%% library-book.tex -- a standalone wrapper around the GENERATED book of theorems
%% and definitions by library (reference/library-book.tex, written at every load by
%% (write-library-book-tex! "reference/library-book.tex"), library-book.scm).
%%
%% One part per library of libraries.sexp, in an order in which every theorem
%% follows the theorems and definitions it depends on.  No proofs: a proposition's
%% proof is (display-proof 'NAME "file.pdf") in the prover.
%%
%% Build:  make library-book     (pdflatex twice, for the table of contents)

\documentclass[11pt]{report}
\usepackage[T1]{fontenc}
\usepackage[utf8]{inputenc}
\usepackage[fleqn]{amsmath}
\setlength{\mathindent}{1em}
\usepackage{amssymb,amsthm}
\usepackage[margin=1in]{geometry}
\usepackage{graphicx}
\usepackage{enumitem}
\usepackage[hidelinks]{hyperref}
\setcounter{tocdepth}{1}
\setlength{\parindent}{0pt}
\setlength{\parskip}{2pt}
\setlength{\emergencystretch}{3em}
\sloppy

\title{VNB: the theorems and definitions, by library}
\author{}
\date{}

\begin{document}
\maketitle

\noindent
Every definition, axiom, assumption and theorem the VNB library installs, as a
statement, without proofs, in seven books that are the libraries of
\texttt{libraries.sexp}.  The books are in import order and the entries within a
book in load order, so a theorem always follows the theorems and definitions it
depends on.  A section is a source file (its path is a link).  An
\textbf{Assumption} is a fact accepted without a machine proof, with the tier of
its warrant; a \textbf{Proposition} is proven, and its proof is
\texttt{(display-proof 'NAME "file.pdf")} in the prover.

\noindent
A \textbf{Definition} of a structure (a field, a metric space, a measure space)
transcribes its \texttt{declare-structure} declaration (manual, section
``Mathematical structures'').  It lists the components, numbered as in the
declaration; each is a projection, so that, for a field, $\mathrm{add}(s)$ is
$\mathrm{nth}(2, s)$ for every $s$; each line names the clause that declared it
(\texttt{carriers}, \texttt{op}, \texttt{constant}, \texttt{derived},
\texttt{substructure}, \texttt{family}, \texttt{family-fun}, \texttt{relation})
and gives its typing.  The laws follow (the \texttt{property} and \texttt{law}
clauses).  The structure's predicate holds of $s$ exactly when $s$ has the stated
length, every component has its typing, and every law holds.

% library-book.tex -- GENERATED by (write-library-book-tex! ...) at load.  Do not hand-edit.
% The theorems and definitions of the tree, no proofs, one part per library.
% \vnbstmt -- a displayed statement that shrinks to the text width only
% when it would overflow it, so short statements are set at normal size and
% a wide one never runs into the margin.  Needs graphicx.
%
% The leading \leavevmode is not decoration.  amsthm defers the theorem head
% (`Theorem 1.1 (Title).') into \everypar, so it is injected at the start of
% the first paragraph of the body; without \leavevmode the opening \par is a
% no-op on an empty paragraph, the head lands in the SAME paragraph as the
% box below, and head-width + \linewidth overfills every statement by 300 pt.
% \leavevmode starts a paragraph for the head to attach to, and the \par then
% closes it before the box.
\providecommand{\vnbstmt}[1]{%
  \leavevmode\par\nobreak\smallskip\noindent
  \resizebox{\ifdim\width>\linewidth\linewidth\else\width\fi}{!}{$\displaystyle #1$}%
  \par\smallskip}
\providecommand{\vnbmissing}[1]{\textit{No theorem \texttt{#1} is installed in the library.}}
% \vnbprf -- one paragraph of the proof block.  Set RAGGED RIGHT: these are
% comma-separated lists of \texttt identifiers, which neither hyphenate nor
% break at their hyphens, so justified setting either overfills the measure or
% (under \sloppy) buys the fit with very loose interword glue.
\providecommand{\vnbprf}[1]{{\raggedright #1\par}}
% \vnbpart -- one PART of a statement that the library gives under several names:
% its own paragraph, the name as a run-in label (the user, 2026-09-21).
\providecommand{\vnbpart}[1]{\par\medskip{\raggedright\noindent\textup{\textbf{Part} #1.}\par}\nopagebreak\smallskip}
\providecommand{\vnbproofpart}[1]{\par\medskip\noindent\textbf{Proof of} #1\textbf{.}\par\nopagebreak}
\providecommand{\vnbstructdef}[2]{\par\medskip\noindent\textbf{Structure definition} (\texttt{#1}).\par\nopagebreak#2\par}
\providecommand{\vnbdef}[2]{\par\medskip\noindent\textbf{Definition} (\texttt{#1}).\par\nopagebreak#2\par}
\providecommand{\vnbaxiom}[2]{\par\medskip\noindent\textbf{Axiom} (\texttt{#1}).\par\nopagebreak#2\par}
\providecommand{\vnbassume}[3]{\par\medskip\noindent\textbf{Assumption} (\texttt{#1}; warrant: #2).\par\nopagebreak#3\par}
\providecommand{\vnbprop}[2]{\par\medskip\noindent\textbf{Proposition} (\texttt{#1}).\par\nopagebreak#2\par}
\providecommand{\vnbbookfile}[2]{\section*{\href{file:../#2.scm}{\texttt{#1}}}\addcontentsline{toc}{section}{\texttt{#1}}}
\tableofcontents

\part{Foundations}
\noindent Logic and sets; the number systems NN, ZZ, QQ, RR and CC with their arithmetic and order; tuples and lists; functions, bijections and intervals; finite sums; cardinality and ordinals; choice (Zermelo, Zorn); the combinatorial lemmas. Every other book stands on this one.
\par\noindent{\small 141 files.}

\vnbbookfile{structure-library/operation-properties}{structure-library/operation-properties}
\noindent{\small operation-properties.scm -- named properties of operations.}\par
\par\medskip\noindent\textbf{Definition} (\texttt{is-associative}).\;For all $\mathit{op}$ and $\mathit{crr}$, 
\begin{equation*}
\begin{array}{@{}l@{}} \operatorname{is\text{-}associative}(\mathit{op}, \mathit{crr}) \\ \Leftrightarrow\; \forall u, v, w \in \mathit{crr}.\; (\operatorname{op}(\operatorname{op}(u, v), w) = \operatorname{op}(u, \operatorname{op}(v, w))) \end{array}
\end{equation*}

\par\medskip\noindent\textbf{Definition} (\texttt{is-commutative}).\;For all $\mathit{op}$ and $\mathit{crr}$, 
\begin{equation*}
\operatorname{is\text{-}commutative}(\mathit{op}, \mathit{crr}) \Leftrightarrow \forall u, v \in \mathit{crr}.\; (\operatorname{op}(u, v) = \operatorname{op}(v, u))
\end{equation*}

\par\medskip\noindent\textbf{Definition} (\texttt{is-identity}).\;For all $\mathit{op}$, $\mathit{unit}$ and $\mathit{crr}$, 
\begin{equation*}
\operatorname{is\text{-}identity}(\mathit{op}, \mathit{unit}, \mathit{crr}) \Leftrightarrow \forall u \in \mathit{crr}.\; ((\operatorname{op}(\mathit{unit}, u) = u) \wedge (\operatorname{op}(u, \mathit{unit}) = u))
\end{equation*}

\par\medskip\noindent\textbf{Definition} (\texttt{has-inverses}).\;For all $\mathit{op}$, $\mathit{unit}$, $\mathit{invop}$ and $\mathit{crr}$, 
\begin{equation*}
\begin{array}{@{}l@{}} \operatorname{has\text{-}inverses}(\mathit{op}, \mathit{unit}, \mathit{invop}, \mathit{crr}) \\ \Leftrightarrow\; \forall u \in \mathit{crr}.\; ((\operatorname{op}(\operatorname{invop}(u), u) = \mathit{unit}) \wedge (\operatorname{op}(u, \operatorname{invop}(u)) = \mathit{unit})) \end{array}
\end{equation*}

\par\medskip\noindent\textbf{Definition} (\texttt{is-distributive}).\;For all $\mathit{addop}$, $\mathit{mulop}$ and $\mathit{crr}$, 
\begin{equation*}
\begin{array}{@{}l@{}} \operatorname{is\text{-}distributive}(\mathit{addop}, \mathit{mulop}, \mathit{crr}) \\ \Leftrightarrow\; \forall u, v, w \in \mathit{crr}. \\ \quad \operatorname{mulop}(u, \operatorname{addop}(v, w)) = \operatorname{addop}(\operatorname{mulop}(u, v), \operatorname{mulop}(u, w)) \wedge \\ \quad \operatorname{mulop}(\operatorname{addop}(u, v), w) = \operatorname{addop}(\operatorname{mulop}(u, w), \operatorname{mulop}(v, w)) \end{array}
\end{equation*}

\par\medskip\noindent\textbf{Definition} (\texttt{is-norm}).\;For all $\mathit{nm}$, $\mathit{addop}$, $\mathit{mulop}$, $\mathit{zr}$ and $\mathit{crr}$, 
\begin{equation*}
\begin{array}{@{}l@{}} \operatorname{is\text{-}norm}(\mathit{nm}, \mathit{addop}, \mathit{mulop}, \mathit{zr}, \mathit{crr}) \\ \Leftrightarrow\; \mathit{nm} \in (\mathit{crr} \to \mathbb{R}) \wedge \\ \forall a \in \mathit{crr}. \\ \quad 0 \leq \operatorname{nm}(a) \wedge \\ \quad (\operatorname{nm}(a) = 0) \Leftrightarrow (a = \mathit{zr}) \wedge \\ \quad \forall b \in \mathit{crr}. \\ \quad \quad (\operatorname{nm}(\operatorname{mulop}(a, b)) = \cdot(\operatorname{nm}(a), \operatorname{nm}(b))) \wedge (\operatorname{nm}(\operatorname{addop}(a, b)) \leq +(\operatorname{nm}(a), \operatorname{nm}(b))) \end{array}
\end{equation*}

\par\medskip\noindent\textbf{Definition} (\texttt{is-group-norm}).\;For all $\mathit{nm}$, $\mathit{op}$, $\mathit{invop}$, $\mathit{unit}$ and $\mathit{crr}$, 
\begin{equation*}
\begin{array}{@{}l@{}} \operatorname{is\text{-}group\text{-}norm}(\mathit{nm}, \mathit{op}, \mathit{invop}, \mathit{unit}, \mathit{crr}) \\ \Leftrightarrow\; \mathit{nm} \in (\mathit{crr} \to \mathbb{R}) \wedge \\ \forall u \in \mathit{crr}. \\ \quad 0 \leq \operatorname{nm}(u) \wedge \\ \quad (\operatorname{nm}(u) = 0) \Leftrightarrow (u = \mathit{unit}) \wedge \\ \quad \operatorname{nm}(\operatorname{invop}(u)) = \operatorname{nm}(u) \wedge \\ \quad \forall v \in \mathit{crr}.\; (\operatorname{nm}(\operatorname{op}(u, v)) \leq +(\operatorname{nm}(u), \operatorname{nm}(v))) \end{array}
\end{equation*}

\par\medskip\noindent\textbf{Definition} (\texttt{is-metric}).\;For all $\mathit{dst}$ and $\mathit{crr}$, 
\begin{equation*}
\begin{array}{@{}l@{}} \operatorname{is\text{-}metric}(\mathit{dst}, \mathit{crr}) \\ \Leftrightarrow\; \forall u \in \mathit{crr}. \\ \quad \operatorname{dst}(u, u) = 0 \wedge \\ \quad \forall v \in \mathit{crr}. \\ \quad \quad 0 \leq \operatorname{dst}(u, v) \wedge \\ \quad \quad (\operatorname{dst}(u, v) = 0) \Rightarrow (u = v) \wedge \\ \quad \quad \operatorname{dst}(u, v) = \operatorname{dst}(v, u) \wedge \\ \quad \quad \forall w \in \mathit{crr}.\; (\operatorname{dst}(u, w) \leq +(\operatorname{dst}(u, v), \operatorname{dst}(v, w))) \end{array}
\end{equation*}

\par\medskip\noindent\textbf{Definition} (\texttt{is-pseudometric}).\;For all $\mathit{dst}$ and $\mathit{crr}$, 
\begin{equation*}
\begin{array}{@{}l@{}} \operatorname{is\text{-}pseudometric}(\mathit{dst}, \mathit{crr}) \\ \Leftrightarrow\; \forall u \in \mathit{crr}. \\ \quad \operatorname{dst}(u, u) = 0 \wedge \\ \quad \forall v \in \mathit{crr}. \\ \quad \quad ((0 \leq \operatorname{dst}(u, v)) \wedge ((\operatorname{dst}(u, v) = \operatorname{dst}(v, u)) \wedge \forall w \in \mathit{crr}.\; (\operatorname{dst}(u, w) \leq +(\operatorname{dst}(u, v), \operatorname{dst}(v, w))))) \end{array}
\end{equation*}


\vnbbookfile{theorem-library/axioms}{theorem-library/axioms}
\noindent{\small axioms.scm -- axioms not yet derivable from primitive inferences}\par
\vnbaxiom{equality-symmetry}{\vnbstmt{\forall a, b.\; ((a = b) \Rightarrow (b = a))}}
\vnbaxiom{quasi-eq-subst-membership}{\vnbstmt{\forall a, b, s.\; (((a \simeq b) \wedge (b \in s)) \Rightarrow (a \in s))}}
\vnbaxiom{quasi-eq-def}{\vnbstmt{\forall a, b.\; ((a = a) \Rightarrow ((a \simeq b) \Leftrightarrow (a = b)))}}
\vnbaxiom{equality-transitivity}{\vnbstmt{\forall a, b, c.\; (((a = b) \wedge (b = c)) \Rightarrow (a = c))}}
\vnbaxiom{apply-tupling-2}{\vnbstmt{\forall f, a, b.\; (f(a, b) \simeq f(\langle a, b \rangle))}}
\vnbaxiom{quasi-eq-reflexivity}{\vnbstmt{\forall a.\; (a \simeq a)}}
\vnbaxiom{eq-subst-membership}{\vnbstmt{\forall a, b, s.\; (((a = b) \wedge (b \in s)) \Rightarrow (a \in s))}}
\vnbaxiom{quasi-eq-symmetry}{\vnbstmt{\forall a, b.\; ((a \simeq b) \Rightarrow (b \simeq a))}}
\vnbaxiom{apply-tupling-3}{\vnbstmt{\forall f, a, b, c.\; (f(a, b, c) \simeq f(\langle a, b, c \rangle))}}
\vnbaxiom{list-sethood}{\vnbstmt{\forall a, l.\; (((a \in \mathbf{Set}) \wedge (l \in \operatorname{tuples}(a))) \Rightarrow (l \in \mathbf{Set}))}}
\vnbaxiom{quasi-eq-transitivity}{\vnbstmt{\forall a, b, c.\; (((a \simeq b) \wedge (b \simeq c)) \Rightarrow (a \simeq c))}}
\vnbaxiom{fun-apply-type}{\vnbstmt{\forall f, a, b, x.\; (((f \in (a \to b)) \wedge (x \in a)) \Rightarrow (f(x) \in b))}}

\vnbbookfile{structure-library/set-vocabulary}{structure-library/set-vocabulary}
\noindent{\small set-vocabulary.scm -- two set constructors that the rest of the library is STATED with, and that nothing but the kernel's own constructors is needed to define.}\par
\par\medskip\noindent\textbf{Definition} (\texttt{difference}).\;$\operatorname{difference}(\operatorname{u\_}, \operatorname{b\_}) :=$ $\operatorname{complement\text{-}in}(\operatorname{u\_}, \operatorname{b\_})$
\par\medskip\noindent\textbf{Definition} (\texttt{singleton}).\;$\operatorname{singleton}(\operatorname{y\_}) :=$ $\operatorname{make\text{-}set}(\langle \operatorname{y\_} \rangle)$

\vnbbookfile{structure-library/intersection-of}{structure-library/intersection-of}
\noindent{\small intersection-of.scm -- INTERSECTION-OF(c), the intersection of the members of a family c, as a DEFINED constructor.}\par
\par\medskip\noindent\textbf{Definition} (\texttt{intersection-of}).\;$\operatorname{intersection\text{-}of}(c) :=$
\begin{equation*}
\{\operatorname{iov\_} \in \operatorname{big\text{-}union}(\operatorname{iow\_}, c, \operatorname{iow\_}) : \forall \operatorname{iou\_} \in c.\; (\operatorname{iov\_} \in \operatorname{iou\_})\}
\end{equation*}


\vnbbookfile{structure-library/compose}{structure-library/compose}
\noindent{\small compose.scm -- COMPOSE: composition of functions, f o g.}\par
\par\medskip\noindent\textbf{Definition} (\texttt{compose}).\;$\operatorname{compose}(f, g) :=$ $(\lambda \operatorname{z\_} \in \operatorname{dom}(g).\; f(g(\operatorname{z\_})))$

\vnbbookfile{structure-library/id-fun}{structure-library/id-fun}
\noindent{\small id-fun.scm -- ID-FUN: the identity map of a set.}\par
\par\medskip\noindent\textbf{Definition} (\texttt{id-fun}).\;$\operatorname{id\text{-}fun}(\operatorname{a\_}) :=$ $(\lambda \operatorname{idfx\_} \in \operatorname{a\_}.\; \operatorname{idfx\_})$

\vnbbookfile{number-systems}{number-systems}
\noindent{\small number-systems.scm -- number system constants and axioms}\par
\par\medskip\noindent\textbf{Definition} (\texttt{rr-upper-bound}).\;For all $s$ and $b$, 
\begin{equation*}
\operatorname{rr\text{-}upper\text{-}bound}(s, b) \Leftrightarrow ((b \in \mathbb{R}) \wedge \forall x \in s.\; (x \leq b))
\end{equation*}

\par\medskip\noindent\textbf{Definition} (\texttt{rr-bounded-above}).\;For every $s$,  $\operatorname{rr\text{-}bounded\text{-}above}(s) \Leftrightarrow \exists b.\; \operatorname{rr\text{-}upper\text{-}bound}(s, b)$
\vnbdef{binary-minus-def}{\vnbstmt{\forall a, b.\; ((a - b) \simeq (a + -b))}}
\vnbdef{binary-divide-def}{\vnbstmt{\forall a, b.\; (/(a, b) \simeq (a \cdot \operatorname{recip}(b)))}}
\vnbdef{rr-abs-def}{\vnbstmt{\begin{array}{@{}l@{}} \forall x \in \mathbb{R}. \\ \quad ((0 \leq x) \Rightarrow (\lvert x \rvert = x)) \wedge (\neg (0 \leq x) \Rightarrow (\lvert x \rvert = -x)) \end{array}}}
\vnbdef{rr-max-def}{\vnbstmt{\begin{array}{@{}l@{}} \forall x, y \in \mathbb{R}. \\ \quad ((y \leq x) \Rightarrow (\operatorname{max}(x, y) = x)) \wedge (\neg (y \leq x) \Rightarrow (\operatorname{max}(x, y) = y)) \end{array}}}
\vnbdef{rr-min-def}{\vnbstmt{\begin{array}{@{}l@{}} \forall x, y \in \mathbb{R}. \\ \quad ((x \leq y) \Rightarrow (\operatorname{min}(x, y) = x)) \wedge (\neg (x \leq y) \Rightarrow (\operatorname{min}(x, y) = y)) \end{array}}}
\vnbdef{real-part-def}{\vnbstmt{\forall z.\; (\operatorname{real\text{-}part}(z) \simeq ((z + \operatorname{conjugate}(z)) \cdot \operatorname{recip}(2)))}}
\vnbdef{imag-part-def}{\vnbstmt{\forall z.\; (\operatorname{imag\text{-}part}(z) \simeq ((z - \operatorname{conjugate}(z)) \cdot \operatorname{recip}(+2i)))}}
\vnbaxiom{rr-add-assoc}{\vnbstmt{\begin{array}{@{}l@{}} \forall a, b, c. \\ \quad ((a \in \mathbb{R}) \wedge ((b \in \mathbb{R}) \wedge (c \in \mathbb{R}))) \Rightarrow (((a + b) + c) = (a + (b + c))) \end{array}}}
\vnbaxiom{cc-conjugate-fixes-rr}{\vnbstmt{\forall a \in \mathbb{R}.\; (\operatorname{conjugate}(a) = a)}}
\vnbaxiom{cc-neg-closed}{\vnbstmt{\forall a \in \mathbb{C}.\; (-a \in \mathbb{C})}}
\vnbaxiom{cc-mul-closed}{\vnbstmt{\forall a, b.\; (((a \in \mathbb{C}) \wedge (b \in \mathbb{C})) \Rightarrow ((a \cdot b) \in \mathbb{C}))}}
\vnbaxiom{qq-subset-rr}{\vnbstmt{\forall n \in \mathbb{Q}.\; (n \in \mathbb{R})}}
\vnbaxiom{qq-recip-closed}{\vnbstmt{\forall a \in \mathbb{Q}.\; (\neg (a = 0) \Rightarrow (\operatorname{recip}(a) \in \mathbb{Q}))}}
\vnbaxiom{rr-add-zero}{\vnbstmt{\forall a \in \mathbb{R}.\; (a = a)}}
\vnbaxiom{cc-add-closed}{\vnbstmt{\forall a, b.\; (((a \in \mathbb{C}) \wedge (b \in \mathbb{C})) \Rightarrow ((a + b) \in \mathbb{C}))}}
\vnbaxiom{nn-succ-closed}{\vnbstmt{\forall n \in \mathbb{N}.\; ((n + 1) \in \mathbb{N})}}
\vnbaxiom{cc-conjugate-closed}{\vnbstmt{\forall a \in \mathbb{C}.\; (\operatorname{conjugate}(a) \in \mathbb{C})}}
\vnbaxiom{zz-add-comm}{\vnbstmt{\forall a, b.\; (((a \in \mathbb{Z}) \wedge (b \in \mathbb{Z})) \Rightarrow ((a + b) = (b + a)))}}
\vnbaxiom{rr-leq-reflexive}{\vnbstmt{\forall a \in \mathbb{R}.\; (a \leq a)}}
\vnbaxiom{qq-recip-inverse}{\vnbstmt{\forall a \in \mathbb{Q}.\; (\neg (a = 0) \Rightarrow ((a \cdot \operatorname{recip}(a)) = 1))}}
\vnbaxiom{qq-mul-comm}{\vnbstmt{\forall a, b.\; (((a \in \mathbb{Q}) \wedge (b \in \mathbb{Q})) \Rightarrow ((a \cdot b) = (b \cdot a)))}}
\vnbaxiom{zz-add-closed}{\vnbstmt{\forall a, b.\; (((a \in \mathbb{Z}) \wedge (b \in \mathbb{Z})) \Rightarrow ((a + b) \in \mathbb{Z}))}}
\vnbaxiom{cc-mul-assoc}{\vnbstmt{\begin{array}{@{}l@{}} \forall a, b, c. \\ \quad ((a \in \mathbb{C}) \wedge ((b \in \mathbb{C}) \wedge (c \in \mathbb{C}))) \Rightarrow (((a \cdot b) \cdot c) = (a \cdot (b \cdot c))) \end{array}}}
\vnbaxiom{zz-mul-comm}{\vnbstmt{\forall a, b.\; (((a \in \mathbb{Z}) \wedge (b \in \mathbb{Z})) \Rightarrow ((a \cdot b) = (b \cdot a)))}}
\vnbaxiom{zz-subset-qq}{\vnbstmt{\forall n \in \mathbb{Z}.\; (n \in \mathbb{Q})}}
\vnbaxiom{rr-zero-in}{\vnbstmt{0 \in \mathbb{R}}}
\vnbaxiom{rr-leq-add-compat}{\vnbstmt{\begin{array}{@{}l@{}} \forall a, b, c. \\ \quad (((a \in \mathbb{R}) \wedge ((b \in \mathbb{R}) \wedge (c \in \mathbb{R}))) \wedge (a \leq b)) \Rightarrow ((a + c) \leq (b + c)) \end{array}}}
\vnbaxiom{power-neg}{\vnbstmt{\begin{array}{@{}l@{}} \forall x, n. \\ \quad ((x \in \mathbb{C}) \wedge ((n \in \mathbb{N}) \wedge \neg (x = 0))) \Rightarrow (\operatorname{power}(x, -n) = \operatorname{recip}(\operatorname{power}(x, n))) \end{array}}}
\vnbaxiom{cc-is-set}{\vnbstmt{\mathbb{C} \in \mathbf{Set}}}
\vnbaxiom{power-succ}{\vnbstmt{\forall x, n.\; (((x \in \mathbb{C}) \wedge (n \in \mathbb{N})) \Rightarrow (\operatorname{power}(x, (n + 1)) = (x \cdot \operatorname{power}(x, n))))}}
\vnbaxiom{qq-mul-closed}{\vnbstmt{\forall a, b.\; (((a \in \mathbb{Q}) \wedge (b \in \mathbb{Q})) \Rightarrow ((a \cdot b) \in \mathbb{Q}))}}
\vnbaxiom{rr-add-comm}{\vnbstmt{\forall a, b.\; (((a \in \mathbb{R}) \wedge (b \in \mathbb{R})) \Rightarrow ((a + b) = (b + a)))}}
\vnbaxiom{zz-mul-closed}{\vnbstmt{\forall a, b.\; (((a \in \mathbb{Z}) \wedge (b \in \mathbb{Z})) \Rightarrow ((a \cdot b) \in \mathbb{Z}))}}
\vnbaxiom{zz-distributive}{\vnbstmt{\begin{array}{@{}l@{}} \forall a, b, c. \\ \quad ((a \in \mathbb{Z}) \wedge ((b \in \mathbb{Z}) \wedge (c \in \mathbb{Z}))) \Rightarrow ((a \cdot (b + c)) = ((a \cdot b) + (a \cdot c))) \end{array}}}
\vnbaxiom{cc-generated-by-rr}{\vnbstmt{\forall z \in \mathbb{C}.\; \exists x, y \in \mathbb{R}.\; (z = (x + (y \cdot +i)))}}
\vnbaxiom{qq-one-mul}{\vnbstmt{\forall a \in \mathbb{Q}.\; (a = a)}}
\vnbaxiom{rr-subset-cc}{\vnbstmt{\forall n \in \mathbb{R}.\; (n \in \mathbb{C})}}
\vnbaxiom{cc-conjugate-involution}{\vnbstmt{\forall a \in \mathbb{C}.\; (\operatorname{conjugate}(\operatorname{conjugate}(a)) = a)}}
\vnbaxiom{qq-is-fraction}{\vnbstmt{\forall q \in \mathbb{Q}.\; \exists a, b \in \mathbb{Z}.\; (\neg (b = 0) \wedge (q = (a \cdot \operatorname{recip}(b))))}}
\vnbaxiom{qq-is-set}{\vnbstmt{\mathbb{Q} \in \mathbf{Set}}}
\vnbaxiom{rr-leq-total}{\vnbstmt{\forall a, b.\; (((a \in \mathbb{R}) \wedge (b \in \mathbb{R})) \Rightarrow ((a \leq b) \vee (b \leq a)))}}
\vnbaxiom{cc-distributive}{\vnbstmt{\begin{array}{@{}l@{}} \forall a, b, c. \\ \quad ((a \in \mathbb{C}) \wedge ((b \in \mathbb{C}) \wedge (c \in \mathbb{C}))) \Rightarrow ((a \cdot (b + c)) = ((a \cdot b) + (a \cdot c))) \end{array}}}
\vnbaxiom{qq-add-assoc}{\vnbstmt{\begin{array}{@{}l@{}} \forall a, b, c. \\ \quad ((a \in \mathbb{Q}) \wedge ((b \in \mathbb{Q}) \wedge (c \in \mathbb{Q}))) \Rightarrow (((a + b) + c) = (a + (b + c))) \end{array}}}
\vnbaxiom{rr-one-in}{\vnbstmt{1 \in \mathbb{R}}}
\vnbaxiom{rr-recip-closed}{\vnbstmt{\forall a \in \mathbb{R}.\; (\neg (a = 0) \Rightarrow (\operatorname{recip}(a) \in \mathbb{R}))}}
\vnbaxiom{qq-mul-assoc}{\vnbstmt{\begin{array}{@{}l@{}} \forall a, b, c. \\ \quad ((a \in \mathbb{Q}) \wedge ((b \in \mathbb{Q}) \wedge (c \in \mathbb{Q}))) \Rightarrow (((a \cdot b) \cdot c) = (a \cdot (b \cdot c))) \end{array}}}
\vnbaxiom{cc-self-conj-real}{\vnbstmt{\forall a \in \mathbb{C}.\; ((a \cdot \operatorname{conjugate}(a)) \in \mathbb{R})}}
\vnbaxiom{qq-one-in}{\vnbstmt{1 \in \mathbb{Q}}}
\vnbaxiom{zz-zero-in}{\vnbstmt{0 \in \mathbb{Z}}}
\vnbaxiom{cc-add-zero}{\vnbstmt{\forall a \in \mathbb{C}.\; (a = a)}}
\vnbaxiom{cc-zero-in}{\vnbstmt{0 \in \mathbb{C}}}
\vnbaxiom{zz-mul-assoc}{\vnbstmt{\begin{array}{@{}l@{}} \forall a, b, c. \\ \quad ((a \in \mathbb{Z}) \wedge ((b \in \mathbb{Z}) \wedge (c \in \mathbb{Z}))) \Rightarrow (((a \cdot b) \cdot c) = (a \cdot (b \cdot c))) \end{array}}}
\vnbaxiom{nn-add-comm}{\vnbstmt{\forall a, b.\; (((a \in \mathbb{N}) \wedge (b \in \mathbb{N})) \Rightarrow ((a + b) = (b + a)))}}
\vnbaxiom{rr-neg-inverse}{\vnbstmt{\forall a \in \mathbb{R}.\; ((a + -a) = 0)}}
\vnbaxiom{cc-add-comm}{\vnbstmt{\forall a, b.\; (((a \in \mathbb{C}) \wedge (b \in \mathbb{C})) \Rightarrow ((a + b) = (b + a)))}}
\vnbaxiom{rr-add-closed}{\vnbstmt{\forall a, b.\; (((a \in \mathbb{R}) \wedge (b \in \mathbb{R})) \Rightarrow ((a + b) \in \mathbb{R}))}}
\vnbaxiom{cc-conjugate-mul}{\vnbstmt{\begin{array}{@{}l@{}} \forall a, b. \\ \quad ((a \in \mathbb{C}) \wedge (b \in \mathbb{C})) \Rightarrow (\operatorname{conjugate}((a \cdot b)) = (\operatorname{conjugate}(a) \cdot \operatorname{conjugate}(b))) \end{array}}}
\vnbaxiom{cc-conjugate-i}{\vnbstmt{\operatorname{conjugate}(+i) = -+i}}
\vnbaxiom{rr-leq-mul-nonneg}{\vnbstmt{\forall a, b.\; ((((a \in \mathbb{R}) \wedge (b \in \mathbb{R})) \wedge ((0 \leq a) \wedge (0 \leq b))) \Rightarrow (0 \leq (a \cdot b)))}}
\vnbaxiom{zz-one-mul}{\vnbstmt{\forall a \in \mathbb{Z}.\; (a = a)}}
\vnbaxiom{cc-neg-inverse}{\vnbstmt{\forall a \in \mathbb{C}.\; ((a + -a) = 0)}}
\vnbaxiom{rr-mul-assoc}{\vnbstmt{\begin{array}{@{}l@{}} \forall a, b, c. \\ \quad ((a \in \mathbb{R}) \wedge ((b \in \mathbb{R}) \wedge (c \in \mathbb{R}))) \Rightarrow (((a \cdot b) \cdot c) = (a \cdot (b \cdot c))) \end{array}}}
\vnbaxiom{qq-add-closed}{\vnbstmt{\forall a, b.\; (((a \in \mathbb{Q}) \wedge (b \in \mathbb{Q})) \Rightarrow ((a + b) \in \mathbb{Q}))}}
\vnbaxiom{cc-conjugate-add}{\vnbstmt{\begin{array}{@{}l@{}} \forall a, b. \\ \quad ((a \in \mathbb{C}) \wedge (b \in \mathbb{C})) \Rightarrow (\operatorname{conjugate}((a + b)) = (\operatorname{conjugate}(a) + \operatorname{conjugate}(b))) \end{array}}}
\vnbaxiom{rr-one-mul}{\vnbstmt{\forall a \in \mathbb{R}.\; (a = a)}}
\vnbaxiom{rr-mul-closed}{\vnbstmt{\forall a, b.\; (((a \in \mathbb{R}) \wedge (b \in \mathbb{R})) \Rightarrow ((a \cdot b) \in \mathbb{R}))}}
\vnbaxiom{rr-leq-antisymmetric}{\vnbstmt{\forall a, b.\; ((((a \in \mathbb{R}) \wedge (b \in \mathbb{R})) \wedge ((a \leq b) \wedge (b \leq a))) \Rightarrow (a = b))}}
\vnbaxiom{power-zero}{\vnbstmt{\forall x \in \mathbb{C}.\; (\operatorname{power}(x, 0) = 1)}}
\vnbaxiom{nn-add-assoc}{\vnbstmt{\begin{array}{@{}l@{}} \forall a, b, c. \\ \quad ((a \in \mathbb{N}) \wedge ((b \in \mathbb{N}) \wedge (c \in \mathbb{N}))) \Rightarrow (((a + b) + c) = (a + (b + c))) \end{array}}}
\vnbaxiom{cc-one-mul}{\vnbstmt{\forall a \in \mathbb{C}.\; (a = a)}}
\vnbaxiom{cc-add-assoc}{\vnbstmt{\begin{array}{@{}l@{}} \forall a, b, c. \\ \quad ((a \in \mathbb{C}) \wedge ((b \in \mathbb{C}) \wedge (c \in \mathbb{C}))) \Rightarrow (((a + b) + c) = (a + (b + c))) \end{array}}}
\vnbaxiom{qq-zero-in}{\vnbstmt{0 \in \mathbb{Q}}}
\vnbaxiom{nn-distributive}{\vnbstmt{\begin{array}{@{}l@{}} \forall a, b, c. \\ \quad ((a \in \mathbb{N}) \wedge ((b \in \mathbb{N}) \wedge (c \in \mathbb{N}))) \Rightarrow ((a \cdot (b + c)) = ((a \cdot b) + (a \cdot c))) \end{array}}}
\vnbaxiom{nn-one-mul}{\vnbstmt{\forall a \in \mathbb{N}.\; (a = a)}}
\vnbaxiom{nn-subset-zz}{\vnbstmt{\forall n \in \mathbb{N}.\; (n \in \mathbb{Z})}}
\vnbaxiom{nn-mul-assoc}{\vnbstmt{\begin{array}{@{}l@{}} \forall a, b, c. \\ \quad ((a \in \mathbb{N}) \wedge ((b \in \mathbb{N}) \wedge (c \in \mathbb{N}))) \Rightarrow (((a \cdot b) \cdot c) = (a \cdot (b \cdot c))) \end{array}}}
\vnbaxiom{cc-recip-inverse}{\vnbstmt{\forall a \in \mathbb{C}.\; (\neg (a = 0) \Rightarrow ((a \cdot \operatorname{recip}(a)) = 1))}}
\vnbaxiom{cc-self-conj-nonneg}{\vnbstmt{\forall a \in \mathbb{C}.\; (((a \cdot \operatorname{conjugate}(a)) \in \mathbb{R}) \wedge (0 \leq (a \cdot \operatorname{conjugate}(a))))}}
\vnbaxiom{nn-mul-comm}{\vnbstmt{\forall a, b.\; (((a \in \mathbb{N}) \wedge (b \in \mathbb{N})) \Rightarrow ((a \cdot b) = (b \cdot a)))}}
\vnbaxiom{rr-distributive}{\vnbstmt{\begin{array}{@{}l@{}} \forall a, b, c. \\ \quad ((a \in \mathbb{R}) \wedge ((b \in \mathbb{R}) \wedge (c \in \mathbb{R}))) \Rightarrow ((a \cdot (b + c)) = ((a \cdot b) + (a \cdot c))) \end{array}}}
\vnbaxiom{zz-neg-inverse}{\vnbstmt{\forall a \in \mathbb{Z}.\; ((a + -a) = 0)}}
\vnbaxiom{zz-add-assoc}{\vnbstmt{\begin{array}{@{}l@{}} \forall a, b, c. \\ \quad ((a \in \mathbb{Z}) \wedge ((b \in \mathbb{Z}) \wedge (c \in \mathbb{Z}))) \Rightarrow (((a + b) + c) = (a + (b + c))) \end{array}}}
\vnbaxiom{qq-neg-closed}{\vnbstmt{\forall a \in \mathbb{Q}.\; (-a \in \mathbb{Q})}}
\vnbaxiom{cc-mul-comm}{\vnbstmt{\forall a, b.\; (((a \in \mathbb{C}) \wedge (b \in \mathbb{C})) \Rightarrow ((a \cdot b) = (b \cdot a)))}}
\vnbaxiom{nn-induction}{\vnbstmt{\begin{array}{@{}l@{}} \forall c.; \forall n \in \mathbb{N}. \\ \quad ((0 \in c) \wedge \forall n \in \mathbb{N}.\; ((n \in c) \Rightarrow ((n + 1) \in c))) \Rightarrow (n \in c) \end{array}}}
\vnbaxiom{zz-is-set}{\vnbstmt{\mathbb{Z} \in \mathbf{Set}}}
\vnbaxiom{nn-add-closed}{\vnbstmt{\forall a, b.\; (((a \in \mathbb{N}) \wedge (b \in \mathbb{N})) \Rightarrow ((a + b) \in \mathbb{N}))}}
\vnbaxiom{cc-one-in}{\vnbstmt{1 \in \mathbb{C}}}
\vnbaxiom{qq-add-comm}{\vnbstmt{\forall a, b.\; (((a \in \mathbb{Q}) \wedge (b \in \mathbb{Q})) \Rightarrow ((a + b) = (b + a)))}}
\vnbaxiom{zz-neg-closed}{\vnbstmt{\forall a \in \mathbb{Z}.\; (-a \in \mathbb{Z})}}
\vnbaxiom{rr-is-set}{\vnbstmt{\mathbb{R} \in \mathbf{Set}}}
\vnbaxiom{nn-is-set}{\vnbstmt{\mathbb{N} \in \mathbf{Set}}}
\vnbaxiom{rr-mul-comm}{\vnbstmt{\forall a, b.\; (((a \in \mathbb{R}) \wedge (b \in \mathbb{R})) \Rightarrow ((a \cdot b) = (b \cdot a)))}}
\vnbaxiom{cc-i-squared}{\vnbstmt{(+i \cdot +i) = -1}}
\vnbaxiom{cc-recip-closed}{\vnbstmt{\forall a \in \mathbb{C}.\; (\neg (a = 0) \Rightarrow (\operatorname{recip}(a) \in \mathbb{C}))}}
\vnbaxiom{rr-leq-transitive}{\vnbstmt{\begin{array}{@{}l@{}} \forall a, b, c. \\ \quad (((a \in \mathbb{R}) \wedge ((b \in \mathbb{R}) \wedge (c \in \mathbb{R}))) \wedge ((a \leq b) \wedge (b \leq c))) \Rightarrow (a \leq c) \end{array}}}
\vnbaxiom{rr-recip-inverse}{\vnbstmt{\forall a \in \mathbb{R}.\; (\neg (a = 0) \Rightarrow ((a \cdot \operatorname{recip}(a)) = 1))}}
\vnbaxiom{qq-neg-inverse}{\vnbstmt{\forall a \in \mathbb{Q}.\; ((a + -a) = 0)}}
\vnbaxiom{nn-mul-closed}{\vnbstmt{\forall a, b.\; (((a \in \mathbb{N}) \wedge (b \in \mathbb{N})) \Rightarrow ((a \cdot b) \in \mathbb{N}))}}
\vnbaxiom{rr-neg-closed}{\vnbstmt{\forall a \in \mathbb{R}.\; (-a \in \mathbb{R})}}
\vnbaxiom{nn-zero-in}{\vnbstmt{0 \in \mathbb{N}}}
\vnbaxiom{rr-sup-in}{\vnbstmt{\begin{array}{@{}l@{}} \forall s. \\ \quad ((s \subseteq \mathbb{R}) \wedge (\exists x.\; (x \in s) \wedge \operatorname{rr\text{-}bounded\text{-}above}(s))) \Rightarrow  \\ \quad \quad \operatorname{sup}(s) \in \mathbb{R} \end{array}}}
\vnbaxiom{zz-add-zero}{\vnbstmt{\forall a \in \mathbb{Z}.\; (a = a)}}
\vnbaxiom{zz-one-in}{\vnbstmt{1 \in \mathbb{Z}}}
\vnbaxiom{qq-add-zero}{\vnbstmt{\forall a \in \mathbb{Q}.\; (a = a)}}
\vnbaxiom{zz-generated-by-nn}{\vnbstmt{\forall a \in \mathbb{Z}.\; \exists n \in \mathbb{N}.\; ((a = n) \vee (a = -n))}}
\vnbaxiom{power-typing-nonneg}{\vnbstmt{\forall x, n.\; (((x \in \mathbb{C}) \wedge (n \in \mathbb{N})) \Rightarrow (\operatorname{power}(x, n) \in \mathbb{C}))}}
\vnbaxiom{nn-add-zero}{\vnbstmt{\forall a \in \mathbb{N}.\; (a = a)}}
\vnbaxiom{cc-i-in}{\vnbstmt{+i \in \mathbb{C}}}
\vnbaxiom{qq-distributive}{\vnbstmt{\begin{array}{@{}l@{}} \forall a, b, c. \\ \quad ((a \in \mathbb{Q}) \wedge ((b \in \mathbb{Q}) \wedge (c \in \mathbb{Q}))) \Rightarrow ((a \cdot (b + c)) = ((a \cdot b) + (a \cdot c))) \end{array}}}
\vnbaxiom{rr-sup-least}{\vnbstmt{\begin{array}{@{}l@{}} \forall s, b. \\ \quad ((s \subseteq \mathbb{R}) \wedge (\exists x.\; (x \in s) \wedge (\operatorname{rr\text{-}bounded\text{-}above}(s) \wedge \operatorname{rr\text{-}upper\text{-}bound}(s, b)))) \Rightarrow  \\ \quad \quad \operatorname{sup}(s) \leq b \end{array}}}
\vnbaxiom{rr-sup-upper}{\vnbstmt{\begin{array}{@{}l@{}} \forall s. \\ \quad ((s \subseteq \mathbb{R}) \wedge (\exists x.\; (x \in s) \wedge \operatorname{rr\text{-}bounded\text{-}above}(s))) \Rightarrow  \\ \quad \quad \operatorname{rr\text{-}upper\text{-}bound}(s, \operatorname{sup}(s)) \end{array}}}

\vnbbookfile{structure-library/nn-arith}{structure-library/nn-arith}
\noindent{\small nn-arith.scm -- the recursion equations for + and * on NN.}\par
\vnbdef{nn-add-succ}{\vnbstmt{\forall a, b \in \mathbb{N}.\; ((a + (b + 1)) = ((a + b) + 1))}}

\vnbbookfile{structure-library/list-recursion}{structure-library/list-recursion}
\noindent{\small list-recursion.scm -- CONS, and the recursive characterization of LENGTH.}\par
\vnbdef{empty-in-tuples}{\vnbstmt{\forall a.\; (\langle  \rangle \in \operatorname{tuples}(a))}}
\vnbdef{cons-in-tuples}{\vnbstmt{\forall a, x, l.\; (((x \in a) \wedge (l \in \operatorname{tuples}(a))) \Rightarrow (\operatorname{cons}(x, l) \in \operatorname{tuples}(a)))}}
\vnbdef{nth-cons-1}{\vnbstmt{\forall x \in \mathbf{Set}, l.\; (\operatorname{nth}(1, \operatorname{cons}(x, l)) = x)}}
\vnbdef{nth-cons-succ}{\vnbstmt{\begin{array}{@{}l@{}} \forall i, x, l. \\ \quad ((x \in \mathbf{Set}) \wedge ((i \in \mathbb{N}) \wedge ((1 \leq i) \wedge (i \leq \operatorname{length}(l))))) \Rightarrow  \\ \quad \quad \operatorname{nth}((i + 1), \operatorname{cons}(x, l)) = \operatorname{nth}(i, l) \end{array}}}
\vnbdef{length-cons}{\vnbstmt{\forall x \in \mathbf{Set}, l.\; (\operatorname{length}(\operatorname{cons}(x, l)) = (\operatorname{length}(l) + 1))}}
\vnbdef{makeset-cons}{\vnbstmt{\begin{array}{@{}l@{}} \forall x \in \mathbf{Set}, l. \\ \quad \operatorname{make\text{-}set}(\operatorname{cons}(x, l)) = (\{x\} \cup \operatorname{make\text{-}set}(l)) \end{array}}}
\vnbaxiom{tuple-cons-decompose}{\vnbstmt{\begin{array}{@{}l@{}} \forall a, n, l. \\ \quad ((n \in \mathbb{N}) \wedge ((l \in \operatorname{tuples}(a)) \wedge (\operatorname{length}(l) = (n + 1)))) \Rightarrow  \\ \quad \quad \exists x \in a, m \in \operatorname{tuples}(a).\; ((\operatorname{length}(m) = n) \wedge (l = \operatorname{cons}(x, m))) \end{array}}}
\vnbaxiom{tuple-length-zero}{\vnbstmt{\forall a.; \forall l \in \operatorname{tuples}(a).\; ((\operatorname{length}(l) = 0) \Rightarrow (l = \langle  \rangle))}}

\vnbbookfile{structure-library/sqn}{structure-library/sqn}
\noindent{\small sqn.scm -- SQN(a): the sequences in a, i.e. FUN(NN, a).}\par
\par\medskip\noindent\textbf{Definition} (\texttt{sqn}).\;$a^{\mathbb{N}} :=$ $(\mathbb{N} \to a)$
\vnbdef{sqn-membership}{\vnbstmt{\forall a, f.\; ((f \in a^{\mathbb{N}}) \Leftrightarrow (f \in (\mathbb{N} \to a)))}}
\vnbdef{sqn-sethood}{\vnbstmt{\forall a \in \mathbf{Set}.\; (a^{\mathbb{N}} \in \mathbf{Set})}}

\vnbbookfile{structure-library/order-predicates}{structure-library/order-predicates}
\noindent{\small order-predicates.scm -- derived order predicates on the kernel $<$=.}\par
\par\medskip\noindent\textbf{Definition} (\texttt{<}).\;For all $x$ and $y$,  $(x < y) \Leftrightarrow ((x \leq y) \wedge \neg (x = y))$
\par\medskip\noindent\textbf{Definition} (\texttt{pos-rr}).\;For every $r$,  $\operatorname{pos\text{-}rr}(r) \Leftrightarrow ((r \in \mathbb{R}) \wedge ((0 \leq r) \wedge \neg (0 = r)))$

\vnbbookfile{structure-library/nn-minus}{structure-library/nn-minus}
\noindent{\small nn-minus.scm -- truncated natural subtraction (monus), as VOCABULARY.}\par
\par\medskip\noindent\textbf{Definition} (\texttt{nn-minus}).\;Suppose $k \in \mathbb{N}$ and $l \in \mathbb{N}$. Then $\operatorname{nn\text{-}minus}(k, l) = \left(\text{if }(l \leq k)\text{ then }-(k, l)\text{ else }0\right)$

\vnbbookfile{structure-library/order-lemmas}{structure-library/order-lemmas}
\noindent{\small RETIRED 2026-09-17 (proven): nn-minus-succ-1, one-in-interval, succ-not-one, nn-minus-1-inj, nn-succ-le-cancel -- theorem-library/rake-intervals.scm, modulo 0.  succ-in-interval and pred-in-interval were FALSE as stated (the interval's upper bound was not typed: q = 3/2 with succ(3/2) interpreted as 0 refutes the ...}\par
\vnbaxiom{recip-defined-nonzero}{\vnbstmt{\forall q.\; ((\operatorname{recip}(q) = \operatorname{recip}(q)) \Rightarrow \neg (q = 0))}}

\vnbbookfile{structure-library/real-powers}{structure-library/real-powers}
\noindent{\small real-powers.scm -- the real power RPOW and the square root SQRT.}\par
\par\medskip\noindent\textbf{Definition} (\texttt{sqrt}).\;$\sqrt{\operatorname{a\_}} :=$ $(\iota\, \operatorname{x\_}.\;((\operatorname{x\_} \in \mathbb{R}) \wedge ((0 \leq \operatorname{x\_}) \wedge ((\operatorname{x\_} \cdot \operatorname{x\_}) = \operatorname{a\_}))))$

\vnbbookfile{structure-library/complex}{structure-library/complex}
\noindent{\small complex.scm -- CC as a metric space; completeness axiom}\par
\par\medskip\noindent\textbf{Definition} (\texttt{cc-ms}).\;
\begin{equation*}
\operatorname{cc\text{-}ms} = \langle \mathbb{C}, (\lambda \langle x, y \rangle \in (\mathbb{C} \times \mathbb{C}).\; \operatorname{magnitude}(-(x, y))) \rangle
\end{equation*}

\vnbdef{magnitude-def}{\vnbstmt{\begin{array}{@{}l@{}} \forall z. \\ \quad (\operatorname{magnitude}(z) \simeq \sqrt{((\operatorname{real\text{-}part}(z) \cdot \operatorname{real\text{-}part}(z)) + (\operatorname{imag\text{-}part}(z) \cdot \operatorname{imag\text{-}part}(z)))}) \end{array}}\par\noindent{\small\textit{Comment.} The right-hand side is defined iff $z$ is a complex number, so $\operatorname{magnitude}(z)$ is defined iff $z$ is a complex number.}\par
}
\vnbdef{cc-ms@dist}{\vnbstmt{(\operatorname{dist}(\operatorname{cc\text{-}ms}) \simeq (\lambda \langle x, y \rangle \in (\mathbb{C} \times \mathbb{C}).\; \operatorname{magnitude}((x - y))))}}
\vnbdef{cc-ms@pts}{\vnbstmt{(\operatorname{pts}(\operatorname{cc\text{-}ms}) \simeq \mathbb{C})}}

\vnbbookfile{structure-library/reduce}{structure-library/reduce}
\noindent{\small reduce.scm -- REDUCE: left-fold of a binary op over a finite indexed family.  Bridges the kiddie n-ary surface form (+ x y z ...) and the adult finite-sum operator SUM-AG.}\par
\par\medskip\noindent\textbf{Definition} (\texttt{reduce}).\;For all $\mathit{op}$ and $f$,  $(\operatorname{reduce}(\mathit{op}, f, 1) \simeq f(0))$\par
Let $\mathit{op}$ and $f$ be arbitrary. Let $n \text{ be s.t. } ((n \in \mathbb{N}) \wedge (1 \leq n))$. Then $(\operatorname{reduce}(\mathit{op}, f, \operatorname{succ}(n)) \simeq \operatorname{op}(\operatorname{reduce}(\mathit{op}, f, n), f(n)))$
\par\medskip\noindent\textbf{Definition} (\texttt{fam-of-list}).\;Let $l$ be arbitrary. Let $i \text{ be s.t. } ((i \in \mathbb{N}) \wedge (\operatorname{succ}(i) \leq \operatorname{length}(l)))$. Then $\operatorname{fam\text{-}of\text{-}list}(l)(i) = \operatorname{nth}(\operatorname{succ}(i), l)$

\vnbbookfile{structure-library/ordinals}{structure-library/ordinals}
\noindent{\small ordinals.scm -- ordinal numbers: axioms, ordering, induction, recursion}\par
\vnbdef{ord-lt-iff}{\vnbstmt{\forall a, b.\; (\operatorname{ord\text{-}lt}(a, b) \Leftrightarrow (\operatorname{ord\text{-}le}(a, b) \wedge \neg (a = b)))}}
\vnbdef{limit-ord-iff}{\vnbstmt{\begin{array}{@{}l@{}} \forall a. \\ \quad \operatorname{limit\text{-}ord}(a) \Leftrightarrow ((a \in \mathrm{Ord}) \wedge (\neg (a = 0) \wedge \neg \exists b \in \mathrm{Ord}.\; (a = (b + 1)))) \end{array}}}
\vnbaxiom{limit-ord-is-sup}{\vnbstmt{\forall a.\; (\operatorname{limit\text{-}ord}(a) \Rightarrow (\operatorname{sup\text{-}ord}(\mathrm{OS}(a)) = a))}}
\vnbaxiom{ord-le-total}{\vnbstmt{\forall a, b.\; (((a \in \mathrm{Ord}) \wedge (b \in \mathrm{Ord})) \Rightarrow (\operatorname{ord\text{-}le}(a, b) \vee \operatorname{ord\text{-}le}(b, a)))}}
\vnbaxiom{ord-succ-in}{\vnbstmt{\forall a \in \mathrm{Ord}.\; ((a + 1) \in \mathrm{Ord})}}
\vnbaxiom{ord-succ-nn}{\vnbstmt{\forall n \in \mathbb{N}.\; ((n + 1) = (n + 1))}}
\vnbaxiom{ord-le-antisymm}{\vnbstmt{\forall a, b.\; ((\operatorname{ord\text{-}le}(a, b) \wedge \operatorname{ord\text{-}le}(b, a)) \Rightarrow (a = b))}}
\vnbaxiom{transfinite-induction}{\vnbstmt{\begin{array}{@{}l@{}} \forall c.; \forall a \in \mathrm{Ord}. \\ \quad \forall a \in \mathrm{Ord}.\; (\forall b.\; (\operatorname{ord\text{-}lt}(b, a) \Rightarrow (b \in c)) \Rightarrow (a \in c)) \Rightarrow  \\ \quad \quad a \in c \end{array}}}
\vnbaxiom{ord-segment-limit}{\vnbstmt{\begin{array}{@{}l@{}} \forall a, x. \\ \quad (\operatorname{limit\text{-}ord}(a) \wedge \operatorname{ord\text{-}lt}(x, a)) \Rightarrow  \\ \quad \quad \exists b.\; (\operatorname{ord\text{-}lt}(x, b) \wedge \operatorname{ord\text{-}lt}(b, a)) \end{array}}}
\vnbaxiom{burali-forti}{\vnbstmt{\neg (\mathrm{Ord} \in \mathbf{Set})}}
\vnbaxiom{sup-ord-upper}{\vnbstmt{\begin{array}{@{}l@{}} \forall a \in \mathbf{Set}, b \in a. \\ \quad \forall x \in a.\; (x \in \mathrm{Ord}) \Rightarrow \operatorname{ord\text{-}le}(b, \operatorname{sup\text{-}ord}(a)) \end{array}}}
\vnbaxiom{sup-ord-in}{\vnbstmt{\forall a \in \mathbf{Set}.\; (\forall x \in a.\; (x \in \mathrm{Ord}) \Rightarrow (\operatorname{sup\text{-}ord}(a) \in \mathrm{Ord}))}}
\vnbaxiom{ord-segment-succ}{\vnbstmt{\begin{array}{@{}l@{}} \forall a \in \mathrm{Ord}, x. \\ \quad (x \in \mathrm{OS}((a + 1))) \Leftrightarrow ((x \in \mathrm{OS}(a)) \vee (x = a)) \end{array}}}
\vnbaxiom{ord-segment-membership}{\vnbstmt{\forall a \in \mathrm{Ord}, x.\; ((x \in \mathrm{OS}(a)) \Leftrightarrow \operatorname{ord\text{-}lt}(x, a))}}
\vnbaxiom{ord-segment-is-set}{\vnbstmt{\forall a \in \mathrm{Ord}.\; (\mathrm{OS}(a) \in \mathbf{Set})}}
\vnbaxiom{ord-succ-above}{\vnbstmt{\forall a \in \mathrm{Ord}.\; \operatorname{ord\text{-}lt}(a, (a + 1))}}
\vnbaxiom{ord-le-closure}{\vnbstmt{\forall a, b.\; (\operatorname{ord\text{-}le}(a, b) \Rightarrow ((a \in \mathrm{Ord}) \wedge (b \in \mathrm{Ord})))}}
\vnbaxiom{ord-succ-immediate}{\vnbstmt{\begin{array}{@{}l@{}} \forall a, b. \\ \quad ((a \in \mathrm{Ord}) \wedge ((b \in \mathrm{Ord}) \wedge \operatorname{ord\text{-}lt}(a, b))) \Rightarrow \operatorname{ord\text{-}le}((a + 1), b) \end{array}}}
\vnbaxiom{nn-subset-ord}{\vnbstmt{\forall n \in \mathbb{N}.\; (n \in \mathrm{Ord})}}
\vnbaxiom{sup-ord-empty}{\vnbstmt{\operatorname{sup\text{-}ord}(\emptyset) = 0}}
\vnbaxiom{sup-ord-succ-segment}{\vnbstmt{\forall a \in \mathrm{Ord}.\; (\operatorname{sup\text{-}ord}(\mathrm{OS}((a + 1))) = a)}}
\vnbaxiom{ord-le-nn-compat}{\vnbstmt{\forall m, n.\; (((m \in \mathbb{N}) \wedge (n \in \mathbb{N})) \Rightarrow (\operatorname{ord\text{-}le}(m, n) \Leftrightarrow (m \leq n)))}}
\vnbaxiom{ord-zero-least}{\vnbstmt{\forall a \in \mathrm{Ord}.\; \operatorname{ord\text{-}le}(0, a)}}
\vnbaxiom{ord-succ-injective}{\vnbstmt{\begin{array}{@{}l@{}} \forall a, b. \\ \quad ((a \in \mathrm{Ord}) \wedge ((b \in \mathrm{Ord}) \wedge ((a + 1) = (b + 1)))) \Rightarrow (a = b) \end{array}}}
\vnbaxiom{sup-ord-least}{\vnbstmt{\begin{array}{@{}l@{}} \forall a \in \mathbf{Set}, b \in \mathrm{Ord}. \\ \quad (\forall x \in a.\; (x \in \mathrm{Ord}) \wedge \forall c \in a.\; \operatorname{ord\text{-}le}(c, b)) \Rightarrow  \\ \quad \quad \operatorname{ord\text{-}le}(\operatorname{sup\text{-}ord}(a), b) \end{array}}}
\vnbaxiom{ord-le-trans}{\vnbstmt{\forall a, b, c.\; ((\operatorname{ord\text{-}le}(a, b) \wedge \operatorname{ord\text{-}le}(b, c)) \Rightarrow \operatorname{ord\text{-}le}(a, c))}}
\vnbaxiom{ord-segment-zero}{\vnbstmt{\mathrm{OS}(0) = \emptyset}}
\vnbaxiom{ord-le-refl}{\vnbstmt{\forall a \in \mathrm{Ord}.\; \operatorname{ord\text{-}le}(a, a)}}

\vnbbookfile{structure-library/bijection}{structure-library/bijection}
\noindent{\small RETIRED 2026-09-17 (proven): inverse-bij-in-fun, inverse-bij-left, inverse-bij-right, inverse-bij-is-bijection -- were UNWARRANTED axioms; theorem-library/rake-inverse-bij.scm, modulo 0. RETIRED 2026-09-17 (proven): bijection-compose (was an `informal' axiom) -- theorem-library/rake-analysis2.scm, via the new ...}\par
\par\medskip\noindent\textbf{Definition} (\texttt{inverse-bij}).\;$\varphi^{-1} :=$ $(\lambda \operatorname{y\_} \in y.\; \varepsilon(\{\operatorname{x\_} \in x : (\varphi(\operatorname{x\_}) = \operatorname{y\_})\}))$
\vnbdef{bijection-membership-iff}{\vnbstmt{\begin{array}{@{}l@{}} \forall x, y, i. \\ \quad i \in \mathrm{Bij}(x, y) \\ \quad \Leftrightarrow\; i \in (x \to y) \wedge \\ \quad \forall a, b \in x.\; ((i(a) = i(b)) \Rightarrow (a = b)) \wedge \\ \quad \forall w \in y.\; \exists z \in x.\; (i(z) = w) \end{array}}}

\vnbbookfile{structure-library/cardinality}{structure-library/cardinality}
\noindent{\small cardinality.scm -- CARD, the cardinality of a set, DEFINED.}\par
\par\medskip\noindent\textbf{Definition} (\texttt{card}).\;$|\operatorname{a\_}| :=$
\begin{equation*}
(\iota\, \alpha.\;((\alpha \in \mathrm{Ord}) \wedge (\exists \varphi.\; (\varphi \in \mathrm{Bij}(\operatorname{a\_}, \mathrm{OS}(\alpha))) \wedge \forall \beta.\; (\operatorname{ord\text{-}lt}(\beta, \alpha) \Rightarrow \neg \exists \psi.\; (\psi \in \mathrm{Bij}(\operatorname{a\_}, \mathrm{OS}(\beta)))))))
\end{equation*}


\vnbbookfile{structure-library/injection}{structure-library/injection}
\noindent{\small RETIRED 2026-09-17 (proven): choose-n-0, choose-0-succ -- theorem-library/rake-combinatorics.scm. choose-succ (Pascal) stays: needs the CHOOSE-SET split + card-image-injection (\textasciitilde{}150 lines). RETIRED 2026-09-17 (proven): falling-in-nn -- theorem-library/rake-finsum-typing-nn.scm. choose-in-nn stays: it bottoms out in the ...}\par
\par\medskip\noindent\textbf{Definition} (\texttt{choose-set}).\;$\operatorname{choose\text{-}set}(n, m) :=$ $\{a \in \operatorname{power}(\mathrm{OS}(n)) : (|a| = m)\}$
\par\medskip\noindent\textbf{Definition} (\texttt{injective-star}).\;For every $f$, 
\begin{equation*}
\operatorname{injective\text{-}star}(f) \Leftrightarrow \forall \operatorname{u\_}, \operatorname{v\_}.\; ((f(\operatorname{u\_}) = f(\operatorname{v\_})) \Rightarrow (\operatorname{u\_} = \operatorname{v\_}))
\end{equation*}

\par\medskip\noindent\textbf{Definition} (\texttt{permutations}).\;$\operatorname{permutations}(n) :=$ $\operatorname{injection}(\mathrm{OS}(n), \mathrm{OS}(n))$
\par\medskip\noindent\textbf{Definition} (\texttt{factorial}).\; $\operatorname{factorial}(0) = 1$\par
Suppose $k \in \mathbb{N}$. Then $\operatorname{factorial}(\operatorname{succ}(k)) = \cdot(\operatorname{succ}(k), \operatorname{factorial}(k))$
\par\medskip\noindent\textbf{Definition} (\texttt{choose}).\;$\operatorname{choose}(n, m) :=$ $|\operatorname{choose\text{-}set}(n, m)|$
\par\medskip\noindent\textbf{Definition} (\texttt{falling}).\;For every $n$,  $(\operatorname{falling}(n, 0) \simeq 1)$\par
Let $n$ be arbitrary. Let $k \in \mathbb{N}$. Then $(\operatorname{falling}(n, \operatorname{succ}(k)) \simeq \cdot(-(n, k), \operatorname{falling}(n, k)))$
\vnbdef{injection-membership-iff}{\vnbstmt{\begin{array}{@{}l@{}} \forall x, y, i. \\ \quad (i \in \operatorname{injection}(x, y)) \Leftrightarrow ((i \in (x \to y)) \wedge \forall a, b \in x.\; ((i(a) = i(b)) \Rightarrow (a = b))) \end{array}}}
\vnbdef{image-membership-iff}{\vnbstmt{\forall i, s, w.\; ((w \in \operatorname{image}(i, s)) \Leftrightarrow ((w \in \mathbf{Set}) \wedge \exists x \in s.\; (i(x) = w)))}}
\vnbaxiom{image-set}{\vnbstmt{\forall i.; \forall s \in \mathbf{Set}.\; (\operatorname{image}(i, s) \in \mathbf{Set})}}

\vnbbookfile{structure-library/compose-typing}{structure-library/compose-typing}
\noindent{\small RETIRED 2026-09-17 (proven): range-membership, fun-range-membership, ran-subset-codomain, compose-type-2/3/4/5 -- theorem-library/rake-compose-typing.scm, all modulo 0 from dom-membership / dom-fun-membership / fun-codomain-iff / image-membership-iff.  They had been `proof'-warranted supports here; the}\par
\par\medskip\noindent\textbf{Definition} (\texttt{ran}).\;$\operatorname{ran}(f) :=$ $\operatorname{image}(f, \operatorname{dom}(f))$

\vnbbookfile{structure-library/inf-subsets}{structure-library/inf-subsets}
\noindent{\small inf-subsets.scm -- the class of infinite subsets of a given set.}\par
\par\medskip\noindent\textbf{Definition} (\texttt{inf-subsets}).\;$\mathrm{Inf}(a) :=$ $\{s \in \operatorname{power}(a) : \neg (|s| \in \mathbb{N})\}$
\vnbdef{inf-subsets-membership}{\vnbstmt{\forall a, s.\; ((s \in \mathrm{Inf}(a)) \Leftrightarrow ((s \subseteq a) \wedge \neg (|s| \in \mathbb{N})))}}

\vnbbookfile{structure-library/finsum}{structure-library/finsum}
\noindent{\small finsum.scm -- FINSUM: summation of an abelian-group-valued function over an arbitrary finite set.}\par
\par\medskip\noindent\textbf{Definition} (\texttt{finsum}).\;$\operatorname{finsum}(\mathit{ag}, f, s) :=$ $\operatorname{sum\text{-}ag}(\mathit{ag}, \operatorname{enum\text{-}fam}(\mathit{ag}, f, \operatorname{fin\text{-}enum}(s), |s|), |s|)$
\par\medskip\noindent\textbf{Definition} (\texttt{fin-enum}).\;$\operatorname{fin\text{-}enum}(s) :=$ $\varepsilon(\mathrm{Bij}(\mathrm{OS}(|s|), s))$
\par\medskip\noindent\textbf{Definition} (\texttt{enum-fam}).\;$\operatorname{enum\text{-}fam}(\mathit{ag}, f, \varphi, n) :=$
\begin{equation*}
(\lambda i \in \mathbb{N}.\; \left(\text{if }(i \in \mathrm{OS}(n))\text{ then }f(\varphi(i))\text{ else }\operatorname{iden}(\mathit{ag})\right))
\end{equation*}

\par\medskip\noindent\textbf{Definition} (\texttt{insert-last}).\;$\operatorname{insert\text{-}last}(\varphi, x, n) :=$ $(\lambda i \in \mathbb{N}.\; \left(\text{if }(i \in \mathrm{OS}(n))\text{ then }\varphi(i)\text{ else }x\right))$
\par\medskip\noindent\textbf{Definition} (\texttt{sum-set}).\;$\operatorname{sum\text{-}set}(r, s, f) :=$ $\operatorname{finsum}(\operatorname{ring\text{-}additive\text{-}ag}(r), f, s)$

\vnbbookfile{structure-library/matrix}{structure-library/matrix}
\noindent{\small RETIRED 2026-09-17 (proven): finsum-two-support -- theorem-library/rake-finsum-laws2.scm RETIRED 2026-09-17 (proven): identmat-left-identity, identmat-right-identity -- theorem-library/rake-identmat.scm (matmul-entry + finsum-single-support + extensionality; no guard needed). RETIRED 2026-09-17 (proven): ...}\par
\par\medskip\noindent\textbf{Definition} (\texttt{size}).\;$\operatorname{size}(m) :=$ $\langle \operatorname{length}(m), \left(\text{if }(\operatorname{length}(m) = 0)\text{ then }0\text{ else }\operatorname{length}(\operatorname{nth}(1, m))\right) \rangle$
\par\medskip\noindent\textbf{Definition} (\texttt{entry}).\;$\operatorname{entry}(m, i, j) :=$ $\operatorname{nth}(j, \operatorname{nth}(i, m))$
\par\medskip\noindent\textbf{Definition} (\texttt{mat}).\;$\operatorname{mat}(m, n, x) :=$
\begin{equation*}
\{p \in \operatorname{matrix}(x) : ((n \in \mathbb{N}) \wedge (((\operatorname{length}(p) = 0) \Rightarrow (m = 0)) \wedge (\neg (\operatorname{length}(p) = 0) \Rightarrow (\operatorname{size}(p) = \langle m, n \rangle))))\}
\end{equation*}

\par\medskip\noindent\textbf{Definition} (\texttt{identmat}).\;$\operatorname{identmat}(a, n) :=$
\begin{equation*}
\begin{array}{@{}l@{}} \operatorname{matof}(n, n, \\ \quad (\lambda \langle i, j \rangle \in (\operatorname{interval}(1, n) \times \operatorname{interval}(1, n)).\; \left(\text{if }(i = j)\text{ then }\operatorname{one}(a)\text{ else }\operatorname{zero}(a)\right))) \end{array}
\end{equation*}

\par\medskip\noindent\textbf{Definition} (\texttt{interval}).\;$\operatorname{interval}(a, b) :=$ $\{i \in \mathbb{N} : ((a \leq i) \wedge (i \leq b))\}$
\par\medskip\noindent\textbf{Definition} (\texttt{matof}).\;$\operatorname{matof}(m, n, g) :=$
\begin{equation*}
(\iota\, p.\;((p \in \operatorname{mat}(m, n, \operatorname{image}(g, (\operatorname{interval}(1, m) \times \operatorname{interval}(1, n))))) \wedge \forall i \in \operatorname{interval}(1, m), j \in \operatorname{interval}(1, n).\; (\operatorname{entry}(p, i, j) = g(i, j))))
\end{equation*}

\par\medskip\noindent\textbf{Definition} (\texttt{matmul}).\;$\operatorname{matmul}(a, p, q) :=$
\begin{equation*}
\begin{array}{@{}l@{}} \operatorname{matof}(\operatorname{nth}(1, \operatorname{size}(p)), \operatorname{nth}(2, \operatorname{size}(q)), \\ \quad (\lambda \langle i, k \rangle \in (\operatorname{interval}(1, \operatorname{nth}(1, \operatorname{size}(p))) \times \operatorname{interval}(1, \operatorname{nth}(2, \operatorname{size}(q)))).\; \\ \quad \quad \operatorname{finsum}(\operatorname{ring\text{-}additive\text{-}ag}(a), \\ \quad \quad \quad (\lambda j \in \operatorname{interval}(1, \operatorname{nth}(2, \operatorname{size}(p))).\; \operatorname{mul}(a)(\operatorname{entry}(p, i, j), \operatorname{entry}(q, j, k))), \\ \quad \quad \quad \operatorname{interval}(1, \operatorname{nth}(2, \operatorname{size}(p)))))) \end{array}
\end{equation*}

\par\medskip\noindent\textbf{Definition} (\texttt{mat-ring}).\;$\operatorname{mat\text{-}ring}(a, n) :=$
\begin{equation*}
\langle \operatorname{mat}(n, n, \operatorname{carr}(a)), (\lambda \langle p, q \rangle \in (\operatorname{mat}(n, n, \operatorname{carr}(a)) \times \operatorname{mat}(n, n, \operatorname{carr}(a))).\; \operatorname{matadd}(a, p, q)), (\lambda \langle p, q \rangle \in (\operatorname{mat}(n, n, \operatorname{carr}(a)) \times \operatorname{mat}(n, n, \operatorname{carr}(a))).\; \operatorname{matmul}(a, p, q)), (\lambda p \in \operatorname{mat}(n, n, \operatorname{carr}(a)).\; \operatorname{matneg}(a, p)), \operatorname{zeromat}(a, n, n), \operatorname{identmat}(a, n) \rangle
\end{equation*}

\par\medskip\noindent\textbf{Definition} (\texttt{zeromat}).\;$\operatorname{zeromat}(a, m, n) :=$
\begin{equation*}
\operatorname{matof}(m, n, (\lambda \langle i, j \rangle \in (\operatorname{interval}(1, m) \times \operatorname{interval}(1, n)).\; \operatorname{zero}(a)))
\end{equation*}

\par\medskip\noindent\textbf{Definition} (\texttt{matadd}).\;$\operatorname{matadd}(a, p, q) :=$
\begin{equation*}
\begin{array}{@{}l@{}} \operatorname{matof}(\operatorname{nth}(1, \operatorname{size}(p)), \operatorname{nth}(2, \operatorname{size}(p)), \\ \quad (\lambda \langle i, j \rangle \in (\operatorname{interval}(1, \operatorname{nth}(1, \operatorname{size}(p))) \times \operatorname{interval}(1, \operatorname{nth}(2, \operatorname{size}(p)))).\; \\ \quad \quad \operatorname{add}(a)(\operatorname{entry}(p, i, j), \operatorname{entry}(q, i, j)))) \end{array}
\end{equation*}

\par\medskip\noindent\textbf{Definition} (\texttt{matneg}).\;$\operatorname{matneg}(a, p) :=$
\begin{equation*}
\begin{array}{@{}l@{}} \operatorname{matof}(\operatorname{nth}(1, \operatorname{size}(p)), \operatorname{nth}(2, \operatorname{size}(p)), \\ \quad (\lambda \langle i, j \rangle \in (\operatorname{interval}(1, \operatorname{nth}(1, \operatorname{size}(p))) \times \operatorname{interval}(1, \operatorname{nth}(2, \operatorname{size}(p)))).\; \\ \quad \quad \operatorname{neg}(a)(\operatorname{entry}(p, i, j)))) \end{array}
\end{equation*}

\par\medskip\noindent\textbf{Definition} (\texttt{matscale}).\;$\operatorname{matscale}(a, r, p) :=$
\begin{equation*}
\begin{array}{@{}l@{}} \operatorname{matof}(\operatorname{nth}(1, \operatorname{size}(p)), \operatorname{nth}(2, \operatorname{size}(p)), \\ \quad (\lambda \langle i, j \rangle \in (\operatorname{interval}(1, \operatorname{nth}(1, \operatorname{size}(p))) \times \operatorname{interval}(1, \operatorname{nth}(2, \operatorname{size}(p)))).\; \\ \quad \quad \operatorname{mul}(a)(r, \operatorname{entry}(p, i, j)))) \end{array}
\end{equation*}

\par\medskip\noindent\textbf{Definition} (\texttt{unitrow}).\;$\operatorname{unitrow}(a, n, i) :=$
\begin{equation*}
\begin{array}{@{}l@{}} \operatorname{matof}(1, n, \\ \quad (\lambda \langle \mathit{rw}, \mathit{cl} \rangle \in (\operatorname{interval}(1, 1) \times \operatorname{interval}(1, n)).\; \left(\text{if }(\mathit{cl} = i)\text{ then }\operatorname{one}(a)\text{ else }\operatorname{zero}(a)\right))) \end{array}
\end{equation*}

\par\medskip\noindent\textbf{Definition} (\texttt{block}).\;$\operatorname{block}(p, k, l) :=$
\begin{equation*}
\operatorname{matof}(k, l, (\lambda \langle i, j \rangle \in (\operatorname{interval}(1, k) \times \operatorname{interval}(1, l)).\; \operatorname{entry}(p, i, j)))
\end{equation*}

\par\medskip\noindent\textbf{Definition} (\texttt{snoc-col}).\;$\operatorname{snoc\text{-}col}(w, n, x) :=$
\begin{equation*}
\begin{array}{@{}l@{}} \operatorname{matof}((n + 1), 1, \\ \quad (\lambda \langle \operatorname{i\_}, \operatorname{j\_} \rangle \in (\operatorname{interval}(1, (n + 1)) \times \operatorname{interval}(1, 1)).\; \\ \quad \quad \left(\text{if }(\operatorname{i\_} = (n + 1))\text{ then }x\text{ else }\operatorname{entry}(w, \operatorname{i\_}, 1)\right))) \end{array}
\end{equation*}

\par\medskip\noindent\textbf{Definition} (\texttt{snoc-row}).\;$\operatorname{snoc\text{-}row}(c, n, r) :=$
\begin{equation*}
\begin{array}{@{}l@{}} \operatorname{matof}(1, (n + 1), \\ \quad (\lambda \langle \operatorname{i\_}, \operatorname{j\_} \rangle \in (\operatorname{interval}(1, 1) \times \operatorname{interval}(1, (n + 1))).\; \\ \quad \quad \left(\text{if }(\operatorname{j\_} = (n + 1))\text{ then }r\text{ else }\operatorname{entry}(c, 1, \operatorname{j\_})\right))) \end{array}
\end{equation*}

\vnbdef{matrix-membership}{\vnbstmt{\begin{array}{@{}l@{}} \forall s, m. \\ \quad m \in \operatorname{matrix}(s) \\ \quad \Leftrightarrow\; m \in \operatorname{tuples}(\operatorname{tuples}(s)) \wedge \\ \quad \forall i, j \in \mathbb{N}. \\ \quad \quad (((1 \leq i) \wedge (i \leq \operatorname{length}(m))) \wedge ((1 \leq j) \wedge (j \leq \operatorname{length}(m)))) \Rightarrow  \\ \quad \quad \quad \operatorname{length}(\operatorname{nth}(i, m)) = \operatorname{length}(\operatorname{nth}(j, m)) \end{array}}}

\vnbbookfile{theorem-library/ord-well-ordered-proof}{theorem-library/ord-well-ordered-proof}
\noindent{\small ord-well-ordered-proof.scm -- the well-ordering of ORD, PROVEN.}\par
\vnbprop{ord-well-ordered}{\vnbstmt{\begin{array}{@{}l@{}} \forall l. \\ \quad (\forall x \in l.\; (x \in \mathrm{Ord}) \wedge \exists w.\; (w \in l)) \Rightarrow  \\ \quad \quad \exists m \in l.\; \forall k \in l.\; \operatorname{ord\text{-}le}(m, k) \end{array}}}

\vnbbookfile{theorem-library/equality-basics}{theorem-library/equality-basics}
\noindent{\small equality-basics.scm -- eq-sym, eq-trans, neq-sym, PROVEN.}\par
\vnbprop{eq-sym}{\vnbstmt{\forall a, b.\; ((a = b) \Rightarrow (b = a))}}
\vnbprop{eq-trans}{\vnbstmt{\forall a, b, c.\; (((a = b) \wedge (b = c)) \Rightarrow (a = c))}}
\vnbprop{neq-sym}{\vnbstmt{\forall a, b.\; (\neg (a = b) \Rightarrow \neg (b = a))}}

\vnbbookfile{theorem-library/factorial-in-nn}{theorem-library/factorial-in-nn}
\noindent{\small factorial-in-nn.scm -- n! is a natural number, PROVEN. Batch 27-A follow-on, 2026-09-24.}\par
\vnbprop{factorial-in-nn}{\vnbstmt{\forall n \in \mathbb{N}.\; (\operatorname{factorial}(n) \in \mathbb{N})}}

\vnbbookfile{theorem-library/rake-compose-typing}{theorem-library/rake-compose-typing}
\noindent{\small rake-compose-typing.scm -- the seven RAN / composition typing supports of structure-library/compose-typing.scm, PROVEN modulo 0.}\par
\vnbprop{range-membership}{\vnbstmt{\forall f.; \forall a \in \operatorname{dom}(f).\; (f(a) \in \operatorname{ran}(f))}}
\vnbprop{fun-range-membership}{\vnbstmt{\forall x, f, a.\; (((f \in \mathrm{Fun}(x)) \wedge (a \in x)) \Rightarrow (f(a) \in \operatorname{ran}(f)))}}
\vnbprop{ran-subset-codomain}{\vnbstmt{\forall x, y.; \forall f \in (x \to y).\; (\operatorname{ran}(f) \subseteq y)}}
\vnbprop{compose-type-2}{\vnbstmt{\begin{array}{@{}l@{}} \forall x, f, g, a. \\ \quad ((f \in \mathrm{Fun}(x)) \wedge ((g \in \mathrm{Fun}(\operatorname{ran}(f))) \wedge (a \in x))) \Rightarrow (g(f(a)) \in \operatorname{ran}(g)) \end{array}}}
\vnbprop{compose-type-3}{\vnbstmt{\begin{array}{@{}l@{}} \forall x, f, g, h, a. \\ \quad ((f \in \mathrm{Fun}(x)) \wedge ((g \in \mathrm{Fun}(\operatorname{ran}(f))) \wedge ((h \in \mathrm{Fun}(\operatorname{ran}(g))) \wedge (a \in x)))) \Rightarrow  \\ \quad \quad h(g(f(a))) \in \operatorname{ran}(h) \end{array}}}
\vnbprop{compose-type-4}{\vnbstmt{\begin{array}{@{}l@{}} \forall x, f, g, h, k, a. \\ \quad ((f \in \mathrm{Fun}(x)) \wedge ((g \in \mathrm{Fun}(\operatorname{ran}(f))) \wedge ((h \in \mathrm{Fun}(\operatorname{ran}(g))) \wedge ((k \in \mathrm{Fun}(\operatorname{ran}(h))) \wedge (a \in x))))) \Rightarrow  \\ \quad \quad k(h(g(f(a)))) \in \operatorname{ran}(k) \end{array}}}
\vnbprop{compose-type-5}{\vnbstmt{\begin{array}{@{}l@{}} \forall x, f, g, h, k, m, a. \\ \quad f \in \mathrm{Fun}(x) \wedge \\ \quad g \in \mathrm{Fun}(\operatorname{ran}(f)) \wedge \\ \quad h \in \mathrm{Fun}(\operatorname{ran}(g)) \wedge \\ \quad k \in \mathrm{Fun}(\operatorname{ran}(h)) \wedge \\ \quad m \in \mathrm{Fun}(\operatorname{ran}(k)) \wedge \\ \quad a \in x \Rightarrow \\ \quad \quad m(k(h(g(f(a))))) \in \operatorname{ran}(m) \end{array}}}

\vnbbookfile{theorem-library/co-eq-lt-trans}{theorem-library/co-eq-lt-trans}
\noindent{\small co-eq-lt-trans.scm -- co-eq-lt-trans and co-lt-eq-trans, PROVEN.}\par
\vnbprop{co-eq-lt-trans}{\vnbstmt{\forall a, b, c.\; (((a = b) \wedge (b < c)) \Rightarrow (a < c))}}
\vnbprop{co-lt-eq-trans}{\vnbstmt{\forall a, b, c.\; (((a < b) \wedge (b = c)) \Rightarrow (a < c))}}

\vnbbookfile{theorem-library/interval-basics}{theorem-library/interval-basics}
\noindent{\small interval-basics.scm -- INTERVAL's read-offs, PROVEN from its definition.}\par
\vnbprop{interval-unfold}{\vnbstmt{\forall a, b.\; (\{i \in \mathbb{N} : ((a \leq i) \wedge (i \leq b))\} \simeq \operatorname{interval}(a, b))}}
\vnbprop{interval-in-set}{\vnbstmt{\forall a, b.\; (\operatorname{interval}(a, b) \in \mathbf{Set})}}
\vnbprop{interval-elt-in-nn}{\vnbstmt{\forall a, b.; \forall i \in \operatorname{interval}(a, b).\; (i \in \mathbb{N})}}
\vnbprop{interval-lo}{\vnbstmt{\forall a, b.; \forall i \in \operatorname{interval}(a, b).\; (a \leq i)}}
\vnbprop{interval-hi}{\vnbstmt{\forall a, b.; \forall i \in \operatorname{interval}(a, b).\; (i \leq b)}}

\vnbbookfile{theorem-library/interval-mem-intro}{theorem-library/interval-mem-intro}
\noindent{\small interval-mem-intro.scm -- membership INTRODUCTION for INTERVAL, proven.}\par
\vnbprop{interval-mem-intro}{\vnbstmt{\forall a, b.; \forall i \in \mathbb{N}.\; (((a \leq i) \wedge (i \leq b)) \Rightarrow (i \in \operatorname{interval}(a, b)))}}

\vnbbookfile{theorem-library/ord-segment-nn-succ-proof}{theorem-library/ord-segment-nn-succ-proof}
\noindent{\small ord-segment-nn-succ-proof.scm -- the NN-flavoured ORD-SEGMENT successor characterisation, proven from three ordinal primitives.}\par
\vnbprop{ord-segment-nn-succ}{\vnbstmt{\begin{array}{@{}l@{}} \forall n \in \mathbb{N}, k. \\ \quad (k \in \mathrm{OS}((n + 1))) \Leftrightarrow ((k \in \mathrm{OS}(n)) \vee (k = n)) \end{array}}}

\vnbbookfile{theorem-library/ord-segment-nn-subset-proof}{theorem-library/ord-segment-nn-subset-proof}
\noindent{\small ord-segment-nn-subset-proof.scm -- a member of ORD-SEGMENT(m), m in NN, is a natural number.  PROVEN, retiring the `informal' support of the same name (theorem-library/ord-segment-nn-subset.scm:9).}\par
\vnbprop{ord-segment-zero-no-members}{\vnbstmt{\forall k.\; \neg (k \in \mathrm{OS}(0))}}
\vnbprop{ord-segment-nn-subset}{\vnbstmt{\forall m \in \mathbb{N}, k \in \mathrm{OS}(m).\; (k \in \mathbb{N})}}

\vnbbookfile{theorem-library/inf-subsets-is-set}{theorem-library/inf-subsets-is-set}
\noindent{\small inf-subsets-is-set.scm -- INF-SUBSETS(A) is a set when A is, proven.}\par
\vnbprop{inf-subsets-is-set}{\vnbstmt{\forall a \in \mathbf{Set}.\; (\mathrm{Inf}(a) \in \mathbf{Set})}}

\vnbbookfile{theorem-library/binary-minus-laws}{theorem-library/binary-minus-laws}
\noindent{\small binary-minus-laws.scm -- what follows from `binary-minus-def'.}\par
\vnbprop{rr-sub-in-rr}{\vnbstmt{\forall u, v.\; (((u \in \mathbb{R}) \wedge (v \in \mathbb{R})) \Rightarrow ((u - v) \in \mathbb{R}))}}
\vnbprop{zz-sub-in-zz}{\vnbstmt{\forall u, v.\; (((u \in \mathbb{Z}) \wedge (v \in \mathbb{Z})) \Rightarrow ((u - v) \in \mathbb{Z}))}}
\vnbprop{rr-add-in-rr}{\vnbstmt{\forall a, b.\; (((a \in \mathbb{R}) \wedge (b \in \mathbb{R})) \Rightarrow ((a + b) \in \mathbb{R}))}}
\vnbprop{rr-mul-in-rr}{\vnbstmt{\forall a, b.\; (((a \in \mathbb{R}) \wedge (b \in \mathbb{R})) \Rightarrow ((a \cdot b) \in \mathbb{R}))}}

\vnbbookfile{theorem-library/fun-apply-type-proof}{theorem-library/fun-apply-type-proof}
\noindent{\small fun-apply-type-proof.scm -- f : A -$>$ B and x in A give f(x) in B, PROVEN.}\par
\vnbprop{fun-apply-type-c}{\vnbstmt{\forall f, a, b, x.\; (((f \in (a \to b)) \wedge (x \in a)) \Rightarrow (f(x) \in b))}}

\vnbbookfile{theorem-library/image-subset-codomain}{theorem-library/image-subset-codomain}
\noindent{\small image-subset-codomain.scm -- the image of phi : dm -$>$ cod lies in cod.}\par
\vnbprop{image-subset-codomain}{\vnbstmt{\forall m, d.; \forall i \in (m \to d), w \in \operatorname{image}(i, m).\; (w \in d)}}

\vnbbookfile{theorem-library/rake-dc-on-nn}{theorem-library/rake-dc-on-nn}
\noindent{\small rake-dc-on-nn.scm -- dependent choice on NN, PROVEN.}\par
\par\medskip\noindent\textbf{Definition} (\texttt{dc-iter}).\;For all $s$, $a$ and $\mathit{nxt}$,  $(\operatorname{dc\text{-}iter}(s, a, \mathit{nxt}, 0) \simeq a)$\par
Let $s$, $a$ and $\mathit{nxt}$ be arbitrary. Let $n \in \mathbb{N}$. Then
\begin{equation*}
(\operatorname{dc\text{-}iter}(s, a, \mathit{nxt}, \operatorname{succ}(n)) \simeq \varepsilon(\{\operatorname{y\_} \in s : (\operatorname{y\_} \in \operatorname{nxt}(n, \operatorname{dc\text{-}iter}(s, a, \mathit{nxt}, n)))\}))
\end{equation*}

\vnbprop{dc-iter-in}{\vnbstmt{\begin{array}{@{}l@{}} \forall n \in \mathbb{N}, s.; \forall a \in s, x. \\ \quad \forall k \in \mathbb{N}, u \in s.\; \exists y \in s.\; (y \in x(k, u)) \Rightarrow  \\ \quad \quad \operatorname{dc\text{-}iter}(s, a, x, n) \in s \end{array}}}
\vnbprop{dc-iter-step}{\vnbstmt{\begin{array}{@{}l@{}} \forall n \in \mathbb{N}, s.; \forall a \in s, x. \\ \quad \forall k \in \mathbb{N}, u \in s.\; \exists y \in s.\; (y \in x(k, u)) \Rightarrow  \\ \quad \quad \operatorname{dc\text{-}iter}(s, a, x, (n + 1)) \in x(n, \operatorname{dc\text{-}iter}(s, a, x, n)) \end{array}}}
\vnbprop{dc-on-nn-pred}{\vnbstmt{\begin{array}{@{}l@{}} \forall x \in \mathbf{Set}, a \in x, t. \\ \quad \forall k \in \mathbb{N}, u \in x.\; \exists y \in x.\; (y \in t(k, u)) \Rightarrow  \\ \quad \quad \exists f \in (\mathbb{N} \to x). \\ \quad \quad \quad (f(0) = a) \wedge \forall k \in \mathbb{N}.\; (f((k + 1)) \in t(k, f(k))) \end{array}}}
\vnbprop{dc-on-nn}{\vnbstmt{\begin{array}{@{}l@{}} \forall x \in \mathbf{Set}, a \in x, r \in \mathbf{Set}. \\ \quad \forall k \in \mathbb{N}, u \in x.\; \exists y \in x.\; (\langle k, u, y \rangle \in r) \Rightarrow  \\ \quad \quad \exists f \in (\mathbb{N} \to x). \\ \quad \quad \quad (f(0) = a) \wedge \forall k \in \mathbb{N}.\; (\langle k, f(k), f((k + 1)) \rangle \in r) \end{array}}}

\vnbbookfile{theorem-library/compose-apply-proof}{theorem-library/compose-apply-proof}
\noindent{\small compose-apply-proof.scm -- (COMPOSE f g)(x) = f(g(x)), PROVEN modulo 0.}\par
\vnbprop{compose-apply}{\vnbstmt{\begin{array}{@{}l@{}} \forall a, b, c.; \forall f \in (b \to c), g \in (a \to b), x \in a. \\ \quad \operatorname{compose}(f, g)(x) = f(g(x)) \end{array}}}
\vnbprop{compose-type}{\vnbstmt{\begin{array}{@{}l@{}} \forall a, b, c, f, g. \\ \quad ((a \in \mathbf{Set}) \wedge ((g \in (a \to b)) \wedge (f \in (b \to c)))) \Rightarrow  \\ \quad \quad \operatorname{compose}(f, g) \in (a \to c) \end{array}}}

\vnbbookfile{theorem-library/id-fun-laws}{theorem-library/id-fun-laws}
\noindent{\small id-fun-laws.scm -- the identity map and the composition laws over FUN.}\par
\vnbprop{id-fun-type}{\vnbstmt{\forall a \in \mathbf{Set}.\; (\operatorname{id\text{-}fun}(a) \in (a \to a))}}
\vnbprop{id-fun-apply}{\vnbstmt{\forall a.; \forall x \in a.\; (\operatorname{id\text{-}fun}(a)(x) = x)}}
\vnbprop{compose-id-left}{\vnbstmt{\begin{array}{@{}l@{}} \forall a, b, f. \\ \quad ((a \in \mathbf{Set}) \wedge ((b \in \mathbf{Set}) \wedge (f \in (a \to b)))) \Rightarrow (\operatorname{compose}(\operatorname{id\text{-}fun}(b), f) = f) \end{array}}}
\vnbprop{compose-id-right}{\vnbstmt{\forall a, b, f.\; (((a \in \mathbf{Set}) \wedge (f \in (a \to b))) \Rightarrow (\operatorname{compose}(f, \operatorname{id\text{-}fun}(a)) = f))}}
\vnbprop{compose-assoc}{\vnbstmt{\begin{array}{@{}l@{}} \forall a, b, c, u, f, g, h. \\ \quad ((a \in \mathbf{Set}) \wedge ((b \in \mathbf{Set}) \wedge ((f \in (a \to b)) \wedge ((g \in (b \to c)) \wedge (h \in (c \to u)))))) \Rightarrow  \\ \quad \quad \operatorname{compose}(h, \operatorname{compose}(g, f)) = \operatorname{compose}(\operatorname{compose}(h, g), f) \end{array}}}

\vnbbookfile{theorem-library/pair-tuple-sethood}{theorem-library/pair-tuple-sethood}
\noindent{\small pair-tuple-sethood.scm -- a pair [u,v] is a set, PROVEN modulo 0.}\par
\vnbprop{pair-in-cartesian}{\vnbstmt{\forall x, y, u, v.\; (((u \in x) \wedge (v \in y)) \Rightarrow (\langle u, v \rangle \in (x \times y)))}}
\vnbprop{pair-tuple-is-set}{\vnbstmt{\forall x, y, u, v.\; (((u \in x) \wedge (v \in y)) \Rightarrow (\langle u, v \rangle \in \mathbf{Set}))}}

\vnbbookfile{theorem-library/nn-order-ord}{theorem-library/nn-order-ord}
\noindent{\small nn-order-ord.scm -- the elementary NN order block, PROVEN.  Seven facts that were `well-known' supports in structure-library/order-lemmas.scm are theorems here, every one billing `modulo 0'.}\par
\vnbprop{nn-one-in}{\vnbstmt{1 \in \mathbb{N}}}
\vnbprop{nn-zero-le}{\vnbstmt{\forall n \in \mathbb{N}.\; (0 \leq n)}}
\vnbprop{nn-le-succ}{\vnbstmt{\forall k \in \mathbb{N}.\; (k \leq (k + 1))}}
\vnbprop{nn-le-imp-neq-succ}{\vnbstmt{\forall k, j \in \mathbb{N}.\; ((j \leq k) \Rightarrow \neg (j = (k + 1)))}}
\vnbprop{nn-succ-mono}{\vnbstmt{\forall a, b \in \mathbb{N}.\; ((a \leq b) \Rightarrow ((a + 1) \leq (b + 1)))}}
\vnbprop{nn-one-le-succ}{\vnbstmt{\forall n \in \mathbb{N}.\; (1 \leq (n + 1))}}
\vnbprop{nn-le-succ-cases}{\vnbstmt{\forall k, j \in \mathbb{N}.\; ((j \leq (k + 1)) \Rightarrow ((j \leq k) \vee (j = (k + 1))))}}

\vnbbookfile{theorem-library/nn-not-le-succ-le}{theorem-library/nn-not-le-succ-le}
\noindent{\small nn-not-le-succ-le.scm -- NN is totally ordered and discrete, PROVEN.}\par
\vnbprop{nn-not-le-succ-le}{\vnbstmt{\forall m, n \in \mathbb{N}.\; (\neg (m \leq n) \Rightarrow ((n + 1) \leq m))}}

\vnbbookfile{theorem-library/nn-order-basics}{theorem-library/nn-order-basics}
\noindent{\small nn-order-basics.scm -- the elementary NN order facts, PROVEN.}\par
\vnbprop{nn-in-rr}{\vnbstmt{\forall k \in \mathbb{N}.\; (k \in \mathbb{R})}}
\vnbprop{rr-le-trans}{\vnbstmt{\forall x, y, z \in \mathbb{R}.\; (((x \leq y) \wedge (y \leq z)) \Rightarrow (x \leq z))}}
\vnbprop{rr-le-trans-c}{\vnbstmt{\forall x, y, z \in \mathbb{R}.\; (((x \leq y) \wedge (y \leq z)) \Rightarrow (x \leq z))}}
\vnbprop{nn-le-trans-guarded}{\vnbstmt{\forall a, b, c \in \mathbb{N}.\; (((a \leq b) \wedge (b \leq c)) \Rightarrow (a \leq c))}}
\vnbprop{nn-le-refl}{\vnbstmt{\forall k \in \mathbb{N}.\; (k \leq k)}}
\vnbprop{nn-le-add-right}{\vnbstmt{\forall n, m \in \mathbb{N}.\; (m \leq (m + n))}}
\vnbprop{nn-pair-upper-bound}{\vnbstmt{\begin{array}{@{}l@{}} \forall a, b. \\ \quad ((a \in \mathbb{N}) \wedge (b \in \mathbb{N})) \Rightarrow \exists c \in \mathbb{N}.\; ((a \leq c) \wedge (b \leq c)) \end{array}}}
\vnbprop{zz-in-rr}{\vnbstmt{\forall k \in \mathbb{Z}.\; (k \in \mathbb{R})}}

\vnbbookfile{theorem-library/mat-basics}{theorem-library/mat-basics}
\noindent{\small mat-basics.scm -- MAT's read-offs, PROVEN from its definition.}\par
\vnbprop{mat-unfold}{\vnbstmt{\begin{array}{@{}l@{}} \forall m, n, x. \\ \quad (\{p \in \operatorname{matrix}(x) : ((n \in \mathbb{N}) \wedge (((\operatorname{length}(p) = 0) \Rightarrow (m = 0)) \wedge (\neg (\operatorname{length}(p) = 0) \Rightarrow (\operatorname{size}(p) = \langle m, n \rangle))))\} \simeq \operatorname{mat}(m, n, x)) \end{array}}}
\vnbprop{size-unfold}{\vnbstmt{\forall m.\; (\langle \operatorname{length}(m), \left(\text{if }(\operatorname{length}(m) = 0)\text{ then }0\text{ else }\operatorname{length}(\operatorname{nth}(1, m))\right) \rangle \simeq \operatorname{size}(m))}}
\vnbprop{size-nonempty}{\vnbstmt{\begin{array}{@{}l@{}} \forall m. \\ \quad \neg (\operatorname{length}(m) = 0) \Rightarrow (\langle \operatorname{length}(m), \operatorname{length}(\operatorname{nth}(1, m)) \rangle \simeq \operatorname{size}(m)) \end{array}}}
\vnbprop{nth1-pair}{\vnbstmt{\forall a, b.\; (((a \in \mathbf{Set}) \wedge (b \in \mathbf{Set})) \Rightarrow (\operatorname{nth}(1, \langle a, b \rangle) \simeq a))}}
\vnbprop{mat-in-matrix}{\vnbstmt{\forall m, n, x.; \forall q \in \operatorname{mat}(m, n, x).\; (q \in \operatorname{matrix}(x))}}
\vnbprop{mat-cols-in-nn}{\vnbstmt{\forall m, n, x.; \forall q \in \operatorname{mat}(m, n, x).\; (n \in \mathbb{N})}}
\vnbprop{mat-length}{\vnbstmt{\forall m, n, x.; \forall q \in \operatorname{mat}(m, n, x).\; (\operatorname{length}(q) = m)}}
\vnbprop{mat-rows-in-nn}{\vnbstmt{\forall m, n, x.; \forall q \in \operatorname{mat}(m, n, x).\; (m \in \mathbb{N})}}
\vnbprop{mat-size}{\vnbstmt{\forall m, n, x.; \forall q \in \operatorname{mat}(m, n, x).\; ((1 \leq m) \Rightarrow (\operatorname{size}(q) = \langle m, n \rangle))}}
\vnbprop{mat-cols}{\vnbstmt{\forall m, n, x.; \forall q \in \operatorname{mat}(m, n, x).\; ((1 \leq m) \Rightarrow (\operatorname{length}(\operatorname{nth}(1, q)) = n))}}
\vnbprop{mat-size-rows}{\vnbstmt{\forall m, n, x.; \forall q \in \operatorname{mat}(m, n, x).\; (\operatorname{nth}(1, \operatorname{size}(q)) = m)}}
\vnbprop{mat-size-cols}{\vnbstmt{\forall m, n, x.; \forall q \in \operatorname{mat}(m, n, x).\; ((1 \leq m) \Rightarrow (\operatorname{nth}(2, \operatorname{size}(q)) = n))}}
\vnbprop{mat-intro}{\vnbstmt{\begin{array}{@{}l@{}} \forall m, n, x.; \forall q \in \operatorname{matrix}(x). \\ \quad ((n \in \mathbb{N}) \wedge ((1 \leq m) \wedge ((\operatorname{length}(q) = m) \wedge (\operatorname{length}(\operatorname{nth}(1, q)) = n)))) \Rightarrow  \\ \quad \quad q \in \operatorname{mat}(m, n, x) \end{array}}}
\vnbprop{mat-intro-empty}{\vnbstmt{\begin{array}{@{}l@{}} \forall n, x.; \forall q \in \operatorname{matrix}(x). \\ \quad ((n \in \mathbb{N}) \wedge (\operatorname{length}(q) = 0)) \Rightarrow (q \in \operatorname{mat}(0, n, x)) \end{array}}}
\vnbprop{nil-in-mat}{\vnbstmt{\forall n \in \mathbb{N}, x.\; (\langle  \rangle \in \operatorname{mat}(0, n, x))}}
\vnbprop{mat-size-cols-in-nn}{\vnbstmt{\forall m, n, x.; \forall q \in \operatorname{mat}(m, n, x).\; (\operatorname{nth}(2, \operatorname{size}(q)) \in \mathbb{N})}}
\vnbprop{mat-size-cols-pos}{\vnbstmt{\begin{array}{@{}l@{}} \forall m, n, x.; \forall q \in \operatorname{mat}(m, n, x). \\ \quad (1 \leq \operatorname{nth}(2, \operatorname{size}(q))) \Rightarrow (\operatorname{nth}(2, \operatorname{size}(q)) = n) \end{array}}}
\vnbprop{mat-empty-any-cols}{\vnbstmt{\forall c, x, q, k.\; (((q \in \operatorname{mat}(0, c, x)) \wedge (k \in \mathbb{N})) \Rightarrow (q \in \operatorname{mat}(0, k, x)))}}
\vnbprop{interval-1-0-empty}{\vnbstmt{\operatorname{interval}(1, 0) = \emptyset}}

\vnbbookfile{theorem-library/nn-pos-of-nonzero}{theorem-library/nn-pos-of-nonzero}
\noindent{\small nn-pos-of-nonzero.scm -- a nonzero natural is strictly positive.}\par
\vnbprop{nn-pos-of-nonzero}{\vnbstmt{\forall k \in \mathbb{N}.\; (\neg (k = 0) \Rightarrow (0 < k))}}

\vnbbookfile{theorem-library/interval-membership}{theorem-library/interval-membership}
\noindent{\small interval-membership.scm -- INTERVAL's membership IFF and the [1,1] instance, both PROVEN, both by citation of the read-offs already in the tree.}\par
\vnbprop{interval-membership}{\vnbstmt{\forall a, b, i.\; ((i \in \operatorname{interval}(a, b)) \Leftrightarrow ((i \in \mathbb{N}) \wedge ((a \leq i) \wedge (i \leq b))))}}
\vnbprop{one-in-interval-1}{\vnbstmt{1 \in \operatorname{interval}(1, 1)}}

\vnbbookfile{theorem-library/nn-not-lt-le-proof}{theorem-library/nn-not-lt-le-proof}
\noindent{\small nn-not-lt-le-proof.scm -- NOT (k $<$ j) =$>$ j $<$= k on NN, `modulo 0'.}\par
\vnbprop{nn-not-lt-le}{\vnbstmt{\forall k, j \in \mathbb{N}.\; (\neg (k < j) \Rightarrow (j \leq k))}}

\vnbbookfile{theorem-library/rr-order-basics}{theorem-library/rr-order-basics}
\noindent{\small rr-order-basics.scm -- the elementary RR order facts, PROVEN.}\par
\vnbprop{rr-lt-trans}{\vnbstmt{\forall x, y, z \in \mathbb{R}.\; (((x < y) \wedge (y < z)) \Rightarrow (x < z))}}
\vnbprop{rr-lt-le-trans}{\vnbstmt{\forall x, y, z \in \mathbb{R}.\; (((x < y) \wedge (y \leq z)) \Rightarrow (x < z))}}
\vnbprop{rr-le-lt-trans}{\vnbstmt{\forall x, y, z \in \mathbb{R}.\; (((x \leq y) \wedge (y < z)) \Rightarrow (x < z))}}
\vnbprop{rr-lt-implies-le}{\vnbstmt{\forall x, y \in \mathbb{R}.\; ((x < y) \Rightarrow (x \leq y))}}
\vnbprop{rr-le-diff-nonpos}{\vnbstmt{\forall u, v \in \mathbb{R}.\; ((u \leq v) \Rightarrow ((u - v) \leq 0))}}
\vnbprop{rr-lt-diff-pos}{\vnbstmt{\forall u, v \in \mathbb{R}.\; ((u < v) \Rightarrow (0 < (v - u)))}}
\vnbprop{rr-lt-diff-neg}{\vnbstmt{\forall u, v \in \mathbb{R}.\; ((u < v) \Rightarrow ((u - v) < 0))}}
\vnbprop{rr-diff-zero-eq}{\vnbstmt{\forall u, v \in \mathbb{R}.\; ((0 = (u - v)) \Rightarrow (u = v))}}
\vnbprop{rr-le-neg}{\vnbstmt{\forall u, v \in \mathbb{R}.\; ((u \leq v) \Rightarrow (-v \leq -u))}}
\vnbprop{rr-neg-eq-zero}{\vnbstmt{\forall u \in \mathbb{R}.\; ((-u = 0) \Rightarrow (u = 0))}}
\vnbprop{rr-sub-ne-zero}{\vnbstmt{\forall u, v \in \mathbb{R}.\; (\neg (u = v) \Rightarrow \neg ((u - v) = 0))}}
\vnbprop{rr-pos-ne-zero}{\vnbstmt{\forall c \in \mathbb{R}.\; ((0 < c) \Rightarrow \neg (c = 0))}}
\vnbprop{rr-lt-irrefl}{\vnbstmt{\forall a \in \mathbb{R}.\; \neg (a < a)}}
\vnbprop{rr-le-cases}{\vnbstmt{\forall u, v \in \mathbb{R}.\; ((u \leq v) \Rightarrow ((u < v) \vee (u = v)))}}
\vnbprop{rr-le-ne-lt}{\vnbstmt{\forall u, v \in \mathbb{R}.\; (((u \leq v) \wedge \neg (u = v)) \Rightarrow (u < v))}}
\vnbprop{rr-le-add}{\vnbstmt{\begin{array}{@{}l@{}} \forall x, y, u, v. \\ \quad ((x \in \mathbb{R}) \wedge ((y \in \mathbb{R}) \wedge ((u \in \mathbb{R}) \wedge ((v \in \mathbb{R}) \wedge ((x \leq y) \wedge (u \leq v)))))) \Rightarrow  \\ \quad \quad (x + u) \leq (y + v) \end{array}}}
\vnbprop{rr-lt-add}{\vnbstmt{\begin{array}{@{}l@{}} \forall x, y, u, v. \\ \quad ((x \in \mathbb{R}) \wedge ((y \in \mathbb{R}) \wedge ((u \in \mathbb{R}) \wedge ((v \in \mathbb{R}) \wedge ((x < y) \wedge (u \leq v)))))) \Rightarrow  \\ \quad \quad (x + u) < (y + v) \end{array}}}
\vnbprop{rr-add-nonneg}{\vnbstmt{\forall x, y.\; (((x \in \mathbb{R}) \wedge ((y \in \mathbb{R}) \wedge ((0 \leq x) \wedge (0 \leq y)))) \Rightarrow (0 \leq (x + y)))}}
\vnbprop{rr-le-from-diff-nonneg}{\vnbstmt{\forall x, y.\; (((x \in \mathbb{R}) \wedge ((y \in \mathbb{R}) \wedge (0 \leq (y - x)))) \Rightarrow (x \leq y))}}
\vnbprop{rr-double-nonneg}{\vnbstmt{\forall x \in \mathbb{R}.\; ((0 \leq (x + x)) \Rightarrow (0 \leq x))}}
\vnbprop{rr-no-zero-divisors}{\vnbstmt{\forall a, b.\; (((a \in \mathbb{R}) \wedge ((b \in \mathbb{R}) \wedge ((a \cdot b) = 0))) \Rightarrow ((a = 0) \vee (b = 0)))}}
\vnbprop{rr-cancel-mul-right}{\vnbstmt{\begin{array}{@{}l@{}} \forall u, v, c. \\ \quad ((u \in \mathbb{R}) \wedge ((v \in \mathbb{R}) \wedge ((c \in \mathbb{R}) \wedge (\neg (c = 0) \wedge ((u \cdot c) = (v \cdot c)))))) \Rightarrow (u = v) \end{array}}}
\vnbprop{rr-lt-scale-pos}{\vnbstmt{\begin{array}{@{}l@{}} \forall c, x, y. \\ \quad ((c \in \mathbb{R}) \wedge ((x \in \mathbb{R}) \wedge ((y \in \mathbb{R}) \wedge ((0 < c) \wedge (x < y))))) \Rightarrow ((c \cdot x) < (c \cdot y)) \end{array}}}
\vnbprop{rr-le-scale-nonneg}{\vnbstmt{\begin{array}{@{}l@{}} \forall c, x, y. \\ \quad ((c \in \mathbb{R}) \wedge ((x \in \mathbb{R}) \wedge ((y \in \mathbb{R}) \wedge ((0 \leq c) \wedge (x \leq y))))) \Rightarrow ((c \cdot x) \leq (c \cdot y)) \end{array}}}
\vnbprop{rr-sq-nonneg}{\vnbstmt{\forall x \in \mathbb{R}.\; (0 \leq (x \cdot x))}}
\vnbprop{rr-lt-trichotomy}{\vnbstmt{\forall u, v \in \mathbb{R}.\; ((u < v) \vee ((u = v) \vee (v < u)))}}
\vnbprop{rr-upper-of-two}{\vnbstmt{\begin{array}{@{}l@{}} \forall u, v. \\ \quad ((u \in \mathbb{R}) \wedge (v \in \mathbb{R})) \Rightarrow \exists w \in \mathbb{R}.\; ((u \leq w) \wedge (v \leq w)) \end{array}}}
\vnbprop{rr-upper-of-two-below}{\vnbstmt{\begin{array}{@{}l@{}} \forall u, v, m. \\ \quad ((u \in \mathbb{R}) \wedge ((v \in \mathbb{R}) \wedge ((m \in \mathbb{R}) \wedge ((u < m) \wedge (v < m))))) \Rightarrow  \\ \quad \quad \exists w \in \mathbb{R}.\; ((u \leq w) \wedge ((v \leq w) \wedge (w < m))) \end{array}}}
\vnbprop{rr-min-pos}{\vnbstmt{\begin{array}{@{}l@{}} \forall u, v. \\ \quad ((u \in \mathbb{R}) \wedge ((v \in \mathbb{R}) \wedge ((0 < u) \wedge (0 < v)))) \Rightarrow  \\ \quad \quad \exists w \in \mathbb{R}.\; ((0 < w) \wedge ((w \leq u) \wedge (w \leq v))) \end{array}}}
\vnbprop{rr-prod-le-prod}{\vnbstmt{\begin{array}{@{}l@{}} \forall a, b, u, v. \\ \quad ((a \in \mathbb{R}) \wedge ((b \in \mathbb{R}) \wedge ((u \in \mathbb{R}) \wedge ((v \in \mathbb{R}) \wedge (((0 \leq a) \wedge (a \leq b)) \wedge ((0 \leq u) \wedge (u \leq v))))))) \Rightarrow  \\ \quad \quad (a \cdot u) \leq (b \cdot v) \end{array}}}
\vnbprop{rr-nonneg-cancel-pos}{\vnbstmt{\forall c, d.\; (((c \in \mathbb{R}) \wedge ((d \in \mathbb{R}) \wedge ((0 < c) \wedge (0 \leq (c \cdot d))))) \Rightarrow (0 \leq d))}}
\vnbprop{rr-le-total}{\vnbstmt{\forall u, v \in \mathbb{R}.\; ((u \leq v) \vee (v \leq u))}}
\vnbprop{rr-lt-not-le}{\vnbstmt{\forall a, b \in \mathbb{R}.\; ((a < b) \Rightarrow \neg (b \leq a))}}
\vnbprop{rr-not-le-lt}{\vnbstmt{\forall a, b \in \mathbb{R}.\; (\neg (a \leq b) \Rightarrow (b < a))}}

\vnbbookfile{theorem-library/co-order-trans-guarded}{theorem-library/co-order-trans-guarded}
\noindent{\small co-order-trans-guarded.scm -- the six `co-' order composition lemmas, PROVEN.  Four of them GUARDED on RR; two of them exactly as stated.}\par
\vnbprop{co-lt-trans}{\vnbstmt{\forall a, b, c \in \mathbb{R}.\; (((a < b) \wedge (b < c)) \Rightarrow (a < c))}}
\vnbprop{co-le-lt-trans}{\vnbstmt{\forall a, b, c \in \mathbb{R}.\; (((a \leq b) \wedge (b < c)) \Rightarrow (a < c))}}
\vnbprop{co-lt-le-trans}{\vnbstmt{\forall a, b, c \in \mathbb{R}.\; (((a < b) \wedge (b \leq c)) \Rightarrow (a < c))}}
\vnbprop{co-le-eq-trans}{\vnbstmt{\forall a, b, c.\; (((a \leq b) \wedge (b = c)) \Rightarrow (a \leq c))}}
\vnbprop{co-eq-le-trans}{\vnbstmt{\forall a, b, c.\; (((a = b) \wedge (b \leq c)) \Rightarrow (a \leq c))}}

\vnbbookfile{theorem-library/pos-rr-bridges}{theorem-library/pos-rr-bridges}
\noindent{\small pos-rr-bridges.scm -- the two projections of POS-RR, and one more positivity fact, PROVEN modulo 0.}\par
\vnbprop{rr-pos-rr-in-rr}{\vnbstmt{\forall t.\; (\operatorname{pos\text{-}rr}(t) \Rightarrow (t \in \mathbb{R}))}}
\vnbprop{rr-lt-of-pos-rr}{\vnbstmt{\forall t.\; (\operatorname{pos\text{-}rr}(t) \Rightarrow (0 < t))}}
\vnbprop{rr-one-plus-nonneg-pos}{\vnbstmt{\forall d \in \mathbb{R}.\; ((0 \leq d) \Rightarrow (0 < (1 + d)))}}

\vnbbookfile{theorem-library/rr-abs-basics}{theorem-library/rr-abs-basics}
\noindent{\small rr-abs-basics.scm -- absolute value on RR, PROVEN from its DEFINITION.}\par
\vnbprop{rr-abs-of-nonneg}{\vnbstmt{\forall x \in \mathbb{R}.\; ((0 \leq x) \Rightarrow (\lvert x \rvert = x))}}
\vnbprop{rr-abs-of-nonpos}{\vnbstmt{\forall x \in \mathbb{R}.\; ((x \leq 0) \Rightarrow (\lvert x \rvert = -x))}}
\vnbprop{rr-abs-of-neg}{\vnbstmt{\forall x \in \mathbb{R}.\; ((x < 0) \Rightarrow (\lvert x \rvert = -x))}}
\vnbprop{rr-abs-cases}{\vnbstmt{\forall x \in \mathbb{R}.\; (((0 \leq x) \wedge (\lvert x \rvert = x)) \vee (\neg (0 \leq x) \wedge (\lvert x \rvert = -x)))}}
\vnbprop{rr-abs-closed}{\vnbstmt{\forall a \in \mathbb{R}.\; (\lvert a \rvert \in \mathbb{R})}}
\vnbprop{rr-abs-nonneg}{\vnbstmt{\forall a \in \mathbb{R}.\; (0 \leq \lvert a \rvert)}}
\vnbprop{rr-abs-zero}{\vnbstmt{\forall a \in \mathbb{R}.\; ((\lvert a \rvert = 0) \Leftrightarrow (a = 0))}}
\vnbprop{rr-le-abs}{\vnbstmt{\forall x \in \mathbb{R}.\; (x \leq \lvert x \rvert)}}
\vnbprop{rr-neg-abs-le}{\vnbstmt{\forall x \in \mathbb{R}.\; (-\lvert x \rvert \leq x)}}
\vnbprop{rr-abs-bound}{\vnbstmt{\forall x, c \in \mathbb{R}.\; ((\lvert x \rvert \leq c) \Leftrightarrow ((-c \leq x) \wedge (x \leq c)))}}
\vnbprop{rr-abs-triangle}{\vnbstmt{\forall a, b.\; (((a \in \mathbb{R}) \wedge (b \in \mathbb{R})) \Rightarrow (\lvert (a + b) \rvert \leq (\lvert a \rvert + \lvert b \rvert)))}}
\vnbprop{rr-abs-triangle-c}{\vnbstmt{\forall a, b.\; (((a \in \mathbb{R}) \wedge (b \in \mathbb{R})) \Rightarrow (\lvert (a + b) \rvert \leq (\lvert a \rvert + \lvert b \rvert)))}}
\vnbprop{rr-abs-neg}{\vnbstmt{\forall a \in \mathbb{R}.\; (\lvert -a \rvert = \lvert a \rvert)}}
\vnbprop{rr-abs-sub-sym}{\vnbstmt{\forall u, v.\; (((u \in \mathbb{R}) \wedge (v \in \mathbb{R})) \Rightarrow (\lvert (u - v) \rvert = \lvert (v - u) \rvert))}}
\vnbprop{rr-abs-mult}{\vnbstmt{\forall a, b.\; (((a \in \mathbb{R}) \wedge (b \in \mathbb{R})) \Rightarrow (\lvert (a \cdot b) \rvert = (\lvert a \rvert \cdot \lvert b \rvert)))}}
\vnbprop{rr-abs-reverse-triangle}{\vnbstmt{\forall x, y \in \mathbb{R}.\; (\lvert (\lvert x \rvert - \lvert y \rvert) \rvert \leq \lvert (x - y) \rvert)}}
\vnbprop{rr-abs-sum-bound}{\vnbstmt{\begin{array}{@{}l@{}} \forall u, v, p, q, h, e. \\ \quad u \in \mathbb{R} \wedge \\ \quad v \in \mathbb{R} \wedge \\ \quad p \in \mathbb{R} \wedge \\ \quad q \in \mathbb{R} \wedge \\ \quad h \in \mathbb{R} \wedge \\ \quad e \in \mathbb{R} \wedge \\ \quad \lvert (u - p) \rvert \leq h \wedge \\ \quad \lvert (v - q) \rvert \leq h \wedge \\ \quad (h + h) = e \Rightarrow \\ \quad \quad \lvert ((u + v) - (p + q)) \rvert \leq e \end{array}}}
\vnbprop{rr-abs-prod-bound}{\vnbstmt{\begin{array}{@{}l@{}} \forall u, v, p, q, m, k, t, e. \\ \quad u \in \mathbb{R} \wedge \\ \quad v \in \mathbb{R} \wedge \\ \quad p \in \mathbb{R} \wedge \\ \quad q \in \mathbb{R} \wedge \\ \quad m \in \mathbb{R} \wedge \\ \quad k \in \mathbb{R} \wedge \\ \quad t \in \mathbb{R} \wedge \\ \quad e \in \mathbb{R} \wedge \\ \quad \lvert u \rvert \leq m \wedge \\ \quad \lvert q \rvert \leq k \wedge \\ \quad \lvert (u - p) \rvert \leq t \wedge \\ \quad \lvert (v - q) \rvert \leq t \wedge \\ \quad ((m + k) \cdot t) \leq e \Rightarrow \\ \quad \quad \lvert ((u \cdot v) - (p \cdot q)) \rvert \leq e \end{array}}}

\vnbbookfile{theorem-library/rr-max-basics}{theorem-library/rr-max-basics}
\noindent{\small rr-max-basics.scm -- the usable facts about MAX on RR, PROVEN from its defining equation.}\par
\vnbprop{rr-max-closed}{\vnbstmt{\forall x, y \in \mathbb{R}.\; (\operatorname{max}(x, y) \in \mathbb{R})}}
\vnbprop{rr-max-cases}{\vnbstmt{\forall x, y \in \mathbb{R}.\; ((\operatorname{max}(x, y) = x) \vee (\operatorname{max}(x, y) = y))}}
\vnbprop{rr-le-max-left}{\vnbstmt{\forall x, y \in \mathbb{R}.\; (x \leq \operatorname{max}(x, y))}}
\vnbprop{rr-le-max-right}{\vnbstmt{\forall x, y \in \mathbb{R}.\; (y \leq \operatorname{max}(x, y))}}
\vnbprop{rr-max-le}{\vnbstmt{\forall x, y, z \in \mathbb{R}.\; (((x \leq z) \wedge (y \leq z)) \Rightarrow (\operatorname{max}(x, y) \leq z))}}

\vnbbookfile{theorem-library/rr-min-basics}{theorem-library/rr-min-basics}
\noindent{\small rr-min-basics.scm -- the usable facts about MIN on RR, PROVEN from its defining equation, and the NN closure of both MIN and MAX.}\par
\vnbprop{rr-min-closed}{\vnbstmt{\forall x, y \in \mathbb{R}.\; (\operatorname{min}(x, y) \in \mathbb{R})}}
\vnbprop{rr-min-cases}{\vnbstmt{\forall x, y \in \mathbb{R}.\; ((\operatorname{min}(x, y) = x) \vee (\operatorname{min}(x, y) = y))}}
\vnbprop{rr-min-le-left}{\vnbstmt{\forall x, y \in \mathbb{R}.\; (\operatorname{min}(x, y) \leq x)}}
\vnbprop{rr-min-le-right}{\vnbstmt{\forall x, y \in \mathbb{R}.\; (\operatorname{min}(x, y) \leq y)}}
\vnbprop{rr-le-min}{\vnbstmt{\forall x, y, z \in \mathbb{R}.\; (((z \leq x) \wedge (z \leq y)) \Rightarrow (z \leq \operatorname{min}(x, y)))}}
\vnbprop{nn-max-closed}{\vnbstmt{\forall x, y \in \mathbb{N}.\; (\operatorname{max}(x, y) \in \mathbb{N})}}
\vnbprop{nn-min-closed}{\vnbstmt{\forall x, y \in \mathbb{N}.\; (\operatorname{min}(x, y) \in \mathbb{N})}}
\vnbprop{rr-min-one-subadditive}{\vnbstmt{\begin{array}{@{}l@{}} \forall x, y \in \mathbb{R}. \\ \quad ((0 \leq x) \wedge (0 \leq y)) \Rightarrow (\operatorname{min}(1, (x + y)) \leq (\operatorname{min}(1, x) + \operatorname{min}(1, y))) \end{array}}}

\vnbbookfile{theorem-library/cc-real-imag}{theorem-library/cc-real-imag}
\noindent{\small cc-real-imag.scm -- the real and imaginary parts of a complex number.}\par
\vnbprop{cc-re-im-of}{\vnbstmt{\begin{array}{@{}l@{}} \forall z \in \mathbb{C}, a, b \in \mathbb{R}. \\ \quad (z = (a + (b \cdot +i))) \Rightarrow ((\operatorname{real\text{-}part}(z) = a) \wedge (\operatorname{imag\text{-}part}(z) = b)) \end{array}}}
\vnbprop{real-part-in-rr}{\vnbstmt{\forall z \in \mathbb{C}.\; (\operatorname{real\text{-}part}(z) \in \mathbb{R})}}
\vnbprop{imag-part-in-rr}{\vnbstmt{\forall z \in \mathbb{C}.\; (\operatorname{imag\text{-}part}(z) \in \mathbb{R})}}
\vnbprop{cc-re-im-decompose}{\vnbstmt{\forall z \in \mathbb{C}.\; (z = (\operatorname{real\text{-}part}(z) + (+i \cdot \operatorname{imag\text{-}part}(z))))}}
\vnbprop{cc-plus-conj-is-2re}{\vnbstmt{\forall z \in \mathbb{C}.\; ((z + \operatorname{conjugate}(z)) = (2 \cdot \operatorname{real\text{-}part}(z)))}}
\vnbprop{cc-re-sq-le-mod-sq}{\vnbstmt{\forall z \in \mathbb{C}.\; ((\operatorname{real\text{-}part}(z) \cdot \operatorname{real\text{-}part}(z)) \leq (z \cdot \operatorname{conjugate}(z)))}}

\vnbbookfile{theorem-library/rr-recip-order}{theorem-library/rr-recip-order}
\noindent{\small rr-recip-order.scm -- the L1 rung of the "attach the loose ends" program: the reciprocal's SIGN, and the two field facts it stands on.}\par
\vnbprop{rr-mul-zero}{\vnbstmt{\forall a \in \mathbb{R}.\; ((a \cdot 0) = 0)}}
\vnbprop{rr-zero-lt-one}{\vnbstmt{0 < 1}}
\vnbprop{rr-mul-pos}{\vnbstmt{\forall a, b \in \mathbb{R}.\; (((0 < a) \wedge (0 < b)) \Rightarrow (0 < (a \cdot b)))}}
\vnbprop{rr-recip-pos}{\vnbstmt{\forall a \in \mathbb{R}.\; ((0 < a) \Rightarrow (0 < \operatorname{recip}(a)))}}

\vnbbookfile{theorem-library/pos-rr-of-lt}{theorem-library/pos-rr-of-lt}
\noindent{\small pos-rr-of-lt.scm -- 0 $<$ x IMPLIES POS-RR(x), the bridge the tree did not have, and `nn-recip-succ-pos' retired as a four-line consequence of it.}\par
\vnbprop{rr-pos-rr-of-lt}{\vnbstmt{\forall x \in \mathbb{R}.\; ((0 < x) \Rightarrow \operatorname{pos\text{-}rr}(x))}}
\vnbprop{nn-recip-succ-pos}{\vnbstmt{\forall n \in \mathbb{N}.\; \operatorname{pos\text{-}rr}(\operatorname{recip}((n + 1)))}}

\vnbbookfile{theorem-library/rr-halving}{theorem-library/rr-halving}
\noindent{\small rr-halving.scm -- every positive real splits into two equal positive halves, PROVEN.}\par
\vnbprop{rr-pos-halvable}{\vnbstmt{\forall s.\; (\operatorname{pos\text{-}rr}(s) \Rightarrow \exists d.\; (\operatorname{pos\text{-}rr}(d) \wedge ((d + d) = s)))}}

\vnbbookfile{theorem-library/rr-order-bundle}{theorem-library/rr-order-bundle}
\noindent{\small rr-order-bundle.scm -- four elementary real-order facts, PROVEN modulo 0.}\par
\vnbprop{rr-prod-nonpos-pos}{\vnbstmt{\forall u, v \in \mathbb{R}.\; (((0 < v) \wedge ((u \cdot v) \leq 0)) \Rightarrow (u \leq 0))}}
\vnbprop{rr-prod-nonpos-neg}{\vnbstmt{\forall u, v \in \mathbb{R}.\; (((v < 0) \wedge ((u \cdot v) \leq 0)) \Rightarrow (0 \leq u))}}
\vnbprop{rr-midpoint-between}{\vnbstmt{\forall u, v \in \mathbb{R}.\; ((u < v) \Rightarrow \exists w \in \mathbb{R}.\; ((u < w) \wedge (w < v)))}}
\vnbprop{rr-le-scale-nonneg-right}{\vnbstmt{\forall x, y, c \in \mathbb{R}.\; (((0 \leq c) \wedge (x \leq y)) \Rightarrow ((x \cdot c) \leq (y \cdot c)))}}

\vnbbookfile{theorem-library/rr-le-all-pos}{theorem-library/rr-le-all-pos}
\noindent{\small rr-le-all-pos.scm -- a real below every positive real is non-positive, PROVEN.}\par
\vnbprop{rr-le-all-pos-nonpos}{\vnbstmt{\forall x \in \mathbb{R}.\; (\forall s.\; (\operatorname{pos\text{-}rr}(s) \Rightarrow (x \leq s)) \Rightarrow (x \leq 0))}}

\vnbbookfile{theorem-library/subset-lemmas}{theorem-library/subset-lemmas}
\noindent{\small subset-lemmas.scm -- the four inclusion facts, PROVEN.}\par
\vnbprop{subset-mem-fwd}{\vnbstmt{\forall a, b, x.\; (((a \subseteq b) \wedge (x \in a)) \Rightarrow (x \in b))}}
\vnbprop{subset-trans}{\vnbstmt{\forall a, b, c.\; (((a \subseteq b) \wedge (b \subseteq c)) \Rightarrow (a \subseteq c))}}
\vnbprop{subclass-of-set-is-set}{\vnbstmt{\forall a.; \forall b \in \mathbf{Set}.\; ((a \subseteq b) \Rightarrow (a \in \mathbf{Set}))}}
\vnbprop{subset-refl}{\vnbstmt{\forall a.\; (a \subseteq a)}}

\vnbbookfile{theorem-library/fun-image-set}{theorem-library/fun-image-set}
\noindent{\small fun-image-set.scm -- the image of ANY class under a function into a SET is a set, proven with no citation of `image-set' (replacement).}\par
\vnbprop{fun-image-set}{\vnbstmt{\forall a, b, f, s.\; (((f \in (a \to b)) \wedge (b \in \mathbf{Set})) \Rightarrow (\operatorname{image}(f, s) \in \mathbf{Set}))}}

\vnbbookfile{theorem-library/rake-analysis2}{theorem-library/rake-analysis2}
\noindent{\small A file shared across books; the entries that belong to this one.}\par
\noindent{\small rake-analysis2.scm -- rake batch L (2026-09-17): five asserted leaves of the analysis / sets area, plus the general lemma that unblocked one of them.}\par
\vnbprop{fun-domain-in-set}{\vnbstmt{\forall a, b.; \forall f \in (a \to b).\; (a \in \mathbf{Set})}}
\vnbprop{bijection-compose}{\vnbstmt{\begin{array}{@{}l@{}} \forall x, y, z, i, a. \\ \quad ((i \in \mathrm{Bij}(x, y)) \wedge (a \in \mathrm{Bij}(y, z))) \Rightarrow  \\ \quad \quad (\lambda b \in x.\; a(i(b))) \in \mathrm{Bij}(x, z) \end{array}}}
\vnbprop{le-bound-mono}{\vnbstmt{\begin{array}{@{}l@{}} \forall p, a, b, n. \\ \quad ((p \in \mathbb{R}) \wedge ((a \in \mathbb{R}) \wedge ((b \in \mathbb{R}) \wedge ((n \in \mathbb{R}) \wedge ((0 \leq n) \wedge ((p \leq (a \cdot n)) \wedge (a \leq b))))))) \Rightarrow  \\ \quad \quad p \leq (b \cdot n) \end{array}}}

\vnbbookfile{theorem-library/rake-inverse-bij}{theorem-library/rake-inverse-bij}
\noindent{\small theorem-library/rake-inverse-bij.scm -- rake batch N (2026-09-17), part 1 of 2. --------------------------------------------------------------------------- NOTE, 2026-09-20 (batch 9-B).  This file was written while cardinality was AXIOMATISED under the name CARD and the defined constant was its companion}\par
\vnbprop{inverse-bij-right}{\vnbstmt{\forall m, d.; \forall i \in \mathrm{Bij}(m, d), y \in d.\; (i(i^{-1}(y)) = y)}}
\vnbprop{inverse-bij-left}{\vnbstmt{\forall m, d.; \forall i \in \mathrm{Bij}(m, d), x \in m.\; (i^{-1}(i(x)) = x)}}
\vnbprop{inverse-bij-in-fun}{\vnbstmt{\forall x, y.; \forall i \in \mathrm{Bij}(x, y).\; (i^{-1} \in (y \to x))}}
\vnbprop{inverse-bij-is-bijection}{\vnbstmt{\forall x, y.; \forall i \in \mathrm{Bij}(x, y).\; (i^{-1} \in \mathrm{Bij}(y, x))}}
\vnbprop{ord-segment-self}{\vnbstmt{\forall n \in \mathbb{N}.\; \neg (n \in \mathrm{OS}(n))}}

\vnbbookfile{theorem-library/nn-parity-proof}{theorem-library/nn-parity-proof}
\noindent{\small nn-parity-proof.scm -- Peano arithmetic on NN: from the two recursion equations (structure-library/nn-arith.scm) to parity and even-square.}\par
\vnbprop{nn-one-is-succ-zero}{\vnbstmt{1 = (0 + 1)}}
\vnbprop{nn-succ-inj}{\vnbstmt{\forall a, b \in \mathbb{N}.\; (((a + 1) = (b + 1)) \Rightarrow (a = b))}}
\vnbprop{nn-succ-nonzero}{\vnbstmt{\forall n \in \mathbb{N}.\; \neg ((n + 1) = 0)}}
\vnbprop{nn-succ-plus-one}{\vnbstmt{\forall n \in \mathbb{N}.\; ((n + 1) = (n + 1))}}
\vnbprop{nn-mul-succ}{\vnbstmt{\forall a, b \in \mathbb{N}.\; ((a \cdot (b + 1)) = ((a \cdot b) + a))}}
\vnbprop{nn-two-is-succ-succ-zero}{\vnbstmt{2 = ((0 + 1) + 1)}}
\vnbprop{nn-plus-two}{\vnbstmt{\forall x \in \mathbb{N}.\; ((x + 2) = ((x + 1) + 1))}}
\vnbprop{nn-two-mul-succ}{\vnbstmt{\forall k \in \mathbb{N}.\; ((2 \cdot (k + 1)) = (((2 \cdot k) + 1) + 1))}}
\vnbprop{nn-nonzero-is-succ}{\vnbstmt{\forall n \in \mathbb{N}.\; (\neg (n = 0) \Rightarrow \exists q \in \mathbb{N}.\; (n = (q + 1)))}}
\vnbprop{nn-parity}{\vnbstmt{\forall n \in \mathbb{N}.\; \exists k \in \mathbb{N}.\; ((n = (2 \cdot k)) \vee (n = ((2 \cdot k) + 1)))}}
\vnbprop{nn-zero-or-succ}{\vnbstmt{\forall y \in \mathbb{N}.\; ((y = 0) \vee \exists q \in \mathbb{N}.\; (y = (q + 1)))}}
\vnbprop{nn-parity-exclusive}{\vnbstmt{\forall x, y \in \mathbb{N}.\; \neg ((2 \cdot x) = ((2 \cdot y) + 1))}}
\vnbprop{nn-even-square}{\vnbstmt{\begin{array}{@{}l@{}} \forall p \in \mathbb{N}. \\ \quad \exists m \in \mathbb{N}.\; ((p \cdot p) = (2 \cdot m)) \Rightarrow \exists k \in \mathbb{N}.\; (p = (2 \cdot k)) \end{array}}}
\vnbprop{bt-succ-minus-1}{\vnbstmt{\forall n \in \mathbb{N}.\; (((n + 1) - 1) = n)}}

\vnbbookfile{theorem-library/nn-pred}{theorem-library/nn-pred}
\noindent{\small nn-pred.scm -- the predecessor on NN, DEFINED by definite description.}\par
\par\medskip\noindent\textbf{Definition} (\texttt{pred}).\;$\operatorname{pred}(\operatorname{n\_}) :=$ $(\iota\, \operatorname{m\_}.\;((\operatorname{m\_} \in \mathbb{N}) \wedge ((\operatorname{m\_} + 1) = \operatorname{n\_})))$
\vnbprop{pred-property}{\vnbstmt{\forall n \in \mathbb{N}.\; (\neg (n = 0) \Rightarrow ((\operatorname{pred}(n) \in \mathbb{N}) \wedge ((\operatorname{pred}(n) + 1) = n)))}}
\vnbprop{pred-in-nn}{\vnbstmt{\forall n \in \mathbb{N}.\; (\neg (n = 0) \Rightarrow (\operatorname{pred}(n) \in \mathbb{N}))}}
\vnbprop{pred-succ}{\vnbstmt{\forall n \in \mathbb{N}.\; (\neg (n = 0) \Rightarrow ((\operatorname{pred}(n) + 1) = n))}}
\vnbprop{pred-of-succ}{\vnbstmt{\forall m \in \mathbb{N}.\; (\operatorname{pred}((m + 1)) = m)}}

\vnbbookfile{theorem-library/ord-segment-arith}{theorem-library/ord-segment-arith}
\noindent{\small ord-segment-arith.scm -- the two bridges between segment membership and arithmetic on NN.}\par
\vnbprop{seg-mem-lt}{\vnbstmt{\forall n, j \in \mathbb{N}.\; ((j \in \mathrm{OS}(n)) \Leftrightarrow (j < n))}}
\vnbprop{seg-mem-succ-le}{\vnbstmt{\forall n, j \in \mathbb{N}.\; ((j \in \mathrm{OS}((n + 1))) \Leftrightarrow (j \leq n))}}

\vnbbookfile{theorem-library/finite-surgery}{theorem-library/finite-surgery}
\noindent{\small finite-surgery.scm -- the finite-surgery kit, first member: COLLAPSE-AT.}\par
\par\medskip\noindent\textbf{Definition} (\texttt{collapse-at}).\;$\operatorname{collapse\text{-}at}(\operatorname{k\_}) :=$ $(\lambda \operatorname{z\_} \in \mathbb{N}.\; \left(\text{if }(\operatorname{k\_} < \operatorname{z\_})\text{ then }\operatorname{pred}(\operatorname{z\_})\text{ else }\operatorname{z\_}\right))$
\vnbprop{nn-lt-succ-le}{\vnbstmt{\forall k, j \in \mathbb{N}.\; ((k < j) \Rightarrow ((k + 1) \leq j))}}
\vnbprop{collapse-at-lo}{\vnbstmt{\forall k, j \in \mathbb{N}.\; (\neg (k < j) \Rightarrow (\operatorname{collapse\text{-}at}(k)(j) = j))}}
\vnbprop{collapse-at-hi}{\vnbstmt{\forall k, j \in \mathbb{N}.\; ((k < j) \Rightarrow (\operatorname{collapse\text{-}at}(k)(j) = \operatorname{pred}(j)))}}
\vnbprop{collapse-at-in-nn}{\vnbstmt{\forall k, j \in \mathbb{N}.\; (\operatorname{collapse\text{-}at}(k)(j) \in \mathbb{N})}}
\vnbprop{collapse-at-lt}{\vnbstmt{\begin{array}{@{}l@{}} \forall n, k, j. \\ \quad ((n \in \mathbb{N}) \wedge ((k \in \mathbb{N}) \wedge ((j \in \mathbb{N}) \wedge ((k \leq n) \wedge ((j \leq n) \wedge \neg (j = k)))))) \Rightarrow  \\ \quad \quad \operatorname{collapse\text{-}at}(k)(j) < n \end{array}}}
\vnbprop{collapse-at-inj}{\vnbstmt{\begin{array}{@{}l@{}} \forall k, a, b. \\ \quad k \in \mathbb{N} \wedge \\ \quad a \in \mathbb{N} \wedge \\ \quad b \in \mathbb{N} \wedge \\ \quad \neg (a = k) \wedge \\ \quad \neg (b = k) \wedge \\ \quad \operatorname{collapse\text{-}at}(k)(a) = \operatorname{collapse\text{-}at}(k)(b) \Rightarrow \\ \quad \quad a = b \end{array}}}
\vnbprop{collapse-at-in-seg}{\vnbstmt{\begin{array}{@{}l@{}} \forall n, k, j. \\ \quad n \in \mathbb{N} \wedge \\ \quad k \in \mathbb{N} \wedge \\ \quad j \in \mathbb{N} \wedge \\ \quad k \in \mathrm{OS}((n + 1)) \wedge \\ \quad j \in \mathrm{OS}((n + 1)) \wedge \\ \quad \neg (j = k) \Rightarrow \\ \quad \quad \operatorname{collapse\text{-}at}(k)(j) \in \mathrm{OS}(n) \end{array}}}
\vnbprop{nn-le-ne-succ}{\vnbstmt{\forall a, k \in \mathbb{N}.\; (((a \leq k) \wedge \neg (k = a)) \Rightarrow ((a + 1) \leq k))}}

\vnbbookfile{theorem-library/pigeonhole-segments}{theorem-library/pigeonhole-segments}
\noindent{\small pigeonhole-segments.scm -- Track A of the CARD plan: finite pigeonhole for ordinal segments, the one piece of real mathematics the plan needs.}\par
\vnbprop{fun-into-seg-zero-absurd}{\vnbstmt{\forall a, f, x.\; (((f \in (a \to \mathrm{OS}(0))) \wedge (x \in a)) \Rightarrow \mathit{falsity})}}
\vnbprop{no-injection-seg-1-into-seg-0}{\vnbstmt{\forall f.\; \neg (f \in \operatorname{injection}(\mathrm{OS}(1), \mathrm{OS}(0)))}}
\vnbprop{pigeonhole-step}{\vnbstmt{\begin{array}{@{}l@{}} \forall n \in \mathbb{N}, g. \\ \quad \forall f.\; \neg (f \in \operatorname{injection}(\mathrm{OS}((n + 1)), \mathrm{OS}(n))) \Rightarrow  \\ \quad \quad \neg (g \in \operatorname{injection}(\mathrm{OS}(((n + 1) + 1)), \mathrm{OS}((n + 1)))) \end{array}}}
\vnbprop{pigeonhole-segments}{\vnbstmt{\forall n \in \mathbb{N}, f.\; \neg (f \in \operatorname{injection}(\mathrm{OS}((n + 1)), \mathrm{OS}(n)))}}
\vnbprop{pigeonhole-segments-gen}{\vnbstmt{\begin{array}{@{}l@{}} \forall m, n, f. \\ \quad ((m \in \mathbb{N}) \wedge ((n \in \mathbb{N}) \wedge (n < m))) \Rightarrow  \\ \quad \quad \neg (f \in \operatorname{injection}(\mathrm{OS}(m), \mathrm{OS}(n))) \end{array}}}

\vnbbookfile{theorem-library/enum-append}{theorem-library/enum-append}
\noindent{\small theorem-library/enum-append.scm -- enum-append-is-bijection, moved out of theorem-library/finsum-insert.scm on 2026-09-20 (batch 9-B, CARD := CARD-STAR).}\par
\vnbprop{enum-append-is-bijection}{\vnbstmt{\begin{array}{@{}l@{}} \forall n \in \mathbb{N}, g, t \in \mathbf{Set}, m \in \mathrm{Bij}(\mathrm{OS}(n), g). \\ \quad \neg (t \in g) \Rightarrow  \\ \quad \quad (\lambda i \in \mathrm{OS}((n + 1)).\; \left(\text{if }(i \in \mathrm{OS}(n))\text{ then }m(i)\text{ else }t\right)) \\ \quad \quad \in\; \mathrm{Bij}(\mathrm{OS}((n + 1)), (g \cup \{t\})) \end{array}}}

\vnbbookfile{theorem-library/union-laws}{theorem-library/union-laws}
\noindent{\small theorem-library/union-laws.scm -- union-empty-right and union-assoc, moved out of theorem-library/fin-subsets.scm on 2026-09-20 (batch 9-B, CARD := CARD-STAR).}\par
\vnbprop{union-empty-right}{\vnbstmt{\forall a.\; ((a \cup \emptyset) = a)}}
\vnbprop{union-assoc}{\vnbstmt{\forall a, b, c.\; (((a \cup b) \cup c) = (a \cup (b \cup c)))}}

\vnbbookfile{structure-library/bijection-derived}{structure-library/bijection-derived}
\noindent{\small bijection-derived.scm -- the three BIJECTION projection lemmas, PROVEN.}\par
\vnbprop{bijection-in-fun}{\vnbstmt{\forall x, y.; \forall i \in \mathrm{Bij}(x, y).\; (i \in (x \to y))}}
\vnbprop{bijection-injective}{\vnbstmt{\forall x, y.; \forall i \in \mathrm{Bij}(x, y), a, b \in x.\; ((i(a) = i(b)) \Rightarrow (a = b))}}
\vnbprop{bijection-surjective}{\vnbstmt{\forall x, y.; \forall i \in \mathrm{Bij}(x, y), w \in y.\; \exists z \in x.\; (i(z) = w)}}

\vnbbookfile{theorem-library/bijection-identity-proof}{theorem-library/bijection-identity-proof}
\noindent{\small bijection-identity-proof.scm -- the identity lambda on a SET is a bijection.}\par
\vnbprop{bijection-identity}{\vnbstmt{\forall x \in \mathbf{Set}.\; ((\lambda a \in x.\; a) \in \mathrm{Bij}(x, x))}}

\vnbbookfile{theorem-library/card-defined}{theorem-library/card-defined}
\noindent{\small card-defined.scm -- CARD, cardinality DEFINED rather than axiomatised.}\par
\vnbprop{cd-body-unique}{\vnbstmt{\begin{array}{@{}l@{}} \forall a, l, e. \\ \quad l \in \mathrm{Ord} \wedge \\ \quad \exists i.\; (i \in \mathrm{Bij}(a, \mathrm{OS}(l))) \wedge \\ \quad \forall b. \\ \quad \quad \operatorname{ord\text{-}lt}(b, l) \Rightarrow \neg \exists c.\; (c \in \mathrm{Bij}(a, \mathrm{OS}(b))) \wedge \\ \quad e \in \mathrm{Ord} \wedge \\ \quad \exists i.\; (i \in \mathrm{Bij}(a, \mathrm{OS}(e))) \wedge \\ \quad \forall b. \\ \quad \quad \operatorname{ord\text{-}lt}(b, e) \Rightarrow \neg \exists c.\; (c \in \mathrm{Bij}(a, \mathrm{OS}(b))) \Rightarrow \\ \quad \quad l = e \end{array}}}
\vnbprop{cd-seg-body}{\vnbstmt{\begin{array}{@{}l@{}} \forall n \in \mathbb{N}. \\ \quad n \in \mathrm{Ord} \wedge \\ \quad \exists i.\; (i \in \mathrm{Bij}(\mathrm{OS}(n), \mathrm{OS}(n))) \wedge \\ \quad \forall a. \\ \quad \quad \operatorname{ord\text{-}lt}(a, n) \Rightarrow  \\ \quad \quad \quad \neg \exists b.\; (b \in \mathrm{Bij}(\mathrm{OS}(n), \mathrm{OS}(a))) \end{array}}}
\vnbprop{card-segment}{\vnbstmt{\forall n \in \mathbb{N}.\; (|\mathrm{OS}(n)| = n)}}

\vnbbookfile{theorem-library/card-finite}{theorem-library/card-finite}
\noindent{\small card-finite.scm -- the finite layer of the DEFINED cardinal, CARD. --------------------------------------------------------------------------- NOTE, 2026-09-20 (batch 9-B).  This file was written while cardinality was AXIOMATISED under the name CARD and the defined constant was its companion}\par
\vnbprop{card-from-body}{\vnbstmt{\begin{array}{@{}l@{}} \forall a.; \forall l \in \mathrm{Ord}. \\ \quad \exists i.\; (i \in \mathrm{Bij}(a, \mathrm{OS}(l))) \wedge \\ \quad \forall b. \\ \quad \quad \operatorname{ord\text{-}lt}(b, l) \Rightarrow \neg \exists c.\; (c \in \mathrm{Bij}(a, \mathrm{OS}(b))) \Rightarrow \\ \quad \quad |a| = l \end{array}}}
\vnbprop{card-bij}{\vnbstmt{\begin{array}{@{}l@{}} \forall n \in \mathbb{N}, a, h, t. \\ \quad ((h \in \mathrm{Bij}(a, \mathrm{OS}(n))) \wedge (t \in \mathrm{Bij}(\mathrm{OS}(n), a))) \Rightarrow  \\ \quad \quad |a| = n \end{array}}}
\vnbprop{seg0-empty}{\vnbstmt{\emptyset = \mathrm{OS}(0)}}
\vnbprop{card-empty}{\vnbstmt{|\emptyset| = 0}}

\vnbbookfile{theorem-library/ord-no-injection}{theorem-library/ord-no-injection}
\noindent{\small ord-no-injection.scm -- NO SET RECEIVES AN INJECTIVE CLASS FUNCTION FROM ORD.}\par
\vnbprop{ord-no-injection-into-set}{\vnbstmt{\begin{array}{@{}l@{}} \forall d, i. \\ \quad d \in \mathbf{Set} \wedge \\ \quad \forall a \in \mathrm{Ord}.\; (i(a) \in d) \wedge \\ \quad \forall a_{1}, a_{2} \in \mathrm{Ord}.\; ((i(a_{1}) = i(a_{2})) \Rightarrow (a_{1} = a_{2})) \Rightarrow \\ \quad \quad \mathit{falsity} \end{array}}}

\vnbbookfile{theorem-library/zen-step}{theorem-library/zen-step}
\noindent{\small zen-step.scm -- ZEN's three recursion equations, collapsed into ONE.}\par
\par\medskip\noindent\textbf{Definition} (\texttt{zen}).\;For every $x$,  $\operatorname{zen}(x, 0) = \varepsilon(x)$\par
Let $x$ be arbitrary. Let $\alpha \in \mathrm{Ord}$. Then
\begin{equation*}
\begin{array}{@{}l@{}} \operatorname{zen}(x, (\alpha + 1)) \\ =\; \varepsilon(\{\operatorname{y\_} \in x : \forall \operatorname{c\_}.\; (\operatorname{ord\text{-}le}(\operatorname{c\_}, \alpha) \Rightarrow \neg (\operatorname{zen}(x, \operatorname{c\_}) = \operatorname{y\_}))\}) \end{array}
\end{equation*}
\par
Let $x$ be arbitrary. Let $\mathit{lam} \text{ be s.t. } \operatorname{limit\text{-}ord}(\mathit{lam})$. Then
\begin{equation*}
\operatorname{zen}(x, \mathit{lam}) = \varepsilon(\{\operatorname{y\_} \in x : \forall \operatorname{c\_}.\; (\operatorname{ord\text{-}lt}(\operatorname{c\_}, \mathit{lam}) \Rightarrow \neg (\operatorname{zen}(x, \operatorname{c\_}) = \operatorname{y\_}))\})
\end{equation*}

\vnbprop{ord-lt-succ-iff-le}{\vnbstmt{\forall l, z \in \mathrm{Ord}.\; (\operatorname{ord\text{-}lt}(z, (l + 1)) \Leftrightarrow \operatorname{ord\text{-}le}(z, l))}}
\vnbprop{zen-step}{\vnbstmt{\begin{array}{@{}l@{}} \forall a \in \mathrm{Ord}, x. \\ \quad \operatorname{zen}(x, a) = \varepsilon(\{y \in x : \forall c.\; (\operatorname{ord\text{-}lt}(c, a) \Rightarrow \neg (\operatorname{zen}(x, c) = y))\}) \end{array}}}
\vnbprop{zen-hits}{\vnbstmt{\begin{array}{@{}l@{}} \forall a \in \mathrm{Ord}, s. \\ \quad \exists x.\; (x \in \{y \in s : \forall c.\; (\operatorname{ord\text{-}lt}(c, a) \Rightarrow \neg (\operatorname{zen}(s, c) = y))\}) \Rightarrow  \\ \quad \quad \operatorname{zen}(s, a) \in \{y \in s : \forall c.\; (\operatorname{ord\text{-}lt}(c, a) \Rightarrow \neg (\operatorname{zen}(s, c) = y))\} \end{array}}}
\vnbprop{zen-exhausts}{\vnbstmt{\begin{array}{@{}l@{}} \forall s \in \mathbf{Set}. \exists a \in \mathrm{Ord}. \\ \quad \neg \exists x.\; (x \in \{y \in s : \forall c.\; (\operatorname{ord\text{-}lt}(c, a) \Rightarrow \neg (\operatorname{zen}(s, c) = y))\}) \end{array}}}

\vnbbookfile{theorem-library/rake-zermelo}{theorem-library/rake-zermelo}
\noindent{\small rake-zermelo.scm -- CARD-free Zermelo: every set bijects with an ordinal --------------------------------------------------------------------------- NOTE, 2026-09-20 (batch 9-B).  This file was written while cardinality was AXIOMATISED under the name CARD and the defined constant was its companion}\par
\vnbprop{zermelo-least-ordinal}{\vnbstmt{\begin{array}{@{}l@{}} \forall s \in \mathbf{Set}. \exists l \in \mathrm{Ord}. \\ \quad \neg \exists x.\; (x \in \{y \in s : \forall c.\; (\operatorname{ord\text{-}lt}(c, l) \Rightarrow \neg (\operatorname{zen}(s, c) = y))\}) \wedge \\ \quad \forall z. \\ \quad \quad \operatorname{ord\text{-}lt}(z, l) \Rightarrow  \\ \quad \quad \quad \exists x. \\ \quad \quad \quad \quad x \in \{y \in s : \forall c.\; (\operatorname{ord\text{-}lt}(c, z) \Rightarrow \neg (\operatorname{zen}(s, c) = y))\} \end{array}}}
\vnbprop{zen-enum-in-fun}{\vnbstmt{\begin{array}{@{}l@{}} \forall l \in \mathrm{Ord}, s \in \mathbf{Set}. \\ \quad \forall z. \\ \quad \quad \operatorname{ord\text{-}lt}(z, l) \Rightarrow  \\ \quad \quad \quad \exists x. \\ \quad \quad \quad \quad x \in \{y \in s : \forall c.\; (\operatorname{ord\text{-}lt}(c, z) \Rightarrow \neg (\operatorname{zen}(s, c) = y))\} \Rightarrow \\ \quad \quad (\lambda a \in \mathrm{OS}(l).\; \operatorname{zen}(s, a)) \in (\mathrm{OS}(l) \to s) \end{array}}}
\vnbprop{zen-enum-injective}{\vnbstmt{\begin{array}{@{}l@{}} \forall l \in \mathrm{Ord}, s \in \mathbf{Set}, a_{1}, a_{2} \in \mathrm{OS}(l). \\ \quad \forall z. \\ \quad \quad \operatorname{ord\text{-}lt}(z, l) \Rightarrow  \\ \quad \quad \quad \exists x. \\ \quad \quad \quad \quad x \in \{y \in s : \forall c.\; (\operatorname{ord\text{-}lt}(c, z) \Rightarrow \neg (\operatorname{zen}(s, c) = y))\} \wedge \\ \quad (\lambda a \in \mathrm{OS}(l).\; \operatorname{zen}(s, a))(a_{1}) = (\lambda a \in \mathrm{OS}(l).\; \operatorname{zen}(s, a))(a_{2}) \Rightarrow \\ \quad \quad a_{1} = a_{2} \end{array}}}
\vnbprop{zen-enum-surjective}{\vnbstmt{\begin{array}{@{}l@{}} \forall l \in \mathrm{Ord}, s \in \mathbf{Set}, w \in s. \\ \quad \neg \exists x.\; (x \in \{y \in s : \forall c.\; (\operatorname{ord\text{-}lt}(c, l) \Rightarrow \neg (\operatorname{zen}(s, c) = y))\}) \Rightarrow  \\ \quad \quad \exists z \in \mathrm{OS}(l). \\ \quad \quad \quad (\lambda a \in \mathrm{OS}(l).\; \operatorname{zen}(s, a))(z) = w \end{array}}}
\vnbprop{zermelo-bijection}{\vnbstmt{\forall s \in \mathbf{Set}.\; \exists l \in \mathrm{Ord}, i.\; (i \in \mathrm{Bij}(\mathrm{OS}(l), s))}}

\vnbbookfile{theorem-library/rake-ord-pigeonhole}{theorem-library/rake-ord-pigeonhole}
\noindent{\small rake-ord-pigeonhole.scm -- rake batch 5c, assignment 5c-U (2026-09-18). --------------------------------------------------------------------------- NOTE, 2026-09-20 (batch 9-B).  This file was written while cardinality was AXIOMATISED under the name CARD and the defined constant was its companion}\par
\vnbprop{ord-segment-pigeonhole-guarded}{\vnbstmt{\begin{array}{@{}l@{}} \forall l, e \in \mathbb{N}, f. \\ \quad \operatorname{ord\text{-}lt}(e, l) \Rightarrow \neg (f \in \operatorname{injection}(\mathrm{OS}(l), \mathrm{OS}(e))) \end{array}}}
\vnbprop{card-body-exists}{\vnbstmt{\begin{array}{@{}l@{}} \forall s \in \mathbf{Set}. \exists l \in \mathrm{Ord}. \\ \quad \exists i.\; (i \in \mathrm{Bij}(s, \mathrm{OS}(l))) \wedge \\ \quad \forall a. \\ \quad \quad \operatorname{ord\text{-}lt}(a, l) \Rightarrow \neg \exists b.\; (b \in \mathrm{Bij}(s, \mathrm{OS}(a))) \end{array}}}
\vnbprop{card-zermelo}{\vnbstmt{\begin{array}{@{}l@{}} \forall s \in \mathbf{Set}. \\ \quad (|s| \in \mathrm{Ord}) \wedge \exists i.\; (i \in \mathrm{Bij}(s, \mathrm{OS}(|s|))) \end{array}}}
\vnbprop{well-ordering-principle}{\vnbstmt{\forall s \in \mathbf{Set}.\; \exists i.\; (i \in \mathrm{Bij}(\mathrm{OS}(|s|), s))}}

\vnbbookfile{theorem-library/rake-card-star-laws}{theorem-library/rake-card-star-laws}
\noindent{\small rake-card-star-laws.scm -- rake batch 5c, assignment 5c-X (2026-09-18). --------------------------------------------------------------------------- NOTE, 2026-09-20 (batch 9-B).  This file was written while cardinality was AXIOMATISED under the name CARD and the defined constant was its companion}\par
\vnbprop{card-insert-curried}{\vnbstmt{\begin{array}{@{}l@{}} \forall a, x \in \mathbf{Set}. \\ \quad ((|a| \in \mathbb{N}) \wedge \neg (x \in a)) \Rightarrow  \\ \quad \quad |(a \cup \{x\})| = (|a| + 1) \end{array}}}
\vnbprop{seg-restrict-is-bijection}{\vnbstmt{\begin{array}{@{}l@{}} \forall n \in \mathbb{N}, b.; \forall g \in \mathrm{Bij}(\mathrm{OS}((n + 1)), b). \\ \quad (\lambda i \in \mathrm{OS}(n).\; g(i)) \in \mathrm{Bij}(\mathrm{OS}(n), \operatorname{image}(g, \mathrm{OS}(n))) \end{array}}}
\vnbprop{card-seg-image}{\vnbstmt{\begin{array}{@{}l@{}} \forall n \in \mathbb{N}, b.; \forall g \in \mathrm{Bij}(\mathrm{OS}((n + 1)), b). \\ \quad |\operatorname{image}(g, \mathrm{OS}(n))| = n \end{array}}}
\vnbprop{seg-image-subset}{\vnbstmt{\begin{array}{@{}l@{}} \forall n \in \mathbb{N}, b.; \forall g \in \mathrm{Bij}(\mathrm{OS}((n + 1)), b), w \in \operatorname{image}(g, \mathrm{OS}(n)). \\ \quad w \in b \end{array}}}
\vnbprop{seg-image-omits-top}{\vnbstmt{\begin{array}{@{}l@{}} \forall n \in \mathbb{N}, b.; \forall g \in \mathrm{Bij}(\mathrm{OS}((n + 1)), b). \\ \quad \neg (g(n) \in \operatorname{image}(g, \mathrm{OS}(n))) \end{array}}}
\vnbprop{seg-image-insert}{\vnbstmt{\begin{array}{@{}l@{}} \forall n \in \mathbb{N}, b.; \forall g \in \mathrm{Bij}(\mathrm{OS}((n + 1)), b). \\ \quad b = (\operatorname{image}(g, \mathrm{OS}(n)) \cup \{g(n)\}) \end{array}}}
\vnbprop{card-zero-is-empty}{\vnbstmt{\forall b \in \mathbf{Set}.\; ((|b| = 0) \Rightarrow (b = \emptyset))}}
\vnbprop{intersection-empty-transfer}{\vnbstmt{\begin{array}{@{}l@{}} \forall a \in \mathbf{Set}, b, c. \\ \quad (((a \cap b) = \emptyset) \wedge \forall w \in c.\; (w \in b)) \Rightarrow  \\ \quad \quad (a \cap c) = \emptyset \end{array}}}
\vnbprop{card-union-disjoint-ind}{\vnbstmt{\begin{array}{@{}l@{}} \forall n \in \mathbb{N}, a, b \in \mathbf{Set}. \\ \quad ((|a| \in \mathbb{N}) \wedge ((|b| = n) \wedge ((a \cap b) = \emptyset))) \Rightarrow  \\ \quad \quad |(a \cup b)| = (|a| + n) \end{array}}}
\vnbprop{card-union-disjoint-curried}{\vnbstmt{\begin{array}{@{}l@{}} \forall a, b \in \mathbf{Set}. \\ \quad ((|a| \in \mathbb{N}) \wedge ((|b| \in \mathbb{N}) \wedge ((a \cap b) = \emptyset))) \Rightarrow  \\ \quad \quad |(a \cup b)| = (|a| + |b|) \end{array}}}
\vnbprop{card-finite-induction-aux}{\vnbstmt{\begin{array}{@{}l@{}} \forall n \in \mathbb{N}, c.; \forall s \in \mathbf{Set}. \\ \quad \emptyset \in c \wedge \\ \quad \forall s, x \in \mathbf{Set}. \\ \quad \quad (((|s| \in \mathbb{N}) \wedge (s \in c)) \wedge \neg (x \in s)) \Rightarrow ((s \cup \{x\}) \in c) \wedge \\ \quad |s| = n \Rightarrow \\ \quad \quad s \in c \end{array}}}
\vnbprop{finite-set-induction}{\vnbstmt{\begin{array}{@{}l@{}} \forall c.; \forall s \in \mathbf{Set}. \\ \quad \emptyset \in c \wedge \\ \quad \forall s, x \in \mathbf{Set}. \\ \quad \quad (((|s| \in \mathbb{N}) \wedge (s \in c)) \wedge \neg (x \in s)) \Rightarrow ((s \cup \{x\}) \in c) \wedge \\ \quad |s| \in \mathbb{N} \Rightarrow \\ \quad \quad s \in c \end{array}}}
\vnbprop{injection-image-is-bijection}{\vnbstmt{\begin{array}{@{}l@{}} \forall m \in \mathbf{Set}, d.; \forall h \in \operatorname{injection}(m, d). \\ \quad (\lambda z \in m.\; h(z)) \in \mathrm{Bij}(m, \operatorname{image}(h, m)) \end{array}}}
\vnbprop{card-image-injection-curried}{\vnbstmt{\begin{array}{@{}l@{}} \forall m \in \mathbf{Set}, d.; \forall h \in \operatorname{injection}(m, d). \\ \quad (|m| \in \mathbb{N}) \Rightarrow (|\operatorname{image}(h, m)| = |m|) \end{array}}}

\vnbbookfile{theorem-library/card-laws}{theorem-library/card-laws}
\noindent{\small theorem-library/card-laws.scm -- the former CARD axioms whose counterpart has a different SHAPE, restated character for character as the axioms stood and proven (2026-09-20, batch 9-B: CARD := CARD-STAR, the user's decision).}\par
\vnbprop{card-in-ord}{\vnbstmt{\forall a \in \mathbf{Set}.\; (|a| \in \mathrm{Ord})}}
\vnbprop{card-insert}{\vnbstmt{\begin{array}{@{}l@{}} \forall a, x \in \mathbf{Set}. \\ \quad ((|a| \in \mathbb{N}) \wedge \neg (x \in a)) \Rightarrow  \\ \quad \quad |(a \cup \{x\})| = (|a| + 1) \end{array}}}
\vnbprop{card-finite-bij}{\vnbstmt{\begin{array}{@{}l@{}} \forall a \in \mathbf{Set}. \\ \quad (|a| \in \mathbb{N}) \Rightarrow \exists i.\; (i \in \mathrm{Bij}(\mathrm{OS}(|a|), a)) \end{array}}}
\vnbprop{card-union-disjoint}{\vnbstmt{\begin{array}{@{}l@{}} \forall a, b \in \mathbf{Set}. \\ \quad ((|a| \in \mathbb{N}) \wedge ((|b| \in \mathbb{N}) \wedge ((a \cap b) = \emptyset))) \Rightarrow  \\ \quad \quad |(a \cup b)| = (|a| + |b|) \end{array}}}
\vnbprop{card-image-injection}{\vnbstmt{\begin{array}{@{}l@{}} \forall m, d.; \forall i \in \operatorname{injection}(m, d). \\ \quad ((m \in \mathbf{Set}) \wedge (|m| \in \mathbb{N})) \Rightarrow (|\operatorname{image}(i, m)| = |m|) \end{array}}}

\vnbbookfile{theorem-library/card-singleton-proof}{theorem-library/card-singleton-proof}
\noindent{\small card-singleton-proof.scm -- the cardinality of a singleton is 1, proven.}\par
\vnbprop{card-singleton}{\vnbstmt{\forall x \in \mathbf{Set}.\; (|\{x\}| = 1)}}

\vnbbookfile{theorem-library/rake-fin-enum}{theorem-library/rake-fin-enum}
\noindent{\small theorem-library/rake-fin-enum.scm -- fin-enum-is-bijection, moved out of theorem-library/rake-inverse-bij.scm on 2026-09-20 (batch 9-B, CARD := CARD-STAR).}\par
\vnbprop{fin-enum-is-bijection}{\vnbstmt{\begin{array}{@{}l@{}} \forall s \in \mathbf{Set}. \\ \quad (|s| \in \mathbb{N}) \Rightarrow (\operatorname{fin\text{-}enum}(s) \in \mathrm{Bij}(\mathrm{OS}(|s|), s)) \end{array}}}

\vnbbookfile{theorem-library/difference-laws}{theorem-library/difference-laws}
\noindent{\small difference-laws.scm -- the two DIFFERENCE laws, PROVEN.}\par
\vnbprop{difference-membership}{\vnbstmt{\forall u, b, x.\; ((x \in \operatorname{difference}(u, b)) \Leftrightarrow ((x \in u) \wedge \neg (x \in b)))}}
\vnbprop{difference-set}{\vnbstmt{\forall x \in \mathbf{Set}, b.\; (\operatorname{difference}(x, b) \in \mathbf{Set})}}

\vnbbookfile{theorem-library/intersection-of-laws}{theorem-library/intersection-of-laws}
\noindent{\small intersection-of-laws.scm -- the laws of INTERSECTION-OF, PROVEN.}\par
\vnbprop{intersection-of-unfold}{\vnbstmt{\begin{array}{@{}l@{}} \forall c. \\ \quad (\operatorname{intersection\text{-}of}(c) \simeq \{v \in \operatorname{big\text{-}union}(w, c, w) : \forall u \in c.\; (v \in u)\}) \end{array}}}
\vnbprop{intersection-of-membership}{\vnbstmt{\begin{array}{@{}l@{}} \forall c, x. \\ \quad (x \in \operatorname{intersection\text{-}of}(c)) \Leftrightarrow (\exists u.\; (u \in c) \wedge \forall u \in c.\; (x \in u)) \end{array}}}
\vnbprop{intersection-of-subset-member}{\vnbstmt{\forall c.; \forall a \in c.\; (\operatorname{intersection\text{-}of}(c) \subseteq a)}}
\vnbprop{intersection-of-empty-family}{\vnbstmt{\forall c.\; (\neg \exists u.\; (u \in c) \Rightarrow (\operatorname{intersection\text{-}of}(c) = \emptyset))}}
\vnbprop{intersection-of-is-set}{\vnbstmt{\forall c.\; (\operatorname{intersection\text{-}of}(c) \in \mathbf{Set})}}
\vnbprop{intersection-of-antitone}{\vnbstmt{\begin{array}{@{}l@{}} \forall c, d. \\ \quad ((c \subseteq d) \wedge \exists u.\; (u \in c)) \Rightarrow  \\ \quad \quad \operatorname{intersection\text{-}of}(d) \subseteq \operatorname{intersection\text{-}of}(c) \end{array}}}
\vnbprop{complement-in-big-union}{\vnbstmt{\begin{array}{@{}l@{}} \forall y, c, x. \\ \quad (x \in \operatorname{complement\text{-}in}(y, \operatorname{big\text{-}union}(w, c, w))) \Leftrightarrow ((x \in y) \wedge \forall u \in c.\; \neg (x \in u)) \end{array}}}
\vnbprop{complement-in-intersection-of}{\vnbstmt{\begin{array}{@{}l@{}} \forall y, c. \\ \quad \exists u.\; (u \in c) \Rightarrow  \\ \quad \quad \operatorname{complement\text{-}in}(y, \operatorname{intersection\text{-}of}(c)) = \operatorname{big\text{-}union}(w, c, \operatorname{complement\text{-}in}(y, w)) \end{array}}}

\vnbbookfile{theorem-library/rake-complement-laws}{theorem-library/rake-complement-laws}
\noindent{\small rake-complement-laws.scm -- two set-theoretic bricks the compactness duality needs, PROVEN (2026-09-20, batch 11-E; helper prefix `rcl-').}\par
\vnbprop{complement-in-double}{\vnbstmt{\forall y, a.\; ((a \subseteq y) \Rightarrow (\operatorname{complement\text{-}in}(y, \operatorname{complement\text{-}in}(y, a)) = a))}}
\vnbprop{family-of-subsets-is-set}{\vnbstmt{\forall y, c.\; (((y \in \mathbf{Set}) \wedge \forall u \in c.\; (u \subseteq y)) \Rightarrow (c \in \mathbf{Set}))}}

\vnbbookfile{theorem-library/rake-fun-codomain}{theorem-library/rake-fun-codomain}
\noindent{\small theorem-library/rake-fun-codomain.scm -- codomain widening for FUN, PROVEN.}\par
\vnbprop{fun-codomain-superset}{\vnbstmt{\forall f, a, b, c.\; (((f \in (a \to b)) \wedge (b \subseteq c)) \Rightarrow (f \in (a \to c)))}}

\vnbbookfile{theorem-library/rake-rr-order-leaves}{theorem-library/rake-rr-order-leaves}
\noindent{\small theorem-library/rake-rr-order-leaves.scm -- the two asserted order leaves of `deriv-pos-strictly-increasing', PROVEN.}\par
\vnbprop{rr-lt-from-diff-pos}{\vnbstmt{\forall x, y \in \mathbb{R}.\; ((0 < (y - x)) \Rightarrow (x < y))}}

\vnbbookfile{theorem-library/rr-sup-approx}{theorem-library/rr-sup-approx}
\noindent{\small rr-sup-approx.scm -- the SUP APPROXIMATION lemma, PROVEN.}\par
\vnbprop{rr-sup-approx}{\vnbstmt{\begin{array}{@{}l@{}} \forall s, d. \\ \quad ((s \subseteq \mathbb{R}) \wedge (\exists x.\; (x \in s) \wedge (\operatorname{rr\text{-}bounded\text{-}above}(s) \wedge ((d \in \mathbb{R}) \wedge (0 < d))))) \Rightarrow  \\ \quad \quad \exists w \in s.\; ((\operatorname{sup}(s) - d) < w) \end{array}}}

\vnbbookfile{theorem-library/nn-unbounded-in-rr}{theorem-library/nn-unbounded-in-rr}
\noindent{\small nn-unbounded-in-rr.scm -- the ARCHIMEDEAN PROPERTY, PROVEN.}\par
\vnbprop{nn-unbounded-in-rr}{\vnbstmt{\forall x \in \mathbb{R}.\; \exists n \in \mathbb{N}.\; (x < n)}}

\vnbbookfile{theorem-library/nn-recip-succ-small}{theorem-library/nn-recip-succ-small}
\noindent{\small nn-recip-succ-small.scm -- 1/(n+1) gets below any positive eps, PROVEN.}\par
\vnbprop{nn-recip-succ-small}{\vnbstmt{\forall s.\; (\operatorname{pos\text{-}rr}(s) \Rightarrow \exists n \in \mathbb{N}.\; (\operatorname{recip}((n + 1)) < s))}}

\vnbbookfile{theorem-library/discrete-space}{theorem-library/discrete-space}
\noindent{\small A file shared across books; the entries that belong to this one.}\par
\noindent{\small discrete-space.scm -- THE FIRST WITNESS OF A MEASURABLE SPACE, and the first outright inhabitant of TOP-SPACE.}\par
\vnbprop{power-mem-sethood}{\vnbstmt{\forall a.; \forall b \in \operatorname{power}(a).\; (b \in \mathbf{Set})}}
\vnbprop{power-mem-in}{\vnbstmt{\forall a, b, z.\; (((b \in \operatorname{power}(a)) \wedge (z \in b)) \Rightarrow (z \in a))}}
\vnbprop{power-mem-intro}{\vnbstmt{\forall a.; \forall b \in \mathbf{Set}.\; (\forall z \in b.\; (z \in a) \Rightarrow (b \in \operatorname{power}(a)))}}
\vnbprop{power-whole-in}{\vnbstmt{\forall a \in \mathbf{Set}.\; (a \in \operatorname{power}(a))}}
\vnbprop{power-empty-in}{\vnbstmt{\forall a.\; (\emptyset \in \operatorname{power}(a))}}
\vnbprop{power-inter-closed}{\vnbstmt{\begin{array}{@{}l@{}} \forall a, u, v. \\ \quad ((u \in \operatorname{power}(a)) \wedge (v \in \operatorname{power}(a))) \Rightarrow ((u \cap v) \in \operatorname{power}(a)) \end{array}}}
\vnbprop{power-complement-closed}{\vnbstmt{\forall a \in \mathbf{Set}, b.\; (\operatorname{complement\text{-}in}(a, b) \in \operatorname{power}(a))}}
\vnbprop{power-big-union-closed}{\vnbstmt{\begin{array}{@{}l@{}} \forall a, m. \\ \quad ((a \in \mathbf{Set}) \wedge (m \subseteq \operatorname{power}(a))) \Rightarrow (\operatorname{big\text{-}union}(u, m, u) \in \operatorname{power}(a)) \end{array}}}
\vnbprop{power-fun-union-closed}{\vnbstmt{\forall a.; \forall f \in (\mathbb{N} \to \operatorname{power}(a)).\; (\operatorname{big\text{-}union}(n, \mathbb{N}, f(n)) \in \operatorname{power}(a))}}

\vnbbookfile{structure-library/subtype-laws}{structure-library/subtype-laws}
\noindent{\small A file shared across books; the entries that belong to this one.}\par
\noindent{\small subtype-laws.scm -- the trivial subtype-subsumption laws, PROVEN.}\par
\vnbprop{bijection-is-injection}{\vnbstmt{\forall x, y.; \forall i \in \mathrm{Bij}(x, y).\; (i \in \operatorname{injection}(x, y))}}

\vnbbookfile{theorem-library/lambda-slot-apply}{theorem-library/lambda-slot-apply}
\noindent{\small lambda-slot-apply.scm -- taking an operation SLOT down to the surface.}\par
\vnbprop{lam-slot-add-apply}{\vnbstmt{\begin{array}{@{}l@{}} \forall m, a, b. \\ \quad ((a \in m) \wedge (b \in m)) \Rightarrow  \\ \quad \quad ((\lambda \langle \operatorname{x\_}, \operatorname{y\_} \rangle \in (m \times m).\; (\operatorname{x\_} + \operatorname{y\_}))(a, b) \simeq (a + b)) \end{array}}}
\vnbprop{lam-slot-mul-apply}{\vnbstmt{\begin{array}{@{}l@{}} \forall m, a, b. \\ \quad ((a \in m) \wedge (b \in m)) \Rightarrow  \\ \quad \quad ((\lambda \langle \operatorname{x\_}, \operatorname{y\_} \rangle \in (m \times m).\; (\operatorname{x\_} \cdot \operatorname{y\_}))(a, b) \simeq (a \cdot b)) \end{array}}}
\vnbprop{lam-slot-neg-apply}{\vnbstmt{\forall m.; \forall a \in m.\; ((\lambda x \in m.\; -x)(a) \simeq -a)}}

\vnbbookfile{theorem-library/nn-cancel}{theorem-library/nn-cancel}
\noindent{\small nn-cancel.scm -- CANCELLATION ON NN, proven in the core (2026-10-07).}\par
\vnbprop{nn-add-cancel}{\vnbstmt{\forall a, b, c \in \mathbb{N}.\; (((a + c) = (b + c)) \Rightarrow (a = b))}}
\vnbprop{nn-mul-cancel}{\vnbstmt{\forall a, b, c \in \mathbb{N}.\; ((\neg (c = 0) \wedge ((a \cdot c) = (b \cdot c))) \Rightarrow (a = b))}}

\vnbbookfile{theorem-library/nn-mod3-proof}{theorem-library/nn-mod3-proof}
\noindent{\small nn-mod3-proof.scm -- mod-3 arithmetic on NN, the mirror of nn-parity-proof.scm.}\par
\vnbprop{nn-three-is-succ3-zero}{\vnbstmt{3 = (((0 + 1) + 1) + 1)}}
\vnbprop{nn-plus-three}{\vnbstmt{\forall x \in \mathbb{N}.\; ((x + 3) = (((x + 1) + 1) + 1))}}
\vnbprop{nn-three-mul-succ}{\vnbstmt{\forall k \in \mathbb{N}.\; ((3 \cdot (k + 1)) = ((((3 \cdot k) + 1) + 1) + 1))}}
\vnbprop{nn-trichotomy-3}{\vnbstmt{\begin{array}{@{}l@{}} \forall n \in \mathbb{N}. \exists k \in \mathbb{N}. \\ \quad (n = (3 \cdot k)) \vee ((n = ((3 \cdot k) + 1)) \vee (n = (((3 \cdot k) + 1) + 1))) \end{array}}}
\vnbprop{nn-3-not-succ}{\vnbstmt{\forall x, y \in \mathbb{N}.\; \neg ((3 \cdot x) = ((3 \cdot y) + 1))}}
\vnbprop{nn-3-div-square}{\vnbstmt{\begin{array}{@{}l@{}} \forall p \in \mathbb{N}. \\ \quad \exists m \in \mathbb{N}.\; ((p \cdot p) = (3 \cdot m)) \Rightarrow \exists k \in \mathbb{N}.\; (p = (3 \cdot k)) \end{array}}}
\vnbprop{nn-3-cancel}{\vnbstmt{\forall x, y \in \mathbb{N}.\; (((3 \cdot x) = (3 \cdot y)) \Rightarrow (x = y))}}
\vnbprop{nn-lt-triple}{\vnbstmt{\forall k \in \mathbb{N}.\; (\neg (k = 0) \Rightarrow (k < (3 \cdot k)))}}

\vnbbookfile{theorem-library/nn-order-via-rr}{theorem-library/nn-order-via-rr}
\noindent{\small nn-order-via-rr.scm -- elementary NN order facts, PROVEN by taking the contradiction in RR.}\par
\vnbprop{nn-not-le-zero-pos}{\vnbstmt{\forall j \in \mathbb{N}.\; ((1 \leq j) \Rightarrow \neg (j \leq 0))}}
\vnbprop{nn-succ-le-antisym}{\vnbstmt{\forall a, b \in \mathbb{N}.\; (((a + 1) \leq b) \Rightarrow \neg (b \leq a))}}

\vnbbookfile{theorem-library/rake-intervals}{theorem-library/rake-intervals}
\noindent{\small rake-intervals.scm -- the interval / monus-predecessor supports of structure-library/order-lemmas.scm, PROVEN (batch K, 2026-09-17).}\par
\vnbprop{nn-succ-le-cancel}{\vnbstmt{\forall a, b \in \mathbb{N}.\; (((a + 1) \leq (b + 1)) \Rightarrow (a \leq b))}}
\vnbprop{nn-minus-succ-1}{\vnbstmt{\forall z \in \mathbb{N}.\; (\operatorname{nn\text{-}minus}((z + 1), 1) = z)}}
\vnbprop{one-in-interval}{\vnbstmt{\forall n \in \mathbb{N}.\; (1 \in \operatorname{interval}(1, (n + 1)))}}
\vnbprop{succ-not-one}{\vnbstmt{\forall q.; \forall z \in \operatorname{interval}(1, q).\; \neg ((z + 1) = 1)}}
\vnbprop{nn-minus-1-inj}{\vnbstmt{\begin{array}{@{}l@{}} \forall i, j. \\ \quad ((i \in \mathbb{N}) \wedge ((j \in \mathbb{N}) \wedge ((1 \leq i) \wedge ((1 \leq j) \wedge \neg (i = j))))) \Rightarrow  \\ \quad \quad \neg (\operatorname{nn\text{-}minus}(i, 1) = \operatorname{nn\text{-}minus}(j, 1)) \end{array}}}
\vnbprop{succ-in-interval}{\vnbstmt{\forall q \in \mathbb{N}, z \in \operatorname{interval}(1, q).\; ((z + 1) \in \operatorname{interval}(1, (q + 1)))}}
\vnbprop{pred-in-interval}{\vnbstmt{\begin{array}{@{}l@{}} \forall p \in \mathbb{N}, i \in \operatorname{interval}(1, (p + 1)). \\ \quad \neg (i = 1) \Rightarrow (\operatorname{nn\text{-}minus}(i, 1) \in \operatorname{interval}(1, p)) \end{array}}}

\vnbbookfile{theorem-library/zz-order}{theorem-library/zz-order}
\noindent{\small zz-order.scm -- the ORDER THEORY OF ZZ, proven.}\par
\vnbprop{zz-abs-in-nn}{\vnbstmt{\forall a \in \mathbb{Z}.\; (\lvert a \rvert \in \mathbb{N})}}
\vnbprop{zz-abs-cases}{\vnbstmt{\forall a \in \mathbb{Z}.\; ((a = \lvert a \rvert) \vee (a = -\lvert a \rvert))}}
\vnbprop{zz-trichotomy}{\vnbstmt{\forall a \in \mathbb{Z}.\; ((a < 0) \vee ((a = 0) \vee (0 < a)))}}
\vnbprop{zz-discrete}{\vnbstmt{\forall a \in \mathbb{Z}.\; ((0 < a) \Rightarrow (1 \leq a))}}
\vnbprop{zz-lt-succ-le}{\vnbstmt{\forall a, b \in \mathbb{Z}.\; ((a < b) \Rightarrow ((a + 1) \leq b))}}
\vnbprop{zz-nonneg-in-nn}{\vnbstmt{\forall a \in \mathbb{Z}.\; ((0 \leq a) \Rightarrow (a \in \mathbb{N}))}}

\vnbbookfile{theorem-library/zz-parity-proof}{theorem-library/zz-parity-proof}
\noindent{\small zz-parity-proof.scm -- parity on the INTEGERS, lifted from NN.}\par
\par\medskip\noindent\textbf{Definition} (\texttt{even}).\;For every $x$,  $\operatorname{even}(x) \Leftrightarrow \exists k \in \mathbb{Z}.\; (x = \cdot(2, k))$
\par\medskip\noindent\textbf{Definition} (\texttt{odd}).\;For every $x$,  $\operatorname{odd}(x) \Leftrightarrow \exists k \in \mathbb{Z}.\; (x = +(\cdot(2, k), 1))$
\vnbdef{even-def}{\vnbstmt{\forall x.\; (\operatorname{even}(x) \Leftrightarrow \exists k \in \mathbb{Z}.\; (x = (2 \cdot k)))}}
\vnbdef{odd-def}{\vnbstmt{\forall x.\; (\operatorname{odd}(x) \Leftrightarrow \exists k \in \mathbb{Z}.\; (x = ((2 \cdot k) + 1)))}}
\vnbprop{zz-parity}{\vnbstmt{\forall a \in \mathbb{Z}.\; (\operatorname{even}(a) \vee \operatorname{odd}(a))}}
\vnbprop{nn-2p-neq-1}{\vnbstmt{\forall p \in \mathbb{N}.\; \neg ((2 \cdot p) = 1)}}
\vnbprop{nn-2p-neq-neg1}{\vnbstmt{\forall p \in \mathbb{N}.\; \neg ((2 \cdot p) = -1)}}
\vnbprop{zz-2m-neq-1}{\vnbstmt{\forall m \in \mathbb{Z}.\; \neg ((2 \cdot m) = 1)}}
\vnbprop{zz-not-even-and-odd}{\vnbstmt{\forall x \in \mathbb{Z}.\; \neg (\operatorname{even}(x) \wedge \operatorname{odd}(x))}}
\vnbprop{zz-even-iff-not-odd}{\vnbstmt{\forall x \in \mathbb{Z}.\; (\operatorname{even}(x) \Leftrightarrow \neg \operatorname{odd}(x))}}
\vnbprop{zz-even-mul}{\vnbstmt{\forall x, y \in \mathbb{Z}.\; (\operatorname{even}(x) \Rightarrow \operatorname{even}((x \cdot y)))}}
\vnbprop{zz-odd-mul-odd}{\vnbstmt{\forall x, y \in \mathbb{Z}.\; ((\operatorname{odd}(x) \wedge \operatorname{odd}(y)) \Rightarrow \operatorname{odd}((x \cdot y)))}}
\vnbprop{zz-odd-sq-odd}{\vnbstmt{\forall x \in \mathbb{Z}.\; (\operatorname{odd}(x) \Rightarrow \operatorname{odd}((x \cdot x)))}}

\vnbbookfile{theorem-library/nn-pos-is-succ}{theorem-library/nn-pos-is-succ}
\noindent{\small nn-pos-is-succ.scm -- a natural at least 1 is a successor.}\par
\vnbprop{nn-pos-is-succ}{\vnbstmt{\forall n \in \mathbb{N}.\; ((1 \leq n) \Rightarrow \exists q \in \mathbb{N}.\; (n = (q + 1)))}}

\vnbbookfile{theorem-library/nn-order-proof}{theorem-library/nn-order-proof}
\noindent{\small nn-order-proof.scm -- three elementary facts about $<$= and + on NN, PROVEN.}\par
\vnbprop{nn-le-zero-is-zero}{\vnbstmt{\forall a \in \mathbb{N}.\; ((a \leq 0) \Rightarrow (a = 0))}}
\vnbprop{nn-add-le-mono}{\vnbstmt{\forall c, a, b \in \mathbb{N}.\; ((b \leq c) \Rightarrow ((a + b) \leq (a + c)))}}
\vnbprop{nn-le-add-left}{\vnbstmt{\forall a, b.\; (((a \in \mathbb{N}) \wedge (b \in \mathbb{N})) \Rightarrow (b \leq (a + b)))}}
\vnbprop{nn-le-antisym}{\vnbstmt{\forall a, b \in \mathbb{N}.\; (((a \leq b) \wedge (b \leq a)) \Rightarrow (a = b))}}
\vnbprop{nn-succ-not-le}{\vnbstmt{\forall n \in \mathbb{N}.\; \neg ((n + 1) \leq n)}}
\vnbprop{nn-le-gap}{\vnbstmt{\forall k, m \in \mathbb{N}.\; ((m \leq k) \Rightarrow \exists d \in \mathbb{N}.\; (k = (m + d)))}}

\vnbbookfile{theorem-library/rake-finsum-typing}{theorem-library/rake-finsum-typing}
\noindent{\small A file shared across books; the entries that belong to this one.}\par
\noindent{\small rake-finsum-typing.scm -- the finsum / monoid-power typing leaves of the rake (batch G), PROVEN, together with the guarded form of the one statement in the batch that is NOT provable as stated.}\par
\vnbprop{funcomp-succ-type}{\vnbstmt{\begin{array}{@{}l@{}} \forall x.; \forall q \in \mathbb{N}, f \in (\operatorname{interval}(1, (q + 1)) \to x). \\ \quad (\lambda z \in \operatorname{interval}(1, q).\; f((z + 1))) \in (\operatorname{interval}(1, q) \to x) \end{array}}}

\vnbbookfile{theorem-library/hom-laws}{theorem-library/hom-laws}
\noindent{\small A file shared across books; the entries that belong to this one.}\par
\noindent{\small hom-laws.scm -- the models of a structure form a CATEGORY: identity arrows, composition, and the hom-set as a set.}\par
\vnbprop{hom-semigroup-compose}{\vnbstmt{\begin{array}{@{}l@{}} \forall a, b, c, f, g. \\ \quad (\operatorname{is\text{-}hom\text{-}semigroup}(a, b, f) \wedge \operatorname{is\text{-}hom\text{-}semigroup}(b, c, g)) \Rightarrow  \\ \quad \quad \operatorname{is\text{-}hom\text{-}semigroup}(a, c, \operatorname{compose}(g, f)) \end{array}}}
\vnbprop{hom-monoid-compose}{\vnbstmt{\begin{array}{@{}l@{}} \forall a, b, c, f, g. \\ \quad (\operatorname{is\text{-}hom\text{-}monoid}(a, b, f) \wedge \operatorname{is\text{-}hom\text{-}monoid}(b, c, g)) \Rightarrow  \\ \quad \quad \operatorname{is\text{-}hom\text{-}monoid}(a, c, \operatorname{compose}(g, f)) \end{array}}}
\vnbprop{hom-comm-monoid-compose}{\vnbstmt{\begin{array}{@{}l@{}} \forall a, b, c, f, g. \\ \quad (\operatorname{is\text{-}hom\text{-}comm\text{-}monoid}(a, b, f) \wedge \operatorname{is\text{-}hom\text{-}comm\text{-}monoid}(b, c, g)) \Rightarrow  \\ \quad \quad \operatorname{is\text{-}hom\text{-}comm\text{-}monoid}(a, c, \operatorname{compose}(g, f)) \end{array}}}
\vnbprop{hom-group-compose}{\vnbstmt{\begin{array}{@{}l@{}} \forall a, b, c, f, g. \\ \quad (\operatorname{is\text{-}hom\text{-}group}(a, b, f) \wedge \operatorname{is\text{-}hom\text{-}group}(b, c, g)) \Rightarrow  \\ \quad \quad \operatorname{is\text{-}hom\text{-}group}(a, c, \operatorname{compose}(g, f)) \end{array}}}
\vnbprop{hom-abelian-group-compose}{\vnbstmt{\begin{array}{@{}l@{}} \forall a, b, c, f, g. \\ \quad (\operatorname{is\text{-}hom\text{-}abelian\text{-}group}(a, b, f) \wedge \operatorname{is\text{-}hom\text{-}abelian\text{-}group}(b, c, g)) \Rightarrow  \\ \quad \quad \operatorname{is\text{-}hom\text{-}abelian\text{-}group}(a, c, \operatorname{compose}(g, f)) \end{array}}}
\vnbprop{hom-ring-compose}{\vnbstmt{\begin{array}{@{}l@{}} \forall a, b, c, f, g. \\ \quad (\operatorname{is\text{-}hom\text{-}ring}(a, b, f) \wedge \operatorname{is\text{-}hom\text{-}ring}(b, c, g)) \Rightarrow  \\ \quad \quad \operatorname{is\text{-}hom\text{-}ring}(a, c, \operatorname{compose}(g, f)) \end{array}}}
\vnbprop{hom-commutative-ring-compose}{\vnbstmt{\begin{array}{@{}l@{}} \forall a, b, c, f, g. \\ \quad (\operatorname{is\text{-}hom\text{-}commutative\text{-}ring}(a, b, f) \wedge \operatorname{is\text{-}hom\text{-}commutative\text{-}ring}(b, c, g)) \Rightarrow  \\ \quad \quad \operatorname{is\text{-}hom\text{-}commutative\text{-}ring}(a, c, \operatorname{compose}(g, f)) \end{array}}}
\vnbprop{hom-integral-domain-compose}{\vnbstmt{\begin{array}{@{}l@{}} \forall a, b, c, f, g. \\ \quad (\operatorname{is\text{-}hom\text{-}integral\text{-}domain}(a, b, f) \wedge \operatorname{is\text{-}hom\text{-}integral\text{-}domain}(b, c, g)) \Rightarrow  \\ \quad \quad \operatorname{is\text{-}hom\text{-}integral\text{-}domain}(a, c, \operatorname{compose}(g, f)) \end{array}}}
\vnbprop{hom-euclidean-ring-compose}{\vnbstmt{\begin{array}{@{}l@{}} \forall a, b, c, f, g. \\ \quad (\operatorname{is\text{-}hom\text{-}euclidean\text{-}ring}(a, b, f) \wedge \operatorname{is\text{-}hom\text{-}euclidean\text{-}ring}(b, c, g)) \Rightarrow  \\ \quad \quad \operatorname{is\text{-}hom\text{-}euclidean\text{-}ring}(a, c, \operatorname{compose}(g, f)) \end{array}}}
\vnbprop{hom-pid-compose}{\vnbstmt{\begin{array}{@{}l@{}} \forall a, b, c, f, g. \\ \quad (\operatorname{is\text{-}hom\text{-}pid}(a, b, f) \wedge \operatorname{is\text{-}hom\text{-}pid}(b, c, g)) \Rightarrow  \\ \quad \quad \operatorname{is\text{-}hom\text{-}pid}(a, c, \operatorname{compose}(g, f)) \end{array}}}
\vnbprop{hom-field-ring-compose}{\vnbstmt{\begin{array}{@{}l@{}} \forall a, b, c, f, g. \\ \quad (\operatorname{is\text{-}hom\text{-}field\text{-}ring}(a, b, f) \wedge \operatorname{is\text{-}hom\text{-}field\text{-}ring}(b, c, g)) \Rightarrow  \\ \quad \quad \operatorname{is\text{-}hom\text{-}field\text{-}ring}(a, c, \operatorname{compose}(g, f)) \end{array}}}
\vnbprop{hom-normed-ag-compose}{\vnbstmt{\begin{array}{@{}l@{}} \forall a, b, c, f, g. \\ \quad (\operatorname{is\text{-}hom\text{-}normed\text{-}ag}(a, b, f) \wedge \operatorname{is\text{-}hom\text{-}normed\text{-}ag}(b, c, g)) \Rightarrow  \\ \quad \quad \operatorname{is\text{-}hom\text{-}normed\text{-}ag}(a, c, \operatorname{compose}(g, f)) \end{array}}}
\vnbprop{hom-normed-field-compose}{\vnbstmt{\begin{array}{@{}l@{}} \forall a, b, c, f, g. \\ \quad (\operatorname{is\text{-}hom\text{-}normed\text{-}field}(a, b, f) \wedge \operatorname{is\text{-}hom\text{-}normed\text{-}field}(b, c, g)) \Rightarrow  \\ \quad \quad \operatorname{is\text{-}hom\text{-}normed\text{-}field}(a, c, \operatorname{compose}(g, f)) \end{array}}}
\vnbprop{hom-module-compose}{\vnbstmt{\begin{array}{@{}l@{}} \forall a, b, c, f, g. \\ \quad (\operatorname{is\text{-}hom\text{-}module}(a, b, f) \wedge \operatorname{is\text{-}hom\text{-}module}(b, c, g)) \Rightarrow  \\ \quad \quad \operatorname{is\text{-}hom\text{-}module}(a, c, \operatorname{compose}(g, f)) \end{array}}}
\vnbprop{hom-vector-space-compose}{\vnbstmt{\begin{array}{@{}l@{}} \forall a, b, c, f, g. \\ \quad (\operatorname{is\text{-}hom\text{-}vector\text{-}space}(a, b, f) \wedge \operatorname{is\text{-}hom\text{-}vector\text{-}space}(b, c, g)) \Rightarrow  \\ \quad \quad \operatorname{is\text{-}hom\text{-}vector\text{-}space}(a, c, \operatorname{compose}(g, f)) \end{array}}}
\vnbprop{hom-normed-vector-space-compose}{\vnbstmt{\begin{array}{@{}l@{}} \forall a, b, c, f, g. \\ \quad (\operatorname{is\text{-}hom\text{-}normed\text{-}vector\text{-}space}(a, b, f) \wedge \operatorname{is\text{-}hom\text{-}normed\text{-}vector\text{-}space}(b, c, g)) \Rightarrow  \\ \quad \quad \operatorname{is\text{-}hom\text{-}normed\text{-}vector\text{-}space}(a, c, \operatorname{compose}(g, f)) \end{array}}}
\vnbprop{hom-complex-inner-product-space-compose}{\vnbstmt{\begin{array}{@{}l@{}} \forall a, b, c, f, g. \\ \quad (\operatorname{is\text{-}hom\text{-}complex\text{-}inner\text{-}product\text{-}space}(a, b, f) \wedge \operatorname{is\text{-}hom\text{-}complex\text{-}inner\text{-}product\text{-}space}(b, c, g)) \Rightarrow  \\ \quad \quad \operatorname{is\text{-}hom\text{-}complex\text{-}inner\text{-}product\text{-}space}(a, c, \operatorname{compose}(g, f)) \end{array}}}
\vnbprop{hom-metric-space-compose}{\vnbstmt{\begin{array}{@{}l@{}} \forall a, b, c, f, g. \\ \quad (\operatorname{is\text{-}hom\text{-}metric\text{-}space}(a, b, f) \wedge \operatorname{is\text{-}hom\text{-}metric\text{-}space}(b, c, g)) \Rightarrow  \\ \quad \quad \operatorname{is\text{-}hom\text{-}metric\text{-}space}(a, c, \operatorname{compose}(g, f)) \end{array}}}
\vnbprop{hom-pseudometric-space-compose}{\vnbstmt{\begin{array}{@{}l@{}} \forall a, b, c, f, g. \\ \quad (\operatorname{is\text{-}hom\text{-}pseudometric\text{-}space}(a, b, f) \wedge \operatorname{is\text{-}hom\text{-}pseudometric\text{-}space}(b, c, g)) \Rightarrow  \\ \quad \quad \operatorname{is\text{-}hom\text{-}pseudometric\text{-}space}(a, c, \operatorname{compose}(g, f)) \end{array}}}
\vnbprop{hom-c-metric-space-compose}{\vnbstmt{\begin{array}{@{}l@{}} \forall a, b, c, f, g. \\ \quad (\operatorname{is\text{-}hom\text{-}c\text{-}metric\text{-}space}(a, b, f) \wedge \operatorname{is\text{-}hom\text{-}c\text{-}metric\text{-}space}(b, c, g)) \Rightarrow  \\ \quad \quad \operatorname{is\text{-}hom\text{-}c\text{-}metric\text{-}space}(a, c, \operatorname{compose}(g, f)) \end{array}}}
\vnbprop{hom-top-space-compose}{\vnbstmt{\begin{array}{@{}l@{}} \forall s, t, w, f, g. \\ \quad (\operatorname{is\text{-}hom\text{-}top\text{-}space}(s, t, f) \wedge \operatorname{is\text{-}hom\text{-}top\text{-}space}(t, w, g)) \Rightarrow  \\ \quad \quad \operatorname{is\text{-}hom\text{-}top\text{-}space}(s, w, \operatorname{compose}(g, f)) \end{array}}}
\vnbprop{hom-metrizable-top-space-compose}{\vnbstmt{\begin{array}{@{}l@{}} \forall s, t, w, f, g. \\ \quad (\operatorname{is\text{-}hom\text{-}metrizable\text{-}top\text{-}space}(s, t, f) \wedge \operatorname{is\text{-}hom\text{-}metrizable\text{-}top\text{-}space}(t, w, g)) \Rightarrow  \\ \quad \quad \operatorname{is\text{-}hom\text{-}metrizable\text{-}top\text{-}space}(s, w, \operatorname{compose}(g, f)) \end{array}}}
\vnbprop{hom-ringoid-compose}{\vnbstmt{\begin{array}{@{}l@{}} \forall a, b, c, f, g. \\ \quad (\operatorname{is\text{-}hom\text{-}ringoid}(a, b, f) \wedge \operatorname{is\text{-}hom\text{-}ringoid}(b, c, g)) \Rightarrow  \\ \quad \quad \operatorname{is\text{-}hom\text{-}ringoid}(a, c, \operatorname{compose}(g, f)) \end{array}}}

\vnbbookfile{theorem-library/hom-kinds}{theorem-library/hom-kinds}
\noindent{\small A file shared across books; the entries that belong to this one.}\par
\noindent{\small hom-kinds.scm -- the category laws for the arrows the generic driver of hom-laws.scm does not reach: the FAMILY slot kinds, and the RELATION kind (SETOID, since batch 40; two carriers in batch 39).}\par
\vnbprop{hom-measurable-space-compose}{\vnbstmt{\begin{array}{@{}l@{}} \forall a, b, c, f, g. \\ \quad (\operatorname{is\text{-}hom\text{-}measurable\text{-}space}(a, b, f) \wedge \operatorname{is\text{-}hom\text{-}measurable\text{-}space}(b, c, g)) \Rightarrow  \\ \quad \quad \operatorname{is\text{-}hom\text{-}measurable\text{-}space}(a, c, \operatorname{compose}(g, f)) \end{array}}}
\vnbprop{hom-measure-space-compose}{\vnbstmt{\begin{array}{@{}l@{}} \forall a, b, c, f, g. \\ \quad (\operatorname{is\text{-}hom\text{-}measure\text{-}space}(a, b, f) \wedge \operatorname{is\text{-}hom\text{-}measure\text{-}space}(b, c, g)) \Rightarrow  \\ \quad \quad \operatorname{is\text{-}hom\text{-}measure\text{-}space}(a, c, \operatorname{compose}(g, f)) \end{array}}}
\vnbprop{hom-setoid-compose}{\vnbstmt{\begin{array}{@{}l@{}} \forall a, b, c, f, g. \\ \quad (\operatorname{is\text{-}hom\text{-}setoid}(a, b, f) \wedge \operatorname{is\text{-}hom\text{-}setoid}(b, c, g)) \Rightarrow  \\ \quad \quad \operatorname{is\text{-}hom\text{-}setoid}(a, c, \operatorname{compose}(g, f)) \end{array}}}
\vnbprop{measure-space-as-measurable-space-functorial}{\vnbstmt{\begin{array}{@{}l@{}} \forall a, b, f. \\ \quad \operatorname{is\text{-}hom\text{-}measure\text{-}space}(a, b, f) \Rightarrow  \\ \quad \quad \operatorname{is\text{-}hom\text{-}measurable\text{-}space}(\operatorname{measure\text{-}space\text{-}as\text{-}measurable\text{-}space}(a), \\ \quad \quad \quad \operatorname{measure\text{-}space\text{-}as\text{-}measurable\text{-}space}(b), f) \end{array}}}

\vnbbookfile{theorem-library/hom-functors}{theorem-library/hom-functors}
\noindent{\small A file shared across books; the entries that belong to this one.}\par
\noindent{\small hom-functors.scm -- the Hom functors: Hom(a, -) covariant, Hom(-, b) contravariant, as maps on hom-sets.}\par
\vnbprop{post-id}{\vnbstmt{\begin{array}{@{}l@{}} \forall a, b, s, f. \\ \quad ((a \in \mathbf{Set}) \wedge ((b \in \mathbf{Set}) \wedge ((f \in s) \wedge (f \in (a \to b))))) \Rightarrow  \\ \quad \quad (\lambda \mathit{hmpf} \in s.\; \operatorname{compose}(\operatorname{id\text{-}fun}(b), \mathit{hmpf}))(f) = f \end{array}}}
\vnbprop{pre-id}{\vnbstmt{\begin{array}{@{}l@{}} \forall a, b, s, f. \\ \quad ((a \in \mathbf{Set}) \wedge ((f \in s) \wedge (f \in (a \to b)))) \Rightarrow  \\ \quad \quad (\lambda \mathit{hmpf} \in s.\; \operatorname{compose}(\mathit{hmpf}, \operatorname{id\text{-}fun}(a)))(f) = f \end{array}}}
\vnbprop{post-compose}{\vnbstmt{\begin{array}{@{}l@{}} \forall a, b, c, u, s, d, f, h, k. \\ \quad a \in \mathbf{Set} \wedge \\ \quad b \in \mathbf{Set} \wedge \\ \quad f \in (a \to b) \wedge \\ \quad h \in (b \to c) \wedge \\ \quad k \in (c \to u) \wedge \\ \quad f \in s \wedge \\ \quad \operatorname{compose}(h, f) \in d \Rightarrow \\ \quad \quad (\lambda \mathit{hmpf} \in s.\; \operatorname{compose}(\operatorname{compose}(k, h), \mathit{hmpf}))(f) \\ \quad \quad =\; (\lambda \mathit{hmpf} \in d.\; \operatorname{compose}(k, \mathit{hmpf}))((\lambda \mathit{hmpf} \in s.\; \operatorname{compose}(h, \mathit{hmpf}))(f)) \end{array}}}
\vnbprop{pre-compose}{\vnbstmt{\begin{array}{@{}l@{}} \forall a, b, c, u, s, d, f, h, k. \\ \quad a \in \mathbf{Set} \wedge \\ \quad b \in \mathbf{Set} \wedge \\ \quad k \in (a \to b) \wedge \\ \quad h \in (b \to c) \wedge \\ \quad f \in (c \to u) \wedge \\ \quad f \in s \wedge \\ \quad \operatorname{compose}(f, h) \in d \Rightarrow \\ \quad \quad (\lambda \mathit{hmpf} \in s.\; \operatorname{compose}(\mathit{hmpf}, \operatorname{compose}(h, k)))(f) \\ \quad \quad =\; (\lambda \mathit{hmpf} \in d.\; \operatorname{compose}(\mathit{hmpf}, k))((\lambda \mathit{hmpf} \in s.\; \operatorname{compose}(\mathit{hmpf}, h))(f)) \end{array}}}

\vnbbookfile{theorem-library/hom-views}{theorem-library/hom-views}
\noindent{\small A file shared across books; the entries that belong to this one.}\par
\noindent{\small hom-views.scm -- a view's action on arrows.}\par
\vnbprop{ringoid-as-ring-functorial}{\vnbstmt{\begin{array}{@{}l@{}} \forall a, b, f. \\ \quad \operatorname{is\text{-}hom\text{-}ringoid}(a, b, f) \Rightarrow  \\ \quad \quad \operatorname{is\text{-}hom\text{-}ring}(\operatorname{ringoid\text{-}as\text{-}ring}(a), \operatorname{ringoid\text{-}as\text{-}ring}(b), f) \end{array}}}

\vnbbookfile{theorem-library/makeset-basics}{theorem-library/makeset-basics}
\noindent{\small makeset-basics.scm -- reading membership out of a literal brace set.}\par
\vnbprop{ms-eq-symm}{\vnbstmt{\forall x, y.\; ((x = y) \Leftrightarrow (y = x))}}
\vnbprop{ms-index-2}{\vnbstmt{\forall i \in \mathbb{N}.\; (((1 \leq i) \wedge (i \leq 2)) \Rightarrow ((i = 1) \vee (i = 2)))}}
\vnbprop{ms-nth-2}{\vnbstmt{\begin{array}{@{}l@{}} \forall a, b, i. \\ \quad ((a \in \mathbf{Set}) \wedge ((b \in \mathbf{Set}) \wedge ((i \in \mathbb{N}) \wedge ((1 \leq i) \wedge (i \leq 2))))) \Rightarrow  \\ \quad \quad (\operatorname{nth}(i, \langle a, b \rangle) = a) \vee (\operatorname{nth}(i, \langle a, b \rangle) = b) \end{array}}}
\vnbprop{makeset2-membership}{\vnbstmt{\begin{array}{@{}l@{}} \forall x, a, b. \\ \quad ((a \in \mathbf{Set}) \wedge (b \in \mathbf{Set})) \Rightarrow ((x \in \operatorname{make\text{-}set}(\langle a, b \rangle)) \Leftrightarrow ((x \in \mathbf{Set}) \wedge ((x = a) \vee (x = b)))) \end{array}}}
\vnbprop{makeset2-split}{\vnbstmt{\forall a, b \in \mathbf{Set}.\; (\operatorname{make\text{-}set}(\langle a, b \rangle) = (\{a\} \cup \{b\}))}}
\vnbprop{card-pair}{\vnbstmt{\forall a, b \in \mathbf{Set}.\; (\neg (a = b) \Rightarrow (|\operatorname{make\text{-}set}(\langle a, b \rangle)| = 2))}}

\vnbbookfile{theorem-library/empty-cardinality}{theorem-library/empty-cardinality}
\noindent{\small empty-cardinality.scm -- the two BACKWARD rungs for "this set is empty, so its cardinal is 0".}\par
\vnbprop{card-empty-le}{\vnbstmt{\forall x.\; ((x = \emptyset) \Rightarrow (|x| \leq 0))}}
\vnbprop{makeset-of-empty-tuple}{\vnbstmt{\forall l.\; ((l = \langle  \rangle) \Rightarrow (\operatorname{make\text{-}set}(l) = \emptyset))}}

\vnbbookfile{theorem-library/tuples-induction}{theorem-library/tuples-induction}
\noindent{\small tuples-induction.scm -- list induction, DERIVED.}\par
\vnbprop{tuples-induction}{\vnbstmt{\begin{array}{@{}l@{}} \forall a, c.; \forall l \in \operatorname{tuples}(a). \\ \quad ((\langle  \rangle \in c) \wedge \forall x, m.\; (((x \in a) \wedge ((m \in \operatorname{tuples}(a)) \wedge (m \in c))) \Rightarrow (\operatorname{cons}(x, m) \in c))) \Rightarrow  \\ \quad \quad l \in c \end{array}}}

\vnbbookfile{theorem-library/tuple-extensionality}{theorem-library/tuple-extensionality}
\noindent{\small tuple-extensionality.scm -- two tuples with the same length and the same entry at every index are equal; and, on top of it, matrix-entry-extensionality.}\par
\vnbprop{nth2-pair}{\vnbstmt{\forall a, b.\; (((a \in \mathbf{Set}) \wedge (b \in \mathbf{Set})) \Rightarrow (\operatorname{nth}(2, \langle a, b \rangle) \simeq b))}}
\vnbprop{entry-unfold}{\vnbstmt{\forall m, i, j.\; (\operatorname{entry}(m, i, j) \simeq \operatorname{nth}(j, \operatorname{nth}(i, m)))}}
\vnbprop{tuple-extensionality}{\vnbstmt{\begin{array}{@{}l@{}} \forall a, l, m. \\ \quad l \in \operatorname{tuples}(a) \wedge \\ \quad m \in \operatorname{tuples}(a) \wedge \\ \quad \operatorname{length}(l) = \operatorname{length}(m) \wedge \\ \quad \forall k \in \mathbb{N}.\; (((1 \leq k) \wedge (k \leq \operatorname{length}(l))) \Rightarrow (\operatorname{nth}(k, l) = \operatorname{nth}(k, m))) \Rightarrow \\ \quad \quad l = m \end{array}}}
\vnbprop{matrix-entry-extensionality}{\vnbstmt{\begin{array}{@{}l@{}} \forall m, n, x, p, q. \\ \quad p \in \operatorname{mat}(m, n, x) \wedge \\ \quad q \in \operatorname{mat}(m, n, x) \wedge \\ \quad \forall i \in \operatorname{interval}(1, m), j \in \operatorname{interval}(1, n). \\ \quad \quad \operatorname{entry}(p, i, j) = \operatorname{entry}(q, i, j) \Rightarrow \\ \quad \quad p = q \end{array}}}

\vnbbookfile{theorem-library/tuple-tabulation}{theorem-library/tuple-tabulation}
\noindent{\small tuple-tabulation.scm -- TABULATION: the missing generation principle for tuples and matrices, and the honest replacement for `matof-exists'.}\par
\vnbprop{tab-succ-add}{\vnbstmt{\forall a, b \in \mathbb{N}.\; (((a + 1) + b) = ((a + b) + 1))}}
\vnbprop{tuple-tabulation-offset}{\vnbstmt{\begin{array}{@{}l@{}} \forall n \in \mathbb{N}, a, f.; \forall k \in \mathbb{N}. \\ \quad \forall i \in \mathbb{N}.\; (((1 \leq i) \wedge (i \leq n)) \Rightarrow (f((k + i)) \in a)) \Rightarrow  \\ \quad \quad \exists l \in \operatorname{tuples}(a). \\ \quad \quad \quad ((\operatorname{length}(l) = n) \wedge \forall i \in \mathbb{N}.\; (((1 \leq i) \wedge (i \leq n)) \Rightarrow (\operatorname{nth}(i, l) = f((k + i))))) \end{array}}}
\vnbprop{row-tabulation-offset}{\vnbstmt{\begin{array}{@{}l@{}} \forall n \in \mathbb{N}, a, g, c.; \forall k \in \mathbb{N}. \\ \quad \forall j \in \mathbb{N}.\; (((1 \leq j) \wedge (j \leq n)) \Rightarrow (g(c, (k + j)) \in a)) \Rightarrow  \\ \quad \quad \exists l \in \operatorname{tuples}(a). \\ \quad \quad \quad ((\operatorname{length}(l) = n) \wedge \forall j \in \mathbb{N}.\; (((1 \leq j) \wedge (j \leq n)) \Rightarrow (\operatorname{nth}(j, l) = g(c, (k + j))))) \end{array}}}
\vnbprop{tuple-tabulation}{\vnbstmt{\begin{array}{@{}l@{}} \forall n \in \mathbb{N}, a, f. \\ \quad \forall i \in \mathbb{N}.\; (((1 \leq i) \wedge (i \leq n)) \Rightarrow (f(i) \in a)) \Rightarrow  \\ \quad \quad \exists l \in \operatorname{tuples}(a). \\ \quad \quad \quad (\operatorname{length}(l) = n) \wedge \forall i \in \mathbb{N}.\; (((1 \leq i) \wedge (i \leq n)) \Rightarrow (\operatorname{nth}(i, l) = f(i))) \end{array}}}
\vnbprop{row-tabulation}{\vnbstmt{\begin{array}{@{}l@{}} \forall n \in \mathbb{N}, a, g, c. \\ \quad \forall j \in \mathbb{N}.\; (((1 \leq j) \wedge (j \leq n)) \Rightarrow (g(c, j) \in a)) \Rightarrow  \\ \quad \quad \exists l \in \operatorname{tuples}(a). \\ \quad \quad \quad ((\operatorname{length}(l) = n) \wedge \forall j \in \mathbb{N}.\; (((1 \leq j) \wedge (j \leq n)) \Rightarrow (\operatorname{nth}(j, l) = g(c, j)))) \end{array}}}
\vnbprop{matrix-tabulation}{\vnbstmt{\begin{array}{@{}l@{}} \forall m \in \mathbb{N}, a, g, n.; \forall k \in \mathbb{N}. \\ \quad n \in \mathbb{N} \wedge \\ \quad \forall i, j \in \mathbb{N}. \\ \quad \quad (((1 \leq i) \wedge (i \leq m)) \wedge ((1 \leq j) \wedge (j \leq n))) \Rightarrow (g((k + i), j) \in a) \Rightarrow \\ \quad \quad \exists p \in \operatorname{tuples}(\operatorname{tuples}(a)). \\ \quad \quad \quad \operatorname{length}(p) = m \wedge \\ \quad \quad \quad \forall i \in \mathbb{N}.\; (((1 \leq i) \wedge (i \leq m)) \Rightarrow (\operatorname{length}(\operatorname{nth}(i, p)) = n)) \wedge \\ \quad \quad \quad \forall i, j \in \mathbb{N}. \\ \quad \quad \quad \quad (((1 \leq i) \wedge (i \leq m)) \wedge ((1 \leq j) \wedge (j \leq n))) \Rightarrow  \\ \quad \quad \quad \quad \quad \operatorname{nth}(j, \operatorname{nth}(i, p)) = g((k + i), j) \end{array}}}
\vnbprop{mat-tabulation-exists}{\vnbstmt{\begin{array}{@{}l@{}} \forall m, n, x, g. \\ \quad ((m \in \mathbb{N}) \wedge ((n \in \mathbb{N}) \wedge ((1 \leq m) \wedge \forall i \in \operatorname{interval}(1, m), j \in \operatorname{interval}(1, n).\; (g(i, j) \in x)))) \Rightarrow  \\ \quad \quad \exists p \in \operatorname{mat}(m, n, x). \forall i \in \operatorname{interval}(1, m), j \in \operatorname{interval}(1, n). \\ \quad \quad \quad \operatorname{entry}(p, i, j) = g(i, j) \end{array}}}
\vnbprop{matof-exists-image}{\vnbstmt{\begin{array}{@{}l@{}} \forall m, n, g. \\ \quad m \in \mathbb{N} \wedge \\ \quad n \in \mathbb{N} \wedge \\ \quad 1 \leq m \wedge \\ \quad \forall i \in \operatorname{interval}(1, m), j \in \operatorname{interval}(1, n).\; (g(i, j) \in \mathbf{Set}) \Rightarrow \\ \quad \quad \exists p \in \operatorname{mat}(m, n, \operatorname{image}(g, (\operatorname{interval}(1, m) \times \operatorname{interval}(1, n)))). \forall i \in \operatorname{interval}(1, m), j \in \operatorname{interval}(1, n). \\ \quad \quad \quad \operatorname{entry}(p, i, j) = g(i, j) \end{array}}}
\vnbprop{entry-of-matof-guarded}{\vnbstmt{\begin{array}{@{}l@{}} \forall m, n, g.; \forall u \in \operatorname{interval}(1, m), v \in \operatorname{interval}(1, n). \\ \quad m \in \mathbb{N} \wedge \\ \quad n \in \mathbb{N} \wedge \\ \quad 1 \leq m \wedge \\ \quad \forall i \in \operatorname{interval}(1, m), j \in \operatorname{interval}(1, n).\; (g(i, j) \in \mathbf{Set}) \Rightarrow \\ \quad \quad \operatorname{entry}(\operatorname{matof}(m, n, g), u, v) = g(u, v) \end{array}}}
\vnbprop{matof-exists}{\vnbstmt{\begin{array}{@{}l@{}} \forall m, n, g. \\ \quad ((m \in \mathbb{N}) \wedge ((n \in \mathbb{N}) \wedge \forall a \in \operatorname{interval}(1, m), b \in \operatorname{interval}(1, n).\; (g(a, b) \in \mathbf{Set}))) \Rightarrow  \\ \quad \quad \exists p \in \operatorname{mat}(m, n, \operatorname{image}(g, (\operatorname{interval}(1, m) \times \operatorname{interval}(1, n)))). \forall i \in \operatorname{interval}(1, m), j \in \operatorname{interval}(1, n). \\ \quad \quad \quad \operatorname{entry}(p, i, j) = g(i, j) \end{array}}}
\vnbprop{entry-of-matof}{\vnbstmt{\begin{array}{@{}l@{}} \forall m, n, g, i, j. \\ \quad m \in \mathbb{N} \wedge \\ \quad n \in \mathbb{N} \wedge \\ \quad \forall a \in \operatorname{interval}(1, m), b \in \operatorname{interval}(1, n).\; (g(a, b) \in \mathbf{Set}) \wedge \\ \quad i \in \operatorname{interval}(1, m) \wedge \\ \quad j \in \operatorname{interval}(1, n) \Rightarrow \\ \quad \quad \operatorname{entry}(\operatorname{matof}(m, n, g), i, j) = g(i, j) \end{array}}}

\vnbbookfile{theorem-library/matof-in-mat}{theorem-library/matof-in-mat}
\noindent{\small matof-in-mat.scm -- the tabulated matrix MATOF(m,n,g) is an m-by-n matrix over X when every value g(i,j) on the index box [1,m] x [1,n] lies in X.}\par
\vnbprop{tuples-mono}{\vnbstmt{\forall a, b.; \forall l \in \operatorname{tuples}(a).\; (\forall w \in a.\; (w \in b) \Rightarrow (l \in \operatorname{tuples}(b)))}}
\vnbprop{matrix-mono}{\vnbstmt{\forall a, b.; \forall m \in \operatorname{matrix}(a).\; (\forall w \in a.\; (w \in b) \Rightarrow (m \in \operatorname{matrix}(b)))}}
\vnbprop{mat-mono}{\vnbstmt{\begin{array}{@{}l@{}} \forall m, n, a, b.; \forall q \in \operatorname{mat}(m, n, a). \\ \quad \forall w \in a.\; (w \in b) \Rightarrow (q \in \operatorname{mat}(m, n, b)) \end{array}}}
\vnbprop{matof-in-mat}{\vnbstmt{\begin{array}{@{}l@{}} \forall m, n, x, g. \\ \quad ((m \in \mathbb{N}) \wedge ((n \in \mathbb{N}) \wedge \forall i \in \operatorname{interval}(1, m), j \in \operatorname{interval}(1, n).\; (g(i, j) \in x))) \Rightarrow  \\ \quad \quad \operatorname{matof}(m, n, g) \in \operatorname{mat}(m, n, x) \end{array}}}

\vnbbookfile{theorem-library/makeset-card-bound}{theorem-library/makeset-card-bound}
\noindent{\small makeset-card-bound.scm -- a tuple of length n has at most n distinct entries.}\par
\vnbprop{union-comm}{\vnbstmt{\forall a, b.\; ((a \cup b) = (b \cup a))}}
\vnbprop{union-singleton-absorb}{\vnbstmt{\forall s.; \forall y \in s.\; ((\{y\} \cup s) = s)}}
\vnbprop{card-union-singleton-bound}{\vnbstmt{\begin{array}{@{}l@{}} \forall s, y \in \mathbf{Set}. \\ \quad (|s| \in \mathbb{N}) \Rightarrow  \\ \quad \quad (|(\{y\} \cup s)| \in \mathbb{N}) \wedge (|(\{y\} \cup s)| \leq (|s| + 1)) \end{array}}}
\vnbprop{makeset-card-bound-strong}{\vnbstmt{\begin{array}{@{}l@{}} \forall n \in \mathbb{N}, a.; \forall l \in \operatorname{tuples}(a). \\ \quad (\operatorname{length}(l) = n) \Rightarrow  \\ \quad \quad (\operatorname{make\text{-}set}(l) \in \mathbf{Set}) \wedge ((|\operatorname{make\text{-}set}(l)| \in \mathbb{N}) \wedge (|\operatorname{make\text{-}set}(l)| \leq n)) \end{array}}}
\vnbprop{makeset-card-bound}{\vnbstmt{\begin{array}{@{}l@{}} \forall n \in \mathbb{N}, a.; \forall l \in \operatorname{tuples}(a). \\ \quad (\operatorname{length}(l) = n) \Rightarrow (|\operatorname{make\text{-}set}(l)| \leq n) \end{array}}}

\vnbbookfile{theorem-library/card-subset-nn}{theorem-library/card-subset-nn}
\noindent{\small card-subset-nn.scm -- a subset of a finite set is finite.}\par
\vnbprop{card-subset-nn}{\vnbstmt{\begin{array}{@{}l@{}} \forall x, s \in \mathbf{Set}. \\ \quad ((|x| \in \mathbb{N}) \wedge \forall z \in s.\; (z \in x)) \Rightarrow (|s| \in \mathbb{N}) \end{array}}}

\vnbbookfile{theorem-library/card-image-finite}{theorem-library/card-image-finite}
\noindent{\small card-image-finite.scm -- the image of a finite set is finite.}\par
\vnbprop{card-image-finite}{\vnbstmt{\forall i.; \forall x \in \mathbf{Set}.\; ((|x| \in \mathbb{N}) \Rightarrow (|\operatorname{image}(i, x)| \in \mathbb{N}))}}

\vnbbookfile{theorem-library/interval-card-in-nn}{theorem-library/interval-card-in-nn}
\noindent{\small interval-card-in-nn.scm -- an interval with a natural upper bound is finite.}\par
\vnbprop{interval-card-in-nn}{\vnbstmt{\forall a.; \forall b \in \mathbb{N}.\; (|\operatorname{interval}(a, b)| \in \mathbb{N})}}

\vnbbookfile{theorem-library/rake-finsum-laws2}{theorem-library/rake-finsum-laws2}
\noindent{\small A file shared across books; the entries that belong to this one.}\par
\noindent{\small theorem-library/rake-finsum-laws2.scm -- rake batch J, part 2: the three DISTRIBUTION laws of FINSUM (a ring element on either side, a module scalar) and the interval back-peel.  Part 1 is theorem-library/rake-finsum-laws.scm, which must load FIRST: this file cites its `enum-fam-value' and}\par
\vnbprop{interval-succ-insert}{\vnbstmt{\begin{array}{@{}l@{}} \forall n \in \mathbb{N}. \\ \quad \operatorname{interval}(1, (n + 1)) = (\operatorname{interval}(1, n) \cup \{(n + 1)\}) \end{array}}}

\vnbbookfile{theorem-library/rake-finsum-core}{theorem-library/rake-finsum-core}
\noindent{\small A file shared across books; the entries that belong to this one.}\par
\noindent{\small theorem-library/rake-finsum-core.scm -- the FINSUM CORE, proven. Rake batch M (2026-09-17).  Nine of the batch's ten leaves; the tenth (finsum-embed) is BLOCKED and its obstacle is stated at the end of this header.}\par
\vnbprop{cra-is-rag}{\vnbstmt{\forall g.\; (\operatorname{commutative\text{-}ring\text{-}additive\text{-}ag}(g) \simeq \operatorname{ring\text{-}additive\text{-}ag}(g))}}
\vnbprop{ord-segment-insert}{\vnbstmt{\forall n \in \mathbb{N}.\; (\mathrm{OS}((n + 1)) = (\mathrm{OS}(n) \cup \{n\}))}}

\vnbbookfile{theorem-library/nn-finite-subset-bounded}{theorem-library/nn-finite-subset-bounded}
\noindent{\small nn-finite-subset-bounded.scm -- a finite subset of NN is bounded.}\par
\vnbprop{nn-finite-subset-bounded}{\vnbstmt{\begin{array}{@{}l@{}} \forall t. \\ \quad ((t \subseteq \mathbb{N}) \wedge (|t| \in \mathbb{N})) \Rightarrow  \\ \quad \quad \exists n \in \mathbb{N}.\; \forall y \in \mathbb{N}.\; ((n \leq y) \Rightarrow \neg (y \in t)) \end{array}}}

\vnbbookfile{theorem-library/card-inequalities}{theorem-library/card-inequalities}
\noindent{\small card-inequalities.scm -- the two missing CARD inequalities. --------------------------------------------------------------------------- NOTE, 2026-09-20 (batch 9-B).  This file was written while cardinality was AXIOMATISED under the name CARD and the defined constant was its companion}\par
\vnbprop{complement-in-subset}{\vnbstmt{\forall a, b.\; (\operatorname{complement\text{-}in}(b, a) \subseteq b)}}
\vnbprop{union-complement-in}{\vnbstmt{\forall a, b.\; ((a \subseteq b) \Rightarrow (b = (a \cup \operatorname{complement\text{-}in}(b, a))))}}
\vnbprop{union-as-disjoint}{\vnbstmt{\forall a, b.\; ((a \cup b) = (a \cup \operatorname{complement\text{-}in}(b, a)))}}
\vnbprop{intersection-complement-in-empty}{\vnbstmt{\forall a, b.\; ((a \cap \operatorname{complement\text{-}in}(b, a)) = \emptyset)}}
\vnbprop{card-subset-mono}{\vnbstmt{\forall b \in \mathbf{Set}, a.\; (((|b| \in \mathbb{N}) \wedge (a \subseteq b)) \Rightarrow (|a| \leq |b|))}}
\vnbprop{card-union-bound}{\vnbstmt{\begin{array}{@{}l@{}} \forall a, b \in \mathbf{Set}. \\ \quad ((|a| \in \mathbb{N}) \wedge (|b| \in \mathbb{N})) \Rightarrow (|(a \cup b)| \leq (|a| + |b|)) \end{array}}}
\vnbprop{card-union-nn}{\vnbstmt{\begin{array}{@{}l@{}} \forall a, b \in \mathbf{Set}. \\ \quad ((|a| \in \mathbb{N}) \wedge (|b| \in \mathbb{N})) \Rightarrow (|(a \cup b)| \in \mathbb{N}) \end{array}}}

\vnbbookfile{theorem-library/rake-combinatorics}{theorem-library/rake-combinatorics}
\noindent{\small A file shared across books; the entries that belong to this one.}\par
\noindent{\small rake-combinatorics.scm -- BATCH P of the 2026-09-17 rake: NN combinatorics, the ring powers, the powerset insert-split, and the two CHOOSE boundary values.  Seventeen asserted supports, PROVEN; every statement is copied LITERALLY from its site.}\par
\vnbprop{succ-nn-ord}{\vnbstmt{\forall n \in \mathbb{N}.\; ((n + 1) = (n + 1))}}
\vnbprop{union-empty-left}{\vnbstmt{\forall a \in \mathbf{Set}.\; ((\emptyset \cup a) = a)}}
\vnbprop{ord-segment-trans}{\vnbstmt{\begin{array}{@{}l@{}} \forall m \in \mathbb{N}, k \in \mathrm{OS}(m), i \in \mathrm{OS}(k). \\ \quad i \in \mathrm{OS}(m) \end{array}}}
\vnbprop{power-zero-base}{\vnbstmt{\forall n \in \mathbb{N}.\; (\operatorname{power}(0, (n + 1)) = 0)}}
\vnbprop{power-insert-disjoint}{\vnbstmt{\begin{array}{@{}l@{}} \forall x, k \in \mathbf{Set}. \\ \quad \neg (k \in x) \Rightarrow  \\ \quad \quad (\operatorname{power}(x) \cap \operatorname{image}((\lambda s \in \operatorname{power}(x).\; (s \cup \{k\})), \operatorname{power}(x))) = \emptyset \end{array}}}
\vnbprop{choose-set-unfold}{\vnbstmt{\forall n, m.\; (\operatorname{choose\text{-}set}(n, m) \simeq \{a \in \operatorname{power}(\mathrm{OS}(n)) : (|a| = m)\})}}
\vnbprop{choose-0-succ}{\vnbstmt{\forall k \in \mathbb{N}.\; (\operatorname{choose}(0, (k + 1)) = 0)}}
\vnbprop{choose-n-0}{\vnbstmt{\forall n \in \mathbb{N}.\; (\operatorname{choose}(n, 0) = 1)}}
\vnbprop{power-insert-cover}{\vnbstmt{\begin{array}{@{}l@{}} \forall x, k \in \mathbf{Set}. \\ \quad \neg (k \in x) \Rightarrow  \\ \quad \quad \operatorname{power}((x \cup \{k\})) \\ \quad \quad =\; (\operatorname{power}(x) \cup \operatorname{image}((\lambda s \in \operatorname{power}(x).\; (s \cup \{k\})), \operatorname{power}(x))) \end{array}}}
\vnbprop{subset-of-empty-is-empty}{\vnbstmt{\forall w \in \mathbf{Set}.\; (\forall z \in w.\; (z \in \emptyset) \Rightarrow (w = \emptyset))}}
\vnbprop{power-of-empty}{\vnbstmt{\operatorname{power}(\emptyset) = \{\emptyset\}}}
\vnbprop{card-power-nn}{\vnbstmt{\forall x \in \mathbf{Set}.\; ((|x| \in \mathbb{N}) \Rightarrow (|\operatorname{power}(x)| \in \mathbb{N}))}}

\vnbbookfile{theorem-library/rake-combinatorics2}{theorem-library/rake-combinatorics2}
\noindent{\small rake-combinatorics2.scm -- BATCH U of the 2026-09-18 rake (batch 5, the off-bill supports): finite combinatorics.  Every statement is copied LITERALLY from its support site; the sites are named beside each proof. Six results, all `proven modulo 0'.}\par
\vnbprop{interval-card}{\vnbstmt{\forall n \in \mathbb{N}.\; (|\operatorname{interval}(1, n)| = n)}}
\vnbprop{choose-in-nn}{\vnbstmt{\forall n, m \in \mathbb{N}.\; (\operatorname{choose}(n, m) \in \mathbb{N})}}
\vnbprop{injection-from-empty}{\vnbstmt{\forall c \in \mathbf{Set}.\; (|\operatorname{injection}(\emptyset, c)| = 1)}}
\vnbprop{permutations-zero}{\vnbstmt{|\operatorname{permutations}(0)| = 1}}
\vnbprop{fibrewise-finite}{\vnbstmt{\begin{array}{@{}l@{}} \forall p.; \forall g, t \in \mathbf{Set}. \\ \quad ((|g| \in \mathbb{N}) \wedge (\forall y \in t.\; (p(y) \in g) \wedge \forall c \in g.\; (|\{x \in t : (p(x) = c)\}| \in \mathbb{N}))) \Rightarrow  \\ \quad \quad |t| \in \mathbb{N} \end{array}}}
\vnbprop{pigeonhole-infinite}{\vnbstmt{\begin{array}{@{}l@{}} \forall s, f \in \mathbf{Set}, p \in (s \to f). \\ \quad (\neg (|s| \in \mathbb{N}) \wedge (|f| \in \mathbb{N})) \Rightarrow  \\ \quad \quad \exists c \in f.\; \neg (|\{x \in s : (p(x) = c)\}| \in \mathbb{N}) \end{array}}}

\vnbbookfile{theorem-library/rake-choose-succ}{theorem-library/rake-choose-succ}
\noindent{\small rake-choose-succ.scm -- PASCAL'S RULE for the binomial coefficient, and the CHOOSE-SET lemmas it is assembled from.  Rake batch 5b, assignment 5b-G.}\par
\vnbprop{choose-card-unfold}{\vnbstmt{\forall n, m.\; (|\operatorname{choose\text{-}set}(n, m)| \simeq \operatorname{choose}(n, m))}}
\vnbprop{choose-set-mem-iff}{\vnbstmt{\begin{array}{@{}l@{}} \forall n, m, a. \\ \quad (a \in \operatorname{choose\text{-}set}(n, m)) \Leftrightarrow ((a \in \operatorname{power}(\mathrm{OS}(n))) \wedge (|a| = m)) \end{array}}}
\vnbprop{choose-set-avoids-point}{\vnbstmt{\forall n \in \mathbb{N}, m.; \forall a \in \operatorname{choose\text{-}set}(n, m).\; \neg (n \in a)}}
\vnbprop{choose-set-widen}{\vnbstmt{\forall n \in \mathbb{N}, m.; \forall a \in \operatorname{choose\text{-}set}(n, m).\; (a \in \operatorname{choose\text{-}set}((n + 1), m))}}
\vnbprop{choose-set-insert-in}{\vnbstmt{\begin{array}{@{}l@{}} \forall n, k \in \mathbb{N}, s \in \operatorname{choose\text{-}set}(n, k). \\ \quad (s \cup \{n\}) \in \operatorname{choose\text{-}set}((n + 1), (k + 1)) \end{array}}}
\vnbprop{choose-set-remove-restores}{\vnbstmt{\begin{array}{@{}l@{}} \forall n, a \in \mathbf{Set}. \\ \quad (n \in a) \Rightarrow ((\operatorname{difference}(a, \{n\}) \cup \{n\}) = a) \end{array}}}
\vnbprop{choose-set-remove-in}{\vnbstmt{\begin{array}{@{}l@{}} \forall n, k \in \mathbb{N}, a \in \operatorname{choose\text{-}set}((n + 1), (k + 1)). \\ \quad (n \in a) \Rightarrow (\operatorname{difference}(a, \{n\}) \in \operatorname{choose\text{-}set}(n, k)) \end{array}}}
\vnbprop{choose-set-avoid-in}{\vnbstmt{\begin{array}{@{}l@{}} \forall n \in \mathbb{N}, m.; \forall a \in \operatorname{choose\text{-}set}((n + 1), m). \\ \quad \neg (n \in a) \Rightarrow (a \in \operatorname{choose\text{-}set}(n, m)) \end{array}}}
\vnbprop{choose-set-insert-injective}{\vnbstmt{\begin{array}{@{}l@{}} \forall n \in \mathbb{N}, m.; \forall s, t \in \operatorname{choose\text{-}set}(n, m). \\ \quad ((s \cup \{n\}) = (t \cup \{n\})) \Rightarrow (s = t) \end{array}}}
\vnbprop{choose-set-split}{\vnbstmt{\begin{array}{@{}l@{}} \forall n, k \in \mathbb{N}. \\ \quad \operatorname{choose\text{-}set}((n + 1), (k + 1)) \\ \quad =\; (\operatorname{choose\text{-}set}(n, (k + 1)) \cup \operatorname{image}((\lambda s \in \operatorname{choose\text{-}set}(n, k).\; (s \cup \{n\})), \operatorname{choose\text{-}set}(n, k))) \end{array}}}
\vnbprop{choose-set-split-disjoint}{\vnbstmt{\begin{array}{@{}l@{}} \forall n, k \in \mathbb{N}. \\ \quad (\operatorname{choose\text{-}set}(n, (k + 1)) \cap \operatorname{image}((\lambda s \in \operatorname{choose\text{-}set}(n, k).\; (s \cup \{n\})), \operatorname{choose\text{-}set}(n, k))) \\ \quad =\; \emptyset \end{array}}}
\vnbprop{choose-succ}{\vnbstmt{\begin{array}{@{}l@{}} \forall n, k \in \mathbb{N}. \\ \quad \operatorname{choose}((n + 1), (k + 1)) = (\operatorname{choose}(n, k) + \operatorname{choose}(n, (k + 1))) \end{array}}}

\vnbbookfile{theorem-library/nn-infinite}{theorem-library/nn-infinite}
\noindent{\small nn-infinite.scm -- NN is infinite, and the two leaves that hang on it.}\par
\vnbprop{inf-subsets-unfold}{\vnbstmt{\forall a.\; (\mathrm{Inf}(a) \simeq \{s \in \operatorname{power}(a) : \neg (|s| \in \mathbb{N})\})}}
\vnbprop{nn-not-finite}{\vnbstmt{\neg (|\mathbb{N}| \in \mathbb{N})}}
\vnbprop{nn-in-inf-subsets}{\vnbstmt{\mathbb{N} \in \mathrm{Inf}(\mathbb{N})}}
\vnbprop{inf-subset-nn-unbounded}{\vnbstmt{\forall t \in \mathrm{Inf}(\mathbb{N}), u \in \mathbb{N}.\; \exists y \in \mathbb{N}.\; ((y \in t) \wedge (u < y))}}

\vnbbookfile{theorem-library/fin-subsets}{theorem-library/fin-subsets}
\noindent{\small fin-subsets.scm -- FIN-SUBSETS(a), the finite subsets of a, and the --------------------------------------------------------------------------- NOTE, 2026-09-20 (batch 9-B).  This file was written while cardinality was AXIOMATISED under the name CARD and the defined constant was its companion}\par
\par\medskip\noindent\textbf{Definition} (\texttt{fin-subsets}).\;$\operatorname{fin\text{-}subsets}(\operatorname{a\_}) :=$ $\{\operatorname{t\_} \in \operatorname{power}(\operatorname{a\_}) : (|\operatorname{t\_}| \in \mathbb{N})\}$
\vnbprop{fin-subsets-membership}{\vnbstmt{\begin{array}{@{}l@{}} \forall a, s. \\ \quad (s \in \operatorname{fin\text{-}subsets}(a)) \Leftrightarrow ((s \in \mathbf{Set}) \wedge (\forall z \in s.\; (z \in a) \wedge (|s| \in \mathbb{N}))) \end{array}}}
\vnbprop{fin-subsets-is-set}{\vnbstmt{\forall a \in \mathbf{Set}.\; (\operatorname{fin\text{-}subsets}(a) \in \mathbf{Set})}}
\vnbprop{fin-subsets-has-empty}{\vnbstmt{\forall a.\; (\emptyset \in \operatorname{fin\text{-}subsets}(a))}}
\vnbprop{fin-subsets-union-closed}{\vnbstmt{\begin{array}{@{}l@{}} \forall a, s, u. \\ \quad ((s \in \operatorname{fin\text{-}subsets}(a)) \wedge (u \in \operatorname{fin\text{-}subsets}(a))) \Rightarrow  \\ \quad \quad (s \cup u) \in \operatorname{fin\text{-}subsets}(a) \end{array}}}

\vnbbookfile{theorem-library/nn-pairing}{theorem-library/nn-pairing}
\noindent{\small nn-pairing.scm -- the Cantor pairing NN x NN -$>$ NN, and the arithmetic it needs.  Step one of the missing re-indexing mechanism.}\par
\par\medskip\noindent\textbf{Definition} (\texttt{trinum}).\; $(\operatorname{trinum}(0) \simeq 0)$\par
Suppose $\operatorname{n\_} \in \mathbb{N}$. Then $(\operatorname{trinum}(\operatorname{succ}(\operatorname{n\_})) \simeq +(\operatorname{trinum}(\operatorname{n\_}), \operatorname{succ}(\operatorname{n\_})))$
\par\medskip\noindent\textbf{Definition} (\texttt{nnpair}).\;$\operatorname{nnpair}(\operatorname{i\_}, \operatorname{j\_}) :=$ $\operatorname{trinum}((\operatorname{i\_} + \operatorname{j\_})) + \operatorname{j\_}$
\par\medskip\noindent\textbf{Definition} (\texttt{nnfst}).\;$\operatorname{nnfst}(\operatorname{n\_}) :=$
\begin{equation*}
(\iota\, \operatorname{a\_}.\;((\operatorname{a\_} \in \mathbb{N}) \wedge \exists \operatorname{b\_} \in \mathbb{N}.\; (\operatorname{nnpair}(\operatorname{a\_}, \operatorname{b\_}) = \operatorname{n\_})))
\end{equation*}

\par\medskip\noindent\textbf{Definition} (\texttt{nnsnd}).\;$\operatorname{nnsnd}(\operatorname{n\_}) :=$
\begin{equation*}
(\iota\, \operatorname{b\_}.\;((\operatorname{b\_} \in \mathbb{N}) \wedge \exists \operatorname{a\_} \in \mathbb{N}.\; (\operatorname{nnpair}(\operatorname{a\_}, \operatorname{b\_}) = \operatorname{n\_})))
\end{equation*}

\vnbprop{trinum-type}{\vnbstmt{\forall n \in \mathbb{N}.\; (\operatorname{trinum}(n) \in \mathbb{N})}}
\vnbprop{trinum-mono}{\vnbstmt{\forall b, a \in \mathbb{N}.\; ((a \leq b) \Rightarrow (\operatorname{trinum}(a) \leq \operatorname{trinum}(b)))}}
\vnbprop{nnpair-onto}{\vnbstmt{\forall n \in \mathbb{N}.\; \exists i, j \in \mathbb{N}.\; (\operatorname{nnpair}(i, j) = n)}}
\vnbprop{nnpair-diag-bound}{\vnbstmt{\begin{array}{@{}l@{}} \forall s, j. \\ \quad ((s \in \mathbb{N}) \wedge ((j \in \mathbb{N}) \wedge (j \leq s))) \Rightarrow (((\operatorname{trinum}(s) + j) + 1) \leq \operatorname{trinum}((s + 1))) \end{array}}}
\vnbprop{nnpair-cross}{\vnbstmt{\begin{array}{@{}l@{}} \forall s, t, j, k. \\ \quad ((s \in \mathbb{N}) \wedge ((t \in \mathbb{N}) \wedge ((j \in \mathbb{N}) \wedge ((k \in \mathbb{N}) \wedge ((j \leq s) \wedge ((s + 1) \leq t)))))) \Rightarrow  \\ \quad \quad ((\operatorname{trinum}(s) + j) + 1) \leq (\operatorname{trinum}(t) + k) \end{array}}}
\vnbprop{nnpair-inj}{\vnbstmt{\begin{array}{@{}l@{}} \forall i, j, \mathit{ii}, \mathit{jj}. \\ \quad ((i \in \mathbb{N}) \wedge ((j \in \mathbb{N}) \wedge ((\mathit{ii} \in \mathbb{N}) \wedge ((\mathit{jj} \in \mathbb{N}) \wedge (\operatorname{nnpair}(i, j) = \operatorname{nnpair}(\mathit{ii}, \mathit{jj})))))) \Rightarrow  \\ \quad \quad (i = \mathit{ii}) \wedge (j = \mathit{jj}) \end{array}}}
\vnbprop{nnpair-type}{\vnbstmt{\forall i, j.\; (((i \in \mathbb{N}) \wedge (j \in \mathbb{N})) \Rightarrow (\operatorname{nnpair}(i, j) \in \mathbb{N}))}}
\vnbprop{nnfst-nnpair}{\vnbstmt{\forall i, j.\; (((i \in \mathbb{N}) \wedge (j \in \mathbb{N})) \Rightarrow (\operatorname{nnfst}(\operatorname{nnpair}(i, j)) = i))}}
\vnbprop{nnsnd-nnpair}{\vnbstmt{\forall i, j.\; (((i \in \mathbb{N}) \wedge (j \in \mathbb{N})) \Rightarrow (\operatorname{nnsnd}(\operatorname{nnpair}(i, j)) = j))}}
\vnbprop{nnfst-type}{\vnbstmt{\forall n \in \mathbb{N}.\; (\operatorname{nnfst}(n) \in \mathbb{N})}}
\vnbprop{nnsnd-type}{\vnbstmt{\forall n \in \mathbb{N}.\; (\operatorname{nnsnd}(n) \in \mathbb{N})}}
\vnbprop{nn-flatten}{\vnbstmt{\begin{array}{@{}l@{}} \forall a.; \forall h \in (\mathbb{N} \to (\mathbb{N} \to a)). \exists e \in (\mathbb{N} \to a). \forall i, j. \\ \quad ((i \in \mathbb{N}) \wedge (j \in \mathbb{N})) \Rightarrow \exists m \in \mathbb{N}.\; (e(m) = h(i)(j)) \end{array}}}

\vnbbookfile{theorem-library/nn-nested-subset-chain-proof}{theorem-library/nn-nested-subset-chain-proof}
\noindent{\small nn-nested-subset-chain-proof.scm -- a descending NN-indexed family is a chain.}\par
\vnbprop{nn-nested-subset-chain-j}{\vnbstmt{\begin{array}{@{}l@{}} \forall m.; \forall j, k \in \mathbb{N}. \\ \quad (\forall k \in \mathbb{N}.\; (m(k) \in \mathbf{Set}) \wedge (\forall k \in \mathbb{N}.\; (m((k + 1)) \subseteq m(k)) \wedge (k \leq j))) \Rightarrow  \\ \quad \quad m(j) \subseteq m(k) \end{array}}}
\vnbprop{nn-nested-subset-chain}{\vnbstmt{\begin{array}{@{}l@{}} \forall t.; \forall k, j \in \mathbb{N}. \\ \quad (\forall k \in \mathbb{N}.\; (t(k) \in \mathbf{Set}) \wedge (\forall k \in \mathbb{N}.\; (t((k + 1)) \subseteq t(k)) \wedge (k \leq j))) \Rightarrow  \\ \quad \quad t(j) \subseteq t(k) \end{array}}}

\vnbbookfile{structure-library/order-zorn}{structure-library/order-zorn}
\noindent{\small MOVED 2026-09-20 (batch 12-A) from theorem-library/: this file is VOCABULARY -- definitions, notation and warranted supports, not one proof -- and every theorem stated with it had to load below it.  Its load.scm slot is unchanged. order-zorn.scm -- abstract partial orders and ZORN'S LEMMA (statement only).}\par
\par\medskip\noindent\textbf{Definition} (\texttt{is-partial-order}).\;For all $\mathit{grd}$ and $\mathit{porel}$, 
\begin{equation*}
\begin{array}{@{}l@{}} \operatorname{is\text{-}partial\text{-}order}(\mathit{grd}, \mathit{porel}) \\ \Leftrightarrow\; \mathit{grd} \in \mathbf{Set} \wedge \\ \mathit{porel} \in \mathbf{Set} \wedge \\ \forall x \in \mathit{grd}.\; (\langle x, x \rangle \in \mathit{porel}) \wedge \\ \forall x, y \in \mathit{grd}.\; (((\langle x, y \rangle \in \mathit{porel}) \wedge (\langle y, x \rangle \in \mathit{porel})) \Rightarrow (x = y)) \wedge \\ \forall x, y, z \in \mathit{grd}. \\ \quad ((\langle x, y \rangle \in \mathit{porel}) \wedge (\langle y, z \rangle \in \mathit{porel})) \Rightarrow (\langle x, z \rangle \in \mathit{porel}) \end{array}
\end{equation*}

\par\medskip\noindent\textbf{Definition} (\texttt{is-chain}).\;For all $\mathit{grd}$, $\mathit{porel}$ and $\mathit{ch}$, 
\begin{equation*}
\begin{array}{@{}l@{}} \operatorname{is\text{-}chain}(\mathit{grd}, \mathit{porel}, \mathit{ch}) \\ \Leftrightarrow\; (\mathit{ch} \subseteq \mathit{grd}) \wedge \forall x, y \in \mathit{ch}.\; ((\langle x, y \rangle \in \mathit{porel}) \vee (\langle y, x \rangle \in \mathit{porel})) \end{array}
\end{equation*}

\par\medskip\noindent\textbf{Definition} (\texttt{is-upper-bound}).\;For all $\mathit{grd}$, $\mathit{porel}$, $\mathit{ch}$ and $b$, 
\begin{equation*}
\operatorname{is\text{-}upper\text{-}bound}(\mathit{grd}, \mathit{porel}, \mathit{ch}, b) \Leftrightarrow ((b \in \mathit{grd}) \wedge \forall x \in \mathit{ch}.\; (\langle x, b \rangle \in \mathit{porel}))
\end{equation*}

\par\medskip\noindent\textbf{Definition} (\texttt{is-maximal}).\;For all $\mathit{grd}$, $\mathit{porel}$ and $\mathit{mx}$, 
\begin{equation*}
\begin{array}{@{}l@{}} \operatorname{is\text{-}maximal}(\mathit{grd}, \mathit{porel}, \mathit{mx}) \\ \Leftrightarrow\; (\mathit{mx} \in \mathit{grd}) \wedge \forall y \in \mathit{grd}.\; ((\langle \mathit{mx}, y \rangle \in \mathit{porel}) \Rightarrow (\langle y, \mathit{mx} \rangle \in \mathit{porel})) \end{array}
\end{equation*}


\vnbbookfile{theorem-library/injective-star}{theorem-library/injective-star}
\noindent{\small injective-star.scm -- the bridge from the set-function vocabulary to the class-function one.}\par
\vnbprop{injection-is-injective-star}{\vnbstmt{\forall x, y.; \forall f \in \operatorname{injection}(x, y).\; \operatorname{injective\text{-}star}(f)}}

\vnbbookfile{theorem-library/zorn-proof}{theorem-library/zorn-proof}
\noindent{\small zorn-proof.scm -- proving ZORN'S LEMMA from the well-ordering principle, via the GREEDY-KEPT-CHAIN construction (avoids the cardinality/Hartogs comparison lemma the base lacks).  Multi-session build; this file grows as lemmas land. Started 2026-07-22.}\par
\par\medskip\noindent\textbf{Definition} (\texttt{zkept}).\;For all $\varphi$, $\mathit{grd}$ and $\mathit{porel}$,  $\operatorname{zkept}(\varphi, \mathit{grd}, \mathit{porel}, 0) = \emptyset$\par
Let $\varphi$, $\mathit{grd}$ and $\mathit{porel}$ be arbitrary. Let $\alpha \in \mathrm{Ord}$. Then
\begin{equation*}
\begin{array}{@{}l@{}} \operatorname{zkept}(\varphi, \mathit{grd}, \mathit{porel}, (\alpha + 1)) \\ =\; (\operatorname{zkept}(\varphi, \mathit{grd}, \mathit{porel}, \alpha) \cup \operatorname{keep\text{-}set}(\varphi, \mathit{grd}, \mathit{porel}, \operatorname{zkept}(\varphi, \mathit{grd}, \mathit{porel}, \alpha), \alpha)) \end{array}
\end{equation*}
\par
Let $\varphi$, $\mathit{grd}$ and $\mathit{porel}$ be arbitrary. Let $\mathit{lam} \text{ be s.t. } \operatorname{limit\text{-}ord}(\mathit{lam})$. Then
\begin{equation*}
\operatorname{zkept}(\varphi, \mathit{grd}, \mathit{porel}, \mathit{lam}) = \operatorname{big\text{-}union}(\beta, \mathrm{OS}(\mathit{lam}), \operatorname{zkept}(\varphi, \mathit{grd}, \mathit{porel}, \beta))
\end{equation*}

\par\medskip\noindent\textbf{Definition} (\texttt{greedy-chain}).\;$\operatorname{greedy\text{-}chain}(\varphi, \mathit{grd}, \mathit{porel}) :=$ $\operatorname{zkept}(\varphi, \mathit{grd}, \mathit{porel}, |\mathit{grd}|)$
\par\medskip\noindent\textbf{Definition} (\texttt{keep-set}).\;$\operatorname{keep\text{-}set}(\varphi, \mathit{grd}, \mathit{porel}, \mathit{kset}, \alpha) :=$
\begin{equation*}
\{y \in \mathit{grd} : ((y = \varphi(\alpha)) \wedge \forall z \in \mathit{kset}.\; (\langle z, y \rangle \in \mathit{porel}))\}
\end{equation*}

\vnbdef{keep-set-membership}{\vnbstmt{\begin{array}{@{}l@{}} \forall i, d, l, t, a, y. \\ \quad (y \in \operatorname{keep\text{-}set}(i, d, l, t, a)) \Leftrightarrow ((y \in d) \wedge ((y = i(a)) \wedge \forall z \in t.\; (\langle z, y \rangle \in l))) \end{array}}}
\vnbprop{zorn-empty-is-chain}{\vnbstmt{\forall d \in \mathbf{Set}, l.\; \operatorname{is\text{-}chain}(d, l, \emptyset)}}
\vnbprop{zorn-chain-plus-keepset}{\vnbstmt{\begin{array}{@{}l@{}} \forall i, d, l, t, a. \\ \quad (\operatorname{is\text{-}partial\text{-}order}(d, l) \wedge \operatorname{is\text{-}chain}(d, l, t)) \Rightarrow  \\ \quad \quad \operatorname{is\text{-}chain}(d, l, (t \cup \operatorname{keep\text{-}set}(i, d, l, t, a))) \end{array}}}
\vnbprop{zkept-defined}{\vnbstmt{\forall a \in \mathrm{Ord}, i, d, l.\; (\operatorname{zkept}(i, d, l, a) = \operatorname{zkept}(i, d, l, a))}}
\vnbprop{zorn-zkept-monotone}{\vnbstmt{\begin{array}{@{}l@{}} \forall b \in \mathrm{Ord}, i, d, l.; \forall a \in \mathrm{Ord}. \\ \quad \operatorname{ord\text{-}le}(a, b) \Rightarrow (\operatorname{zkept}(i, d, l, a) \subseteq \operatorname{zkept}(i, d, l, b)) \end{array}}}
\vnbprop{zorn-zkept-limit-is-chain}{\vnbstmt{\begin{array}{@{}l@{}} \forall i, d, l, m. \\ \quad \operatorname{is\text{-}partial\text{-}order}(d, l) \wedge \\ \quad \operatorname{limit\text{-}ord}(m) \wedge \\ \quad \forall a \in \mathrm{OS}(m).\; \operatorname{is\text{-}chain}(d, l, \operatorname{zkept}(i, d, l, a)) \Rightarrow \\ \quad \quad \operatorname{is\text{-}chain}(d, l, \operatorname{big\text{-}union}(a_{2}, \mathrm{OS}(m), \operatorname{zkept}(i, d, l, a_{2}))) \end{array}}}
\vnbprop{zorn-zkept-is-chain}{\vnbstmt{\begin{array}{@{}l@{}} \forall a \in \mathrm{Ord}, i, d, l. \\ \quad \operatorname{is\text{-}partial\text{-}order}(d, l) \Rightarrow \operatorname{is\text{-}chain}(d, l, \operatorname{zkept}(i, d, l, a)) \end{array}}}
\vnbprop{greedy-chain-is-chain}{\vnbstmt{\begin{array}{@{}l@{}} \forall i, d, l. \\ \quad \operatorname{is\text{-}partial\text{-}order}(d, l) \Rightarrow \operatorname{is\text{-}chain}(d, l, \operatorname{greedy\text{-}chain}(i, d, l)) \end{array}}}

\vnbbookfile{theorem-library/zorn-route-two}{theorem-library/zorn-route-two}
\noindent{\small zorn-route-two.scm -- ZORN'S LEMMA, PROVED.}\par
\par\medskip\noindent\textbf{Definition} (\texttt{zup}).\;For all $\mathit{grd}$ and $\mathit{porel}$,  $\operatorname{zup}(\mathit{grd}, \mathit{porel}, 0) = \varepsilon(\mathit{grd})$\par
Let $\mathit{grd}$ and $\mathit{porel}$ be arbitrary. Let $\alpha \in \mathrm{Ord}$. Then
\begin{equation*}
\begin{array}{@{}l@{}} \operatorname{zup}(\mathit{grd}, \mathit{porel}, (\alpha + 1)) \\ =\; \varepsilon(\{\operatorname{y\_} \in \mathit{grd} : \operatorname{is\text{-}strictly\text{-}below}(\mathit{porel}, \operatorname{zup}(\mathit{grd}, \mathit{porel}, \alpha), \operatorname{y\_})\}) \end{array}
\end{equation*}
\par
Let $\mathit{grd}$ and $\mathit{porel}$ be arbitrary. Let $\mathit{lam} \text{ be s.t. } \operatorname{limit\text{-}ord}(\mathit{lam})$. Then
\begin{equation*}
\begin{array}{@{}l@{}} \operatorname{zup}(\mathit{grd}, \mathit{porel}, \mathit{lam}) \\ =\; \varepsilon(\{\operatorname{y\_} \in \mathit{grd} : \forall \operatorname{c\_}.\; (\operatorname{ord\text{-}lt}(\operatorname{c\_}, \mathit{lam}) \Rightarrow (\langle \operatorname{zup}(\mathit{grd}, \mathit{porel}, \operatorname{c\_}), \operatorname{y\_} \rangle \in \mathit{porel}))\}) \end{array}
\end{equation*}

\par\medskip\noindent\textbf{Definition} (\texttt{is-strictly-below}).\;For all $\mathit{porel}$, $x$ and $y$, 
\begin{equation*}
\operatorname{is\text{-}strictly\text{-}below}(\mathit{porel}, x, y) \Leftrightarrow ((\langle x, y \rangle \in \mathit{porel}) \wedge \neg (\langle y, x \rangle \in \mathit{porel}))
\end{equation*}

\vnbprop{strictly-below-irreflexive}{\vnbstmt{\forall l, x.\; \neg \operatorname{is\text{-}strictly\text{-}below}(l, x, x)}}
\vnbprop{strictly-below-transitive}{\vnbstmt{\begin{array}{@{}l@{}} \forall d, l, x, y, z. \\ \quad \operatorname{is\text{-}partial\text{-}order}(d, l) \wedge \\ \quad x \in d \wedge \\ \quad y \in d \wedge \\ \quad z \in d \wedge \\ \quad \operatorname{is\text{-}strictly\text{-}below}(l, x, y) \wedge \\ \quad \operatorname{is\text{-}strictly\text{-}below}(l, y, z) \Rightarrow \\ \quad \quad \operatorname{is\text{-}strictly\text{-}below}(l, x, z) \end{array}}}
\vnbprop{strictly-below-not-equal}{\vnbstmt{\forall l, x, y.\; ((x = y) \Rightarrow \neg \operatorname{is\text{-}strictly\text{-}below}(l, x, y))}}
\vnbprop{no-maximal-strictly-above}{\vnbstmt{\begin{array}{@{}l@{}} \forall d, l.; \forall x \in d. \\ \quad \neg \exists \mathit{mx}.\; \operatorname{is\text{-}maximal}(d, l, \mathit{mx}) \Rightarrow  \\ \quad \quad \exists y \in d.\; \operatorname{is\text{-}strictly\text{-}below}(l, x, y) \end{array}}}
\vnbprop{po-reflexive}{\vnbstmt{\forall d, l, x.\; ((\operatorname{is\text{-}partial\text{-}order}(d, l) \wedge (x \in d)) \Rightarrow (\langle x, x \rangle \in l))}}
\vnbprop{po-transitive}{\vnbstmt{\begin{array}{@{}l@{}} \forall d, l, x, y, z. \\ \quad (\operatorname{is\text{-}partial\text{-}order}(d, l) \wedge ((x \in d) \wedge ((y \in d) \wedge ((z \in d) \wedge ((\langle x, y \rangle \in l) \wedge (\langle y, z \rangle \in l)))))) \Rightarrow  \\ \quad \quad \langle x, z \rangle \in l \end{array}}}
\vnbprop{ord-lt-succ-cases}{\vnbstmt{\begin{array}{@{}l@{}} \forall a, b. \\ \quad ((a \in \mathrm{Ord}) \wedge ((b \in \mathrm{Ord}) \wedge \operatorname{ord\text{-}lt}(b, (a + 1)))) \Rightarrow (\operatorname{ord\text{-}lt}(b, a) \vee (b = a)) \end{array}}}
\vnbprop{ord-succ-lt-limit}{\vnbstmt{\begin{array}{@{}l@{}} \forall m, b. \\ \quad (\operatorname{limit\text{-}ord}(m) \wedge ((b \in \mathrm{Ord}) \wedge \operatorname{ord\text{-}lt}(b, m))) \Rightarrow \operatorname{ord\text{-}lt}((b + 1), m) \end{array}}}
\vnbprop{zup-tower}{\vnbstmt{\begin{array}{@{}l@{}} \forall a \in \mathrm{Ord}, d, l. \\ \quad \operatorname{is\text{-}partial\text{-}order}(d, l) \wedge \\ \quad \exists w.\; (w \in d) \wedge \\ \quad \forall h.\; (\operatorname{is\text{-}chain}(d, l, h) \Rightarrow \exists u.\; \operatorname{is\text{-}upper\text{-}bound}(d, l, h, u)) \wedge \\ \quad \forall x \in d.\; \exists y \in d.\; \operatorname{is\text{-}strictly\text{-}below}(l, x, y) \Rightarrow \\ \quad \quad \operatorname{zup}(d, l, a) \in d \wedge \\ \quad \quad \forall b. \\ \quad \quad \quad \operatorname{ord\text{-}lt}(b, a) \Rightarrow  \\ \quad \quad \quad \quad (\operatorname{zup}(d, l, b) \in d) \wedge \operatorname{is\text{-}strictly\text{-}below}(l, \operatorname{zup}(d, l, b), \operatorname{zup}(d, l, a)) \end{array}}}
\vnbprop{zup-injective}{\vnbstmt{\begin{array}{@{}l@{}} \forall d, l, a_{1}, a_{2}. \\ \quad \operatorname{is\text{-}partial\text{-}order}(d, l) \wedge \\ \quad \exists w.\; (w \in d) \wedge \\ \quad \forall h.\; (\operatorname{is\text{-}chain}(d, l, h) \Rightarrow \exists u.\; \operatorname{is\text{-}upper\text{-}bound}(d, l, h, u)) \wedge \\ \quad \forall x \in d.\; \exists y \in d.\; \operatorname{is\text{-}strictly\text{-}below}(l, x, y) \wedge \\ \quad a_{1} \in \mathrm{Ord} \wedge \\ \quad a_{2} \in \mathrm{Ord} \wedge \\ \quad \operatorname{zup}(d, l, a_{1}) = \operatorname{zup}(d, l, a_{2}) \Rightarrow \\ \quad \quad a_{1} = a_{2} \end{array}}}
\vnbprop{zorn-lemma}{\vnbstmt{\begin{array}{@{}l@{}} \forall d, l. \\ \quad \operatorname{is\text{-}partial\text{-}order}(d, l) \wedge \\ \quad \exists w.\; (w \in d) \wedge \\ \quad \forall h.\; (\operatorname{is\text{-}chain}(d, l, h) \Rightarrow \exists u.\; \operatorname{is\text{-}upper\text{-}bound}(d, l, h, u)) \Rightarrow \\ \quad \quad \exists x.\; \operatorname{is\text{-}maximal}(d, l, x) \end{array}}}

\vnbbookfile{theorem-library/nn-least-element}{theorem-library/nn-least-element}
\noindent{\small theorem-library/nn-least-element.scm =================================================================== The well-ordering of NN, PROVEN from the well-ordering of ORD.}\par
\vnbprop{nn-least-element}{\vnbstmt{\begin{array}{@{}l@{}} \forall t. \\ \quad ((t \subseteq \mathbb{N}) \wedge \exists n.\; (n \in t)) \Rightarrow  \\ \quad \quad \exists m \in t.\; \forall k \in t.\; (m \leq k) \end{array}}}

\vnbbookfile{theorem-library/pair-self-cartesian3}{theorem-library/pair-self-cartesian3}
\noindent{\small pair-self-cartesian3.scm -- three generic set lemmas: a set is the one member of its own singleton PAIR(a, a) (pair-self-member, pair-self-eq), and the ternary product of three sets is a set (cartesian3-set).}\par
\vnbprop{pair-self-member}{\vnbstmt{\forall a \in \mathbf{Set}.\; (a \in \{a\})}}
\vnbprop{pair-self-eq}{\vnbstmt{\forall a, x.\; (((a \in \mathbf{Set}) \wedge (x \in \{a\})) \Rightarrow (x = a))}}
\vnbprop{cartesian3-set}{\vnbstmt{\begin{array}{@{}l@{}} \forall u, v, w. \\ \quad ((u \in \mathbf{Set}) \wedge ((v \in \mathbf{Set}) \wedge (w \in \mathbf{Set}))) \Rightarrow ((u \times v) \in \mathbf{Set}) \end{array}}}

\vnbbookfile{theorem-library/interval-widen}{theorem-library/interval-widen}
\noindent{\small interval-widen.scm -- an interval membership survives widening the top end.}\par
\vnbprop{interval-widen}{\vnbstmt{\begin{array}{@{}l@{}} \forall b, c \in \mathbb{N}, a.; \forall i \in \operatorname{interval}(a, b). \\ \quad (b \leq c) \Rightarrow (i \in \operatorname{interval}(a, c)) \end{array}}}

\vnbbookfile{theorem-library/rake-finsum-typing-nn}{theorem-library/rake-finsum-typing-nn}
\noindent{\small rake-finsum-typing-nn.scm -- two NN typing leaves of the rake, PROVEN.}\par
\vnbprop{nn-minus-in-nn}{\vnbstmt{\forall k, l \in \mathbb{N}.\; (\operatorname{nn\text{-}minus}(k, l) \in \mathbb{N})}}
\vnbprop{falling-in-nn-ind}{\vnbstmt{\begin{array}{@{}l@{}} \forall m, n \in \mathbb{N}. \\ \quad (\operatorname{falling}(n, m) \in \mathbb{N}) \wedge (((n - m) \in \mathbb{N}) \vee (\operatorname{falling}(n, m) = 0)) \end{array}}}
\vnbprop{falling-in-nn}{\vnbstmt{\forall n, m \in \mathbb{N}.\; (\operatorname{falling}(n, m) \in \mathbb{N})}}

\vnbbookfile{theorem-library/rake-asm-2}{theorem-library/rake-asm-2}
\noindent{\small A file shared across books; the entries that belong to this one.}\par
\noindent{\small rake-asm-2.scm -- RAKE-ASM agent RA-2 (2026-10-12): elementary laws that stood as ASSUMPTIONS (add-axiom! / support), each proven here under its EXACT name and statement (copied character for character from its definition site).}\par
\vnbprop{falling-zero-succ}{\vnbstmt{\forall k \in \mathbb{N}.\; (\operatorname{falling}(0, (k + 1)) = 0)}}
\vnbprop{falling-succ-succ}{\vnbstmt{\forall k, n \in \mathbb{N}.\; (\operatorname{falling}((n + 1), (k + 1)) = ((n + 1) \cdot \operatorname{falling}(n, k)))}}
\vnbprop{choose-times-factorial-all}{\vnbstmt{\forall n, m \in \mathbb{N}.\; ((\operatorname{choose}(n, m) \cdot \operatorname{factorial}(m)) = \operatorname{falling}(n, m))}}
\vnbprop{choose-times-factorial}{\vnbstmt{\begin{array}{@{}l@{}} \forall n, m \in \mathbb{N}. \\ \quad (m \leq n) \Rightarrow ((\operatorname{choose}(n, m) \cdot \operatorname{factorial}(m)) = \operatorname{falling}(n, m)) \end{array}}}

\vnbbookfile{theorem-library/rake-setoid2}{theorem-library/rake-setoid2}
\noindent{\small A file shared across books; the entries that belong to this one.}\par
\noindent{\small rake-setoid2.scm -- BATCH O of the 2026-09-17 rake: the SETOID quotient and the METRIC COMPLETION.  Twelve of the thirteen asserted leaves of the batch are PROVEN here; the thirteenth is a STATEMENT DEFECT and is recorded below. Every statement is its support's statement UNCHANGED (copied from the}\par
\vnbprop{rko2-dist-seq-at}{\vnbstmt{\forall m, f, g.; \forall n \in \mathbb{N}.\; (\operatorname{dist\text{-}seq}(m, f, g)(n) \simeq \operatorname{dist}(m)(f(n), g(n)))}}

\vnbbookfile{theorem-library/rake-asm-1}{theorem-library/rake-asm-1}
\noindent{\small A file shared across books; the entries that belong to this one.}\par
\noindent{\small rake-asm-1.scm -- RAKE-ASM agent RA-1 (2026-10-12): elementary assumptions PROVEN.}\par
\vnbprop{injection-in-fun}{\vnbstmt{\forall x, y.; \forall i \in \operatorname{injection}(x, y).\; (i \in (x \to y))}}
\vnbprop{injection-injective}{\vnbstmt{\forall x, y.; \forall i \in \operatorname{injection}(x, y), a, b \in x.\; ((i(a) = i(b)) \Rightarrow (a = b))}}
\vnbprop{injection-set-iff}{\vnbstmt{\forall x, y.\; (((x \in \mathbf{Set}) \wedge (y \in \mathbf{Set})) \Rightarrow (\operatorname{injection}(x, y) \in \mathbf{Set}))}}
\vnbprop{bijection-set-iff}{\vnbstmt{\forall x, y.\; (((x \in \mathbf{Set}) \wedge (y \in \mathbf{Set})) \Rightarrow (\mathrm{Bij}(x, y) \in \mathbf{Set}))}}

\vnbbookfile{theorem-library/rake-asm-b5}{theorem-library/rake-asm-b5}
\noindent{\small A file shared across books; the entries that belong to this one.}\par
\noindent{\small rake-asm-b5.scm -- RAKE-ASM-B, agent RB-5 (2026-10-12).  Helper prefix `rb5-'.}\par
\vnbprop{rb5-nary-plus-3-cc}{\vnbstmt{\forall x, y, z \in \mathbb{C}.\; ((x + y + z) = ((x + y) + z))}}
\vnbprop{rb5-nary-plus-4-cc}{\vnbstmt{\begin{array}{@{}l@{}} \forall w, x, y, z \in \mathbb{C}. \\ \quad (w + x + y + z) = (((w + x) + y) + z) \end{array}}}
\vnbprop{rb5-nary-plus-5-cc}{\vnbstmt{\begin{array}{@{}l@{}} \forall v, w, x, y, z \in \mathbb{C}. \\ \quad (v + w + x + y + z) = ((((v + w) + x) + y) + z) \end{array}}}
\vnbprop{rb5-nary-times-3-cc}{\vnbstmt{\forall x, y, z \in \mathbb{C}.\; ((x \cdot y \cdot z) = ((x \cdot y) \cdot z))}}
\vnbprop{rb5-nary-times-4-cc}{\vnbstmt{\begin{array}{@{}l@{}} \forall w, x, y, z \in \mathbb{C}. \\ \quad (w \cdot x \cdot y \cdot z) = (((w \cdot x) \cdot y) \cdot z) \end{array}}}
\vnbprop{rb5-nary-times-5-cc}{\vnbstmt{\begin{array}{@{}l@{}} \forall v, w, x, y, z \in \mathbb{C}. \\ \quad (v \cdot w \cdot x \cdot y \cdot z) = ((((v \cdot w) \cdot x) \cdot y) \cdot z) \end{array}}}
\vnbprop{rb5-nary-plus-2-list-cc}{\vnbstmt{\forall x, y \in \mathbb{C}.\; ((x + y) = \operatorname{reduce}(\mathit{binplus}, \operatorname{fam\text{-}of\text{-}list}(\langle x, y \rangle), 2))}}
\vnbprop{rb5-nary-plus-3-list-cc}{\vnbstmt{\begin{array}{@{}l@{}} \forall x, y, z \in \mathbb{C}. \\ \quad (x + y + z) = \operatorname{reduce}(\mathit{binplus}, \operatorname{fam\text{-}of\text{-}list}(\langle x, y, z \rangle), 3) \end{array}}}
\vnbprop{rb5-nary-plus-4-list-cc}{\vnbstmt{\begin{array}{@{}l@{}} \forall w, x, y, z \in \mathbb{C}. \\ \quad (w + x + y + z) = \operatorname{reduce}(\mathit{binplus}, \operatorname{fam\text{-}of\text{-}list}(\langle w, x, y, z \rangle), 4) \end{array}}}
\vnbprop{rb5-nary-plus-5-list-cc}{\vnbstmt{\begin{array}{@{}l@{}} \forall v, w, x, y, z \in \mathbb{C}. \\ \quad (v + w + x + y + z) = \operatorname{reduce}(\mathit{binplus}, \operatorname{fam\text{-}of\text{-}list}(\langle v, w, x, y, z \rangle), 5) \end{array}}}
\vnbprop{rb5-nary-times-2-list-cc}{\vnbstmt{\forall x, y \in \mathbb{C}.\; ((x \cdot y) = \operatorname{reduce}(\mathit{bintimes}, \operatorname{fam\text{-}of\text{-}list}(\langle x, y \rangle), 2))}}
\vnbprop{rb5-nary-times-3-list-cc}{\vnbstmt{\begin{array}{@{}l@{}} \forall x, y, z \in \mathbb{C}. \\ \quad (x \cdot y \cdot z) = \operatorname{reduce}(\mathit{bintimes}, \operatorname{fam\text{-}of\text{-}list}(\langle x, y, z \rangle), 3) \end{array}}}
\vnbprop{rb5-nary-times-4-list-cc}{\vnbstmt{\begin{array}{@{}l@{}} \forall w, x, y, z \in \mathbb{C}. \\ \quad (w \cdot x \cdot y \cdot z) = \operatorname{reduce}(\mathit{bintimes}, \operatorname{fam\text{-}of\text{-}list}(\langle w, x, y, z \rangle), 4) \end{array}}}
\vnbprop{rb5-nary-times-5-list-cc}{\vnbstmt{\begin{array}{@{}l@{}} \forall v, w, x, y, z \in \mathbb{C}. \\ \quad (v \cdot w \cdot x \cdot y \cdot z) = \operatorname{reduce}(\mathit{bintimes}, \operatorname{fam\text{-}of\text{-}list}(\langle v, w, x, y, z \rangle), 5) \end{array}}}
\vnbprop{rb5-fibre-count}{\vnbstmt{\begin{array}{@{}l@{}} \forall p.; \forall m \in \mathbb{N}, g, t \in \mathbf{Set}. \\ \quad ((|g| \in \mathbb{N}) \wedge (\forall y \in t.\; (p(y) \in g) \wedge \forall c \in g.\; (|\{x \in t : (p(x) = c)\}| = m))) \Rightarrow  \\ \quad \quad |t| = (m \cdot |g|) \end{array}}}
\vnbprop{rb5-restrict-injection}{\vnbstmt{\begin{array}{@{}l@{}} \forall a, d, c, f. \\ \quad ((a \in \mathbf{Set}) \wedge ((f \in \operatorname{injection}(d, c)) \wedge (a \subseteq d))) \Rightarrow  \\ \quad \quad (\lambda x \in a.\; f(x)) \in \operatorname{injection}(a, c) \end{array}}}
\vnbprop{rb5-ext-at-b}{\vnbstmt{\begin{array}{@{}l@{}} \forall a, c, b, g, v. \\ \quad ((b \in \mathbf{Set}) \wedge (v \in c)) \Rightarrow  \\ \quad \quad (\lambda z \in (a \cup \{b\}).\; \left(\text{if }(z = b)\text{ then }v\text{ else }g(z)\right))(b) = v \end{array}}}
\vnbprop{rb5-ext-off-b}{\vnbstmt{\begin{array}{@{}l@{}} \forall a, c, b, g, v.; \forall w \in a. \\ \quad ((b \in \mathbf{Set}) \wedge (\neg (b \in a) \wedge (g \in (a \to c)))) \Rightarrow  \\ \quad \quad (\lambda z \in (a \cup \{b\}).\; \left(\text{if }(z = b)\text{ then }v\text{ else }g(z)\right))(w) = g(w) \end{array}}}
\vnbprop{rb5-ext-injection}{\vnbstmt{\begin{array}{@{}l@{}} \forall a, c, b, g, v. \\ \quad ((a \in \mathbf{Set}) \wedge ((b \in \mathbf{Set}) \wedge (\neg (b \in a) \wedge ((g \in \operatorname{injection}(a, c)) \wedge ((v \in c) \wedge \neg (v \in \operatorname{image}(g, a))))))) \Rightarrow  \\ \quad \quad (\lambda z \in (a \cup \{b\}).\; \left(\text{if }(z = b)\text{ then }v\text{ else }g(z)\right)) \in \operatorname{injection}((a \cup \{b\}), c) \end{array}}}
\vnbprop{rb5-ext-restrict}{\vnbstmt{\begin{array}{@{}l@{}} \forall a, c, b, g, v. \\ \quad ((a \in \mathbf{Set}) \wedge ((b \in \mathbf{Set}) \wedge (\neg (b \in a) \wedge ((g \in \operatorname{injection}(a, c)) \wedge ((v \in c) \wedge \neg (v \in \operatorname{image}(g, a))))))) \Rightarrow  \\ \quad \quad (\lambda x \in a.\; (\lambda z \in (a \cup \{b\}).\; \left(\text{if }(z = b)\text{ then }v\text{ else }g(z)\right))(x)) = g \end{array}}}
\vnbprop{rb5-free-value}{\vnbstmt{\begin{array}{@{}l@{}} \forall a, c, b, g, f. \\ \quad a \in \mathbf{Set} \wedge \\ \quad b \in \mathbf{Set} \wedge \\ \quad \neg (b \in a) \wedge \\ \quad f \in \operatorname{injection}((a \cup \{b\}), c) \wedge \\ \quad (\lambda x \in a.\; f(x)) = g \Rightarrow \\ \quad \quad (f(b) \in c) \wedge \neg (f(b) \in \operatorname{image}(g, a)) \end{array}}}
\vnbprop{rb5-ext-recover}{\vnbstmt{\begin{array}{@{}l@{}} \forall a, c, b, g, f. \\ \quad a \in \mathbf{Set} \wedge \\ \quad b \in \mathbf{Set} \wedge \\ \quad \neg (b \in a) \wedge \\ \quad g \in \operatorname{injection}(a, c) \wedge \\ \quad f \in \operatorname{injection}((a \cup \{b\}), c) \wedge \\ \quad (\lambda x \in a.\; f(x)) = g \Rightarrow \\ \quad \quad f = (\lambda z \in (a \cup \{b\}).\; \left(\text{if }(z = b)\text{ then }f(b)\text{ else }g(z)\right)) \end{array}}}
\vnbprop{rb5-complement-card}{\vnbstmt{\begin{array}{@{}l@{}} \forall a, c, g, m. \\ \quad a \in \mathbf{Set} \wedge \\ \quad c \in \mathbf{Set} \wedge \\ \quad |a| \in \mathbb{N} \wedge \\ \quad m \in \mathbb{N} \wedge \\ \quad |c| = (|a| + m) \wedge \\ \quad g \in \operatorname{injection}(a, c) \Rightarrow \\ \quad \quad |\operatorname{difference}(c, \operatorname{image}(g, a))| = m \end{array}}}
\vnbprop{rb5-fibre-image}{\vnbstmt{\begin{array}{@{}l@{}} \forall a, c, b, g. \\ \quad ((a \in \mathbf{Set}) \wedge ((c \in \mathbf{Set}) \wedge ((b \in \mathbf{Set}) \wedge (\neg (b \in a) \wedge (g \in \operatorname{injection}(a, c)))))) \Rightarrow  \\ \quad \quad \operatorname{image}((\lambda \mathit{rbc} \in \operatorname{difference}(c, \operatorname{image}(g, a)).\; \\ \quad \quad \quad \quad (\lambda z \in (a \cup \{b\}).\; \left(\text{if }(z = b)\text{ then }\mathit{rbc}\text{ else }g(z)\right))), \\ \quad \quad \quad \operatorname{difference}(c, \operatorname{image}(g, a))) \\ \quad \quad =\; \{x \in \operatorname{injection}((a \cup \{b\}), c) : ((\lambda f \in \operatorname{injection}((a \cup \{b\}), c).\; (\lambda \mathit{rsx} \in a.\; f(\mathit{rsx})))(x) = g)\} \end{array}}}
\vnbprop{rb5-fibre-card}{\vnbstmt{\begin{array}{@{}l@{}} \forall a, c, b, m, g. \\ \quad a \in \mathbf{Set} \wedge \\ \quad c \in \mathbf{Set} \wedge \\ \quad b \in \mathbf{Set} \wedge \\ \quad \neg (b \in a) \wedge \\ \quad |a| \in \mathbb{N} \wedge \\ \quad m \in \mathbb{N} \wedge \\ \quad |c| = (|a| + m) \wedge \\ \quad g \in \operatorname{injection}(a, c) \Rightarrow \\ \quad \quad |\{x \in \operatorname{injection}((a \cup \{b\}), c) : ((\lambda f \in \operatorname{injection}((a \cup \{b\}), c).\; (\lambda \mathit{rsx} \in a.\; f(\mathit{rsx})))(x) = g)\}| \\ \quad \quad =\; m \end{array}}}
\vnbprop{rb5-injection-finite}{\vnbstmt{\begin{array}{@{}l@{}} \forall a, c. \\ \quad ((a \in \mathbf{Set}) \wedge ((|a| \in \mathbb{N}) \wedge ((c \in \mathbf{Set}) \wedge (|c| \in \mathbb{N})))) \Rightarrow  \\ \quad \quad |\operatorname{injection}(a, c)| \in \mathbb{N} \end{array}}}
\vnbprop{injection-extension-recurrence}{\vnbstmt{\begin{array}{@{}l@{}} \forall a, c, b, m. \\ \quad \neg (b \in a) \wedge \\ \quad a \in \mathbf{Set} \wedge \\ \quad c \in \mathbf{Set} \wedge \\ \quad |a| \in \mathbb{N} \wedge \\ \quad m \in \mathbb{N} \wedge \\ \quad |c| = (|a| + m) \wedge \\ \quad b \in \mathbf{Set} \Rightarrow \\ \quad \quad |\operatorname{injection}((a \cup \{b\}), c)| = (m \cdot |\operatorname{injection}(a, c)|) \end{array}}}
\vnbprop{rb5-count-falling-aux}{\vnbstmt{\begin{array}{@{}l@{}} \forall m, n \in \mathbb{N}. \\ \quad (m \leq n) \Rightarrow (|\operatorname{injection}(\mathrm{OS}(m), \mathrm{OS}(n))| = \operatorname{falling}(n, m)) \end{array}}}
\vnbprop{injection-count-falling}{\vnbstmt{\begin{array}{@{}l@{}} \forall n, m \in \mathbb{N}. \\ \quad (m \leq n) \Rightarrow (|\operatorname{injection}(\mathrm{OS}(m), \mathrm{OS}(n))| = \operatorname{falling}(n, m)) \end{array}}}
\vnbprop{permutation-recurrence}{\vnbstmt{\forall n \in \mathbb{N}.\; (|\operatorname{permutations}((n + 1))| = ((n + 1) \cdot |\operatorname{permutations}(n)|))}}

\vnbbookfile{structure-library/functoriality}{structure-library/functoriality}
\noindent{\small A file shared across books; the entries that belong to this one.}\par
\noindent{\small functoriality.scm -- every view-as is a FUNCTOR, and here is the proof.}\par
\vnbprop{abelian-group-as-monoid-functorial}{\vnbstmt{\begin{array}{@{}l@{}} \forall a, b, f. \\ \quad \operatorname{is\text{-}hom\text{-}abelian\text{-}group}(a, b, f) \Rightarrow  \\ \quad \quad \operatorname{is\text{-}hom\text{-}monoid}(\operatorname{abelian\text{-}group\text{-}as\text{-}monoid}(a), \operatorname{abelian\text{-}group\text{-}as\text{-}monoid}(b), f) \end{array}}}
\vnbprop{commutative-ring-additive-ag-functorial}{\vnbstmt{\begin{array}{@{}l@{}} \forall a, b, f. \\ \quad \operatorname{is\text{-}hom\text{-}commutative\text{-}ring}(a, b, f) \Rightarrow  \\ \quad \quad \operatorname{is\text{-}hom\text{-}abelian\text{-}group}(\operatorname{commutative\text{-}ring\text{-}additive\text{-}ag}(a), \operatorname{commutative\text{-}ring\text{-}additive\text{-}ag}(b), f) \end{array}}}
\vnbprop{commutative-ring-multiplicative-cm-functorial}{\vnbstmt{\begin{array}{@{}l@{}} \forall a, b, f. \\ \quad \operatorname{is\text{-}hom\text{-}commutative\text{-}ring}(a, b, f) \Rightarrow  \\ \quad \quad \operatorname{is\text{-}hom\text{-}comm\text{-}monoid}(\operatorname{commutative\text{-}ring\text{-}multiplicative\text{-}cm}(a), \operatorname{commutative\text{-}ring\text{-}multiplicative\text{-}cm}(b), f) \end{array}}}
\vnbprop{complex-inner-product-space-as-module-functorial}{\vnbstmt{\begin{array}{@{}l@{}} \forall a, b, f. \\ \quad \operatorname{is\text{-}hom\text{-}complex\text{-}inner\text{-}product\text{-}space}(a, b, f) \Rightarrow  \\ \quad \quad \operatorname{is\text{-}hom\text{-}module}(\operatorname{complex\text{-}inner\text{-}product\text{-}space\text{-}as\text{-}module}(a), \operatorname{complex\text{-}inner\text{-}product\text{-}space\text{-}as\text{-}module}(b), f) \end{array}}}
\vnbprop{field-additive-ag-functorial}{\vnbstmt{\begin{array}{@{}l@{}} \forall a, b, f. \\ \quad \operatorname{is\text{-}hom\text{-}field}(a, b, f) \Rightarrow  \\ \quad \quad \operatorname{is\text{-}hom\text{-}abelian\text{-}group}(\operatorname{field\text{-}additive\text{-}ag}(a), \operatorname{field\text{-}additive\text{-}ag}(b), f) \end{array}}}
\vnbprop{module-vector-ag-functorial}{\vnbstmt{\begin{array}{@{}l@{}} \forall a, b, f. \\ \quad \operatorname{is\text{-}hom\text{-}module}(a, b, f) \Rightarrow  \\ \quad \quad \operatorname{is\text{-}hom\text{-}abelian\text{-}group}(\operatorname{module\text{-}vector\text{-}ag}(a), \operatorname{module\text{-}vector\text{-}ag}(b), f) \end{array}}}
\vnbprop{normed-ag-as-abelian-group-functorial}{\vnbstmt{\begin{array}{@{}l@{}} \forall a, b, f. \\ \quad \operatorname{is\text{-}hom\text{-}normed\text{-}ag}(a, b, f) \Rightarrow  \\ \quad \quad \operatorname{is\text{-}hom\text{-}abelian\text{-}group}(\operatorname{normed\text{-}ag\text{-}as\text{-}abelian\text{-}group}(a), \operatorname{normed\text{-}ag\text{-}as\text{-}abelian\text{-}group}(b), f) \end{array}}}
\vnbprop{normed-field-additive-ag-functorial}{\vnbstmt{\begin{array}{@{}l@{}} \forall a, b, f. \\ \quad \operatorname{is\text{-}hom\text{-}normed\text{-}field}(a, b, f) \Rightarrow  \\ \quad \quad \operatorname{is\text{-}hom\text{-}abelian\text{-}group}(\operatorname{normed\text{-}field\text{-}additive\text{-}ag}(a), \operatorname{normed\text{-}field\text{-}additive\text{-}ag}(b), f) \end{array}}}
\vnbprop{normed-vector-space-as-module-functorial}{\vnbstmt{\begin{array}{@{}l@{}} \forall a, b, f. \\ \quad \operatorname{is\text{-}hom\text{-}normed\text{-}vector\text{-}space}(a, b, f) \Rightarrow  \\ \quad \quad \operatorname{is\text{-}hom\text{-}module}(\operatorname{normed\text{-}vector\text{-}space\text{-}as\text{-}module}(a), \operatorname{normed\text{-}vector\text{-}space\text{-}as\text{-}module}(b), f) \end{array}}}
\vnbprop{normed-vector-space-as-normed-ag-functorial}{\vnbstmt{\begin{array}{@{}l@{}} \forall a, b, f. \\ \quad \operatorname{is\text{-}hom\text{-}normed\text{-}vector\text{-}space}(a, b, f) \Rightarrow  \\ \quad \quad \operatorname{is\text{-}hom\text{-}normed\text{-}ag}(\operatorname{normed\text{-}vector\text{-}space\text{-}as\text{-}normed\text{-}ag}(a), \operatorname{normed\text{-}vector\text{-}space\text{-}as\text{-}normed\text{-}ag}(b), f) \end{array}}}
\vnbprop{ring-additive-ag-functorial}{\vnbstmt{\begin{array}{@{}l@{}} \forall a, b, f. \\ \quad \operatorname{is\text{-}hom\text{-}ring}(a, b, f) \Rightarrow  \\ \quad \quad \operatorname{is\text{-}hom\text{-}abelian\text{-}group}(\operatorname{ring\text{-}additive\text{-}ag}(a), \operatorname{ring\text{-}additive\text{-}ag}(b), f) \end{array}}}
\vnbprop{ring-multiplicative-monoid-functorial}{\vnbstmt{\begin{array}{@{}l@{}} \forall a, b, f. \\ \quad \operatorname{is\text{-}hom\text{-}ring}(a, b, f) \Rightarrow  \\ \quad \quad \operatorname{is\text{-}hom\text{-}monoid}(\operatorname{ring\text{-}multiplicative\text{-}monoid}(a), \operatorname{ring\text{-}multiplicative\text{-}monoid}(b), f) \end{array}}}

\part{Algebra}
\noindent Semigroups, monoids, groups, rings, integral domains and fields; ideals, Euclidean rings and Bezout; polynomials and monoid algebras; setoids, quotients and ringoids; homomorphisms, categories and views.
\par\noindent\textit{Builds on:} Foundations.
\par\noindent{\small 87 files.}

\vnbbookfile{structure-library/semigroup}{structure-library/semigroup}
\noindent{\small semigroup.scm -- SEMIGROUP structure}\par
\vnbstructdef{semigroup}{A semigroup $s$ is a 2-tuple whose components are listed below.\par
\begin{enumerate}[nosep]
\item $\operatorname{carr}(s) = \operatorname{nth}(1, s)$, a carrier (\texttt{carriers}): $\operatorname{carr}(s) \in \mathbf{Set}$.
\item $\operatorname{opr}(s) = \operatorname{nth}(2, s)$, an operation (\texttt{op}): $\operatorname{opr}(s) \in ((\operatorname{carr}(s) \times \operatorname{carr}(s)) \to \operatorname{carr}(s))$.
\end{enumerate}
The laws (the \texttt{property} and \texttt{law} clauses):
\begin{itemize}[nosep]
\item $\operatorname{is\text{-}associative}(\operatorname{opr}(s), \operatorname{carr}(s))$
\end{itemize}
The characteristic law (the axiom \texttt{is-semigroup}): $\operatorname{is\text{-}semigroup}(s)$ holds if and only if $\operatorname{length}(s) = 2$, every component has the typing listed with it, and every law holds.
}
\par\medskip\noindent\textbf{Definition} (\texttt{hom-semigroup}).\;$\operatorname{hom\text{-}semigroup}(a, b) :=$
\begin{equation*}
\{\operatorname{homf\_} \in (\operatorname{carr}(a) \to \operatorname{carr}(b)) : \operatorname{is\text{-}hom\text{-}semigroup}(a, b, \operatorname{homf\_})\}
\end{equation*}

\vnbdef{semigroup-class}{\vnbstmt{\forall s.\; ((s \in \mathit{semigroup}) \Leftrightarrow \operatorname{is\text{-}semigroup}(s))}}
\vnbdef{is-hom-semigroup-def}{\vnbstmt{\begin{array}{@{}l@{}} \forall a, b, f. \\ \quad \operatorname{is\text{-}hom\text{-}semigroup}(a, b, f) \\ \quad \Leftrightarrow\; \operatorname{is\text{-}semigroup}(a) \wedge \\ \quad \operatorname{is\text{-}semigroup}(b) \wedge \\ \quad f \in (\operatorname{carr}(a) \to \operatorname{carr}(b)) \wedge \\ \quad \forall x_{1}, x_{2} \in \operatorname{carr}(a). \\ \quad \quad f(\operatorname{opr}(a)(x_{1}, x_{2})) = \operatorname{opr}(b)(f(x_{1}), f(x_{2})) \end{array}}}

\vnbbookfile{structure-library/monoid}{structure-library/monoid}
\noindent{\small RETIRED 2026-09-17 (proven): comm-monoid-opr-comm (was an unwarranted axiom, no citers) -- theorem-library/rake-finsum-core.scm RETIRED 2026-09-17 (proven): monoid-identity-in, monoid-carrier-closed-opr -- were UNWARRANTED axioms here; theorem-library/rake-algebra2.scm, projections of is-monoid-def, modulo 0.}\par
\vnbstructdef{monoid}{A monoid $s$ is a 3-tuple whose components are listed below.\par
\begin{enumerate}[nosep]
\item $\operatorname{carr}(s) = \operatorname{nth}(1, s)$, a carrier (\texttt{carriers}): $\operatorname{carr}(s) \in \mathbf{Set}$.
\item $\operatorname{opr}(s) = \operatorname{nth}(2, s)$, an operation (\texttt{op}): $\operatorname{opr}(s) \in ((\operatorname{carr}(s) \times \operatorname{carr}(s)) \to \operatorname{carr}(s))$.
\item $\operatorname{iden}(s) = \operatorname{nth}(3, s)$, a constant (\texttt{constant}): $\operatorname{iden}(s) \in \operatorname{carr}(s)$.
\end{enumerate}
The laws (the \texttt{property} and \texttt{law} clauses):
\begin{itemize}[nosep]
\item $\operatorname{is\text{-}associative}(\operatorname{opr}(s), \operatorname{carr}(s))$
\item $\operatorname{is\text{-}identity}(\operatorname{opr}(s), \operatorname{iden}(s), \operatorname{carr}(s))$
\end{itemize}
The characteristic law (the axiom \texttt{is-monoid}): $\operatorname{is\text{-}monoid}(s)$ holds if and only if $\operatorname{length}(s) = 3$, every component has the typing listed with it, and every law holds.
}
\vnbstructdef{comm-monoid}{A comm monoid $s$ is a 3-tuple whose components are listed below.\par
\begin{enumerate}[nosep]
\item $\operatorname{carr}(s) = \operatorname{nth}(1, s)$, a carrier (\texttt{carriers}): $\operatorname{carr}(s) \in \mathbf{Set}$.
\item $\operatorname{opr}(s) = \operatorname{nth}(2, s)$, an operation (\texttt{op}): $\operatorname{opr}(s) \in ((\operatorname{carr}(s) \times \operatorname{carr}(s)) \to \operatorname{carr}(s))$.
\item $\operatorname{iden}(s) = \operatorname{nth}(3, s)$, a constant (\texttt{constant}): $\operatorname{iden}(s) \in \operatorname{carr}(s)$.
\end{enumerate}
The laws (the \texttt{property} and \texttt{law} clauses):
\begin{itemize}[nosep]
\item $\operatorname{is\text{-}associative}(\operatorname{opr}(s), \operatorname{carr}(s))$
\item $\operatorname{is\text{-}identity}(\operatorname{opr}(s), \operatorname{iden}(s), \operatorname{carr}(s))$
\item $\operatorname{is\text{-}commutative}(\operatorname{opr}(s), \operatorname{carr}(s))$
\end{itemize}
The characteristic law (the axiom \texttt{is-comm-monoid}): $\operatorname{is\text{-}comm\text{-}monoid}(s)$ holds if and only if $\operatorname{length}(s) = 3$, every component has the typing listed with it, and every law holds.
}
\par\medskip\noindent\textbf{Definition} (\texttt{hom-comm-monoid}).\;$\operatorname{hom\text{-}comm\text{-}monoid}(a, b) :=$
\begin{equation*}
\{\operatorname{homf\_} \in (\operatorname{carr}(a) \to \operatorname{carr}(b)) : \operatorname{is\text{-}hom\text{-}comm\text{-}monoid}(a, b, \operatorname{homf\_})\}
\end{equation*}

\par\medskip\noindent\textbf{Definition} (\texttt{hom-monoid}).\;$\operatorname{hom\text{-}monoid}(a, b) :=$ $\{\operatorname{homf\_} \in (\operatorname{carr}(a) \to \operatorname{carr}(b)) : \operatorname{is\text{-}hom\text{-}monoid}(a, b, \operatorname{homf\_})\}$
\vnbdef{monoid-class}{\vnbstmt{\forall s.\; ((s \in \mathit{monoid}) \Leftrightarrow \operatorname{is\text{-}monoid}(s))}}
\vnbdef{is-hom-monoid-def}{\vnbstmt{\begin{array}{@{}l@{}} \forall a, b, f. \\ \quad \operatorname{is\text{-}hom\text{-}monoid}(a, b, f) \\ \quad \Leftrightarrow\; \operatorname{is\text{-}monoid}(a) \wedge \\ \quad \operatorname{is\text{-}monoid}(b) \wedge \\ \quad f \in (\operatorname{carr}(a) \to \operatorname{carr}(b)) \wedge \\ \quad \forall x_{1}, x_{2} \in \operatorname{carr}(a). \\ \quad \quad f(\operatorname{opr}(a)(x_{1}, x_{2})) = \operatorname{opr}(b)(f(x_{1}), f(x_{2})) \wedge \\ \quad f(\operatorname{iden}(a)) = \operatorname{iden}(b) \end{array}}}
\vnbdef{is-hom-comm-monoid-def}{\vnbstmt{\begin{array}{@{}l@{}} \forall a, b, f. \\ \quad \operatorname{is\text{-}hom\text{-}comm\text{-}monoid}(a, b, f) \\ \quad \Leftrightarrow\; \operatorname{is\text{-}comm\text{-}monoid}(a) \wedge \\ \quad \operatorname{is\text{-}comm\text{-}monoid}(b) \wedge \\ \quad f \in (\operatorname{carr}(a) \to \operatorname{carr}(b)) \wedge \\ \quad \forall x_{1}, x_{2} \in \operatorname{carr}(a). \\ \quad \quad f(\operatorname{opr}(a)(x_{1}, x_{2})) = \operatorname{opr}(b)(f(x_{1}), f(x_{2})) \wedge \\ \quad f(\operatorname{iden}(a)) = \operatorname{iden}(b) \end{array}}}
\vnbdef{comm-monoid-class}{\vnbstmt{\forall s.\; ((s \in \operatorname{comm\text{-}monoid}) \Leftrightarrow \operatorname{is\text{-}comm\text{-}monoid}(s))}}

\vnbbookfile{structure-library/group}{structure-library/group}
\noindent{\small group.scm -- GROUP structure}\par
\vnbstructdef{group}{A group $s$ is a 4-tuple whose components are listed below.\par
\begin{enumerate}[nosep]
\item $\operatorname{carr}(s) = \operatorname{nth}(1, s)$, a carrier (\texttt{carriers}): $\operatorname{carr}(s) \in \mathbf{Set}$.
\item $\operatorname{opr}(s) = \operatorname{nth}(2, s)$, an operation (\texttt{op}): $\operatorname{opr}(s) \in ((\operatorname{carr}(s) \times \operatorname{carr}(s)) \to \operatorname{carr}(s))$.
\item $\operatorname{iden}(s) = \operatorname{nth}(3, s)$, a constant (\texttt{constant}): $\operatorname{iden}(s) \in \operatorname{carr}(s)$.
\item $\operatorname{inv}(s) = \operatorname{nth}(4, s)$, an operation (\texttt{op}): $\operatorname{inv}(s) \in (\operatorname{carr}(s) \to \operatorname{carr}(s))$.
\end{enumerate}
The laws (the \texttt{property} and \texttt{law} clauses):
\begin{itemize}[nosep]
\item $\operatorname{is\text{-}associative}(\operatorname{opr}(s), \operatorname{carr}(s))$
\item $\operatorname{is\text{-}identity}(\operatorname{opr}(s), \operatorname{iden}(s), \operatorname{carr}(s))$
\item $\operatorname{has\text{-}inverses}(\operatorname{opr}(s), \operatorname{iden}(s), \operatorname{inv}(s), \operatorname{carr}(s))$
\end{itemize}
The characteristic law (the axiom \texttt{is-group}): $\operatorname{is\text{-}group}(s)$ holds if and only if $\operatorname{length}(s) = 4$, every component has the typing listed with it, and every law holds.
}
\par\medskip\noindent\textbf{Definition} (\texttt{hom-group}).\;$\operatorname{hom\text{-}group}(a, b) :=$ $\{\operatorname{homf\_} \in (\operatorname{carr}(a) \to \operatorname{carr}(b)) : \operatorname{is\text{-}hom\text{-}group}(a, b, \operatorname{homf\_})\}$
\vnbdef{is-hom-group-def}{\vnbstmt{\begin{array}{@{}l@{}} \forall a, b, f. \\ \quad \operatorname{is\text{-}hom\text{-}group}(a, b, f) \\ \quad \Leftrightarrow\; \operatorname{is\text{-}group}(a) \wedge \\ \quad \operatorname{is\text{-}group}(b) \wedge \\ \quad f \in (\operatorname{carr}(a) \to \operatorname{carr}(b)) \wedge \\ \quad \forall x_{1}, x_{2} \in \operatorname{carr}(a). \\ \quad \quad f(\operatorname{opr}(a)(x_{1}, x_{2})) = \operatorname{opr}(b)(f(x_{1}), f(x_{2})) \wedge \\ \quad f(\operatorname{iden}(a)) = \operatorname{iden}(b) \wedge \\ \quad \forall x_{1} \in \operatorname{carr}(a).\; (f(\operatorname{inv}(a)(x_{1})) = \operatorname{inv}(b)(f(x_{1}))) \end{array}}}
\vnbdef{group-class}{\vnbstmt{\forall s.\; ((s \in \mathit{group}) \Leftrightarrow \operatorname{is\text{-}group}(s))}}

\vnbbookfile{structure-library/abelian-group}{structure-library/abelian-group}
\noindent{\small RETIRED 2026-09-17 (proven): abelian-group-assoc, abelian-group-right-id, abelian-group-inverse-unique -- theorem-library/rake-algebra2.scm (which also rebuilds the three MODULE-VECTOR-AG companions that proofs cite by name). abelian-group.scm -- ABELIAN-GROUP structure}\par
\vnbstructdef{abelian-group}{An abelian group $s$ is a 4-tuple whose components are listed below.\par
\begin{enumerate}[nosep]
\item $\operatorname{carr}(s) = \operatorname{nth}(1, s)$, a carrier (\texttt{carriers}): $\operatorname{carr}(s) \in \mathbf{Set}$.
\item $\operatorname{opr}(s) = \operatorname{nth}(2, s)$, an operation (\texttt{op}): $\operatorname{opr}(s) \in ((\operatorname{carr}(s) \times \operatorname{carr}(s)) \to \operatorname{carr}(s))$.
\item $\operatorname{iden}(s) = \operatorname{nth}(3, s)$, a constant (\texttt{constant}): $\operatorname{iden}(s) \in \operatorname{carr}(s)$.
\item $\operatorname{inv}(s) = \operatorname{nth}(4, s)$, an operation (\texttt{op}): $\operatorname{inv}(s) \in (\operatorname{carr}(s) \to \operatorname{carr}(s))$.
\end{enumerate}
The laws (the \texttt{property} and \texttt{law} clauses):
\begin{itemize}[nosep]
\item $\operatorname{is\text{-}associative}(\operatorname{opr}(s), \operatorname{carr}(s))$
\item $\operatorname{is\text{-}identity}(\operatorname{opr}(s), \operatorname{iden}(s), \operatorname{carr}(s))$
\item $\operatorname{has\text{-}inverses}(\operatorname{opr}(s), \operatorname{iden}(s), \operatorname{inv}(s), \operatorname{carr}(s))$
\item $\operatorname{is\text{-}commutative}(\operatorname{opr}(s), \operatorname{carr}(s))$
\end{itemize}
The characteristic law (the axiom \texttt{is-abelian-group}): $\operatorname{is\text{-}abelian\text{-}group}(s)$ holds if and only if $\operatorname{length}(s) = 4$, every component has the typing listed with it, and every law holds.
}
\par\medskip\noindent\textbf{Definition} (\texttt{hom-abelian-group}).\;$\operatorname{hom\text{-}abelian\text{-}group}(a, b) :=$
\begin{equation*}
\{\operatorname{homf\_} \in (\operatorname{carr}(a) \to \operatorname{carr}(b)) : \operatorname{is\text{-}hom\text{-}abelian\text{-}group}(a, b, \operatorname{homf\_})\}
\end{equation*}

\vnbdef{abelian-group-class}{\vnbstmt{\forall s.\; ((s \in \operatorname{abelian\text{-}group}) \Leftrightarrow \operatorname{is\text{-}abelian\text{-}group}(s))}}
\vnbdef{is-hom-abelian-group-def}{\vnbstmt{\begin{array}{@{}l@{}} \forall a, b, f. \\ \quad \operatorname{is\text{-}hom\text{-}abelian\text{-}group}(a, b, f) \\ \quad \Leftrightarrow\; \operatorname{is\text{-}abelian\text{-}group}(a) \wedge \\ \quad \operatorname{is\text{-}abelian\text{-}group}(b) \wedge \\ \quad f \in (\operatorname{carr}(a) \to \operatorname{carr}(b)) \wedge \\ \quad \forall x_{1}, x_{2} \in \operatorname{carr}(a). \\ \quad \quad f(\operatorname{opr}(a)(x_{1}, x_{2})) = \operatorname{opr}(b)(f(x_{1}), f(x_{2})) \wedge \\ \quad f(\operatorname{iden}(a)) = \operatorname{iden}(b) \wedge \\ \quad \forall x_{1} \in \operatorname{carr}(a).\; (f(\operatorname{inv}(a)(x_{1})) = \operatorname{inv}(b)(f(x_{1}))) \end{array}}}

\vnbbookfile{structure-library/ring}{structure-library/ring}
\noindent{\small RETIRED 2026-09-17 (proven): ring-neg-mul-left -- theorem-library/rake-algebra2.scm RETIRED 2026-09-17 (proven): ring-carr-in-set -- theorem-library/rake-algebra.scm ring.scm -- RING, RING-PROD, ZERO-RING}\par
\vnbstructdef{ring}{A ring $s$ is a 6-tuple whose components are listed below.\par
\begin{enumerate}[nosep]
\item $\operatorname{carr}(s) = \operatorname{nth}(1, s)$, a carrier (\texttt{carriers}): $\operatorname{carr}(s) \in \mathbf{Set}$.
\item $\operatorname{add}(s) = \operatorname{nth}(2, s)$, an operation (\texttt{op}): $\operatorname{add}(s) \in ((\operatorname{carr}(s) \times \operatorname{carr}(s)) \to \operatorname{carr}(s))$.
\item $\operatorname{mul}(s) = \operatorname{nth}(3, s)$, an operation (\texttt{op}): $\operatorname{mul}(s) \in ((\operatorname{carr}(s) \times \operatorname{carr}(s)) \to \operatorname{carr}(s))$.
\item $\operatorname{neg}(s) = \operatorname{nth}(4, s)$, an operation (\texttt{op}): $\operatorname{neg}(s) \in (\operatorname{carr}(s) \to \operatorname{carr}(s))$.
\item $\operatorname{zero}(s) = \operatorname{nth}(5, s)$, a constant (\texttt{constant}): $\operatorname{zero}(s) \in \operatorname{carr}(s)$.
\item $\operatorname{one}(s) = \operatorname{nth}(6, s)$, a constant (\texttt{constant}): $\operatorname{one}(s) \in \operatorname{carr}(s)$.
\end{enumerate}
The laws (the \texttt{property} and \texttt{law} clauses):
\begin{itemize}[nosep]
\item $\operatorname{is\text{-}associative}(\operatorname{add}(s), \operatorname{carr}(s))$
\item $\operatorname{is\text{-}commutative}(\operatorname{add}(s), \operatorname{carr}(s))$
\item $\operatorname{is\text{-}identity}(\operatorname{add}(s), \operatorname{zero}(s), \operatorname{carr}(s))$
\item $\operatorname{has\text{-}inverses}(\operatorname{add}(s), \operatorname{zero}(s), \operatorname{neg}(s), \operatorname{carr}(s))$
\item $\operatorname{is\text{-}associative}(\operatorname{mul}(s), \operatorname{carr}(s))$
\item $\operatorname{is\text{-}identity}(\operatorname{mul}(s), \operatorname{one}(s), \operatorname{carr}(s))$
\item $\operatorname{is\text{-}distributive}(\operatorname{add}(s), \operatorname{mul}(s), \operatorname{carr}(s))$
\end{itemize}
The characteristic law (the axiom \texttt{is-ring}): $\operatorname{is\text{-}ring}(s)$ holds if and only if $\operatorname{length}(s) = 6$, every component has the typing listed with it, and every law holds.
}
\par\medskip\noindent\textbf{Definition} (\texttt{ring-prod}).\;$\operatorname{ring\text{-}prod}(x, y) :=$
\begin{equation*}
\langle (\operatorname{carr}(x) \times \operatorname{carr}(y)), (\lambda \langle p, q \rangle \in ((\operatorname{carr}(x) \times \operatorname{carr}(y)) \times (\operatorname{carr}(x) \times \operatorname{carr}(y))).\; \langle \operatorname{add}(x)(\operatorname{nth}(1, p), \operatorname{nth}(1, q)), \operatorname{add}(y)(\operatorname{nth}(2, p), \operatorname{nth}(2, q)) \rangle), (\lambda \langle p, q \rangle \in ((\operatorname{carr}(x) \times \operatorname{carr}(y)) \times (\operatorname{carr}(x) \times \operatorname{carr}(y))).\; \langle \operatorname{mul}(x)(\operatorname{nth}(1, p), \operatorname{nth}(1, q)), \operatorname{mul}(y)(\operatorname{nth}(2, p), \operatorname{nth}(2, q)) \rangle), (\lambda p \in (\operatorname{carr}(x) \times \operatorname{carr}(y)).\; \langle \operatorname{neg}(x)(\operatorname{nth}(1, p)), \operatorname{neg}(y)(\operatorname{nth}(2, p)) \rangle), \langle \operatorname{zero}(x), \operatorname{zero}(y) \rangle, \langle \operatorname{one}(x), \operatorname{one}(y) \rangle \rangle
\end{equation*}

\par\medskip\noindent\textbf{Definition} (\texttt{zero-ring}).\;
\begin{equation*}
\begin{array}{@{}l@{}} \operatorname{zero\text{-}ring} \\ =\; \langle \operatorname{make\text{-}set}(\langle 0 \rangle), (\lambda \langle p, q \rangle \in (\operatorname{make\text{-}set}(\langle 0 \rangle) \times \operatorname{make\text{-}set}(\langle 0 \rangle)).\; 0), (\lambda \langle p, q \rangle \in (\operatorname{make\text{-}set}(\langle 0 \rangle) \times \operatorname{make\text{-}set}(\langle 0 \rangle)).\; 0), (\lambda p \in \operatorname{make\text{-}set}(\langle 0 \rangle).\; 0), 0, 0 \rangle \end{array}
\end{equation*}

\par\medskip\noindent\textbf{Definition} (\texttt{hom-ring}).\;$\operatorname{hom\text{-}ring}(a, b) :=$ $\{\operatorname{homf\_} \in (\operatorname{carr}(a) \to \operatorname{carr}(b)) : \operatorname{is\text{-}hom\text{-}ring}(a, b, \operatorname{homf\_})\}$
\vnbdef{ring-add-assoc}{\vnbstmt{\begin{array}{@{}l@{}} \forall s.; \forall a, b, c \in \operatorname{carr}(s). \\ \quad \operatorname{is\text{-}ring}(s) \Rightarrow (\operatorname{add}(s)(\operatorname{add}(s)(a, b), c) = \operatorname{add}(s)(a, \operatorname{add}(s)(b, c))) \end{array}}}
\vnbdef{ring-add-comm}{\vnbstmt{\begin{array}{@{}l@{}} \forall s.; \forall a, b \in \operatorname{carr}(s). \\ \quad \operatorname{is\text{-}ring}(s) \Rightarrow (\operatorname{add}(s)(a, b) = \operatorname{add}(s)(b, a)) \end{array}}}
\vnbdef{ring-add-left-id}{\vnbstmt{\forall s.; \forall a \in \operatorname{carr}(s).\; (\operatorname{is\text{-}ring}(s) \Rightarrow (\operatorname{add}(s)(\operatorname{zero}(s), a) = a))}}
\vnbdef{ring-add-left-inv}{\vnbstmt{\forall s.; \forall a \in \operatorname{carr}(s).\; (\operatorname{is\text{-}ring}(s) \Rightarrow (\operatorname{add}(s)(\operatorname{neg}(s)(a), a) = \operatorname{zero}(s)))}}
\vnbdef{ring-mul-assoc}{\vnbstmt{\begin{array}{@{}l@{}} \forall s.; \forall a, b, c \in \operatorname{carr}(s). \\ \quad \operatorname{is\text{-}ring}(s) \Rightarrow (\operatorname{mul}(s)(\operatorname{mul}(s)(a, b), c) = \operatorname{mul}(s)(a, \operatorname{mul}(s)(b, c))) \end{array}}}
\vnbdef{ring-mul-left-id}{\vnbstmt{\forall s.; \forall a \in \operatorname{carr}(s).\; (\operatorname{is\text{-}ring}(s) \Rightarrow (\operatorname{mul}(s)(\operatorname{one}(s), a) = a))}}
\vnbdef{ring-mul-right-id}{\vnbstmt{\forall s.; \forall a \in \operatorname{carr}(s).\; (\operatorname{is\text{-}ring}(s) \Rightarrow (\operatorname{mul}(s)(a, \operatorname{one}(s)) = a))}}
\vnbdef{ring-left-dist}{\vnbstmt{\begin{array}{@{}l@{}} \forall s.; \forall a, b, c \in \operatorname{carr}(s). \\ \quad \operatorname{is\text{-}ring}(s) \Rightarrow  \\ \quad \quad \operatorname{mul}(s)(a, \operatorname{add}(s)(b, c)) = \operatorname{add}(s)(\operatorname{mul}(s)(a, b), \operatorname{mul}(s)(a, c)) \end{array}}}
\vnbdef{ring-right-dist}{\vnbstmt{\begin{array}{@{}l@{}} \forall s.; \forall a, b, c \in \operatorname{carr}(s). \\ \quad \operatorname{is\text{-}ring}(s) \Rightarrow  \\ \quad \quad \operatorname{mul}(s)(\operatorname{add}(s)(a, b), c) = \operatorname{add}(s)(\operatorname{mul}(s)(a, c), \operatorname{mul}(s)(b, c)) \end{array}}}
\vnbdef{ring-zero-in}{\vnbstmt{\forall r.\; (\operatorname{is\text{-}ring}(r) \Rightarrow (\operatorname{zero}(r) \in \operatorname{carr}(r)))}}
\vnbdef{ring-carrier-closed-add}{\vnbstmt{\begin{array}{@{}l@{}} \forall r, a, b. \\ \quad (\operatorname{is\text{-}ring}(r) \wedge ((a \in \operatorname{carr}(r)) \wedge (b \in \operatorname{carr}(r)))) \Rightarrow (\operatorname{add}(r)(a, b) \in \operatorname{carr}(r)) \end{array}}}
\vnbdef{is-hom-ring-def}{\vnbstmt{\begin{array}{@{}l@{}} \forall a, b, f. \\ \quad \operatorname{is\text{-}hom\text{-}ring}(a, b, f) \\ \quad \Leftrightarrow\; \operatorname{is\text{-}ring}(a) \wedge \\ \quad \operatorname{is\text{-}ring}(b) \wedge \\ \quad f \in (\operatorname{carr}(a) \to \operatorname{carr}(b)) \wedge \\ \quad \forall x_{1}, x_{2} \in \operatorname{carr}(a). \\ \quad \quad f(\operatorname{add}(a)(x_{1}, x_{2})) = \operatorname{add}(b)(f(x_{1}), f(x_{2})) \wedge \\ \quad \forall x_{1}, x_{2} \in \operatorname{carr}(a). \\ \quad \quad f(\operatorname{mul}(a)(x_{1}, x_{2})) = \operatorname{mul}(b)(f(x_{1}), f(x_{2})) \wedge \\ \quad \forall x_{1} \in \operatorname{carr}(a).\; (f(\operatorname{neg}(a)(x_{1})) = \operatorname{neg}(b)(f(x_{1}))) \wedge \\ \quad f(\operatorname{zero}(a)) = \operatorname{zero}(b) \wedge \\ \quad f(\operatorname{one}(a)) = \operatorname{one}(b) \end{array}}}
\vnbdef{ring-class}{\vnbstmt{\forall s.\; ((s \in \mathit{ring}) \Leftrightarrow \operatorname{is\text{-}ring}(s))}}

\vnbbookfile{structure-library/setoid}{structure-library/setoid}
\noindent{\small RETIRED 2026-09-17 (proven, rake batch O): class-self, class-subset-carrier, class-eq-iff, class-disjoint, class-in-quotient, quotient-rep, descend-computes, quotient-universal, descend2-computes -- theorem-library/rake-setoid2.scm, all modulo 0.  This file now holds no support. RETIRED 2026-09-17 (proven): ...}\par
\vnbstructdef{setoid}{A setoid $s$ is a 2-tuple whose components are listed below.\par
\begin{enumerate}[nosep]
\item $\operatorname{pts}(s) = \operatorname{nth}(1, s)$, a carrier (\texttt{carriers}): $\operatorname{pts}(s) \in \mathbf{Set}$.
\item $\operatorname{rel}(s) = \operatorname{nth}(2, s)$, a relation (\texttt{relation}): $\operatorname{rel}(s) \in \operatorname{power}((\operatorname{pts}(s) \times \operatorname{pts}(s)))$.
\end{enumerate}
The laws (the \texttt{property} and \texttt{law} clauses):
\begin{itemize}[nosep]
\item $\operatorname{is\text{-}equivalence}(\operatorname{rel}(s), \operatorname{pts}(s))$
\end{itemize}
The characteristic law (the axiom \texttt{is-setoid}): $\operatorname{is\text{-}setoid}(s)$ holds if and only if $\operatorname{length}(s) = 2$, every component has the typing listed with it, and every law holds.
}
\par\medskip\noindent\textbf{Definition} (\texttt{quotient}).\;$\operatorname{quotient}(s) :=$ $\operatorname{image}(\operatorname{proj}(s), \operatorname{pts}(s))$
\par\medskip\noindent\textbf{Definition} (\texttt{class}).\;$\operatorname{class}(s, a) :=$ $\{b \in \operatorname{pts}(s) : \operatorname{related}(s, a, b)\}$
\par\medskip\noindent\textbf{Definition} (\texttt{related}).\;$\operatorname{related}(s, a, b) :=$ $\langle a, b \rangle \in \operatorname{rel}(s)$
\par\medskip\noindent\textbf{Definition} (\texttt{is-equivalence}).\;For all $\rho$ and $\mathit{crr}$, 
\begin{equation*}
\begin{array}{@{}l@{}} \operatorname{is\text{-}equivalence}(\rho, \mathit{crr}) \\ \Leftrightarrow\; \rho \in \operatorname{power}((\mathit{crr} \times \mathit{crr})) \wedge \\ \forall u \in \mathit{crr}.\; (\langle u, u \rangle \in \rho) \wedge \\ \forall u, v \in \mathit{crr}.\; ((\langle u, v \rangle \in \rho) \Rightarrow (\langle v, u \rangle \in \rho)) \wedge \\ \forall u, v, w \in \mathit{crr}. \\ \quad ((\langle u, v \rangle \in \rho) \wedge (\langle v, w \rangle \in \rho)) \Rightarrow (\langle u, w \rangle \in \rho) \end{array}
\end{equation*}

\par\medskip\noindent\textbf{Definition} (\texttt{proj}).\;$\operatorname{proj}(s) :=$ $(\lambda a \in \operatorname{pts}(s).\; \operatorname{class}(s, a))$
\par\medskip\noindent\textbf{Definition} (\texttt{respects}).\;$\operatorname{respects}(s, f) :=$
\begin{equation*}
\forall a, b \in \operatorname{pts}(s).\; (\operatorname{related}(s, a, b) \Rightarrow (f(a) = f(b)))
\end{equation*}

\par\medskip\noindent\textbf{Definition} (\texttt{descend}).\;$\operatorname{descend}(s, f) :=$
\begin{equation*}
(\lambda c \in \operatorname{quotient}(s).\; (\iota\, z.\;\exists a \in c.\; (z = f(a))))
\end{equation*}

\par\medskip\noindent\textbf{Definition} (\texttt{respects2}).\;$\operatorname{respects2}(s, f) :=$
\begin{equation*}
\begin{array}{@{}l@{}} \forall a, b, \operatorname{a\_}, \operatorname{b\_} \in \operatorname{pts}(s). \\ \quad (\operatorname{related}(s, a, \operatorname{a\_}) \wedge \operatorname{related}(s, b, \operatorname{b\_})) \Rightarrow (f(a, b) = f(\operatorname{a\_}, \operatorname{b\_})) \end{array}
\end{equation*}

\par\medskip\noindent\textbf{Definition} (\texttt{descend2}).\;$\operatorname{descend2}(s, f) :=$
\begin{equation*}
\begin{array}{@{}l@{}} (\lambda \langle c, d \rangle \in (\operatorname{quotient}(s) \times \operatorname{quotient}(s)).\; \\ \quad (\iota\, z.\;\exists a \in c, b \in d.\; (z = f(a, b)))) \end{array}
\end{equation*}

\par\medskip\noindent\textbf{Definition} (\texttt{hom-setoid}).\;$\operatorname{hom\text{-}setoid}(a, b) :=$ $\{\operatorname{homf\_} \in (\operatorname{pts}(a) \to \operatorname{pts}(b)) : \operatorname{is\text{-}hom\text{-}setoid}(a, b, \operatorname{homf\_})\}$
\vnbdef{setoid-class}{\vnbstmt{\forall s.\; ((s \in \mathit{setoid}) \Leftrightarrow \operatorname{is\text{-}setoid}(s))}}
\vnbdef{is-hom-setoid-def}{\vnbstmt{\begin{array}{@{}l@{}} \forall a, b, f. \\ \quad \operatorname{is\text{-}hom\text{-}setoid}(a, b, f) \\ \quad \Leftrightarrow\; \operatorname{is\text{-}setoid}(a) \wedge \\ \quad \operatorname{is\text{-}setoid}(b) \wedge \\ \quad f \in (\operatorname{pts}(a) \to \operatorname{pts}(b)) \wedge \\ \quad \forall x_{1}, x_{2} \in \operatorname{pts}(a). \\ \quad \quad (\langle x_{1}, x_{2} \rangle \in \operatorname{rel}(a)) \Rightarrow (\langle f(x_{1}), f(x_{2}) \rangle \in \operatorname{rel}(b)) \end{array}}}

\vnbbookfile{structure-library/ringoid}{structure-library/ringoid}
\noindent{\small ringoid.scm -- a ring with a distinguished two-sided ideal.}\par
\vnbstructdef{ringoid}{A ringoid $s$ is a 7-tuple whose components are listed below.\par
\begin{enumerate}[nosep]
\item $\operatorname{carr}(s) = \operatorname{nth}(1, s)$, a carrier (\texttt{carriers}): $\operatorname{carr}(s) \in \mathbf{Set}$.
\item $\operatorname{add}(s) = \operatorname{nth}(2, s)$, an operation (\texttt{op}): $\operatorname{add}(s) \in ((\operatorname{carr}(s) \times \operatorname{carr}(s)) \to \operatorname{carr}(s))$.
\item $\operatorname{mul}(s) = \operatorname{nth}(3, s)$, an operation (\texttt{op}): $\operatorname{mul}(s) \in ((\operatorname{carr}(s) \times \operatorname{carr}(s)) \to \operatorname{carr}(s))$.
\item $\operatorname{neg}(s) = \operatorname{nth}(4, s)$, an operation (\texttt{op}): $\operatorname{neg}(s) \in (\operatorname{carr}(s) \to \operatorname{carr}(s))$.
\item $\operatorname{zero}(s) = \operatorname{nth}(5, s)$, a constant (\texttt{constant}): $\operatorname{zero}(s) \in \operatorname{carr}(s)$.
\item $\operatorname{one}(s) = \operatorname{nth}(6, s)$, a constant (\texttt{constant}): $\operatorname{one}(s) \in \operatorname{carr}(s)$.
\item $\operatorname{idl}(s) = \operatorname{nth}(7, s)$, a constant (\texttt{constant}): $\operatorname{idl}(s) \in \operatorname{power}(\operatorname{carr}(s))$.
\end{enumerate}
The laws (the \texttt{property} and \texttt{law} clauses):
\begin{itemize}[nosep]
\item $\operatorname{is\text{-}associative}(\operatorname{add}(s), \operatorname{carr}(s))$
\item $\operatorname{is\text{-}commutative}(\operatorname{add}(s), \operatorname{carr}(s))$
\item $\operatorname{is\text{-}identity}(\operatorname{add}(s), \operatorname{zero}(s), \operatorname{carr}(s))$
\item $\operatorname{has\text{-}inverses}(\operatorname{add}(s), \operatorname{zero}(s), \operatorname{neg}(s), \operatorname{carr}(s))$
\item $\operatorname{is\text{-}associative}(\operatorname{mul}(s), \operatorname{carr}(s))$
\item $\operatorname{is\text{-}identity}(\operatorname{mul}(s), \operatorname{one}(s), \operatorname{carr}(s))$
\item $\operatorname{is\text{-}distributive}(\operatorname{add}(s), \operatorname{mul}(s), \operatorname{carr}(s))$
\item $\operatorname{is\text{-}ideal\text{-}in}(\operatorname{idl}(s), \operatorname{carr}(s), \operatorname{add}(s), \operatorname{mul}(s), \operatorname{neg}(s), \operatorname{zero}(s))$
\end{itemize}
The characteristic law (the axiom \texttt{is-ringoid}): $\operatorname{is\text{-}ringoid}(s)$ holds if and only if $\operatorname{length}(s) = 7$, every component has the typing listed with it, and every law holds.
}
\par\medskip\noindent\textbf{Definition} (\texttt{ringoid-quotient}).\;$\operatorname{ringoid\text{-}quotient}(r) :=$ $\operatorname{quotient}(\operatorname{ringoid\text{-}setoid}(r))$
\par\medskip\noindent\textbf{Definition} (\texttt{is-ideal-in}).\;For all $\operatorname{i\_}$, $\operatorname{carr\_}$, $\operatorname{add\_}$, $\operatorname{mul\_}$, $\operatorname{neg\_}$ and $\operatorname{zero\_}$, 
\begin{equation*}
\begin{array}{@{}l@{}} \operatorname{is\text{-}ideal\text{-}in}(\operatorname{i\_}, \operatorname{carr\_}, \operatorname{add\_}, \operatorname{mul\_}, \operatorname{neg\_}, \operatorname{zero\_}) \\ \Leftrightarrow\; \operatorname{i\_} \subseteq \operatorname{carr\_} \wedge \\ \operatorname{zero\_} \in \operatorname{i\_} \wedge \\ \forall a, b \in \operatorname{i\_}.\; (\operatorname{add\_}(a, b) \in \operatorname{i\_}) \wedge \\ \forall a \in \operatorname{i\_}.\; (\operatorname{neg\_}(a) \in \operatorname{i\_}) \wedge \\ \forall \operatorname{x\_} \in \operatorname{carr\_}, a \in \operatorname{i\_}.\; (\operatorname{mul\_}(\operatorname{x\_}, a) \in \operatorname{i\_}) \wedge \\ \forall \operatorname{x\_} \in \operatorname{carr\_}, a \in \operatorname{i\_}.\; (\operatorname{mul\_}(a, \operatorname{x\_}) \in \operatorname{i\_}) \end{array}
\end{equation*}

\par\medskip\noindent\textbf{Definition} (\texttt{ringoid-rel}).\;$\operatorname{ringoid\text{-}rel}(r) :=$
\begin{equation*}
\{p \in (\operatorname{carr}(r) \times \operatorname{carr}(r)) : (\operatorname{add}(r)(\operatorname{nth}(1, p), \operatorname{neg}(r)(\operatorname{nth}(2, p))) \in \operatorname{idl}(r))\}
\end{equation*}

\par\medskip\noindent\textbf{Definition} (\texttt{ringoid-setoid}).\;$\operatorname{ringoid\text{-}setoid}(r) :=$ $\langle \operatorname{carr}(r), \operatorname{ringoid\text{-}rel}(r) \rangle$
\par\medskip\noindent\textbf{Definition} (\texttt{ringoid-quotient-ring}).\;$\operatorname{ringoid\text{-}quotient\text{-}ring}(r) :=$
\begin{equation*}
\langle \operatorname{quotient}(\operatorname{ringoid\text{-}setoid}(r)), \operatorname{descend2}(\operatorname{ringoid\text{-}setoid}(r), (\lambda \langle a, b \rangle \in (\operatorname{carr}(r) \times \operatorname{carr}(r)).\; \operatorname{class}(\operatorname{ringoid\text{-}setoid}(r), \operatorname{add}(r)(a, b)))), \operatorname{descend2}(\operatorname{ringoid\text{-}setoid}(r), (\lambda \langle a, b \rangle \in (\operatorname{carr}(r) \times \operatorname{carr}(r)).\; \operatorname{class}(\operatorname{ringoid\text{-}setoid}(r), \operatorname{mul}(r)(a, b)))), \operatorname{descend}(\operatorname{ringoid\text{-}setoid}(r), (\lambda a \in \operatorname{carr}(r).\; \operatorname{class}(\operatorname{ringoid\text{-}setoid}(r), \operatorname{neg}(r)(a)))), \operatorname{class}(\operatorname{ringoid\text{-}setoid}(r), \operatorname{zero}(r)), \operatorname{class}(\operatorname{ringoid\text{-}setoid}(r), \operatorname{one}(r)) \rangle
\end{equation*}

\par\medskip\noindent\textbf{Definition} (\texttt{hom-ringoid}).\;$\operatorname{hom\text{-}ringoid}(a, b) :=$
\begin{equation*}
\{\operatorname{homf\_} \in (\operatorname{carr}(a) \to \operatorname{carr}(b)) : \operatorname{is\text{-}hom\text{-}ringoid}(a, b, \operatorname{homf\_})\}
\end{equation*}

\vnbdef{rq-carr}{\vnbstmt{\forall r.\; (\operatorname{carr}(\operatorname{ringoid\text{-}quotient\text{-}ring}(r)) \simeq \operatorname{quotient}(\operatorname{ringoid\text{-}setoid}(r)))}}
\vnbdef{rq-add}{\vnbstmt{\begin{array}{@{}l@{}} \forall r. \\ \quad (\operatorname{add}(\operatorname{ringoid\text{-}quotient\text{-}ring}(r)) \simeq \operatorname{descend2}(\operatorname{ringoid\text{-}setoid}(r), (\lambda \langle a, b \rangle \in (\operatorname{carr}(r) \times \operatorname{carr}(r)).\; \operatorname{class}(\operatorname{ringoid\text{-}setoid}(r), \operatorname{add}(r)(a, b))))) \end{array}}}
\vnbdef{rq-mul}{\vnbstmt{\begin{array}{@{}l@{}} \forall r. \\ \quad (\operatorname{mul}(\operatorname{ringoid\text{-}quotient\text{-}ring}(r)) \simeq \operatorname{descend2}(\operatorname{ringoid\text{-}setoid}(r), (\lambda \langle a, b \rangle \in (\operatorname{carr}(r) \times \operatorname{carr}(r)).\; \operatorname{class}(\operatorname{ringoid\text{-}setoid}(r), \operatorname{mul}(r)(a, b))))) \end{array}}}
\vnbdef{rq-neg}{\vnbstmt{\begin{array}{@{}l@{}} \forall r. \\ \quad (\operatorname{neg}(\operatorname{ringoid\text{-}quotient\text{-}ring}(r)) \simeq \operatorname{descend}(\operatorname{ringoid\text{-}setoid}(r), (\lambda a \in \operatorname{carr}(r).\; \operatorname{class}(\operatorname{ringoid\text{-}setoid}(r), \operatorname{neg}(r)(a))))) \end{array}}}
\vnbdef{rq-zero}{\vnbstmt{\forall r.\; (\operatorname{zero}(\operatorname{ringoid\text{-}quotient\text{-}ring}(r)) \simeq \operatorname{class}(\operatorname{ringoid\text{-}setoid}(r), \operatorname{zero}(r)))}}
\vnbdef{rq-one}{\vnbstmt{\forall r.\; (\operatorname{one}(\operatorname{ringoid\text{-}quotient\text{-}ring}(r)) \simeq \operatorname{class}(\operatorname{ringoid\text{-}setoid}(r), \operatorname{one}(r)))}}
\vnbdef{ringoid-setoid-pts}{\vnbstmt{\forall r.\; (\operatorname{pts}(\operatorname{ringoid\text{-}setoid}(r)) \simeq \operatorname{carr}(r))}}
\vnbdef{ringoid-rel-mem}{\vnbstmt{\begin{array}{@{}l@{}} \forall r, a, b. \\ \quad \langle a, b \rangle \in \operatorname{ringoid\text{-}rel}(r) \\ \quad \Leftrightarrow\; (\langle a, b \rangle \in (\operatorname{carr}(r) \times \operatorname{carr}(r))) \wedge (\operatorname{add}(r)(a, \operatorname{neg}(r)(b)) \in \operatorname{idl}(r)) \end{array}}}
\vnbdef{ringoid-carr-in-set}{\vnbstmt{\forall r.\; (\operatorname{is\text{-}ringoid}(r) \Rightarrow (\operatorname{carr}(r) \in \mathbf{Set}))}}
\vnbdef{ringoid-add-right-inv}{\vnbstmt{\begin{array}{@{}l@{}} \forall r, a. \\ \quad (\operatorname{is\text{-}ringoid}(r) \wedge (a \in \operatorname{carr}(r))) \Rightarrow (\operatorname{add}(r)(a, \operatorname{neg}(r)(a)) = \operatorname{zero}(r)) \end{array}}}
\vnbdef{ringoid-ideal-subset}{\vnbstmt{\forall r.\; (\operatorname{is\text{-}ringoid}(r) \Rightarrow (\operatorname{idl}(r) \subseteq \operatorname{carr}(r)))}}
\vnbdef{ringoid-ideal-zero}{\vnbstmt{\forall r.\; (\operatorname{is\text{-}ringoid}(r) \Rightarrow (\operatorname{zero}(r) \in \operatorname{idl}(r)))}}
\vnbdef{ringoid-ideal-add}{\vnbstmt{\begin{array}{@{}l@{}} \forall r, a, b. \\ \quad (\operatorname{is\text{-}ringoid}(r) \wedge ((a \in \operatorname{idl}(r)) \wedge (b \in \operatorname{idl}(r)))) \Rightarrow (\operatorname{add}(r)(a, b) \in \operatorname{idl}(r)) \end{array}}}
\vnbdef{ringoid-ideal-neg}{\vnbstmt{\forall r, a.\; ((\operatorname{is\text{-}ringoid}(r) \wedge (a \in \operatorname{idl}(r))) \Rightarrow (\operatorname{neg}(r)(a) \in \operatorname{idl}(r)))}}
\vnbdef{ringoid-ideal-mul-left}{\vnbstmt{\begin{array}{@{}l@{}} \forall r, x, a. \\ \quad (\operatorname{is\text{-}ringoid}(r) \wedge ((x \in \operatorname{carr}(r)) \wedge (a \in \operatorname{idl}(r)))) \Rightarrow (\operatorname{mul}(r)(x, a) \in \operatorname{idl}(r)) \end{array}}}
\vnbdef{ringoid-ideal-mul-right}{\vnbstmt{\begin{array}{@{}l@{}} \forall r, a, x. \\ \quad (\operatorname{is\text{-}ringoid}(r) \wedge ((a \in \operatorname{idl}(r)) \wedge (x \in \operatorname{carr}(r)))) \Rightarrow (\operatorname{mul}(r)(a, x) \in \operatorname{idl}(r)) \end{array}}}
\vnbdef{ringoid-as-ring-is-ring}{\vnbstmt{\forall r.\; (\operatorname{is\text{-}ringoid}(r) \Rightarrow \operatorname{is\text{-}ring}(\operatorname{ringoid\text{-}as\text{-}ring}(r)))}}
\vnbdef{ringoid-class}{\vnbstmt{\forall s.\; ((s \in \mathit{ringoid}) \Leftrightarrow \operatorname{is\text{-}ringoid}(s))}}
\vnbdef{is-hom-ringoid-def}{\vnbstmt{\begin{array}{@{}l@{}} \forall a, b, f. \\ \quad \operatorname{is\text{-}hom\text{-}ringoid}(a, b, f) \\ \quad \Leftrightarrow\; \operatorname{is\text{-}ringoid}(a) \wedge \\ \quad \operatorname{is\text{-}ringoid}(b) \wedge \\ \quad f \in (\operatorname{carr}(a) \to \operatorname{carr}(b)) \wedge \\ \quad \operatorname{is\text{-}hom\text{-}ring}(\operatorname{ringoid\text{-}as\text{-}ring}(a), \operatorname{ringoid\text{-}as\text{-}ring}(b), f) \wedge \\ \quad \forall x \in \operatorname{idl}(a).\; (f(x) \in \operatorname{idl}(b)) \end{array}}}

\vnbbookfile{structure-library/commutative-ring}{structure-library/commutative-ring}
\noindent{\small commutative-ring.scm -- COMMUTATIVE-RING: a ring whose MUL is commutative.}\par
\par\medskip\noindent\textbf{Definition} (\texttt{hom-commutative-ring}).\;$\operatorname{hom\text{-}commutative\text{-}ring}(a, b) :=$
\begin{equation*}
\{\operatorname{homf\_} \in (\operatorname{carr}(a) \to \operatorname{carr}(b)) : \operatorname{is\text{-}hom\text{-}commutative\text{-}ring}(a, b, \operatorname{homf\_})\}
\end{equation*}

\vnbdef{is-commutative-ring-def}{\vnbstmt{\begin{array}{@{}l@{}} \forall s. \\ \quad \operatorname{is\text{-}commutative\text{-}ring}(s) \\ \quad \Leftrightarrow\; (\operatorname{is\text{-}ring}(s) \wedge \forall a, b \in \operatorname{carr}(s).\; (\operatorname{mul}(s)(a, b) = \operatorname{mul}(s)(b, a))) \end{array}}}
\vnbdef{is-hom-commutative-ring-def}{\vnbstmt{\begin{array}{@{}l@{}} \forall a, b, f. \\ \quad \operatorname{is\text{-}hom\text{-}commutative\text{-}ring}(a, b, f) \\ \quad \Leftrightarrow\; \operatorname{is\text{-}commutative\text{-}ring}(a) \wedge (\operatorname{is\text{-}commutative\text{-}ring}(b) \wedge \operatorname{is\text{-}hom\text{-}ring}(a, b, f)) \end{array}}}
\vnbdef{commutative-ring-class}{\vnbstmt{\forall s.\; ((s \in \operatorname{commutative\text{-}ring}) \Leftrightarrow \operatorname{is\text{-}commutative\text{-}ring}(s))}}

\vnbbookfile{structure-library/integral-domain}{structure-library/integral-domain}
\noindent{\small integral-domain.scm -- INTEGRAL-DOMAIN: a commutative ring with ONE /= ZERO and no zero divisors.  Reuses the RING shape and accessors.}\par
\par\medskip\noindent\textbf{Definition} (\texttt{hom-integral-domain}).\;$\operatorname{hom\text{-}integral\text{-}domain}(a, b) :=$
\begin{equation*}
\{\operatorname{homf\_} \in (\operatorname{carr}(a) \to \operatorname{carr}(b)) : \operatorname{is\text{-}hom\text{-}integral\text{-}domain}(a, b, \operatorname{homf\_})\}
\end{equation*}

\vnbdef{is-integral-domain-def}{\vnbstmt{\begin{array}{@{}l@{}} \forall s. \\ \quad \operatorname{is\text{-}integral\text{-}domain}(s) \\ \quad \Leftrightarrow\; \operatorname{is\text{-}commutative\text{-}ring}(s) \wedge \\ \quad \neg (\operatorname{one}(s) = \operatorname{zero}(s)) \wedge \\ \quad \forall a, b \in \operatorname{carr}(s). \\ \quad \quad (\operatorname{mul}(s)(a, b) = \operatorname{zero}(s)) \Rightarrow ((a = \operatorname{zero}(s)) \vee (b = \operatorname{zero}(s))) \end{array}}}
\vnbdef{integral-domain-class}{\vnbstmt{\forall s.\; ((s \in \operatorname{integral\text{-}domain}) \Leftrightarrow \operatorname{is\text{-}integral\text{-}domain}(s))}}
\vnbdef{is-hom-integral-domain-def}{\vnbstmt{\begin{array}{@{}l@{}} \forall a, b, f. \\ \quad \operatorname{is\text{-}hom\text{-}integral\text{-}domain}(a, b, f) \\ \quad \Leftrightarrow\; (\operatorname{is\text{-}integral\text{-}domain}(a) \wedge (\operatorname{is\text{-}integral\text{-}domain}(b) \wedge \operatorname{is\text{-}hom\text{-}commutative\text{-}ring}(a, b, f))) \end{array}}}

\vnbbookfile{structure-library/field}{structure-library/field}
\noindent{\small field.scm -- FIELD: a commutative ring in which every nonzero element has a multiplicative inverse.}\par
\vnbstructdef{field}{A field $s$ is an 8-tuple whose components are listed below.\par
\begin{enumerate}[nosep]
\item $\operatorname{carr}(s) = \operatorname{nth}(1, s)$, a carrier (\texttt{carriers}): $\operatorname{carr}(s) \in \mathbf{Set}$.
\item $\operatorname{add}(s) = \operatorname{nth}(2, s)$, an operation (\texttt{op}): $\operatorname{add}(s) \in ((\operatorname{carr}(s) \times \operatorname{carr}(s)) \to \operatorname{carr}(s))$.
\item $\operatorname{mul}(s) = \operatorname{nth}(3, s)$, an operation (\texttt{op}): $\operatorname{mul}(s) \in ((\operatorname{carr}(s) \times \operatorname{carr}(s)) \to \operatorname{carr}(s))$.
\item $\operatorname{neg}(s) = \operatorname{nth}(4, s)$, an operation (\texttt{op}): $\operatorname{neg}(s) \in (\operatorname{carr}(s) \to \operatorname{carr}(s))$.
\item $\operatorname{zero}(s) = \operatorname{nth}(5, s)$, a constant (\texttt{constant}): $\operatorname{zero}(s) \in \operatorname{carr}(s)$.
\item $\operatorname{one}(s) = \operatorname{nth}(6, s)$, a constant (\texttt{constant}): $\operatorname{one}(s) \in \operatorname{carr}(s)$.
\item $\operatorname{non\text{-}zero}(s) = \operatorname{nth}(7, s)$, a derived carrier, fixed by the other components (\texttt{derived}): $\operatorname{non\text{-}zero}(s) \in \mathbf{Set}$ and $\operatorname{non\text{-}zero}(s) = \operatorname{difference}(\operatorname{carr}(s), \operatorname{singleton}(\operatorname{zero}(s)))$.
\item $\operatorname{mul\text{-}inv}(s) = \operatorname{nth}(8, s)$, an operation (\texttt{op}): $\operatorname{mul\text{-}inv}(s) \in (\operatorname{non\text{-}zero}(s) \to \operatorname{non\text{-}zero}(s))$.
\end{enumerate}
The laws (the \texttt{property} and \texttt{law} clauses):
\begin{itemize}[nosep]
\item $\operatorname{is\text{-}associative}(\operatorname{add}(s), \operatorname{carr}(s))$
\item $\operatorname{is\text{-}commutative}(\operatorname{add}(s), \operatorname{carr}(s))$
\item $\operatorname{is\text{-}identity}(\operatorname{add}(s), \operatorname{zero}(s), \operatorname{carr}(s))$
\item $\operatorname{has\text{-}inverses}(\operatorname{add}(s), \operatorname{zero}(s), \operatorname{neg}(s), \operatorname{carr}(s))$
\item $\operatorname{is\text{-}associative}(\operatorname{mul}(s), \operatorname{carr}(s))$
\item $\operatorname{is\text{-}commutative}(\operatorname{mul}(s), \operatorname{carr}(s))$
\item $\operatorname{is\text{-}identity}(\operatorname{mul}(s), \operatorname{one}(s), \operatorname{carr}(s))$
\item $\operatorname{is\text{-}distributive}(\operatorname{add}(s), \operatorname{mul}(s), \operatorname{carr}(s))$
\item $\neg (\operatorname{zero}(s) = \operatorname{one}(s))$
\item $\forall \operatorname{a\_} \in \operatorname{non\text{-}zero}(s).\; (\operatorname{mul}(s)(\operatorname{a\_}, \operatorname{mul\text{-}inv}(s)(\operatorname{a\_})) = \operatorname{one}(s))$
\end{itemize}
The characteristic law (the axiom \texttt{is-field}): $\operatorname{is\text{-}field}(s)$ holds if and only if $\operatorname{length}(s) = 8$, every component has the typing listed with it, and every law holds.
}
\par\medskip\noindent\textbf{Definition} (\texttt{hom-field}).\;$\operatorname{hom\text{-}field}(a, b) :=$ $\{\operatorname{homf\_} \in (\operatorname{carr}(a) \to \operatorname{carr}(b)) : \operatorname{is\text{-}hom\text{-}field}(a, b, \operatorname{homf\_})\}$
\par\medskip\noindent\textbf{Definition} (\texttt{hom-field-ring}).\;$\operatorname{hom\text{-}field\text{-}ring}(a, b) :=$
\begin{equation*}
\{\operatorname{homf\_} \in (\operatorname{carr}(a) \to \operatorname{carr}(b)) : \operatorname{is\text{-}hom\text{-}field\text{-}ring}(a, b, \operatorname{homf\_})\}
\end{equation*}

\vnbdef{field-non-zero-carrier}{\vnbstmt{\begin{array}{@{}l@{}} \forall s. \\ \quad \operatorname{is\text{-}field}(s) \Rightarrow (\operatorname{non\text{-}zero}(s) = \operatorname{difference}(\operatorname{carr}(s), \operatorname{singleton}(\operatorname{zero}(s)))) \end{array}}}
\vnbdef{is-field-ring-def}{\vnbstmt{\begin{array}{@{}l@{}} \forall s. \\ \quad \operatorname{is\text{-}field\text{-}ring}(s) \\ \quad \Leftrightarrow\; \operatorname{is\text{-}commutative\text{-}ring}(s) \wedge \\ \quad \neg (\operatorname{one}(s) = \operatorname{zero}(s)) \wedge \\ \quad \forall a \in \operatorname{carr}(s). \\ \quad \quad \neg (a = \operatorname{zero}(s)) \Rightarrow \exists b \in \operatorname{carr}(s).\; (\operatorname{mul}(s)(a, b) = \operatorname{one}(s)) \end{array}}}
\vnbdef{is-hom-field-ring-def}{\vnbstmt{\begin{array}{@{}l@{}} \forall a, b, f. \\ \quad \operatorname{is\text{-}hom\text{-}field\text{-}ring}(a, b, f) \\ \quad \Leftrightarrow\; \operatorname{is\text{-}field\text{-}ring}(a) \wedge (\operatorname{is\text{-}field\text{-}ring}(b) \wedge \operatorname{is\text{-}hom\text{-}commutative\text{-}ring}(a, b, f)) \end{array}}}
\vnbdef{field-class}{\vnbstmt{\forall s.\; ((s \in \mathit{field}) \Leftrightarrow \operatorname{is\text{-}field}(s))}}
\vnbdef{is-hom-field-def}{\vnbstmt{\begin{array}{@{}l@{}} \forall a, b, f. \\ \quad \operatorname{is\text{-}hom\text{-}field}(a, b, f) \\ \quad \Leftrightarrow\; \operatorname{is\text{-}field}(a) \wedge \\ \quad \operatorname{is\text{-}field}(b) \wedge \\ \quad f \in (\operatorname{carr}(a) \to \operatorname{carr}(b)) \wedge \\ \quad \forall x_{1}, x_{2} \in \operatorname{carr}(a). \\ \quad \quad f(\operatorname{add}(a)(x_{1}, x_{2})) = \operatorname{add}(b)(f(x_{1}), f(x_{2})) \wedge \\ \quad \forall x_{1}, x_{2} \in \operatorname{carr}(a). \\ \quad \quad f(\operatorname{mul}(a)(x_{1}, x_{2})) = \operatorname{mul}(b)(f(x_{1}), f(x_{2})) \wedge \\ \quad \forall x_{1} \in \operatorname{carr}(a).\; (f(\operatorname{neg}(a)(x_{1})) = \operatorname{neg}(b)(f(x_{1}))) \wedge \\ \quad f(\operatorname{zero}(a)) = \operatorname{zero}(b) \wedge \\ \quad f(\operatorname{one}(a)) = \operatorname{one}(b) \wedge \\ \quad \forall x_{1} \in \operatorname{non\text{-}zero}(a).\; (f(\operatorname{mul\text{-}inv}(a)(x_{1})) = \operatorname{mul\text{-}inv}(b)(f(x_{1}))) \end{array}}}
\vnbdef{field-ring-class}{\vnbstmt{\forall s.\; ((s \in \operatorname{field\text{-}ring}) \Leftrightarrow \operatorname{is\text{-}field\text{-}ring}(s))}}

\vnbbookfile{structure-library/euclidean-ring}{structure-library/euclidean-ring}
\noindent{\small euclidean-ring.scm -- EUCLIDEAN-RING: an integral domain admitting a Euclidean (degree) function with division-with-remainder.}\par
\par\medskip\noindent\textbf{Definition} (\texttt{euclidean-gauges}).\;$\operatorname{euclidean\text{-}gauges}(s) :=$ $\{\mathit{dg} \in (\operatorname{carr}(s) \to \mathbb{N}) : \operatorname{has\text{-}div\text{-}remainder}(s, \mathit{dg})\}$
\par\medskip\noindent\textbf{Definition} (\texttt{gauge}).\;$\operatorname{gauge}(s) :=$ $\varepsilon(\operatorname{euclidean\text{-}gauges}(s))$
\par\medskip\noindent\textbf{Definition} (\texttt{has-div-remainder}).\;For all $s$ and $\mathit{dg}$, 
\begin{equation*}
\begin{array}{@{}l@{}} \operatorname{has\text{-}div\text{-}remainder}(s, \mathit{dg}) \\ \Leftrightarrow\; \forall \operatorname{a\_}, b \in \operatorname{carr}(s). \\ \quad \neg (b = \operatorname{zero}(s)) \Rightarrow  \\ \quad \quad \exists q, r \in \operatorname{carr}(s). \\ \quad \quad \quad (\operatorname{a\_} = \operatorname{add}(s)(\operatorname{mul}(s)(q, b), r)) \wedge ((r = \operatorname{zero}(s)) \vee (\operatorname{succ}(\operatorname{dg}(r)) \leq \operatorname{dg}(b))) \end{array}
\end{equation*}

\par\medskip\noindent\textbf{Definition} (\texttt{hom-euclidean-ring}).\;$\operatorname{hom\text{-}euclidean\text{-}ring}(a, b) :=$
\begin{equation*}
\{\operatorname{homf\_} \in (\operatorname{carr}(a) \to \operatorname{carr}(b)) : \operatorname{is\text{-}hom\text{-}euclidean\text{-}ring}(a, b, \operatorname{homf\_})\}
\end{equation*}

\vnbdef{is-euclidean-ring-def}{\vnbstmt{\begin{array}{@{}l@{}} \forall s. \\ \quad \operatorname{is\text{-}euclidean\text{-}ring}(s) \\ \quad \Leftrightarrow\; \operatorname{is\text{-}integral\text{-}domain}(s) \wedge \\ \quad \exists g \in (\operatorname{carr}(s) \to \mathbb{N}). \forall a, b \in \operatorname{carr}(s). \\ \quad \quad \neg (b = \operatorname{zero}(s)) \Rightarrow  \\ \quad \quad \quad \exists q, r \in \operatorname{carr}(s). \\ \quad \quad \quad \quad (a = \operatorname{add}(s)(\operatorname{mul}(s)(q, b), r)) \wedge ((r = \operatorname{zero}(s)) \vee ((g(r) + 1) \leq g(b))) \end{array}}}
\vnbdef{is-hom-euclidean-ring-def}{\vnbstmt{\begin{array}{@{}l@{}} \forall a, b, f. \\ \quad \operatorname{is\text{-}hom\text{-}euclidean\text{-}ring}(a, b, f) \\ \quad \Leftrightarrow\; (\operatorname{is\text{-}euclidean\text{-}ring}(a) \wedge (\operatorname{is\text{-}euclidean\text{-}ring}(b) \wedge \operatorname{is\text{-}hom\text{-}integral\text{-}domain}(a, b, f))) \end{array}}}
\vnbdef{euclidean-ring-class}{\vnbstmt{\forall s.\; ((s \in \operatorname{euclidean\text{-}ring}) \Leftrightarrow \operatorname{is\text{-}euclidean\text{-}ring}(s))}}

\vnbbookfile{structure-library/ideal}{structure-library/ideal}
\noindent{\small RETIRED 2026-09-17 (proven): ideal-elt-in-carrier, principal-ideal-in-ideal -- theorem-library/rake-algebra.scm ideal.scm -- ideals, principal ideals, and principal-ideal domains over a commutative ring.  The vocabulary the Euclidean-ring =$>$ PID stress test needs (theorem-library/euclidean-ring-pid.scm).}\par
\par\medskip\noindent\textbf{Definition} (\texttt{principal-ideal}).\;$\operatorname{principal\text{-}ideal}(s, a) :=$ $\{x \in \operatorname{carr}(s) : \exists r \in \operatorname{carr}(s).\; (x = \operatorname{mul}(s)(r, a))\}$
\par\medskip\noindent\textbf{Definition} (\texttt{is-ideal}).\;For all $s$ and $i$, 
\begin{equation*}
\begin{array}{@{}l@{}} \operatorname{is\text{-}ideal}(s, i) \\ \Leftrightarrow\; \operatorname{is\text{-}commutative\text{-}ring}(s) \wedge \\ i \subseteq \operatorname{carr}(s) \wedge \\ \operatorname{zero}(s) \in i \wedge \\ \forall a, b \in i.\; (\operatorname{add}(s)(a, b) \in i) \wedge \\ \forall a \in i.\; (\operatorname{neg}(s)(a) \in i) \wedge \\ \forall r \in \operatorname{carr}(s), a \in i.\; (\operatorname{mul}(s)(r, a) \in i) \end{array}
\end{equation*}

\par\medskip\noindent\textbf{Definition} (\texttt{hom-pid}).\;$\operatorname{hom\text{-}pid}(a, b) :=$ $\{\operatorname{homf\_} \in (\operatorname{carr}(a) \to \operatorname{carr}(b)) : \operatorname{is\text{-}hom\text{-}pid}(a, b, \operatorname{homf\_})\}$
\vnbdef{principal-ideal-membership}{\vnbstmt{\begin{array}{@{}l@{}} \forall s, a, x. \\ \quad (x \in \operatorname{principal\text{-}ideal}(s, a)) \Leftrightarrow ((x \in \operatorname{carr}(s)) \wedge \exists r \in \operatorname{carr}(s).\; (x = \operatorname{mul}(s)(r, a))) \end{array}}}
\vnbdef{pid-class}{\vnbstmt{\forall s.\; ((s \in \mathit{pid}) \Leftrightarrow \operatorname{is\text{-}pid}(s))}}
\vnbdef{is-hom-pid-def}{\vnbstmt{\begin{array}{@{}l@{}} \forall a, b, f. \\ \quad \operatorname{is\text{-}hom\text{-}pid}(a, b, f) \Leftrightarrow (\operatorname{is\text{-}pid}(a) \wedge (\operatorname{is\text{-}pid}(b) \wedge \operatorname{is\text{-}hom\text{-}integral\text{-}domain}(a, b, f))) \end{array}}}
\vnbdef{is-pid-def}{\vnbstmt{\begin{array}{@{}l@{}} \forall s. \\ \quad \operatorname{is\text{-}pid}(s) \\ \quad \Leftrightarrow\; \operatorname{is\text{-}integral\text{-}domain}(s) \wedge \\ \quad \forall i. \\ \quad \quad \operatorname{is\text{-}ideal}(s, i) \Rightarrow \exists a \in \operatorname{carr}(s).\; (i = \operatorname{principal\text{-}ideal}(s, a)) \end{array}}}

\vnbbookfile{structure-library/views}{structure-library/views}
\noindent{\small views.scm -- view-as declarations between structures}\par
\vnbdef{ring-additive-ag-is-abelian-group}{\vnbstmt{\forall r.\; (\operatorname{is\text{-}ring}(r) \Rightarrow \operatorname{is\text{-}abelian\text{-}group}(\operatorname{ring\text{-}additive\text{-}ag}(r)))}}
\vnbdef{normed-field-additive-ag-is-abelian-group}{\vnbstmt{\begin{array}{@{}l@{}} \forall r. \\ \quad \operatorname{is\text{-}normed\text{-}field}(r) \Rightarrow \operatorname{is\text{-}abelian\text{-}group}(\operatorname{normed\text{-}field\text{-}additive\text{-}ag}(r)) \end{array}}}
\vnbdef{ring-multiplicative-monoid-is-monoid}{\vnbstmt{\forall r.\; (\operatorname{is\text{-}ring}(r) \Rightarrow \operatorname{is\text{-}monoid}(\operatorname{ring\text{-}multiplicative\text{-}monoid}(r)))}}
\vnbdef{commutative-ring-multiplicative-cm-is-comm-monoid}{\vnbstmt{\begin{array}{@{}l@{}} \forall r. \\ \quad \operatorname{is\text{-}commutative\text{-}ring}(r) \Rightarrow  \\ \quad \quad \operatorname{is\text{-}comm\text{-}monoid}(\operatorname{commutative\text{-}ring\text{-}multiplicative\text{-}cm}(r)) \end{array}}}
\vnbdef{commutative-ring-additive-ag-is-abelian-group}{\vnbstmt{\begin{array}{@{}l@{}} \forall r. \\ \quad \operatorname{is\text{-}commutative\text{-}ring}(r) \Rightarrow \operatorname{is\text{-}abelian\text{-}group}(\operatorname{commutative\text{-}ring\text{-}additive\text{-}ag}(r)) \end{array}}}
\vnbdef{field-multiplicative-group-is-group}{\vnbstmt{\forall r.\; (\operatorname{is\text{-}field}(r) \Rightarrow \operatorname{is\text{-}group}(\operatorname{field\text{-}multiplicative\text{-}group}(r)))}}
\vnbdef{normed-field-as-integral-domain-is-integral-domain}{\vnbstmt{\begin{array}{@{}l@{}} \forall r. \\ \quad \operatorname{is\text{-}normed\text{-}field}(r) \Rightarrow  \\ \quad \quad \operatorname{is\text{-}integral\text{-}domain}(\operatorname{normed\text{-}field\text{-}as\text{-}integral\text{-}domain}(r)) \end{array}}}
\vnbdef{module-vector-ag-is-abelian-group}{\vnbstmt{\forall r.\; (\operatorname{is\text{-}module}(r) \Rightarrow \operatorname{is\text{-}abelian\text{-}group}(\operatorname{module\text{-}vector\text{-}ag}(r)))}}
\vnbdef{abelian-group-as-monoid-is-monoid}{\vnbstmt{\forall r.\; (\operatorname{is\text{-}abelian\text{-}group}(r) \Rightarrow \operatorname{is\text{-}monoid}(\operatorname{abelian\text{-}group\text{-}as\text{-}monoid}(r)))}}
\vnbdef{field-as-integral-domain-is-integral-domain}{\vnbstmt{\forall r.\; (\operatorname{is\text{-}field}(r) \Rightarrow \operatorname{is\text{-}integral\text{-}domain}(\operatorname{field\text{-}as\text{-}integral\text{-}domain}(r)))}}
\vnbdef{normed-ag-as-abelian-group-is-abelian-group}{\vnbstmt{\begin{array}{@{}l@{}} \forall r. \\ \quad \operatorname{is\text{-}normed\text{-}ag}(r) \Rightarrow \operatorname{is\text{-}abelian\text{-}group}(\operatorname{normed\text{-}ag\text{-}as\text{-}abelian\text{-}group}(r)) \end{array}}}
\vnbdef{field-as-euclidean-ring-is-euclidean-ring}{\vnbstmt{\forall r.\; (\operatorname{is\text{-}field}(r) \Rightarrow \operatorname{is\text{-}euclidean\text{-}ring}(\operatorname{field\text{-}as\text{-}euclidean\text{-}ring}(r)))}}
\vnbdef{field-additive-ag-is-abelian-group}{\vnbstmt{\forall r.\; (\operatorname{is\text{-}field}(r) \Rightarrow \operatorname{is\text{-}abelian\text{-}group}(\operatorname{field\text{-}additive\text{-}ag}(r)))}}
\vnbdef{normed-field-as-commutative-ring-is-commutative-ring}{\vnbstmt{\begin{array}{@{}l@{}} \forall r. \\ \quad \operatorname{is\text{-}normed\text{-}field}(r) \Rightarrow  \\ \quad \quad \operatorname{is\text{-}commutative\text{-}ring}(\operatorname{normed\text{-}field\text{-}as\text{-}commutative\text{-}ring}(r)) \end{array}}}

\vnbbookfile{structure-library/numeric-instances}{structure-library/numeric-instances}
\noindent{\small basic-rings.scm -- "officialize" NN/ZZ/QQ/RR/CC as algebraic structures}\par
\par\medskip\noindent\textbf{Definition} (\texttt{zz-ring}).\;
\begin{equation*}
\begin{array}{@{}l@{}} \operatorname{zz\text{-}ring} \\ =\; \langle \mathbb{Z}, (\lambda \langle \operatorname{x\_}, \operatorname{y\_} \rangle \in (\mathbb{Z} \times \mathbb{Z}).\; +(\operatorname{x\_}, \operatorname{y\_})), (\lambda \langle \operatorname{x\_}, \operatorname{y\_} \rangle \in (\mathbb{Z} \times \mathbb{Z}).\; \cdot(\operatorname{x\_}, \operatorname{y\_})), (\lambda \operatorname{x\_} \in \mathbb{Z}.\; -(\operatorname{x\_})), 0, 1 \rangle \end{array}
\end{equation*}

\par\medskip\noindent\textbf{Definition} (\texttt{qq-ring}).\;
\begin{equation*}
\begin{array}{@{}l@{}} \operatorname{qq\text{-}ring} \\ =\; \langle \mathbb{Q}, (\lambda \langle \operatorname{x\_}, \operatorname{y\_} \rangle \in (\mathbb{Q} \times \mathbb{Q}).\; +(\operatorname{x\_}, \operatorname{y\_})), (\lambda \langle \operatorname{x\_}, \operatorname{y\_} \rangle \in (\mathbb{Q} \times \mathbb{Q}).\; \cdot(\operatorname{x\_}, \operatorname{y\_})), (\lambda \operatorname{x\_} \in \mathbb{Q}.\; -(\operatorname{x\_})), 0, 1 \rangle \end{array}
\end{equation*}

\par\medskip\noindent\textbf{Definition} (\texttt{rr-normed-field}).\;
\begin{equation*}
\begin{array}{@{}l@{}} \operatorname{rr\text{-}normed\text{-}field} \\ =\; \langle \mathbb{R}, (\lambda \langle \operatorname{x\_}, \operatorname{y\_} \rangle \in (\mathbb{R} \times \mathbb{R}).\; +(\operatorname{x\_}, \operatorname{y\_})), (\lambda \langle \operatorname{x\_}, \operatorname{y\_} \rangle \in (\mathbb{R} \times \mathbb{R}).\; \cdot(\operatorname{x\_}, \operatorname{y\_})), (\lambda \operatorname{x\_} \in \mathbb{R}.\; -(\operatorname{x\_})), 0, 1, (\lambda \operatorname{x\_} \in \mathbb{R}.\; \lvert \operatorname{x\_} \rvert) \rangle \end{array}
\end{equation*}

\par\medskip\noindent\textbf{Definition} (\texttt{cc-normed-field}).\;
\begin{equation*}
\begin{array}{@{}l@{}} \operatorname{cc\text{-}normed\text{-}field} \\ =\; \langle \mathbb{C}, (\lambda \langle \operatorname{x\_}, \operatorname{y\_} \rangle \in (\mathbb{C} \times \mathbb{C}).\; +(\operatorname{x\_}, \operatorname{y\_})), (\lambda \langle \operatorname{x\_}, \operatorname{y\_} \rangle \in (\mathbb{C} \times \mathbb{C}).\; \cdot(\operatorname{x\_}, \operatorname{y\_})), (\lambda \operatorname{x\_} \in \mathbb{C}.\; -(\operatorname{x\_})), 0, 1, (\lambda \operatorname{x\_} \in \mathbb{C}.\; \operatorname{magnitude}(\operatorname{x\_})) \rangle \end{array}
\end{equation*}

\par\medskip\noindent\textbf{Definition} (\texttt{nn-add-monoid}).\;
\begin{equation*}
\operatorname{nn\text{-}add\text{-}monoid} = \langle \mathbb{N}, (\lambda \langle \operatorname{x\_}, \operatorname{y\_} \rangle \in (\mathbb{N} \times \mathbb{N}).\; +(\operatorname{x\_}, \operatorname{y\_})), 0 \rangle
\end{equation*}

\par\medskip\noindent\textbf{Definition} (\texttt{rr-ms}).\;
\begin{equation*}
\operatorname{rr\text{-}ms} = \langle \mathbb{R}, (\lambda \langle x, y \rangle \in (\mathbb{R} \times \mathbb{R}).\; \lvert -(x, y) \rvert) \rangle
\end{equation*}

\par\medskip\noindent\textbf{Definition} (\texttt{qq-field}).\;
\begin{equation*}
\begin{array}{@{}l@{}} \operatorname{qq\text{-}field} \\ =\; \langle \mathbb{Q}, (\lambda \langle \operatorname{x\_}, \operatorname{y\_} \rangle \in (\mathbb{Q} \times \mathbb{Q}).\; +(\operatorname{x\_}, \operatorname{y\_})), (\lambda \langle \operatorname{x\_}, \operatorname{y\_} \rangle \in (\mathbb{Q} \times \mathbb{Q}).\; \cdot(\operatorname{x\_}, \operatorname{y\_})), (\lambda \operatorname{x\_} \in \mathbb{Q}.\; -(\operatorname{x\_})), 0, 1, \operatorname{difference}(\mathbb{Q}, \operatorname{singleton}(0)), (\lambda \operatorname{x\_} \in \operatorname{difference}(\mathbb{Q}, \operatorname{singleton}(0)).\; \operatorname{recip}(\operatorname{x\_})) \rangle \end{array}
\end{equation*}

\vnbdef{binplus-apply}{\vnbstmt{\forall x, y.\; ((x + y) \simeq (x + y))}}
\vnbdef{bintimes-apply}{\vnbstmt{\forall x, y.\; ((x \cdot y) \simeq (x \cdot y))}}
\vnbdef{binneg-apply}{\vnbstmt{\forall x.\; (\operatorname{binneg}(x) \simeq -x)}}
\vnbdef{nary-plus-2}{\vnbstmt{\forall x, y.\; ((x + y) \simeq (x + y))}}
\vnbdef{nary-times-2}{\vnbstmt{\forall x, y.\; ((x \cdot y) \simeq (x \cdot y))}}
\vnbdef{nary-neg-1}{\vnbstmt{\forall x.\; (-x \simeq \operatorname{binneg}(x))}}
\vnbdef{nary-minus-2}{\vnbstmt{\forall x, y.\; ((x - y) \simeq (x + \operatorname{binneg}(y)))}}
\vnbdef{nary-plus-3}{\vnbstmt{\forall x, y, z.\; ((x + y + z) \simeq ((x + y) + z))}}
\vnbdef{nary-plus-4}{\vnbstmt{\forall w, x, y, z.\; ((w + x + y + z) \simeq (((w + x) + y) + z))}}
\vnbdef{nary-plus-5}{\vnbstmt{\begin{array}{@{}l@{}} \forall v, w, x, y, z. \\ \quad ((v + w + x + y + z) \simeq ((((v + w) + x) + y) + z)) \end{array}}}
\vnbdef{nary-times-3}{\vnbstmt{\forall x, y, z.\; ((x \cdot y \cdot z) \simeq ((x \cdot y) \cdot z))}}
\vnbdef{nary-times-4}{\vnbstmt{\forall w, x, y, z.\; ((w \cdot x \cdot y \cdot z) \simeq (((w \cdot x) \cdot y) \cdot z))}}
\vnbdef{nary-times-5}{\vnbstmt{\begin{array}{@{}l@{}} \forall v, w, x, y, z. \\ \quad ((v \cdot w \cdot x \cdot y \cdot z) \simeq ((((v \cdot w) \cdot x) \cdot y) \cdot z)) \end{array}}}
\vnbdef{rr-normed-field@add}{\vnbstmt{(\operatorname{add}(\operatorname{rr\text{-}normed\text{-}field}) \simeq (\lambda \langle \operatorname{x\_}, \operatorname{y\_} \rangle \in (\mathbb{R} \times \mathbb{R}).\; (\operatorname{x\_} + \operatorname{y\_})))}}
\vnbdef{qq-field@add}{\vnbstmt{(\operatorname{add}(\operatorname{qq\text{-}field}) \simeq (\lambda \langle \operatorname{x\_}, \operatorname{y\_} \rangle \in (\mathbb{Q} \times \mathbb{Q}).\; (\operatorname{x\_} + \operatorname{y\_})))}}
\vnbdef{qq-ring@neg}{\vnbstmt{(\operatorname{neg}(\operatorname{qq\text{-}ring}) \simeq (\lambda x \in \mathbb{Q}.\; -x))}}
\vnbdef{rr-normed-field@one}{\vnbstmt{(\operatorname{one}(\operatorname{rr\text{-}normed\text{-}field}) \simeq 1)}}
\vnbdef{qq-field@one}{\vnbstmt{(\operatorname{one}(\operatorname{qq\text{-}field}) \simeq 1)}}
\vnbdef{qq-ring@add}{\vnbstmt{(\operatorname{add}(\operatorname{qq\text{-}ring}) \simeq (\lambda \langle \operatorname{x\_}, \operatorname{y\_} \rangle \in (\mathbb{Q} \times \mathbb{Q}).\; (\operatorname{x\_} + \operatorname{y\_})))}}
\vnbdef{qq-field@zero}{\vnbstmt{(\operatorname{zero}(\operatorname{qq\text{-}field}) \simeq 0)}}
\vnbdef{cc-normed-field@carr}{\vnbstmt{(\operatorname{carr}(\operatorname{cc\text{-}normed\text{-}field}) \simeq \mathbb{C})}}
\vnbdef{rr-normed-field@neg}{\vnbstmt{(\operatorname{neg}(\operatorname{rr\text{-}normed\text{-}field}) \simeq (\lambda x \in \mathbb{R}.\; -x))}}
\vnbdef{cc-normed-field@zero}{\vnbstmt{(\operatorname{zero}(\operatorname{cc\text{-}normed\text{-}field}) \simeq 0)}}
\vnbdef{qq-field@carr}{\vnbstmt{(\operatorname{carr}(\operatorname{qq\text{-}field}) \simeq \mathbb{Q})}}
\vnbdef{rr-normed-field@carr}{\vnbstmt{(\operatorname{carr}(\operatorname{rr\text{-}normed\text{-}field}) \simeq \mathbb{R})}}
\vnbdef{nn-add-monoid@opr}{\vnbstmt{(\operatorname{opr}(\operatorname{nn\text{-}add\text{-}monoid}) \simeq (\lambda \langle \operatorname{x\_}, \operatorname{y\_} \rangle \in (\mathbb{N} \times \mathbb{N}).\; (\operatorname{x\_} + \operatorname{y\_})))}}
\vnbdef{qq-field@non-zero}{\vnbstmt{(\operatorname{non\text{-}zero}(\operatorname{qq\text{-}field}) \simeq \operatorname{difference}(\mathbb{Q}, \operatorname{singleton}(0)))}}
\vnbdef{zz-ring@mul}{\vnbstmt{(\operatorname{mul}(\operatorname{zz\text{-}ring}) \simeq (\lambda \langle \operatorname{x\_}, \operatorname{y\_} \rangle \in (\mathbb{Z} \times \mathbb{Z}).\; (\operatorname{x\_} \cdot \operatorname{y\_})))}}
\vnbdef{nn-add-monoid@iden}{\vnbstmt{(\operatorname{iden}(\operatorname{nn\text{-}add\text{-}monoid}) \simeq 0)}}
\vnbdef{zz-ring@neg}{\vnbstmt{(\operatorname{neg}(\operatorname{zz\text{-}ring}) \simeq (\lambda x \in \mathbb{Z}.\; -x))}}
\vnbdef{qq-field@mul-inv}{\vnbstmt{(\operatorname{mul\text{-}inv}(\operatorname{qq\text{-}field}) \simeq (\lambda x \in \operatorname{difference}(\mathbb{Q}, \operatorname{singleton}(0)).\; \operatorname{recip}(x)))}}
\vnbdef{qq-ring@one}{\vnbstmt{(\operatorname{one}(\operatorname{qq\text{-}ring}) \simeq 1)}}
\vnbdef{zz-ring@one}{\vnbstmt{(\operatorname{one}(\operatorname{zz\text{-}ring}) \simeq 1)}}
\vnbdef{nn-add-monoid@carr}{\vnbstmt{(\operatorname{carr}(\operatorname{nn\text{-}add\text{-}monoid}) \simeq \mathbb{N})}}
\vnbdef{cc-normed-field@neg}{\vnbstmt{(\operatorname{neg}(\operatorname{cc\text{-}normed\text{-}field}) \simeq (\lambda x \in \mathbb{C}.\; -x))}}
\vnbdef{rr-ms@dist}{\vnbstmt{(\operatorname{dist}(\operatorname{rr\text{-}ms}) \simeq (\lambda \langle x, y \rangle \in (\mathbb{R} \times \mathbb{R}).\; \lvert (x - y) \rvert))}}
\vnbdef{cc-normed-field@fnrm}{\vnbstmt{(\operatorname{fnrm}(\operatorname{cc\text{-}normed\text{-}field}) \simeq (\lambda x \in \mathbb{C}.\; \operatorname{magnitude}(x)))}}
\vnbdef{qq-field@mul}{\vnbstmt{(\operatorname{mul}(\operatorname{qq\text{-}field}) \simeq (\lambda \langle \operatorname{x\_}, \operatorname{y\_} \rangle \in (\mathbb{Q} \times \mathbb{Q}).\; (\operatorname{x\_} \cdot \operatorname{y\_})))}}
\vnbdef{zz-ring@zero}{\vnbstmt{(\operatorname{zero}(\operatorname{zz\text{-}ring}) \simeq 0)}}
\vnbdef{zz-ring@add}{\vnbstmt{(\operatorname{add}(\operatorname{zz\text{-}ring}) \simeq (\lambda \langle \operatorname{x\_}, \operatorname{y\_} \rangle \in (\mathbb{Z} \times \mathbb{Z}).\; (\operatorname{x\_} + \operatorname{y\_})))}}
\vnbdef{cc-normed-field@one}{\vnbstmt{(\operatorname{one}(\operatorname{cc\text{-}normed\text{-}field}) \simeq 1)}}
\vnbdef{qq-ring@carr}{\vnbstmt{(\operatorname{carr}(\operatorname{qq\text{-}ring}) \simeq \mathbb{Q})}}
\vnbdef{cc-normed-field@mul}{\vnbstmt{(\operatorname{mul}(\operatorname{cc\text{-}normed\text{-}field}) \simeq (\lambda \langle \operatorname{x\_}, \operatorname{y\_} \rangle \in (\mathbb{C} \times \mathbb{C}).\; (\operatorname{x\_} \cdot \operatorname{y\_})))}}
\vnbdef{cc-normed-field@add}{\vnbstmt{(\operatorname{add}(\operatorname{cc\text{-}normed\text{-}field}) \simeq (\lambda \langle \operatorname{x\_}, \operatorname{y\_} \rangle \in (\mathbb{C} \times \mathbb{C}).\; (\operatorname{x\_} + \operatorname{y\_})))}}
\vnbdef{zz-ring@carr}{\vnbstmt{(\operatorname{carr}(\operatorname{zz\text{-}ring}) \simeq \mathbb{Z})}}
\vnbdef{qq-ring@zero}{\vnbstmt{(\operatorname{zero}(\operatorname{qq\text{-}ring}) \simeq 0)}}
\vnbdef{rr-normed-field@mul}{\vnbstmt{(\operatorname{mul}(\operatorname{rr\text{-}normed\text{-}field}) \simeq (\lambda \langle \operatorname{x\_}, \operatorname{y\_} \rangle \in (\mathbb{R} \times \mathbb{R}).\; (\operatorname{x\_} \cdot \operatorname{y\_})))}}
\vnbdef{qq-ring@mul}{\vnbstmt{(\operatorname{mul}(\operatorname{qq\text{-}ring}) \simeq (\lambda \langle \operatorname{x\_}, \operatorname{y\_} \rangle \in (\mathbb{Q} \times \mathbb{Q}).\; (\operatorname{x\_} \cdot \operatorname{y\_})))}}
\vnbdef{rr-ms@pts}{\vnbstmt{(\operatorname{pts}(\operatorname{rr\text{-}ms}) \simeq \mathbb{R})}}
\vnbdef{qq-field@neg}{\vnbstmt{(\operatorname{neg}(\operatorname{qq\text{-}field}) \simeq (\lambda x \in \mathbb{Q}.\; -x))}}
\vnbdef{rr-normed-field@zero}{\vnbstmt{(\operatorname{zero}(\operatorname{rr\text{-}normed\text{-}field}) \simeq 0)}}
\vnbdef{rr-normed-field@fnrm}{\vnbstmt{(\operatorname{fnrm}(\operatorname{rr\text{-}normed\text{-}field}) \simeq (\lambda x \in \mathbb{R}.\; \lvert x \rvert))}}

\vnbbookfile{structure-library/zz-divisibility}{structure-library/zz-divisibility}
\noindent{\small zz-divisibility.scm -- divisibility, Bezout combinations, gcd and coprimality on the INTEGERS.  Vocabulary only; nothing is asserted here.}\par
\par\medskip\noindent\textbf{Definition} (\texttt{zz-divides}).\;For all $a$ and $b$, 
\begin{equation*}
\operatorname{zz\text{-}divides}(a, b) \Leftrightarrow ((a \in \mathbb{Z}) \wedge ((b \in \mathbb{Z}) \wedge \exists \operatorname{c\_} \in \mathbb{Z}.\; (b = \cdot(a, \operatorname{c\_}))))
\end{equation*}

\par\medskip\noindent\textbf{Definition} (\texttt{zz-bezout-set}).\;$\operatorname{zz\text{-}bezout\text{-}set}(a, b) :=$
\begin{equation*}
\{z \in \mathbb{Z} : \exists \operatorname{x\_}, \operatorname{y\_} \in \mathbb{Z}.\; (z = ((\operatorname{x\_} \cdot a) + (\operatorname{y\_} \cdot b)))\}
\end{equation*}

\par\medskip\noindent\textbf{Definition} (\texttt{zz-is-gcd}).\;For all $a$, $b$ and $d$, 
\begin{equation*}
\begin{array}{@{}l@{}} \operatorname{zz\text{-}is\text{-}gcd}(a, b, d) \\ \Leftrightarrow\; \operatorname{zz\text{-}divides}(d, a) \wedge \\ \operatorname{zz\text{-}divides}(d, b) \wedge \\ \forall \operatorname{e\_} \in \mathbb{Z}. \\ \quad \operatorname{zz\text{-}divides}(\operatorname{e\_}, a) \Rightarrow (\operatorname{zz\text{-}divides}(\operatorname{e\_}, b) \Rightarrow \operatorname{zz\text{-}divides}(\operatorname{e\_}, d)) \end{array}
\end{equation*}

\par\medskip\noindent\textbf{Definition} (\texttt{zz-coprime}).\;For all $a$ and $b$,  $\operatorname{zz\text{-}coprime}(a, b) \Leftrightarrow \operatorname{zz\text{-}is\text{-}gcd}(a, b, 1)$
\vnbdef{zz-bezout-set-membership}{\vnbstmt{\begin{array}{@{}l@{}} \forall a, b, z. \\ \quad (z \in \operatorname{zz\text{-}bezout\text{-}set}(a, b)) \Leftrightarrow ((z \in \mathbb{Z}) \wedge \exists x, y \in \mathbb{Z}.\; (z = ((x \cdot a) + (y \cdot b)))) \end{array}}}

\vnbbookfile{structure-library/monoid-power}{structure-library/monoid-power}
\noindent{\small RETIRED 2026-09-17 (proven): mpow-type -- was an unwarranted axiom; theorem-library/rake-algebra2.scm (mpow-type-ind, exponent outermost, then one fact). monoid-power.scm -- MPOW: the natural-number power x\textasciicircum{}n in a monoid.}\par
\par\medskip\noindent\textbf{Definition} (\texttt{mpow}).\;For all $m$ and $x$,  $(\operatorname{mpow}(m, x, 0) \simeq \operatorname{iden}(m))$\par
Let $m$ and $x$ be arbitrary. Let $n \in \mathbb{N}$. Then $(\operatorname{mpow}(m, x, \operatorname{succ}(n)) \simeq \operatorname{opr}(m)(x, \operatorname{mpow}(m, x, n)))$

\vnbbookfile{structure-library/zz-action}{structure-library/zz-action}
\noindent{\small zz-action.scm -- ZZ-ACT: the integer action on an abelian group.}\par
\par\medskip\noindent\textbf{Definition} (\texttt{zz-act}).\;$\operatorname{zz\text{-}act}(\operatorname{g\_}, \operatorname{k\_}, \operatorname{a\_}) :=$
\begin{equation*}
\left(\text{if }(\operatorname{k\_} \in \mathbb{N})\text{ then }\operatorname{mpow}(\operatorname{abelian\text{-}group\text{-}as\text{-}monoid}(\operatorname{g\_}), \operatorname{a\_}, \operatorname{k\_})\text{ else }\operatorname{inv}(\operatorname{g\_})(\operatorname{mpow}(\operatorname{abelian\text{-}group\text{-}as\text{-}monoid}(\operatorname{g\_}), \operatorname{a\_}, -\operatorname{k\_}))\right)
\end{equation*}


\vnbbookfile{structure-library/ring-power}{structure-library/ring-power}
\noindent{\small RETIRED 2026-09-17 (proven, rake batch P): ring-power-one, ring-power-add, ring-power-mult -- theorem-library/rake-combinatorics.scm, modulo 0 (mpow laws through the multiplicative view). RETIRED 2026-09-17 (proven): ring-power-type -- theorem-library/rake-finsum-typing.scm ring-power.scm -- RING-POWER: the ...}\par
\par\medskip\noindent\textbf{Definition} (\texttt{ring-power}).\;$\operatorname{ring\text{-}power}(r, x, n) :=$ $\operatorname{mpow}(\operatorname{commutative\text{-}ring\text{-}multiplicative\text{-}cm}(r), x, n)$

\vnbbookfile{structure-library/sequences}{structure-library/sequences}
\noindent{\small sequences.scm -- PROD-ORD, SUM, SUM-AG: products and sums over initial NN-segments}\par
\par\medskip\noindent\textbf{Definition} (\texttt{prod-ord}).\;For all $m$ and $f$,  $(\operatorname{prod\text{-}ord}(m, f, 0) \simeq \operatorname{iden}(m))$\par
Let $m$ and $f$ be arbitrary. Let $n \in \mathbb{N}$. Then $(\operatorname{prod\text{-}ord}(m, f, \operatorname{succ}(n)) \simeq \operatorname{opr}(m)(\operatorname{prod\text{-}ord}(m, f, n), f(n)))$
\par\medskip\noindent\textbf{Definition} (\texttt{sum}).\;For all $r$ and $f$,  $(\operatorname{sum}(r, f, 0) \simeq \operatorname{zero}(r))$\par
Let $r$ and $f$ be arbitrary. Let $n \in \mathbb{N}$. Then $(\operatorname{sum}(r, f, \operatorname{succ}(n)) \simeq \operatorname{add}(r)(\operatorname{sum}(r, f, n), f(n)))$
\par\medskip\noindent\textbf{Definition} (\texttt{sum-ag}).\;For all $\mathit{ag}$ and $f$,  $(\operatorname{sum\text{-}ag}(\mathit{ag}, f, 0) \simeq \operatorname{iden}(\mathit{ag}))$\par
Let $\mathit{ag}$ and $f$ be arbitrary. Let $n \in \mathbb{N}$. Then $(\operatorname{sum\text{-}ag}(\mathit{ag}, f, \operatorname{succ}(n)) \simeq \operatorname{opr}(\mathit{ag})(\operatorname{sum\text{-}ag}(\mathit{ag}, f, n), f(n)))$
\par\medskip\noindent\textbf{Definition} (\texttt{ring-prod-n}).\;For every $f$,  $(\operatorname{ring\text{-}prod\text{-}n}(f, 0) \simeq \operatorname{zero\text{-}ring})$\par
Let $f$ be arbitrary. Let $n \in \mathbb{N}$. Then
\begin{equation*}
(\operatorname{ring\text{-}prod\text{-}n}(f, \operatorname{succ}(n)) \simeq \operatorname{ring\text{-}prod}(\operatorname{ring\text{-}prod\text{-}n}(f, n), f(n)))
\end{equation*}


\vnbbookfile{structure-library/finprod}{structure-library/finprod}
\noindent{\small finprod.scm -- FINPROD / PROD-RING: finite PRODUCT of a function over a finite set.}\par
\par\medskip\noindent\textbf{Definition} (\texttt{finprod}).\;$\operatorname{finprod}(m, f, s) :=$ $\operatorname{finsum}(m, f, s)$
\par\medskip\noindent\textbf{Definition} (\texttt{prod-ring}).\;$\operatorname{prod\text{-}ring}(r, f, s) :=$ $\operatorname{finprod}(\operatorname{commutative\text{-}ring\text{-}multiplicative\text{-}cm}(r), f, s)$
\par\medskip\noindent\textbf{Definition} (\texttt{prod-set}).\;$\operatorname{prod\text{-}set}(\mathit{cm}, s, f) :=$ $\operatorname{finsum}(\mathit{cm}, f, s)$

\vnbbookfile{structure-library/polynomial}{structure-library/polynomial}
\noindent{\small polynomial.scm -- the monoid-algebra construction A[M], and polynomials as its specialisation A[NN].}\par
\par\medskip\noindent\textbf{Definition} (\texttt{supp}).\;$\operatorname{supp}(a, m, f) :=$ $\{\operatorname{x\_} \in \operatorname{carr}(m) : \neg (f(\operatorname{x\_}) = \operatorname{zero}(a))\}$
\par\medskip\noindent\textbf{Definition} (\texttt{finsupp}).\;$\operatorname{finsupp}(a, m) :=$ $\{f \in (\operatorname{carr}(m) \to \operatorname{carr}(a)) : (|\operatorname{supp}(a, m, f)| \in \mathbb{N})\}$
\par\medskip\noindent\textbf{Definition} (\texttt{monalg-add}).\;$\operatorname{monalg\text{-}add}(a, m, f, g) :=$ $(\lambda \operatorname{x\_} \in \operatorname{carr}(m).\; \operatorname{add}(a)(f(\operatorname{x\_}), g(\operatorname{x\_})))$
\par\medskip\noindent\textbf{Definition} (\texttt{monalg-neg}).\;$\operatorname{monalg\text{-}neg}(a, m, f) :=$ $(\lambda \operatorname{x\_} \in \operatorname{carr}(m).\; \operatorname{neg}(a)(f(\operatorname{x\_})))$
\par\medskip\noindent\textbf{Definition} (\texttt{monalg-zero}).\;$\operatorname{monalg\text{-}zero}(a, m) :=$ $(\lambda \operatorname{x\_} \in \operatorname{carr}(m).\; \operatorname{zero}(a))$
\par\medskip\noindent\textbf{Definition} (\texttt{monalg-one}).\;$\operatorname{monalg\text{-}one}(a, m) :=$ $(\lambda \operatorname{x\_} \in \operatorname{carr}(m).\; \left(\text{if }(\operatorname{x\_} = \operatorname{iden}(m))\text{ then }\operatorname{one}(a)\text{ else }\operatorname{zero}(a)\right))$
\par\medskip\noindent\textbf{Definition} (\texttt{monalg-mul}).\;$\operatorname{monalg\text{-}mul}(a, m, f, g) :=$
\begin{equation*}
\begin{array}{@{}l@{}} (\lambda \operatorname{x\_} \in \operatorname{carr}(m).\; \\ \quad \operatorname{finsum}(\operatorname{ring\text{-}additive\text{-}ag}(a), \\ \quad \quad (\lambda p \in \{p \in (\operatorname{supp}(a, m, f) \times \operatorname{supp}(a, m, g)) : (\operatorname{opr}(m)(\operatorname{nth}(1, p), \operatorname{nth}(2, p)) = \operatorname{x\_})\}.\; \\ \quad \quad \quad \operatorname{mul}(a)(f(\operatorname{nth}(1, p)), g(\operatorname{nth}(2, p)))), \\ \quad \quad \{p \in (\operatorname{supp}(a, m, f) \times \operatorname{supp}(a, m, g)) : (\operatorname{opr}(m)(\operatorname{nth}(1, p), \operatorname{nth}(2, p)) = \operatorname{x\_})\})) \end{array}
\end{equation*}

\par\medskip\noindent\textbf{Definition} (\texttt{monalg}).\;$\operatorname{monalg}(a, m) :=$
\begin{equation*}
\langle \operatorname{finsupp}(a, m), (\lambda \langle f, g \rangle \in (\operatorname{finsupp}(a, m) \times \operatorname{finsupp}(a, m)).\; \operatorname{monalg\text{-}add}(a, m, f, g)), (\lambda \langle f, g \rangle \in (\operatorname{finsupp}(a, m) \times \operatorname{finsupp}(a, m)).\; \operatorname{monalg\text{-}mul}(a, m, f, g)), (\lambda f \in \operatorname{finsupp}(a, m).\; \operatorname{monalg\text{-}neg}(a, m, f)), \operatorname{monalg\text{-}zero}(a, m), \operatorname{monalg\text{-}one}(a, m) \rangle
\end{equation*}

\par\medskip\noindent\textbf{Definition} (\texttt{poly}).\;$\operatorname{poly}(a) :=$ $\operatorname{monalg}(a, \operatorname{nn\text{-}add\text{-}monoid})$

\vnbbookfile{structure-library/binomial}{structure-library/binomial}
\noindent{\small MOVED 2026-09-20 (batch 12-A) from theorem-library/: this file is VOCABULARY -- definitions, notation and warranted supports, not one proof -- and every theorem stated with it had to load below it.  Its load.scm slot is unchanged. RETIRED 2026-09-17 (proven): bt-mul-comm -- theorem-library/rake-combinatorics.scm}\par
\par\medskip\noindent\textbf{Definition} (\texttt{comb-kk}).\;For all $r$, $x$ and $y$, 
\begin{equation*}
(\operatorname{comb\text{-}kk}(r, x, y, 0) \simeq (\lambda k \in \mathbb{Z}.\; \left(\text{if }(k = 0)\text{ then }\operatorname{one}(r)\text{ else }\operatorname{zero}(r)\right)))
\end{equation*}


\noindent\emph{(Here $\operatorname{comb\text{-}kk}(r, x, y, 0)$ is the internal representation of the binomial coefficient term; it is used throughout the proof.)}
\par
Let $r$, $x$ and $y$ be arbitrary. Let $m \in \mathbb{N}$. Then
\begin{equation*}
(\operatorname{comb\text{-}kk}(r, x, y, \operatorname{succ}(m)) \simeq (\lambda k \in \mathbb{Z}.\; \operatorname{add}(r)(\operatorname{mul}(r)(x, \operatorname{comb\text{-}kk}(r, x, y, m)(-(k, 1))), \operatorname{mul}(r)(y, \operatorname{comb\text{-}kk}(r, x, y, m)(k)))))
\end{equation*}


\noindent\emph{(Here $\operatorname{comb\text{-}kk}(r, x, y, \operatorname{succ}(m))$ is the internal representation of the binomial coefficient term; it is used throughout the proof.)}


\vnbbookfile{theorem-library/poly-membership}{theorem-library/poly-membership}
\noindent{\small poly-membership.scm -- SUPP / FINSUPP / POLY read-offs, PROVEN.}\par
\vnbprop{supp-unfold}{\vnbstmt{\forall a, m, f.\; (\operatorname{supp}(a, m, f) \simeq \{x \in \operatorname{carr}(m) : \neg (f(x) = \operatorname{zero}(a))\})}}
\vnbprop{supp-membership}{\vnbstmt{\begin{array}{@{}l@{}} \forall a, m, f, x. \\ \quad (x \in \operatorname{supp}(a, m, f)) \Leftrightarrow ((x \in \operatorname{carr}(m)) \wedge \neg (f(x) = \operatorname{zero}(a))) \end{array}}}
\vnbprop{supp-in-set}{\vnbstmt{\forall a, m, f.\; ((\operatorname{carr}(m) \in \mathbf{Set}) \Rightarrow (\operatorname{supp}(a, m, f) \in \mathbf{Set}))}}
\vnbprop{finsupp-unfold}{\vnbstmt{\begin{array}{@{}l@{}} \forall a, m. \\ \quad (\operatorname{finsupp}(a, m) \simeq \{f \in (\operatorname{carr}(m) \to \operatorname{carr}(a)) : (|\operatorname{supp}(a, m, f)| \in \mathbb{N})\}) \end{array}}}
\vnbprop{finsupp-membership}{\vnbstmt{\begin{array}{@{}l@{}} \forall a, m, f. \\ \quad (f \in \operatorname{finsupp}(a, m)) \Leftrightarrow ((f \in (\operatorname{carr}(m) \to \operatorname{carr}(a))) \wedge (|\operatorname{supp}(a, m, f)| \in \mathbb{N})) \end{array}}}
\vnbprop{poly-carrier}{\vnbstmt{\forall a.\; (\operatorname{is\text{-}ring}(a) \Rightarrow (\operatorname{carr}(\operatorname{poly}(a)) \simeq \operatorname{finsupp}(a, \operatorname{nn\text{-}add\text{-}monoid})))}}
\vnbprop{poly-membership}{\vnbstmt{\begin{array}{@{}l@{}} \forall a, f. \\ \quad \operatorname{is\text{-}ring}(a) \Rightarrow  \\ \quad \quad (f \in \operatorname{carr}(\operatorname{poly}(a))) \Leftrightarrow ((f \in \operatorname{carr}(a)^{\mathbb{N}}) \wedge (|\operatorname{supp}(a, \operatorname{nn\text{-}add\text{-}monoid}, f)| \in \mathbb{N})) \end{array}}}

\vnbbookfile{theorem-library/ag-view-read-offs}{theorem-library/ag-view-read-offs}
\noindent{\small ag-view-read-offs.scm -- THE ADDITIVE-GROUP VIEW READ-OFFS, PROVEN.}\par
\vnbprop{ras-carr}{\vnbstmt{\forall a.\; (\operatorname{is\text{-}ring}(a) \Rightarrow (\operatorname{carr}(\operatorname{ring\text{-}additive\text{-}ag}(a)) = \operatorname{carr}(a)))}}
\vnbprop{ras-op}{\vnbstmt{\forall a.\; (\operatorname{is\text{-}ring}(a) \Rightarrow (\operatorname{opr}(\operatorname{ring\text{-}additive\text{-}ag}(a)) = \operatorname{add}(a)))}}
\vnbprop{ras-id}{\vnbstmt{\forall a.\; (\operatorname{is\text{-}ring}(a) \Rightarrow (\operatorname{iden}(\operatorname{ring\text{-}additive\text{-}ag}(a)) = \operatorname{zero}(a)))}}

\vnbbookfile{theorem-library/rake-setoid}{theorem-library/rake-setoid}
\noindent{\small rake-setoid.scm -- BATCH F of the 2026-09-17 rake: setoids and carriers. Seven asserted leaves of the PSS, PROVEN.  Every statement here is its support's statement UNCHANGED.}\par
\vnbprop{class-unfold}{\vnbstmt{\forall s, a.\; (\operatorname{class}(s, a) \simeq \{b \in \operatorname{pts}(s) : \operatorname{related}(s, a, b)\})}}
\vnbprop{proj-unfold}{\vnbstmt{\forall s.\; (\operatorname{proj}(s) \simeq (\lambda a \in \operatorname{pts}(s).\; \operatorname{class}(s, a)))}}
\vnbprop{quotient-unfold}{\vnbstmt{\forall s.\; (\operatorname{quotient}(s) \simeq \operatorname{image}(\operatorname{proj}(s), \operatorname{pts}(s)))}}
\vnbprop{respects-unfold}{\vnbstmt{\begin{array}{@{}l@{}} \forall s, f. \\ \quad \operatorname{respects}(s, f) \Leftrightarrow \forall a, b \in \operatorname{pts}(s).\; (\operatorname{related}(s, a, b) \Rightarrow (f(a) = f(b))) \end{array}}}
\vnbprop{respects2-unfold}{\vnbstmt{\begin{array}{@{}l@{}} \forall s, f. \\ \quad \operatorname{respects2}(s, f) \\ \quad \Leftrightarrow\; \forall a, b, c, d. \\ \quad \quad a \in \operatorname{pts}(s) \wedge \\ \quad \quad b \in \operatorname{pts}(s) \wedge \\ \quad \quad c \in \operatorname{pts}(s) \wedge \\ \quad \quad d \in \operatorname{pts}(s) \wedge \\ \quad \quad \operatorname{related}(s, a, c) \wedge \\ \quad \quad \operatorname{related}(s, b, d) \Rightarrow \\ \quad \quad \quad f(a, b) = f(c, d) \end{array}}}
\vnbprop{class-is-set}{\vnbstmt{\forall s.; \forall a \in \operatorname{pts}(s).\; (\operatorname{is\text{-}setoid}(s) \Rightarrow (\operatorname{class}(s, a) \in \mathbf{Set}))}}
\vnbprop{quotient-is-set}{\vnbstmt{\forall s.\; (\operatorname{is\text{-}setoid}(s) \Rightarrow (\operatorname{quotient}(s) \in \mathbf{Set}))}}
\vnbprop{proj-in-fun}{\vnbstmt{\forall s.\; (\operatorname{is\text{-}setoid}(s) \Rightarrow (\operatorname{proj}(s) \in (\operatorname{pts}(s) \to \operatorname{quotient}(s))))}}
\vnbprop{descend-in-fun}{\vnbstmt{\begin{array}{@{}l@{}} \forall s, z.; \forall f \in (\operatorname{pts}(s) \to z). \\ \quad (\operatorname{is\text{-}setoid}(s) \wedge \operatorname{respects}(s, f)) \Rightarrow (\operatorname{descend}(s, f) \in (\operatorname{quotient}(s) \to z)) \end{array}}}
\vnbprop{descend2-in-fun}{\vnbstmt{\begin{array}{@{}l@{}} \forall s, z.; \forall f \in ((\operatorname{pts}(s) \times \operatorname{pts}(s)) \to z). \\ \quad (\operatorname{is\text{-}setoid}(s) \wedge \operatorname{respects2}(s, f)) \Rightarrow  \\ \quad \quad \operatorname{descend2}(s, f) \in ((\operatorname{quotient}(s) \times \operatorname{quotient}(s)) \to z) \end{array}}}

\vnbbookfile{theorem-library/bt-shims}{theorem-library/bt-shims}
\noindent{\small bt-shims.scm -- the five "Bernstein typing" shims of theorem-library/binomial.scm, PROVEN.}\par
\vnbprop{bt-nn-in-zz}{\vnbstmt{\forall n \in \mathbb{N}.\; (n \in \mathbb{Z})}}
\vnbprop{bt-succ-in-nn}{\vnbstmt{\forall n \in \mathbb{N}.\; ((n + 1) \in \mathbb{N})}}
\vnbprop{bt-neg1-in-zz}{\vnbstmt{-1 \in \mathbb{Z}}}
\vnbprop{bt-lt-succ}{\vnbstmt{\forall n \in \mathbb{N}.\; (n < (n + 1))}}
\vnbprop{bt-neg1-neg}{\vnbstmt{-1 < 0}}

\vnbbookfile{structure-library/subtype-laws}{structure-library/subtype-laws}
\noindent{\small subtype-laws.scm -- the trivial subtype-subsumption laws, PROVEN.}\par
\vnbprop{abelian-group-is-group}{\vnbstmt{\forall s.\; (\operatorname{is\text{-}abelian\text{-}group}(s) \Rightarrow \operatorname{is\text{-}group}(s))}}
\vnbprop{commutative-ring-is-ring}{\vnbstmt{\forall s.\; (\operatorname{is\text{-}commutative\text{-}ring}(s) \Rightarrow \operatorname{is\text{-}ring}(s))}}
\vnbprop{integral-domain-is-commutative-ring}{\vnbstmt{\forall s.\; (\operatorname{is\text{-}integral\text{-}domain}(s) \Rightarrow \operatorname{is\text{-}commutative\text{-}ring}(s))}}
\vnbprop{euclidean-ring-is-integral-domain}{\vnbstmt{\forall s.\; (\operatorname{is\text{-}euclidean\text{-}ring}(s) \Rightarrow \operatorname{is\text{-}integral\text{-}domain}(s))}}
\vnbprop{comm-monoid-is-monoid}{\vnbstmt{\forall s.\; (\operatorname{is\text{-}comm\text{-}monoid}(s) \Rightarrow \operatorname{is\text{-}monoid}(s))}}
\vnbprop{abelian-group-opr-comm}{\vnbstmt{\begin{array}{@{}l@{}} \forall s.; \forall a, b \in \operatorname{carr}(s). \\ \quad \operatorname{is\text{-}abelian\text{-}group}(s) \Rightarrow (\operatorname{opr}(s)(a, b) = \operatorname{opr}(s)(b, a)) \end{array}}}
\vnbprop{group-assoc}{\vnbstmt{\begin{array}{@{}l@{}} \forall s.; \forall a, b, c \in \operatorname{carr}(s). \\ \quad \operatorname{is\text{-}group}(s) \Rightarrow (\operatorname{opr}(s)(\operatorname{opr}(s)(a, b), c) = \operatorname{opr}(s)(a, \operatorname{opr}(s)(b, c))) \end{array}}}
\vnbprop{group-left-id}{\vnbstmt{\forall s.; \forall a \in \operatorname{carr}(s).\; (\operatorname{is\text{-}group}(s) \Rightarrow (\operatorname{opr}(s)(\operatorname{iden}(s), a) = a))}}
\vnbprop{group-left-inv}{\vnbstmt{\begin{array}{@{}l@{}} \forall s.; \forall a \in \operatorname{carr}(s). \\ \quad \operatorname{is\text{-}group}(s) \Rightarrow (\operatorname{opr}(s)(\operatorname{inv}(s)(a), a) = \operatorname{iden}(s)) \end{array}}}
\vnbprop{group-identity-in}{\vnbstmt{\forall s.\; (\operatorname{is\text{-}group}(s) \Rightarrow (\operatorname{iden}(s) \in \operatorname{carr}(s)))}}

\vnbbookfile{theorem-library/op-typing}{theorem-library/op-typing}
\noindent{\small A file shared across books; the entries that belong to this one.}\par
\noindent{\small op-typing.scm -- the codomain typing of a structure operation, in APPLIED form, PROVEN.  One driver, seven theorems.}\par
\vnbprop{ring-add-closed}{\vnbstmt{\begin{array}{@{}l@{}} \forall s.; \forall a, b \in \operatorname{carr}(s). \\ \quad \operatorname{is\text{-}ring}(s) \Rightarrow (\operatorname{add}(s)(a, b) \in \operatorname{carr}(s)) \end{array}}}
\vnbprop{ring-carrier-closed-mul}{\vnbstmt{\begin{array}{@{}l@{}} \forall r.; \forall a, b \in \operatorname{carr}(r). \\ \quad \operatorname{is\text{-}ring}(r) \Rightarrow (\operatorname{mul}(r)(a, b) \in \operatorname{carr}(r)) \end{array}}}
\vnbprop{ring-neg-in-carr}{\vnbstmt{\forall s.; \forall a \in \operatorname{carr}(s).\; (\operatorname{is\text{-}ring}(s) \Rightarrow (\operatorname{neg}(s)(a) \in \operatorname{carr}(s)))}}
\vnbprop{bt-add-in-carr}{\vnbstmt{\begin{array}{@{}l@{}} \forall r.; \forall a, b \in \operatorname{carr}(r). \\ \quad \operatorname{is\text{-}commutative\text{-}ring}(r) \Rightarrow (\operatorname{add}(r)(a, b) \in \operatorname{carr}(r)) \end{array}}}

\vnbbookfile{theorem-library/finsum-type-proof}{theorem-library/finsum-type-proof}
\noindent{\small theorem-library/finsum-type-proof.scm -- the FINSUM typing floor, PROVEN.}\par
\vnbprop{group-carrier-closed-opr}{\vnbstmt{\begin{array}{@{}l@{}} \forall s.; \forall a, b \in \operatorname{carr}(s). \\ \quad \operatorname{is\text{-}group}(s) \Rightarrow (\operatorname{opr}(s)(a, b) \in \operatorname{carr}(s)) \end{array}}}
\vnbprop{sum-ag-type-ind}{\vnbstmt{\begin{array}{@{}l@{}} \forall n \in \mathbb{N}, g.; \forall f \in (\mathbb{N} \to \operatorname{carr}(g)). \\ \quad \operatorname{is\text{-}abelian\text{-}group}(g) \Rightarrow (\operatorname{sum\text{-}ag}(g, f, n) \in \operatorname{carr}(g)) \end{array}}}
\vnbprop{sum-ag-type}{\vnbstmt{\begin{array}{@{}l@{}} \forall g.; \forall f \in (\mathbb{N} \to \operatorname{carr}(g)), n \in \mathbb{N}. \\ \quad \operatorname{is\text{-}abelian\text{-}group}(g) \Rightarrow (\operatorname{sum\text{-}ag}(g, f, n) \in \operatorname{carr}(g)) \end{array}}}
\vnbprop{enum-fam-in-fun}{\vnbstmt{\begin{array}{@{}l@{}} \forall n \in \mathbb{N}, g, s.; \forall i \in (\mathrm{OS}(n) \to s), f \in (s \to \operatorname{carr}(g)). \\ \quad \operatorname{is\text{-}group}(g) \Rightarrow (\operatorname{enum\text{-}fam}(g, f, i, n) \in (\mathbb{N} \to \operatorname{carr}(g))) \end{array}}}
\vnbprop{finsum-type}{\vnbstmt{\begin{array}{@{}l@{}} \forall g.; \forall s \in \mathbf{Set}, f \in (s \to \operatorname{carr}(g)). \\ \quad (\operatorname{is\text{-}abelian\text{-}group}(g) \wedge (|s| \in \mathbb{N})) \Rightarrow (\operatorname{finsum}(g, f, s) \in \operatorname{carr}(g)) \end{array}}}
\vnbprop{sum-ag-all-id-ind}{\vnbstmt{\begin{array}{@{}l@{}} \forall n \in \mathbb{N}, g, a. \\ \quad (\operatorname{is\text{-}abelian\text{-}group}(g) \wedge \forall i \in \mathrm{OS}(n).\; (a(i) = \operatorname{iden}(g))) \Rightarrow  \\ \quad \quad \operatorname{sum\text{-}ag}(g, a, n) = \operatorname{iden}(g) \end{array}}}
\vnbprop{finsum-all-id}{\vnbstmt{\begin{array}{@{}l@{}} \forall g.; \forall s \in \mathbf{Set}, f \in (s \to \operatorname{carr}(g)). \\ \quad (\operatorname{is\text{-}abelian\text{-}group}(g) \wedge ((|s| \in \mathbb{N}) \wedge \forall z \in s.\; (f(z) = \operatorname{iden}(g)))) \Rightarrow  \\ \quad \quad \operatorname{finsum}(g, f, s) = \operatorname{iden}(g) \end{array}}}

\vnbbookfile{theorem-library/finsum-single-support}{theorem-library/finsum-single-support}
\noindent{\small finsum-single-support.scm -- a finite sum whose summand vanishes off a single index equals its value there.  Both theorems `modulo 0'.}\par
\vnbprop{sum-ag-single-ind}{\vnbstmt{\begin{array}{@{}l@{}} \forall n \in \mathbb{N}, g, a.; \forall p \in \mathrm{OS}(n). \\ \quad \operatorname{is\text{-}abelian\text{-}group}(g) \wedge \\ \quad a(p) \in \operatorname{carr}(g) \wedge \\ \quad \forall i \in \mathrm{OS}(n).\; (\neg (i = p) \Rightarrow (a(i) = \operatorname{iden}(g))) \Rightarrow \\ \quad \quad \operatorname{sum\text{-}ag}(g, a, n) = a(p) \end{array}}}
\vnbprop{finsum-single-support}{\vnbstmt{\begin{array}{@{}l@{}} \forall g.; \forall s \in \mathbf{Set}, f \in (s \to \operatorname{carr}(g)), a \in s. \\ \quad (\operatorname{is\text{-}abelian\text{-}group}(g) \wedge ((|s| \in \mathbb{N}) \wedge \forall j \in s.\; (\neg (j = a) \Rightarrow (f(j) = \operatorname{iden}(g))))) \Rightarrow  \\ \quad \quad \operatorname{finsum}(g, f, s) = f(a) \end{array}}}

\vnbbookfile{structure-library/integral-domain-laws}{structure-library/integral-domain-laws}
\noindent{\small integral-domain-laws.scm -- integral-domain-cancel-zero, PROVEN from the definition.}\par
\vnbprop{integral-domain-cancel-zero}{\vnbstmt{\begin{array}{@{}l@{}} \forall s.; \forall a, b \in \operatorname{carr}(s). \\ \quad (\operatorname{is\text{-}integral\text{-}domain}(s) \wedge ((\operatorname{mul}(s)(a, b) = \operatorname{zero}(s)) \wedge \neg (b = \operatorname{zero}(s)))) \Rightarrow  \\ \quad \quad a = \operatorname{zero}(s) \end{array}}}
\vnbprop{integral-domain-nontrivial}{\vnbstmt{\forall s.\; (\operatorname{is\text{-}integral\text{-}domain}(s) \Rightarrow \neg (\operatorname{one}(s) = \operatorname{zero}(s)))}}

\vnbbookfile{theorem-library/zz-ring-is-ring}{theorem-library/zz-ring-is-ring}
\noindent{\small zz-ring-is-ring.scm -- IS-RING(ZZ-RING), PROVED.}\par
\vnbprop{zz-is-ring}{\vnbstmt{\operatorname{is\text{-}ring}(\operatorname{zz\text{-}ring})}}
\vnbprop{qq-is-ring}{\vnbstmt{\operatorname{is\text{-}ring}(\operatorname{qq\text{-}ring})}}

\vnbbookfile{theorem-library/zz-integral-domain}{theorem-library/zz-integral-domain}
\noindent{\small zz-integral-domain.scm -- IS-COMMUTATIVE-RING(ZZ-RING) and IS-INTEGRAL-DOMAIN(ZZ-RING), PROVED.}\par
\vnbprop{zz-is-commutative-ring}{\vnbstmt{\operatorname{is\text{-}commutative\text{-}ring}(\operatorname{zz\text{-}ring})}}
\vnbprop{zz-no-zero-divisors}{\vnbstmt{\forall a, b \in \mathbb{Z}.\; (((a \cdot b) = 0) \Rightarrow ((a = 0) \vee (b = 0)))}}
\vnbprop{zz-is-integral-domain}{\vnbstmt{\operatorname{is\text{-}integral\text{-}domain}(\operatorname{zz\text{-}ring})}}

\vnbbookfile{theorem-library/nn-add-monoid}{theorem-library/nn-add-monoid}
\noindent{\small nn-add-monoid.scm -- IS-COMM-MONOID(NN-ADD-MONOID), PROVED.}\par
\vnbprop{nn-add-monoid-is-comm-monoid}{\vnbstmt{\operatorname{is\text{-}comm\text{-}monoid}(\operatorname{nn\text{-}add\text{-}monoid})}}

\vnbbookfile{theorem-library/cancellation}{theorem-library/cancellation}
\noindent{\small cancellation.scm -- cancellation, proved ONCE where it lives, then carried.}\par
\vnbprop{group-inv-in}{\vnbstmt{\forall s.; \forall a \in \operatorname{carr}(s).\; (\operatorname{is\text{-}group}(s) \Rightarrow (\operatorname{inv}(s)(a) \in \operatorname{carr}(s)))}}
\vnbprop{group-cancel-left}{\vnbstmt{\begin{array}{@{}l@{}} \forall s.; \forall a, b, c \in \operatorname{carr}(s). \\ \quad (\operatorname{is\text{-}group}(s) \wedge (\operatorname{opr}(s)(c, a) = \operatorname{opr}(s)(c, b))) \Rightarrow (a = b) \end{array}}}
\vnbprop{ag-cancel-right}{\vnbstmt{\begin{array}{@{}l@{}} \forall s.; \forall a, b, c \in \operatorname{carr}(s). \\ \quad (\operatorname{is\text{-}abelian\text{-}group}(s) \wedge (\operatorname{opr}(s)(a, c) = \operatorname{opr}(s)(b, c))) \Rightarrow (a = b) \end{array}}}
\vnbprop{abelian-group-idempotent-is-id}{\vnbstmt{\begin{array}{@{}l@{}} \forall s.; \forall a \in \operatorname{carr}(s). \\ \quad (\operatorname{is\text{-}abelian\text{-}group}(s) \wedge (\operatorname{opr}(s)(a, a) = a)) \Rightarrow (a = \operatorname{iden}(s)) \end{array}}}
\vnbprop{zz-add-cancel}{\vnbstmt{\forall a, b, c \in \mathbb{Z}.\; (((a + c) = (b + c)) \Rightarrow (a = b))}}
\vnbprop{qq-add-cancel}{\vnbstmt{\forall a, b, c \in \mathbb{Q}.\; (((a + c) = (b + c)) \Rightarrow (a = b))}}

\vnbbookfile{theorem-library/ring-zero-one-power}{theorem-library/ring-zero-one-power}
\noindent{\small ring-zero-one-power.scm -- four ring facts that were asserted, PROVEN.}\par
\vnbprop{ring-one-in}{\vnbstmt{\forall s.\; (\operatorname{is\text{-}ring}(s) \Rightarrow (\operatorname{one}(s) \in \operatorname{carr}(s)))}}
\vnbprop{ring-mul-zero-left}{\vnbstmt{\forall s.; \forall a \in \operatorname{carr}(s).\; (\operatorname{is\text{-}ring}(s) \Rightarrow (\operatorname{mul}(s)(\operatorname{zero}(s), a) = \operatorname{zero}(s)))}}
\vnbprop{ring-mul-zero-right}{\vnbstmt{\forall s.; \forall a \in \operatorname{carr}(s).\; (\operatorname{is\text{-}ring}(s) \Rightarrow (\operatorname{mul}(s)(a, \operatorname{zero}(s)) = \operatorname{zero}(s)))}}
\vnbprop{ring-power-zero}{\vnbstmt{\forall r, x.\; (\operatorname{is\text{-}ring}(r) \Rightarrow (\operatorname{ring\text{-}power}(r, x, 0) \simeq \operatorname{one}(r)))}}

\vnbbookfile{theorem-library/nn-integral}{theorem-library/nn-integral}
\noindent{\small nn-integral.scm -- the NN arithmetic support layer.}\par
\vnbprop{zz-mul-cancel-zero}{\vnbstmt{\forall a, b \in \mathbb{Z}.\; ((((a \cdot b) = 0) \wedge \neg (b = 0)) \Rightarrow (a = 0))}}
\vnbprop{nn-mul-nonzero}{\vnbstmt{\forall a, b \in \mathbb{N}.\; ((\neg (a = 0) \wedge \neg (b = 0)) \Rightarrow \neg ((a \cdot b) = 0))}}
\vnbprop{nn-2-cancel}{\vnbstmt{\forall x, y \in \mathbb{N}.\; (((2 \cdot x) = (2 \cdot y)) \Rightarrow (x = y))}}
\vnbprop{nn-lt-double}{\vnbstmt{\forall k \in \mathbb{N}.\; (\neg (k = 0) \Rightarrow (k < (2 \cdot k)))}}

\vnbbookfile{theorem-library/comb-kk-laws}{theorem-library/comb-kk-laws}
\noindent{\small comb-kk-laws.scm -- the six COMB-KK / Bernstein plumbing supports of theorem-library/binomial.scm, PROVEN.}\par
\vnbprop{sum-in-carr-zz-ind}{\vnbstmt{\begin{array}{@{}l@{}} \forall n \in \mathbb{N}, r.; \forall g \in (\mathbb{Z} \to \operatorname{carr}(r)). \\ \quad \operatorname{is\text{-}commutative\text{-}ring}(r) \Rightarrow (\operatorname{sum}(r, g, n) \in \operatorname{carr}(r)) \end{array}}}
\vnbprop{bt-sum-in-carr-zz}{\vnbstmt{\begin{array}{@{}l@{}} \forall r.; \forall g \in (\mathbb{Z} \to \operatorname{carr}(r)), n \in \mathbb{N}. \\ \quad \operatorname{is\text{-}commutative\text{-}ring}(r) \Rightarrow (\operatorname{sum}(r, g, n) \in \operatorname{carr}(r)) \end{array}}}
\vnbprop{comb-kk-0-0}{\vnbstmt{\begin{array}{@{}l@{}} \forall r.; \forall x, y \in \operatorname{carr}(r). \\ \quad \operatorname{is\text{-}commutative\text{-}ring}(r) \Rightarrow (\operatorname{comb\text{-}kk}(r, x, y, 0)(0) = \operatorname{one}(r)) \end{array}}}
\vnbprop{comb-kk-in-fun-ind}{\vnbstmt{\begin{array}{@{}l@{}} \forall m \in \mathbb{N}, r.; \forall x, y \in \operatorname{carr}(r). \\ \quad \operatorname{is\text{-}commutative\text{-}ring}(r) \Rightarrow (\operatorname{comb\text{-}kk}(r, x, y, m) \in (\mathbb{Z} \to \operatorname{carr}(r))) \end{array}}}
\vnbprop{comb-kk-in-fun}{\vnbstmt{\begin{array}{@{}l@{}} \forall r.; \forall x, y \in \operatorname{carr}(r), m \in \mathbb{N}. \\ \quad \operatorname{is\text{-}commutative\text{-}ring}(r) \Rightarrow (\operatorname{comb\text{-}kk}(r, x, y, m) \in (\mathbb{Z} \to \operatorname{carr}(r))) \end{array}}}
\vnbprop{comb-kk-null-ind}{\vnbstmt{\begin{array}{@{}l@{}} \forall m \in \mathbb{N}, r.; \forall x, y \in \operatorname{carr}(r), k \in \mathbb{Z}. \\ \quad (\operatorname{is\text{-}commutative\text{-}ring}(r) \wedge (k < 0)) \Rightarrow (\operatorname{comb\text{-}kk}(r, x, y, m)(k) = \operatorname{zero}(r)) \end{array}}}
\vnbprop{comb-kk-null}{\vnbstmt{\begin{array}{@{}l@{}} \forall r.; \forall x, y \in \operatorname{carr}(r), m \in \mathbb{N}, k \in \mathbb{Z}. \\ \quad (\operatorname{is\text{-}commutative\text{-}ring}(r) \wedge (k < 0)) \Rightarrow (\operatorname{comb\text{-}kk}(r, x, y, m)(k) = \operatorname{zero}(r)) \end{array}}}
\vnbprop{comb-kk-above-ind}{\vnbstmt{\begin{array}{@{}l@{}} \forall m \in \mathbb{N}, r.; \forall x, y \in \operatorname{carr}(r), k \in \mathbb{Z}. \\ \quad (\operatorname{is\text{-}commutative\text{-}ring}(r) \wedge (m < k)) \Rightarrow (\operatorname{comb\text{-}kk}(r, x, y, m)(k) = \operatorname{zero}(r)) \end{array}}}
\vnbprop{comb-kk-above}{\vnbstmt{\begin{array}{@{}l@{}} \forall r.; \forall x, y \in \operatorname{carr}(r), m \in \mathbb{N}, k \in \mathbb{Z}. \\ \quad (\operatorname{is\text{-}commutative\text{-}ring}(r) \wedge (m < k)) \Rightarrow (\operatorname{comb\text{-}kk}(r, x, y, m)(k) = \operatorname{zero}(r)) \end{array}}}

\vnbbookfile{theorem-library/zz-division}{theorem-library/zz-division}
\noindent{\small zz-division.scm -- DIVISION WITH REMAINDER, and ZZ as a Euclidean ring.}\par
\vnbprop{nn-division}{\vnbstmt{\begin{array}{@{}l@{}} \forall a, b \in \mathbb{N}. \\ \quad \neg (b = 0) \Rightarrow \exists q, r \in \mathbb{N}.\; ((a = ((q \cdot b) + r)) \wedge ((r + 1) \leq b)) \end{array}}}
\vnbprop{zz-division}{\vnbstmt{\begin{array}{@{}l@{}} \forall a, b \in \mathbb{Z}. \\ \quad \neg (b = 0) \Rightarrow  \\ \quad \quad \exists q, r \in \mathbb{Z}.\; ((a = ((q \cdot b) + r)) \wedge ((r = 0) \vee ((\lvert r \rvert + 1) \leq \lvert b \rvert))) \end{array}}}
\vnbprop{zz-is-euclidean-ring}{\vnbstmt{\operatorname{is\text{-}euclidean\text{-}ring}(\operatorname{zz\text{-}ring})}}

\vnbbookfile{theorem-library/rake-algebra}{theorem-library/rake-algebra}
\noindent{\small rake-algebra.scm -- BATCH D of the 2026-09-17 rake: seven asserted leaves of the PSS, PROVEN.  Every statement here is its support's statement UNCHANGED.}\par
\vnbprop{ring-carr-in-set}{\vnbstmt{\forall r.\; (\operatorname{is\text{-}ring}(r) \Rightarrow (\operatorname{carr}(r) \in \mathbf{Set}))}}
\vnbprop{ideal-elt-in-carrier}{\vnbstmt{\forall s, i, x.\; ((\operatorname{is\text{-}ideal}(s, i) \wedge (x \in i)) \Rightarrow (x \in \operatorname{carr}(s)))}}
\vnbprop{principal-ideal-in-ideal}{\vnbstmt{\begin{array}{@{}l@{}} \forall s, i, b.; \forall x \in \operatorname{principal\text{-}ideal}(s, b). \\ \quad (\operatorname{is\text{-}ideal}(s, i) \wedge (b \in i)) \Rightarrow (x \in i) \end{array}}}
\vnbprop{ring-add-right-id}{\vnbstmt{\forall s.; \forall a \in \operatorname{carr}(s).\; (\operatorname{is\text{-}ring}(s) \Rightarrow (\operatorname{add}(s)(a, \operatorname{zero}(s)) = a))}}

\vnbbookfile{theorem-library/rake-algebra2}{theorem-library/rake-algebra2}
\noindent{\small rake-algebra2.scm -- BATCH I of the 2026-09-17 rake: the algebra two-liners and the four unwarranted algebra axioms, PROVEN.  Every statement is its site's statement UNCHANGED (extracted from the site, not retyped).}\par
\vnbprop{abelian-group-assoc}{\vnbstmt{\begin{array}{@{}l@{}} \forall s.; \forall a, b, c \in \operatorname{carr}(s). \\ \quad \operatorname{is\text{-}abelian\text{-}group}(s) \Rightarrow  \\ \quad \quad \operatorname{opr}(s)(\operatorname{opr}(s)(a, b), c) = \operatorname{opr}(s)(a, \operatorname{opr}(s)(b, c)) \end{array}}}
\vnbprop{abelian-group-right-id}{\vnbstmt{\forall s.; \forall a \in \operatorname{carr}(s).\; (\operatorname{is\text{-}abelian\text{-}group}(s) \Rightarrow (\operatorname{opr}(s)(a, \operatorname{iden}(s)) = a))}}
\vnbprop{abelian-group-inverse-unique}{\vnbstmt{\begin{array}{@{}l@{}} \forall s.; \forall a, b \in \operatorname{carr}(s). \\ \quad (\operatorname{is\text{-}abelian\text{-}group}(s) \wedge (\operatorname{opr}(s)(a, b) = \operatorname{iden}(s))) \Rightarrow (b = \operatorname{inv}(s)(a)) \end{array}}}
\vnbprop{ring-add-right-inv}{\vnbstmt{\forall s.; \forall a \in \operatorname{carr}(s).\; (\operatorname{is\text{-}ring}(s) \Rightarrow (\operatorname{add}(s)(a, \operatorname{neg}(s)(a)) = \operatorname{zero}(s)))}}
\vnbprop{ring-add-inverse-unique}{\vnbstmt{\begin{array}{@{}l@{}} \forall s.; \forall a, b \in \operatorname{carr}(s). \\ \quad (\operatorname{is\text{-}ring}(s) \wedge (\operatorname{add}(s)(a, b) = \operatorname{zero}(s))) \Rightarrow (b = \operatorname{neg}(s)(a)) \end{array}}}
\vnbprop{ring-neg-neg}{\vnbstmt{\forall s.; \forall r \in \operatorname{carr}(s).\; (\operatorname{is\text{-}ring}(s) \Rightarrow (\operatorname{neg}(s)(\operatorname{neg}(s)(r)) = r))}}
\vnbprop{ring-neg-mul-left}{\vnbstmt{\begin{array}{@{}l@{}} \forall s.; \forall a, b \in \operatorname{carr}(s). \\ \quad \operatorname{is\text{-}ring}(s) \Rightarrow (\operatorname{mul}(s)(\operatorname{neg}(s)(a), b) = \operatorname{neg}(s)(\operatorname{mul}(s)(a, b))) \end{array}}}
\vnbprop{matrix-sethood}{\vnbstmt{\forall s \in \mathbf{Set}.\; (\operatorname{matrix}(s) \in \mathbf{Set})}}
\vnbprop{monoid-identity-in}{\vnbstmt{\forall m.\; (\operatorname{is\text{-}monoid}(m) \Rightarrow (\operatorname{iden}(m) \in \operatorname{carr}(m)))}}
\vnbprop{monoid-carrier-closed-opr}{\vnbstmt{\begin{array}{@{}l@{}} \forall m, a, b. \\ \quad (\operatorname{is\text{-}monoid}(m) \wedge ((a \in \operatorname{carr}(m)) \wedge (b \in \operatorname{carr}(m)))) \Rightarrow (\operatorname{opr}(m)(a, b) \in \operatorname{carr}(m)) \end{array}}}
\vnbprop{mpow-type-ind}{\vnbstmt{\begin{array}{@{}l@{}} \forall n \in \mathbb{N}, m.; \forall x \in \operatorname{carr}(m). \\ \quad \operatorname{is\text{-}monoid}(m) \Rightarrow (\operatorname{mpow}(m, x, n) \in \operatorname{carr}(m)) \end{array}}}
\vnbprop{mpow-type}{\vnbstmt{\begin{array}{@{}l@{}} \forall m.; \forall x \in \operatorname{carr}(m), n \in \mathbb{N}. \\ \quad \operatorname{is\text{-}monoid}(m) \Rightarrow (\operatorname{mpow}(m, x, n) \in \operatorname{carr}(m)) \end{array}}}

\vnbbookfile{theorem-library/rake-monoid-laws}{theorem-library/rake-monoid-laws}
\noindent{\small rake-monoid-laws.scm -- BATCH 5c-P of the 2026-09-18 rake: the three MONOID shape projections, PROVEN.}\par
\vnbprop{monoid-assoc}{\vnbstmt{\begin{array}{@{}l@{}} \forall s.; \forall a, b, c \in \operatorname{carr}(s). \\ \quad \operatorname{is\text{-}monoid}(s) \Rightarrow (\operatorname{opr}(s)(\operatorname{opr}(s)(a, b), c) = \operatorname{opr}(s)(a, \operatorname{opr}(s)(b, c))) \end{array}}}
\vnbprop{monoid-left-id}{\vnbstmt{\forall s.; \forall a \in \operatorname{carr}(s).\; (\operatorname{is\text{-}monoid}(s) \Rightarrow (\operatorname{opr}(s)(\operatorname{iden}(s), a) = a))}}
\vnbprop{monoid-right-id}{\vnbstmt{\forall s.; \forall a \in \operatorname{carr}(s).\; (\operatorname{is\text{-}monoid}(s) \Rightarrow (\operatorname{opr}(s)(a, \operatorname{iden}(s)) = a))}}

\vnbbookfile{theorem-library/rake-zz-act}{theorem-library/rake-zz-act}
\noindent{\small rake-zz-act.scm -- BATCH 6 of the rake: ZZ-ACT made a DEFINITION, and the five billed ZZ-ACT laws proven from it, each `modulo 0'.}\par
\vnbprop{ag-as-monoid-carr}{\vnbstmt{\begin{array}{@{}l@{}} \forall g. \\ \quad \operatorname{is\text{-}abelian\text{-}group}(g) \Rightarrow (\operatorname{carr}(g) \simeq \operatorname{carr}(\operatorname{abelian\text{-}group\text{-}as\text{-}monoid}(g))) \end{array}}}
\vnbprop{ag-as-monoid-iden}{\vnbstmt{\begin{array}{@{}l@{}} \forall g. \\ \quad \operatorname{is\text{-}abelian\text{-}group}(g) \Rightarrow (\operatorname{iden}(g) \simeq \operatorname{iden}(\operatorname{abelian\text{-}group\text{-}as\text{-}monoid}(g))) \end{array}}}
\vnbprop{ag-as-monoid-opr}{\vnbstmt{\begin{array}{@{}l@{}} \forall g. \\ \quad \operatorname{is\text{-}abelian\text{-}group}(g) \Rightarrow (\operatorname{opr}(g) \simeq \operatorname{opr}(\operatorname{abelian\text{-}group\text{-}as\text{-}monoid}(g))) \end{array}}}
\vnbprop{abelian-group-right-inv}{\vnbstmt{\begin{array}{@{}l@{}} \forall s.; \forall a \in \operatorname{carr}(s). \\ \quad \operatorname{is\text{-}abelian\text{-}group}(s) \Rightarrow (\operatorname{opr}(s)(a, \operatorname{inv}(s)(a)) = \operatorname{iden}(s)) \end{array}}}
\vnbprop{abelian-group-inv-iden}{\vnbstmt{\forall g.\; (\operatorname{is\text{-}abelian\text{-}group}(g) \Rightarrow (\operatorname{inv}(g)(\operatorname{iden}(g)) = \operatorname{iden}(g)))}}
\vnbprop{abelian-group-inv-inv}{\vnbstmt{\forall s.; \forall a \in \operatorname{carr}(s).\; (\operatorname{is\text{-}abelian\text{-}group}(s) \Rightarrow (\operatorname{inv}(s)(\operatorname{inv}(s)(a)) = a))}}
\vnbprop{abelian-group-inv-opr}{\vnbstmt{\begin{array}{@{}l@{}} \forall s.; \forall a, b \in \operatorname{carr}(s). \\ \quad \operatorname{is\text{-}abelian\text{-}group}(s) \Rightarrow  \\ \quad \quad \operatorname{inv}(s)(\operatorname{opr}(s)(a, b)) = \operatorname{opr}(s)(\operatorname{inv}(s)(a), \operatorname{inv}(s)(b)) \end{array}}}
\vnbprop{mpow-one}{\vnbstmt{\forall m.; \forall x \in \operatorname{carr}(m).\; (\operatorname{is\text{-}monoid}(m) \Rightarrow (\operatorname{mpow}(m, x, 1) = x))}}
\vnbprop{mpow-add-ind}{\vnbstmt{\begin{array}{@{}l@{}} \forall j \in \mathbb{N}, m.; \forall x \in \operatorname{carr}(m), k \in \mathbb{N}. \\ \quad \operatorname{is\text{-}monoid}(m) \Rightarrow (\operatorname{mpow}(m, x, (j + k)) = \operatorname{opr}(m)(\operatorname{mpow}(m, x, j), \operatorname{mpow}(m, x, k))) \end{array}}}
\vnbprop{mpow-add}{\vnbstmt{\begin{array}{@{}l@{}} \forall m.; \forall x \in \operatorname{carr}(m), j, k \in \mathbb{N}. \\ \quad \operatorname{is\text{-}monoid}(m) \Rightarrow (\operatorname{mpow}(m, x, (j + k)) = \operatorname{opr}(m)(\operatorname{mpow}(m, x, j), \operatorname{mpow}(m, x, k))) \end{array}}}
\vnbprop{ag-mpow-diff-nonneg}{\vnbstmt{\begin{array}{@{}l@{}} \forall g.; \forall d, n, m \in \mathbb{N}, a \in \operatorname{carr}(g). \\ \quad (\operatorname{is\text{-}abelian\text{-}group}(g) \wedge ((d + n) = m)) \Rightarrow  \\ \quad \quad \operatorname{opr}(g)(\operatorname{mpow}(\operatorname{abelian\text{-}group\text{-}as\text{-}monoid}(g), a, m), \operatorname{inv}(g)(\operatorname{mpow}(\operatorname{abelian\text{-}group\text{-}as\text{-}monoid}(g), a, n))) \\ \quad \quad =\; \operatorname{mpow}(\operatorname{abelian\text{-}group\text{-}as\text{-}monoid}(g), a, d) \end{array}}}
\vnbprop{ag-mpow-diff-neg}{\vnbstmt{\begin{array}{@{}l@{}} \forall g.; \forall p, m, n \in \mathbb{N}, a \in \operatorname{carr}(g). \\ \quad (\operatorname{is\text{-}abelian\text{-}group}(g) \wedge ((p + m) = n)) \Rightarrow  \\ \quad \quad \operatorname{opr}(g)(\operatorname{mpow}(\operatorname{abelian\text{-}group\text{-}as\text{-}monoid}(g), a, m), \operatorname{inv}(g)(\operatorname{mpow}(\operatorname{abelian\text{-}group\text{-}as\text{-}monoid}(g), a, n))) \\ \quad \quad =\; \operatorname{inv}(g)(\operatorname{mpow}(\operatorname{abelian\text{-}group\text{-}as\text{-}monoid}(g), a, p)) \end{array}}}
\vnbprop{zz-act-nonneg}{\vnbstmt{\begin{array}{@{}l@{}} \forall g.; \forall k \in \mathbb{N}, a \in \operatorname{carr}(g). \\ \quad \operatorname{is\text{-}abelian\text{-}group}(g) \Rightarrow  \\ \quad \quad \operatorname{zz\text{-}act}(g, k, a) = \operatorname{mpow}(\operatorname{abelian\text{-}group\text{-}as\text{-}monoid}(g), a, k) \end{array}}}
\vnbprop{zz-act-neg}{\vnbstmt{\begin{array}{@{}l@{}} \forall g.; \forall k \in \mathbb{N}, a \in \operatorname{carr}(g). \\ \quad \operatorname{is\text{-}abelian\text{-}group}(g) \Rightarrow  \\ \quad \quad \operatorname{zz\text{-}act}(g, -k, a) = \operatorname{inv}(g)(\operatorname{mpow}(\operatorname{abelian\text{-}group\text{-}as\text{-}monoid}(g), a, k)) \end{array}}}
\vnbprop{zz-act-zero}{\vnbstmt{\forall g.; \forall a \in \operatorname{carr}(g).\; (\operatorname{is\text{-}abelian\text{-}group}(g) \Rightarrow (\operatorname{zz\text{-}act}(g, 0, a) = \operatorname{iden}(g)))}}
\vnbprop{zz-act-one}{\vnbstmt{\forall g.; \forall a \in \operatorname{carr}(g).\; (\operatorname{is\text{-}abelian\text{-}group}(g) \Rightarrow (\operatorname{zz\text{-}act}(g, 1, a) = a))}}
\vnbprop{zz-act-type}{\vnbstmt{\begin{array}{@{}l@{}} \forall g.; \forall k \in \mathbb{Z}, a \in \operatorname{carr}(g). \\ \quad \operatorname{is\text{-}abelian\text{-}group}(g) \Rightarrow (\operatorname{zz\text{-}act}(g, k, a) \in \operatorname{carr}(g)) \end{array}}}
\vnbprop{zz-act-neg-sign}{\vnbstmt{\begin{array}{@{}l@{}} \forall g.; \forall k \in \mathbb{Z}, a \in \operatorname{carr}(g). \\ \quad \operatorname{is\text{-}abelian\text{-}group}(g) \Rightarrow (\operatorname{zz\text{-}act}(g, -k, a) = \operatorname{inv}(g)(\operatorname{zz\text{-}act}(g, k, a))) \end{array}}}
\vnbprop{zz-act-add}{\vnbstmt{\begin{array}{@{}l@{}} \forall g.; \forall j, k \in \mathbb{Z}, a \in \operatorname{carr}(g). \\ \quad \operatorname{is\text{-}abelian\text{-}group}(g) \Rightarrow  \\ \quad \quad \operatorname{zz\text{-}act}(g, (j + k), a) = \operatorname{opr}(g)(\operatorname{zz\text{-}act}(g, j, a), \operatorname{zz\text{-}act}(g, k, a)) \end{array}}}

\vnbbookfile{theorem-library/rake-ringoid-additive}{theorem-library/rake-ringoid-additive}
\noindent{\small rake-ringoid-additive.scm -- the additive group of a RINGOID, proven.}\par
\vnbprop{r6b-radd-type}{\vnbstmt{\begin{array}{@{}l@{}} \forall r, a, b. \\ \quad (\operatorname{is\text{-}ringoid}(r) \wedge ((a \in \operatorname{carr}(r)) \wedge (b \in \operatorname{carr}(r)))) \Rightarrow  \\ \quad \quad \operatorname{add}(r)(a, b) \in \operatorname{carr}(r) \end{array}}}
\vnbprop{r6b-rneg-type}{\vnbstmt{\forall r.; \forall a \in \operatorname{carr}(r).\; (\operatorname{is\text{-}ringoid}(r) \Rightarrow (\operatorname{neg}(r)(a) \in \operatorname{carr}(r)))}}
\vnbprop{r6b-radd-right-id}{\vnbstmt{\forall r.; \forall a \in \operatorname{carr}(r).\; (\operatorname{is\text{-}ringoid}(r) \Rightarrow (\operatorname{add}(r)(a, \operatorname{zero}(r)) = a))}}
\vnbprop{r6b-radd-interchange}{\vnbstmt{\begin{array}{@{}l@{}} \forall r, a, b, c, d. \\ \quad (\operatorname{is\text{-}ringoid}(r) \wedge ((a \in \operatorname{carr}(r)) \wedge ((b \in \operatorname{carr}(r)) \wedge ((c \in \operatorname{carr}(r)) \wedge (d \in \operatorname{carr}(r)))))) \Rightarrow  \\ \quad \quad \operatorname{add}(r)(\operatorname{add}(r)(a, b), \operatorname{add}(r)(c, d)) = \operatorname{add}(r)(\operatorname{add}(r)(a, c), \operatorname{add}(r)(b, d)) \end{array}}}
\vnbprop{r6b-rinv-unique}{\vnbstmt{\begin{array}{@{}l@{}} \forall r, a, b. \\ \quad (\operatorname{is\text{-}ringoid}(r) \wedge ((a \in \operatorname{carr}(r)) \wedge ((b \in \operatorname{carr}(r)) \wedge (\operatorname{add}(r)(a, b) = \operatorname{zero}(r))))) \Rightarrow  \\ \quad \quad b = \operatorname{neg}(r)(a) \end{array}}}
\vnbprop{r6b-rneg-add}{\vnbstmt{\begin{array}{@{}l@{}} \forall r, a, b. \\ \quad (\operatorname{is\text{-}ringoid}(r) \wedge ((a \in \operatorname{carr}(r)) \wedge (b \in \operatorname{carr}(r)))) \Rightarrow  \\ \quad \quad \operatorname{neg}(r)(\operatorname{add}(r)(a, b)) = \operatorname{add}(r)(\operatorname{neg}(r)(a), \operatorname{neg}(r)(b)) \end{array}}}
\vnbprop{r6b-rneg-neg}{\vnbstmt{\forall r.; \forall a \in \operatorname{carr}(r).\; (\operatorname{is\text{-}ringoid}(r) \Rightarrow (\operatorname{neg}(r)(\operatorname{neg}(r)(a)) = a))}}
\vnbprop{r6b-rdiff-add}{\vnbstmt{\begin{array}{@{}l@{}} \forall r, a, b, u, v. \\ \quad (\operatorname{is\text{-}ringoid}(r) \wedge ((a \in \operatorname{carr}(r)) \wedge ((b \in \operatorname{carr}(r)) \wedge ((u \in \operatorname{carr}(r)) \wedge (v \in \operatorname{carr}(r)))))) \Rightarrow  \\ \quad \quad \operatorname{add}(r)(\operatorname{add}(r)(a, b), \operatorname{neg}(r)(\operatorname{add}(r)(u, v))) \\ \quad \quad =\; \operatorname{add}(r)(\operatorname{add}(r)(a, \operatorname{neg}(r)(u)), \operatorname{add}(r)(b, \operatorname{neg}(r)(v))) \end{array}}}
\vnbprop{ringoid-neg-diff}{\vnbstmt{\begin{array}{@{}l@{}} \forall r, a, b. \\ \quad (\operatorname{is\text{-}ringoid}(r) \wedge ((a \in \operatorname{carr}(r)) \wedge (b \in \operatorname{carr}(r)))) \Rightarrow  \\ \quad \quad \operatorname{neg}(r)(\operatorname{add}(r)(a, \operatorname{neg}(r)(b))) = \operatorname{add}(r)(b, \operatorname{neg}(r)(a)) \end{array}}}
\vnbprop{ringoid-diff-telescope}{\vnbstmt{\begin{array}{@{}l@{}} \forall r, a, b, c. \\ \quad (\operatorname{is\text{-}ringoid}(r) \wedge ((a \in \operatorname{carr}(r)) \wedge ((b \in \operatorname{carr}(r)) \wedge (c \in \operatorname{carr}(r))))) \Rightarrow  \\ \quad \quad \operatorname{add}(r)(\operatorname{add}(r)(a, \operatorname{neg}(r)(b)), \operatorname{add}(r)(b, \operatorname{neg}(r)(c))) = \operatorname{add}(r)(a, \operatorname{neg}(r)(c)) \end{array}}}

\vnbbookfile{theorem-library/rake-finsum-typing}{theorem-library/rake-finsum-typing}
\noindent{\small rake-finsum-typing.scm -- the finsum / monoid-power typing leaves of the rake (batch G), PROVEN, together with the guarded form of the one statement in the batch that is NOT provable as stated.}\par
\vnbprop{comm-monoid-carrier-closed-opr}{\vnbstmt{\begin{array}{@{}l@{}} \forall s.; \forall a, b \in \operatorname{carr}(s). \\ \quad \operatorname{is\text{-}comm\text{-}monoid}(s) \Rightarrow (\operatorname{opr}(s)(a, b) \in \operatorname{carr}(s)) \end{array}}}
\vnbprop{comm-monoid-identity-in-carr}{\vnbstmt{\forall s.\; (\operatorname{is\text{-}comm\text{-}monoid}(s) \Rightarrow (\operatorname{iden}(s) \in \operatorname{carr}(s)))}}
\vnbprop{sum-ag-comm-monoid-type-ind}{\vnbstmt{\begin{array}{@{}l@{}} \forall n \in \mathbb{N}, m.; \forall f \in (\mathbb{N} \to \operatorname{carr}(m)). \\ \quad \operatorname{is\text{-}comm\text{-}monoid}(m) \Rightarrow (\operatorname{sum\text{-}ag}(m, f, n) \in \operatorname{carr}(m)) \end{array}}}
\vnbprop{enum-fam-comm-monoid-in-fun}{\vnbstmt{\begin{array}{@{}l@{}} \forall n \in \mathbb{N}, m, s.; \forall i \in (\mathrm{OS}(n) \to s), f \in (s \to \operatorname{carr}(m)). \\ \quad \operatorname{is\text{-}comm\text{-}monoid}(m) \Rightarrow (\operatorname{enum\text{-}fam}(m, f, i, n) \in (\mathbb{N} \to \operatorname{carr}(m))) \end{array}}}
\vnbprop{finsum-comm-monoid-type}{\vnbstmt{\begin{array}{@{}l@{}} \forall m.; \forall s \in \mathbf{Set}, f \in (s \to \operatorname{carr}(m)). \\ \quad (\operatorname{is\text{-}comm\text{-}monoid}(m) \wedge (|s| \in \mathbb{N})) \Rightarrow (\operatorname{finsum}(m, f, s) \in \operatorname{carr}(m)) \end{array}}}
\vnbprop{crmcm-carr}{\vnbstmt{\begin{array}{@{}l@{}} \forall r. \\ \quad \operatorname{is\text{-}commutative\text{-}ring}(r) \Rightarrow  \\ \quad \quad \operatorname{carr}(\operatorname{commutative\text{-}ring\text{-}multiplicative\text{-}cm}(r)) = \operatorname{carr}(r) \end{array}}}
\vnbprop{prod-ring-type}{\vnbstmt{\begin{array}{@{}l@{}} \forall r.; \forall x \in \mathbf{Set}, f \in (x \to \operatorname{carr}(r)). \\ \quad (\operatorname{is\text{-}commutative\text{-}ring}(r) \wedge (|x| \in \mathbb{N})) \Rightarrow (\operatorname{prod\text{-}ring}(r, f, x) \in \operatorname{carr}(r)) \end{array}}}
\vnbprop{mpow-comm-monoid-type-ind}{\vnbstmt{\begin{array}{@{}l@{}} \forall n \in \mathbb{N}, m.; \forall x \in \operatorname{carr}(m). \\ \quad \operatorname{is\text{-}comm\text{-}monoid}(m) \Rightarrow (\operatorname{mpow}(m, x, n) \in \operatorname{carr}(m)) \end{array}}}
\vnbprop{ring-power-type}{\vnbstmt{\begin{array}{@{}l@{}} \forall r.; \forall x \in \operatorname{carr}(r), n \in \mathbb{N}. \\ \quad \operatorname{is\text{-}commutative\text{-}ring}(r) \Rightarrow (\operatorname{ring\text{-}power}(r, x, n) \in \operatorname{carr}(r)) \end{array}}}
\vnbprop{sum-ag-in-subset-ind}{\vnbstmt{\begin{array}{@{}l@{}} \forall n \in \mathbb{N}, g, m, a. \\ \quad \operatorname{is\text{-}abelian\text{-}group}(g) \wedge \\ \quad \operatorname{iden}(g) \in m \wedge \\ \quad \forall x, y \in m.\; (\operatorname{opr}(g)(x, y) \in m) \wedge \\ \quad \forall i \in \mathrm{OS}(n).\; (a(i) \in m) \Rightarrow \\ \quad \quad \operatorname{sum\text{-}ag}(g, a, n) \in m \end{array}}}
\vnbprop{finsum-in-subset}{\vnbstmt{\begin{array}{@{}l@{}} \forall g, m.; \forall s \in \mathbf{Set}, f \in (s \to \operatorname{carr}(g)). \\ \quad \operatorname{is\text{-}abelian\text{-}group}(g) \wedge \\ \quad \operatorname{iden}(g) \in m \wedge \\ \quad \forall x, y \in m.\; (\operatorname{opr}(g)(x, y) \in m) \wedge \\ \quad |s| \in \mathbb{N} \wedge \\ \quad \forall z \in s.\; (f(z) \in m) \Rightarrow \\ \quad \quad \operatorname{finsum}(g, f, s) \in m \end{array}}}

\vnbbookfile{theorem-library/rake-finsum-welldef}{theorem-library/rake-finsum-welldef}
\noindent{\small theorem-library/rake-finsum-welldef.scm -- rake batch N (2026-09-17), part 2 of 2 --------------------------------------------------------------------------- NOTE, 2026-09-20 (batch 9-B).  This file was written while cardinality was AXIOMATISED under the name CARD and the defined constant was its companion}\par
\vnbprop{opr-rearrange-laws}{\vnbstmt{\begin{array}{@{}l@{}} \forall s.; \forall a, b, c \in \operatorname{carr}(s). \\ \quad \operatorname{iden}(s) \in \operatorname{carr}(s) \wedge \\ \quad \forall \alpha, \beta \in \operatorname{carr}(s).\; (\operatorname{opr}(s)(\alpha, \beta) \in \operatorname{carr}(s)) \wedge \\ \quad \forall \alpha, \beta, \gamma \in \operatorname{carr}(s). \\ \quad \quad \operatorname{opr}(s)(\operatorname{opr}(s)(\alpha, \beta), \gamma) = \operatorname{opr}(s)(\alpha, \operatorname{opr}(s)(\beta, \gamma)) \wedge \\ \quad \forall \alpha, \beta \in \operatorname{carr}(s). \\ \quad \quad \operatorname{opr}(s)(\alpha, \beta) = \operatorname{opr}(s)(\beta, \alpha) \Rightarrow \\ \quad \quad \operatorname{opr}(s)(\operatorname{opr}(s)(a, b), c) = \operatorname{opr}(s)(\operatorname{opr}(s)(a, c), b) \end{array}}}
\vnbprop{sum-ag-type-laws}{\vnbstmt{\begin{array}{@{}l@{}} \forall n \in \mathbb{N}, s.; \forall g \in (\mathbb{N} \to \operatorname{carr}(s)). \\ \quad \operatorname{iden}(s) \in \operatorname{carr}(s) \wedge \\ \quad \forall a, b \in \operatorname{carr}(s).\; (\operatorname{opr}(s)(a, b) \in \operatorname{carr}(s)) \wedge \\ \quad \forall a, b, c \in \operatorname{carr}(s). \\ \quad \quad \operatorname{opr}(s)(\operatorname{opr}(s)(a, b), c) = \operatorname{opr}(s)(a, \operatorname{opr}(s)(b, c)) \wedge \\ \quad \forall a, b \in \operatorname{carr}(s).\; (\operatorname{opr}(s)(a, b) = \operatorname{opr}(s)(b, a)) \Rightarrow \\ \quad \quad \operatorname{sum\text{-}ag}(s, g, n) \in \operatorname{carr}(s) \end{array}}}
\vnbprop{sum-ag-segment-congruence}{\vnbstmt{\begin{array}{@{}l@{}} \forall n \in \mathbb{N}, a, g, h. \\ \quad \forall i \in \mathrm{OS}(n).\; (g(i) \simeq h(i)) \Rightarrow  \\ \quad \quad (\operatorname{sum\text{-}ag}(a, g, n) \simeq \operatorname{sum\text{-}ag}(a, h, n)) \end{array}}}
\vnbprop{sum-ag-replace-entry-laws}{\vnbstmt{\begin{array}{@{}l@{}} \forall n \in \mathbb{N}, s.; \forall g, p \in (\mathbb{N} \to \operatorname{carr}(s)), k \in \mathrm{OS}(n). \\ \quad \operatorname{iden}(s) \in \operatorname{carr}(s) \wedge \\ \quad \forall a, b \in \operatorname{carr}(s).\; (\operatorname{opr}(s)(a, b) \in \operatorname{carr}(s)) \wedge \\ \quad \forall a, b, c \in \operatorname{carr}(s). \\ \quad \quad \operatorname{opr}(s)(\operatorname{opr}(s)(a, b), c) = \operatorname{opr}(s)(a, \operatorname{opr}(s)(b, c)) \wedge \\ \quad \forall a, b \in \operatorname{carr}(s).\; (\operatorname{opr}(s)(a, b) = \operatorname{opr}(s)(b, a)) \wedge \\ \quad \forall i \in \mathrm{OS}(n).\; (\neg (i = k) \Rightarrow (g(i) \simeq p(i))) \Rightarrow \\ \quad \quad \operatorname{opr}(s)(\operatorname{sum\text{-}ag}(s, g, n), p(k)) = \operatorname{opr}(s)(\operatorname{sum\text{-}ag}(s, p, n), g(k)) \end{array}}}
\vnbprop{seg-perm-restrict-is-bijection}{\vnbstmt{\begin{array}{@{}l@{}} \forall n \in \mathbb{N}, i \in \mathrm{Bij}(\mathrm{OS}((n + 1)), \mathrm{OS}((n + 1))), k. \\ \quad (i(n) = k) \Rightarrow  \\ \quad \quad (\lambda a \in \mathrm{OS}(n).\; \left(\text{if }(i(a) = n)\text{ then }k\text{ else }i(a)\right)) \in \mathrm{Bij}(\mathrm{OS}(n), \mathrm{OS}(n)) \end{array}}}
\vnbprop{sum-ag-permutation-invariance-laws}{\vnbstmt{\begin{array}{@{}l@{}} \forall n \in \mathbb{N}, s.; \forall g, h \in (\mathbb{N} \to \operatorname{carr}(s)), a \in \mathrm{Bij}(\mathrm{OS}(n), \mathrm{OS}(n)). \\ \quad \operatorname{iden}(s) \in \operatorname{carr}(s) \wedge \\ \quad \forall \alpha, b \in \operatorname{carr}(s).\; (\operatorname{opr}(s)(\alpha, b) \in \operatorname{carr}(s)) \wedge \\ \quad \forall \alpha, b, c \in \operatorname{carr}(s). \\ \quad \quad \operatorname{opr}(s)(\operatorname{opr}(s)(\alpha, b), c) = \operatorname{opr}(s)(\alpha, \operatorname{opr}(s)(b, c)) \wedge \\ \quad \forall \alpha, b \in \operatorname{carr}(s). \\ \quad \quad \operatorname{opr}(s)(\alpha, b) = \operatorname{opr}(s)(b, \alpha) \wedge \\ \quad \forall i \in \mathrm{OS}(n).\; (h(i) = g(a(i))) \Rightarrow \\ \quad \quad \operatorname{sum\text{-}ag}(s, g, n) = \operatorname{sum\text{-}ag}(s, h, n) \end{array}}}
\vnbprop{sum-ag-permutation-invariance}{\vnbstmt{\begin{array}{@{}l@{}} \forall n \in \mathbb{N}, a.; \forall g, h \in (\mathbb{N} \to \operatorname{carr}(a)), b \in \mathrm{Bij}(\mathrm{OS}(n), \mathrm{OS}(n)). \\ \quad (\operatorname{is\text{-}abelian\text{-}group}(a) \wedge \forall i \in \mathrm{OS}(n).\; (h(i) = g(b(i)))) \Rightarrow  \\ \quad \quad \operatorname{sum\text{-}ag}(a, g, n) = \operatorname{sum\text{-}ag}(a, h, n) \end{array}}}
\vnbprop{finsum-comm-monoid-permutation-invariance}{\vnbstmt{\begin{array}{@{}l@{}} \forall n \in \mathbb{N}, m.; \forall g, h \in (\mathbb{N} \to \operatorname{carr}(m)), a \in \mathrm{Bij}(\mathrm{OS}(n), \mathrm{OS}(n)). \\ \quad (\operatorname{is\text{-}comm\text{-}monoid}(m) \wedge \forall i \in \mathrm{OS}(n).\; (h(i) = g(a(i)))) \Rightarrow  \\ \quad \quad \operatorname{sum\text{-}ag}(m, g, n) = \operatorname{sum\text{-}ag}(m, h, n) \end{array}}}
\vnbprop{finsum-well-defined}{\vnbstmt{\begin{array}{@{}l@{}} \forall s \in \mathbf{Set}, g.; \forall f \in (s \to \operatorname{carr}(g)), m \in \mathrm{Bij}(\mathrm{OS}(|s|), s). \\ \quad ((|s| \in \mathbb{N}) \wedge \operatorname{is\text{-}abelian\text{-}group}(g)) \Rightarrow  \\ \quad \quad \operatorname{finsum}(g, f, s) = \operatorname{sum\text{-}ag}(g, \operatorname{enum\text{-}fam}(g, f, m, |s|), |s|) \end{array}}}
\vnbprop{finsum-comm-monoid-well-defined}{\vnbstmt{\begin{array}{@{}l@{}} \forall s \in \mathbf{Set}, m.; \forall f \in (s \to \operatorname{carr}(m)), a \in \mathrm{Bij}(\mathrm{OS}(|s|), s). \\ \quad ((|s| \in \mathbb{N}) \wedge \operatorname{is\text{-}comm\text{-}monoid}(m)) \Rightarrow  \\ \quad \quad \operatorname{finsum}(m, f, s) = \operatorname{sum\text{-}ag}(m, \operatorname{enum\text{-}fam}(m, f, a, |s|), |s|) \end{array}}}

\vnbbookfile{theorem-library/finsum-insert}{theorem-library/finsum-insert}
\noindent{\small theorem-library/finsum-insert.scm -- the insertion recurrence for FINSUM. --------------------------------------------------------------------------- NOTE, 2026-09-20 (batch 9-B).  This file was written while cardinality was AXIOMATISED under the name CARD and the defined constant was its companion}\par
\vnbprop{finsum-insert-from-enum}{\vnbstmt{\begin{array}{@{}l@{}} \forall m, f, x, k, a. \\ \quad |x| \in \mathbb{N} \wedge \\ \quad |(x \cup \{k\})| = (|x| + 1) \wedge \\ \quad (a(|x|) \simeq k) \wedge \\ \quad \forall i \in \mathrm{OS}(|x|).\; (a(i) \simeq \operatorname{fin\text{-}enum}(x)(i)) \wedge \\ \quad \operatorname{finsum}(m, f, (x \cup \{k\})) \\ \quad =\; \operatorname{sum\text{-}ag}(m, \operatorname{enum\text{-}fam}(m, f, a, |(x \cup \{k\})|), |(x \cup \{k\})|) \Rightarrow \\ \quad \quad \operatorname{finsum}(m, f, (x \cup \{k\})) = \operatorname{opr}(m)(\operatorname{finsum}(m, f, x), f(k)) \end{array}}}
\vnbprop{finsum-insert}{\vnbstmt{\begin{array}{@{}l@{}} \forall m.; \forall x, k \in \mathbf{Set}, f \in ((x \cup \{k\}) \to \operatorname{carr}(m)). \\ \quad (\operatorname{is\text{-}comm\text{-}monoid}(m) \wedge ((|x| \in \mathbb{N}) \wedge \neg (k \in x))) \Rightarrow  \\ \quad \quad \operatorname{finsum}(m, f, (x \cup \{k\})) = \operatorname{opr}(m)(\operatorname{finsum}(m, f, x), f(k)) \end{array}}}
\vnbprop{finsum-insert-ag}{\vnbstmt{\begin{array}{@{}l@{}} \forall g.; \forall x, k \in \mathbf{Set}, f \in ((x \cup \{k\}) \to \operatorname{carr}(g)). \\ \quad (\operatorname{is\text{-}abelian\text{-}group}(g) \wedge ((|x| \in \mathbb{N}) \wedge \neg (k \in x))) \Rightarrow  \\ \quad \quad \operatorname{finsum}(g, f, (x \cup \{k\})) = \operatorname{opr}(g)(\operatorname{finsum}(g, f, x), f(k)) \end{array}}}
\vnbprop{finsum-empty}{\vnbstmt{\forall g, f.\; (\operatorname{finsum}(g, f, \emptyset) \simeq \operatorname{iden}(g))}}

\vnbbookfile{theorem-library/field-ring-view}{theorem-library/field-ring-view}
\noindent{\small field-ring-view.scm -- EVERY FIELD IS A FIELD-RING.}\par
\vnbprop{singleton-unfold}{\vnbstmt{\forall y.\; (\operatorname{singleton}(y) \simeq \operatorname{make\text{-}set}(\langle y \rangle))}}
\vnbprop{singleton-membership}{\vnbstmt{\forall y, x.\; ((x \in \operatorname{singleton}(y)) \Leftrightarrow ((x \in \mathbf{Set}) \wedge (x = y)))}}
\vnbprop{field-id-carr}{\vnbstmt{\forall f.\; (\operatorname{is\text{-}field}(f) \Rightarrow (\operatorname{carr}(\operatorname{field\text{-}as\text{-}integral\text{-}domain}(f)) \simeq \operatorname{carr}(f)))}}
\vnbprop{field-id-mul}{\vnbstmt{\forall f.\; (\operatorname{is\text{-}field}(f) \Rightarrow (\operatorname{mul}(\operatorname{field\text{-}as\text{-}integral\text{-}domain}(f)) \simeq \operatorname{mul}(f)))}}
\vnbprop{field-id-one}{\vnbstmt{\forall f.\; (\operatorname{is\text{-}field}(f) \Rightarrow (\operatorname{one}(\operatorname{field\text{-}as\text{-}integral\text{-}domain}(f)) \simeq \operatorname{one}(f)))}}
\vnbprop{field-id-zero}{\vnbstmt{\forall f.\; (\operatorname{is\text{-}field}(f) \Rightarrow (\operatorname{zero}(\operatorname{field\text{-}as\text{-}integral\text{-}domain}(f)) \simeq \operatorname{zero}(f)))}}
\vnbprop{field-id-add}{\vnbstmt{\forall f.\; (\operatorname{is\text{-}field}(f) \Rightarrow (\operatorname{add}(\operatorname{field\text{-}as\text{-}integral\text{-}domain}(f)) \simeq \operatorname{add}(f)))}}
\vnbprop{field-id-neg}{\vnbstmt{\forall f.\; (\operatorname{is\text{-}field}(f) \Rightarrow (\operatorname{neg}(\operatorname{field\text{-}as\text{-}integral\text{-}domain}(f)) \simeq \operatorname{neg}(f)))}}
\vnbprop{field-is-field-ring}{\vnbstmt{\forall s.\; (\operatorname{is\text{-}field}(s) \Rightarrow \operatorname{is\text{-}field\text{-}ring}(\operatorname{field\text{-}as\text{-}integral\text{-}domain}(s)))}}

\vnbbookfile{theorem-library/hom-laws}{theorem-library/hom-laws}
\noindent{\small hom-laws.scm -- the models of a structure form a CATEGORY: identity arrows, composition, and the hom-set as a set.}\par
\vnbprop{hom-semigroup-member-iff}{\vnbstmt{\begin{array}{@{}l@{}} \forall a, b, f. \\ \quad (f \in \operatorname{hom\text{-}semigroup}(a, b)) \Leftrightarrow ((f \in (\operatorname{carr}(a) \to \operatorname{carr}(b))) \wedge \operatorname{is\text{-}hom\text{-}semigroup}(a, b, f)) \end{array}}}
\vnbprop{hom-semigroup-id}{\vnbstmt{\forall a.\; (\operatorname{is\text{-}semigroup}(a) \Rightarrow \operatorname{is\text{-}hom\text{-}semigroup}(a, a, \operatorname{id\text{-}fun}(\operatorname{carr}(a))))}}
\vnbprop{hom-semigroup-in-set}{\vnbstmt{\begin{array}{@{}l@{}} \forall a, b. \\ \quad (\operatorname{is\text{-}semigroup}(a) \wedge \operatorname{is\text{-}semigroup}(b)) \Rightarrow (\operatorname{hom\text{-}semigroup}(a, b) \in \mathbf{Set}) \end{array}}}
\vnbprop{hom-monoid-member-iff}{\vnbstmt{\begin{array}{@{}l@{}} \forall a, b, f. \\ \quad (f \in \operatorname{hom\text{-}monoid}(a, b)) \Leftrightarrow ((f \in (\operatorname{carr}(a) \to \operatorname{carr}(b))) \wedge \operatorname{is\text{-}hom\text{-}monoid}(a, b, f)) \end{array}}}
\vnbprop{hom-monoid-id}{\vnbstmt{\forall a.\; (\operatorname{is\text{-}monoid}(a) \Rightarrow \operatorname{is\text{-}hom\text{-}monoid}(a, a, \operatorname{id\text{-}fun}(\operatorname{carr}(a))))}}
\vnbprop{hom-monoid-in-set}{\vnbstmt{\forall a, b.\; ((\operatorname{is\text{-}monoid}(a) \wedge \operatorname{is\text{-}monoid}(b)) \Rightarrow (\operatorname{hom\text{-}monoid}(a, b) \in \mathbf{Set}))}}
\vnbprop{hom-comm-monoid-member-iff}{\vnbstmt{\begin{array}{@{}l@{}} \forall a, b, f. \\ \quad (f \in \operatorname{hom\text{-}comm\text{-}monoid}(a, b)) \Leftrightarrow ((f \in (\operatorname{carr}(a) \to \operatorname{carr}(b))) \wedge \operatorname{is\text{-}hom\text{-}comm\text{-}monoid}(a, b, f)) \end{array}}}
\vnbprop{hom-comm-monoid-id}{\vnbstmt{\forall a.\; (\operatorname{is\text{-}comm\text{-}monoid}(a) \Rightarrow \operatorname{is\text{-}hom\text{-}comm\text{-}monoid}(a, a, \operatorname{id\text{-}fun}(\operatorname{carr}(a))))}}
\vnbprop{hom-comm-monoid-in-set}{\vnbstmt{\begin{array}{@{}l@{}} \forall a, b. \\ \quad (\operatorname{is\text{-}comm\text{-}monoid}(a) \wedge \operatorname{is\text{-}comm\text{-}monoid}(b)) \Rightarrow (\operatorname{hom\text{-}comm\text{-}monoid}(a, b) \in \mathbf{Set}) \end{array}}}
\vnbprop{hom-group-member-iff}{\vnbstmt{\begin{array}{@{}l@{}} \forall a, b, f. \\ \quad (f \in \operatorname{hom\text{-}group}(a, b)) \Leftrightarrow ((f \in (\operatorname{carr}(a) \to \operatorname{carr}(b))) \wedge \operatorname{is\text{-}hom\text{-}group}(a, b, f)) \end{array}}}
\vnbprop{hom-group-id}{\vnbstmt{\forall a.\; (\operatorname{is\text{-}group}(a) \Rightarrow \operatorname{is\text{-}hom\text{-}group}(a, a, \operatorname{id\text{-}fun}(\operatorname{carr}(a))))}}
\vnbprop{hom-group-in-set}{\vnbstmt{\forall a, b.\; ((\operatorname{is\text{-}group}(a) \wedge \operatorname{is\text{-}group}(b)) \Rightarrow (\operatorname{hom\text{-}group}(a, b) \in \mathbf{Set}))}}
\vnbprop{hom-abelian-group-member-iff}{\vnbstmt{\begin{array}{@{}l@{}} \forall a, b, f. \\ \quad (f \in \operatorname{hom\text{-}abelian\text{-}group}(a, b)) \Leftrightarrow ((f \in (\operatorname{carr}(a) \to \operatorname{carr}(b))) \wedge \operatorname{is\text{-}hom\text{-}abelian\text{-}group}(a, b, f)) \end{array}}}
\vnbprop{hom-abelian-group-id}{\vnbstmt{\begin{array}{@{}l@{}} \forall a. \\ \quad \operatorname{is\text{-}abelian\text{-}group}(a) \Rightarrow \operatorname{is\text{-}hom\text{-}abelian\text{-}group}(a, a, \operatorname{id\text{-}fun}(\operatorname{carr}(a))) \end{array}}}
\vnbprop{hom-abelian-group-in-set}{\vnbstmt{\begin{array}{@{}l@{}} \forall a, b. \\ \quad (\operatorname{is\text{-}abelian\text{-}group}(a) \wedge \operatorname{is\text{-}abelian\text{-}group}(b)) \Rightarrow  \\ \quad \quad \operatorname{hom\text{-}abelian\text{-}group}(a, b) \in \mathbf{Set} \end{array}}}
\vnbprop{hom-ring-member-iff}{\vnbstmt{\begin{array}{@{}l@{}} \forall a, b, f. \\ \quad (f \in \operatorname{hom\text{-}ring}(a, b)) \Leftrightarrow ((f \in (\operatorname{carr}(a) \to \operatorname{carr}(b))) \wedge \operatorname{is\text{-}hom\text{-}ring}(a, b, f)) \end{array}}}
\vnbprop{hom-ring-id}{\vnbstmt{\forall a.\; (\operatorname{is\text{-}ring}(a) \Rightarrow \operatorname{is\text{-}hom\text{-}ring}(a, a, \operatorname{id\text{-}fun}(\operatorname{carr}(a))))}}
\vnbprop{hom-ring-in-set}{\vnbstmt{\forall a, b.\; ((\operatorname{is\text{-}ring}(a) \wedge \operatorname{is\text{-}ring}(b)) \Rightarrow (\operatorname{hom\text{-}ring}(a, b) \in \mathbf{Set}))}}
\vnbprop{hom-commutative-ring-member-iff}{\vnbstmt{\begin{array}{@{}l@{}} \forall a, b, f. \\ \quad (f \in \operatorname{hom\text{-}commutative\text{-}ring}(a, b)) \Leftrightarrow ((f \in (\operatorname{carr}(a) \to \operatorname{carr}(b))) \wedge \operatorname{is\text{-}hom\text{-}commutative\text{-}ring}(a, b, f)) \end{array}}}
\vnbprop{hom-commutative-ring-id}{\vnbstmt{\begin{array}{@{}l@{}} \forall a. \\ \quad \operatorname{is\text{-}commutative\text{-}ring}(a) \Rightarrow \operatorname{is\text{-}hom\text{-}commutative\text{-}ring}(a, a, \operatorname{id\text{-}fun}(\operatorname{carr}(a))) \end{array}}}
\vnbprop{hom-commutative-ring-in-set}{\vnbstmt{\begin{array}{@{}l@{}} \forall a, b. \\ \quad (\operatorname{is\text{-}commutative\text{-}ring}(a) \wedge \operatorname{is\text{-}commutative\text{-}ring}(b)) \Rightarrow  \\ \quad \quad \operatorname{hom\text{-}commutative\text{-}ring}(a, b) \in \mathbf{Set} \end{array}}}
\vnbprop{hom-integral-domain-member-iff}{\vnbstmt{\begin{array}{@{}l@{}} \forall a, b, f. \\ \quad (f \in \operatorname{hom\text{-}integral\text{-}domain}(a, b)) \Leftrightarrow ((f \in (\operatorname{carr}(a) \to \operatorname{carr}(b))) \wedge \operatorname{is\text{-}hom\text{-}integral\text{-}domain}(a, b, f)) \end{array}}}
\vnbprop{hom-integral-domain-id}{\vnbstmt{\begin{array}{@{}l@{}} \forall a. \\ \quad \operatorname{is\text{-}integral\text{-}domain}(a) \Rightarrow \operatorname{is\text{-}hom\text{-}integral\text{-}domain}(a, a, \operatorname{id\text{-}fun}(\operatorname{carr}(a))) \end{array}}}
\vnbprop{hom-integral-domain-in-set}{\vnbstmt{\begin{array}{@{}l@{}} \forall a, b. \\ \quad (\operatorname{is\text{-}integral\text{-}domain}(a) \wedge \operatorname{is\text{-}integral\text{-}domain}(b)) \Rightarrow  \\ \quad \quad \operatorname{hom\text{-}integral\text{-}domain}(a, b) \in \mathbf{Set} \end{array}}}
\vnbprop{hom-euclidean-ring-member-iff}{\vnbstmt{\begin{array}{@{}l@{}} \forall a, b, f. \\ \quad (f \in \operatorname{hom\text{-}euclidean\text{-}ring}(a, b)) \Leftrightarrow ((f \in (\operatorname{carr}(a) \to \operatorname{carr}(b))) \wedge \operatorname{is\text{-}hom\text{-}euclidean\text{-}ring}(a, b, f)) \end{array}}}
\vnbprop{hom-euclidean-ring-id}{\vnbstmt{\begin{array}{@{}l@{}} \forall a. \\ \quad \operatorname{is\text{-}euclidean\text{-}ring}(a) \Rightarrow \operatorname{is\text{-}hom\text{-}euclidean\text{-}ring}(a, a, \operatorname{id\text{-}fun}(\operatorname{carr}(a))) \end{array}}}
\vnbprop{hom-euclidean-ring-in-set}{\vnbstmt{\begin{array}{@{}l@{}} \forall a, b. \\ \quad (\operatorname{is\text{-}euclidean\text{-}ring}(a) \wedge \operatorname{is\text{-}euclidean\text{-}ring}(b)) \Rightarrow  \\ \quad \quad \operatorname{hom\text{-}euclidean\text{-}ring}(a, b) \in \mathbf{Set} \end{array}}}
\vnbprop{hom-pid-member-iff}{\vnbstmt{\begin{array}{@{}l@{}} \forall a, b, f. \\ \quad (f \in \operatorname{hom\text{-}pid}(a, b)) \Leftrightarrow ((f \in (\operatorname{carr}(a) \to \operatorname{carr}(b))) \wedge \operatorname{is\text{-}hom\text{-}pid}(a, b, f)) \end{array}}}
\vnbprop{hom-pid-id}{\vnbstmt{\forall a.\; (\operatorname{is\text{-}pid}(a) \Rightarrow \operatorname{is\text{-}hom\text{-}pid}(a, a, \operatorname{id\text{-}fun}(\operatorname{carr}(a))))}}
\vnbprop{hom-pid-in-set}{\vnbstmt{\forall a, b.\; ((\operatorname{is\text{-}pid}(a) \wedge \operatorname{is\text{-}pid}(b)) \Rightarrow (\operatorname{hom\text{-}pid}(a, b) \in \mathbf{Set}))}}
\vnbprop{hom-field-ring-member-iff}{\vnbstmt{\begin{array}{@{}l@{}} \forall a, b, f. \\ \quad (f \in \operatorname{hom\text{-}field\text{-}ring}(a, b)) \Leftrightarrow ((f \in (\operatorname{carr}(a) \to \operatorname{carr}(b))) \wedge \operatorname{is\text{-}hom\text{-}field\text{-}ring}(a, b, f)) \end{array}}}
\vnbprop{hom-field-ring-id}{\vnbstmt{\forall a.\; (\operatorname{is\text{-}field\text{-}ring}(a) \Rightarrow \operatorname{is\text{-}hom\text{-}field\text{-}ring}(a, a, \operatorname{id\text{-}fun}(\operatorname{carr}(a))))}}
\vnbprop{hom-field-ring-in-set}{\vnbstmt{\begin{array}{@{}l@{}} \forall a, b. \\ \quad (\operatorname{is\text{-}field\text{-}ring}(a) \wedge \operatorname{is\text{-}field\text{-}ring}(b)) \Rightarrow (\operatorname{hom\text{-}field\text{-}ring}(a, b) \in \mathbf{Set}) \end{array}}}
\vnbprop{hom-field-non-zero}{\vnbstmt{\begin{array}{@{}l@{}} \forall a, b, f.; \forall x \in \operatorname{non\text{-}zero}(a). \\ \quad \operatorname{is\text{-}hom\text{-}field}(a, b, f) \Rightarrow (f(x) \in \operatorname{non\text{-}zero}(b)) \end{array}}}
\vnbprop{hom-field-member-iff}{\vnbstmt{\begin{array}{@{}l@{}} \forall a, b, f. \\ \quad (f \in \operatorname{hom\text{-}field}(a, b)) \Leftrightarrow ((f \in (\operatorname{carr}(a) \to \operatorname{carr}(b))) \wedge \operatorname{is\text{-}hom\text{-}field}(a, b, f)) \end{array}}}
\vnbprop{hom-field-id}{\vnbstmt{\forall a.\; (\operatorname{is\text{-}field}(a) \Rightarrow \operatorname{is\text{-}hom\text{-}field}(a, a, \operatorname{id\text{-}fun}(\operatorname{carr}(a))))}}
\vnbprop{hom-field-compose}{\vnbstmt{\begin{array}{@{}l@{}} \forall a, b, c, f, g. \\ \quad (\operatorname{is\text{-}hom\text{-}field}(a, b, f) \wedge \operatorname{is\text{-}hom\text{-}field}(b, c, g)) \Rightarrow  \\ \quad \quad \operatorname{is\text{-}hom\text{-}field}(a, c, \operatorname{compose}(g, f)) \end{array}}}
\vnbprop{hom-field-in-set}{\vnbstmt{\forall a, b.\; ((\operatorname{is\text{-}field}(a) \wedge \operatorname{is\text{-}field}(b)) \Rightarrow (\operatorname{hom\text{-}field}(a, b) \in \mathbf{Set}))}}
\vnbprop{hom-ringoid-member-iff}{\vnbstmt{\begin{array}{@{}l@{}} \forall a, b, f. \\ \quad (f \in \operatorname{hom\text{-}ringoid}(a, b)) \Leftrightarrow ((f \in (\operatorname{carr}(a) \to \operatorname{carr}(b))) \wedge \operatorname{is\text{-}hom\text{-}ringoid}(a, b, f)) \end{array}}}
\vnbprop{hom-ringoid-id}{\vnbstmt{\forall a.\; (\operatorname{is\text{-}ringoid}(a) \Rightarrow \operatorname{is\text{-}hom\text{-}ringoid}(a, a, \operatorname{id\text{-}fun}(\operatorname{carr}(a))))}}
\vnbprop{hom-ringoid-in-set}{\vnbstmt{\forall a, b.\; ((\operatorname{is\text{-}ringoid}(a) \wedge \operatorname{is\text{-}ringoid}(b)) \Rightarrow (\operatorname{hom\text{-}ringoid}(a, b) \in \mathbf{Set}))}}

\vnbbookfile{theorem-library/hom-kinds}{theorem-library/hom-kinds}
\noindent{\small hom-kinds.scm -- the category laws for the arrows the generic driver of hom-laws.scm does not reach: the FAMILY slot kinds, and the RELATION kind (SETOID, since batch 40; two carriers in batch 39).}\par
\vnbprop{hom-setoid-member-iff}{\vnbstmt{\begin{array}{@{}l@{}} \forall a, b, f. \\ \quad (f \in \operatorname{hom\text{-}setoid}(a, b)) \Leftrightarrow ((f \in (\operatorname{pts}(a) \to \operatorname{pts}(b))) \wedge \operatorname{is\text{-}hom\text{-}setoid}(a, b, f)) \end{array}}}
\vnbprop{hom-setoid-id}{\vnbstmt{\forall a.\; (\operatorname{is\text{-}setoid}(a) \Rightarrow \operatorname{is\text{-}hom\text{-}setoid}(a, a, \operatorname{id\text{-}fun}(\operatorname{pts}(a))))}}
\vnbprop{hom-setoid-in-set}{\vnbstmt{\forall a, b.\; ((\operatorname{is\text{-}setoid}(a) \wedge \operatorname{is\text{-}setoid}(b)) \Rightarrow (\operatorname{hom\text{-}setoid}(a, b) \in \mathbf{Set}))}}

\vnbbookfile{theorem-library/hom-functors}{theorem-library/hom-functors}
\noindent{\small hom-functors.scm -- the Hom functors: Hom(a, -) covariant, Hom(-, b) contravariant, as maps on hom-sets.}\par
\vnbprop{hom-semigroup-post-type}{\vnbstmt{\begin{array}{@{}l@{}} \forall a, b, c, h. \\ \quad (\operatorname{is\text{-}semigroup}(a) \wedge (h \in \operatorname{hom\text{-}semigroup}(b, c))) \Rightarrow  \\ \quad \quad (\lambda f \in \operatorname{hom\text{-}semigroup}(a, b).\; \operatorname{compose}(h, f)) \in (\operatorname{hom\text{-}semigroup}(a, b) \to \operatorname{hom\text{-}semigroup}(a, c)) \end{array}}}
\vnbprop{hom-semigroup-pre-type}{\vnbstmt{\begin{array}{@{}l@{}} \forall a, b, c, k. \\ \quad (\operatorname{is\text{-}semigroup}(b) \wedge (k \in \operatorname{hom\text{-}semigroup}(c, a))) \Rightarrow  \\ \quad \quad (\lambda f \in \operatorname{hom\text{-}semigroup}(a, b).\; \operatorname{compose}(f, k)) \in (\operatorname{hom\text{-}semigroup}(a, b) \to \operatorname{hom\text{-}semigroup}(c, b)) \end{array}}}
\vnbprop{hom-monoid-post-type}{\vnbstmt{\begin{array}{@{}l@{}} \forall a, b, c, h. \\ \quad (\operatorname{is\text{-}monoid}(a) \wedge (h \in \operatorname{hom\text{-}monoid}(b, c))) \Rightarrow  \\ \quad \quad (\lambda f \in \operatorname{hom\text{-}monoid}(a, b).\; \operatorname{compose}(h, f)) \in (\operatorname{hom\text{-}monoid}(a, b) \to \operatorname{hom\text{-}monoid}(a, c)) \end{array}}}
\vnbprop{hom-monoid-pre-type}{\vnbstmt{\begin{array}{@{}l@{}} \forall a, b, c, k. \\ \quad (\operatorname{is\text{-}monoid}(b) \wedge (k \in \operatorname{hom\text{-}monoid}(c, a))) \Rightarrow  \\ \quad \quad (\lambda f \in \operatorname{hom\text{-}monoid}(a, b).\; \operatorname{compose}(f, k)) \in (\operatorname{hom\text{-}monoid}(a, b) \to \operatorname{hom\text{-}monoid}(c, b)) \end{array}}}
\vnbprop{hom-comm-monoid-post-type}{\vnbstmt{\begin{array}{@{}l@{}} \forall a, b, c, h. \\ \quad (\operatorname{is\text{-}comm\text{-}monoid}(a) \wedge (h \in \operatorname{hom\text{-}comm\text{-}monoid}(b, c))) \Rightarrow  \\ \quad \quad (\lambda f \in \operatorname{hom\text{-}comm\text{-}monoid}(a, b).\; \operatorname{compose}(h, f)) \\ \quad \quad \in\; (\operatorname{hom\text{-}comm\text{-}monoid}(a, b) \to \operatorname{hom\text{-}comm\text{-}monoid}(a, c)) \end{array}}}
\vnbprop{hom-comm-monoid-pre-type}{\vnbstmt{\begin{array}{@{}l@{}} \forall a, b, c, k. \\ \quad (\operatorname{is\text{-}comm\text{-}monoid}(b) \wedge (k \in \operatorname{hom\text{-}comm\text{-}monoid}(c, a))) \Rightarrow  \\ \quad \quad (\lambda f \in \operatorname{hom\text{-}comm\text{-}monoid}(a, b).\; \operatorname{compose}(f, k)) \\ \quad \quad \in\; (\operatorname{hom\text{-}comm\text{-}monoid}(a, b) \to \operatorname{hom\text{-}comm\text{-}monoid}(c, b)) \end{array}}}
\vnbprop{hom-group-post-type}{\vnbstmt{\begin{array}{@{}l@{}} \forall a, b, c, h. \\ \quad (\operatorname{is\text{-}group}(a) \wedge (h \in \operatorname{hom\text{-}group}(b, c))) \Rightarrow  \\ \quad \quad (\lambda f \in \operatorname{hom\text{-}group}(a, b).\; \operatorname{compose}(h, f)) \in (\operatorname{hom\text{-}group}(a, b) \to \operatorname{hom\text{-}group}(a, c)) \end{array}}}
\vnbprop{hom-group-pre-type}{\vnbstmt{\begin{array}{@{}l@{}} \forall a, b, c, k. \\ \quad (\operatorname{is\text{-}group}(b) \wedge (k \in \operatorname{hom\text{-}group}(c, a))) \Rightarrow  \\ \quad \quad (\lambda f \in \operatorname{hom\text{-}group}(a, b).\; \operatorname{compose}(f, k)) \in (\operatorname{hom\text{-}group}(a, b) \to \operatorname{hom\text{-}group}(c, b)) \end{array}}}
\vnbprop{hom-abelian-group-post-type}{\vnbstmt{\begin{array}{@{}l@{}} \forall a, b, c, h. \\ \quad (\operatorname{is\text{-}abelian\text{-}group}(a) \wedge (h \in \operatorname{hom\text{-}abelian\text{-}group}(b, c))) \Rightarrow  \\ \quad \quad (\lambda f \in \operatorname{hom\text{-}abelian\text{-}group}(a, b).\; \operatorname{compose}(h, f)) \\ \quad \quad \in\; (\operatorname{hom\text{-}abelian\text{-}group}(a, b) \to \operatorname{hom\text{-}abelian\text{-}group}(a, c)) \end{array}}}
\vnbprop{hom-abelian-group-pre-type}{\vnbstmt{\begin{array}{@{}l@{}} \forall a, b, c, k. \\ \quad (\operatorname{is\text{-}abelian\text{-}group}(b) \wedge (k \in \operatorname{hom\text{-}abelian\text{-}group}(c, a))) \Rightarrow  \\ \quad \quad (\lambda f \in \operatorname{hom\text{-}abelian\text{-}group}(a, b).\; \operatorname{compose}(f, k)) \\ \quad \quad \in\; (\operatorname{hom\text{-}abelian\text{-}group}(a, b) \to \operatorname{hom\text{-}abelian\text{-}group}(c, b)) \end{array}}}
\vnbprop{hom-ring-post-type}{\vnbstmt{\begin{array}{@{}l@{}} \forall a, b, c, h. \\ \quad (\operatorname{is\text{-}ring}(a) \wedge (h \in \operatorname{hom\text{-}ring}(b, c))) \Rightarrow  \\ \quad \quad (\lambda f \in \operatorname{hom\text{-}ring}(a, b).\; \operatorname{compose}(h, f)) \in (\operatorname{hom\text{-}ring}(a, b) \to \operatorname{hom\text{-}ring}(a, c)) \end{array}}}
\vnbprop{hom-ring-pre-type}{\vnbstmt{\begin{array}{@{}l@{}} \forall a, b, c, k. \\ \quad (\operatorname{is\text{-}ring}(b) \wedge (k \in \operatorname{hom\text{-}ring}(c, a))) \Rightarrow  \\ \quad \quad (\lambda f \in \operatorname{hom\text{-}ring}(a, b).\; \operatorname{compose}(f, k)) \in (\operatorname{hom\text{-}ring}(a, b) \to \operatorname{hom\text{-}ring}(c, b)) \end{array}}}
\vnbprop{hom-commutative-ring-post-type}{\vnbstmt{\begin{array}{@{}l@{}} \forall a, b, c, h. \\ \quad (\operatorname{is\text{-}commutative\text{-}ring}(a) \wedge (h \in \operatorname{hom\text{-}commutative\text{-}ring}(b, c))) \Rightarrow  \\ \quad \quad (\lambda f \in \operatorname{hom\text{-}commutative\text{-}ring}(a, b).\; \operatorname{compose}(h, f)) \\ \quad \quad \in\; (\operatorname{hom\text{-}commutative\text{-}ring}(a, b) \to \operatorname{hom\text{-}commutative\text{-}ring}(a, c)) \end{array}}}
\vnbprop{hom-commutative-ring-pre-type}{\vnbstmt{\begin{array}{@{}l@{}} \forall a, b, c, k. \\ \quad (\operatorname{is\text{-}commutative\text{-}ring}(b) \wedge (k \in \operatorname{hom\text{-}commutative\text{-}ring}(c, a))) \Rightarrow  \\ \quad \quad (\lambda f \in \operatorname{hom\text{-}commutative\text{-}ring}(a, b).\; \operatorname{compose}(f, k)) \\ \quad \quad \in\; (\operatorname{hom\text{-}commutative\text{-}ring}(a, b) \to \operatorname{hom\text{-}commutative\text{-}ring}(c, b)) \end{array}}}
\vnbprop{hom-integral-domain-post-type}{\vnbstmt{\begin{array}{@{}l@{}} \forall a, b, c, h. \\ \quad (\operatorname{is\text{-}integral\text{-}domain}(a) \wedge (h \in \operatorname{hom\text{-}integral\text{-}domain}(b, c))) \Rightarrow  \\ \quad \quad (\lambda f \in \operatorname{hom\text{-}integral\text{-}domain}(a, b).\; \operatorname{compose}(h, f)) \\ \quad \quad \in\; (\operatorname{hom\text{-}integral\text{-}domain}(a, b) \to \operatorname{hom\text{-}integral\text{-}domain}(a, c)) \end{array}}}
\vnbprop{hom-integral-domain-pre-type}{\vnbstmt{\begin{array}{@{}l@{}} \forall a, b, c, k. \\ \quad (\operatorname{is\text{-}integral\text{-}domain}(b) \wedge (k \in \operatorname{hom\text{-}integral\text{-}domain}(c, a))) \Rightarrow  \\ \quad \quad (\lambda f \in \operatorname{hom\text{-}integral\text{-}domain}(a, b).\; \operatorname{compose}(f, k)) \\ \quad \quad \in\; (\operatorname{hom\text{-}integral\text{-}domain}(a, b) \to \operatorname{hom\text{-}integral\text{-}domain}(c, b)) \end{array}}}
\vnbprop{hom-euclidean-ring-post-type}{\vnbstmt{\begin{array}{@{}l@{}} \forall a, b, c, h. \\ \quad (\operatorname{is\text{-}euclidean\text{-}ring}(a) \wedge (h \in \operatorname{hom\text{-}euclidean\text{-}ring}(b, c))) \Rightarrow  \\ \quad \quad (\lambda f \in \operatorname{hom\text{-}euclidean\text{-}ring}(a, b).\; \operatorname{compose}(h, f)) \\ \quad \quad \in\; (\operatorname{hom\text{-}euclidean\text{-}ring}(a, b) \to \operatorname{hom\text{-}euclidean\text{-}ring}(a, c)) \end{array}}}
\vnbprop{hom-euclidean-ring-pre-type}{\vnbstmt{\begin{array}{@{}l@{}} \forall a, b, c, k. \\ \quad (\operatorname{is\text{-}euclidean\text{-}ring}(b) \wedge (k \in \operatorname{hom\text{-}euclidean\text{-}ring}(c, a))) \Rightarrow  \\ \quad \quad (\lambda f \in \operatorname{hom\text{-}euclidean\text{-}ring}(a, b).\; \operatorname{compose}(f, k)) \\ \quad \quad \in\; (\operatorname{hom\text{-}euclidean\text{-}ring}(a, b) \to \operatorname{hom\text{-}euclidean\text{-}ring}(c, b)) \end{array}}}
\vnbprop{hom-pid-post-type}{\vnbstmt{\begin{array}{@{}l@{}} \forall a, b, c, h. \\ \quad (\operatorname{is\text{-}pid}(a) \wedge (h \in \operatorname{hom\text{-}pid}(b, c))) \Rightarrow  \\ \quad \quad (\lambda f \in \operatorname{hom\text{-}pid}(a, b).\; \operatorname{compose}(h, f)) \in (\operatorname{hom\text{-}pid}(a, b) \to \operatorname{hom\text{-}pid}(a, c)) \end{array}}}
\vnbprop{hom-pid-pre-type}{\vnbstmt{\begin{array}{@{}l@{}} \forall a, b, c, k. \\ \quad (\operatorname{is\text{-}pid}(b) \wedge (k \in \operatorname{hom\text{-}pid}(c, a))) \Rightarrow  \\ \quad \quad (\lambda f \in \operatorname{hom\text{-}pid}(a, b).\; \operatorname{compose}(f, k)) \in (\operatorname{hom\text{-}pid}(a, b) \to \operatorname{hom\text{-}pid}(c, b)) \end{array}}}
\vnbprop{hom-field-ring-post-type}{\vnbstmt{\begin{array}{@{}l@{}} \forall a, b, c, h. \\ \quad (\operatorname{is\text{-}field\text{-}ring}(a) \wedge (h \in \operatorname{hom\text{-}field\text{-}ring}(b, c))) \Rightarrow  \\ \quad \quad (\lambda f \in \operatorname{hom\text{-}field\text{-}ring}(a, b).\; \operatorname{compose}(h, f)) \\ \quad \quad \in\; (\operatorname{hom\text{-}field\text{-}ring}(a, b) \to \operatorname{hom\text{-}field\text{-}ring}(a, c)) \end{array}}}
\vnbprop{hom-field-ring-pre-type}{\vnbstmt{\begin{array}{@{}l@{}} \forall a, b, c, k. \\ \quad (\operatorname{is\text{-}field\text{-}ring}(b) \wedge (k \in \operatorname{hom\text{-}field\text{-}ring}(c, a))) \Rightarrow  \\ \quad \quad (\lambda f \in \operatorname{hom\text{-}field\text{-}ring}(a, b).\; \operatorname{compose}(f, k)) \\ \quad \quad \in\; (\operatorname{hom\text{-}field\text{-}ring}(a, b) \to \operatorname{hom\text{-}field\text{-}ring}(c, b)) \end{array}}}
\vnbprop{hom-field-post-type}{\vnbstmt{\begin{array}{@{}l@{}} \forall a, b, c, h. \\ \quad (\operatorname{is\text{-}field}(a) \wedge (h \in \operatorname{hom\text{-}field}(b, c))) \Rightarrow  \\ \quad \quad (\lambda f \in \operatorname{hom\text{-}field}(a, b).\; \operatorname{compose}(h, f)) \in (\operatorname{hom\text{-}field}(a, b) \to \operatorname{hom\text{-}field}(a, c)) \end{array}}}
\vnbprop{hom-field-pre-type}{\vnbstmt{\begin{array}{@{}l@{}} \forall a, b, c, k. \\ \quad (\operatorname{is\text{-}field}(b) \wedge (k \in \operatorname{hom\text{-}field}(c, a))) \Rightarrow  \\ \quad \quad (\lambda f \in \operatorname{hom\text{-}field}(a, b).\; \operatorname{compose}(f, k)) \in (\operatorname{hom\text{-}field}(a, b) \to \operatorname{hom\text{-}field}(c, b)) \end{array}}}
\vnbprop{hom-ringoid-post-type}{\vnbstmt{\begin{array}{@{}l@{}} \forall a, b, c, h. \\ \quad (\operatorname{is\text{-}ringoid}(a) \wedge (h \in \operatorname{hom\text{-}ringoid}(b, c))) \Rightarrow  \\ \quad \quad (\lambda f \in \operatorname{hom\text{-}ringoid}(a, b).\; \operatorname{compose}(h, f)) \in (\operatorname{hom\text{-}ringoid}(a, b) \to \operatorname{hom\text{-}ringoid}(a, c)) \end{array}}}
\vnbprop{hom-ringoid-pre-type}{\vnbstmt{\begin{array}{@{}l@{}} \forall a, b, c, k. \\ \quad (\operatorname{is\text{-}ringoid}(b) \wedge (k \in \operatorname{hom\text{-}ringoid}(c, a))) \Rightarrow  \\ \quad \quad (\lambda f \in \operatorname{hom\text{-}ringoid}(a, b).\; \operatorname{compose}(f, k)) \in (\operatorname{hom\text{-}ringoid}(a, b) \to \operatorname{hom\text{-}ringoid}(c, b)) \end{array}}}
\vnbprop{hom-setoid-post-type}{\vnbstmt{\begin{array}{@{}l@{}} \forall a, b, c, h. \\ \quad (\operatorname{is\text{-}setoid}(a) \wedge (h \in \operatorname{hom\text{-}setoid}(b, c))) \Rightarrow  \\ \quad \quad (\lambda f \in \operatorname{hom\text{-}setoid}(a, b).\; \operatorname{compose}(h, f)) \in (\operatorname{hom\text{-}setoid}(a, b) \to \operatorname{hom\text{-}setoid}(a, c)) \end{array}}}
\vnbprop{hom-setoid-pre-type}{\vnbstmt{\begin{array}{@{}l@{}} \forall a, b, c, k. \\ \quad (\operatorname{is\text{-}setoid}(b) \wedge (k \in \operatorname{hom\text{-}setoid}(c, a))) \Rightarrow  \\ \quad \quad (\lambda f \in \operatorname{hom\text{-}setoid}(a, b).\; \operatorname{compose}(f, k)) \in (\operatorname{hom\text{-}setoid}(a, b) \to \operatorname{hom\text{-}setoid}(c, b)) \end{array}}}

\vnbbookfile{theorem-library/qq-field-is-field}{theorem-library/qq-field-is-field}
\noindent{\small qq-field-is-field.scm -- IS-FIELD(QQ-FIELD), PROVED.}\par
\vnbprop{qq-recip-nonzero}{\vnbstmt{\forall a \in \mathbb{Q}.\; (\neg (a = 0) \Rightarrow \neg (\operatorname{recip}(a) = 0))}}
\vnbprop{qq-field-mul-inv-type}{\vnbstmt{\begin{array}{@{}l@{}} (\lambda x \in \operatorname{difference}(\mathbb{Q}, \operatorname{singleton}(0)).\; \operatorname{recip}(x)) \\ \in\; (\operatorname{difference}(\mathbb{Q}, \operatorname{singleton}(0)) \to \operatorname{difference}(\mathbb{Q}, \operatorname{singleton}(0))) \end{array}}}
\vnbprop{qq-field-is-field}{\vnbstmt{\operatorname{is\text{-}field}(\operatorname{qq\text{-}field})}}

\vnbbookfile{theorem-library/rake-finsum-laws}{theorem-library/rake-finsum-laws}
\noindent{\small theorem-library/rake-finsum-laws.scm -- the ADDITIVE algebra of FINSUM, proven.  Rake batch J (2026-09-17), part 1 of 2; part 2 is theorem-library/rake-finsum-laws2.scm (the three distribution laws and the interval back-peel), which loads AFTER this file and cites two of its}\par
\vnbprop{enum-fam-value}{\vnbstmt{\begin{array}{@{}l@{}} \forall n \in \mathbb{N}, g, u, a.; \forall i \in \mathrm{OS}(n). \\ \quad (\operatorname{enum\text{-}fam}(g, u, a, n)(i) \simeq u(a(i))) \end{array}}}
\vnbprop{sum-ag-type-ptwise}{\vnbstmt{\begin{array}{@{}l@{}} \forall n \in \mathbb{N}, g, a. \\ \quad (\operatorname{is\text{-}abelian\text{-}group}(g) \wedge \forall i \in \mathrm{OS}(n).\; (a(i) \in \operatorname{carr}(g))) \Rightarrow  \\ \quad \quad \operatorname{sum\text{-}ag}(g, a, n) \in \operatorname{carr}(g) \end{array}}}
\vnbprop{abelian-group-opr-interchange}{\vnbstmt{\begin{array}{@{}l@{}} \forall s.; \forall a, b, c, d \in \operatorname{carr}(s). \\ \quad \operatorname{is\text{-}abelian\text{-}group}(s) \Rightarrow  \\ \quad \quad \operatorname{opr}(s)(\operatorname{opr}(s)(a, b), \operatorname{opr}(s)(c, d)) = \operatorname{opr}(s)(\operatorname{opr}(s)(a, c), \operatorname{opr}(s)(b, d)) \end{array}}}
\vnbprop{finsum-congruence-q}{\vnbstmt{\begin{array}{@{}l@{}} \forall a.; \forall s \in \mathbf{Set}, f, g. \\ \quad (\operatorname{is\text{-}abelian\text{-}group}(a) \wedge ((|s| \in \mathbb{N}) \wedge \forall z \in s.\; (f(z) = g(z)))) \Rightarrow  \\ \quad \quad (\operatorname{finsum}(a, f, s) \simeq \operatorname{finsum}(a, g, s)) \end{array}}}
\vnbprop{finsum-congruence-guarded}{\vnbstmt{\begin{array}{@{}l@{}} \forall a.; \forall s \in \mathbf{Set}, f \in (s \to \operatorname{carr}(a)), g. \\ \quad (\operatorname{is\text{-}abelian\text{-}group}(a) \wedge ((|s| \in \mathbb{N}) \wedge \forall z \in s.\; (f(z) = g(z)))) \Rightarrow  \\ \quad \quad \operatorname{finsum}(a, f, s) = \operatorname{finsum}(a, g, s) \end{array}}}
\vnbprop{finsum-type-ptwise}{\vnbstmt{\begin{array}{@{}l@{}} \forall g.; \forall s \in \mathbf{Set}, f. \\ \quad (\operatorname{is\text{-}abelian\text{-}group}(g) \wedge ((|s| \in \mathbb{N}) \wedge \forall z \in s.\; (f(z) \in \operatorname{carr}(g)))) \Rightarrow  \\ \quad \quad \operatorname{finsum}(g, f, s) \in \operatorname{carr}(g) \end{array}}}
\vnbprop{finsum-congruence}{\vnbstmt{\begin{array}{@{}l@{}} \forall a.; \forall s \in \mathbf{Set}, f, g. \\ \quad \operatorname{is\text{-}abelian\text{-}group}(a) \wedge \\ \quad |s| \in \mathbb{N} \wedge \\ \quad \forall z \in s.\; (f(z) = g(z)) \wedge \\ \quad \forall z \in s.\; (f(z) \in \operatorname{carr}(a)) \Rightarrow \\ \quad \quad \operatorname{finsum}(a, f, s) = \operatorname{finsum}(a, g, s) \end{array}}}
\vnbprop{sum-ag-add-ind}{\vnbstmt{\begin{array}{@{}l@{}} \forall n \in \mathbb{N}, g, a, b, c. \\ \quad \operatorname{is\text{-}abelian\text{-}group}(g) \wedge \\ \quad \forall i \in \mathrm{OS}(n).\; (a(i) \in \operatorname{carr}(g)) \wedge \\ \quad \forall i \in \mathrm{OS}(n).\; (b(i) \in \operatorname{carr}(g)) \wedge \\ \quad \forall i \in \mathrm{OS}(n).\; (c(i) \simeq \operatorname{opr}(g)(a(i), b(i))) \Rightarrow \\ \quad \quad \operatorname{sum\text{-}ag}(g, c, n) = \operatorname{opr}(g)(\operatorname{sum\text{-}ag}(g, a, n), \operatorname{sum\text{-}ag}(g, b, n)) \end{array}}}
\vnbprop{finsum-add-ag}{\vnbstmt{\begin{array}{@{}l@{}} \forall g.; \forall s \in \mathbf{Set}, f, h \in (s \to \operatorname{carr}(g)). \\ \quad (\operatorname{is\text{-}abelian\text{-}group}(g) \wedge (|s| \in \mathbb{N})) \Rightarrow  \\ \quad \quad \operatorname{finsum}(g, (\lambda z \in s.\; \operatorname{opr}(g)(f(z), h(z))), s) = \operatorname{opr}(g)(\operatorname{finsum}(g, f, s), \operatorname{finsum}(g, h, s)) \end{array}}}
\vnbprop{sum-ag-two-support-ind}{\vnbstmt{\begin{array}{@{}l@{}} \forall n \in \mathbb{N}, g, a, p, q. \\ \quad \operatorname{is\text{-}abelian\text{-}group}(g) \wedge \\ \quad p \in \mathrm{OS}(n) \wedge \\ \quad q \in \mathrm{OS}(n) \wedge \\ \quad \neg (p = q) \wedge \\ \quad a(p) \in \operatorname{carr}(g) \wedge \\ \quad a(q) \in \operatorname{carr}(g) \wedge \\ \quad \forall i \in \mathrm{OS}(n).\; ((\neg (i = p) \wedge \neg (i = q)) \Rightarrow (a(i) = \operatorname{iden}(g))) \Rightarrow \\ \quad \quad \operatorname{sum\text{-}ag}(g, a, n) = \operatorname{opr}(g)(a(p), a(q)) \end{array}}}
\vnbprop{finsum-two-support}{\vnbstmt{\begin{array}{@{}l@{}} \forall g.; \forall s \in \mathbf{Set}, f \in (s \to \operatorname{carr}(g)), a, b \in s. \\ \quad \operatorname{is\text{-}abelian\text{-}group}(g) \wedge \\ \quad |s| \in \mathbb{N} \wedge \\ \quad \neg (a = b) \wedge \\ \quad \forall j \in s.\; ((\neg (j = a) \wedge \neg (j = b)) \Rightarrow (f(j) = \operatorname{iden}(g))) \Rightarrow \\ \quad \quad \operatorname{finsum}(g, f, s) = \operatorname{opr}(g)(f(a), f(b)) \end{array}}}

\vnbbookfile{theorem-library/rake-finsum-laws2}{theorem-library/rake-finsum-laws2}
\noindent{\small theorem-library/rake-finsum-laws2.scm -- rake batch J, part 2: the three DISTRIBUTION laws of FINSUM (a ring element on either side, a module scalar) and the interval back-peel.  Part 1 is theorem-library/rake-finsum-laws.scm, which must load FIRST: this file cites its `enum-fam-value' and}\par
\vnbprop{sum-ag-ring-left-ind}{\vnbstmt{\begin{array}{@{}l@{}} \forall n \in \mathbb{N}, g.; \forall r \in \operatorname{carr}(g), a, b. \\ \quad \operatorname{is\text{-}ring}(g) \wedge \\ \quad \forall i \in \mathrm{OS}(n).\; (a(i) \in \operatorname{carr}(g)) \wedge \\ \quad \forall i \in \mathrm{OS}(n).\; (b(i) \simeq \operatorname{mul}(g)(r, a(i))) \Rightarrow \\ \quad \quad \operatorname{sum\text{-}ag}(\operatorname{ring\text{-}additive\text{-}ag}(g), b, n) = \operatorname{mul}(g)(r, \operatorname{sum\text{-}ag}(\operatorname{ring\text{-}additive\text{-}ag}(g), a, n)) \end{array}}}
\vnbprop{finsum-ring-distrib-left-gen}{\vnbstmt{\begin{array}{@{}l@{}} \forall g.; \forall r \in \operatorname{carr}(g), s \in \mathbf{Set}, f \in (s \to \operatorname{carr}(g)). \\ \quad (\operatorname{is\text{-}ring}(g) \wedge (|s| \in \mathbb{N})) \Rightarrow  \\ \quad \quad \operatorname{mul}(g)(r, \operatorname{finsum}(\operatorname{ring\text{-}additive\text{-}ag}(g), f, s)) \\ \quad \quad =\; \operatorname{finsum}(\operatorname{ring\text{-}additive\text{-}ag}(g), (\lambda z \in s.\; \operatorname{mul}(g)(r, f(z))), s) \end{array}}}
\vnbprop{sum-ag-ring-right-ind}{\vnbstmt{\begin{array}{@{}l@{}} \forall n \in \mathbb{N}, g.; \forall r \in \operatorname{carr}(g), a, b. \\ \quad \operatorname{is\text{-}ring}(g) \wedge \\ \quad \forall i \in \mathrm{OS}(n).\; (a(i) \in \operatorname{carr}(g)) \wedge \\ \quad \forall i \in \mathrm{OS}(n).\; (b(i) \simeq \operatorname{mul}(g)(a(i), r)) \Rightarrow \\ \quad \quad \operatorname{sum\text{-}ag}(\operatorname{ring\text{-}additive\text{-}ag}(g), b, n) = \operatorname{mul}(g)(\operatorname{sum\text{-}ag}(\operatorname{ring\text{-}additive\text{-}ag}(g), a, n), r) \end{array}}}
\vnbprop{finsum-ring-distrib-right-gen}{\vnbstmt{\begin{array}{@{}l@{}} \forall g.; \forall r \in \operatorname{carr}(g), s \in \mathbf{Set}, f \in (s \to \operatorname{carr}(g)). \\ \quad (\operatorname{is\text{-}ring}(g) \wedge (|s| \in \mathbb{N})) \Rightarrow  \\ \quad \quad \operatorname{mul}(g)(\operatorname{finsum}(\operatorname{ring\text{-}additive\text{-}ag}(g), f, s), r) \\ \quad \quad =\; \operatorname{finsum}(\operatorname{ring\text{-}additive\text{-}ag}(g), (\lambda z \in s.\; \operatorname{mul}(g)(f(z), r)), s) \end{array}}}
\vnbprop{finsum-interval-peel}{\vnbstmt{\begin{array}{@{}l@{}} \forall g.; \forall n \in \mathbb{N}, f \in (\operatorname{interval}(1, (n + 1)) \to \operatorname{carr}(g)). \\ \quad \operatorname{is\text{-}abelian\text{-}group}(g) \Rightarrow  \\ \quad \quad \operatorname{finsum}(g, f, \operatorname{interval}(1, (n + 1))) = \operatorname{opr}(g)(\operatorname{finsum}(g, f, \operatorname{interval}(1, n)), f((n + 1))) \end{array}}}

\vnbbookfile{theorem-library/rake-finsum-core}{theorem-library/rake-finsum-core}
\noindent{\small theorem-library/rake-finsum-core.scm -- the FINSUM CORE, proven. Rake batch M (2026-09-17).  Nine of the batch's ten leaves; the tenth (finsum-embed) is BLOCKED and its obstacle is stated at the end of this header.}\par
\vnbprop{finsum-ring-distrib-left}{\vnbstmt{\begin{array}{@{}l@{}} \forall g.; \forall r \in \operatorname{carr}(g), s \in \mathbf{Set}, f \in (s \to \operatorname{carr}(g)). \\ \quad (\operatorname{is\text{-}commutative\text{-}ring}(g) \wedge (|s| \in \mathbb{N})) \Rightarrow  \\ \quad \quad \operatorname{mul}(g)(r, \operatorname{finsum}(\operatorname{commutative\text{-}ring\text{-}additive\text{-}ag}(g), f, s)) \\ \quad \quad =\; \operatorname{finsum}(\operatorname{commutative\text{-}ring\text{-}additive\text{-}ag}(g), (\lambda z \in s.\; \operatorname{mul}(g)(r, f(z))), s) \end{array}}}
\vnbprop{finsum-ring-distrib-right}{\vnbstmt{\begin{array}{@{}l@{}} \forall g.; \forall r \in \operatorname{carr}(g), s \in \mathbf{Set}, f \in (s \to \operatorname{carr}(g)). \\ \quad (\operatorname{is\text{-}commutative\text{-}ring}(g) \wedge (|s| \in \mathbb{N})) \Rightarrow  \\ \quad \quad \operatorname{mul}(g)(\operatorname{finsum}(\operatorname{commutative\text{-}ring\text{-}additive\text{-}ag}(g), f, s), r) \\ \quad \quad =\; \operatorname{finsum}(\operatorname{commutative\text{-}ring\text{-}additive\text{-}ag}(g), (\lambda z \in s.\; \operatorname{mul}(g)(f(z), r)), s) \end{array}}}
\vnbprop{finsum-singleton}{\vnbstmt{\begin{array}{@{}l@{}} \forall g.; \forall x \in \mathbf{Set}, f \in (\{x\} \to \operatorname{carr}(g)). \\ \quad \operatorname{is\text{-}abelian\text{-}group}(g) \Rightarrow (\operatorname{finsum}(g, f, \{x\}) = f(x)) \end{array}}}
\vnbprop{finsum-ord-peel}{\vnbstmt{\begin{array}{@{}l@{}} \forall g.; \forall n \in \mathbb{N}, f \in (\mathrm{OS}((n + 1)) \to \operatorname{carr}(g)). \\ \quad \operatorname{is\text{-}abelian\text{-}group}(g) \Rightarrow  \\ \quad \quad \operatorname{finsum}(g, f, \mathrm{OS}((n + 1))) = \operatorname{opr}(g)(\operatorname{finsum}(g, f, \mathrm{OS}(n)), f(n)) \end{array}}}
\vnbprop{comm-monoid-opr-comm}{\vnbstmt{\begin{array}{@{}l@{}} \forall s.; \forall a, b \in \operatorname{carr}(s). \\ \quad \operatorname{is\text{-}comm\text{-}monoid}(s) \Rightarrow (\operatorname{opr}(s)(a, b) = \operatorname{opr}(s)(b, a)) \end{array}}}
\vnbprop{comm-monoid-assoc}{\vnbstmt{\begin{array}{@{}l@{}} \forall s.; \forall a, b, c \in \operatorname{carr}(s). \\ \quad \operatorname{is\text{-}comm\text{-}monoid}(s) \Rightarrow  \\ \quad \quad \operatorname{opr}(s)(\operatorname{opr}(s)(a, b), c) = \operatorname{opr}(s)(a, \operatorname{opr}(s)(b, c)) \end{array}}}
\vnbprop{comm-monoid-left-id}{\vnbstmt{\forall s.; \forall a \in \operatorname{carr}(s).\; (\operatorname{is\text{-}comm\text{-}monoid}(s) \Rightarrow (\operatorname{opr}(s)(\operatorname{iden}(s), a) = a))}}
\vnbprop{comm-monoid-opr-interchange}{\vnbstmt{\begin{array}{@{}l@{}} \forall s.; \forall a, b, c, d \in \operatorname{carr}(s). \\ \quad \operatorname{is\text{-}comm\text{-}monoid}(s) \Rightarrow  \\ \quad \quad \operatorname{opr}(s)(\operatorname{opr}(s)(a, b), \operatorname{opr}(s)(c, d)) = \operatorname{opr}(s)(\operatorname{opr}(s)(a, c), \operatorname{opr}(s)(b, d)) \end{array}}}
\vnbprop{sum-ag-comm-monoid-type-ptwise}{\vnbstmt{\begin{array}{@{}l@{}} \forall n \in \mathbb{N}, m, g. \\ \quad (\operatorname{is\text{-}comm\text{-}monoid}(m) \wedge \forall i \in \mathrm{OS}(n).\; (g(i) \in \operatorname{carr}(m))) \Rightarrow  \\ \quad \quad \operatorname{sum\text{-}ag}(m, g, n) \in \operatorname{carr}(m) \end{array}}}
\vnbprop{sum-ag-comm-monoid-add-ind}{\vnbstmt{\begin{array}{@{}l@{}} \forall n \in \mathbb{N}, m, a, b, c. \\ \quad \operatorname{is\text{-}comm\text{-}monoid}(m) \wedge \\ \quad \forall i \in \mathrm{OS}(n).\; (a(i) \in \operatorname{carr}(m)) \wedge \\ \quad \forall i \in \mathrm{OS}(n).\; (b(i) \in \operatorname{carr}(m)) \wedge \\ \quad \forall i \in \mathrm{OS}(n).\; (c(i) \simeq \operatorname{opr}(m)(a(i), b(i))) \Rightarrow \\ \quad \quad \operatorname{sum\text{-}ag}(m, c, n) = \operatorname{opr}(m)(\operatorname{sum\text{-}ag}(m, a, n), \operatorname{sum\text{-}ag}(m, b, n)) \end{array}}}
\vnbprop{finsum-add}{\vnbstmt{\begin{array}{@{}l@{}} \forall m.; \forall s \in \mathbf{Set}, f, h \in (s \to \operatorname{carr}(m)). \\ \quad (\operatorname{is\text{-}comm\text{-}monoid}(m) \wedge (|s| \in \mathbb{N})) \Rightarrow  \\ \quad \quad \operatorname{finsum}(m, (\lambda z \in s.\; \operatorname{opr}(m)(f(z), h(z))), s) = \operatorname{opr}(m)(\operatorname{finsum}(m, f, s), \operatorname{finsum}(m, h, s)) \end{array}}}
\vnbprop{sum-ag-fubini-ind}{\vnbstmt{\begin{array}{@{}l@{}} \forall m \in \mathbb{N}, g.; \forall n \in \mathbb{N}, f, \mathit{gg}, s, \mathit{cs}. \\ \quad \operatorname{is\text{-}abelian\text{-}group}(g) \wedge \\ \quad \forall i \in \mathrm{OS}(m), j \in \mathrm{OS}(n).\; (f(i)(j) \in \operatorname{carr}(g)) \wedge \\ \quad \forall i \in \mathrm{OS}(m), j \in \mathrm{OS}(n).\; (\operatorname{gg}(j)(i) \simeq f(i)(j)) \wedge \\ \quad \forall i \in \mathrm{OS}(m).\; (s(i) \simeq \operatorname{sum\text{-}ag}(g, f(i), n)) \wedge \\ \quad \forall j \in \mathrm{OS}(n).\; (\operatorname{cs}(j) \simeq \operatorname{sum\text{-}ag}(g, \operatorname{gg}(j), m)) \wedge \\ \quad \forall \mathit{fbi} \in \mathrm{OS}(m).\; (s(\mathit{fbi}) \in \operatorname{carr}(g)) \wedge \\ \quad \forall \mathit{fbj} \in \mathrm{OS}(n).\; (\operatorname{cs}(\mathit{fbj}) \in \operatorname{carr}(g)) \Rightarrow \\ \quad \quad \operatorname{sum\text{-}ag}(g, s, m) = \operatorname{sum\text{-}ag}(g, \mathit{cs}, n) \end{array}}}
\vnbprop{finsum-fubini}{\vnbstmt{\begin{array}{@{}l@{}} \forall g.; \forall x, y \in \mathbf{Set}, f \in ((x \times y) \to \operatorname{carr}(g)). \\ \quad (\operatorname{is\text{-}abelian\text{-}group}(g) \wedge ((|x| \in \mathbb{N}) \wedge (|y| \in \mathbb{N}))) \Rightarrow  \\ \quad \quad \operatorname{finsum}(g, (\lambda i \in x.\; \operatorname{finsum}(g, (\lambda j \in y.\; f(\langle i, j \rangle)), y)), x) \\ \quad \quad =\; \operatorname{finsum}(g, (\lambda j \in y.\; \operatorname{finsum}(g, (\lambda i \in x.\; f(\langle i, j \rangle)), x)), y) \end{array}}}
\vnbprop{ras-inv}{\vnbstmt{\forall a.\; (\operatorname{is\text{-}ring}(a) \Rightarrow (\operatorname{inv}(\operatorname{ring\text{-}additive\text{-}ag}(a)) = \operatorname{neg}(a)))}}
\vnbprop{commutative-ring-mul-comm}{\vnbstmt{\begin{array}{@{}l@{}} \forall s.; \forall a, b \in \operatorname{carr}(s). \\ \quad \operatorname{is\text{-}commutative\text{-}ring}(s) \Rightarrow (\operatorname{mul}(s)(a, b) = \operatorname{mul}(s)(b, a)) \end{array}}}
\vnbprop{ring-scalar-zz-nn}{\vnbstmt{\begin{array}{@{}l@{}} \forall k \in \mathbb{N}, g.; \forall r, a \in \operatorname{carr}(g). \\ \quad \operatorname{is\text{-}commutative\text{-}ring}(g) \Rightarrow  \\ \quad \quad \operatorname{mul}(g)(r, \operatorname{zz\text{-}act}(\operatorname{ring\text{-}additive\text{-}ag}(g), k, a)) = \operatorname{zz\text{-}act}(\operatorname{ring\text{-}additive\text{-}ag}(g), k, \operatorname{mul}(g)(r, a)) \end{array}}}
\vnbprop{finsum-ring-scalar-zz}{\vnbstmt{\begin{array}{@{}l@{}} \forall g.; \forall r \in \operatorname{carr}(g), c \in \mathbb{Z}, a \in \operatorname{carr}(g). \\ \quad \operatorname{is\text{-}commutative\text{-}ring}(g) \Rightarrow  \\ \quad \quad \operatorname{mul}(g)(r, \operatorname{zz\text{-}act}(\operatorname{commutative\text{-}ring\text{-}additive\text{-}ag}(g), c, a)) \\ \quad \quad =\; \operatorname{zz\text{-}act}(\operatorname{commutative\text{-}ring\text{-}additive\text{-}ag}(g), c, \operatorname{mul}(g)(r, a)) \end{array}}}

\vnbbookfile{theorem-library/rake-finsum-cm-ptwise}{theorem-library/rake-finsum-cm-ptwise}
\noindent{\small theorem-library/rake-finsum-cm-ptwise.scm -- finsum-comm-monoid-type-ptwise, the COMMUTATIVE-MONOID twin of `finsum-type-ptwise'.  Rake batch 5c-K.}\par
\vnbprop{finsum-comm-monoid-type-ptwise}{\vnbstmt{\begin{array}{@{}l@{}} \forall m.; \forall s \in \mathbf{Set}, f. \\ \quad (\operatorname{is\text{-}comm\text{-}monoid}(m) \wedge ((|s| \in \mathbb{N}) \wedge \forall z \in s.\; (f(z) \in \operatorname{carr}(m)))) \Rightarrow  \\ \quad \quad \operatorname{finsum}(m, f, s) \in \operatorname{carr}(m) \end{array}}}

\vnbbookfile{theorem-library/rake-finsum-fubini-c}{theorem-library/rake-finsum-fubini-c}
\noindent{\small rake-finsum-fubini-c.scm -- finsum-fubini-c, moved out of rake-finsum-laws.scm on 2026-09-17 (rake batch M) because it cites finsum-fubini, now PROVEN in rake-finsum-core.scm, which loads after rake-finsum-laws.  Window: after rake-finsum-core, before finsum-fiber (its earliest citer).  Helper prefix rkj-}\par
\vnbprop{finsum-fubini-c}{\vnbstmt{\begin{array}{@{}l@{}} \forall g.; \forall x, y \in \mathbf{Set}, f \in ((x \times y) \to \operatorname{carr}(g)). \\ \quad (\operatorname{is\text{-}abelian\text{-}group}(g) \wedge ((|x| \in \mathbb{N}) \wedge (|y| \in \mathbb{N}))) \Rightarrow  \\ \quad \quad \operatorname{finsum}(g, (\lambda i \in x.\; \operatorname{finsum}(g, (\lambda j \in y.\; f(\langle i, j \rangle)), y)), x) \\ \quad \quad =\; \operatorname{finsum}(g, (\lambda j \in y.\; \operatorname{finsum}(g, (\lambda i \in x.\; f(\langle i, j \rangle)), x)), y) \end{array}}}

\vnbbookfile{theorem-library/rake-combinatorics}{theorem-library/rake-combinatorics}
\noindent{\small rake-combinatorics.scm -- BATCH P of the 2026-09-17 rake: NN combinatorics, the ring powers, the powerset insert-split, and the two CHOOSE boundary values.  Seventeen asserted supports, PROVEN; every statement is copied LITERALLY from its site.}\par
\vnbprop{crmcm-op}{\vnbstmt{\begin{array}{@{}l@{}} \forall r. \\ \quad \operatorname{is\text{-}commutative\text{-}ring}(r) \Rightarrow  \\ \quad \quad \operatorname{opr}(\operatorname{commutative\text{-}ring\text{-}multiplicative\text{-}cm}(r)) = \operatorname{mul}(r) \end{array}}}
\vnbprop{crmcm-iden}{\vnbstmt{\begin{array}{@{}l@{}} \forall r. \\ \quad \operatorname{is\text{-}commutative\text{-}ring}(r) \Rightarrow  \\ \quad \quad \operatorname{iden}(\operatorname{commutative\text{-}ring\text{-}multiplicative\text{-}cm}(r)) = \operatorname{one}(r) \end{array}}}
\vnbprop{ring-power-succ-left}{\vnbstmt{\begin{array}{@{}l@{}} \forall r.; \forall x \in \operatorname{carr}(r), n \in \mathbb{N}. \\ \quad \operatorname{is\text{-}commutative\text{-}ring}(r) \Rightarrow  \\ \quad \quad \operatorname{ring\text{-}power}(r, x, (n + 1)) = \operatorname{mul}(r)(x, \operatorname{ring\text{-}power}(r, x, n)) \end{array}}}
\vnbprop{ring-power-one}{\vnbstmt{\begin{array}{@{}l@{}} \forall r.; \forall x \in \operatorname{carr}(r). \\ \quad \operatorname{is\text{-}commutative\text{-}ring}(r) \Rightarrow (\operatorname{ring\text{-}power}(r, x, 1) = x) \end{array}}}
\vnbprop{ring-power-succ}{\vnbstmt{\begin{array}{@{}l@{}} \forall r.; \forall x \in \operatorname{carr}(r), n \in \mathbb{N}. \\ \quad \operatorname{is\text{-}commutative\text{-}ring}(r) \Rightarrow  \\ \quad \quad \operatorname{ring\text{-}power}(r, x, (n + 1)) = \operatorname{mul}(r)(\operatorname{ring\text{-}power}(r, x, n), x) \end{array}}}
\vnbprop{ring-mul-interchange}{\vnbstmt{\begin{array}{@{}l@{}} \forall r.; \forall a, b, c, d \in \operatorname{carr}(r). \\ \quad \operatorname{is\text{-}commutative\text{-}ring}(r) \Rightarrow  \\ \quad \quad \operatorname{mul}(r)(\operatorname{mul}(r)(a, b), \operatorname{mul}(r)(c, d)) = \operatorname{mul}(r)(\operatorname{mul}(r)(a, c), \operatorname{mul}(r)(b, d)) \end{array}}}
\vnbprop{ring-power-add-ind}{\vnbstmt{\begin{array}{@{}l@{}} \forall k \in \mathbb{N}, r.; \forall x \in \operatorname{carr}(r), j \in \mathbb{N}. \\ \quad \operatorname{is\text{-}commutative\text{-}ring}(r) \Rightarrow  \\ \quad \quad \operatorname{ring\text{-}power}(r, x, (j + k)) = \operatorname{mul}(r)(\operatorname{ring\text{-}power}(r, x, j), \operatorname{ring\text{-}power}(r, x, k)) \end{array}}}
\vnbprop{ring-power-add}{\vnbstmt{\begin{array}{@{}l@{}} \forall r.; \forall x \in \operatorname{carr}(r), j, k \in \mathbb{N}. \\ \quad \operatorname{is\text{-}commutative\text{-}ring}(r) \Rightarrow  \\ \quad \quad \operatorname{ring\text{-}power}(r, x, (j + k)) = \operatorname{mul}(r)(\operatorname{ring\text{-}power}(r, x, j), \operatorname{ring\text{-}power}(r, x, k)) \end{array}}}
\vnbprop{ring-power-mult-ind}{\vnbstmt{\begin{array}{@{}l@{}} \forall n \in \mathbb{N}, r.; \forall x, y \in \operatorname{carr}(r). \\ \quad \operatorname{is\text{-}commutative\text{-}ring}(r) \Rightarrow  \\ \quad \quad \operatorname{ring\text{-}power}(r, \operatorname{mul}(r)(x, y), n) = \operatorname{mul}(r)(\operatorname{ring\text{-}power}(r, x, n), \operatorname{ring\text{-}power}(r, y, n)) \end{array}}}
\vnbprop{ring-power-mult}{\vnbstmt{\begin{array}{@{}l@{}} \forall r.; \forall x, y \in \operatorname{carr}(r), n \in \mathbb{N}. \\ \quad \operatorname{is\text{-}commutative\text{-}ring}(r) \Rightarrow  \\ \quad \quad \operatorname{ring\text{-}power}(r, \operatorname{mul}(r)(x, y), n) = \operatorname{mul}(r)(\operatorname{ring\text{-}power}(r, x, n), \operatorname{ring\text{-}power}(r, y, n)) \end{array}}}
\vnbprop{prod-ring-empty}{\vnbstmt{\begin{array}{@{}l@{}} \forall r, f. \\ \quad \operatorname{is\text{-}commutative\text{-}ring}(r) \Rightarrow (\operatorname{prod\text{-}ring}(r, f, \emptyset) = \operatorname{one}(r)) \end{array}}}
\vnbprop{prod-ring-insert}{\vnbstmt{\begin{array}{@{}l@{}} \forall r.; \forall x, k \in \mathbf{Set}, f \in ((x \cup \{k\}) \to \operatorname{carr}(r)). \\ \quad (\operatorname{is\text{-}commutative\text{-}ring}(r) \wedge ((|x| \in \mathbb{N}) \wedge \neg (k \in x))) \Rightarrow  \\ \quad \quad \operatorname{prod\text{-}ring}(r, f, (x \cup \{k\})) = \operatorname{mul}(r)(\operatorname{prod\text{-}ring}(r, f, x), f(k)) \end{array}}}
\vnbprop{prod-ring-singleton}{\vnbstmt{\begin{array}{@{}l@{}} \forall r.; \forall x \in \mathbf{Set}, f \in (\{x\} \to \operatorname{carr}(r)). \\ \quad \operatorname{is\text{-}commutative\text{-}ring}(r) \Rightarrow (\operatorname{prod\text{-}ring}(r, f, \{x\}) = f(x)) \end{array}}}

\vnbbookfile{theorem-library/rake-finsum-union}{theorem-library/rake-finsum-union}
\noindent{\small theorem-library/rake-finsum-union.scm -- the two set/card bricks that three rake batches built inline, the POINTWISE finsum recurrences they make possible, the DISJOINT-UNION law for FINSUM, and `finsum-embed'. Rake batch 5c, assignment 5c-M.  Helper prefix `r7m-'; every helper is}\par
\vnbprop{card-insert-converse}{\vnbstmt{\begin{array}{@{}l@{}} \forall a, x \in \mathbf{Set}, k \in \mathbb{N}. \\ \quad (\neg (x \in a) \wedge (|(a \cup \{x\})| = (k + 1))) \Rightarrow (|a| = k) \end{array}}}
\vnbprop{remove-restore}{\vnbstmt{\forall a \in \mathbf{Set}, x \in a.\; ((\operatorname{difference}(a, \{x\}) \cup \{x\}) = a)}}
\vnbprop{finsum-insert-ptwise}{\vnbstmt{\begin{array}{@{}l@{}} \forall g.; \forall s, k \in \mathbf{Set}, f. \\ \quad \operatorname{is\text{-}abelian\text{-}group}(g) \wedge \\ \quad |s| \in \mathbb{N} \wedge \\ \quad \neg (k \in s) \wedge \\ \quad \forall z \in (s \cup \{k\}).\; (f(z) \in \operatorname{carr}(g)) \Rightarrow \\ \quad \quad \operatorname{finsum}(g, f, (s \cup \{k\})) = \operatorname{opr}(g)(\operatorname{finsum}(g, f, s), f(k)) \end{array}}}
\vnbprop{finsum-union-disjoint}{\vnbstmt{\begin{array}{@{}l@{}} \forall g.; \forall s, t \in \mathbf{Set}, f. \\ \quad \operatorname{is\text{-}abelian\text{-}group}(g) \wedge \\ \quad |s| \in \mathbb{N} \wedge \\ \quad |t| \in \mathbb{N} \wedge \\ \quad \forall z \in t.\; \neg (z \in s) \wedge \\ \quad \forall z \in (s \cup t).\; (f(z) \in \operatorname{carr}(g)) \Rightarrow \\ \quad \quad \operatorname{finsum}(g, f, (s \cup t)) = \operatorname{opr}(g)(\operatorname{finsum}(g, f, s), \operatorname{finsum}(g, f, t)) \end{array}}}
\vnbprop{finsum-union-disjoint-fun}{\vnbstmt{\begin{array}{@{}l@{}} \forall g.; \forall s, t \in \mathbf{Set}, f \in ((s \cup t) \to \operatorname{carr}(g)). \\ \quad (\operatorname{is\text{-}abelian\text{-}group}(g) \wedge ((|s| \in \mathbb{N}) \wedge ((|t| \in \mathbb{N}) \wedge ((s \cap t) = \emptyset)))) \Rightarrow  \\ \quad \quad \operatorname{finsum}(g, f, (s \cup t)) = \operatorname{opr}(g)(\operatorname{finsum}(g, f, s), \operatorname{finsum}(g, f, t)) \end{array}}}
\vnbprop{finsum-all-id-ptwise}{\vnbstmt{\begin{array}{@{}l@{}} \forall g.; \forall s \in \mathbf{Set}, f. \\ \quad (\operatorname{is\text{-}abelian\text{-}group}(g) \wedge ((|s| \in \mathbb{N}) \wedge \forall z \in s.\; (f(z) = \operatorname{iden}(g)))) \Rightarrow  \\ \quad \quad \operatorname{finsum}(g, f, s) = \operatorname{iden}(g) \end{array}}}
\vnbprop{finsum-embed}{\vnbstmt{\begin{array}{@{}l@{}} \forall g.; \forall s, a \in \mathbf{Set}, f \in (a \to \operatorname{carr}(g)). \\ \quad \operatorname{is\text{-}abelian\text{-}group}(g) \wedge \\ \quad |s| \in \mathbb{N} \wedge \\ \quad |a| \in \mathbb{N} \wedge \\ \quad s \subseteq a \wedge \\ \quad \forall z \in a.\; (\neg (z \in s) \Rightarrow (f(z) = \operatorname{iden}(g))) \Rightarrow \\ \quad \quad \operatorname{finsum}(g, f, a) = \operatorname{finsum}(g, f, s) \end{array}}}

\vnbbookfile{theorem-library/rake-finsum-cm-union}{theorem-library/rake-finsum-cm-union}
\noindent{\small theorem-library/rake-finsum-cm-union.scm -- the COMMUTATIVE-MONOID mirror of the FINSUM congruence / insertion / disjoint-union block, plus the two set identities that a peeling induction over POWER(X) needs. Rake batch 5c, assignment 5c-T.  Helper prefix `r7t-'; every helper is}\par
\vnbprop{finsum-cm-congruence-q}{\vnbstmt{\begin{array}{@{}l@{}} \forall m.; \forall s \in \mathbf{Set}, f, g. \\ \quad (\operatorname{is\text{-}comm\text{-}monoid}(m) \wedge ((|s| \in \mathbb{N}) \wedge \forall z \in s.\; (f(z) = g(z)))) \Rightarrow  \\ \quad \quad (\operatorname{finsum}(m, f, s) \simeq \operatorname{finsum}(m, g, s)) \end{array}}}
\vnbprop{finsum-cm-insert-ptwise}{\vnbstmt{\begin{array}{@{}l@{}} \forall m.; \forall s, k \in \mathbf{Set}, f. \\ \quad \operatorname{is\text{-}comm\text{-}monoid}(m) \wedge \\ \quad |s| \in \mathbb{N} \wedge \\ \quad \neg (k \in s) \wedge \\ \quad \forall z \in (s \cup \{k\}).\; (f(z) \in \operatorname{carr}(m)) \Rightarrow \\ \quad \quad \operatorname{finsum}(m, f, (s \cup \{k\})) = \operatorname{opr}(m)(\operatorname{finsum}(m, f, s), f(k)) \end{array}}}
\vnbprop{finsum-cm-union-disjoint}{\vnbstmt{\begin{array}{@{}l@{}} \forall m.; \forall s, t \in \mathbf{Set}, f. \\ \quad \operatorname{is\text{-}comm\text{-}monoid}(m) \wedge \\ \quad |s| \in \mathbb{N} \wedge \\ \quad |t| \in \mathbb{N} \wedge \\ \quad \forall z \in t.\; \neg (z \in s) \wedge \\ \quad \forall z \in (s \cup t).\; (f(z) \in \operatorname{carr}(m)) \Rightarrow \\ \quad \quad \operatorname{finsum}(m, f, (s \cup t)) = \operatorname{opr}(m)(\operatorname{finsum}(m, f, s), \operatorname{finsum}(m, f, t)) \end{array}}}
\vnbprop{insert-difference-outside}{\vnbstmt{\begin{array}{@{}l@{}} \forall x, s.; \forall k \in \mathbf{Set}. \\ \quad \neg (k \in s) \Rightarrow  \\ \quad \quad \operatorname{difference}((x \cup \{k\}), s) = (\operatorname{difference}(x, s) \cup \{k\}) \end{array}}}
\vnbprop{insert-difference-cancel}{\vnbstmt{\begin{array}{@{}l@{}} \forall x, s.; \forall k \in \mathbf{Set}. \\ \quad \neg (k \in x) \Rightarrow  \\ \quad \quad \operatorname{difference}((x \cup \{k\}), (s \cup \{k\})) = \operatorname{difference}(x, s) \end{array}}}

\vnbbookfile{theorem-library/rake-sum-set-defined}{theorem-library/rake-sum-set-defined}
\noindent{\small theorem-library/rake-sum-set-defined.scm -- SUM-SET DEFINED, and its four characterising laws PROVEN from the definition.  Rake batch 5c, assignment 5c-S (the user's decision (B) of 2026-09-18).  Helper prefix `r7s-'.}\par
\vnbprop{sum-set-empty}{\vnbstmt{\forall r, f.\; (\operatorname{is\text{-}ring}(r) \Rightarrow (\operatorname{sum\text{-}set}(r, \emptyset, f) \simeq \operatorname{zero}(r)))}}
\vnbprop{finsum-singleton-ptwise}{\vnbstmt{\begin{array}{@{}l@{}} \forall g.; \forall x \in \mathbf{Set}, f. \\ \quad (\operatorname{is\text{-}abelian\text{-}group}(g) \wedge (f(x) \in \operatorname{carr}(g))) \Rightarrow (\operatorname{finsum}(g, f, \{x\}) = f(x)) \end{array}}}
\vnbprop{sum-set-singleton-defined}{\vnbstmt{\begin{array}{@{}l@{}} \forall r.; \forall x \in \mathbf{Set}, f. \\ \quad (\operatorname{is\text{-}ring}(r) \wedge (f(x) \in \operatorname{carr}(r))) \Rightarrow (\operatorname{sum\text{-}set}(r, \{x\}, f) = f(x)) \end{array}}}
\vnbprop{sum-set-type-defined}{\vnbstmt{\begin{array}{@{}l@{}} \forall r, x, s, f. \\ \quad (\operatorname{is\text{-}ring}(r) \wedge ((x \in \mathbf{Set}) \wedge ((s \subseteq x) \wedge ((f \in (x \to \operatorname{carr}(r))) \wedge (|s| \in \mathbb{N}))))) \Rightarrow  \\ \quad \quad \operatorname{sum\text{-}set}(r, s, f) \in \operatorname{carr}(r) \end{array}}}
\vnbprop{sum-set-disjoint-union-defined}{\vnbstmt{\begin{array}{@{}l@{}} \forall r, x, a, b, f. \\ \quad \operatorname{is\text{-}ring}(r) \wedge \\ \quad x \in \mathbf{Set} \wedge \\ \quad a \in \mathbf{Set} \wedge \\ \quad b \in \mathbf{Set} \wedge \\ \quad (a \cap b) = \emptyset \wedge \\ \quad (a \cup b) \subseteq x \wedge \\ \quad f \in (x \to \operatorname{carr}(r)) \wedge \\ \quad |a| \in \mathbb{N} \wedge \\ \quad |b| \in \mathbb{N} \Rightarrow \\ \quad \quad \operatorname{sum\text{-}set}(r, (a \cup b), f) = \operatorname{add}(r)(\operatorname{sum\text{-}set}(r, a, f), \operatorname{sum\text{-}set}(r, b, f)) \end{array}}}

\vnbbookfile{theorem-library/rake-prod-set-defined}{theorem-library/rake-prod-set-defined}
\noindent{\small theorem-library/rake-prod-set-defined.scm -- PROD-SET DEFINED, and its four characterising laws PROVEN from the definition.  Rake batch 5c, assignment 5c-P (the user's decision of 2026-09-19).  Helper prefix `psd-'.}\par
\vnbprop{prod-set-empty}{\vnbstmt{\forall m, f.\; (\operatorname{prod\text{-}set}(m, \emptyset, f) \simeq \operatorname{iden}(m))}}
\vnbprop{finsum-cm-singleton-ptwise}{\vnbstmt{\begin{array}{@{}l@{}} \forall m.; \forall x \in \mathbf{Set}, f. \\ \quad (\operatorname{is\text{-}comm\text{-}monoid}(m) \wedge (f(x) \in \operatorname{carr}(m))) \Rightarrow (\operatorname{finsum}(m, f, \{x\}) = f(x)) \end{array}}}
\vnbprop{prod-set-singleton-defined}{\vnbstmt{\begin{array}{@{}l@{}} \forall m.; \forall x \in \mathbf{Set}, f. \\ \quad (\operatorname{is\text{-}comm\text{-}monoid}(m) \wedge (f(x) \in \operatorname{carr}(m))) \Rightarrow (\operatorname{prod\text{-}set}(m, \{x\}, f) = f(x)) \end{array}}}
\vnbprop{prod-set-type-defined}{\vnbstmt{\begin{array}{@{}l@{}} \forall m, x, s, f. \\ \quad (\operatorname{is\text{-}comm\text{-}monoid}(m) \wedge ((x \in \mathbf{Set}) \wedge ((s \subseteq x) \wedge ((f \in (x \to \operatorname{carr}(m))) \wedge (|s| \in \mathbb{N}))))) \Rightarrow  \\ \quad \quad \operatorname{prod\text{-}set}(m, s, f) \in \operatorname{carr}(m) \end{array}}}
\vnbprop{prod-set-disjoint-union-defined}{\vnbstmt{\begin{array}{@{}l@{}} \forall m, a, b, f. \\ \quad \operatorname{is\text{-}comm\text{-}monoid}(m) \wedge \\ \quad a \in \mathbf{Set} \wedge \\ \quad b \in \mathbf{Set} \wedge \\ \quad (a \cap b) = \emptyset \wedge \\ \quad f \in ((a \cup b) \to \operatorname{carr}(m)) \wedge \\ \quad |a| \in \mathbb{N} \wedge \\ \quad |b| \in \mathbb{N} \Rightarrow \\ \quad \quad \operatorname{prod\text{-}set}(m, (a \cup b), f) = \operatorname{opr}(m)(\operatorname{prod\text{-}set}(m, a, f), \operatorname{prod\text{-}set}(m, b, f)) \end{array}}}

\vnbbookfile{theorem-library/rake-sum-set-scalar}{theorem-library/rake-sum-set-scalar}
\noindent{\small rake-sum-set-scalar.scm -- the two scalar pull-outs for SUM-SET, PROVEN.}\par
\vnbprop{sum-set-left-scalar}{\vnbstmt{\begin{array}{@{}l@{}} \forall s.; \forall a \in \operatorname{carr}(s), x \in \mathbf{Set}, f \in (x \to \operatorname{carr}(s)). \\ \quad (\operatorname{is\text{-}ring}(s) \wedge (|x| \in \mathbb{N})) \Rightarrow  \\ \quad \quad \operatorname{mul}(s)(a, \operatorname{sum\text{-}set}(s, x, f)) = \operatorname{sum\text{-}set}(s, x, (\lambda z \in x.\; \operatorname{mul}(s)(a, f(z)))) \end{array}}}
\vnbprop{sum-set-right-scalar}{\vnbstmt{\begin{array}{@{}l@{}} \forall s.; \forall b \in \operatorname{carr}(s), x \in \mathbf{Set}, f \in (x \to \operatorname{carr}(s)). \\ \quad (\operatorname{is\text{-}ring}(s) \wedge (|x| \in \mathbb{N})) \Rightarrow  \\ \quad \quad \operatorname{mul}(s)(\operatorname{sum\text{-}set}(s, x, f), b) = \operatorname{sum\text{-}set}(s, x, (\lambda z \in x.\; \operatorname{mul}(s)(f(z), b))) \end{array}}}

\vnbbookfile{theorem-library/fin-subset-monoid}{theorem-library/fin-subset-monoid}
\noindent{\small fin-subset-monoid.scm -- FIN-SUBSETS(a) is a commutative monoid under union.}\par
\par\medskip\noindent\textbf{Definition} (\texttt{fin-subset-monoid}).\;$\operatorname{fin\text{-}subset\text{-}monoid}(\operatorname{a\_}) :=$
\begin{equation*}
\langle \operatorname{fin\text{-}subsets}(\operatorname{a\_}), (\lambda \langle \operatorname{x\_}, \operatorname{y\_} \rangle \in (\operatorname{fin\text{-}subsets}(\operatorname{a\_}) \times \operatorname{fin\text{-}subsets}(\operatorname{a\_})).\; (\operatorname{x\_} \cup \operatorname{y\_})), \emptyset \rangle
\end{equation*}

\vnbprop{fsm-carr}{\vnbstmt{\forall a \in \mathbf{Set}.\; (\operatorname{carr}(\operatorname{fin\text{-}subset\text{-}monoid}(a)) \simeq \operatorname{fin\text{-}subsets}(a))}}
\vnbprop{fsm-opr}{\vnbstmt{\begin{array}{@{}l@{}} \forall a \in \mathbf{Set}. \\ \quad (\operatorname{opr}(\operatorname{fin\text{-}subset\text{-}monoid}(a)) \simeq (\lambda \langle \operatorname{x\_}, \operatorname{y\_} \rangle \in (\operatorname{fin\text{-}subsets}(a) \times \operatorname{fin\text{-}subsets}(a)).\; (\operatorname{x\_} \cup \operatorname{y\_}))) \end{array}}}
\vnbprop{fsm-iden}{\vnbstmt{\forall a \in \mathbf{Set}.\; (\operatorname{iden}(\operatorname{fin\text{-}subset\text{-}monoid}(a)) \simeq \emptyset)}}
\vnbprop{fin-subset-monoid-is-comm-monoid}{\vnbstmt{\forall a \in \mathbf{Set}.\; \operatorname{is\text{-}comm\text{-}monoid}(\operatorname{fin\text{-}subset\text{-}monoid}(a))}}

\vnbbookfile{theorem-library/rake-algebra3}{theorem-library/rake-algebra3}
\noindent{\small theorem-library/rake-algebra3.scm -- rake batch 5, batch V (algebra leftovers).}\par
\vnbprop{qq-nonneg-is-fraction}{\vnbstmt{\begin{array}{@{}l@{}} \forall q \in \mathbb{Q}. \\ \quad (0 \leq q) \Rightarrow \exists a, b \in \mathbb{N}.\; (\neg (b = 0) \wedge ((q \cdot b) = a)) \end{array}}}
\vnbprop{card-bijection-eq}{\vnbstmt{\begin{array}{@{}l@{}} \forall m, d, i. \\ \quad ((m \in \mathbf{Set}) \wedge ((|m| \in \mathbb{N}) \wedge (i \in \mathrm{Bij}(m, d)))) \Rightarrow (|d| = |m|) \end{array}}}
\vnbprop{finsum-reindex}{\vnbstmt{\begin{array}{@{}l@{}} \forall m.; \forall s, t \in \mathbf{Set}, i \in \mathrm{Bij}(t, s), f \in (s \to \operatorname{carr}(m)). \\ \quad (\operatorname{is\text{-}comm\text{-}monoid}(m) \wedge ((|s| \in \mathbb{N}) \wedge (|t| \in \mathbb{N}))) \Rightarrow  \\ \quad \quad \operatorname{finsum}(m, f, s) = \operatorname{finsum}(m, (\lambda z \in t.\; f(i(z))), t) \end{array}}}
\vnbprop{finsum-reindex-ag}{\vnbstmt{\begin{array}{@{}l@{}} \forall g.; \forall s, t \in \mathbf{Set}, i \in \mathrm{Bij}(t, s), f \in (s \to \operatorname{carr}(g)). \\ \quad (\operatorname{is\text{-}abelian\text{-}group}(g) \wedge ((|s| \in \mathbb{N}) \wedge (|t| \in \mathbb{N}))) \Rightarrow  \\ \quad \quad \operatorname{finsum}(g, f, s) = \operatorname{finsum}(g, (\lambda z \in t.\; f(i(z))), t) \end{array}}}

\vnbbookfile{theorem-library/poly-tail-zero}{theorem-library/poly-tail-zero}
\noindent{\small poly-tail-zero.scm -- "a polynomial is a sequence that is eventually zero".}\par
\vnbprop{poly-tail-zero-fwd}{\vnbstmt{\begin{array}{@{}l@{}} \forall a, f. \\ \quad (\operatorname{is\text{-}ring}(a) \wedge (f \in \operatorname{carr}(\operatorname{poly}(a)))) \Rightarrow  \\ \quad \quad \exists n \in \mathbb{N}.\; \forall k \in \mathbb{N}.\; ((n \leq k) \Rightarrow (f(k) = \operatorname{zero}(a))) \end{array}}}
\vnbprop{poly-tail-zero-bwd}{\vnbstmt{\begin{array}{@{}l@{}} \forall a, f. \\ \quad \operatorname{is\text{-}ring}(a) \wedge \\ \quad f \in \operatorname{carr}(a)^{\mathbb{N}} \wedge \\ \quad \exists n \in \mathbb{N}.\; \forall k \in \mathbb{N}.\; ((n \leq k) \Rightarrow (f(k) = \operatorname{zero}(a))) \Rightarrow \\ \quad \quad f \in \operatorname{carr}(\operatorname{poly}(a)) \end{array}}}
\vnbprop{poly-tail-zero}{\vnbstmt{\begin{array}{@{}l@{}} \forall a, f. \\ \quad \operatorname{is\text{-}ring}(a) \Rightarrow  \\ \quad \quad f \in \operatorname{carr}(\operatorname{poly}(a)) \\ \quad \quad \Leftrightarrow\; ((f \in \operatorname{carr}(a)^{\mathbb{N}}) \wedge \exists n \in \mathbb{N}.\; \forall k \in \mathbb{N}.\; ((n \leq k) \Rightarrow (f(k) = \operatorname{zero}(a)))) \end{array}}}

\vnbbookfile{theorem-library/finsum-fiber}{theorem-library/finsum-fiber}
\noindent{\small finsum-fiber.scm -- FIBERED FUBINI: group the terms of a finite sum by the fibers of an arbitrary map, and the CARTESIAN projection brick it needs.}\par
\vnbprop{cartesian-nth}{\vnbstmt{\begin{array}{@{}l@{}} \forall x, a, b. \\ \quad (x \in (a \times b)) \Rightarrow ((\operatorname{nth}(1, x) \in a) \wedge (\operatorname{nth}(2, x) \in b)) \end{array}}}
\vnbprop{finsum-fiber}{\vnbstmt{\begin{array}{@{}l@{}} \forall g.; \forall s, t \in \mathbf{Set}, i \in (s \to t), f \in (s \to \operatorname{carr}(g)). \\ \quad (\operatorname{is\text{-}abelian\text{-}group}(g) \wedge ((|s| \in \mathbb{N}) \wedge (|t| \in \mathbb{N}))) \Rightarrow  \\ \quad \quad \operatorname{finsum}(g, f, s) = \operatorname{finsum}(g, (\lambda a \in t.\; \operatorname{finsum}(g, f, \{b \in s : (i(b) = a)\})), t) \end{array}}}

\vnbbookfile{theorem-library/monalg-laws}{theorem-library/monalg-laws}
\noindent{\small monalg-laws.scm -- the monoid-algebra laws that go through without reindexing, PROVEN.}\par
\vnbprop{monoid-carrier-is-set}{\vnbstmt{\forall m.\; (\operatorname{is\text{-}monoid}(m) \Rightarrow (\operatorname{carr}(m) \in \mathbf{Set}))}}
\vnbprop{monalg-add-apply}{\vnbstmt{\begin{array}{@{}l@{}} \forall a, m, f, g.; \forall x \in \operatorname{carr}(m). \\ \quad (\operatorname{monalg\text{-}add}(a, m, f, g)(x) \simeq \operatorname{add}(a)(f(x), g(x))) \end{array}}}
\vnbprop{monalg-add-fun}{\vnbstmt{\begin{array}{@{}l@{}} \forall a, m, f, g. \\ \quad (\operatorname{is\text{-}ring}(a) \wedge (\operatorname{is\text{-}monoid}(m) \wedge ((f \in \operatorname{finsupp}(a, m)) \wedge (g \in \operatorname{finsupp}(a, m))))) \Rightarrow  \\ \quad \quad \operatorname{monalg\text{-}add}(a, m, f, g) \in \operatorname{finsupp}(a, m) \end{array}}}

\vnbbookfile{theorem-library/poly-zero}{theorem-library/poly-zero}
\noindent{\small poly-zero.scm -- the ZERO of a monoid algebra, pointwise, and the EXTENSIONALITY BRIDGE for polynomials.  Nothing here is asserted; every statement is `proven modulo 0'.}\par
\vnbprop{monalg-zero-apply}{\vnbstmt{\forall a, m.; \forall x \in \operatorname{carr}(m).\; (\operatorname{monalg\text{-}zero}(a, m)(x) \simeq \operatorname{zero}(a))}}
\vnbprop{monalg-zero-fun}{\vnbstmt{\begin{array}{@{}l@{}} \forall a, m. \\ \quad ((\operatorname{carr}(m) \in \mathbf{Set}) \wedge (\operatorname{zero}(a) \in \operatorname{carr}(a))) \Rightarrow  \\ \quad \quad \operatorname{monalg\text{-}zero}(a, m) \in (\operatorname{carr}(m) \to \operatorname{carr}(a)) \end{array}}}
\vnbprop{poly-zero-unfold}{\vnbstmt{\forall a.\; (\operatorname{is\text{-}ring}(a) \Rightarrow (\operatorname{zero}(\operatorname{poly}(a)) \simeq \operatorname{monalg\text{-}zero}(a, \operatorname{nn\text{-}add\text{-}monoid})))}}
\vnbprop{poly-zero-apply}{\vnbstmt{\forall a, k.\; ((\operatorname{is\text{-}ring}(a) \wedge (k \in \mathbb{N})) \Rightarrow (\operatorname{zero}(\operatorname{poly}(a))(k) = \operatorname{zero}(a)))}}
\vnbprop{poly-zero-fun}{\vnbstmt{\forall a.\; (\operatorname{is\text{-}ring}(a) \Rightarrow (\operatorname{zero}(\operatorname{poly}(a)) \in (\mathbb{N} \to \operatorname{carr}(a))))}}
\vnbprop{poly-zero-iff}{\vnbstmt{\begin{array}{@{}l@{}} \forall a, p. \\ \quad (\operatorname{is\text{-}ring}(a) \wedge (p \in \operatorname{carr}(\operatorname{poly}(a)))) \Rightarrow  \\ \quad \quad (p = \operatorname{zero}(\operatorname{poly}(a))) \Leftrightarrow \forall k \in \mathbb{N}.\; (p(k) = \operatorname{zero}(a)) \end{array}}}

\vnbbookfile{theorem-library/monalg-is-ring}{theorem-library/monalg-is-ring}
\noindent{\small monalg-is-ring.scm -- the monoid algebra A[M] is a ring (structure-library/polynomial.scm's umbrella seed), PROVEN, with every additive law, both distributive laws, both unit laws and the typing of the convolution proven on the way.  The one law still chained to its seed is the associativity of convolution, ...}\par
\vnbprop{monalg-carr}{\vnbstmt{\begin{array}{@{}l@{}} \forall a, m. \\ \quad (\operatorname{is\text{-}ring}(a) \wedge \operatorname{is\text{-}monoid}(m)) \Rightarrow (\operatorname{carr}(\operatorname{monalg}(a, m)) \simeq \operatorname{finsupp}(a, m)) \end{array}}}
\vnbprop{monalg-add-op}{\vnbstmt{\begin{array}{@{}l@{}} \forall a, m. \\ \quad (\operatorname{is\text{-}ring}(a) \wedge \operatorname{is\text{-}monoid}(m)) \Rightarrow  \\ \quad \quad (\operatorname{add}(\operatorname{monalg}(a, m)) \simeq (\lambda \langle f, g \rangle \in (\operatorname{finsupp}(a, m) \times \operatorname{finsupp}(a, m)).\; \operatorname{monalg\text{-}add}(a, m, f, g))) \end{array}}}
\vnbprop{monalg-mul-op}{\vnbstmt{\begin{array}{@{}l@{}} \forall a, m. \\ \quad (\operatorname{is\text{-}ring}(a) \wedge \operatorname{is\text{-}monoid}(m)) \Rightarrow  \\ \quad \quad (\operatorname{mul}(\operatorname{monalg}(a, m)) \simeq (\lambda \langle f, g \rangle \in (\operatorname{finsupp}(a, m) \times \operatorname{finsupp}(a, m)).\; \operatorname{monalg\text{-}mul}(a, m, f, g))) \end{array}}}
\vnbprop{monalg-neg-op}{\vnbstmt{\begin{array}{@{}l@{}} \forall a, m. \\ \quad (\operatorname{is\text{-}ring}(a) \wedge \operatorname{is\text{-}monoid}(m)) \Rightarrow  \\ \quad \quad (\operatorname{neg}(\operatorname{monalg}(a, m)) \simeq (\lambda f \in \operatorname{finsupp}(a, m).\; \operatorname{monalg\text{-}neg}(a, m, f))) \end{array}}}
\vnbprop{monalg-zero-slot}{\vnbstmt{\begin{array}{@{}l@{}} \forall a, m. \\ \quad (\operatorname{is\text{-}ring}(a) \wedge \operatorname{is\text{-}monoid}(m)) \Rightarrow (\operatorname{zero}(\operatorname{monalg}(a, m)) \simeq \operatorname{monalg\text{-}zero}(a, m)) \end{array}}}
\vnbprop{monalg-one-slot}{\vnbstmt{\begin{array}{@{}l@{}} \forall a, m. \\ \quad (\operatorname{is\text{-}ring}(a) \wedge \operatorname{is\text{-}monoid}(m)) \Rightarrow (\operatorname{one}(\operatorname{monalg}(a, m)) \simeq \operatorname{monalg\text{-}one}(a, m)) \end{array}}}
\vnbprop{finsupp-in-set}{\vnbstmt{\forall a, m.\; ((\operatorname{is\text{-}ring}(a) \wedge \operatorname{is\text{-}monoid}(m)) \Rightarrow (\operatorname{finsupp}(a, m) \in \mathbf{Set}))}}
\vnbprop{monalg-ext}{\vnbstmt{\begin{array}{@{}l@{}} \forall a, m, f, g. \\ \quad ((f \in \operatorname{finsupp}(a, m)) \wedge ((g \in \operatorname{finsupp}(a, m)) \wedge \forall x \in \operatorname{carr}(m).\; (f(x) = g(x)))) \Rightarrow  \\ \quad \quad f = g \end{array}}}
\vnbprop{monalg-neg-apply}{\vnbstmt{\forall a, m, f.; \forall x \in \operatorname{carr}(m).\; (\operatorname{monalg\text{-}neg}(a, m, f)(x) \simeq \operatorname{neg}(a)(f(x)))}}
\vnbprop{monalg-one-apply}{\vnbstmt{\begin{array}{@{}l@{}} \forall a, m.; \forall x \in \operatorname{carr}(m). \\ \quad (\operatorname{monalg\text{-}one}(a, m)(x) \simeq \left(\text{if }(x = \operatorname{iden}(m))\text{ then }\operatorname{one}(a)\text{ else }\operatorname{zero}(a)\right)) \end{array}}}
\vnbprop{ring-neg-zero}{\vnbstmt{\forall a.\; (\operatorname{is\text{-}ring}(a) \Rightarrow (\operatorname{neg}(a)(\operatorname{zero}(a)) = \operatorname{zero}(a)))}}
\vnbprop{monalg-neg-fun}{\vnbstmt{\begin{array}{@{}l@{}} \forall a, m, f. \\ \quad (\operatorname{is\text{-}ring}(a) \wedge (\operatorname{is\text{-}monoid}(m) \wedge (f \in \operatorname{finsupp}(a, m)))) \Rightarrow  \\ \quad \quad \operatorname{monalg\text{-}neg}(a, m, f) \in \operatorname{finsupp}(a, m) \end{array}}}
\vnbprop{monalg-zero-in}{\vnbstmt{\begin{array}{@{}l@{}} \forall a, m. \\ \quad (\operatorname{is\text{-}ring}(a) \wedge \operatorname{is\text{-}monoid}(m)) \Rightarrow (\operatorname{monalg\text{-}zero}(a, m) \in \operatorname{finsupp}(a, m)) \end{array}}}
\vnbprop{monalg-one-in}{\vnbstmt{\begin{array}{@{}l@{}} \forall a, m. \\ \quad (\operatorname{is\text{-}ring}(a) \wedge \operatorname{is\text{-}monoid}(m)) \Rightarrow (\operatorname{monalg\text{-}one}(a, m) \in \operatorname{finsupp}(a, m)) \end{array}}}
\vnbprop{monalg-add-assoc}{\vnbstmt{\begin{array}{@{}l@{}} \forall a, m, f, g, h. \\ \quad (\operatorname{is\text{-}ring}(a) \wedge (\operatorname{is\text{-}monoid}(m) \wedge ((f \in \operatorname{finsupp}(a, m)) \wedge ((g \in \operatorname{finsupp}(a, m)) \wedge (h \in \operatorname{finsupp}(a, m)))))) \Rightarrow  \\ \quad \quad \operatorname{monalg\text{-}add}(a, m, \operatorname{monalg\text{-}add}(a, m, f, g), h) = \operatorname{monalg\text{-}add}(a, m, f, \operatorname{monalg\text{-}add}(a, m, g, h)) \end{array}}}
\vnbprop{monalg-add-comm}{\vnbstmt{\begin{array}{@{}l@{}} \forall a, m, f, g. \\ \quad (\operatorname{is\text{-}ring}(a) \wedge (\operatorname{is\text{-}monoid}(m) \wedge ((f \in \operatorname{finsupp}(a, m)) \wedge (g \in \operatorname{finsupp}(a, m))))) \Rightarrow  \\ \quad \quad \operatorname{monalg\text{-}add}(a, m, f, g) = \operatorname{monalg\text{-}add}(a, m, g, f) \end{array}}}
\vnbprop{monalg-add-zero-left}{\vnbstmt{\begin{array}{@{}l@{}} \forall a, m, f. \\ \quad (\operatorname{is\text{-}ring}(a) \wedge (\operatorname{is\text{-}monoid}(m) \wedge (f \in \operatorname{finsupp}(a, m)))) \Rightarrow  \\ \quad \quad \operatorname{monalg\text{-}add}(a, m, \operatorname{monalg\text{-}zero}(a, m), f) = f \end{array}}}
\vnbprop{monalg-add-zero-right}{\vnbstmt{\begin{array}{@{}l@{}} \forall a, m, f. \\ \quad (\operatorname{is\text{-}ring}(a) \wedge (\operatorname{is\text{-}monoid}(m) \wedge (f \in \operatorname{finsupp}(a, m)))) \Rightarrow  \\ \quad \quad \operatorname{monalg\text{-}add}(a, m, f, \operatorname{monalg\text{-}zero}(a, m)) = f \end{array}}}
\vnbprop{monalg-add-neg-left}{\vnbstmt{\begin{array}{@{}l@{}} \forall a, m, f. \\ \quad (\operatorname{is\text{-}ring}(a) \wedge (\operatorname{is\text{-}monoid}(m) \wedge (f \in \operatorname{finsupp}(a, m)))) \Rightarrow  \\ \quad \quad \operatorname{monalg\text{-}add}(a, m, \operatorname{monalg\text{-}neg}(a, m, f), f) = \operatorname{monalg\text{-}zero}(a, m) \end{array}}}
\vnbprop{monalg-add-neg-right}{\vnbstmt{\begin{array}{@{}l@{}} \forall a, m, f. \\ \quad (\operatorname{is\text{-}ring}(a) \wedge (\operatorname{is\text{-}monoid}(m) \wedge (f \in \operatorname{finsupp}(a, m)))) \Rightarrow  \\ \quad \quad \operatorname{monalg\text{-}add}(a, m, f, \operatorname{monalg\text{-}neg}(a, m, f)) = \operatorname{monalg\text{-}zero}(a, m) \end{array}}}
\vnbprop{card-cartesian-nn}{\vnbstmt{\begin{array}{@{}l@{}} \forall x, y. \\ \quad ((x \in \mathbf{Set}) \wedge ((|x| \in \mathbb{N}) \wedge ((y \in \mathbf{Set}) \wedge (|y| \in \mathbb{N})))) \Rightarrow  \\ \quad \quad |(x \times y)| \in \mathbb{N} \end{array}}}
\vnbprop{raag-carr}{\vnbstmt{\forall a.\; (\operatorname{is\text{-}ring}(a) \Rightarrow (\operatorname{carr}(\operatorname{ring\text{-}additive\text{-}ag}(a)) \simeq \operatorname{carr}(a)))}}
\vnbprop{raag-opr}{\vnbstmt{\forall a.\; (\operatorname{is\text{-}ring}(a) \Rightarrow (\operatorname{opr}(\operatorname{ring\text{-}additive\text{-}ag}(a)) \simeq \operatorname{add}(a)))}}
\vnbprop{raag-iden}{\vnbstmt{\forall a.\; (\operatorname{is\text{-}ring}(a) \Rightarrow (\operatorname{iden}(\operatorname{ring\text{-}additive\text{-}ag}(a)) \simeq \operatorname{zero}(a)))}}
\vnbprop{monalg-mul-apply}{\vnbstmt{\begin{array}{@{}l@{}} \forall a, m, f, g.; \forall x \in \operatorname{carr}(m). \\ \quad (\operatorname{monalg\text{-}mul}(a, m, f, g)(x) \simeq \operatorname{finsum}(\operatorname{ring\text{-}additive\text{-}ag}(a), (\lambda p \in \{p \in (\operatorname{supp}(a, m, f) \times \operatorname{supp}(a, m, g)) : (\operatorname{opr}(m)(\operatorname{nth}(1, p), \operatorname{nth}(2, p)) = x)\}.\; \operatorname{mul}(a)(f(\operatorname{nth}(1, p)), g(\operatorname{nth}(2, p)))), \{p \in (\operatorname{supp}(a, m, f) \times \operatorname{supp}(a, m, g)) : (\operatorname{opr}(m)(\operatorname{nth}(1, p), \operatorname{nth}(2, p)) = x)\})) \end{array}}}
\vnbprop{monalg-mul-value-in-carr}{\vnbstmt{\begin{array}{@{}l@{}} \forall a, m, f, g, x. \\ \quad (\operatorname{is\text{-}ring}(a) \wedge (\operatorname{is\text{-}monoid}(m) \wedge ((f \in \operatorname{finsupp}(a, m)) \wedge ((g \in \operatorname{finsupp}(a, m)) \wedge (x \in \operatorname{carr}(m)))))) \Rightarrow  \\ \quad \quad \operatorname{monalg\text{-}mul}(a, m, f, g)(x) \in \operatorname{carr}(a) \end{array}}}
\vnbprop{monalg-mul-fun}{\vnbstmt{\begin{array}{@{}l@{}} \forall a, m, f, g. \\ \quad (\operatorname{is\text{-}ring}(a) \wedge (\operatorname{is\text{-}monoid}(m) \wedge ((f \in \operatorname{finsupp}(a, m)) \wedge (g \in \operatorname{finsupp}(a, m))))) \Rightarrow  \\ \quad \quad \operatorname{monalg\text{-}mul}(a, m, f, g) \in \operatorname{finsupp}(a, m) \end{array}}}
\vnbprop{cartesian-pair-eq}{\vnbstmt{\forall x, a, b.\; ((x \in (a \times b)) \Rightarrow (x = \langle \operatorname{nth}(1, x), \operatorname{nth}(2, x) \rangle))}}
\vnbprop{supp-add-subset}{\vnbstmt{\begin{array}{@{}l@{}} \forall a, m, g, h, z. \\ \quad (\operatorname{is\text{-}ring}(a) \wedge ((g \in \operatorname{finsupp}(a, m)) \wedge ((h \in \operatorname{finsupp}(a, m)) \wedge (z \in \operatorname{supp}(a, m, \operatorname{monalg\text{-}add}(a, m, g, h)))))) \Rightarrow  \\ \quad \quad z \in (\operatorname{supp}(a, m, g) \cup \operatorname{supp}(a, m, h)) \end{array}}}
\vnbprop{monalg-distrib-left-at}{\vnbstmt{\begin{array}{@{}l@{}} \forall a, m, f, g, h, x. \\ \quad \operatorname{is\text{-}ring}(a) \wedge \\ \quad \operatorname{is\text{-}monoid}(m) \wedge \\ \quad f \in \operatorname{finsupp}(a, m) \wedge \\ \quad g \in \operatorname{finsupp}(a, m) \wedge \\ \quad h \in \operatorname{finsupp}(a, m) \wedge \\ \quad x \in \operatorname{carr}(m) \Rightarrow \\ \quad \quad \operatorname{finsum}(\operatorname{ring\text{-}additive\text{-}ag}(a), \\ \quad \quad \quad (\lambda p \in \{p \in (\operatorname{supp}(a, m, f) \times \operatorname{supp}(a, m, \operatorname{monalg\text{-}add}(a, m, g, h))) : (\operatorname{opr}(m)(\operatorname{nth}(1, p), \operatorname{nth}(2, p)) = x)\}.\; \\ \quad \quad \quad \quad \operatorname{mul}(a)(f(\operatorname{nth}(1, p)), \operatorname{monalg\text{-}add}(a, m, g, h)(\operatorname{nth}(2, p)))), \\ \quad \quad \quad \{p \in (\operatorname{supp}(a, m, f) \times \operatorname{supp}(a, m, \operatorname{monalg\text{-}add}(a, m, g, h))) : (\operatorname{opr}(m)(\operatorname{nth}(1, p), \operatorname{nth}(2, p)) = x)\}) \\ \quad \quad =\; \operatorname{add}(a)(\operatorname{finsum}(\operatorname{ring\text{-}additive\text{-}ag}(a), \\ \quad \quad \quad \quad (\lambda p \in \{p \in (\operatorname{supp}(a, m, f) \times \operatorname{supp}(a, m, g)) : (\operatorname{opr}(m)(\operatorname{nth}(1, p), \operatorname{nth}(2, p)) = x)\}.\; \\ \quad \quad \quad \quad \quad \operatorname{mul}(a)(f(\operatorname{nth}(1, p)), g(\operatorname{nth}(2, p)))), \\ \quad \quad \quad \quad \{p \in (\operatorname{supp}(a, m, f) \times \operatorname{supp}(a, m, g)) : (\operatorname{opr}(m)(\operatorname{nth}(1, p), \operatorname{nth}(2, p)) = x)\}), \\ \quad \quad \quad \operatorname{finsum}(\operatorname{ring\text{-}additive\text{-}ag}(a), \\ \quad \quad \quad \quad (\lambda p \in \{p \in (\operatorname{supp}(a, m, f) \times \operatorname{supp}(a, m, h)) : (\operatorname{opr}(m)(\operatorname{nth}(1, p), \operatorname{nth}(2, p)) = x)\}.\; \\ \quad \quad \quad \quad \quad \operatorname{mul}(a)(f(\operatorname{nth}(1, p)), h(\operatorname{nth}(2, p)))), \\ \quad \quad \quad \quad \{p \in (\operatorname{supp}(a, m, f) \times \operatorname{supp}(a, m, h)) : (\operatorname{opr}(m)(\operatorname{nth}(1, p), \operatorname{nth}(2, p)) = x)\})) \end{array}}}
\vnbprop{monalg-distrib-right-at}{\vnbstmt{\begin{array}{@{}l@{}} \forall a, m, f, g, h, x. \\ \quad \operatorname{is\text{-}ring}(a) \wedge \\ \quad \operatorname{is\text{-}monoid}(m) \wedge \\ \quad f \in \operatorname{finsupp}(a, m) \wedge \\ \quad g \in \operatorname{finsupp}(a, m) \wedge \\ \quad h \in \operatorname{finsupp}(a, m) \wedge \\ \quad x \in \operatorname{carr}(m) \Rightarrow \\ \quad \quad \operatorname{finsum}(\operatorname{ring\text{-}additive\text{-}ag}(a), \\ \quad \quad \quad (\lambda p \in \{p \in (\operatorname{supp}(a, m, \operatorname{monalg\text{-}add}(a, m, f, g)) \times \operatorname{supp}(a, m, h)) : (\operatorname{opr}(m)(\operatorname{nth}(1, p), \operatorname{nth}(2, p)) = x)\}.\; \\ \quad \quad \quad \quad \operatorname{mul}(a)(\operatorname{monalg\text{-}add}(a, m, f, g)(\operatorname{nth}(1, p)), h(\operatorname{nth}(2, p)))), \\ \quad \quad \quad \{p \in (\operatorname{supp}(a, m, \operatorname{monalg\text{-}add}(a, m, f, g)) \times \operatorname{supp}(a, m, h)) : (\operatorname{opr}(m)(\operatorname{nth}(1, p), \operatorname{nth}(2, p)) = x)\}) \\ \quad \quad =\; \operatorname{add}(a)(\operatorname{finsum}(\operatorname{ring\text{-}additive\text{-}ag}(a), \\ \quad \quad \quad \quad (\lambda p \in \{p \in (\operatorname{supp}(a, m, f) \times \operatorname{supp}(a, m, h)) : (\operatorname{opr}(m)(\operatorname{nth}(1, p), \operatorname{nth}(2, p)) = x)\}.\; \\ \quad \quad \quad \quad \quad \operatorname{mul}(a)(f(\operatorname{nth}(1, p)), h(\operatorname{nth}(2, p)))), \\ \quad \quad \quad \quad \{p \in (\operatorname{supp}(a, m, f) \times \operatorname{supp}(a, m, h)) : (\operatorname{opr}(m)(\operatorname{nth}(1, p), \operatorname{nth}(2, p)) = x)\}), \\ \quad \quad \quad \operatorname{finsum}(\operatorname{ring\text{-}additive\text{-}ag}(a), \\ \quad \quad \quad \quad (\lambda p \in \{p \in (\operatorname{supp}(a, m, g) \times \operatorname{supp}(a, m, h)) : (\operatorname{opr}(m)(\operatorname{nth}(1, p), \operatorname{nth}(2, p)) = x)\}.\; \\ \quad \quad \quad \quad \quad \operatorname{mul}(a)(g(\operatorname{nth}(1, p)), h(\operatorname{nth}(2, p)))), \\ \quad \quad \quad \quad \{p \in (\operatorname{supp}(a, m, g) \times \operatorname{supp}(a, m, h)) : (\operatorname{opr}(m)(\operatorname{nth}(1, p), \operatorname{nth}(2, p)) = x)\})) \end{array}}}
\vnbprop{monalg-one-supp-subset}{\vnbstmt{\begin{array}{@{}l@{}} \forall a, m, z. \\ \quad (\operatorname{is\text{-}ring}(a) \wedge (\operatorname{is\text{-}monoid}(m) \wedge (z \in \operatorname{supp}(a, m, \operatorname{monalg\text{-}one}(a, m))))) \Rightarrow  \\ \quad \quad z \in \{\operatorname{iden}(m)\} \end{array}}}
\vnbprop{monalg-one-left-at}{\vnbstmt{\begin{array}{@{}l@{}} \forall a, m, f, x. \\ \quad (\operatorname{is\text{-}ring}(a) \wedge (\operatorname{is\text{-}monoid}(m) \wedge ((f \in \operatorname{finsupp}(a, m)) \wedge (x \in \operatorname{carr}(m))))) \Rightarrow  \\ \quad \quad \operatorname{finsum}(\operatorname{ring\text{-}additive\text{-}ag}(a), \\ \quad \quad \quad (\lambda p \in \{p \in (\operatorname{supp}(a, m, \operatorname{monalg\text{-}one}(a, m)) \times \operatorname{supp}(a, m, f)) : (\operatorname{opr}(m)(\operatorname{nth}(1, p), \operatorname{nth}(2, p)) = x)\}.\; \\ \quad \quad \quad \quad \operatorname{mul}(a)(\operatorname{monalg\text{-}one}(a, m)(\operatorname{nth}(1, p)), f(\operatorname{nth}(2, p)))), \\ \quad \quad \quad \{p \in (\operatorname{supp}(a, m, \operatorname{monalg\text{-}one}(a, m)) \times \operatorname{supp}(a, m, f)) : (\operatorname{opr}(m)(\operatorname{nth}(1, p), \operatorname{nth}(2, p)) = x)\}) \\ \quad \quad =\; f(x) \end{array}}}
\vnbprop{monalg-one-right-at}{\vnbstmt{\begin{array}{@{}l@{}} \forall a, m, f, x. \\ \quad (\operatorname{is\text{-}ring}(a) \wedge (\operatorname{is\text{-}monoid}(m) \wedge ((f \in \operatorname{finsupp}(a, m)) \wedge (x \in \operatorname{carr}(m))))) \Rightarrow  \\ \quad \quad \operatorname{finsum}(\operatorname{ring\text{-}additive\text{-}ag}(a), \\ \quad \quad \quad (\lambda p \in \{p \in (\operatorname{supp}(a, m, f) \times \operatorname{supp}(a, m, \operatorname{monalg\text{-}one}(a, m))) : (\operatorname{opr}(m)(\operatorname{nth}(1, p), \operatorname{nth}(2, p)) = x)\}.\; \\ \quad \quad \quad \quad \operatorname{mul}(a)(f(\operatorname{nth}(1, p)), \operatorname{monalg\text{-}one}(a, m)(\operatorname{nth}(2, p)))), \\ \quad \quad \quad \{p \in (\operatorname{supp}(a, m, f) \times \operatorname{supp}(a, m, \operatorname{monalg\text{-}one}(a, m))) : (\operatorname{opr}(m)(\operatorname{nth}(1, p), \operatorname{nth}(2, p)) = x)\}) \\ \quad \quad =\; f(x) \end{array}}}
\vnbprop{monalg-distrib-left}{\vnbstmt{\begin{array}{@{}l@{}} \forall a, m, f, g, h. \\ \quad (\operatorname{is\text{-}ring}(a) \wedge (\operatorname{is\text{-}monoid}(m) \wedge ((f \in \operatorname{finsupp}(a, m)) \wedge ((g \in \operatorname{finsupp}(a, m)) \wedge (h \in \operatorname{finsupp}(a, m)))))) \Rightarrow  \\ \quad \quad \operatorname{monalg\text{-}mul}(a, m, f, \operatorname{monalg\text{-}add}(a, m, g, h)) \\ \quad \quad =\; \operatorname{monalg\text{-}add}(a, m, \operatorname{monalg\text{-}mul}(a, m, f, g), \operatorname{monalg\text{-}mul}(a, m, f, h)) \end{array}}}
\vnbprop{monalg-distrib-right}{\vnbstmt{\begin{array}{@{}l@{}} \forall a, m, f, g, h. \\ \quad (\operatorname{is\text{-}ring}(a) \wedge (\operatorname{is\text{-}monoid}(m) \wedge ((f \in \operatorname{finsupp}(a, m)) \wedge ((g \in \operatorname{finsupp}(a, m)) \wedge (h \in \operatorname{finsupp}(a, m)))))) \Rightarrow  \\ \quad \quad \operatorname{monalg\text{-}mul}(a, m, \operatorname{monalg\text{-}add}(a, m, f, g), h) \\ \quad \quad =\; \operatorname{monalg\text{-}add}(a, m, \operatorname{monalg\text{-}mul}(a, m, f, h), \operatorname{monalg\text{-}mul}(a, m, g, h)) \end{array}}}
\vnbprop{monalg-one-left-fn}{\vnbstmt{\begin{array}{@{}l@{}} \forall a, m, f. \\ \quad (\operatorname{is\text{-}ring}(a) \wedge (\operatorname{is\text{-}monoid}(m) \wedge (f \in \operatorname{finsupp}(a, m)))) \Rightarrow  \\ \quad \quad \operatorname{monalg\text{-}mul}(a, m, \operatorname{monalg\text{-}one}(a, m), f) = f \end{array}}}
\vnbprop{monalg-one-right-fn}{\vnbstmt{\begin{array}{@{}l@{}} \forall a, m, f. \\ \quad (\operatorname{is\text{-}ring}(a) \wedge (\operatorname{is\text{-}monoid}(m) \wedge (f \in \operatorname{finsupp}(a, m)))) \Rightarrow  \\ \quad \quad \operatorname{monalg\text{-}mul}(a, m, f, \operatorname{monalg\text{-}one}(a, m)) = f \end{array}}}
\vnbprop{monalg-one-left}{\vnbstmt{\begin{array}{@{}l@{}} \forall a, m, f, x. \\ \quad (\operatorname{is\text{-}ring}(a) \wedge (\operatorname{is\text{-}monoid}(m) \wedge ((f \in \operatorname{finsupp}(a, m)) \wedge (x \in \operatorname{carr}(m))))) \Rightarrow  \\ \quad \quad \operatorname{monalg\text{-}mul}(a, m, \operatorname{monalg\text{-}one}(a, m), f)(x) = f(x) \end{array}}}
\vnbprop{bijection-from-inverse}{\vnbstmt{\begin{array}{@{}l@{}} \forall x, y, i, a. \\ \quad i \in (x \to y) \wedge \\ \quad a \in (y \to x) \wedge \\ \quad \forall u \in x.\; (a(i(u)) = u) \wedge \\ \quad \forall v \in y.\; (i(a(v)) = v) \Rightarrow \\ \quad \quad i \in \mathrm{Bij}(x, y) \end{array}}}
\vnbprop{lambda-compose-value}{\vnbstmt{\forall f, h, m.; \forall t \in m.\; ((\lambda z \in m.\; f(h(z)))(t) \simeq f(h(t)))}}
\vnbprop{finsum-reindex-inverse}{\vnbstmt{\begin{array}{@{}l@{}} \forall g.; \forall s, t \in \mathbf{Set}, i \in (t \to s), a \in (s \to t), f \in (s \to \operatorname{carr}(g)). \\ \quad \operatorname{is\text{-}abelian\text{-}group}(g) \wedge \\ \quad |s| \in \mathbb{N} \wedge \\ \quad |t| \in \mathbb{N} \wedge \\ \quad \forall u \in t.\; (a(i(u)) = u) \wedge \\ \quad \forall v \in s.\; (i(a(v)) = v) \Rightarrow \\ \quad \quad \operatorname{finsum}(g, f, s) = \operatorname{finsum}(g, (\lambda z \in t.\; f(i(z))), t) \end{array}}}
\vnbprop{finsum-reindex-inverse-ptwise}{\vnbstmt{\begin{array}{@{}l@{}} \forall g.; \forall s, t \in \mathbf{Set}, i \in (t \to s), a \in (s \to t), f. \\ \quad \operatorname{is\text{-}abelian\text{-}group}(g) \wedge \\ \quad |s| \in \mathbb{N} \wedge \\ \quad |t| \in \mathbb{N} \wedge \\ \quad \forall u \in t.\; (a(i(u)) = u) \wedge \\ \quad \forall v \in s.\; (i(a(v)) = v) \wedge \\ \quad \forall b \in s.\; (f(b) \in \operatorname{carr}(g)) \Rightarrow \\ \quad \quad \operatorname{finsum}(g, f, s) = \operatorname{finsum}(g, (\lambda z \in t.\; f(i(z))), t) \end{array}}}
\vnbprop{finsum-fiber-value}{\vnbstmt{\begin{array}{@{}l@{}} \forall g, f, h, s, t.; \forall a \in t. \\ \quad ((\lambda b \in t.\; \operatorname{finsum}(g, f, \{c \in s : (h(c) = b)\}))(a) \simeq \operatorname{finsum}(g, f, \{c \in s : (h(c) = a)\})) \end{array}}}
\vnbprop{finsum-fiber-slice}{\vnbstmt{\begin{array}{@{}l@{}} \forall g.; \forall s, k \in \mathbf{Set}, h, p, f, l, e. \\ \quad \operatorname{is\text{-}abelian\text{-}group}(g) \wedge \\ \quad |s| \in \mathbb{N} \wedge \\ \quad |k| \in \mathbb{N} \wedge \\ \quad l \in (k \to s) \wedge \\ \quad \forall a \in k.\; (h(l(a)) = p) \wedge \\ \quad \forall w \in s.\; ((h(w) = p) \Rightarrow (e(w) \in k)) \wedge \\ \quad \forall a \in k.\; (e(l(a)) = a) \wedge \\ \quad \forall w \in s.\; ((h(w) = p) \Rightarrow (l(e(w)) = w)) \wedge \\ \quad \forall b \in s.\; (f(b) \in \operatorname{carr}(g)) \Rightarrow \\ \quad \quad \operatorname{finsum}(g, f, \{c \in s : (h(c) = p)\}) = \operatorname{finsum}(g, (\lambda z \in k.\; f(l(z))), k) \end{array}}}
\vnbprop{monalg-mul-supp-in-image}{\vnbstmt{\begin{array}{@{}l@{}} \forall a, m, f, g, z. \\ \quad \operatorname{is\text{-}ring}(a) \wedge \\ \quad \operatorname{is\text{-}monoid}(m) \wedge \\ \quad f \in \operatorname{finsupp}(a, m) \wedge \\ \quad g \in \operatorname{finsupp}(a, m) \wedge \\ \quad z \in \operatorname{supp}(a, m, \operatorname{monalg\text{-}mul}(a, m, f, g)) \Rightarrow \\ \quad \quad z \\ \quad \quad \in\; \operatorname{image}((\lambda q \in (\operatorname{supp}(a, m, f) \times \operatorname{supp}(a, m, g)).\; \operatorname{opr}(m)(\operatorname{nth}(1, q), \operatorname{nth}(2, q))), \\ \quad \quad \quad (\operatorname{supp}(a, m, f) \times \operatorname{supp}(a, m, g))) \end{array}}}
\vnbprop{monalg-mul-assoc-left-core}{\vnbstmt{\begin{array}{@{}l@{}} \forall a, m, f, g, h, x. \\ \quad \operatorname{is\text{-}ring}(a) \wedge \\ \quad \operatorname{is\text{-}monoid}(m) \wedge \\ \quad f \in \operatorname{finsupp}(a, m) \wedge \\ \quad g \in \operatorname{finsupp}(a, m) \wedge \\ \quad h \in \operatorname{finsupp}(a, m) \wedge \\ \quad x \in \operatorname{carr}(m) \Rightarrow \\ \quad \quad \operatorname{finsum}(\operatorname{ring\text{-}additive\text{-}ag}(a), \\ \quad \quad \quad (\lambda p \in \{p \in (\operatorname{image}((\lambda q \in (\operatorname{supp}(a, m, f) \times \operatorname{supp}(a, m, g)).\; \operatorname{opr}(m)(\operatorname{nth}(1, q), \operatorname{nth}(2, q))), (\operatorname{supp}(a, m, f) \times \operatorname{supp}(a, m, g))) \times \operatorname{supp}(a, m, h)) : (\operatorname{opr}(m)(\operatorname{nth}(1, p), \operatorname{nth}(2, p)) = x)\}.\; \\ \quad \quad \quad \quad \operatorname{mul}(a)(\operatorname{monalg\text{-}mul}(a, m, f, g)(\operatorname{nth}(1, p)), h(\operatorname{nth}(2, p)))), \\ \quad \quad \quad \{p \in (\operatorname{image}((\lambda q \in (\operatorname{supp}(a, m, f) \times \operatorname{supp}(a, m, g)).\; \operatorname{opr}(m)(\operatorname{nth}(1, q), \operatorname{nth}(2, q))), (\operatorname{supp}(a, m, f) \times \operatorname{supp}(a, m, g))) \times \operatorname{supp}(a, m, h)) : (\operatorname{opr}(m)(\operatorname{nth}(1, p), \operatorname{nth}(2, p)) = x)\}) \\ \quad \quad =\; \operatorname{finsum}(\operatorname{ring\text{-}additive\text{-}ag}(a), \\ \quad \quad \quad (\lambda w \in \{w \in ((\operatorname{supp}(a, m, f) \times \operatorname{supp}(a, m, g)) \times \operatorname{supp}(a, m, h)) : (\operatorname{opr}(m)(\operatorname{opr}(m)(\operatorname{nth}(1, \operatorname{nth}(1, w)), \operatorname{nth}(2, \operatorname{nth}(1, w))), \operatorname{nth}(2, w)) = x)\}.\; \\ \quad \quad \quad \quad \operatorname{mul}(a)(\operatorname{mul}(a)(f(\operatorname{nth}(1, \operatorname{nth}(1, w))), g(\operatorname{nth}(2, \operatorname{nth}(1, w)))), h(\operatorname{nth}(2, w)))), \\ \quad \quad \quad \{w \in ((\operatorname{supp}(a, m, f) \times \operatorname{supp}(a, m, g)) \times \operatorname{supp}(a, m, h)) : (\operatorname{opr}(m)(\operatorname{opr}(m)(\operatorname{nth}(1, \operatorname{nth}(1, w)), \operatorname{nth}(2, \operatorname{nth}(1, w))), \operatorname{nth}(2, w)) = x)\}) \end{array}}}
\vnbprop{monalg-mul-assoc-right-core}{\vnbstmt{\begin{array}{@{}l@{}} \forall a, m, f, g, h, x. \\ \quad \operatorname{is\text{-}ring}(a) \wedge \\ \quad \operatorname{is\text{-}monoid}(m) \wedge \\ \quad f \in \operatorname{finsupp}(a, m) \wedge \\ \quad g \in \operatorname{finsupp}(a, m) \wedge \\ \quad h \in \operatorname{finsupp}(a, m) \wedge \\ \quad x \in \operatorname{carr}(m) \Rightarrow \\ \quad \quad \operatorname{finsum}(\operatorname{ring\text{-}additive\text{-}ag}(a), \\ \quad \quad \quad (\lambda p \in \{p \in (\operatorname{supp}(a, m, f) \times \operatorname{image}((\lambda q \in (\operatorname{supp}(a, m, g) \times \operatorname{supp}(a, m, h)).\; \operatorname{opr}(m)(\operatorname{nth}(1, q), \operatorname{nth}(2, q))), (\operatorname{supp}(a, m, g) \times \operatorname{supp}(a, m, h)))) : (\operatorname{opr}(m)(\operatorname{nth}(1, p), \operatorname{nth}(2, p)) = x)\}.\; \\ \quad \quad \quad \quad \operatorname{mul}(a)(f(\operatorname{nth}(1, p)), \operatorname{monalg\text{-}mul}(a, m, g, h)(\operatorname{nth}(2, p)))), \\ \quad \quad \quad \{p \in (\operatorname{supp}(a, m, f) \times \operatorname{image}((\lambda q \in (\operatorname{supp}(a, m, g) \times \operatorname{supp}(a, m, h)).\; \operatorname{opr}(m)(\operatorname{nth}(1, q), \operatorname{nth}(2, q))), (\operatorname{supp}(a, m, g) \times \operatorname{supp}(a, m, h)))) : (\operatorname{opr}(m)(\operatorname{nth}(1, p), \operatorname{nth}(2, p)) = x)\}) \\ \quad \quad =\; \operatorname{finsum}(\operatorname{ring\text{-}additive\text{-}ag}(a), \\ \quad \quad \quad (\lambda w \in \{w \in ((\operatorname{supp}(a, m, f) \times \operatorname{supp}(a, m, g)) \times \operatorname{supp}(a, m, h)) : (\operatorname{opr}(m)(\operatorname{opr}(m)(\operatorname{nth}(1, \operatorname{nth}(1, w)), \operatorname{nth}(2, \operatorname{nth}(1, w))), \operatorname{nth}(2, w)) = x)\}.\; \\ \quad \quad \quad \quad \operatorname{mul}(a)(f(\operatorname{nth}(1, \operatorname{nth}(1, w))), \operatorname{mul}(a)(g(\operatorname{nth}(2, \operatorname{nth}(1, w))), h(\operatorname{nth}(2, w))))), \\ \quad \quad \quad \{w \in ((\operatorname{supp}(a, m, f) \times \operatorname{supp}(a, m, g)) \times \operatorname{supp}(a, m, h)) : (\operatorname{opr}(m)(\operatorname{opr}(m)(\operatorname{nth}(1, \operatorname{nth}(1, w)), \operatorname{nth}(2, \operatorname{nth}(1, w))), \operatorname{nth}(2, w)) = x)\}) \end{array}}}
\vnbprop{monalg-mul-assoc-at}{\vnbstmt{\begin{array}{@{}l@{}} \forall a, m, f, g, h, x. \\ \quad \operatorname{is\text{-}ring}(a) \wedge \\ \quad \operatorname{is\text{-}monoid}(m) \wedge \\ \quad f \in \operatorname{finsupp}(a, m) \wedge \\ \quad g \in \operatorname{finsupp}(a, m) \wedge \\ \quad h \in \operatorname{finsupp}(a, m) \wedge \\ \quad x \in \operatorname{carr}(m) \Rightarrow \\ \quad \quad \operatorname{monalg\text{-}mul}(a, m, \operatorname{monalg\text{-}mul}(a, m, f, g), h)(x) = \operatorname{monalg\text{-}mul}(a, m, f, \operatorname{monalg\text{-}mul}(a, m, g, h))(x) \end{array}}}
\vnbprop{monalg-mul-assoc}{\vnbstmt{\begin{array}{@{}l@{}} \forall a, m, f, g, h. \\ \quad (\operatorname{is\text{-}ring}(a) \wedge (\operatorname{is\text{-}monoid}(m) \wedge ((f \in \operatorname{finsupp}(a, m)) \wedge ((g \in \operatorname{finsupp}(a, m)) \wedge (h \in \operatorname{finsupp}(a, m)))))) \Rightarrow  \\ \quad \quad \operatorname{monalg\text{-}mul}(a, m, \operatorname{monalg\text{-}mul}(a, m, f, g), h) = \operatorname{monalg\text{-}mul}(a, m, f, \operatorname{monalg\text{-}mul}(a, m, g, h)) \end{array}}}
\vnbprop{monalg-is-ring}{\vnbstmt{\forall a, m.\; ((\operatorname{is\text{-}ring}(a) \wedge \operatorname{is\text{-}monoid}(m)) \Rightarrow \operatorname{is\text{-}ring}(\operatorname{monalg}(a, m)))}}

\vnbbookfile{theorem-library/rake-prod-of-sums}{theorem-library/rake-prod-of-sums}
\noindent{\small rake-prod-of-sums.scm -- the product-of-sums expansion, proven.}\par
\vnbprop{difference-empty-left}{\vnbstmt{\forall s.\; (\operatorname{difference}(\emptyset, s) = \emptyset)}}
\vnbprop{insert-remove}{\vnbstmt{\begin{array}{@{}l@{}} \forall s.; \forall k \in \mathbf{Set}. \\ \quad \neg (k \in s) \Rightarrow (\operatorname{difference}((s \cup \{k\}), \{k\}) = s) \end{array}}}
\vnbprop{prod-ring-type-ptwise}{\vnbstmt{\begin{array}{@{}l@{}} \forall r.; \forall x \in \mathbf{Set}, f. \\ \quad (\operatorname{is\text{-}commutative\text{-}ring}(r) \wedge ((|x| \in \mathbb{N}) \wedge \forall z \in x.\; (f(z) \in \operatorname{carr}(r)))) \Rightarrow  \\ \quad \quad \operatorname{prod\text{-}ring}(r, f, x) \in \operatorname{carr}(r) \end{array}}}
\vnbprop{prod-ring-insert-ptwise}{\vnbstmt{\begin{array}{@{}l@{}} \forall r.; \forall x, k \in \mathbf{Set}, f. \\ \quad \operatorname{is\text{-}commutative\text{-}ring}(r) \wedge \\ \quad |x| \in \mathbb{N} \wedge \\ \quad \neg (k \in x) \wedge \\ \quad \forall z \in (x \cup \{k\}).\; (f(z) \in \operatorname{carr}(r)) \Rightarrow \\ \quad \quad \operatorname{prod\text{-}ring}(r, f, (x \cup \{k\})) = \operatorname{mul}(r)(\operatorname{prod\text{-}ring}(r, f, x), f(k)) \end{array}}}
\vnbprop{prod-ring-congruence-q}{\vnbstmt{\begin{array}{@{}l@{}} \forall r.; \forall x \in \mathbf{Set}, f, g. \\ \quad (\operatorname{is\text{-}commutative\text{-}ring}(r) \wedge ((|x| \in \mathbb{N}) \wedge \forall z \in x.\; (f(z) = g(z)))) \Rightarrow  \\ \quad \quad (\operatorname{prod\text{-}ring}(r, f, x) \simeq \operatorname{prod\text{-}ring}(r, g, x)) \end{array}}}
\vnbprop{cra-carr}{\vnbstmt{\begin{array}{@{}l@{}} \forall r. \\ \quad \operatorname{is\text{-}commutative\text{-}ring}(r) \Rightarrow (\operatorname{carr}(\operatorname{commutative\text{-}ring\text{-}additive\text{-}ag}(r)) = \operatorname{carr}(r)) \end{array}}}
\vnbprop{cra-op}{\vnbstmt{\begin{array}{@{}l@{}} \forall r. \\ \quad \operatorname{is\text{-}commutative\text{-}ring}(r) \Rightarrow (\operatorname{opr}(\operatorname{commutative\text{-}ring\text{-}additive\text{-}ag}(r)) = \operatorname{add}(r)) \end{array}}}
\vnbprop{difference-in-power}{\vnbstmt{\forall w \in \mathbf{Set}, s.\; (\operatorname{difference}(w, s) \in \operatorname{power}(w))}}
\vnbprop{prod-ring-type-subset}{\vnbstmt{\begin{array}{@{}l@{}} \forall r.; \forall w \in \mathbf{Set}, f.; \forall s \in \operatorname{power}(w). \\ \quad (\operatorname{is\text{-}commutative\text{-}ring}(r) \wedge ((|w| \in \mathbb{N}) \wedge \forall z \in w.\; (f(z) \in \operatorname{carr}(r)))) \Rightarrow  \\ \quad \quad \operatorname{prod\text{-}ring}(r, f, s) \in \operatorname{carr}(r) \end{array}}}
\vnbprop{prod-of-sums-summand-type}{\vnbstmt{\begin{array}{@{}l@{}} \forall r.; \forall w \in \mathbf{Set}, a, b.; \forall s \in \operatorname{power}(w). \\ \quad \operatorname{is\text{-}commutative\text{-}ring}(r) \wedge \\ \quad |w| \in \mathbb{N} \wedge \\ \quad \forall z \in w.\; (a(z) \in \operatorname{carr}(r)) \wedge \\ \quad \forall z \in w.\; (b(z) \in \operatorname{carr}(r)) \Rightarrow \\ \quad \quad \operatorname{mul}(r)(\operatorname{prod\text{-}ring}(r, a, s), \operatorname{prod\text{-}ring}(r, b, \operatorname{difference}(w, s))) \in \operatorname{carr}(r) \end{array}}}
\vnbprop{power-insert-injective}{\vnbstmt{\begin{array}{@{}l@{}} \forall x, k \in \mathbf{Set}. \\ \quad \neg (k \in x) \Rightarrow  \\ \quad \quad (\lambda s \in \operatorname{power}(x).\; (s \cup \{k\})) \in \operatorname{injection}(\operatorname{power}(x), \operatorname{power}((x \cup \{k\}))) \end{array}}}
\vnbprop{lambda-eta-value}{\vnbstmt{\forall f, m.; \forall t \in m.\; ((\lambda z \in m.\; f(z))(t) \simeq f(t))}}
\vnbprop{prod-of-sums-ptwise}{\vnbstmt{\begin{array}{@{}l@{}} \forall r.; \forall w \in \mathbf{Set}, a, b. \\ \quad \operatorname{is\text{-}commutative\text{-}ring}(r) \wedge \\ \quad |w| \in \mathbb{N} \wedge \\ \quad \forall z \in w.\; (a(z) \in \operatorname{carr}(r)) \wedge \\ \quad \forall z \in w.\; (b(z) \in \operatorname{carr}(r)) \Rightarrow \\ \quad \quad \operatorname{prod\text{-}ring}(r, (\lambda v \in w.\; \operatorname{add}(r)(a(v), b(v))), w) \\ \quad \quad =\; \operatorname{finsum}(\operatorname{commutative\text{-}ring\text{-}additive\text{-}ag}(r), \\ \quad \quad \quad (\lambda \mathit{sv} \in \operatorname{power}(w).\; \operatorname{mul}(r)(\operatorname{prod\text{-}ring}(r, a, \mathit{sv}), \operatorname{prod\text{-}ring}(r, b, \operatorname{difference}(w, \mathit{sv})))), \\ \quad \quad \quad \operatorname{power}(w)) \end{array}}}
\vnbprop{prod-of-sums-expansion}{\vnbstmt{\begin{array}{@{}l@{}} \forall r.; \forall x \in \mathbf{Set}, a, b \in (x \to \operatorname{carr}(r)). \\ \quad (\operatorname{is\text{-}commutative\text{-}ring}(r) \wedge (|x| \in \mathbb{N})) \Rightarrow  \\ \quad \quad \operatorname{prod\text{-}ring}(r, (\lambda k \in x.\; \operatorname{add}(r)(a(k), b(k))), x) \\ \quad \quad =\; \operatorname{finsum}(\operatorname{commutative\text{-}ring\text{-}additive\text{-}ag}(r), \\ \quad \quad \quad (\lambda s \in \operatorname{power}(x).\; \operatorname{mul}(r)(\operatorname{prod\text{-}ring}(r, a, s), \operatorname{prod\text{-}ring}(r, b, \operatorname{difference}(x, s)))), \\ \quad \quad \quad \operatorname{power}(x)) \end{array}}}

\vnbbookfile{theorem-library/rake-monalg-comm}{theorem-library/rake-monalg-comm}
\noindent{\small theorem-library/rake-monalg-comm.scm -- rake batch 5b, assignment 5b-A.}\par
\vnbprop{monalg-index-swap-fun}{\vnbstmt{\begin{array}{@{}l@{}} \forall a, m, f, g, x. \\ \quad \operatorname{is\text{-}comm\text{-}monoid}(m) \Rightarrow  \\ \quad \quad (\lambda q \in \{p \in (\operatorname{supp}(a, m, f) \times \operatorname{supp}(a, m, g)) : (\operatorname{opr}(m)(\operatorname{nth}(1, p), \operatorname{nth}(2, p)) = x)\}.\; \\ \quad \quad \quad \langle \operatorname{nth}(2, q), \operatorname{nth}(1, q) \rangle) \\ \quad \quad \in\; (\{p \in (\operatorname{supp}(a, m, f) \times \operatorname{supp}(a, m, g)) : (\operatorname{opr}(m)(\operatorname{nth}(1, p), \operatorname{nth}(2, p)) = x)\} \to \{p \in (\operatorname{supp}(a, m, g) \times \operatorname{supp}(a, m, f)) : (\operatorname{opr}(m)(\operatorname{nth}(1, p), \operatorname{nth}(2, p)) = x)\}) \end{array}}}
\vnbprop{monalg-index-swap-inverse}{\vnbstmt{\begin{array}{@{}l@{}} \forall a, m, f, g, x.; \forall u \in \{p \in (\operatorname{supp}(a, m, f) \times \operatorname{supp}(a, m, g)) : (\operatorname{opr}(m)(\operatorname{nth}(1, p), \operatorname{nth}(2, p)) = x)\}. \\ \quad \operatorname{is\text{-}comm\text{-}monoid}(m) \Rightarrow  \\ \quad \quad (\lambda q \in \{p \in (\operatorname{supp}(a, m, g) \times \operatorname{supp}(a, m, f)) : (\operatorname{opr}(m)(\operatorname{nth}(1, p), \operatorname{nth}(2, p)) = x)\}.\; \langle \operatorname{nth}(2, q), \operatorname{nth}(1, q) \rangle)((\lambda q \in \{p \in (\operatorname{supp}(a, m, f) \times \operatorname{supp}(a, m, g)) : (\operatorname{opr}(m)(\operatorname{nth}(1, p), \operatorname{nth}(2, p)) = x)\}.\; \langle \operatorname{nth}(2, q), \operatorname{nth}(1, q) \rangle)(u)) \\ \quad \quad =\; u \end{array}}}
\vnbprop{monalg-index-swap-bijection}{\vnbstmt{\begin{array}{@{}l@{}} \forall a, m, f, g, x. \\ \quad \operatorname{is\text{-}comm\text{-}monoid}(m) \Rightarrow  \\ \quad \quad (\lambda q \in \{p \in (\operatorname{supp}(a, m, f) \times \operatorname{supp}(a, m, g)) : (\operatorname{opr}(m)(\operatorname{nth}(1, p), \operatorname{nth}(2, p)) = x)\}.\; \\ \quad \quad \quad \langle \operatorname{nth}(2, q), \operatorname{nth}(1, q) \rangle) \\ \quad \quad \in\; \mathrm{Bij}(\{p \in (\operatorname{supp}(a, m, f) \times \operatorname{supp}(a, m, g)) : (\operatorname{opr}(m)(\operatorname{nth}(1, p), \operatorname{nth}(2, p)) = x)\}, \{p \in (\operatorname{supp}(a, m, g) \times \operatorname{supp}(a, m, f)) : (\operatorname{opr}(m)(\operatorname{nth}(1, p), \operatorname{nth}(2, p)) = x)\}) \end{array}}}
\vnbprop{monalg-mul-comm-at}{\vnbstmt{\begin{array}{@{}l@{}} \forall a, m, f, g, x. \\ \quad \operatorname{is\text{-}commutative\text{-}ring}(a) \wedge \\ \quad \operatorname{is\text{-}comm\text{-}monoid}(m) \wedge \\ \quad f \in \operatorname{finsupp}(a, m) \wedge \\ \quad g \in \operatorname{finsupp}(a, m) \wedge \\ \quad x \in \operatorname{carr}(m) \Rightarrow \\ \quad \quad \operatorname{monalg\text{-}mul}(a, m, f, g)(x) = \operatorname{monalg\text{-}mul}(a, m, g, f)(x) \end{array}}}
\vnbprop{monalg-mul-comm}{\vnbstmt{\begin{array}{@{}l@{}} \forall a, m, f, g. \\ \quad (\operatorname{is\text{-}commutative\text{-}ring}(a) \wedge (\operatorname{is\text{-}comm\text{-}monoid}(m) \wedge ((f \in \operatorname{finsupp}(a, m)) \wedge (g \in \operatorname{finsupp}(a, m))))) \Rightarrow  \\ \quad \quad \operatorname{monalg\text{-}mul}(a, m, f, g) = \operatorname{monalg\text{-}mul}(a, m, g, f) \end{array}}}
\vnbprop{monalg-comm}{\vnbstmt{\begin{array}{@{}l@{}} \forall a, m. \\ \quad (\operatorname{is\text{-}commutative\text{-}ring}(a) \wedge \operatorname{is\text{-}comm\text{-}monoid}(m)) \Rightarrow  \\ \quad \quad \operatorname{is\text{-}commutative\text{-}ring}(\operatorname{monalg}(a, m)) \end{array}}}

\vnbbookfile{structure-library/poly-degree}{structure-library/poly-degree}
\noindent{\small poly-degree.scm -- DEG, LEADCOEF and MONOMIAL for the monoid algebra's polynomial specialisation A[NN].}\par
\par\medskip\noindent\textbf{Definition} (\texttt{deg}).\;$\operatorname{deg}(a, p) :=$
\begin{equation*}
(\iota\, \operatorname{n\_}.\;((\operatorname{n\_} \in \mathbb{N}) \wedge (\operatorname{deg\text{-}bound}(a, p, \operatorname{n\_}) \wedge \forall \operatorname{j\_} \in \mathbb{N}.\; (\operatorname{deg\text{-}bound}(a, p, \operatorname{j\_}) \Rightarrow (\operatorname{n\_} \leq \operatorname{j\_})))))
\end{equation*}

\par\medskip\noindent\textbf{Definition} (\texttt{leadcoef}).\;$\operatorname{leadcoef}(a, p) :=$ $p(\operatorname{deg}(a, p))$
\par\medskip\noindent\textbf{Definition} (\texttt{monomial}).\;$\operatorname{monomial}(a, c, n) :=$ $(\lambda \operatorname{x\_} \in \mathbb{N}.\; \left(\text{if }(\operatorname{x\_} = n)\text{ then }c\text{ else }\operatorname{zero}(a)\right))$
\par\medskip\noindent\textbf{Definition} (\texttt{deg-bound}).\;For all $a$, $p$ and $n$, 
\begin{equation*}
\operatorname{deg\text{-}bound}(a, p, n) \Leftrightarrow \forall \operatorname{k\_} \in \mathbb{N}.\; ((n < \operatorname{k\_}) \Rightarrow (p(\operatorname{k\_}) = \operatorname{zero}(a)))
\end{equation*}


\vnbbookfile{theorem-library/poly-degree-laws}{theorem-library/poly-degree-laws}
\noindent{\small poly-degree-laws.scm -- DEG / LEADCOEF / MONOMIAL, PROVEN.}\par
\vnbprop{deg-well-defined}{\vnbstmt{\begin{array}{@{}l@{}} \forall a, p. \\ \quad (\operatorname{is\text{-}ring}(a) \wedge (p \in \operatorname{carr}(\operatorname{poly}(a)))) \Rightarrow  \\ \quad \quad \operatorname{deg}(a, p) \in \mathbb{N} \wedge \\ \quad \quad \operatorname{deg\text{-}bound}(a, p, \operatorname{deg}(a, p)) \wedge \\ \quad \quad \forall j \in \mathbb{N}.\; (\operatorname{deg\text{-}bound}(a, p, j) \Rightarrow (\operatorname{deg}(a, p) \leq j)) \end{array}}}
\vnbprop{deg-bound-apply}{\vnbstmt{\begin{array}{@{}l@{}} \forall a, p, n, k. \\ \quad (\operatorname{deg\text{-}bound}(a, p, n) \wedge ((k \in \mathbb{N}) \wedge (n < k))) \Rightarrow (p(k) = \operatorname{zero}(a)) \end{array}}}
\vnbprop{deg-unfold}{\vnbstmt{\begin{array}{@{}l@{}} \forall a, p. \\ \quad (\operatorname{deg}(a, p) \simeq (\iota\, n.\;((n \in \mathbb{N}) \wedge (\operatorname{deg\text{-}bound}(a, p, n) \wedge \forall j \in \mathbb{N}.\; (\operatorname{deg\text{-}bound}(a, p, j) \Rightarrow (n \leq j)))))) \end{array}}}
\vnbprop{leadcoef-unfold}{\vnbstmt{\forall a, p.\; (\operatorname{leadcoef}(a, p) \simeq p(\operatorname{deg}(a, p)))}}
\vnbprop{leadcoef-nonzero-ptwise}{\vnbstmt{\begin{array}{@{}l@{}} \forall a, p. \\ \quad (\operatorname{is\text{-}ring}(a) \wedge ((p \in \operatorname{carr}(\operatorname{poly}(a))) \wedge \exists k \in \mathbb{N}.\; \neg (p(k) = \operatorname{zero}(a)))) \Rightarrow  \\ \quad \quad \neg (\operatorname{leadcoef}(a, p) = \operatorname{zero}(a)) \end{array}}}
\vnbprop{leadcoef-nonzero}{\vnbstmt{\begin{array}{@{}l@{}} \forall a, p. \\ \quad (\operatorname{is\text{-}ring}(a) \wedge ((p \in \operatorname{carr}(\operatorname{poly}(a))) \wedge \neg (p = \operatorname{zero}(\operatorname{poly}(a))))) \Rightarrow  \\ \quad \quad \neg (\operatorname{leadcoef}(a, p) = \operatorname{zero}(a)) \end{array}}}
\vnbprop{monomial-at}{\vnbstmt{\forall a, c, n.\; (((c \in \operatorname{carr}(a)) \wedge (n \in \mathbb{N})) \Rightarrow (\operatorname{monomial}(a, c, n)(n) = c))}}
\vnbprop{monomial-off}{\vnbstmt{\begin{array}{@{}l@{}} \forall a, c, n, k. \\ \quad (\operatorname{is\text{-}ring}(a) \wedge ((c \in \operatorname{carr}(a)) \wedge ((k \in \mathbb{N}) \wedge \neg (k = n)))) \Rightarrow  \\ \quad \quad \operatorname{monomial}(a, c, n)(k) = \operatorname{zero}(a) \end{array}}}
\vnbprop{monomial-in-poly}{\vnbstmt{\begin{array}{@{}l@{}} \forall a, c, n. \\ \quad (\operatorname{is\text{-}ring}(a) \wedge ((c \in \operatorname{carr}(a)) \wedge (n \in \mathbb{N}))) \Rightarrow  \\ \quad \quad \operatorname{monomial}(a, c, n) \in \operatorname{carr}(\operatorname{poly}(a)) \end{array}}}
\vnbprop{monomial-support}{\vnbstmt{\begin{array}{@{}l@{}} \forall a, c, n, k. \\ \quad (\operatorname{is\text{-}ring}(a) \wedge ((c \in \operatorname{carr}(a)) \wedge (\neg (c = \operatorname{zero}(a)) \wedge (n \in \mathbb{N})))) \Rightarrow  \\ \quad \quad (k \in \operatorname{supp}(a, \operatorname{nn\text{-}add\text{-}monoid}, \operatorname{monomial}(a, c, n))) \Leftrightarrow ((k \in \mathbb{N}) \wedge (k = n)) \end{array}}}
\vnbprop{deg-is}{\vnbstmt{\begin{array}{@{}l@{}} \forall a, p, m. \\ \quad \operatorname{is\text{-}ring}(a) \wedge \\ \quad p \in \operatorname{carr}(\operatorname{poly}(a)) \wedge \\ \quad m \in \mathbb{N} \wedge \\ \quad \operatorname{deg\text{-}bound}(a, p, m) \wedge \\ \quad \forall j \in \mathbb{N}.\; (\operatorname{deg\text{-}bound}(a, p, j) \Rightarrow (m \leq j)) \Rightarrow \\ \quad \quad \operatorname{deg}(a, p) = m \end{array}}}
\vnbprop{monomial-deg}{\vnbstmt{\begin{array}{@{}l@{}} \forall a, c, n. \\ \quad (\operatorname{is\text{-}ring}(a) \wedge ((c \in \operatorname{carr}(a)) \wedge (\neg (c = \operatorname{zero}(a)) \wedge (n \in \mathbb{N})))) \Rightarrow  \\ \quad \quad \operatorname{deg}(a, \operatorname{monomial}(a, c, n)) = n \end{array}}}
\vnbprop{monomial-leadcoef}{\vnbstmt{\begin{array}{@{}l@{}} \forall a, c, n. \\ \quad (\operatorname{is\text{-}ring}(a) \wedge ((c \in \operatorname{carr}(a)) \wedge (\neg (c = \operatorname{zero}(a)) \wedge (n \in \mathbb{N})))) \Rightarrow  \\ \quad \quad \operatorname{leadcoef}(a, \operatorname{monomial}(a, c, n)) = c \end{array}}}
\vnbprop{deg-add-le}{\vnbstmt{\begin{array}{@{}l@{}} \forall a, f, g, n. \\ \quad \operatorname{is\text{-}ring}(a) \wedge \\ \quad f \in \operatorname{carr}(\operatorname{poly}(a)) \wedge \\ \quad g \in \operatorname{carr}(\operatorname{poly}(a)) \wedge \\ \quad n \in \mathbb{N} \wedge \\ \quad \operatorname{deg}(a, f) \leq n \wedge \\ \quad \operatorname{deg}(a, g) \leq n \Rightarrow \\ \quad \quad \operatorname{deg}(a, \operatorname{monalg\text{-}add}(a, \operatorname{nn\text{-}add\text{-}monoid}, f, g)) \leq n \end{array}}}

\vnbbookfile{theorem-library/gauge-is-degree}{theorem-library/gauge-is-degree}
\noindent{\small gauge-is-degree.scm -- the named Euclidean gauge GAUGE(s) is a degree function, PROVEN, together with the four supports its warrant cites.}\par
\vnbprop{euclidean-gauges-unfold}{\vnbstmt{\begin{array}{@{}l@{}} \forall s. \\ \quad (\operatorname{euclidean\text{-}gauges}(s) \simeq \{g \in (\operatorname{carr}(s) \to \mathbb{N}) : \operatorname{has\text{-}div\text{-}remainder}(s, g)\}) \end{array}}}
\vnbprop{euclidean-ring-has-gauge}{\vnbstmt{\begin{array}{@{}l@{}} \forall s. \\ \quad \operatorname{is\text{-}euclidean\text{-}ring}(s) \Rightarrow  \\ \quad \quad \exists g \in (\operatorname{carr}(s) \to \mathbb{N}).\; \operatorname{has\text{-}div\text{-}remainder}(s, g) \end{array}}}
\vnbprop{gauges-mem-build}{\vnbstmt{\begin{array}{@{}l@{}} \forall s.; \forall g \in (\operatorname{carr}(s) \to \mathbb{N}). \\ \quad \operatorname{has\text{-}div\text{-}remainder}(s, g) \Rightarrow (g \in \operatorname{euclidean\text{-}gauges}(s)) \end{array}}}
\vnbprop{gauges-in-fun}{\vnbstmt{\forall s.; \forall g \in \operatorname{euclidean\text{-}gauges}(s).\; (g \in (\operatorname{carr}(s) \to \mathbb{N}))}}
\vnbprop{gauges-spec}{\vnbstmt{\forall s.; \forall g \in \operatorname{euclidean\text{-}gauges}(s).\; \operatorname{has\text{-}div\text{-}remainder}(s, g)}}
\vnbprop{gauge-is-degree}{\vnbstmt{\begin{array}{@{}l@{}} \forall s. \\ \quad \operatorname{is\text{-}euclidean\text{-}ring}(s) \Rightarrow  \\ \quad \quad (\operatorname{gauge}(s) \in (\operatorname{carr}(s) \to \mathbb{N})) \wedge \operatorname{has\text{-}div\text{-}remainder}(s, \operatorname{gauge}(s)) \end{array}}}

\vnbbookfile{theorem-library/euclidean-ideal-generator-proof}{theorem-library/euclidean-ideal-generator-proof}
\noindent{\small euclidean-ideal-generator-proof.scm -- every ideal of a Euclidean ring has a generator: some b in I with I subset (b).}\par
\vnbprop{euclidean-ideal-has-generator}{\vnbstmt{\begin{array}{@{}l@{}} \forall s, i. \\ \quad (\operatorname{is\text{-}euclidean\text{-}ring}(s) \wedge \operatorname{is\text{-}ideal}(s, i)) \Rightarrow  \\ \quad \quad \exists b \in i.\; \forall a \in i.\; (a \in \operatorname{principal\text{-}ideal}(s, b)) \end{array}}}

\vnbbookfile{theorem-library/zz-bezout-proof}{theorem-library/zz-bezout-proof}
\noindent{\small zz-bezout-proof.scm -- BEZOUT on the integers, from the Euclidean-ring core.}\par
\vnbprop{zz-bezout-set-is-ideal}{\vnbstmt{\begin{array}{@{}l@{}} \forall a, b. \\ \quad ((a \in \mathbb{Z}) \wedge (b \in \mathbb{Z})) \Rightarrow \operatorname{is\text{-}ideal}(\operatorname{zz\text{-}ring}, \operatorname{zz\text{-}bezout\text{-}set}(a, b)) \end{array}}}
\vnbprop{zz-bezout}{\vnbstmt{\begin{array}{@{}l@{}} \forall a, b. \\ \quad ((a \in \mathbb{Z}) \wedge (b \in \mathbb{Z})) \Rightarrow  \\ \quad \quad \exists d \in \operatorname{zz\text{-}bezout\text{-}set}(a, b).\; (\operatorname{zz\text{-}divides}(d, a) \wedge \operatorname{zz\text{-}divides}(d, b)) \end{array}}}

\vnbbookfile{theorem-library/euclidean-division-proof}{theorem-library/euclidean-division-proof}
\noindent{\small euclidean-division-proof.scm -- Brick B2 (LA Phase B, algebraic-numbers.pdf ch.3): division-with-remainder in usable form over any euclidean ring.  For a euclidean ring s, a in CARR, b in CARR, b/=0, there are q,r with a = q.b + r and (r = 0 or deg(r) $<$ deg(b)), where deg = GAUGE(s) (the named degree).  This}\par
\vnbprop{euclidean-division}{\vnbstmt{\begin{array}{@{}l@{}} \forall s.; \forall a, b \in \operatorname{carr}(s). \\ \quad (\operatorname{is\text{-}euclidean\text{-}ring}(s) \wedge \neg (b = \operatorname{zero}(s))) \Rightarrow  \\ \quad \quad \exists q, r \in \operatorname{carr}(s). \\ \quad \quad \quad ((a = \operatorname{add}(s)(\operatorname{mul}(s)(q, b), r)) \wedge ((r = \operatorname{zero}(s)) \vee ((\operatorname{gauge}(s)(r) + 1) \leq \operatorname{gauge}(s)(b)))) \end{array}}}

\vnbbookfile{theorem-library/binomial-proof}{theorem-library/binomial-proof}
\noindent{\small binomial-proof.scm -- machine proof of the Binomial Theorem (SUM form), PROVEN by ordinary induction, IMPS-style (docs/algebra.t).}\par
\vnbprop{sum-expansion}{\vnbstmt{\begin{array}{@{}l@{}} \forall n \in \mathbb{N}, r.; \forall x, y \in \operatorname{carr}(r), g \in (\mathbb{Z} \to \operatorname{carr}(r)). \\ \quad \operatorname{is\text{-}commutative\text{-}ring}(r) \Rightarrow  \\ \quad \quad \operatorname{sum}(r, (\lambda k \in \mathbb{Z}.\; \operatorname{add}(r)(\operatorname{mul}(r)(x, g((k - 1))), \operatorname{mul}(r)(y, g(k)))), (n + 1)) \\ \quad \quad =\; \operatorname{add}(r)(\operatorname{add}(r)(\operatorname{mul}(r)(\operatorname{add}(r)(x, y), \operatorname{sum}(r, g, n)), \operatorname{mul}(r)(x, g(-1))), \operatorname{mul}(r)(y, g(n))) \end{array}}}
\vnbprop{binomial-theorem}{\vnbstmt{\begin{array}{@{}l@{}} \forall n \in \mathbb{N}, r.; \forall x, y \in \operatorname{carr}(r). \\ \quad \operatorname{is\text{-}commutative\text{-}ring}(r) \Rightarrow  \\ \quad \quad \operatorname{ring\text{-}power}(r, \operatorname{add}(r)(x, y), n) = \operatorname{sum}(r, \operatorname{comb\text{-}kk}(r, x, y, n), (n + 1)) \end{array}}}

\vnbbookfile{theorem-library/poly-is-ring-proof}{theorem-library/poly-is-ring-proof}
\noindent{\small poly-is-ring-proof.scm -- POLY(A) is a ring when A is a ring.}\par
\vnbprop{poly-is-ring}{\vnbstmt{\forall a.\; (\operatorname{is\text{-}ring}(a) \Rightarrow \operatorname{is\text{-}ring}(\operatorname{poly}(a)))}}

\vnbbookfile{theorem-library/rake-asm-2}{theorem-library/rake-asm-2}
\noindent{\small rake-asm-2.scm -- RAKE-ASM agent RA-2 (2026-10-12): elementary laws that stood as ASSUMPTIONS (add-axiom! / support), each proven here under its EXACT name and statement (copied character for character from its definition site).}\par
\vnbprop{prod-ord-singleton}{\vnbstmt{\begin{array}{@{}l@{}} \forall m.; \forall f \in (\mathbb{N} \to \operatorname{carr}(m)). \\ \quad \operatorname{is\text{-}monoid}(m) \Rightarrow (\operatorname{prod\text{-}ord}(m, f, 1) = f(0)) \end{array}}}
\vnbprop{sum-singleton}{\vnbstmt{\forall r.; \forall f \in (\mathbb{N} \to \operatorname{carr}(r)).\; (\operatorname{is\text{-}ring}(r) \Rightarrow (\operatorname{sum}(r, f, 1) = f(0)))}}
\vnbprop{sum-ag-singleton}{\vnbstmt{\begin{array}{@{}l@{}} \forall g.; \forall f \in (\mathbb{N} \to \operatorname{carr}(g)). \\ \quad \operatorname{is\text{-}abelian\text{-}group}(g) \Rightarrow (\operatorname{sum\text{-}ag}(g, f, 1) = f(0)) \end{array}}}
\vnbprop{prod-ord-type-ind}{\vnbstmt{\begin{array}{@{}l@{}} \forall n \in \mathbb{N}, m.; \forall f \in (\mathbb{N} \to \operatorname{carr}(m)). \\ \quad \operatorname{is\text{-}monoid}(m) \Rightarrow (\operatorname{prod\text{-}ord}(m, f, n) \in \operatorname{carr}(m)) \end{array}}}
\vnbprop{prod-ord-type}{\vnbstmt{\begin{array}{@{}l@{}} \forall m.; \forall f \in (\mathbb{N} \to \operatorname{carr}(m)), n \in \mathbb{N}. \\ \quad \operatorname{is\text{-}monoid}(m) \Rightarrow (\operatorname{prod\text{-}ord}(m, f, n) \in \operatorname{carr}(m)) \end{array}}}
\vnbprop{sum-type-ind}{\vnbstmt{\begin{array}{@{}l@{}} \forall n \in \mathbb{N}, r.; \forall f \in (\mathbb{N} \to \operatorname{carr}(r)). \\ \quad \operatorname{is\text{-}ring}(r) \Rightarrow (\operatorname{sum}(r, f, n) \in \operatorname{carr}(r)) \end{array}}}
\vnbprop{sum-type}{\vnbstmt{\begin{array}{@{}l@{}} \forall r.; \forall f \in (\mathbb{N} \to \operatorname{carr}(r)), n \in \mathbb{N}. \\ \quad \operatorname{is\text{-}ring}(r) \Rightarrow (\operatorname{sum}(r, f, n) \in \operatorname{carr}(r)) \end{array}}}
\vnbprop{sum-ag-as-reduce-ind}{\vnbstmt{\begin{array}{@{}l@{}} \forall k \in \mathbb{N}, g.; \forall f \in (\mathbb{N} \to \operatorname{carr}(g)). \\ \quad \operatorname{is\text{-}abelian\text{-}group}(g) \Rightarrow (\operatorname{sum\text{-}ag}(g, f, (k + 1)) = \operatorname{reduce}(\operatorname{opr}(g), f, (k + 1))) \end{array}}}
\vnbprop{sum-ag-as-reduce}{\vnbstmt{\begin{array}{@{}l@{}} \forall g.; \forall f \in (\mathbb{N} \to \operatorname{carr}(g)), n \in \mathbb{N}. \\ \quad (\operatorname{is\text{-}abelian\text{-}group}(g) \wedge (1 \leq n)) \Rightarrow (\operatorname{sum\text{-}ag}(g, f, n) = \operatorname{reduce}(\operatorname{opr}(g), f, n)) \end{array}}}
\vnbprop{mpow-mult-ind}{\vnbstmt{\begin{array}{@{}l@{}} \forall n \in \mathbb{N}, m.; \forall x, y \in \operatorname{carr}(m). \\ \quad \operatorname{is\text{-}comm\text{-}monoid}(m) \Rightarrow  \\ \quad \quad \operatorname{mpow}(m, \operatorname{opr}(m)(x, y), n) = \operatorname{opr}(m)(\operatorname{mpow}(m, x, n), \operatorname{mpow}(m, y, n)) \end{array}}}
\vnbprop{mpow-mult}{\vnbstmt{\begin{array}{@{}l@{}} \forall m.; \forall x, y \in \operatorname{carr}(m), n \in \mathbb{N}. \\ \quad \operatorname{is\text{-}comm\text{-}monoid}(m) \Rightarrow  \\ \quad \quad \operatorname{mpow}(m, \operatorname{opr}(m)(x, y), n) = \operatorname{opr}(m)(\operatorname{mpow}(m, x, n), \operatorname{mpow}(m, y, n)) \end{array}}}
\vnbprop{sum-left-scalar-ind}{\vnbstmt{\begin{array}{@{}l@{}} \forall n \in \mathbb{N}, r.; \forall a \in \operatorname{carr}(r), f \in (\mathbb{N} \to \operatorname{carr}(r)). \\ \quad \operatorname{is\text{-}ring}(r) \Rightarrow  \\ \quad \quad \operatorname{sum}(r, (\lambda i \in \mathbb{N}.\; \operatorname{mul}(r)(a, f(i))), n) = \operatorname{mul}(r)(a, \operatorname{sum}(r, f, n)) \end{array}}}
\vnbprop{sum-left-scalar}{\vnbstmt{\begin{array}{@{}l@{}} \forall r.; \forall a \in \operatorname{carr}(r), f \in (\mathbb{N} \to \operatorname{carr}(r)), n \in \mathbb{N}. \\ \quad \operatorname{is\text{-}ring}(r) \Rightarrow  \\ \quad \quad \operatorname{sum}(r, (\lambda i \in \mathbb{N}.\; \operatorname{mul}(r)(a, f(i))), n) = \operatorname{mul}(r)(a, \operatorname{sum}(r, f, n)) \end{array}}}
\vnbprop{zz-act-distrib-nn}{\vnbstmt{\begin{array}{@{}l@{}} \forall k \in \mathbb{N}, g.; \forall a, b \in \operatorname{carr}(g). \\ \quad \operatorname{is\text{-}abelian\text{-}group}(g) \Rightarrow  \\ \quad \quad \operatorname{zz\text{-}act}(g, k, \operatorname{opr}(g)(a, b)) = \operatorname{opr}(g)(\operatorname{zz\text{-}act}(g, k, a), \operatorname{zz\text{-}act}(g, k, b)) \end{array}}}
\vnbprop{zz-act-distrib}{\vnbstmt{\begin{array}{@{}l@{}} \forall g.; \forall k \in \mathbb{Z}, a, b \in \operatorname{carr}(g). \\ \quad \operatorname{is\text{-}abelian\text{-}group}(g) \Rightarrow  \\ \quad \quad \operatorname{zz\text{-}act}(g, k, \operatorname{opr}(g)(a, b)) = \operatorname{opr}(g)(\operatorname{zz\text{-}act}(g, k, a), \operatorname{zz\text{-}act}(g, k, b)) \end{array}}}
\vnbprop{zz-act-assoc-nn}{\vnbstmt{\begin{array}{@{}l@{}} \forall j \in \mathbb{N}, g.; \forall k \in \mathbb{Z}, a \in \operatorname{carr}(g). \\ \quad \operatorname{is\text{-}abelian\text{-}group}(g) \Rightarrow (\operatorname{zz\text{-}act}(g, (j \cdot k), a) = \operatorname{zz\text{-}act}(g, j, \operatorname{zz\text{-}act}(g, k, a))) \end{array}}}
\vnbprop{zz-act-assoc}{\vnbstmt{\begin{array}{@{}l@{}} \forall g.; \forall j, k \in \mathbb{Z}, a \in \operatorname{carr}(g). \\ \quad \operatorname{is\text{-}abelian\text{-}group}(g) \Rightarrow (\operatorname{zz\text{-}act}(g, (j \cdot k), a) = \operatorname{zz\text{-}act}(g, j, \operatorname{zz\text{-}act}(g, k, a))) \end{array}}}

\vnbbookfile{theorem-library/rake-setoid2}{theorem-library/rake-setoid2}
\noindent{\small A file shared across books; the entries that belong to this one.}\par
\noindent{\small rake-setoid2.scm -- BATCH O of the 2026-09-17 rake: the SETOID quotient and the METRIC COMPLETION.  Twelve of the thirteen asserted leaves of the batch are PROVEN here; the thirteenth is a STATEMENT DEFECT and is recorded below. Every statement is its support's statement UNCHANGED (copied from the}\par
\vnbprop{rko2-related-unfold}{\vnbstmt{\forall s, a, b.\; (\operatorname{related}(s, a, b) \Leftrightarrow (\langle a, b \rangle \in \operatorname{rel}(s)))}}
\vnbprop{rko2-refl}{\vnbstmt{\forall s.; \forall a \in \operatorname{pts}(s).\; (\operatorname{is\text{-}setoid}(s) \Rightarrow \operatorname{related}(s, a, a))}}
\vnbprop{rko2-sym}{\vnbstmt{\begin{array}{@{}l@{}} \forall s, a, b. \\ \quad (\operatorname{is\text{-}setoid}(s) \wedge ((a \in \operatorname{pts}(s)) \wedge ((b \in \operatorname{pts}(s)) \wedge \operatorname{related}(s, a, b)))) \Rightarrow  \\ \quad \quad \operatorname{related}(s, b, a) \end{array}}}
\vnbprop{rko2-trans}{\vnbstmt{\begin{array}{@{}l@{}} \forall s, a, b, c. \\ \quad \operatorname{is\text{-}setoid}(s) \wedge \\ \quad a \in \operatorname{pts}(s) \wedge \\ \quad b \in \operatorname{pts}(s) \wedge \\ \quad c \in \operatorname{pts}(s) \wedge \\ \quad \operatorname{related}(s, a, b) \wedge \\ \quad \operatorname{related}(s, b, c) \Rightarrow \\ \quad \quad \operatorname{related}(s, a, c) \end{array}}}
\vnbprop{class-self}{\vnbstmt{\forall s.; \forall a \in \operatorname{pts}(s).\; (\operatorname{is\text{-}setoid}(s) \Rightarrow (a \in \operatorname{class}(s, a)))}}
\vnbprop{class-in-quotient}{\vnbstmt{\forall s.; \forall a \in \operatorname{pts}(s).\; (\operatorname{is\text{-}setoid}(s) \Rightarrow (\operatorname{class}(s, a) \in \operatorname{quotient}(s)))}}
\vnbprop{quotient-rep}{\vnbstmt{\begin{array}{@{}l@{}} \forall s.; \forall x \in \operatorname{quotient}(s). \\ \quad \operatorname{is\text{-}setoid}(s) \Rightarrow \exists a \in \operatorname{pts}(s).\; (x = \operatorname{class}(s, a)) \end{array}}}
\vnbprop{class-subset-carrier}{\vnbstmt{\forall s.; \forall a \in \operatorname{pts}(s).\; (\operatorname{is\text{-}setoid}(s) \Rightarrow (\operatorname{class}(s, a) \subseteq \operatorname{pts}(s)))}}
\vnbprop{class-eq-iff}{\vnbstmt{\begin{array}{@{}l@{}} \forall s.; \forall a, b \in \operatorname{pts}(s). \\ \quad \operatorname{is\text{-}setoid}(s) \Rightarrow ((\operatorname{class}(s, a) = \operatorname{class}(s, b)) \Leftrightarrow \operatorname{related}(s, a, b)) \end{array}}}
\vnbprop{class-disjoint}{\vnbstmt{\begin{array}{@{}l@{}} \forall s.; \forall a, b \in \operatorname{pts}(s). \\ \quad \operatorname{is\text{-}setoid}(s) \Rightarrow  \\ \quad \quad (\operatorname{class}(s, a) = \operatorname{class}(s, b)) \vee ((\operatorname{class}(s, a) \cap \operatorname{class}(s, b)) = \emptyset) \end{array}}}
\vnbprop{descend-computes}{\vnbstmt{\begin{array}{@{}l@{}} \forall s, z.; \forall f \in (\operatorname{pts}(s) \to z), a \in \operatorname{pts}(s). \\ \quad (\operatorname{is\text{-}setoid}(s) \wedge \operatorname{respects}(s, f)) \Rightarrow (\operatorname{descend}(s, f)(\operatorname{class}(s, a)) = f(a)) \end{array}}}
\vnbprop{descend2-computes}{\vnbstmt{\begin{array}{@{}l@{}} \forall s, z.; \forall f \in ((\operatorname{pts}(s) \times \operatorname{pts}(s)) \to z), a, b. \\ \quad (\operatorname{is\text{-}setoid}(s) \wedge (\operatorname{respects2}(s, f) \wedge ((a \in \operatorname{pts}(s)) \wedge (b \in \operatorname{pts}(s))))) \Rightarrow  \\ \quad \quad \operatorname{descend2}(s, f)(\operatorname{class}(s, a), \operatorname{class}(s, b)) = f(a, b) \end{array}}}
\vnbprop{quotient-universal}{\vnbstmt{\begin{array}{@{}l@{}} \forall s, z.; \forall f \in (\operatorname{pts}(s) \to z). \\ \quad (\operatorname{is\text{-}setoid}(s) \wedge \operatorname{respects}(s, f)) \Rightarrow  \\ \quad \quad \exists g \in (\operatorname{quotient}(s) \to z). \\ \quad \quad \quad \forall a \in \operatorname{pts}(s).\; (g(\operatorname{class}(s, a)) = f(a)) \wedge \\ \quad \quad \quad \forall b \in (\operatorname{quotient}(s) \to z). \\ \quad \quad \quad \quad \forall a \in \operatorname{pts}(s).\; (b(\operatorname{class}(s, a)) = f(a)) \Rightarrow (b = g) \end{array}}}

\vnbbookfile{theorem-library/ringoid-setoid-proof}{theorem-library/ringoid-setoid-proof}
\noindent{\small ringoid-setoid-proof.scm}\par
\vnbprop{ringoid-rel-is-equivalence}{\vnbstmt{\forall r.\; (\operatorname{is\text{-}ringoid}(r) \Rightarrow \operatorname{is\text{-}equivalence}(\operatorname{ringoid\text{-}rel}(r), \operatorname{carr}(r)))}}
\vnbprop{ringoid-setoid-is-setoid}{\vnbstmt{\forall r.\; (\operatorname{is\text{-}ringoid}(r) \Rightarrow \operatorname{is\text{-}setoid}(\operatorname{ringoid\text{-}setoid}(r)))}}

\vnbbookfile{theorem-library/rake-ringoid}{theorem-library/rake-ringoid}
\noindent{\small rake-ringoid.scm -- rq-add-computes: addition on R/I computes on classes.}\par
\vnbprop{r6b-setoid-rel}{\vnbstmt{\forall r.\; (\operatorname{is\text{-}ringoid}(r) \Rightarrow (\operatorname{rel}(\operatorname{ringoid\text{-}setoid}(r)) \simeq \operatorname{ringoid\text{-}rel}(r)))}}
\vnbprop{r6b-related-out}{\vnbstmt{\begin{array}{@{}l@{}} \forall r, a, b. \\ \quad (\operatorname{is\text{-}ringoid}(r) \wedge ((a \in \operatorname{carr}(r)) \wedge ((b \in \operatorname{carr}(r)) \wedge \operatorname{related}(\operatorname{ringoid\text{-}setoid}(r), a, b)))) \Rightarrow  \\ \quad \quad \operatorname{add}(r)(a, \operatorname{neg}(r)(b)) \in \operatorname{idl}(r) \end{array}}}
\vnbprop{r6b-related-in}{\vnbstmt{\begin{array}{@{}l@{}} \forall r, a, b. \\ \quad (\operatorname{is\text{-}ringoid}(r) \wedge ((a \in \operatorname{carr}(r)) \wedge ((b \in \operatorname{carr}(r)) \wedge (\operatorname{add}(r)(a, \operatorname{neg}(r)(b)) \in \operatorname{idl}(r))))) \Rightarrow  \\ \quad \quad \operatorname{related}(\operatorname{ringoid\text{-}setoid}(r), a, b) \end{array}}}
\vnbprop{r6b-rq-add-lam-type}{\vnbstmt{\begin{array}{@{}l@{}} \forall r. \\ \quad \operatorname{is\text{-}ringoid}(r) \Rightarrow  \\ \quad \quad (\lambda \langle a, b \rangle \in (\operatorname{carr}(r) \times \operatorname{carr}(r)).\; \operatorname{class}(\operatorname{ringoid\text{-}setoid}(r), \operatorname{add}(r)(a, b))) \\ \quad \quad \in\; ((\operatorname{pts}(\operatorname{ringoid\text{-}setoid}(r)) \times \operatorname{pts}(\operatorname{ringoid\text{-}setoid}(r))) \to \operatorname{quotient}(\operatorname{ringoid\text{-}setoid}(r))) \end{array}}}
\vnbprop{r6b-add-respects2}{\vnbstmt{\begin{array}{@{}l@{}} \forall r. \\ \quad \operatorname{is\text{-}ringoid}(r) \Rightarrow  \\ \quad \quad \operatorname{respects2}(\operatorname{ringoid\text{-}setoid}(r), \\ \quad \quad \quad (\lambda \langle a, b \rangle \in (\operatorname{carr}(r) \times \operatorname{carr}(r)).\; \operatorname{class}(\operatorname{ringoid\text{-}setoid}(r), \operatorname{add}(r)(a, b)))) \end{array}}}
\vnbprop{rq-add-computes}{\vnbstmt{\begin{array}{@{}l@{}} \forall r, a, b. \\ \quad (\operatorname{is\text{-}ringoid}(r) \wedge ((a \in \operatorname{carr}(r)) \wedge (b \in \operatorname{carr}(r)))) \Rightarrow  \\ \quad \quad \operatorname{add}(\operatorname{ringoid\text{-}quotient\text{-}ring}(r))(\operatorname{class}(\operatorname{ringoid\text{-}setoid}(r), a), \operatorname{class}(\operatorname{ringoid\text{-}setoid}(r), b)) \\ \quad \quad =\; \operatorname{class}(\operatorname{ringoid\text{-}setoid}(r), \operatorname{add}(r)(a, b)) \end{array}}}

\vnbbookfile{theorem-library/rake-ringoid-quotient}{theorem-library/rake-ringoid-quotient}
\noindent{\small rake-ringoid-quotient.scm -- MUL and NEG compute on R/I.}\par
\vnbprop{r7n-rmul-type}{\vnbstmt{\begin{array}{@{}l@{}} \forall r, a, b. \\ \quad (\operatorname{is\text{-}ringoid}(r) \wedge ((a \in \operatorname{carr}(r)) \wedge (b \in \operatorname{carr}(r)))) \Rightarrow  \\ \quad \quad \operatorname{mul}(r)(a, b) \in \operatorname{carr}(r) \end{array}}}
\vnbprop{r7n-rone-type}{\vnbstmt{\forall r.\; (\operatorname{is\text{-}ringoid}(r) \Rightarrow (\operatorname{one}(r) \in \operatorname{carr}(r)))}}
\vnbprop{r7n-radd-idem-zero}{\vnbstmt{\begin{array}{@{}l@{}} \forall r.; \forall a \in \operatorname{carr}(r). \\ \quad (\operatorname{is\text{-}ringoid}(r) \wedge (\operatorname{add}(r)(a, a) = a)) \Rightarrow (a = \operatorname{zero}(r)) \end{array}}}
\vnbprop{r7n-rzero-mul-left}{\vnbstmt{\forall r.; \forall b \in \operatorname{carr}(r).\; (\operatorname{is\text{-}ringoid}(r) \Rightarrow (\operatorname{mul}(r)(\operatorname{zero}(r), b) = \operatorname{zero}(r)))}}
\vnbprop{r7n-rmul-zero-right}{\vnbstmt{\forall r.; \forall a \in \operatorname{carr}(r).\; (\operatorname{is\text{-}ringoid}(r) \Rightarrow (\operatorname{mul}(r)(a, \operatorname{zero}(r)) = \operatorname{zero}(r)))}}
\vnbprop{r7n-rneg-mul-left}{\vnbstmt{\begin{array}{@{}l@{}} \forall r, a, b. \\ \quad (\operatorname{is\text{-}ringoid}(r) \wedge ((a \in \operatorname{carr}(r)) \wedge (b \in \operatorname{carr}(r)))) \Rightarrow  \\ \quad \quad \operatorname{mul}(r)(\operatorname{neg}(r)(a), b) = \operatorname{neg}(r)(\operatorname{mul}(r)(a, b)) \end{array}}}
\vnbprop{r7n-rmul-neg-right}{\vnbstmt{\begin{array}{@{}l@{}} \forall r, a, b. \\ \quad (\operatorname{is\text{-}ringoid}(r) \wedge ((a \in \operatorname{carr}(r)) \wedge (b \in \operatorname{carr}(r)))) \Rightarrow  \\ \quad \quad \operatorname{mul}(r)(a, \operatorname{neg}(r)(b)) = \operatorname{neg}(r)(\operatorname{mul}(r)(a, b)) \end{array}}}
\vnbprop{r7n-rdiff-mul}{\vnbstmt{\begin{array}{@{}l@{}} \forall r, a, b, u, v. \\ \quad (\operatorname{is\text{-}ringoid}(r) \wedge ((a \in \operatorname{carr}(r)) \wedge ((b \in \operatorname{carr}(r)) \wedge ((u \in \operatorname{carr}(r)) \wedge (v \in \operatorname{carr}(r)))))) \Rightarrow  \\ \quad \quad \operatorname{add}(r)(\operatorname{mul}(r)(a, b), \operatorname{neg}(r)(\operatorname{mul}(r)(u, v))) \\ \quad \quad =\; \operatorname{add}(r)(\operatorname{mul}(r)(\operatorname{add}(r)(a, \operatorname{neg}(r)(u)), b), \operatorname{mul}(r)(u, \operatorname{add}(r)(b, \operatorname{neg}(r)(v)))) \end{array}}}
\vnbprop{r7n-rq-mul-lam-type}{\vnbstmt{\begin{array}{@{}l@{}} \forall r. \\ \quad \operatorname{is\text{-}ringoid}(r) \Rightarrow  \\ \quad \quad (\lambda \langle a, b \rangle \in (\operatorname{carr}(r) \times \operatorname{carr}(r)).\; \operatorname{class}(\operatorname{ringoid\text{-}setoid}(r), \operatorname{mul}(r)(a, b))) \\ \quad \quad \in\; ((\operatorname{pts}(\operatorname{ringoid\text{-}setoid}(r)) \times \operatorname{pts}(\operatorname{ringoid\text{-}setoid}(r))) \to \operatorname{quotient}(\operatorname{ringoid\text{-}setoid}(r))) \end{array}}}
\vnbprop{r7n-mul-respects2}{\vnbstmt{\begin{array}{@{}l@{}} \forall r. \\ \quad \operatorname{is\text{-}ringoid}(r) \Rightarrow  \\ \quad \quad \operatorname{respects2}(\operatorname{ringoid\text{-}setoid}(r), \\ \quad \quad \quad (\lambda \langle a, b \rangle \in (\operatorname{carr}(r) \times \operatorname{carr}(r)).\; \operatorname{class}(\operatorname{ringoid\text{-}setoid}(r), \operatorname{mul}(r)(a, b)))) \end{array}}}
\vnbprop{rq-mul-computes}{\vnbstmt{\begin{array}{@{}l@{}} \forall r, a, b. \\ \quad (\operatorname{is\text{-}ringoid}(r) \wedge ((a \in \operatorname{carr}(r)) \wedge (b \in \operatorname{carr}(r)))) \Rightarrow  \\ \quad \quad \operatorname{mul}(\operatorname{ringoid\text{-}quotient\text{-}ring}(r))(\operatorname{class}(\operatorname{ringoid\text{-}setoid}(r), a), \operatorname{class}(\operatorname{ringoid\text{-}setoid}(r), b)) \\ \quad \quad =\; \operatorname{class}(\operatorname{ringoid\text{-}setoid}(r), \operatorname{mul}(r)(a, b)) \end{array}}}
\vnbprop{r7n-rq-neg-lam-type}{\vnbstmt{\begin{array}{@{}l@{}} \forall r. \\ \quad \operatorname{is\text{-}ringoid}(r) \Rightarrow  \\ \quad \quad (\lambda a \in \operatorname{carr}(r).\; \operatorname{class}(\operatorname{ringoid\text{-}setoid}(r), \operatorname{neg}(r)(a))) \\ \quad \quad \in\; (\operatorname{pts}(\operatorname{ringoid\text{-}setoid}(r)) \to \operatorname{quotient}(\operatorname{ringoid\text{-}setoid}(r))) \end{array}}}
\vnbprop{r7n-neg-respects}{\vnbstmt{\begin{array}{@{}l@{}} \forall r. \\ \quad \operatorname{is\text{-}ringoid}(r) \Rightarrow  \\ \quad \quad \operatorname{respects}(\operatorname{ringoid\text{-}setoid}(r), (\lambda a \in \operatorname{carr}(r).\; \operatorname{class}(\operatorname{ringoid\text{-}setoid}(r), \operatorname{neg}(r)(a)))) \end{array}}}
\vnbprop{rq-neg-computes}{\vnbstmt{\begin{array}{@{}l@{}} \forall r.; \forall a \in \operatorname{carr}(r). \\ \quad \operatorname{is\text{-}ringoid}(r) \Rightarrow  \\ \quad \quad \operatorname{neg}(\operatorname{ringoid\text{-}quotient\text{-}ring}(r))(\operatorname{class}(\operatorname{ringoid\text{-}setoid}(r), a)) = \operatorname{class}(\operatorname{ringoid\text{-}setoid}(r), \operatorname{neg}(r)(a)) \end{array}}}
\vnbprop{ringoid-quotient-is-ring}{\vnbstmt{\forall r.\; (\operatorname{is\text{-}ringoid}(r) \Rightarrow \operatorname{is\text{-}ring}(\operatorname{ringoid\text{-}quotient\text{-}ring}(r)))}}

\vnbbookfile{theorem-library/rake-asm-1}{theorem-library/rake-asm-1}
\noindent{\small A file shared across books; the entries that belong to this one.}\par
\noindent{\small rake-asm-1.scm -- RAKE-ASM agent RA-1 (2026-10-12): elementary assumptions PROVEN.}\par
\vnbprop{semigroup-assoc}{\vnbstmt{\begin{array}{@{}l@{}} \forall s.; \forall a, b, c \in \operatorname{carr}(s). \\ \quad \operatorname{is\text{-}semigroup}(s) \Rightarrow  \\ \quad \quad \operatorname{opr}(s)(\operatorname{opr}(s)(a, b), c) = \operatorname{opr}(s)(a, \operatorname{opr}(s)(b, c)) \end{array}}}
\vnbprop{qq-is-commutative-ring}{\vnbstmt{\operatorname{is\text{-}commutative\text{-}ring}(\operatorname{qq\text{-}ring})}}
\vnbprop{qq-is-integral-domain}{\vnbstmt{\operatorname{is\text{-}integral\text{-}domain}(\operatorname{qq\text{-}ring})}}
\vnbprop{zero-ring-is-ring}{\vnbstmt{\operatorname{is\text{-}ring}(\operatorname{zero\text{-}ring})}}
\vnbprop{ring-prod-is-ring}{\vnbstmt{\forall x, y.\; ((\operatorname{is\text{-}ring}(x) \wedge \operatorname{is\text{-}ring}(y)) \Rightarrow \operatorname{is\text{-}ring}(\operatorname{ring\text{-}prod}(x, y)))}}
\vnbprop{ring-prod-n-is-ring-ind}{\vnbstmt{\begin{array}{@{}l@{}} \forall n \in \mathbb{N}, f. \\ \quad \forall i \in \mathbb{N}.\; \operatorname{is\text{-}ring}(f(i)) \Rightarrow \operatorname{is\text{-}ring}(\operatorname{ring\text{-}prod\text{-}n}(f, n)) \end{array}}}
\vnbprop{ring-prod-n-is-ring}{\vnbstmt{\begin{array}{@{}l@{}} \forall f.; \forall n \in \mathbb{N}. \\ \quad \forall i \in \mathbb{N}.\; \operatorname{is\text{-}ring}(f(i)) \Rightarrow \operatorname{is\text{-}ring}(\operatorname{ring\text{-}prod\text{-}n}(f, n)) \end{array}}}

\vnbbookfile{structure-library/functoriality}{structure-library/functoriality}
\noindent{\small functoriality.scm -- every view-as is a FUNCTOR, and here is the proof.}\par
\vnbprop{field-as-euclidean-ring-functorial}{\vnbstmt{\begin{array}{@{}l@{}} \forall a, b, f. \\ \quad \operatorname{is\text{-}hom\text{-}field}(a, b, f) \Rightarrow  \\ \quad \quad \operatorname{is\text{-}hom\text{-}euclidean\text{-}ring}(\operatorname{field\text{-}as\text{-}euclidean\text{-}ring}(a), \operatorname{field\text{-}as\text{-}euclidean\text{-}ring}(b), f) \end{array}}}
\vnbprop{field-as-integral-domain-functorial}{\vnbstmt{\begin{array}{@{}l@{}} \forall a, b, f. \\ \quad \operatorname{is\text{-}hom\text{-}field}(a, b, f) \Rightarrow  \\ \quad \quad \operatorname{is\text{-}hom\text{-}integral\text{-}domain}(\operatorname{field\text{-}as\text{-}integral\text{-}domain}(a), \operatorname{field\text{-}as\text{-}integral\text{-}domain}(b), f) \end{array}}}
\vnbprop{normed-field-as-commutative-ring-functorial}{\vnbstmt{\begin{array}{@{}l@{}} \forall a, b, f. \\ \quad \operatorname{is\text{-}hom\text{-}normed\text{-}field}(a, b, f) \Rightarrow  \\ \quad \quad \operatorname{is\text{-}hom\text{-}commutative\text{-}ring}(\operatorname{normed\text{-}field\text{-}as\text{-}commutative\text{-}ring}(a), \operatorname{normed\text{-}field\text{-}as\text{-}commutative\text{-}ring}(b), f) \end{array}}}
\vnbprop{normed-field-as-integral-domain-functorial}{\vnbstmt{\begin{array}{@{}l@{}} \forall a, b, f. \\ \quad \operatorname{is\text{-}hom\text{-}normed\text{-}field}(a, b, f) \Rightarrow  \\ \quad \quad \operatorname{is\text{-}hom\text{-}integral\text{-}domain}(\operatorname{normed\text{-}field\text{-}as\text{-}integral\text{-}domain}(a), \operatorname{normed\text{-}field\text{-}as\text{-}integral\text{-}domain}(b), f) \end{array}}}

\part{Linear Algebra}
\noindent Modules and vector spaces; matrices, elementary operations, determinants and the Smith normal form; spans, submodules, bases and rank; finite-dimensional spaces.
\par\noindent\textit{Builds on:} Algebra.
\par\noindent{\small 72 files.}

\vnbbookfile{structure-library/module}{structure-library/module}
\noindent{\small module.scm -- MODULE over a ring, with the scalar action ACT.}\par
\vnbstructdef{module}{A module $s$ is a 6-tuple whose components are listed below.\par
\begin{enumerate}[nosep]
\item $\operatorname{scal}(s) = \operatorname{nth}(1, s)$, a substructure (\texttt{substructure}): $\operatorname{is\text{-}ring}(\operatorname{scal}(s))$.
\item $\operatorname{vec}(s) = \operatorname{nth}(2, s)$, a carrier (\texttt{carriers}): $\operatorname{vec}(s) \in \mathbf{Set}$.
\item $\operatorname{vadd}(s) = \operatorname{nth}(3, s)$, an operation (\texttt{op}): $\operatorname{vadd}(s) \in ((\operatorname{vec}(s) \times \operatorname{vec}(s)) \to \operatorname{vec}(s))$.
\item $\operatorname{vzero}(s) = \operatorname{nth}(4, s)$, a constant (\texttt{constant}): $\operatorname{vzero}(s) \in \operatorname{vec}(s)$.
\item $\operatorname{vneg}(s) = \operatorname{nth}(5, s)$, an operation (\texttt{op}): $\operatorname{vneg}(s) \in (\operatorname{vec}(s) \to \operatorname{vec}(s))$.
\item $\operatorname{act}(s) = \operatorname{nth}(6, s)$, an operation (\texttt{op}): $\operatorname{act}(s) \in ((\operatorname{carr}(\operatorname{scal}(s)) \times \operatorname{vec}(s)) \to \operatorname{vec}(s))$.
\end{enumerate}
The laws (the \texttt{property} and \texttt{law} clauses):
\begin{itemize}[nosep]
\item $\operatorname{is\text{-}associative}(\operatorname{vadd}(s), \operatorname{vec}(s))$
\item $\operatorname{is\text{-}commutative}(\operatorname{vadd}(s), \operatorname{vec}(s))$
\item $\operatorname{is\text{-}identity}(\operatorname{vadd}(s), \operatorname{vzero}(s), \operatorname{vec}(s))$
\item $\operatorname{has\text{-}inverses}(\operatorname{vadd}(s), \operatorname{vzero}(s), \operatorname{vneg}(s), \operatorname{vec}(s))$
\item
\begin{equation*}
\begin{array}{@{}l@{}} \forall \operatorname{r\_} \in \operatorname{carr}(\operatorname{scal}(s)), \operatorname{x\_}, \operatorname{y\_} \in \operatorname{vec}(s). \\ \quad \operatorname{act}(s)(\operatorname{r\_}, \operatorname{vadd}(s)(\operatorname{x\_}, \operatorname{y\_})) = \operatorname{vadd}(s)(\operatorname{act}(s)(\operatorname{r\_}, \operatorname{x\_}), \operatorname{act}(s)(\operatorname{r\_}, \operatorname{y\_})) \end{array}
\end{equation*}

\item
\begin{equation*}
\begin{array}{@{}l@{}} \forall \operatorname{r\_}, \operatorname{s\_} \in \operatorname{carr}(\operatorname{scal}(s)), \operatorname{x\_} \in \operatorname{vec}(s). \\ \quad \operatorname{act}(s)(\operatorname{add}(\operatorname{scal}(s))(\operatorname{r\_}, \operatorname{s\_}), \operatorname{x\_}) = \operatorname{vadd}(s)(\operatorname{act}(s)(\operatorname{r\_}, \operatorname{x\_}), \operatorname{act}(s)(\operatorname{s\_}, \operatorname{x\_})) \end{array}
\end{equation*}

\item
\begin{equation*}
\begin{array}{@{}l@{}} \forall \operatorname{r\_}, \operatorname{s\_} \in \operatorname{carr}(\operatorname{scal}(s)), \operatorname{x\_} \in \operatorname{vec}(s). \\ \quad \operatorname{act}(s)(\operatorname{mul}(\operatorname{scal}(s))(\operatorname{r\_}, \operatorname{s\_}), \operatorname{x\_}) = \operatorname{act}(s)(\operatorname{r\_}, \operatorname{act}(s)(\operatorname{s\_}, \operatorname{x\_})) \end{array}
\end{equation*}

\item $\forall \operatorname{x\_} \in \operatorname{vec}(s).\; (\operatorname{act}(s)(\operatorname{one}(\operatorname{scal}(s)), \operatorname{x\_}) = \operatorname{x\_})$
\end{itemize}
The characteristic law (the axiom \texttt{is-module}): $\operatorname{is\text{-}module}(s)$ holds if and only if $\operatorname{length}(s) = 6$, every component has the typing listed with it, and every law holds.
}
\par\medskip\noindent\textbf{Definition} (\texttt{hom-module}).\;$\operatorname{hom\text{-}module}(a, b) :=$ $\{\operatorname{homf\_} \in (\operatorname{vec}(a) \to \operatorname{vec}(b)) : \operatorname{is\text{-}hom\text{-}module}(a, b, \operatorname{homf\_})\}$
\vnbdef{module-scalar-ring}{\vnbstmt{\forall m.\; (\operatorname{is\text{-}module}(m) \Rightarrow \operatorname{is\text{-}ring}(\operatorname{scal}(m)))}}
\vnbdef{module-scalar-zero-in}{\vnbstmt{\forall m.\; (\operatorname{is\text{-}module}(m) \Rightarrow (\operatorname{zero}(\operatorname{scal}(m)) \in \operatorname{carr}(\operatorname{scal}(m))))}}
\vnbdef{module-vzero-in}{\vnbstmt{\forall m.\; (\operatorname{is\text{-}module}(m) \Rightarrow (\operatorname{vzero}(m) \in \operatorname{vec}(m)))}}
\vnbdef{module-act-type}{\vnbstmt{\begin{array}{@{}l@{}} \forall m.; \forall r \in \operatorname{carr}(\operatorname{scal}(m)), x \in \operatorname{vec}(m). \\ \quad \operatorname{is\text{-}module}(m) \Rightarrow (\operatorname{act}(m)(r, x) \in \operatorname{vec}(m)) \end{array}}}
\vnbdef{module-vadd-type}{\vnbstmt{\begin{array}{@{}l@{}} \forall m.; \forall x, y \in \operatorname{vec}(m). \\ \quad \operatorname{is\text{-}module}(m) \Rightarrow (\operatorname{vadd}(m)(x, y) \in \operatorname{vec}(m)) \end{array}}}
\vnbdef{module-vneg-type}{\vnbstmt{\forall m.; \forall x \in \operatorname{vec}(m).\; (\operatorname{is\text{-}module}(m) \Rightarrow (\operatorname{vneg}(m)(x) \in \operatorname{vec}(m)))}}
\vnbdef{module-act-distrib-vec}{\vnbstmt{\begin{array}{@{}l@{}} \forall m.; \forall r \in \operatorname{carr}(\operatorname{scal}(m)), x, y \in \operatorname{vec}(m). \\ \quad \operatorname{is\text{-}module}(m) \Rightarrow  \\ \quad \quad \operatorname{act}(m)(r, \operatorname{vadd}(m)(x, y)) = \operatorname{vadd}(m)(\operatorname{act}(m)(r, x), \operatorname{act}(m)(r, y)) \end{array}}}
\vnbdef{module-act-distrib-scalar}{\vnbstmt{\begin{array}{@{}l@{}} \forall m.; \forall r, s \in \operatorname{carr}(\operatorname{scal}(m)), x \in \operatorname{vec}(m). \\ \quad \operatorname{is\text{-}module}(m) \Rightarrow  \\ \quad \quad \operatorname{act}(m)(\operatorname{add}(\operatorname{scal}(m))(r, s), x) = \operatorname{vadd}(m)(\operatorname{act}(m)(r, x), \operatorname{act}(m)(s, x)) \end{array}}}
\vnbdef{module-act-mul-compat}{\vnbstmt{\begin{array}{@{}l@{}} \forall m.; \forall r, s \in \operatorname{carr}(\operatorname{scal}(m)), x \in \operatorname{vec}(m). \\ \quad \operatorname{is\text{-}module}(m) \Rightarrow  \\ \quad \quad \operatorname{act}(m)(\operatorname{mul}(\operatorname{scal}(m))(r, s), x) = \operatorname{act}(m)(r, \operatorname{act}(m)(s, x)) \end{array}}}
\vnbdef{module-act-unital}{\vnbstmt{\forall m.; \forall x \in \operatorname{vec}(m).\; (\operatorname{is\text{-}module}(m) \Rightarrow (\operatorname{act}(m)(\operatorname{one}(\operatorname{scal}(m)), x) = x))}}
\vnbdef{is-hom-module-def}{\vnbstmt{\begin{array}{@{}l@{}} \forall a, b, f. \\ \quad \operatorname{is\text{-}hom\text{-}module}(a, b, f) \\ \quad \Leftrightarrow\; \operatorname{is\text{-}module}(a) \wedge \\ \quad \operatorname{is\text{-}module}(b) \wedge \\ \quad f \in (\operatorname{vec}(a) \to \operatorname{vec}(b)) \wedge \\ \quad \operatorname{scal}(a) = \operatorname{scal}(b) \wedge \\ \quad \forall x_{1}, x_{2} \in \operatorname{vec}(a). \\ \quad \quad f(\operatorname{vadd}(a)(x_{1}, x_{2})) = \operatorname{vadd}(b)(f(x_{1}), f(x_{2})) \wedge \\ \quad f(\operatorname{vzero}(a)) = \operatorname{vzero}(b) \wedge \\ \quad \forall x_{1} \in \operatorname{vec}(a).\; (f(\operatorname{vneg}(a)(x_{1})) = \operatorname{vneg}(b)(f(x_{1}))) \wedge \\ \quad \forall x_{1} \in \operatorname{carr}(\operatorname{scal}(a)), x_{2} \in \operatorname{vec}(a). \\ \quad \quad f(\operatorname{act}(a)(x_{1}, x_{2})) = \operatorname{act}(b)(x_{1}, f(x_{2})) \end{array}}}
\vnbdef{module-class}{\vnbstmt{\forall s.\; ((s \in \mathit{module}) \Leftrightarrow \operatorname{is\text{-}module}(s))}}

\vnbbookfile{structure-library/finite-dimensional}{structure-library/finite-dimensional}
\noindent{\small finite-dimensional.scm -- finite-dimensional vector spaces, Zorn-free.}\par
\par\medskip\noindent\textbf{Definition} (\texttt{is-finite-dimensional}).\;For every $m$, 
\begin{equation*}
\operatorname{is\text{-}finite\text{-}dimensional}(m) \Leftrightarrow (\operatorname{is\text{-}vector\text{-}space}(m) \wedge \operatorname{is\text{-}noetherian}(m))
\end{equation*}

\par\medskip\noindent\textbf{Definition} (\texttt{is-submodule}).\;For all $m$ and $s$, 
\begin{equation*}
\begin{array}{@{}l@{}} \operatorname{is\text{-}submodule}(m, s) \\ \Leftrightarrow\; s \subseteq \operatorname{vec}(m) \wedge \\ \operatorname{vzero}(m) \in s \wedge \\ \forall \operatorname{x\_}, \operatorname{y\_} \in s.\; (\operatorname{vadd}(m)(\operatorname{x\_}, \operatorname{y\_}) \in s) \wedge \\ \forall \operatorname{x\_} \in s.\; (\operatorname{vneg}(m)(\operatorname{x\_}) \in s) \wedge \\ \forall \operatorname{r\_} \in \operatorname{carr}(\operatorname{scal}(m)), \operatorname{x\_} \in s.\; (\operatorname{act}(m)(\operatorname{r\_}, \operatorname{x\_}) \in s) \end{array}
\end{equation*}

\par\medskip\noindent\textbf{Definition} (\texttt{is-noetherian}).\;For every $m$, 
\begin{equation*}
\begin{array}{@{}l@{}} \operatorname{is\text{-}noetherian}(m) \\ \Leftrightarrow\; \operatorname{is\text{-}module}(m) \wedge \\ \forall \operatorname{f\_} \in (\mathbb{N} \to \operatorname{power}(\operatorname{vec}(m))). \\ \quad (\forall \operatorname{n\_} \in \mathbb{N}.\; \operatorname{is\text{-}submodule}(m, \operatorname{f\_}(\operatorname{n\_})) \wedge \forall \operatorname{n\_} \in \mathbb{N}.\; (\operatorname{f\_}(\operatorname{n\_}) \subseteq \operatorname{f\_}(\operatorname{succ}(\operatorname{n\_})))) \Rightarrow  \\ \quad \quad \exists \operatorname{k\_} \in \mathbb{N}.\; \forall \operatorname{n\_}.\; (((\operatorname{n\_} \in \mathbb{N}) \wedge (\operatorname{k\_} \leq \operatorname{n\_})) \Rightarrow (\operatorname{f\_}(\operatorname{n\_}) = \operatorname{f\_}(\operatorname{k\_}))) \end{array}
\end{equation*}

\par\medskip\noindent\textbf{Definition} (\texttt{is-subspace}).\;For all $m$ and $s$,  $\operatorname{is\text{-}subspace}(m, s) \Leftrightarrow \operatorname{is\text{-}submodule}(m, s)$
\par\medskip\noindent\textbf{Definition} (\texttt{span-add-one}).\;$\operatorname{span\text{-}add\text{-}one}(m, s, v) :=$
\begin{equation*}
\{\operatorname{y\_} \in \operatorname{vec}(m) : \exists \operatorname{x\_} \in s, \operatorname{r\_} \in \mathbb{R}.\; (\operatorname{y\_} = \operatorname{vadd}(m)(\operatorname{x\_}, \operatorname{act}(m)(\operatorname{r\_}, v)))\}
\end{equation*}

\par\medskip\noindent\textbf{Definition} (\texttt{hom-vector-space}).\;$\operatorname{hom\text{-}vector\text{-}space}(a, b) :=$
\begin{equation*}
\{\operatorname{homf\_} \in (\operatorname{vec}(a) \to \operatorname{vec}(b)) : \operatorname{is\text{-}hom\text{-}vector\text{-}space}(a, b, \operatorname{homf\_})\}
\end{equation*}

\vnbdef{submodule-subset}{\vnbstmt{\forall m, s.\; (\operatorname{is\text{-}submodule}(m, s) \Rightarrow (s \subseteq \operatorname{vec}(m)))}}
\vnbdef{submodule-vzero-in}{\vnbstmt{\forall m, s.\; (\operatorname{is\text{-}submodule}(m, s) \Rightarrow (\operatorname{vzero}(m) \in s))}}
\vnbdef{submodule-vadd-closed}{\vnbstmt{\forall m, s.; \forall x, y \in s.\; (\operatorname{is\text{-}submodule}(m, s) \Rightarrow (\operatorname{vadd}(m)(x, y) \in s))}}
\vnbdef{submodule-vneg-closed}{\vnbstmt{\forall m, s.; \forall x \in s.\; (\operatorname{is\text{-}submodule}(m, s) \Rightarrow (\operatorname{vneg}(m)(x) \in s))}}
\vnbdef{submodule-act-closed}{\vnbstmt{\begin{array}{@{}l@{}} \forall m, s.; \forall r \in \operatorname{carr}(\operatorname{scal}(m)), x \in s. \\ \quad \operatorname{is\text{-}submodule}(m, s) \Rightarrow (\operatorname{act}(m)(r, x) \in s) \end{array}}}
\vnbdef{is-hom-vector-space-def}{\vnbstmt{\begin{array}{@{}l@{}} \forall a, b, f. \\ \quad \operatorname{is\text{-}hom\text{-}vector\text{-}space}(a, b, f) \\ \quad \Leftrightarrow\; \operatorname{is\text{-}vector\text{-}space}(a) \wedge (\operatorname{is\text{-}vector\text{-}space}(b) \wedge \operatorname{is\text{-}hom\text{-}module}(a, b, f)) \end{array}}}
\vnbdef{vector-space-class}{\vnbstmt{\forall s.\; ((s \in \operatorname{vector\text{-}space}) \Leftrightarrow \operatorname{is\text{-}vector\text{-}space}(s))}}
\vnbdef{is-vector-space-def}{\vnbstmt{\forall s.\; (\operatorname{is\text{-}vector\text{-}space}(s) \Leftrightarrow (\operatorname{is\text{-}module}(s) \wedge \operatorname{is\text{-}field\text{-}ring}(\operatorname{scal}(s))))}}

\vnbbookfile{structure-library/transpose}{structure-library/transpose}
\noindent{\small transpose.scm -- the TRANSPOSE of a matrix (batch 32, 2026-09-25; the user's decision of 2026-09-25, on batch 25-C's proposal).}\par
\par\medskip\noindent\textbf{Definition} (\texttt{transpose}).\;$\operatorname{transpose}(a, m, n) :=$
\begin{equation*}
\begin{array}{@{}l@{}} \operatorname{matof}(n, m, \\ \quad (\lambda \langle \operatorname{tri\_}, \operatorname{trj\_} \rangle \in (\operatorname{interval}(1, n) \times \operatorname{interval}(1, m)).\; \operatorname{entry}(a, \operatorname{trj\_}, \operatorname{tri\_}))) \end{array}
\end{equation*}


\vnbbookfile{structure-library/elementary-matrix}{structure-library/elementary-matrix}
\noindent{\small elementary-matrix.scm -- matrix units and elementary column operations (algebraic-numbers.pdf ch.3, Def 3.2-3.7).  The whole elementary-operation theory is over a COMMUTATIVE ring here; the euclidean assumption only enters later, at the reduction-to-diagonal (Smith) step, NOT here.}\par
\par\medskip\noindent\textbf{Definition} (\texttt{matunit}).\;$\operatorname{matunit}(a, n, k, l) :=$
\begin{equation*}
\begin{array}{@{}l@{}} \operatorname{matof}(n, n, \\ \quad (\lambda \langle i, j \rangle \in (\operatorname{interval}(1, n) \times \operatorname{interval}(1, n)).\; \left(\text{if }((i = k) \wedge (j = l))\text{ then }\operatorname{one}(a)\text{ else }\operatorname{zero}(a)\right))) \end{array}
\end{equation*}

\par\medskip\noindent\textbf{Definition} (\texttt{elem-f}).\;$\operatorname{elem\text{-}f}(a, n, k, l) :=$
\begin{equation*}
\begin{array}{@{}l@{}} \operatorname{matof}(n, n, \\ \quad (\lambda \langle i, j \rangle \in (\operatorname{interval}(1, n) \times \operatorname{interval}(1, n)).\; \\ \quad \quad \left(\text{if }(((i = k) \wedge (j = l)) \vee (((i = l) \wedge (j = k)) \vee ((i = j) \wedge (\neg (i = k) \wedge \neg (i = l)))))\text{ then }\operatorname{one}(a)\text{ else }\operatorname{zero}(a)\right))) \end{array}
\end{equation*}

\par\medskip\noindent\textbf{Definition} (\texttt{elem-g}).\;$\operatorname{elem\text{-}g}(a, n, r, k, l) :=$
\begin{equation*}
\begin{array}{@{}l@{}} \operatorname{matof}(n, n, \\ \quad (\lambda \langle i, j \rangle \in (\operatorname{interval}(1, n) \times \operatorname{interval}(1, n)).\; \\ \quad \quad \left(\text{if }(i = j)\text{ then }\operatorname{one}(a)\text{ else }\left(\text{if }((i = k) \wedge (j = l))\text{ then }r\text{ else }\operatorname{zero}(a)\right)\right))) \end{array}
\end{equation*}

\par\medskip\noindent\textbf{Definition} (\texttt{elem-h}).\;$\operatorname{elem\text{-}h}(a, n, r, k) :=$
\begin{equation*}
\begin{array}{@{}l@{}} \operatorname{matof}(n, n, \\ \quad (\lambda \langle i, j \rangle \in (\operatorname{interval}(1, n) \times \operatorname{interval}(1, n)).\; \\ \quad \quad \left(\text{if }(i = j)\text{ then }\left(\text{if }(i = k)\text{ then }r\text{ else }\operatorname{one}(a)\right)\text{ else }\operatorname{zero}(a)\right))) \end{array}
\end{equation*}


\vnbbookfile{structure-library/mat-equiv}{structure-library/mat-equiv}
\noindent{\small RETIRED 2026-09-17 (proven): ring-neg-neg -- theorem-library/rake-algebra2.scm RETIRED 2026-09-17 (proven): submat-type, border-type -- theorem-library/rake-mat-typing.scm RETIRED 2026-09-17 (proven): diagonal-off-entry -- theorem-library/rake-algebra.scm mat-equiv.scm -- the matrix EQUIVALENCE relation \textasciitilde{} ...}\par
\par\medskip\noindent\textbf{Definition} (\texttt{is-invertible-mat}).\;For all $a$, $n$ and $u$, 
\begin{equation*}
\begin{array}{@{}l@{}} \operatorname{is\text{-}invertible\text{-}mat}(a, n, u) \\ \Leftrightarrow\; u \in \operatorname{mat}(n, n, \operatorname{carr}(a)) \wedge \\ \exists v \in \operatorname{mat}(n, n, \operatorname{carr}(a)). \\ \quad (\operatorname{matmul}(a, u, v) = \operatorname{identmat}(a, n)) \wedge (\operatorname{matmul}(a, v, u) = \operatorname{identmat}(a, n)) \end{array}
\end{equation*}

\par\medskip\noindent\textbf{Definition} (\texttt{mat-equiv}).\;For all $a$, $m$, $n$, $c$ and $d$, 
\begin{equation*}
\begin{array}{@{}l@{}} \operatorname{mat\text{-}equiv}(a, m, n, c, d) \\ \Leftrightarrow\; \exists u. \\ \quad \operatorname{is\text{-}invertible\text{-}mat}(a, m, u) \wedge \\ \quad \exists v.\; (\operatorname{is\text{-}invertible\text{-}mat}(a, n, v) \wedge (d = \operatorname{matmul}(a, \operatorname{matmul}(a, u, c), v))) \end{array}
\end{equation*}

\par\medskip\noindent\textbf{Definition} (\texttt{submat}).\;$\operatorname{submat}(s, p, q) :=$
\begin{equation*}
\begin{array}{@{}l@{}} \operatorname{matof}(p, q, \\ \quad (\lambda \langle i, j \rangle \in (\operatorname{interval}(1, p) \times \operatorname{interval}(1, q)).\; \operatorname{entry}(s, (i + 1), (j + 1)))) \end{array}
\end{equation*}

\par\medskip\noindent\textbf{Definition} (\texttt{border}).\;$\operatorname{border}(a, b, m, p, q) :=$
\begin{equation*}
\begin{array}{@{}l@{}} \operatorname{matof}((p + 1), (q + 1), \\ \quad (\lambda \langle i, j \rangle \in (\operatorname{interval}(1, (p + 1)) \times \operatorname{interval}(1, (q + 1))).\; \\ \quad \quad \left(\text{if }(i = 1)\text{ then }\left(\text{if }(j = 1)\text{ then }b\text{ else }\operatorname{zero}(a)\right)\text{ else }\left(\text{if }(j = 1)\text{ then }\operatorname{zero}(a)\text{ else }\operatorname{entry}(m, \operatorname{nn\text{-}minus}(i, 1), \operatorname{nn\text{-}minus}(j, 1))\right)\right))) \end{array}
\end{equation*}

\par\medskip\noindent\textbf{Definition} (\texttt{is-diagonal}).\;For all $a$, $m$, $n$ and $d$, 
\begin{equation*}
\begin{array}{@{}l@{}} \operatorname{is\text{-}diagonal}(a, m, n, d) \\ \Leftrightarrow\; \forall i \in \operatorname{interval}(1, m), j \in \operatorname{interval}(1, n). \\ \quad \neg (i = j) \Rightarrow (\operatorname{entry}(d, i, j) = \operatorname{zero}(a)) \end{array}
\end{equation*}

\par\medskip\noindent\textbf{Definition} (\texttt{smith-staircase}).\;For all $a$, $m$, $n$, $d$ and $k$, 
\begin{equation*}
\begin{array}{@{}l@{}} \operatorname{smith\text{-}staircase}(a, m, n, d, k) \\ \Leftrightarrow\; \operatorname{is\text{-}diagonal}(a, m, n, d) \wedge \\ k \leq m \wedge \\ k \leq n \wedge \\ \forall \operatorname{i\_} \in \operatorname{interval}(1, k).\; \neg (\operatorname{entry}(d, \operatorname{i\_}, \operatorname{i\_}) = \operatorname{zero}(a)) \wedge \\ \forall \operatorname{i\_} \in \operatorname{interval}(1, m). \\ \quad \neg (\operatorname{i\_} \leq k) \Rightarrow \forall \operatorname{j\_} \in \operatorname{interval}(1, n).\; (\operatorname{entry}(d, \operatorname{i\_}, \operatorname{j\_}) = \operatorname{zero}(a)) \end{array}
\end{equation*}


\vnbbookfile{structure-library/determinant}{structure-library/determinant}
\noindent{\small RETIRED 2026-09-17 (proven): minor-type (modulo 0), det-in-carrier (modulo mpow-type) -- theorem-library/rake-mat-typing.scm, by det-in-carrier-ind with the NN binder outermost. determinant.scm -- the determinant of a square matrix over a ring, by RECURSIVE COFACTOR (Laplace) EXPANSION along the first row.  The ...}\par
\par\medskip\noindent\textbf{Definition} (\texttt{det}).\;For all $r$ and $a$,  $(\operatorname{det}(r, 0, a) \simeq \operatorname{one}(r))$\par
Suppose $r \text{ s.t. } \operatorname{is\text{-}ring}(r)$. Let $n \in \mathbb{N}$ and $a \in \operatorname{mat}(\operatorname{succ}(n), \operatorname{succ}(n), \operatorname{carr}(r))$. Then
\begin{equation*}
(\operatorname{det}(r, \operatorname{succ}(n), a) \simeq \operatorname{finsum}(\operatorname{ring\text{-}additive\text{-}ag}(r), (\lambda j \in \operatorname{interval}(1, \operatorname{succ}(n)).\; \operatorname{mul}(r)(\operatorname{mpow}(\operatorname{ring\text{-}multiplicative\text{-}monoid}(r), \operatorname{neg}(r)(\operatorname{one}(r)), \operatorname{succ}(j)), \operatorname{mul}(r)(\operatorname{entry}(a, 1, j), \operatorname{det}(r, n, \operatorname{minor}(a, 1, j, n))))), \operatorname{interval}(1, \operatorname{succ}(n))))
\end{equation*}

\par\medskip\noindent\textbf{Definition} (\texttt{minor}).\;$\operatorname{minor}(s, r, c, n) :=$
\begin{equation*}
\begin{array}{@{}l@{}} \operatorname{matof}(n, n, \\ \quad (\lambda \langle i, j \rangle \in (\operatorname{interval}(1, n) \times \operatorname{interval}(1, n)).\; \\ \quad \quad \operatorname{entry}(s, \left(\text{if }(i < r)\text{ then }i\text{ else }(i + 1)\right), \left(\text{if }(j < c)\text{ then }j\text{ else }(j + 1)\right)))) \end{array}
\end{equation*}


\vnbbookfile{structure-library/mod-seq}{structure-library/mod-seq}
\noindent{\small RETIRED 2026-09-17 (proven): finsum-act-collect-gen -- theorem-library/rake-finsum-core.scm RETIRED 2026-09-17 (proven): finsum-act-distrib-gen -- theorem-library/rake-finsum-laws2.scm RETIRED 2026-09-17 (proven): matact-entry -- theorem-library/rake-identmat.scm (the module twin of matmul-entry). RETIRED 2026-09-17 ...}\par
\par\medskip\noindent\textbf{Definition} (\texttt{matact}).\;$\operatorname{matact}(\mathit{md}, p, u) :=$
\begin{equation*}
\begin{array}{@{}l@{}} \operatorname{matof}(\operatorname{nth}(1, \operatorname{size}(p)), \operatorname{nth}(2, \operatorname{size}(u)), \\ \quad (\lambda \langle i, c \rangle \in (\operatorname{interval}(1, \operatorname{nth}(1, \operatorname{size}(p))) \times \operatorname{interval}(1, \operatorname{nth}(2, \operatorname{size}(u)))).\; \\ \quad \quad \operatorname{finsum}(\operatorname{module\text{-}vector\text{-}ag}(\mathit{md}), \\ \quad \quad \quad (\lambda j \in \operatorname{interval}(1, \operatorname{nth}(2, \operatorname{size}(p))).\; \operatorname{act}(\mathit{md})(\operatorname{entry}(p, i, j), \operatorname{entry}(u, j, c))), \\ \quad \quad \quad \operatorname{interval}(1, \operatorname{nth}(2, \operatorname{size}(p)))))) \end{array}
\end{equation*}

\par\medskip\noindent\textbf{Definition} (\texttt{lincomb}).\;$\operatorname{lincomb}(\mathit{md}, n, c, u) :=$
\begin{equation*}
\begin{array}{@{}l@{}} \operatorname{finsum}(\operatorname{module\text{-}vector\text{-}ag}(\mathit{md}), \\ \quad (\lambda j \in \operatorname{interval}(1, n).\; \operatorname{act}(\mathit{md})(\operatorname{entry}(c, 1, j), \operatorname{entry}(u, j, 1))), \\ \quad \operatorname{interval}(1, n)) \end{array}
\end{equation*}

\par\medskip\noindent\textbf{Definition} (\texttt{generates}).\;For all $\mathit{md}$, $n$ and $u$, 
\begin{equation*}
\begin{array}{@{}l@{}} \operatorname{generates}(\mathit{md}, n, u) \\ \Leftrightarrow\; \forall \operatorname{x\_} \in \operatorname{vec}(\mathit{md}). \exists \operatorname{c\_} \in \operatorname{mat}(1, n, \operatorname{carr}(\operatorname{scal}(\mathit{md}))). \\ \quad \operatorname{x\_} = \operatorname{lincomb}(\mathit{md}, n, \operatorname{c\_}, u) \end{array}
\end{equation*}

\par\medskip\noindent\textbf{Definition} (\texttt{rel-free}).\;For all $\mathit{md}$, $n$ and $u$, 
\begin{equation*}
\begin{array}{@{}l@{}} \operatorname{rel\text{-}free}(\mathit{md}, n, u) \\ \Leftrightarrow\; \forall \operatorname{c\_} \in \operatorname{mat}(1, n, \operatorname{carr}(\operatorname{scal}(\mathit{md}))). \\ \quad (\operatorname{lincomb}(\mathit{md}, n, \operatorname{c\_}, u) = \operatorname{vzero}(\mathit{md})) \Rightarrow  \\ \quad \quad \forall \operatorname{j\_} \in \operatorname{interval}(1, n).\; (\operatorname{entry}(\operatorname{c\_}, 1, \operatorname{j\_}) = \operatorname{zero}(\operatorname{scal}(\mathit{md}))) \end{array}
\end{equation*}

\par\medskip\noindent\textbf{Definition} (\texttt{spans}).\;For all $\mathit{md}$, $n$, $u$ and $\mathit{sm}$, 
\begin{equation*}
\begin{array}{@{}l@{}} \operatorname{spans}(\mathit{md}, n, u, \mathit{sm}) \\ \Leftrightarrow\; \forall \operatorname{j\_} \in \operatorname{interval}(1, n).\; (\operatorname{entry}(u, \operatorname{j\_}, 1) \in \mathit{sm}) \wedge \\ \forall \operatorname{x\_} \in \mathit{sm}. \exists \operatorname{c\_} \in \operatorname{mat}(1, n, \operatorname{carr}(\operatorname{scal}(\mathit{md}))). \\ \quad \operatorname{x\_} = \operatorname{lincomb}(\mathit{md}, n, \operatorname{c\_}, u) \end{array}
\end{equation*}

\par\medskip\noindent\textbf{Definition} (\texttt{span}).\;$\operatorname{span}(\mathit{md}, n, u) :=$
\begin{equation*}
\{\operatorname{x\_} \in \operatorname{vec}(\mathit{md}) : \exists \operatorname{c\_} \in \operatorname{mat}(1, n, \operatorname{carr}(\operatorname{scal}(\mathit{md}))).\; (\operatorname{x\_} = \operatorname{lincomb}(\mathit{md}, n, \operatorname{c\_}, u))\}
\end{equation*}

\par\medskip\noindent\textbf{Definition} (\texttt{lastcoeff-set}).\;$\operatorname{lastcoeff\text{-}set}(\mathit{md}, p, u, \mathit{sm}) :=$
\begin{equation*}
\{\operatorname{r\_} \in \operatorname{carr}(\operatorname{scal}(\mathit{md})) : \exists \operatorname{c\_} \in \operatorname{mat}(1, (p + 1), \operatorname{carr}(\operatorname{scal}(\mathit{md}))).\; ((\operatorname{entry}(\operatorname{c\_}, 1, (p + 1)) = \operatorname{r\_}) \wedge (\operatorname{lincomb}(\mathit{md}, (p + 1), \operatorname{c\_}, u) \in \mathit{sm}))\}
\end{equation*}

\vnbdef{span-membership}{\vnbstmt{\begin{array}{@{}l@{}} \forall d, n, u, x. \\ \quad x \in \operatorname{span}(d, n, u) \\ \quad \Leftrightarrow\; ((x \in \operatorname{vec}(d)) \wedge \exists c \in \operatorname{mat}(1, n, \operatorname{carr}(\operatorname{scal}(d))).\; (x = \operatorname{lincomb}(d, n, c, u))) \end{array}}}
\vnbdef{lastcoeff-set-membership}{\vnbstmt{\begin{array}{@{}l@{}} \forall d, p, u, m, r. \\ \quad r \in \operatorname{lastcoeff\text{-}set}(d, p, u, m) \\ \quad \Leftrightarrow\; r \in \operatorname{carr}(\operatorname{scal}(d)) \wedge \\ \quad \exists c \in \operatorname{mat}(1, (p + 1), \operatorname{carr}(\operatorname{scal}(d))). \\ \quad \quad (\operatorname{entry}(c, 1, (p + 1)) = r) \wedge (\operatorname{lincomb}(d, (p + 1), c, u) \in m) \end{array}}}

\vnbbookfile{theorem-library/ag-view-read-offs}{theorem-library/ag-view-read-offs}
\noindent{\small A file shared across books; the entries that belong to this one.}\par
\noindent{\small ag-view-read-offs.scm -- THE ADDITIVE-GROUP VIEW READ-OFFS, PROVEN.}\par
\vnbprop{mvag-carr}{\vnbstmt{\forall d.\; (\operatorname{is\text{-}module}(d) \Rightarrow (\operatorname{carr}(\operatorname{module\text{-}vector\text{-}ag}(d)) = \operatorname{vec}(d)))}}
\vnbprop{mvag-op}{\vnbstmt{\forall d.\; (\operatorname{is\text{-}module}(d) \Rightarrow (\operatorname{opr}(\operatorname{module\text{-}vector\text{-}ag}(d)) = \operatorname{vadd}(d)))}}
\vnbprop{mvag-id}{\vnbstmt{\forall d.\; (\operatorname{is\text{-}module}(d) \Rightarrow (\operatorname{iden}(\operatorname{module\text{-}vector\text{-}ag}(d)) = \operatorname{vzero}(d)))}}

\vnbbookfile{theorem-library/rake-setoid}{theorem-library/rake-setoid}
\noindent{\small A file shared across books; the entries that belong to this one.}\par
\noindent{\small rake-setoid.scm -- BATCH F of the 2026-09-17 rake: setoids and carriers. Seven asserted leaves of the PSS, PROVEN.  Every statement here is its support's statement UNCHANGED.}\par
\vnbprop{is-vector-space-unfold}{\vnbstmt{\forall m.\; (\operatorname{is\text{-}vector\text{-}space}(m) \Leftrightarrow (\operatorname{is\text{-}module}(m) \wedge \operatorname{is\text{-}field\text{-}ring}(\operatorname{scal}(m))))}}
\vnbprop{vspace-vec-is-set}{\vnbstmt{\forall m.\; (\operatorname{is\text{-}vector\text{-}space}(m) \Rightarrow (\operatorname{vec}(m) \in \mathbf{Set}))}}

\vnbbookfile{theorem-library/entry-in-carrier}{theorem-library/entry-in-carrier}
\noindent{\small entry-in-carrier.scm -- the entries of an m-by-n matrix over X lie in X.}\par
\vnbprop{entry-in-carrier}{\vnbstmt{\begin{array}{@{}l@{}} \forall m, n, x, p, i, j. \\ \quad ((p \in \operatorname{mat}(m, n, x)) \wedge ((i \in \operatorname{interval}(1, m)) \wedge (j \in \operatorname{interval}(1, n)))) \Rightarrow  \\ \quad \quad \operatorname{entry}(p, i, j) \in x \end{array}}}

\vnbbookfile{theorem-library/rake-algebra}{theorem-library/rake-algebra}
\noindent{\small A file shared across books; the entries that belong to this one.}\par
\noindent{\small rake-algebra.scm -- BATCH D of the 2026-09-17 rake: seven asserted leaves of the PSS, PROVEN.  Every statement here is its support's statement UNCHANGED.}\par
\vnbprop{diagonal-off-entry}{\vnbstmt{\begin{array}{@{}l@{}} \forall a, m, n, d.; \forall i \in \operatorname{interval}(1, m), j \in \operatorname{interval}(1, n). \\ \quad (\operatorname{is\text{-}diagonal}(a, m, n, d) \wedge \neg (i = j)) \Rightarrow (\operatorname{entry}(d, i, j) = \operatorname{zero}(a)) \end{array}}}
\vnbprop{span-add-one-unfold}{\vnbstmt{\begin{array}{@{}l@{}} \forall m, s, v. \\ \quad (\operatorname{span\text{-}add\text{-}one}(m, s, v) \simeq \{y \in \operatorname{vec}(m) : \exists x \in s, r \in \mathbb{R}.\; (y = \operatorname{vadd}(m)(x, \operatorname{act}(m)(r, v)))\}) \end{array}}}
\vnbprop{span-add-one-membership}{\vnbstmt{\begin{array}{@{}l@{}} \forall m, s, v, w. \\ \quad w \in \operatorname{span\text{-}add\text{-}one}(m, s, v) \\ \quad \Leftrightarrow\; (w \in \operatorname{vec}(m)) \wedge \exists y \in s, r \in \mathbb{R}.\; (w = \operatorname{vadd}(m)(y, \operatorname{act}(m)(r, v))) \end{array}}}

\vnbbookfile{theorem-library/rake-finsum-typing}{theorem-library/rake-finsum-typing}
\noindent{\small A file shared across books; the entries that belong to this one.}\par
\noindent{\small rake-finsum-typing.scm -- the finsum / monoid-power typing leaves of the rake (batch G), PROVEN, together with the guarded form of the one statement in the batch that is NOT provable as stated.}\par
\vnbprop{submodule-finsum-closed}{\vnbstmt{\begin{array}{@{}l@{}} \forall d, m.; \forall s \in \mathbf{Set}, f \in (s \to \operatorname{carr}(\operatorname{module\text{-}vector\text{-}ag}(d))). \\ \quad (\operatorname{is\text{-}module}(d) \wedge (\operatorname{is\text{-}submodule}(d, m) \wedge ((|s| \in \mathbb{N}) \wedge \forall z \in s.\; (f(z) \in m)))) \Rightarrow  \\ \quad \quad \operatorname{finsum}(\operatorname{module\text{-}vector\text{-}ag}(d), f, s) \in m \end{array}}}

\vnbbookfile{theorem-library/hom-laws}{theorem-library/hom-laws}
\noindent{\small A file shared across books; the entries that belong to this one.}\par
\noindent{\small hom-laws.scm -- the models of a structure form a CATEGORY: identity arrows, composition, and the hom-set as a set.}\par
\vnbprop{hom-module-member-iff}{\vnbstmt{\begin{array}{@{}l@{}} \forall a, b, f. \\ \quad (f \in \operatorname{hom\text{-}module}(a, b)) \Leftrightarrow ((f \in (\operatorname{vec}(a) \to \operatorname{vec}(b))) \wedge \operatorname{is\text{-}hom\text{-}module}(a, b, f)) \end{array}}}
\vnbprop{hom-module-id}{\vnbstmt{\forall a.\; (\operatorname{is\text{-}module}(a) \Rightarrow \operatorname{is\text{-}hom\text{-}module}(a, a, \operatorname{id\text{-}fun}(\operatorname{vec}(a))))}}
\vnbprop{hom-module-in-set}{\vnbstmt{\forall a, b.\; ((\operatorname{is\text{-}module}(a) \wedge \operatorname{is\text{-}module}(b)) \Rightarrow (\operatorname{hom\text{-}module}(a, b) \in \mathbf{Set}))}}
\vnbprop{hom-vector-space-member-iff}{\vnbstmt{\begin{array}{@{}l@{}} \forall a, b, f. \\ \quad (f \in \operatorname{hom\text{-}vector\text{-}space}(a, b)) \Leftrightarrow ((f \in (\operatorname{vec}(a) \to \operatorname{vec}(b))) \wedge \operatorname{is\text{-}hom\text{-}vector\text{-}space}(a, b, f)) \end{array}}}
\vnbprop{hom-vector-space-id}{\vnbstmt{\begin{array}{@{}l@{}} \forall a. \\ \quad \operatorname{is\text{-}vector\text{-}space}(a) \Rightarrow \operatorname{is\text{-}hom\text{-}vector\text{-}space}(a, a, \operatorname{id\text{-}fun}(\operatorname{vec}(a))) \end{array}}}
\vnbprop{hom-vector-space-in-set}{\vnbstmt{\begin{array}{@{}l@{}} \forall a, b. \\ \quad (\operatorname{is\text{-}vector\text{-}space}(a) \wedge \operatorname{is\text{-}vector\text{-}space}(b)) \Rightarrow (\operatorname{hom\text{-}vector\text{-}space}(a, b) \in \mathbf{Set}) \end{array}}}

\vnbbookfile{theorem-library/hom-functors}{theorem-library/hom-functors}
\noindent{\small A file shared across books; the entries that belong to this one.}\par
\noindent{\small hom-functors.scm -- the Hom functors: Hom(a, -) covariant, Hom(-, b) contravariant, as maps on hom-sets.}\par
\vnbprop{hom-module-post-type}{\vnbstmt{\begin{array}{@{}l@{}} \forall a, b, c, h. \\ \quad (\operatorname{is\text{-}module}(a) \wedge (h \in \operatorname{hom\text{-}module}(b, c))) \Rightarrow  \\ \quad \quad (\lambda f \in \operatorname{hom\text{-}module}(a, b).\; \operatorname{compose}(h, f)) \in (\operatorname{hom\text{-}module}(a, b) \to \operatorname{hom\text{-}module}(a, c)) \end{array}}}
\vnbprop{hom-module-pre-type}{\vnbstmt{\begin{array}{@{}l@{}} \forall a, b, c, k. \\ \quad (\operatorname{is\text{-}module}(b) \wedge (k \in \operatorname{hom\text{-}module}(c, a))) \Rightarrow  \\ \quad \quad (\lambda f \in \operatorname{hom\text{-}module}(a, b).\; \operatorname{compose}(f, k)) \in (\operatorname{hom\text{-}module}(a, b) \to \operatorname{hom\text{-}module}(c, b)) \end{array}}}
\vnbprop{hom-vector-space-post-type}{\vnbstmt{\begin{array}{@{}l@{}} \forall a, b, c, h. \\ \quad (\operatorname{is\text{-}vector\text{-}space}(a) \wedge (h \in \operatorname{hom\text{-}vector\text{-}space}(b, c))) \Rightarrow  \\ \quad \quad (\lambda f \in \operatorname{hom\text{-}vector\text{-}space}(a, b).\; \operatorname{compose}(h, f)) \\ \quad \quad \in\; (\operatorname{hom\text{-}vector\text{-}space}(a, b) \to \operatorname{hom\text{-}vector\text{-}space}(a, c)) \end{array}}}
\vnbprop{hom-vector-space-pre-type}{\vnbstmt{\begin{array}{@{}l@{}} \forall a, b, c, k. \\ \quad (\operatorname{is\text{-}vector\text{-}space}(b) \wedge (k \in \operatorname{hom\text{-}vector\text{-}space}(c, a))) \Rightarrow  \\ \quad \quad (\lambda f \in \operatorname{hom\text{-}vector\text{-}space}(a, b).\; \operatorname{compose}(f, k)) \\ \quad \quad \in\; (\operatorname{hom\text{-}vector\text{-}space}(a, b) \to \operatorname{hom\text{-}vector\text{-}space}(c, b)) \end{array}}}

\vnbbookfile{theorem-library/qq-vs-exemplification}{theorem-library/qq-vs-exemplification}
\noindent{\small qq-vs-exemplification.scm -- THE FIRST WITNESS OF A VECTOR SPACE.}\par
\par\medskip\noindent\textbf{Definition} (\texttt{qq-line}).\;
\begin{equation*}
\begin{array}{@{}l@{}} \operatorname{qq\text{-}line} \\ =\; \langle \operatorname{field\text{-}as\text{-}integral\text{-}domain}(\operatorname{qq\text{-}field}), \mathbb{Q}, (\lambda \langle \operatorname{x\_}, \operatorname{y\_} \rangle \in (\mathbb{Q} \times \mathbb{Q}).\; +(\operatorname{x\_}, \operatorname{y\_})), 0, (\lambda \operatorname{x\_} \in \mathbb{Q}.\; -(\operatorname{x\_})), (\lambda \langle \operatorname{r\_}, \operatorname{x\_} \rangle \in (\mathbb{Q} \times \mathbb{Q}).\; \cdot(\operatorname{r\_}, \operatorname{x\_})) \rangle \end{array}
\end{equation*}

\vnbdef{qq-line@vzero}{\vnbstmt{(\operatorname{vzero}(\operatorname{qq\text{-}line}) \simeq 0)}}
\vnbdef{qq-line@scal}{\vnbstmt{(\operatorname{scal}(\operatorname{qq\text{-}line}) \simeq \operatorname{field\text{-}as\text{-}integral\text{-}domain}(\operatorname{qq\text{-}field}))}}
\vnbdef{qq-line@vec}{\vnbstmt{(\operatorname{vec}(\operatorname{qq\text{-}line}) \simeq \mathbb{Q})}}
\vnbdef{qq-line@vadd}{\vnbstmt{(\operatorname{vadd}(\operatorname{qq\text{-}line}) \simeq (\lambda \langle \operatorname{x\_}, \operatorname{y\_} \rangle \in (\mathbb{Q} \times \mathbb{Q}).\; (\operatorname{x\_} + \operatorname{y\_})))}}
\vnbdef{qq-line@act}{\vnbstmt{(\operatorname{act}(\operatorname{qq\text{-}line}) \simeq (\lambda \langle \operatorname{r\_}, \operatorname{x\_} \rangle \in (\mathbb{Q} \times \mathbb{Q}).\; (\operatorname{r\_} \cdot \operatorname{x\_})))}}
\vnbdef{qq-line@vneg}{\vnbstmt{(\operatorname{vneg}(\operatorname{qq\text{-}line}) \simeq (\lambda x \in \mathbb{Q}.\; -x))}}
\vnbprop{qq-line-is-module}{\vnbstmt{\operatorname{is\text{-}module}(\operatorname{qq\text{-}line})}}
\vnbprop{qq-line-is-vector-space}{\vnbstmt{\operatorname{is\text{-}vector\text{-}space}(\operatorname{qq\text{-}line})}}

\vnbbookfile{theorem-library/rake-border-entry}{theorem-library/rake-border-entry}
\noindent{\small rake-border-entry.scm -- the BORDER block entry read-offs, PROVEN.}\par
\vnbprop{border-entries-defined}{\vnbstmt{\begin{array}{@{}l@{}} \forall a, b, m, p, q.; \forall c \in \operatorname{interval}(1, (p + 1)), d \in \operatorname{interval}(1, (q + 1)). \\ \quad (\operatorname{is\text{-}ring}(a) \wedge ((p \in \mathbb{N}) \wedge ((q \in \mathbb{N}) \wedge ((b \in \operatorname{carr}(a)) \wedge (m \in \operatorname{mat}(p, q, \operatorname{carr}(a))))))) \Rightarrow  \\ \quad \quad (\lambda \langle i, j \rangle \in (\operatorname{interval}(1, (p + 1)) \times \operatorname{interval}(1, (q + 1))).\; \left(\text{if }(i = 1)\text{ then }\left(\text{if }(j = 1)\text{ then }b\text{ else }\operatorname{zero}(a)\right)\text{ else }\left(\text{if }(j = 1)\text{ then }\operatorname{zero}(a)\text{ else }\operatorname{entry}(m, \operatorname{nn\text{-}minus}(i, 1), \operatorname{nn\text{-}minus}(j, 1))\right)\right))(c, d) \\ \quad \quad \in\; \mathbf{Set} \end{array}}}
\vnbprop{border-entry-block2}{\vnbstmt{\begin{array}{@{}l@{}} \forall a, b, m, p, q, i, j. \\ \quad \operatorname{is\text{-}ring}(a) \wedge \\ \quad p \in \mathbb{N} \wedge \\ \quad q \in \mathbb{N} \wedge \\ \quad b \in \operatorname{carr}(a) \wedge \\ \quad m \in \operatorname{mat}(p, q, \operatorname{carr}(a)) \wedge \\ \quad i \in \operatorname{interval}(1, (p + 1)) \wedge \\ \quad \neg (i = 1) \wedge \\ \quad j \in \operatorname{interval}(1, (q + 1)) \wedge \\ \quad \neg (j = 1) \Rightarrow \\ \quad \quad \operatorname{entry}(\operatorname{border}(a, b, m, p, q), i, j) = \operatorname{entry}(m, \operatorname{nn\text{-}minus}(i, 1), \operatorname{nn\text{-}minus}(j, 1)) \end{array}}}
\vnbprop{border-entry-block}{\vnbstmt{\begin{array}{@{}l@{}} \forall a, b, m, p, q, i, j. \\ \quad \operatorname{is\text{-}ring}(a) \wedge \\ \quad p \in \mathbb{N} \wedge \\ \quad q \in \mathbb{N} \wedge \\ \quad b \in \operatorname{carr}(a) \wedge \\ \quad m \in \operatorname{mat}(p, q, \operatorname{carr}(a)) \wedge \\ \quad i \in \operatorname{interval}(1, p) \wedge \\ \quad j \in \operatorname{interval}(1, q) \Rightarrow \\ \quad \quad \operatorname{entry}(\operatorname{border}(a, b, m, p, q), (i + 1), (j + 1)) = \operatorname{entry}(m, i, j) \end{array}}}

\vnbbookfile{theorem-library/rake-border-siblings}{theorem-library/rake-border-siblings}
\noindent{\small rake-border-siblings.scm -- the three remaining BORDER entry read-offs, PROVEN: border-entry-11, border-entry-1j, border-entry-i1.}\par
\vnbprop{border-entry-11}{\vnbstmt{\begin{array}{@{}l@{}} \forall a, b, m, p, q. \\ \quad (\operatorname{is\text{-}ring}(a) \wedge ((p \in \mathbb{N}) \wedge ((q \in \mathbb{N}) \wedge ((b \in \operatorname{carr}(a)) \wedge (m \in \operatorname{mat}(p, q, \operatorname{carr}(a))))))) \Rightarrow  \\ \quad \quad \operatorname{entry}(\operatorname{border}(a, b, m, p, q), 1, 1) = b \end{array}}}
\vnbprop{border-entry-1j}{\vnbstmt{\begin{array}{@{}l@{}} \forall a, b, m, p, q, j. \\ \quad \operatorname{is\text{-}ring}(a) \wedge \\ \quad p \in \mathbb{N} \wedge \\ \quad q \in \mathbb{N} \wedge \\ \quad b \in \operatorname{carr}(a) \wedge \\ \quad m \in \operatorname{mat}(p, q, \operatorname{carr}(a)) \wedge \\ \quad j \in \operatorname{interval}(1, (q + 1)) \wedge \\ \quad \neg (j = 1) \Rightarrow \\ \quad \quad \operatorname{entry}(\operatorname{border}(a, b, m, p, q), 1, j) = \operatorname{zero}(a) \end{array}}}
\vnbprop{border-entry-i1}{\vnbstmt{\begin{array}{@{}l@{}} \forall a, b, m, p, q, i. \\ \quad \operatorname{is\text{-}ring}(a) \wedge \\ \quad p \in \mathbb{N} \wedge \\ \quad q \in \mathbb{N} \wedge \\ \quad b \in \operatorname{carr}(a) \wedge \\ \quad m \in \operatorname{mat}(p, q, \operatorname{carr}(a)) \wedge \\ \quad i \in \operatorname{interval}(1, (p + 1)) \wedge \\ \quad \neg (i = 1) \Rightarrow \\ \quad \quad \operatorname{entry}(\operatorname{border}(a, b, m, p, q), i, 1) = \operatorname{zero}(a) \end{array}}}
\vnbprop{succ-nn-minus-1}{\vnbstmt{\forall i \in \mathbb{N}.\; ((1 \leq i) \Rightarrow ((\operatorname{nn\text{-}minus}(i, 1) + 1) = i))}}

\vnbbookfile{theorem-library/rake-finsum-laws2}{theorem-library/rake-finsum-laws2}
\noindent{\small A file shared across books; the entries that belong to this one.}\par
\noindent{\small theorem-library/rake-finsum-laws2.scm -- rake batch J, part 2: the three DISTRIBUTION laws of FINSUM (a ring element on either side, a module scalar) and the interval back-peel.  Part 1 is theorem-library/rake-finsum-laws.scm, which must load FIRST: this file cites its `enum-fam-value' and}\par
\vnbprop{module-act-zero-vec}{\vnbstmt{\begin{array}{@{}l@{}} \forall d.; \forall r \in \operatorname{carr}(\operatorname{scal}(d)). \\ \quad \operatorname{is\text{-}module}(d) \Rightarrow (\operatorname{act}(d)(r, \operatorname{vzero}(d)) = \operatorname{vzero}(d)) \end{array}}}
\vnbprop{sum-ag-act-ind}{\vnbstmt{\begin{array}{@{}l@{}} \forall n \in \mathbb{N}, d.; \forall r \in \operatorname{carr}(\operatorname{scal}(d)), a, b. \\ \quad \operatorname{is\text{-}module}(d) \wedge \\ \quad \forall i \in \mathrm{OS}(n).\; (a(i) \in \operatorname{carr}(\operatorname{module\text{-}vector\text{-}ag}(d))) \wedge \\ \quad \forall i \in \mathrm{OS}(n).\; (b(i) \simeq \operatorname{act}(d)(r, a(i))) \Rightarrow \\ \quad \quad \operatorname{sum\text{-}ag}(\operatorname{module\text{-}vector\text{-}ag}(d), b, n) = \operatorname{act}(d)(r, \operatorname{sum\text{-}ag}(\operatorname{module\text{-}vector\text{-}ag}(d), a, n)) \end{array}}}
\vnbprop{finsum-act-distrib-gen}{\vnbstmt{\begin{array}{@{}l@{}} \forall d.; \forall r \in \operatorname{carr}(\operatorname{scal}(d)), s \in \mathbf{Set}, f \in (s \to \operatorname{carr}(\operatorname{module\text{-}vector\text{-}ag}(d))). \\ \quad (\operatorname{is\text{-}module}(d) \wedge (|s| \in \mathbb{N})) \Rightarrow  \\ \quad \quad \operatorname{act}(d)(r, \operatorname{finsum}(\operatorname{module\text{-}vector\text{-}ag}(d), f, s)) \\ \quad \quad =\; \operatorname{finsum}(\operatorname{module\text{-}vector\text{-}ag}(d), (\lambda z \in s.\; \operatorname{act}(d)(r, f(z))), s) \end{array}}}

\vnbbookfile{theorem-library/rake-finsum-core}{theorem-library/rake-finsum-core}
\noindent{\small A file shared across books; the entries that belong to this one.}\par
\noindent{\small theorem-library/rake-finsum-core.scm -- the FINSUM CORE, proven. Rake batch M (2026-09-17).  Nine of the batch's ten leaves; the tenth (finsum-embed) is BLOCKED and its obstacle is stated at the end of this header.}\par
\vnbprop{sum-ag-act-collect-ind}{\vnbstmt{\begin{array}{@{}l@{}} \forall n \in \mathbb{N}, d.; \forall x \in \operatorname{vec}(d), c, v. \\ \quad \operatorname{is\text{-}module}(d) \wedge \\ \quad \forall i \in \mathrm{OS}(n).\; (c(i) \in \operatorname{carr}(\operatorname{scal}(d))) \wedge \\ \quad \forall i \in \mathrm{OS}(n).\; (v(i) \simeq \operatorname{act}(d)(c(i), x)) \Rightarrow \\ \quad \quad \operatorname{sum\text{-}ag}(\operatorname{module\text{-}vector\text{-}ag}(d), v, n) = \operatorname{act}(d)(\operatorname{sum\text{-}ag}(\operatorname{ring\text{-}additive\text{-}ag}(\operatorname{scal}(d)), c, n), x) \end{array}}}
\vnbprop{finsum-act-collect-gen}{\vnbstmt{\begin{array}{@{}l@{}} \forall d.; \forall x \in \operatorname{vec}(d), s \in \mathbf{Set}, c \in (s \to \operatorname{carr}(\operatorname{scal}(d))). \\ \quad (\operatorname{is\text{-}module}(d) \wedge (|s| \in \mathbb{N})) \Rightarrow  \\ \quad \quad \operatorname{act}(d)(\operatorname{finsum}(\operatorname{ring\text{-}additive\text{-}ag}(\operatorname{scal}(d)), c, s), x) \\ \quad \quad =\; \operatorname{finsum}(\operatorname{module\text{-}vector\text{-}ag}(d), (\lambda z \in s.\; \operatorname{act}(d)(c(z), x)), s) \end{array}}}

\vnbbookfile{theorem-library/mat-typing-bundle}{theorem-library/mat-typing-bundle}
\noindent{\small mat-typing-bundle.scm -- three MATOF typings, PROVEN from the definitions.}\par
\vnbprop{mat-colcount-transfer}{\vnbstmt{\begin{array}{@{}l@{}} \forall n, k, y, q, m, x, r. \\ \quad ((q \in \operatorname{mat}(n, k, y)) \wedge (((n = 0) \Rightarrow ((m = 0) \vee (k = 0))) \wedge (r \in \operatorname{mat}(m, \operatorname{nth}(2, \operatorname{size}(q)), x)))) \Rightarrow  \\ \quad \quad r \in \operatorname{mat}(m, k, x) \end{array}}}
\vnbprop{identmat-type}{\vnbstmt{\forall a.; \forall n \in \mathbb{N}.\; (\operatorname{is\text{-}ring}(a) \Rightarrow (\operatorname{identmat}(a, n) \in \operatorname{mat}(n, n, \operatorname{carr}(a))))}}
\vnbprop{elem-g-type}{\vnbstmt{\begin{array}{@{}l@{}} \forall a, n, r, k, l. \\ \quad (\operatorname{is\text{-}ring}(a) \wedge ((n \in \mathbb{N}) \wedge (r \in \operatorname{carr}(a)))) \Rightarrow  \\ \quad \quad \operatorname{elem\text{-}g}(a, n, r, k, l) \in \operatorname{mat}(n, n, \operatorname{carr}(a)) \end{array}}}
\vnbprop{zeromat-type}{\vnbstmt{\begin{array}{@{}l@{}} \forall a, m, n. \\ \quad (\operatorname{is\text{-}ring}(a) \wedge ((m \in \mathbb{N}) \wedge (n \in \mathbb{N}))) \Rightarrow  \\ \quad \quad \operatorname{zeromat}(a, m, n) \in \operatorname{mat}(m, n, \operatorname{carr}(a)) \end{array}}}
\vnbprop{unitrow-type}{\vnbstmt{\begin{array}{@{}l@{}} \forall a, n, i. \\ \quad (\operatorname{is\text{-}ring}(a) \wedge (n \in \mathbb{N})) \Rightarrow (\operatorname{unitrow}(a, n, i) \in \operatorname{mat}(1, n, \operatorname{carr}(a))) \end{array}}}
\vnbprop{elem-h-type}{\vnbstmt{\begin{array}{@{}l@{}} \forall a, n, r, k. \\ \quad (\operatorname{is\text{-}ring}(a) \wedge ((n \in \mathbb{N}) \wedge (r \in \operatorname{carr}(a)))) \Rightarrow  \\ \quad \quad \operatorname{elem\text{-}h}(a, n, r, k) \in \operatorname{mat}(n, n, \operatorname{carr}(a)) \end{array}}}
\vnbprop{matmul-type}{\vnbstmt{\begin{array}{@{}l@{}} \forall a, m, n, k, p, q. \\ \quad \operatorname{is\text{-}ring}(a) \wedge \\ \quad p \in \operatorname{mat}(m, n, \operatorname{carr}(a)) \wedge \\ \quad q \in \operatorname{mat}(n, k, \operatorname{carr}(a)) \wedge \\ \quad (n = 0) \Rightarrow ((m = 0) \vee (k = 0)) \Rightarrow \\ \quad \quad \operatorname{matmul}(a, p, q) \in \operatorname{mat}(m, k, \operatorname{carr}(a)) \end{array}}}

\vnbbookfile{theorem-library/rake-mat-typing}{theorem-library/rake-mat-typing}
\noindent{\small rake-mat-typing.scm -- the ELEVEN matrix typing / sethood supports, PROVEN.}\par
\vnbprop{mat-is-set}{\vnbstmt{\forall x \in \mathbf{Set}, m, n.\; (\operatorname{mat}(m, n, x) \in \mathbf{Set})}}
\vnbprop{block-type}{\vnbstmt{\begin{array}{@{}l@{}} \forall m, n, x, p, k, l. \\ \quad ((k \in \mathbb{N}) \wedge ((l \in \mathbb{N}) \wedge ((p \in \operatorname{mat}(m, n, x)) \wedge ((k \leq m) \wedge (l \leq n))))) \Rightarrow  \\ \quad \quad \operatorname{block}(p, k, l) \in \operatorname{mat}(k, l, x) \end{array}}}
\vnbprop{snoc-col-type}{\vnbstmt{\begin{array}{@{}l@{}} \forall x, n, w, v. \\ \quad ((n \in \mathbb{N}) \wedge ((w \in \operatorname{mat}(n, 1, x)) \wedge (v \in x))) \Rightarrow  \\ \quad \quad \operatorname{snoc\text{-}col}(w, n, v) \in \operatorname{mat}((n + 1), 1, x) \end{array}}}
\vnbprop{snoc-row-type}{\vnbstmt{\begin{array}{@{}l@{}} \forall x, n, c, r. \\ \quad ((n \in \mathbb{N}) \wedge ((c \in \operatorname{mat}(1, n, x)) \wedge (r \in x))) \Rightarrow  \\ \quad \quad \operatorname{snoc\text{-}row}(c, n, r) \in \operatorname{mat}(1, (n + 1), x) \end{array}}}
\vnbprop{matadd-type}{\vnbstmt{\begin{array}{@{}l@{}} \forall a, m, n, p, q. \\ \quad (\operatorname{is\text{-}ring}(a) \wedge ((p \in \operatorname{mat}(m, n, \operatorname{carr}(a))) \wedge (q \in \operatorname{mat}(m, n, \operatorname{carr}(a))))) \Rightarrow  \\ \quad \quad \operatorname{matadd}(a, p, q) \in \operatorname{mat}(m, n, \operatorname{carr}(a)) \end{array}}}
\vnbprop{matscale-type}{\vnbstmt{\begin{array}{@{}l@{}} \forall a, m, n, r, p. \\ \quad (\operatorname{is\text{-}ring}(a) \wedge ((r \in \operatorname{carr}(a)) \wedge (p \in \operatorname{mat}(m, n, \operatorname{carr}(a))))) \Rightarrow  \\ \quad \quad \operatorname{matscale}(a, r, p) \in \operatorname{mat}(m, n, \operatorname{carr}(a)) \end{array}}}
\vnbprop{matneg-type}{\vnbstmt{\begin{array}{@{}l@{}} \forall a, m, n.; \forall p \in \operatorname{mat}(m, n, \operatorname{carr}(a)). \\ \quad \operatorname{is\text{-}ring}(a) \Rightarrow (\operatorname{matneg}(a, p) \in \operatorname{mat}(m, n, \operatorname{carr}(a))) \end{array}}}
\vnbprop{submat-type}{\vnbstmt{\begin{array}{@{}l@{}} \forall a, p, q, s. \\ \quad ((p \in \mathbb{N}) \wedge ((q \in \mathbb{N}) \wedge (s \in \operatorname{mat}((p + 1), (q + 1), \operatorname{carr}(a))))) \Rightarrow  \\ \quad \quad \operatorname{submat}(s, p, q) \in \operatorname{mat}(p, q, \operatorname{carr}(a)) \end{array}}}
\vnbprop{border-type}{\vnbstmt{\begin{array}{@{}l@{}} \forall a, b, m, p, q. \\ \quad (\operatorname{is\text{-}ring}(a) \wedge ((p \in \mathbb{N}) \wedge ((q \in \mathbb{N}) \wedge ((b \in \operatorname{carr}(a)) \wedge (m \in \operatorname{mat}(p, q, \operatorname{carr}(a))))))) \Rightarrow  \\ \quad \quad \operatorname{border}(a, b, m, p, q) \in \operatorname{mat}((p + 1), (q + 1), \operatorname{carr}(a)) \end{array}}}
\vnbprop{minor-type}{\vnbstmt{\begin{array}{@{}l@{}} \forall r, s, p, q, n. \\ \quad (\operatorname{is\text{-}ring}(r) \wedge ((p \in \mathbb{N}) \wedge ((q \in \mathbb{N}) \wedge ((n \in \mathbb{N}) \wedge (s \in \operatorname{mat}((n + 1), (n + 1), \operatorname{carr}(r))))))) \Rightarrow  \\ \quad \quad \operatorname{minor}(s, p, q, n) \in \operatorname{mat}(n, n, \operatorname{carr}(r)) \end{array}}}
\vnbprop{rmm-carr}{\vnbstmt{\forall r.\; (\operatorname{is\text{-}ring}(r) \Rightarrow (\operatorname{carr}(\operatorname{ring\text{-}multiplicative\text{-}monoid}(r)) = \operatorname{carr}(r)))}}
\vnbprop{det-in-carrier-ind}{\vnbstmt{\begin{array}{@{}l@{}} \forall r.; \forall n \in \mathbb{N}, a \in \operatorname{mat}(n, n, \operatorname{carr}(r)). \\ \quad \operatorname{is\text{-}ring}(r) \Rightarrow (\operatorname{det}(r, n, a) \in \operatorname{carr}(r)) \end{array}}}
\vnbprop{det-in-carrier}{\vnbstmt{\begin{array}{@{}l@{}} \forall r, n, a. \\ \quad (\operatorname{is\text{-}ring}(r) \wedge ((n \in \mathbb{N}) \wedge (a \in \operatorname{mat}(n, n, \operatorname{carr}(r))))) \Rightarrow  \\ \quad \quad \operatorname{det}(r, n, a) \in \operatorname{carr}(r) \end{array}}}

\vnbbookfile{theorem-library/rake-det-small}{theorem-library/rake-det-small}
\noindent{\small rake-det-small.scm -- the small determinants, PROVEN (rake batch 5b-C, 2026-09-18).}\par
\vnbprop{interval-1-1}{\vnbstmt{\operatorname{interval}(1, 1) = \{1\}}}
\vnbprop{rmm-op}{\vnbstmt{\forall r.\; (\operatorname{is\text{-}ring}(r) \Rightarrow (\operatorname{opr}(\operatorname{ring\text{-}multiplicative\text{-}monoid}(r)) = \operatorname{mul}(r)))}}
\vnbprop{rmm-iden}{\vnbstmt{\forall r.\; (\operatorname{is\text{-}ring}(r) \Rightarrow (\operatorname{iden}(\operatorname{ring\text{-}multiplicative\text{-}monoid}(r)) = \operatorname{one}(r)))}}
\vnbprop{ring-mpow-one}{\vnbstmt{\begin{array}{@{}l@{}} \forall r.; \forall x \in \operatorname{carr}(r). \\ \quad \operatorname{is\text{-}ring}(r) \Rightarrow (\operatorname{mpow}(\operatorname{ring\text{-}multiplicative\text{-}monoid}(r), x, 1) = x) \end{array}}}
\vnbprop{minor-entry}{\vnbstmt{\begin{array}{@{}l@{}} \forall r, s, p, q, n, u, v. \\ \quad \operatorname{is\text{-}ring}(r) \wedge \\ \quad n \in \mathbb{N} \wedge \\ \quad s \in \operatorname{mat}((n + 1), (n + 1), \operatorname{carr}(r)) \wedge \\ \quad u \in \operatorname{interval}(1, n) \wedge \\ \quad v \in \operatorname{interval}(1, n) \Rightarrow \\ \quad \quad \operatorname{entry}(\operatorname{minor}(s, p, q, n), u, v) = \operatorname{entry}(s, \left(\text{if }(u < p)\text{ then }u\text{ else }(u + 1)\right), \left(\text{if }(v < q)\text{ then }v\text{ else }(v + 1)\right)) \end{array}}}
\vnbprop{det-1x1}{\vnbstmt{\begin{array}{@{}l@{}} \forall r, a. \\ \quad (\operatorname{is\text{-}ring}(r) \wedge (a \in \operatorname{mat}(1, 1, \operatorname{carr}(r)))) \Rightarrow (\operatorname{det}(r, 1, a) = \operatorname{entry}(a, 1, 1)) \end{array}}}
\vnbprop{det-2x2}{\vnbstmt{\begin{array}{@{}l@{}} \forall r, a. \\ \quad (\operatorname{is\text{-}ring}(r) \wedge (a \in \operatorname{mat}(2, 2, \operatorname{carr}(r)))) \Rightarrow  \\ \quad \quad \operatorname{det}(r, 2, a) \\ \quad \quad =\; \operatorname{add}(r)(\operatorname{mul}(r)(\operatorname{entry}(a, 1, 1), \operatorname{entry}(a, 2, 2)), \\ \quad \quad \quad \operatorname{neg}(r)(\operatorname{mul}(r)(\operatorname{entry}(a, 1, 2), \operatorname{entry}(a, 2, 1)))) \end{array}}}
\vnbprop{identmat-minor}{\vnbstmt{\begin{array}{@{}l@{}} \forall r, n. \\ \quad (\operatorname{is\text{-}ring}(r) \wedge (n \in \mathbb{N})) \Rightarrow  \\ \quad \quad \operatorname{minor}(\operatorname{identmat}(r, (n + 1)), 1, 1, n) = \operatorname{identmat}(r, n) \end{array}}}
\vnbprop{mat-ring-one}{\vnbstmt{\begin{array}{@{}l@{}} \forall a, n. \\ \quad (\operatorname{is\text{-}ring}(a) \wedge (n \in \mathbb{N})) \Rightarrow (\operatorname{one}(\operatorname{mat\text{-}ring}(a, n)) = \operatorname{identmat}(a, n)) \end{array}}}
\vnbprop{det-identity-ind}{\vnbstmt{\forall r.; \forall n \in \mathbb{N}.\; (\operatorname{is\text{-}ring}(r) \Rightarrow (\operatorname{det}(r, n, \operatorname{identmat}(r, n)) = \operatorname{one}(r)))}}
\vnbprop{det-identity}{\vnbstmt{\begin{array}{@{}l@{}} \forall r, n. \\ \quad (\operatorname{is\text{-}ring}(r) \wedge (n \in \mathbb{N})) \Rightarrow (\operatorname{det}(r, n, \operatorname{one}(\operatorname{mat\text{-}ring}(r, n))) = \operatorname{one}(r)) \end{array}}}

\vnbbookfile{theorem-library/rake-mat-ring-readoffs}{theorem-library/rake-mat-ring-readoffs}
\noindent{\small rake-mat-ring-readoffs.scm -- the MAT-RING slot read-offs and the two curried-operation FUN typings, PROVEN (rake, 2026-09-19).}\par
\vnbprop{mat-ring-carr}{\vnbstmt{\begin{array}{@{}l@{}} \forall a, n. \\ \quad (\operatorname{is\text{-}ring}(a) \wedge (n \in \mathbb{N})) \Rightarrow (\operatorname{carr}(\operatorname{mat\text{-}ring}(a, n)) = \operatorname{mat}(n, n, \operatorname{carr}(a))) \end{array}}}
\vnbprop{mat-ring-add}{\vnbstmt{\begin{array}{@{}l@{}} \forall a, n. \\ \quad (\operatorname{is\text{-}ring}(a) \wedge (n \in \mathbb{N})) \Rightarrow  \\ \quad \quad \operatorname{add}(\operatorname{mat\text{-}ring}(a, n)) \\ \quad \quad =\; (\lambda \langle p, q \rangle \in (\operatorname{mat}(n, n, \operatorname{carr}(a)) \times \operatorname{mat}(n, n, \operatorname{carr}(a))).\; \operatorname{matadd}(a, p, q)) \end{array}}}
\vnbprop{mat-ring-mul}{\vnbstmt{\begin{array}{@{}l@{}} \forall a, n. \\ \quad (\operatorname{is\text{-}ring}(a) \wedge (n \in \mathbb{N})) \Rightarrow  \\ \quad \quad \operatorname{mul}(\operatorname{mat\text{-}ring}(a, n)) \\ \quad \quad =\; (\lambda \langle p, q \rangle \in (\operatorname{mat}(n, n, \operatorname{carr}(a)) \times \operatorname{mat}(n, n, \operatorname{carr}(a))).\; \operatorname{matmul}(a, p, q)) \end{array}}}
\vnbprop{mat-ring-neg}{\vnbstmt{\begin{array}{@{}l@{}} \forall a, n. \\ \quad (\operatorname{is\text{-}ring}(a) \wedge (n \in \mathbb{N})) \Rightarrow  \\ \quad \quad \operatorname{neg}(\operatorname{mat\text{-}ring}(a, n)) = (\lambda p \in \operatorname{mat}(n, n, \operatorname{carr}(a)).\; \operatorname{matneg}(a, p)) \end{array}}}
\vnbprop{mat-ring-zero}{\vnbstmt{\begin{array}{@{}l@{}} \forall a, n. \\ \quad (\operatorname{is\text{-}ring}(a) \wedge (n \in \mathbb{N})) \Rightarrow (\operatorname{zero}(\operatorname{mat\text{-}ring}(a, n)) = \operatorname{zeromat}(a, n, n)) \end{array}}}
\vnbprop{mat-ring-add-fun}{\vnbstmt{\begin{array}{@{}l@{}} \forall a.; \forall n \in \mathbb{N}. \\ \quad \operatorname{is\text{-}ring}(a) \Rightarrow  \\ \quad \quad (\lambda \langle p, q \rangle \in (\operatorname{mat}(n, n, \operatorname{carr}(a)) \times \operatorname{mat}(n, n, \operatorname{carr}(a))).\; \operatorname{matadd}(a, p, q)) \\ \quad \quad \in\; ((\operatorname{mat}(n, n, \operatorname{carr}(a)) \times \operatorname{mat}(n, n, \operatorname{carr}(a))) \to \operatorname{mat}(n, n, \operatorname{carr}(a))) \end{array}}}
\vnbprop{mat-ring-mul-fun}{\vnbstmt{\begin{array}{@{}l@{}} \forall a.; \forall n \in \mathbb{N}. \\ \quad \operatorname{is\text{-}ring}(a) \Rightarrow  \\ \quad \quad (\lambda \langle p, q \rangle \in (\operatorname{mat}(n, n, \operatorname{carr}(a)) \times \operatorname{mat}(n, n, \operatorname{carr}(a))).\; \operatorname{matmul}(a, p, q)) \\ \quad \quad \in\; ((\operatorname{mat}(n, n, \operatorname{carr}(a)) \times \operatorname{mat}(n, n, \operatorname{carr}(a))) \to \operatorname{mat}(n, n, \operatorname{carr}(a))) \end{array}}}

\vnbbookfile{theorem-library/matmul-entry-proof}{theorem-library/matmul-entry-proof}
\noindent{\small matmul-entry-proof.scm -- the matrix-product entry formula, PROVEN.}\par
\vnbprop{matmul-entry}{\vnbstmt{\begin{array}{@{}l@{}} \forall a, m, n, k, p, q.; \forall i \in \operatorname{interval}(1, m), c \in \operatorname{interval}(1, k). \\ \quad (\operatorname{is\text{-}ring}(a) \wedge ((p \in \operatorname{mat}(m, n, \operatorname{carr}(a))) \wedge ((q \in \operatorname{mat}(n, k, \operatorname{carr}(a))) \wedge (1 \leq n)))) \Rightarrow  \\ \quad \quad \operatorname{entry}(\operatorname{matmul}(a, p, q), i, c) \\ \quad \quad =\; \operatorname{finsum}(\operatorname{ring\text{-}additive\text{-}ag}(a), \\ \quad \quad \quad (\lambda j \in \operatorname{interval}(1, n).\; \operatorname{mul}(a)(\operatorname{entry}(p, i, j), \operatorname{entry}(q, j, c))), \\ \quad \quad \quad \operatorname{interval}(1, n)) \end{array}}}

\vnbbookfile{theorem-library/matunit-matact-type}{theorem-library/matunit-matact-type}
\noindent{\small matunit-matact-type.scm -- two MATOF typings, PROVEN from the definitions.}\par
\vnbprop{matunit-type}{\vnbstmt{\begin{array}{@{}l@{}} \forall a, n, k, l. \\ \quad (\operatorname{is\text{-}ring}(a) \wedge (n \in \mathbb{N})) \Rightarrow (\operatorname{matunit}(a, n, k, l) \in \operatorname{mat}(n, n, \operatorname{carr}(a))) \end{array}}}
\vnbprop{matact-type}{\vnbstmt{\begin{array}{@{}l@{}} \forall d, m, n, q, p, u. \\ \quad \operatorname{is\text{-}module}(d) \wedge \\ \quad p \in \operatorname{mat}(m, n, \operatorname{carr}(\operatorname{scal}(d))) \wedge \\ \quad u \in \operatorname{mat}(n, q, \operatorname{vec}(d)) \wedge \\ \quad (n = 0) \Rightarrow ((m = 0) \vee (q = 0)) \Rightarrow \\ \quad \quad \operatorname{matact}(d, p, u) \in \operatorname{mat}(m, q, \operatorname{vec}(d)) \end{array}}}

\vnbbookfile{theorem-library/lam-fun-bricks}{theorem-library/lam-fun-bricks}
\noindent{\small lam-fun-bricks.scm -- the lambda FUN-typing bricks of the matrix arc, PROVEN.}\par
\vnbprop{matprod-summand-type}{\vnbstmt{\begin{array}{@{}l@{}} \forall a, m, n, b, p, q, i, c. \\ \quad \operatorname{is\text{-}ring}(a) \wedge \\ \quad p \in \operatorname{mat}(m, n, \operatorname{carr}(a)) \wedge \\ \quad q \in \operatorname{mat}(n, b, \operatorname{carr}(a)) \wedge \\ \quad i \in \operatorname{interval}(1, m) \wedge \\ \quad c \in \operatorname{interval}(1, b) \Rightarrow \\ \quad \quad (\lambda j \in \operatorname{interval}(1, n).\; \operatorname{mul}(a)(\operatorname{entry}(p, i, j), \operatorname{entry}(q, j, c))) \\ \quad \quad \in\; (\operatorname{interval}(1, n) \to \operatorname{carr}(\operatorname{ring\text{-}additive\text{-}ag}(a))) \end{array}}}
\vnbprop{matmul-assoc-summand-type}{\vnbstmt{\begin{array}{@{}l@{}} \forall a, m, n, k, l, p, q, r, w, b. \\ \quad \operatorname{is\text{-}ring}(a) \wedge \\ \quad p \in \operatorname{mat}(m, n, \operatorname{carr}(a)) \wedge \\ \quad q \in \operatorname{mat}(n, k, \operatorname{carr}(a)) \wedge \\ \quad r \in \operatorname{mat}(k, l, \operatorname{carr}(a)) \wedge \\ \quad w \in \operatorname{interval}(1, m) \wedge \\ \quad b \in \operatorname{interval}(1, l) \Rightarrow \\ \quad \quad (\lambda z \in (\operatorname{interval}(1, k) \times \operatorname{interval}(1, n)).\; \\ \quad \quad \quad \operatorname{mul}(a)(\operatorname{mul}(a)(\operatorname{entry}(p, w, \operatorname{nth}(2, z)), \operatorname{entry}(q, \operatorname{nth}(2, z), \operatorname{nth}(1, z))), \operatorname{entry}(r, \operatorname{nth}(1, z), b))) \\ \quad \quad \in\; ((\operatorname{interval}(1, k) \times \operatorname{interval}(1, n)) \to \operatorname{carr}(\operatorname{ring\text{-}additive\text{-}ag}(a))) \end{array}}}
\vnbprop{matact-summand-type}{\vnbstmt{\begin{array}{@{}l@{}} \forall d, m, n, q, p, u, i, c. \\ \quad \operatorname{is\text{-}module}(d) \wedge \\ \quad p \in \operatorname{mat}(m, n, \operatorname{carr}(\operatorname{scal}(d))) \wedge \\ \quad u \in \operatorname{mat}(n, q, \operatorname{vec}(d)) \wedge \\ \quad i \in \operatorname{interval}(1, m) \wedge \\ \quad c \in \operatorname{interval}(1, q) \Rightarrow \\ \quad \quad (\lambda j \in \operatorname{interval}(1, n).\; \operatorname{act}(d)(\operatorname{entry}(p, i, j), \operatorname{entry}(u, j, c))) \\ \quad \quad \in\; (\operatorname{interval}(1, n) \to \operatorname{carr}(\operatorname{module\text{-}vector\text{-}ag}(d))) \end{array}}}
\vnbprop{matunit-summand-type}{\vnbstmt{\begin{array}{@{}l@{}} \forall a, m, n, p, k, l, i, c. \\ \quad (\operatorname{is\text{-}ring}(a) \wedge ((p \in \operatorname{mat}(m, n, \operatorname{carr}(a))) \wedge ((i \in \operatorname{interval}(1, m)) \wedge (c \in \operatorname{interval}(1, n))))) \Rightarrow  \\ \quad \quad (\lambda j \in \operatorname{interval}(1, n).\; \operatorname{mul}(a)(\operatorname{entry}(p, i, j), \operatorname{entry}(\operatorname{matunit}(a, n, k, l), j, c))) \\ \quad \quad \in\; (\operatorname{interval}(1, n) \to \operatorname{carr}(\operatorname{ring\text{-}additive\text{-}ag}(a))) \end{array}}}
\vnbprop{tel-outf-type}{\vnbstmt{\begin{array}{@{}l@{}} \forall a, m, n, k, l, p, q, r, w, b. \\ \quad \operatorname{is\text{-}ring}(a) \wedge \\ \quad p \in \operatorname{mat}(m, n, \operatorname{carr}(a)) \wedge \\ \quad q \in \operatorname{mat}(n, k, \operatorname{carr}(a)) \wedge \\ \quad r \in \operatorname{mat}(k, l, \operatorname{carr}(a)) \wedge \\ \quad w \in \operatorname{interval}(1, m) \wedge \\ \quad b \in \operatorname{interval}(1, l) \wedge \\ \quad 1 \leq n \Rightarrow \\ \quad \quad (\lambda j \in \operatorname{interval}(1, k).\; \operatorname{mul}(a)(\operatorname{entry}(\operatorname{matmul}(a, p, q), w, j), \operatorname{entry}(r, j, b))) \\ \quad \quad \in\; (\operatorname{interval}(1, k) \to \operatorname{carr}(\operatorname{ring\text{-}additive\text{-}ag}(a))) \end{array}}}
\vnbprop{tel-tout-type}{\vnbstmt{\begin{array}{@{}l@{}} \forall a, m, n, k, l, p, q, r, w, b. \\ \quad \operatorname{is\text{-}ring}(a) \wedge \\ \quad p \in \operatorname{mat}(m, n, \operatorname{carr}(a)) \wedge \\ \quad q \in \operatorname{mat}(n, k, \operatorname{carr}(a)) \wedge \\ \quad r \in \operatorname{mat}(k, l, \operatorname{carr}(a)) \wedge \\ \quad w \in \operatorname{interval}(1, m) \wedge \\ \quad b \in \operatorname{interval}(1, l) \Rightarrow \\ \quad \quad (\lambda c \in \operatorname{interval}(1, k).\; \\ \quad \quad \quad \operatorname{finsum}(\operatorname{ring\text{-}additive\text{-}ag}(a), \\ \quad \quad \quad \quad (\lambda j \in \operatorname{interval}(1, n).\; \\ \quad \quad \quad \quad \quad (\lambda z \in (\operatorname{interval}(1, k) \times \operatorname{interval}(1, n)).\; \operatorname{mul}(a)(\operatorname{mul}(a)(\operatorname{entry}(p, w, \operatorname{nth}(2, z)), \operatorname{entry}(q, \operatorname{nth}(2, z), \operatorname{nth}(1, z))), \operatorname{entry}(r, \operatorname{nth}(1, z), b)))(\langle c, j \rangle)), \\ \quad \quad \quad \quad \operatorname{interval}(1, n))) \\ \quad \quad \in\; (\operatorname{interval}(1, k) \to \operatorname{carr}(\operatorname{ring\text{-}additive\text{-}ag}(a))) \end{array}}}
\vnbprop{tel-inf-type}{\vnbstmt{\begin{array}{@{}l@{}} \forall a, m, n, k, l, p, q, r, w, b, x. \\ \quad \operatorname{is\text{-}ring}(a) \wedge \\ \quad p \in \operatorname{mat}(m, n, \operatorname{carr}(a)) \wedge \\ \quad q \in \operatorname{mat}(n, k, \operatorname{carr}(a)) \wedge \\ \quad r \in \operatorname{mat}(k, l, \operatorname{carr}(a)) \wedge \\ \quad w \in \operatorname{interval}(1, m) \wedge \\ \quad b \in \operatorname{interval}(1, l) \wedge \\ \quad x \in \operatorname{interval}(1, k) \Rightarrow \\ \quad \quad (\lambda j \in \operatorname{interval}(1, n).\; \operatorname{mul}(a)(\operatorname{entry}(p, w, j), \operatorname{entry}(q, j, x))) \\ \quad \quad \in\; (\operatorname{interval}(1, n) \to \operatorname{carr}(a)) \end{array}}}
\vnbprop{tel-dist-type}{\vnbstmt{\begin{array}{@{}l@{}} \forall a, m, n, k, l, p, q, r, w, b, x. \\ \quad \operatorname{is\text{-}ring}(a) \wedge \\ \quad p \in \operatorname{mat}(m, n, \operatorname{carr}(a)) \wedge \\ \quad q \in \operatorname{mat}(n, k, \operatorname{carr}(a)) \wedge \\ \quad r \in \operatorname{mat}(k, l, \operatorname{carr}(a)) \wedge \\ \quad w \in \operatorname{interval}(1, m) \wedge \\ \quad b \in \operatorname{interval}(1, l) \wedge \\ \quad x \in \operatorname{interval}(1, k) \Rightarrow \\ \quad \quad (\lambda z \in \operatorname{interval}(1, n).\; \\ \quad \quad \quad \operatorname{mul}(a)((\lambda j \in \operatorname{interval}(1, n).\; \operatorname{mul}(a)(\operatorname{entry}(p, w, j), \operatorname{entry}(q, j, x)))(z), \\ \quad \quad \quad \quad \operatorname{entry}(r, x, b))) \\ \quad \quad \in\; (\operatorname{interval}(1, n) \to \operatorname{carr}(\operatorname{ring\text{-}additive\text{-}ag}(a))) \end{array}}}
\vnbprop{tel-red-type}{\vnbstmt{\begin{array}{@{}l@{}} \forall a, m, n, k, l, p, q, r, w, b, x. \\ \quad \operatorname{is\text{-}ring}(a) \wedge \\ \quad p \in \operatorname{mat}(m, n, \operatorname{carr}(a)) \wedge \\ \quad q \in \operatorname{mat}(n, k, \operatorname{carr}(a)) \wedge \\ \quad r \in \operatorname{mat}(k, l, \operatorname{carr}(a)) \wedge \\ \quad w \in \operatorname{interval}(1, m) \wedge \\ \quad b \in \operatorname{interval}(1, l) \wedge \\ \quad x \in \operatorname{interval}(1, k) \Rightarrow \\ \quad \quad (\lambda j \in \operatorname{interval}(1, n).\; \operatorname{mul}(a)(\operatorname{mul}(a)(\operatorname{entry}(p, w, j), \operatorname{entry}(q, j, x)), \operatorname{entry}(r, x, b))) \\ \quad \quad \in\; (\operatorname{interval}(1, n) \to \operatorname{carr}(\operatorname{ring\text{-}additive\text{-}ag}(a))) \end{array}}}
\vnbprop{ter-outf-type}{\vnbstmt{\begin{array}{@{}l@{}} \forall a, m, n, k, l, p, q, r, w, b. \\ \quad \operatorname{is\text{-}ring}(a) \wedge \\ \quad p \in \operatorname{mat}(m, n, \operatorname{carr}(a)) \wedge \\ \quad q \in \operatorname{mat}(n, k, \operatorname{carr}(a)) \wedge \\ \quad r \in \operatorname{mat}(k, l, \operatorname{carr}(a)) \wedge \\ \quad w \in \operatorname{interval}(1, m) \wedge \\ \quad b \in \operatorname{interval}(1, l) \wedge \\ \quad 1 \leq k \Rightarrow \\ \quad \quad (\lambda j \in \operatorname{interval}(1, n).\; \operatorname{mul}(a)(\operatorname{entry}(p, w, j), \operatorname{entry}(\operatorname{matmul}(a, q, r), j, b))) \\ \quad \quad \in\; (\operatorname{interval}(1, n) \to \operatorname{carr}(\operatorname{ring\text{-}additive\text{-}ag}(a))) \end{array}}}
\vnbprop{ter-tout-type}{\vnbstmt{\begin{array}{@{}l@{}} \forall a, m, n, k, l, p, q, r, w, b. \\ \quad \operatorname{is\text{-}ring}(a) \wedge \\ \quad p \in \operatorname{mat}(m, n, \operatorname{carr}(a)) \wedge \\ \quad q \in \operatorname{mat}(n, k, \operatorname{carr}(a)) \wedge \\ \quad r \in \operatorname{mat}(k, l, \operatorname{carr}(a)) \wedge \\ \quad w \in \operatorname{interval}(1, m) \wedge \\ \quad b \in \operatorname{interval}(1, l) \Rightarrow \\ \quad \quad (\lambda j \in \operatorname{interval}(1, n).\; \\ \quad \quad \quad \operatorname{finsum}(\operatorname{ring\text{-}additive\text{-}ag}(a), \\ \quad \quad \quad \quad (\lambda c \in \operatorname{interval}(1, k).\; \\ \quad \quad \quad \quad \quad (\lambda z \in (\operatorname{interval}(1, k) \times \operatorname{interval}(1, n)).\; \operatorname{mul}(a)(\operatorname{mul}(a)(\operatorname{entry}(p, w, \operatorname{nth}(2, z)), \operatorname{entry}(q, \operatorname{nth}(2, z), \operatorname{nth}(1, z))), \operatorname{entry}(r, \operatorname{nth}(1, z), b)))(\langle c, j \rangle)), \\ \quad \quad \quad \quad \operatorname{interval}(1, k))) \\ \quad \quad \in\; (\operatorname{interval}(1, n) \to \operatorname{carr}(\operatorname{ring\text{-}additive\text{-}ag}(a))) \end{array}}}
\vnbprop{ter-gj-type}{\vnbstmt{\begin{array}{@{}l@{}} \forall a, m, n, k, l, p, q, r, w, b, x. \\ \quad \operatorname{is\text{-}ring}(a) \wedge \\ \quad p \in \operatorname{mat}(m, n, \operatorname{carr}(a)) \wedge \\ \quad q \in \operatorname{mat}(n, k, \operatorname{carr}(a)) \wedge \\ \quad r \in \operatorname{mat}(k, l, \operatorname{carr}(a)) \wedge \\ \quad w \in \operatorname{interval}(1, m) \wedge \\ \quad b \in \operatorname{interval}(1, l) \wedge \\ \quad x \in \operatorname{interval}(1, n) \Rightarrow \\ \quad \quad (\lambda j \in \operatorname{interval}(1, k).\; \operatorname{mul}(a)(\operatorname{entry}(q, x, j), \operatorname{entry}(r, j, b))) \\ \quad \quad \in\; (\operatorname{interval}(1, k) \to \operatorname{carr}(a)) \end{array}}}
\vnbprop{ter-dist-type}{\vnbstmt{\begin{array}{@{}l@{}} \forall a, m, n, k, l, p, q, r, w, b, x. \\ \quad \operatorname{is\text{-}ring}(a) \wedge \\ \quad p \in \operatorname{mat}(m, n, \operatorname{carr}(a)) \wedge \\ \quad q \in \operatorname{mat}(n, k, \operatorname{carr}(a)) \wedge \\ \quad r \in \operatorname{mat}(k, l, \operatorname{carr}(a)) \wedge \\ \quad w \in \operatorname{interval}(1, m) \wedge \\ \quad b \in \operatorname{interval}(1, l) \wedge \\ \quad x \in \operatorname{interval}(1, n) \Rightarrow \\ \quad \quad (\lambda z \in \operatorname{interval}(1, k).\; \\ \quad \quad \quad \operatorname{mul}(a)(\operatorname{entry}(p, w, x), \\ \quad \quad \quad \quad (\lambda j \in \operatorname{interval}(1, k).\; \operatorname{mul}(a)(\operatorname{entry}(q, x, j), \operatorname{entry}(r, j, b)))(z))) \\ \quad \quad \in\; (\operatorname{interval}(1, k) \to \operatorname{carr}(\operatorname{ring\text{-}additive\text{-}ag}(a))) \end{array}}}
\vnbprop{ter-red-type}{\vnbstmt{\begin{array}{@{}l@{}} \forall a, m, n, k, l, p, q, r, w, b, x. \\ \quad \operatorname{is\text{-}ring}(a) \wedge \\ \quad p \in \operatorname{mat}(m, n, \operatorname{carr}(a)) \wedge \\ \quad q \in \operatorname{mat}(n, k, \operatorname{carr}(a)) \wedge \\ \quad r \in \operatorname{mat}(k, l, \operatorname{carr}(a)) \wedge \\ \quad w \in \operatorname{interval}(1, m) \wedge \\ \quad b \in \operatorname{interval}(1, l) \wedge \\ \quad x \in \operatorname{interval}(1, n) \Rightarrow \\ \quad \quad (\lambda c \in \operatorname{interval}(1, k).\; \operatorname{mul}(a)(\operatorname{mul}(a)(\operatorname{entry}(p, w, x), \operatorname{entry}(q, x, c)), \operatorname{entry}(r, c, b))) \\ \quad \quad \in\; (\operatorname{interval}(1, k) \to \operatorname{carr}(\operatorname{ring\text{-}additive\text{-}ag}(a))) \end{array}}}
\vnbprop{mal-outf-type}{\vnbstmt{\begin{array}{@{}l@{}} \forall d, m, n, k, l, p, q, u, w, a. \\ \quad \operatorname{is\text{-}module}(d) \wedge \\ \quad p \in \operatorname{mat}(m, n, \operatorname{carr}(\operatorname{scal}(d))) \wedge \\ \quad q \in \operatorname{mat}(n, k, \operatorname{carr}(\operatorname{scal}(d))) \wedge \\ \quad u \in \operatorname{mat}(k, l, \operatorname{vec}(d)) \wedge \\ \quad w \in \operatorname{interval}(1, m) \wedge \\ \quad a \in \operatorname{interval}(1, l) \wedge \\ \quad 1 \leq n \Rightarrow \\ \quad \quad (\lambda j \in \operatorname{interval}(1, k).\; \operatorname{act}(d)(\operatorname{entry}(\operatorname{matmul}(\operatorname{scal}(d), p, q), w, j), \operatorname{entry}(u, j, a))) \\ \quad \quad \in\; (\operatorname{interval}(1, k) \to \operatorname{carr}(\operatorname{module\text{-}vector\text{-}ag}(d))) \end{array}}}
\vnbprop{mal-tout-type}{\vnbstmt{\begin{array}{@{}l@{}} \forall d, m, n, k, l, p, q, u, w, a. \\ \quad \operatorname{is\text{-}module}(d) \wedge \\ \quad p \in \operatorname{mat}(m, n, \operatorname{carr}(\operatorname{scal}(d))) \wedge \\ \quad q \in \operatorname{mat}(n, k, \operatorname{carr}(\operatorname{scal}(d))) \wedge \\ \quad u \in \operatorname{mat}(k, l, \operatorname{vec}(d)) \wedge \\ \quad w \in \operatorname{interval}(1, m) \wedge \\ \quad a \in \operatorname{interval}(1, l) \Rightarrow \\ \quad \quad (\lambda c \in \operatorname{interval}(1, k).\; \\ \quad \quad \quad \operatorname{finsum}(\operatorname{module\text{-}vector\text{-}ag}(d), \\ \quad \quad \quad \quad (\lambda j \in \operatorname{interval}(1, n).\; \\ \quad \quad \quad \quad \quad (\lambda z \in (\operatorname{interval}(1, k) \times \operatorname{interval}(1, n)).\; \operatorname{act}(d)(\operatorname{mul}(\operatorname{scal}(d))(\operatorname{entry}(p, w, \operatorname{nth}(2, z)), \operatorname{entry}(q, \operatorname{nth}(2, z), \operatorname{nth}(1, z))), \operatorname{entry}(u, \operatorname{nth}(1, z), a)))(\langle c, j \rangle)), \\ \quad \quad \quad \quad \operatorname{interval}(1, n))) \\ \quad \quad \in\; (\operatorname{interval}(1, k) \to \operatorname{carr}(\operatorname{module\text{-}vector\text{-}ag}(d))) \end{array}}}
\vnbprop{mal-inf-type}{\vnbstmt{\begin{array}{@{}l@{}} \forall d, m, n, k, l, p, q, u, w, a, x. \\ \quad \operatorname{is\text{-}module}(d) \wedge \\ \quad p \in \operatorname{mat}(m, n, \operatorname{carr}(\operatorname{scal}(d))) \wedge \\ \quad q \in \operatorname{mat}(n, k, \operatorname{carr}(\operatorname{scal}(d))) \wedge \\ \quad u \in \operatorname{mat}(k, l, \operatorname{vec}(d)) \wedge \\ \quad w \in \operatorname{interval}(1, m) \wedge \\ \quad a \in \operatorname{interval}(1, l) \wedge \\ \quad x \in \operatorname{interval}(1, k) \Rightarrow \\ \quad \quad (\lambda j \in \operatorname{interval}(1, n).\; \operatorname{mul}(\operatorname{scal}(d))(\operatorname{entry}(p, w, j), \operatorname{entry}(q, j, x))) \\ \quad \quad \in\; (\operatorname{interval}(1, n) \to \operatorname{carr}(\operatorname{scal}(d))) \end{array}}}
\vnbprop{mal-dist-type}{\vnbstmt{\begin{array}{@{}l@{}} \forall d, m, n, k, l, p, q, u, w, a, x. \\ \quad \operatorname{is\text{-}module}(d) \wedge \\ \quad p \in \operatorname{mat}(m, n, \operatorname{carr}(\operatorname{scal}(d))) \wedge \\ \quad q \in \operatorname{mat}(n, k, \operatorname{carr}(\operatorname{scal}(d))) \wedge \\ \quad u \in \operatorname{mat}(k, l, \operatorname{vec}(d)) \wedge \\ \quad w \in \operatorname{interval}(1, m) \wedge \\ \quad a \in \operatorname{interval}(1, l) \wedge \\ \quad x \in \operatorname{interval}(1, k) \Rightarrow \\ \quad \quad (\lambda z \in \operatorname{interval}(1, n).\; \\ \quad \quad \quad \operatorname{act}(d)((\lambda j \in \operatorname{interval}(1, n).\; \operatorname{mul}(\operatorname{scal}(d))(\operatorname{entry}(p, w, j), \operatorname{entry}(q, j, x)))(z), \\ \quad \quad \quad \quad \operatorname{entry}(u, x, a))) \\ \quad \quad \in\; (\operatorname{interval}(1, n) \to \operatorname{carr}(\operatorname{module\text{-}vector\text{-}ag}(d))) \end{array}}}
\vnbprop{mal-red-type}{\vnbstmt{\begin{array}{@{}l@{}} \forall d, m, n, k, l, p, q, u, w, a, x. \\ \quad \operatorname{is\text{-}module}(d) \wedge \\ \quad p \in \operatorname{mat}(m, n, \operatorname{carr}(\operatorname{scal}(d))) \wedge \\ \quad q \in \operatorname{mat}(n, k, \operatorname{carr}(\operatorname{scal}(d))) \wedge \\ \quad u \in \operatorname{mat}(k, l, \operatorname{vec}(d)) \wedge \\ \quad w \in \operatorname{interval}(1, m) \wedge \\ \quad a \in \operatorname{interval}(1, l) \wedge \\ \quad x \in \operatorname{interval}(1, k) \Rightarrow \\ \quad \quad (\lambda j \in \operatorname{interval}(1, n).\; \\ \quad \quad \quad \operatorname{act}(d)(\operatorname{mul}(\operatorname{scal}(d))(\operatorname{entry}(p, w, j), \operatorname{entry}(q, j, x)), \operatorname{entry}(u, x, a))) \\ \quad \quad \in\; (\operatorname{interval}(1, n) \to \operatorname{carr}(\operatorname{module\text{-}vector\text{-}ag}(d))) \end{array}}}
\vnbprop{mar-outf-type}{\vnbstmt{\begin{array}{@{}l@{}} \forall d, m, n, k, l, p, q, u, w, a. \\ \quad \operatorname{is\text{-}module}(d) \wedge \\ \quad p \in \operatorname{mat}(m, n, \operatorname{carr}(\operatorname{scal}(d))) \wedge \\ \quad q \in \operatorname{mat}(n, k, \operatorname{carr}(\operatorname{scal}(d))) \wedge \\ \quad u \in \operatorname{mat}(k, l, \operatorname{vec}(d)) \wedge \\ \quad w \in \operatorname{interval}(1, m) \wedge \\ \quad a \in \operatorname{interval}(1, l) \wedge \\ \quad 1 \leq k \Rightarrow \\ \quad \quad (\lambda j \in \operatorname{interval}(1, n).\; \operatorname{act}(d)(\operatorname{entry}(p, w, j), \operatorname{entry}(\operatorname{matact}(d, q, u), j, a))) \\ \quad \quad \in\; (\operatorname{interval}(1, n) \to \operatorname{carr}(\operatorname{module\text{-}vector\text{-}ag}(d))) \end{array}}}
\vnbprop{mar-tout-type}{\vnbstmt{\begin{array}{@{}l@{}} \forall d, m, n, k, l, p, q, u, w, a. \\ \quad \operatorname{is\text{-}module}(d) \wedge \\ \quad p \in \operatorname{mat}(m, n, \operatorname{carr}(\operatorname{scal}(d))) \wedge \\ \quad q \in \operatorname{mat}(n, k, \operatorname{carr}(\operatorname{scal}(d))) \wedge \\ \quad u \in \operatorname{mat}(k, l, \operatorname{vec}(d)) \wedge \\ \quad w \in \operatorname{interval}(1, m) \wedge \\ \quad a \in \operatorname{interval}(1, l) \Rightarrow \\ \quad \quad (\lambda j \in \operatorname{interval}(1, n).\; \\ \quad \quad \quad \operatorname{finsum}(\operatorname{module\text{-}vector\text{-}ag}(d), \\ \quad \quad \quad \quad (\lambda c \in \operatorname{interval}(1, k).\; \\ \quad \quad \quad \quad \quad (\lambda z \in (\operatorname{interval}(1, k) \times \operatorname{interval}(1, n)).\; \operatorname{act}(d)(\operatorname{mul}(\operatorname{scal}(d))(\operatorname{entry}(p, w, \operatorname{nth}(2, z)), \operatorname{entry}(q, \operatorname{nth}(2, z), \operatorname{nth}(1, z))), \operatorname{entry}(u, \operatorname{nth}(1, z), a)))(\langle c, j \rangle)), \\ \quad \quad \quad \quad \operatorname{interval}(1, k))) \\ \quad \quad \in\; (\operatorname{interval}(1, n) \to \operatorname{carr}(\operatorname{module\text{-}vector\text{-}ag}(d))) \end{array}}}
\vnbprop{mar-gj-type}{\vnbstmt{\begin{array}{@{}l@{}} \forall d, m, n, k, l, p, q, u, w, a, x. \\ \quad \operatorname{is\text{-}module}(d) \wedge \\ \quad p \in \operatorname{mat}(m, n, \operatorname{carr}(\operatorname{scal}(d))) \wedge \\ \quad q \in \operatorname{mat}(n, k, \operatorname{carr}(\operatorname{scal}(d))) \wedge \\ \quad u \in \operatorname{mat}(k, l, \operatorname{vec}(d)) \wedge \\ \quad w \in \operatorname{interval}(1, m) \wedge \\ \quad a \in \operatorname{interval}(1, l) \wedge \\ \quad x \in \operatorname{interval}(1, n) \Rightarrow \\ \quad \quad (\lambda j \in \operatorname{interval}(1, k).\; \operatorname{act}(d)(\operatorname{entry}(q, x, j), \operatorname{entry}(u, j, a))) \\ \quad \quad \in\; (\operatorname{interval}(1, k) \to \operatorname{carr}(\operatorname{module\text{-}vector\text{-}ag}(d))) \end{array}}}
\vnbprop{mar-dist-type}{\vnbstmt{\begin{array}{@{}l@{}} \forall d, m, n, k, l, p, q, u, w, a, x. \\ \quad \operatorname{is\text{-}module}(d) \wedge \\ \quad p \in \operatorname{mat}(m, n, \operatorname{carr}(\operatorname{scal}(d))) \wedge \\ \quad q \in \operatorname{mat}(n, k, \operatorname{carr}(\operatorname{scal}(d))) \wedge \\ \quad u \in \operatorname{mat}(k, l, \operatorname{vec}(d)) \wedge \\ \quad w \in \operatorname{interval}(1, m) \wedge \\ \quad a \in \operatorname{interval}(1, l) \wedge \\ \quad x \in \operatorname{interval}(1, n) \Rightarrow \\ \quad \quad (\lambda z \in \operatorname{interval}(1, k).\; \\ \quad \quad \quad \operatorname{act}(d)(\operatorname{entry}(p, w, x), \\ \quad \quad \quad \quad (\lambda j \in \operatorname{interval}(1, k).\; \operatorname{act}(d)(\operatorname{entry}(q, x, j), \operatorname{entry}(u, j, a)))(z))) \\ \quad \quad \in\; (\operatorname{interval}(1, k) \to \operatorname{carr}(\operatorname{module\text{-}vector\text{-}ag}(d))) \end{array}}}
\vnbprop{mar-red-type}{\vnbstmt{\begin{array}{@{}l@{}} \forall d, m, n, k, l, p, q, u, w, a, x. \\ \quad \operatorname{is\text{-}module}(d) \wedge \\ \quad p \in \operatorname{mat}(m, n, \operatorname{carr}(\operatorname{scal}(d))) \wedge \\ \quad q \in \operatorname{mat}(n, k, \operatorname{carr}(\operatorname{scal}(d))) \wedge \\ \quad u \in \operatorname{mat}(k, l, \operatorname{vec}(d)) \wedge \\ \quad w \in \operatorname{interval}(1, m) \wedge \\ \quad a \in \operatorname{interval}(1, l) \wedge \\ \quad x \in \operatorname{interval}(1, n) \Rightarrow \\ \quad \quad (\lambda c \in \operatorname{interval}(1, k).\; \\ \quad \quad \quad \operatorname{act}(d)(\operatorname{mul}(\operatorname{scal}(d))(\operatorname{entry}(p, w, x), \operatorname{entry}(q, x, c)), \operatorname{entry}(u, c, a))) \\ \quad \quad \in\; (\operatorname{interval}(1, k) \to \operatorname{carr}(\operatorname{module\text{-}vector\text{-}ag}(d))) \end{array}}}
\vnbprop{matact-assoc-summand-type}{\vnbstmt{\begin{array}{@{}l@{}} \forall d, m, n, k, l, p, q, u, w, a. \\ \quad \operatorname{is\text{-}module}(d) \wedge \\ \quad p \in \operatorname{mat}(m, n, \operatorname{carr}(\operatorname{scal}(d))) \wedge \\ \quad q \in \operatorname{mat}(n, k, \operatorname{carr}(\operatorname{scal}(d))) \wedge \\ \quad u \in \operatorname{mat}(k, l, \operatorname{vec}(d)) \wedge \\ \quad w \in \operatorname{interval}(1, m) \wedge \\ \quad a \in \operatorname{interval}(1, l) \Rightarrow \\ \quad \quad (\lambda z \in (\operatorname{interval}(1, k) \times \operatorname{interval}(1, n)).\; \\ \quad \quad \quad \operatorname{act}(d)(\operatorname{mul}(\operatorname{scal}(d))(\operatorname{entry}(p, w, \operatorname{nth}(2, z)), \operatorname{entry}(q, \operatorname{nth}(2, z), \operatorname{nth}(1, z))), \\ \quad \quad \quad \quad \operatorname{entry}(u, \operatorname{nth}(1, z), a))) \\ \quad \quad \in\; ((\operatorname{interval}(1, k) \times \operatorname{interval}(1, n)) \to \operatorname{carr}(\operatorname{module\text{-}vector\text{-}ag}(d))) \end{array}}}
\vnbprop{mra-combined-summand-type}{\vnbstmt{\begin{array}{@{}l@{}} \forall d, n, a, b, u. \\ \quad \operatorname{is\text{-}module}(d) \wedge \\ \quad a \in \operatorname{mat}(1, n, \operatorname{carr}(\operatorname{scal}(d))) \wedge \\ \quad b \in \operatorname{mat}(1, n, \operatorname{carr}(\operatorname{scal}(d))) \wedge \\ \quad u \in \operatorname{mat}(n, 1, \operatorname{vec}(d)) \Rightarrow \\ \quad \quad (\lambda z \in \operatorname{interval}(1, n).\; \\ \quad \quad \quad \operatorname{opr}(\operatorname{module\text{-}vector\text{-}ag}(d))((\lambda j \in \operatorname{interval}(1, n).\; \operatorname{act}(d)(\operatorname{entry}(a, 1, j), \operatorname{entry}(u, j, 1)))(z), \\ \quad \quad \quad \quad (\lambda j \in \operatorname{interval}(1, n).\; \operatorname{act}(d)(\operatorname{entry}(b, 1, j), \operatorname{entry}(u, j, 1)))(z))) \\ \quad \quad \in\; (\operatorname{interval}(1, n) \to \operatorname{carr}(\operatorname{module\text{-}vector\text{-}ag}(d))) \end{array}}}
\vnbprop{mrs-scaled-summand-type}{\vnbstmt{\begin{array}{@{}l@{}} \forall d, n, r, c, u. \\ \quad (\operatorname{is\text{-}module}(d) \wedge ((r \in \operatorname{carr}(\operatorname{scal}(d))) \wedge ((c \in \operatorname{mat}(1, n, \operatorname{carr}(\operatorname{scal}(d)))) \wedge (u \in \operatorname{mat}(n, 1, \operatorname{vec}(d)))))) \Rightarrow  \\ \quad \quad (\lambda z \in \operatorname{interval}(1, n).\; \\ \quad \quad \quad \operatorname{act}(d)(r, (\lambda j \in \operatorname{interval}(1, n).\; \operatorname{act}(d)(\operatorname{entry}(c, 1, j), \operatorname{entry}(u, j, 1)))(z))) \\ \quad \quad \in\; (\operatorname{interval}(1, n) \to \operatorname{carr}(\operatorname{module\text{-}vector\text{-}ag}(d))) \end{array}}}

\vnbbookfile{theorem-library/rake-identmat}{theorem-library/rake-identmat}
\noindent{\small rake-identmat.scm -- the two IDENTMAT identity laws and the module twin of matmul-entry, PROVEN (batch K, 2026-09-17).}\par
\vnbprop{identmat-left-identity}{\vnbstmt{\begin{array}{@{}l@{}} \forall a, m, n.; \forall p \in \operatorname{mat}(m, n, \operatorname{carr}(a)). \\ \quad \operatorname{is\text{-}ring}(a) \Rightarrow (\operatorname{matmul}(a, \operatorname{identmat}(a, m), p) = p) \end{array}}}
\vnbprop{identmat-right-identity}{\vnbstmt{\begin{array}{@{}l@{}} \forall a, m, n.; \forall p \in \operatorname{mat}(m, n, \operatorname{carr}(a)). \\ \quad \operatorname{is\text{-}ring}(a) \Rightarrow (\operatorname{matmul}(a, p, \operatorname{identmat}(a, n)) = p) \end{array}}}
\vnbprop{matact-entry}{\vnbstmt{\begin{array}{@{}l@{}} \forall d, m, n, q, p, u.; \forall i \in \operatorname{interval}(1, m), c \in \operatorname{interval}(1, q). \\ \quad (\operatorname{is\text{-}module}(d) \wedge ((p \in \operatorname{mat}(m, n, \operatorname{carr}(\operatorname{scal}(d)))) \wedge ((u \in \operatorname{mat}(n, q, \operatorname{vec}(d))) \wedge (1 \leq n)))) \Rightarrow  \\ \quad \quad \operatorname{entry}(\operatorname{matact}(d, p, u), i, c) \\ \quad \quad =\; \operatorname{finsum}(\operatorname{module\text{-}vector\text{-}ag}(d), \\ \quad \quad \quad (\lambda j \in \operatorname{interval}(1, n).\; \operatorname{act}(d)(\operatorname{entry}(p, i, j), \operatorname{entry}(u, j, c))), \\ \quad \quad \quad \operatorname{interval}(1, n)) \end{array}}}

\vnbbookfile{theorem-library/rake-generates-coeff}{theorem-library/rake-generates-coeff}
\noindent{\small rake-generates-coeff.scm -- generates-coeff-matrix, proven.}\par
\vnbprop{generates-coeff-matrix}{\vnbstmt{\begin{array}{@{}l@{}} \forall d, n, m, u, v. \\ \quad \operatorname{is\text{-}module}(d) \wedge \\ \quad u \in \operatorname{mat}(n, 1, \operatorname{vec}(d)) \wedge \\ \quad v \in \operatorname{mat}(m, 1, \operatorname{vec}(d)) \wedge \\ \quad (n = 0) \Rightarrow (m = 0) \wedge \\ \quad \operatorname{generates}(d, n, u) \Rightarrow \\ \quad \quad \exists a \in \operatorname{mat}(m, n, \operatorname{carr}(\operatorname{scal}(d))).\; (v = \operatorname{matact}(d, a, u)) \end{array}}}

\vnbbookfile{theorem-library/rake-smith-leaves}{theorem-library/rake-smith-leaves}
\noindent{\small rake-smith-leaves.scm -- finsum-interval-shift, PROVEN.}\par
\vnbprop{finsum-interval-shift}{\vnbstmt{\begin{array}{@{}l@{}} \forall g.; \forall n \in \mathbb{N}, f \in (\operatorname{interval}(1, (n + 1)) \to \operatorname{carr}(g)). \\ \quad \operatorname{is\text{-}abelian\text{-}group}(g) \Rightarrow  \\ \quad \quad \operatorname{finsum}(g, f, \operatorname{interval}(1, (n + 1))) \\ \quad \quad =\; \operatorname{opr}(g)(f(1), \operatorname{finsum}(g, (\lambda z \in \operatorname{interval}(1, n).\; f((z + 1))), \operatorname{interval}(1, n))) \end{array}}}

\vnbbookfile{theorem-library/module-zero-act}{theorem-library/module-zero-act}
\noindent{\small module-zero-act.scm -- 0 . x = 0\_V : the ring zero annihilates every vector.}\par
\vnbprop{module-zero-act}{\vnbstmt{\begin{array}{@{}l@{}} \forall s.; \forall x \in \operatorname{vec}(s). \\ \quad \operatorname{is\text{-}module}(s) \Rightarrow (\operatorname{act}(s)(\operatorname{zero}(\operatorname{scal}(s)), x) = \operatorname{vzero}(s)) \end{array}}}

\vnbbookfile{theorem-library/triple-entry-proof}{theorem-library/triple-entry-proof}
\noindent{\small triple-entry-proof.scm -- the two matrix-product entry-expansion lemmas that matmul-assoc rests on, PROVEN (were warranted in matrix.scm) from the (B) finite-sum bricks (finsum-additive.scm): finsum-congruence, and general-ring distribution of a scalar into a finite sum (finsum-ring-distrib-left/right-gen).}\par
\vnbprop{triple-entry-left}{\vnbstmt{\begin{array}{@{}l@{}} \forall a, m, n, k, l, p, q, r, w, b. \\ \quad \operatorname{is\text{-}ring}(a) \wedge \\ \quad p \in \operatorname{mat}(m, n, \operatorname{carr}(a)) \wedge \\ \quad q \in \operatorname{mat}(n, k, \operatorname{carr}(a)) \wedge \\ \quad r \in \operatorname{mat}(k, l, \operatorname{carr}(a)) \wedge \\ \quad 1 \leq n \wedge \\ \quad 1 \leq k \wedge \\ \quad w \in \operatorname{interval}(1, m) \wedge \\ \quad b \in \operatorname{interval}(1, l) \Rightarrow \\ \quad \quad \operatorname{entry}(\operatorname{matmul}(a, \operatorname{matmul}(a, p, q), r), w, b) \\ \quad \quad =\; \operatorname{finsum}(\operatorname{ring\text{-}additive\text{-}ag}(a), \\ \quad \quad \quad (\lambda c \in \operatorname{interval}(1, k).\; \\ \quad \quad \quad \quad \operatorname{finsum}(\operatorname{ring\text{-}additive\text{-}ag}(a), \\ \quad \quad \quad \quad \quad (\lambda j \in \operatorname{interval}(1, n).\; \\ \quad \quad \quad \quad \quad \quad (\lambda z \in (\operatorname{interval}(1, k) \times \operatorname{interval}(1, n)).\; \operatorname{mul}(a)(\operatorname{mul}(a)(\operatorname{entry}(p, w, \operatorname{nth}(2, z)), \operatorname{entry}(q, \operatorname{nth}(2, z), \operatorname{nth}(1, z))), \operatorname{entry}(r, \operatorname{nth}(1, z), b)))(\langle c, j \rangle)), \\ \quad \quad \quad \quad \quad \operatorname{interval}(1, n))), \\ \quad \quad \quad \operatorname{interval}(1, k)) \end{array}}}
\vnbprop{triple-entry-right}{\vnbstmt{\begin{array}{@{}l@{}} \forall a, m, n, k, l, p, q, r, w, b. \\ \quad \operatorname{is\text{-}ring}(a) \wedge \\ \quad p \in \operatorname{mat}(m, n, \operatorname{carr}(a)) \wedge \\ \quad q \in \operatorname{mat}(n, k, \operatorname{carr}(a)) \wedge \\ \quad r \in \operatorname{mat}(k, l, \operatorname{carr}(a)) \wedge \\ \quad 1 \leq n \wedge \\ \quad 1 \leq k \wedge \\ \quad w \in \operatorname{interval}(1, m) \wedge \\ \quad b \in \operatorname{interval}(1, l) \Rightarrow \\ \quad \quad \operatorname{entry}(\operatorname{matmul}(a, p, \operatorname{matmul}(a, q, r)), w, b) \\ \quad \quad =\; \operatorname{finsum}(\operatorname{ring\text{-}additive\text{-}ag}(a), \\ \quad \quad \quad (\lambda j \in \operatorname{interval}(1, n).\; \\ \quad \quad \quad \quad \operatorname{finsum}(\operatorname{ring\text{-}additive\text{-}ag}(a), \\ \quad \quad \quad \quad \quad (\lambda c \in \operatorname{interval}(1, k).\; \\ \quad \quad \quad \quad \quad \quad (\lambda z \in (\operatorname{interval}(1, k) \times \operatorname{interval}(1, n)).\; \operatorname{mul}(a)(\operatorname{mul}(a)(\operatorname{entry}(p, w, \operatorname{nth}(2, z)), \operatorname{entry}(q, \operatorname{nth}(2, z), \operatorname{nth}(1, z))), \operatorname{entry}(r, \operatorname{nth}(1, z), b)))(\langle c, j \rangle)), \\ \quad \quad \quad \quad \quad \operatorname{interval}(1, k))), \\ \quad \quad \quad \operatorname{interval}(1, n)) \end{array}}}

\vnbbookfile{theorem-library/matmul-assoc-proof}{theorem-library/matmul-assoc-proof}
\noindent{\small matmul-assoc-proof.scm -- matrix multiplication is associative, via Fubini.}\par
\vnbprop{matmul-assoc}{\vnbstmt{\begin{array}{@{}l@{}} \forall a, m, n, k, l, p, q, r. \\ \quad \operatorname{is\text{-}ring}(a) \wedge \\ \quad p \in \operatorname{mat}(m, n, \operatorname{carr}(a)) \wedge \\ \quad q \in \operatorname{mat}(n, k, \operatorname{carr}(a)) \wedge \\ \quad r \in \operatorname{mat}(k, l, \operatorname{carr}(a)) \wedge \\ \quad (n = 0) \Rightarrow ((m = 0) \vee ((k = 0) \wedge (l = 0))) \wedge \\ \quad (k = 0) \Rightarrow ((l = 0) \vee ((n = 0) \wedge (m = 0))) \Rightarrow \\ \quad \quad \operatorname{matmul}(a, \operatorname{matmul}(a, p, q), r) = \operatorname{matmul}(a, p, \operatorname{matmul}(a, q, r)) \end{array}}}

\vnbbookfile{theorem-library/elem-entry-readoffs}{theorem-library/elem-entry-readoffs}
\noindent{\small elem-entry-readoffs.scm -- the elementary-matrix ENTRY read-offs, PROVEN.}\par
\vnbprop{matunit-entries-defined}{\vnbstmt{\begin{array}{@{}l@{}} \forall a, n, k, l.; \forall b, c \in \operatorname{interval}(1, n). \\ \quad (\operatorname{is\text{-}ring}(a) \wedge (n \in \mathbb{N})) \Rightarrow  \\ \quad \quad (\lambda \langle i, j \rangle \in (\operatorname{interval}(1, n) \times \operatorname{interval}(1, n)).\; \left(\text{if }((i = k) \wedge (j = l))\text{ then }\operatorname{one}(a)\text{ else }\operatorname{zero}(a)\right))(b, c) \\ \quad \quad \in\; \mathbf{Set} \end{array}}}
\vnbprop{elem-f-entries-defined}{\vnbstmt{\begin{array}{@{}l@{}} \forall a, n, k, l.; \forall b, c \in \operatorname{interval}(1, n). \\ \quad (\operatorname{is\text{-}ring}(a) \wedge (n \in \mathbb{N})) \Rightarrow  \\ \quad \quad (\lambda \langle i, j \rangle \in (\operatorname{interval}(1, n) \times \operatorname{interval}(1, n)).\; \left(\text{if }(((i = k) \wedge (j = l)) \vee (((i = l) \wedge (j = k)) \vee ((i = j) \wedge (\neg (i = k) \wedge \neg (i = l)))))\text{ then }\operatorname{one}(a)\text{ else }\operatorname{zero}(a)\right))(b, c) \\ \quad \quad \in\; \mathbf{Set} \end{array}}}
\vnbprop{elem-g-entries-defined}{\vnbstmt{\begin{array}{@{}l@{}} \forall a, n, r, k, l.; \forall b, c \in \operatorname{interval}(1, n). \\ \quad (\operatorname{is\text{-}ring}(a) \wedge ((n \in \mathbb{N}) \wedge (r \in \operatorname{carr}(a)))) \Rightarrow  \\ \quad \quad (\lambda \langle i, j \rangle \in (\operatorname{interval}(1, n) \times \operatorname{interval}(1, n)).\; \left(\text{if }(i = j)\text{ then }\operatorname{one}(a)\text{ else }\left(\text{if }((i = k) \wedge (j = l))\text{ then }r\text{ else }\operatorname{zero}(a)\right)\right))(b, c) \\ \quad \quad \in\; \mathbf{Set} \end{array}}}
\vnbprop{elem-h-entries-defined}{\vnbstmt{\begin{array}{@{}l@{}} \forall a, n, r, k.; \forall b, c \in \operatorname{interval}(1, n). \\ \quad (\operatorname{is\text{-}ring}(a) \wedge ((n \in \mathbb{N}) \wedge (r \in \operatorname{carr}(a)))) \Rightarrow  \\ \quad \quad (\lambda \langle i, j \rangle \in (\operatorname{interval}(1, n) \times \operatorname{interval}(1, n)).\; \left(\text{if }(i = j)\text{ then }\left(\text{if }(i = k)\text{ then }r\text{ else }\operatorname{one}(a)\right)\text{ else }\operatorname{zero}(a)\right))(b, c) \\ \quad \quad \in\; \mathbf{Set} \end{array}}}
\vnbprop{entry-of-matunit}{\vnbstmt{\begin{array}{@{}l@{}} \forall a, n, k, l, i, j. \\ \quad (\operatorname{is\text{-}ring}(a) \wedge ((n \in \mathbb{N}) \wedge ((i \in \operatorname{interval}(1, n)) \wedge (j \in \operatorname{interval}(1, n))))) \Rightarrow  \\ \quad \quad \operatorname{entry}(\operatorname{matunit}(a, n, k, l), i, j) = \left(\text{if }((i = k) \wedge (j = l))\text{ then }\operatorname{one}(a)\text{ else }\operatorname{zero}(a)\right) \end{array}}}
\vnbprop{matunit-entry-off-row}{\vnbstmt{\begin{array}{@{}l@{}} \forall a, n, k, l, j, c. \\ \quad (\operatorname{is\text{-}ring}(a) \wedge ((n \in \mathbb{N}) \wedge ((j \in \operatorname{interval}(1, n)) \wedge ((c \in \operatorname{interval}(1, n)) \wedge \neg (j = k))))) \Rightarrow  \\ \quad \quad \operatorname{entry}(\operatorname{matunit}(a, n, k, l), j, c) = \operatorname{zero}(a) \end{array}}}
\vnbprop{matunit-entry-k-row}{\vnbstmt{\begin{array}{@{}l@{}} \forall a, n, k, l, c. \\ \quad (\operatorname{is\text{-}ring}(a) \wedge ((n \in \mathbb{N}) \wedge ((k \in \operatorname{interval}(1, n)) \wedge (c \in \operatorname{interval}(1, n))))) \Rightarrow  \\ \quad \quad \operatorname{entry}(\operatorname{matunit}(a, n, k, l), k, c) = \left(\text{if }(c = l)\text{ then }\operatorname{one}(a)\text{ else }\operatorname{zero}(a)\right) \end{array}}}
\vnbprop{entry-of-elem-f}{\vnbstmt{\begin{array}{@{}l@{}} \forall a, n, k, l, i, j. \\ \quad (\operatorname{is\text{-}ring}(a) \wedge ((n \in \mathbb{N}) \wedge ((i \in \operatorname{interval}(1, n)) \wedge (j \in \operatorname{interval}(1, n))))) \Rightarrow  \\ \quad \quad \operatorname{entry}(\operatorname{elem\text{-}f}(a, n, k, l), i, j) \\ \quad \quad =\; \left(\text{if }(((i = k) \wedge (j = l)) \vee (((i = l) \wedge (j = k)) \vee ((i = j) \wedge (\neg (i = k) \wedge \neg (i = l)))))\text{ then }\operatorname{one}(a)\text{ else }\operatorname{zero}(a)\right) \end{array}}}
\vnbprop{entry-of-elem-g}{\vnbstmt{\begin{array}{@{}l@{}} \forall a, n, r, k, l, i, j. \\ \quad (\operatorname{is\text{-}ring}(a) \wedge ((n \in \mathbb{N}) \wedge ((r \in \operatorname{carr}(a)) \wedge ((i \in \operatorname{interval}(1, n)) \wedge (j \in \operatorname{interval}(1, n)))))) \Rightarrow  \\ \quad \quad \operatorname{entry}(\operatorname{elem\text{-}g}(a, n, r, k, l), i, j) = \left(\text{if }(i = j)\text{ then }\operatorname{one}(a)\text{ else }\left(\text{if }((i = k) \wedge (j = l))\text{ then }r\text{ else }\operatorname{zero}(a)\right)\right) \end{array}}}
\vnbprop{entry-of-elem-h}{\vnbstmt{\begin{array}{@{}l@{}} \forall a, n, r, k, i, j. \\ \quad (\operatorname{is\text{-}ring}(a) \wedge ((n \in \mathbb{N}) \wedge ((r \in \operatorname{carr}(a)) \wedge ((i \in \operatorname{interval}(1, n)) \wedge (j \in \operatorname{interval}(1, n)))))) \Rightarrow  \\ \quad \quad \operatorname{entry}(\operatorname{elem\text{-}h}(a, n, r, k), i, j) = \left(\text{if }(i = j)\text{ then }\left(\text{if }(i = k)\text{ then }r\text{ else }\operatorname{one}(a)\right)\text{ else }\operatorname{zero}(a)\right) \end{array}}}
\vnbprop{elem-g-entry-off-l-off-diag}{\vnbstmt{\begin{array}{@{}l@{}} \forall a, n, r, k, l, j, c. \\ \quad \operatorname{is\text{-}ring}(a) \wedge \\ \quad n \in \mathbb{N} \wedge \\ \quad r \in \operatorname{carr}(a) \wedge \\ \quad j \in \operatorname{interval}(1, n) \wedge \\ \quad c \in \operatorname{interval}(1, n) \wedge \\ \quad \neg (c = l) \wedge \\ \quad \neg (j = c) \Rightarrow \\ \quad \quad \operatorname{entry}(\operatorname{elem\text{-}g}(a, n, r, k, l), j, c) = \operatorname{zero}(a) \end{array}}}
\vnbprop{elem-g-entry-off-l-diag}{\vnbstmt{\begin{array}{@{}l@{}} \forall a, n, r, k, l, c. \\ \quad (\operatorname{is\text{-}ring}(a) \wedge ((n \in \mathbb{N}) \wedge ((r \in \operatorname{carr}(a)) \wedge ((c \in \operatorname{interval}(1, n)) \wedge \neg (c = l))))) \Rightarrow  \\ \quad \quad \operatorname{entry}(\operatorname{elem\text{-}g}(a, n, r, k, l), c, c) = \operatorname{one}(a) \end{array}}}
\vnbprop{elem-g-entry-l-vanish}{\vnbstmt{\begin{array}{@{}l@{}} \forall a, n, r, k, l, j, c. \\ \quad \operatorname{is\text{-}ring}(a) \wedge \\ \quad n \in \mathbb{N} \wedge \\ \quad r \in \operatorname{carr}(a) \wedge \\ \quad j \in \operatorname{interval}(1, n) \wedge \\ \quad c \in \operatorname{interval}(1, n) \wedge \\ \quad c = l \wedge \\ \quad \neg (j = l) \wedge \\ \quad \neg (j = k) \Rightarrow \\ \quad \quad \operatorname{entry}(\operatorname{elem\text{-}g}(a, n, r, k, l), j, c) = \operatorname{zero}(a) \end{array}}}
\vnbprop{elem-g-entry-l-at-l}{\vnbstmt{\begin{array}{@{}l@{}} \forall a, n, r, k, l, c. \\ \quad (\operatorname{is\text{-}ring}(a) \wedge ((n \in \mathbb{N}) \wedge ((r \in \operatorname{carr}(a)) \wedge ((c \in \operatorname{interval}(1, n)) \wedge (c = l))))) \Rightarrow  \\ \quad \quad \operatorname{entry}(\operatorname{elem\text{-}g}(a, n, r, k, l), l, c) = \operatorname{one}(a) \end{array}}}
\vnbprop{elem-f-ck-off}{\vnbstmt{\begin{array}{@{}l@{}} \forall a, n, k, l, j, c. \\ \quad \operatorname{is\text{-}ring}(a) \wedge \\ \quad n \in \mathbb{N} \wedge \\ \quad j \in \operatorname{interval}(1, n) \wedge \\ \quad c \in \operatorname{interval}(1, n) \wedge \\ \quad c = k \wedge \\ \quad \neg (k = l) \wedge \\ \quad \neg (j = l) \Rightarrow \\ \quad \quad \operatorname{entry}(\operatorname{elem\text{-}f}(a, n, k, l), j, c) = \operatorname{zero}(a) \end{array}}}
\vnbprop{elem-f-cl-off}{\vnbstmt{\begin{array}{@{}l@{}} \forall a, n, k, l, j, c. \\ \quad \operatorname{is\text{-}ring}(a) \wedge \\ \quad n \in \mathbb{N} \wedge \\ \quad j \in \operatorname{interval}(1, n) \wedge \\ \quad c \in \operatorname{interval}(1, n) \wedge \\ \quad c = l \wedge \\ \quad \neg (k = l) \wedge \\ \quad \neg (j = k) \Rightarrow \\ \quad \quad \operatorname{entry}(\operatorname{elem\text{-}f}(a, n, k, l), j, c) = \operatorname{zero}(a) \end{array}}}
\vnbprop{elem-f-co-off}{\vnbstmt{\begin{array}{@{}l@{}} \forall a, n, k, l, j, c. \\ \quad \operatorname{is\text{-}ring}(a) \wedge \\ \quad n \in \mathbb{N} \wedge \\ \quad j \in \operatorname{interval}(1, n) \wedge \\ \quad c \in \operatorname{interval}(1, n) \wedge \\ \quad \neg (c = k) \wedge \\ \quad \neg (c = l) \wedge \\ \quad \neg (j = c) \Rightarrow \\ \quad \quad \operatorname{entry}(\operatorname{elem\text{-}f}(a, n, k, l), j, c) = \operatorname{zero}(a) \end{array}}}
\vnbprop{elem-f-co-at}{\vnbstmt{\begin{array}{@{}l@{}} \forall a, n, k, l, c. \\ \quad (\operatorname{is\text{-}ring}(a) \wedge ((n \in \mathbb{N}) \wedge ((c \in \operatorname{interval}(1, n)) \wedge (\neg (c = k) \wedge \neg (c = l))))) \Rightarrow  \\ \quad \quad \operatorname{entry}(\operatorname{elem\text{-}f}(a, n, k, l), c, c) = \operatorname{one}(a) \end{array}}}
\vnbprop{elem-h-entry-off-diag}{\vnbstmt{\begin{array}{@{}l@{}} \forall a, n, r, k, j, c. \\ \quad \operatorname{is\text{-}ring}(a) \wedge \\ \quad n \in \mathbb{N} \wedge \\ \quad r \in \operatorname{carr}(a) \wedge \\ \quad j \in \operatorname{interval}(1, n) \wedge \\ \quad c \in \operatorname{interval}(1, n) \wedge \\ \quad \neg (j = c) \Rightarrow \\ \quad \quad \operatorname{entry}(\operatorname{elem\text{-}h}(a, n, r, k), j, c) = \operatorname{zero}(a) \end{array}}}
\vnbprop{elem-h-entry-diag}{\vnbstmt{\begin{array}{@{}l@{}} \forall a, n, r, k, c. \\ \quad (\operatorname{is\text{-}ring}(a) \wedge ((n \in \mathbb{N}) \wedge ((r \in \operatorname{carr}(a)) \wedge (c \in \operatorname{interval}(1, n))))) \Rightarrow  \\ \quad \quad \operatorname{entry}(\operatorname{elem\text{-}h}(a, n, r, k), c, c) = \left(\text{if }(c = k)\text{ then }r\text{ else }\operatorname{one}(a)\right) \end{array}}}
\vnbprop{elem-f-rk-off}{\vnbstmt{\begin{array}{@{}l@{}} \forall a, n, k, l, i, j. \\ \quad \operatorname{is\text{-}ring}(a) \wedge \\ \quad n \in \mathbb{N} \wedge \\ \quad i \in \operatorname{interval}(1, n) \wedge \\ \quad j \in \operatorname{interval}(1, n) \wedge \\ \quad i = k \wedge \\ \quad \neg (k = l) \wedge \\ \quad \neg (j = l) \Rightarrow \\ \quad \quad \operatorname{entry}(\operatorname{elem\text{-}f}(a, n, k, l), i, j) = \operatorname{zero}(a) \end{array}}}
\vnbprop{elem-f-rl-off}{\vnbstmt{\begin{array}{@{}l@{}} \forall a, n, k, l, i, j. \\ \quad \operatorname{is\text{-}ring}(a) \wedge \\ \quad n \in \mathbb{N} \wedge \\ \quad i \in \operatorname{interval}(1, n) \wedge \\ \quad j \in \operatorname{interval}(1, n) \wedge \\ \quad i = l \wedge \\ \quad \neg (k = l) \wedge \\ \quad \neg (j = k) \Rightarrow \\ \quad \quad \operatorname{entry}(\operatorname{elem\text{-}f}(a, n, k, l), i, j) = \operatorname{zero}(a) \end{array}}}
\vnbprop{elem-f-ro-off}{\vnbstmt{\begin{array}{@{}l@{}} \forall a, n, k, l, i, j. \\ \quad \operatorname{is\text{-}ring}(a) \wedge \\ \quad n \in \mathbb{N} \wedge \\ \quad i \in \operatorname{interval}(1, n) \wedge \\ \quad j \in \operatorname{interval}(1, n) \wedge \\ \quad \neg (i = k) \wedge \\ \quad \neg (i = l) \wedge \\ \quad \neg (j = i) \Rightarrow \\ \quad \quad \operatorname{entry}(\operatorname{elem\text{-}f}(a, n, k, l), i, j) = \operatorname{zero}(a) \end{array}}}
\vnbprop{elem-g-rk-vanish}{\vnbstmt{\begin{array}{@{}l@{}} \forall a, n, r, k, l, i, j. \\ \quad \operatorname{is\text{-}ring}(a) \wedge \\ \quad n \in \mathbb{N} \wedge \\ \quad r \in \operatorname{carr}(a) \wedge \\ \quad i \in \operatorname{interval}(1, n) \wedge \\ \quad j \in \operatorname{interval}(1, n) \wedge \\ \quad i = k \wedge \\ \quad \neg (j = k) \wedge \\ \quad \neg (j = l) \Rightarrow \\ \quad \quad \operatorname{entry}(\operatorname{elem\text{-}g}(a, n, r, k, l), i, j) = \operatorname{zero}(a) \end{array}}}
\vnbprop{elem-g-ro-off}{\vnbstmt{\begin{array}{@{}l@{}} \forall a, n, r, k, l, i, j. \\ \quad \operatorname{is\text{-}ring}(a) \wedge \\ \quad n \in \mathbb{N} \wedge \\ \quad r \in \operatorname{carr}(a) \wedge \\ \quad i \in \operatorname{interval}(1, n) \wedge \\ \quad j \in \operatorname{interval}(1, n) \wedge \\ \quad \neg (i = k) \wedge \\ \quad \neg (j = i) \Rightarrow \\ \quad \quad \operatorname{entry}(\operatorname{elem\text{-}g}(a, n, r, k, l), i, j) = \operatorname{zero}(a) \end{array}}}
\vnbprop{elem-h-ro-off}{\vnbstmt{\begin{array}{@{}l@{}} \forall a, n, r, k, i, j. \\ \quad \operatorname{is\text{-}ring}(a) \wedge \\ \quad n \in \mathbb{N} \wedge \\ \quad r \in \operatorname{carr}(a) \wedge \\ \quad i \in \operatorname{interval}(1, n) \wedge \\ \quad j \in \operatorname{interval}(1, n) \wedge \\ \quad \neg (j = i) \Rightarrow \\ \quad \quad \operatorname{entry}(\operatorname{elem\text{-}h}(a, n, r, k), i, j) = \operatorname{zero}(a) \end{array}}}
\vnbprop{elem-f-ck-at}{\vnbstmt{\begin{array}{@{}l@{}} \forall a, n, k, l, c. \\ \quad (\operatorname{is\text{-}ring}(a) \wedge ((n \in \mathbb{N}) \wedge ((l \in \operatorname{interval}(1, n)) \wedge ((c \in \operatorname{interval}(1, n)) \wedge (c = k))))) \Rightarrow  \\ \quad \quad \operatorname{entry}(\operatorname{elem\text{-}f}(a, n, k, l), l, c) = \operatorname{one}(a) \end{array}}}
\vnbprop{elem-f-cl-at}{\vnbstmt{\begin{array}{@{}l@{}} \forall a, n, k, l, c. \\ \quad (\operatorname{is\text{-}ring}(a) \wedge ((n \in \mathbb{N}) \wedge ((k \in \operatorname{interval}(1, n)) \wedge ((c \in \operatorname{interval}(1, n)) \wedge (c = l))))) \Rightarrow  \\ \quad \quad \operatorname{entry}(\operatorname{elem\text{-}f}(a, n, k, l), k, c) = \operatorname{one}(a) \end{array}}}
\vnbprop{elem-f-col-k}{\vnbstmt{\begin{array}{@{}l@{}} \forall a, n, k, l, i. \\ \quad (\operatorname{is\text{-}ring}(a) \wedge ((n \in \mathbb{N}) \wedge ((i \in \operatorname{interval}(1, n)) \wedge ((k \in \operatorname{interval}(1, n)) \wedge \neg (k = l))))) \Rightarrow  \\ \quad \quad \operatorname{entry}(\operatorname{elem\text{-}f}(a, n, k, l), i, k) = \left(\text{if }(i = l)\text{ then }\operatorname{one}(a)\text{ else }\operatorname{zero}(a)\right) \end{array}}}
\vnbprop{elem-f-col-l}{\vnbstmt{\begin{array}{@{}l@{}} \forall a, n, k, l, i. \\ \quad (\operatorname{is\text{-}ring}(a) \wedge ((n \in \mathbb{N}) \wedge ((i \in \operatorname{interval}(1, n)) \wedge ((l \in \operatorname{interval}(1, n)) \wedge \neg (k = l))))) \Rightarrow  \\ \quad \quad \operatorname{entry}(\operatorname{elem\text{-}f}(a, n, k, l), i, l) = \left(\text{if }(i = k)\text{ then }\operatorname{one}(a)\text{ else }\operatorname{zero}(a)\right) \end{array}}}
\vnbprop{elem-f-col-other}{\vnbstmt{\begin{array}{@{}l@{}} \forall a, n, k, l, i, c. \\ \quad (\operatorname{is\text{-}ring}(a) \wedge ((n \in \mathbb{N}) \wedge ((i \in \operatorname{interval}(1, n)) \wedge ((c \in \operatorname{interval}(1, n)) \wedge (\neg (c = k) \wedge \neg (c = l)))))) \Rightarrow  \\ \quad \quad \operatorname{entry}(\operatorname{elem\text{-}f}(a, n, k, l), i, c) = \left(\text{if }(i = c)\text{ then }\operatorname{one}(a)\text{ else }\operatorname{zero}(a)\right) \end{array}}}
\vnbprop{elem-f-rk-at}{\vnbstmt{\begin{array}{@{}l@{}} \forall a, n, k, l, i. \\ \quad (\operatorname{is\text{-}ring}(a) \wedge ((n \in \mathbb{N}) \wedge ((i \in \operatorname{interval}(1, n)) \wedge ((l \in \operatorname{interval}(1, n)) \wedge ((i = k) \wedge \neg (k = l)))))) \Rightarrow  \\ \quad \quad \operatorname{entry}(\operatorname{elem\text{-}f}(a, n, k, l), i, l) = \operatorname{one}(a) \end{array}}}
\vnbprop{elem-f-rl-at}{\vnbstmt{\begin{array}{@{}l@{}} \forall a, n, k, l, i. \\ \quad (\operatorname{is\text{-}ring}(a) \wedge ((n \in \mathbb{N}) \wedge ((i \in \operatorname{interval}(1, n)) \wedge ((k \in \operatorname{interval}(1, n)) \wedge ((i = l) \wedge \neg (k = l)))))) \Rightarrow  \\ \quad \quad \operatorname{entry}(\operatorname{elem\text{-}f}(a, n, k, l), i, k) = \operatorname{one}(a) \end{array}}}
\vnbprop{elem-g-entry-l-at-k}{\vnbstmt{\begin{array}{@{}l@{}} \forall a, n, r, k, l, c. \\ \quad \operatorname{is\text{-}ring}(a) \wedge \\ \quad n \in \mathbb{N} \wedge \\ \quad r \in \operatorname{carr}(a) \wedge \\ \quad k \in \operatorname{interval}(1, n) \wedge \\ \quad c \in \operatorname{interval}(1, n) \wedge \\ \quad c = l \wedge \\ \quad \neg (k = l) \Rightarrow \\ \quad \quad \operatorname{entry}(\operatorname{elem\text{-}g}(a, n, r, k, l), k, c) = r \end{array}}}
\vnbprop{elem-g-col-k}{\vnbstmt{\begin{array}{@{}l@{}} \forall a, n, r, k, l, i. \\ \quad \operatorname{is\text{-}ring}(a) \wedge \\ \quad n \in \mathbb{N} \wedge \\ \quad r \in \operatorname{carr}(a) \wedge \\ \quad i \in \operatorname{interval}(1, n) \wedge \\ \quad k \in \operatorname{interval}(1, n) \wedge \\ \quad \neg (k = l) \Rightarrow \\ \quad \quad \operatorname{entry}(\operatorname{elem\text{-}g}(a, n, r, k, l), i, k) = \left(\text{if }(i = k)\text{ then }\operatorname{one}(a)\text{ else }\operatorname{zero}(a)\right) \end{array}}}
\vnbprop{elem-g-col-l}{\vnbstmt{\begin{array}{@{}l@{}} \forall a, n, r, k, l, i. \\ \quad (\operatorname{is\text{-}ring}(a) \wedge ((n \in \mathbb{N}) \wedge ((r \in \operatorname{carr}(a)) \wedge ((i \in \operatorname{interval}(1, n)) \wedge (l \in \operatorname{interval}(1, n)))))) \Rightarrow  \\ \quad \quad \operatorname{entry}(\operatorname{elem\text{-}g}(a, n, r, k, l), i, l) = \left(\text{if }(i = l)\text{ then }\operatorname{one}(a)\text{ else }\left(\text{if }(i = k)\text{ then }r\text{ else }\operatorname{zero}(a)\right)\right) \end{array}}}
\vnbprop{elem-g-col-other}{\vnbstmt{\begin{array}{@{}l@{}} \forall a, n, r, k, l, i, c. \\ \quad \operatorname{is\text{-}ring}(a) \wedge \\ \quad n \in \mathbb{N} \wedge \\ \quad r \in \operatorname{carr}(a) \wedge \\ \quad i \in \operatorname{interval}(1, n) \wedge \\ \quad c \in \operatorname{interval}(1, n) \wedge \\ \quad \neg (c = l) \Rightarrow \\ \quad \quad \operatorname{entry}(\operatorname{elem\text{-}g}(a, n, r, k, l), i, c) = \left(\text{if }(i = c)\text{ then }\operatorname{one}(a)\text{ else }\operatorname{zero}(a)\right) \end{array}}}
\vnbprop{elem-g-rk-at-k}{\vnbstmt{\begin{array}{@{}l@{}} \forall a, n, r, k, l, i. \\ \quad (\operatorname{is\text{-}ring}(a) \wedge ((n \in \mathbb{N}) \wedge ((r \in \operatorname{carr}(a)) \wedge ((i \in \operatorname{interval}(1, n)) \wedge ((k \in \operatorname{interval}(1, n)) \wedge (i = k)))))) \Rightarrow  \\ \quad \quad \operatorname{entry}(\operatorname{elem\text{-}g}(a, n, r, k, l), i, k) = \operatorname{one}(a) \end{array}}}
\vnbprop{elem-g-rk-at-l}{\vnbstmt{\begin{array}{@{}l@{}} \forall a, n, r, k, l, i. \\ \quad \operatorname{is\text{-}ring}(a) \wedge \\ \quad n \in \mathbb{N} \wedge \\ \quad r \in \operatorname{carr}(a) \wedge \\ \quad i \in \operatorname{interval}(1, n) \wedge \\ \quad l \in \operatorname{interval}(1, n) \wedge \\ \quad i = k \wedge \\ \quad \neg (k = l) \Rightarrow \\ \quad \quad \operatorname{entry}(\operatorname{elem\text{-}g}(a, n, r, k, l), i, l) = r \end{array}}}
\vnbprop{elem-g-ro-at}{\vnbstmt{\begin{array}{@{}l@{}} \forall a, n, r, k, l, i. \\ \quad (\operatorname{is\text{-}ring}(a) \wedge ((n \in \mathbb{N}) \wedge ((r \in \operatorname{carr}(a)) \wedge ((i \in \operatorname{interval}(1, n)) \wedge \neg (i = k))))) \Rightarrow  \\ \quad \quad \operatorname{entry}(\operatorname{elem\text{-}g}(a, n, r, k, l), i, i) = \operatorname{one}(a) \end{array}}}
\vnbprop{identmat-entries-defined}{\vnbstmt{\begin{array}{@{}l@{}} \forall a, n.; \forall b, c \in \operatorname{interval}(1, n). \\ \quad (\operatorname{is\text{-}ring}(a) \wedge (n \in \mathbb{N})) \Rightarrow  \\ \quad \quad (\lambda \langle i, j \rangle \in (\operatorname{interval}(1, n) \times \operatorname{interval}(1, n)).\; \left(\text{if }(i = j)\text{ then }\operatorname{one}(a)\text{ else }\operatorname{zero}(a)\right))(b, c) \\ \quad \quad \in\; \mathbf{Set} \end{array}}}
\vnbprop{unitrow-entries-defined}{\vnbstmt{\begin{array}{@{}l@{}} \forall a, n, i.; \forall b \in \operatorname{interval}(1, 1), j \in \operatorname{interval}(1, n). \\ \quad (\operatorname{is\text{-}ring}(a) \wedge (n \in \mathbb{N})) \Rightarrow  \\ \quad \quad (\lambda \langle \mathit{rw}, \mathit{cl} \rangle \in (\operatorname{interval}(1, 1) \times \operatorname{interval}(1, n)).\; \left(\text{if }(\mathit{cl} = i)\text{ then }\operatorname{one}(a)\text{ else }\operatorname{zero}(a)\right))(b, j) \\ \quad \quad \in\; \mathbf{Set} \end{array}}}
\vnbprop{zeromat-entries-defined}{\vnbstmt{\begin{array}{@{}l@{}} \forall a, m, n.; \forall b \in \operatorname{interval}(1, m), c \in \operatorname{interval}(1, n). \\ \quad (\operatorname{is\text{-}ring}(a) \wedge ((m \in \mathbb{N}) \wedge (n \in \mathbb{N}))) \Rightarrow  \\ \quad \quad (\lambda \langle i, j \rangle \in (\operatorname{interval}(1, m) \times \operatorname{interval}(1, n)).\; \operatorname{zero}(a))(b, c) \in \mathbf{Set} \end{array}}}
\vnbprop{entry-of-identmat}{\vnbstmt{\begin{array}{@{}l@{}} \forall a, n, i, j. \\ \quad (\operatorname{is\text{-}ring}(a) \wedge ((n \in \mathbb{N}) \wedge ((i \in \operatorname{interval}(1, n)) \wedge (j \in \operatorname{interval}(1, n))))) \Rightarrow  \\ \quad \quad \operatorname{entry}(\operatorname{identmat}(a, n), i, j) = \left(\text{if }(i = j)\text{ then }\operatorname{one}(a)\text{ else }\operatorname{zero}(a)\right) \end{array}}}
\vnbprop{identmat-entry-diag}{\vnbstmt{\begin{array}{@{}l@{}} \forall a, n, i. \\ \quad (\operatorname{is\text{-}ring}(a) \wedge ((n \in \mathbb{N}) \wedge (i \in \operatorname{interval}(1, n)))) \Rightarrow  \\ \quad \quad \operatorname{entry}(\operatorname{identmat}(a, n), i, i) = \operatorname{one}(a) \end{array}}}
\vnbprop{identmat-entry-off}{\vnbstmt{\begin{array}{@{}l@{}} \forall a, n, i, j. \\ \quad (\operatorname{is\text{-}ring}(a) \wedge ((n \in \mathbb{N}) \wedge ((i \in \operatorname{interval}(1, n)) \wedge ((j \in \operatorname{interval}(1, n)) \wedge \neg (i = j))))) \Rightarrow  \\ \quad \quad \operatorname{entry}(\operatorname{identmat}(a, n), i, j) = \operatorname{zero}(a) \end{array}}}
\vnbprop{unitrow-entry-at}{\vnbstmt{\begin{array}{@{}l@{}} \forall a, n, i. \\ \quad (\operatorname{is\text{-}ring}(a) \wedge ((n \in \mathbb{N}) \wedge (i \in \operatorname{interval}(1, n)))) \Rightarrow  \\ \quad \quad \operatorname{entry}(\operatorname{unitrow}(a, n, i), 1, i) = \operatorname{one}(a) \end{array}}}
\vnbprop{unitrow-entry-off}{\vnbstmt{\begin{array}{@{}l@{}} \forall a, n, i, j. \\ \quad (\operatorname{is\text{-}ring}(a) \wedge ((n \in \mathbb{N}) \wedge ((i \in \operatorname{interval}(1, n)) \wedge ((j \in \operatorname{interval}(1, n)) \wedge \neg (j = i))))) \Rightarrow  \\ \quad \quad \operatorname{entry}(\operatorname{unitrow}(a, n, i), 1, j) = \operatorname{zero}(a) \end{array}}}
\vnbprop{entry-of-zeromat}{\vnbstmt{\begin{array}{@{}l@{}} \forall a, m, n, i, j. \\ \quad (\operatorname{is\text{-}ring}(a) \wedge ((m \in \mathbb{N}) \wedge ((n \in \mathbb{N}) \wedge ((i \in \operatorname{interval}(1, m)) \wedge (j \in \operatorname{interval}(1, n)))))) \Rightarrow  \\ \quad \quad \operatorname{entry}(\operatorname{zeromat}(a, m, n), i, j) = \operatorname{zero}(a) \end{array}}}
\vnbprop{elem-f-type}{\vnbstmt{\begin{array}{@{}l@{}} \forall a, n, k, l. \\ \quad (\operatorname{is\text{-}ring}(a) \wedge (n \in \mathbb{N})) \Rightarrow (\operatorname{elem\text{-}f}(a, n, k, l) \in \operatorname{mat}(n, n, \operatorname{carr}(a))) \end{array}}}
\vnbprop{elem-f-symmetric}{\vnbstmt{\begin{array}{@{}l@{}} \forall a, n, k, l. \\ \quad (\operatorname{is\text{-}ring}(a) \wedge (n \in \mathbb{N})) \Rightarrow (\operatorname{elem\text{-}f}(a, n, k, l) = \operatorname{elem\text{-}f}(a, n, l, k)) \end{array}}}
\vnbprop{entry-of-block}{\vnbstmt{\begin{array}{@{}l@{}} \forall p, k, l, i, j, m, n, x. \\ \quad k \in \mathbb{N} \wedge \\ \quad l \in \mathbb{N} \wedge \\ \quad p \in \operatorname{mat}(m, n, x) \wedge \\ \quad k \leq m \wedge \\ \quad l \leq n \wedge \\ \quad i \in \operatorname{interval}(1, k) \wedge \\ \quad j \in \operatorname{interval}(1, l) \Rightarrow \\ \quad \quad \operatorname{entry}(\operatorname{block}(p, k, l), i, j) = \operatorname{entry}(p, i, j) \end{array}}}
\vnbprop{entry-of-snoc-col}{\vnbstmt{\begin{array}{@{}l@{}} \forall w, n, v, i, x. \\ \quad ((n \in \mathbb{N}) \wedge ((w \in \operatorname{mat}(n, 1, x)) \wedge ((v \in x) \wedge (i \in \operatorname{interval}(1, n))))) \Rightarrow  \\ \quad \quad \operatorname{entry}(\operatorname{snoc\text{-}col}(w, n, v), i, 1) = \operatorname{entry}(w, i, 1) \end{array}}}
\vnbprop{snoc-col-last}{\vnbstmt{\begin{array}{@{}l@{}} \forall w, n, v, x. \\ \quad ((n \in \mathbb{N}) \wedge ((w \in \operatorname{mat}(n, 1, x)) \wedge (v \in x))) \Rightarrow  \\ \quad \quad \operatorname{entry}(\operatorname{snoc\text{-}col}(w, n, v), (n + 1), 1) = v \end{array}}}
\vnbprop{entry-of-snoc-row}{\vnbstmt{\begin{array}{@{}l@{}} \forall c, n, r, j, x. \\ \quad ((n \in \mathbb{N}) \wedge ((c \in \operatorname{mat}(1, n, x)) \wedge ((r \in x) \wedge (j \in \operatorname{interval}(1, n))))) \Rightarrow  \\ \quad \quad \operatorname{entry}(\operatorname{snoc\text{-}row}(c, n, r), 1, j) = \operatorname{entry}(c, 1, j) \end{array}}}
\vnbprop{snoc-row-last}{\vnbstmt{\begin{array}{@{}l@{}} \forall c, n, r, x. \\ \quad ((n \in \mathbb{N}) \wedge ((c \in \operatorname{mat}(1, n, x)) \wedge (r \in x))) \Rightarrow  \\ \quad \quad \operatorname{entry}(\operatorname{snoc\text{-}row}(c, n, r), 1, (n + 1)) = r \end{array}}}
\vnbprop{matadd-entry}{\vnbstmt{\begin{array}{@{}l@{}} \forall a, m, n, p, q.; \forall i \in \operatorname{interval}(1, m), j \in \operatorname{interval}(1, n). \\ \quad (\operatorname{is\text{-}ring}(a) \wedge ((p \in \operatorname{mat}(m, n, \operatorname{carr}(a))) \wedge (q \in \operatorname{mat}(m, n, \operatorname{carr}(a))))) \Rightarrow  \\ \quad \quad \operatorname{entry}(\operatorname{matadd}(a, p, q), i, j) = \operatorname{add}(a)(\operatorname{entry}(p, i, j), \operatorname{entry}(q, i, j)) \end{array}}}
\vnbprop{matscale-entry}{\vnbstmt{\begin{array}{@{}l@{}} \forall a, m, n, r, p.; \forall i \in \operatorname{interval}(1, m), j \in \operatorname{interval}(1, n). \\ \quad (\operatorname{is\text{-}ring}(a) \wedge ((r \in \operatorname{carr}(a)) \wedge (p \in \operatorname{mat}(m, n, \operatorname{carr}(a))))) \Rightarrow  \\ \quad \quad \operatorname{entry}(\operatorname{matscale}(a, r, p), i, j) = \operatorname{mul}(a)(r, \operatorname{entry}(p, i, j)) \end{array}}}
\vnbprop{entry-of-submat}{\vnbstmt{\begin{array}{@{}l@{}} \forall s, p, q, i, j, x. \\ \quad p \in \mathbb{N} \wedge \\ \quad q \in \mathbb{N} \wedge \\ \quad s \in \operatorname{mat}((p + 1), (q + 1), x) \wedge \\ \quad i \in \operatorname{interval}(1, p) \wedge \\ \quad j \in \operatorname{interval}(1, q) \Rightarrow \\ \quad \quad \operatorname{entry}(\operatorname{submat}(s, p, q), i, j) = \operatorname{entry}(s, (i + 1), (j + 1)) \end{array}}}

\vnbbookfile{theorem-library/rake-matrix-laws}{theorem-library/rake-matrix-laws}
\noindent{\small rake-matrix-laws.scm -- the eight matrix-algebra laws of `mat-ring-is-ring', PROVEN by ONE entrywise driver.}\par
\vnbprop{matneg-entry}{\vnbstmt{\begin{array}{@{}l@{}} \forall a, m, n, p.; \forall i \in \operatorname{interval}(1, m), j \in \operatorname{interval}(1, n). \\ \quad (\operatorname{is\text{-}ring}(a) \wedge (p \in \operatorname{mat}(m, n, \operatorname{carr}(a)))) \Rightarrow  \\ \quad \quad \operatorname{entry}(\operatorname{matneg}(a, p), i, j) = \operatorname{neg}(a)(\operatorname{entry}(p, i, j)) \end{array}}}
\vnbprop{matadd-comm}{\vnbstmt{\begin{array}{@{}l@{}} \forall a, m, n, p, q. \\ \quad (\operatorname{is\text{-}ring}(a) \wedge ((p \in \operatorname{mat}(m, n, \operatorname{carr}(a))) \wedge (q \in \operatorname{mat}(m, n, \operatorname{carr}(a))))) \Rightarrow  \\ \quad \quad \operatorname{matadd}(a, p, q) = \operatorname{matadd}(a, q, p) \end{array}}}
\vnbprop{matadd-assoc}{\vnbstmt{\begin{array}{@{}l@{}} \forall a, m, n, p, q, r. \\ \quad (\operatorname{is\text{-}ring}(a) \wedge ((p \in \operatorname{mat}(m, n, \operatorname{carr}(a))) \wedge ((q \in \operatorname{mat}(m, n, \operatorname{carr}(a))) \wedge (r \in \operatorname{mat}(m, n, \operatorname{carr}(a)))))) \Rightarrow  \\ \quad \quad \operatorname{matadd}(a, \operatorname{matadd}(a, p, q), r) = \operatorname{matadd}(a, p, \operatorname{matadd}(a, q, r)) \end{array}}}
\vnbprop{matadd-zero-left}{\vnbstmt{\begin{array}{@{}l@{}} \forall a, m, n.; \forall p \in \operatorname{mat}(m, n, \operatorname{carr}(a)). \\ \quad \operatorname{is\text{-}ring}(a) \Rightarrow (\operatorname{matadd}(a, \operatorname{zeromat}(a, m, n), p) = p) \end{array}}}
\vnbprop{matadd-zero-right}{\vnbstmt{\begin{array}{@{}l@{}} \forall a, m, n.; \forall p \in \operatorname{mat}(m, n, \operatorname{carr}(a)). \\ \quad \operatorname{is\text{-}ring}(a) \Rightarrow (\operatorname{matadd}(a, p, \operatorname{zeromat}(a, m, n)) = p) \end{array}}}
\vnbprop{matadd-neg-left}{\vnbstmt{\begin{array}{@{}l@{}} \forall a, m, n.; \forall p \in \operatorname{mat}(m, n, \operatorname{carr}(a)). \\ \quad \operatorname{is\text{-}ring}(a) \Rightarrow (\operatorname{matadd}(a, \operatorname{matneg}(a, p), p) = \operatorname{zeromat}(a, m, n)) \end{array}}}
\vnbprop{matadd-neg-right}{\vnbstmt{\begin{array}{@{}l@{}} \forall a, m, n.; \forall p \in \operatorname{mat}(m, n, \operatorname{carr}(a)). \\ \quad \operatorname{is\text{-}ring}(a) \Rightarrow (\operatorname{matadd}(a, p, \operatorname{matneg}(a, p)) = \operatorname{zeromat}(a, m, n)) \end{array}}}
\vnbprop{matmul-left-dist-guarded}{\vnbstmt{\begin{array}{@{}l@{}} \forall a, m, n, k, p, q, r. \\ \quad \operatorname{is\text{-}ring}(a) \wedge \\ \quad p \in \operatorname{mat}(m, n, \operatorname{carr}(a)) \wedge \\ \quad q \in \operatorname{mat}(n, k, \operatorname{carr}(a)) \wedge \\ \quad r \in \operatorname{mat}(n, k, \operatorname{carr}(a)) \wedge \\ \quad (n = 0) \Rightarrow ((m = 0) \vee (k = 0)) \Rightarrow \\ \quad \quad \operatorname{matmul}(a, p, \operatorname{matadd}(a, q, r)) = \operatorname{matadd}(a, \operatorname{matmul}(a, p, q), \operatorname{matmul}(a, p, r)) \end{array}}}
\vnbprop{matmul-right-dist-guarded}{\vnbstmt{\begin{array}{@{}l@{}} \forall a, m, n, k, p, q, r. \\ \quad \operatorname{is\text{-}ring}(a) \wedge \\ \quad p \in \operatorname{mat}(m, n, \operatorname{carr}(a)) \wedge \\ \quad q \in \operatorname{mat}(m, n, \operatorname{carr}(a)) \wedge \\ \quad r \in \operatorname{mat}(n, k, \operatorname{carr}(a)) \wedge \\ \quad (n = 0) \Rightarrow ((m = 0) \vee (k = 0)) \Rightarrow \\ \quad \quad \operatorname{matmul}(a, \operatorname{matadd}(a, p, q), r) = \operatorname{matadd}(a, \operatorname{matmul}(a, p, r), \operatorname{matmul}(a, q, r)) \end{array}}}

\vnbbookfile{theorem-library/mat-ring-proof}{theorem-library/mat-ring-proof}
\noindent{\small mat-ring-proof.scm -- the n-by-n matrices over a ring form a ring.}\par
\vnbprop{mat-ring-is-ring}{\vnbstmt{\forall a.; \forall n \in \mathbb{N}.\; (\operatorname{is\text{-}ring}(a) \Rightarrow \operatorname{is\text{-}ring}(\operatorname{mat\text{-}ring}(a, n)))}}

\vnbbookfile{theorem-library/matunit-shift-proof}{theorem-library/matunit-shift-proof}
\noindent{\small matunit-shift-proof.scm -- Lemma 3.3 (algebraic-numbers.pdf ch.3): the column-shift formula for right-multiplication by a matrix unit.}\par
\vnbprop{matunit-col-shift}{\vnbstmt{\begin{array}{@{}l@{}} \forall a, m, n, p, k, l.; \forall i \in \operatorname{interval}(1, m), c \in \operatorname{interval}(1, n). \\ \quad (\operatorname{is\text{-}ring}(a) \wedge ((p \in \operatorname{mat}(m, n, \operatorname{carr}(a))) \wedge ((k \in \operatorname{interval}(1, n)) \wedge (l \in \operatorname{interval}(1, n))))) \Rightarrow  \\ \quad \quad \operatorname{entry}(\operatorname{matmul}(a, p, \operatorname{matunit}(a, n, k, l)), i, c) = \left(\text{if }(c = l)\text{ then }\operatorname{entry}(p, i, k)\text{ else }\operatorname{zero}(a)\right) \end{array}}}

\vnbbookfile{theorem-library/elem-actions-proof}{theorem-library/elem-actions-proof}
\noindent{\small elem-actions-proof.scm -- Prop 3.5 (algebraic-numbers.pdf ch.3): the action of the elementary column matrices on a matrix by right-multiplication.  Each is the matmul-entry expansion + a finsum collapse (single-support for H, and for the off-l column of G; two-support for G's l-column), closed by an}\par
\vnbprop{elem-h-action}{\vnbstmt{\begin{array}{@{}l@{}} \forall a, m, n, p, r, k.; \forall i \in \operatorname{interval}(1, m), c \in \operatorname{interval}(1, n). \\ \quad (\operatorname{is\text{-}ring}(a) \wedge ((p \in \operatorname{mat}(m, n, \operatorname{carr}(a))) \wedge ((r \in \operatorname{carr}(a)) \wedge (k \in \operatorname{interval}(1, n))))) \Rightarrow  \\ \quad \quad \operatorname{entry}(\operatorname{matmul}(a, p, \operatorname{elem\text{-}h}(a, n, r, k)), i, c) \\ \quad \quad =\; \left(\text{if }(c = k)\text{ then }\operatorname{mul}(a)(\operatorname{entry}(p, i, k), r)\text{ else }\operatorname{entry}(p, i, c)\right) \end{array}}}
\vnbprop{elem-g-action}{\vnbstmt{\begin{array}{@{}l@{}} \forall a, m, n, p, r, k, l.; \forall i \in \operatorname{interval}(1, m), c \in \operatorname{interval}(1, n). \\ \quad \operatorname{is\text{-}ring}(a) \wedge \\ \quad p \in \operatorname{mat}(m, n, \operatorname{carr}(a)) \wedge \\ \quad r \in \operatorname{carr}(a) \wedge \\ \quad k \in \operatorname{interval}(1, n) \wedge \\ \quad l \in \operatorname{interval}(1, n) \wedge \\ \quad \neg (k = l) \Rightarrow \\ \quad \quad \operatorname{entry}(\operatorname{matmul}(a, p, \operatorname{elem\text{-}g}(a, n, r, k, l)), i, c) \\ \quad \quad =\; \left(\text{if }(c = l)\text{ then }\operatorname{add}(a)(\operatorname{entry}(p, i, l), \operatorname{mul}(a)(\operatorname{entry}(p, i, k), r))\text{ else }\operatorname{entry}(p, i, c)\right) \end{array}}}
\vnbprop{elem-f-action}{\vnbstmt{\begin{array}{@{}l@{}} \forall a, m, n, p, k, l.; \forall i \in \operatorname{interval}(1, m), c \in \operatorname{interval}(1, n). \\ \quad \operatorname{is\text{-}ring}(a) \wedge \\ \quad p \in \operatorname{mat}(m, n, \operatorname{carr}(a)) \wedge \\ \quad k \in \operatorname{interval}(1, n) \wedge \\ \quad l \in \operatorname{interval}(1, n) \wedge \\ \quad \neg (k = l) \Rightarrow \\ \quad \quad \operatorname{entry}(\operatorname{matmul}(a, p, \operatorname{elem\text{-}f}(a, n, k, l)), i, c) \\ \quad \quad =\; \left(\text{if }(c = k)\text{ then }\operatorname{entry}(p, i, l)\text{ else }\left(\text{if }(c = l)\text{ then }\operatorname{entry}(p, i, k)\text{ else }\operatorname{entry}(p, i, c)\right)\right) \end{array}}}

\vnbbookfile{theorem-library/elem-inverses-proof}{theorem-library/elem-inverses-proof}
\noindent{\small elem-inverses-proof.scm -- Cor 3.6 (algebraic-numbers.pdf ch.3): the elementary column matrices are invertible, with elementary inverses.  Each reduces, via matrix-entry-extensionality, to an entry identity whose LHS is the PROVEN Prop 3.5 action (with P = the would-be inverse) and whose RHS is entry-of-identmat;}\par
\vnbprop{elem-h-inverse}{\vnbstmt{\begin{array}{@{}l@{}} \forall a, n, r, s, k. \\ \quad \operatorname{is\text{-}ring}(a) \wedge \\ \quad n \in \mathbb{N} \wedge \\ \quad r \in \operatorname{carr}(a) \wedge \\ \quad s \in \operatorname{carr}(a) \wedge \\ \quad \operatorname{mul}(a)(r, s) = \operatorname{one}(a) \wedge \\ \quad k \in \operatorname{interval}(1, n) \Rightarrow \\ \quad \quad \operatorname{matmul}(a, \operatorname{elem\text{-}h}(a, n, r, k), \operatorname{elem\text{-}h}(a, n, s, k)) = \operatorname{identmat}(a, n) \end{array}}}
\vnbprop{elem-g-inverse}{\vnbstmt{\begin{array}{@{}l@{}} \forall a, n, r, k, l. \\ \quad \operatorname{is\text{-}ring}(a) \wedge \\ \quad n \in \mathbb{N} \wedge \\ \quad r \in \operatorname{carr}(a) \wedge \\ \quad k \in \operatorname{interval}(1, n) \wedge \\ \quad l \in \operatorname{interval}(1, n) \wedge \\ \quad \neg (k = l) \Rightarrow \\ \quad \quad \operatorname{matmul}(a, \operatorname{elem\text{-}g}(a, n, r, k, l), \operatorname{elem\text{-}g}(a, n, \operatorname{neg}(a)(r), k, l)) = \operatorname{identmat}(a, n) \end{array}}}
\vnbprop{elem-f-inverse}{\vnbstmt{\begin{array}{@{}l@{}} \forall a, n, k, l. \\ \quad (\operatorname{is\text{-}ring}(a) \wedge ((n \in \mathbb{N}) \wedge ((k \in \operatorname{interval}(1, n)) \wedge ((l \in \operatorname{interval}(1, n)) \wedge \neg (k = l))))) \Rightarrow  \\ \quad \quad \operatorname{matmul}(a, \operatorname{elem\text{-}f}(a, n, k, l), \operatorname{elem\text{-}f}(a, n, l, k)) = \operatorname{identmat}(a, n) \end{array}}}

\vnbbookfile{theorem-library/elem-row-actions-proof}{theorem-library/elem-row-actions-proof}
\noindent{\small elem-row-actions-proof.scm -- Prop 3.29 (row form, algebraic-numbers.pdf ch.3): the action of the elementary matrices by LEFT-multiplication (row operations). Each mirrors theorem-library/elem-actions-proof.scm with the elementary matrix as the LEFT factor of matmul-entry: (E.P)\_\{ic\} = FINSUM\_j E\_\{ij\}.P\_\{jc\}, so a}\par
\vnbprop{elem-h-row-action}{\vnbstmt{\begin{array}{@{}l@{}} \forall a, m, n, p, r, k.; \forall i \in \operatorname{interval}(1, m), c \in \operatorname{interval}(1, n). \\ \quad (\operatorname{is\text{-}ring}(a) \wedge ((p \in \operatorname{mat}(m, n, \operatorname{carr}(a))) \wedge ((r \in \operatorname{carr}(a)) \wedge (k \in \operatorname{interval}(1, m))))) \Rightarrow  \\ \quad \quad \operatorname{entry}(\operatorname{matmul}(a, \operatorname{elem\text{-}h}(a, m, r, k), p), i, c) \\ \quad \quad =\; \left(\text{if }(i = k)\text{ then }\operatorname{mul}(a)(r, \operatorname{entry}(p, k, c))\text{ else }\operatorname{entry}(p, i, c)\right) \end{array}}}
\vnbprop{elem-g-row-action}{\vnbstmt{\begin{array}{@{}l@{}} \forall a, m, n, p, r, k, l.; \forall i \in \operatorname{interval}(1, m), c \in \operatorname{interval}(1, n). \\ \quad \operatorname{is\text{-}ring}(a) \wedge \\ \quad p \in \operatorname{mat}(m, n, \operatorname{carr}(a)) \wedge \\ \quad r \in \operatorname{carr}(a) \wedge \\ \quad k \in \operatorname{interval}(1, m) \wedge \\ \quad l \in \operatorname{interval}(1, m) \wedge \\ \quad \neg (k = l) \Rightarrow \\ \quad \quad \operatorname{entry}(\operatorname{matmul}(a, \operatorname{elem\text{-}g}(a, m, r, k, l), p), i, c) \\ \quad \quad =\; \left(\text{if }(i = k)\text{ then }\operatorname{add}(a)(\operatorname{entry}(p, k, c), \operatorname{mul}(a)(r, \operatorname{entry}(p, l, c)))\text{ else }\operatorname{entry}(p, i, c)\right) \end{array}}}
\vnbprop{elem-f-row-action}{\vnbstmt{\begin{array}{@{}l@{}} \forall a, m, n, p, k, l.; \forall i \in \operatorname{interval}(1, m), c \in \operatorname{interval}(1, n). \\ \quad \operatorname{is\text{-}ring}(a) \wedge \\ \quad p \in \operatorname{mat}(m, n, \operatorname{carr}(a)) \wedge \\ \quad k \in \operatorname{interval}(1, m) \wedge \\ \quad l \in \operatorname{interval}(1, m) \wedge \\ \quad \neg (k = l) \Rightarrow \\ \quad \quad \operatorname{entry}(\operatorname{matmul}(a, \operatorname{elem\text{-}f}(a, m, k, l), p), i, c) \\ \quad \quad =\; \left(\text{if }(i = k)\text{ then }\operatorname{entry}(p, l, c)\text{ else }\left(\text{if }(i = l)\text{ then }\operatorname{entry}(p, k, c)\text{ else }\operatorname{entry}(p, i, c)\right)\right) \end{array}}}

\vnbbookfile{theorem-library/det-rows}{theorem-library/det-rows}
\noindent{\small det-rows.scm -- the determinant is LINEAR in each row and ALTERNATING; the interchange of two rows, the elementary row operations, and the expansion along any row (batch 25-C, 2026-09-24).  Linear-algebra 3.18 (row clauses), the row form of 3.19, and 3.22 (row half).  All over a COMMUTATIVE ring.}\par
\vnbprop{dtr-ring-lin1}{\vnbstmt{\begin{array}{@{}l@{}} \forall r, s, a, b, c, d. \\ \quad \operatorname{is\text{-}commutative\text{-}ring}(r) \wedge \\ \quad s \in \operatorname{carr}(r) \wedge \\ \quad a \in \operatorname{carr}(r) \wedge \\ \quad b \in \operatorname{carr}(r) \wedge \\ \quad c \in \operatorname{carr}(r) \wedge \\ \quad d \in \operatorname{carr}(r) \Rightarrow \\ \quad \quad \operatorname{mul}(r)(s, \operatorname{mul}(r)(\operatorname{add}(r)(\operatorname{mul}(r)(c, a), b), d)) \\ \quad \quad =\; \operatorname{add}(r)(\operatorname{mul}(r)(c, \operatorname{mul}(r)(s, \operatorname{mul}(r)(a, d))), \operatorname{mul}(r)(s, \operatorname{mul}(r)(b, d))) \end{array}}}
\vnbprop{dtr-ring-lin2}{\vnbstmt{\begin{array}{@{}l@{}} \forall r, s, a, c, x, y. \\ \quad \operatorname{is\text{-}commutative\text{-}ring}(r) \wedge \\ \quad s \in \operatorname{carr}(r) \wedge \\ \quad a \in \operatorname{carr}(r) \wedge \\ \quad c \in \operatorname{carr}(r) \wedge \\ \quad x \in \operatorname{carr}(r) \wedge \\ \quad y \in \operatorname{carr}(r) \Rightarrow \\ \quad \quad \operatorname{mul}(r)(s, \operatorname{mul}(r)(a, \operatorname{add}(r)(\operatorname{mul}(r)(c, x), y))) \\ \quad \quad =\; \operatorname{add}(r)(\operatorname{mul}(r)(c, \operatorname{mul}(r)(s, \operatorname{mul}(r)(a, x))), \operatorname{mul}(r)(s, \operatorname{mul}(r)(a, y))) \end{array}}}
\vnbprop{dtr-finsum-lin}{\vnbstmt{\begin{array}{@{}l@{}} \forall r, s, c, f, g, h. \\ \quad \operatorname{is\text{-}ring}(r) \wedge \\ \quad s \in \mathbf{Set} \wedge \\ \quad |s| \in \mathbb{N} \wedge \\ \quad c \in \operatorname{carr}(r) \wedge \\ \quad \forall z \in s.\; (f(z) \in \operatorname{carr}(r)) \wedge \\ \quad \forall z \in s.\; (g(z) \in \operatorname{carr}(r)) \wedge \\ \quad \forall z \in s.\; (h(z) = \operatorname{add}(r)(\operatorname{mul}(r)(c, f(z)), g(z))) \Rightarrow \\ \quad \quad \operatorname{finsum}(\operatorname{ring\text{-}additive\text{-}ag}(r), h, s) \\ \quad \quad =\; \operatorname{add}(r)(\operatorname{mul}(r)(c, \operatorname{finsum}(\operatorname{ring\text{-}additive\text{-}ag}(r), f, s)), \operatorname{finsum}(\operatorname{ring\text{-}additive\text{-}ag}(r), g, s)) \end{array}}}
\vnbprop{dtr-succ-in}{\vnbstmt{\begin{array}{@{}l@{}} \forall n, u. \\ \quad ((n \in \mathbb{N}) \wedge (u \in \operatorname{interval}(1, n))) \Rightarrow ((u + 1) \in \operatorname{interval}(1, (n + 1))) \end{array}}}
\vnbprop{dtr-skip-in}{\vnbstmt{\begin{array}{@{}l@{}} \forall n, j, v. \\ \quad ((n \in \mathbb{N}) \wedge (v \in \operatorname{interval}(1, n))) \Rightarrow  \\ \quad \quad \left(\text{if }(v < j)\text{ then }v\text{ else }(v + 1)\right) \in \operatorname{interval}(1, (n + 1)) \end{array}}}
\vnbprop{dtr-minor1-entry}{\vnbstmt{\begin{array}{@{}l@{}} \forall r, n, s, j, u, v. \\ \quad \operatorname{is\text{-}ring}(r) \wedge \\ \quad n \in \mathbb{N} \wedge \\ \quad s \in \operatorname{mat}((n + 1), (n + 1), \operatorname{carr}(r)) \wedge \\ \quad u \in \operatorname{interval}(1, n) \wedge \\ \quad v \in \operatorname{interval}(1, n) \Rightarrow \\ \quad \quad \operatorname{entry}(\operatorname{minor}(s, 1, j, n), u, v) = \operatorname{entry}(s, (u + 1), \left(\text{if }(v < j)\text{ then }v\text{ else }(v + 1)\right)) \end{array}}}
\vnbprop{dtr-sign-type}{\vnbstmt{\begin{array}{@{}l@{}} \forall r, j. \\ \quad (\operatorname{is\text{-}ring}(r) \wedge (j \in \mathbb{N})) \Rightarrow  \\ \quad \quad \operatorname{mpow}(\operatorname{ring\text{-}multiplicative\text{-}monoid}(r), \operatorname{neg}(r)(\operatorname{one}(r)), j) \in \operatorname{carr}(r) \end{array}}}
\vnbprop{dtr-minor1-congr}{\vnbstmt{\begin{array}{@{}l@{}} \forall r, n, x, y, p, j. \\ \quad \operatorname{is\text{-}ring}(r) \wedge \\ \quad n \in \mathbb{N} \wedge \\ \quad x \in \operatorname{mat}((n + 1), (n + 1), \operatorname{carr}(r)) \wedge \\ \quad y \in \operatorname{mat}((n + 1), (n + 1), \operatorname{carr}(r)) \wedge \\ \quad j \in \mathbb{N} \wedge \\ \quad p = 1 \wedge \\ \quad \forall i, k \in \operatorname{interval}(1, (n + 1)). \\ \quad \quad \neg (i = p) \Rightarrow (\operatorname{entry}(x, i, k) = \operatorname{entry}(y, i, k)) \Rightarrow \\ \quad \quad \operatorname{minor}(x, 1, j, n) = \operatorname{minor}(y, 1, j, n) \end{array}}}
\vnbprop{det-row-linear-ind}{\vnbstmt{\begin{array}{@{}l@{}} \forall n \in \mathbb{N}, r, a, b, c, p, \mathit{dsc}. \\ \quad \operatorname{is\text{-}commutative\text{-}ring}(r) \wedge \\ \quad a \in \operatorname{mat}(n, n, \operatorname{carr}(r)) \wedge \\ \quad b \in \operatorname{mat}(n, n, \operatorname{carr}(r)) \wedge \\ \quad c \in \operatorname{mat}(n, n, \operatorname{carr}(r)) \wedge \\ \quad p \in \operatorname{interval}(1, n) \wedge \\ \quad \mathit{dsc} \in \operatorname{carr}(r) \wedge \\ \quad \forall i, k \in \operatorname{interval}(1, n). \\ \quad \quad \neg (i = p) \Rightarrow (\operatorname{entry}(b, i, k) = \operatorname{entry}(a, i, k)) \wedge \\ \quad \forall i, k \in \operatorname{interval}(1, n). \\ \quad \quad \neg (i = p) \Rightarrow (\operatorname{entry}(c, i, k) = \operatorname{entry}(a, i, k)) \wedge \\ \quad \forall k \in \operatorname{interval}(1, n). \\ \quad \quad \operatorname{entry}(c, p, k) = \operatorname{add}(r)(\operatorname{mul}(r)(\mathit{dsc}, \operatorname{entry}(a, p, k)), \operatorname{entry}(b, p, k)) \Rightarrow \\ \quad \quad \operatorname{det}(r, n, c) = \operatorname{add}(r)(\operatorname{mul}(r)(\mathit{dsc}, \operatorname{det}(r, n, a)), \operatorname{det}(r, n, b)) \end{array}}}
\vnbprop{det-row-linear}{\vnbstmt{\begin{array}{@{}l@{}} \forall r, n, a, b, c, p, d. \\ \quad \operatorname{is\text{-}commutative\text{-}ring}(r) \wedge \\ \quad n \in \mathbb{N} \wedge \\ \quad a \in \operatorname{mat}(n, n, \operatorname{carr}(r)) \wedge \\ \quad b \in \operatorname{mat}(n, n, \operatorname{carr}(r)) \wedge \\ \quad c \in \operatorname{mat}(n, n, \operatorname{carr}(r)) \wedge \\ \quad p \in \operatorname{interval}(1, n) \wedge \\ \quad d \in \operatorname{carr}(r) \wedge \\ \quad \forall i, k \in \operatorname{interval}(1, n). \\ \quad \quad \neg (i = p) \Rightarrow (\operatorname{entry}(b, i, k) = \operatorname{entry}(a, i, k)) \wedge \\ \quad \forall i, k \in \operatorname{interval}(1, n). \\ \quad \quad \neg (i = p) \Rightarrow (\operatorname{entry}(c, i, k) = \operatorname{entry}(a, i, k)) \wedge \\ \quad \forall k \in \operatorname{interval}(1, n). \\ \quad \quad \operatorname{entry}(c, p, k) = \operatorname{add}(r)(\operatorname{mul}(r)(d, \operatorname{entry}(a, p, k)), \operatorname{entry}(b, p, k)) \Rightarrow \\ \quad \quad \operatorname{det}(r, n, c) = \operatorname{add}(r)(\operatorname{mul}(r)(d, \operatorname{det}(r, n, a)), \operatorname{det}(r, n, b)) \end{array}}}
\vnbprop{dtr-sign-succ}{\vnbstmt{\begin{array}{@{}l@{}} \forall r, x. \\ \quad (\operatorname{is\text{-}ring}(r) \wedge (x \in \mathbb{N})) \Rightarrow  \\ \quad \quad \operatorname{mpow}(\operatorname{ring\text{-}multiplicative\text{-}monoid}(r), \operatorname{neg}(r)(\operatorname{one}(r)), (x + 1)) \\ \quad \quad =\; \operatorname{mul}(r)(\operatorname{neg}(r)(\operatorname{one}(r)), \operatorname{mpow}(\operatorname{ring\text{-}multiplicative\text{-}monoid}(r), \operatorname{neg}(r)(\operatorname{one}(r)), x)) \end{array}}}
\vnbprop{dtr-minor-minor}{\vnbstmt{\begin{array}{@{}l@{}} \forall r, n, a, j, k. \\ \quad \operatorname{is\text{-}ring}(r) \wedge \\ \quad n \in \mathbb{N} \wedge \\ \quad a \in \operatorname{mat}(((n + 1) + 1), ((n + 1) + 1), \operatorname{carr}(r)) \wedge \\ \quad j \in \mathbb{N} \wedge \\ \quad k \in \mathbb{N} \wedge \\ \quad 1 \leq j \wedge \\ \quad j \leq k \wedge \\ \quad k \leq (n + 1) \Rightarrow \\ \quad \quad \operatorname{minor}(\operatorname{minor}(a, 1, j, (n + 1)), 1, k, n) = \operatorname{minor}(\operatorname{minor}(a, 1, (k + 1), (n + 1)), 1, j, n) \end{array}}}
\vnbprop{dtr-ring-zero-add}{\vnbstmt{\begin{array}{@{}l@{}} \forall r, y, z. \\ \quad (\operatorname{is\text{-}commutative\text{-}ring}(r) \wedge ((y \in \operatorname{carr}(r)) \wedge (z \in \operatorname{carr}(r)))) \Rightarrow  \\ \quad \quad \operatorname{add}(r)(\operatorname{add}(r)(\operatorname{zero}(r), y), z) = \operatorname{add}(r)(y, z) \end{array}}}
\vnbprop{dtr-dsum-anti}{\vnbstmt{\begin{array}{@{}l@{}} \forall n \in \mathbb{N}, r, f. \\ \quad \operatorname{is\text{-}commutative\text{-}ring}(r) \wedge \\ \quad \forall j \in \operatorname{interval}(1, (n + 1)), k \in \operatorname{interval}(1, n).\; (f(j, k) \in \operatorname{carr}(r)) \wedge \\ \quad \forall j, k \in \mathbb{N}. \\ \quad \quad ((1 \leq j) \wedge ((j \leq k) \wedge (k \leq n))) \Rightarrow (f((k + 1), j) = \operatorname{neg}(r)(f(j, k))) \Rightarrow \\ \quad \quad \operatorname{finsum}(\operatorname{ring\text{-}additive\text{-}ag}(r), \\ \quad \quad \quad (\lambda a \in \operatorname{interval}(1, (n + 1)).\; \\ \quad \quad \quad \quad \operatorname{finsum}(\operatorname{ring\text{-}additive\text{-}ag}(r), (\lambda b \in \operatorname{interval}(1, n).\; f(a, b)), \operatorname{interval}(1, n))), \\ \quad \quad \quad \operatorname{interval}(1, (n + 1))) \\ \quad \quad =\; \operatorname{zero}(r) \end{array}}}
\vnbprop{dtr-ring-anti}{\vnbstmt{\begin{array}{@{}l@{}} \forall r, s, t, a, b, d. \\ \quad \operatorname{is\text{-}commutative\text{-}ring}(r) \wedge \\ \quad s \in \operatorname{carr}(r) \wedge \\ \quad t \in \operatorname{carr}(r) \wedge \\ \quad a \in \operatorname{carr}(r) \wedge \\ \quad b \in \operatorname{carr}(r) \wedge \\ \quad d \in \operatorname{carr}(r) \Rightarrow \\ \quad \quad \operatorname{mul}(r)(\operatorname{mul}(r)(\operatorname{neg}(r)(\operatorname{one}(r)), t), \operatorname{mul}(r)(b, \operatorname{mul}(r)(s, \operatorname{mul}(r)(a, d)))) \\ \quad \quad =\; \operatorname{neg}(r)(\operatorname{mul}(r)(s, \operatorname{mul}(r)(a, \operatorname{mul}(r)(t, \operatorname{mul}(r)(b, d))))) \end{array}}}
\vnbprop{dtr-det-rows12}{\vnbstmt{\begin{array}{@{}l@{}} \forall n, r, a. \\ \quad n \in \mathbb{N} \wedge \\ \quad \operatorname{is\text{-}commutative\text{-}ring}(r) \wedge \\ \quad a \in \operatorname{mat}(((n + 1) + 1), ((n + 1) + 1), \operatorname{carr}(r)) \wedge \\ \quad \forall k \in \operatorname{interval}(1, ((n + 1) + 1)).\; (\operatorname{entry}(a, 1, k) = \operatorname{entry}(a, 2, k)) \Rightarrow \\ \quad \quad \operatorname{det}(r, ((n + 1) + 1), a) = \operatorname{zero}(r) \end{array}}}
\vnbprop{dtr-det-rows12b}{\vnbstmt{\begin{array}{@{}l@{}} \forall m, n, r, a. \\ \quad n \in \mathbb{N} \wedge \\ \quad m = ((n + 1) + 1) \wedge \\ \quad \operatorname{is\text{-}commutative\text{-}ring}(r) \wedge \\ \quad a \in \operatorname{mat}(m, m, \operatorname{carr}(r)) \wedge \\ \quad \forall k \in \operatorname{interval}(1, m).\; (\operatorname{entry}(a, 1, k) = \operatorname{entry}(a, 2, k)) \Rightarrow \\ \quad \quad \operatorname{det}(r, m, a) = \operatorname{zero}(r) \end{array}}}
\vnbprop{dtr-alt-low}{\vnbstmt{\begin{array}{@{}l@{}} \forall n \in \mathbb{N}, r, a, p, q. \\ \quad \forall r, a, p, q. \\ \quad \quad \operatorname{is\text{-}commutative\text{-}ring}(r) \wedge \\ \quad \quad a \in \operatorname{mat}(n, n, \operatorname{carr}(r)) \wedge \\ \quad \quad p \in \operatorname{interval}(1, n) \wedge \\ \quad \quad q \in \operatorname{interval}(1, n) \wedge \\ \quad \quad \neg (p = q) \wedge \\ \quad \quad \forall k \in \operatorname{interval}(1, n).\; (\operatorname{entry}(a, p, k) = \operatorname{entry}(a, q, k)) \Rightarrow \\ \quad \quad \quad \operatorname{det}(r, n, a) = \operatorname{zero}(r) \wedge \\ \quad \operatorname{is\text{-}commutative\text{-}ring}(r) \wedge \\ \quad a \in \operatorname{mat}((n + 1), (n + 1), \operatorname{carr}(r)) \wedge \\ \quad p \in \operatorname{interval}(1, (n + 1)) \wedge \\ \quad q \in \operatorname{interval}(1, (n + 1)) \wedge \\ \quad \neg (p = q) \wedge \\ \quad \neg (p = 1) \wedge \\ \quad \neg (q = 1) \wedge \\ \quad \forall k \in \operatorname{interval}(1, (n + 1)).\; (\operatorname{entry}(a, p, k) = \operatorname{entry}(a, q, k)) \Rightarrow \\ \quad \quad \operatorname{det}(r, (n + 1), a) = \operatorname{zero}(r) \end{array}}}
\vnbprop{dtr-ring-final}{\vnbstmt{\begin{array}{@{}l@{}} \forall r, t, u_{3}, u_{4}, u_{5}, u_{6}, u_{7}, u_{8}. \\ \quad \operatorname{is\text{-}commutative\text{-}ring}(r) \wedge \\ \quad t \in \operatorname{carr}(r) \wedge \\ \quad u_{3} \in \operatorname{carr}(r) \wedge \\ \quad u_{4} \in \operatorname{carr}(r) \wedge \\ \quad u_{5} \in \operatorname{carr}(r) \wedge \\ \quad u_{6} \in \operatorname{carr}(r) \wedge \\ \quad u_{7} \in \operatorname{carr}(r) \wedge \\ \quad u_{8} \in \operatorname{carr}(r) \wedge \\ \quad u_{3} = \operatorname{add}(r)(\operatorname{mul}(r)(\operatorname{one}(r), u_{4}), u_{5}) \wedge \\ \quad u_{4} = \operatorname{add}(r)(\operatorname{mul}(r)(\operatorname{one}(r), u_{6}), t) \wedge \\ \quad u_{5} = \operatorname{add}(r)(\operatorname{mul}(r)(\operatorname{one}(r), u_{7}), u_{8}) \wedge \\ \quad u_{3} = \operatorname{zero}(r) \wedge \\ \quad u_{6} = \operatorname{zero}(r) \wedge \\ \quad u_{7} = \operatorname{zero}(r) \wedge \\ \quad u_{8} = \operatorname{zero}(r) \Rightarrow \\ \quad \quad t = \operatorname{zero}(r) \end{array}}}
\vnbprop{dtr-alt-row1}{\vnbstmt{\begin{array}{@{}l@{}} \forall n \in \mathbb{N}, r, a, q. \\ \quad \forall r, a, p, q. \\ \quad \quad \operatorname{is\text{-}commutative\text{-}ring}(r) \wedge \\ \quad \quad a \in \operatorname{mat}(n, n, \operatorname{carr}(r)) \wedge \\ \quad \quad p \in \operatorname{interval}(1, n) \wedge \\ \quad \quad q \in \operatorname{interval}(1, n) \wedge \\ \quad \quad \neg (p = q) \wedge \\ \quad \quad \forall k \in \operatorname{interval}(1, n).\; (\operatorname{entry}(a, p, k) = \operatorname{entry}(a, q, k)) \Rightarrow \\ \quad \quad \quad \operatorname{det}(r, n, a) = \operatorname{zero}(r) \wedge \\ \quad \operatorname{is\text{-}commutative\text{-}ring}(r) \wedge \\ \quad a \in \operatorname{mat}((n + 1), (n + 1), \operatorname{carr}(r)) \wedge \\ \quad q \in \operatorname{interval}(1, (n + 1)) \wedge \\ \quad \neg (q = 1) \wedge \\ \quad \forall k \in \operatorname{interval}(1, (n + 1)).\; (\operatorname{entry}(a, 1, k) = \operatorname{entry}(a, q, k)) \Rightarrow \\ \quad \quad \operatorname{det}(r, (n + 1), a) = \operatorname{zero}(r) \end{array}}}
\vnbprop{det-alternating-ind}{\vnbstmt{\begin{array}{@{}l@{}} \forall n \in \mathbb{N}, r, a, p, q. \\ \quad \operatorname{is\text{-}commutative\text{-}ring}(r) \wedge \\ \quad a \in \operatorname{mat}(n, n, \operatorname{carr}(r)) \wedge \\ \quad p \in \operatorname{interval}(1, n) \wedge \\ \quad q \in \operatorname{interval}(1, n) \wedge \\ \quad \neg (p = q) \wedge \\ \quad \forall k \in \operatorname{interval}(1, n).\; (\operatorname{entry}(a, p, k) = \operatorname{entry}(a, q, k)) \Rightarrow \\ \quad \quad \operatorname{det}(r, n, a) = \operatorname{zero}(r) \end{array}}}
\vnbprop{det-alternating-rows}{\vnbstmt{\begin{array}{@{}l@{}} \forall r, n, a, p, q. \\ \quad \operatorname{is\text{-}commutative\text{-}ring}(r) \wedge \\ \quad n \in \mathbb{N} \wedge \\ \quad a \in \operatorname{mat}(n, n, \operatorname{carr}(r)) \wedge \\ \quad p \in \operatorname{interval}(1, n) \wedge \\ \quad q \in \operatorname{interval}(1, n) \wedge \\ \quad \neg (p = q) \wedge \\ \quad \forall k \in \operatorname{interval}(1, n).\; (\operatorname{entry}(a, p, k) = \operatorname{entry}(a, q, k)) \Rightarrow \\ \quad \quad \operatorname{det}(r, n, a) = \operatorname{zero}(r) \end{array}}}
\vnbprop{dtr-ring-swap}{\vnbstmt{\begin{array}{@{}l@{}} \forall r, a, b, u_{1}, u_{2}, u_{3}, u_{4}, u_{5}. \\ \quad \operatorname{is\text{-}commutative\text{-}ring}(r) \wedge \\ \quad a \in \operatorname{carr}(r) \wedge \\ \quad b \in \operatorname{carr}(r) \wedge \\ \quad u_{1} \in \operatorname{carr}(r) \wedge \\ \quad u_{2} \in \operatorname{carr}(r) \wedge \\ \quad u_{3} \in \operatorname{carr}(r) \wedge \\ \quad u_{4} \in \operatorname{carr}(r) \wedge \\ \quad u_{5} \in \operatorname{carr}(r) \wedge \\ \quad u_{1} = \operatorname{add}(r)(\operatorname{mul}(r)(\operatorname{one}(r), u_{2}), u_{3}) \wedge \\ \quad u_{2} = \operatorname{add}(r)(\operatorname{mul}(r)(\operatorname{one}(r), u_{4}), a) \wedge \\ \quad u_{3} = \operatorname{add}(r)(\operatorname{mul}(r)(\operatorname{one}(r), b), u_{5}) \wedge \\ \quad u_{1} = \operatorname{zero}(r) \wedge \\ \quad u_{4} = \operatorname{zero}(r) \wedge \\ \quad u_{5} = \operatorname{zero}(r) \Rightarrow \\ \quad \quad b = \operatorname{neg}(r)(a) \end{array}}}
\vnbprop{dtr-swap-s}{\vnbstmt{\begin{array}{@{}l@{}} \forall n, r, a, b, p, q. \\ \quad n \in \mathbb{N} \wedge \\ \quad \operatorname{is\text{-}commutative\text{-}ring}(r) \wedge \\ \quad a \in \operatorname{mat}((n + 1), (n + 1), \operatorname{carr}(r)) \wedge \\ \quad b \in \operatorname{mat}((n + 1), (n + 1), \operatorname{carr}(r)) \wedge \\ \quad p \in \operatorname{interval}(1, (n + 1)) \wedge \\ \quad q \in \operatorname{interval}(1, (n + 1)) \wedge \\ \quad \neg (p = q) \wedge \\ \quad \forall i, k \in \operatorname{interval}(1, (n + 1)). \\ \quad \quad (\neg (i = p) \wedge \neg (i = q)) \Rightarrow (\operatorname{entry}(b, i, k) = \operatorname{entry}(a, i, k)) \wedge \\ \quad \forall k \in \operatorname{interval}(1, (n + 1)).\; (\operatorname{entry}(b, p, k) = \operatorname{entry}(a, q, k)) \wedge \\ \quad \forall k \in \operatorname{interval}(1, (n + 1)).\; (\operatorname{entry}(b, q, k) = \operatorname{entry}(a, p, k)) \Rightarrow \\ \quad \quad \operatorname{det}(r, (n + 1), b) = \operatorname{neg}(r)(\operatorname{det}(r, (n + 1), a)) \end{array}}}
\vnbprop{det-swap-rows}{\vnbstmt{\begin{array}{@{}l@{}} \forall r, n, a, b, p, q. \\ \quad \operatorname{is\text{-}commutative\text{-}ring}(r) \wedge \\ \quad n \in \mathbb{N} \wedge \\ \quad a \in \operatorname{mat}(n, n, \operatorname{carr}(r)) \wedge \\ \quad b \in \operatorname{mat}(n, n, \operatorname{carr}(r)) \wedge \\ \quad p \in \operatorname{interval}(1, n) \wedge \\ \quad q \in \operatorname{interval}(1, n) \wedge \\ \quad \neg (p = q) \wedge \\ \quad \forall i, k \in \operatorname{interval}(1, n). \\ \quad \quad (\neg (i = p) \wedge \neg (i = q)) \Rightarrow (\operatorname{entry}(b, i, k) = \operatorname{entry}(a, i, k)) \wedge \\ \quad \forall k \in \operatorname{interval}(1, n).\; (\operatorname{entry}(b, p, k) = \operatorname{entry}(a, q, k)) \wedge \\ \quad \forall k \in \operatorname{interval}(1, n).\; (\operatorname{entry}(b, q, k) = \operatorname{entry}(a, p, k)) \Rightarrow \\ \quad \quad \operatorname{det}(r, n, b) = \operatorname{neg}(r)(\operatorname{det}(r, n, a)) \end{array}}}
\vnbprop{dtr-ring-negself}{\vnbstmt{\begin{array}{@{}l@{}} \forall r.; \forall u \in \operatorname{carr}(r). \\ \quad (\operatorname{is\text{-}commutative\text{-}ring}(r) \wedge (u = \operatorname{add}(r)(\operatorname{mul}(r)(\operatorname{neg}(r)(\operatorname{one}(r)), u), u))) \Rightarrow  \\ \quad \quad u = \operatorname{zero}(r) \end{array}}}
\vnbprop{dtr-zero-row-s}{\vnbstmt{\begin{array}{@{}l@{}} \forall n, r, a, k. \\ \quad n \in \mathbb{N} \wedge \\ \quad \operatorname{is\text{-}commutative\text{-}ring}(r) \wedge \\ \quad a \in \operatorname{mat}((n + 1), (n + 1), \operatorname{carr}(r)) \wedge \\ \quad k \in \operatorname{interval}(1, (n + 1)) \wedge \\ \quad \forall \mathit{drk} \in \operatorname{interval}(1, (n + 1)).\; (\operatorname{entry}(a, k, \mathit{drk}) = \operatorname{zero}(r)) \Rightarrow \\ \quad \quad \operatorname{det}(r, (n + 1), a) = \operatorname{zero}(r) \end{array}}}
\vnbprop{dtr-elem-h-s}{\vnbstmt{\begin{array}{@{}l@{}} \forall n, r, a, \mathit{dsr}, k. \\ \quad n \in \mathbb{N} \wedge \\ \quad \operatorname{is\text{-}commutative\text{-}ring}(r) \wedge \\ \quad a \in \operatorname{mat}((n + 1), (n + 1), \operatorname{carr}(r)) \wedge \\ \quad \mathit{dsr} \in \operatorname{carr}(r) \wedge \\ \quad k \in \operatorname{interval}(1, (n + 1)) \Rightarrow \\ \quad \quad \operatorname{det}(r, (n + 1), \operatorname{matmul}(r, \operatorname{elem\text{-}h}(r, (n + 1), \mathit{dsr}, k), a)) = \operatorname{mul}(r)(\mathit{dsr}, \operatorname{det}(r, (n + 1), a)) \end{array}}}
\vnbprop{dtr-elem-g-s}{\vnbstmt{\begin{array}{@{}l@{}} \forall n, r, a, \mathit{dsr}, k, l. \\ \quad n \in \mathbb{N} \wedge \\ \quad \operatorname{is\text{-}commutative\text{-}ring}(r) \wedge \\ \quad a \in \operatorname{mat}((n + 1), (n + 1), \operatorname{carr}(r)) \wedge \\ \quad \mathit{dsr} \in \operatorname{carr}(r) \wedge \\ \quad k \in \operatorname{interval}(1, (n + 1)) \wedge \\ \quad l \in \operatorname{interval}(1, (n + 1)) \wedge \\ \quad \neg (k = l) \Rightarrow \\ \quad \quad \operatorname{det}(r, (n + 1), \operatorname{matmul}(r, \operatorname{elem\text{-}g}(r, (n + 1), \mathit{dsr}, k, l), a)) = \operatorname{det}(r, (n + 1), a) \end{array}}}
\vnbprop{dtr-elem-f-s}{\vnbstmt{\begin{array}{@{}l@{}} \forall n, r, a, k, l. \\ \quad n \in \mathbb{N} \wedge \\ \quad \operatorname{is\text{-}commutative\text{-}ring}(r) \wedge \\ \quad a \in \operatorname{mat}((n + 1), (n + 1), \operatorname{carr}(r)) \wedge \\ \quad k \in \operatorname{interval}(1, (n + 1)) \wedge \\ \quad l \in \operatorname{interval}(1, (n + 1)) \wedge \\ \quad \neg (k = l) \Rightarrow \\ \quad \quad \operatorname{det}(r, (n + 1), \operatorname{matmul}(r, \operatorname{elem\text{-}f}(r, (n + 1), k, l), a)) = \operatorname{neg}(r)(\operatorname{det}(r, (n + 1), a)) \end{array}}}
\vnbprop{det-elem-h-row}{\vnbstmt{\begin{array}{@{}l@{}} \forall r, n, a, c, k. \\ \quad \operatorname{is\text{-}commutative\text{-}ring}(r) \wedge \\ \quad n \in \mathbb{N} \wedge \\ \quad a \in \operatorname{mat}(n, n, \operatorname{carr}(r)) \wedge \\ \quad c \in \operatorname{carr}(r) \wedge \\ \quad k \in \operatorname{interval}(1, n) \Rightarrow \\ \quad \quad \operatorname{det}(r, n, \operatorname{matmul}(r, \operatorname{elem\text{-}h}(r, n, c, k), a)) = \operatorname{mul}(r)(c, \operatorname{det}(r, n, a)) \end{array}}}
\vnbprop{det-elem-g-row}{\vnbstmt{\begin{array}{@{}l@{}} \forall r, n, a, c, k, l. \\ \quad \operatorname{is\text{-}commutative\text{-}ring}(r) \wedge \\ \quad n \in \mathbb{N} \wedge \\ \quad a \in \operatorname{mat}(n, n, \operatorname{carr}(r)) \wedge \\ \quad c \in \operatorname{carr}(r) \wedge \\ \quad k \in \operatorname{interval}(1, n) \wedge \\ \quad l \in \operatorname{interval}(1, n) \wedge \\ \quad \neg (k = l) \Rightarrow \\ \quad \quad \operatorname{det}(r, n, \operatorname{matmul}(r, \operatorname{elem\text{-}g}(r, n, c, k, l), a)) = \operatorname{det}(r, n, a) \end{array}}}
\vnbprop{det-elem-f-row}{\vnbstmt{\begin{array}{@{}l@{}} \forall r, n, a, k, l. \\ \quad \operatorname{is\text{-}commutative\text{-}ring}(r) \wedge \\ \quad n \in \mathbb{N} \wedge \\ \quad a \in \operatorname{mat}(n, n, \operatorname{carr}(r)) \wedge \\ \quad k \in \operatorname{interval}(1, n) \wedge \\ \quad l \in \operatorname{interval}(1, n) \wedge \\ \quad \neg (k = l) \Rightarrow \\ \quad \quad \operatorname{det}(r, n, \operatorname{matmul}(r, \operatorname{elem\text{-}f}(r, n, k, l), a)) = \operatorname{neg}(r)(\operatorname{det}(r, n, a)) \end{array}}}
\vnbprop{dtr-ring-rx}{\vnbstmt{\begin{array}{@{}l@{}} \forall r, d, e, t. \\ \quad \operatorname{is\text{-}commutative\text{-}ring}(r) \wedge \\ \quad d \in \operatorname{carr}(r) \wedge \\ \quad e \in \operatorname{carr}(r) \wedge \\ \quad t \in \operatorname{carr}(r) \wedge \\ \quad e = \operatorname{neg}(r)(d) \wedge \\ \quad e = t \Rightarrow \\ \quad \quad d = \operatorname{add}(r)(\operatorname{mul}(r)(\operatorname{neg}(r)(\operatorname{one}(r)), t), \operatorname{zero}(r)) \end{array}}}
\vnbprop{dtr-ring-rx2}{\vnbstmt{\begin{array}{@{}l@{}} \forall r, s, a, d. \\ \quad (\operatorname{is\text{-}commutative\text{-}ring}(r) \wedge ((s \in \operatorname{carr}(r)) \wedge ((a \in \operatorname{carr}(r)) \wedge (d \in \operatorname{carr}(r))))) \Rightarrow  \\ \quad \quad \operatorname{mul}(r)(\operatorname{mul}(r)(\operatorname{neg}(r)(\operatorname{one}(r)), s), \operatorname{mul}(r)(a, d)) \\ \quad \quad =\; \operatorname{add}(r)(\operatorname{mul}(r)(\operatorname{neg}(r)(\operatorname{one}(r)), \operatorname{mul}(r)(s, \operatorname{mul}(r)(a, d))), \operatorname{zero}(r)) \end{array}}}
\vnbprop{det-expand-row-ind}{\vnbstmt{\begin{array}{@{}l@{}} \forall i \in \mathbb{N}, n, r, a. \\ \quad n \in \mathbb{N} \wedge \\ \quad \operatorname{is\text{-}commutative\text{-}ring}(r) \wedge \\ \quad a \in \operatorname{mat}((n + 1), (n + 1), \operatorname{carr}(r)) \wedge \\ \quad 1 \leq i \wedge \\ \quad i \leq (n + 1) \Rightarrow \\ \quad \quad \operatorname{det}(r, (n + 1), a) \\ \quad \quad =\; \operatorname{finsum}(\operatorname{ring\text{-}additive\text{-}ag}(r), \\ \quad \quad \quad (\lambda j \in \operatorname{interval}(1, (n + 1)).\; \\ \quad \quad \quad \quad \operatorname{mul}(r)(\operatorname{mpow}(\operatorname{ring\text{-}multiplicative\text{-}monoid}(r), \operatorname{neg}(r)(\operatorname{one}(r)), (i + j)), \\ \quad \quad \quad \quad \quad \operatorname{mul}(r)(\operatorname{entry}(a, i, j), \operatorname{det}(r, n, \operatorname{minor}(a, i, j, n))))), \\ \quad \quad \quad \operatorname{interval}(1, (n + 1))) \end{array}}}
\vnbprop{det-expand-row}{\vnbstmt{\begin{array}{@{}l@{}} \forall r, n, a, i. \\ \quad \operatorname{is\text{-}commutative\text{-}ring}(r) \wedge \\ \quad n \in \mathbb{N} \wedge \\ \quad a \in \operatorname{mat}((n + 1), (n + 1), \operatorname{carr}(r)) \wedge \\ \quad i \in \operatorname{interval}(1, (n + 1)) \Rightarrow \\ \quad \quad \operatorname{det}(r, (n + 1), a) \\ \quad \quad =\; \operatorname{finsum}(\operatorname{ring\text{-}additive\text{-}ag}(r), \\ \quad \quad \quad (\lambda j \in \operatorname{interval}(1, (n + 1)).\; \\ \quad \quad \quad \quad \operatorname{mul}(r)(\operatorname{mpow}(\operatorname{ring\text{-}multiplicative\text{-}monoid}(r), \operatorname{neg}(r)(\operatorname{one}(r)), (i + j)), \\ \quad \quad \quad \quad \quad \operatorname{mul}(r)(\operatorname{entry}(a, i, j), \operatorname{det}(r, n, \operatorname{minor}(a, i, j, n))))), \\ \quad \quad \quad \operatorname{interval}(1, (n + 1))) \end{array}}}

\vnbbookfile{theorem-library/det-mul}{theorem-library/det-mul}
\noindent{\small det-mul.scm -- det(E A) = det(E) det(A) for an ELEMENTARY matrix E (batch 25-C, 2026-09-24).  The multiplicativity of the determinant as far as the tree carries it; see the report (scratchpad/triage/BATCH25-REPORTS.md, 25-C) for what stands between this and linear-algebra 3.20 / 3.27.}\par
\vnbprop{det-elem-h}{\vnbstmt{\begin{array}{@{}l@{}} \forall r, n, c, k. \\ \quad (\operatorname{is\text{-}commutative\text{-}ring}(r) \wedge ((n \in \mathbb{N}) \wedge ((c \in \operatorname{carr}(r)) \wedge (k \in \operatorname{interval}(1, n))))) \Rightarrow  \\ \quad \quad \operatorname{det}(r, n, \operatorname{elem\text{-}h}(r, n, c, k)) = c \end{array}}}
\vnbprop{det-elem-g}{\vnbstmt{\begin{array}{@{}l@{}} \forall r, n, c, k, l. \\ \quad \operatorname{is\text{-}commutative\text{-}ring}(r) \wedge \\ \quad n \in \mathbb{N} \wedge \\ \quad c \in \operatorname{carr}(r) \wedge \\ \quad k \in \operatorname{interval}(1, n) \wedge \\ \quad l \in \operatorname{interval}(1, n) \wedge \\ \quad \neg (k = l) \Rightarrow \\ \quad \quad \operatorname{det}(r, n, \operatorname{elem\text{-}g}(r, n, c, k, l)) = \operatorname{one}(r) \end{array}}}
\vnbprop{det-elem-f}{\vnbstmt{\begin{array}{@{}l@{}} \forall r, n, k, l. \\ \quad (\operatorname{is\text{-}commutative\text{-}ring}(r) \wedge ((n \in \mathbb{N}) \wedge ((k \in \operatorname{interval}(1, n)) \wedge ((l \in \operatorname{interval}(1, n)) \wedge \neg (k = l))))) \Rightarrow  \\ \quad \quad \operatorname{det}(r, n, \operatorname{elem\text{-}f}(r, n, k, l)) = \operatorname{neg}(r)(\operatorname{one}(r)) \end{array}}}
\vnbprop{det-mul-elem-h}{\vnbstmt{\begin{array}{@{}l@{}} \forall r, n, a, c, k. \\ \quad \operatorname{is\text{-}commutative\text{-}ring}(r) \wedge \\ \quad n \in \mathbb{N} \wedge \\ \quad a \in \operatorname{mat}(n, n, \operatorname{carr}(r)) \wedge \\ \quad c \in \operatorname{carr}(r) \wedge \\ \quad k \in \operatorname{interval}(1, n) \Rightarrow \\ \quad \quad \operatorname{det}(r, n, \operatorname{matmul}(r, \operatorname{elem\text{-}h}(r, n, c, k), a)) = \operatorname{mul}(r)(\operatorname{det}(r, n, \operatorname{elem\text{-}h}(r, n, c, k)), \operatorname{det}(r, n, a)) \end{array}}}
\vnbprop{det-mul-elem-g}{\vnbstmt{\begin{array}{@{}l@{}} \forall r, n, a, c, k, l. \\ \quad \operatorname{is\text{-}commutative\text{-}ring}(r) \wedge \\ \quad n \in \mathbb{N} \wedge \\ \quad a \in \operatorname{mat}(n, n, \operatorname{carr}(r)) \wedge \\ \quad c \in \operatorname{carr}(r) \wedge \\ \quad k \in \operatorname{interval}(1, n) \wedge \\ \quad l \in \operatorname{interval}(1, n) \wedge \\ \quad \neg (k = l) \Rightarrow \\ \quad \quad \operatorname{det}(r, n, \operatorname{matmul}(r, \operatorname{elem\text{-}g}(r, n, c, k, l), a)) \\ \quad \quad =\; \operatorname{mul}(r)(\operatorname{det}(r, n, \operatorname{elem\text{-}g}(r, n, c, k, l)), \operatorname{det}(r, n, a)) \end{array}}}
\vnbprop{det-mul-elem-f}{\vnbstmt{\begin{array}{@{}l@{}} \forall r, n, a, k, l. \\ \quad \operatorname{is\text{-}commutative\text{-}ring}(r) \wedge \\ \quad n \in \mathbb{N} \wedge \\ \quad a \in \operatorname{mat}(n, n, \operatorname{carr}(r)) \wedge \\ \quad k \in \operatorname{interval}(1, n) \wedge \\ \quad l \in \operatorname{interval}(1, n) \wedge \\ \quad \neg (k = l) \Rightarrow \\ \quad \quad \operatorname{det}(r, n, \operatorname{matmul}(r, \operatorname{elem\text{-}f}(r, n, k, l), a)) = \operatorname{mul}(r)(\operatorname{det}(r, n, \operatorname{elem\text{-}f}(r, n, k, l)), \operatorname{det}(r, n, a)) \end{array}}}

\vnbbookfile{theorem-library/pivot-col-reduce-proof}{theorem-library/pivot-col-reduce-proof}
\noindent{\small pivot-col-reduce-proof.scm -- Brick B2b (LA Phase B, algebraic-numbers.pdf ch.3): ONE elementary column operation reduces a first-row entry against the pivot.  Over a euclidean ring A, for P with nonzero pivot P\_\{1,1\} and target column j/=1, there are q,r with (P . G[-q,1,j])\_\{1,j\} = r and (r=0 or}\par
\vnbprop{pivot-col-reduce}{\vnbstmt{\begin{array}{@{}l@{}} \forall a, m, n, p, j. \\ \quad \operatorname{is\text{-}euclidean\text{-}ring}(a) \wedge \\ \quad p \in \operatorname{mat}(m, n, \operatorname{carr}(a)) \wedge \\ \quad 1 \in \operatorname{interval}(1, m) \wedge \\ \quad 1 \in \operatorname{interval}(1, n) \wedge \\ \quad j \in \operatorname{interval}(1, n) \wedge \\ \quad \neg (1 = j) \wedge \\ \quad \neg (\operatorname{entry}(p, 1, 1) = \operatorname{zero}(a)) \Rightarrow \\ \quad \quad \exists q, r \in \operatorname{carr}(a). \\ \quad \quad \quad \operatorname{entry}(\operatorname{matmul}(a, p, \operatorname{elem\text{-}g}(a, n, \operatorname{neg}(a)(q), 1, j)), 1, j) = r \wedge \\ \quad \quad \quad (r = \operatorname{zero}(a)) \vee ((\operatorname{gauge}(a)(r) + 1) \leq \operatorname{gauge}(a)(\operatorname{entry}(p, 1, 1))) \end{array}}}

\vnbbookfile{theorem-library/pivot-row-reduce-proof}{theorem-library/pivot-row-reduce-proof}
\noindent{\small pivot-row-reduce-proof.scm -- Brick B3a (LA Phase B, algebraic-numbers.pdf ch.3): the ROW-reduction step, mirror of pivot-col-reduce.  Over a euclidean ring A, for P with nonzero pivot P\_\{1,1\} and target row i/=1, there are q,r with (ELEM-G[-q,i,1] . P)\_\{i,1\} = r and (r=0 or deg(r) $<$ deg(P\_\{1,1\})) -- one}\par
\vnbprop{pivot-row-reduce}{\vnbstmt{\begin{array}{@{}l@{}} \forall a, m, n, p, i. \\ \quad \operatorname{is\text{-}euclidean\text{-}ring}(a) \wedge \\ \quad p \in \operatorname{mat}(m, n, \operatorname{carr}(a)) \wedge \\ \quad 1 \in \operatorname{interval}(1, m) \wedge \\ \quad 1 \in \operatorname{interval}(1, n) \wedge \\ \quad i \in \operatorname{interval}(1, m) \wedge \\ \quad \neg (i = 1) \wedge \\ \quad \neg (\operatorname{entry}(p, 1, 1) = \operatorname{zero}(a)) \Rightarrow \\ \quad \quad \exists q, r \in \operatorname{carr}(a). \\ \quad \quad \quad \operatorname{entry}(\operatorname{matmul}(a, \operatorname{elem\text{-}g}(a, m, \operatorname{neg}(a)(q), i, 1), p), i, 1) = r \wedge \\ \quad \quad \quad (r = \operatorname{zero}(a)) \vee ((\operatorname{gauge}(a)(r) + 1) \leq \operatorname{gauge}(a)(\operatorname{entry}(p, 1, 1))) \end{array}}}

\vnbbookfile{theorem-library/mat-equiv-proof}{theorem-library/mat-equiv-proof}
\noindent{\small mat-equiv-proof.scm -- Brick B3 infrastructure: basic laws of \textasciitilde{} (mat-equiv.scm). invertible-mat-is-mat, identmat-invertible, mat-equiv-refl, mat-equiv-right-mult, mat-equiv-left-mult, elem-f-invertible.  \textasciitilde{} is reflexive and each elementary swap (row/col) yields an equivalent matrix.  elem-g-invertible + transitivity + ...}\par
\vnbprop{invertible-mat-is-mat}{\vnbstmt{\forall a, n, u.\; (\operatorname{is\text{-}invertible\text{-}mat}(a, n, u) \Rightarrow (u \in \operatorname{mat}(n, n, \operatorname{carr}(a))))}}
\vnbprop{identmat-invertible}{\vnbstmt{\forall a.; \forall n \in \mathbb{N}.\; (\operatorname{is\text{-}ring}(a) \Rightarrow \operatorname{is\text{-}invertible\text{-}mat}(a, n, \operatorname{identmat}(a, n)))}}
\vnbprop{mat-equiv-refl}{\vnbstmt{\begin{array}{@{}l@{}} \forall a, m, n.; \forall c \in \operatorname{mat}(m, n, \operatorname{carr}(a)). \\ \quad \operatorname{is\text{-}ring}(a) \Rightarrow \operatorname{mat\text{-}equiv}(a, m, n, c, c) \end{array}}}
\vnbprop{mat-equiv-right-mult}{\vnbstmt{\begin{array}{@{}l@{}} \forall a, m, n, c, v. \\ \quad (\operatorname{is\text{-}ring}(a) \wedge ((c \in \operatorname{mat}(m, n, \operatorname{carr}(a))) \wedge \operatorname{is\text{-}invertible\text{-}mat}(a, n, v))) \Rightarrow  \\ \quad \quad \operatorname{mat\text{-}equiv}(a, m, n, c, \operatorname{matmul}(a, c, v)) \end{array}}}
\vnbprop{mat-equiv-left-mult}{\vnbstmt{\begin{array}{@{}l@{}} \forall a, m, n, c, u. \\ \quad (\operatorname{is\text{-}ring}(a) \wedge ((c \in \operatorname{mat}(m, n, \operatorname{carr}(a))) \wedge \operatorname{is\text{-}invertible\text{-}mat}(a, m, u))) \Rightarrow  \\ \quad \quad \operatorname{mat\text{-}equiv}(a, m, n, c, \operatorname{matmul}(a, u, c)) \end{array}}}
\vnbprop{elem-f-invertible}{\vnbstmt{\begin{array}{@{}l@{}} \forall a, n, k, l. \\ \quad (\operatorname{is\text{-}ring}(a) \wedge ((n \in \mathbb{N}) \wedge ((k \in \operatorname{interval}(1, n)) \wedge ((l \in \operatorname{interval}(1, n)) \wedge \neg (k = l))))) \Rightarrow  \\ \quad \quad \operatorname{is\text{-}invertible\text{-}mat}(a, n, \operatorname{elem\text{-}f}(a, n, k, l)) \end{array}}}
\vnbprop{elem-g-param-cong}{\vnbstmt{\begin{array}{@{}l@{}} \forall a, n, r, s, k, l. \\ \quad (\operatorname{is\text{-}ring}(a) \wedge ((n \in \mathbb{N}) \wedge ((s \in \operatorname{carr}(a)) \wedge (r = s)))) \Rightarrow  \\ \quad \quad \operatorname{elem\text{-}g}(a, n, r, k, l) = \operatorname{elem\text{-}g}(a, n, s, k, l) \end{array}}}
\vnbprop{elem-g-invertible}{\vnbstmt{\begin{array}{@{}l@{}} \forall a, n, r, k, l. \\ \quad \operatorname{is\text{-}ring}(a) \wedge \\ \quad n \in \mathbb{N} \wedge \\ \quad r \in \operatorname{carr}(a) \wedge \\ \quad k \in \operatorname{interval}(1, n) \wedge \\ \quad l \in \operatorname{interval}(1, n) \wedge \\ \quad \neg (k = l) \Rightarrow \\ \quad \quad \operatorname{is\text{-}invertible\text{-}mat}(a, n, \operatorname{elem\text{-}g}(a, n, r, k, l)) \end{array}}}
\vnbprop{product-of-invertibles-is-invertible}{\vnbstmt{\begin{array}{@{}l@{}} \forall a, n, b, c. \\ \quad (\operatorname{is\text{-}ring}(a) \wedge (\operatorname{is\text{-}invertible\text{-}mat}(a, n, b) \wedge \operatorname{is\text{-}invertible\text{-}mat}(a, n, c))) \Rightarrow  \\ \quad \quad \operatorname{is\text{-}invertible\text{-}mat}(a, n, \operatorname{matmul}(a, b, c)) \end{array}}}
\vnbprop{mat-equiv-trans}{\vnbstmt{\begin{array}{@{}l@{}} \forall a, m, n, c, d, e. \\ \quad (\operatorname{is\text{-}ring}(a) \wedge ((c \in \operatorname{mat}(m, n, \operatorname{carr}(a))) \wedge (\operatorname{mat\text{-}equiv}(a, m, n, c, d) \wedge \operatorname{mat\text{-}equiv}(a, m, n, d, e)))) \Rightarrow  \\ \quad \quad \operatorname{mat\text{-}equiv}(a, m, n, c, e) \end{array}}}

\vnbbookfile{theorem-library/min-degree-entry-proof}{theorem-library/min-degree-entry-proof}
\noindent{\small min-degree-entry-proof.scm -- a matrix with a nonzero entry has a nonzero entry of MINIMAL Euclidean degree.  The "mu is achieved" fact that starts the Smith reduction (Prop 3.36).}\par
\vnbprop{min-degree-entry}{\vnbstmt{\begin{array}{@{}l@{}} \forall a, m, n.; \forall p \in \operatorname{mat}(m, n, \operatorname{carr}(a)). \\ \quad \operatorname{is\text{-}euclidean\text{-}ring}(a) \wedge \\ \quad \exists b \in \operatorname{interval}(1, m), c \in \operatorname{interval}(1, n). \\ \quad \quad \neg (\operatorname{entry}(p, b, c) = \operatorname{zero}(a)) \Rightarrow \\ \quad \quad \exists s \in \operatorname{interval}(1, m), d \in \operatorname{interval}(1, n). \\ \quad \quad \quad \neg (\operatorname{entry}(p, s, d) = \operatorname{zero}(a)) \wedge \\ \quad \quad \quad \forall i, j. \\ \quad \quad \quad \quad ((i \in \operatorname{interval}(1, m)) \wedge ((j \in \operatorname{interval}(1, n)) \wedge \neg (\operatorname{entry}(p, i, j) = \operatorname{zero}(a)))) \Rightarrow  \\ \quad \quad \quad \quad \quad \operatorname{gauge}(a)(\operatorname{entry}(p, s, d)) \leq \operatorname{gauge}(a)(\operatorname{entry}(p, i, j)) \end{array}}}

\vnbbookfile{theorem-library/smith-proof}{theorem-library/smith-proof}
\noindent{\small smith-proof.scm -- Brick B3 (LA Phase B): the Smith normal-form reduction, built on the \textasciitilde{} equivalence infrastructure (mat-equiv-proof.scm). equiv-mul-both  -- A \textasciitilde{} U.A.V for U,V invertible (sandwiching preserves \textasciitilde{}); swap-to-corner  -- bring any entry (i0,j0) to position (1,1) preserving \textasciitilde{}}\par
\vnbprop{equiv-mul-both}{\vnbstmt{\begin{array}{@{}l@{}} \forall r, m, n, a, u, v. \\ \quad (\operatorname{is\text{-}ring}(r) \wedge ((a \in \operatorname{mat}(m, n, \operatorname{carr}(r))) \wedge (\operatorname{is\text{-}invertible\text{-}mat}(r, m, u) \wedge \operatorname{is\text{-}invertible\text{-}mat}(r, n, v)))) \Rightarrow  \\ \quad \quad \operatorname{mat\text{-}equiv}(r, m, n, a, \operatorname{matmul}(r, \operatorname{matmul}(r, u, a), v)) \end{array}}}
\vnbprop{swap-to-corner}{\vnbstmt{\begin{array}{@{}l@{}} \forall r, m, n, a, c, d. \\ \quad \operatorname{is\text{-}ring}(r) \wedge \\ \quad a \in \operatorname{mat}(m, n, \operatorname{carr}(r)) \wedge \\ \quad 1 \in \operatorname{interval}(1, m) \wedge \\ \quad 1 \in \operatorname{interval}(1, n) \wedge \\ \quad c \in \operatorname{interval}(1, m) \wedge \\ \quad d \in \operatorname{interval}(1, n) \wedge \\ \quad \neg (1 = c) \wedge \\ \quad \neg (1 = d) \Rightarrow \\ \quad \quad \exists b.\; (\operatorname{mat\text{-}equiv}(r, m, n, a, b) \wedge (\operatorname{entry}(b, 1, 1) = \operatorname{entry}(a, c, d))) \end{array}}}
\vnbprop{swap-to-corner-gen}{\vnbstmt{\begin{array}{@{}l@{}} \forall r, m, n, a, c, d. \\ \quad \operatorname{is\text{-}ring}(r) \wedge \\ \quad a \in \operatorname{mat}(m, n, \operatorname{carr}(r)) \wedge \\ \quad 1 \in \operatorname{interval}(1, m) \wedge \\ \quad 1 \in \operatorname{interval}(1, n) \wedge \\ \quad c \in \operatorname{interval}(1, m) \wedge \\ \quad d \in \operatorname{interval}(1, n) \Rightarrow \\ \quad \quad \exists b.\; (\operatorname{mat\text{-}equiv}(r, m, n, a, b) \wedge (\operatorname{entry}(b, 1, 1) = \operatorname{entry}(a, c, d))) \end{array}}}
\vnbprop{place-min-pivot}{\vnbstmt{\begin{array}{@{}l@{}} \forall a, m, n.; \forall p \in \operatorname{mat}(m, n, \operatorname{carr}(a)). \\ \quad \operatorname{is\text{-}euclidean\text{-}ring}(a) \wedge \\ \quad 1 \in \operatorname{interval}(1, m) \wedge \\ \quad 1 \in \operatorname{interval}(1, n) \wedge \\ \quad \exists c \in \operatorname{interval}(1, m), d \in \operatorname{interval}(1, n). \\ \quad \quad \neg (\operatorname{entry}(p, c, d) = \operatorname{zero}(a)) \Rightarrow \\ \quad \quad \exists b. \\ \quad \quad \quad \operatorname{mat\text{-}equiv}(a, m, n, p, b) \wedge \\ \quad \quad \quad \neg (\operatorname{entry}(b, 1, 1) = \operatorname{zero}(a)) \wedge \\ \quad \quad \quad \forall i, j. \\ \quad \quad \quad \quad ((i \in \operatorname{interval}(1, m)) \wedge ((j \in \operatorname{interval}(1, n)) \wedge \neg (\operatorname{entry}(p, i, j) = \operatorname{zero}(a)))) \Rightarrow  \\ \quad \quad \quad \quad \quad \operatorname{gauge}(a)(\operatorname{entry}(b, 1, 1)) \leq \operatorname{gauge}(a)(\operatorname{entry}(p, i, j)) \end{array}}}

\vnbbookfile{theorem-library/class-min-pivot-proof}{theorem-library/class-min-pivot-proof}
\noindent{\small class-min-pivot-proof.scm -- the minimal-degree pivot over the whole \textasciitilde{}-equivalence class, and the typing lemma it needs.}\par
\vnbprop{mat-equiv-target-is-mat}{\vnbstmt{\begin{array}{@{}l@{}} \forall a, m, n, c, d. \\ \quad (\operatorname{is\text{-}ring}(a) \wedge ((c \in \operatorname{mat}(m, n, \operatorname{carr}(a))) \wedge \operatorname{mat\text{-}equiv}(a, m, n, c, d))) \Rightarrow  \\ \quad \quad d \in \operatorname{mat}(m, n, \operatorname{carr}(a)) \end{array}}}
\vnbprop{class-min-pivot}{\vnbstmt{\begin{array}{@{}l@{}} \forall a, m, n.; \forall p \in \operatorname{mat}(m, n, \operatorname{carr}(a)). \\ \quad \operatorname{is\text{-}euclidean\text{-}ring}(a) \wedge \\ \quad 1 \in \operatorname{interval}(1, m) \wedge \\ \quad 1 \in \operatorname{interval}(1, n) \wedge \\ \quad \exists d \in \operatorname{interval}(1, m), f \in \operatorname{interval}(1, n). \\ \quad \quad \neg (\operatorname{entry}(p, d, f) = \operatorname{zero}(a)) \Rightarrow \\ \quad \quad \exists b. \\ \quad \quad \quad \operatorname{mat\text{-}equiv}(a, m, n, p, b) \wedge \\ \quad \quad \quad \neg (\operatorname{entry}(b, 1, 1) = \operatorname{zero}(a)) \wedge \\ \quad \quad \quad \forall c, i, j. \\ \quad \quad \quad \quad \operatorname{mat\text{-}equiv}(a, m, n, p, c) \wedge \\ \quad \quad \quad \quad i \in \operatorname{interval}(1, m) \wedge \\ \quad \quad \quad \quad j \in \operatorname{interval}(1, n) \wedge \\ \quad \quad \quad \quad \neg (\operatorname{entry}(c, i, j) = \operatorname{zero}(a)) \Rightarrow \\ \quad \quad \quad \quad \quad \operatorname{gauge}(a)(\operatorname{entry}(b, 1, 1)) \leq \operatorname{gauge}(a)(\operatorname{entry}(c, i, j)) \end{array}}}

\vnbbookfile{theorem-library/smith-clear-proof}{theorem-library/smith-clear-proof}
\noindent{\small smith-clear-proof.scm -- Brick B3 (LA Phase B): the DESCENT step of the Smith reduction.  pivot-clears-col: over a euclidean ring, if P's pivot P\_\{1,1\} is nonzero and of MINIMAL degree over the whole equivalence class of P (class-min-pivot supplies such a P), then a single elementary COLUMN op zeroes}\par
\vnbprop{pivot-clears-col}{\vnbstmt{\begin{array}{@{}l@{}} \forall a, m, n, p, j. \\ \quad \operatorname{is\text{-}euclidean\text{-}ring}(a) \wedge \\ \quad p \in \operatorname{mat}(m, n, \operatorname{carr}(a)) \wedge \\ \quad 1 \in \operatorname{interval}(1, m) \wedge \\ \quad 1 \in \operatorname{interval}(1, n) \wedge \\ \quad j \in \operatorname{interval}(1, n) \wedge \\ \quad \neg (1 = j) \wedge \\ \quad \neg (\operatorname{entry}(p, 1, 1) = \operatorname{zero}(a)) \wedge \\ \quad \forall c, i, b. \\ \quad \quad \operatorname{mat\text{-}equiv}(a, m, n, p, c) \wedge \\ \quad \quad i \in \operatorname{interval}(1, m) \wedge \\ \quad \quad b \in \operatorname{interval}(1, n) \wedge \\ \quad \quad \neg (\operatorname{entry}(c, i, b) = \operatorname{zero}(a)) \Rightarrow \\ \quad \quad \quad \operatorname{gauge}(a)(\operatorname{entry}(p, 1, 1)) \leq \operatorname{gauge}(a)(\operatorname{entry}(c, i, b)) \Rightarrow \\ \quad \quad \exists q. \\ \quad \quad \quad \operatorname{mat\text{-}equiv}(a, m, n, p, q) \wedge \\ \quad \quad \quad \operatorname{entry}(q, 1, 1) = \operatorname{entry}(p, 1, 1) \wedge \\ \quad \quad \quad \operatorname{entry}(q, 1, j) = \operatorname{zero}(a) \wedge \\ \quad \quad \quad \forall i, c. \\ \quad \quad \quad \quad ((i \in \operatorname{interval}(1, m)) \wedge ((c \in \operatorname{interval}(1, n)) \wedge \neg (c = j))) \Rightarrow  \\ \quad \quad \quad \quad \quad \operatorname{entry}(q, i, c) = \operatorname{entry}(p, i, c) \end{array}}}

\vnbbookfile{theorem-library/clear-first-row-proof}{theorem-library/clear-first-row-proof}
\noindent{\small clear-first-row-proof.scm -- Brick B3 (LA Phase B): clear the whole first row. Iterates pivot-clears-col across columns 2..n by NN-induction on the number of cleared columns, then instantiates at n.  With a class-minimal pivot at (1,1), the result is a matrix \textasciitilde{} P, same pivot, whose first row is zero off the pivot.}\par
\vnbprop{classmin-transport}{\vnbstmt{\begin{array}{@{}l@{}} \forall a, m, n, p, q, c, i, j. \\ \quad \operatorname{is\text{-}ring}(a) \wedge \\ \quad p \in \operatorname{mat}(m, n, \operatorname{carr}(a)) \wedge \\ \quad \operatorname{mat\text{-}equiv}(a, m, n, p, q) \wedge \\ \quad \operatorname{entry}(q, 1, 1) = \operatorname{entry}(p, 1, 1) \wedge \\ \quad \forall c, i, j. \\ \quad \quad \operatorname{mat\text{-}equiv}(a, m, n, p, c) \wedge \\ \quad \quad i \in \operatorname{interval}(1, m) \wedge \\ \quad \quad j \in \operatorname{interval}(1, n) \wedge \\ \quad \quad \neg (\operatorname{entry}(c, i, j) = \operatorname{zero}(a)) \Rightarrow \\ \quad \quad \quad \operatorname{gauge}(a)(\operatorname{entry}(p, 1, 1)) \leq \operatorname{gauge}(a)(\operatorname{entry}(c, i, j)) \wedge \\ \quad \operatorname{mat\text{-}equiv}(a, m, n, q, c) \wedge \\ \quad i \in \operatorname{interval}(1, m) \wedge \\ \quad j \in \operatorname{interval}(1, n) \wedge \\ \quad \neg (\operatorname{entry}(c, i, j) = \operatorname{zero}(a)) \Rightarrow \\ \quad \quad \operatorname{gauge}(a)(\operatorname{entry}(q, 1, 1)) \leq \operatorname{gauge}(a)(\operatorname{entry}(c, i, j)) \end{array}}}
\vnbprop{clear-row-upto}{\vnbstmt{\begin{array}{@{}l@{}} \forall a, m, n.; \forall p \in \operatorname{mat}(m, n, \operatorname{carr}(a)), k \in \mathbb{N}. \\ \quad \operatorname{is\text{-}euclidean\text{-}ring}(a) \wedge \\ \quad 1 \in \operatorname{interval}(1, m) \wedge \\ \quad 1 \in \operatorname{interval}(1, n) \wedge \\ \quad \neg (\operatorname{entry}(p, 1, 1) = \operatorname{zero}(a)) \wedge \\ \quad \forall c, i, b. \\ \quad \quad \operatorname{mat\text{-}equiv}(a, m, n, p, c) \wedge \\ \quad \quad i \in \operatorname{interval}(1, m) \wedge \\ \quad \quad b \in \operatorname{interval}(1, n) \wedge \\ \quad \quad \neg (\operatorname{entry}(c, i, b) = \operatorname{zero}(a)) \Rightarrow \\ \quad \quad \quad \operatorname{gauge}(a)(\operatorname{entry}(p, 1, 1)) \leq \operatorname{gauge}(a)(\operatorname{entry}(c, i, b)) \Rightarrow \\ \quad \quad \exists q. \\ \quad \quad \quad \operatorname{mat\text{-}equiv}(a, m, n, p, q) \wedge \\ \quad \quad \quad \operatorname{entry}(q, 1, 1) = \operatorname{entry}(p, 1, 1) \wedge \\ \quad \quad \quad \forall j \in \operatorname{interval}(1, n). \\ \quad \quad \quad \quad (\neg (j = 1) \wedge (j \leq k)) \Rightarrow (\operatorname{entry}(q, 1, j) = \operatorname{zero}(a)) \end{array}}}
\vnbprop{clear-first-row}{\vnbstmt{\begin{array}{@{}l@{}} \forall a, m, n.; \forall p \in \operatorname{mat}(m, n, \operatorname{carr}(a)). \\ \quad \operatorname{is\text{-}euclidean\text{-}ring}(a) \wedge \\ \quad m \in \mathbb{N} \wedge \\ \quad n \in \mathbb{N} \wedge \\ \quad 1 \in \operatorname{interval}(1, m) \wedge \\ \quad 1 \in \operatorname{interval}(1, n) \wedge \\ \quad \neg (\operatorname{entry}(p, 1, 1) = \operatorname{zero}(a)) \wedge \\ \quad \forall c, i, b. \\ \quad \quad \operatorname{mat\text{-}equiv}(a, m, n, p, c) \wedge \\ \quad \quad i \in \operatorname{interval}(1, m) \wedge \\ \quad \quad b \in \operatorname{interval}(1, n) \wedge \\ \quad \quad \neg (\operatorname{entry}(c, i, b) = \operatorname{zero}(a)) \Rightarrow \\ \quad \quad \quad \operatorname{gauge}(a)(\operatorname{entry}(p, 1, 1)) \leq \operatorname{gauge}(a)(\operatorname{entry}(c, i, b)) \Rightarrow \\ \quad \quad \exists q. \\ \quad \quad \quad \operatorname{mat\text{-}equiv}(a, m, n, p, q) \wedge \\ \quad \quad \quad \operatorname{entry}(q, 1, 1) = \operatorname{entry}(p, 1, 1) \wedge \\ \quad \quad \quad \forall j \in \operatorname{interval}(1, n).\; (\neg (j = 1) \Rightarrow (\operatorname{entry}(q, 1, j) = \operatorname{zero}(a))) \end{array}}}

\vnbbookfile{theorem-library/clear-first-col-proof}{theorem-library/clear-first-col-proof}
\noindent{\small clear-first-col-proof.scm -- Brick B3 (LA Phase B): clear the whole first column. The ROW-op mirror of clear-first-row: left-multiply by elementary row matrices. pivot-clears-row -- one row op zeroes an off-pivot column-1 entry (i,1); the euclidean remainder is forced to vanish by class-minimality}\par
\vnbprop{pivot-clears-row}{\vnbstmt{\begin{array}{@{}l@{}} \forall a, m, n, p, i. \\ \quad \operatorname{is\text{-}euclidean\text{-}ring}(a) \wedge \\ \quad p \in \operatorname{mat}(m, n, \operatorname{carr}(a)) \wedge \\ \quad 1 \in \operatorname{interval}(1, m) \wedge \\ \quad 1 \in \operatorname{interval}(1, n) \wedge \\ \quad i \in \operatorname{interval}(1, m) \wedge \\ \quad \neg (i = 1) \wedge \\ \quad \neg (\operatorname{entry}(p, 1, 1) = \operatorname{zero}(a)) \wedge \\ \quad \forall c, b, j. \\ \quad \quad \operatorname{mat\text{-}equiv}(a, m, n, p, c) \wedge \\ \quad \quad b \in \operatorname{interval}(1, m) \wedge \\ \quad \quad j \in \operatorname{interval}(1, n) \wedge \\ \quad \quad \neg (\operatorname{entry}(c, b, j) = \operatorname{zero}(a)) \Rightarrow \\ \quad \quad \quad \operatorname{gauge}(a)(\operatorname{entry}(p, 1, 1)) \leq \operatorname{gauge}(a)(\operatorname{entry}(c, b, j)) \Rightarrow \\ \quad \quad \exists q. \\ \quad \quad \quad \operatorname{mat\text{-}equiv}(a, m, n, p, q) \wedge \\ \quad \quad \quad \operatorname{entry}(q, 1, 1) = \operatorname{entry}(p, 1, 1) \wedge \\ \quad \quad \quad \operatorname{entry}(q, i, 1) = \operatorname{zero}(a) \wedge \\ \quad \quad \quad \forall w, c. \\ \quad \quad \quad \quad ((w \in \operatorname{interval}(1, m)) \wedge ((c \in \operatorname{interval}(1, n)) \wedge \neg (w = i))) \Rightarrow  \\ \quad \quad \quad \quad \quad \operatorname{entry}(q, w, c) = \operatorname{entry}(p, w, c) \end{array}}}
\vnbprop{clear-col-upto}{\vnbstmt{\begin{array}{@{}l@{}} \forall a, m, n.; \forall p \in \operatorname{mat}(m, n, \operatorname{carr}(a)), k \in \mathbb{N}. \\ \quad \operatorname{is\text{-}euclidean\text{-}ring}(a) \wedge \\ \quad 1 \in \operatorname{interval}(1, m) \wedge \\ \quad 1 \in \operatorname{interval}(1, n) \wedge \\ \quad \neg (\operatorname{entry}(p, 1, 1) = \operatorname{zero}(a)) \wedge \\ \quad \forall c, b, j. \\ \quad \quad \operatorname{mat\text{-}equiv}(a, m, n, p, c) \wedge \\ \quad \quad b \in \operatorname{interval}(1, m) \wedge \\ \quad \quad j \in \operatorname{interval}(1, n) \wedge \\ \quad \quad \neg (\operatorname{entry}(c, b, j) = \operatorname{zero}(a)) \Rightarrow \\ \quad \quad \quad \operatorname{gauge}(a)(\operatorname{entry}(p, 1, 1)) \leq \operatorname{gauge}(a)(\operatorname{entry}(c, b, j)) \Rightarrow \\ \quad \quad \exists q. \\ \quad \quad \quad \operatorname{mat\text{-}equiv}(a, m, n, p, q) \wedge \\ \quad \quad \quad \operatorname{entry}(q, 1, 1) = \operatorname{entry}(p, 1, 1) \wedge \\ \quad \quad \quad \forall i \in \operatorname{interval}(1, m). \\ \quad \quad \quad \quad (\neg (i = 1) \wedge (i \leq k)) \Rightarrow (\operatorname{entry}(q, i, 1) = \operatorname{zero}(a)) \wedge \\ \quad \quad \quad \forall c \in \operatorname{interval}(1, n).\; (\operatorname{entry}(q, 1, c) = \operatorname{entry}(p, 1, c)) \end{array}}}
\vnbprop{clear-first-col}{\vnbstmt{\begin{array}{@{}l@{}} \forall a, m, n.; \forall p \in \operatorname{mat}(m, n, \operatorname{carr}(a)). \\ \quad \operatorname{is\text{-}euclidean\text{-}ring}(a) \wedge \\ \quad m \in \mathbb{N} \wedge \\ \quad n \in \mathbb{N} \wedge \\ \quad 1 \in \operatorname{interval}(1, m) \wedge \\ \quad 1 \in \operatorname{interval}(1, n) \wedge \\ \quad \neg (\operatorname{entry}(p, 1, 1) = \operatorname{zero}(a)) \wedge \\ \quad \forall c, b, j. \\ \quad \quad \operatorname{mat\text{-}equiv}(a, m, n, p, c) \wedge \\ \quad \quad b \in \operatorname{interval}(1, m) \wedge \\ \quad \quad j \in \operatorname{interval}(1, n) \wedge \\ \quad \quad \neg (\operatorname{entry}(c, b, j) = \operatorname{zero}(a)) \Rightarrow \\ \quad \quad \quad \operatorname{gauge}(a)(\operatorname{entry}(p, 1, 1)) \leq \operatorname{gauge}(a)(\operatorname{entry}(c, b, j)) \Rightarrow \\ \quad \quad \exists q. \\ \quad \quad \quad \operatorname{mat\text{-}equiv}(a, m, n, p, q) \wedge \\ \quad \quad \quad \operatorname{entry}(q, 1, 1) = \operatorname{entry}(p, 1, 1) \wedge \\ \quad \quad \quad \forall i \in \operatorname{interval}(1, m).\; (\neg (i = 1) \Rightarrow (\operatorname{entry}(q, i, 1) = \operatorname{zero}(a))) \wedge \\ \quad \quad \quad \forall c \in \operatorname{interval}(1, n).\; (\operatorname{entry}(q, 1, c) = \operatorname{entry}(p, 1, c)) \end{array}}}

\vnbbookfile{theorem-library/matact-summand-type-le-proof}{theorem-library/matact-summand-type-le-proof}
\noindent{\small matact-summand-type-le-proof.scm -- matact-summand-type-le, GUARDED and PROVEN.}\par
\vnbprop{matact-summand-type-le-guarded}{\vnbstmt{\begin{array}{@{}l@{}} \forall d, m, n, q, p, u, i, c, k. \\ \quad \operatorname{is\text{-}module}(d) \wedge \\ \quad p \in \operatorname{mat}(m, n, \operatorname{carr}(\operatorname{scal}(d))) \wedge \\ \quad u \in \operatorname{mat}(n, q, \operatorname{vec}(d)) \wedge \\ \quad i \in \operatorname{interval}(1, m) \wedge \\ \quad c \in \operatorname{interval}(1, q) \wedge \\ \quad k \in \mathbb{N} \wedge \\ \quad k \leq n \Rightarrow \\ \quad \quad (\lambda j \in \operatorname{interval}(1, k).\; \operatorname{act}(d)(\operatorname{entry}(p, i, j), \operatorname{entry}(u, j, c))) \\ \quad \quad \in\; (\operatorname{interval}(1, k) \to \operatorname{carr}(\operatorname{module\text{-}vector\text{-}ag}(d))) \end{array}}}

\vnbbookfile{theorem-library/border-mult-proof}{theorem-library/border-mult-proof}
\noindent{\small border-mult-proof.scm -- Brick B3 (LA Phase B): block-matrix multiplication. BORDER(A,b,X,p,q) . BORDER(A,d,Y,q,r) = BORDER(A, b*d, X.Y, p, r)   [trust:none] The [[b,0],[0,X]] . [[d,0],[0,Y]] = [[bd,0],[0,XY]] direct-sum law -- the crux that lets the Smith recursion border a block equivalence back to the full}\par
\vnbprop{border-mult}{\vnbstmt{\begin{array}{@{}l@{}} \forall a, p, q, r, x, y, b, d. \\ \quad \operatorname{is\text{-}ring}(a) \wedge \\ \quad p \in \mathbb{N} \wedge \\ \quad q \in \mathbb{N} \wedge \\ \quad r \in \mathbb{N} \wedge \\ \quad x \in \operatorname{mat}(p, q, \operatorname{carr}(a)) \wedge \\ \quad y \in \operatorname{mat}(q, r, \operatorname{carr}(a)) \wedge \\ \quad b \in \operatorname{carr}(a) \wedge \\ \quad d \in \operatorname{carr}(a) \wedge \\ \quad (q = 0) \Rightarrow ((p = 0) \vee (r = 0)) \Rightarrow \\ \quad \quad \operatorname{matmul}(a, \operatorname{border}(a, b, x, p, q), \operatorname{border}(a, d, y, q, r)) \\ \quad \quad =\; \operatorname{border}(a, \operatorname{mul}(a)(b, d), \operatorname{matmul}(a, x, y), p, r) \end{array}}}

\vnbbookfile{theorem-library/border-assembly-proof}{theorem-library/border-assembly-proof}
\vnbprop{border-identity}{\vnbstmt{\begin{array}{@{}l@{}} \forall a.; \forall p \in \mathbb{N}. \\ \quad \operatorname{is\text{-}ring}(a) \Rightarrow  \\ \quad \quad \operatorname{border}(a, \operatorname{one}(a), \operatorname{identmat}(a, p), p, p) = \operatorname{identmat}(a, (p + 1)) \end{array}}}
\vnbprop{border-is-diagonal}{\vnbstmt{\begin{array}{@{}l@{}} \forall a, b, e, p, q. \\ \quad \operatorname{is\text{-}ring}(a) \wedge \\ \quad p \in \mathbb{N} \wedge \\ \quad q \in \mathbb{N} \wedge \\ \quad \operatorname{is\text{-}diagonal}(a, p, q, e) \wedge \\ \quad b \in \operatorname{carr}(a) \wedge \\ \quad e \in \operatorname{mat}(p, q, \operatorname{carr}(a)) \Rightarrow \\ \quad \quad \operatorname{is\text{-}diagonal}(a, (p + 1), (q + 1), \operatorname{border}(a, b, e, p, q)) \end{array}}}
\vnbprop{border-invertible}{\vnbstmt{\begin{array}{@{}l@{}} \forall a, n, u. \\ \quad (\operatorname{is\text{-}ring}(a) \wedge ((n \in \mathbb{N}) \wedge \operatorname{is\text{-}invertible\text{-}mat}(a, n, u))) \Rightarrow  \\ \quad \quad \operatorname{is\text{-}invertible\text{-}mat}(a, (n + 1), \operatorname{border}(a, \operatorname{one}(a), u, n, n)) \end{array}}}
\vnbprop{bordering}{\vnbstmt{\begin{array}{@{}l@{}} \forall a, p, q, b, w, g, e. \\ \quad \operatorname{is\text{-}ring}(a) \wedge \\ \quad p \in \mathbb{N} \wedge \\ \quad q \in \mathbb{N} \wedge \\ \quad b \in \operatorname{carr}(a) \wedge \\ \quad w \in \operatorname{mat}(p, q, \operatorname{carr}(a)) \wedge \\ \quad g \in \operatorname{mat}((p + 1), (q + 1), \operatorname{carr}(a)) \wedge \\ \quad g = \operatorname{border}(a, b, w, p, q) \wedge \\ \quad \operatorname{mat\text{-}equiv}(a, p, q, w, e) \Rightarrow \\ \quad \quad \operatorname{mat\text{-}equiv}(a, (p + 1), (q + 1), g, \operatorname{border}(a, b, e, p, q)) \end{array}}}

\vnbbookfile{theorem-library/bordered-eq-border-proof}{theorem-library/bordered-eq-border-proof}
\noindent{\small bordered-eq-border -- a matrix C whose first row and first column are cleared off the pivot IS literally BORDER(C\_\{1,1\}, SUBMAT(C)): C = [[C11,0],[0,C']]. The bridge from "cleared cross" (clear-first-row + clear-first-col) to the BORDER block form, so the Smith recursion can border a sub-block equivalence}\par
\vnbprop{bordered-eq-border}{\vnbstmt{\begin{array}{@{}l@{}} \forall a, p, q, c. \\ \quad \operatorname{is\text{-}ring}(a) \wedge \\ \quad p \in \mathbb{N} \wedge \\ \quad q \in \mathbb{N} \wedge \\ \quad c \in \operatorname{mat}((p + 1), (q + 1), \operatorname{carr}(a)) \wedge \\ \quad \forall j \in \operatorname{interval}(1, (q + 1)).\; (\neg (j = 1) \Rightarrow (\operatorname{entry}(c, 1, j) = \operatorname{zero}(a))) \wedge \\ \quad \forall i \in \operatorname{interval}(1, (p + 1)).\; (\neg (i = 1) \Rightarrow (\operatorname{entry}(c, i, 1) = \operatorname{zero}(a))) \Rightarrow \\ \quad \quad c = \operatorname{border}(a, \operatorname{entry}(c, 1, 1), \operatorname{submat}(c, p, q), p, q) \end{array}}}

\vnbbookfile{theorem-library/det-alternating-form}{theorem-library/det-alternating-form}
\noindent{\small det-alternating-form.scm -- Hoffman-Kunze 5.3 Theorem 2 (batch 26-B, 2026-09-24): an ALTERNATING n-LINEAR function D on the n x n matrices over a commutative ring R satisfies  D(A) = det(A) D(I)  for every A.  Proven WITHOUT permutations.}\par
\vnbprop{daf-ring-rdist}{\vnbstmt{\begin{array}{@{}l@{}} \forall r, s, x, y, z. \\ \quad (\operatorname{is\text{-}commutative\text{-}ring}(r) \wedge ((s \in \operatorname{carr}(r)) \wedge ((x \in \operatorname{carr}(r)) \wedge ((y \in \operatorname{carr}(r)) \wedge (z \in \operatorname{carr}(r)))))) \Rightarrow  \\ \quad \quad \operatorname{mul}(r)(\operatorname{add}(r)(\operatorname{mul}(r)(s, x), y), z) = \operatorname{add}(r)(\operatorname{mul}(r)(s, \operatorname{mul}(r)(x, z)), \operatorname{mul}(r)(y, z)) \end{array}}}
\vnbprop{daf-ring-cancel0}{\vnbstmt{\begin{array}{@{}l@{}} \forall r, s, x, y. \\ \quad \operatorname{is\text{-}commutative\text{-}ring}(r) \wedge \\ \quad s \in \operatorname{carr}(r) \wedge \\ \quad x \in \operatorname{carr}(r) \wedge \\ \quad y \in \operatorname{carr}(r) \wedge \\ \quad x = \operatorname{add}(r)(\operatorname{mul}(r)(s, \operatorname{zero}(r)), y) \Rightarrow \\ \quad \quad y = x \end{array}}}
\vnbprop{daf-det-fun}{\vnbstmt{\begin{array}{@{}l@{}} \forall r, n. \\ \quad (\operatorname{is\text{-}commutative\text{-}ring}(r) \wedge (n \in \mathbb{N})) \Rightarrow  \\ \quad \quad (\lambda l \in \operatorname{mat}(n, n, \operatorname{carr}(r)).\; \operatorname{det}(r, n, l)) \in (\operatorname{mat}(n, n, \operatorname{carr}(r)) \to \operatorname{carr}(r)) \end{array}}}
\vnbprop{daf-det-lin}{\vnbstmt{\begin{array}{@{}l@{}} \forall r, n, a, b, c, p, s. \\ \quad \operatorname{is\text{-}commutative\text{-}ring}(r) \wedge \\ \quad n \in \mathbb{N} \wedge \\ \quad a \in \operatorname{mat}(n, n, \operatorname{carr}(r)) \wedge \\ \quad b \in \operatorname{mat}(n, n, \operatorname{carr}(r)) \wedge \\ \quad c \in \operatorname{mat}(n, n, \operatorname{carr}(r)) \wedge \\ \quad p \in \operatorname{interval}(1, n) \wedge \\ \quad s \in \operatorname{carr}(r) \wedge \\ \quad \forall i, k \in \operatorname{interval}(1, n). \\ \quad \quad \neg (i = p) \Rightarrow (\operatorname{entry}(b, i, k) = \operatorname{entry}(a, i, k)) \wedge \\ \quad \forall i, k \in \operatorname{interval}(1, n). \\ \quad \quad \neg (i = p) \Rightarrow (\operatorname{entry}(c, i, k) = \operatorname{entry}(a, i, k)) \wedge \\ \quad \forall k \in \operatorname{interval}(1, n). \\ \quad \quad \operatorname{entry}(c, p, k) = \operatorname{add}(r)(\operatorname{mul}(r)(s, \operatorname{entry}(a, p, k)), \operatorname{entry}(b, p, k)) \Rightarrow \\ \quad \quad (\lambda l \in \operatorname{mat}(n, n, \operatorname{carr}(r)).\; \operatorname{det}(r, n, l))(c) \\ \quad \quad =\; \operatorname{add}(r)(\operatorname{mul}(r)(s, (\lambda l \in \operatorname{mat}(n, n, \operatorname{carr}(r)).\; \operatorname{det}(r, n, l))(a)), \\ \quad \quad \quad (\lambda l \in \operatorname{mat}(n, n, \operatorname{carr}(r)).\; \operatorname{det}(r, n, l))(b)) \end{array}}}
\vnbprop{daf-det-alt}{\vnbstmt{\begin{array}{@{}l@{}} \forall r, n, a, p, q. \\ \quad \operatorname{is\text{-}commutative\text{-}ring}(r) \wedge \\ \quad n \in \mathbb{N} \wedge \\ \quad a \in \operatorname{mat}(n, n, \operatorname{carr}(r)) \wedge \\ \quad p \in \operatorname{interval}(1, n) \wedge \\ \quad q \in \operatorname{interval}(1, n) \wedge \\ \quad \neg (p = q) \wedge \\ \quad \forall k \in \operatorname{interval}(1, n).\; (\operatorname{entry}(a, p, k) = \operatorname{entry}(a, q, k)) \Rightarrow \\ \quad \quad (\lambda l \in \operatorname{mat}(n, n, \operatorname{carr}(r)).\; \operatorname{det}(r, n, l))(a) = \operatorname{zero}(r) \end{array}}}
\vnbprop{daf-mm-row-eq}{\vnbstmt{\begin{array}{@{}l@{}} \forall r, n, x, y, q, i, j, k. \\ \quad \operatorname{is\text{-}commutative\text{-}ring}(r) \wedge \\ \quad n \in \mathbb{N} \wedge \\ \quad x \in \operatorname{mat}(n, n, \operatorname{carr}(r)) \wedge \\ \quad y \in \operatorname{mat}(n, n, \operatorname{carr}(r)) \wedge \\ \quad q \in \operatorname{mat}(n, n, \operatorname{carr}(r)) \wedge \\ \quad i \in \operatorname{interval}(1, n) \wedge \\ \quad j \in \operatorname{interval}(1, n) \wedge \\ \quad k \in \operatorname{interval}(1, n) \wedge \\ \quad \forall \mathit{drk} \in \operatorname{interval}(1, n).\; (\operatorname{entry}(x, i, \mathit{drk}) = \operatorname{entry}(y, j, \mathit{drk})) \Rightarrow \\ \quad \quad \operatorname{entry}(\operatorname{matmul}(r, x, q), i, k) = \operatorname{entry}(\operatorname{matmul}(r, y, q), j, k) \end{array}}}
\vnbprop{daf-mm-row-lin}{\vnbstmt{\begin{array}{@{}l@{}} \forall r, n, a, b, c, q, p, s, k. \\ \quad \operatorname{is\text{-}commutative\text{-}ring}(r) \wedge \\ \quad n \in \mathbb{N} \wedge \\ \quad a \in \operatorname{mat}(n, n, \operatorname{carr}(r)) \wedge \\ \quad b \in \operatorname{mat}(n, n, \operatorname{carr}(r)) \wedge \\ \quad c \in \operatorname{mat}(n, n, \operatorname{carr}(r)) \wedge \\ \quad q \in \operatorname{mat}(n, n, \operatorname{carr}(r)) \wedge \\ \quad p \in \operatorname{interval}(1, n) \wedge \\ \quad s \in \operatorname{carr}(r) \wedge \\ \quad k \in \operatorname{interval}(1, n) \wedge \\ \quad \forall \mathit{drk} \in \operatorname{interval}(1, n). \\ \quad \quad \operatorname{entry}(c, p, \mathit{drk}) = \operatorname{add}(r)(\operatorname{mul}(r)(s, \operatorname{entry}(a, p, \mathit{drk})), \operatorname{entry}(b, p, \mathit{drk})) \Rightarrow \\ \quad \quad \operatorname{entry}(\operatorname{matmul}(r, c, q), p, k) \\ \quad \quad =\; \operatorname{add}(r)(\operatorname{mul}(r)(s, \operatorname{entry}(\operatorname{matmul}(r, a, q), p, k)), \operatorname{entry}(\operatorname{matmul}(r, b, q), p, k)) \end{array}}}
\vnbprop{daf-rm-fun}{\vnbstmt{\begin{array}{@{}l@{}} \forall r, n, d, q. \\ \quad \operatorname{is\text{-}commutative\text{-}ring}(r) \wedge \\ \quad n \in \mathbb{N} \wedge \\ \quad d \in (\operatorname{mat}(n, n, \operatorname{carr}(r)) \to \operatorname{carr}(r)) \wedge \\ \quad q \in \operatorname{mat}(n, n, \operatorname{carr}(r)) \Rightarrow \\ \quad \quad (\lambda l \in \operatorname{mat}(n, n, \operatorname{carr}(r)).\; d(\operatorname{matmul}(r, l, q))) \in (\operatorname{mat}(n, n, \operatorname{carr}(r)) \to \operatorname{carr}(r)) \end{array}}}
\vnbprop{daf-rm-lin}{\vnbstmt{\begin{array}{@{}l@{}} \forall r, n, d, q, a, b, c, p, s. \\ \quad \operatorname{is\text{-}commutative\text{-}ring}(r) \wedge \\ \quad n \in \mathbb{N} \wedge \\ \quad d \in (\operatorname{mat}(n, n, \operatorname{carr}(r)) \to \operatorname{carr}(r)) \wedge \\ \quad \forall a, b, c, p, s. \\ \quad \quad a \in \operatorname{mat}(n, n, \operatorname{carr}(r)) \wedge \\ \quad \quad b \in \operatorname{mat}(n, n, \operatorname{carr}(r)) \wedge \\ \quad \quad c \in \operatorname{mat}(n, n, \operatorname{carr}(r)) \wedge \\ \quad \quad p \in \operatorname{interval}(1, n) \wedge \\ \quad \quad s \in \operatorname{carr}(r) \wedge \\ \quad \quad \forall i, k \in \operatorname{interval}(1, n). \\ \quad \quad \quad \neg (i = p) \Rightarrow (\operatorname{entry}(b, i, k) = \operatorname{entry}(a, i, k)) \wedge \\ \quad \quad \forall i, k \in \operatorname{interval}(1, n). \\ \quad \quad \quad \neg (i = p) \Rightarrow (\operatorname{entry}(c, i, k) = \operatorname{entry}(a, i, k)) \wedge \\ \quad \quad \forall k \in \operatorname{interval}(1, n). \\ \quad \quad \quad \operatorname{entry}(c, p, k) = \operatorname{add}(r)(\operatorname{mul}(r)(s, \operatorname{entry}(a, p, k)), \operatorname{entry}(b, p, k)) \Rightarrow \\ \quad \quad \quad d(c) = \operatorname{add}(r)(\operatorname{mul}(r)(s, d(a)), d(b)) \wedge \\ \quad q \in \operatorname{mat}(n, n, \operatorname{carr}(r)) \wedge \\ \quad a \in \operatorname{mat}(n, n, \operatorname{carr}(r)) \wedge \\ \quad b \in \operatorname{mat}(n, n, \operatorname{carr}(r)) \wedge \\ \quad c \in \operatorname{mat}(n, n, \operatorname{carr}(r)) \wedge \\ \quad p \in \operatorname{interval}(1, n) \wedge \\ \quad s \in \operatorname{carr}(r) \wedge \\ \quad \forall i, k \in \operatorname{interval}(1, n). \\ \quad \quad \neg (i = p) \Rightarrow (\operatorname{entry}(b, i, k) = \operatorname{entry}(a, i, k)) \wedge \\ \quad \forall i, k \in \operatorname{interval}(1, n). \\ \quad \quad \neg (i = p) \Rightarrow (\operatorname{entry}(c, i, k) = \operatorname{entry}(a, i, k)) \wedge \\ \quad \forall k \in \operatorname{interval}(1, n). \\ \quad \quad \operatorname{entry}(c, p, k) = \operatorname{add}(r)(\operatorname{mul}(r)(s, \operatorname{entry}(a, p, k)), \operatorname{entry}(b, p, k)) \Rightarrow \\ \quad \quad (\lambda l \in \operatorname{mat}(n, n, \operatorname{carr}(r)).\; d(\operatorname{matmul}(r, l, q)))(c) \\ \quad \quad =\; \operatorname{add}(r)(\operatorname{mul}(r)(s, (\lambda l \in \operatorname{mat}(n, n, \operatorname{carr}(r)).\; d(\operatorname{matmul}(r, l, q)))(a)), \\ \quad \quad \quad (\lambda l \in \operatorname{mat}(n, n, \operatorname{carr}(r)).\; d(\operatorname{matmul}(r, l, q)))(b)) \end{array}}}
\vnbprop{daf-rm-alt}{\vnbstmt{\begin{array}{@{}l@{}} \forall r, n, d, q, a, p, \mathit{dfq}. \\ \quad \operatorname{is\text{-}commutative\text{-}ring}(r) \wedge \\ \quad n \in \mathbb{N} \wedge \\ \quad d \in (\operatorname{mat}(n, n, \operatorname{carr}(r)) \to \operatorname{carr}(r)) \wedge \\ \quad \forall a, p, \mathit{dfq}. \\ \quad \quad a \in \operatorname{mat}(n, n, \operatorname{carr}(r)) \wedge \\ \quad \quad p \in \operatorname{interval}(1, n) \wedge \\ \quad \quad \mathit{dfq} \in \operatorname{interval}(1, n) \wedge \\ \quad \quad \neg (p = \mathit{dfq}) \wedge \\ \quad \quad \forall k \in \operatorname{interval}(1, n).\; (\operatorname{entry}(a, p, k) = \operatorname{entry}(a, \mathit{dfq}, k)) \Rightarrow \\ \quad \quad \quad d(a) = \operatorname{zero}(r) \wedge \\ \quad q \in \operatorname{mat}(n, n, \operatorname{carr}(r)) \wedge \\ \quad a \in \operatorname{mat}(n, n, \operatorname{carr}(r)) \wedge \\ \quad p \in \operatorname{interval}(1, n) \wedge \\ \quad \mathit{dfq} \in \operatorname{interval}(1, n) \wedge \\ \quad \neg (p = \mathit{dfq}) \wedge \\ \quad \forall k \in \operatorname{interval}(1, n).\; (\operatorname{entry}(a, p, k) = \operatorname{entry}(a, \mathit{dfq}, k)) \Rightarrow \\ \quad \quad (\lambda l \in \operatorname{mat}(n, n, \operatorname{carr}(r)).\; d(\operatorname{matmul}(r, l, q)))(a) = \operatorname{zero}(r) \end{array}}}
\vnbprop{daf-det-rm-fun}{\vnbstmt{\begin{array}{@{}l@{}} \forall r, n, q. \\ \quad (\operatorname{is\text{-}commutative\text{-}ring}(r) \wedge ((n \in \mathbb{N}) \wedge (q \in \operatorname{mat}(n, n, \operatorname{carr}(r))))) \Rightarrow  \\ \quad \quad (\lambda l \in \operatorname{mat}(n, n, \operatorname{carr}(r)).\; \operatorname{det}(r, n, \operatorname{matmul}(r, l, q))) \in (\operatorname{mat}(n, n, \operatorname{carr}(r)) \to \operatorname{carr}(r)) \end{array}}}
\vnbprop{daf-det-rm-lin}{\vnbstmt{\begin{array}{@{}l@{}} \forall r, n, q, a, b, c, p, s. \\ \quad \operatorname{is\text{-}commutative\text{-}ring}(r) \wedge \\ \quad n \in \mathbb{N} \wedge \\ \quad q \in \operatorname{mat}(n, n, \operatorname{carr}(r)) \wedge \\ \quad a \in \operatorname{mat}(n, n, \operatorname{carr}(r)) \wedge \\ \quad b \in \operatorname{mat}(n, n, \operatorname{carr}(r)) \wedge \\ \quad c \in \operatorname{mat}(n, n, \operatorname{carr}(r)) \wedge \\ \quad p \in \operatorname{interval}(1, n) \wedge \\ \quad s \in \operatorname{carr}(r) \wedge \\ \quad \forall i, k \in \operatorname{interval}(1, n). \\ \quad \quad \neg (i = p) \Rightarrow (\operatorname{entry}(b, i, k) = \operatorname{entry}(a, i, k)) \wedge \\ \quad \forall i, k \in \operatorname{interval}(1, n). \\ \quad \quad \neg (i = p) \Rightarrow (\operatorname{entry}(c, i, k) = \operatorname{entry}(a, i, k)) \wedge \\ \quad \forall k \in \operatorname{interval}(1, n). \\ \quad \quad \operatorname{entry}(c, p, k) = \operatorname{add}(r)(\operatorname{mul}(r)(s, \operatorname{entry}(a, p, k)), \operatorname{entry}(b, p, k)) \Rightarrow \\ \quad \quad (\lambda l \in \operatorname{mat}(n, n, \operatorname{carr}(r)).\; \operatorname{det}(r, n, \operatorname{matmul}(r, l, q)))(c) \\ \quad \quad =\; \operatorname{add}(r)(\operatorname{mul}(r)(s, (\lambda l \in \operatorname{mat}(n, n, \operatorname{carr}(r)).\; \operatorname{det}(r, n, \operatorname{matmul}(r, l, q)))(a)), \\ \quad \quad \quad (\lambda l \in \operatorname{mat}(n, n, \operatorname{carr}(r)).\; \operatorname{det}(r, n, \operatorname{matmul}(r, l, q)))(b)) \end{array}}}
\vnbprop{daf-det-rm-alt}{\vnbstmt{\begin{array}{@{}l@{}} \forall r, n, q, a, p, \mathit{dfq}. \\ \quad \operatorname{is\text{-}commutative\text{-}ring}(r) \wedge \\ \quad n \in \mathbb{N} \wedge \\ \quad q \in \operatorname{mat}(n, n, \operatorname{carr}(r)) \wedge \\ \quad a \in \operatorname{mat}(n, n, \operatorname{carr}(r)) \wedge \\ \quad p \in \operatorname{interval}(1, n) \wedge \\ \quad \mathit{dfq} \in \operatorname{interval}(1, n) \wedge \\ \quad \neg (p = \mathit{dfq}) \wedge \\ \quad \forall k \in \operatorname{interval}(1, n).\; (\operatorname{entry}(a, p, k) = \operatorname{entry}(a, \mathit{dfq}, k)) \Rightarrow \\ \quad \quad (\lambda l \in \operatorname{mat}(n, n, \operatorname{carr}(r)).\; \operatorname{det}(r, n, \operatorname{matmul}(r, l, q)))(a) = \operatorname{zero}(r) \end{array}}}
\vnbprop{daf-zero-row}{\vnbstmt{\begin{array}{@{}l@{}} \forall r, n, d, a, k. \\ \quad \operatorname{is\text{-}commutative\text{-}ring}(r) \wedge \\ \quad n \in \mathbb{N} \wedge \\ \quad d \in (\operatorname{mat}(n, n, \operatorname{carr}(r)) \to \operatorname{carr}(r)) \wedge \\ \quad \forall \mathit{dfa}, b, c, p, s. \\ \quad \quad \mathit{dfa} \in \operatorname{mat}(n, n, \operatorname{carr}(r)) \wedge \\ \quad \quad b \in \operatorname{mat}(n, n, \operatorname{carr}(r)) \wedge \\ \quad \quad c \in \operatorname{mat}(n, n, \operatorname{carr}(r)) \wedge \\ \quad \quad p \in \operatorname{interval}(1, n) \wedge \\ \quad \quad s \in \operatorname{carr}(r) \wedge \\ \quad \quad \forall i, \mathit{drk} \in \operatorname{interval}(1, n). \\ \quad \quad \quad \neg (i = p) \Rightarrow (\operatorname{entry}(b, i, \mathit{drk}) = \operatorname{entry}(\mathit{dfa}, i, \mathit{drk})) \wedge \\ \quad \quad \forall i, \mathit{drk} \in \operatorname{interval}(1, n). \\ \quad \quad \quad \neg (i = p) \Rightarrow (\operatorname{entry}(c, i, \mathit{drk}) = \operatorname{entry}(\mathit{dfa}, i, \mathit{drk})) \wedge \\ \quad \quad \forall \mathit{drk} \in \operatorname{interval}(1, n). \\ \quad \quad \quad \operatorname{entry}(c, p, \mathit{drk}) = \operatorname{add}(r)(\operatorname{mul}(r)(s, \operatorname{entry}(\mathit{dfa}, p, \mathit{drk})), \operatorname{entry}(b, p, \mathit{drk})) \Rightarrow \\ \quad \quad \quad d(c) = \operatorname{add}(r)(\operatorname{mul}(r)(s, d(\mathit{dfa})), d(b)) \wedge \\ \quad a \in \operatorname{mat}(n, n, \operatorname{carr}(r)) \wedge \\ \quad k \in \operatorname{interval}(1, n) \wedge \\ \quad \forall \mathit{drk} \in \operatorname{interval}(1, n).\; (\operatorname{entry}(a, k, \mathit{drk}) = \operatorname{zero}(r)) \Rightarrow \\ \quad \quad d(a) = \operatorname{zero}(r) \end{array}}}
\vnbprop{daf-bd-fun}{\vnbstmt{\begin{array}{@{}l@{}} \forall r, n, d. \\ \quad (\operatorname{is\text{-}commutative\text{-}ring}(r) \wedge ((n \in \mathbb{N}) \wedge (d \in (\operatorname{mat}((n + 1), (n + 1), \operatorname{carr}(r)) \to \operatorname{carr}(r))))) \Rightarrow  \\ \quad \quad (\lambda l \in \operatorname{mat}(n, n, \operatorname{carr}(r)).\; d(\operatorname{border}(r, \operatorname{one}(r), l, n, n))) \in (\operatorname{mat}(n, n, \operatorname{carr}(r)) \to \operatorname{carr}(r)) \end{array}}}
\vnbprop{daf-bd-lin}{\vnbstmt{\begin{array}{@{}l@{}} \forall r, n, d, a, b, c, p, s. \\ \quad \operatorname{is\text{-}commutative\text{-}ring}(r) \wedge \\ \quad n \in \mathbb{N} \wedge \\ \quad d \in (\operatorname{mat}((n + 1), (n + 1), \operatorname{carr}(r)) \to \operatorname{carr}(r)) \wedge \\ \quad \forall a, b, c, p, s. \\ \quad \quad a \in \operatorname{mat}((n + 1), (n + 1), \operatorname{carr}(r)) \wedge \\ \quad \quad b \in \operatorname{mat}((n + 1), (n + 1), \operatorname{carr}(r)) \wedge \\ \quad \quad c \in \operatorname{mat}((n + 1), (n + 1), \operatorname{carr}(r)) \wedge \\ \quad \quad p \in \operatorname{interval}(1, (n + 1)) \wedge \\ \quad \quad s \in \operatorname{carr}(r) \wedge \\ \quad \quad \forall i, k \in \operatorname{interval}(1, (n + 1)). \\ \quad \quad \quad \neg (i = p) \Rightarrow (\operatorname{entry}(b, i, k) = \operatorname{entry}(a, i, k)) \wedge \\ \quad \quad \forall i, k \in \operatorname{interval}(1, (n + 1)). \\ \quad \quad \quad \neg (i = p) \Rightarrow (\operatorname{entry}(c, i, k) = \operatorname{entry}(a, i, k)) \wedge \\ \quad \quad \forall k \in \operatorname{interval}(1, (n + 1)). \\ \quad \quad \quad \operatorname{entry}(c, p, k) = \operatorname{add}(r)(\operatorname{mul}(r)(s, \operatorname{entry}(a, p, k)), \operatorname{entry}(b, p, k)) \Rightarrow \\ \quad \quad \quad d(c) = \operatorname{add}(r)(\operatorname{mul}(r)(s, d(a)), d(b)) \wedge \\ \quad a \in \operatorname{mat}(n, n, \operatorname{carr}(r)) \wedge \\ \quad b \in \operatorname{mat}(n, n, \operatorname{carr}(r)) \wedge \\ \quad c \in \operatorname{mat}(n, n, \operatorname{carr}(r)) \wedge \\ \quad p \in \operatorname{interval}(1, n) \wedge \\ \quad s \in \operatorname{carr}(r) \wedge \\ \quad \forall i, k \in \operatorname{interval}(1, n). \\ \quad \quad \neg (i = p) \Rightarrow (\operatorname{entry}(b, i, k) = \operatorname{entry}(a, i, k)) \wedge \\ \quad \forall i, k \in \operatorname{interval}(1, n). \\ \quad \quad \neg (i = p) \Rightarrow (\operatorname{entry}(c, i, k) = \operatorname{entry}(a, i, k)) \wedge \\ \quad \forall k \in \operatorname{interval}(1, n). \\ \quad \quad \operatorname{entry}(c, p, k) = \operatorname{add}(r)(\operatorname{mul}(r)(s, \operatorname{entry}(a, p, k)), \operatorname{entry}(b, p, k)) \Rightarrow \\ \quad \quad (\lambda l \in \operatorname{mat}(n, n, \operatorname{carr}(r)).\; d(\operatorname{border}(r, \operatorname{one}(r), l, n, n)))(c) \\ \quad \quad =\; \operatorname{add}(r)(\operatorname{mul}(r)(s, (\lambda l \in \operatorname{mat}(n, n, \operatorname{carr}(r)).\; d(\operatorname{border}(r, \operatorname{one}(r), l, n, n)))(a)), \\ \quad \quad \quad (\lambda l \in \operatorname{mat}(n, n, \operatorname{carr}(r)).\; d(\operatorname{border}(r, \operatorname{one}(r), l, n, n)))(b)) \end{array}}}
\vnbprop{daf-bd-alt}{\vnbstmt{\begin{array}{@{}l@{}} \forall r, n, d, a, p, q. \\ \quad \operatorname{is\text{-}commutative\text{-}ring}(r) \wedge \\ \quad n \in \mathbb{N} \wedge \\ \quad d \in (\operatorname{mat}((n + 1), (n + 1), \operatorname{carr}(r)) \to \operatorname{carr}(r)) \wedge \\ \quad \forall a, p, q. \\ \quad \quad a \in \operatorname{mat}((n + 1), (n + 1), \operatorname{carr}(r)) \wedge \\ \quad \quad p \in \operatorname{interval}(1, (n + 1)) \wedge \\ \quad \quad q \in \operatorname{interval}(1, (n + 1)) \wedge \\ \quad \quad \neg (p = q) \wedge \\ \quad \quad \forall k \in \operatorname{interval}(1, (n + 1)).\; (\operatorname{entry}(a, p, k) = \operatorname{entry}(a, q, k)) \Rightarrow \\ \quad \quad \quad d(a) = \operatorname{zero}(r) \wedge \\ \quad a \in \operatorname{mat}(n, n, \operatorname{carr}(r)) \wedge \\ \quad p \in \operatorname{interval}(1, n) \wedge \\ \quad q \in \operatorname{interval}(1, n) \wedge \\ \quad \neg (p = q) \wedge \\ \quad \forall k \in \operatorname{interval}(1, n).\; (\operatorname{entry}(a, p, k) = \operatorname{entry}(a, q, k)) \Rightarrow \\ \quad \quad (\lambda l \in \operatorname{mat}(n, n, \operatorname{carr}(r)).\; d(\operatorname{border}(r, \operatorname{one}(r), l, n, n)))(a) = \operatorname{zero}(r) \end{array}}}
\vnbprop{daf-swap}{\vnbstmt{\begin{array}{@{}l@{}} \forall r, n, d, a, b, p, q. \\ \quad \operatorname{is\text{-}commutative\text{-}ring}(r) \wedge \\ \quad n \in \mathbb{N} \wedge \\ \quad d \in (\operatorname{mat}(n, n, \operatorname{carr}(r)) \to \operatorname{carr}(r)) \wedge \\ \quad \forall \mathit{dfa}, \mathit{dfb}, c, \mathit{dfp}, s. \\ \quad \quad \mathit{dfa} \in \operatorname{mat}(n, n, \operatorname{carr}(r)) \wedge \\ \quad \quad \mathit{dfb} \in \operatorname{mat}(n, n, \operatorname{carr}(r)) \wedge \\ \quad \quad c \in \operatorname{mat}(n, n, \operatorname{carr}(r)) \wedge \\ \quad \quad \mathit{dfp} \in \operatorname{interval}(1, n) \wedge \\ \quad \quad s \in \operatorname{carr}(r) \wedge \\ \quad \quad \forall i, k \in \operatorname{interval}(1, n). \\ \quad \quad \quad \neg (i = \mathit{dfp}) \Rightarrow (\operatorname{entry}(\mathit{dfb}, i, k) = \operatorname{entry}(\mathit{dfa}, i, k)) \wedge \\ \quad \quad \forall i, k \in \operatorname{interval}(1, n). \\ \quad \quad \quad \neg (i = \mathit{dfp}) \Rightarrow (\operatorname{entry}(c, i, k) = \operatorname{entry}(\mathit{dfa}, i, k)) \wedge \\ \quad \quad \forall k \in \operatorname{interval}(1, n). \\ \quad \quad \quad \operatorname{entry}(c, \mathit{dfp}, k) = \operatorname{add}(r)(\operatorname{mul}(r)(s, \operatorname{entry}(\mathit{dfa}, \mathit{dfp}, k)), \operatorname{entry}(\mathit{dfb}, \mathit{dfp}, k)) \Rightarrow \\ \quad \quad \quad d(c) = \operatorname{add}(r)(\operatorname{mul}(r)(s, d(\mathit{dfa})), d(\mathit{dfb})) \wedge \\ \quad \forall \mathit{dfa}, \mathit{dfp}, \mathit{dfq}. \\ \quad \quad \mathit{dfa} \in \operatorname{mat}(n, n, \operatorname{carr}(r)) \wedge \\ \quad \quad \mathit{dfp} \in \operatorname{interval}(1, n) \wedge \\ \quad \quad \mathit{dfq} \in \operatorname{interval}(1, n) \wedge \\ \quad \quad \neg (\mathit{dfp} = \mathit{dfq}) \wedge \\ \quad \quad \forall k \in \operatorname{interval}(1, n).\; (\operatorname{entry}(\mathit{dfa}, \mathit{dfp}, k) = \operatorname{entry}(\mathit{dfa}, \mathit{dfq}, k)) \Rightarrow \\ \quad \quad \quad d(\mathit{dfa}) = \operatorname{zero}(r) \wedge \\ \quad a \in \operatorname{mat}(n, n, \operatorname{carr}(r)) \wedge \\ \quad b \in \operatorname{mat}(n, n, \operatorname{carr}(r)) \wedge \\ \quad p \in \operatorname{interval}(1, n) \wedge \\ \quad q \in \operatorname{interval}(1, n) \wedge \\ \quad \neg (p = q) \wedge \\ \quad \forall i, k \in \operatorname{interval}(1, n). \\ \quad \quad (\neg (i = p) \wedge \neg (i = q)) \Rightarrow (\operatorname{entry}(b, i, k) = \operatorname{entry}(a, i, k)) \wedge \\ \quad \forall k \in \operatorname{interval}(1, n).\; (\operatorname{entry}(b, p, k) = \operatorname{entry}(a, q, k)) \wedge \\ \quad \forall k \in \operatorname{interval}(1, n).\; (\operatorname{entry}(b, q, k) = \operatorname{entry}(a, p, k)) \Rightarrow \\ \quad \quad d(b) = \operatorname{neg}(r)(d(a)) \end{array}}}
\vnbprop{daf-elem-f}{\vnbstmt{\begin{array}{@{}l@{}} \forall r, n, d, k, l. \\ \quad \operatorname{is\text{-}commutative\text{-}ring}(r) \wedge \\ \quad n \in \mathbb{N} \wedge \\ \quad d \in (\operatorname{mat}(n, n, \operatorname{carr}(r)) \to \operatorname{carr}(r)) \wedge \\ \quad \forall a, b, c, p, s. \\ \quad \quad a \in \operatorname{mat}(n, n, \operatorname{carr}(r)) \wedge \\ \quad \quad b \in \operatorname{mat}(n, n, \operatorname{carr}(r)) \wedge \\ \quad \quad c \in \operatorname{mat}(n, n, \operatorname{carr}(r)) \wedge \\ \quad \quad p \in \operatorname{interval}(1, n) \wedge \\ \quad \quad s \in \operatorname{carr}(r) \wedge \\ \quad \quad \forall i, \mathit{drk} \in \operatorname{interval}(1, n). \\ \quad \quad \quad \neg (i = p) \Rightarrow (\operatorname{entry}(b, i, \mathit{drk}) = \operatorname{entry}(a, i, \mathit{drk})) \wedge \\ \quad \quad \forall i, \mathit{drk} \in \operatorname{interval}(1, n). \\ \quad \quad \quad \neg (i = p) \Rightarrow (\operatorname{entry}(c, i, \mathit{drk}) = \operatorname{entry}(a, i, \mathit{drk})) \wedge \\ \quad \quad \forall \mathit{drk} \in \operatorname{interval}(1, n). \\ \quad \quad \quad \operatorname{entry}(c, p, \mathit{drk}) = \operatorname{add}(r)(\operatorname{mul}(r)(s, \operatorname{entry}(a, p, \mathit{drk})), \operatorname{entry}(b, p, \mathit{drk})) \Rightarrow \\ \quad \quad \quad d(c) = \operatorname{add}(r)(\operatorname{mul}(r)(s, d(a)), d(b)) \wedge \\ \quad \forall a, p, q. \\ \quad \quad a \in \operatorname{mat}(n, n, \operatorname{carr}(r)) \wedge \\ \quad \quad p \in \operatorname{interval}(1, n) \wedge \\ \quad \quad q \in \operatorname{interval}(1, n) \wedge \\ \quad \quad \neg (p = q) \wedge \\ \quad \quad \forall \mathit{drk} \in \operatorname{interval}(1, n).\; (\operatorname{entry}(a, p, \mathit{drk}) = \operatorname{entry}(a, q, \mathit{drk})) \Rightarrow \\ \quad \quad \quad d(a) = \operatorname{zero}(r) \wedge \\ \quad k \in \operatorname{interval}(1, n) \wedge \\ \quad l \in \operatorname{interval}(1, n) \wedge \\ \quad \neg (k = l) \Rightarrow \\ \quad \quad d(\operatorname{elem\text{-}f}(r, n, k, l)) = \operatorname{neg}(r)(d(\operatorname{identmat}(r, n))) \end{array}}}
\vnbprop{daf-clear-ind}{\vnbstmt{\begin{array}{@{}l@{}} \forall m \in \mathbb{N}, r, n, d, y. \\ \quad \operatorname{is\text{-}commutative\text{-}ring}(r) \wedge \\ \quad n \in \mathbb{N} \wedge \\ \quad d \in (\operatorname{mat}((n + 1), (n + 1), \operatorname{carr}(r)) \to \operatorname{carr}(r)) \wedge \\ \quad \forall a, b, c, p, s. \\ \quad \quad a \in \operatorname{mat}((n + 1), (n + 1), \operatorname{carr}(r)) \wedge \\ \quad \quad b \in \operatorname{mat}((n + 1), (n + 1), \operatorname{carr}(r)) \wedge \\ \quad \quad c \in \operatorname{mat}((n + 1), (n + 1), \operatorname{carr}(r)) \wedge \\ \quad \quad p \in \operatorname{interval}(1, (n + 1)) \wedge \\ \quad \quad s \in \operatorname{carr}(r) \wedge \\ \quad \quad \forall i, k \in \operatorname{interval}(1, (n + 1)). \\ \quad \quad \quad \neg (i = p) \Rightarrow (\operatorname{entry}(b, i, k) = \operatorname{entry}(a, i, k)) \wedge \\ \quad \quad \forall i, k \in \operatorname{interval}(1, (n + 1)). \\ \quad \quad \quad \neg (i = p) \Rightarrow (\operatorname{entry}(c, i, k) = \operatorname{entry}(a, i, k)) \wedge \\ \quad \quad \forall k \in \operatorname{interval}(1, (n + 1)). \\ \quad \quad \quad \operatorname{entry}(c, p, k) = \operatorname{add}(r)(\operatorname{mul}(r)(s, \operatorname{entry}(a, p, k)), \operatorname{entry}(b, p, k)) \Rightarrow \\ \quad \quad \quad d(c) = \operatorname{add}(r)(\operatorname{mul}(r)(s, d(a)), d(b)) \wedge \\ \quad \forall a, p, q. \\ \quad \quad a \in \operatorname{mat}((n + 1), (n + 1), \operatorname{carr}(r)) \wedge \\ \quad \quad p \in \operatorname{interval}(1, (n + 1)) \wedge \\ \quad \quad q \in \operatorname{interval}(1, (n + 1)) \wedge \\ \quad \quad \neg (p = q) \wedge \\ \quad \quad \forall k \in \operatorname{interval}(1, (n + 1)).\; (\operatorname{entry}(a, p, k) = \operatorname{entry}(a, q, k)) \Rightarrow \\ \quad \quad \quad d(a) = \operatorname{zero}(r) \wedge \\ \quad y \in \operatorname{mat}((n + 1), (n + 1), \operatorname{carr}(r)) \wedge \\ \quad \operatorname{entry}(y, 1, 1) = \operatorname{one}(r) \wedge \\ \quad \forall k \in \operatorname{interval}(1, (n + 1)).\; (\neg (k = 1) \Rightarrow (\operatorname{entry}(y, 1, k) = \operatorname{zero}(r))) \wedge \\ \quad m \leq n \Rightarrow \\ \quad \quad d(\operatorname{matof}((n + 1), (n + 1), \\ \quad \quad \quad \quad (\lambda \langle \operatorname{dav\_}, \operatorname{daw\_} \rangle \in (\operatorname{interval}(1, (n + 1)) \times \operatorname{interval}(1, (n + 1))).\; \\ \quad \quad \quad \quad \quad \left(\text{if }(\operatorname{daw\_} = 1)\text{ then }\left(\text{if }(\operatorname{dav\_} = 1)\text{ then }\operatorname{entry}(y, \operatorname{dav\_}, \operatorname{daw\_})\text{ else }\left(\text{if }(\operatorname{dav\_} \leq (m + 1))\text{ then }\operatorname{zero}(r)\text{ else }\operatorname{entry}(y, \operatorname{dav\_}, \operatorname{daw\_})\right)\right)\text{ else }\operatorname{entry}(y, \operatorname{dav\_}, \operatorname{daw\_})\right)))) \\ \quad \quad =\; d(y) \end{array}}}
\vnbprop{daf-clear}{\vnbstmt{\begin{array}{@{}l@{}} \forall r, n, d, y. \\ \quad \operatorname{is\text{-}commutative\text{-}ring}(r) \wedge \\ \quad n \in \mathbb{N} \wedge \\ \quad d \in (\operatorname{mat}((n + 1), (n + 1), \operatorname{carr}(r)) \to \operatorname{carr}(r)) \wedge \\ \quad \forall a, b, c, p, s. \\ \quad \quad a \in \operatorname{mat}((n + 1), (n + 1), \operatorname{carr}(r)) \wedge \\ \quad \quad b \in \operatorname{mat}((n + 1), (n + 1), \operatorname{carr}(r)) \wedge \\ \quad \quad c \in \operatorname{mat}((n + 1), (n + 1), \operatorname{carr}(r)) \wedge \\ \quad \quad p \in \operatorname{interval}(1, (n + 1)) \wedge \\ \quad \quad s \in \operatorname{carr}(r) \wedge \\ \quad \quad \forall i, k \in \operatorname{interval}(1, (n + 1)). \\ \quad \quad \quad \neg (i = p) \Rightarrow (\operatorname{entry}(b, i, k) = \operatorname{entry}(a, i, k)) \wedge \\ \quad \quad \forall i, k \in \operatorname{interval}(1, (n + 1)). \\ \quad \quad \quad \neg (i = p) \Rightarrow (\operatorname{entry}(c, i, k) = \operatorname{entry}(a, i, k)) \wedge \\ \quad \quad \forall k \in \operatorname{interval}(1, (n + 1)). \\ \quad \quad \quad \operatorname{entry}(c, p, k) = \operatorname{add}(r)(\operatorname{mul}(r)(s, \operatorname{entry}(a, p, k)), \operatorname{entry}(b, p, k)) \Rightarrow \\ \quad \quad \quad d(c) = \operatorname{add}(r)(\operatorname{mul}(r)(s, d(a)), d(b)) \wedge \\ \quad \forall a, p, q. \\ \quad \quad a \in \operatorname{mat}((n + 1), (n + 1), \operatorname{carr}(r)) \wedge \\ \quad \quad p \in \operatorname{interval}(1, (n + 1)) \wedge \\ \quad \quad q \in \operatorname{interval}(1, (n + 1)) \wedge \\ \quad \quad \neg (p = q) \wedge \\ \quad \quad \forall k \in \operatorname{interval}(1, (n + 1)).\; (\operatorname{entry}(a, p, k) = \operatorname{entry}(a, q, k)) \Rightarrow \\ \quad \quad \quad d(a) = \operatorname{zero}(r) \wedge \\ \quad y \in \operatorname{mat}((n + 1), (n + 1), \operatorname{carr}(r)) \wedge \\ \quad \operatorname{entry}(y, 1, 1) = \operatorname{one}(r) \wedge \\ \quad \forall k \in \operatorname{interval}(1, (n + 1)).\; (\neg (k = 1) \Rightarrow (\operatorname{entry}(y, 1, k) = \operatorname{zero}(r))) \Rightarrow \\ \quad \quad d(\operatorname{matof}((n + 1), (n + 1), \\ \quad \quad \quad \quad (\lambda \langle \operatorname{dav\_}, \operatorname{daw\_} \rangle \in (\operatorname{interval}(1, (n + 1)) \times \operatorname{interval}(1, (n + 1))).\; \\ \quad \quad \quad \quad \quad \left(\text{if }(\operatorname{daw\_} = 1)\text{ then }\left(\text{if }(\operatorname{dav\_} = 1)\text{ then }\operatorname{entry}(y, \operatorname{dav\_}, \operatorname{daw\_})\text{ else }\left(\text{if }(\operatorname{dav\_} \leq (n + 1))\text{ then }\operatorname{zero}(r)\text{ else }\operatorname{entry}(y, \operatorname{dav\_}, \operatorname{daw\_})\right)\right)\text{ else }\operatorname{entry}(y, \operatorname{dav\_}, \operatorname{daw\_})\right)))) \\ \quad \quad =\; d(y) \end{array}}}
\vnbprop{daf-bordered}{\vnbstmt{\begin{array}{@{}l@{}} \forall r, n, y. \\ \quad \operatorname{is\text{-}commutative\text{-}ring}(r) \wedge \\ \quad n \in \mathbb{N} \wedge \\ \quad y \in \operatorname{mat}((n + 1), (n + 1), \operatorname{carr}(r)) \wedge \\ \quad \operatorname{entry}(y, 1, 1) = \operatorname{one}(r) \wedge \\ \quad \forall k \in \operatorname{interval}(1, (n + 1)).\; (\neg (k = 1) \Rightarrow (\operatorname{entry}(y, 1, k) = \operatorname{zero}(r))) \Rightarrow \\ \quad \quad \operatorname{matof}((n + 1), (n + 1), \\ \quad \quad \quad (\lambda \langle \operatorname{dav\_}, \operatorname{daw\_} \rangle \in (\operatorname{interval}(1, (n + 1)) \times \operatorname{interval}(1, (n + 1))).\; \\ \quad \quad \quad \quad \left(\text{if }(\operatorname{daw\_} = 1)\text{ then }\left(\text{if }(\operatorname{dav\_} = 1)\text{ then }\operatorname{entry}(y, \operatorname{dav\_}, \operatorname{daw\_})\text{ else }\left(\text{if }(\operatorname{dav\_} \leq (n + 1))\text{ then }\operatorname{zero}(r)\text{ else }\operatorname{entry}(y, \operatorname{dav\_}, \operatorname{daw\_})\right)\right)\text{ else }\operatorname{entry}(y, \operatorname{dav\_}, \operatorname{daw\_})\right))) \\ \quad \quad =\; \operatorname{border}(r, \operatorname{one}(r), \\ \quad \quad \quad \operatorname{submat}(\operatorname{matof}((n + 1), (n + 1), \\ \quad \quad \quad \quad \quad (\lambda \langle \operatorname{dav\_}, \operatorname{daw\_} \rangle \in (\operatorname{interval}(1, (n + 1)) \times \operatorname{interval}(1, (n + 1))).\; \\ \quad \quad \quad \quad \quad \quad \left(\text{if }(\operatorname{daw\_} = 1)\text{ then }\left(\text{if }(\operatorname{dav\_} = 1)\text{ then }\operatorname{entry}(y, \operatorname{dav\_}, \operatorname{daw\_})\text{ else }\left(\text{if }(\operatorname{dav\_} \leq (n + 1))\text{ then }\operatorname{zero}(r)\text{ else }\operatorname{entry}(y, \operatorname{dav\_}, \operatorname{daw\_})\right)\right)\text{ else }\operatorname{entry}(y, \operatorname{dav\_}, \operatorname{daw\_})\right))), \\ \quad \quad \quad \quad n, n), \\ \quad \quad \quad n, n) \end{array}}}
\vnbprop{daf-cl}{\vnbstmt{\begin{array}{@{}l@{}} \forall n \in \mathbb{N}, r, d, y. \\ \quad \forall r, d.; \forall x \in \operatorname{mat}(n, n, \operatorname{carr}(r)). \\ \quad \quad \operatorname{is\text{-}commutative\text{-}ring}(r) \wedge \\ \quad \quad d \in (\operatorname{mat}(n, n, \operatorname{carr}(r)) \to \operatorname{carr}(r)) \wedge \\ \quad \quad \forall a, b, c, p, s. \\ \quad \quad \quad a \in \operatorname{mat}(n, n, \operatorname{carr}(r)) \wedge \\ \quad \quad \quad b \in \operatorname{mat}(n, n, \operatorname{carr}(r)) \wedge \\ \quad \quad \quad c \in \operatorname{mat}(n, n, \operatorname{carr}(r)) \wedge \\ \quad \quad \quad p \in \operatorname{interval}(1, n) \wedge \\ \quad \quad \quad s \in \operatorname{carr}(r) \wedge \\ \quad \quad \quad \forall i, k \in \operatorname{interval}(1, n). \\ \quad \quad \quad \quad \neg (i = p) \Rightarrow (\operatorname{entry}(b, i, k) = \operatorname{entry}(a, i, k)) \wedge \\ \quad \quad \quad \forall i, k \in \operatorname{interval}(1, n). \\ \quad \quad \quad \quad \neg (i = p) \Rightarrow (\operatorname{entry}(c, i, k) = \operatorname{entry}(a, i, k)) \wedge \\ \quad \quad \quad \forall k \in \operatorname{interval}(1, n). \\ \quad \quad \quad \quad \operatorname{entry}(c, p, k) = \operatorname{add}(r)(\operatorname{mul}(r)(s, \operatorname{entry}(a, p, k)), \operatorname{entry}(b, p, k)) \Rightarrow \\ \quad \quad \quad \quad d(c) = \operatorname{add}(r)(\operatorname{mul}(r)(s, d(a)), d(b)) \wedge \\ \quad \quad \forall a, p, q. \\ \quad \quad \quad a \in \operatorname{mat}(n, n, \operatorname{carr}(r)) \wedge \\ \quad \quad \quad p \in \operatorname{interval}(1, n) \wedge \\ \quad \quad \quad q \in \operatorname{interval}(1, n) \wedge \\ \quad \quad \quad \neg (p = q) \wedge \\ \quad \quad \quad \forall k \in \operatorname{interval}(1, n).\; (\operatorname{entry}(a, p, k) = \operatorname{entry}(a, q, k)) \Rightarrow \\ \quad \quad \quad \quad d(a) = \operatorname{zero}(r) \Rightarrow \\ \quad \quad \quad d(x) = \operatorname{mul}(r)(\operatorname{det}(r, n, x), d(\operatorname{identmat}(r, n))) \wedge \\ \quad \operatorname{is\text{-}commutative\text{-}ring}(r) \wedge \\ \quad d \in (\operatorname{mat}((n + 1), (n + 1), \operatorname{carr}(r)) \to \operatorname{carr}(r)) \wedge \\ \quad \forall a, b, c, p, s. \\ \quad \quad a \in \operatorname{mat}((n + 1), (n + 1), \operatorname{carr}(r)) \wedge \\ \quad \quad b \in \operatorname{mat}((n + 1), (n + 1), \operatorname{carr}(r)) \wedge \\ \quad \quad c \in \operatorname{mat}((n + 1), (n + 1), \operatorname{carr}(r)) \wedge \\ \quad \quad p \in \operatorname{interval}(1, (n + 1)) \wedge \\ \quad \quad s \in \operatorname{carr}(r) \wedge \\ \quad \quad \forall i, k \in \operatorname{interval}(1, (n + 1)). \\ \quad \quad \quad \neg (i = p) \Rightarrow (\operatorname{entry}(b, i, k) = \operatorname{entry}(a, i, k)) \wedge \\ \quad \quad \forall i, k \in \operatorname{interval}(1, (n + 1)). \\ \quad \quad \quad \neg (i = p) \Rightarrow (\operatorname{entry}(c, i, k) = \operatorname{entry}(a, i, k)) \wedge \\ \quad \quad \forall k \in \operatorname{interval}(1, (n + 1)). \\ \quad \quad \quad \operatorname{entry}(c, p, k) = \operatorname{add}(r)(\operatorname{mul}(r)(s, \operatorname{entry}(a, p, k)), \operatorname{entry}(b, p, k)) \Rightarrow \\ \quad \quad \quad d(c) = \operatorname{add}(r)(\operatorname{mul}(r)(s, d(a)), d(b)) \wedge \\ \quad \forall a, p, q. \\ \quad \quad a \in \operatorname{mat}((n + 1), (n + 1), \operatorname{carr}(r)) \wedge \\ \quad \quad p \in \operatorname{interval}(1, (n + 1)) \wedge \\ \quad \quad q \in \operatorname{interval}(1, (n + 1)) \wedge \\ \quad \quad \neg (p = q) \wedge \\ \quad \quad \forall k \in \operatorname{interval}(1, (n + 1)).\; (\operatorname{entry}(a, p, k) = \operatorname{entry}(a, q, k)) \Rightarrow \\ \quad \quad \quad d(a) = \operatorname{zero}(r) \wedge \\ \quad y \in \operatorname{mat}((n + 1), (n + 1), \operatorname{carr}(r)) \wedge \\ \quad \operatorname{entry}(y, 1, 1) = \operatorname{one}(r) \wedge \\ \quad \forall k \in \operatorname{interval}(1, (n + 1)).\; (\neg (k = 1) \Rightarrow (\operatorname{entry}(y, 1, k) = \operatorname{zero}(r))) \Rightarrow \\ \quad \quad d(y) \\ \quad \quad =\; \operatorname{mul}(r)(\operatorname{det}(r, n, \\ \quad \quad \quad \quad \operatorname{submat}(\operatorname{matof}((n + 1), (n + 1), \\ \quad \quad \quad \quad \quad \quad (\lambda \langle \operatorname{dav\_}, \operatorname{daw\_} \rangle \in (\operatorname{interval}(1, (n + 1)) \times \operatorname{interval}(1, (n + 1))).\; \\ \quad \quad \quad \quad \quad \quad \quad \left(\text{if }(\operatorname{daw\_} = 1)\text{ then }\left(\text{if }(\operatorname{dav\_} = 1)\text{ then }\operatorname{entry}(y, \operatorname{dav\_}, \operatorname{daw\_})\text{ else }\left(\text{if }(\operatorname{dav\_} \leq (n + 1))\text{ then }\operatorname{zero}(r)\text{ else }\operatorname{entry}(y, \operatorname{dav\_}, \operatorname{daw\_})\right)\right)\text{ else }\operatorname{entry}(y, \operatorname{dav\_}, \operatorname{daw\_})\right))), \\ \quad \quad \quad \quad \quad n, n)), \\ \quad \quad \quad d(\operatorname{identmat}(r, (n + 1)))) \end{array}}}
\vnbprop{daf-basis}{\vnbstmt{\begin{array}{@{}l@{}} \forall n \in \mathbb{N}, r, d, x, j. \\ \quad \forall r, d.; \forall \mathit{dfx} \in \operatorname{mat}(n, n, \operatorname{carr}(r)). \\ \quad \quad \operatorname{is\text{-}commutative\text{-}ring}(r) \wedge \\ \quad \quad d \in (\operatorname{mat}(n, n, \operatorname{carr}(r)) \to \operatorname{carr}(r)) \wedge \\ \quad \quad \forall a, b, c, p, s. \\ \quad \quad \quad a \in \operatorname{mat}(n, n, \operatorname{carr}(r)) \wedge \\ \quad \quad \quad b \in \operatorname{mat}(n, n, \operatorname{carr}(r)) \wedge \\ \quad \quad \quad c \in \operatorname{mat}(n, n, \operatorname{carr}(r)) \wedge \\ \quad \quad \quad p \in \operatorname{interval}(1, n) \wedge \\ \quad \quad \quad s \in \operatorname{carr}(r) \wedge \\ \quad \quad \quad \forall i, k \in \operatorname{interval}(1, n). \\ \quad \quad \quad \quad \neg (i = p) \Rightarrow (\operatorname{entry}(b, i, k) = \operatorname{entry}(a, i, k)) \wedge \\ \quad \quad \quad \forall i, k \in \operatorname{interval}(1, n). \\ \quad \quad \quad \quad \neg (i = p) \Rightarrow (\operatorname{entry}(c, i, k) = \operatorname{entry}(a, i, k)) \wedge \\ \quad \quad \quad \forall k \in \operatorname{interval}(1, n). \\ \quad \quad \quad \quad \operatorname{entry}(c, p, k) = \operatorname{add}(r)(\operatorname{mul}(r)(s, \operatorname{entry}(a, p, k)), \operatorname{entry}(b, p, k)) \Rightarrow \\ \quad \quad \quad \quad d(c) = \operatorname{add}(r)(\operatorname{mul}(r)(s, d(a)), d(b)) \wedge \\ \quad \quad \forall a, p, q. \\ \quad \quad \quad a \in \operatorname{mat}(n, n, \operatorname{carr}(r)) \wedge \\ \quad \quad \quad p \in \operatorname{interval}(1, n) \wedge \\ \quad \quad \quad q \in \operatorname{interval}(1, n) \wedge \\ \quad \quad \quad \neg (p = q) \wedge \\ \quad \quad \quad \forall k \in \operatorname{interval}(1, n).\; (\operatorname{entry}(a, p, k) = \operatorname{entry}(a, q, k)) \Rightarrow \\ \quad \quad \quad \quad d(a) = \operatorname{zero}(r) \Rightarrow \\ \quad \quad \quad d(\mathit{dfx}) = \operatorname{mul}(r)(\operatorname{det}(r, n, \mathit{dfx}), d(\operatorname{identmat}(r, n))) \wedge \\ \quad \operatorname{is\text{-}commutative\text{-}ring}(r) \wedge \\ \quad d \in (\operatorname{mat}((n + 1), (n + 1), \operatorname{carr}(r)) \to \operatorname{carr}(r)) \wedge \\ \quad \forall a, b, c, p, s. \\ \quad \quad a \in \operatorname{mat}((n + 1), (n + 1), \operatorname{carr}(r)) \wedge \\ \quad \quad b \in \operatorname{mat}((n + 1), (n + 1), \operatorname{carr}(r)) \wedge \\ \quad \quad c \in \operatorname{mat}((n + 1), (n + 1), \operatorname{carr}(r)) \wedge \\ \quad \quad p \in \operatorname{interval}(1, (n + 1)) \wedge \\ \quad \quad s \in \operatorname{carr}(r) \wedge \\ \quad \quad \forall i, k \in \operatorname{interval}(1, (n + 1)). \\ \quad \quad \quad \neg (i = p) \Rightarrow (\operatorname{entry}(b, i, k) = \operatorname{entry}(a, i, k)) \wedge \\ \quad \quad \forall i, k \in \operatorname{interval}(1, (n + 1)). \\ \quad \quad \quad \neg (i = p) \Rightarrow (\operatorname{entry}(c, i, k) = \operatorname{entry}(a, i, k)) \wedge \\ \quad \quad \forall k \in \operatorname{interval}(1, (n + 1)). \\ \quad \quad \quad \operatorname{entry}(c, p, k) = \operatorname{add}(r)(\operatorname{mul}(r)(s, \operatorname{entry}(a, p, k)), \operatorname{entry}(b, p, k)) \Rightarrow \\ \quad \quad \quad d(c) = \operatorname{add}(r)(\operatorname{mul}(r)(s, d(a)), d(b)) \wedge \\ \quad \forall a, p, q. \\ \quad \quad a \in \operatorname{mat}((n + 1), (n + 1), \operatorname{carr}(r)) \wedge \\ \quad \quad p \in \operatorname{interval}(1, (n + 1)) \wedge \\ \quad \quad q \in \operatorname{interval}(1, (n + 1)) \wedge \\ \quad \quad \neg (p = q) \wedge \\ \quad \quad \forall k \in \operatorname{interval}(1, (n + 1)).\; (\operatorname{entry}(a, p, k) = \operatorname{entry}(a, q, k)) \Rightarrow \\ \quad \quad \quad d(a) = \operatorname{zero}(r) \wedge \\ \quad x \in \operatorname{mat}((n + 1), (n + 1), \operatorname{carr}(r)) \wedge \\ \quad j \in \operatorname{interval}(1, (n + 1)) \wedge \\ \quad \forall k \in \operatorname{interval}(1, (n + 1)).\; (\operatorname{entry}(x, 1, k) = \left(\text{if }(k = j)\text{ then }\operatorname{one}(r)\text{ else }\operatorname{zero}(r)\right)) \Rightarrow \\ \quad \quad d(x) = \operatorname{mul}(r)(\operatorname{det}(r, (n + 1), x), d(\operatorname{identmat}(r, (n + 1)))) \end{array}}}
\vnbprop{daf-psum-ind}{\vnbstmt{\begin{array}{@{}l@{}} \forall m \in \mathbb{N}, r, n, d, a. \\ \quad \operatorname{is\text{-}commutative\text{-}ring}(r) \wedge \\ \quad n \in \mathbb{N} \wedge \\ \quad d \in (\operatorname{mat}((n + 1), (n + 1), \operatorname{carr}(r)) \to \operatorname{carr}(r)) \wedge \\ \quad \forall \mathit{dfa}, b, c, p, s. \\ \quad \quad \mathit{dfa} \in \operatorname{mat}((n + 1), (n + 1), \operatorname{carr}(r)) \wedge \\ \quad \quad b \in \operatorname{mat}((n + 1), (n + 1), \operatorname{carr}(r)) \wedge \\ \quad \quad c \in \operatorname{mat}((n + 1), (n + 1), \operatorname{carr}(r)) \wedge \\ \quad \quad p \in \operatorname{interval}(1, (n + 1)) \wedge \\ \quad \quad s \in \operatorname{carr}(r) \wedge \\ \quad \quad \forall i, k \in \operatorname{interval}(1, (n + 1)). \\ \quad \quad \quad \neg (i = p) \Rightarrow (\operatorname{entry}(b, i, k) = \operatorname{entry}(\mathit{dfa}, i, k)) \wedge \\ \quad \quad \forall i, k \in \operatorname{interval}(1, (n + 1)). \\ \quad \quad \quad \neg (i = p) \Rightarrow (\operatorname{entry}(c, i, k) = \operatorname{entry}(\mathit{dfa}, i, k)) \wedge \\ \quad \quad \forall k \in \operatorname{interval}(1, (n + 1)). \\ \quad \quad \quad \operatorname{entry}(c, p, k) = \operatorname{add}(r)(\operatorname{mul}(r)(s, \operatorname{entry}(\mathit{dfa}, p, k)), \operatorname{entry}(b, p, k)) \Rightarrow \\ \quad \quad \quad d(c) = \operatorname{add}(r)(\operatorname{mul}(r)(s, d(\mathit{dfa})), d(b)) \wedge \\ \quad \forall \mathit{dfa}, p, q. \\ \quad \quad \mathit{dfa} \in \operatorname{mat}((n + 1), (n + 1), \operatorname{carr}(r)) \wedge \\ \quad \quad p \in \operatorname{interval}(1, (n + 1)) \wedge \\ \quad \quad q \in \operatorname{interval}(1, (n + 1)) \wedge \\ \quad \quad \neg (p = q) \wedge \\ \quad \quad \forall k \in \operatorname{interval}(1, (n + 1)).\; (\operatorname{entry}(\mathit{dfa}, p, k) = \operatorname{entry}(\mathit{dfa}, q, k)) \Rightarrow \\ \quad \quad \quad d(\mathit{dfa}) = \operatorname{zero}(r) \wedge \\ \quad a \in \operatorname{mat}((n + 1), (n + 1), \operatorname{carr}(r)) \wedge \\ \quad \forall j \in \operatorname{interval}(1, (n + 1)). \\ \quad \quad d(\operatorname{matof}((n + 1), (n + 1), \\ \quad \quad \quad \quad (\lambda \langle \operatorname{dav\_}, \operatorname{daw\_} \rangle \in (\operatorname{interval}(1, (n + 1)) \times \operatorname{interval}(1, (n + 1))).\; \\ \quad \quad \quad \quad \quad \left(\text{if }(\operatorname{dav\_} = 1)\text{ then }\left(\text{if }(\operatorname{daw\_} = j)\text{ then }\operatorname{one}(r)\text{ else }\operatorname{zero}(r)\right)\text{ else }\operatorname{entry}(a, \operatorname{dav\_}, \operatorname{daw\_})\right)))) \\ \quad \quad =\; \operatorname{mul}(r)(\operatorname{det}(r, (n + 1), \\ \quad \quad \quad \quad \operatorname{matof}((n + 1), (n + 1), \\ \quad \quad \quad \quad \quad (\lambda \langle \operatorname{dav\_}, \operatorname{daw\_} \rangle \in (\operatorname{interval}(1, (n + 1)) \times \operatorname{interval}(1, (n + 1))).\; \\ \quad \quad \quad \quad \quad \quad \left(\text{if }(\operatorname{dav\_} = 1)\text{ then }\left(\text{if }(\operatorname{daw\_} = j)\text{ then }\operatorname{one}(r)\text{ else }\operatorname{zero}(r)\right)\text{ else }\operatorname{entry}(a, \operatorname{dav\_}, \operatorname{daw\_})\right)))), \\ \quad \quad \quad d(\operatorname{identmat}(r, (n + 1)))) \wedge \\ \quad m \leq (n + 1) \Rightarrow \\ \quad \quad d(\operatorname{matof}((n + 1), (n + 1), \\ \quad \quad \quad \quad (\lambda \langle \operatorname{dav\_}, \operatorname{daw\_} \rangle \in (\operatorname{interval}(1, (n + 1)) \times \operatorname{interval}(1, (n + 1))).\; \\ \quad \quad \quad \quad \quad \left(\text{if }(\operatorname{dav\_} = 1)\text{ then }\left(\text{if }(\operatorname{daw\_} \leq m)\text{ then }\operatorname{entry}(a, 1, \operatorname{daw\_})\text{ else }\operatorname{zero}(r)\right)\text{ else }\operatorname{entry}(a, \operatorname{dav\_}, \operatorname{daw\_})\right)))) \\ \quad \quad =\; \operatorname{mul}(r)(\operatorname{det}(r, (n + 1), \\ \quad \quad \quad \quad \operatorname{matof}((n + 1), (n + 1), \\ \quad \quad \quad \quad \quad (\lambda \langle \operatorname{dav\_}, \operatorname{daw\_} \rangle \in (\operatorname{interval}(1, (n + 1)) \times \operatorname{interval}(1, (n + 1))).\; \\ \quad \quad \quad \quad \quad \quad \left(\text{if }(\operatorname{dav\_} = 1)\text{ then }\left(\text{if }(\operatorname{daw\_} \leq m)\text{ then }\operatorname{entry}(a, 1, \operatorname{daw\_})\text{ else }\operatorname{zero}(r)\right)\text{ else }\operatorname{entry}(a, \operatorname{dav\_}, \operatorname{daw\_})\right)))), \\ \quad \quad \quad d(\operatorname{identmat}(r, (n + 1)))) \end{array}}}
\vnbprop{daf-psum}{\vnbstmt{\begin{array}{@{}l@{}} \forall r, n, d, a. \\ \quad \operatorname{is\text{-}commutative\text{-}ring}(r) \wedge \\ \quad n \in \mathbb{N} \wedge \\ \quad d \in (\operatorname{mat}((n + 1), (n + 1), \operatorname{carr}(r)) \to \operatorname{carr}(r)) \wedge \\ \quad \forall \mathit{dfa}, b, c, p, s. \\ \quad \quad \mathit{dfa} \in \operatorname{mat}((n + 1), (n + 1), \operatorname{carr}(r)) \wedge \\ \quad \quad b \in \operatorname{mat}((n + 1), (n + 1), \operatorname{carr}(r)) \wedge \\ \quad \quad c \in \operatorname{mat}((n + 1), (n + 1), \operatorname{carr}(r)) \wedge \\ \quad \quad p \in \operatorname{interval}(1, (n + 1)) \wedge \\ \quad \quad s \in \operatorname{carr}(r) \wedge \\ \quad \quad \forall i, k \in \operatorname{interval}(1, (n + 1)). \\ \quad \quad \quad \neg (i = p) \Rightarrow (\operatorname{entry}(b, i, k) = \operatorname{entry}(\mathit{dfa}, i, k)) \wedge \\ \quad \quad \forall i, k \in \operatorname{interval}(1, (n + 1)). \\ \quad \quad \quad \neg (i = p) \Rightarrow (\operatorname{entry}(c, i, k) = \operatorname{entry}(\mathit{dfa}, i, k)) \wedge \\ \quad \quad \forall k \in \operatorname{interval}(1, (n + 1)). \\ \quad \quad \quad \operatorname{entry}(c, p, k) = \operatorname{add}(r)(\operatorname{mul}(r)(s, \operatorname{entry}(\mathit{dfa}, p, k)), \operatorname{entry}(b, p, k)) \Rightarrow \\ \quad \quad \quad d(c) = \operatorname{add}(r)(\operatorname{mul}(r)(s, d(\mathit{dfa})), d(b)) \wedge \\ \quad \forall \mathit{dfa}, p, q. \\ \quad \quad \mathit{dfa} \in \operatorname{mat}((n + 1), (n + 1), \operatorname{carr}(r)) \wedge \\ \quad \quad p \in \operatorname{interval}(1, (n + 1)) \wedge \\ \quad \quad q \in \operatorname{interval}(1, (n + 1)) \wedge \\ \quad \quad \neg (p = q) \wedge \\ \quad \quad \forall k \in \operatorname{interval}(1, (n + 1)).\; (\operatorname{entry}(\mathit{dfa}, p, k) = \operatorname{entry}(\mathit{dfa}, q, k)) \Rightarrow \\ \quad \quad \quad d(\mathit{dfa}) = \operatorname{zero}(r) \wedge \\ \quad a \in \operatorname{mat}((n + 1), (n + 1), \operatorname{carr}(r)) \wedge \\ \quad \forall j \in \operatorname{interval}(1, (n + 1)). \\ \quad \quad d(\operatorname{matof}((n + 1), (n + 1), \\ \quad \quad \quad \quad (\lambda \langle \operatorname{dav\_}, \operatorname{daw\_} \rangle \in (\operatorname{interval}(1, (n + 1)) \times \operatorname{interval}(1, (n + 1))).\; \\ \quad \quad \quad \quad \quad \left(\text{if }(\operatorname{dav\_} = 1)\text{ then }\left(\text{if }(\operatorname{daw\_} = j)\text{ then }\operatorname{one}(r)\text{ else }\operatorname{zero}(r)\right)\text{ else }\operatorname{entry}(a, \operatorname{dav\_}, \operatorname{daw\_})\right)))) \\ \quad \quad =\; \operatorname{mul}(r)(\operatorname{det}(r, (n + 1), \\ \quad \quad \quad \quad \operatorname{matof}((n + 1), (n + 1), \\ \quad \quad \quad \quad \quad (\lambda \langle \operatorname{dav\_}, \operatorname{daw\_} \rangle \in (\operatorname{interval}(1, (n + 1)) \times \operatorname{interval}(1, (n + 1))).\; \\ \quad \quad \quad \quad \quad \quad \left(\text{if }(\operatorname{dav\_} = 1)\text{ then }\left(\text{if }(\operatorname{daw\_} = j)\text{ then }\operatorname{one}(r)\text{ else }\operatorname{zero}(r)\right)\text{ else }\operatorname{entry}(a, \operatorname{dav\_}, \operatorname{daw\_})\right)))), \\ \quad \quad \quad d(\operatorname{identmat}(r, (n + 1)))) \Rightarrow \\ \quad \quad d(a) = \operatorname{mul}(r)(\operatorname{det}(r, (n + 1), a), d(\operatorname{identmat}(r, (n + 1)))) \end{array}}}
\vnbprop{det-alternating-form-ind}{\vnbstmt{\begin{array}{@{}l@{}} \forall n \in \mathbb{N}, r, d.; \forall x \in \operatorname{mat}(n, n, \operatorname{carr}(r)). \\ \quad \operatorname{is\text{-}commutative\text{-}ring}(r) \wedge \\ \quad d \in (\operatorname{mat}(n, n, \operatorname{carr}(r)) \to \operatorname{carr}(r)) \wedge \\ \quad \forall a, b, c, p, s. \\ \quad \quad a \in \operatorname{mat}(n, n, \operatorname{carr}(r)) \wedge \\ \quad \quad b \in \operatorname{mat}(n, n, \operatorname{carr}(r)) \wedge \\ \quad \quad c \in \operatorname{mat}(n, n, \operatorname{carr}(r)) \wedge \\ \quad \quad p \in \operatorname{interval}(1, n) \wedge \\ \quad \quad s \in \operatorname{carr}(r) \wedge \\ \quad \quad \forall i, k \in \operatorname{interval}(1, n). \\ \quad \quad \quad \neg (i = p) \Rightarrow (\operatorname{entry}(b, i, k) = \operatorname{entry}(a, i, k)) \wedge \\ \quad \quad \forall i, k \in \operatorname{interval}(1, n). \\ \quad \quad \quad \neg (i = p) \Rightarrow (\operatorname{entry}(c, i, k) = \operatorname{entry}(a, i, k)) \wedge \\ \quad \quad \forall k \in \operatorname{interval}(1, n). \\ \quad \quad \quad \operatorname{entry}(c, p, k) = \operatorname{add}(r)(\operatorname{mul}(r)(s, \operatorname{entry}(a, p, k)), \operatorname{entry}(b, p, k)) \Rightarrow \\ \quad \quad \quad d(c) = \operatorname{add}(r)(\operatorname{mul}(r)(s, d(a)), d(b)) \wedge \\ \quad \forall a, p, q. \\ \quad \quad a \in \operatorname{mat}(n, n, \operatorname{carr}(r)) \wedge \\ \quad \quad p \in \operatorname{interval}(1, n) \wedge \\ \quad \quad q \in \operatorname{interval}(1, n) \wedge \\ \quad \quad \neg (p = q) \wedge \\ \quad \quad \forall k \in \operatorname{interval}(1, n).\; (\operatorname{entry}(a, p, k) = \operatorname{entry}(a, q, k)) \Rightarrow \\ \quad \quad \quad d(a) = \operatorname{zero}(r) \Rightarrow \\ \quad \quad d(x) = \operatorname{mul}(r)(\operatorname{det}(r, n, x), d(\operatorname{identmat}(r, n))) \end{array}}}
\vnbprop{det-alternating-form}{\vnbstmt{\begin{array}{@{}l@{}} \forall r, n, d, a. \\ \quad \operatorname{is\text{-}commutative\text{-}ring}(r) \wedge \\ \quad n \in \mathbb{N} \wedge \\ \quad d \in (\operatorname{mat}(n, n, \operatorname{carr}(r)) \to \operatorname{carr}(r)) \wedge \\ \quad \forall \mathit{dfa}, b, c, p, s. \\ \quad \quad \mathit{dfa} \in \operatorname{mat}(n, n, \operatorname{carr}(r)) \wedge \\ \quad \quad b \in \operatorname{mat}(n, n, \operatorname{carr}(r)) \wedge \\ \quad \quad c \in \operatorname{mat}(n, n, \operatorname{carr}(r)) \wedge \\ \quad \quad p \in \operatorname{interval}(1, n) \wedge \\ \quad \quad s \in \operatorname{carr}(r) \wedge \\ \quad \quad \forall i, k \in \operatorname{interval}(1, n). \\ \quad \quad \quad \neg (i = p) \Rightarrow (\operatorname{entry}(b, i, k) = \operatorname{entry}(\mathit{dfa}, i, k)) \wedge \\ \quad \quad \forall i, k \in \operatorname{interval}(1, n). \\ \quad \quad \quad \neg (i = p) \Rightarrow (\operatorname{entry}(c, i, k) = \operatorname{entry}(\mathit{dfa}, i, k)) \wedge \\ \quad \quad \forall k \in \operatorname{interval}(1, n). \\ \quad \quad \quad \operatorname{entry}(c, p, k) = \operatorname{add}(r)(\operatorname{mul}(r)(s, \operatorname{entry}(\mathit{dfa}, p, k)), \operatorname{entry}(b, p, k)) \Rightarrow \\ \quad \quad \quad d(c) = \operatorname{add}(r)(\operatorname{mul}(r)(s, d(\mathit{dfa})), d(b)) \wedge \\ \quad \forall \mathit{dfa}, p, q. \\ \quad \quad \mathit{dfa} \in \operatorname{mat}(n, n, \operatorname{carr}(r)) \wedge \\ \quad \quad p \in \operatorname{interval}(1, n) \wedge \\ \quad \quad q \in \operatorname{interval}(1, n) \wedge \\ \quad \quad \neg (p = q) \wedge \\ \quad \quad \forall k \in \operatorname{interval}(1, n).\; (\operatorname{entry}(\mathit{dfa}, p, k) = \operatorname{entry}(\mathit{dfa}, q, k)) \Rightarrow \\ \quad \quad \quad d(\mathit{dfa}) = \operatorname{zero}(r) \wedge \\ \quad a \in \operatorname{mat}(n, n, \operatorname{carr}(r)) \Rightarrow \\ \quad \quad d(a) = \operatorname{mul}(r)(\operatorname{det}(r, n, a), d(\operatorname{identmat}(r, n))) \end{array}}}

\vnbbookfile{theorem-library/det-multiplicative}{theorem-library/det-multiplicative}
\noindent{\small det-multiplicative.scm -- det(P Q) = det(P) det(Q) over a commutative ring (Hoffman-Kunze 5.3 Theorem 3; the notes' proposition "det is multiplicative"), and the notes' corollary 3.20 (R a field, V invertible: det(A V) = det A det V). Batch 26-B, 2026-09-24.}\par
\vnbprop{det-multiplicative}{\vnbstmt{\begin{array}{@{}l@{}} \forall r, n, p, q. \\ \quad (\operatorname{is\text{-}commutative\text{-}ring}(r) \wedge ((n \in \mathbb{N}) \wedge ((p \in \operatorname{mat}(n, n, \operatorname{carr}(r))) \wedge (q \in \operatorname{mat}(n, n, \operatorname{carr}(r)))))) \Rightarrow  \\ \quad \quad \operatorname{det}(r, n, \operatorname{matmul}(r, p, q)) = \operatorname{mul}(r)(\operatorname{det}(r, n, p), \operatorname{det}(r, n, q)) \end{array}}}
\vnbprop{det-mul-invertible}{\vnbstmt{\begin{array}{@{}l@{}} \forall r, n, a, v. \\ \quad (\operatorname{is\text{-}field\text{-}ring}(r) \wedge ((n \in \mathbb{N}) \wedge ((a \in \operatorname{mat}(n, n, \operatorname{carr}(r))) \wedge \operatorname{is\text{-}invertible\text{-}mat}(r, n, v)))) \Rightarrow  \\ \quad \quad \operatorname{det}(r, n, \operatorname{matmul}(r, a, v)) = \operatorname{mul}(r)(\operatorname{det}(r, n, a), \operatorname{det}(r, n, v)) \end{array}}}

\vnbbookfile{theorem-library/det-transpose}{theorem-library/det-transpose}
\noindent{\small det-transpose.scm -- the transpose of a matrix and det(A\textasciicircum{}t) = det A; the COLUMN clauses of the linear-algebra notes' 3.18 (linearity in each column, two equal columns give 0, interchanging two columns negates det), 3.19 (det under an elementary COLUMN operation) and 3.22 (expansion along any column).  Batch 32}\par
\vnbprop{transpose-in-mat}{\vnbstmt{\begin{array}{@{}l@{}} \forall a, m, n, x. \\ \quad (a \in \operatorname{mat}(m, n, x)) \Rightarrow (\operatorname{transpose}(a, m, n) \in \operatorname{mat}(n, m, x)) \end{array}}}
\vnbprop{transpose-entry}{\vnbstmt{\begin{array}{@{}l@{}} \forall a, m, n, x, i, j. \\ \quad ((a \in \operatorname{mat}(m, n, x)) \wedge ((i \in \operatorname{interval}(1, n)) \wedge (j \in \operatorname{interval}(1, m)))) \Rightarrow  \\ \quad \quad \operatorname{entry}(\operatorname{transpose}(a, m, n), i, j) = \operatorname{entry}(a, j, i) \end{array}}}
\vnbprop{transpose-transpose}{\vnbstmt{\begin{array}{@{}l@{}} \forall a, m, n, x. \\ \quad (a \in \operatorname{mat}(m, n, x)) \Rightarrow (\operatorname{transpose}(\operatorname{transpose}(a, m, n), n, m) = a) \end{array}}}
\vnbprop{transpose-identity}{\vnbstmt{\begin{array}{@{}l@{}} \forall r, n. \\ \quad (\operatorname{is\text{-}ring}(r) \wedge (n \in \mathbb{N})) \Rightarrow (\operatorname{transpose}(\operatorname{identmat}(r, n), n, n) = \operatorname{identmat}(r, n)) \end{array}}}
\vnbprop{transpose-matmul}{\vnbstmt{\begin{array}{@{}l@{}} \forall r, m, n, c, p, q. \\ \quad \operatorname{is\text{-}commutative\text{-}ring}(r) \wedge \\ \quad p \in \operatorname{mat}(m, n, \operatorname{carr}(r)) \wedge \\ \quad q \in \operatorname{mat}(n, c, \operatorname{carr}(r)) \wedge \\ \quad (n = 0) \Rightarrow ((m = 0) \vee (c = 0)) \Rightarrow \\ \quad \quad \operatorname{transpose}(\operatorname{matmul}(r, p, q), m, c) = \operatorname{matmul}(r, \operatorname{transpose}(q, n, c), \operatorname{transpose}(p, m, n)) \end{array}}}
\vnbprop{dtt-ring-mul-neg1}{\vnbstmt{\begin{array}{@{}l@{}} \forall r.; \forall x \in \operatorname{carr}(r). \\ \quad \operatorname{is\text{-}commutative\text{-}ring}(r) \Rightarrow (\operatorname{mul}(r)(x, \operatorname{neg}(r)(\operatorname{one}(r))) = \operatorname{neg}(r)(x)) \end{array}}}
\vnbprop{dtt-ring-mul-one}{\vnbstmt{\begin{array}{@{}l@{}} \forall r.; \forall x \in \operatorname{carr}(r). \\ \quad \operatorname{is\text{-}commutative\text{-}ring}(r) \Rightarrow (\operatorname{mul}(r)(x, \operatorname{one}(r)) = x) \end{array}}}
\vnbprop{dtt-ring-neg-lin}{\vnbstmt{\begin{array}{@{}l@{}} \forall r, s, a, b. \\ \quad (\operatorname{is\text{-}commutative\text{-}ring}(r) \wedge ((s \in \operatorname{carr}(r)) \wedge ((a \in \operatorname{carr}(r)) \wedge (b \in \operatorname{carr}(r))))) \Rightarrow  \\ \quad \quad \operatorname{neg}(r)(\operatorname{add}(r)(\operatorname{mul}(r)(s, \operatorname{neg}(r)(a)), \operatorname{neg}(r)(b))) = \operatorname{add}(r)(\operatorname{mul}(r)(s, a), b) \end{array}}}
\vnbprop{dtt-ring-neg-neg}{\vnbstmt{\begin{array}{@{}l@{}} \forall r.; \forall x \in \operatorname{carr}(r). \\ \quad \operatorname{is\text{-}commutative\text{-}ring}(r) \Rightarrow (\operatorname{neg}(r)(\operatorname{neg}(r)(x)) = x) \end{array}}}
\vnbprop{dtt-ring-negx-x}{\vnbstmt{\begin{array}{@{}l@{}} \forall r.; \forall x \in \operatorname{carr}(r). \\ \quad \operatorname{is\text{-}commutative\text{-}ring}(r) \Rightarrow  \\ \quad \quad \operatorname{add}(r)(\operatorname{mul}(r)(\operatorname{neg}(r)(\operatorname{one}(r)), x), x) = \operatorname{zero}(r) \end{array}}}
\vnbprop{dtt-ring-zero-lin}{\vnbstmt{\begin{array}{@{}l@{}} \forall r. \\ \quad \operatorname{is\text{-}commutative\text{-}ring}(r) \Rightarrow  \\ \quad \quad \operatorname{zero}(r) = \operatorname{add}(r)(\operatorname{mul}(r)(\operatorname{neg}(r)(\operatorname{one}(r)), \operatorname{zero}(r)), \operatorname{zero}(r)) \end{array}}}
\vnbprop{dtt-ring-x-xneg}{\vnbstmt{\begin{array}{@{}l@{}} \forall r.; \forall x \in \operatorname{carr}(r). \\ \quad \operatorname{is\text{-}commutative\text{-}ring}(r) \Rightarrow  \\ \quad \quad \operatorname{add}(r)(x, \operatorname{mul}(r)(x, \operatorname{neg}(r)(\operatorname{one}(r)))) = \operatorname{zero}(r) \end{array}}}
\vnbprop{det-mul-elem-f-col}{\vnbstmt{\begin{array}{@{}l@{}} \forall r, n, a, k, l. \\ \quad \operatorname{is\text{-}commutative\text{-}ring}(r) \wedge \\ \quad n \in \mathbb{N} \wedge \\ \quad a \in \operatorname{mat}(n, n, \operatorname{carr}(r)) \wedge \\ \quad k \in \operatorname{interval}(1, n) \wedge \\ \quad l \in \operatorname{interval}(1, n) \wedge \\ \quad \neg (k = l) \Rightarrow \\ \quad \quad \operatorname{det}(r, n, \operatorname{matmul}(r, a, \operatorname{elem\text{-}f}(r, n, k, l))) = \operatorname{mul}(r)(\operatorname{det}(r, n, a), \operatorname{det}(r, n, \operatorname{elem\text{-}f}(r, n, k, l))) \end{array}}}
\vnbprop{det-elem-f-col}{\vnbstmt{\begin{array}{@{}l@{}} \forall r, n, a, k, l. \\ \quad \operatorname{is\text{-}commutative\text{-}ring}(r) \wedge \\ \quad n \in \mathbb{N} \wedge \\ \quad a \in \operatorname{mat}(n, n, \operatorname{carr}(r)) \wedge \\ \quad k \in \operatorname{interval}(1, n) \wedge \\ \quad l \in \operatorname{interval}(1, n) \wedge \\ \quad \neg (k = l) \Rightarrow \\ \quad \quad \operatorname{det}(r, n, \operatorname{matmul}(r, a, \operatorname{elem\text{-}f}(r, n, k, l))) = \operatorname{neg}(r)(\operatorname{det}(r, n, a)) \end{array}}}
\vnbprop{det-mul-elem-g-col}{\vnbstmt{\begin{array}{@{}l@{}} \forall r, n, a, c, k, l. \\ \quad \operatorname{is\text{-}commutative\text{-}ring}(r) \wedge \\ \quad n \in \mathbb{N} \wedge \\ \quad a \in \operatorname{mat}(n, n, \operatorname{carr}(r)) \wedge \\ \quad c \in \operatorname{carr}(r) \wedge \\ \quad k \in \operatorname{interval}(1, n) \wedge \\ \quad l \in \operatorname{interval}(1, n) \wedge \\ \quad \neg (k = l) \Rightarrow \\ \quad \quad \operatorname{det}(r, n, \operatorname{matmul}(r, a, \operatorname{elem\text{-}g}(r, n, c, k, l))) \\ \quad \quad =\; \operatorname{mul}(r)(\operatorname{det}(r, n, a), \operatorname{det}(r, n, \operatorname{elem\text{-}g}(r, n, c, k, l))) \end{array}}}
\vnbprop{det-elem-g-col}{\vnbstmt{\begin{array}{@{}l@{}} \forall r, n, a, c, k, l. \\ \quad \operatorname{is\text{-}commutative\text{-}ring}(r) \wedge \\ \quad n \in \mathbb{N} \wedge \\ \quad a \in \operatorname{mat}(n, n, \operatorname{carr}(r)) \wedge \\ \quad c \in \operatorname{carr}(r) \wedge \\ \quad k \in \operatorname{interval}(1, n) \wedge \\ \quad l \in \operatorname{interval}(1, n) \wedge \\ \quad \neg (k = l) \Rightarrow \\ \quad \quad \operatorname{det}(r, n, \operatorname{matmul}(r, a, \operatorname{elem\text{-}g}(r, n, c, k, l))) = \operatorname{det}(r, n, a) \end{array}}}
\vnbprop{det-mul-elem-h-col}{\vnbstmt{\begin{array}{@{}l@{}} \forall r, n, a, c, k. \\ \quad \operatorname{is\text{-}commutative\text{-}ring}(r) \wedge \\ \quad n \in \mathbb{N} \wedge \\ \quad a \in \operatorname{mat}(n, n, \operatorname{carr}(r)) \wedge \\ \quad c \in \operatorname{carr}(r) \wedge \\ \quad k \in \operatorname{interval}(1, n) \Rightarrow \\ \quad \quad \operatorname{det}(r, n, \operatorname{matmul}(r, a, \operatorname{elem\text{-}h}(r, n, c, k))) = \operatorname{mul}(r)(\operatorname{det}(r, n, a), \operatorname{det}(r, n, \operatorname{elem\text{-}h}(r, n, c, k))) \end{array}}}
\vnbprop{det-elem-h-col}{\vnbstmt{\begin{array}{@{}l@{}} \forall r, n, a, c, k. \\ \quad \operatorname{is\text{-}commutative\text{-}ring}(r) \wedge \\ \quad n \in \mathbb{N} \wedge \\ \quad a \in \operatorname{mat}(n, n, \operatorname{carr}(r)) \wedge \\ \quad c \in \operatorname{carr}(r) \wedge \\ \quad k \in \operatorname{interval}(1, n) \Rightarrow \\ \quad \quad \operatorname{det}(r, n, \operatorname{matmul}(r, a, \operatorname{elem\text{-}h}(r, n, c, k))) = \operatorname{mul}(r)(c, \operatorname{det}(r, n, a)) \end{array}}}
\vnbprop{dtt-minor1-colcongr}{\vnbstmt{\begin{array}{@{}l@{}} \forall r, n, x, y, p. \\ \quad \operatorname{is\text{-}ring}(r) \wedge \\ \quad n \in \mathbb{N} \wedge \\ \quad x \in \operatorname{mat}((n + 1), (n + 1), \operatorname{carr}(r)) \wedge \\ \quad y \in \operatorname{mat}((n + 1), (n + 1), \operatorname{carr}(r)) \wedge \\ \quad p = 1 \wedge \\ \quad \forall j, i \in \operatorname{interval}(1, (n + 1)). \\ \quad \quad \neg (j = p) \Rightarrow (\operatorname{entry}(x, i, j) = \operatorname{entry}(y, i, j)) \Rightarrow \\ \quad \quad \operatorname{minor}(x, 1, p, n) = \operatorname{minor}(y, 1, p, n) \end{array}}}
\vnbprop{dtt-col1-ind}{\vnbstmt{\begin{array}{@{}l@{}} \forall n \in \mathbb{N}, r, a, b, c, p, \mathit{dsc}. \\ \quad \operatorname{is\text{-}commutative\text{-}ring}(r) \wedge \\ \quad a \in \operatorname{mat}(n, n, \operatorname{carr}(r)) \wedge \\ \quad b \in \operatorname{mat}(n, n, \operatorname{carr}(r)) \wedge \\ \quad c \in \operatorname{mat}(n, n, \operatorname{carr}(r)) \wedge \\ \quad p \in \operatorname{interval}(1, n) \wedge \\ \quad p = 1 \wedge \\ \quad \mathit{dsc} \in \operatorname{carr}(r) \wedge \\ \quad \forall j, i \in \operatorname{interval}(1, n). \\ \quad \quad \neg (j = p) \Rightarrow (\operatorname{entry}(b, i, j) = \operatorname{entry}(a, i, j)) \wedge \\ \quad \forall j, i \in \operatorname{interval}(1, n). \\ \quad \quad \neg (j = p) \Rightarrow (\operatorname{entry}(c, i, j) = \operatorname{entry}(a, i, j)) \wedge \\ \quad \forall i \in \operatorname{interval}(1, n). \\ \quad \quad \operatorname{entry}(c, i, p) = \operatorname{add}(r)(\operatorname{mul}(r)(\mathit{dsc}, \operatorname{entry}(a, i, p)), \operatorname{entry}(b, i, p)) \Rightarrow \\ \quad \quad \operatorname{det}(r, n, c) = \operatorname{add}(r)(\operatorname{mul}(r)(\mathit{dsc}, \operatorname{det}(r, n, a)), \operatorname{det}(r, n, b)) \end{array}}}
\vnbprop{det-col-linear}{\vnbstmt{\begin{array}{@{}l@{}} \forall r, n, a, b, c, p, d. \\ \quad \operatorname{is\text{-}commutative\text{-}ring}(r) \wedge \\ \quad n \in \mathbb{N} \wedge \\ \quad a \in \operatorname{mat}(n, n, \operatorname{carr}(r)) \wedge \\ \quad b \in \operatorname{mat}(n, n, \operatorname{carr}(r)) \wedge \\ \quad c \in \operatorname{mat}(n, n, \operatorname{carr}(r)) \wedge \\ \quad p \in \operatorname{interval}(1, n) \wedge \\ \quad d \in \operatorname{carr}(r) \wedge \\ \quad \forall j, i \in \operatorname{interval}(1, n). \\ \quad \quad \neg (j = p) \Rightarrow (\operatorname{entry}(b, i, j) = \operatorname{entry}(a, i, j)) \wedge \\ \quad \forall j, i \in \operatorname{interval}(1, n). \\ \quad \quad \neg (j = p) \Rightarrow (\operatorname{entry}(c, i, j) = \operatorname{entry}(a, i, j)) \wedge \\ \quad \forall i \in \operatorname{interval}(1, n). \\ \quad \quad \operatorname{entry}(c, i, p) = \operatorname{add}(r)(\operatorname{mul}(r)(d, \operatorname{entry}(a, i, p)), \operatorname{entry}(b, i, p)) \Rightarrow \\ \quad \quad \operatorname{det}(r, n, c) = \operatorname{add}(r)(\operatorname{mul}(r)(d, \operatorname{det}(r, n, a)), \operatorname{det}(r, n, b)) \end{array}}}
\vnbprop{dtt-zero-col}{\vnbstmt{\begin{array}{@{}l@{}} \forall r, n, c, q. \\ \quad \operatorname{is\text{-}commutative\text{-}ring}(r) \wedge \\ \quad n \in \mathbb{N} \wedge \\ \quad c \in \operatorname{mat}(n, n, \operatorname{carr}(r)) \wedge \\ \quad q \in \operatorname{interval}(1, n) \wedge \\ \quad \forall i \in \operatorname{interval}(1, n).\; (\operatorname{entry}(c, i, q) = \operatorname{zero}(r)) \Rightarrow \\ \quad \quad \operatorname{det}(r, n, c) = \operatorname{zero}(r) \end{array}}}
\vnbprop{det-equal-cols-zero}{\vnbstmt{\begin{array}{@{}l@{}} \forall r, n, a, p, q. \\ \quad \operatorname{is\text{-}commutative\text{-}ring}(r) \wedge \\ \quad n \in \mathbb{N} \wedge \\ \quad a \in \operatorname{mat}(n, n, \operatorname{carr}(r)) \wedge \\ \quad p \in \operatorname{interval}(1, n) \wedge \\ \quad q \in \operatorname{interval}(1, n) \wedge \\ \quad \neg (p = q) \wedge \\ \quad \forall i \in \operatorname{interval}(1, n).\; (\operatorname{entry}(a, i, p) = \operatorname{entry}(a, i, q)) \Rightarrow \\ \quad \quad \operatorname{det}(r, n, a) = \operatorname{zero}(r) \end{array}}}
\vnbprop{det-swap-cols}{\vnbstmt{\begin{array}{@{}l@{}} \forall r, n, a, b, p, q. \\ \quad \operatorname{is\text{-}commutative\text{-}ring}(r) \wedge \\ \quad n \in \mathbb{N} \wedge \\ \quad a \in \operatorname{mat}(n, n, \operatorname{carr}(r)) \wedge \\ \quad b \in \operatorname{mat}(n, n, \operatorname{carr}(r)) \wedge \\ \quad p \in \operatorname{interval}(1, n) \wedge \\ \quad q \in \operatorname{interval}(1, n) \wedge \\ \quad \neg (p = q) \wedge \\ \quad \forall j, i \in \operatorname{interval}(1, n). \\ \quad \quad (\neg (j = p) \wedge \neg (j = q)) \Rightarrow (\operatorname{entry}(b, i, j) = \operatorname{entry}(a, i, j)) \wedge \\ \quad \forall i \in \operatorname{interval}(1, n).\; (\operatorname{entry}(b, i, p) = \operatorname{entry}(a, i, q)) \wedge \\ \quad \forall i \in \operatorname{interval}(1, n).\; (\operatorname{entry}(b, i, q) = \operatorname{entry}(a, i, p)) \Rightarrow \\ \quad \quad \operatorname{det}(r, n, b) = \operatorname{neg}(r)(\operatorname{det}(r, n, a)) \end{array}}}
\vnbprop{dtt-tr-fun}{\vnbstmt{\begin{array}{@{}l@{}} \forall r, n. \\ \quad (\operatorname{is\text{-}commutative\text{-}ring}(r) \wedge (n \in \mathbb{N})) \Rightarrow  \\ \quad \quad (\lambda x \in \operatorname{mat}(n, n, \operatorname{carr}(r)).\; \operatorname{det}(r, n, \operatorname{transpose}(x, n, n))) \in (\operatorname{mat}(n, n, \operatorname{carr}(r)) \to \operatorname{carr}(r)) \end{array}}}
\vnbprop{dtt-tr-lin}{\vnbstmt{\begin{array}{@{}l@{}} \forall r, n, a, b, c, p, s. \\ \quad \operatorname{is\text{-}commutative\text{-}ring}(r) \wedge \\ \quad n \in \mathbb{N} \wedge \\ \quad a \in \operatorname{mat}(n, n, \operatorname{carr}(r)) \wedge \\ \quad b \in \operatorname{mat}(n, n, \operatorname{carr}(r)) \wedge \\ \quad c \in \operatorname{mat}(n, n, \operatorname{carr}(r)) \wedge \\ \quad p \in \operatorname{interval}(1, n) \wedge \\ \quad s \in \operatorname{carr}(r) \wedge \\ \quad \forall i, k \in \operatorname{interval}(1, n). \\ \quad \quad \neg (i = p) \Rightarrow (\operatorname{entry}(b, i, k) = \operatorname{entry}(a, i, k)) \wedge \\ \quad \forall i, k \in \operatorname{interval}(1, n). \\ \quad \quad \neg (i = p) \Rightarrow (\operatorname{entry}(c, i, k) = \operatorname{entry}(a, i, k)) \wedge \\ \quad \forall k \in \operatorname{interval}(1, n). \\ \quad \quad \operatorname{entry}(c, p, k) = \operatorname{add}(r)(\operatorname{mul}(r)(s, \operatorname{entry}(a, p, k)), \operatorname{entry}(b, p, k)) \Rightarrow \\ \quad \quad (\lambda x \in \operatorname{mat}(n, n, \operatorname{carr}(r)).\; \operatorname{det}(r, n, \operatorname{transpose}(x, n, n)))(c) \\ \quad \quad =\; \operatorname{add}(r)(\operatorname{mul}(r)(s, (\lambda x \in \operatorname{mat}(n, n, \operatorname{carr}(r)).\; \operatorname{det}(r, n, \operatorname{transpose}(x, n, n)))(a)), \\ \quad \quad \quad (\lambda x \in \operatorname{mat}(n, n, \operatorname{carr}(r)).\; \operatorname{det}(r, n, \operatorname{transpose}(x, n, n)))(b)) \end{array}}}
\vnbprop{dtt-tr-alt}{\vnbstmt{\begin{array}{@{}l@{}} \forall r, n, a, p, q. \\ \quad \operatorname{is\text{-}commutative\text{-}ring}(r) \wedge \\ \quad n \in \mathbb{N} \wedge \\ \quad a \in \operatorname{mat}(n, n, \operatorname{carr}(r)) \wedge \\ \quad p \in \operatorname{interval}(1, n) \wedge \\ \quad q \in \operatorname{interval}(1, n) \wedge \\ \quad \neg (p = q) \wedge \\ \quad \forall k \in \operatorname{interval}(1, n).\; (\operatorname{entry}(a, p, k) = \operatorname{entry}(a, q, k)) \Rightarrow \\ \quad \quad (\lambda x \in \operatorname{mat}(n, n, \operatorname{carr}(r)).\; \operatorname{det}(r, n, \operatorname{transpose}(x, n, n)))(a) = \operatorname{zero}(r) \end{array}}}
\vnbprop{det-transpose}{\vnbstmt{\begin{array}{@{}l@{}} \forall r, n, a. \\ \quad (\operatorname{is\text{-}commutative\text{-}ring}(r) \wedge ((n \in \mathbb{N}) \wedge (a \in \operatorname{mat}(n, n, \operatorname{carr}(r))))) \Rightarrow  \\ \quad \quad \operatorname{det}(r, n, \operatorname{transpose}(a, n, n)) = \operatorname{det}(r, n, a) \end{array}}}
\vnbprop{dtt-minor-transpose}{\vnbstmt{\begin{array}{@{}l@{}} \forall r, n, s, p, q. \\ \quad (\operatorname{is\text{-}ring}(r) \wedge ((n \in \mathbb{N}) \wedge ((s \in \operatorname{mat}((n + 1), (n + 1), \operatorname{carr}(r))) \wedge ((p \in \mathbb{N}) \wedge (q \in \mathbb{N}))))) \Rightarrow  \\ \quad \quad \operatorname{minor}(\operatorname{transpose}(s, (n + 1), (n + 1)), q, p, n) = \operatorname{transpose}(\operatorname{minor}(s, p, q, n), n, n) \end{array}}}
\vnbprop{det-expand-col}{\vnbstmt{\begin{array}{@{}l@{}} \forall r, n, a, j. \\ \quad \operatorname{is\text{-}commutative\text{-}ring}(r) \wedge \\ \quad n \in \mathbb{N} \wedge \\ \quad a \in \operatorname{mat}((n + 1), (n + 1), \operatorname{carr}(r)) \wedge \\ \quad j \in \operatorname{interval}(1, (n + 1)) \Rightarrow \\ \quad \quad \operatorname{det}(r, (n + 1), a) \\ \quad \quad =\; \operatorname{finsum}(\operatorname{ring\text{-}additive\text{-}ag}(r), \\ \quad \quad \quad (\lambda e \in \operatorname{interval}(1, (n + 1)).\; \\ \quad \quad \quad \quad \operatorname{mul}(r)(\operatorname{mpow}(\operatorname{ring\text{-}multiplicative\text{-}monoid}(r), \operatorname{neg}(r)(\operatorname{one}(r)), (e + j)), \\ \quad \quad \quad \quad \quad \operatorname{mul}(r)(\operatorname{entry}(a, e, j), \operatorname{det}(r, n, \operatorname{minor}(a, e, j, n))))), \\ \quad \quad \quad \operatorname{interval}(1, (n + 1))) \end{array}}}

\vnbbookfile{theorem-library/clear-pivot-cross-proof}{theorem-library/clear-pivot-cross-proof}
\noindent{\small clear-pivot-cross -- over a euclidean ring, a matrix P (in MAT(succ p, succ q)) with a nonzero entry is \textasciitilde{}-equivalent to a matrix C2 in BORDER block form [[b,0],[0,C']] with b = C2\_\{1,1\} nonzero.  This is the per-level Smith step: place a class-minimal pivot at (1,1), clear its row and column, and read the}\par
\vnbprop{clear-pivot-cross}{\vnbstmt{\begin{array}{@{}l@{}} \forall a, p, q, m. \\ \quad \operatorname{is\text{-}euclidean\text{-}ring}(a) \wedge \\ \quad p \in \mathbb{N} \wedge \\ \quad q \in \mathbb{N} \wedge \\ \quad m \in \operatorname{mat}((p + 1), (q + 1), \operatorname{carr}(a)) \wedge \\ \quad \exists b \in \operatorname{interval}(1, (p + 1)), c \in \operatorname{interval}(1, (q + 1)). \\ \quad \quad \neg (\operatorname{entry}(m, b, c) = \operatorname{zero}(a)) \Rightarrow \\ \quad \quad \exists d. \\ \quad \quad \quad \operatorname{mat\text{-}equiv}(a, (p + 1), (q + 1), m, d) \wedge \\ \quad \quad \quad d = \operatorname{border}(a, \operatorname{entry}(d, 1, 1), \operatorname{submat}(d, p, q), p, q) \wedge \\ \quad \quad \quad \neg (\operatorname{entry}(d, 1, 1) = \operatorname{zero}(a)) \end{array}}}

\vnbbookfile{theorem-library/smith-diagonalization-proof}{theorem-library/smith-diagonalization-proof}
\noindent{\small smith-diagonalization -- every matrix P in MAT(m,n) over a euclidean ring is \textasciitilde{}-equivalent to a DIAGONAL matrix (Smith normal form, DIAGONALIZATION ONLY -- no divisibility chain).  The capstone of the LA arc.  Induction (ni) on the row dim k, inner statement INNER(k) = forall n P over MAT(k,n) exists diagonal}\par
\vnbprop{smith-diagonalization}{\vnbstmt{\begin{array}{@{}l@{}} \forall a.; \forall k \in \mathbb{N}, n, p. \\ \quad (\operatorname{is\text{-}euclidean\text{-}ring}(a) \wedge ((n \in \mathbb{N}) \wedge (p \in \operatorname{mat}(k, n, \operatorname{carr}(a))))) \Rightarrow  \\ \quad \quad \exists d.\; (\operatorname{mat\text{-}equiv}(a, k, n, p, d) \wedge \operatorname{is\text{-}diagonal}(a, k, n, d)) \end{array}}}

\vnbbookfile{theorem-library/smith-staircase-proof}{theorem-library/smith-staircase-proof}
\noindent{\small smith-staircase-proof.scm -- the Smith normal form with its nonzero diagonal entries recorded as an INITIAL SEGMENT.}\par
\vnbprop{border-staircase}{\vnbstmt{\begin{array}{@{}l@{}} \forall a, b, w, p, q, c. \\ \quad \operatorname{is\text{-}ring}(a) \wedge \\ \quad p \in \mathbb{N} \wedge \\ \quad q \in \mathbb{N} \wedge \\ \quad c \in \mathbb{N} \wedge \\ \quad b \in \operatorname{carr}(a) \wedge \\ \quad \neg (b = \operatorname{zero}(a)) \wedge \\ \quad \operatorname{smith\text{-}staircase}(a, p, q, w, c) \wedge \\ \quad w \in \operatorname{mat}(p, q, \operatorname{carr}(a)) \Rightarrow \\ \quad \quad \operatorname{smith\text{-}staircase}(a, (p + 1), (q + 1), \operatorname{border}(a, b, w, p, q), (c + 1)) \end{array}}}
\vnbprop{smith-normal-form}{\vnbstmt{\begin{array}{@{}l@{}} \forall a.; \forall k \in \mathbb{N}, n, p. \\ \quad (\operatorname{is\text{-}euclidean\text{-}ring}(a) \wedge ((n \in \mathbb{N}) \wedge (p \in \operatorname{mat}(k, n, \operatorname{carr}(a))))) \Rightarrow  \\ \quad \quad \exists d.; \exists b \in \mathbb{N}. \\ \quad \quad \quad \operatorname{mat\text{-}equiv}(a, k, n, p, d) \wedge \operatorname{smith\text{-}staircase}(a, k, n, d, b) \end{array}}}

\vnbbookfile{theorem-library/matact-row-linear-proof}{theorem-library/matact-row-linear-proof}
\noindent{\small matact-row-linear-proof.scm -- the coefficient row acts LINEARLY.}\par
\vnbprop{lincomb-unfold}{\vnbstmt{\begin{array}{@{}l@{}} \forall d, n, c, u. \\ \quad (\operatorname{lincomb}(d, n, c, u) \simeq \operatorname{finsum}(\operatorname{module\text{-}vector\text{-}ag}(d), (\lambda j \in \operatorname{interval}(1, n).\; \operatorname{act}(d)(\operatorname{entry}(c, 1, j), \operatorname{entry}(u, j, 1))), \operatorname{interval}(1, n))) \end{array}}}
\vnbprop{lincomb-type}{\vnbstmt{\begin{array}{@{}l@{}} \forall d, n, c, u. \\ \quad (\operatorname{is\text{-}module}(d) \wedge ((c \in \operatorname{mat}(1, n, \operatorname{carr}(\operatorname{scal}(d)))) \wedge (u \in \operatorname{mat}(n, 1, \operatorname{vec}(d))))) \Rightarrow  \\ \quad \quad \operatorname{lincomb}(d, n, c, u) \in \operatorname{vec}(d) \end{array}}}
\vnbprop{matact-lincomb}{\vnbstmt{\begin{array}{@{}l@{}} \forall d, n, c, u. \\ \quad (\operatorname{is\text{-}module}(d) \wedge ((c \in \operatorname{mat}(1, n, \operatorname{carr}(\operatorname{scal}(d)))) \wedge ((u \in \operatorname{mat}(n, 1, \operatorname{vec}(d))) \wedge (1 \leq n)))) \Rightarrow  \\ \quad \quad \operatorname{entry}(\operatorname{matact}(d, c, u), 1, 1) = \operatorname{lincomb}(d, n, c, u) \end{array}}}
\vnbprop{lincomb-row-add}{\vnbstmt{\begin{array}{@{}l@{}} \forall d, n, a, b, u. \\ \quad \operatorname{is\text{-}module}(d) \wedge \\ \quad a \in \operatorname{mat}(1, n, \operatorname{carr}(\operatorname{scal}(d))) \wedge \\ \quad b \in \operatorname{mat}(1, n, \operatorname{carr}(\operatorname{scal}(d))) \wedge \\ \quad u \in \operatorname{mat}(n, 1, \operatorname{vec}(d)) \Rightarrow \\ \quad \quad \operatorname{lincomb}(d, n, \operatorname{matadd}(\operatorname{scal}(d), a, b), u) = \operatorname{vadd}(d)(\operatorname{lincomb}(d, n, a, u), \operatorname{lincomb}(d, n, b, u)) \end{array}}}
\vnbprop{lincomb-row-scale}{\vnbstmt{\begin{array}{@{}l@{}} \forall d, n, r, c, u. \\ \quad (\operatorname{is\text{-}module}(d) \wedge ((r \in \operatorname{carr}(\operatorname{scal}(d))) \wedge ((c \in \operatorname{mat}(1, n, \operatorname{carr}(\operatorname{scal}(d)))) \wedge (u \in \operatorname{mat}(n, 1, \operatorname{vec}(d)))))) \Rightarrow  \\ \quad \quad \operatorname{lincomb}(d, n, \operatorname{matscale}(\operatorname{scal}(d), r, c), u) = \operatorname{act}(d)(r, \operatorname{lincomb}(d, n, c, u)) \end{array}}}

\vnbbookfile{theorem-library/span-bricks-proof}{theorem-library/span-bricks-proof}
\noindent{\small span-bricks-proof.scm -- the module-action bricks under spans-submodule-fg.}\par
\vnbprop{module-act-neg-one}{\vnbstmt{\begin{array}{@{}l@{}} \forall d.; \forall x \in \operatorname{vec}(d). \\ \quad \operatorname{is\text{-}module}(d) \Rightarrow (\operatorname{act}(d)(\operatorname{neg}(\operatorname{scal}(d))(\operatorname{one}(\operatorname{scal}(d))), x) = \operatorname{vneg}(d)(x)) \end{array}}}
\vnbprop{lincomb-zerorow}{\vnbstmt{\begin{array}{@{}l@{}} \forall d, n.; \forall u \in \operatorname{mat}(n, 1, \operatorname{vec}(d)). \\ \quad \operatorname{is\text{-}module}(d) \Rightarrow (\operatorname{lincomb}(d, n, \operatorname{zeromat}(\operatorname{scal}(d), 1, n), u) = \operatorname{vzero}(d)) \end{array}}}
\vnbprop{lincomb-unitrow}{\vnbstmt{\begin{array}{@{}l@{}} \forall d, n, i.; \forall u \in \operatorname{mat}(n, 1, \operatorname{vec}(d)). \\ \quad (\operatorname{is\text{-}module}(d) \wedge (i \in \operatorname{interval}(1, n))) \Rightarrow  \\ \quad \quad \operatorname{lincomb}(d, n, \operatorname{unitrow}(\operatorname{scal}(d), n, i), u) = \operatorname{entry}(u, i, 1) \end{array}}}
\vnbprop{lincomb-empty}{\vnbstmt{\forall d, c, u.\; (\operatorname{is\text{-}module}(d) \Rightarrow (\operatorname{lincomb}(d, 0, c, u) = \operatorname{vzero}(d)))}}
\vnbprop{span-is-submodule}{\vnbstmt{\begin{array}{@{}l@{}} \forall d, n.; \forall u \in \operatorname{mat}(n, 1, \operatorname{vec}(d)). \\ \quad \operatorname{is\text{-}module}(d) \Rightarrow \operatorname{is\text{-}submodule}(d, \operatorname{span}(d, n, u)) \end{array}}}
\vnbprop{spans-span}{\vnbstmt{\begin{array}{@{}l@{}} \forall d, n.; \forall u \in \operatorname{mat}(n, 1, \operatorname{vec}(d)). \\ \quad \operatorname{is\text{-}module}(d) \Rightarrow \operatorname{spans}(d, n, u, \operatorname{span}(d, n, u)) \end{array}}}

\vnbbookfile{theorem-library/span-bricks2-proof}{theorem-library/span-bricks2-proof}
\noindent{\small span-bricks2-proof.scm -- bricks 4, 5, 6 under spans-submodule-fg.}\par
\vnbprop{lincomb-row-peel}{\vnbstmt{\begin{array}{@{}l@{}} \forall d.; \forall n \in \mathbb{N}, c \in \operatorname{mat}(1, (n + 1), \operatorname{carr}(\operatorname{scal}(d))), u \in \operatorname{mat}((n + 1), 1, \operatorname{vec}(d)). \\ \quad \operatorname{is\text{-}module}(d) \Rightarrow  \\ \quad \quad \operatorname{lincomb}(d, (n + 1), c, u) \\ \quad \quad =\; \operatorname{vadd}(d)(\operatorname{lincomb}(d, n, \operatorname{block}(c, 1, n), \operatorname{block}(u, n, 1)), \\ \quad \quad \quad \operatorname{act}(d)(\operatorname{entry}(c, 1, (n + 1)), \operatorname{entry}(u, (n + 1), 1))) \end{array}}}
\vnbprop{lincomb-snoc}{\vnbstmt{\begin{array}{@{}l@{}} \forall d.; \forall n \in \mathbb{N}, c \in \operatorname{mat}(1, n, \operatorname{carr}(\operatorname{scal}(d))), u \in \operatorname{mat}(n, 1, \operatorname{vec}(d)), r \in \operatorname{carr}(\operatorname{scal}(d)), x \in \operatorname{vec}(d). \\ \quad \operatorname{is\text{-}module}(d) \Rightarrow  \\ \quad \quad \operatorname{lincomb}(d, (n + 1), \operatorname{snoc\text{-}row}(c, n, r), \operatorname{snoc\text{-}col}(u, n, x)) \\ \quad \quad =\; \operatorname{vadd}(d)(\operatorname{lincomb}(d, n, c, u), \operatorname{act}(d)(r, x)) \end{array}}}
\vnbprop{submodule-intersection}{\vnbstmt{\begin{array}{@{}l@{}} \forall d, a, b. \\ \quad (\operatorname{is\text{-}module}(d) \wedge (\operatorname{is\text{-}submodule}(d, a) \wedge \operatorname{is\text{-}submodule}(d, b))) \Rightarrow  \\ \quad \quad \operatorname{is\text{-}submodule}(d, (a \cap b)) \end{array}}}
\vnbprop{mat-1-0-nonempty}{\vnbstmt{\forall x.\; \exists p.\; (p \in \operatorname{mat}(1, 0, x))}}
\vnbprop{mat-0-1-nonempty}{\vnbstmt{\forall x.\; \exists p.\; (p \in \operatorname{mat}(0, 1, x))}}

\vnbbookfile{theorem-library/lastcoeff-ideal-proof}{theorem-library/lastcoeff-ideal-proof}
\noindent{\small lastcoeff-ideal-proof.scm -- the two self-contained lemmas the spans-submodule-fg descent leans on (no induction here).}\par
\vnbprop{lastcoeff-zero-in-span}{\vnbstmt{\begin{array}{@{}l@{}} \forall d.; \forall p \in \mathbb{N}, u \in \operatorname{mat}((p + 1), 1, \operatorname{vec}(d)), c \in \operatorname{mat}(1, (p + 1), \operatorname{carr}(\operatorname{scal}(d))). \\ \quad (\operatorname{is\text{-}module}(d) \wedge (\operatorname{entry}(c, 1, (p + 1)) = \operatorname{zero}(\operatorname{scal}(d)))) \Rightarrow  \\ \quad \quad \operatorname{lincomb}(d, (p + 1), c, u) \in \operatorname{span}(d, p, \operatorname{block}(u, p, 1)) \end{array}}}
\vnbprop{lastcoeff-set-is-ideal}{\vnbstmt{\begin{array}{@{}l@{}} \forall d.; \forall p \in \mathbb{N}, u \in \operatorname{mat}((p + 1), 1, \operatorname{vec}(d)), m. \\ \quad (\operatorname{is\text{-}module}(d) \wedge (\operatorname{is\text{-}commutative\text{-}ring}(\operatorname{scal}(d)) \wedge \operatorname{is\text{-}submodule}(d, m))) \Rightarrow  \\ \quad \quad \operatorname{is\text{-}ideal}(\operatorname{scal}(d), \operatorname{lastcoeff\text{-}set}(d, p, u, m)) \end{array}}}
\vnbprop{descent-remainder}{\vnbstmt{\begin{array}{@{}l@{}} \forall d.; \forall n \in \mathbb{N}, u \in \operatorname{mat}((n + 1), 1, \operatorname{vec}(d)), c, a \in \operatorname{mat}(1, (n + 1), \operatorname{carr}(\operatorname{scal}(d))), q, b \in \operatorname{carr}(\operatorname{scal}(d)). \\ \quad \operatorname{is\text{-}module}(d) \wedge \\ \quad \operatorname{is\text{-}commutative\text{-}ring}(\operatorname{scal}(d)) \wedge \\ \quad \operatorname{entry}(a, 1, (n + 1)) = b \wedge \\ \quad \operatorname{entry}(c, 1, (n + 1)) = \operatorname{mul}(\operatorname{scal}(d))(q, b) \Rightarrow \\ \quad \quad \exists y \in \operatorname{span}(d, n, \operatorname{block}(u, n, 1)). \\ \quad \quad \quad \operatorname{lincomb}(d, (n + 1), c, u) = \operatorname{vadd}(d)(y, \operatorname{act}(d)(q, \operatorname{lincomb}(d, (n + 1), a, u))) \wedge \\ \quad \quad \quad y = \operatorname{vadd}(d)(\operatorname{lincomb}(d, (n + 1), c, u), \operatorname{act}(d)(\operatorname{neg}(\operatorname{scal}(d))(q), \operatorname{lincomb}(d, (n + 1), a, u))) \end{array}}}

\vnbbookfile{theorem-library/spans-submodule-fg-proof}{theorem-library/spans-submodule-fg-proof}
\noindent{\small spans-submodule-fg-proof.scm -- THE DESCENT.}\par
\vnbprop{spans-fg-base}{\vnbstmt{\begin{array}{@{}l@{}} \forall d.; \forall u \in \operatorname{mat}(0, 1, \operatorname{vec}(d)), m, a. \\ \quad \operatorname{is\text{-}module}(d) \wedge \\ \quad \operatorname{is\text{-}euclidean\text{-}ring}(\operatorname{scal}(d)) \wedge \\ \quad \operatorname{is\text{-}submodule}(d, m) \wedge \\ \quad \operatorname{spans}(d, 0, u, m) \wedge \\ \quad \operatorname{is\text{-}submodule}(d, a) \wedge \\ \quad a \subseteq m \Rightarrow \\ \quad \quad \exists k. \\ \quad \quad \quad ((k \in \mathbb{N}) \wedge (k \leq 0)) \wedge \exists w \in \operatorname{mat}(k, 1, \operatorname{vec}(d)).\; \operatorname{spans}(d, k, w, a) \end{array}}}
\vnbprop{spans-fg-step}{\vnbstmt{\begin{array}{@{}l@{}} \forall d.; \forall n \in \mathbb{N}, u \in \operatorname{mat}((n + 1), 1, \operatorname{vec}(d)), m, a. \\ \quad \operatorname{is\text{-}module}(d) \wedge \\ \quad \operatorname{is\text{-}euclidean\text{-}ring}(\operatorname{scal}(d)) \wedge \\ \quad \forall u \in \operatorname{mat}(n, 1, \operatorname{vec}(d)), m, a. \\ \quad \quad (\operatorname{is\text{-}submodule}(d, m) \wedge (\operatorname{spans}(d, n, u, m) \wedge (\operatorname{is\text{-}submodule}(d, a) \wedge (a \subseteq m)))) \Rightarrow  \\ \quad \quad \quad \exists k. \\ \quad \quad \quad \quad ((k \in \mathbb{N}) \wedge (k \leq n)) \wedge \exists w \in \operatorname{mat}(k, 1, \operatorname{vec}(d)).\; \operatorname{spans}(d, k, w, a) \wedge \\ \quad \operatorname{is\text{-}submodule}(d, m) \wedge \\ \quad \operatorname{spans}(d, (n + 1), u, m) \wedge \\ \quad \operatorname{is\text{-}submodule}(d, a) \wedge \\ \quad a \subseteq m \Rightarrow \\ \quad \quad \exists k. \\ \quad \quad \quad (((k \in \mathbb{N}) \wedge (k \leq (n + 1))) \wedge \exists w \in \operatorname{mat}(k, 1, \operatorname{vec}(d)).\; \operatorname{spans}(d, k, w, a)) \end{array}}}
\vnbprop{spans-submodule-fg}{\vnbstmt{\begin{array}{@{}l@{}} \forall d.; \forall n \in \mathbb{N}, u \in \operatorname{mat}(n, 1, \operatorname{vec}(d)), m, a. \\ \quad \operatorname{is\text{-}module}(d) \wedge \\ \quad \operatorname{is\text{-}euclidean\text{-}ring}(\operatorname{scal}(d)) \wedge \\ \quad \operatorname{is\text{-}submodule}(d, m) \wedge \\ \quad \operatorname{spans}(d, n, u, m) \wedge \\ \quad \operatorname{is\text{-}submodule}(d, a) \wedge \\ \quad a \subseteq m \Rightarrow \\ \quad \quad \exists k. \\ \quad \quad \quad ((k \in \mathbb{N}) \wedge (k \leq n)) \wedge \exists w \in \operatorname{mat}(k, 1, \operatorname{vec}(d)).\; \operatorname{spans}(d, k, w, a) \end{array}}}

\vnbbookfile{theorem-library/submodule-fg-proof}{theorem-library/submodule-fg-proof}
\noindent{\small submodule-fg-proof.scm -- submodule-fg, PROVEN as a corollary of the SPANS-relativized statement it should always have had.}\par
\vnbprop{whole-module-is-submodule}{\vnbstmt{\forall d.\; (\operatorname{is\text{-}module}(d) \Rightarrow \operatorname{is\text{-}submodule}(d, \operatorname{vec}(d)))}}
\vnbprop{generates-implies-spans-vec}{\vnbstmt{\begin{array}{@{}l@{}} \forall d.; \forall n \in \mathbb{N}, u \in \operatorname{mat}(n, 1, \operatorname{vec}(d)). \\ \quad (\operatorname{is\text{-}module}(d) \wedge \operatorname{generates}(d, n, u)) \Rightarrow \operatorname{spans}(d, n, u, \operatorname{vec}(d)) \end{array}}}
\vnbprop{submodule-fg}{\vnbstmt{\begin{array}{@{}l@{}} \forall d.; \forall n \in \mathbb{N}, u \in \operatorname{mat}(n, 1, \operatorname{vec}(d)), m. \\ \quad ((\operatorname{is\text{-}module}(d) \wedge \operatorname{is\text{-}euclidean\text{-}ring}(\operatorname{scal}(d))) \wedge (\operatorname{generates}(d, n, u) \wedge \operatorname{is\text{-}submodule}(d, m))) \Rightarrow  \\ \quad \quad \exists k. \\ \quad \quad \quad ((k \in \mathbb{N}) \wedge (k \leq n)) \wedge \exists w \in \operatorname{mat}(k, 1, \operatorname{vec}(d)).\; \operatorname{spans}(d, k, w, m) \end{array}}}

\vnbbookfile{theorem-library/matact-assoc-proof}{theorem-library/matact-assoc-proof}
\noindent{\small matact-assoc-proof.scm -- Remark 3.39: the matrix action on a sequence of module elements is associative.}\par
\vnbprop{matact-triple-left}{\vnbstmt{\begin{array}{@{}l@{}} \forall d, m, n, k, l, p, q, u, w, a. \\ \quad \operatorname{is\text{-}module}(d) \wedge \\ \quad p \in \operatorname{mat}(m, n, \operatorname{carr}(\operatorname{scal}(d))) \wedge \\ \quad q \in \operatorname{mat}(n, k, \operatorname{carr}(\operatorname{scal}(d))) \wedge \\ \quad u \in \operatorname{mat}(k, l, \operatorname{vec}(d)) \wedge \\ \quad 1 \leq n \wedge \\ \quad 1 \leq k \wedge \\ \quad w \in \operatorname{interval}(1, m) \wedge \\ \quad a \in \operatorname{interval}(1, l) \Rightarrow \\ \quad \quad \operatorname{entry}(\operatorname{matact}(d, \operatorname{matmul}(\operatorname{scal}(d), p, q), u), w, a) \\ \quad \quad =\; \operatorname{finsum}(\operatorname{module\text{-}vector\text{-}ag}(d), \\ \quad \quad \quad (\lambda c \in \operatorname{interval}(1, k).\; \\ \quad \quad \quad \quad \operatorname{finsum}(\operatorname{module\text{-}vector\text{-}ag}(d), \\ \quad \quad \quad \quad \quad (\lambda j \in \operatorname{interval}(1, n).\; \\ \quad \quad \quad \quad \quad \quad (\lambda z \in (\operatorname{interval}(1, k) \times \operatorname{interval}(1, n)).\; \operatorname{act}(d)(\operatorname{mul}(\operatorname{scal}(d))(\operatorname{entry}(p, w, \operatorname{nth}(2, z)), \operatorname{entry}(q, \operatorname{nth}(2, z), \operatorname{nth}(1, z))), \operatorname{entry}(u, \operatorname{nth}(1, z), a)))(\langle c, j \rangle)), \\ \quad \quad \quad \quad \quad \operatorname{interval}(1, n))), \\ \quad \quad \quad \operatorname{interval}(1, k)) \end{array}}}
\vnbprop{matact-triple-right}{\vnbstmt{\begin{array}{@{}l@{}} \forall d, m, n, k, l, p, q, u, w, a. \\ \quad \operatorname{is\text{-}module}(d) \wedge \\ \quad p \in \operatorname{mat}(m, n, \operatorname{carr}(\operatorname{scal}(d))) \wedge \\ \quad q \in \operatorname{mat}(n, k, \operatorname{carr}(\operatorname{scal}(d))) \wedge \\ \quad u \in \operatorname{mat}(k, l, \operatorname{vec}(d)) \wedge \\ \quad 1 \leq n \wedge \\ \quad 1 \leq k \wedge \\ \quad w \in \operatorname{interval}(1, m) \wedge \\ \quad a \in \operatorname{interval}(1, l) \Rightarrow \\ \quad \quad \operatorname{entry}(\operatorname{matact}(d, p, \operatorname{matact}(d, q, u)), w, a) \\ \quad \quad =\; \operatorname{finsum}(\operatorname{module\text{-}vector\text{-}ag}(d), \\ \quad \quad \quad (\lambda j \in \operatorname{interval}(1, n).\; \\ \quad \quad \quad \quad \operatorname{finsum}(\operatorname{module\text{-}vector\text{-}ag}(d), \\ \quad \quad \quad \quad \quad (\lambda c \in \operatorname{interval}(1, k).\; \\ \quad \quad \quad \quad \quad \quad (\lambda z \in (\operatorname{interval}(1, k) \times \operatorname{interval}(1, n)).\; \operatorname{act}(d)(\operatorname{mul}(\operatorname{scal}(d))(\operatorname{entry}(p, w, \operatorname{nth}(2, z)), \operatorname{entry}(q, \operatorname{nth}(2, z), \operatorname{nth}(1, z))), \operatorname{entry}(u, \operatorname{nth}(1, z), a)))(\langle c, j \rangle)), \\ \quad \quad \quad \quad \quad \operatorname{interval}(1, k))), \\ \quad \quad \quad \operatorname{interval}(1, n)) \end{array}}}
\vnbprop{matact-assoc}{\vnbstmt{\begin{array}{@{}l@{}} \forall d, m, n, k, l, p, q, u. \\ \quad \operatorname{is\text{-}module}(d) \wedge \\ \quad p \in \operatorname{mat}(m, n, \operatorname{carr}(\operatorname{scal}(d))) \wedge \\ \quad q \in \operatorname{mat}(n, k, \operatorname{carr}(\operatorname{scal}(d))) \wedge \\ \quad u \in \operatorname{mat}(k, l, \operatorname{vec}(d)) \wedge \\ \quad (n = 0) \Rightarrow ((m = 0) \vee ((k = 0) \wedge (l = 0))) \wedge \\ \quad (k = 0) \Rightarrow ((l = 0) \vee ((n = 0) \wedge (m = 0))) \Rightarrow \\ \quad \quad \operatorname{matact}(d, \operatorname{matmul}(\operatorname{scal}(d), p, q), u) = \operatorname{matact}(d, p, \operatorname{matact}(d, q, u)) \end{array}}}

\vnbbookfile{theorem-library/mod-basis-proof}{theorem-library/mod-basis-proof}
\noindent{\small mod-basis-proof.scm -- Lemma 3.40 (algebraic-numbers.pdf ch.3 sec 8.2): an invertible matrix carries generating sequences to generating sequences and relation-free sequences to relation-free sequences.  Plus the identity action I . u = u that says MATACT is a genuine action.}\par
\vnbprop{matact-identmat}{\vnbstmt{\begin{array}{@{}l@{}} \forall d, n, q.; \forall u \in \operatorname{mat}(n, q, \operatorname{vec}(d)). \\ \quad \operatorname{is\text{-}module}(d) \Rightarrow (\operatorname{matact}(d, \operatorname{identmat}(\operatorname{scal}(d), n), u) = u) \end{array}}}
\vnbprop{generates-transport}{\vnbstmt{\begin{array}{@{}l@{}} \forall d, n, u, m. \\ \quad (\operatorname{is\text{-}module}(d) \wedge ((u \in \operatorname{mat}(n, 1, \operatorname{vec}(d))) \wedge (\operatorname{is\text{-}invertible\text{-}mat}(\operatorname{scal}(d), n, m) \wedge \operatorname{generates}(d, n, u)))) \Rightarrow  \\ \quad \quad \operatorname{generates}(d, n, \operatorname{matact}(d, m, u)) \end{array}}}
\vnbprop{free-transport}{\vnbstmt{\begin{array}{@{}l@{}} \forall d, m, v, a. \\ \quad (\operatorname{is\text{-}module}(d) \wedge ((v \in \operatorname{mat}(m, 1, \operatorname{vec}(d))) \wedge (\operatorname{is\text{-}invertible\text{-}mat}(\operatorname{scal}(d), m, a) \wedge \operatorname{rel\text{-}free}(d, m, v)))) \Rightarrow  \\ \quad \quad \operatorname{rel\text{-}free}(d, m, \operatorname{matact}(d, a, v)) \end{array}}}

\vnbbookfile{theorem-library/inverse-invertible-proof}{theorem-library/inverse-invertible-proof}
\noindent{\small inverse-invertible-proof.scm -- the inverse of an invertible matrix is invertible.}\par
\vnbprop{inverse-is-invertible}{\vnbstmt{\begin{array}{@{}l@{}} \forall a, n, m, b. \\ \quad \operatorname{is\text{-}ring}(a) \wedge \\ \quad m \in \operatorname{mat}(n, n, \operatorname{carr}(a)) \wedge \\ \quad b \in \operatorname{mat}(n, n, \operatorname{carr}(a)) \wedge \\ \quad \operatorname{matmul}(a, m, b) = \operatorname{identmat}(a, n) \wedge \\ \quad \operatorname{matmul}(a, b, m) = \operatorname{identmat}(a, n) \Rightarrow \\ \quad \quad \operatorname{is\text{-}invertible\text{-}mat}(a, n, b) \end{array}}}

\vnbbookfile{theorem-library/spans-transport-proof}{theorem-library/spans-transport-proof}
\noindent{\small spans-transport-proof.scm -- the SPANS-relativized generates-transport.}\par
\vnbprop{spans-transport}{\vnbstmt{\begin{array}{@{}l@{}} \forall d, n, u, m, a. \\ \quad \operatorname{is\text{-}module}(d) \wedge \\ \quad u \in \operatorname{mat}(n, 1, \operatorname{vec}(d)) \wedge \\ \quad \operatorname{is\text{-}invertible\text{-}mat}(\operatorname{scal}(d), n, m) \wedge \\ \quad \operatorname{is\text{-}submodule}(d, a) \wedge \\ \quad \operatorname{spans}(d, n, u, a) \Rightarrow \\ \quad \quad \operatorname{spans}(d, n, \operatorname{matact}(d, m, u), a) \end{array}}}

\vnbbookfile{theorem-library/rake-asm-b4}{theorem-library/rake-asm-b4}
\noindent{\small theorem-library/rake-asm-b4.scm  (RB-4, rake-asm-b, 2026-10-12) ==================================================================== submodule-of-free-is-free (Cor 3.46, linear-algebra.tex cor:22195-45888-207248): over a euclidean ring, a submodule sm of a module}\par
\vnbprop{rb4-finsum-ag-pad-ptwise}{\vnbstmt{\begin{array}{@{}l@{}} \forall a.; \forall t \in \mathbf{Set}, r, g. \\ \quad \operatorname{is\text{-}abelian\text{-}group}(a) \wedge \\ \quad |t| \in \mathbb{N} \wedge \\ \quad \forall z \in r.\; (z \in t) \wedge \\ \quad \forall z \in t.\; (g(z) \in \operatorname{carr}(a)) \wedge \\ \quad \forall z \in t.\; (\neg (z \in r) \Rightarrow (g(z) = \operatorname{iden}(a))) \Rightarrow \\ \quad \quad \operatorname{finsum}(a, g, t) = \operatorname{finsum}(a, g, r) \end{array}}}
\vnbprop{rb4-lincomb-trunc}{\vnbstmt{\begin{array}{@{}l@{}} \forall d, m, k, c, y. \\ \quad \operatorname{is\text{-}module}(d) \wedge \\ \quad c \in \operatorname{mat}(1, m, \operatorname{carr}(\operatorname{scal}(d))) \wedge \\ \quad y \in \operatorname{mat}(m, 1, \operatorname{vec}(d)) \wedge \\ \quad k \in \mathbb{N} \wedge \\ \quad k \leq m \wedge \\ \quad \forall j \in \operatorname{interval}(1, m). \\ \quad \quad \neg (j \leq k) \Rightarrow (\operatorname{act}(d)(\operatorname{entry}(c, 1, j), \operatorname{entry}(y, j, 1)) = \operatorname{vzero}(d)) \Rightarrow \\ \quad \quad \operatorname{lincomb}(d, m, c, y) = \operatorname{lincomb}(d, k, c, y) \end{array}}}
\vnbprop{rb4-act-entry-single}{\vnbstmt{\begin{array}{@{}l@{}} \forall d, m, n, p, e, i. \\ \quad \operatorname{is\text{-}module}(d) \wedge \\ \quad p \in \operatorname{mat}(m, n, \operatorname{carr}(\operatorname{scal}(d))) \wedge \\ \quad e \in \operatorname{mat}(n, 1, \operatorname{vec}(d)) \wedge \\ \quad i \in \operatorname{interval}(1, m) \wedge \\ \quad i \in \operatorname{interval}(1, n) \wedge \\ \quad \forall j \in \operatorname{interval}(1, n).\; (\neg (j = i) \Rightarrow (\operatorname{entry}(p, i, j) = \operatorname{zero}(\operatorname{scal}(d)))) \Rightarrow \\ \quad \quad \operatorname{entry}(\operatorname{matact}(d, p, e), i, 1) = \operatorname{act}(d)(\operatorname{entry}(p, i, i), \operatorname{entry}(e, i, 1)) \end{array}}}
\vnbprop{rb4-act-entry-zero-row}{\vnbstmt{\begin{array}{@{}l@{}} \forall d, m, n, p, e, i. \\ \quad \operatorname{is\text{-}module}(d) \wedge \\ \quad p \in \operatorname{mat}(m, n, \operatorname{carr}(\operatorname{scal}(d))) \wedge \\ \quad e \in \operatorname{mat}(n, 1, \operatorname{vec}(d)) \wedge \\ \quad i \in \operatorname{interval}(1, m) \wedge \\ \quad 1 \leq n \wedge \\ \quad \forall j \in \operatorname{interval}(1, n).\; (\operatorname{entry}(p, i, j) = \operatorname{zero}(\operatorname{scal}(d))) \Rightarrow \\ \quad \quad \operatorname{entry}(\operatorname{matact}(d, p, e), i, 1) = \operatorname{vzero}(d) \end{array}}}
\vnbprop{rb4-matmul-entry-single}{\vnbstmt{\begin{array}{@{}l@{}} \forall a, m, n, k, p, q, w, l. \\ \quad \operatorname{is\text{-}ring}(a) \wedge \\ \quad p \in \operatorname{mat}(m, n, \operatorname{carr}(a)) \wedge \\ \quad q \in \operatorname{mat}(n, k, \operatorname{carr}(a)) \wedge \\ \quad w \in \operatorname{interval}(1, m) \wedge \\ \quad l \in \operatorname{interval}(1, n) \wedge \\ \quad l \in \operatorname{interval}(1, k) \wedge \\ \quad \forall i \in \operatorname{interval}(1, n).\; (\neg (i = l) \Rightarrow (\operatorname{entry}(q, i, l) = \operatorname{zero}(a))) \Rightarrow \\ \quad \quad \operatorname{entry}(\operatorname{matmul}(a, p, q), w, l) = \operatorname{mul}(a)(\operatorname{entry}(p, w, l), \operatorname{entry}(q, l, l)) \end{array}}}
\vnbprop{rb4-diag-rel-free}{\vnbstmt{\begin{array}{@{}l@{}} \forall a, m, n, k, d, e, w. \\ \quad \operatorname{is\text{-}module}(a) \wedge \\ \quad \operatorname{is\text{-}integral\text{-}domain}(\operatorname{scal}(a)) \wedge \\ \quad m \in \mathbb{N} \wedge \\ \quad n \in \mathbb{N} \wedge \\ \quad k \in \mathbb{N} \wedge \\ \quad k \leq m \wedge \\ \quad k \leq n \wedge \\ \quad d \in \operatorname{mat}(m, n, \operatorname{carr}(\operatorname{scal}(a))) \wedge \\ \quad \operatorname{is\text{-}diagonal}(\operatorname{scal}(a), m, n, d) \wedge \\ \quad \forall i \in \operatorname{interval}(1, k).\; \neg (\operatorname{entry}(d, i, i) = \operatorname{zero}(\operatorname{scal}(a))) \wedge \\ \quad e \in \operatorname{mat}(n, 1, \operatorname{vec}(a)) \wedge \\ \quad \operatorname{rel\text{-}free}(a, n, e) \wedge \\ \quad w \in \operatorname{mat}(k, 1, \operatorname{vec}(a)) \wedge \\ \quad \forall i \in \operatorname{interval}(1, k). \\ \quad \quad \operatorname{entry}(w, i, 1) = \operatorname{act}(a)(\operatorname{entry}(d, i, i), \operatorname{entry}(e, i, 1)) \Rightarrow \\ \quad \quad \operatorname{rel\text{-}free}(a, k, w) \end{array}}}
\vnbprop{rb4-span-trunc}{\vnbstmt{\begin{array}{@{}l@{}} \forall d, m, k, y, a. \\ \quad \operatorname{is\text{-}module}(d) \wedge \\ \quad m \in \mathbb{N} \wedge \\ \quad k \in \mathbb{N} \wedge \\ \quad k \leq m \wedge \\ \quad y \in \operatorname{mat}(m, 1, \operatorname{vec}(d)) \wedge \\ \quad \operatorname{spans}(d, m, y, a) \wedge \\ \quad \forall j \in \operatorname{interval}(1, m).\; (\neg (j \leq k) \Rightarrow (\operatorname{entry}(y, j, 1) = \operatorname{vzero}(d))) \Rightarrow \\ \quad \quad \operatorname{spans}(d, k, \operatorname{block}(y, k, 1), a) \end{array}}}
\vnbprop{rb4-equiv-act}{\vnbstmt{\begin{array}{@{}l@{}} \forall d, m, n, a, b, s. \\ \quad \operatorname{is\text{-}module}(d) \wedge \\ \quad m \in \mathbb{N} \wedge \\ \quad n \in \mathbb{N} \wedge \\ \quad 1 \leq m \wedge \\ \quad 1 \leq n \wedge \\ \quad a \in \operatorname{mat}(m, n, \operatorname{carr}(\operatorname{scal}(d))) \wedge \\ \quad s \in \operatorname{mat}(n, 1, \operatorname{vec}(d)) \wedge \\ \quad \operatorname{mat\text{-}equiv}(\operatorname{scal}(d), m, n, a, b) \Rightarrow \\ \quad \quad b \in \operatorname{mat}(m, n, \operatorname{carr}(\operatorname{scal}(d))) \wedge \\ \quad \quad \exists u. \\ \quad \quad \quad \operatorname{is\text{-}invertible\text{-}mat}(\operatorname{scal}(d), m, u) \wedge \\ \quad \quad \quad \exists w. \\ \quad \quad \quad \quad (\operatorname{is\text{-}invertible\text{-}mat}(\operatorname{scal}(d), n, w) \wedge (\operatorname{matact}(d, u, \operatorname{matact}(d, a, s)) = \operatorname{matact}(d, b, \operatorname{matact}(d, w, s)))) \end{array}}}
\vnbprop{submodule-of-free-is-free}{\vnbstmt{\begin{array}{@{}l@{}} \forall d.; \forall n \in \mathbb{N}, u \in \operatorname{mat}(n, 1, \operatorname{vec}(d)), m. \\ \quad \operatorname{is\text{-}module}(d) \wedge \\ \quad \operatorname{is\text{-}euclidean\text{-}ring}(\operatorname{scal}(d)) \wedge \\ \quad \operatorname{generates}(d, n, u) \wedge \\ \quad \operatorname{rel\text{-}free}(d, n, u) \wedge \\ \quad \operatorname{is\text{-}submodule}(d, m) \Rightarrow \\ \quad \quad \exists k. \\ \quad \quad \quad (((k \in \mathbb{N}) \wedge (k \leq n)) \wedge \exists w \in \operatorname{mat}(k, 1, \operatorname{vec}(d)).\; (\operatorname{spans}(d, k, w, m) \wedge \operatorname{rel\text{-}free}(d, k, w))) \end{array}}}

\vnbbookfile{theorem-library/rank-bound-proof}{theorem-library/rank-bound-proof}
\noindent{\small rank-bound-proof.scm -- Proposition 3.41 (algebraic-numbers.pdf ch.3 sec 8.2), the payoff of the whole Smith arc:}\par
\vnbprop{free-length-le-generators}{\vnbstmt{\begin{array}{@{}l@{}} \forall d, n, m, u, v. \\ \quad \operatorname{is\text{-}module}(d) \wedge \\ \quad \operatorname{is\text{-}euclidean\text{-}ring}(\operatorname{scal}(d)) \wedge \\ \quad n \in \mathbb{N} \wedge \\ \quad m \in \mathbb{N} \wedge \\ \quad u \in \operatorname{mat}(n, 1, \operatorname{vec}(d)) \wedge \\ \quad v \in \operatorname{mat}(m, 1, \operatorname{vec}(d)) \wedge \\ \quad \operatorname{generates}(d, n, u) \wedge \\ \quad \operatorname{rel\text{-}free}(d, m, v) \Rightarrow \\ \quad \quad m \leq n \end{array}}}

\part{Analysis}
\noindent Metric spaces and their topology, completeness, compactness and continuity; normed fields and normed vector spaces; sequences and series of real and complex numbers; the calculus of one real variable on an interval: derivatives, the mean value theorem, Taylor, antiderivatives, regulated functions and the integral; roads and line integrals; Ascoli, Baire, Hahn-Banach, Bernstein.
\par\noindent\textit{Builds on:} Linear Algebra.
\par\noindent{\small 331 files.}

\vnbbookfile{structure-library/metric-space}{structure-library/metric-space}
\noindent{\small metric-space.scm -- METRIC-SPACE structure}\par
\vnbstructdef{metric-space}{A metric space $s$ is a 2-tuple whose components are listed below.\par
\begin{enumerate}[nosep]
\item $\operatorname{pts}(s) = \operatorname{nth}(1, s)$, a carrier (\texttt{carriers}): $\operatorname{pts}(s) \in \mathbf{Set}$.
\item $\operatorname{dist}(s) = \operatorname{nth}(2, s)$, an operation (\texttt{op}): $\operatorname{dist}(s) \in ((\operatorname{pts}(s) \times \operatorname{pts}(s)) \to \mathbb{R})$.
\end{enumerate}
The laws (the \texttt{property} and \texttt{law} clauses):
\begin{itemize}[nosep]
\item $\operatorname{is\text{-}metric}(\operatorname{dist}(s), \operatorname{pts}(s))$
\end{itemize}
The characteristic law (the axiom \texttt{is-metric-space}): $\operatorname{is\text{-}metric\text{-}space}(s)$ holds if and only if $\operatorname{length}(s) = 2$, every component has the typing listed with it, and every law holds.
}
\par\medskip\noindent\textbf{Definition} (\texttt{hom-metric-space}).\;$\operatorname{hom\text{-}metric\text{-}space}(a, b) :=$
\begin{equation*}
\{\operatorname{homf\_} \in (\operatorname{pts}(a) \to \operatorname{pts}(b)) : \operatorname{is\text{-}hom\text{-}metric\text{-}space}(a, b, \operatorname{homf\_})\}
\end{equation*}

\vnbdef{metric-space-class}{\vnbstmt{\forall s.\; ((s \in \operatorname{metric\text{-}space}) \Leftrightarrow \operatorname{is\text{-}metric\text{-}space}(s))}}
\vnbdef{is-hom-metric-space-def}{\vnbstmt{\begin{array}{@{}l@{}} \forall a, b, f. \\ \quad \operatorname{is\text{-}hom\text{-}metric\text{-}space}(a, b, f) \\ \quad \Leftrightarrow\; \operatorname{is\text{-}metric\text{-}space}(a) \wedge \\ \quad \operatorname{is\text{-}metric\text{-}space}(b) \wedge \\ \quad f \in (\operatorname{pts}(a) \to \operatorname{pts}(b)) \wedge \\ \quad \forall x_{1}, x_{2} \in \operatorname{pts}(a). \\ \quad \quad \operatorname{dist}(a)(x_{1}, x_{2}) = \operatorname{dist}(b)(f(x_{1}), f(x_{2})) \end{array}}}

\vnbbookfile{structure-library/metric-topology}{structure-library/metric-topology}
\noindent{\small metric-topology.scm -- balls, r-nets, total-boundedness.}\par
\par\medskip\noindent\textbf{Definition} (\texttt{ball}).\;$\operatorname{ball}(s, c, r) :=$
\begin{equation*}
\{y \in \operatorname{pts}(s) : ((\operatorname{dist}(s)(c, y) \leq r) \wedge \neg (\operatorname{dist}(s)(c, y) = r))\}
\end{equation*}

\par\medskip\noindent\textbf{Definition} (\texttt{is-r-net}).\;For all $s$, $f$, $a$ and $r$, 
\begin{equation*}
\begin{array}{@{}l@{}} \operatorname{is\text{-}r\text{-}net}(s, f, a, r) \\ \Leftrightarrow\; f \subseteq a \wedge \\ \forall p \in a. \exists c \in f. \\ \quad (\operatorname{dist}(s)(c, p) \leq r) \wedge \neg (\operatorname{dist}(s)(c, p) = r) \end{array}
\end{equation*}

\par\medskip\noindent\textbf{Definition} (\texttt{totally-bounded}).\;For every $s$, 
\begin{equation*}
\begin{array}{@{}l@{}} \operatorname{totally\text{-}bounded}(s) \\ \Leftrightarrow\; \operatorname{is\text{-}metric\text{-}space}(s) \wedge \\ \forall r. \\ \quad ((r \in \mathbb{R}) \wedge ((0 \leq r) \wedge \neg (0 = r))) \Rightarrow  \\ \quad \quad \exists f.\; ((|f| \in \mathbb{N}) \wedge \operatorname{is\text{-}r\text{-}net}(s, f, \operatorname{pts}(s), r)) \end{array}
\end{equation*}

\par\medskip\noindent\textbf{Definition} (\texttt{closed-ball}).\;$\operatorname{closed\text{-}ball}(s, c, r) :=$ $\{y \in \operatorname{pts}(s) : (\operatorname{dist}(s)(c, y) \leq r)\}$

\vnbbookfile{structure-library/metric-completeness}{structure-library/metric-completeness}
\noindent{\small metric-completeness.scm -- Cauchy sequences, convergence, completeness on a generic METRIC-SPACE.}\par
\par\medskip\noindent\textbf{Definition} (\texttt{is-cauchy-seq}).\;For all $s$ and $f$, 
\begin{equation*}
\begin{array}{@{}l@{}} \operatorname{is\text{-}cauchy\text{-}seq}(s, f) \\ \Leftrightarrow\; \operatorname{is\text{-}metric\text{-}space}(s) \wedge \\ f \in (\mathbb{N} \to \operatorname{pts}(s)) \wedge \\ \forall \mathit{eps}. \\ \quad \operatorname{pos\text{-}rr}(\mathit{eps}) \Rightarrow  \\ \quad \quad \exists n \in \mathbb{N}. \forall m, \operatorname{n\_} \in \mathbb{N}. \\ \quad \quad \quad ((n \leq m) \wedge (n \leq \operatorname{n\_})) \Rightarrow (\operatorname{dist}(s)(f(m), f(\operatorname{n\_})) \leq \mathit{eps}) \end{array}
\end{equation*}

\par\medskip\noindent\textbf{Definition} (\texttt{converges-to}).\;For all $s$, $f$ and $l$, 
\begin{equation*}
\begin{array}{@{}l@{}} \operatorname{converges\text{-}to}(s, f, l) \\ \Leftrightarrow\; \operatorname{is\text{-}metric\text{-}space}(s) \wedge \\ f \in (\mathbb{N} \to \operatorname{pts}(s)) \wedge \\ l \in \operatorname{pts}(s) \wedge \\ \forall \mathit{eps}. \\ \quad \operatorname{pos\text{-}rr}(\mathit{eps}) \Rightarrow  \\ \quad \quad \exists n \in \mathbb{N}. \forall \operatorname{n\_} \in \mathbb{N}. \\ \quad \quad \quad (n \leq \operatorname{n\_}) \Rightarrow (\operatorname{dist}(s)(f(\operatorname{n\_}), l) \leq \mathit{eps}) \end{array}
\end{equation*}

\par\medskip\noindent\textbf{Definition} (\texttt{converges}).\;For all $s$ and $f$,  $\operatorname{converges}(s, f) \Leftrightarrow \exists l.\; \operatorname{converges\text{-}to}(s, f, l)$
\par\medskip\noindent\textbf{Definition} (\texttt{is-complete}).\;For every $s$, 
\begin{equation*}
\operatorname{is\text{-}complete}(s) \Leftrightarrow (\operatorname{is\text{-}metric\text{-}space}(s) \wedge \forall f.\; (\operatorname{is\text{-}cauchy\text{-}seq}(s, f) \Rightarrow \operatorname{converges}(s, f)))
\end{equation*}

\vnbdef{complete-cauchy-converges}{\vnbstmt{\forall s, f.\; ((\operatorname{is\text{-}complete}(s) \wedge \operatorname{is\text{-}cauchy\text{-}seq}(s, f)) \Rightarrow \operatorname{converges}(s, f))}}

\vnbbookfile{structure-library/metric-continuity}{structure-library/metric-continuity}
\noindent{\small RETIRED 2026-09-17 (proven): metric-hom-is-continuous -- theorem-library/rake-analysis2.scm metric-continuity.scm -- continuous maps between metric spaces.}\par
\par\medskip\noindent\textbf{Definition} (\texttt{is-continuous-at}).\;For all $s$, $t$, $f$ and $a$, 
\begin{equation*}
\begin{array}{@{}l@{}} \operatorname{is\text{-}continuous\text{-}at}(s, t, f, a) \\ \Leftrightarrow\; \operatorname{is\text{-}metric\text{-}space}(s) \wedge \\ \operatorname{is\text{-}metric\text{-}space}(t) \wedge \\ f \in (\operatorname{pts}(s) \to \operatorname{pts}(t)) \wedge \\ a \in \operatorname{pts}(s) \wedge \\ \forall \mathit{eps}. \\ \quad \operatorname{pos\text{-}rr}(\mathit{eps}) \Rightarrow  \\ \quad \quad \exists \delta. \\ \quad \quad \quad \operatorname{pos\text{-}rr}(\delta) \wedge \\ \quad \quad \quad \forall b \in \operatorname{pts}(s). \\ \quad \quad \quad \quad (\operatorname{dist}(s)(a, b) \leq \delta) \Rightarrow (\operatorname{dist}(t)(f(a), f(b)) \leq \mathit{eps}) \end{array}
\end{equation*}

\par\medskip\noindent\textbf{Definition} (\texttt{is-continuous}).\;For all $s$, $t$ and $f$, 
\begin{equation*}
\begin{array}{@{}l@{}} \operatorname{is\text{-}continuous}(s, t, f) \\ \Leftrightarrow\; \operatorname{is\text{-}metric\text{-}space}(s) \wedge \\ \operatorname{is\text{-}metric\text{-}space}(t) \wedge \\ f \in (\operatorname{pts}(s) \to \operatorname{pts}(t)) \wedge \\ \forall a \in \operatorname{pts}(s).\; \operatorname{is\text{-}continuous\text{-}at}(s, t, f, a) \end{array}
\end{equation*}

\par\medskip\noindent\textbf{Definition} (\texttt{is-uniformly-continuous}).\;For all $s$, $t$ and $f$, 
\begin{equation*}
\begin{array}{@{}l@{}} \operatorname{is\text{-}uniformly\text{-}continuous}(s, t, f) \\ \Leftrightarrow\; \operatorname{is\text{-}metric\text{-}space}(s) \wedge \\ \operatorname{is\text{-}metric\text{-}space}(t) \wedge \\ f \in (\operatorname{pts}(s) \to \operatorname{pts}(t)) \wedge \\ \forall \mathit{eps}. \\ \quad \operatorname{pos\text{-}rr}(\mathit{eps}) \Rightarrow  \\ \quad \quad \exists \delta. \\ \quad \quad \quad \operatorname{pos\text{-}rr}(\delta) \wedge \\ \quad \quad \quad \forall a, b \in \operatorname{pts}(s). \\ \quad \quad \quad \quad (\operatorname{dist}(s)(a, b) \leq \delta) \Rightarrow (\operatorname{dist}(t)(f(a), f(b)) \leq \mathit{eps}) \end{array}
\end{equation*}

\par\medskip\noindent\textbf{Definition} (\texttt{is-isometry}).\;For all $s$, $t$ and $f$, 
\begin{equation*}
\begin{array}{@{}l@{}} \operatorname{is\text{-}isometry}(s, t, f) \\ \Leftrightarrow\; \operatorname{is\text{-}metric\text{-}space}(s) \wedge \\ \operatorname{is\text{-}metric\text{-}space}(t) \wedge \\ f \in (\operatorname{pts}(s) \to \operatorname{pts}(t)) \wedge \\ \forall x, y \in \operatorname{pts}(s).\; (\operatorname{dist}(t)(f(x), f(y)) = \operatorname{dist}(s)(x, y)) \end{array}
\end{equation*}


\vnbbookfile{structure-library/metric-open-sets}{structure-library/metric-open-sets}
\noindent{\small RETIRED 2026-09-17 (proven): continuous-implies-open-preimage -- theorem-library/rake-cont-preimage.scm (modulo 0, via dk-halve!; this file now holds no support at all). RETIRED 2026-09-17 (proven): empty-is-open, carrier-is-open, union-of-opens-open, inter-of-opens-open, preimage-complement, ...}\par
\par\medskip\noindent\textbf{Definition} (\texttt{is-open}).\;For all $s$ and $u$, 
\begin{equation*}
\begin{array}{@{}l@{}} \operatorname{is\text{-}open}(s, u) \\ \Leftrightarrow\; \operatorname{is\text{-}metric\text{-}space}(s) \wedge \\ u \subseteq \operatorname{pts}(s) \wedge \\ \forall y \in u.\; \exists r.\; (\operatorname{pos\text{-}rr}(r) \wedge (\operatorname{ball}(s, y, r) \subseteq u)) \end{array}
\end{equation*}

\par\medskip\noindent\textbf{Definition} (\texttt{preimage}).\;$\operatorname{preimage}(s, f, v) :=$ $\{a \in \operatorname{pts}(s) : (f(a) \in v)\}$
\par\medskip\noindent\textbf{Definition} (\texttt{is-closed}).\;For all $s$ and $a$, 
\begin{equation*}
\operatorname{is\text{-}closed}(s, a) \Leftrightarrow (\operatorname{is\text{-}metric\text{-}space}(s) \wedge ((a \subseteq \operatorname{pts}(s)) \wedge \operatorname{is\text{-}open}(s, \operatorname{complement\text{-}in}(\operatorname{pts}(s), a))))
\end{equation*}

\vnbdef{preimage-membership}{\vnbstmt{\forall s, f, v, a.\; ((a \in \operatorname{preimage}(s, f, v)) \Leftrightarrow ((a \in \operatorname{pts}(s)) \wedge (f(a) \in v)))}}

\vnbbookfile{structure-library/metric-connected}{structure-library/metric-connected}
\noindent{\small metric-connected.scm -- CONNECTEDNESS OF A SUBSET OF A METRIC SPACE. DEFINITION ONLY; the laws are PROVEN in theorem-library/metric-connected-laws.scm. Agent CA-5 of week 3 of the October roadmap, 2026-10-04 (a NEW DEFINITION, the user's to confirm).}\par
\par\medskip\noindent\textbf{Definition} (\texttt{is-connected}).\;For all $\operatorname{cns\_}$ and $\operatorname{cna\_}$, 
\begin{equation*}
\begin{array}{@{}l@{}} \operatorname{is\text{-}connected}(\operatorname{cns\_}, \operatorname{cna\_}) \\ \Leftrightarrow\; \operatorname{is\text{-}metric\text{-}space}(\operatorname{cns\_}) \wedge \\ \operatorname{cna\_} \subseteq \operatorname{pts}(\operatorname{cns\_}) \wedge \\ \forall \operatorname{cnu\_}. \\ \quad \operatorname{is\text{-}open}(\operatorname{cns\_}, \operatorname{cnu\_}) \Rightarrow  \\ \quad \quad \forall \operatorname{cnv\_}. \\ \quad \quad \quad \operatorname{is\text{-}open}(\operatorname{cns\_}, \operatorname{cnv\_}) \Rightarrow  \\ \quad \quad \quad (\operatorname{cna\_} \subseteq (\operatorname{cnu\_} \cup \operatorname{cnv\_})) \Rightarrow  \\ \quad \quad \quad \forall \operatorname{cnx\_} \in \operatorname{cna\_}.\; \neg ((\operatorname{cnx\_} \in \operatorname{cnu\_}) \wedge (\operatorname{cnx\_} \in \operatorname{cnv\_})) \Rightarrow  \\ \quad \quad \quad \exists \operatorname{cnx\_} \in \operatorname{cna\_}.\; (\operatorname{cnx\_} \in \operatorname{cnu\_}) \Rightarrow  \\ \quad \quad \quad \quad \neg \exists \operatorname{cnx\_} \in \operatorname{cna\_}.\; (\operatorname{cnx\_} \in \operatorname{cnv\_}) \end{array}
\end{equation*}


\vnbbookfile{structure-library/top-space}{structure-library/top-space}
\noindent{\small top-space.scm -- TOPOLOGICAL SPACES, and the functor Met -$>$ Top}\par
\vnbstructdef{top-space}{A top space $s$ is a 2-tuple whose components are listed below.\par
\begin{enumerate}[nosep]
\item $\operatorname{pts}(s) = \operatorname{nth}(1, s)$, a carrier (\texttt{carriers}): $\operatorname{pts}(s) \in \mathbf{Set}$.
\item $\operatorname{opens}(s) = \operatorname{nth}(2, s)$, a family of subsets (\texttt{family}): $\operatorname{opens}(s) \in \operatorname{power}(\operatorname{power}(\operatorname{pts}(s)))$.
\end{enumerate}
The laws (the \texttt{property} and \texttt{law} clauses):
\begin{itemize}[nosep]
\item $\emptyset \in \operatorname{opens}(s)$
\item $\operatorname{pts}(s) \in \operatorname{opens}(s)$
\item $\forall u, v \in \operatorname{opens}(s).\; ((u \cap v) \in \operatorname{opens}(s))$
\item $\forall \mathit{fam}.\; ((\mathit{fam} \subseteq \operatorname{opens}(s)) \Rightarrow (\operatorname{big\text{-}union}(u, \mathit{fam}, u) \in \operatorname{opens}(s)))$
\end{itemize}
The characteristic law (the axiom \texttt{is-top-space}): $\operatorname{is\text{-}top\text{-}space}(s)$ holds if and only if $\operatorname{length}(s) = 2$, every component has the typing listed with it, and every law holds.
}
\par\medskip\noindent\textbf{Definition} (\texttt{metric-top}).\;$\operatorname{metric\text{-}top}(\mathit{md}) :=$ $\langle \operatorname{pts}(\mathit{md}), \{u \in \operatorname{power}(\operatorname{pts}(\mathit{md})) : \operatorname{is\text{-}open}(\mathit{md}, u)\} \rangle$
\par\medskip\noindent\textbf{Definition} (\texttt{is-hausdorff}).\;For every $s$, 
\begin{equation*}
\begin{array}{@{}l@{}} \operatorname{is\text{-}hausdorff}(s) \\ \Leftrightarrow\; \operatorname{is\text{-}top\text{-}space}(s) \wedge \\ \forall a, b \in \operatorname{pts}(s). \\ \quad \neg (a \simeq b) \Rightarrow  \\ \quad \quad \exists u, v \in \operatorname{opens}(s). \\ \quad \quad \quad (a \in u) \wedge ((b \in v) \wedge ((u \cap v) \simeq \emptyset)) \end{array}
\end{equation*}

\par\medskip\noindent\textbf{Definition} (\texttt{hom-top-space}).\;$\operatorname{hom\text{-}top\text{-}space}(a, b) :=$
\begin{equation*}
\{\operatorname{homf\_} \in (\operatorname{pts}(a) \to \operatorname{pts}(b)) : \operatorname{is\text{-}hom\text{-}top\text{-}space}(a, b, \operatorname{homf\_})\}
\end{equation*}

\par\medskip\noindent\textbf{Definition} (\texttt{hom-metrizable-top-space}).\;$\operatorname{hom\text{-}metrizable\text{-}top\text{-}space}(a, b) :=$
\begin{equation*}
\{\operatorname{homf\_} \in (\operatorname{pts}(a) \to \operatorname{pts}(b)) : \operatorname{is\text{-}hom\text{-}metrizable\text{-}top\text{-}space}(a, b, \operatorname{homf\_})\}
\end{equation*}

\vnbdef{top-space-class}{\vnbstmt{\forall s.\; ((s \in \operatorname{top\text{-}space}) \Leftrightarrow \operatorname{is\text{-}top\text{-}space}(s))}}
\vnbdef{metrizable-top-space-class}{\vnbstmt{\forall s.\; ((s \in \operatorname{metrizable\text{-}top\text{-}space}) \Leftrightarrow \operatorname{is\text{-}metrizable\text{-}top\text{-}space}(s))}}
\vnbdef{is-metrizable-top-space-def}{\vnbstmt{\begin{array}{@{}l@{}} \forall s. \\ \quad \operatorname{is\text{-}metrizable\text{-}top\text{-}space}(s) \\ \quad \Leftrightarrow\; \operatorname{is\text{-}top\text{-}space}(s) \wedge \\ \quad \exists d. \\ \quad \quad \operatorname{is\text{-}metric\text{-}space}(d) \wedge (s \simeq \langle \operatorname{pts}(d), \{u \in \operatorname{power}(\operatorname{pts}(d)) : \operatorname{is\text{-}open}(d, u)\} \rangle) \end{array}}}
\vnbdef{is-hom-top-space-def}{\vnbstmt{\begin{array}{@{}l@{}} \forall s, t, f. \\ \quad \operatorname{is\text{-}hom\text{-}top\text{-}space}(s, t, f) \\ \quad \Leftrightarrow\; \operatorname{is\text{-}top\text{-}space}(s) \wedge \\ \quad \operatorname{is\text{-}top\text{-}space}(t) \wedge \\ \quad f \in (\operatorname{pts}(s) \to \operatorname{pts}(t)) \wedge \\ \quad \forall u \in \operatorname{opens}(t).\; (\operatorname{preimage}(s, f, u) \in \operatorname{opens}(s)) \end{array}}}
\vnbdef{is-hom-metrizable-top-space-def}{\vnbstmt{\begin{array}{@{}l@{}} \forall s, t, f. \\ \quad \operatorname{is\text{-}hom\text{-}metrizable\text{-}top\text{-}space}(s, t, f) \\ \quad \Leftrightarrow\; \operatorname{is\text{-}metrizable\text{-}top\text{-}space}(s) \wedge \\ \quad \operatorname{is\text{-}metrizable\text{-}top\text{-}space}(t) \wedge \\ \quad f \in (\operatorname{pts}(s) \to \operatorname{pts}(t)) \wedge \\ \quad \forall u \in \operatorname{opens}(t).\; (\operatorname{preimage}(s, f, u) \in \operatorname{opens}(s)) \end{array}}}

\vnbbookfile{structure-library/compactness}{structure-library/compactness}
\noindent{\small compactness.scm -- compactness of a metric space and the four-way characterization (calculus.pdf Prop 3.12).  Supplies the missing vocabulary (IS-COMPACT, IS-OPEN-COVER, CLUSTER-POINT, HAS-FIP) and states the characterization as warranted supports (library phase); the proof}\par
\par\medskip\noindent\textbf{Definition} (\texttt{is-compact}).\;For every $s$, 
\begin{equation*}
\begin{array}{@{}l@{}} \operatorname{is\text{-}compact}(s) \\ \Leftrightarrow\; \operatorname{is\text{-}metric\text{-}space}(s) \wedge \\ \forall c. \\ \quad \operatorname{is\text{-}open\text{-}cover}(s, c) \Rightarrow  \\ \quad \quad \exists f.\; ((f \subseteq c) \wedge ((|f| \in \mathbb{N}) \wedge \operatorname{is\text{-}open\text{-}cover}(s, f))) \end{array}
\end{equation*}

\par\medskip\noindent\textbf{Definition} (\texttt{is-open-cover}).\;For all $s$ and $c$, 
\begin{equation*}
\begin{array}{@{}l@{}} \operatorname{is\text{-}open\text{-}cover}(s, c) \\ \Leftrightarrow\; (\operatorname{is\text{-}metric\text{-}space}(s) \wedge (\forall u \in c.\; \operatorname{is\text{-}open}(s, u) \wedge (\operatorname{big\text{-}union}(u, c, u) \simeq \operatorname{pts}(s)))) \end{array}
\end{equation*}

\par\medskip\noindent\textbf{Definition} (\texttt{cluster-point}).\;For all $s$, $f$ and $x$, 
\begin{equation*}
\begin{array}{@{}l@{}} \operatorname{cluster\text{-}point}(s, f, x) \\ \Leftrightarrow\; \operatorname{is\text{-}metric\text{-}space}(s) \wedge \\ f \in (\mathbb{N} \to \operatorname{pts}(s)) \wedge \\ x \in \operatorname{pts}(s) \wedge \\ \forall \mathit{eps}. \\ \quad \operatorname{pos\text{-}rr}(\mathit{eps}) \Rightarrow  \\ \quad \quad \forall m \in \mathbb{N}.\; \exists n \in \mathbb{N}.\; ((m \leq n) \wedge (\operatorname{dist}(s)(f(n), x) < \mathit{eps})) \end{array}
\end{equation*}

\par\medskip\noindent\textbf{Definition} (\texttt{has-fip}).\;For all $s$ and $c$, 
\begin{equation*}
\begin{array}{@{}l@{}} \operatorname{has\text{-}fip}(s, c) \\ \Leftrightarrow\; \operatorname{is\text{-}metric\text{-}space}(s) \wedge \\ \exists a_{0}.\; (a_{0} \in c) \wedge \\ \forall a \in c.\; \operatorname{is\text{-}closed}(s, a) \wedge \\ \forall f. \\ \quad (f \subseteq c) \Rightarrow  \\ \quad (|f| \in \mathbb{N}) \Rightarrow  \\ \quad \exists a_{0}.\; (a_{0} \in f) \Rightarrow \exists p.\; (p \in \operatorname{intersection\text{-}of}(f)) \end{array}
\end{equation*}

\par\medskip\noindent\textbf{Definition} (\texttt{ball-cover}).\;$\operatorname{ball\text{-}cover}(s, r) :=$ $\operatorname{image}((\lambda c \in \operatorname{pts}(s).\; \operatorname{ball}(s, c, r)), \operatorname{pts}(s))$
\par\medskip\noindent\textbf{Definition} (\texttt{centres}).\;$\operatorname{centres}(s, b, r) :=$ $\{c \in \operatorname{pts}(s) : (\operatorname{ball}(s, c, r) = b)\}$
\par\medskip\noindent\textbf{Definition} (\texttt{centre-set}).\;$\operatorname{centre\text{-}set}(s, r, f) :=$ $\operatorname{image}((\lambda b \in f.\; \varepsilon(\operatorname{centres}(s, b, r))), f)$

\vnbbookfile{structure-library/top-continuity}{structure-library/top-continuity}
\noindent{\small top-continuity.scm -- COMPACTNESS of a topological space, and the SUBSPACE TOPOLOGY.  DEFINITIONS AND NOTATION ONLY; every law is proven in theorem-library/metrizable-compact-image.scm.}\par
\par\medskip\noindent\textbf{Definition} (\texttt{is-compact-t}).\;For every $s$, 
\begin{equation*}
\begin{array}{@{}l@{}} \operatorname{is\text{-}compact\text{-}t}(s) \\ \Leftrightarrow\; \operatorname{is\text{-}top\text{-}space}(s) \wedge \\ \forall c. \\ \quad \forall u \in c.\; (u \in \operatorname{opens}(s)) \Rightarrow  \\ \quad (\operatorname{big\text{-}union}(u, c, u) \simeq \operatorname{pts}(s)) \Rightarrow  \\ \quad \quad \exists f.\; ((f \subseteq c) \wedge ((|f| \in \mathbb{N}) \wedge (\operatorname{big\text{-}union}(u, f, u) \simeq \operatorname{pts}(s)))) \end{array}
\end{equation*}

\par\medskip\noindent\textbf{Definition} (\texttt{subspace-top}).\;$\operatorname{subspace\text{-}top}(t, a) :=$
\begin{equation*}
\langle a, \{\operatorname{stw\_} \in \operatorname{power}(a) : \exists \operatorname{stu\_} \in \operatorname{opens}(t).\; (\operatorname{stw\_} = (\operatorname{stu\_} \cap a))\} \rangle
\end{equation*}


\vnbbookfile{structure-library/separable}{structure-library/separable}
\noindent{\small separable.scm -- separability of a metric space (a dense sequence), the vocabulary and the one support the bridge compact =$>$ separable still needs. Added 2026-07-22 for the Ascoli-Arzela arc.}\par
\par\medskip\noindent\textbf{Definition} (\texttt{is-separable}).\;For every $s$, 
\begin{equation*}
\begin{array}{@{}l@{}} \operatorname{is\text{-}separable}(s) \\ \Leftrightarrow\; \operatorname{is\text{-}metric\text{-}space}(s) \wedge \\ \exists \mathit{dseq} \in (\mathbb{N} \to \operatorname{pts}(s)). \forall x \in \operatorname{pts}(s), \mathit{eps}. \\ \quad \operatorname{pos\text{-}rr}(\mathit{eps}) \Rightarrow \exists n \in \mathbb{N}.\; (\operatorname{dist}(s)(x, \operatorname{dseq}(n)) < \mathit{eps}) \end{array}
\end{equation*}


\vnbbookfile{structure-library/normed-field}{structure-library/normed-field}
\noindent{\small normed-field.scm -- NORMED-FIELD: a commutative ring with multiplicative inverses and a structurally-carried norm into RR.}\par
\vnbstructdef{normed-field}{A normed field $s$ is a 7-tuple whose components are listed below.\par
\begin{enumerate}[nosep]
\item $\operatorname{carr}(s) = \operatorname{nth}(1, s)$, a carrier (\texttt{carriers}): $\operatorname{carr}(s) \in \mathbf{Set}$.
\item $\operatorname{add}(s) = \operatorname{nth}(2, s)$, an operation (\texttt{op}): $\operatorname{add}(s) \in ((\operatorname{carr}(s) \times \operatorname{carr}(s)) \to \operatorname{carr}(s))$.
\item $\operatorname{mul}(s) = \operatorname{nth}(3, s)$, an operation (\texttt{op}): $\operatorname{mul}(s) \in ((\operatorname{carr}(s) \times \operatorname{carr}(s)) \to \operatorname{carr}(s))$.
\item $\operatorname{neg}(s) = \operatorname{nth}(4, s)$, an operation (\texttt{op}): $\operatorname{neg}(s) \in (\operatorname{carr}(s) \to \operatorname{carr}(s))$.
\item $\operatorname{zero}(s) = \operatorname{nth}(5, s)$, a constant (\texttt{constant}): $\operatorname{zero}(s) \in \operatorname{carr}(s)$.
\item $\operatorname{one}(s) = \operatorname{nth}(6, s)$, a constant (\texttt{constant}): $\operatorname{one}(s) \in \operatorname{carr}(s)$.
\item $\operatorname{fnrm}(s) = \operatorname{nth}(7, s)$, an operation (\texttt{op}): $\operatorname{fnrm}(s) \in (\operatorname{carr}(s) \to \mathbb{R})$.
\end{enumerate}
The laws (the \texttt{property} and \texttt{law} clauses):
\begin{itemize}[nosep]
\item $\operatorname{is\text{-}associative}(\operatorname{add}(s), \operatorname{carr}(s))$
\item $\operatorname{is\text{-}commutative}(\operatorname{add}(s), \operatorname{carr}(s))$
\item $\operatorname{is\text{-}identity}(\operatorname{add}(s), \operatorname{zero}(s), \operatorname{carr}(s))$
\item $\operatorname{has\text{-}inverses}(\operatorname{add}(s), \operatorname{zero}(s), \operatorname{neg}(s), \operatorname{carr}(s))$
\item $\operatorname{is\text{-}associative}(\operatorname{mul}(s), \operatorname{carr}(s))$
\item $\operatorname{is\text{-}commutative}(\operatorname{mul}(s), \operatorname{carr}(s))$
\item $\operatorname{is\text{-}identity}(\operatorname{mul}(s), \operatorname{one}(s), \operatorname{carr}(s))$
\item $\operatorname{is\text{-}distributive}(\operatorname{add}(s), \operatorname{mul}(s), \operatorname{carr}(s))$
\item $\operatorname{is\text{-}norm}(\operatorname{fnrm}(s), \operatorname{add}(s), \operatorname{mul}(s), \operatorname{zero}(s), \operatorname{carr}(s))$
\item $\neg (\operatorname{zero}(s) = \operatorname{one}(s))$
\item
\begin{equation*}
\begin{array}{@{}l@{}} \forall \operatorname{a\_} \in \operatorname{carr}(s). \\ \quad \neg (\operatorname{a\_} = \operatorname{zero}(s)) \Rightarrow  \\ \quad \quad \exists \operatorname{b\_} \in \operatorname{carr}(s).\; (\neg (\operatorname{b\_} = \operatorname{zero}(s)) \wedge (\operatorname{mul}(s)(\operatorname{a\_}, \operatorname{b\_}) = \operatorname{one}(s))) \end{array}
\end{equation*}

\end{itemize}
The characteristic law (the axiom \texttt{is-normed-field}): $\operatorname{is\text{-}normed\text{-}field}(s)$ holds if and only if $\operatorname{length}(s) = 7$, every component has the typing listed with it, and every law holds.
}
\par\medskip\noindent\textbf{Definition} (\texttt{hom-normed-field}).\;$\operatorname{hom\text{-}normed\text{-}field}(a, b) :=$
\begin{equation*}
\{\operatorname{homf\_} \in (\operatorname{carr}(a) \to \operatorname{carr}(b)) : \operatorname{is\text{-}hom\text{-}normed\text{-}field}(a, b, \operatorname{homf\_})\}
\end{equation*}

\vnbdef{normed-field-class}{\vnbstmt{\forall s.\; ((s \in \operatorname{normed\text{-}field}) \Leftrightarrow \operatorname{is\text{-}normed\text{-}field}(s))}}
\vnbdef{is-hom-normed-field-def}{\vnbstmt{\begin{array}{@{}l@{}} \forall a, b, f. \\ \quad \operatorname{is\text{-}hom\text{-}normed\text{-}field}(a, b, f) \\ \quad \Leftrightarrow\; \operatorname{is\text{-}normed\text{-}field}(a) \wedge \\ \quad \operatorname{is\text{-}normed\text{-}field}(b) \wedge \\ \quad f \in (\operatorname{carr}(a) \to \operatorname{carr}(b)) \wedge \\ \quad \forall x_{1}, x_{2} \in \operatorname{carr}(a). \\ \quad \quad f(\operatorname{add}(a)(x_{1}, x_{2})) = \operatorname{add}(b)(f(x_{1}), f(x_{2})) \wedge \\ \quad \forall x_{1}, x_{2} \in \operatorname{carr}(a). \\ \quad \quad f(\operatorname{mul}(a)(x_{1}, x_{2})) = \operatorname{mul}(b)(f(x_{1}), f(x_{2})) \wedge \\ \quad \forall x_{1} \in \operatorname{carr}(a).\; (f(\operatorname{neg}(a)(x_{1})) = \operatorname{neg}(b)(f(x_{1}))) \wedge \\ \quad f(\operatorname{zero}(a)) = \operatorname{zero}(b) \wedge \\ \quad f(\operatorname{one}(a)) = \operatorname{one}(b) \wedge \\ \quad \forall x_{1} \in \operatorname{carr}(a).\; (\operatorname{fnrm}(a)(x_{1}) = \operatorname{fnrm}(b)(f(x_{1}))) \end{array}}}

\vnbbookfile{structure-library/normed-ag}{structure-library/normed-ag}
\noindent{\small normed-ag.scm -- NORMED-AG structure (normed abelian group)}\par
\vnbstructdef{normed-ag}{A normed ag $s$ is a 5-tuple whose components are listed below.\par
\begin{enumerate}[nosep]
\item $\operatorname{carr}(s) = \operatorname{nth}(1, s)$, a carrier (\texttt{carriers}): $\operatorname{carr}(s) \in \mathbf{Set}$.
\item $\operatorname{opr}(s) = \operatorname{nth}(2, s)$, an operation (\texttt{op}): $\operatorname{opr}(s) \in ((\operatorname{carr}(s) \times \operatorname{carr}(s)) \to \operatorname{carr}(s))$.
\item $\operatorname{iden}(s) = \operatorname{nth}(3, s)$, a constant (\texttt{constant}): $\operatorname{iden}(s) \in \operatorname{carr}(s)$.
\item $\operatorname{inv}(s) = \operatorname{nth}(4, s)$, an operation (\texttt{op}): $\operatorname{inv}(s) \in (\operatorname{carr}(s) \to \operatorname{carr}(s))$.
\item $\operatorname{nrm}(s) = \operatorname{nth}(5, s)$, an operation (\texttt{op}): $\operatorname{nrm}(s) \in (\operatorname{carr}(s) \to \mathbb{R})$.
\end{enumerate}
The laws (the \texttt{property} and \texttt{law} clauses):
\begin{itemize}[nosep]
\item $\operatorname{is\text{-}associative}(\operatorname{opr}(s), \operatorname{carr}(s))$
\item $\operatorname{is\text{-}identity}(\operatorname{opr}(s), \operatorname{iden}(s), \operatorname{carr}(s))$
\item $\operatorname{has\text{-}inverses}(\operatorname{opr}(s), \operatorname{iden}(s), \operatorname{inv}(s), \operatorname{carr}(s))$
\item $\operatorname{is\text{-}commutative}(\operatorname{opr}(s), \operatorname{carr}(s))$
\item $\operatorname{is\text{-}group\text{-}norm}(\operatorname{nrm}(s), \operatorname{opr}(s), \operatorname{inv}(s), \operatorname{iden}(s), \operatorname{carr}(s))$
\end{itemize}
The characteristic law (the axiom \texttt{is-normed-ag}): $\operatorname{is\text{-}normed\text{-}ag}(s)$ holds if and only if $\operatorname{length}(s) = 5$, every component has the typing listed with it, and every law holds.
}
\par\medskip\noindent\textbf{Definition} (\texttt{hom-normed-ag}).\;$\operatorname{hom\text{-}normed\text{-}ag}(a, b) :=$
\begin{equation*}
\{\operatorname{homf\_} \in (\operatorname{carr}(a) \to \operatorname{carr}(b)) : \operatorname{is\text{-}hom\text{-}normed\text{-}ag}(a, b, \operatorname{homf\_})\}
\end{equation*}

\vnbdef{is-hom-normed-ag-def}{\vnbstmt{\begin{array}{@{}l@{}} \forall a, b, f. \\ \quad \operatorname{is\text{-}hom\text{-}normed\text{-}ag}(a, b, f) \\ \quad \Leftrightarrow\; \operatorname{is\text{-}normed\text{-}ag}(a) \wedge \\ \quad \operatorname{is\text{-}normed\text{-}ag}(b) \wedge \\ \quad f \in (\operatorname{carr}(a) \to \operatorname{carr}(b)) \wedge \\ \quad \forall x_{1}, x_{2} \in \operatorname{carr}(a). \\ \quad \quad f(\operatorname{opr}(a)(x_{1}, x_{2})) = \operatorname{opr}(b)(f(x_{1}), f(x_{2})) \wedge \\ \quad f(\operatorname{iden}(a)) = \operatorname{iden}(b) \wedge \\ \quad \forall x_{1} \in \operatorname{carr}(a).\; (f(\operatorname{inv}(a)(x_{1})) = \operatorname{inv}(b)(f(x_{1}))) \wedge \\ \quad \forall x_{1} \in \operatorname{carr}(a).\; (\operatorname{nrm}(a)(x_{1}) = \operatorname{nrm}(b)(f(x_{1}))) \end{array}}}
\vnbdef{normed-ag-class}{\vnbstmt{\forall s.\; ((s \in \operatorname{normed\text{-}ag}) \Leftrightarrow \operatorname{is\text{-}normed\text{-}ag}(s))}}

\vnbbookfile{structure-library/normed-field-metric}{structure-library/normed-field-metric}
\noindent{\small normed-field-metric.scm -- the metric space underlying a normed field.}\par
\par\medskip\noindent\textbf{Definition} (\texttt{nf-metric-space}).\;$\operatorname{nf\text{-}metric\text{-}space}(\mathit{nf}) :=$
\begin{equation*}
\langle \operatorname{carr}(\mathit{nf}), (\lambda \langle x, y \rangle \in (\operatorname{carr}(\mathit{nf}) \times \operatorname{carr}(\mathit{nf})).\; \operatorname{fnrm}(\mathit{nf})(\operatorname{add}(\mathit{nf})(x, \operatorname{neg}(\mathit{nf})(y)))) \rangle
\end{equation*}

\vnbdef{nf-metric-distance}{\vnbstmt{\begin{array}{@{}l@{}} \forall f.; \forall x, y \in \operatorname{carr}(f). \\ \quad \operatorname{is\text{-}normed\text{-}field}(f) \Rightarrow  \\ \quad \quad \operatorname{dist}(\operatorname{nf\text{-}metric\text{-}space}(f))(x, y) = \operatorname{fnrm}(f)(\operatorname{add}(f)(x, \operatorname{neg}(f)(y))) \end{array}}}
\vnbdef{nf-metric-carrier}{\vnbstmt{\forall f.\; (\operatorname{pts}(\operatorname{nf\text{-}metric\text{-}space}(f)) \simeq \operatorname{carr}(f))}}

\vnbbookfile{structure-library/normed-ag-metric}{structure-library/normed-ag-metric}
\noindent{\small normed-ag-metric.scm -- the metric space underlying a normed abelian group.}\par
\par\medskip\noindent\textbf{Definition} (\texttt{nag-metric-space}).\;$\operatorname{nag\text{-}metric\text{-}space}(\mathit{nag}) :=$
\begin{equation*}
\langle \operatorname{carr}(\mathit{nag}), (\lambda \langle u, v \rangle \in (\operatorname{carr}(\mathit{nag}) \times \operatorname{carr}(\mathit{nag})).\; \operatorname{nrm}(\mathit{nag})(\operatorname{opr}(\mathit{nag})(u, \operatorname{inv}(\mathit{nag})(v)))) \rangle
\end{equation*}


\vnbbookfile{structure-library/normed-vector-space}{structure-library/normed-vector-space}
\noindent{\small normed-vector-space.scm -- NORMED-VECTOR-SPACE: a real normed vector space. Vocabulary only (no proofs).  Three decisions, as agreed:}\par
\vnbstructdef{normed-vector-space}{A normed vector space $s$ is a 7-tuple whose components are listed below.\par
\begin{enumerate}[nosep]
\item $\operatorname{scal}(s) = \operatorname{nth}(1, s)$, a substructure (\texttt{substructure}): $\operatorname{is\text{-}ring}(\operatorname{scal}(s))$.
\item $\operatorname{vec}(s) = \operatorname{nth}(2, s)$, a carrier (\texttt{carriers}): $\operatorname{vec}(s) \in \mathbf{Set}$.
\item $\operatorname{vadd}(s) = \operatorname{nth}(3, s)$, an operation (\texttt{op}): $\operatorname{vadd}(s) \in ((\operatorname{vec}(s) \times \operatorname{vec}(s)) \to \operatorname{vec}(s))$.
\item $\operatorname{vzero}(s) = \operatorname{nth}(4, s)$, a constant (\texttt{constant}): $\operatorname{vzero}(s) \in \operatorname{vec}(s)$.
\item $\operatorname{vneg}(s) = \operatorname{nth}(5, s)$, an operation (\texttt{op}): $\operatorname{vneg}(s) \in (\operatorname{vec}(s) \to \operatorname{vec}(s))$.
\item $\operatorname{act}(s) = \operatorname{nth}(6, s)$, an operation (\texttt{op}): $\operatorname{act}(s) \in ((\operatorname{carr}(\operatorname{scal}(s)) \times \operatorname{vec}(s)) \to \operatorname{vec}(s))$.
\item $\operatorname{vnrm}(s) = \operatorname{nth}(7, s)$, an operation (\texttt{op}): $\operatorname{vnrm}(s) \in (\operatorname{vec}(s) \to \mathbb{R})$.
\end{enumerate}
The laws (the \texttt{property} and \texttt{law} clauses):
\begin{itemize}[nosep]
\item $\operatorname{is\text{-}associative}(\operatorname{vadd}(s), \operatorname{vec}(s))$
\item $\operatorname{is\text{-}commutative}(\operatorname{vadd}(s), \operatorname{vec}(s))$
\item $\operatorname{is\text{-}identity}(\operatorname{vadd}(s), \operatorname{vzero}(s), \operatorname{vec}(s))$
\item $\operatorname{has\text{-}inverses}(\operatorname{vadd}(s), \operatorname{vzero}(s), \operatorname{vneg}(s), \operatorname{vec}(s))$
\item $\operatorname{scal}(s) = \operatorname{normed\text{-}field\text{-}as\text{-}commutative\text{-}ring}(\operatorname{rr\text{-}normed\text{-}field})$
\item
\begin{equation*}
\begin{array}{@{}l@{}} \forall \operatorname{r\_} \in \operatorname{carr}(\operatorname{scal}(s)), \operatorname{x\_}, \operatorname{y\_} \in \operatorname{vec}(s). \\ \quad \operatorname{act}(s)(\operatorname{r\_}, \operatorname{vadd}(s)(\operatorname{x\_}, \operatorname{y\_})) = \operatorname{vadd}(s)(\operatorname{act}(s)(\operatorname{r\_}, \operatorname{x\_}), \operatorname{act}(s)(\operatorname{r\_}, \operatorname{y\_})) \end{array}
\end{equation*}

\item
\begin{equation*}
\begin{array}{@{}l@{}} \forall \operatorname{r\_}, \operatorname{s\_} \in \operatorname{carr}(\operatorname{scal}(s)), \operatorname{x\_} \in \operatorname{vec}(s). \\ \quad \operatorname{act}(s)(\operatorname{add}(\operatorname{scal}(s))(\operatorname{r\_}, \operatorname{s\_}), \operatorname{x\_}) = \operatorname{vadd}(s)(\operatorname{act}(s)(\operatorname{r\_}, \operatorname{x\_}), \operatorname{act}(s)(\operatorname{s\_}, \operatorname{x\_})) \end{array}
\end{equation*}

\item
\begin{equation*}
\begin{array}{@{}l@{}} \forall \operatorname{r\_}, \operatorname{s\_} \in \operatorname{carr}(\operatorname{scal}(s)), \operatorname{x\_} \in \operatorname{vec}(s). \\ \quad \operatorname{act}(s)(\operatorname{mul}(\operatorname{scal}(s))(\operatorname{r\_}, \operatorname{s\_}), \operatorname{x\_}) = \operatorname{act}(s)(\operatorname{r\_}, \operatorname{act}(s)(\operatorname{s\_}, \operatorname{x\_})) \end{array}
\end{equation*}

\item $\forall \operatorname{x\_} \in \operatorname{vec}(s).\; (\operatorname{act}(s)(\operatorname{one}(\operatorname{scal}(s)), \operatorname{x\_}) = \operatorname{x\_})$
\item $\forall \operatorname{x\_} \in \operatorname{vec}(s).\; (0 \leq \operatorname{vnrm}(s)(\operatorname{x\_}))$
\item $\forall \operatorname{x\_} \in \operatorname{vec}(s).\; ((\operatorname{vnrm}(s)(\operatorname{x\_}) = 0) \Leftrightarrow (\operatorname{x\_} = \operatorname{vzero}(s)))$
\item
\begin{equation*}
\begin{array}{@{}l@{}} \forall \operatorname{r\_} \in \operatorname{carr}(\operatorname{scal}(s)), \operatorname{x\_} \in \operatorname{vec}(s). \\ \quad \operatorname{vnrm}(s)(\operatorname{act}(s)(\operatorname{r\_}, \operatorname{x\_})) = (\lvert \operatorname{r\_} \rvert \cdot \operatorname{vnrm}(s)(\operatorname{x\_})) \end{array}
\end{equation*}

\item
\begin{equation*}
\begin{array}{@{}l@{}} \forall \operatorname{x\_}, \operatorname{y\_} \in \operatorname{vec}(s). \\ \quad \operatorname{vnrm}(s)(\operatorname{vadd}(s)(\operatorname{x\_}, \operatorname{y\_})) \leq (\operatorname{vnrm}(s)(\operatorname{x\_}) + \operatorname{vnrm}(s)(\operatorname{y\_})) \end{array}
\end{equation*}

\end{itemize}
The characteristic law (the axiom \texttt{is-normed-vector-space}): $\operatorname{is\text{-}normed\text{-}vector\text{-}space}(s)$ holds if and only if $\operatorname{length}(s) = 7$, every component has the typing listed with it, and every law holds.
}
\par\medskip\noindent\textbf{Definition} (\texttt{hom-normed-vector-space}).\;$\operatorname{hom\text{-}normed\text{-}vector\text{-}space}(a, b) :=$
\begin{equation*}
\{\operatorname{homf\_} \in (\operatorname{vec}(a) \to \operatorname{vec}(b)) : \operatorname{is\text{-}hom\text{-}normed\text{-}vector\text{-}space}(a, b, \operatorname{homf\_})\}
\end{equation*}

\vnbdef{normed-vector-space-class}{\vnbstmt{\forall s.\; ((s \in \operatorname{normed\text{-}vector\text{-}space}) \Leftrightarrow \operatorname{is\text{-}normed\text{-}vector\text{-}space}(s))}}
\vnbdef{normed-vector-space-as-module-is-module}{\vnbstmt{\begin{array}{@{}l@{}} \forall r. \\ \quad \operatorname{is\text{-}normed\text{-}vector\text{-}space}(r) \Rightarrow \operatorname{is\text{-}module}(\operatorname{normed\text{-}vector\text{-}space\text{-}as\text{-}module}(r)) \end{array}}}
\vnbdef{normed-vector-space-as-normed-ag-is-normed-ag}{\vnbstmt{\begin{array}{@{}l@{}} \forall r. \\ \quad \operatorname{is\text{-}normed\text{-}vector\text{-}space}(r) \Rightarrow  \\ \quad \quad \operatorname{is\text{-}normed\text{-}ag}(\operatorname{normed\text{-}vector\text{-}space\text{-}as\text{-}normed\text{-}ag}(r)) \end{array}}}
\vnbdef{is-hom-normed-vector-space-def}{\vnbstmt{\begin{array}{@{}l@{}} \forall a, b, f. \\ \quad \operatorname{is\text{-}hom\text{-}normed\text{-}vector\text{-}space}(a, b, f) \\ \quad \Leftrightarrow\; \operatorname{is\text{-}normed\text{-}vector\text{-}space}(a) \wedge \\ \quad \operatorname{is\text{-}normed\text{-}vector\text{-}space}(b) \wedge \\ \quad f \in (\operatorname{vec}(a) \to \operatorname{vec}(b)) \wedge \\ \quad \operatorname{scal}(a) = \operatorname{scal}(b) \wedge \\ \quad \forall x_{1}, x_{2} \in \operatorname{vec}(a). \\ \quad \quad f(\operatorname{vadd}(a)(x_{1}, x_{2})) = \operatorname{vadd}(b)(f(x_{1}), f(x_{2})) \wedge \\ \quad f(\operatorname{vzero}(a)) = \operatorname{vzero}(b) \wedge \\ \quad \forall x_{1} \in \operatorname{vec}(a).\; (f(\operatorname{vneg}(a)(x_{1})) = \operatorname{vneg}(b)(f(x_{1}))) \wedge \\ \quad \forall x_{1} \in \operatorname{carr}(\operatorname{scal}(a)), x_{2} \in \operatorname{vec}(a). \\ \quad \quad f(\operatorname{act}(a)(x_{1}, x_{2})) = \operatorname{act}(b)(x_{1}, f(x_{2})) \wedge \\ \quad \forall x_{1} \in \operatorname{vec}(a).\; (\operatorname{vnrm}(a)(x_{1}) = \operatorname{vnrm}(b)(f(x_{1}))) \end{array}}}

\vnbbookfile{structure-library/nvs-metric}{structure-library/nvs-metric}
\noindent{\small nvs-metric.scm -- the metric induced by the norm of a real normed vector space, as VOCABULARY.}\par
\par\medskip\noindent\textbf{Definition} (\texttt{nvs-metric-space}).\;$\operatorname{nvs\text{-}metric\text{-}space}(m) :=$
\begin{equation*}
\langle \operatorname{vec}(m), (\lambda \langle x, y \rangle \in (\operatorname{vec}(m) \times \operatorname{vec}(m)).\; \operatorname{vnrm}(m)(\operatorname{vadd}(m)(x, \operatorname{vneg}(m)(y)))) \rangle
\end{equation*}


\vnbbookfile{structure-library/linear-functional}{structure-library/linear-functional}
\noindent{\small linear-functional.scm -- bounded linear functionals and the dual norm on a real NORMED-VECTOR-SPACE.  Vocabulary only (no proofs).}\par
\par\medskip\noindent\textbf{Definition} (\texttt{is-linear-functional}).\;For all $m$ and $f$, 
\begin{equation*}
\begin{array}{@{}l@{}} \operatorname{is\text{-}linear\text{-}functional}(m, f) \\ \Leftrightarrow\; f \in (\operatorname{vec}(m) \to \mathbb{R}) \wedge \\ \forall \operatorname{x\_}, \operatorname{y\_} \in \operatorname{vec}(m).\; (f(\operatorname{vadd}(m)(\operatorname{x\_}, \operatorname{y\_})) = +(f(\operatorname{x\_}), f(\operatorname{y\_}))) \wedge \\ \forall \operatorname{r\_} \in \mathbb{R}, \operatorname{x\_} \in \operatorname{vec}(m).\; (f(\operatorname{act}(m)(\operatorname{r\_}, \operatorname{x\_})) = \cdot(\operatorname{r\_}, f(\operatorname{x\_}))) \end{array}
\end{equation*}

\par\medskip\noindent\textbf{Definition} (\texttt{is-bounded-linear-functional}).\;For all $m$ and $f$, 
\begin{equation*}
\begin{array}{@{}l@{}} \operatorname{is\text{-}bounded\text{-}linear\text{-}functional}(m, f) \\ \Leftrightarrow\; \operatorname{is\text{-}linear\text{-}functional}(m, f) \wedge \\ \exists \operatorname{c\_} \in \mathbb{R}. \\ \quad (0 \leq \operatorname{c\_}) \wedge \forall \operatorname{x\_} \in \operatorname{vec}(m).\; (\lvert f(\operatorname{x\_}) \rvert \leq \cdot(\operatorname{c\_}, \operatorname{vnrm}(m)(\operatorname{x\_}))) \end{array}
\end{equation*}

\par\medskip\noindent\textbf{Definition} (\texttt{dual-norm}).\;$\operatorname{dual\text{-}norm}(m, f) :=$
\begin{equation*}
(\iota\, \operatorname{c\_}.\;((\operatorname{c\_} \in \mathbb{R}) \wedge ((0 \leq \operatorname{c\_}) \wedge (\forall \operatorname{x\_} \in \operatorname{vec}(m).\; (\lvert f(\operatorname{x\_}) \rvert \leq (\operatorname{c\_} \cdot \operatorname{vnrm}(m)(\operatorname{x\_}))) \wedge \forall \operatorname{d\_}.\; (((\operatorname{d\_} \in \mathbb{R}) \wedge ((0 \leq \operatorname{d\_}) \wedge \forall \operatorname{x\_} \in \operatorname{vec}(m).\; (\lvert f(\operatorname{x\_}) \rvert \leq (\operatorname{d\_} \cdot \operatorname{vnrm}(m)(\operatorname{x\_}))))) \Rightarrow (\operatorname{c\_} \leq \operatorname{d\_}))))))
\end{equation*}

\par\medskip\noindent\textbf{Definition} (\texttt{is-linear-functional-on}).\;For all $m$, $s$ and $f$, 
\begin{equation*}
\begin{array}{@{}l@{}} \operatorname{is\text{-}linear\text{-}functional\text{-}on}(m, s, f) \\ \Leftrightarrow\; f \in (s \to \mathbb{R}) \wedge \\ \forall \operatorname{x\_}, \operatorname{y\_} \in s.\; (f(\operatorname{vadd}(m)(\operatorname{x\_}, \operatorname{y\_})) = +(f(\operatorname{x\_}), f(\operatorname{y\_}))) \wedge \\ \forall \operatorname{r\_} \in \mathbb{R}, \operatorname{x\_} \in s.\; (f(\operatorname{act}(m)(\operatorname{r\_}, \operatorname{x\_})) = \cdot(\operatorname{r\_}, f(\operatorname{x\_}))) \end{array}
\end{equation*}

\par\medskip\noindent\textbf{Definition} (\texttt{is-bounded-linear-functional-on}).\;For all $m$, $s$ and $f$, 
\begin{equation*}
\begin{array}{@{}l@{}} \operatorname{is\text{-}bounded\text{-}linear\text{-}functional\text{-}on}(m, s, f) \\ \Leftrightarrow\; \operatorname{is\text{-}linear\text{-}functional\text{-}on}(m, s, f) \wedge \\ \exists \operatorname{c\_} \in \mathbb{R}. \\ \quad (0 \leq \operatorname{c\_}) \wedge \forall \operatorname{x\_} \in s.\; (\lvert f(\operatorname{x\_}) \rvert \leq \cdot(\operatorname{c\_}, \operatorname{vnrm}(m)(\operatorname{x\_}))) \end{array}
\end{equation*}

\par\medskip\noindent\textbf{Definition} (\texttt{dual-norm-on}).\;$\operatorname{dual\text{-}norm\text{-}on}(m, s, f) :=$
\begin{equation*}
(\iota\, \operatorname{c\_}.\;((\operatorname{c\_} \in \mathbb{R}) \wedge ((0 \leq \operatorname{c\_}) \wedge (\forall \operatorname{x\_} \in s.\; (\lvert f(\operatorname{x\_}) \rvert \leq (\operatorname{c\_} \cdot \operatorname{vnrm}(m)(\operatorname{x\_}))) \wedge \forall \operatorname{d\_}.\; (((\operatorname{d\_} \in \mathbb{R}) \wedge ((0 \leq \operatorname{d\_}) \wedge \forall \operatorname{x\_} \in s.\; (\lvert f(\operatorname{x\_}) \rvert \leq (\operatorname{d\_} \cdot \operatorname{vnrm}(m)(\operatorname{x\_}))))) \Rightarrow (\operatorname{c\_} \leq \operatorname{d\_}))))))
\end{equation*}

\par\medskip\noindent\textbf{Definition} (\texttt{extends-on}).\;For all $s$, $g$ and $f$,  $\operatorname{extends\text{-}on}(s, g, f) \Leftrightarrow \forall \operatorname{x\_} \in s.\; (g(\operatorname{x\_}) = f(\operatorname{x\_}))$
\par\medskip\noindent\textbf{Definition} (\texttt{npe}).\;For all $m$, $s$, $f$, $t$ and $g$, 
\begin{equation*}
\begin{array}{@{}l@{}} \operatorname{npe}(m, s, f, t, g) \\ \Leftrightarrow\; \operatorname{is\text{-}submodule}(m, t) \wedge \\ s \subseteq t \wedge \\ \operatorname{is\text{-}linear\text{-}functional\text{-}on}(m, t, g) \wedge \\ \operatorname{extends\text{-}on}(s, g, f) \wedge \\ \forall \operatorname{w\_} \in t.\; (\lvert g(\operatorname{w\_}) \rvert \leq \cdot(\operatorname{dual\text{-}norm\text{-}on}(m, s, f), \operatorname{vnrm}(m)(\operatorname{w\_}))) \end{array}
\end{equation*}

\par\medskip\noindent\textbf{Definition} (\texttt{good-sub}).\;For all $m$, $s$, $f$ and $t$,  $\operatorname{good\text{-}sub}(m, s, f, t) \Leftrightarrow \exists \operatorname{g\_}.\; \operatorname{npe}(m, s, f, t, \operatorname{g\_})$
\par\medskip\noindent\textbf{Definition} (\texttt{line}).\;$\operatorname{line}(m, v) :=$ $\{\operatorname{y\_} \in \operatorname{vec}(m) : \exists \operatorname{r\_} \in \mathbb{R}.\; (\operatorname{y\_} = \operatorname{act}(m)(\operatorname{r\_}, v))\}$

\vnbbookfile{structure-library/dual-space}{structure-library/dual-space}
\noindent{\small dual-space.scm -- THE DUAL of a real normed vector space, as an OBJECT.}\par
\par\medskip\noindent\textbf{Definition} (\texttt{dual}).\;${m}^{*} :=$
\begin{equation*}
\langle \operatorname{rr\text{-}normed\text{-}field}, \operatorname{dual\text{-}vec}(m), (\lambda \langle \operatorname{f\_}, \operatorname{g\_} \rangle \in (\operatorname{dual\text{-}vec}(m) \times \operatorname{dual\text{-}vec}(m)).\; (\lambda \operatorname{x\_} \in \operatorname{vec}(m).\; (\operatorname{f\_}(\operatorname{x\_}) + \operatorname{g\_}(\operatorname{x\_})))), (\lambda \operatorname{x\_} \in \operatorname{vec}(m).\; 0), (\lambda \operatorname{f\_} \in \operatorname{dual\text{-}vec}(m).\; (\lambda \operatorname{x\_} \in \operatorname{vec}(m).\; -\operatorname{f\_}(\operatorname{x\_}))), (\lambda \langle \operatorname{r\_}, \operatorname{f\_} \rangle \in (\mathbb{R} \times \operatorname{dual\text{-}vec}(m)).\; (\lambda \operatorname{x\_} \in \operatorname{vec}(m).\; (\operatorname{r\_} \cdot \operatorname{f\_}(\operatorname{x\_})))), (\lambda \operatorname{f\_} \in \operatorname{dual\text{-}vec}(m).\; \operatorname{dual\text{-}norm}(m, \operatorname{f\_})) \rangle
\end{equation*}

\par\medskip\noindent\textbf{Definition} (\texttt{dual-vec}).\;$\operatorname{dual\text{-}vec}(m) :=$ $\{\operatorname{f\_} \in (\operatorname{vec}(m) \to \mathbb{R}) : \operatorname{is\text{-}bounded\text{-}linear\text{-}functional}(m, \operatorname{f\_})\}$
\vnbdef{dual-vec-membership}{\vnbstmt{\begin{array}{@{}l@{}} \forall m, f. \\ \quad (f \in \operatorname{dual\text{-}vec}(m)) \Leftrightarrow ((f \in (\operatorname{vec}(m) \to \mathbb{R})) \wedge \operatorname{is\text{-}bounded\text{-}linear\text{-}functional}(m, f)) \end{array}}}
\vnbdef{dual-scal}{\vnbstmt{\forall m.\; (\operatorname{scal}({m}^{*}) \simeq \operatorname{rr\text{-}normed\text{-}field})}}
\vnbdef{dual-carrier}{\vnbstmt{\forall m.\; (\operatorname{vec}({m}^{*}) \simeq \operatorname{dual\text{-}vec}(m))}}
\vnbdef{dual-norm-value}{\vnbstmt{\forall m.; \forall f \in \operatorname{dual\text{-}vec}(m).\; (\operatorname{vnrm}({m}^{*})(f) \simeq \operatorname{dual\text{-}norm}(m, f))}}

\vnbbookfile{structure-library/metric-subspace}{structure-library/metric-subspace}
\noindent{\small metric-subspace.scm -- the METRIC SUBSPACE of a metric space, and the RESTRICTION of a function to a subset.  DEFINITIONS ONLY; every law is PROVEN in theorem-library/metric-subspace-laws.scm.}\par
\par\medskip\noindent\textbf{Definition} (\texttt{subspace-ms}).\;$\operatorname{subspace\text{-}ms}(s, a) :=$
\begin{equation*}
\langle a, (\lambda \langle \operatorname{sbu\_}, \operatorname{sbv\_} \rangle \in (a \times a).\; \operatorname{dist}(s)(\operatorname{sbu\_}, \operatorname{sbv\_})) \rangle
\end{equation*}

\par\medskip\noindent\textbf{Definition} (\texttt{restrict}).\;$\operatorname{restrict}(f, a) :=$ $(\lambda \operatorname{rsx\_} \in a.\; f(\operatorname{rsx\_}))$

\vnbbookfile{structure-library/diff-on}{structure-library/diff-on}
\noindent{\small diff-on.scm -- THE DERIVATIVE ON AN OPEN SET, OVER A NORMED FIELD. DEFINITIONS ONLY; every law is PROVEN in theorem-library/diff-on-laws.scm.}\par
\par\medskip\noindent\textbf{Definition} (\texttt{is-diff-on}).\;For all $k$, $u$, $f$, $a$ and $l$, 
\begin{equation*}
\begin{array}{@{}l@{}} \operatorname{is\text{-}diff\text{-}on}(k, u, f, a, l) \\ \Leftrightarrow\; \operatorname{is\text{-}normed\text{-}field}(k) \wedge \\ \operatorname{is\text{-}open}(\operatorname{nf\text{-}metric\text{-}space}(k), u) \wedge \\ f \in (u \to \operatorname{carr}(k)) \wedge \\ a \in u \wedge \\ l \in \operatorname{carr}(k) \wedge \\ \exists \varphi \in (u \to \operatorname{carr}(k)). \\ \quad \operatorname{is\text{-}continuous\text{-}at}(\operatorname{subspace\text{-}ms}(\operatorname{nf\text{-}metric\text{-}space}(k), u), \operatorname{nf\text{-}metric\text{-}space}(k), \varphi, a) \wedge \\ \quad \varphi(a) = l \wedge \\ \quad \forall x \in u. \\ \quad \quad \operatorname{add}(k)(f(x), \operatorname{neg}(k)(f(a))) = \operatorname{mul}(k)(\varphi(x), \operatorname{add}(k)(x, \operatorname{neg}(k)(a))) \end{array}
\end{equation*}

\par\medskip\noindent\textbf{Definition} (\texttt{deriv-on}).\;$\operatorname{deriv\text{-}on}(k, u, f, a) :=$ $(\iota\, l.\;\operatorname{is\text{-}diff\text{-}on}(k, u, f, a, l))$
\par\medskip\noindent\textbf{Definition} (\texttt{holomorphic-on}).\;For all $u$ and $f$, 
\begin{equation*}
\begin{array}{@{}l@{}} \operatorname{holomorphic\text{-}on}(u, f) \\ \Leftrightarrow\; \operatorname{is\text{-}open}(\operatorname{nf\text{-}metric\text{-}space}(\operatorname{cc\text{-}normed\text{-}field}), u) \wedge \\ f \in (u \to \mathbb{C}) \wedge \\ \forall a \in u.\; \exists l.\; \operatorname{is\text{-}diff\text{-}on}(\operatorname{cc\text{-}normed\text{-}field}, u, f, a, l) \end{array}
\end{equation*}


\vnbbookfile{structure-library/extreme-value}{structure-library/extreme-value}
\noindent{\small MOVED 2026-09-20 (batch 12-A) from theorem-library/, and MOVED UP in load.scm to the calculus-vocabulary block just after numeric-instances: this file is VOCABULARY -- one def-functoid, not one proof -- while 43 files are STATED with CCINT and every one of them had to load below it.  CCINT is a SEP over RR and has no ...}\par
\par\medskip\noindent\textbf{Definition} (\texttt{ccint}).\;$\operatorname{ccint}(a, b) :=$ $\{x \in \mathbb{R} : ((a \leq x) \wedge (x \leq b))\}$

\vnbbookfile{structure-library/derivative}{structure-library/derivative}
\noindent{\small derivative.scm -- the Caratheodory derivative, as VOCABULARY.}\par
\par\medskip\noindent\textbf{Definition} (\texttt{is-diff-at}).\;For all $f$, $a$ and $l$, 
\begin{equation*}
\begin{array}{@{}l@{}} \operatorname{is\text{-}diff\text{-}at}(f, a, l) \\ \Leftrightarrow\; f \in (\mathbb{R} \to \mathbb{R}) \wedge \\ a \in \mathbb{R} \wedge \\ l \in \mathbb{R} \wedge \\ \exists \varphi \in (\mathbb{R} \to \mathbb{R}). \\ \quad \operatorname{is\text{-}continuous\text{-}at}(\operatorname{rr\text{-}ms}, \operatorname{rr\text{-}ms}, \varphi, a) \wedge \\ \quad \varphi(a) = l \wedge \\ \quad \forall x \in \mathbb{R}.\; (-(f(x), f(a)) = \cdot(\varphi(x), -(x, a))) \end{array}
\end{equation*}

\par\medskip\noindent\textbf{Definition} (\texttt{deriv}).\;$\operatorname{deriv}(f, a) :=$ $(\iota\, l.\;\operatorname{is\text{-}diff\text{-}at}(f, a, l))$

\vnbbookfile{structure-library/cc-coords}{structure-library/cc-coords}
\noindent{\small cc-coords.scm -- CC and cartesian(RR, RR), as VOCABULARY.}\par
\par\medskip\noindent\textbf{Definition} (\texttt{cc-coords}).\;
\begin{equation*}
(\operatorname{cc\text{-}coords} \simeq (\lambda \operatorname{zv\_} \in \mathbb{C}.\; \langle \operatorname{real\text{-}part}(\operatorname{zv\_}), \operatorname{imag\text{-}part}(\operatorname{zv\_}) \rangle))
\end{equation*}

\par\medskip\noindent\textbf{Definition} (\texttt{cc-of-pair}).\;
\begin{equation*}
(\operatorname{cc\text{-}of\text{-}pair} \simeq (\lambda \operatorname{pv\_} \in (\mathbb{R} \times \mathbb{R}).\; +(\operatorname{nth}(1, \operatorname{pv\_}), \cdot(\operatorname{nth}(2, \operatorname{pv\_}), +i))))
\end{equation*}

\par\medskip\noindent\textbf{Definition} (\texttt{is-cc-diff-at}).\;For all $g$, $a$ and $l$, 
\begin{equation*}
\begin{array}{@{}l@{}} \operatorname{is\text{-}cc\text{-}diff\text{-}at}(g, a, l) \\ \Leftrightarrow\; g \in (\mathbb{R} \to \mathbb{C}) \wedge \\ a \in \mathbb{R} \wedge \\ l \in \mathbb{C} \wedge \\ \operatorname{is\text{-}diff\text{-}at}((\lambda \operatorname{sv\_} \in \mathbb{R}.\; \operatorname{real\text{-}part}(g(\operatorname{sv\_}))), a, \operatorname{real\text{-}part}(l)) \wedge \\ \operatorname{is\text{-}diff\text{-}at}((\lambda \operatorname{sv\_} \in \mathbb{R}.\; \operatorname{imag\text{-}part}(g(\operatorname{sv\_}))), a, \operatorname{imag\text{-}part}(l)) \end{array}
\end{equation*}


\vnbbookfile{structure-library/antiderivative}{structure-library/antiderivative}
\noindent{\small antiderivative.scm -- Definition 4.6 (calculus.pdf), as VOCABULARY.}\par
\par\medskip\noindent\textbf{Definition} (\texttt{is-antiderivative}).\;For all $f$, $\varphi$, $a$ and $b$, 
\begin{equation*}
\begin{array}{@{}l@{}} \operatorname{is\text{-}antiderivative}(f, \varphi, a, b) \\ \Leftrightarrow\; f \in (\mathbb{R} \to \mathbb{R}) \wedge \\ \varphi \in (\mathbb{R} \to \mathbb{R}) \wedge \\ a \in \mathbb{R} \wedge \\ b \in \mathbb{R} \wedge \\ a < b \wedge \\ \forall \operatorname{x\_} \in \operatorname{ccint}(a, b).\; \operatorname{is\text{-}continuous\text{-}at}(\operatorname{rr\text{-}ms}, \operatorname{rr\text{-}ms}, f, \operatorname{x\_}) \wedge \\ \forall \operatorname{th\_}. \\ \quad ((\operatorname{th\_} \in \mathbb{R}) \wedge ((a < \operatorname{th\_}) \wedge (\operatorname{th\_} < b))) \Rightarrow \operatorname{is\text{-}diff\text{-}at}(f, \operatorname{th\_}, \varphi(\operatorname{th\_})) \end{array}
\end{equation*}

\par\medskip\noindent\textbf{Definition} (\texttt{is-antiderivable}).\;For all $\varphi$, $a$ and $b$, 
\begin{equation*}
\operatorname{is\text{-}antiderivable}(\varphi, a, b) \Leftrightarrow \exists \operatorname{f\_}.\; \operatorname{is\text{-}antiderivative}(\operatorname{f\_}, \varphi, a, b)
\end{equation*}


\vnbbookfile{structure-library/complex-inner-product}{structure-library/complex-inner-product}
\noindent{\small complex-inner-product.scm -- COMPLEX-INNER-PRODUCT-SPACE: a vector space over CC carrying a sesquilinear, conjugate-symmetric, positive-definite form IP. Vocabulary only (no proofs); the two inequalities it exists for are proved in theorem-library/inner-product-inequalities.scm.}\par
\vnbstructdef{complex-inner-product-space}{A complex inner product space $s$ is a 7-tuple whose components are listed below.\par
\begin{enumerate}[nosep]
\item $\operatorname{scal}(s) = \operatorname{nth}(1, s)$, a substructure (\texttt{substructure}): $\operatorname{is\text{-}ring}(\operatorname{scal}(s))$.
\item $\operatorname{vec}(s) = \operatorname{nth}(2, s)$, a carrier (\texttt{carriers}): $\operatorname{vec}(s) \in \mathbf{Set}$.
\item $\operatorname{vadd}(s) = \operatorname{nth}(3, s)$, an operation (\texttt{op}): $\operatorname{vadd}(s) \in ((\operatorname{vec}(s) \times \operatorname{vec}(s)) \to \operatorname{vec}(s))$.
\item $\operatorname{vzero}(s) = \operatorname{nth}(4, s)$, a constant (\texttt{constant}): $\operatorname{vzero}(s) \in \operatorname{vec}(s)$.
\item $\operatorname{vneg}(s) = \operatorname{nth}(5, s)$, an operation (\texttt{op}): $\operatorname{vneg}(s) \in (\operatorname{vec}(s) \to \operatorname{vec}(s))$.
\item $\operatorname{act}(s) = \operatorname{nth}(6, s)$, an operation (\texttt{op}): $\operatorname{act}(s) \in ((\operatorname{carr}(\operatorname{scal}(s)) \times \operatorname{vec}(s)) \to \operatorname{vec}(s))$.
\item $\langle\cdot,\cdot\rangle_{s} = \operatorname{nth}(7, s)$, an operation (\texttt{op}): $\langle\cdot,\cdot\rangle_{s} \in ((\operatorname{vec}(s) \times \operatorname{vec}(s)) \to \mathbb{C})$.
\end{enumerate}
The laws (the \texttt{property} and \texttt{law} clauses):
\begin{itemize}[nosep]
\item $\operatorname{is\text{-}associative}(\operatorname{vadd}(s), \operatorname{vec}(s))$
\item $\operatorname{is\text{-}commutative}(\operatorname{vadd}(s), \operatorname{vec}(s))$
\item $\operatorname{is\text{-}identity}(\operatorname{vadd}(s), \operatorname{vzero}(s), \operatorname{vec}(s))$
\item $\operatorname{has\text{-}inverses}(\operatorname{vadd}(s), \operatorname{vzero}(s), \operatorname{vneg}(s), \operatorname{vec}(s))$
\item $\operatorname{scal}(s) = \operatorname{normed\text{-}field\text{-}as\text{-}commutative\text{-}ring}(\operatorname{cc\text{-}normed\text{-}field})$
\item $\operatorname{carr}(\operatorname{scal}(s)) = \mathbb{C}$
\item
\begin{equation*}
\begin{array}{@{}l@{}} \forall \operatorname{r\_} \in \operatorname{carr}(\operatorname{scal}(s)), \operatorname{x\_}, \operatorname{y\_} \in \operatorname{vec}(s). \\ \quad \operatorname{act}(s)(\operatorname{r\_}, \operatorname{vadd}(s)(\operatorname{x\_}, \operatorname{y\_})) = \operatorname{vadd}(s)(\operatorname{act}(s)(\operatorname{r\_}, \operatorname{x\_}), \operatorname{act}(s)(\operatorname{r\_}, \operatorname{y\_})) \end{array}
\end{equation*}

\item
\begin{equation*}
\begin{array}{@{}l@{}} \forall \operatorname{r\_}, \operatorname{s\_} \in \operatorname{carr}(\operatorname{scal}(s)), \operatorname{x\_} \in \operatorname{vec}(s). \\ \quad \operatorname{act}(s)(\operatorname{add}(\operatorname{scal}(s))(\operatorname{r\_}, \operatorname{s\_}), \operatorname{x\_}) = \operatorname{vadd}(s)(\operatorname{act}(s)(\operatorname{r\_}, \operatorname{x\_}), \operatorname{act}(s)(\operatorname{s\_}, \operatorname{x\_})) \end{array}
\end{equation*}

\item
\begin{equation*}
\begin{array}{@{}l@{}} \forall \operatorname{r\_}, \operatorname{s\_} \in \operatorname{carr}(\operatorname{scal}(s)), \operatorname{x\_} \in \operatorname{vec}(s). \\ \quad \operatorname{act}(s)(\operatorname{mul}(\operatorname{scal}(s))(\operatorname{r\_}, \operatorname{s\_}), \operatorname{x\_}) = \operatorname{act}(s)(\operatorname{r\_}, \operatorname{act}(s)(\operatorname{s\_}, \operatorname{x\_})) \end{array}
\end{equation*}

\item $\forall \operatorname{x\_} \in \operatorname{vec}(s).\; (\operatorname{act}(s)(\operatorname{one}(\operatorname{scal}(s)), \operatorname{x\_}) = \operatorname{x\_})$
\item
\begin{equation*}
\begin{array}{@{}l@{}} \forall \operatorname{x\_}, \operatorname{y\_}, \operatorname{z\_} \in \operatorname{vec}(s). \\ \quad \langle\cdot,\cdot\rangle_{s}(\operatorname{vadd}(s)(\operatorname{x\_}, \operatorname{y\_}), \operatorname{z\_}) = (\langle\cdot,\cdot\rangle_{s}(\operatorname{x\_}, \operatorname{z\_}) + \langle\cdot,\cdot\rangle_{s}(\operatorname{y\_}, \operatorname{z\_})) \end{array}
\end{equation*}

\item
\begin{equation*}
\begin{array}{@{}l@{}} \forall \operatorname{a\_} \in \mathbb{C}, \operatorname{x\_}, \operatorname{y\_} \in \operatorname{vec}(s). \\ \quad \langle\cdot,\cdot\rangle_{s}(\operatorname{act}(s)(\operatorname{a\_}, \operatorname{x\_}), \operatorname{y\_}) = (\operatorname{a\_} \cdot \langle\cdot,\cdot\rangle_{s}(\operatorname{x\_}, \operatorname{y\_})) \end{array}
\end{equation*}

\item $\forall \operatorname{x\_}, \operatorname{y\_} \in \operatorname{vec}(s).\; (\langle\cdot,\cdot\rangle_{s}(\operatorname{y\_}, \operatorname{x\_}) = \operatorname{conjugate}(\langle\cdot,\cdot\rangle_{s}(\operatorname{x\_}, \operatorname{y\_})))$
\item $\forall \operatorname{x\_} \in \operatorname{vec}(s).\; (\langle\cdot,\cdot\rangle_{s}(\operatorname{x\_}, \operatorname{x\_}) \in \mathbb{R})$
\item $\forall \operatorname{x\_} \in \operatorname{vec}(s).\; (0 \leq \langle\cdot,\cdot\rangle_{s}(\operatorname{x\_}, \operatorname{x\_}))$
\item $\forall \operatorname{x\_} \in \operatorname{vec}(s).\; ((\langle\cdot,\cdot\rangle_{s}(\operatorname{x\_}, \operatorname{x\_}) = 0) \Leftrightarrow (\operatorname{x\_} = \operatorname{vzero}(s)))$
\end{itemize}
The characteristic law (the axiom \texttt{is-complex-inner-product-space}): $\operatorname{is\text{-}complex\text{-}inner\text{-}product\text{-}space}(s)$ holds if and only if $\operatorname{length}(s) = 7$, every component has the typing listed with it, and every law holds.
}
\par\medskip\noindent\textbf{Definition} (\texttt{ip-normed-ag}).\;$\operatorname{ip\text{-}normed\text{-}ag}(v) :=$
\begin{equation*}
\langle \operatorname{vec}(v), \operatorname{vadd}(v), \operatorname{vzero}(v), \operatorname{vneg}(v), (\lambda \operatorname{x\_} \in \operatorname{vec}(v).\; \lVert \operatorname{x\_} \rVert) \rangle
\end{equation*}

\par\medskip\noindent\textbf{Definition} (\texttt{ip-norm}).\;$\lVert x \rVert :=$ $\sqrt{\langle\cdot,\cdot\rangle_{v}(x, x)}$
\par\medskip\noindent\textbf{Definition} (\texttt{hom-complex-inner-product-space}).\;$\operatorname{hom\text{-}complex\text{-}inner\text{-}product\text{-}space}(a, b) :=$
\begin{equation*}
\{\operatorname{homf\_} \in (\operatorname{vec}(a) \to \operatorname{vec}(b)) : \operatorname{is\text{-}hom\text{-}complex\text{-}inner\text{-}product\text{-}space}(a, b, \operatorname{homf\_})\}
\end{equation*}

\vnbdef{cips-act-type}{\vnbstmt{\begin{array}{@{}l@{}} \forall v.; \forall a \in \mathbb{C}, x \in \operatorname{vec}(v). \\ \quad \operatorname{is\text{-}complex\text{-}inner\text{-}product\text{-}space}(v) \Rightarrow (\operatorname{act}(v)(a, x) \in \operatorname{vec}(v)) \end{array}}}
\vnbdef{cips-scal-carrier}{\vnbstmt{\forall v.\; (\operatorname{is\text{-}complex\text{-}inner\text{-}product\text{-}space}(v) \Rightarrow (\operatorname{carr}(\operatorname{scal}(v)) = \mathbb{C}))}}
\vnbdef{cips-vzero-in}{\vnbstmt{\forall v.\; (\operatorname{is\text{-}complex\text{-}inner\text{-}product\text{-}space}(v) \Rightarrow (\operatorname{vzero}(v) \in \operatorname{vec}(v)))}}
\vnbdef{cips-vadd-vzero}{\vnbstmt{\begin{array}{@{}l@{}} \forall v.; \forall x \in \operatorname{vec}(v). \\ \quad \operatorname{is\text{-}complex\text{-}inner\text{-}product\text{-}space}(v) \Rightarrow (\operatorname{vadd}(v)(\operatorname{vzero}(v), x) = x) \end{array}}}
\vnbdef{cips-vadd-type}{\vnbstmt{\begin{array}{@{}l@{}} \forall v.; \forall x, y \in \operatorname{vec}(v). \\ \quad \operatorname{is\text{-}complex\text{-}inner\text{-}product\text{-}space}(v) \Rightarrow (\operatorname{vadd}(v)(x, y) \in \operatorname{vec}(v)) \end{array}}}
\vnbdef{cips-ip-type}{\vnbstmt{\begin{array}{@{}l@{}} \forall v.; \forall x, y \in \operatorname{vec}(v). \\ \quad \operatorname{is\text{-}complex\text{-}inner\text{-}product\text{-}space}(v) \Rightarrow (\langle\cdot,\cdot\rangle_{v}(x, y) \in \mathbb{C}) \end{array}}}
\vnbdef{cips-ip-add-left}{\vnbstmt{\begin{array}{@{}l@{}} \forall v.; \forall x, y, z \in \operatorname{vec}(v). \\ \quad \operatorname{is\text{-}complex\text{-}inner\text{-}product\text{-}space}(v) \Rightarrow  \\ \quad \quad \langle\cdot,\cdot\rangle_{v}(\operatorname{vadd}(v)(x, y), z) = (\langle\cdot,\cdot\rangle_{v}(x, z) + \langle\cdot,\cdot\rangle_{v}(y, z)) \end{array}}}
\vnbdef{cips-ip-homog-left}{\vnbstmt{\begin{array}{@{}l@{}} \forall v.; \forall a \in \mathbb{C}, x, y \in \operatorname{vec}(v). \\ \quad \operatorname{is\text{-}complex\text{-}inner\text{-}product\text{-}space}(v) \Rightarrow  \\ \quad \quad \langle\cdot,\cdot\rangle_{v}(\operatorname{act}(v)(a, x), y) = (a \cdot \langle\cdot,\cdot\rangle_{v}(x, y)) \end{array}}}
\vnbdef{cips-ip-conj-sym}{\vnbstmt{\begin{array}{@{}l@{}} \forall v.; \forall x, y \in \operatorname{vec}(v). \\ \quad \operatorname{is\text{-}complex\text{-}inner\text{-}product\text{-}space}(v) \Rightarrow  \\ \quad \quad \langle\cdot,\cdot\rangle_{v}(y, x) = \operatorname{conjugate}(\langle\cdot,\cdot\rangle_{v}(x, y)) \end{array}}}
\vnbdef{cips-ip-self-real}{\vnbstmt{\begin{array}{@{}l@{}} \forall v.; \forall x \in \operatorname{vec}(v). \\ \quad \operatorname{is\text{-}complex\text{-}inner\text{-}product\text{-}space}(v) \Rightarrow (\langle\cdot,\cdot\rangle_{v}(x, x) \in \mathbb{R}) \end{array}}}
\vnbdef{cips-ip-self-nonneg}{\vnbstmt{\begin{array}{@{}l@{}} \forall v.; \forall x \in \operatorname{vec}(v). \\ \quad \operatorname{is\text{-}complex\text{-}inner\text{-}product\text{-}space}(v) \Rightarrow (0 \leq \langle\cdot,\cdot\rangle_{v}(x, x)) \end{array}}}
\vnbdef{cips-ip-zero-vector}{\vnbstmt{\begin{array}{@{}l@{}} \forall v.; \forall x \in \operatorname{vec}(v). \\ \quad (\operatorname{is\text{-}complex\text{-}inner\text{-}product\text{-}space}(v) \wedge (\langle\cdot,\cdot\rangle_{v}(x, x) = 0)) \Rightarrow (x = \operatorname{vzero}(v)) \end{array}}}
\vnbdef{cips-ip-definite}{\vnbstmt{\begin{array}{@{}l@{}} \forall v.; \forall x \in \operatorname{vec}(v). \\ \quad \operatorname{is\text{-}complex\text{-}inner\text{-}product\text{-}space}(v) \Rightarrow ((\langle\cdot,\cdot\rangle_{v}(x, x) = 0) \Leftrightarrow (x = \operatorname{vzero}(v))) \end{array}}}
\vnbdef{ip-normed-ag-carrier}{\vnbstmt{\begin{array}{@{}l@{}} \forall v. \\ \quad \operatorname{is\text{-}complex\text{-}inner\text{-}product\text{-}space}(v) \Rightarrow (\operatorname{carr}(\operatorname{ip\text{-}normed\text{-}ag}(v)) \simeq \operatorname{vec}(v)) \end{array}}}
\vnbdef{ip-normed-ag-norm-value}{\vnbstmt{\forall v.; \forall x \in \operatorname{vec}(v).\; (\operatorname{nrm}(\operatorname{ip\text{-}normed\text{-}ag}(v))(x) \simeq \lVert x \rVert)}}
\vnbdef{is-hom-complex-inner-product-space-def}{\vnbstmt{\begin{array}{@{}l@{}} \forall a, b, f. \\ \quad \operatorname{is\text{-}hom\text{-}complex\text{-}inner\text{-}product\text{-}space}(a, b, f) \\ \quad \Leftrightarrow\; \operatorname{is\text{-}complex\text{-}inner\text{-}product\text{-}space}(a) \wedge \\ \quad \operatorname{is\text{-}complex\text{-}inner\text{-}product\text{-}space}(b) \wedge \\ \quad f \in (\operatorname{vec}(a) \to \operatorname{vec}(b)) \wedge \\ \quad \operatorname{scal}(a) = \operatorname{scal}(b) \wedge \\ \quad \forall x_{1}, x_{2} \in \operatorname{vec}(a). \\ \quad \quad f(\operatorname{vadd}(a)(x_{1}, x_{2})) = \operatorname{vadd}(b)(f(x_{1}), f(x_{2})) \wedge \\ \quad f(\operatorname{vzero}(a)) = \operatorname{vzero}(b) \wedge \\ \quad \forall x_{1} \in \operatorname{vec}(a).\; (f(\operatorname{vneg}(a)(x_{1})) = \operatorname{vneg}(b)(f(x_{1}))) \wedge \\ \quad \forall x_{1} \in \operatorname{carr}(\operatorname{scal}(a)), x_{2} \in \operatorname{vec}(a). \\ \quad \quad f(\operatorname{act}(a)(x_{1}, x_{2})) = \operatorname{act}(b)(x_{1}, f(x_{2})) \wedge \\ \quad \forall x_{1}, x_{2} \in \operatorname{vec}(a).\; (\langle\cdot,\cdot\rangle_{a}(x_{1}, x_{2}) = \langle\cdot,\cdot\rangle_{b}(f(x_{1}), f(x_{2}))) \end{array}}}
\vnbdef{complex-inner-product-space-as-module-is-module}{\vnbstmt{\begin{array}{@{}l@{}} \forall r. \\ \quad \operatorname{is\text{-}complex\text{-}inner\text{-}product\text{-}space}(r) \Rightarrow  \\ \quad \quad \operatorname{is\text{-}module}(\operatorname{complex\text{-}inner\text{-}product\text{-}space\text{-}as\text{-}module}(r)) \end{array}}}
\vnbdef{complex-inner-product-space-class}{\vnbstmt{\begin{array}{@{}l@{}} \forall s. \\ \quad (s \in \operatorname{complex\text{-}inner\text{-}product\text{-}space}) \Leftrightarrow \operatorname{is\text{-}complex\text{-}inner\text{-}product\text{-}space}(s) \end{array}}}

\vnbbookfile{structure-library/bounded-metric}{structure-library/bounded-metric}
\noindent{\small bounded-metric.scm -- the bounded metric  rho = d/(1+d)  of a metric space, and the fact that it is topologically equivalent to the original (the identity map is bicontinuous).  Library phase: the construction is a def-functoid, the metric/equivalence facts are warranted supports.}\par
\par\medskip\noindent\textbf{Definition} (\texttt{bdd-metric}).\;$\operatorname{bdd\text{-}metric}(s) :=$
\begin{equation*}
\langle \operatorname{pts}(s), (\lambda \langle u, v \rangle \in (\operatorname{pts}(s) \times \operatorname{pts}(s)).\; /(\operatorname{dist}(s)(u, v), (1 + \operatorname{dist}(s)(u, v)))) \rangle
\end{equation*}

\par\medskip\noindent\textbf{Definition} (\texttt{rr-bounded-ms}).\; $(\operatorname{rr\text{-}bounded\text{-}ms} \simeq \operatorname{bdd\text{-}metric}(\operatorname{rr\text{-}ms}))$
\par\medskip\noindent\textbf{Definition} (\texttt{is-bounded-metric-space}).\;For every $s$, 
\begin{equation*}
\begin{array}{@{}l@{}} \operatorname{is\text{-}bounded\text{-}metric\text{-}space}(s) \\ \Leftrightarrow\; \operatorname{is\text{-}metric\text{-}space}(s) \wedge \\ \exists b \in \mathbb{R}. \forall x, y \in \operatorname{pts}(s). \\ \quad \operatorname{dist}(s)(x, y) \leq b \end{array}
\end{equation*}


\vnbbookfile{structure-library/trunc-metric}{structure-library/trunc-metric}
\noindent{\small trunc-metric.scm -- the TRUNCATED METRIC  d'(x,y) = min(1, d(x,y))  of a metric space, as an ENDOFUNCTOR of METRIC-SPACE.}\par
\par\medskip\noindent\textbf{Definition} (\texttt{trunc-metric}).\;$\operatorname{trunc\text{-}metric}(s) :=$
\begin{equation*}
\langle \operatorname{pts}(s), (\lambda \langle u, v \rangle \in (\operatorname{pts}(s) \times \operatorname{pts}(s)).\; \operatorname{min}(1, \operatorname{dist}(s)(u, v))) \rangle
\end{equation*}


\vnbbookfile{structure-library/pseudometric}{structure-library/pseudometric}
\noindent{\small pseudometric.scm -- PSEUDOMETRIC-SPACE, the topology GENERATED by a countable family of pseudometrics (the gauge topology), and the metrizability cluster T2/T3/T4 stated as warranted PSS supports.}\par
\vnbstructdef{pseudometric-space}{A pseudometric space $s$ is a 2-tuple whose components are listed below.\par
\begin{enumerate}[nosep]
\item $\operatorname{pts}(s) = \operatorname{nth}(1, s)$, a carrier (\texttt{carriers}): $\operatorname{pts}(s) \in \mathbf{Set}$.
\item $\operatorname{dist}(s) = \operatorname{nth}(2, s)$, an operation (\texttt{op}): $\operatorname{dist}(s) \in ((\operatorname{pts}(s) \times \operatorname{pts}(s)) \to \mathbb{R})$.
\end{enumerate}
The laws (the \texttt{property} and \texttt{law} clauses):
\begin{itemize}[nosep]
\item $\operatorname{is\text{-}pseudometric}(\operatorname{dist}(s), \operatorname{pts}(s))$
\end{itemize}
The characteristic law (the axiom \texttt{is-pseudometric-space}): $\operatorname{is\text{-}pseudometric\text{-}space}(s)$ holds if and only if $\operatorname{length}(s) = 2$, every component has the typing listed with it, and every law holds.
}
\par\medskip\noindent\textbf{Definition} (\texttt{pseudo-gauge-top}).\;$\operatorname{pseudo\text{-}gauge\text{-}top}(\mathit{fam}, \mathit{ground}) :=$
\begin{equation*}
\langle \mathit{ground}, \{u \in \operatorname{power}(\mathit{ground}) : \operatorname{is\text{-}open\text{-}gauge}(\mathit{fam}, \mathit{ground}, u)\} \rangle
\end{equation*}

\par\medskip\noindent\textbf{Definition} (\texttt{is-countable-pseudometric-family}).\;For all $\mathit{fam}$ and $\mathit{ground}$, 
\begin{equation*}
\begin{array}{@{}l@{}} \operatorname{is\text{-}countable\text{-}pseudometric\text{-}family}(\mathit{fam}, \mathit{ground}) \\ \Leftrightarrow\; \forall \operatorname{n\_} \in \mathbb{N}.\; (\operatorname{is\text{-}pseudometric\text{-}space}(\operatorname{fam}(\operatorname{n\_})) \wedge (\operatorname{pts}(\operatorname{fam}(\operatorname{n\_})) \simeq \mathit{ground})) \end{array}
\end{equation*}

\par\medskip\noindent\textbf{Definition} (\texttt{is-open-gauge}).\;For all $\mathit{fam}$, $\mathit{ground}$ and $u$, 
\begin{equation*}
\begin{array}{@{}l@{}} \operatorname{is\text{-}open\text{-}gauge}(\mathit{fam}, \mathit{ground}, u) \\ \Leftrightarrow\; u \subseteq \mathit{ground} \wedge \\ \forall y \in u. \exists \operatorname{k\_} \in \mathbb{N}, r. \\ \quad \operatorname{pos\text{-}rr}(r) \wedge \\ \quad \forall x \in \mathit{ground}. \\ \quad \quad \forall \operatorname{n\_} \in \mathbb{N}.\; ((\operatorname{n\_} \leq \operatorname{k\_}) \Rightarrow (\operatorname{dist}(\operatorname{fam}(\operatorname{n\_}))(y, x) < r)) \Rightarrow (x \in u) \end{array}
\end{equation*}

\par\medskip\noindent\textbf{Definition} (\texttt{is-gauge-countable}).\;For every $s$, 
\begin{equation*}
\begin{array}{@{}l@{}} \operatorname{is\text{-}gauge\text{-}countable}(s) \\ \Leftrightarrow\; \operatorname{is\text{-}top\text{-}space}(s) \wedge \\ \exists \mathit{fam}. \\ \quad (\operatorname{is\text{-}countable\text{-}pseudometric\text{-}family}(\mathit{fam}, \operatorname{pts}(s)) \wedge (\operatorname{pseudo\text{-}gauge\text{-}top}(\mathit{fam}, \operatorname{pts}(s)) \simeq s)) \end{array}
\end{equation*}

\par\medskip\noindent\textbf{Definition} (\texttt{hom-pseudometric-space}).\;$\operatorname{hom\text{-}pseudometric\text{-}space}(a, b) :=$
\begin{equation*}
\{\operatorname{homf\_} \in (\operatorname{pts}(a) \to \operatorname{pts}(b)) : \operatorname{is\text{-}hom\text{-}pseudometric\text{-}space}(a, b, \operatorname{homf\_})\}
\end{equation*}

\vnbdef{is-hom-pseudometric-space-def}{\vnbstmt{\begin{array}{@{}l@{}} \forall a, b, f. \\ \quad \operatorname{is\text{-}hom\text{-}pseudometric\text{-}space}(a, b, f) \\ \quad \Leftrightarrow\; \operatorname{is\text{-}pseudometric\text{-}space}(a) \wedge \\ \quad \operatorname{is\text{-}pseudometric\text{-}space}(b) \wedge \\ \quad f \in (\operatorname{pts}(a) \to \operatorname{pts}(b)) \wedge \\ \quad \forall x_{1}, x_{2} \in \operatorname{pts}(a). \\ \quad \quad \operatorname{dist}(a)(x_{1}, x_{2}) = \operatorname{dist}(b)(f(x_{1}), f(x_{2})) \end{array}}}
\vnbdef{pseudometric-space-class}{\vnbstmt{\forall s.\; ((s \in \operatorname{pseudometric\text{-}space}) \Leftrightarrow \operatorname{is\text{-}pseudometric\text{-}space}(s))}}

\vnbbookfile{structure-library/metric-completion}{structure-library/metric-completion}
\noindent{\small RETIRED 2026-09-17 (proven, rake batch O): cauchy-setoid-is-setoid; embed-isometry -- theorem-library/rake-setoid2.scm, modulo 0.  embed-isometry was MISSPELLED here: `(EMBED M u)' applies the one-parameter functoid EMBED to the PAIR [M,u] (apply-tupling-2), which is a function on PTS([M,u]), not a point of the ...}\par
\par\medskip\noindent\textbf{Definition} (\texttt{embed}).\;$\operatorname{embed}(m) :=$
\begin{equation*}
(\lambda u \in \operatorname{pts}(m).\; \operatorname{class}(\operatorname{cauchy\text{-}setoid}(m), \operatorname{embed\text{-}seq}(m, u)))
\end{equation*}

\par\medskip\noindent\textbf{Definition} (\texttt{completion}).\;$\operatorname{completion}(m) :=$ $\langle \operatorname{quotient}(\operatorname{cauchy\text{-}setoid}(m)), \operatorname{completion\text{-}dist}(m) \rangle$
\par\medskip\noindent\textbf{Definition} (\texttt{dist-seq}).\;$\operatorname{dist\text{-}seq}(m, f, g) :=$ $(\lambda n \in \mathbb{N}.\; \operatorname{dist}(m)(f(n), g(n)))$
\par\medskip\noindent\textbf{Definition} (\texttt{cauchy-setoid}).\;$\operatorname{cauchy\text{-}setoid}(m) :=$ $\langle \operatorname{cseq}(m), \operatorname{crel}(m) \rangle$
\par\medskip\noindent\textbf{Definition} (\texttt{cseq-equiv}).\;For all $m$, $f$ and $g$, 
\begin{equation*}
\operatorname{cseq\text{-}equiv}(m, f, g) \Leftrightarrow \operatorname{converges\text{-}to}(\operatorname{rr\text{-}ms}, \operatorname{dist\text{-}seq}(m, f, g), 0)
\end{equation*}

\par\medskip\noindent\textbf{Definition} (\texttt{cseq}).\;$\operatorname{cseq}(m) :=$ $\{f \in (\mathbb{N} \to \operatorname{pts}(m)) : \operatorname{is\text{-}cauchy\text{-}seq}(m, f)\}$
\par\medskip\noindent\textbf{Definition} (\texttt{crel}).\;$\operatorname{crel}(m) :=$
\begin{equation*}
\{p \in (\operatorname{cseq}(m) \times \operatorname{cseq}(m)) : \operatorname{cseq\text{-}equiv}(m, \operatorname{nth}(1, p), \operatorname{nth}(2, p))\}
\end{equation*}

\par\medskip\noindent\textbf{Definition} (\texttt{completion-dist}).\;$\operatorname{completion\text{-}dist}(m) :=$
\begin{equation*}
\begin{array}{@{}l@{}} (\lambda p \in (\operatorname{quotient}(\operatorname{cauchy\text{-}setoid}(m)) \times \operatorname{quotient}(\operatorname{cauchy\text{-}setoid}(m))).\; \\ \quad (\iota\, \mathit{dval}.\;\exists f, g \in \operatorname{cseq}(m).\; ((\operatorname{nth}(1, p) = \operatorname{class}(\operatorname{cauchy\text{-}setoid}(m), f)) \wedge ((\operatorname{nth}(2, p) = \operatorname{class}(\operatorname{cauchy\text{-}setoid}(m), g)) \wedge \operatorname{converges\text{-}to}(\operatorname{rr\text{-}ms}, \operatorname{dist\text{-}seq}(m, f, g), \mathit{dval}))))) \end{array}
\end{equation*}

\par\medskip\noindent\textbf{Definition} (\texttt{embed-seq}).\;$\operatorname{embed\text{-}seq}(m, u) :=$ $(\lambda n \in \mathbb{N}.\; u)$

\vnbbookfile{structure-library/extended-reals}{structure-library/extended-reals}
\noindent{\small extended-reals.scm -- the extended real line RR-STAR = [-inf, +inf]}\par
\vnbdef{rr-star-membership}{\vnbstmt{\forall x.\; ((x \in \operatorname{rr\text{-}star}) \Leftrightarrow ((x \in \mathbb{R}) \vee ((x = \operatorname{pos\text{-}inf}) \vee (x = \operatorname{neg\text{-}inf}))))}}
\vnbaxiom{pos-inf-above-reals}{\vnbstmt{\forall x \in \mathbb{R}.\; \neg (\operatorname{pos\text{-}inf} \leq x)}}
\vnbaxiom{neg-inf-lower-bound}{\vnbstmt{\forall x \in \operatorname{rr\text{-}star}.\; (\operatorname{neg\text{-}inf} \leq x)}}
\vnbaxiom{neg-inf-not-in-rr}{\vnbstmt{\neg (\operatorname{neg\text{-}inf} \in \mathbb{R})}}
\vnbaxiom{pos-inf-upper-bound}{\vnbstmt{\forall x \in \operatorname{rr\text{-}star}.\; (x \leq \operatorname{pos\text{-}inf})}}

\vnbbookfile{structure-library/extended-reals-pos}{structure-library/extended-reals-pos}
\noindent{\small extended-reals-pos.scm -- RR-POS-STAR = [0, +inf], the nonnegative extended reals}\par
\par\medskip\noindent\textbf{Definition} (\texttt{esup}).\;$\operatorname{esup}(\operatorname{s\_}) :=$
\begin{equation*}
(\iota\, \operatorname{bb\_}.\;((\operatorname{bb\_} \in \operatorname{rr\text{-}pos\text{-}star}) \wedge (\forall x \in \operatorname{s\_}.\; (x \leq \operatorname{bb\_}) \wedge \forall b.\; (((b \in \operatorname{rr\text{-}pos\text{-}star}) \wedge \forall x \in \operatorname{s\_}.\; (x \leq b)) \Rightarrow (\operatorname{bb\_} \leq b)))))
\end{equation*}

\par\medskip\noindent\textbf{Definition} (\texttt{rr-pos-star-add-monoid}).\; $\operatorname{rr\text{-}pos\text{-}star\text{-}add\text{-}monoid} = \langle \operatorname{rr\text{-}pos\text{-}star}, \mathit{eplus}, 0 \rangle$
\vnbdef{rr-pos-star-membership}{\vnbstmt{\forall x.\; ((x \in \operatorname{rr\text{-}pos\text{-}star}) \Leftrightarrow (((x \in \mathbb{R}) \wedge (0 \leq x)) \vee (x = \operatorname{pos\text{-}inf})))}}
\vnbdef{eplus-def}{\vnbstmt{(\mathit{eplus} \simeq (\lambda \langle \operatorname{x\_}, \operatorname{y\_} \rangle \in (\operatorname{rr\text{-}pos\text{-}star} \times \operatorname{rr\text{-}pos\text{-}star}).\; \left(\text{if }((\operatorname{x\_} = \operatorname{pos\text{-}inf}) \vee (\operatorname{y\_} = \operatorname{pos\text{-}inf}))\text{ then }\operatorname{pos\text{-}inf}\text{ else }(\operatorname{x\_} + \operatorname{y\_})\right)))}}
\vnbdef{rr-pos-star-add-monoid@iden}{\vnbstmt{(\operatorname{iden}(\operatorname{rr\text{-}pos\text{-}star\text{-}add\text{-}monoid}) \simeq 0)}}
\vnbdef{rr-pos-star-add-monoid@opr}{\vnbstmt{(\operatorname{opr}(\operatorname{rr\text{-}pos\text{-}star\text{-}add\text{-}monoid}) \simeq \mathit{eplus})}}
\vnbdef{rr-pos-star-add-monoid@carr}{\vnbstmt{(\operatorname{carr}(\operatorname{rr\text{-}pos\text{-}star\text{-}add\text{-}monoid}) \simeq \operatorname{rr\text{-}pos\text{-}star})}}

\vnbbookfile{structure-library/nth-deriv}{structure-library/nth-deriv}
\noindent{\small nth-deriv.scm -- THE ITERATED DERIVATIVE ON AN OPEN SET, OVER A NORMED FIELD. DEFINITION ONLY; the laws are PROVEN in theorem-library/nth-deriv-laws.scm. Batch 36, 2026-09-26.}\par
\par\medskip\noindent\textbf{Definition} (\texttt{nth-deriv-on}).\;For all $\operatorname{ndvk\_}$, $\operatorname{ndvu\_}$ and $\operatorname{ndvf\_}$,  $(\operatorname{ndvf\_}^{(0)} \simeq \operatorname{ndvf\_})$\par
Let $\operatorname{ndvk\_}$, $\operatorname{ndvu\_}$ and $\operatorname{ndvf\_}$ be arbitrary. Let $\operatorname{ndvn\_} \in \mathbb{N}$. Then
\begin{equation*}
(\operatorname{ndvf\_}^{(\operatorname{succ}(\operatorname{ndvn\_}))} \simeq (\lambda \operatorname{ndvz\_} \in \operatorname{ndvu\_}.\; \operatorname{deriv\text{-}on}(\operatorname{ndvk\_}, \operatorname{ndvu\_}, \operatorname{ndvf\_}^{(\operatorname{ndvn\_})}, \operatorname{ndvz\_})))
\end{equation*}


\vnbbookfile{structure-library/block-family-combinatorial}{structure-library/block-family-combinatorial}
\noindent{\small MOVED 2026-09-20 (batch 12-A) from theorem-library/: this file is VOCABULARY -- definitions, notation and warranted supports, not one proof -- and every theorem stated with it had to load below it.  Its load.scm slot is unchanged. structure-library/block-family-combinatorial.scm}\par
\par\medskip\noindent\textbf{Definition} (\texttt{is-finite-cover}).\;For all $c$ and $a$, 
\begin{equation*}
\operatorname{is\text{-}finite\text{-}cover}(c, a) \Leftrightarrow ((c \in \mathbf{Set}) \wedge ((|c| \in \mathbb{N}) \wedge (a \subseteq \operatorname{big\text{-}union}(u, c, u))))
\end{equation*}


\vnbbookfile{structure-library/subsequence-capture}{structure-library/subsequence-capture}
\noindent{\small MOVED 2026-09-20 (batch 12-A) from theorem-library/: this file is VOCABULARY -- definitions, notation and warranted supports, not one proof -- and every theorem stated with it had to load below it.  Its load.scm slot is unchanged. structure-library/subsequence-capture.scm}\par
\par\medskip\noindent\textbf{Definition} (\texttt{nn-enum}).\;$\operatorname{nn\text{-}enum}(s) :=$
\begin{equation*}
\varepsilon(\{f \in (\mathbb{N} \to s) : \forall m, n \in \mathbb{N}.\; ((m < n) \Rightarrow (f(m) < f(n)))\})
\end{equation*}


\vnbbookfile{structure-library/cauchy-subsequence}{structure-library/cauchy-subsequence}
\noindent{\small MOVED 2026-09-20 (batch 12-A) from theorem-library/: this file is VOCABULARY -- definitions, notation and warranted supports, not one proof -- and every theorem stated with it had to load below it.  Its load.scm slot is unchanged. RETIRED 2026-09-17 (proven): subseq-is-fun -- theorem-library/rake-analysis2.scm}\par
\par\medskip\noindent\textbf{Definition} (\texttt{strictly-mono-nn}).\;For every $\varphi$, 
\begin{equation*}
\begin{array}{@{}l@{}} \operatorname{strictly\text{-}mono\text{-}nn}(\varphi) \\ \Leftrightarrow\; ((\varphi \in (\mathbb{N} \to \mathbb{N})) \wedge \forall m, \operatorname{n\_} \in \mathbb{N}.\; ((m < \operatorname{n\_}) \Rightarrow (\varphi(m) < \varphi(\operatorname{n\_})))) \end{array}
\end{equation*}

\par\medskip\noindent\textbf{Definition} (\texttt{subseq}).\;$\operatorname{subseq}(f, \varphi) :=$ $(\lambda k \in \mathbb{N}.\; f(\varphi(k)))$
\par\medskip\noindent\textbf{Definition} (\texttt{is-subsequence}).\;For all $s$, $y$ and $f$, 
\begin{equation*}
\begin{array}{@{}l@{}} \operatorname{is\text{-}subsequence}(s, y, f) \\ \Leftrightarrow\; ((f \in (\mathbb{N} \to \operatorname{pts}(s))) \wedge \exists \varphi.\; (\operatorname{strictly\text{-}mono\text{-}nn}(\varphi) \wedge (y = \operatorname{subseq}(f, \varphi)))) \end{array}
\end{equation*}

\par\medskip\noindent\textbf{Definition} (\texttt{subsqn}).\;For all $y$ and $f$, 
\begin{equation*}
\operatorname{subsqn}(y, f) \Leftrightarrow \exists \varphi.\; (\operatorname{strictly\text{-}mono\text{-}nn}(\varphi) \wedge (y = \operatorname{subseq}(f, \varphi)))
\end{equation*}

\par\medskip\noindent\textbf{Definition} (\texttt{is-eps-cauchy-seq}).\;For all $s$, $\mathit{eps}$ and $y$, 
\begin{equation*}
\begin{array}{@{}l@{}} \operatorname{is\text{-}eps\text{-}cauchy\text{-}seq}(s, \mathit{eps}, y) \\ \Leftrightarrow\; \operatorname{is\text{-}metric\text{-}space}(s) \wedge \\ y \in (\mathbb{N} \to \operatorname{pts}(s)) \wedge \\ \forall m, \operatorname{n\_} \in \mathbb{N}.\; (\operatorname{dist}(s)(y(m), y(\operatorname{n\_})) \leq \mathit{eps}) \end{array}
\end{equation*}

\par\medskip\noindent\textbf{Definition} (\texttt{null-rr-seq}).\;For every $\mathit{rad}$, 
\begin{equation*}
\begin{array}{@{}l@{}} \operatorname{null\text{-}rr\text{-}seq}(\mathit{rad}) \\ \Leftrightarrow\; \mathit{rad} \in (\mathbb{N} \to \mathbb{R}) \wedge \\ \forall k \in \mathbb{N}.\; \operatorname{pos\text{-}rr}(\operatorname{rad}(k)) \wedge \\ \forall \mathit{eps}. \\ \quad \operatorname{pos\text{-}rr}(\mathit{eps}) \Rightarrow  \\ \quad \quad \exists n \in \mathbb{N}.\; \forall k.\; (((k \in \mathbb{N}) \wedge (n \leq k)) \Rightarrow (\operatorname{rad}(k) \leq \mathit{eps})) \end{array}
\end{equation*}


\vnbbookfile{structure-library/extended-sum}{structure-library/extended-sum}
\noindent{\small MOVED 2026-09-20 (batch 12-A) from theorem-library/: this file is VOCABULARY -- definitions, notation and warranted supports, not one proof -- and every theorem stated with it had to load below it.  Its load.scm slot is unchanged. structure-library/extended-sum.scm -- ESUM: the unordered sum in RR-POS-STAR}\par
\par\medskip\noindent\textbf{Definition} (\texttt{esum}).\;$\operatorname{esum}(\operatorname{f\_}) :=$
\begin{equation*}
\operatorname{esup}(\{\operatorname{v\_} \in \operatorname{rr\text{-}pos\text{-}star} : \exists \operatorname{t\_}.\; (((\operatorname{t\_} \in \mathbf{Set}) \wedge ((|\operatorname{t\_}| \in \mathbb{N}) \wedge (\operatorname{t\_} \subseteq \operatorname{dom}(\operatorname{f\_})))) \wedge (\operatorname{v\_} = \operatorname{finsum}(\operatorname{rr\text{-}pos\text{-}star\text{-}add\text{-}monoid}, \operatorname{f\_}, \operatorname{t\_})))\})
\end{equation*}


\vnbbookfile{structure-library/extended-arith}{structure-library/extended-arith}
\noindent{\small extended-arith.scm -- the rest of the arithmetic and order vocabulary of RR-POS-STAR = [0, +inf]: multiplication, infima, and the limit operations of a sequence.}\par
\par\medskip\noindent\textbf{Definition} (\texttt{einf}).\;$\operatorname{einf}(s) :=$ $\operatorname{esup}(\operatorname{elower\text{-}bounds}(s))$
\par\medskip\noindent\textbf{Definition} (\texttt{etail}).\;$\operatorname{etail}(f, n) :=$ $\operatorname{image}(f, \{\operatorname{p\_} \in \mathbb{N} : (n \leq \operatorname{p\_})\})$
\par\medskip\noindent\textbf{Definition} (\texttt{elower-bounds}).\;$\operatorname{elower\text{-}bounds}(s) :=$ $\{\operatorname{b\_} \in \operatorname{rr\text{-}pos\text{-}star} : \forall \operatorname{x\_} \in s.\; (\operatorname{b\_} \leq \operatorname{x\_})\}$
\par\medskip\noindent\textbf{Definition} (\texttt{eliminf}).\;$\operatorname{eliminf}(f) :=$ $\operatorname{esup}(\operatorname{image}((\lambda \operatorname{n\_} \in \mathbb{N}.\; \operatorname{einf}(\operatorname{etail}(f, \operatorname{n\_}))), \mathbb{N}))$
\par\medskip\noindent\textbf{Definition} (\texttt{elimsup}).\;$\operatorname{elimsup}(f) :=$ $\operatorname{einf}(\operatorname{image}((\lambda \operatorname{n\_} \in \mathbb{N}.\; \operatorname{esup}(\operatorname{etail}(f, \operatorname{n\_}))), \mathbb{N}))$
\par\medskip\noindent\textbf{Definition} (\texttt{econverges-to}).\;For all $f$ and $l$, 
\begin{equation*}
\begin{array}{@{}l@{}} \operatorname{econverges\text{-}to}(f, l) \\ \Leftrightarrow\; ((f \in (\mathbb{N} \to \operatorname{rr\text{-}pos\text{-}star})) \wedge ((l \in \operatorname{rr\text{-}pos\text{-}star}) \wedge ((\operatorname{eliminf}(f) = l) \wedge (\operatorname{elimsup}(f) = l)))) \end{array}
\end{equation*}

\vnbdef{etimes-def}{\vnbstmt{(\mathit{etimes} \simeq (\lambda \langle \operatorname{x\_}, \operatorname{y\_} \rangle \in (\operatorname{rr\text{-}pos\text{-}star} \times \operatorname{rr\text{-}pos\text{-}star}).\; \left(\text{if }((\operatorname{x\_} = \operatorname{pos\text{-}inf}) \vee (\operatorname{y\_} = \operatorname{pos\text{-}inf}))\text{ then }\left(\text{if }((\operatorname{x\_} = 0) \vee (\operatorname{y\_} = 0))\text{ then }0\text{ else }\operatorname{pos\text{-}inf}\right)\text{ else }(\operatorname{x\_} \cdot \operatorname{y\_})\right)))}}

\vnbbookfile{structure-library/summability}{structure-library/summability}
\noindent{\small MOVED 2026-09-20 (batch 12-A) from theorem-library/: this file is VOCABULARY -- definitions, notation and warranted supports, not one proof -- and every theorem stated with it had to load below it.  Its load.scm slot is unchanged. summability.scm -- unconditional summability of a normed-AG-valued function}\par
\par\medskip\noindent\textbf{Definition} (\texttt{sums-to}).\;For all $\mathit{grp}$, $f$ and $r$, 
\begin{equation*}
\begin{array}{@{}l@{}} \operatorname{sums\text{-}to}(\mathit{grp}, f, r) \\ \Leftrightarrow\; \forall \mathit{eps}. \\ \quad \operatorname{pos\text{-}rr}(\mathit{eps}) \Rightarrow  \\ \quad \quad \exists \mathit{fin} \in \mathbf{Set}. \\ \quad \quad \quad |\mathit{fin}| \in \mathbb{N} \wedge \\ \quad \quad \quad \mathit{fin} \subseteq \operatorname{dom}(f) \wedge \\ \quad \quad \quad \forall \mathit{ext}. \\ \quad \quad \quad \quad ((\mathit{ext} \in \mathbf{Set}) \wedge ((|\mathit{ext}| \in \mathbb{N}) \wedge ((\mathit{fin} \subseteq \mathit{ext}) \wedge (\mathit{ext} \subseteq \operatorname{dom}(f))))) \Rightarrow  \\ \quad \quad \quad \quad \quad \operatorname{nrm}(\mathit{grp})(\operatorname{opr}(\mathit{grp})(r, \operatorname{inv}(\mathit{grp})(\operatorname{finsum}(\operatorname{normed\text{-}ag\text{-}as\text{-}abelian\text{-}group}(\mathit{grp}), f, \mathit{ext})))) < \mathit{eps} \end{array}
\end{equation*}

\par\medskip\noindent\textbf{Definition} (\texttt{is-summable}).\;For all $\mathit{grp}$ and $f$, 
\begin{equation*}
\operatorname{is\text{-}summable}(\mathit{grp}, f) \Leftrightarrow \exists r \in \operatorname{carr}(\mathit{grp}).\; \operatorname{sums\text{-}to}(\mathit{grp}, f, r)
\end{equation*}

\par\medskip\noindent\textbf{Definition} (\texttt{is-absolutely-summable}).\;For all $\mathit{grp}$ and $f$, 
\begin{equation*}
\operatorname{is\text{-}absolutely\text{-}summable}(\mathit{grp}, f) \Leftrightarrow (\operatorname{esum}((\lambda i \in \operatorname{dom}(f).\; \operatorname{nrm}(\mathit{grp})(f(i)))) \in \mathbb{R})
\end{equation*}


\vnbbookfile{structure-library/power-series}{structure-library/power-series}
\noindent{\small MOVED 2026-09-20 (batch 12-A) from theorem-library/: this file is VOCABULARY -- definitions, notation and warranted supports, not one proof -- and every theorem stated with it had to load below it.  Its load.scm slot is unchanged. power-series.scm -- real power series  Sum\_\{n$>$=0\} coef(n) x\textasciicircum{}n  in RR.}\par
\par\medskip\noindent\textbf{Definition} (\texttt{ps-partial-sum}).\;$\operatorname{ps\text{-}partial\text{-}sum}(\mathit{coef}, x, k) :=$
\begin{equation*}
\operatorname{sum\text{-}ag}(\operatorname{normed\text{-}field\text{-}additive\text{-}ag}(\operatorname{rr\text{-}normed\text{-}field}), (\lambda n \in \mathbb{N}.\; (\operatorname{coef}(n) \cdot \operatorname{power}(x, n))), k)
\end{equation*}

\par\medskip\noindent\textbf{Definition} (\texttt{ps-converges-to-at}).\;For all $\mathit{coef}$, $x$ and $l$, 
\begin{equation*}
\begin{array}{@{}l@{}} \operatorname{ps\text{-}converges\text{-}to\text{-}at}(\mathit{coef}, x, l) \\ \Leftrightarrow\; \operatorname{converges\text{-}to}(\operatorname{rr\text{-}ms}, (\lambda k \in \mathbb{N}.\; \operatorname{ps\text{-}partial\text{-}sum}(\mathit{coef}, x, k)), l) \end{array}
\end{equation*}

\par\medskip\noindent\textbf{Definition} (\texttt{ps-converges-at}).\;For all $\mathit{coef}$ and $x$, 
\begin{equation*}
\operatorname{ps\text{-}converges\text{-}at}(\mathit{coef}, x) \Leftrightarrow \operatorname{converges}(\operatorname{rr\text{-}ms}, (\lambda k \in \mathbb{N}.\; \operatorname{ps\text{-}partial\text{-}sum}(\mathit{coef}, x, k)))
\end{equation*}

\par\medskip\noindent\textbf{Definition} (\texttt{ps-ratio-limit}).\;For all $\mathit{coef}$ and $l$, 
\begin{equation*}
\begin{array}{@{}l@{}} \operatorname{ps\text{-}ratio\text{-}limit}(\mathit{coef}, l) \\ \Leftrightarrow\; \operatorname{converges\text{-}to}(\operatorname{rr\text{-}ms}, (\lambda n \in \mathbb{N}.\; \lvert \cdot(\operatorname{coef}(\operatorname{succ}(n)), \operatorname{recip}(\operatorname{coef}(n))) \rvert), l) \end{array}
\end{equation*}

\par\medskip\noindent\textbf{Definition} (\texttt{ps-absolutely-converges-at}).\;For all $\mathit{coef}$ and $x$, 
\begin{equation*}
\operatorname{ps\text{-}absolutely\text{-}converges\text{-}at}(\mathit{coef}, x) \Leftrightarrow \operatorname{ps\text{-}converges\text{-}at}((\lambda n \in \mathbb{N}.\; \lvert \operatorname{coef}(n) \rvert), \lvert x \rvert)
\end{equation*}

\par\medskip\noindent\textbf{Definition} (\texttt{series-partial-sum}).\;$\operatorname{series\text{-}partial\text{-}sum}(f, k) :=$ $\operatorname{sum\text{-}ag}(\operatorname{normed\text{-}field\text{-}additive\text{-}ag}(\operatorname{rr\text{-}normed\text{-}field}), f, k)$
\par\medskip\noindent\textbf{Definition} (\texttt{series-converges-to}).\;For all $f$ and $l$, 
\begin{equation*}
\operatorname{series\text{-}converges\text{-}to}(f, l) \Leftrightarrow \operatorname{converges\text{-}to}(\operatorname{rr\text{-}ms}, (\lambda k \in \mathbb{N}.\; \operatorname{series\text{-}partial\text{-}sum}(f, k)), l)
\end{equation*}

\par\medskip\noindent\textbf{Definition} (\texttt{series-converges}).\;For every $f$, 
\begin{equation*}
\operatorname{series\text{-}converges}(f) \Leftrightarrow \operatorname{converges}(\operatorname{rr\text{-}ms}, (\lambda k \in \mathbb{N}.\; \operatorname{series\text{-}partial\text{-}sum}(f, k)))
\end{equation*}


\vnbbookfile{structure-library/product-metric}{structure-library/product-metric}
\noindent{\small product-metric.scm -- the countable product of metric spaces, as a metric space carrying the PRODUCT TOPOLOGY.  Library phase: functoids for the construction, warranted supports for the metric/topology facts.}\par
\par\medskip\noindent\textbf{Definition} (\texttt{product-metric-w}).\;$\operatorname{product\text{-}metric\text{-}w}(\mathit{ms}, w) :=$
\begin{equation*}
\langle \operatorname{product\text{-}carrier}(\mathit{ms}), (\lambda \langle x, y \rangle \in (\operatorname{product\text{-}carrier}(\mathit{ms}) \times \operatorname{product\text{-}carrier}(\mathit{ms})).\; (\iota\, l.\;\operatorname{series\text{-}converges\text{-}to}((\lambda n \in \mathbb{N}.\; (w(n) \cdot \operatorname{dist}(\operatorname{bdd\text{-}metric}(\operatorname{ms}(n)))(x(n), y(n)))), l))) \rangle
\end{equation*}

\par\medskip\noindent\textbf{Definition} (\texttt{product-metric}).\;$\operatorname{product\text{-}metric}(\mathit{ms}) :=$ $\operatorname{product\text{-}metric\text{-}w}(\mathit{ms}, (\lambda n \in \mathbb{N}.\; /(1, \operatorname{power}(2, (n + 1)))))$
\par\medskip\noindent\textbf{Definition} (\texttt{is-ms-sequence}).\;For every $\mathit{ms}$, 
\begin{equation*}
\operatorname{is\text{-}ms\text{-}sequence}(\mathit{ms}) \Leftrightarrow \forall n \in \mathbb{N}.\; \operatorname{is\text{-}metric\text{-}space}(\operatorname{ms}(n))
\end{equation*}

\par\medskip\noindent\textbf{Definition} (\texttt{summable-weight}).\;For every $w$, 
\begin{equation*}
\operatorname{summable\text{-}weight}(w) \Leftrightarrow ((w \in (\mathbb{N} \to \mathbb{R})) \wedge (\forall n \in \mathbb{N}.\; (0 < w(n)) \wedge \operatorname{series\text{-}converges}(w)))
\end{equation*}

\par\medskip\noindent\textbf{Definition} (\texttt{product-carrier}).\;$\operatorname{product\text{-}carrier}(\mathit{ms}) :=$
\begin{equation*}
\{x \in (\mathbb{N} \to \operatorname{big\text{-}union}(n, \mathbb{N}, \operatorname{pts}(\operatorname{ms}(n)))) : \forall n \in \mathbb{N}.\; (x(n) \in \operatorname{pts}(\operatorname{ms}(n)))\}
\end{equation*}

\par\medskip\noindent\textbf{Definition} (\texttt{product-proj}).\;$\operatorname{product\text{-}proj}(\mathit{ms}, n) :=$ $(\lambda x \in \operatorname{product\text{-}carrier}(\mathit{ms}).\; x(n))$

\vnbbookfile{structure-library/c-metric-space}{structure-library/c-metric-space}
\noindent{\small c-metric-space.scm -- C-METRIC-SPACE, a set carrying a COUNTABLE FAMILY OF PSEUDOMETRICS that separates points, and the CANONICAL METRIC it induces.}\par
\vnbstructdef{c-metric-space}{A c metric space $s$ is a 2-tuple whose components are listed below.\par
\begin{enumerate}[nosep]
\item $\operatorname{pts}(s) = \operatorname{nth}(1, s)$, a carrier (\texttt{carriers}): $\operatorname{pts}(s) \in \mathbf{Set}$.
\item $\operatorname{dists}(s) = \operatorname{nth}(2, s)$, an operation (\texttt{op}): $\operatorname{dists}(s) \in (\mathbb{N} \to ((\operatorname{pts}(s) \times \operatorname{pts}(s)) \to \mathbb{R}))$.
\end{enumerate}
The laws (the \texttt{property} and \texttt{law} clauses):
\begin{itemize}[nosep]
\item $\forall \operatorname{k\_} \in \mathbb{N}.\; \operatorname{is\text{-}pseudometric}(\operatorname{dists}(s)(\operatorname{k\_}), \operatorname{pts}(s))$
\item
\begin{equation*}
\begin{array}{@{}l@{}} \forall u, v \in \operatorname{pts}(s). \\ \quad \forall \operatorname{k\_} \in \mathbb{N}.\; (\operatorname{dists}(s)(\operatorname{k\_})(u, v) = 0) \Rightarrow (u = v) \end{array}
\end{equation*}

\end{itemize}
The characteristic law (the axiom \texttt{is-c-metric-space}): $\operatorname{is\text{-}c\text{-}metric\text{-}space}(s)$ holds if and only if $\operatorname{length}(s) = 2$, every component has the typing listed with it, and every law holds.
}
\par\medskip\noindent\textbf{Definition} (\texttt{c-metric}).\;$\operatorname{c\text{-}metric}(s) :=$ $\operatorname{c\text{-}metric\text{-}w}(s, (\lambda n \in \mathbb{N}.\; /(1, \operatorname{power}(2, (n + 1)))))$
\par\medskip\noindent\textbf{Definition} (\texttt{c-metric-w}).\;$\operatorname{c\text{-}metric\text{-}w}(s, w) :=$
\begin{equation*}
\langle \operatorname{pts}(s), (\lambda \langle u, v \rangle \in (\operatorname{pts}(s) \times \operatorname{pts}(s)).\; (\iota\, l.\;\operatorname{series\text{-}converges\text{-}to}((\lambda k \in \mathbb{N}.\; (w(k) \cdot \operatorname{min}(1, \operatorname{dists}(s)(k)(u, v)))), l))) \rangle
\end{equation*}

\par\medskip\noindent\textbf{Definition} (\texttt{c-pseudo-family}).\;$\operatorname{c\text{-}pseudo\text{-}family}(s) :=$ $(\lambda k \in \mathbb{N}.\; \langle \operatorname{pts}(s), \operatorname{dists}(s)(k) \rangle)$
\par\medskip\noindent\textbf{Definition} (\texttt{hom-c-metric-space}).\;$\operatorname{hom\text{-}c\text{-}metric\text{-}space}(a, b) :=$
\begin{equation*}
\{\operatorname{homf\_} \in (\operatorname{pts}(a) \to \operatorname{pts}(b)) : \operatorname{is\text{-}hom\text{-}c\text{-}metric\text{-}space}(a, b, \operatorname{homf\_})\}
\end{equation*}

\vnbdef{c-metric-space-class}{\vnbstmt{\forall s.\; ((s \in \operatorname{c\text{-}metric\text{-}space}) \Leftrightarrow \operatorname{is\text{-}c\text{-}metric\text{-}space}(s))}}
\vnbdef{is-hom-c-metric-space-def}{\vnbstmt{\begin{array}{@{}l@{}} \forall a, b, f. \\ \quad \operatorname{is\text{-}hom\text{-}c\text{-}metric\text{-}space}(a, b, f) \\ \quad \Leftrightarrow\; \operatorname{is\text{-}c\text{-}metric\text{-}space}(a) \wedge \\ \quad \operatorname{is\text{-}c\text{-}metric\text{-}space}(b) \wedge \\ \quad f \in (\operatorname{pts}(a) \to \operatorname{pts}(b)) \wedge \\ \quad \forall x_{1} \in \mathbb{N}.\; (\operatorname{dists}(a)(x_{1}) = \operatorname{dists}(b)(x_{1})) \end{array}}}

\vnbbookfile{structure-library/seq-compact-product}{structure-library/seq-compact-product}
\noindent{\small MOVED 2026-09-20 (batch 12-A) from theorem-library/: this file is VOCABULARY -- definitions, notation and warranted supports, not one proof -- and every theorem stated with it had to load below it.  Its load.scm slot is unchanged. structure-library/seq-compact-product.scm}\par
\par\medskip\noindent\textbf{Definition} (\texttt{seq-compact}).\;For every $s$, 
\begin{equation*}
\begin{array}{@{}l@{}} \operatorname{seq\text{-}compact}(s) \\ \Leftrightarrow\; \operatorname{is\text{-}metric\text{-}space}(s) \wedge \\ \forall f \in (\mathbb{N} \to \operatorname{pts}(s)). \exists \varphi. \\ \quad (\operatorname{strictly\text{-}mono\text{-}nn}(\varphi) \wedge \exists l \in \operatorname{pts}(s).\; \operatorname{converges\text{-}to}(s, \operatorname{subseq}(f, \varphi), l)) \end{array}
\end{equation*}

\par\medskip\noindent\textbf{Definition} (\texttt{converges-along}).\;For all $s$, $g$, $b$ and $p$, 
\begin{equation*}
\begin{array}{@{}l@{}} \operatorname{converges\text{-}along}(s, g, b, p) \\ \Leftrightarrow\; \operatorname{is\text{-}metric\text{-}space}(s) \wedge \\ g \in (\mathbb{N} \to \operatorname{pts}(s)) \wedge \\ p \in \operatorname{pts}(s) \wedge \\ \forall \mathit{eps}. \\ \quad \operatorname{pos\text{-}rr}(\mathit{eps}) \Rightarrow  \\ \quad \quad \exists n \in \mathbb{N}. \forall i \in \mathbb{N}. \\ \quad \quad \quad (i \in b) \Rightarrow ((n \leq i) \Rightarrow (\operatorname{dist}(s)(g(i), p) \leq \mathit{eps})) \end{array}
\end{equation*}


\vnbbookfile{theorem-library/subseq-apply}{theorem-library/subseq-apply}
\noindent{\small subseq-apply.scm -- the VALUE of a subsequence at one TYPED index.}\par
\vnbprop{subseq-apply}{\vnbstmt{\forall f, i.; \forall k \in \mathbb{N}.\; (\operatorname{subseq}(f, i)(k) \simeq f(i(k)))}}

\vnbbookfile{theorem-library/bdd-metric-carrier}{theorem-library/bdd-metric-carrier}
\noindent{\small bdd-metric-carrier.scm -- BDD-METRIC keeps the point set, proven.}\par
\vnbprop{bdd-metric-carrier}{\vnbstmt{\forall s.\; (\operatorname{is\text{-}metric\text{-}space}(s) \Rightarrow (\operatorname{pts}(\operatorname{bdd\text{-}metric}(s)) \simeq \operatorname{pts}(s)))}}

\vnbbookfile{theorem-library/normed-field-ring-view}{theorem-library/normed-field-ring-view}
\noindent{\small normed-field-ring-view.scm -- THE SCALAR BRIDGE, PROVEN.}\par
\vnbprop{normed-field-ring-view-carr}{\vnbstmt{\begin{array}{@{}l@{}} \forall f. \\ \quad \operatorname{is\text{-}normed\text{-}field}(f) \Rightarrow (\operatorname{carr}(\operatorname{normed\text{-}field\text{-}as\text{-}commutative\text{-}ring}(f)) \simeq \operatorname{carr}(f)) \end{array}}}
\vnbprop{normed-field-ring-view-add}{\vnbstmt{\begin{array}{@{}l@{}} \forall f. \\ \quad \operatorname{is\text{-}normed\text{-}field}(f) \Rightarrow (\operatorname{add}(\operatorname{normed\text{-}field\text{-}as\text{-}commutative\text{-}ring}(f)) \simeq \operatorname{add}(f)) \end{array}}}
\vnbprop{normed-field-ring-view-mul}{\vnbstmt{\begin{array}{@{}l@{}} \forall f. \\ \quad \operatorname{is\text{-}normed\text{-}field}(f) \Rightarrow (\operatorname{mul}(\operatorname{normed\text{-}field\text{-}as\text{-}commutative\text{-}ring}(f)) \simeq \operatorname{mul}(f)) \end{array}}}
\vnbprop{normed-field-ring-view-neg}{\vnbstmt{\begin{array}{@{}l@{}} \forall f. \\ \quad \operatorname{is\text{-}normed\text{-}field}(f) \Rightarrow (\operatorname{neg}(\operatorname{normed\text{-}field\text{-}as\text{-}commutative\text{-}ring}(f)) \simeq \operatorname{neg}(f)) \end{array}}}
\vnbprop{normed-field-ring-view-zero}{\vnbstmt{\begin{array}{@{}l@{}} \forall f. \\ \quad \operatorname{is\text{-}normed\text{-}field}(f) \Rightarrow (\operatorname{zero}(\operatorname{normed\text{-}field\text{-}as\text{-}commutative\text{-}ring}(f)) \simeq \operatorname{zero}(f)) \end{array}}}
\vnbprop{normed-field-ring-view-one}{\vnbstmt{\begin{array}{@{}l@{}} \forall f. \\ \quad \operatorname{is\text{-}normed\text{-}field}(f) \Rightarrow (\operatorname{one}(\operatorname{normed\text{-}field\text{-}as\text{-}commutative\text{-}ring}(f)) \simeq \operatorname{one}(f)) \end{array}}}

\vnbbookfile{theorem-library/rake-setoid}{theorem-library/rake-setoid}
\noindent{\small A file shared across books; the entries that belong to this one.}\par
\noindent{\small rake-setoid.scm -- BATCH F of the 2026-09-17 rake: setoids and carriers. Seven asserted leaves of the PSS, PROVEN.  Every statement here is its support's statement UNCHANGED.}\par
\vnbprop{vec-is-set}{\vnbstmt{\forall m.\; (\operatorname{is\text{-}normed\text{-}vector\text{-}space}(m) \Rightarrow (\operatorname{vec}(m) \in \mathbf{Set}))}}

\vnbbookfile{theorem-library/rake-analysis2}{theorem-library/rake-analysis2}
\noindent{\small rake-analysis2.scm -- rake batch L (2026-09-17): five asserted leaves of the analysis / sets area, plus the general lemma that unblocked one of them.}\par
\vnbprop{subseq-is-fun}{\vnbstmt{\begin{array}{@{}l@{}} \forall s, f, i. \\ \quad ((f \in (\mathbb{N} \to \operatorname{pts}(s))) \wedge \operatorname{strictly\text{-}mono\text{-}nn}(i)) \Rightarrow  \\ \quad \quad \operatorname{subseq}(f, i) \in (\mathbb{N} \to \operatorname{pts}(s)) \end{array}}}
\vnbprop{metric-hom-is-continuous}{\vnbstmt{\forall s, t, f.\; (\operatorname{is\text{-}hom\text{-}metric\text{-}space}(s, t, f) \Rightarrow \operatorname{is\text{-}continuous}(s, t, f))}}
\vnbprop{centre-set-contains-choice}{\vnbstmt{\begin{array}{@{}l@{}} \forall s, r, f.; \forall u \in f. \\ \quad \exists c.\; (c \in \operatorname{centres}(s, u, r)) \Rightarrow  \\ \quad \quad \varepsilon(\operatorname{centres}(s, u, r)) \in \operatorname{centre\text{-}set}(s, r, f) \end{array}}}

\vnbbookfile{theorem-library/rake-infinity-points}{theorem-library/rake-infinity-points}
\noindent{\small theorem-library/rake-infinity-points.scm -- the five DERIVABLE axioms of structure-library/extended-reals.scm, proven.  Rake "infinity points".}\par
\vnbprop{rr-subset-rr-star}{\vnbstmt{\mathbb{R} \subseteq \operatorname{rr\text{-}star}}}
\vnbprop{pos-inf-in-rr-star}{\vnbstmt{\operatorname{pos\text{-}inf} \in \operatorname{rr\text{-}star}}}
\vnbprop{neg-inf-in-rr-star}{\vnbstmt{\operatorname{neg\text{-}inf} \in \operatorname{rr\text{-}star}}}
\vnbprop{pos-inf-not-in-rr}{\vnbstmt{\neg (\operatorname{pos\text{-}inf} \in \mathbb{R})}}
\vnbprop{pos-inf-neq-neg-inf}{\vnbstmt{\neg (\operatorname{pos\text{-}inf} = \operatorname{neg\text{-}inf})}}

\vnbbookfile{theorem-library/rake-rr-pos-star}{theorem-library/rake-rr-pos-star}
\noindent{\small theorem-library/rake-rr-pos-star.scm -- the six CARRIER facts of RR-POS-STAR = [0, +inf], proven.  Rake batch 5c-K.}\par
\vnbprop{zero-in-rr-pos-star}{\vnbstmt{0 \in \operatorname{rr\text{-}pos\text{-}star}}}
\vnbprop{pos-inf-in-rr-pos-star}{\vnbstmt{\operatorname{pos\text{-}inf} \in \operatorname{rr\text{-}pos\text{-}star}}}
\vnbprop{neg-inf-not-in-rr-pos-star}{\vnbstmt{\neg (\operatorname{neg\text{-}inf} \in \operatorname{rr\text{-}pos\text{-}star})}}
\vnbprop{rr-pos-star-nonneg}{\vnbstmt{\forall x \in \operatorname{rr\text{-}pos\text{-}star}.\; (0 \leq x)}}
\vnbprop{rr-pos-star-below-pos-inf}{\vnbstmt{\forall x \in \operatorname{rr\text{-}pos\text{-}star}.\; (x \leq \operatorname{pos\text{-}inf})}}
\vnbprop{rr-pos-star-subset-rr-star}{\vnbstmt{\operatorname{rr\text{-}pos\text{-}star} \subseteq \operatorname{rr\text{-}star}}}

\vnbbookfile{theorem-library/rake-eplus-defined}{theorem-library/rake-eplus-defined}
\noindent{\small theorem-library/rake-eplus-defined.scm -- extended addition on [0,+inf] DEFINED, and its laws proven.  Rake batch 5c-V.}\par
\vnbprop{apply-congruence-2}{\vnbstmt{\forall f, g, u, v.\; ((f \simeq g) \Rightarrow (f(u, v) \simeq g(u, v)))}}
\vnbprop{rr-pos-star-is-set}{\vnbstmt{\operatorname{rr\text{-}pos\text{-}star} \in \mathbf{Set}}}
\vnbprop{eplus-in-fun}{\vnbstmt{\mathit{eplus} \in ((\operatorname{rr\text{-}pos\text{-}star} \times \operatorname{rr\text{-}pos\text{-}star}) \to \operatorname{rr\text{-}pos\text{-}star})}}
\vnbprop{eplus-pos-inf-left}{\vnbstmt{\forall y \in \operatorname{rr\text{-}pos\text{-}star}.\; (\operatorname{eplus}(\operatorname{pos\text{-}inf}, y) = \operatorname{pos\text{-}inf})}}
\vnbprop{eplus-pos-inf-right}{\vnbstmt{\forall x \in \operatorname{rr\text{-}pos\text{-}star}.\; (\operatorname{eplus}(x, \operatorname{pos\text{-}inf}) = \operatorname{pos\text{-}inf})}}
\vnbprop{eplus-real-defined}{\vnbstmt{\begin{array}{@{}l@{}} \forall x, y. \\ \quad ((x \in \mathbb{R}) \wedge ((0 \leq x) \wedge ((y \in \mathbb{R}) \wedge (0 \leq y)))) \Rightarrow (\operatorname{eplus}(x, y) = (x + y)) \end{array}}}

\vnbbookfile{theorem-library/rake-etimes-defined}{theorem-library/rake-etimes-defined}
\noindent{\small theorem-library/rake-etimes-defined.scm -- extended MULTIPLICATION on [0,+inf] DEFINED, and its laws proven.  Rake batch 5c-Y.}\par
\vnbprop{etimes-in-fun}{\vnbstmt{\mathit{etimes} \in ((\operatorname{rr\text{-}pos\text{-}star} \times \operatorname{rr\text{-}pos\text{-}star}) \to \operatorname{rr\text{-}pos\text{-}star})}}
\vnbprop{etimes-zero-pos-inf}{\vnbstmt{\operatorname{etimes}(0, \operatorname{pos\text{-}inf}) = 0}}
\vnbprop{etimes-pos-inf-zero}{\vnbstmt{\operatorname{etimes}(\operatorname{pos\text{-}inf}, 0) = 0}}
\vnbprop{etimes-pos-inf-left}{\vnbstmt{\forall y \in \operatorname{rr\text{-}pos\text{-}star}.\; (\neg (y = 0) \Rightarrow (\operatorname{etimes}(\operatorname{pos\text{-}inf}, y) = \operatorname{pos\text{-}inf}))}}
\vnbprop{etimes-pos-inf-right}{\vnbstmt{\forall x \in \operatorname{rr\text{-}pos\text{-}star}.\; (\neg (x = 0) \Rightarrow (\operatorname{etimes}(x, \operatorname{pos\text{-}inf}) = \operatorname{pos\text{-}inf}))}}
\vnbprop{etimes-real-defined}{\vnbstmt{\begin{array}{@{}l@{}} \forall x, y. \\ \quad ((x \in \mathbb{R}) \wedge ((0 \leq x) \wedge ((y \in \mathbb{R}) \wedge (0 \leq y)))) \Rightarrow (\operatorname{etimes}(x, y) = (x \cdot y)) \end{array}}}

\vnbbookfile{theorem-library/discrete-space}{theorem-library/discrete-space}
\noindent{\small discrete-space.scm -- THE FIRST WITNESS OF A MEASURABLE SPACE, and the first outright inhabitant of TOP-SPACE.}\par
\par\medskip\noindent\textbf{Definition} (\texttt{discrete-space}).\;$\operatorname{discrete\text{-}space}(\operatorname{a\_}) :=$ $\langle \operatorname{a\_}, \operatorname{power}(\operatorname{a\_}) \rangle$
\vnbprop{ds-pts}{\vnbstmt{\forall a \in \mathbf{Set}.\; (\operatorname{pts}(\operatorname{discrete\text{-}space}(a)) \simeq a)}}
\vnbprop{ds-opens}{\vnbstmt{\forall a \in \mathbf{Set}.\; (\operatorname{opens}(\operatorname{discrete\text{-}space}(a)) \simeq \operatorname{power}(a))}}
\vnbprop{ds-sigma}{\vnbstmt{\forall a \in \mathbf{Set}.\; (\sigma(\operatorname{discrete\text{-}space}(a)) \simeq \operatorname{power}(a))}}
\vnbprop{discrete-space-is-top-space}{\vnbstmt{\forall a \in \mathbf{Set}.\; \operatorname{is\text{-}top\text{-}space}(\operatorname{discrete\text{-}space}(a))}}

\vnbbookfile{theorem-library/op-typing}{theorem-library/op-typing}
\noindent{\small op-typing.scm -- the codomain typing of a structure operation, in APPLIED form, PROVEN.  One driver, seven theorems.}\par
\vnbprop{metric-dist-real}{\vnbstmt{\begin{array}{@{}l@{}} \forall s.; \forall x, y \in \operatorname{pts}(s). \\ \quad \operatorname{is\text{-}metric\text{-}space}(s) \Rightarrow (\operatorname{dist}(s)(x, y) \in \mathbb{R}) \end{array}}}
\vnbprop{vnrm-real}{\vnbstmt{\begin{array}{@{}l@{}} \forall m, w. \\ \quad (\operatorname{is\text{-}normed\text{-}vector\text{-}space}(m) \wedge (w \in \operatorname{vec}(m))) \Rightarrow (\operatorname{vnrm}(m)(w) \in \mathbb{R}) \end{array}}}
\vnbprop{nvs-vadd-in-vec}{\vnbstmt{\begin{array}{@{}l@{}} \forall m.; \forall x, y \in \operatorname{vec}(m). \\ \quad \operatorname{is\text{-}normed\text{-}vector\text{-}space}(m) \Rightarrow (\operatorname{vadd}(m)(x, y) \in \operatorname{vec}(m)) \end{array}}}

\vnbbookfile{theorem-library/rr-is-normed-field}{theorem-library/rr-is-normed-field}
\noindent{\small rr-is-normed-field.scm -- IS-NORMED-FIELD(RR-NORMED-FIELD), PROVED.}\par
\vnbprop{rr-is-normed-field}{\vnbstmt{\operatorname{is\text{-}normed\text{-}field}(\operatorname{rr\text{-}normed\text{-}field})}}
\vnbprop{rr-scalar-ring-carr}{\vnbstmt{(\operatorname{carr}(\operatorname{normed\text{-}field\text{-}as\text{-}commutative\text{-}ring}(\operatorname{rr\text{-}normed\text{-}field})) \simeq \mathbb{R})}}
\vnbprop{rr-scalar-ring-add}{\vnbstmt{(\operatorname{add}(\operatorname{normed\text{-}field\text{-}as\text{-}commutative\text{-}ring}(\operatorname{rr\text{-}normed\text{-}field})) \simeq (\lambda \langle \operatorname{x\_}, \operatorname{y\_} \rangle \in (\mathbb{R} \times \mathbb{R}).\; (\operatorname{x\_} + \operatorname{y\_})))}}
\vnbprop{rr-scalar-ring-mul}{\vnbstmt{(\operatorname{mul}(\operatorname{normed\text{-}field\text{-}as\text{-}commutative\text{-}ring}(\operatorname{rr\text{-}normed\text{-}field})) \simeq (\lambda \langle \operatorname{x\_}, \operatorname{y\_} \rangle \in (\mathbb{R} \times \mathbb{R}).\; (\operatorname{x\_} \cdot \operatorname{y\_})))}}
\vnbprop{rr-scalar-ring-neg}{\vnbstmt{(\operatorname{neg}(\operatorname{normed\text{-}field\text{-}as\text{-}commutative\text{-}ring}(\operatorname{rr\text{-}normed\text{-}field})) \simeq (\lambda x \in \mathbb{R}.\; -x))}}
\vnbprop{rr-scalar-ring-zero}{\vnbstmt{(\operatorname{zero}(\operatorname{normed\text{-}field\text{-}as\text{-}commutative\text{-}ring}(\operatorname{rr\text{-}normed\text{-}field})) \simeq 0)}}
\vnbprop{rr-scalar-ring-one}{\vnbstmt{(\operatorname{one}(\operatorname{normed\text{-}field\text{-}as\text{-}commutative\text{-}ring}(\operatorname{rr\text{-}normed\text{-}field})) \simeq 1)}}

\vnbbookfile{theorem-library/rr-nvs-exemplification}{theorem-library/rr-nvs-exemplification}
\noindent{\small rr-nvs-exemplification.scm -- THE FIRST WITNESS OF A NORMED VECTOR SPACE.}\par
\par\medskip\noindent\textbf{Definition} (\texttt{rr-nvs}).\;
\begin{equation*}
\begin{array}{@{}l@{}} \operatorname{rr\text{-}nvs} \\ =\; \langle \operatorname{normed\text{-}field\text{-}as\text{-}commutative\text{-}ring}(\operatorname{rr\text{-}normed\text{-}field}), \mathbb{R}, (\lambda \langle \operatorname{x\_}, \operatorname{y\_} \rangle \in (\mathbb{R} \times \mathbb{R}).\; +(\operatorname{x\_}, \operatorname{y\_})), 0, (\lambda \operatorname{x\_} \in \mathbb{R}.\; -(\operatorname{x\_})), (\lambda \langle \operatorname{r\_}, \operatorname{x\_} \rangle \in (\mathbb{R} \times \mathbb{R}).\; \cdot(\operatorname{r\_}, \operatorname{x\_})), (\lambda \operatorname{x\_} \in \mathbb{R}.\; \lvert \operatorname{x\_} \rvert) \rangle \end{array}
\end{equation*}

\vnbdef{rr-nvs@vnrm}{\vnbstmt{(\operatorname{vnrm}(\operatorname{rr\text{-}nvs}) \simeq (\lambda x \in \mathbb{R}.\; \lvert x \rvert))}}
\vnbdef{rr-nvs@vec}{\vnbstmt{(\operatorname{vec}(\operatorname{rr\text{-}nvs}) \simeq \mathbb{R})}}
\vnbdef{rr-nvs@act}{\vnbstmt{(\operatorname{act}(\operatorname{rr\text{-}nvs}) \simeq (\lambda \langle \operatorname{r\_}, \operatorname{x\_} \rangle \in (\mathbb{R} \times \mathbb{R}).\; (\operatorname{r\_} \cdot \operatorname{x\_})))}}
\vnbdef{rr-nvs@vadd}{\vnbstmt{(\operatorname{vadd}(\operatorname{rr\text{-}nvs}) \simeq (\lambda \langle \operatorname{x\_}, \operatorname{y\_} \rangle \in (\mathbb{R} \times \mathbb{R}).\; (\operatorname{x\_} + \operatorname{y\_})))}}
\vnbdef{rr-nvs@vzero}{\vnbstmt{(\operatorname{vzero}(\operatorname{rr\text{-}nvs}) \simeq 0)}}
\vnbdef{rr-nvs@scal}{\vnbstmt{(\operatorname{scal}(\operatorname{rr\text{-}nvs}) \simeq \operatorname{normed\text{-}field\text{-}as\text{-}commutative\text{-}ring}(\operatorname{rr\text{-}normed\text{-}field}))}}
\vnbdef{rr-nvs@vneg}{\vnbstmt{(\operatorname{vneg}(\operatorname{rr\text{-}nvs}) \simeq (\lambda x \in \mathbb{R}.\; -x))}}
\vnbprop{rr-nvs-is-normed-vector-space}{\vnbstmt{\operatorname{is\text{-}normed\text{-}vector\text{-}space}(\operatorname{rr\text{-}nvs})}}

\vnbbookfile{theorem-library/rake-algebra}{theorem-library/rake-algebra}
\noindent{\small A file shared across books; the entries that belong to this one.}\par
\noindent{\small rake-algebra.scm -- BATCH D of the 2026-09-17 rake: seven asserted leaves of the PSS, PROVEN.  Every statement here is its support's statement UNCHANGED.}\par
\vnbprop{strictly-mono-ge-id-ind}{\vnbstmt{\forall k \in \mathbb{N}, i.\; (\operatorname{strictly\text{-}mono\text{-}nn}(i) \Rightarrow (k \leq i(k)))}}
\vnbprop{strictly-mono-ge-id}{\vnbstmt{\forall i.; \forall k \in \mathbb{N}.\; (\operatorname{strictly\text{-}mono\text{-}nn}(i) \Rightarrow (k \leq i(k)))}}

\vnbbookfile{theorem-library/rake-extended-order}{theorem-library/rake-extended-order}
\noindent{\small theorem-library/rake-extended-order.scm -- the order and addition theory of RR-POS-STAR = [0, +inf], built on the user's new axiom `pos-inf-above-reals' (structure-library/extended-reals.scm, stamped `primitive', 2026-09-18).}\par
\vnbprop{rr-pos-star-le-refl}{\vnbstmt{\forall x \in \operatorname{rr\text{-}pos\text{-}star}.\; (x \leq x)}}
\vnbprop{rr-pos-star-le-real-in-rr}{\vnbstmt{\forall x, y.\; ((((x \in \operatorname{rr\text{-}pos\text{-}star}) \wedge (y \in \mathbb{R})) \wedge (x \leq y)) \Rightarrow (x \in \mathbb{R}))}}
\vnbprop{rr-pos-star-le-trans}{\vnbstmt{\begin{array}{@{}l@{}} \forall x, y, z. \\ \quad (((x \in \operatorname{rr\text{-}pos\text{-}star}) \wedge ((y \in \operatorname{rr\text{-}pos\text{-}star}) \wedge (z \in \operatorname{rr\text{-}pos\text{-}star}))) \wedge ((x \leq y) \wedge (y \leq z))) \Rightarrow  \\ \quad \quad x \leq z \end{array}}}
\vnbprop{rr-pos-star-le-antisymm}{\vnbstmt{\begin{array}{@{}l@{}} \forall x, y. \\ \quad (((x \in \operatorname{rr\text{-}pos\text{-}star}) \wedge (y \in \operatorname{rr\text{-}pos\text{-}star})) \wedge ((x \leq y) \wedge (y \leq x))) \Rightarrow (x = y) \end{array}}}
\vnbprop{eplus-closed}{\vnbstmt{\begin{array}{@{}l@{}} \forall x, y. \\ \quad ((x \in \operatorname{rr\text{-}pos\text{-}star}) \wedge (y \in \operatorname{rr\text{-}pos\text{-}star})) \Rightarrow (\operatorname{eplus}(x, y) \in \operatorname{rr\text{-}pos\text{-}star}) \end{array}}}
\vnbprop{eplus-zero-left}{\vnbstmt{\forall x \in \operatorname{rr\text{-}pos\text{-}star}.\; (\operatorname{eplus}(0, x) = x)}}
\vnbprop{eplus-zero-right}{\vnbstmt{\forall x \in \operatorname{rr\text{-}pos\text{-}star}.\; (\operatorname{eplus}(x, 0) = x)}}
\vnbprop{eplus-comm}{\vnbstmt{\begin{array}{@{}l@{}} \forall x, y. \\ \quad ((x \in \operatorname{rr\text{-}pos\text{-}star}) \wedge (y \in \operatorname{rr\text{-}pos\text{-}star})) \Rightarrow (\operatorname{eplus}(x, y) = \operatorname{eplus}(y, x)) \end{array}}}
\vnbprop{eplus-assoc}{\vnbstmt{\begin{array}{@{}l@{}} \forall x, y, z. \\ \quad ((x \in \operatorname{rr\text{-}pos\text{-}star}) \wedge ((y \in \operatorname{rr\text{-}pos\text{-}star}) \wedge (z \in \operatorname{rr\text{-}pos\text{-}star}))) \Rightarrow  \\ \quad \quad \operatorname{eplus}(\operatorname{eplus}(x, y), z) = \operatorname{eplus}(x, \operatorname{eplus}(y, z)) \end{array}}}
\vnbprop{eplus-le-left}{\vnbstmt{\forall x, y.\; (((x \in \operatorname{rr\text{-}pos\text{-}star}) \wedge (y \in \operatorname{rr\text{-}pos\text{-}star})) \Rightarrow (x \leq \operatorname{eplus}(x, y)))}}
\vnbprop{eplus-mono}{\vnbstmt{\begin{array}{@{}l@{}} \forall x, y, u, v. \\ \quad x \in \operatorname{rr\text{-}pos\text{-}star} \wedge \\ \quad y \in \operatorname{rr\text{-}pos\text{-}star} \wedge \\ \quad u \in \operatorname{rr\text{-}pos\text{-}star} \wedge \\ \quad v \in \operatorname{rr\text{-}pos\text{-}star} \wedge \\ \quad x \leq u \wedge \\ \quad y \leq v \Rightarrow \\ \quad \quad \operatorname{eplus}(x, y) \leq \operatorname{eplus}(u, v) \end{array}}}

\vnbbookfile{theorem-library/rake-esup-defined}{theorem-library/rake-esup-defined}
\noindent{\small theorem-library/rake-esup-defined.scm -- ESUP DEFINED as a definite description, and its four characterising axioms PROVEN from the definition. Assignment "ESUP and ESUM defined" (the user's decision of 2026-09-19). Helper prefix `esd-'.}\par
\vnbprop{esup-exists}{\vnbstmt{\begin{array}{@{}l@{}} \forall s. \\ \quad (s \subseteq \operatorname{rr\text{-}pos\text{-}star}) \Rightarrow  \\ \quad \quad \exists \mathit{bb} \in \operatorname{rr\text{-}pos\text{-}star}. \\ \quad \quad \quad (\forall x \in s.\; (x \leq \mathit{bb}) \wedge \forall b \in \operatorname{rr\text{-}pos\text{-}star}.\; (\forall x \in s.\; (x \leq b) \Rightarrow (\mathit{bb} \leq b))) \end{array}}}
\vnbprop{esup-unique}{\vnbstmt{\begin{array}{@{}l@{}} \forall s, b_{1}, b_{2}. \\ \quad b_{1} \in \operatorname{rr\text{-}pos\text{-}star} \wedge \\ \quad \forall x \in s.\; (x \leq b_{1}) \wedge \\ \quad \forall b \in \operatorname{rr\text{-}pos\text{-}star}.\; (\forall x \in s.\; (x \leq b) \Rightarrow (b_{1} \leq b)) \wedge \\ \quad b_{2} \in \operatorname{rr\text{-}pos\text{-}star} \wedge \\ \quad \forall x \in s.\; (x \leq b_{2}) \wedge \\ \quad \forall b \in \operatorname{rr\text{-}pos\text{-}star}.\; (\forall x \in s.\; (x \leq b) \Rightarrow (b_{2} \leq b)) \Rightarrow \\ \quad \quad b_{1} = b_{2} \end{array}}}
\vnbprop{esup-prop}{\vnbstmt{\begin{array}{@{}l@{}} \forall s. \\ \quad (s \subseteq \operatorname{rr\text{-}pos\text{-}star}) \Rightarrow  \\ \quad \quad \operatorname{esup}(s) \in \operatorname{rr\text{-}pos\text{-}star} \wedge \\ \quad \quad \forall x \in s.\; (x \leq \operatorname{esup}(s)) \wedge \\ \quad \quad \forall b \in \operatorname{rr\text{-}pos\text{-}star}.\; (\forall x \in s.\; (x \leq b) \Rightarrow (\operatorname{esup}(s) \leq b)) \end{array}}}
\vnbprop{esup-in}{\vnbstmt{\forall s.\; ((s \subseteq \operatorname{rr\text{-}pos\text{-}star}) \Rightarrow (\operatorname{esup}(s) \in \operatorname{rr\text{-}pos\text{-}star}))}}
\vnbprop{esup-upper}{\vnbstmt{\forall s.; \forall x \in s.\; ((s \subseteq \operatorname{rr\text{-}pos\text{-}star}) \Rightarrow (x \leq \operatorname{esup}(s)))}}
\vnbprop{esup-least}{\vnbstmt{\begin{array}{@{}l@{}} \forall s.; \forall b \in \operatorname{rr\text{-}pos\text{-}star}. \\ \quad ((s \subseteq \operatorname{rr\text{-}pos\text{-}star}) \wedge \forall x \in s.\; (x \leq b)) \Rightarrow (\operatorname{esup}(s) \leq b) \end{array}}}
\vnbprop{esup-empty}{\vnbstmt{\operatorname{esup}(\emptyset) = 0}}

\vnbbookfile{theorem-library/rake-rr-pos-star-monoid}{theorem-library/rake-rr-pos-star-monoid}
\noindent{\small theorem-library/rake-rr-pos-star-monoid.scm -- IS-COMM-MONOID(RR-POS-STAR-ADD-MONOID), PROVEN.  Rake batch 6-C.  Helper prefix `rpm-'.}\par
\vnbprop{rr-pos-star-is-comm-monoid}{\vnbstmt{\operatorname{is\text{-}comm\text{-}monoid}(\operatorname{rr\text{-}pos\text{-}star\text{-}add\text{-}monoid})}}

\vnbbookfile{theorem-library/hom-laws}{theorem-library/hom-laws}
\noindent{\small A file shared across books; the entries that belong to this one.}\par
\noindent{\small hom-laws.scm -- the models of a structure form a CATEGORY: identity arrows, composition, and the hom-set as a set.}\par
\vnbprop{hom-normed-ag-member-iff}{\vnbstmt{\begin{array}{@{}l@{}} \forall a, b, f. \\ \quad (f \in \operatorname{hom\text{-}normed\text{-}ag}(a, b)) \Leftrightarrow ((f \in (\operatorname{carr}(a) \to \operatorname{carr}(b))) \wedge \operatorname{is\text{-}hom\text{-}normed\text{-}ag}(a, b, f)) \end{array}}}
\vnbprop{hom-normed-ag-id}{\vnbstmt{\forall a.\; (\operatorname{is\text{-}normed\text{-}ag}(a) \Rightarrow \operatorname{is\text{-}hom\text{-}normed\text{-}ag}(a, a, \operatorname{id\text{-}fun}(\operatorname{carr}(a))))}}
\vnbprop{hom-normed-ag-in-set}{\vnbstmt{\begin{array}{@{}l@{}} \forall a, b. \\ \quad (\operatorname{is\text{-}normed\text{-}ag}(a) \wedge \operatorname{is\text{-}normed\text{-}ag}(b)) \Rightarrow (\operatorname{hom\text{-}normed\text{-}ag}(a, b) \in \mathbf{Set}) \end{array}}}
\vnbprop{hom-normed-field-member-iff}{\vnbstmt{\begin{array}{@{}l@{}} \forall a, b, f. \\ \quad (f \in \operatorname{hom\text{-}normed\text{-}field}(a, b)) \Leftrightarrow ((f \in (\operatorname{carr}(a) \to \operatorname{carr}(b))) \wedge \operatorname{is\text{-}hom\text{-}normed\text{-}field}(a, b, f)) \end{array}}}
\vnbprop{hom-normed-field-id}{\vnbstmt{\begin{array}{@{}l@{}} \forall a. \\ \quad \operatorname{is\text{-}normed\text{-}field}(a) \Rightarrow \operatorname{is\text{-}hom\text{-}normed\text{-}field}(a, a, \operatorname{id\text{-}fun}(\operatorname{carr}(a))) \end{array}}}
\vnbprop{hom-normed-field-in-set}{\vnbstmt{\begin{array}{@{}l@{}} \forall a, b. \\ \quad (\operatorname{is\text{-}normed\text{-}field}(a) \wedge \operatorname{is\text{-}normed\text{-}field}(b)) \Rightarrow (\operatorname{hom\text{-}normed\text{-}field}(a, b) \in \mathbf{Set}) \end{array}}}
\vnbprop{hom-normed-vector-space-member-iff}{\vnbstmt{\begin{array}{@{}l@{}} \forall a, b, f. \\ \quad f \in \operatorname{hom\text{-}normed\text{-}vector\text{-}space}(a, b) \\ \quad \Leftrightarrow\; (f \in (\operatorname{vec}(a) \to \operatorname{vec}(b))) \wedge \operatorname{is\text{-}hom\text{-}normed\text{-}vector\text{-}space}(a, b, f) \end{array}}}
\vnbprop{hom-normed-vector-space-id}{\vnbstmt{\begin{array}{@{}l@{}} \forall a. \\ \quad \operatorname{is\text{-}normed\text{-}vector\text{-}space}(a) \Rightarrow  \\ \quad \quad \operatorname{is\text{-}hom\text{-}normed\text{-}vector\text{-}space}(a, a, \operatorname{id\text{-}fun}(\operatorname{vec}(a))) \end{array}}}
\vnbprop{hom-normed-vector-space-in-set}{\vnbstmt{\begin{array}{@{}l@{}} \forall a, b. \\ \quad (\operatorname{is\text{-}normed\text{-}vector\text{-}space}(a) \wedge \operatorname{is\text{-}normed\text{-}vector\text{-}space}(b)) \Rightarrow  \\ \quad \quad \operatorname{hom\text{-}normed\text{-}vector\text{-}space}(a, b) \in \mathbf{Set} \end{array}}}
\vnbprop{hom-complex-inner-product-space-member-iff}{\vnbstmt{\begin{array}{@{}l@{}} \forall a, b, f. \\ \quad f \in \operatorname{hom\text{-}complex\text{-}inner\text{-}product\text{-}space}(a, b) \\ \quad \Leftrightarrow\; (f \in (\operatorname{vec}(a) \to \operatorname{vec}(b))) \wedge \operatorname{is\text{-}hom\text{-}complex\text{-}inner\text{-}product\text{-}space}(a, b, f) \end{array}}}
\vnbprop{hom-complex-inner-product-space-id}{\vnbstmt{\begin{array}{@{}l@{}} \forall a. \\ \quad \operatorname{is\text{-}complex\text{-}inner\text{-}product\text{-}space}(a) \Rightarrow  \\ \quad \quad \operatorname{is\text{-}hom\text{-}complex\text{-}inner\text{-}product\text{-}space}(a, a, \operatorname{id\text{-}fun}(\operatorname{vec}(a))) \end{array}}}
\vnbprop{hom-complex-inner-product-space-in-set}{\vnbstmt{\begin{array}{@{}l@{}} \forall a, b. \\ \quad (\operatorname{is\text{-}complex\text{-}inner\text{-}product\text{-}space}(a) \wedge \operatorname{is\text{-}complex\text{-}inner\text{-}product\text{-}space}(b)) \Rightarrow  \\ \quad \quad \operatorname{hom\text{-}complex\text{-}inner\text{-}product\text{-}space}(a, b) \in \mathbf{Set} \end{array}}}
\vnbprop{hom-metric-space-member-iff}{\vnbstmt{\begin{array}{@{}l@{}} \forall a, b, f. \\ \quad (f \in \operatorname{hom\text{-}metric\text{-}space}(a, b)) \Leftrightarrow ((f \in (\operatorname{pts}(a) \to \operatorname{pts}(b))) \wedge \operatorname{is\text{-}hom\text{-}metric\text{-}space}(a, b, f)) \end{array}}}
\vnbprop{hom-metric-space-id}{\vnbstmt{\begin{array}{@{}l@{}} \forall a. \\ \quad \operatorname{is\text{-}metric\text{-}space}(a) \Rightarrow \operatorname{is\text{-}hom\text{-}metric\text{-}space}(a, a, \operatorname{id\text{-}fun}(\operatorname{pts}(a))) \end{array}}}
\vnbprop{hom-metric-space-in-set}{\vnbstmt{\begin{array}{@{}l@{}} \forall a, b. \\ \quad (\operatorname{is\text{-}metric\text{-}space}(a) \wedge \operatorname{is\text{-}metric\text{-}space}(b)) \Rightarrow (\operatorname{hom\text{-}metric\text{-}space}(a, b) \in \mathbf{Set}) \end{array}}}
\vnbprop{hom-pseudometric-space-member-iff}{\vnbstmt{\begin{array}{@{}l@{}} \forall a, b, f. \\ \quad f \in \operatorname{hom\text{-}pseudometric\text{-}space}(a, b) \\ \quad \Leftrightarrow\; (f \in (\operatorname{pts}(a) \to \operatorname{pts}(b))) \wedge \operatorname{is\text{-}hom\text{-}pseudometric\text{-}space}(a, b, f) \end{array}}}
\vnbprop{hom-pseudometric-space-id}{\vnbstmt{\begin{array}{@{}l@{}} \forall a. \\ \quad \operatorname{is\text{-}pseudometric\text{-}space}(a) \Rightarrow \operatorname{is\text{-}hom\text{-}pseudometric\text{-}space}(a, a, \operatorname{id\text{-}fun}(\operatorname{pts}(a))) \end{array}}}
\vnbprop{hom-pseudometric-space-in-set}{\vnbstmt{\begin{array}{@{}l@{}} \forall a, b. \\ \quad (\operatorname{is\text{-}pseudometric\text{-}space}(a) \wedge \operatorname{is\text{-}pseudometric\text{-}space}(b)) \Rightarrow  \\ \quad \quad \operatorname{hom\text{-}pseudometric\text{-}space}(a, b) \in \mathbf{Set} \end{array}}}
\vnbprop{hom-c-metric-space-member-iff}{\vnbstmt{\begin{array}{@{}l@{}} \forall a, b, f. \\ \quad (f \in \operatorname{hom\text{-}c\text{-}metric\text{-}space}(a, b)) \Leftrightarrow ((f \in (\operatorname{pts}(a) \to \operatorname{pts}(b))) \wedge \operatorname{is\text{-}hom\text{-}c\text{-}metric\text{-}space}(a, b, f)) \end{array}}}
\vnbprop{hom-c-metric-space-id}{\vnbstmt{\begin{array}{@{}l@{}} \forall a. \\ \quad \operatorname{is\text{-}c\text{-}metric\text{-}space}(a) \Rightarrow \operatorname{is\text{-}hom\text{-}c\text{-}metric\text{-}space}(a, a, \operatorname{id\text{-}fun}(\operatorname{pts}(a))) \end{array}}}
\vnbprop{hom-c-metric-space-in-set}{\vnbstmt{\begin{array}{@{}l@{}} \forall a, b. \\ \quad (\operatorname{is\text{-}c\text{-}metric\text{-}space}(a) \wedge \operatorname{is\text{-}c\text{-}metric\text{-}space}(b)) \Rightarrow  \\ \quad \quad \operatorname{hom\text{-}c\text{-}metric\text{-}space}(a, b) \in \mathbf{Set} \end{array}}}
\vnbprop{hom-top-space-member-iff}{\vnbstmt{\begin{array}{@{}l@{}} \forall a, b, f. \\ \quad (f \in \operatorname{hom\text{-}top\text{-}space}(a, b)) \Leftrightarrow ((f \in (\operatorname{pts}(a) \to \operatorname{pts}(b))) \wedge \operatorname{is\text{-}hom\text{-}top\text{-}space}(a, b, f)) \end{array}}}
\vnbprop{hom-top-space-id}{\vnbstmt{\forall s.\; (\operatorname{is\text{-}top\text{-}space}(s) \Rightarrow \operatorname{is\text{-}hom\text{-}top\text{-}space}(s, s, \operatorname{id\text{-}fun}(\operatorname{pts}(s))))}}
\vnbprop{hom-top-space-in-set}{\vnbstmt{\begin{array}{@{}l@{}} \forall a, b. \\ \quad (\operatorname{is\text{-}top\text{-}space}(a) \wedge \operatorname{is\text{-}top\text{-}space}(b)) \Rightarrow (\operatorname{hom\text{-}top\text{-}space}(a, b) \in \mathbf{Set}) \end{array}}}
\vnbprop{hom-metrizable-top-space-member-iff}{\vnbstmt{\begin{array}{@{}l@{}} \forall a, b, f. \\ \quad f \in \operatorname{hom\text{-}metrizable\text{-}top\text{-}space}(a, b) \\ \quad \Leftrightarrow\; (f \in (\operatorname{pts}(a) \to \operatorname{pts}(b))) \wedge \operatorname{is\text{-}hom\text{-}metrizable\text{-}top\text{-}space}(a, b, f) \end{array}}}
\vnbprop{hom-metrizable-top-space-id}{\vnbstmt{\begin{array}{@{}l@{}} \forall s. \\ \quad \operatorname{is\text{-}metrizable\text{-}top\text{-}space}(s) \Rightarrow  \\ \quad \quad \operatorname{is\text{-}hom\text{-}metrizable\text{-}top\text{-}space}(s, s, \operatorname{id\text{-}fun}(\operatorname{pts}(s))) \end{array}}}
\vnbprop{hom-metrizable-top-space-in-set}{\vnbstmt{\begin{array}{@{}l@{}} \forall a, b. \\ \quad (\operatorname{is\text{-}metrizable\text{-}top\text{-}space}(a) \wedge \operatorname{is\text{-}metrizable\text{-}top\text{-}space}(b)) \Rightarrow  \\ \quad \quad \operatorname{hom\text{-}metrizable\text{-}top\text{-}space}(a, b) \in \mathbf{Set} \end{array}}}

\vnbbookfile{theorem-library/hom-functors}{theorem-library/hom-functors}
\noindent{\small A file shared across books; the entries that belong to this one.}\par
\noindent{\small hom-functors.scm -- the Hom functors: Hom(a, -) covariant, Hom(-, b) contravariant, as maps on hom-sets.}\par
\vnbprop{hom-normed-ag-post-type}{\vnbstmt{\begin{array}{@{}l@{}} \forall a, b, c, h. \\ \quad (\operatorname{is\text{-}normed\text{-}ag}(a) \wedge (h \in \operatorname{hom\text{-}normed\text{-}ag}(b, c))) \Rightarrow  \\ \quad \quad (\lambda f \in \operatorname{hom\text{-}normed\text{-}ag}(a, b).\; \operatorname{compose}(h, f)) \in (\operatorname{hom\text{-}normed\text{-}ag}(a, b) \to \operatorname{hom\text{-}normed\text{-}ag}(a, c)) \end{array}}}
\vnbprop{hom-normed-ag-pre-type}{\vnbstmt{\begin{array}{@{}l@{}} \forall a, b, c, k. \\ \quad (\operatorname{is\text{-}normed\text{-}ag}(b) \wedge (k \in \operatorname{hom\text{-}normed\text{-}ag}(c, a))) \Rightarrow  \\ \quad \quad (\lambda f \in \operatorname{hom\text{-}normed\text{-}ag}(a, b).\; \operatorname{compose}(f, k)) \in (\operatorname{hom\text{-}normed\text{-}ag}(a, b) \to \operatorname{hom\text{-}normed\text{-}ag}(c, b)) \end{array}}}
\vnbprop{hom-normed-field-post-type}{\vnbstmt{\begin{array}{@{}l@{}} \forall a, b, c, h. \\ \quad (\operatorname{is\text{-}normed\text{-}field}(a) \wedge (h \in \operatorname{hom\text{-}normed\text{-}field}(b, c))) \Rightarrow  \\ \quad \quad (\lambda f \in \operatorname{hom\text{-}normed\text{-}field}(a, b).\; \operatorname{compose}(h, f)) \\ \quad \quad \in\; (\operatorname{hom\text{-}normed\text{-}field}(a, b) \to \operatorname{hom\text{-}normed\text{-}field}(a, c)) \end{array}}}
\vnbprop{hom-normed-field-pre-type}{\vnbstmt{\begin{array}{@{}l@{}} \forall a, b, c, k. \\ \quad (\operatorname{is\text{-}normed\text{-}field}(b) \wedge (k \in \operatorname{hom\text{-}normed\text{-}field}(c, a))) \Rightarrow  \\ \quad \quad (\lambda f \in \operatorname{hom\text{-}normed\text{-}field}(a, b).\; \operatorname{compose}(f, k)) \\ \quad \quad \in\; (\operatorname{hom\text{-}normed\text{-}field}(a, b) \to \operatorname{hom\text{-}normed\text{-}field}(c, b)) \end{array}}}
\vnbprop{hom-normed-vector-space-post-type}{\vnbstmt{\begin{array}{@{}l@{}} \forall a, b, c, h. \\ \quad (\operatorname{is\text{-}normed\text{-}vector\text{-}space}(a) \wedge (h \in \operatorname{hom\text{-}normed\text{-}vector\text{-}space}(b, c))) \Rightarrow  \\ \quad \quad (\lambda f \in \operatorname{hom\text{-}normed\text{-}vector\text{-}space}(a, b).\; \operatorname{compose}(h, f)) \\ \quad \quad \in\; (\operatorname{hom\text{-}normed\text{-}vector\text{-}space}(a, b) \to \operatorname{hom\text{-}normed\text{-}vector\text{-}space}(a, c)) \end{array}}}
\vnbprop{hom-normed-vector-space-pre-type}{\vnbstmt{\begin{array}{@{}l@{}} \forall a, b, c, k. \\ \quad (\operatorname{is\text{-}normed\text{-}vector\text{-}space}(b) \wedge (k \in \operatorname{hom\text{-}normed\text{-}vector\text{-}space}(c, a))) \Rightarrow  \\ \quad \quad (\lambda f \in \operatorname{hom\text{-}normed\text{-}vector\text{-}space}(a, b).\; \operatorname{compose}(f, k)) \\ \quad \quad \in\; (\operatorname{hom\text{-}normed\text{-}vector\text{-}space}(a, b) \to \operatorname{hom\text{-}normed\text{-}vector\text{-}space}(c, b)) \end{array}}}
\vnbprop{hom-complex-inner-product-space-post-type}{\vnbstmt{\begin{array}{@{}l@{}} \forall a, b, c, h. \\ \quad (\operatorname{is\text{-}complex\text{-}inner\text{-}product\text{-}space}(a) \wedge (h \in \operatorname{hom\text{-}complex\text{-}inner\text{-}product\text{-}space}(b, c))) \Rightarrow  \\ \quad \quad (\lambda f \in \operatorname{hom\text{-}complex\text{-}inner\text{-}product\text{-}space}(a, b).\; \operatorname{compose}(h, f)) \\ \quad \quad \in\; (\operatorname{hom\text{-}complex\text{-}inner\text{-}product\text{-}space}(a, b) \to \operatorname{hom\text{-}complex\text{-}inner\text{-}product\text{-}space}(a, c)) \end{array}}}
\vnbprop{hom-complex-inner-product-space-pre-type}{\vnbstmt{\begin{array}{@{}l@{}} \forall a, b, c, k. \\ \quad (\operatorname{is\text{-}complex\text{-}inner\text{-}product\text{-}space}(b) \wedge (k \in \operatorname{hom\text{-}complex\text{-}inner\text{-}product\text{-}space}(c, a))) \Rightarrow  \\ \quad \quad (\lambda f \in \operatorname{hom\text{-}complex\text{-}inner\text{-}product\text{-}space}(a, b).\; \operatorname{compose}(f, k)) \\ \quad \quad \in\; (\operatorname{hom\text{-}complex\text{-}inner\text{-}product\text{-}space}(a, b) \to \operatorname{hom\text{-}complex\text{-}inner\text{-}product\text{-}space}(c, b)) \end{array}}}
\vnbprop{hom-metric-space-post-type}{\vnbstmt{\begin{array}{@{}l@{}} \forall a, b, c, h. \\ \quad (\operatorname{is\text{-}metric\text{-}space}(a) \wedge (h \in \operatorname{hom\text{-}metric\text{-}space}(b, c))) \Rightarrow  \\ \quad \quad (\lambda f \in \operatorname{hom\text{-}metric\text{-}space}(a, b).\; \operatorname{compose}(h, f)) \\ \quad \quad \in\; (\operatorname{hom\text{-}metric\text{-}space}(a, b) \to \operatorname{hom\text{-}metric\text{-}space}(a, c)) \end{array}}}
\vnbprop{hom-metric-space-pre-type}{\vnbstmt{\begin{array}{@{}l@{}} \forall a, b, c, k. \\ \quad (\operatorname{is\text{-}metric\text{-}space}(b) \wedge (k \in \operatorname{hom\text{-}metric\text{-}space}(c, a))) \Rightarrow  \\ \quad \quad (\lambda f \in \operatorname{hom\text{-}metric\text{-}space}(a, b).\; \operatorname{compose}(f, k)) \\ \quad \quad \in\; (\operatorname{hom\text{-}metric\text{-}space}(a, b) \to \operatorname{hom\text{-}metric\text{-}space}(c, b)) \end{array}}}
\vnbprop{hom-pseudometric-space-post-type}{\vnbstmt{\begin{array}{@{}l@{}} \forall a, b, c, h. \\ \quad (\operatorname{is\text{-}pseudometric\text{-}space}(a) \wedge (h \in \operatorname{hom\text{-}pseudometric\text{-}space}(b, c))) \Rightarrow  \\ \quad \quad (\lambda f \in \operatorname{hom\text{-}pseudometric\text{-}space}(a, b).\; \operatorname{compose}(h, f)) \\ \quad \quad \in\; (\operatorname{hom\text{-}pseudometric\text{-}space}(a, b) \to \operatorname{hom\text{-}pseudometric\text{-}space}(a, c)) \end{array}}}
\vnbprop{hom-pseudometric-space-pre-type}{\vnbstmt{\begin{array}{@{}l@{}} \forall a, b, c, k. \\ \quad (\operatorname{is\text{-}pseudometric\text{-}space}(b) \wedge (k \in \operatorname{hom\text{-}pseudometric\text{-}space}(c, a))) \Rightarrow  \\ \quad \quad (\lambda f \in \operatorname{hom\text{-}pseudometric\text{-}space}(a, b).\; \operatorname{compose}(f, k)) \\ \quad \quad \in\; (\operatorname{hom\text{-}pseudometric\text{-}space}(a, b) \to \operatorname{hom\text{-}pseudometric\text{-}space}(c, b)) \end{array}}}
\vnbprop{hom-c-metric-space-post-type}{\vnbstmt{\begin{array}{@{}l@{}} \forall a, b, c, h. \\ \quad (\operatorname{is\text{-}c\text{-}metric\text{-}space}(a) \wedge (h \in \operatorname{hom\text{-}c\text{-}metric\text{-}space}(b, c))) \Rightarrow  \\ \quad \quad (\lambda f \in \operatorname{hom\text{-}c\text{-}metric\text{-}space}(a, b).\; \operatorname{compose}(h, f)) \\ \quad \quad \in\; (\operatorname{hom\text{-}c\text{-}metric\text{-}space}(a, b) \to \operatorname{hom\text{-}c\text{-}metric\text{-}space}(a, c)) \end{array}}}
\vnbprop{hom-c-metric-space-pre-type}{\vnbstmt{\begin{array}{@{}l@{}} \forall a, b, c, k. \\ \quad (\operatorname{is\text{-}c\text{-}metric\text{-}space}(b) \wedge (k \in \operatorname{hom\text{-}c\text{-}metric\text{-}space}(c, a))) \Rightarrow  \\ \quad \quad (\lambda f \in \operatorname{hom\text{-}c\text{-}metric\text{-}space}(a, b).\; \operatorname{compose}(f, k)) \\ \quad \quad \in\; (\operatorname{hom\text{-}c\text{-}metric\text{-}space}(a, b) \to \operatorname{hom\text{-}c\text{-}metric\text{-}space}(c, b)) \end{array}}}
\vnbprop{hom-top-space-post-type}{\vnbstmt{\begin{array}{@{}l@{}} \forall a, b, c, h. \\ \quad (\operatorname{is\text{-}top\text{-}space}(a) \wedge (h \in \operatorname{hom\text{-}top\text{-}space}(b, c))) \Rightarrow  \\ \quad \quad (\lambda f \in \operatorname{hom\text{-}top\text{-}space}(a, b).\; \operatorname{compose}(h, f)) \in (\operatorname{hom\text{-}top\text{-}space}(a, b) \to \operatorname{hom\text{-}top\text{-}space}(a, c)) \end{array}}}
\vnbprop{hom-top-space-pre-type}{\vnbstmt{\begin{array}{@{}l@{}} \forall a, b, c, k. \\ \quad (\operatorname{is\text{-}top\text{-}space}(b) \wedge (k \in \operatorname{hom\text{-}top\text{-}space}(c, a))) \Rightarrow  \\ \quad \quad (\lambda f \in \operatorname{hom\text{-}top\text{-}space}(a, b).\; \operatorname{compose}(f, k)) \in (\operatorname{hom\text{-}top\text{-}space}(a, b) \to \operatorname{hom\text{-}top\text{-}space}(c, b)) \end{array}}}
\vnbprop{hom-metrizable-top-space-post-type}{\vnbstmt{\begin{array}{@{}l@{}} \forall a, b, c, h. \\ \quad (\operatorname{is\text{-}metrizable\text{-}top\text{-}space}(a) \wedge (h \in \operatorname{hom\text{-}metrizable\text{-}top\text{-}space}(b, c))) \Rightarrow  \\ \quad \quad (\lambda f \in \operatorname{hom\text{-}metrizable\text{-}top\text{-}space}(a, b).\; \operatorname{compose}(h, f)) \\ \quad \quad \in\; (\operatorname{hom\text{-}metrizable\text{-}top\text{-}space}(a, b) \to \operatorname{hom\text{-}metrizable\text{-}top\text{-}space}(a, c)) \end{array}}}
\vnbprop{hom-metrizable-top-space-pre-type}{\vnbstmt{\begin{array}{@{}l@{}} \forall a, b, c, k. \\ \quad (\operatorname{is\text{-}metrizable\text{-}top\text{-}space}(b) \wedge (k \in \operatorname{hom\text{-}metrizable\text{-}top\text{-}space}(c, a))) \Rightarrow  \\ \quad \quad (\lambda f \in \operatorname{hom\text{-}metrizable\text{-}top\text{-}space}(a, b).\; \operatorname{compose}(f, k)) \\ \quad \quad \in\; (\operatorname{hom\text{-}metrizable\text{-}top\text{-}space}(a, b) \to \operatorname{hom\text{-}metrizable\text{-}top\text{-}space}(c, b)) \end{array}}}

\vnbbookfile{theorem-library/rake-esum-defined}{theorem-library/rake-esum-defined}
\noindent{\small theorem-library/rake-esum-defined.scm -- ESUM DEFINED as the ESUP of the set of finite partial sums, and its three characterising axioms PROVEN from the definition.  Rake batch 6-C, the ESUP half of which is theorem-library/rake-esup-defined.scm (agent 6-B).  Helper prefix `esm-'.}\par
\vnbprop{esum-partial-sum-type}{\vnbstmt{\begin{array}{@{}l@{}} \forall f, s. \\ \quad ((f \in (\operatorname{dom}(f) \to \operatorname{rr\text{-}pos\text{-}star})) \wedge ((s \in \mathbf{Set}) \wedge ((|s| \in \mathbb{N}) \wedge (s \subseteq \operatorname{dom}(f))))) \Rightarrow  \\ \quad \quad \operatorname{finsum}(\operatorname{rr\text{-}pos\text{-}star\text{-}add\text{-}monoid}, f, s) \in \operatorname{rr\text{-}pos\text{-}star} \end{array}}}
\vnbprop{esum-in}{\vnbstmt{\forall f \in (\operatorname{dom}(f) \to \operatorname{rr\text{-}pos\text{-}star}).\; (\operatorname{esum}(f) \in \operatorname{rr\text{-}pos\text{-}star})}}
\vnbprop{esum-upper}{\vnbstmt{\begin{array}{@{}l@{}} \forall f \in (\operatorname{dom}(f) \to \operatorname{rr\text{-}pos\text{-}star}), s \in \mathbf{Set}. \\ \quad ((|s| \in \mathbb{N}) \wedge (s \subseteq \operatorname{dom}(f))) \Rightarrow  \\ \quad \quad \operatorname{finsum}(\operatorname{rr\text{-}pos\text{-}star\text{-}add\text{-}monoid}, f, s) \leq \operatorname{esum}(f) \end{array}}}
\vnbprop{esum-least}{\vnbstmt{\begin{array}{@{}l@{}} \forall f \in (\operatorname{dom}(f) \to \operatorname{rr\text{-}pos\text{-}star}), b \in \operatorname{rr\text{-}pos\text{-}star}. \\ \quad \forall s \in \mathbf{Set}. \\ \quad \quad ((|s| \in \mathbb{N}) \wedge (s \subseteq \operatorname{dom}(f))) \Rightarrow  \\ \quad \quad \quad \operatorname{finsum}(\operatorname{rr\text{-}pos\text{-}star\text{-}add\text{-}monoid}, f, s) \leq b \Rightarrow \\ \quad \quad \operatorname{esum}(f) \leq b \end{array}}}

\vnbbookfile{theorem-library/rake-esum-finite}{theorem-library/rake-esum-finite}
\noindent{\small rake-esum-finite.scm -- ESUM: the finiteness dichotomy.}\par
\vnbprop{esum-bounded-implies-finite}{\vnbstmt{\begin{array}{@{}l@{}} \forall f \in (\operatorname{dom}(f) \to \operatorname{rr\text{-}pos\text{-}star}). \\ \quad \exists m \in \mathbb{R}. \forall s \in \mathbf{Set}. \\ \quad \quad ((|s| \in \mathbb{N}) \wedge (s \subseteq \operatorname{dom}(f))) \Rightarrow  \\ \quad \quad \quad \operatorname{finsum}(\operatorname{rr\text{-}pos\text{-}star\text{-}add\text{-}monoid}, f, s) \leq m \Rightarrow \\ \quad \quad \operatorname{esum}(f) \in \mathbb{R} \end{array}}}
\vnbprop{esum-finite-iff-bounded}{\vnbstmt{\begin{array}{@{}l@{}} \forall f \in (\operatorname{dom}(f) \to \operatorname{rr\text{-}pos\text{-}star}). \\ \quad \operatorname{esum}(f) \in \mathbb{R} \\ \quad \Leftrightarrow\; \exists m \in \mathbb{R}. \forall s \in \mathbf{Set}. \\ \quad \quad ((|s| \in \mathbb{N}) \wedge (s \subseteq \operatorname{dom}(f))) \Rightarrow  \\ \quad \quad \quad \operatorname{finsum}(\operatorname{rr\text{-}pos\text{-}star\text{-}add\text{-}monoid}, f, s) \leq m \end{array}}}

\vnbbookfile{theorem-library/ball-cover-lemmas}{theorem-library/ball-cover-lemmas}
\noindent{\small ball-cover-lemmas.scm -- the membership plumbing under finite-ball-subcover-r-net, PROVEN (2026-09-14).}\par
\vnbprop{centres-unfold}{\vnbstmt{\forall s, b, r.\; (\operatorname{centres}(s, b, r) \simeq \{c \in \operatorname{pts}(s) : (\operatorname{ball}(s, c, r) = b)\})}}
\vnbprop{ball-cover-unfold}{\vnbstmt{\begin{array}{@{}l@{}} \forall s, r. \\ \quad (\operatorname{ball\text{-}cover}(s, r) \simeq \operatorname{image}((\lambda c \in \operatorname{pts}(s).\; \operatorname{ball}(s, c, r)), \operatorname{pts}(s))) \end{array}}}
\vnbprop{centre-set-unfold}{\vnbstmt{\begin{array}{@{}l@{}} \forall s, r, f. \\ \quad (\operatorname{centre\text{-}set}(s, r, f) \simeq \operatorname{image}((\lambda b \in f.\; \varepsilon(\operatorname{centres}(s, b, r))), f)) \end{array}}}
\vnbprop{ball-sep-unfold}{\vnbstmt{\begin{array}{@{}l@{}} \forall s, c, r. \\ \quad (\operatorname{ball}(s, c, r) \simeq \{y \in \operatorname{pts}(s) : ((\operatorname{dist}(s)(c, y) \leq r) \wedge \neg (\operatorname{dist}(s)(c, y) = r))\}) \end{array}}}
\vnbprop{centres-mem-build}{\vnbstmt{\forall s, b, r.; \forall c \in \operatorname{pts}(s).\; ((\operatorname{ball}(s, c, r) = b) \Rightarrow (c \in \operatorname{centres}(s, b, r)))}}
\vnbprop{centres-in-carrier}{\vnbstmt{\forall s, b, r.; \forall c \in \operatorname{centres}(s, b, r).\; (c \in \operatorname{pts}(s))}}
\vnbprop{centres-ball-eq}{\vnbstmt{\forall s, b, r.; \forall c \in \operatorname{centres}(s, b, r).\; (\operatorname{ball}(s, c, r) = b)}}
\vnbprop{ball-cover-mem-fwd}{\vnbstmt{\forall s, r.; \forall u \in \operatorname{ball\text{-}cover}(s, r).\; \exists c \in \operatorname{pts}(s).\; (\operatorname{ball}(s, c, r) = u)}}
\vnbprop{open-cover-covers-point}{\vnbstmt{\begin{array}{@{}l@{}} \forall s, f.; \forall p \in \operatorname{pts}(s). \\ \quad \operatorname{is\text{-}open\text{-}cover}(s, f) \Rightarrow \exists u \in f.\; (p \in u) \end{array}}}
\vnbprop{ball-point-le}{\vnbstmt{\forall s, c, r, w.; \forall p \in w.\; ((\operatorname{ball}(s, c, r) = w) \Rightarrow (\operatorname{dist}(s)(c, p) \leq r))}}
\vnbprop{ball-point-ne}{\vnbstmt{\forall s, c, r, w.; \forall p \in w.\; ((\operatorname{ball}(s, c, r) = w) \Rightarrow \neg (\operatorname{dist}(s)(c, p) = r))}}
\vnbprop{centre-set-finite-guarded}{\vnbstmt{\forall s, r.; \forall f \in \mathbf{Set}.\; ((|f| \in \mathbb{N}) \Rightarrow (|\operatorname{centre\text{-}set}(s, r, f)| \in \mathbb{N}))}}

\vnbbookfile{structure-library/metric-laws}{structure-library/metric-laws}
\noindent{\small metric-laws.scm -- the five metric laws, PROVEN from the definition.}\par
\vnbprop{metric-pos}{\vnbstmt{\begin{array}{@{}l@{}} \forall s.; \forall x, y \in \operatorname{pts}(s). \\ \quad \operatorname{is\text{-}metric\text{-}space}(s) \Rightarrow (0 \leq \operatorname{dist}(s)(x, y)) \end{array}}}
\vnbprop{metric-self-zero}{\vnbstmt{\forall s.; \forall x \in \operatorname{pts}(s).\; (\operatorname{is\text{-}metric\text{-}space}(s) \Rightarrow (\operatorname{dist}(s)(x, x) = 0))}}
\vnbprop{metric-zero-eq}{\vnbstmt{\begin{array}{@{}l@{}} \forall s.; \forall x, y \in \operatorname{pts}(s). \\ \quad (\operatorname{is\text{-}metric\text{-}space}(s) \wedge (\operatorname{dist}(s)(x, y) = 0)) \Rightarrow (x = y) \end{array}}}
\vnbprop{metric-sym}{\vnbstmt{\begin{array}{@{}l@{}} \forall s.; \forall x, y \in \operatorname{pts}(s). \\ \quad \operatorname{is\text{-}metric\text{-}space}(s) \Rightarrow (\operatorname{dist}(s)(x, y) = \operatorname{dist}(s)(y, x)) \end{array}}}
\vnbprop{metric-triangle}{\vnbstmt{\begin{array}{@{}l@{}} \forall s.; \forall x, y, z \in \operatorname{pts}(s). \\ \quad \operatorname{is\text{-}metric\text{-}space}(s) \Rightarrow (\operatorname{dist}(s)(x, z) \leq (\operatorname{dist}(s)(x, y) + \operatorname{dist}(s)(y, z))) \end{array}}}

\vnbbookfile{theorem-library/ball-is-open}{theorem-library/ball-is-open}
\noindent{\small ball-is-open.scm -- an open ball is an open set.}\par
\vnbprop{ball-is-open}{\vnbstmt{\begin{array}{@{}l@{}} \forall s.; \forall x \in \operatorname{pts}(s), r \in \mathbb{R}. \\ \quad (\operatorname{is\text{-}metric\text{-}space}(s) \wedge ((0 \leq r) \wedge \neg (0 = r))) \Rightarrow \operatorname{is\text{-}open}(s, \operatorname{ball}(s, x, r)) \end{array}}}

\vnbbookfile{theorem-library/rake-compact-complete}{theorem-library/rake-compact-complete}
\noindent{\small rake-compact-complete.scm -- two of the three asserted leaves on the bill of compact-implies-complete (calculus/compact-complete-proof.scm), PROVEN.}\par
\vnbprop{cauchy-seq-is-fun}{\vnbstmt{\forall s, f.\; (\operatorname{is\text{-}cauchy\text{-}seq}(s, f) \Rightarrow (f \in (\mathbb{N} \to \operatorname{pts}(s))))}}
\vnbprop{cauchy-cluster-converges}{\vnbstmt{\begin{array}{@{}l@{}} \forall s, f. \\ \quad (\operatorname{is\text{-}cauchy\text{-}seq}(s, f) \wedge \exists x.\; \operatorname{cluster\text{-}point}(s, f, x)) \Rightarrow  \\ \quad \quad \operatorname{converges}(s, f) \end{array}}}

\vnbbookfile{theorem-library/rake-compact-cluster}{theorem-library/rake-compact-cluster}
\noindent{\small rake-compact-cluster.scm -- a compact metric space is sequentially compact.}\par
\vnbprop{compact-seq-has-cluster}{\vnbstmt{\begin{array}{@{}l@{}} \forall s, f. \\ \quad (\operatorname{is\text{-}compact}(s) \wedge (f \in (\mathbb{N} \to \operatorname{pts}(s)))) \Rightarrow  \\ \quad \quad \exists x.\; \operatorname{cluster\text{-}point}(s, f, x) \end{array}}}

\vnbbookfile{theorem-library/rake-balls}{theorem-library/rake-balls}
\noindent{\small rake-balls.scm -- the open-ball read-offs of structure-library/metric-topology.scm, PROVEN.  Five facts that stood there as `support' + `warrant! ... 'proof' -- a tier that claims a machine-checked proof exists -- and had none:}\par
\vnbprop{ball-membership}{\vnbstmt{\begin{array}{@{}l@{}} \forall s, x, r, y. \\ \quad (y \in \operatorname{ball}(s, x, r)) \Leftrightarrow ((y \in \operatorname{pts}(s)) \wedge ((\operatorname{dist}(s)(x, y) \leq r) \wedge \neg (\operatorname{dist}(s)(x, y) = r))) \end{array}}}
\vnbprop{ball-mem-from-le}{\vnbstmt{\begin{array}{@{}l@{}} \forall s.; \forall x \in \operatorname{pts}(s), y, d, r. \\ \quad (\operatorname{is\text{-}metric\text{-}space}(s) \wedge ((y \in \operatorname{pts}(s)) \wedge ((d \in \mathbb{R}) \wedge ((r \in \mathbb{R}) \wedge ((\operatorname{dist}(s)(x, y) \leq d) \wedge (d < r)))))) \Rightarrow  \\ \quad \quad y \in \operatorname{ball}(s, x, r) \end{array}}}
\vnbprop{ball-is-set}{\vnbstmt{\forall s, x, r.\; (\operatorname{is\text{-}metric\text{-}space}(s) \Rightarrow (\operatorname{ball}(s, x, r) \in \mathbf{Set}))}}
\vnbprop{ball-center-in}{\vnbstmt{\begin{array}{@{}l@{}} \forall s.; \forall x \in \operatorname{pts}(s), r \in \mathbb{R}. \\ \quad (\operatorname{is\text{-}metric\text{-}space}(s) \wedge ((0 \leq r) \wedge \neg (0 = r))) \Rightarrow (x \in \operatorname{ball}(s, x, r)) \end{array}}}
\vnbprop{ball-2r-triangle}{\vnbstmt{\begin{array}{@{}l@{}} \forall s.; \forall x \in \operatorname{pts}(s), r \in \mathbb{R}, y, z \in \operatorname{ball}(s, x, r). \\ \quad (\operatorname{is\text{-}metric\text{-}space}(s) \wedge ((0 \leq r) \wedge \neg (0 = r))) \Rightarrow  \\ \quad \quad (\operatorname{dist}(s)(y, z) \leq (r + r)) \wedge \neg (\operatorname{dist}(s)(y, z) = (r + r)) \end{array}}}
\vnbprop{cauchy-block-estimate}{\vnbstmt{\begin{array}{@{}l@{}} \forall s.; \forall c \in \operatorname{pts}(s), r, u.; \forall y, z \in u, d.; \forall a \in \mathbb{R}. \\ \quad (\operatorname{is\text{-}metric\text{-}space}(s) \wedge (\operatorname{pos\text{-}rr}(r) \wedge ((u = \operatorname{ball}(s, c, r)) \wedge (\operatorname{pos\text{-}rr}(d) \wedge ((r \leq d) \wedge ((d + d) = a)))))) \Rightarrow  \\ \quad \quad \operatorname{dist}(s)(y, z) \leq a \end{array}}}

\vnbbookfile{theorem-library/metric-connected-laws}{theorem-library/metric-connected-laws}
\noindent{\small metric-connected-laws.scm -- A LOCALLY CONSTANT FUNCTION ON A CONNECTED SET IS CONSTANT.  Agent CA-5 of week 3 of the October roadmap, 2026-10-04. Definition: structure-library/metric-connected.scm (IS-CONNECTED, new this batch).}\par
\vnbprop{connected-locally-constant}{\vnbstmt{\begin{array}{@{}l@{}} \forall s, a, f.; \forall p, q \in a. \\ \quad \operatorname{is\text{-}connected}(s, a) \wedge \\ \quad \forall x \in a. \exists r. \\ \quad \quad \operatorname{pos\text{-}rr}(r) \wedge \forall y \in a.\; ((y \in \operatorname{ball}(s, x, r)) \Rightarrow (f(y) = f(x))) \Rightarrow \\ \quad \quad f(p) = f(q) \end{array}}}

\vnbbookfile{theorem-library/metric-subspace-laws}{theorem-library/metric-subspace-laws}
\noindent{\small metric-subspace-laws.scm -- the laws of SUBSPACE-MS and RESTRICT (structure-library/metric-subspace.scm).  Everything here is PROVEN; the definition file asserts nothing.}\par
\vnbprop{subspace-pts}{\vnbstmt{\forall s.; \forall a \in \mathbf{Set}.\; (\operatorname{pts}(\operatorname{subspace\text{-}ms}(s, a)) \simeq a)}}
\vnbprop{subspace-dist-slot}{\vnbstmt{\begin{array}{@{}l@{}} \forall s.; \forall a \in \mathbf{Set}. \\ \quad (\operatorname{dist}(\operatorname{subspace\text{-}ms}(s, a)) \simeq (\lambda \langle \operatorname{sbu\_}, \operatorname{sbv\_} \rangle \in (a \times a).\; \operatorname{dist}(s)(\operatorname{sbu\_}, \operatorname{sbv\_}))) \end{array}}}
\vnbprop{subspace-dist}{\vnbstmt{\begin{array}{@{}l@{}} \forall s.; \forall a \in \mathbf{Set}, u, v. \\ \quad ((u \in a) \wedge (v \in a)) \Rightarrow (\operatorname{dist}(\operatorname{subspace\text{-}ms}(s, a))(u, v) \simeq \operatorname{dist}(s)(u, v)) \end{array}}}
\vnbprop{subspace-is-metric-space}{\vnbstmt{\begin{array}{@{}l@{}} \forall s, a. \\ \quad (\operatorname{is\text{-}metric\text{-}space}(s) \wedge (a \subseteq \operatorname{pts}(s))) \Rightarrow  \\ \quad \quad \operatorname{is\text{-}metric\text{-}space}(\operatorname{subspace\text{-}ms}(s, a)) \end{array}}}
\vnbprop{restrict-apply}{\vnbstmt{\forall f, a.; \forall u \in a.\; (\operatorname{restrict}(f, a)(u) \simeq f(u))}}
\vnbprop{restrict-in-fun}{\vnbstmt{\begin{array}{@{}l@{}} \forall f, m, d, a. \\ \quad ((f \in (m \to d)) \wedge (a \subseteq m)) \Rightarrow (\operatorname{restrict}(f, a) \in (a \to d)) \end{array}}}
\vnbprop{subspace-whole}{\vnbstmt{\begin{array}{@{}l@{}} \forall s. \\ \quad \operatorname{is\text{-}metric\text{-}space}(s) \Rightarrow  \\ \quad \quad (\operatorname{pts}(\operatorname{subspace\text{-}ms}(s, \operatorname{pts}(s))) \simeq \operatorname{pts}(s)) \wedge \\ \quad \quad \forall u, v. \\ \quad \quad \quad ((u \in \operatorname{pts}(s)) \wedge (v \in \operatorname{pts}(s))) \Rightarrow  \\ \quad \quad \quad \quad (\operatorname{dist}(\operatorname{subspace\text{-}ms}(s, \operatorname{pts}(s)))(u, v) \simeq \operatorname{dist}(s)(u, v)) \end{array}}}
\vnbprop{subspace-ball-membership}{\vnbstmt{\begin{array}{@{}l@{}} \forall s, a.; \forall c \in a, r, y. \\ \quad (\operatorname{is\text{-}metric\text{-}space}(s) \wedge (a \subseteq \operatorname{pts}(s))) \Rightarrow  \\ \quad \quad (y \in \operatorname{ball}(\operatorname{subspace\text{-}ms}(s, a), c, r)) \Leftrightarrow ((y \in \operatorname{ball}(s, c, r)) \wedge (y \in a)) \end{array}}}
\vnbprop{subspace-ball}{\vnbstmt{\begin{array}{@{}l@{}} \forall s, a.; \forall c \in a, r. \\ \quad (\operatorname{is\text{-}metric\text{-}space}(s) \wedge (a \subseteq \operatorname{pts}(s))) \Rightarrow  \\ \quad \quad \operatorname{ball}(\operatorname{subspace\text{-}ms}(s, a), c, r) = (\operatorname{ball}(s, c, r) \cap a) \end{array}}}
\vnbprop{subspace-open-trace}{\vnbstmt{\begin{array}{@{}l@{}} \forall s, a, u. \\ \quad (\operatorname{is\text{-}metric\text{-}space}(s) \wedge ((a \subseteq \operatorname{pts}(s)) \wedge \operatorname{is\text{-}open}(s, u))) \Rightarrow  \\ \quad \quad \operatorname{is\text{-}open}(\operatorname{subspace\text{-}ms}(s, a), (u \cap a)) \end{array}}}
\vnbprop{restrict-continuous-at}{\vnbstmt{\begin{array}{@{}l@{}} \forall s, a, t, f.; \forall \mathit{pa} \in a. \\ \quad (\operatorname{is\text{-}metric\text{-}space}(s) \wedge ((a \subseteq \operatorname{pts}(s)) \wedge \operatorname{is\text{-}continuous\text{-}at}(s, t, f, \mathit{pa}))) \Rightarrow  \\ \quad \quad \operatorname{is\text{-}continuous\text{-}at}(\operatorname{subspace\text{-}ms}(s, a), t, \operatorname{restrict}(f, a), \mathit{pa}) \end{array}}}
\vnbprop{subspace-converges-to}{\vnbstmt{\begin{array}{@{}l@{}} \forall s, a.; \forall g \in (\mathbb{N} \to a), l \in a. \\ \quad (\operatorname{is\text{-}metric\text{-}space}(s) \wedge (a \subseteq \operatorname{pts}(s))) \Rightarrow  \\ \quad \quad \operatorname{converges\text{-}to}(\operatorname{subspace\text{-}ms}(s, a), g, l) \Leftrightarrow \operatorname{converges\text{-}to}(s, g, l) \end{array}}}
\vnbprop{subspace-open-is-trace}{\vnbstmt{\begin{array}{@{}l@{}} \forall s, a, \mathit{vs}. \\ \quad (\operatorname{is\text{-}metric\text{-}space}(s) \wedge ((a \subseteq \operatorname{pts}(s)) \wedge \operatorname{is\text{-}open}(\operatorname{subspace\text{-}ms}(s, a), \mathit{vs}))) \Rightarrow  \\ \quad \quad \exists \mathit{us}.\; (\operatorname{is\text{-}open}(s, \mathit{us}) \wedge (\mathit{vs} = (\mathit{us} \cap a))) \end{array}}}
\vnbprop{closed-limit-in}{\vnbstmt{\begin{array}{@{}l@{}} \forall s, a.; \forall f \in (\mathbb{N} \to a), l. \\ \quad (\operatorname{is\text{-}metric\text{-}space}(s) \wedge (\operatorname{is\text{-}closed}(s, a) \wedge \operatorname{converges\text{-}to}(s, f, l))) \Rightarrow (l \in a) \end{array}}}
\vnbprop{subspace-is-cauchy-seq}{\vnbstmt{\begin{array}{@{}l@{}} \forall s, a.; \forall f \in (\mathbb{N} \to a). \\ \quad (\operatorname{is\text{-}metric\text{-}space}(s) \wedge (a \subseteq \operatorname{pts}(s))) \Rightarrow  \\ \quad \quad \operatorname{is\text{-}cauchy\text{-}seq}(\operatorname{subspace\text{-}ms}(s, a), f) \Leftrightarrow \operatorname{is\text{-}cauchy\text{-}seq}(s, f) \end{array}}}
\vnbprop{subspace-complete}{\vnbstmt{\begin{array}{@{}l@{}} \forall s, a. \\ \quad (\operatorname{is\text{-}complete}(s) \wedge \operatorname{is\text{-}closed}(s, a)) \Rightarrow \operatorname{is\text{-}complete}(\operatorname{subspace\text{-}ms}(s, a)) \end{array}}}

\vnbbookfile{theorem-library/compact-subspace}{theorem-library/compact-subspace}
\noindent{\small compact-subspace.scm -- a CLOSED SUBSET of a COMPACT space is COMPACT as a subspace (structure-library/metric-subspace.scm, theorem-library/metric-subspace-laws.scm).}\par
\vnbprop{open-in-power}{\vnbstmt{\forall s, u.\; ((\operatorname{is\text{-}metric\text{-}space}(s) \wedge \operatorname{is\text{-}open}(s, u)) \Rightarrow (u \in \operatorname{power}(\operatorname{pts}(s))))}}
\vnbprop{closed-subset-of-compact-is-compact}{\vnbstmt{\begin{array}{@{}l@{}} \forall s, a. \\ \quad (\operatorname{is\text{-}compact}(s) \wedge \operatorname{is\text{-}closed}(s, a)) \Rightarrow \operatorname{is\text{-}compact}(\operatorname{subspace\text{-}ms}(s, a)) \end{array}}}

\vnbbookfile{theorem-library/rake-analysis-typing}{theorem-library/rake-analysis-typing}
\noindent{\small rake-analysis-typing.scm -- BATCH H of the leaf rake (2026-09-17): the four ANALYSIS typing leaves.  Two are proven as stated; two are UNDERDETERMINED as stated -- the same defect that made `nth-deriv-v-in-vec' false -- and are proven here with the guard their scalar twins already carry.}\par
\vnbprop{rkt-cseq-is-set}{\vnbstmt{\forall m.\; (\operatorname{is\text{-}metric\text{-}space}(m) \Rightarrow (\operatorname{cseq}(m) \in \mathbf{Set}))}}
\vnbprop{rkt-completion-pts}{\vnbstmt{\begin{array}{@{}l@{}} \forall m. \\ \quad \operatorname{is\text{-}metric\text{-}space}(m) \Rightarrow (\operatorname{pts}(\operatorname{completion}(m)) \simeq \operatorname{quotient}(\operatorname{cauchy\text{-}setoid}(m))) \end{array}}}
\vnbprop{rkt-cauchy-setoid-pts}{\vnbstmt{\forall m.\; (\operatorname{is\text{-}metric\text{-}space}(m) \Rightarrow (\operatorname{pts}(\operatorname{cauchy\text{-}setoid}(m)) \simeq \operatorname{cseq}(m)))}}
\vnbprop{rkt-class-is-set}{\vnbstmt{\forall m, w.\; (\operatorname{is\text{-}metric\text{-}space}(m) \Rightarrow (\operatorname{class}(\operatorname{cauchy\text{-}setoid}(m), w) \in \mathbf{Set}))}}
\vnbprop{rkt-embed-seq-in-fun}{\vnbstmt{\begin{array}{@{}l@{}} \forall m, u. \\ \quad (\operatorname{is\text{-}metric\text{-}space}(m) \wedge (u \in \operatorname{pts}(m))) \Rightarrow (\operatorname{embed\text{-}seq}(m, u) \in (\mathbb{N} \to \operatorname{pts}(m))) \end{array}}}
\vnbprop{rkt-embed-seq-cauchy}{\vnbstmt{\begin{array}{@{}l@{}} \forall m, u. \\ \quad (\operatorname{is\text{-}metric\text{-}space}(m) \wedge (u \in \operatorname{pts}(m))) \Rightarrow \operatorname{is\text{-}cauchy\text{-}seq}(m, \operatorname{embed\text{-}seq}(m, u)) \end{array}}}
\vnbprop{rkt-embed-seq-in-cseq}{\vnbstmt{\begin{array}{@{}l@{}} \forall m, u. \\ \quad (\operatorname{is\text{-}metric\text{-}space}(m) \wedge (u \in \operatorname{pts}(m))) \Rightarrow (\operatorname{embed\text{-}seq}(m, u) \in \operatorname{cseq}(m)) \end{array}}}
\vnbprop{embed-in-fun}{\vnbstmt{\begin{array}{@{}l@{}} \forall m. \\ \quad \operatorname{is\text{-}metric\text{-}space}(m) \Rightarrow (\operatorname{embed}(m) \in (\operatorname{pts}(m) \to \operatorname{pts}(\operatorname{completion}(m)))) \end{array}}}

\vnbbookfile{theorem-library/rake-cont-preimage}{theorem-library/rake-cont-preimage}
\noindent{\small rake-cont-preimage.scm -- continuous-implies-open-preimage, PROVEN.}\par
\vnbprop{continuous-implies-open-preimage}{\vnbstmt{\begin{array}{@{}l@{}} \forall s, t, f, v. \\ \quad (\operatorname{is\text{-}continuous}(s, t, f) \wedge \operatorname{is\text{-}open}(t, v)) \Rightarrow \operatorname{is\text{-}open}(s, \operatorname{preimage}(s, f, v)) \end{array}}}

\vnbbookfile{theorem-library/pseudometric-laws}{theorem-library/pseudometric-laws}
\noindent{\small pseudometric-laws.scm -- the FOUR PSEUDOMETRIC LAWS, proven by projecting the `is-pseudometric' property, exactly as structure-library/metric-laws.scm projects `is-metric'.}\par
\vnbprop{pseudometric-self-zero}{\vnbstmt{\forall t, r.; \forall u \in r.\; (\operatorname{is\text{-}pseudometric}(t, r) \Rightarrow (t(u, u) = 0))}}
\vnbprop{pseudometric-pos}{\vnbstmt{\forall t, r.; \forall u, v \in r.\; (\operatorname{is\text{-}pseudometric}(t, r) \Rightarrow (0 \leq t(u, v)))}}
\vnbprop{pseudometric-sym}{\vnbstmt{\forall t, r.; \forall u, v \in r.\; (\operatorname{is\text{-}pseudometric}(t, r) \Rightarrow (t(u, v) = t(v, u)))}}
\vnbprop{pseudometric-triangle}{\vnbstmt{\begin{array}{@{}l@{}} \forall t, r.; \forall u, v, w \in r. \\ \quad \operatorname{is\text{-}pseudometric}(t, r) \Rightarrow (t(u, w) \leq (t(u, v) + t(v, w))) \end{array}}}

\vnbbookfile{theorem-library/bdd-metric-distance}{theorem-library/bdd-metric-distance}
\noindent{\small bdd-metric-distance.scm -- the BDD-METRIC distance formula, proven.}\par
\vnbprop{bdd-metric-distance}{\vnbstmt{\begin{array}{@{}l@{}} \forall s.; \forall x, y \in \operatorname{pts}(s). \\ \quad \operatorname{is\text{-}metric\text{-}space}(s) \Rightarrow  \\ \quad \quad \operatorname{dist}(\operatorname{bdd\text{-}metric}(s))(x, y) = /(\operatorname{dist}(s)(x, y), (1 + \operatorname{dist}(s)(x, y))) \end{array}}}

\vnbbookfile{theorem-library/bdd-metric-basics}{theorem-library/bdd-metric-basics}
\noindent{\small bdd-metric-basics.scm -- the t/(1+t) scalar family and the two BDD-METRIC statements that read off it, PROVEN.}\par
\vnbprop{bdd-fn-nonneg}{\vnbstmt{\forall t \in \mathbb{R}.\; ((0 \leq t) \Rightarrow (0 \leq /(t, (1 + t))))}}
\vnbprop{bdd-fn-lt-one}{\vnbstmt{\forall t \in \mathbb{R}.\; ((0 \leq t) \Rightarrow (/(t, (1 + t)) < 1))}}
\vnbprop{bdd-fn-le-arg}{\vnbstmt{\forall t \in \mathbb{R}.\; ((0 \leq t) \Rightarrow (/(t, (1 + t)) \leq t))}}
\vnbprop{bdd-fn-zero}{\vnbstmt{\forall t \in \mathbb{R}.\; (((0 \leq t) \wedge (t = 0)) \Rightarrow (/(t, (1 + t)) = 0))}}
\vnbprop{bdd-fn-zero-eq}{\vnbstmt{\forall t \in \mathbb{R}.\; (((0 \leq t) \wedge (/(t, (1 + t)) = 0)) \Rightarrow (t = 0))}}
\vnbprop{bdd-fn-mono}{\vnbstmt{\begin{array}{@{}l@{}} \forall s, t \in \mathbb{R}. \\ \quad ((0 \leq s) \wedge (s \leq t)) \Rightarrow (/(s, (1 + s)) \leq /(t, (1 + t))) \end{array}}}
\vnbprop{bdd-fn-subadd}{\vnbstmt{\begin{array}{@{}l@{}} \forall a, b \in \mathbb{R}. \\ \quad ((0 \leq a) \wedge (0 \leq b)) \Rightarrow  \\ \quad \quad /((a + b), (1 + (a + b))) \leq (/(a, (1 + a)) + /(b, (1 + b))) \end{array}}}
\vnbprop{bdd-metric-bounded}{\vnbstmt{\begin{array}{@{}l@{}} \forall s.; \forall x, y \in \operatorname{pts}(s). \\ \quad \operatorname{is\text{-}metric\text{-}space}(s) \Rightarrow (\operatorname{dist}(\operatorname{bdd\text{-}metric}(s))(x, y) < 1) \end{array}}}
\vnbprop{bdd-metric-is-metric-space}{\vnbstmt{\forall s.\; (\operatorname{is\text{-}metric\text{-}space}(s) \Rightarrow \operatorname{is\text{-}metric\text{-}space}(\operatorname{bdd\text{-}metric}(s)))}}

\vnbbookfile{theorem-library/rake-subseq-leaves}{theorem-library/rake-subseq-leaves}
\noindent{\small rake-subseq-leaves.scm -- the remaining asserted leaves of the subsequence / diagonalization cluster, PROVEN.}\par
\vnbprop{nn-step-mono-ptwise}{\vnbstmt{\begin{array}{@{}l@{}} \forall n \in \mathbb{N}, g.; \forall m \in \mathbb{N}. \\ \quad (\forall k \in \mathbb{N}.\; (g(k) \in \mathbb{N}) \wedge (\forall k \in \mathbb{N}.\; (g(k) < g((k + 1))) \wedge (m < n))) \Rightarrow  \\ \quad \quad g(m) < g(n) \end{array}}}
\vnbprop{nn-step-strictly-mono}{\vnbstmt{\begin{array}{@{}l@{}} \forall g \in (\mathbb{N} \to \mathbb{N}). \\ \quad \forall k \in \mathbb{N}.\; (g(k) < g((k + 1))) \Rightarrow \operatorname{strictly\text{-}mono\text{-}nn}(g) \end{array}}}
\vnbprop{subsequence-capture}{\vnbstmt{\begin{array}{@{}l@{}} \forall s \in \mathrm{Inf}(\mathbb{N}). \exists f \in (\mathbb{N} \to s). \forall m, n \in \mathbb{N}. \\ \quad (m < n) \Rightarrow (f(m) < f(n)) \end{array}}}
\vnbprop{nn-recip-succ-small-at}{\vnbstmt{\forall s.; \forall k \in \mathbb{N}.\; ((\operatorname{pos\text{-}rr}(s) \wedge (\operatorname{recip}(s) < (k + 1))) \Rightarrow (\operatorname{recip}((k + 1)) < s))}}
\vnbprop{null-rr-seq-exists}{\vnbstmt{\exists d.\; \operatorname{null\text{-}rr\text{-}seq}(d)}}

\vnbbookfile{theorem-library/rake-dc-consumers}{theorem-library/rake-dc-consumers}
\noindent{\small rake-dc-consumers.scm -- the consumers of recursive choice (rake batch 8-C).}\par
\vnbprop{cauchy-rapid-subsequence}{\vnbstmt{\begin{array}{@{}l@{}} \forall s, f, d. \\ \quad (\operatorname{is\text{-}cauchy\text{-}seq}(s, f) \wedge ((d \in (\mathbb{N} \to \mathbb{R})) \wedge \forall k \in \mathbb{N}.\; \operatorname{pos\text{-}rr}(d(k)))) \Rightarrow  \\ \quad \quad \exists i \in (\mathbb{N} \to \mathbb{N}). \\ \quad \quad \quad \forall m, n \in \mathbb{N}.\; ((m < n) \Rightarrow (i(m) < i(n))) \wedge \\ \quad \quad \quad \forall k \in \mathbb{N}.\; (\operatorname{dist}(s)(f(i(k)), f(i((k + 1)))) \leq d(k)) \end{array}}}

\vnbbookfile{theorem-library/rake-compact-iff-seq-compact}{theorem-library/rake-compact-iff-seq-compact}
\noindent{\small rake-compact-iff-seq-compact.scm -- the FORWARD half of Prop 3.12 (1)=$>$(3'): a compact metric space is sequentially compact.}\par
\vnbprop{cluster-point-has-convergent-subseq}{\vnbstmt{\begin{array}{@{}l@{}} \forall s, f, x. \\ \quad \operatorname{cluster\text{-}point}(s, f, x) \Rightarrow  \\ \quad \quad \exists i.\; (\operatorname{strictly\text{-}mono\text{-}nn}(i) \wedge \operatorname{converges\text{-}to}(s, \operatorname{subseq}(f, i), x)) \end{array}}}
\vnbprop{compact-implies-seq-compact}{\vnbstmt{\forall s.\; ((\operatorname{is\text{-}metric\text{-}space}(s) \wedge \operatorname{is\text{-}compact}(s)) \Rightarrow \operatorname{seq\text{-}compact}(s))}}

\vnbbookfile{theorem-library/rake-seq-compact-tb}{theorem-library/rake-seq-compact-tb}
\noindent{\small rake-seq-compact-tb.scm -- a sequentially compact metric space is totally bounded (the (C1) leg of Prop 3.12's missing direction; see the closing block of theorem-library/rake-compact-iff-seq-compact.scm).}\par
\vnbprop{seq-compact-implies-totally-bounded}{\vnbstmt{\forall s.\; ((\operatorname{is\text{-}metric\text{-}space}(s) \wedge \operatorname{seq\text{-}compact}(s)) \Rightarrow \operatorname{totally\text{-}bounded}(s))}}

\vnbbookfile{theorem-library/rake-lebesgue-number}{theorem-library/rake-lebesgue-number}
\noindent{\small rake-lebesgue-number.scm -- the Lebesgue number lemma, and the CONVERSE half of Prop 3.12 (1)$<$=$>$(3'): a sequentially compact metric space is compact.}\par
\vnbprop{lebesgue-number}{\vnbstmt{\begin{array}{@{}l@{}} \forall s, v. \\ \quad (\operatorname{is\text{-}metric\text{-}space}(s) \wedge (\operatorname{seq\text{-}compact}(s) \wedge \operatorname{is\text{-}open\text{-}cover}(s, v))) \Rightarrow  \\ \quad \quad \exists l. \\ \quad \quad \quad \operatorname{pos\text{-}rr}(l) \wedge \forall p \in \operatorname{pts}(s).\; \exists u \in v.\; (\operatorname{ball}(s, p, l) \subseteq u) \end{array}}}
\vnbprop{seq-compact-implies-compact}{\vnbstmt{\forall s.\; ((\operatorname{is\text{-}metric\text{-}space}(s) \wedge \operatorname{seq\text{-}compact}(s)) \Rightarrow \operatorname{is\text{-}compact}(s))}}
\vnbprop{compact-iff-seq-compact}{\vnbstmt{\forall s.\; (\operatorname{is\text{-}metric\text{-}space}(s) \Rightarrow (\operatorname{is\text{-}compact}(s) \Leftrightarrow \operatorname{seq\text{-}compact}(s)))}}

\vnbbookfile{theorem-library/rake-nn-enum}{theorem-library/rake-nn-enum}
\noindent{\small theorem-library/rake-nn-enum.scm ==================================================================== The chosen enumeration NN-ENUM and the two strictly-monotone bricks the refinement tower (convergence-block-tower) wants.}\par
\vnbprop{nn-enum-spec}{\vnbstmt{\begin{array}{@{}l@{}} \forall s \in \mathrm{Inf}(\mathbb{N}). \\ \quad \operatorname{nn\text{-}enum}(s) \in (\mathbb{N} \to s) \wedge \\ \quad \forall m, n \in \mathbb{N}.\; ((m < n) \Rightarrow (\operatorname{nn\text{-}enum}(s)(m) < \operatorname{nn\text{-}enum}(s)(n))) \end{array}}}
\vnbprop{strictly-mono-image-infinite}{\vnbstmt{\begin{array}{@{}l@{}} \forall i, a. \\ \quad (\operatorname{strictly\text{-}mono\text{-}nn}(i) \wedge ((a \subseteq \mathbb{N}) \wedge \forall k \in \mathbb{N}.\; (i(k) \in a))) \Rightarrow  \\ \quad \quad \{m \in a : \exists j \in \mathbb{N}.\; (m = i(j))\} \in \mathrm{Inf}(\mathbb{N}) \end{array}}}
\vnbprop{strictly-mono-le-reflect}{\vnbstmt{\begin{array}{@{}l@{}} \forall i.; \forall a, j \in \mathbb{N}. \\ \quad (\operatorname{strictly\text{-}mono\text{-}nn}(i) \wedge (i(a) \leq i(j))) \Rightarrow (a \leq j) \end{array}}}

\vnbbookfile{theorem-library/rake-tb-leaves}{theorem-library/rake-tb-leaves}
\noindent{\small rake-tb-leaves.scm -- the totally-bounded / pigeonhole leaves of the Cauchy-subsequence and separability arcs.  Two theorems, both modulo 0:}\par
\vnbprop{cover-block-step}{\vnbstmt{\begin{array}{@{}l@{}} \forall v \in \mathbf{Set}, f \in (\mathbb{N} \to v), c.; \forall j \in \mathrm{Inf}(\mathbb{N}). \\ \quad \operatorname{is\text{-}finite\text{-}cover}(c, v) \Rightarrow  \\ \quad \quad \exists a \in \mathrm{Inf}(\mathbb{N}). \\ \quad \quad \quad (a \subseteq j) \wedge \exists u \in c.\; \forall i \in a.\; (f(i) \in u) \end{array}}}
\vnbprop{dist-le-implies-in-carrier}{\vnbstmt{\begin{array}{@{}l@{}} \forall s, c, p, r. \\ \quad (\operatorname{is\text{-}metric\text{-}space}(s) \wedge (\operatorname{dist}(s)(c, p) \leq r)) \Rightarrow (c \in \operatorname{pts}(s)) \end{array}}}

\vnbbookfile{theorem-library/rake-tb-leaves-2}{theorem-library/rake-tb-leaves-2}
\noindent{\small rake-tb-leaves-2.scm -- the two totally-bounded leaves that the IS-R-NET repair of 2026-09-19 unblocked.}\par
\vnbprop{tb-scale-dense-seq}{\vnbstmt{\begin{array}{@{}l@{}} \forall s.; \forall a \in \operatorname{pts}(s), r. \\ \quad (\operatorname{totally\text{-}bounded}(s) \wedge \operatorname{pos\text{-}rr}(r)) \Rightarrow  \\ \quad \quad \exists g \in (\mathbb{N} \to \operatorname{pts}(s)). \forall p \in \operatorname{pts}(s). \exists j \in \mathbb{N}. \\ \quad \quad \quad \operatorname{dist}(s)(p, g(j)) < r \end{array}}}
\vnbprop{tb-rad-ball-cover}{\vnbstmt{\begin{array}{@{}l@{}} \forall s, d. \\ \quad (\operatorname{totally\text{-}bounded}(s) \wedge \forall k \in \mathbb{N}.\; \operatorname{pos\text{-}rr}(d(k))) \Rightarrow  \\ \quad \quad \operatorname{pts}(s) \in \mathbf{Set} \wedge \\ \quad \quad \exists v. \\ \quad \quad \quad \forall k \in \mathbb{N}.\; \operatorname{is\text{-}finite\text{-}cover}(v(k), \operatorname{pts}(s)) \wedge \\ \quad \quad \quad \forall k \in \mathbb{N}, u \in v(k).\; \exists c \in \operatorname{pts}(s).\; (u = \operatorname{ball}(s, c, d(k))) \end{array}}}

\vnbbookfile{theorem-library/diagonalization}{theorem-library/diagonalization}
\noindent{\small theorem-library/diagonalization.scm}\par
\vnbprop{diagonalization}{\vnbstmt{\begin{array}{@{}l@{}} \forall s \in (\mathbb{N} \to \mathrm{Inf}(\mathbb{N})). \\ \quad \forall k \in \mathbb{N}.\; (s((k + 1)) \subseteq s(k)) \Rightarrow  \\ \quad \quad \exists f. \\ \quad \quad \quad \operatorname{strictly\text{-}mono\text{-}nn}(f) \wedge \forall k, j \in \mathbb{N}.\; ((k \leq j) \Rightarrow (f(j) \in s(k))) \end{array}}}

\vnbbookfile{theorem-library/block-family-combinatorial-proof}{theorem-library/block-family-combinatorial-proof}
\noindent{\small theorem-library/block-family-combinatorial-proof.scm}\par
\vnbprop{block-family-combinatorial}{\vnbstmt{\begin{array}{@{}l@{}} \forall v \in \mathbf{Set}, f \in (\mathbb{N} \to v), a. \\ \quad \forall k \in \mathbb{N}.\; \operatorname{is\text{-}finite\text{-}cover}(a(k), v) \Rightarrow  \\ \quad \quad \exists b \in (\mathbb{N} \to \mathrm{Inf}(\mathbb{N})). \\ \quad \quad \quad \forall k \in \mathbb{N}.\; (b((k + 1)) \subseteq b(k)) \wedge \\ \quad \quad \quad \forall k \in \mathbb{N}.\; \exists u \in a(k).\; \forall i \in b(k).\; (f(i) \in u) \end{array}}}

\vnbbookfile{theorem-library/cauchy-subseq-proof}{theorem-library/cauchy-subseq-proof}
\noindent{\small theorem-library/cauchy-subseq-proof.scm ==================================================================== totally-bounded-has-cauchy-subsequence -- PROVEN (was asserted).}\par
\vnbprop{totally-bounded-has-cauchy-subsequence}{\vnbstmt{\begin{array}{@{}l@{}} \forall s.; \forall f \in (\mathbb{N} \to \operatorname{pts}(s)). \\ \quad \operatorname{totally\text{-}bounded}(s) \Rightarrow  \\ \quad \quad \exists i.\; (\operatorname{strictly\text{-}mono\text{-}nn}(i) \wedge \operatorname{is\text{-}cauchy\text{-}seq}(s, \operatorname{subseq}(f, i))) \end{array}}}

\vnbbookfile{theorem-library/rr-ms-dist}{theorem-library/rr-ms-dist}
\noindent{\small rr-ms-dist.scm -- the distance of RR-MS, on the surface.}\par
\vnbprop{rr-ms-dist}{\vnbstmt{\forall u, v.\; (((u \in \mathbb{R}) \wedge (v \in \mathbb{R})) \Rightarrow (\operatorname{dist}(\operatorname{rr\text{-}ms})(u, v) \simeq \lvert (u - v) \rvert))}}

\vnbbookfile{theorem-library/rr-metric-space-proof}{theorem-library/rr-metric-space-proof}
\noindent{\small rr-metric-space-proof.scm -- IS-METRIC-SPACE(RR-MS), PROVEN.}\par
\vnbprop{rr-is-metric-space}{\vnbstmt{\operatorname{is\text{-}metric\text{-}space}(\operatorname{rr\text{-}ms})}}

\vnbbookfile{theorem-library/trunc-metric-proof}{theorem-library/trunc-metric-proof}
\noindent{\small trunc-metric-proof.scm -- the truncated metric d'(x,y) = min(1, d(x,y)) is a metric, is bounded by 1, and is UNIFORMLY equivalent to d.}\par
\vnbprop{rr-min-one-mono}{\vnbstmt{\forall a, b \in \mathbb{R}.\; ((a \leq b) \Rightarrow (\operatorname{min}(1, a) \leq \operatorname{min}(1, b)))}}
\vnbprop{trunc-metric-dist}{\vnbstmt{\begin{array}{@{}l@{}} \forall s, u, v. \\ \quad ((u \in \operatorname{pts}(s)) \wedge (v \in \operatorname{pts}(s))) \Rightarrow  \\ \quad \quad (\operatorname{dist}(\operatorname{trunc\text{-}metric}(s))(u, v) \simeq \operatorname{min}(1, \operatorname{dist}(s)(u, v))) \end{array}}}
\vnbprop{trunc-metric-is-metric-space}{\vnbstmt{\forall s.\; (\operatorname{is\text{-}metric\text{-}space}(s) \Rightarrow \operatorname{is\text{-}metric\text{-}space}(\operatorname{trunc\text{-}metric}(s)))}}
\vnbprop{trunc-metric-bounded}{\vnbstmt{\begin{array}{@{}l@{}} \forall s, u, v. \\ \quad (\operatorname{is\text{-}metric\text{-}space}(s) \wedge ((u \in \operatorname{pts}(s)) \wedge (v \in \operatorname{pts}(s)))) \Rightarrow  \\ \quad \quad \operatorname{dist}(\operatorname{trunc\text{-}metric}(s))(u, v) \leq 1 \end{array}}}
\vnbprop{trunc-metric-id-uniformly-continuous}{\vnbstmt{\begin{array}{@{}l@{}} \forall s. \\ \quad \operatorname{is\text{-}metric\text{-}space}(s) \Rightarrow  \\ \quad \quad \operatorname{is\text{-}uniformly\text{-}continuous}(s, \operatorname{trunc\text{-}metric}(s), (\lambda a \in \operatorname{pts}(s).\; a)) \end{array}}}
\vnbprop{trunc-metric-id-uniformly-continuous-back}{\vnbstmt{\begin{array}{@{}l@{}} \forall s. \\ \quad \operatorname{is\text{-}metric\text{-}space}(s) \Rightarrow  \\ \quad \quad \operatorname{is\text{-}uniformly\text{-}continuous}(\operatorname{trunc\text{-}metric}(s), s, (\lambda a \in \operatorname{pts}(s).\; a)) \end{array}}}

\vnbbookfile{theorem-library/rr-complete-proof}{theorem-library/rr-complete-proof}
\noindent{\small rr-complete-proof.scm -- RR is a COMPLETE METRIC SPACE.  Proven 2026-08-01/02.}\par
\vnbprop{rr-complete}{\vnbstmt{\operatorname{is\text{-}complete}(\operatorname{rr\text{-}ms})}}

\vnbbookfile{theorem-library/metric-limit-unique}{theorem-library/metric-limit-unique}
\noindent{\small metric-limit-unique.scm -- A SEQUENCE IN A METRIC SPACE HAS AT MOST ONE LIMIT, proven once for an arbitrary metric space, with the RR and CC cases taken as instances.}\par
\vnbprop{metric-limit-unique}{\vnbstmt{\begin{array}{@{}l@{}} \forall s, f, \ell_{1}, \ell_{2}. \\ \quad (\operatorname{converges\text{-}to}(s, f, \ell_{1}) \wedge \operatorname{converges\text{-}to}(s, f, \ell_{2})) \Rightarrow (\ell_{1} = \ell_{2}) \end{array}}}
\vnbprop{rr-limit-unique}{\vnbstmt{\begin{array}{@{}l@{}} \forall f \in (\mathbb{N} \to \mathbb{R}), \ell_{1}, \ell_{2} \in \mathbb{R}. \\ \quad (\operatorname{converges\text{-}to}(\operatorname{rr\text{-}ms}, f, \ell_{1}) \wedge \operatorname{converges\text{-}to}(\operatorname{rr\text{-}ms}, f, \ell_{2})) \Rightarrow  \\ \quad \quad \ell_{1} = \ell_{2} \end{array}}}
\vnbprop{cc-limit-unique}{\vnbstmt{\begin{array}{@{}l@{}} \forall f \in (\mathbb{N} \to \mathbb{C}), \ell_{1}, \ell_{2} \in \mathbb{C}. \\ \quad (\operatorname{converges\text{-}to}(\operatorname{cc\text{-}ms}, f, \ell_{1}) \wedge \operatorname{converges\text{-}to}(\operatorname{cc\text{-}ms}, f, \ell_{2})) \Rightarrow  \\ \quad \quad \ell_{1} = \ell_{2} \end{array}}}

\vnbbookfile{theorem-library/seq-limit-core}{theorem-library/seq-limit-core}
\noindent{\small theorem-library/seq-limit-core.scm -- THE LIMIT OF A REAL SEQUENCE AS A TERM (SEQ-LIMIT), its three characterising theorems, the TAIL bound rr-limit-tail-abs-le, and rr-cauchy-converges.}\par
\par\medskip\noindent\textbf{Definition} (\texttt{seq-limit}).\;$\operatorname{seq\text{-}limit}(f) :=$
\begin{equation*}
\left(\text{if }\operatorname{converges}(\operatorname{rr\text{-}ms}, f)\text{ then }(\iota\, \operatorname{lm\_}.\;\operatorname{converges\text{-}to}(\operatorname{rr\text{-}ms}, f, \operatorname{lm\_}))\text{ else }0\right)
\end{equation*}

\vnbprop{seq-limit-converges-to}{\vnbstmt{\forall f.\; (\operatorname{converges}(\operatorname{rr\text{-}ms}, f) \Rightarrow \operatorname{converges\text{-}to}(\operatorname{rr\text{-}ms}, f, \operatorname{seq\text{-}limit}(f)))}}
\vnbprop{seq-limit-in-rr}{\vnbstmt{\forall f.\; (\operatorname{seq\text{-}limit}(f) \in \mathbb{R})}}
\vnbprop{seq-limit-value}{\vnbstmt{\forall f, v.\; (\operatorname{converges\text{-}to}(\operatorname{rr\text{-}ms}, f, v) \Rightarrow (\operatorname{seq\text{-}limit}(f) = v))}}
\vnbprop{rr-limit-tail-abs-le}{\vnbstmt{\begin{array}{@{}l@{}} \forall f \in (\mathbb{N} \to \mathbb{R}), v, a, c \in \mathbb{R}, n \in \mathbb{N}. \\ \quad (\operatorname{converges\text{-}to}(\operatorname{rr\text{-}ms}, f, v) \wedge \forall k \in \mathbb{N}.\; ((n \leq k) \Rightarrow (\lvert (f(k) - a) \rvert \leq c))) \Rightarrow  \\ \quad \quad \lvert (v - a) \rvert \leq c \end{array}}}
\vnbprop{rr-cauchy-converges}{\vnbstmt{\begin{array}{@{}l@{}} \forall f \in (\mathbb{N} \to \mathbb{R}). \\ \quad \forall s. \\ \quad \quad \operatorname{pos\text{-}rr}(s) \Rightarrow  \\ \quad \quad \quad \exists n \in \mathbb{N}. \forall m, p \in \mathbb{N}. \\ \quad \quad \quad \quad ((n \leq m) \wedge (n \leq p)) \Rightarrow (\lvert (f(m) - f(p)) \rvert \leq s) \Rightarrow \\ \quad \quad \operatorname{converges}(\operatorname{rr\text{-}ms}, f) \end{array}}}

\vnbbookfile{structure-library/ascoli-arzela-statement}{structure-library/ascoli-arzela-statement}
\noindent{\small NOTE 2026-09-20 (batch 12-B).  Where this file says the tree has no metric SUBSPACE structure, that is history: SUBSPACE-MS(s, A) and RESTRICT(f, A) are defined in structure-library/metric-subspace.scm, their laws are in theorem-library/metric-subspace-laws.scm and compact-subspace.scm, and}\par
\par\medskip\noindent\textbf{Definition} (\texttt{pointwise-bounded}).\;For all $s$ and $\mathit{fam}$, 
\begin{equation*}
\begin{array}{@{}l@{}} \operatorname{pointwise\text{-}bounded}(s, \mathit{fam}) \\ \Leftrightarrow\; \operatorname{is\text{-}metric\text{-}space}(s) \wedge \\ \mathit{fam} \in (\mathbb{N} \to (\operatorname{pts}(s) \to \mathbb{R})) \wedge \\ \forall x \in \operatorname{pts}(s). \exists \mathit{bd} \in \mathbb{R}. \forall k \in \mathbb{N}. \\ \quad \lvert \operatorname{fam}(k)(x) \rvert \leq \mathit{bd} \end{array}
\end{equation*}

\par\medskip\noindent\textbf{Definition} (\texttt{is-equicontinuous}).\;For all $s$, $t$ and $\mathit{fam}$, 
\begin{equation*}
\begin{array}{@{}l@{}} \operatorname{is\text{-}equicontinuous}(s, t, \mathit{fam}) \\ \Leftrightarrow\; \operatorname{is\text{-}metric\text{-}space}(s) \wedge \\ \operatorname{is\text{-}metric\text{-}space}(t) \wedge \\ \mathit{fam} \in (\mathbb{N} \to (\operatorname{pts}(s) \to \operatorname{pts}(t))) \wedge \\ \forall x \in \operatorname{pts}(s), \mathit{eps}. \\ \quad \operatorname{pos\text{-}rr}(\mathit{eps}) \Rightarrow  \\ \quad \quad \exists \mathit{del}. \\ \quad \quad \quad \operatorname{pos\text{-}rr}(\mathit{del}) \wedge \\ \quad \quad \quad \forall k \in \mathbb{N}, y \in \operatorname{pts}(s). \\ \quad \quad \quad \quad (\operatorname{dist}(s)(x, y) < \mathit{del}) \Rightarrow (\operatorname{dist}(t)(\operatorname{fam}(k)(x), \operatorname{fam}(k)(y)) < \mathit{eps}) \end{array}
\end{equation*}

\par\medskip\noindent\textbf{Definition} (\texttt{converges-uniformly}).\;For all $s$, $\mathit{seq}$ and $g$, 
\begin{equation*}
\begin{array}{@{}l@{}} \operatorname{converges\text{-}uniformly}(s, \mathit{seq}, g) \\ \Leftrightarrow\; \operatorname{is\text{-}metric\text{-}space}(s) \wedge \\ \mathit{seq} \in (\mathbb{N} \to (\operatorname{pts}(s) \to \mathbb{R})) \wedge \\ g \in (\operatorname{pts}(s) \to \mathbb{R}) \wedge \\ \forall \mathit{eps}. \\ \quad \operatorname{pos\text{-}rr}(\mathit{eps}) \Rightarrow  \\ \quad \quad \exists \mathit{cap} \in \mathbb{N}. \forall k. \\ \quad \quad \quad ((k \in \mathbb{N}) \wedge (\mathit{cap} \leq k)) \Rightarrow  \\ \quad \quad \quad \quad \forall x \in \operatorname{pts}(s).\; (\lvert -(\operatorname{seq}(k)(x), g(x)) \rvert < \mathit{eps}) \end{array}
\end{equation*}


\vnbbookfile{theorem-library/ascoli-analytic-cores}{theorem-library/ascoli-analytic-cores}
\noindent{\small ascoli-analytic-cores.scm -- the analytic content of Ascoli-Arzelà, machine-proved. bridge, MACHINE-PROVED, and the one rung it still stands on.}\par
\par\medskip\noindent\textbf{Definition} (\texttt{is-unif-cauchy}).\;For all $s$ and $\mathit{fam}$, 
\begin{equation*}
\begin{array}{@{}l@{}} \operatorname{is\text{-}unif\text{-}cauchy}(s, \mathit{fam}) \\ \Leftrightarrow\; \operatorname{is\text{-}metric\text{-}space}(s) \wedge \\ \mathit{fam} \in (\mathbb{N} \to (\operatorname{pts}(s) \to \mathbb{R})) \wedge \\ \forall \mathit{eps}. \\ \quad \operatorname{pos\text{-}rr}(\mathit{eps}) \Rightarrow  \\ \quad \quad \exists \mathit{cap} \in \mathbb{N}. \forall k. \\ \quad \quad \quad ((k \in \mathbb{N}) \wedge (\mathit{cap} \leq k)) \Rightarrow  \\ \quad \quad \quad \quad \forall l. \\ \quad \quad \quad \quad \quad ((l \in \mathbb{N}) \wedge (\mathit{cap} \leq l)) \Rightarrow  \\ \quad \quad \quad \quad \quad \quad \forall x \in \operatorname{pts}(s).\; (\lvert -(\operatorname{fam}(k)(x), \operatorname{fam}(l)(x)) \rvert < \mathit{eps}) \end{array}
\end{equation*}

\par\medskip\noindent\textbf{Definition} (\texttt{is-ptwise-cauchy}).\;For all $s$ and $\mathit{fam}$, 
\begin{equation*}
\begin{array}{@{}l@{}} \operatorname{is\text{-}ptwise\text{-}cauchy}(s, \mathit{fam}) \\ \Leftrightarrow\; \operatorname{is\text{-}metric\text{-}space}(s) \wedge \\ \mathit{fam} \in (\mathbb{N} \to (\operatorname{pts}(s) \to \mathbb{R})) \wedge \\ \forall x \in \operatorname{pts}(s), \mathit{eps}. \\ \quad \operatorname{pos\text{-}rr}(\mathit{eps}) \Rightarrow  \\ \quad \quad \exists \mathit{cap} \in \mathbb{N}. \forall k. \\ \quad \quad \quad ((k \in \mathbb{N}) \wedge (\mathit{cap} \leq k)) \Rightarrow  \\ \quad \quad \quad \quad \forall l.\; (((l \in \mathbb{N}) \wedge (\mathit{cap} \leq l)) \Rightarrow (\lvert -(\operatorname{fam}(k)(x), \operatorname{fam}(l)(x)) \rvert < \mathit{eps})) \end{array}
\end{equation*}

\par\medskip\noindent\textbf{Definition} (\texttt{is-dense-seq}).\;For all $s$ and $\mathit{dseq}$, 
\begin{equation*}
\begin{array}{@{}l@{}} \operatorname{is\text{-}dense\text{-}seq}(s, \mathit{dseq}) \\ \Leftrightarrow\; \operatorname{is\text{-}metric\text{-}space}(s) \wedge \\ \mathit{dseq} \in (\mathbb{N} \to \operatorname{pts}(s)) \wedge \\ \forall x \in \operatorname{pts}(s), \mathit{eps}. \\ \quad \operatorname{pos\text{-}rr}(\mathit{eps}) \Rightarrow \exists m \in \mathbb{N}.\; (\operatorname{dist}(s)(x, \operatorname{dseq}(m)) < \mathit{eps}) \end{array}
\end{equation*}

\par\medskip\noindent\textbf{Definition} (\texttt{converges-on}).\;For all $s$, $\mathit{fam}$ and $\mathit{dseq}$, 
\begin{equation*}
\begin{array}{@{}l@{}} \operatorname{converges\text{-}on}(s, \mathit{fam}, \mathit{dseq}) \\ \Leftrightarrow\; \operatorname{is\text{-}metric\text{-}space}(s) \wedge \\ \mathit{fam} \in (\mathbb{N} \to (\operatorname{pts}(s) \to \mathbb{R})) \wedge \\ \mathit{dseq} \in (\mathbb{N} \to \operatorname{pts}(s)) \wedge \\ \forall m \in \mathbb{N}.\; \operatorname{converges}(\operatorname{rr\text{-}ms}, (\lambda k \in \mathbb{N}.\; \operatorname{fam}(k)(\operatorname{dseq}(m)))) \end{array}
\end{equation*}

\vnbprop{equicont-dense-conv-ptwise-cauchy}{\vnbstmt{\begin{array}{@{}l@{}} \forall s.; \forall m \in (\mathbb{N} \to (\operatorname{pts}(s) \to \mathbb{R})). \\ \quad (\operatorname{is\text{-}equicontinuous}(s, \operatorname{rr\text{-}ms}, m) \wedge \exists q.\; (\operatorname{is\text{-}dense\text{-}seq}(s, q) \wedge \operatorname{converges\text{-}on}(s, m, q))) \Rightarrow  \\ \quad \quad \operatorname{is\text{-}ptwise\text{-}cauchy}(s, m) \end{array}}}
\vnbprop{uniform-limit-continuous}{\vnbstmt{\begin{array}{@{}l@{}} \forall s, m, g. \\ \quad m \in (\mathbb{N} \to (\operatorname{pts}(s) \to \mathbb{R})) \wedge \\ \quad g \in (\operatorname{pts}(s) \to \mathbb{R}) \wedge \\ \quad \forall k \in \mathbb{N}.\; \operatorname{is\text{-}continuous}(s, \operatorname{rr\text{-}ms}, m(k)) \wedge \\ \quad \operatorname{converges\text{-}uniformly}(s, m, g) \Rightarrow \\ \quad \quad \operatorname{is\text{-}continuous}(s, \operatorname{rr\text{-}ms}, g) \end{array}}}

\vnbbookfile{theorem-library/rake-open-sets}{theorem-library/rake-open-sets}
\noindent{\small rake-open-sets.scm -- the metric topology axioms, PROVEN.}\par
\vnbprop{empty-is-open}{\vnbstmt{\forall s.\; (\operatorname{is\text{-}metric\text{-}space}(s) \Rightarrow \operatorname{is\text{-}open}(s, \emptyset))}}
\vnbprop{carrier-is-open}{\vnbstmt{\forall s.\; (\operatorname{is\text{-}metric\text{-}space}(s) \Rightarrow \operatorname{is\text{-}open}(s, \operatorname{pts}(s)))}}
\vnbprop{union-of-opens-open}{\vnbstmt{\begin{array}{@{}l@{}} \forall s, a, g. \\ \quad (\operatorname{is\text{-}metric\text{-}space}(s) \wedge \forall i \in a.\; \operatorname{is\text{-}open}(s, g(i))) \Rightarrow  \\ \quad \quad \operatorname{is\text{-}open}(s, \operatorname{big\text{-}union}(i, a, g(i))) \end{array}}}
\vnbprop{inter-of-opens-open}{\vnbstmt{\begin{array}{@{}l@{}} \forall s, u, w. \\ \quad (\operatorname{is\text{-}metric\text{-}space}(s) \wedge (\operatorname{is\text{-}open}(s, u) \wedge \operatorname{is\text{-}open}(s, w))) \Rightarrow  \\ \quad \quad \operatorname{is\text{-}open}(s, (u \cap w)) \end{array}}}
\vnbprop{preimage-complement}{\vnbstmt{\begin{array}{@{}l@{}} \forall s, t, f, u. \\ \quad (\operatorname{is\text{-}metric\text{-}space}(s) \wedge (\operatorname{is\text{-}metric\text{-}space}(t) \wedge (f \in (\operatorname{pts}(s) \to \operatorname{pts}(t))))) \Rightarrow  \\ \quad \quad \operatorname{preimage}(s, f, \operatorname{complement\text{-}in}(\operatorname{pts}(t), u)) = \operatorname{complement\text{-}in}(\operatorname{pts}(s), \operatorname{preimage}(s, f, u)) \end{array}}}
\vnbprop{open-preimage-implies-continuous}{\vnbstmt{\begin{array}{@{}l@{}} \forall s, t, f. \\ \quad \operatorname{is\text{-}metric\text{-}space}(s) \wedge \\ \quad \operatorname{is\text{-}metric\text{-}space}(t) \wedge \\ \quad f \in (\operatorname{pts}(s) \to \operatorname{pts}(t)) \wedge \\ \quad \forall v.\; (\operatorname{is\text{-}open}(t, v) \Rightarrow \operatorname{is\text{-}open}(s, \operatorname{preimage}(s, f, v))) \Rightarrow \\ \quad \quad \operatorname{is\text{-}continuous}(s, t, f) \end{array}}}
\vnbprop{continuous-implies-closed-preimage}{\vnbstmt{\begin{array}{@{}l@{}} \forall s, t, f, a. \\ \quad (\operatorname{is\text{-}continuous}(s, t, f) \wedge \operatorname{is\text{-}closed}(t, a)) \Rightarrow  \\ \quad \quad \operatorname{is\text{-}closed}(s, \operatorname{preimage}(s, f, a)) \end{array}}}
\vnbprop{closed-preimage-implies-continuous}{\vnbstmt{\begin{array}{@{}l@{}} \forall s, t, f. \\ \quad \operatorname{is\text{-}metric\text{-}space}(s) \wedge \\ \quad \operatorname{is\text{-}metric\text{-}space}(t) \wedge \\ \quad f \in (\operatorname{pts}(s) \to \operatorname{pts}(t)) \wedge \\ \quad \forall a.\; (\operatorname{is\text{-}closed}(t, a) \Rightarrow \operatorname{is\text{-}closed}(s, \operatorname{preimage}(s, f, a))) \Rightarrow \\ \quad \quad \operatorname{is\text{-}continuous}(s, t, f) \end{array}}}

\vnbbookfile{theorem-library/ptwise-cauchy-unif}{theorem-library/ptwise-cauchy-unif}
\noindent{\small ptwise-cauchy-unif.scm -- pointwise Cauchy + equicontinuous on a COMPACT space  =$>$  uniformly Cauchy  (the user's notes, Prop 3.33, Cauchy spelling).}\par
\vnbprop{ptwise-cauchy-compact-equicont-unif}{\vnbstmt{\begin{array}{@{}l@{}} \forall s, m. \\ \quad \operatorname{is\text{-}compact}(s) \wedge \\ \quad m \in (\mathbb{N} \to (\operatorname{pts}(s) \to \mathbb{R})) \wedge \\ \quad \operatorname{is\text{-}equicontinuous}(s, \operatorname{rr\text{-}ms}, m) \wedge \\ \quad \operatorname{is\text{-}ptwise\text{-}cauchy}(s, m) \Rightarrow \\ \quad \quad \operatorname{is\text{-}unif\text{-}cauchy}(s, m) \end{array}}}

\vnbbookfile{theorem-library/ascoli-assembly}{theorem-library/ascoli-assembly}
\noindent{\small theorem-library/ascoli-assembly.scm -- equicont-dense-conv-implies-unif-cauchy.}\par
\vnbprop{equicont-dense-conv-implies-unif-cauchy}{\vnbstmt{\begin{array}{@{}l@{}} \forall s, m. \\ \quad \operatorname{is\text{-}compact}(s) \wedge \\ \quad m \in (\mathbb{N} \to (\operatorname{pts}(s) \to \mathbb{R})) \wedge \\ \quad \operatorname{is\text{-}equicontinuous}(s, \operatorname{rr\text{-}ms}, m) \wedge \\ \quad \exists q.\; (\operatorname{is\text{-}dense\text{-}seq}(s, q) \wedge \operatorname{converges\text{-}on}(s, m, q)) \Rightarrow \\ \quad \quad \operatorname{is\text{-}unif\text{-}cauchy}(s, m) \end{array}}}

\vnbbookfile{theorem-library/unif-cauchy-limit}{theorem-library/unif-cauchy-limit}
\noindent{\small unif-cauchy-limit.scm -- unif-cauchy-has-uniform-limit, PROVEN (2026-09-14).}\par
\vnbprop{unif-cauchy-has-uniform-limit}{\vnbstmt{\begin{array}{@{}l@{}} \forall s.; \forall m \in (\mathbb{N} \to (\operatorname{pts}(s) \to \mathbb{R})). \\ \quad \operatorname{is\text{-}unif\text{-}cauchy}(s, m) \Rightarrow  \\ \quad \quad \exists g \in (\operatorname{pts}(s) \to \mathbb{R}).\; \operatorname{converges\text{-}uniformly}(s, m, g) \end{array}}}

\vnbbookfile{theorem-library/ascoli-bridge}{theorem-library/ascoli-bridge}
\noindent{\small ascoli-bridge.scm -- the dense-to-uniform bridge for Ascoli-Arzelà.}\par
\vnbprop{unif-cauchy-cont-implies-uniform-limit}{\vnbstmt{\begin{array}{@{}l@{}} \forall s.; \forall m \in (\mathbb{N} \to (\operatorname{pts}(s) \to \mathbb{R})). \\ \quad (\forall k \in \mathbb{N}.\; \operatorname{is\text{-}continuous}(s, \operatorname{rr\text{-}ms}, m(k)) \wedge \operatorname{is\text{-}unif\text{-}cauchy}(s, m)) \Rightarrow  \\ \quad \quad \exists g \in (\operatorname{pts}(s) \to \mathbb{R}). \\ \quad \quad \quad \operatorname{is\text{-}continuous}(s, \operatorname{rr\text{-}ms}, g) \wedge \operatorname{converges\text{-}uniformly}(s, m, g) \end{array}}}
\vnbprop{ascoli-dense-bridge}{\vnbstmt{\begin{array}{@{}l@{}} \forall s, m. \\ \quad \operatorname{is\text{-}compact}(s) \wedge \\ \quad m \in (\mathbb{N} \to (\operatorname{pts}(s) \to \mathbb{R})) \wedge \\ \quad \forall k \in \mathbb{N}.\; \operatorname{is\text{-}continuous}(s, \operatorname{rr\text{-}ms}, m(k)) \wedge \\ \quad \operatorname{is\text{-}equicontinuous}(s, \operatorname{rr\text{-}ms}, m) \wedge \\ \quad \exists q.\; (\operatorname{is\text{-}dense\text{-}seq}(s, q) \wedge \operatorname{converges\text{-}on}(s, m, q)) \Rightarrow \\ \quad \quad \exists g \in (\operatorname{pts}(s) \to \mathbb{R}). \\ \quad \quad \quad \operatorname{is\text{-}continuous}(s, \operatorname{rr\text{-}ms}, g) \wedge \operatorname{converges\text{-}uniformly}(s, m, g) \end{array}}}

\vnbbookfile{structure-library/seminorm-hahn-banach}{structure-library/seminorm-hahn-banach}
\noindent{\small MOVED 2026-09-20 (batch 12-A) from theorem-library/: this file is VOCABULARY -- definitions, notation and warranted supports, not one proof -- and every theorem stated with it had to load below it.  Its load.scm slot is unchanged. seminorm-hahn-banach.scm -- seminorms on a K-vector space and the}\par
\par\medskip\noindent\textbf{Definition} (\texttt{is-seminorm}).\;For all $\mathit{fld}$, $m$ and $p$, 
\begin{equation*}
\begin{array}{@{}l@{}} \operatorname{is\text{-}seminorm}(\mathit{fld}, m, p) \\ \Leftrightarrow\; \operatorname{is\text{-}normed\text{-}field}(\mathit{fld}) \wedge \\ \operatorname{is\text{-}module}(m) \wedge \\ \operatorname{scal}(m) = \operatorname{normed\text{-}field\text{-}as\text{-}commutative\text{-}ring}(\mathit{fld}) \wedge \\ p \in (\operatorname{vec}(m) \to \mathbb{R}) \wedge \\ \forall x \in \operatorname{vec}(m).\; (0 \leq p(x)) \wedge \\ \forall x, y \in \operatorname{vec}(m).\; (p(\operatorname{vadd}(m)(x, y)) \leq +(p(x), p(y))) \wedge \\ \forall \mathit{lam} \in \operatorname{carr}(\operatorname{scal}(m)), x \in \operatorname{vec}(m). \\ \quad p(\operatorname{act}(m)(\mathit{lam}, x)) = \cdot(\operatorname{fnrm}(\mathit{fld})(\mathit{lam}), p(x)) \end{array}
\end{equation*}

\par\medskip\noindent\textbf{Definition} (\texttt{is-k-linear-on}).\;For all $m$, $s$ and $f$, 
\begin{equation*}
\begin{array}{@{}l@{}} \operatorname{is\text{-}k\text{-}linear\text{-}on}(m, s, f) \\ \Leftrightarrow\; \operatorname{is\text{-}module}(m) \wedge \\ \operatorname{is\text{-}submodule}(m, s) \wedge \\ f \in (s \to \operatorname{carr}(\operatorname{scal}(m))) \wedge \\ \forall x, y \in s.\; (f(\operatorname{vadd}(m)(x, y)) = \operatorname{add}(\operatorname{scal}(m))(f(x), f(y))) \wedge \\ \forall \mathit{lam} \in \operatorname{carr}(\operatorname{scal}(m)), x \in s. \\ \quad f(\operatorname{act}(m)(\mathit{lam}, x)) = \operatorname{mul}(\operatorname{scal}(m))(\mathit{lam}, f(x)) \end{array}
\end{equation*}

\par\medskip\noindent\textbf{Definition} (\texttt{is-k-linear}).\;For all $m$ and $f$, 
\begin{equation*}
\begin{array}{@{}l@{}} \operatorname{is\text{-}k\text{-}linear}(m, f) \\ \Leftrightarrow\; \operatorname{is\text{-}module}(m) \wedge \\ f \in (\operatorname{vec}(m) \to \operatorname{carr}(\operatorname{scal}(m))) \wedge \\ \forall x, y \in \operatorname{vec}(m). \\ \quad f(\operatorname{vadd}(m)(x, y)) = \operatorname{add}(\operatorname{scal}(m))(f(x), f(y)) \wedge \\ \forall \mathit{lam} \in \operatorname{carr}(\operatorname{scal}(m)), x \in \operatorname{vec}(m). \\ \quad f(\operatorname{act}(m)(\mathit{lam}, x)) = \operatorname{mul}(\operatorname{scal}(m))(\mathit{lam}, f(x)) \end{array}
\end{equation*}

\par\medskip\noindent\textbf{Definition} (\texttt{is-nontrivially-valued}).\;For every $\mathit{fld}$, 
\begin{equation*}
\begin{array}{@{}l@{}} \operatorname{is\text{-}nontrivially\text{-}valued}(\mathit{fld}) \\ \Leftrightarrow\; \operatorname{is\text{-}normed\text{-}field}(\mathit{fld}) \wedge \\ \exists c \in \operatorname{carr}(\mathit{fld}). \\ \quad (1 < \operatorname{fnrm}(\mathit{fld})(c)) \wedge \exists b \in \operatorname{carr}(\mathit{fld}).\; (\operatorname{mul}(\mathit{fld})(c, b) = \operatorname{one}(\mathit{fld})) \end{array}
\end{equation*}


\vnbbookfile{structure-library/baire-category}{structure-library/baire-category}
\noindent{\small MOVED 2026-09-20 (batch 12-A) from theorem-library/: this file is VOCABULARY -- definitions, notation and warranted supports, not one proof -- and every theorem stated with it had to load below it.  Its load.scm slot is unchanged. baire-category.scm -- interior/closure, nowhere-dense, meager/nonmeager, and}\par
\par\medskip\noindent\textbf{Definition} (\texttt{interior}).\;$\operatorname{interior}(s, a) :=$
\begin{equation*}
\{x \in \operatorname{pts}(s) : \exists u.\; (\operatorname{is\text{-}open}(s, u) \wedge ((x \in u) \wedge (u \subseteq a)))\}
\end{equation*}

\par\medskip\noindent\textbf{Definition} (\texttt{closure}).\;$\operatorname{closure}(s, a) :=$
\begin{equation*}
\{x \in \operatorname{pts}(s) : \forall u.\; ((\operatorname{is\text{-}open}(s, u) \wedge (x \in u)) \Rightarrow \exists y \in u.\; (y \in a))\}
\end{equation*}

\par\medskip\noindent\textbf{Definition} (\texttt{is-nowhere-dense}).\;For all $s$ and $a$, 
\begin{equation*}
\begin{array}{@{}l@{}} \operatorname{is\text{-}nowhere\text{-}dense}(s, a) \\ \Leftrightarrow\; (\operatorname{is\text{-}metric\text{-}space}(s) \wedge ((a \subseteq \operatorname{pts}(s)) \wedge (\operatorname{interior}(s, \operatorname{closure}(s, a)) = \emptyset))) \end{array}
\end{equation*}

\par\medskip\noindent\textbf{Definition} (\texttt{is-meager}).\;For all $s$ and $a$, 
\begin{equation*}
\begin{array}{@{}l@{}} \operatorname{is\text{-}meager}(s, a) \\ \Leftrightarrow\; \operatorname{is\text{-}metric\text{-}space}(s) \wedge \\ a \subseteq \operatorname{pts}(s) \wedge \\ \exists \mathit{ee} \in (\mathbb{N} \to \operatorname{power}(\operatorname{pts}(s))). \\ \quad (\forall n \in \mathbb{N}.\; \operatorname{is\text{-}nowhere\text{-}dense}(s, \operatorname{ee}(n)) \wedge (a \subseteq \operatorname{big\text{-}union}(n, \mathbb{N}, \operatorname{ee}(n)))) \end{array}
\end{equation*}

\par\medskip\noindent\textbf{Definition} (\texttt{is-nonmeager}).\;For all $s$ and $a$, 
\begin{equation*}
\operatorname{is\text{-}nonmeager}(s, a) \Leftrightarrow (\operatorname{is\text{-}metric\text{-}space}(s) \wedge ((a \subseteq \operatorname{pts}(s)) \wedge \neg \operatorname{is\text{-}meager}(s, a)))
\end{equation*}


\vnbbookfile{theorem-library/metric-closure-laws}{theorem-library/metric-closure-laws}
\noindent{\small metric-closure-laws.scm -- CLOSURE, INTERIOR and the CLOSED BALL.}\par
\vnbprop{closed-ball-sep-unfold}{\vnbstmt{\forall s, c, r.\; (\operatorname{closed\text{-}ball}(s, c, r) \simeq \{y \in \operatorname{pts}(s) : (\operatorname{dist}(s)(c, y) \leq r)\})}}
\vnbprop{closure-sep-unfold}{\vnbstmt{\begin{array}{@{}l@{}} \forall s, a. \\ \quad (\operatorname{closure}(s, a) \simeq \{x \in \operatorname{pts}(s) : \forall u.\; ((\operatorname{is\text{-}open}(s, u) \wedge (x \in u)) \Rightarrow \exists y \in u.\; (y \in a))\}) \end{array}}}
\vnbprop{interior-sep-unfold}{\vnbstmt{\begin{array}{@{}l@{}} \forall s, a. \\ \quad (\operatorname{interior}(s, a) \simeq \{x \in \operatorname{pts}(s) : \exists u.\; (\operatorname{is\text{-}open}(s, u) \wedge ((x \in u) \wedge (u \subseteq a)))\}) \end{array}}}
\vnbprop{closed-ball-membership}{\vnbstmt{\begin{array}{@{}l@{}} \forall s, c, r, y. \\ \quad (y \in \operatorname{closed\text{-}ball}(s, c, r)) \Leftrightarrow ((y \in \operatorname{pts}(s)) \wedge (\operatorname{dist}(s)(c, y) \leq r)) \end{array}}}
\vnbprop{closure-membership}{\vnbstmt{\begin{array}{@{}l@{}} \forall s, a, x. \\ \quad x \in \operatorname{closure}(s, a) \\ \quad \Leftrightarrow\; ((x \in \operatorname{pts}(s)) \wedge \forall u.\; ((\operatorname{is\text{-}open}(s, u) \wedge (x \in u)) \Rightarrow \exists y \in u.\; (y \in a))) \end{array}}}
\vnbprop{interior-membership}{\vnbstmt{\begin{array}{@{}l@{}} \forall s, a, x. \\ \quad (x \in \operatorname{interior}(s, a)) \Leftrightarrow ((x \in \operatorname{pts}(s)) \wedge \exists u.\; (\operatorname{is\text{-}open}(s, u) \wedge ((x \in u) \wedge (u \subseteq a)))) \end{array}}}
\vnbprop{closed-ball-is-set}{\vnbstmt{\forall s, c, r.\; (\operatorname{is\text{-}metric\text{-}space}(s) \Rightarrow (\operatorname{closed\text{-}ball}(s, c, r) \in \mathbf{Set}))}}
\vnbprop{closure-in-carrier}{\vnbstmt{\forall s, a.\; (\operatorname{is\text{-}metric\text{-}space}(s) \Rightarrow (\operatorname{closure}(s, a) \subseteq \operatorname{pts}(s)))}}
\vnbprop{interior-in-carrier}{\vnbstmt{\forall s, a.\; (\operatorname{is\text{-}metric\text{-}space}(s) \Rightarrow (\operatorname{interior}(s, a) \subseteq \operatorname{pts}(s)))}}
\vnbprop{interior-in-set}{\vnbstmt{\forall s, a.\; (\operatorname{is\text{-}metric\text{-}space}(s) \Rightarrow (\operatorname{interior}(s, a) \subseteq a))}}
\vnbprop{closure-contains}{\vnbstmt{\begin{array}{@{}l@{}} \forall s, a. \\ \quad (\operatorname{is\text{-}metric\text{-}space}(s) \wedge (a \subseteq \operatorname{pts}(s))) \Rightarrow (a \subseteq \operatorname{closure}(s, a)) \end{array}}}
\vnbprop{ball-in-closed-ball}{\vnbstmt{\begin{array}{@{}l@{}} \forall s, c, r. \\ \quad \operatorname{is\text{-}metric\text{-}space}(s) \Rightarrow (\operatorname{ball}(s, c, r) \subseteq \operatorname{closed\text{-}ball}(s, c, r)) \end{array}}}
\vnbprop{closed-ball-in-ball}{\vnbstmt{\begin{array}{@{}l@{}} \forall s.; \forall c \in \operatorname{pts}(s), r, q \in \mathbb{R}. \\ \quad (\operatorname{is\text{-}metric\text{-}space}(s) \wedge (r < q)) \Rightarrow (\operatorname{closed\text{-}ball}(s, c, r) \subseteq \operatorname{ball}(s, c, q)) \end{array}}}
\vnbprop{interior-mem-intro}{\vnbstmt{\begin{array}{@{}l@{}} \forall s, a, u, x. \\ \quad (\operatorname{is\text{-}open}(s, u) \wedge ((u \subseteq a) \wedge (x \in u))) \Rightarrow (x \in \operatorname{interior}(s, a)) \end{array}}}
\vnbprop{interior-is-open}{\vnbstmt{\forall s, a.\; (\operatorname{is\text{-}metric\text{-}space}(s) \Rightarrow \operatorname{is\text{-}open}(s, \operatorname{interior}(s, a)))}}
\vnbprop{closure-is-closed}{\vnbstmt{\forall s, a.\; (\operatorname{is\text{-}metric\text{-}space}(s) \Rightarrow \operatorname{is\text{-}closed}(s, \operatorname{closure}(s, a)))}}
\vnbprop{closed-ball-is-closed}{\vnbstmt{\begin{array}{@{}l@{}} \forall s.; \forall c \in \operatorname{pts}(s), r \in \mathbb{R}. \\ \quad \operatorname{is\text{-}metric\text{-}space}(s) \Rightarrow \operatorname{is\text{-}closed}(s, \operatorname{closed\text{-}ball}(s, c, r)) \end{array}}}
\vnbprop{closure-in-closed}{\vnbstmt{\forall s, a, t.\; ((\operatorname{is\text{-}closed}(s, t) \wedge (a \subseteq t)) \Rightarrow (\operatorname{closure}(s, a) \subseteq t))}}
\vnbprop{closure-of-ball-in-closed-ball}{\vnbstmt{\begin{array}{@{}l@{}} \forall s.; \forall c \in \operatorname{pts}(s), r \in \mathbb{R}. \\ \quad \operatorname{is\text{-}metric\text{-}space}(s) \Rightarrow (\operatorname{closure}(s, \operatorname{ball}(s, c, r)) \subseteq \operatorname{closed\text{-}ball}(s, c, r)) \end{array}}}
\vnbprop{closed-set-limit-in}{\vnbstmt{\begin{array}{@{}l@{}} \forall s, a, f, l. \\ \quad (\operatorname{is\text{-}closed}(s, a) \wedge (\operatorname{converges\text{-}to}(s, f, l) \wedge \forall n \in \mathbb{N}.\; (f(n) \in a))) \Rightarrow  \\ \quad \quad l \in a \end{array}}}
\vnbprop{closure-is-set}{\vnbstmt{\forall s, a.\; (\operatorname{is\text{-}metric\text{-}space}(s) \Rightarrow (\operatorname{closure}(s, a) \in \mathbf{Set}))}}
\vnbprop{interior-is-set}{\vnbstmt{\forall s, a.\; (\operatorname{is\text{-}metric\text{-}space}(s) \Rightarrow (\operatorname{interior}(s, a) \in \mathbf{Set}))}}
\vnbprop{nowhere-dense-ball-point}{\vnbstmt{\begin{array}{@{}l@{}} \forall s, a.; \forall c \in \operatorname{pts}(s), r. \\ \quad (\operatorname{is\text{-}nowhere\text{-}dense}(s, a) \wedge \operatorname{pos\text{-}rr}(r)) \Rightarrow  \\ \quad \quad \exists x \in \operatorname{ball}(s, c, r).\; \neg (x \in \operatorname{closure}(s, a)) \end{array}}}
\vnbprop{nowhere-dense-closed-ball}{\vnbstmt{\begin{array}{@{}l@{}} \forall s, a.; \forall c \in \operatorname{pts}(s), r, d. \\ \quad (\operatorname{is\text{-}nowhere\text{-}dense}(s, a) \wedge (\operatorname{pos\text{-}rr}(r) \wedge \operatorname{pos\text{-}rr}(d))) \Rightarrow  \\ \quad \quad \exists x \in \operatorname{pts}(s), q. \\ \quad \quad \quad \operatorname{pos\text{-}rr}(q) \wedge \\ \quad \quad \quad q \leq d \wedge \\ \quad \quad \quad \operatorname{closed\text{-}ball}(s, x, q) \subseteq \operatorname{ball}(s, c, r) \wedge \\ \quad \quad \quad \forall z \in \operatorname{closed\text{-}ball}(s, x, q).\; \neg (z \in \operatorname{closure}(s, a)) \end{array}}}

\vnbbookfile{theorem-library/rake-baire}{theorem-library/rake-baire}
\noindent{\small rake-baire.scm -- towards the BAIRE CATEGORY theorem (batch 12-F, 2026-09-20).}\par
\vnbprop{nested-family-monotone}{\vnbstmt{\begin{array}{@{}l@{}} \forall m \in \mathbb{N}, w.; \forall g \in (\mathbb{N} \to w), n \in \mathbb{N}. \\ \quad (\forall n \in \mathbb{N}.\; (g((n + 1)) \subseteq g(n)) \wedge (n \leq m)) \Rightarrow (g(m) \subseteq g(n)) \end{array}}}

\vnbbookfile{structure-library/frechet-open-mapping}{structure-library/frechet-open-mapping}
\noindent{\small MOVED 2026-09-20 (batch 12-A) from theorem-library/: this file is VOCABULARY -- definitions, notation and warranted supports, not one proof -- and every theorem stated with it had to load below it.  Its load.scm slot is unchanged. frechet-open-mapping.scm -- Frechet spaces, continuous linear maps, and the}\par
\par\medskip\noindent\textbf{Definition} (\texttt{fr-ball}).\;$\operatorname{fr\text{-}ball}(m, \mathit{fam}, n, \mathit{eps}) :=$
\begin{equation*}
\{x \in \operatorname{vec}(m) : \forall k.\; (((k \in \mathbb{N}) \wedge (k \leq n)) \Rightarrow (\operatorname{fam}(k)(x) < \mathit{eps}))\}
\end{equation*}
\par\noindent{\small\textit{Comment.} The pair $(m, \mathit{fam})$ is a seminormed space: $m$ a module, its vectors $\operatorname{vec}(m)$, and $\mathit{fam}(k)$ the $k$-th seminorm of a family enumerated by $\mathbb{N}$. The ball is cut out by the first $n + 1$ seminorms, radius $\varepsilon$.}\par

\par\medskip\noindent\textbf{Definition} (\texttt{has-closed-graph}).\;For all $m_{1}$, $\mathit{fam}_{1}$, $m_{2}$, $\mathit{fam}_{2}$ and $\mathit{tt}$, 
\begin{equation*}
\begin{array}{@{}l@{}} \operatorname{has\text{-}closed\text{-}graph}(m_{1}, \mathit{fam}_{1}, m_{2}, \mathit{fam}_{2}, \mathit{tt}) \\ \Leftrightarrow\; \operatorname{is\text{-}linear\text{-}map}(m_{1}, m_{2}, \mathit{tt}) \wedge \\ \forall \mathit{seq} \in (\mathbb{N} \to \operatorname{vec}(m_{1})), p_{1} \in \operatorname{vec}(m_{1}), p_{2} \in \operatorname{vec}(m_{2}). \\ \quad (\operatorname{fr\text{-}conv}(m_{1}, \mathit{fam}_{1}, \mathit{seq}, p_{1}) \wedge \operatorname{fr\text{-}conv}(m_{2}, \mathit{fam}_{2}, (\lambda i \in \mathbb{N}.\; \operatorname{tt}(\operatorname{seq}(i))), p_{2})) \Rightarrow  \\ \quad \quad p_{2} = \operatorname{tt}(p_{1}) \end{array}
\end{equation*}
\par\noindent{\small\textit{Comment.} $(m_1, \mathit{fam}_1)$ is the domain and $(m_2, \mathit{fam}_2)$ the target, each a module with an $\mathbb{N}$-enumerated seminorm family; $\mathit{tt}$ is the map.}\par

\par\medskip\noindent\textbf{Definition} (\texttt{is-seminorm-family}).\;For all $\mathit{fld}$, $m$ and $\mathit{fam}$, 
\begin{equation*}
\begin{array}{@{}l@{}} \operatorname{is\text{-}seminorm\text{-}family}(\mathit{fld}, m, \mathit{fam}) \\ \Leftrightarrow\; \operatorname{is\text{-}normed\text{-}field}(\mathit{fld}) \wedge \\ \operatorname{is\text{-}module}(m) \wedge \\ \operatorname{scal}(m) = \operatorname{normed\text{-}field\text{-}as\text{-}commutative\text{-}ring}(\mathit{fld}) \wedge \\ \mathit{fam} \in (\mathbb{N} \to (\operatorname{vec}(m) \to \mathbb{R})) \wedge \\ \forall k \in \mathbb{N}.\; \operatorname{is\text{-}seminorm}(\mathit{fld}, m, \operatorname{fam}(k)) \wedge \\ \forall x \in \operatorname{vec}(m). \\ \quad \neg (x = \operatorname{vzero}(m)) \Rightarrow \exists k \in \mathbb{N}.\; \neg (\operatorname{fam}(k)(x) = 0) \end{array}
\end{equation*}
\par\noindent{\small\textit{Comment.} $\mathit{fld}$ is the normed field of scalars, $m$ a module over it, and $\mathit{fam} : \mathbb{N} \to (\operatorname{vec}(m) \to \mathbb{R})$ the family of seminorms; separating: a nonzero vector has a nonzero seminorm.}\par

\par\medskip\noindent\textbf{Definition} (\texttt{is-frechet-structure}).\;For all $\mathit{fld}$, $m$ and $\mathit{fam}$, 
\begin{equation*}
\begin{array}{@{}l@{}} \operatorname{is\text{-}frechet\text{-}structure}(\mathit{fld}, m, \mathit{fam}) \\ \Leftrightarrow\; \operatorname{is\text{-}seminorm\text{-}family}(\mathit{fld}, m, \mathit{fam}) \wedge \\ \forall \mathit{seq} \in (\mathbb{N} \to \operatorname{vec}(m)). \\ \quad \operatorname{fr\text{-}cauchy}(m, \mathit{fam}, \mathit{seq}) \Rightarrow  \\ \quad \quad \exists \mathit{lim} \in \operatorname{vec}(m).\; \operatorname{fr\text{-}conv}(m, \mathit{fam}, \mathit{seq}, \mathit{lim}) \end{array}
\end{equation*}
\par\noindent{\small\textit{Comment.} $\mathit{fld}$ is the normed field of scalars; $(m, \mathit{fam})$ is the space, $m$ a module and $\mathit{fam}(k)$ its $k$-th seminorm.}\par

\par\medskip\noindent\textbf{Definition} (\texttt{is-linear-map}).\;For all $m_{1}$, $m_{2}$ and $\mathit{tt}$, 
\begin{equation*}
\begin{array}{@{}l@{}} \operatorname{is\text{-}linear\text{-}map}(m_{1}, m_{2}, \mathit{tt}) \\ \Leftrightarrow\; \operatorname{is\text{-}module}(m_{1}) \wedge \\ \operatorname{is\text{-}module}(m_{2}) \wedge \\ \operatorname{scal}(m_{1}) = \operatorname{scal}(m_{2}) \wedge \\ \mathit{tt} \in (\operatorname{vec}(m_{1}) \to \operatorname{vec}(m_{2})) \wedge \\ \forall x, y \in \operatorname{vec}(m_{1}). \\ \quad \operatorname{tt}(\operatorname{vadd}(m_{1})(x, y)) = \operatorname{vadd}(m_{2})(\operatorname{tt}(x), \operatorname{tt}(y)) \wedge \\ \forall \mathit{lam} \in \operatorname{carr}(\operatorname{scal}(m_{1})), x \in \operatorname{vec}(m_{1}). \\ \quad \operatorname{tt}(\operatorname{act}(m_{1})(\mathit{lam}, x)) = \operatorname{act}(m_{2})(\mathit{lam}, \operatorname{tt}(x)) \end{array}
\end{equation*}
\par\noindent{\small\textit{Comment.} $m_1$ and $m_2$ are modules over one scalar ring, $\mathit{tt} : \operatorname{vec}(m_1) \to \operatorname{vec}(m_2)$ the map.}\par

\par\medskip\noindent\textbf{Definition} (\texttt{fr-conv}).\;For all $m$, $\mathit{fam}$, $\mathit{seq}$ and $\mathit{lim}$, 
\begin{equation*}
\begin{array}{@{}l@{}} \operatorname{fr\text{-}conv}(m, \mathit{fam}, \mathit{seq}, \mathit{lim}) \\ \Leftrightarrow\; \forall k \in \mathbb{N}, \mathit{eps}. \\ \quad \operatorname{pos\text{-}rr}(\mathit{eps}) \Rightarrow  \\ \quad \quad \exists \mathit{cap} \in \mathbb{N}. \forall i. \\ \quad \quad \quad ((i \in \mathbb{N}) \wedge (\mathit{cap} \leq i)) \Rightarrow (\operatorname{fam}(k)(\operatorname{vadd}(m)(\operatorname{seq}(i), \operatorname{vneg}(m)(\mathit{lim}))) < \mathit{eps}) \end{array}
\end{equation*}
\par\noindent{\small\textit{Comment.} The pair $(m, \mathit{fam})$ is a seminormed space: $m$ a module, its vectors $\operatorname{vec}(m)$, and $\mathit{fam}(k)$ the $k$-th seminorm of a family enumerated by $\mathbb{N}$. $\mathit{seq}$ is a sequence of vectors and $\mathit{lim}$ the limit.}\par

\par\medskip\noindent\textbf{Definition} (\texttt{fr-cauchy}).\;For all $m$, $\mathit{fam}$ and $\mathit{seq}$, 
\begin{equation*}
\begin{array}{@{}l@{}} \operatorname{fr\text{-}cauchy}(m, \mathit{fam}, \mathit{seq}) \\ \Leftrightarrow\; \forall k \in \mathbb{N}, \mathit{eps}. \\ \quad \operatorname{pos\text{-}rr}(\mathit{eps}) \Rightarrow  \\ \quad \quad \exists \mathit{cap} \in \mathbb{N}. \forall i. \\ \quad \quad \quad ((i \in \mathbb{N}) \wedge (\mathit{cap} \leq i)) \Rightarrow  \\ \quad \quad \quad \quad \forall j. \\ \quad \quad \quad \quad \quad ((j \in \mathbb{N}) \wedge (\mathit{cap} \leq j)) \Rightarrow  \\ \quad \quad \quad \quad \quad \quad \operatorname{fam}(k)(\operatorname{vadd}(m)(\operatorname{seq}(i), \operatorname{vneg}(m)(\operatorname{seq}(j)))) < \mathit{eps} \end{array}
\end{equation*}
\par\noindent{\small\textit{Comment.} The pair $(m, \mathit{fam})$ is a seminormed space: $m$ a module, its vectors $\operatorname{vec}(m)$, and $\mathit{fam}(k)$ the $k$-th seminorm of a family enumerated by $\mathbb{N}$. $\mathit{seq}$ is a sequence of vectors.}\par

\par\medskip\noindent\textbf{Definition} (\texttt{is-cont-lin}).\;For all $m_{1}$, $\mathit{fam}_{1}$, $m_{2}$, $\mathit{fam}_{2}$ and $\mathit{tt}$, 
\begin{equation*}
\begin{array}{@{}l@{}} \operatorname{is\text{-}cont\text{-}lin}(m_{1}, \mathit{fam}_{1}, m_{2}, \mathit{fam}_{2}, \mathit{tt}) \\ \Leftrightarrow\; \operatorname{is\text{-}linear\text{-}map}(m_{1}, m_{2}, \mathit{tt}) \wedge \\ \forall \mathit{seq} \in (\mathbb{N} \to \operatorname{vec}(m_{1})), \mathit{lim} \in \operatorname{vec}(m_{1}). \\ \quad \operatorname{fr\text{-}conv}(m_{1}, \mathit{fam}_{1}, \mathit{seq}, \mathit{lim}) \Rightarrow  \\ \quad \quad \operatorname{fr\text{-}conv}(m_{2}, \mathit{fam}_{2}, (\lambda i \in \mathbb{N}.\; \operatorname{tt}(\operatorname{seq}(i))), \operatorname{tt}(\mathit{lim})) \end{array}
\end{equation*}
\par\noindent{\small\textit{Comment.} $(m_1, \mathit{fam}_1)$ is the domain and $(m_2, \mathit{fam}_2)$ the target, each a module with an $\mathbb{N}$-enumerated seminorm family; $\mathit{tt}$ is the map. Continuity is sequential.}\par

\par\medskip\noindent\textbf{Definition} (\texttt{is-open-lin-map}).\;For all $m_{1}$, $\mathit{fam}_{1}$, $m_{2}$, $\mathit{fam}_{2}$ and $\mathit{tt}$, 
\begin{equation*}
\begin{array}{@{}l@{}} \operatorname{is\text{-}open\text{-}lin\text{-}map}(m_{1}, \mathit{fam}_{1}, m_{2}, \mathit{fam}_{2}, \mathit{tt}) \\ \Leftrightarrow\; \operatorname{is\text{-}linear\text{-}map}(m_{1}, m_{2}, \mathit{tt}) \wedge \\ \forall n \in \mathbb{N}, \mathit{eps}. \\ \quad \operatorname{pos\text{-}rr}(\mathit{eps}) \Rightarrow  \\ \quad \quad \exists \operatorname{n\_} \in \mathbb{N}, \operatorname{eps\_}. \\ \quad \quad \quad (\operatorname{pos\text{-}rr}(\operatorname{eps\_}) \wedge (\operatorname{fr\text{-}ball}(m_{2}, \mathit{fam}_{2}, \operatorname{n\_}, \operatorname{eps\_}) \subseteq \operatorname{image}(\mathit{tt}, \operatorname{fr\text{-}ball}(m_{1}, \mathit{fam}_{1}, n, \mathit{eps})))) \end{array}
\end{equation*}
\par\noindent{\small\textit{Comment.} $(m_1, \mathit{fam}_1)$ is the domain and $(m_2, \mathit{fam}_2)$ the target, each a module with an $\mathbb{N}$-enumerated seminorm family; $\mathit{tt}$ is the map.}\par


\vnbbookfile{structure-library/frechet-metric}{structure-library/frechet-metric}
\noindent{\small frechet-metric.scm -- the METRIC of a seminormed space whose seminorm family is COUNTABLE (Thayer, "Analise Funcional e Teoria Espectral", Ch. 8 s.4, Prop. 8.24). Added 2026-10-10 (batch TVS-1).  Vocabulary only: the theorems are in theorem-library/frechet-metric.scm.}\par
\par\medskip\noindent\textbf{Definition} (\texttt{fr-c-space}).\;$\operatorname{fr\text{-}c\text{-}space}(m, \mathit{fam}) :=$
\begin{equation*}
\langle \operatorname{vec}(m), (\lambda \operatorname{frk\_} \in \mathbb{N}.\; (\lambda \langle \operatorname{fru\_}, \operatorname{frv\_} \rangle \in (\operatorname{vec}(m) \times \operatorname{vec}(m)).\; \operatorname{fam}(\operatorname{frk\_})(\operatorname{vadd}(m)(\operatorname{fru\_}, \operatorname{vneg}(m)(\operatorname{frv\_}))))) \rangle
\end{equation*}

\par\medskip\noindent\textbf{Definition} (\texttt{fr-metric-space}).\;$\operatorname{fr\text{-}metric\text{-}space}(m, \mathit{fam}) :=$ $\operatorname{c\text{-}metric}(\operatorname{fr\text{-}c\text{-}space}(m, \mathit{fam}))$
\par\medskip\noindent\textbf{Definition} (\texttt{fr-metric}).\;$\operatorname{fr\text{-}metric}(m, \mathit{fam}) :=$ $\operatorname{dist}(\operatorname{fr\text{-}metric\text{-}space}(m, \mathit{fam}))$

\vnbbookfile{structure-library/frechet-banach-steinhaus}{structure-library/frechet-banach-steinhaus}
\noindent{\small frechet-banach-steinhaus.scm -- BOUNDED SETS and EQUICONTINUOUS FAMILIES of linear maps in a space with a countable family of seminorms (Thayer, "Analise Funcional e Teoria Espectral", Ch. 8 s.4, Def. 8.31 and the equicontinuity of Thm 8.32, p. 257). Added 2026-10-11 (batch TVS-2B).  Vocabulary only: the theorems ...}\par
\par\medskip\noindent\textbf{Definition} (\texttt{is-fr-bounded}).\;For all $m$, $\mathit{fam}$ and $\mathit{aa}$, 
\begin{equation*}
\begin{array}{@{}l@{}} \operatorname{is\text{-}fr\text{-}bounded}(m, \mathit{fam}, \mathit{aa}) \\ \Leftrightarrow\; ((\mathit{aa} \subseteq \operatorname{vec}(m)) \wedge \forall k \in \mathbb{N}.\; \exists \mathit{bd} \in \mathbb{R}.\; \forall y \in \mathit{aa}.\; (\operatorname{fam}(k)(y) \leq \mathit{bd})) \end{array}
\end{equation*}
\par\noindent{\small\textit{Comment.} $(m, \mathit{fam})$ is the seminormed space ($m$ a module, $\mathit{fam}(k)$ its $k$-th seminorm); $\mathit{aa}$ is the set, a subset of $\operatorname{vec}(m)$.}\par

\par\medskip\noindent\textbf{Definition} (\texttt{is-fr-equicont}).\;For all $m_{1}$, $\mathit{fam}_{1}$, $m_{2}$, $\mathit{fam}_{2}$, $\mathit{ix}$ and $\mathit{tf}$, 
\begin{equation*}
\begin{array}{@{}l@{}} \operatorname{is\text{-}fr\text{-}equicont}(m_{1}, \mathit{fam}_{1}, m_{2}, \mathit{fam}_{2}, \mathit{ix}, \mathit{tf}) \\ \Leftrightarrow\; \mathit{tf} \in (\mathit{ix} \to (\operatorname{vec}(m_{1}) \to \operatorname{vec}(m_{2}))) \wedge \\ \forall k \in \mathbb{N}. \exists \mathit{nb} \in \mathbb{N}, \mathit{ck}. \\ \quad \operatorname{pos\text{-}rr}(\mathit{ck}) \wedge \\ \quad \forall i \in \mathit{ix}, x \in \operatorname{vec}(m_{1}), r \in \mathbb{R}. \\ \quad \quad \forall j \in \mathbb{N}.\; ((j \leq \mathit{nb}) \Rightarrow (\operatorname{fam1}(j)(x) \leq r)) \Rightarrow  \\ \quad \quad \quad \operatorname{fam2}(k)(\operatorname{tf}(i)(x)) \leq \cdot(\mathit{ck}, r) \end{array}
\end{equation*}
\par\noindent{\small\textit{Comment.} $(m_1, \mathit{fam}_1)$ is the domain and $(m_2, \mathit{fam}_2)$ the target, each a module with an $\mathbb{N}$-enumerated seminorm family; $\mathit{ix}$ is the index class and $\mathit{tf}(i)$, for $i$ in $\mathit{ix}$, the $i$-th map of the family.}\par


\vnbbookfile{theorem-library/subseq-convergence-proof}{theorem-library/subseq-convergence-proof}
\noindent{\small theorem-library/subseq-convergence-proof.scm ==================================================================== subseq-of-convergent -- PROVEN.}\par
\vnbprop{subseq-of-convergent}{\vnbstmt{\begin{array}{@{}l@{}} \forall s, f, i, l. \\ \quad (\operatorname{converges\text{-}to}(s, f, l) \wedge \operatorname{strictly\text{-}mono\text{-}nn}(i)) \Rightarrow  \\ \quad \quad \operatorname{converges\text{-}to}(s, \operatorname{subseq}(f, i), l) \end{array}}}

\vnbbookfile{theorem-library/subsqn-basics}{theorem-library/subsqn-basics}
\noindent{\small subsqn-basics.scm -- the SUBSQN relation and what is proven about it.}\par
\vnbprop{is-subsequence-iff-typed-subsqn}{\vnbstmt{\begin{array}{@{}l@{}} \forall s, y, f. \\ \quad \operatorname{is\text{-}subsequence}(s, y, f) \Leftrightarrow ((f \in (\mathbb{N} \to \operatorname{pts}(s))) \wedge \operatorname{subsqn}(y, f)) \end{array}}}
\vnbprop{subsqn-preserves-sqn}{\vnbstmt{\forall a, f, y.\; (((f \in a^{\mathbb{N}}) \wedge \operatorname{subsqn}(y, f)) \Rightarrow (y \in a^{\mathbb{N}}))}}

\vnbbookfile{theorem-library/subsequence-principle}{theorem-library/subsequence-principle}
\noindent{\small theorem-library/subsequence-principle.scm ====================================================================== subsequence-principle -- PROVEN.}\par
\vnbprop{subsequence-principle}{\vnbstmt{\begin{array}{@{}l@{}} \forall s, f, l. \\ \quad \operatorname{is\text{-}metric\text{-}space}(s) \wedge \\ \quad f \in (\mathbb{N} \to \operatorname{pts}(s)) \wedge \\ \quad l \in \operatorname{pts}(s) \wedge \\ \quad \forall i. \\ \quad \quad \operatorname{strictly\text{-}mono\text{-}nn}(i) \Rightarrow  \\ \quad \quad \quad \exists a. \\ \quad \quad \quad \quad \operatorname{strictly\text{-}mono\text{-}nn}(a) \wedge \operatorname{converges\text{-}to}(s, \operatorname{subseq}(\operatorname{subseq}(f, i), a), l) \Rightarrow \\ \quad \quad \operatorname{converges\text{-}to}(s, f, l) \end{array}}}

\vnbbookfile{theorem-library/rake-offbill-combinatorial}{theorem-library/rake-offbill-combinatorial}
\noindent{\small theorem-library/rake-offbill-combinatorial.scm ==================================================================== Batch 8-D: the off-bill `combinatorial' PSS entries that batch 6/7 made cheap.  Three theorems, in dependency order:}\par
\vnbprop{tb-block-step}{\vnbstmt{\begin{array}{@{}l@{}} \forall s.; \forall f \in (\mathbb{N} \to \operatorname{pts}(s)), r.; \forall j \in \mathrm{Inf}(\mathbb{N}). \\ \quad (\operatorname{totally\text{-}bounded}(s) \wedge \operatorname{pos\text{-}rr}(r)) \Rightarrow  \\ \quad \quad \exists a \in \mathrm{Inf}(\mathbb{N}). \\ \quad \quad \quad (a \subseteq j) \wedge \exists c \in \operatorname{pts}(s).\; \forall i \in a.\; (f(i) \in \operatorname{ball}(s, c, r)) \end{array}}}
\vnbprop{tb-has-eps-cauchy-subseq}{\vnbstmt{\begin{array}{@{}l@{}} \forall s.; \forall f \in (\mathbb{N} \to \operatorname{pts}(s)), a. \\ \quad (\operatorname{totally\text{-}bounded}(s) \wedge \operatorname{pos\text{-}rr}(a)) \Rightarrow  \\ \quad \quad \exists i.\; (\operatorname{strictly\text{-}mono\text{-}nn}(i) \wedge \operatorname{is\text{-}eps\text{-}cauchy\text{-}seq}(s, a, \operatorname{subseq}(f, i))) \end{array}}}
\vnbprop{block-family}{\vnbstmt{\begin{array}{@{}l@{}} \forall s.; \forall f \in (\mathbb{N} \to \operatorname{pts}(s)), d. \\ \quad (\operatorname{totally\text{-}bounded}(s) \wedge \operatorname{null\text{-}rr\text{-}seq}(d)) \Rightarrow  \\ \quad \quad \exists a \in (\mathbb{N} \to \mathrm{Inf}(\mathbb{N})). \\ \quad \quad \quad \forall k \in \mathbb{N}.\; (a((k + 1)) \subseteq a(k)) \wedge \\ \quad \quad \quad \forall k \in \mathbb{N}. \exists c \in \operatorname{pts}(s). \forall i \in a(k). \\ \quad \quad \quad \quad f(i) \in \operatorname{ball}(s, c, d(k)) \end{array}}}

\vnbbookfile{theorem-library/rake-baire-2}{theorem-library/rake-baire-2}
\noindent{\small rake-baire-2.scm -- the BAIRE CATEGORY theorem (batch 13-C, 2026-09-20).}\par
\vnbprop{closed-ball-in-carrier}{\vnbstmt{\begin{array}{@{}l@{}} \forall s, v, \mathit{qv}. \\ \quad \operatorname{is\text{-}metric\text{-}space}(s) \Rightarrow (\operatorname{closed\text{-}ball}(s, v, \mathit{qv}) \subseteq \operatorname{pts}(s)) \end{array}}}
\vnbprop{closed-ball-center-in}{\vnbstmt{\begin{array}{@{}l@{}} \forall s.; \forall v \in \operatorname{pts}(s), \mathit{qv} \in \mathbb{R}. \\ \quad (\operatorname{is\text{-}metric\text{-}space}(s) \wedge (0 \leq \mathit{qv})) \Rightarrow (v \in \operatorname{closed\text{-}ball}(s, v, \mathit{qv})) \end{array}}}
\vnbprop{closed-ball-in-power}{\vnbstmt{\begin{array}{@{}l@{}} \forall s, v, \mathit{qv}. \\ \quad \operatorname{is\text{-}metric\text{-}space}(s) \Rightarrow (\operatorname{closed\text{-}ball}(s, v, \mathit{qv}) \in \operatorname{power}(\operatorname{pts}(s))) \end{array}}}
\vnbprop{closed-set-limit-in-tail}{\vnbstmt{\begin{array}{@{}l@{}} \forall s, v, \mathit{fv}, \mathit{lv}, b. \\ \quad \operatorname{is\text{-}closed}(s, v) \wedge \\ \quad \operatorname{converges\text{-}to}(s, \mathit{fv}, \mathit{lv}) \wedge \\ \quad b \in \mathbb{N} \wedge \\ \quad \forall n \in \mathbb{N}.\; ((b \leq n) \Rightarrow (\operatorname{fv}(n) \in v)) \Rightarrow \\ \quad \quad \mathit{lv} \in v \end{array}}}
\vnbprop{nested-closed-balls-point}{\vnbstmt{\begin{array}{@{}l@{}} \forall s.; \forall r \in (\mathbb{N} \to \operatorname{pts}(s)), \mathit{rds} \in (\mathbb{N} \to \mathbb{R}). \\ \quad \operatorname{is\text{-}complete}(s) \wedge \\ \quad \forall n \in \mathbb{N}.\; (0 \leq \operatorname{rds}(n)) \wedge \\ \quad \forall n \in \mathbb{N}. \\ \quad \quad \operatorname{closed\text{-}ball}(s, r((n + 1)), \operatorname{rds}((n + 1))) \subseteq \operatorname{closed\text{-}ball}(s, r(n), \operatorname{rds}(n)) \wedge \\ \quad \forall n \in \mathbb{N}.\; (\operatorname{rds}((n + 1)) \leq \operatorname{recip}((n + 1))) \Rightarrow \\ \quad \quad \exists v \in \operatorname{pts}(s).\; \forall n \in \mathbb{N}.\; (v \in \operatorname{closed\text{-}ball}(s, r(n), \operatorname{rds}(n))) \end{array}}}
\vnbprop{baire-category}{\vnbstmt{\begin{array}{@{}l@{}} \forall s. \\ \quad (\operatorname{is\text{-}complete}(s) \wedge \exists x.\; (x \in \operatorname{pts}(s))) \Rightarrow \operatorname{is\text{-}nonmeager}(s, \operatorname{pts}(s)) \end{array}}}

\vnbbookfile{theorem-library/coord-block-estimate-proof}{theorem-library/coord-block-estimate-proof}
\noindent{\small theorem-library/coord-block-estimate-proof.scm ==================================================================== coord-block-estimate -- PROVEN.}\par
\vnbprop{coord-block-estimate}{\vnbstmt{\begin{array}{@{}l@{}} \forall s, g, b, p, a, c. \\ \quad \operatorname{converges\text{-}along}(s, g, b, p) \wedge \\ \quad \operatorname{strictly\text{-}mono\text{-}nn}(a) \wedge \\ \quad c \in \mathbb{N} \wedge \\ \quad \forall j \in \mathbb{N}.\; ((c \leq j) \Rightarrow (a(j) \in b)) \Rightarrow \\ \quad \quad \operatorname{converges\text{-}to}(s, \operatorname{subseq}(g, a), p) \end{array}}}

\vnbbookfile{theorem-library/sqrt2-proof}{theorem-library/sqrt2-proof}
\noindent{\small sqrt2-proof.scm -- sqrt(2) is irrational, stated on the naturals:}\par
\vnbprop{sqrt2-irrational}{\vnbstmt{\forall p, q \in \mathbb{N}.\; (\neg (q = 0) \Rightarrow \neg ((p \cdot p) = (2 \cdot (q \cdot q))))}}

\vnbbookfile{theorem-library/sqrt3-proof}{theorem-library/sqrt3-proof}
\noindent{\small sqrt3-proof.scm -- sqrt(3) is irrational, on the naturals: forall p,q in NN.  q /= 0  =$>$  p*p /= 3*(q*q).}\par
\vnbprop{sqrt3-irrational}{\vnbstmt{\forall p, q \in \mathbb{N}.\; (\neg (q = 0) \Rightarrow \neg ((p \cdot p) = (3 \cdot (q \cdot q))))}}

\vnbbookfile{theorem-library/continuity-basics}{theorem-library/continuity-basics}
\noindent{\small continuity-basics.scm -- the constant and identity maps are continuous, PROVEN, with their FUN typings.}\par
\vnbprop{const-lam-in-fun}{\vnbstmt{\forall c \in \mathbb{R}.\; ((\lambda x \in \mathbb{R}.\; c) \in (\mathbb{R} \to \mathbb{R}))}}
\vnbprop{ident-lam-in-fun}{\vnbstmt{(\lambda x \in \mathbb{R}.\; x) \in (\mathbb{R} \to \mathbb{R})}}
\vnbprop{const-continuous-at}{\vnbstmt{\begin{array}{@{}l@{}} \forall c, a \in \mathbb{R}. \\ \quad \operatorname{is\text{-}continuous\text{-}at}(\operatorname{rr\text{-}ms}, \operatorname{rr\text{-}ms}, (\lambda x \in \mathbb{R}.\; c), a) \end{array}}}
\vnbprop{identity-continuous-at}{\vnbstmt{\forall a \in \mathbb{R}.\; \operatorname{is\text{-}continuous\text{-}at}(\operatorname{rr\text{-}ms}, \operatorname{rr\text{-}ms}, (\lambda x \in \mathbb{R}.\; x), a)}}

\vnbbookfile{theorem-library/neg-continuous}{theorem-library/neg-continuous}
\noindent{\small neg-continuous.scm -- NEGATION preserves membership in FUN(RR,RR) and continuity, PROVEN.}\par
\vnbprop{neg-fun-in-fun}{\vnbstmt{\forall f \in (\mathbb{R} \to \mathbb{R}).\; ((\lambda z \in \mathbb{R}.\; -f(z)) \in (\mathbb{R} \to \mathbb{R}))}}
\vnbprop{neg-continuous-at}{\vnbstmt{\begin{array}{@{}l@{}} \forall f, x. \\ \quad ((f \in (\mathbb{R} \to \mathbb{R})) \wedge ((x \in \mathbb{R}) \wedge \operatorname{is\text{-}continuous\text{-}at}(\operatorname{rr\text{-}ms}, \operatorname{rr\text{-}ms}, f, x))) \Rightarrow  \\ \quad \quad \operatorname{is\text{-}continuous\text{-}at}(\operatorname{rr\text{-}ms}, \operatorname{rr\text{-}ms}, (\lambda z \in \mathbb{R}.\; -f(z)), x) \end{array}}}

\vnbbookfile{theorem-library/continuity-sum}{theorem-library/continuity-sum}
\noindent{\small continuity-sum.scm -- the pointwise SUM of two maps continuous at a point is continuous there, PROVEN, with its FUN typing.}\par
\vnbprop{sum-lam-in-fun}{\vnbstmt{\begin{array}{@{}l@{}} \forall g, h. \\ \quad ((g \in (\mathbb{R} \to \mathbb{R})) \wedge (h \in (\mathbb{R} \to \mathbb{R}))) \Rightarrow  \\ \quad \quad (\lambda x \in \mathbb{R}.\; (g(x) + h(x))) \in (\mathbb{R} \to \mathbb{R}) \end{array}}}
\vnbprop{sum-continuous-at}{\vnbstmt{\begin{array}{@{}l@{}} \forall g, h, a. \\ \quad (\operatorname{is\text{-}continuous\text{-}at}(\operatorname{rr\text{-}ms}, \operatorname{rr\text{-}ms}, g, a) \wedge \operatorname{is\text{-}continuous\text{-}at}(\operatorname{rr\text{-}ms}, \operatorname{rr\text{-}ms}, h, a)) \Rightarrow  \\ \quad \quad \operatorname{is\text{-}continuous\text{-}at}(\operatorname{rr\text{-}ms}, \operatorname{rr\text{-}ms}, (\lambda x \in \mathbb{R}.\; (g(x) + h(x))), a) \end{array}}}

\vnbbookfile{theorem-library/continuity-product}{theorem-library/continuity-product}
\noindent{\small continuity-product.scm -- the pointwise PRODUCT of two maps continuous at a point is continuous there, PROVEN, with its FUN typing.}\par
\vnbprop{prod-lam-in-fun}{\vnbstmt{\begin{array}{@{}l@{}} \forall g, h. \\ \quad ((g \in (\mathbb{R} \to \mathbb{R})) \wedge (h \in (\mathbb{R} \to \mathbb{R}))) \Rightarrow  \\ \quad \quad (\lambda x \in \mathbb{R}.\; (g(x) \cdot h(x))) \in (\mathbb{R} \to \mathbb{R}) \end{array}}}
\vnbprop{product-continuous-at}{\vnbstmt{\begin{array}{@{}l@{}} \forall g, h, a. \\ \quad (\operatorname{is\text{-}continuous\text{-}at}(\operatorname{rr\text{-}ms}, \operatorname{rr\text{-}ms}, g, a) \wedge \operatorname{is\text{-}continuous\text{-}at}(\operatorname{rr\text{-}ms}, \operatorname{rr\text{-}ms}, h, a)) \Rightarrow  \\ \quad \quad \operatorname{is\text{-}continuous\text{-}at}(\operatorname{rr\text{-}ms}, \operatorname{rr\text{-}ms}, (\lambda x \in \mathbb{R}.\; (g(x) \cdot h(x))), a) \end{array}}}

\vnbbookfile{theorem-library/continuity-transfer}{theorem-library/continuity-transfer}
\noindent{\small continuity-transfer.scm -- continuity reads only the VALUES: a map that agrees pointwise with a map continuous at a is itself continuous at a. PROVEN.}\par
\vnbprop{cont-transfer-ptwise-eq}{\vnbstmt{\begin{array}{@{}l@{}} \forall f, g, a. \\ \quad ((f \in (\mathbb{R} \to \mathbb{R})) \wedge (\operatorname{is\text{-}continuous\text{-}at}(\operatorname{rr\text{-}ms}, \operatorname{rr\text{-}ms}, g, a) \wedge \forall x \in \mathbb{R}.\; (f(x) = g(x)))) \Rightarrow  \\ \quad \quad \operatorname{is\text{-}continuous\text{-}at}(\operatorname{rr\text{-}ms}, \operatorname{rr\text{-}ms}, f, a) \end{array}}}

\vnbbookfile{theorem-library/continuity-sub}{theorem-library/continuity-sub}
\noindent{\small continuity-sub.scm -- the pointwise DIFFERENCE of two maps continuous at a point is continuous there, PROVEN, with its FUN typing.}\par
\vnbprop{sub-lam-in-fun}{\vnbstmt{\begin{array}{@{}l@{}} \forall g, h. \\ \quad ((g \in (\mathbb{R} \to \mathbb{R})) \wedge (h \in (\mathbb{R} \to \mathbb{R}))) \Rightarrow  \\ \quad \quad (\lambda x \in \mathbb{R}.\; (g(x) - h(x))) \in (\mathbb{R} \to \mathbb{R}) \end{array}}}
\vnbprop{sub-continuous-at}{\vnbstmt{\begin{array}{@{}l@{}} \forall g, h, a. \\ \quad (\operatorname{is\text{-}continuous\text{-}at}(\operatorname{rr\text{-}ms}, \operatorname{rr\text{-}ms}, g, a) \wedge \operatorname{is\text{-}continuous\text{-}at}(\operatorname{rr\text{-}ms}, \operatorname{rr\text{-}ms}, h, a)) \Rightarrow  \\ \quad \quad \operatorname{is\text{-}continuous\text{-}at}(\operatorname{rr\text{-}ms}, \operatorname{rr\text{-}ms}, (\lambda x \in \mathbb{R}.\; (g(x) - h(x))), a) \end{array}}}

\vnbbookfile{theorem-library/continuity-scale}{theorem-library/continuity-scale}
\noindent{\small continuity-scale.scm -- the SCALAR MULTIPLE of a real map: its FUN typing and its continuity, both PROVEN `modulo 0'.}\par
\vnbprop{scale-lam-in-fun}{\vnbstmt{\begin{array}{@{}l@{}} \forall c, f. \\ \quad ((c \in \mathbb{R}) \wedge (f \in (\mathbb{R} \to \mathbb{R}))) \Rightarrow ((\lambda x \in \mathbb{R}.\; (c \cdot f(x))) \in (\mathbb{R} \to \mathbb{R})) \end{array}}}
\vnbprop{scale-continuous-at}{\vnbstmt{\begin{array}{@{}l@{}} \forall c \in \mathbb{R}, f, a. \\ \quad \operatorname{is\text{-}continuous\text{-}at}(\operatorname{rr\text{-}ms}, \operatorname{rr\text{-}ms}, f, a) \Rightarrow  \\ \quad \quad \operatorname{is\text{-}continuous\text{-}at}(\operatorname{rr\text{-}ms}, \operatorname{rr\text{-}ms}, (\lambda x \in \mathbb{R}.\; (c \cdot f(x))), a) \end{array}}}

\vnbbookfile{theorem-library/continuity-compose}{theorem-library/continuity-compose}
\noindent{\small continuity-compose.scm -- COMPOSITION OF CONTINUOUS MAPS, PROVEN.}\par
\vnbprop{compose-continuous-at}{\vnbstmt{\begin{array}{@{}l@{}} \forall g, f, a. \\ \quad (\operatorname{is\text{-}continuous\text{-}at}(\operatorname{rr\text{-}ms}, \operatorname{rr\text{-}ms}, f, a) \wedge \operatorname{is\text{-}continuous\text{-}at}(\operatorname{rr\text{-}ms}, \operatorname{rr\text{-}ms}, g, f(a))) \Rightarrow  \\ \quad \quad \operatorname{is\text{-}continuous\text{-}at}(\operatorname{rr\text{-}ms}, \operatorname{rr\text{-}ms}, \operatorname{compose}(g, f), a) \end{array}}}

\vnbbookfile{theorem-library/continuity-recip}{theorem-library/continuity-recip}
\noindent{\small continuity-recip.scm -- THE RECIPROCAL OF A NON-VANISHING CONTINUOUS MAP, PROVEN `modulo 0'.}\par
\vnbprop{recip-lam-in-fun}{\vnbstmt{\begin{array}{@{}l@{}} \forall h \in (\mathbb{R} \to \mathbb{R}). \\ \quad \forall x \in \mathbb{R}.\; \neg (h(x) = 0) \Rightarrow  \\ \quad \quad (\lambda x \in \mathbb{R}.\; \operatorname{recip}(h(x))) \in (\mathbb{R} \to \mathbb{R}) \end{array}}}
\vnbprop{recip-continuous-at}{\vnbstmt{\begin{array}{@{}l@{}} \forall h, p. \\ \quad ((h \in (\mathbb{R} \to \mathbb{R})) \wedge (\forall x \in \mathbb{R}.\; \neg (h(x) = 0) \wedge \operatorname{is\text{-}continuous\text{-}at}(\operatorname{rr\text{-}ms}, \operatorname{rr\text{-}ms}, h, p))) \Rightarrow  \\ \quad \quad \operatorname{is\text{-}continuous\text{-}at}(\operatorname{rr\text{-}ms}, \operatorname{rr\text{-}ms}, (\lambda x \in \mathbb{R}.\; \operatorname{recip}(h(x))), p) \end{array}}}

\vnbbookfile{theorem-library/cont-agree-off-pt}{theorem-library/cont-agree-off-pt}
\noindent{\small cont-agree-off-pt.scm -- two functions continuous at a point and agreeing AWAY from it agree AT it, PROVEN.}\par
\vnbprop{cont-agree-off-pt}{\vnbstmt{\begin{array}{@{}l@{}} \forall a, b.; \forall t \in \mathbb{R}. \\ \quad \operatorname{is\text{-}continuous\text{-}at}(\operatorname{rr\text{-}ms}, \operatorname{rr\text{-}ms}, a, t) \wedge \\ \quad \operatorname{is\text{-}continuous\text{-}at}(\operatorname{rr\text{-}ms}, \operatorname{rr\text{-}ms}, b, t) \wedge \\ \quad \forall x \in \mathbb{R}.\; (\neg (x = t) \Rightarrow (a(x) = b(x))) \Rightarrow \\ \quad \quad a(t) = b(t) \end{array}}}

\vnbbookfile{theorem-library/continuous-one-sided-sign}{theorem-library/continuous-one-sided-sign}
\noindent{\small continuous-one-sided-sign.scm -- sign preservation under continuity, from ONE side, PROVEN.  Two theorems, one driver parameterised by the side:}\par
\vnbprop{continuous-nonneg-left}{\vnbstmt{\begin{array}{@{}l@{}} \forall g, a, h. \\ \quad g \in (\mathbb{R} \to \mathbb{R}) \wedge \\ \quad a \in \mathbb{R} \wedge \\ \quad h \in \mathbb{R} \wedge \\ \quad a < h \wedge \\ \quad \operatorname{is\text{-}continuous\text{-}at}(\operatorname{rr\text{-}ms}, \operatorname{rr\text{-}ms}, g, h) \wedge \\ \quad \forall x \in \mathbb{R}.\; (((a < x) \wedge (x < h)) \Rightarrow (0 \leq g(x))) \Rightarrow \\ \quad \quad 0 \leq g(h) \end{array}}}
\vnbprop{continuous-nonpos-right}{\vnbstmt{\begin{array}{@{}l@{}} \forall g, h, b. \\ \quad g \in (\mathbb{R} \to \mathbb{R}) \wedge \\ \quad h \in \mathbb{R} \wedge \\ \quad b \in \mathbb{R} \wedge \\ \quad h < b \wedge \\ \quad \operatorname{is\text{-}continuous\text{-}at}(\operatorname{rr\text{-}ms}, \operatorname{rr\text{-}ms}, g, h) \wedge \\ \quad \forall x \in \mathbb{R}.\; (((h < x) \wedge (x < b)) \Rightarrow (g(x) \leq 0)) \Rightarrow \\ \quad \quad g(h) \leq 0 \end{array}}}

\vnbbookfile{theorem-library/differentiation}{theorem-library/differentiation}
\noindent{\small differentiation.scm -- Chapter 2 (Differentiation) of calculus.pdf, founded on the CARATHEODORY / o(h) ALGEBRAIC definition (user's choice 2026-06-27).}\par
\vnbprop{diff-implies-continuous}{\vnbstmt{\begin{array}{@{}l@{}} \forall f, a, l. \\ \quad \operatorname{is\text{-}diff\text{-}at}(f, a, l) \Rightarrow \operatorname{is\text{-}continuous\text{-}at}(\operatorname{rr\text{-}ms}, \operatorname{rr\text{-}ms}, f, a) \end{array}}}
\vnbprop{derivative-unique}{\vnbstmt{\forall f, a, l, m.\; ((\operatorname{is\text{-}diff\text{-}at}(f, a, l) \wedge \operatorname{is\text{-}diff\text{-}at}(f, a, m)) \Rightarrow (l = m))}}
\vnbprop{deriv-const}{\vnbstmt{\begin{array}{@{}l@{}} \forall c, a. \\ \quad ((c \in \mathbb{R}) \wedge (a \in \mathbb{R})) \Rightarrow \operatorname{is\text{-}diff\text{-}at}((\lambda x \in \mathbb{R}.\; c), a, 0) \end{array}}}
\vnbprop{deriv-identity}{\vnbstmt{\forall a \in \mathbb{R}.\; \operatorname{is\text{-}diff\text{-}at}((\lambda x \in \mathbb{R}.\; x), a, 1)}}

\vnbbookfile{theorem-library/diff-transfer}{theorem-library/diff-transfer}
\noindent{\small diff-transfer.scm -- differentiability reads only the VALUES: a map that agrees pointwise with a map differentiable at a is itself differentiable at a, with the same derivative.  PROVEN, `modulo 0'.}\par
\vnbprop{diff-transfer-ptwise-eq}{\vnbstmt{\begin{array}{@{}l@{}} \forall f, g, a, l. \\ \quad ((f \in (\mathbb{R} \to \mathbb{R})) \wedge (\operatorname{is\text{-}diff\text{-}at}(g, a, l) \wedge \forall x \in \mathbb{R}.\; (f(x) \simeq g(x)))) \Rightarrow  \\ \quad \quad \operatorname{is\text{-}diff\text{-}at}(f, a, l) \end{array}}}

\vnbbookfile{theorem-library/deriv-sum-product}{theorem-library/deriv-sum-product}
\noindent{\small deriv-sum-product.scm -- calculus.pdf Prop 2.5 (the SUM rule) and Prop 2.6 (the PRODUCT rule), PROVEN, both `modulo 0'.}\par
\vnbprop{deriv-sum}{\vnbstmt{\begin{array}{@{}l@{}} \forall f, g, a, l, m. \\ \quad (\operatorname{is\text{-}diff\text{-}at}(f, a, l) \wedge \operatorname{is\text{-}diff\text{-}at}(g, a, m)) \Rightarrow  \\ \quad \quad \operatorname{is\text{-}diff\text{-}at}((\lambda x \in \mathbb{R}.\; (f(x) + g(x))), a, (l + m)) \end{array}}}
\vnbprop{deriv-product}{\vnbstmt{\begin{array}{@{}l@{}} \forall f, g, a, l, m. \\ \quad (\operatorname{is\text{-}diff\text{-}at}(f, a, l) \wedge \operatorname{is\text{-}diff\text{-}at}(g, a, m)) \Rightarrow  \\ \quad \quad \operatorname{is\text{-}diff\text{-}at}((\lambda x \in \mathbb{R}.\; (f(x) \cdot g(x))), a, ((l \cdot g(a)) + (f(a) \cdot m))) \end{array}}}

\vnbbookfile{theorem-library/chain-rule}{theorem-library/chain-rule}
\noindent{\small chain-rule.scm -- THE CHAIN RULE (calculus.pdf Prop 2.8), PROVEN.}\par
\vnbprop{deriv-chain}{\vnbstmt{\begin{array}{@{}l@{}} \forall f, g, a, l, m. \\ \quad (\operatorname{is\text{-}diff\text{-}at}(f, a, l) \wedge \operatorname{is\text{-}diff\text{-}at}(g, f(a), m)) \Rightarrow  \\ \quad \quad \operatorname{is\text{-}diff\text{-}at}(\operatorname{compose}(g, f), a, (m \cdot l)) \end{array}}}
\vnbprop{deriv-of-is-diff-at}{\vnbstmt{\forall f, a, v.\; (\operatorname{is\text{-}diff\text{-}at}(f, a, v) \Rightarrow (\operatorname{deriv}(f, a) = v))}}
\vnbprop{deriv-chain-value}{\vnbstmt{\begin{array}{@{}l@{}} \forall f, g, a, l, m. \\ \quad (\operatorname{is\text{-}diff\text{-}at}(f, a, l) \wedge \operatorname{is\text{-}diff\text{-}at}(g, f(a), m)) \Rightarrow  \\ \quad \quad \operatorname{deriv}(\operatorname{compose}(g, f), a) = (\operatorname{deriv}(g, f(a)) \cdot \operatorname{deriv}(f, a)) \end{array}}}

\vnbbookfile{theorem-library/inverse-function}{theorem-library/inverse-function}
\noindent{\small inverse-function.scm -- THE INVERSE FUNCTION THEOREM ON THE LINE, Caratheodory form, PROVEN:}\par
\vnbprop{rr-recip-solve}{\vnbstmt{\begin{array}{@{}l@{}} \forall c, u, v. \\ \quad ((c \in \mathbb{R}) \wedge ((u \in \mathbb{R}) \wedge ((v \in \mathbb{R}) \wedge (\neg (c = 0) \wedge (u = (c \cdot v)))))) \Rightarrow  \\ \quad \quad v = (\operatorname{recip}(c) \cdot u) \end{array}}}
\vnbprop{deriv-inverse}{\vnbstmt{\begin{array}{@{}l@{}} \forall f, g, a, l. \\ \quad \operatorname{is\text{-}diff\text{-}at}(f, a, l) \wedge \\ \quad \neg (l = 0) \wedge \\ \quad g \in (\mathbb{R} \to \mathbb{R}) \wedge \\ \quad \forall x \in \mathbb{R}.\; (g(f(x)) = x) \wedge \\ \quad \forall y \in \mathbb{R}.\; (f(g(y)) = y) \wedge \\ \quad \operatorname{is\text{-}continuous\text{-}at}(\operatorname{rr\text{-}ms}, \operatorname{rr\text{-}ms}, g, f(a)) \Rightarrow \\ \quad \quad \operatorname{is\text{-}diff\text{-}at}(g, f(a), \operatorname{recip}(l)) \end{array}}}

\vnbbookfile{theorem-library/higher-derivatives}{theorem-library/higher-derivatives}
\noindent{\small higher-derivatives.scm -- calculus.pdf Def 2.2: the n-th derivative f\textasciicircum{}(n), by induction on n, built on DERIV / IS-DIFF-AT (differentiation.scm).}\par
\par\medskip\noindent\textbf{Definition} (\texttt{nth-deriv}).\;For every $f$,  $(\operatorname{nth\text{-}deriv}(f, 0) \simeq f)$\par
Let $f$ be arbitrary. Let $n \in \mathbb{N}$. Then
\begin{equation*}
(\operatorname{nth\text{-}deriv}(f, \operatorname{succ}(n)) \simeq (\lambda x \in \mathbb{R}.\; \operatorname{deriv}(\operatorname{nth\text{-}deriv}(f, n), x)))
\end{equation*}

\vnbprop{nth-deriv-one}{\vnbstmt{\forall f.\; (\operatorname{nth\text{-}deriv}(f, 1) \simeq (\lambda x \in \mathbb{R}.\; \operatorname{deriv}(f, x)))}}

\vnbbookfile{structure-library/little-o}{structure-library/little-o}
\noindent{\small MOVED 2026-09-20 (batch 12-A) from theorem-library/: this file is VOCABULARY -- definitions, notation and warranted supports, not one proof -- and every theorem stated with it had to load below it.  Its load.scm slot is unchanged. little-o.scm -- calculus.pdf Section 2.2, the calculus of o/O notation, in a}\par
\par\medskip\noindent\textbf{Definition} (\texttt{little-o-at}).\;For all $g$ and $a$, 
\begin{equation*}
\begin{array}{@{}l@{}} \operatorname{little\text{-}o\text{-}at}(g, a) \\ \Leftrightarrow\; g \in (\mathbb{R} \to \mathbb{R}) \wedge \\ a \in \mathbb{R} \wedge \\ \exists \mathit{eps} \in (\mathbb{R} \to \mathbb{R}). \\ \quad (\operatorname{is\text{-}continuous\text{-}at}(\operatorname{rr\text{-}ms}, \operatorname{rr\text{-}ms}, \mathit{eps}, a) \wedge ((\operatorname{eps}(a) = 0) \wedge \forall x \in \mathbb{R}.\; (g(x) = \cdot(\operatorname{eps}(x), -(x, a))))) \end{array}
\end{equation*}


\vnbbookfile{theorem-library/ccint-basics}{theorem-library/ccint-basics}
\noindent{\small ccint-basics.scm -- CCINT's membership law, PROVEN from its definition.}\par
\vnbprop{ccint-membership}{\vnbstmt{\forall a, b, x.\; ((x \in \operatorname{ccint}(a, b)) \Leftrightarrow ((x \in \mathbb{R}) \wedge ((a \leq x) \wedge (x \leq b))))}}
\vnbprop{ccint-elt-in-rr}{\vnbstmt{\forall a, b.; \forall x \in \operatorname{ccint}(a, b).\; (x \in \mathbb{R})}}

\vnbbookfile{theorem-library/ivt-proof}{theorem-library/ivt-proof}
\noindent{\small ivt-proof.scm -- the INTERMEDIATE VALUE THEOREM, PROVEN.}\par
\vnbprop{ivt}{\vnbstmt{\begin{array}{@{}l@{}} \forall f, a, b, c. \\ \quad f \in (\mathbb{R} \to \mathbb{R}) \wedge \\ \quad a \in \mathbb{R} \wedge \\ \quad b \in \mathbb{R} \wedge \\ \quad c \in \mathbb{R} \wedge \\ \quad a \leq b \wedge \\ \quad \forall x \in \operatorname{ccint}(a, b).\; \operatorname{is\text{-}continuous\text{-}at}(\operatorname{rr\text{-}ms}, \operatorname{rr\text{-}ms}, f, x) \wedge \\ \quad f(a) \leq c \wedge \\ \quad c \leq f(b) \Rightarrow \\ \quad \quad \exists p \in \mathbb{R}.\; ((p \in \operatorname{ccint}(a, b)) \wedge (f(p) = c)) \end{array}}}

\vnbbookfile{theorem-library/monotone-inverse}{theorem-library/monotone-inverse}
\noindent{\small monotone-inverse.scm -- MONOTONICITY VOCABULARY FOR REAL FUNCTIONS, and the EXISTENCE HALF of the inverse function theorem on a closed interval. Everything here is `modulo 0'.}\par
\par\medskip\noindent\textbf{Definition} (\texttt{is-strictly-increasing-on}).\;For all $f$ and $s$, 
\begin{equation*}
\begin{array}{@{}l@{}} \operatorname{is\text{-}strictly\text{-}increasing\text{-}on}(f, s) \\ \Leftrightarrow\; ((f \in (\mathbb{R} \to \mathbb{R})) \wedge ((s \subseteq \mathbb{R}) \wedge \forall u, v \in s.\; ((u < v) \Rightarrow (f(u) < f(v))))) \end{array}
\end{equation*}

\vnbprop{ccint-subset-rr}{\vnbstmt{\forall a, b.\; (\operatorname{ccint}(a, b) \subseteq \mathbb{R})}}
\vnbprop{strict-increase-injective}{\vnbstmt{\begin{array}{@{}l@{}} \forall f, s.; \forall u, v \in s. \\ \quad (\operatorname{is\text{-}strictly\text{-}increasing\text{-}on}(f, s) \wedge (f(u) = f(v))) \Rightarrow (u = v) \end{array}}}
\vnbprop{monotone-ccint-inverse-exists}{\vnbstmt{\begin{array}{@{}l@{}} \forall f, a, b.; \forall y \in \operatorname{ccint}(f(a), f(b)). \\ \quad f \in (\mathbb{R} \to \mathbb{R}) \wedge \\ \quad a \in \mathbb{R} \wedge \\ \quad b \in \mathbb{R} \wedge \\ \quad a \leq b \wedge \\ \quad \forall x \in \operatorname{ccint}(a, b).\; \operatorname{is\text{-}continuous\text{-}at}(\operatorname{rr\text{-}ms}, \operatorname{rr\text{-}ms}, f, x) \wedge \\ \quad \operatorname{is\text{-}strictly\text{-}increasing\text{-}on}(f, \operatorname{ccint}(a, b)) \Rightarrow \\ \quad \quad \exists p \in \operatorname{ccint}(a, b). \\ \quad \quad \quad (f(p) = y) \wedge \forall q \in \operatorname{ccint}(a, b).\; ((f(q) = y) \Rightarrow (p = q)) \end{array}}}
\vnbprop{monotone-ccint-inverse-order}{\vnbstmt{\begin{array}{@{}l@{}} \forall f, a, b.; \forall u, v \in \operatorname{ccint}(a, b). \\ \quad (\operatorname{is\text{-}strictly\text{-}increasing\text{-}on}(f, \operatorname{ccint}(a, b)) \wedge (f(u) < f(v))) \Rightarrow (u < v) \end{array}}}

\vnbbookfile{theorem-library/sq-continuous}{theorem-library/sq-continuous}
\noindent{\small sq-continuous.scm -- z |-$>$ z*z is a function RR -$>$ RR and is CONTINUOUS at every real point, PROVEN.}\par
\vnbprop{sq-fun-in-fun}{\vnbstmt{(\lambda z \in \mathbb{R}.\; (z \cdot z)) \in (\mathbb{R} \to \mathbb{R})}}
\vnbprop{sq-continuous-at}{\vnbstmt{\forall t \in \mathbb{R}.\; \operatorname{is\text{-}continuous\text{-}at}(\operatorname{rr\text{-}ms}, \operatorname{rr\text{-}ms}, (\lambda z \in \mathbb{R}.\; (z \cdot z)), t)}}

\vnbbookfile{theorem-library/sqrt-defined}{theorem-library/sqrt-defined}
\noindent{\small sqrt-defined.scm -- SQRT DEFINED by the intermediate value theorem, and the five supports that used to characterise it RETIRED as theorems.}\par
\vnbprop{sqrt-exists}{\vnbstmt{\forall a \in \mathbb{R}.\; ((0 \leq a) \Rightarrow \exists x \in \mathbb{R}.\; ((0 \leq x) \wedge ((x \cdot x) = a)))}}
\vnbprop{sqrt-unique}{\vnbstmt{\begin{array}{@{}l@{}} \forall x, y. \\ \quad (((x \in \mathbb{R}) \wedge (y \in \mathbb{R})) \wedge (((0 \leq x) \wedge (0 \leq y)) \wedge ((x \cdot x) = (y \cdot y)))) \Rightarrow (x = y) \end{array}}}
\vnbprop{sqrt-prop}{\vnbstmt{\begin{array}{@{}l@{}} \forall a \in \mathbb{R}. \\ \quad (0 \leq a) \Rightarrow ((\sqrt{a} \in \mathbb{R}) \wedge ((0 \leq \sqrt{a}) \wedge ((\sqrt{a} \cdot \sqrt{a}) = a))) \end{array}}}
\vnbprop{sqrt-char}{\vnbstmt{\begin{array}{@{}l@{}} \forall a, x. \\ \quad (((a \in \mathbb{R}) \wedge (0 \leq a)) \wedge (((x \in \mathbb{R}) \wedge (0 \leq x)) \wedge ((x \cdot x) = a))) \Rightarrow (\sqrt{a} = x) \end{array}}}
\vnbprop{sqrt-nonneg}{\vnbstmt{\forall a \in \mathbb{R}.\; ((0 \leq a) \Rightarrow ((\sqrt{a} \in \mathbb{R}) \wedge (0 \leq \sqrt{a})))}}
\vnbprop{sqrt-sq}{\vnbstmt{\forall a \in \mathbb{R}.\; ((0 \leq a) \Rightarrow ((\sqrt{a} \cdot \sqrt{a}) = a))}}
\vnbprop{sqrt-of-sq}{\vnbstmt{\forall a \in \mathbb{R}.\; (\sqrt{(a \cdot a)} = \lvert a \rvert)}}
\vnbprop{sqrt-mono}{\vnbstmt{\forall a, b \in \mathbb{R}.\; (((0 \leq a) \wedge (a \leq b)) \Rightarrow (\sqrt{a} \leq \sqrt{b}))}}
\vnbprop{sqrt-mul}{\vnbstmt{\begin{array}{@{}l@{}} \forall a, b \in \mathbb{R}. \\ \quad ((0 \leq a) \wedge (0 \leq b)) \Rightarrow (\sqrt{(a \cdot b)} = (\sqrt{a} \cdot \sqrt{b})) \end{array}}}

\vnbbookfile{theorem-library/cc-magnitude}{theorem-library/cc-magnitude}
\noindent{\small cc-magnitude.scm -- the complex modulus, PROVEN from its definition.}\par
\vnbprop{cc-magnitude-closed}{\vnbstmt{\forall z \in \mathbb{C}.\; (\operatorname{magnitude}(z) \in \mathbb{R})}}
\vnbprop{cc-magnitude-nonneg}{\vnbstmt{\forall z \in \mathbb{C}.\; (0 \leq \operatorname{magnitude}(z))}}
\vnbprop{sqrt-zero}{\vnbstmt{\sqrt{0} = 0}}
\vnbprop{rr-le-of-sq-le}{\vnbstmt{\forall t, s.\; (((t \in \mathbb{R}) \wedge ((s \in \mathbb{R}) \wedge ((0 \leq s) \wedge ((t \cdot t) \leq (s \cdot s))))) \Rightarrow (t \leq s))}}
\vnbprop{cc-magnitude-zero-iff}{\vnbstmt{\forall z \in \mathbb{C}.\; ((\operatorname{magnitude}(z) = 0) \Leftrightarrow (z = 0))}}
\vnbprop{cc-magnitude-neg}{\vnbstmt{\forall z \in \mathbb{C}.\; (\operatorname{magnitude}(-z) = \operatorname{magnitude}(z))}}
\vnbprop{rr-magnitude-is-abs}{\vnbstmt{\forall a \in \mathbb{R}.\; (\operatorname{magnitude}(a) = \lvert a \rvert)}}
\vnbprop{cc-magnitude-mul}{\vnbstmt{\begin{array}{@{}l@{}} \forall a, b. \\ \quad ((a \in \mathbb{C}) \wedge (b \in \mathbb{C})) \Rightarrow (\operatorname{magnitude}((a \cdot b)) = (\operatorname{magnitude}(a) \cdot \operatorname{magnitude}(b))) \end{array}}}
\vnbprop{cc-magnitude-triangle}{\vnbstmt{\begin{array}{@{}l@{}} \forall a, b. \\ \quad ((a \in \mathbb{C}) \wedge (b \in \mathbb{C})) \Rightarrow (\operatorname{magnitude}((a + b)) \leq (\operatorname{magnitude}(a) + \operatorname{magnitude}(b))) \end{array}}}

\vnbbookfile{theorem-library/cc-normed-field}{theorem-library/cc-normed-field}
\noindent{\small cc-normed-field.scm -- IS-NORMED-FIELD(CC-NORMED-FIELD), PROVED.}\par
\vnbprop{cc-is-normed-field}{\vnbstmt{\operatorname{is\text{-}normed\text{-}field}(\operatorname{cc\text{-}normed\text{-}field})}}

\vnbbookfile{theorem-library/cc-metric-space-proof}{theorem-library/cc-metric-space-proof}
\noindent{\small cc-metric-space-proof.scm -- IS-METRIC-SPACE(CC-MS), PROVEN.}\par
\vnbprop{cc-sub-in-cc}{\vnbstmt{\forall a, b.\; (((a \in \mathbb{C}) \wedge (b \in \mathbb{C})) \Rightarrow ((a - b) \in \mathbb{C}))}}
\vnbprop{cc-sub-zero-eq}{\vnbstmt{\forall a, b.\; (((a \in \mathbb{C}) \wedge ((b \in \mathbb{C}) \wedge ((a - b) = 0))) \Rightarrow (a = b))}}
\vnbprop{cc-is-metric-space}{\vnbstmt{\operatorname{is\text{-}metric\text{-}space}(\operatorname{cc\text{-}ms})}}

\vnbbookfile{theorem-library/complex-inner-product-laws}{theorem-library/complex-inner-product-laws}
\noindent{\small complex-inner-product-laws.scm -- what a sesquilinear form does in its SECOND argument, and the two expansion identities that follow.  All PROVEN modulo 0.}\par
\vnbprop{cc-conjugate-neg}{\vnbstmt{\forall a \in \mathbb{C}.\; (\operatorname{conjugate}(-a) = -\operatorname{conjugate}(a))}}
\vnbprop{cips-ip-add-right}{\vnbstmt{\begin{array}{@{}l@{}} \forall v.; \forall x, y, z \in \operatorname{vec}(v). \\ \quad \operatorname{is\text{-}complex\text{-}inner\text{-}product\text{-}space}(v) \Rightarrow  \\ \quad \quad \langle\cdot,\cdot\rangle_{v}(x, \operatorname{vadd}(v)(y, z)) = (\langle\cdot,\cdot\rangle_{v}(x, y) + \langle\cdot,\cdot\rangle_{v}(x, z)) \end{array}}}
\vnbprop{cips-ip-conj-homog-right}{\vnbstmt{\begin{array}{@{}l@{}} \forall v.; \forall a \in \mathbb{C}, x, y \in \operatorname{vec}(v). \\ \quad \operatorname{is\text{-}complex\text{-}inner\text{-}product\text{-}space}(v) \Rightarrow  \\ \quad \quad \langle\cdot,\cdot\rangle_{v}(x, \operatorname{act}(v)(a, y)) = (\operatorname{conjugate}(a) \cdot \langle\cdot,\cdot\rangle_{v}(x, y)) \end{array}}}
\vnbprop{cips-ip-zero-right}{\vnbstmt{\begin{array}{@{}l@{}} \forall v.; \forall x \in \operatorname{vec}(v). \\ \quad \operatorname{is\text{-}complex\text{-}inner\text{-}product\text{-}space}(v) \Rightarrow (\langle\cdot,\cdot\rangle_{v}(x, \operatorname{vzero}(v)) = 0) \end{array}}}
\vnbprop{cips-expand-sum}{\vnbstmt{\begin{array}{@{}l@{}} \forall v.; \forall x, y \in \operatorname{vec}(v). \\ \quad \operatorname{is\text{-}complex\text{-}inner\text{-}product\text{-}space}(v) \Rightarrow  \\ \quad \quad \langle\cdot,\cdot\rangle_{v}(\operatorname{vadd}(v)(x, y), \operatorname{vadd}(v)(x, y)) \\ \quad \quad =\; (\langle\cdot,\cdot\rangle_{v}(x, x) + \langle\cdot,\cdot\rangle_{v}(x, y)) + (\operatorname{conjugate}(\langle\cdot,\cdot\rangle_{v}(x, y)) + \langle\cdot,\cdot\rangle_{v}(y, y)) \end{array}}}
\vnbprop{cips-expand-two}{\vnbstmt{\begin{array}{@{}l@{}} \forall v, a, b, x, y. \\ \quad (\operatorname{is\text{-}complex\text{-}inner\text{-}product\text{-}space}(v) \wedge ((a \in \mathbb{C}) \wedge ((b \in \mathbb{C}) \wedge ((x \in \operatorname{vec}(v)) \wedge (y \in \operatorname{vec}(v)))))) \Rightarrow  \\ \quad \quad \langle\cdot,\cdot\rangle_{v}(\operatorname{vadd}(v)(\operatorname{act}(v)(a, x), \operatorname{act}(v)(b, y)), \operatorname{vadd}(v)(\operatorname{act}(v)(a, x), \operatorname{act}(v)(b, y))) \\ \quad \quad =\; (((a \cdot \operatorname{conjugate}(a)) \cdot \langle\cdot,\cdot\rangle_{v}(x, x)) + ((a \cdot \operatorname{conjugate}(b)) \cdot \langle\cdot,\cdot\rangle_{v}(x, y))) + (((b \cdot \operatorname{conjugate}(a)) \cdot \operatorname{conjugate}(\langle\cdot,\cdot\rangle_{v}(x, y))) + ((b \cdot \operatorname{conjugate}(b)) \cdot \langle\cdot,\cdot\rangle_{v}(y, y))) \end{array}}}

\vnbbookfile{theorem-library/inner-product-inequalities}{theorem-library/inner-product-inequalities}
\noindent{\small inner-product-inequalities.scm -- CAUCHY-SCHWARZ and MINKOWSKI at the INNER-PRODUCT level, over a COMPLEX-INNER-PRODUCT-SPACE.  Both PROVEN modulo 0.}\par
\vnbprop{cips-schwarz}{\vnbstmt{\begin{array}{@{}l@{}} \forall v, x, y. \\ \quad (\operatorname{is\text{-}complex\text{-}inner\text{-}product\text{-}space}(v) \wedge ((x \in \operatorname{vec}(v)) \wedge (y \in \operatorname{vec}(v)))) \Rightarrow  \\ \quad \quad (\langle\cdot,\cdot\rangle_{v}(x, y) \cdot \operatorname{conjugate}(\langle\cdot,\cdot\rangle_{v}(x, y))) \leq (\langle\cdot,\cdot\rangle_{v}(x, x) \cdot \langle\cdot,\cdot\rangle_{v}(y, y)) \end{array}}}
\vnbprop{cips-minkowski}{\vnbstmt{\begin{array}{@{}l@{}} \forall v, x, y. \\ \quad (\operatorname{is\text{-}complex\text{-}inner\text{-}product\text{-}space}(v) \wedge ((x \in \operatorname{vec}(v)) \wedge (y \in \operatorname{vec}(v)))) \Rightarrow  \\ \quad \quad \lVert \operatorname{vadd}(v)(x, y) \rVert \leq (\lVert x \rVert + \lVert y \rVert) \end{array}}}

\vnbbookfile{theorem-library/rake-ip-normed-ag}{theorem-library/rake-ip-normed-ag}
\noindent{\small rake-ip-normed-ag.scm -- rake batch 5b E: the induced norm of a complex inner product makes the additive group of its vectors a NORMED ABELIAN GROUP.  The support proved here is}\par
\vnbprop{ip-norm-unfold}{\vnbstmt{\forall v, x.\; (\lVert x \rVert \simeq \sqrt{\langle\cdot,\cdot\rangle_{v}(x, x)})}}
\vnbprop{cips-vneg-type}{\vnbstmt{\begin{array}{@{}l@{}} \forall v.; \forall x \in \operatorname{vec}(v). \\ \quad \operatorname{is\text{-}complex\text{-}inner\text{-}product\text{-}space}(v) \Rightarrow (\operatorname{vneg}(v)(x) \in \operatorname{vec}(v)) \end{array}}}
\vnbprop{cips-vadd-vneg}{\vnbstmt{\begin{array}{@{}l@{}} \forall v.; \forall x \in \operatorname{vec}(v). \\ \quad \operatorname{is\text{-}complex\text{-}inner\text{-}product\text{-}space}(v) \Rightarrow (\operatorname{vadd}(v)(x, \operatorname{vneg}(v)(x)) = \operatorname{vzero}(v)) \end{array}}}
\vnbprop{cips-ip-zero-left}{\vnbstmt{\begin{array}{@{}l@{}} \forall v.; \forall z \in \operatorname{vec}(v). \\ \quad \operatorname{is\text{-}complex\text{-}inner\text{-}product\text{-}space}(v) \Rightarrow (\langle\cdot,\cdot\rangle_{v}(\operatorname{vzero}(v), z) = 0) \end{array}}}
\vnbprop{cips-ip-neg-left}{\vnbstmt{\begin{array}{@{}l@{}} \forall v.; \forall x, z \in \operatorname{vec}(v). \\ \quad \operatorname{is\text{-}complex\text{-}inner\text{-}product\text{-}space}(v) \Rightarrow  \\ \quad \quad \langle\cdot,\cdot\rangle_{v}(\operatorname{vneg}(v)(x), z) = -\langle\cdot,\cdot\rangle_{v}(x, z) \end{array}}}
\vnbprop{cips-ip-neg-self}{\vnbstmt{\begin{array}{@{}l@{}} \forall v.; \forall x \in \operatorname{vec}(v). \\ \quad \operatorname{is\text{-}complex\text{-}inner\text{-}product\text{-}space}(v) \Rightarrow  \\ \quad \quad \langle\cdot,\cdot\rangle_{v}(\operatorname{vneg}(v)(x), \operatorname{vneg}(v)(x)) = \langle\cdot,\cdot\rangle_{v}(x, x) \end{array}}}
\vnbprop{ip-normed-ag-is-normed-ag}{\vnbstmt{\begin{array}{@{}l@{}} \forall v. \\ \quad \operatorname{is\text{-}complex\text{-}inner\text{-}product\text{-}space}(v) \Rightarrow \operatorname{is\text{-}normed\text{-}ag}(\operatorname{ip\text{-}normed\text{-}ag}(v)) \end{array}}}

\vnbbookfile{theorem-library/monotone-convergence-proof}{theorem-library/monotone-convergence-proof}
\noindent{\small monotone-convergence-proof.scm -- MONOTONE CONVERGENCE on RR, PROVEN.}\par
\vnbprop{nn-monotone-step-implies-le}{\vnbstmt{\begin{array}{@{}l@{}} \forall f \in (\mathbb{N} \to \mathbb{R}), n, m \in \mathbb{N}. \\ \quad (\forall k \in \mathbb{N}.\; (f(k) \leq f((k + 1))) \wedge (m \leq n)) \Rightarrow (f(m) \leq f(n)) \end{array}}}
\vnbprop{monotone-convergence-rr}{\vnbstmt{\begin{array}{@{}l@{}} \forall f \in (\mathbb{N} \to \mathbb{R}). \\ \quad (\forall k \in \mathbb{N}.\; (f(k) \leq f((k + 1))) \wedge \exists d \in \mathbb{R}.\; \forall k \in \mathbb{N}.\; (f(k) \leq d)) \Rightarrow  \\ \quad \quad \operatorname{converges}(\operatorname{rr\text{-}ms}, f) \end{array}}}

\vnbbookfile{theorem-library/comparison-test-proof}{theorem-library/comparison-test-proof}
\noindent{\small comparison-test-proof.scm -- THE COMPARISON TEST for real series, PROVEN, together with the two order lemmas it rests on and the readouts that make a real partial sum computable.}\par
\vnbprop{rr-additive-ag-opr}{\vnbstmt{(\operatorname{opr}(\operatorname{normed\text{-}field\text{-}additive\text{-}ag}(\operatorname{rr\text{-}normed\text{-}field})) \simeq (\lambda \langle \operatorname{x\_}, \operatorname{y\_} \rangle \in (\mathbb{R} \times \mathbb{R}).\; (\operatorname{x\_} + \operatorname{y\_})))}}
\vnbprop{rr-additive-ag-iden}{\vnbstmt{(\operatorname{iden}(\operatorname{normed\text{-}field\text{-}additive\text{-}ag}(\operatorname{rr\text{-}normed\text{-}field})) \simeq 0)}}
\vnbprop{series-partial-sum-zero}{\vnbstmt{\forall f.\; (\operatorname{series\text{-}partial\text{-}sum}(f, 0) \simeq 0)}}
\vnbprop{sum-ag-rr-in-rr}{\vnbstmt{\begin{array}{@{}l@{}} \forall k \in \mathbb{N}, f \in (\mathbb{N} \to \mathbb{R}). \\ \quad \operatorname{sum\text{-}ag}(\operatorname{normed\text{-}field\text{-}additive\text{-}ag}(\operatorname{rr\text{-}normed\text{-}field}), f, k) \in \mathbb{R} \end{array}}}
\vnbprop{series-partial-sum-in-rr}{\vnbstmt{\forall k \in \mathbb{N}, f \in (\mathbb{N} \to \mathbb{R}).\; (\operatorname{series\text{-}partial\text{-}sum}(f, k) \in \mathbb{R})}}
\vnbprop{series-partial-sum-unfold}{\vnbstmt{\begin{array}{@{}l@{}} \forall f, k. \\ \quad (\operatorname{series\text{-}partial\text{-}sum}(f, k) \simeq \operatorname{sum\text{-}ag}(\operatorname{normed\text{-}field\text{-}additive\text{-}ag}(\operatorname{rr\text{-}normed\text{-}field}), f, k)) \end{array}}}
\vnbprop{series-partial-sum-succ}{\vnbstmt{\begin{array}{@{}l@{}} \forall f.; \forall k \in \mathbb{N}. \\ \quad ((\operatorname{series\text{-}partial\text{-}sum}(f, k) \in \mathbb{R}) \wedge (f(k) \in \mathbb{R})) \Rightarrow  \\ \quad \quad (\operatorname{series\text{-}partial\text{-}sum}(f, (k + 1)) \simeq (\operatorname{series\text{-}partial\text{-}sum}(f, k) + f(k))) \end{array}}}
\vnbprop{series-partial-sum-monotone-nonneg}{\vnbstmt{\begin{array}{@{}l@{}} \forall f \in (\mathbb{N} \to \mathbb{R}), k \in \mathbb{N}. \\ \quad \forall n \in \mathbb{N}.\; (0 \leq f(n)) \Rightarrow  \\ \quad \quad \operatorname{series\text{-}partial\text{-}sum}(f, k) \leq \operatorname{series\text{-}partial\text{-}sum}(f, (k + 1)) \end{array}}}
\vnbprop{series-partial-sum-le-termwise}{\vnbstmt{\begin{array}{@{}l@{}} \forall f, g \in (\mathbb{N} \to \mathbb{R}), k \in \mathbb{N}. \\ \quad \forall n \in \mathbb{N}.\; (f(n) \leq g(n)) \Rightarrow  \\ \quad \quad \operatorname{series\text{-}partial\text{-}sum}(f, k) \leq \operatorname{series\text{-}partial\text{-}sum}(g, k) \end{array}}}
\vnbprop{monotone-convergent-bounded-above}{\vnbstmt{\begin{array}{@{}l@{}} \forall f \in (\mathbb{N} \to \mathbb{R}). \\ \quad (\forall k \in \mathbb{N}.\; (f(k) \leq f((k + 1))) \wedge \operatorname{converges}(\operatorname{rr\text{-}ms}, f)) \Rightarrow  \\ \quad \quad \exists d \in \mathbb{R}.\; \forall k \in \mathbb{N}.\; (f(k) \leq d) \end{array}}}
\vnbprop{series-partial-sum-seq-apply}{\vnbstmt{\begin{array}{@{}l@{}} \forall f.; \forall j \in \mathbb{N}. \\ \quad ((\lambda k \in \mathbb{N}.\; \operatorname{series\text{-}partial\text{-}sum}(f, k))(j) \simeq \operatorname{series\text{-}partial\text{-}sum}(f, j)) \end{array}}}
\vnbprop{series-partial-sum-seq-in-fun}{\vnbstmt{\begin{array}{@{}l@{}} \forall f \in (\mathbb{N} \to \mathbb{R}). \\ \quad (\lambda k \in \mathbb{N}.\; \operatorname{series\text{-}partial\text{-}sum}(f, k)) \in (\mathbb{N} \to \mathbb{R}) \end{array}}}
\vnbprop{abs-seq-in-fun}{\vnbstmt{\forall f \in (\mathbb{N} \to \mathbb{R}).\; ((\lambda n \in \mathbb{N}.\; \lvert f(n) \rvert) \in (\mathbb{N} \to \mathbb{R}))}}
\vnbprop{comparison-test}{\vnbstmt{\begin{array}{@{}l@{}} \forall f, g \in (\mathbb{N} \to \mathbb{R}). \\ \quad (\forall n \in \mathbb{N}.\; ((0 \leq f(n)) \wedge (f(n) \leq g(n))) \wedge \operatorname{series\text{-}converges}(g)) \Rightarrow  \\ \quad \quad \operatorname{series\text{-}converges}(f) \end{array}}}
\vnbprop{series-partial-sum-abs-le}{\vnbstmt{\begin{array}{@{}l@{}} \forall k \in \mathbb{N}, f \in (\mathbb{N} \to \mathbb{R}). \\ \quad \lvert \operatorname{series\text{-}partial\text{-}sum}(f, k) \rvert \leq \operatorname{series\text{-}partial\text{-}sum}((\lambda n \in \mathbb{N}.\; \lvert f(n) \rvert), k) \end{array}}}

\vnbbookfile{theorem-library/series-abs-triangle}{theorem-library/series-abs-triangle}
\noindent{\small series-abs-triangle.scm -- |Sum\_\{n$<$k\} f(n)| $<$= Sum\_\{n$<$k\} |f(n)|, PROVEN, and `abs(0) = 0', which the tree did not state.}\par
\vnbprop{rr-abs-zero-value}{\vnbstmt{\lvert 0 \rvert = 0}}
\vnbprop{series-abs-triangle}{\vnbstmt{\begin{array}{@{}l@{}} \forall k \in \mathbb{N}, f, g \in (\mathbb{N} \to \mathbb{R}). \\ \quad \forall j \in \mathbb{N}.\; (g(j) = \lvert f(j) \rvert) \Rightarrow  \\ \quad \quad \lvert \operatorname{series\text{-}partial\text{-}sum}(f, k) \rvert \leq \operatorname{series\text{-}partial\text{-}sum}(g, k) \end{array}}}

\vnbbookfile{theorem-library/series-block-abs}{theorem-library/series-block-abs}
\noindent{\small series-block-abs.scm -- THE BLOCK TRIANGLE INEQUALITY, in the form a Cauchy argument can use, and the NN gap lemma under it.}\par
\vnbprop{series-abs-triangle-block}{\vnbstmt{\begin{array}{@{}l@{}} \forall d, m \in \mathbb{N}, f, g \in (\mathbb{N} \to \mathbb{R}). \\ \quad \forall j \in \mathbb{N}.\; (g(j) = \lvert f(j) \rvert) \Rightarrow  \\ \quad \quad \lvert (\operatorname{series\text{-}partial\text{-}sum}(f, (m + d)) - \operatorname{series\text{-}partial\text{-}sum}(f, m)) \rvert \\ \quad \quad \leq\; (\operatorname{series\text{-}partial\text{-}sum}(g, (m + d)) - \operatorname{series\text{-}partial\text{-}sum}(g, m)) \end{array}}}
\vnbprop{series-abs-triangle-le}{\vnbstmt{\begin{array}{@{}l@{}} \forall k, m \in \mathbb{N}, f, g \in (\mathbb{N} \to \mathbb{R}). \\ \quad (\forall j \in \mathbb{N}.\; (g(j) = \lvert f(j) \rvert) \wedge (m \leq k)) \Rightarrow  \\ \quad \quad \lvert (\operatorname{series\text{-}partial\text{-}sum}(f, k) - \operatorname{series\text{-}partial\text{-}sum}(f, m)) \rvert \\ \quad \quad \leq\; (\operatorname{series\text{-}partial\text{-}sum}(g, k) - \operatorname{series\text{-}partial\text{-}sum}(g, m)) \end{array}}}

\vnbbookfile{theorem-library/series-linearity}{theorem-library/series-linearity}
\noindent{\small series-linearity.scm -- LINEARITY OF THE REAL PARTIAL SUM, in POINTWISE form.}\par
\vnbprop{series-partial-sum-in-rr-ptwise}{\vnbstmt{\begin{array}{@{}l@{}} \forall n \in \mathbb{N}, f. \\ \quad \forall k \in \mathbb{N}.\; (f(k) \in \mathbb{R}) \Rightarrow (\operatorname{series\text{-}partial\text{-}sum}(f, n) \in \mathbb{R}) \end{array}}}
\vnbprop{series-partial-sum-add-ptwise}{\vnbstmt{\begin{array}{@{}l@{}} \forall n \in \mathbb{N}, f, g, h. \\ \quad \forall k \in \mathbb{N}.\; (f(k) \in \mathbb{R}) \wedge \\ \quad \forall k \in \mathbb{N}.\; (g(k) \in \mathbb{R}) \wedge \\ \quad \forall k \in \mathbb{N}.\; (h(k) = (f(k) + g(k))) \Rightarrow \\ \quad \quad \operatorname{series\text{-}partial\text{-}sum}(h, n) = (\operatorname{series\text{-}partial\text{-}sum}(f, n) + \operatorname{series\text{-}partial\text{-}sum}(g, n)) \end{array}}}
\vnbprop{series-partial-sum-scale-ptwise}{\vnbstmt{\begin{array}{@{}l@{}} \forall n \in \mathbb{N}, c, f, h. \\ \quad ((c \in \mathbb{R}) \wedge (\forall k \in \mathbb{N}.\; (f(k) \in \mathbb{R}) \wedge \forall k \in \mathbb{N}.\; (h(k) = (c \cdot f(k))))) \Rightarrow  \\ \quad \quad \operatorname{series\text{-}partial\text{-}sum}(h, n) = (c \cdot \operatorname{series\text{-}partial\text{-}sum}(f, n)) \end{array}}}
\vnbprop{series-partial-sum-weighted-expansion}{\vnbstmt{\begin{array}{@{}l@{}} \forall n \in \mathbb{N}, x, y, g, w, p, q, s. \\ \quad x \in \mathbb{R} \wedge \\ \quad y \in \mathbb{R} \wedge \\ \quad \forall k \in \mathbb{Z}.\; (g(k) \in \mathbb{R}) \wedge \\ \quad \forall k \in \mathbb{Z}.\; (w(k) \in \mathbb{R}) \wedge \\ \quad \forall k \in \mathbb{N}.\; (q(k) \in \mathbb{R}) \wedge \\ \quad \forall k \in \mathbb{N}.\; (s(k) \in \mathbb{R}) \wedge \\ \quad \forall k \in \mathbb{N}.\; (p(k) = (w(k) \cdot ((x \cdot g((k - 1))) + (y \cdot g(k))))) \wedge \\ \quad \forall k \in \mathbb{N}.\; (q(k) = (w((k + 1)) \cdot g(k))) \wedge \\ \quad \forall k \in \mathbb{N}.\; (s(k) = (w(k) \cdot g(k))) \Rightarrow \\ \quad \quad \operatorname{series\text{-}partial\text{-}sum}(p, (n + 1)) \\ \quad \quad =\; (x \cdot ((w(0) \cdot g(-1)) + \operatorname{series\text{-}partial\text{-}sum}(q, n))) + (y \cdot (\operatorname{series\text{-}partial\text{-}sum}(s, n) + (w(n) \cdot g(n)))) \end{array}}}
\vnbprop{series-partial-sum-le-termwise-ptwise}{\vnbstmt{\begin{array}{@{}l@{}} \forall n \in \mathbb{N}, f, g. \\ \quad (\forall k \in \mathbb{N}.\; (f(k) \in \mathbb{R}) \wedge (\forall k \in \mathbb{N}.\; (g(k) \in \mathbb{R}) \wedge \forall k \in \mathbb{N}.\; (f(k) \leq g(k)))) \Rightarrow  \\ \quad \quad \operatorname{series\text{-}partial\text{-}sum}(f, n) \leq \operatorname{series\text{-}partial\text{-}sum}(g, n) \end{array}}}

\vnbbookfile{theorem-library/bernstein-basis}{theorem-library/bernstein-basis}
\noindent{\small bernstein-basis.scm -- the Bernstein basis polynomials, and the first of the two identities Proposition 5.4 of docs/calculus.pdf asks for.}\par
\par\medskip\noindent\textbf{Definition} (\texttt{bernstein-basis}).\;$\operatorname{bernstein\text{-}basis}(n, x) :=$
\begin{equation*}
\operatorname{comb\text{-}kk}(\operatorname{normed\text{-}field\text{-}as\text{-}commutative\text{-}ring}(\operatorname{rr\text{-}normed\text{-}field}), x, (1 - x), n)
\end{equation*}

\vnbprop{ring-power-of-one}{\vnbstmt{\begin{array}{@{}l@{}} \forall n \in \mathbb{N}, r. \\ \quad \operatorname{is\text{-}commutative\text{-}ring}(r) \Rightarrow (\operatorname{ring\text{-}power}(r, \operatorname{one}(r), n) = \operatorname{one}(r)) \end{array}}}
\vnbprop{series-partial-sum-is-ring-sum}{\vnbstmt{\begin{array}{@{}l@{}} \forall n \in \mathbb{N}, f. \\ \quad (\operatorname{series\text{-}partial\text{-}sum}(f, n) \simeq \operatorname{sum}(\operatorname{normed\text{-}field\text{-}as\text{-}commutative\text{-}ring}(\operatorname{rr\text{-}normed\text{-}field}), f, n)) \end{array}}}
\vnbprop{binomial-sum-value}{\vnbstmt{\begin{array}{@{}l@{}} \forall n, r, x, y. \\ \quad ((n \in \mathbb{N}) \wedge (\operatorname{is\text{-}commutative\text{-}ring}(r) \wedge ((x \in \operatorname{carr}(r)) \wedge (y \in \operatorname{carr}(r))))) \Rightarrow  \\ \quad \quad \operatorname{sum}(r, \operatorname{comb\text{-}kk}(r, x, y, n), (n + 1)) = \operatorname{ring\text{-}power}(r, \operatorname{add}(r)(x, y), n) \end{array}}}
\vnbprop{bernstein-basis-unfold}{\vnbstmt{\begin{array}{@{}l@{}} \forall n, x. \\ \quad (\operatorname{bernstein\text{-}basis}(n, x) \simeq \operatorname{comb\text{-}kk}(\operatorname{normed\text{-}field\text{-}as\text{-}commutative\text{-}ring}(\operatorname{rr\text{-}normed\text{-}field}), x, (1 - x), n)) \end{array}}}
\vnbprop{bernstein-basis-in-fun}{\vnbstmt{\forall n \in \mathbb{N}, x \in \mathbb{R}.\; (\operatorname{bernstein\text{-}basis}(n, x) \in (\mathbb{Z} \to \mathbb{R}))}}
\vnbprop{bernstein-partition}{\vnbstmt{\begin{array}{@{}l@{}} \forall n \in \mathbb{N}, x \in \mathbb{R}. \\ \quad \operatorname{series\text{-}partial\text{-}sum}(\operatorname{bernstein\text{-}basis}(n, x), (n + 1)) = 1 \end{array}}}

\vnbbookfile{theorem-library/bernstein-moments}{theorem-library/bernstein-moments}
\noindent{\small bernstein-moments.scm -- the moments of the Bernstein basis, and identity (76) of docs/calculus.pdf Proposition 5.4.}\par
\vnbprop{bernstein-basis-ptwise-in-rr}{\vnbstmt{\begin{array}{@{}l@{}} \forall n, x, k. \\ \quad ((n \in \mathbb{N}) \wedge ((x \in \mathbb{R}) \wedge (k \in \mathbb{Z}))) \Rightarrow (\operatorname{bernstein\text{-}basis}(n, x)(k) \in \mathbb{R}) \end{array}}}
\vnbprop{bernstein-basis-succ}{\vnbstmt{\begin{array}{@{}l@{}} \forall n, x, k. \\ \quad ((n \in \mathbb{N}) \wedge ((x \in \mathbb{R}) \wedge (k \in \mathbb{N}))) \Rightarrow  \\ \quad \quad \operatorname{bernstein\text{-}basis}((n + 1), x)(k) \\ \quad \quad =\; (x \cdot \operatorname{bernstein\text{-}basis}(n, x)((k - 1))) + ((1 - x) \cdot \operatorname{bernstein\text{-}basis}(n, x)(k)) \end{array}}}
\vnbprop{bernstein-basis-null}{\vnbstmt{\begin{array}{@{}l@{}} \forall n, x, k. \\ \quad ((n \in \mathbb{N}) \wedge ((x \in \mathbb{R}) \wedge ((k \in \mathbb{Z}) \wedge (k < 0)))) \Rightarrow (\operatorname{bernstein\text{-}basis}(n, x)(k) = 0) \end{array}}}
\vnbprop{bernstein-basis-above}{\vnbstmt{\begin{array}{@{}l@{}} \forall n, x, k. \\ \quad ((n \in \mathbb{N}) \wedge ((x \in \mathbb{R}) \wedge ((k \in \mathbb{Z}) \wedge (n < k)))) \Rightarrow (\operatorname{bernstein\text{-}basis}(n, x)(k) = 0) \end{array}}}
\vnbprop{bernstein-weighted-step}{\vnbstmt{\begin{array}{@{}l@{}} \forall n, x, w, p, q, s. \\ \quad n \in \mathbb{N} \wedge \\ \quad x \in \mathbb{R} \wedge \\ \quad \forall k \in \mathbb{Z}.\; (w(k) \in \mathbb{R}) \wedge \\ \quad \forall k \in \mathbb{N}.\; (q(k) \in \mathbb{R}) \wedge \\ \quad \forall k \in \mathbb{N}.\; (s(k) \in \mathbb{R}) \wedge \\ \quad \forall k \in \mathbb{N}.\; (p(k) = (w(k) \cdot \operatorname{bernstein\text{-}basis}((n + 1), x)(k))) \wedge \\ \quad \forall k \in \mathbb{N}.\; (q(k) = (w((k + 1)) \cdot \operatorname{bernstein\text{-}basis}(n, x)(k))) \wedge \\ \quad \forall k \in \mathbb{N}.\; (s(k) = (w(k) \cdot \operatorname{bernstein\text{-}basis}(n, x)(k))) \Rightarrow \\ \quad \quad \operatorname{series\text{-}partial\text{-}sum}(p, ((n + 1) + 1)) \\ \quad \quad =\; (x \cdot \operatorname{series\text{-}partial\text{-}sum}(q, (n + 1))) + ((1 - x) \cdot \operatorname{series\text{-}partial\text{-}sum}(s, (n + 1))) \end{array}}}
\vnbprop{bernstein-moment-1}{\vnbstmt{\begin{array}{@{}l@{}} \forall n \in \mathbb{N}, x, h. \\ \quad ((x \in \mathbb{R}) \wedge \forall k \in \mathbb{N}.\; (h(k) = (k \cdot \operatorname{bernstein\text{-}basis}(n, x)(k)))) \Rightarrow  \\ \quad \quad \operatorname{series\text{-}partial\text{-}sum}(h, (n + 1)) = (n \cdot x) \end{array}}}
\vnbprop{bernstein-moment-2}{\vnbstmt{\begin{array}{@{}l@{}} \forall n \in \mathbb{N}, x, h. \\ \quad ((x \in \mathbb{R}) \wedge \forall k \in \mathbb{N}.\; (h(k) = ((k \cdot (k - 1)) \cdot \operatorname{bernstein\text{-}basis}(n, x)(k)))) \Rightarrow  \\ \quad \quad \operatorname{series\text{-}partial\text{-}sum}(h, (n + 1)) = ((n \cdot (n - 1)) \cdot (x \cdot x)) \end{array}}}
\vnbprop{bernstein-variance}{\vnbstmt{\begin{array}{@{}l@{}} \forall n, x, h. \\ \quad n \in \mathbb{N} \wedge \\ \quad x \in \mathbb{R} \wedge \\ \quad \forall k \in \mathbb{N}. \\ \quad \quad h(k) = (((k - (n \cdot x)) \cdot (k - (n \cdot x))) \cdot \operatorname{bernstein\text{-}basis}(n, x)(k)) \Rightarrow \\ \quad \quad \operatorname{series\text{-}partial\text{-}sum}(h, (n + 1)) = ((n \cdot x) \cdot (1 - x)) \end{array}}}
\vnbprop{bernstein-variance-bound}{\vnbstmt{\begin{array}{@{}l@{}} \forall n, x, h. \\ \quad n \in \mathbb{N} \wedge \\ \quad \neg (n = 0) \wedge \\ \quad x \in \mathbb{R} \wedge \\ \quad \forall k \in \mathbb{N}. \\ \quad \quad h(k) = (((k - (n \cdot x)) \cdot (k - (n \cdot x))) \cdot \operatorname{bernstein\text{-}basis}(n, x)(k)) \Rightarrow \\ \quad \quad (/(1, n) \cdot \operatorname{series\text{-}partial\text{-}sum}(h, (n + 1))) \leq 1 \end{array}}}

\vnbbookfile{theorem-library/series-cauchy-proof}{theorem-library/series-cauchy-proof}
\noindent{\small series-cauchy-proof.scm -- THE CAUCHY CRITERION FOR REAL SERIES, PROVEN.}\par
\vnbprop{series-cauchy-criterion}{\vnbstmt{\begin{array}{@{}l@{}} \forall f \in (\mathbb{N} \to \mathbb{R}), s. \\ \quad (\operatorname{series\text{-}converges}(f) \wedge \operatorname{pos\text{-}rr}(s)) \Rightarrow  \\ \quad \quad \exists d \in \mathbb{N}. \forall m, n \in \mathbb{N}. \\ \quad \quad \quad ((d \leq m) \wedge (m \leq n)) \Rightarrow  \\ \quad \quad \quad \quad \lvert (\operatorname{series\text{-}partial\text{-}sum}(f, n) - \operatorname{series\text{-}partial\text{-}sum}(f, m)) \rvert \leq s \end{array}}}

\vnbbookfile{theorem-library/dominated-convergence}{theorem-library/dominated-convergence}
\noindent{\small dominated-convergence.scm -- DOMINATED CONVERGENCE, PROVEN, in two forms: for real series (the content) and in ELL-ONE (the statement the user asked for), together with the twenty limit / partial-sum / complex-magnitude lemmas the two are assembled from.}\par
\par\medskip\noindent\textbf{Definition} (\texttt{ell-one}).\;
\begin{equation*}
\operatorname{ell\text{-}one} = \{\operatorname{x\_} \in \mathbb{C}^{\mathbb{N}} : \operatorname{series\text{-}converges}((\lambda \operatorname{k\_} \in \mathbb{N}.\; \operatorname{magnitude}(\operatorname{x\_}(\operatorname{k\_}))))\}
\end{equation*}

\par\medskip\noindent\textbf{Definition} (\texttt{series-limit}).\;$\operatorname{series\text{-}limit}(f) :=$ $(\iota\, \operatorname{lm\_}.\;\operatorname{series\text{-}converges\text{-}to}(f, \operatorname{lm\_}))$
\vnbprop{rr-null-sum}{\vnbstmt{\begin{array}{@{}l@{}} \forall f, g, h \in (\mathbb{N} \to \mathbb{R}). \\ \quad (\operatorname{converges\text{-}to}(\operatorname{rr\text{-}ms}, f, 0) \wedge (\operatorname{converges\text{-}to}(\operatorname{rr\text{-}ms}, g, 0) \wedge \forall j \in \mathbb{N}.\; (h(j) = (f(j) + g(j))))) \Rightarrow  \\ \quad \quad \operatorname{converges\text{-}to}(\operatorname{rr\text{-}ms}, h, 0) \end{array}}}
\vnbprop{rr-limit-abs-le}{\vnbstmt{\begin{array}{@{}l@{}} \forall f \in (\mathbb{N} \to \mathbb{R}), v, d \in \mathbb{R}. \\ \quad (\operatorname{converges\text{-}to}(\operatorname{rr\text{-}ms}, f, v) \wedge \forall n \in \mathbb{N}.\; (\lvert f(n) \rvert \leq d)) \Rightarrow  \\ \quad \quad \lvert v \rvert \leq d \end{array}}}
\vnbprop{series-block-le}{\vnbstmt{\begin{array}{@{}l@{}} \forall k, n \in \mathbb{N}, f, g \in (\mathbb{N} \to \mathbb{R}). \\ \quad (\forall i \in \mathbb{N}.\; ((0 \leq f(i)) \wedge (f(i) \leq g(i))) \wedge (n \leq k)) \Rightarrow  \\ \quad \quad (\operatorname{series\text{-}partial\text{-}sum}(f, k) - \operatorname{series\text{-}partial\text{-}sum}(f, n)) \\ \quad \quad \leq\; (\operatorname{series\text{-}partial\text{-}sum}(g, k) - \operatorname{series\text{-}partial\text{-}sum}(g, n)) \end{array}}}
\vnbprop{series-head-null}{\vnbstmt{\begin{array}{@{}l@{}} \forall n \in \mathbb{N}, u.; \forall d \in (\mathbb{N} \to \mathbb{R}). \\ \quad \forall j \in \mathbb{N}.\; (u(j) \in (\mathbb{N} \to \mathbb{R})) \wedge \\ \quad \forall i \in \mathbb{N}.\; \operatorname{converges\text{-}to}(\operatorname{rr\text{-}ms}, (\lambda j \in \mathbb{N}.\; u(j)(i)), 0) \wedge \\ \quad \forall j \in \mathbb{N}.\; (d(j) = \operatorname{series\text{-}partial\text{-}sum}(u(j), n)) \Rightarrow \\ \quad \quad \operatorname{converges\text{-}to}(\operatorname{rr\text{-}ms}, d, 0) \end{array}}}
\vnbprop{series-partial-sum-nonneg}{\vnbstmt{\begin{array}{@{}l@{}} \forall k \in \mathbb{N}, f \in (\mathbb{N} \to \mathbb{R}). \\ \quad \forall i \in \mathbb{N}.\; (0 \leq f(i)) \Rightarrow (0 \leq \operatorname{series\text{-}partial\text{-}sum}(f, k)) \end{array}}}
\vnbprop{series-partial-sum-mono}{\vnbstmt{\begin{array}{@{}l@{}} \forall k, m \in \mathbb{N}, f \in (\mathbb{N} \to \mathbb{R}). \\ \quad (\forall i \in \mathbb{N}.\; (0 \leq f(i)) \wedge (m \leq k)) \Rightarrow  \\ \quad \quad \operatorname{series\text{-}partial\text{-}sum}(f, m) \leq \operatorname{series\text{-}partial\text{-}sum}(f, k) \end{array}}}
\vnbprop{series-tail-small}{\vnbstmt{\begin{array}{@{}l@{}} \forall g \in (\mathbb{N} \to \mathbb{R}), s. \\ \quad (\operatorname{series\text{-}converges}(g) \wedge \operatorname{pos\text{-}rr}(s)) \Rightarrow  \\ \quad \quad \exists d \in \mathbb{N}. \forall k \in \mathbb{N}. \\ \quad \quad \quad (d \leq k) \Rightarrow ((\operatorname{series\text{-}partial\text{-}sum}(g, k) - \operatorname{series\text{-}partial\text{-}sum}(g, d)) \leq s) \end{array}}}
\vnbprop{dominated-null-series}{\vnbstmt{\begin{array}{@{}l@{}} \forall g, e \in (\mathbb{N} \to \mathbb{R}), u. \\ \quad \forall j \in \mathbb{N}.\; (u(j) \in (\mathbb{N} \to \mathbb{R})) \wedge \\ \quad \forall j, i \in \mathbb{N}.\; ((0 \leq u(j)(i)) \wedge (u(j)(i) \leq g(i))) \wedge \\ \quad \operatorname{series\text{-}converges}(g) \wedge \\ \quad \forall i \in \mathbb{N}.\; \operatorname{converges\text{-}to}(\operatorname{rr\text{-}ms}, (\lambda j \in \mathbb{N}.\; u(j)(i)), 0) \wedge \\ \quad \forall j \in \mathbb{N}.\; \operatorname{series\text{-}converges\text{-}to}(u(j), e(j)) \Rightarrow \\ \quad \quad \operatorname{converges\text{-}to}(\operatorname{rr\text{-}ms}, e, 0) \end{array}}}
\vnbprop{series-limit-converges-to}{\vnbstmt{\begin{array}{@{}l@{}} \forall f \in (\mathbb{N} \to \mathbb{R}). \\ \quad \operatorname{series\text{-}converges}(f) \Rightarrow \operatorname{series\text{-}converges\text{-}to}(f, \operatorname{series\text{-}limit}(f)) \end{array}}}
\vnbprop{series-limit-in-rr}{\vnbstmt{\forall f \in (\mathbb{N} \to \mathbb{R}).\; (\operatorname{series\text{-}converges}(f) \Rightarrow (\operatorname{series\text{-}limit}(f) \in \mathbb{R}))}}
\vnbprop{series-partial-sum-add}{\vnbstmt{\begin{array}{@{}l@{}} \forall k \in \mathbb{N}, f, g, h \in (\mathbb{N} \to \mathbb{R}). \\ \quad \forall i \in \mathbb{N}.\; (h(i) = (f(i) + g(i))) \Rightarrow  \\ \quad \quad \operatorname{series\text{-}partial\text{-}sum}(h, k) = (\operatorname{series\text{-}partial\text{-}sum}(f, k) + \operatorname{series\text{-}partial\text{-}sum}(g, k)) \end{array}}}
\vnbprop{series-partial-sum-bounded}{\vnbstmt{\begin{array}{@{}l@{}} \forall f \in (\mathbb{N} \to \mathbb{R}). \\ \quad (\forall i \in \mathbb{N}.\; (0 \leq f(i)) \wedge \operatorname{series\text{-}converges}(f)) \Rightarrow  \\ \quad \quad \exists d \in \mathbb{R}.\; \forall x \in \mathbb{N}.\; (\operatorname{series\text{-}partial\text{-}sum}(f, x) \leq d) \end{array}}}
\vnbprop{series-converges-sum}{\vnbstmt{\begin{array}{@{}l@{}} \forall f, g, h \in (\mathbb{N} \to \mathbb{R}). \\ \quad \forall i \in \mathbb{N}.\; (0 \leq f(i)) \wedge \\ \quad \forall i \in \mathbb{N}.\; (0 \leq g(i)) \wedge \\ \quad \forall i \in \mathbb{N}.\; (h(i) = (f(i) + g(i))) \wedge \\ \quad \operatorname{series\text{-}converges}(f) \wedge \\ \quad \operatorname{series\text{-}converges}(g) \Rightarrow \\ \quad \quad \operatorname{series\text{-}converges}(h) \end{array}}}
\vnbprop{cc-ms-dist}{\vnbstmt{\begin{array}{@{}l@{}} \forall u, v. \\ \quad ((u \in \mathbb{C}) \wedge (v \in \mathbb{C})) \Rightarrow (\operatorname{dist}(\operatorname{cc\text{-}ms})(u, v) \simeq \operatorname{magnitude}((u - v))) \end{array}}}
\vnbprop{cc-magnitude-sub-triangle}{\vnbstmt{\forall p, q \in \mathbb{C}.\; (\operatorname{magnitude}((p - q)) \leq (\operatorname{magnitude}(p) + \operatorname{magnitude}(q)))}}
\vnbprop{cc-magnitude-sub-sym}{\vnbstmt{\forall p, q \in \mathbb{C}.\; (\operatorname{magnitude}((p - q)) = \operatorname{magnitude}((q - p)))}}
\vnbprop{cc-magnitude-le-add}{\vnbstmt{\forall p, q \in \mathbb{C}.\; (\operatorname{magnitude}(p) \leq (\operatorname{magnitude}((p - q)) + \operatorname{magnitude}(q)))}}
\vnbprop{cc-limit-magnitude-le}{\vnbstmt{\begin{array}{@{}l@{}} \forall s \in (\mathbb{N} \to \mathbb{C}), v \in \mathbb{C}, d \in \mathbb{R}. \\ \quad (\operatorname{converges\text{-}to}(\operatorname{cc\text{-}ms}, s, v) \wedge \forall n \in \mathbb{N}.\; (\operatorname{magnitude}(s(n)) \leq d)) \Rightarrow  \\ \quad \quad \operatorname{magnitude}(v) \leq d \end{array}}}
\vnbprop{cc-converges-magnitude-null}{\vnbstmt{\begin{array}{@{}l@{}} \forall s \in (\mathbb{N} \to \mathbb{C}), v \in \mathbb{C}, m \in (\mathbb{N} \to \mathbb{R}). \\ \quad (\operatorname{converges\text{-}to}(\operatorname{cc\text{-}ms}, s, v) \wedge \forall j \in \mathbb{N}.\; (m(j) = \operatorname{magnitude}((s(j) - v)))) \Rightarrow  \\ \quad \quad \operatorname{converges\text{-}to}(\operatorname{rr\text{-}ms}, m, 0) \end{array}}}
\vnbprop{ell-one-membership}{\vnbstmt{\begin{array}{@{}l@{}} \forall y. \\ \quad (y \in \operatorname{ell\text{-}one}) \Leftrightarrow ((y \in \mathbb{C}^{\mathbb{N}}) \wedge \operatorname{series\text{-}converges}((\lambda k \in \mathbb{N}.\; \operatorname{magnitude}(y(k))))) \end{array}}}
\vnbprop{ell-one-dominated-convergence}{\vnbstmt{\begin{array}{@{}l@{}} \forall a \in (\mathbb{N} \to \mathbb{R}), x \in \mathbb{C}^{\mathbb{N}}, e \in (\mathbb{N} \to \mathbb{R}), s. \\ \quad \forall j \in \mathbb{N}.\; (s(j) \in \mathbb{C}^{\mathbb{N}}) \wedge \\ \quad \operatorname{series\text{-}converges}(a) \wedge \\ \quad \forall j, k \in \mathbb{N}.\; (\operatorname{magnitude}(s(j)(k)) \leq a(k)) \wedge \\ \quad \forall k \in \mathbb{N}.\; \operatorname{converges\text{-}to}(\operatorname{cc\text{-}ms}, (\lambda j \in \mathbb{N}.\; s(j)(k)), x(k)) \wedge \\ \quad \forall j \in \mathbb{N}. \\ \quad \quad e(j) = \operatorname{series\text{-}limit}((\lambda k \in \mathbb{N}.\; \operatorname{magnitude}((s(j)(k) - x(k))))) \Rightarrow \\ \quad \quad (x \in \operatorname{ell\text{-}one}) \wedge \operatorname{converges\text{-}to}(\operatorname{rr\text{-}ms}, e, 0) \end{array}}}

\vnbbookfile{theorem-library/cc-complete-proof}{theorem-library/cc-complete-proof}
\noindent{\small cc-complete-proof.scm -- CC IS A COMPLETE METRIC SPACE.}\par
\vnbprop{cc-re-sub}{\vnbstmt{\begin{array}{@{}l@{}} \forall z, w. \\ \quad ((z \in \mathbb{C}) \wedge (w \in \mathbb{C})) \Rightarrow (\operatorname{real\text{-}part}((z - w)) = (\operatorname{real\text{-}part}(z) - \operatorname{real\text{-}part}(w))) \end{array}}}
\vnbprop{cc-im-sub}{\vnbstmt{\begin{array}{@{}l@{}} \forall z, w. \\ \quad ((z \in \mathbb{C}) \wedge (w \in \mathbb{C})) \Rightarrow (\operatorname{imag\text{-}part}((z - w)) = (\operatorname{imag\text{-}part}(z) - \operatorname{imag\text{-}part}(w))) \end{array}}}
\vnbprop{cc-abs-re-le-magnitude}{\vnbstmt{\forall z \in \mathbb{C}.\; (\lvert \operatorname{real\text{-}part}(z) \rvert \leq \operatorname{magnitude}(z))}}
\vnbprop{cc-abs-im-le-magnitude}{\vnbstmt{\forall z \in \mathbb{C}.\; (\lvert \operatorname{imag\text{-}part}(z) \rvert \leq \operatorname{magnitude}(z))}}
\vnbprop{cc-magnitude-le-re-im}{\vnbstmt{\forall z \in \mathbb{C}.\; (\operatorname{magnitude}(z) \leq (\lvert \operatorname{real\text{-}part}(z) \rvert + \lvert \operatorname{imag\text{-}part}(z) \rvert))}}
\vnbprop{cc-cauchy-of-dominated}{\vnbstmt{\begin{array}{@{}l@{}} \forall f \in (\mathbb{N} \to \mathbb{C}), p \in (\mathbb{N} \to \mathbb{R}). \\ \quad (\operatorname{is\text{-}cauchy\text{-}seq}(\operatorname{cc\text{-}ms}, f) \wedge \forall m, n \in \mathbb{N}.\; (\lvert (p(m) - p(n)) \rvert \leq \operatorname{magnitude}((f(m) - f(n))))) \Rightarrow  \\ \quad \quad \operatorname{is\text{-}cauchy\text{-}seq}(\operatorname{rr\text{-}ms}, p) \end{array}}}
\vnbprop{cc-converges-of-coords}{\vnbstmt{\begin{array}{@{}l@{}} \forall f \in (\mathbb{N} \to \mathbb{C}), g, q \in (\mathbb{N} \to \mathbb{R}), v \in \mathbb{C}. \\ \quad \forall n \in \mathbb{N}.\; (\operatorname{real\text{-}part}(f(n)) = g(n)) \wedge \\ \quad \forall n \in \mathbb{N}.\; (\operatorname{imag\text{-}part}(f(n)) = q(n)) \wedge \\ \quad \operatorname{converges\text{-}to}(\operatorname{rr\text{-}ms}, g, \operatorname{real\text{-}part}(v)) \wedge \\ \quad \operatorname{converges\text{-}to}(\operatorname{rr\text{-}ms}, q, \operatorname{imag\text{-}part}(v)) \Rightarrow \\ \quad \quad \operatorname{converges\text{-}to}(\operatorname{cc\text{-}ms}, f, v) \end{array}}}
\vnbprop{cc-complete}{\vnbstmt{\operatorname{is\text{-}complete}(\operatorname{cc\text{-}ms})}}

\vnbbookfile{theorem-library/series-abs-converges}{theorem-library/series-abs-converges}
\noindent{\small series-abs-converges.scm -- ABSOLUTE CONVERGENCE IMPLIES CONVERGENCE, for real series.}\par
\vnbprop{series-abs-converges}{\vnbstmt{\begin{array}{@{}l@{}} \forall f, g \in (\mathbb{N} \to \mathbb{R}). \\ \quad (\forall j \in \mathbb{N}.\; (g(j) = \lvert f(j) \rvert) \wedge \operatorname{series\text{-}converges}(g)) \Rightarrow  \\ \quad \quad \operatorname{series\text{-}converges}(f) \end{array}}}

\vnbbookfile{theorem-library/mono-le-limit}{theorem-library/mono-le-limit}
\noindent{\small mono-le-limit.scm -- A TERM OF A NONDECREASING CONVERGENT SEQUENCE IS AT MOST ITS LIMIT, and the corollary the product metric's zero law needs.}\par
\vnbprop{rr-mono-le-limit}{\vnbstmt{\begin{array}{@{}l@{}} \forall f \in (\mathbb{N} \to \mathbb{R}), \ell.; \forall m \in \mathbb{N}. \\ \quad (\forall k \in \mathbb{N}.\; (f(k) \leq f((k + 1))) \wedge \operatorname{converges\text{-}to}(\operatorname{rr\text{-}ms}, f, \ell)) \Rightarrow  \\ \quad \quad f(m) \leq \ell \end{array}}}
\vnbprop{series-term-le-sum}{\vnbstmt{\begin{array}{@{}l@{}} \forall f \in (\mathbb{N} \to \mathbb{R}), \ell.; \forall n \in \mathbb{N}. \\ \quad (\forall n \in \mathbb{N}.\; (0 \leq f(n)) \wedge \operatorname{series\text{-}converges\text{-}to}(f, \ell)) \Rightarrow (f(n) \leq \ell) \end{array}}}
\vnbprop{series-nonneg-zero-sum}{\vnbstmt{\begin{array}{@{}l@{}} \forall f \in (\mathbb{N} \to \mathbb{R}), n \in \mathbb{N}. \\ \quad (\forall n \in \mathbb{N}.\; (0 \leq f(n)) \wedge \operatorname{series\text{-}converges\text{-}to}(f, 0)) \Rightarrow (f(n) = 0) \end{array}}}

\vnbbookfile{theorem-library/limit-arithmetic}{theorem-library/limit-arithmetic}
\noindent{\small limit-arithmetic.scm -- THE TWO LIMIT LAWS THE SERIES AND PRODUCT LANES WERE MISSING: `$<$=' carried across TWO limits, and the limit of a pointwise sum.  Plus the series readings of the second.}\par
\vnbprop{rr-limit-le}{\vnbstmt{\begin{array}{@{}l@{}} \forall f, g \in (\mathbb{N} \to \mathbb{R}), v, a \in \mathbb{R}. \\ \quad (\forall n \in \mathbb{N}.\; (f(n) \leq g(n)) \wedge (\operatorname{converges\text{-}to}(\operatorname{rr\text{-}ms}, f, v) \wedge \operatorname{converges\text{-}to}(\operatorname{rr\text{-}ms}, g, a))) \Rightarrow  \\ \quad \quad v \leq a \end{array}}}
\vnbprop{rr-limit-add}{\vnbstmt{\begin{array}{@{}l@{}} \forall f, g, h \in (\mathbb{N} \to \mathbb{R}), v, a \in \mathbb{R}. \\ \quad (\forall j \in \mathbb{N}.\; (h(j) = (f(j) + g(j))) \wedge (\operatorname{converges\text{-}to}(\operatorname{rr\text{-}ms}, f, v) \wedge \operatorname{converges\text{-}to}(\operatorname{rr\text{-}ms}, g, a))) \Rightarrow  \\ \quad \quad \operatorname{converges\text{-}to}(\operatorname{rr\text{-}ms}, h, (v + a)) \end{array}}}
\vnbprop{series-converges-to-converges}{\vnbstmt{\begin{array}{@{}l@{}} \forall f \in (\mathbb{N} \to \mathbb{R}), v \in \mathbb{R}. \\ \quad \operatorname{series\text{-}converges\text{-}to}(f, v) \Rightarrow \operatorname{series\text{-}converges}(f) \end{array}}}
\vnbprop{series-converges-to-add}{\vnbstmt{\begin{array}{@{}l@{}} \forall f, g, h \in (\mathbb{N} \to \mathbb{R}), v, a \in \mathbb{R}. \\ \quad (\forall i \in \mathbb{N}.\; (h(i) = (f(i) + g(i))) \wedge (\operatorname{series\text{-}converges\text{-}to}(f, v) \wedge \operatorname{series\text{-}converges\text{-}to}(g, a))) \Rightarrow  \\ \quad \quad \operatorname{series\text{-}converges\text{-}to}(h, (v + a)) \end{array}}}
\vnbprop{series-limit-add}{\vnbstmt{\begin{array}{@{}l@{}} \forall f, g, h \in (\mathbb{N} \to \mathbb{R}). \\ \quad (\forall i \in \mathbb{N}.\; (h(i) = (f(i) + g(i))) \wedge (\operatorname{series\text{-}converges}(f) \wedge \operatorname{series\text{-}converges}(g))) \Rightarrow  \\ \quad \quad \operatorname{series\text{-}limit}(h) = (\operatorname{series\text{-}limit}(f) + \operatorname{series\text{-}limit}(g)) \end{array}}}

\vnbbookfile{theorem-library/ell-two}{theorem-library/ell-two}
\noindent{\small ell-two.scm -- ELL-TWO, the square-summable complex sequences, and the closure of that carrier under the pointwise vector operations.}\par
\par\medskip\noindent\textbf{Definition} (\texttt{ell-two}).\;
\begin{equation*}
\begin{array}{@{}l@{}} \operatorname{ell\text{-}two} \\ =\; \{\operatorname{x\_} \in \mathbb{C}^{\mathbb{N}} : \operatorname{series\text{-}converges}((\lambda \operatorname{k\_} \in \mathbb{N}.\; \cdot(\operatorname{magnitude}(\operatorname{x\_}(\operatorname{k\_})), \operatorname{magnitude}(\operatorname{x\_}(\operatorname{k\_})))))\} \end{array}
\end{equation*}

\vnbprop{rr-sq-add-le}{\vnbstmt{\begin{array}{@{}l@{}} \forall u, v \in \mathbb{R}. \\ \quad ((u + v) \cdot (u + v)) \leq (((u \cdot u) + (u \cdot u)) + ((v \cdot v) + (v \cdot v))) \end{array}}}
\vnbprop{cc-magnitude-sq-add-le}{\vnbstmt{\begin{array}{@{}l@{}} \forall p, q \in \mathbb{C}. \\ \quad \operatorname{magnitude}((p + q)) \cdot \operatorname{magnitude}((p + q)) \\ \quad \leq\; ((\operatorname{magnitude}(p) \cdot \operatorname{magnitude}(p)) + (\operatorname{magnitude}(p) \cdot \operatorname{magnitude}(p))) + ((\operatorname{magnitude}(q) \cdot \operatorname{magnitude}(q)) + (\operatorname{magnitude}(q) \cdot \operatorname{magnitude}(q))) \end{array}}}
\vnbprop{ell-two-membership}{\vnbstmt{\begin{array}{@{}l@{}} \forall y. \\ \quad y \in \operatorname{ell\text{-}two} \\ \quad \Leftrightarrow\; ((y \in \mathbb{C}^{\mathbb{N}}) \wedge \operatorname{series\text{-}converges}((\lambda k \in \mathbb{N}.\; (\operatorname{magnitude}(y(k)) \cdot \operatorname{magnitude}(y(k)))))) \end{array}}}
\vnbprop{ell-two-add-closed}{\vnbstmt{\begin{array}{@{}l@{}} \forall x, y \in \operatorname{ell\text{-}two}, z \in \mathbb{C}^{\mathbb{N}}. \\ \quad \forall k \in \mathbb{N}.\; (z(k) = (x(k) + y(k))) \Rightarrow (z \in \operatorname{ell\text{-}two}) \end{array}}}
\vnbprop{zero-seq-in-fun}{\vnbstmt{(\lambda n \in \mathbb{N}.\; 0) \in (\mathbb{N} \to \mathbb{R})}}
\vnbprop{series-partial-sum-zero-seq}{\vnbstmt{\forall k \in \mathbb{N}.\; (\operatorname{series\text{-}partial\text{-}sum}((\lambda n \in \mathbb{N}.\; 0), k) = 0)}}
\vnbprop{zero-series-converges}{\vnbstmt{\operatorname{series\text{-}converges}((\lambda n \in \mathbb{N}.\; 0))}}
\vnbprop{ell-two-zero-in}{\vnbstmt{(\lambda n \in \mathbb{N}.\; 0) \in \operatorname{ell\text{-}two}}}
\vnbprop{ell-two-neg-closed}{\vnbstmt{\begin{array}{@{}l@{}} \forall x \in \operatorname{ell\text{-}two}, z \in \mathbb{C}^{\mathbb{N}}. \\ \quad \forall k \in \mathbb{N}.\; (z(k) = -x(k)) \Rightarrow (z \in \operatorname{ell\text{-}two}) \end{array}}}

\vnbbookfile{theorem-library/dyadic-weights}{theorem-library/dyadic-weights}
\noindent{\small dyadic-weights.scm -- THE CANONICAL PRODUCT-METRIC WEIGHTS ARE SUMMABLE, proven -- retiring `product-metric-default-summable' (structure-library/product-metric.scm), the last asserted support of that file that is not blocked on a bridge somebody else is proving:}\par
\vnbprop{power-closed-at}{\vnbstmt{\forall n \in \mathbb{N}, x \in \mathbb{R}.\; (\operatorname{power}(x, n) \in \mathbb{R})}}
\vnbprop{power-real-closed}{\vnbstmt{\forall x \in \mathbb{R}, n \in \mathbb{N}.\; (\operatorname{power}(x, n) \in \mathbb{R})}}
\vnbprop{rr-zero-lt-two}{\vnbstmt{0 < 2}}
\vnbprop{power-two-pos}{\vnbstmt{\forall k \in \mathbb{N}.\; (0 < \operatorname{power}(2, k))}}
\vnbprop{rr-scaled-inverse-unique}{\vnbstmt{\begin{array}{@{}l@{}} \forall a, u, v \in \mathbb{R}. \\ \quad (((a \cdot u) = 1) \wedge (((2 \cdot a) \cdot v) = 1)) \Rightarrow (u = (2 \cdot v)) \end{array}}}
\vnbprop{recip-power-two-halves}{\vnbstmt{\forall k \in \mathbb{N}.\; (\operatorname{recip}(\operatorname{power}(2, k)) = (2 \cdot \operatorname{recip}((2 \cdot \operatorname{power}(2, k)))))}}
\vnbprop{dyadic-weight-apply}{\vnbstmt{\begin{array}{@{}l@{}} \forall k \in \mathbb{N}. \\ \quad (\lambda n \in \mathbb{N}.\; /(1, \operatorname{power}(2, (n + 1))))(k) = \operatorname{recip}(\operatorname{power}(2, (k + 1))) \end{array}}}
\vnbprop{rr-halving-identity}{\vnbstmt{\forall b \in \mathbb{R}.\; (((1 - (2 \cdot b)) + b) = (1 - b))}}
\vnbprop{dyadic-weight-in-fun}{\vnbstmt{(\lambda n \in \mathbb{N}.\; /(1, \operatorname{power}(2, (n + 1)))) \in (\mathbb{N} \to \mathbb{R})}}
\vnbprop{dyadic-weight-pos}{\vnbstmt{\forall a \in \mathbb{N}.\; (0 < (\lambda n \in \mathbb{N}.\; /(1, \operatorname{power}(2, (n + 1))))(a))}}
\vnbprop{dyadic-partial-sum}{\vnbstmt{\begin{array}{@{}l@{}} \forall k \in \mathbb{N}. \\ \quad \operatorname{series\text{-}partial\text{-}sum}((\lambda n \in \mathbb{N}.\; /(1, \operatorname{power}(2, (n + 1)))), k) = (1 - \operatorname{recip}(\operatorname{power}(2, k))) \end{array}}}
\vnbprop{product-metric-default-summable}{\vnbstmt{\operatorname{summable\text{-}weight}((\lambda n \in \mathbb{N}.\; /(1, \operatorname{power}(2, (n + 1)))))}}

\vnbbookfile{theorem-library/rake-inequalities}{theorem-library/rake-inequalities}
\noindent{\small rake-inequalities.scm -- the real-inequality leaves of the rake (batch 5, batch Q), PROVEN.  Nine asserted supports, and the bridge that four of them turned out to need.}\par
\vnbprop{rr-young-2}{\vnbstmt{\forall a, b \in \mathbb{R}.\; ((a \cdot b) \leq /(((a \cdot a) + (b \cdot b)), 2))}}
\vnbprop{rr-amgm-2}{\vnbstmt{\forall a, b \in \mathbb{R}.\; ((4 \cdot (a \cdot b)) \leq ((a + b) \cdot (a + b)))}}
\vnbprop{rr-qm-am-2}{\vnbstmt{\forall a, b \in \mathbb{R}.\; (((a + b) \cdot (a + b)) \leq (2 \cdot ((a \cdot a) + (b \cdot b))))}}
\vnbprop{rr-cauchy-schwarz-2}{\vnbstmt{\begin{array}{@{}l@{}} \forall a, b, c, d \in \mathbb{R}. \\ \quad (((a \cdot c) + (b \cdot d)) \cdot ((a \cdot c) + (b \cdot d))) \leq (((a \cdot a) + (b \cdot b)) \cdot ((c \cdot c) + (d \cdot d))) \end{array}}}
\vnbprop{rr-pos-shrink}{\vnbstmt{\forall s.\; (\operatorname{pos\text{-}rr}(s) \Rightarrow \exists d.\; (\operatorname{pos\text{-}rr}(d) \wedge (d < s)))}}
\vnbprop{rr-bernoulli-ind}{\vnbstmt{\forall n \in \mathbb{N}, x \in \mathbb{R}.\; ((-1 \leq x) \Rightarrow ((1 + (n \cdot x)) \leq \operatorname{power}((1 + x), n)))}}
\vnbprop{rr-bernoulli}{\vnbstmt{\forall x \in \mathbb{R}, n \in \mathbb{N}.\; ((-1 \leq x) \Rightarrow ((1 + (n \cdot x)) \leq \operatorname{power}((1 + x), n)))}}
\vnbprop{cra-ag-rr-carr}{\vnbstmt{(\operatorname{carr}(\operatorname{commutative\text{-}ring\text{-}additive\text{-}ag}(\operatorname{rr\text{-}normed\text{-}field})) \simeq \mathbb{R})}}
\vnbprop{cra-ag-rr-iden}{\vnbstmt{(\operatorname{iden}(\operatorname{commutative\text{-}ring\text{-}additive\text{-}ag}(\operatorname{rr\text{-}normed\text{-}field})) \simeq 0)}}
\vnbprop{cra-ag-rr-opr}{\vnbstmt{(\operatorname{opr}(\operatorname{commutative\text{-}ring\text{-}additive\text{-}ag}(\operatorname{rr\text{-}normed\text{-}field})) \simeq (\lambda \langle \operatorname{x\_}, \operatorname{y\_} \rangle \in (\mathbb{R} \times \mathbb{R}).\; (\operatorname{x\_} + \operatorname{y\_})))}}
\vnbprop{cra-ag-is-nf-ag}{\vnbstmt{\forall r.\; (\operatorname{commutative\text{-}ring\text{-}additive\text{-}ag}(r) \simeq \operatorname{normed\text{-}field\text{-}additive\text{-}ag}(r))}}
\vnbprop{cra-ag-rr-is-abelian-group}{\vnbstmt{\operatorname{is\text{-}abelian\text{-}group}(\operatorname{commutative\text{-}ring\text{-}additive\text{-}ag}(\operatorname{rr\text{-}normed\text{-}field}))}}
\vnbprop{cra-ag-rr-opr-apply}{\vnbstmt{\begin{array}{@{}l@{}} \forall u, v \in \mathbb{R}. \\ \quad (\operatorname{opr}(\operatorname{commutative\text{-}ring\text{-}additive\text{-}ag}(\operatorname{rr\text{-}normed\text{-}field}))(u, v) \simeq (u + v)) \end{array}}}
\vnbprop{rr-nonneg-has-zero}{\vnbstmt{\operatorname{iden}(\operatorname{commutative\text{-}ring\text{-}additive\text{-}ag}(\operatorname{rr\text{-}normed\text{-}field})) \in \{t \in \mathbb{R} : (0 \leq t)\}}}
\vnbprop{rr-nonneg-opr-closed}{\vnbstmt{\begin{array}{@{}l@{}} \forall x, y \in \{t \in \mathbb{R} : (0 \leq t)\}. \\ \quad \operatorname{opr}(\operatorname{commutative\text{-}ring\text{-}additive\text{-}ag}(\operatorname{rr\text{-}normed\text{-}field}))(x, y) \in \{t \in \mathbb{R} : (0 \leq t)\} \end{array}}}
\vnbprop{finsum-rr-in-rr}{\vnbstmt{\begin{array}{@{}l@{}} \forall s \in \mathbf{Set}, f \in (s \to \mathbb{R}). \\ \quad (|s| \in \mathbb{N}) \Rightarrow  \\ \quad \quad \operatorname{finsum}(\operatorname{commutative\text{-}ring\text{-}additive\text{-}ag}(\operatorname{rr\text{-}normed\text{-}field}), f, s) \in \mathbb{R} \end{array}}}
\vnbprop{finsum-rr-nonneg}{\vnbstmt{\begin{array}{@{}l@{}} \forall s \in \mathbf{Set}, f \in (s \to \mathbb{R}). \\ \quad ((|s| \in \mathbb{N}) \wedge \forall z \in s.\; (0 \leq f(z))) \Rightarrow  \\ \quad \quad 0 \leq \operatorname{finsum}(\operatorname{commutative\text{-}ring\text{-}additive\text{-}ag}(\operatorname{rr\text{-}normed\text{-}field}), f, s) \end{array}}}
\vnbprop{finsum-sq-nonneg}{\vnbstmt{\begin{array}{@{}l@{}} \forall s \in \mathbf{Set}, a \in (s \to \mathbb{R}). \\ \quad (|s| \in \mathbb{N}) \Rightarrow  \\ \quad \quad 0 \leq \operatorname{finsum}(\operatorname{commutative\text{-}ring\text{-}additive\text{-}ag}(\operatorname{rr\text{-}normed\text{-}field}), (\lambda i \in s.\; (a(i) \cdot a(i))), s) \end{array}}}
\vnbprop{finsum-le-termwise}{\vnbstmt{\begin{array}{@{}l@{}} \forall s \in \mathbf{Set}, a, b \in (s \to \mathbb{R}). \\ \quad ((|s| \in \mathbb{N}) \wedge \forall i \in s.\; (a(i) \leq b(i))) \Rightarrow  \\ \quad \quad \operatorname{finsum}(\operatorname{commutative\text{-}ring\text{-}additive\text{-}ag}(\operatorname{rr\text{-}normed\text{-}field}), a, s) \\ \quad \quad \leq\; \operatorname{finsum}(\operatorname{commutative\text{-}ring\text{-}additive\text{-}ag}(\operatorname{rr\text{-}normed\text{-}field}), b, s) \end{array}}}
\vnbprop{finsum-abs-triangle}{\vnbstmt{\begin{array}{@{}l@{}} \forall s \in \mathbf{Set}, a \in (s \to \mathbb{R}). \\ \quad (|s| \in \mathbb{N}) \Rightarrow  \\ \quad \quad \lvert \operatorname{finsum}(\operatorname{commutative\text{-}ring\text{-}additive\text{-}ag}(\operatorname{rr\text{-}normed\text{-}field}), a, s) \rvert \\ \quad \quad \leq\; \operatorname{finsum}(\operatorname{commutative\text{-}ring\text{-}additive\text{-}ag}(\operatorname{rr\text{-}normed\text{-}field}), (\lambda i \in s.\; \lvert a(i) \rvert), s) \end{array}}}

\vnbbookfile{theorem-library/deriv-power}{theorem-library/deriv-power}
\noindent{\small deriv-power.scm -- the POWER RULE for the real monomial, PROVEN `modulo 0'.}\par
\vnbprop{pw-step-identity}{\vnbstmt{\begin{array}{@{}l@{}} \forall c, u, v. \\ \quad ((c \in \mathbb{R}) \wedge ((u \in \mathbb{R}) \wedge (v \in \mathbb{R}))) \Rightarrow  \\ \quad \quad ((c + 1) \cdot (u \cdot v)) = ((u \cdot v) + (u \cdot (c \cdot v))) \end{array}}}
\vnbprop{pow-lam-in-fun}{\vnbstmt{\forall n \in \mathbb{N}.\; ((\lambda x \in \mathbb{R}.\; \operatorname{power}(x, n)) \in (\mathbb{R} \to \mathbb{R}))}}
\vnbprop{deriv-power}{\vnbstmt{\begin{array}{@{}l@{}} \forall n \in \mathbb{N}, a \in \mathbb{R}. \\ \quad \operatorname{is\text{-}diff\text{-}at}((\lambda x \in \mathbb{R}.\; \operatorname{power}(x, (n + 1))), a, ((n + 1) \cdot \operatorname{power}(a, n))) \end{array}}}

\vnbbookfile{theorem-library/deriv-polynomial}{theorem-library/deriv-polynomial}
\noindent{\small deriv-polynomial.scm -- rung 1 of the integration arc: the DERIVATIVE OF A POLYNOMIAL FUNCTION, in the coefficient-lambda representation.}\par
\vnbprop{poly-term-lam-in-fun}{\vnbstmt{\begin{array}{@{}l@{}} \forall f \in (\mathbb{N} \to \mathbb{R}), x \in \mathbb{R}. \\ \quad (\lambda k \in \mathbb{N}.\; (f(k) \cdot \operatorname{power}(x, k))) \in (\mathbb{N} \to \mathbb{R}) \end{array}}}
\vnbprop{poly-dseq-lam-in-fun}{\vnbstmt{\begin{array}{@{}l@{}} \forall f \in (\mathbb{N} \to \mathbb{R}), t \in \mathbb{R}. \\ \quad (\lambda j \in \mathbb{N}.\; (((j + 1) \cdot f((j + 1))) \cdot \operatorname{power}(t, j))) \in (\mathbb{N} \to \mathbb{R}) \end{array}}}
\vnbprop{poly-lam-in-fun}{\vnbstmt{\begin{array}{@{}l@{}} \forall f \in (\mathbb{N} \to \mathbb{R}), n \in \mathbb{N}. \\ \quad (\lambda x \in \mathbb{R}.\; \operatorname{series\text{-}partial\text{-}sum}((\lambda k \in \mathbb{N}.\; (f(k) \cdot \operatorname{power}(x, k))), n)) \in (\mathbb{R} \to \mathbb{R}) \end{array}}}
\vnbprop{poly-step-identity}{\vnbstmt{\begin{array}{@{}l@{}} \forall c, u, v. \\ \quad ((c \in \mathbb{R}) \wedge ((u \in \mathbb{R}) \wedge (v \in \mathbb{R}))) \Rightarrow (((c \cdot u) \cdot v) = (u \cdot (c \cdot v))) \end{array}}}
\vnbprop{deriv-scalar-mult}{\vnbstmt{\begin{array}{@{}l@{}} \forall c \in \mathbb{R}, f, t, l. \\ \quad \operatorname{is\text{-}diff\text{-}at}(f, t, l) \Rightarrow \operatorname{is\text{-}diff\text{-}at}((\lambda x \in \mathbb{R}.\; (c \cdot f(x))), t, (c \cdot l)) \end{array}}}
\vnbprop{deriv-coef-monomial}{\vnbstmt{\begin{array}{@{}l@{}} \forall c \in \mathbb{R}, g \in \mathbb{N}, t \in \mathbb{R}. \\ \quad \operatorname{is\text{-}diff\text{-}at}((\lambda x \in \mathbb{R}.\; (c \cdot \operatorname{power}(x, (g + 1)))), t, (((g + 1) \cdot c) \cdot \operatorname{power}(t, g))) \end{array}}}
\vnbprop{deriv-polynomial}{\vnbstmt{\begin{array}{@{}l@{}} \forall n \in \mathbb{N}, f \in (\mathbb{N} \to \mathbb{R}), t \in \mathbb{R}. \\ \quad \operatorname{is\text{-}diff\text{-}at}((\lambda x \in \mathbb{R}.\; \operatorname{series\text{-}partial\text{-}sum}((\lambda k \in \mathbb{N}.\; (f(k) \cdot \operatorname{power}(x, k))), (n + 1))), \\ \quad \quad t, \\ \quad \quad \operatorname{series\text{-}partial\text{-}sum}((\lambda j \in \mathbb{N}.\; (((j + 1) \cdot f((j + 1))) \cdot \operatorname{power}(t, j))), n)) \end{array}}}

\vnbbookfile{theorem-library/poly-antiderivative}{theorem-library/poly-antiderivative}
\noindent{\small poly-antiderivative.scm -- calculus.pdf EXAMPLE 4.7, "every polynomial function is antiderivable", with the antiderivative written out:}\par
\vnbprop{anti-term-in-rr}{\vnbstmt{\begin{array}{@{}l@{}} \forall f \in (\mathbb{N} \to \mathbb{R}), k \in \mathbb{N}, x \in \mathbb{R}. \\ \quad ((\operatorname{recip}((k + 1)) \cdot f(k)) \cdot \operatorname{power}(x, (k + 1))) \in \mathbb{R} \end{array}}}
\vnbprop{anti-term-lam-in-fun}{\vnbstmt{\begin{array}{@{}l@{}} \forall f \in (\mathbb{N} \to \mathbb{R}), x \in \mathbb{R}. \\ \quad (\lambda k \in \mathbb{N}.\; ((\operatorname{recip}((k + 1)) \cdot f(k)) \cdot \operatorname{power}(x, (k + 1)))) \in (\mathbb{N} \to \mathbb{R}) \end{array}}}
\vnbprop{anti-lam-in-fun}{\vnbstmt{\begin{array}{@{}l@{}} \forall f \in (\mathbb{N} \to \mathbb{R}), m \in \mathbb{N}. \\ \quad (\lambda x \in \mathbb{R}.\; \operatorname{series\text{-}partial\text{-}sum}((\lambda k \in \mathbb{N}.\; ((\operatorname{recip}((k + 1)) \cdot f(k)) \cdot \operatorname{power}(x, (k + 1)))), m)) \\ \quad \in\; (\mathbb{R} \to \mathbb{R}) \end{array}}}
\vnbprop{rr-zero-plus}{\vnbstmt{\forall t \in \mathbb{R}.\; (t = t)}}
\vnbprop{recip-succ-cancel}{\vnbstmt{\begin{array}{@{}l@{}} \forall m \in \mathbb{N}, c, v \in \mathbb{R}. \\ \quad (c \cdot v) = (((m + 1) \cdot (\operatorname{recip}((m + 1)) \cdot c)) \cdot v) \end{array}}}
\vnbprop{deriv-anti-monomial}{\vnbstmt{\begin{array}{@{}l@{}} \forall c \in \mathbb{R}, m \in \mathbb{N}, t \in \mathbb{R}. \\ \quad \operatorname{is\text{-}diff\text{-}at}((\lambda x \in \mathbb{R}.\; ((\operatorname{recip}((m + 1)) \cdot c) \cdot \operatorname{power}(x, (m + 1)))), t, (c \cdot \operatorname{power}(t, m))) \end{array}}}
\vnbprop{poly-antiderivative}{\vnbstmt{\begin{array}{@{}l@{}} \forall n \in \mathbb{N}, f \in (\mathbb{N} \to \mathbb{R}), t \in \mathbb{R}. \\ \quad \operatorname{is\text{-}diff\text{-}at}((\lambda x \in \mathbb{R}.\; \\ \quad \quad \quad \operatorname{series\text{-}partial\text{-}sum}((\lambda k \in \mathbb{N}.\; ((\operatorname{recip}((k + 1)) \cdot f(k)) \cdot \operatorname{power}(x, (k + 1)))), (n + 1))), \\ \quad \quad t, \\ \quad \quad \operatorname{series\text{-}partial\text{-}sum}((\lambda k \in \mathbb{N}.\; (f(k) \cdot \operatorname{power}(t, k))), (n + 1))) \end{array}}}

\vnbbookfile{theorem-library/product-summable}{theorem-library/product-summable}
\noindent{\small product-summable.scm -- THE DEFINING SERIES OF THE PRODUCT METRIC CONVERGES, proven -- retiring the asserted `product-weighted-summable' (structure-library/product-metric.scm), which is what makes the IOTA in PRODUCT-METRIC-W a well-defined distance, and with it the second asserted}\par
\vnbprop{product-carrier-unfold}{\vnbstmt{\begin{array}{@{}l@{}} \forall s. \\ \quad (\operatorname{product\text{-}carrier}(s) \simeq \{x \in (\mathbb{N} \to \operatorname{big\text{-}union}(n, \mathbb{N}, \operatorname{pts}(s(n)))) : \forall n \in \mathbb{N}.\; (x(n) \in \operatorname{pts}(s(n)))\}) \end{array}}}
\vnbprop{product-carrier-coord}{\vnbstmt{\forall s.; \forall x \in \operatorname{product\text{-}carrier}(s), n \in \mathbb{N}.\; (x(n) \in \operatorname{pts}(s(n)))}}
\vnbprop{bdd-metric-weight-bound}{\vnbstmt{\begin{array}{@{}l@{}} \forall s.; \forall u, v \in \operatorname{pts}(s), n \in \mathbb{R}. \\ \quad (\operatorname{is\text{-}metric\text{-}space}(s) \wedge (0 \leq n)) \Rightarrow  \\ \quad \quad (n \cdot \operatorname{dist}(\operatorname{bdd\text{-}metric}(s))(u, v)) \in \mathbb{R} \wedge \\ \quad \quad 0 \leq (n \cdot \operatorname{dist}(\operatorname{bdd\text{-}metric}(s))(u, v)) \wedge \\ \quad \quad (n \cdot \operatorname{dist}(\operatorname{bdd\text{-}metric}(s))(u, v)) \leq n \end{array}}}
\vnbprop{product-metric-carrier}{\vnbstmt{\forall s, w.\; (\operatorname{pts}(\operatorname{product\text{-}metric\text{-}w}(s, w)) \simeq \operatorname{product\text{-}carrier}(s))}}
\vnbprop{product-term-seq-in-fun}{\vnbstmt{\begin{array}{@{}l@{}} \forall s, w.; \forall x, y \in \operatorname{product\text{-}carrier}(s). \\ \quad (\operatorname{is\text{-}ms\text{-}sequence}(s) \wedge \operatorname{summable\text{-}weight}(w)) \Rightarrow  \\ \quad \quad (\lambda n \in \mathbb{N}.\; (w(n) \cdot \operatorname{dist}(\operatorname{bdd\text{-}metric}(s(n)))(x(n), y(n)))) \in (\mathbb{N} \to \mathbb{R}) \end{array}}}
\vnbprop{product-weighted-summable}{\vnbstmt{\begin{array}{@{}l@{}} \forall s, w.; \forall x, y \in \operatorname{product\text{-}carrier}(s). \\ \quad (\operatorname{is\text{-}ms\text{-}sequence}(s) \wedge \operatorname{summable\text{-}weight}(w)) \Rightarrow  \\ \quad \quad \operatorname{series\text{-}converges}((\lambda n \in \mathbb{N}.\; (w(n) \cdot \operatorname{dist}(\operatorname{bdd\text{-}metric}(s(n)))(x(n), y(n))))) \end{array}}}
\vnbprop{product-metric-distance}{\vnbstmt{\begin{array}{@{}l@{}} \forall s, w.; \forall x, y \in \operatorname{product\text{-}carrier}(s). \\ \quad (\operatorname{dist}(\operatorname{product\text{-}metric\text{-}w}(s, w))(x, y) \simeq \operatorname{series\text{-}limit}((\lambda n \in \mathbb{N}.\; (w(n) \cdot \operatorname{dist}(\operatorname{bdd\text{-}metric}(s(n)))(x(n), y(n)))))) \end{array}}}
\vnbprop{product-metric-dist-converges}{\vnbstmt{\begin{array}{@{}l@{}} \forall s, w.; \forall x, y \in \operatorname{product\text{-}carrier}(s). \\ \quad (\operatorname{is\text{-}ms\text{-}sequence}(s) \wedge \operatorname{summable\text{-}weight}(w)) \Rightarrow  \\ \quad \quad \operatorname{series\text{-}converges\text{-}to}((\lambda n \in \mathbb{N}.\; (w(n) \cdot \operatorname{dist}(\operatorname{bdd\text{-}metric}(s(n)))(x(n), y(n)))), \\ \quad \quad \quad \operatorname{dist}(\operatorname{product\text{-}metric\text{-}w}(s, w))(x, y)) \end{array}}}

\vnbbookfile{theorem-library/product-is-metric-space}{theorem-library/product-is-metric-space}
\noindent{\small product-is-metric-space.scm -- THE COUNTABLE PRODUCT OF METRIC SPACES IS A METRIC SPACE, proven -- retiring the asserted `product-is-metric-space' of structure-library/product-metric.scm (line 79).}\par
\vnbprop{series-limit-le}{\vnbstmt{\begin{array}{@{}l@{}} \forall f, g \in (\mathbb{N} \to \mathbb{R}), v, a \in \mathbb{R}. \\ \quad (\forall n \in \mathbb{N}.\; (f(n) \leq g(n)) \wedge (\operatorname{series\text{-}converges\text{-}to}(f, v) \wedge \operatorname{series\text{-}converges\text{-}to}(g, a))) \Rightarrow  \\ \quad \quad v \leq a \end{array}}}
\vnbprop{series-limit-iota}{\vnbstmt{\forall f.\; ((\iota\, l.\;\operatorname{series\text{-}converges\text{-}to}(f, l)) \simeq \operatorname{series\text{-}limit}(f))}}
\vnbprop{product-dist-in-rr}{\vnbstmt{\begin{array}{@{}l@{}} \forall s, w.; \forall x, y \in \operatorname{product\text{-}carrier}(s). \\ \quad (\operatorname{is\text{-}ms\text{-}sequence}(s) \wedge \operatorname{summable\text{-}weight}(w)) \Rightarrow  \\ \quad \quad \operatorname{dist}(\operatorname{product\text{-}metric\text{-}w}(s, w))(x, y) \in \mathbb{R} \end{array}}}
\vnbprop{product-term-nonneg}{\vnbstmt{\begin{array}{@{}l@{}} \forall s, w.; \forall x, y \in \operatorname{product\text{-}carrier}(s), a \in \mathbb{N}. \\ \quad (\operatorname{is\text{-}ms\text{-}sequence}(s) \wedge \operatorname{summable\text{-}weight}(w)) \Rightarrow  \\ \quad \quad 0 \leq (\lambda n \in \mathbb{N}.\; (w(n) \cdot \operatorname{dist}(\operatorname{bdd\text{-}metric}(s(n)))(x(n), y(n))))(a) \end{array}}}
\vnbprop{product-term-self-zero}{\vnbstmt{\begin{array}{@{}l@{}} \forall s, w.; \forall x \in \operatorname{product\text{-}carrier}(s), a \in \mathbb{N}. \\ \quad (\operatorname{is\text{-}ms\text{-}sequence}(s) \wedge \operatorname{summable\text{-}weight}(w)) \Rightarrow  \\ \quad \quad (\lambda n \in \mathbb{N}.\; (w(n) \cdot \operatorname{dist}(\operatorname{bdd\text{-}metric}(s(n)))(x(n), x(n))))(a) = 0 \end{array}}}
\vnbprop{product-term-sym}{\vnbstmt{\begin{array}{@{}l@{}} \forall s, w.; \forall x, y \in \operatorname{product\text{-}carrier}(s), a \in \mathbb{N}. \\ \quad (\operatorname{is\text{-}ms\text{-}sequence}(s) \wedge \operatorname{summable\text{-}weight}(w)) \Rightarrow  \\ \quad \quad (\lambda n \in \mathbb{N}.\; (w(n) \cdot \operatorname{dist}(\operatorname{bdd\text{-}metric}(s(n)))(x(n), y(n))))(a) \\ \quad \quad \leq\; (\lambda n \in \mathbb{N}.\; (w(n) \cdot \operatorname{dist}(\operatorname{bdd\text{-}metric}(s(n)))(y(n), x(n))))(a) \end{array}}}
\vnbprop{product-sum-seq-in-fun}{\vnbstmt{\begin{array}{@{}l@{}} \forall s, w.; \forall x, y, z \in \operatorname{product\text{-}carrier}(s). \\ \quad (\operatorname{is\text{-}ms\text{-}sequence}(s) \wedge \operatorname{summable\text{-}weight}(w)) \Rightarrow  \\ \quad \quad (\lambda n \in \mathbb{N}.\; \\ \quad \quad \quad (w(n) \cdot \operatorname{dist}(\operatorname{bdd\text{-}metric}(s(n)))(x(n), y(n))) + (w(n) \cdot \operatorname{dist}(\operatorname{bdd\text{-}metric}(s(n)))(y(n), z(n)))) \\ \quad \quad \in\; (\mathbb{N} \to \mathbb{R}) \end{array}}}
\vnbprop{product-sum-seq-split}{\vnbstmt{\begin{array}{@{}l@{}} \forall s, w.; \forall x, y, z \in \operatorname{product\text{-}carrier}(s), i \in \mathbb{N}. \\ \quad (\operatorname{is\text{-}ms\text{-}sequence}(s) \wedge \operatorname{summable\text{-}weight}(w)) \Rightarrow  \\ \quad \quad (\lambda n \in \mathbb{N}.\; ((w(n) \cdot \operatorname{dist}(\operatorname{bdd\text{-}metric}(s(n)))(x(n), y(n))) + (w(n) \cdot \operatorname{dist}(\operatorname{bdd\text{-}metric}(s(n)))(y(n), z(n)))))(i) \\ \quad \quad =\; (\lambda n \in \mathbb{N}.\; (w(n) \cdot \operatorname{dist}(\operatorname{bdd\text{-}metric}(s(n)))(x(n), y(n))))(i) + (\lambda n \in \mathbb{N}.\; (w(n) \cdot \operatorname{dist}(\operatorname{bdd\text{-}metric}(s(n)))(y(n), z(n))))(i) \end{array}}}
\vnbprop{product-term-triangle}{\vnbstmt{\begin{array}{@{}l@{}} \forall s, w.; \forall x, y, z \in \operatorname{product\text{-}carrier}(s), a \in \mathbb{N}. \\ \quad (\operatorname{is\text{-}ms\text{-}sequence}(s) \wedge \operatorname{summable\text{-}weight}(w)) \Rightarrow  \\ \quad \quad (\lambda n \in \mathbb{N}.\; (w(n) \cdot \operatorname{dist}(\operatorname{bdd\text{-}metric}(s(n)))(x(n), z(n))))(a) \\ \quad \quad \leq\; (\lambda n \in \mathbb{N}.\; ((w(n) \cdot \operatorname{dist}(\operatorname{bdd\text{-}metric}(s(n)))(x(n), y(n))) + (w(n) \cdot \operatorname{dist}(\operatorname{bdd\text{-}metric}(s(n)))(y(n), z(n)))))(a) \end{array}}}
\vnbprop{product-is-metric-space}{\vnbstmt{\begin{array}{@{}l@{}} \forall s, w. \\ \quad (\operatorname{is\text{-}ms\text{-}sequence}(s) \wedge \operatorname{summable\text{-}weight}(w)) \Rightarrow  \\ \quad \quad \operatorname{is\text{-}metric\text{-}space}(\operatorname{product\text{-}metric\text{-}w}(s, w)) \end{array}}}

\vnbbookfile{theorem-library/rake-block-tower}{theorem-library/rake-block-tower}
\noindent{\small theorem-library/rake-block-tower.scm ==================================================================== convergence-block-tower (theorem-library/seq-compact-product.scm:91) -- the refinement-tower core of the cross-coordinate diagonal argument, and}\par
\vnbprop{block-step-converges}{\vnbstmt{\begin{array}{@{}l@{}} \forall s.; \forall h \in (\mathbb{N} \to \operatorname{pts}(s)), j \in \mathrm{Inf}(\mathbb{N}). \\ \quad \operatorname{seq\text{-}compact}(s) \Rightarrow  \\ \quad \quad \exists b \in \mathrm{Inf}(\mathbb{N}). \\ \quad \quad \quad (b \subseteq j) \wedge \exists p \in \operatorname{pts}(s).\; \operatorname{converges\text{-}along}(s, h, b, p) \end{array}}}
\vnbprop{convergence-block-tower}{\vnbstmt{\begin{array}{@{}l@{}} \forall a.; \forall q \in (\mathbb{N} \to \operatorname{product\text{-}carrier}(a)). \\ \quad (\operatorname{is\text{-}ms\text{-}sequence}(a) \wedge \forall n \in \mathbb{N}.\; \operatorname{seq\text{-}compact}(a(n))) \Rightarrow  \\ \quad \quad \exists s \in (\mathbb{N} \to \mathrm{Inf}(\mathbb{N})). \\ \quad \quad \quad \forall k \in \mathbb{N}.\; (s((k + 1)) \subseteq s(k)) \wedge \\ \quad \quad \quad \exists l \in \operatorname{product\text{-}carrier}(a). \forall n \in \mathbb{N}. \\ \quad \quad \quad \quad \operatorname{converges\text{-}along}(a(n), (\lambda i \in \mathbb{N}.\; q(i)(n)), s(n), l(n)) \end{array}}}

\vnbbookfile{theorem-library/rr-null-scale}{theorem-library/rr-null-scale}
\noindent{\small rr-null-scale.scm -- a NONNEGATIVE CONSTANT MULTIPLE OF A NULL SEQUENCE IS NULL, proven, together with the eps/delta lemma that carries it.}\par
\vnbprop{rr-scale-eps}{\vnbstmt{\begin{array}{@{}l@{}} \forall c, s \in \mathbb{R}. \\ \quad ((0 \leq c) \wedge \operatorname{pos\text{-}rr}(s)) \Rightarrow  \\ \quad \quad \exists d. \\ \quad \quad \quad \operatorname{pos\text{-}rr}(d) \wedge \forall t \in \mathbb{R}.\; (((0 \leq t) \wedge (t \leq d)) \Rightarrow ((c \cdot t) \leq s)) \end{array}}}
\vnbprop{rr-null-scale}{\vnbstmt{\begin{array}{@{}l@{}} \forall c \in \mathbb{R}, f, h \in (\mathbb{N} \to \mathbb{R}). \\ \quad ((0 \leq c) \wedge (\operatorname{converges\text{-}to}(\operatorname{rr\text{-}ms}, f, 0) \wedge \forall j \in \mathbb{N}.\; (h(j) = (c \cdot f(j))))) \Rightarrow  \\ \quad \quad \operatorname{converges\text{-}to}(\operatorname{rr\text{-}ms}, h, 0) \end{array}}}

\vnbbookfile{theorem-library/bdd-metric-convergence}{theorem-library/bdd-metric-convergence}
\noindent{\small bdd-metric-convergence.scm -- CONVERGENCE IS THE SAME IN d AND IN d/(1+d), proven: a sequence converges to L in BDD-METRIC(s) if and only if it converges to L in s.}\par
\vnbprop{bdd-fn-reflect}{\vnbstmt{\begin{array}{@{}l@{}} \forall a, b \in \mathbb{R}. \\ \quad ((0 \leq a) \wedge ((0 \leq b) \wedge (/(a, (1 + a)) \leq /(b, (1 + b))))) \Rightarrow (a \leq b) \end{array}}}
\vnbprop{bdd-metric-dist-le}{\vnbstmt{\begin{array}{@{}l@{}} \forall s.; \forall x, y \in \operatorname{pts}(s). \\ \quad \operatorname{is\text{-}metric\text{-}space}(s) \Rightarrow (\operatorname{dist}(\operatorname{bdd\text{-}metric}(s))(x, y) \leq \operatorname{dist}(s)(x, y)) \end{array}}}
\vnbprop{bdd-metric-dist-reflect}{\vnbstmt{\begin{array}{@{}l@{}} \forall s.; \forall x, y \in \operatorname{pts}(s), e \in \mathbb{R}. \\ \quad (\operatorname{is\text{-}metric\text{-}space}(s) \wedge ((0 \leq e) \wedge (\operatorname{dist}(\operatorname{bdd\text{-}metric}(s))(x, y) \leq /(e, (1 + e))))) \Rightarrow  \\ \quad \quad \operatorname{dist}(s)(x, y) \leq e \end{array}}}
\vnbprop{bdd-metric-converges-fwd}{\vnbstmt{\begin{array}{@{}l@{}} \forall s.; \forall f \in (\mathbb{N} \to \operatorname{pts}(s)), \ell \in \operatorname{pts}(s). \\ \quad (\operatorname{is\text{-}metric\text{-}space}(s) \wedge \operatorname{converges\text{-}to}(s, f, \ell)) \Rightarrow  \\ \quad \quad \operatorname{converges\text{-}to}(\operatorname{bdd\text{-}metric}(s), f, \ell) \end{array}}}
\vnbprop{bdd-metric-converges-bwd}{\vnbstmt{\begin{array}{@{}l@{}} \forall s.; \forall f \in (\mathbb{N} \to \operatorname{pts}(s)), \ell \in \operatorname{pts}(s). \\ \quad (\operatorname{is\text{-}metric\text{-}space}(s) \wedge \operatorname{converges\text{-}to}(\operatorname{bdd\text{-}metric}(s), f, \ell)) \Rightarrow  \\ \quad \quad \operatorname{converges\text{-}to}(s, f, \ell) \end{array}}}
\vnbprop{bdd-metric-converges-iff}{\vnbstmt{\begin{array}{@{}l@{}} \forall s.; \forall f \in (\mathbb{N} \to \operatorname{pts}(s)), \ell \in \operatorname{pts}(s). \\ \quad \operatorname{is\text{-}metric\text{-}space}(s) \Rightarrow  \\ \quad \quad \operatorname{converges\text{-}to}(\operatorname{bdd\text{-}metric}(s), f, \ell) \Leftrightarrow \operatorname{converges\text{-}to}(s, f, \ell) \end{array}}}

\vnbbookfile{theorem-library/rake-metric-constructions}{theorem-library/rake-metric-constructions}
\noindent{\small rake-metric-constructions.scm -- BATCH S of the 2026-09-18 rake: the metric CONSTRUCTIONS (bounded metric, ball cover, countable product).  Every statement below is its support's statement UNCHANGED, copied from the support site, except where the header says otherwise.}\par
\vnbprop{bdd-metric-id-bicontinuous}{\vnbstmt{\begin{array}{@{}l@{}} \forall s. \\ \quad \operatorname{is\text{-}metric\text{-}space}(s) \Rightarrow  \\ \quad \quad \operatorname{is\text{-}continuous}(s, \operatorname{bdd\text{-}metric}(s), (\lambda x \in \operatorname{pts}(s).\; x)) \wedge \\ \quad \quad \operatorname{is\text{-}continuous}(\operatorname{bdd\text{-}metric}(s), s, (\lambda x \in \operatorname{pts}(\operatorname{bdd\text{-}metric}(s)).\; x)) \end{array}}}
\vnbprop{bdd-rr-ms-is-metric-space}{\vnbstmt{\operatorname{is\text{-}metric\text{-}space}(\operatorname{bdd\text{-}metric}(\operatorname{rr\text{-}ms}))}}
\vnbprop{bdd-rr-ms-bounded}{\vnbstmt{\forall x, y \in \mathbb{R}.\; (\operatorname{dist}(\operatorname{bdd\text{-}metric}(\operatorname{rr\text{-}ms}))(x, y) < 1)}}
\vnbprop{bdd-rr-equivalent}{\vnbstmt{\begin{array}{@{}l@{}} \operatorname{is\text{-}continuous}(\operatorname{rr\text{-}ms}, \operatorname{bdd\text{-}metric}(\operatorname{rr\text{-}ms}), (\lambda x \in \mathbb{R}.\; x)) \wedge \\ \operatorname{is\text{-}continuous}(\operatorname{bdd\text{-}metric}(\operatorname{rr\text{-}ms}), \operatorname{rr\text{-}ms}, (\lambda x \in \mathbb{R}.\; x)) \end{array}}}
\vnbprop{ball-cover-is-open-cover}{\vnbstmt{\begin{array}{@{}l@{}} \forall s.; \forall r \in \mathbb{R}. \\ \quad (\operatorname{is\text{-}metric\text{-}space}(s) \wedge ((0 \leq r) \wedge \neg (0 = r))) \Rightarrow  \\ \quad \quad \operatorname{is\text{-}open\text{-}cover}(s, \operatorname{ball\text{-}cover}(s, r)) \end{array}}}
\vnbprop{product-projection-continuous}{\vnbstmt{\begin{array}{@{}l@{}} \forall s, w.; \forall n \in \mathbb{N}. \\ \quad (\operatorname{is\text{-}ms\text{-}sequence}(s) \wedge \operatorname{summable\text{-}weight}(w)) \Rightarrow  \\ \quad \quad \operatorname{is\text{-}continuous}(\operatorname{product\text{-}metric\text{-}w}(s, w), s(n), \operatorname{product\text{-}proj}(s, n)) \end{array}}}

\vnbbookfile{theorem-library/rake-rr-bounded-ms}{theorem-library/rake-rr-bounded-ms}
\noindent{\small rake-rr-bounded-ms.scm -- the three facts about the constant RR-BOUNDED-MS (rake batch 5, integrator, 2026-09-18).}\par
\vnbprop{rr-bounded-ms-is-metric-space}{\vnbstmt{\operatorname{is\text{-}metric\text{-}space}(\operatorname{rr\text{-}bounded\text{-}ms})}}
\vnbprop{rr-bounded-ms-bounded}{\vnbstmt{\forall x, y \in \mathbb{R}.\; (\operatorname{dist}(\operatorname{rr\text{-}bounded\text{-}ms})(x, y) < 1)}}
\vnbprop{rr-bounded-equivalent}{\vnbstmt{\begin{array}{@{}l@{}} \operatorname{is\text{-}continuous}(\operatorname{rr\text{-}ms}, \operatorname{rr\text{-}bounded\text{-}ms}, (\lambda x \in \mathbb{R}.\; x)) \wedge \\ \operatorname{is\text{-}continuous}(\operatorname{rr\text{-}bounded\text{-}ms}, \operatorname{rr\text{-}ms}, (\lambda x \in \mathbb{R}.\; x)) \end{array}}}

\vnbbookfile{theorem-library/rake-compact-equivalences}{theorem-library/rake-compact-equivalences}
\noindent{\small rake-compact-equivalences.scm -- Prop 3.12, the two remaining equivalences (batch 8, assignment 8-A; helper prefix `r9a-').}\par
\vnbprop{compact-iff-cluster-point}{\vnbstmt{\begin{array}{@{}l@{}} \forall s. \\ \quad \operatorname{is\text{-}metric\text{-}space}(s) \Rightarrow  \\ \quad \quad \operatorname{is\text{-}compact}(s) \Leftrightarrow \forall f \in (\mathbb{N} \to \operatorname{pts}(s)).\; \exists x.\; \operatorname{cluster\text{-}point}(s, f, x) \end{array}}}
\vnbprop{compact-iff-tb-complete}{\vnbstmt{\begin{array}{@{}l@{}} \forall s. \\ \quad \operatorname{is\text{-}metric\text{-}space}(s) \Rightarrow  \\ \quad \quad \operatorname{is\text{-}compact}(s) \Leftrightarrow (\operatorname{totally\text{-}bounded}(s) \wedge \operatorname{is\text{-}complete}(s)) \end{array}}}

\vnbbookfile{theorem-library/rake-compact-fip}{theorem-library/rake-compact-fip}
\noindent{\small rake-compact-fip.scm -- Prop 3.12 (1) $<$=$>$ (2), the finite-intersection characterisation of compactness, PROVEN (2026-09-20, batch 11-E; helper prefix `rcf-').}\par
\vnbprop{compact-implies-fip}{\vnbstmt{\begin{array}{@{}l@{}} \forall s, c. \\ \quad (\operatorname{is\text{-}metric\text{-}space}(s) \wedge (\operatorname{is\text{-}compact}(s) \wedge \operatorname{has\text{-}fip}(s, c))) \Rightarrow  \\ \quad \quad \exists p.\; (p \in \operatorname{intersection\text{-}of}(c)) \end{array}}}
\vnbprop{fip-implies-compact}{\vnbstmt{\begin{array}{@{}l@{}} \forall s. \\ \quad (\operatorname{is\text{-}metric\text{-}space}(s) \wedge \forall c.\; (\operatorname{has\text{-}fip}(s, c) \Rightarrow \exists p.\; (p \in \operatorname{intersection\text{-}of}(c)))) \Rightarrow  \\ \quad \quad \operatorname{is\text{-}compact}(s) \end{array}}}
\vnbprop{compact-iff-fip}{\vnbstmt{\begin{array}{@{}l@{}} \forall s. \\ \quad \operatorname{is\text{-}metric\text{-}space}(s) \Rightarrow  \\ \quad \quad \operatorname{is\text{-}compact}(s) \Leftrightarrow \forall c.\; (\operatorname{has\text{-}fip}(s, c) \Rightarrow \exists p.\; (p \in \operatorname{intersection\text{-}of}(c))) \end{array}}}

\vnbbookfile{theorem-library/compact-separable-proof}{theorem-library/compact-separable-proof}
\noindent{\small theorem-library/compact-separable-proof.scm}\par
\vnbprop{compact-metric-is-separable}{\vnbstmt{\forall s.\; ((\operatorname{is\text{-}compact}(s) \wedge \exists x.\; (x \in \operatorname{pts}(s))) \Rightarrow \operatorname{is\text{-}separable}(s))}}

\vnbbookfile{theorem-library/rake-ascoli}{theorem-library/rake-ascoli}
\noindent{\small theorem-library/rake-ascoli.scm ==================================================================== Ascoli-Arzela (sequential, RR-valued) -- batch 8-B.}\par
\vnbprop{equicont-subseq}{\vnbstmt{\begin{array}{@{}l@{}} \forall s, t, m, i. \\ \quad (\operatorname{is\text{-}equicontinuous}(s, t, m) \wedge \operatorname{strictly\text{-}mono\text{-}nn}(i)) \Rightarrow  \\ \quad \quad \operatorname{is\text{-}equicontinuous}(s, t, \operatorname{subseq}(m, i)) \end{array}}}
\vnbprop{ascoli-sequential-from-diagonal}{\vnbstmt{\begin{array}{@{}l@{}} \forall s, m. \\ \quad \operatorname{is\text{-}compact}(s) \wedge \\ \quad \exists x.\; (x \in \operatorname{pts}(s)) \wedge \\ \quad m \in (\mathbb{N} \to (\operatorname{pts}(s) \to \mathbb{R})) \wedge \\ \quad \forall k \in \mathbb{N}.\; \operatorname{is\text{-}continuous}(s, \operatorname{rr\text{-}ms}, m(k)) \wedge \\ \quad \operatorname{is\text{-}equicontinuous}(s, \operatorname{rr\text{-}ms}, m) \wedge \\ \quad \operatorname{pointwise\text{-}bounded}(s, m) \wedge \\ \quad \forall q. \\ \quad \quad \operatorname{is\text{-}dense\text{-}seq}(s, q) \Rightarrow  \\ \quad \quad \quad \exists l.\; (\operatorname{strictly\text{-}mono\text{-}nn}(l) \wedge \operatorname{converges\text{-}on}(s, \operatorname{subseq}(m, l), q)) \Rightarrow \\ \quad \quad \exists l. \\ \quad \quad \quad \operatorname{strictly\text{-}mono\text{-}nn}(l) \wedge \\ \quad \quad \quad \exists g \in (\operatorname{pts}(s) \to \mathbb{R}). \\ \quad \quad \quad \quad \operatorname{is\text{-}continuous}(s, \operatorname{rr\text{-}ms}, g) \wedge \operatorname{converges\text{-}uniformly}(s, \operatorname{subseq}(m, l), g) \end{array}}}

\vnbbookfile{theorem-library/converges-dist-null}{theorem-library/converges-dist-null}
\noindent{\small converges-dist-null.scm -- CONVERGENCE IN A METRIC SPACE IS A NULL SEQUENCE OF REAL DISTANCES, proven both ways:}\par
\vnbprop{dist-seq-apply}{\vnbstmt{\begin{array}{@{}l@{}} \forall t, f, v.; \forall n \in \mathbb{N}. \\ \quad ((\lambda j \in \mathbb{N}.\; \operatorname{dist}(t)(f(j), v))(n) \simeq \operatorname{dist}(t)(f(n), v)) \end{array}}}
\vnbprop{dist-seq-in-fun}{\vnbstmt{\begin{array}{@{}l@{}} \forall t.; \forall f \in (\mathbb{N} \to \operatorname{pts}(t)), v \in \operatorname{pts}(t). \\ \quad \operatorname{is\text{-}metric\text{-}space}(t) \Rightarrow ((\lambda j \in \mathbb{N}.\; \operatorname{dist}(t)(f(j), v)) \in (\mathbb{N} \to \mathbb{R})) \end{array}}}
\vnbprop{converges-dist-null-fwd}{\vnbstmt{\begin{array}{@{}l@{}} \forall t.; \forall f \in (\mathbb{N} \to \operatorname{pts}(t)), v \in \operatorname{pts}(t). \\ \quad (\operatorname{is\text{-}metric\text{-}space}(t) \wedge \operatorname{converges\text{-}to}(t, f, v)) \Rightarrow  \\ \quad \quad \operatorname{converges\text{-}to}(\operatorname{rr\text{-}ms}, (\lambda j \in \mathbb{N}.\; \operatorname{dist}(t)(f(j), v)), 0) \end{array}}}
\vnbprop{converges-dist-null-bwd}{\vnbstmt{\begin{array}{@{}l@{}} \forall t.; \forall f \in (\mathbb{N} \to \operatorname{pts}(t)), v \in \operatorname{pts}(t). \\ \quad (\operatorname{is\text{-}metric\text{-}space}(t) \wedge \operatorname{converges\text{-}to}(\operatorname{rr\text{-}ms}, (\lambda j \in \mathbb{N}.\; \operatorname{dist}(t)(f(j), v)), 0)) \Rightarrow  \\ \quad \quad \operatorname{converges\text{-}to}(t, f, v) \end{array}}}
\vnbprop{converges-iff-dist-null}{\vnbstmt{\begin{array}{@{}l@{}} \forall t.; \forall f \in (\mathbb{N} \to \operatorname{pts}(t)), v \in \operatorname{pts}(t). \\ \quad \operatorname{is\text{-}metric\text{-}space}(t) \Rightarrow  \\ \quad \quad \operatorname{converges\text{-}to}(t, f, v) \Leftrightarrow \operatorname{converges\text{-}to}(\operatorname{rr\text{-}ms}, (\lambda j \in \mathbb{N}.\; \operatorname{dist}(t)(f(j), v)), 0) \end{array}}}

\vnbbookfile{theorem-library/sequential-continuity}{theorem-library/sequential-continuity}
\noindent{\small sequential-continuity.scm -- THE SEQUENTIAL CHARACTERIZATION OF CONTINUITY, both directions PROVEN:}\par
\vnbprop{continuous-at-implies-sequential}{\vnbstmt{\begin{array}{@{}l@{}} \forall s, t, f, a.; \forall q \in \operatorname{pts}(s)^{\mathbb{N}}. \\ \quad \operatorname{is\text{-}metric\text{-}space}(s) \wedge \\ \quad \operatorname{is\text{-}metric\text{-}space}(t) \wedge \\ \quad f \in (\operatorname{pts}(s) \to \operatorname{pts}(t)) \wedge \\ \quad a \in \operatorname{pts}(s) \wedge \\ \quad \operatorname{is\text{-}continuous\text{-}at}(s, t, f, a) \wedge \\ \quad \operatorname{converges\text{-}to}(s, q, a) \Rightarrow \\ \quad \quad \operatorname{converges\text{-}to}(t, \operatorname{compose}(f, q), f(a)) \end{array}}}
\vnbprop{not-continuous-witness}{\vnbstmt{\begin{array}{@{}l@{}} \forall s, t, f, a, b, l. \\ \quad \neg \forall \beta \in \operatorname{pts}(s).\; ((\operatorname{dist}(s)(a, \beta) \leq l) \Rightarrow (\operatorname{dist}(t)(f(a), f(\beta)) \leq b)) \Rightarrow  \\ \quad \quad \exists \beta \in \operatorname{pts}(s). \\ \quad \quad \quad (\operatorname{dist}(s)(a, \beta) \leq l) \wedge \neg (\operatorname{dist}(t)(f(a), f(\beta)) \leq b) \end{array}}}
\vnbprop{rr-recip-antitone}{\vnbstmt{\forall u, v \in \mathbb{R}.\; (((0 < u) \wedge (u \leq v)) \Rightarrow (\operatorname{recip}(v) \leq \operatorname{recip}(u)))}}
\vnbprop{nn-recip-succ-antitone}{\vnbstmt{\forall m, n \in \mathbb{N}.\; ((m \leq n) \Rightarrow (\operatorname{recip}((n + 1)) \leq \operatorname{recip}((m + 1))))}}
\vnbprop{sequential-implies-continuous-at}{\vnbstmt{\begin{array}{@{}l@{}} \forall s, t, f, a. \\ \quad \operatorname{is\text{-}metric\text{-}space}(s) \wedge \\ \quad \operatorname{is\text{-}metric\text{-}space}(t) \wedge \\ \quad f \in (\operatorname{pts}(s) \to \operatorname{pts}(t)) \wedge \\ \quad a \in \operatorname{pts}(s) \wedge \\ \quad \forall q \in \operatorname{pts}(s)^{\mathbb{N}}. \\ \quad \quad \operatorname{converges\text{-}to}(s, q, a) \Rightarrow \operatorname{converges\text{-}to}(t, \operatorname{compose}(f, q), f(a)) \Rightarrow \\ \quad \quad \operatorname{is\text{-}continuous\text{-}at}(s, t, f, a) \end{array}}}
\vnbprop{continuous-at-iff-sequential}{\vnbstmt{\begin{array}{@{}l@{}} \forall s, t, f, a. \\ \quad (\operatorname{is\text{-}metric\text{-}space}(s) \wedge (\operatorname{is\text{-}metric\text{-}space}(t) \wedge ((f \in (\operatorname{pts}(s) \to \operatorname{pts}(t))) \wedge (a \in \operatorname{pts}(s))))) \Rightarrow  \\ \quad \quad \operatorname{is\text{-}continuous\text{-}at}(s, t, f, a) \\ \quad \quad \Leftrightarrow\; \forall q \in \operatorname{pts}(s)^{\mathbb{N}}. \\ \quad \quad \quad \operatorname{converges\text{-}to}(s, q, a) \Rightarrow \operatorname{converges\text{-}to}(t, \operatorname{compose}(f, q), f(a)) \end{array}}}

\vnbbookfile{theorem-library/recip-succ-null}{theorem-library/recip-succ-null}
\noindent{\small recip-succ-null.scm -- 1/(n+1) -$>$ 0, PROVEN.}\par
\vnbprop{rr-recip-succ-converges-to-zero}{\vnbstmt{\operatorname{converges\text{-}to}(\operatorname{rr\text{-}ms}, (\lambda k \in \mathbb{N}.\; \operatorname{recip}((k + 1))), 0)}}

\vnbbookfile{theorem-library/product-convergence}{theorem-library/product-convergence}
\noindent{\small product-convergence.scm -- CONVERGENCE IN THE COUNTABLE PRODUCT IS EXACTLY COORDINATEWISE CONVERGENCE, proven both ways, retiring the asserted `product-convergence-coordinatewise' of structure-library/product-metric.scm -- the defining property of the PRODUCT TOPOLOGY, and the rung every later}\par
\vnbprop{rr-null-squeeze}{\vnbstmt{\begin{array}{@{}l@{}} \forall f, h \in (\mathbb{N} \to \mathbb{R}). \\ \quad (\operatorname{converges\text{-}to}(\operatorname{rr\text{-}ms}, f, 0) \wedge \forall j \in \mathbb{N}.\; ((0 \leq h(j)) \wedge (h(j) \leq f(j)))) \Rightarrow  \\ \quad \quad \operatorname{converges\text{-}to}(\operatorname{rr\text{-}ms}, h, 0) \end{array}}}
\vnbprop{converges-to-transfer}{\vnbstmt{\begin{array}{@{}l@{}} \forall t.; \forall f, h \in (\mathbb{N} \to \operatorname{pts}(t)), v \in \operatorname{pts}(t). \\ \quad (\operatorname{is\text{-}metric\text{-}space}(t) \wedge (\operatorname{converges\text{-}to}(t, f, v) \wedge \forall j \in \mathbb{N}.\; (h(j) = f(j)))) \Rightarrow  \\ \quad \quad \operatorname{converges\text{-}to}(t, h, v) \end{array}}}
\vnbprop{product-coord-weighted-null}{\vnbstmt{\begin{array}{@{}l@{}} \forall s, w.; \forall q \in (\mathbb{N} \to \operatorname{product\text{-}carrier}(s)), v \in \operatorname{product\text{-}carrier}(s), n \in \mathbb{N}. \\ \quad (\operatorname{is\text{-}ms\text{-}sequence}(s) \wedge (\operatorname{summable\text{-}weight}(w) \wedge \operatorname{converges\text{-}to}(s(n), (\lambda k \in \mathbb{N}.\; q(k)(n)), v(n)))) \Rightarrow  \\ \quad \quad \operatorname{converges\text{-}to}(\operatorname{rr\text{-}ms}, (\lambda j \in \mathbb{N}.\; (w(n) \cdot \operatorname{dist}(\operatorname{bdd\text{-}metric}(s(n)))(q(j)(n), v(n)))), 0) \end{array}}}
\vnbprop{product-coord-from-weighted-null}{\vnbstmt{\begin{array}{@{}l@{}} \forall s, w.; \forall q \in (\mathbb{N} \to \operatorname{product\text{-}carrier}(s)), v \in \operatorname{product\text{-}carrier}(s), n \in \mathbb{N}. \\ \quad \operatorname{is\text{-}ms\text{-}sequence}(s) \wedge \\ \quad \operatorname{summable\text{-}weight}(w) \wedge \\ \quad \operatorname{converges\text{-}to}(\operatorname{rr\text{-}ms}, (\lambda j \in \mathbb{N}.\; (w(n) \cdot \operatorname{dist}(\operatorname{bdd\text{-}metric}(s(n)))(q(j)(n), v(n)))), 0) \Rightarrow \\ \quad \quad \operatorname{converges\text{-}to}(s(n), (\lambda k \in \mathbb{N}.\; q(k)(n)), v(n)) \end{array}}}
\vnbprop{product-convergence-coordinatewise-bwd}{\vnbstmt{\begin{array}{@{}l@{}} \forall s, w.; \forall q \in (\mathbb{N} \to \operatorname{product\text{-}carrier}(s)), v \in \operatorname{product\text{-}carrier}(s). \\ \quad \operatorname{is\text{-}ms\text{-}sequence}(s) \wedge \\ \quad \operatorname{summable\text{-}weight}(w) \wedge \\ \quad \forall n \in \mathbb{N}.\; \operatorname{converges\text{-}to}(s(n), (\lambda k \in \mathbb{N}.\; q(k)(n)), v(n)) \Rightarrow \\ \quad \quad \operatorname{converges\text{-}to}(\operatorname{product\text{-}metric\text{-}w}(s, w), q, v) \end{array}}}
\vnbprop{product-convergence-coordinatewise-fwd}{\vnbstmt{\begin{array}{@{}l@{}} \forall s, w.; \forall q \in (\mathbb{N} \to \operatorname{product\text{-}carrier}(s)), v \in \operatorname{product\text{-}carrier}(s), n \in \mathbb{N}. \\ \quad (\operatorname{is\text{-}ms\text{-}sequence}(s) \wedge (\operatorname{summable\text{-}weight}(w) \wedge \operatorname{converges\text{-}to}(\operatorname{product\text{-}metric\text{-}w}(s, w), q, v))) \Rightarrow  \\ \quad \quad \operatorname{converges\text{-}to}(s(n), (\lambda k \in \mathbb{N}.\; q(k)(n)), v(n)) \end{array}}}
\vnbprop{product-convergence-coordinatewise}{\vnbstmt{\begin{array}{@{}l@{}} \forall s, w.; \forall q \in (\mathbb{N} \to \operatorname{product\text{-}carrier}(s)), v \in \operatorname{product\text{-}carrier}(s). \\ \quad (\operatorname{is\text{-}ms\text{-}sequence}(s) \wedge \operatorname{summable\text{-}weight}(w)) \Rightarrow  \\ \quad \quad \operatorname{converges\text{-}to}(\operatorname{product\text{-}metric\text{-}w}(s, w), q, v) \\ \quad \quad \Leftrightarrow\; \forall n \in \mathbb{N}.\; \operatorname{converges\text{-}to}(s(n), (\lambda k \in \mathbb{N}.\; q(k)(n)), v(n)) \end{array}}}

\vnbbookfile{theorem-library/rake-diagonal-subseq}{theorem-library/rake-diagonal-subseq}
\noindent{\small theorem-library/rake-diagonal-subseq.scm ==================================================================== The cross-coordinate DIAGONAL argument, and the sequential Tychonoff assembly on top of it.  Two supports of theorem-library/seq-compact-product}\par
\vnbprop{coordinatewise-diagonal-subseq}{\vnbstmt{\begin{array}{@{}l@{}} \forall s.; \forall q \in (\mathbb{N} \to \operatorname{product\text{-}carrier}(s)). \\ \quad (\operatorname{is\text{-}ms\text{-}sequence}(s) \wedge \forall n \in \mathbb{N}.\; \operatorname{seq\text{-}compact}(s(n))) \Rightarrow  \\ \quad \quad \exists a. \\ \quad \quad \quad \operatorname{strictly\text{-}mono\text{-}nn}(a) \wedge \\ \quad \quad \quad \exists l \in \operatorname{product\text{-}carrier}(s). \forall n \in \mathbb{N}. \\ \quad \quad \quad \quad \operatorname{converges\text{-}to}(s(n), (\lambda k \in \mathbb{N}.\; \operatorname{subseq}(q, a)(k)(n)), l(n)) \end{array}}}
\vnbprop{seq-compact-countable-product}{\vnbstmt{\begin{array}{@{}l@{}} \forall s. \\ \quad (\operatorname{is\text{-}ms\text{-}sequence}(s) \wedge \forall n \in \mathbb{N}.\; \operatorname{seq\text{-}compact}(s(n))) \Rightarrow  \\ \quad \quad \operatorname{seq\text{-}compact}(\operatorname{product\text{-}metric}(s)) \end{array}}}

\vnbbookfile{theorem-library/tychonoff-proof}{theorem-library/tychonoff-proof}
\noindent{\small theorem-library/tychonoff-proof.scm ==================================================================== compact-countable-product (countable Tychonoff headline) -- PROVEN.}\par
\vnbprop{compact-countable-product}{\vnbstmt{\begin{array}{@{}l@{}} \forall s. \\ \quad (\operatorname{is\text{-}ms\text{-}sequence}(s) \wedge \forall n \in \mathbb{N}.\; \operatorname{is\text{-}compact}(s(n))) \Rightarrow  \\ \quad \quad \operatorname{is\text{-}compact}(\operatorname{product\text{-}metric}(s)) \end{array}}}

\vnbbookfile{theorem-library/rake-setoid2}{theorem-library/rake-setoid2}
\noindent{\small rake-setoid2.scm -- BATCH O of the 2026-09-17 rake: the SETOID quotient and the METRIC COMPLETION.  Twelve of the thirteen asserted leaves of the batch are PROVEN here; the thirteenth is a STATEMENT DEFECT and is recorded below. Every statement is its support's statement UNCHANGED (copied from the}\par
\vnbprop{rko2-cseq-unfold}{\vnbstmt{\forall m.\; (\operatorname{cseq}(m) \simeq \{f \in (\mathbb{N} \to \operatorname{pts}(m)) : \operatorname{is\text{-}cauchy\text{-}seq}(m, f)\})}}
\vnbprop{rko2-cseq-in-fun}{\vnbstmt{\forall m.; \forall f \in \operatorname{cseq}(m).\; (f \in (\mathbb{N} \to \operatorname{pts}(m)))}}
\vnbprop{rko2-completion-dist}{\vnbstmt{\begin{array}{@{}l@{}} \forall m. \\ \quad \operatorname{is\text{-}metric\text{-}space}(m) \Rightarrow (\operatorname{dist}(\operatorname{completion}(m)) \simeq \operatorname{completion\text{-}dist}(m)) \end{array}}}
\vnbprop{rko2-cauchy-setoid-rel}{\vnbstmt{\forall m.\; (\operatorname{is\text{-}metric\text{-}space}(m) \Rightarrow (\operatorname{rel}(\operatorname{cauchy\text{-}setoid}(m)) \simeq \operatorname{crel}(m)))}}
\vnbprop{rko2-crel-unfold}{\vnbstmt{\begin{array}{@{}l@{}} \forall m. \\ \quad (\operatorname{crel}(m) \simeq \{p \in (\operatorname{cseq}(m) \times \operatorname{cseq}(m)) : \operatorname{cseq\text{-}equiv}(m, \operatorname{nth}(1, p), \operatorname{nth}(2, p))\}) \end{array}}}
\vnbprop{rko2-crel-intro}{\vnbstmt{\begin{array}{@{}l@{}} \forall m, p, q. \\ \quad ((p \in \operatorname{cseq}(m)) \wedge ((q \in \operatorname{cseq}(m)) \wedge \operatorname{cseq\text{-}equiv}(m, p, q))) \Rightarrow (\langle p, q \rangle \in \operatorname{crel}(m)) \end{array}}}
\vnbprop{rko2-crel-elim}{\vnbstmt{\forall m, p, q.\; ((\langle p, q \rangle \in \operatorname{crel}(m)) \Rightarrow \operatorname{cseq\text{-}equiv}(m, p, q))}}
\vnbprop{rko2-cauchy-setoid-length}{\vnbstmt{\forall m.\; (\operatorname{is\text{-}metric\text{-}space}(m) \Rightarrow (\operatorname{length}(\operatorname{cauchy\text{-}setoid}(m)) = 2))}}
\vnbprop{rko2-crel-is-set}{\vnbstmt{\forall m.\; (\operatorname{is\text{-}metric\text{-}space}(m) \Rightarrow (\operatorname{crel}(m) \in \mathbf{Set}))}}
\vnbprop{rko2-dist-seq-unfold}{\vnbstmt{\forall m, f, g.\; (\operatorname{dist\text{-}seq}(m, f, g) \simeq (\lambda n \in \mathbb{N}.\; \operatorname{dist}(m)(f(n), g(n))))}}
\vnbprop{rko2-dist-seq-type}{\vnbstmt{\begin{array}{@{}l@{}} \forall m, f, g. \\ \quad (\operatorname{is\text{-}metric\text{-}space}(m) \wedge ((f \in (\mathbb{N} \to \operatorname{pts}(m))) \wedge (g \in (\mathbb{N} \to \operatorname{pts}(m))))) \Rightarrow  \\ \quad \quad \operatorname{dist\text{-}seq}(m, f, g) \in (\mathbb{N} \to \mathbb{R}) \end{array}}}
\vnbprop{rko2-null-refl}{\vnbstmt{\begin{array}{@{}l@{}} \forall m, f. \\ \quad (\operatorname{is\text{-}metric\text{-}space}(m) \wedge (f \in (\mathbb{N} \to \operatorname{pts}(m)))) \Rightarrow  \\ \quad \quad \operatorname{converges\text{-}to}(\operatorname{rr\text{-}ms}, \operatorname{dist\text{-}seq}(m, f, f), 0) \end{array}}}
\vnbprop{rko2-null-sym}{\vnbstmt{\begin{array}{@{}l@{}} \forall m, f, g. \\ \quad \operatorname{is\text{-}metric\text{-}space}(m) \wedge \\ \quad f \in (\mathbb{N} \to \operatorname{pts}(m)) \wedge \\ \quad g \in (\mathbb{N} \to \operatorname{pts}(m)) \wedge \\ \quad \operatorname{converges\text{-}to}(\operatorname{rr\text{-}ms}, \operatorname{dist\text{-}seq}(m, f, g), 0) \Rightarrow \\ \quad \quad \operatorname{converges\text{-}to}(\operatorname{rr\text{-}ms}, \operatorname{dist\text{-}seq}(m, g, f), 0) \end{array}}}
\vnbprop{rko2-null-trans}{\vnbstmt{\begin{array}{@{}l@{}} \forall m, f, g, h. \\ \quad \operatorname{is\text{-}metric\text{-}space}(m) \wedge \\ \quad f \in (\mathbb{N} \to \operatorname{pts}(m)) \wedge \\ \quad g \in (\mathbb{N} \to \operatorname{pts}(m)) \wedge \\ \quad h \in (\mathbb{N} \to \operatorname{pts}(m)) \wedge \\ \quad \operatorname{converges\text{-}to}(\operatorname{rr\text{-}ms}, \operatorname{dist\text{-}seq}(m, f, g), 0) \wedge \\ \quad \operatorname{converges\text{-}to}(\operatorname{rr\text{-}ms}, \operatorname{dist\text{-}seq}(m, g, h), 0) \Rightarrow \\ \quad \quad \operatorname{converges\text{-}to}(\operatorname{rr\text{-}ms}, \operatorname{dist\text{-}seq}(m, f, h), 0) \end{array}}}
\vnbprop{cauchy-setoid-is-setoid}{\vnbstmt{\forall m.\; (\operatorname{is\text{-}metric\text{-}space}(m) \Rightarrow \operatorname{is\text{-}setoid}(\operatorname{cauchy\text{-}setoid}(m)))}}
\vnbprop{rko2-quad}{\vnbstmt{\begin{array}{@{}l@{}} \forall m, a, b, c, e. \\ \quad (\operatorname{is\text{-}metric\text{-}space}(m) \wedge ((a \in \operatorname{pts}(m)) \wedge ((b \in \operatorname{pts}(m)) \wedge ((c \in \operatorname{pts}(m)) \wedge (e \in \operatorname{pts}(m)))))) \Rightarrow  \\ \quad \quad \lvert (\operatorname{dist}(m)(a, b) - \operatorname{dist}(m)(c, e)) \rvert \leq (\operatorname{dist}(m)(a, c) + \operatorname{dist}(m)(b, e)) \end{array}}}
\vnbprop{rko2-dist-converges}{\vnbstmt{\begin{array}{@{}l@{}} \forall m, f, g, a, b. \\ \quad \operatorname{is\text{-}metric\text{-}space}(m) \wedge \\ \quad f \in (\mathbb{N} \to \operatorname{pts}(m)) \wedge \\ \quad g \in (\mathbb{N} \to \operatorname{pts}(m)) \wedge \\ \quad a \in \operatorname{pts}(m) \wedge \\ \quad b \in \operatorname{pts}(m) \wedge \\ \quad \operatorname{converges\text{-}to}(m, f, a) \wedge \\ \quad \operatorname{converges\text{-}to}(m, g, b) \Rightarrow \\ \quad \quad \operatorname{converges\text{-}to}(\operatorname{rr\text{-}ms}, \operatorname{dist\text{-}seq}(m, f, g), \operatorname{dist}(m)(a, b)) \end{array}}}
\vnbprop{rko2-const-seq-converges}{\vnbstmt{\begin{array}{@{}l@{}} \forall m, u. \\ \quad (\operatorname{is\text{-}metric\text{-}space}(m) \wedge (u \in \operatorname{pts}(m))) \Rightarrow \operatorname{converges\text{-}to}(m, \operatorname{embed\text{-}seq}(m, u), u) \end{array}}}
\vnbprop{rko2-class-eq-null}{\vnbstmt{\begin{array}{@{}l@{}} \forall m, p, q. \\ \quad \operatorname{is\text{-}metric\text{-}space}(m) \wedge \\ \quad p \in \operatorname{cseq}(m) \wedge \\ \quad q \in \operatorname{cseq}(m) \wedge \\ \quad \operatorname{class}(\operatorname{cauchy\text{-}setoid}(m), p) = \operatorname{class}(\operatorname{cauchy\text{-}setoid}(m), q) \Rightarrow \\ \quad \quad \operatorname{converges\text{-}to}(\operatorname{rr\text{-}ms}, \operatorname{dist\text{-}seq}(m, p, q), 0) \end{array}}}
\vnbprop{rko2-const-equiv-converges}{\vnbstmt{\begin{array}{@{}l@{}} \forall m, u, f. \\ \quad \operatorname{is\text{-}metric\text{-}space}(m) \wedge \\ \quad u \in \operatorname{pts}(m) \wedge \\ \quad f \in (\mathbb{N} \to \operatorname{pts}(m)) \wedge \\ \quad \operatorname{converges\text{-}to}(\operatorname{rr\text{-}ms}, \operatorname{dist\text{-}seq}(m, \operatorname{embed\text{-}seq}(m, u), f), 0) \Rightarrow \\ \quad \quad \operatorname{converges\text{-}to}(m, f, u) \end{array}}}
\vnbprop{embed-isometry}{\vnbstmt{\begin{array}{@{}l@{}} \forall m.; \forall u, v \in \operatorname{pts}(m). \\ \quad \operatorname{is\text{-}metric\text{-}space}(m) \Rightarrow  \\ \quad \quad \operatorname{dist}(\operatorname{completion}(m))(\operatorname{embed}(m)(u), \operatorname{embed}(m)(v)) = \operatorname{dist}(m)(u, v) \end{array}}}

\vnbbookfile{theorem-library/rake-completion-ms}{theorem-library/rake-completion-ms}
\noindent{\small rake-completion-ms.scm -- BATCH 5c-Q of the 2026-09-18 rake: completion-is-metric-space (structure-library/metric-completion.scm:107, asserted `well-known').  PROVEN `modulo 0'.  The statement is the support's, copied from the support site and unchanged:}\par
\vnbprop{r7q-cseq-cauchy}{\vnbstmt{\forall m.; \forall f \in \operatorname{cseq}(m).\; \operatorname{is\text{-}cauchy\text{-}seq}(m, f)}}
\vnbprop{r7q-dist-seq-cauchy}{\vnbstmt{\begin{array}{@{}l@{}} \forall m.; \forall f, g \in \operatorname{cseq}(m). \\ \quad \operatorname{is\text{-}metric\text{-}space}(m) \Rightarrow \operatorname{is\text{-}cauchy\text{-}seq}(\operatorname{rr\text{-}ms}, \operatorname{dist\text{-}seq}(m, f, g)) \end{array}}}
\vnbprop{r7q-dist-seq-converges}{\vnbstmt{\begin{array}{@{}l@{}} \forall m.; \forall f, g \in \operatorname{cseq}(m). \\ \quad \operatorname{is\text{-}metric\text{-}space}(m) \Rightarrow \operatorname{converges}(\operatorname{rr\text{-}ms}, \operatorname{dist\text{-}seq}(m, f, g)) \end{array}}}
\vnbprop{r7q-dist-seq-sym}{\vnbstmt{\begin{array}{@{}l@{}} \forall m.; \forall f, g \in (\mathbb{N} \to \operatorname{pts}(m)), v. \\ \quad (\operatorname{is\text{-}metric\text{-}space}(m) \wedge \operatorname{converges\text{-}to}(\operatorname{rr\text{-}ms}, \operatorname{dist\text{-}seq}(m, f, g), v)) \Rightarrow  \\ \quad \quad \operatorname{converges\text{-}to}(\operatorname{rr\text{-}ms}, \operatorname{dist\text{-}seq}(m, g, f), v) \end{array}}}
\vnbprop{r7q-null-transfer}{\vnbstmt{\begin{array}{@{}l@{}} \forall t.; \forall f, h \in (\mathbb{N} \to \operatorname{pts}(t)), v \in \operatorname{pts}(t). \\ \quad (\operatorname{is\text{-}metric\text{-}space}(t) \wedge (\operatorname{converges\text{-}to}(t, f, v) \wedge \operatorname{converges\text{-}to}(\operatorname{rr\text{-}ms}, \operatorname{dist\text{-}seq}(t, f, h), 0))) \Rightarrow  \\ \quad \quad \operatorname{converges\text{-}to}(t, h, v) \end{array}}}
\vnbprop{r7q-quad-null}{\vnbstmt{\begin{array}{@{}l@{}} \forall m.; \forall f, g, p, q \in (\mathbb{N} \to \operatorname{pts}(m)). \\ \quad \operatorname{is\text{-}metric\text{-}space}(m) \wedge \\ \quad \operatorname{converges\text{-}to}(\operatorname{rr\text{-}ms}, \operatorname{dist\text{-}seq}(m, f, p), 0) \wedge \\ \quad \operatorname{converges\text{-}to}(\operatorname{rr\text{-}ms}, \operatorname{dist\text{-}seq}(m, g, q), 0) \Rightarrow \\ \quad \quad \operatorname{converges\text{-}to}(\operatorname{rr\text{-}ms}, \operatorname{dist\text{-}seq}(\operatorname{rr\text{-}ms}, \operatorname{dist\text{-}seq}(m, f, g), \operatorname{dist\text{-}seq}(m, p, q)), 0) \end{array}}}
\vnbprop{r7q-limit-indep}{\vnbstmt{\begin{array}{@{}l@{}} \forall m.; \forall f, g, p, q \in (\mathbb{N} \to \operatorname{pts}(m)), v. \\ \quad \operatorname{is\text{-}metric\text{-}space}(m) \wedge \\ \quad \operatorname{converges\text{-}to}(\operatorname{rr\text{-}ms}, \operatorname{dist\text{-}seq}(m, f, p), 0) \wedge \\ \quad \operatorname{converges\text{-}to}(\operatorname{rr\text{-}ms}, \operatorname{dist\text{-}seq}(m, g, q), 0) \wedge \\ \quad \operatorname{converges\text{-}to}(\operatorname{rr\text{-}ms}, \operatorname{dist\text{-}seq}(m, f, g), v) \Rightarrow \\ \quad \quad \operatorname{converges\text{-}to}(\operatorname{rr\text{-}ms}, \operatorname{dist\text{-}seq}(m, p, q), v) \end{array}}}
\vnbprop{r7q-iota-dhat}{\vnbstmt{\begin{array}{@{}l@{}} \forall m, x, y. \\ \quad (\operatorname{is\text{-}metric\text{-}space}(m) \wedge ((x \in \operatorname{quotient}(\operatorname{cauchy\text{-}setoid}(m))) \wedge (y \in \operatorname{quotient}(\operatorname{cauchy\text{-}setoid}(m))))) \Rightarrow  \\ \quad \quad ((\iota\, l.\;\exists f, g \in \operatorname{cseq}(m).\; ((x = \operatorname{class}(\operatorname{cauchy\text{-}setoid}(m), f)) \wedge ((y = \operatorname{class}(\operatorname{cauchy\text{-}setoid}(m), g)) \wedge \operatorname{converges\text{-}to}(\operatorname{rr\text{-}ms}, \operatorname{dist\text{-}seq}(m, f, g), l)))) \simeq \operatorname{dist}(\operatorname{completion}(m))(x, y)) \end{array}}}
\vnbprop{r7q-dhat-iota}{\vnbstmt{\begin{array}{@{}l@{}} \forall m, x, y. \\ \quad (\operatorname{is\text{-}metric\text{-}space}(m) \wedge ((x \in \operatorname{quotient}(\operatorname{cauchy\text{-}setoid}(m))) \wedge (y \in \operatorname{quotient}(\operatorname{cauchy\text{-}setoid}(m))))) \Rightarrow  \\ \quad \quad (\operatorname{dist}(\operatorname{completion}(m))(x, y) \simeq (\iota\, l.\;\exists f, g \in \operatorname{cseq}(m).\; ((x = \operatorname{class}(\operatorname{cauchy\text{-}setoid}(m), f)) \wedge ((y = \operatorname{class}(\operatorname{cauchy\text{-}setoid}(m), g)) \wedge \operatorname{converges\text{-}to}(\operatorname{rr\text{-}ms}, \operatorname{dist\text{-}seq}(m, f, g), l))))) \end{array}}}
\vnbprop{r7q-dist-value}{\vnbstmt{\begin{array}{@{}l@{}} \forall m.; \forall f, g \in \operatorname{cseq}(m), v. \\ \quad (\operatorname{is\text{-}metric\text{-}space}(m) \wedge \operatorname{converges\text{-}to}(\operatorname{rr\text{-}ms}, \operatorname{dist\text{-}seq}(m, f, g), v)) \Rightarrow  \\ \quad \quad \operatorname{dist}(\operatorname{completion}(m))(\operatorname{class}(\operatorname{cauchy\text{-}setoid}(m), f), \operatorname{class}(\operatorname{cauchy\text{-}setoid}(m), g)) = v \end{array}}}
\vnbprop{r7q-rr-pts}{\vnbstmt{(\mathbb{R} \simeq \operatorname{pts}(\operatorname{rr\text{-}ms}))}}
\vnbprop{r7q-dhat-in-rr}{\vnbstmt{\begin{array}{@{}l@{}} \forall m.; \forall x, y \in \operatorname{quotient}(\operatorname{cauchy\text{-}setoid}(m)). \\ \quad \operatorname{is\text{-}metric\text{-}space}(m) \Rightarrow (\operatorname{dist}(\operatorname{completion}(m))(x, y) \in \mathbb{R}) \end{array}}}
\vnbprop{r7q-dhat-self-zero}{\vnbstmt{\begin{array}{@{}l@{}} \forall m.; \forall x \in \operatorname{quotient}(\operatorname{cauchy\text{-}setoid}(m)). \\ \quad \operatorname{is\text{-}metric\text{-}space}(m) \Rightarrow (\operatorname{dist}(\operatorname{completion}(m))(x, x) = 0) \end{array}}}
\vnbprop{r7q-dhat-nonneg}{\vnbstmt{\begin{array}{@{}l@{}} \forall m.; \forall x, y \in \operatorname{quotient}(\operatorname{cauchy\text{-}setoid}(m)). \\ \quad \operatorname{is\text{-}metric\text{-}space}(m) \Rightarrow (0 \leq \operatorname{dist}(\operatorname{completion}(m))(x, y)) \end{array}}}
\vnbprop{r7q-dhat-zero-eq}{\vnbstmt{\begin{array}{@{}l@{}} \forall m.; \forall x, y \in \operatorname{quotient}(\operatorname{cauchy\text{-}setoid}(m)). \\ \quad (\operatorname{is\text{-}metric\text{-}space}(m) \wedge (\operatorname{dist}(\operatorname{completion}(m))(x, y) = 0)) \Rightarrow (x = y) \end{array}}}
\vnbprop{r7q-dhat-sym}{\vnbstmt{\begin{array}{@{}l@{}} \forall m.; \forall x, y \in \operatorname{quotient}(\operatorname{cauchy\text{-}setoid}(m)). \\ \quad \operatorname{is\text{-}metric\text{-}space}(m) \Rightarrow  \\ \quad \quad \operatorname{dist}(\operatorname{completion}(m))(x, y) = \operatorname{dist}(\operatorname{completion}(m))(y, x) \end{array}}}
\vnbprop{r7q-dhat-triangle}{\vnbstmt{\begin{array}{@{}l@{}} \forall m.; \forall x, y, z \in \operatorname{quotient}(\operatorname{cauchy\text{-}setoid}(m)). \\ \quad \operatorname{is\text{-}metric\text{-}space}(m) \Rightarrow  \\ \quad \quad \operatorname{dist}(\operatorname{completion}(m))(x, z) \leq (\operatorname{dist}(\operatorname{completion}(m))(x, y) + \operatorname{dist}(\operatorname{completion}(m))(y, z)) \end{array}}}
\vnbprop{completion-is-metric-space}{\vnbstmt{\forall m.\; (\operatorname{is\text{-}metric\text{-}space}(m) \Rightarrow \operatorname{is\text{-}metric\text{-}space}(\operatorname{completion}(m)))}}

\vnbbookfile{theorem-library/rake-completion-complete}{theorem-library/rake-completion-complete}
\noindent{\small rake-completion-complete.scm -- BATCH 12-D (2026-09-20): `completion-is-complete' (structure-library/metric-completion.scm:110, asserted `well-known') and the general bricks its diagonal argument needs.}\par
\vnbprop{rr-limit-tail-le}{\vnbstmt{\begin{array}{@{}l@{}} \forall f, v, c. \\ \quad ((f \in (\mathbb{N} \to \mathbb{R})) \wedge ((v \in \mathbb{R}) \wedge ((c \in \mathbb{R}) \wedge (\operatorname{converges\text{-}to}(\operatorname{rr\text{-}ms}, f, v) \wedge \operatorname{pos\text{-}rr}(c))))) \Rightarrow  \\ \quad \quad \exists t \in \mathbb{N}.\; \forall k \in \mathbb{N}.\; ((t \leq k) \Rightarrow (f(k) \leq (v + c))) \end{array}}}
\vnbprop{rr-limit-le-of-tail}{\vnbstmt{\begin{array}{@{}l@{}} \forall f, v, c, t. \\ \quad f \in (\mathbb{N} \to \mathbb{R}) \wedge \\ \quad v \in \mathbb{R} \wedge \\ \quad c \in \mathbb{R} \wedge \\ \quad t \in \mathbb{N} \wedge \\ \quad \operatorname{converges\text{-}to}(\operatorname{rr\text{-}ms}, f, v) \wedge \\ \quad \forall k \in \mathbb{N}.\; ((t \leq k) \Rightarrow (f(k) \leq c)) \Rightarrow \\ \quad \quad v \leq c \end{array}}}
\vnbprop{r12d-repfam}{\vnbstmt{\begin{array}{@{}l@{}} \forall m.; \forall f \in (\mathbb{N} \to \operatorname{quotient}(\operatorname{cauchy\text{-}setoid}(m))). \\ \quad \operatorname{is\text{-}metric\text{-}space}(m) \Rightarrow  \\ \quad \quad \exists a \in (\mathbb{N} \to \operatorname{cseq}(m)). \forall p \in \mathbb{N}. \\ \quad \quad \quad f(p) = \operatorname{class}(\operatorname{cauchy\text{-}setoid}(m), a(p)) \end{array}}}
\vnbprop{r12d-rapid-indices}{\vnbstmt{\begin{array}{@{}l@{}} \forall m, f, d. \\ \quad \operatorname{is\text{-}metric\text{-}space}(m) \wedge \\ \quad f \in (\mathbb{N} \to \operatorname{cseq}(m)) \wedge \\ \quad d \in (\mathbb{N} \to \mathbb{R}) \wedge \\ \quad \forall p \in \mathbb{N}.\; \operatorname{pos\text{-}rr}(d(p)) \Rightarrow \\ \quad \quad \exists a \in (\mathbb{N} \to \mathbb{N}). \forall p \in \mathbb{N}, i, j. \\ \quad \quad \quad ((i \in \mathbb{N}) \wedge ((j \in \mathbb{N}) \wedge ((a(p) \leq i) \wedge (a(p) \leq j)))) \Rightarrow  \\ \quad \quad \quad \quad \operatorname{dist}(m)(f(p)(i), f(p)(j)) \leq d(p) \end{array}}}
\vnbprop{r12d-key}{\vnbstmt{\begin{array}{@{}l@{}} \forall m, f, a, p, q, b, c. \\ \quad \operatorname{is\text{-}metric\text{-}space}(m) \wedge \\ \quad f \in (\mathbb{N} \to \operatorname{quotient}(\operatorname{cauchy\text{-}setoid}(m))) \wedge \\ \quad a \in (\mathbb{N} \to \operatorname{cseq}(m)) \wedge \\ \quad \forall p \in \mathbb{N}.\; (f(p) = \operatorname{class}(\operatorname{cauchy\text{-}setoid}(m), a(p))) \wedge \\ \quad p \in \mathbb{N} \wedge \\ \quad q \in \mathbb{N} \wedge \\ \quad b \in \mathbb{N} \wedge \\ \quad c \in \mathbb{R} \wedge \\ \quad \operatorname{pos\text{-}rr}(c) \Rightarrow \\ \quad \quad \exists n \in \mathbb{N}. \\ \quad \quad \quad ((b \leq n) \wedge (\operatorname{dist}(m)(a(p)(n), a(q)(n)) \leq (\operatorname{dist}(\operatorname{completion}(m))(f(p), f(q)) + c))) \end{array}}}
\vnbprop{r12d-diag-at}{\vnbstmt{\forall f, m.; \forall p \in \mathbb{N}.\; ((\lambda v \in \mathbb{N}.\; f(v)(m(v)))(p) \simeq f(p)(m(p)))}}
\vnbprop{r12d-diag-in-fun}{\vnbstmt{\begin{array}{@{}l@{}} \forall m, f, a. \\ \quad (\operatorname{is\text{-}metric\text{-}space}(m) \wedge ((f \in (\mathbb{N} \to \operatorname{cseq}(m))) \wedge (a \in (\mathbb{N} \to \mathbb{N})))) \Rightarrow  \\ \quad \quad (\lambda v \in \mathbb{N}.\; f(v)(a(v))) \in (\mathbb{N} \to \operatorname{pts}(m)) \end{array}}}
\vnbprop{r12d-rad-tail}{\vnbstmt{\begin{array}{@{}l@{}} \forall d, v. \\ \quad (\operatorname{null\text{-}rr\text{-}seq}(d) \wedge \operatorname{pos\text{-}rr}(v)) \Rightarrow  \\ \quad \quad \exists r \in \mathbb{N}.\; \forall k \in \mathbb{N}.\; ((r \leq k) \Rightarrow (d(k) \leq v)) \end{array}}}
\vnbprop{r12d-diag-est}{\vnbstmt{\begin{array}{@{}l@{}} \forall m, f, a, d, b, p, q, c. \\ \quad \operatorname{is\text{-}metric\text{-}space}(m) \wedge \\ \quad f \in (\mathbb{N} \to \operatorname{quotient}(\operatorname{cauchy\text{-}setoid}(m))) \wedge \\ \quad a \in (\mathbb{N} \to \operatorname{cseq}(m)) \wedge \\ \quad d \in (\mathbb{N} \to \mathbb{R}) \wedge \\ \quad b \in (\mathbb{N} \to \mathbb{N}) \wedge \\ \quad \forall p \in \mathbb{N}.\; (f(p) = \operatorname{class}(\operatorname{cauchy\text{-}setoid}(m), a(p))) \wedge \\ \quad \forall p \in \mathbb{N}, i, j. \\ \quad \quad ((i \in \mathbb{N}) \wedge ((j \in \mathbb{N}) \wedge ((b(p) \leq i) \wedge (b(p) \leq j)))) \Rightarrow  \\ \quad \quad \quad \operatorname{dist}(m)(a(p)(i), a(p)(j)) \leq d(p) \wedge \\ \quad p \in \mathbb{N} \wedge \\ \quad q \in \mathbb{N} \wedge \\ \quad c \in \mathbb{R} \wedge \\ \quad \operatorname{pos\text{-}rr}(c) \Rightarrow \\ \quad \quad \operatorname{dist}(m)(a(p)(b(p)), a(q)(b(q))) \leq ((d(p) + d(q)) + (\operatorname{dist}(\operatorname{completion}(m))(f(p), f(q)) + c)) \end{array}}}
\vnbprop{r12d-cross-est}{\vnbstmt{\begin{array}{@{}l@{}} \forall m, f, a, d, b, p, q, c. \\ \quad \operatorname{is\text{-}metric\text{-}space}(m) \wedge \\ \quad f \in (\mathbb{N} \to \operatorname{quotient}(\operatorname{cauchy\text{-}setoid}(m))) \wedge \\ \quad a \in (\mathbb{N} \to \operatorname{cseq}(m)) \wedge \\ \quad d \in (\mathbb{N} \to \mathbb{R}) \wedge \\ \quad b \in (\mathbb{N} \to \mathbb{N}) \wedge \\ \quad \forall p \in \mathbb{N}.\; (f(p) = \operatorname{class}(\operatorname{cauchy\text{-}setoid}(m), a(p))) \wedge \\ \quad \forall p \in \mathbb{N}, i, j. \\ \quad \quad ((i \in \mathbb{N}) \wedge ((j \in \mathbb{N}) \wedge ((b(p) \leq i) \wedge (b(p) \leq j)))) \Rightarrow  \\ \quad \quad \quad \operatorname{dist}(m)(a(p)(i), a(p)(j)) \leq d(p) \wedge \\ \quad p \in \mathbb{N} \wedge \\ \quad q \in \mathbb{N} \wedge \\ \quad c \in \mathbb{R} \wedge \\ \quad \operatorname{pos\text{-}rr}(c) \wedge \\ \quad b(p) \leq q \Rightarrow \\ \quad \quad \operatorname{dist}(m)(a(q)(b(q)), a(p)(q)) \leq ((d(q) + d(p)) + (\operatorname{dist}(\operatorname{completion}(m))(f(q), f(p)) + c)) \end{array}}}
\vnbprop{r12d-diag-cauchy}{\vnbstmt{\begin{array}{@{}l@{}} \forall m, f, a, d, b. \\ \quad \operatorname{is\text{-}metric\text{-}space}(m) \wedge \\ \quad f \in (\mathbb{N} \to \operatorname{quotient}(\operatorname{cauchy\text{-}setoid}(m))) \wedge \\ \quad a \in (\mathbb{N} \to \operatorname{cseq}(m)) \wedge \\ \quad \operatorname{null\text{-}rr\text{-}seq}(d) \wedge \\ \quad b \in (\mathbb{N} \to \mathbb{N}) \wedge \\ \quad \forall p \in \mathbb{N}.\; (f(p) = \operatorname{class}(\operatorname{cauchy\text{-}setoid}(m), a(p))) \wedge \\ \quad \forall p \in \mathbb{N}, i, j. \\ \quad \quad ((i \in \mathbb{N}) \wedge ((j \in \mathbb{N}) \wedge ((b(p) \leq i) \wedge (b(p) \leq j)))) \Rightarrow  \\ \quad \quad \quad \operatorname{dist}(m)(a(p)(i), a(p)(j)) \leq d(p) \wedge \\ \quad \operatorname{is\text{-}cauchy\text{-}seq}(\operatorname{completion}(m), f) \Rightarrow \\ \quad \quad \operatorname{is\text{-}cauchy\text{-}seq}(m, (\lambda v \in \mathbb{N}.\; a(v)(b(v)))) \end{array}}}
\vnbprop{r12d-diag-in-cseq}{\vnbstmt{\begin{array}{@{}l@{}} \forall m, f, a, d, b. \\ \quad \operatorname{is\text{-}metric\text{-}space}(m) \wedge \\ \quad f \in (\mathbb{N} \to \operatorname{quotient}(\operatorname{cauchy\text{-}setoid}(m))) \wedge \\ \quad a \in (\mathbb{N} \to \operatorname{cseq}(m)) \wedge \\ \quad \operatorname{null\text{-}rr\text{-}seq}(d) \wedge \\ \quad b \in (\mathbb{N} \to \mathbb{N}) \wedge \\ \quad \forall p \in \mathbb{N}.\; (f(p) = \operatorname{class}(\operatorname{cauchy\text{-}setoid}(m), a(p))) \wedge \\ \quad \forall p \in \mathbb{N}, i, j. \\ \quad \quad ((i \in \mathbb{N}) \wedge ((j \in \mathbb{N}) \wedge ((b(p) \leq i) \wedge (b(p) \leq j)))) \Rightarrow  \\ \quad \quad \quad \operatorname{dist}(m)(a(p)(i), a(p)(j)) \leq d(p) \wedge \\ \quad \operatorname{is\text{-}cauchy\text{-}seq}(\operatorname{completion}(m), f) \Rightarrow \\ \quad \quad (\lambda v \in \mathbb{N}.\; a(v)(b(v))) \in \operatorname{cseq}(m) \end{array}}}
\vnbprop{r12d-diag-limit}{\vnbstmt{\begin{array}{@{}l@{}} \forall m, f, a, d, b. \\ \quad \operatorname{is\text{-}metric\text{-}space}(m) \wedge \\ \quad f \in (\mathbb{N} \to \operatorname{quotient}(\operatorname{cauchy\text{-}setoid}(m))) \wedge \\ \quad a \in (\mathbb{N} \to \operatorname{cseq}(m)) \wedge \\ \quad \operatorname{null\text{-}rr\text{-}seq}(d) \wedge \\ \quad b \in (\mathbb{N} \to \mathbb{N}) \wedge \\ \quad \forall p \in \mathbb{N}.\; (f(p) = \operatorname{class}(\operatorname{cauchy\text{-}setoid}(m), a(p))) \wedge \\ \quad \forall p \in \mathbb{N}, i, j. \\ \quad \quad ((i \in \mathbb{N}) \wedge ((j \in \mathbb{N}) \wedge ((b(p) \leq i) \wedge (b(p) \leq j)))) \Rightarrow  \\ \quad \quad \quad \operatorname{dist}(m)(a(p)(i), a(p)(j)) \leq d(p) \wedge \\ \quad \operatorname{is\text{-}cauchy\text{-}seq}(\operatorname{completion}(m), f) \Rightarrow \\ \quad \quad \operatorname{converges\text{-}to}(\operatorname{completion}(m), f, \operatorname{class}(\operatorname{cauchy\text{-}setoid}(m), (\lambda v \in \mathbb{N}.\; a(v)(b(v))))) \end{array}}}
\vnbprop{completion-is-complete}{\vnbstmt{\forall m.\; (\operatorname{is\text{-}metric\text{-}space}(m) \Rightarrow \operatorname{is\text{-}complete}(\operatorname{completion}(m)))}}

\vnbbookfile{theorem-library/product-weights}{theorem-library/product-weights}
\noindent{\small product-weights.scm -- ANY TWO SUMMABLE WEIGHT SEQUENCES GIVE THE SAME TOPOLOGY ON THE COUNTABLE PRODUCT, proven -- retiring the asserted `product-weights-equivalent' of structure-library/product-metric.scm, the last statement of that file about the product's topology.}\par
\vnbprop{product-identity-continuous}{\vnbstmt{\begin{array}{@{}l@{}} \forall s, w, a. \\ \quad (\operatorname{is\text{-}ms\text{-}sequence}(s) \wedge (\operatorname{summable\text{-}weight}(w) \wedge \operatorname{summable\text{-}weight}(a))) \Rightarrow  \\ \quad \quad \operatorname{is\text{-}continuous}(\operatorname{product\text{-}metric\text{-}w}(s, w), \operatorname{product\text{-}metric\text{-}w}(s, a), \\ \quad \quad \quad (\lambda x \in \operatorname{pts}(\operatorname{product\text{-}metric\text{-}w}(s, w)).\; x)) \end{array}}}
\vnbprop{product-weights-equivalent}{\vnbstmt{\begin{array}{@{}l@{}} \forall s, w, a. \\ \quad (\operatorname{is\text{-}ms\text{-}sequence}(s) \wedge (\operatorname{summable\text{-}weight}(w) \wedge \operatorname{summable\text{-}weight}(a))) \Rightarrow  \\ \quad \quad \operatorname{is\text{-}continuous}(\operatorname{product\text{-}metric\text{-}w}(s, w), \operatorname{product\text{-}metric\text{-}w}(s, a), \\ \quad \quad \quad (\lambda x \in \operatorname{pts}(\operatorname{product\text{-}metric\text{-}w}(s, w)).\; x)) \wedge \\ \quad \quad \operatorname{is\text{-}continuous}(\operatorname{product\text{-}metric\text{-}w}(s, a), \operatorname{product\text{-}metric\text{-}w}(s, w), \\ \quad \quad \quad (\lambda x \in \operatorname{pts}(\operatorname{product\text{-}metric\text{-}w}(s, a)).\; x)) \end{array}}}

\vnbbookfile{theorem-library/c-metric-summable}{theorem-library/c-metric-summable}
\noindent{\small c-metric-summable.scm -- THE SERIES DEFINING THE CANONICAL METRIC OF A COUNTABLY-METRISED SPACE CONVERGES, proven, together with the plumbing every later argument about C-METRIC needs.}\par
\vnbprop{c-metric-pseudometric-at}{\vnbstmt{\begin{array}{@{}l@{}} \forall s.; \forall k \in \mathbb{N}. \\ \quad \operatorname{is\text{-}c\text{-}metric\text{-}space}(s) \Rightarrow \operatorname{is\text{-}pseudometric}(\operatorname{dists}(s)(k), \operatorname{pts}(s)) \end{array}}}
\vnbprop{c-metric-dist-real}{\vnbstmt{\begin{array}{@{}l@{}} \forall s.; \forall k \in \mathbb{N}, u, v \in \operatorname{pts}(s). \\ \quad \operatorname{is\text{-}c\text{-}metric\text{-}space}(s) \Rightarrow (\operatorname{dists}(s)(k)(u, v) \in \mathbb{R}) \end{array}}}
\vnbprop{c-metric-term-bound}{\vnbstmt{\begin{array}{@{}l@{}} \forall s, k, u, v, a. \\ \quad (\operatorname{is\text{-}c\text{-}metric\text{-}space}(s) \wedge ((k \in \mathbb{N}) \wedge ((u \in \operatorname{pts}(s)) \wedge ((v \in \operatorname{pts}(s)) \wedge ((a \in \mathbb{R}) \wedge (0 \leq a)))))) \Rightarrow  \\ \quad \quad (a \cdot \operatorname{min}(1, \operatorname{dists}(s)(k)(u, v))) \in \mathbb{R} \wedge \\ \quad \quad 0 \leq (a \cdot \operatorname{min}(1, \operatorname{dists}(s)(k)(u, v))) \wedge \\ \quad \quad (a \cdot \operatorname{min}(1, \operatorname{dists}(s)(k)(u, v))) \leq a \end{array}}}
\vnbprop{c-metric-term-seq-in-fun}{\vnbstmt{\begin{array}{@{}l@{}} \forall s, w, u, v. \\ \quad (\operatorname{is\text{-}c\text{-}metric\text{-}space}(s) \wedge (\operatorname{summable\text{-}weight}(w) \wedge ((u \in \operatorname{pts}(s)) \wedge (v \in \operatorname{pts}(s))))) \Rightarrow  \\ \quad \quad (\lambda k \in \mathbb{N}.\; (w(k) \cdot \operatorname{min}(1, \operatorname{dists}(s)(k)(u, v)))) \in (\mathbb{N} \to \mathbb{R}) \end{array}}}
\vnbprop{c-metric-summable}{\vnbstmt{\begin{array}{@{}l@{}} \forall s, w, u, v. \\ \quad (\operatorname{is\text{-}c\text{-}metric\text{-}space}(s) \wedge (\operatorname{summable\text{-}weight}(w) \wedge ((u \in \operatorname{pts}(s)) \wedge (v \in \operatorname{pts}(s))))) \Rightarrow  \\ \quad \quad \operatorname{series\text{-}converges}((\lambda k \in \mathbb{N}.\; (w(k) \cdot \operatorname{min}(1, \operatorname{dists}(s)(k)(u, v))))) \end{array}}}
\vnbprop{c-metric-carrier}{\vnbstmt{\forall s, w.\; (\operatorname{is\text{-}c\text{-}metric\text{-}space}(s) \Rightarrow (\operatorname{pts}(\operatorname{c\text{-}metric\text{-}w}(s, w)) \simeq \operatorname{pts}(s)))}}
\vnbprop{c-metric-distance}{\vnbstmt{\begin{array}{@{}l@{}} \forall s, w, u, v. \\ \quad ((u \in \operatorname{pts}(s)) \wedge (v \in \operatorname{pts}(s))) \Rightarrow  \\ \quad \quad (\operatorname{dist}(\operatorname{c\text{-}metric\text{-}w}(s, w))(u, v) \simeq \operatorname{series\text{-}limit}((\lambda k \in \mathbb{N}.\; (w(k) \cdot \operatorname{min}(1, \operatorname{dists}(s)(k)(u, v)))))) \end{array}}}
\vnbprop{c-metric-dist-converges}{\vnbstmt{\begin{array}{@{}l@{}} \forall s, w, u, v. \\ \quad (\operatorname{is\text{-}c\text{-}metric\text{-}space}(s) \wedge (\operatorname{summable\text{-}weight}(w) \wedge ((u \in \operatorname{pts}(s)) \wedge (v \in \operatorname{pts}(s))))) \Rightarrow  \\ \quad \quad \operatorname{series\text{-}converges\text{-}to}((\lambda k \in \mathbb{N}.\; (w(k) \cdot \operatorname{min}(1, \operatorname{dists}(s)(k)(u, v)))), \\ \quad \quad \quad \operatorname{dist}(\operatorname{c\text{-}metric\text{-}w}(s, w))(u, v)) \end{array}}}
\vnbprop{c-metric-canonical-summable}{\vnbstmt{\begin{array}{@{}l@{}} \forall s, u, v. \\ \quad (\operatorname{is\text{-}c\text{-}metric\text{-}space}(s) \wedge ((u \in \operatorname{pts}(s)) \wedge (v \in \operatorname{pts}(s)))) \Rightarrow  \\ \quad \quad \operatorname{series\text{-}converges}((\lambda k \in \mathbb{N}.\; ((\lambda n \in \mathbb{N}.\; /(1, \operatorname{power}(2, (n + 1))))(k) \cdot \operatorname{min}(1, \operatorname{dists}(s)(k)(u, v))))) \end{array}}}

\vnbbookfile{theorem-library/c-metric-is-metric}{theorem-library/c-metric-is-metric}
\noindent{\small c-metric-is-metric.scm -- THE CANONICAL METRIC OF A COUNTABLY-METRISED SPACE IS A METRIC SPACE, proven.}\par
\vnbprop{cmim-dist-in-rr}{\vnbstmt{\begin{array}{@{}l@{}} \forall s, w.; \forall x, y \in \operatorname{pts}(s). \\ \quad (\operatorname{is\text{-}c\text{-}metric\text{-}space}(s) \wedge \operatorname{summable\text{-}weight}(w)) \Rightarrow  \\ \quad \quad \operatorname{dist}(\operatorname{c\text{-}metric\text{-}w}(s, w))(x, y) \in \mathbb{R} \end{array}}}
\vnbprop{cmim-term-nonneg}{\vnbstmt{\begin{array}{@{}l@{}} \forall s, w.; \forall x, y \in \operatorname{pts}(s), n \in \mathbb{N}. \\ \quad (\operatorname{is\text{-}c\text{-}metric\text{-}space}(s) \wedge \operatorname{summable\text{-}weight}(w)) \Rightarrow  \\ \quad \quad 0 \leq (\lambda k \in \mathbb{N}.\; (w(k) \cdot \operatorname{min}(1, \operatorname{dists}(s)(k)(x, y))))(n) \end{array}}}
\vnbprop{cmim-dist-nonneg}{\vnbstmt{\begin{array}{@{}l@{}} \forall s, w.; \forall x, y \in \operatorname{pts}(s). \\ \quad (\operatorname{is\text{-}c\text{-}metric\text{-}space}(s) \wedge \operatorname{summable\text{-}weight}(w)) \Rightarrow  \\ \quad \quad 0 \leq \operatorname{dist}(\operatorname{c\text{-}metric\text{-}w}(s, w))(x, y) \end{array}}}
\vnbprop{cmim-term-self-zero}{\vnbstmt{\begin{array}{@{}l@{}} \forall s, w.; \forall x \in \operatorname{pts}(s), n \in \mathbb{N}. \\ \quad (\operatorname{is\text{-}c\text{-}metric\text{-}space}(s) \wedge \operatorname{summable\text{-}weight}(w)) \Rightarrow  \\ \quad \quad (\lambda k \in \mathbb{N}.\; (w(k) \cdot \operatorname{min}(1, \operatorname{dists}(s)(k)(x, x))))(n) = 0 \end{array}}}
\vnbprop{cmim-dist-self-zero}{\vnbstmt{\begin{array}{@{}l@{}} \forall s, w.; \forall x \in \operatorname{pts}(s). \\ \quad (\operatorname{is\text{-}c\text{-}metric\text{-}space}(s) \wedge \operatorname{summable\text{-}weight}(w)) \Rightarrow  \\ \quad \quad \operatorname{dist}(\operatorname{c\text{-}metric\text{-}w}(s, w))(x, x) = 0 \end{array}}}
\vnbprop{cmim-term-zero}{\vnbstmt{\begin{array}{@{}l@{}} \forall s, w.; \forall x, y \in \operatorname{pts}(s), n \in \mathbb{N}. \\ \quad \operatorname{is\text{-}c\text{-}metric\text{-}space}(s) \wedge \\ \quad \operatorname{summable\text{-}weight}(w) \wedge \\ \quad (\lambda k \in \mathbb{N}.\; (w(k) \cdot \operatorname{min}(1, \operatorname{dists}(s)(k)(x, y))))(n) = 0 \Rightarrow \\ \quad \quad \operatorname{dists}(s)(n)(x, y) = 0 \end{array}}}
\vnbprop{cmim-dist-zero-eq}{\vnbstmt{\begin{array}{@{}l@{}} \forall s, w.; \forall x, y \in \operatorname{pts}(s). \\ \quad (\operatorname{is\text{-}c\text{-}metric\text{-}space}(s) \wedge (\operatorname{summable\text{-}weight}(w) \wedge (\operatorname{dist}(\operatorname{c\text{-}metric\text{-}w}(s, w))(x, y) = 0))) \Rightarrow  \\ \quad \quad x = y \end{array}}}
\vnbprop{cmim-term-sym}{\vnbstmt{\begin{array}{@{}l@{}} \forall s, w.; \forall x, y \in \operatorname{pts}(s), n \in \mathbb{N}. \\ \quad (\operatorname{is\text{-}c\text{-}metric\text{-}space}(s) \wedge \operatorname{summable\text{-}weight}(w)) \Rightarrow  \\ \quad \quad (\lambda k \in \mathbb{N}.\; (w(k) \cdot \operatorname{min}(1, \operatorname{dists}(s)(k)(x, y))))(n) \\ \quad \quad \leq\; (\lambda k \in \mathbb{N}.\; (w(k) \cdot \operatorname{min}(1, \operatorname{dists}(s)(k)(y, x))))(n) \end{array}}}
\vnbprop{cmim-dist-sym}{\vnbstmt{\begin{array}{@{}l@{}} \forall s, w.; \forall x, y \in \operatorname{pts}(s). \\ \quad (\operatorname{is\text{-}c\text{-}metric\text{-}space}(s) \wedge \operatorname{summable\text{-}weight}(w)) \Rightarrow  \\ \quad \quad \operatorname{dist}(\operatorname{c\text{-}metric\text{-}w}(s, w))(x, y) = \operatorname{dist}(\operatorname{c\text{-}metric\text{-}w}(s, w))(y, x) \end{array}}}
\vnbprop{cmim-sum-seq-in-fun}{\vnbstmt{\begin{array}{@{}l@{}} \forall s, w.; \forall x, y, z \in \operatorname{pts}(s). \\ \quad (\operatorname{is\text{-}c\text{-}metric\text{-}space}(s) \wedge \operatorname{summable\text{-}weight}(w)) \Rightarrow  \\ \quad \quad (\lambda k \in \mathbb{N}.\; ((w(k) \cdot \operatorname{min}(1, \operatorname{dists}(s)(k)(x, y))) + (w(k) \cdot \operatorname{min}(1, \operatorname{dists}(s)(k)(y, z))))) \\ \quad \quad \in\; (\mathbb{N} \to \mathbb{R}) \end{array}}}
\vnbprop{cmim-sum-seq-split}{\vnbstmt{\begin{array}{@{}l@{}} \forall s, w.; \forall x, y, z \in \operatorname{pts}(s), i \in \mathbb{N}. \\ \quad (\operatorname{is\text{-}c\text{-}metric\text{-}space}(s) \wedge \operatorname{summable\text{-}weight}(w)) \Rightarrow  \\ \quad \quad (\lambda k \in \mathbb{N}.\; ((w(k) \cdot \operatorname{min}(1, \operatorname{dists}(s)(k)(x, y))) + (w(k) \cdot \operatorname{min}(1, \operatorname{dists}(s)(k)(y, z)))))(i) \\ \quad \quad =\; (\lambda k \in \mathbb{N}.\; (w(k) \cdot \operatorname{min}(1, \operatorname{dists}(s)(k)(x, y))))(i) + (\lambda k \in \mathbb{N}.\; (w(k) \cdot \operatorname{min}(1, \operatorname{dists}(s)(k)(y, z))))(i) \end{array}}}
\vnbprop{cmim-term-triangle}{\vnbstmt{\begin{array}{@{}l@{}} \forall s, w.; \forall x, y, z \in \operatorname{pts}(s), n \in \mathbb{N}. \\ \quad (\operatorname{is\text{-}c\text{-}metric\text{-}space}(s) \wedge \operatorname{summable\text{-}weight}(w)) \Rightarrow  \\ \quad \quad (\lambda k \in \mathbb{N}.\; (w(k) \cdot \operatorname{min}(1, \operatorname{dists}(s)(k)(x, z))))(n) \\ \quad \quad \leq\; (\lambda k \in \mathbb{N}.\; ((w(k) \cdot \operatorname{min}(1, \operatorname{dists}(s)(k)(x, y))) + (w(k) \cdot \operatorname{min}(1, \operatorname{dists}(s)(k)(y, z)))))(n) \end{array}}}
\vnbprop{cmim-dist-triangle}{\vnbstmt{\begin{array}{@{}l@{}} \forall s, w.; \forall x, y, z \in \operatorname{pts}(s). \\ \quad (\operatorname{is\text{-}c\text{-}metric\text{-}space}(s) \wedge \operatorname{summable\text{-}weight}(w)) \Rightarrow  \\ \quad \quad \operatorname{dist}(\operatorname{c\text{-}metric\text{-}w}(s, w))(x, z) \leq (\operatorname{dist}(\operatorname{c\text{-}metric\text{-}w}(s, w))(x, y) + \operatorname{dist}(\operatorname{c\text{-}metric\text{-}w}(s, w))(y, z)) \end{array}}}
\vnbprop{c-metric-w-is-metric-space}{\vnbstmt{\begin{array}{@{}l@{}} \forall s, w. \\ \quad (\operatorname{is\text{-}c\text{-}metric\text{-}space}(s) \wedge \operatorname{summable\text{-}weight}(w)) \Rightarrow  \\ \quad \quad \operatorname{is\text{-}metric\text{-}space}(\operatorname{c\text{-}metric\text{-}w}(s, w)) \end{array}}}
\vnbprop{c-metric-is-metric-space}{\vnbstmt{\forall s.\; (\operatorname{is\text{-}c\text{-}metric\text{-}space}(s) \Rightarrow \operatorname{is\text{-}metric\text{-}space}(\operatorname{c\text{-}metric}(s)))}}

\vnbbookfile{theorem-library/ccint-creep}{theorem-library/ccint-creep}
\noindent{\small ccint-creep.scm -- the CREEPING PRINCIPLE on a closed interval, PROVEN.}\par
\vnbprop{ccint-creep}{\vnbstmt{\begin{array}{@{}l@{}} \forall a, b, g. \\ \quad a \in \mathbb{R} \wedge \\ \quad b \in \mathbb{R} \wedge \\ \quad a \leq b \wedge \\ \quad g \subseteq \mathbb{R} \wedge \\ \quad a \in g \wedge \\ \quad \operatorname{rr\text{-}upper\text{-}bound}(g, b) \wedge \\ \quad \forall t \in \operatorname{ccint}(a, b). \exists d \in \mathbb{R}. \\ \quad \quad ((0 < d) \wedge \forall y \in \operatorname{ccint}(a, b).\; ((\exists w \in g.\; ((t - d) < w) \wedge (y \leq (t + d))) \Rightarrow (y \in g))) \Rightarrow \\ \quad \quad b \in g \end{array}}}

\vnbbookfile{theorem-library/ccint-bounded}{theorem-library/ccint-bounded}
\noindent{\small ccint-bounded.scm -- a continuous function on [a,b] is BOUNDED ABOVE, PROVEN.}\par
\vnbprop{continuous-bounded-above-on-ccint}{\vnbstmt{\begin{array}{@{}l@{}} \forall f, a, b. \\ \quad f \in (\mathbb{R} \to \mathbb{R}) \wedge \\ \quad a \in \mathbb{R} \wedge \\ \quad b \in \mathbb{R} \wedge \\ \quad a \leq b \wedge \\ \quad \forall x \in \operatorname{ccint}(a, b).\; \operatorname{is\text{-}continuous\text{-}at}(\operatorname{rr\text{-}ms}, \operatorname{rr\text{-}ms}, f, x) \Rightarrow \\ \quad \quad \exists k \in \mathbb{R}.\; \forall z \in \operatorname{ccint}(a, b).\; (f(z) \leq k) \end{array}}}

\vnbbookfile{theorem-library/uniform-continuity-ccint}{theorem-library/uniform-continuity-ccint}
\noindent{\small uniform-continuity-ccint.scm -- a continuous function on a CLOSED BOUNDED INTERVAL is UNIFORMLY continuous there, PROVEN.}\par
\vnbprop{continuous-uniformly-continuous-on-ccint}{\vnbstmt{\begin{array}{@{}l@{}} \forall f, a, b, s. \\ \quad f \in (\mathbb{R} \to \mathbb{R}) \wedge \\ \quad a \in \mathbb{R} \wedge \\ \quad b \in \mathbb{R} \wedge \\ \quad a \leq b \wedge \\ \quad \operatorname{pos\text{-}rr}(s) \wedge \\ \quad \forall x \in \operatorname{ccint}(a, b).\; \operatorname{is\text{-}continuous\text{-}at}(\operatorname{rr\text{-}ms}, \operatorname{rr\text{-}ms}, f, x) \Rightarrow \\ \quad \quad \exists d \in \mathbb{R}. \\ \quad \quad \quad 0 < d \wedge \\ \quad \quad \quad \forall p, q \in \operatorname{ccint}(a, b). \\ \quad \quad \quad \quad (\lvert (p - q) \rvert \leq d) \Rightarrow (\lvert (f(p) - f(q)) \rvert \leq s) \end{array}}}

\vnbbookfile{theorem-library/evt-proof}{theorem-library/evt-proof}
\noindent{\small evt-proof.scm -- the EXTREME VALUE THEOREM (max form), PROVEN.}\par
\vnbprop{extreme-value-max}{\vnbstmt{\begin{array}{@{}l@{}} \forall f, a, b. \\ \quad f \in (\mathbb{R} \to \mathbb{R}) \wedge \\ \quad a \in \mathbb{R} \wedge \\ \quad b \in \mathbb{R} \wedge \\ \quad a \leq b \wedge \\ \quad \forall x \in \operatorname{ccint}(a, b).\; \operatorname{is\text{-}continuous\text{-}at}(\operatorname{rr\text{-}ms}, \operatorname{rr\text{-}ms}, f, x) \Rightarrow \\ \quad \quad \exists c \in \operatorname{ccint}(a, b).\; \forall x \in \operatorname{ccint}(a, b).\; (f(x) \leq f(c)) \end{array}}}

\vnbbookfile{theorem-library/heine-borel-baby}{theorem-library/heine-borel-baby}
\noindent{\small NOTE 2026-09-20 (batch 12-B).  Where this file says the tree has no metric SUBSPACE structure, that is history: SUBSPACE-MS(s, A) and RESTRICT(f, A) are defined in structure-library/metric-subspace.scm, their laws are in theorem-library/metric-subspace-laws.scm and compact-subspace.scm, and}\par
\vnbprop{empty-subset-any}{\vnbstmt{\forall b.\; (\emptyset \subseteq b)}}
\vnbprop{union-singleton-mem}{\vnbstmt{\forall y \in \mathbf{Set}, s.\; (y \in (\{y\} \cup s))}}
\vnbprop{union-right-subset}{\vnbstmt{\forall a, s.\; (s \subseteq (a \cup s))}}
\vnbprop{union-singleton-subset}{\vnbstmt{\forall b, y, s.\; (((y \in b) \wedge (s \subseteq b)) \Rightarrow ((\{y\} \cup s) \subseteq b))}}
\vnbprop{big-union-mono}{\vnbstmt{\forall a, b.\; ((a \subseteq b) \Rightarrow (\operatorname{big\text{-}union}(u, a, u) \subseteq \operatorname{big\text{-}union}(u, b, u)))}}
\vnbprop{card-union-singleton-nn}{\vnbstmt{\begin{array}{@{}l@{}} \forall s, y \in \mathbf{Set}. \\ \quad (|s| \in \mathbb{N}) \Rightarrow (|(\{y\} \cup s)| \in \mathbb{N}) \end{array}}}
\vnbprop{heine-borel-baby}{\vnbstmt{\begin{array}{@{}l@{}} \forall a, b, f. \\ \quad a \in \mathbb{R} \wedge \\ \quad b \in \mathbb{R} \wedge \\ \quad a \leq b \wedge \\ \quad \forall v \in f.\; \operatorname{is\text{-}open}(\operatorname{rr\text{-}ms}, v) \wedge \\ \quad \operatorname{ccint}(a, b) \subseteq \operatorname{big\text{-}union}(u, f, u) \Rightarrow \\ \quad \quad \exists g. \\ \quad \quad \quad ((g \subseteq f) \wedge ((g \in \mathbf{Set}) \wedge ((|g| \in \mathbb{N}) \wedge (\operatorname{ccint}(a, b) \subseteq \operatorname{big\text{-}union}(u, g, u))))) \end{array}}}

\vnbbookfile{theorem-library/evt-min-proof}{theorem-library/evt-min-proof}
\noindent{\small evt-min-proof.scm -- the EXTREME VALUE THEOREM (min form), PROVEN.}\par
\vnbprop{extreme-value-min}{\vnbstmt{\begin{array}{@{}l@{}} \forall f, a, b. \\ \quad f \in (\mathbb{R} \to \mathbb{R}) \wedge \\ \quad a \in \mathbb{R} \wedge \\ \quad b \in \mathbb{R} \wedge \\ \quad a \leq b \wedge \\ \quad \forall x \in \operatorname{ccint}(a, b).\; \operatorname{is\text{-}continuous\text{-}at}(\operatorname{rr\text{-}ms}, \operatorname{rr\text{-}ms}, f, x) \Rightarrow \\ \quad \quad \exists c \in \operatorname{ccint}(a, b).\; \forall x \in \operatorname{ccint}(a, b).\; (f(c) \leq f(x)) \end{array}}}

\vnbbookfile{theorem-library/ccint-abs-bounded}{theorem-library/ccint-abs-bounded}
\noindent{\small ccint-abs-bounded.scm -- a continuous function on a CLOSED BOUNDED INTERVAL is bounded in ABSOLUTE VALUE there, PROVEN modulo 0.}\par
\vnbprop{continuous-abs-bounded-on-ccint}{\vnbstmt{\begin{array}{@{}l@{}} \forall f, a, b. \\ \quad f \in (\mathbb{R} \to \mathbb{R}) \wedge \\ \quad a \in \mathbb{R} \wedge \\ \quad b \in \mathbb{R} \wedge \\ \quad a \leq b \wedge \\ \quad \forall x \in \operatorname{ccint}(a, b).\; \operatorname{is\text{-}continuous\text{-}at}(\operatorname{rr\text{-}ms}, \operatorname{rr\text{-}ms}, f, x) \Rightarrow \\ \quad \quad \exists m \in \mathbb{R}.\; ((0 \leq m) \wedge \forall z \in \operatorname{ccint}(a, b).\; (\lvert f(z) \rvert \leq m)) \end{array}}}

\vnbbookfile{theorem-library/bernstein-density}{theorem-library/bernstein-density}
\noindent{\small bernstein-density.scm -- THEOREM 5.2 of docs/calculus.pdf: a continuous function on [0,1] is the uniform limit of its Bernstein polynomials.}\par
\par\medskip\noindent\textbf{Definition} (\texttt{bernstein-poly}).\;$\operatorname{bernstein\text{-}poly}(f, n, x) :=$
\begin{equation*}
\operatorname{series\text{-}partial\text{-}sum}((\lambda \operatorname{l\_} \in \mathbb{N}.\; (f(/(\operatorname{l\_}, n)) \cdot \operatorname{bernstein\text{-}basis}(n, x)(\operatorname{l\_}))), (n + 1))
\end{equation*}

\vnbprop{bernstein-basis-nonneg}{\vnbstmt{\begin{array}{@{}l@{}} \forall n \in \mathbb{N}, x, k. \\ \quad ((x \in \mathbb{R}) \wedge ((0 \leq x) \wedge ((x \leq 1) \wedge (k \in \mathbb{N})))) \Rightarrow  \\ \quad \quad 0 \leq \operatorname{bernstein\text{-}basis}(n, x)(k) \end{array}}}
\vnbprop{bernstein-term-bound}{\vnbstmt{\begin{array}{@{}l@{}} \forall f, n, x, h, a, m, c, l. \\ \quad f \in (\mathbb{R} \to \mathbb{R}) \wedge \\ \quad n \in \mathbb{N} \wedge \\ \quad \neg (n = 0) \wedge \\ \quad x \in \operatorname{ccint}(0, 1) \wedge \\ \quad h \in \mathbb{R} \wedge \\ \quad 0 < h \wedge \\ \quad a \in \mathbb{R} \wedge \\ \quad 0 < a \wedge \\ \quad m \in \mathbb{R} \wedge \\ \quad 0 \leq m \wedge \\ \quad c \in \mathbb{R} \wedge \\ \quad 0 \leq c \wedge \\ \quad ((n \cdot n) \cdot ((a \cdot a) \cdot c)) = (m + m) \wedge \\ \quad \forall p, q \in \operatorname{ccint}(0, 1). \\ \quad \quad (\lvert (p - q) \rvert \leq a) \Rightarrow (\lvert (f(p) - f(q)) \rvert \leq h) \wedge \\ \quad \forall z \in \operatorname{ccint}(0, 1).\; (\lvert f(z) \rvert \leq m) \wedge \\ \quad l \in \mathbb{N} \Rightarrow \\ \quad \quad \operatorname{bernstein\text{-}basis}(n, x)(l) \cdot (f(x) - f(/(l, n))) \\ \quad \quad \leq\; \operatorname{bernstein\text{-}basis}(n, x)(l) \cdot (h + (c \cdot ((l - (n \cdot x)) \cdot (l - (n \cdot x))))) \wedge \\ \quad \quad \operatorname{bernstein\text{-}basis}(n, x)(l) \cdot (-1 \cdot (h + (c \cdot ((l - (n \cdot x)) \cdot (l - (n \cdot x)))))) \\ \quad \quad \leq\; \operatorname{bernstein\text{-}basis}(n, x)(l) \cdot (f(x) - f(/(l, n))) \end{array}}}
\vnbprop{bernstein-approx-bound}{\vnbstmt{\begin{array}{@{}l@{}} \forall f, n, x, h, a, m, c. \\ \quad f \in (\mathbb{R} \to \mathbb{R}) \wedge \\ \quad n \in \mathbb{N} \wedge \\ \quad \neg (n = 0) \wedge \\ \quad x \in \operatorname{ccint}(0, 1) \wedge \\ \quad h \in \mathbb{R} \wedge \\ \quad 0 < h \wedge \\ \quad a \in \mathbb{R} \wedge \\ \quad 0 < a \wedge \\ \quad m \in \mathbb{R} \wedge \\ \quad 0 \leq m \wedge \\ \quad c \in \mathbb{R} \wedge \\ \quad 0 \leq c \wedge \\ \quad ((n \cdot n) \cdot ((a \cdot a) \cdot c)) = (m + m) \wedge \\ \quad (c \cdot n) \leq h \wedge \\ \quad \forall p, q \in \operatorname{ccint}(0, 1). \\ \quad \quad (\lvert (p - q) \rvert \leq a) \Rightarrow (\lvert (f(p) - f(q)) \rvert \leq h) \wedge \\ \quad \forall z \in \operatorname{ccint}(0, 1).\; (\lvert f(z) \rvert \leq m) \Rightarrow \\ \quad \quad \lvert (f(x) - \operatorname{bernstein\text{-}poly}(f, n, x)) \rvert \leq (h + h) \end{array}}}
\vnbprop{bernstein-uniform-approximation}{\vnbstmt{\begin{array}{@{}l@{}} \forall f, s. \\ \quad ((f \in (\mathbb{R} \to \mathbb{R})) \wedge (\forall x \in \operatorname{ccint}(0, 1).\; \operatorname{is\text{-}continuous\text{-}at}(\operatorname{rr\text{-}ms}, \operatorname{rr\text{-}ms}, f, x) \wedge \operatorname{pos\text{-}rr}(s))) \Rightarrow  \\ \quad \quad \exists p \in \mathbb{N}. \forall n, a. \\ \quad \quad \quad ((n \in \mathbb{N}) \wedge ((p \leq n) \wedge (a \in \operatorname{ccint}(0, 1)))) \Rightarrow  \\ \quad \quad \quad \quad \lvert (f(a) - \operatorname{bernstein\text{-}poly}(f, n, a)) \rvert \leq s \end{array}}}

\vnbbookfile{theorem-library/mvt-cluster-readoffs}{theorem-library/mvt-cluster-readoffs}
\noindent{\small theorem-library/mvt-cluster-readoffs.scm}\par
\vnbprop{diff-value-real}{\vnbstmt{\forall f, a, l.\; (\operatorname{is\text{-}diff\text{-}at}(f, a, l) \Rightarrow (l \in \mathbb{R}))}}
\vnbprop{deriv-neg}{\vnbstmt{\begin{array}{@{}l@{}} \forall f, a, l. \\ \quad \operatorname{is\text{-}diff\text{-}at}(f, a, l) \Rightarrow \operatorname{is\text{-}diff\text{-}at}((\lambda z \in \mathbb{R}.\; -f(z)), a, -l) \end{array}}}
\vnbprop{max-val-strict-interior}{\vnbstmt{\begin{array}{@{}l@{}} \forall h, a, b.; \forall c \in \operatorname{ccint}(a, b). \\ \quad ((h(a) = h(b)) \wedge (h(a) < h(c))) \Rightarrow ((a < c) \wedge (c < b)) \end{array}}}
\vnbprop{min-val-strict-interior}{\vnbstmt{\begin{array}{@{}l@{}} \forall h, a, b.; \forall c \in \operatorname{ccint}(a, b). \\ \quad ((h(a) = h(b)) \wedge (h(c) < h(a))) \Rightarrow ((a < c) \wedge (c < b)) \end{array}}}

\vnbbookfile{theorem-library/deriv-difference}{theorem-library/deriv-difference}
\noindent{\small deriv-difference.scm -- the DIFFERENCE rule for the Caratheodory derivative IS-DIFF-AT: if f has derivative L at a and g has derivative M at a, then x |-$>$ f(x) - g(x) has derivative L - M at a.  PROVEN, `modulo 0'.}\par
\vnbprop{deriv-difference}{\vnbstmt{\begin{array}{@{}l@{}} \forall f, g, a, l, m. \\ \quad (\operatorname{is\text{-}diff\text{-}at}(f, a, l) \wedge \operatorname{is\text{-}diff\text{-}at}(g, a, m)) \Rightarrow  \\ \quad \quad \operatorname{is\text{-}diff\text{-}at}((\lambda x \in \mathbb{R}.\; (f(x) - g(x))), a, (l - m)) \end{array}}}

\vnbbookfile{theorem-library/mvt-aux-guarded}{theorem-library/mvt-aux-guarded}
\noindent{\small theorem-library/mvt-aux-guarded.scm -- the four MVT auxiliary-function supports, GUARDED and PROVEN.}\par
\vnbprop{mvt-aux-diff}{\vnbstmt{\begin{array}{@{}l@{}} \forall f, a, b, x, l. \\ \quad ((a \in \mathbb{R}) \wedge ((b \in \mathbb{R}) \wedge \operatorname{is\text{-}diff\text{-}at}(f, x, l))) \Rightarrow  \\ \quad \quad \operatorname{is\text{-}diff\text{-}at}((\lambda z \in \mathbb{R}.\; ((f(z) \cdot (b - a)) - (z \cdot (f(b) - f(a))))), x, ((l \cdot (b - a)) - (f(b) - f(a)))) \end{array}}}
\vnbprop{mvt-aux-cont}{\vnbstmt{\begin{array}{@{}l@{}} \forall f, a, b, x. \\ \quad ((a \in \mathbb{R}) \wedge ((b \in \mathbb{R}) \wedge \operatorname{is\text{-}continuous\text{-}at}(\operatorname{rr\text{-}ms}, \operatorname{rr\text{-}ms}, f, x))) \Rightarrow  \\ \quad \quad \operatorname{is\text{-}continuous\text{-}at}(\operatorname{rr\text{-}ms}, \operatorname{rr\text{-}ms}, (\lambda z \in \mathbb{R}.\; ((f(z) \cdot (b - a)) - (z \cdot (f(b) - f(a))))), x) \end{array}}}
\vnbprop{gmvt-aux-diff}{\vnbstmt{\begin{array}{@{}l@{}} \forall f, g, a, b, x, l, m. \\ \quad ((a \in \mathbb{R}) \wedge ((b \in \mathbb{R}) \wedge (\operatorname{is\text{-}diff\text{-}at}(f, x, l) \wedge \operatorname{is\text{-}diff\text{-}at}(g, x, m)))) \Rightarrow  \\ \quad \quad \operatorname{is\text{-}diff\text{-}at}((\lambda z \in \mathbb{R}.\; ((f(z) \cdot (g(b) - g(a))) - (g(z) \cdot (f(b) - f(a))))), \\ \quad \quad \quad x, ((l \cdot (g(b) - g(a))) - (m \cdot (f(b) - f(a))))) \end{array}}}
\vnbprop{gmvt-aux-cont}{\vnbstmt{\begin{array}{@{}l@{}} \forall f, g, a, b, x. \\ \quad a \in \mathbb{R} \wedge \\ \quad b \in \mathbb{R} \wedge \\ \quad \operatorname{is\text{-}continuous\text{-}at}(\operatorname{rr\text{-}ms}, \operatorname{rr\text{-}ms}, f, x) \wedge \\ \quad \operatorname{is\text{-}continuous\text{-}at}(\operatorname{rr\text{-}ms}, \operatorname{rr\text{-}ms}, g, x) \Rightarrow \\ \quad \quad \operatorname{is\text{-}continuous\text{-}at}(\operatorname{rr\text{-}ms}, \operatorname{rr\text{-}ms}, (\lambda z \in \mathbb{R}.\; ((f(z) \cdot (g(b) - g(a))) - (g(z) \cdot (f(b) - f(a))))), x) \end{array}}}

\vnbbookfile{theorem-library/interior-extremum-proof}{theorem-library/interior-extremum-proof}
\noindent{\small theorem-library/interior-extremum-proof.scm ==================================================================== Prop 2.10 (interior max): interior-max-deriv-zero, MACHINE-PROVEN.}\par
\vnbprop{interior-max-deriv-zero}{\vnbstmt{\begin{array}{@{}l@{}} \forall f, a, b, c, l. \\ \quad f \in (\mathbb{R} \to \mathbb{R}) \wedge \\ \quad a \in \mathbb{R} \wedge \\ \quad b \in \mathbb{R} \wedge \\ \quad c \in \mathbb{R} \wedge \\ \quad a < c \wedge \\ \quad c < b \wedge \\ \quad \forall x \in \operatorname{ccint}(a, b).\; (f(x) \leq f(c)) \wedge \\ \quad \operatorname{is\text{-}diff\text{-}at}(f, c, l) \Rightarrow \\ \quad \quad l = 0 \end{array}}}
\vnbprop{interior-min-deriv-zero}{\vnbstmt{\begin{array}{@{}l@{}} \forall f, a, b, c, l. \\ \quad f \in (\mathbb{R} \to \mathbb{R}) \wedge \\ \quad a \in \mathbb{R} \wedge \\ \quad b \in \mathbb{R} \wedge \\ \quad c \in \mathbb{R} \wedge \\ \quad a < c \wedge \\ \quad c < b \wedge \\ \quad \forall x \in \operatorname{ccint}(a, b).\; (f(c) \leq f(x)) \wedge \\ \quad \operatorname{is\text{-}diff\text{-}at}(f, c, l) \Rightarrow \\ \quad \quad l = 0 \end{array}}}

\vnbbookfile{theorem-library/rolle-proof}{theorem-library/rolle-proof}
\noindent{\small theorem-library/rolle-proof.scm -- Rolle's theorem, MACHINE-PROVEN. EVT gives argmax c / argmin d on [a,b]; if either is interior, the proven Fermat theorem (interior-max/min-deriv-zero) gives h'=0; else h is constant and a midpoint (rr-midpoint-between) is an interior max.}\par
\vnbprop{rolle}{\vnbstmt{\begin{array}{@{}l@{}} \forall h, a, b. \\ \quad h \in (\mathbb{R} \to \mathbb{R}) \wedge \\ \quad a \in \mathbb{R} \wedge \\ \quad b \in \mathbb{R} \wedge \\ \quad a < b \wedge \\ \quad \forall x \in \operatorname{ccint}(a, b).\; \operatorname{is\text{-}continuous\text{-}at}(\operatorname{rr\text{-}ms}, \operatorname{rr\text{-}ms}, h, x) \wedge \\ \quad \forall x \in \mathbb{R}.\; (((a < x) \wedge (x < b)) \Rightarrow \exists l.\; \operatorname{is\text{-}diff\text{-}at}(h, x, l)) \wedge \\ \quad h(a) = h(b) \Rightarrow \\ \quad \quad \exists c \in \mathbb{R}.\; ((a < c) \wedge ((c < b) \wedge \operatorname{is\text{-}diff\text{-}at}(h, c, 0))) \end{array}}}

\vnbbookfile{theorem-library/mvt-proof}{theorem-library/mvt-proof}
\noindent{\small theorem-library/mvt-proof.scm -- the Mean Value Theorem, MACHINE-PROVEN. Strategy: apply Rolle's theorem to the auxiliary h(z) = f(z)*(b-a) - z*(f(b)-f(a)) which satisfies h(a)=h(b); Rolle gives an interior theta with h'(theta)=0,}\par
\vnbprop{mvt}{\vnbstmt{\begin{array}{@{}l@{}} \forall f, a, b. \\ \quad f \in (\mathbb{R} \to \mathbb{R}) \wedge \\ \quad a \in \mathbb{R} \wedge \\ \quad b \in \mathbb{R} \wedge \\ \quad a < b \wedge \\ \quad \forall x \in \operatorname{ccint}(a, b).\; \operatorname{is\text{-}continuous\text{-}at}(\operatorname{rr\text{-}ms}, \operatorname{rr\text{-}ms}, f, x) \wedge \\ \quad \forall x \in \mathbb{R}.\; (((a < x) \wedge (x < b)) \Rightarrow \exists l.\; \operatorname{is\text{-}diff\text{-}at}(f, x, l)) \Rightarrow \\ \quad \quad \exists c \in \mathbb{R}. \\ \quad \quad \quad ((a < c) \wedge ((c < b) \wedge \exists l.\; (\operatorname{is\text{-}diff\text{-}at}(f, c, l) \wedge ((l \cdot (b - a)) = (f(b) - f(a)))))) \end{array}}}

\vnbbookfile{theorem-library/deriv-constant-proof}{theorem-library/deriv-constant-proof}
\noindent{\small theorem-library/deriv-constant-proof.scm -- Cor 2.15, MACHINE-PROVEN. f continuous on [a,b], f'=0 on (a,b)  =$>$  f constant on [a,b]. Strategy: for u$<$v in [a,b], the proven MVT on [u,v] gives an interior theta with f'(theta)(v-u)=f(v)-f(u); the zero-derivative hypothesis + derivative-}\par
\vnbprop{deriv-zero-const-up}{\vnbstmt{\begin{array}{@{}l@{}} \forall f, a, b, u, v. \\ \quad f \in (\mathbb{R} \to \mathbb{R}) \wedge \\ \quad a \in \mathbb{R} \wedge \\ \quad b \in \mathbb{R} \wedge \\ \quad a < b \wedge \\ \quad \forall x \in \operatorname{ccint}(a, b).\; \operatorname{is\text{-}continuous\text{-}at}(\operatorname{rr\text{-}ms}, \operatorname{rr\text{-}ms}, f, x) \wedge \\ \quad \forall x \in \mathbb{R}.\; (((a < x) \wedge (x < b)) \Rightarrow \operatorname{is\text{-}diff\text{-}at}(f, x, 0)) \wedge \\ \quad u \in \operatorname{ccint}(a, b) \wedge \\ \quad v \in \operatorname{ccint}(a, b) \wedge \\ \quad u < v \Rightarrow \\ \quad \quad f(u) = f(v) \end{array}}}
\vnbprop{deriv-zero-implies-constant}{\vnbstmt{\begin{array}{@{}l@{}} \forall f, a, b, u, v. \\ \quad f \in (\mathbb{R} \to \mathbb{R}) \wedge \\ \quad a \in \mathbb{R} \wedge \\ \quad b \in \mathbb{R} \wedge \\ \quad a < b \wedge \\ \quad \forall x \in \operatorname{ccint}(a, b).\; \operatorname{is\text{-}continuous\text{-}at}(\operatorname{rr\text{-}ms}, \operatorname{rr\text{-}ms}, f, x) \wedge \\ \quad \forall x \in \mathbb{R}.\; (((a < x) \wedge (x < b)) \Rightarrow \operatorname{is\text{-}diff\text{-}at}(f, x, 0)) \wedge \\ \quad u \in \operatorname{ccint}(a, b) \wedge \\ \quad v \in \operatorname{ccint}(a, b) \Rightarrow \\ \quad \quad f(u) = f(v) \end{array}}}

\vnbbookfile{theorem-library/mvt-bounds-proof}{theorem-library/mvt-bounds-proof}
\noindent{\small theorem-library/mvt-bounds-proof.scm -- Cor 2.14, MACHINE-PROVEN. f' $<$= M on (a,b)  =$>$  f(b)-f(a) $<$= M(b-a)     (mvt-upper-bound) m $<$= f' on (a,b)  =$>$  m(b-a) $<$= f(b)-f(a)     (mvt-lower-bound) Strategy: MVT gives theta with f(b)-f(a) = f'(theta)(b-a) = L(b-a); the}\par
\vnbprop{mvt-upper-bound}{\vnbstmt{\begin{array}{@{}l@{}} \forall f, a, b, m. \\ \quad f \in (\mathbb{R} \to \mathbb{R}) \wedge \\ \quad a \in \mathbb{R} \wedge \\ \quad b \in \mathbb{R} \wedge \\ \quad m \in \mathbb{R} \wedge \\ \quad a < b \wedge \\ \quad \forall x \in \operatorname{ccint}(a, b).\; \operatorname{is\text{-}continuous\text{-}at}(\operatorname{rr\text{-}ms}, \operatorname{rr\text{-}ms}, f, x) \wedge \\ \quad \forall x \in \mathbb{R}. \\ \quad \quad ((a < x) \wedge (x < b)) \Rightarrow \exists l.\; (\operatorname{is\text{-}diff\text{-}at}(f, x, l) \wedge (l \leq m)) \Rightarrow \\ \quad \quad (f(b) - f(a)) \leq (m \cdot (b - a)) \end{array}}}
\vnbprop{mvt-lower-bound}{\vnbstmt{\begin{array}{@{}l@{}} \forall f, a, b, m. \\ \quad f \in (\mathbb{R} \to \mathbb{R}) \wedge \\ \quad a \in \mathbb{R} \wedge \\ \quad b \in \mathbb{R} \wedge \\ \quad m \in \mathbb{R} \wedge \\ \quad a < b \wedge \\ \quad \forall x \in \operatorname{ccint}(a, b).\; \operatorname{is\text{-}continuous\text{-}at}(\operatorname{rr\text{-}ms}, \operatorname{rr\text{-}ms}, f, x) \wedge \\ \quad \forall x \in \mathbb{R}. \\ \quad \quad ((a < x) \wedge (x < b)) \Rightarrow \exists l.\; (\operatorname{is\text{-}diff\text{-}at}(f, x, l) \wedge (m \leq l)) \Rightarrow \\ \quad \quad (m \cdot (b - a)) \leq (f(b) - f(a)) \end{array}}}

\vnbbookfile{theorem-library/mvt-abs-bound}{theorem-library/mvt-abs-bound}
\noindent{\small mvt-abs-bound.scm -- the TWO-SIDED corollary of the Mean Value Theorem:}\par
\vnbprop{mvt-abs-bound}{\vnbstmt{\begin{array}{@{}l@{}} \forall f, a, b, m. \\ \quad f \in (\mathbb{R} \to \mathbb{R}) \wedge \\ \quad a \in \mathbb{R} \wedge \\ \quad b \in \mathbb{R} \wedge \\ \quad m \in \mathbb{R} \wedge \\ \quad a < b \wedge \\ \quad \forall x \in \operatorname{ccint}(a, b).\; \operatorname{is\text{-}continuous\text{-}at}(\operatorname{rr\text{-}ms}, \operatorname{rr\text{-}ms}, f, x) \wedge \\ \quad \forall x \in \mathbb{R}. \\ \quad \quad ((a < x) \wedge (x < b)) \Rightarrow \exists l \in \mathbb{R}.\; (\operatorname{is\text{-}diff\text{-}at}(f, x, l) \wedge (\lvert l \rvert \leq m)) \Rightarrow \\ \quad \quad \lvert (f(b) - f(a)) \rvert \leq (m \cdot (b - a)) \end{array}}}

\vnbbookfile{theorem-library/continuity-local}{theorem-library/continuity-local}
\noindent{\small continuity-local.scm -- CONTINUITY AT A POINT READS ONLY A NEIGHBOURHOOD OF IT: a map agreeing with a continuous map on a BALL about the point, and saying nothing whatever elsewhere, is continuous at that point.  PROVEN.}\par
\vnbprop{continuous-at-local}{\vnbstmt{\begin{array}{@{}l@{}} \forall f, g \in (\mathbb{R} \to \mathbb{R}), t, o \in \mathbb{R}. \\ \quad 0 < o \wedge \\ \quad \operatorname{is\text{-}continuous\text{-}at}(\operatorname{rr\text{-}ms}, \operatorname{rr\text{-}ms}, g, t) \wedge \\ \quad \forall x \in \mathbb{R}.\; ((\lvert (x - t) \rvert \leq o) \Rightarrow (f(x) = g(x))) \Rightarrow \\ \quad \quad \operatorname{is\text{-}continuous\text{-}at}(\operatorname{rr\text{-}ms}, \operatorname{rr\text{-}ms}, f, t) \end{array}}}

\vnbbookfile{theorem-library/recip-star}{theorem-library/recip-star}
\noindent{\small recip-star.scm -- THE TOTALISED RECIPROCAL, and the reason it exists.}\par
\par\medskip\noindent\textbf{Definition} (\texttt{recip-star}).\;${\operatorname{u\_}}^{\ast\!-1} :=$ $\left(\text{if }(\operatorname{u\_} = 0)\text{ then }0\text{ else }\operatorname{recip}(\operatorname{u\_})\right)$
\vnbprop{recip-star-in-rr}{\vnbstmt{\forall u \in \mathbb{R}.\; ({u}^{\ast\!-1} \in \mathbb{R})}}
\vnbprop{recip-star-lam-in-fun}{\vnbstmt{(\lambda z \in \mathbb{R}.\; {z}^{\ast\!-1}) \in (\mathbb{R} \to \mathbb{R})}}
\vnbprop{recip-star-value}{\vnbstmt{\forall u \in \mathbb{R}.\; (\neg (u = 0) \Rightarrow ({u}^{\ast\!-1} = \operatorname{recip}(u)))}}
\vnbprop{recip-star-zero}{\vnbstmt{{0}^{\ast\!-1} = 0}}
\vnbprop{recip-star-inverse}{\vnbstmt{\forall u \in \mathbb{R}.\; (\neg (u = 0) \Rightarrow ((u \cdot {u}^{\ast\!-1}) = 1))}}
\vnbprop{rr-half-ball-nonzero}{\vnbstmt{\begin{array}{@{}l@{}} \forall c, m, y \in \mathbb{R}. \\ \quad ((0 < m) \wedge (((m + m) = \lvert c \rvert) \wedge (\lvert (y - c) \rvert \leq m))) \Rightarrow \neg (y = 0) \end{array}}}
\vnbprop{nonzero-ident-lam-in-fun}{\vnbstmt{(\lambda w \in \mathbb{R}.\; \left(\text{if }(w = 0)\text{ then }1\text{ else }w\right)) \in (\mathbb{R} \to \mathbb{R})}}
\vnbprop{nonzero-ident-nowhere-zero}{\vnbstmt{\forall x \in \mathbb{R}.\; \neg ((\lambda w \in \mathbb{R}.\; \left(\text{if }(w = 0)\text{ then }1\text{ else }w\right))(x) = 0)}}
\vnbprop{nonzero-ident-continuous-at}{\vnbstmt{\begin{array}{@{}l@{}} \forall c \in \mathbb{R}. \\ \quad \neg (c = 0) \Rightarrow  \\ \quad \quad \operatorname{is\text{-}continuous\text{-}at}(\operatorname{rr\text{-}ms}, \operatorname{rr\text{-}ms}, (\lambda w \in \mathbb{R}.\; \left(\text{if }(w = 0)\text{ then }1\text{ else }w\right)), c) \end{array}}}
\vnbprop{recip-star-continuous-at}{\vnbstmt{\begin{array}{@{}l@{}} \forall c \in \mathbb{R}. \\ \quad \neg (c = 0) \Rightarrow  \\ \quad \quad \operatorname{is\text{-}continuous\text{-}at}(\operatorname{rr\text{-}ms}, \operatorname{rr\text{-}ms}, (\lambda z \in \mathbb{R}.\; {z}^{\ast\!-1}), c) \end{array}}}
\vnbprop{recip-star-continuous-on-ccint}{\vnbstmt{\begin{array}{@{}l@{}} \forall a, b, x \in \mathbb{R}. \\ \quad ((0 < a) \wedge (x \in \operatorname{ccint}(a, b))) \Rightarrow  \\ \quad \quad \operatorname{is\text{-}continuous\text{-}at}(\operatorname{rr\text{-}ms}, \operatorname{rr\text{-}ms}, (\lambda z \in \mathbb{R}.\; {z}^{\ast\!-1}), x) \end{array}}}

\vnbbookfile{theorem-library/clamp}{theorem-library/clamp}
\noindent{\small clamp.scm -- CLAMP(a,b,x) = min(max(x,a),b), the retraction of the line onto the interval [a,b], and the ONE mechanism that dissolves the endpoint problem of docs/calculus.pdf Prop 4.16.}\par
\par\medskip\noindent\textbf{Definition} (\texttt{clamp}).\;$\operatorname{clamp}(a, b, x) :=$ $\operatorname{min}(\operatorname{max}(x, a), b)$
\vnbprop{clamp-in-rr}{\vnbstmt{\forall a, b, x \in \mathbb{R}.\; (\operatorname{clamp}(a, b, x) \in \mathbb{R})}}
\vnbprop{clamp-in-ccint}{\vnbstmt{\begin{array}{@{}l@{}} \forall a, b, x \in \mathbb{R}. \\ \quad (a \leq b) \Rightarrow (\operatorname{clamp}(a, b, x) \in \operatorname{ccint}(a, b)) \end{array}}}
\vnbprop{clamp-fixes}{\vnbstmt{\begin{array}{@{}l@{}} \forall a, b, x \in \mathbb{R}. \\ \quad ((a \leq x) \wedge (x \leq b)) \Rightarrow (\operatorname{clamp}(a, b, x) = x) \end{array}}}
\vnbprop{rr-max-lipschitz}{\vnbstmt{\forall u, v, c \in \mathbb{R}.\; (\lvert (\operatorname{max}(u, c) - \operatorname{max}(v, c)) \rvert \leq \lvert (u - v) \rvert)}}
\vnbprop{rr-min-lipschitz}{\vnbstmt{\forall u, v, c \in \mathbb{R}.\; (\lvert (\operatorname{min}(u, c) - \operatorname{min}(v, c)) \rvert \leq \lvert (u - v) \rvert)}}
\vnbprop{clamp-lipschitz}{\vnbstmt{\begin{array}{@{}l@{}} \forall a, b, y, x \in \mathbb{R}. \\ \quad \lvert (\operatorname{clamp}(a, b, y) - \operatorname{clamp}(a, b, x)) \rvert \leq \lvert (y - x) \rvert \end{array}}}
\vnbprop{clamp-compose-lam-in-fun}{\vnbstmt{\begin{array}{@{}l@{}} \forall f \in (\mathbb{R} \to \mathbb{R}), a, b \in \mathbb{R}. \\ \quad (\lambda z \in \mathbb{R}.\; f(\operatorname{clamp}(a, b, z))) \in (\mathbb{R} \to \mathbb{R}) \end{array}}}
\vnbprop{clamp-compose-fixes}{\vnbstmt{\begin{array}{@{}l@{}} \forall f \in (\mathbb{R} \to \mathbb{R}), a, b, x \in \mathbb{R}. \\ \quad ((a \leq x) \wedge (x \leq b)) \Rightarrow ((\lambda z \in \mathbb{R}.\; f(\operatorname{clamp}(a, b, z)))(x) = f(x)) \end{array}}}
\vnbprop{clamp-compose-continuous-at}{\vnbstmt{\begin{array}{@{}l@{}} \forall f \in (\mathbb{R} \to \mathbb{R}), a, b, t \in \mathbb{R}. \\ \quad ((a \leq b) \wedge \forall x \in \mathbb{R}.\; ((x \in \operatorname{ccint}(a, b)) \Rightarrow \operatorname{is\text{-}continuous\text{-}at}(\operatorname{rr\text{-}ms}, \operatorname{rr\text{-}ms}, f, x))) \Rightarrow  \\ \quad \quad \operatorname{is\text{-}continuous\text{-}at}(\operatorname{rr\text{-}ms}, \operatorname{rr\text{-}ms}, (\lambda z \in \mathbb{R}.\; f(\operatorname{clamp}(a, b, z))), t) \end{array}}}

\vnbbookfile{theorem-library/seq-limit}{theorem-library/seq-limit}
\noindent{\small theorem-library/seq-limit.scm -- THE LIMIT OF A REAL SEQUENCE AS A TERM, and the two pieces of limit arithmetic the tree was missing.}\par
\vnbprop{rr-limit-scale}{\vnbstmt{\begin{array}{@{}l@{}} \forall c \in \mathbb{R}, f, h \in (\mathbb{N} \to \mathbb{R}), v \in \mathbb{R}. \\ \quad (\forall j \in \mathbb{N}.\; (h(j) = (f(j) \cdot c)) \wedge \operatorname{converges\text{-}to}(\operatorname{rr\text{-}ms}, f, v)) \Rightarrow  \\ \quad \quad \operatorname{converges\text{-}to}(\operatorname{rr\text{-}ms}, h, (v \cdot c)) \end{array}}}
\vnbprop{rr-limit-sub}{\vnbstmt{\begin{array}{@{}l@{}} \forall f, g, h \in (\mathbb{N} \to \mathbb{R}), v, a \in \mathbb{R}. \\ \quad (\forall j \in \mathbb{N}.\; (h(j) = (f(j) - g(j))) \wedge (\operatorname{converges\text{-}to}(\operatorname{rr\text{-}ms}, f, v) \wedge \operatorname{converges\text{-}to}(\operatorname{rr\text{-}ms}, g, a))) \Rightarrow  \\ \quad \quad \operatorname{converges\text{-}to}(\operatorname{rr\text{-}ms}, h, (v - a)) \end{array}}}

\vnbbookfile{structure-library/regulated}{structure-library/regulated}
\noindent{\small MOVED 2026-09-20 (batch 12-A) from theorem-library/: this file is VOCABULARY -- definitions, notation and warranted supports, not one proof -- and every theorem stated with it had to load below it.  Its load.scm slot is unchanged. regulated.scm -- docs/calculus.pdf CHAPTER 4 SECTION 1: the vocabulary of}\par
\par\medskip\noindent\textbf{Definition} (\texttt{is-right-limit}).\;For all $f$, $x$ and $l$, 
\begin{equation*}
\begin{array}{@{}l@{}} \operatorname{is\text{-}right\text{-}limit}(f, x, l) \\ \Leftrightarrow\; f \in (\mathbb{R} \to \mathbb{R}) \wedge \\ x \in \mathbb{R} \wedge \\ l \in \mathbb{R} \wedge \\ \forall \operatorname{eps\_}. \\ \quad \operatorname{pos\text{-}rr}(\operatorname{eps\_}) \Rightarrow  \\ \quad \quad \exists \operatorname{delta\_}. \\ \quad \quad \quad \operatorname{pos\text{-}rr}(\operatorname{delta\_}) \wedge \\ \quad \quad \quad \forall \operatorname{t\_}. \\ \quad \quad \quad \quad ((\operatorname{t\_} \in \mathbb{R}) \wedge ((x < \operatorname{t\_}) \wedge (\operatorname{t\_} \leq +(x, \operatorname{delta\_})))) \Rightarrow (\lvert -(f(\operatorname{t\_}), l) \rvert \leq \operatorname{eps\_}) \end{array}
\end{equation*}

\par\medskip\noindent\textbf{Definition} (\texttt{is-regulated}).\;For all $f$, $a$ and $b$, 
\begin{equation*}
\begin{array}{@{}l@{}} \operatorname{is\text{-}regulated}(f, a, b) \\ \Leftrightarrow\; f \in (\mathbb{R} \to \mathbb{R}) \wedge \\ a \in \mathbb{R} \wedge \\ b \in \mathbb{R} \wedge \\ a < b \wedge \\ \forall \operatorname{x\_} \in \operatorname{ccint}(a, b). \\ \quad (\operatorname{x\_} < b) \Rightarrow \exists \operatorname{l\_}.\; \operatorname{is\text{-}right\text{-}limit}(f, \operatorname{x\_}, \operatorname{l\_}) \wedge \\ \quad (a < \operatorname{x\_}) \Rightarrow \exists \operatorname{l\_}.\; \operatorname{is\text{-}left\text{-}limit}(f, \operatorname{x\_}, \operatorname{l\_}) \end{array}
\end{equation*}

\par\medskip\noindent\textbf{Definition} (\texttt{is-unif-limit-on}).\;For all $\mathit{fam}$, $g$, $a$ and $b$, 
\begin{equation*}
\begin{array}{@{}l@{}} \operatorname{is\text{-}unif\text{-}limit\text{-}on}(\mathit{fam}, g, a, b) \\ \Leftrightarrow\; \mathit{fam} \in (\mathbb{N} \to (\mathbb{R} \to \mathbb{R})) \wedge \\ g \in (\mathbb{R} \to \mathbb{R}) \wedge \\ a \in \mathbb{R} \wedge \\ b \in \mathbb{R} \wedge \\ a < b \wedge \\ \forall \operatorname{eps\_}. \\ \quad \operatorname{pos\text{-}rr}(\operatorname{eps\_}) \Rightarrow  \\ \quad \quad \exists \operatorname{n\_} \in \mathbb{N}. \forall \operatorname{k\_}. \\ \quad \quad \quad ((\operatorname{k\_} \in \mathbb{N}) \wedge (\operatorname{n\_} \leq \operatorname{k\_})) \Rightarrow  \\ \quad \quad \quad \quad \forall \operatorname{x\_} \in \operatorname{ccint}(a, b).\; (\lvert -(g(\operatorname{x\_}), \operatorname{fam}(\operatorname{k\_})(\operatorname{x\_})) \rvert \leq \operatorname{eps\_}) \end{array}
\end{equation*}

\par\medskip\noindent\textbf{Definition} (\texttt{is-left-limit}).\;For all $f$, $x$ and $l$, 
\begin{equation*}
\begin{array}{@{}l@{}} \operatorname{is\text{-}left\text{-}limit}(f, x, l) \\ \Leftrightarrow\; f \in (\mathbb{R} \to \mathbb{R}) \wedge \\ x \in \mathbb{R} \wedge \\ l \in \mathbb{R} \wedge \\ \forall \operatorname{eps\_}. \\ \quad \operatorname{pos\text{-}rr}(\operatorname{eps\_}) \Rightarrow  \\ \quad \quad \exists \operatorname{delta\_}. \\ \quad \quad \quad \operatorname{pos\text{-}rr}(\operatorname{delta\_}) \wedge \\ \quad \quad \quad \forall \operatorname{t\_}. \\ \quad \quad \quad \quad ((\operatorname{t\_} \in \mathbb{R}) \wedge ((-(x, \operatorname{delta\_}) \leq \operatorname{t\_}) \wedge (\operatorname{t\_} < x))) \Rightarrow (\lvert -(f(\operatorname{t\_}), l) \rvert \leq \operatorname{eps\_}) \end{array}
\end{equation*}

\par\medskip\noindent\textbf{Definition} (\texttt{is-partition}).\;For all $p$, $n$, $a$ and $b$, 
\begin{equation*}
\begin{array}{@{}l@{}} \operatorname{is\text{-}partition}(p, n, a, b) \\ \Leftrightarrow\; n \in \mathbb{N} \wedge \\ 1 \leq n \wedge \\ a \in \mathbb{R} \wedge \\ b \in \mathbb{R} \wedge \\ a < b \wedge \\ p \in (\mathbb{N} \to \mathbb{R}) \wedge \\ p(0) = a \wedge \\ p(n) = b \wedge \\ \forall \operatorname{i\_}, \operatorname{j\_}. \\ \quad ((\operatorname{i\_} \in \operatorname{interval}(0, n)) \wedge ((\operatorname{j\_} \in \operatorname{interval}(0, n)) \wedge (\operatorname{i\_} < \operatorname{j\_}))) \Rightarrow (p(\operatorname{i\_}) < p(\operatorname{j\_})) \end{array}
\end{equation*}

\par\medskip\noindent\textbf{Definition} (\texttt{is-step-fn}).\;For all $f$, $a$ and $b$, 
\begin{equation*}
\begin{array}{@{}l@{}} \operatorname{is\text{-}step\text{-}fn}(f, a, b) \\ \Leftrightarrow\; f \in (\mathbb{R} \to \mathbb{R}) \wedge \\ a \in \mathbb{R} \wedge \\ b \in \mathbb{R} \wedge \\ a < b \wedge \\ \exists \operatorname{n\_}, \operatorname{p\_}. \\ \quad \operatorname{is\text{-}partition}(\operatorname{p\_}, \operatorname{n\_}, a, b) \wedge \\ \quad \forall \operatorname{i\_}. \\ \quad \quad ((\operatorname{i\_} \in \operatorname{interval}(0, \operatorname{n\_})) \wedge (\operatorname{i\_} < \operatorname{n\_})) \Rightarrow  \\ \quad \quad \quad \exists \operatorname{c\_} \in \mathbb{R}. \forall \operatorname{t\_}. \\ \quad \quad \quad \quad ((\operatorname{t\_} \in \mathbb{R}) \wedge ((\operatorname{p\_}(\operatorname{i\_}) < \operatorname{t\_}) \wedge (\operatorname{t\_} < \operatorname{p\_}(\operatorname{succ}(\operatorname{i\_}))))) \Rightarrow (f(\operatorname{t\_}) = \operatorname{c\_}) \end{array}
\end{equation*}

\par\medskip\noindent\textbf{Definition} (\texttt{is-piecewise-continuous}).\;For all $f$, $a$ and $b$, 
\begin{equation*}
\begin{array}{@{}l@{}} \operatorname{is\text{-}piecewise\text{-}continuous}(f, a, b) \\ \Leftrightarrow\; f \in (\mathbb{R} \to \mathbb{R}) \wedge \\ a \in \mathbb{R} \wedge \\ b \in \mathbb{R} \wedge \\ a < b \wedge \\ \exists \operatorname{n\_}, \operatorname{p\_}. \\ \quad \operatorname{is\text{-}partition}(\operatorname{p\_}, \operatorname{n\_}, a, b) \wedge \\ \quad \forall \operatorname{i\_}. \\ \quad \quad ((\operatorname{i\_} \in \operatorname{interval}(0, \operatorname{n\_})) \wedge (\operatorname{i\_} < \operatorname{n\_})) \Rightarrow  \\ \quad \quad \quad \forall \operatorname{t\_}. \\ \quad \quad \quad \quad ((\operatorname{t\_} \in \mathbb{R}) \wedge ((\operatorname{p\_}(\operatorname{i\_}) < \operatorname{t\_}) \wedge (\operatorname{t\_} < \operatorname{p\_}(\operatorname{succ}(\operatorname{i\_}))))) \Rightarrow  \\ \quad \quad \quad \quad \quad \operatorname{is\text{-}continuous\text{-}at}(\operatorname{rr\text{-}ms}, \operatorname{rr\text{-}ms}, f, \operatorname{t\_}) \end{array}
\end{equation*}


\vnbbookfile{theorem-library/onesided-limits}{theorem-library/onesided-limits}
\noindent{\small onesided-limits.scm -- docs/calculus.pdf Ch 4 Sec 1. Continuity at x gives BOTH one-sided limits at x, with value f(x).  Both `modulo 0'.  The two drivers are the same text with IS-RIGHT-LIMIT swapped for IS-LEFT-LIMIT: the interval side is the only difference, and `ineq!'}\par
\vnbprop{continuous-at-right-limit}{\vnbstmt{\begin{array}{@{}l@{}} \forall f, x. \\ \quad ((f \in (\mathbb{R} \to \mathbb{R})) \wedge ((x \in \mathbb{R}) \wedge \operatorname{is\text{-}continuous\text{-}at}(\operatorname{rr\text{-}ms}, \operatorname{rr\text{-}ms}, f, x))) \Rightarrow  \\ \quad \quad \operatorname{is\text{-}right\text{-}limit}(f, x, f(x)) \end{array}}}
\vnbprop{continuous-at-left-limit}{\vnbstmt{\begin{array}{@{}l@{}} \forall f, x. \\ \quad ((f \in (\mathbb{R} \to \mathbb{R})) \wedge ((x \in \mathbb{R}) \wedge \operatorname{is\text{-}continuous\text{-}at}(\operatorname{rr\text{-}ms}, \operatorname{rr\text{-}ms}, f, x))) \Rightarrow  \\ \quad \quad \operatorname{is\text{-}left\text{-}limit}(f, x, f(x)) \end{array}}}

\vnbbookfile{theorem-library/recip-succ-small}{theorem-library/recip-succ-small}
\noindent{\small recip-succ-small.scm -- the eventual form of the Archimedean bound.}\par
\vnbprop{recip-succ-small}{\vnbstmt{\begin{array}{@{}l@{}} \forall d. \\ \quad \operatorname{pos\text{-}rr}(d) \Rightarrow  \\ \quad \quad \exists n \in \mathbb{N}.\; \forall m \in \mathbb{N}.\; ((n \leq m) \Rightarrow (\operatorname{recip}((m + 1)) \leq d)) \end{array}}}

\vnbbookfile{theorem-library/continuous-is-regulated}{theorem-library/continuous-is-regulated}
\noindent{\small continuous-is-regulated.scm -- docs/calculus.pdf Ch 4, the bridge from Section 1 to the finished Section 2/3 arc.}\par
\vnbprop{continuous-is-regulated}{\vnbstmt{\begin{array}{@{}l@{}} \forall f, a, b. \\ \quad f \in (\mathbb{R} \to \mathbb{R}) \wedge \\ \quad a \in \mathbb{R} \wedge \\ \quad b \in \mathbb{R} \wedge \\ \quad a < b \wedge \\ \quad \forall x \in \operatorname{ccint}(a, b).\; \operatorname{is\text{-}continuous\text{-}at}(\operatorname{rr\text{-}ms}, \operatorname{rr\text{-}ms}, f, x) \Rightarrow \\ \quad \quad \operatorname{is\text{-}regulated}(f, a, b) \end{array}}}

\vnbbookfile{theorem-library/cauchy-criterion-left}{theorem-library/cauchy-criterion-left}
\noindent{\small cauchy-criterion-left.scm -- docs/calculus.pdf Ch 4 Sec 1.}\par
\vnbprop{cauchy-criterion-left}{\vnbstmt{\begin{array}{@{}l@{}} \forall f, x. \\ \quad f \in (\mathbb{R} \to \mathbb{R}) \wedge \\ \quad x \in \mathbb{R} \wedge \\ \quad \forall s. \\ \quad \quad \operatorname{pos\text{-}rr}(s) \Rightarrow  \\ \quad \quad \quad \exists a. \\ \quad \quad \quad \quad \operatorname{pos\text{-}rr}(a) \wedge \\ \quad \quad \quad \quad \forall b, t. \\ \quad \quad \quad \quad \quad (((b \in \mathbb{R}) \wedge (((x - a) \leq b) \wedge (b < x))) \wedge ((t \in \mathbb{R}) \wedge (((x - a) \leq t) \wedge (t < x)))) \Rightarrow  \\ \quad \quad \quad \quad \quad \quad \lvert (f(b) - f(t)) \rvert \leq s \Rightarrow \\ \quad \quad \exists l.\; \operatorname{is\text{-}left\text{-}limit}(f, x, l) \end{array}}}

\vnbbookfile{theorem-library/cauchy-criterion-right}{theorem-library/cauchy-criterion-right}
\noindent{\small cauchy-criterion-right.scm -- docs/calculus.pdf Ch 4 Sec 1.}\par
\vnbprop{cauchy-criterion-right}{\vnbstmt{\begin{array}{@{}l@{}} \forall f, x. \\ \quad f \in (\mathbb{R} \to \mathbb{R}) \wedge \\ \quad x \in \mathbb{R} \wedge \\ \quad \forall s. \\ \quad \quad \operatorname{pos\text{-}rr}(s) \Rightarrow  \\ \quad \quad \quad \exists a. \\ \quad \quad \quad \quad \operatorname{pos\text{-}rr}(a) \wedge \\ \quad \quad \quad \quad \forall b, t. \\ \quad \quad \quad \quad \quad (((b \in \mathbb{R}) \wedge ((x < b) \wedge (b \leq (x + a)))) \wedge ((t \in \mathbb{R}) \wedge ((x < t) \wedge (t \leq (x + a))))) \Rightarrow  \\ \quad \quad \quad \quad \quad \quad \lvert (f(b) - f(t)) \rvert \leq s \Rightarrow \\ \quad \quad \exists l.\; \operatorname{is\text{-}right\text{-}limit}(f, x, l) \end{array}}}

\vnbbookfile{theorem-library/partition-locates-point}{theorem-library/partition-locates-point}
\noindent{\small partition-locates-point.scm -- docs/calculus.pdf Ch 4 Sec 1.}\par
\vnbprop{partition-locates-point}{\vnbstmt{\begin{array}{@{}l@{}} \forall n \in \mathbb{N}, p, a, b, x. \\ \quad (\operatorname{is\text{-}partition}(p, n, a, b) \wedge (x \in \operatorname{ccint}(a, b))) \Rightarrow  \\ \quad \quad \exists i \in \operatorname{interval}(0, n).\; ((i < n) \wedge ((p(i) \leq x) \wedge (x \leq p((i + 1))))) \end{array}}}

\vnbbookfile{theorem-library/uniform-limit-regulated}{theorem-library/uniform-limit-regulated}
\noindent{\small uniform-limit-regulated.scm -- a uniform limit of regulated functions is regulated.  docs/calculus.pdf Prop 4.5, the sufficiency half.}\par
\vnbprop{uniform-limit-of-regulated}{\vnbstmt{\begin{array}{@{}l@{}} \forall m, g, a, b. \\ \quad (\operatorname{is\text{-}unif\text{-}limit\text{-}on}(m, g, a, b) \wedge \forall k \in \mathbb{N}.\; \operatorname{is\text{-}regulated}(m(k), a, b)) \Rightarrow  \\ \quad \quad \operatorname{is\text{-}regulated}(g, a, b) \end{array}}}

\vnbbookfile{theorem-library/uniform-limit-local}{theorem-library/uniform-limit-local}
\noindent{\small theorem-library/uniform-limit-local.scm -- A UNIFORM LIMIT IS CONTINUOUS AT A POINT, from uniform convergence on a BALL about that point only.}\par
\vnbprop{uniform-limit-continuous-at}{\vnbstmt{\begin{array}{@{}l@{}} \forall m \in (\mathbb{N} \to (\mathbb{R} \to \mathbb{R})), g \in (\mathbb{R} \to \mathbb{R}), t, o \in \mathbb{R}. \\ \quad 0 < o \wedge \\ \quad \forall k \in \mathbb{N}.\; \operatorname{is\text{-}continuous\text{-}at}(\operatorname{rr\text{-}ms}, \operatorname{rr\text{-}ms}, m(k), t) \wedge \\ \quad \forall s. \\ \quad \quad \operatorname{pos\text{-}rr}(s) \Rightarrow  \\ \quad \quad \quad \exists p \in \mathbb{N}. \forall k \in \mathbb{N}, x \in \mathbb{R}. \\ \quad \quad \quad \quad ((p \leq k) \wedge (\lvert (x - t) \rvert \leq o)) \Rightarrow (\lvert (g(x) - m(k)(x)) \rvert \leq s) \Rightarrow \\ \quad \quad \operatorname{is\text{-}continuous\text{-}at}(\operatorname{rr\text{-}ms}, \operatorname{rr\text{-}ms}, g, t) \end{array}}}

\vnbbookfile{theorem-library/diff-at-local}{theorem-library/diff-at-local}
\noindent{\small theorem-library/diff-at-local.scm -- DIFFERENTIABILITY IS LOCAL, and the Caratheodory transfer of a POINTWISE LIMIT.}\par
\vnbprop{diff-at-local}{\vnbstmt{\begin{array}{@{}l@{}} \forall f, i \in (\mathbb{R} \to \mathbb{R}), t, v, o \in \mathbb{R}. \\ \quad 0 < o \wedge \\ \quad \operatorname{is\text{-}continuous\text{-}at}(\operatorname{rr\text{-}ms}, \operatorname{rr\text{-}ms}, i, t) \wedge \\ \quad i(t) = v \wedge \\ \quad \forall x \in \mathbb{R}.\; ((\lvert (x - t) \rvert \leq o) \Rightarrow ((f(x) - f(t)) = (i(x) \cdot (x - t)))) \Rightarrow \\ \quad \quad \operatorname{is\text{-}diff\text{-}at}(f, t, v) \end{array}}}
\vnbprop{diff-at-ptwise-limit}{\vnbstmt{\begin{array}{@{}l@{}} \forall m, a \in (\mathbb{N} \to (\mathbb{R} \to \mathbb{R})), f, i \in (\mathbb{R} \to \mathbb{R}), t, o \in \mathbb{R}. \\ \quad 0 < o \wedge \\ \quad \operatorname{is\text{-}continuous\text{-}at}(\operatorname{rr\text{-}ms}, \operatorname{rr\text{-}ms}, i, t) \wedge \\ \quad \forall k \in \mathbb{N}, y \in \mathbb{R}. \\ \quad \quad (\lvert (y - t) \rvert \leq o) \Rightarrow ((m(k)(y) - m(k)(t)) = (a(k)(y) \cdot (y - t))) \wedge \\ \quad \forall y \in \mathbb{R}. \\ \quad \quad (\lvert (y - t) \rvert \leq o) \Rightarrow \operatorname{converges\text{-}to}(\operatorname{rr\text{-}ms}, (\lambda b \in \mathbb{N}.\; m(b)(y)), f(y)) \wedge \\ \quad \forall y \in \mathbb{R}. \\ \quad \quad (\lvert (y - t) \rvert \leq o) \Rightarrow \operatorname{converges\text{-}to}(\operatorname{rr\text{-}ms}, (\lambda b \in \mathbb{N}.\; a(b)(y)), i(y)) \Rightarrow \\ \quad \quad \operatorname{is\text{-}diff\text{-}at}(f, t, i(t)) \end{array}}}

\vnbbookfile{theorem-library/generalized-mvt-proof}{theorem-library/generalized-mvt-proof}
\noindent{\small theorem-library/generalized-mvt-proof.scm -- Cauchy MVT (Thm 2.11), PROVEN. L*(g(b)-g(a)) = M*(f(b)-f(a))  for some interior theta with f'=L, g'=M. Strategy: apply Rolle to the auxiliary h(z) = f(z)*(g(b)-g(a)) - g(z)*(f(b)-f(a))}\par
\vnbprop{generalized-mvt}{\vnbstmt{\begin{array}{@{}l@{}} \forall f, g, a, b. \\ \quad f \in (\mathbb{R} \to \mathbb{R}) \wedge \\ \quad g \in (\mathbb{R} \to \mathbb{R}) \wedge \\ \quad a \in \mathbb{R} \wedge \\ \quad b \in \mathbb{R} \wedge \\ \quad a < b \wedge \\ \quad \forall x \in \operatorname{ccint}(a, b). \\ \quad \quad \operatorname{is\text{-}continuous\text{-}at}(\operatorname{rr\text{-}ms}, \operatorname{rr\text{-}ms}, f, x) \wedge \operatorname{is\text{-}continuous\text{-}at}(\operatorname{rr\text{-}ms}, \operatorname{rr\text{-}ms}, g, x) \wedge \\ \quad \forall x \in \mathbb{R}. \\ \quad \quad ((a < x) \wedge (x < b)) \Rightarrow  \\ \quad \quad \quad \exists l.\; \operatorname{is\text{-}diff\text{-}at}(f, x, l) \wedge \exists m.\; \operatorname{is\text{-}diff\text{-}at}(g, x, m) \Rightarrow \\ \quad \quad \exists c \in \mathbb{R}. \\ \quad \quad \quad a < c \wedge \\ \quad \quad \quad c < b \wedge \\ \quad \quad \quad \exists l, m. \\ \quad \quad \quad \quad (\operatorname{is\text{-}diff\text{-}at}(f, c, l) \wedge (\operatorname{is\text{-}diff\text{-}at}(g, c, m) \wedge ((l \cdot (g(b) - g(a))) = (m \cdot (f(b) - f(a)))))) \end{array}}}

\vnbbookfile{theorem-library/deriv-monotone-proof}{theorem-library/deriv-monotone-proof}
\noindent{\small theorem-library/deriv-monotone-proof.scm -- increasing function theorem. f continuous on [a,b], f' $>$ 0 on (a,b)  =$>$  f strictly increasing on [a,b]. Strategy: for u$<$v in [a,b], MVT on [u,v] gives theta with f(v)-f(u) = f'(theta)(v-u);  f'(theta) $>$ 0 and v-u $>$ 0, so f(v)-f(u) $>$ 0.}\par
\vnbprop{deriv-pos-strictly-increasing}{\vnbstmt{\begin{array}{@{}l@{}} \forall f, a, b, u, v. \\ \quad f \in (\mathbb{R} \to \mathbb{R}) \wedge \\ \quad a \in \mathbb{R} \wedge \\ \quad b \in \mathbb{R} \wedge \\ \quad a < b \wedge \\ \quad \forall x \in \operatorname{ccint}(a, b).\; \operatorname{is\text{-}continuous\text{-}at}(\operatorname{rr\text{-}ms}, \operatorname{rr\text{-}ms}, f, x) \wedge \\ \quad \forall x \in \mathbb{R}. \\ \quad \quad ((a < x) \wedge (x < b)) \Rightarrow \exists l.\; (\operatorname{is\text{-}diff\text{-}at}(f, x, l) \wedge (0 < l)) \wedge \\ \quad u \in \operatorname{ccint}(a, b) \wedge \\ \quad v \in \operatorname{ccint}(a, b) \wedge \\ \quad u < v \Rightarrow \\ \quad \quad f(u) < f(v) \end{array}}}

\vnbbookfile{theorem-library/taylor-proof}{theorem-library/taylor-proof}
\noindent{\small RETIRED 2026-09-17 (proven): power-zero-base (was an add-to-pss here) -- theorem-library/rake-combinatorics.scm theorem-library/taylor-proof.scm -- Taylor's theorem, Lagrange remainder. f\textasciicircum{}(n) C\textasciicircum{}0 on [a,x], f\textasciicircum{}(n+1) exists on (a,x), a$<$x  =$>$  there is theta in (a,x) with  (n+1)! * R\_n(x) = f\textasciicircum{}(n+1)(theta) * (x-a)\textasciicircum{}(n+1),}\par
\par\medskip\noindent\textbf{Definition} (\texttt{taylor-poly}).\;$\operatorname{taylor\text{-}poly}(f, a, n, x) :=$
\begin{equation*}
\begin{array}{@{}l@{}} \operatorname{series\text{-}partial\text{-}sum}((\lambda k \in \mathbb{N}.\; ((\operatorname{nth\text{-}deriv}(f, k)(a) \cdot \operatorname{power}((x - a), k)) \cdot \operatorname{recip}(\operatorname{factorial}(k)))), \\ \quad (n + 1)) \end{array}
\end{equation*}

\par\medskip\noindent\textbf{Definition} (\texttt{taylor-differentiable}).\;For all $f$, $a$, $x$ and $n$, 
\begin{equation*}
\begin{array}{@{}l@{}} \operatorname{taylor\text{-}differentiable}(f, a, x, n) \\ \Leftrightarrow\; \forall k. \\ \quad ((k \in \mathbb{N}) \wedge (k \leq n)) \Rightarrow  \\ \quad \quad \forall t \in \operatorname{ccint}(a, x).\; \operatorname{is\text{-}continuous\text{-}at}(\operatorname{rr\text{-}ms}, \operatorname{rr\text{-}ms}, \operatorname{nth\text{-}deriv}(f, k), t) \wedge \\ \forall k. \\ \quad ((k \in \mathbb{N}) \wedge (k \leq n)) \Rightarrow  \\ \quad \quad \forall t. \\ \quad \quad \quad ((a < t) \wedge (t < x)) \Rightarrow  \\ \quad \quad \quad \quad \operatorname{is\text{-}diff\text{-}at}(\operatorname{nth\text{-}deriv}(f, k), t, \operatorname{nth\text{-}deriv}(f, \operatorname{succ}(k))(t)) \end{array}
\end{equation*}

\vnbprop{rr-recip-one}{\vnbstmt{\operatorname{recip}(1) = 1}}
\vnbprop{rr-zero-mul}{\vnbstmt{\forall v \in \mathbb{R}.\; ((0 \cdot v) = 0)}}
\vnbprop{rr-mul-zero-mid}{\vnbstmt{\forall u, v \in \mathbb{R}.\; (((u \cdot 0) \cdot v) = 0)}}
\vnbprop{rr-mul-one-mid}{\vnbstmt{\forall u \in \mathbb{R}.\; (u = u)}}
\vnbprop{factorial-real-pos}{\vnbstmt{\forall n \in \mathbb{N}.\; ((\operatorname{factorial}(n) \in \mathbb{R}) \wedge (0 < \operatorname{factorial}(n)))}}
\vnbprop{recip-factorial-in-rr}{\vnbstmt{\forall n \in \mathbb{N}.\; (\operatorname{recip}(\operatorname{factorial}(n)) \in \mathbb{R})}}
\vnbprop{taylor-center-term-at-zero}{\vnbstmt{\begin{array}{@{}l@{}} \forall f.; \forall x \in \mathbb{R}. \\ \quad (f(x) \in \mathbb{R}) \Rightarrow  \\ \quad \quad (\lambda k \in \mathbb{N}.\; ((\operatorname{nth\text{-}deriv}(f, k)(x) \cdot \operatorname{power}((x - x), k)) \cdot \operatorname{recip}(\operatorname{factorial}(k))))(0) = f(x) \end{array}}}
\vnbprop{taylor-center-term-at-succ}{\vnbstmt{\begin{array}{@{}l@{}} \forall f.; \forall x \in \mathbb{R}, m \in \mathbb{N}. \\ \quad (\operatorname{nth\text{-}deriv}(f, (m + 1))(x) \in \mathbb{R}) \Rightarrow  \\ \quad \quad (\lambda k \in \mathbb{N}.\; ((\operatorname{nth\text{-}deriv}(f, k)(x) \cdot \operatorname{power}((x - x), k)) \cdot \operatorname{recip}(\operatorname{factorial}(k))))((m + 1)) = 0 \end{array}}}
\vnbprop{taylor-center-partial-sum}{\vnbstmt{\begin{array}{@{}l@{}} \forall n \in \mathbb{N}, f.; \forall x \in \mathbb{R}. \\ \quad \forall k \in \mathbb{N}.\; ((k \leq n) \Rightarrow (\operatorname{nth\text{-}deriv}(f, k)(x) \in \mathbb{R})) \Rightarrow  \\ \quad \quad \operatorname{series\text{-}partial\text{-}sum}((\lambda k \in \mathbb{N}.\; ((\operatorname{nth\text{-}deriv}(f, k)(x) \cdot \operatorname{power}((x - x), k)) \cdot \operatorname{recip}(\operatorname{factorial}(k)))), \\ \quad \quad \quad (n + 1)) \\ \quad \quad =\; f(x) \end{array}}}
\vnbprop{taylor-poly-at-center}{\vnbstmt{\begin{array}{@{}l@{}} \forall f, x, n. \\ \quad ((x \in \mathbb{R}) \wedge ((n \in \mathbb{N}) \wedge \forall k \in \mathbb{N}.\; ((k \leq n) \Rightarrow (\operatorname{nth\text{-}deriv}(f, k)(x) \in \mathbb{R})))) \Rightarrow  \\ \quad \quad \operatorname{taylor\text{-}poly}(f, x, n, x) = f(x) \end{array}}}
\vnbprop{taylor-derivs-in-fun}{\vnbstmt{\begin{array}{@{}l@{}} \forall f, a, x, n.; \forall k \in \mathbb{N}. \\ \quad (\operatorname{taylor\text{-}differentiable}(f, a, x, n) \wedge ((a \in \mathbb{R}) \wedge ((x \in \mathbb{R}) \wedge ((a < x) \wedge (k \leq n))))) \Rightarrow  \\ \quad \quad \operatorname{nth\text{-}deriv}(f, k) \in (\mathbb{R} \to \mathbb{R}) \end{array}}}
\vnbprop{taylor-derivs-values-real}{\vnbstmt{\begin{array}{@{}l@{}} \forall f, n.; \forall t \in \mathbb{R}, k \in \mathbb{N}. \\ \quad (\forall k \in \mathbb{N}.\; ((k \leq n) \Rightarrow (\operatorname{nth\text{-}deriv}(f, k) \in (\mathbb{R} \to \mathbb{R}))) \wedge (k \leq n)) \Rightarrow  \\ \quad \quad \operatorname{nth\text{-}deriv}(f, k)(t) \in \mathbb{R} \end{array}}}
\vnbprop{taylor-poly-in-rr-at}{\vnbstmt{\begin{array}{@{}l@{}} \forall n \in \mathbb{N}, f.; \forall a, x \in \mathbb{R}. \\ \quad \forall k \in \mathbb{N}.\; ((k \leq n) \Rightarrow (\operatorname{nth\text{-}deriv}(f, k)(a) \in \mathbb{R})) \Rightarrow  \\ \quad \quad \operatorname{series\text{-}partial\text{-}sum}((\lambda k \in \mathbb{N}.\; ((\operatorname{nth\text{-}deriv}(f, k)(a) \cdot \operatorname{power}((x - a), k)) \cdot \operatorname{recip}(\operatorname{factorial}(k)))), \\ \quad \quad \quad (n + 1)) \\ \quad \quad \in\; \mathbb{R} \end{array}}}
\vnbprop{taylor-poly-in-rr}{\vnbstmt{\begin{array}{@{}l@{}} \forall f \in (\mathbb{R} \to \mathbb{R}), a \in \mathbb{R}, n \in \mathbb{N}, x \in \mathbb{R}. \\ \quad \forall k \in \mathbb{N}.\; ((k \leq n) \Rightarrow (\operatorname{nth\text{-}deriv}(f, k) \in (\mathbb{R} \to \mathbb{R}))) \Rightarrow  \\ \quad \quad \operatorname{taylor\text{-}poly}(f, a, n, x) \in \mathbb{R} \end{array}}}
\vnbprop{taylor-g-in-fun}{\vnbstmt{\begin{array}{@{}l@{}} \forall n \in \mathbb{N}, f \in (\mathbb{R} \to \mathbb{R}), x \in \mathbb{R}. \\ \quad \forall k \in \mathbb{N}.\; ((k \leq n) \Rightarrow (\operatorname{nth\text{-}deriv}(f, k) \in (\mathbb{R} \to \mathbb{R}))) \Rightarrow  \\ \quad \quad (\lambda z \in \mathbb{R}.\; (f(x) - \operatorname{taylor\text{-}poly}(f, z, n, x))) \in (\mathbb{R} \to \mathbb{R}) \end{array}}}
\vnbprop{taylor-h-in-fun}{\vnbstmt{\forall n \in \mathbb{N}, x \in \mathbb{R}.\; ((\lambda z \in \mathbb{R}.\; \operatorname{power}((x - z), (n + 1))) \in (\mathbb{R} \to \mathbb{R}))}}
\vnbprop{taylor-h-at-a}{\vnbstmt{\begin{array}{@{}l@{}} \forall a, x, n. \\ \quad ((a \in \mathbb{R}) \wedge ((x \in \mathbb{R}) \wedge (n \in \mathbb{N}))) \Rightarrow  \\ \quad \quad (\lambda z \in \mathbb{R}.\; \operatorname{power}((x - z), (n + 1)))(a) = \operatorname{power}((x - a), (n + 1)) \end{array}}}
\vnbprop{taylor-h-at-x}{\vnbstmt{\begin{array}{@{}l@{}} \forall x \in \mathbb{R}, n \in \mathbb{N}. \\ \quad (\lambda z \in \mathbb{R}.\; \operatorname{power}((x - z), (n + 1)))(x) = 0 \end{array}}}
\vnbprop{taylor-g-at-a}{\vnbstmt{\begin{array}{@{}l@{}} \forall f, a, x, n. \\ \quad f \in (\mathbb{R} \to \mathbb{R}) \wedge \\ \quad a \in \mathbb{R} \wedge \\ \quad x \in \mathbb{R} \wedge \\ \quad n \in \mathbb{N} \wedge \\ \quad \forall k \in \mathbb{N}.\; ((k \leq n) \Rightarrow (\operatorname{nth\text{-}deriv}(f, k) \in (\mathbb{R} \to \mathbb{R}))) \Rightarrow \\ \quad \quad (\lambda z \in \mathbb{R}.\; (f(x) - \operatorname{taylor\text{-}poly}(f, z, n, x)))(a) = (f(x) - \operatorname{taylor\text{-}poly}(f, a, n, x)) \end{array}}}
\vnbprop{taylor-g-at-x}{\vnbstmt{\begin{array}{@{}l@{}} \forall f, x.; \forall n \in \mathbb{N}. \\ \quad ((f \in (\mathbb{R} \to \mathbb{R})) \wedge ((x \in \mathbb{R}) \wedge \forall k \in \mathbb{N}.\; ((k \leq n) \Rightarrow (\operatorname{nth\text{-}deriv}(f, k) \in (\mathbb{R} \to \mathbb{R}))))) \Rightarrow  \\ \quad \quad (\lambda z \in \mathbb{R}.\; (f(x) - \operatorname{taylor\text{-}poly}(f, z, n, x)))(x) = 0 \end{array}}}
\vnbprop{taylor-deriv-real}{\vnbstmt{\begin{array}{@{}l@{}} \forall f, a, x, n, t. \\ \quad (\operatorname{taylor\text{-}differentiable}(f, a, x, n) \wedge ((a < t) \wedge ((t < x) \wedge (n \in \mathbb{N})))) \Rightarrow  \\ \quad \quad \operatorname{nth\text{-}deriv}(f, (n + 1))(t) \in \mathbb{R} \end{array}}}
\vnbprop{tg-neg-scale}{\vnbstmt{\forall s, p \in \mathbb{R}.\; ((0 - (s \cdot p)) = ((s \cdot p) \cdot -1))}}
\vnbprop{taylor-h-diff}{\vnbstmt{\begin{array}{@{}l@{}} \forall x, n, t. \\ \quad ((x \in \mathbb{R}) \wedge ((t \in \mathbb{R}) \wedge (n \in \mathbb{N}))) \Rightarrow  \\ \quad \quad \operatorname{is\text{-}diff\text{-}at}((\lambda z \in \mathbb{R}.\; \operatorname{power}((x - z), (n + 1))), t, (0 - ((n + 1) \cdot \operatorname{power}((x - t), n)))) \end{array}}}
\vnbprop{tgd-recip-factorial-succ}{\vnbstmt{\forall m \in \mathbb{N}.\; (\operatorname{recip}(\operatorname{factorial}(m)) = ((m + 1) \cdot \operatorname{recip}(\operatorname{factorial}((m + 1)))))}}
\vnbprop{tgd-sub-diff}{\vnbstmt{\forall x, t \in \mathbb{R}.\; \operatorname{is\text{-}diff\text{-}at}((\lambda z \in \mathbb{R}.\; (x - z)), t, -1)}}
\vnbprop{tgd-power-sub-diff}{\vnbstmt{\begin{array}{@{}l@{}} \forall m \in \mathbb{N}, x, t \in \mathbb{R}. \\ \quad \operatorname{is\text{-}diff\text{-}at}((\lambda z \in \mathbb{R}.\; \operatorname{power}((x - z), (m + 1))), t, (((m + 1) \cdot \operatorname{power}((x - t), m)) \cdot -1)) \end{array}}}
\vnbprop{tgd-step-identity}{\vnbstmt{\begin{array}{@{}l@{}} \forall s, r, a, b, p, q \in \mathbb{R}. \\ \quad ((((b \cdot q) + (a \cdot ((s \cdot p) \cdot -1))) \cdot r) + ((a \cdot q) \cdot 0)) = (((r \cdot b) \cdot q) - (((s \cdot r) \cdot a) \cdot p)) \end{array}}}
\vnbprop{tgd-telescope-identity}{\vnbstmt{\forall u, v \in \mathbb{R}.\; (v = (u + (v - u)))}}
\vnbprop{tgd-term-diff}{\vnbstmt{\begin{array}{@{}l@{}} \forall m \in \mathbb{N}, f.; \forall x, t \in \mathbb{R}. \\ \quad \operatorname{is\text{-}diff\text{-}at}(\operatorname{nth\text{-}deriv}(f, (m + 1)), t, \operatorname{nth\text{-}deriv}(f, ((m + 1) + 1))(t)) \Rightarrow  \\ \quad \quad \operatorname{is\text{-}diff\text{-}at}((\lambda z \in \mathbb{R}.\; ((\operatorname{nth\text{-}deriv}(f, (m + 1))(z) \cdot \operatorname{power}((x - z), (m + 1))) \cdot \operatorname{recip}(\operatorname{factorial}((m + 1))))), \\ \quad \quad \quad t, \\ \quad \quad \quad (((\operatorname{recip}(\operatorname{factorial}((m + 1))) \cdot \operatorname{nth\text{-}deriv}(f, ((m + 1) + 1))(t)) \cdot \operatorname{power}((x - t), (m + 1))) - ((\operatorname{recip}(\operatorname{factorial}(m)) \cdot \operatorname{nth\text{-}deriv}(f, (m + 1))(t)) \cdot \operatorname{power}((x - t), m)))) \end{array}}}
\vnbprop{tgd-base-sum-value}{\vnbstmt{\begin{array}{@{}l@{}} \forall f \in (\mathbb{R} \to \mathbb{R}), x, z \in \mathbb{R}. \\ \quad (\operatorname{series\text{-}partial\text{-}sum}((\lambda k \in \mathbb{N}.\; ((\operatorname{nth\text{-}deriv}(f, k)(z) \cdot \operatorname{power}((x - z), k)) \cdot \operatorname{recip}(\operatorname{factorial}(k)))), 1) \simeq f(z)) \end{array}}}
\vnbprop{tgd-sum-diff}{\vnbstmt{\begin{array}{@{}l@{}} \forall m \in \mathbb{N}, f.; \forall x, t \in \mathbb{R}. \\ \quad \forall k \in \mathbb{N}.\; ((k \leq m) \Rightarrow \operatorname{is\text{-}diff\text{-}at}(\operatorname{nth\text{-}deriv}(f, k), t, \operatorname{nth\text{-}deriv}(f, (k + 1))(t))) \Rightarrow  \\ \quad \quad \operatorname{is\text{-}diff\text{-}at}((\lambda z \in \mathbb{R}.\; \\ \quad \quad \quad \quad \operatorname{series\text{-}partial\text{-}sum}((\lambda k \in \mathbb{N}.\; ((\operatorname{nth\text{-}deriv}(f, k)(z) \cdot \operatorname{power}((x - z), k)) \cdot \operatorname{recip}(\operatorname{factorial}(k)))), \\ \quad \quad \quad \quad \quad (m + 1))), \\ \quad \quad \quad t, \\ \quad \quad \quad (\operatorname{recip}(\operatorname{factorial}(m)) \cdot \operatorname{nth\text{-}deriv}(f, (m + 1))(t)) \cdot \operatorname{power}((x - t), m)) \end{array}}}
\vnbprop{taylor-g-diff}{\vnbstmt{\begin{array}{@{}l@{}} \forall f, a, x, n, t. \\ \quad ((n \in \mathbb{N}) \wedge ((x \in \mathbb{R}) \wedge (\operatorname{taylor\text{-}differentiable}(f, a, x, n) \wedge ((a < t) \wedge (t < x))))) \Rightarrow  \\ \quad \quad \operatorname{is\text{-}diff\text{-}at}((\lambda z \in \mathbb{R}.\; (f(x) - \operatorname{taylor\text{-}poly}(f, z, n, x))), t, \\ \quad \quad \quad (0 - ((\operatorname{recip}(\operatorname{factorial}(n)) \cdot \operatorname{nth\text{-}deriv}(f, (n + 1))(t)) \cdot \operatorname{power}((x - t), n)))) \end{array}}}
\vnbprop{taylor-h-cont}{\vnbstmt{\begin{array}{@{}l@{}} \forall x, n, t. \\ \quad ((x \in \mathbb{R}) \wedge ((t \in \mathbb{R}) \wedge (n \in \mathbb{N}))) \Rightarrow  \\ \quad \quad \operatorname{is\text{-}continuous\text{-}at}(\operatorname{rr\text{-}ms}, \operatorname{rr\text{-}ms}, (\lambda z \in \mathbb{R}.\; \operatorname{power}((x - z), (n + 1))), t) \end{array}}}
\vnbprop{tgc-term-cont}{\vnbstmt{\begin{array}{@{}l@{}} \forall m \in \mathbb{N}, f.; \forall x, t \in \mathbb{R}. \\ \quad \operatorname{is\text{-}continuous\text{-}at}(\operatorname{rr\text{-}ms}, \operatorname{rr\text{-}ms}, \operatorname{nth\text{-}deriv}(f, (m + 1)), t) \Rightarrow  \\ \quad \quad \operatorname{is\text{-}continuous\text{-}at}(\operatorname{rr\text{-}ms}, \operatorname{rr\text{-}ms}, \\ \quad \quad \quad (\lambda z \in \mathbb{R}.\; ((\operatorname{nth\text{-}deriv}(f, (m + 1))(z) \cdot \operatorname{power}((x - z), (m + 1))) \cdot \operatorname{recip}(\operatorname{factorial}((m + 1))))), \\ \quad \quad \quad t) \end{array}}}
\vnbprop{tgc-sum-cont}{\vnbstmt{\begin{array}{@{}l@{}} \forall m \in \mathbb{N}, f.; \forall x, t \in \mathbb{R}. \\ \quad \forall k \in \mathbb{N}.\; ((k \leq m) \Rightarrow \operatorname{is\text{-}continuous\text{-}at}(\operatorname{rr\text{-}ms}, \operatorname{rr\text{-}ms}, \operatorname{nth\text{-}deriv}(f, k), t)) \Rightarrow  \\ \quad \quad \operatorname{is\text{-}continuous\text{-}at}(\operatorname{rr\text{-}ms}, \operatorname{rr\text{-}ms}, \\ \quad \quad \quad (\lambda z \in \mathbb{R}.\; \\ \quad \quad \quad \quad \operatorname{series\text{-}partial\text{-}sum}((\lambda k \in \mathbb{N}.\; ((\operatorname{nth\text{-}deriv}(f, k)(z) \cdot \operatorname{power}((x - z), k)) \cdot \operatorname{recip}(\operatorname{factorial}(k)))), \\ \quad \quad \quad \quad \quad (m + 1))), \\ \quad \quad \quad t) \end{array}}}
\vnbprop{taylor-g-cont}{\vnbstmt{\begin{array}{@{}l@{}} \forall f, a, x, n, t. \\ \quad ((n \in \mathbb{N}) \wedge ((x \in \mathbb{R}) \wedge (\operatorname{taylor\text{-}differentiable}(f, a, x, n) \wedge (t \in \operatorname{ccint}(a, x))))) \Rightarrow  \\ \quad \quad \operatorname{is\text{-}continuous\text{-}at}(\operatorname{rr\text{-}ms}, \operatorname{rr\text{-}ms}, (\lambda z \in \mathbb{R}.\; (f(x) - \operatorname{taylor\text{-}poly}(f, z, n, x))), t) \end{array}}}
\vnbprop{taylor-gmvt-cont}{\vnbstmt{\begin{array}{@{}l@{}} \forall f, a, x, n.; \forall t \in \operatorname{ccint}(a, x). \\ \quad (\operatorname{taylor\text{-}differentiable}(f, a, x, n) \wedge ((x \in \mathbb{R}) \wedge (n \in \mathbb{N}))) \Rightarrow  \\ \quad \quad \operatorname{is\text{-}continuous\text{-}at}(\operatorname{rr\text{-}ms}, \operatorname{rr\text{-}ms}, (\lambda z \in \mathbb{R}.\; (f(x) - \operatorname{taylor\text{-}poly}(f, z, n, x))), t) \wedge \\ \quad \quad \operatorname{is\text{-}continuous\text{-}at}(\operatorname{rr\text{-}ms}, \operatorname{rr\text{-}ms}, (\lambda z \in \mathbb{R}.\; \operatorname{power}((x - z), (n + 1))), t) \end{array}}}
\vnbprop{taylor-gmvt-diff}{\vnbstmt{\begin{array}{@{}l@{}} \forall f, a, x, n.; \forall t \in \mathbb{R}. \\ \quad (\operatorname{taylor\text{-}differentiable}(f, a, x, n) \wedge ((n \in \mathbb{N}) \wedge ((x \in \mathbb{R}) \wedge ((a < t) \wedge (t < x))))) \Rightarrow  \\ \quad \quad \exists l. \\ \quad \quad \quad \operatorname{is\text{-}diff\text{-}at}((\lambda z \in \mathbb{R}.\; (f(x) - \operatorname{taylor\text{-}poly}(f, z, n, x))), t, l) \wedge \\ \quad \quad \exists m.\; \operatorname{is\text{-}diff\text{-}at}((\lambda z \in \mathbb{R}.\; \operatorname{power}((x - z), (n + 1))), t, m) \end{array}}}
\vnbprop{rr-power-pos}{\vnbstmt{\forall d \in \mathbb{R}, n \in \mathbb{N}.\; ((0 < d) \Rightarrow (0 < \operatorname{power}(d, n)))}}
\vnbprop{rr-cancel-mul-left}{\vnbstmt{\forall c, u, v \in \mathbb{R}.\; ((\neg (c = 0) \wedge ((c \cdot u) = (c \cdot v))) \Rightarrow (u = v))}}
\vnbprop{rr-recip-factorial}{\vnbstmt{\forall n \in \mathbb{N}.\; ((\operatorname{factorial}(n) \cdot \operatorname{recip}(\operatorname{factorial}(n))) = 1)}}
\vnbprop{taylor-clear}{\vnbstmt{\begin{array}{@{}l@{}} \forall n_{1}, w, d, r \in \mathbb{R}, x, a, b, c.; \forall n \in \mathbb{N}. \\ \quad 0 < w \wedge \\ \quad x = 0 \wedge \\ \quad b = 0 \wedge \\ \quad c = d \wedge \\ \quad a = r \wedge \\ \quad ((0 - ((\operatorname{recip}(\operatorname{factorial}(n)) \cdot n_{1}) \cdot w)) \cdot (b - c)) = ((0 - ((n + 1) \cdot w)) \cdot (x - a)) \Rightarrow \\ \quad \quad (\operatorname{factorial}((n + 1)) \cdot r) = (n_{1} \cdot d) \end{array}}}
\vnbprop{taylor-lagrange}{\vnbstmt{\begin{array}{@{}l@{}} \forall f, a, x, n. \\ \quad f \in (\mathbb{R} \to \mathbb{R}) \wedge \\ \quad a \in \mathbb{R} \wedge \\ \quad x \in \mathbb{R} \wedge \\ \quad n \in \mathbb{N} \wedge \\ \quad a < x \wedge \\ \quad \operatorname{taylor\text{-}differentiable}(f, a, x, n) \Rightarrow \\ \quad \quad \exists b \in \mathbb{R}. \\ \quad \quad \quad a < b \wedge \\ \quad \quad \quad b < x \wedge \\ \quad \quad \quad \operatorname{factorial}((n + 1)) \cdot (f(x) - \operatorname{taylor\text{-}poly}(f, a, n, x)) \\ \quad \quad \quad =\; \operatorname{nth\text{-}deriv}(f, (n + 1))(b) \cdot \operatorname{power}((x - a), (n + 1)) \end{array}}}

\vnbbookfile{theorem-library/rake-hb-leaves}{theorem-library/rake-hb-leaves}
\noindent{\small theorem-library/rake-hb-leaves.scm -- nine asserted leaves of the Hahn-Banach cluster, PROVEN.}\par
\vnbprop{extends-on-trans}{\vnbstmt{\begin{array}{@{}l@{}} \forall s, t, a, b, f. \\ \quad ((s \subseteq t) \wedge (\operatorname{extends\text{-}on}(t, a, b) \wedge \operatorname{extends\text{-}on}(s, b, f))) \Rightarrow  \\ \quad \quad \operatorname{extends\text{-}on}(s, a, f) \end{array}}}
\vnbprop{vnrm-nonneg}{\vnbstmt{\begin{array}{@{}l@{}} \forall m, w. \\ \quad (\operatorname{is\text{-}normed\text{-}vector\text{-}space}(m) \wedge (w \in \operatorname{vec}(m))) \Rightarrow (0 \leq \operatorname{vnrm}(m)(w)) \end{array}}}
\vnbprop{linfun-on-vec-is-linfun}{\vnbstmt{\begin{array}{@{}l@{}} \forall m, f. \\ \quad \operatorname{is\text{-}linear\text{-}functional\text{-}on}(m, \operatorname{vec}(m), f) \Rightarrow \operatorname{is\text{-}linear\text{-}functional}(m, f) \end{array}}}
\vnbprop{blf-app-real}{\vnbstmt{\begin{array}{@{}l@{}} \forall m, g, v. \\ \quad (\operatorname{is\text{-}bounded\text{-}linear\text{-}functional}(m, g) \wedge (v \in \operatorname{vec}(m))) \Rightarrow (g(v) \in \mathbb{R}) \end{array}}}
\vnbprop{bdd-linfun-abs-real}{\vnbstmt{\begin{array}{@{}l@{}} \forall m, f, x. \\ \quad (\operatorname{is\text{-}bounded\text{-}linear\text{-}functional}(m, f) \wedge (x \in \operatorname{vec}(m))) \Rightarrow (\lvert f(x) \rvert \in \mathbb{R}) \end{array}}}
\vnbprop{linfun-app-abs-real}{\vnbstmt{\begin{array}{@{}l@{}} \forall m, t, g, w. \\ \quad (\operatorname{is\text{-}linear\text{-}functional\text{-}on}(m, t, g) \wedge (w \in t)) \Rightarrow (\lvert g(w) \rvert \in \mathbb{R}) \end{array}}}
\vnbprop{abs-nonneg-le}{\vnbstmt{\forall a, c.\; (((a \in \mathbb{R}) \wedge ((0 \leq a) \wedge (\lvert a \rvert \leq c))) \Rightarrow (a \leq c))}}
\vnbprop{rr-mul-le-right}{\vnbstmt{\forall a, b, c \in \mathbb{R}.\; (((a \leq b) \wedge (0 \leq c)) \Rightarrow ((a \cdot c) \leq (b \cdot c)))}}
\vnbprop{proper-subset-witness}{\vnbstmt{\forall a, b.\; (((a \subseteq b) \wedge \neg (a = b)) \Rightarrow \exists x \in b.\; \neg (x \in a))}}

\vnbbookfile{theorem-library/rake-dual-norm-spec}{theorem-library/rake-dual-norm-spec}
\noindent{\small rake-dual-norm-spec.scm -- the DUAL-NORM / DUAL-NORM-ON descriptions DENOTE, and the corollaries that follow from them.  Rake batch 7, assignment 7-A. Helper prefix `r8a-'.}\par
\vnbprop{rr-le-of-le-add-all-pos}{\vnbstmt{\begin{array}{@{}l@{}} \forall a, b \in \mathbb{R}. \\ \quad \forall d.\; (\operatorname{pos\text{-}rr}(d) \Rightarrow (a \leq (b + d))) \Rightarrow (a \leq b) \end{array}}}
\vnbprop{rr-recip-cancel-right}{\vnbstmt{\forall a, n \in \mathbb{R}.\; (\neg (n = 0) \Rightarrow (a = ((a \cdot \operatorname{recip}(n)) \cdot n)))}}
\vnbprop{rr-mul-recip-cancel}{\vnbstmt{\forall a, n \in \mathbb{R}.\; (\neg (n = 0) \Rightarrow (a = ((a \cdot n) \cdot \operatorname{recip}(n))))}}
\vnbprop{linfun-is-linfun-on-vec}{\vnbstmt{\begin{array}{@{}l@{}} \forall m, f. \\ \quad \operatorname{is\text{-}linear\text{-}functional}(m, f) \Rightarrow \operatorname{is\text{-}linear\text{-}functional\text{-}on}(m, \operatorname{vec}(m), f) \end{array}}}
\vnbprop{blf-is-blf-on-vec}{\vnbstmt{\begin{array}{@{}l@{}} \forall m, f. \\ \quad \operatorname{is\text{-}bounded\text{-}linear\text{-}functional}(m, f) \Rightarrow  \\ \quad \quad \operatorname{is\text{-}bounded\text{-}linear\text{-}functional\text{-}on}(m, \operatorname{vec}(m), f) \end{array}}}
\vnbprop{dual-norm-on-lub-exists}{\vnbstmt{\begin{array}{@{}l@{}} \forall m, s, f. \\ \quad (\operatorname{is\text{-}normed\text{-}vector\text{-}space}(m) \wedge ((s \subseteq \operatorname{vec}(m)) \wedge \operatorname{is\text{-}bounded\text{-}linear\text{-}functional\text{-}on}(m, s, f))) \Rightarrow  \\ \quad \quad \exists c \in \mathbb{R}. \\ \quad \quad \quad 0 \leq c \wedge \\ \quad \quad \quad \forall x \in s.\; (\lvert f(x) \rvert \leq (c \cdot \operatorname{vnrm}(m)(x))) \wedge \\ \quad \quad \quad \forall d \in \mathbb{R}. \\ \quad \quad \quad \quad ((0 \leq d) \wedge \forall x \in s.\; (\lvert f(x) \rvert \leq (d \cdot \operatorname{vnrm}(m)(x)))) \Rightarrow (c \leq d) \end{array}}}
\vnbprop{dual-norm-on-lub-unique}{\vnbstmt{\begin{array}{@{}l@{}} \forall m, s, f, c_{1}, c_{2}. \\ \quad c_{1} \in \mathbb{R} \wedge \\ \quad 0 \leq c_{1} \wedge \\ \quad \forall x \in s.\; (\lvert f(x) \rvert \leq (c_{1} \cdot \operatorname{vnrm}(m)(x))) \wedge \\ \quad \forall d \in \mathbb{R}. \\ \quad \quad ((0 \leq d) \wedge \forall x \in s.\; (\lvert f(x) \rvert \leq (d \cdot \operatorname{vnrm}(m)(x)))) \Rightarrow (c_{1} \leq d) \wedge \\ \quad c_{2} \in \mathbb{R} \wedge \\ \quad 0 \leq c_{2} \wedge \\ \quad \forall x \in s.\; (\lvert f(x) \rvert \leq (c_{2} \cdot \operatorname{vnrm}(m)(x))) \wedge \\ \quad \forall d \in \mathbb{R}. \\ \quad \quad ((0 \leq d) \wedge \forall x \in s.\; (\lvert f(x) \rvert \leq (d \cdot \operatorname{vnrm}(m)(x)))) \Rightarrow (c_{2} \leq d) \Rightarrow \\ \quad \quad c_{1} = c_{2} \end{array}}}
\vnbprop{dual-norm-on-spec}{\vnbstmt{\begin{array}{@{}l@{}} \forall m, s, f. \\ \quad (\operatorname{is\text{-}normed\text{-}vector\text{-}space}(m) \wedge ((s \subseteq \operatorname{vec}(m)) \wedge \operatorname{is\text{-}bounded\text{-}linear\text{-}functional\text{-}on}(m, s, f))) \Rightarrow  \\ \quad \quad \operatorname{dual\text{-}norm\text{-}on}(m, s, f) \in \mathbb{R} \wedge \\ \quad \quad 0 \leq \operatorname{dual\text{-}norm\text{-}on}(m, s, f) \wedge \\ \quad \quad \forall x \in s.\; (\lvert f(x) \rvert \leq (\operatorname{dual\text{-}norm\text{-}on}(m, s, f) \cdot \operatorname{vnrm}(m)(x))) \wedge \\ \quad \quad \forall d \in \mathbb{R}. \\ \quad \quad \quad ((0 \leq d) \wedge \forall x \in s.\; (\lvert f(x) \rvert \leq (d \cdot \operatorname{vnrm}(m)(x)))) \Rightarrow  \\ \quad \quad \quad \quad \operatorname{dual\text{-}norm\text{-}on}(m, s, f) \leq d \end{array}}}
\vnbprop{dual-norm-spec}{\vnbstmt{\begin{array}{@{}l@{}} \forall m, f. \\ \quad (\operatorname{is\text{-}normed\text{-}vector\text{-}space}(m) \wedge \operatorname{is\text{-}bounded\text{-}linear\text{-}functional}(m, f)) \Rightarrow  \\ \quad \quad \operatorname{dual\text{-}norm}(m, f) \in \mathbb{R} \wedge \\ \quad \quad 0 \leq \operatorname{dual\text{-}norm}(m, f) \wedge \\ \quad \quad \forall x \in \operatorname{vec}(m).\; (\lvert f(x) \rvert \leq (\operatorname{dual\text{-}norm}(m, f) \cdot \operatorname{vnrm}(m)(x))) \wedge \\ \quad \quad \forall d \in \mathbb{R}. \\ \quad \quad \quad ((0 \leq d) \wedge \forall x \in \operatorname{vec}(m).\; (\lvert f(x) \rvert \leq (d \cdot \operatorname{vnrm}(m)(x)))) \Rightarrow  \\ \quad \quad \quad \quad \operatorname{dual\text{-}norm}(m, f) \leq d \end{array}}}
\vnbprop{dual-norm-on-nonneg-nvs}{\vnbstmt{\begin{array}{@{}l@{}} \forall m, s, f. \\ \quad (\operatorname{is\text{-}normed\text{-}vector\text{-}space}(m) \wedge ((s \subseteq \operatorname{vec}(m)) \wedge \operatorname{is\text{-}bounded\text{-}linear\text{-}functional\text{-}on}(m, s, f))) \Rightarrow  \\ \quad \quad (\operatorname{dual\text{-}norm\text{-}on}(m, s, f) \in \mathbb{R}) \wedge (0 \leq \operatorname{dual\text{-}norm\text{-}on}(m, s, f)) \end{array}}}
\vnbprop{dual-norm-on-le-bound-nvs}{\vnbstmt{\begin{array}{@{}l@{}} \forall m, s, g, c. \\ \quad \operatorname{is\text{-}normed\text{-}vector\text{-}space}(m) \wedge \\ \quad s \subseteq \operatorname{vec}(m) \wedge \\ \quad \operatorname{is\text{-}linear\text{-}functional\text{-}on}(m, s, g) \wedge \\ \quad c \in \mathbb{R} \wedge \\ \quad 0 \leq c \wedge \\ \quad \forall w \in s.\; (\lvert g(w) \rvert \leq (c \cdot \operatorname{vnrm}(m)(w))) \Rightarrow \\ \quad \quad \operatorname{dual\text{-}norm\text{-}on}(m, s, g) \leq c \end{array}}}
\vnbprop{dual-norm-on-is-bound-nvs}{\vnbstmt{\begin{array}{@{}l@{}} \forall m, s, f, x. \\ \quad \operatorname{is\text{-}normed\text{-}vector\text{-}space}(m) \wedge \\ \quad s \subseteq \operatorname{vec}(m) \wedge \\ \quad \operatorname{is\text{-}bounded\text{-}linear\text{-}functional\text{-}on}(m, s, f) \wedge \\ \quad x \in s \Rightarrow \\ \quad \quad \lvert f(x) \rvert \leq (\operatorname{dual\text{-}norm\text{-}on}(m, s, f) \cdot \operatorname{vnrm}(m)(x)) \end{array}}}
\vnbprop{dual-norm-nonneg-nvs}{\vnbstmt{\begin{array}{@{}l@{}} \forall m, f. \\ \quad (\operatorname{is\text{-}normed\text{-}vector\text{-}space}(m) \wedge \operatorname{is\text{-}bounded\text{-}linear\text{-}functional}(m, f)) \Rightarrow  \\ \quad \quad (\operatorname{dual\text{-}norm}(m, f) \in \mathbb{R}) \wedge (0 \leq \operatorname{dual\text{-}norm}(m, f)) \end{array}}}
\vnbprop{dual-norm-is-bound-nvs}{\vnbstmt{\begin{array}{@{}l@{}} \forall m, f, x. \\ \quad (\operatorname{is\text{-}normed\text{-}vector\text{-}space}(m) \wedge (\operatorname{is\text{-}bounded\text{-}linear\text{-}functional}(m, f) \wedge (x \in \operatorname{vec}(m)))) \Rightarrow  \\ \quad \quad \lvert f(x) \rvert \leq (\operatorname{dual\text{-}norm}(m, f) \cdot \operatorname{vnrm}(m)(x)) \end{array}}}
\vnbprop{dual-norm-le-bound-nvs}{\vnbstmt{\begin{array}{@{}l@{}} \forall m, f, c. \\ \quad \operatorname{is\text{-}normed\text{-}vector\text{-}space}(m) \wedge \\ \quad \operatorname{is\text{-}linear\text{-}functional}(m, f) \wedge \\ \quad c \in \mathbb{R} \wedge \\ \quad 0 \leq c \wedge \\ \quad \forall w \in \operatorname{vec}(m).\; (\lvert f(w) \rvert \leq (c \cdot \operatorname{vnrm}(m)(w))) \Rightarrow \\ \quad \quad \operatorname{dual\text{-}norm}(m, f) \leq c \end{array}}}
\vnbprop{good-sub-self-nvs}{\vnbstmt{\begin{array}{@{}l@{}} \forall m, s, f. \\ \quad (\operatorname{is\text{-}normed\text{-}vector\text{-}space}(m) \wedge (\operatorname{is\text{-}submodule}(m, s) \wedge \operatorname{is\text{-}bounded\text{-}linear\text{-}functional\text{-}on}(m, s, f))) \Rightarrow  \\ \quad \quad \operatorname{good\text{-}sub}(m, s, f, s) \end{array}}}

\vnbbookfile{theorem-library/nvs-module-view}{theorem-library/nvs-module-view}
\noindent{\small nvs-module-view.scm -- THE MODULE BRIDGE OF A NORMED VECTOR SPACE, PROVEN.}\par
\vnbprop{nvs-module-view-scal}{\vnbstmt{\begin{array}{@{}l@{}} \forall m. \\ \quad \operatorname{is\text{-}normed\text{-}vector\text{-}space}(m) \Rightarrow  \\ \quad \quad (\operatorname{scal}(\operatorname{normed\text{-}vector\text{-}space\text{-}as\text{-}module}(m)) \simeq \operatorname{scal}(m)) \end{array}}}
\vnbprop{nvs-module-view-vec}{\vnbstmt{\begin{array}{@{}l@{}} \forall m. \\ \quad \operatorname{is\text{-}normed\text{-}vector\text{-}space}(m) \Rightarrow  \\ \quad \quad (\operatorname{vec}(\operatorname{normed\text{-}vector\text{-}space\text{-}as\text{-}module}(m)) \simeq \operatorname{vec}(m)) \end{array}}}
\vnbprop{nvs-module-view-vadd}{\vnbstmt{\begin{array}{@{}l@{}} \forall m. \\ \quad \operatorname{is\text{-}normed\text{-}vector\text{-}space}(m) \Rightarrow  \\ \quad \quad (\operatorname{vadd}(\operatorname{normed\text{-}vector\text{-}space\text{-}as\text{-}module}(m)) \simeq \operatorname{vadd}(m)) \end{array}}}
\vnbprop{nvs-module-view-vzero}{\vnbstmt{\begin{array}{@{}l@{}} \forall m. \\ \quad \operatorname{is\text{-}normed\text{-}vector\text{-}space}(m) \Rightarrow  \\ \quad \quad (\operatorname{vzero}(\operatorname{normed\text{-}vector\text{-}space\text{-}as\text{-}module}(m)) \simeq \operatorname{vzero}(m)) \end{array}}}
\vnbprop{nvs-module-view-vneg}{\vnbstmt{\begin{array}{@{}l@{}} \forall m. \\ \quad \operatorname{is\text{-}normed\text{-}vector\text{-}space}(m) \Rightarrow  \\ \quad \quad (\operatorname{vneg}(\operatorname{normed\text{-}vector\text{-}space\text{-}as\text{-}module}(m)) \simeq \operatorname{vneg}(m)) \end{array}}}
\vnbprop{nvs-module-view-act}{\vnbstmt{\begin{array}{@{}l@{}} \forall m. \\ \quad \operatorname{is\text{-}normed\text{-}vector\text{-}space}(m) \Rightarrow  \\ \quad \quad (\operatorname{act}(\operatorname{normed\text{-}vector\text{-}space\text{-}as\text{-}module}(m)) \simeq \operatorname{act}(m)) \end{array}}}
\vnbprop{submodule-nvs-module-view}{\vnbstmt{\begin{array}{@{}l@{}} \forall m, t. \\ \quad (\operatorname{is\text{-}normed\text{-}vector\text{-}space}(m) \wedge \operatorname{is\text{-}submodule}(m, t)) \Rightarrow  \\ \quad \quad \operatorname{is\text{-}submodule}(\operatorname{normed\text{-}vector\text{-}space\text{-}as\text{-}module}(m), t) \end{array}}}

\vnbbookfile{theorem-library/nvs-act-laws}{theorem-library/nvs-act-laws}
\noindent{\small nvs-act-laws.scm -- THE GROUP AND SCALAR-ACTION LAWS OF A REAL NORMED VECTOR SPACE.  Since 2026-09-19 (rake batch 8, assignment 8-K1) this file is their ONE home: section (4) below holds the laws that three to five agents had each proved separately on 2026-09-19 under their own prefixes.  The NORM}\par
\vnbprop{nvs-scal-carr}{\vnbstmt{\forall m.\; (\operatorname{is\text{-}normed\text{-}vector\text{-}space}(m) \Rightarrow (\operatorname{carr}(\operatorname{scal}(m)) \simeq \mathbb{R}))}}
\vnbprop{nvs-scal-add}{\vnbstmt{\begin{array}{@{}l@{}} \forall m. \\ \quad \operatorname{is\text{-}normed\text{-}vector\text{-}space}(m) \Rightarrow  \\ \quad \quad (\operatorname{add}(\operatorname{scal}(m)) \simeq (\lambda \langle \operatorname{x\_}, \operatorname{y\_} \rangle \in (\mathbb{R} \times \mathbb{R}).\; (\operatorname{x\_} + \operatorname{y\_}))) \end{array}}}
\vnbprop{nvs-scal-mul}{\vnbstmt{\begin{array}{@{}l@{}} \forall m. \\ \quad \operatorname{is\text{-}normed\text{-}vector\text{-}space}(m) \Rightarrow  \\ \quad \quad (\operatorname{mul}(\operatorname{scal}(m)) \simeq (\lambda \langle \operatorname{x\_}, \operatorname{y\_} \rangle \in (\mathbb{R} \times \mathbb{R}).\; (\operatorname{x\_} \cdot \operatorname{y\_}))) \end{array}}}
\vnbprop{nvs-scal-zero}{\vnbstmt{\forall m.\; (\operatorname{is\text{-}normed\text{-}vector\text{-}space}(m) \Rightarrow (\operatorname{zero}(\operatorname{scal}(m)) \simeq 0))}}
\vnbprop{nvs-scal-one}{\vnbstmt{\forall m.\; (\operatorname{is\text{-}normed\text{-}vector\text{-}space}(m) \Rightarrow (\operatorname{one}(\operatorname{scal}(m)) \simeq 1))}}
\vnbprop{nvs-act-distrib-scalar}{\vnbstmt{\begin{array}{@{}l@{}} \forall m.; \forall r, s \in \operatorname{carr}(\operatorname{scal}(m)), x \in \operatorname{vec}(m). \\ \quad \operatorname{is\text{-}normed\text{-}vector\text{-}space}(m) \Rightarrow  \\ \quad \quad \operatorname{act}(m)(\operatorname{add}(\operatorname{scal}(m))(r, s), x) = \operatorname{vadd}(m)(\operatorname{act}(m)(r, x), \operatorname{act}(m)(s, x)) \end{array}}}
\vnbprop{nvs-act-mul-compat}{\vnbstmt{\begin{array}{@{}l@{}} \forall m.; \forall r, s \in \operatorname{carr}(\operatorname{scal}(m)), x \in \operatorname{vec}(m). \\ \quad \operatorname{is\text{-}normed\text{-}vector\text{-}space}(m) \Rightarrow  \\ \quad \quad \operatorname{act}(m)(\operatorname{mul}(\operatorname{scal}(m))(r, s), x) = \operatorname{act}(m)(r, \operatorname{act}(m)(s, x)) \end{array}}}
\vnbprop{nvs-act-unital}{\vnbstmt{\begin{array}{@{}l@{}} \forall m.; \forall x \in \operatorname{vec}(m). \\ \quad \operatorname{is\text{-}normed\text{-}vector\text{-}space}(m) \Rightarrow (\operatorname{act}(m)(\operatorname{one}(\operatorname{scal}(m)), x) = x) \end{array}}}
\vnbprop{nvs-vadd-assoc}{\vnbstmt{\begin{array}{@{}l@{}} \forall m.; \forall u, v, w \in \operatorname{vec}(m). \\ \quad \operatorname{is\text{-}normed\text{-}vector\text{-}space}(m) \Rightarrow  \\ \quad \quad \operatorname{vadd}(m)(\operatorname{vadd}(m)(u, v), w) = \operatorname{vadd}(m)(u, \operatorname{vadd}(m)(v, w)) \end{array}}}
\vnbprop{nvs-vzero-right}{\vnbstmt{\begin{array}{@{}l@{}} \forall m.; \forall u \in \operatorname{vec}(m). \\ \quad \operatorname{is\text{-}normed\text{-}vector\text{-}space}(m) \Rightarrow (\operatorname{vadd}(m)(u, \operatorname{vzero}(m)) = u) \end{array}}}
\vnbprop{nvs-act-type}{\vnbstmt{\begin{array}{@{}l@{}} \forall m.; \forall r \in \operatorname{carr}(\operatorname{scal}(m)), x \in \operatorname{vec}(m). \\ \quad \operatorname{is\text{-}normed\text{-}vector\text{-}space}(m) \Rightarrow (\operatorname{act}(m)(r, x) \in \operatorname{vec}(m)) \end{array}}}
\vnbprop{nvs-act-in-vec}{\vnbstmt{\begin{array}{@{}l@{}} \forall m.; \forall r \in \mathbb{R}, x \in \operatorname{vec}(m). \\ \quad \operatorname{is\text{-}normed\text{-}vector\text{-}space}(m) \Rightarrow (\operatorname{act}(m)(r, x) \in \operatorname{vec}(m)) \end{array}}}
\vnbprop{nvs-act-scale-assoc}{\vnbstmt{\begin{array}{@{}l@{}} \forall m.; \forall r, s \in \mathbb{R}, x \in \operatorname{vec}(m). \\ \quad \operatorname{is\text{-}normed\text{-}vector\text{-}space}(m) \Rightarrow  \\ \quad \quad \operatorname{act}(m)(r, \operatorname{act}(m)(s, x)) = \operatorname{act}(m)((r \cdot s), x) \end{array}}}
\vnbprop{nvs-act-collect}{\vnbstmt{\begin{array}{@{}l@{}} \forall m.; \forall r, s \in \mathbb{R}, x, y \in \operatorname{vec}(m). \\ \quad \operatorname{is\text{-}normed\text{-}vector\text{-}space}(m) \Rightarrow  \\ \quad \quad \operatorname{vadd}(m)(\operatorname{vadd}(m)(y, \operatorname{act}(m)(r, x)), \operatorname{act}(m)(s, x)) = \operatorname{vadd}(m)(y, \operatorname{act}(m)((r + s), x)) \end{array}}}
\vnbprop{nvs-act-zero}{\vnbstmt{\begin{array}{@{}l@{}} \forall m.; \forall x, y \in \operatorname{vec}(m). \\ \quad \operatorname{is\text{-}normed\text{-}vector\text{-}space}(m) \Rightarrow (\operatorname{vadd}(m)(y, \operatorname{act}(m)(0, x)) = y) \end{array}}}
\vnbprop{nvs-act-one}{\vnbstmt{\begin{array}{@{}l@{}} \forall m.; \forall x, y \in \operatorname{vec}(m). \\ \quad \operatorname{is\text{-}normed\text{-}vector\text{-}space}(m) \Rightarrow (\operatorname{vadd}(m)(y, \operatorname{act}(m)(1, x)) = \operatorname{vadd}(m)(y, x)) \end{array}}}
\vnbprop{nvs-scal-neg}{\vnbstmt{\begin{array}{@{}l@{}} \forall m. \\ \quad \operatorname{is\text{-}normed\text{-}vector\text{-}space}(m) \Rightarrow (\operatorname{neg}(\operatorname{scal}(m)) \simeq (\lambda x \in \mathbb{R}.\; -x)) \end{array}}}
\vnbprop{nvs-vzero-left}{\vnbstmt{\begin{array}{@{}l@{}} \forall m.; \forall u \in \operatorname{vec}(m). \\ \quad \operatorname{is\text{-}normed\text{-}vector\text{-}space}(m) \Rightarrow (\operatorname{vadd}(m)(\operatorname{vzero}(m), u) = u) \end{array}}}
\vnbprop{nvs-vadd-comm}{\vnbstmt{\begin{array}{@{}l@{}} \forall m.; \forall u, w \in \operatorname{vec}(m). \\ \quad \operatorname{is\text{-}normed\text{-}vector\text{-}space}(m) \Rightarrow (\operatorname{vadd}(m)(u, w) = \operatorname{vadd}(m)(w, u)) \end{array}}}
\vnbprop{nvs-act-distrib-vec}{\vnbstmt{\begin{array}{@{}l@{}} \forall m.; \forall r \in \operatorname{carr}(\operatorname{scal}(m)), x, y \in \operatorname{vec}(m). \\ \quad \operatorname{is\text{-}normed\text{-}vector\text{-}space}(m) \Rightarrow  \\ \quad \quad \operatorname{act}(m)(r, \operatorname{vadd}(m)(x, y)) = \operatorname{vadd}(m)(\operatorname{act}(m)(r, x), \operatorname{act}(m)(r, y)) \end{array}}}
\vnbprop{nvs-zero-act}{\vnbstmt{\begin{array}{@{}l@{}} \forall m.; \forall x \in \operatorname{vec}(m). \\ \quad \operatorname{is\text{-}normed\text{-}vector\text{-}space}(m) \Rightarrow (\operatorname{act}(m)(0, x) = \operatorname{vzero}(m)) \end{array}}}
\vnbprop{nvs-neg-one-act}{\vnbstmt{\begin{array}{@{}l@{}} \forall m.; \forall x \in \operatorname{vec}(m). \\ \quad \operatorname{is\text{-}normed\text{-}vector\text{-}space}(m) \Rightarrow (\operatorname{act}(m)(-1, x) = \operatorname{vneg}(m)(x)) \end{array}}}
\vnbprop{nvs-act-add-scalars}{\vnbstmt{\begin{array}{@{}l@{}} \forall m.; \forall r, s \in \mathbb{R}, x \in \operatorname{vec}(m). \\ \quad \operatorname{is\text{-}normed\text{-}vector\text{-}space}(m) \Rightarrow  \\ \quad \quad \operatorname{vadd}(m)(\operatorname{act}(m)(r, x), \operatorname{act}(m)(s, x)) = \operatorname{act}(m)((r + s), x) \end{array}}}
\vnbprop{nvs-act-vneg}{\vnbstmt{\begin{array}{@{}l@{}} \forall m.; \forall r \in \mathbb{R}, x \in \operatorname{vec}(m). \\ \quad \operatorname{is\text{-}normed\text{-}vector\text{-}space}(m) \Rightarrow  \\ \quad \quad \operatorname{vneg}(m)(\operatorname{act}(m)(r, x)) = \operatorname{act}(m)((-1 \cdot r), x) \end{array}}}
\vnbprop{nvs-vneg-vadd}{\vnbstmt{\begin{array}{@{}l@{}} \forall m.; \forall x, y \in \operatorname{vec}(m). \\ \quad \operatorname{is\text{-}normed\text{-}vector\text{-}space}(m) \Rightarrow  \\ \quad \quad \operatorname{vneg}(m)(\operatorname{vadd}(m)(x, y)) = \operatorname{vadd}(m)(\operatorname{vneg}(m)(x), \operatorname{vneg}(m)(y)) \end{array}}}
\vnbprop{nvs-vadd-shuffle}{\vnbstmt{\begin{array}{@{}l@{}} \forall m.; \forall a, b, c, d \in \operatorname{vec}(m). \\ \quad \operatorname{is\text{-}normed\text{-}vector\text{-}space}(m) \Rightarrow  \\ \quad \quad \operatorname{vadd}(m)(\operatorname{vadd}(m)(a, c), \operatorname{vadd}(m)(b, d)) = \operatorname{vadd}(m)(\operatorname{vadd}(m)(a, b), \operatorname{vadd}(m)(c, d)) \end{array}}}
\vnbprop{nvs-vzero-in-vec}{\vnbstmt{\forall m.\; (\operatorname{is\text{-}normed\text{-}vector\text{-}space}(m) \Rightarrow (\operatorname{vzero}(m) \in \operatorname{vec}(m)))}}
\vnbprop{nvs-vneg-in-vec}{\vnbstmt{\begin{array}{@{}l@{}} \forall m.; \forall x \in \operatorname{vec}(m). \\ \quad \operatorname{is\text{-}normed\text{-}vector\text{-}space}(m) \Rightarrow (\operatorname{vneg}(m)(x) \in \operatorname{vec}(m)) \end{array}}}
\vnbprop{nvs-vneg-left}{\vnbstmt{\begin{array}{@{}l@{}} \forall m.; \forall u \in \operatorname{vec}(m). \\ \quad \operatorname{is\text{-}normed\text{-}vector\text{-}space}(m) \Rightarrow (\operatorname{vadd}(m)(\operatorname{vneg}(m)(u), u) = \operatorname{vzero}(m)) \end{array}}}
\vnbprop{nvs-act-distrib-vec-rr}{\vnbstmt{\begin{array}{@{}l@{}} \forall m.; \forall r \in \mathbb{R}, x, y \in \operatorname{vec}(m). \\ \quad \operatorname{is\text{-}normed\text{-}vector\text{-}space}(m) \Rightarrow  \\ \quad \quad \operatorname{act}(m)(r, \operatorname{vadd}(m)(x, y)) = \operatorname{vadd}(m)(\operatorname{act}(m)(r, x), \operatorname{act}(m)(r, y)) \end{array}}}
\vnbprop{nvs-one-act}{\vnbstmt{\forall m.; \forall x \in \operatorname{vec}(m).\; (\operatorname{is\text{-}normed\text{-}vector\text{-}space}(m) \Rightarrow (\operatorname{act}(m)(1, x) = x))}}
\vnbprop{nvs-vadd-cancel}{\vnbstmt{\begin{array}{@{}l@{}} \forall m.; \forall x, z, w \in \operatorname{vec}(m). \\ \quad (\operatorname{is\text{-}normed\text{-}vector\text{-}space}(m) \wedge (\operatorname{vadd}(m)(x, z) = \operatorname{vadd}(m)(x, w))) \Rightarrow (z = w) \end{array}}}
\vnbprop{nvs-vadd-cancel-right}{\vnbstmt{\begin{array}{@{}l@{}} \forall m.; \forall x, z, w \in \operatorname{vec}(m). \\ \quad (\operatorname{is\text{-}normed\text{-}vector\text{-}space}(m) \wedge (\operatorname{vadd}(m)(z, x) = \operatorname{vadd}(m)(w, x))) \Rightarrow (z = w) \end{array}}}
\vnbprop{nvs-vneg-right}{\vnbstmt{\begin{array}{@{}l@{}} \forall m.; \forall u \in \operatorname{vec}(m). \\ \quad \operatorname{is\text{-}normed\text{-}vector\text{-}space}(m) \Rightarrow (\operatorname{vadd}(m)(u, \operatorname{vneg}(m)(u)) = \operatorname{vzero}(m)) \end{array}}}
\vnbprop{nvs-vadd-cancel-zero}{\vnbstmt{\begin{array}{@{}l@{}} \forall m.; \forall x, z \in \operatorname{vec}(m). \\ \quad (\operatorname{is\text{-}normed\text{-}vector\text{-}space}(m) \wedge (\operatorname{vadd}(m)(x, z) = x)) \Rightarrow (z = \operatorname{vzero}(m)) \end{array}}}
\vnbprop{nvs-neg-one-act-sub}{\vnbstmt{\begin{array}{@{}l@{}} \forall m.; \forall u \in \operatorname{vec}(m). \\ \quad \operatorname{is\text{-}normed\text{-}vector\text{-}space}(m) \Rightarrow (\operatorname{act}(m)(-1, u) = \operatorname{vneg}(m)(u)) \end{array}}}

\vnbbookfile{theorem-library/rake-norm-metrics}{theorem-library/rake-norm-metrics}
\noindent{\small rake-norm-metrics.scm -- rake batch 5 R: continuity projections and the norm metrics, PROVEN.}\par
\vnbprop{ag-norm-metric-is-metric-space}{\vnbstmt{\begin{array}{@{}l@{}} \forall g, m. \\ \quad (\operatorname{is\text{-}abelian\text{-}group}(g) \wedge \operatorname{is\text{-}group\text{-}norm}(m, \operatorname{opr}(g), \operatorname{inv}(g), \operatorname{iden}(g), \operatorname{carr}(g))) \Rightarrow  \\ \quad \quad \operatorname{is\text{-}metric\text{-}space}(\langle \operatorname{carr}(g), (\lambda \langle u, v \rangle \in (\operatorname{carr}(g) \times \operatorname{carr}(g)).\; m(\operatorname{opr}(g)(u, \operatorname{inv}(g)(v)))) \rangle) \end{array}}}
\vnbprop{continuous-is-continuous-at}{\vnbstmt{\begin{array}{@{}l@{}} \forall s, t, f.; \forall a \in \operatorname{pts}(s). \\ \quad \operatorname{is\text{-}continuous}(s, t, f) \Rightarrow \operatorname{is\text{-}continuous\text{-}at}(s, t, f, a) \end{array}}}
\vnbprop{uniformly-continuous-is-continuous}{\vnbstmt{\begin{array}{@{}l@{}} \forall s, t, f. \\ \quad \operatorname{is\text{-}uniformly\text{-}continuous}(s, t, f) \Rightarrow \operatorname{is\text{-}continuous}(s, t, f) \end{array}}}
\vnbprop{nag-metric-carrier}{\vnbstmt{\forall g.\; (\operatorname{is\text{-}normed\text{-}ag}(g) \Rightarrow (\operatorname{pts}(\operatorname{nag\text{-}metric\text{-}space}(g)) \simeq \operatorname{carr}(g)))}}
\vnbprop{nag-metric-distance}{\vnbstmt{\begin{array}{@{}l@{}} \forall g.; \forall u, v \in \operatorname{carr}(g). \\ \quad \operatorname{is\text{-}normed\text{-}ag}(g) \Rightarrow  \\ \quad \quad \operatorname{dist}(\operatorname{nag\text{-}metric\text{-}space}(g))(u, v) = \operatorname{nrm}(g)(\operatorname{opr}(g)(u, \operatorname{inv}(g)(v))) \end{array}}}
\vnbprop{nag-metric-space-is-metric-space}{\vnbstmt{\forall g.\; (\operatorname{is\text{-}normed\text{-}ag}(g) \Rightarrow \operatorname{is\text{-}metric\text{-}space}(\operatorname{nag\text{-}metric\text{-}space}(g)))}}
\vnbprop{nvs-metric-is-ms}{\vnbstmt{\begin{array}{@{}l@{}} \forall m. \\ \quad \operatorname{is\text{-}normed\text{-}vector\text{-}space}(m) \Rightarrow \operatorname{is\text{-}metric\text{-}space}(\operatorname{nvs\text{-}metric\text{-}space}(m)) \end{array}}}

\vnbbookfile{theorem-library/nvs-norm-laws}{theorem-library/nvs-norm-laws}
\noindent{\small nvs-norm-laws.scm -- THE NORM LAWS OF A REAL NORMED VECTOR SPACE, in ONE place.  Rake batch 8, assignment 8-K1.}\par
\vnbprop{nvs-vnrm-triangle}{\vnbstmt{\begin{array}{@{}l@{}} \forall m.; \forall x, y \in \operatorname{vec}(m). \\ \quad \operatorname{is\text{-}normed\text{-}vector\text{-}space}(m) \Rightarrow  \\ \quad \quad \operatorname{vnrm}(m)(\operatorname{vadd}(m)(x, y)) \leq (\operatorname{vnrm}(m)(x) + \operatorname{vnrm}(m)(y)) \end{array}}}
\vnbprop{nvs-vnrm-homog-c}{\vnbstmt{\begin{array}{@{}l@{}} \forall m.; \forall r \in \operatorname{carr}(\operatorname{scal}(m)), x \in \operatorname{vec}(m). \\ \quad \operatorname{is\text{-}normed\text{-}vector\text{-}space}(m) \Rightarrow  \\ \quad \quad \operatorname{vnrm}(m)(\operatorname{act}(m)(r, x)) = (\lvert r \rvert \cdot \operatorname{vnrm}(m)(x)) \end{array}}}
\vnbprop{nvs-vnrm-homog}{\vnbstmt{\begin{array}{@{}l@{}} \forall m.; \forall r \in \mathbb{R}, x \in \operatorname{vec}(m). \\ \quad \operatorname{is\text{-}normed\text{-}vector\text{-}space}(m) \Rightarrow  \\ \quad \quad \operatorname{vnrm}(m)(\operatorname{act}(m)(r, x)) = (\lvert r \rvert \cdot \operatorname{vnrm}(m)(x)) \end{array}}}
\vnbprop{nvs-vnrm-vzero}{\vnbstmt{\forall m.\; (\operatorname{is\text{-}normed\text{-}vector\text{-}space}(m) \Rightarrow (\operatorname{vnrm}(m)(\operatorname{vzero}(m)) = 0))}}
\vnbprop{nvs-ms-pts}{\vnbstmt{\begin{array}{@{}l@{}} \forall m. \\ \quad \operatorname{is\text{-}metric\text{-}space}(\operatorname{nvs\text{-}metric\text{-}space}(m)) \Rightarrow (\operatorname{pts}(\operatorname{nvs\text{-}metric\text{-}space}(m)) \simeq \operatorname{vec}(m)) \end{array}}}
\vnbprop{nvs-metric-distance}{\vnbstmt{\begin{array}{@{}l@{}} \forall m.; \forall u, v \in \operatorname{vec}(m). \\ \quad \operatorname{is\text{-}normed\text{-}vector\text{-}space}(m) \Rightarrow  \\ \quad \quad \operatorname{dist}(\operatorname{nvs\text{-}metric\text{-}space}(m))(u, v) = \operatorname{vnrm}(m)(\operatorname{vadd}(m)(u, \operatorname{vneg}(m)(v))) \end{array}}}
\vnbprop{nvs-metric-distance-q}{\vnbstmt{\begin{array}{@{}l@{}} \forall m.; \forall u, v \in \operatorname{vec}(m). \\ \quad \operatorname{is\text{-}normed\text{-}vector\text{-}space}(m) \Rightarrow  \\ \quad \quad (\operatorname{dist}(\operatorname{nvs\text{-}metric\text{-}space}(m))(u, v) \simeq \operatorname{vnrm}(m)(\operatorname{vadd}(m)(u, \operatorname{vneg}(m)(v)))) \end{array}}}
\vnbprop{nvs-dist-shift}{\vnbstmt{\begin{array}{@{}l@{}} \forall m.; \forall u, w \in \operatorname{vec}(m). \\ \quad \operatorname{is\text{-}normed\text{-}vector\text{-}space}(m) \Rightarrow  \\ \quad \quad \operatorname{dist}(\operatorname{nvs\text{-}metric\text{-}space}(m))(u, \operatorname{vadd}(m)(u, w)) = \operatorname{vnrm}(m)(w) \end{array}}}

\vnbbookfile{theorem-library/rake-hb-submodules}{theorem-library/rake-hb-submodules}
\noindent{\small rake-hb-submodules.scm -- the two SUBMODULE leaves of the Hahn-Banach cluster, PROVEN, and the normed-vector-space action bricks they need.}\par
\vnbprop{line-unfold}{\vnbstmt{\begin{array}{@{}l@{}} \forall m, v. \\ \quad (\operatorname{line}(m, v) \simeq \{y \in \operatorname{vec}(m) : \exists r \in \mathbb{R}.\; (y = \operatorname{act}(m)(r, v))\}) \end{array}}}
\vnbprop{line-is-submodule}{\vnbstmt{\begin{array}{@{}l@{}} \forall m, v. \\ \quad (\operatorname{is\text{-}normed\text{-}vector\text{-}space}(m) \wedge (v \in \operatorname{vec}(m))) \Rightarrow \operatorname{is\text{-}submodule}(m, \operatorname{line}(m, v)) \end{array}}}
\vnbprop{span-add-one-submodule}{\vnbstmt{\begin{array}{@{}l@{}} \forall m, t, v. \\ \quad (\operatorname{is\text{-}normed\text{-}vector\text{-}space}(m) \wedge (\operatorname{is\text{-}submodule}(m, t) \wedge (v \in \operatorname{vec}(m)))) \Rightarrow  \\ \quad \quad \operatorname{is\text{-}submodule}(m, \operatorname{span\text{-}add\text{-}one}(m, t, v)) \end{array}}}

\vnbbookfile{theorem-library/rake-span-add-one-guarded}{theorem-library/rake-span-add-one-guarded}
\noindent{\small rake-span-add-one-guarded.scm -- the GUARDED forms of span-add-one-superset and span-add-one-has-v, PROVEN.}\par
\vnbprop{span-add-one-superset-nvs}{\vnbstmt{\begin{array}{@{}l@{}} \forall m, t, v. \\ \quad (\operatorname{is\text{-}normed\text{-}vector\text{-}space}(m) \wedge (\operatorname{is\text{-}submodule}(m, t) \wedge (v \in \operatorname{vec}(m)))) \Rightarrow  \\ \quad \quad t \subseteq \operatorname{span\text{-}add\text{-}one}(m, t, v) \end{array}}}
\vnbprop{span-add-one-has-v-nvs}{\vnbstmt{\begin{array}{@{}l@{}} \forall m, t, v. \\ \quad (\operatorname{is\text{-}normed\text{-}vector\text{-}space}(m) \wedge (\operatorname{is\text{-}submodule}(m, t) \wedge (v \in \operatorname{vec}(m)))) \Rightarrow  \\ \quad \quad v \in \operatorname{span\text{-}add\text{-}one}(m, t, v) \end{array}}}

\vnbbookfile{theorem-library/rake-hb-leaves-2}{theorem-library/rake-hb-leaves-2}
\noindent{\small theorem-library/rake-hb-leaves-2.scm -- three Hahn-Banach leaves, PROVEN modulo 0: good-sub-submodule, good-sub-in-power, line-has-v.}\par
\vnbprop{good-sub-submodule}{\vnbstmt{\forall m, s, f, t.\; (\operatorname{good\text{-}sub}(m, s, f, t) \Rightarrow \operatorname{is\text{-}submodule}(m, t))}}
\vnbprop{good-sub-in-power}{\vnbstmt{\begin{array}{@{}l@{}} \forall m, s, f, t. \\ \quad (\operatorname{is\text{-}normed\text{-}vector\text{-}space}(m) \wedge \operatorname{good\text{-}sub}(m, s, f, t)) \Rightarrow (t \in \operatorname{power}(\operatorname{vec}(m))) \end{array}}}
\vnbprop{line-has-v}{\vnbstmt{\begin{array}{@{}l@{}} \forall m, v. \\ \quad (\operatorname{is\text{-}normed\text{-}vector\text{-}space}(m) \wedge (v \in \operatorname{vec}(m))) \Rightarrow (v \in \operatorname{line}(m, v)) \end{array}}}

\vnbbookfile{theorem-library/rake-hb-gap}{theorem-library/rake-hb-gap}
\noindent{\small rake-hb-gap.scm -- hb-gap (theorem-library/hahn-banach-proof.scm:24), the analytic core of the one-dimension Hahn-Banach step, PROVEN.}\par
\vnbprop{hbg-scale-pos}{\vnbstmt{\begin{array}{@{}l@{}} \forall d, t, u, a, o \in \mathbb{R}. \\ \quad ((0 \leq o) \wedge (((-u - (d \cdot t)) \leq a) \wedge (a \leq ((d \cdot t) - u)))) \Rightarrow  \\ \quad \quad \lvert ((o \cdot u) + (o \cdot a)) \rvert \leq (d \cdot (o \cdot t)) \end{array}}}
\vnbprop{hbg-scale-neg}{\vnbstmt{\begin{array}{@{}l@{}} \forall d, t, u, a, o \in \mathbb{R}. \\ \quad ((o \leq 0) \wedge (((-u - (d \cdot t)) \leq a) \wedge (a \leq ((d \cdot t) - u)))) \Rightarrow  \\ \quad \quad \lvert ((o \cdot u) + (o \cdot a)) \rvert \leq (d \cdot (-o \cdot t)) \end{array}}}
\vnbprop{hbg-neg-one-cancel}{\vnbstmt{\begin{array}{@{}l@{}} \forall m.; \forall y, v \in \operatorname{vec}(m). \\ \quad \operatorname{is\text{-}normed\text{-}vector\text{-}space}(m) \Rightarrow  \\ \quad \quad \operatorname{vadd}(m)(\operatorname{vadd}(m)(y, v), \operatorname{act}(m)(-1, v)) = y \end{array}}}
\vnbprop{hbg-shift}{\vnbstmt{\begin{array}{@{}l@{}} \forall m.; \forall w, z, v \in \operatorname{vec}(m). \\ \quad \operatorname{is\text{-}normed\text{-}vector\text{-}space}(m) \Rightarrow  \\ \quad \quad \operatorname{vadd}(m)(w, \operatorname{act}(m)(-1, z)) = \operatorname{vadd}(m)(\operatorname{vadd}(m)(w, v), \operatorname{act}(m)(-1, \operatorname{vadd}(m)(z, v))) \end{array}}}
\vnbprop{hbg-diff-norm}{\vnbstmt{\begin{array}{@{}l@{}} \forall m.; \forall w, z, v \in \operatorname{vec}(m). \\ \quad \operatorname{is\text{-}normed\text{-}vector\text{-}space}(m) \Rightarrow  \\ \quad \quad \operatorname{vnrm}(m)(\operatorname{vadd}(m)(w, \operatorname{act}(m)(-1, z))) \leq (\operatorname{vnrm}(m)(\operatorname{vadd}(m)(w, v)) + \operatorname{vnrm}(m)(\operatorname{vadd}(m)(z, v))) \end{array}}}
\vnbprop{hbg-scale-back}{\vnbstmt{\begin{array}{@{}l@{}} \forall m.; \forall y \in \operatorname{vec}(m), r \in \mathbb{R}. \\ \quad (\operatorname{is\text{-}normed\text{-}vector\text{-}space}(m) \wedge \neg (r = 0)) \Rightarrow  \\ \quad \quad \operatorname{act}(m)(r, \operatorname{act}(m)(\operatorname{recip}(r), y)) = y \end{array}}}
\vnbprop{hbg-scale-decomp}{\vnbstmt{\begin{array}{@{}l@{}} \forall m.; \forall y, v \in \operatorname{vec}(m), r \in \mathbb{R}. \\ \quad (\operatorname{is\text{-}normed\text{-}vector\text{-}space}(m) \wedge \neg (r = 0)) \Rightarrow  \\ \quad \quad \operatorname{vadd}(m)(y, \operatorname{act}(m)(r, v)) = \operatorname{act}(m)(r, \operatorname{vadd}(m)(\operatorname{act}(m)(\operatorname{recip}(r), y), v)) \end{array}}}
\vnbprop{hbg-linfun-app-real}{\vnbstmt{\begin{array}{@{}l@{}} \forall m, s, f, x. \\ \quad (\operatorname{is\text{-}linear\text{-}functional\text{-}on}(m, s, f) \wedge (x \in s)) \Rightarrow (f(x) \in \mathbb{R}) \end{array}}}
\vnbprop{hbg-sep}{\vnbstmt{\begin{array}{@{}l@{}} \forall m, s, f, v, c.; \forall z, w \in s. \\ \quad \operatorname{is\text{-}normed\text{-}vector\text{-}space}(m) \wedge \\ \quad \operatorname{is\text{-}submodule}(m, s) \wedge \\ \quad \operatorname{is\text{-}linear\text{-}functional\text{-}on}(m, s, f) \wedge \\ \quad v \in \operatorname{vec}(m) \wedge \\ \quad c \in \mathbb{R} \wedge \\ \quad \forall x \in s.\; (\lvert f(x) \rvert \leq (c \cdot \operatorname{vnrm}(m)(x))) \wedge \\ \quad 0 \leq c \Rightarrow \\ \quad \quad (-f(z) - (c \cdot \operatorname{vnrm}(m)(\operatorname{vadd}(m)(z, v)))) \leq ((c \cdot \operatorname{vnrm}(m)(\operatorname{vadd}(m)(w, v))) - f(w)) \end{array}}}
\vnbprop{hbg-core}{\vnbstmt{\begin{array}{@{}l@{}} \forall m, s, f, v, c, a.; \forall y \in s, r \in \mathbb{R}. \\ \quad \operatorname{is\text{-}normed\text{-}vector\text{-}space}(m) \wedge \\ \quad \operatorname{is\text{-}submodule}(m, s) \wedge \\ \quad \operatorname{is\text{-}linear\text{-}functional\text{-}on}(m, s, f) \wedge \\ \quad v \in \operatorname{vec}(m) \wedge \\ \quad c \in \mathbb{R} \wedge \\ \quad a \in \mathbb{R} \wedge \\ \quad \forall x \in s.\; (\lvert f(x) \rvert \leq (c \cdot \operatorname{vnrm}(m)(x))) \wedge \\ \quad \forall u \in s.\; ((-f(u) - (c \cdot \operatorname{vnrm}(m)(\operatorname{vadd}(m)(u, v)))) \leq a) \wedge \\ \quad \forall u \in s.\; (a \leq ((c \cdot \operatorname{vnrm}(m)(\operatorname{vadd}(m)(u, v))) - f(u))) \Rightarrow \\ \quad \quad \lvert (f(y) + (r \cdot a)) \rvert \leq (c \cdot \operatorname{vnrm}(m)(\operatorname{vadd}(m)(y, \operatorname{act}(m)(r, v)))) \end{array}}}
\vnbprop{hb-gap-bound}{\vnbstmt{\begin{array}{@{}l@{}} \forall m, s, f, v, c. \\ \quad \operatorname{is\text{-}normed\text{-}vector\text{-}space}(m) \wedge \\ \quad \operatorname{is\text{-}submodule}(m, s) \wedge \\ \quad \operatorname{is\text{-}linear\text{-}functional\text{-}on}(m, s, f) \wedge \\ \quad v \in \operatorname{vec}(m) \wedge \\ \quad c \in \mathbb{R} \wedge \\ \quad \forall x \in s.\; (\lvert f(x) \rvert \leq (c \cdot \operatorname{vnrm}(m)(x))) \wedge \\ \quad 0 \leq c \Rightarrow \\ \quad \quad \exists a \in \mathbb{R}. \forall y \in s, r \in \mathbb{R}. \\ \quad \quad \quad \lvert (f(y) + (r \cdot a)) \rvert \leq (c \cdot \operatorname{vnrm}(m)(\operatorname{vadd}(m)(y, \operatorname{act}(m)(r, v)))) \end{array}}}
\vnbprop{hb-gap}{\vnbstmt{\begin{array}{@{}l@{}} \forall m, s, f, v. \\ \quad \operatorname{is\text{-}normed\text{-}vector\text{-}space}(m) \wedge \\ \quad \operatorname{is\text{-}submodule}(m, s) \wedge \\ \quad \operatorname{is\text{-}bounded\text{-}linear\text{-}functional\text{-}on}(m, s, f) \wedge \\ \quad v \in \operatorname{vec}(m) \wedge \\ \quad \neg (v \in s) \Rightarrow \\ \quad \quad \exists a \in \mathbb{R}. \forall y \in s, r \in \mathbb{R}. \\ \quad \quad \quad \lvert (f(y) + (r \cdot a)) \rvert \leq (\operatorname{dual\text{-}norm\text{-}on}(m, s, f) \cdot \operatorname{vnrm}(m)(\operatorname{vadd}(m)(y, \operatorname{act}(m)(r, v)))) \end{array}}}

\vnbbookfile{theorem-library/rake-hb-extend-construct}{theorem-library/rake-hb-extend-construct}
\noindent{\small rake-hb-extend-construct.scm -- hb-extend-construct (theorem-library/ hahn-banach-proof.scm:44), rake batch 7, assignment 7-C, second leaf.}\par
\vnbprop{hbx-s-in-vec}{\vnbstmt{\begin{array}{@{}l@{}} \forall m, s, p. \\ \quad (\operatorname{is\text{-}normed\text{-}vector\text{-}space}(m) \wedge (\operatorname{is\text{-}submodule}(m, s) \wedge (p \in s))) \Rightarrow (p \in \operatorname{vec}(m)) \end{array}}}
\vnbprop{hbx-span-in-vec}{\vnbstmt{\forall m, s, v.; \forall w \in \operatorname{span\text{-}add\text{-}one}(m, s, v).\; (w \in \operatorname{vec}(m))}}
\vnbprop{hbx-span-decomp}{\vnbstmt{\begin{array}{@{}l@{}} \forall m, s, v.; \forall w \in \operatorname{span\text{-}add\text{-}one}(m, s, v). \exists y \in s, r \in \mathbb{R}. \\ \quad w = \operatorname{vadd}(m)(y, \operatorname{act}(m)(r, v)) \end{array}}}
\vnbprop{hbx-span-mem-eq}{\vnbstmt{\begin{array}{@{}l@{}} \forall m, s, v, w, p, c. \\ \quad \operatorname{is\text{-}normed\text{-}vector\text{-}space}(m) \wedge \\ \quad \operatorname{is\text{-}submodule}(m, s) \wedge \\ \quad v \in \operatorname{vec}(m) \wedge \\ \quad p \in s \wedge \\ \quad c \in \mathbb{R} \wedge \\ \quad w = \operatorname{vadd}(m)(p, \operatorname{act}(m)(c, v)) \Rightarrow \\ \quad \quad w \in \operatorname{span\text{-}add\text{-}one}(m, s, v) \end{array}}}
\vnbprop{hbx-decomp-unique}{\vnbstmt{\begin{array}{@{}l@{}} \forall m, s, v, p_{1}, c_{1}, p_{2}, c_{2}. \\ \quad \operatorname{is\text{-}normed\text{-}vector\text{-}space}(m) \wedge \\ \quad \operatorname{is\text{-}submodule}(m, s) \wedge \\ \quad v \in \operatorname{vec}(m) \wedge \\ \quad \neg (v \in s) \wedge \\ \quad p_{1} \in s \wedge \\ \quad c_{1} \in \mathbb{R} \wedge \\ \quad p_{2} \in s \wedge \\ \quad c_{2} \in \mathbb{R} \wedge \\ \quad \operatorname{vadd}(m)(p_{1}, \operatorname{act}(m)(c_{1}, v)) = \operatorname{vadd}(m)(p_{2}, \operatorname{act}(m)(c_{2}, v)) \Rightarrow \\ \quad \quad (c_{1} = c_{2}) \wedge (p_{1} = p_{2}) \end{array}}}
\vnbprop{hbx-vadd-exchange}{\vnbstmt{\begin{array}{@{}l@{}} \forall m, a, b, c, d. \\ \quad (\operatorname{is\text{-}normed\text{-}vector\text{-}space}(m) \wedge ((a \in \operatorname{vec}(m)) \wedge ((b \in \operatorname{vec}(m)) \wedge ((c \in \operatorname{vec}(m)) \wedge (d \in \operatorname{vec}(m)))))) \Rightarrow  \\ \quad \quad \operatorname{vadd}(m)(\operatorname{vadd}(m)(a, b), \operatorname{vadd}(m)(c, d)) = \operatorname{vadd}(m)(\operatorname{vadd}(m)(a, c), \operatorname{vadd}(m)(b, d)) \end{array}}}
\vnbprop{hbx-add-decomp}{\vnbstmt{\begin{array}{@{}l@{}} \forall m, v, y_{1}, c_{1}, y_{2}, c_{2}. \\ \quad \operatorname{is\text{-}normed\text{-}vector\text{-}space}(m) \wedge \\ \quad v \in \operatorname{vec}(m) \wedge \\ \quad y_{1} \in \operatorname{vec}(m) \wedge \\ \quad c_{1} \in \mathbb{R} \wedge \\ \quad y_{2} \in \operatorname{vec}(m) \wedge \\ \quad c_{2} \in \mathbb{R} \Rightarrow \\ \quad \quad \operatorname{vadd}(m)(\operatorname{vadd}(m)(y_{1}, \operatorname{act}(m)(c_{1}, v)), \operatorname{vadd}(m)(y_{2}, \operatorname{act}(m)(c_{2}, v))) \\ \quad \quad =\; \operatorname{vadd}(m)(\operatorname{vadd}(m)(y_{1}, y_{2}), \operatorname{act}(m)((c_{1} + c_{2}), v)) \end{array}}}
\vnbprop{hbx-f-add}{\vnbstmt{\begin{array}{@{}l@{}} \forall m, s, f, x, y. \\ \quad (\operatorname{is\text{-}linear\text{-}functional\text{-}on}(m, s, f) \wedge ((x \in s) \wedge (y \in s))) \Rightarrow  \\ \quad \quad f(\operatorname{vadd}(m)(x, y)) = (f(x) + f(y)) \end{array}}}
\vnbprop{hbx-f-act}{\vnbstmt{\begin{array}{@{}l@{}} \forall m, s, f, r, x. \\ \quad (\operatorname{is\text{-}linear\text{-}functional\text{-}on}(m, s, f) \wedge ((r \in \mathbb{R}) \wedge (x \in s))) \Rightarrow  \\ \quad \quad f(\operatorname{act}(m)(r, x)) = (r \cdot f(x)) \end{array}}}
\vnbprop{hbx-value}{\vnbstmt{\begin{array}{@{}l@{}} \forall m, s, f, v, a, u, p, c. \\ \quad \operatorname{is\text{-}normed\text{-}vector\text{-}space}(m) \wedge \\ \quad \operatorname{is\text{-}submodule}(m, s) \wedge \\ \quad \operatorname{is\text{-}linear\text{-}functional\text{-}on}(m, s, f) \wedge \\ \quad v \in \operatorname{vec}(m) \wedge \\ \quad \neg (v \in s) \wedge \\ \quad a \in \mathbb{R} \wedge \\ \quad p \in s \wedge \\ \quad c \in \mathbb{R} \wedge \\ \quad u = \operatorname{vadd}(m)(p, \operatorname{act}(m)(c, v)) \Rightarrow \\ \quad \quad (\iota\, \mathit{tv}.\;((\mathit{tv} \in \mathbb{R}) \wedge \exists \mathit{pv} \in s, \mathit{cv} \in \mathbb{R}.\; ((u = \operatorname{vadd}(m)(\mathit{pv}, \operatorname{act}(m)(\mathit{cv}, v))) \wedge (\mathit{tv} = (f(\mathit{pv}) + (\mathit{cv} \cdot a)))))) \\ \quad \quad =\; f(p) + (c \cdot a) \end{array}}}
\vnbprop{hb-extend-construct}{\vnbstmt{\begin{array}{@{}l@{}} \forall m, s, f, v, a. \\ \quad \operatorname{is\text{-}normed\text{-}vector\text{-}space}(m) \wedge \\ \quad \operatorname{is\text{-}submodule}(m, s) \wedge \\ \quad \operatorname{is\text{-}linear\text{-}functional\text{-}on}(m, s, f) \wedge \\ \quad v \in \operatorname{vec}(m) \wedge \\ \quad \neg (v \in s) \wedge \\ \quad a \in \mathbb{R} \Rightarrow \\ \quad \quad \exists g. \\ \quad \quad \quad \operatorname{is\text{-}linear\text{-}functional\text{-}on}(m, \operatorname{span\text{-}add\text{-}one}(m, s, v), g) \wedge \\ \quad \quad \quad \operatorname{extends\text{-}on}(s, g, f) \wedge \\ \quad \quad \quad \forall y \in s, r \in \mathbb{R}.\; (g(\operatorname{vadd}(m)(y, \operatorname{act}(m)(r, v))) = (f(y) + (r \cdot a))) \end{array}}}

\vnbbookfile{theorem-library/rake-line-functional}{theorem-library/rake-line-functional}
\noindent{\small rake-line-functional.scm -- line-functional-exists (theorem-library/ norm-as-sup-proof.scm:43), rake batch 7, assignment 7-C.}\par
\vnbprop{lfe-line-membership}{\vnbstmt{\begin{array}{@{}l@{}} \forall m, v, w. \\ \quad (w \in \operatorname{line}(m, v)) \Leftrightarrow ((w \in \operatorname{vec}(m)) \wedge \exists r \in \mathbb{R}.\; (w = \operatorname{act}(m)(r, v))) \end{array}}}
\vnbprop{lfe-line-in-vec}{\vnbstmt{\forall m, v.; \forall w \in \operatorname{line}(m, v).\; (w \in \operatorname{vec}(m))}}
\vnbprop{lfe-line-subset}{\vnbstmt{\forall m, v.\; (\operatorname{is\text{-}normed\text{-}vector\text{-}space}(m) \Rightarrow (\operatorname{line}(m, v) \subseteq \operatorname{vec}(m)))}}
\vnbprop{lfe-line-coef}{\vnbstmt{\forall m, v.; \forall w \in \operatorname{line}(m, v).\; \exists r \in \mathbb{R}.\; (w = \operatorname{act}(m)(r, v))}}
\vnbprop{lfe-line-mem-eq}{\vnbstmt{\begin{array}{@{}l@{}} \forall m, v, w, r. \\ \quad (\operatorname{is\text{-}normed\text{-}vector\text{-}space}(m) \wedge ((v \in \operatorname{vec}(m)) \wedge ((r \in \mathbb{R}) \wedge (w = \operatorname{act}(m)(r, v))))) \Rightarrow  \\ \quad \quad w \in \operatorname{line}(m, v) \end{array}}}
\vnbprop{lfe-coef-unique}{\vnbstmt{\begin{array}{@{}l@{}} \forall m, v.; \forall a, b \in \mathbb{R}. \\ \quad (\operatorname{is\text{-}normed\text{-}vector\text{-}space}(m) \wedge ((v \in \operatorname{vec}(m)) \wedge (\operatorname{act}(m)(a, v) = \operatorname{act}(m)(b, v)))) \Rightarrow  \\ \quad \quad (a \cdot \operatorname{vnrm}(m)(v)) = (b \cdot \operatorname{vnrm}(m)(v)) \end{array}}}
\vnbprop{lfe-coef-value}{\vnbstmt{\begin{array}{@{}l@{}} \forall m, v, u, c. \\ \quad (\operatorname{is\text{-}normed\text{-}vector\text{-}space}(m) \wedge ((v \in \operatorname{vec}(m)) \wedge ((c \in \mathbb{R}) \wedge (u = \operatorname{act}(m)(c, v))))) \Rightarrow  \\ \quad \quad (\iota\, t.\;((t \in \mathbb{R}) \wedge \exists q \in \mathbb{R}.\; ((u = \operatorname{act}(m)(q, v)) \wedge (t = (q \cdot \operatorname{vnrm}(m)(v)))))) \\ \quad \quad =\; c \cdot \operatorname{vnrm}(m)(v) \end{array}}}
\vnbprop{line-functional-exists}{\vnbstmt{\begin{array}{@{}l@{}} \forall m, v. \\ \quad (\operatorname{is\text{-}normed\text{-}vector\text{-}space}(m) \wedge (v \in \operatorname{vec}(m))) \Rightarrow  \\ \quad \quad \exists f. \\ \quad \quad \quad \operatorname{is\text{-}bounded\text{-}linear\text{-}functional\text{-}on}(m, \operatorname{line}(m, v), f) \wedge \\ \quad \quad \quad \operatorname{dual\text{-}norm\text{-}on}(m, \operatorname{line}(m, v), f) \leq 1 \wedge \\ \quad \quad \quad f(v) = \operatorname{vnrm}(m)(v) \end{array}}}

\vnbbookfile{theorem-library/hahn-banach-proof}{theorem-library/hahn-banach-proof}
\noindent{\small hahn-banach-proof.scm -- Hahn-Banach, the one-dimension-higher extension step (real normed vector space), MACHINE-PROVEN modulo the warranted core.}\par
\vnbprop{hahn-banach-extend-one}{\vnbstmt{\begin{array}{@{}l@{}} \forall m, s, f, v. \\ \quad \operatorname{is\text{-}normed\text{-}vector\text{-}space}(m) \wedge \\ \quad \operatorname{is\text{-}submodule}(m, s) \wedge \\ \quad \operatorname{is\text{-}bounded\text{-}linear\text{-}functional\text{-}on}(m, s, f) \wedge \\ \quad v \in \operatorname{vec}(m) \wedge \\ \quad \neg (v \in s) \Rightarrow \\ \quad \quad \exists g. \\ \quad \quad \quad \operatorname{is\text{-}linear\text{-}functional\text{-}on}(m, \operatorname{span\text{-}add\text{-}one}(m, s, v), g) \wedge \\ \quad \quad \quad \operatorname{extends\text{-}on}(s, g, f) \wedge \\ \quad \quad \quad \forall w \in \operatorname{span\text{-}add\text{-}one}(m, s, v). \\ \quad \quad \quad \quad \lvert g(w) \rvert \leq (\operatorname{dual\text{-}norm\text{-}on}(m, s, f) \cdot \operatorname{vnrm}(m)(w)) \end{array}}}

\vnbbookfile{theorem-library/noetherian-maximal-proof}{theorem-library/noetherian-maximal-proof}
\noindent{\small noetherian-maximal-proof.scm -- the noetherian maximal-element principle, MACHINE-PROVEN, and its specialisation hb-good-has-maximal (the input that makes finite-dimensional Hahn-Banach terminate).}\par
\vnbprop{noetherian-set-has-maximal}{\vnbstmt{\begin{array}{@{}l@{}} \forall m.; \forall g \in \mathbf{Set}. \\ \quad \operatorname{is\text{-}noetherian}(m) \wedge \\ \quad g \subseteq \operatorname{power}(\operatorname{vec}(m)) \wedge \\ \quad \forall t \in g.\; \operatorname{is\text{-}submodule}(m, t) \wedge \\ \quad \exists a.\; (a \in g) \Rightarrow \\ \quad \quad \exists t \in g.\; \forall u \in g.\; ((t \subseteq u) \Rightarrow (t = u)) \end{array}}}
\vnbprop{hb-good-has-maximal}{\vnbstmt{\begin{array}{@{}l@{}} \forall m, s, f. \\ \quad \operatorname{is\text{-}normed\text{-}vector\text{-}space}(m) \wedge \\ \quad \operatorname{is\text{-}finite\text{-}dimensional}(\operatorname{normed\text{-}vector\text{-}space\text{-}as\text{-}module}(m)) \wedge \\ \quad \operatorname{is\text{-}submodule}(m, s) \wedge \\ \quad \operatorname{is\text{-}bounded\text{-}linear\text{-}functional\text{-}on}(m, s, f) \Rightarrow \\ \quad \quad \exists t. \\ \quad \quad \quad (\operatorname{good\text{-}sub}(m, s, f, t) \wedge \forall u.\; ((\operatorname{good\text{-}sub}(m, s, f, u) \wedge (t \subseteq u)) \Rightarrow (t = u))) \end{array}}}

\vnbbookfile{theorem-library/hahn-banach-full-proof}{theorem-library/hahn-banach-full-proof}
\noindent{\small RETIRED 2026-09-17 (proven): le-bound-mono (was an add-to-pss here) -- theorem-library/rake-analysis2.scm hahn-banach-full-proof.scm -- FULL Hahn-Banach for a finite-dimensional real normed vector space, MACHINE-PROVEN modulo a warranted plumbing core.}\par
\vnbprop{good-step}{\vnbstmt{\begin{array}{@{}l@{}} \forall m, s, f, t, x. \\ \quad \operatorname{is\text{-}normed\text{-}vector\text{-}space}(m) \wedge \\ \quad \operatorname{is\text{-}bounded\text{-}linear\text{-}functional\text{-}on}(m, s, f) \wedge \\ \quad \operatorname{good\text{-}sub}(m, s, f, t) \wedge \\ \quad x \in \operatorname{vec}(m) \wedge \\ \quad \neg (x \in t) \Rightarrow \\ \quad \quad \operatorname{good\text{-}sub}(m, s, f, \operatorname{span\text{-}add\text{-}one}(m, t, x)) \end{array}}}
\vnbprop{hahn-banach}{\vnbstmt{\begin{array}{@{}l@{}} \forall m, s, f. \\ \quad \operatorname{is\text{-}normed\text{-}vector\text{-}space}(m) \wedge \\ \quad \operatorname{is\text{-}finite\text{-}dimensional}(\operatorname{normed\text{-}vector\text{-}space\text{-}as\text{-}module}(m)) \wedge \\ \quad \operatorname{is\text{-}submodule}(m, s) \wedge \\ \quad \operatorname{is\text{-}bounded\text{-}linear\text{-}functional\text{-}on}(m, s, f) \Rightarrow \\ \quad \quad \exists g. \\ \quad \quad \quad \operatorname{is\text{-}linear\text{-}functional\text{-}on}(m, \operatorname{vec}(m), g) \wedge \\ \quad \quad \quad \operatorname{extends\text{-}on}(s, g, f) \wedge \\ \quad \quad \quad \forall w \in \operatorname{vec}(m).\; (\lvert g(w) \rvert \leq (\operatorname{dual\text{-}norm\text{-}on}(m, s, f) \cdot \operatorname{vnrm}(m)(w))) \end{array}}}

\vnbbookfile{structure-library/hahn-banach-sublinear}{structure-library/hahn-banach-sublinear}
\noindent{\small structure-library/hahn-banach-sublinear.scm -- the vocabulary of the Hahn-Banach theorem for a SUBLINEAR functional (RM-3A, 2026-10-10). DEFINITIONS ONLY; every law is PROVEN in theorem-library/hahn-banach-sublinear.scm (prefix hbs-).}\par
\par\medskip\noindent\textbf{Definition} (\texttt{is-sublinear}).\;For all $m$ and $p$, 
\begin{equation*}
\begin{array}{@{}l@{}} \operatorname{is\text{-}sublinear}(m, p) \\ \Leftrightarrow\; p \in (\operatorname{vec}(m) \to \mathbb{R}) \wedge \\ \forall \operatorname{hsx\_}, \operatorname{hsy\_} \in \operatorname{vec}(m). \\ \quad p(\operatorname{vadd}(m)(\operatorname{hsx\_}, \operatorname{hsy\_})) \leq +(p(\operatorname{hsx\_}), p(\operatorname{hsy\_})) \wedge \\ \forall \operatorname{hst\_} \in \mathbb{R}. \\ \quad (0 \leq \operatorname{hst\_}) \Rightarrow  \\ \quad \quad \forall \operatorname{hsx\_} \in \operatorname{vec}(m).\; (p(\operatorname{act}(m)(\operatorname{hst\_}, \operatorname{hsx\_})) = \cdot(\operatorname{hst\_}, p(\operatorname{hsx\_}))) \end{array}
\end{equation*}


\vnbbookfile{theorem-library/hahn-banach-sublinear}{theorem-library/hahn-banach-sublinear}
\noindent{\small theorem-library/hahn-banach-sublinear.scm -- HAHN-BANACH FOR A SUBLINEAR FUNCTIONAL on a real vector space (RM-3A, 2026-10-10).  Prefix hbs-.}\par
\vnbprop{hbs-carr}{\vnbstmt{\begin{array}{@{}l@{}} \forall m. \\ \quad (\operatorname{is\text{-}module}(m) \wedge (\operatorname{scal}(m) = \operatorname{normed\text{-}field\text{-}as\text{-}commutative\text{-}ring}(\operatorname{rr\text{-}normed\text{-}field}))) \Rightarrow  \\ \quad \quad \operatorname{carr}(\operatorname{scal}(m)) = \mathbb{R} \end{array}}}
\vnbprop{hbs-in-carr}{\vnbstmt{\begin{array}{@{}l@{}} \forall m.; \forall r \in \mathbb{R}. \\ \quad (\operatorname{is\text{-}module}(m) \wedge (\operatorname{scal}(m) = \operatorname{normed\text{-}field\text{-}as\text{-}commutative\text{-}ring}(\operatorname{rr\text{-}normed\text{-}field}))) \Rightarrow  \\ \quad \quad r \in \operatorname{carr}(\operatorname{scal}(m)) \end{array}}}
\vnbprop{hbs-add-eq}{\vnbstmt{\begin{array}{@{}l@{}} \forall m.; \forall r, s \in \mathbb{R}. \\ \quad (\operatorname{is\text{-}module}(m) \wedge (\operatorname{scal}(m) = \operatorname{normed\text{-}field\text{-}as\text{-}commutative\text{-}ring}(\operatorname{rr\text{-}normed\text{-}field}))) \Rightarrow  \\ \quad \quad \operatorname{add}(\operatorname{scal}(m))(r, s) = (r + s) \end{array}}}
\vnbprop{hbs-mul-eq}{\vnbstmt{\begin{array}{@{}l@{}} \forall m.; \forall r, s \in \mathbb{R}. \\ \quad (\operatorname{is\text{-}module}(m) \wedge (\operatorname{scal}(m) = \operatorname{normed\text{-}field\text{-}as\text{-}commutative\text{-}ring}(\operatorname{rr\text{-}normed\text{-}field}))) \Rightarrow  \\ \quad \quad \operatorname{mul}(\operatorname{scal}(m))(r, s) = (r \cdot s) \end{array}}}
\vnbprop{hbs-one-eq}{\vnbstmt{\begin{array}{@{}l@{}} \forall m. \\ \quad (\operatorname{is\text{-}module}(m) \wedge (\operatorname{scal}(m) = \operatorname{normed\text{-}field\text{-}as\text{-}commutative\text{-}ring}(\operatorname{rr\text{-}normed\text{-}field}))) \Rightarrow  \\ \quad \quad \operatorname{one}(\operatorname{scal}(m)) = 1 \end{array}}}
\vnbprop{hbs-zero-eq}{\vnbstmt{\begin{array}{@{}l@{}} \forall m. \\ \quad (\operatorname{is\text{-}module}(m) \wedge (\operatorname{scal}(m) = \operatorname{normed\text{-}field\text{-}as\text{-}commutative\text{-}ring}(\operatorname{rr\text{-}normed\text{-}field}))) \Rightarrow  \\ \quad \quad \operatorname{zero}(\operatorname{scal}(m)) = 0 \end{array}}}
\vnbprop{hbs-act-in}{\vnbstmt{\begin{array}{@{}l@{}} \forall m.; \forall r \in \mathbb{R}, x \in \operatorname{vec}(m). \\ \quad (\operatorname{is\text{-}module}(m) \wedge (\operatorname{scal}(m) = \operatorname{normed\text{-}field\text{-}as\text{-}commutative\text{-}ring}(\operatorname{rr\text{-}normed\text{-}field}))) \Rightarrow  \\ \quad \quad \operatorname{act}(m)(r, x) \in \operatorname{vec}(m) \end{array}}}
\vnbprop{hbs-act-add}{\vnbstmt{\begin{array}{@{}l@{}} \forall m.; \forall r, s \in \mathbb{R}, x \in \operatorname{vec}(m). \\ \quad (\operatorname{is\text{-}module}(m) \wedge (\operatorname{scal}(m) = \operatorname{normed\text{-}field\text{-}as\text{-}commutative\text{-}ring}(\operatorname{rr\text{-}normed\text{-}field}))) \Rightarrow  \\ \quad \quad \operatorname{act}(m)((r + s), x) = \operatorname{vadd}(m)(\operatorname{act}(m)(r, x), \operatorname{act}(m)(s, x)) \end{array}}}
\vnbprop{hbs-act-mul}{\vnbstmt{\begin{array}{@{}l@{}} \forall m.; \forall r, s \in \mathbb{R}, x \in \operatorname{vec}(m). \\ \quad (\operatorname{is\text{-}module}(m) \wedge (\operatorname{scal}(m) = \operatorname{normed\text{-}field\text{-}as\text{-}commutative\text{-}ring}(\operatorname{rr\text{-}normed\text{-}field}))) \Rightarrow  \\ \quad \quad \operatorname{act}(m)((r \cdot s), x) = \operatorname{act}(m)(r, \operatorname{act}(m)(s, x)) \end{array}}}
\vnbprop{hbs-act-one}{\vnbstmt{\begin{array}{@{}l@{}} \forall m.; \forall x \in \operatorname{vec}(m). \\ \quad (\operatorname{is\text{-}module}(m) \wedge (\operatorname{scal}(m) = \operatorname{normed\text{-}field\text{-}as\text{-}commutative\text{-}ring}(\operatorname{rr\text{-}normed\text{-}field}))) \Rightarrow  \\ \quad \quad \operatorname{act}(m)(1, x) = x \end{array}}}
\vnbprop{hbs-act-zero}{\vnbstmt{\begin{array}{@{}l@{}} \forall m.; \forall x \in \operatorname{vec}(m). \\ \quad (\operatorname{is\text{-}module}(m) \wedge (\operatorname{scal}(m) = \operatorname{normed\text{-}field\text{-}as\text{-}commutative\text{-}ring}(\operatorname{rr\text{-}normed\text{-}field}))) \Rightarrow  \\ \quad \quad \operatorname{act}(m)(0, x) = \operatorname{vzero}(m) \end{array}}}
\vnbprop{hbs-act-vadd}{\vnbstmt{\begin{array}{@{}l@{}} \forall m.; \forall r \in \mathbb{R}, x, y \in \operatorname{vec}(m). \\ \quad (\operatorname{is\text{-}module}(m) \wedge (\operatorname{scal}(m) = \operatorname{normed\text{-}field\text{-}as\text{-}commutative\text{-}ring}(\operatorname{rr\text{-}normed\text{-}field}))) \Rightarrow  \\ \quad \quad \operatorname{act}(m)(r, \operatorname{vadd}(m)(x, y)) = \operatorname{vadd}(m)(\operatorname{act}(m)(r, x), \operatorname{act}(m)(r, y)) \end{array}}}
\vnbprop{hbs-vadd-assoc}{\vnbstmt{\begin{array}{@{}l@{}} \forall m.; \forall x, y, z \in \operatorname{vec}(m). \\ \quad \operatorname{is\text{-}module}(m) \Rightarrow  \\ \quad \quad \operatorname{vadd}(m)(\operatorname{vadd}(m)(x, y), z) = \operatorname{vadd}(m)(x, \operatorname{vadd}(m)(y, z)) \end{array}}}
\vnbprop{hbs-vadd-comm}{\vnbstmt{\begin{array}{@{}l@{}} \forall m.; \forall x, y \in \operatorname{vec}(m). \\ \quad \operatorname{is\text{-}module}(m) \Rightarrow (\operatorname{vadd}(m)(x, y) = \operatorname{vadd}(m)(y, x)) \end{array}}}
\vnbprop{hbs-vzero-lr}{\vnbstmt{\begin{array}{@{}l@{}} \forall m.; \forall x \in \operatorname{vec}(m). \\ \quad \operatorname{is\text{-}module}(m) \Rightarrow ((\operatorname{vadd}(m)(\operatorname{vzero}(m), x) = x) \wedge (\operatorname{vadd}(m)(x, \operatorname{vzero}(m)) = x)) \end{array}}}
\vnbprop{hbs-vzero-r}{\vnbstmt{\forall m.; \forall x \in \operatorname{vec}(m).\; (\operatorname{is\text{-}module}(m) \Rightarrow (\operatorname{vadd}(m)(x, \operatorname{vzero}(m)) = x))}}
\vnbprop{hbs-vzero-l}{\vnbstmt{\forall m.; \forall x \in \operatorname{vec}(m).\; (\operatorname{is\text{-}module}(m) \Rightarrow (\operatorname{vadd}(m)(\operatorname{vzero}(m), x) = x))}}
\vnbprop{hbs-neg-cancel}{\vnbstmt{\begin{array}{@{}l@{}} \forall m.; \forall x \in \operatorname{vec}(m). \\ \quad (\operatorname{is\text{-}module}(m) \wedge (\operatorname{scal}(m) = \operatorname{normed\text{-}field\text{-}as\text{-}commutative\text{-}ring}(\operatorname{rr\text{-}normed\text{-}field}))) \Rightarrow  \\ \quad \quad \operatorname{vadd}(m)(x, \operatorname{act}(m)(-1, x)) = \operatorname{vzero}(m) \end{array}}}
\vnbprop{hbs-neg-cancel-l}{\vnbstmt{\begin{array}{@{}l@{}} \forall m.; \forall x \in \operatorname{vec}(m). \\ \quad (\operatorname{is\text{-}module}(m) \wedge (\operatorname{scal}(m) = \operatorname{normed\text{-}field\text{-}as\text{-}commutative\text{-}ring}(\operatorname{rr\text{-}normed\text{-}field}))) \Rightarrow  \\ \quad \quad \operatorname{vadd}(m)(\operatorname{act}(m)(-1, x), x) = \operatorname{vzero}(m) \end{array}}}
\vnbprop{hbs-cancel-r}{\vnbstmt{\begin{array}{@{}l@{}} \forall m.; \forall x, y \in \operatorname{vec}(m). \\ \quad (\operatorname{is\text{-}module}(m) \wedge (\operatorname{scal}(m) = \operatorname{normed\text{-}field\text{-}as\text{-}commutative\text{-}ring}(\operatorname{rr\text{-}normed\text{-}field}))) \Rightarrow  \\ \quad \quad \operatorname{vadd}(m)(\operatorname{vadd}(m)(x, y), \operatorname{act}(m)(-1, y)) = x \end{array}}}
\vnbprop{hbs-cancel-l}{\vnbstmt{\begin{array}{@{}l@{}} \forall m.; \forall x, y \in \operatorname{vec}(m). \\ \quad (\operatorname{is\text{-}module}(m) \wedge (\operatorname{scal}(m) = \operatorname{normed\text{-}field\text{-}as\text{-}commutative\text{-}ring}(\operatorname{rr\text{-}normed\text{-}field}))) \Rightarrow  \\ \quad \quad \operatorname{vadd}(m)(\operatorname{vadd}(m)(x, y), \operatorname{act}(m)(-1, x)) = y \end{array}}}
\vnbprop{hbs-interchange}{\vnbstmt{\begin{array}{@{}l@{}} \forall m.; \forall a, b, c, d \in \operatorname{vec}(m). \\ \quad \operatorname{is\text{-}module}(m) \Rightarrow  \\ \quad \quad \operatorname{vadd}(m)(\operatorname{vadd}(m)(a, b), \operatorname{vadd}(m)(c, d)) = \operatorname{vadd}(m)(\operatorname{vadd}(m)(a, c), \operatorname{vadd}(m)(b, d)) \end{array}}}
\vnbprop{hbs-inf-point}{\vnbstmt{\begin{array}{@{}l@{}} \forall i, a, x. \\ \quad \exists i_{0}.\; (i_{0} \in i) \wedge \\ \quad \forall j \in i.\; (a(x, j) \in \mathbb{R}) \wedge \\ \quad \exists b \in \mathbb{R}.\; \forall j \in i.\; (b \leq a(x, j)) \Rightarrow \\ \quad \quad -\operatorname{sup}(\{w \in \mathbb{R} : \exists k \in i.\; (w = -a(x, k))\}) \in \mathbb{R} \wedge \\ \quad \quad \forall j \in i.\; (-\operatorname{sup}(\{w \in \mathbb{R} : \exists k \in i.\; (w = -a(x, k))\}) \leq a(x, j)) \wedge \\ \quad \quad \forall b \in \mathbb{R}. \\ \quad \quad \quad \forall j \in i.\; (b \leq a(x, j)) \Rightarrow  \\ \quad \quad \quad \quad b \leq -\operatorname{sup}(\{w \in \mathbb{R} : \exists k \in i.\; (w = -a(x, k))\}) \end{array}}}
\vnbprop{hbs-inf-family}{\vnbstmt{\begin{array}{@{}l@{}} \forall d, f, i. \\ \quad d \in \mathbf{Set} \wedge \\ \quad \forall x \in d.\; \exists i_{0}.\; (i_{0} \in f(x)) \wedge \\ \quad \forall x \in d, j \in f(x).\; (i(x, j) \in \mathbb{R}) \wedge \\ \quad \forall x \in d.\; \exists b \in \mathbb{R}.\; \forall j \in f(x).\; (b \leq i(x, j)) \Rightarrow \\ \quad \quad \exists r \in (d \to \mathbb{R}). \\ \quad \quad \quad \forall x \in d, j \in f(x).\; (r(x) \leq i(x, j)) \wedge \\ \quad \quad \quad \forall x \in d, b \in \mathbb{R}.\; (\forall j \in f(x).\; (b \leq i(x, j)) \Rightarrow (b \leq r(x))) \end{array}}}
\vnbprop{hbs-inf-family-c}{\vnbstmt{\begin{array}{@{}l@{}} \forall d, i, a. \\ \quad d \in \mathbf{Set} \wedge \\ \quad \forall x \in d.\; \exists i_{0}.\; (i_{0} \in i) \wedge \\ \quad \forall x \in d, j \in i.\; (a(x, j) \in \mathbb{R}) \wedge \\ \quad \forall x \in d.\; \exists b \in \mathbb{R}.\; \forall j \in i.\; (b \leq a(x, j)) \Rightarrow \\ \quad \quad \exists r \in (d \to \mathbb{R}). \\ \quad \quad \quad \forall x \in d, j \in i.\; (r(x) \leq a(x, j)) \wedge \\ \quad \quad \quad \forall x \in d, b \in \mathbb{R}.\; (\forall j \in i.\; (b \leq a(x, j)) \Rightarrow (b \leq r(x))) \end{array}}}
\vnbprop{hbs-pos-inv}{\vnbstmt{\forall t \in \mathbb{R}.\; ((0 < t) \Rightarrow \exists u \in \mathbb{R}.\; ((0 < u) \wedge ((t \cdot u) = 1)))}}
\vnbprop{hbs-rr-div-le}{\vnbstmt{\begin{array}{@{}l@{}} \forall t, u, c, a \in \mathbb{R}. \\ \quad ((0 < t) \wedge ((0 < u) \wedge (((t \cdot u) = 1) \wedge (c \leq (t \cdot a))))) \Rightarrow ((u \cdot c) \leq a) \end{array}}}
\vnbprop{hbs-rr-mul-back}{\vnbstmt{\begin{array}{@{}l@{}} \forall t, u, c, a \in \mathbb{R}. \\ \quad ((0 < t) \wedge ((0 < u) \wedge (((t \cdot u) = 1) \wedge ((u \cdot c) \leq a)))) \Rightarrow (c \leq (t \cdot a)) \end{array}}}
\vnbprop{hbs-sl-fun}{\vnbstmt{\forall m, q.\; (\operatorname{is\text{-}sublinear}(m, q) \Rightarrow (q \in (\operatorname{vec}(m) \to \mathbb{R})))}}
\vnbprop{hbs-sl-sub}{\vnbstmt{\begin{array}{@{}l@{}} \forall m, q.; \forall x, y \in \operatorname{vec}(m). \\ \quad \operatorname{is\text{-}sublinear}(m, q) \Rightarrow (q(\operatorname{vadd}(m)(x, y)) \leq (q(x) + q(y))) \end{array}}}
\vnbprop{hbs-sl-hom}{\vnbstmt{\begin{array}{@{}l@{}} \forall m, q.; \forall t \in \mathbb{R}, x \in \operatorname{vec}(m). \\ \quad (\operatorname{is\text{-}sublinear}(m, q) \wedge (0 \leq t)) \Rightarrow (q(\operatorname{act}(m)(t, x)) = (t \cdot q(x))) \end{array}}}
\vnbprop{hbs-sl-val}{\vnbstmt{\forall m, q.; \forall x \in \operatorname{vec}(m).\; (\operatorname{is\text{-}sublinear}(m, q) \Rightarrow (q(x) \in \mathbb{R}))}}
\vnbprop{hbs-sl-zero}{\vnbstmt{\begin{array}{@{}l@{}} \forall m, q. \\ \quad (\operatorname{is\text{-}module}(m) \wedge ((\operatorname{scal}(m) = \operatorname{normed\text{-}field\text{-}as\text{-}commutative\text{-}ring}(\operatorname{rr\text{-}normed\text{-}field})) \wedge \operatorname{is\text{-}sublinear}(m, q))) \Rightarrow  \\ \quad \quad q(\operatorname{vzero}(m)) = 0 \end{array}}}
\vnbprop{hbs-sl-negsum}{\vnbstmt{\begin{array}{@{}l@{}} \forall m, q.; \forall x \in \operatorname{vec}(m). \\ \quad (\operatorname{is\text{-}module}(m) \wedge ((\operatorname{scal}(m) = \operatorname{normed\text{-}field\text{-}as\text{-}commutative\text{-}ring}(\operatorname{rr\text{-}normed\text{-}field})) \wedge \operatorname{is\text{-}sublinear}(m, q))) \Rightarrow  \\ \quad \quad 0 \leq (q(x) + q(\operatorname{act}(m)(-1, x))) \end{array}}}
\vnbprop{hbs-sl-intro-pos}{\vnbstmt{\begin{array}{@{}l@{}} \forall m.; \forall q \in (\operatorname{vec}(m) \to \mathbb{R}). \\ \quad \operatorname{is\text{-}module}(m) \wedge \\ \quad \operatorname{scal}(m) = \operatorname{normed\text{-}field\text{-}as\text{-}commutative\text{-}ring}(\operatorname{rr\text{-}normed\text{-}field}) \wedge \\ \quad \forall x, y \in \operatorname{vec}(m).\; (q(\operatorname{vadd}(m)(x, y)) \leq (q(x) + q(y))) \wedge \\ \quad \forall t \in \mathbb{R}, x \in \operatorname{vec}(m).\; ((0 < t) \Rightarrow (q(\operatorname{act}(m)(t, x)) \leq (t \cdot q(x)))) \wedge \\ \quad q(\operatorname{vzero}(m)) = 0 \Rightarrow \\ \quad \quad \operatorname{is\text{-}sublinear}(m, q) \end{array}}}
\vnbprop{hbs-sl-odd-linear}{\vnbstmt{\begin{array}{@{}l@{}} \forall m, q. \\ \quad \operatorname{is\text{-}module}(m) \wedge \\ \quad \operatorname{scal}(m) = \operatorname{normed\text{-}field\text{-}as\text{-}commutative\text{-}ring}(\operatorname{rr\text{-}normed\text{-}field}) \wedge \\ \quad \operatorname{is\text{-}sublinear}(m, q) \wedge \\ \quad \forall y \in \operatorname{vec}(m).\; (q(\operatorname{act}(m)(-1, y)) \leq -q(y)) \Rightarrow \\ \quad \quad \operatorname{is\text{-}linear\text{-}functional\text{-}on}(m, \operatorname{vec}(m), q) \end{array}}}
\vnbprop{hbs-tilt-exists}{\vnbstmt{\begin{array}{@{}l@{}} \forall m, q.; \forall y \in \operatorname{vec}(m). \\ \quad (\operatorname{is\text{-}module}(m) \wedge ((\operatorname{scal}(m) = \operatorname{normed\text{-}field\text{-}as\text{-}commutative\text{-}ring}(\operatorname{rr\text{-}normed\text{-}field})) \wedge \operatorname{is\text{-}sublinear}(m, q))) \Rightarrow  \\ \quad \quad \exists r \in (\operatorname{vec}(m) \to \mathbb{R}). \\ \quad \quad \quad \forall x \in \operatorname{vec}(m), t \in \mathbb{R}. \\ \quad \quad \quad \quad (0 \leq t) \Rightarrow (r(x) \leq (q(\operatorname{vadd}(m)(x, \operatorname{act}(m)(t, y))) - (t \cdot q(y)))) \wedge \\ \quad \quad \quad \forall x \in \operatorname{vec}(m), b \in \mathbb{R}. \\ \quad \quad \quad \quad \forall t \in \mathbb{R}.\; ((0 \leq t) \Rightarrow (b \leq (q(\operatorname{vadd}(m)(x, \operatorname{act}(m)(t, y))) - (t \cdot q(y))))) \Rightarrow  \\ \quad \quad \quad \quad \quad b \leq r(x) \end{array}}}
\vnbprop{hbs-tilt-props}{\vnbstmt{\begin{array}{@{}l@{}} \forall m, q.; \forall y \in \operatorname{vec}(m), t \in (\operatorname{vec}(m) \to \mathbb{R}). \\ \quad \operatorname{is\text{-}module}(m) \wedge \\ \quad \operatorname{scal}(m) = \operatorname{normed\text{-}field\text{-}as\text{-}commutative\text{-}ring}(\operatorname{rr\text{-}normed\text{-}field}) \wedge \\ \quad \operatorname{is\text{-}sublinear}(m, q) \wedge \\ \quad \forall x \in \operatorname{vec}(m), a \in \mathbb{R}. \\ \quad \quad (0 \leq a) \Rightarrow (t(x) \leq (q(\operatorname{vadd}(m)(x, \operatorname{act}(m)(a, y))) - (a \cdot q(y)))) \wedge \\ \quad \forall x \in \operatorname{vec}(m), b \in \mathbb{R}. \\ \quad \quad \forall a \in \mathbb{R}.\; ((0 \leq a) \Rightarrow (b \leq (q(\operatorname{vadd}(m)(x, \operatorname{act}(m)(a, y))) - (a \cdot q(y))))) \Rightarrow  \\ \quad \quad \quad b \leq t(x) \Rightarrow \\ \quad \quad (\operatorname{is\text{-}sublinear}(m, t) \wedge (\forall x \in \operatorname{vec}(m).\; (t(x) \leq q(x)) \wedge (t(\operatorname{act}(m)(-1, y)) \leq -q(y)))) \end{array}}}
\vnbprop{hbs-chain-exists}{\vnbstmt{\begin{array}{@{}l@{}} \forall m, h. \\ \quad \operatorname{is\text{-}module}(m) \wedge \\ \quad \operatorname{scal}(m) = \operatorname{normed\text{-}field\text{-}as\text{-}commutative\text{-}ring}(\operatorname{rr\text{-}normed\text{-}field}) \wedge \\ \quad \exists q_{0}.\; (q_{0} \in h) \wedge \\ \quad \forall q \in h.\; (q \in (\operatorname{vec}(m) \to \mathbb{R})) \wedge \\ \quad \forall x \in \operatorname{vec}(m).\; \exists b \in \mathbb{R}.\; \forall q \in h.\; (b \leq q(x)) \Rightarrow \\ \quad \quad \exists t \in (\operatorname{vec}(m) \to \mathbb{R}). \\ \quad \quad \quad \forall x \in \operatorname{vec}(m), q \in h.\; (t(x) \leq q(x)) \wedge \\ \quad \quad \quad \forall x \in \operatorname{vec}(m), b \in \mathbb{R}.\; (\forall q \in h.\; (b \leq q(x)) \Rightarrow (b \leq t(x))) \end{array}}}
\vnbprop{hbs-chain-props}{\vnbstmt{\begin{array}{@{}l@{}} \forall m, h.; \forall t \in (\operatorname{vec}(m) \to \mathbb{R}). \\ \quad \operatorname{is\text{-}module}(m) \wedge \\ \quad \operatorname{scal}(m) = \operatorname{normed\text{-}field\text{-}as\text{-}commutative\text{-}ring}(\operatorname{rr\text{-}normed\text{-}field}) \wedge \\ \quad \exists q_{0}.\; (q_{0} \in h) \wedge \\ \quad \forall q \in h.\; \operatorname{is\text{-}sublinear}(m, q) \wedge \\ \quad \forall q_{1}, q_{2} \in h. \\ \quad \quad \forall x \in \operatorname{vec}(m).\; (\operatorname{q1}(x) \leq \operatorname{q2}(x)) \vee \forall x \in \operatorname{vec}(m).\; (\operatorname{q2}(x) \leq \operatorname{q1}(x)) \wedge \\ \quad \forall x \in \operatorname{vec}(m), q \in h.\; (t(x) \leq q(x)) \wedge \\ \quad \forall x \in \operatorname{vec}(m), b \in \mathbb{R}.\; (\forall q \in h.\; (b \leq q(x)) \Rightarrow (b \leq t(x))) \Rightarrow \\ \quad \quad \operatorname{is\text{-}sublinear}(m, t) \end{array}}}
\vnbprop{hbs-lf-fun}{\vnbstmt{\forall m, s, f.\; (\operatorname{is\text{-}linear\text{-}functional\text{-}on}(m, s, f) \Rightarrow (f \in (s \to \mathbb{R})))}}
\vnbprop{hbs-lf-add}{\vnbstmt{\begin{array}{@{}l@{}} \forall m, s, f.; \forall x, y \in s. \\ \quad \operatorname{is\text{-}linear\text{-}functional\text{-}on}(m, s, f) \Rightarrow (f(\operatorname{vadd}(m)(x, y)) = (f(x) + f(y))) \end{array}}}
\vnbprop{hbs-lf-hom}{\vnbstmt{\begin{array}{@{}l@{}} \forall m, s, f.; \forall r \in \mathbb{R}, x \in s. \\ \quad \operatorname{is\text{-}linear\text{-}functional\text{-}on}(m, s, f) \Rightarrow (f(\operatorname{act}(m)(r, x)) = (r \cdot f(x))) \end{array}}}
\vnbprop{hbs-lf-val}{\vnbstmt{\forall m, s, f.; \forall x \in s.\; (\operatorname{is\text{-}linear\text{-}functional\text{-}on}(m, s, f) \Rightarrow (f(x) \in \mathbb{R}))}}
\vnbprop{hbs-sub-vec}{\vnbstmt{\forall m, s.; \forall x \in s.\; ((\operatorname{is\text{-}module}(m) \wedge \operatorname{is\text{-}submodule}(m, s)) \Rightarrow (x \in \operatorname{vec}(m)))}}
\vnbprop{hbs-sub-act}{\vnbstmt{\begin{array}{@{}l@{}} \forall m, s.; \forall r \in \mathbb{R}, x \in s. \\ \quad (\operatorname{is\text{-}module}(m) \wedge ((\operatorname{scal}(m) = \operatorname{normed\text{-}field\text{-}as\text{-}commutative\text{-}ring}(\operatorname{rr\text{-}normed\text{-}field})) \wedge \operatorname{is\text{-}submodule}(m, s))) \Rightarrow  \\ \quad \quad \operatorname{act}(m)(r, x) \in s \end{array}}}
\vnbprop{hbs-lf-zero}{\vnbstmt{\begin{array}{@{}l@{}} \forall m, s, f. \\ \quad (\operatorname{is\text{-}module}(m) \wedge (\operatorname{is\text{-}submodule}(m, s) \wedge \operatorname{is\text{-}linear\text{-}functional\text{-}on}(m, s, f))) \Rightarrow  \\ \quad \quad f(\operatorname{vzero}(m)) = 0 \end{array}}}
\vnbprop{hbs-red-exists}{\vnbstmt{\begin{array}{@{}l@{}} \forall m, p, s, f. \\ \quad \operatorname{is\text{-}module}(m) \wedge \\ \quad \operatorname{scal}(m) = \operatorname{normed\text{-}field\text{-}as\text{-}commutative\text{-}ring}(\operatorname{rr\text{-}normed\text{-}field}) \wedge \\ \quad \operatorname{is\text{-}sublinear}(m, p) \wedge \\ \quad \operatorname{is\text{-}submodule}(m, s) \wedge \\ \quad \operatorname{is\text{-}linear\text{-}functional\text{-}on}(m, s, f) \wedge \\ \quad \forall z \in s.\; (f(z) \leq p(z)) \Rightarrow \\ \quad \quad \exists a \in (\operatorname{vec}(m) \to \mathbb{R}). \\ \quad \quad \quad \forall x \in \operatorname{vec}(m), z \in s.\; (a(x) \leq (p(\operatorname{vadd}(m)(x, z)) - f(z))) \wedge \\ \quad \quad \quad \forall x \in \operatorname{vec}(m), b \in \mathbb{R}. \\ \quad \quad \quad \quad \forall z \in s.\; (b \leq (p(\operatorname{vadd}(m)(x, z)) - f(z))) \Rightarrow (b \leq a(x)) \end{array}}}
\vnbprop{hbs-red-props}{\vnbstmt{\begin{array}{@{}l@{}} \forall m, p, s, f.; \forall z \in (\operatorname{vec}(m) \to \mathbb{R}). \\ \quad \operatorname{is\text{-}module}(m) \wedge \\ \quad \operatorname{scal}(m) = \operatorname{normed\text{-}field\text{-}as\text{-}commutative\text{-}ring}(\operatorname{rr\text{-}normed\text{-}field}) \wedge \\ \quad \operatorname{is\text{-}sublinear}(m, p) \wedge \\ \quad \operatorname{is\text{-}submodule}(m, s) \wedge \\ \quad \operatorname{is\text{-}linear\text{-}functional\text{-}on}(m, s, f) \wedge \\ \quad \forall a \in s.\; (f(a) \leq p(a)) \wedge \\ \quad \forall x \in \operatorname{vec}(m), a \in s.\; (z(x) \leq (p(\operatorname{vadd}(m)(x, a)) - f(a))) \wedge \\ \quad \forall x \in \operatorname{vec}(m), b \in \mathbb{R}. \\ \quad \quad \forall a \in s.\; (b \leq (p(\operatorname{vadd}(m)(x, a)) - f(a))) \Rightarrow (b \leq z(x)) \Rightarrow \\ \quad \quad (\operatorname{is\text{-}sublinear}(m, z) \wedge (\forall x \in \operatorname{vec}(m).\; (z(x) \leq p(x)) \wedge \forall a \in s.\; (z(a) \leq f(a)))) \end{array}}}
\vnbprop{hbs-grd-in}{\vnbstmt{\begin{array}{@{}l@{}} \forall m, z.; \forall q \in (\operatorname{vec}(m) \to \mathbb{R}). \\ \quad (\operatorname{is\text{-}sublinear}(m, q) \wedge \forall x \in \operatorname{vec}(m).\; (q(x) \leq z(x))) \Rightarrow  \\ \quad \quad q \in \{\mathit{hgq} \in (\operatorname{vec}(m) \to \mathbb{R}) : (\operatorname{is\text{-}sublinear}(m, \mathit{hgq}) \wedge \forall \mathit{hgx} \in \operatorname{vec}(m).\; (\operatorname{hgq}(\mathit{hgx}) \leq z(\mathit{hgx})))\} \end{array}}}
\vnbprop{hbs-grd-parts}{\vnbstmt{\begin{array}{@{}l@{}} \forall m, z.; \forall q \in \{\mathit{hgq} \in (\operatorname{vec}(m) \to \mathbb{R}) : (\operatorname{is\text{-}sublinear}(m, \mathit{hgq}) \wedge \forall x \in \operatorname{vec}(m).\; (\operatorname{hgq}(x) \leq z(x)))\}. \\ \quad ((q \in (\operatorname{vec}(m) \to \mathbb{R})) \wedge (\operatorname{is\text{-}sublinear}(m, q) \wedge \forall a \in \operatorname{vec}(m).\; (q(a) \leq z(a)))) \end{array}}}
\vnbprop{hbs-porel-in}{\vnbstmt{\begin{array}{@{}l@{}} \forall m, z, a, b. \\ \quad a \in \{q \in (\operatorname{vec}(m) \to \mathbb{R}) : (\operatorname{is\text{-}sublinear}(m, q) \wedge \forall x \in \operatorname{vec}(m).\; (q(x) \leq z(x)))\} \wedge \\ \quad b \in \{q \in (\operatorname{vec}(m) \to \mathbb{R}) : (\operatorname{is\text{-}sublinear}(m, q) \wedge \forall x \in \operatorname{vec}(m).\; (q(x) \leq z(x)))\} \wedge \\ \quad \forall c \in \operatorname{vec}(m).\; (b(c) \leq a(c)) \Rightarrow \\ \quad \quad \langle a, b \rangle \\ \quad \quad \in\; \{w \in (\{q \in (\operatorname{vec}(m) \to \mathbb{R}) : (\operatorname{is\text{-}sublinear}(m, q) \wedge \forall x \in \operatorname{vec}(m).\; (q(x) \leq z(x)))\} \times \{q \in (\operatorname{vec}(m) \to \mathbb{R}) : (\operatorname{is\text{-}sublinear}(m, q) \wedge \forall x \in \operatorname{vec}(m).\; (q(x) \leq z(x)))\}) : ((\operatorname{nth}(1, w) \in \{q \in (\operatorname{vec}(m) \to \mathbb{R}) : (\operatorname{is\text{-}sublinear}(m, q) \wedge \forall x \in \operatorname{vec}(m).\; (q(x) \leq z(x)))\}) \wedge ((\operatorname{nth}(2, w) \in \{q \in (\operatorname{vec}(m) \to \mathbb{R}) : (\operatorname{is\text{-}sublinear}(m, q) \wedge \forall x \in \operatorname{vec}(m).\; (q(x) \leq z(x)))\}) \wedge \forall x \in \operatorname{vec}(m).\; (\operatorname{nth}(2, w)(x) \leq \operatorname{nth}(1, w)(x))))\} \end{array}}}
\vnbprop{hbs-porel-parts}{\vnbstmt{\begin{array}{@{}l@{}} \forall m, z, a, b. \\ \quad \langle a, b \rangle \\ \quad \in\; \{w \in (\{q \in (\operatorname{vec}(m) \to \mathbb{R}) : (\operatorname{is\text{-}sublinear}(m, q) \wedge \forall x \in \operatorname{vec}(m).\; (q(x) \leq z(x)))\} \times \{q \in (\operatorname{vec}(m) \to \mathbb{R}) : (\operatorname{is\text{-}sublinear}(m, q) \wedge \forall x \in \operatorname{vec}(m).\; (q(x) \leq z(x)))\}) : ((\operatorname{nth}(1, w) \in \{q \in (\operatorname{vec}(m) \to \mathbb{R}) : (\operatorname{is\text{-}sublinear}(m, q) \wedge \forall x \in \operatorname{vec}(m).\; (q(x) \leq z(x)))\}) \wedge ((\operatorname{nth}(2, w) \in \{q \in (\operatorname{vec}(m) \to \mathbb{R}) : (\operatorname{is\text{-}sublinear}(m, q) \wedge \forall x \in \operatorname{vec}(m).\; (q(x) \leq z(x)))\}) \wedge \forall x \in \operatorname{vec}(m).\; (\operatorname{nth}(2, w)(x) \leq \operatorname{nth}(1, w)(x))))\} \Rightarrow \\ \quad \quad a \in \{q \in (\operatorname{vec}(m) \to \mathbb{R}) : (\operatorname{is\text{-}sublinear}(m, q) \wedge \forall x \in \operatorname{vec}(m).\; (q(x) \leq z(x)))\} \wedge \\ \quad \quad b \in \{q \in (\operatorname{vec}(m) \to \mathbb{R}) : (\operatorname{is\text{-}sublinear}(m, q) \wedge \forall x \in \operatorname{vec}(m).\; (q(x) \leq z(x)))\} \wedge \\ \quad \quad \forall c \in \operatorname{vec}(m).\; (b(c) \leq a(c)) \end{array}}}
\vnbprop{hbs-dominated-linear}{\vnbstmt{\begin{array}{@{}l@{}} \forall m, z. \\ \quad (\operatorname{is\text{-}module}(m) \wedge ((\operatorname{scal}(m) = \operatorname{normed\text{-}field\text{-}as\text{-}commutative\text{-}ring}(\operatorname{rr\text{-}normed\text{-}field})) \wedge \operatorname{is\text{-}sublinear}(m, z))) \Rightarrow  \\ \quad \quad \exists f. \\ \quad \quad \quad \operatorname{is\text{-}linear\text{-}functional\text{-}on}(m, \operatorname{vec}(m), f) \wedge \forall x \in \operatorname{vec}(m).\; (f(x) \leq z(x)) \end{array}}}
\vnbprop{hbs-hahn-banach-sublinear}{\vnbstmt{\begin{array}{@{}l@{}} \forall m, p, s, f. \\ \quad \operatorname{is\text{-}module}(m) \wedge \\ \quad \operatorname{scal}(m) = \operatorname{normed\text{-}field\text{-}as\text{-}commutative\text{-}ring}(\operatorname{rr\text{-}normed\text{-}field}) \wedge \\ \quad \operatorname{is\text{-}sublinear}(m, p) \wedge \\ \quad \operatorname{is\text{-}submodule}(m, s) \wedge \\ \quad \operatorname{is\text{-}linear\text{-}functional\text{-}on}(m, s, f) \wedge \\ \quad \forall z \in s.\; (f(z) \leq p(z)) \Rightarrow \\ \quad \quad \exists a. \\ \quad \quad \quad \operatorname{is\text{-}linear\text{-}functional\text{-}on}(m, \operatorname{vec}(m), a) \wedge \\ \quad \quad \quad \operatorname{extends\text{-}on}(s, a, f) \wedge \\ \quad \quad \quad \forall x \in \operatorname{vec}(m).\; (a(x) \leq p(x)) \end{array}}}
\vnbprop{hbs-rr-fnrm}{\vnbstmt{\forall t \in \mathbb{R}.\; (\operatorname{fnrm}(\operatorname{rr\text{-}normed\text{-}field})(t) = \lvert t \rvert)}}
\vnbprop{hahn-banach-seminorm-rr}{\vnbstmt{\begin{array}{@{}l@{}} \forall m, p, s, f. \\ \quad \operatorname{is\text{-}module}(m) \wedge \\ \quad \operatorname{scal}(m) = \operatorname{normed\text{-}field\text{-}as\text{-}commutative\text{-}ring}(\operatorname{rr\text{-}normed\text{-}field}) \wedge \\ \quad \operatorname{is\text{-}seminorm}(\operatorname{rr\text{-}normed\text{-}field}, m, p) \wedge \\ \quad \operatorname{is\text{-}submodule}(m, s) \wedge \\ \quad \operatorname{is\text{-}k\text{-}linear\text{-}on}(m, s, f) \wedge \\ \quad \forall x \in s.\; (\operatorname{fnrm}(\operatorname{rr\text{-}normed\text{-}field})(f(x)) \leq p(x)) \Rightarrow \\ \quad \quad \exists a. \\ \quad \quad \quad \operatorname{is\text{-}k\text{-}linear}(m, a) \wedge \\ \quad \quad \quad \operatorname{extends\text{-}on}(s, a, f) \wedge \\ \quad \quad \quad \forall x \in \operatorname{vec}(m).\; (\operatorname{fnrm}(\operatorname{rr\text{-}normed\text{-}field})(a(x)) \leq p(x)) \end{array}}}

\vnbbookfile{theorem-library/norm-as-sup-proof}{theorem-library/norm-as-sup-proof}
\noindent{\small norm-as-sup-proof.scm -- the norm as a supremum of bounded functionals: for a finite-dimensional real NVS m and x in VEC(m), ||x|| = sup \{ |f(x)| : f a bounded linear functional, ||f|| $<$= 1 \}.}\par
\vnbprop{norm-bounded-by-functionals}{\vnbstmt{\begin{array}{@{}l@{}} \forall m, f, x. \\ \quad \operatorname{is\text{-}normed\text{-}vector\text{-}space}(m) \wedge \\ \quad \operatorname{is\text{-}bounded\text{-}linear\text{-}functional}(m, f) \wedge \\ \quad x \in \operatorname{vec}(m) \wedge \\ \quad \operatorname{dual\text{-}norm}(m, f) \leq 1 \Rightarrow \\ \quad \quad \lvert f(x) \rvert \leq \operatorname{vnrm}(m)(x) \end{array}}}
\vnbprop{norm-attained-by-functional}{\vnbstmt{\begin{array}{@{}l@{}} \forall m, x. \\ \quad \operatorname{is\text{-}normed\text{-}vector\text{-}space}(m) \wedge \\ \quad \operatorname{is\text{-}finite\text{-}dimensional}(\operatorname{normed\text{-}vector\text{-}space\text{-}as\text{-}module}(m)) \wedge \\ \quad x \in \operatorname{vec}(m) \Rightarrow \\ \quad \quad \exists g. \\ \quad \quad \quad (\operatorname{is\text{-}bounded\text{-}linear\text{-}functional}(m, g) \wedge ((\operatorname{dual\text{-}norm}(m, g) \leq 1) \wedge (g(x) = \operatorname{vnrm}(m)(x)))) \end{array}}}
\vnbprop{norm-as-sup}{\vnbstmt{\begin{array}{@{}l@{}} \forall m, x. \\ \quad \operatorname{is\text{-}normed\text{-}vector\text{-}space}(m) \wedge \\ \quad \operatorname{is\text{-}finite\text{-}dimensional}(\operatorname{normed\text{-}vector\text{-}space\text{-}as\text{-}module}(m)) \wedge \\ \quad x \in \operatorname{vec}(m) \Rightarrow \\ \quad \quad \forall f. \\ \quad \quad \quad (\operatorname{is\text{-}bounded\text{-}linear\text{-}functional}(m, f) \wedge (\operatorname{dual\text{-}norm}(m, f) \leq 1)) \Rightarrow  \\ \quad \quad \quad \quad \lvert f(x) \rvert \leq \operatorname{vnrm}(m)(x) \wedge \\ \quad \quad \forall d \in \mathbb{R}. \\ \quad \quad \quad \forall f.\; ((\operatorname{is\text{-}bounded\text{-}linear\text{-}functional}(m, f) \wedge (\operatorname{dual\text{-}norm}(m, f) \leq 1)) \Rightarrow (\lvert f(x) \rvert \leq d)) \Rightarrow  \\ \quad \quad \quad \quad \operatorname{vnrm}(m)(x) \leq d \end{array}}}

\vnbbookfile{theorem-library/vector-taylor-proof}{theorem-library/vector-taylor-proof}
\noindent{\small vector-taylor-proof.scm -- vector-valued Taylor with a norm remainder bound, REDUCED TO THE SCALAR CASE via a norm-attaining functional (Hahn-Banach payoff).}\par
\par\medskip\noindent\textbf{Definition} (\texttt{taylor-poly-v}).\;For all $m$, $f$, $a$ and $x$, 
\begin{equation*}
(\operatorname{taylor\text{-}poly\text{-}v}(m, f, a, x, 0) \simeq \operatorname{act}(m)(\cdot(\operatorname{power}(-(x, a), 0), \operatorname{recip}(\operatorname{factorial}(0))), \operatorname{nth\text{-}deriv\text{-}v}(m, f, 0)(a)))
\end{equation*}
\par
Let $m$, $f$, $a$ and $x$ be arbitrary. Let $n \in \mathbb{N}$. Then
\begin{equation*}
(\operatorname{taylor\text{-}poly\text{-}v}(m, f, a, x, \operatorname{succ}(n)) \simeq \operatorname{vadd}(m)(\operatorname{taylor\text{-}poly\text{-}v}(m, f, a, x, n), \operatorname{act}(m)(\cdot(\operatorname{power}(-(x, a), \operatorname{succ}(n)), \operatorname{recip}(\operatorname{factorial}(\operatorname{succ}(n)))), \operatorname{nth\text{-}deriv\text{-}v}(m, f, \operatorname{succ}(n))(a))))
\end{equation*}

\par\medskip\noindent\textbf{Definition} (\texttt{is-diff-at-v}).\;For all $m$, $f$, $a$ and $l$, 
\begin{equation*}
\begin{array}{@{}l@{}} \operatorname{is\text{-}diff\text{-}at\text{-}v}(m, f, a, l) \\ \Leftrightarrow\; \operatorname{is\text{-}normed\text{-}vector\text{-}space}(m) \wedge \\ f \in (\mathbb{R} \to \operatorname{vec}(m)) \wedge \\ a \in \mathbb{R} \wedge \\ l \in \operatorname{vec}(m) \wedge \\ \exists \varphi \in (\mathbb{R} \to \operatorname{vec}(m)). \\ \quad \operatorname{is\text{-}continuous\text{-}at}(\operatorname{rr\text{-}ms}, \operatorname{nvs\text{-}metric\text{-}space}(m), \varphi, a) \wedge \\ \quad \varphi(a) = l \wedge \\ \quad \forall \operatorname{x\_} \in \mathbb{R}. \\ \quad \quad \operatorname{vadd}(m)(f(\operatorname{x\_}), \operatorname{vneg}(m)(f(a))) = \operatorname{act}(m)(-(\operatorname{x\_}, a), \varphi(\operatorname{x\_})) \end{array}
\end{equation*}

\par\medskip\noindent\textbf{Definition} (\texttt{deriv-v}).\;$\operatorname{deriv\text{-}v}(m, f, a) :=$ $(\iota\, l.\;\operatorname{is\text{-}diff\text{-}at\text{-}v}(m, f, a, l))$
\par\medskip\noindent\textbf{Definition} (\texttt{nth-deriv-v}).\;For all $m$ and $f$,  $(\operatorname{nth\text{-}deriv\text{-}v}(m, f, 0) \simeq f)$\par
Let $m$ and $f$ be arbitrary. Let $n \in \mathbb{N}$. Then
\begin{equation*}
(\operatorname{nth\text{-}deriv\text{-}v}(m, f, \operatorname{succ}(n)) \simeq (\lambda x \in \mathbb{R}.\; \operatorname{deriv\text{-}v}(m, \operatorname{nth\text{-}deriv\text{-}v}(m, f, n), x)))
\end{equation*}

\par\medskip\noindent\textbf{Definition} (\texttt{taylor-differentiable-v}).\;For all $m$, $f$, $a$, $x$ and $n$, 
\begin{equation*}
\begin{array}{@{}l@{}} \operatorname{taylor\text{-}differentiable\text{-}v}(m, f, a, x, n) \\ \Leftrightarrow\; \forall k. \\ \quad ((k \in \mathbb{N}) \wedge (k \leq n)) \Rightarrow  \\ \quad \quad \forall t \in \operatorname{ccint}(a, x). \\ \quad \quad \quad \operatorname{is\text{-}continuous\text{-}at}(\operatorname{rr\text{-}ms}, \operatorname{nvs\text{-}metric\text{-}space}(m), \operatorname{nth\text{-}deriv\text{-}v}(m, f, k), t) \wedge \\ \forall k. \\ \quad ((k \in \mathbb{N}) \wedge (k \leq n)) \Rightarrow  \\ \quad \quad \forall t. \\ \quad \quad \quad ((a < t) \wedge (t < x)) \Rightarrow  \\ \quad \quad \quad \quad \operatorname{is\text{-}diff\text{-}at\text{-}v}(m, \operatorname{nth\text{-}deriv\text{-}v}(m, f, k), t, \operatorname{nth\text{-}deriv\text{-}v}(m, f, \operatorname{succ}(k))(t)) \end{array}
\end{equation*}

\vnbprop{diff-v-value-in-vec}{\vnbstmt{\forall m, f, a, l.\; (\operatorname{is\text{-}diff\text{-}at\text{-}v}(m, f, a, l) \Rightarrow (l \in \operatorname{vec}(m)))}}
\vnbprop{taylor-v-derivs-in-fun}{\vnbstmt{\begin{array}{@{}l@{}} \forall m, f, a, x, n.; \forall j \in \mathbb{N}. \\ \quad (\operatorname{taylor\text{-}differentiable\text{-}v}(m, f, a, x, n) \wedge ((a \in \mathbb{R}) \wedge ((x \in \mathbb{R}) \wedge ((a < x) \wedge (j \leq n))))) \Rightarrow  \\ \quad \quad \operatorname{nth\text{-}deriv\text{-}v}(m, f, j) \in (\mathbb{R} \to \operatorname{vec}(m)) \end{array}}}
\vnbprop{dfun-v-mono}{\vnbstmt{\begin{array}{@{}l@{}} \forall m, f, n, k.; \forall j \in \mathbb{N}. \\ \quad \forall j \in \mathbb{N}.\; ((j \leq n) \Rightarrow (\operatorname{nth\text{-}deriv\text{-}v}(m, f, j) \in (\mathbb{R} \to \operatorname{vec}(m)))) \wedge \\ \quad n \in \mathbb{N} \wedge \\ \quad k \in \mathbb{N} \wedge \\ \quad k \leq n \wedge \\ \quad j \leq k \Rightarrow \\ \quad \quad \operatorname{nth\text{-}deriv\text{-}v}(m, f, j) \in (\mathbb{R} \to \operatorname{vec}(m)) \end{array}}}
\vnbprop{nth-deriv-v-in-vec}{\vnbstmt{\begin{array}{@{}l@{}} \forall m, f, k, t. \\ \quad \operatorname{is\text{-}normed\text{-}vector\text{-}space}(m) \wedge \\ \quad f \in (\mathbb{R} \to \operatorname{vec}(m)) \wedge \\ \quad k \in \mathbb{N} \wedge \\ \quad t \in \mathbb{R} \wedge \\ \quad \forall j \in \mathbb{N}.\; ((j \leq k) \Rightarrow (\operatorname{nth\text{-}deriv\text{-}v}(m, f, j) \in (\mathbb{R} \to \operatorname{vec}(m)))) \Rightarrow \\ \quad \quad \operatorname{nth\text{-}deriv\text{-}v}(m, f, k)(t) \in \operatorname{vec}(m) \end{array}}}
\vnbprop{taylor-v-deriv-in-vec}{\vnbstmt{\begin{array}{@{}l@{}} \forall m, f, a, x, n, t. \\ \quad (\operatorname{taylor\text{-}differentiable\text{-}v}(m, f, a, x, n) \wedge ((a < t) \wedge ((t < x) \wedge (n \in \mathbb{N})))) \Rightarrow  \\ \quad \quad \operatorname{nth\text{-}deriv\text{-}v}(m, f, (n + 1))(t) \in \operatorname{vec}(m) \end{array}}}
\vnbprop{gof-in-fun}{\vnbstmt{\begin{array}{@{}l@{}} \forall m, f, g. \\ \quad (\operatorname{is\text{-}bounded\text{-}linear\text{-}functional}(m, g) \wedge (f \in (\mathbb{R} \to \operatorname{vec}(m)))) \Rightarrow  \\ \quad \quad \operatorname{compose}(g, f) \in (\mathbb{R} \to \mathbb{R}) \end{array}}}
\vnbprop{rkt-vtaylor-poly-in-vec-ind}{\vnbstmt{\begin{array}{@{}l@{}} \forall n \in \mathbb{N}, m, f, a, x. \\ \quad \operatorname{is\text{-}normed\text{-}vector\text{-}space}(m) \wedge \\ \quad f \in (\mathbb{R} \to \operatorname{vec}(m)) \wedge \\ \quad a \in \mathbb{R} \wedge \\ \quad x \in \mathbb{R} \wedge \\ \quad \forall j \in \mathbb{N}.\; ((j \leq n) \Rightarrow (\operatorname{nth\text{-}deriv\text{-}v}(m, f, j) \in (\mathbb{R} \to \operatorname{vec}(m)))) \Rightarrow \\ \quad \quad \operatorname{taylor\text{-}poly\text{-}v}(m, f, a, x, n) \in \operatorname{vec}(m) \end{array}}}
\vnbprop{vtaylor-poly-in-vec}{\vnbstmt{\begin{array}{@{}l@{}} \forall m, f, a, x, n. \\ \quad \operatorname{is\text{-}normed\text{-}vector\text{-}space}(m) \wedge \\ \quad f \in (\mathbb{R} \to \operatorname{vec}(m)) \wedge \\ \quad a \in \mathbb{R} \wedge \\ \quad x \in \mathbb{R} \wedge \\ \quad n \in \mathbb{N} \wedge \\ \quad \forall j \in \mathbb{N}.\; ((j \leq n) \Rightarrow (\operatorname{nth\text{-}deriv\text{-}v}(m, f, j) \in (\mathbb{R} \to \operatorname{vec}(m)))) \Rightarrow \\ \quad \quad \operatorname{taylor\text{-}poly\text{-}v}(m, f, a, x, n) \in \operatorname{vec}(m) \end{array}}}
\vnbprop{vtaylor-remainder-in-vec}{\vnbstmt{\begin{array}{@{}l@{}} \forall m, f, a, x, n. \\ \quad \operatorname{is\text{-}normed\text{-}vector\text{-}space}(m) \wedge \\ \quad f \in (\mathbb{R} \to \operatorname{vec}(m)) \wedge \\ \quad a \in \mathbb{R} \wedge \\ \quad x \in \mathbb{R} \wedge \\ \quad n \in \mathbb{N} \wedge \\ \quad \forall j \in \mathbb{N}.\; ((j \leq n) \Rightarrow (\operatorname{nth\text{-}deriv\text{-}v}(m, f, j) \in (\mathbb{R} \to \operatorname{vec}(m)))) \Rightarrow \\ \quad \quad \operatorname{vadd}(m)(f(x), \operatorname{vneg}(m)(\operatorname{taylor\text{-}poly\text{-}v}(m, f, a, x, n))) \in \operatorname{vec}(m) \end{array}}}
\vnbprop{r8e-blf-vneg}{\vnbstmt{\begin{array}{@{}l@{}} \forall m, g.; \forall u \in \operatorname{vec}(m). \\ \quad (\operatorname{is\text{-}normed\text{-}vector\text{-}space}(m) \wedge \operatorname{is\text{-}bounded\text{-}linear\text{-}functional}(m, g)) \Rightarrow  \\ \quad \quad g(\operatorname{vneg}(m)(u)) = (0 - g(u)) \end{array}}}
\vnbprop{r8e-blf-sub}{\vnbstmt{\begin{array}{@{}l@{}} \forall m, g.; \forall u, v \in \operatorname{vec}(m). \\ \quad (\operatorname{is\text{-}normed\text{-}vector\text{-}space}(m) \wedge \operatorname{is\text{-}bounded\text{-}linear\text{-}functional}(m, g)) \Rightarrow  \\ \quad \quad g(\operatorname{vadd}(m)(u, \operatorname{vneg}(m)(v))) = (g(u) - g(v)) \end{array}}}
\vnbprop{r8e-eps-over}{\vnbstmt{\begin{array}{@{}l@{}} \forall c \in \mathbb{R}, e. \\ \quad ((0 \leq c) \wedge \operatorname{pos\text{-}rr}(e)) \Rightarrow  \\ \quad \quad \exists d.\; (\operatorname{pos\text{-}rr}(d) \wedge \forall z \in \mathbb{R}.\; ((z \leq d) \Rightarrow ((c \cdot z) \leq e))) \end{array}}}
\vnbprop{r8e-blf-compose-continuous}{\vnbstmt{\begin{array}{@{}l@{}} \forall m, g, h, t. \\ \quad \operatorname{is\text{-}normed\text{-}vector\text{-}space}(m) \wedge \\ \quad \operatorname{is\text{-}bounded\text{-}linear\text{-}functional}(m, g) \wedge \\ \quad h \in (\mathbb{R} \to \operatorname{vec}(m)) \wedge \\ \quad \operatorname{is\text{-}continuous\text{-}at}(\operatorname{rr\text{-}ms}, \operatorname{nvs\text{-}metric\text{-}space}(m), h, t) \Rightarrow \\ \quad \quad \operatorname{is\text{-}continuous\text{-}at}(\operatorname{rr\text{-}ms}, \operatorname{rr\text{-}ms}, \operatorname{compose}(g, h), t) \end{array}}}
\vnbprop{r8e-blf-diff-commute}{\vnbstmt{\begin{array}{@{}l@{}} \forall m, g, h, t, l. \\ \quad (\operatorname{is\text{-}bounded\text{-}linear\text{-}functional}(m, g) \wedge \operatorname{is\text{-}diff\text{-}at\text{-}v}(m, h, t, l)) \Rightarrow  \\ \quad \quad \operatorname{is\text{-}diff\text{-}at}(\operatorname{compose}(g, h), t, g(l)) \end{array}}}
\vnbprop{r8e-nth-deriv-v-diff}{\vnbstmt{\begin{array}{@{}l@{}} \forall m, f, j, x. \\ \quad j \in \mathbb{N} \wedge \\ \quad x \in \mathbb{R} \wedge \\ \quad \operatorname{nth\text{-}deriv\text{-}v}(m, f, j) \in (\mathbb{R} \to \operatorname{vec}(m)) \wedge \\ \quad \operatorname{nth\text{-}deriv\text{-}v}(m, f, (j + 1)) \in (\mathbb{R} \to \operatorname{vec}(m)) \Rightarrow \\ \quad \quad \operatorname{is\text{-}diff\text{-}at\text{-}v}(m, \operatorname{nth\text{-}deriv\text{-}v}(m, f, j), x, \operatorname{nth\text{-}deriv\text{-}v}(m, f, (j + 1))(x)) \end{array}}}
\vnbprop{r8e-gof-nth-deriv-fn}{\vnbstmt{\begin{array}{@{}l@{}} \forall k \in \mathbb{N}, m, f, g. \\ \quad \operatorname{is\text{-}bounded\text{-}linear\text{-}functional}(m, g) \wedge \\ \quad f \in (\mathbb{R} \to \operatorname{vec}(m)) \wedge \\ \quad \forall j \in \mathbb{N}.\; ((j \leq k) \Rightarrow (\operatorname{nth\text{-}deriv\text{-}v}(m, f, j) \in (\mathbb{R} \to \operatorname{vec}(m)))) \Rightarrow \\ \quad \quad \operatorname{nth\text{-}deriv}(\operatorname{compose}(g, f), k) = \operatorname{compose}(g, \operatorname{nth\text{-}deriv\text{-}v}(m, f, k)) \end{array}}}
\vnbprop{gof-nth-deriv-dfun}{\vnbstmt{\begin{array}{@{}l@{}} \forall m, f, g, n, t. \\ \quad \operatorname{is\text{-}bounded\text{-}linear\text{-}functional}(m, g) \wedge \\ \quad f \in (\mathbb{R} \to \operatorname{vec}(m)) \wedge \\ \quad n \in \mathbb{N} \wedge \\ \quad t \in \mathbb{R} \wedge \\ \quad \forall j \in \mathbb{N}.\; ((j \leq n) \Rightarrow (\operatorname{nth\text{-}deriv\text{-}v}(m, f, j) \in (\mathbb{R} \to \operatorname{vec}(m)))) \wedge \\ \quad \operatorname{is\text{-}diff\text{-}at\text{-}v}(m, \operatorname{nth\text{-}deriv\text{-}v}(m, f, n), t, \operatorname{nth\text{-}deriv\text{-}v}(m, f, (n + 1))(t)) \Rightarrow \\ \quad \quad \operatorname{nth\text{-}deriv}(\operatorname{compose}(g, f), (n + 1))(t) = g(\operatorname{nth\text{-}deriv\text{-}v}(m, f, (n + 1))(t)) \end{array}}}
\vnbprop{gof-taylor-diff-dfun}{\vnbstmt{\begin{array}{@{}l@{}} \forall m, f, g, a, x, n. \\ \quad \operatorname{is\text{-}normed\text{-}vector\text{-}space}(m) \wedge \\ \quad \operatorname{is\text{-}bounded\text{-}linear\text{-}functional}(m, g) \wedge \\ \quad f \in (\mathbb{R} \to \operatorname{vec}(m)) \wedge \\ \quad n \in \mathbb{N} \wedge \\ \quad \operatorname{taylor\text{-}differentiable\text{-}v}(m, f, a, x, n) \wedge \\ \quad \forall j \in \mathbb{N}.\; ((j \leq n) \Rightarrow (\operatorname{nth\text{-}deriv\text{-}v}(m, f, j) \in (\mathbb{R} \to \operatorname{vec}(m)))) \Rightarrow \\ \quad \quad \operatorname{taylor\text{-}differentiable}(\operatorname{compose}(g, f), a, x, n) \end{array}}}
\vnbprop{r8e-td-v-diff-at}{\vnbstmt{\begin{array}{@{}l@{}} \forall m, f, a, x, n, t. \\ \quad (\operatorname{taylor\text{-}differentiable\text{-}v}(m, f, a, x, n) \wedge ((n \in \mathbb{N}) \wedge ((a < t) \wedge (t < x)))) \Rightarrow  \\ \quad \quad \operatorname{is\text{-}diff\text{-}at\text{-}v}(m, \operatorname{nth\text{-}deriv\text{-}v}(m, f, n), t, \operatorname{nth\text{-}deriv\text{-}v}(m, f, (n + 1))(t)) \end{array}}}
\vnbprop{r8e-tp-unfold}{\vnbstmt{\begin{array}{@{}l@{}} \forall f, a, d, x. \\ \quad (\operatorname{taylor\text{-}poly}(f, a, d, x) \simeq \operatorname{series\text{-}partial\text{-}sum}((\lambda k \in \mathbb{N}.\; ((\operatorname{nth\text{-}deriv}(f, k)(a) \cdot \operatorname{power}((x - a), k)) \cdot \operatorname{recip}(\operatorname{factorial}(k)))), (d + 1))) \end{array}}}
\vnbprop{r8e-coef-comm}{\vnbstmt{\forall p, q, v \in \mathbb{R}.\; (((p \cdot q) \cdot v) = ((v \cdot p) \cdot q))}}
\vnbprop{r8e-g-of-poly-ind}{\vnbstmt{\begin{array}{@{}l@{}} \forall n \in \mathbb{N}, m, f, g, a, x. \\ \quad \operatorname{is\text{-}normed\text{-}vector\text{-}space}(m) \wedge \\ \quad \operatorname{is\text{-}bounded\text{-}linear\text{-}functional}(m, g) \wedge \\ \quad f \in (\mathbb{R} \to \operatorname{vec}(m)) \wedge \\ \quad a \in \mathbb{R} \wedge \\ \quad x \in \mathbb{R} \wedge \\ \quad \forall j \in \mathbb{N}.\; ((j \leq n) \Rightarrow (\operatorname{nth\text{-}deriv\text{-}v}(m, f, j) \in (\mathbb{R} \to \operatorname{vec}(m)))) \Rightarrow \\ \quad \quad g(\operatorname{taylor\text{-}poly\text{-}v}(m, f, a, x, n)) = \operatorname{taylor\text{-}poly}(\operatorname{compose}(g, f), a, n, x) \end{array}}}
\vnbprop{g-of-remainder-dfun}{\vnbstmt{\begin{array}{@{}l@{}} \forall m, f, g, a, x, n. \\ \quad \operatorname{is\text{-}normed\text{-}vector\text{-}space}(m) \wedge \\ \quad \operatorname{is\text{-}bounded\text{-}linear\text{-}functional}(m, g) \wedge \\ \quad f \in (\mathbb{R} \to \operatorname{vec}(m)) \wedge \\ \quad a \in \mathbb{R} \wedge \\ \quad x \in \mathbb{R} \wedge \\ \quad n \in \mathbb{N} \wedge \\ \quad \forall j \in \mathbb{N}.\; ((j \leq n) \Rightarrow (\operatorname{nth\text{-}deriv\text{-}v}(m, f, j) \in (\mathbb{R} \to \operatorname{vec}(m)))) \Rightarrow \\ \quad \quad g(\operatorname{vadd}(m)(f(x), \operatorname{vneg}(m)(\operatorname{taylor\text{-}poly\text{-}v}(m, f, a, x, n)))) \\ \quad \quad =\; (\operatorname{compose}(g, f)(x) - \operatorname{taylor\text{-}poly}(\operatorname{compose}(g, f), a, n, x)) \end{array}}}
\vnbprop{vector-taylor-remainder-bound}{\vnbstmt{\begin{array}{@{}l@{}} \forall m, f, a, x, n. \\ \quad \operatorname{is\text{-}normed\text{-}vector\text{-}space}(m) \wedge \\ \quad \operatorname{is\text{-}finite\text{-}dimensional}(\operatorname{normed\text{-}vector\text{-}space\text{-}as\text{-}module}(m)) \wedge \\ \quad f \in (\mathbb{R} \to \operatorname{vec}(m)) \wedge \\ \quad a \in \mathbb{R} \wedge \\ \quad x \in \mathbb{R} \wedge \\ \quad n \in \mathbb{N} \wedge \\ \quad a < x \wedge \\ \quad \operatorname{taylor\text{-}differentiable\text{-}v}(m, f, a, x, n) \Rightarrow \\ \quad \quad \exists b \in \mathbb{R}. \\ \quad \quad \quad a < b \wedge \\ \quad \quad \quad b < x \wedge \\ \quad \quad \quad \operatorname{factorial}((n + 1)) \cdot \operatorname{vnrm}(m)(\operatorname{vadd}(m)(f(x), \operatorname{vneg}(m)(\operatorname{taylor\text{-}poly\text{-}v}(m, f, a, x, n)))) \\ \quad \quad \quad \leq\; \operatorname{vnrm}(m)(\operatorname{nth\text{-}deriv\text{-}v}(m, f, (n + 1))(b)) \cdot \operatorname{power}((x - a), (n + 1)) \end{array}}}

\vnbbookfile{theorem-library/rake-nf-norm}{theorem-library/rake-nf-norm}
\noindent{\small rake-nf-norm.scm -- rake batch 5b-D: the norm of a negative, and the metric space underlying a normed field.}\par
\vnbprop{nf-norm-neg}{\vnbstmt{\begin{array}{@{}l@{}} \forall f.; \forall a \in \operatorname{carr}(f). \\ \quad \operatorname{is\text{-}normed\text{-}field}(f) \Rightarrow (\operatorname{fnrm}(f)(\operatorname{neg}(f)(a)) = \operatorname{fnrm}(f)(a)) \end{array}}}
\vnbprop{nf-metric-space-is-metric-space}{\vnbstmt{\forall f.\; (\operatorname{is\text{-}normed\text{-}field}(f) \Rightarrow \operatorname{is\text{-}metric\text{-}space}(\operatorname{nf\text{-}metric\text{-}space}(f)))}}

\vnbbookfile{theorem-library/diff-on-laws}{theorem-library/diff-on-laws}
\noindent{\small diff-on-laws.scm -- the laws of IS-DIFF-ON / DERIV-ON / HOLOMORPHIC-ON (structure-library/diff-on.scm).  Everything here is PROVEN; the definition file asserts nothing.}\par
\vnbprop{diff-on-normed-field}{\vnbstmt{\forall k, u, f, a, l.\; (\operatorname{is\text{-}diff\text{-}on}(k, u, f, a, l) \Rightarrow \operatorname{is\text{-}normed\text{-}field}(k))}}
\vnbprop{diff-on-open}{\vnbstmt{\begin{array}{@{}l@{}} \forall k, u, f, a, l. \\ \quad \operatorname{is\text{-}diff\text{-}on}(k, u, f, a, l) \Rightarrow \operatorname{is\text{-}open}(\operatorname{nf\text{-}metric\text{-}space}(k), u) \end{array}}}
\vnbprop{diff-on-in-fun}{\vnbstmt{\forall k, u, f, a, l.\; (\operatorname{is\text{-}diff\text{-}on}(k, u, f, a, l) \Rightarrow (f \in (u \to \operatorname{carr}(k))))}}
\vnbprop{diff-on-pt-in}{\vnbstmt{\forall k, u, f, a, l.\; (\operatorname{is\text{-}diff\text{-}on}(k, u, f, a, l) \Rightarrow (a \in u))}}
\vnbprop{diff-on-deriv-in-carr}{\vnbstmt{\forall k, u, f, a, l.\; (\operatorname{is\text{-}diff\text{-}on}(k, u, f, a, l) \Rightarrow (l \in \operatorname{carr}(k)))}}
\vnbprop{nf-open-subset-carr}{\vnbstmt{\forall k, u.\; (\operatorname{is\text{-}open}(\operatorname{nf\text{-}metric\text{-}space}(k), u) \Rightarrow (u \subseteq \operatorname{carr}(k)))}}
\vnbprop{diff-on-value-in-carr}{\vnbstmt{\begin{array}{@{}l@{}} \forall k, u, f, a, l.; \forall x \in u. \\ \quad \operatorname{is\text{-}diff\text{-}on}(k, u, f, a, l) \Rightarrow (f(x) \in \operatorname{carr}(k)) \end{array}}}
\vnbprop{rr-nf-carr}{\vnbstmt{(\operatorname{carr}(\operatorname{rr\text{-}normed\text{-}field}) \simeq \mathbb{R})}}
\vnbprop{cc-nf-carr}{\vnbstmt{(\operatorname{carr}(\operatorname{cc\text{-}normed\text{-}field}) \simeq \mathbb{C})}}
\vnbprop{rr-nf-add}{\vnbstmt{(\operatorname{add}(\operatorname{rr\text{-}normed\text{-}field}) \simeq (\lambda \langle \operatorname{x\_}, \operatorname{y\_} \rangle \in (\mathbb{R} \times \mathbb{R}).\; (\operatorname{x\_} + \operatorname{y\_})))}}
\vnbprop{rr-nf-mul}{\vnbstmt{(\operatorname{mul}(\operatorname{rr\text{-}normed\text{-}field}) \simeq (\lambda \langle \operatorname{x\_}, \operatorname{y\_} \rangle \in (\mathbb{R} \times \mathbb{R}).\; (\operatorname{x\_} \cdot \operatorname{y\_})))}}
\vnbprop{rr-nf-neg}{\vnbstmt{(\operatorname{neg}(\operatorname{rr\text{-}normed\text{-}field}) \simeq (\lambda x \in \mathbb{R}.\; -x))}}
\vnbprop{rr-nf-fnrm}{\vnbstmt{(\operatorname{fnrm}(\operatorname{rr\text{-}normed\text{-}field}) \simeq (\lambda x \in \mathbb{R}.\; \lvert x \rvert))}}
\vnbprop{cc-nf-fnrm}{\vnbstmt{(\operatorname{fnrm}(\operatorname{cc\text{-}normed\text{-}field}) \simeq (\lambda x \in \mathbb{C}.\; \operatorname{magnitude}(x)))}}
\vnbprop{cc-nf-zero}{\vnbstmt{(\operatorname{zero}(\operatorname{cc\text{-}normed\text{-}field}) \simeq 0)}}
\vnbprop{rr-nf-zero}{\vnbstmt{(\operatorname{zero}(\operatorname{rr\text{-}normed\text{-}field}) \simeq 0)}}
\vnbprop{rr-nf-add-apply}{\vnbstmt{\begin{array}{@{}l@{}} \forall u, v. \\ \quad ((u \in \mathbb{R}) \wedge (v \in \mathbb{R})) \Rightarrow (\operatorname{add}(\operatorname{rr\text{-}normed\text{-}field})(u, v) \simeq (u + v)) \end{array}}}
\vnbprop{rr-nf-mul-apply}{\vnbstmt{\begin{array}{@{}l@{}} \forall u, v. \\ \quad ((u \in \mathbb{R}) \wedge (v \in \mathbb{R})) \Rightarrow (\operatorname{mul}(\operatorname{rr\text{-}normed\text{-}field})(u, v) \simeq (u \cdot v)) \end{array}}}
\vnbprop{rr-nf-neg-apply}{\vnbstmt{\forall u \in \mathbb{R}.\; (\operatorname{neg}(\operatorname{rr\text{-}normed\text{-}field})(u) \simeq -u)}}
\vnbprop{rr-nf-fnrm-apply}{\vnbstmt{\forall u \in \mathbb{R}.\; (\operatorname{fnrm}(\operatorname{rr\text{-}normed\text{-}field})(u) \simeq \lvert u \rvert)}}
\vnbprop{cc-nf-fnrm-apply}{\vnbstmt{\forall u \in \mathbb{C}.\; (\operatorname{fnrm}(\operatorname{cc\text{-}normed\text{-}field})(u) \simeq \operatorname{magnitude}(u))}}
\vnbprop{ms-agree-ball}{\vnbstmt{\begin{array}{@{}l@{}} \forall a, b.; \forall c \in \operatorname{pts}(a), r. \\ \quad \operatorname{is\text{-}metric\text{-}space}(a) \wedge \\ \quad \operatorname{is\text{-}metric\text{-}space}(b) \wedge \\ \quad (\operatorname{pts}(b) \simeq \operatorname{pts}(a)) \wedge \\ \quad \forall u, v \in \operatorname{pts}(a).\; (\operatorname{dist}(b)(u, v) \simeq \operatorname{dist}(a)(u, v)) \Rightarrow \\ \quad \quad \operatorname{ball}(b, c, r) = \operatorname{ball}(a, c, r) \end{array}}}
\vnbprop{ms-agree-open}{\vnbstmt{\begin{array}{@{}l@{}} \forall a, b, w. \\ \quad \operatorname{is\text{-}metric\text{-}space}(b) \wedge \\ \quad (\operatorname{pts}(b) \simeq \operatorname{pts}(a)) \wedge \\ \quad \forall u, v \in \operatorname{pts}(a).\; (\operatorname{dist}(b)(u, v) \simeq \operatorname{dist}(a)(u, v)) \wedge \\ \quad \operatorname{is\text{-}open}(a, w) \Rightarrow \\ \quad \quad \operatorname{is\text{-}open}(b, w) \end{array}}}
\vnbprop{ms-agree-continuous-at}{\vnbstmt{\begin{array}{@{}l@{}} \forall a, b, \mathit{mta}, \mathit{mtb}, g, p. \\ \quad \operatorname{is\text{-}metric\text{-}space}(b) \wedge \\ \quad \operatorname{is\text{-}metric\text{-}space}(\mathit{mtb}) \wedge \\ \quad (\operatorname{pts}(b) \simeq \operatorname{pts}(a)) \wedge \\ \quad (\operatorname{pts}(\mathit{mtb}) \simeq \operatorname{pts}(\mathit{mta})) \wedge \\ \quad \forall u, v \in \operatorname{pts}(a).\; (\operatorname{dist}(b)(u, v) \simeq \operatorname{dist}(a)(u, v)) \wedge \\ \quad \forall u, v \in \operatorname{pts}(\mathit{mta}).\; (\operatorname{dist}(\mathit{mtb})(u, v) \simeq \operatorname{dist}(\mathit{mta})(u, v)) \wedge \\ \quad \operatorname{is\text{-}continuous\text{-}at}(a, \mathit{mta}, g, p) \Rightarrow \\ \quad \quad \operatorname{is\text{-}continuous\text{-}at}(b, \mathit{mtb}, g, p) \end{array}}}
\vnbprop{ms-agree-pts-sym}{\vnbstmt{\forall a, b.\; ((\operatorname{pts}(b) \simeq \operatorname{pts}(a)) \Rightarrow (\operatorname{pts}(a) \simeq \operatorname{pts}(b)))}}
\vnbprop{ms-agree-dist-sym}{\vnbstmt{\begin{array}{@{}l@{}} \forall a, b.; \forall u, v \in \operatorname{pts}(b). \\ \quad ((\operatorname{pts}(b) \simeq \operatorname{pts}(a)) \wedge \forall u, v \in \operatorname{pts}(a).\; (\operatorname{dist}(b)(u, v) \simeq \operatorname{dist}(a)(u, v))) \Rightarrow  \\ \quad \quad (\operatorname{dist}(a)(u, v) \simeq \operatorname{dist}(b)(u, v)) \end{array}}}
\vnbprop{nf-rr-is-metric-space}{\vnbstmt{\operatorname{is\text{-}metric\text{-}space}(\operatorname{nf\text{-}metric\text{-}space}(\operatorname{rr\text{-}normed\text{-}field}))}}
\vnbprop{nf-rr-agree-pts}{\vnbstmt{(\operatorname{pts}(\operatorname{nf\text{-}metric\text{-}space}(\operatorname{rr\text{-}normed\text{-}field})) \simeq \operatorname{pts}(\operatorname{rr\text{-}ms}))}}
\vnbprop{nf-rr-agree-dist}{\vnbstmt{\begin{array}{@{}l@{}} \forall u, v \in \operatorname{pts}(\operatorname{rr\text{-}ms}). \\ \quad (\operatorname{dist}(\operatorname{nf\text{-}metric\text{-}space}(\operatorname{rr\text{-}normed\text{-}field}))(u, v) \simeq \operatorname{dist}(\operatorname{rr\text{-}ms})(u, v)) \end{array}}}
\vnbprop{rr-open-in-nf}{\vnbstmt{\operatorname{is\text{-}open}(\operatorname{nf\text{-}metric\text{-}space}(\operatorname{rr\text{-}normed\text{-}field}), \mathbb{R})}}
\vnbprop{rr-subset-nf-pts}{\vnbstmt{\mathbb{R} \subseteq \operatorname{pts}(\operatorname{nf\text{-}metric\text{-}space}(\operatorname{rr\text{-}normed\text{-}field}))}}
\vnbprop{sub-rr-is-metric-space}{\vnbstmt{\operatorname{is\text{-}metric\text{-}space}(\operatorname{subspace\text{-}ms}(\operatorname{nf\text{-}metric\text{-}space}(\operatorname{rr\text{-}normed\text{-}field}), \mathbb{R}))}}
\vnbprop{sub-rr-agree-pts}{\vnbstmt{(\operatorname{pts}(\operatorname{subspace\text{-}ms}(\operatorname{nf\text{-}metric\text{-}space}(\operatorname{rr\text{-}normed\text{-}field}), \mathbb{R})) \simeq \operatorname{pts}(\operatorname{rr\text{-}ms}))}}
\vnbprop{sub-rr-agree-dist}{\vnbstmt{\begin{array}{@{}l@{}} \forall u, v \in \operatorname{pts}(\operatorname{rr\text{-}ms}). \\ \quad (\operatorname{dist}(\operatorname{subspace\text{-}ms}(\operatorname{nf\text{-}metric\text{-}space}(\operatorname{rr\text{-}normed\text{-}field}), \mathbb{R}))(u, v) \simeq \operatorname{dist}(\operatorname{rr\text{-}ms})(u, v)) \end{array}}}
\vnbprop{rr-nf-sub}{\vnbstmt{\begin{array}{@{}l@{}} \forall u, v. \\ \quad ((u \in \mathbb{R}) \wedge (v \in \mathbb{R})) \Rightarrow  \\ \quad \quad (\operatorname{add}(\operatorname{rr\text{-}normed\text{-}field})(u, \operatorname{neg}(\operatorname{rr\text{-}normed\text{-}field})(v)) \simeq (u - v)) \end{array}}}
\vnbprop{rr-nf-fun-class}{\vnbstmt{((\mathbb{R} \to \operatorname{carr}(\operatorname{rr\text{-}normed\text{-}field})) \simeq (\mathbb{R} \to \mathbb{R}))}}
\vnbprop{diff-at-iff-diff-on}{\vnbstmt{\begin{array}{@{}l@{}} \forall f, a, l. \\ \quad \operatorname{is\text{-}diff\text{-}at}(f, a, l) \Leftrightarrow \operatorname{is\text{-}diff\text{-}on}(\operatorname{rr\text{-}normed\text{-}field}, \mathbb{R}, f, a, l) \end{array}}}
\vnbprop{nf-norm-in-rr}{\vnbstmt{\forall k.; \forall u \in \operatorname{carr}(k).\; (\operatorname{is\text{-}normed\text{-}field}(k) \Rightarrow (\operatorname{fnrm}(k)(u) \in \mathbb{R}))}}
\vnbprop{nf-norm-nonneg}{\vnbstmt{\forall k.; \forall u \in \operatorname{carr}(k).\; (\operatorname{is\text{-}normed\text{-}field}(k) \Rightarrow (0 \leq \operatorname{fnrm}(k)(u)))}}
\vnbprop{nf-norm-mult}{\vnbstmt{\begin{array}{@{}l@{}} \forall k, u, v. \\ \quad (\operatorname{is\text{-}normed\text{-}field}(k) \wedge ((u \in \operatorname{carr}(k)) \wedge (v \in \operatorname{carr}(k)))) \Rightarrow  \\ \quad \quad \operatorname{fnrm}(k)(\operatorname{mul}(k)(u, v)) = (\operatorname{fnrm}(k)(u) \cdot \operatorname{fnrm}(k)(v)) \end{array}}}
\vnbprop{nf-norm-subadd}{\vnbstmt{\begin{array}{@{}l@{}} \forall k, u, v. \\ \quad (\operatorname{is\text{-}normed\text{-}field}(k) \wedge ((u \in \operatorname{carr}(k)) \wedge (v \in \operatorname{carr}(k)))) \Rightarrow  \\ \quad \quad \operatorname{fnrm}(k)(\operatorname{add}(k)(u, v)) \leq (\operatorname{fnrm}(k)(u) + \operatorname{fnrm}(k)(v)) \end{array}}}
\vnbprop{nf-add-in-carr}{\vnbstmt{\begin{array}{@{}l@{}} \forall k, u, v. \\ \quad (\operatorname{is\text{-}normed\text{-}field}(k) \wedge ((u \in \operatorname{carr}(k)) \wedge (v \in \operatorname{carr}(k)))) \Rightarrow  \\ \quad \quad \operatorname{add}(k)(u, v) \in \operatorname{carr}(k) \end{array}}}
\vnbprop{nf-mul-in-carr}{\vnbstmt{\begin{array}{@{}l@{}} \forall k, u, v. \\ \quad (\operatorname{is\text{-}normed\text{-}field}(k) \wedge ((u \in \operatorname{carr}(k)) \wedge (v \in \operatorname{carr}(k)))) \Rightarrow  \\ \quad \quad \operatorname{mul}(k)(u, v) \in \operatorname{carr}(k) \end{array}}}
\vnbprop{nf-neg-in-carr}{\vnbstmt{\forall k.; \forall u \in \operatorname{carr}(k).\; (\operatorname{is\text{-}normed\text{-}field}(k) \Rightarrow (\operatorname{neg}(k)(u) \in \operatorname{carr}(k)))}}
\vnbprop{nf-sub-add-back}{\vnbstmt{\begin{array}{@{}l@{}} \forall k, u, v. \\ \quad (\operatorname{is\text{-}normed\text{-}field}(k) \wedge ((u \in \operatorname{carr}(k)) \wedge (v \in \operatorname{carr}(k)))) \Rightarrow  \\ \quad \quad \operatorname{add}(k)(\operatorname{add}(k)(u, \operatorname{neg}(k)(v)), v) = u \end{array}}}
\vnbprop{nf-norm-le-add}{\vnbstmt{\begin{array}{@{}l@{}} \forall k, u, v. \\ \quad (\operatorname{is\text{-}normed\text{-}field}(k) \wedge ((u \in \operatorname{carr}(k)) \wedge (v \in \operatorname{carr}(k)))) \Rightarrow  \\ \quad \quad \operatorname{fnrm}(k)(u) \leq (\operatorname{fnrm}(k)(\operatorname{add}(k)(u, \operatorname{neg}(k)(v))) + \operatorname{fnrm}(k)(v)) \end{array}}}
\vnbprop{diff-on-implies-continuous}{\vnbstmt{\begin{array}{@{}l@{}} \forall k, u, f, a, l. \\ \quad \operatorname{is\text{-}diff\text{-}on}(k, u, f, a, l) \Rightarrow  \\ \quad \quad \operatorname{is\text{-}continuous\text{-}at}(\operatorname{subspace\text{-}ms}(\operatorname{nf\text{-}metric\text{-}space}(k), u), \operatorname{nf\text{-}metric\text{-}space}(k), f, a) \end{array}}}
\vnbprop{deriv-on-unfold}{\vnbstmt{\forall k, u, f, a.\; (\operatorname{deriv\text{-}on}(k, u, f, a) \simeq (\iota\, l.\;\operatorname{is\text{-}diff\text{-}on}(k, u, f, a, l)))}}
\vnbprop{holomorphic-on-open}{\vnbstmt{\begin{array}{@{}l@{}} \forall u, f. \\ \quad \operatorname{holomorphic\text{-}on}(u, f) \Rightarrow \operatorname{is\text{-}open}(\operatorname{nf\text{-}metric\text{-}space}(\operatorname{cc\text{-}normed\text{-}field}), u) \end{array}}}
\vnbprop{holomorphic-on-in-fun}{\vnbstmt{\forall u, f.\; (\operatorname{holomorphic\text{-}on}(u, f) \Rightarrow (f \in (u \to \mathbb{C})))}}
\vnbprop{holomorphic-on-diff}{\vnbstmt{\begin{array}{@{}l@{}} \forall u, f.; \forall a \in u. \\ \quad \operatorname{holomorphic\text{-}on}(u, f) \Rightarrow  \\ \quad \quad \exists l.\; \operatorname{is\text{-}diff\text{-}on}(\operatorname{cc\text{-}normed\text{-}field}, u, f, a, l) \end{array}}}

\vnbbookfile{theorem-library/ms-continuity-algebra}{theorem-library/ms-continuity-algebra}
\noindent{\small ms-continuity-algebra.scm -- THE CONTINUITY ALGEBRA OFF THE REAL LINE.}\par
\vnbprop{ms-pts-is-set}{\vnbstmt{\forall s.\; (\operatorname{is\text{-}metric\text{-}space}(s) \Rightarrow (\operatorname{pts}(s) \in \mathbf{Set}))}}
\vnbprop{ms-const-lam-in-fun}{\vnbstmt{\begin{array}{@{}l@{}} \forall s, t.; \forall c \in \operatorname{pts}(t). \\ \quad \operatorname{is\text{-}metric\text{-}space}(s) \Rightarrow ((\lambda z \in \operatorname{pts}(s).\; c) \in (\operatorname{pts}(s) \to \operatorname{pts}(t))) \end{array}}}
\vnbprop{ms-const-continuous-at}{\vnbstmt{\begin{array}{@{}l@{}} \forall s, t.; \forall c \in \operatorname{pts}(t), a \in \operatorname{pts}(s). \\ \quad (\operatorname{is\text{-}metric\text{-}space}(s) \wedge \operatorname{is\text{-}metric\text{-}space}(t)) \Rightarrow  \\ \quad \quad \operatorname{is\text{-}continuous\text{-}at}(s, t, (\lambda z \in \operatorname{pts}(s).\; c), a) \end{array}}}
\vnbprop{ms-ident-lam-in-fun}{\vnbstmt{\begin{array}{@{}l@{}} \forall s. \\ \quad \operatorname{is\text{-}metric\text{-}space}(s) \Rightarrow ((\lambda z \in \operatorname{pts}(s).\; z) \in (\operatorname{pts}(s) \to \operatorname{pts}(s))) \end{array}}}
\vnbprop{ms-identity-continuous-at}{\vnbstmt{\begin{array}{@{}l@{}} \forall s.; \forall a \in \operatorname{pts}(s). \\ \quad \operatorname{is\text{-}metric\text{-}space}(s) \Rightarrow \operatorname{is\text{-}continuous\text{-}at}(s, s, (\lambda z \in \operatorname{pts}(s).\; z), a) \end{array}}}
\vnbprop{ms-cont-transfer-ptwise-eq}{\vnbstmt{\begin{array}{@{}l@{}} \forall s, t, g, f, a. \\ \quad (\operatorname{is\text{-}continuous\text{-}at}(s, t, g, a) \wedge ((f \in (\operatorname{pts}(s) \to \operatorname{pts}(t))) \wedge \forall z \in \operatorname{pts}(s).\; (f(z) = g(z)))) \Rightarrow  \\ \quad \quad \operatorname{is\text{-}continuous\text{-}at}(s, t, f, a) \end{array}}}
\vnbprop{ms-compose-continuous-at}{\vnbstmt{\begin{array}{@{}l@{}} \forall s, t, u, f, g, a. \\ \quad (\operatorname{is\text{-}continuous\text{-}at}(s, t, f, a) \wedge \operatorname{is\text{-}continuous\text{-}at}(t, u, g, f(a))) \Rightarrow  \\ \quad \quad \operatorname{is\text{-}continuous\text{-}at}(s, u, (\lambda z \in \operatorname{pts}(s).\; g(f(z))), a) \end{array}}}
\vnbprop{ms-cont-agree-off-pt}{\vnbstmt{\begin{array}{@{}l@{}} \forall s, t, f, g, p. \\ \quad \operatorname{is\text{-}continuous\text{-}at}(s, t, f, p) \wedge \\ \quad \operatorname{is\text{-}continuous\text{-}at}(s, t, g, p) \wedge \\ \quad \forall r. \\ \quad \quad \operatorname{pos\text{-}rr}(r) \Rightarrow \exists b \in \operatorname{pts}(s).\; (\neg (b = p) \wedge (\operatorname{dist}(s)(p, b) \leq r)) \wedge \\ \quad \forall y \in \operatorname{pts}(s).\; (\neg (y = p) \Rightarrow (f(y) = g(y))) \Rightarrow \\ \quad \quad f(p) = g(p) \end{array}}}
\vnbprop{nf-sub-split}{\vnbstmt{\begin{array}{@{}l@{}} \forall k, u, v, p, q. \\ \quad (\operatorname{is\text{-}normed\text{-}field}(k) \wedge ((u \in \operatorname{carr}(k)) \wedge ((v \in \operatorname{carr}(k)) \wedge ((p \in \operatorname{carr}(k)) \wedge (q \in \operatorname{carr}(k)))))) \Rightarrow  \\ \quad \quad \operatorname{add}(k)(\operatorname{add}(k)(u, v), \operatorname{neg}(k)(\operatorname{add}(k)(p, q))) \\ \quad \quad =\; \operatorname{add}(k)(\operatorname{add}(k)(u, \operatorname{neg}(k)(p)), \operatorname{add}(k)(v, \operatorname{neg}(k)(q))) \end{array}}}
\vnbprop{nf-mul-sub-split}{\vnbstmt{\begin{array}{@{}l@{}} \forall k, u, v, p, q. \\ \quad (\operatorname{is\text{-}normed\text{-}field}(k) \wedge ((u \in \operatorname{carr}(k)) \wedge ((v \in \operatorname{carr}(k)) \wedge ((p \in \operatorname{carr}(k)) \wedge (q \in \operatorname{carr}(k)))))) \Rightarrow  \\ \quad \quad \operatorname{add}(k)(\operatorname{mul}(k)(u, v), \operatorname{neg}(k)(\operatorname{mul}(k)(p, q))) \\ \quad \quad =\; \operatorname{add}(k)(\operatorname{mul}(k)(u, \operatorname{add}(k)(v, \operatorname{neg}(k)(q))), \operatorname{mul}(k)(q, \operatorname{add}(k)(u, \operatorname{neg}(k)(p)))) \end{array}}}
\vnbprop{nf-neg-sub-split}{\vnbstmt{\begin{array}{@{}l@{}} \forall k, u, v. \\ \quad (\operatorname{is\text{-}normed\text{-}field}(k) \wedge ((u \in \operatorname{carr}(k)) \wedge (v \in \operatorname{carr}(k)))) \Rightarrow  \\ \quad \quad \operatorname{add}(k)(\operatorname{neg}(k)(u), \operatorname{neg}(k)(\operatorname{neg}(k)(v))) = \operatorname{neg}(k)(\operatorname{add}(k)(u, \operatorname{neg}(k)(v))) \end{array}}}
\vnbprop{nf-sub-anti}{\vnbstmt{\begin{array}{@{}l@{}} \forall k, u, v. \\ \quad (\operatorname{is\text{-}normed\text{-}field}(k) \wedge ((u \in \operatorname{carr}(k)) \wedge (v \in \operatorname{carr}(k)))) \Rightarrow  \\ \quad \quad \operatorname{neg}(k)(\operatorname{add}(k)(u, \operatorname{neg}(k)(v))) = \operatorname{add}(k)(v, \operatorname{neg}(k)(u)) \end{array}}}
\vnbprop{nf-norm-sub-sym}{\vnbstmt{\begin{array}{@{}l@{}} \forall k, u, v. \\ \quad (\operatorname{is\text{-}normed\text{-}field}(k) \wedge ((u \in \operatorname{carr}(k)) \wedge (v \in \operatorname{carr}(k)))) \Rightarrow  \\ \quad \quad \operatorname{fnrm}(k)(\operatorname{add}(k)(u, \operatorname{neg}(k)(v))) = \operatorname{fnrm}(k)(\operatorname{add}(k)(v, \operatorname{neg}(k)(u))) \end{array}}}
\vnbprop{nf-norm-neg-diff}{\vnbstmt{\begin{array}{@{}l@{}} \forall k, u, v. \\ \quad (\operatorname{is\text{-}normed\text{-}field}(k) \wedge ((u \in \operatorname{carr}(k)) \wedge (v \in \operatorname{carr}(k)))) \Rightarrow  \\ \quad \quad \operatorname{fnrm}(k)(\operatorname{add}(k)(\operatorname{neg}(k)(u), \operatorname{neg}(k)(\operatorname{neg}(k)(v)))) = \operatorname{fnrm}(k)(\operatorname{add}(k)(u, \operatorname{neg}(k)(v))) \end{array}}}
\vnbprop{nf-norm-sum-diff}{\vnbstmt{\begin{array}{@{}l@{}} \forall k, u, v, p, q. \\ \quad (\operatorname{is\text{-}normed\text{-}field}(k) \wedge ((u \in \operatorname{carr}(k)) \wedge ((v \in \operatorname{carr}(k)) \wedge ((p \in \operatorname{carr}(k)) \wedge (q \in \operatorname{carr}(k)))))) \Rightarrow  \\ \quad \quad \operatorname{fnrm}(k)(\operatorname{add}(k)(\operatorname{add}(k)(u, v), \operatorname{neg}(k)(\operatorname{add}(k)(p, q)))) \\ \quad \quad \leq\; \operatorname{fnrm}(k)(\operatorname{add}(k)(u, \operatorname{neg}(k)(p))) + \operatorname{fnrm}(k)(\operatorname{add}(k)(v, \operatorname{neg}(k)(q))) \end{array}}}
\vnbprop{nf-norm-prod-bound}{\vnbstmt{\begin{array}{@{}l@{}} \forall k, u, v, p, q, m, \mathit{nfk}, t, e. \\ \quad \operatorname{is\text{-}normed\text{-}field}(k) \wedge \\ \quad u \in \operatorname{carr}(k) \wedge \\ \quad v \in \operatorname{carr}(k) \wedge \\ \quad p \in \operatorname{carr}(k) \wedge \\ \quad q \in \operatorname{carr}(k) \wedge \\ \quad m \in \mathbb{R} \wedge \\ \quad \mathit{nfk} \in \mathbb{R} \wedge \\ \quad t \in \mathbb{R} \wedge \\ \quad e \in \mathbb{R} \wedge \\ \quad \operatorname{fnrm}(k)(u) \leq m \wedge \\ \quad \operatorname{fnrm}(k)(q) \leq \mathit{nfk} \wedge \\ \quad \operatorname{fnrm}(k)(\operatorname{add}(k)(u, \operatorname{neg}(k)(p))) \leq t \wedge \\ \quad \operatorname{fnrm}(k)(\operatorname{add}(k)(v, \operatorname{neg}(k)(q))) \leq t \wedge \\ \quad ((m + \mathit{nfk}) \cdot t) \leq e \Rightarrow \\ \quad \quad \operatorname{fnrm}(k)(\operatorname{add}(k)(\operatorname{mul}(k)(u, v), \operatorname{neg}(k)(\operatorname{mul}(k)(p, q)))) \leq e \end{array}}}
\vnbprop{nf-sum-lam-in-fun}{\vnbstmt{\begin{array}{@{}l@{}} \forall k, s.; \forall g, h \in (\operatorname{pts}(s) \to \operatorname{pts}(\operatorname{nf\text{-}metric\text{-}space}(k))). \\ \quad (\operatorname{is\text{-}normed\text{-}field}(k) \wedge \operatorname{is\text{-}metric\text{-}space}(s)) \Rightarrow  \\ \quad \quad (\lambda z \in \operatorname{pts}(s).\; \operatorname{add}(k)(g(z), h(z))) \in (\operatorname{pts}(s) \to \operatorname{pts}(\operatorname{nf\text{-}metric\text{-}space}(k))) \end{array}}}
\vnbprop{nf-prod-lam-in-fun}{\vnbstmt{\begin{array}{@{}l@{}} \forall k, s.; \forall g, h \in (\operatorname{pts}(s) \to \operatorname{pts}(\operatorname{nf\text{-}metric\text{-}space}(k))). \\ \quad (\operatorname{is\text{-}normed\text{-}field}(k) \wedge \operatorname{is\text{-}metric\text{-}space}(s)) \Rightarrow  \\ \quad \quad (\lambda z \in \operatorname{pts}(s).\; \operatorname{mul}(k)(g(z), h(z))) \in (\operatorname{pts}(s) \to \operatorname{pts}(\operatorname{nf\text{-}metric\text{-}space}(k))) \end{array}}}
\vnbprop{nf-neg-lam-in-fun}{\vnbstmt{\begin{array}{@{}l@{}} \forall k, s.; \forall g \in (\operatorname{pts}(s) \to \operatorname{pts}(\operatorname{nf\text{-}metric\text{-}space}(k))). \\ \quad (\operatorname{is\text{-}normed\text{-}field}(k) \wedge \operatorname{is\text{-}metric\text{-}space}(s)) \Rightarrow  \\ \quad \quad (\lambda z \in \operatorname{pts}(s).\; \operatorname{neg}(k)(g(z))) \in (\operatorname{pts}(s) \to \operatorname{pts}(\operatorname{nf\text{-}metric\text{-}space}(k))) \end{array}}}
\vnbprop{nf-sum-continuous-at}{\vnbstmt{\begin{array}{@{}l@{}} \forall k, s, g, h, a. \\ \quad \operatorname{is\text{-}normed\text{-}field}(k) \wedge \\ \quad \operatorname{is\text{-}continuous\text{-}at}(s, \operatorname{nf\text{-}metric\text{-}space}(k), g, a) \wedge \\ \quad \operatorname{is\text{-}continuous\text{-}at}(s, \operatorname{nf\text{-}metric\text{-}space}(k), h, a) \Rightarrow \\ \quad \quad \operatorname{is\text{-}continuous\text{-}at}(s, \operatorname{nf\text{-}metric\text{-}space}(k), (\lambda z \in \operatorname{pts}(s).\; \operatorname{add}(k)(g(z), h(z))), a) \end{array}}}
\vnbprop{nf-neg-continuous-at}{\vnbstmt{\begin{array}{@{}l@{}} \forall k, s, g, a. \\ \quad (\operatorname{is\text{-}normed\text{-}field}(k) \wedge \operatorname{is\text{-}continuous\text{-}at}(s, \operatorname{nf\text{-}metric\text{-}space}(k), g, a)) \Rightarrow  \\ \quad \quad \operatorname{is\text{-}continuous\text{-}at}(s, \operatorname{nf\text{-}metric\text{-}space}(k), (\lambda z \in \operatorname{pts}(s).\; \operatorname{neg}(k)(g(z))), a) \end{array}}}
\vnbprop{nf-product-continuous-at}{\vnbstmt{\begin{array}{@{}l@{}} \forall k, s, g, h, a. \\ \quad \operatorname{is\text{-}normed\text{-}field}(k) \wedge \\ \quad \operatorname{is\text{-}continuous\text{-}at}(s, \operatorname{nf\text{-}metric\text{-}space}(k), g, a) \wedge \\ \quad \operatorname{is\text{-}continuous\text{-}at}(s, \operatorname{nf\text{-}metric\text{-}space}(k), h, a) \Rightarrow \\ \quad \quad \operatorname{is\text{-}continuous\text{-}at}(s, \operatorname{nf\text{-}metric\text{-}space}(k), (\lambda z \in \operatorname{pts}(s).\; \operatorname{mul}(k)(g(z), h(z))), a) \end{array}}}

\vnbbookfile{structure-library/riesz-markov-model}{structure-library/riesz-markov-model}
\noindent{\small structure-library/riesz-markov-model.scm -- the Cantor model of a compact metric space (Riesz-Markov, stage 1, RM-1, 2026-10-10).  DEFINITIONS ONLY; every law is PROVEN in theorem-library/riesz-markov-model.scm (prefix rmm-).}\par
\par\medskip\noindent\textbf{Definition} (\texttt{rm-ball}).\;$\operatorname{rm\text{-}ball}(s, \mathit{dseq}, n) :=$ $\operatorname{closed\text{-}ball}(s, \operatorname{dseq}(\operatorname{nnfst}(n)), \operatorname{recip}((\operatorname{nnsnd}(n) + 1)))$
\par\medskip\noindent\textbf{Definition} (\texttt{rm-shell}).\;$\operatorname{rm\text{-}shell}(s, \mathit{dseq}, n) :=$
\begin{equation*}
\{\operatorname{rsy\_} \in \operatorname{pts}(s) : (\operatorname{recip}((\operatorname{nnsnd}(n) + 1)) \leq \operatorname{dist}(s)(\operatorname{dseq}(\operatorname{nnfst}(n)), \operatorname{rsy\_}))\}
\end{equation*}

\par\medskip\noindent\textbf{Definition} (\texttt{rm-two}).\; $(\operatorname{rm\text{-}two} \simeq \{0, 1\})$
\par\medskip\noindent\textbf{Definition} (\texttt{rm-two-ms}).\; $(\operatorname{rm\text{-}two\text{-}ms} \simeq \operatorname{subspace\text{-}ms}(\operatorname{rr\text{-}ms}, \operatorname{rm\text{-}two}))$
\par\medskip\noindent\textbf{Definition} (\texttt{cantor-ms}).\; $(\operatorname{cantor\text{-}ms} \simeq \operatorname{product\text{-}metric}((\lambda \operatorname{ctk\_} \in \mathbb{N}.\; \operatorname{rm\text{-}two\text{-}ms})))$
\par\medskip\noindent\textbf{Definition} (\texttt{rm-realizes}).\;For all $s$, $\mathit{dseq}$, $\mathit{es}$, $\mathit{nb}$ and $y$, 
\begin{equation*}
\begin{array}{@{}l@{}} \operatorname{rm\text{-}realizes}(s, \mathit{dseq}, \mathit{es}, \mathit{nb}, y) \\ \Leftrightarrow\; ((y \in \operatorname{pts}(s)) \wedge \forall \operatorname{rri\_} \in \mathbb{N}.\; ((\operatorname{rri\_} < \mathit{nb}) \Rightarrow ((y \in \operatorname{rm\text{-}ball}(s, \mathit{dseq}, \operatorname{rri\_})) \Leftrightarrow (\operatorname{es}(\operatorname{rri\_}) = 1)))) \end{array}
\end{equation*}

\par\medskip\noindent\textbf{Definition} (\texttt{rm-k}).\;$\operatorname{rm\text{-}k}(s, \mathit{dseq}) :=$
\begin{equation*}
\{\operatorname{rke\_} \in \operatorname{pts}(\operatorname{cantor\text{-}ms}) : \forall \operatorname{rkn\_} \in \mathbb{N}.\; \exists \operatorname{rky\_}.\; \operatorname{rm\text{-}realizes}(s, \mathit{dseq}, \operatorname{rke\_}, \operatorname{rkn\_}, \operatorname{rky\_})\}
\end{equation*}

\par\medskip\noindent\textbf{Definition} (\texttt{rm-kms}).\;$\operatorname{rm\text{-}kms}(s, \mathit{dseq}) :=$ $\operatorname{subspace\text{-}ms}(\operatorname{cantor\text{-}ms}, \operatorname{rm\text{-}k}(s, \mathit{dseq}))$
\par\medskip\noindent\textbf{Definition} (\texttt{rm-address}).\;For all $s$, $\mathit{dseq}$, $\mathit{es}$ and $y$, 
\begin{equation*}
\begin{array}{@{}l@{}} \operatorname{rm\text{-}address}(s, \mathit{dseq}, \mathit{es}, y) \\ \Leftrightarrow\; y \in \operatorname{pts}(s) \wedge \\ \forall \operatorname{ran\_} \in \mathbb{N}. \\ \quad (\operatorname{es}(\operatorname{ran\_}) = 1) \Rightarrow (y \in \operatorname{rm\text{-}ball}(s, \mathit{dseq}, \operatorname{ran\_})) \wedge \\ \quad \neg (\operatorname{es}(\operatorname{ran\_}) = 1) \Rightarrow (y \in \operatorname{rm\text{-}shell}(s, \mathit{dseq}, \operatorname{ran\_})) \end{array}
\end{equation*}

\par\medskip\noindent\textbf{Definition} (\texttt{rm-address-below}).\;For all $s$, $\mathit{dseq}$, $\mathit{es}$, $\mathit{nb}$ and $y$, 
\begin{equation*}
\begin{array}{@{}l@{}} \operatorname{rm\text{-}address\text{-}below}(s, \mathit{dseq}, \mathit{es}, \mathit{nb}, y) \\ \Leftrightarrow\; y \in \operatorname{pts}(s) \wedge \\ \forall \operatorname{rbn\_} \in \mathbb{N}. \\ \quad (\operatorname{rbn\_} < \mathit{nb}) \Rightarrow  \\ \quad \quad (\operatorname{es}(\operatorname{rbn\_}) = 1) \Rightarrow (y \in \operatorname{rm\text{-}ball}(s, \mathit{dseq}, \operatorname{rbn\_})) \wedge \\ \quad \quad \neg (\operatorname{es}(\operatorname{rbn\_}) = 1) \Rightarrow (y \in \operatorname{rm\text{-}shell}(s, \mathit{dseq}, \operatorname{rbn\_})) \end{array}
\end{equation*}

\par\medskip\noindent\textbf{Definition} (\texttt{rm-pi}).\;$\operatorname{rm\text{-}pi}(s, \mathit{dseq}) :=$
\begin{equation*}
(\lambda \operatorname{rpe\_} \in \operatorname{rm\text{-}k}(s, \mathit{dseq}).\; (\iota\, \operatorname{rpy\_}.\;\operatorname{rm\text{-}address}(s, \mathit{dseq}, \operatorname{rpe\_}, \operatorname{rpy\_})))
\end{equation*}

\par\medskip\noindent\textbf{Definition} (\texttt{rm-code}).\;$\operatorname{rm\text{-}code}(s, \mathit{dseq}, x) :=$ $(\lambda \operatorname{rox\_} \in \mathbb{N}.\; \left(\text{if }(x \in \operatorname{rm\text{-}ball}(s, \mathit{dseq}, \operatorname{rox\_}))\text{ then }1\text{ else }0\right))$
\par\medskip\noindent\textbf{Definition} (\texttt{rm-cyl}).\;$\operatorname{rm\text{-}cyl}(s, \mathit{dseq}, a, \mathit{nb}) :=$
\begin{equation*}
\{\operatorname{rce\_} \in \operatorname{rm\text{-}k}(s, \mathit{dseq}) : \forall \operatorname{rci\_} \in \mathbb{N}.\; ((\operatorname{rci\_} < \mathit{nb}) \Rightarrow (\operatorname{rce\_}(\operatorname{rci\_}) = a(\operatorname{rci\_})))\}
\end{equation*}


\vnbbookfile{theorem-library/riesz-markov-model}{theorem-library/riesz-markov-model}
\noindent{\small theorem-library/riesz-markov-model.scm -- RIESZ-MARKOV, STAGE 1 (RM-1, 2026-10-10): the Cantor model K of a compact metric space and the map pi : K -$>$ s.  Vocabulary: structure-library/riesz-markov-model.scm. Source: the user's notes-63, concrete form (b).  Helper prefix: rmm-.}\par
\vnbprop{rmm-two-mem}{\vnbstmt{\forall x.\; ((x \in \operatorname{rm\text{-}two}) \Leftrightarrow ((x = 0) \vee (x = 1)))}}
\vnbprop{rmm-two-cases}{\vnbstmt{\forall x \in \operatorname{rm\text{-}two}.\; ((x = 0) \vee (x = 1))}}
\vnbprop{rmm-two-makeset}{\vnbstmt{\operatorname{rm\text{-}two} = \operatorname{make\text{-}set}(\langle 0, 1 \rangle)}}
\vnbprop{rmm-two-sub-rr}{\vnbstmt{\operatorname{rm\text{-}two} \subseteq \mathbb{R}}}
\vnbprop{rmm-two-ms-metric}{\vnbstmt{\operatorname{is\text{-}metric\text{-}space}(\operatorname{rm\text{-}two\text{-}ms})}}
\vnbprop{rmm-two-ms-pts}{\vnbstmt{(\operatorname{pts}(\operatorname{rm\text{-}two\text{-}ms}) \simeq \operatorname{rm\text{-}two})}}
\vnbprop{rmm-two-ms-dist}{\vnbstmt{\forall u, v \in \operatorname{rm\text{-}two}.\; (\operatorname{dist}(\operatorname{rm\text{-}two\text{-}ms})(u, v) \simeq \lvert (u - v) \rvert)}}
\vnbprop{rmm-two-near-eq}{\vnbstmt{\forall u, v \in \operatorname{rm\text{-}two}.\; ((\operatorname{dist}(\operatorname{rm\text{-}two\text{-}ms})(u, v) < 1) \Rightarrow (u = v))}}
\vnbprop{rmm-two-ms-set}{\vnbstmt{\operatorname{rm\text{-}two\text{-}ms} \in \mathbf{Set}}}
\vnbprop{rmm-ms-seq}{\vnbstmt{\operatorname{is\text{-}ms\text{-}sequence}((\lambda k \in \mathbb{N}.\; \operatorname{rm\text{-}two\text{-}ms}))}}
\vnbprop{rmm-cantor-metric}{\vnbstmt{\operatorname{is\text{-}metric\text{-}space}(\operatorname{cantor\text{-}ms})}}
\vnbprop{rmm-cantor-pts}{\vnbstmt{(\operatorname{pts}(\operatorname{cantor\text{-}ms}) \simeq \operatorname{product\text{-}carrier}((\lambda k \in \mathbb{N}.\; \operatorname{rm\text{-}two\text{-}ms})))}}
\vnbprop{rmm-cantor-coord}{\vnbstmt{\forall s \in \operatorname{pts}(\operatorname{cantor\text{-}ms}), n \in \mathbb{N}.\; (s(n) \in \operatorname{rm\text{-}two})}}
\vnbprop{rmm-two-seq-compact}{\vnbstmt{\operatorname{seq\text{-}compact}(\operatorname{rm\text{-}two\text{-}ms})}}
\vnbprop{rmm-two-compact}{\vnbstmt{\operatorname{is\text{-}compact}(\operatorname{rm\text{-}two\text{-}ms})}}
\vnbprop{rmm-cantor-compact}{\vnbstmt{\operatorname{is\text{-}compact}(\operatorname{cantor\text{-}ms})}}
\vnbprop{rmm-cantor-unfold}{\vnbstmt{(\operatorname{cantor\text{-}ms} \simeq \operatorname{product\text{-}metric\text{-}w}((\lambda k \in \mathbb{N}.\; \operatorname{rm\text{-}two\text{-}ms}), (\lambda n \in \mathbb{N}.\; /(1, \operatorname{power}(2, (n + 1))))))}}
\vnbprop{rmm-proj-value}{\vnbstmt{\begin{array}{@{}l@{}} \forall n \in \mathbb{N}, s \in \operatorname{pts}(\operatorname{cantor\text{-}ms}). \\ \quad \operatorname{product\text{-}proj}((\lambda k \in \mathbb{N}.\; \operatorname{rm\text{-}two\text{-}ms}), n)(s) = s(n) \end{array}}}
\vnbprop{rmm-coord-continuous}{\vnbstmt{\begin{array}{@{}l@{}} \forall n \in \mathbb{N}. \\ \quad \operatorname{is\text{-}continuous}(\operatorname{cantor\text{-}ms}, \operatorname{rm\text{-}two\text{-}ms}, \operatorname{product\text{-}proj}((\lambda k \in \mathbb{N}.\; \operatorname{rm\text{-}two\text{-}ms}), n)) \end{array}}}
\vnbprop{rmm-coord-near}{\vnbstmt{\begin{array}{@{}l@{}} \forall n \in \mathbb{N}, s \in \operatorname{pts}(\operatorname{cantor\text{-}ms}). \exists d. \\ \quad (\operatorname{pos\text{-}rr}(d) \wedge \forall b \in \operatorname{pts}(\operatorname{cantor\text{-}ms}).\; ((\operatorname{dist}(\operatorname{cantor\text{-}ms})(s, b) \leq d) \Rightarrow (b(n) = s(n)))) \end{array}}}
\vnbprop{rmm-agree-near}{\vnbstmt{\begin{array}{@{}l@{}} \forall b \in \mathbb{N}, s \in \operatorname{pts}(\operatorname{cantor\text{-}ms}). \exists d. \\ \quad \operatorname{pos\text{-}rr}(d) \wedge \\ \quad \forall \mathit{rab} \in \operatorname{pts}(\operatorname{cantor\text{-}ms}), i \in \mathbb{N}. \\ \quad \quad ((\operatorname{dist}(\operatorname{cantor\text{-}ms})(s, \mathit{rab}) \leq d) \wedge (i < b)) \Rightarrow (\operatorname{rab}(i) = s(i)) \end{array}}}
\vnbprop{rmm-dense-parts}{\vnbstmt{\begin{array}{@{}l@{}} \forall s, q. \\ \quad \operatorname{is\text{-}dense\text{-}seq}(s, q) \Rightarrow (\operatorname{is\text{-}metric\text{-}space}(s) \wedge (q \in (\mathbb{N} \to \operatorname{pts}(s)))) \end{array}}}
\vnbprop{rmm-dense-near}{\vnbstmt{\begin{array}{@{}l@{}} \forall s, q.; \forall x \in \operatorname{pts}(s), a. \\ \quad (\operatorname{is\text{-}dense\text{-}seq}(s, q) \wedge \operatorname{pos\text{-}rr}(a)) \Rightarrow  \\ \quad \quad \exists m \in \mathbb{N}.\; (\operatorname{dist}(s)(x, q(m)) < a) \end{array}}}
\vnbprop{rmm-centre-in}{\vnbstmt{\forall s, q.; \forall n \in \mathbb{N}.\; (\operatorname{is\text{-}dense\text{-}seq}(s, q) \Rightarrow (q(\operatorname{nnfst}(n)) \in \operatorname{pts}(s)))}}
\vnbprop{rmm-rad-pos}{\vnbstmt{\forall n \in \mathbb{N}.\; \operatorname{pos\text{-}rr}(\operatorname{recip}((\operatorname{nnsnd}(n) + 1)))}}
\vnbprop{rmm-ball-mem}{\vnbstmt{\begin{array}{@{}l@{}} \forall s, q, n, y. \\ \quad (y \in \operatorname{rm\text{-}ball}(s, q, n)) \Leftrightarrow ((y \in \operatorname{pts}(s)) \wedge (\operatorname{dist}(s)(q(\operatorname{nnfst}(n)), y) \leq \operatorname{recip}((\operatorname{nnsnd}(n) + 1)))) \end{array}}}
\vnbprop{rmm-shell-mem}{\vnbstmt{\begin{array}{@{}l@{}} \forall s, q, n, y. \\ \quad (y \in \operatorname{rm\text{-}shell}(s, q, n)) \Leftrightarrow ((y \in \operatorname{pts}(s)) \wedge (\operatorname{recip}((\operatorname{nnsnd}(n) + 1)) \leq \operatorname{dist}(s)(q(\operatorname{nnfst}(n)), y))) \end{array}}}
\vnbprop{rmm-near-centre}{\vnbstmt{\begin{array}{@{}l@{}} \forall s, q.; \forall x \in \operatorname{pts}(s), j \in \mathbb{N}. \\ \quad \operatorname{is\text{-}dense\text{-}seq}(s, q) \Rightarrow  \\ \quad \quad \exists n \in \mathbb{N}. \\ \quad \quad \quad (\operatorname{nnsnd}(n) = j) \wedge (\operatorname{dist}(s)(q(\operatorname{nnfst}(n)), x) < \operatorname{recip}((\operatorname{nnsnd}(n) + 1))) \end{array}}}
\vnbprop{rmm-inside-ball}{\vnbstmt{\begin{array}{@{}l@{}} \forall s, q.; \forall n \in \mathbb{N}, x \in \operatorname{pts}(s). \\ \quad (\operatorname{is\text{-}dense\text{-}seq}(s, q) \wedge (\operatorname{dist}(s)(q(\operatorname{nnfst}(n)), x) < \operatorname{recip}((\operatorname{nnsnd}(n) + 1)))) \Rightarrow  \\ \quad \quad (x \in \operatorname{rm\text{-}ball}(s, q, n)) \wedge \neg (x \in \operatorname{rm\text{-}shell}(s, q, n)) \end{array}}}
\vnbprop{rmm-strong-separation}{\vnbstmt{\begin{array}{@{}l@{}} \forall s, q.; \forall x, y \in \operatorname{pts}(s). \\ \quad (\operatorname{is\text{-}dense\text{-}seq}(s, q) \wedge \neg (x = y)) \Rightarrow  \\ \quad \quad \exists n, m \in \mathbb{N}. \\ \quad \quad \quad x \in \operatorname{rm\text{-}ball}(s, q, n) \wedge \\ \quad \quad \quad \neg (x \in \operatorname{rm\text{-}shell}(s, q, n)) \wedge \\ \quad \quad \quad y \in \operatorname{rm\text{-}ball}(s, q, m) \wedge \\ \quad \quad \quad \neg (y \in \operatorname{rm\text{-}shell}(s, q, m)) \wedge \\ \quad \quad \quad \forall z.\; \neg ((z \in \operatorname{rm\text{-}ball}(s, q, n)) \wedge (z \in \operatorname{rm\text{-}ball}(s, q, m))) \end{array}}}
\vnbprop{rmm-k-mem}{\vnbstmt{\begin{array}{@{}l@{}} \forall s, q, \mathit{es}. \\ \quad \mathit{es} \in \operatorname{rm\text{-}k}(s, q) \\ \quad \Leftrightarrow\; ((\mathit{es} \in \operatorname{pts}(\operatorname{cantor\text{-}ms})) \wedge \forall n \in \mathbb{N}.\; \exists y.\; \operatorname{rm\text{-}realizes}(s, q, \mathit{es}, n, y)) \end{array}}}
\vnbprop{rmm-k-sub}{\vnbstmt{\forall s, q.\; (\operatorname{rm\text{-}k}(s, q) \subseteq \operatorname{pts}(\operatorname{cantor\text{-}ms}))}}
\vnbprop{rmm-realizes-agree}{\vnbstmt{\begin{array}{@{}l@{}} \forall s, q, \mathit{es}, b, a, y. \\ \quad (\forall i \in \mathbb{N}.\; ((i < a) \Rightarrow (b(i) = \operatorname{es}(i))) \wedge \operatorname{rm\text{-}realizes}(s, q, \mathit{es}, a, y)) \Rightarrow  \\ \quad \quad \operatorname{rm\text{-}realizes}(s, q, b, a, y) \end{array}}}
\vnbprop{rmm-k-out}{\vnbstmt{\begin{array}{@{}l@{}} \forall s, q.; \forall \mathit{es} \in \operatorname{pts}(\operatorname{cantor\text{-}ms}). \\ \quad \neg (\mathit{es} \in \operatorname{rm\text{-}k}(s, q)) \Rightarrow  \\ \quad \quad \exists n \in \mathbb{N}.\; \forall y.\; \neg \operatorname{rm\text{-}realizes}(s, q, \mathit{es}, n, y) \end{array}}}
\vnbprop{rmm-k-closed}{\vnbstmt{\forall s, q.\; \operatorname{is\text{-}closed}(\operatorname{cantor\text{-}ms}, \operatorname{rm\text{-}k}(s, q))}}
\vnbprop{rmm-k-set}{\vnbstmt{\forall s, q.\; (\operatorname{rm\text{-}k}(s, q) \in \mathbf{Set})}}
\vnbprop{rmm-k-compact}{\vnbstmt{\forall s, q.\; \operatorname{is\text{-}compact}(\operatorname{rm\text{-}kms}(s, q))}}
\vnbprop{rmm-kms-metric}{\vnbstmt{\forall s, q.\; \operatorname{is\text{-}metric\text{-}space}(\operatorname{rm\text{-}kms}(s, q))}}
\vnbprop{rmm-kms-pts}{\vnbstmt{\forall s, q.\; (\operatorname{pts}(\operatorname{rm\text{-}kms}(s, q)) \simeq \operatorname{rm\text{-}k}(s, q))}}
\vnbprop{rmm-nn-choice}{\vnbstmt{\begin{array}{@{}l@{}} \forall a, m. \\ \quad \forall k \in \mathbb{N}.\; \exists w \in a.\; (w \in m(k)) \Rightarrow  \\ \quad \quad \exists g \in (\mathbb{N} \to a).\; \forall k \in \mathbb{N}.\; (g(k) \in m(k)) \end{array}}}
\vnbprop{rmm-address-unique}{\vnbstmt{\begin{array}{@{}l@{}} \forall s, q, \mathit{es}, y_{1}, y_{2}. \\ \quad (\operatorname{is\text{-}dense\text{-}seq}(s, q) \wedge (\operatorname{rm\text{-}address}(s, q, \mathit{es}, y_{1}) \wedge \operatorname{rm\text{-}address}(s, q, \mathit{es}, y_{2}))) \Rightarrow  \\ \quad \quad y_{1} = y_{2} \end{array}}}
\vnbprop{rmm-shell-closed}{\vnbstmt{\begin{array}{@{}l@{}} \forall s, q.; \forall n \in \mathbb{N}. \\ \quad \operatorname{is\text{-}dense\text{-}seq}(s, q) \Rightarrow \operatorname{is\text{-}closed}(s, \operatorname{rm\text{-}shell}(s, q, n)) \end{array}}}
\vnbprop{rmm-realizes-pts}{\vnbstmt{\forall s, q, \mathit{es}, b, y.\; (\operatorname{rm\text{-}realizes}(s, q, \mathit{es}, b, y) \Rightarrow (y \in \operatorname{pts}(s)))}}
\vnbprop{rmm-ball-closed}{\vnbstmt{\begin{array}{@{}l@{}} \forall s, q.; \forall n \in \mathbb{N}. \\ \quad \operatorname{is\text{-}dense\text{-}seq}(s, q) \Rightarrow \operatorname{is\text{-}closed}(s, \operatorname{rm\text{-}ball}(s, q, n)) \end{array}}}
\vnbprop{rmm-address-exists}{\vnbstmt{\begin{array}{@{}l@{}} \forall s, q.; \forall \mathit{es} \in \operatorname{rm\text{-}k}(s, q). \\ \quad (\operatorname{is\text{-}compact}(s) \wedge \operatorname{is\text{-}dense\text{-}seq}(s, q)) \Rightarrow  \\ \quad \quad \exists y.\; \operatorname{rm\text{-}address}(s, q, \mathit{es}, y) \end{array}}}
\vnbprop{rmm-address-eu}{\vnbstmt{\begin{array}{@{}l@{}} \forall s, q.; \forall \mathit{es} \in \operatorname{rm\text{-}k}(s, q). \\ \quad (\operatorname{is\text{-}compact}(s) \wedge \operatorname{is\text{-}dense\text{-}seq}(s, q)) \Rightarrow  \\ \quad \quad \exists y. \\ \quad \quad \quad \operatorname{rm\text{-}address}(s, q, \mathit{es}, y) \wedge \forall z.\; (\operatorname{rm\text{-}address}(s, q, \mathit{es}, z) \Rightarrow (y = z)) \end{array}}}
\vnbprop{rmm-pi-spec}{\vnbstmt{\begin{array}{@{}l@{}} \forall s, q.; \forall \mathit{es} \in \operatorname{rm\text{-}k}(s, q). \\ \quad (\operatorname{is\text{-}compact}(s) \wedge \operatorname{is\text{-}dense\text{-}seq}(s, q)) \Rightarrow  \\ \quad \quad \operatorname{rm\text{-}address}(s, q, \mathit{es}, \operatorname{rm\text{-}pi}(s, q)(\mathit{es})) \end{array}}}
\vnbprop{rmm-pi-value}{\vnbstmt{\begin{array}{@{}l@{}} \forall s, q.; \forall \mathit{es} \in \operatorname{rm\text{-}k}(s, q), y. \\ \quad (\operatorname{is\text{-}compact}(s) \wedge (\operatorname{is\text{-}dense\text{-}seq}(s, q) \wedge \operatorname{rm\text{-}address}(s, q, \mathit{es}, y))) \Rightarrow  \\ \quad \quad \operatorname{rm\text{-}pi}(s, q)(\mathit{es}) = y \end{array}}}
\vnbprop{rmm-pi-fun}{\vnbstmt{\begin{array}{@{}l@{}} \forall s, q. \\ \quad (\operatorname{is\text{-}compact}(s) \wedge \operatorname{is\text{-}dense\text{-}seq}(s, q)) \Rightarrow  \\ \quad \quad \operatorname{rm\text{-}pi}(s, q) \in (\operatorname{rm\text{-}k}(s, q) \to \operatorname{pts}(s)) \end{array}}}
\vnbprop{rmm-code-value}{\vnbstmt{\forall s, q, x.; \forall n \in \mathbb{N}.\; ((\operatorname{rm\text{-}code}(s, q, x)(n) = 1) \Leftrightarrow (x \in \operatorname{rm\text{-}ball}(s, q, n)))}}
\vnbprop{rmm-code-two}{\vnbstmt{\forall s, q, x.\; (\operatorname{rm\text{-}code}(s, q, x) \in (\mathbb{N} \to \operatorname{rm\text{-}two}))}}
\vnbprop{rmm-cantor-pts-intro}{\vnbstmt{\forall s \in (\mathbb{N} \to \operatorname{rm\text{-}two}).\; (s \in \operatorname{pts}(\operatorname{cantor\text{-}ms}))}}
\vnbprop{rmm-code-in-k}{\vnbstmt{\begin{array}{@{}l@{}} \forall s, q.; \forall x \in \operatorname{pts}(s). \\ \quad \operatorname{is\text{-}dense\text{-}seq}(s, q) \Rightarrow (\operatorname{rm\text{-}code}(s, q, x) \in \operatorname{rm\text{-}k}(s, q)) \end{array}}}
\vnbprop{rmm-code-address}{\vnbstmt{\begin{array}{@{}l@{}} \forall s, q.; \forall x \in \operatorname{pts}(s). \\ \quad \operatorname{is\text{-}dense\text{-}seq}(s, q) \Rightarrow \operatorname{rm\text{-}address}(s, q, \operatorname{rm\text{-}code}(s, q, x), x) \end{array}}}
\vnbprop{rmm-pi-code}{\vnbstmt{\begin{array}{@{}l@{}} \forall s, q.; \forall x \in \operatorname{pts}(s). \\ \quad (\operatorname{is\text{-}compact}(s) \wedge \operatorname{is\text{-}dense\text{-}seq}(s, q)) \Rightarrow (\operatorname{rm\text{-}pi}(s, q)(\operatorname{rm\text{-}code}(s, q, x)) = x) \end{array}}}
\vnbprop{rmm-pi-onto}{\vnbstmt{\begin{array}{@{}l@{}} \forall s, q.; \forall x \in \operatorname{pts}(s). \\ \quad (\operatorname{is\text{-}compact}(s) \wedge \operatorname{is\text{-}dense\text{-}seq}(s, q)) \Rightarrow  \\ \quad \quad \exists \mathit{es} \in \operatorname{rm\text{-}k}(s, q).\; (\operatorname{rm\text{-}pi}(s, q)(\mathit{es}) = x) \end{array}}}
\vnbprop{rmm-address-below}{\vnbstmt{\begin{array}{@{}l@{}} \forall s, q, \mathit{es}, y.; \forall b \in \mathbb{N}. \\ \quad \operatorname{rm\text{-}address}(s, q, \mathit{es}, y) \Rightarrow \operatorname{rm\text{-}address\text{-}below}(s, q, \mathit{es}, b, y) \end{array}}}
\vnbprop{rmm-below-agree}{\vnbstmt{\begin{array}{@{}l@{}} \forall s, q, \mathit{es}, b, a, y. \\ \quad (\forall i \in \mathbb{N}.\; ((i < a) \Rightarrow (b(i) = \operatorname{es}(i))) \wedge \operatorname{rm\text{-}address\text{-}below}(s, q, b, a, y)) \Rightarrow  \\ \quad \quad \operatorname{rm\text{-}address\text{-}below}(s, q, \mathit{es}, a, y) \end{array}}}
\vnbprop{rmm-shrink}{\vnbstmt{\begin{array}{@{}l@{}} \forall s, q.; \forall \mathit{es} \in \operatorname{rm\text{-}k}(s, q), a. \\ \quad (\operatorname{is\text{-}compact}(s) \wedge (\operatorname{is\text{-}dense\text{-}seq}(s, q) \wedge \operatorname{pos\text{-}rr}(a))) \Rightarrow  \\ \quad \quad \exists n \in \mathbb{N}. \forall y. \\ \quad \quad \quad \operatorname{rm\text{-}address\text{-}below}(s, q, \mathit{es}, n, y) \Rightarrow (\operatorname{dist}(s)(\operatorname{rm\text{-}pi}(s, q)(\mathit{es}), y) \leq a) \end{array}}}
\vnbprop{rmm-kms-dist}{\vnbstmt{\begin{array}{@{}l@{}} \forall s, q.; \forall u, v \in \operatorname{rm\text{-}k}(s, q). \\ \quad (\operatorname{dist}(\operatorname{rm\text{-}kms}(s, q))(u, v) \simeq \operatorname{dist}(\operatorname{cantor\text{-}ms})(u, v)) \end{array}}}
\vnbprop{rmm-pi-continuous}{\vnbstmt{\begin{array}{@{}l@{}} \forall s, q. \\ \quad (\operatorname{is\text{-}compact}(s) \wedge \operatorname{is\text{-}dense\text{-}seq}(s, q)) \Rightarrow  \\ \quad \quad \operatorname{is\text{-}continuous}(\operatorname{rm\text{-}kms}(s, q), s, \operatorname{rm\text{-}pi}(s, q)) \end{array}}}
\vnbprop{rmm-cyl-mem}{\vnbstmt{\begin{array}{@{}l@{}} \forall s, q, a, b, \mathit{es}. \\ \quad (\mathit{es} \in \operatorname{rm\text{-}cyl}(s, q, a, b)) \Leftrightarrow ((\mathit{es} \in \operatorname{rm\text{-}k}(s, q)) \wedge \forall i \in \mathbb{N}.\; ((i < b) \Rightarrow (\operatorname{es}(i) = a(i)))) \end{array}}}
\vnbprop{rmm-cyl-sub}{\vnbstmt{\forall s, q, a, b.\; (\operatorname{rm\text{-}cyl}(s, q, a, b) \subseteq \operatorname{rm\text{-}k}(s, q))}}
\vnbprop{rmm-cyl-open}{\vnbstmt{\forall s, q, a.; \forall b \in \mathbb{N}.\; \operatorname{is\text{-}open}(\operatorname{rm\text{-}kms}(s, q), \operatorname{rm\text{-}cyl}(s, q, a, b))}}
\vnbprop{rmm-cyl-closed}{\vnbstmt{\forall s, q, a.; \forall b \in \mathbb{N}.\; \operatorname{is\text{-}closed}(\operatorname{rm\text{-}kms}(s, q), \operatorname{rm\text{-}cyl}(s, q, a, b))}}
\vnbprop{rmm-kms-unfold}{\vnbstmt{\forall s, q.\; (\operatorname{rm\text{-}kms}(s, q) \simeq \operatorname{subspace\text{-}ms}(\operatorname{cantor\text{-}ms}, \operatorname{rm\text{-}k}(s, q)))}}
\vnbprop{rmm-cyl-base}{\vnbstmt{\begin{array}{@{}l@{}} \forall s, q, u.; \forall \mathit{es} \in u. \\ \quad \operatorname{is\text{-}open}(\operatorname{rm\text{-}kms}(s, q), u) \Rightarrow  \\ \quad \quad \exists k \in \mathbb{N}.\; (\operatorname{rm\text{-}cyl}(s, q, \mathit{es}, k) \subseteq u) \end{array}}}
\vnbprop{rmm-cyl-self}{\vnbstmt{\forall s, q.; \forall \mathit{es} \in \operatorname{rm\text{-}k}(s, q), b.\; (\mathit{es} \in \operatorname{rm\text{-}cyl}(s, q, \mathit{es}, b))}}
\vnbprop{rmm-dense-seq-exists}{\vnbstmt{\begin{array}{@{}l@{}} \forall s. \\ \quad (\operatorname{is\text{-}compact}(s) \wedge \exists x.\; (x \in \operatorname{pts}(s))) \Rightarrow  \\ \quad \quad \exists q.\; \operatorname{is\text{-}dense\text{-}seq}(s, q) \end{array}}}

\vnbbookfile{theorem-library/categories}{theorem-library/categories}
\noindent{\small categories.scm -- FOUR CATEGORIES BESIDE THE DEFAULTS: Lipschitz, uniformly continuous and continuous maps of metric spaces, bounded linear maps of normed vector spaces; their laws, their Hom functors, and the inclusions.}\par
\par\medskip\noindent\textbf{Definition} (\texttt{hom-lipschitz}).\;$\operatorname{hom\text{-}lipschitz}(a, b) :=$
\begin{equation*}
\{\operatorname{homf\_} \in (\operatorname{pts}(a) \to \operatorname{pts}(b)) : \operatorname{is\text{-}lipschitz\text{-}arrow}(a, b, \operatorname{homf\_})\}
\end{equation*}

\par\medskip\noindent\textbf{Definition} (\texttt{hom-uniformly-continuous}).\;$\operatorname{hom\text{-}uniformly\text{-}continuous}(a, b) :=$
\begin{equation*}
\{\operatorname{homf\_} \in (\operatorname{pts}(a) \to \operatorname{pts}(b)) : \operatorname{is\text{-}uniformly\text{-}continuous\text{-}arrow}(a, b, \operatorname{homf\_})\}
\end{equation*}

\par\medskip\noindent\textbf{Definition} (\texttt{hom-continuous}).\;$\operatorname{hom\text{-}continuous}(a, b) :=$
\begin{equation*}
\{\operatorname{homf\_} \in (\operatorname{pts}(a) \to \operatorname{pts}(b)) : \operatorname{is\text{-}continuous\text{-}arrow}(a, b, \operatorname{homf\_})\}
\end{equation*}

\par\medskip\noindent\textbf{Definition} (\texttt{hom-bounded-linear}).\;$\operatorname{hom\text{-}bounded\text{-}linear}(a, b) :=$
\begin{equation*}
\{\operatorname{homf\_} \in (\operatorname{vec}(a) \to \operatorname{vec}(b)) : \operatorname{is\text{-}bounded\text{-}linear\text{-}arrow}(a, b, \operatorname{homf\_})\}
\end{equation*}

\vnbdef{is-lipschitz-arrow-def}{\vnbstmt{\begin{array}{@{}l@{}} \forall a, b, f. \\ \quad \operatorname{is\text{-}lipschitz\text{-}arrow}(a, b, f) \\ \quad \Leftrightarrow\; \operatorname{is\text{-}metric\text{-}space}(a) \wedge \\ \quad \operatorname{is\text{-}metric\text{-}space}(b) \wedge \\ \quad f \in (\operatorname{pts}(a) \to \operatorname{pts}(b)) \wedge \\ \quad \exists k \in \mathbb{R}. \\ \quad \quad ((0 \leq k) \wedge \forall x, y \in \operatorname{pts}(a).\; (\operatorname{dist}(b)(f(x), f(y)) \leq (k \cdot \operatorname{dist}(a)(x, y)))) \end{array}}}
\vnbdef{is-uniformly-continuous-arrow-def}{\vnbstmt{\begin{array}{@{}l@{}} \forall a, b, f. \\ \quad \operatorname{is\text{-}uniformly\text{-}continuous\text{-}arrow}(a, b, f) \\ \quad \Leftrightarrow\; \operatorname{is\text{-}metric\text{-}space}(a) \wedge \\ \quad \operatorname{is\text{-}metric\text{-}space}(b) \wedge \\ \quad f \in (\operatorname{pts}(a) \to \operatorname{pts}(b)) \wedge \\ \quad \operatorname{is\text{-}uniformly\text{-}continuous}(a, b, f) \end{array}}}
\vnbdef{is-continuous-arrow-def}{\vnbstmt{\begin{array}{@{}l@{}} \forall a, b, f. \\ \quad \operatorname{is\text{-}continuous\text{-}arrow}(a, b, f) \\ \quad \Leftrightarrow\; (\operatorname{is\text{-}metric\text{-}space}(a) \wedge (\operatorname{is\text{-}metric\text{-}space}(b) \wedge ((f \in (\operatorname{pts}(a) \to \operatorname{pts}(b))) \wedge \operatorname{is\text{-}continuous}(a, b, f)))) \end{array}}}
\vnbdef{is-bounded-linear-arrow-def}{\vnbstmt{\begin{array}{@{}l@{}} \forall a, b, f. \\ \quad \operatorname{is\text{-}bounded\text{-}linear\text{-}arrow}(a, b, f) \\ \quad \Leftrightarrow\; \operatorname{is\text{-}normed\text{-}vector\text{-}space}(a) \wedge \\ \quad \operatorname{is\text{-}normed\text{-}vector\text{-}space}(b) \wedge \\ \quad f \in (\operatorname{vec}(a) \to \operatorname{vec}(b)) \wedge \\ \quad \operatorname{scal}(a) = \operatorname{scal}(b) \wedge \\ \quad \forall x_{1}, x_{2} \in \operatorname{vec}(a). \\ \quad \quad f(\operatorname{vadd}(a)(x_{1}, x_{2})) = \operatorname{vadd}(b)(f(x_{1}), f(x_{2})) \wedge \\ \quad \forall x_{1} \in \operatorname{carr}(\operatorname{scal}(a)), x_{2} \in \operatorname{vec}(a). \\ \quad \quad f(\operatorname{act}(a)(x_{1}, x_{2})) = \operatorname{act}(b)(x_{1}, f(x_{2})) \wedge \\ \quad \exists k \in \mathbb{R}. \\ \quad \quad (0 \leq k) \wedge \forall x_{1} \in \operatorname{vec}(a).\; (\operatorname{vnrm}(b)(f(x_{1})) \leq (k \cdot \operatorname{vnrm}(a)(x_{1}))) \end{array}}}
\vnbprop{hom-lipschitz-member-iff}{\vnbstmt{\begin{array}{@{}l@{}} \forall a, b, f. \\ \quad (f \in \operatorname{hom\text{-}lipschitz}(a, b)) \Leftrightarrow ((f \in (\operatorname{pts}(a) \to \operatorname{pts}(b))) \wedge \operatorname{is\text{-}lipschitz\text{-}arrow}(a, b, f)) \end{array}}}
\vnbprop{hom-lipschitz-in-set}{\vnbstmt{\begin{array}{@{}l@{}} \forall a, b. \\ \quad (\operatorname{is\text{-}metric\text{-}space}(a) \wedge \operatorname{is\text{-}metric\text{-}space}(b)) \Rightarrow (\operatorname{hom\text{-}lipschitz}(a, b) \in \mathbf{Set}) \end{array}}}
\vnbprop{hom-uniformly-continuous-member-iff}{\vnbstmt{\begin{array}{@{}l@{}} \forall a, b, f. \\ \quad f \in \operatorname{hom\text{-}uniformly\text{-}continuous}(a, b) \\ \quad \Leftrightarrow\; (f \in (\operatorname{pts}(a) \to \operatorname{pts}(b))) \wedge \operatorname{is\text{-}uniformly\text{-}continuous\text{-}arrow}(a, b, f) \end{array}}}
\vnbprop{hom-uniformly-continuous-in-set}{\vnbstmt{\begin{array}{@{}l@{}} \forall a, b. \\ \quad (\operatorname{is\text{-}metric\text{-}space}(a) \wedge \operatorname{is\text{-}metric\text{-}space}(b)) \Rightarrow  \\ \quad \quad \operatorname{hom\text{-}uniformly\text{-}continuous}(a, b) \in \mathbf{Set} \end{array}}}
\vnbprop{hom-continuous-member-iff}{\vnbstmt{\begin{array}{@{}l@{}} \forall a, b, f. \\ \quad (f \in \operatorname{hom\text{-}continuous}(a, b)) \Leftrightarrow ((f \in (\operatorname{pts}(a) \to \operatorname{pts}(b))) \wedge \operatorname{is\text{-}continuous\text{-}arrow}(a, b, f)) \end{array}}}
\vnbprop{hom-continuous-in-set}{\vnbstmt{\begin{array}{@{}l@{}} \forall a, b. \\ \quad (\operatorname{is\text{-}metric\text{-}space}(a) \wedge \operatorname{is\text{-}metric\text{-}space}(b)) \Rightarrow (\operatorname{hom\text{-}continuous}(a, b) \in \mathbf{Set}) \end{array}}}
\vnbprop{hom-bounded-linear-member-iff}{\vnbstmt{\begin{array}{@{}l@{}} \forall a, b, f. \\ \quad (f \in \operatorname{hom\text{-}bounded\text{-}linear}(a, b)) \Leftrightarrow ((f \in (\operatorname{vec}(a) \to \operatorname{vec}(b))) \wedge \operatorname{is\text{-}bounded\text{-}linear\text{-}arrow}(a, b, f)) \end{array}}}
\vnbprop{hom-bounded-linear-in-set}{\vnbstmt{\begin{array}{@{}l@{}} \forall a, b. \\ \quad (\operatorname{is\text{-}normed\text{-}vector\text{-}space}(a) \wedge \operatorname{is\text{-}normed\text{-}vector\text{-}space}(b)) \Rightarrow  \\ \quad \quad \operatorname{hom\text{-}bounded\text{-}linear}(a, b) \in \mathbf{Set} \end{array}}}
\vnbprop{hom-metric-space-is-lipschitz-arrow}{\vnbstmt{\begin{array}{@{}l@{}} \forall a, b, f. \\ \quad \operatorname{is\text{-}hom\text{-}metric\text{-}space}(a, b, f) \Rightarrow \operatorname{is\text{-}lipschitz\text{-}arrow}(a, b, f) \end{array}}}
\vnbprop{lipschitz-arrow-is-uniformly-continuous-arrow}{\vnbstmt{\begin{array}{@{}l@{}} \forall a, b, f. \\ \quad \operatorname{is\text{-}lipschitz\text{-}arrow}(a, b, f) \Rightarrow \operatorname{is\text{-}uniformly\text{-}continuous\text{-}arrow}(a, b, f) \end{array}}}
\vnbprop{uniformly-continuous-arrow-is-continuous-arrow}{\vnbstmt{\begin{array}{@{}l@{}} \forall a, b, f. \\ \quad \operatorname{is\text{-}uniformly\text{-}continuous\text{-}arrow}(a, b, f) \Rightarrow \operatorname{is\text{-}continuous\text{-}arrow}(a, b, f) \end{array}}}
\vnbprop{hom-normed-vector-space-is-bounded-linear-arrow}{\vnbstmt{\begin{array}{@{}l@{}} \forall a, b, f. \\ \quad \operatorname{is\text{-}hom\text{-}normed\text{-}vector\text{-}space}(a, b, f) \Rightarrow \operatorname{is\text{-}bounded\text{-}linear\text{-}arrow}(a, b, f) \end{array}}}
\vnbprop{hom-lipschitz-id}{\vnbstmt{\forall a.\; (\operatorname{is\text{-}metric\text{-}space}(a) \Rightarrow \operatorname{is\text{-}lipschitz\text{-}arrow}(a, a, \operatorname{id\text{-}fun}(\operatorname{pts}(a))))}}
\vnbprop{hom-uniformly-continuous-id}{\vnbstmt{\begin{array}{@{}l@{}} \forall a. \\ \quad \operatorname{is\text{-}metric\text{-}space}(a) \Rightarrow \operatorname{is\text{-}uniformly\text{-}continuous\text{-}arrow}(a, a, \operatorname{id\text{-}fun}(\operatorname{pts}(a))) \end{array}}}
\vnbprop{hom-continuous-id}{\vnbstmt{\begin{array}{@{}l@{}} \forall a. \\ \quad \operatorname{is\text{-}metric\text{-}space}(a) \Rightarrow \operatorname{is\text{-}continuous\text{-}arrow}(a, a, \operatorname{id\text{-}fun}(\operatorname{pts}(a))) \end{array}}}
\vnbprop{hom-bounded-linear-id}{\vnbstmt{\begin{array}{@{}l@{}} \forall a. \\ \quad \operatorname{is\text{-}normed\text{-}vector\text{-}space}(a) \Rightarrow \operatorname{is\text{-}bounded\text{-}linear\text{-}arrow}(a, a, \operatorname{id\text{-}fun}(\operatorname{vec}(a))) \end{array}}}
\vnbprop{hom-lipschitz-compose}{\vnbstmt{\begin{array}{@{}l@{}} \forall a, b, c, f, g. \\ \quad (\operatorname{is\text{-}lipschitz\text{-}arrow}(a, b, f) \wedge \operatorname{is\text{-}lipschitz\text{-}arrow}(b, c, g)) \Rightarrow  \\ \quad \quad \operatorname{is\text{-}lipschitz\text{-}arrow}(a, c, \operatorname{compose}(g, f)) \end{array}}}
\vnbprop{hom-uniformly-continuous-compose}{\vnbstmt{\begin{array}{@{}l@{}} \forall a, b, c, f, g. \\ \quad (\operatorname{is\text{-}uniformly\text{-}continuous\text{-}arrow}(a, b, f) \wedge \operatorname{is\text{-}uniformly\text{-}continuous\text{-}arrow}(b, c, g)) \Rightarrow  \\ \quad \quad \operatorname{is\text{-}uniformly\text{-}continuous\text{-}arrow}(a, c, \operatorname{compose}(g, f)) \end{array}}}
\vnbprop{hom-continuous-compose}{\vnbstmt{\begin{array}{@{}l@{}} \forall a, b, c, f, g. \\ \quad (\operatorname{is\text{-}continuous\text{-}arrow}(a, b, f) \wedge \operatorname{is\text{-}continuous\text{-}arrow}(b, c, g)) \Rightarrow  \\ \quad \quad \operatorname{is\text{-}continuous\text{-}arrow}(a, c, \operatorname{compose}(g, f)) \end{array}}}
\vnbprop{hom-bounded-linear-compose}{\vnbstmt{\begin{array}{@{}l@{}} \forall a, b, c, f, g. \\ \quad (\operatorname{is\text{-}bounded\text{-}linear\text{-}arrow}(a, b, f) \wedge \operatorname{is\text{-}bounded\text{-}linear\text{-}arrow}(b, c, g)) \Rightarrow  \\ \quad \quad \operatorname{is\text{-}bounded\text{-}linear\text{-}arrow}(a, c, \operatorname{compose}(g, f)) \end{array}}}
\vnbprop{hom-lipschitz-post-type}{\vnbstmt{\begin{array}{@{}l@{}} \forall a, b, c, h. \\ \quad (\operatorname{is\text{-}metric\text{-}space}(a) \wedge (h \in \operatorname{hom\text{-}lipschitz}(b, c))) \Rightarrow  \\ \quad \quad (\lambda f \in \operatorname{hom\text{-}lipschitz}(a, b).\; \operatorname{compose}(h, f)) \in (\operatorname{hom\text{-}lipschitz}(a, b) \to \operatorname{hom\text{-}lipschitz}(a, c)) \end{array}}}
\vnbprop{hom-lipschitz-pre-type}{\vnbstmt{\begin{array}{@{}l@{}} \forall a, b, c, k. \\ \quad (\operatorname{is\text{-}metric\text{-}space}(b) \wedge (k \in \operatorname{hom\text{-}lipschitz}(c, a))) \Rightarrow  \\ \quad \quad (\lambda f \in \operatorname{hom\text{-}lipschitz}(a, b).\; \operatorname{compose}(f, k)) \in (\operatorname{hom\text{-}lipschitz}(a, b) \to \operatorname{hom\text{-}lipschitz}(c, b)) \end{array}}}
\vnbprop{hom-uniformly-continuous-post-type}{\vnbstmt{\begin{array}{@{}l@{}} \forall a, b, c, h. \\ \quad (\operatorname{is\text{-}metric\text{-}space}(a) \wedge (h \in \operatorname{hom\text{-}uniformly\text{-}continuous}(b, c))) \Rightarrow  \\ \quad \quad (\lambda f \in \operatorname{hom\text{-}uniformly\text{-}continuous}(a, b).\; \operatorname{compose}(h, f)) \\ \quad \quad \in\; (\operatorname{hom\text{-}uniformly\text{-}continuous}(a, b) \to \operatorname{hom\text{-}uniformly\text{-}continuous}(a, c)) \end{array}}}
\vnbprop{hom-uniformly-continuous-pre-type}{\vnbstmt{\begin{array}{@{}l@{}} \forall a, b, c, k. \\ \quad (\operatorname{is\text{-}metric\text{-}space}(b) \wedge (k \in \operatorname{hom\text{-}uniformly\text{-}continuous}(c, a))) \Rightarrow  \\ \quad \quad (\lambda f \in \operatorname{hom\text{-}uniformly\text{-}continuous}(a, b).\; \operatorname{compose}(f, k)) \\ \quad \quad \in\; (\operatorname{hom\text{-}uniformly\text{-}continuous}(a, b) \to \operatorname{hom\text{-}uniformly\text{-}continuous}(c, b)) \end{array}}}
\vnbprop{hom-continuous-post-type}{\vnbstmt{\begin{array}{@{}l@{}} \forall a, b, c, h. \\ \quad (\operatorname{is\text{-}metric\text{-}space}(a) \wedge (h \in \operatorname{hom\text{-}continuous}(b, c))) \Rightarrow  \\ \quad \quad (\lambda f \in \operatorname{hom\text{-}continuous}(a, b).\; \operatorname{compose}(h, f)) \\ \quad \quad \in\; (\operatorname{hom\text{-}continuous}(a, b) \to \operatorname{hom\text{-}continuous}(a, c)) \end{array}}}
\vnbprop{hom-continuous-pre-type}{\vnbstmt{\begin{array}{@{}l@{}} \forall a, b, c, k. \\ \quad (\operatorname{is\text{-}metric\text{-}space}(b) \wedge (k \in \operatorname{hom\text{-}continuous}(c, a))) \Rightarrow  \\ \quad \quad (\lambda f \in \operatorname{hom\text{-}continuous}(a, b).\; \operatorname{compose}(f, k)) \\ \quad \quad \in\; (\operatorname{hom\text{-}continuous}(a, b) \to \operatorname{hom\text{-}continuous}(c, b)) \end{array}}}
\vnbprop{hom-bounded-linear-post-type}{\vnbstmt{\begin{array}{@{}l@{}} \forall a, b, c, h. \\ \quad (\operatorname{is\text{-}normed\text{-}vector\text{-}space}(a) \wedge (h \in \operatorname{hom\text{-}bounded\text{-}linear}(b, c))) \Rightarrow  \\ \quad \quad (\lambda f \in \operatorname{hom\text{-}bounded\text{-}linear}(a, b).\; \operatorname{compose}(h, f)) \\ \quad \quad \in\; (\operatorname{hom\text{-}bounded\text{-}linear}(a, b) \to \operatorname{hom\text{-}bounded\text{-}linear}(a, c)) \end{array}}}
\vnbprop{hom-bounded-linear-pre-type}{\vnbstmt{\begin{array}{@{}l@{}} \forall a, b, c, k. \\ \quad (\operatorname{is\text{-}normed\text{-}vector\text{-}space}(b) \wedge (k \in \operatorname{hom\text{-}bounded\text{-}linear}(c, a))) \Rightarrow  \\ \quad \quad (\lambda f \in \operatorname{hom\text{-}bounded\text{-}linear}(a, b).\; \operatorname{compose}(f, k)) \\ \quad \quad \in\; (\operatorname{hom\text{-}bounded\text{-}linear}(a, b) \to \operatorname{hom\text{-}bounded\text{-}linear}(c, b)) \end{array}}}
\vnbprop{hom-metric-space-subset-hom-lipschitz}{\vnbstmt{\begin{array}{@{}l@{}} \forall a, b. \\ \quad (\operatorname{is\text{-}metric\text{-}space}(a) \wedge \operatorname{is\text{-}metric\text{-}space}(b)) \Rightarrow  \\ \quad \quad \operatorname{hom\text{-}metric\text{-}space}(a, b) \subseteq \operatorname{hom\text{-}lipschitz}(a, b) \end{array}}}
\vnbprop{hom-lipschitz-subset-hom-uniformly-continuous}{\vnbstmt{\begin{array}{@{}l@{}} \forall a, b. \\ \quad (\operatorname{is\text{-}metric\text{-}space}(a) \wedge \operatorname{is\text{-}metric\text{-}space}(b)) \Rightarrow  \\ \quad \quad \operatorname{hom\text{-}lipschitz}(a, b) \subseteq \operatorname{hom\text{-}uniformly\text{-}continuous}(a, b) \end{array}}}
\vnbprop{hom-uniformly-continuous-subset-hom-continuous}{\vnbstmt{\begin{array}{@{}l@{}} \forall a, b. \\ \quad (\operatorname{is\text{-}metric\text{-}space}(a) \wedge \operatorname{is\text{-}metric\text{-}space}(b)) \Rightarrow  \\ \quad \quad \operatorname{hom\text{-}uniformly\text{-}continuous}(a, b) \subseteq \operatorname{hom\text{-}continuous}(a, b) \end{array}}}
\vnbprop{hom-normed-vector-space-subset-hom-bounded-linear}{\vnbstmt{\begin{array}{@{}l@{}} \forall a, b. \\ \quad (\operatorname{is\text{-}normed\text{-}vector\text{-}space}(a) \wedge \operatorname{is\text{-}normed\text{-}vector\text{-}space}(b)) \Rightarrow  \\ \quad \quad \operatorname{hom\text{-}normed\text{-}vector\text{-}space}(a, b) \subseteq \operatorname{hom\text{-}bounded\text{-}linear}(a, b) \end{array}}}

\vnbbookfile{theorem-library/c-continuous-category}{theorem-library/c-continuous-category}
\noindent{\small c-continuous-category.scm -- THE CATEGORY OF COUNTABLY-METRISED SPACES AND CONTINUOUS MAPS: objects the C-METRIC-SPACEs, arrows the maps f : PTS(a) -$>$ PTS(b) continuous for the canonical metrics C-METRIC(a), C-METRIC(b).}\par
\par\medskip\noindent\textbf{Definition} (\texttt{hom-c-continuous}).\;$\operatorname{hom\text{-}c\text{-}continuous}(a, b) :=$
\begin{equation*}
\{\operatorname{homf\_} \in (\operatorname{pts}(a) \to \operatorname{pts}(b)) : \operatorname{is\text{-}c\text{-}continuous\text{-}arrow}(a, b, \operatorname{homf\_})\}
\end{equation*}

\vnbdef{is-c-continuous-arrow-def}{\vnbstmt{\begin{array}{@{}l@{}} \forall a, b, f. \\ \quad \operatorname{is\text{-}c\text{-}continuous\text{-}arrow}(a, b, f) \\ \quad \Leftrightarrow\; \operatorname{is\text{-}c\text{-}metric\text{-}space}(a) \wedge \\ \quad \operatorname{is\text{-}c\text{-}metric\text{-}space}(b) \wedge \\ \quad f \in (\operatorname{pts}(a) \to \operatorname{pts}(b)) \wedge \\ \quad \operatorname{is\text{-}continuous}(\operatorname{c\text{-}metric}(a), \operatorname{c\text{-}metric}(b), f) \end{array}}}
\vnbprop{c-metric-canonical-carrier}{\vnbstmt{\forall s.\; (\operatorname{is\text{-}c\text{-}metric\text{-}space}(s) \Rightarrow (\operatorname{pts}(\operatorname{c\text{-}metric}(s)) \simeq \operatorname{pts}(s)))}}
\vnbprop{hom-c-continuous-member-iff}{\vnbstmt{\begin{array}{@{}l@{}} \forall a, b, f. \\ \quad (f \in \operatorname{hom\text{-}c\text{-}continuous}(a, b)) \Leftrightarrow ((f \in (\operatorname{pts}(a) \to \operatorname{pts}(b))) \wedge \operatorname{is\text{-}c\text{-}continuous\text{-}arrow}(a, b, f)) \end{array}}}
\vnbprop{hom-c-continuous-in-set}{\vnbstmt{\begin{array}{@{}l@{}} \forall a, b. \\ \quad (\operatorname{is\text{-}c\text{-}metric\text{-}space}(a) \wedge \operatorname{is\text{-}c\text{-}metric\text{-}space}(b)) \Rightarrow  \\ \quad \quad \operatorname{hom\text{-}c\text{-}continuous}(a, b) \in \mathbf{Set} \end{array}}}
\vnbprop{hom-c-continuous-compose}{\vnbstmt{\begin{array}{@{}l@{}} \forall a, b, c, f, g. \\ \quad (\operatorname{is\text{-}c\text{-}continuous\text{-}arrow}(a, b, f) \wedge \operatorname{is\text{-}c\text{-}continuous\text{-}arrow}(b, c, g)) \Rightarrow  \\ \quad \quad \operatorname{is\text{-}c\text{-}continuous\text{-}arrow}(a, c, \operatorname{compose}(g, f)) \end{array}}}
\vnbprop{hom-c-continuous-id-of-metric}{\vnbstmt{\begin{array}{@{}l@{}} \forall a. \\ \quad (\operatorname{is\text{-}c\text{-}metric\text{-}space}(a) \wedge \operatorname{is\text{-}metric\text{-}space}(\operatorname{c\text{-}metric}(a))) \Rightarrow  \\ \quad \quad \operatorname{is\text{-}c\text{-}continuous\text{-}arrow}(a, a, \operatorname{id\text{-}fun}(\operatorname{pts}(a))) \end{array}}}
\vnbprop{hom-c-continuous-post-type}{\vnbstmt{\begin{array}{@{}l@{}} \forall a, b, c, h. \\ \quad (\operatorname{is\text{-}c\text{-}metric\text{-}space}(a) \wedge (h \in \operatorname{hom\text{-}c\text{-}continuous}(b, c))) \Rightarrow  \\ \quad \quad (\lambda f \in \operatorname{hom\text{-}c\text{-}continuous}(a, b).\; \operatorname{compose}(h, f)) \\ \quad \quad \in\; (\operatorname{hom\text{-}c\text{-}continuous}(a, b) \to \operatorname{hom\text{-}c\text{-}continuous}(a, c)) \end{array}}}
\vnbprop{hom-c-continuous-pre-type}{\vnbstmt{\begin{array}{@{}l@{}} \forall a, b, c, k. \\ \quad (\operatorname{is\text{-}c\text{-}metric\text{-}space}(b) \wedge (k \in \operatorname{hom\text{-}c\text{-}continuous}(c, a))) \Rightarrow  \\ \quad \quad (\lambda f \in \operatorname{hom\text{-}c\text{-}continuous}(a, b).\; \operatorname{compose}(f, k)) \\ \quad \quad \in\; (\operatorname{hom\text{-}c\text{-}continuous}(a, b) \to \operatorname{hom\text{-}c\text{-}continuous}(c, b)) \end{array}}}
\vnbprop{hom-c-continuous-id}{\vnbstmt{\begin{array}{@{}l@{}} \forall a. \\ \quad \operatorname{is\text{-}c\text{-}metric\text{-}space}(a) \Rightarrow \operatorname{is\text{-}c\text{-}continuous\text{-}arrow}(a, a, \operatorname{id\text{-}fun}(\operatorname{pts}(a))) \end{array}}}

\vnbbookfile{theorem-library/diff-on-laws-2}{theorem-library/diff-on-laws-2}
\noindent{\small diff-on-laws-2.scm -- THE DERIVATIVE LAWS ON AN OPEN SET, over a normed field: uniqueness, the constant and the identity, the sum, the product and the chain rule.  Continues theorem-library/diff-on-laws.scm (batch 12-G), which proved the read-offs, the bridge `diff-at-iff-diff-on' and}\par
\vnbprop{ms-const-continuous-on}{\vnbstmt{\begin{array}{@{}l@{}} \forall s, d, t.; \forall c \in \operatorname{pts}(t), a \in \operatorname{pts}(s). \\ \quad ((\operatorname{pts}(s) \simeq d) \wedge (\operatorname{is\text{-}metric\text{-}space}(s) \wedge \operatorname{is\text{-}metric\text{-}space}(t))) \Rightarrow  \\ \quad \quad \operatorname{is\text{-}continuous\text{-}at}(s, t, (\lambda z \in d.\; c), a) \end{array}}}
\vnbprop{ms-identity-continuous-on}{\vnbstmt{\begin{array}{@{}l@{}} \forall s, d.; \forall a \in \operatorname{pts}(s). \\ \quad ((\operatorname{pts}(s) \simeq d) \wedge \operatorname{is\text{-}metric\text{-}space}(s)) \Rightarrow  \\ \quad \quad \operatorname{is\text{-}continuous\text{-}at}(s, s, (\lambda z \in d.\; z), a) \end{array}}}
\vnbprop{nf-sum-continuous-on}{\vnbstmt{\begin{array}{@{}l@{}} \forall k, s, d, g, h, a. \\ \quad \operatorname{is\text{-}normed\text{-}field}(k) \wedge \\ \quad (\operatorname{pts}(s) \simeq d) \wedge \\ \quad \operatorname{is\text{-}continuous\text{-}at}(s, \operatorname{nf\text{-}metric\text{-}space}(k), g, a) \wedge \\ \quad \operatorname{is\text{-}continuous\text{-}at}(s, \operatorname{nf\text{-}metric\text{-}space}(k), h, a) \Rightarrow \\ \quad \quad \operatorname{is\text{-}continuous\text{-}at}(s, \operatorname{nf\text{-}metric\text{-}space}(k), (\lambda z \in d.\; \operatorname{add}(k)(g(z), h(z))), a) \end{array}}}
\vnbprop{nf-product-continuous-on}{\vnbstmt{\begin{array}{@{}l@{}} \forall k, s, d, g, h, a. \\ \quad \operatorname{is\text{-}normed\text{-}field}(k) \wedge \\ \quad (\operatorname{pts}(s) \simeq d) \wedge \\ \quad \operatorname{is\text{-}continuous\text{-}at}(s, \operatorname{nf\text{-}metric\text{-}space}(k), g, a) \wedge \\ \quad \operatorname{is\text{-}continuous\text{-}at}(s, \operatorname{nf\text{-}metric\text{-}space}(k), h, a) \Rightarrow \\ \quad \quad \operatorname{is\text{-}continuous\text{-}at}(s, \operatorname{nf\text{-}metric\text{-}space}(k), (\lambda z \in d.\; \operatorname{mul}(k)(g(z), h(z))), a) \end{array}}}
\vnbprop{ms-compose-continuous-on}{\vnbstmt{\begin{array}{@{}l@{}} \forall s, d, t, u, f, g, a. \\ \quad ((\operatorname{pts}(s) \simeq d) \wedge (\operatorname{is\text{-}continuous\text{-}at}(s, t, f, a) \wedge \operatorname{is\text{-}continuous\text{-}at}(t, u, g, f(a)))) \Rightarrow  \\ \quad \quad \operatorname{is\text{-}continuous\text{-}at}(s, u, (\lambda z \in d.\; g(f(z))), a) \end{array}}}
\vnbprop{ms-cont-transfer-on}{\vnbstmt{\begin{array}{@{}l@{}} \forall s, t, d, c, g, f, a. \\ \quad (\operatorname{pts}(s) \simeq d) \wedge \\ \quad (\operatorname{pts}(t) \simeq c) \wedge \\ \quad \operatorname{is\text{-}continuous\text{-}at}(s, t, g, a) \wedge \\ \quad f \in (d \to c) \wedge \\ \quad \forall y \in d.\; (f(y) = g(y)) \Rightarrow \\ \quad \quad \operatorname{is\text{-}continuous\text{-}at}(s, t, f, a) \end{array}}}
\vnbprop{nf-lam-const-in-fun}{\vnbstmt{\begin{array}{@{}l@{}} \forall k.; \forall u \in \mathbf{Set}, c \in \operatorname{carr}(k). \\ \quad \operatorname{is\text{-}normed\text{-}field}(k) \Rightarrow ((\lambda x \in u.\; c) \in (u \to \operatorname{carr}(k))) \end{array}}}
\vnbprop{nf-lam-id-in-fun}{\vnbstmt{\begin{array}{@{}l@{}} \forall k.; \forall u \in \mathbf{Set}. \\ \quad (\operatorname{is\text{-}normed\text{-}field}(k) \wedge (u \subseteq \operatorname{carr}(k))) \Rightarrow  \\ \quad \quad (\lambda x \in u.\; x) \in (u \to \operatorname{carr}(k)) \end{array}}}
\vnbprop{nf-lam-sum-in-fun}{\vnbstmt{\begin{array}{@{}l@{}} \forall k.; \forall u \in \mathbf{Set}, f, g \in (u \to \operatorname{carr}(k)). \\ \quad \operatorname{is\text{-}normed\text{-}field}(k) \Rightarrow  \\ \quad \quad (\lambda x \in u.\; \operatorname{add}(k)(f(x), g(x))) \in (u \to \operatorname{carr}(k)) \end{array}}}
\vnbprop{nf-lam-prod-in-fun}{\vnbstmt{\begin{array}{@{}l@{}} \forall k.; \forall u \in \mathbf{Set}, f, g \in (u \to \operatorname{carr}(k)). \\ \quad \operatorname{is\text{-}normed\text{-}field}(k) \Rightarrow  \\ \quad \quad (\lambda x \in u.\; \operatorname{mul}(k)(f(x), g(x))) \in (u \to \operatorname{carr}(k)) \end{array}}}
\vnbprop{nf-norm-zero-iff}{\vnbstmt{\begin{array}{@{}l@{}} \forall k.; \forall u \in \operatorname{carr}(k). \\ \quad \operatorname{is\text{-}normed\text{-}field}(k) \Rightarrow ((\operatorname{fnrm}(k)(u) = 0) \Leftrightarrow (u = \operatorname{zero}(k))) \end{array}}}
\vnbprop{nf-norm-zero-implies}{\vnbstmt{\begin{array}{@{}l@{}} \forall k.; \forall u \in \operatorname{carr}(k). \\ \quad (\operatorname{is\text{-}normed\text{-}field}(k) \wedge (\operatorname{fnrm}(k)(u) = 0)) \Rightarrow (u = \operatorname{zero}(k)) \end{array}}}
\vnbprop{rr-nf-small-elements}{\vnbstmt{\begin{array}{@{}l@{}} \forall e. \\ \quad \operatorname{pos\text{-}rr}(e) \Rightarrow  \\ \quad \quad \exists h \in \operatorname{carr}(\operatorname{rr\text{-}normed\text{-}field}). \\ \quad \quad \quad \neg (h = \operatorname{zero}(\operatorname{rr\text{-}normed\text{-}field})) \wedge (\operatorname{fnrm}(\operatorname{rr\text{-}normed\text{-}field})(h) < e) \end{array}}}
\vnbprop{nf-open-limit-point}{\vnbstmt{\begin{array}{@{}l@{}} \forall k, u.; \forall a \in u, r. \\ \quad \operatorname{is\text{-}normed\text{-}field}(k) \wedge \\ \quad \forall e. \\ \quad \quad \operatorname{pos\text{-}rr}(e) \Rightarrow \exists h \in \operatorname{carr}(k).\; (\neg (h = \operatorname{zero}(k)) \wedge (\operatorname{fnrm}(k)(h) < e)) \wedge \\ \quad \operatorname{is\text{-}open}(\operatorname{nf\text{-}metric\text{-}space}(k), u) \wedge \\ \quad \operatorname{pos\text{-}rr}(r) \Rightarrow \\ \quad \quad \exists b \in u. \\ \quad \quad \quad \neg (b = a) \wedge (\operatorname{dist}(\operatorname{subspace\text{-}ms}(\operatorname{nf\text{-}metric\text{-}space}(k), u))(a, b) \leq r) \end{array}}}
\vnbprop{nf-zero-in-carr}{\vnbstmt{\forall k.\; (\operatorname{is\text{-}normed\text{-}field}(k) \Rightarrow (\operatorname{zero}(k) \in \operatorname{carr}(k)))}}
\vnbprop{nf-norm-of-zero}{\vnbstmt{\forall k.\; (\operatorname{is\text{-}normed\text{-}field}(k) \Rightarrow (\operatorname{fnrm}(k)(\operatorname{zero}(k)) = 0))}}
\vnbprop{nf-sub-zero-eq}{\vnbstmt{\begin{array}{@{}l@{}} \forall k, u, v. \\ \quad (\operatorname{is\text{-}normed\text{-}field}(k) \wedge ((u \in \operatorname{carr}(k)) \wedge ((v \in \operatorname{carr}(k)) \wedge (\operatorname{add}(k)(u, \operatorname{neg}(k)(v)) = \operatorname{zero}(k))))) \Rightarrow  \\ \quad \quad u = v \end{array}}}
\vnbprop{nf-mul-cancel}{\vnbstmt{\begin{array}{@{}l@{}} \forall k, u, v, w. \\ \quad \operatorname{is\text{-}normed\text{-}field}(k) \wedge \\ \quad u \in \operatorname{carr}(k) \wedge \\ \quad v \in \operatorname{carr}(k) \wedge \\ \quad w \in \operatorname{carr}(k) \wedge \\ \quad \neg (w = \operatorname{zero}(k)) \wedge \\ \quad \operatorname{mul}(k)(u, w) = \operatorname{mul}(k)(v, w) \Rightarrow \\ \quad \quad u = v \end{array}}}
\vnbprop{diff-on-unique}{\vnbstmt{\begin{array}{@{}l@{}} \forall k, u, f, a, l, m. \\ \quad \operatorname{is\text{-}normed\text{-}field}(k) \wedge \\ \quad \forall e. \\ \quad \quad \operatorname{pos\text{-}rr}(e) \Rightarrow \exists h \in \operatorname{carr}(k).\; (\neg (h = \operatorname{zero}(k)) \wedge (\operatorname{fnrm}(k)(h) < e)) \wedge \\ \quad \operatorname{is\text{-}diff\text{-}on}(k, u, f, a, l) \wedge \\ \quad \operatorname{is\text{-}diff\text{-}on}(k, u, f, a, m) \Rightarrow \\ \quad \quad l = m \end{array}}}
\vnbprop{nf-one-in-carr}{\vnbstmt{\forall k.\; (\operatorname{is\text{-}normed\text{-}field}(k) \Rightarrow (\operatorname{one}(k) \in \operatorname{carr}(k)))}}
\vnbprop{nf-open-is-set}{\vnbstmt{\begin{array}{@{}l@{}} \forall k, u. \\ \quad (\operatorname{is\text{-}normed\text{-}field}(k) \wedge \operatorname{is\text{-}open}(\operatorname{nf\text{-}metric\text{-}space}(k), u)) \Rightarrow (u \in \mathbf{Set}) \end{array}}}
\vnbprop{diff-on-const}{\vnbstmt{\begin{array}{@{}l@{}} \forall k, u.; \forall c \in \operatorname{carr}(k), a \in u. \\ \quad (\operatorname{is\text{-}normed\text{-}field}(k) \wedge \operatorname{is\text{-}open}(\operatorname{nf\text{-}metric\text{-}space}(k), u)) \Rightarrow  \\ \quad \quad \operatorname{is\text{-}diff\text{-}on}(k, u, (\lambda x \in u.\; c), a, \operatorname{zero}(k)) \end{array}}}
\vnbprop{diff-on-identity}{\vnbstmt{\begin{array}{@{}l@{}} \forall k, u.; \forall a \in u. \\ \quad (\operatorname{is\text{-}normed\text{-}field}(k) \wedge \operatorname{is\text{-}open}(\operatorname{nf\text{-}metric\text{-}space}(k), u)) \Rightarrow  \\ \quad \quad \operatorname{is\text{-}diff\text{-}on}(k, u, (\lambda x \in u.\; x), a, \operatorname{one}(k)) \end{array}}}
\vnbprop{diff-on-sum}{\vnbstmt{\begin{array}{@{}l@{}} \forall k, u, f, g, a, l, m. \\ \quad (\operatorname{is\text{-}normed\text{-}field}(k) \wedge (\operatorname{is\text{-}diff\text{-}on}(k, u, f, a, l) \wedge \operatorname{is\text{-}diff\text{-}on}(k, u, g, a, m))) \Rightarrow  \\ \quad \quad \operatorname{is\text{-}diff\text{-}on}(k, u, (\lambda x \in u.\; \operatorname{add}(k)(f(x), g(x))), a, \operatorname{add}(k)(l, m)) \end{array}}}
\vnbprop{nf-lam-const-z-in-fun}{\vnbstmt{\begin{array}{@{}l@{}} \forall k.; \forall u \in \mathbf{Set}, c \in \operatorname{carr}(k). \\ \quad \operatorname{is\text{-}normed\text{-}field}(k) \Rightarrow ((\lambda z \in u.\; c) \in (u \to \operatorname{carr}(k))) \end{array}}}
\vnbprop{nf-lam-prod-v-in-fun}{\vnbstmt{\begin{array}{@{}l@{}} \forall k.; \forall u \in \mathbf{Set}, f, g \in (u \to \operatorname{carr}(k)). \\ \quad \operatorname{is\text{-}normed\text{-}field}(k) \Rightarrow  \\ \quad \quad (\lambda v \in u.\; \operatorname{mul}(k)(f(v), g(v))) \in (u \to \operatorname{carr}(k)) \end{array}}}
\vnbprop{nf-product-continuous-v}{\vnbstmt{\begin{array}{@{}l@{}} \forall k, s, d, g, h, a. \\ \quad \operatorname{is\text{-}normed\text{-}field}(k) \wedge \\ \quad (\operatorname{pts}(s) \simeq d) \wedge \\ \quad \operatorname{is\text{-}continuous\text{-}at}(s, \operatorname{nf\text{-}metric\text{-}space}(k), g, a) \wedge \\ \quad \operatorname{is\text{-}continuous\text{-}at}(s, \operatorname{nf\text{-}metric\text{-}space}(k), h, a) \Rightarrow \\ \quad \quad \operatorname{is\text{-}continuous\text{-}at}(s, \operatorname{nf\text{-}metric\text{-}space}(k), (\lambda v \in d.\; \operatorname{mul}(k)(g(v), h(v))), a) \end{array}}}
\vnbprop{diff-on-product}{\vnbstmt{\begin{array}{@{}l@{}} \forall k, u, f, g, a, l, m. \\ \quad (\operatorname{is\text{-}normed\text{-}field}(k) \wedge (\operatorname{is\text{-}diff\text{-}on}(k, u, f, a, l) \wedge \operatorname{is\text{-}diff\text{-}on}(k, u, g, a, m))) \Rightarrow  \\ \quad \quad \operatorname{is\text{-}diff\text{-}on}(k, u, (\lambda x \in u.\; \operatorname{mul}(k)(f(x), g(x))), a, \\ \quad \quad \quad \operatorname{add}(k)(\operatorname{mul}(k)(l, g(a)), \operatorname{mul}(k)(m, f(a)))) \end{array}}}
\vnbprop{ms-corestrict-continuous}{\vnbstmt{\begin{array}{@{}l@{}} \forall s, t, \mathit{d2s}, f, a. \\ \quad \operatorname{is\text{-}metric\text{-}space}(t) \wedge \\ \quad \mathit{d2s} \subseteq \operatorname{pts}(t) \wedge \\ \quad \operatorname{is\text{-}continuous\text{-}at}(s, t, f, a) \wedge \\ \quad \forall y \in \operatorname{pts}(s).\; (f(y) \in \mathit{d2s}) \Rightarrow \\ \quad \quad \operatorname{is\text{-}continuous\text{-}at}(s, \operatorname{subspace\text{-}ms}(t, \mathit{d2s}), f, a) \end{array}}}
\vnbprop{nf-lam-comp-z-in-fun}{\vnbstmt{\begin{array}{@{}l@{}} \forall k.; \forall u \in \mathbf{Set}, s, f.; \forall g \in (s \to \operatorname{carr}(k)). \\ \quad (\operatorname{is\text{-}normed\text{-}field}(k) \wedge \forall y \in u.\; (f(y) \in s)) \Rightarrow  \\ \quad \quad (\lambda z \in u.\; g(f(z))) \in (u \to \operatorname{carr}(k)) \end{array}}}
\vnbprop{diff-on-chain}{\vnbstmt{\begin{array}{@{}l@{}} \forall k, u, s, f, g, a, l, m. \\ \quad \operatorname{is\text{-}normed\text{-}field}(k) \wedge \\ \quad \operatorname{is\text{-}diff\text{-}on}(k, u, f, a, l) \wedge \\ \quad \operatorname{is\text{-}diff\text{-}on}(k, s, g, f(a), m) \wedge \\ \quad \forall y \in u.\; (f(y) \in s) \Rightarrow \\ \quad \quad \operatorname{is\text{-}diff\text{-}on}(k, u, (\lambda x \in u.\; g(f(x))), a, \operatorname{mul}(k)(m, l)) \end{array}}}

\vnbbookfile{theorem-library/holomorphic-basics}{theorem-library/holomorphic-basics}
\noindent{\small holomorphic-basics.scm -- the COMPLEX side of IS-DIFF-ON / HOLOMORPHIC-ON.}\par
\vnbprop{cc-nf-add-apply}{\vnbstmt{\begin{array}{@{}l@{}} \forall u, v. \\ \quad ((u \in \mathbb{C}) \wedge (v \in \mathbb{C})) \Rightarrow (\operatorname{add}(\operatorname{cc\text{-}normed\text{-}field})(u, v) \simeq (u + v)) \end{array}}}
\vnbprop{cc-nf-mul-apply}{\vnbstmt{\begin{array}{@{}l@{}} \forall u, v. \\ \quad ((u \in \mathbb{C}) \wedge (v \in \mathbb{C})) \Rightarrow (\operatorname{mul}(\operatorname{cc\text{-}normed\text{-}field})(u, v) \simeq (u \cdot v)) \end{array}}}
\vnbprop{cc-nf-neg-apply}{\vnbstmt{\forall u \in \mathbb{C}.\; (\operatorname{neg}(\operatorname{cc\text{-}normed\text{-}field})(u) \simeq -u)}}
\vnbprop{cc-nf-sub}{\vnbstmt{\begin{array}{@{}l@{}} \forall u, v. \\ \quad ((u \in \mathbb{C}) \wedge (v \in \mathbb{C})) \Rightarrow  \\ \quad \quad (\operatorname{add}(\operatorname{cc\text{-}normed\text{-}field})(u, \operatorname{neg}(\operatorname{cc\text{-}normed\text{-}field})(v)) \simeq (u - v)) \end{array}}}
\vnbprop{nf-cc-is-metric-space}{\vnbstmt{\operatorname{is\text{-}metric\text{-}space}(\operatorname{nf\text{-}metric\text{-}space}(\operatorname{cc\text{-}normed\text{-}field}))}}
\vnbprop{nf-cc-agree-pts}{\vnbstmt{(\operatorname{pts}(\operatorname{nf\text{-}metric\text{-}space}(\operatorname{cc\text{-}normed\text{-}field})) \simeq \operatorname{pts}(\operatorname{cc\text{-}ms}))}}
\vnbprop{nf-cc-agree-dist}{\vnbstmt{\begin{array}{@{}l@{}} \forall u, v \in \operatorname{pts}(\operatorname{cc\text{-}ms}). \\ \quad (\operatorname{dist}(\operatorname{nf\text{-}metric\text{-}space}(\operatorname{cc\text{-}normed\text{-}field}))(u, v) \simeq \operatorname{dist}(\operatorname{cc\text{-}ms})(u, v)) \end{array}}}
\vnbprop{cc-ms-open-iff}{\vnbstmt{\forall w.\; (\operatorname{is\text{-}open}(\operatorname{nf\text{-}metric\text{-}space}(\operatorname{cc\text{-}normed\text{-}field}), w) \Leftrightarrow \operatorname{is\text{-}open}(\operatorname{cc\text{-}ms}, w))}}
\vnbprop{cc-nf-small-elements}{\vnbstmt{\begin{array}{@{}l@{}} \forall s. \\ \quad \operatorname{pos\text{-}rr}(s) \Rightarrow  \\ \quad \quad \exists h \in \operatorname{carr}(\operatorname{cc\text{-}normed\text{-}field}). \\ \quad \quad \quad \neg (h = \operatorname{zero}(\operatorname{cc\text{-}normed\text{-}field})) \wedge (\operatorname{fnrm}(\operatorname{cc\text{-}normed\text{-}field})(h) < s) \end{array}}}
\vnbprop{holomorphic-on-open-cc-ms}{\vnbstmt{\forall u, f.\; (\operatorname{holomorphic\text{-}on}(u, f) \Rightarrow \operatorname{is\text{-}open}(\operatorname{cc\text{-}ms}, u))}}
\vnbprop{holomorphic-on-iff-cc-ms}{\vnbstmt{\begin{array}{@{}l@{}} \forall u, f. \\ \quad \operatorname{holomorphic\text{-}on}(u, f) \\ \quad \Leftrightarrow\; \operatorname{is\text{-}open}(\operatorname{cc\text{-}ms}, u) \wedge \\ \quad f \in (u \to \mathbb{C}) \wedge \\ \quad \forall a \in u.\; \exists l.\; \operatorname{is\text{-}diff\text{-}on}(\operatorname{cc\text{-}normed\text{-}field}, u, f, a, l) \end{array}}}

\vnbbookfile{theorem-library/deriv-on-read-off}{theorem-library/deriv-on-read-off}
\noindent{\small deriv-on-read-off.scm -- DERIV-ON reads off IS-DIFF-ON. Batch 27-A follow-on, 2026-09-24.}\par
\vnbprop{deriv-on-of-is-diff-on}{\vnbstmt{\begin{array}{@{}l@{}} \forall k, u, f, a, l. \\ \quad \operatorname{is\text{-}normed\text{-}field}(k) \wedge \\ \quad \forall e. \\ \quad \quad \operatorname{pos\text{-}rr}(e) \Rightarrow \exists h \in \operatorname{carr}(k).\; (\neg (h = \operatorname{zero}(k)) \wedge (\operatorname{fnrm}(k)(h) < e)) \wedge \\ \quad \operatorname{is\text{-}diff\text{-}on}(k, u, f, a, l) \Rightarrow \\ \quad \quad \operatorname{deriv\text{-}on}(k, u, f, a) = l \end{array}}}
\vnbprop{cc-deriv-on-of-is-diff-on}{\vnbstmt{\begin{array}{@{}l@{}} \forall u, f, a, l. \\ \quad \operatorname{is\text{-}diff\text{-}on}(\operatorname{cc\text{-}normed\text{-}field}, u, f, a, l) \Rightarrow  \\ \quad \quad \operatorname{deriv\text{-}on}(\operatorname{cc\text{-}normed\text{-}field}, u, f, a) = l \end{array}}}

\vnbbookfile{theorem-library/holomorphic-basics-2}{theorem-library/holomorphic-basics-2}
\noindent{\small holomorphic-basics-2.scm -- section 2.1 of the user's complex-analysis notes, finished: the constant and the identity are holomorphic on any open set, sums and products of holomorphic functions are holomorphic, and so is z |-$>$ z\textasciicircum{}n for every natural n.}\par
\vnbprop{diff-on-transfer-ptwise-eq}{\vnbstmt{\begin{array}{@{}l@{}} \forall k, u, f, g, a, l. \\ \quad ((f \in (u \to \operatorname{carr}(k))) \wedge (\operatorname{is\text{-}diff\text{-}on}(k, u, g, a, l) \wedge \forall x \in u.\; (f(x) \simeq g(x)))) \Rightarrow  \\ \quad \quad \operatorname{is\text{-}diff\text{-}on}(k, u, f, a, l) \end{array}}}
\vnbprop{holomorphic-transfer-ptwise-eq}{\vnbstmt{\begin{array}{@{}l@{}} \forall u, f, g. \\ \quad (\operatorname{holomorphic\text{-}on}(u, g) \wedge ((f \in (u \to \mathbb{C})) \wedge \forall x \in u.\; (f(x) \simeq g(x)))) \Rightarrow  \\ \quad \quad \operatorname{holomorphic\text{-}on}(u, f) \end{array}}}
\vnbprop{holomorphic-const}{\vnbstmt{\begin{array}{@{}l@{}} \forall u.; \forall c \in \mathbb{C}. \\ \quad \operatorname{is\text{-}open}(\operatorname{cc\text{-}ms}, u) \Rightarrow \operatorname{holomorphic\text{-}on}(u, (\lambda x \in u.\; c)) \end{array}}}
\vnbprop{holomorphic-identity}{\vnbstmt{\forall u.\; (\operatorname{is\text{-}open}(\operatorname{cc\text{-}ms}, u) \Rightarrow \operatorname{holomorphic\text{-}on}(u, (\lambda x \in u.\; x)))}}
\vnbprop{holomorphic-sum}{\vnbstmt{\begin{array}{@{}l@{}} \forall u, f, g, m. \\ \quad \operatorname{holomorphic\text{-}on}(u, f) \wedge \\ \quad \operatorname{holomorphic\text{-}on}(u, g) \wedge \\ \quad m \in (u \to \mathbb{C}) \wedge \\ \quad \forall x \in u.\; (m(x) \simeq (f(x) + g(x))) \Rightarrow \\ \quad \quad \operatorname{holomorphic\text{-}on}(u, m) \end{array}}}
\vnbprop{holomorphic-product}{\vnbstmt{\begin{array}{@{}l@{}} \forall u, f, g, m. \\ \quad \operatorname{holomorphic\text{-}on}(u, f) \wedge \\ \quad \operatorname{holomorphic\text{-}on}(u, g) \wedge \\ \quad m \in (u \to \mathbb{C}) \wedge \\ \quad \forall x \in u.\; (m(x) \simeq (f(x) \cdot g(x))) \Rightarrow \\ \quad \quad \operatorname{holomorphic\text{-}on}(u, m) \end{array}}}
\vnbprop{holomorphic-power}{\vnbstmt{\begin{array}{@{}l@{}} \forall n \in \mathbb{N}, u.; \forall w \in (u \to \mathbb{C}). \\ \quad (\operatorname{is\text{-}open}(\operatorname{cc\text{-}ms}, u) \wedge \forall z \in u.\; (w(z) \simeq \operatorname{power}(z, n))) \Rightarrow  \\ \quad \quad \operatorname{holomorphic\text{-}on}(u, w) \end{array}}}

\vnbbookfile{theorem-library/cc-reciprocal}{theorem-library/cc-reciprocal}
\noindent{\small cc-reciprocal.scm -- z |-$>$ 1/z: the punctured plane, the continuity of the reciprocal, its derivative, and the reciprocal and quotient rules for HOLOMORPHIC-ON.}\par
\vnbprop{cc-nonzero-membership}{\vnbstmt{\forall w.\; ((w \in \operatorname{difference}(\mathbb{C}, \operatorname{singleton}(0))) \Leftrightarrow ((w \in \mathbb{C}) \wedge \neg (w = 0)))}}
\vnbprop{cc-nonzero-subset}{\vnbstmt{\operatorname{difference}(\mathbb{C}, \operatorname{singleton}(0)) \subseteq \mathbb{C}}}
\vnbprop{cc-nonzero-is-set}{\vnbstmt{\operatorname{difference}(\mathbb{C}, \operatorname{singleton}(0)) \in \mathbf{Set}}}
\vnbprop{cc-nonzero-open}{\vnbstmt{\operatorname{is\text{-}open}(\operatorname{cc\text{-}ms}, \operatorname{difference}(\mathbb{C}, \operatorname{singleton}(0)))}}
\vnbprop{cc-nonzero-open-nf}{\vnbstmt{\operatorname{is\text{-}open}(\operatorname{nf\text{-}metric\text{-}space}(\operatorname{cc\text{-}normed\text{-}field}), \operatorname{difference}(\mathbb{C}, \operatorname{singleton}(0)))}}
\vnbprop{cc-recip-lam-in-fun}{\vnbstmt{(\lambda z \in \operatorname{difference}(\mathbb{C}, \operatorname{singleton}(0)).\; \operatorname{recip}(z)) \in (\operatorname{difference}(\mathbb{C}, \operatorname{singleton}(0)) \to \mathbb{C})}}
\vnbprop{cc-recip-continuous-at}{\vnbstmt{\begin{array}{@{}l@{}} \forall a \in \operatorname{difference}(\mathbb{C}, \operatorname{singleton}(0)). \\ \quad \operatorname{is\text{-}continuous\text{-}at}(\operatorname{subspace\text{-}ms}(\operatorname{nf\text{-}metric\text{-}space}(\operatorname{cc\text{-}normed\text{-}field}), \operatorname{difference}(\mathbb{C}, \operatorname{singleton}(0))), \\ \quad \quad \operatorname{nf\text{-}metric\text{-}space}(\operatorname{cc\text{-}normed\text{-}field}), \\ \quad \quad (\lambda z \in \operatorname{difference}(\mathbb{C}, \operatorname{singleton}(0)).\; \operatorname{recip}(z)), \\ \quad \quad a) \end{array}}}
\vnbprop{cc-mul-nonzero}{\vnbstmt{\forall u, v \in \mathbb{C}.\; ((\neg (u = 0) \wedge \neg (v = 0)) \Rightarrow \neg ((u \cdot v) = 0))}}
\vnbprop{cc-recip-mul}{\vnbstmt{\begin{array}{@{}l@{}} \forall u, v \in \mathbb{C}. \\ \quad (\neg (u = 0) \wedge \neg (v = 0)) \Rightarrow (\operatorname{recip}((u \cdot v)) = (\operatorname{recip}(u) \cdot \operatorname{recip}(v))) \end{array}}}
\vnbprop{cc-recip-diff-on}{\vnbstmt{\begin{array}{@{}l@{}} \forall a \in \operatorname{difference}(\mathbb{C}, \operatorname{singleton}(0)). \\ \quad \operatorname{is\text{-}diff\text{-}on}(\operatorname{cc\text{-}normed\text{-}field}, \operatorname{difference}(\mathbb{C}, \operatorname{singleton}(0)), \\ \quad \quad (\lambda z \in \operatorname{difference}(\mathbb{C}, \operatorname{singleton}(0)).\; \operatorname{recip}(z)), \\ \quad \quad a, -\operatorname{recip}((a \cdot a))) \end{array}}}
\vnbprop{holomorphic-recip}{\vnbstmt{\operatorname{holomorphic\text{-}on}(\operatorname{difference}(\mathbb{C}, \operatorname{singleton}(0)), (\lambda z \in \operatorname{difference}(\mathbb{C}, \operatorname{singleton}(0)).\; \operatorname{recip}(z)))}}
\vnbprop{diff-on-recip-cc}{\vnbstmt{\begin{array}{@{}l@{}} \forall u, f, a, l, q. \\ \quad \operatorname{is\text{-}diff\text{-}on}(\operatorname{cc\text{-}normed\text{-}field}, u, f, a, l) \wedge \\ \quad \forall y \in u.\; \neg (f(y) = 0) \wedge \\ \quad q \in (u \to \mathbb{C}) \wedge \\ \quad \forall y \in u.\; (q(y) \simeq \operatorname{recip}(f(y))) \Rightarrow \\ \quad \quad \operatorname{is\text{-}diff\text{-}on}(\operatorname{cc\text{-}normed\text{-}field}, u, q, a, -(l \cdot \operatorname{recip}((f(a) \cdot f(a))))) \end{array}}}
\vnbprop{holomorphic-recip-of}{\vnbstmt{\begin{array}{@{}l@{}} \forall u, f, q. \\ \quad \operatorname{holomorphic\text{-}on}(u, f) \wedge \\ \quad \forall y \in u.\; \neg (f(y) = 0) \wedge \\ \quad q \in (u \to \mathbb{C}) \wedge \\ \quad \forall y \in u.\; (q(y) \simeq \operatorname{recip}(f(y))) \Rightarrow \\ \quad \quad \operatorname{holomorphic\text{-}on}(u, q) \end{array}}}
\vnbprop{holomorphic-quotient}{\vnbstmt{\begin{array}{@{}l@{}} \forall u, f, g, q. \\ \quad \operatorname{holomorphic\text{-}on}(u, f) \wedge \\ \quad \operatorname{holomorphic\text{-}on}(u, g) \wedge \\ \quad \forall y \in u.\; \neg (g(y) = 0) \wedge \\ \quad q \in (u \to \mathbb{C}) \wedge \\ \quad \forall y \in u.\; (q(y) \simeq /(f(y), g(y))) \Rightarrow \\ \quad \quad \operatorname{holomorphic\text{-}on}(u, q) \end{array}}}

\vnbbookfile{structure-library/interval-calculus}{structure-library/interval-calculus}
\noindent{\small interval-calculus.scm -- THE VOCABULARY OF THE CALCULUS ON AN INTERVAL. DEFINITIONS ONLY; every law is PROVEN in theorem-library/interval-calculus-laws.scm and interval-mvt.scm.}\par
\par\medskip\noindent\textbf{Definition} (\texttt{ooint}).\;$\operatorname{ooint}(a, b) :=$ $\{\operatorname{icx\_} \in \mathbb{R} : ((a < \operatorname{icx\_}) \wedge (\operatorname{icx\_} < b))\}$
\par\medskip\noindent\textbf{Definition} (\texttt{is-continuous-on}).\;For all $f$ and $d$, 
\begin{equation*}
\begin{array}{@{}l@{}} \operatorname{is\text{-}continuous\text{-}on}(f, d) \\ \Leftrightarrow\; ((f \in (d \to \mathbb{R})) \wedge \forall \operatorname{icp\_} \in d.\; \operatorname{is\text{-}continuous\text{-}at}(\operatorname{subspace\text{-}ms}(\operatorname{rr\text{-}ms}, d), \operatorname{rr\text{-}ms}, f, \operatorname{icp\_})) \end{array}
\end{equation*}

\par\medskip\noindent\textbf{Definition} (\texttt{has-deriv-at}).\;For all $f$, $a$ and $l$, 
\begin{equation*}
\begin{array}{@{}l@{}} \operatorname{has\text{-}deriv\text{-}at}(f, a, l) \\ \Leftrightarrow\; \exists \operatorname{ice\_}. \\ \quad \operatorname{pos\text{-}rr}(\operatorname{ice\_}) \wedge \\ \quad \operatorname{is\text{-}diff\text{-}on}(\operatorname{rr\text{-}normed\text{-}field}, \operatorname{ooint}(-(a, \operatorname{ice\_}), +(a, \operatorname{ice\_})), \operatorname{restrict}(f, \operatorname{ooint}(-(a, \operatorname{ice\_}), +(a, \operatorname{ice\_}))), a, l) \end{array}
\end{equation*}

\par\medskip\noindent\textbf{Definition} (\texttt{extend-const}).\;$\operatorname{extend\text{-}const}(f, a, b) :=$ $(\lambda \operatorname{ecx\_} \in \mathbb{R}.\; f(\operatorname{clamp}(a, b, \operatorname{ecx\_})))$

\vnbbookfile{structure-library/regulated-primitive}{structure-library/regulated-primitive}
\noindent{\small regulated-primitive.scm -- COUNTABLE SETS, ONE-SIDED LIMITS WITHIN A SET, REGULATED FUNCTIONS ON AN INTERVAL, PRIMITIVES.  DEFINITIONS ONLY; the laws are PROVEN in theorem-library/regulated-primitive-laws.scm.}\par
\par\medskip\noindent\textbf{Definition} (\texttt{is-countable}).\;For every $\operatorname{rpc\_}$, 
\begin{equation*}
\begin{array}{@{}l@{}} \operatorname{is\text{-}countable}(\operatorname{rpc\_}) \\ \Leftrightarrow\; \operatorname{rpc\_} \in \mathbf{Set} \wedge \\ \operatorname{rpc\_} = \emptyset \vee \\ \exists \operatorname{rpen\_} \in (\mathbb{N} \to \operatorname{rpc\_}). \forall \operatorname{rpy\_} \in \operatorname{rpc\_}. \exists \operatorname{rpn\_} \in \mathbb{N}. \\ \quad \operatorname{rpy\_} = \operatorname{rpen\_}(\operatorname{rpn\_}) \end{array}
\end{equation*}

\par\medskip\noindent\textbf{Definition} (\texttt{is-right-limit-within}).\;For all $\operatorname{rpf\_}$, $\operatorname{rpw\_}$, $\operatorname{rpx\_}$ and $\operatorname{rpl\_}$, 
\begin{equation*}
\begin{array}{@{}l@{}} \operatorname{is\text{-}right\text{-}limit\text{-}within}(\operatorname{rpf\_}, \operatorname{rpw\_}, \operatorname{rpx\_}, \operatorname{rpl\_}) \\ \Leftrightarrow\; \operatorname{rpw\_} \subseteq \mathbb{R} \wedge \\ \operatorname{rpf\_} \in (\operatorname{rpw\_} \to \mathbb{R}) \wedge \\ \operatorname{rpx\_} \in \mathbb{R} \wedge \\ \operatorname{rpl\_} \in \mathbb{R} \wedge \\ \forall \operatorname{rpe\_}. \\ \quad \operatorname{pos\text{-}rr}(\operatorname{rpe\_}) \Rightarrow  \\ \quad \quad \exists \operatorname{rpdl\_}. \\ \quad \quad \quad \operatorname{pos\text{-}rr}(\operatorname{rpdl\_}) \wedge \\ \quad \quad \quad \forall \operatorname{rpt\_} \in \operatorname{rpw\_}. \\ \quad \quad \quad \quad (\operatorname{rpx\_} < \operatorname{rpt\_}) \Rightarrow ((\operatorname{rpt\_} < +(\operatorname{rpx\_}, \operatorname{rpdl\_})) \Rightarrow (\lvert -(\operatorname{rpf\_}(\operatorname{rpt\_}), \operatorname{rpl\_}) \rvert < \operatorname{rpe\_})) \end{array}
\end{equation*}

\par\medskip\noindent\textbf{Definition} (\texttt{is-left-limit-within}).\;For all $\operatorname{rpf\_}$, $\operatorname{rpw\_}$, $\operatorname{rpx\_}$ and $\operatorname{rpl\_}$, 
\begin{equation*}
\begin{array}{@{}l@{}} \operatorname{is\text{-}left\text{-}limit\text{-}within}(\operatorname{rpf\_}, \operatorname{rpw\_}, \operatorname{rpx\_}, \operatorname{rpl\_}) \\ \Leftrightarrow\; \operatorname{rpw\_} \subseteq \mathbb{R} \wedge \\ \operatorname{rpf\_} \in (\operatorname{rpw\_} \to \mathbb{R}) \wedge \\ \operatorname{rpx\_} \in \mathbb{R} \wedge \\ \operatorname{rpl\_} \in \mathbb{R} \wedge \\ \forall \operatorname{rple\_}. \\ \quad \operatorname{pos\text{-}rr}(\operatorname{rple\_}) \Rightarrow  \\ \quad \quad \exists \operatorname{rpld\_}. \\ \quad \quad \quad \operatorname{pos\text{-}rr}(\operatorname{rpld\_}) \wedge \\ \quad \quad \quad \forall \operatorname{rplt\_} \in \operatorname{rpw\_}. \\ \quad \quad \quad \quad (-(\operatorname{rpx\_}, \operatorname{rpld\_}) < \operatorname{rplt\_}) \Rightarrow  \\ \quad \quad \quad \quad (\operatorname{rplt\_} < \operatorname{rpx\_}) \Rightarrow (\lvert -(\operatorname{rpf\_}(\operatorname{rplt\_}), \operatorname{rpl\_}) \rvert < \operatorname{rple\_}) \end{array}
\end{equation*}

\par\medskip\noindent\textbf{Definition} (\texttt{is-regulated-on}).\;For all $\operatorname{rpf\_}$, $\operatorname{rpa\_}$ and $\operatorname{rpb\_}$, 
\begin{equation*}
\begin{array}{@{}l@{}} \operatorname{is\text{-}regulated\text{-}on}(\operatorname{rpf\_}, \operatorname{rpa\_}, \operatorname{rpb\_}) \\ \Leftrightarrow\; \operatorname{rpa\_} \in \mathbb{R} \wedge \\ \operatorname{rpb\_} \in \mathbb{R} \wedge \\ \operatorname{rpa\_} < \operatorname{rpb\_} \wedge \\ \operatorname{rpf\_} \in (\operatorname{ccint}(\operatorname{rpa\_}, \operatorname{rpb\_}) \to \mathbb{R}) \wedge \\ \forall \operatorname{rpx\_} \in \operatorname{ccint}(\operatorname{rpa\_}, \operatorname{rpb\_}). \\ \quad (\operatorname{rpx\_} < \operatorname{rpb\_}) \Rightarrow  \\ \quad \quad \exists \operatorname{rpl\_}.\; \operatorname{is\text{-}right\text{-}limit\text{-}within}(\operatorname{rpf\_}, \operatorname{ccint}(\operatorname{rpa\_}, \operatorname{rpb\_}), \operatorname{rpx\_}, \operatorname{rpl\_}) \wedge \\ \forall \operatorname{rpx\_} \in \operatorname{ccint}(\operatorname{rpa\_}, \operatorname{rpb\_}). \\ \quad (\operatorname{rpa\_} < \operatorname{rpx\_}) \Rightarrow  \\ \quad \quad \exists \operatorname{rpl\_}.\; \operatorname{is\text{-}left\text{-}limit\text{-}within}(\operatorname{rpf\_}, \operatorname{ccint}(\operatorname{rpa\_}, \operatorname{rpb\_}), \operatorname{rpx\_}, \operatorname{rpl\_}) \end{array}
\end{equation*}

\par\medskip\noindent\textbf{Definition} (\texttt{is-primitive}).\;For all $\operatorname{rpg\_}$, $\operatorname{rpf\_}$, $\operatorname{rpa\_}$ and $\operatorname{rpb\_}$, 
\begin{equation*}
\begin{array}{@{}l@{}} \operatorname{is\text{-}primitive}(\operatorname{rpg\_}, \operatorname{rpf\_}, \operatorname{rpa\_}, \operatorname{rpb\_}) \\ \Leftrightarrow\; \operatorname{rpa\_} \in \mathbb{R} \wedge \\ \operatorname{rpb\_} \in \mathbb{R} \wedge \\ \operatorname{rpa\_} < \operatorname{rpb\_} \wedge \\ \operatorname{rpg\_} \in (\operatorname{ccint}(\operatorname{rpa\_}, \operatorname{rpb\_}) \to \mathbb{R}) \wedge \\ \operatorname{rpf\_} \in (\operatorname{ccint}(\operatorname{rpa\_}, \operatorname{rpb\_}) \to \mathbb{R}) \wedge \\ \operatorname{is\text{-}continuous\text{-}on}(\operatorname{rpg\_}, \operatorname{ccint}(\operatorname{rpa\_}, \operatorname{rpb\_})) \wedge \\ \exists \operatorname{rps\_}. \\ \quad \operatorname{is\text{-}countable}(\operatorname{rps\_}) \wedge \\ \quad \operatorname{rps\_} \subseteq \operatorname{ccint}(\operatorname{rpa\_}, \operatorname{rpb\_}) \wedge \\ \quad \forall \operatorname{rpt\_} \in \operatorname{ooint}(\operatorname{rpa\_}, \operatorname{rpb\_}). \\ \quad \quad \neg (\operatorname{rpt\_} \in \operatorname{rps\_}) \Rightarrow \operatorname{has\text{-}deriv\text{-}at}(\operatorname{rpg\_}, \operatorname{rpt\_}, \operatorname{rpf\_}(\operatorname{rpt\_})) \end{array}
\end{equation*}


\vnbbookfile{theorem-library/interval-calculus-laws}{theorem-library/interval-calculus-laws}
\noindent{\small interval-calculus-laws.scm -- the laws of the vocabulary of structure-library/interval-calculus.scm: OOINT, IS-CONTINUOUS-ON, HAS-DERIV-AT and EXTEND-CONST.  Everything here is PROVEN; the definition file asserts nothing.}\par
\vnbprop{ooint-membership}{\vnbstmt{\forall a, b, y.\; ((y \in \operatorname{ooint}(a, b)) \Leftrightarrow ((y \in \mathbb{R}) \wedge ((a < y) \wedge (y < b))))}}
\vnbprop{ooint-elt-in-rr}{\vnbstmt{\forall a, b.; \forall y \in \operatorname{ooint}(a, b).\; (y \in \mathbb{R})}}
\vnbprop{ooint-subset-rr}{\vnbstmt{\forall a, b.\; (\operatorname{ooint}(a, b) \subseteq \mathbb{R})}}
\vnbprop{ooint-subset-ccint}{\vnbstmt{\forall a, b.\; (((a \in \mathbb{R}) \wedge (b \in \mathbb{R})) \Rightarrow (\operatorname{ooint}(a, b) \subseteq \operatorname{ccint}(a, b)))}}
\vnbprop{ooint-center}{\vnbstmt{\forall y, e.\; (((y \in \mathbb{R}) \wedge \operatorname{pos\text{-}rr}(e)) \Rightarrow (y \in \operatorname{ooint}((y - e), (y + e))))}}
\vnbprop{ooint-open}{\vnbstmt{\forall a, b.\; (((a \in \mathbb{R}) \wedge (b \in \mathbb{R})) \Rightarrow \operatorname{is\text{-}open}(\operatorname{rr\text{-}ms}, \operatorname{ooint}(a, b)))}}
\vnbprop{ooint-open-nf}{\vnbstmt{\begin{array}{@{}l@{}} \forall a, b. \\ \quad ((a \in \mathbb{R}) \wedge (b \in \mathbb{R})) \Rightarrow  \\ \quad \quad \operatorname{is\text{-}open}(\operatorname{nf\text{-}metric\text{-}space}(\operatorname{rr\text{-}normed\text{-}field}), \operatorname{ooint}(a, b)) \end{array}}}
\vnbprop{subspace-restrict-continuous-at}{\vnbstmt{\begin{array}{@{}l@{}} \forall s, u, v, t, f.; \forall z \in v. \\ \quad \operatorname{is\text{-}metric\text{-}space}(s) \wedge \\ \quad u \subseteq \operatorname{pts}(s) \wedge \\ \quad v \subseteq u \wedge \\ \quad \operatorname{is\text{-}metric\text{-}space}(t) \wedge \\ \quad \operatorname{is\text{-}continuous\text{-}at}(\operatorname{subspace\text{-}ms}(s, u), t, f, z) \Rightarrow \\ \quad \quad \operatorname{is\text{-}continuous\text{-}at}(\operatorname{subspace\text{-}ms}(s, v), t, \operatorname{restrict}(f, v), z) \end{array}}}
\vnbprop{diff-on-shrink}{\vnbstmt{\begin{array}{@{}l@{}} \forall k, u, f, a, l, v. \\ \quad (\operatorname{is\text{-}diff\text{-}on}(k, u, f, a, l) \wedge (\operatorname{is\text{-}open}(\operatorname{nf\text{-}metric\text{-}space}(k), v) \wedge ((v \subseteq u) \wedge (a \in v)))) \Rightarrow  \\ \quad \quad \operatorname{is\text{-}diff\text{-}on}(k, v, \operatorname{restrict}(f, v), a, l) \end{array}}}
\vnbprop{restrict-in-fun-of-restrict}{\vnbstmt{\begin{array}{@{}l@{}} \forall f, u, c, v. \\ \quad ((\operatorname{restrict}(f, u) \in (u \to c)) \wedge (v \subseteq u)) \Rightarrow (\operatorname{restrict}(f, v) \in (v \to c)) \end{array}}}
\vnbprop{diff-on-shrink-restrict}{\vnbstmt{\begin{array}{@{}l@{}} \forall k, u, f, a, l, v. \\ \quad (\operatorname{is\text{-}diff\text{-}on}(k, u, \operatorname{restrict}(f, u), a, l) \wedge (\operatorname{is\text{-}open}(\operatorname{nf\text{-}metric\text{-}space}(k), v) \wedge ((v \subseteq u) \wedge (a \in v)))) \Rightarrow  \\ \quad \quad \operatorname{is\text{-}diff\text{-}on}(k, v, \operatorname{restrict}(f, v), a, l) \end{array}}}
\vnbprop{diff-on-ooint-shrink}{\vnbstmt{\begin{array}{@{}l@{}} \forall a \in \mathbb{R}, e, w, f, l. \\ \quad \operatorname{pos\text{-}rr}(e) \wedge \\ \quad \operatorname{pos\text{-}rr}(w) \wedge \\ \quad w \leq e \wedge \\ \quad \operatorname{is\text{-}diff\text{-}on}(\operatorname{rr\text{-}normed\text{-}field}, \operatorname{ooint}((a - e), (a + e)), \operatorname{restrict}(f, \operatorname{ooint}((a - e), (a + e))), a, l) \Rightarrow \\ \quad \quad \operatorname{is\text{-}diff\text{-}on}(\operatorname{rr\text{-}normed\text{-}field}, \operatorname{ooint}((a - w), (a + w)), \operatorname{restrict}(f, \operatorname{ooint}((a - w), (a + w))), a, l) \end{array}}}
\vnbprop{has-deriv-at-intro}{\vnbstmt{\begin{array}{@{}l@{}} \forall f, a, l, e. \\ \quad \operatorname{pos\text{-}rr}(e) \wedge \\ \quad \operatorname{is\text{-}diff\text{-}on}(\operatorname{rr\text{-}normed\text{-}field}, \operatorname{ooint}((a - e), (a + e)), \operatorname{restrict}(f, \operatorname{ooint}((a - e), (a + e))), a, l) \Rightarrow \\ \quad \quad \operatorname{has\text{-}deriv\text{-}at}(f, a, l) \end{array}}}
\vnbprop{has-deriv-at-shrink}{\vnbstmt{\begin{array}{@{}l@{}} \forall f, a, l, r. \\ \quad (\operatorname{has\text{-}deriv\text{-}at}(f, a, l) \wedge \operatorname{pos\text{-}rr}(r)) \Rightarrow  \\ \quad \quad \exists w. \\ \quad \quad \quad \operatorname{pos\text{-}rr}(w) \wedge \\ \quad \quad \quad w \leq r \wedge \\ \quad \quad \quad \operatorname{is\text{-}diff\text{-}on}(\operatorname{rr\text{-}normed\text{-}field}, \operatorname{ooint}((a - w), (a + w)), \operatorname{restrict}(f, \operatorname{ooint}((a - w), (a + w))), a, l) \end{array}}}
\vnbprop{has-deriv-at-pt-in-rr}{\vnbstmt{\forall f, a, l.\; (\operatorname{has\text{-}deriv\text{-}at}(f, a, l) \Rightarrow (a \in \mathbb{R}))}}
\vnbprop{has-deriv-at-in-rr}{\vnbstmt{\forall f, a, l.\; (\operatorname{has\text{-}deriv\text{-}at}(f, a, l) \Rightarrow (l \in \mathbb{R}))}}
\vnbprop{has-deriv-at-unique}{\vnbstmt{\begin{array}{@{}l@{}} \forall f, a, l, m. \\ \quad (\operatorname{has\text{-}deriv\text{-}at}(f, a, l) \wedge \operatorname{has\text{-}deriv\text{-}at}(f, a, m)) \Rightarrow (l = m) \end{array}}}
\vnbprop{diff-on-nbhd-implies-diff-at}{\vnbstmt{\begin{array}{@{}l@{}} \forall f \in (\mathbb{R} \to \mathbb{R}), a \in \mathbb{R}, e, l. \\ \quad \operatorname{pos\text{-}rr}(e) \wedge \\ \quad \operatorname{is\text{-}diff\text{-}on}(\operatorname{rr\text{-}normed\text{-}field}, \operatorname{ooint}((a - e), (a + e)), \operatorname{restrict}(f, \operatorname{ooint}((a - e), (a + e))), a, l) \Rightarrow \\ \quad \quad \operatorname{is\text{-}diff\text{-}at}(f, a, l) \end{array}}}
\vnbprop{has-deriv-at-iff-diff-at}{\vnbstmt{\forall f \in (\mathbb{R} \to \mathbb{R}), a, l.\; (\operatorname{has\text{-}deriv\text{-}at}(f, a, l) \Leftrightarrow \operatorname{is\text{-}diff\text{-}at}(f, a, l))}}
\vnbprop{extend-const-value}{\vnbstmt{\forall f, a, b.; \forall z \in \mathbb{R}.\; (\operatorname{extend\text{-}const}(f, a, b)(z) \simeq f(\operatorname{clamp}(a, b, z)))}}
\vnbprop{extend-const-in-fun}{\vnbstmt{\begin{array}{@{}l@{}} \forall a, b \in \mathbb{R}, f \in (\operatorname{ccint}(a, b) \to \mathbb{R}). \\ \quad (a \leq b) \Rightarrow (\operatorname{extend\text{-}const}(f, a, b) \in (\mathbb{R} \to \mathbb{R})) \end{array}}}
\vnbprop{extend-const-fixes}{\vnbstmt{\begin{array}{@{}l@{}} \forall a, b, z \in \mathbb{R}, f. \\ \quad ((a \leq z) \wedge (z \leq b)) \Rightarrow (\operatorname{extend\text{-}const}(f, a, b)(z) \simeq f(z)) \end{array}}}
\vnbprop{extend-const-ccint-fixes}{\vnbstmt{\begin{array}{@{}l@{}} \forall a, b \in \mathbb{R}, z \in \operatorname{ccint}(a, b), f. \\ \quad (\operatorname{extend\text{-}const}(f, a, b)(z) \simeq f(z)) \end{array}}}
\vnbprop{extend-const-continuous-at}{\vnbstmt{\begin{array}{@{}l@{}} \forall a, b \in \mathbb{R}, f.; \forall t \in \mathbb{R}. \\ \quad ((a \leq b) \wedge \operatorname{is\text{-}continuous\text{-}on}(f, \operatorname{ccint}(a, b))) \Rightarrow  \\ \quad \quad \operatorname{is\text{-}continuous\text{-}at}(\operatorname{rr\text{-}ms}, \operatorname{rr\text{-}ms}, \operatorname{extend\text{-}const}(f, a, b), t) \end{array}}}
\vnbprop{extend-const-deriv-fwd}{\vnbstmt{\begin{array}{@{}l@{}} \forall a, b \in \mathbb{R}, f \in (\operatorname{ccint}(a, b) \to \mathbb{R}), x \in \operatorname{ooint}(a, b), l. \\ \quad ((a \leq b) \wedge \operatorname{has\text{-}deriv\text{-}at}(f, x, l)) \Rightarrow \operatorname{is\text{-}diff\text{-}at}(\operatorname{extend\text{-}const}(f, a, b), x, l) \end{array}}}
\vnbprop{extend-const-deriv-bwd}{\vnbstmt{\begin{array}{@{}l@{}} \forall a, b \in \mathbb{R}, f \in (\operatorname{ccint}(a, b) \to \mathbb{R}), x \in \operatorname{ooint}(a, b), l. \\ \quad ((a \leq b) \wedge \operatorname{is\text{-}diff\text{-}at}(\operatorname{extend\text{-}const}(f, a, b), x, l)) \Rightarrow \operatorname{has\text{-}deriv\text{-}at}(f, x, l) \end{array}}}
\vnbprop{extend-const-deriv-iff}{\vnbstmt{\begin{array}{@{}l@{}} \forall a, b \in \mathbb{R}, f \in (\operatorname{ccint}(a, b) \to \mathbb{R}), x \in \operatorname{ooint}(a, b), l. \\ \quad (a \leq b) \Rightarrow (\operatorname{has\text{-}deriv\text{-}at}(f, x, l) \Leftrightarrow \operatorname{is\text{-}diff\text{-}at}(\operatorname{extend\text{-}const}(f, a, b), x, l)) \end{array}}}
\vnbprop{has-deriv-at-local}{\vnbstmt{\begin{array}{@{}l@{}} \forall r, f, g, x, l. \\ \quad \operatorname{pos\text{-}rr}(r) \wedge \\ \quad \operatorname{restrict}(f, \operatorname{ooint}((x - r), (x + r))) \in (\operatorname{ooint}((x - r), (x + r)) \to \mathbb{R}) \wedge \\ \quad \forall \mathit{hbx} \in \operatorname{ooint}((x - r), (x + r)).\; (f(\mathit{hbx}) \simeq g(\mathit{hbx})) \wedge \\ \quad \operatorname{has\text{-}deriv\text{-}at}(g, x, l) \Rightarrow \\ \quad \quad \operatorname{has\text{-}deriv\text{-}at}(f, x, l) \end{array}}}
\vnbprop{has-deriv-at-sum}{\vnbstmt{\begin{array}{@{}l@{}} \forall r, f, g, h, x, l, m. \\ \quad \operatorname{pos\text{-}rr}(r) \wedge \\ \quad \operatorname{restrict}(h, \operatorname{ooint}((x - r), (x + r))) \in (\operatorname{ooint}((x - r), (x + r)) \to \mathbb{R}) \wedge \\ \quad \forall \mathit{hbx} \in \operatorname{ooint}((x - r), (x + r)).\; (h(\mathit{hbx}) \simeq (f(\mathit{hbx}) + g(\mathit{hbx}))) \wedge \\ \quad \operatorname{has\text{-}deriv\text{-}at}(f, x, l) \wedge \\ \quad \operatorname{has\text{-}deriv\text{-}at}(g, x, m) \Rightarrow \\ \quad \quad \operatorname{has\text{-}deriv\text{-}at}(h, x, (l + m)) \end{array}}}
\vnbprop{has-deriv-at-product}{\vnbstmt{\begin{array}{@{}l@{}} \forall r, f, g, h, x, l, m. \\ \quad \operatorname{pos\text{-}rr}(r) \wedge \\ \quad \operatorname{restrict}(h, \operatorname{ooint}((x - r), (x + r))) \in (\operatorname{ooint}((x - r), (x + r)) \to \mathbb{R}) \wedge \\ \quad \forall \mathit{hbx} \in \operatorname{ooint}((x - r), (x + r)).\; (h(\mathit{hbx}) \simeq (f(\mathit{hbx}) \cdot g(\mathit{hbx}))) \wedge \\ \quad \operatorname{has\text{-}deriv\text{-}at}(f, x, l) \wedge \\ \quad \operatorname{has\text{-}deriv\text{-}at}(g, x, m) \Rightarrow \\ \quad \quad \operatorname{has\text{-}deriv\text{-}at}(h, x, ((l \cdot g(x)) + (m \cdot f(x)))) \end{array}}}

\vnbbookfile{theorem-library/interval-mvt}{theorem-library/interval-mvt}
\noindent{\small interval-mvt.scm -- ROLLE, THE TWO MEAN VALUE THEOREMS AND THE THREE BOUNDS, STATED FOR A FUNCTION ON ITS INTERVAL.}\par
\vnbprop{rolle-on-interval}{\vnbstmt{\begin{array}{@{}l@{}} \forall a, b \in \mathbb{R}, h \in (\operatorname{ccint}(a, b) \to \mathbb{R}). \\ \quad a < b \wedge \\ \quad \operatorname{is\text{-}continuous\text{-}on}(h, \operatorname{ccint}(a, b)) \wedge \\ \quad \forall x \in \operatorname{ooint}(a, b).\; \exists l.\; \operatorname{has\text{-}deriv\text{-}at}(h, x, l) \wedge \\ \quad h(a) = h(b) \Rightarrow \\ \quad \quad \exists c \in \operatorname{ooint}(a, b).\; \operatorname{has\text{-}deriv\text{-}at}(h, c, 0) \end{array}}}
\vnbprop{generalized-mvt-on-interval}{\vnbstmt{\begin{array}{@{}l@{}} \forall a, b \in \mathbb{R}, f, g. \\ \quad a < b \wedge \\ \quad f \in (\operatorname{ccint}(a, b) \to \mathbb{R}) \wedge \\ \quad g \in (\operatorname{ccint}(a, b) \to \mathbb{R}) \wedge \\ \quad \operatorname{is\text{-}continuous\text{-}on}(f, \operatorname{ccint}(a, b)) \wedge \\ \quad \operatorname{is\text{-}continuous\text{-}on}(g, \operatorname{ccint}(a, b)) \wedge \\ \quad \forall x \in \operatorname{ooint}(a, b).\; \exists l.\; \operatorname{has\text{-}deriv\text{-}at}(f, x, l) \wedge \\ \quad \forall x \in \operatorname{ooint}(a, b).\; \exists m.\; \operatorname{has\text{-}deriv\text{-}at}(g, x, m) \Rightarrow \\ \quad \quad \exists c \in \operatorname{ooint}(a, b), l, m. \\ \quad \quad \quad (\operatorname{has\text{-}deriv\text{-}at}(f, c, l) \wedge (\operatorname{has\text{-}deriv\text{-}at}(g, c, m) \wedge ((l \cdot (g(b) - g(a))) = (m \cdot (f(b) - f(a)))))) \end{array}}}
\vnbprop{mvt-on-interval}{\vnbstmt{\begin{array}{@{}l@{}} \forall a, b \in \mathbb{R}, f \in (\operatorname{ccint}(a, b) \to \mathbb{R}). \\ \quad a < b \wedge \\ \quad \operatorname{is\text{-}continuous\text{-}on}(f, \operatorname{ccint}(a, b)) \wedge \\ \quad \forall x \in \operatorname{ooint}(a, b).\; \exists l.\; \operatorname{has\text{-}deriv\text{-}at}(f, x, l) \Rightarrow \\ \quad \quad \exists c \in \operatorname{ooint}(a, b), l. \\ \quad \quad \quad \operatorname{has\text{-}deriv\text{-}at}(f, c, l) \wedge (l = ((f(b) - f(a)) \cdot \operatorname{recip}((b - a)))) \end{array}}}
\vnbprop{mvt-upper-bound-on-interval}{\vnbstmt{\begin{array}{@{}l@{}} \forall a, b, m \in \mathbb{R}, f \in (\operatorname{ccint}(a, b) \to \mathbb{R}). \\ \quad a < b \wedge \\ \quad \operatorname{is\text{-}continuous\text{-}on}(f, \operatorname{ccint}(a, b)) \wedge \\ \quad \forall x \in \operatorname{ooint}(a, b).\; \exists l.\; (\operatorname{has\text{-}deriv\text{-}at}(f, x, l) \wedge (l \leq m)) \Rightarrow \\ \quad \quad (f(b) - f(a)) \leq (m \cdot (b - a)) \end{array}}}
\vnbprop{mvt-lower-bound-on-interval}{\vnbstmt{\begin{array}{@{}l@{}} \forall a, b, m \in \mathbb{R}, f \in (\operatorname{ccint}(a, b) \to \mathbb{R}). \\ \quad a < b \wedge \\ \quad \operatorname{is\text{-}continuous\text{-}on}(f, \operatorname{ccint}(a, b)) \wedge \\ \quad \forall x \in \operatorname{ooint}(a, b).\; \exists l.\; (\operatorname{has\text{-}deriv\text{-}at}(f, x, l) \wedge (m \leq l)) \Rightarrow \\ \quad \quad (m \cdot (b - a)) \leq (f(b) - f(a)) \end{array}}}
\vnbprop{mvt-abs-bound-on-interval}{\vnbstmt{\begin{array}{@{}l@{}} \forall a, b, m \in \mathbb{R}, f \in (\operatorname{ccint}(a, b) \to \mathbb{R}). \\ \quad a < b \wedge \\ \quad \operatorname{is\text{-}continuous\text{-}on}(f, \operatorname{ccint}(a, b)) \wedge \\ \quad \forall x \in \operatorname{ooint}(a, b).\; \exists l.\; (\operatorname{has\text{-}deriv\text{-}at}(f, x, l) \wedge (\lvert l \rvert \leq m)) \Rightarrow \\ \quad \quad \lvert (f(b) - f(a)) \rvert \leq (m \cdot (b - a)) \end{array}}}

\vnbbookfile{structure-library/interval-extremum}{structure-library/interval-extremum}
\noindent{\small interval-extremum.scm -- LOCAL MAXIMUM AND LOCAL MINIMUM OF A FUNCTION ON A SET.  DEFINITIONS ONLY; Proposition 2.10 is proven in theorem-library/interval-extremum.scm.}\par
\par\medskip\noindent\textbf{Definition} (\texttt{is-local-max-at}).\;For all $f$, $d$ and $\operatorname{lmt\_}$, 
\begin{equation*}
\begin{array}{@{}l@{}} \operatorname{is\text{-}local\text{-}max\text{-}at}(f, d, \operatorname{lmt\_}) \\ \Leftrightarrow\; f \in (d \to \mathbb{R}) \wedge \\ \operatorname{lmt\_} \in d \wedge \\ \exists \operatorname{lma\_}. \\ \quad (\operatorname{pos\text{-}rr}(\operatorname{lma\_}) \wedge \forall \operatorname{lmx\_} \in d.\; ((\lvert -(\operatorname{lmx\_}, \operatorname{lmt\_}) \rvert \leq \operatorname{lma\_}) \Rightarrow (f(\operatorname{lmx\_}) \leq f(\operatorname{lmt\_})))) \end{array}
\end{equation*}

\par\medskip\noindent\textbf{Definition} (\texttt{is-local-min-at}).\;For all $f$, $d$ and $\operatorname{lmt\_}$, 
\begin{equation*}
\begin{array}{@{}l@{}} \operatorname{is\text{-}local\text{-}min\text{-}at}(f, d, \operatorname{lmt\_}) \\ \Leftrightarrow\; f \in (d \to \mathbb{R}) \wedge \\ \operatorname{lmt\_} \in d \wedge \\ \exists \operatorname{lma\_}. \\ \quad (\operatorname{pos\text{-}rr}(\operatorname{lma\_}) \wedge \forall \operatorname{lmx\_} \in d.\; ((\lvert -(\operatorname{lmx\_}, \operatorname{lmt\_}) \rvert \leq \operatorname{lma\_}) \Rightarrow (f(\operatorname{lmt\_}) \leq f(\operatorname{lmx\_})))) \end{array}
\end{equation*}


\vnbbookfile{theorem-library/interval-extremum}{theorem-library/interval-extremum}
\noindent{\small interval-extremum.scm -- PROPOSITION 2.10 FOR A FUNCTION ON ITS INTERVAL.}\par
\vnbprop{interval-local-max-deriv-zero}{\vnbstmt{\begin{array}{@{}l@{}} \forall a, b \in \mathbb{R}, f.; \forall c \in \operatorname{ooint}(a, b), l. \\ \quad (\operatorname{is\text{-}local\text{-}max\text{-}at}(f, \operatorname{ccint}(a, b), c) \wedge \operatorname{has\text{-}deriv\text{-}at}(f, c, l)) \Rightarrow (l = 0) \end{array}}}
\vnbprop{interval-local-min-deriv-zero}{\vnbstmt{\begin{array}{@{}l@{}} \forall a, b \in \mathbb{R}, f.; \forall c \in \operatorname{ooint}(a, b), l. \\ \quad (\operatorname{is\text{-}local\text{-}min\text{-}at}(f, \operatorname{ccint}(a, b), c) \wedge \operatorname{has\text{-}deriv\text{-}at}(f, c, l)) \Rightarrow (l = 0) \end{array}}}

\vnbbookfile{structure-library/one-sided-derivative}{structure-library/one-sided-derivative}
\noindent{\small one-sided-derivative.scm -- THE ONE-SIDED DERIVATIVES OF THE USER'S NOTES, section 2.1.  DEFINITIONS ONLY; every law is PROVEN in theorem-library/one-sided-derivative-laws.scm.}\par
\par\medskip\noindent\textbf{Definition} (\texttt{coint}).\;$\operatorname{coint}(a, b) :=$ $\{\operatorname{osq\_} \in \mathbb{R} : ((a \leq \operatorname{osq\_}) \wedge (\operatorname{osq\_} < b))\}$
\par\medskip\noindent\textbf{Definition} (\texttt{ocint}).\;$\operatorname{ocint}(a, b) :=$ $\{\operatorname{osq\_} \in \mathbb{R} : ((a < \operatorname{osq\_}) \wedge (\operatorname{osq\_} \leq b))\}$
\par\medskip\noindent\textbf{Definition} (\texttt{has-deriv-at-within}).\;For all $f$, $w$, $x$ and $l$, 
\begin{equation*}
\begin{array}{@{}l@{}} \operatorname{has\text{-}deriv\text{-}at\text{-}within}(f, w, x, l) \\ \Leftrightarrow\; \operatorname{restrict}(f, w) \in (w \to \mathbb{R}) \wedge \\ x \in w \wedge \\ l \in \mathbb{R} \wedge \\ \exists \operatorname{osp\_} \in (w \to \mathbb{R}). \\ \quad \operatorname{is\text{-}continuous\text{-}at}(\operatorname{subspace\text{-}ms}(\operatorname{rr\text{-}ms}, w), \operatorname{rr\text{-}ms}, \operatorname{osp\_}, x) \wedge \\ \quad \operatorname{osp\_}(x) = l \wedge \\ \quad \forall \operatorname{ost\_} \in w.\; (-(f(\operatorname{ost\_}), f(x)) = \cdot(\operatorname{osp\_}(\operatorname{ost\_}), -(\operatorname{ost\_}, x))) \end{array}
\end{equation*}

\par\medskip\noindent\textbf{Definition} (\texttt{has-right-deriv-at}).\;For all $f$, $x$ and $l$, 
\begin{equation*}
\begin{array}{@{}l@{}} \operatorname{has\text{-}right\text{-}deriv\text{-}at}(f, x, l) \\ \Leftrightarrow\; \exists \operatorname{ose\_}. \\ \quad \operatorname{pos\text{-}rr}(\operatorname{ose\_}) \wedge \operatorname{has\text{-}deriv\text{-}at\text{-}within}(f, \operatorname{coint}(x, +(x, \operatorname{ose\_})), x, l) \end{array}
\end{equation*}

\par\medskip\noindent\textbf{Definition} (\texttt{has-left-deriv-at}).\;For all $f$, $x$ and $l$, 
\begin{equation*}
\begin{array}{@{}l@{}} \operatorname{has\text{-}left\text{-}deriv\text{-}at}(f, x, l) \\ \Leftrightarrow\; \exists \operatorname{ose\_}. \\ \quad \operatorname{pos\text{-}rr}(\operatorname{ose\_}) \wedge \operatorname{has\text{-}deriv\text{-}at\text{-}within}(f, \operatorname{ocint}(-(x, \operatorname{ose\_}), x), x, l) \end{array}
\end{equation*}


\vnbbookfile{theorem-library/one-sided-derivative-laws}{theorem-library/one-sided-derivative-laws}
\noindent{\small one-sided-derivative-laws.scm -- the laws of structure-library/one-sided-derivative.scm: COINT, OCINT, HAS-DERIV-AT-WITHIN, HAS-RIGHT-DERIV-AT and HAS-LEFT-DERIV-AT.  Everything here is PROVEN; the definition file asserts nothing.}\par
\vnbprop{coint-membership}{\vnbstmt{\forall a, b, y.\; ((y \in \operatorname{coint}(a, b)) \Leftrightarrow ((y \in \mathbb{R}) \wedge ((a \leq y) \wedge (y < b))))}}
\vnbprop{ocint-membership}{\vnbstmt{\forall a, b, y.\; ((y \in \operatorname{ocint}(a, b)) \Leftrightarrow ((y \in \mathbb{R}) \wedge ((a < y) \wedge (y \leq b))))}}
\vnbprop{coint-elt-in-rr}{\vnbstmt{\forall a, b.; \forall y \in \operatorname{coint}(a, b).\; (y \in \mathbb{R})}}
\vnbprop{ocint-elt-in-rr}{\vnbstmt{\forall a, b.; \forall y \in \operatorname{ocint}(a, b).\; (y \in \mathbb{R})}}
\vnbprop{coint-subset-rr}{\vnbstmt{\forall a, b.\; (\operatorname{coint}(a, b) \subseteq \mathbb{R})}}
\vnbprop{ocint-subset-rr}{\vnbstmt{\forall a, b.\; (\operatorname{ocint}(a, b) \subseteq \mathbb{R})}}
\vnbprop{coint-left-endpoint}{\vnbstmt{\forall x, r.\; (((x \in \mathbb{R}) \wedge \operatorname{pos\text{-}rr}(r)) \Rightarrow (x \in \operatorname{coint}(x, (x + r))))}}
\vnbprop{ocint-right-endpoint}{\vnbstmt{\forall x, r.\; (((x \in \mathbb{R}) \wedge \operatorname{pos\text{-}rr}(r)) \Rightarrow (x \in \operatorname{ocint}((x - r), x)))}}
\vnbprop{coint-subset-ooint}{\vnbstmt{\begin{array}{@{}l@{}} \forall x, r. \\ \quad ((x \in \mathbb{R}) \wedge \operatorname{pos\text{-}rr}(r)) \Rightarrow (\operatorname{coint}(x, (x + r)) \subseteq \operatorname{ooint}((x - r), (x + r))) \end{array}}}
\vnbprop{ocint-subset-ooint}{\vnbstmt{\begin{array}{@{}l@{}} \forall x, r. \\ \quad ((x \in \mathbb{R}) \wedge \operatorname{pos\text{-}rr}(r)) \Rightarrow (\operatorname{ocint}((x - r), x) \subseteq \operatorname{ooint}((x - r), (x + r))) \end{array}}}
\vnbprop{coint-shrink}{\vnbstmt{\begin{array}{@{}l@{}} \forall x, r, v. \\ \quad ((x \in \mathbb{R}) \wedge (\operatorname{pos\text{-}rr}(r) \wedge (\operatorname{pos\text{-}rr}(v) \wedge (v \leq r)))) \Rightarrow  \\ \quad \quad \operatorname{coint}(x, (x + v)) \subseteq \operatorname{coint}(x, (x + r)) \end{array}}}
\vnbprop{ocint-shrink}{\vnbstmt{\begin{array}{@{}l@{}} \forall x, r, v. \\ \quad ((x \in \mathbb{R}) \wedge (\operatorname{pos\text{-}rr}(r) \wedge (\operatorname{pos\text{-}rr}(v) \wedge (v \leq r)))) \Rightarrow  \\ \quad \quad \operatorname{ocint}((x - v), x) \subseteq \operatorname{ocint}((x - r), x) \end{array}}}
\vnbprop{coint-left-limit-point}{\vnbstmt{\begin{array}{@{}l@{}} \forall x, d, r. \\ \quad ((x \in \mathbb{R}) \wedge (\operatorname{pos\text{-}rr}(d) \wedge \operatorname{pos\text{-}rr}(r))) \Rightarrow  \\ \quad \quad \exists b \in \operatorname{coint}(x, (x + d)).\; (\neg (b = x) \wedge (\lvert (x - b) \rvert \leq r)) \end{array}}}
\vnbprop{ocint-right-limit-point}{\vnbstmt{\begin{array}{@{}l@{}} \forall x, d, r. \\ \quad ((x \in \mathbb{R}) \wedge (\operatorname{pos\text{-}rr}(d) \wedge \operatorname{pos\text{-}rr}(r))) \Rightarrow  \\ \quad \quad \exists b \in \operatorname{ocint}((x - d), x).\; (\neg (b = x) \wedge (\lvert (x - b) \rvert \leq r)) \end{array}}}
\vnbprop{restrict-fun-domain-subset}{\vnbstmt{\begin{array}{@{}l@{}} \forall f, d, c, w, b. \\ \quad ((f \in (d \to c)) \wedge (\operatorname{restrict}(f, w) \in (w \to b))) \Rightarrow (w \subseteq d) \end{array}}}
\vnbprop{deriv-within-in-rr}{\vnbstmt{\forall f, w, x, l.\; (\operatorname{has\text{-}deriv\text{-}at\text{-}within}(f, w, x, l) \Rightarrow (l \in \mathbb{R}))}}
\vnbprop{deriv-within-pt-in}{\vnbstmt{\forall f, w, x, l.\; (\operatorname{has\text{-}deriv\text{-}at\text{-}within}(f, w, x, l) \Rightarrow (x \in w))}}
\vnbprop{deriv-within-restrict-in-fun}{\vnbstmt{\begin{array}{@{}l@{}} \forall f, w, x, l. \\ \quad \operatorname{has\text{-}deriv\text{-}at\text{-}within}(f, w, x, l) \Rightarrow (\operatorname{restrict}(f, w) \in (w \to \mathbb{R})) \end{array}}}
\vnbprop{deriv-within-shrink}{\vnbstmt{\begin{array}{@{}l@{}} \forall f, w, v, x, l. \\ \quad ((w \subseteq \mathbb{R}) \wedge (\operatorname{has\text{-}deriv\text{-}at\text{-}within}(f, w, x, l) \wedge ((v \subseteq w) \wedge (x \in v)))) \Rightarrow  \\ \quad \quad \operatorname{has\text{-}deriv\text{-}at\text{-}within}(f, v, x, l) \end{array}}}
\vnbprop{deriv-within-local}{\vnbstmt{\begin{array}{@{}l@{}} \forall f, g, w, x, l. \\ \quad ((\operatorname{restrict}(f, w) \in (w \to \mathbb{R})) \wedge (\forall y \in w.\; (f(y) \simeq g(y)) \wedge \operatorname{has\text{-}deriv\text{-}at\text{-}within}(g, w, x, l))) \Rightarrow  \\ \quad \quad \operatorname{has\text{-}deriv\text{-}at\text{-}within}(f, w, x, l) \end{array}}}
\vnbprop{deriv-within-unique}{\vnbstmt{\begin{array}{@{}l@{}} \forall f, w, x, l, m. \\ \quad w \subseteq \mathbb{R} \wedge \\ \quad \forall r.\; (\operatorname{pos\text{-}rr}(r) \Rightarrow \exists b \in w.\; (\neg (b = x) \wedge (\lvert (x - b) \rvert \leq r))) \wedge \\ \quad \operatorname{has\text{-}deriv\text{-}at\text{-}within}(f, w, x, l) \wedge \\ \quad \operatorname{has\text{-}deriv\text{-}at\text{-}within}(f, w, x, m) \Rightarrow \\ \quad \quad l = m \end{array}}}
\vnbprop{rr-mul-neg-neg-pos}{\vnbstmt{\forall y, b.\; (((y \in \mathbb{R}) \wedge ((b \in \mathbb{R}) \wedge ((y < 0) \wedge (b < 0)))) \Rightarrow (0 < (y \cdot b)))}}
\vnbprop{has-deriv-at-of-within-ooint}{\vnbstmt{\begin{array}{@{}l@{}} \forall x, r, f, l. \\ \quad (\operatorname{pos\text{-}rr}(r) \wedge \operatorname{has\text{-}deriv\text{-}at\text{-}within}(f, \operatorname{ooint}((x - r), (x + r)), x, l)) \Rightarrow  \\ \quad \quad \operatorname{has\text{-}deriv\text{-}at}(f, x, l) \end{array}}}
\vnbprop{has-deriv-at-to-within-ooint}{\vnbstmt{\begin{array}{@{}l@{}} \forall f, x, l, r. \\ \quad (\operatorname{has\text{-}deriv\text{-}at}(f, x, l) \wedge \operatorname{pos\text{-}rr}(r)) \Rightarrow  \\ \quad \quad \exists v. \\ \quad \quad \quad \operatorname{pos\text{-}rr}(v) \wedge ((v \leq r) \wedge \operatorname{has\text{-}deriv\text{-}at\text{-}within}(f, \operatorname{ooint}((x - v), (x + v)), x, l)) \end{array}}}
\vnbprop{has-deriv-at-iff-within-ooint}{\vnbstmt{\begin{array}{@{}l@{}} \forall f, x, l. \\ \quad \operatorname{has\text{-}deriv\text{-}at}(f, x, l) \\ \quad \Leftrightarrow\; \exists v.\; (\operatorname{pos\text{-}rr}(v) \wedge \operatorname{has\text{-}deriv\text{-}at\text{-}within}(f, \operatorname{ooint}((x - v), (x + v)), x, l)) \end{array}}}
\vnbprop{has-right-deriv-at-intro}{\vnbstmt{\begin{array}{@{}l@{}} \forall f, x, l, r. \\ \quad (\operatorname{pos\text{-}rr}(r) \wedge \operatorname{has\text{-}deriv\text{-}at\text{-}within}(f, \operatorname{coint}(x, (x + r)), x, l)) \Rightarrow  \\ \quad \quad \operatorname{has\text{-}right\text{-}deriv\text{-}at}(f, x, l) \end{array}}}
\vnbprop{has-left-deriv-at-intro}{\vnbstmt{\begin{array}{@{}l@{}} \forall f, x, l, r. \\ \quad (\operatorname{pos\text{-}rr}(r) \wedge \operatorname{has\text{-}deriv\text{-}at\text{-}within}(f, \operatorname{ocint}((x - r), x), x, l)) \Rightarrow  \\ \quad \quad \operatorname{has\text{-}left\text{-}deriv\text{-}at}(f, x, l) \end{array}}}
\vnbprop{has-right-deriv-at-pt-in-rr}{\vnbstmt{\forall f, x, l.\; (\operatorname{has\text{-}right\text{-}deriv\text{-}at}(f, x, l) \Rightarrow (x \in \mathbb{R}))}}
\vnbprop{has-right-deriv-at-in-rr}{\vnbstmt{\forall f, x, l.\; (\operatorname{has\text{-}right\text{-}deriv\text{-}at}(f, x, l) \Rightarrow (l \in \mathbb{R}))}}
\vnbprop{has-left-deriv-at-pt-in-rr}{\vnbstmt{\forall f, x, l.\; (\operatorname{has\text{-}left\text{-}deriv\text{-}at}(f, x, l) \Rightarrow (x \in \mathbb{R}))}}
\vnbprop{has-left-deriv-at-in-rr}{\vnbstmt{\forall f, x, l.\; (\operatorname{has\text{-}left\text{-}deriv\text{-}at}(f, x, l) \Rightarrow (l \in \mathbb{R}))}}
\vnbprop{has-right-deriv-at-shrink}{\vnbstmt{\begin{array}{@{}l@{}} \forall f, x, l, r. \\ \quad (\operatorname{has\text{-}right\text{-}deriv\text{-}at}(f, x, l) \wedge \operatorname{pos\text{-}rr}(r)) \Rightarrow  \\ \quad \quad \exists v. \\ \quad \quad \quad \operatorname{pos\text{-}rr}(v) \wedge ((v \leq r) \wedge \operatorname{has\text{-}deriv\text{-}at\text{-}within}(f, \operatorname{coint}(x, (x + v)), x, l)) \end{array}}}
\vnbprop{has-left-deriv-at-shrink}{\vnbstmt{\begin{array}{@{}l@{}} \forall f, x, l, r. \\ \quad (\operatorname{has\text{-}left\text{-}deriv\text{-}at}(f, x, l) \wedge \operatorname{pos\text{-}rr}(r)) \Rightarrow  \\ \quad \quad \exists v. \\ \quad \quad \quad \operatorname{pos\text{-}rr}(v) \wedge ((v \leq r) \wedge \operatorname{has\text{-}deriv\text{-}at\text{-}within}(f, \operatorname{ocint}((x - v), x), x, l)) \end{array}}}
\vnbprop{has-right-deriv-at-unique}{\vnbstmt{\begin{array}{@{}l@{}} \forall f, x, l, m. \\ \quad (\operatorname{has\text{-}right\text{-}deriv\text{-}at}(f, x, l) \wedge \operatorname{has\text{-}right\text{-}deriv\text{-}at}(f, x, m)) \Rightarrow (l = m) \end{array}}}
\vnbprop{has-left-deriv-at-unique}{\vnbstmt{\begin{array}{@{}l@{}} \forall f, x, l, m. \\ \quad (\operatorname{has\text{-}left\text{-}deriv\text{-}at}(f, x, l) \wedge \operatorname{has\text{-}left\text{-}deriv\text{-}at}(f, x, m)) \Rightarrow (l = m) \end{array}}}
\vnbprop{has-right-deriv-at-local}{\vnbstmt{\begin{array}{@{}l@{}} \forall r, f, g, x, l. \\ \quad \operatorname{pos\text{-}rr}(r) \wedge \\ \quad \operatorname{restrict}(f, \operatorname{coint}(x, (x + r))) \in (\operatorname{coint}(x, (x + r)) \to \mathbb{R}) \wedge \\ \quad \forall y \in \operatorname{coint}(x, (x + r)).\; (f(y) \simeq g(y)) \wedge \\ \quad \operatorname{has\text{-}right\text{-}deriv\text{-}at}(g, x, l) \Rightarrow \\ \quad \quad \operatorname{has\text{-}right\text{-}deriv\text{-}at}(f, x, l) \end{array}}}
\vnbprop{has-left-deriv-at-local}{\vnbstmt{\begin{array}{@{}l@{}} \forall r, f, g, x, l. \\ \quad \operatorname{pos\text{-}rr}(r) \wedge \\ \quad \operatorname{restrict}(f, \operatorname{ocint}((x - r), x)) \in (\operatorname{ocint}((x - r), x) \to \mathbb{R}) \wedge \\ \quad \forall y \in \operatorname{ocint}((x - r), x).\; (f(y) \simeq g(y)) \wedge \\ \quad \operatorname{has\text{-}left\text{-}deriv\text{-}at}(g, x, l) \Rightarrow \\ \quad \quad \operatorname{has\text{-}left\text{-}deriv\text{-}at}(f, x, l) \end{array}}}
\vnbprop{has-deriv-at-implies-right}{\vnbstmt{\forall f, x, l.\; (\operatorname{has\text{-}deriv\text{-}at}(f, x, l) \Rightarrow \operatorname{has\text{-}right\text{-}deriv\text{-}at}(f, x, l))}}
\vnbprop{has-deriv-at-implies-left}{\vnbstmt{\forall f, x, l.\; (\operatorname{has\text{-}deriv\text{-}at}(f, x, l) \Rightarrow \operatorname{has\text{-}left\text{-}deriv\text{-}at}(f, x, l))}}
\vnbprop{local-max-right-deriv-nonpos}{\vnbstmt{\begin{array}{@{}l@{}} \forall f, d, x, l. \\ \quad (\operatorname{is\text{-}local\text{-}max\text{-}at}(f, d, x) \wedge \operatorname{has\text{-}right\text{-}deriv\text{-}at}(f, x, l)) \Rightarrow (l \leq 0) \end{array}}}
\vnbprop{local-max-left-deriv-nonneg}{\vnbstmt{\begin{array}{@{}l@{}} \forall f, d, x, l. \\ \quad (\operatorname{is\text{-}local\text{-}max\text{-}at}(f, d, x) \wedge \operatorname{has\text{-}left\text{-}deriv\text{-}at}(f, x, l)) \Rightarrow (0 \leq l) \end{array}}}
\vnbprop{interval-local-max-right-deriv-nonpos}{\vnbstmt{\begin{array}{@{}l@{}} \forall a, b \in \mathbb{R}, f, x, l. \\ \quad (\operatorname{is\text{-}local\text{-}max\text{-}at}(f, \operatorname{ccint}(a, b), x) \wedge \operatorname{has\text{-}right\text{-}deriv\text{-}at}(f, x, l)) \Rightarrow  \\ \quad \quad l \leq 0 \end{array}}}
\vnbprop{interval-local-max-left-deriv-nonneg}{\vnbstmt{\begin{array}{@{}l@{}} \forall a, b \in \mathbb{R}, f, x, l. \\ \quad (\operatorname{is\text{-}local\text{-}max\text{-}at}(f, \operatorname{ccint}(a, b), x) \wedge \operatorname{has\text{-}left\text{-}deriv\text{-}at}(f, x, l)) \Rightarrow (0 \leq l) \end{array}}}
\vnbprop{rr-max-left-le}{\vnbstmt{\forall y, b.\; (((y \in \mathbb{R}) \wedge ((b \in \mathbb{R}) \wedge (b \leq y))) \Rightarrow (\operatorname{max}(y, b) = y))}}
\vnbprop{rr-min-left-le}{\vnbstmt{\forall y, b.\; (((y \in \mathbb{R}) \wedge ((b \in \mathbb{R}) \wedge (y \leq b))) \Rightarrow (\operatorname{min}(y, b) = y))}}
\vnbprop{rr-max-right-le}{\vnbstmt{\forall y, b.\; (((y \in \mathbb{R}) \wedge ((b \in \mathbb{R}) \wedge (y \leq b))) \Rightarrow (\operatorname{max}(y, b) = b))}}
\vnbprop{rr-min-right-le}{\vnbstmt{\forall y, b.\; (((y \in \mathbb{R}) \wedge ((b \in \mathbb{R}) \wedge (b \leq y))) \Rightarrow (\operatorname{min}(y, b) = b))}}
\vnbprop{rr-max-lt}{\vnbstmt{\begin{array}{@{}l@{}} \forall y, b, c. \\ \quad ((y \in \mathbb{R}) \wedge ((b \in \mathbb{R}) \wedge ((c \in \mathbb{R}) \wedge ((y < c) \wedge (b < c))))) \Rightarrow (\operatorname{max}(y, b) < c) \end{array}}}
\vnbprop{rr-lt-min}{\vnbstmt{\begin{array}{@{}l@{}} \forall c, y, b. \\ \quad ((c \in \mathbb{R}) \wedge ((y \in \mathbb{R}) \wedge ((b \in \mathbb{R}) \wedge ((c < y) \wedge (c < b))))) \Rightarrow (c < \operatorname{min}(y, b)) \end{array}}}
\vnbprop{rr-max-retract-close}{\vnbstmt{\forall y, c.\; (((y \in \mathbb{R}) \wedge (c \in \mathbb{R})) \Rightarrow (\lvert (c - \operatorname{max}(y, c)) \rvert \leq \lvert (c - y) \rvert))}}
\vnbprop{rr-min-retract-close}{\vnbstmt{\forall y, c.\; (((y \in \mathbb{R}) \wedge (c \in \mathbb{R})) \Rightarrow (\lvert (c - \operatorname{min}(y, c)) \rvert \leq \lvert (c - y) \rvert))}}
\vnbprop{coint-of-ooint}{\vnbstmt{\begin{array}{@{}l@{}} \forall x, r.; \forall y \in \operatorname{ooint}((x - r), (x + r)). \\ \quad ((x \in \mathbb{R}) \wedge (\operatorname{pos\text{-}rr}(r) \wedge (x \leq y))) \Rightarrow (y \in \operatorname{coint}(x, (x + r))) \end{array}}}
\vnbprop{ocint-of-ooint}{\vnbstmt{\begin{array}{@{}l@{}} \forall x, r.; \forall y \in \operatorname{ooint}((x - r), (x + r)). \\ \quad ((x \in \mathbb{R}) \wedge (\operatorname{pos\text{-}rr}(r) \wedge (y \leq x))) \Rightarrow (y \in \operatorname{ocint}((x - r), x)) \end{array}}}
\vnbprop{has-deriv-at-of-one-sided}{\vnbstmt{\begin{array}{@{}l@{}} \forall f, x, l. \\ \quad (\operatorname{has\text{-}right\text{-}deriv\text{-}at}(f, x, l) \wedge \operatorname{has\text{-}left\text{-}deriv\text{-}at}(f, x, l)) \Rightarrow  \\ \quad \quad \operatorname{has\text{-}deriv\text{-}at}(f, x, l) \end{array}}}
\vnbprop{has-deriv-at-iff-one-sided}{\vnbstmt{\begin{array}{@{}l@{}} \forall f, x, l. \\ \quad \operatorname{has\text{-}deriv\text{-}at}(f, x, l) \Leftrightarrow (\operatorname{has\text{-}right\text{-}deriv\text{-}at}(f, x, l) \wedge \operatorname{has\text{-}left\text{-}deriv\text{-}at}(f, x, l)) \end{array}}}
\vnbprop{deriv-within-sum}{\vnbstmt{\begin{array}{@{}l@{}} \forall f, g, h, w, x, l, m. \\ \quad w \subseteq \mathbb{R} \wedge \\ \quad \operatorname{restrict}(h, w) \in (w \to \mathbb{R}) \wedge \\ \quad \forall y \in w.\; (h(y) \simeq (f(y) + g(y))) \wedge \\ \quad \operatorname{has\text{-}deriv\text{-}at\text{-}within}(f, w, x, l) \wedge \\ \quad \operatorname{has\text{-}deriv\text{-}at\text{-}within}(g, w, x, m) \Rightarrow \\ \quad \quad \operatorname{has\text{-}deriv\text{-}at\text{-}within}(h, w, x, (l + m)) \end{array}}}
\vnbprop{has-right-deriv-at-sum}{\vnbstmt{\begin{array}{@{}l@{}} \forall r, f, g, h, x, l, m. \\ \quad \operatorname{pos\text{-}rr}(r) \wedge \\ \quad \operatorname{restrict}(h, \operatorname{coint}(x, (x + r))) \in (\operatorname{coint}(x, (x + r)) \to \mathbb{R}) \wedge \\ \quad \forall y \in \operatorname{coint}(x, (x + r)).\; (h(y) \simeq (f(y) + g(y))) \wedge \\ \quad \operatorname{has\text{-}right\text{-}deriv\text{-}at}(f, x, l) \wedge \\ \quad \operatorname{has\text{-}right\text{-}deriv\text{-}at}(g, x, m) \Rightarrow \\ \quad \quad \operatorname{has\text{-}right\text{-}deriv\text{-}at}(h, x, (l + m)) \end{array}}}
\vnbprop{has-left-deriv-at-sum}{\vnbstmt{\begin{array}{@{}l@{}} \forall r, f, g, h, x, l, m. \\ \quad \operatorname{pos\text{-}rr}(r) \wedge \\ \quad \operatorname{restrict}(h, \operatorname{ocint}((x - r), x)) \in (\operatorname{ocint}((x - r), x) \to \mathbb{R}) \wedge \\ \quad \forall y \in \operatorname{ocint}((x - r), x).\; (h(y) \simeq (f(y) + g(y))) \wedge \\ \quad \operatorname{has\text{-}left\text{-}deriv\text{-}at}(f, x, l) \wedge \\ \quad \operatorname{has\text{-}left\text{-}deriv\text{-}at}(g, x, m) \Rightarrow \\ \quad \quad \operatorname{has\text{-}left\text{-}deriv\text{-}at}(h, x, (l + m)) \end{array}}}

\vnbbookfile{structure-library/interval-taylor}{structure-library/interval-taylor}
\noindent{\small interval-taylor.scm -- THE VOCABULARY OF TAYLOR'S FORMULA ON AN INTERVAL. DEFINITIONS ONLY; every theorem is in theorem-library/interval-taylor.scm.}\par
\par\medskip\noindent\textbf{Definition} (\texttt{taylor-sum}).\;$\operatorname{taylor\text{-}sum}(d, n, b, \operatorname{tv\_}) :=$ $\operatorname{series\text{-}partial\text{-}sum}(\operatorname{taylor\text{-}term}(d, b, \operatorname{tv\_}), n)$
\par\medskip\noindent\textbf{Definition} (\texttt{taylor-aux}).\;$\operatorname{taylor\text{-}aux}(d, n, a, b) :=$ $(\lambda \operatorname{tu\_} \in \operatorname{ccint}(a, b).\; \operatorname{taylor\text{-}sum}(d, n, b, \operatorname{tu\_}))$
\par\medskip\noindent\textbf{Definition} (\texttt{taylor-rem}).\;$\operatorname{taylor\text{-}rem}(d, n, a, b) :=$ $(d(0)(b) - \operatorname{taylor\text{-}sum}(d, n, b, a))$
\par\medskip\noindent\textbf{Definition} (\texttt{taylor-term}).\;$\operatorname{taylor\text{-}term}(d, b, \operatorname{tv\_}) :=$
\begin{equation*}
(\lambda \operatorname{tk\_} \in \mathbb{N}.\; ((d(\operatorname{tk\_})(\operatorname{tv\_}) \cdot \operatorname{power}((b - \operatorname{tv\_}), \operatorname{tk\_})) \cdot \operatorname{recip}(\operatorname{factorial}(\operatorname{tk\_}))))
\end{equation*}

\par\medskip\noindent\textbf{Definition} (\texttt{is-taylor-family}).\;For all $d$, $n$, $a$ and $b$, 
\begin{equation*}
\begin{array}{@{}l@{}} \operatorname{is\text{-}taylor\text{-}family}(d, n, a, b) \\ \Leftrightarrow\; n \in \mathbb{N} \wedge \\ 1 \leq n \wedge \\ a \in \mathbb{R} \wedge \\ b \in \mathbb{R} \wedge \\ a < b \wedge \\ d(n) \in (\operatorname{ooint}(a, b) \to \mathbb{R}) \wedge \\ \forall \operatorname{tk\_} \in \mathbb{N}. \\ \quad (\operatorname{tk\_} < n) \Rightarrow  \\ \quad \quad d(\operatorname{tk\_}) \in (\operatorname{ccint}(a, b) \to \mathbb{R}) \wedge \\ \quad \quad \operatorname{is\text{-}continuous\text{-}on}(d(\operatorname{tk\_}), \operatorname{ccint}(a, b)) \wedge \\ \quad \quad \forall \operatorname{tx\_} \in \operatorname{ooint}(a, b).\; \operatorname{has\text{-}deriv\text{-}at}(d(\operatorname{tk\_}), \operatorname{tx\_}, d(\operatorname{succ}(\operatorname{tk\_}))(\operatorname{tx\_})) \end{array}
\end{equation*}


\vnbbookfile{theorem-library/directional-derivative}{theorem-library/directional-derivative}
\noindent{\small directional-derivative.scm -- calculus.pdf Chapter 2 section 8 "Multivariable Calculus", Def 2.23 through Prop 2.27, over an arbitrary real normed vector space.}\par
\par\medskip\noindent\textbf{Definition} (\texttt{seg-curve}).\;$\operatorname{seg\text{-}curve}(m, f, a, \eta) :=$ $(\lambda \operatorname{u\_} \in \mathbb{R}.\; f(\operatorname{vadd}(m)(a, \operatorname{act}(m)(\operatorname{u\_}, \eta))))$
\par\medskip\noindent\textbf{Definition} (\texttt{dir-deriv}).\;$\operatorname{dir\text{-}deriv}(m, f, a, \eta) :=$ $\operatorname{deriv}(\operatorname{seg\text{-}curve}(m, f, a, \eta), 0)$
\par\medskip\noindent\textbf{Definition} (\texttt{is-dir-diff-at}).\;For all $m$, $f$, $a$, $\eta$ and $\mathit{dl}$, 
\begin{equation*}
\operatorname{is\text{-}dir\text{-}diff\text{-}at}(m, f, a, \eta, \mathit{dl}) \Leftrightarrow \operatorname{is\text{-}diff\text{-}at}(\operatorname{seg\text{-}curve}(m, f, a, \eta), 0, \mathit{dl})
\end{equation*}

\vnbprop{seg-curve-apply}{\vnbstmt{\begin{array}{@{}l@{}} \forall m, f, a, b.; \forall t \in \mathbb{R}. \\ \quad (\operatorname{seg\text{-}curve}(m, f, a, b)(t) \simeq f(\operatorname{vadd}(m)(a, \operatorname{act}(m)(t, b)))) \end{array}}}
\vnbprop{diff-at-in-fun}{\vnbstmt{\forall f, a, l.\; (\operatorname{is\text{-}diff\text{-}at}(f, a, l) \Rightarrow (f \in (\mathbb{R} \to \mathbb{R})))}}
\vnbprop{affine-lam-in-fun}{\vnbstmt{\begin{array}{@{}l@{}} \forall c, m. \\ \quad ((c \in \mathbb{R}) \wedge (m \in \mathbb{R})) \Rightarrow ((\lambda x \in \mathbb{R}.\; (c + (m \cdot x))) \in (\mathbb{R} \to \mathbb{R})) \end{array}}}
\vnbprop{deriv-affine}{\vnbstmt{\begin{array}{@{}l@{}} \forall c, m, a. \\ \quad ((c \in \mathbb{R}) \wedge ((m \in \mathbb{R}) \wedge (a \in \mathbb{R}))) \Rightarrow  \\ \quad \quad \operatorname{is\text{-}diff\text{-}at}((\lambda x \in \mathbb{R}.\; (c + (m \cdot x))), a, m) \end{array}}}
\vnbprop{dir-diff-reparam}{\vnbstmt{\begin{array}{@{}l@{}} \forall f, g, c, m, l. \\ \quad f \in (\mathbb{R} \to \mathbb{R}) \wedge \\ \quad c \in \mathbb{R} \wedge \\ \quad m \in \mathbb{R} \wedge \\ \quad \operatorname{is\text{-}diff\text{-}at}(g, c, l) \wedge \\ \quad \forall x \in \mathbb{R}.\; (f(x) \simeq g((c + (m \cdot x)))) \Rightarrow \\ \quad \quad \operatorname{is\text{-}diff\text{-}at}(f, 0, (l \cdot m)) \end{array}}}
\vnbprop{dir-diff-scale}{\vnbstmt{\begin{array}{@{}l@{}} \forall f, g, m, l. \\ \quad ((f \in (\mathbb{R} \to \mathbb{R})) \wedge ((m \in \mathbb{R}) \wedge (\operatorname{is\text{-}diff\text{-}at}(g, 0, l) \wedge \forall x \in \mathbb{R}.\; (f(x) \simeq g((m \cdot x)))))) \Rightarrow  \\ \quad \quad \operatorname{is\text{-}diff\text{-}at}(f, 0, (m \cdot l)) \end{array}}}
\vnbprop{dir-diff-shift}{\vnbstmt{\begin{array}{@{}l@{}} \forall f, g, c, l. \\ \quad ((f \in (\mathbb{R} \to \mathbb{R})) \wedge ((c \in \mathbb{R}) \wedge (\operatorname{is\text{-}diff\text{-}at}(g, c, l) \wedge \forall x \in \mathbb{R}.\; (f(x) \simeq g((c + x)))))) \Rightarrow  \\ \quad \quad \operatorname{is\text{-}diff\text{-}at}(f, 0, l) \end{array}}}
\vnbprop{seg-curve-in-fun}{\vnbstmt{\begin{array}{@{}l@{}} \forall m, f, a, b. \\ \quad (\operatorname{is\text{-}normed\text{-}vector\text{-}space}(m) \wedge ((f \in (\operatorname{vec}(m) \to \mathbb{R})) \wedge ((a \in \operatorname{vec}(m)) \wedge (b \in \operatorname{vec}(m))))) \Rightarrow  \\ \quad \quad \operatorname{seg\text{-}curve}(m, f, a, b) \in (\mathbb{R} \to \mathbb{R}) \end{array}}}
\vnbprop{dir-deriv-homogeneous}{\vnbstmt{\begin{array}{@{}l@{}} \forall m, f, a, b, c, l. \\ \quad \operatorname{is\text{-}normed\text{-}vector\text{-}space}(m) \wedge \\ \quad f \in (\operatorname{vec}(m) \to \mathbb{R}) \wedge \\ \quad a \in \operatorname{vec}(m) \wedge \\ \quad b \in \operatorname{vec}(m) \wedge \\ \quad c \in \mathbb{R} \wedge \\ \quad \operatorname{is\text{-}dir\text{-}diff\text{-}at}(m, f, a, b, l) \Rightarrow \\ \quad \quad \operatorname{is\text{-}dir\text{-}diff\text{-}at}(m, f, a, \operatorname{act}(m)(c, b), (c \cdot l)) \end{array}}}
\vnbprop{dir-deriv-mvt}{\vnbstmt{\begin{array}{@{}l@{}} \forall m, f, a, b. \\ \quad \operatorname{is\text{-}normed\text{-}vector\text{-}space}(m) \wedge \\ \quad f \in (\operatorname{vec}(m) \to \mathbb{R}) \wedge \\ \quad a \in \operatorname{vec}(m) \wedge \\ \quad b \in \operatorname{vec}(m) \wedge \\ \quad \forall x \in \operatorname{ccint}(0, 1). \\ \quad \quad \operatorname{is\text{-}continuous\text{-}at}(\operatorname{rr\text{-}ms}, \operatorname{rr\text{-}ms}, \operatorname{seg\text{-}curve}(m, f, a, b), x) \wedge \\ \quad \forall x \in \mathbb{R}. \\ \quad \quad ((0 < x) \wedge (x < 1)) \Rightarrow \exists l.\; \operatorname{is\text{-}diff\text{-}at}(\operatorname{seg\text{-}curve}(m, f, a, b), x, l) \Rightarrow \\ \quad \quad \exists h \in \mathbb{R}. \\ \quad \quad \quad 0 < h \wedge \\ \quad \quad \quad h < 1 \wedge \\ \quad \quad \quad \exists v. \\ \quad \quad \quad \quad (\operatorname{is\text{-}dir\text{-}diff\text{-}at}(m, f, \operatorname{vadd}(m)(a, \operatorname{act}(m)(h, b)), b, v) \wedge (v = (f(\operatorname{vadd}(m)(a, b)) - f(a)))) \end{array}}}
\vnbprop{dir-deriv-value}{\vnbstmt{\begin{array}{@{}l@{}} \forall m, f, a, b, l. \\ \quad \operatorname{is\text{-}dir\text{-}diff\text{-}at}(m, f, a, b, l) \Rightarrow (\operatorname{dir\text{-}deriv}(m, f, a, b) = l) \end{array}}}
\vnbprop{diff-at-shift-back}{\vnbstmt{\begin{array}{@{}l@{}} \forall f, g, c, l. \\ \quad ((f \in (\mathbb{R} \to \mathbb{R})) \wedge ((c \in \mathbb{R}) \wedge (\operatorname{is\text{-}diff\text{-}at}(g, 0, l) \wedge \forall x \in \mathbb{R}.\; (g(x) \simeq f((c + x)))))) \Rightarrow  \\ \quad \quad \operatorname{is\text{-}diff\text{-}at}(f, c, l) \end{array}}}
\vnbprop{seg-curve-diff-at}{\vnbstmt{\begin{array}{@{}l@{}} \forall m.; \forall f \in (\operatorname{vec}(m) \to \mathbb{R}), b, a \in \operatorname{vec}(m), s \in \mathbb{R}, k. \\ \quad (\operatorname{is\text{-}normed\text{-}vector\text{-}space}(m) \wedge \operatorname{is\text{-}dir\text{-}diff\text{-}at}(m, f, \operatorname{vadd}(m)(b, \operatorname{act}(m)(s, a)), a, k)) \Rightarrow  \\ \quad \quad \operatorname{is\text{-}diff\text{-}at}(\operatorname{seg\text{-}curve}(m, f, b, a), s, k) \end{array}}}
\vnbprop{dir-deriv-increment}{\vnbstmt{\begin{array}{@{}l@{}} \forall m.; \forall f \in (\operatorname{vec}(m) \to \mathbb{R}), b, a \in \operatorname{vec}(m). \\ \quad (\operatorname{is\text{-}normed\text{-}vector\text{-}space}(m) \wedge \forall x \in \operatorname{vec}(m).\; \exists k.\; \operatorname{is\text{-}dir\text{-}diff\text{-}at}(m, f, x, a, k)) \Rightarrow  \\ \quad \quad \exists h \in \mathbb{R}. \\ \quad \quad \quad ((0 < h) \wedge ((h < 1) \wedge ((f(\operatorname{vadd}(m)(b, a)) - f(b)) = \operatorname{dir\text{-}deriv}(m, f, \operatorname{vadd}(m)(b, \operatorname{act}(m)(h, a)), a)))) \end{array}}}
\vnbprop{dir-diff-value-in-rr}{\vnbstmt{\forall m, f, a, b, l.\; (\operatorname{is\text{-}dir\text{-}diff\text{-}at}(m, f, a, b, l) \Rightarrow (l \in \mathbb{R}))}}
\vnbprop{dir-deriv-increment-scaled}{\vnbstmt{\begin{array}{@{}l@{}} \forall m.; \forall f \in (\operatorname{vec}(m) \to \mathbb{R}), b, a \in \operatorname{vec}(m), t \in \mathbb{R}. \\ \quad (\operatorname{is\text{-}normed\text{-}vector\text{-}space}(m) \wedge \forall x \in \operatorname{vec}(m).\; \exists k.\; \operatorname{is\text{-}dir\text{-}diff\text{-}at}(m, f, x, a, k)) \Rightarrow  \\ \quad \quad \exists h \in \mathbb{R}. \\ \quad \quad \quad 0 < h \wedge \\ \quad \quad \quad h < 1 \wedge \\ \quad \quad \quad (f(\operatorname{vadd}(m)(b, \operatorname{act}(m)(t, a))) - f(b)) = (t \cdot \operatorname{dir\text{-}deriv}(m, f, \operatorname{vadd}(m)(b, \operatorname{act}(m)((h \cdot t), a)), a)) \end{array}}}
\vnbprop{rr-abs-unit-scale}{\vnbstmt{\forall s, t \in \mathbb{R}.\; (((0 < s) \wedge (s < 1)) \Rightarrow (\lvert (s \cdot t) \rvert \leq \lvert t \rvert))}}
\vnbprop{nvs-norm-two-term-bound}{\vnbstmt{\begin{array}{@{}l@{}} \forall m.; \forall a, b \in \operatorname{vec}(m), c, d, f \in \mathbb{R}. \\ \quad (\operatorname{is\text{-}normed\text{-}vector\text{-}space}(m) \wedge ((0 \leq f) \wedge ((\lvert c \rvert \leq f) \wedge (\lvert d \rvert \leq f)))) \Rightarrow  \\ \quad \quad \operatorname{vnrm}(m)(\operatorname{vadd}(m)(\operatorname{act}(m)(c, a), \operatorname{act}(m)(d, b))) \leq (f \cdot (1 + (\operatorname{vnrm}(m)(b) + \operatorname{vnrm}(m)(a)))) \end{array}}}
\vnbprop{nvs-norm-one-term-bound}{\vnbstmt{\begin{array}{@{}l@{}} \forall m.; \forall a, b \in \operatorname{vec}(m), c, d \in \mathbb{R}. \\ \quad (\operatorname{is\text{-}normed\text{-}vector\text{-}space}(m) \wedge ((0 \leq d) \wedge (\lvert c \rvert \leq d))) \Rightarrow  \\ \quad \quad \operatorname{vnrm}(m)(\operatorname{act}(m)(c, a)) \leq (d \cdot (1 + (\operatorname{vnrm}(m)(b) + \operatorname{vnrm}(m)(a)))) \end{array}}}
\vnbprop{rr-diff-split-mid}{\vnbstmt{\forall x, y, z \in \mathbb{R}.\; ((x - z) = ((x - y) + (y - z)))}}
\vnbprop{dir-deriv-corner-split}{\vnbstmt{\begin{array}{@{}l@{}} \forall m.; \forall f \in (\operatorname{vec}(m) \to \mathbb{R}), a, b, i \in \operatorname{vec}(m), t \in \mathbb{R}. \\ \quad \operatorname{is\text{-}normed\text{-}vector\text{-}space}(m) \wedge \\ \quad \forall x \in \operatorname{vec}(m).\; \exists k.\; \operatorname{is\text{-}dir\text{-}diff\text{-}at}(m, f, x, b, k) \wedge \\ \quad \forall x \in \operatorname{vec}(m).\; \exists k.\; \operatorname{is\text{-}dir\text{-}diff\text{-}at}(m, f, x, i, k) \Rightarrow \\ \quad \quad \exists c, d \in \operatorname{vec}(m). \\ \quad \quad \quad \operatorname{dist}(\operatorname{nvs\text{-}metric\text{-}space}(m))(a, c) \leq (\lvert t \rvert \cdot (1 + (\operatorname{vnrm}(m)(b) + \operatorname{vnrm}(m)(i)))) \wedge \\ \quad \quad \quad \operatorname{dist}(\operatorname{nvs\text{-}metric\text{-}space}(m))(a, d) \leq (\lvert t \rvert \cdot (1 + (\operatorname{vnrm}(m)(b) + \operatorname{vnrm}(m)(i)))) \wedge \\ \quad \quad \quad (\operatorname{seg\text{-}curve}(m, f, a, \operatorname{vadd}(m)(b, i))(t) - \operatorname{seg\text{-}curve}(m, f, a, \operatorname{vadd}(m)(b, i))(0)) \\ \quad \quad \quad =\; t \cdot (\operatorname{dir\text{-}deriv}(m, f, c, b) + \operatorname{dir\text{-}deriv}(m, f, d, i)) \end{array}}}
\vnbprop{rr-factor-shift}{\vnbstmt{\forall d, u, r \in \mathbb{R}.\; (((d \cdot r) \cdot u) = (d \cdot (u \cdot r)))}}
\vnbprop{dir-quotient-in-fun}{\vnbstmt{\begin{array}{@{}l@{}} \forall m.; \forall f \in (\operatorname{vec}(m) \to \mathbb{R}), a, b, i \in \operatorname{vec}(m), l, c \in \mathbb{R}. \\ \quad \operatorname{is\text{-}normed\text{-}vector\text{-}space}(m) \Rightarrow  \\ \quad \quad (\lambda z \in \mathbb{R}.\; \\ \quad \quad \quad \left(\text{if }(z = 0)\text{ then }(l + c)\text{ else }((\operatorname{seg\text{-}curve}(m, f, a, \operatorname{vadd}(m)(b, i))(z) - \operatorname{seg\text{-}curve}(m, f, a, \operatorname{vadd}(m)(b, i))(0)) \cdot \operatorname{recip}(z))\right)) \\ \quad \quad \in\; (\mathbb{R} \to \mathbb{R}) \end{array}}}
\vnbprop{dir-quotient-at-zero}{\vnbstmt{\begin{array}{@{}l@{}} \forall m.; \forall f \in (\operatorname{vec}(m) \to \mathbb{R}), a, b, i \in \operatorname{vec}(m), l, c \in \mathbb{R}. \\ \quad \operatorname{is\text{-}normed\text{-}vector\text{-}space}(m) \Rightarrow  \\ \quad \quad (\lambda z \in \mathbb{R}.\; \left(\text{if }(z = 0)\text{ then }(l + c)\text{ else }((\operatorname{seg\text{-}curve}(m, f, a, \operatorname{vadd}(m)(b, i))(z) - \operatorname{seg\text{-}curve}(m, f, a, \operatorname{vadd}(m)(b, i))(0)) \cdot \operatorname{recip}(z))\right))(0) \\ \quad \quad =\; l + c \end{array}}}
\vnbprop{dir-quotient-identity}{\vnbstmt{\begin{array}{@{}l@{}} \forall m.; \forall f \in (\operatorname{vec}(m) \to \mathbb{R}), a, b, i \in \operatorname{vec}(m), l, c, x \in \mathbb{R}. \\ \quad \operatorname{is\text{-}normed\text{-}vector\text{-}space}(m) \Rightarrow  \\ \quad \quad (\operatorname{seg\text{-}curve}(m, f, a, \operatorname{vadd}(m)(b, i))(x) - \operatorname{seg\text{-}curve}(m, f, a, \operatorname{vadd}(m)(b, i))(0)) \\ \quad \quad =\; (\lambda z \in \mathbb{R}.\; \left(\text{if }(z = 0)\text{ then }(l + c)\text{ else }((\operatorname{seg\text{-}curve}(m, f, a, \operatorname{vadd}(m)(b, i))(z) - \operatorname{seg\text{-}curve}(m, f, a, \operatorname{vadd}(m)(b, i))(0)) \cdot \operatorname{recip}(z))\right))(x) \cdot x \end{array}}}
\vnbprop{rr-mul-rotate}{\vnbstmt{\forall s, u, r \in \mathbb{R}.\; (((u \cdot s) \cdot r) = (s \cdot (u \cdot r)))}}
\vnbprop{dir-quotient-continuous}{\vnbstmt{\begin{array}{@{}l@{}} \forall m.; \forall f \in (\operatorname{vec}(m) \to \mathbb{R}), a, b, i \in \operatorname{vec}(m), l, c. \\ \quad \operatorname{is\text{-}normed\text{-}vector\text{-}space}(m) \wedge \\ \quad \operatorname{is\text{-}dir\text{-}diff\text{-}at}(m, f, a, b, l) \wedge \\ \quad \operatorname{is\text{-}dir\text{-}diff\text{-}at}(m, f, a, i, c) \wedge \\ \quad \forall x \in \operatorname{vec}(m).\; \exists k.\; \operatorname{is\text{-}dir\text{-}diff\text{-}at}(m, f, x, b, k) \wedge \\ \quad \forall x \in \operatorname{vec}(m).\; \exists k.\; \operatorname{is\text{-}dir\text{-}diff\text{-}at}(m, f, x, i, k) \wedge \\ \quad \operatorname{is\text{-}continuous\text{-}at}(\operatorname{nvs\text{-}metric\text{-}space}(m), \operatorname{rr\text{-}ms}, (\lambda x \in \operatorname{vec}(m).\; \operatorname{dir\text{-}deriv}(m, f, x, b)), a) \wedge \\ \quad \operatorname{is\text{-}continuous\text{-}at}(\operatorname{nvs\text{-}metric\text{-}space}(m), \operatorname{rr\text{-}ms}, (\lambda x \in \operatorname{vec}(m).\; \operatorname{dir\text{-}deriv}(m, f, x, i)), a) \Rightarrow \\ \quad \quad \operatorname{is\text{-}continuous\text{-}at}(\operatorname{rr\text{-}ms}, \operatorname{rr\text{-}ms}, \\ \quad \quad \quad (\lambda z \in \mathbb{R}.\; \\ \quad \quad \quad \quad \left(\text{if }(z = 0)\text{ then }(l + c)\text{ else }((\operatorname{seg\text{-}curve}(m, f, a, \operatorname{vadd}(m)(b, i))(z) - \operatorname{seg\text{-}curve}(m, f, a, \operatorname{vadd}(m)(b, i))(0)) \cdot \operatorname{recip}(z))\right)), \\ \quad \quad \quad 0) \end{array}}}
\vnbprop{dir-deriv-additive}{\vnbstmt{\begin{array}{@{}l@{}} \forall m, f, a, b, i, l, c. \\ \quad \operatorname{is\text{-}normed\text{-}vector\text{-}space}(m) \wedge \\ \quad f \in (\operatorname{vec}(m) \to \mathbb{R}) \wedge \\ \quad a \in \operatorname{vec}(m) \wedge \\ \quad b \in \operatorname{vec}(m) \wedge \\ \quad i \in \operatorname{vec}(m) \wedge \\ \quad \operatorname{is\text{-}dir\text{-}diff\text{-}at}(m, f, a, b, l) \wedge \\ \quad \operatorname{is\text{-}dir\text{-}diff\text{-}at}(m, f, a, i, c) \wedge \\ \quad \forall x \in \operatorname{vec}(m).\; \exists k.\; \operatorname{is\text{-}dir\text{-}diff\text{-}at}(m, f, x, b, k) \wedge \\ \quad \forall x \in \operatorname{vec}(m).\; \exists k.\; \operatorname{is\text{-}dir\text{-}diff\text{-}at}(m, f, x, i, k) \wedge \\ \quad \operatorname{is\text{-}continuous\text{-}at}(\operatorname{nvs\text{-}metric\text{-}space}(m), \operatorname{rr\text{-}ms}, (\lambda x \in \operatorname{vec}(m).\; \operatorname{dir\text{-}deriv}(m, f, x, b)), a) \wedge \\ \quad \operatorname{is\text{-}continuous\text{-}at}(\operatorname{nvs\text{-}metric\text{-}space}(m), \operatorname{rr\text{-}ms}, (\lambda x \in \operatorname{vec}(m).\; \operatorname{dir\text{-}deriv}(m, f, x, i)), a) \Rightarrow \\ \quad \quad \operatorname{is\text{-}dir\text{-}diff\text{-}at}(m, f, a, \operatorname{vadd}(m)(b, i), (l + c)) \end{array}}}
\vnbprop{dir-deriv-linear}{\vnbstmt{\begin{array}{@{}l@{}} \forall m, f, a, b, i, r, s, l, c. \\ \quad \operatorname{is\text{-}normed\text{-}vector\text{-}space}(m) \wedge \\ \quad f \in (\operatorname{vec}(m) \to \mathbb{R}) \wedge \\ \quad a \in \operatorname{vec}(m) \wedge \\ \quad b \in \operatorname{vec}(m) \wedge \\ \quad i \in \operatorname{vec}(m) \wedge \\ \quad r \in \mathbb{R} \wedge \\ \quad s \in \mathbb{R} \wedge \\ \quad \operatorname{is\text{-}dir\text{-}diff\text{-}at}(m, f, a, b, l) \wedge \\ \quad \operatorname{is\text{-}dir\text{-}diff\text{-}at}(m, f, a, i, c) \wedge \\ \quad \forall x \in \operatorname{vec}(m).\; \exists k.\; \operatorname{is\text{-}dir\text{-}diff\text{-}at}(m, f, x, \operatorname{act}(m)(r, b), k) \wedge \\ \quad \forall x \in \operatorname{vec}(m).\; \exists k.\; \operatorname{is\text{-}dir\text{-}diff\text{-}at}(m, f, x, \operatorname{act}(m)(s, i), k) \wedge \\ \quad \operatorname{is\text{-}continuous\text{-}at}(\operatorname{nvs\text{-}metric\text{-}space}(m), \operatorname{rr\text{-}ms}, \\ \quad \quad (\lambda x \in \operatorname{vec}(m).\; \operatorname{dir\text{-}deriv}(m, f, x, \operatorname{act}(m)(r, b))), \\ \quad \quad a) \wedge \\ \quad \operatorname{is\text{-}continuous\text{-}at}(\operatorname{nvs\text{-}metric\text{-}space}(m), \operatorname{rr\text{-}ms}, \\ \quad \quad (\lambda x \in \operatorname{vec}(m).\; \operatorname{dir\text{-}deriv}(m, f, x, \operatorname{act}(m)(s, i))), \\ \quad \quad a) \Rightarrow \\ \quad \quad \operatorname{is\text{-}dir\text{-}diff\text{-}at}(m, f, a, \operatorname{vadd}(m)(\operatorname{act}(m)(r, b), \operatorname{act}(m)(s, i)), ((r \cdot l) + (s \cdot c))) \end{array}}}

\vnbbookfile{theorem-library/nvs-taylor-proof}{theorem-library/nvs-taylor-proof}
\noindent{\small nvs-taylor-proof.scm -- Taylor's theorem with remainder bound for a map between two finite-dimensional real normed vector spaces.  PROVED.}\par
\vnbprop{rr-power-one}{\vnbstmt{\forall k \in \mathbb{N}.\; (\operatorname{power}(1, k) = 1)}}
\vnbprop{nt-unit-factor}{\vnbstmt{\forall y, k.\; (((y \in \mathbb{R}) \wedge (k \in \mathbb{N})) \Rightarrow ((y \cdot \operatorname{power}(1, (k + 1))) = y))}}
\vnbprop{nvs-taylor-remainder-bound}{\vnbstmt{\begin{array}{@{}l@{}} \forall m, b, f, a, h, n. \\ \quad \operatorname{is\text{-}normed\text{-}vector\text{-}space}(m) \wedge \\ \quad \operatorname{is\text{-}finite\text{-}dimensional}(\operatorname{normed\text{-}vector\text{-}space\text{-}as\text{-}module}(m)) \wedge \\ \quad \operatorname{is\text{-}normed\text{-}vector\text{-}space}(b) \wedge \\ \quad \operatorname{is\text{-}finite\text{-}dimensional}(\operatorname{normed\text{-}vector\text{-}space\text{-}as\text{-}module}(b)) \wedge \\ \quad f \in (\operatorname{vec}(m) \to \operatorname{vec}(b)) \wedge \\ \quad a \in \operatorname{vec}(m) \wedge \\ \quad h \in \operatorname{vec}(m) \wedge \\ \quad n \in \mathbb{N} \wedge \\ \quad \operatorname{taylor\text{-}differentiable\text{-}v}(b, (\lambda t \in \mathbb{R}.\; f(\operatorname{vadd}(m)(a, \operatorname{act}(m)(t, h)))), 0, 1, n) \Rightarrow \\ \quad \quad \exists c. \\ \quad \quad \quad 0 < c \wedge \\ \quad \quad \quad c < 1 \wedge \\ \quad \quad \quad \operatorname{factorial}((n + 1)) \cdot \operatorname{vnrm}(b)(\operatorname{vadd}(b)(f(\operatorname{vadd}(m)(a, h)), \operatorname{vneg}(b)(\operatorname{taylor\text{-}poly\text{-}v}(b, (\lambda t \in \mathbb{R}.\; f(\operatorname{vadd}(m)(a, \operatorname{act}(m)(t, h)))), 0, 1, n)))) \\ \quad \quad \quad \leq\; \operatorname{vnrm}(b)(\operatorname{nth\text{-}deriv\text{-}v}(b, (\lambda t \in \mathbb{R}.\; f(\operatorname{vadd}(m)(a, \operatorname{act}(m)(t, h)))), (n + 1))(c)) \end{array}}}

\vnbbookfile{theorem-library/antiderivative}{theorem-library/antiderivative}
\noindent{\small antiderivative.scm -- docs/calculus.pdf CHAPTER 4 SECTION 2: Definition 4.6 (antiderivative / antiderivable), Proposition 4.8 (closed under sums; the scalar half is handed off), Proposition 4.10 (two antiderivatives of the same function differ by a constant) and Corollary 4.11 (so f(b) - f(a) does}\par
\vnbprop{antiderivative-map-in-fun}{\vnbstmt{\forall f, i, a, b.\; (\operatorname{is\text{-}antiderivative}(f, i, a, b) \Rightarrow (f \in (\mathbb{R} \to \mathbb{R})))}}
\vnbprop{antiderivative-fn-in-fun}{\vnbstmt{\forall f, i, a, b.\; (\operatorname{is\text{-}antiderivative}(f, i, a, b) \Rightarrow (i \in (\mathbb{R} \to \mathbb{R})))}}
\vnbprop{antiderivative-endpoints}{\vnbstmt{\begin{array}{@{}l@{}} \forall f, i, a, b. \\ \quad \operatorname{is\text{-}antiderivative}(f, i, a, b) \Rightarrow ((a \in \mathbb{R}) \wedge ((b \in \mathbb{R}) \wedge (a < b))) \end{array}}}
\vnbprop{deriv-sub}{\vnbstmt{\begin{array}{@{}l@{}} \forall f, g, a, l, m. \\ \quad (\operatorname{is\text{-}diff\text{-}at}(f, a, l) \wedge \operatorname{is\text{-}diff\text{-}at}(g, a, m)) \Rightarrow  \\ \quad \quad \operatorname{is\text{-}diff\text{-}at}((\lambda x \in \mathbb{R}.\; (f(x) - g(x))), a, (l - m)) \end{array}}}
\vnbprop{antiderivative-differ-by-constant}{\vnbstmt{\begin{array}{@{}l@{}} \forall f, g, i, a, b. \\ \quad (\operatorname{is\text{-}antiderivative}(f, i, a, b) \wedge \operatorname{is\text{-}antiderivative}(g, i, a, b)) \Rightarrow  \\ \quad \quad \exists t \in \mathbb{R}.\; \forall s \in \operatorname{ccint}(a, b).\; (f(s) = (g(s) + t)) \end{array}}}
\vnbprop{antiderivative-endpoint-difference}{\vnbstmt{\begin{array}{@{}l@{}} \forall f, g, i, a, b. \\ \quad (\operatorname{is\text{-}antiderivative}(f, i, a, b) \wedge \operatorname{is\text{-}antiderivative}(g, i, a, b)) \Rightarrow  \\ \quad \quad (f(b) - f(a)) = (g(b) - g(a)) \end{array}}}
\vnbprop{poly-is-antiderivable}{\vnbstmt{\begin{array}{@{}l@{}} \forall n \in \mathbb{N}, f \in (\mathbb{N} \to \mathbb{R}), a, b. \\ \quad ((a \in \mathbb{R}) \wedge ((b \in \mathbb{R}) \wedge (a < b))) \Rightarrow  \\ \quad \quad \operatorname{is\text{-}antiderivable}((\lambda x \in \mathbb{R}.\; \operatorname{series\text{-}partial\text{-}sum}((\lambda k \in \mathbb{N}.\; (f(k) \cdot \operatorname{power}(x, k))), (n + 1))), \\ \quad \quad \quad a, b) \end{array}}}
\vnbprop{antiderivative-add}{\vnbstmt{\begin{array}{@{}l@{}} \forall f, g, i, c, a, b. \\ \quad (\operatorname{is\text{-}antiderivative}(f, i, a, b) \wedge \operatorname{is\text{-}antiderivative}(g, c, a, b)) \Rightarrow  \\ \quad \quad \operatorname{is\text{-}antiderivative}((\lambda x \in \mathbb{R}.\; (f(x) + g(x))), (\lambda x \in \mathbb{R}.\; (i(x) + c(x))), a, b) \end{array}}}
\vnbprop{antiderivable-add}{\vnbstmt{\begin{array}{@{}l@{}} \forall i, c, a, b. \\ \quad (\operatorname{is\text{-}antiderivable}(i, a, b) \wedge \operatorname{is\text{-}antiderivable}(c, a, b)) \Rightarrow  \\ \quad \quad \operatorname{is\text{-}antiderivable}((\lambda x \in \mathbb{R}.\; (i(x) + c(x))), a, b) \end{array}}}
\vnbprop{antiderivative-scale}{\vnbstmt{\begin{array}{@{}l@{}} \forall c, f, i, a, b. \\ \quad ((c \in \mathbb{R}) \wedge \operatorname{is\text{-}antiderivative}(f, i, a, b)) \Rightarrow  \\ \quad \quad \operatorname{is\text{-}antiderivative}((\lambda x \in \mathbb{R}.\; (c \cdot f(x))), (\lambda x \in \mathbb{R}.\; (c \cdot i(x))), a, b) \end{array}}}
\vnbprop{antiderivable-scale}{\vnbstmt{\begin{array}{@{}l@{}} \forall c, i, a, b. \\ \quad ((c \in \mathbb{R}) \wedge \operatorname{is\text{-}antiderivable}(i, a, b)) \Rightarrow  \\ \quad \quad \operatorname{is\text{-}antiderivable}((\lambda x \in \mathbb{R}.\; (c \cdot i(x))), a, b) \end{array}}}

\vnbbookfile{theorem-library/antiderivative-transfer}{theorem-library/antiderivative-transfer}
\noindent{\small antiderivative-transfer.scm -- the POINTWISE TRANSFER for Def 4.6, and the DIFFERENCE of two antiderivatives.  All three `modulo 0'.}\par
\vnbprop{antiderivative-transfer-ptwise-eq}{\vnbstmt{\begin{array}{@{}l@{}} \forall f, i \in (\mathbb{R} \to \mathbb{R}), g, c, a, b. \\ \quad (\operatorname{is\text{-}antiderivative}(g, c, a, b) \wedge (\forall x \in \mathbb{R}.\; (f(x) = g(x)) \wedge \forall x \in \mathbb{R}.\; (i(x) = c(x)))) \Rightarrow  \\ \quad \quad \operatorname{is\text{-}antiderivative}(f, i, a, b) \end{array}}}
\vnbprop{antiderivative-sub}{\vnbstmt{\begin{array}{@{}l@{}} \forall f, g, i, c, a, b. \\ \quad (\operatorname{is\text{-}antiderivative}(f, i, a, b) \wedge \operatorname{is\text{-}antiderivative}(g, c, a, b)) \Rightarrow  \\ \quad \quad \operatorname{is\text{-}antiderivative}((\lambda x \in \mathbb{R}.\; (f(x) - g(x))), (\lambda x \in \mathbb{R}.\; (i(x) - c(x))), a, b) \end{array}}}
\vnbprop{antiderivable-sub}{\vnbstmt{\begin{array}{@{}l@{}} \forall i, c, a, b. \\ \quad (\operatorname{is\text{-}antiderivable}(i, a, b) \wedge \operatorname{is\text{-}antiderivable}(c, a, b)) \Rightarrow  \\ \quad \quad \operatorname{is\text{-}antiderivable}((\lambda x \in \mathbb{R}.\; (i(x) - c(x))), a, b) \end{array}}}

\vnbbookfile{theorem-library/antiderivative-lipschitz}{theorem-library/antiderivative-lipschitz}
\noindent{\small antiderivative-lipschitz.scm -- THE WITNESS ESTIMATE of docs/calculus.pdf Prop 4.16: a bound on the INTEGRAND bounds the CARATHEODORY WITNESS, at every point of the interval and by the same constant.}\par
\vnbprop{antiderivative-increment-bound}{\vnbstmt{\begin{array}{@{}l@{}} \forall h, i \in (\mathbb{R} \to \mathbb{R}), a, b, m, u, v \in \mathbb{R}. \\ \quad \operatorname{is\text{-}antiderivative}(h, i, a, b) \wedge \\ \quad \forall x \in \mathbb{R}.\; ((x \in \operatorname{ccint}(a, b)) \Rightarrow (\lvert i(x) \rvert \leq m)) \wedge \\ \quad u \in \operatorname{ccint}(a, b) \wedge \\ \quad v \in \operatorname{ccint}(a, b) \wedge \\ \quad u < v \Rightarrow \\ \quad \quad \lvert (h(v) - h(u)) \rvert \leq (m \cdot (v - u)) \end{array}}}
\vnbprop{antiderivative-lipschitz}{\vnbstmt{\begin{array}{@{}l@{}} \forall h, i \in (\mathbb{R} \to \mathbb{R}), a, b, m, u, v \in \mathbb{R}. \\ \quad \operatorname{is\text{-}antiderivative}(h, i, a, b) \wedge \\ \quad \forall x \in \mathbb{R}.\; ((x \in \operatorname{ccint}(a, b)) \Rightarrow (\lvert i(x) \rvert \leq m)) \wedge \\ \quad u \in \operatorname{ccint}(a, b) \wedge \\ \quad v \in \operatorname{ccint}(a, b) \Rightarrow \\ \quad \quad \lvert (h(v) - h(u)) \rvert \leq (m \cdot \lvert (v - u) \rvert) \end{array}}}
\vnbprop{caratheodory-witness-bound}{\vnbstmt{\begin{array}{@{}l@{}} \forall h, i, c \in (\mathbb{R} \to \mathbb{R}), a, b, t, m, x \in \mathbb{R}. \\ \quad \operatorname{is\text{-}antiderivative}(h, i, a, b) \wedge \\ \quad \forall x \in \mathbb{R}.\; ((x \in \operatorname{ccint}(a, b)) \Rightarrow (\lvert i(x) \rvert \leq m)) \wedge \\ \quad t \in \operatorname{ccint}(a, b) \wedge \\ \quad c(t) = i(t) \wedge \\ \quad \forall x \in \mathbb{R}.\; ((h(x) - h(t)) = (c(x) \cdot (x - t))) \wedge \\ \quad x \in \operatorname{ccint}(a, b) \Rightarrow \\ \quad \quad \lvert c(x) \rvert \leq m \end{array}}}
\vnbprop{antiderivative-pair-close}{\vnbstmt{\begin{array}{@{}l@{}} \forall f, g, i, c \in (\mathbb{R} \to \mathbb{R}), a, b, m, x \in \mathbb{R}. \\ \quad \operatorname{is\text{-}antiderivative}(f, i, a, b) \wedge \\ \quad \operatorname{is\text{-}antiderivative}(g, c, a, b) \wedge \\ \quad f(a) = g(a) \wedge \\ \quad \forall x \in \mathbb{R}.\; ((x \in \operatorname{ccint}(a, b)) \Rightarrow (\lvert (i(x) - c(x)) \rvert \leq m)) \wedge \\ \quad x \in \operatorname{ccint}(a, b) \Rightarrow \\ \quad \quad \lvert (f(x) - g(x)) \rvert \leq (m \cdot (b - a)) \end{array}}}

\vnbbookfile{theorem-library/witness-family-choice}{theorem-library/witness-family-choice}
\noindent{\small witness-family-choice.scm -- TWO FAMILIES BUILT BY CHOICE, and nothing else.}\par
\vnbprop{antiderivable-family-choice}{\vnbstmt{\begin{array}{@{}l@{}} \forall m \in (\mathbb{N} \to (\mathbb{R} \to \mathbb{R})), a, b \in \mathbb{R}. \\ \quad \forall k \in \mathbb{N}.\; \operatorname{is\text{-}antiderivable}(m(k), a, b) \Rightarrow  \\ \quad \quad \exists c \in (\mathbb{N} \to (\mathbb{R} \to \mathbb{R})). \forall k \in \mathbb{N}. \\ \quad \quad \quad \operatorname{is\text{-}antiderivative}(c(k), m(k), a, b) \end{array}}}
\vnbprop{diff-at-witness-family}{\vnbstmt{\begin{array}{@{}l@{}} \forall m, a, t. \\ \quad m \in (\mathbb{N} \to (\mathbb{R} \to \mathbb{R})) \wedge \\ \quad a \in (\mathbb{N} \to \mathbb{R}) \wedge \\ \quad t \in \mathbb{R} \wedge \\ \quad \forall k \in \mathbb{N}.\; \operatorname{is\text{-}diff\text{-}at}(m(k), t, a(k)) \Rightarrow \\ \quad \quad \exists b \in (\mathbb{N} \to (\mathbb{R} \to \mathbb{R})). \forall k \in \mathbb{N}. \\ \quad \quad \quad \operatorname{is\text{-}continuous\text{-}at}(\operatorname{rr\text{-}ms}, \operatorname{rr\text{-}ms}, b(k), t) \wedge \\ \quad \quad \quad b(k)(t) = a(k) \wedge \\ \quad \quad \quad \forall x \in \mathbb{R}.\; ((m(k)(x) - m(k)(t)) = (b(k)(x) \cdot (x - t))) \end{array}}}

\vnbbookfile{theorem-library/antiderivative-normalize}{theorem-library/antiderivative-normalize}
\noindent{\small antiderivative-normalize.scm -- SHIFTING AN ANTIDERIVATIVE BY A CONSTANT, and the NORMALISED antiderivative family that docs/calculus.pdf Cor 4.17 needs and that nothing in the tree supplied.}\par
\vnbprop{antiderivative-shift-const}{\vnbstmt{\begin{array}{@{}l@{}} \forall f, i, a, b, c. \\ \quad f \in (\mathbb{R} \to \mathbb{R}) \wedge \\ \quad i \in (\mathbb{R} \to \mathbb{R}) \wedge \\ \quad a \in \mathbb{R} \wedge \\ \quad b \in \mathbb{R} \wedge \\ \quad c \in \mathbb{R} \wedge \\ \quad \operatorname{is\text{-}antiderivative}(f, i, a, b) \Rightarrow \\ \quad \quad \operatorname{is\text{-}antiderivative}((\lambda x \in \mathbb{R}.\; (f(x) - c)), i, a, b) \end{array}}}
\vnbprop{antiderivable-family-normalized}{\vnbstmt{\begin{array}{@{}l@{}} \forall m \in (\mathbb{N} \to (\mathbb{R} \to \mathbb{R})), a, b \in \mathbb{R}. \\ \quad \forall k \in \mathbb{N}.\; \operatorname{is\text{-}antiderivable}(m(k), a, b) \Rightarrow  \\ \quad \quad \exists c \in (\mathbb{N} \to (\mathbb{R} \to \mathbb{R})). \forall k \in \mathbb{N}. \\ \quad \quad \quad \operatorname{is\text{-}antiderivative}(c(k), m(k), a, b) \wedge (c(k)(a) = 0) \end{array}}}

\vnbbookfile{theorem-library/antiderivable-uniform-limit}{theorem-library/antiderivable-uniform-limit}
\noindent{\small theorem-library/antiderivable-uniform-limit.scm -- docs/calculus.pdf PROP 4.16 and COR 4.17: THE ANTIDERIVABLE FUNCTIONS ARE CLOSED UNDER UNIFORM LIMITS.  Rung 3 of the integration arc.}\par
\vnbprop{rr-converges-to-abs}{\vnbstmt{\begin{array}{@{}l@{}} \forall f \in (\mathbb{N} \to \mathbb{R}), v \in \mathbb{R}. \\ \quad \forall s. \\ \quad \quad \operatorname{pos\text{-}rr}(s) \Rightarrow  \\ \quad \quad \quad \exists n \in \mathbb{N}.\; \forall k \in \mathbb{N}.\; ((n \leq k) \Rightarrow (\lvert (f(k) - v) \rvert \leq s)) \Rightarrow \\ \quad \quad \operatorname{converges\text{-}to}(\operatorname{rr\text{-}ms}, f, v) \end{array}}}
\vnbprop{rr-limit-ptwise-eq}{\vnbstmt{\begin{array}{@{}l@{}} \forall f, h \in (\mathbb{N} \to \mathbb{R}), v \in \mathbb{R}. \\ \quad (\forall j \in \mathbb{N}.\; (h(j) = f(j)) \wedge \operatorname{converges\text{-}to}(\operatorname{rr\text{-}ms}, f, v)) \Rightarrow  \\ \quad \quad \operatorname{converges\text{-}to}(\operatorname{rr\text{-}ms}, h, v) \end{array}}}
\vnbprop{uniform-cauchy-limit}{\vnbstmt{\begin{array}{@{}l@{}} \forall g \in (\mathbb{N} \to (\mathbb{R} \to \mathbb{R})), s. \\ \quad \forall a. \\ \quad \quad \operatorname{pos\text{-}rr}(a) \Rightarrow  \\ \quad \quad \quad \exists n \in \mathbb{N}. \forall k, l \in \mathbb{N}, x \in \mathbb{R}. \\ \quad \quad \quad \quad ((n \leq k) \wedge ((n \leq l) \wedge (x \in s))) \Rightarrow (\lvert (g(k)(x) - g(l)(x)) \rvert \leq a) \Rightarrow \\ \quad \quad \exists h \in (\mathbb{R} \to \mathbb{R}). \\ \quad \quad \quad \forall x \in \mathbb{R}. \\ \quad \quad \quad \quad (x \in s) \Rightarrow \operatorname{converges\text{-}to}(\operatorname{rr\text{-}ms}, (\lambda k \in \mathbb{N}.\; g(k)(x)), h(x)) \wedge \\ \quad \quad \quad \forall a. \\ \quad \quad \quad \quad \operatorname{pos\text{-}rr}(a) \Rightarrow  \\ \quad \quad \quad \quad \quad \exists n \in \mathbb{N}. \forall k \in \mathbb{N}, x \in \mathbb{R}. \\ \quad \quad \quad \quad \quad \quad ((n \leq k) \wedge (x \in s)) \Rightarrow (\lvert (h(x) - g(k)(x)) \rvert \leq a) \end{array}}}
\vnbprop{caratheodory-witness-pair-bound}{\vnbstmt{\begin{array}{@{}l@{}} \forall c, d, f, g, h, k \in (\mathbb{R} \to \mathbb{R}), a, b, t, m, x \in \mathbb{R}. \\ \quad \operatorname{is\text{-}antiderivative}(c, f, a, b) \wedge \\ \quad \operatorname{is\text{-}antiderivative}(d, g, a, b) \wedge \\ \quad \forall x \in \mathbb{R}.\; ((x \in \operatorname{ccint}(a, b)) \Rightarrow (\lvert (f(x) - g(x)) \rvert \leq m)) \wedge \\ \quad t \in \operatorname{ccint}(a, b) \wedge \\ \quad h(t) = f(t) \wedge \\ \quad k(t) = g(t) \wedge \\ \quad \forall x \in \mathbb{R}.\; ((c(x) - c(t)) = (h(x) \cdot (x - t))) \wedge \\ \quad \forall x \in \mathbb{R}.\; ((d(x) - d(t)) = (k(x) \cdot (x - t))) \wedge \\ \quad x \in \operatorname{ccint}(a, b) \Rightarrow \\ \quad \quad \lvert (h(x) - k(x)) \rvert \leq m \end{array}}}
\vnbprop{uniform-witness-diff-at}{\vnbstmt{\begin{array}{@{}l@{}} \forall m, c \in (\mathbb{N} \to (\mathbb{R} \to \mathbb{R})), d \in (\mathbb{R} \to \mathbb{R}), v, a, b, t, o \in \mathbb{R}. \\ \quad 0 < o \wedge \\ \quad \forall y \in \mathbb{R}.\; ((\lvert (y - t) \rvert \leq o) \Rightarrow (y \in \operatorname{ccint}(a, b))) \wedge \\ \quad t \in \operatorname{ccint}(a, b) \wedge \\ \quad \forall k \in \mathbb{N}.\; \operatorname{is\text{-}continuous\text{-}at}(\operatorname{rr\text{-}ms}, \operatorname{rr\text{-}ms}, c(k), t) \wedge \\ \quad \forall k \in \mathbb{N}, y \in \mathbb{R}. \\ \quad \quad (y \in \operatorname{ccint}(a, b)) \Rightarrow ((m(k)(y) - m(k)(t)) = (c(k)(y) \cdot (y - t))) \wedge \\ \quad \forall y \in \mathbb{R}. \\ \quad \quad (y \in \operatorname{ccint}(a, b)) \Rightarrow \operatorname{converges\text{-}to}(\operatorname{rr\text{-}ms}, (\lambda k \in \mathbb{N}.\; m(k)(y)), d(y)) \wedge \\ \quad \forall s. \\ \quad \quad \operatorname{pos\text{-}rr}(s) \Rightarrow  \\ \quad \quad \quad \exists n \in \mathbb{N}. \forall k, l \in \mathbb{N}, x \in \mathbb{R}. \\ \quad \quad \quad \quad ((n \leq k) \wedge ((n \leq l) \wedge (x \in \operatorname{ccint}(a, b)))) \Rightarrow (\lvert (c(k)(x) - c(l)(x)) \rvert \leq s) \wedge \\ \quad \operatorname{converges\text{-}to}(\operatorname{rr\text{-}ms}, (\lambda k \in \mathbb{N}.\; c(k)(t)), v) \Rightarrow \\ \quad \quad \operatorname{is\text{-}diff\text{-}at}(d, t, v) \end{array}}}
\vnbprop{antiderivative-family-uniform-limit}{\vnbstmt{\begin{array}{@{}l@{}} \forall m, c \in (\mathbb{N} \to (\mathbb{R} \to \mathbb{R})), i \in (\mathbb{R} \to \mathbb{R}), a, b \in \mathbb{R}. \\ \quad a < b \wedge \\ \quad \forall k \in \mathbb{N}.\; (\operatorname{is\text{-}antiderivative}(m(k), c(k), a, b) \wedge (m(k)(a) = 0)) \wedge \\ \quad \forall s. \\ \quad \quad \operatorname{pos\text{-}rr}(s) \Rightarrow  \\ \quad \quad \quad \exists n \in \mathbb{N}. \forall k \in \mathbb{N}, x \in \mathbb{R}. \\ \quad \quad \quad \quad ((n \leq k) \wedge (x \in \operatorname{ccint}(a, b))) \Rightarrow (\lvert (i(x) - c(k)(x)) \rvert \leq s) \Rightarrow \\ \quad \quad \exists f \in (\mathbb{R} \to \mathbb{R}). \\ \quad \quad \quad \operatorname{is\text{-}antiderivative}(f, i, a, b) \wedge \\ \quad \quad \quad \forall s. \\ \quad \quad \quad \quad \operatorname{pos\text{-}rr}(s) \Rightarrow  \\ \quad \quad \quad \quad \quad \exists n \in \mathbb{N}. \forall k \in \mathbb{N}, x \in \mathbb{R}. \\ \quad \quad \quad \quad \quad \quad ((n \leq k) \wedge (x \in \operatorname{ccint}(a, b))) \Rightarrow (\lvert (f(x) - m(k)(x)) \rvert \leq s) \end{array}}}
\vnbprop{antiderivable-uniform-limit}{\vnbstmt{\begin{array}{@{}l@{}} \forall m \in (\mathbb{N} \to (\mathbb{R} \to \mathbb{R})), i \in (\mathbb{R} \to \mathbb{R}), a, b \in \mathbb{R}. \\ \quad \forall k \in \mathbb{N}.\; \operatorname{is\text{-}antiderivable}(m(k), a, b) \wedge \\ \quad \forall s. \\ \quad \quad \operatorname{pos\text{-}rr}(s) \Rightarrow  \\ \quad \quad \quad \exists n \in \mathbb{N}. \forall k \in \mathbb{N}, x \in \mathbb{R}. \\ \quad \quad \quad \quad ((n \leq k) \wedge (x \in \operatorname{ccint}(a, b))) \Rightarrow (\lvert (i(x) - m(k)(x)) \rvert \leq s) \Rightarrow \\ \quad \quad \operatorname{is\text{-}antiderivable}(i, a, b) \end{array}}}

\vnbbookfile{theorem-library/series-antiderivable}{theorem-library/series-antiderivable}
\noindent{\small series-antiderivable.scm -- A FINITE SUM OF ANTIDERIVABLE TERMS IS ANTIDERIVABLE, and the vocabulary that makes such a statement usable.}\par
\vnbprop{antiderivable-fn-in-fun}{\vnbstmt{\forall i, a, b.\; (\operatorname{is\text{-}antiderivable}(i, a, b) \Rightarrow (i \in (\mathbb{R} \to \mathbb{R})))}}
\vnbprop{antiderivative-integrand-transfer}{\vnbstmt{\begin{array}{@{}l@{}} \forall f, i, c, a, b. \\ \quad ((c \in (\mathbb{R} \to \mathbb{R})) \wedge (\operatorname{is\text{-}antiderivative}(f, i, a, b) \wedge \forall x \in \mathbb{R}.\; (c(x) \simeq i(x)))) \Rightarrow  \\ \quad \quad \operatorname{is\text{-}antiderivative}(f, c, a, b) \end{array}}}
\vnbprop{antiderivable-integrand-transfer}{\vnbstmt{\begin{array}{@{}l@{}} \forall i, c, a, b. \\ \quad ((c \in (\mathbb{R} \to \mathbb{R})) \wedge (\operatorname{is\text{-}antiderivable}(i, a, b) \wedge \forall x \in \mathbb{R}.\; (c(x) \simeq i(x)))) \Rightarrow  \\ \quad \quad \operatorname{is\text{-}antiderivable}(c, a, b) \end{array}}}
\vnbprop{antiderivable-from-deriv}{\vnbstmt{\begin{array}{@{}l@{}} \forall f, i, a, b. \\ \quad f \in (\mathbb{R} \to \mathbb{R}) \wedge \\ \quad i \in (\mathbb{R} \to \mathbb{R}) \wedge \\ \quad a \in \mathbb{R} \wedge \\ \quad b \in \mathbb{R} \wedge \\ \quad a < b \wedge \\ \quad \forall x \in \mathbb{R}.\; \operatorname{is\text{-}diff\text{-}at}(f, x, i(x)) \Rightarrow \\ \quad \quad \operatorname{is\text{-}antiderivable}(i, a, b) \end{array}}}
\vnbprop{zero-is-antiderivable}{\vnbstmt{\begin{array}{@{}l@{}} \forall a, b. \\ \quad ((a \in \mathbb{R}) \wedge ((b \in \mathbb{R}) \wedge (a < b))) \Rightarrow  \\ \quad \quad \operatorname{is\text{-}antiderivable}((\lambda x \in \mathbb{R}.\; 0), a, b) \end{array}}}
\vnbprop{series-partial-sum-antiderivable}{\vnbstmt{\begin{array}{@{}l@{}} \forall n \in \mathbb{N}, s, h, a, b. \\ \quad a \in \mathbb{R} \wedge \\ \quad b \in \mathbb{R} \wedge \\ \quad a < b \wedge \\ \quad h \in (\mathbb{R} \to \mathbb{R}) \wedge \\ \quad \forall j \in \mathbb{N}.\; \operatorname{is\text{-}antiderivable}(s(j), a, b) \wedge \\ \quad \forall x \in \mathbb{R}.\; (h(x) \simeq \operatorname{series\text{-}partial\text{-}sum}((\lambda j \in \mathbb{N}.\; s(j)(x)), n)) \Rightarrow \\ \quad \quad \operatorname{is\text{-}antiderivable}(h, a, b) \end{array}}}
\vnbprop{power-antiderivable}{\vnbstmt{\begin{array}{@{}l@{}} \forall m, a, b. \\ \quad ((m \in \mathbb{N}) \wedge ((a \in \mathbb{R}) \wedge ((b \in \mathbb{R}) \wedge (a < b)))) \Rightarrow  \\ \quad \quad \operatorname{is\text{-}antiderivable}((\lambda x \in \mathbb{R}.\; \operatorname{power}(x, m)), a, b) \end{array}}}
\vnbprop{poly-fn-antiderivable}{\vnbstmt{\begin{array}{@{}l@{}} \forall n \in \mathbb{N}, f, h, a, b. \\ \quad f \in (\mathbb{N} \to \mathbb{R}) \wedge \\ \quad a \in \mathbb{R} \wedge \\ \quad b \in \mathbb{R} \wedge \\ \quad a < b \wedge \\ \quad h \in (\mathbb{R} \to \mathbb{R}) \wedge \\ \quad \forall x \in \mathbb{R}. \\ \quad \quad (h(x) \simeq \operatorname{series\text{-}partial\text{-}sum}((\lambda k \in \mathbb{N}.\; (f(k) \cdot \operatorname{power}(x, k))), n)) \Rightarrow \\ \quad \quad \operatorname{is\text{-}antiderivable}(h, a, b) \end{array}}}

\vnbbookfile{theorem-library/c-int}{theorem-library/c-int}
\noindent{\small c-int.scm -- docs/calculus.pdf CHAPTER 4, equation (64): the DEFINITE INTEGRAL of a continuous function, `C-INT'.  Rung 4.}\par
\par\medskip\noindent\textbf{Definition} (\texttt{c-int}).\;$\int_{a}^{b} \varphi :=$
\begin{equation*}
(\iota\, \operatorname{v\_}.\;((\operatorname{v\_} \in \mathbb{R}) \wedge \exists \operatorname{f\_}.\; (\operatorname{is\text{-}antiderivative}(\operatorname{f\_}, \varphi, a, b) \wedge (\operatorname{v\_} = (\operatorname{f\_}(b) - \operatorname{f\_}(a))))))
\end{equation*}

\vnbprop{c-int-value}{\vnbstmt{\begin{array}{@{}l@{}} \forall f, i, a, b. \\ \quad \operatorname{is\text{-}antiderivative}(f, i, a, b) \Rightarrow (\int_{a}^{b} i = (f(b) - f(a))) \end{array}}}
\vnbprop{c-int-in-rr}{\vnbstmt{\forall i, a, b.\; (\operatorname{is\text{-}antiderivable}(i, a, b) \Rightarrow (\int_{a}^{b} i \in \mathbb{R}))}}
\vnbprop{c-int-add}{\vnbstmt{\begin{array}{@{}l@{}} \forall i, c, a, b. \\ \quad (\operatorname{is\text{-}antiderivable}(i, a, b) \wedge \operatorname{is\text{-}antiderivable}(c, a, b)) \Rightarrow  \\ \quad \quad \int_{a}^{b} (\lambda x \in \mathbb{R}.\; (i(x) + c(x))) = (\int_{a}^{b} i + \int_{a}^{b} c) \end{array}}}
\vnbprop{c-int-scale}{\vnbstmt{\begin{array}{@{}l@{}} \forall c, i, a, b. \\ \quad ((c \in \mathbb{R}) \wedge \operatorname{is\text{-}antiderivable}(i, a, b)) \Rightarrow  \\ \quad \quad \int_{a}^{b} (\lambda x \in \mathbb{R}.\; (c \cdot i(x))) = (c \cdot \int_{a}^{b} i) \end{array}}}

\vnbbookfile{theorem-library/c-int-oriented}{theorem-library/c-int-oriented}
\noindent{\small c-int-oriented.scm -- THE ORIENTED DEFINITE INTEGRAL, the sub-interval restriction of Def 4.6 (docs/calculus.pdf REMARK 4.9), and ADDITIVITY IN THE BOUNDS.}\par
\par\medskip\noindent\textbf{Definition} (\texttt{c-int-or}).\;$\int_{a}^{b} \varphi :=$
\begin{equation*}
\left(\text{if }\operatorname{is\text{-}antiderivable}(\varphi, a, b)\text{ then }\int_{a}^{b} \varphi\text{ else }\left(\text{if }\operatorname{is\text{-}antiderivable}(\varphi, b, a)\text{ then }-\int_{b}^{a} \varphi\text{ else }0\right)\right)
\end{equation*}

\vnbprop{rr-lt-asymm}{\vnbstmt{\forall a, b \in \mathbb{R}.\; ((a < b) \Rightarrow \neg (b < a))}}
\vnbprop{antiderivable-endpoints}{\vnbstmt{\begin{array}{@{}l@{}} \forall i, a, b. \\ \quad \operatorname{is\text{-}antiderivable}(i, a, b) \Rightarrow ((a \in \mathbb{R}) \wedge ((b \in \mathbb{R}) \wedge (a < b))) \end{array}}}
\vnbprop{antiderivable-not-reverse}{\vnbstmt{\begin{array}{@{}l@{}} \forall i, a, b. \\ \quad \operatorname{is\text{-}antiderivable}(i, a, b) \Rightarrow \neg \operatorname{is\text{-}antiderivable}(i, b, a) \end{array}}}
\vnbprop{antiderivable-not-degenerate}{\vnbstmt{\forall i, a.\; \neg \operatorname{is\text{-}antiderivable}(i, a, a)}}
\vnbprop{c-int-or-anti}{\vnbstmt{\begin{array}{@{}l@{}} \forall i, a, b. \\ \quad \operatorname{is\text{-}antiderivable}(i, a, b) \Rightarrow (\int_{a}^{b} i = \int_{a}^{b} i) \end{array}}}
\vnbprop{c-int-or-flip}{\vnbstmt{\begin{array}{@{}l@{}} \forall i, a, b. \\ \quad \operatorname{is\text{-}antiderivable}(i, a, b) \Rightarrow (\int_{b}^{a} i = -\int_{a}^{b} i) \end{array}}}
\vnbprop{c-int-or-none}{\vnbstmt{\begin{array}{@{}l@{}} \forall i, a, b. \\ \quad (\neg \operatorname{is\text{-}antiderivable}(i, a, b) \wedge \neg \operatorname{is\text{-}antiderivable}(i, b, a)) \Rightarrow  \\ \quad \quad \int_{a}^{b} i = 0 \end{array}}}
\vnbprop{c-int-or-degenerate}{\vnbstmt{\forall i, a.\; (\int_{a}^{a} i = 0)}}
\vnbprop{c-int-or-in-rr}{\vnbstmt{\forall i, a, b.\; (\int_{a}^{b} i \in \mathbb{R})}}
\vnbprop{c-int-or-reverse}{\vnbstmt{\forall i, a, b.\; (\int_{b}^{a} i = -\int_{a}^{b} i)}}
\vnbprop{antiderivative-subinterval}{\vnbstmt{\begin{array}{@{}l@{}} \forall f, i, a, b, o, c. \\ \quad (\operatorname{is\text{-}antiderivative}(f, i, a, b) \wedge ((o \in \mathbb{R}) \wedge ((c \in \mathbb{R}) \wedge ((a \leq o) \wedge ((o < c) \wedge (c \leq b)))))) \Rightarrow  \\ \quad \quad \operatorname{is\text{-}antiderivative}(f, i, o, c) \end{array}}}
\vnbprop{antiderivable-subinterval}{\vnbstmt{\begin{array}{@{}l@{}} \forall i, a, b, o, c. \\ \quad (\operatorname{is\text{-}antiderivable}(i, a, b) \wedge ((o \in \mathbb{R}) \wedge ((c \in \mathbb{R}) \wedge ((a \leq o) \wedge ((o < c) \wedge (c \leq b)))))) \Rightarrow  \\ \quad \quad \operatorname{is\text{-}antiderivable}(i, o, c) \end{array}}}
\vnbprop{c-int-or-value}{\vnbstmt{\begin{array}{@{}l@{}} \forall f, i, p, q, a, b. \\ \quad (\operatorname{is\text{-}antiderivative}(f, i, p, q) \wedge ((a \in \operatorname{ccint}(p, q)) \wedge (b \in \operatorname{ccint}(p, q)))) \Rightarrow  \\ \quad \quad \int_{a}^{b} i = (f(b) - f(a)) \end{array}}}
\vnbprop{c-int-or-additive}{\vnbstmt{\begin{array}{@{}l@{}} \forall i, p, q, a, b, c. \\ \quad (\operatorname{is\text{-}antiderivable}(i, p, q) \wedge ((a \in \operatorname{ccint}(p, q)) \wedge ((b \in \operatorname{ccint}(p, q)) \wedge (c \in \operatorname{ccint}(p, q))))) \Rightarrow  \\ \quad \quad \int_{a}^{c} i = (\int_{a}^{b} i + \int_{b}^{c} i) \end{array}}}
\vnbprop{c-int-add-bounds}{\vnbstmt{\begin{array}{@{}l@{}} \forall i, a, b, c. \\ \quad (\operatorname{is\text{-}antiderivable}(i, a, c) \wedge ((b \in \mathbb{R}) \wedge ((a < b) \wedge (b < c)))) \Rightarrow  \\ \quad \quad \int_{a}^{c} i = (\int_{a}^{b} i + \int_{b}^{c} i) \end{array}}}

\vnbbookfile{theorem-library/bernstein-ccint}{theorem-library/bernstein-ccint}
\noindent{\small bernstein-ccint.scm -- THEOREM 5.2 of docs/calculus.pdf moved off [0,1] and onto an arbitrary nondegenerate [a,b], by the affine change of variable}\par
\vnbprop{ccint-parts}{\vnbstmt{\forall o, i.; \forall t \in \operatorname{ccint}(o, i).\; ((t \in \mathbb{R}) \wedge ((o \leq t) \wedge (t \leq i)))}}
\vnbprop{ccint-affine-in}{\vnbstmt{\begin{array}{@{}l@{}} \forall a, b, t. \\ \quad (((a \in \mathbb{R}) \wedge ((b \in \mathbb{R}) \wedge (a < b))) \wedge (t \in \operatorname{ccint}(0, 1))) \Rightarrow  \\ \quad \quad (a + ((b - a) \cdot t)) \in \operatorname{ccint}(a, b) \end{array}}}
\vnbprop{ccint-coaffine-in}{\vnbstmt{\begin{array}{@{}l@{}} \forall a, b, x. \\ \quad (((a \in \mathbb{R}) \wedge ((b \in \mathbb{R}) \wedge (a < b))) \wedge (x \in \operatorname{ccint}(a, b))) \Rightarrow  \\ \quad \quad ((x - a) \cdot \operatorname{recip}((b - a))) \in \operatorname{ccint}(0, 1) \end{array}}}
\vnbprop{ccint-affine-inverse}{\vnbstmt{\begin{array}{@{}l@{}} \forall a, b, x. \\ \quad ((a \in \mathbb{R}) \wedge ((b \in \mathbb{R}) \wedge ((a < b) \wedge (x \in \mathbb{R})))) \Rightarrow  \\ \quad \quad (a + ((b - a) \cdot ((x - a) \cdot \operatorname{recip}((b - a))))) = x \end{array}}}
\vnbprop{affine-continuous-at}{\vnbstmt{\begin{array}{@{}l@{}} \forall c, m, t. \\ \quad ((c \in \mathbb{R}) \wedge ((m \in \mathbb{R}) \wedge (t \in \mathbb{R}))) \Rightarrow  \\ \quad \quad \operatorname{is\text{-}continuous\text{-}at}(\operatorname{rr\text{-}ms}, \operatorname{rr\text{-}ms}, (\lambda x \in \mathbb{R}.\; (c + (m \cdot x))), t) \end{array}}}
\vnbprop{bernstein-uniform-approximation-ccint}{\vnbstmt{\begin{array}{@{}l@{}} \forall n, a, b, p. \\ \quad n \in (\mathbb{R} \to \mathbb{R}) \wedge \\ \quad a \in \mathbb{R} \wedge \\ \quad b \in \mathbb{R} \wedge \\ \quad a < b \wedge \\ \quad \forall c \in \operatorname{ccint}(a, b).\; \operatorname{is\text{-}continuous\text{-}at}(\operatorname{rr\text{-}ms}, \operatorname{rr\text{-}ms}, n, c) \wedge \\ \quad \operatorname{pos\text{-}rr}(p) \Rightarrow \\ \quad \quad \exists \mathit{cap} \in \mathbb{N}. \forall d \in \mathbb{N}, c \in \operatorname{ccint}(a, b). \\ \quad \quad \quad (\mathit{cap} \leq d) \Rightarrow  \\ \quad \quad \quad \quad \lvert (n(c) - \operatorname{bernstein\text{-}poly}(\operatorname{compose}(n, (\lambda x \in \mathbb{R}.\; (a + ((b - a) \cdot x)))), d, ((c - a) \cdot \operatorname{recip}((b - a))))) \rvert \\ \quad \quad \quad \quad \leq\; p \end{array}}}

\vnbbookfile{theorem-library/rake-bolzano-weierstrass}{theorem-library/rake-bolzano-weierstrass}
\noindent{\small theorem-library/rake-bolzano-weierstrass.scm ==================================================================== The eps-GRID FINITE COVER of a closed interval -- batch 8-E, item (1) -- and the scaled archimedean property it needs.}\par
\vnbprop{rr-archimedean-multiple}{\vnbstmt{\forall v.; \forall \mathit{xv} \in \mathbb{R}.\; (\operatorname{pos\text{-}rr}(v) \Rightarrow \exists \mathit{nv} \in \mathbb{N}.\; (\mathit{xv} \leq (\mathit{nv} \cdot v)))}}
\vnbprop{ccint-eps-grid-cover-at}{\vnbstmt{\begin{array}{@{}l@{}} \forall v \in \mathbb{N}, \mathit{av}, \mathit{bv} \in \mathbb{R}, \mathit{ev}. \\ \quad (\operatorname{pos\text{-}rr}(\mathit{ev}) \wedge (\mathit{bv} \leq (\mathit{av} + (v \cdot \mathit{ev})))) \Rightarrow  \\ \quad \quad \exists \mathit{cv}. \\ \quad \quad \quad \operatorname{is\text{-}finite\text{-}cover}(\mathit{cv}, \operatorname{ccint}(\mathit{av}, \mathit{bv})) \wedge \\ \quad \quad \quad \forall \mathit{uv} \in \mathit{cv}, \mathit{pv}, \mathit{qv} \in \mathit{uv}. \\ \quad \quad \quad \quad (\mathit{pv} \in \mathbb{R}) \wedge ((\mathit{qv} \in \mathbb{R}) \wedge (\lvert (\mathit{pv} - \mathit{qv}) \rvert \leq \mathit{ev})) \end{array}}}
\vnbprop{ccint-eps-grid-cover}{\vnbstmt{\begin{array}{@{}l@{}} \forall v, \mathit{bv} \in \mathbb{R}, \mathit{ev}. \\ \quad \operatorname{pos\text{-}rr}(\mathit{ev}) \Rightarrow  \\ \quad \quad \exists \mathit{cv}. \\ \quad \quad \quad \operatorname{is\text{-}finite\text{-}cover}(\mathit{cv}, \operatorname{ccint}(v, \mathit{bv})) \wedge \\ \quad \quad \quad \forall \mathit{uv} \in \mathit{cv}, \mathit{pv}, \mathit{qv} \in \mathit{uv}. \\ \quad \quad \quad \quad (\mathit{pv} \in \mathbb{R}) \wedge ((\mathit{qv} \in \mathbb{R}) \wedge (\lvert (\mathit{pv} - \mathit{qv}) \rvert \leq \mathit{ev})) \end{array}}}
\vnbprop{ccint-grid-cover-seq}{\vnbstmt{\begin{array}{@{}l@{}} \forall v, \mathit{bv} \in \mathbb{R}, d. \\ \quad \forall \mathit{kv} \in \mathbb{N}.\; \operatorname{pos\text{-}rr}(d(\mathit{kv})) \Rightarrow  \\ \quad \quad \exists w. \forall \mathit{kv} \in \mathbb{N}. \\ \quad \quad \quad \operatorname{is\text{-}finite\text{-}cover}(w(\mathit{kv}), \operatorname{ccint}(v, \mathit{bv})) \wedge \\ \quad \quad \quad \forall \mathit{uv} \in w(\mathit{kv}), \mathit{pv}, \mathit{qv} \in \mathit{uv}. \\ \quad \quad \quad \quad (\mathit{pv} \in \mathbb{R}) \wedge ((\mathit{qv} \in \mathbb{R}) \wedge (\lvert (\mathit{pv} - \mathit{qv}) \rvert \leq d(\mathit{kv}))) \end{array}}}

\vnbbookfile{theorem-library/rake-bolzano-weierstrass-2}{theorem-library/rake-bolzano-weierstrass-2}
\noindent{\small NOTE 2026-09-20 (batch 12-B).  Where this file says the tree has no metric SUBSPACE structure, that is history: SUBSPACE-MS(s, A) and RESTRICT(f, A) are defined in structure-library/metric-subspace.scm, their laws are in theorem-library/metric-subspace-laws.scm and compact-subspace.scm, and}\par
\vnbprop{rr-interval-seq-has-convergent-subseq}{\vnbstmt{\begin{array}{@{}l@{}} \forall v, \mathit{hiv} \in \mathbb{R}, \mathit{hqv} \in (\mathbb{N} \to \operatorname{ccint}(v, \mathit{hiv})). \exists \mathit{phv}. \\ \quad (\operatorname{strictly\text{-}mono\text{-}nn}(\mathit{phv}) \wedge \exists \mathit{ptv} \in \mathbb{R}.\; \operatorname{converges\text{-}to}(\operatorname{rr\text{-}ms}, \operatorname{subseq}(\mathit{hqv}, \mathit{phv}), \mathit{ptv})) \end{array}}}
\vnbprop{rr-bolzano-weierstrass}{\vnbstmt{\begin{array}{@{}l@{}} \forall v \in (\mathbb{N} \to \mathbb{R}), \mathit{bdv} \in \mathbb{R}. \\ \quad \forall \mathit{iqv} \in \mathbb{N}.\; (\lvert v(\mathit{iqv}) \rvert \leq \mathit{bdv}) \Rightarrow  \\ \quad \quad \exists \mathit{phv}. \\ \quad \quad \quad (\operatorname{strictly\text{-}mono\text{-}nn}(\mathit{phv}) \wedge \exists \mathit{ptv} \in \mathbb{R}.\; \operatorname{converges\text{-}to}(\operatorname{rr\text{-}ms}, \operatorname{subseq}(v, \mathit{phv}), \mathit{ptv})) \end{array}}}
\vnbprop{bounded-block-converges}{\vnbstmt{\begin{array}{@{}l@{}} \forall v \in (\mathbb{N} \to \mathbb{R}), \mathit{bdv} \in \mathbb{R}, \mathit{jqv} \in \mathrm{Inf}(\mathbb{N}). \\ \quad \forall \mathit{iqv} \in \mathbb{N}.\; (\lvert v(\mathit{iqv}) \rvert \leq \mathit{bdv}) \Rightarrow  \\ \quad \quad \exists \mathit{bqv} \in \mathrm{Inf}(\mathbb{N}). \\ \quad \quad \quad (\mathit{bqv} \subseteq \mathit{jqv}) \wedge \exists \mathit{ptv} \in \mathbb{R}.\; \operatorname{converges\text{-}along}(\operatorname{rr\text{-}ms}, v, \mathit{bqv}, \mathit{ptv}) \end{array}}}

\vnbbookfile{theorem-library/rake-ascoli-diagonal}{theorem-library/rake-ascoli-diagonal}
\noindent{\small theorem-library/rake-ascoli-diagonal.scm ==================================================================== ASCOLI-ARZELA, the last stretch -- batch 11-F.  Two theorems:}\par
\vnbprop{ascoli-pointwise-diagonal}{\vnbstmt{\begin{array}{@{}l@{}} \forall s, m, q. \\ \quad (\operatorname{pointwise\text{-}bounded}(s, m) \wedge \operatorname{is\text{-}dense\text{-}seq}(s, q)) \Rightarrow  \\ \quad \quad \exists l.\; (\operatorname{strictly\text{-}mono\text{-}nn}(l) \wedge \operatorname{converges\text{-}on}(s, \operatorname{subseq}(m, l), q)) \end{array}}}
\vnbprop{ascoli-arzela-sequential}{\vnbstmt{\begin{array}{@{}l@{}} \forall s, m. \\ \quad \operatorname{is\text{-}compact}(s) \wedge \\ \quad \exists x.\; (x \in \operatorname{pts}(s)) \wedge \\ \quad m \in (\mathbb{N} \to (\operatorname{pts}(s) \to \mathbb{R})) \wedge \\ \quad \forall k \in \mathbb{N}.\; \operatorname{is\text{-}continuous}(s, \operatorname{rr\text{-}ms}, m(k)) \wedge \\ \quad \operatorname{is\text{-}equicontinuous}(s, \operatorname{rr\text{-}ms}, m) \wedge \\ \quad \operatorname{pointwise\text{-}bounded}(s, m) \Rightarrow \\ \quad \quad \exists l. \\ \quad \quad \quad \operatorname{strictly\text{-}mono\text{-}nn}(l) \wedge \\ \quad \quad \quad \exists g \in (\operatorname{pts}(s) \to \mathbb{R}). \\ \quad \quad \quad \quad \operatorname{is\text{-}continuous}(s, \operatorname{rr\text{-}ms}, g) \wedge \operatorname{converges\text{-}uniformly}(s, \operatorname{subseq}(m, l), g) \end{array}}}

\vnbbookfile{theorem-library/bernstein-antiderivable}{theorem-library/bernstein-antiderivable}
\noindent{\small bernstein-antiderivable.scm -- EVERY BERNSTEIN APPROXIMANT IS ANTIDERIVABLE.}\par
\vnbprop{bernstein-basis-degree-zero-diag}{\vnbstmt{\begin{array}{@{}l@{}} \forall x, k. \\ \quad ((x \in \mathbb{R}) \wedge ((k \in \mathbb{N}) \wedge (k = 0))) \Rightarrow (\operatorname{bernstein\text{-}basis}(0, x)(k) = 1) \end{array}}}
\vnbprop{bernstein-basis-degree-zero-off}{\vnbstmt{\begin{array}{@{}l@{}} \forall x, k. \\ \quad ((x \in \mathbb{R}) \wedge ((k \in \mathbb{N}) \wedge \neg (k = 0))) \Rightarrow (\operatorname{bernstein\text{-}basis}(0, x)(k) = 0) \end{array}}}
\vnbprop{bernstein-term-antiderivable}{\vnbstmt{\begin{array}{@{}l@{}} \forall n \in \mathbb{N}, a, b, k, j. \\ \quad (((a \in \mathbb{R}) \wedge ((b \in \mathbb{R}) \wedge (a < b))) \wedge ((k \in \mathbb{N}) \wedge (j \in \mathbb{N}))) \Rightarrow  \\ \quad \quad \operatorname{is\text{-}antiderivable}((\lambda x \in \mathbb{R}.\; (\operatorname{power}(x, j) \cdot \operatorname{bernstein\text{-}basis}(n, x)(k))), a, b) \end{array}}}
\vnbprop{bernstein-basis-antiderivable}{\vnbstmt{\begin{array}{@{}l@{}} \forall n, k, a, b. \\ \quad ((n \in \mathbb{N}) \wedge ((k \in \mathbb{N}) \wedge ((a \in \mathbb{R}) \wedge ((b \in \mathbb{R}) \wedge (a < b))))) \Rightarrow  \\ \quad \quad \operatorname{is\text{-}antiderivable}((\lambda x \in \mathbb{R}.\; \operatorname{bernstein\text{-}basis}(n, x)(k)), a, b) \end{array}}}
\vnbprop{bernstein-basis-scaled-antiderivable}{\vnbstmt{\begin{array}{@{}l@{}} \forall c, n, k, a, b. \\ \quad ((c \in \mathbb{R}) \wedge ((n \in \mathbb{N}) \wedge ((k \in \mathbb{N}) \wedge ((a \in \mathbb{R}) \wedge ((b \in \mathbb{R}) \wedge (a < b)))))) \Rightarrow  \\ \quad \quad \operatorname{is\text{-}antiderivable}((\lambda x \in \mathbb{R}.\; (c \cdot \operatorname{bernstein\text{-}basis}(n, x)(k))), a, b) \end{array}}}
\vnbprop{bernstein-poly-antiderivable}{\vnbstmt{\begin{array}{@{}l@{}} \forall f, n, a, b. \\ \quad ((f \in (\mathbb{R} \to \mathbb{R})) \wedge ((n \in \mathbb{N}) \wedge (\neg (n = 0) \wedge ((a \in \mathbb{R}) \wedge ((b \in \mathbb{R}) \wedge (a < b)))))) \Rightarrow  \\ \quad \quad \operatorname{is\text{-}antiderivable}((\lambda x \in \mathbb{R}.\; \operatorname{bernstein\text{-}poly}(f, n, x)), a, b) \end{array}}}

\vnbbookfile{theorem-library/antiderivable-affine-subst}{theorem-library/antiderivable-affine-subst}
\noindent{\small antiderivable-affine-subst.scm -- ANTIDERIVABILITY SURVIVES AN AFFINE CHANGE OF VARIABLE, and with it the rung-3 form of the Bernstein approximants.}\par
\vnbprop{antiderivable-affine-subst}{\vnbstmt{\begin{array}{@{}l@{}} \forall i, k, m, p, q, a, b. \\ \quad i \in (\mathbb{R} \to \mathbb{R}) \wedge \\ \quad k \in \mathbb{R} \wedge \\ \quad m \in \mathbb{R} \wedge \\ \quad \neg (m = 0) \wedge \\ \quad p \in \mathbb{R} \wedge \\ \quad q \in \mathbb{R} \wedge \\ \quad p < q \wedge \\ \quad \operatorname{is\text{-}antiderivable}(i, a, b) \wedge \\ \quad \forall t \in \operatorname{ccint}(p, q).\; ((k + (m \cdot t)) \in \operatorname{ccint}(a, b)) \wedge \\ \quad \forall t \in \mathbb{R}.\; (((p < t) \wedge (t < q)) \Rightarrow ((a < (k + (m \cdot t))) \wedge ((k + (m \cdot t)) < b))) \Rightarrow \\ \quad \quad \operatorname{is\text{-}antiderivable}((\lambda x \in \mathbb{R}.\; i((k + (m \cdot x)))), p, q) \end{array}}}
\vnbprop{antiderivable-affine-subst-pos}{\vnbstmt{\begin{array}{@{}l@{}} \forall i, k, m, p, q. \\ \quad i \in (\mathbb{R} \to \mathbb{R}) \wedge \\ \quad k \in \mathbb{R} \wedge \\ \quad m \in \mathbb{R} \wedge \\ \quad 0 < m \wedge \\ \quad p \in \mathbb{R} \wedge \\ \quad q \in \mathbb{R} \wedge \\ \quad p < q \wedge \\ \quad \operatorname{is\text{-}antiderivable}(i, (k + (m \cdot p)), (k + (m \cdot q))) \Rightarrow \\ \quad \quad \operatorname{is\text{-}antiderivable}((\lambda x \in \mathbb{R}.\; i((k + (m \cdot x)))), p, q) \end{array}}}
\vnbprop{antiderivable-affine-subst-pos-transfer}{\vnbstmt{\begin{array}{@{}l@{}} \forall i, a, k, m, p, q. \\ \quad i \in (\mathbb{R} \to \mathbb{R}) \wedge \\ \quad a \in (\mathbb{R} \to \mathbb{R}) \wedge \\ \quad k \in \mathbb{R} \wedge \\ \quad m \in \mathbb{R} \wedge \\ \quad 0 < m \wedge \\ \quad p \in \mathbb{R} \wedge \\ \quad q \in \mathbb{R} \wedge \\ \quad p < q \wedge \\ \quad \operatorname{is\text{-}antiderivable}(i, (k + (m \cdot p)), (k + (m \cdot q))) \wedge \\ \quad \forall x \in \mathbb{R}.\; (a(x) \simeq i((k + (m \cdot x)))) \Rightarrow \\ \quad \quad \operatorname{is\text{-}antiderivable}(a, p, q) \end{array}}}
\vnbprop{bernstein-poly-unit-coordinate-antiderivable}{\vnbstmt{\begin{array}{@{}l@{}} \forall g, n, a, b. \\ \quad ((g \in (\mathbb{R} \to \mathbb{R})) \wedge ((n \in \mathbb{N}) \wedge (\neg (n = 0) \wedge ((a \in \mathbb{R}) \wedge ((b \in \mathbb{R}) \wedge (a < b)))))) \Rightarrow  \\ \quad \quad \operatorname{is\text{-}antiderivable}((\lambda x \in \mathbb{R}.\; \operatorname{bernstein\text{-}poly}(g, n, ((x - a) \cdot \operatorname{recip}((b - a))))), a, b) \end{array}}}
\vnbprop{bernstein-ccint-approximant-antiderivable}{\vnbstmt{\begin{array}{@{}l@{}} \forall f, n, a, b. \\ \quad ((f \in (\mathbb{R} \to \mathbb{R})) \wedge ((n \in \mathbb{N}) \wedge (\neg (n = 0) \wedge ((a \in \mathbb{R}) \wedge ((b \in \mathbb{R}) \wedge (a < b)))))) \Rightarrow  \\ \quad \quad \operatorname{is\text{-}antiderivable}((\lambda x \in \mathbb{R}.\; \\ \quad \quad \quad \quad \operatorname{bernstein\text{-}poly}(\operatorname{compose}(f, (\lambda c \in \mathbb{R}.\; (a + ((b - a) \cdot c)))), n, ((x - a) \cdot \operatorname{recip}((b - a))))), \\ \quad \quad \quad a, b) \end{array}}}

\vnbbookfile{theorem-library/continuous-antiderivable}{theorem-library/continuous-antiderivable}
\noindent{\small continuous-antiderivable.scm -- RUNG 3, AND THE INTEGRAL THAT FOLLOWS.}\par
\vnbprop{continuous-is-antiderivable}{\vnbstmt{\begin{array}{@{}l@{}} \forall f, a, b. \\ \quad f \in (\mathbb{R} \to \mathbb{R}) \wedge \\ \quad a \in \mathbb{R} \wedge \\ \quad b \in \mathbb{R} \wedge \\ \quad a < b \wedge \\ \quad \forall x \in \operatorname{ccint}(a, b).\; \operatorname{is\text{-}continuous\text{-}at}(\operatorname{rr\text{-}ms}, \operatorname{rr\text{-}ms}, f, x) \Rightarrow \\ \quad \quad \operatorname{is\text{-}antiderivable}(f, a, b) \end{array}}}
\vnbprop{c-int-defined-for-continuous}{\vnbstmt{\begin{array}{@{}l@{}} \forall f, a, b. \\ \quad f \in (\mathbb{R} \to \mathbb{R}) \wedge \\ \quad a \in \mathbb{R} \wedge \\ \quad b \in \mathbb{R} \wedge \\ \quad a < b \wedge \\ \quad \forall x \in \operatorname{ccint}(a, b).\; \operatorname{is\text{-}continuous\text{-}at}(\operatorname{rr\text{-}ms}, \operatorname{rr\text{-}ms}, f, x) \Rightarrow \\ \quad \quad \int_{a}^{b} f \in \mathbb{R} \end{array}}}
\vnbprop{c-int-value-continuous}{\vnbstmt{\begin{array}{@{}l@{}} \forall f, a, b. \\ \quad f \in (\mathbb{R} \to \mathbb{R}) \wedge \\ \quad a \in \mathbb{R} \wedge \\ \quad b \in \mathbb{R} \wedge \\ \quad a < b \wedge \\ \quad \forall x \in \operatorname{ccint}(a, b).\; \operatorname{is\text{-}continuous\text{-}at}(\operatorname{rr\text{-}ms}, \operatorname{rr\text{-}ms}, f, x) \Rightarrow \\ \quad \quad \exists g.\; (\operatorname{is\text{-}antiderivative}(g, f, a, b) \wedge (\int_{a}^{b} f = (g(b) - g(a)))) \end{array}}}
\vnbprop{c-int-add-continuous}{\vnbstmt{\begin{array}{@{}l@{}} \forall i, c, a, b. \\ \quad i \in (\mathbb{R} \to \mathbb{R}) \wedge \\ \quad c \in (\mathbb{R} \to \mathbb{R}) \wedge \\ \quad a \in \mathbb{R} \wedge \\ \quad b \in \mathbb{R} \wedge \\ \quad a < b \wedge \\ \quad \forall d \in \operatorname{ccint}(a, b).\; \operatorname{is\text{-}continuous\text{-}at}(\operatorname{rr\text{-}ms}, \operatorname{rr\text{-}ms}, i, d) \wedge \\ \quad \forall d \in \operatorname{ccint}(a, b).\; \operatorname{is\text{-}continuous\text{-}at}(\operatorname{rr\text{-}ms}, \operatorname{rr\text{-}ms}, c, d) \Rightarrow \\ \quad \quad \int_{a}^{b} (\lambda x \in \mathbb{R}.\; (i(x) + c(x))) = (\int_{a}^{b} i + \int_{a}^{b} c) \end{array}}}
\vnbprop{c-int-scale-continuous}{\vnbstmt{\begin{array}{@{}l@{}} \forall c, i, a, b. \\ \quad c \in \mathbb{R} \wedge \\ \quad i \in (\mathbb{R} \to \mathbb{R}) \wedge \\ \quad a \in \mathbb{R} \wedge \\ \quad b \in \mathbb{R} \wedge \\ \quad a < b \wedge \\ \quad \forall d \in \operatorname{ccint}(a, b).\; \operatorname{is\text{-}continuous\text{-}at}(\operatorname{rr\text{-}ms}, \operatorname{rr\text{-}ms}, i, d) \Rightarrow \\ \quad \quad \int_{a}^{b} (\lambda x \in \mathbb{R}.\; (c \cdot i(x))) = (c \cdot \int_{a}^{b} i) \end{array}}}

\vnbbookfile{theorem-library/c-int-change-of-variable}{theorem-library/c-int-change-of-variable}
\noindent{\small c-int-change-of-variable.scm -- CHANGE OF VARIABLES FOR `C-INT', in the INCREASING case, docs/calculus.pdf's substitution rule:}\par
\vnbprop{change-of-variable-antiderivative}{\vnbstmt{\begin{array}{@{}l@{}} \forall w, i, p, r, a, q. \\ \quad p \in (\mathbb{R} \to \mathbb{R}) \wedge \\ \quad r \in (\mathbb{R} \to \mathbb{R}) \wedge \\ \quad a \in \mathbb{R} \wedge \\ \quad q \in \mathbb{R} \wedge \\ \quad a < q \wedge \\ \quad \operatorname{is\text{-}antiderivative}(w, i, p(a), p(q)) \wedge \\ \quad \forall u \in \operatorname{ccint}(a, q).\; \operatorname{is\text{-}continuous\text{-}at}(\operatorname{rr\text{-}ms}, \operatorname{rr\text{-}ms}, p, u) \wedge \\ \quad \forall b \in \mathbb{R}.\; (((a < b) \wedge (b < q)) \Rightarrow \operatorname{is\text{-}diff\text{-}at}(p, b, r(b))) \wedge \\ \quad \forall u \in \operatorname{ccint}(a, q).\; (p(u) \in \operatorname{ccint}(p(a), p(q))) \wedge \\ \quad \forall b \in \mathbb{R}.\; (((a < b) \wedge (b < q)) \Rightarrow ((p(a) < p(b)) \wedge (p(b) < p(q)))) \Rightarrow \\ \quad \quad \operatorname{is\text{-}antiderivative}(\operatorname{compose}(w, p), (\lambda x \in \mathbb{R}.\; (i(p(x)) \cdot r(x))), a, q) \end{array}}}
\vnbprop{change-of-variable-antiderivable}{\vnbstmt{\begin{array}{@{}l@{}} \forall i, p, r, a, q. \\ \quad p \in (\mathbb{R} \to \mathbb{R}) \wedge \\ \quad r \in (\mathbb{R} \to \mathbb{R}) \wedge \\ \quad a \in \mathbb{R} \wedge \\ \quad q \in \mathbb{R} \wedge \\ \quad a < q \wedge \\ \quad \operatorname{is\text{-}antiderivable}(i, p(a), p(q)) \wedge \\ \quad \forall u \in \operatorname{ccint}(a, q).\; \operatorname{is\text{-}continuous\text{-}at}(\operatorname{rr\text{-}ms}, \operatorname{rr\text{-}ms}, p, u) \wedge \\ \quad \forall w \in \mathbb{R}.\; (((a < w) \wedge (w < q)) \Rightarrow \operatorname{is\text{-}diff\text{-}at}(p, w, r(w))) \wedge \\ \quad \forall u \in \operatorname{ccint}(a, q).\; (p(u) \in \operatorname{ccint}(p(a), p(q))) \wedge \\ \quad \forall w \in \mathbb{R}.\; (((a < w) \wedge (w < q)) \Rightarrow ((p(a) < p(w)) \wedge (p(w) < p(q)))) \Rightarrow \\ \quad \quad \operatorname{is\text{-}antiderivable}((\lambda x \in \mathbb{R}.\; (i(p(x)) \cdot r(x))), a, q) \end{array}}}
\vnbprop{c-int-change-of-variable}{\vnbstmt{\begin{array}{@{}l@{}} \forall i, p, r, a, q. \\ \quad p \in (\mathbb{R} \to \mathbb{R}) \wedge \\ \quad r \in (\mathbb{R} \to \mathbb{R}) \wedge \\ \quad a \in \mathbb{R} \wedge \\ \quad q \in \mathbb{R} \wedge \\ \quad a < q \wedge \\ \quad \operatorname{is\text{-}antiderivable}(i, p(a), p(q)) \wedge \\ \quad \forall u \in \operatorname{ccint}(a, q).\; \operatorname{is\text{-}continuous\text{-}at}(\operatorname{rr\text{-}ms}, \operatorname{rr\text{-}ms}, p, u) \wedge \\ \quad \forall w \in \mathbb{R}.\; (((a < w) \wedge (w < q)) \Rightarrow \operatorname{is\text{-}diff\text{-}at}(p, w, r(w))) \wedge \\ \quad \forall u \in \operatorname{ccint}(a, q).\; (p(u) \in \operatorname{ccint}(p(a), p(q))) \wedge \\ \quad \forall w \in \mathbb{R}.\; (((a < w) \wedge (w < q)) \Rightarrow ((p(a) < p(w)) \wedge (p(w) < p(q)))) \Rightarrow \\ \quad \quad \int_{a}^{q} (\lambda x \in \mathbb{R}.\; (i(p(x)) \cdot r(x))) = \int_{p(a)}^{p(q)} i \end{array}}}
\vnbprop{c-int-change-of-variable-transfer}{\vnbstmt{\begin{array}{@{}l@{}} \forall i, a, p, r, b, q. \\ \quad p \in (\mathbb{R} \to \mathbb{R}) \wedge \\ \quad r \in (\mathbb{R} \to \mathbb{R}) \wedge \\ \quad b \in \mathbb{R} \wedge \\ \quad q \in \mathbb{R} \wedge \\ \quad b < q \wedge \\ \quad a \in (\mathbb{R} \to \mathbb{R}) \wedge \\ \quad \operatorname{is\text{-}antiderivable}(i, p(b), p(q)) \wedge \\ \quad \forall u \in \operatorname{ccint}(b, q).\; \operatorname{is\text{-}continuous\text{-}at}(\operatorname{rr\text{-}ms}, \operatorname{rr\text{-}ms}, p, u) \wedge \\ \quad \forall w \in \mathbb{R}.\; (((b < w) \wedge (w < q)) \Rightarrow \operatorname{is\text{-}diff\text{-}at}(p, w, r(w))) \wedge \\ \quad \forall u \in \operatorname{ccint}(b, q).\; (p(u) \in \operatorname{ccint}(p(b), p(q))) \wedge \\ \quad \forall w \in \mathbb{R}.\; (((b < w) \wedge (w < q)) \Rightarrow ((p(b) < p(w)) \wedge (p(w) < p(q)))) \wedge \\ \quad \forall x \in \mathbb{R}.\; (a(x) \simeq (i(p(x)) \cdot r(x))) \Rightarrow \\ \quad \quad \int_{b}^{q} a = \int_{p(b)}^{p(q)} i \end{array}}}

\vnbbookfile{theorem-library/log}{theorem-library/log}
\noindent{\small theorem-library/log.scm -- THE LOGARITHM, DEFINED AS AN INTEGRAL, and its basic laws.  The one that costs anything is the functional equation, and what it costs is a CHANGE OF VARIABLES.}\par
\par\medskip\noindent\textbf{Definition} (\texttt{log}).\;$\log\left(\operatorname{x\_}\right) :=$ $\int_{1}^{\operatorname{x\_}} (\lambda \operatorname{z\_} \in \mathbb{R}.\; {\operatorname{z\_}}^{\ast\!-1})$
\vnbprop{log-unfold}{\vnbstmt{\forall x.\; (\log\left(x\right) \simeq \int_{1}^{x} (\lambda z \in \mathbb{R}.\; {z}^{\ast\!-1}))}}
\vnbprop{log-in-rr}{\vnbstmt{\forall x.\; (\log\left(x\right) \in \mathbb{R})}}
\vnbprop{log-lam-in-fun}{\vnbstmt{(\lambda x \in \mathbb{R}.\; \log\left(x\right)) \in (\mathbb{R} \to \mathbb{R})}}
\vnbprop{log-one}{\vnbstmt{\log\left(1\right) = 0}}
\vnbprop{recip-star-antiderivable-pos}{\vnbstmt{\begin{array}{@{}l@{}} \forall a, b. \\ \quad ((a \in \mathbb{R}) \wedge ((b \in \mathbb{R}) \wedge ((0 < a) \wedge (a < b)))) \Rightarrow  \\ \quad \quad \operatorname{is\text{-}antiderivable}((\lambda z \in \mathbb{R}.\; {z}^{\ast\!-1}), a, b) \end{array}}}
\vnbprop{recip-star-scale}{\vnbstmt{\begin{array}{@{}l@{}} \forall m, u. \\ \quad ((m \in \mathbb{R}) \wedge ((u \in \mathbb{R}) \wedge \neg (m = 0))) \Rightarrow (({(m \cdot u)}^{\ast\!-1} \cdot m) = {u}^{\ast\!-1}) \end{array}}}
\vnbprop{c-int-or-scale-ordered}{\vnbstmt{\begin{array}{@{}l@{}} \forall m, a, b. \\ \quad ((m \in \mathbb{R}) \wedge ((a \in \mathbb{R}) \wedge ((b \in \mathbb{R}) \wedge ((0 < m) \wedge ((0 < a) \wedge (a < b)))))) \Rightarrow  \\ \quad \quad \int_{(m \cdot a)}^{(m \cdot b)} (\lambda z \in \mathbb{R}.\; {z}^{\ast\!-1}) \\ \quad \quad =\; \int_{a}^{b} (\lambda z \in \mathbb{R}.\; {z}^{\ast\!-1}) \end{array}}}
\vnbprop{c-int-or-scale}{\vnbstmt{\begin{array}{@{}l@{}} \forall m, a, b. \\ \quad ((m \in \mathbb{R}) \wedge ((a \in \mathbb{R}) \wedge ((b \in \mathbb{R}) \wedge ((0 < m) \wedge ((0 < a) \wedge (0 < b)))))) \Rightarrow  \\ \quad \quad \int_{(m \cdot a)}^{(m \cdot b)} (\lambda z \in \mathbb{R}.\; {z}^{\ast\!-1}) \\ \quad \quad =\; \int_{a}^{b} (\lambda z \in \mathbb{R}.\; {z}^{\ast\!-1}) \end{array}}}
\vnbprop{recip-star-int-additive}{\vnbstmt{\begin{array}{@{}l@{}} \forall u, v, w. \\ \quad ((u \in \mathbb{R}) \wedge ((v \in \mathbb{R}) \wedge ((w \in \mathbb{R}) \wedge ((0 < u) \wedge ((0 < v) \wedge (0 < w)))))) \Rightarrow  \\ \quad \quad \int_{u}^{w} (\lambda z \in \mathbb{R}.\; {z}^{\ast\!-1}) \\ \quad \quad =\; \int_{u}^{v} (\lambda z \in \mathbb{R}.\; {z}^{\ast\!-1}) + \int_{v}^{w} (\lambda z \in \mathbb{R}.\; {z}^{\ast\!-1}) \end{array}}}
\vnbprop{log-diff}{\vnbstmt{\begin{array}{@{}l@{}} \forall u, v. \\ \quad ((u \in \mathbb{R}) \wedge ((v \in \mathbb{R}) \wedge ((0 < u) \wedge (0 < v)))) \Rightarrow  \\ \quad \quad \int_{u}^{v} (\lambda z \in \mathbb{R}.\; {z}^{\ast\!-1}) = (\log\left(v\right) - \log\left(u\right)) \end{array}}}
\vnbprop{log-mul}{\vnbstmt{\begin{array}{@{}l@{}} \forall x, y. \\ \quad ((x \in \mathbb{R}) \wedge ((y \in \mathbb{R}) \wedge ((0 < x) \wedge (0 < y)))) \Rightarrow (\log\left((x \cdot y)\right) = (\log\left(x\right) + \log\left(y\right))) \end{array}}}
\vnbprop{log-strictly-increasing}{\vnbstmt{\forall u, v.\; (((u \in \mathbb{R}) \wedge ((v \in \mathbb{R}) \wedge ((0 < u) \wedge (u < v)))) \Rightarrow (\log\left(u\right) < \log\left(v\right)))}}
\vnbprop{log-neg-below-one}{\vnbstmt{\forall x \in \mathbb{R}.\; (((0 < x) \wedge (x < 1)) \Rightarrow (\log\left(x\right) < 0))}}
\vnbprop{log-pos-above-one}{\vnbstmt{\forall x \in \mathbb{R}.\; ((1 < x) \Rightarrow (0 < \log\left(x\right)))}}
\vnbprop{log-recip}{\vnbstmt{\forall x \in \mathbb{R}.\; ((0 < x) \Rightarrow (\log\left({x}^{\ast\!-1}\right) = -\log\left(x\right)))}}
\vnbprop{log-deriv}{\vnbstmt{\begin{array}{@{}l@{}} \forall c \in \mathbb{R}. \\ \quad (0 < c) \Rightarrow \operatorname{is\text{-}diff\text{-}at}((\lambda x \in \mathbb{R}.\; \log\left(x\right)), c, {c}^{\ast\!-1}) \end{array}}}

\vnbbookfile{theorem-library/r-exp}{theorem-library/r-exp}
\noindent{\small theorem-library/r-exp.scm -- THE REAL EXPONENTIAL, DEFINED AS THE INVERSE OF THE LOGARITHM, and its laws.}\par
\par\medskip\noindent\textbf{Definition} (\texttt{r-exp}).\;$\exp\left(\operatorname{y\_}\right) :=$ $(\iota\, \operatorname{x\_}.\;((\operatorname{x\_} \in \mathbb{R}) \wedge ((0 < \operatorname{x\_}) \wedge (\log\left(\operatorname{x\_}\right) = \operatorname{y\_}))))$
\vnbprop{log-two-pos}{\vnbstmt{0 < \log\left(2\right)}}
\vnbprop{log-unbounded-above}{\vnbstmt{\forall n \in \mathbb{N}.\; \exists b \in \mathbb{R}.\; ((0 < b) \wedge ((n \cdot \log\left(2\right)) \leq \log\left(b\right)))}}
\vnbprop{log-above-any}{\vnbstmt{\forall y \in \mathbb{R}.\; \exists b \in \mathbb{R}.\; ((0 < b) \wedge (y \leq \log\left(b\right)))}}
\vnbprop{log-below-any}{\vnbstmt{\forall y \in \mathbb{R}.\; \exists a \in \mathbb{R}.\; ((0 < a) \wedge (\log\left(a\right) \leq y))}}
\vnbprop{log-surjective}{\vnbstmt{\forall y \in \mathbb{R}.\; \exists x \in \mathbb{R}.\; ((0 < x) \wedge (\log\left(x\right) = y))}}
\vnbprop{log-injective-pos}{\vnbstmt{\begin{array}{@{}l@{}} \forall u, v. \\ \quad ((u \in \mathbb{R}) \wedge ((v \in \mathbb{R}) \wedge ((0 < u) \wedge ((0 < v) \wedge (\log\left(u\right) = \log\left(v\right)))))) \Rightarrow (u = v) \end{array}}}
\vnbprop{r-exp-prop}{\vnbstmt{\forall y \in \mathbb{R}.\; ((\exp\left(y\right) \in \mathbb{R}) \wedge ((0 < \exp\left(y\right)) \wedge (\log\left(\exp\left(y\right)\right) = y)))}}
\vnbprop{r-exp-in-rr}{\vnbstmt{\forall y \in \mathbb{R}.\; (\exp\left(y\right) \in \mathbb{R})}}
\vnbprop{r-exp-pos}{\vnbstmt{\forall y \in \mathbb{R}.\; (0 < \exp\left(y\right))}}
\vnbprop{log-r-exp}{\vnbstmt{\forall y \in \mathbb{R}.\; (\log\left(\exp\left(y\right)\right) = y)}}
\vnbprop{r-exp-char}{\vnbstmt{\begin{array}{@{}l@{}} \forall y, x. \\ \quad ((y \in \mathbb{R}) \wedge ((x \in \mathbb{R}) \wedge ((0 < x) \wedge (\log\left(x\right) = y)))) \Rightarrow (\exp\left(y\right) = x) \end{array}}}
\vnbprop{r-exp-lam-in-fun}{\vnbstmt{(\lambda y \in \mathbb{R}.\; \exp\left(y\right)) \in (\mathbb{R} \to \mathbb{R})}}
\vnbprop{r-exp-zero}{\vnbstmt{\exp\left(0\right) = 1}}
\vnbprop{r-exp-log}{\vnbstmt{\forall x \in \mathbb{R}.\; ((0 < x) \Rightarrow (\exp\left(\log\left(x\right)\right) = x))}}
\vnbprop{r-exp-add}{\vnbstmt{\forall u, v.\; (((u \in \mathbb{R}) \wedge (v \in \mathbb{R})) \Rightarrow (\exp\left((u + v)\right) = (\exp\left(u\right) \cdot \exp\left(v\right))))}}
\vnbprop{r-exp-strictly-increasing}{\vnbstmt{\forall u, v.\; (((u \in \mathbb{R}) \wedge ((v \in \mathbb{R}) \wedge (u < v))) \Rightarrow (\exp\left(u\right) < \exp\left(v\right)))}}
\vnbprop{r-exp-mono}{\vnbstmt{\forall u, v.\; (((u \in \mathbb{R}) \wedge ((v \in \mathbb{R}) \wedge (u \leq v))) \Rightarrow (\exp\left(u\right) \leq \exp\left(v\right)))}}
\vnbprop{r-exp-neg}{\vnbstmt{\forall y \in \mathbb{R}.\; (\exp\left(-y\right) = {\exp\left(y\right)}^{\ast\!-1})}}
\vnbprop{r-exp-log-not-global}{\vnbstmt{\neg \forall x \in \mathbb{R}.\; (\exp\left(\log\left(x\right)\right) = x)}}
\vnbprop{r-exp-continuous-at}{\vnbstmt{\begin{array}{@{}l@{}} \forall t \in \mathbb{R}. \\ \quad \operatorname{is\text{-}continuous\text{-}at}(\operatorname{rr\text{-}ms}, \operatorname{rr\text{-}ms}, (\lambda y \in \mathbb{R}.\; \exp\left(y\right)), t) \end{array}}}
\vnbprop{rr-recip-recip-pos}{\vnbstmt{\forall c \in \mathbb{R}.\; ((0 < c) \Rightarrow (\operatorname{recip}(\operatorname{recip}(c)) = c))}}
\vnbprop{deriv-right-inverse}{\vnbstmt{\begin{array}{@{}l@{}} \forall m, \mathit{gm}, \mathit{am}, v. \\ \quad \operatorname{is\text{-}diff\text{-}at}(m, \mathit{am}, v) \wedge \\ \quad \neg (v = 0) \wedge \\ \quad \mathit{gm} \in (\mathbb{R} \to \mathbb{R}) \wedge \\ \quad \forall z \in \mathbb{R}.\; (m(\operatorname{gm}(z)) = z) \wedge \\ \quad \operatorname{gm}(m(\mathit{am})) = \mathit{am} \wedge \\ \quad \operatorname{is\text{-}continuous\text{-}at}(\operatorname{rr\text{-}ms}, \operatorname{rr\text{-}ms}, \mathit{gm}, m(\mathit{am})) \Rightarrow \\ \quad \quad \operatorname{is\text{-}diff\text{-}at}(\mathit{gm}, m(\mathit{am}), \operatorname{recip}(v)) \end{array}}}
\vnbprop{r-exp-diff-at-log}{\vnbstmt{\begin{array}{@{}l@{}} \forall c, w. \\ \quad ((c \in \mathbb{R}) \wedge ((w \in \mathbb{R}) \wedge ((0 < c) \wedge (w = \log\left(c\right))))) \Rightarrow  \\ \quad \quad \operatorname{is\text{-}diff\text{-}at}((\lambda y \in \mathbb{R}.\; \exp\left(y\right)), w, c) \end{array}}}
\vnbprop{r-exp-deriv}{\vnbstmt{\forall t \in \mathbb{R}.\; \operatorname{is\text{-}diff\text{-}at}((\lambda y \in \mathbb{R}.\; \exp\left(y\right)), t, \exp\left(t\right))}}

\vnbbookfile{theorem-library/rpow-star}{theorem-library/rpow-star}
\noindent{\small theorem-library/rpow-star.scm -- THE REAL POWER, DEFINED, with the base-zero conventions carried:}\par
\par\medskip\noindent\textbf{Definition} (\texttt{rpow-star}).\;${\operatorname{x\_}}^{\operatorname{s\_}} :=$ $\left(\text{if }(0 < \operatorname{x\_})\text{ then }\exp\left((\operatorname{s\_} \cdot \log\left(\operatorname{x\_}\right))\right)\text{ else }\left(\text{if }(0 < \operatorname{s\_})\text{ then }0\text{ else }1\right)\right)$
\vnbprop{rpow-star-unfold}{\vnbstmt{\forall x, s.\; ({x}^{s} \simeq \left(\text{if }(0 < x)\text{ then }\exp\left((s \cdot \log\left(x\right))\right)\text{ else }\left(\text{if }(0 < s)\text{ then }0\text{ else }1\right)\right))}}
\vnbprop{rpow-star-value}{\vnbstmt{\forall x.; \forall s \in \mathbb{R}.\; ((0 < x) \Rightarrow ({x}^{s} = \exp\left((s \cdot \log\left(x\right))\right)))}}
\vnbprop{rpow-star-zero-base}{\vnbstmt{\forall s \in \mathbb{R}.\; ((0 < s) \Rightarrow ({0}^{s} = 0))}}
\vnbprop{rpow-star-zero-zero}{\vnbstmt{{0}^{0} = 1}}
\vnbprop{rpow-star-in-rr}{\vnbstmt{\forall x.; \forall s \in \mathbb{R}.\; ((0 < x) \Rightarrow ({x}^{s} \in \mathbb{R}))}}
\vnbprop{rpow-star-pos}{\vnbstmt{\forall x.; \forall s \in \mathbb{R}.\; ((0 < x) \Rightarrow (0 < {x}^{s}))}}
\vnbprop{log-rpow-star}{\vnbstmt{\forall x.; \forall s \in \mathbb{R}.\; ((0 < x) \Rightarrow (\log\left({x}^{s}\right) = (s \cdot \log\left(x\right))))}}
\vnbprop{rpow-star-zero}{\vnbstmt{\forall x \in \mathbb{R}.\; ((0 < x) \Rightarrow ({x}^{0} = 1))}}
\vnbprop{rpow-star-one}{\vnbstmt{\forall x \in \mathbb{R}.\; ((0 < x) \Rightarrow ({x}^{1} = x))}}
\vnbprop{rpow-star-add}{\vnbstmt{\begin{array}{@{}l@{}} \forall x, s, t. \\ \quad ((x \in \mathbb{R}) \wedge ((s \in \mathbb{R}) \wedge ((t \in \mathbb{R}) \wedge (0 < x)))) \Rightarrow  \\ \quad \quad {x}^{(s + t)} = ({x}^{s} \cdot {x}^{t}) \end{array}}}
\vnbprop{rpow-star-mul-base}{\vnbstmt{\begin{array}{@{}l@{}} \forall x, y, s. \\ \quad ((x \in \mathbb{R}) \wedge ((y \in \mathbb{R}) \wedge ((s \in \mathbb{R}) \wedge ((0 < x) \wedge (0 < y))))) \Rightarrow  \\ \quad \quad {(x \cdot y)}^{s} = ({x}^{s} \cdot {y}^{s}) \end{array}}}
\vnbprop{rpow-star-pow}{\vnbstmt{\begin{array}{@{}l@{}} \forall x, s, t. \\ \quad ((x \in \mathbb{R}) \wedge ((s \in \mathbb{R}) \wedge ((t \in \mathbb{R}) \wedge (0 < x)))) \Rightarrow  \\ \quad \quad {{x}^{s}}^{t} = {x}^{(s \cdot t)} \end{array}}}
\vnbprop{rpow-star-neg}{\vnbstmt{\begin{array}{@{}l@{}} \forall x, s. \\ \quad ((x \in \mathbb{R}) \wedge ((s \in \mathbb{R}) \wedge (0 < x))) \Rightarrow  \\ \quad \quad {x}^{-s} = {{x}^{s}}^{\ast\!-1} \end{array}}}
\vnbprop{rpow-star-mono-exp}{\vnbstmt{\begin{array}{@{}l@{}} \forall x, s, t. \\ \quad ((x \in \mathbb{R}) \wedge ((s \in \mathbb{R}) \wedge ((t \in \mathbb{R}) \wedge ((1 < x) \wedge (s < t))))) \Rightarrow  \\ \quad \quad {x}^{s} < {x}^{t} \end{array}}}
\vnbprop{rpow-star-mono-base}{\vnbstmt{\begin{array}{@{}l@{}} \forall x, y, s. \\ \quad ((x \in \mathbb{R}) \wedge ((y \in \mathbb{R}) \wedge ((s \in \mathbb{R}) \wedge ((0 < x) \wedge ((x < y) \wedge (0 < s)))))) \Rightarrow  \\ \quad \quad {x}^{s} < {y}^{s} \end{array}}}

\vnbbookfile{theorem-library/rpow-star-convexity}{theorem-library/rpow-star-convexity}
\noindent{\small theorem-library/rpow-star-convexity.scm -- the two REAL inequalities of the axiomatic RPOW, proved for the DEFINED power RPOW-STAR (rpow-star.scm):}\par
\vnbprop{r-exp-ge-one-plus}{\vnbstmt{\forall z \in \mathbb{R}.\; ((1 + z) \leq \exp\left(z\right))}}
\vnbprop{log-le-minus-one}{\vnbstmt{\forall w \in \mathbb{R}.\; ((0 < w) \Rightarrow (\log\left(w\right) \leq (w - 1)))}}
\vnbprop{log-jensen-two}{\vnbstmt{\begin{array}{@{}l@{}} \forall u, v, s, t \in \mathbb{R}. \\ \quad ((0 < u) \wedge ((0 < v) \wedge ((0 < s) \wedge ((0 < t) \wedge ((s + t) = 1))))) \Rightarrow  \\ \quad \quad ((s \cdot \log\left(u\right)) + (t \cdot \log\left(v\right))) \leq \log\left(((s \cdot u) + (t \cdot v))\right) \end{array}}}
\vnbprop{young-inequality-star}{\vnbstmt{\begin{array}{@{}l@{}} \forall a, b, p, q \in \mathbb{R}. \\ \quad ((0 < a) \wedge ((0 < b) \wedge ((1 < p) \wedge ((1 < q) \wedge ((/(1, p) + /(1, q)) = 1))))) \Rightarrow  \\ \quad \quad (a \cdot b) \leq (/({a}^{p}, p) + /({b}^{q}, q)) \end{array}}}
\vnbprop{bernoulli-rpow-star}{\vnbstmt{\begin{array}{@{}l@{}} \forall x, b \in \mathbb{R}. \\ \quad ((0 < (1 + x)) \wedge (1 \leq b)) \Rightarrow ((1 + (b \cdot x)) \leq {(1 + x)}^{b}) \end{array}}}

\vnbbookfile{structure-library/functor-invariance}{structure-library/functor-invariance}
\noindent{\small functor-invariance.scm -- a construction that cannot see the functor.}\par
\vnbprop{metric-top@preimage}{\vnbstmt{\forall s, f, v.\; (\operatorname{preimage}(\operatorname{metric\text{-}top}(s), f, v) \simeq \operatorname{preimage}(s, f, v))}}
\vnbprop{trunc-metric@preimage}{\vnbstmt{\forall s, f, v.\; (\operatorname{preimage}(\operatorname{trunc\text{-}metric}(s), f, v) \simeq \operatorname{preimage}(s, f, v))}}

\vnbbookfile{theorem-library/metric-top-proof}{theorem-library/metric-top-proof}
\noindent{\small metric-top-proof.scm -- the metric open sets form a topology.}\par
\vnbprop{metric-top-is-top-space}{\vnbstmt{\forall d.\; (\operatorname{is\text{-}metric\text{-}space}(d) \Rightarrow \operatorname{is\text{-}top\text{-}space}(\operatorname{metric\text{-}top}(d)))}}
\vnbprop{metric-top-is-metrizable-top-space}{\vnbstmt{\forall d.\; (\operatorname{is\text{-}metric\text{-}space}(d) \Rightarrow \operatorname{is\text{-}metrizable\text{-}top\text{-}space}(\operatorname{metric\text{-}top}(d)))}}

\vnbbookfile{theorem-library/metric-top-functorial-proof}{theorem-library/metric-top-functorial-proof}
\noindent{\small metric-top-functorial-proof.scm -- the Met -$>$ Top functor acts on morphisms.}\par
\vnbprop{metric-top-functorial}{\vnbstmt{\begin{array}{@{}l@{}} \forall a, b, f. \\ \quad \operatorname{is\text{-}hom\text{-}metric\text{-}space}(a, b, f) \Rightarrow  \\ \quad \quad \operatorname{is\text{-}hom\text{-}metrizable\text{-}top\text{-}space}(\operatorname{metric\text{-}top}(a), \operatorname{metric\text{-}top}(b), f) \end{array}}}

\vnbbookfile{theorem-library/rake-bdd-metric}{theorem-library/rake-bdd-metric}
\noindent{\small rake-bdd-metric.scm -- the three asserted leaves still on the bill of `metrizable-iff-bounded-metrizable' (theorem-library/metrizable-bounded-proof.scm), PROVEN.  Statements copied VERBATIM from their support sites:}\par
\vnbprop{metrizable-has-metric-top}{\vnbstmt{\begin{array}{@{}l@{}} \forall s. \\ \quad \operatorname{is\text{-}metrizable\text{-}top\text{-}space}(s) \Rightarrow  \\ \quad \quad \exists d.\; (\operatorname{is\text{-}metric\text{-}space}(d) \wedge (\operatorname{metric\text{-}top}(d) \simeq s)) \end{array}}}
\vnbprop{bdd-metric-is-bounded-metric-space}{\vnbstmt{\forall s.\; (\operatorname{is\text{-}metric\text{-}space}(s) \Rightarrow \operatorname{is\text{-}bounded\text{-}metric\text{-}space}(\operatorname{bdd\text{-}metric}(s)))}}
\vnbprop{bdd-metric-open-bwd}{\vnbstmt{\begin{array}{@{}l@{}} \forall s, u. \\ \quad (\operatorname{is\text{-}metric\text{-}space}(s) \wedge \operatorname{is\text{-}open}(\operatorname{bdd\text{-}metric}(s), u)) \Rightarrow \operatorname{is\text{-}open}(s, u) \end{array}}}
\vnbprop{bdd-metric-open-fwd}{\vnbstmt{\begin{array}{@{}l@{}} \forall s, u. \\ \quad (\operatorname{is\text{-}metric\text{-}space}(s) \wedge \operatorname{is\text{-}open}(s, u)) \Rightarrow \operatorname{is\text{-}open}(\operatorname{bdd\text{-}metric}(s), u) \end{array}}}
\vnbprop{bdd-metric-preserves-metric-top}{\vnbstmt{\begin{array}{@{}l@{}} \forall s. \\ \quad \operatorname{is\text{-}metric\text{-}space}(s) \Rightarrow (\operatorname{metric\text{-}top}(\operatorname{bdd\text{-}metric}(s)) \simeq \operatorname{metric\text{-}top}(s)) \end{array}}}

\vnbbookfile{theorem-library/metrizable-compact-image}{theorem-library/metrizable-compact-image}
\noindent{\small metrizable-compact-image.scm -- THE CONTINUOUS IMAGE OF A COMPACT SPACE IS COMPACT, IN THE CATEGORY OF METRIZABLE TOPOLOGICAL SPACES, PROVED THROUGH SEQUENCES (the user's request, 2026-10-01: "in that category ... it can be done using sequences.  I like using sequences to prove things whenever}\par
\vnbprop{metric-top-opens-iff}{\vnbstmt{\begin{array}{@{}l@{}} \forall d, u. \\ \quad \operatorname{is\text{-}metric\text{-}space}(d) \Rightarrow ((u \in \operatorname{opens}(\operatorname{metric\text{-}top}(d))) \Leftrightarrow \operatorname{is\text{-}open}(d, u)) \end{array}}}
\vnbprop{metric-top-compact}{\vnbstmt{\begin{array}{@{}l@{}} \forall d. \\ \quad \operatorname{is\text{-}metric\text{-}space}(d) \Rightarrow (\operatorname{is\text{-}compact\text{-}t}(\operatorname{metric\text{-}top}(d)) \Leftrightarrow \operatorname{is\text{-}compact}(d)) \end{array}}}
\vnbprop{metric-top-hom}{\vnbstmt{\begin{array}{@{}l@{}} \forall s, t, f. \\ \quad (\operatorname{is\text{-}metric\text{-}space}(s) \wedge \operatorname{is\text{-}metric\text{-}space}(t)) \Rightarrow  \\ \quad \quad \operatorname{is\text{-}hom\text{-}top\text{-}space}(\operatorname{metric\text{-}top}(s), \operatorname{metric\text{-}top}(t), f) \Leftrightarrow \operatorname{is\text{-}continuous}(s, t, f) \end{array}}}
\vnbprop{seq-compact-image}{\vnbstmt{\begin{array}{@{}l@{}} \forall d, t, f. \\ \quad (\operatorname{is\text{-}metric\text{-}space}(d) \wedge (\operatorname{is\text{-}metric\text{-}space}(t) \wedge (\operatorname{is\text{-}continuous}(d, t, f) \wedge \operatorname{seq\text{-}compact}(d)))) \Rightarrow  \\ \quad \quad \operatorname{seq\text{-}compact}(\operatorname{subspace\text{-}ms}(t, \operatorname{image}(f, \operatorname{pts}(d)))) \end{array}}}
\vnbprop{metrizable-compact-t-iff-seq-compact}{\vnbstmt{\begin{array}{@{}l@{}} \forall s, d. \\ \quad (\operatorname{is\text{-}metric\text{-}space}(d) \wedge (\operatorname{metric\text{-}top}(d) \simeq s)) \Rightarrow  \\ \quad \quad \operatorname{is\text{-}compact\text{-}t}(s) \Leftrightarrow \operatorname{seq\text{-}compact}(d) \end{array}}}
\vnbprop{metrizable-hom-iff-sequential}{\vnbstmt{\begin{array}{@{}l@{}} \forall s, t, d, b, f. \\ \quad (\operatorname{is\text{-}metric\text{-}space}(d) \wedge (\operatorname{is\text{-}metric\text{-}space}(b) \wedge ((\operatorname{metric\text{-}top}(d) \simeq s) \wedge (\operatorname{metric\text{-}top}(b) \simeq t)))) \Rightarrow  \\ \quad \quad \operatorname{is\text{-}hom\text{-}top\text{-}space}(s, t, f) \\ \quad \quad \Leftrightarrow\; f \in (\operatorname{pts}(d) \to \operatorname{pts}(b)) \wedge \\ \quad \quad \forall a \in \operatorname{pts}(d), q \in \operatorname{pts}(d)^{\mathbb{N}}. \\ \quad \quad \quad \operatorname{converges\text{-}to}(d, q, a) \Rightarrow \operatorname{converges\text{-}to}(b, \operatorname{compose}(f, q), f(a)) \end{array}}}
\vnbprop{metrizable-compact-image-ms}{\vnbstmt{\begin{array}{@{}l@{}} \forall s, t, f, a. \\ \quad \operatorname{is\text{-}metrizable\text{-}top\text{-}space}(s) \wedge \\ \quad \operatorname{is\text{-}metrizable\text{-}top\text{-}space}(t) \wedge \\ \quad \operatorname{is\text{-}hom\text{-}top\text{-}space}(s, t, f) \wedge \\ \quad \operatorname{is\text{-}compact\text{-}t}(s) \wedge \\ \quad \operatorname{is\text{-}metric\text{-}space}(a) \wedge \\ \quad (\operatorname{metric\text{-}top}(a) \simeq t) \Rightarrow \\ \quad \quad \operatorname{is\text{-}compact}(\operatorname{subspace\text{-}ms}(a, \operatorname{image}(f, \operatorname{pts}(s)))) \end{array}}}
\vnbprop{subspace-top-pts}{\vnbstmt{\forall t.; \forall a \in \mathbf{Set}.\; (\operatorname{pts}(\operatorname{subspace\text{-}top}(t, a)) \simeq a)}}
\vnbprop{subspace-top-opens}{\vnbstmt{\begin{array}{@{}l@{}} \forall t.; \forall a \in \mathbf{Set}. \\ \quad (\operatorname{opens}(\operatorname{subspace\text{-}top}(t, a)) \simeq \{w \in \operatorname{power}(a) : \exists u \in \operatorname{opens}(t).\; (w = (u \cap a))\}) \end{array}}}
\vnbprop{subspace-top-is-top-space}{\vnbstmt{\begin{array}{@{}l@{}} \forall t, a. \\ \quad (\operatorname{is\text{-}top\text{-}space}(t) \wedge (a \subseteq \operatorname{pts}(t))) \Rightarrow \operatorname{is\text{-}top\text{-}space}(\operatorname{subspace\text{-}top}(t, a)) \end{array}}}
\vnbprop{subspace-top-of-metric-top}{\vnbstmt{\begin{array}{@{}l@{}} \forall t, a. \\ \quad (\operatorname{is\text{-}metric\text{-}space}(t) \wedge (a \subseteq \operatorname{pts}(t))) \Rightarrow  \\ \quad \quad (\operatorname{subspace\text{-}top}(\operatorname{metric\text{-}top}(t), a) \simeq \operatorname{metric\text{-}top}(\operatorname{subspace\text{-}ms}(t, a))) \end{array}}}
\vnbprop{metrizable-compact-image}{\vnbstmt{\begin{array}{@{}l@{}} \forall s, t, f. \\ \quad \operatorname{is\text{-}metrizable\text{-}top\text{-}space}(s) \wedge \\ \quad \operatorname{is\text{-}metrizable\text{-}top\text{-}space}(t) \wedge \\ \quad \operatorname{is\text{-}hom\text{-}top\text{-}space}(s, t, f) \wedge \\ \quad \operatorname{is\text{-}compact\text{-}t}(s) \Rightarrow \\ \quad \quad \operatorname{is\text{-}compact\text{-}t}(\operatorname{subspace\text{-}top}(t, \operatorname{image}(f, \operatorname{pts}(s)))) \end{array}}}

\vnbbookfile{theorem-library/metrizable-bounded-proof}{theorem-library/metrizable-bounded-proof}
\noindent{\small metrizable-bounded-proof.scm -- T1: a topological space is metrizable iff it is metrizable by a BOUNDED metric.}\par
\vnbprop{metrizable-iff-bounded-metrizable}{\vnbstmt{\begin{array}{@{}l@{}} \forall s. \\ \quad \operatorname{is\text{-}top\text{-}space}(s) \Rightarrow  \\ \quad \quad \operatorname{is\text{-}metrizable\text{-}top\text{-}space}(s) \Leftrightarrow \exists d.\; (\operatorname{is\text{-}bounded\text{-}metric\text{-}space}(d) \wedge (\operatorname{metric\text{-}top}(d) \simeq s)) \end{array}}}

\vnbbookfile{theorem-library/geometric-series}{theorem-library/geometric-series}
\noindent{\small geometric-series.scm -- THE GEOMETRIC MAJORANT, PROVEN, and the four small facts about powers it rests on.}\par
\vnbprop{rr-abs-power}{\vnbstmt{\forall k \in \mathbb{N}, r \in \mathbb{R}.\; (\lvert \operatorname{power}(r, k) \rvert = \operatorname{power}(\lvert r \rvert, k))}}
\vnbprop{geometric-partial-sum-mul}{\vnbstmt{\begin{array}{@{}l@{}} \forall k \in \mathbb{N}, r \in \mathbb{R}. \\ \quad ((1 - r) \cdot \operatorname{series\text{-}partial\text{-}sum}((\lambda n \in \mathbb{N}.\; \operatorname{power}(r, n)), k)) = (1 - \operatorname{power}(r, k)) \end{array}}}
\vnbprop{rr-power-nonneg}{\vnbstmt{\forall k \in \mathbb{N}, r \in \mathbb{R}.\; ((0 \leq r) \Rightarrow (0 \leq \operatorname{power}(r, k)))}}
\vnbprop{geometric-partial-sum}{\vnbstmt{\begin{array}{@{}l@{}} \forall r \in \mathbb{R}, k \in \mathbb{N}. \\ \quad \neg (r = 1) \Rightarrow  \\ \quad \quad \operatorname{series\text{-}partial\text{-}sum}((\lambda n \in \mathbb{N}.\; \operatorname{power}(r, n)), k) = ((1 - \operatorname{power}(r, k)) \cdot \operatorname{recip}((1 - r))) \end{array}}}
\vnbprop{geometric-series-converges}{\vnbstmt{\begin{array}{@{}l@{}} \forall r \in \mathbb{R}. \\ \quad ((0 \leq r) \wedge (r < 1)) \Rightarrow \operatorname{series\text{-}converges}((\lambda n \in \mathbb{N}.\; \operatorname{power}(r, n))) \end{array}}}

\vnbbookfile{theorem-library/limsup-tests}{theorem-library/limsup-tests}
\noindent{\small limsup-tests.scm -- THE UPPER LIMIT OF A SEQUENCE IN [0, +inf], THE ROOT TEST AND THE RATIO TEST, as the notes state them. Batch 25-A, 2026-09-24.}\par
\vnbprop{einf-lower}{\vnbstmt{\forall s.; \forall x \in s.\; ((s \subseteq \operatorname{rr\text{-}pos\text{-}star}) \Rightarrow (\operatorname{einf}(s) \leq x))}}
\vnbprop{einf-greatest}{\vnbstmt{\begin{array}{@{}l@{}} \forall s.; \forall b \in \operatorname{rr\text{-}pos\text{-}star}. \\ \quad ((s \subseteq \operatorname{rr\text{-}pos\text{-}star}) \wedge \forall x \in s.\; (b \leq x)) \Rightarrow (b \leq \operatorname{einf}(s)) \end{array}}}
\vnbprop{einf-in}{\vnbstmt{\forall s.\; ((s \subseteq \operatorname{rr\text{-}pos\text{-}star}) \Rightarrow (\operatorname{einf}(s) \in \operatorname{rr\text{-}pos\text{-}star}))}}
\vnbprop{etail-subset-rr-pos-star}{\vnbstmt{\begin{array}{@{}l@{}} \forall f, n. \\ \quad (\forall k \in \mathbb{N}.\; (f(k) \in \mathbb{R}) \wedge \forall k \in \mathbb{N}.\; (0 \leq f(k))) \Rightarrow  \\ \quad \quad \operatorname{etail}(f, n) \subseteq \operatorname{rr\text{-}pos\text{-}star} \end{array}}}
\vnbprop{elimsup-eventually-below}{\vnbstmt{\begin{array}{@{}l@{}} \forall f.; \forall r \in \mathbb{R}. \\ \quad (\forall k \in \mathbb{N}.\; (f(k) \in \mathbb{R}) \wedge (\forall k \in \mathbb{N}.\; (0 \leq f(k)) \wedge (\operatorname{elimsup}(f) < r))) \Rightarrow  \\ \quad \quad \exists n \in \mathbb{N}.\; \forall k \in \mathbb{N}.\; ((n \leq k) \Rightarrow (f(k) < r)) \end{array}}}
\vnbprop{elimsup-frequently-above}{\vnbstmt{\begin{array}{@{}l@{}} \forall f.; \forall r \in \mathbb{R}, n \in \mathbb{N}. \\ \quad (\forall k \in \mathbb{N}.\; (f(k) \in \mathbb{R}) \wedge (\forall k \in \mathbb{N}.\; (0 \leq f(k)) \wedge ((0 \leq r) \wedge (r < \operatorname{elimsup}(f))))) \Rightarrow  \\ \quad \quad \exists k \in \mathbb{N}.\; ((n \leq k) \wedge (r < f(k))) \end{array}}}
\vnbprop{eliminf-eventually-above}{\vnbstmt{\begin{array}{@{}l@{}} \forall f.; \forall r \in \mathbb{R}. \\ \quad (\forall k \in \mathbb{N}.\; (f(k) \in \mathbb{R}) \wedge (\forall k \in \mathbb{N}.\; (0 \leq f(k)) \wedge ((0 \leq r) \wedge (r < \operatorname{eliminf}(f))))) \Rightarrow  \\ \quad \quad \exists n \in \mathbb{N}.\; \forall k \in \mathbb{N}.\; ((n \leq k) \Rightarrow (r < f(k))) \end{array}}}
\vnbprop{elimsup-in-rr-pos-star}{\vnbstmt{\begin{array}{@{}l@{}} \forall f. \\ \quad (\forall k \in \mathbb{N}.\; (f(k) \in \mathbb{R}) \wedge \forall k \in \mathbb{N}.\; (0 \leq f(k))) \Rightarrow  \\ \quad \quad \operatorname{elimsup}(f) \in \operatorname{rr\text{-}pos\text{-}star} \end{array}}}
\vnbprop{eliminf-in-rr-pos-star}{\vnbstmt{\begin{array}{@{}l@{}} \forall f. \\ \quad (\forall k \in \mathbb{N}.\; (f(k) \in \mathbb{R}) \wedge \forall k \in \mathbb{N}.\; (0 \leq f(k))) \Rightarrow  \\ \quad \quad \operatorname{eliminf}(f) \in \operatorname{rr\text{-}pos\text{-}star} \end{array}}}
\vnbprop{rr-power-mono-base}{\vnbstmt{\forall k \in \mathbb{N}, x, y \in \mathbb{R}.\; (((0 \leq x) \wedge (x \leq y)) \Rightarrow (\operatorname{power}(x, k) \leq \operatorname{power}(y, k)))}}
\vnbprop{rpow-star-nat}{\vnbstmt{\forall k \in \mathbb{N}, x \in \mathbb{R}.\; ((0 < x) \Rightarrow ({x}^{k} = \operatorname{power}(x, k)))}}
\vnbprop{recip-star-nn-pos}{\vnbstmt{\forall k \in \mathbb{N}.\; ((1 \leq k) \Rightarrow (0 < {k}^{\ast\!-1}))}}
\vnbprop{recip-star-nn-nonneg}{\vnbstmt{\forall k \in \mathbb{N}.\; (0 \leq {k}^{\ast\!-1})}}
\vnbprop{rpow-star-root-power}{\vnbstmt{\begin{array}{@{}l@{}} \forall k \in \mathbb{N}, x \in \mathbb{R}. \\ \quad ((1 \leq k) \wedge (0 < x)) \Rightarrow (\operatorname{power}({x}^{{k}^{\ast\!-1}}, k) = x) \end{array}}}
\vnbprop{rpow-star-nonneg-in-rr}{\vnbstmt{\forall x, s \in \mathbb{R}.\; (((0 \leq x) \wedge (0 \leq s)) \Rightarrow ({x}^{s} \in \mathbb{R}))}}
\vnbprop{rpow-star-nonneg-nonneg}{\vnbstmt{\forall x, s \in \mathbb{R}.\; (((0 \leq x) \wedge (0 \leq s)) \Rightarrow (0 \leq {x}^{s}))}}
\vnbprop{series-bounded-converges}{\vnbstmt{\begin{array}{@{}l@{}} \forall f \in (\mathbb{N} \to \mathbb{R}), b \in \mathbb{R}. \\ \quad (\forall k \in \mathbb{N}.\; (0 \leq f(k)) \wedge \forall k \in \mathbb{N}.\; (\operatorname{series\text{-}partial\text{-}sum}(f, k) \leq b)) \Rightarrow  \\ \quad \quad \operatorname{series\text{-}converges}(f) \end{array}}}
\vnbprop{series-eventually-geometric-converges}{\vnbstmt{\begin{array}{@{}l@{}} \forall f \in (\mathbb{N} \to \mathbb{R}), q, c \in \mathbb{R}, n \in \mathbb{N}. \\ \quad \forall k \in \mathbb{N}.\; (0 \leq f(k)) \wedge \\ \quad 0 \leq q \wedge \\ \quad q < 1 \wedge \\ \quad 0 \leq c \wedge \\ \quad \forall k \in \mathbb{N}.\; ((n \leq k) \Rightarrow (f(k) \leq (c \cdot \operatorname{power}(q, k)))) \Rightarrow \\ \quad \quad \operatorname{series\text{-}converges}(f) \end{array}}}
\vnbprop{series-frequently-large-diverges}{\vnbstmt{\begin{array}{@{}l@{}} \forall f \in (\mathbb{N} \to \mathbb{R}), c \in \mathbb{R}. \\ \quad ((0 < c) \wedge \forall n \in \mathbb{N}.\; \exists k \in \mathbb{N}.\; ((n \leq k) \wedge (c \leq f(k)))) \Rightarrow  \\ \quad \quad \neg \operatorname{series\text{-}converges}(f) \end{array}}}
\vnbprop{root-test-converges}{\vnbstmt{\begin{array}{@{}l@{}} \forall h \in (\mathbb{N} \to \mathbb{R}). \\ \quad (\forall k \in \mathbb{N}.\; (0 \leq h(k)) \wedge (\operatorname{elimsup}((\lambda j \in \mathbb{N}.\; {h(j)}^{{j}^{\ast\!-1}})) < 1)) \Rightarrow  \\ \quad \quad \operatorname{series\text{-}converges}(h) \end{array}}}
\vnbprop{root-test-diverges}{\vnbstmt{\begin{array}{@{}l@{}} \forall h \in (\mathbb{N} \to \mathbb{R}). \\ \quad (\forall k \in \mathbb{N}.\; (0 \leq h(k)) \wedge (1 < \operatorname{elimsup}((\lambda j \in \mathbb{N}.\; {h(j)}^{{j}^{\ast\!-1}})))) \Rightarrow  \\ \quad \quad \neg \operatorname{series\text{-}converges}(h) \end{array}}}
\vnbprop{root-test}{\vnbstmt{\begin{array}{@{}l@{}} \forall h \in (\mathbb{N} \to \mathbb{R}). \\ \quad \forall k \in \mathbb{N}.\; (0 \leq h(k)) \Rightarrow  \\ \quad \quad (\operatorname{elimsup}((\lambda j \in \mathbb{N}.\; {h(j)}^{{j}^{\ast\!-1}})) < 1) \Rightarrow  \\ \quad \quad \quad \operatorname{series\text{-}converges}(h) \wedge \\ \quad \quad (1 < \operatorname{elimsup}((\lambda j \in \mathbb{N}.\; {h(j)}^{{j}^{\ast\!-1}}))) \Rightarrow  \\ \quad \quad \quad \neg \operatorname{series\text{-}converges}(h) \end{array}}}
\vnbprop{ratio-test-converges-limsup}{\vnbstmt{\begin{array}{@{}l@{}} \forall h \in (\mathbb{N} \to \mathbb{R}). \\ \quad (\forall k \in \mathbb{N}.\; (0 < h(k)) \wedge (\operatorname{elimsup}((\lambda j \in \mathbb{N}.\; (h((j + 1)) \cdot \operatorname{recip}(h(j))))) < 1)) \Rightarrow  \\ \quad \quad \operatorname{series\text{-}converges}(h) \end{array}}}
\vnbprop{ratio-test-diverges-liminf}{\vnbstmt{\begin{array}{@{}l@{}} \forall h \in (\mathbb{N} \to \mathbb{R}). \\ \quad (\forall k \in \mathbb{N}.\; (0 < h(k)) \wedge (1 < \operatorname{eliminf}((\lambda j \in \mathbb{N}.\; (h((j + 1)) \cdot \operatorname{recip}(h(j))))))) \Rightarrow  \\ \quad \quad \neg \operatorname{series\text{-}converges}(h) \end{array}}}
\vnbprop{ratio-test}{\vnbstmt{\begin{array}{@{}l@{}} \forall h \in (\mathbb{N} \to \mathbb{R}). \\ \quad \forall k \in \mathbb{N}.\; (0 < h(k)) \Rightarrow  \\ \quad \quad (\operatorname{elimsup}((\lambda j \in \mathbb{N}.\; (h((j + 1)) \cdot \operatorname{recip}(h(j))))) < 1) \Rightarrow  \\ \quad \quad \quad \operatorname{series\text{-}converges}(h) \wedge \\ \quad \quad (1 < \operatorname{eliminf}((\lambda j \in \mathbb{N}.\; (h((j + 1)) \cdot \operatorname{recip}(h(j)))))) \Rightarrow  \\ \quad \quad \quad \neg \operatorname{series\text{-}converges}(h) \end{array}}}

\vnbbookfile{theorem-library/rpow-defined}{theorem-library/rpow-defined}
\noindent{\small theorem-library/rpow-defined.scm -- THE AXIOMATIC REAL POWER RETIRED: RPOW is DEFINED as RPOW-STAR, and its laws are THEOREMS under their old names and with their old statements (the user's decision, 2026-10-03).}\par
\vnbdef{rpow-def}{\vnbstmt{\forall a, b.\; (\operatorname{rpow}(a, b) \simeq {a}^{b})}}
\vnbprop{rpow-pos}{\vnbstmt{\forall a \in \mathbb{R}, b \in \mathbb{Q}.\; ((0 < a) \Rightarrow ((\operatorname{rpow}(a, b) \in \mathbb{R}) \wedge (0 < \operatorname{rpow}(a, b))))}}
\vnbprop{rpow-zero}{\vnbstmt{\forall a \in \mathbb{R}.\; ((0 < a) \Rightarrow (\operatorname{rpow}(a, 0) = 1))}}
\vnbprop{rpow-one}{\vnbstmt{\forall a \in \mathbb{R}.\; ((0 < a) \Rightarrow (\operatorname{rpow}(a, 1) = a))}}
\vnbprop{rpow-add}{\vnbstmt{\begin{array}{@{}l@{}} \forall a \in \mathbb{R}, b, d \in \mathbb{Q}. \\ \quad (0 < a) \Rightarrow (\operatorname{rpow}(a, (b + d)) = (\operatorname{rpow}(a, b) \cdot \operatorname{rpow}(a, d))) \end{array}}}
\vnbprop{rpow-mul-base}{\vnbstmt{\begin{array}{@{}l@{}} \forall a, c \in \mathbb{R}, b \in \mathbb{Q}. \\ \quad ((0 < a) \wedge (0 < c)) \Rightarrow (\operatorname{rpow}((a \cdot c), b) = (\operatorname{rpow}(a, b) \cdot \operatorname{rpow}(c, b))) \end{array}}}
\vnbprop{rpow-pow}{\vnbstmt{\begin{array}{@{}l@{}} \forall a \in \mathbb{R}, b, d \in \mathbb{Q}. \\ \quad (0 < a) \Rightarrow (\operatorname{rpow}(\operatorname{rpow}(a, b), d) = \operatorname{rpow}(a, (b \cdot d))) \end{array}}}
\vnbprop{rpow-neg}{\vnbstmt{\forall a \in \mathbb{R}, b \in \mathbb{Q}.\; ((0 < a) \Rightarrow (\operatorname{rpow}(a, (0 - b)) = \operatorname{recip}(\operatorname{rpow}(a, b))))}}
\vnbprop{rpow-nat}{\vnbstmt{\forall a \in \mathbb{R}, n \in \mathbb{N}.\; ((0 < a) \Rightarrow (\operatorname{rpow}(a, n) = \operatorname{power}(a, n)))}}
\vnbprop{rpow-zero-base}{\vnbstmt{\forall b \in \mathbb{Q}.\; ((0 < b) \Rightarrow (\operatorname{rpow}(0, b) = 0))}}
\vnbprop{rpow-zero-zero}{\vnbstmt{\operatorname{rpow}(0, 0) = 1}}
\vnbprop{rpow-mono-base}{\vnbstmt{\begin{array}{@{}l@{}} \forall a, c \in \mathbb{R}, b \in \mathbb{Q}. \\ \quad ((0 < a) \wedge ((a \leq c) \wedge (0 \leq b))) \Rightarrow (\operatorname{rpow}(a, b) \leq \operatorname{rpow}(c, b)) \end{array}}}
\vnbprop{rpow-mono-exp-ge1}{\vnbstmt{\begin{array}{@{}l@{}} \forall a \in \mathbb{R}, b, d \in \mathbb{Q}. \\ \quad ((1 \leq a) \wedge (b \leq d)) \Rightarrow (\operatorname{rpow}(a, b) \leq \operatorname{rpow}(a, d)) \end{array}}}
\vnbprop{rpow-mono-exp-le1}{\vnbstmt{\begin{array}{@{}l@{}} \forall a \in \mathbb{R}, b, d \in \mathbb{Q}. \\ \quad (((0 < a) \wedge (a \leq 1)) \wedge (b \leq d)) \Rightarrow (\operatorname{rpow}(a, d) \leq \operatorname{rpow}(a, b)) \end{array}}}
\vnbprop{sqrt-rpow}{\vnbstmt{\forall a \in \mathbb{R}.\; ((0 < a) \Rightarrow (\sqrt{a} = \operatorname{rpow}(a, /(1, 2))))}}
\vnbprop{amgm-2-sqrt}{\vnbstmt{\begin{array}{@{}l@{}} \forall a, b \in \mathbb{R}. \\ \quad ((0 \leq a) \wedge (0 \leq b)) \Rightarrow (\sqrt{(a \cdot b)} \leq /((a + b), 2)) \end{array}}}
\vnbprop{young-inequality}{\vnbstmt{\begin{array}{@{}l@{}} \forall a, b \in \mathbb{R}, p, q \in \mathbb{Q}. \\ \quad ((0 < a) \wedge ((0 < b) \wedge ((1 < p) \wedge ((1 < q) \wedge ((/(1, p) + /(1, q)) = 1))))) \Rightarrow  \\ \quad \quad (a \cdot b) \leq (/(\operatorname{rpow}(a, p), p) + /(\operatorname{rpow}(b, q), q)) \end{array}}}
\vnbprop{bernoulli-rpow}{\vnbstmt{\begin{array}{@{}l@{}} \forall x \in \mathbb{R}, b \in \mathbb{Q}. \\ \quad ((0 < (1 + x)) \wedge (1 \leq b)) \Rightarrow ((1 + (b \cdot x)) \leq \operatorname{rpow}((1 + x), b)) \end{array}}}

\vnbbookfile{theorem-library/rake-asm-b1}{theorem-library/rake-asm-b1}
\noindent{\small theorem-library/rake-asm-b1.scm -- THE FOUR FINITE-SUM INEQUALITIES OF analysis-inequalities.scm, PROVEN under their own names and statements (RAKE-ASM-B, agent RB-1, 2026-10-12):}\par
\vnbprop{rb1-fs-empty}{\vnbstmt{\begin{array}{@{}l@{}} \forall f. \\ \quad \operatorname{finsum}(\operatorname{commutative\text{-}ring\text{-}additive\text{-}ag}(\operatorname{rr\text{-}normed\text{-}field}), f, \emptyset) = 0 \end{array}}}
\vnbprop{rb1-fs-real}{\vnbstmt{\begin{array}{@{}l@{}} \forall s \in \mathbf{Set}, f. \\ \quad ((|s| \in \mathbb{N}) \wedge \forall z \in s.\; (f(z) \in \mathbb{R})) \Rightarrow  \\ \quad \quad \operatorname{finsum}(\operatorname{commutative\text{-}ring\text{-}additive\text{-}ag}(\operatorname{rr\text{-}normed\text{-}field}), f, s) \in \mathbb{R} \end{array}}}
\vnbprop{rb1-fs-ins}{\vnbstmt{\begin{array}{@{}l@{}} \forall s, k \in \mathbf{Set}, f. \\ \quad ((|s| \in \mathbb{N}) \wedge (\neg (k \in s) \wedge \forall z \in (s \cup \{k\}).\; (f(z) \in \mathbb{R}))) \Rightarrow  \\ \quad \quad \operatorname{finsum}(\operatorname{commutative\text{-}ring\text{-}additive\text{-}ag}(\operatorname{rr\text{-}normed\text{-}field}), f, (s \cup \{k\})) \\ \quad \quad =\; \operatorname{finsum}(\operatorname{commutative\text{-}ring\text{-}additive\text{-}ag}(\operatorname{rr\text{-}normed\text{-}field}), f, s) + f(k) \end{array}}}
\vnbprop{rb1-fs-lin}{\vnbstmt{\begin{array}{@{}l@{}} \forall s \in \mathbf{Set}, a, b \in \mathbb{R}, f, g, h. \\ \quad |s| \in \mathbb{N} \wedge \\ \quad \forall z \in s.\; (f(z) \in \mathbb{R}) \wedge \\ \quad \forall z \in s.\; (g(z) \in \mathbb{R}) \wedge \\ \quad \forall z \in s.\; (h(z) = ((a \cdot f(z)) + (b \cdot g(z)))) \Rightarrow \\ \quad \quad \operatorname{finsum}(\operatorname{commutative\text{-}ring\text{-}additive\text{-}ag}(\operatorname{rr\text{-}normed\text{-}field}), h, s) \\ \quad \quad =\; (a \cdot \operatorname{finsum}(\operatorname{commutative\text{-}ring\text{-}additive\text{-}ag}(\operatorname{rr\text{-}normed\text{-}field}), f, s)) + (b \cdot \operatorname{finsum}(\operatorname{commutative\text{-}ring\text{-}additive\text{-}ag}(\operatorname{rr\text{-}normed\text{-}field}), g, s)) \end{array}}}
\vnbprop{rb1-disc}{\vnbstmt{\begin{array}{@{}l@{}} \forall a, b, c \in \mathbb{R}. \\ \quad \forall l, m \in \mathbb{R}.\; (0 \leq ((((l \cdot l) \cdot a) + (-(2 \cdot (l \cdot m)) \cdot c)) + ((m \cdot m) \cdot b))) \Rightarrow  \\ \quad \quad (c \cdot c) \leq (a \cdot b) \end{array}}}
\vnbprop{cauchy-schwarz-finite}{\vnbstmt{\begin{array}{@{}l@{}} \forall s \in \mathbf{Set}, a, b \in (s \to \mathbb{R}). \\ \quad (|s| \in \mathbb{N}) \Rightarrow  \\ \quad \quad \operatorname{finsum}(\operatorname{commutative\text{-}ring\text{-}additive\text{-}ag}(\operatorname{rr\text{-}normed\text{-}field}), (\lambda i \in s.\; (a(i) \cdot b(i))), s) \cdot \operatorname{finsum}(\operatorname{commutative\text{-}ring\text{-}additive\text{-}ag}(\operatorname{rr\text{-}normed\text{-}field}), (\lambda i \in s.\; (a(i) \cdot b(i))), s) \\ \quad \quad \leq\; \operatorname{finsum}(\operatorname{commutative\text{-}ring\text{-}additive\text{-}ag}(\operatorname{rr\text{-}normed\text{-}field}), (\lambda i \in s.\; (a(i) \cdot a(i))), s) \cdot \operatorname{finsum}(\operatorname{commutative\text{-}ring\text{-}additive\text{-}ag}(\operatorname{rr\text{-}normed\text{-}field}), (\lambda i \in s.\; (b(i) \cdot b(i))), s) \end{array}}}
\vnbprop{cauchy-schwarz-sqrt}{\vnbstmt{\begin{array}{@{}l@{}} \forall s \in \mathbf{Set}, a, b \in (s \to \mathbb{R}). \\ \quad (|s| \in \mathbb{N}) \Rightarrow  \\ \quad \quad \lvert \operatorname{finsum}(\operatorname{commutative\text{-}ring\text{-}additive\text{-}ag}(\operatorname{rr\text{-}normed\text{-}field}), (\lambda i \in s.\; (a(i) \cdot b(i))), s) \rvert \\ \quad \quad \leq\; \sqrt{\operatorname{finsum}(\operatorname{commutative\text{-}ring\text{-}additive\text{-}ag}(\operatorname{rr\text{-}normed\text{-}field}), (\lambda i \in s.\; (a(i) \cdot a(i))), s)} \cdot \sqrt{\operatorname{finsum}(\operatorname{commutative\text{-}ring\text{-}additive\text{-}ag}(\operatorname{rr\text{-}normed\text{-}field}), (\lambda i \in s.\; (b(i) \cdot b(i))), s)} \end{array}}}
\vnbprop{minkowski-l2}{\vnbstmt{\begin{array}{@{}l@{}} \forall s \in \mathbf{Set}, a, b \in (s \to \mathbb{R}). \\ \quad (|s| \in \mathbb{N}) \Rightarrow  \\ \quad \quad \sqrt{\operatorname{finsum}(\operatorname{commutative\text{-}ring\text{-}additive\text{-}ag}(\operatorname{rr\text{-}normed\text{-}field}), (\lambda i \in s.\; ((a(i) + b(i)) \cdot (a(i) + b(i)))), s)} \\ \quad \quad \leq\; \sqrt{\operatorname{finsum}(\operatorname{commutative\text{-}ring\text{-}additive\text{-}ag}(\operatorname{rr\text{-}normed\text{-}field}), (\lambda i \in s.\; (a(i) \cdot a(i))), s)} + \sqrt{\operatorname{finsum}(\operatorname{commutative\text{-}ring\text{-}additive\text{-}ag}(\operatorname{rr\text{-}normed\text{-}field}), (\lambda i \in s.\; (b(i) \cdot b(i))), s)} \end{array}}}
\vnbprop{rb1-rpow-nn}{\vnbstmt{\begin{array}{@{}l@{}} \forall x \in \mathbb{R}, r \in \mathbb{Q}. \\ \quad ((0 \leq x) \wedge (0 < r)) \Rightarrow ((\operatorname{rpow}(x, r) \in \mathbb{R}) \wedge (0 \leq \operatorname{rpow}(x, r))) \end{array}}}
\vnbprop{rb1-rpow-zero}{\vnbstmt{\forall x \in \mathbb{R}, r \in \mathbb{Q}.\; (((0 \leq x) \wedge ((0 < r) \wedge (\operatorname{rpow}(x, r) = 0))) \Rightarrow (x = 0))}}
\vnbprop{rb1-rpow-scale}{\vnbstmt{\begin{array}{@{}l@{}} \forall x, c \in \mathbb{R}, r \in \mathbb{Q}. \\ \quad ((0 \leq x) \wedge ((0 < c) \wedge (0 < r))) \Rightarrow  \\ \quad \quad \operatorname{rpow}((x \cdot \operatorname{recip}(c)), r) = (\operatorname{rpow}(x, r) \cdot \operatorname{recip}(\operatorname{rpow}(c, r))) \end{array}}}
\vnbprop{rb1-young0}{\vnbstmt{\begin{array}{@{}l@{}} \forall p, q \in \mathbb{Q}, u, v \in \mathbb{R}. \\ \quad ((1 < p) \wedge ((1 < q) \wedge (((/(1, p) + /(1, q)) = 1) \wedge ((0 \leq u) \wedge (0 \leq v))))) \Rightarrow  \\ \quad \quad (u \cdot v) \leq (/(\operatorname{rpow}(u, p), p) + /(\operatorname{rpow}(v, q), q)) \end{array}}}
\vnbprop{rb1-young-scaled}{\vnbstmt{\begin{array}{@{}l@{}} \forall p, q \in \mathbb{Q}, s, t, x, y \in \mathbb{R}. \\ \quad ((1 < p) \wedge ((1 < q) \wedge (((/(1, p) + /(1, q)) = 1) \wedge ((0 < s) \wedge ((0 < t) \wedge ((0 \leq x) \wedge (0 \leq y))))))) \Rightarrow  \\ \quad \quad x \cdot y \\ \quad \quad \leq\; ((((s \cdot t) \cdot \operatorname{recip}(\operatorname{rpow}(s, p))) \cdot \operatorname{recip}(p)) \cdot \operatorname{rpow}(x, p)) + ((((s \cdot t) \cdot \operatorname{recip}(\operatorname{rpow}(t, q))) \cdot \operatorname{recip}(q)) \cdot \operatorname{rpow}(y, q)) \end{array}}}
\vnbprop{rb1-fs-term-le}{\vnbstmt{\begin{array}{@{}l@{}} \forall s \in \mathbf{Set}, f. \\ \quad ((|s| \in \mathbb{N}) \wedge (\forall z \in s.\; (f(z) \in \mathbb{R}) \wedge \forall z \in s.\; (0 \leq f(z)))) \Rightarrow  \\ \quad \quad 0 \leq \operatorname{finsum}(\operatorname{commutative\text{-}ring\text{-}additive\text{-}ag}(\operatorname{rr\text{-}normed\text{-}field}), f, s) \wedge \\ \quad \quad \forall k \in s. \\ \quad \quad \quad f(k) \leq \operatorname{finsum}(\operatorname{commutative\text{-}ring\text{-}additive\text{-}ag}(\operatorname{rr\text{-}normed\text{-}field}), f, s) \end{array}}}
\vnbprop{holder-finite}{\vnbstmt{\begin{array}{@{}l@{}} \forall s \in \mathbf{Set}, a, b \in (s \to \mathbb{R}), p, q \in \mathbb{Q}. \\ \quad ((|s| \in \mathbb{N}) \wedge ((1 < p) \wedge ((1 < q) \wedge ((/(1, p) + /(1, q)) = 1)))) \Rightarrow  \\ \quad \quad \operatorname{finsum}(\operatorname{commutative\text{-}ring\text{-}additive\text{-}ag}(\operatorname{rr\text{-}normed\text{-}field}), (\lambda i \in s.\; \lvert (a(i) \cdot b(i)) \rvert), s) \\ \quad \quad \leq\; \operatorname{rpow}(\operatorname{finsum}(\operatorname{commutative\text{-}ring\text{-}additive\text{-}ag}(\operatorname{rr\text{-}normed\text{-}field}), (\lambda i \in s.\; \operatorname{rpow}(\lvert a(i) \rvert, p)), s), /(1, p)) \cdot \operatorname{rpow}(\operatorname{finsum}(\operatorname{commutative\text{-}ring\text{-}additive\text{-}ag}(\operatorname{rr\text{-}normed\text{-}field}), (\lambda i \in s.\; \operatorname{rpow}(\lvert b(i) \rvert, q)), s), /(1, q)) \end{array}}}

\vnbbookfile{theorem-library/ps-series-bridges}{theorem-library/ps-series-bridges}
\noindent{\small ps-series-bridges.scm -- THE THREE PS/SERIES BRIDGES, PROVEN, and the CONVERGES twin of rr-limit-ptwise-eq that the third one needed.}\par
\vnbprop{ps-abs-term}{\vnbstmt{\begin{array}{@{}l@{}} \forall f \in (\mathbb{N} \to \mathbb{R}), x \in \mathbb{R}, n \in \mathbb{N}. \\ \quad \lvert (f(n) \cdot \operatorname{power}(x, n)) \rvert = (\lvert f(n) \rvert \cdot \operatorname{power}(\lvert x \rvert, n)) \end{array}}}
\vnbprop{ps-partial-sum-as-series}{\vnbstmt{\begin{array}{@{}l@{}} \forall f \in (\mathbb{N} \to \mathbb{R}), x \in \mathbb{R}, k \in \mathbb{N}. \\ \quad \operatorname{ps\text{-}partial\text{-}sum}(f, x, k) = \operatorname{series\text{-}partial\text{-}sum}((\lambda n \in \mathbb{N}.\; (f(n) \cdot \operatorname{power}(x, n))), k) \end{array}}}
\vnbprop{rr-converges-ptwise-eq}{\vnbstmt{\begin{array}{@{}l@{}} \forall f, h \in (\mathbb{N} \to \mathbb{R}). \\ \quad (\forall j \in \mathbb{N}.\; (h(j) = f(j)) \wedge \operatorname{converges}(\operatorname{rr\text{-}ms}, f)) \Rightarrow  \\ \quad \quad \operatorname{converges}(\operatorname{rr\text{-}ms}, h) \end{array}}}
\vnbprop{ps-converges-as-series}{\vnbstmt{\begin{array}{@{}l@{}} \forall f \in (\mathbb{N} \to \mathbb{R}), x \in \mathbb{R}. \\ \quad \operatorname{ps\text{-}converges\text{-}at}(f, x) \Leftrightarrow \operatorname{series\text{-}converges}((\lambda n \in \mathbb{N}.\; (f(n) \cdot \operatorname{power}(x, n)))) \end{array}}}

\vnbbookfile{theorem-library/ratio-test-proof}{theorem-library/ratio-test-proof}
\noindent{\small ratio-test-proof.scm -- THE RATIO TEST FOR A REAL POWER SERIES (ratio-test-converges, structure-library/power-series.scm), PROVEN. Batch 27-A, 2026-09-24.}\par
\vnbprop{elimsup-eventually-le}{\vnbstmt{\begin{array}{@{}l@{}} \forall f.; \forall r \in \mathbb{R}, n \in \mathbb{N}. \\ \quad \forall k \in \mathbb{N}.\; (f(k) \in \mathbb{R}) \wedge \\ \quad \forall k \in \mathbb{N}.\; (0 \leq f(k)) \wedge \\ \quad \forall a \in \mathbb{N}.\; ((n \leq a) \Rightarrow (f(a) \leq r)) \Rightarrow \\ \quad \quad \operatorname{elimsup}(f) \leq r \end{array}}}
\vnbprop{ratio-test-converges}{\vnbstmt{\begin{array}{@{}l@{}} \forall f \in (\mathbb{N} \to \mathbb{R}), l, x \in \mathbb{R}. \\ \quad (((0 \leq l) \wedge (\forall n \in \mathbb{N}.\; \neg (f(n) = 0) \wedge \operatorname{ps\text{-}ratio\text{-}limit}(f, l))) \wedge ((\lvert x \rvert \cdot l) < 1)) \Rightarrow  \\ \quad \quad \operatorname{ps\text{-}converges\text{-}at}(f, x) \end{array}}}

\vnbbookfile{theorem-library/rake-series}{theorem-library/rake-series}
\noindent{\small rake-series.scm -- rake batch 5 T (2026-09-18): the little-o algebra, the power-series absolute-convergence bridge, and the Taylor clearing step.}\par
\vnbprop{vtaylor-clear}{\vnbstmt{\begin{array}{@{}l@{}} \forall v, a, b, w \in \mathbb{R}, n \in \mathbb{N}. \\ \quad (((\operatorname{factorial}((n + 1)) \cdot v) = (a \cdot w)) \wedge ((a \leq b) \wedge (0 \leq w))) \Rightarrow  \\ \quad \quad (\operatorname{factorial}((n + 1)) \cdot v) \leq (b \cdot w) \end{array}}}
\vnbprop{little-o-scalar}{\vnbstmt{\begin{array}{@{}l@{}} \forall c, g, a. \\ \quad ((c \in \mathbb{R}) \wedge \operatorname{little\text{-}o\text{-}at}(g, a)) \Rightarrow  \\ \quad \quad \operatorname{little\text{-}o\text{-}at}((\lambda x \in \mathbb{R}.\; (c \cdot g(x))), a) \end{array}}}
\vnbprop{little-o-sum}{\vnbstmt{\begin{array}{@{}l@{}} \forall g, h, a. \\ \quad (\operatorname{little\text{-}o\text{-}at}(g, a) \wedge \operatorname{little\text{-}o\text{-}at}(h, a)) \Rightarrow  \\ \quad \quad \operatorname{little\text{-}o\text{-}at}((\lambda x \in \mathbb{R}.\; (g(x) + h(x))), a) \end{array}}}
\vnbprop{shift-lam-in-fun}{\vnbstmt{\forall k \in \mathbb{R}, f \in (\mathbb{R} \to \mathbb{R}).\; ((\lambda x \in \mathbb{R}.\; (f(x) + k)) \in (\mathbb{R} \to \mathbb{R}))}}
\vnbprop{shift-continuous-at}{\vnbstmt{\begin{array}{@{}l@{}} \forall k \in \mathbb{R}, f, a. \\ \quad \operatorname{is\text{-}continuous\text{-}at}(\operatorname{rr\text{-}ms}, \operatorname{rr\text{-}ms}, f, a) \Rightarrow  \\ \quad \quad \operatorname{is\text{-}continuous\text{-}at}(\operatorname{rr\text{-}ms}, \operatorname{rr\text{-}ms}, (\lambda x \in \mathbb{R}.\; (f(x) + k)), a) \end{array}}}
\vnbprop{diff-iff-little-o}{\vnbstmt{\begin{array}{@{}l@{}} \forall f, a, l. \\ \quad \operatorname{is\text{-}diff\text{-}at}(f, a, l) \\ \quad \Leftrightarrow\; f \in (\mathbb{R} \to \mathbb{R}) \wedge \\ \quad a \in \mathbb{R} \wedge \\ \quad l \in \mathbb{R} \wedge \\ \quad \operatorname{little\text{-}o\text{-}at}((\lambda x \in \mathbb{R}.\; ((f(x) - f(a)) - (l \cdot (x - a)))), a) \end{array}}}
\vnbprop{ps-absolute-implies-convergent}{\vnbstmt{\begin{array}{@{}l@{}} \forall f \in (\mathbb{N} \to \mathbb{R}), x \in \mathbb{R}. \\ \quad \operatorname{ps\text{-}absolutely\text{-}converges\text{-}at}(f, x) \Rightarrow \operatorname{ps\text{-}converges\text{-}at}(f, x) \end{array}}}
\vnbprop{esum-finite-implies-bounded}{\vnbstmt{\begin{array}{@{}l@{}} \forall f \in (\operatorname{dom}(f) \to \operatorname{rr\text{-}pos\text{-}star}). \\ \quad (\operatorname{esum}(f) \in \mathbb{R}) \Rightarrow  \\ \quad \quad \exists m \in \mathbb{R}. \forall s \in \mathbf{Set}. \\ \quad \quad \quad ((|s| \in \mathbb{N}) \wedge (s \subseteq \operatorname{dom}(f))) \Rightarrow  \\ \quad \quad \quad \quad \operatorname{finsum}(\operatorname{rr\text{-}pos\text{-}star\text{-}add\text{-}monoid}, f, s) \leq m \end{array}}}

\vnbbookfile{theorem-library/rake-series2}{theorem-library/rake-series2}
\noindent{\small rake-series2.scm -- rake batch 5c (2026-09-18), assignment 5c-R: the two series leaves batch T left BLOCKED, and the four bricks they were missing.}\par
\vnbprop{metric-chain-bound}{\vnbstmt{\begin{array}{@{}l@{}} \forall n \in \mathbb{N}, s, f, d, m. \\ \quad \operatorname{is\text{-}metric\text{-}space}(s) \wedge \\ \quad f \in (\mathbb{N} \to \operatorname{pts}(s)) \wedge \\ \quad d \in (\mathbb{N} \to \mathbb{R}) \wedge \\ \quad \forall k \in \mathbb{N}.\; (\operatorname{dist}(s)(f(k), f((k + 1))) \leq d(k)) \wedge \\ \quad m \in \mathbb{N} \wedge \\ \quad m \leq n \Rightarrow \\ \quad \quad \operatorname{dist}(s)(f(m), f(n)) \leq (\operatorname{series\text{-}partial\text{-}sum}(d, n) - \operatorname{series\text{-}partial\text{-}sum}(d, m)) \end{array}}}
\vnbprop{metric-chain-triangle}{\vnbstmt{\begin{array}{@{}l@{}} \forall s, f.; \forall m, n \in \mathbb{N}. \\ \quad ((\operatorname{is\text{-}metric\text{-}space}(s) \wedge (f \in (\mathbb{N} \to \operatorname{pts}(s)))) \wedge (m \leq n)) \Rightarrow  \\ \quad \quad \operatorname{dist}(s)(f(m), f(n)) \\ \quad \quad \leq\; (\operatorname{series\text{-}partial\text{-}sum}((\lambda k \in \mathbb{N}.\; \operatorname{dist}(s)(f(k), f((k + 1)))), n) - \operatorname{series\text{-}partial\text{-}sum}((\lambda k \in \mathbb{N}.\; \operatorname{dist}(s)(f(k), f((k + 1)))), m)) \end{array}}}
\vnbprop{summable-bound-implies-cauchy}{\vnbstmt{\begin{array}{@{}l@{}} \forall s, f, d. \\ \quad \operatorname{is\text{-}metric\text{-}space}(s) \wedge \\ \quad f \in (\mathbb{N} \to \operatorname{pts}(s)) \wedge \\ \quad d \in (\mathbb{N} \to \mathbb{R}) \wedge \\ \quad \operatorname{series\text{-}converges}(d) \wedge \\ \quad \forall k \in \mathbb{N}.\; (\operatorname{dist}(s)(f(k), f((k + 1))) \leq d(k)) \Rightarrow \\ \quad \quad \operatorname{is\text{-}cauchy\text{-}seq}(s, f) \end{array}}}
\vnbprop{rr-power-le-one}{\vnbstmt{\forall k \in \mathbb{N}, r \in \mathbb{R}.\; (((0 \leq r) \wedge (r \leq 1)) \Rightarrow (\operatorname{power}(r, k) \leq 1))}}
\vnbprop{rr-power-shrink-bound}{\vnbstmt{\forall k \in \mathbb{N}, r \in \mathbb{R}.\; (((0 \leq r) \wedge (r \leq 1)) \Rightarrow (((k \cdot (1 - r)) \cdot \operatorname{power}(r, k)) \leq 1))}}
\vnbprop{power-small-threshold}{\vnbstmt{\begin{array}{@{}l@{}} \forall r \in \mathbb{R}, s. \\ \quad (((0 \leq r) \wedge (r < 1)) \wedge \operatorname{pos\text{-}rr}(s)) \Rightarrow  \\ \quad \quad \exists p \in \mathbb{N}.\; \forall k \in \mathbb{N}.\; ((p \leq k) \Rightarrow (\operatorname{power}(r, k) \leq s)) \end{array}}}
\vnbprop{power-tends-to-zero}{\vnbstmt{\begin{array}{@{}l@{}} \forall r \in \mathbb{R}. \\ \quad (\lvert r \rvert < 1) \Rightarrow \operatorname{converges\text{-}to}(\operatorname{rr\text{-}ms}, (\lambda n \in \mathbb{N}.\; \operatorname{power}(r, n)), 0) \end{array}}}
\vnbprop{rr-const-converges-to}{\vnbstmt{\forall c \in \mathbb{R}.\; \operatorname{converges\text{-}to}(\operatorname{rr\text{-}ms}, (\lambda n \in \mathbb{N}.\; c), c)}}
\vnbprop{geometric-series-converges-to}{\vnbstmt{\begin{array}{@{}l@{}} \forall r \in \mathbb{R}. \\ \quad (\lvert r \rvert < 1) \Rightarrow \operatorname{series\text{-}converges\text{-}to}((\lambda n \in \mathbb{N}.\; \operatorname{power}(r, n)), \operatorname{recip}((1 - r))) \end{array}}}

\vnbbookfile{theorem-library/heine-borel-interval}{theorem-library/heine-borel-interval}
\noindent{\small heine-borel-interval.scm -- HEINE-BOREL for a closed interval AS A SPACE:}\par
\vnbprop{rr-limit-in-ccint}{\vnbstmt{\begin{array}{@{}l@{}} \forall o, i \in \mathbb{R}, q \in (\mathbb{N} \to \operatorname{ccint}(o, i)), t \in \mathbb{R}. \\ \quad \operatorname{converges\text{-}to}(\operatorname{rr\text{-}ms}, q, t) \Rightarrow (t \in \operatorname{ccint}(o, i)) \end{array}}}
\vnbprop{heine-borel-ccint}{\vnbstmt{\forall v, \mathit{bv} \in \mathbb{R}.\; \operatorname{is\text{-}compact}(\operatorname{subspace\text{-}ms}(\operatorname{rr\text{-}ms}, \operatorname{ccint}(v, \mathit{bv})))}}

\vnbbookfile{theorem-library/cc-coords-laws}{theorem-library/cc-coords-laws}
\noindent{\small cc-coords-laws.scm -- CC and cartesian(RR, RR), explicitly.}\par
\vnbprop{apply-congruence-1}{\vnbstmt{\forall f, g, u.\; ((f \simeq g) \Rightarrow (f(u) \simeq g(u)))}}
\vnbprop{cc-coords-apply}{\vnbstmt{\forall z \in \mathbb{C}.\; (\operatorname{cc\text{-}coords}(z) \simeq \langle \operatorname{real\text{-}part}(z), \operatorname{imag\text{-}part}(z) \rangle)}}
\vnbprop{cc-coords-in-fun}{\vnbstmt{\operatorname{cc\text{-}coords} \in (\mathbb{C} \to (\mathbb{R} \times \mathbb{R}))}}
\vnbprop{cc-of-pair-apply}{\vnbstmt{\forall a, b \in \mathbb{R}.\; (\operatorname{cc\text{-}of\text{-}pair}(\langle a, b \rangle) \simeq (a + (b \cdot +i)))}}
\vnbprop{cc-of-pair-in-fun}{\vnbstmt{\operatorname{cc\text{-}of\text{-}pair} \in ((\mathbb{R} \times \mathbb{R}) \to \mathbb{C})}}
\vnbprop{cc-of-pair-of-coords}{\vnbstmt{\forall z \in \mathbb{C}.\; (\operatorname{cc\text{-}of\text{-}pair}(\operatorname{cc\text{-}coords}(z)) = z)}}
\vnbprop{cc-coords-of-pair}{\vnbstmt{\forall r \in (\mathbb{R} \times \mathbb{R}).\; (\operatorname{cc\text{-}coords}(\operatorname{cc\text{-}of\text{-}pair}(r)) = r)}}
\vnbprop{cc-coords-is-bijection}{\vnbstmt{\operatorname{cc\text{-}coords} \in \mathrm{Bij}(\mathbb{C}, (\mathbb{R} \times \mathbb{R}))}}
\vnbprop{cc-of-pair-is-bijection}{\vnbstmt{\operatorname{cc\text{-}of\text{-}pair} \in \mathrm{Bij}((\mathbb{R} \times \mathbb{R}), \mathbb{C})}}
\vnbprop{cc-re-add}{\vnbstmt{\forall z, w \in \mathbb{C}.\; (\operatorname{real\text{-}part}((z + w)) = (\operatorname{real\text{-}part}(z) + \operatorname{real\text{-}part}(w)))}}
\vnbprop{cc-im-add}{\vnbstmt{\forall z, w \in \mathbb{C}.\; (\operatorname{imag\text{-}part}((z + w)) = (\operatorname{imag\text{-}part}(z) + \operatorname{imag\text{-}part}(w)))}}
\vnbprop{cc-re-real-mul}{\vnbstmt{\forall c \in \mathbb{R}, z \in \mathbb{C}.\; (\operatorname{real\text{-}part}((c \cdot z)) = (c \cdot \operatorname{real\text{-}part}(z)))}}
\vnbprop{cc-im-real-mul}{\vnbstmt{\forall c \in \mathbb{R}, z \in \mathbb{C}.\; (\operatorname{imag\text{-}part}((c \cdot z)) = (c \cdot \operatorname{imag\text{-}part}(z)))}}
\vnbprop{cc-re-mul}{\vnbstmt{\begin{array}{@{}l@{}} \forall z, w \in \mathbb{C}. \\ \quad \operatorname{real\text{-}part}((z \cdot w)) = ((\operatorname{real\text{-}part}(z) \cdot \operatorname{real\text{-}part}(w)) - (\operatorname{imag\text{-}part}(z) \cdot \operatorname{imag\text{-}part}(w))) \end{array}}}
\vnbprop{cc-im-mul}{\vnbstmt{\begin{array}{@{}l@{}} \forall z, w \in \mathbb{C}. \\ \quad \operatorname{imag\text{-}part}((z \cdot w)) = ((\operatorname{real\text{-}part}(z) \cdot \operatorname{imag\text{-}part}(w)) + (\operatorname{imag\text{-}part}(z) \cdot \operatorname{real\text{-}part}(w))) \end{array}}}
\vnbprop{cc-coords-add}{\vnbstmt{\begin{array}{@{}l@{}} \forall z, w \in \mathbb{C}. \\ \quad \operatorname{cc\text{-}coords}((z + w)) = \langle (\operatorname{real\text{-}part}(z) + \operatorname{real\text{-}part}(w)), (\operatorname{imag\text{-}part}(z) + \operatorname{imag\text{-}part}(w)) \rangle \end{array}}}
\vnbprop{cc-coords-real-mul}{\vnbstmt{\begin{array}{@{}l@{}} \forall c \in \mathbb{R}, z \in \mathbb{C}. \\ \quad \operatorname{cc\text{-}coords}((c \cdot z)) = \langle (c \cdot \operatorname{real\text{-}part}(z)), (c \cdot \operatorname{imag\text{-}part}(z)) \rangle \end{array}}}
\vnbprop{cc-of-pair-add}{\vnbstmt{\begin{array}{@{}l@{}} \forall a_{1}, b_{1}, a_{2}, b_{2} \in \mathbb{R}. \\ \quad \operatorname{cc\text{-}of\text{-}pair}(\langle (a_{1} + a_{2}), (b_{1} + b_{2}) \rangle) = (\operatorname{cc\text{-}of\text{-}pair}(\langle a_{1}, b_{1} \rangle) + \operatorname{cc\text{-}of\text{-}pair}(\langle a_{2}, b_{2} \rangle)) \end{array}}}
\vnbprop{cc-of-pair-real-mul}{\vnbstmt{\begin{array}{@{}l@{}} \forall c, a, b \in \mathbb{R}. \\ \quad \operatorname{cc\text{-}of\text{-}pair}(\langle (c \cdot a), (c \cdot b) \rangle) = (c \cdot \operatorname{cc\text{-}of\text{-}pair}(\langle a, b \rangle)) \end{array}}}
\vnbprop{cc-re-dist-le}{\vnbstmt{\forall z, w \in \mathbb{C}.\; (\lvert (\operatorname{real\text{-}part}(z) - \operatorname{real\text{-}part}(w)) \rvert \leq \operatorname{magnitude}((z - w)))}}
\vnbprop{cc-im-dist-le}{\vnbstmt{\forall z, w \in \mathbb{C}.\; (\lvert (\operatorname{imag\text{-}part}(z) - \operatorname{imag\text{-}part}(w)) \rvert \leq \operatorname{magnitude}((z - w)))}}
\vnbprop{cc-dist-le-coords}{\vnbstmt{\begin{array}{@{}l@{}} \forall z, w \in \mathbb{C}. \\ \quad \operatorname{magnitude}((z - w)) \leq (\lvert (\operatorname{real\text{-}part}(z) - \operatorname{real\text{-}part}(w)) \rvert + \lvert (\operatorname{imag\text{-}part}(z) - \operatorname{imag\text{-}part}(w)) \rvert) \end{array}}}
\vnbprop{cc-coords-of-converges}{\vnbstmt{\begin{array}{@{}l@{}} \forall f \in (\mathbb{N} \to \mathbb{C}), g, q \in (\mathbb{N} \to \mathbb{R}), v \in \mathbb{C}. \\ \quad \forall n \in \mathbb{N}.\; (\operatorname{real\text{-}part}(f(n)) = g(n)) \wedge \\ \quad \forall n \in \mathbb{N}.\; (\operatorname{imag\text{-}part}(f(n)) = q(n)) \wedge \\ \quad \operatorname{converges\text{-}to}(\operatorname{cc\text{-}ms}, f, v) \Rightarrow \\ \quad \quad \operatorname{converges\text{-}to}(\operatorname{rr\text{-}ms}, g, \operatorname{real\text{-}part}(v)) \wedge \operatorname{converges\text{-}to}(\operatorname{rr\text{-}ms}, q, \operatorname{imag\text{-}part}(v)) \end{array}}}
\vnbprop{cc-converges-iff-coords}{\vnbstmt{\begin{array}{@{}l@{}} \forall f \in (\mathbb{N} \to \mathbb{C}), g, q \in (\mathbb{N} \to \mathbb{R}), v \in \mathbb{C}. \\ \quad (\forall n \in \mathbb{N}.\; (\operatorname{real\text{-}part}(f(n)) = g(n)) \wedge \forall n \in \mathbb{N}.\; (\operatorname{imag\text{-}part}(f(n)) = q(n))) \Rightarrow  \\ \quad \quad \operatorname{converges\text{-}to}(\operatorname{cc\text{-}ms}, f, v) \\ \quad \quad \Leftrightarrow\; \operatorname{converges\text{-}to}(\operatorname{rr\text{-}ms}, g, \operatorname{real\text{-}part}(v)) \wedge \operatorname{converges\text{-}to}(\operatorname{rr\text{-}ms}, q, \operatorname{imag\text{-}part}(v)) \end{array}}}
\vnbprop{cc-continuous-at-iff-coords}{\vnbstmt{\begin{array}{@{}l@{}} \forall s, f, r, i, a. \\ \quad \operatorname{is\text{-}metric\text{-}space}(s) \wedge \\ \quad f \in (\operatorname{pts}(s) \to \mathbb{C}) \wedge \\ \quad r \in (\operatorname{pts}(s) \to \mathbb{R}) \wedge \\ \quad i \in (\operatorname{pts}(s) \to \mathbb{R}) \wedge \\ \quad a \in \operatorname{pts}(s) \wedge \\ \quad \forall u \in \operatorname{pts}(s).\; (\operatorname{real\text{-}part}(f(u)) = r(u)) \wedge \\ \quad \forall u \in \operatorname{pts}(s).\; (\operatorname{imag\text{-}part}(f(u)) = i(u)) \Rightarrow \\ \quad \quad \operatorname{is\text{-}continuous\text{-}at}(s, \operatorname{cc\text{-}ms}, f, a) \\ \quad \quad \Leftrightarrow\; \operatorname{is\text{-}continuous\text{-}at}(s, \operatorname{rr\text{-}ms}, r, a) \wedge \operatorname{is\text{-}continuous\text{-}at}(s, \operatorname{rr\text{-}ms}, i, a) \end{array}}}
\vnbprop{cc-diff-at-iff-coords}{\vnbstmt{\begin{array}{@{}l@{}} \forall g \in (\mathbb{R} \to \mathbb{C}), r, i \in (\mathbb{R} \to \mathbb{R}), v \in \mathbb{R}, a \in \mathbb{C}. \\ \quad (\forall x \in \mathbb{R}.\; (\operatorname{real\text{-}part}(g(x)) = r(x)) \wedge \forall x \in \mathbb{R}.\; (\operatorname{imag\text{-}part}(g(x)) = i(x))) \Rightarrow  \\ \quad \quad \operatorname{is\text{-}cc\text{-}diff\text{-}at}(g, v, a) \Leftrightarrow (\operatorname{is\text{-}diff\text{-}at}(r, v, \operatorname{real\text{-}part}(a)) \wedge \operatorname{is\text{-}diff\text{-}at}(i, v, \operatorname{imag\text{-}part}(a))) \end{array}}}
\vnbprop{cc-diff-at-sum}{\vnbstmt{\begin{array}{@{}l@{}} \forall g, h, m \in (\mathbb{R} \to \mathbb{C}), v \in \mathbb{R}, a, b \in \mathbb{C}. \\ \quad (\forall x \in \mathbb{R}.\; (m(x) = (g(x) + h(x))) \wedge (\operatorname{is\text{-}cc\text{-}diff\text{-}at}(g, v, a) \wedge \operatorname{is\text{-}cc\text{-}diff\text{-}at}(h, v, b))) \Rightarrow  \\ \quad \quad \operatorname{is\text{-}cc\text{-}diff\text{-}at}(m, v, (a + b)) \end{array}}}
\vnbprop{cc-diff-at-real-mul}{\vnbstmt{\begin{array}{@{}l@{}} \forall c \in \mathbb{R}, g, d \in (\mathbb{R} \to \mathbb{C}), v \in \mathbb{R}, a \in \mathbb{C}. \\ \quad (\forall x \in \mathbb{R}.\; (d(x) = (c \cdot g(x))) \wedge \operatorname{is\text{-}cc\text{-}diff\text{-}at}(g, v, a)) \Rightarrow  \\ \quad \quad \operatorname{is\text{-}cc\text{-}diff\text{-}at}(d, v, (c \cdot a)) \end{array}}}
\vnbprop{cc-diff-at-product}{\vnbstmt{\begin{array}{@{}l@{}} \forall g, h, r \in (\mathbb{R} \to \mathbb{C}), v \in \mathbb{R}, a, b \in \mathbb{C}. \\ \quad (\forall x \in \mathbb{R}.\; (r(x) = (g(x) \cdot h(x))) \wedge (\operatorname{is\text{-}cc\text{-}diff\text{-}at}(g, v, a) \wedge \operatorname{is\text{-}cc\text{-}diff\text{-}at}(h, v, b))) \Rightarrow  \\ \quad \quad \operatorname{is\text{-}cc\text{-}diff\text{-}at}(r, v, ((a \cdot h(v)) + (g(v) \cdot b))) \end{array}}}

\vnbbookfile{structure-library/path-integral}{structure-library/path-integral}
\noindent{\small path-integral.scm -- THE VOCABULARY OF THE INTEGRAL ALONG A PATH. DEFINITIONS ONLY; every law is PROVEN in theorem-library/pw-antiderivative-laws.scm, cc-int-laws.scm and road-laws.scm.}\par
\par\medskip\noindent\textbf{Definition} (\texttt{is-pw-continuous-on}).\;For all $\operatorname{pphi\_}$, $a$ and $b$, 
\begin{equation*}
\begin{array}{@{}l@{}} \operatorname{is\text{-}pw\text{-}continuous\text{-}on}(\operatorname{pphi\_}, a, b) \\ \Leftrightarrow\; a \in \mathbb{R} \wedge \\ b \in \mathbb{R} \wedge \\ a < b \wedge \\ \operatorname{pphi\_} \in (\operatorname{ccint}(a, b) \to \mathbb{R}) \wedge \\ \exists \operatorname{pas\_} \in \mathbf{Set}. \\ \quad |\operatorname{pas\_}| \in \mathbb{N} \wedge \\ \quad \operatorname{pas\_} \subseteq \operatorname{ccint}(a, b) \wedge \\ \quad \forall \operatorname{pat\_} \in \operatorname{ooint}(a, b). \\ \quad \quad \neg (\operatorname{pat\_} \in \operatorname{pas\_}) \Rightarrow  \\ \quad \quad \quad \operatorname{is\text{-}continuous\text{-}at}(\operatorname{subspace\text{-}ms}(\operatorname{rr\text{-}ms}, \operatorname{ccint}(a, b)), \operatorname{rr\text{-}ms}, \operatorname{pphi\_}, \operatorname{pat\_}) \end{array}
\end{equation*}

\par\medskip\noindent\textbf{Definition} (\texttt{pw-int}).\;$\int_{a}^{b} \operatorname{pphi\_} :=$
\begin{equation*}
(\iota\, \operatorname{pav\_}.\;((\operatorname{pav\_} \in \mathbb{R}) \wedge \exists \operatorname{paw\_}.\; (\operatorname{is\text{-}primitive}(\operatorname{paw\_}, \operatorname{pphi\_}, a, b) \wedge (\operatorname{pav\_} = (\operatorname{paw\_}(b) - \operatorname{paw\_}(a))))))
\end{equation*}

\par\medskip\noindent\textbf{Definition} (\texttt{cc-int}).\;$\int_{a}^{b} \operatorname{pphi\_} :=$
\begin{equation*}
\int_{a}^{b} (\lambda \operatorname{pat\_} \in \operatorname{ccint}(a, b).\; \operatorname{real\text{-}part}(\operatorname{pphi\_}(\operatorname{pat\_}))) + (\int_{a}^{b} (\lambda \operatorname{pat\_} \in \operatorname{ccint}(a, b).\; \operatorname{imag\text{-}part}(\operatorname{pphi\_}(\operatorname{pat\_}))) \cdot +i)
\end{equation*}

\par\medskip\noindent\textbf{Definition} (\texttt{is-path}).\;For all $\mathit{pgam}$, $a$ and $b$, 
\begin{equation*}
\begin{array}{@{}l@{}} \operatorname{is\text{-}path}(\mathit{pgam}, a, b) \\ \Leftrightarrow\; a \in \mathbb{R} \wedge \\ b \in \mathbb{R} \wedge \\ a < b \wedge \\ \mathit{pgam} \in (\operatorname{ccint}(a, b) \to \mathbb{C}) \wedge \\ \forall \operatorname{pax\_} \in \operatorname{ccint}(a, b). \\ \quad \operatorname{is\text{-}continuous\text{-}at}(\operatorname{subspace\text{-}ms}(\operatorname{rr\text{-}ms}, \operatorname{ccint}(a, b)), \operatorname{cc\text{-}ms}, \mathit{pgam}, \operatorname{pax\_}) \end{array}
\end{equation*}

\par\medskip\noindent\textbf{Definition} (\texttt{is-road}).\;For all $\mathit{pgam}$, $\mathit{dgam}$, $a$ and $b$, 
\begin{equation*}
\begin{array}{@{}l@{}} \operatorname{is\text{-}road}(\mathit{pgam}, \mathit{dgam}, a, b) \\ \Leftrightarrow\; \operatorname{is\text{-}path}(\mathit{pgam}, a, b) \wedge \\ \mathit{dgam} \in (\operatorname{ccint}(a, b) \to \mathbb{C}) \wedge \\ \operatorname{is\text{-}regulated\text{-}on}((\lambda \operatorname{pat\_} \in \operatorname{ccint}(a, b).\; \operatorname{real\text{-}part}(\operatorname{dgam}(\operatorname{pat\_}))), a, b) \wedge \\ \operatorname{is\text{-}regulated\text{-}on}((\lambda \operatorname{pat\_} \in \operatorname{ccint}(a, b).\; \operatorname{imag\text{-}part}(\operatorname{dgam}(\operatorname{pat\_}))), a, b) \wedge \\ \operatorname{is\text{-}primitive}((\lambda \operatorname{pat\_} \in \operatorname{ccint}(a, b).\; \operatorname{real\text{-}part}(\operatorname{pgam}(\operatorname{pat\_}))), \\ \quad (\lambda \operatorname{pat\_} \in \operatorname{ccint}(a, b).\; \operatorname{real\text{-}part}(\operatorname{dgam}(\operatorname{pat\_}))), \\ \quad a, b) \wedge \\ \operatorname{is\text{-}primitive}((\lambda \operatorname{pat\_} \in \operatorname{ccint}(a, b).\; \operatorname{imag\text{-}part}(\operatorname{pgam}(\operatorname{pat\_}))), \\ \quad (\lambda \operatorname{pat\_} \in \operatorname{ccint}(a, b).\; \operatorname{imag\text{-}part}(\operatorname{dgam}(\operatorname{pat\_}))), \\ \quad a, b) \end{array}
\end{equation*}

\par\medskip\noindent\textbf{Definition} (\texttt{trace}).\;$\operatorname{trace}(\mathit{pgam}, a, b) :=$ $\operatorname{image}(\mathit{pgam}, \operatorname{ccint}(a, b))$
\par\medskip\noindent\textbf{Definition} (\texttt{line-int}).\;$\operatorname{line\text{-}int}(\mathit{pf}, \mathit{pgam}, \mathit{dgam}, a, b) :=$
\begin{equation*}
\int_{a}^{b} (\lambda \operatorname{pat\_} \in \operatorname{ccint}(a, b).\; (\operatorname{pf}(\operatorname{pgam}(\operatorname{pat\_})) \cdot \operatorname{dgam}(\operatorname{pat\_})))
\end{equation*}

\par\medskip\noindent\textbf{Definition} (\texttt{is-pw-antiderivable}).\;For all $\operatorname{pphi\_}$, $a$ and $b$, 
\begin{equation*}
\operatorname{is\text{-}pw\text{-}antiderivable}(\operatorname{pphi\_}, a, b) \Leftrightarrow \exists \operatorname{paw\_}.\; \operatorname{is\text{-}primitive}(\operatorname{paw\_}, \operatorname{pphi\_}, a, b)
\end{equation*}


\vnbbookfile{structure-library/segments-triangles}{structure-library/segments-triangles}
\noindent{\small segments-triangles.scm -- THE VOCABULARY OF CAUCHY-GOURSAT (complex-analysis.pdf ch. 3, s. 3.2).  DEFINITIONS ONLY; every law is PROVEN in theorem-library/goursat.scm.  Batch 30, 2026-09-25.}\par
\par\medskip\noindent\textbf{Definition} (\texttt{seg-path}).\;$\operatorname{seg\text{-}path}(\operatorname{sgp\_}, \operatorname{sgq\_}) :=$ $(\lambda \operatorname{sgs\_} \in \operatorname{ccint}(0, 1).\; (\operatorname{sgp\_} + (\operatorname{sgs\_} \cdot (\operatorname{sgq\_} - \operatorname{sgp\_}))))$
\par\medskip\noindent\textbf{Definition} (\texttt{seg-deriv}).\;$\operatorname{seg\text{-}deriv}(\operatorname{sgp\_}, \operatorname{sgq\_}) :=$ $(\lambda \operatorname{sgs\_} \in \operatorname{ccint}(0, 1).\; (\operatorname{sgq\_} - \operatorname{sgp\_}))$
\par\medskip\noindent\textbf{Definition} (\texttt{seg-int}).\;$\operatorname{seg\text{-}int}(\operatorname{sgf\_}, \operatorname{sgp\_}, \operatorname{sgq\_}) :=$
\begin{equation*}
\operatorname{line\text{-}int}(\operatorname{sgf\_}, \operatorname{seg\text{-}path}(\operatorname{sgp\_}, \operatorname{sgq\_}), \operatorname{seg\text{-}deriv}(\operatorname{sgp\_}, \operatorname{sgq\_}), 0, 1)
\end{equation*}

\par\medskip\noindent\textbf{Definition} (\texttt{tri-int}).\;$\operatorname{tri\text{-}int}(\operatorname{sgf\_}, \operatorname{sga\_}, \operatorname{sgb\_}, \operatorname{sgc\_}) :=$
\begin{equation*}
(\operatorname{seg\text{-}int}(\operatorname{sgf\_}, \operatorname{sga\_}, \operatorname{sgb\_}) + \operatorname{seg\text{-}int}(\operatorname{sgf\_}, \operatorname{sgb\_}, \operatorname{sgc\_})) + \operatorname{seg\text{-}int}(\operatorname{sgf\_}, \operatorname{sgc\_}, \operatorname{sga\_})
\end{equation*}

\par\medskip\noindent\textbf{Definition} (\texttt{conv3}).\;$\operatorname{conv3}(\operatorname{cva\_}, \operatorname{cvb\_}, \operatorname{cvc\_}) :=$
\begin{equation*}
\{\operatorname{cvz\_} \in \mathbb{C} : \exists \operatorname{cvl\_}, \operatorname{cvm\_}, \operatorname{cvn\_} \in \operatorname{ccint}(0, 1).\; ((((\operatorname{cvl\_} + \operatorname{cvm\_}) + \operatorname{cvn\_}) = 1) \wedge (\operatorname{cvz\_} = (((\operatorname{cvl\_} \cdot \operatorname{cva\_}) + (\operatorname{cvm\_} \cdot \operatorname{cvb\_})) + (\operatorname{cvn\_} \cdot \operatorname{cvc\_}))))\}
\end{equation*}

\par\medskip\noindent\textbf{Definition} (\texttt{cc-mid}).\;$\operatorname{cc\text{-}mid}(\operatorname{cmx\_}, \operatorname{cmy\_}) :=$ $(\operatorname{cmx\_} + \operatorname{cmy\_}) \cdot \operatorname{recip}(2)$
\par\medskip\noindent\textbf{Definition} (\texttt{is-convex-cc}).\;For every $\operatorname{cvu\_}$, 
\begin{equation*}
\begin{array}{@{}l@{}} \operatorname{is\text{-}convex\text{-}cc}(\operatorname{cvu\_}) \\ \Leftrightarrow\; \operatorname{cvu\_} \subseteq \mathbb{C} \wedge \\ \forall \operatorname{cvx\_}, \operatorname{cvy\_} \in \operatorname{cvu\_}, \operatorname{cvt\_} \in \operatorname{ccint}(0, 1). \\ \quad +(\operatorname{cvx\_}, \cdot(\operatorname{cvt\_}, -(\operatorname{cvy\_}, \operatorname{cvx\_}))) \in \operatorname{cvu\_} \end{array}
\end{equation*}


\vnbbookfile{structure-library/road-juxtaposition}{structure-library/road-juxtaposition}
\noindent{\small road-juxtaposition.scm -- THE JUXTAPOSITION OF ROADS, THE TRIANGULAR ROAD AND THE LENGTH OF A ROAD.  DEFINITIONS ONLY; every law is PROVEN in theorem-library/road-juxtaposition-laws.scm.  Batch 37, 2026-09-26.}\par
\par\medskip\noindent\textbf{Definition} (\texttt{juxta}).\;$\operatorname{juxta}(\operatorname{jxp\_}, \operatorname{jxq\_}, \operatorname{jxa\_}, \operatorname{jxb\_}, \operatorname{jxc\_}) :=$
\begin{equation*}
(\lambda \operatorname{jxt\_} \in \operatorname{ccint}(\operatorname{jxa\_}, \operatorname{jxc\_}).\; \left(\text{if }(\operatorname{jxt\_} \leq \operatorname{jxb\_})\text{ then }\operatorname{jxp\_}(\operatorname{jxt\_})\text{ else }\operatorname{jxq\_}(\operatorname{jxt\_})\right))
\end{equation*}

\par\medskip\noindent\textbf{Definition} (\texttt{road-length}).\;$\operatorname{road\text{-}length}(\operatorname{rld\_}, \operatorname{rla\_}, \operatorname{rlb\_}) :=$
\begin{equation*}
\int_{\operatorname{rla\_}}^{\operatorname{rlb\_}} (\lambda \operatorname{rlt\_} \in \operatorname{ccint}(\operatorname{rla\_}, \operatorname{rlb\_}).\; \operatorname{magnitude}(\operatorname{rld\_}(\operatorname{rlt\_})))
\end{equation*}

\par\medskip\noindent\textbf{Definition} (\texttt{tri-road}).\;$\operatorname{tri\text{-}road}(\operatorname{tra\_}, \operatorname{trb\_}, \operatorname{trc\_}) :=$
\begin{equation*}
\begin{array}{@{}l@{}} \operatorname{juxta}(\operatorname{juxta}(\operatorname{seg\text{-}path}(\operatorname{tra\_}, \operatorname{trb\_}), \\ \quad \quad (\lambda \operatorname{trs\_} \in \operatorname{ccint}(1, 2).\; \operatorname{seg\text{-}path}(\operatorname{trb\_}, \operatorname{trc\_})((\operatorname{trs\_} - 1))), \\ \quad \quad 0, 1, 2), \\ \quad (\lambda \operatorname{trs\_} \in \operatorname{ccint}(2, 3).\; \operatorname{seg\text{-}path}(\operatorname{trc\_}, \operatorname{tra\_})((\operatorname{trs\_} - 2))), \\ \quad 0, 2, 3) \end{array}
\end{equation*}

\par\medskip\noindent\textbf{Definition} (\texttt{tri-road-deriv}).\;$\operatorname{tri\text{-}road\text{-}deriv}(\operatorname{tra\_}, \operatorname{trb\_}, \operatorname{trc\_}) :=$
\begin{equation*}
\begin{array}{@{}l@{}} \operatorname{juxta}(\operatorname{juxta}(\operatorname{seg\text{-}deriv}(\operatorname{tra\_}, \operatorname{trb\_}), \\ \quad \quad (\lambda \operatorname{trs\_} \in \operatorname{ccint}(1, 2).\; \operatorname{seg\text{-}deriv}(\operatorname{trb\_}, \operatorname{trc\_})((\operatorname{trs\_} - 1))), \\ \quad \quad 0, 1, 2), \\ \quad (\lambda \operatorname{trs\_} \in \operatorname{ccint}(2, 3).\; \operatorname{seg\text{-}deriv}(\operatorname{trc\_}, \operatorname{tra\_})((\operatorname{trs\_} - 2))), \\ \quad 0, 2, 3) \end{array}
\end{equation*}


\vnbbookfile{theorem-library/regulated-primitive-laws}{theorem-library/regulated-primitive-laws}
\noindent{\small regulated-primitive-laws.scm -- the laws of IS-COUNTABLE, IS-RIGHT-LIMIT-WITHIN / IS-LEFT-LIMIT-WITHIN, IS-REGULATED-ON and IS-PRIMITIVE (structure-library/regulated-primitive.scm).  Batch 20-D, 2026-09-23.  THE SPECIFICATION is docs/roads-regulated-primitives-2026-09-23.md.}\par
\vnbprop{countable-of-enum}{\vnbstmt{\begin{array}{@{}l@{}} \forall d, e. \\ \quad ((d \in \mathbf{Set}) \wedge ((e \in (\mathbb{N} \to d)) \wedge \forall y \in d.\; \exists n \in \mathbb{N}.\; (y = e(n)))) \Rightarrow  \\ \quad \quad \operatorname{is\text{-}countable}(d) \end{array}}}
\vnbprop{countable-empty}{\vnbstmt{\operatorname{is\text{-}countable}(\emptyset)}}
\vnbprop{rr-abs-lt-parts}{\vnbstmt{\forall x, c \in \mathbb{R}.\; ((\lvert x \rvert < c) \Rightarrow ((-c < x) \wedge (x < c)))}}
\vnbprop{countable-singleton}{\vnbstmt{\forall x \in \mathbf{Set}.\; \operatorname{is\text{-}countable}(\{x\})}}
\vnbprop{countable-subset}{\vnbstmt{\forall a, d.\; ((\operatorname{is\text{-}countable}(d) \wedge (a \subseteq d)) \Rightarrow \operatorname{is\text{-}countable}(a))}}
\vnbprop{countable-cases}{\vnbstmt{\begin{array}{@{}l@{}} \forall d. \\ \quad \operatorname{is\text{-}countable}(d) \Rightarrow  \\ \quad \quad d \in \mathbf{Set} \wedge \\ \quad \quad ((d = \emptyset) \vee \exists e \in (\mathbb{N} \to d).\; \forall y \in d.\; \exists n \in \mathbb{N}.\; (y = e(n))) \end{array}}}
\vnbprop{countable-of-two-enum}{\vnbstmt{\begin{array}{@{}l@{}} \forall u, f, g. \\ \quad u \in \mathbf{Set} \wedge \\ \quad f \in (\mathbb{N} \to u) \wedge \\ \quad g \in (\mathbb{N} \to u) \wedge \\ \quad \forall y \in u.\; (\exists j \in \mathbb{N}.\; (y = f(j)) \vee \exists j \in \mathbb{N}.\; (y = g(j))) \Rightarrow \\ \quad \quad \operatorname{is\text{-}countable}(u) \end{array}}}
\vnbprop{countable-union-2}{\vnbstmt{\begin{array}{@{}l@{}} \forall a, b. \\ \quad (\operatorname{is\text{-}countable}(a) \wedge \operatorname{is\text{-}countable}(b)) \Rightarrow \operatorname{is\text{-}countable}((a \cup b)) \end{array}}}
\vnbprop{finite-implies-countable}{\vnbstmt{\forall d \in \mathbf{Set}.\; ((|d| \in \mathbb{N}) \Rightarrow \operatorname{is\text{-}countable}(d))}}
\vnbprop{countable-image}{\vnbstmt{\begin{array}{@{}l@{}} \forall d, h, a, b. \\ \quad (\operatorname{is\text{-}countable}(d) \wedge ((h \in (a \to b)) \wedge (d \subseteq a))) \Rightarrow  \\ \quad \quad \operatorname{is\text{-}countable}(\operatorname{image}(h, d)) \end{array}}}
\vnbprop{rr-abs-le-parts}{\vnbstmt{\forall x, c \in \mathbb{R}.\; ((\lvert x \rvert \leq c) \Rightarrow ((-c \leq x) \wedge (x \leq c)))}}
\vnbprop{rr-abs-le-of-parts}{\vnbstmt{\forall x, c \in \mathbb{R}.\; (((-c \leq x) \wedge (x \leq c)) \Rightarrow (\lvert x \rvert \leq c))}}
\vnbprop{rr-abs-lt-of-parts}{\vnbstmt{\forall x, c \in \mathbb{R}.\; (((-c < x) \wedge (x < c)) \Rightarrow (\lvert x \rvert < c))}}
\vnbprop{continuous-on-right-limit-within}{\vnbstmt{\begin{array}{@{}l@{}} \forall f, w.; \forall x \in w. \\ \quad ((w \subseteq \mathbb{R}) \wedge \operatorname{is\text{-}continuous\text{-}on}(f, w)) \Rightarrow  \\ \quad \quad \operatorname{is\text{-}right\text{-}limit\text{-}within}(f, w, x, f(x)) \end{array}}}
\vnbprop{continuous-on-left-limit-within}{\vnbstmt{\begin{array}{@{}l@{}} \forall f, w.; \forall x \in w. \\ \quad ((w \subseteq \mathbb{R}) \wedge \operatorname{is\text{-}continuous\text{-}on}(f, w)) \Rightarrow  \\ \quad \quad \operatorname{is\text{-}left\text{-}limit\text{-}within}(f, w, x, f(x)) \end{array}}}
\vnbprop{continuous-on-implies-regulated-on}{\vnbstmt{\begin{array}{@{}l@{}} \forall f, a, b. \\ \quad ((a \in \mathbb{R}) \wedge ((b \in \mathbb{R}) \wedge ((a < b) \wedge \operatorname{is\text{-}continuous\text{-}on}(f, \operatorname{ccint}(a, b))))) \Rightarrow  \\ \quad \quad \operatorname{is\text{-}regulated\text{-}on}(f, a, b) \end{array}}}
\vnbprop{right-limit-within-le}{\vnbstmt{\begin{array}{@{}l@{}} \forall f, w, x, l, m. \\ \quad \operatorname{is\text{-}right\text{-}limit\text{-}within}(f, w, x, l) \wedge \\ \quad \operatorname{is\text{-}right\text{-}limit\text{-}within}(f, w, x, m) \wedge \\ \quad \forall r.\; (\operatorname{pos\text{-}rr}(r) \Rightarrow \exists t \in w.\; ((x < t) \wedge (t < (x + r)))) \Rightarrow \\ \quad \quad (l - m) \leq 0 \end{array}}}
\vnbprop{left-limit-within-le}{\vnbstmt{\begin{array}{@{}l@{}} \forall f, w, x, l, m. \\ \quad \operatorname{is\text{-}left\text{-}limit\text{-}within}(f, w, x, l) \wedge \\ \quad \operatorname{is\text{-}left\text{-}limit\text{-}within}(f, w, x, m) \wedge \\ \quad \forall r.\; (\operatorname{pos\text{-}rr}(r) \Rightarrow \exists t \in w.\; (((x - r) < t) \wedge (t < x))) \Rightarrow \\ \quad \quad (l - m) \leq 0 \end{array}}}
\vnbprop{right-limit-within-unique}{\vnbstmt{\begin{array}{@{}l@{}} \forall f, w, x, l, m. \\ \quad \operatorname{is\text{-}right\text{-}limit\text{-}within}(f, w, x, l) \wedge \\ \quad \operatorname{is\text{-}right\text{-}limit\text{-}within}(f, w, x, m) \wedge \\ \quad \forall r.\; (\operatorname{pos\text{-}rr}(r) \Rightarrow \exists t \in w.\; ((x < t) \wedge (t < (x + r)))) \Rightarrow \\ \quad \quad l = m \end{array}}}
\vnbprop{left-limit-within-unique}{\vnbstmt{\begin{array}{@{}l@{}} \forall f, w, x, l, m. \\ \quad \operatorname{is\text{-}left\text{-}limit\text{-}within}(f, w, x, l) \wedge \\ \quad \operatorname{is\text{-}left\text{-}limit\text{-}within}(f, w, x, m) \wedge \\ \quad \forall r.\; (\operatorname{pos\text{-}rr}(r) \Rightarrow \exists t \in w.\; (((x - r) < t) \wedge (t < x))) \Rightarrow \\ \quad \quad l = m \end{array}}}
\vnbprop{ccint-right-approach}{\vnbstmt{\begin{array}{@{}l@{}} \forall a, b, x \in \mathbb{R}, r. \\ \quad ((a \leq x) \wedge ((x < b) \wedge \operatorname{pos\text{-}rr}(r))) \Rightarrow  \\ \quad \quad \exists t \in \operatorname{ccint}(a, b).\; ((x < t) \wedge (t < (x + r))) \end{array}}}
\vnbprop{ccint-left-approach}{\vnbstmt{\begin{array}{@{}l@{}} \forall a, b, x \in \mathbb{R}, r. \\ \quad ((a < x) \wedge ((x \leq b) \wedge \operatorname{pos\text{-}rr}(r))) \Rightarrow  \\ \quad \quad \exists t \in \operatorname{ccint}(a, b).\; (((x - r) < t) \wedge (t < x)) \end{array}}}
\vnbprop{right-limit-within-ccint-unique}{\vnbstmt{\begin{array}{@{}l@{}} \forall f.; \forall a, b, x \in \mathbb{R}, l, m. \\ \quad a \leq x \wedge \\ \quad x < b \wedge \\ \quad \operatorname{is\text{-}right\text{-}limit\text{-}within}(f, \operatorname{ccint}(a, b), x, l) \wedge \\ \quad \operatorname{is\text{-}right\text{-}limit\text{-}within}(f, \operatorname{ccint}(a, b), x, m) \Rightarrow \\ \quad \quad l = m \end{array}}}
\vnbprop{left-limit-within-ccint-unique}{\vnbstmt{\begin{array}{@{}l@{}} \forall f.; \forall a, b, x \in \mathbb{R}, l, m. \\ \quad a < x \wedge \\ \quad x \leq b \wedge \\ \quad \operatorname{is\text{-}left\text{-}limit\text{-}within}(f, \operatorname{ccint}(a, b), x, l) \wedge \\ \quad \operatorname{is\text{-}left\text{-}limit\text{-}within}(f, \operatorname{ccint}(a, b), x, m) \Rightarrow \\ \quad \quad l = m \end{array}}}
\vnbprop{right-limit-within-subset}{\vnbstmt{\begin{array}{@{}l@{}} \forall f, w, x, l, g, v. \\ \quad \operatorname{is\text{-}right\text{-}limit\text{-}within}(f, w, x, l) \wedge \\ \quad v \subseteq w \wedge \\ \quad g \in (v \to \mathbb{R}) \wedge \\ \quad \forall z \in v.\; (g(z) = f(z)) \Rightarrow \\ \quad \quad \operatorname{is\text{-}right\text{-}limit\text{-}within}(g, v, x, l) \end{array}}}
\vnbprop{left-limit-within-subset}{\vnbstmt{\begin{array}{@{}l@{}} \forall f, w, x, l, g, v. \\ \quad \operatorname{is\text{-}left\text{-}limit\text{-}within}(f, w, x, l) \wedge \\ \quad v \subseteq w \wedge \\ \quad g \in (v \to \mathbb{R}) \wedge \\ \quad \forall z \in v.\; (g(z) = f(z)) \Rightarrow \\ \quad \quad \operatorname{is\text{-}left\text{-}limit\text{-}within}(g, v, x, l) \end{array}}}
\vnbprop{regulated-on-restrict}{\vnbstmt{\begin{array}{@{}l@{}} \forall f, a, b, c, d, g. \\ \quad \operatorname{is\text{-}regulated\text{-}on}(f, a, b) \wedge \\ \quad c \in \mathbb{R} \wedge \\ \quad d \in \mathbb{R} \wedge \\ \quad a \leq c \wedge \\ \quad c < d \wedge \\ \quad d \leq b \wedge \\ \quad g \in (\operatorname{ccint}(c, d) \to \mathbb{R}) \wedge \\ \quad \forall z \in \operatorname{ccint}(c, d).\; (g(z) = f(z)) \Rightarrow \\ \quad \quad \operatorname{is\text{-}regulated\text{-}on}(g, c, d) \end{array}}}
\vnbprop{primitive-endpoints}{\vnbstmt{\begin{array}{@{}l@{}} \forall f, i, a, b. \\ \quad \operatorname{is\text{-}primitive}(f, i, a, b) \Rightarrow ((a \in \mathbb{R}) \wedge ((b \in \mathbb{R}) \wedge (a < b))) \end{array}}}
\vnbprop{primitive-in-fun}{\vnbstmt{\forall f, i, a, b.\; (\operatorname{is\text{-}primitive}(f, i, a, b) \Rightarrow (f \in (\operatorname{ccint}(a, b) \to \mathbb{R})))}}
\vnbprop{primitive-integrand-in-fun}{\vnbstmt{\forall f, i, a, b.\; (\operatorname{is\text{-}primitive}(f, i, a, b) \Rightarrow (i \in (\operatorname{ccint}(a, b) \to \mathbb{R})))}}
\vnbprop{primitive-continuous}{\vnbstmt{\begin{array}{@{}l@{}} \forall f, i, a, b. \\ \quad \operatorname{is\text{-}primitive}(f, i, a, b) \Rightarrow \operatorname{is\text{-}continuous\text{-}on}(f, \operatorname{ccint}(a, b)) \end{array}}}
\vnbprop{primitive-exceptional-set}{\vnbstmt{\begin{array}{@{}l@{}} \forall f, i, a, b. \\ \quad \operatorname{is\text{-}primitive}(f, i, a, b) \Rightarrow  \\ \quad \quad \exists s. \\ \quad \quad \quad \operatorname{is\text{-}countable}(s) \wedge \\ \quad \quad \quad s \subseteq \operatorname{ccint}(a, b) \wedge \\ \quad \quad \quad \forall t \in \operatorname{ooint}(a, b).\; (\neg (t \in s) \Rightarrow \operatorname{has\text{-}deriv\text{-}at}(f, t, i(t))) \end{array}}}

\vnbbookfile{theorem-library/regulated-step-approx}{theorem-library/regulated-step-approx}
\noindent{\small regulated-step-approx.scm -- Dieudonne (7.6.1), NECESSITY: a function regulated on [a, b] is the uniform limit on [a, b] of step functions. (Batch 22-A, 2026-09-23.)}\par
\vnbprop{strict-seq-append-point}{\vnbstmt{\begin{array}{@{}l@{}} \forall p, n, d. \\ \quad p \in (\mathbb{N} \to \mathbb{R}) \wedge \\ \quad n \in \mathbb{N} \wedge \\ \quad d \in \mathbb{R} \wedge \\ \quad \forall i, j. \\ \quad \quad ((i \in \operatorname{interval}(0, n)) \wedge ((j \in \operatorname{interval}(0, n)) \wedge (i < j))) \Rightarrow (p(i) < p(j)) \wedge \\ \quad p(n) < d \Rightarrow \\ \quad \quad \exists q \in (\mathbb{N} \to \mathbb{R}). \\ \quad \quad \quad q((n + 1)) = d \wedge \\ \quad \quad \quad \forall k \in \mathbb{N}.\; ((k \leq n) \Rightarrow (q(k) = p(k))) \wedge \\ \quad \quad \quad \forall i, j. \\ \quad \quad \quad \quad ((i \in \operatorname{interval}(0, (n + 1))) \wedge ((j \in \operatorname{interval}(0, (n + 1))) \wedge (i < j))) \Rightarrow  \\ \quad \quad \quad \quad \quad q(i) < q(j) \end{array}}}
\vnbprop{fun-append-constant-piece}{\vnbstmt{\begin{array}{@{}l@{}} \forall h \in (\mathbb{R} \to \mathbb{R}), c, d, k, v \in \mathbb{R}. \\ \quad (c < d) \Rightarrow  \\ \quad \quad \exists g \in (\mathbb{R} \to \mathbb{R}). \\ \quad \quad \quad \forall t \in \mathbb{R}.\; ((t \leq c) \Rightarrow (g(t) = h(t))) \wedge \\ \quad \quad \quad \forall t \in \mathbb{R}.\; (((c < t) \wedge (t < d)) \Rightarrow (g(t) = k)) \wedge \\ \quad \quad \quad g(d) = v \end{array}}}
\vnbprop{step-pieces-append}{\vnbstmt{\begin{array}{@{}l@{}} \forall p, n, h, q, g, c, d, k.; \forall i \in \operatorname{interval}(0, (n + 1)). \\ \quad p \in (\mathbb{N} \to \mathbb{R}) \wedge \\ \quad n \in \mathbb{N} \wedge \\ \quad h \in (\mathbb{R} \to \mathbb{R}) \wedge \\ \quad q \in (\mathbb{N} \to \mathbb{R}) \wedge \\ \quad g \in (\mathbb{R} \to \mathbb{R}) \wedge \\ \quad c \in \mathbb{R} \wedge \\ \quad d \in \mathbb{R} \wedge \\ \quad k \in \mathbb{R} \wedge \\ \quad \forall \mathit{rsi}, j. \\ \quad \quad ((\mathit{rsi} \in \operatorname{interval}(0, n)) \wedge ((j \in \operatorname{interval}(0, n)) \wedge (\mathit{rsi} < j))) \Rightarrow (p(\mathit{rsi}) < p(j)) \wedge \\ \quad \forall i \in \operatorname{interval}(0, n). \\ \quad \quad (i < n) \Rightarrow  \\ \quad \quad \quad \exists \mathit{rspc} \in \mathbb{R}. \forall t \in \mathbb{R}. \\ \quad \quad \quad \quad ((p(i) < t) \wedge (t < p((i + 1)))) \Rightarrow (h(t) = \mathit{rspc}) \wedge \\ \quad p(n) = c \wedge \\ \quad \forall j \in \mathbb{N}.\; ((j \leq n) \Rightarrow (q(j) = p(j))) \wedge \\ \quad q((n + 1)) = d \wedge \\ \quad \forall \mathit{rst} \in \mathbb{R}.\; ((\mathit{rst} \leq c) \Rightarrow (g(\mathit{rst}) = h(\mathit{rst}))) \wedge \\ \quad \forall \mathit{rst} \in \mathbb{R}.\; (((c < \mathit{rst}) \wedge (\mathit{rst} < d)) \Rightarrow (g(\mathit{rst}) = k)) \wedge \\ \quad i < (n + 1) \Rightarrow \\ \quad \quad \exists \mathit{rspc} \in \mathbb{R}. \forall t \in \mathbb{R}. \\ \quad \quad \quad ((q(i) < t) \wedge (t < q((i + 1)))) \Rightarrow (g(t) = \mathit{rspc}) \end{array}}}
\vnbprop{step-approx-start}{\vnbstmt{\begin{array}{@{}l@{}} \forall a, f, e. \\ \quad ((a \in \mathbb{R}) \wedge ((f \in (\mathbb{R} \to \mathbb{R})) \wedge ((e \in \mathbb{R}) \wedge (0 \leq e)))) \Rightarrow  \\ \quad \quad \exists n \in \mathbb{N}, p \in (\mathbb{N} \to \mathbb{R}), h \in (\mathbb{R} \to \mathbb{R}). \\ \quad \quad \quad p(0) = a \wedge \\ \quad \quad \quad p(n) = a \wedge \\ \quad \quad \quad \forall i, j. \\ \quad \quad \quad \quad ((i \in \operatorname{interval}(0, n)) \wedge ((j \in \operatorname{interval}(0, n)) \wedge (i < j))) \Rightarrow (p(i) < p(j)) \wedge \\ \quad \quad \quad \forall \mathit{rspi} \in \operatorname{interval}(0, n). \\ \quad \quad \quad \quad (\mathit{rspi} < n) \Rightarrow  \\ \quad \quad \quad \quad \quad \exists c \in \mathbb{R}. \forall t \in \mathbb{R}. \\ \quad \quad \quad \quad \quad \quad ((p(\mathit{rspi}) < t) \wedge (t < p((\mathit{rspi} + 1)))) \Rightarrow (h(t) = c) \wedge \\ \quad \quad \quad \forall z \in \operatorname{ccint}(a, a).\; (\lvert (f(z) - h(z)) \rvert \leq e) \end{array}}}
\vnbprop{step-approx-append}{\vnbstmt{\begin{array}{@{}l@{}} \forall a, x, y, f, e, k. \\ \quad a \in \mathbb{R} \wedge \\ \quad x \in \mathbb{R} \wedge \\ \quad y \in \mathbb{R} \wedge \\ \quad f \in (\mathbb{R} \to \mathbb{R}) \wedge \\ \quad e \in \mathbb{R} \wedge \\ \quad k \in \mathbb{R} \wedge \\ \quad 0 \leq e \wedge \\ \quad \exists n \in \mathbb{N}, p \in (\mathbb{N} \to \mathbb{R}), h \in (\mathbb{R} \to \mathbb{R}). \\ \quad \quad p(0) = a \wedge \\ \quad \quad p(n) = x \wedge \\ \quad \quad \forall i, j. \\ \quad \quad \quad ((i \in \operatorname{interval}(0, n)) \wedge ((j \in \operatorname{interval}(0, n)) \wedge (i < j))) \Rightarrow (p(i) < p(j)) \wedge \\ \quad \quad \forall \mathit{rspi} \in \operatorname{interval}(0, n). \\ \quad \quad \quad (\mathit{rspi} < n) \Rightarrow  \\ \quad \quad \quad \quad \exists c \in \mathbb{R}. \forall t \in \mathbb{R}. \\ \quad \quad \quad \quad \quad ((p(\mathit{rspi}) < t) \wedge (t < p((\mathit{rspi} + 1)))) \Rightarrow (h(t) = c) \wedge \\ \quad \quad \forall z \in \operatorname{ccint}(a, x).\; (\lvert (f(z) - h(z)) \rvert \leq e) \wedge \\ \quad x < y \wedge \\ \quad \forall s \in \mathbb{R}.\; (((x < s) \wedge (s < y)) \Rightarrow (\lvert (f(s) - k) \rvert \leq e)) \Rightarrow \\ \quad \quad \exists n \in \mathbb{N}, p \in (\mathbb{N} \to \mathbb{R}), h \in (\mathbb{R} \to \mathbb{R}). \\ \quad \quad \quad p(0) = a \wedge \\ \quad \quad \quad p(n) = y \wedge \\ \quad \quad \quad \forall i, j. \\ \quad \quad \quad \quad ((i \in \operatorname{interval}(0, n)) \wedge ((j \in \operatorname{interval}(0, n)) \wedge (i < j))) \Rightarrow (p(i) < p(j)) \wedge \\ \quad \quad \quad \forall \mathit{rspi} \in \operatorname{interval}(0, n). \\ \quad \quad \quad \quad (\mathit{rspi} < n) \Rightarrow  \\ \quad \quad \quad \quad \quad \exists c \in \mathbb{R}. \forall t \in \mathbb{R}. \\ \quad \quad \quad \quad \quad \quad ((p(\mathit{rspi}) < t) \wedge (t < p((\mathit{rspi} + 1)))) \Rightarrow (h(t) = c) \wedge \\ \quad \quad \quad \forall z \in \operatorname{ccint}(a, y).\; (\lvert (f(z) - h(z)) \rvert \leq e) \end{array}}}
\vnbprop{regulated-on-right-window}{\vnbstmt{\begin{array}{@{}l@{}} \forall a, b, f, g, e, t. \\ \quad \operatorname{is\text{-}regulated\text{-}on}(f, a, b) \wedge \\ \quad g \in (\mathbb{R} \to \mathbb{R}) \wedge \\ \quad \forall z \in \operatorname{ccint}(a, b).\; (g(z) = f(z)) \wedge \\ \quad \operatorname{pos\text{-}rr}(e) \wedge \\ \quad t \in \operatorname{ccint}(a, b) \Rightarrow \\ \quad \quad \exists d \in \mathbb{R}. \\ \quad \quad \quad 0 < d \wedge \\ \quad \quad \quad \exists l \in \mathbb{R}. \forall s \in \mathbb{R}. \\ \quad \quad \quad \quad ((t < s) \wedge ((s < (t + d)) \wedge (s \leq b))) \Rightarrow (\lvert (g(s) - l) \rvert \leq e) \end{array}}}
\vnbprop{regulated-on-left-window}{\vnbstmt{\begin{array}{@{}l@{}} \forall a, b, f, g, e, t. \\ \quad \operatorname{is\text{-}regulated\text{-}on}(f, a, b) \wedge \\ \quad g \in (\mathbb{R} \to \mathbb{R}) \wedge \\ \quad \forall z \in \operatorname{ccint}(a, b).\; (g(z) = f(z)) \wedge \\ \quad \operatorname{pos\text{-}rr}(e) \wedge \\ \quad t \in \operatorname{ccint}(a, b) \Rightarrow \\ \quad \quad \exists d \in \mathbb{R}. \\ \quad \quad \quad 0 < d \wedge \\ \quad \quad \quad \exists l \in \mathbb{R}. \forall s \in \mathbb{R}. \\ \quad \quad \quad \quad ((a \leq s) \wedge (((t - d) < s) \wedge (s < t))) \Rightarrow (\lvert (g(s) - l) \rvert \leq e) \end{array}}}
\vnbprop{regulated-on-step-approx-creep}{\vnbstmt{\begin{array}{@{}l@{}} \forall a, b, f, g, e. \\ \quad \operatorname{is\text{-}regulated\text{-}on}(f, a, b) \wedge \\ \quad g \in (\mathbb{R} \to \mathbb{R}) \wedge \\ \quad \forall z \in \operatorname{ccint}(a, b).\; (g(z) = f(z)) \wedge \\ \quad \operatorname{pos\text{-}rr}(e) \Rightarrow \\ \quad \quad \exists n \in \mathbb{N}, p \in (\mathbb{N} \to \mathbb{R}), h \in (\mathbb{R} \to \mathbb{R}). \\ \quad \quad \quad p(0) = a \wedge \\ \quad \quad \quad p(n) = b \wedge \\ \quad \quad \quad \forall i, j. \\ \quad \quad \quad \quad ((i \in \operatorname{interval}(0, n)) \wedge ((j \in \operatorname{interval}(0, n)) \wedge (i < j))) \Rightarrow (p(i) < p(j)) \wedge \\ \quad \quad \quad \forall \mathit{rspi} \in \operatorname{interval}(0, n). \\ \quad \quad \quad \quad (\mathit{rspi} < n) \Rightarrow  \\ \quad \quad \quad \quad \quad \exists c \in \mathbb{R}. \forall t \in \mathbb{R}. \\ \quad \quad \quad \quad \quad \quad ((p(\mathit{rspi}) < t) \wedge (t < p((\mathit{rspi} + 1)))) \Rightarrow (h(t) = c) \wedge \\ \quad \quad \quad \forall \mathit{rsqz} \in \operatorname{ccint}(a, b).\; (\lvert (g(\mathit{rsqz}) - h(\mathit{rsqz})) \rvert \leq e) \end{array}}}
\vnbprop{step-approx-q-step-fn}{\vnbstmt{\begin{array}{@{}l@{}} \forall a, b, f, e. \\ \quad a \in \mathbb{R} \wedge \\ \quad b \in \mathbb{R} \wedge \\ \quad a < b \wedge \\ \quad \exists n \in \mathbb{N}, p \in (\mathbb{N} \to \mathbb{R}), h \in (\mathbb{R} \to \mathbb{R}). \\ \quad \quad p(0) = a \wedge \\ \quad \quad p(n) = b \wedge \\ \quad \quad \forall i, j. \\ \quad \quad \quad ((i \in \operatorname{interval}(0, n)) \wedge ((j \in \operatorname{interval}(0, n)) \wedge (i < j))) \Rightarrow (p(i) < p(j)) \wedge \\ \quad \quad \forall \mathit{rspi} \in \operatorname{interval}(0, n). \\ \quad \quad \quad (\mathit{rspi} < n) \Rightarrow  \\ \quad \quad \quad \quad \exists c \in \mathbb{R}. \forall t \in \mathbb{R}. \\ \quad \quad \quad \quad \quad ((p(\mathit{rspi}) < t) \wedge (t < p((\mathit{rspi} + 1)))) \Rightarrow (h(t) = c) \wedge \\ \quad \quad \forall z \in \operatorname{ccint}(a, b).\; (\lvert (f(z) - h(z)) \rvert \leq e) \Rightarrow \\ \quad \quad \exists \mathit{rsth} \in (\mathbb{R} \to \mathbb{R}). \\ \quad \quad \quad \operatorname{is\text{-}step\text{-}fn}(\mathit{rsth}, a, b) \wedge \forall z \in \operatorname{ccint}(a, b).\; (\lvert (f(z) - \operatorname{rsth}(z)) \rvert \leq e) \end{array}}}
\vnbprop{regulated-on-step-approx}{\vnbstmt{\begin{array}{@{}l@{}} \forall a, b, f, e. \\ \quad (\operatorname{is\text{-}regulated\text{-}on}(f, a, b) \wedge \operatorname{pos\text{-}rr}(e)) \Rightarrow  \\ \quad \quad \exists h \in (\mathbb{R} \to \mathbb{R}). \\ \quad \quad \quad (\operatorname{is\text{-}step\text{-}fn}(h, a, b) \wedge \forall z \in \operatorname{ccint}(a, b).\; (\lvert (\operatorname{extend\text{-}const}(f, a, b)(z) - h(z)) \rvert \leq e)) \end{array}}}
\vnbprop{regulated-on-is-step-limit}{\vnbstmt{\begin{array}{@{}l@{}} \forall a, b \in \mathbb{R}, f. \\ \quad \operatorname{is\text{-}regulated\text{-}on}(f, a, b) \Rightarrow  \\ \quad \quad \exists m \in (\mathbb{N} \to (\mathbb{R} \to \mathbb{R})). \\ \quad \quad \quad (\forall k \in \mathbb{N}.\; \operatorname{is\text{-}step\text{-}fn}(m(k), a, b) \wedge \operatorname{is\text{-}unif\text{-}limit\text{-}on}(m, \operatorname{extend\text{-}const}(f, a, b), a, b)) \end{array}}}

\vnbbookfile{theorem-library/has-deriv-at-more}{theorem-library/has-deriv-at-more}
\noindent{\small has-deriv-at-more.scm -- the four HAS-DERIV-AT laws that theorem-library/interval-calculus-laws.scm did not prove: the VALUE typing, the AFFINE function (and with it the constant and the identity), the REAL MULTIPLE and the DIFFERENCE.}\par
\vnbprop{has-deriv-at-value-in-rr}{\vnbstmt{\forall f, x, l.\; (\operatorname{has\text{-}deriv\text{-}at}(f, x, l) \Rightarrow (f(x) \in \mathbb{R}))}}
\vnbprop{has-deriv-at-affine}{\vnbstmt{\begin{array}{@{}l@{}} \forall r, f.; \forall c, m, x \in \mathbb{R}. \\ \quad \operatorname{pos\text{-}rr}(r) \wedge \\ \quad \operatorname{restrict}(f, \operatorname{ooint}((x - r), (x + r))) \in (\operatorname{ooint}((x - r), (x + r)) \to \mathbb{R}) \wedge \\ \quad \forall \mathit{hbx} \in \operatorname{ooint}((x - r), (x + r)).\; (f(\mathit{hbx}) \simeq (c + (m \cdot \mathit{hbx}))) \Rightarrow \\ \quad \quad \operatorname{has\text{-}deriv\text{-}at}(f, x, m) \end{array}}}
\vnbprop{has-deriv-at-const}{\vnbstmt{\begin{array}{@{}l@{}} \forall r, f.; \forall c, x \in \mathbb{R}. \\ \quad \operatorname{pos\text{-}rr}(r) \wedge \\ \quad \operatorname{restrict}(f, \operatorname{ooint}((x - r), (x + r))) \in (\operatorname{ooint}((x - r), (x + r)) \to \mathbb{R}) \wedge \\ \quad \forall \mathit{hbx} \in \operatorname{ooint}((x - r), (x + r)).\; (f(\mathit{hbx}) \simeq c) \Rightarrow \\ \quad \quad \operatorname{has\text{-}deriv\text{-}at}(f, x, 0) \end{array}}}
\vnbprop{has-deriv-at-identity}{\vnbstmt{\begin{array}{@{}l@{}} \forall r, f.; \forall x \in \mathbb{R}. \\ \quad \operatorname{pos\text{-}rr}(r) \wedge \\ \quad \operatorname{restrict}(f, \operatorname{ooint}((x - r), (x + r))) \in (\operatorname{ooint}((x - r), (x + r)) \to \mathbb{R}) \wedge \\ \quad \forall \mathit{hbx} \in \operatorname{ooint}((x - r), (x + r)).\; (f(\mathit{hbx}) \simeq \mathit{hbx}) \Rightarrow \\ \quad \quad \operatorname{has\text{-}deriv\text{-}at}(f, x, 1) \end{array}}}
\vnbprop{has-deriv-at-real-mul}{\vnbstmt{\begin{array}{@{}l@{}} \forall r, g, h.; \forall c \in \mathbb{R}, x, m. \\ \quad \operatorname{pos\text{-}rr}(r) \wedge \\ \quad \operatorname{restrict}(h, \operatorname{ooint}((x - r), (x + r))) \in (\operatorname{ooint}((x - r), (x + r)) \to \mathbb{R}) \wedge \\ \quad \forall \mathit{hbx} \in \operatorname{ooint}((x - r), (x + r)).\; (h(\mathit{hbx}) \simeq (c \cdot g(\mathit{hbx}))) \wedge \\ \quad \operatorname{has\text{-}deriv\text{-}at}(g, x, m) \Rightarrow \\ \quad \quad \operatorname{has\text{-}deriv\text{-}at}(h, x, (c \cdot m)) \end{array}}}
\vnbprop{has-deriv-at-difference}{\vnbstmt{\begin{array}{@{}l@{}} \forall r, f, g, h, x, l, m. \\ \quad \operatorname{pos\text{-}rr}(r) \wedge \\ \quad \operatorname{restrict}(h, \operatorname{ooint}((x - r), (x + r))) \in (\operatorname{ooint}((x - r), (x + r)) \to \mathbb{R}) \wedge \\ \quad \operatorname{restrict}(g, \operatorname{ooint}((x - r), (x + r))) \in (\operatorname{ooint}((x - r), (x + r)) \to \mathbb{R}) \wedge \\ \quad \forall \mathit{hbx} \in \operatorname{ooint}((x - r), (x + r)).\; (h(\mathit{hbx}) \simeq (f(\mathit{hbx}) - g(\mathit{hbx}))) \wedge \\ \quad \operatorname{has\text{-}deriv\text{-}at}(f, x, l) \wedge \\ \quad \operatorname{has\text{-}deriv\text{-}at}(g, x, m) \Rightarrow \\ \quad \quad \operatorname{has\text{-}deriv\text{-}at}(h, x, (l - m)) \end{array}}}

\vnbbookfile{theorem-library/has-deriv-at-chain}{theorem-library/has-deriv-at-chain}
\noindent{\small has-deriv-at-chain.scm -- THE CHAIN RULE FOR THE LOCAL DERIVATIVE, and the corollary at an affine inner map.}\par
\vnbprop{has-deriv-at-maps-into}{\vnbstmt{\begin{array}{@{}l@{}} \forall f, x, l, d, b. \\ \quad (\operatorname{has\text{-}deriv\text{-}at}(f, x, l) \wedge (\operatorname{pos\text{-}rr}(d) \wedge \operatorname{pos\text{-}rr}(b))) \Rightarrow  \\ \quad \quad \exists w. \\ \quad \quad \quad (\operatorname{pos\text{-}rr}(w) \wedge ((w \leq b) \wedge \forall y \in \operatorname{ooint}((x - w), (x + w)).\; (f(y) \in \operatorname{ooint}((f(x) - d), (f(x) + d))))) \end{array}}}
\vnbprop{has-deriv-at-chain}{\vnbstmt{\begin{array}{@{}l@{}} \forall r, f, g, h, x, l, m. \\ \quad \operatorname{pos\text{-}rr}(r) \wedge \\ \quad \operatorname{restrict}(h, \operatorname{ooint}((x - r), (x + r))) \in (\operatorname{ooint}((x - r), (x + r)) \to \mathbb{R}) \wedge \\ \quad \forall y \in \operatorname{ooint}((x - r), (x + r)).\; (h(y) \simeq g(f(y))) \wedge \\ \quad \operatorname{has\text{-}deriv\text{-}at}(f, x, l) \wedge \\ \quad \operatorname{has\text{-}deriv\text{-}at}(g, f(x), m) \Rightarrow \\ \quad \quad \operatorname{has\text{-}deriv\text{-}at}(h, x, (m \cdot l)) \end{array}}}
\vnbprop{has-deriv-at-affine-chain}{\vnbstmt{\begin{array}{@{}l@{}} \forall r, g, h.; \forall p, q, x \in \mathbb{R}, m. \\ \quad \operatorname{pos\text{-}rr}(r) \wedge \\ \quad \operatorname{restrict}(h, \operatorname{ooint}((x - r), (x + r))) \in (\operatorname{ooint}((x - r), (x + r)) \to \mathbb{R}) \wedge \\ \quad \forall y \in \operatorname{ooint}((x - r), (x + r)).\; (h(y) \simeq g((p + (q \cdot y)))) \wedge \\ \quad \operatorname{has\text{-}deriv\text{-}at}(g, (p + (q \cdot x)), m) \Rightarrow \\ \quad \quad \operatorname{has\text{-}deriv\text{-}at}(h, x, (m \cdot q)) \end{array}}}

\vnbbookfile{theorem-library/interval-taylor}{theorem-library/interval-taylor}
\noindent{\small interval-taylor.scm -- TAYLOR'S FORMULA (2.18) FOR A FUNCTION ON ITS INTERVAL, AND ITS TWO COROLLARIES.}\par
\vnbprop{gmvt-with-known-derivative}{\vnbstmt{\begin{array}{@{}l@{}} \forall a, b \in \mathbb{R}, f, g.; \forall v \in (\operatorname{ooint}(a, b) \to \mathbb{R}). \\ \quad a < b \wedge \\ \quad f \in (\operatorname{ccint}(a, b) \to \mathbb{R}) \wedge \\ \quad \operatorname{is\text{-}continuous\text{-}on}(f, \operatorname{ccint}(a, b)) \wedge \\ \quad \forall x \in \operatorname{ooint}(a, b).\; \operatorname{has\text{-}deriv\text{-}at}(f, x, v(x)) \wedge \\ \quad g \in (\operatorname{ccint}(a, b) \to \mathbb{R}) \wedge \\ \quad \operatorname{is\text{-}continuous\text{-}on}(g, \operatorname{ccint}(a, b)) \wedge \\ \quad \forall x \in \operatorname{ooint}(a, b).\; \exists m.\; \operatorname{has\text{-}deriv\text{-}at}(g, x, m) \Rightarrow \\ \quad \quad \exists i \in \operatorname{ooint}(a, b), m. \\ \quad \quad \quad \operatorname{has\text{-}deriv\text{-}at}(g, i, m) \wedge ((v(i) \cdot (g(b) - g(a))) = (m \cdot (f(b) - f(a)))) \end{array}}}
\vnbprop{taylor-family-order-in-nn}{\vnbstmt{\forall d, n, a, b.\; (\operatorname{is\text{-}taylor\text{-}family}(d, n, a, b) \Rightarrow (n \in \mathbb{N}))}}
\vnbprop{taylor-family-order-pos}{\vnbstmt{\forall d, n, a, b.\; (\operatorname{is\text{-}taylor\text{-}family}(d, n, a, b) \Rightarrow (1 \leq n))}}
\vnbprop{taylor-family-left-in-rr}{\vnbstmt{\forall d, n, a, b.\; (\operatorname{is\text{-}taylor\text{-}family}(d, n, a, b) \Rightarrow (a \in \mathbb{R}))}}
\vnbprop{taylor-family-right-in-rr}{\vnbstmt{\forall d, n, a, b.\; (\operatorname{is\text{-}taylor\text{-}family}(d, n, a, b) \Rightarrow (b \in \mathbb{R}))}}
\vnbprop{taylor-family-endpoints-lt}{\vnbstmt{\forall d, n, a, b.\; (\operatorname{is\text{-}taylor\text{-}family}(d, n, a, b) \Rightarrow (a < b))}}
\vnbprop{taylor-family-top-in-fun}{\vnbstmt{\begin{array}{@{}l@{}} \forall d, n, a, b. \\ \quad \operatorname{is\text{-}taylor\text{-}family}(d, n, a, b) \Rightarrow (d(n) \in (\operatorname{ooint}(a, b) \to \mathbb{R})) \end{array}}}
\vnbprop{taylor-family-in-fun}{\vnbstmt{\begin{array}{@{}l@{}} \forall d, n, a, b.; \forall k \in \mathbb{N}. \\ \quad (\operatorname{is\text{-}taylor\text{-}family}(d, n, a, b) \wedge (k < n)) \Rightarrow (d(k) \in (\operatorname{ccint}(a, b) \to \mathbb{R})) \end{array}}}
\vnbprop{taylor-family-continuous}{\vnbstmt{\begin{array}{@{}l@{}} \forall d, n, a, b.; \forall k \in \mathbb{N}. \\ \quad (\operatorname{is\text{-}taylor\text{-}family}(d, n, a, b) \wedge (k < n)) \Rightarrow  \\ \quad \quad \operatorname{is\text{-}continuous\text{-}on}(d(k), \operatorname{ccint}(a, b)) \end{array}}}
\vnbprop{taylor-family-deriv}{\vnbstmt{\begin{array}{@{}l@{}} \forall d, n, a, b.; \forall k \in \mathbb{N}, x \in \operatorname{ooint}(a, b). \\ \quad (\operatorname{is\text{-}taylor\text{-}family}(d, n, a, b) \wedge (k < n)) \Rightarrow  \\ \quad \quad \operatorname{has\text{-}deriv\text{-}at}(d(k), x, d((k + 1))(x)) \end{array}}}
\vnbprop{taylor-aux-value}{\vnbstmt{\begin{array}{@{}l@{}} \forall d, n, a, b.; \forall z \in \operatorname{ccint}(a, b). \\ \quad (\operatorname{taylor\text{-}aux}(d, n, a, b)(z) \simeq \operatorname{taylor\text{-}sum}(d, n, b, z)) \end{array}}}
\vnbprop{taylor-clear-recip}{\vnbstmt{\begin{array}{@{}l@{}} \forall p, q, r, s, u, w \in \mathbb{R}. \\ \quad (((s \cdot r) = 1) \wedge ((((p \cdot q) \cdot r) \cdot w) = u)) \Rightarrow  \\ \quad \quad ((u \cdot s) = ((p \cdot q) \cdot w)) \wedge (u = (((p \cdot r) \cdot q) \cdot w)) \end{array}}}
\vnbprop{series-partial-sum-in-rr-below}{\vnbstmt{\begin{array}{@{}l@{}} \forall n \in \mathbb{N}, f. \\ \quad \forall k \in \mathbb{N}.\; ((k < n) \Rightarrow (f(k) \in \mathbb{R})) \Rightarrow  \\ \quad \quad \operatorname{series\text{-}partial\text{-}sum}(f, n) \in \mathbb{R} \end{array}}}
\vnbprop{taylor-term-value}{\vnbstmt{\begin{array}{@{}l@{}} \forall d, b, v.; \forall k \in \mathbb{N}. \\ \quad (\operatorname{taylor\text{-}term}(d, b, v)(k) \simeq ((d(k)(v) \cdot \operatorname{power}((b - v), k)) \cdot \operatorname{recip}(\operatorname{factorial}(k)))) \end{array}}}
\vnbprop{taylor-sum-unfold}{\vnbstmt{\begin{array}{@{}l@{}} \forall d, n, b, v. \\ \quad (\operatorname{taylor\text{-}sum}(d, n, b, v) \simeq \operatorname{series\text{-}partial\text{-}sum}(\operatorname{taylor\text{-}term}(d, b, v), n)) \end{array}}}
\vnbprop{taylor-sum-in-rr}{\vnbstmt{\begin{array}{@{}l@{}} \forall d, n, a, b.; \forall j \in \mathbb{N}, z \in \operatorname{ccint}(a, b). \\ \quad (\operatorname{is\text{-}taylor\text{-}family}(d, n, a, b) \wedge (j \leq n)) \Rightarrow (\operatorname{taylor\text{-}sum}(d, j, b, z) \in \mathbb{R}) \end{array}}}
\vnbprop{taylor-aux-in-fun}{\vnbstmt{\begin{array}{@{}l@{}} \forall d, n, a, b.; \forall j \in \mathbb{N}. \\ \quad (\operatorname{is\text{-}taylor\text{-}family}(d, n, a, b) \wedge (j \leq n)) \Rightarrow  \\ \quad \quad \operatorname{taylor\text{-}aux}(d, j, a, b) \in (\operatorname{ccint}(a, b) \to \mathbb{R}) \end{array}}}
\vnbprop{taylor-sum-one}{\vnbstmt{\begin{array}{@{}l@{}} \forall d.; \forall b, v \in \mathbb{R}. \\ \quad (d(0)(v) \in \mathbb{R}) \Rightarrow (\operatorname{taylor\text{-}sum}(d, 1, b, v) = d(0)(v)) \end{array}}}
\vnbprop{taylor-sum-at-endpoint}{\vnbstmt{\begin{array}{@{}l@{}} \forall m \in \mathbb{N}, d.; \forall b \in \mathbb{R}. \\ \quad \forall k \in \mathbb{N}.\; ((k \leq m) \Rightarrow (d(k)(b) \in \mathbb{R})) \Rightarrow  \\ \quad \quad \operatorname{taylor\text{-}sum}(d, (m + 1), b, b) = d(0)(b) \end{array}}}
\vnbprop{taylor-family-sum-at-endpoint}{\vnbstmt{\begin{array}{@{}l@{}} \forall d, n, a, b.; \forall m \in \mathbb{N}. \\ \quad (\operatorname{is\text{-}taylor\text{-}family}(d, n, a, b) \wedge ((m + 1) \leq n)) \Rightarrow  \\ \quad \quad \operatorname{taylor\text{-}sum}(d, (m + 1), b, b) = d(0)(b) \end{array}}}
\vnbprop{cont-on-in-fun}{\vnbstmt{\forall f, m.\; (\operatorname{is\text{-}continuous\text{-}on}(f, m) \Rightarrow (f \in (m \to \mathbb{R})))}}
\vnbprop{cont-on-transfer-ptwise-eq}{\vnbstmt{\begin{array}{@{}l@{}} \forall m, f, g. \\ \quad ((m \subseteq \mathbb{R}) \wedge ((f \in (m \to \mathbb{R})) \wedge (\operatorname{is\text{-}continuous\text{-}on}(g, m) \wedge \forall z \in m.\; (f(z) = g(z))))) \Rightarrow  \\ \quad \quad \operatorname{is\text{-}continuous\text{-}on}(f, m) \end{array}}}
\vnbprop{cont-on-of-total}{\vnbstmt{\begin{array}{@{}l@{}} \forall m.; \forall g \in (\mathbb{R} \to \mathbb{R}), f \in (m \to \mathbb{R}). \\ \quad m \subseteq \mathbb{R} \wedge \\ \quad \forall y \in \mathbb{R}.\; \operatorname{is\text{-}continuous\text{-}at}(\operatorname{rr\text{-}ms}, \operatorname{rr\text{-}ms}, g, y) \wedge \\ \quad \forall z \in m.\; (f(z) = g(z)) \Rightarrow \\ \quad \quad \operatorname{is\text{-}continuous\text{-}on}(f, m) \end{array}}}
\vnbprop{power-sub-succ-continuous}{\vnbstmt{\begin{array}{@{}l@{}} \forall b \in \mathbb{R}, k \in \mathbb{N}, y \in \mathbb{R}. \\ \quad \operatorname{is\text{-}continuous\text{-}at}(\operatorname{rr\text{-}ms}, \operatorname{rr\text{-}ms}, (\lambda z \in \mathbb{R}.\; \operatorname{power}((b - z), (k + 1))), y) \end{array}}}
\vnbprop{taylor-step-continuous}{\vnbstmt{\begin{array}{@{}l@{}} \forall b, c \in \mathbb{R}, k \in \mathbb{N}, g, f \in (\mathbb{R} \to \mathbb{R}), y \in \mathbb{R}. \\ \quad (\operatorname{is\text{-}continuous\text{-}at}(\operatorname{rr\text{-}ms}, \operatorname{rr\text{-}ms}, g, y) \wedge \operatorname{is\text{-}continuous\text{-}at}(\operatorname{rr\text{-}ms}, \operatorname{rr\text{-}ms}, f, y)) \Rightarrow  \\ \quad \quad \operatorname{is\text{-}continuous\text{-}at}(\operatorname{rr\text{-}ms}, \operatorname{rr\text{-}ms}, (\lambda w \in \mathbb{R}.\; (g(w) + ((f(w) \cdot \operatorname{power}((b - w), (k + 1))) \cdot c))), y) \end{array}}}
\vnbprop{taylor-aux-step}{\vnbstmt{\begin{array}{@{}l@{}} \forall d, n, a, b.; \forall j \in \mathbb{N}, z \in \operatorname{ccint}(a, b). \\ \quad (\operatorname{is\text{-}taylor\text{-}family}(d, n, a, b) \wedge ((j + 1) \leq n)) \Rightarrow  \\ \quad \quad \operatorname{taylor\text{-}aux}(d, (j + 1), a, b)(z) \\ \quad \quad =\; \operatorname{taylor\text{-}aux}(d, j, a, b)(z) + ((d(j)(z) \cdot \operatorname{power}((b - z), j)) \cdot \operatorname{recip}(\operatorname{factorial}(j))) \end{array}}}
\vnbprop{taylor-aux-continuous-on}{\vnbstmt{\begin{array}{@{}l@{}} \forall m \in \mathbb{N}, d, n, a, b. \\ \quad (\operatorname{is\text{-}taylor\text{-}family}(d, n, a, b) \wedge ((m + 1) \leq n)) \Rightarrow  \\ \quad \quad \operatorname{is\text{-}continuous\text{-}on}(\operatorname{taylor\text{-}aux}(d, (m + 1), a, b), \operatorname{ccint}(a, b)) \end{array}}}
\vnbprop{ooint-interior-radius}{\vnbstmt{\begin{array}{@{}l@{}} \forall a, b \in \mathbb{R}, x \in \operatorname{ooint}(a, b). \exists r. \\ \quad \operatorname{pos\text{-}rr}(r) \wedge (\operatorname{ooint}((x - r), (x + r)) \subseteq \operatorname{ooint}(a, b)) \end{array}}}
\vnbprop{power-sub-succ-has-deriv}{\vnbstmt{\begin{array}{@{}l@{}} \forall b \in \mathbb{R}, k \in \mathbb{N}, x \in \mathbb{R}. \\ \quad \operatorname{has\text{-}deriv\text{-}at}((\lambda z \in \mathbb{R}.\; \operatorname{power}((b - z), (k + 1))), x, (((k + 1) \cdot \operatorname{power}((b - x), k)) \cdot -1)) \end{array}}}
\vnbprop{taylor-step-identity}{\vnbstmt{\begin{array}{@{}l@{}} \forall a, b, p, q, r, s \in \mathbb{R}. \\ \quad (((a \cdot p) \cdot (s \cdot r)) + (r \cdot ((b \cdot q) + (((s \cdot p) \cdot -1) \cdot a)))) = ((b \cdot q) \cdot r) \end{array}}}
\vnbprop{taylor-aux-has-deriv}{\vnbstmt{\begin{array}{@{}l@{}} \forall m \in \mathbb{N}, d, n, a, b.; \forall x \in \operatorname{ooint}(a, b). \\ \quad (\operatorname{is\text{-}taylor\text{-}family}(d, n, a, b) \wedge ((m + 1) \leq n)) \Rightarrow  \\ \quad \quad \operatorname{has\text{-}deriv\text{-}at}(\operatorname{taylor\text{-}aux}(d, (m + 1), a, b), x, ((d((m + 1))(x) \cdot \operatorname{power}((b - x), m)) \cdot \operatorname{recip}(\operatorname{factorial}(m)))) \end{array}}}
\vnbprop{taylor-formula-on-interval-from-aux}{\vnbstmt{\begin{array}{@{}l@{}} \forall a, b \in \mathbb{R}, d.; \forall m \in \mathbb{N}, g \in (\operatorname{ccint}(a, b) \to \mathbb{R}). \\ \quad \operatorname{is\text{-}taylor\text{-}family}(d, (m + 1), a, b) \wedge \\ \quad \operatorname{is\text{-}continuous\text{-}on}(\operatorname{taylor\text{-}aux}(d, (m + 1), a, b), \operatorname{ccint}(a, b)) \wedge \\ \quad \forall x \in \operatorname{ooint}(a, b). \\ \quad \quad \operatorname{has\text{-}deriv\text{-}at}(\operatorname{taylor\text{-}aux}(d, (m + 1), a, b), x, ((d((m + 1))(x) \cdot \operatorname{power}((b - x), m)) \cdot \operatorname{recip}(\operatorname{factorial}(m)))) \wedge \\ \quad \operatorname{taylor\text{-}sum}(d, (m + 1), b, b) = d(0)(b) \wedge \\ \quad \operatorname{is\text{-}continuous\text{-}on}(g, \operatorname{ccint}(a, b)) \wedge \\ \quad \forall x \in \operatorname{ooint}(a, b).\; \exists l.\; \operatorname{has\text{-}deriv\text{-}at}(g, x, l) \Rightarrow \\ \quad \quad \exists i \in \operatorname{ooint}(a, b), u. \\ \quad \quad \quad \operatorname{has\text{-}deriv\text{-}at}(g, i, u) \wedge \\ \quad \quad \quad ((u \cdot \operatorname{taylor\text{-}rem}(d, (m + 1), a, b)) \cdot \operatorname{factorial}(m)) = ((d((m + 1))(i) \cdot \operatorname{power}((b - i), m)) \cdot (g(b) - g(a))) \wedge \\ \quad \quad \quad u \cdot \operatorname{taylor\text{-}rem}(d, (m + 1), a, b) \\ \quad \quad \quad =\; ((d((m + 1))(i) \cdot \operatorname{recip}(\operatorname{factorial}(m))) \cdot \operatorname{power}((b - i), m)) \cdot (g(b) - g(a)) \end{array}}}
\vnbprop{taylor-formula-on-interval}{\vnbstmt{\begin{array}{@{}l@{}} \forall a, b \in \mathbb{R}, d.; \forall m \in \mathbb{N}, g \in (\operatorname{ccint}(a, b) \to \mathbb{R}). \\ \quad \operatorname{is\text{-}taylor\text{-}family}(d, (m + 1), a, b) \wedge \\ \quad \operatorname{is\text{-}continuous\text{-}on}(g, \operatorname{ccint}(a, b)) \wedge \\ \quad \forall x \in \operatorname{ooint}(a, b).\; \exists l.\; \operatorname{has\text{-}deriv\text{-}at}(g, x, l) \Rightarrow \\ \quad \quad \exists i \in \operatorname{ooint}(a, b), u. \\ \quad \quad \quad \operatorname{has\text{-}deriv\text{-}at}(g, i, u) \wedge \\ \quad \quad \quad ((u \cdot \operatorname{taylor\text{-}rem}(d, (m + 1), a, b)) \cdot \operatorname{factorial}(m)) = ((d((m + 1))(i) \cdot \operatorname{power}((b - i), m)) \cdot (g(b) - g(a))) \wedge \\ \quad \quad \quad u \cdot \operatorname{taylor\text{-}rem}(d, (m + 1), a, b) \\ \quad \quad \quad =\; ((d((m + 1))(i) \cdot \operatorname{recip}(\operatorname{factorial}(m))) \cdot \operatorname{power}((b - i), m)) \cdot (g(b) - g(a)) \end{array}}}
\vnbprop{taylor-lagrange-on-interval}{\vnbstmt{\begin{array}{@{}l@{}} \forall a, b \in \mathbb{R}, d.; \forall m \in \mathbb{N}. \\ \quad \operatorname{is\text{-}taylor\text{-}family}(d, (m + 1), a, b) \Rightarrow  \\ \quad \quad \exists i \in \operatorname{ooint}(a, b). \\ \quad \quad \quad \operatorname{taylor\text{-}rem}(d, (m + 1), a, b) = ((d((m + 1))(i) \cdot \operatorname{recip}(\operatorname{factorial}((m + 1)))) \cdot \operatorname{power}((b - a), (m + 1))) \end{array}}}

\vnbbookfile{theorem-library/zero-deriv-off-countable}{theorem-library/zero-deriv-off-countable}
\noindent{\small zero-deriv-off-countable.scm -- CONSTANCY AND THE MEAN VALUE INEQUALITY WITH A COUNTABLE EXCEPTIONAL SET (batch 20-A, 2026-09-23).}\par
\vnbprop{pw-lt-ne}{\vnbstmt{\forall u, v \in \mathbb{R}.\; ((u < v) \Rightarrow (\neg (u = v) \wedge \neg (v = u)))}}
\vnbprop{nn-antitone-step-implies-le}{\vnbstmt{\begin{array}{@{}l@{}} \forall f \in (\mathbb{N} \to \mathbb{R}), n, m \in \mathbb{N}. \\ \quad (\forall k \in \mathbb{N}.\; (f((k + 1)) \leq f(k)) \wedge (m \leq n)) \Rightarrow (f(n) \leq f(m)) \end{array}}}
\vnbprop{rr-nested-intervals}{\vnbstmt{\begin{array}{@{}l@{}} \forall o, i \in (\mathbb{N} \to \mathbb{R}). \\ \quad \forall k \in \mathbb{N}.\; (o(k) \leq o((k + 1))) \wedge \\ \quad \forall k \in \mathbb{N}.\; (i((k + 1)) \leq i(k)) \wedge \\ \quad \forall k \in \mathbb{N}.\; (o(k) \leq i(k)) \Rightarrow \\ \quad \quad \exists y \in \mathbb{R}.\; \forall n \in \mathbb{N}.\; ((o(n) \leq y) \wedge (y \leq i(n))) \end{array}}}
\vnbprop{rr-interval-avoids-sequence}{\vnbstmt{\begin{array}{@{}l@{}} \forall e \in (\mathbb{N} \to \mathbb{R}), p, q \in \mathbb{R}. \\ \quad (p < q) \Rightarrow  \\ \quad \quad \exists y \in \mathbb{R}.\; ((p < y) \wedge ((y < q) \wedge \forall n \in \mathbb{N}.\; \neg (e(n) = y))) \end{array}}}
\vnbprop{countable-image-avoids-interval}{\vnbstmt{\begin{array}{@{}l@{}} \forall c, m, h.; \forall p, q \in \mathbb{R}. \\ \quad (\operatorname{is\text{-}countable}(c) \wedge ((c \subseteq m) \wedge ((h \in (m \to \mathbb{R})) \wedge (p < q)))) \Rightarrow  \\ \quad \quad \exists y \in \mathbb{R}.\; ((p < y) \wedge ((y < q) \wedge \forall x \in c.\; \neg (h(x) = y))) \end{array}}}
\vnbprop{neg-deriv-off-countable-nonincreasing}{\vnbstmt{\begin{array}{@{}l@{}} \forall h \in (\mathbb{R} \to \mathbb{R}), u, v \in \mathbb{R}, c. \\ \quad u < v \wedge \\ \quad \forall x \in \operatorname{ccint}(u, v).\; \operatorname{is\text{-}continuous\text{-}at}(\operatorname{rr\text{-}ms}, \operatorname{rr\text{-}ms}, h, x) \wedge \\ \quad \operatorname{is\text{-}countable}(c) \wedge \\ \quad c \subseteq \mathbb{R} \wedge \\ \quad \forall t \in \operatorname{ooint}(u, v). \\ \quad \quad \neg (t \in c) \Rightarrow \exists l.\; (\operatorname{is\text{-}diff\text{-}at}(h, t, l) \wedge (l < 0)) \Rightarrow \\ \quad \quad h(v) \leq h(u) \end{array}}}
\vnbprop{mvi-off-countable}{\vnbstmt{\begin{array}{@{}l@{}} \forall h \in (\mathbb{R} \to \mathbb{R}), m, u, v \in \mathbb{R}, c. \\ \quad u < v \wedge \\ \quad \forall x \in \operatorname{ccint}(u, v).\; \operatorname{is\text{-}continuous\text{-}at}(\operatorname{rr\text{-}ms}, \operatorname{rr\text{-}ms}, h, x) \wedge \\ \quad \operatorname{is\text{-}countable}(c) \wedge \\ \quad c \subseteq \mathbb{R} \wedge \\ \quad \forall t \in \operatorname{ooint}(u, v). \\ \quad \quad \neg (t \in c) \Rightarrow \exists l.\; (\operatorname{is\text{-}diff\text{-}at}(h, t, l) \wedge (l \leq m)) \Rightarrow \\ \quad \quad (h(v) - h(u)) \leq (m \cdot (v - u)) \end{array}}}
\vnbprop{mvi-abs-off-countable}{\vnbstmt{\begin{array}{@{}l@{}} \forall h \in (\mathbb{R} \to \mathbb{R}), m, u, v \in \mathbb{R}, c. \\ \quad u < v \wedge \\ \quad \forall x \in \operatorname{ccint}(u, v).\; \operatorname{is\text{-}continuous\text{-}at}(\operatorname{rr\text{-}ms}, \operatorname{rr\text{-}ms}, h, x) \wedge \\ \quad \operatorname{is\text{-}countable}(c) \wedge \\ \quad c \subseteq \mathbb{R} \wedge \\ \quad \forall t \in \operatorname{ooint}(u, v). \\ \quad \quad \neg (t \in c) \Rightarrow \exists l.\; (\operatorname{is\text{-}diff\text{-}at}(h, t, l) \wedge (\lvert l \rvert \leq m)) \Rightarrow \\ \quad \quad \lvert (h(v) - h(u)) \rvert \leq (m \cdot (v - u)) \end{array}}}
\vnbprop{zero-deriv-off-countable}{\vnbstmt{\begin{array}{@{}l@{}} \forall s.; \forall h \in (\mathbb{R} \to \mathbb{R}), u, v \in \mathbb{R}. \\ \quad \operatorname{is\text{-}countable}(s) \wedge \\ \quad s \subseteq \mathbb{R} \wedge \\ \quad \forall x \in \mathbb{R}.\; \operatorname{is\text{-}continuous\text{-}at}(\operatorname{rr\text{-}ms}, \operatorname{rr\text{-}ms}, h, x) \wedge \\ \quad u < v \wedge \\ \quad \forall t \in \operatorname{ooint}(u, v).\; (\neg (t \in s) \Rightarrow \operatorname{is\text{-}diff\text{-}at}(h, t, 0)) \Rightarrow \\ \quad \quad h(u) = h(v) \end{array}}}
\vnbprop{mvi-abs-off-countable-on-interval}{\vnbstmt{\begin{array}{@{}l@{}} \forall a, b, m \in \mathbb{R}, f \in (\operatorname{ccint}(a, b) \to \mathbb{R}), d. \\ \quad a < b \wedge \\ \quad \operatorname{is\text{-}continuous\text{-}on}(f, \operatorname{ccint}(a, b)) \wedge \\ \quad \operatorname{is\text{-}countable}(d) \wedge \\ \quad d \subseteq \operatorname{ccint}(a, b) \wedge \\ \quad \forall t \in \operatorname{ooint}(a, b). \\ \quad \quad \neg (t \in d) \Rightarrow \exists l.\; (\operatorname{has\text{-}deriv\text{-}at}(f, t, l) \wedge (\lvert l \rvert \leq m)) \Rightarrow \\ \quad \quad \lvert (f(b) - f(a)) \rvert \leq (m \cdot (b - a)) \end{array}}}
\vnbprop{zero-deriv-off-countable-constant}{\vnbstmt{\begin{array}{@{}l@{}} \forall a, b \in \mathbb{R}, g \in (\operatorname{ccint}(a, b) \to \mathbb{R}), d.; \forall s, w \in \operatorname{ccint}(a, b). \\ \quad a < b \wedge \\ \quad \operatorname{is\text{-}continuous\text{-}on}(g, \operatorname{ccint}(a, b)) \wedge \\ \quad \operatorname{is\text{-}countable}(d) \wedge \\ \quad d \subseteq \operatorname{ccint}(a, b) \wedge \\ \quad \forall t \in \operatorname{ooint}(a, b).\; (\neg (t \in d) \Rightarrow \operatorname{has\text{-}deriv\text{-}at}(g, t, 0)) \Rightarrow \\ \quad \quad g(s) = g(w) \end{array}}}

\vnbbookfile{theorem-library/pw-antiderivative-laws}{theorem-library/pw-antiderivative-laws}
\noindent{\small pw-antiderivative-laws.scm -- the laws of IS-PW-ANTIDERIVATIVE and PW-INT (structure-library/path-integral.scm), with the exceptional set a FINITE SET.}\par
\vnbprop{pw-antiderivative-value-in-rr}{\vnbstmt{\begin{array}{@{}l@{}} \forall f, i, a, b.; \forall y \in \operatorname{ccint}(a, b). \\ \quad \operatorname{is\text{-}primitive}(f, i, a, b) \Rightarrow (f(y) \in \mathbb{R}) \end{array}}}
\vnbprop{pw-not-in-insert}{\vnbstmt{\begin{array}{@{}l@{}} \forall y \in \mathbf{Set}, w, t. \\ \quad (\neg (t \in w) \wedge \neg (t = y)) \Rightarrow \neg (t \in (w \cup \{y\})) \end{array}}}
\vnbprop{ooint-inner-radius}{\vnbstmt{\begin{array}{@{}l@{}} \forall u, v \in \mathbb{R}, t \in \operatorname{ooint}(u, v). \exists r. \\ \quad \operatorname{pos\text{-}rr}(r) \wedge (\operatorname{ooint}((t - r), (t + r)) \subseteq \operatorname{ccint}(u, v)) \end{array}}}
\vnbprop{pw-zero-deriv-total}{\vnbstmt{\begin{array}{@{}l@{}} \forall h \in (\mathbb{R} \to \mathbb{R}), u, v \in \mathbb{R}. \\ \quad \forall x \in \mathbb{R}.\; \operatorname{is\text{-}continuous\text{-}at}(\operatorname{rr\text{-}ms}, \operatorname{rr\text{-}ms}, h, x) \wedge \\ \quad u < v \wedge \\ \quad \forall t \in \operatorname{ooint}(u, v).\; \operatorname{is\text{-}diff\text{-}at}(h, t, 0) \Rightarrow \\ \quad \quad h(u) = h(v) \end{array}}}
\vnbprop{deriv-zero-constant-on-interval}{\vnbstmt{\begin{array}{@{}l@{}} \forall a, b \in \mathbb{R}, h \in (\operatorname{ccint}(a, b) \to \mathbb{R}), u, v \in \operatorname{ccint}(a, b). \\ \quad ((a < b) \wedge (\operatorname{is\text{-}continuous\text{-}on}(h, \operatorname{ccint}(a, b)) \wedge \forall t \in \operatorname{ooint}(a, b).\; \operatorname{has\text{-}deriv\text{-}at}(h, t, 0))) \Rightarrow  \\ \quad \quad h(u) = h(v) \end{array}}}
\vnbprop{pw-zero-deriv-off-finite-set}{\vnbstmt{\begin{array}{@{}l@{}} \forall s \in \mathbf{Set}, h \in (\mathbb{R} \to \mathbb{R}), u, v \in \mathbb{R}. \\ \quad |s| \in \mathbb{N} \wedge \\ \quad \forall x \in \mathbb{R}.\; \operatorname{is\text{-}continuous\text{-}at}(\operatorname{rr\text{-}ms}, \operatorname{rr\text{-}ms}, h, x) \wedge \\ \quad u < v \wedge \\ \quad \forall t \in \operatorname{ooint}(u, v).\; (\neg (t \in s) \Rightarrow \operatorname{is\text{-}diff\text{-}at}(h, t, 0)) \Rightarrow \\ \quad \quad h(u) = h(v) \end{array}}}
\vnbprop{pw-antiderivative-endpoint-difference}{\vnbstmt{\begin{array}{@{}l@{}} \forall f, w, i, a, b. \\ \quad (\operatorname{is\text{-}primitive}(f, i, a, b) \wedge \operatorname{is\text{-}primitive}(w, i, a, b)) \Rightarrow  \\ \quad \quad (f(b) - f(a)) = (w(b) - w(a)) \end{array}}}
\vnbprop{pw-int-value}{\vnbstmt{\begin{array}{@{}l@{}} \forall f, i, a, b. \\ \quad \operatorname{is\text{-}primitive}(f, i, a, b) \Rightarrow (\int_{a}^{b} i = (f(b) - f(a))) \end{array}}}
\vnbprop{pw-int-in-rr}{\vnbstmt{\forall f, i, a, b.\; (\operatorname{is\text{-}primitive}(f, i, a, b) \Rightarrow (\int_{a}^{b} i \in \mathbb{R}))}}
\vnbprop{pw-antiderivative-integrand-congruence}{\vnbstmt{\begin{array}{@{}l@{}} \forall f, i, \mathit{ppsi}, a, b. \\ \quad \operatorname{is\text{-}primitive}(f, i, a, b) \wedge \\ \quad \mathit{ppsi} \in (\operatorname{ccint}(a, b) \to \mathbb{R}) \wedge \\ \quad \forall y \in \operatorname{ccint}(a, b).\; (i(y) \simeq \operatorname{ppsi}(y)) \Rightarrow \\ \quad \quad \operatorname{is\text{-}primitive}(f, \mathit{ppsi}, a, b) \end{array}}}
\vnbprop{pw-int-congruence}{\vnbstmt{\begin{array}{@{}l@{}} \forall f, i, \mathit{ppsi}, a, b. \\ \quad \operatorname{is\text{-}primitive}(f, i, a, b) \wedge \\ \quad \mathit{ppsi} \in (\operatorname{ccint}(a, b) \to \mathbb{R}) \wedge \\ \quad \forall y \in \operatorname{ccint}(a, b).\; (i(y) \simeq \operatorname{ppsi}(y)) \Rightarrow \\ \quad \quad \int_{a}^{b} i = \int_{a}^{b} \mathit{ppsi} \end{array}}}
\vnbprop{pw-antiderivative-one-piece}{\vnbstmt{\begin{array}{@{}l@{}} \forall a, b \in \mathbb{R}, f, i. \\ \quad a < b \wedge \\ \quad f \in (\operatorname{ccint}(a, b) \to \mathbb{R}) \wedge \\ \quad i \in (\operatorname{ccint}(a, b) \to \mathbb{R}) \wedge \\ \quad \operatorname{is\text{-}continuous\text{-}on}(f, \operatorname{ccint}(a, b)) \wedge \\ \quad \forall t \in \operatorname{ooint}(a, b).\; \operatorname{has\text{-}deriv\text{-}at}(f, t, i(t)) \Rightarrow \\ \quad \quad \operatorname{is\text{-}primitive}(f, i, a, b) \end{array}}}
\vnbprop{pw-restrict-continuous-on}{\vnbstmt{\begin{array}{@{}l@{}} \forall a, b \in \mathbb{R}, f, \mathit{pwf}. \\ \quad a < b \wedge \\ \quad f \in (\mathbb{R} \to \mathbb{R}) \wedge \\ \quad \forall y \in \mathbb{R}.\; \operatorname{is\text{-}continuous\text{-}at}(\operatorname{rr\text{-}ms}, \operatorname{rr\text{-}ms}, f, y) \wedge \\ \quad \mathit{pwf} \in (\operatorname{ccint}(a, b) \to \mathbb{R}) \wedge \\ \quad \forall y \in \operatorname{ccint}(a, b).\; (\operatorname{pwf}(y) \simeq f(y)) \Rightarrow \\ \quad \quad \operatorname{is\text{-}continuous\text{-}on}(\mathit{pwf}, \operatorname{ccint}(a, b)) \end{array}}}
\vnbprop{pw-affine-antiderivative}{\vnbstmt{\begin{array}{@{}l@{}} \forall a, b, c, m \in \mathbb{R}, f, i. \\ \quad a < b \wedge \\ \quad f \in (\operatorname{ccint}(a, b) \to \mathbb{R}) \wedge \\ \quad i \in (\operatorname{ccint}(a, b) \to \mathbb{R}) \wedge \\ \quad \forall y \in \operatorname{ccint}(a, b).\; (f(y) \simeq (c + (m \cdot y))) \wedge \\ \quad \forall y \in \operatorname{ccint}(a, b).\; (i(y) \simeq m) \Rightarrow \\ \quad \quad \operatorname{is\text{-}primitive}(f, i, a, b) \end{array}}}
\vnbprop{pw-int-const}{\vnbstmt{\begin{array}{@{}l@{}} \forall a, b, c \in \mathbb{R}, i \in (\operatorname{ccint}(a, b) \to \mathbb{R}). \\ \quad ((a < b) \wedge \forall y \in \operatorname{ccint}(a, b).\; (i(y) \simeq c)) \Rightarrow  \\ \quad \quad \int_{a}^{b} i = (c \cdot (b - a)) \end{array}}}
\vnbprop{pw-antiderivative-sum}{\vnbstmt{\begin{array}{@{}l@{}} \forall a, b \in \mathbb{R}, f, w, i, \mathit{ppsi}.; \forall h, \mathit{pchi} \in (\operatorname{ccint}(a, b) \to \mathbb{R}). \\ \quad \operatorname{is\text{-}primitive}(f, i, a, b) \wedge \\ \quad \operatorname{is\text{-}primitive}(w, \mathit{ppsi}, a, b) \wedge \\ \quad \forall y \in \operatorname{ccint}(a, b).\; (h(y) \simeq (f(y) + w(y))) \wedge \\ \quad \forall y \in \operatorname{ccint}(a, b).\; (\operatorname{pchi}(y) \simeq (i(y) + \operatorname{ppsi}(y))) \Rightarrow \\ \quad \quad (\operatorname{is\text{-}primitive}(h, \mathit{pchi}, a, b) \wedge (\int_{a}^{b} \mathit{pchi} = (\int_{a}^{b} i + \int_{a}^{b} \mathit{ppsi}))) \end{array}}}

\vnbbookfile{theorem-library/cc-int-laws}{theorem-library/cc-int-laws}
\noindent{\small cc-int-laws.scm -- the read-offs of CC-INT, TRACE and LINE-INT (structure-library/path-integral.scm), sections 3.3 and 3.4 of docs/paths-and-line-integrals-2026-09-21.md.}\par
\vnbprop{cc-int-unfold}{\vnbstmt{\begin{array}{@{}l@{}} \forall i, a, b. \\ \quad (\int_{a}^{b} i \simeq (\int_{a}^{b} (\lambda t \in \operatorname{ccint}(a, b).\; \operatorname{real\text{-}part}(i(t))) + (\int_{a}^{b} (\lambda t \in \operatorname{ccint}(a, b).\; \operatorname{imag\text{-}part}(i(t))) \cdot +i))) \end{array}}}
\vnbprop{line-int-unfold}{\vnbstmt{\begin{array}{@{}l@{}} \forall f, m, c, a, b. \\ \quad (\operatorname{line\text{-}int}(f, m, c, a, b) \simeq \int_{a}^{b} (\lambda t \in \operatorname{ccint}(a, b).\; (f(m(t)) \cdot c(t)))) \end{array}}}
\vnbprop{trace-unfold}{\vnbstmt{\forall m, a, b.\; (\operatorname{trace}(m, a, b) \simeq \operatorname{image}(m, \operatorname{ccint}(a, b)))}}
\vnbprop{trace-value-in}{\vnbstmt{\begin{array}{@{}l@{}} \forall m, a, b.; \forall t \in \operatorname{ccint}(a, b). \\ \quad (m \in (\operatorname{ccint}(a, b) \to \mathbb{C})) \Rightarrow (m(t) \in \operatorname{trace}(m, a, b)) \end{array}}}
\vnbprop{trace-subset-cc}{\vnbstmt{\forall m, a, b.\; ((m \in (\operatorname{ccint}(a, b) \to \mathbb{C})) \Rightarrow (\operatorname{trace}(m, a, b) \subseteq \mathbb{C}))}}
\vnbprop{trace-is-set}{\vnbstmt{\forall m, a, b.\; (((a \in \mathbb{R}) \wedge (b \in \mathbb{R})) \Rightarrow (\operatorname{trace}(m, a, b) \in \mathbf{Set}))}}
\vnbprop{line-int-integrand-in-fun}{\vnbstmt{\begin{array}{@{}l@{}} \forall a, b, m, c, f, d. \\ \quad a \in \mathbb{R} \wedge \\ \quad b \in \mathbb{R} \wedge \\ \quad m \in (\operatorname{ccint}(a, b) \to \mathbb{C}) \wedge \\ \quad c \in (\operatorname{ccint}(a, b) \to \mathbb{C}) \wedge \\ \quad f \in (d \to \mathbb{C}) \wedge \\ \quad \operatorname{trace}(m, a, b) \subseteq d \Rightarrow \\ \quad \quad (\lambda t \in \operatorname{ccint}(a, b).\; (f(m(t)) \cdot c(t))) \in (\operatorname{ccint}(a, b) \to \mathbb{C}) \end{array}}}

\vnbbookfile{theorem-library/road-laws}{theorem-library/road-laws}
\noindent{\small road-laws.scm -- the laws of IS-PW-CONTINUOUS-ON, IS-PATH and IS-ROAD (structure-library/path-integral.scm), section 3.4 of docs/paths-and-line-integrals-2026-09-21.md and Dieudonne 9.6.}\par
\vnbprop{pw-continuous-endpoints}{\vnbstmt{\begin{array}{@{}l@{}} \forall i, a, b. \\ \quad \operatorname{is\text{-}pw\text{-}continuous\text{-}on}(i, a, b) \Rightarrow ((a \in \mathbb{R}) \wedge ((b \in \mathbb{R}) \wedge (a < b))) \end{array}}}
\vnbprop{pw-continuous-in-fun}{\vnbstmt{\begin{array}{@{}l@{}} \forall i, a, b. \\ \quad \operatorname{is\text{-}pw\text{-}continuous\text{-}on}(i, a, b) \Rightarrow (i \in (\operatorname{ccint}(a, b) \to \mathbb{R})) \end{array}}}
\vnbprop{is-path-endpoints}{\vnbstmt{\forall m, a, b.\; (\operatorname{is\text{-}path}(m, a, b) \Rightarrow ((a \in \mathbb{R}) \wedge ((b \in \mathbb{R}) \wedge (a < b))))}}
\vnbprop{is-path-in-fun}{\vnbstmt{\forall m, a, b.\; (\operatorname{is\text{-}path}(m, a, b) \Rightarrow (m \in (\operatorname{ccint}(a, b) \to \mathbb{C})))}}
\vnbprop{is-path-continuous}{\vnbstmt{\begin{array}{@{}l@{}} \forall m, a, b.; \forall x \in \operatorname{ccint}(a, b). \\ \quad \operatorname{is\text{-}path}(m, a, b) \Rightarrow  \\ \quad \quad \operatorname{is\text{-}continuous\text{-}at}(\operatorname{subspace\text{-}ms}(\operatorname{rr\text{-}ms}, \operatorname{ccint}(a, b)), \operatorname{cc\text{-}ms}, m, x) \end{array}}}
\vnbprop{is-road-is-path}{\vnbstmt{\forall m, c, a, b.\; (\operatorname{is\text{-}road}(m, c, a, b) \Rightarrow \operatorname{is\text{-}path}(m, a, b))}}
\vnbprop{is-road-dgam-in-fun}{\vnbstmt{\forall m, c, a, b.\; (\operatorname{is\text{-}road}(m, c, a, b) \Rightarrow (c \in (\operatorname{ccint}(a, b) \to \mathbb{C})))}}
\vnbprop{is-road-re-antiderivative}{\vnbstmt{\begin{array}{@{}l@{}} \forall m, c, a, b. \\ \quad \operatorname{is\text{-}road}(m, c, a, b) \Rightarrow  \\ \quad \quad \operatorname{is\text{-}primitive}((\lambda t \in \operatorname{ccint}(a, b).\; \operatorname{real\text{-}part}(m(t))), \\ \quad \quad \quad (\lambda t \in \operatorname{ccint}(a, b).\; \operatorname{real\text{-}part}(c(t))), a, b) \end{array}}}
\vnbprop{is-road-im-antiderivative}{\vnbstmt{\begin{array}{@{}l@{}} \forall m, c, a, b. \\ \quad \operatorname{is\text{-}road}(m, c, a, b) \Rightarrow  \\ \quad \quad \operatorname{is\text{-}primitive}((\lambda t \in \operatorname{ccint}(a, b).\; \operatorname{imag\text{-}part}(m(t))), \\ \quad \quad \quad (\lambda t \in \operatorname{ccint}(a, b).\; \operatorname{imag\text{-}part}(c(t))), a, b) \end{array}}}
\vnbprop{is-road-re-regulated}{\vnbstmt{\begin{array}{@{}l@{}} \forall m, c, a, b. \\ \quad \operatorname{is\text{-}road}(m, c, a, b) \Rightarrow  \\ \quad \quad \operatorname{is\text{-}regulated\text{-}on}((\lambda t \in \operatorname{ccint}(a, b).\; \operatorname{real\text{-}part}(c(t))), a, b) \end{array}}}
\vnbprop{is-road-im-regulated}{\vnbstmt{\begin{array}{@{}l@{}} \forall m, c, a, b. \\ \quad \operatorname{is\text{-}road}(m, c, a, b) \Rightarrow  \\ \quad \quad \operatorname{is\text{-}regulated\text{-}on}((\lambda t \in \operatorname{ccint}(a, b).\; \operatorname{imag\text{-}part}(c(t))), a, b) \end{array}}}
\vnbprop{pw-continuous-one-piece}{\vnbstmt{\begin{array}{@{}l@{}} \forall a, b \in \mathbb{R}, i \in (\operatorname{ccint}(a, b) \to \mathbb{R}). \\ \quad ((a < b) \wedge \operatorname{is\text{-}continuous\text{-}on}(i, \operatorname{ccint}(a, b))) \Rightarrow \operatorname{is\text{-}pw\text{-}continuous\text{-}on}(i, a, b) \end{array}}}
\vnbprop{is-path-of-coords}{\vnbstmt{\begin{array}{@{}l@{}} \forall a, b \in \mathbb{R}, m, r, i. \\ \quad a < b \wedge \\ \quad m \in (\operatorname{ccint}(a, b) \to \mathbb{C}) \wedge \\ \quad r \in (\operatorname{ccint}(a, b) \to \mathbb{R}) \wedge \\ \quad i \in (\operatorname{ccint}(a, b) \to \mathbb{R}) \wedge \\ \quad \forall y \in \operatorname{ccint}(a, b).\; (\operatorname{real\text{-}part}(m(y)) = r(y)) \wedge \\ \quad \forall y \in \operatorname{ccint}(a, b).\; (\operatorname{imag\text{-}part}(m(y)) = i(y)) \wedge \\ \quad \operatorname{is\text{-}continuous\text{-}on}(r, \operatorname{ccint}(a, b)) \wedge \\ \quad \operatorname{is\text{-}continuous\text{-}on}(i, \operatorname{ccint}(a, b)) \Rightarrow \\ \quad \quad \operatorname{is\text{-}path}(m, a, b) \end{array}}}
\vnbprop{segment-is-road}{\vnbstmt{\begin{array}{@{}l@{}} \forall z_{0}, w. \\ \quad ((z_{0} \in \mathbb{C}) \wedge (w \in \mathbb{C})) \Rightarrow  \\ \quad \quad \operatorname{is\text{-}road}((\lambda x \in \operatorname{ccint}(0, 1).\; (z_{0} + (x \cdot w))), (\lambda x \in \operatorname{ccint}(0, 1).\; w), 0, 1) \end{array}}}
\vnbprop{line-int-of-one}{\vnbstmt{\begin{array}{@{}l@{}} \forall m, c, a, b, f, d. \\ \quad (\operatorname{is\text{-}road}(m, c, a, b) \wedge ((f \in (d \to \mathbb{C})) \wedge ((\operatorname{trace}(m, a, b) \subseteq d) \wedge \forall y \in d.\; (f(y) = 1)))) \Rightarrow  \\ \quad \quad \operatorname{line\text{-}int}(f, m, c, a, b) = (m(b) - m(a)) \end{array}}}

\vnbbookfile{theorem-library/pw-int-laws-2}{theorem-library/pw-int-laws-2}
\noindent{\small pw-int-laws-2.scm -- the LINEARITY and the ADDITIVITY of the integral on [a, b], continuing theorem-library/pw-antiderivative-laws.scm.}\par
\vnbprop{pw-antiderivative-real-mul}{\vnbstmt{\begin{array}{@{}l@{}} \forall a, b \in \mathbb{R}, f, i.; \forall c \in \mathbb{R}, h, \mathit{pchi} \in (\operatorname{ccint}(a, b) \to \mathbb{R}). \\ \quad \operatorname{is\text{-}primitive}(f, i, a, b) \wedge \\ \quad \forall y \in \operatorname{ccint}(a, b).\; (h(y) \simeq (c \cdot f(y))) \wedge \\ \quad \forall y \in \operatorname{ccint}(a, b).\; (\operatorname{pchi}(y) \simeq (c \cdot i(y))) \Rightarrow \\ \quad \quad \operatorname{is\text{-}primitive}(h, \mathit{pchi}, a, b) \wedge (\int_{a}^{b} \mathit{pchi} = (c \cdot \int_{a}^{b} i)) \end{array}}}
\vnbprop{pw-antiderivative-differ-by-constant}{\vnbstmt{\begin{array}{@{}l@{}} \forall f, w, i, a, b.; \forall s \in \operatorname{ccint}(a, b). \\ \quad (\operatorname{is\text{-}primitive}(f, i, a, b) \wedge \operatorname{is\text{-}primitive}(w, i, a, b)) \Rightarrow  \\ \quad \quad f(s) = (w(s) + (f(a) - w(a))) \end{array}}}
\vnbprop{pw-antiderivative-differ-by-some-constant}{\vnbstmt{\begin{array}{@{}l@{}} \forall f, w, i, a, b. \\ \quad (\operatorname{is\text{-}primitive}(f, i, a, b) \wedge \operatorname{is\text{-}primitive}(w, i, a, b)) \Rightarrow  \\ \quad \quad \exists c \in \mathbb{R}.\; \forall s \in \operatorname{ccint}(a, b).\; (f(s) = (w(s) + c)) \end{array}}}
\vnbprop{ccint-subset-ccint}{\vnbstmt{\begin{array}{@{}l@{}} \forall a, b, v, \mathit{pdv} \in \mathbb{R}. \\ \quad ((a \leq v) \wedge (\mathit{pdv} \leq b)) \Rightarrow (\operatorname{ccint}(v, \mathit{pdv}) \subseteq \operatorname{ccint}(a, b)) \end{array}}}
\vnbprop{pw-antiderivative-restrict}{\vnbstmt{\begin{array}{@{}l@{}} \forall a, b, v, \mathit{pdv} \in \mathbb{R}, f, i. \\ \quad (\operatorname{is\text{-}primitive}(f, i, a, b) \wedge ((a \leq v) \wedge ((v < \mathit{pdv}) \wedge (\mathit{pdv} \leq b)))) \Rightarrow  \\ \quad \quad \operatorname{is\text{-}primitive}(\operatorname{restrict}(f, \operatorname{ccint}(v, \mathit{pdv})), \operatorname{restrict}(i, \operatorname{ccint}(v, \mathit{pdv})), v, \mathit{pdv}) \end{array}}}
\vnbprop{pw-int-adjacent}{\vnbstmt{\begin{array}{@{}l@{}} \forall a, b \in \mathbb{R}, v \in \operatorname{ooint}(a, b), f, i. \\ \quad \operatorname{is\text{-}primitive}(f, i, a, b) \Rightarrow  \\ \quad \quad \int_{a}^{b} i = (\int_{a}^{v} \operatorname{restrict}(i, \operatorname{ccint}(a, v)) + \int_{v}^{b} \operatorname{restrict}(i, \operatorname{ccint}(v, b))) \end{array}}}
\vnbprop{cc-int-value}{\vnbstmt{\begin{array}{@{}l@{}} \forall a, b, i, f, w. \\ \quad a \in \mathbb{R} \wedge \\ \quad b \in \mathbb{R} \wedge \\ \quad \operatorname{is\text{-}primitive}(f, (\lambda t \in \operatorname{ccint}(a, b).\; \operatorname{real\text{-}part}(i(t))), a, b) \wedge \\ \quad \operatorname{is\text{-}primitive}(w, (\lambda t \in \operatorname{ccint}(a, b).\; \operatorname{imag\text{-}part}(i(t))), a, b) \Rightarrow \\ \quad \quad \int_{a}^{b} i = ((f(b) - f(a)) + ((w(b) - w(a)) \cdot +i)) \end{array}}}
\vnbprop{cc-int-in-cc}{\vnbstmt{\begin{array}{@{}l@{}} \forall a, b, i, f, w. \\ \quad a \in \mathbb{R} \wedge \\ \quad b \in \mathbb{R} \wedge \\ \quad \operatorname{is\text{-}primitive}(f, (\lambda t \in \operatorname{ccint}(a, b).\; \operatorname{real\text{-}part}(i(t))), a, b) \wedge \\ \quad \operatorname{is\text{-}primitive}(w, (\lambda t \in \operatorname{ccint}(a, b).\; \operatorname{imag\text{-}part}(i(t))), a, b) \Rightarrow \\ \quad \quad \int_{a}^{b} i \in \mathbb{C} \end{array}}}
\vnbprop{cc-int-sum}{\vnbstmt{\begin{array}{@{}l@{}} \forall a, b, i, \mathit{ppsi}, \mathit{pchi}, f, w, h, \mathit{pbw}. \\ \quad a \in \mathbb{R} \wedge \\ \quad b \in \mathbb{R} \wedge \\ \quad i \in (\operatorname{ccint}(a, b) \to \mathbb{C}) \wedge \\ \quad \mathit{ppsi} \in (\operatorname{ccint}(a, b) \to \mathbb{C}) \wedge \\ \quad \mathit{pchi} \in (\operatorname{ccint}(a, b) \to \mathbb{C}) \wedge \\ \quad \forall y \in \operatorname{ccint}(a, b).\; (\operatorname{pchi}(y) \simeq (i(y) + \operatorname{ppsi}(y))) \wedge \\ \quad \operatorname{is\text{-}primitive}(f, (\lambda t \in \operatorname{ccint}(a, b).\; \operatorname{real\text{-}part}(i(t))), a, b) \wedge \\ \quad \operatorname{is\text{-}primitive}(w, (\lambda t \in \operatorname{ccint}(a, b).\; \operatorname{imag\text{-}part}(i(t))), a, b) \wedge \\ \quad \operatorname{is\text{-}primitive}(h, (\lambda t \in \operatorname{ccint}(a, b).\; \operatorname{real\text{-}part}(\operatorname{ppsi}(t))), a, b) \wedge \\ \quad \operatorname{is\text{-}primitive}(\mathit{pbw}, (\lambda t \in \operatorname{ccint}(a, b).\; \operatorname{imag\text{-}part}(\operatorname{ppsi}(t))), a, b) \Rightarrow \\ \quad \quad \int_{a}^{b} \mathit{pchi} = (\int_{a}^{b} i + \int_{a}^{b} \mathit{ppsi}) \end{array}}}
\vnbprop{cc-int-real-mul}{\vnbstmt{\begin{array}{@{}l@{}} \forall a, b, i, \mathit{pchi}.; \forall c \in \mathbb{R}, f, w. \\ \quad a \in \mathbb{R} \wedge \\ \quad b \in \mathbb{R} \wedge \\ \quad i \in (\operatorname{ccint}(a, b) \to \mathbb{C}) \wedge \\ \quad \mathit{pchi} \in (\operatorname{ccint}(a, b) \to \mathbb{C}) \wedge \\ \quad \forall y \in \operatorname{ccint}(a, b).\; (\operatorname{pchi}(y) \simeq (c \cdot i(y))) \wedge \\ \quad \operatorname{is\text{-}primitive}(f, (\lambda t \in \operatorname{ccint}(a, b).\; \operatorname{real\text{-}part}(i(t))), a, b) \wedge \\ \quad \operatorname{is\text{-}primitive}(w, (\lambda t \in \operatorname{ccint}(a, b).\; \operatorname{imag\text{-}part}(i(t))), a, b) \Rightarrow \\ \quad \quad \int_{a}^{b} \mathit{pchi} = (c \cdot \int_{a}^{b} i) \end{array}}}

\vnbbookfile{theorem-library/taylor-cluster}{theorem-library/taylor-cluster}
\noindent{\small taylor-cluster.scm -- THE TAYLOR CLUSTER OF THE CALCULUS NOTES (calculus.pdf, ch. 2, sections 1, 6 and 7): 2.3, 2.16, 2.20, 2.21, 2.22.  Batch 25-B, 2026-09-24.}\par
\vnbprop{nth-deriv-succ-inner}{\vnbstmt{\begin{array}{@{}l@{}} \forall k \in \mathbb{N}, f. \\ \quad (\operatorname{nth\text{-}deriv}(f, (k + 1)) \simeq \operatorname{nth\text{-}deriv}((\lambda x \in \mathbb{R}.\; \operatorname{deriv}(f, x)), k)) \end{array}}}
\vnbprop{nth-deriv-monomial-core}{\vnbstmt{\begin{array}{@{}l@{}} \forall k, m \in \mathbb{N}. \\ \quad (\operatorname{nth\text{-}deriv}((\lambda x \in \mathbb{R}.\; \operatorname{power}(x, (m + k))), k) \simeq (\lambda x \in \mathbb{R}.\; ((\operatorname{factorial}((m + k)) \cdot \operatorname{recip}(\operatorname{factorial}(m))) \cdot \operatorname{power}(x, m)))) \end{array}}}
\vnbprop{nth-deriv-monomial-gap}{\vnbstmt{\begin{array}{@{}l@{}} \forall k, m \in \mathbb{N}, p. \\ \quad (p = (m + k)) \Rightarrow  \\ \quad \quad (\operatorname{nth\text{-}deriv}((\lambda x \in \mathbb{R}.\; \operatorname{power}(x, p)), k) \simeq (\lambda x \in \mathbb{R}.\; ((\operatorname{factorial}(p) \cdot \operatorname{recip}(\operatorname{factorial}(m))) \cdot \operatorname{power}(x, m)))) \end{array}}}
\vnbprop{nth-deriv-monomial-beyond}{\vnbstmt{\begin{array}{@{}l@{}} \forall j, n \in \mathbb{N}. \\ \quad (\operatorname{nth\text{-}deriv}((\lambda x \in \mathbb{R}.\; \operatorname{power}(x, n)), ((n + j) + 1)) \simeq (\lambda x \in \mathbb{R}.\; 0)) \end{array}}}
\vnbprop{nth-deriv-monomial-vanish}{\vnbstmt{\begin{array}{@{}l@{}} \forall n, k \in \mathbb{N}. \\ \quad (n < k) \Rightarrow (\operatorname{nth\text{-}deriv}((\lambda x \in \mathbb{R}.\; \operatorname{power}(x, n)), k) \simeq (\lambda x \in \mathbb{R}.\; 0)) \end{array}}}
\vnbprop{nth-deriv-monomial}{\vnbstmt{\begin{array}{@{}l@{}} \forall n, k \in \mathbb{N}. \\ \quad (k \leq n) \Rightarrow  \\ \quad \quad (\operatorname{nth\text{-}deriv}((\lambda x \in \mathbb{R}.\; \operatorname{power}(x, n)), k) \simeq (\lambda x \in \mathbb{R}.\; ((\operatorname{factorial}(n) \cdot \operatorname{recip}(\operatorname{factorial}(\operatorname{nn\text{-}minus}(n, k)))) \cdot \operatorname{power}(x, \operatorname{nn\text{-}minus}(n, k))))) \end{array}}}
\vnbprop{poly-deriv-fn}{\vnbstmt{\begin{array}{@{}l@{}} \forall n \in \mathbb{N}, f, d \in (\mathbb{N} \to \mathbb{R}). \\ \quad \forall j \in \mathbb{N}.\; (d(j) = ((j + 1) \cdot f((j + 1)))) \Rightarrow  \\ \quad \quad ((\lambda \mathit{tcx} \in \mathbb{R}.\; \operatorname{deriv}((\lambda x \in \mathbb{R}.\; \operatorname{series\text{-}partial\text{-}sum}((\lambda k \in \mathbb{N}.\; (f(k) \cdot \operatorname{power}(x, k))), (n + 1))), \mathit{tcx})) \simeq (\lambda x \in \mathbb{R}.\; \operatorname{series\text{-}partial\text{-}sum}((\lambda k \in \mathbb{N}.\; (d(k) \cdot \operatorname{power}(x, k))), n))) \end{array}}}
\vnbprop{poly-deriv-coef-exists}{\vnbstmt{\begin{array}{@{}l@{}} \forall f \in (\mathbb{N} \to \mathbb{R}). \exists d \in (\mathbb{N} \to \mathbb{R}). \forall j \in \mathbb{N}. \\ \quad d(j) = ((j + 1) \cdot f((j + 1))) \end{array}}}
\vnbprop{poly-deriv-fn-zero}{\vnbstmt{\begin{array}{@{}l@{}} \forall f \in (\mathbb{N} \to \mathbb{R}). \\ \quad ((\lambda \mathit{tcx} \in \mathbb{R}.\; \operatorname{deriv}((\lambda x \in \mathbb{R}.\; \operatorname{series\text{-}partial\text{-}sum}((\lambda k \in \mathbb{N}.\; (f(k) \cdot \operatorname{power}(x, k))), 0)), \mathit{tcx})) \simeq (\lambda x \in \mathbb{R}.\; \operatorname{series\text{-}partial\text{-}sum}((\lambda k \in \mathbb{N}.\; (f(k) \cdot \operatorname{power}(x, k))), 0))) \end{array}}}
\vnbprop{poly-nth-deriv-in-fun}{\vnbstmt{\begin{array}{@{}l@{}} \forall \mathit{tck}, n \in \mathbb{N}, f \in (\mathbb{N} \to \mathbb{R}). \\ \quad \operatorname{nth\text{-}deriv}((\lambda x \in \mathbb{R}.\; \operatorname{series\text{-}partial\text{-}sum}((\lambda k \in \mathbb{N}.\; (f(k) \cdot \operatorname{power}(x, k))), n)), \mathit{tck}) \\ \quad \in\; (\mathbb{R} \to \mathbb{R}) \end{array}}}
\vnbprop{poly-value-at-zero}{\vnbstmt{\begin{array}{@{}l@{}} \forall n \in \mathbb{N}, f \in (\mathbb{N} \to \mathbb{R}). \\ \quad \operatorname{series\text{-}partial\text{-}sum}((\lambda k \in \mathbb{N}.\; (f(k) \cdot \operatorname{power}(0, k))), (n + 1)) = f(0) \end{array}}}
\vnbprop{poly-nth-deriv-at-zero}{\vnbstmt{\begin{array}{@{}l@{}} \forall \mathit{tck}, n \in \mathbb{N}, f \in (\mathbb{N} \to \mathbb{R}). \\ \quad (\mathit{tck} \leq n) \Rightarrow  \\ \quad \quad \operatorname{nth\text{-}deriv}((\lambda x \in \mathbb{R}.\; \operatorname{series\text{-}partial\text{-}sum}((\lambda k \in \mathbb{N}.\; (f(k) \cdot \operatorname{power}(x, k))), (n + 1))), \mathit{tck})(0) \\ \quad \quad =\; \operatorname{factorial}(\mathit{tck}) \cdot f(\mathit{tck}) \end{array}}}
\vnbprop{taylor-coef-from-derivs}{\vnbstmt{\begin{array}{@{}l@{}} \forall n \in \mathbb{N}, f.; \forall a \in (\mathbb{N} \to \mathbb{R}), \mathit{tck} \in \mathbb{N}. \\ \quad \forall \mathit{tck} \in \mathbb{N}. \\ \quad \quad (\mathit{tck} \leq n) \Rightarrow  \\ \quad \quad \quad \operatorname{nth\text{-}deriv}((\lambda x \in \mathbb{R}.\; \operatorname{series\text{-}partial\text{-}sum}((\lambda k \in \mathbb{N}.\; (a(k) \cdot \operatorname{power}(x, k))), (n + 1))), \mathit{tck})(0) \\ \quad \quad \quad =\; \operatorname{nth\text{-}deriv}(f, \mathit{tck})(0) \wedge \\ \quad \mathit{tck} \leq n \Rightarrow \\ \quad \quad a(\mathit{tck}) = (\operatorname{nth\text{-}deriv}(f, \mathit{tck})(0) \cdot \operatorname{recip}(\operatorname{factorial}(\mathit{tck}))) \end{array}}}
\vnbprop{taylor-derivs-from-coef}{\vnbstmt{\begin{array}{@{}l@{}} \forall n \in \mathbb{N}, f.; \forall a \in (\mathbb{N} \to \mathbb{R}), \mathit{tck} \in \mathbb{N}. \\ \quad \forall \mathit{tck} \in \mathbb{N}.\; ((\mathit{tck} \leq n) \Rightarrow (\operatorname{nth\text{-}deriv}(f, \mathit{tck})(0) \in \mathbb{R})) \wedge \\ \quad \forall \mathit{tck} \in \mathbb{N}. \\ \quad \quad (\mathit{tck} \leq n) \Rightarrow (a(\mathit{tck}) = (\operatorname{nth\text{-}deriv}(f, \mathit{tck})(0) \cdot \operatorname{recip}(\operatorname{factorial}(\mathit{tck})))) \wedge \\ \quad \mathit{tck} \leq n \Rightarrow \\ \quad \quad \operatorname{nth\text{-}deriv}((\lambda x \in \mathbb{R}.\; \operatorname{series\text{-}partial\text{-}sum}((\lambda k \in \mathbb{N}.\; (a(k) \cdot \operatorname{power}(x, k))), (n + 1))), \mathit{tck})(0) \\ \quad \quad =\; \operatorname{nth\text{-}deriv}(f, \mathit{tck})(0) \end{array}}}
\vnbprop{taylor-interpolant-at-zero}{\vnbstmt{\begin{array}{@{}l@{}} \forall n \in \mathbb{N}, f. \\ \quad \forall \mathit{tck} \in \mathbb{N}.\; ((\mathit{tck} \leq n) \Rightarrow (\operatorname{nth\text{-}deriv}(f, \mathit{tck})(0) \in \mathbb{R})) \Rightarrow  \\ \quad \quad \exists a \in (\mathbb{N} \to \mathbb{R}). \\ \quad \quad \quad \forall \mathit{tck} \in \mathbb{N}. \\ \quad \quad \quad \quad (\mathit{tck} \leq n) \Rightarrow (a(\mathit{tck}) = (\operatorname{nth\text{-}deriv}(f, \mathit{tck})(0) \cdot \operatorname{recip}(\operatorname{factorial}(\mathit{tck})))) \wedge \\ \quad \quad \quad \forall \mathit{tck} \in \mathbb{N}. \\ \quad \quad \quad \quad (\mathit{tck} \leq n) \Rightarrow  \\ \quad \quad \quad \quad \quad \operatorname{nth\text{-}deriv}((\lambda x \in \mathbb{R}.\; \operatorname{series\text{-}partial\text{-}sum}((\lambda k \in \mathbb{N}.\; (a(k) \cdot \operatorname{power}(x, k))), (n + 1))), \mathit{tck})(0) \\ \quad \quad \quad \quad \quad =\; \operatorname{nth\text{-}deriv}(f, \mathit{tck})(0) \wedge \\ \quad \quad \quad \forall g \in (\mathbb{N} \to \mathbb{R}), \mathit{tck} \in \mathbb{N}. \\ \quad \quad \quad \quad \forall \mathit{tck} \in \mathbb{N}. \\ \quad \quad \quad \quad \quad (\mathit{tck} \leq n) \Rightarrow  \\ \quad \quad \quad \quad \quad \quad \operatorname{nth\text{-}deriv}((\lambda x \in \mathbb{R}.\; \operatorname{series\text{-}partial\text{-}sum}((\lambda k \in \mathbb{N}.\; (g(k) \cdot \operatorname{power}(x, k))), (n + 1))), \mathit{tck})(0) \\ \quad \quad \quad \quad \quad \quad =\; \operatorname{nth\text{-}deriv}(f, \mathit{tck})(0) \wedge \\ \quad \quad \quad \quad \mathit{tck} \leq n \Rightarrow \\ \quad \quad \quad \quad \quad g(\mathit{tck}) = a(\mathit{tck}) \end{array}}}
\vnbprop{series-partial-sum-shift}{\vnbstmt{\begin{array}{@{}l@{}} \forall n \in \mathbb{N}, u \in (\mathbb{N} \to \mathbb{R}). \\ \quad \operatorname{series\text{-}partial\text{-}sum}(u, (n + 1)) = (u(0) + \operatorname{series\text{-}partial\text{-}sum}((\lambda i \in \mathbb{N}.\; u((i + 1))), n)) \end{array}}}
\vnbprop{poly-degree-zero-value}{\vnbstmt{\begin{array}{@{}l@{}} \forall f \in (\mathbb{N} \to \mathbb{R}), v \in \mathbb{R}. \\ \quad \operatorname{series\text{-}partial\text{-}sum}((\lambda k \in \mathbb{N}.\; (f(k) \cdot \operatorname{power}(v, k))), 1) = f(0) \end{array}}}
\vnbprop{taylor-exact-polynomial}{\vnbstmt{\begin{array}{@{}l@{}} \forall m \in \mathbb{N}, f \in (\mathbb{N} \to \mathbb{R}), a, h \in \mathbb{R}. \\ \quad (\lambda x \in \mathbb{R}.\; \operatorname{series\text{-}partial\text{-}sum}((\lambda k \in \mathbb{N}.\; (f(k) \cdot \operatorname{power}(x, k))), (m + 1)))((a + h)) \\ \quad =\; \operatorname{series\text{-}partial\text{-}sum}((\lambda i \in \mathbb{N}.\; \\ \quad \quad \quad (\operatorname{nth\text{-}deriv}((\lambda x \in \mathbb{R}.\; \operatorname{series\text{-}partial\text{-}sum}((\lambda k \in \mathbb{N}.\; (f(k) \cdot \operatorname{power}(x, k))), (m + 1))), i)(a) \cdot \operatorname{power}(h, i)) \cdot \operatorname{recip}(\operatorname{factorial}(i))), \\ \quad \quad (m + 1)) \end{array}}}
\vnbprop{taylor-remainder-sup-bound}{\vnbstmt{\begin{array}{@{}l@{}} \forall a, b \in \mathbb{R}, f.; \forall m \in \mathbb{N}. \\ \quad \operatorname{is\text{-}taylor\text{-}family}(f, (m + 1), a, b) \wedge \\ \quad \operatorname{rr\text{-}bounded\text{-}above}(\operatorname{image}((\lambda v \in \operatorname{ooint}(a, b).\; \lvert f((m + 1))(v) \rvert), \operatorname{ooint}(a, b))) \Rightarrow \\ \quad \quad \lvert \operatorname{taylor\text{-}rem}(f, (m + 1), a, b) \rvert \\ \quad \quad \leq\; (\operatorname{power}((b - a), (m + 1)) \cdot \operatorname{recip}(\operatorname{factorial}((m + 1)))) \cdot \operatorname{sup}(\operatorname{image}((\lambda v \in \operatorname{ooint}(a, b).\; \lvert f((m + 1))(v) \rvert), \operatorname{ooint}(a, b))) \end{array}}}
\vnbprop{series-partial-sum-agree-below}{\vnbstmt{\begin{array}{@{}l@{}} \forall n \in \mathbb{N}, u, g. \\ \quad (\forall k \in \mathbb{N}.\; ((k < n) \Rightarrow (g(k) \in \mathbb{R})) \wedge \forall k \in \mathbb{N}.\; ((k < n) \Rightarrow (u(k) = g(k)))) \Rightarrow  \\ \quad \quad \operatorname{series\text{-}partial\text{-}sum}(u, n) = \operatorname{series\text{-}partial\text{-}sum}(g, n) \end{array}}}
\vnbprop{taylor-lagrange-subinterval}{\vnbstmt{\begin{array}{@{}l@{}} \forall a, b \in \mathbb{R}, f.; \forall m \in \mathbb{N}, t \in \mathbb{R}. \\ \quad \forall k \in \mathbb{N}.\; ((k \leq (m + 1)) \Rightarrow (f(k) \in (\operatorname{ccint}(a, b) \to \mathbb{R}))) \wedge \\ \quad \forall k \in \mathbb{N}.\; ((k < (m + 1)) \Rightarrow \operatorname{is\text{-}continuous\text{-}on}(f(k), \operatorname{ccint}(a, b))) \wedge \\ \quad \forall k \in \mathbb{N}, x \in \operatorname{ooint}(a, b). \\ \quad \quad (k < (m + 1)) \Rightarrow \operatorname{has\text{-}deriv\text{-}at}(f(k), x, f((k + 1))(x)) \wedge \\ \quad a < t \wedge \\ \quad t \leq b \Rightarrow \\ \quad \quad \exists i \in \operatorname{ooint}(a, t). \\ \quad \quad \quad \operatorname{taylor\text{-}rem}(f, (m + 1), a, t) = ((f((m + 1))(i) \cdot \operatorname{recip}(\operatorname{factorial}((m + 1)))) \cdot \operatorname{power}((t - a), (m + 1))) \end{array}}}
\vnbprop{taylor-sum-at-base}{\vnbstmt{\begin{array}{@{}l@{}} \forall j \in \mathbb{N}, f.; \forall a \in \mathbb{R}. \\ \quad \forall k \in \mathbb{N}.\; ((k \leq j) \Rightarrow (f(k)(a) \in \mathbb{R})) \Rightarrow  \\ \quad \quad \operatorname{series\text{-}partial\text{-}sum}((\lambda \mathit{tk} \in \mathbb{N}.\; ((f(\mathit{tk})(a) \cdot \operatorname{power}((a - a), \mathit{tk})) \cdot \operatorname{recip}(\operatorname{factorial}(\mathit{tk})))), (j + 1)) \\ \quad \quad =\; f(0)(a) \end{array}}}
\vnbprop{taylor-factor-value}{\vnbstmt{\begin{array}{@{}l@{}} \forall a, b \in \mathbb{R}, f.; \forall m \in \mathbb{N}, t \in \mathbb{R}. \\ \quad \forall k \in \mathbb{N}.\; ((k \leq (m + 1)) \Rightarrow (f(k) \in (\operatorname{ccint}(a, b) \to \mathbb{R}))) \wedge \\ \quad \forall k \in \mathbb{N}.\; ((k < (m + 1)) \Rightarrow \operatorname{is\text{-}continuous\text{-}on}(f(k), \operatorname{ccint}(a, b))) \wedge \\ \quad \forall k \in \mathbb{N}, x \in \operatorname{ooint}(a, b). \\ \quad \quad (k < (m + 1)) \Rightarrow \operatorname{has\text{-}deriv\text{-}at}(f(k), x, f((k + 1))(x)) \wedge \\ \quad a < t \wedge \\ \quad t \leq b \Rightarrow \\ \quad \quad \exists i \in \operatorname{ooint}(a, t). \\ \quad \quad \quad ((\operatorname{taylor\text{-}rem}(f, (m + 1), a, t) \cdot \operatorname{factorial}((m + 1))) \cdot \operatorname{recip}(\operatorname{power}((t - a), (m + 1)))) = f((m + 1))(i) \end{array}}}
\vnbprop{taylor-sum-continuous}{\vnbstmt{\begin{array}{@{}l@{}} \forall a \in \mathbb{R}, f.; \forall m \in \mathbb{N}, p \in \mathbb{R}. \\ \quad \forall k \in \mathbb{N}.\; ((k < (m + 1)) \Rightarrow (f(k)(a) \in \mathbb{R})) \Rightarrow  \\ \quad \quad \operatorname{is\text{-}continuous\text{-}at}(\operatorname{rr\text{-}ms}, \operatorname{rr\text{-}ms}, (\lambda s \in \mathbb{R}.\; \operatorname{taylor\text{-}sum}(f, (m + 1), s, a)), p) \end{array}}}
\vnbprop{taylor-factor-fn-continuous}{\vnbstmt{\begin{array}{@{}l@{}} \forall a, b \in \mathbb{R}, f.; \forall m \in \mathbb{N}, p \in \mathbb{R}. \\ \quad a \leq b \wedge \\ \quad \operatorname{is\text{-}continuous\text{-}on}(f(0), \operatorname{ccint}(a, b)) \wedge \\ \quad \forall k \in \mathbb{N}.\; ((k < (m + 1)) \Rightarrow (f(k)(a) \in \mathbb{R})) \wedge \\ \quad a < p \Rightarrow \\ \quad \quad \operatorname{is\text{-}continuous\text{-}at}(\operatorname{rr\text{-}ms}, \operatorname{rr\text{-}ms}, \\ \quad \quad \quad (\lambda s \in \mathbb{R}.\; \\ \quad \quad \quad \quad ((\operatorname{extend\text{-}const}(f(0), a, b)(s) - \operatorname{taylor\text{-}sum}(f, (m + 1), s, a)) \cdot \operatorname{factorial}((m + 1))) \cdot {\operatorname{power}((s - a), (m + 1))}^{\ast\!-1}), \\ \quad \quad \quad p) \end{array}}}
\vnbprop{taylor-factor-fn-value}{\vnbstmt{\begin{array}{@{}l@{}} \forall a, b \in \mathbb{R}, f.; \forall m \in \mathbb{N}, t \in \mathbb{R}. \\ \quad \forall k \in \mathbb{N}.\; ((k \leq (m + 1)) \Rightarrow (f(k) \in (\operatorname{ccint}(a, b) \to \mathbb{R}))) \wedge \\ \quad \forall k \in \mathbb{N}.\; ((k < (m + 1)) \Rightarrow \operatorname{is\text{-}continuous\text{-}on}(f(k), \operatorname{ccint}(a, b))) \wedge \\ \quad \forall k \in \mathbb{N}, x \in \operatorname{ooint}(a, b). \\ \quad \quad (k < (m + 1)) \Rightarrow \operatorname{has\text{-}deriv\text{-}at}(f(k), x, f((k + 1))(x)) \wedge \\ \quad a < t \wedge \\ \quad t \leq b \Rightarrow \\ \quad \quad \exists i \in \operatorname{ooint}(a, t). \\ \quad \quad \quad (\lambda s \in \mathbb{R}.\; (((\operatorname{extend\text{-}const}(f(0), a, b)(s) - \operatorname{taylor\text{-}sum}(f, (m + 1), s, a)) \cdot \operatorname{factorial}((m + 1))) \cdot {\operatorname{power}((s - a), (m + 1))}^{\ast\!-1}))(t) \\ \quad \quad \quad =\; f((m + 1))(i) \end{array}}}
\vnbprop{taylor-factor-glue-continuous}{\vnbstmt{\begin{array}{@{}l@{}} \forall a, b \in \mathbb{R}, f.; \forall m \in \mathbb{N}, p \in \mathbb{R}. \\ \quad a < b \wedge \\ \quad \forall k \in \mathbb{N}.\; ((k \leq (m + 1)) \Rightarrow (f(k) \in (\operatorname{ccint}(a, b) \to \mathbb{R}))) \wedge \\ \quad \forall k \in \mathbb{N}.\; ((k < (m + 1)) \Rightarrow \operatorname{is\text{-}continuous\text{-}on}(f(k), \operatorname{ccint}(a, b))) \wedge \\ \quad \forall k \in \mathbb{N}, x \in \operatorname{ooint}(a, b). \\ \quad \quad (k < (m + 1)) \Rightarrow \operatorname{has\text{-}deriv\text{-}at}(f(k), x, f((k + 1))(x)) \wedge \\ \quad \operatorname{is\text{-}continuous\text{-}on}(f((m + 1)), \operatorname{ccint}(a, b)) \Rightarrow \\ \quad \quad \operatorname{is\text{-}continuous\text{-}at}(\operatorname{rr\text{-}ms}, \operatorname{rr\text{-}ms}, \\ \quad \quad \quad (\lambda s \in \mathbb{R}.\; \\ \quad \quad \quad \quad \left(\text{if }(s \leq a)\text{ then }f((m + 1))(a)\text{ else }(((\operatorname{extend\text{-}const}(f(0), a, b)(s) - \operatorname{taylor\text{-}sum}(f, (m + 1), s, a)) \cdot \operatorname{factorial}((m + 1))) \cdot {\operatorname{power}((s - a), (m + 1))}^{\ast\!-1})\right)), \\ \quad \quad \quad p) \end{array}}}
\vnbprop{taylor-remainder-factor}{\vnbstmt{\begin{array}{@{}l@{}} \forall a, b \in \mathbb{R}, f.; \forall m \in \mathbb{N}. \\ \quad a < b \wedge \\ \quad \forall k \in \mathbb{N}.\; ((k \leq (m + 1)) \Rightarrow (f(k) \in (\operatorname{ccint}(a, b) \to \mathbb{R}))) \wedge \\ \quad \forall k \in \mathbb{N}.\; ((k < (m + 1)) \Rightarrow \operatorname{is\text{-}continuous\text{-}on}(f(k), \operatorname{ccint}(a, b))) \wedge \\ \quad \forall k \in \mathbb{N}, x \in \operatorname{ooint}(a, b). \\ \quad \quad (k < (m + 1)) \Rightarrow \operatorname{has\text{-}deriv\text{-}at}(f(k), x, f((k + 1))(x)) \wedge \\ \quad \operatorname{is\text{-}continuous\text{-}on}(f((m + 1)), \operatorname{ccint}(a, b)) \Rightarrow \\ \quad \quad \exists s \in (\operatorname{ccint}(a, b) \to \mathbb{R}). \\ \quad \quad \quad \operatorname{is\text{-}continuous\text{-}on}(s, \operatorname{ccint}(a, b)) \wedge \\ \quad \quad \quad s(a) = f((m + 1))(a) \wedge \\ \quad \quad \quad \forall t \in \operatorname{ccint}(a, b). \\ \quad \quad \quad \quad \operatorname{taylor\text{-}rem}(f, (m + 1), a, t) = ((s(t) \cdot \operatorname{power}((t - a), (m + 1))) \cdot \operatorname{recip}(\operatorname{factorial}((m + 1)))) \wedge \\ \quad \quad \quad \forall t \in \operatorname{ccint}(a, b). \\ \quad \quad \quad \quad \lvert s(t) \rvert \leq \operatorname{sup}(\operatorname{image}((\lambda v \in \operatorname{ccint}(a, b).\; \lvert f((m + 1))(v) \rvert), \operatorname{ccint}(a, b))) \end{array}}}

\vnbbookfile{theorem-library/frechet-metric}{theorem-library/frechet-metric}
\noindent{\small frechet-metric.scm -- a seminormed space with a COUNTABLE separating family of seminorms as a METRIC space: Thayer, "Analise Funcional e Teoria Espectral", Ch. 8 s.4, Props. 8.24-8.26 (pp. 254-256).  Batch TVS-1, 2026-10-10. Vocabulary: structure-library/frechet-metric.scm (FR-C-SPACE, FR-METRIC-SPACE,}\par
\vnbprop{frm-mvag-inv}{\vnbstmt{\forall d.\; (\operatorname{is\text{-}module}(d) \Rightarrow (\operatorname{inv}(\operatorname{module\text{-}vector\text{-}ag}(d)) = \operatorname{vneg}(d)))}}
\vnbprop{frm-v-right-inv}{\vnbstmt{\begin{array}{@{}l@{}} \forall m.; \forall x \in \operatorname{vec}(m). \\ \quad \operatorname{is\text{-}module}(m) \Rightarrow (\operatorname{vadd}(m)(x, \operatorname{vneg}(m)(x)) = \operatorname{vzero}(m)) \end{array}}}
\vnbprop{frm-v-neg-add}{\vnbstmt{\begin{array}{@{}l@{}} \forall m.; \forall x, y \in \operatorname{vec}(m). \\ \quad \operatorname{is\text{-}module}(m) \Rightarrow  \\ \quad \quad \operatorname{vneg}(m)(\operatorname{vadd}(m)(x, y)) = \operatorname{vadd}(m)(\operatorname{vneg}(m)(x), \operatorname{vneg}(m)(y)) \end{array}}}
\vnbprop{frm-v-neg-neg}{\vnbstmt{\forall m.; \forall x \in \operatorname{vec}(m).\; (\operatorname{is\text{-}module}(m) \Rightarrow (\operatorname{vneg}(m)(\operatorname{vneg}(m)(x)) = x))}}
\vnbprop{frm-v-sub-anti}{\vnbstmt{\begin{array}{@{}l@{}} \forall m.; \forall x, y \in \operatorname{vec}(m). \\ \quad \operatorname{is\text{-}module}(m) \Rightarrow  \\ \quad \quad \operatorname{vneg}(m)(\operatorname{vadd}(m)(x, \operatorname{vneg}(m)(y))) = \operatorname{vadd}(m)(y, \operatorname{vneg}(m)(x)) \end{array}}}
\vnbprop{frm-v-sub-chain}{\vnbstmt{\begin{array}{@{}l@{}} \forall m.; \forall x, y, z \in \operatorname{vec}(m). \\ \quad \operatorname{is\text{-}module}(m) \Rightarrow  \\ \quad \quad \operatorname{vadd}(m)(\operatorname{vadd}(m)(x, \operatorname{vneg}(m)(y)), \operatorname{vadd}(m)(y, \operatorname{vneg}(m)(z))) = \operatorname{vadd}(m)(x, \operatorname{vneg}(m)(z)) \end{array}}}
\vnbprop{frm-v-sub-zero-eq}{\vnbstmt{\begin{array}{@{}l@{}} \forall m.; \forall x, y \in \operatorname{vec}(m). \\ \quad (\operatorname{is\text{-}module}(m) \wedge (\operatorname{vadd}(m)(x, \operatorname{vneg}(m)(y)) = \operatorname{vzero}(m))) \Rightarrow (x = y) \end{array}}}
\vnbprop{frm-v-sub-shift}{\vnbstmt{\begin{array}{@{}l@{}} \forall m.; \forall x, y, r \in \operatorname{vec}(m). \\ \quad \operatorname{is\text{-}module}(m) \Rightarrow  \\ \quad \quad \operatorname{vadd}(m)(\operatorname{vadd}(m)(x, r), \operatorname{vneg}(m)(\operatorname{vadd}(m)(y, r))) = \operatorname{vadd}(m)(x, \operatorname{vneg}(m)(y)) \end{array}}}
\vnbprop{frm-v-sub-add}{\vnbstmt{\begin{array}{@{}l@{}} \forall m.; \forall x, y, u, v \in \operatorname{vec}(m). \\ \quad \operatorname{is\text{-}module}(m) \Rightarrow  \\ \quad \quad \operatorname{vadd}(m)(\operatorname{vadd}(m)(x, y), \operatorname{vneg}(m)(\operatorname{vadd}(m)(u, v))) \\ \quad \quad =\; \operatorname{vadd}(m)(\operatorname{vadd}(m)(x, \operatorname{vneg}(m)(u)), \operatorname{vadd}(m)(y, \operatorname{vneg}(m)(v))) \end{array}}}
\vnbprop{frm-c-space-pts}{\vnbstmt{\forall m, a.\; ((\operatorname{vec}(m) \in \mathbf{Set}) \Rightarrow (\operatorname{pts}(\operatorname{fr\text{-}c\text{-}space}(m, a)) \simeq \operatorname{vec}(m)))}}
\vnbprop{frm-c-space-dists}{\vnbstmt{\begin{array}{@{}l@{}} \forall m, a. \\ \quad (\operatorname{vec}(m) \in \mathbf{Set}) \Rightarrow  \\ \quad \quad (\operatorname{dists}(\operatorname{fr\text{-}c\text{-}space}(m, a)) \simeq (\lambda k \in \mathbb{N}.\; (\lambda \langle \operatorname{fru\_}, \operatorname{frv\_} \rangle \in (\operatorname{vec}(m) \times \operatorname{vec}(m)).\; a(k)(\operatorname{vadd}(m)(\operatorname{fru\_}, \operatorname{vneg}(m)(\operatorname{frv\_})))))) \end{array}}}
\vnbprop{frm-c-space-dists-apply}{\vnbstmt{\begin{array}{@{}l@{}} \forall m, a, k, x, y. \\ \quad ((\operatorname{vec}(m) \in \mathbf{Set}) \wedge ((k \in \mathbb{N}) \wedge ((x \in \operatorname{vec}(m)) \wedge (y \in \operatorname{vec}(m))))) \Rightarrow  \\ \quad \quad (\operatorname{dists}(\operatorname{fr\text{-}c\text{-}space}(m, a))(k)(x, y) \simeq a(k)(\operatorname{vadd}(m)(x, \operatorname{vneg}(m)(y)))) \end{array}}}
\vnbprop{frm-sn-module}{\vnbstmt{\forall d, m, p.\; (\operatorname{is\text{-}seminorm}(d, m, p) \Rightarrow \operatorname{is\text{-}module}(m))}}
\vnbprop{frm-sn-type}{\vnbstmt{\forall d, m, p.\; (\operatorname{is\text{-}seminorm}(d, m, p) \Rightarrow (p \in (\operatorname{vec}(m) \to \mathbb{R})))}}
\vnbprop{frm-sn-nonneg}{\vnbstmt{\forall d, m, p.; \forall x \in \operatorname{vec}(m).\; (\operatorname{is\text{-}seminorm}(d, m, p) \Rightarrow (0 \leq p(x)))}}
\vnbprop{frm-sn-subadd}{\vnbstmt{\begin{array}{@{}l@{}} \forall d, m, p.; \forall x, y \in \operatorname{vec}(m). \\ \quad \operatorname{is\text{-}seminorm}(d, m, p) \Rightarrow (p(\operatorname{vadd}(m)(x, y)) \leq (p(x) + p(y))) \end{array}}}
\vnbprop{frm-sn-zero}{\vnbstmt{\forall d, m, p.\; (\operatorname{is\text{-}seminorm}(d, m, p) \Rightarrow (p(\operatorname{vzero}(m)) = 0))}}
\vnbprop{frm-sn-neg}{\vnbstmt{\begin{array}{@{}l@{}} \forall d, m, p.; \forall x \in \operatorname{vec}(m). \\ \quad \operatorname{is\text{-}seminorm}(d, m, p) \Rightarrow (p(\operatorname{vneg}(m)(x)) = p(x)) \end{array}}}
\vnbprop{frm-fam-module}{\vnbstmt{\forall d, m, a.\; (\operatorname{is\text{-}seminorm\text{-}family}(d, m, a) \Rightarrow \operatorname{is\text{-}module}(m))}}
\vnbprop{frm-fam-type}{\vnbstmt{\begin{array}{@{}l@{}} \forall d, m, a. \\ \quad \operatorname{is\text{-}seminorm\text{-}family}(d, m, a) \Rightarrow (a \in (\mathbb{N} \to (\operatorname{vec}(m) \to \mathbb{R}))) \end{array}}}
\vnbprop{frm-fam-seminorm}{\vnbstmt{\begin{array}{@{}l@{}} \forall d, m, a.; \forall k \in \mathbb{N}. \\ \quad \operatorname{is\text{-}seminorm\text{-}family}(d, m, a) \Rightarrow \operatorname{is\text{-}seminorm}(d, m, a(k)) \end{array}}}
\vnbprop{frm-module-vec-set}{\vnbstmt{\forall m.\; (\operatorname{is\text{-}module}(m) \Rightarrow (\operatorname{vec}(m) \in \mathbf{Set}))}}
\vnbprop{frm-fam-separates}{\vnbstmt{\begin{array}{@{}l@{}} \forall d, m, a.; \forall x \in \operatorname{vec}(m). \\ \quad (\operatorname{is\text{-}seminorm\text{-}family}(d, m, a) \wedge \forall k \in \mathbb{N}.\; (a(k)(x) = 0)) \Rightarrow  \\ \quad \quad x = \operatorname{vzero}(m) \end{array}}}
\vnbprop{frm-c-space-is-c-metric-space}{\vnbstmt{\begin{array}{@{}l@{}} \forall d, m, a. \\ \quad \operatorname{is\text{-}seminorm\text{-}family}(d, m, a) \Rightarrow \operatorname{is\text{-}c\text{-}metric\text{-}space}(\operatorname{fr\text{-}c\text{-}space}(m, a)) \end{array}}}
\vnbprop{frm-metric-space-is-metric-space}{\vnbstmt{\begin{array}{@{}l@{}} \forall d, m, a. \\ \quad \operatorname{is\text{-}seminorm\text{-}family}(d, m, a) \Rightarrow \operatorname{is\text{-}metric\text{-}space}(\operatorname{fr\text{-}metric\text{-}space}(m, a)) \end{array}}}
\vnbprop{frm-metric-space-pts}{\vnbstmt{\begin{array}{@{}l@{}} \forall d, m, a. \\ \quad \operatorname{is\text{-}seminorm\text{-}family}(d, m, a) \Rightarrow (\operatorname{pts}(\operatorname{fr\text{-}metric\text{-}space}(m, a)) \simeq \operatorname{vec}(m)) \end{array}}}
\vnbprop{frm-c-ball-prefix}{\vnbstmt{\begin{array}{@{}l@{}} \forall s.; \forall k \in \mathbb{N}, a. \\ \quad (\operatorname{is\text{-}c\text{-}metric\text{-}space}(s) \wedge \operatorname{pos\text{-}rr}(a)) \Rightarrow  \\ \quad \quad \exists l. \\ \quad \quad \quad \operatorname{pos\text{-}rr}(l) \wedge \\ \quad \quad \quad \forall x, y \in \operatorname{pts}(s). \\ \quad \quad \quad \quad (\operatorname{dist}(\operatorname{c\text{-}metric}(s))(x, y) < l) \Rightarrow (\operatorname{dists}(s)(k)(x, y) < a) \end{array}}}
\vnbprop{frm-c-head-small}{\vnbstmt{\begin{array}{@{}l@{}} \forall s.; \forall x, y \in \operatorname{pts}(s), b \in \mathbb{N}, e \in \mathbb{R}. \\ \quad (\operatorname{is\text{-}c\text{-}metric\text{-}space}(s) \wedge ((0 \leq e) \wedge \forall j \in \mathbb{N}.\; ((j \leq b) \Rightarrow (\operatorname{dists}(s)(j)(x, y) < e)))) \Rightarrow  \\ \quad \quad \operatorname{series\text{-}partial\text{-}sum}((\lambda k \in \mathbb{N}.\; ((\lambda n \in \mathbb{N}.\; /(1, \operatorname{power}(2, (n + 1))))(k) \cdot \operatorname{min}(1, \operatorname{dists}(s)(k)(x, y)))), \\ \quad \quad \quad b) \\ \quad \quad \leq\; e \end{array}}}
\vnbprop{frm-c-term-bounds}{\vnbstmt{\begin{array}{@{}l@{}} \forall s, x, y.; \forall i \in \mathbb{N}. \\ \quad (\operatorname{is\text{-}c\text{-}metric\text{-}space}(s) \wedge ((x \in \operatorname{pts}(s)) \wedge (y \in \operatorname{pts}(s)))) \Rightarrow  \\ \quad \quad 0 \\ \quad \quad \leq\; (\lambda k \in \mathbb{N}.\; ((\lambda n \in \mathbb{N}.\; /(1, \operatorname{power}(2, (n + 1))))(k) \cdot \operatorname{min}(1, \operatorname{dists}(s)(k)(x, y))))(i) \wedge \\ \quad \quad (\lambda k \in \mathbb{N}.\; ((\lambda n \in \mathbb{N}.\; /(1, \operatorname{power}(2, (n + 1))))(k) \cdot \operatorname{min}(1, \operatorname{dists}(s)(k)(x, y))))(i) \\ \quad \quad \leq\; (\lambda n \in \mathbb{N}.\; /(1, \operatorname{power}(2, (n + 1))))(i) \end{array}}}
\vnbprop{frm-c-prefix-ball}{\vnbstmt{\begin{array}{@{}l@{}} \forall s, a. \\ \quad (\operatorname{is\text{-}c\text{-}metric\text{-}space}(s) \wedge \operatorname{pos\text{-}rr}(a)) \Rightarrow  \\ \quad \quad \exists b \in \mathbb{N}, l. \\ \quad \quad \quad \operatorname{pos\text{-}rr}(l) \wedge \\ \quad \quad \quad \forall x, y \in \operatorname{pts}(s). \\ \quad \quad \quad \quad \forall j \in \mathbb{N}.\; ((j \leq b) \Rightarrow (\operatorname{dists}(s)(j)(x, y) < l)) \Rightarrow  \\ \quad \quad \quad \quad \quad \operatorname{dist}(\operatorname{c\text{-}metric}(s))(x, y) \leq a \end{array}}}
\vnbprop{frm-ball-seminorm}{\vnbstmt{\begin{array}{@{}l@{}} \forall d, m, a.; \forall k \in \mathbb{N}, s. \\ \quad (\operatorname{is\text{-}seminorm\text{-}family}(d, m, a) \wedge \operatorname{pos\text{-}rr}(s)) \Rightarrow  \\ \quad \quad \exists l. \\ \quad \quad \quad \operatorname{pos\text{-}rr}(l) \wedge \\ \quad \quad \quad \forall x, y. \\ \quad \quad \quad \quad ((x \in \operatorname{vec}(m)) \wedge ((y \in \operatorname{vec}(m)) \wedge (\operatorname{dist}(\operatorname{fr\text{-}metric\text{-}space}(m, a))(x, y) < l))) \Rightarrow  \\ \quad \quad \quad \quad \quad a(k)(\operatorname{vadd}(m)(x, \operatorname{vneg}(m)(y))) < s \end{array}}}
\vnbprop{frm-seminorm-ball}{\vnbstmt{\begin{array}{@{}l@{}} \forall d, m, a, s. \\ \quad (\operatorname{is\text{-}seminorm\text{-}family}(d, m, a) \wedge \operatorname{pos\text{-}rr}(s)) \Rightarrow  \\ \quad \quad \exists b \in \mathbb{N}, l. \\ \quad \quad \quad \operatorname{pos\text{-}rr}(l) \wedge \\ \quad \quad \quad \forall x, y. \\ \quad \quad \quad \quad x \in \operatorname{vec}(m) \wedge \\ \quad \quad \quad \quad y \in \operatorname{vec}(m) \wedge \\ \quad \quad \quad \quad \forall j \in \mathbb{N}.\; ((j \leq b) \Rightarrow (a(j)(\operatorname{vadd}(m)(x, \operatorname{vneg}(m)(y))) < l)) \Rightarrow \\ \quad \quad \quad \quad \quad \operatorname{dist}(\operatorname{fr\text{-}metric\text{-}space}(m, a))(x, y) \leq s \end{array}}}
\vnbprop{frm-prefix-eventually}{\vnbstmt{\begin{array}{@{}l@{}} \forall b \in \mathbb{N}, u, l. \\ \quad \forall j \in \mathbb{N}. \exists p \in \mathbb{N}. \forall i, a \in \mathbb{N}. \\ \quad \quad ((p \leq i) \wedge (p \leq a)) \Rightarrow (u(j)(i)(a) < l) \Rightarrow \\ \quad \quad \exists t \in \mathbb{N}. \forall j, i, a \in \mathbb{N}. \\ \quad \quad \quad ((j \leq b) \wedge ((t \leq i) \wedge (t \leq a))) \Rightarrow (u(j)(i)(a) < l) \end{array}}}
\vnbprop{frm-conv-of-converges}{\vnbstmt{\begin{array}{@{}l@{}} \forall d, m, a, q, b. \\ \quad \operatorname{is\text{-}seminorm\text{-}family}(d, m, a) \wedge \\ \quad q \in (\mathbb{N} \to \operatorname{vec}(m)) \wedge \\ \quad b \in \operatorname{vec}(m) \wedge \\ \quad \operatorname{converges\text{-}to}(\operatorname{fr\text{-}metric\text{-}space}(m, a), q, b) \Rightarrow \\ \quad \quad \operatorname{fr\text{-}conv}(m, a, q, b) \end{array}}}
\vnbprop{frm-converges-of-conv}{\vnbstmt{\begin{array}{@{}l@{}} \forall d, m, a, q, b. \\ \quad (\operatorname{is\text{-}seminorm\text{-}family}(d, m, a) \wedge ((q \in (\mathbb{N} \to \operatorname{vec}(m))) \wedge ((b \in \operatorname{vec}(m)) \wedge \operatorname{fr\text{-}conv}(m, a, q, b)))) \Rightarrow  \\ \quad \quad \operatorname{converges\text{-}to}(\operatorname{fr\text{-}metric\text{-}space}(m, a), q, b) \end{array}}}
\vnbprop{frm-cauchy-of-metric}{\vnbstmt{\begin{array}{@{}l@{}} \forall d, m, a.; \forall q \in (\mathbb{N} \to \operatorname{vec}(m)). \\ \quad (\operatorname{is\text{-}seminorm\text{-}family}(d, m, a) \wedge \operatorname{is\text{-}cauchy\text{-}seq}(\operatorname{fr\text{-}metric\text{-}space}(m, a), q)) \Rightarrow  \\ \quad \quad \operatorname{fr\text{-}cauchy}(m, a, q) \end{array}}}
\vnbprop{frm-metric-cauchy-of-cauchy}{\vnbstmt{\begin{array}{@{}l@{}} \forall d, m, a.; \forall q \in (\mathbb{N} \to \operatorname{vec}(m)). \\ \quad (\operatorname{is\text{-}seminorm\text{-}family}(d, m, a) \wedge \operatorname{fr\text{-}cauchy}(m, a, q)) \Rightarrow  \\ \quad \quad \operatorname{is\text{-}cauchy\text{-}seq}(\operatorname{fr\text{-}metric\text{-}space}(m, a), q) \end{array}}}
\vnbprop{frm-converges-iff}{\vnbstmt{\begin{array}{@{}l@{}} \forall d, m, a, q, b. \\ \quad (\operatorname{is\text{-}seminorm\text{-}family}(d, m, a) \wedge ((q \in (\mathbb{N} \to \operatorname{vec}(m))) \wedge (b \in \operatorname{vec}(m)))) \Rightarrow  \\ \quad \quad \operatorname{converges\text{-}to}(\operatorname{fr\text{-}metric\text{-}space}(m, a), q, b) \Leftrightarrow \operatorname{fr\text{-}conv}(m, a, q, b) \end{array}}}
\vnbprop{frm-cauchy-iff}{\vnbstmt{\begin{array}{@{}l@{}} \forall d, m, a.; \forall q \in (\mathbb{N} \to \operatorname{vec}(m)). \\ \quad \operatorname{is\text{-}seminorm\text{-}family}(d, m, a) \Rightarrow  \\ \quad \quad \operatorname{is\text{-}cauchy\text{-}seq}(\operatorname{fr\text{-}metric\text{-}space}(m, a), q) \Leftrightarrow \operatorname{fr\text{-}cauchy}(m, a, q) \end{array}}}
\vnbprop{frm-frechet-iff-complete}{\vnbstmt{\begin{array}{@{}l@{}} \forall d, m, a. \\ \quad \operatorname{is\text{-}frechet\text{-}structure}(d, m, a) \Leftrightarrow (\operatorname{is\text{-}seminorm\text{-}family}(d, m, a) \wedge \operatorname{is\text{-}complete}(\operatorname{fr\text{-}metric\text{-}space}(m, a))) \end{array}}}
\vnbprop{frm-c-dist-congr}{\vnbstmt{\begin{array}{@{}l@{}} \forall s, x, y, a, b. \\ \quad \operatorname{is\text{-}c\text{-}metric\text{-}space}(s) \wedge \\ \quad x \in \operatorname{pts}(s) \wedge \\ \quad y \in \operatorname{pts}(s) \wedge \\ \quad a \in \operatorname{pts}(s) \wedge \\ \quad b \in \operatorname{pts}(s) \wedge \\ \quad \forall k \in \mathbb{N}.\; (\operatorname{dists}(s)(k)(x, y) = \operatorname{dists}(s)(k)(a, b)) \Rightarrow \\ \quad \quad \operatorname{dist}(\operatorname{c\text{-}metric}(s))(x, y) = \operatorname{dist}(\operatorname{c\text{-}metric}(s))(a, b) \end{array}}}
\vnbprop{frm-metric-translation-invariant}{\vnbstmt{\begin{array}{@{}l@{}} \forall d, m, a, x, y, r. \\ \quad (\operatorname{is\text{-}seminorm\text{-}family}(d, m, a) \wedge ((x \in \operatorname{vec}(m)) \wedge ((y \in \operatorname{vec}(m)) \wedge (r \in \operatorname{vec}(m))))) \Rightarrow  \\ \quad \quad \operatorname{dist}(\operatorname{fr\text{-}metric\text{-}space}(m, a))(\operatorname{vadd}(m)(x, r), \operatorname{vadd}(m)(y, r)) \\ \quad \quad =\; \operatorname{dist}(\operatorname{fr\text{-}metric\text{-}space}(m, a))(x, y) \end{array}}}
\vnbprop{frm-metric-value}{\vnbstmt{\begin{array}{@{}l@{}} \forall d, m, a, x, y. \\ \quad (\operatorname{is\text{-}seminorm\text{-}family}(d, m, a) \wedge ((x \in \operatorname{vec}(m)) \wedge (y \in \operatorname{vec}(m)))) \Rightarrow  \\ \quad \quad (\operatorname{fr\text{-}metric}(m, a)(x, y) \simeq \operatorname{series\text{-}limit}((\lambda k \in \mathbb{N}.\; (/(1, \operatorname{power}(2, (k + 1))) \cdot \operatorname{min}(1, a(k)(\operatorname{vadd}(m)(x, \operatorname{vneg}(m)(y)))))))) \end{array}}}
\vnbprop{frm-add-continuous}{\vnbstmt{\begin{array}{@{}l@{}} \forall d, m, c, f, g, a, b. \\ \quad \operatorname{is\text{-}seminorm\text{-}family}(d, m, c) \wedge \\ \quad f \in (\mathbb{N} \to \operatorname{vec}(m)) \wedge \\ \quad g \in (\mathbb{N} \to \operatorname{vec}(m)) \wedge \\ \quad a \in \operatorname{vec}(m) \wedge \\ \quad b \in \operatorname{vec}(m) \wedge \\ \quad \operatorname{converges\text{-}to}(\operatorname{fr\text{-}metric\text{-}space}(m, c), f, a) \wedge \\ \quad \operatorname{converges\text{-}to}(\operatorname{fr\text{-}metric\text{-}space}(m, c), g, b) \Rightarrow \\ \quad \quad \operatorname{converges\text{-}to}(\operatorname{fr\text{-}metric\text{-}space}(m, c), (\lambda n \in \mathbb{N}.\; \operatorname{vadd}(m)(f(n), g(n))), \operatorname{vadd}(m)(a, b)) \end{array}}}
\vnbprop{frm-linear-sub}{\vnbstmt{\begin{array}{@{}l@{}} \forall c, d, t, a, b. \\ \quad (\operatorname{is\text{-}linear\text{-}map}(c, d, t) \wedge ((a \in \operatorname{vec}(c)) \wedge (b \in \operatorname{vec}(c)))) \Rightarrow  \\ \quad \quad t(\operatorname{vadd}(c)(a, \operatorname{vneg}(c)(b))) = \operatorname{vadd}(d)(t(a), \operatorname{vneg}(d)(t(b))) \end{array}}}
\vnbprop{frm-cont-lin-of-bound}{\vnbstmt{\begin{array}{@{}l@{}} \forall d, a, b, c, f, t. \\ \quad \operatorname{is\text{-}seminorm\text{-}family}(d, a, b) \wedge \\ \quad \operatorname{is\text{-}seminorm\text{-}family}(d, c, f) \wedge \\ \quad \operatorname{is\text{-}linear\text{-}map}(a, c, t) \wedge \\ \quad \forall k \in \mathbb{N}. \exists g \in \mathbb{N}, h. \\ \quad \quad \operatorname{pos\text{-}rr}(h) \wedge \\ \quad \quad \forall x \in \operatorname{vec}(a), r \in \mathbb{R}. \\ \quad \quad \quad \forall j \in \mathbb{N}.\; ((j \leq g) \Rightarrow (b(j)(x) \leq r)) \Rightarrow (f(k)(t(x)) \leq (h \cdot r)) \Rightarrow \\ \quad \quad \operatorname{is\text{-}cont\text{-}lin}(a, b, c, f, t) \end{array}}}
\vnbprop{frm-ti-uniform}{\vnbstmt{\begin{array}{@{}l@{}} \forall m, a, b. \\ \quad \operatorname{is\text{-}module}(m) \wedge \\ \quad \operatorname{is\text{-}metric\text{-}space}(a) \wedge \\ \quad \operatorname{is\text{-}metric\text{-}space}(b) \wedge \\ \quad (\operatorname{pts}(a) \simeq \operatorname{vec}(m)) \wedge \\ \quad (\operatorname{pts}(b) \simeq \operatorname{vec}(m)) \wedge \\ \quad \forall x, y, r. \\ \quad \quad ((x \in \operatorname{vec}(m)) \wedge ((y \in \operatorname{vec}(m)) \wedge (r \in \operatorname{vec}(m)))) \Rightarrow  \\ \quad \quad \quad \operatorname{dist}(a)(\operatorname{vadd}(m)(x, r), \operatorname{vadd}(m)(y, r)) = \operatorname{dist}(a)(x, y) \wedge \\ \quad \forall x, y, r. \\ \quad \quad ((x \in \operatorname{vec}(m)) \wedge ((y \in \operatorname{vec}(m)) \wedge (r \in \operatorname{vec}(m)))) \Rightarrow  \\ \quad \quad \quad \operatorname{dist}(b)(\operatorname{vadd}(m)(x, r), \operatorname{vadd}(m)(y, r)) = \operatorname{dist}(b)(x, y) \wedge \\ \quad \forall q \in (\mathbb{N} \to \operatorname{vec}(m)). \\ \quad \quad \operatorname{converges\text{-}to}(a, q, \operatorname{vzero}(m)) \Rightarrow \operatorname{converges\text{-}to}(b, q, \operatorname{vzero}(m)) \Rightarrow \\ \quad \quad \operatorname{is\text{-}uniformly\text{-}continuous}(a, b, \operatorname{id\text{-}fun}(\operatorname{vec}(m))) \end{array}}}
\vnbprop{frm-metrics-uniformly-equivalent}{\vnbstmt{\begin{array}{@{}l@{}} \forall d, m, a, t. \\ \quad \operatorname{is\text{-}seminorm\text{-}family}(d, m, a) \wedge \\ \quad \operatorname{is\text{-}metric\text{-}space}(t) \wedge \\ \quad (\operatorname{pts}(t) \simeq \operatorname{vec}(m)) \wedge \\ \quad \forall x, y, r. \\ \quad \quad ((x \in \operatorname{vec}(m)) \wedge ((y \in \operatorname{vec}(m)) \wedge (r \in \operatorname{vec}(m)))) \Rightarrow  \\ \quad \quad \quad \operatorname{dist}(t)(\operatorname{vadd}(m)(x, r), \operatorname{vadd}(m)(y, r)) = \operatorname{dist}(t)(x, y) \wedge \\ \quad \forall q, b. \\ \quad \quad ((q \in (\mathbb{N} \to \operatorname{vec}(m))) \wedge (b \in \operatorname{vec}(m))) \Rightarrow  \\ \quad \quad \quad \operatorname{converges\text{-}to}(t, q, b) \Leftrightarrow \operatorname{fr\text{-}conv}(m, a, q, b) \Rightarrow \\ \quad \quad \operatorname{is\text{-}uniformly\text{-}continuous}(\operatorname{fr\text{-}metric\text{-}space}(m, a), t, \operatorname{id\text{-}fun}(\operatorname{vec}(m))) \wedge \\ \quad \quad \operatorname{is\text{-}uniformly\text{-}continuous}(t, \operatorname{fr\text{-}metric\text{-}space}(m, a), \operatorname{id\text{-}fun}(\operatorname{vec}(m))) \end{array}}}
\vnbprop{frm-sn-homog}{\vnbstmt{\begin{array}{@{}l@{}} \forall d, m, p.; \forall a \in \operatorname{carr}(\operatorname{scal}(m)), x \in \operatorname{vec}(m). \\ \quad \operatorname{is\text{-}seminorm}(d, m, p) \Rightarrow (p(\operatorname{act}(m)(a, x)) = (\operatorname{fnrm}(d)(a) \cdot p(x))) \end{array}}}
\vnbprop{frm-fam-field}{\vnbstmt{\forall d, m, a.\; (\operatorname{is\text{-}seminorm\text{-}family}(d, m, a) \Rightarrow \operatorname{is\text{-}normed\text{-}field}(d))}}
\vnbprop{frm-fam-scal}{\vnbstmt{\begin{array}{@{}l@{}} \forall d, m, a. \\ \quad \operatorname{is\text{-}seminorm\text{-}family}(d, m, a) \Rightarrow  \\ \quad \quad \operatorname{scal}(m) = \operatorname{normed\text{-}field\text{-}as\text{-}commutative\text{-}ring}(d) \end{array}}}
\vnbprop{frm-act-neg-vec}{\vnbstmt{\begin{array}{@{}l@{}} \forall m, r, x. \\ \quad (\operatorname{is\text{-}module}(m) \wedge ((r \in \operatorname{carr}(\operatorname{scal}(m))) \wedge (x \in \operatorname{vec}(m)))) \Rightarrow  \\ \quad \quad \operatorname{act}(m)(r, \operatorname{vneg}(m)(x)) = \operatorname{vneg}(m)(\operatorname{act}(m)(r, x)) \end{array}}}
\vnbprop{frm-act-neg-scal}{\vnbstmt{\begin{array}{@{}l@{}} \forall m, r, x. \\ \quad (\operatorname{is\text{-}module}(m) \wedge ((r \in \operatorname{carr}(\operatorname{scal}(m))) \wedge (x \in \operatorname{vec}(m)))) \Rightarrow  \\ \quad \quad \operatorname{act}(m)(\operatorname{neg}(\operatorname{scal}(m))(r), x) = \operatorname{vneg}(m)(\operatorname{act}(m)(r, x)) \end{array}}}
\vnbprop{frm-act-sub-vec}{\vnbstmt{\begin{array}{@{}l@{}} \forall m, r, x, y. \\ \quad (\operatorname{is\text{-}module}(m) \wedge ((r \in \operatorname{carr}(\operatorname{scal}(m))) \wedge ((x \in \operatorname{vec}(m)) \wedge (y \in \operatorname{vec}(m))))) \Rightarrow  \\ \quad \quad \operatorname{act}(m)(r, \operatorname{vadd}(m)(x, \operatorname{vneg}(m)(y))) = \operatorname{vadd}(m)(\operatorname{act}(m)(r, x), \operatorname{vneg}(m)(\operatorname{act}(m)(r, y))) \end{array}}}
\vnbprop{frm-act-sub-scal}{\vnbstmt{\begin{array}{@{}l@{}} \forall m, r, s, x. \\ \quad (\operatorname{is\text{-}module}(m) \wedge ((r \in \operatorname{carr}(\operatorname{scal}(m))) \wedge ((s \in \operatorname{carr}(\operatorname{scal}(m))) \wedge (x \in \operatorname{vec}(m))))) \Rightarrow  \\ \quad \quad \operatorname{act}(m)(\operatorname{add}(\operatorname{scal}(m))(r, \operatorname{neg}(\operatorname{scal}(m))(s)), x) \\ \quad \quad =\; \operatorname{vadd}(m)(\operatorname{act}(m)(r, x), \operatorname{vneg}(m)(\operatorname{act}(m)(s, x))) \end{array}}}
\vnbprop{frm-act-diff}{\vnbstmt{\begin{array}{@{}l@{}} \forall m, r, s, x, y. \\ \quad (\operatorname{is\text{-}module}(m) \wedge ((r \in \operatorname{carr}(\operatorname{scal}(m))) \wedge ((s \in \operatorname{carr}(\operatorname{scal}(m))) \wedge ((x \in \operatorname{vec}(m)) \wedge (y \in \operatorname{vec}(m)))))) \Rightarrow  \\ \quad \quad \operatorname{vadd}(m)(\operatorname{act}(m)(r, x), \operatorname{vneg}(m)(\operatorname{act}(m)(s, y))) \\ \quad \quad =\; \operatorname{vadd}(m)(\operatorname{act}(m)(r, \operatorname{vadd}(m)(x, \operatorname{vneg}(m)(y))), \\ \quad \quad \quad \operatorname{act}(m)(\operatorname{add}(\operatorname{scal}(m))(r, \operatorname{neg}(\operatorname{scal}(m))(s)), y)) \end{array}}}
\vnbprop{frm-scal-continuous}{\vnbstmt{\begin{array}{@{}l@{}} \forall d, m, a.; \forall s \in (\mathbb{N} \to \operatorname{carr}(d)), b \in (\mathbb{N} \to \operatorname{vec}(m)), l \in \operatorname{carr}(d), x \in \operatorname{vec}(m). \\ \quad \operatorname{is\text{-}seminorm\text{-}family}(d, m, a) \wedge \\ \quad \operatorname{converges\text{-}to}(\operatorname{nf\text{-}metric\text{-}space}(d), s, l) \wedge \\ \quad \operatorname{converges\text{-}to}(\operatorname{fr\text{-}metric\text{-}space}(m, a), b, x) \Rightarrow \\ \quad \quad \operatorname{converges\text{-}to}(\operatorname{fr\text{-}metric\text{-}space}(m, a), (\lambda n \in \mathbb{N}.\; \operatorname{act}(m)(s(n), b(n))), \operatorname{act}(m)(l, x)) \end{array}}}
\vnbprop{frm-v-sub-zero}{\vnbstmt{\begin{array}{@{}l@{}} \forall m.; \forall x \in \operatorname{vec}(m). \\ \quad \operatorname{is\text{-}module}(m) \Rightarrow (\operatorname{vadd}(m)(x, \operatorname{vneg}(m)(\operatorname{vzero}(m))) = x) \end{array}}}
\vnbprop{frm-linear-zero}{\vnbstmt{\forall a, b, t.\; (\operatorname{is\text{-}linear\text{-}map}(a, b, t) \Rightarrow (t(\operatorname{vzero}(a)) = \operatorname{vzero}(b)))}}
\vnbprop{frm-linear-homog}{\vnbstmt{\begin{array}{@{}l@{}} \forall a, b, t.; \forall m \in \operatorname{carr}(\operatorname{scal}(a)), x \in \operatorname{vec}(a). \\ \quad \operatorname{is\text{-}linear\text{-}map}(a, b, t) \Rightarrow (t(\operatorname{act}(a)(m, x)) = \operatorname{act}(b)(m, t(x))) \end{array}}}
\vnbprop{frm-cont-lin-seq}{\vnbstmt{\begin{array}{@{}l@{}} \forall a, b, c, d, t.; \forall q \in (\mathbb{N} \to \operatorname{vec}(a)), m \in \operatorname{vec}(a). \\ \quad (\operatorname{is\text{-}cont\text{-}lin}(a, b, c, d, t) \wedge \operatorname{fr\text{-}conv}(a, b, q, m)) \Rightarrow  \\ \quad \quad \operatorname{fr\text{-}conv}(c, d, (\lambda i \in \mathbb{N}.\; t(q(i))), t(m)) \end{array}}}
\vnbprop{frm-cont-lin-nbhd}{\vnbstmt{\begin{array}{@{}l@{}} \forall d, a, b, c, f, t.; \forall k \in \mathbb{N}. \\ \quad \operatorname{is\text{-}seminorm\text{-}family}(d, a, b) \wedge \\ \quad \operatorname{is\text{-}seminorm\text{-}family}(d, c, f) \wedge \\ \quad \operatorname{is\text{-}linear\text{-}map}(a, c, t) \wedge \\ \quad \operatorname{is\text{-}cont\text{-}lin}(a, b, c, f, t) \Rightarrow \\ \quad \quad \exists g \in \mathbb{N}, l. \\ \quad \quad \quad \operatorname{pos\text{-}rr}(l) \wedge \\ \quad \quad \quad \forall x \in \operatorname{vec}(a). \\ \quad \quad \quad \quad \forall j \in \mathbb{N}.\; ((j \leq g) \Rightarrow (b(j)(x) < l)) \Rightarrow (f(k)(t(x)) < 1) \end{array}}}
\vnbprop{frm-nf-norm-one}{\vnbstmt{\begin{array}{@{}l@{}} \forall d.; \forall c \in \operatorname{carr}(d). \\ \quad (\operatorname{is\text{-}normed\text{-}field}(d) \wedge (0 < \operatorname{fnrm}(d)(c))) \Rightarrow (\operatorname{fnrm}(d)(\operatorname{one}(d)) = 1) \end{array}}}
\vnbprop{frm-nf-ladder}{\vnbstmt{\begin{array}{@{}l@{}} \forall n \in \mathbb{N}, d.; \forall c, b \in \operatorname{carr}(d). \\ \quad (\operatorname{is\text{-}normed\text{-}field}(d) \wedge ((1 < \operatorname{fnrm}(d)(c)) \wedge (\operatorname{mul}(d)(c, b) = \operatorname{one}(d)))) \Rightarrow  \\ \quad \quad \exists \mathit{fsc}, \mathit{fsb} \in \operatorname{carr}(d). \\ \quad \quad \quad (\operatorname{fnrm}(d)(\mathit{fsc}) \cdot \operatorname{fnrm}(d)(\mathit{fsb})) = 1 \wedge \\ \quad \quad \quad 1 \leq \operatorname{fnrm}(d)(\mathit{fsc}) \wedge \\ \quad \quad \quad (1 + ((n \cdot \operatorname{fnrm}(d)(c)) - n)) \leq \operatorname{fnrm}(d)(\mathit{fsc}) \wedge \\ \quad \quad \quad \forall u, v \in \mathbb{R}. \\ \quad \quad \quad \quad ((0 < u) \wedge ((0 < v) \wedge ((v < u) \wedge (u \leq (\operatorname{fnrm}(d)(\mathit{fsc}) \cdot v))))) \Rightarrow  \\ \quad \quad \quad \quad \quad \exists s \in \operatorname{carr}(d). \\ \quad \quad \quad \quad \quad \quad ((\operatorname{fnrm}(d)(s) \cdot v) < u) \wedge (u \leq (\operatorname{fnrm}(d)(c) \cdot (\operatorname{fnrm}(d)(s) \cdot v))) \wedge \\ \quad \quad \quad \forall u, v \in \mathbb{R}. \\ \quad \quad \quad \quad ((0 < u) \wedge ((0 < v) \wedge (((\operatorname{fnrm}(d)(\mathit{fsb}) \cdot v) < u) \wedge (u \leq v)))) \Rightarrow  \\ \quad \quad \quad \quad \quad \exists s \in \operatorname{carr}(d). \\ \quad \quad \quad \quad \quad \quad ((\operatorname{fnrm}(d)(s) \cdot v) < u) \wedge (u \leq (\operatorname{fnrm}(d)(c) \cdot (\operatorname{fnrm}(d)(s) \cdot v))) \end{array}}}
\vnbprop{frm-nf-bracket}{\vnbstmt{\begin{array}{@{}l@{}} \forall d.; \forall c, b \in \operatorname{carr}(d), u, v \in \mathbb{R}. \\ \quad (\operatorname{is\text{-}normed\text{-}field}(d) \wedge ((1 < \operatorname{fnrm}(d)(c)) \wedge ((\operatorname{mul}(d)(c, b) = \operatorname{one}(d)) \wedge ((0 < u) \wedge (0 < v))))) \Rightarrow  \\ \quad \quad \exists s \in \operatorname{carr}(d). \\ \quad \quad \quad ((\operatorname{fnrm}(d)(s) \cdot v) < u) \wedge (u \leq (\operatorname{fnrm}(d)(c) \cdot (\operatorname{fnrm}(d)(s) \cdot v))) \end{array}}}
\vnbprop{frm-bound-of-cont-lin}{\vnbstmt{\begin{array}{@{}l@{}} \forall d, a, b, c, f, t.; \forall k \in \mathbb{N}. \\ \quad \operatorname{is\text{-}seminorm\text{-}family}(d, a, b) \wedge \\ \quad \operatorname{is\text{-}seminorm\text{-}family}(d, c, f) \wedge \\ \quad \operatorname{is\text{-}linear\text{-}map}(a, c, t) \wedge \\ \quad \operatorname{is\text{-}nontrivially\text{-}valued}(d) \wedge \\ \quad \operatorname{is\text{-}cont\text{-}lin}(a, b, c, f, t) \Rightarrow \\ \quad \quad \exists g \in \mathbb{N}, h. \\ \quad \quad \quad \operatorname{pos\text{-}rr}(h) \wedge \\ \quad \quad \quad \forall x \in \operatorname{vec}(a), r \in \mathbb{R}. \\ \quad \quad \quad \quad \forall j \in \mathbb{N}.\; ((j \leq g) \Rightarrow (b(j)(x) \leq r)) \Rightarrow (f(k)(t(x)) \leq (h \cdot r)) \end{array}}}
\vnbprop{frm-cont-lin-iff-bound}{\vnbstmt{\begin{array}{@{}l@{}} \forall d, a, b, c, f, t. \\ \quad \operatorname{is\text{-}seminorm\text{-}family}(d, a, b) \wedge \\ \quad \operatorname{is\text{-}seminorm\text{-}family}(d, c, f) \wedge \\ \quad \operatorname{is\text{-}linear\text{-}map}(a, c, t) \wedge \\ \quad \operatorname{is\text{-}nontrivially\text{-}valued}(d) \Rightarrow \\ \quad \quad \operatorname{is\text{-}cont\text{-}lin}(a, b, c, f, t) \\ \quad \quad \Leftrightarrow\; \forall k \in \mathbb{N}. \exists g \in \mathbb{N}, h. \\ \quad \quad \quad \operatorname{pos\text{-}rr}(h) \wedge \\ \quad \quad \quad \forall x \in \operatorname{vec}(a), r \in \mathbb{R}. \\ \quad \quad \quad \quad \forall j \in \mathbb{N}.\; ((j \leq g) \Rightarrow (b(j)(x) \leq r)) \Rightarrow (f(k)(t(x)) \leq (h \cdot r)) \end{array}}}
\vnbprop{frm-rr-nontrivially-valued}{\vnbstmt{\operatorname{is\text{-}nontrivially\text{-}valued}(\operatorname{rr\text{-}normed\text{-}field})}}

\vnbbookfile{theorem-library/frechet-open-mapping-proof}{theorem-library/frechet-open-mapping-proof}
\noindent{\small frechet-open-mapping-proof.scm -- the OPEN MAPPING THEOREM for Frechet spaces and its bounded-inverse corollary: Thayer, "Analise Funcional e Teoria Espectral", Ch. 8 s.4, Thm 8.30, proved by the roadmap of exercise 8.9 (a)-(i).  Batch TVS-2A, 2026-10-11/12.  Helper prefix: fom-.  Every theorem here is proven modulo 0; ...}\par
\vnbprop{fom-prefix-bound}{\vnbstmt{\begin{array}{@{}l@{}} \forall b \in \mathbb{N}, m, x. \\ \quad \forall i \in \mathbb{N}.\; (m(i)(x) \in \mathbb{R}) \Rightarrow  \\ \quad \quad \exists p \in \mathbb{N}.\; \forall i \in \mathbb{N}.\; ((i \leq b) \Rightarrow (m(i)(x) < p)) \end{array}}}
\vnbprop{fom-dist-prefix}{\vnbstmt{\begin{array}{@{}l@{}} \forall k \in \mathbb{N}, d, m, a, s. \\ \quad (\operatorname{is\text{-}seminorm\text{-}family}(d, m, a) \wedge \operatorname{pos\text{-}rr}(s)) \Rightarrow  \\ \quad \quad \exists l. \\ \quad \quad \quad \operatorname{pos\text{-}rr}(l) \wedge \\ \quad \quad \quad \forall x, y.; \forall i \in \mathbb{N}. \\ \quad \quad \quad \quad ((x \in \operatorname{vec}(m)) \wedge ((y \in \operatorname{vec}(m)) \wedge ((\operatorname{dist}(\operatorname{fr\text{-}metric\text{-}space}(m, a))(x, y) < l) \wedge (i \leq k)))) \Rightarrow  \\ \quad \quad \quad \quad \quad a(i)(\operatorname{vadd}(m)(x, \operatorname{vneg}(m)(y))) < s \end{array}}}
\vnbprop{fom-unit-ladder}{\vnbstmt{\begin{array}{@{}l@{}} \forall n \in \mathbb{N}, d.; \forall c, b \in \operatorname{carr}(d). \\ \quad (\operatorname{is\text{-}normed\text{-}field}(d) \wedge ((1 < \operatorname{fnrm}(d)(c)) \wedge (\operatorname{mul}(d)(c, b) = \operatorname{one}(d)))) \Rightarrow  \\ \quad \quad \exists \mathit{fsc}, \mathit{fsb} \in \operatorname{carr}(d). \\ \quad \quad \quad ((\operatorname{mul}(d)(\mathit{fsc}, \mathit{fsb}) = \operatorname{one}(d)) \wedge ((1 \leq \operatorname{fnrm}(d)(\mathit{fsc})) \wedge ((1 + ((n \cdot \operatorname{fnrm}(d)(c)) - n)) \leq \operatorname{fnrm}(d)(\mathit{fsc})))) \end{array}}}
\vnbprop{fom-small-unit}{\vnbstmt{\begin{array}{@{}l@{}} \forall d.; \forall r \in \mathbb{R}, l. \\ \quad (\operatorname{is\text{-}nontrivially\text{-}valued}(d) \wedge ((0 \leq r) \wedge \operatorname{pos\text{-}rr}(l))) \Rightarrow  \\ \quad \quad \exists g, a \in \operatorname{carr}(d). \\ \quad \quad \quad \operatorname{mul}(d)(g, a) = \operatorname{one}(d) \wedge \\ \quad \quad \quad (\operatorname{fnrm}(d)(g) \cdot \operatorname{fnrm}(d)(a)) = 1 \wedge \\ \quad \quad \quad 0 < \operatorname{fnrm}(d)(g) \wedge \\ \quad \quad \quad 0 < \operatorname{fnrm}(d)(a) \wedge \\ \quad \quad \quad (\operatorname{fnrm}(d)(g) \cdot r) < l \end{array}}}
\vnbprop{fom-closure-near}{\vnbstmt{\begin{array}{@{}l@{}} \forall s, a.; \forall z \in \operatorname{closure}(s, a), o. \\ \quad (\operatorname{is\text{-}metric\text{-}space}(s) \wedge \operatorname{pos\text{-}rr}(o)) \Rightarrow \exists t \in a.\; (\operatorname{dist}(s)(z, t) < o) \end{array}}}
\vnbprop{fom-baire-ball}{\vnbstmt{\begin{array}{@{}l@{}} \forall s.; \forall a \in \operatorname{pts}(s), e \in (\mathbb{N} \to \operatorname{power}(\operatorname{pts}(s))). \\ \quad (\operatorname{is\text{-}complete}(s) \wedge \forall y \in \operatorname{pts}(s).\; \exists j \in \mathbb{N}.\; (y \in e(j))) \Rightarrow  \\ \quad \quad \exists j \in \mathbb{N}, c \in \operatorname{pts}(s), r. \\ \quad \quad \quad (\operatorname{pos\text{-}rr}(r) \wedge \forall z \in \operatorname{pts}(s).\; ((\operatorname{dist}(s)(c, z) < r) \Rightarrow (z \in \operatorname{closure}(s, e(j))))) \end{array}}}
\vnbprop{fom-v-shift-back}{\vnbstmt{\begin{array}{@{}l@{}} \forall m.; \forall w, c \in \operatorname{vec}(m). \\ \quad \operatorname{is\text{-}module}(m) \Rightarrow (\operatorname{vadd}(m)(\operatorname{vadd}(m)(w, c), \operatorname{vneg}(m)(c)) = w) \end{array}}}
\vnbprop{fom-v-quad}{\vnbstmt{\begin{array}{@{}l@{}} \forall m.; \forall p, q, a, b \in \operatorname{vec}(m). \\ \quad \operatorname{is\text{-}module}(m) \Rightarrow  \\ \quad \quad \operatorname{vadd}(m)(\operatorname{vadd}(m)(p, \operatorname{vneg}(m)(q)), \operatorname{vneg}(m)(\operatorname{vadd}(m)(a, \operatorname{vneg}(m)(b)))) \\ \quad \quad =\; \operatorname{vadd}(m)(\operatorname{vadd}(m)(p, \operatorname{vneg}(m)(a)), \operatorname{vneg}(m)(\operatorname{vadd}(m)(q, \operatorname{vneg}(m)(b)))) \end{array}}}
\vnbprop{fom-v-sub-sum}{\vnbstmt{\begin{array}{@{}l@{}} \forall m.; \forall y, a, b \in \operatorname{vec}(m). \\ \quad \operatorname{is\text{-}module}(m) \Rightarrow  \\ \quad \quad \operatorname{vadd}(m)(y, \operatorname{vneg}(m)(\operatorname{vadd}(m)(a, b))) = \operatorname{vadd}(m)(\operatorname{vadd}(m)(y, \operatorname{vneg}(m)(b)), \operatorname{vneg}(m)(a)) \end{array}}}
\vnbprop{fom-lin-m1}{\vnbstmt{\forall a, b, t.\; (\operatorname{is\text{-}linear\text{-}map}(a, b, t) \Rightarrow \operatorname{is\text{-}module}(a))}}
\vnbprop{fom-lin-m2}{\vnbstmt{\forall a, b, t.\; (\operatorname{is\text{-}linear\text{-}map}(a, b, t) \Rightarrow \operatorname{is\text{-}module}(b))}}
\vnbprop{fom-lin-scal}{\vnbstmt{\forall a, b, t.\; (\operatorname{is\text{-}linear\text{-}map}(a, b, t) \Rightarrow (\operatorname{scal}(a) = \operatorname{scal}(b)))}}
\vnbprop{fom-lin-type}{\vnbstmt{\forall a, b, t.\; (\operatorname{is\text{-}linear\text{-}map}(a, b, t) \Rightarrow (t \in (\operatorname{vec}(a) \to \operatorname{vec}(b))))}}
\vnbprop{fom-lin-add}{\vnbstmt{\begin{array}{@{}l@{}} \forall a, b, t.; \forall x, y \in \operatorname{vec}(a). \\ \quad \operatorname{is\text{-}linear\text{-}map}(a, b, t) \Rightarrow (t(\operatorname{vadd}(a)(x, y)) = \operatorname{vadd}(b)(t(x), t(y))) \end{array}}}
\vnbprop{fom-sn-sub-le}{\vnbstmt{\begin{array}{@{}l@{}} \forall d, m, p.; \forall u, v \in \operatorname{vec}(m). \\ \quad \operatorname{is\text{-}seminorm}(d, m, p) \Rightarrow (p(\operatorname{vadd}(m)(u, \operatorname{vneg}(m)(v))) \leq (p(u) + p(v))) \end{array}}}
\vnbprop{fom-act-cancel}{\vnbstmt{\begin{array}{@{}l@{}} \forall d, m.; \forall g, a \in \operatorname{carr}(d), z \in \operatorname{vec}(m). \\ \quad \operatorname{is\text{-}normed\text{-}field}(d) \wedge \\ \quad \operatorname{is\text{-}module}(m) \wedge \\ \quad \operatorname{scal}(m) = \operatorname{normed\text{-}field\text{-}as\text{-}commutative\text{-}ring}(d) \wedge \\ \quad \operatorname{mul}(d)(g, a) = \operatorname{one}(d) \Rightarrow \\ \quad \quad \operatorname{act}(m)(g, \operatorname{act}(m)(a, z)) = z \end{array}}}
\vnbprop{fom-mul-lt-left}{\vnbstmt{\forall a, b, c \in \mathbb{R}.\; (((0 < c) \wedge (a < b)) \Rightarrow ((c \cdot a) < (c \cdot b)))}}
\vnbprop{fom-baire-approx}{\vnbstmt{\begin{array}{@{}l@{}} \forall d, a, b, c, f, t.; \forall n \in \mathbb{N}, l. \\ \quad \operatorname{is\text{-}nontrivially\text{-}valued}(d) \wedge \\ \quad \operatorname{is\text{-}seminorm\text{-}family}(d, a, b) \wedge \\ \quad \operatorname{is\text{-}frechet\text{-}structure}(d, c, f) \wedge \\ \quad \operatorname{is\text{-}linear\text{-}map}(a, c, t) \wedge \\ \quad \forall y \in \operatorname{vec}(c).\; \exists x \in \operatorname{vec}(a).\; (y = t(x)) \wedge \\ \quad \operatorname{pos\text{-}rr}(l) \Rightarrow \\ \quad \quad \exists \mathit{fnb} \in \mathbb{N}, \mathit{feta}. \\ \quad \quad \quad \operatorname{pos\text{-}rr}(\mathit{feta}) \wedge \\ \quad \quad \quad \forall z \in \operatorname{vec}(c), k \in \mathbb{N}, h. \\ \quad \quad \quad \quad (\forall i \in \mathbb{N}.\; ((i \leq \mathit{fnb}) \Rightarrow (f(i)(z) < \mathit{feta})) \wedge \operatorname{pos\text{-}rr}(h)) \Rightarrow  \\ \quad \quad \quad \quad \quad \exists \mathit{fx} \in \operatorname{vec}(a). \\ \quad \quad \quad \quad \quad \quad \forall j \in \mathbb{N}.\; ((j \leq n) \Rightarrow (b(j)(\mathit{fx}) < l)) \wedge \\ \quad \quad \quad \quad \quad \quad \forall i \in \mathbb{N}.\; ((i \leq k) \Rightarrow (f(i)(\operatorname{vadd}(c)(z, \operatorname{vneg}(c)(t(\mathit{fx})))) < h)) \end{array}}}
\vnbprop{fom-a-facts}{\vnbstmt{\begin{array}{@{}l@{}} \forall a, k \in \mathbb{N}. \\ \quad \operatorname{recip}(((a + k) + 1)) \in \mathbb{R} \wedge \\ \quad 0 < \operatorname{recip}(((a + k) + 1)) \wedge \\ \quad \operatorname{recip}(((a + (k + 1)) + 1)) < \operatorname{recip}(((a + k) + 1)) \end{array}}}
\vnbprop{fom-a-small}{\vnbstmt{\begin{array}{@{}l@{}} \forall d. \\ \quad \operatorname{pos\text{-}rr}(d) \Rightarrow  \\ \quad \quad \exists p \in \mathbb{N}. \forall a, k \in \mathbb{N}. \\ \quad \quad \quad (p \leq k) \Rightarrow (\operatorname{recip}(((a + k) + 1)) \leq d) \end{array}}}
\vnbprop{fom-telescope}{\vnbstmt{\begin{array}{@{}l@{}} \forall l \in \mathbb{N}, d, m, a, f, v, b.; \forall k, j \in \mathbb{N}. \\ \quad \operatorname{is\text{-}seminorm\text{-}family}(d, m, a) \wedge \\ \quad f \in (\mathbb{N} \to \operatorname{vec}(m)) \wedge \\ \quad v \in \mathbb{N} \wedge \\ \quad b \in \mathbb{N} \wedge \\ \quad \forall i, j \in \mathbb{N}. \\ \quad \quad (j \leq (v + i)) \Rightarrow  \\ \quad \quad \quad a(j)(\operatorname{vadd}(m)(f((i + 1)), \operatorname{vneg}(m)(f(i)))) < (\operatorname{recip}(((b + i) + 1)) - \operatorname{recip}(((b + (i + 1)) + 1))) \wedge \\ \quad k \leq l \wedge \\ \quad j \leq (v + k) \Rightarrow \\ \quad \quad a(j)(\operatorname{vadd}(m)(f(l), \operatorname{vneg}(m)(f(k)))) \leq (\operatorname{recip}(((b + k) + 1)) - \operatorname{recip}(((b + l) + 1))) \end{array}}}
\vnbprop{fom-fr-ball-unfold}{\vnbstmt{\begin{array}{@{}l@{}} \forall m, a, n, s. \\ \quad (\operatorname{fr\text{-}ball}(m, a, n, s) \simeq \{x \in \operatorname{vec}(m) : \forall k \in \mathbb{N}.\; ((k \leq n) \Rightarrow (a(k)(x) < s))\}) \end{array}}}
\vnbprop{fom-sn-swap}{\vnbstmt{\begin{array}{@{}l@{}} \forall d, m, p.; \forall a, b \in \operatorname{vec}(m). \\ \quad \operatorname{is\text{-}seminorm}(d, m, p) \Rightarrow  \\ \quad \quad p(\operatorname{vadd}(m)(a, \operatorname{vneg}(m)(b))) = p(\operatorname{vadd}(m)(b, \operatorname{vneg}(m)(a))) \end{array}}}
\vnbprop{fom-open-of-approx}{\vnbstmt{\begin{array}{@{}l@{}} \forall d, a, b, c, f, t. \\ \quad \operatorname{is\text{-}frechet\text{-}structure}(d, a, b) \wedge \\ \quad \operatorname{is\text{-}seminorm\text{-}family}(d, c, f) \wedge \\ \quad \operatorname{is\text{-}cont\text{-}lin}(a, b, c, f, t) \wedge \\ \quad \forall n \in \mathbb{N}, l. \\ \quad \quad \operatorname{pos\text{-}rr}(l) \Rightarrow  \\ \quad \quad \quad \exists \mathit{fnb} \in \mathbb{N}, \mathit{feta}. \\ \quad \quad \quad \quad \operatorname{pos\text{-}rr}(\mathit{feta}) \wedge \\ \quad \quad \quad \quad \forall z \in \operatorname{vec}(c), k \in \mathbb{N}, h. \\ \quad \quad \quad \quad \quad (\forall i \in \mathbb{N}.\; ((i \leq \mathit{fnb}) \Rightarrow (f(i)(z) < \mathit{feta})) \wedge \operatorname{pos\text{-}rr}(h)) \Rightarrow  \\ \quad \quad \quad \quad \quad \quad \exists x \in \operatorname{vec}(a). \\ \quad \quad \quad \quad \quad \quad \quad \forall j \in \mathbb{N}.\; ((j \leq n) \Rightarrow (b(j)(x) < l)) \wedge \\ \quad \quad \quad \quad \quad \quad \quad \forall i \in \mathbb{N}.\; ((i \leq k) \Rightarrow (f(i)(\operatorname{vadd}(c)(z, \operatorname{vneg}(c)(t(x)))) < h)) \Rightarrow \\ \quad \quad \operatorname{is\text{-}open\text{-}lin\text{-}map}(a, b, c, f, t) \end{array}}}
\vnbprop{open-mapping-theorem}{\vnbstmt{\begin{array}{@{}l@{}} \forall d, a, b, c, f, t. \\ \quad \operatorname{is\text{-}nontrivially\text{-}valued}(d) \wedge \\ \quad \operatorname{is\text{-}frechet\text{-}structure}(d, a, b) \wedge \\ \quad \operatorname{is\text{-}frechet\text{-}structure}(d, c, f) \wedge \\ \quad \operatorname{is\text{-}cont\text{-}lin}(a, b, c, f, t) \wedge \\ \quad \forall y \in \operatorname{vec}(c).\; \exists x \in \operatorname{vec}(a).\; (y = t(x)) \Rightarrow \\ \quad \quad \operatorname{is\text{-}open\text{-}lin\text{-}map}(a, b, c, f, t) \end{array}}}
\vnbprop{open-mapping-bounded-inverse}{\vnbstmt{\begin{array}{@{}l@{}} \forall d, a, b, c, f, t, g. \\ \quad \operatorname{is\text{-}nontrivially\text{-}valued}(d) \wedge \\ \quad \operatorname{is\text{-}frechet\text{-}structure}(d, a, b) \wedge \\ \quad \operatorname{is\text{-}frechet\text{-}structure}(d, c, f) \wedge \\ \quad \operatorname{is\text{-}cont\text{-}lin}(a, b, c, f, t) \wedge \\ \quad \operatorname{is\text{-}linear\text{-}map}(c, a, g) \wedge \\ \quad \forall x \in \operatorname{vec}(a).\; (g(t(x)) = x) \wedge \\ \quad \forall y \in \operatorname{vec}(c).\; (t(g(y)) = y) \Rightarrow \\ \quad \quad \operatorname{is\text{-}cont\text{-}lin}(c, f, a, b, g) \end{array}}}

\vnbbookfile{theorem-library/frechet-banach-steinhaus}{theorem-library/frechet-banach-steinhaus}
\noindent{\small frechet-banach-steinhaus.scm -- Thayer, "Analise Funcional e Teoria Espectral", Ch. 8 s.4, pp. 256-257: Thm 8.29 (closed graph bis) and Thm 8.32 (Banach-Steinhaus bis, exercise 8.10).  Batch TVS-2B, 2026-10-11/12.  Vocabulary (Def. 8.31, the equicontinuity of 8.32): structure-library/frechet-banach-steinhaus.scm.}\par
\vnbprop{fbs-conv-dom}{\vnbstmt{\begin{array}{@{}l@{}} \forall a, b, f, c, s, d, g, h. \\ \quad \operatorname{is\text{-}module}(a) \wedge \\ \quad \operatorname{is\text{-}module}(b) \wedge \\ \quad f \in (\mathbb{N} \to (\operatorname{vec}(a) \to \mathbb{R})) \wedge \\ \quad c \in (\mathbb{N} \to (\operatorname{vec}(b) \to \mathbb{R})) \wedge \\ \quad s \in (\mathbb{N} \to \operatorname{vec}(a)) \wedge \\ \quad d \in (\mathbb{N} \to \operatorname{vec}(b)) \wedge \\ \quad g \in \operatorname{vec}(a) \wedge \\ \quad h \in \operatorname{vec}(b) \wedge \\ \quad \forall k, i. \\ \quad \quad ((k \in \mathbb{N}) \wedge (i \in \mathbb{N})) \Rightarrow  \\ \quad \quad \quad c(k)(\operatorname{vadd}(b)(d(i), \operatorname{vneg}(b)(h))) \leq f(k)(\operatorname{vadd}(a)(s(i), \operatorname{vneg}(a)(g))) \wedge \\ \quad \operatorname{fr\text{-}conv}(a, f, s, g) \Rightarrow \\ \quad \quad \operatorname{fr\text{-}conv}(b, c, d, h) \end{array}}}
\vnbprop{fbs-cauchy-dom}{\vnbstmt{\begin{array}{@{}l@{}} \forall a, b, f, c, s, d. \\ \quad \operatorname{is\text{-}module}(a) \wedge \\ \quad \operatorname{is\text{-}module}(b) \wedge \\ \quad f \in (\mathbb{N} \to (\operatorname{vec}(a) \to \mathbb{R})) \wedge \\ \quad c \in (\mathbb{N} \to (\operatorname{vec}(b) \to \mathbb{R})) \wedge \\ \quad s \in (\mathbb{N} \to \operatorname{vec}(a)) \wedge \\ \quad d \in (\mathbb{N} \to \operatorname{vec}(b)) \wedge \\ \quad \forall k, i, j. \\ \quad \quad ((k \in \mathbb{N}) \wedge ((i \in \mathbb{N}) \wedge (j \in \mathbb{N}))) \Rightarrow  \\ \quad \quad \quad c(k)(\operatorname{vadd}(b)(d(i), \operatorname{vneg}(b)(d(j)))) \leq f(k)(\operatorname{vadd}(a)(s(i), \operatorname{vneg}(a)(s(j)))) \wedge \\ \quad \operatorname{fr\text{-}cauchy}(a, f, s) \Rightarrow \\ \quad \quad \operatorname{fr\text{-}cauchy}(b, c, d) \end{array}}}
\vnbprop{fbs-conv-sum}{\vnbstmt{\begin{array}{@{}l@{}} \forall a, b, f, c, d, s, g, h, m. \\ \quad \operatorname{is\text{-}module}(a) \wedge \\ \quad \operatorname{is\text{-}module}(b) \wedge \\ \quad f \in (\mathbb{N} \to (\operatorname{vec}(a) \to \mathbb{R})) \wedge \\ \quad c \in (\mathbb{N} \to (\operatorname{vec}(a) \to \mathbb{R})) \wedge \\ \quad d \in (\mathbb{N} \to (\operatorname{vec}(b) \to \mathbb{R})) \wedge \\ \quad s \in (\mathbb{N} \to \operatorname{vec}(a)) \wedge \\ \quad g \in (\mathbb{N} \to \operatorname{vec}(b)) \wedge \\ \quad h \in \operatorname{vec}(a) \wedge \\ \quad m \in \operatorname{vec}(b) \wedge \\ \quad \forall k, i. \\ \quad \quad ((k \in \mathbb{N}) \wedge (i \in \mathbb{N})) \Rightarrow  \\ \quad \quad \quad f(k)(\operatorname{vadd}(a)(s(i), \operatorname{vneg}(a)(h))) \\ \quad \quad \quad \leq\; c(k)(\operatorname{vadd}(a)(s(i), \operatorname{vneg}(a)(h))) + d(k)(\operatorname{vadd}(b)(g(i), \operatorname{vneg}(b)(m))) \wedge \\ \quad \operatorname{fr\text{-}conv}(a, c, s, h) \wedge \\ \quad \operatorname{fr\text{-}conv}(b, d, g, m) \Rightarrow \\ \quad \quad \operatorname{fr\text{-}conv}(a, f, s, h) \end{array}}}
\vnbprop{fbs-graph-fam-exists}{\vnbstmt{\begin{array}{@{}l@{}} \forall d, a, b, c, f, t. \\ \quad (\operatorname{is\text{-}seminorm\text{-}family}(d, a, b) \wedge (\operatorname{is\text{-}seminorm\text{-}family}(d, c, f) \wedge (t \in (\operatorname{vec}(a) \to \operatorname{vec}(c))))) \Rightarrow  \\ \quad \quad \exists g \in (\mathbb{N} \to (\operatorname{vec}(a) \to \mathbb{R})). \forall k \in \mathbb{N}, x \in \operatorname{vec}(a). \\ \quad \quad \quad g(k)(x) = (b(k)(x) + f(k)(t(x))) \end{array}}}
\vnbprop{fbs-graph-seminorm}{\vnbstmt{\begin{array}{@{}l@{}} \forall d, a, b, c, f, t, g.; \forall k \in \mathbb{N}. \\ \quad \operatorname{is\text{-}seminorm\text{-}family}(d, a, b) \wedge \\ \quad \operatorname{is\text{-}seminorm\text{-}family}(d, c, f) \wedge \\ \quad \operatorname{is\text{-}linear\text{-}map}(a, c, t) \wedge \\ \quad g \in (\mathbb{N} \to (\operatorname{vec}(a) \to \mathbb{R})) \wedge \\ \quad \forall k \in \mathbb{N}, x \in \operatorname{vec}(a).\; (g(k)(x) = (b(k)(x) + f(k)(t(x)))) \Rightarrow \\ \quad \quad \operatorname{is\text{-}seminorm}(d, a, g(k)) \end{array}}}
\vnbprop{fbs-graph-family}{\vnbstmt{\begin{array}{@{}l@{}} \forall d, a, b, c, f, t, g. \\ \quad \operatorname{is\text{-}seminorm\text{-}family}(d, a, b) \wedge \\ \quad \operatorname{is\text{-}seminorm\text{-}family}(d, c, f) \wedge \\ \quad \operatorname{is\text{-}linear\text{-}map}(a, c, t) \wedge \\ \quad g \in (\mathbb{N} \to (\operatorname{vec}(a) \to \mathbb{R})) \wedge \\ \quad \forall k \in \mathbb{N}, x \in \operatorname{vec}(a).\; (g(k)(x) = (b(k)(x) + f(k)(t(x)))) \Rightarrow \\ \quad \quad \operatorname{is\text{-}seminorm\text{-}family}(d, a, g) \end{array}}}
\vnbprop{fbs-graph-frechet}{\vnbstmt{\begin{array}{@{}l@{}} \forall d, a, b, c, f, t, g. \\ \quad \operatorname{is\text{-}frechet\text{-}structure}(d, a, b) \wedge \\ \quad \operatorname{is\text{-}frechet\text{-}structure}(d, c, f) \wedge \\ \quad \operatorname{has\text{-}closed\text{-}graph}(a, b, c, f, t) \wedge \\ \quad g \in (\mathbb{N} \to (\operatorname{vec}(a) \to \mathbb{R})) \wedge \\ \quad \forall k \in \mathbb{N}, x \in \operatorname{vec}(a).\; (g(k)(x) = (b(k)(x) + f(k)(t(x)))) \Rightarrow \\ \quad \quad \operatorname{is\text{-}frechet\text{-}structure}(d, a, g) \end{array}}}
\vnbprop{fbs-id-linear}{\vnbstmt{\forall a.\; (\operatorname{is\text{-}module}(a) \Rightarrow \operatorname{is\text{-}linear\text{-}map}(a, a, (\lambda v \in \operatorname{vec}(a).\; v)))}}
\vnbprop{closed-graph-theorem}{\vnbstmt{\begin{array}{@{}l@{}} \forall d, a, b, c, f, t. \\ \quad \operatorname{is\text{-}nontrivially\text{-}valued}(d) \wedge \\ \quad \operatorname{is\text{-}frechet\text{-}structure}(d, a, b) \wedge \\ \quad \operatorname{is\text{-}frechet\text{-}structure}(d, c, f) \wedge \\ \quad \operatorname{has\text{-}closed\text{-}graph}(a, b, c, f, t) \Rightarrow \\ \quad \quad \operatorname{is\text{-}cont\text{-}lin}(a, b, c, f, t) \end{array}}}
\vnbprop{fbs-baire-ball}{\vnbstmt{\begin{array}{@{}l@{}} \forall d, a, b, c, f, g, h.; \forall k \in \mathbb{N}. \\ \quad \operatorname{is\text{-}frechet\text{-}structure}(d, a, b) \wedge \\ \quad \operatorname{is\text{-}seminorm\text{-}family}(d, c, f) \wedge \\ \quad h \in (g \to (\operatorname{vec}(a) \to \operatorname{vec}(c))) \wedge \\ \quad \forall x \in \operatorname{vec}(a). \\ \quad \quad \operatorname{is\text{-}fr\text{-}bounded}(c, f, \{y \in \operatorname{vec}(c) : \exists i \in g.\; (y = h(i)(x))\}) \Rightarrow \\ \quad \quad \exists n \in \mathbb{N}, m \in \operatorname{vec}(a), p. \\ \quad \quad \quad \operatorname{pos\text{-}rr}(p) \wedge \\ \quad \quad \quad \operatorname{ball}(\operatorname{fr\text{-}metric\text{-}space}(a, b), m, p) \\ \quad \quad \quad \subseteq\; \operatorname{closure}(\operatorname{fr\text{-}metric\text{-}space}(a, b), \{\mathit{bx} \in \operatorname{vec}(a) : \forall \mathit{bi} \in g.\; (f(k)(h(\mathit{bi})(\mathit{bx})) \leq n)\}) \end{array}}}
\vnbprop{fbs-v-sub-add}{\vnbstmt{\begin{array}{@{}l@{}} \forall m.; \forall x, y \in \operatorname{vec}(m). \\ \quad \operatorname{is\text{-}module}(m) \Rightarrow (\operatorname{vadd}(m)(\operatorname{vadd}(m)(x, \operatorname{vneg}(m)(y)), y) = x) \end{array}}}
\vnbprop{fbs-v-add-sub}{\vnbstmt{\begin{array}{@{}l@{}} \forall m.; \forall x, y \in \operatorname{vec}(m). \\ \quad \operatorname{is\text{-}module}(m) \Rightarrow (\operatorname{vadd}(m)(\operatorname{vadd}(m)(x, y), \operatorname{vneg}(m)(y)) = x) \end{array}}}
\vnbprop{fbs-q-cont}{\vnbstmt{\begin{array}{@{}l@{}} \forall d, a, b, c, f, t, k, y, s. \\ \quad \operatorname{is\text{-}seminorm\text{-}family}(d, a, b) \wedge \\ \quad \operatorname{is\text{-}seminorm\text{-}family}(d, c, f) \wedge \\ \quad \operatorname{is\text{-}cont\text{-}lin}(a, b, c, f, t) \wedge \\ \quad k \in \mathbb{N} \wedge \\ \quad y \in \operatorname{vec}(a) \wedge \\ \quad \operatorname{pos\text{-}rr}(s) \Rightarrow \\ \quad \quad \exists e. \\ \quad \quad \quad \operatorname{pos\text{-}rr}(e) \wedge \\ \quad \quad \quad \forall z \in \operatorname{vec}(a). \\ \quad \quad \quad \quad (\operatorname{dist}(\operatorname{fr\text{-}metric\text{-}space}(a, b))(y, z) \leq e) \Rightarrow  \\ \quad \quad \quad \quad \quad f(k)(t(y)) \leq (f(k)(t(z)) + s) \end{array}}}
\vnbprop{fbs-closure-bound}{\vnbstmt{\begin{array}{@{}l@{}} \forall d, a, b, c, f, x, g, k, n, z.; \forall j \in x. \\ \quad \operatorname{is\text{-}seminorm\text{-}family}(d, a, b) \wedge \\ \quad \operatorname{is\text{-}seminorm\text{-}family}(d, c, f) \wedge \\ \quad g \in (x \to (\operatorname{vec}(a) \to \operatorname{vec}(c))) \wedge \\ \quad \forall i \in x.\; \operatorname{is\text{-}cont\text{-}lin}(a, b, c, f, g(i)) \wedge \\ \quad k \in \mathbb{N} \wedge \\ \quad n \in \mathbb{N} \wedge \\ \quad z \in \operatorname{vec}(a) \wedge \\ \quad z \in \operatorname{closure}(\operatorname{fr\text{-}metric\text{-}space}(a, b), \{\mathit{bx} \in \operatorname{vec}(a) : \forall \mathit{bi} \in x.\; (f(k)(g(\mathit{bi})(\mathit{bx})) \leq n)\}) \Rightarrow \\ \quad \quad f(k)(g(j)(z)) \leq n \end{array}}}
\vnbprop{fbs-ball-bound}{\vnbstmt{\begin{array}{@{}l@{}} \forall d, a, f, g, h, t.; \forall c, b \in \operatorname{carr}(d), k, m, l, n.; \forall x \in \operatorname{vec}(a), r \in \mathbb{R}. \\ \quad \operatorname{is\text{-}seminorm\text{-}family}(d, a, f) \wedge \\ \quad \operatorname{is\text{-}seminorm\text{-}family}(d, g, h) \wedge \\ \quad \operatorname{is\text{-}linear\text{-}map}(a, g, t) \wedge \\ \quad 1 < \operatorname{fnrm}(d)(c) \wedge \\ \quad \operatorname{mul}(d)(c, b) = \operatorname{one}(d) \wedge \\ \quad k \in \mathbb{N} \wedge \\ \quad m \in \mathbb{N} \wedge \\ \quad \operatorname{pos\text{-}rr}(l) \wedge \\ \quad \operatorname{pos\text{-}rr}(n) \wedge \\ \quad \forall x \in \operatorname{vec}(a). \\ \quad \quad \forall j \in \mathbb{N}.\; ((j \leq m) \Rightarrow (f(j)(x) < l)) \Rightarrow (h(k)(t(x)) \leq n) \wedge \\ \quad \forall j \in \mathbb{N}.\; ((j \leq m) \Rightarrow (f(j)(x) \leq r)) \Rightarrow \\ \quad \quad h(k)(t(x)) \leq (((\operatorname{fnrm}(d)(c) \cdot n) \cdot \operatorname{recip}(l)) \cdot r) \end{array}}}
\vnbprop{fbs-uniform-ball}{\vnbstmt{\begin{array}{@{}l@{}} \forall d, a, b, c, f, g, h.; \forall k \in \mathbb{N}. \\ \quad \operatorname{is\text{-}frechet\text{-}structure}(d, a, b) \wedge \\ \quad \operatorname{is\text{-}seminorm\text{-}family}(d, c, f) \wedge \\ \quad h \in (g \to (\operatorname{vec}(a) \to \operatorname{vec}(c))) \wedge \\ \quad \forall i \in g.\; \operatorname{is\text{-}cont\text{-}lin}(a, b, c, f, h(i)) \wedge \\ \quad \forall x \in \operatorname{vec}(a). \\ \quad \quad \operatorname{is\text{-}fr\text{-}bounded}(c, f, \{y \in \operatorname{vec}(c) : \exists i \in g.\; (y = h(i)(x))\}) \Rightarrow \\ \quad \quad \exists m \in \mathbb{N}, l. \\ \quad \quad \quad \operatorname{pos\text{-}rr}(l) \wedge \\ \quad \quad \quad \exists n. \\ \quad \quad \quad \quad \operatorname{pos\text{-}rr}(n) \wedge \\ \quad \quad \quad \quad \forall i \in g, x \in \operatorname{vec}(a). \\ \quad \quad \quad \quad \quad \forall j \in \mathbb{N}.\; ((j \leq m) \Rightarrow (b(j)(x) < l)) \Rightarrow (f(k)(h(i)(x)) \leq n) \end{array}}}
\vnbprop{frechet-banach-steinhaus}{\vnbstmt{\begin{array}{@{}l@{}} \forall d, a, b, c, f, g, h. \\ \quad \operatorname{is\text{-}nontrivially\text{-}valued}(d) \wedge \\ \quad \operatorname{is\text{-}frechet\text{-}structure}(d, a, b) \wedge \\ \quad \operatorname{is\text{-}seminorm\text{-}family}(d, c, f) \wedge \\ \quad h \in (g \to (\operatorname{vec}(a) \to \operatorname{vec}(c))) \wedge \\ \quad \forall i \in g.\; \operatorname{is\text{-}cont\text{-}lin}(a, b, c, f, h(i)) \wedge \\ \quad \forall x \in \operatorname{vec}(a). \\ \quad \quad \operatorname{is\text{-}fr\text{-}bounded}(c, f, \{y \in \operatorname{vec}(c) : \exists i \in g.\; (y = h(i)(x))\}) \Rightarrow \\ \quad \quad \operatorname{is\text{-}fr\text{-}equicont}(a, b, c, f, g, h) \end{array}}}

\vnbbookfile{theorem-library/rake-asm-1}{theorem-library/rake-asm-1}
\noindent{\small rake-asm-1.scm -- RAKE-ASM agent RA-1 (2026-10-12): elementary assumptions PROVEN.}\par
\vnbprop{normed-ag-nrm-type}{\vnbstmt{\forall s.\; (\operatorname{is\text{-}normed\text{-}ag}(s) \Rightarrow (\operatorname{nrm}(s) \in (\operatorname{carr}(s) \to \mathbb{R})))}}
\vnbprop{normed-ag-nrm-nonneg}{\vnbstmt{\forall s.; \forall a \in \operatorname{carr}(s).\; (\operatorname{is\text{-}normed\text{-}ag}(s) \Rightarrow (0 \leq \operatorname{nrm}(s)(a)))}}
\vnbprop{normed-ag-nrm-definite}{\vnbstmt{\begin{array}{@{}l@{}} \forall s.; \forall a \in \operatorname{carr}(s). \\ \quad \operatorname{is\text{-}normed\text{-}ag}(s) \Rightarrow ((\operatorname{nrm}(s)(a) = 0) \Leftrightarrow (a = \operatorname{iden}(s))) \end{array}}}
\vnbprop{normed-ag-nrm-subadditive}{\vnbstmt{\begin{array}{@{}l@{}} \forall s.; \forall a, b \in \operatorname{carr}(s). \\ \quad \operatorname{is\text{-}normed\text{-}ag}(s) \Rightarrow (\operatorname{nrm}(s)(\operatorname{opr}(s)(a, b)) \leq (\operatorname{nrm}(s)(a) + \operatorname{nrm}(s)(b))) \end{array}}}
\vnbprop{rr-archimedean-floor}{\vnbstmt{\begin{array}{@{}l@{}} \forall r.; \forall s \in \mathbb{R}. \\ \quad (\operatorname{pos\text{-}rr}(r) \wedge (0 \leq s)) \Rightarrow \exists w \in \mathbb{N}.\; ((s < (w \cdot r)) \wedge (((w \cdot r) - r) \leq s)) \end{array}}}
\vnbprop{qq-dense-in-rr}{\vnbstmt{\begin{array}{@{}l@{}} \forall s, x. \\ \quad (\operatorname{pos\text{-}rr}(s) \wedge (x \in \mathbb{R})) \Rightarrow \exists a \in \mathbb{Q}.\; (((a - s) < x) \wedge (x < (a + s))) \end{array}}}

\vnbbookfile{theorem-library/rake-asm-b2}{theorem-library/rake-asm-b2}
\noindent{\small rake-asm-b2.scm -- RAKE-ASM-B agent RB-2 (2026-10-12): four assumptions PROVEN.}\par
\vnbprop{totally-bounded-has-cauchy-subseq-rad}{\vnbstmt{\begin{array}{@{}l@{}} \forall s.; \forall f \in (\mathbb{N} \to \operatorname{pts}(s)), d. \\ \quad (\operatorname{totally\text{-}bounded}(s) \wedge \operatorname{null\text{-}rr\text{-}seq}(d)) \Rightarrow  \\ \quad \quad \exists i.\; (\operatorname{strictly\text{-}mono\text{-}nn}(i) \wedge \operatorname{is\text{-}cauchy\text{-}seq}(s, \operatorname{subseq}(f, i))) \end{array}}}
\vnbprop{summable-bound-converges}{\vnbstmt{\begin{array}{@{}l@{}} \forall s, f, d. \\ \quad \operatorname{is\text{-}complete}(s) \wedge \\ \quad f \in (\mathbb{N} \to \operatorname{pts}(s)) \wedge \\ \quad d \in (\mathbb{N} \to \mathbb{R}) \wedge \\ \quad \operatorname{series\text{-}converges}(d) \wedge \\ \quad \forall k \in \mathbb{N}.\; (\operatorname{dist}(s)(f(k), f((k + 1))) \leq d(k)) \Rightarrow \\ \quad \quad \operatorname{converges}(s, f) \end{array}}}
\vnbprop{sums-to-unique}{\vnbstmt{\begin{array}{@{}l@{}} \forall p.; \forall f \in (\operatorname{dom}(f) \to \operatorname{carr}(p)), a, b \in \operatorname{carr}(p). \\ \quad (\operatorname{is\text{-}normed\text{-}ag}(p) \wedge (\operatorname{sums\text{-}to}(p, f, a) \wedge \operatorname{sums\text{-}to}(p, f, b))) \Rightarrow (a = b) \end{array}}}
\vnbprop{rb2-nrm-finsum-le}{\vnbstmt{\begin{array}{@{}l@{}} \forall p, f, g, d.; \forall w \in \mathbf{Set}. \\ \quad \operatorname{is\text{-}normed\text{-}ag}(p) \wedge \\ \quad \forall z \in d.\; (f(z) \in \operatorname{carr}(p)) \wedge \\ \quad \forall z \in d.\; (g(z) = \operatorname{nrm}(p)(f(z))) \wedge \\ \quad |w| \in \mathbb{N} \wedge \\ \quad w \subseteq d \Rightarrow \\ \quad \quad \operatorname{finsum}(\operatorname{rr\text{-}pos\text{-}star\text{-}add\text{-}monoid}, g, w) \in \mathbb{R} \wedge \\ \quad \quad \operatorname{nrm}(p)(\operatorname{finsum}(\operatorname{normed\text{-}ag\text{-}as\text{-}abelian\text{-}group}(p), f, w)) \leq \operatorname{finsum}(\operatorname{rr\text{-}pos\text{-}star\text{-}add\text{-}monoid}, g, w) \end{array}}}
\vnbprop{rb2-dist-union}{\vnbstmt{\begin{array}{@{}l@{}} \forall p, f, g, d.; \forall a, c \in \mathbf{Set}. \\ \quad \operatorname{is\text{-}normed\text{-}ag}(p) \wedge \\ \quad \forall z \in d.\; (f(z) \in \operatorname{carr}(p)) \wedge \\ \quad \forall z \in d.\; (g(z) = \operatorname{nrm}(p)(f(z))) \wedge \\ \quad |a| \in \mathbb{N} \wedge \\ \quad |c| \in \mathbb{N} \wedge \\ \quad \forall z \in c.\; \neg (z \in a) \wedge \\ \quad a \subseteq d \wedge \\ \quad c \subseteq d \Rightarrow \\ \quad \quad \operatorname{finsum}(\operatorname{rr\text{-}pos\text{-}star\text{-}add\text{-}monoid}, g, (a \cup c)) \\ \quad \quad =\; \operatorname{finsum}(\operatorname{rr\text{-}pos\text{-}star\text{-}add\text{-}monoid}, g, a) + \operatorname{finsum}(\operatorname{rr\text{-}pos\text{-}star\text{-}add\text{-}monoid}, g, c) \wedge \\ \quad \quad \operatorname{dist}(\operatorname{nag\text{-}metric\text{-}space}(p))(\operatorname{finsum}(\operatorname{normed\text{-}ag\text{-}as\text{-}abelian\text{-}group}(p), f, a), \\ \quad \quad \quad \operatorname{finsum}(\operatorname{normed\text{-}ag\text{-}as\text{-}abelian\text{-}group}(p), f, (a \cup c))) \\ \quad \quad \leq\; \operatorname{finsum}(\operatorname{rr\text{-}pos\text{-}star\text{-}add\text{-}monoid}, g, c) \end{array}}}
\vnbprop{rb2-dist-subset}{\vnbstmt{\begin{array}{@{}l@{}} \forall p, f, g, d.; \forall a, b \in \mathbf{Set}. \\ \quad \operatorname{is\text{-}normed\text{-}ag}(p) \wedge \\ \quad \forall z \in d.\; (f(z) \in \operatorname{carr}(p)) \wedge \\ \quad \forall z \in d.\; (g(z) = \operatorname{nrm}(p)(f(z))) \wedge \\ \quad |a| \in \mathbb{N} \wedge \\ \quad |b| \in \mathbb{N} \wedge \\ \quad a \subseteq b \wedge \\ \quad b \subseteq d \Rightarrow \\ \quad \quad \operatorname{dist}(\operatorname{nag\text{-}metric\text{-}space}(p))(\operatorname{finsum}(\operatorname{normed\text{-}ag\text{-}as\text{-}abelian\text{-}group}(p), f, a), \\ \quad \quad \quad \operatorname{finsum}(\operatorname{normed\text{-}ag\text{-}as\text{-}abelian\text{-}group}(p), f, b)) \\ \quad \quad \leq\; (\operatorname{finsum}(\operatorname{rr\text{-}pos\text{-}star\text{-}add\text{-}monoid}, g, b) - \operatorname{finsum}(\operatorname{rr\text{-}pos\text{-}star\text{-}add\text{-}monoid}, g, a)) \end{array}}}
\vnbprop{rb2-esum-approx}{\vnbstmt{\begin{array}{@{}l@{}} \forall g \in (\operatorname{dom}(g) \to \operatorname{rr\text{-}pos\text{-}star}), l. \\ \quad ((\operatorname{esum}(g) \in \mathbb{R}) \wedge \operatorname{pos\text{-}rr}(l)) \Rightarrow  \\ \quad \quad \exists t \in \mathbf{Set}. \\ \quad \quad \quad ((|t| \in \mathbb{N}) \wedge ((t \subseteq \operatorname{dom}(g)) \wedge (\operatorname{esum}(g) < (\operatorname{finsum}(\operatorname{rr\text{-}pos\text{-}star\text{-}add\text{-}monoid}, g, t) + l)))) \end{array}}}
\vnbprop{rb2-choice-seq}{\vnbstmt{\begin{array}{@{}l@{}} \forall g \in (\operatorname{dom}(g) \to \operatorname{rr\text{-}pos\text{-}star}), d. \\ \quad ((\operatorname{esum}(g) \in \mathbb{R}) \wedge \forall k \in \mathbb{N}.\; \operatorname{pos\text{-}rr}(d(k))) \Rightarrow  \\ \quad \quad \exists s. \forall n \in \mathbb{N}. \\ \quad \quad \quad s(n) \in \mathbf{Set} \wedge \\ \quad \quad \quad |s(n)| \in \mathbb{N} \wedge \\ \quad \quad \quad s(n) \subseteq \operatorname{dom}(g) \wedge \\ \quad \quad \quad \operatorname{esum}(g) < (\operatorname{finsum}(\operatorname{rr\text{-}pos\text{-}star\text{-}add\text{-}monoid}, g, s(n)) + d(n)) \end{array}}}
\vnbprop{absolute-summable-implies-summable}{\vnbstmt{\begin{array}{@{}l@{}} \forall p.; \forall f \in (\operatorname{dom}(f) \to \operatorname{carr}(p)). \\ \quad ((\operatorname{is\text{-}normed\text{-}ag}(p) \wedge \operatorname{is\text{-}complete}(\operatorname{nag\text{-}metric\text{-}space}(p))) \wedge \operatorname{is\text{-}absolutely\text{-}summable}(p, f)) \Rightarrow  \\ \quad \quad \operatorname{is\text{-}summable}(p, f) \end{array}}}

\vnbbookfile{theorem-library/primitive-glue}{theorem-library/primitive-glue}
\noindent{\small primitive-glue.scm -- PRIMITIVES ON ADJACENT INTERVALS GLUE (batch 22-C).}\par
\vnbprop{clamp-left-const}{\vnbstmt{\begin{array}{@{}l@{}} \forall a, b, x \in \mathbb{R}. \\ \quad ((x \leq a) \wedge (a \leq b)) \Rightarrow (\operatorname{clamp}(a, b, x) = a) \end{array}}}
\vnbprop{clamp-right-const}{\vnbstmt{\begin{array}{@{}l@{}} \forall a, b, x \in \mathbb{R}. \\ \quad ((a \leq b) \wedge (b \leq x)) \Rightarrow (\operatorname{clamp}(a, b, x) = b) \end{array}}}
\vnbprop{primitive-integrand-interior-congruence}{\vnbstmt{\begin{array}{@{}l@{}} \forall f, i, \mathit{ppsi}, a, b. \\ \quad \operatorname{is\text{-}primitive}(f, i, a, b) \wedge \\ \quad \mathit{ppsi} \in (\operatorname{ccint}(a, b) \to \mathbb{R}) \wedge \\ \quad \forall y \in \operatorname{ooint}(a, b).\; (i(y) \simeq \operatorname{ppsi}(y)) \Rightarrow \\ \quad \quad \operatorname{is\text{-}primitive}(f, \mathit{ppsi}, a, b) \end{array}}}
\vnbprop{primitive-glue}{\vnbstmt{\begin{array}{@{}l@{}} \forall a, c, b \in \mathbb{R}, d, h, k, m.; \forall g, f \in (\operatorname{ccint}(a, b) \to \mathbb{R}). \\ \quad a < c \wedge \\ \quad c < b \wedge \\ \quad \operatorname{is\text{-}primitive}(d, h, a, c) \wedge \\ \quad \operatorname{is\text{-}primitive}(k, m, c, b) \wedge \\ \quad d(c) = k(c) \wedge \\ \quad \forall t \in \operatorname{ccint}(a, c).\; (g(t) \simeq d(t)) \wedge \\ \quad \forall t \in \operatorname{ccint}(c, b).\; (g(t) \simeq k(t)) \wedge \\ \quad \forall t \in \operatorname{ooint}(a, c).\; (f(t) \simeq h(t)) \wedge \\ \quad \forall t \in \operatorname{ooint}(c, b).\; (f(t) \simeq m(t)) \Rightarrow \\ \quad \quad \operatorname{is\text{-}primitive}(g, f, a, b) \end{array}}}
\vnbprop{primitive-of-constant-piece}{\vnbstmt{\begin{array}{@{}l@{}} \forall u, v \in \mathbb{R}, f \in (\mathbb{R} \to \mathbb{R}), a, k \in \mathbb{R}. \\ \quad ((u < v) \wedge \forall t \in \mathbb{R}.\; (((u < t) \wedge (t < v)) \Rightarrow (f(t) = a))) \Rightarrow  \\ \quad \quad \operatorname{is\text{-}primitive}((\lambda s \in \operatorname{ccint}(u, v).\; (k + (a \cdot (s - u)))), \operatorname{restrict}(f, \operatorname{ccint}(u, v)), u, v) \end{array}}}
\vnbprop{step-fn-has-primitive-ind}{\vnbstmt{\begin{array}{@{}l@{}} \forall n \in \mathbb{N}, p, a, b.; \forall f \in (\mathbb{R} \to \mathbb{R}). \\ \quad \operatorname{is\text{-}partition}(p, n, a, b) \wedge \\ \quad \forall i \in \operatorname{interval}(0, n). \\ \quad \quad (i < n) \Rightarrow  \\ \quad \quad \quad \exists c \in \mathbb{R}. \forall t \in \mathbb{R}. \\ \quad \quad \quad \quad ((p(i) < t) \wedge (t < p((i + 1)))) \Rightarrow (f(t) = c) \Rightarrow \\ \quad \quad \exists g.\; \operatorname{is\text{-}primitive}(g, \operatorname{restrict}(f, \operatorname{ccint}(a, b)), a, b) \end{array}}}
\vnbprop{step-fn-has-primitive}{\vnbstmt{\begin{array}{@{}l@{}} \forall a, b \in \mathbb{R}, f. \\ \quad \operatorname{is\text{-}step\text{-}fn}(f, a, b) \Rightarrow  \\ \quad \quad \exists g.\; \operatorname{is\text{-}primitive}(g, \operatorname{restrict}(f, \operatorname{ccint}(a, b)), a, b) \end{array}}}

\vnbbookfile{theorem-library/pw-int-laws-3}{theorem-library/pw-int-laws-3}
\noindent{\small pw-int-laws-3.scm -- LINEARITY over the reals as ONE law, the COMPLEX scalar, and the additivity of the complex integral over adjacent intervals.  Continues theorem-library/pw-int-laws-2.scm.}\par
\vnbprop{pw-antiderivative-linear}{\vnbstmt{\begin{array}{@{}l@{}} \forall a, b, c, d \in \mathbb{R}, f, w, i, \mathit{ppsi}.; \forall h, \mathit{pchi} \in (\operatorname{ccint}(a, b) \to \mathbb{R}). \\ \quad \operatorname{is\text{-}primitive}(f, i, a, b) \wedge \\ \quad \operatorname{is\text{-}primitive}(w, \mathit{ppsi}, a, b) \wedge \\ \quad \forall y \in \operatorname{ccint}(a, b).\; (h(y) \simeq ((c \cdot f(y)) + (d \cdot w(y)))) \wedge \\ \quad \forall y \in \operatorname{ccint}(a, b).\; (\operatorname{pchi}(y) \simeq ((c \cdot i(y)) + (d \cdot \operatorname{ppsi}(y)))) \Rightarrow \\ \quad \quad (\operatorname{is\text{-}primitive}(h, \mathit{pchi}, a, b) \wedge (\int_{a}^{b} \mathit{pchi} = ((c \cdot \int_{a}^{b} i) + (d \cdot \int_{a}^{b} \mathit{ppsi})))) \end{array}}}
\vnbprop{cc-int-complex-mul}{\vnbstmt{\begin{array}{@{}l@{}} \forall a, b, i, \mathit{pchi}.; \forall w \in \mathbb{C}, f, \mathit{paw}. \\ \quad a \in \mathbb{R} \wedge \\ \quad b \in \mathbb{R} \wedge \\ \quad i \in (\operatorname{ccint}(a, b) \to \mathbb{C}) \wedge \\ \quad \mathit{pchi} \in (\operatorname{ccint}(a, b) \to \mathbb{C}) \wedge \\ \quad \forall y \in \operatorname{ccint}(a, b).\; (\operatorname{pchi}(y) \simeq (w \cdot i(y))) \wedge \\ \quad \operatorname{is\text{-}primitive}(f, (\lambda t \in \operatorname{ccint}(a, b).\; \operatorname{real\text{-}part}(i(t))), a, b) \wedge \\ \quad \operatorname{is\text{-}primitive}(\mathit{paw}, (\lambda t \in \operatorname{ccint}(a, b).\; \operatorname{imag\text{-}part}(i(t))), a, b) \Rightarrow \\ \quad \quad \int_{a}^{b} \mathit{pchi} = (w \cdot \int_{a}^{b} i) \end{array}}}
\vnbprop{cc-int-adjacent}{\vnbstmt{\begin{array}{@{}l@{}} \forall a, b.; \forall v \in \operatorname{ooint}(a, b), i, f, w. \\ \quad a \in \mathbb{R} \wedge \\ \quad b \in \mathbb{R} \wedge \\ \quad i \in (\operatorname{ccint}(a, b) \to \mathbb{C}) \wedge \\ \quad \operatorname{is\text{-}primitive}(f, (\lambda t \in \operatorname{ccint}(a, b).\; \operatorname{real\text{-}part}(i(t))), a, b) \wedge \\ \quad \operatorname{is\text{-}primitive}(w, (\lambda t \in \operatorname{ccint}(a, b).\; \operatorname{imag\text{-}part}(i(t))), a, b) \Rightarrow \\ \quad \quad \int_{a}^{b} i = (\int_{a}^{v} \operatorname{restrict}(i, \operatorname{ccint}(a, v)) + \int_{v}^{b} \operatorname{restrict}(i, \operatorname{ccint}(v, b))) \end{array}}}

\vnbbookfile{theorem-library/primitive-uniform-limit}{theorem-library/primitive-uniform-limit}
\noindent{\small theorem-library/primitive-uniform-limit.scm -- Dieudonne (8.6.4) on a compact interval [a,b], with COUNTABLE exceptional sets: A UNIFORM LIMIT OF PRIMITIVES IS A PRIMITIVE.  Batch 22-B, 2026-09-23.}\par
\vnbprop{countable-union-nn}{\vnbstmt{\begin{array}{@{}l@{}} \forall d.; \forall u \in \mathbf{Set}. \\ \quad (\forall k \in \mathbb{N}.\; \operatorname{is\text{-}countable}(d(k)) \wedge \forall y \in u.\; \exists k \in \mathbb{N}.\; (y \in d(k))) \Rightarrow  \\ \quad \quad \operatorname{is\text{-}countable}(u) \end{array}}}
\vnbprop{primitive-pair-lipschitz}{\vnbstmt{\begin{array}{@{}l@{}} \forall a, b, g_{1}, f_{1}, g_{2}, f_{2}.; \forall m \in \mathbb{R}, x, y \in \operatorname{ccint}(a, b). \\ \quad \operatorname{is\text{-}primitive}(g_{1}, f_{1}, a, b) \wedge \\ \quad \operatorname{is\text{-}primitive}(g_{2}, f_{2}, a, b) \wedge \\ \quad \forall x \in \operatorname{ccint}(a, b).\; (\lvert (\operatorname{f1}(x) - \operatorname{f2}(x)) \rvert \leq m) \wedge \\ \quad x < y \Rightarrow \\ \quad \quad \lvert ((\operatorname{g1}(y) - \operatorname{g2}(y)) - (\operatorname{g1}(x) - \operatorname{g2}(x))) \rvert \leq (m \cdot (y - x)) \end{array}}}
\vnbprop{rr-mul-le-cancel-pos}{\vnbstmt{\forall a, c, m \in \mathbb{R}.\; (((0 < c) \wedge ((a \cdot c) \leq (m \cdot c))) \Rightarrow (a \leq m))}}
\vnbprop{primitive-pair-lipschitz-abs}{\vnbstmt{\begin{array}{@{}l@{}} \forall a, b, g_{1}, f_{1}, g_{2}, f_{2}.; \forall m \in \mathbb{R}, x, y \in \operatorname{ccint}(a, b). \\ \quad \operatorname{is\text{-}primitive}(g_{1}, f_{1}, a, b) \wedge \\ \quad \operatorname{is\text{-}primitive}(g_{2}, f_{2}, a, b) \wedge \\ \quad \forall x \in \operatorname{ccint}(a, b).\; (\lvert (\operatorname{f1}(x) - \operatorname{f2}(x)) \rvert \leq m) \Rightarrow \\ \quad \quad \lvert ((\operatorname{g1}(y) - \operatorname{g2}(y)) - (\operatorname{g1}(x) - \operatorname{g2}(x))) \rvert \leq (m \cdot \lvert (y - x) \rvert) \end{array}}}
\vnbprop{primitive-family-limit}{\vnbstmt{\begin{array}{@{}l@{}} \forall a, b \in \mathbb{R}, m, c \in (\mathbb{N} \to (\operatorname{ccint}(a, b) \to \mathbb{R})). \\ \quad a < b \wedge \\ \quad \forall k \in \mathbb{N}.\; \operatorname{is\text{-}primitive}(m(k), c(k), a, b) \wedge \\ \quad \forall s. \\ \quad \quad \operatorname{pos\text{-}rr}(s) \Rightarrow  \\ \quad \quad \quad \exists n \in \mathbb{N}. \forall k, l \in \mathbb{N}, x \in \operatorname{ccint}(a, b). \\ \quad \quad \quad \quad ((n \leq k) \wedge (n \leq l)) \Rightarrow (\lvert (c(k)(x) - c(l)(x)) \rvert \leq s) \wedge \\ \quad \forall k \in \mathbb{N}.\; (m(k)(a) = 0) \Rightarrow \\ \quad \quad \exists \mathit{pufc} \in (\mathbb{R} \to \mathbb{R}). \\ \quad \quad \quad \forall x \in \operatorname{ccint}(a, b). \\ \quad \quad \quad \quad \operatorname{converges\text{-}to}(\operatorname{rr\text{-}ms}, (\lambda k \in \mathbb{N}.\; m(k)(x)), \operatorname{pufc}(x)) \wedge \\ \quad \quad \quad \forall t \in \mathbb{R}.\; \operatorname{is\text{-}continuous\text{-}at}(\operatorname{rr\text{-}ms}, \operatorname{rr\text{-}ms}, \mathit{pufc}, t) \end{array}}}
\vnbprop{primitive-family-deriv-at}{\vnbstmt{\begin{array}{@{}l@{}} \forall a, b \in \mathbb{R}, m, c \in (\mathbb{N} \to (\operatorname{ccint}(a, b) \to \mathbb{R})), \mathit{pufc} \in (\mathbb{R} \to \mathbb{R}), t \in \operatorname{ooint}(a, b), v \in \mathbb{R}. \\ \quad a < b \wedge \\ \quad \forall k \in \mathbb{N}.\; \operatorname{is\text{-}primitive}(m(k), c(k), a, b) \wedge \\ \quad \forall s. \\ \quad \quad \operatorname{pos\text{-}rr}(s) \Rightarrow  \\ \quad \quad \quad \exists n \in \mathbb{N}. \forall k, l \in \mathbb{N}, x \in \operatorname{ccint}(a, b). \\ \quad \quad \quad \quad ((n \leq k) \wedge (n \leq l)) \Rightarrow (\lvert (c(k)(x) - c(l)(x)) \rvert \leq s) \wedge \\ \quad \forall x \in \operatorname{ccint}(a, b). \\ \quad \quad \operatorname{converges\text{-}to}(\operatorname{rr\text{-}ms}, (\lambda k \in \mathbb{N}.\; m(k)(x)), \operatorname{pufc}(x)) \wedge \\ \quad \forall k \in \mathbb{N}.\; \operatorname{has\text{-}deriv\text{-}at}(m(k), t, c(k)(t)) \wedge \\ \quad \operatorname{converges\text{-}to}(\operatorname{rr\text{-}ms}, (\lambda k \in \mathbb{N}.\; c(k)(t)), v) \Rightarrow \\ \quad \quad \operatorname{is\text{-}diff\text{-}at}(\mathit{pufc}, t, v) \end{array}}}
\vnbprop{primitive-uniform-limit}{\vnbstmt{\begin{array}{@{}l@{}} \forall a, b \in \mathbb{R}, m, c \in (\mathbb{N} \to (\operatorname{ccint}(a, b) \to \mathbb{R})), f \in (\operatorname{ccint}(a, b) \to \mathbb{R}). \\ \quad a < b \wedge \\ \quad \forall k \in \mathbb{N}.\; \operatorname{is\text{-}primitive}(m(k), c(k), a, b) \wedge \\ \quad \forall s. \\ \quad \quad \operatorname{pos\text{-}rr}(s) \Rightarrow  \\ \quad \quad \quad \exists n \in \mathbb{N}. \forall k \in \mathbb{N}, x \in \operatorname{ccint}(a, b). \\ \quad \quad \quad \quad (n \leq k) \Rightarrow (\lvert (c(k)(x) - f(x)) \rvert \leq s) \wedge \\ \quad \forall k \in \mathbb{N}.\; (m(k)(a) = 0) \Rightarrow \\ \quad \quad \exists g. \\ \quad \quad \quad \operatorname{is\text{-}primitive}(g, f, a, b) \wedge \\ \quad \quad \quad \forall x \in \operatorname{ccint}(a, b). \\ \quad \quad \quad \quad \operatorname{converges\text{-}to}(\operatorname{rr\text{-}ms}, (\lambda k \in \mathbb{N}.\; m(k)(x)), g(x)) \end{array}}}

\vnbbookfile{structure-library/hilbert-space}{structure-library/hilbert-space}
\noindent{\small hilbert-space.scm -- HILBERT SPACES: the metric space of a complex inner product space, completeness, and the orthogonal complement.  Vocabulary only: no proof, no support.  Agent HILBERT-1, 2026-10-10 (scratchpad/triage/HILBERT-1.md).  The proofs are}\par
\par\medskip\noindent\textbf{Definition} (\texttt{perp}).\;${\mathit{ks}}^{\perp} :=$
\begin{equation*}
\{\operatorname{prp\_} \in \operatorname{vec}(v) : \forall \operatorname{prk\_} \in \mathit{ks}.\; (\langle\cdot,\cdot\rangle_{v}(\operatorname{prp\_}, \operatorname{prk\_}) = 0)\}
\end{equation*}

\par\medskip\noindent\textbf{Definition} (\texttt{ip-metric-space}).\;$\operatorname{ip\text{-}metric\text{-}space}(v) :=$ $\operatorname{nag\text{-}metric\text{-}space}(\operatorname{ip\text{-}normed\text{-}ag}(v))$
\par\medskip\noindent\textbf{Definition} (\texttt{is-hilbert-space}).\;For every $v$, 
\begin{equation*}
\operatorname{is\text{-}hilbert\text{-}space}(v) \Leftrightarrow (\operatorname{is\text{-}complex\text{-}inner\text{-}product\text{-}space}(v) \wedge \operatorname{is\text{-}complete}(\operatorname{ip\text{-}metric\text{-}space}(v)))
\end{equation*}

\par\medskip\noindent\textbf{Definition} (\texttt{is-orthogonal-projection}).\;For all $v$, $\mathit{ks}$ and $\mathit{pk}$, 
\begin{equation*}
\begin{array}{@{}l@{}} \operatorname{is\text{-}orthogonal\text{-}projection}(v, \mathit{ks}, \mathit{pk}) \\ \Leftrightarrow\; \mathit{pk} \in (\operatorname{vec}(v) \to \operatorname{vec}(v)) \wedge \\ \forall \operatorname{opz\_} \in \operatorname{vec}(v). \\ \quad (\operatorname{pk}(\operatorname{opz\_}) \in \mathit{ks}) \wedge (\operatorname{vadd}(v)(\operatorname{opz\_}, \operatorname{vneg}(v)(\operatorname{pk}(\operatorname{opz\_}))) \in {\mathit{ks}}^{\perp}) \end{array}
\end{equation*}

\par\medskip\noindent\textbf{Definition} (\texttt{is-cc-linear-functional}).\;For all $v$ and $f$, 
\begin{equation*}
\begin{array}{@{}l@{}} \operatorname{is\text{-}cc\text{-}linear\text{-}functional}(v, f) \\ \Leftrightarrow\; f \in (\operatorname{vec}(v) \to \mathbb{C}) \wedge \\ \forall \operatorname{hlx\_}, \operatorname{hly\_} \in \operatorname{vec}(v). \\ \quad f(\operatorname{vadd}(v)(\operatorname{hlx\_}, \operatorname{hly\_})) = +(f(\operatorname{hlx\_}), f(\operatorname{hly\_})) \wedge \\ \forall \operatorname{hlc\_} \in \mathbb{C}, \operatorname{hlx\_} \in \operatorname{vec}(v).\; (f(\operatorname{act}(v)(\operatorname{hlc\_}, \operatorname{hlx\_})) = \cdot(\operatorname{hlc\_}, f(\operatorname{hlx\_}))) \end{array}
\end{equation*}


\vnbbookfile{theorem-library/hilbert-projection}{theorem-library/hilbert-projection}
\noindent{\small hilbert-projection.scm -- HILBERT SPACES: the parallelogram law, K-perp, the orthogonal projection theorem (Prop. 2.19), the Riesz representation theorem (Thm. 2.22) and the density criterion (Cor. 2.21).  All PROVEN modulo 0. Agent HILBERT-1, 2026-10-10 (scratchpad/triage/HILBERT-1.md).  Vocabulary:}\par
\vnbprop{ip-metric-space-unfold}{\vnbstmt{\forall v.\; (\operatorname{ip\text{-}metric\text{-}space}(v) \simeq \operatorname{nag\text{-}metric\text{-}space}(\operatorname{ip\text{-}normed\text{-}ag}(v)))}}
\vnbprop{perp-unfold}{\vnbstmt{\forall v, s.\; ({s}^{\perp} \simeq \{p \in \operatorname{vec}(v) : \forall k \in s.\; (\langle\cdot,\cdot\rangle_{v}(p, k) = 0)\})}}
\vnbprop{perp-membership}{\vnbstmt{\begin{array}{@{}l@{}} \forall v, s, z. \\ \quad (z \in {s}^{\perp}) \Leftrightarrow ((z \in \operatorname{vec}(v)) \wedge \forall k \in s.\; (\langle\cdot,\cdot\rangle_{v}(z, k) = 0)) \end{array}}}
\vnbprop{hilbert-space-is-cips}{\vnbstmt{\forall v.\; (\operatorname{is\text{-}hilbert\text{-}space}(v) \Rightarrow \operatorname{is\text{-}complex\text{-}inner\text{-}product\text{-}space}(v))}}
\vnbprop{hilbert-space-complete}{\vnbstmt{\forall v.\; (\operatorname{is\text{-}hilbert\text{-}space}(v) \Rightarrow \operatorname{is\text{-}complete}(\operatorname{ip\text{-}metric\text{-}space}(v)))}}
\vnbprop{ip-metric-space-is-metric-space}{\vnbstmt{\begin{array}{@{}l@{}} \forall v. \\ \quad \operatorname{is\text{-}complex\text{-}inner\text{-}product\text{-}space}(v) \Rightarrow \operatorname{is\text{-}metric\text{-}space}(\operatorname{ip\text{-}metric\text{-}space}(v)) \end{array}}}
\vnbprop{ip-metric-space-pts}{\vnbstmt{\begin{array}{@{}l@{}} \forall v. \\ \quad \operatorname{is\text{-}complex\text{-}inner\text{-}product\text{-}space}(v) \Rightarrow (\operatorname{pts}(\operatorname{ip\text{-}metric\text{-}space}(v)) \simeq \operatorname{vec}(v)) \end{array}}}
\vnbprop{ip-normed-ag-opr}{\vnbstmt{\begin{array}{@{}l@{}} \forall v. \\ \quad \operatorname{is\text{-}complex\text{-}inner\text{-}product\text{-}space}(v) \Rightarrow (\operatorname{opr}(\operatorname{ip\text{-}normed\text{-}ag}(v)) \simeq \operatorname{vadd}(v)) \end{array}}}
\vnbprop{ip-normed-ag-inv}{\vnbstmt{\begin{array}{@{}l@{}} \forall v. \\ \quad \operatorname{is\text{-}complex\text{-}inner\text{-}product\text{-}space}(v) \Rightarrow (\operatorname{inv}(\operatorname{ip\text{-}normed\text{-}ag}(v)) \simeq \operatorname{vneg}(v)) \end{array}}}
\vnbprop{ip-metric-space-dist}{\vnbstmt{\begin{array}{@{}l@{}} \forall v.; \forall x, y \in \operatorname{vec}(v). \\ \quad \operatorname{is\text{-}complex\text{-}inner\text{-}product\text{-}space}(v) \Rightarrow  \\ \quad \quad \operatorname{dist}(\operatorname{ip\text{-}metric\text{-}space}(v))(x, y) \\ \quad \quad =\; \sqrt{\langle\cdot,\cdot\rangle_{v}(\operatorname{vadd}(v)(x, \operatorname{vneg}(v)(y)), \operatorname{vadd}(v)(x, \operatorname{vneg}(v)(y)))} \end{array}}}
\vnbprop{cips-ip-neg-right}{\vnbstmt{\begin{array}{@{}l@{}} \forall v.; \forall x, y \in \operatorname{vec}(v). \\ \quad \operatorname{is\text{-}complex\text{-}inner\text{-}product\text{-}space}(v) \Rightarrow  \\ \quad \quad \langle\cdot,\cdot\rangle_{v}(x, \operatorname{vneg}(v)(y)) = -\langle\cdot,\cdot\rangle_{v}(x, y) \end{array}}}
\vnbprop{cips-parallelogram-ip}{\vnbstmt{\begin{array}{@{}l@{}} \forall v.; \forall x, y \in \operatorname{vec}(v). \\ \quad \operatorname{is\text{-}complex\text{-}inner\text{-}product\text{-}space}(v) \Rightarrow  \\ \quad \quad \langle\cdot,\cdot\rangle_{v}(\operatorname{vadd}(v)(x, y), \operatorname{vadd}(v)(x, y)) + \langle\cdot,\cdot\rangle_{v}(\operatorname{vadd}(v)(x, \operatorname{vneg}(v)(y)), \operatorname{vadd}(v)(x, \operatorname{vneg}(v)(y))) \\ \quad \quad =\; (2 \cdot \langle\cdot,\cdot\rangle_{v}(x, x)) + (2 \cdot \langle\cdot,\cdot\rangle_{v}(y, y)) \end{array}}}
\vnbprop{cips-parallelogram}{\vnbstmt{\begin{array}{@{}l@{}} \forall v, x, y. \\ \quad (\operatorname{is\text{-}complex\text{-}inner\text{-}product\text{-}space}(v) \wedge ((x \in \operatorname{vec}(v)) \wedge (y \in \operatorname{vec}(v)))) \Rightarrow  \\ \quad \quad (\lVert \operatorname{vadd}(v)(x, y) \rVert \cdot \lVert \operatorname{vadd}(v)(x, y) \rVert) + (\lVert \operatorname{vadd}(v)(x, \operatorname{vneg}(v)(y)) \rVert \cdot \lVert \operatorname{vadd}(v)(x, \operatorname{vneg}(v)(y)) \rVert) \\ \quad \quad =\; (2 \cdot (\lVert x \rVert \cdot \lVert x \rVert)) + (2 \cdot (\lVert y \rVert \cdot \lVert y \rVert)) \end{array}}}
\vnbprop{cips-pythagoras}{\vnbstmt{\begin{array}{@{}l@{}} \forall v, x, y. \\ \quad (\operatorname{is\text{-}complex\text{-}inner\text{-}product\text{-}space}(v) \wedge ((x \in \operatorname{vec}(v)) \wedge ((y \in \operatorname{vec}(v)) \wedge (\langle\cdot,\cdot\rangle_{v}(x, y) = 0)))) \Rightarrow  \\ \quad \quad \lVert \operatorname{vadd}(v)(x, y) \rVert \cdot \lVert \operatorname{vadd}(v)(x, y) \rVert \\ \quad \quad =\; (\lVert x \rVert \cdot \lVert x \rVert) + (\lVert y \rVert \cdot \lVert y \rVert) \end{array}}}
\vnbprop{cc-zero-of-self-conj-zero}{\vnbstmt{\forall a \in \mathbb{C}.\; (((a \cdot \operatorname{conjugate}(a)) = 0) \Rightarrow (a = 0))}}
\vnbprop{cc-magnitude-sqrt-self-conj}{\vnbstmt{\forall z \in \mathbb{C}.\; (\operatorname{magnitude}(z) = \sqrt{(z \cdot \operatorname{conjugate}(z))})}}
\vnbprop{ip-orth-closure}{\vnbstmt{\begin{array}{@{}l@{}} \forall v, a, w, x. \\ \quad \operatorname{is\text{-}complex\text{-}inner\text{-}product\text{-}space}(v) \wedge \\ \quad w \in \operatorname{vec}(v) \wedge \\ \quad x \in \operatorname{closure}(\operatorname{ip\text{-}metric\text{-}space}(v), a) \wedge \\ \quad \forall y \in a.\; (\langle\cdot,\cdot\rangle_{v}(y, w) = 0) \Rightarrow \\ \quad \quad \langle\cdot,\cdot\rangle_{v}(x, w) = 0 \end{array}}}
\vnbprop{perp-orth}{\vnbstmt{\forall v, s, z, k.\; (((z \in {s}^{\perp}) \wedge (k \in s)) \Rightarrow (\langle\cdot,\cdot\rangle_{v}(z, k) = 0))}}
\vnbprop{perp-subset}{\vnbstmt{\begin{array}{@{}l@{}} \forall v, s. \\ \quad \operatorname{is\text{-}complex\text{-}inner\text{-}product\text{-}space}(v) \Rightarrow ({s}^{\perp} \subseteq \operatorname{vec}(v)) \end{array}}}
\vnbprop{closure-subset-closed}{\vnbstmt{\begin{array}{@{}l@{}} \forall s, a. \\ \quad (\operatorname{is\text{-}metric\text{-}space}(s) \wedge ((a \subseteq \operatorname{pts}(s)) \wedge (\operatorname{closure}(s, a) \subseteq a))) \Rightarrow  \\ \quad \quad \operatorname{is\text{-}closed}(s, a) \end{array}}}
\vnbprop{perp-is-submodule}{\vnbstmt{\begin{array}{@{}l@{}} \forall v, s. \\ \quad (\operatorname{is\text{-}complex\text{-}inner\text{-}product\text{-}space}(v) \wedge (s \subseteq \operatorname{vec}(v))) \Rightarrow  \\ \quad \quad \operatorname{is\text{-}submodule}(v, {s}^{\perp}) \end{array}}}
\vnbprop{perp-is-closed}{\vnbstmt{\begin{array}{@{}l@{}} \forall v, s. \\ \quad (\operatorname{is\text{-}complex\text{-}inner\text{-}product\text{-}space}(v) \wedge (s \subseteq \operatorname{vec}(v))) \Rightarrow  \\ \quad \quad \operatorname{is\text{-}closed}(\operatorname{ip\text{-}metric\text{-}space}(v), {s}^{\perp}) \end{array}}}
\vnbprop{cips-dist-sq-inf}{\vnbstmt{\begin{array}{@{}l@{}} \forall v, k, i. \\ \quad (\operatorname{is\text{-}complex\text{-}inner\text{-}product\text{-}space}(v) \wedge (\operatorname{is\text{-}submodule}(v, k) \wedge (i \in \operatorname{vec}(v)))) \Rightarrow  \\ \quad \quad \exists a \in \mathbb{R}. \\ \quad \quad \quad 0 \leq a \wedge \\ \quad \quad \quad \forall p \in k. \\ \quad \quad \quad \quad a \leq \langle\cdot,\cdot\rangle_{v}(\operatorname{vadd}(v)(p, \operatorname{vneg}(v)(i)), \operatorname{vadd}(v)(p, \operatorname{vneg}(v)(i))) \wedge \\ \quad \quad \quad \forall n \in \mathbb{N}. \exists q \in k. \\ \quad \quad \quad \quad \langle\cdot,\cdot\rangle_{v}(\operatorname{vadd}(v)(q, \operatorname{vneg}(v)(i)), \operatorname{vadd}(v)(q, \operatorname{vneg}(v)(i))) \leq (a + \operatorname{recip}((n + 1))) \end{array}}}
\vnbprop{cips-midpoint-estimate}{\vnbstmt{\begin{array}{@{}l@{}} \forall v, k, i, a, p_{1}, p_{2}. \\ \quad \operatorname{is\text{-}complex\text{-}inner\text{-}product\text{-}space}(v) \wedge \\ \quad \operatorname{is\text{-}submodule}(v, k) \wedge \\ \quad i \in \operatorname{vec}(v) \wedge \\ \quad a \in \mathbb{R} \wedge \\ \quad \forall p \in k. \\ \quad \quad a \leq \langle\cdot,\cdot\rangle_{v}(\operatorname{vadd}(v)(p, \operatorname{vneg}(v)(i)), \operatorname{vadd}(v)(p, \operatorname{vneg}(v)(i))) \wedge \\ \quad p_{1} \in k \wedge \\ \quad p_{2} \in k \Rightarrow \\ \quad \quad \langle\cdot,\cdot\rangle_{v}(\operatorname{vadd}(v)(p_{1}, \operatorname{vneg}(v)(p_{2})), \operatorname{vadd}(v)(p_{1}, \operatorname{vneg}(v)(p_{2}))) \\ \quad \quad \leq\; (((2 \cdot \langle\cdot,\cdot\rangle_{v}(\operatorname{vadd}(v)(p_{1}, \operatorname{vneg}(v)(i)), \operatorname{vadd}(v)(p_{1}, \operatorname{vneg}(v)(i)))) + (2 \cdot \langle\cdot,\cdot\rangle_{v}(\operatorname{vadd}(v)(p_{2}, \operatorname{vneg}(v)(i)), \operatorname{vadd}(v)(p_{2}, \operatorname{vneg}(v)(i))))) - (4 \cdot a)) \end{array}}}
\vnbprop{cips-eq-of-dist-sq-zero}{\vnbstmt{\begin{array}{@{}l@{}} \forall v, x, y. \\ \quad \operatorname{is\text{-}complex\text{-}inner\text{-}product\text{-}space}(v) \wedge \\ \quad x \in \operatorname{vec}(v) \wedge \\ \quad y \in \operatorname{vec}(v) \wedge \\ \quad \langle\cdot,\cdot\rangle_{v}(\operatorname{vadd}(v)(x, \operatorname{vneg}(v)(y)), \operatorname{vadd}(v)(x, \operatorname{vneg}(v)(y))) = 0 \Rightarrow \\ \quad \quad x = y \end{array}}}
\vnbprop{hilbert-closest-point}{\vnbstmt{\begin{array}{@{}l@{}} \forall v, k, i. \\ \quad (\operatorname{is\text{-}hilbert\text{-}space}(v) \wedge (\operatorname{is\text{-}submodule}(v, k) \wedge (\operatorname{is\text{-}closed}(\operatorname{ip\text{-}metric\text{-}space}(v), k) \wedge (i \in \operatorname{vec}(v))))) \Rightarrow  \\ \quad \quad \exists p \in k. \forall \mathit{hdp} \in k. \\ \quad \quad \quad \langle\cdot,\cdot\rangle_{v}(\operatorname{vadd}(v)(p, \operatorname{vneg}(v)(i)), \operatorname{vadd}(v)(p, \operatorname{vneg}(v)(i))) \\ \quad \quad \quad \leq\; \langle\cdot,\cdot\rangle_{v}(\operatorname{vadd}(v)(\mathit{hdp}, \operatorname{vneg}(v)(i)), \operatorname{vadd}(v)(\mathit{hdp}, \operatorname{vneg}(v)(i))) \end{array}}}
\vnbprop{cips-variational}{\vnbstmt{\begin{array}{@{}l@{}} \forall v, w, e. \\ \quad \operatorname{is\text{-}complex\text{-}inner\text{-}product\text{-}space}(v) \wedge \\ \quad w \in \operatorname{vec}(v) \wedge \\ \quad e \in \operatorname{vec}(v) \wedge \\ \quad \forall t \in \mathbb{C}. \\ \quad \quad \langle\cdot,\cdot\rangle_{v}(w, w) \leq \langle\cdot,\cdot\rangle_{v}(\operatorname{vadd}(v)(w, \operatorname{act}(v)(t, e)), \operatorname{vadd}(v)(w, \operatorname{act}(v)(t, e))) \Rightarrow \\ \quad \quad \langle\cdot,\cdot\rangle_{v}(w, e) = 0 \end{array}}}
\vnbprop{hilbert-projection-exists}{\vnbstmt{\begin{array}{@{}l@{}} \forall v, k, i. \\ \quad (\operatorname{is\text{-}hilbert\text{-}space}(v) \wedge (\operatorname{is\text{-}submodule}(v, k) \wedge (\operatorname{is\text{-}closed}(\operatorname{ip\text{-}metric\text{-}space}(v), k) \wedge (i \in \operatorname{vec}(v))))) \Rightarrow  \\ \quad \quad \exists a \in k, b \in {k}^{\perp}.\; (i = \operatorname{vadd}(v)(a, b)) \end{array}}}
\vnbprop{hilbert-decomposition-unique}{\vnbstmt{\begin{array}{@{}l@{}} \forall v, k, a_{1}, a_{2}, b_{1}, b_{2}. \\ \quad \operatorname{is\text{-}complex\text{-}inner\text{-}product\text{-}space}(v) \wedge \\ \quad \operatorname{is\text{-}submodule}(v, k) \wedge \\ \quad a_{1} \in k \wedge \\ \quad a_{2} \in k \wedge \\ \quad b_{1} \in {k}^{\perp} \wedge \\ \quad b_{2} \in {k}^{\perp} \wedge \\ \quad \operatorname{vadd}(v)(a_{1}, b_{1}) = \operatorname{vadd}(v)(a_{2}, b_{2}) \Rightarrow \\ \quad \quad (a_{1} = a_{2}) \wedge (b_{1} = b_{2}) \end{array}}}
\vnbprop{hilbert-projection-point}{\vnbstmt{\begin{array}{@{}l@{}} \forall v, k, i. \\ \quad (\operatorname{is\text{-}hilbert\text{-}space}(v) \wedge (\operatorname{is\text{-}submodule}(v, k) \wedge (\operatorname{is\text{-}closed}(\operatorname{ip\text{-}metric\text{-}space}(v), k) \wedge (i \in \operatorname{vec}(v))))) \Rightarrow  \\ \quad \quad \exists a \in k.\; (\operatorname{vadd}(v)(i, \operatorname{vneg}(v)(a)) \in {k}^{\perp}) \end{array}}}
\vnbprop{orthogonal-projection-apply}{\vnbstmt{\begin{array}{@{}l@{}} \forall v, k, a, x. \\ \quad (\operatorname{is\text{-}orthogonal\text{-}projection}(v, k, a) \wedge (x \in \operatorname{vec}(v))) \Rightarrow  \\ \quad \quad (a(x) \in k) \wedge (\operatorname{vadd}(v)(x, \operatorname{vneg}(v)(a(x))) \in {k}^{\perp}) \end{array}}}
\vnbprop{orthogonal-projection-type}{\vnbstmt{\begin{array}{@{}l@{}} \forall v, k, a. \\ \quad \operatorname{is\text{-}orthogonal\text{-}projection}(v, k, a) \Rightarrow (a \in (\operatorname{vec}(v) \to \operatorname{vec}(v))) \end{array}}}
\vnbprop{orthogonal-projection-exists}{\vnbstmt{\begin{array}{@{}l@{}} \forall v, k. \\ \quad (\operatorname{is\text{-}hilbert\text{-}space}(v) \wedge (\operatorname{is\text{-}submodule}(v, k) \wedge \operatorname{is\text{-}closed}(\operatorname{ip\text{-}metric\text{-}space}(v), k))) \Rightarrow  \\ \quad \quad \exists \mathit{hpk}.\; \operatorname{is\text{-}orthogonal\text{-}projection}(v, k, \mathit{hpk}) \end{array}}}
\vnbprop{orthogonal-projection-unique}{\vnbstmt{\begin{array}{@{}l@{}} \forall v, k, a, b. \\ \quad \operatorname{is\text{-}complex\text{-}inner\text{-}product\text{-}space}(v) \wedge \\ \quad \operatorname{is\text{-}submodule}(v, k) \wedge \\ \quad \operatorname{is\text{-}orthogonal\text{-}projection}(v, k, a) \wedge \\ \quad \operatorname{is\text{-}orthogonal\text{-}projection}(v, k, b) \Rightarrow \\ \quad \quad a = b \end{array}}}
\vnbprop{orthogonal-projection-additive}{\vnbstmt{\begin{array}{@{}l@{}} \forall v, k, a, x, y. \\ \quad \operatorname{is\text{-}complex\text{-}inner\text{-}product\text{-}space}(v) \wedge \\ \quad \operatorname{is\text{-}submodule}(v, k) \wedge \\ \quad \operatorname{is\text{-}orthogonal\text{-}projection}(v, k, a) \wedge \\ \quad x \in \operatorname{vec}(v) \wedge \\ \quad y \in \operatorname{vec}(v) \Rightarrow \\ \quad \quad a(\operatorname{vadd}(v)(x, y)) = \operatorname{vadd}(v)(a(x), a(y)) \end{array}}}
\vnbprop{orthogonal-projection-homogeneous}{\vnbstmt{\begin{array}{@{}l@{}} \forall v, k, a, c, x. \\ \quad \operatorname{is\text{-}complex\text{-}inner\text{-}product\text{-}space}(v) \wedge \\ \quad \operatorname{is\text{-}submodule}(v, k) \wedge \\ \quad \operatorname{is\text{-}orthogonal\text{-}projection}(v, k, a) \wedge \\ \quad c \in \mathbb{C} \wedge \\ \quad x \in \operatorname{vec}(v) \Rightarrow \\ \quad \quad a(\operatorname{act}(v)(c, x)) = \operatorname{act}(v)(c, a(x)) \end{array}}}
\vnbprop{orthogonal-projection-norm-le}{\vnbstmt{\begin{array}{@{}l@{}} \forall v, k, a, x. \\ \quad \operatorname{is\text{-}complex\text{-}inner\text{-}product\text{-}space}(v) \wedge \\ \quad \operatorname{is\text{-}submodule}(v, k) \wedge \\ \quad \operatorname{is\text{-}orthogonal\text{-}projection}(v, k, a) \wedge \\ \quad x \in \operatorname{vec}(v) \Rightarrow \\ \quad \quad \lVert a(x) \rVert \leq \lVert x \rVert \end{array}}}
\vnbprop{orthogonal-projection-fixes}{\vnbstmt{\begin{array}{@{}l@{}} \forall v, k, a, y. \\ \quad \operatorname{is\text{-}complex\text{-}inner\text{-}product\text{-}space}(v) \wedge \\ \quad \operatorname{is\text{-}submodule}(v, k) \wedge \\ \quad \operatorname{is\text{-}orthogonal\text{-}projection}(v, k, a) \wedge \\ \quad y \in k \Rightarrow \\ \quad \quad a(y) = y \end{array}}}
\vnbprop{orthogonal-projection-idempotent}{\vnbstmt{\begin{array}{@{}l@{}} \forall v, k, a, x. \\ \quad \operatorname{is\text{-}complex\text{-}inner\text{-}product\text{-}space}(v) \wedge \\ \quad \operatorname{is\text{-}submodule}(v, k) \wedge \\ \quad \operatorname{is\text{-}orthogonal\text{-}projection}(v, k, a) \wedge \\ \quad x \in \operatorname{vec}(v) \Rightarrow \\ \quad \quad a(a(x)) = a(x) \end{array}}}
\vnbprop{cc-linear-functional-laws}{\vnbstmt{\begin{array}{@{}l@{}} \forall v, f. \\ \quad \operatorname{is\text{-}cc\text{-}linear\text{-}functional}(v, f) \Rightarrow  \\ \quad \quad f \in (\operatorname{vec}(v) \to \mathbb{C}) \wedge \\ \quad \quad \forall x, y \in \operatorname{vec}(v).\; (f(\operatorname{vadd}(v)(x, y)) = (f(x) + f(y))) \wedge \\ \quad \quad \forall c \in \mathbb{C}, x \in \operatorname{vec}(v).\; (f(\operatorname{act}(v)(c, x)) = (c \cdot f(x))) \end{array}}}
\vnbprop{cc-linear-functional-zero-neg}{\vnbstmt{\begin{array}{@{}l@{}} \forall v, f. \\ \quad (\operatorname{is\text{-}complex\text{-}inner\text{-}product\text{-}space}(v) \wedge \operatorname{is\text{-}cc\text{-}linear\text{-}functional}(v, f)) \Rightarrow  \\ \quad \quad (f(\operatorname{vzero}(v)) = 0) \wedge \forall x \in \operatorname{vec}(v).\; (f(\operatorname{vneg}(v)(x)) = -f(x)) \end{array}}}
\vnbprop{cc-linear-functional-kernel-membership}{\vnbstmt{\begin{array}{@{}l@{}} \forall v, f, x. \\ \quad (\operatorname{is\text{-}complex\text{-}inner\text{-}product\text{-}space}(v) \wedge (\operatorname{is\text{-}cc\text{-}linear\text{-}functional}(v, f) \wedge (x \in \operatorname{vec}(v)))) \Rightarrow  \\ \quad \quad (x \in \operatorname{preimage}(\operatorname{ip\text{-}metric\text{-}space}(v), f, \operatorname{closed\text{-}ball}(\operatorname{cc\text{-}ms}, 0, 0))) \Leftrightarrow (f(x) = 0) \end{array}}}
\vnbprop{cc-linear-functional-kernel-submodule}{\vnbstmt{\begin{array}{@{}l@{}} \forall v, f. \\ \quad (\operatorname{is\text{-}complex\text{-}inner\text{-}product\text{-}space}(v) \wedge \operatorname{is\text{-}cc\text{-}linear\text{-}functional}(v, f)) \Rightarrow  \\ \quad \quad \operatorname{is\text{-}submodule}(v, \operatorname{preimage}(\operatorname{ip\text{-}metric\text{-}space}(v), f, \operatorname{closed\text{-}ball}(\operatorname{cc\text{-}ms}, 0, 0))) \end{array}}}
\vnbprop{cc-linear-functional-kernel-closed}{\vnbstmt{\begin{array}{@{}l@{}} \forall v, f. \\ \quad (\operatorname{is\text{-}complex\text{-}inner\text{-}product\text{-}space}(v) \wedge \operatorname{is\text{-}continuous}(\operatorname{ip\text{-}metric\text{-}space}(v), \operatorname{cc\text{-}ms}, f)) \Rightarrow  \\ \quad \quad \operatorname{is\text{-}closed}(\operatorname{ip\text{-}metric\text{-}space}(v), \operatorname{preimage}(\operatorname{ip\text{-}metric\text{-}space}(v), f, \operatorname{closed\text{-}ball}(\operatorname{cc\text{-}ms}, 0, 0))) \end{array}}}
\vnbprop{cips-ip-separates}{\vnbstmt{\begin{array}{@{}l@{}} \forall v, e_{1}, e_{2}. \\ \quad \operatorname{is\text{-}complex\text{-}inner\text{-}product\text{-}space}(v) \wedge \\ \quad e_{1} \in \operatorname{vec}(v) \wedge \\ \quad e_{2} \in \operatorname{vec}(v) \wedge \\ \quad \forall x \in \operatorname{vec}(v).\; (\langle\cdot,\cdot\rangle_{v}(x, e_{1}) = \langle\cdot,\cdot\rangle_{v}(x, e_{2})) \Rightarrow \\ \quad \quad e_{1} = e_{2} \end{array}}}
\vnbprop{cips-ip-magnitude-le}{\vnbstmt{\begin{array}{@{}l@{}} \forall v, x, e. \\ \quad (\operatorname{is\text{-}complex\text{-}inner\text{-}product\text{-}space}(v) \wedge ((x \in \operatorname{vec}(v)) \wedge (e \in \operatorname{vec}(v)))) \Rightarrow  \\ \quad \quad \operatorname{magnitude}(\langle\cdot,\cdot\rangle_{v}(x, e)) \leq (\lVert x \rVert \cdot \lVert e \rVert) \end{array}}}
\vnbprop{cips-ip-bound-least}{\vnbstmt{\begin{array}{@{}l@{}} \forall v, e, c. \\ \quad \operatorname{is\text{-}complex\text{-}inner\text{-}product\text{-}space}(v) \wedge \\ \quad e \in \operatorname{vec}(v) \wedge \\ \quad c \in \mathbb{R} \wedge \\ \quad 0 \leq c \wedge \\ \quad \forall x \in \operatorname{vec}(v).\; (\operatorname{magnitude}(\langle\cdot,\cdot\rangle_{v}(x, e)) \leq (c \cdot \lVert x \rVert)) \Rightarrow \\ \quad \quad \lVert e \rVert \leq c \end{array}}}
\vnbprop{riesz-representation-of-kernel}{\vnbstmt{\begin{array}{@{}l@{}} \forall v, f, k. \\ \quad \operatorname{is\text{-}hilbert\text{-}space}(v) \wedge \\ \quad \operatorname{is\text{-}cc\text{-}linear\text{-}functional}(v, f) \wedge \\ \quad \operatorname{is\text{-}submodule}(v, k) \wedge \\ \quad \operatorname{is\text{-}closed}(\operatorname{ip\text{-}metric\text{-}space}(v), k) \wedge \\ \quad \forall x \in \operatorname{vec}(v).\; ((x \in k) \Leftrightarrow (f(x) = 0)) \Rightarrow \\ \quad \quad \exists e \in \operatorname{vec}(v).\; \forall \mathit{hrx} \in \operatorname{vec}(v).\; (f(\mathit{hrx}) = \langle\cdot,\cdot\rangle_{v}(\mathit{hrx}, e)) \end{array}}}
\vnbprop{riesz-representation-exists}{\vnbstmt{\begin{array}{@{}l@{}} \forall v, f. \\ \quad \operatorname{is\text{-}hilbert\text{-}space}(v) \wedge \\ \quad \operatorname{is\text{-}cc\text{-}linear\text{-}functional}(v, f) \wedge \\ \quad \operatorname{is\text{-}continuous}(\operatorname{ip\text{-}metric\text{-}space}(v), \operatorname{cc\text{-}ms}, f) \Rightarrow \\ \quad \quad \exists e \in \operatorname{vec}(v).\; \forall x \in \operatorname{vec}(v).\; (f(x) = \langle\cdot,\cdot\rangle_{v}(x, e)) \end{array}}}
\vnbprop{riesz-representation}{\vnbstmt{\begin{array}{@{}l@{}} \forall v, f. \\ \quad \operatorname{is\text{-}hilbert\text{-}space}(v) \wedge \\ \quad \operatorname{is\text{-}cc\text{-}linear\text{-}functional}(v, f) \wedge \\ \quad \operatorname{is\text{-}continuous}(\operatorname{ip\text{-}metric\text{-}space}(v), \operatorname{cc\text{-}ms}, f) \Rightarrow \\ \quad \quad \exists e \in \operatorname{vec}(v). \\ \quad \quad \quad \forall x \in \operatorname{vec}(v).\; (f(x) = \langle\cdot,\cdot\rangle_{v}(x, e)) \wedge \\ \quad \quad \quad \forall x \in \operatorname{vec}(v).\; (\operatorname{magnitude}(f(x)) \leq (\lVert x \rVert \cdot \lVert e \rVert)) \wedge \\ \quad \quad \quad \forall c \in \mathbb{R}. \\ \quad \quad \quad \quad ((0 \leq c) \wedge \forall x \in \operatorname{vec}(v).\; (\operatorname{magnitude}(f(x)) \leq (c \cdot \lVert x \rVert))) \Rightarrow  \\ \quad \quad \quad \quad \quad \lVert e \rVert \leq c \end{array}}}
\vnbprop{riesz-representation-unique}{\vnbstmt{\begin{array}{@{}l@{}} \forall v, f, e_{1}, e_{2}. \\ \quad \operatorname{is\text{-}complex\text{-}inner\text{-}product\text{-}space}(v) \wedge \\ \quad e_{1} \in \operatorname{vec}(v) \wedge \\ \quad e_{2} \in \operatorname{vec}(v) \wedge \\ \quad \forall x \in \operatorname{vec}(v).\; (f(x) = \langle\cdot,\cdot\rangle_{v}(x, e_{1})) \wedge \\ \quad \forall x \in \operatorname{vec}(v).\; (f(x) = \langle\cdot,\cdot\rangle_{v}(x, e_{2})) \Rightarrow \\ \quad \quad e_{1} = e_{2} \end{array}}}
\vnbprop{ip-closure-elim}{\vnbstmt{\begin{array}{@{}l@{}} \forall v, a, x, e. \\ \quad (\operatorname{is\text{-}complex\text{-}inner\text{-}product\text{-}space}(v) \wedge ((x \in \operatorname{closure}(\operatorname{ip\text{-}metric\text{-}space}(v), a)) \wedge \operatorname{pos\text{-}rr}(e))) \Rightarrow  \\ \quad \quad \exists y \in a. \\ \quad \quad \quad \langle\cdot,\cdot\rangle_{v}(\operatorname{vadd}(v)(x, \operatorname{vneg}(v)(y)), \operatorname{vadd}(v)(x, \operatorname{vneg}(v)(y))) \leq e \end{array}}}
\vnbprop{ip-closure-intro}{\vnbstmt{\begin{array}{@{}l@{}} \forall v, a, x. \\ \quad \operatorname{is\text{-}complex\text{-}inner\text{-}product\text{-}space}(v) \wedge \\ \quad a \subseteq \operatorname{vec}(v) \wedge \\ \quad x \in \operatorname{vec}(v) \wedge \\ \quad \forall e. \\ \quad \quad \operatorname{pos\text{-}rr}(e) \Rightarrow  \\ \quad \quad \quad \exists y \in a. \\ \quad \quad \quad \quad \langle\cdot,\cdot\rangle_{v}(\operatorname{vadd}(v)(x, \operatorname{vneg}(v)(y)), \operatorname{vadd}(v)(x, \operatorname{vneg}(v)(y))) \leq e \Rightarrow \\ \quad \quad x \in \operatorname{closure}(\operatorname{ip\text{-}metric\text{-}space}(v), a) \end{array}}}
\vnbprop{ip-closure-submodule}{\vnbstmt{\begin{array}{@{}l@{}} \forall v, k. \\ \quad (\operatorname{is\text{-}complex\text{-}inner\text{-}product\text{-}space}(v) \wedge \operatorname{is\text{-}submodule}(v, k)) \Rightarrow  \\ \quad \quad \operatorname{is\text{-}submodule}(v, \operatorname{closure}(\operatorname{ip\text{-}metric\text{-}space}(v), k)) \end{array}}}
\vnbprop{hilbert-dense-iff-perp-zero}{\vnbstmt{\begin{array}{@{}l@{}} \forall v, k. \\ \quad (\operatorname{is\text{-}hilbert\text{-}space}(v) \wedge \operatorname{is\text{-}submodule}(v, k)) \Rightarrow  \\ \quad \quad (\operatorname{closure}(\operatorname{ip\text{-}metric\text{-}space}(v), k) = \operatorname{vec}(v)) \Leftrightarrow \forall z \in {k}^{\perp}.\; (z = \operatorname{vzero}(v)) \end{array}}}

\vnbbookfile{theorem-library/rake-asm-b3}{theorem-library/rake-asm-b3}
\noindent{\small theorem-library/rake-asm-b3.scm -- RAKE-ASM-B, agent RB-3 (2026-10-12).  Prefix rb3-.}\par
\vnbprop{rb3-cc-ring-carr}{\vnbstmt{(\operatorname{carr}(\operatorname{normed\text{-}field\text{-}as\text{-}commutative\text{-}ring}(\operatorname{cc\text{-}normed\text{-}field})) \simeq \mathbb{C})}}
\vnbprop{rb3-cc-ring-add}{\vnbstmt{(\operatorname{add}(\operatorname{normed\text{-}field\text{-}as\text{-}commutative\text{-}ring}(\operatorname{cc\text{-}normed\text{-}field})) \simeq (\lambda \langle \operatorname{x\_}, \operatorname{y\_} \rangle \in (\mathbb{C} \times \mathbb{C}).\; (\operatorname{x\_} + \operatorname{y\_})))}}
\vnbprop{rb3-cc-ring-mul}{\vnbstmt{(\operatorname{mul}(\operatorname{normed\text{-}field\text{-}as\text{-}commutative\text{-}ring}(\operatorname{cc\text{-}normed\text{-}field})) \simeq (\lambda \langle \operatorname{x\_}, \operatorname{y\_} \rangle \in (\mathbb{C} \times \mathbb{C}).\; (\operatorname{x\_} \cdot \operatorname{y\_})))}}
\vnbprop{rb3-cc-ring-one}{\vnbstmt{(\operatorname{one}(\operatorname{normed\text{-}field\text{-}as\text{-}commutative\text{-}ring}(\operatorname{cc\text{-}normed\text{-}field})) \simeq 1)}}
\vnbprop{rb3-cc-ring-zero}{\vnbstmt{(\operatorname{zero}(\operatorname{normed\text{-}field\text{-}as\text{-}commutative\text{-}ring}(\operatorname{cc\text{-}normed\text{-}field})) \simeq 0)}}
\vnbprop{rb3-c-carr}{\vnbstmt{\begin{array}{@{}l@{}} \forall m. \\ \quad (\operatorname{is\text{-}module}(m) \wedge (\operatorname{scal}(m) = \operatorname{normed\text{-}field\text{-}as\text{-}commutative\text{-}ring}(\operatorname{cc\text{-}normed\text{-}field}))) \Rightarrow  \\ \quad \quad \operatorname{carr}(\operatorname{scal}(m)) = \mathbb{C} \end{array}}}
\vnbprop{rb3-c-in-carr}{\vnbstmt{\begin{array}{@{}l@{}} \forall m.; \forall r \in \mathbb{C}. \\ \quad (\operatorname{is\text{-}module}(m) \wedge (\operatorname{scal}(m) = \operatorname{normed\text{-}field\text{-}as\text{-}commutative\text{-}ring}(\operatorname{cc\text{-}normed\text{-}field}))) \Rightarrow  \\ \quad \quad r \in \operatorname{carr}(\operatorname{scal}(m)) \end{array}}}
\vnbprop{rb3-c-add-eq}{\vnbstmt{\begin{array}{@{}l@{}} \forall m.; \forall r, s \in \mathbb{C}. \\ \quad (\operatorname{is\text{-}module}(m) \wedge (\operatorname{scal}(m) = \operatorname{normed\text{-}field\text{-}as\text{-}commutative\text{-}ring}(\operatorname{cc\text{-}normed\text{-}field}))) \Rightarrow  \\ \quad \quad \operatorname{add}(\operatorname{scal}(m))(r, s) = (r + s) \end{array}}}
\vnbprop{rb3-c-mul-eq}{\vnbstmt{\begin{array}{@{}l@{}} \forall m.; \forall r, s \in \mathbb{C}. \\ \quad (\operatorname{is\text{-}module}(m) \wedge (\operatorname{scal}(m) = \operatorname{normed\text{-}field\text{-}as\text{-}commutative\text{-}ring}(\operatorname{cc\text{-}normed\text{-}field}))) \Rightarrow  \\ \quad \quad \operatorname{mul}(\operatorname{scal}(m))(r, s) = (r \cdot s) \end{array}}}
\vnbprop{rb3-c-one-eq}{\vnbstmt{\begin{array}{@{}l@{}} \forall m. \\ \quad (\operatorname{is\text{-}module}(m) \wedge (\operatorname{scal}(m) = \operatorname{normed\text{-}field\text{-}as\text{-}commutative\text{-}ring}(\operatorname{cc\text{-}normed\text{-}field}))) \Rightarrow  \\ \quad \quad \operatorname{one}(\operatorname{scal}(m)) = 1 \end{array}}}
\vnbprop{rb3-c-act-in}{\vnbstmt{\begin{array}{@{}l@{}} \forall m.; \forall r \in \mathbb{C}, x \in \operatorname{vec}(m). \\ \quad (\operatorname{is\text{-}module}(m) \wedge (\operatorname{scal}(m) = \operatorname{normed\text{-}field\text{-}as\text{-}commutative\text{-}ring}(\operatorname{cc\text{-}normed\text{-}field}))) \Rightarrow  \\ \quad \quad \operatorname{act}(m)(r, x) \in \operatorname{vec}(m) \end{array}}}
\vnbprop{rb3-c-act-add}{\vnbstmt{\begin{array}{@{}l@{}} \forall m.; \forall r, s \in \mathbb{C}, x \in \operatorname{vec}(m). \\ \quad (\operatorname{is\text{-}module}(m) \wedge (\operatorname{scal}(m) = \operatorname{normed\text{-}field\text{-}as\text{-}commutative\text{-}ring}(\operatorname{cc\text{-}normed\text{-}field}))) \Rightarrow  \\ \quad \quad \operatorname{act}(m)((r + s), x) = \operatorname{vadd}(m)(\operatorname{act}(m)(r, x), \operatorname{act}(m)(s, x)) \end{array}}}
\vnbprop{rb3-c-act-mul}{\vnbstmt{\begin{array}{@{}l@{}} \forall m.; \forall r, s \in \mathbb{C}, x \in \operatorname{vec}(m). \\ \quad (\operatorname{is\text{-}module}(m) \wedge (\operatorname{scal}(m) = \operatorname{normed\text{-}field\text{-}as\text{-}commutative\text{-}ring}(\operatorname{cc\text{-}normed\text{-}field}))) \Rightarrow  \\ \quad \quad \operatorname{act}(m)((r \cdot s), x) = \operatorname{act}(m)(r, \operatorname{act}(m)(s, x)) \end{array}}}
\vnbprop{rb3-c-act-vadd}{\vnbstmt{\begin{array}{@{}l@{}} \forall m.; \forall r \in \mathbb{C}, x, y \in \operatorname{vec}(m). \\ \quad (\operatorname{is\text{-}module}(m) \wedge (\operatorname{scal}(m) = \operatorname{normed\text{-}field\text{-}as\text{-}commutative\text{-}ring}(\operatorname{cc\text{-}normed\text{-}field}))) \Rightarrow  \\ \quad \quad \operatorname{act}(m)(r, \operatorname{vadd}(m)(x, y)) = \operatorname{vadd}(m)(\operatorname{act}(m)(r, x), \operatorname{act}(m)(r, y)) \end{array}}}
\vnbprop{rb3-tuple6}{\vnbstmt{\begin{array}{@{}l@{}} \forall a_{2}, a_{3}, a_{4}, a_{5}, a_{6}. \\ \quad ((a_{2} \in \mathbf{Set}) \wedge ((a_{3} \in \mathbf{Set}) \wedge ((a_{4} \in \mathbf{Set}) \wedge ((a_{5} \in \mathbf{Set}) \wedge (a_{6} \in \mathbf{Set}))))) \Rightarrow  \\ \quad \quad (\operatorname{scal}(\langle \operatorname{normed\text{-}field\text{-}as\text{-}commutative\text{-}ring}(\operatorname{rr\text{-}normed\text{-}field}), a_{2}, a_{3}, a_{4}, a_{5}, a_{6} \rangle) \simeq \operatorname{normed\text{-}field\text{-}as\text{-}commutative\text{-}ring}(\operatorname{rr\text{-}normed\text{-}field})) \wedge \\ \quad \quad (\operatorname{vec}(\langle \operatorname{normed\text{-}field\text{-}as\text{-}commutative\text{-}ring}(\operatorname{rr\text{-}normed\text{-}field}), a_{2}, a_{3}, a_{4}, a_{5}, a_{6} \rangle) \simeq a_{2}) \wedge \\ \quad \quad (\operatorname{vadd}(\langle \operatorname{normed\text{-}field\text{-}as\text{-}commutative\text{-}ring}(\operatorname{rr\text{-}normed\text{-}field}), a_{2}, a_{3}, a_{4}, a_{5}, a_{6} \rangle) \simeq a_{3}) \wedge \\ \quad \quad (\operatorname{vzero}(\langle \operatorname{normed\text{-}field\text{-}as\text{-}commutative\text{-}ring}(\operatorname{rr\text{-}normed\text{-}field}), a_{2}, a_{3}, a_{4}, a_{5}, a_{6} \rangle) \simeq a_{4}) \wedge \\ \quad \quad (\operatorname{vneg}(\langle \operatorname{normed\text{-}field\text{-}as\text{-}commutative\text{-}ring}(\operatorname{rr\text{-}normed\text{-}field}), a_{2}, a_{3}, a_{4}, a_{5}, a_{6} \rangle) \simeq a_{5}) \wedge \\ \quad \quad (\operatorname{act}(\langle \operatorname{normed\text{-}field\text{-}as\text{-}commutative\text{-}ring}(\operatorname{rr\text{-}normed\text{-}field}), a_{2}, a_{3}, a_{4}, a_{5}, a_{6} \rangle) \simeq a_{6}) \wedge \\ \quad \quad \operatorname{length}(\langle \operatorname{normed\text{-}field\text{-}as\text{-}commutative\text{-}ring}(\operatorname{rr\text{-}normed\text{-}field}), a_{2}, a_{3}, a_{4}, a_{5}, a_{6} \rangle) = 6 \end{array}}}
\vnbprop{rb3-mr-exists}{\vnbstmt{\begin{array}{@{}l@{}} \forall m. \\ \quad (\operatorname{is\text{-}module}(m) \wedge (\operatorname{scal}(m) = \operatorname{normed\text{-}field\text{-}as\text{-}commutative\text{-}ring}(\operatorname{cc\text{-}normed\text{-}field}))) \Rightarrow  \\ \quad \quad \exists r. \\ \quad \quad \quad \operatorname{is\text{-}module}(r) \wedge \\ \quad \quad \quad \operatorname{scal}(r) = \operatorname{normed\text{-}field\text{-}as\text{-}commutative\text{-}ring}(\operatorname{rr\text{-}normed\text{-}field}) \wedge \\ \quad \quad \quad \operatorname{vec}(r) = \operatorname{vec}(m) \wedge \\ \quad \quad \quad \operatorname{vadd}(r) = \operatorname{vadd}(m) \wedge \\ \quad \quad \quad \operatorname{vzero}(r) = \operatorname{vzero}(m) \wedge \\ \quad \quad \quad \operatorname{vneg}(r) = \operatorname{vneg}(m) \wedge \\ \quad \quad \quad \forall a \in \mathbb{R}, x \in \operatorname{vec}(m).\; (\operatorname{act}(r)(a, x) = \operatorname{act}(m)(a, x)) \end{array}}}
\vnbprop{rb3-mr-sublinear}{\vnbstmt{\begin{array}{@{}l@{}} \forall m, r, p. \\ \quad \operatorname{is\text{-}module}(m) \wedge \\ \quad \operatorname{scal}(m) = \operatorname{normed\text{-}field\text{-}as\text{-}commutative\text{-}ring}(\operatorname{cc\text{-}normed\text{-}field}) \wedge \\ \quad \operatorname{is\text{-}module}(r) \wedge \\ \quad \operatorname{scal}(r) = \operatorname{normed\text{-}field\text{-}as\text{-}commutative\text{-}ring}(\operatorname{rr\text{-}normed\text{-}field}) \wedge \\ \quad \operatorname{vec}(r) = \operatorname{vec}(m) \wedge \\ \quad \operatorname{vadd}(r) = \operatorname{vadd}(m) \wedge \\ \quad \operatorname{vzero}(r) = \operatorname{vzero}(m) \wedge \\ \quad \operatorname{vneg}(r) = \operatorname{vneg}(m) \wedge \\ \quad \forall a \in \mathbb{R}, x \in \operatorname{vec}(m).\; (\operatorname{act}(r)(a, x) = \operatorname{act}(m)(a, x)) \wedge \\ \quad \operatorname{is\text{-}seminorm}(\operatorname{cc\text{-}normed\text{-}field}, m, p) \Rightarrow \\ \quad \quad \operatorname{is\text{-}sublinear}(r, p) \end{array}}}
\vnbprop{rb3-mr-submodule}{\vnbstmt{\begin{array}{@{}l@{}} \forall m, r, s. \\ \quad \operatorname{is\text{-}module}(m) \wedge \\ \quad \operatorname{scal}(m) = \operatorname{normed\text{-}field\text{-}as\text{-}commutative\text{-}ring}(\operatorname{cc\text{-}normed\text{-}field}) \wedge \\ \quad \operatorname{is\text{-}module}(r) \wedge \\ \quad \operatorname{scal}(r) = \operatorname{normed\text{-}field\text{-}as\text{-}commutative\text{-}ring}(\operatorname{rr\text{-}normed\text{-}field}) \wedge \\ \quad \operatorname{vec}(r) = \operatorname{vec}(m) \wedge \\ \quad \operatorname{vadd}(r) = \operatorname{vadd}(m) \wedge \\ \quad \operatorname{vzero}(r) = \operatorname{vzero}(m) \wedge \\ \quad \operatorname{vneg}(r) = \operatorname{vneg}(m) \wedge \\ \quad \forall a \in \mathbb{R}, x \in \operatorname{vec}(m).\; (\operatorname{act}(r)(a, x) = \operatorname{act}(m)(a, x)) \wedge \\ \quad \operatorname{is\text{-}submodule}(m, s) \Rightarrow \\ \quad \quad \operatorname{is\text{-}submodule}(r, s) \end{array}}}
\vnbprop{rb3-mr-re-linear}{\vnbstmt{\begin{array}{@{}l@{}} \forall m, r, s, f.; \forall a \in (s \to \mathbb{R}). \\ \quad \operatorname{is\text{-}module}(m) \wedge \\ \quad \operatorname{scal}(m) = \operatorname{normed\text{-}field\text{-}as\text{-}commutative\text{-}ring}(\operatorname{cc\text{-}normed\text{-}field}) \wedge \\ \quad \operatorname{is\text{-}module}(r) \wedge \\ \quad \operatorname{scal}(r) = \operatorname{normed\text{-}field\text{-}as\text{-}commutative\text{-}ring}(\operatorname{rr\text{-}normed\text{-}field}) \wedge \\ \quad \operatorname{vec}(r) = \operatorname{vec}(m) \wedge \\ \quad \operatorname{vadd}(r) = \operatorname{vadd}(m) \wedge \\ \quad \operatorname{vzero}(r) = \operatorname{vzero}(m) \wedge \\ \quad \operatorname{vneg}(r) = \operatorname{vneg}(m) \wedge \\ \quad \forall b \in \mathbb{R}, x \in \operatorname{vec}(m).\; (\operatorname{act}(r)(b, x) = \operatorname{act}(m)(b, x)) \wedge \\ \quad \operatorname{is\text{-}submodule}(m, s) \wedge \\ \quad \operatorname{is\text{-}k\text{-}linear\text{-}on}(m, s, f) \wedge \\ \quad \forall z \in s.\; (a(z) \simeq \operatorname{real\text{-}part}(f(z))) \Rightarrow \\ \quad \quad \operatorname{is\text{-}linear\text{-}functional\text{-}on}(r, s, a) \end{array}}}
\vnbprop{rb3-cf-re}{\vnbstmt{\begin{array}{@{}l@{}} \forall m.; \forall n \in (\operatorname{vec}(m) \to \mathbb{R}), a \in (\operatorname{vec}(m) \to \mathbb{C}), y \in \operatorname{vec}(m). \\ \quad \operatorname{is\text{-}module}(m) \wedge \\ \quad \operatorname{scal}(m) = \operatorname{normed\text{-}field\text{-}as\text{-}commutative\text{-}ring}(\operatorname{cc\text{-}normed\text{-}field}) \wedge \\ \quad \forall x, \mathit{ray} \in \operatorname{vec}(m).\; (n(\operatorname{vadd}(m)(x, \mathit{ray})) = (n(x) + n(\mathit{ray}))) \wedge \\ \quad \forall r \in \mathbb{R}, \mathit{rhx} \in \operatorname{vec}(m).\; (n(\operatorname{act}(m)(r, \mathit{rhx})) = (r \cdot n(\mathit{rhx}))) \wedge \\ \quad \forall \mathit{cvy} \in \operatorname{vec}(m).\; (a(\mathit{cvy}) \simeq (n(\mathit{cvy}) + (-n(\operatorname{act}(m)(+i, \mathit{cvy})) \cdot +i))) \Rightarrow \\ \quad \quad \operatorname{real\text{-}part}(a(y)) = n(y) \end{array}}}
\vnbprop{rb3-cf-add}{\vnbstmt{\begin{array}{@{}l@{}} \forall m.; \forall n \in (\operatorname{vec}(m) \to \mathbb{R}), a \in (\operatorname{vec}(m) \to \mathbb{C}), x, y \in \operatorname{vec}(m). \\ \quad \operatorname{is\text{-}module}(m) \wedge \\ \quad \operatorname{scal}(m) = \operatorname{normed\text{-}field\text{-}as\text{-}commutative\text{-}ring}(\operatorname{cc\text{-}normed\text{-}field}) \wedge \\ \quad \forall \mathit{rax}, \mathit{ray} \in \operatorname{vec}(m).\; (n(\operatorname{vadd}(m)(\mathit{rax}, \mathit{ray})) = (n(\mathit{rax}) + n(\mathit{ray}))) \wedge \\ \quad \forall r \in \mathbb{R}, \mathit{rhx} \in \operatorname{vec}(m).\; (n(\operatorname{act}(m)(r, \mathit{rhx})) = (r \cdot n(\mathit{rhx}))) \wedge \\ \quad \forall \mathit{cvy} \in \operatorname{vec}(m).\; (a(\mathit{cvy}) \simeq (n(\mathit{cvy}) + (-n(\operatorname{act}(m)(+i, \mathit{cvy})) \cdot +i))) \Rightarrow \\ \quad \quad a(\operatorname{vadd}(m)(x, y)) = (a(x) + a(y)) \end{array}}}
\vnbprop{rb3-cf-rhom}{\vnbstmt{\begin{array}{@{}l@{}} \forall m.; \forall n \in (\operatorname{vec}(m) \to \mathbb{R}), a \in (\operatorname{vec}(m) \to \mathbb{C}), r \in \mathbb{R}, x \in \operatorname{vec}(m). \\ \quad \operatorname{is\text{-}module}(m) \wedge \\ \quad \operatorname{scal}(m) = \operatorname{normed\text{-}field\text{-}as\text{-}commutative\text{-}ring}(\operatorname{cc\text{-}normed\text{-}field}) \wedge \\ \quad \forall \mathit{rax}, y \in \operatorname{vec}(m).\; (n(\operatorname{vadd}(m)(\mathit{rax}, y)) = (n(\mathit{rax}) + n(y))) \wedge \\ \quad \forall \mathit{rhr} \in \mathbb{R}, \mathit{rhx} \in \operatorname{vec}(m).\; (n(\operatorname{act}(m)(\mathit{rhr}, \mathit{rhx})) = (\mathit{rhr} \cdot n(\mathit{rhx}))) \wedge \\ \quad \forall \mathit{cvy} \in \operatorname{vec}(m).\; (a(\mathit{cvy}) \simeq (n(\mathit{cvy}) + (-n(\operatorname{act}(m)(+i, \mathit{cvy})) \cdot +i))) \Rightarrow \\ \quad \quad a(\operatorname{act}(m)(r, x)) = (r \cdot a(x)) \end{array}}}
\vnbprop{rb3-cf-ihom}{\vnbstmt{\begin{array}{@{}l@{}} \forall m.; \forall n \in (\operatorname{vec}(m) \to \mathbb{R}), a \in (\operatorname{vec}(m) \to \mathbb{C}), x \in \operatorname{vec}(m). \\ \quad \operatorname{is\text{-}module}(m) \wedge \\ \quad \operatorname{scal}(m) = \operatorname{normed\text{-}field\text{-}as\text{-}commutative\text{-}ring}(\operatorname{cc\text{-}normed\text{-}field}) \wedge \\ \quad \forall \mathit{rax}, y \in \operatorname{vec}(m).\; (n(\operatorname{vadd}(m)(\mathit{rax}, y)) = (n(\mathit{rax}) + n(y))) \wedge \\ \quad \forall r \in \mathbb{R}, \mathit{rhx} \in \operatorname{vec}(m).\; (n(\operatorname{act}(m)(r, \mathit{rhx})) = (r \cdot n(\mathit{rhx}))) \wedge \\ \quad \forall \mathit{cvy} \in \operatorname{vec}(m).\; (a(\mathit{cvy}) \simeq (n(\mathit{cvy}) + (-n(\operatorname{act}(m)(+i, \mathit{cvy})) \cdot +i))) \Rightarrow \\ \quad \quad a(\operatorname{act}(m)(+i, x)) = (+i \cdot a(x)) \end{array}}}
\vnbprop{rb3-cf-abhom}{\vnbstmt{\begin{array}{@{}l@{}} \forall m.; \forall n \in (\operatorname{vec}(m) \to \mathbb{R}), a \in (\operatorname{vec}(m) \to \mathbb{C}), \alpha, b \in \mathbb{R}, x \in \operatorname{vec}(m). \\ \quad \operatorname{is\text{-}module}(m) \wedge \\ \quad \operatorname{scal}(m) = \operatorname{normed\text{-}field\text{-}as\text{-}commutative\text{-}ring}(\operatorname{cc\text{-}normed\text{-}field}) \wedge \\ \quad \forall \mathit{rax}, y \in \operatorname{vec}(m).\; (n(\operatorname{vadd}(m)(\mathit{rax}, y)) = (n(\mathit{rax}) + n(y))) \wedge \\ \quad \forall r \in \mathbb{R}, \mathit{rhx} \in \operatorname{vec}(m).\; (n(\operatorname{act}(m)(r, \mathit{rhx})) = (r \cdot n(\mathit{rhx}))) \wedge \\ \quad \forall \mathit{cvy} \in \operatorname{vec}(m).\; (a(\mathit{cvy}) \simeq (n(\mathit{cvy}) + (-n(\operatorname{act}(m)(+i, \mathit{cvy})) \cdot +i))) \Rightarrow \\ \quad \quad a(\operatorname{act}(m)((\alpha + (+i \cdot b)), x)) = ((\alpha + (+i \cdot b)) \cdot a(x)) \end{array}}}
\vnbprop{rb3-cf-chom}{\vnbstmt{\begin{array}{@{}l@{}} \forall m.; \forall n \in (\operatorname{vec}(m) \to \mathbb{R}), a \in (\operatorname{vec}(m) \to \mathbb{C}), b \in \mathbb{C}, x \in \operatorname{vec}(m). \\ \quad \operatorname{is\text{-}module}(m) \wedge \\ \quad \operatorname{scal}(m) = \operatorname{normed\text{-}field\text{-}as\text{-}commutative\text{-}ring}(\operatorname{cc\text{-}normed\text{-}field}) \wedge \\ \quad \forall \mathit{rax}, y \in \operatorname{vec}(m).\; (n(\operatorname{vadd}(m)(\mathit{rax}, y)) = (n(\mathit{rax}) + n(y))) \wedge \\ \quad \forall r \in \mathbb{R}, \mathit{rhx} \in \operatorname{vec}(m).\; (n(\operatorname{act}(m)(r, \mathit{rhx})) = (r \cdot n(\mathit{rhx}))) \wedge \\ \quad \forall \mathit{cvy} \in \operatorname{vec}(m).\; (a(\mathit{cvy}) \simeq (n(\mathit{cvy}) + (-n(\operatorname{act}(m)(+i, \mathit{cvy})) \cdot +i))) \Rightarrow \\ \quad \quad a(\operatorname{act}(m)(b, x)) = (b \cdot a(x)) \end{array}}}
\vnbprop{rb3-cf-klinear}{\vnbstmt{\begin{array}{@{}l@{}} \forall m.; \forall n \in (\operatorname{vec}(m) \to \mathbb{R}), a \in (\operatorname{vec}(m) \to \mathbb{C}). \\ \quad \operatorname{is\text{-}module}(m) \wedge \\ \quad \operatorname{scal}(m) = \operatorname{normed\text{-}field\text{-}as\text{-}commutative\text{-}ring}(\operatorname{cc\text{-}normed\text{-}field}) \wedge \\ \quad \forall x, y \in \operatorname{vec}(m).\; (n(\operatorname{vadd}(m)(x, y)) = (n(x) + n(y))) \wedge \\ \quad \forall r \in \mathbb{R}, \mathit{rhx} \in \operatorname{vec}(m).\; (n(\operatorname{act}(m)(r, \mathit{rhx})) = (r \cdot n(\mathit{rhx}))) \wedge \\ \quad \forall \mathit{cvy} \in \operatorname{vec}(m).\; (a(\mathit{cvy}) \simeq (n(\mathit{cvy}) + (-n(\operatorname{act}(m)(+i, \mathit{cvy})) \cdot +i))) \Rightarrow \\ \quad \quad \operatorname{is\text{-}k\text{-}linear}(m, a) \end{array}}}
\vnbprop{rb3-eq-decomp}{\vnbstmt{\forall z \in \mathbb{C}, w.\; ((w = (\operatorname{real\text{-}part}(z) + (+i \cdot \operatorname{imag\text{-}part}(z)))) \Rightarrow (w = z))}}
\vnbprop{rb3-re-im-i}{\vnbstmt{(\operatorname{real\text{-}part}(+i) = 0) \wedge (\operatorname{imag\text{-}part}(+i) = 1)}}
\vnbprop{rb3-cf-extends}{\vnbstmt{\begin{array}{@{}l@{}} \forall m.; \forall n \in (\operatorname{vec}(m) \to \mathbb{R}), a \in (\operatorname{vec}(m) \to \mathbb{C}), s.; \forall f \in (s \to \mathbb{C}). \\ \quad \operatorname{is\text{-}module}(m) \wedge \\ \quad \operatorname{scal}(m) = \operatorname{normed\text{-}field\text{-}as\text{-}commutative\text{-}ring}(\operatorname{cc\text{-}normed\text{-}field}) \wedge \\ \quad \forall x, y \in \operatorname{vec}(m).\; (n(\operatorname{vadd}(m)(x, y)) = (n(x) + n(y))) \wedge \\ \quad \forall r \in \mathbb{R}, \mathit{rhx} \in \operatorname{vec}(m).\; (n(\operatorname{act}(m)(r, \mathit{rhx})) = (r \cdot n(\mathit{rhx}))) \wedge \\ \quad \forall \mathit{cvy} \in \operatorname{vec}(m).\; (a(\mathit{cvy}) \simeq (n(\mathit{cvy}) + (-n(\operatorname{act}(m)(+i, \mathit{cvy})) \cdot +i))) \wedge \\ \quad \operatorname{is\text{-}submodule}(m, s) \wedge \\ \quad \forall \mathit{eix} \in s.\; (f(\operatorname{act}(m)(+i, \mathit{eix})) = (+i \cdot f(\mathit{eix}))) \wedge \\ \quad \forall \mathit{erx} \in s.\; (n(\mathit{erx}) = \operatorname{real\text{-}part}(f(\mathit{erx}))) \Rightarrow \\ \quad \quad \operatorname{extends\text{-}on}(s, a, f) \end{array}}}
\vnbprop{rb3-mag-sq-conj}{\vnbstmt{\forall z \in \mathbb{C}.\; ((\operatorname{magnitude}(z) \cdot \operatorname{magnitude}(z)) = (z \cdot \operatorname{conjugate}(z)))}}
\vnbprop{rb3-mag-conj}{\vnbstmt{\forall z \in \mathbb{C}.\; (\operatorname{magnitude}(\operatorname{conjugate}(z)) = \operatorname{magnitude}(z))}}
\vnbprop{rb3-cf-bound}{\vnbstmt{\begin{array}{@{}l@{}} \forall m.; \forall n \in (\operatorname{vec}(m) \to \mathbb{R}), a \in (\operatorname{vec}(m) \to \mathbb{C}), p.; \forall y \in \operatorname{vec}(m). \\ \quad \operatorname{is\text{-}module}(m) \wedge \\ \quad \operatorname{scal}(m) = \operatorname{normed\text{-}field\text{-}as\text{-}commutative\text{-}ring}(\operatorname{cc\text{-}normed\text{-}field}) \wedge \\ \quad \forall x, \mathit{ray} \in \operatorname{vec}(m).\; (n(\operatorname{vadd}(m)(x, \mathit{ray})) = (n(x) + n(\mathit{ray}))) \wedge \\ \quad \forall r \in \mathbb{R}, \mathit{rhx} \in \operatorname{vec}(m).\; (n(\operatorname{act}(m)(r, \mathit{rhx})) = (r \cdot n(\mathit{rhx}))) \wedge \\ \quad \forall \mathit{cvy} \in \operatorname{vec}(m).\; (a(\mathit{cvy}) \simeq (n(\mathit{cvy}) + (-n(\operatorname{act}(m)(+i, \mathit{cvy})) \cdot +i))) \wedge \\ \quad \operatorname{is\text{-}seminorm}(\operatorname{cc\text{-}normed\text{-}field}, m, p) \wedge \\ \quad \forall \mathit{bx} \in \operatorname{vec}(m).\; (n(\mathit{bx}) \leq p(\mathit{bx})) \Rightarrow \\ \quad \quad \operatorname{fnrm}(\operatorname{cc\text{-}normed\text{-}field})(a(y)) \leq p(y) \end{array}}}
\vnbprop{hahn-banach-seminorm}{\vnbstmt{\begin{array}{@{}l@{}} \forall m, p, s, f, d. \\ \quad \operatorname{is\text{-}module}(m) \wedge \\ \quad \operatorname{is\text{-}normed\text{-}field}(d) \wedge \\ \quad (d = \operatorname{rr\text{-}normed\text{-}field}) \vee (d = \operatorname{cc\text{-}normed\text{-}field}) \wedge \\ \quad \operatorname{scal}(m) = \operatorname{normed\text{-}field\text{-}as\text{-}commutative\text{-}ring}(d) \wedge \\ \quad \operatorname{is\text{-}seminorm}(d, m, p) \wedge \\ \quad \operatorname{is\text{-}submodule}(m, s) \wedge \\ \quad \operatorname{is\text{-}k\text{-}linear\text{-}on}(m, s, f) \wedge \\ \quad \forall x \in s.\; (\operatorname{fnrm}(d)(f(x)) \leq p(x)) \Rightarrow \\ \quad \quad \exists a. \\ \quad \quad \quad (\operatorname{is\text{-}k\text{-}linear}(m, a) \wedge (\operatorname{extends\text{-}on}(s, a, f) \wedge \forall x \in \operatorname{vec}(m).\; (\operatorname{fnrm}(d)(a(x)) \leq p(x)))) \end{array}}}
\vnbprop{rb3-pgt-pts}{\vnbstmt{\forall m.; \forall g \in \mathbf{Set}.\; (\operatorname{pts}(\operatorname{pseudo\text{-}gauge\text{-}top}(m, g)) \simeq g)}}
\vnbprop{rb3-pgt-opens}{\vnbstmt{\begin{array}{@{}l@{}} \forall m.; \forall g \in \mathbf{Set}. \\ \quad (\operatorname{opens}(\operatorname{pseudo\text{-}gauge\text{-}top}(m, g)) \simeq \{u \in \operatorname{power}(g) : \operatorname{is\text{-}open\text{-}gauge}(m, g, u)\}) \end{array}}}
\vnbprop{rb3-fam-dist-real}{\vnbstmt{\begin{array}{@{}l@{}} \forall m, g.; \forall n \in \mathbb{N}, y, x \in g. \\ \quad \operatorname{is\text{-}countable\text{-}pseudometric\text{-}family}(m, g) \Rightarrow (\operatorname{dist}(m(n))(y, x) \in \mathbb{R}) \end{array}}}
\vnbprop{rb3-fam-pseudo}{\vnbstmt{\begin{array}{@{}l@{}} \forall m, g.; \forall n \in \mathbb{N}. \\ \quad \operatorname{is\text{-}countable\text{-}pseudometric\text{-}family}(m, g) \Rightarrow \operatorname{is\text{-}pseudometric}(\operatorname{dist}(m(n)), g) \end{array}}}
\vnbprop{rb3-og-inter}{\vnbstmt{\begin{array}{@{}l@{}} \forall m, g, u, v. \\ \quad (\operatorname{is\text{-}countable\text{-}pseudometric\text{-}family}(m, g) \wedge (\operatorname{is\text{-}open\text{-}gauge}(m, g, u) \wedge \operatorname{is\text{-}open\text{-}gauge}(m, g, v))) \Rightarrow  \\ \quad \quad \operatorname{is\text{-}open\text{-}gauge}(m, g, (u \cap v)) \end{array}}}
\vnbprop{rb3-og-union}{\vnbstmt{\begin{array}{@{}l@{}} \forall m, g, w. \\ \quad \operatorname{is\text{-}countable\text{-}pseudometric\text{-}family}(m, g) \wedge \\ \quad \forall a \in w.\; \operatorname{is\text{-}open\text{-}gauge}(m, g, a) \wedge \\ \quad \forall a \in w.\; (a \subseteq g) \Rightarrow \\ \quad \quad \operatorname{is\text{-}open\text{-}gauge}(m, g, \operatorname{big\text{-}union}(u, w, u)) \end{array}}}
\vnbprop{rb3-og-empty}{\vnbstmt{\begin{array}{@{}l@{}} \forall m, g. \\ \quad \operatorname{is\text{-}countable\text{-}pseudometric\text{-}family}(m, g) \Rightarrow \operatorname{is\text{-}open\text{-}gauge}(m, g, \emptyset) \end{array}}}
\vnbprop{rb3-og-full}{\vnbstmt{\begin{array}{@{}l@{}} \forall m, g. \\ \quad \operatorname{is\text{-}countable\text{-}pseudometric\text{-}family}(m, g) \Rightarrow \operatorname{is\text{-}open\text{-}gauge}(m, g, g) \end{array}}}
\vnbprop{rb3-gauge-top-is-top-space}{\vnbstmt{\begin{array}{@{}l@{}} \forall m, g. \\ \quad (\operatorname{is\text{-}countable\text{-}pseudometric\text{-}family}(m, g) \wedge (g \in \mathbf{Set})) \Rightarrow  \\ \quad \quad \operatorname{is\text{-}top\text{-}space}(\operatorname{pseudo\text{-}gauge\text{-}top}(m, g)) \end{array}}}
\vnbprop{rb3-og-zero-dist}{\vnbstmt{\begin{array}{@{}l@{}} \forall m, g, u.; \forall y \in u, x \in g. \\ \quad \operatorname{is\text{-}countable\text{-}pseudometric\text{-}family}(m, g) \wedge \\ \quad \operatorname{is\text{-}open\text{-}gauge}(m, g, u) \wedge \\ \quad \forall n \in \mathbb{N}.\; (\operatorname{dist}(m(n))(y, x) = 0) \Rightarrow \\ \quad \quad x \in u \end{array}}}
\vnbprop{rb3-metrizable-hausdorff}{\vnbstmt{\forall s.\; (\operatorname{is\text{-}metrizable\text{-}top\text{-}space}(s) \Rightarrow \operatorname{is\text{-}hausdorff}(s))}}
\vnbprop{rb3-zero-family}{\vnbstmt{\begin{array}{@{}l@{}} \operatorname{is\text{-}countable\text{-}pseudometric\text{-}family}((\lambda n \in \mathbb{N}.\; \langle \mathbb{N}, (\lambda \langle \operatorname{za\_}, \operatorname{zb\_} \rangle \in (\mathbb{N} \times \mathbb{N}).\; 0) \rangle), \\ \quad \mathbb{N}) \end{array}}}
\vnbprop{rb3-zero-gauge-witness}{\vnbstmt{\begin{array}{@{}l@{}} \operatorname{is\text{-}gauge\text{-}countable}(\operatorname{pseudo\text{-}gauge\text{-}top}((\lambda n \in \mathbb{N}.\; \langle \mathbb{N}, (\lambda \langle \operatorname{za\_}, \operatorname{zb\_} \rangle \in (\mathbb{N} \times \mathbb{N}).\; 0) \rangle), \mathbb{N})) \wedge \\ \neg \operatorname{is\text{-}metrizable\text{-}top\text{-}space}(\operatorname{pseudo\text{-}gauge\text{-}top}((\lambda n \in \mathbb{N}.\; \langle \mathbb{N}, (\lambda \langle \operatorname{za\_}, \operatorname{zb\_} \rangle \in (\mathbb{N} \times \mathbb{N}).\; 0) \rangle), \mathbb{N})) \end{array}}}
\vnbprop{gauge-countable-not-metrizable-witness}{\vnbstmt{\exists s.\; (\operatorname{is\text{-}gauge\text{-}countable}(s) \wedge \neg \operatorname{is\text{-}metrizable\text{-}top\text{-}space}(s))}}
\vnbprop{rb3-dual-is-normed-vector-space-false}{\vnbstmt{\neg \forall m.\; (\operatorname{is\text{-}normed\text{-}vector\text{-}space}(m) \Rightarrow \operatorname{is\text{-}normed\text{-}vector\text{-}space}({m}^{*}))}}
\vnbprop{rb3-tuple2}{\vnbstmt{\begin{array}{@{}l@{}} \forall a, b. \\ \quad ((a \in \mathbf{Set}) \wedge (b \in \mathbf{Set})) \Rightarrow ((\operatorname{pts}(\langle a, b \rangle) \simeq a) \wedge (\operatorname{dist}(\langle a, b \rangle) \simeq b)) \end{array}}}
\vnbprop{rb3-pms-of}{\vnbstmt{\begin{array}{@{}l@{}} \forall q, \mathit{dq}. \\ \quad ((q \in \mathbf{Set}) \wedge ((\mathit{dq} \in ((q \times q) \to \mathbb{R})) \wedge \operatorname{is\text{-}pseudometric}(\mathit{dq}, q))) \Rightarrow  \\ \quad \quad \operatorname{is\text{-}pseudometric\text{-}space}(\langle q, \mathit{dq} \rangle) \end{array}}}
\vnbprop{rb3-metric-pseudo}{\vnbstmt{\forall d.\; (\operatorname{is\text{-}metric\text{-}space}(d) \Rightarrow \operatorname{is\text{-}pseudometric}(\operatorname{dist}(d), \operatorname{pts}(d)))}}
\vnbprop{rb3-fam1-dist}{\vnbstmt{\begin{array}{@{}l@{}} \forall d.; \forall n \in \mathbb{N}. \\ \quad \operatorname{is\text{-}metric\text{-}space}(d) \Rightarrow  \\ \quad \quad (\operatorname{dist}((\lambda \mathit{zn} \in \mathbb{N}.\; \langle \operatorname{pts}(d), \operatorname{dist}(d) \rangle)(n)) \simeq \operatorname{dist}(d)) \end{array}}}
\vnbprop{rb3-metric-gauge-iff}{\vnbstmt{\begin{array}{@{}l@{}} \forall d, u. \\ \quad (\operatorname{is\text{-}metric\text{-}space}(d) \wedge (u \subseteq \operatorname{pts}(d))) \Rightarrow  \\ \quad \quad \operatorname{is\text{-}open}(d, u) \Leftrightarrow \operatorname{is\text{-}open\text{-}gauge}((\lambda n \in \mathbb{N}.\; \langle \operatorname{pts}(d), \operatorname{dist}(d) \rangle), \operatorname{pts}(d), u) \end{array}}}
\vnbprop{rb3-mgc-fwd}{\vnbstmt{\begin{array}{@{}l@{}} \forall s. \\ \quad (\operatorname{is\text{-}top\text{-}space}(s) \wedge \operatorname{is\text{-}metrizable\text{-}top\text{-}space}(s)) \Rightarrow  \\ \quad \quad \operatorname{is\text{-}hausdorff}(s) \wedge \operatorname{is\text{-}gauge\text{-}countable}(s) \end{array}}}
\vnbprop{rb3-cm-term-le}{\vnbstmt{\begin{array}{@{}l@{}} \forall s, w, x, y.; \forall n \in \mathbb{N}. \\ \quad (\operatorname{is\text{-}c\text{-}metric\text{-}space}(s) \wedge (\operatorname{summable\text{-}weight}(w) \wedge ((x \in \operatorname{pts}(s)) \wedge (y \in \operatorname{pts}(s))))) \Rightarrow  \\ \quad \quad (w(n) \cdot \operatorname{min}(1, \operatorname{dists}(s)(n)(x, y))) \leq \operatorname{dist}(\operatorname{c\text{-}metric\text{-}w}(s, w))(x, y) \end{array}}}
\vnbprop{rb3-cm-dist-le}{\vnbstmt{\begin{array}{@{}l@{}} \forall s, w, e. \\ \quad \operatorname{is\text{-}c\text{-}metric\text{-}space}(s) \wedge \\ \quad \operatorname{summable\text{-}weight}(w) \wedge \\ \quad \forall k \in \mathbb{N}.\; (\operatorname{series\text{-}partial\text{-}sum}(w, k) \leq 1) \wedge \\ \quad \operatorname{pos\text{-}rr}(e) \Rightarrow \\ \quad \quad \exists d \in \mathbb{N}. \forall x, y \in \operatorname{pts}(s), r \in \mathbb{R}. \\ \quad \quad \quad ((0 < r) \wedge \forall n \in \mathbb{N}.\; ((n \leq d) \Rightarrow (\operatorname{dists}(s)(n)(x, y) < r))) \Rightarrow  \\ \quad \quad \quad \quad \operatorname{dist}(\operatorname{c\text{-}metric\text{-}w}(s, w))(x, y) \leq (r + e) \end{array}}}
\vnbprop{rb3-tuple2d}{\vnbstmt{\begin{array}{@{}l@{}} \forall a, b. \\ \quad ((a \in \mathbf{Set}) \wedge (b \in \mathbf{Set})) \Rightarrow ((\operatorname{pts}(\langle a, b \rangle) \simeq a) \wedge (\operatorname{dists}(\langle a, b \rangle) \simeq b)) \end{array}}}
\vnbprop{rb3-cs-is-cms}{\vnbstmt{\begin{array}{@{}l@{}} \forall m, g. \\ \quad \operatorname{is\text{-}countable\text{-}pseudometric\text{-}family}(m, g) \wedge \\ \quad g \in \mathbf{Set} \wedge \\ \quad \forall u, v \in g. \\ \quad \quad \forall k \in \mathbb{N}.\; (\operatorname{dist}(m(k))(u, v) = 0) \Rightarrow (u = v) \Rightarrow \\ \quad \quad \operatorname{is\text{-}c\text{-}metric\text{-}space}(\langle g, (\lambda \mathit{zk} \in \mathbb{N}.\; \operatorname{dist}(m(\mathit{zk}))) \rangle) \end{array}}}
\vnbprop{rb3-haus-sep}{\vnbstmt{\begin{array}{@{}l@{}} \forall s, m, g.; \forall u, v \in g. \\ \quad \operatorname{is\text{-}hausdorff}(s) \wedge \\ \quad \operatorname{is\text{-}countable\text{-}pseudometric\text{-}family}(m, g) \wedge \\ \quad (\operatorname{pseudo\text{-}gauge\text{-}top}(m, g) \simeq s) \wedge \\ \quad \forall k \in \mathbb{N}.\; (\operatorname{dist}(m(k))(u, v) = 0) \Rightarrow \\ \quad \quad u = v \end{array}}}
\vnbprop{rb3-cm-open-iff}{\vnbstmt{\begin{array}{@{}l@{}} \forall s, w, m, g, u. \\ \quad \operatorname{is\text{-}c\text{-}metric\text{-}space}(s) \wedge \\ \quad \operatorname{summable\text{-}weight}(w) \wedge \\ \quad \forall k \in \mathbb{N}.\; (\operatorname{series\text{-}partial\text{-}sum}(w, k) \leq 1) \wedge \\ \quad \forall n, k \in \mathbb{N}.\; ((n \leq k) \Rightarrow (w(k) \leq w(n))) \wedge \\ \quad (\operatorname{pts}(s) \simeq g) \wedge \\ \quad \forall n \in \mathbb{N}.\; (\operatorname{dists}(s)(n) \simeq \operatorname{dist}(m(n))) \wedge \\ \quad u \subseteq g \Rightarrow \\ \quad \quad \operatorname{is\text{-}open}(\operatorname{c\text{-}metric\text{-}w}(s, w), u) \Leftrightarrow \operatorname{is\text{-}open\text{-}gauge}(m, g, u) \end{array}}}
\vnbprop{rb3-dyadic-le-one}{\vnbstmt{\begin{array}{@{}l@{}} \forall k \in \mathbb{N}. \\ \quad \operatorname{series\text{-}partial\text{-}sum}((\lambda n \in \mathbb{N}.\; /(1, \operatorname{power}(2, (n + 1)))), k) \leq 1 \end{array}}}
\vnbprop{rb3-dyadic-antitone}{\vnbstmt{\begin{array}{@{}l@{}} \forall a, k \in \mathbb{N}. \\ \quad (a \leq k) \Rightarrow  \\ \quad \quad (\lambda n \in \mathbb{N}.\; /(1, \operatorname{power}(2, (n + 1))))(k) \leq (\lambda n \in \mathbb{N}.\; /(1, \operatorname{power}(2, (n + 1))))(a) \end{array}}}
\vnbprop{rb3-mgc-bwd}{\vnbstmt{\begin{array}{@{}l@{}} \forall s, m, g. \\ \quad \operatorname{is\text{-}top\text{-}space}(s) \wedge \\ \quad \operatorname{is\text{-}hausdorff}(s) \wedge \\ \quad \operatorname{is\text{-}countable\text{-}pseudometric\text{-}family}(m, g) \wedge \\ \quad (\operatorname{pseudo\text{-}gauge\text{-}top}(m, g) \simeq s) \Rightarrow \\ \quad \quad \operatorname{is\text{-}metrizable\text{-}top\text{-}space}(s) \end{array}}}
\vnbprop{metrizable-iff-gauge-countable}{\vnbstmt{\begin{array}{@{}l@{}} \forall s. \\ \quad \operatorname{is\text{-}top\text{-}space}(s) \Rightarrow  \\ \quad \quad \operatorname{is\text{-}metrizable\text{-}top\text{-}space}(s) \Leftrightarrow (\operatorname{is\text{-}hausdorff}(s) \wedge \operatorname{is\text{-}gauge\text{-}countable}(s)) \end{array}}}
\vnbprop{bounded-metrizable-iff-gauge-countable}{\vnbstmt{\begin{array}{@{}l@{}} \forall s. \\ \quad \operatorname{is\text{-}top\text{-}space}(s) \Rightarrow  \\ \quad \quad \exists d.\; (\operatorname{is\text{-}bounded\text{-}metric\text{-}space}(d) \wedge (\operatorname{metric\text{-}top}(d) \simeq s)) \\ \quad \quad \Leftrightarrow\; \operatorname{is\text{-}hausdorff}(s) \wedge \operatorname{is\text{-}gauge\text{-}countable}(s) \end{array}}}

\vnbbookfile{structure-library/positive-functional}{structure-library/positive-functional}
\noindent{\small structure-library/positive-functional.scm -- continuous real functions on a metric space, positive linear functionals on them, clopen sets.  DEFINITIONS ONLY (Riesz-Markov, stage 2, RM-2, 2026-10-10); the laws are PROVEN in theorem-library/riesz-markov-cantor.scm (prefix rmc-).}\par
\par\medskip\noindent\textbf{Definition} (\texttt{cont-fns}).\;$\operatorname{cont\text{-}fns}(s) :=$ $\{\operatorname{cfh\_} \in (\operatorname{pts}(s) \to \mathbb{R}) : \operatorname{is\text{-}continuous}(s, \operatorname{rr\text{-}ms}, \operatorname{cfh\_})\}$
\par\medskip\noindent\textbf{Definition} (\texttt{is-positive-functional}).\;For all $s$ and $\mathit{lam}$, 
\begin{equation*}
\begin{array}{@{}l@{}} \operatorname{is\text{-}positive\text{-}functional}(s, \mathit{lam}) \\ \Leftrightarrow\; \mathit{lam} \in (\operatorname{cont\text{-}fns}(s) \to \mathbb{R}) \wedge \\ \forall \operatorname{pff\_}, \operatorname{pfg\_} \in \operatorname{cont\text{-}fns}(s). \\ \quad \operatorname{lam}((\lambda \operatorname{sx\_} \in \operatorname{pts}(s).\; +(\operatorname{pff\_}(\operatorname{sx\_}), \operatorname{pfg\_}(\operatorname{sx\_})))) = +(\operatorname{lam}(\operatorname{pff\_}), \operatorname{lam}(\operatorname{pfg\_})) \wedge \\ \forall \operatorname{pfc\_} \in \mathbb{R}, \operatorname{pff\_} \in \operatorname{cont\text{-}fns}(s). \\ \quad \operatorname{lam}((\lambda \operatorname{sx\_} \in \operatorname{pts}(s).\; \cdot(\operatorname{pfc\_}, \operatorname{pff\_}(\operatorname{sx\_})))) = \cdot(\operatorname{pfc\_}, \operatorname{lam}(\operatorname{pff\_})) \wedge \\ \forall \operatorname{pff\_} \in \operatorname{cont\text{-}fns}(s). \\ \quad \forall \operatorname{pfx\_} \in \operatorname{pts}(s).\; (0 \leq \operatorname{pff\_}(\operatorname{pfx\_})) \Rightarrow (0 \leq \operatorname{lam}(\operatorname{pff\_})) \end{array}
\end{equation*}

\par\medskip\noindent\textbf{Definition} (\texttt{clopen-sets}).\;$\operatorname{clopen\text{-}sets}(s) :=$
\begin{equation*}
\{\operatorname{cou\_} \in \operatorname{power}(\operatorname{pts}(s)) : (\operatorname{is\text{-}open}(s, \operatorname{cou\_}) \wedge \operatorname{is\text{-}closed}(s, \operatorname{cou\_}))\}
\end{equation*}


\vnbbookfile{theorem-library/series-domination}{theorem-library/series-domination}
\noindent{\small series-domination.scm -- SHIFTING AND EVENTUAL DOMINATION FOR REAL SERIES. Batch 26-A, 2026-09-24.  The two lemmas the batch-25 report named as standing in the way of Proposition 2.10 (theorem-library/cc-power-series-deriv.scm):}\par
\vnbprop{series-shift-partial-sum}{\vnbstmt{\begin{array}{@{}l@{}} \forall k \in \mathbb{N}, f, g \in (\mathbb{N} \to \mathbb{R}). \\ \quad \forall j \in \mathbb{N}.\; (g(j) = f((j + 1))) \Rightarrow  \\ \quad \quad \operatorname{series\text{-}partial\text{-}sum}(f, (k + 1)) = (f(0) + \operatorname{series\text{-}partial\text{-}sum}(g, k)) \end{array}}}
\vnbprop{series-shift-converges-iff}{\vnbstmt{\begin{array}{@{}l@{}} \forall f, g \in (\mathbb{N} \to \mathbb{R}). \\ \quad (\forall j \in \mathbb{N}.\; (0 \leq f(j)) \wedge \forall j \in \mathbb{N}.\; (g(j) = f((j + 1)))) \Rightarrow  \\ \quad \quad \operatorname{series\text{-}converges}(f) \Leftrightarrow \operatorname{series\text{-}converges}(g) \end{array}}}
\vnbprop{series-dominated-partial-bound}{\vnbstmt{\begin{array}{@{}l@{}} \forall k \in \mathbb{N}, f, g \in (\mathbb{N} \to \mathbb{R}), n \in \mathbb{N}. \\ \quad \forall j \in \mathbb{N}.\; (0 \leq f(j)) \wedge \\ \quad \forall j \in \mathbb{N}.\; (0 \leq g(j)) \wedge \\ \quad \forall j \in \mathbb{N}.\; ((n \leq j) \Rightarrow (f(j) \leq g(j))) \Rightarrow \\ \quad \quad \operatorname{series\text{-}partial\text{-}sum}(f, k) \leq (\operatorname{series\text{-}partial\text{-}sum}(f, n) + \operatorname{series\text{-}partial\text{-}sum}(g, k)) \end{array}}}
\vnbprop{series-eventually-dominated-converges}{\vnbstmt{\begin{array}{@{}l@{}} \forall f, g \in (\mathbb{N} \to \mathbb{R}), n \in \mathbb{N}. \\ \quad \forall j \in \mathbb{N}.\; (0 \leq f(j)) \wedge \\ \quad \forall j \in \mathbb{N}.\; (0 \leq g(j)) \wedge \\ \quad \forall j \in \mathbb{N}.\; ((n \leq j) \Rightarrow (f(j) \leq g(j))) \wedge \\ \quad \operatorname{series\text{-}converges}(g) \Rightarrow \\ \quad \quad \operatorname{series\text{-}converges}(f) \end{array}}}

\vnbbookfile{theorem-library/regulated-has-primitive}{theorem-library/regulated-has-primitive}
\noindent{\small regulated-has-primitive.scm -- EVERY REGULATED FUNCTION HAS A PRIMITIVE on [a, b] (batch 22, stage 2: the assembly).}\par
\vnbprop{primitive-shift-const}{\vnbstmt{\begin{array}{@{}l@{}} \forall g, f, a, b.; \forall c \in \mathbb{R}. \\ \quad \operatorname{is\text{-}primitive}(g, f, a, b) \Rightarrow  \\ \quad \quad \operatorname{is\text{-}primitive}((\lambda y \in \operatorname{ccint}(a, b).\; (g(y) - c)), f, a, b) \end{array}}}
\vnbprop{primitive-family-normalised-choice}{\vnbstmt{\begin{array}{@{}l@{}} \forall a, b \in \mathbb{R}, f. \\ \quad \forall k \in \mathbb{N}.\; \exists g.\; \operatorname{is\text{-}primitive}(g, f(k), a, b) \Rightarrow  \\ \quad \quad \exists \mathit{rhpgf} \in (\mathbb{N} \to (\operatorname{ccint}(a, b) \to \mathbb{R})). \forall k \in \mathbb{N}. \\ \quad \quad \quad \operatorname{is\text{-}primitive}(\operatorname{rhpgf}(k), f(k), a, b) \wedge (\operatorname{rhpgf}(k)(a) = 0) \end{array}}}
\vnbprop{regulated-on-has-primitive}{\vnbstmt{\begin{array}{@{}l@{}} \forall a, b \in \mathbb{R}, f. \\ \quad \operatorname{is\text{-}regulated\text{-}on}(f, a, b) \Rightarrow \exists g.\; \operatorname{is\text{-}primitive}(g, f, a, b) \end{array}}}

\vnbbookfile{theorem-library/road-laws-2}{theorem-library/road-laws-2}
\noindent{\small road-laws-2.scm -- THE OPPOSITE ROAD AND DIEUDONNE (9.6.1).}\par
\vnbprop{ccint-reflect-in}{\vnbstmt{\forall a, b \in \mathbb{R}, y \in \operatorname{ccint}(a, b).\; (((a + b) - y) \in \operatorname{ccint}(a, b))}}
\vnbprop{ooint-reflect-in}{\vnbstmt{\forall a, b \in \mathbb{R}, y \in \operatorname{ooint}(a, b).\; (((a + b) - y) \in \operatorname{ooint}(a, b))}}
\vnbprop{reflect-continuous-at-sub}{\vnbstmt{\begin{array}{@{}l@{}} \forall a, b \in \mathbb{R}, f.; \forall c \in \mathbb{R}, q \in (\operatorname{ccint}(a, b) \to \mathbb{R}), t \in \operatorname{ccint}(a, b). \\ \quad \forall y \in \operatorname{ccint}(a, b).\; (q(y) \simeq (c \cdot f(((a + b) - y)))) \wedge \\ \quad \operatorname{is\text{-}continuous\text{-}at}(\operatorname{subspace\text{-}ms}(\operatorname{rr\text{-}ms}, \operatorname{ccint}(a, b)), \operatorname{rr\text{-}ms}, f, ((a + b) - t)) \Rightarrow \\ \quad \quad \operatorname{is\text{-}continuous\text{-}at}(\operatorname{subspace\text{-}ms}(\operatorname{rr\text{-}ms}, \operatorname{ccint}(a, b)), \operatorname{rr\text{-}ms}, q, t) \end{array}}}
\vnbprop{reflect-continuous-on}{\vnbstmt{\begin{array}{@{}l@{}} \forall a, b \in \mathbb{R}, f.; \forall c \in \mathbb{R}, q \in (\operatorname{ccint}(a, b) \to \mathbb{R}). \\ \quad (\operatorname{is\text{-}continuous\text{-}on}(f, \operatorname{ccint}(a, b)) \wedge \forall y \in \operatorname{ccint}(a, b).\; (q(y) \simeq (c \cdot f(((a + b) - y))))) \Rightarrow  \\ \quad \quad \operatorname{is\text{-}continuous\text{-}on}(q, \operatorname{ccint}(a, b)) \end{array}}}
\vnbprop{regulated-on-reflect}{\vnbstmt{\begin{array}{@{}l@{}} \forall a, b \in \mathbb{R}, f.; \forall c \in \mathbb{R}, q \in (\operatorname{ccint}(a, b) \to \mathbb{R}). \\ \quad (\operatorname{is\text{-}regulated\text{-}on}(f, a, b) \wedge \forall y \in \operatorname{ccint}(a, b).\; (q(y) \simeq (c \cdot f(((a + b) - y))))) \Rightarrow  \\ \quad \quad \operatorname{is\text{-}regulated\text{-}on}(q, a, b) \end{array}}}
\vnbprop{reflect-image-finite}{\vnbstmt{\begin{array}{@{}l@{}} \forall a, b.; \forall d \in \mathbf{Set}. \\ \quad ((a \in \mathbb{R}) \wedge ((b \in \mathbb{R}) \wedge ((|d| \in \mathbb{N}) \wedge (d \subseteq \operatorname{ccint}(a, b))))) \Rightarrow  \\ \quad \quad \operatorname{image}((\lambda x \in \operatorname{ccint}(a, b).\; ((a + b) - x)), d) \in \mathbf{Set} \wedge \\ \quad \quad |\operatorname{image}((\lambda x \in \operatorname{ccint}(a, b).\; ((a + b) - x)), d)| \in \mathbb{N} \wedge \\ \quad \quad \operatorname{image}((\lambda x \in \operatorname{ccint}(a, b).\; ((a + b) - x)), d) \subseteq \operatorname{ccint}(a, b) \end{array}}}
\vnbprop{reflect-image-countable}{\vnbstmt{\begin{array}{@{}l@{}} \forall a, b, d. \\ \quad ((a \in \mathbb{R}) \wedge ((b \in \mathbb{R}) \wedge (\operatorname{is\text{-}countable}(d) \wedge (d \subseteq \operatorname{ccint}(a, b))))) \Rightarrow  \\ \quad \quad \operatorname{is\text{-}countable}(\operatorname{image}((\lambda x \in \operatorname{ccint}(a, b).\; ((a + b) - x)), d)) \wedge \\ \quad \quad \operatorname{image}((\lambda x \in \operatorname{ccint}(a, b).\; ((a + b) - x)), d) \subseteq \operatorname{ccint}(a, b) \end{array}}}
\vnbprop{reflect-not-in-image}{\vnbstmt{\begin{array}{@{}l@{}} \forall a, b, d.; \forall t \in \operatorname{ccint}(a, b). \\ \quad a \in \mathbb{R} \wedge \\ \quad b \in \mathbb{R} \wedge \\ \quad d \subseteq \operatorname{ccint}(a, b) \wedge \\ \quad \neg (t \in \operatorname{image}((\lambda x \in \operatorname{ccint}(a, b).\; ((a + b) - x)), d)) \Rightarrow \\ \quad \quad \neg (((a + b) - t) \in d) \end{array}}}
\vnbprop{pw-antiderivative-reflect}{\vnbstmt{\begin{array}{@{}l@{}} \forall a, b \in \mathbb{R}, f, i.; \forall g, k \in (\operatorname{ccint}(a, b) \to \mathbb{R}). \\ \quad \operatorname{is\text{-}primitive}(f, i, a, b) \wedge \\ \quad \forall y \in \operatorname{ccint}(a, b).\; (g(y) \simeq f(((a + b) - y))) \wedge \\ \quad \forall y \in \operatorname{ccint}(a, b).\; (k(y) \simeq (-1 \cdot i(((a + b) - y)))) \Rightarrow \\ \quad \quad \operatorname{is\text{-}primitive}(g, k, a, b) \end{array}}}
\vnbprop{pw-int-reflect}{\vnbstmt{\begin{array}{@{}l@{}} \forall a, b \in \mathbb{R}, f, i.; \forall h, s \in (\operatorname{ccint}(a, b) \to \mathbb{R}). \\ \quad \operatorname{is\text{-}primitive}(f, i, a, b) \wedge \\ \quad \forall y \in \operatorname{ccint}(a, b).\; (h(y) \simeq (-1 \cdot f(((a + b) - y)))) \wedge \\ \quad \forall y \in \operatorname{ccint}(a, b).\; (s(y) \simeq i(((a + b) - y))) \Rightarrow \\ \quad \quad \operatorname{is\text{-}primitive}(h, s, a, b) \wedge (\int_{a}^{b} s = \int_{a}^{b} i) \end{array}}}
\vnbprop{pw-continuous-reflect}{\vnbstmt{\begin{array}{@{}l@{}} \forall a, b \in \mathbb{R}, i.; \forall c \in \mathbb{R}, q \in (\operatorname{ccint}(a, b) \to \mathbb{R}). \\ \quad (\operatorname{is\text{-}pw\text{-}continuous\text{-}on}(i, a, b) \wedge \forall y \in \operatorname{ccint}(a, b).\; (q(y) \simeq (c \cdot i(((a + b) - y))))) \Rightarrow  \\ \quad \quad \operatorname{is\text{-}pw\text{-}continuous\text{-}on}(q, a, b) \end{array}}}
\vnbprop{trace-opposite-subset}{\vnbstmt{\begin{array}{@{}l@{}} \forall a, b \in \mathbb{R}, m, \mathit{prgb}. \\ \quad m \in (\operatorname{ccint}(a, b) \to \mathbb{C}) \wedge \\ \quad \mathit{prgb} \in (\operatorname{ccint}(a, b) \to \mathbb{C}) \wedge \\ \quad \forall y \in \operatorname{ccint}(a, b).\; (\operatorname{prgb}(y) \simeq m(((a + b) - y))) \Rightarrow \\ \quad \quad \operatorname{trace}(\mathit{prgb}, a, b) \subseteq \operatorname{trace}(m, a, b) \end{array}}}
\vnbprop{trace-opposite}{\vnbstmt{\begin{array}{@{}l@{}} \forall a, b \in \mathbb{R}, m, \mathit{prgb}. \\ \quad m \in (\operatorname{ccint}(a, b) \to \mathbb{C}) \wedge \\ \quad \mathit{prgb} \in (\operatorname{ccint}(a, b) \to \mathbb{C}) \wedge \\ \quad \forall y \in \operatorname{ccint}(a, b).\; (\operatorname{prgb}(y) \simeq m(((a + b) - y))) \Rightarrow \\ \quad \quad ((\operatorname{trace}(\mathit{prgb}, a, b) \subseteq \operatorname{trace}(m, a, b)) \wedge (\operatorname{trace}(m, a, b) \subseteq \operatorname{trace}(\mathit{prgb}, a, b))) \end{array}}}
\vnbprop{opposite-is-road}{\vnbstmt{\begin{array}{@{}l@{}} \forall a, b \in \mathbb{R}, m, c, \mathit{prgb}, \mathit{prdb}. \\ \quad \operatorname{is\text{-}road}(m, c, a, b) \wedge \\ \quad \mathit{prgb} \in (\operatorname{ccint}(a, b) \to \mathbb{C}) \wedge \\ \quad \mathit{prdb} \in (\operatorname{ccint}(a, b) \to \mathbb{C}) \wedge \\ \quad \forall y \in \operatorname{ccint}(a, b).\; (\operatorname{prgb}(y) \simeq m(((a + b) - y))) \wedge \\ \quad \forall y \in \operatorname{ccint}(a, b).\; (\operatorname{prdb}(y) \simeq (-1 \cdot c(((a + b) - y)))) \Rightarrow \\ \quad \quad \operatorname{is\text{-}road}(\mathit{prgb}, \mathit{prdb}, a, b) \end{array}}}
\vnbprop{cc-int-reflect}{\vnbstmt{\begin{array}{@{}l@{}} \forall a, b, i, \mathit{prb}, f, w. \\ \quad a \in \mathbb{R} \wedge \\ \quad b \in \mathbb{R} \wedge \\ \quad i \in (\operatorname{ccint}(a, b) \to \mathbb{C}) \wedge \\ \quad \mathit{prb} \in (\operatorname{ccint}(a, b) \to \mathbb{C}) \wedge \\ \quad \forall y \in \operatorname{ccint}(a, b).\; (\operatorname{prb}(y) \simeq i(((a + b) - y))) \wedge \\ \quad \operatorname{is\text{-}primitive}(f, (\lambda t \in \operatorname{ccint}(a, b).\; \operatorname{real\text{-}part}(i(t))), a, b) \wedge \\ \quad \operatorname{is\text{-}primitive}(w, (\lambda t \in \operatorname{ccint}(a, b).\; \operatorname{imag\text{-}part}(i(t))), a, b) \Rightarrow \\ \quad \quad \int_{a}^{b} \mathit{prb} = \int_{a}^{b} i \end{array}}}
\vnbprop{line-int-opposite}{\vnbstmt{\begin{array}{@{}l@{}} \forall a, b, m, c, f, d, \mathit{prgb}, \mathit{prdb}, \mathit{pwf}, w. \\ \quad a \in \mathbb{R} \wedge \\ \quad b \in \mathbb{R} \wedge \\ \quad \operatorname{is\text{-}road}(m, c, a, b) \wedge \\ \quad f \in (d \to \mathbb{C}) \wedge \\ \quad \operatorname{trace}(m, a, b) \subseteq d \wedge \\ \quad \mathit{prgb} \in (\operatorname{ccint}(a, b) \to \mathbb{C}) \wedge \\ \quad \mathit{prdb} \in (\operatorname{ccint}(a, b) \to \mathbb{C}) \wedge \\ \quad \forall y \in \operatorname{ccint}(a, b).\; (\operatorname{prgb}(y) \simeq m(((a + b) - y))) \wedge \\ \quad \forall y \in \operatorname{ccint}(a, b).\; (\operatorname{prdb}(y) \simeq (-1 \cdot c(((a + b) - y)))) \wedge \\ \quad \operatorname{is\text{-}primitive}(\mathit{pwf}, \\ \quad \quad (\lambda t \in \operatorname{ccint}(a, b).\; \operatorname{real\text{-}part}((\lambda t \in \operatorname{ccint}(a, b).\; (f(m(t)) \cdot c(t)))(t))), \\ \quad \quad a, b) \wedge \\ \quad \operatorname{is\text{-}primitive}(w, \\ \quad \quad (\lambda t \in \operatorname{ccint}(a, b).\; \operatorname{imag\text{-}part}((\lambda t \in \operatorname{ccint}(a, b).\; (f(m(t)) \cdot c(t)))(t))), \\ \quad \quad a, b) \Rightarrow \\ \quad \quad \operatorname{line\text{-}int}(f, \mathit{prgb}, \mathit{prdb}, a, b) = (-1 \cdot \operatorname{line\text{-}int}(f, m, c, a, b)) \end{array}}}

\vnbbookfile{theorem-library/compact-image}{theorem-library/compact-image}
\noindent{\small compact-image.scm -- THE CONTINUOUS IMAGE OF A COMPACT SPACE IS COMPACT (calculus.pdf, Proposition 3.19), and the COMPACTNESS OF THE TRACE of a path, which follows from it (complex-analysis.pdf 3.1.1: "gamma\_* is sometimes called the trace of gamma.  It is always a compact set").}\par
\vnbprop{preimage-in-power}{\vnbstmt{\forall s, f, u.\; ((\operatorname{pts}(s) \in \mathbf{Set}) \Rightarrow (\operatorname{preimage}(s, f, u) \in \operatorname{power}(\operatorname{pts}(s))))}}
\vnbprop{image-of-preimage-onto}{\vnbstmt{\begin{array}{@{}l@{}} \forall s, f, u. \\ \quad ((\operatorname{pts}(s) \in \mathbf{Set}) \wedge (u \subseteq \operatorname{image}(f, \operatorname{pts}(s)))) \Rightarrow  \\ \quad \quad \operatorname{image}(f, \operatorname{preimage}(s, f, u)) = u \end{array}}}
\vnbprop{preimage-subspace-open}{\vnbstmt{\begin{array}{@{}l@{}} \forall s, t, f, b, u. \\ \quad \operatorname{is\text{-}continuous}(s, t, f) \wedge \\ \quad b \subseteq \operatorname{pts}(t) \wedge \\ \quad \forall p \in \operatorname{pts}(s).\; (f(p) \in b) \wedge \\ \quad \operatorname{is\text{-}open}(\operatorname{subspace\text{-}ms}(t, b), u) \Rightarrow \\ \quad \quad \operatorname{is\text{-}open}(s, \operatorname{preimage}(s, f, u)) \end{array}}}
\vnbprop{compact-image}{\vnbstmt{\begin{array}{@{}l@{}} \forall s, t, f. \\ \quad (\operatorname{is\text{-}continuous}(s, t, f) \wedge \operatorname{is\text{-}compact}(s)) \Rightarrow  \\ \quad \quad \operatorname{is\text{-}compact}(\operatorname{subspace\text{-}ms}(t, \operatorname{image}(f, \operatorname{pts}(s)))) \end{array}}}
\vnbprop{image-restrict}{\vnbstmt{\forall a \in \mathbf{Set}, f.\; (\operatorname{image}(\operatorname{restrict}(f, a), a) = \operatorname{image}(f, a))}}
\vnbprop{continuous-image-of-compact}{\vnbstmt{\begin{array}{@{}l@{}} \forall s, t, f, a. \\ \quad (\operatorname{is\text{-}continuous}(s, t, f) \wedge ((a \subseteq \operatorname{pts}(s)) \wedge \operatorname{is\text{-}compact}(\operatorname{subspace\text{-}ms}(s, a)))) \Rightarrow  \\ \quad \quad \operatorname{is\text{-}compact}(\operatorname{subspace\text{-}ms}(t, \operatorname{image}(f, a))) \end{array}}}
\vnbprop{path-trace-compact}{\vnbstmt{\begin{array}{@{}l@{}} \forall m, v, \mathit{bv}. \\ \quad \operatorname{is\text{-}path}(m, v, \mathit{bv}) \Rightarrow \operatorname{is\text{-}compact}(\operatorname{subspace\text{-}ms}(\operatorname{cc\text{-}ms}, \operatorname{trace}(m, v, \mathit{bv}))) \end{array}}}

\vnbbookfile{theorem-library/pw-int-order}{theorem-library/pw-int-order}
\noindent{\small pw-int-order.scm -- ORDER PROPERTIES OF THE PIECEWISE INTEGRAL.}\par
\vnbprop{nondecreasing-total}{\vnbstmt{\begin{array}{@{}l@{}} \forall h \in (\mathbb{R} \to \mathbb{R}), u, v \in \mathbb{R}. \\ \quad \forall x \in \mathbb{R}.\; \operatorname{is\text{-}continuous\text{-}at}(\operatorname{rr\text{-}ms}, \operatorname{rr\text{-}ms}, h, x) \wedge \\ \quad u < v \wedge \\ \quad \forall t \in \operatorname{ooint}(u, v).\; \exists l.\; (\operatorname{is\text{-}diff\text{-}at}(h, t, l) \wedge (0 \leq l)) \Rightarrow \\ \quad \quad h(u) \leq h(v) \end{array}}}
\vnbprop{nondecreasing-off-finite-set}{\vnbstmt{\begin{array}{@{}l@{}} \forall s \in \mathbf{Set}, h \in (\mathbb{R} \to \mathbb{R}), u, v \in \mathbb{R}. \\ \quad |s| \in \mathbb{N} \wedge \\ \quad \forall x \in \mathbb{R}.\; \operatorname{is\text{-}continuous\text{-}at}(\operatorname{rr\text{-}ms}, \operatorname{rr\text{-}ms}, h, x) \wedge \\ \quad u < v \wedge \\ \quad \forall t \in \operatorname{ooint}(u, v). \\ \quad \quad \neg (t \in s) \Rightarrow \exists l.\; (\operatorname{is\text{-}diff\text{-}at}(h, t, l) \wedge (0 \leq l)) \Rightarrow \\ \quad \quad h(u) \leq h(v) \end{array}}}
\vnbprop{pw-nondecreasing-on-interval}{\vnbstmt{\begin{array}{@{}l@{}} \forall a, b \in \mathbb{R}, h \in (\operatorname{ccint}(a, b) \to \mathbb{R}), s \in \mathbf{Set}. \\ \quad a < b \wedge \\ \quad \operatorname{is\text{-}continuous\text{-}on}(h, \operatorname{ccint}(a, b)) \wedge \\ \quad |s| \in \mathbb{N} \wedge \\ \quad \forall t \in \operatorname{ooint}(a, b). \\ \quad \quad \neg (t \in s) \Rightarrow \exists l.\; (\operatorname{has\text{-}deriv\text{-}at}(h, t, l) \wedge (0 \leq l)) \Rightarrow \\ \quad \quad h(a) \leq h(b) \end{array}}}
\vnbprop{nondecreasing-off-countable}{\vnbstmt{\begin{array}{@{}l@{}} \forall h \in (\mathbb{R} \to \mathbb{R}), u, v \in \mathbb{R}, s. \\ \quad u < v \wedge \\ \quad \forall x \in \operatorname{ccint}(u, v).\; \operatorname{is\text{-}continuous\text{-}at}(\operatorname{rr\text{-}ms}, \operatorname{rr\text{-}ms}, h, x) \wedge \\ \quad \operatorname{is\text{-}countable}(s) \wedge \\ \quad s \subseteq \mathbb{R} \wedge \\ \quad \forall t \in \operatorname{ooint}(u, v). \\ \quad \quad \neg (t \in s) \Rightarrow \exists l.\; (\operatorname{is\text{-}diff\text{-}at}(h, t, l) \wedge (0 \leq l)) \Rightarrow \\ \quad \quad h(u) \leq h(v) \end{array}}}
\vnbprop{nondecreasing-off-countable-on-interval}{\vnbstmt{\begin{array}{@{}l@{}} \forall a, b \in \mathbb{R}, h \in (\operatorname{ccint}(a, b) \to \mathbb{R}), s. \\ \quad a < b \wedge \\ \quad \operatorname{is\text{-}continuous\text{-}on}(h, \operatorname{ccint}(a, b)) \wedge \\ \quad \operatorname{is\text{-}countable}(s) \wedge \\ \quad s \subseteq \operatorname{ccint}(a, b) \wedge \\ \quad \forall t \in \operatorname{ooint}(a, b). \\ \quad \quad \neg (t \in s) \Rightarrow \exists l.\; (\operatorname{has\text{-}deriv\text{-}at}(h, t, l) \wedge (0 \leq l)) \Rightarrow \\ \quad \quad h(a) \leq h(b) \end{array}}}
\vnbprop{pw-int-nonneg}{\vnbstmt{\begin{array}{@{}l@{}} \forall f, i, a, b. \\ \quad (\operatorname{is\text{-}primitive}(f, i, a, b) \wedge \forall y \in \operatorname{ccint}(a, b).\; (0 \leq i(y))) \Rightarrow  \\ \quad \quad 0 \leq \int_{a}^{b} i \end{array}}}
\vnbprop{pw-int-monotone}{\vnbstmt{\begin{array}{@{}l@{}} \forall f, i, w, \mathit{ppsi}, a, b. \\ \quad \operatorname{is\text{-}primitive}(f, i, a, b) \wedge \\ \quad \operatorname{is\text{-}primitive}(w, \mathit{ppsi}, a, b) \wedge \\ \quad \forall y \in \operatorname{ccint}(a, b).\; (i(y) \leq \operatorname{ppsi}(y)) \Rightarrow \\ \quad \quad \int_{a}^{b} i \leq \int_{a}^{b} \mathit{ppsi} \end{array}}}
\vnbprop{pw-int-abs-bound}{\vnbstmt{\begin{array}{@{}l@{}} \forall f, i, w, s, a, b. \\ \quad \operatorname{is\text{-}primitive}(f, i, a, b) \wedge \\ \quad \operatorname{is\text{-}primitive}(w, s, a, b) \wedge \\ \quad \forall y \in \operatorname{ccint}(a, b).\; (s(y) \simeq \lvert i(y) \rvert) \Rightarrow \\ \quad \quad \lvert \int_{a}^{b} i \rvert \leq \int_{a}^{b} s \end{array}}}

\part{Complex Analysis}
\noindent Holomorphic functions and power series over CC; the exponential and the logarithm; line integrals, Goursat and the Cauchy integral formula; winding numbers; Cauchy's estimates, Liouville and the fundamental theorem of algebra; Morera; the Gamma function.
\par\noindent\textit{Builds on:} Analysis.
\par\noindent{\small 43 files.}

\vnbbookfile{theorem-library/cc-series}{theorem-library/cc-series}
\noindent{\small cc-series.scm -- THE COMPLEX SERIES LAYER: ordered partial sums in CC, and their convergence in CC-MS.}\par
\par\medskip\noindent\textbf{Definition} (\texttt{cc-series-partial-sum}).\;$\operatorname{cc\text{-}series\text{-}partial\text{-}sum}(f, k) :=$ $\operatorname{sum\text{-}ag}(\operatorname{normed\text{-}field\text{-}additive\text{-}ag}(\operatorname{cc\text{-}normed\text{-}field}), f, k)$
\par\medskip\noindent\textbf{Definition} (\texttt{cc-series-converges-to}).\;For all $f$ and $l$, 
\begin{equation*}
\begin{array}{@{}l@{}} \operatorname{cc\text{-}series\text{-}converges\text{-}to}(f, l) \\ \Leftrightarrow\; \operatorname{converges\text{-}to}(\operatorname{cc\text{-}ms}, (\lambda k \in \mathbb{N}.\; \operatorname{cc\text{-}series\text{-}partial\text{-}sum}(f, k)), l) \end{array}
\end{equation*}

\par\medskip\noindent\textbf{Definition} (\texttt{cc-series-converges}).\;For every $f$, 
\begin{equation*}
\operatorname{cc\text{-}series\text{-}converges}(f) \Leftrightarrow \operatorname{converges}(\operatorname{cc\text{-}ms}, (\lambda k \in \mathbb{N}.\; \operatorname{cc\text{-}series\text{-}partial\text{-}sum}(f, k)))
\end{equation*}

\vnbprop{cc-additive-ag-opr}{\vnbstmt{(\operatorname{opr}(\operatorname{normed\text{-}field\text{-}additive\text{-}ag}(\operatorname{cc\text{-}normed\text{-}field})) \simeq (\lambda \langle \operatorname{x\_}, \operatorname{y\_} \rangle \in (\mathbb{C} \times \mathbb{C}).\; (\operatorname{x\_} + \operatorname{y\_})))}}
\vnbprop{cc-additive-ag-iden}{\vnbstmt{(\operatorname{iden}(\operatorname{normed\text{-}field\text{-}additive\text{-}ag}(\operatorname{cc\text{-}normed\text{-}field})) \simeq 0)}}
\vnbprop{cc-series-partial-sum-zero}{\vnbstmt{\forall f.\; (\operatorname{cc\text{-}series\text{-}partial\text{-}sum}(f, 0) \simeq 0)}}
\vnbprop{sum-ag-cc-in-cc}{\vnbstmt{\begin{array}{@{}l@{}} \forall k \in \mathbb{N}, f \in (\mathbb{N} \to \mathbb{C}). \\ \quad \operatorname{sum\text{-}ag}(\operatorname{normed\text{-}field\text{-}additive\text{-}ag}(\operatorname{cc\text{-}normed\text{-}field}), f, k) \in \mathbb{C} \end{array}}}
\vnbprop{cc-series-partial-sum-in-cc}{\vnbstmt{\forall k \in \mathbb{N}, f \in (\mathbb{N} \to \mathbb{C}).\; (\operatorname{cc\text{-}series\text{-}partial\text{-}sum}(f, k) \in \mathbb{C})}}
\vnbprop{cc-series-partial-sum-unfold}{\vnbstmt{\begin{array}{@{}l@{}} \forall f, k. \\ \quad (\operatorname{cc\text{-}series\text{-}partial\text{-}sum}(f, k) \simeq \operatorname{sum\text{-}ag}(\operatorname{normed\text{-}field\text{-}additive\text{-}ag}(\operatorname{cc\text{-}normed\text{-}field}), f, k)) \end{array}}}
\vnbprop{cc-series-partial-sum-succ}{\vnbstmt{\begin{array}{@{}l@{}} \forall f.; \forall k \in \mathbb{N}. \\ \quad ((\operatorname{cc\text{-}series\text{-}partial\text{-}sum}(f, k) \in \mathbb{C}) \wedge (f(k) \in \mathbb{C})) \Rightarrow  \\ \quad \quad (\operatorname{cc\text{-}series\text{-}partial\text{-}sum}(f, (k + 1)) \simeq (\operatorname{cc\text{-}series\text{-}partial\text{-}sum}(f, k) + f(k))) \end{array}}}
\vnbprop{cc-magnitude-zero}{\vnbstmt{\operatorname{magnitude}(0) = 0}}
\vnbprop{cc-series-partial-sum-magnitude-le}{\vnbstmt{\begin{array}{@{}l@{}} \forall k \in \mathbb{N}, f \in (\mathbb{N} \to \mathbb{C}), g \in (\mathbb{N} \to \mathbb{R}). \\ \quad \forall j \in \mathbb{N}.\; (g(j) = \operatorname{magnitude}(f(j))) \Rightarrow  \\ \quad \quad \operatorname{magnitude}(\operatorname{cc\text{-}series\text{-}partial\text{-}sum}(f, k)) \leq \operatorname{series\text{-}partial\text{-}sum}(g, k) \end{array}}}

\vnbbookfile{theorem-library/cc-diff-laws-2}{theorem-library/cc-diff-laws-2}
\noindent{\small cc-diff-laws-2.scm -- the rest of the elementary theory of IS-CC-DIFF-AT, the derivative of a map RR -$>$ CC (structure-library/cc-coords.scm). Continues theorem-library/cc-coords-laws.scm (batch 14-A), which defined the predicate coordinatewise and proved the transfer equivalence}\par
\vnbprop{cc-diff-at-implies-continuous}{\vnbstmt{\begin{array}{@{}l@{}} \forall g \in (\mathbb{R} \to \mathbb{C}), v \in \mathbb{R}, a \in \mathbb{C}. \\ \quad \operatorname{is\text{-}cc\text{-}diff\text{-}at}(g, v, a) \Rightarrow \operatorname{is\text{-}continuous\text{-}at}(\operatorname{rr\text{-}ms}, \operatorname{cc\text{-}ms}, g, v) \end{array}}}
\vnbprop{cc-diff-at-const}{\vnbstmt{\begin{array}{@{}l@{}} \forall v \in \mathbb{C}, g \in (\mathbb{R} \to \mathbb{C}), \mathit{tv} \in \mathbb{R}. \\ \quad \forall x \in \mathbb{R}.\; (g(x) = v) \Rightarrow \operatorname{is\text{-}cc\text{-}diff\text{-}at}(g, \mathit{tv}, 0) \end{array}}}
\vnbprop{cc-diff-at-cc-mul}{\vnbstmt{\begin{array}{@{}l@{}} \forall v \in \mathbb{C}, g, d \in (\mathbb{R} \to \mathbb{C}), \mathit{tv} \in \mathbb{R}, a \in \mathbb{C}. \\ \quad (\forall x \in \mathbb{R}.\; (d(x) = (v \cdot g(x))) \wedge \operatorname{is\text{-}cc\text{-}diff\text{-}at}(g, \mathit{tv}, a)) \Rightarrow  \\ \quad \quad \operatorname{is\text{-}cc\text{-}diff\text{-}at}(d, \mathit{tv}, (v \cdot a)) \end{array}}}
\vnbprop{cc-diff-at-difference}{\vnbstmt{\begin{array}{@{}l@{}} \forall g, h, b \in (\mathbb{R} \to \mathbb{C}), v \in \mathbb{R}, a, c \in \mathbb{C}. \\ \quad (\forall x \in \mathbb{R}.\; (b(x) = (g(x) - h(x))) \wedge (\operatorname{is\text{-}cc\text{-}diff\text{-}at}(g, v, a) \wedge \operatorname{is\text{-}cc\text{-}diff\text{-}at}(h, v, c))) \Rightarrow  \\ \quad \quad \operatorname{is\text{-}cc\text{-}diff\text{-}at}(b, v, (a - c)) \end{array}}}
\vnbprop{cc-diff-at-unique}{\vnbstmt{\begin{array}{@{}l@{}} \forall g \in (\mathbb{R} \to \mathbb{C}), v \in \mathbb{R}, a, b \in \mathbb{C}. \\ \quad (\operatorname{is\text{-}cc\text{-}diff\text{-}at}(g, v, a) \wedge \operatorname{is\text{-}cc\text{-}diff\text{-}at}(g, v, b)) \Rightarrow (a = b) \end{array}}}
\vnbprop{cc-diff-at-chain-real}{\vnbstmt{\begin{array}{@{}l@{}} \forall g \in (\mathbb{R} \to \mathbb{C}), h \in (\mathbb{R} \to \mathbb{R}), m \in (\mathbb{R} \to \mathbb{C}), v, a \in \mathbb{R}, b \in \mathbb{C}. \\ \quad (\forall x \in \mathbb{R}.\; (m(x) = g(h(x))) \wedge (\operatorname{is\text{-}diff\text{-}at}(h, v, a) \wedge \operatorname{is\text{-}cc\text{-}diff\text{-}at}(g, h(v), b))) \Rightarrow  \\ \quad \quad \operatorname{is\text{-}cc\text{-}diff\text{-}at}(m, v, (a \cdot b)) \end{array}}}

\vnbbookfile{structure-library/cc-power-series}{structure-library/cc-power-series}
\noindent{\small cc-power-series.scm -- THE VOCABULARY OF COMPLEX POWER SERIES (the notes 2.2). DEFINITIONS ONLY; the laws are PROVEN in theorem-library/cc-power-series-laws.scm (and the root / ratio tests they rest on in theorem-library/limsup-tests.scm). Batch 25-A, 2026-09-24.  Every definition below is a DECISION FOR THE USER.}\par
\par\medskip\noindent\textbf{Definition} (\texttt{cps-radius}).\;$\operatorname{cps\text{-}radius}(\mathit{cf}) :=$
\begin{equation*}
\operatorname{esup}(\{\operatorname{cpsr\_} \in \mathbb{R} : ((0 \leq \operatorname{cpsr\_}) \wedge \operatorname{series\text{-}converges}((\lambda \operatorname{cpsk\_} \in \mathbb{N}.\; (\operatorname{magnitude}(\operatorname{cf}(\operatorname{cpsk\_})) \cdot \operatorname{power}(\operatorname{cpsr\_}, \operatorname{cpsk\_})))))\})
\end{equation*}

\par\medskip\noindent\textbf{Definition} (\texttt{cc-sup-norm}).\;$\operatorname{cc\text{-}sup\text{-}norm}(u, \operatorname{cpsxs\_}) :=$
\begin{equation*}
\operatorname{esup}(\operatorname{image}((\lambda \operatorname{cpsx\_} \in \operatorname{cpsxs\_}.\; \operatorname{magnitude}(u(\operatorname{cpsx\_}))), \operatorname{cpsxs\_}))
\end{equation*}

\par\medskip\noindent\textbf{Definition} (\texttt{is-normally-convergent}).\;For all $u$ and $\operatorname{cpsxs\_}$, 
\begin{equation*}
\begin{array}{@{}l@{}} \operatorname{is\text{-}normally\text{-}convergent}(u, \operatorname{cpsxs\_}) \\ \Leftrightarrow\; \forall \operatorname{cpsi\_} \in \mathbb{N}.\; (\operatorname{cc\text{-}sup\text{-}norm}(u(\operatorname{cpsi\_}), \operatorname{cpsxs\_}) \in \mathbb{R}) \wedge \\ \operatorname{series\text{-}converges}((\lambda \operatorname{cpsi\_} \in \mathbb{N}.\; \operatorname{cc\text{-}sup\text{-}norm}(u(\operatorname{cpsi\_}), \operatorname{cpsxs\_}))) \end{array}
\end{equation*}

\par\medskip\noindent\textbf{Definition} (\texttt{cc-series-limit}).\;$\operatorname{cc\text{-}series\text{-}limit}(f) :=$ $(\iota\, \operatorname{cpsl\_}.\;\operatorname{cc\text{-}series\text{-}converges\text{-}to}(f, \operatorname{cpsl\_}))$

\vnbbookfile{theorem-library/cc-power-series-laws}{theorem-library/cc-power-series-laws}
\noindent{\small cc-power-series-laws.scm -- COMPLEX POWER SERIES: THE RADIUS OF CONVERGENCE (Prop 2.7) AND WHAT FOLLOWS FROM IT, as the notes state them. Batch 25-A, 2026-09-24.}\par
\vnbprop{rr-power-mul-base}{\vnbstmt{\forall k \in \mathbb{N}, x, y \in \mathbb{R}.\; (\operatorname{power}((x \cdot y), k) = (\operatorname{power}(x, k) \cdot \operatorname{power}(y, k)))}}
\vnbprop{cps-root-conv}{\vnbstmt{\begin{array}{@{}l@{}} \forall f \in (\mathbb{N} \to \mathbb{C}), s, r \in \mathbb{R}. \\ \quad ((\operatorname{elimsup}((\lambda j \in \mathbb{N}.\; {\operatorname{magnitude}(f(j))}^{{j}^{\ast\!-1}})) < s) \wedge ((0 \leq r) \wedge ((r \cdot s) < 1))) \Rightarrow  \\ \quad \quad \operatorname{series\text{-}converges}((\lambda k \in \mathbb{N}.\; (\operatorname{magnitude}(f(k)) \cdot \operatorname{power}(r, k)))) \end{array}}}
\vnbprop{cps-root-div}{\vnbstmt{\begin{array}{@{}l@{}} \forall f \in (\mathbb{N} \to \mathbb{C}), s, r \in \mathbb{R}. \\ \quad 0 \leq s \wedge \\ \quad s < \operatorname{elimsup}((\lambda j \in \mathbb{N}.\; {\operatorname{magnitude}(f(j))}^{{j}^{\ast\!-1}})) \wedge \\ \quad 0 \leq r \wedge \\ \quad 1 \leq (r \cdot s) \Rightarrow \\ \quad \quad \neg \operatorname{series\text{-}converges}((\lambda k \in \mathbb{N}.\; (\operatorname{magnitude}(f(k)) \cdot \operatorname{power}(r, k)))) \end{array}}}
\vnbprop{cps-small-factor}{\vnbstmt{\forall r \in \mathbb{R}.\; ((0 \leq r) \Rightarrow \exists s \in \mathbb{R}.\; ((0 < s) \wedge ((r \cdot s) < 1)))}}
\vnbprop{cps-radius-infinite}{\vnbstmt{\begin{array}{@{}l@{}} \forall f \in (\mathbb{N} \to \mathbb{C}). \\ \quad (\operatorname{elimsup}((\lambda j \in \mathbb{N}.\; {\operatorname{magnitude}(f(j))}^{{j}^{\ast\!-1}})) = 0) \Rightarrow  \\ \quad \quad \operatorname{cps\text{-}radius}(f) = \operatorname{pos\text{-}inf} \end{array}}}
\vnbprop{cps-radius-zero}{\vnbstmt{\begin{array}{@{}l@{}} \forall f \in (\mathbb{N} \to \mathbb{C}). \\ \quad (\operatorname{elimsup}((\lambda j \in \mathbb{N}.\; {\operatorname{magnitude}(f(j))}^{{j}^{\ast\!-1}})) = \operatorname{pos\text{-}inf}) \Rightarrow  \\ \quad \quad \operatorname{cps\text{-}radius}(f) = 0 \end{array}}}
\vnbprop{cps-radius-recip}{\vnbstmt{\begin{array}{@{}l@{}} \forall f \in (\mathbb{N} \to \mathbb{C}). \\ \quad \operatorname{elimsup}((\lambda j \in \mathbb{N}.\; {\operatorname{magnitude}(f(j))}^{{j}^{\ast\!-1}})) \in \mathbb{R} \wedge \\ \quad 0 < \operatorname{elimsup}((\lambda j \in \mathbb{N}.\; {\operatorname{magnitude}(f(j))}^{{j}^{\ast\!-1}})) \Rightarrow \\ \quad \quad \operatorname{cps\text{-}radius}(f) = \operatorname{recip}(\operatorname{elimsup}((\lambda j \in \mathbb{N}.\; {\operatorname{magnitude}(f(j))}^{{j}^{\ast\!-1}}))) \end{array}}}
\vnbprop{cps-radius-root}{\vnbstmt{\begin{array}{@{}l@{}} \forall f \in (\mathbb{N} \to \mathbb{C}). \\ \quad (\operatorname{elimsup}((\lambda j \in \mathbb{N}.\; {\operatorname{magnitude}(f(j))}^{{j}^{\ast\!-1}})) = 0) \Rightarrow  \\ \quad \quad \operatorname{cps\text{-}radius}(f) = \operatorname{pos\text{-}inf} \wedge \\ \quad (\operatorname{elimsup}((\lambda j \in \mathbb{N}.\; {\operatorname{magnitude}(f(j))}^{{j}^{\ast\!-1}})) = \operatorname{pos\text{-}inf}) \Rightarrow  \\ \quad \quad \operatorname{cps\text{-}radius}(f) = 0 \wedge \\ \quad \operatorname{elimsup}((\lambda j \in \mathbb{N}.\; {\operatorname{magnitude}(f(j))}^{{j}^{\ast\!-1}})) \in \mathbb{R} \wedge \\ \quad 0 < \operatorname{elimsup}((\lambda j \in \mathbb{N}.\; {\operatorname{magnitude}(f(j))}^{{j}^{\ast\!-1}})) \Rightarrow \\ \quad \quad \operatorname{cps\text{-}radius}(f) = \operatorname{recip}(\operatorname{elimsup}((\lambda j \in \mathbb{N}.\; {\operatorname{magnitude}(f(j))}^{{j}^{\ast\!-1}}))) \end{array}}}
\vnbprop{cps-radius-unfold}{\vnbstmt{\begin{array}{@{}l@{}} \forall f. \\ \quad (\operatorname{cps\text{-}radius}(f) \simeq \operatorname{esup}(\{r \in \mathbb{R} : ((0 \leq r) \wedge \operatorname{series\text{-}converges}((\lambda k \in \mathbb{N}.\; (\operatorname{magnitude}(f(k)) \cdot \operatorname{power}(r, k)))))\})) \end{array}}}
\vnbprop{cc-closed-ball-origin}{\vnbstmt{\forall d.; \forall z \in \operatorname{closed\text{-}ball}(\operatorname{cc\text{-}ms}, 0, d).\; ((z \in \mathbb{C}) \wedge (\operatorname{magnitude}(z) \leq d))}}
\vnbprop{cc-magnitude-power}{\vnbstmt{\forall k \in \mathbb{N}, z \in \mathbb{C}.\; (\operatorname{magnitude}(\operatorname{power}(z, k)) = \operatorname{power}(\operatorname{magnitude}(z), k))}}
\vnbprop{cc-sup-norm-bound}{\vnbstmt{\begin{array}{@{}l@{}} \forall u, s.; \forall m \in \mathbb{R}. \\ \quad ((0 \leq m) \wedge (\forall x \in s.\; (u(x) \in \mathbb{C}) \wedge \forall x \in s.\; (\operatorname{magnitude}(u(x)) \leq m))) \Rightarrow  \\ \quad \quad (\operatorname{cc\text{-}sup\text{-}norm}(u, s) \in \mathbb{R}) \wedge ((0 \leq \operatorname{cc\text{-}sup\text{-}norm}(u, s)) \wedge (\operatorname{cc\text{-}sup\text{-}norm}(u, s) \leq m)) \end{array}}}
\vnbprop{cps-radius-approx}{\vnbstmt{\begin{array}{@{}l@{}} \forall f \in (\mathbb{N} \to \mathbb{C}), d \in \mathbb{R}. \\ \quad ((0 \leq d) \wedge (d < \operatorname{cps\text{-}radius}(f))) \Rightarrow  \\ \quad \quad \exists t \in \mathbb{R}. \\ \quad \quad \quad (d < t) \wedge \operatorname{series\text{-}converges}((\lambda k \in \mathbb{N}.\; (\operatorname{magnitude}(f(k)) \cdot \operatorname{power}(t, k)))) \end{array}}}
\vnbprop{cps-normally-convergent}{\vnbstmt{\begin{array}{@{}l@{}} \forall f \in (\mathbb{N} \to \mathbb{C}), d \in \mathbb{R}. \\ \quad ((0 \leq d) \wedge (d < \operatorname{cps\text{-}radius}(f))) \Rightarrow  \\ \quad \quad \operatorname{is\text{-}normally\text{-}convergent}((\lambda k \in \mathbb{N}.\; (\lambda z \in \operatorname{closed\text{-}ball}(\operatorname{cc\text{-}ms}, 0, d).\; (f(k) \cdot \operatorname{power}(z, k)))), \\ \quad \quad \quad \operatorname{closed\text{-}ball}(\operatorname{cc\text{-}ms}, 0, d)) \end{array}}}
\vnbprop{cps-abs-conv-closed-ball}{\vnbstmt{\begin{array}{@{}l@{}} \forall f \in (\mathbb{N} \to \mathbb{C}), d \in \mathbb{R}, z \in \operatorname{closed\text{-}ball}(\operatorname{cc\text{-}ms}, 0, d). \\ \quad ((0 \leq d) \wedge (d < \operatorname{cps\text{-}radius}(f))) \Rightarrow  \\ \quad \quad \operatorname{series\text{-}converges}((\lambda k \in \mathbb{N}.\; \operatorname{magnitude}((f(k) \cdot \operatorname{power}(z, k))))) \end{array}}}
\vnbprop{cc-series-block-le}{\vnbstmt{\begin{array}{@{}l@{}} \forall k, m \in \mathbb{N}, f \in (\mathbb{N} \to \mathbb{C}), g \in (\mathbb{N} \to \mathbb{R}). \\ \quad (\forall j \in \mathbb{N}.\; (g(j) = \operatorname{magnitude}(f(j))) \wedge (m \leq k)) \Rightarrow  \\ \quad \quad \operatorname{magnitude}((\operatorname{cc\text{-}series\text{-}partial\text{-}sum}(f, k) - \operatorname{cc\text{-}series\text{-}partial\text{-}sum}(f, m))) \\ \quad \quad \leq\; (\operatorname{series\text{-}partial\text{-}sum}(g, k) - \operatorname{series\text{-}partial\text{-}sum}(g, m)) \end{array}}}
\vnbprop{cc-series-abs-converges}{\vnbstmt{\begin{array}{@{}l@{}} \forall f \in (\mathbb{N} \to \mathbb{C}), g \in (\mathbb{N} \to \mathbb{R}). \\ \quad (\forall j \in \mathbb{N}.\; (g(j) = \operatorname{magnitude}(f(j))) \wedge \operatorname{series\text{-}converges}(g)) \Rightarrow  \\ \quad \quad \operatorname{cc\text{-}series\text{-}converges}(f) \end{array}}}
\vnbprop{cc-series-limit-converges-to}{\vnbstmt{\begin{array}{@{}l@{}} \forall f \in (\mathbb{N} \to \mathbb{C}). \\ \quad \operatorname{cc\text{-}series\text{-}converges}(f) \Rightarrow \operatorname{cc\text{-}series\text{-}converges\text{-}to}(f, \operatorname{cc\text{-}series\text{-}limit}(f)) \end{array}}}
\vnbprop{cc-series-limit-in-cc}{\vnbstmt{\begin{array}{@{}l@{}} \forall f \in (\mathbb{N} \to \mathbb{C}). \\ \quad \operatorname{cc\text{-}series\text{-}converges}(f) \Rightarrow (\operatorname{cc\text{-}series\text{-}limit}(f) \in \mathbb{C}) \end{array}}}
\vnbprop{cps-converges-closed-ball}{\vnbstmt{\begin{array}{@{}l@{}} \forall f \in (\mathbb{N} \to \mathbb{C}), d \in \mathbb{R}, z \in \operatorname{closed\text{-}ball}(\operatorname{cc\text{-}ms}, 0, d). \\ \quad ((0 \leq d) \wedge (d < \operatorname{cps\text{-}radius}(f))) \Rightarrow  \\ \quad \quad \operatorname{cc\text{-}series\text{-}converges}((\lambda k \in \mathbb{N}.\; (f(k) \cdot \operatorname{power}(z, k)))) \end{array}}}
\vnbprop{cps-uniform-closed-ball}{\vnbstmt{\begin{array}{@{}l@{}} \forall f \in (\mathbb{N} \to \mathbb{C}), d \in \mathbb{R}, e. \\ \quad ((0 \leq d) \wedge ((d < \operatorname{cps\text{-}radius}(f)) \wedge \operatorname{pos\text{-}rr}(e))) \Rightarrow  \\ \quad \quad \exists n \in \mathbb{N}. \forall m \in \mathbb{N}, z \in \operatorname{closed\text{-}ball}(\operatorname{cc\text{-}ms}, 0, d). \\ \quad \quad \quad (n \leq m) \Rightarrow  \\ \quad \quad \quad \quad \operatorname{magnitude}((\operatorname{cc\text{-}series\text{-}limit}((\lambda k \in \mathbb{N}.\; (f(k) \cdot \operatorname{power}(z, k)))) - \operatorname{cc\text{-}series\text{-}partial\text{-}sum}((\lambda k \in \mathbb{N}.\; (f(k) \cdot \operatorname{power}(z, k))), m))) \\ \quad \quad \quad \quad \leq\; e \end{array}}}
\vnbprop{rr-power-power}{\vnbstmt{\forall k, m \in \mathbb{N}, x \in \mathbb{R}.\; (\operatorname{power}(\operatorname{power}(x, k), m) = \operatorname{power}(\operatorname{power}(x, m), k))}}
\vnbprop{cps-poly-geom-bound}{\vnbstmt{\begin{array}{@{}l@{}} \forall m \in \mathbb{N}, q \in \mathbb{R}. \\ \quad ((0 \leq q) \wedge (q < 1)) \Rightarrow  \\ \quad \quad \exists b \in \mathbb{R}.\; ((0 \leq b) \wedge \forall k \in \mathbb{N}.\; ((\operatorname{power}(k, m) \cdot \operatorname{power}(q, k)) \leq b)) \end{array}}}
\vnbprop{series-scale-converges}{\vnbstmt{\begin{array}{@{}l@{}} \forall f \in (\mathbb{N} \to \mathbb{R}), c \in \mathbb{R}, h \in (\mathbb{N} \to \mathbb{R}). \\ \quad \forall k \in \mathbb{N}.\; (0 \leq f(k)) \wedge \\ \quad 0 \leq c \wedge \\ \quad \forall k \in \mathbb{N}.\; (h(k) = (c \cdot f(k))) \wedge \\ \quad \operatorname{series\text{-}converges}(f) \Rightarrow \\ \quad \quad \operatorname{series\text{-}converges}(h) \end{array}}}
\vnbprop{cps-factor-conv}{\vnbstmt{\begin{array}{@{}l@{}} \forall f, g, a \in (\mathbb{N} \to \mathbb{C}), c \in \mathbb{R}, m \in \mathbb{N}, t, d \in \mathbb{R}. \\ \quad \forall k \in \mathbb{N}.\; (a(k) = (f(k) \cdot g(k))) \wedge \\ \quad \forall k \in \mathbb{N}.\; (\operatorname{magnitude}(g(k)) \leq (c \cdot \operatorname{power}(k, m))) \wedge \\ \quad 0 \leq d \wedge \\ \quad d < t \wedge \\ \quad \operatorname{series\text{-}converges}((\lambda k \in \mathbb{N}.\; (\operatorname{magnitude}(f(k)) \cdot \operatorname{power}(t, k)))) \Rightarrow \\ \quad \quad \operatorname{series\text{-}converges}((\lambda k \in \mathbb{N}.\; (\operatorname{magnitude}(a(k)) \cdot \operatorname{power}(d, k)))) \end{array}}}
\vnbprop{cps-radius-factor}{\vnbstmt{\begin{array}{@{}l@{}} \forall f, g, a \in (\mathbb{N} \to \mathbb{C}), c \in \mathbb{R}, m \in \mathbb{N}. \\ \quad (\forall k \in \mathbb{N}.\; (a(k) = (f(k) \cdot g(k))) \wedge \forall k \in \mathbb{N}.\; (\operatorname{magnitude}(g(k)) \leq (c \cdot \operatorname{power}(k, m)))) \Rightarrow  \\ \quad \quad \operatorname{cps\text{-}radius}(f) \leq \operatorname{cps\text{-}radius}(a) \end{array}}}

\vnbbookfile{theorem-library/cc-power-series-deriv}{theorem-library/cc-power-series-deriv}
\noindent{\small cc-power-series-deriv.scm -- PROPOSITION 2.10: THE SUM OF A COMPLEX POWER SERIES IS HOLOMORPHIC INSIDE THE RADIUS, AND ITS DERIVATIVE IS THE TERMWISE DERIVED SERIES, WHICH HAS THE SAME RADIUS. Batch 26-A, 2026-09-24.}\par
\vnbprop{cpd-cc-squeeze}{\vnbstmt{\begin{array}{@{}l@{}} \forall s \in (\mathbb{N} \to \mathbb{C}), v \in \mathbb{C}, m \in (\mathbb{N} \to \mathbb{R}). \\ \quad (\operatorname{converges\text{-}to}(\operatorname{rr\text{-}ms}, m, 0) \wedge \forall j \in \mathbb{N}.\; (\operatorname{magnitude}((s(j) - v)) \leq m(j))) \Rightarrow  \\ \quad \quad \operatorname{converges\text{-}to}(\operatorname{cc\text{-}ms}, s, v) \end{array}}}
\vnbprop{cc-limit-lincomb}{\vnbstmt{\begin{array}{@{}l@{}} \forall s, t, h \in (\mathbb{N} \to \mathbb{C}), v, a, c, d \in \mathbb{C}. \\ \quad \forall j \in \mathbb{N}.\; (h(j) = ((c \cdot s(j)) + (d \cdot t(j)))) \wedge \\ \quad \operatorname{converges\text{-}to}(\operatorname{cc\text{-}ms}, s, v) \wedge \\ \quad \operatorname{converges\text{-}to}(\operatorname{cc\text{-}ms}, t, a) \Rightarrow \\ \quad \quad \operatorname{converges\text{-}to}(\operatorname{cc\text{-}ms}, h, ((c \cdot v) + (d \cdot a))) \end{array}}}
\vnbprop{cc-limit-shift}{\vnbstmt{\begin{array}{@{}l@{}} \forall s, h \in (\mathbb{N} \to \mathbb{C}), v \in \mathbb{C}. \\ \quad (\forall j \in \mathbb{N}.\; (h(j) = s((j + 1))) \wedge \operatorname{converges\text{-}to}(\operatorname{cc\text{-}ms}, s, v)) \Rightarrow  \\ \quad \quad \operatorname{converges\text{-}to}(\operatorname{cc\text{-}ms}, h, v) \end{array}}}
\vnbprop{cc-series-partial-sum-lincomb}{\vnbstmt{\begin{array}{@{}l@{}} \forall k \in \mathbb{N}, f, g, h \in (\mathbb{N} \to \mathbb{C}), c, d \in \mathbb{C}. \\ \quad \forall j \in \mathbb{N}.\; (h(j) = ((c \cdot f(j)) + (d \cdot g(j)))) \Rightarrow  \\ \quad \quad \operatorname{cc\text{-}series\text{-}partial\text{-}sum}(h, k) = ((c \cdot \operatorname{cc\text{-}series\text{-}partial\text{-}sum}(f, k)) + (d \cdot \operatorname{cc\text{-}series\text{-}partial\text{-}sum}(g, k))) \end{array}}}
\vnbprop{cc-series-limit-linear}{\vnbstmt{\begin{array}{@{}l@{}} \forall f, g, h \in (\mathbb{N} \to \mathbb{C}), c, d \in \mathbb{C}. \\ \quad (\forall j \in \mathbb{N}.\; (h(j) = ((c \cdot f(j)) + (d \cdot g(j)))) \wedge (\operatorname{cc\text{-}series\text{-}converges}(f) \wedge \operatorname{cc\text{-}series\text{-}converges}(g))) \Rightarrow  \\ \quad \quad (\operatorname{cc\text{-}series\text{-}converges}(h) \wedge (\operatorname{cc\text{-}series\text{-}limit}(h) = ((c \cdot \operatorname{cc\text{-}series\text{-}limit}(f)) + (d \cdot \operatorname{cc\text{-}series\text{-}limit}(g))))) \end{array}}}
\vnbprop{cpd-power-second-order}{\vnbstmt{\begin{array}{@{}l@{}} \forall k \in \mathbb{N}, z, u \in \mathbb{C}, r \in \mathbb{R}. \\ \quad ((\operatorname{magnitude}(z) \leq r) \wedge (\operatorname{magnitude}(u) \leq r)) \Rightarrow  \\ \quad \quad r \cdot \operatorname{magnitude}(((\operatorname{power}(z, (k + 1)) - \operatorname{power}(u, (k + 1))) - (((k + 1) \cdot \operatorname{power}(u, k)) \cdot (z - u)))) \\ \quad \quad \leq\; (((k + 1) \cdot (k + 1)) \cdot \operatorname{power}(r, k)) \cdot (\operatorname{magnitude}((z - u)) \cdot \operatorname{magnitude}((z - u))) \end{array}}}
\vnbprop{cc-power-diff-bound}{\vnbstmt{\begin{array}{@{}l@{}} \forall k \in \mathbb{N}, z, u \in \mathbb{C}, r \in \mathbb{R}. \\ \quad ((\operatorname{magnitude}(z) \leq r) \wedge (\operatorname{magnitude}(u) \leq r)) \Rightarrow  \\ \quad \quad \operatorname{magnitude}((\operatorname{power}(z, (k + 1)) - \operatorname{power}(u, (k + 1)))) \leq (((k + 1) \cdot \operatorname{power}(r, k)) \cdot \operatorname{magnitude}((z - u))) \end{array}}}
\vnbprop{cps-shift-conv}{\vnbstmt{\begin{array}{@{}l@{}} \forall f, g \in (\mathbb{N} \to \mathbb{C}), t \in \mathbb{R}. \\ \quad (\forall j \in \mathbb{N}.\; (\operatorname{magnitude}(g(j)) = \operatorname{magnitude}(f((j + 1)))) \wedge (0 < t)) \Rightarrow  \\ \quad \quad \operatorname{series\text{-}converges}((\lambda k \in \mathbb{N}.\; (\operatorname{magnitude}(f(k)) \cdot \operatorname{power}(t, k)))) \\ \quad \quad \Leftrightarrow\; \operatorname{series\text{-}converges}((\lambda k \in \mathbb{N}.\; (\operatorname{magnitude}(g(k)) \cdot \operatorname{power}(t, k)))) \end{array}}}
\vnbprop{cps-radius-in}{\vnbstmt{\forall f.\; (\operatorname{cps\text{-}radius}(f) \in \operatorname{rr\text{-}pos\text{-}star})}}
\vnbprop{cps-radius-mono}{\vnbstmt{\begin{array}{@{}l@{}} \forall f, g \in (\mathbb{N} \to \mathbb{C}). \\ \quad \forall t \in \mathbb{R}. \\ \quad \quad ((0 < t) \wedge \operatorname{series\text{-}converges}((\lambda k \in \mathbb{N}.\; (\operatorname{magnitude}(f(k)) \cdot \operatorname{power}(t, k))))) \Rightarrow  \\ \quad \quad \quad \operatorname{series\text{-}converges}((\lambda k \in \mathbb{N}.\; (\operatorname{magnitude}(g(k)) \cdot \operatorname{power}(t, k)))) \Rightarrow \\ \quad \quad \operatorname{cps\text{-}radius}(f) \leq \operatorname{cps\text{-}radius}(g) \end{array}}}
\vnbprop{cps-radius-derived}{\vnbstmt{\begin{array}{@{}l@{}} \forall f, a \in (\mathbb{N} \to \mathbb{C}). \\ \quad \forall j \in \mathbb{N}.\; (a(j) = ((j + 1) \cdot f((j + 1)))) \Rightarrow  \\ \quad \quad \operatorname{cps\text{-}radius}(a) = \operatorname{cps\text{-}radius}(f) \end{array}}}
\vnbprop{cc-diff-on-quadratic}{\vnbstmt{\begin{array}{@{}l@{}} \forall s.; \forall f \in (s \to \mathbb{C}), a \in s, l \in \mathbb{C}, m \in \mathbb{R}. \\ \quad \operatorname{is\text{-}open}(\operatorname{nf\text{-}metric\text{-}space}(\operatorname{cc\text{-}normed\text{-}field}), s) \wedge \\ \quad 0 \leq m \wedge \\ \quad \forall x \in s. \\ \quad \quad \operatorname{magnitude}(((f(x) - f(a)) - ((x - a) \cdot l))) \leq (m \cdot (\operatorname{magnitude}((x - a)) \cdot \operatorname{magnitude}((x - a)))) \Rightarrow \\ \quad \quad \operatorname{is\text{-}diff\text{-}on}(\operatorname{cc\text{-}normed\text{-}field}, s, f, a, l) \end{array}}}
\vnbprop{cps-partial-error-bound}{\vnbstmt{\begin{array}{@{}l@{}} \forall k \in \mathbb{N}, f, a \in (\mathbb{N} \to \mathbb{C}), z, u \in \mathbb{C}, r \in \mathbb{R}, \mathit{cpdtz}, \mathit{cpdtu}, d \in (\mathbb{N} \to \mathbb{C}), v \in (\mathbb{N} \to \mathbb{R}). \\ \quad \forall j \in \mathbb{N}.\; (a(j) = ((j + 1) \cdot f((j + 1)))) \wedge \\ \quad \operatorname{magnitude}(z) \leq r \wedge \\ \quad \operatorname{magnitude}(u) \leq r \wedge \\ \quad \forall j \in \mathbb{N}.\; (\operatorname{cpdtz}(j) = (f(j) \cdot \operatorname{power}(z, j))) \wedge \\ \quad \forall j \in \mathbb{N}.\; (\operatorname{cpdtu}(j) = (f(j) \cdot \operatorname{power}(u, j))) \wedge \\ \quad \forall j \in \mathbb{N}.\; (d(j) = (a(j) \cdot \operatorname{power}(u, j))) \wedge \\ \quad \forall j \in \mathbb{N}. \\ \quad \quad v(j) = (\operatorname{magnitude}((f((j + 1)) \cdot ((j + 1) \cdot (j + 1)))) \cdot \operatorname{power}(r, (j + 1))) \Rightarrow \\ \quad \quad (r \cdot r) \cdot \operatorname{magnitude}(((\operatorname{cc\text{-}series\text{-}partial\text{-}sum}(\mathit{cpdtz}, (k + 1)) - \operatorname{cc\text{-}series\text{-}partial\text{-}sum}(\mathit{cpdtu}, (k + 1))) - ((z - u) \cdot \operatorname{cc\text{-}series\text{-}partial\text{-}sum}(d, k)))) \\ \quad \quad \leq\; (\operatorname{magnitude}((z - u)) \cdot \operatorname{magnitude}((z - u))) \cdot \operatorname{series\text{-}partial\text{-}sum}(v, k) \end{array}}}
\vnbprop{cc-partial-error-limit}{\vnbstmt{\begin{array}{@{}l@{}} \forall f, g, h \in (\mathbb{N} \to \mathbb{C}), w \in \mathbb{C}. \\ \quad (\operatorname{cc\text{-}series\text{-}converges}(f) \wedge (\operatorname{cc\text{-}series\text{-}converges}(g) \wedge \operatorname{cc\text{-}series\text{-}converges}(h))) \Rightarrow  \\ \quad \quad \operatorname{converges\text{-}to}(\operatorname{cc\text{-}ms}, \\ \quad \quad \quad (\lambda j \in \mathbb{N}.\; \\ \quad \quad \quad \quad ((\operatorname{cc\text{-}series\text{-}partial\text{-}sum}(f, (j + 1)) - \operatorname{cc\text{-}series\text{-}partial\text{-}sum}(g, (j + 1))) - (w \cdot \operatorname{cc\text{-}series\text{-}partial\text{-}sum}(h, j)))), \\ \quad \quad \quad ((\operatorname{cc\text{-}series\text{-}limit}(f) - \operatorname{cc\text{-}series\text{-}limit}(g)) - (w \cdot \operatorname{cc\text{-}series\text{-}limit}(h)))) \end{array}}}
\vnbprop{cps-diff-on-ball}{\vnbstmt{\begin{array}{@{}l@{}} \forall f, a \in (\mathbb{N} \to \mathbb{C}), r \in \mathbb{R}, \mathit{cpda} \in \operatorname{ball}(\operatorname{cc\text{-}ms}, 0, r). \\ \quad (\forall j \in \mathbb{N}.\; (a(j) = ((j + 1) \cdot f((j + 1)))) \wedge ((0 < r) \wedge (r < \operatorname{cps\text{-}radius}(f)))) \Rightarrow  \\ \quad \quad \operatorname{is\text{-}diff\text{-}on}(\operatorname{cc\text{-}normed\text{-}field}, \operatorname{ball}(\operatorname{cc\text{-}ms}, 0, r), \\ \quad \quad \quad (\lambda z \in \operatorname{ball}(\operatorname{cc\text{-}ms}, 0, r).\; \operatorname{cc\text{-}series\text{-}limit}((\lambda k \in \mathbb{N}.\; (f(k) \cdot \operatorname{power}(z, k))))), \\ \quad \quad \quad \mathit{cpda}, \\ \quad \quad \quad \operatorname{cc\text{-}series\text{-}limit}((\lambda k \in \mathbb{N}.\; (a(k) \cdot \operatorname{power}(\mathit{cpda}, k))))) \end{array}}}
\vnbprop{cps-holomorphic-ball}{\vnbstmt{\begin{array}{@{}l@{}} \forall f \in (\mathbb{N} \to \mathbb{C}), r \in \mathbb{R}. \\ \quad ((0 < r) \wedge (r < \operatorname{cps\text{-}radius}(f))) \Rightarrow  \\ \quad \quad \operatorname{holomorphic\text{-}on}(\operatorname{ball}(\operatorname{cc\text{-}ms}, 0, r), \\ \quad \quad \quad (\lambda z \in \operatorname{ball}(\operatorname{cc\text{-}ms}, 0, r).\; \operatorname{cc\text{-}series\text{-}limit}((\lambda k \in \mathbb{N}.\; (f(k) \cdot \operatorname{power}(z, k)))))) \wedge \\ \quad \quad \operatorname{cps\text{-}radius}((\lambda n \in \mathbb{N}.\; ((n + 1) \cdot f((n + 1))))) = \operatorname{cps\text{-}radius}(f) \wedge \\ \quad \quad \forall a \in \operatorname{ball}(\operatorname{cc\text{-}ms}, 0, r). \\ \quad \quad \quad \operatorname{is\text{-}diff\text{-}on}(\operatorname{cc\text{-}normed\text{-}field}, \operatorname{ball}(\operatorname{cc\text{-}ms}, 0, r), \\ \quad \quad \quad \quad (\lambda z \in \operatorname{ball}(\operatorname{cc\text{-}ms}, 0, r).\; \operatorname{cc\text{-}series\text{-}limit}((\lambda k \in \mathbb{N}.\; (f(k) \cdot \operatorname{power}(z, k))))), \\ \quad \quad \quad \quad a, \\ \quad \quad \quad \quad \operatorname{cc\text{-}series\text{-}limit}((\lambda k \in \mathbb{N}.\; (((k + 1) \cdot f((k + 1))) \cdot \operatorname{power}(a, k))))) \end{array}}}

\vnbbookfile{theorem-library/cps-taylor-coefficients}{theorem-library/cps-taylor-coefficients}
\noindent{\small cps-taylor-coefficients.scm -- COROLLARY 2.14: THE TAYLOR COEFFICIENTS OF A COMPLEX POWER SERIES.  f\textasciicircum{}(k)(0) = k! a\_k. Batch 27-A, 2026-09-24.}\par
\vnbprop{cps-kth-derived-step}{\vnbstmt{\begin{array}{@{}l@{}} \forall f \in (\mathbb{N} \to \mathbb{C}), k \in \mathbb{N}, g, h \in (\mathbb{N} \to \mathbb{C}), j \in \mathbb{N}. \\ \quad \forall n \in \mathbb{N}.\; (g(n) = ((\operatorname{factorial}((n + k)) \cdot \operatorname{recip}(\operatorname{factorial}(n))) \cdot f((n + k)))) \wedge \\ \quad \forall n \in \mathbb{N}. \\ \quad \quad h(n) = ((\operatorname{factorial}((n + (k + 1))) \cdot \operatorname{recip}(\operatorname{factorial}(n))) \cdot f((n + (k + 1)))) \Rightarrow \\ \quad \quad h(j) = ((j + 1) \cdot g((j + 1))) \end{array}}}
\vnbprop{cps-kth-derived-radius}{\vnbstmt{\begin{array}{@{}l@{}} \forall f \in (\mathbb{N} \to \mathbb{C}), k \in \mathbb{N}, g \in (\mathbb{N} \to \mathbb{C}). \\ \quad \forall n \in \mathbb{N}.\; (g(n) = ((\operatorname{factorial}((n + k)) \cdot \operatorname{recip}(\operatorname{factorial}(n))) \cdot f((n + k)))) \Rightarrow  \\ \quad \quad \operatorname{cps\text{-}radius}(g) = \operatorname{cps\text{-}radius}(f) \end{array}}}
\vnbprop{cps-sum-at-zero}{\vnbstmt{\begin{array}{@{}l@{}} \forall g \in (\mathbb{N} \to \mathbb{C}). \\ \quad \operatorname{cc\text{-}series\text{-}converges}((\lambda k \in \mathbb{N}.\; (g(k) \cdot \operatorname{power}(0, k)))) \wedge \\ \quad \operatorname{cc\text{-}series\text{-}limit}((\lambda k \in \mathbb{N}.\; (g(k) \cdot \operatorname{power}(0, k)))) = g(0) \end{array}}}
\vnbprop{cps-taylor-coefficients}{\vnbstmt{\begin{array}{@{}l@{}} \forall f \in (\mathbb{N} \to \mathbb{C}), k \in \mathbb{N}, g, h \in (\mathbb{N} \to \mathbb{C}), r \in \mathbb{R}. \\ \quad \forall n \in \mathbb{N}.\; (g(n) = ((\operatorname{factorial}((n + k)) \cdot \operatorname{recip}(\operatorname{factorial}(n))) \cdot f((n + k)))) \wedge \\ \quad \forall n \in \mathbb{N}. \\ \quad \quad h(n) = ((\operatorname{factorial}((n + (k + 1))) \cdot \operatorname{recip}(\operatorname{factorial}(n))) \cdot f((n + (k + 1)))) \wedge \\ \quad 0 < r \wedge \\ \quad r < \operatorname{cps\text{-}radius}(f) \Rightarrow \\ \quad \quad \operatorname{holomorphic\text{-}on}(\operatorname{ball}(\operatorname{cc\text{-}ms}, 0, r), \\ \quad \quad \quad (\lambda z \in \operatorname{ball}(\operatorname{cc\text{-}ms}, 0, r).\; \operatorname{cc\text{-}series\text{-}limit}((\lambda \mathit{cpsk} \in \mathbb{N}.\; (g(\mathit{cpsk}) \cdot \operatorname{power}(z, \mathit{cpsk})))))) \wedge \\ \quad \quad (\lambda z \in \operatorname{ball}(\operatorname{cc\text{-}ms}, 0, r).\; \operatorname{cc\text{-}series\text{-}limit}((\lambda \mathit{cpsk} \in \mathbb{N}.\; (g(\mathit{cpsk}) \cdot \operatorname{power}(z, \mathit{cpsk})))))(0) \\ \quad \quad =\; \operatorname{factorial}(k) \cdot f(k) \wedge \\ \quad \quad \forall a \in \operatorname{ball}(\operatorname{cc\text{-}ms}, 0, r). \\ \quad \quad \quad \operatorname{deriv\text{-}on}(\operatorname{cc\text{-}normed\text{-}field}, \operatorname{ball}(\operatorname{cc\text{-}ms}, 0, r), \\ \quad \quad \quad \quad (\lambda z \in \operatorname{ball}(\operatorname{cc\text{-}ms}, 0, r).\; \operatorname{cc\text{-}series\text{-}limit}((\lambda \mathit{cpsk} \in \mathbb{N}.\; (g(\mathit{cpsk}) \cdot \operatorname{power}(z, \mathit{cpsk}))))), \\ \quad \quad \quad \quad a) \\ \quad \quad \quad =\; (\lambda z \in \operatorname{ball}(\operatorname{cc\text{-}ms}, 0, r).\; \operatorname{cc\text{-}series\text{-}limit}((\lambda \mathit{cpsk} \in \mathbb{N}.\; (h(\mathit{cpsk}) \cdot \operatorname{power}(z, \mathit{cpsk})))))(a) \end{array}}}
\vnbprop{cps-coefficient-formula}{\vnbstmt{\begin{array}{@{}l@{}} \forall f \in (\mathbb{N} \to \mathbb{C}), k \in \mathbb{N}, g \in (\mathbb{N} \to \mathbb{C}), r \in \mathbb{R}. \\ \quad (\forall n \in \mathbb{N}.\; (g(n) = ((\operatorname{factorial}((n + k)) \cdot \operatorname{recip}(\operatorname{factorial}(n))) \cdot f((n + k)))) \wedge (0 < r)) \Rightarrow  \\ \quad \quad f(k) \\ \quad \quad =\; (\lambda z \in \operatorname{ball}(\operatorname{cc\text{-}ms}, 0, r).\; \operatorname{cc\text{-}series\text{-}limit}((\lambda \mathit{cpsk} \in \mathbb{N}.\; (g(\mathit{cpsk}) \cdot \operatorname{power}(z, \mathit{cpsk})))))(0) \cdot \operatorname{recip}(\operatorname{factorial}(k)) \end{array}}}

\vnbbookfile{theorem-library/nth-deriv-laws}{theorem-library/nth-deriv-laws}
\noindent{\small nth-deriv-laws.scm -- THE LAWS OF NTH-DERIV-ON, AND COROLLARY 2.14 RESTATED LITERALLY: f\textasciicircum{}(k)(0) = k! a\_k. Batch 36, 2026-09-26.}\par
\vnbprop{nth-deriv-on-succ-apply}{\vnbstmt{\begin{array}{@{}l@{}} \forall f, u, \mathit{ndlf}.; \forall j \in \mathbb{N}, z \in u. \\ \quad (\mathit{ndlf}^{((j + 1))}(z) \simeq \operatorname{deriv\text{-}on}(f, u, \mathit{ndlf}^{(j)}, z)) \end{array}}}
\vnbprop{nth-deriv-on-of-is-diff-on}{\vnbstmt{\begin{array}{@{}l@{}} \forall f, u, \mathit{ndlf}.; \forall j \in \mathbb{N}, g.; \forall w \in u. \\ \quad \operatorname{is\text{-}normed\text{-}field}(f) \wedge \\ \quad \forall e. \\ \quad \quad \operatorname{pos\text{-}rr}(e) \Rightarrow \exists h \in \operatorname{carr}(f).\; (\neg (h = \operatorname{zero}(f)) \wedge (\operatorname{fnrm}(f)(h) < e)) \wedge \\ \quad \forall z \in u.\; \operatorname{is\text{-}diff\text{-}on}(f, u, \mathit{ndlf}^{(j)}, z, g(z)) \Rightarrow \\ \quad \quad \mathit{ndlf}^{((j + 1))}(w) = g(w) \end{array}}}
\vnbprop{nth-deriv-on-succ-eq}{\vnbstmt{\begin{array}{@{}l@{}} \forall f, u, \mathit{ndlf}.; \forall j \in \mathbb{N}, g \in (u \to \operatorname{carr}(f)). \\ \quad \operatorname{is\text{-}normed\text{-}field}(f) \wedge \\ \quad \forall e. \\ \quad \quad \operatorname{pos\text{-}rr}(e) \Rightarrow \exists h \in \operatorname{carr}(f).\; (\neg (h = \operatorname{zero}(f)) \wedge (\operatorname{fnrm}(f)(h) < e)) \wedge \\ \quad \forall z \in u.\; \operatorname{is\text{-}diff\text{-}on}(f, u, \mathit{ndlf}^{(j)}, z, g(z)) \Rightarrow \\ \quad \quad \mathit{ndlf}^{((j + 1))} = g \end{array}}}
\vnbprop{nth-deriv-on-in-fun}{\vnbstmt{\begin{array}{@{}l@{}} \forall f, u.; \forall \mathit{ndlf} \in (u \to \operatorname{carr}(f)), j \in \mathbb{N}. \\ \quad \operatorname{is\text{-}normed\text{-}field}(f) \wedge \\ \quad \forall e. \\ \quad \quad \operatorname{pos\text{-}rr}(e) \Rightarrow \exists h \in \operatorname{carr}(f).\; (\neg (h = \operatorname{zero}(f)) \wedge (\operatorname{fnrm}(f)(h) < e)) \wedge \\ \quad \forall i \in \mathbb{N}, z \in u. \\ \quad \quad (i < j) \Rightarrow \exists l.\; \operatorname{is\text{-}diff\text{-}on}(f, u, \mathit{ndlf}^{(i)}, z, l) \Rightarrow \\ \quad \quad \mathit{ndlf}^{(j)} \in (u \to \operatorname{carr}(f)) \end{array}}}
\vnbprop{cps-nth-deriv}{\vnbstmt{\begin{array}{@{}l@{}} \forall f \in (\mathbb{N} \to \mathbb{C}), r \in \mathbb{R}, k \in \mathbb{N}, g \in (\mathbb{N} \to \mathbb{C}). \\ \quad 0 < r \wedge \\ \quad r < \operatorname{cps\text{-}radius}(f) \wedge \\ \quad \forall n \in \mathbb{N}.\; (g(n) = ((\operatorname{factorial}((n + k)) \cdot \operatorname{recip}(\operatorname{factorial}(n))) \cdot f((n + k)))) \Rightarrow \\ \quad \quad (\lambda z \in \operatorname{ball}(\operatorname{cc\text{-}ms}, 0, r).\; \operatorname{cc\text{-}series\text{-}limit}((\lambda \mathit{cpsk} \in \mathbb{N}.\; (f(\mathit{cpsk}) \cdot \operatorname{power}(z, \mathit{cpsk})))))^{(k)} \\ \quad \quad =\; (\lambda z \in \operatorname{ball}(\operatorname{cc\text{-}ms}, 0, r).\; \operatorname{cc\text{-}series\text{-}limit}((\lambda \mathit{cpsk} \in \mathbb{N}.\; (g(\mathit{cpsk}) \cdot \operatorname{power}(z, \mathit{cpsk}))))) \end{array}}}
\vnbprop{taylor-coefficients}{\vnbstmt{\begin{array}{@{}l@{}} \forall f \in (\mathbb{N} \to \mathbb{C}), r \in \mathbb{R}, k \in \mathbb{N}. \\ \quad ((0 < r) \wedge (r < \operatorname{cps\text{-}radius}(f))) \Rightarrow  \\ \quad \quad (\lambda z \in \operatorname{ball}(\operatorname{cc\text{-}ms}, 0, r).\; \operatorname{cc\text{-}series\text{-}limit}((\lambda \mathit{cpsk} \in \mathbb{N}.\; (f(\mathit{cpsk}) \cdot \operatorname{power}(z, \mathit{cpsk})))))^{(k)}(0) \\ \quad \quad =\; \operatorname{factorial}(k) \cdot f(k) \end{array}}}
\vnbprop{taylor-coefficient-formula}{\vnbstmt{\begin{array}{@{}l@{}} \forall f \in (\mathbb{N} \to \mathbb{C}), r \in \mathbb{R}, k \in \mathbb{N}. \\ \quad ((0 < r) \wedge (r < \operatorname{cps\text{-}radius}(f))) \Rightarrow  \\ \quad \quad f(k) \\ \quad \quad =\; (\lambda z \in \operatorname{ball}(\operatorname{cc\text{-}ms}, 0, r).\; \operatorname{cc\text{-}series\text{-}limit}((\lambda \mathit{cpsk} \in \mathbb{N}.\; (f(\mathit{cpsk}) \cdot \operatorname{power}(z, \mathit{cpsk})))))^{(k)}(0) \cdot \operatorname{recip}(\operatorname{factorial}(k)) \end{array}}}
\vnbprop{nth-deriv-on-holomorphic-ball}{\vnbstmt{\begin{array}{@{}l@{}} \forall f \in (\mathbb{N} \to \mathbb{C}), r \in \mathbb{R}, k \in \mathbb{N}. \\ \quad ((0 < r) \wedge (r < \operatorname{cps\text{-}radius}(f))) \Rightarrow  \\ \quad \quad \operatorname{holomorphic\text{-}on}(\operatorname{ball}(\operatorname{cc\text{-}ms}, 0, r), \\ \quad \quad \quad (\lambda z \in \operatorname{ball}(\operatorname{cc\text{-}ms}, 0, r).\; \operatorname{cc\text{-}series\text{-}limit}((\lambda \mathit{cpsk} \in \mathbb{N}.\; (f(\mathit{cpsk}) \cdot \operatorname{power}(z, \mathit{cpsk})))))^{(k)}) \end{array}}}

\vnbbookfile{theorem-library/line-int-fundamental}{theorem-library/line-int-fundamental}
\noindent{\small line-int-fundamental.scm -- THE FUNDAMENTAL THEOREM FOR INTEGRALS ALONG A ROAD (the user's notes, \textasciitilde{}/docs/complex-analysis.pdf, section 3.1, Proposition 3.2 and equation (50)), and the chain rule it rests on.}\par
\vnbprop{cc-compose-continuous-at}{\vnbstmt{\begin{array}{@{}l@{}} \forall s, u, g, i, x. \\ \quad \operatorname{is\text{-}metric\text{-}space}(s) \wedge \\ \quad u \subseteq \mathbb{C} \wedge \\ \quad \operatorname{is\text{-}continuous\text{-}at}(s, \operatorname{cc\text{-}ms}, g, x) \wedge \\ \quad \forall y \in \operatorname{pts}(s).\; (g(y) \in u) \wedge \\ \quad \operatorname{is\text{-}continuous\text{-}at}(\operatorname{subspace\text{-}ms}(\operatorname{nf\text{-}metric\text{-}space}(\operatorname{cc\text{-}normed\text{-}field}), u), \\ \quad \quad \operatorname{nf\text{-}metric\text{-}space}(\operatorname{cc\text{-}normed\text{-}field}), i, g(x)) \Rightarrow \\ \quad \quad \operatorname{is\text{-}continuous\text{-}at}(s, \operatorname{cc\text{-}ms}, (\lambda z \in \operatorname{pts}(s).\; i(g(z))), x) \end{array}}}
\vnbprop{cc-product-continuous-at}{\vnbstmt{\begin{array}{@{}l@{}} \forall s, p, q, x. \\ \quad (\operatorname{is\text{-}continuous\text{-}at}(s, \operatorname{cc\text{-}ms}, p, x) \wedge \operatorname{is\text{-}continuous\text{-}at}(s, \operatorname{cc\text{-}ms}, q, x)) \Rightarrow  \\ \quad \quad \operatorname{is\text{-}continuous\text{-}at}(s, \operatorname{cc\text{-}ms}, (\lambda z \in \operatorname{pts}(s).\; (p(z) \cdot q(z))), x) \end{array}}}
\vnbprop{lfn-cc-factor-identity}{\vnbstmt{\begin{array}{@{}l@{}} \forall a, b, c, e, h \in \mathbb{C}. \\ \quad (a \cdot ((b \cdot h) + (e \cdot (c \cdot h)))) = (h \cdot (a \cdot (b + (c \cdot e)))) \end{array}}}
\vnbprop{holomorphic-chain-within}{\vnbstmt{\begin{array}{@{}l@{}} \forall m, w, g, x, u, f, l, d. \\ \quad m \subseteq \mathbb{R} \wedge \\ \quad w \subseteq m \wedge \\ \quad g \in (m \to \mathbb{C}) \wedge \\ \quad \operatorname{is\text{-}continuous\text{-}at}(\operatorname{subspace\text{-}ms}(\operatorname{rr\text{-}ms}, m), \operatorname{cc\text{-}ms}, g, x) \wedge \\ \quad \forall y \in m.\; (g(y) \in u) \wedge \\ \quad \operatorname{is\text{-}diff\text{-}on}(\operatorname{cc\text{-}normed\text{-}field}, u, f, g(x), l) \wedge \\ \quad d \in \mathbb{C} \wedge \\ \quad \operatorname{has\text{-}deriv\text{-}at\text{-}within}((\lambda t \in m.\; \operatorname{real\text{-}part}(g(t))), w, x, \operatorname{real\text{-}part}(d)) \wedge \\ \quad \operatorname{has\text{-}deriv\text{-}at\text{-}within}((\lambda t \in m.\; \operatorname{imag\text{-}part}(g(t))), w, x, \operatorname{imag\text{-}part}(d)) \Rightarrow \\ \quad \quad \operatorname{has\text{-}deriv\text{-}at\text{-}within}((\lambda t \in m.\; \operatorname{real\text{-}part}(f(g(t)))), w, x, \operatorname{real\text{-}part}((l \cdot d))) \wedge \\ \quad \quad \operatorname{has\text{-}deriv\text{-}at\text{-}within}((\lambda t \in m.\; \operatorname{imag\text{-}part}(f(g(t)))), w, x, \operatorname{imag\text{-}part}((l \cdot d))) \end{array}}}
\vnbprop{ooint-center-nested}{\vnbstmt{\begin{array}{@{}l@{}} \forall x \in \mathbb{R}, v.; \forall r \in \mathbb{R}. \\ \quad (\operatorname{pos\text{-}rr}(v) \wedge (v \leq r)) \Rightarrow (\operatorname{ooint}((x - v), (x + v)) \subseteq \operatorname{ooint}((x - r), (x + r))) \end{array}}}
\vnbprop{holomorphic-chain}{\vnbstmt{\begin{array}{@{}l@{}} \forall m, g, x, u, f, l, d. \\ \quad m \subseteq \mathbb{R} \wedge \\ \quad g \in (m \to \mathbb{C}) \wedge \\ \quad \operatorname{is\text{-}continuous\text{-}at}(\operatorname{subspace\text{-}ms}(\operatorname{rr\text{-}ms}, m), \operatorname{cc\text{-}ms}, g, x) \wedge \\ \quad \forall y \in m.\; (g(y) \in u) \wedge \\ \quad \operatorname{is\text{-}diff\text{-}on}(\operatorname{cc\text{-}normed\text{-}field}, u, f, g(x), l) \wedge \\ \quad d \in \mathbb{C} \wedge \\ \quad \operatorname{has\text{-}deriv\text{-}at}((\lambda t \in m.\; \operatorname{real\text{-}part}(g(t))), x, \operatorname{real\text{-}part}(d)) \wedge \\ \quad \operatorname{has\text{-}deriv\text{-}at}((\lambda t \in m.\; \operatorname{imag\text{-}part}(g(t))), x, \operatorname{imag\text{-}part}(d)) \Rightarrow \\ \quad \quad \operatorname{has\text{-}deriv\text{-}at}((\lambda t \in m.\; \operatorname{real\text{-}part}(f(g(t)))), x, \operatorname{real\text{-}part}((l \cdot d))) \wedge \\ \quad \quad \operatorname{has\text{-}deriv\text{-}at}((\lambda t \in m.\; \operatorname{imag\text{-}part}(f(g(t)))), x, \operatorname{imag\text{-}part}((l \cdot d))) \end{array}}}
\vnbprop{line-int-of-derivative}{\vnbstmt{\begin{array}{@{}l@{}} \forall u, f, c, m, d, a, b. \\ \quad \operatorname{holomorphic\text{-}on}(u, f) \wedge \\ \quad c \in (u \to \mathbb{C}) \wedge \\ \quad \forall z \in u.\; \operatorname{is\text{-}diff\text{-}on}(\operatorname{cc\text{-}normed\text{-}field}, u, f, z, c(z)) \wedge \\ \quad \operatorname{is\text{-}road}(m, d, a, b) \wedge \\ \quad \operatorname{trace}(m, a, b) \subseteq u \Rightarrow \\ \quad \quad \operatorname{line\text{-}int}(c, m, d, a, b) = (f(m(b)) - f(m(a))) \end{array}}}

\vnbbookfile{theorem-library/goursat-bricks}{theorem-library/goursat-bricks}
\noindent{\small goursat-bricks.scm -- the definition-free bricks of Cauchy-Goursat (complex-analysis.pdf ch. 3, s. 3.2).  Batch 29-B, 2026-09-25.}\par
\vnbprop{quarter-step}{\vnbstmt{\begin{array}{@{}l@{}} \forall j, a, b, c, d \in \mathbb{R}. \\ \quad ((0 \leq a) \wedge ((0 \leq b) \wedge ((0 \leq c) \wedge ((0 \leq d) \wedge (j \leq (((a + b) + c) + d)))))) \Rightarrow  \\ \quad \quad (j \leq (4 \cdot a)) \vee ((j \leq (4 \cdot b)) \vee ((j \leq (4 \cdot c)) \vee (j \leq (4 \cdot d)))) \end{array}}}
\vnbprop{four-power-halving}{\vnbstmt{\forall n \in \mathbb{N}.\; ((\operatorname{power}(4, n) \cdot (\operatorname{recip}(\operatorname{power}(2, n)) \cdot \operatorname{recip}(\operatorname{power}(2, n)))) = 1)}}
\vnbprop{line-int-closed-road-of-derivative}{\vnbstmt{\begin{array}{@{}l@{}} \forall u, f, c, m, d, a, b. \\ \quad \operatorname{holomorphic\text{-}on}(u, f) \wedge \\ \quad c \in (u \to \mathbb{C}) \wedge \\ \quad \forall z \in u.\; \operatorname{is\text{-}diff\text{-}on}(\operatorname{cc\text{-}normed\text{-}field}, u, f, z, c(z)) \wedge \\ \quad \operatorname{is\text{-}road}(m, d, a, b) \wedge \\ \quad \operatorname{trace}(m, a, b) \subseteq u \wedge \\ \quad m(b) = m(a) \Rightarrow \\ \quad \quad \operatorname{line\text{-}int}(c, m, d, a, b) = 0 \end{array}}}
\vnbprop{affine-has-primitive-cc}{\vnbstmt{\begin{array}{@{}l@{}} \forall p, q, w \in \mathbb{C}, f \in (\mathbb{C} \to \mathbb{C}). \\ \quad \forall z \in \mathbb{C}.\; (f(z) \simeq ((p \cdot z) + ((q \cdot ((z - w) \cdot (z - w))) \cdot \operatorname{recip}(2)))) \Rightarrow  \\ \quad \quad (\operatorname{holomorphic\text{-}on}(\mathbb{C}, f) \wedge \forall z \in \mathbb{C}.\; \operatorname{is\text{-}diff\text{-}on}(\operatorname{cc\text{-}normed\text{-}field}, \mathbb{C}, f, z, (p + (q \cdot (z - w))))) \end{array}}}
\vnbprop{affine-closed-road-integral}{\vnbstmt{\begin{array}{@{}l@{}} \forall p, q, w \in \mathbb{C}, g \in (\mathbb{C} \to \mathbb{C}), m, c, a, b. \\ \quad (\forall z \in \mathbb{C}.\; (g(z) \simeq (p + (q \cdot (z - w)))) \wedge (\operatorname{is\text{-}road}(m, c, a, b) \wedge (m(b) = m(a)))) \Rightarrow  \\ \quad \quad \operatorname{line\text{-}int}(g, m, c, a, b) = 0 \end{array}}}
\vnbprop{caratheodory-remainder-small}{\vnbstmt{\begin{array}{@{}l@{}} \forall s, f, a, l, e. \\ \quad (\operatorname{is\text{-}diff\text{-}on}(\operatorname{cc\text{-}normed\text{-}field}, s, f, a, l) \wedge \operatorname{pos\text{-}rr}(e)) \Rightarrow  \\ \quad \quad \exists r. \\ \quad \quad \quad \operatorname{pos\text{-}rr}(r) \wedge \\ \quad \quad \quad \forall z \in s. \\ \quad \quad \quad \quad (\operatorname{magnitude}((z - a)) < r) \Rightarrow  \\ \quad \quad \quad \quad \quad \operatorname{magnitude}((f(z) - (f(a) + (l \cdot (z - a))))) \leq (e \cdot \operatorname{magnitude}((z - a))) \end{array}}}
\vnbprop{nn-succ-le-power-two}{\vnbstmt{\forall k \in \mathbb{N}.\; ((k + 1) \leq \operatorname{power}(2, k))}}
\vnbprop{cc-telescoping-limit-bound}{\vnbstmt{\begin{array}{@{}l@{}} \forall s \in (\mathbb{N} \to \mathbb{C}), c \in \mathbb{R}. \\ \quad ((0 \leq c) \wedge \forall k \in \mathbb{N}.\; (\operatorname{magnitude}((s((k + 1)) - s(k))) \leq (c \cdot \operatorname{recip}(\operatorname{power}(2, k))))) \Rightarrow  \\ \quad \quad \exists v \in \mathbb{C}. \\ \quad \quad \quad (\operatorname{converges\text{-}to}(\operatorname{cc\text{-}ms}, s, v) \wedge \forall n \in \mathbb{N}.\; (\operatorname{magnitude}((s(n) - v)) \leq ((2 \cdot c) \cdot \operatorname{recip}(\operatorname{power}(2, n))))) \end{array}}}

\vnbbookfile{theorem-library/regulated-algebra}{theorem-library/regulated-algebra}
\noindent{\small regulated-algebra.scm -- THE ALGEBRA OF REGULATED FUNCTIONS ON [a, b], THE INTEGRAND OF A ROAD IS REGULATED, AND THE INTEGRAL ALONG A ROAD EXISTS. Batch 23-A, 2026-09-23.}\par
\vnbprop{right-limit-within-sum}{\vnbstmt{\begin{array}{@{}l@{}} \forall f, g, h, w, x, l, m. \\ \quad \operatorname{is\text{-}right\text{-}limit\text{-}within}(f, w, x, l) \wedge \\ \quad \operatorname{is\text{-}right\text{-}limit\text{-}within}(g, w, x, m) \wedge \\ \quad h \in (w \to \mathbb{R}) \wedge \\ \quad \forall z \in w.\; (h(z) \simeq (f(z) + g(z))) \Rightarrow \\ \quad \quad \operatorname{is\text{-}right\text{-}limit\text{-}within}(h, w, x, (l + m)) \end{array}}}
\vnbprop{left-limit-within-sum}{\vnbstmt{\begin{array}{@{}l@{}} \forall f, g, h, w, x, l, m. \\ \quad \operatorname{is\text{-}left\text{-}limit\text{-}within}(f, w, x, l) \wedge \\ \quad \operatorname{is\text{-}left\text{-}limit\text{-}within}(g, w, x, m) \wedge \\ \quad h \in (w \to \mathbb{R}) \wedge \\ \quad \forall z \in w.\; (h(z) \simeq (f(z) + g(z))) \Rightarrow \\ \quad \quad \operatorname{is\text{-}left\text{-}limit\text{-}within}(h, w, x, (l + m)) \end{array}}}
\vnbprop{right-limit-within-sub}{\vnbstmt{\begin{array}{@{}l@{}} \forall f, g, h, w, x, l, m. \\ \quad \operatorname{is\text{-}right\text{-}limit\text{-}within}(f, w, x, l) \wedge \\ \quad \operatorname{is\text{-}right\text{-}limit\text{-}within}(g, w, x, m) \wedge \\ \quad h \in (w \to \mathbb{R}) \wedge \\ \quad \forall z \in w.\; (h(z) \simeq (f(z) - g(z))) \Rightarrow \\ \quad \quad \operatorname{is\text{-}right\text{-}limit\text{-}within}(h, w, x, (l - m)) \end{array}}}
\vnbprop{left-limit-within-sub}{\vnbstmt{\begin{array}{@{}l@{}} \forall f, g, h, w, x, l, m. \\ \quad \operatorname{is\text{-}left\text{-}limit\text{-}within}(f, w, x, l) \wedge \\ \quad \operatorname{is\text{-}left\text{-}limit\text{-}within}(g, w, x, m) \wedge \\ \quad h \in (w \to \mathbb{R}) \wedge \\ \quad \forall z \in w.\; (h(z) \simeq (f(z) - g(z))) \Rightarrow \\ \quad \quad \operatorname{is\text{-}left\text{-}limit\text{-}within}(h, w, x, (l - m)) \end{array}}}
\vnbprop{right-limit-within-mul}{\vnbstmt{\begin{array}{@{}l@{}} \forall f, g, h, w, x, l, m. \\ \quad \operatorname{is\text{-}right\text{-}limit\text{-}within}(f, w, x, l) \wedge \\ \quad \operatorname{is\text{-}right\text{-}limit\text{-}within}(g, w, x, m) \wedge \\ \quad h \in (w \to \mathbb{R}) \wedge \\ \quad \forall z \in w.\; (h(z) \simeq (f(z) \cdot g(z))) \Rightarrow \\ \quad \quad \operatorname{is\text{-}right\text{-}limit\text{-}within}(h, w, x, (l \cdot m)) \end{array}}}
\vnbprop{left-limit-within-mul}{\vnbstmt{\begin{array}{@{}l@{}} \forall f, g, h, w, x, l, m. \\ \quad \operatorname{is\text{-}left\text{-}limit\text{-}within}(f, w, x, l) \wedge \\ \quad \operatorname{is\text{-}left\text{-}limit\text{-}within}(g, w, x, m) \wedge \\ \quad h \in (w \to \mathbb{R}) \wedge \\ \quad \forall z \in w.\; (h(z) \simeq (f(z) \cdot g(z))) \Rightarrow \\ \quad \quad \operatorname{is\text{-}left\text{-}limit\text{-}within}(h, w, x, (l \cdot m)) \end{array}}}
\vnbprop{right-limit-within-scalar}{\vnbstmt{\begin{array}{@{}l@{}} \forall f, h, w, x, l.; \forall c \in \mathbb{R}. \\ \quad (\operatorname{is\text{-}right\text{-}limit\text{-}within}(f, w, x, l) \wedge ((h \in (w \to \mathbb{R})) \wedge \forall z \in w.\; (h(z) \simeq (c \cdot f(z))))) \Rightarrow  \\ \quad \quad \operatorname{is\text{-}right\text{-}limit\text{-}within}(h, w, x, (c \cdot l)) \end{array}}}
\vnbprop{left-limit-within-scalar}{\vnbstmt{\begin{array}{@{}l@{}} \forall f, h, w, x, l.; \forall c \in \mathbb{R}. \\ \quad (\operatorname{is\text{-}left\text{-}limit\text{-}within}(f, w, x, l) \wedge ((h \in (w \to \mathbb{R})) \wedge \forall z \in w.\; (h(z) \simeq (c \cdot f(z))))) \Rightarrow  \\ \quad \quad \operatorname{is\text{-}left\text{-}limit\text{-}within}(h, w, x, (c \cdot l)) \end{array}}}
\vnbprop{regulated-on-sum}{\vnbstmt{\begin{array}{@{}l@{}} \forall f, g, h, a, b. \\ \quad \operatorname{is\text{-}regulated\text{-}on}(f, a, b) \wedge \\ \quad \operatorname{is\text{-}regulated\text{-}on}(g, a, b) \wedge \\ \quad h \in (\operatorname{ccint}(a, b) \to \mathbb{R}) \wedge \\ \quad \forall z \in \operatorname{ccint}(a, b).\; (h(z) \simeq (f(z) + g(z))) \Rightarrow \\ \quad \quad \operatorname{is\text{-}regulated\text{-}on}(h, a, b) \end{array}}}
\vnbprop{regulated-on-sub}{\vnbstmt{\begin{array}{@{}l@{}} \forall f, g, h, a, b. \\ \quad \operatorname{is\text{-}regulated\text{-}on}(f, a, b) \wedge \\ \quad \operatorname{is\text{-}regulated\text{-}on}(g, a, b) \wedge \\ \quad h \in (\operatorname{ccint}(a, b) \to \mathbb{R}) \wedge \\ \quad \forall z \in \operatorname{ccint}(a, b).\; (h(z) \simeq (f(z) - g(z))) \Rightarrow \\ \quad \quad \operatorname{is\text{-}regulated\text{-}on}(h, a, b) \end{array}}}
\vnbprop{regulated-on-mul}{\vnbstmt{\begin{array}{@{}l@{}} \forall f, g, h, a, b. \\ \quad \operatorname{is\text{-}regulated\text{-}on}(f, a, b) \wedge \\ \quad \operatorname{is\text{-}regulated\text{-}on}(g, a, b) \wedge \\ \quad h \in (\operatorname{ccint}(a, b) \to \mathbb{R}) \wedge \\ \quad \forall z \in \operatorname{ccint}(a, b).\; (h(z) \simeq (f(z) \cdot g(z))) \Rightarrow \\ \quad \quad \operatorname{is\text{-}regulated\text{-}on}(h, a, b) \end{array}}}
\vnbprop{regulated-on-scalar}{\vnbstmt{\begin{array}{@{}l@{}} \forall f, h, a, b.; \forall c \in \mathbb{R}. \\ \quad \operatorname{is\text{-}regulated\text{-}on}(f, a, b) \wedge \\ \quad h \in (\operatorname{ccint}(a, b) \to \mathbb{R}) \wedge \\ \quad \forall z \in \operatorname{ccint}(a, b).\; (h(z) \simeq (c \cdot f(z))) \Rightarrow \\ \quad \quad \operatorname{is\text{-}regulated\text{-}on}(h, a, b) \end{array}}}
\vnbprop{regulated-on-congruence}{\vnbstmt{\begin{array}{@{}l@{}} \forall f, g, a, b. \\ \quad (\operatorname{is\text{-}regulated\text{-}on}(g, a, b) \wedge ((f \in (\operatorname{ccint}(a, b) \to \mathbb{R})) \wedge \forall z \in \operatorname{ccint}(a, b).\; (f(z) \simeq g(z)))) \Rightarrow  \\ \quad \quad \operatorname{is\text{-}regulated\text{-}on}(f, a, b) \end{array}}}
\vnbprop{continuous-on-compose-path}{\vnbstmt{\begin{array}{@{}l@{}} \forall u, f, m, a, b. \\ \quad u \subseteq \mathbb{C} \wedge \\ \quad f \in (u \to \mathbb{C}) \wedge \\ \quad \forall y \in u. \\ \quad \quad \operatorname{is\text{-}continuous\text{-}at}(\operatorname{subspace\text{-}ms}(\operatorname{nf\text{-}metric\text{-}space}(\operatorname{cc\text{-}normed\text{-}field}), u), \\ \quad \quad \quad \operatorname{nf\text{-}metric\text{-}space}(\operatorname{cc\text{-}normed\text{-}field}), f, y) \wedge \\ \quad \operatorname{is\text{-}path}(m, a, b) \wedge \\ \quad \operatorname{trace}(m, a, b) \subseteq u \Rightarrow \\ \quad \quad \operatorname{is\text{-}continuous\text{-}on}((\lambda t \in \operatorname{ccint}(a, b).\; \operatorname{real\text{-}part}(f(m(t)))), \operatorname{ccint}(a, b)) \wedge \\ \quad \quad \operatorname{is\text{-}continuous\text{-}on}((\lambda t \in \operatorname{ccint}(a, b).\; \operatorname{imag\text{-}part}(f(m(t)))), \operatorname{ccint}(a, b)) \end{array}}}
\vnbprop{continuous-on-compose-regulated}{\vnbstmt{\begin{array}{@{}l@{}} \forall u, f, m, a, b. \\ \quad u \subseteq \mathbb{C} \wedge \\ \quad f \in (u \to \mathbb{C}) \wedge \\ \quad \forall y \in u. \\ \quad \quad \operatorname{is\text{-}continuous\text{-}at}(\operatorname{subspace\text{-}ms}(\operatorname{nf\text{-}metric\text{-}space}(\operatorname{cc\text{-}normed\text{-}field}), u), \\ \quad \quad \quad \operatorname{nf\text{-}metric\text{-}space}(\operatorname{cc\text{-}normed\text{-}field}), f, y) \wedge \\ \quad \operatorname{is\text{-}path}(m, a, b) \wedge \\ \quad \operatorname{trace}(m, a, b) \subseteq u \Rightarrow \\ \quad \quad \operatorname{is\text{-}regulated\text{-}on}((\lambda t \in \operatorname{ccint}(a, b).\; \operatorname{real\text{-}part}(f(m(t)))), a, b) \wedge \\ \quad \quad \operatorname{is\text{-}regulated\text{-}on}((\lambda t \in \operatorname{ccint}(a, b).\; \operatorname{imag\text{-}part}(f(m(t)))), a, b) \end{array}}}
\vnbprop{road-integrand-regulated}{\vnbstmt{\begin{array}{@{}l@{}} \forall u, f, m, c, a, b. \\ \quad u \subseteq \mathbb{C} \wedge \\ \quad f \in (u \to \mathbb{C}) \wedge \\ \quad \forall y \in u. \\ \quad \quad \operatorname{is\text{-}continuous\text{-}at}(\operatorname{subspace\text{-}ms}(\operatorname{nf\text{-}metric\text{-}space}(\operatorname{cc\text{-}normed\text{-}field}), u), \\ \quad \quad \quad \operatorname{nf\text{-}metric\text{-}space}(\operatorname{cc\text{-}normed\text{-}field}), f, y) \wedge \\ \quad \operatorname{is\text{-}road}(m, c, a, b) \wedge \\ \quad \operatorname{trace}(m, a, b) \subseteq u \Rightarrow \\ \quad \quad \operatorname{is\text{-}regulated\text{-}on}((\lambda t \in \operatorname{ccint}(a, b).\; \operatorname{real\text{-}part}((f(m(t)) \cdot c(t)))), a, b) \wedge \\ \quad \quad \operatorname{is\text{-}regulated\text{-}on}((\lambda t \in \operatorname{ccint}(a, b).\; \operatorname{imag\text{-}part}((f(m(t)) \cdot c(t)))), a, b) \end{array}}}
\vnbprop{line-int-exists}{\vnbstmt{\begin{array}{@{}l@{}} \forall u, f, m, c, a, b. \\ \quad u \subseteq \mathbb{C} \wedge \\ \quad f \in (u \to \mathbb{C}) \wedge \\ \quad \forall y \in u. \\ \quad \quad \operatorname{is\text{-}continuous\text{-}at}(\operatorname{subspace\text{-}ms}(\operatorname{nf\text{-}metric\text{-}space}(\operatorname{cc\text{-}normed\text{-}field}), u), \\ \quad \quad \quad \operatorname{nf\text{-}metric\text{-}space}(\operatorname{cc\text{-}normed\text{-}field}), f, y) \wedge \\ \quad \operatorname{is\text{-}road}(m, c, a, b) \wedge \\ \quad \operatorname{trace}(m, a, b) \subseteq u \Rightarrow \\ \quad \quad \exists \mathit{pwf}, w. \\ \quad \quad \quad \operatorname{is\text{-}primitive}(\mathit{pwf}, \\ \quad \quad \quad \quad (\lambda t \in \operatorname{ccint}(a, b).\; \operatorname{real\text{-}part}((\lambda t \in \operatorname{ccint}(a, b).\; (f(m(t)) \cdot c(t)))(t))), \\ \quad \quad \quad \quad a, b) \wedge \\ \quad \quad \quad \operatorname{is\text{-}primitive}(w, \\ \quad \quad \quad \quad (\lambda t \in \operatorname{ccint}(a, b).\; \operatorname{imag\text{-}part}((\lambda t \in \operatorname{ccint}(a, b).\; (f(m(t)) \cdot c(t)))(t))), \\ \quad \quad \quad \quad a, b) \end{array}}}
\vnbprop{line-int-in-cc}{\vnbstmt{\begin{array}{@{}l@{}} \forall u, f, m, c, a, b. \\ \quad u \subseteq \mathbb{C} \wedge \\ \quad f \in (u \to \mathbb{C}) \wedge \\ \quad \forall y \in u. \\ \quad \quad \operatorname{is\text{-}continuous\text{-}at}(\operatorname{subspace\text{-}ms}(\operatorname{nf\text{-}metric\text{-}space}(\operatorname{cc\text{-}normed\text{-}field}), u), \\ \quad \quad \quad \operatorname{nf\text{-}metric\text{-}space}(\operatorname{cc\text{-}normed\text{-}field}), f, y) \wedge \\ \quad \operatorname{is\text{-}road}(m, c, a, b) \wedge \\ \quad \operatorname{trace}(m, a, b) \subseteq u \Rightarrow \\ \quad \quad \operatorname{line\text{-}int}(f, m, c, a, b) \in \mathbb{C} \end{array}}}

\vnbbookfile{theorem-library/cc-int-estimate}{theorem-library/cc-int-estimate}
\noindent{\small cc-int-estimate.scm -- THE NOTES' ESTIMATE (45).}\par
\vnbprop{cc-magnitude-sq}{\vnbstmt{\begin{array}{@{}l@{}} \forall z \in \mathbb{C}. \\ \quad (\operatorname{magnitude}(z) \cdot \operatorname{magnitude}(z)) = ((\operatorname{real\text{-}part}(z) \cdot \operatorname{real\text{-}part}(z)) + (\operatorname{imag\text{-}part}(z) \cdot \operatorname{imag\text{-}part}(z))) \end{array}}}
\vnbprop{cc-re-inner-bound}{\vnbstmt{\begin{array}{@{}l@{}} \forall w, z \in \mathbb{C}. \\ \quad ((\operatorname{real\text{-}part}(w) \cdot \operatorname{real\text{-}part}(z)) + (\operatorname{imag\text{-}part}(w) \cdot \operatorname{imag\text{-}part}(z))) \leq (\operatorname{magnitude}(w) \cdot \operatorname{magnitude}(z)) \end{array}}}
\vnbprop{cc-int-abs-bound}{\vnbstmt{\begin{array}{@{}l@{}} \forall i, f, w, s, h, a, b. \\ \quad i \in (\operatorname{ccint}(a, b) \to \mathbb{C}) \wedge \\ \quad \operatorname{is\text{-}primitive}(f, (\lambda t \in \operatorname{ccint}(a, b).\; \operatorname{real\text{-}part}(i(t))), a, b) \wedge \\ \quad \operatorname{is\text{-}primitive}(w, (\lambda t \in \operatorname{ccint}(a, b).\; \operatorname{imag\text{-}part}(i(t))), a, b) \wedge \\ \quad \operatorname{is\text{-}primitive}(h, s, a, b) \wedge \\ \quad \forall y \in \operatorname{ccint}(a, b).\; (s(y) \simeq \operatorname{magnitude}(i(y))) \Rightarrow \\ \quad \quad \operatorname{magnitude}(\int_{a}^{b} i) \leq \int_{a}^{b} s \end{array}}}

\vnbbookfile{theorem-library/line-int-laws}{theorem-library/line-int-laws}
\noindent{\small line-int-laws.scm -- THE LAWS OF THE INTEGRAL ALONG A ROAD, for a road and a CONTINUOUS f, with no primitive in the statement. Batch 24-A, 2026-09-23.}\par
\vnbprop{line-int-value}{\vnbstmt{\begin{array}{@{}l@{}} \forall f, m, c, a, b, g_{1}, g_{2}. \\ \quad \operatorname{is\text{-}primitive}(g_{1}, (\lambda t \in \operatorname{ccint}(a, b).\; \operatorname{real\text{-}part}((f(m(t)) \cdot c(t)))), a, b) \wedge \\ \quad \operatorname{is\text{-}primitive}(g_{2}, (\lambda t \in \operatorname{ccint}(a, b).\; \operatorname{imag\text{-}part}((f(m(t)) \cdot c(t)))), a, b) \Rightarrow \\ \quad \quad \operatorname{line\text{-}int}(f, m, c, a, b) = ((\operatorname{g1}(b) - \operatorname{g1}(a)) + ((\operatorname{g2}(b) - \operatorname{g2}(a)) \cdot +i)) \end{array}}}
\vnbprop{line-int-sum}{\vnbstmt{\begin{array}{@{}l@{}} \forall u, f, g, m, c, a, b.; \forall h \in (u \to \mathbb{C}). \\ \quad u \subseteq \mathbb{C} \wedge \\ \quad f \in (u \to \mathbb{C}) \wedge \\ \quad \forall y \in u. \\ \quad \quad \operatorname{is\text{-}continuous\text{-}at}(\operatorname{subspace\text{-}ms}(\operatorname{nf\text{-}metric\text{-}space}(\operatorname{cc\text{-}normed\text{-}field}), u), \\ \quad \quad \quad \operatorname{nf\text{-}metric\text{-}space}(\operatorname{cc\text{-}normed\text{-}field}), f, y) \wedge \\ \quad g \in (u \to \mathbb{C}) \wedge \\ \quad \forall y \in u. \\ \quad \quad \operatorname{is\text{-}continuous\text{-}at}(\operatorname{subspace\text{-}ms}(\operatorname{nf\text{-}metric\text{-}space}(\operatorname{cc\text{-}normed\text{-}field}), u), \\ \quad \quad \quad \operatorname{nf\text{-}metric\text{-}space}(\operatorname{cc\text{-}normed\text{-}field}), g, y) \wedge \\ \quad \operatorname{is\text{-}road}(m, c, a, b) \wedge \\ \quad \operatorname{trace}(m, a, b) \subseteq u \wedge \\ \quad \forall z \in u.\; (h(z) \simeq (f(z) + g(z))) \Rightarrow \\ \quad \quad \operatorname{line\text{-}int}(h, m, c, a, b) = (\operatorname{line\text{-}int}(f, m, c, a, b) + \operatorname{line\text{-}int}(g, m, c, a, b)) \end{array}}}
\vnbprop{line-int-scalar}{\vnbstmt{\begin{array}{@{}l@{}} \forall u, f, m, c, a, b.; \forall \mathit{lic} \in \mathbb{C}, h \in (u \to \mathbb{C}). \\ \quad u \subseteq \mathbb{C} \wedge \\ \quad f \in (u \to \mathbb{C}) \wedge \\ \quad \forall y \in u. \\ \quad \quad \operatorname{is\text{-}continuous\text{-}at}(\operatorname{subspace\text{-}ms}(\operatorname{nf\text{-}metric\text{-}space}(\operatorname{cc\text{-}normed\text{-}field}), u), \\ \quad \quad \quad \operatorname{nf\text{-}metric\text{-}space}(\operatorname{cc\text{-}normed\text{-}field}), f, y) \wedge \\ \quad \operatorname{is\text{-}road}(m, c, a, b) \wedge \\ \quad \operatorname{trace}(m, a, b) \subseteq u \wedge \\ \quad \forall z \in u.\; (h(z) \simeq (\mathit{lic} \cdot f(z))) \Rightarrow \\ \quad \quad \operatorname{line\text{-}int}(h, m, c, a, b) = (\mathit{lic} \cdot \operatorname{line\text{-}int}(f, m, c, a, b)) \end{array}}}
\vnbprop{line-int-opposite-road}{\vnbstmt{\begin{array}{@{}l@{}} \forall u, f, m, c, a, b, g, d. \\ \quad u \subseteq \mathbb{C} \wedge \\ \quad f \in (u \to \mathbb{C}) \wedge \\ \quad \forall y \in u. \\ \quad \quad \operatorname{is\text{-}continuous\text{-}at}(\operatorname{subspace\text{-}ms}(\operatorname{nf\text{-}metric\text{-}space}(\operatorname{cc\text{-}normed\text{-}field}), u), \\ \quad \quad \quad \operatorname{nf\text{-}metric\text{-}space}(\operatorname{cc\text{-}normed\text{-}field}), f, y) \wedge \\ \quad \operatorname{is\text{-}road}(m, c, a, b) \wedge \\ \quad \operatorname{trace}(m, a, b) \subseteq u \wedge \\ \quad g \in (\operatorname{ccint}(a, b) \to \mathbb{C}) \wedge \\ \quad d \in (\operatorname{ccint}(a, b) \to \mathbb{C}) \wedge \\ \quad \forall \mathit{pay} \in \operatorname{ccint}(a, b).\; (g(\mathit{pay}) \simeq m(((a + b) - \mathit{pay}))) \wedge \\ \quad \forall \mathit{pay} \in \operatorname{ccint}(a, b).\; (d(\mathit{pay}) \simeq (-1 \cdot c(((a + b) - \mathit{pay})))) \Rightarrow \\ \quad \quad \operatorname{line\text{-}int}(f, g, d, a, b) = (-1 \cdot \operatorname{line\text{-}int}(f, m, c, a, b)) \end{array}}}
\vnbprop{road-restrict}{\vnbstmt{\begin{array}{@{}l@{}} \forall m, c, a, b.; \forall v, \mathit{lidv} \in \mathbb{R}. \\ \quad (\operatorname{is\text{-}road}(m, c, a, b) \wedge ((a \leq v) \wedge ((v < \mathit{lidv}) \wedge (\mathit{lidv} \leq b)))) \Rightarrow  \\ \quad \quad \operatorname{is\text{-}road}(\operatorname{restrict}(m, \operatorname{ccint}(v, \mathit{lidv})), \operatorname{restrict}(c, \operatorname{ccint}(v, \mathit{lidv})), v, \mathit{lidv}) \end{array}}}
\vnbprop{cc-int-congruence}{\vnbstmt{\begin{array}{@{}l@{}} \forall i, \mathit{ppsi}, a, b, g_{1}, g_{2}. \\ \quad \operatorname{is\text{-}primitive}(g_{1}, (\lambda t \in \operatorname{ccint}(a, b).\; \operatorname{real\text{-}part}(i(t))), a, b) \wedge \\ \quad \operatorname{is\text{-}primitive}(g_{2}, (\lambda t \in \operatorname{ccint}(a, b).\; \operatorname{imag\text{-}part}(i(t))), a, b) \wedge \\ \quad \mathit{ppsi} \in (\operatorname{ccint}(a, b) \to \mathbb{C}) \wedge \\ \quad \forall y \in \operatorname{ccint}(a, b).\; (i(y) \simeq \operatorname{ppsi}(y)) \Rightarrow \\ \quad \quad \int_{a}^{b} i = \int_{a}^{b} \mathit{ppsi} \end{array}}}
\vnbprop{line-int-adjacent}{\vnbstmt{\begin{array}{@{}l@{}} \forall u, f, m, c, a, b.; \forall v \in \operatorname{ooint}(a, b). \\ \quad u \subseteq \mathbb{C} \wedge \\ \quad f \in (u \to \mathbb{C}) \wedge \\ \quad \forall y \in u. \\ \quad \quad \operatorname{is\text{-}continuous\text{-}at}(\operatorname{subspace\text{-}ms}(\operatorname{nf\text{-}metric\text{-}space}(\operatorname{cc\text{-}normed\text{-}field}), u), \\ \quad \quad \quad \operatorname{nf\text{-}metric\text{-}space}(\operatorname{cc\text{-}normed\text{-}field}), f, y) \wedge \\ \quad \operatorname{is\text{-}road}(m, c, a, b) \wedge \\ \quad \operatorname{trace}(m, a, b) \subseteq u \Rightarrow \\ \quad \quad \operatorname{line\text{-}int}(f, m, c, a, b) \\ \quad \quad =\; \operatorname{line\text{-}int}(f, \operatorname{restrict}(m, \operatorname{ccint}(a, v)), \operatorname{restrict}(c, \operatorname{ccint}(a, v)), a, v) + \operatorname{line\text{-}int}(f, \operatorname{restrict}(m, \operatorname{ccint}(v, b)), \operatorname{restrict}(c, \operatorname{ccint}(v, b)), v, b) \end{array}}}
\vnbprop{right-limit-within-magnitude}{\vnbstmt{\begin{array}{@{}l@{}} \forall f, g, h, w, x, l, k. \\ \quad \operatorname{is\text{-}right\text{-}limit\text{-}within}(f, w, x, l) \wedge \\ \quad \operatorname{is\text{-}right\text{-}limit\text{-}within}(g, w, x, k) \wedge \\ \quad h \in (w \to \mathbb{R}) \wedge \\ \quad \forall z \in w.\; (h(z) \simeq \operatorname{magnitude}((f(z) + (g(z) \cdot +i)))) \Rightarrow \\ \quad \quad \operatorname{is\text{-}right\text{-}limit\text{-}within}(h, w, x, \operatorname{magnitude}((l + (k \cdot +i)))) \end{array}}}
\vnbprop{left-limit-within-magnitude}{\vnbstmt{\begin{array}{@{}l@{}} \forall f, g, h, w, x, l, k. \\ \quad \operatorname{is\text{-}left\text{-}limit\text{-}within}(f, w, x, l) \wedge \\ \quad \operatorname{is\text{-}left\text{-}limit\text{-}within}(g, w, x, k) \wedge \\ \quad h \in (w \to \mathbb{R}) \wedge \\ \quad \forall z \in w.\; (h(z) \simeq \operatorname{magnitude}((f(z) + (g(z) \cdot +i)))) \Rightarrow \\ \quad \quad \operatorname{is\text{-}left\text{-}limit\text{-}within}(h, w, x, \operatorname{magnitude}((l + (k \cdot +i)))) \end{array}}}
\vnbprop{regulated-on-magnitude}{\vnbstmt{\begin{array}{@{}l@{}} \forall f, g, h, a, b. \\ \quad \operatorname{is\text{-}regulated\text{-}on}(f, a, b) \wedge \\ \quad \operatorname{is\text{-}regulated\text{-}on}(g, a, b) \wedge \\ \quad h \in (\operatorname{ccint}(a, b) \to \mathbb{R}) \wedge \\ \quad \forall z \in \operatorname{ccint}(a, b).\; (h(z) \simeq \operatorname{magnitude}((f(z) + (g(z) \cdot +i)))) \Rightarrow \\ \quad \quad \operatorname{is\text{-}regulated\text{-}on}(h, a, b) \end{array}}}
\vnbprop{line-int-abs-bound}{\vnbstmt{\begin{array}{@{}l@{}} \forall u, f, m, c, a, b.; \forall \mathit{liab} \in (\operatorname{ccint}(a, b) \to \mathbb{R}). \\ \quad u \subseteq \mathbb{C} \wedge \\ \quad f \in (u \to \mathbb{C}) \wedge \\ \quad \forall y \in u. \\ \quad \quad \operatorname{is\text{-}continuous\text{-}at}(\operatorname{subspace\text{-}ms}(\operatorname{nf\text{-}metric\text{-}space}(\operatorname{cc\text{-}normed\text{-}field}), u), \\ \quad \quad \quad \operatorname{nf\text{-}metric\text{-}space}(\operatorname{cc\text{-}normed\text{-}field}), f, y) \wedge \\ \quad \operatorname{is\text{-}road}(m, c, a, b) \wedge \\ \quad \operatorname{trace}(m, a, b) \subseteq u \wedge \\ \quad \forall \mathit{liy} \in \operatorname{ccint}(a, b). \\ \quad \quad (\operatorname{liab}(\mathit{liy}) \simeq (\operatorname{magnitude}(f(m(\mathit{liy}))) \cdot \operatorname{magnitude}(c(\mathit{liy})))) \Rightarrow \\ \quad \quad \operatorname{magnitude}(\operatorname{line\text{-}int}(f, m, c, a, b)) \leq \int_{a}^{b} \mathit{liab} \end{array}}}

\vnbbookfile{theorem-library/line-int-reparam}{theorem-library/line-int-reparam}
\noindent{\small line-int-reparam.scm -- THE AFFINE CHANGE OF PARAMETER (the notes' (48), affine case) AND THE SEGMENT LEMMAS OF GOURSAT'S DEFINITION-FREE HALF. Batch 29-A, 2026-09-25.}\par
\vnbprop{affine-reparam-ccint-in}{\vnbstmt{\begin{array}{@{}l@{}} \forall a, b, c, d, l, m \in \mathbb{R}, z \in \operatorname{ccint}(c, d). \\ \quad ((0 < l) \wedge (((m + (l \cdot c)) = a) \wedge ((m + (l \cdot d)) = b))) \Rightarrow ((m + (l \cdot z)) \in \operatorname{ccint}(a, b)) \end{array}}}
\vnbprop{affine-reparam-ooint-in}{\vnbstmt{\begin{array}{@{}l@{}} \forall a, b, c, d, l, m \in \mathbb{R}, z \in \operatorname{ooint}(c, d). \\ \quad ((0 < l) \wedge (((m + (l \cdot c)) = a) \wedge ((m + (l \cdot d)) = b))) \Rightarrow ((m + (l \cdot z)) \in \operatorname{ooint}(a, b)) \end{array}}}
\vnbprop{affine-reparam-solve}{\vnbstmt{\begin{array}{@{}l@{}} \forall a, b, c, d, l, m \in \mathbb{R}, z \in \operatorname{ccint}(a, b), q \in \mathbb{R}. \\ \quad ((0 < l) \wedge (((m + (l \cdot c)) = a) \wedge (((m + (l \cdot d)) = b) \wedge ((z - m) = (q \cdot l))))) \Rightarrow  \\ \quad \quad (q \in \operatorname{ccint}(c, d)) \wedge ((m + (l \cdot q)) = z) \end{array}}}
\vnbprop{affine-reparam-inverse}{\vnbstmt{\begin{array}{@{}l@{}} \forall a, b, c, d, l, m \in \mathbb{R}, z \in \operatorname{ccint}(a, b). \\ \quad ((0 < l) \wedge (((m + (l \cdot c)) = a) \wedge ((m + (l \cdot d)) = b))) \Rightarrow  \\ \quad \quad (((z - m) \cdot \operatorname{recip}(l)) \in \operatorname{ccint}(c, d)) \wedge ((m + (l \cdot ((z - m) \cdot \operatorname{recip}(l)))) = z) \end{array}}}
\vnbprop{affine-reparam-inverse-value}{\vnbstmt{\begin{array}{@{}l@{}} \forall l, m, z \in \mathbb{R}. \\ \quad (0 < l) \Rightarrow ((((m + (l \cdot z)) - m) \cdot \operatorname{recip}(l)) = z) \end{array}}}
\vnbprop{affine-reparam-continuous-at-sub}{\vnbstmt{\begin{array}{@{}l@{}} \forall a, b, c, d, l, m \in \mathbb{R}, f.; \forall g \in (\operatorname{ccint}(c, d) \to \mathbb{R}), t \in \operatorname{ccint}(c, d). \\ \quad 0 < l \wedge \\ \quad (m + (l \cdot c)) = a \wedge \\ \quad (m + (l \cdot d)) = b \wedge \\ \quad \forall z \in \operatorname{ccint}(c, d).\; (g(z) \simeq f((m + (l \cdot z)))) \wedge \\ \quad \operatorname{is\text{-}continuous\text{-}at}(\operatorname{subspace\text{-}ms}(\operatorname{rr\text{-}ms}, \operatorname{ccint}(a, b)), \operatorname{rr\text{-}ms}, f, (m + (l \cdot t))) \Rightarrow \\ \quad \quad \operatorname{is\text{-}continuous\text{-}at}(\operatorname{subspace\text{-}ms}(\operatorname{rr\text{-}ms}, \operatorname{ccint}(c, d)), \operatorname{rr\text{-}ms}, g, t) \end{array}}}
\vnbprop{affine-reparam-continuous-on}{\vnbstmt{\begin{array}{@{}l@{}} \forall a, b, c, d, l, m \in \mathbb{R}, f.; \forall g \in (\operatorname{ccint}(c, d) \to \mathbb{R}). \\ \quad 0 < l \wedge \\ \quad (m + (l \cdot c)) = a \wedge \\ \quad (m + (l \cdot d)) = b \wedge \\ \quad \operatorname{is\text{-}continuous\text{-}on}(f, \operatorname{ccint}(a, b)) \wedge \\ \quad \forall z \in \operatorname{ccint}(c, d).\; (g(z) \simeq f((m + (l \cdot z)))) \Rightarrow \\ \quad \quad \operatorname{is\text{-}continuous\text{-}on}(g, \operatorname{ccint}(c, d)) \end{array}}}
\vnbprop{affine-reparam-image-countable}{\vnbstmt{\begin{array}{@{}l@{}} \forall a, b, c, d, l, m \in \mathbb{R}, e. \\ \quad ((0 < l) \wedge (((m + (l \cdot c)) = a) \wedge (((m + (l \cdot d)) = b) \wedge (\operatorname{is\text{-}countable}(e) \wedge (e \subseteq \operatorname{ccint}(a, b)))))) \Rightarrow  \\ \quad \quad \operatorname{is\text{-}countable}(\operatorname{image}((\lambda v \in \operatorname{ccint}(a, b).\; ((v - m) \cdot \operatorname{recip}(l))), e)) \wedge \\ \quad \quad \operatorname{image}((\lambda v \in \operatorname{ccint}(a, b).\; ((v - m) \cdot \operatorname{recip}(l))), e) \subseteq \operatorname{ccint}(c, d) \end{array}}}
\vnbprop{affine-reparam-not-in-image}{\vnbstmt{\begin{array}{@{}l@{}} \forall a, b, c, d, l, m \in \mathbb{R}, e.; \forall t \in \operatorname{ccint}(c, d). \\ \quad 0 < l \wedge \\ \quad (m + (l \cdot c)) = a \wedge \\ \quad (m + (l \cdot d)) = b \wedge \\ \quad e \subseteq \operatorname{ccint}(a, b) \wedge \\ \quad \neg (t \in \operatorname{image}((\lambda v \in \operatorname{ccint}(a, b).\; ((v - m) \cdot \operatorname{recip}(l))), e)) \Rightarrow \\ \quad \quad \neg ((m + (l \cdot t)) \in e) \end{array}}}
\vnbprop{affine-reparam-regulated}{\vnbstmt{\begin{array}{@{}l@{}} \forall a, b, c, d, l, m \in \mathbb{R}, f.; \forall g \in (\operatorname{ccint}(c, d) \to \mathbb{R}). \\ \quad 0 < l \wedge \\ \quad (m + (l \cdot c)) = a \wedge \\ \quad (m + (l \cdot d)) = b \wedge \\ \quad c < d \wedge \\ \quad \operatorname{is\text{-}regulated\text{-}on}(f, a, b) \wedge \\ \quad \forall z \in \operatorname{ccint}(c, d).\; (g(z) \simeq f((m + (l \cdot z)))) \Rightarrow \\ \quad \quad \operatorname{is\text{-}regulated\text{-}on}(g, c, d) \end{array}}}
\vnbprop{pw-int-affine-subst}{\vnbstmt{\begin{array}{@{}l@{}} \forall a, b, c, d, l, m \in \mathbb{R}, f, i.; \forall h, s \in (\operatorname{ccint}(c, d) \to \mathbb{R}). \\ \quad 0 < l \wedge \\ \quad (m + (l \cdot c)) = a \wedge \\ \quad (m + (l \cdot d)) = b \wedge \\ \quad c < d \wedge \\ \quad \operatorname{is\text{-}primitive}(f, i, a, b) \wedge \\ \quad \forall z \in \operatorname{ccint}(c, d).\; (h(z) \simeq f((m + (l \cdot z)))) \wedge \\ \quad \forall z \in \operatorname{ccint}(c, d).\; (s(z) \simeq (l \cdot i((m + (l \cdot z))))) \Rightarrow \\ \quad \quad \operatorname{is\text{-}primitive}(h, s, c, d) \wedge (\int_{c}^{d} s = \int_{a}^{b} i) \end{array}}}
\vnbprop{cc-int-affine-subst}{\vnbstmt{\begin{array}{@{}l@{}} \forall a, b, c, d, l, m \in \mathbb{R}, i, \mathit{lrpb}, f, w. \\ \quad 0 < l \wedge \\ \quad (m + (l \cdot c)) = a \wedge \\ \quad (m + (l \cdot d)) = b \wedge \\ \quad c < d \wedge \\ \quad i \in (\operatorname{ccint}(a, b) \to \mathbb{C}) \wedge \\ \quad \mathit{lrpb} \in (\operatorname{ccint}(c, d) \to \mathbb{C}) \wedge \\ \quad \forall z \in \operatorname{ccint}(c, d).\; (\operatorname{lrpb}(z) \simeq (l \cdot i((m + (l \cdot z))))) \wedge \\ \quad \operatorname{is\text{-}primitive}(f, (\lambda t \in \operatorname{ccint}(a, b).\; \operatorname{real\text{-}part}(i(t))), a, b) \wedge \\ \quad \operatorname{is\text{-}primitive}(w, (\lambda t \in \operatorname{ccint}(a, b).\; \operatorname{imag\text{-}part}(i(t))), a, b) \Rightarrow \\ \quad \quad \int_{c}^{d} \mathit{lrpb} = \int_{a}^{b} i \end{array}}}
\vnbprop{trace-affine-reparam}{\vnbstmt{\begin{array}{@{}l@{}} \forall a, b, c, d, l, m \in \mathbb{R}, f, r. \\ \quad 0 < l \wedge \\ \quad (m + (l \cdot c)) = a \wedge \\ \quad (m + (l \cdot d)) = b \wedge \\ \quad f \in (\operatorname{ccint}(a, b) \to \mathbb{C}) \wedge \\ \quad r \in (\operatorname{ccint}(c, d) \to \mathbb{C}) \wedge \\ \quad \forall z \in \operatorname{ccint}(c, d).\; (r(z) \simeq f((m + (l \cdot z)))) \Rightarrow \\ \quad \quad (\operatorname{trace}(r, c, d) \subseteq \operatorname{trace}(f, a, b)) \wedge (\operatorname{trace}(f, a, b) \subseteq \operatorname{trace}(r, c, d)) \end{array}}}
\vnbprop{road-affine-reparam}{\vnbstmt{\begin{array}{@{}l@{}} \forall a, b, c, d, l, m \in \mathbb{R}, f, g, r, k. \\ \quad 0 < l \wedge \\ \quad (m + (l \cdot c)) = a \wedge \\ \quad (m + (l \cdot d)) = b \wedge \\ \quad c < d \wedge \\ \quad \operatorname{is\text{-}road}(f, g, a, b) \wedge \\ \quad r \in (\operatorname{ccint}(c, d) \to \mathbb{C}) \wedge \\ \quad k \in (\operatorname{ccint}(c, d) \to \mathbb{C}) \wedge \\ \quad \forall z \in \operatorname{ccint}(c, d).\; (r(z) \simeq f((m + (l \cdot z)))) \wedge \\ \quad \forall z \in \operatorname{ccint}(c, d).\; (k(z) \simeq (l \cdot g((m + (l \cdot z))))) \Rightarrow \\ \quad \quad \operatorname{is\text{-}road}(r, k, c, d) \end{array}}}
\vnbprop{line-int-affine-reparam}{\vnbstmt{\begin{array}{@{}l@{}} \forall u, f, m, c, a, b.; \forall \mathit{lrpc}, d, l, \mathit{lrpm} \in \mathbb{R}, r, k. \\ \quad u \subseteq \mathbb{C} \wedge \\ \quad f \in (u \to \mathbb{C}) \wedge \\ \quad \forall y \in u. \\ \quad \quad \operatorname{is\text{-}continuous\text{-}at}(\operatorname{subspace\text{-}ms}(\operatorname{nf\text{-}metric\text{-}space}(\operatorname{cc\text{-}normed\text{-}field}), u), \\ \quad \quad \quad \operatorname{nf\text{-}metric\text{-}space}(\operatorname{cc\text{-}normed\text{-}field}), f, y) \wedge \\ \quad \operatorname{is\text{-}road}(m, c, a, b) \wedge \\ \quad \operatorname{trace}(m, a, b) \subseteq u \wedge \\ \quad 0 < l \wedge \\ \quad (\mathit{lrpm} + (l \cdot \mathit{lrpc})) = a \wedge \\ \quad (\mathit{lrpm} + (l \cdot d)) = b \wedge \\ \quad \mathit{lrpc} < d \wedge \\ \quad r \in (\operatorname{ccint}(\mathit{lrpc}, d) \to \mathbb{C}) \wedge \\ \quad k \in (\operatorname{ccint}(\mathit{lrpc}, d) \to \mathbb{C}) \wedge \\ \quad \forall z \in \operatorname{ccint}(\mathit{lrpc}, d).\; (r(z) \simeq m((\mathit{lrpm} + (l \cdot z)))) \wedge \\ \quad \forall z \in \operatorname{ccint}(\mathit{lrpc}, d).\; (k(z) \simeq (l \cdot c((\mathit{lrpm} + (l \cdot z))))) \Rightarrow \\ \quad \quad \operatorname{line\text{-}int}(f, r, k, \mathit{lrpc}, d) = \operatorname{line\text{-}int}(f, m, c, a, b) \end{array}}}
\vnbprop{trace-restrict-subset}{\vnbstmt{\begin{array}{@{}l@{}} \forall m, a, b.; \forall c, d \in \mathbb{R}. \\ \quad ((m \in (\operatorname{ccint}(a, b) \to \mathbb{C})) \wedge ((a \in \mathbb{R}) \wedge ((b \in \mathbb{R}) \wedge ((a \leq c) \wedge (d \leq b))))) \Rightarrow  \\ \quad \quad \operatorname{trace}(\operatorname{restrict}(m, \operatorname{ccint}(c, d)), c, d) \subseteq \operatorname{trace}(m, a, b) \end{array}}}
\vnbprop{segment-int-split}{\vnbstmt{\begin{array}{@{}l@{}} \forall u, f.; \forall z_{0}, w \in \mathbb{C}, t \in \operatorname{ooint}(0, 1). \\ \quad u \subseteq \mathbb{C} \wedge \\ \quad f \in (u \to \mathbb{C}) \wedge \\ \quad \forall y \in u. \\ \quad \quad \operatorname{is\text{-}continuous\text{-}at}(\operatorname{subspace\text{-}ms}(\operatorname{nf\text{-}metric\text{-}space}(\operatorname{cc\text{-}normed\text{-}field}), u), \\ \quad \quad \quad \operatorname{nf\text{-}metric\text{-}space}(\operatorname{cc\text{-}normed\text{-}field}), f, y) \wedge \\ \quad \operatorname{trace}((\lambda x \in \operatorname{ccint}(0, 1).\; (z_{0} + (x \cdot w))), 0, 1) \subseteq u \Rightarrow \\ \quad \quad \operatorname{line\text{-}int}(f, (\lambda x \in \operatorname{ccint}(0, 1).\; (z_{0} + (x \cdot w))), (\lambda x \in \operatorname{ccint}(0, 1).\; w), 0, 1) \\ \quad \quad =\; \operatorname{line\text{-}int}(f, (\lambda x \in \operatorname{ccint}(0, 1).\; (z_{0} + (x \cdot (t \cdot w)))), (\lambda x \in \operatorname{ccint}(0, 1).\; (t \cdot w)), 0, 1) + \operatorname{line\text{-}int}(f, (\lambda x \in \operatorname{ccint}(0, 1).\; ((z_{0} + (t \cdot w)) + (x \cdot ((1 - t) \cdot w)))), (\lambda x \in \operatorname{ccint}(0, 1).\; ((1 - t) \cdot w)), 0, 1) \end{array}}}
\vnbprop{segment-int-abs-bound}{\vnbstmt{\begin{array}{@{}l@{}} \forall u, f.; \forall z_{0}, w \in \mathbb{C}, n \in \mathbb{R}. \\ \quad u \subseteq \mathbb{C} \wedge \\ \quad f \in (u \to \mathbb{C}) \wedge \\ \quad \forall y \in u. \\ \quad \quad \operatorname{is\text{-}continuous\text{-}at}(\operatorname{subspace\text{-}ms}(\operatorname{nf\text{-}metric\text{-}space}(\operatorname{cc\text{-}normed\text{-}field}), u), \\ \quad \quad \quad \operatorname{nf\text{-}metric\text{-}space}(\operatorname{cc\text{-}normed\text{-}field}), f, y) \wedge \\ \quad \operatorname{trace}((\lambda x \in \operatorname{ccint}(0, 1).\; (z_{0} + (x \cdot w))), 0, 1) \subseteq u \wedge \\ \quad \forall \mathit{lrpy} \in \operatorname{trace}((\lambda x \in \operatorname{ccint}(0, 1).\; (z_{0} + (x \cdot w))), 0, 1). \\ \quad \quad \operatorname{magnitude}(f(\mathit{lrpy})) \leq n \Rightarrow \\ \quad \quad \operatorname{magnitude}(\operatorname{line\text{-}int}(f, (\lambda x \in \operatorname{ccint}(0, 1).\; (z_{0} + (x \cdot w))), (\lambda x \in \operatorname{ccint}(0, 1).\; w), 0, 1)) \\ \quad \quad \leq\; \operatorname{magnitude}(w) \cdot n \end{array}}}
\vnbprop{segment-int-reverse}{\vnbstmt{\begin{array}{@{}l@{}} \forall u, f.; \forall z_{0}, w \in \mathbb{C}. \\ \quad u \subseteq \mathbb{C} \wedge \\ \quad f \in (u \to \mathbb{C}) \wedge \\ \quad \forall y \in u. \\ \quad \quad \operatorname{is\text{-}continuous\text{-}at}(\operatorname{subspace\text{-}ms}(\operatorname{nf\text{-}metric\text{-}space}(\operatorname{cc\text{-}normed\text{-}field}), u), \\ \quad \quad \quad \operatorname{nf\text{-}metric\text{-}space}(\operatorname{cc\text{-}normed\text{-}field}), f, y) \wedge \\ \quad \operatorname{trace}((\lambda x \in \operatorname{ccint}(0, 1).\; (z_{0} + (x \cdot w))), 0, 1) \subseteq u \Rightarrow \\ \quad \quad \operatorname{line\text{-}int}(f, (\lambda x \in \operatorname{ccint}(0, 1).\; ((z_{0} + w) + (x \cdot (-1 \cdot w)))), \\ \quad \quad \quad (\lambda x \in \operatorname{ccint}(0, 1).\; (-1 \cdot w)), 0, 1) \\ \quad \quad =\; -1 \cdot \operatorname{line\text{-}int}(f, (\lambda x \in \operatorname{ccint}(0, 1).\; (z_{0} + (x \cdot w))), (\lambda x \in \operatorname{ccint}(0, 1).\; w), 0, 1) \end{array}}}

\vnbbookfile{theorem-library/goursat}{theorem-library/goursat}
\noindent{\small goursat.scm -- CAUCHY-GOURSAT FOR A TRIANGLE AND CAUCHY ON A CONVEX SET (complex-analysis.pdf ch. 3, s. 3.2: Lemma 3.4, Lemma 3.5, Prop 3.6, Cor 3.7, Prop 3.8, Cor 3.9, all for f holomorphic on the whole open set).  Batch 30, 2026-09-25.}\par
\vnbprop{seg-int-unfold}{\vnbstmt{\begin{array}{@{}l@{}} \forall f, a, b. \\ \quad (\operatorname{seg\text{-}int}(f, a, b) \simeq \operatorname{line\text{-}int}(f, \operatorname{seg\text{-}path}(a, b), \operatorname{seg\text{-}deriv}(a, b), 0, 1)) \end{array}}}
\vnbprop{tri-int-unfold}{\vnbstmt{\begin{array}{@{}l@{}} \forall f, a, b, c. \\ \quad (\operatorname{tri\text{-}int}(f, a, b, c) \simeq ((\operatorname{seg\text{-}int}(f, a, b) + \operatorname{seg\text{-}int}(f, b, c)) + \operatorname{seg\text{-}int}(f, c, a))) \end{array}}}
\vnbprop{cc-mid-unfold}{\vnbstmt{\forall x, y.\; (\operatorname{cc\text{-}mid}(x, y) \simeq ((x + y) \cdot \operatorname{recip}(2)))}}
\vnbprop{seg-path-bridge}{\vnbstmt{\begin{array}{@{}l@{}} \forall a, b, p, w. \\ \quad ((a = p) \wedge ((b - a) = w)) \Rightarrow  \\ \quad \quad (\operatorname{seg\text{-}path}(a, b) \simeq (\lambda x \in \operatorname{ccint}(0, 1).\; (p + (x \cdot w)))) \end{array}}}
\vnbprop{seg-deriv-bridge}{\vnbstmt{\begin{array}{@{}l@{}} \forall a, b, w. \\ \quad ((b - a) = w) \Rightarrow (\operatorname{seg\text{-}deriv}(a, b) \simeq (\lambda x \in \operatorname{ccint}(0, 1).\; w)) \end{array}}}
\vnbprop{seg-int-bridge}{\vnbstmt{\begin{array}{@{}l@{}} \forall f, a, b, p, w. \\ \quad ((a = p) \wedge ((b - a) = w)) \Rightarrow  \\ \quad \quad (\operatorname{seg\text{-}int}(f, a, b) \simeq \operatorname{line\text{-}int}(f, (\lambda x \in \operatorname{ccint}(0, 1).\; (p + (x \cdot w))), (\lambda x \in \operatorname{ccint}(0, 1).\; w), 0, 1)) \end{array}}}
\vnbprop{seg-trace-bridge}{\vnbstmt{\begin{array}{@{}l@{}} \forall a, b, p, w. \\ \quad ((a = p) \wedge ((b - a) = w)) \Rightarrow  \\ \quad \quad (\operatorname{trace}(\operatorname{seg\text{-}path}(a, b), 0, 1) \simeq \operatorname{trace}((\lambda x \in \operatorname{ccint}(0, 1).\; (p + (x \cdot w))), 0, 1)) \end{array}}}
\vnbprop{seg-path-is-road}{\vnbstmt{\forall a, b \in \mathbb{C}.\; \operatorname{is\text{-}road}(\operatorname{seg\text{-}path}(a, b), \operatorname{seg\text{-}deriv}(a, b), 0, 1)}}
\vnbprop{in-transport-eq}{\vnbstmt{\forall x, y, s.\; (((x \in s) \wedge (x = y)) \Rightarrow (y \in s))}}
\vnbprop{seg-path-apply}{\vnbstmt{\begin{array}{@{}l@{}} \forall a, b \in \mathbb{C}, t \in \operatorname{ccint}(0, 1). \\ \quad \operatorname{seg\text{-}path}(a, b)(t) = (a + (t \cdot (b - a))) \end{array}}}
\vnbprop{seg-trace-intro}{\vnbstmt{\begin{array}{@{}l@{}} \forall a, b \in \mathbb{C}, t \in \operatorname{ccint}(0, 1). \\ \quad (a + (t \cdot (b - a))) \in \operatorname{trace}(\operatorname{seg\text{-}path}(a, b), 0, 1) \end{array}}}
\vnbprop{seg-trace-elim}{\vnbstmt{\begin{array}{@{}l@{}} \forall a, b \in \mathbb{C}, z \in \operatorname{trace}(\operatorname{seg\text{-}path}(a, b), 0, 1). \exists t \in \operatorname{ccint}(0, 1). \\ \quad z = (a + (t \cdot (b - a))) \end{array}}}
\vnbprop{conv3-member-iff}{\vnbstmt{\begin{array}{@{}l@{}} \forall a, b, c, z. \\ \quad z \in \operatorname{conv3}(a, b, c) \\ \quad \Leftrightarrow\; z \in \mathbb{C} \wedge \\ \quad \exists l, m, n \in \operatorname{ccint}(0, 1). \\ \quad \quad (((l + m) + n) = 1) \wedge (z = (((l \cdot a) + (m \cdot b)) + (n \cdot c))) \end{array}}}
\vnbprop{conv3-intro}{\vnbstmt{\begin{array}{@{}l@{}} \forall a, b, c \in \mathbb{C}, l, m, n \in \operatorname{ccint}(0, 1), z. \\ \quad ((((l + m) + n) = 1) \wedge (z = (((l \cdot a) + (m \cdot b)) + (n \cdot c)))) \Rightarrow (z \in \operatorname{conv3}(a, b, c)) \end{array}}}
\vnbprop{is-convex-unfold}{\vnbstmt{\begin{array}{@{}l@{}} \forall u.; \forall x, y \in u, t \in \operatorname{ccint}(0, 1). \\ \quad \operatorname{is\text{-}convex\text{-}cc}(u) \Rightarrow ((x + (t \cdot (y - x))) \in u) \end{array}}}
\vnbprop{conv3-is-convex}{\vnbstmt{\forall a, b, c \in \mathbb{C}.\; \operatorname{is\text{-}convex\text{-}cc}(\operatorname{conv3}(a, b, c))}}
\vnbprop{seg-trace-in-convex}{\vnbstmt{\begin{array}{@{}l@{}} \forall u.; \forall a, b \in \mathbb{C}. \\ \quad (\operatorname{is\text{-}convex\text{-}cc}(u) \wedge ((a \in u) \wedge (b \in u))) \Rightarrow  \\ \quad \quad \operatorname{trace}(\operatorname{seg\text{-}path}(a, b), 0, 1) \subseteq u \end{array}}}
\vnbprop{conv3-in-convex}{\vnbstmt{\begin{array}{@{}l@{}} \forall u.; \forall a, b, c \in \mathbb{C}. \\ \quad (\operatorname{is\text{-}convex\text{-}cc}(u) \wedge ((a \in u) \wedge ((b \in u) \wedge (c \in u)))) \Rightarrow (\operatorname{conv3}(a, b, c) \subseteq u) \end{array}}}
\vnbprop{conv3-vertices-in}{\vnbstmt{\begin{array}{@{}l@{}} \forall a, b, c \in \mathbb{C}. \\ \quad (a \in \operatorname{conv3}(a, b, c)) \wedge ((b \in \operatorname{conv3}(a, b, c)) \wedge (c \in \operatorname{conv3}(a, b, c))) \end{array}}}
\vnbprop{conv3-seg-trace}{\vnbstmt{\begin{array}{@{}l@{}} \forall a, b, c \in \mathbb{C}, x, y \in \operatorname{conv3}(a, b, c). \\ \quad \operatorname{trace}(\operatorname{seg\text{-}path}(x, y), 0, 1) \subseteq \operatorname{conv3}(a, b, c) \end{array}}}
\vnbprop{conv3-sub-hull}{\vnbstmt{\begin{array}{@{}l@{}} \forall a, b, c \in \mathbb{C}, x, y, z \in \operatorname{conv3}(a, b, c). \\ \quad \operatorname{conv3}(x, y, z) \subseteq \operatorname{conv3}(a, b, c) \end{array}}}
\vnbprop{cc-mid-affine}{\vnbstmt{\begin{array}{@{}l@{}} \forall x, y \in \mathbb{C}. \\ \quad (\operatorname{cc\text{-}mid}(x, y) = (x + (1/2 \cdot (y - x)))) \wedge (\operatorname{cc\text{-}mid}(x, y) = (y + (1/2 \cdot (x - y)))) \end{array}}}
\vnbprop{cc-mid-in-seg-trace}{\vnbstmt{\begin{array}{@{}l@{}} \forall x, y \in \mathbb{C}. \\ \quad ((\operatorname{cc\text{-}mid}(x, y) \in \operatorname{trace}(\operatorname{seg\text{-}path}(x, y), 0, 1)) \wedge (\operatorname{cc\text{-}mid}(x, y) \in \operatorname{trace}(\operatorname{seg\text{-}path}(y, x), 0, 1))) \end{array}}}
\vnbprop{conv3-mid-in}{\vnbstmt{\begin{array}{@{}l@{}} \forall a, b, c \in \mathbb{C}, x, y \in \operatorname{conv3}(a, b, c). \\ \quad \operatorname{cc\text{-}mid}(x, y) \in \operatorname{conv3}(a, b, c) \end{array}}}
\vnbprop{seg-trace-in-conv3}{\vnbstmt{\begin{array}{@{}l@{}} \forall a, b, c \in \mathbb{C}. \\ \quad \operatorname{trace}(\operatorname{seg\text{-}path}(a, b), 0, 1) \subseteq \operatorname{conv3}(a, b, c) \wedge \\ \quad \operatorname{trace}(\operatorname{seg\text{-}path}(b, c), 0, 1) \subseteq \operatorname{conv3}(a, b, c) \wedge \\ \quad \operatorname{trace}(\operatorname{seg\text{-}path}(c, a), 0, 1) \subseteq \operatorname{conv3}(a, b, c) \end{array}}}
\vnbprop{conv3-sub-triangle}{\vnbstmt{\begin{array}{@{}l@{}} \forall a, b, c \in \mathbb{C}. \\ \quad \operatorname{conv3}(a, \operatorname{cc\text{-}mid}(a, b), \operatorname{cc\text{-}mid}(a, c)) \subseteq \operatorname{conv3}(a, b, c) \wedge \\ \quad \operatorname{conv3}(\operatorname{cc\text{-}mid}(a, b), b, \operatorname{cc\text{-}mid}(b, c)) \subseteq \operatorname{conv3}(a, b, c) \wedge \\ \quad \operatorname{conv3}(\operatorname{cc\text{-}mid}(a, c), \operatorname{cc\text{-}mid}(b, c), c) \subseteq \operatorname{conv3}(a, b, c) \wedge \\ \quad \operatorname{conv3}(\operatorname{cc\text{-}mid}(a, b), \operatorname{cc\text{-}mid}(b, c), \operatorname{cc\text{-}mid}(a, c)) \subseteq \operatorname{conv3}(a, b, c) \end{array}}}
\vnbprop{seg-int-in-cc}{\vnbstmt{\begin{array}{@{}l@{}} \forall u, f.; \forall a, b \in \mathbb{C}. \\ \quad u \subseteq \mathbb{C} \wedge \\ \quad f \in (u \to \mathbb{C}) \wedge \\ \quad \forall y \in u. \\ \quad \quad \operatorname{is\text{-}continuous\text{-}at}(\operatorname{subspace\text{-}ms}(\operatorname{nf\text{-}metric\text{-}space}(\operatorname{cc\text{-}normed\text{-}field}), u), \\ \quad \quad \quad \operatorname{nf\text{-}metric\text{-}space}(\operatorname{cc\text{-}normed\text{-}field}), f, y) \wedge \\ \quad \operatorname{trace}(\operatorname{seg\text{-}path}(a, b), 0, 1) \subseteq u \Rightarrow \\ \quad \quad \operatorname{seg\text{-}int}(f, a, b) \in \mathbb{C} \end{array}}}
\vnbprop{seg-int-degenerate}{\vnbstmt{\begin{array}{@{}l@{}} \forall u, f.; \forall z \in u. \\ \quad u \subseteq \mathbb{C} \wedge \\ \quad f \in (u \to \mathbb{C}) \wedge \\ \quad \forall y \in u. \\ \quad \quad \operatorname{is\text{-}continuous\text{-}at}(\operatorname{subspace\text{-}ms}(\operatorname{nf\text{-}metric\text{-}space}(\operatorname{cc\text{-}normed\text{-}field}), u), \\ \quad \quad \quad \operatorname{nf\text{-}metric\text{-}space}(\operatorname{cc\text{-}normed\text{-}field}), f, y) \Rightarrow \\ \quad \quad \operatorname{seg\text{-}int}(f, z, z) = 0 \end{array}}}
\vnbprop{seg-int-reverse-at}{\vnbstmt{\begin{array}{@{}l@{}} \forall u, f.; \forall a, b \in \mathbb{C}. \\ \quad u \subseteq \mathbb{C} \wedge \\ \quad f \in (u \to \mathbb{C}) \wedge \\ \quad \forall y \in u. \\ \quad \quad \operatorname{is\text{-}continuous\text{-}at}(\operatorname{subspace\text{-}ms}(\operatorname{nf\text{-}metric\text{-}space}(\operatorname{cc\text{-}normed\text{-}field}), u), \\ \quad \quad \quad \operatorname{nf\text{-}metric\text{-}space}(\operatorname{cc\text{-}normed\text{-}field}), f, y) \wedge \\ \quad \operatorname{trace}(\operatorname{seg\text{-}path}(a, b), 0, 1) \subseteq u \Rightarrow \\ \quad \quad \operatorname{seg\text{-}int}(f, b, a) = (-1 \cdot \operatorname{seg\text{-}int}(f, a, b)) \end{array}}}
\vnbprop{seg-int-split-at}{\vnbstmt{\begin{array}{@{}l@{}} \forall u, f.; \forall a, b \in \mathbb{C}, c \in \operatorname{trace}(\operatorname{seg\text{-}path}(a, b), 0, 1). \\ \quad u \subseteq \mathbb{C} \wedge \\ \quad f \in (u \to \mathbb{C}) \wedge \\ \quad \forall y \in u. \\ \quad \quad \operatorname{is\text{-}continuous\text{-}at}(\operatorname{subspace\text{-}ms}(\operatorname{nf\text{-}metric\text{-}space}(\operatorname{cc\text{-}normed\text{-}field}), u), \\ \quad \quad \quad \operatorname{nf\text{-}metric\text{-}space}(\operatorname{cc\text{-}normed\text{-}field}), f, y) \wedge \\ \quad \operatorname{trace}(\operatorname{seg\text{-}path}(a, b), 0, 1) \subseteq u \Rightarrow \\ \quad \quad \operatorname{seg\text{-}int}(f, a, b) = (\operatorname{seg\text{-}int}(f, a, c) + \operatorname{seg\text{-}int}(f, c, b)) \end{array}}}
\vnbprop{tri-int-subdivision}{\vnbstmt{\begin{array}{@{}l@{}} \forall u, f.; \forall a, b, c \in \mathbb{C}. \\ \quad u \subseteq \mathbb{C} \wedge \\ \quad f \in (u \to \mathbb{C}) \wedge \\ \quad \forall y \in u. \\ \quad \quad \operatorname{is\text{-}continuous\text{-}at}(\operatorname{subspace\text{-}ms}(\operatorname{nf\text{-}metric\text{-}space}(\operatorname{cc\text{-}normed\text{-}field}), u), \\ \quad \quad \quad \operatorname{nf\text{-}metric\text{-}space}(\operatorname{cc\text{-}normed\text{-}field}), f, y) \wedge \\ \quad \operatorname{conv3}(a, b, c) \subseteq u \Rightarrow \\ \quad \quad \operatorname{tri\text{-}int}(f, a, b, c) \\ \quad \quad =\; ((\operatorname{tri\text{-}int}(f, a, \operatorname{cc\text{-}mid}(a, b), \operatorname{cc\text{-}mid}(a, c)) + \operatorname{tri\text{-}int}(f, \operatorname{cc\text{-}mid}(a, b), b, \operatorname{cc\text{-}mid}(b, c))) + \operatorname{tri\text{-}int}(f, \operatorname{cc\text{-}mid}(a, c), \operatorname{cc\text{-}mid}(b, c), c)) + \operatorname{tri\text{-}int}(f, \operatorname{cc\text{-}mid}(a, b), \operatorname{cc\text{-}mid}(b, c), \operatorname{cc\text{-}mid}(a, c)) \end{array}}}
\vnbprop{tri-int-abs-bound}{\vnbstmt{\begin{array}{@{}l@{}} \forall u, f.; \forall a, b, c \in \mathbb{C}, m \in \mathbb{R}. \\ \quad u \subseteq \mathbb{C} \wedge \\ \quad f \in (u \to \mathbb{C}) \wedge \\ \quad \forall y \in u. \\ \quad \quad \operatorname{is\text{-}continuous\text{-}at}(\operatorname{subspace\text{-}ms}(\operatorname{nf\text{-}metric\text{-}space}(\operatorname{cc\text{-}normed\text{-}field}), u), \\ \quad \quad \quad \operatorname{nf\text{-}metric\text{-}space}(\operatorname{cc\text{-}normed\text{-}field}), f, y) \wedge \\ \quad \operatorname{conv3}(a, b, c) \subseteq u \wedge \\ \quad \forall z \in \operatorname{conv3}(a, b, c).\; (\operatorname{magnitude}(f(z)) \leq m) \Rightarrow \\ \quad \quad \operatorname{magnitude}(\operatorname{tri\text{-}int}(f, a, b, c)) \leq (((\operatorname{magnitude}((b - a)) + \operatorname{magnitude}((c - b))) + \operatorname{magnitude}((a - c))) \cdot m) \end{array}}}
\vnbprop{tri-int-sum}{\vnbstmt{\begin{array}{@{}l@{}} \forall u, f, g.; \forall h \in (u \to \mathbb{C}), a, b, c \in \mathbb{C}. \\ \quad u \subseteq \mathbb{C} \wedge \\ \quad f \in (u \to \mathbb{C}) \wedge \\ \quad \forall y \in u. \\ \quad \quad \operatorname{is\text{-}continuous\text{-}at}(\operatorname{subspace\text{-}ms}(\operatorname{nf\text{-}metric\text{-}space}(\operatorname{cc\text{-}normed\text{-}field}), u), \\ \quad \quad \quad \operatorname{nf\text{-}metric\text{-}space}(\operatorname{cc\text{-}normed\text{-}field}), f, y) \wedge \\ \quad g \in (u \to \mathbb{C}) \wedge \\ \quad \forall y \in u. \\ \quad \quad \operatorname{is\text{-}continuous\text{-}at}(\operatorname{subspace\text{-}ms}(\operatorname{nf\text{-}metric\text{-}space}(\operatorname{cc\text{-}normed\text{-}field}), u), \\ \quad \quad \quad \operatorname{nf\text{-}metric\text{-}space}(\operatorname{cc\text{-}normed\text{-}field}), g, y) \wedge \\ \quad \forall z \in u.\; (h(z) \simeq (f(z) + g(z))) \wedge \\ \quad \operatorname{conv3}(a, b, c) \subseteq u \Rightarrow \\ \quad \quad \operatorname{tri\text{-}int}(h, a, b, c) = (\operatorname{tri\text{-}int}(f, a, b, c) + \operatorname{tri\text{-}int}(g, a, b, c)) \end{array}}}
\vnbprop{cc-open-subset-cc}{\vnbstmt{\forall u.\; (\operatorname{is\text{-}open}(\operatorname{nf\text{-}metric\text{-}space}(\operatorname{cc\text{-}normed\text{-}field}), u) \Rightarrow (u \subseteq \mathbb{C}))}}
\vnbprop{affine-primitive-on}{\vnbstmt{\begin{array}{@{}l@{}} \forall u.; \forall p, q, w \in \mathbb{C}, f \in (u \to \mathbb{C}). \\ \quad \operatorname{is\text{-}open}(\operatorname{nf\text{-}metric\text{-}space}(\operatorname{cc\text{-}normed\text{-}field}), u) \wedge \\ \quad \forall z \in u.\; (f(z) \simeq ((p \cdot z) + ((q \cdot ((z - w) \cdot (z - w))) \cdot \operatorname{recip}(2)))) \Rightarrow \\ \quad \quad (\operatorname{holomorphic\text{-}on}(u, f) \wedge \forall z \in u.\; \operatorname{is\text{-}diff\text{-}on}(\operatorname{cc\text{-}normed\text{-}field}, u, f, z, (p + (q \cdot (z - w))))) \end{array}}}
\vnbprop{tri-int-affine-zero}{\vnbstmt{\begin{array}{@{}l@{}} \forall u.; \forall p, q, w \in \mathbb{C}, g \in (u \to \mathbb{C}), a, b, c \in \mathbb{C}. \\ \quad \operatorname{is\text{-}open}(\operatorname{nf\text{-}metric\text{-}space}(\operatorname{cc\text{-}normed\text{-}field}), u) \wedge \\ \quad \forall z \in u.\; (g(z) \simeq (p + (q \cdot (z - w)))) \wedge \\ \quad \operatorname{conv3}(a, b, c) \subseteq u \Rightarrow \\ \quad \quad \operatorname{tri\text{-}int}(g, a, b, c) = 0 \end{array}}}
\vnbprop{holomorphic-on-continuous}{\vnbstmt{\begin{array}{@{}l@{}} \forall u, f.; \forall y \in u. \\ \quad \operatorname{holomorphic\text{-}on}(u, f) \Rightarrow  \\ \quad \quad \operatorname{is\text{-}continuous\text{-}at}(\operatorname{subspace\text{-}ms}(\operatorname{nf\text{-}metric\text{-}space}(\operatorname{cc\text{-}normed\text{-}field}), u), \\ \quad \quad \quad \operatorname{nf\text{-}metric\text{-}space}(\operatorname{cc\text{-}normed\text{-}field}), f, y) \end{array}}}
\vnbprop{cc-mid-coords}{\vnbstmt{\begin{array}{@{}l@{}} \forall x, y \in \mathbb{C}. \\ \quad \operatorname{real\text{-}part}(\operatorname{cc\text{-}mid}(x, y)) = (1/2 \cdot (\operatorname{real\text{-}part}(x) + \operatorname{real\text{-}part}(y))) \wedge \\ \quad \operatorname{imag\text{-}part}(\operatorname{cc\text{-}mid}(x, y)) = (1/2 \cdot (\operatorname{imag\text{-}part}(x) + \operatorname{imag\text{-}part}(y))) \end{array}}}
\vnbprop{std-mid}{\vnbstmt{\begin{array}{@{}l@{}} \forall x, y \in \mathbb{C}. \\ \quad 0 \leq \operatorname{real\text{-}part}(x) \wedge \\ \quad 0 \leq \operatorname{imag\text{-}part}(x) \wedge \\ \quad (\operatorname{real\text{-}part}(x) + \operatorname{imag\text{-}part}(x)) \leq 1 \wedge \\ \quad 0 \leq \operatorname{real\text{-}part}(y) \wedge \\ \quad 0 \leq \operatorname{imag\text{-}part}(y) \wedge \\ \quad (\operatorname{real\text{-}part}(y) + \operatorname{imag\text{-}part}(y)) \leq 1 \Rightarrow \\ \quad \quad 0 \leq \operatorname{real\text{-}part}(\operatorname{cc\text{-}mid}(x, y)) \wedge \\ \quad \quad 0 \leq \operatorname{imag\text{-}part}(\operatorname{cc\text{-}mid}(x, y)) \wedge \\ \quad \quad (\operatorname{real\text{-}part}(\operatorname{cc\text{-}mid}(x, y)) + \operatorname{imag\text{-}part}(\operatorname{cc\text{-}mid}(x, y))) \leq 1 \end{array}}}
\vnbprop{std-dist}{\vnbstmt{\begin{array}{@{}l@{}} \forall x, y \in \mathbb{C}. \\ \quad 0 \leq \operatorname{real\text{-}part}(x) \wedge \\ \quad 0 \leq \operatorname{imag\text{-}part}(x) \wedge \\ \quad (\operatorname{real\text{-}part}(x) + \operatorname{imag\text{-}part}(x)) \leq 1 \wedge \\ \quad 0 \leq \operatorname{real\text{-}part}(y) \wedge \\ \quad 0 \leq \operatorname{imag\text{-}part}(y) \wedge \\ \quad (\operatorname{real\text{-}part}(y) + \operatorname{imag\text{-}part}(y)) \leq 1 \Rightarrow \\ \quad \quad \operatorname{magnitude}((x - y)) \leq 2 \end{array}}}
\vnbprop{std-limit}{\vnbstmt{\begin{array}{@{}l@{}} \forall s \in (\mathbb{N} \to \mathbb{C}), w \in \mathbb{C}. \\ \quad \forall n \in \mathbb{N}. \\ \quad \quad ((0 \leq \operatorname{real\text{-}part}(s(n))) \wedge ((0 \leq \operatorname{imag\text{-}part}(s(n))) \wedge ((\operatorname{real\text{-}part}(s(n)) + \operatorname{imag\text{-}part}(s(n))) \leq 1))) \wedge \\ \quad \operatorname{converges\text{-}to}(\operatorname{cc\text{-}ms}, s, w) \Rightarrow \\ \quad \quad (0 \leq \operatorname{real\text{-}part}(w)) \wedge ((0 \leq \operatorname{imag\text{-}part}(w)) \wedge ((\operatorname{real\text{-}part}(w) + \operatorname{imag\text{-}part}(w)) \leq 1)) \end{array}}}
\vnbprop{bary-mid}{\vnbstmt{\begin{array}{@{}l@{}} \forall a, b, c, x, y \in \mathbb{C}. \\ \quad ((\operatorname{real\text{-}part}(\operatorname{cc\text{-}mid}(x, y)) \cdot a) + (\operatorname{imag\text{-}part}(\operatorname{cc\text{-}mid}(x, y)) \cdot b)) + (((1 - \operatorname{real\text{-}part}(\operatorname{cc\text{-}mid}(x, y))) - \operatorname{imag\text{-}part}(\operatorname{cc\text{-}mid}(x, y))) \cdot c) \\ \quad =\; \operatorname{cc\text{-}mid}(((\operatorname{real\text{-}part}(x) \cdot a) + (\operatorname{imag\text{-}part}(x) \cdot b)) + (((1 - \operatorname{real\text{-}part}(x)) - \operatorname{imag\text{-}part}(x)) \cdot c), \\ \quad \quad ((\operatorname{real\text{-}part}(y) \cdot a) + (\operatorname{imag\text{-}part}(y) \cdot b)) + (((1 - \operatorname{real\text{-}part}(y)) - \operatorname{imag\text{-}part}(y)) \cdot c)) \end{array}}}
\vnbprop{bary-in-conv3}{\vnbstmt{\begin{array}{@{}l@{}} \forall a, b, c, x \in \mathbb{C}. \\ \quad ((0 \leq \operatorname{real\text{-}part}(x)) \wedge ((0 \leq \operatorname{imag\text{-}part}(x)) \wedge ((\operatorname{real\text{-}part}(x) + \operatorname{imag\text{-}part}(x)) \leq 1))) \Rightarrow  \\ \quad \quad (((\operatorname{real\text{-}part}(x) \cdot a) + (\operatorname{imag\text{-}part}(x) \cdot b)) + (((1 - \operatorname{real\text{-}part}(x)) - \operatorname{imag\text{-}part}(x)) \cdot c)) \in \operatorname{conv3}(a, b, c) \end{array}}}
\vnbprop{bary-vertices}{\vnbstmt{\begin{array}{@{}l@{}} \forall a, b, c \in \mathbb{C}. \\ \quad (((\operatorname{real\text{-}part}(1) \cdot a) + (\operatorname{imag\text{-}part}(1) \cdot b)) + (((1 - \operatorname{real\text{-}part}(1)) - \operatorname{imag\text{-}part}(1)) \cdot c)) = a \wedge \\ \quad (((\operatorname{real\text{-}part}(+i) \cdot a) + (\operatorname{imag\text{-}part}(+i) \cdot b)) + (((1 - \operatorname{real\text{-}part}(+i)) - \operatorname{imag\text{-}part}(+i)) \cdot c)) = b \wedge \\ \quad (((\operatorname{real\text{-}part}(0) \cdot a) + (\operatorname{imag\text{-}part}(0) \cdot b)) + (((1 - \operatorname{real\text{-}part}(0)) - \operatorname{imag\text{-}part}(0)) \cdot c)) = c \end{array}}}
\vnbprop{bary-lip}{\vnbstmt{\begin{array}{@{}l@{}} \forall a, b, c, x, y \in \mathbb{C}. \\ \quad \operatorname{magnitude}(((((\operatorname{real\text{-}part}(x) \cdot a) + (\operatorname{imag\text{-}part}(x) \cdot b)) + (((1 - \operatorname{real\text{-}part}(x)) - \operatorname{imag\text{-}part}(x)) \cdot c)) - (((\operatorname{real\text{-}part}(y) \cdot a) + (\operatorname{imag\text{-}part}(y) \cdot b)) + (((1 - \operatorname{real\text{-}part}(y)) - \operatorname{imag\text{-}part}(y)) \cdot c)))) \\ \quad \leq\; (\operatorname{magnitude}((a - c)) + \operatorname{magnitude}((b - c))) \cdot \operatorname{magnitude}((x - y)) \end{array}}}
\vnbprop{half-magnitude}{\vnbstmt{\forall x \in \mathbb{C}.\; (\operatorname{magnitude}((1/2 \cdot x)) = (1/2 \cdot \operatorname{magnitude}(x)))}}
\vnbprop{tri-child-geom-1}{\vnbstmt{\begin{array}{@{}l@{}} \forall x, y, z \in \mathbb{C}. \\ \quad (\operatorname{magnitude}((\operatorname{cc\text{-}mid}(x, y) - x)) + \operatorname{magnitude}((\operatorname{cc\text{-}mid}(x, z) - \operatorname{cc\text{-}mid}(x, y)))) + \operatorname{magnitude}((x - \operatorname{cc\text{-}mid}(x, z))) \\ \quad =\; 1/2 \cdot ((\operatorname{magnitude}((y - x)) + \operatorname{magnitude}((z - y))) + \operatorname{magnitude}((x - z))) \wedge \\ \quad \operatorname{magnitude}((x - x)) \leq ((\operatorname{magnitude}((y - x)) + \operatorname{magnitude}((z - y))) + \operatorname{magnitude}((x - z))) \end{array}}}
\vnbprop{tri-child-geom-2}{\vnbstmt{\begin{array}{@{}l@{}} \forall x, y, z \in \mathbb{C}. \\ \quad (\operatorname{magnitude}((y - \operatorname{cc\text{-}mid}(x, y))) + \operatorname{magnitude}((\operatorname{cc\text{-}mid}(y, z) - y))) + \operatorname{magnitude}((\operatorname{cc\text{-}mid}(x, y) - \operatorname{cc\text{-}mid}(y, z))) \\ \quad =\; 1/2 \cdot ((\operatorname{magnitude}((y - x)) + \operatorname{magnitude}((z - y))) + \operatorname{magnitude}((x - z))) \wedge \\ \quad \operatorname{magnitude}((\operatorname{cc\text{-}mid}(x, y) - x)) \leq ((\operatorname{magnitude}((y - x)) + \operatorname{magnitude}((z - y))) + \operatorname{magnitude}((x - z))) \end{array}}}
\vnbprop{tri-child-geom-3}{\vnbstmt{\begin{array}{@{}l@{}} \forall x, y, z \in \mathbb{C}. \\ \quad (\operatorname{magnitude}((\operatorname{cc\text{-}mid}(y, z) - \operatorname{cc\text{-}mid}(x, z))) + \operatorname{magnitude}((z - \operatorname{cc\text{-}mid}(y, z)))) + \operatorname{magnitude}((\operatorname{cc\text{-}mid}(x, z) - z)) \\ \quad =\; 1/2 \cdot ((\operatorname{magnitude}((y - x)) + \operatorname{magnitude}((z - y))) + \operatorname{magnitude}((x - z))) \wedge \\ \quad \operatorname{magnitude}((\operatorname{cc\text{-}mid}(x, z) - x)) \leq ((\operatorname{magnitude}((y - x)) + \operatorname{magnitude}((z - y))) + \operatorname{magnitude}((x - z))) \end{array}}}
\vnbprop{tri-child-geom-4}{\vnbstmt{\begin{array}{@{}l@{}} \forall x, y, z \in \mathbb{C}. \\ \quad (\operatorname{magnitude}((\operatorname{cc\text{-}mid}(y, z) - \operatorname{cc\text{-}mid}(x, y))) + \operatorname{magnitude}((\operatorname{cc\text{-}mid}(x, z) - \operatorname{cc\text{-}mid}(y, z)))) + \operatorname{magnitude}((\operatorname{cc\text{-}mid}(x, y) - \operatorname{cc\text{-}mid}(x, z))) \\ \quad =\; 1/2 \cdot ((\operatorname{magnitude}((y - x)) + \operatorname{magnitude}((z - y))) + \operatorname{magnitude}((x - z))) \wedge \\ \quad \operatorname{magnitude}((\operatorname{cc\text{-}mid}(x, y) - x)) \leq ((\operatorname{magnitude}((y - x)) + \operatorname{magnitude}((z - y))) + \operatorname{magnitude}((x - z))) \end{array}}}
\vnbprop{tri-int-in-cc}{\vnbstmt{\begin{array}{@{}l@{}} \forall u, f.; \forall a, b, c \in \mathbb{C}. \\ \quad u \subseteq \mathbb{C} \wedge \\ \quad f \in (u \to \mathbb{C}) \wedge \\ \quad \forall y \in u. \\ \quad \quad \operatorname{is\text{-}continuous\text{-}at}(\operatorname{subspace\text{-}ms}(\operatorname{nf\text{-}metric\text{-}space}(\operatorname{cc\text{-}normed\text{-}field}), u), \\ \quad \quad \quad \operatorname{nf\text{-}metric\text{-}space}(\operatorname{cc\text{-}normed\text{-}field}), f, y) \wedge \\ \quad \operatorname{conv3}(a, b, c) \subseteq u \Rightarrow \\ \quad \quad \operatorname{tri\text{-}int}(f, a, b, c) \in \mathbb{C} \end{array}}}
\vnbprop{goursat-nest}{\vnbstmt{\begin{array}{@{}l@{}} \forall u, f.; \forall a, b, c \in \mathbb{C}. \\ \quad u \subseteq \mathbb{C} \wedge \\ \quad f \in (u \to \mathbb{C}) \wedge \\ \quad \forall y \in u. \\ \quad \quad \operatorname{is\text{-}continuous\text{-}at}(\operatorname{subspace\text{-}ms}(\operatorname{nf\text{-}metric\text{-}space}(\operatorname{cc\text{-}normed\text{-}field}), u), \\ \quad \quad \quad \operatorname{nf\text{-}metric\text{-}space}(\operatorname{cc\text{-}normed\text{-}field}), f, y) \wedge \\ \quad \operatorname{conv3}(a, b, c) \subseteq u \Rightarrow \\ \quad \quad \exists w \in \mathbb{C}. \\ \quad \quad \quad 0 \leq \operatorname{real\text{-}part}(w) \wedge \\ \quad \quad \quad 0 \leq \operatorname{imag\text{-}part}(w) \wedge \\ \quad \quad \quad (\operatorname{real\text{-}part}(w) + \operatorname{imag\text{-}part}(w)) \leq 1 \wedge \\ \quad \quad \quad \forall n \in \mathbb{N}. \exists x, \mathit{gouy}, z \in \mathbb{C}. \\ \quad \quad \quad \quad 0 \leq \operatorname{real\text{-}part}(x) \wedge \\ \quad \quad \quad \quad 0 \leq \operatorname{imag\text{-}part}(x) \wedge \\ \quad \quad \quad \quad (\operatorname{real\text{-}part}(x) + \operatorname{imag\text{-}part}(x)) \leq 1 \wedge \\ \quad \quad \quad \quad 0 \leq \operatorname{real\text{-}part}(\mathit{gouy}) \wedge \\ \quad \quad \quad \quad 0 \leq \operatorname{imag\text{-}part}(\mathit{gouy}) \wedge \\ \quad \quad \quad \quad (\operatorname{real\text{-}part}(\mathit{gouy}) + \operatorname{imag\text{-}part}(\mathit{gouy})) \leq 1 \wedge \\ \quad \quad \quad \quad 0 \leq \operatorname{real\text{-}part}(z) \wedge \\ \quad \quad \quad \quad 0 \leq \operatorname{imag\text{-}part}(z) \wedge \\ \quad \quad \quad \quad (\operatorname{real\text{-}part}(z) + \operatorname{imag\text{-}part}(z)) \leq 1 \wedge \\ \quad \quad \quad \quad ((\operatorname{magnitude}((\mathit{gouy} - x)) + \operatorname{magnitude}((z - \mathit{gouy}))) + \operatorname{magnitude}((x - z))) \leq (6 \cdot \operatorname{recip}(\operatorname{power}(2, n))) \wedge \\ \quad \quad \quad \quad \operatorname{magnitude}((x - w)) \leq (12 \cdot \operatorname{recip}(\operatorname{power}(2, n))) \wedge \\ \quad \quad \quad \quad \operatorname{magnitude}(\operatorname{tri\text{-}int}(f, a, b, c)) \\ \quad \quad \quad \quad \leq\; \operatorname{power}(4, n) \cdot \operatorname{magnitude}(\operatorname{tri\text{-}int}(f, (((\operatorname{real\text{-}part}(x) \cdot a) + (\operatorname{imag\text{-}part}(x) \cdot b)) + (((1 - \operatorname{real\text{-}part}(x)) - \operatorname{imag\text{-}part}(x)) \cdot c)), (((\operatorname{real\text{-}part}(\mathit{gouy}) \cdot a) + (\operatorname{imag\text{-}part}(\mathit{gouy}) \cdot b)) + (((1 - \operatorname{real\text{-}part}(\mathit{gouy})) - \operatorname{imag\text{-}part}(\mathit{gouy})) \cdot c)), (((\operatorname{real\text{-}part}(z) \cdot a) + (\operatorname{imag\text{-}part}(z) \cdot b)) + (((1 - \operatorname{real\text{-}part}(z)) - \operatorname{imag\text{-}part}(z)) \cdot c)))) \end{array}}}
\vnbprop{affine-holomorphic}{\vnbstmt{\begin{array}{@{}l@{}} \forall u.; \forall p, q, w \in \mathbb{C}, g \in (u \to \mathbb{C}). \\ \quad (\operatorname{is\text{-}open}(\operatorname{nf\text{-}metric\text{-}space}(\operatorname{cc\text{-}normed\text{-}field}), u) \wedge \forall z \in u.\; (g(z) \simeq (p + (q \cdot (z - w))))) \Rightarrow  \\ \quad \quad \operatorname{holomorphic\text{-}on}(u, g) \end{array}}}
\vnbprop{small-dyadic}{\vnbstmt{\begin{array}{@{}l@{}} \forall m \in \mathbb{R}, r. \\ \quad ((0 \leq m) \wedge \operatorname{pos\text{-}rr}(r)) \Rightarrow \exists n \in \mathbb{N}.\; ((m \cdot \operatorname{recip}(\operatorname{power}(2, n))) < r) \end{array}}}
\vnbprop{bary-side-near}{\vnbstmt{\begin{array}{@{}l@{}} \forall a, b, c, p, q, x, w \in \mathbb{C}, r, m \in \mathbb{R}, z \in \operatorname{trace}(\operatorname{seg\text{-}path}((((\operatorname{real\text{-}part}(p) \cdot a) + (\operatorname{imag\text{-}part}(p) \cdot b)) + (((1 - \operatorname{real\text{-}part}(p)) - \operatorname{imag\text{-}part}(p)) \cdot c)), (((\operatorname{real\text{-}part}(q) \cdot a) + (\operatorname{imag\text{-}part}(q) \cdot b)) + (((1 - \operatorname{real\text{-}part}(q)) - \operatorname{imag\text{-}part}(q)) \cdot c))), 0, 1). \\ \quad ((\operatorname{magnitude}((p - x)) \leq r) \wedge ((\operatorname{magnitude}((q - p)) \leq r) \wedge ((r \leq (6 \cdot m)) \wedge (\operatorname{magnitude}((x - w)) \leq (12 \cdot m))))) \Rightarrow  \\ \quad \quad \operatorname{magnitude}((z - (((\operatorname{real\text{-}part}(w) \cdot a) + (\operatorname{imag\text{-}part}(w) \cdot b)) + (((1 - \operatorname{real\text{-}part}(w)) - \operatorname{imag\text{-}part}(w)) \cdot c)))) \\ \quad \quad \leq\; (24 \cdot (\operatorname{magnitude}((a - c)) + \operatorname{magnitude}((b - c)))) \cdot m \end{array}}}
\vnbprop{tri-int-abs-bound-sides}{\vnbstmt{\begin{array}{@{}l@{}} \forall u, f.; \forall a, b, c \in \mathbb{C}, m \in \mathbb{R}. \\ \quad u \subseteq \mathbb{C} \wedge \\ \quad f \in (u \to \mathbb{C}) \wedge \\ \quad \forall y \in u. \\ \quad \quad \operatorname{is\text{-}continuous\text{-}at}(\operatorname{subspace\text{-}ms}(\operatorname{nf\text{-}metric\text{-}space}(\operatorname{cc\text{-}normed\text{-}field}), u), \\ \quad \quad \quad \operatorname{nf\text{-}metric\text{-}space}(\operatorname{cc\text{-}normed\text{-}field}), f, y) \wedge \\ \quad \operatorname{trace}(\operatorname{seg\text{-}path}(a, b), 0, 1) \subseteq u \wedge \\ \quad \operatorname{trace}(\operatorname{seg\text{-}path}(b, c), 0, 1) \subseteq u \wedge \\ \quad \operatorname{trace}(\operatorname{seg\text{-}path}(c, a), 0, 1) \subseteq u \wedge \\ \quad \forall z \in \operatorname{trace}(\operatorname{seg\text{-}path}(a, b), 0, 1).\; (\operatorname{magnitude}(f(z)) \leq m) \wedge \\ \quad \forall z \in \operatorname{trace}(\operatorname{seg\text{-}path}(b, c), 0, 1).\; (\operatorname{magnitude}(f(z)) \leq m) \wedge \\ \quad \forall z \in \operatorname{trace}(\operatorname{seg\text{-}path}(c, a), 0, 1).\; (\operatorname{magnitude}(f(z)) \leq m) \Rightarrow \\ \quad \quad \operatorname{magnitude}(\operatorname{tri\text{-}int}(f, a, b, c)) \leq (((\operatorname{magnitude}((b - a)) + \operatorname{magnitude}((c - b))) + \operatorname{magnitude}((a - c))) \cdot m) \end{array}}}
\vnbprop{goursat-triangle}{\vnbstmt{\begin{array}{@{}l@{}} \forall u, f.; \forall a, b, c \in u. \\ \quad (\operatorname{holomorphic\text{-}on}(u, f) \wedge (\operatorname{conv3}(a, b, c) \subseteq u)) \Rightarrow (\operatorname{tri\text{-}int}(f, a, b, c) = 0) \end{array}}}
\vnbprop{cc-diff-on-local-quadratic}{\vnbstmt{\begin{array}{@{}l@{}} \forall s.; \forall f \in (s \to \mathbb{C}), a \in s, l \in \mathbb{C}, m \in \mathbb{R}, r. \\ \quad \operatorname{is\text{-}open}(\operatorname{nf\text{-}metric\text{-}space}(\operatorname{cc\text{-}normed\text{-}field}), s) \wedge \\ \quad 0 \leq m \wedge \\ \quad \operatorname{pos\text{-}rr}(r) \wedge \\ \quad \forall x \in s. \\ \quad \quad (\operatorname{magnitude}((x - a)) < r) \Rightarrow  \\ \quad \quad \quad \operatorname{magnitude}(((f(x) - f(a)) - ((x - a) \cdot l))) \leq (m \cdot (\operatorname{magnitude}((x - a)) \cdot \operatorname{magnitude}((x - a)))) \Rightarrow \\ \quad \quad \operatorname{is\text{-}diff\text{-}on}(\operatorname{cc\text{-}normed\text{-}field}, s, f, a, l) \end{array}}}
\vnbprop{cauchy-convex-segments}{\vnbstmt{\begin{array}{@{}l@{}} \forall u, f.; \forall a, b, c \in u. \\ \quad (\operatorname{is\text{-}convex\text{-}cc}(u) \wedge \operatorname{holomorphic\text{-}on}(u, f)) \Rightarrow  \\ \quad \quad \operatorname{seg\text{-}int}(f, a, c) = (\operatorname{seg\text{-}int}(f, a, b) + \operatorname{seg\text{-}int}(f, b, c)) \end{array}}}
\vnbprop{seg-int-sum}{\vnbstmt{\begin{array}{@{}l@{}} \forall u, f, g.; \forall h \in (u \to \mathbb{C}), a, b \in \mathbb{C}. \\ \quad u \subseteq \mathbb{C} \wedge \\ \quad f \in (u \to \mathbb{C}) \wedge \\ \quad \forall y \in u. \\ \quad \quad \operatorname{is\text{-}continuous\text{-}at}(\operatorname{subspace\text{-}ms}(\operatorname{nf\text{-}metric\text{-}space}(\operatorname{cc\text{-}normed\text{-}field}), u), \\ \quad \quad \quad \operatorname{nf\text{-}metric\text{-}space}(\operatorname{cc\text{-}normed\text{-}field}), f, y) \wedge \\ \quad g \in (u \to \mathbb{C}) \wedge \\ \quad \forall y \in u. \\ \quad \quad \operatorname{is\text{-}continuous\text{-}at}(\operatorname{subspace\text{-}ms}(\operatorname{nf\text{-}metric\text{-}space}(\operatorname{cc\text{-}normed\text{-}field}), u), \\ \quad \quad \quad \operatorname{nf\text{-}metric\text{-}space}(\operatorname{cc\text{-}normed\text{-}field}), g, y) \wedge \\ \quad \forall z \in u.\; (h(z) \simeq (f(z) + g(z))) \wedge \\ \quad \operatorname{trace}(\operatorname{seg\text{-}path}(a, b), 0, 1) \subseteq u \Rightarrow \\ \quad \quad \operatorname{seg\text{-}int}(h, a, b) = (\operatorname{seg\text{-}int}(f, a, b) + \operatorname{seg\text{-}int}(g, a, b)) \end{array}}}
\vnbprop{seg-int-abs-bound}{\vnbstmt{\begin{array}{@{}l@{}} \forall u, f.; \forall a, b \in \mathbb{C}, m \in \mathbb{R}. \\ \quad u \subseteq \mathbb{C} \wedge \\ \quad f \in (u \to \mathbb{C}) \wedge \\ \quad \forall y \in u. \\ \quad \quad \operatorname{is\text{-}continuous\text{-}at}(\operatorname{subspace\text{-}ms}(\operatorname{nf\text{-}metric\text{-}space}(\operatorname{cc\text{-}normed\text{-}field}), u), \\ \quad \quad \quad \operatorname{nf\text{-}metric\text{-}space}(\operatorname{cc\text{-}normed\text{-}field}), f, y) \wedge \\ \quad \operatorname{trace}(\operatorname{seg\text{-}path}(a, b), 0, 1) \subseteq u \wedge \\ \quad \forall z \in \operatorname{trace}(\operatorname{seg\text{-}path}(a, b), 0, 1).\; (\operatorname{magnitude}(f(z)) \leq m) \Rightarrow \\ \quad \quad \operatorname{magnitude}(\operatorname{seg\text{-}int}(f, a, b)) \leq (\operatorname{magnitude}((b - a)) \cdot m) \end{array}}}
\vnbprop{seg-int-const}{\vnbstmt{\begin{array}{@{}l@{}} \forall u.; \forall p, w \in \mathbb{C}, g \in (u \to \mathbb{C}), a, b \in \mathbb{C}. \\ \quad \operatorname{is\text{-}open}(\operatorname{nf\text{-}metric\text{-}space}(\operatorname{cc\text{-}normed\text{-}field}), u) \wedge \\ \quad \forall z \in u.\; (g(z) \simeq (p + (0 \cdot (z - w)))) \wedge \\ \quad \operatorname{trace}(\operatorname{seg\text{-}path}(a, b), 0, 1) \subseteq u \Rightarrow \\ \quad \quad \operatorname{seg\text{-}int}(g, a, b) = (p \cdot (b - a)) \end{array}}}
\vnbprop{cauchy-convex-primitive}{\vnbstmt{\begin{array}{@{}l@{}} \forall u, f.; \forall o \in u. \\ \quad (\operatorname{is\text{-}convex\text{-}cc}(u) \wedge \operatorname{holomorphic\text{-}on}(u, f)) \Rightarrow  \\ \quad \quad \exists g \in (u \to \mathbb{C}). \\ \quad \quad \quad (\operatorname{holomorphic\text{-}on}(u, g) \wedge \forall z \in u.\; \operatorname{is\text{-}diff\text{-}on}(\operatorname{cc\text{-}normed\text{-}field}, u, g, z, f(z))) \end{array}}}
\vnbprop{cauchy-convex-closed}{\vnbstmt{\begin{array}{@{}l@{}} \forall u, f, m, c, a, b. \\ \quad \operatorname{is\text{-}convex\text{-}cc}(u) \wedge \\ \quad \operatorname{holomorphic\text{-}on}(u, f) \wedge \\ \quad \operatorname{is\text{-}road}(m, c, a, b) \wedge \\ \quad \operatorname{trace}(m, a, b) \subseteq u \wedge \\ \quad m(b) = m(a) \Rightarrow \\ \quad \quad \operatorname{line\text{-}int}(f, m, c, a, b) = 0 \end{array}}}

\vnbbookfile{theorem-library/goursat-exceptional}{theorem-library/goursat-exceptional}
\noindent{\small goursat-exceptional.scm -- CAUCHY-GOURSAT WITH ONE EXCEPTIONAL POINT (complex-analysis.pdf ch. 3, s. 3.2: Prop 3.6, Cor 3.7, Prop 3.8, Cor 3.9 in the notes' form: f CONTINUOUS on the open set U and HOLOMORPHIC on U \textbackslash{} \{x\} for some x in U).  Batch 35, 2026-09-26.}\par
\vnbprop{seg-int-congruence}{\vnbstmt{\begin{array}{@{}l@{}} \forall u, f, w.; \forall g \in (w \to \mathbb{C}), a, b \in \mathbb{C}. \\ \quad u \subseteq \mathbb{C} \wedge \\ \quad f \in (u \to \mathbb{C}) \wedge \\ \quad \forall y \in u. \\ \quad \quad \operatorname{is\text{-}continuous\text{-}at}(\operatorname{subspace\text{-}ms}(\operatorname{nf\text{-}metric\text{-}space}(\operatorname{cc\text{-}normed\text{-}field}), u), \\ \quad \quad \quad \operatorname{nf\text{-}metric\text{-}space}(\operatorname{cc\text{-}normed\text{-}field}), f, y) \wedge \\ \quad \operatorname{trace}(\operatorname{seg\text{-}path}(a, b), 0, 1) \subseteq u \wedge \\ \quad \operatorname{trace}(\operatorname{seg\text{-}path}(a, b), 0, 1) \subseteq w \wedge \\ \quad \forall z \in \operatorname{trace}(\operatorname{seg\text{-}path}(a, b), 0, 1).\; (g(z) \simeq f(z)) \Rightarrow \\ \quad \quad \operatorname{seg\text{-}int}(g, a, b) = \operatorname{seg\text{-}int}(f, a, b) \end{array}}}
\vnbprop{tri-int-restrict}{\vnbstmt{\begin{array}{@{}l@{}} \forall u, f, w.; \forall a, b, c \in \mathbb{C}. \\ \quad u \subseteq \mathbb{C} \wedge \\ \quad f \in (u \to \mathbb{C}) \wedge \\ \quad \forall y \in u. \\ \quad \quad \operatorname{is\text{-}continuous\text{-}at}(\operatorname{subspace\text{-}ms}(\operatorname{nf\text{-}metric\text{-}space}(\operatorname{cc\text{-}normed\text{-}field}), u), \\ \quad \quad \quad \operatorname{nf\text{-}metric\text{-}space}(\operatorname{cc\text{-}normed\text{-}field}), f, y) \wedge \\ \quad w \subseteq u \wedge \\ \quad \operatorname{conv3}(a, b, c) \subseteq w \Rightarrow \\ \quad \quad \operatorname{tri\text{-}int}(\operatorname{restrict}(f, w), a, b, c) = \operatorname{tri\text{-}int}(f, a, b, c) \end{array}}}
\vnbprop{tri-int-fan}{\vnbstmt{\begin{array}{@{}l@{}} \forall u, f.; \forall a, b, c \in \mathbb{C}, x \in \operatorname{conv3}(a, b, c). \\ \quad u \subseteq \mathbb{C} \wedge \\ \quad f \in (u \to \mathbb{C}) \wedge \\ \quad \forall y \in u. \\ \quad \quad \operatorname{is\text{-}continuous\text{-}at}(\operatorname{subspace\text{-}ms}(\operatorname{nf\text{-}metric\text{-}space}(\operatorname{cc\text{-}normed\text{-}field}), u), \\ \quad \quad \quad \operatorname{nf\text{-}metric\text{-}space}(\operatorname{cc\text{-}normed\text{-}field}), f, y) \wedge \\ \quad \operatorname{conv3}(a, b, c) \subseteq u \Rightarrow \\ \quad \quad \operatorname{tri\text{-}int}(f, a, b, c) = ((\operatorname{tri\text{-}int}(f, a, b, x) + \operatorname{tri\text{-}int}(f, b, c, x)) + \operatorname{tri\text{-}int}(f, c, a, x)) \end{array}}}
\vnbprop{tri-int-collinear}{\vnbstmt{\begin{array}{@{}l@{}} \forall u, f.; \forall a, b, c \in \mathbb{C}. \\ \quad u \subseteq \mathbb{C} \wedge \\ \quad f \in (u \to \mathbb{C}) \wedge \\ \quad \forall y \in u. \\ \quad \quad \operatorname{is\text{-}continuous\text{-}at}(\operatorname{subspace\text{-}ms}(\operatorname{nf\text{-}metric\text{-}space}(\operatorname{cc\text{-}normed\text{-}field}), u), \\ \quad \quad \quad \operatorname{nf\text{-}metric\text{-}space}(\operatorname{cc\text{-}normed\text{-}field}), f, y) \wedge \\ \quad \operatorname{conv3}(a, b, c) \subseteq u \wedge \\ \quad c \in \operatorname{trace}(\operatorname{seg\text{-}path}(a, b), 0, 1) \vee \\ \quad a \in \operatorname{trace}(\operatorname{seg\text{-}path}(b, c), 0, 1) \vee \\ \quad b \in \operatorname{trace}(\operatorname{seg\text{-}path}(c, a), 0, 1) \Rightarrow \\ \quad \quad \operatorname{tri\text{-}int}(f, a, b, c) = 0 \end{array}}}
\vnbprop{cc-cross-coords}{\vnbstmt{\begin{array}{@{}l@{}} \forall p, q \in \mathbb{C}. \\ \quad \operatorname{imag\text{-}part}((p \cdot \operatorname{conjugate}(q))) = ((\operatorname{imag\text{-}part}(p) \cdot \operatorname{real\text{-}part}(q)) - (\operatorname{real\text{-}part}(p) \cdot \operatorname{imag\text{-}part}(q))) \end{array}}}
\vnbprop{cc-affinely-independent}{\vnbstmt{\begin{array}{@{}l@{}} \forall a, b, c \in \mathbb{C}, p, q \in \mathbb{R}. \\ \quad (\neg (\operatorname{imag\text{-}part}(((b - a) \cdot \operatorname{conjugate}((c - a)))) = 0) \wedge (((p \cdot (b - a)) + (q \cdot (c - a))) = 0)) \Rightarrow  \\ \quad \quad (p = 0) \wedge (q = 0) \end{array}}}
\vnbprop{conv3-avoids-vertex}{\vnbstmt{\begin{array}{@{}l@{}} \forall a, b, c \in \mathbb{C}, t \in \mathbb{R}. \\ \quad (\neg (\operatorname{imag\text{-}part}(((b - a) \cdot \operatorname{conjugate}((c - a)))) = 0) \wedge ((0 < t) \wedge (t \leq 1))) \Rightarrow  \\ \quad \quad (\neg (a \in \operatorname{conv3}((a + (t \cdot (b - a))), b, c)) \wedge \neg (a \in \operatorname{conv3}((a + (t \cdot (b - a))), c, (a + (t \cdot (c - a)))))) \end{array}}}
\vnbprop{conv3-near-vertex}{\vnbstmt{\begin{array}{@{}l@{}} \forall a, b, c \in \mathbb{C}, z \in \operatorname{conv3}(a, b, c). \\ \quad \operatorname{magnitude}((z - a)) \leq (\operatorname{magnitude}((b - a)) + \operatorname{magnitude}((c - a))) \end{array}}}
\vnbprop{collinear-cases}{\vnbstmt{\begin{array}{@{}l@{}} \forall a, b, c \in \mathbb{C}. \\ \quad (\operatorname{imag\text{-}part}(((b - a) \cdot \operatorname{conjugate}((c - a)))) = 0) \Rightarrow  \\ \quad \quad c \in \operatorname{trace}(\operatorname{seg\text{-}path}(a, b), 0, 1) \vee \\ \quad \quad a \in \operatorname{trace}(\operatorname{seg\text{-}path}(b, c), 0, 1) \vee \\ \quad \quad b \in \operatorname{trace}(\operatorname{seg\text{-}path}(c, a), 0, 1) \end{array}}}
\vnbprop{conv3-vertex-shrink}{\vnbstmt{\begin{array}{@{}l@{}} \forall a, b, c \in \mathbb{C}, t \in \operatorname{ccint}(0, 1). \\ \quad \operatorname{conv3}(a, (a + (t \cdot (b - a))), (a + (t \cdot (c - a)))) \subseteq \operatorname{conv3}(a, b, c) \wedge \\ \quad \operatorname{conv3}((a + (t \cdot (b - a))), b, c) \subseteq \operatorname{conv3}(a, b, c) \wedge \\ \quad \operatorname{conv3}((a + (t \cdot (b - a))), c, (a + (t \cdot (c - a)))) \subseteq \operatorname{conv3}(a, b, c) \end{array}}}
\vnbprop{tri-int-vertex-split}{\vnbstmt{\begin{array}{@{}l@{}} \forall u, f.; \forall a, b, c \in \mathbb{C}, t \in \operatorname{ccint}(0, 1). \\ \quad u \subseteq \mathbb{C} \wedge \\ \quad f \in (u \to \mathbb{C}) \wedge \\ \quad \forall y \in u. \\ \quad \quad \operatorname{is\text{-}continuous\text{-}at}(\operatorname{subspace\text{-}ms}(\operatorname{nf\text{-}metric\text{-}space}(\operatorname{cc\text{-}normed\text{-}field}), u), \\ \quad \quad \quad \operatorname{nf\text{-}metric\text{-}space}(\operatorname{cc\text{-}normed\text{-}field}), f, y) \wedge \\ \quad \operatorname{conv3}(a, b, c) \subseteq u \Rightarrow \\ \quad \quad \operatorname{tri\text{-}int}(f, a, b, c) \\ \quad \quad =\; (\operatorname{tri\text{-}int}(f, a, (a + (t \cdot (b - a))), (a + (t \cdot (c - a)))) + \operatorname{tri\text{-}int}(f, (a + (t \cdot (b - a))), b, c)) + \operatorname{tri\text{-}int}(f, (a + (t \cdot (b - a))), c, (a + (t \cdot (c - a)))) \end{array}}}
\vnbprop{goursat-vertex}{\vnbstmt{\begin{array}{@{}l@{}} \forall u, f, x.; \forall b, c \in u. \\ \quad \operatorname{is\text{-}open}(\operatorname{nf\text{-}metric\text{-}space}(\operatorname{cc\text{-}normed\text{-}field}), u) \wedge \\ \quad f \in (u \to \mathbb{C}) \wedge \\ \quad \forall y \in u. \\ \quad \quad \operatorname{is\text{-}continuous\text{-}at}(\operatorname{subspace\text{-}ms}(\operatorname{nf\text{-}metric\text{-}space}(\operatorname{cc\text{-}normed\text{-}field}), u), \\ \quad \quad \quad \operatorname{nf\text{-}metric\text{-}space}(\operatorname{cc\text{-}normed\text{-}field}), f, y) \wedge \\ \quad x \in u \wedge \\ \quad \operatorname{holomorphic\text{-}on}(\operatorname{difference}(u, \operatorname{singleton}(x)), \operatorname{restrict}(f, \operatorname{difference}(u, \operatorname{singleton}(x)))) \wedge \\ \quad \operatorname{conv3}(x, b, c) \subseteq u \Rightarrow \\ \quad \quad \operatorname{tri\text{-}int}(f, x, b, c) = 0 \end{array}}}
\vnbprop{goursat-triangle-except}{\vnbstmt{\begin{array}{@{}l@{}} \forall u, f, x.; \forall a, b, c \in u. \\ \quad \operatorname{is\text{-}open}(\operatorname{nf\text{-}metric\text{-}space}(\operatorname{cc\text{-}normed\text{-}field}), u) \wedge \\ \quad f \in (u \to \mathbb{C}) \wedge \\ \quad \forall y \in u. \\ \quad \quad \operatorname{is\text{-}continuous\text{-}at}(\operatorname{subspace\text{-}ms}(\operatorname{nf\text{-}metric\text{-}space}(\operatorname{cc\text{-}normed\text{-}field}), u), \\ \quad \quad \quad \operatorname{nf\text{-}metric\text{-}space}(\operatorname{cc\text{-}normed\text{-}field}), f, y) \wedge \\ \quad x \in u \wedge \\ \quad \operatorname{holomorphic\text{-}on}(\operatorname{difference}(u, \operatorname{singleton}(x)), \operatorname{restrict}(f, \operatorname{difference}(u, \operatorname{singleton}(x)))) \wedge \\ \quad \operatorname{conv3}(a, b, c) \subseteq u \Rightarrow \\ \quad \quad \operatorname{tri\text{-}int}(f, a, b, c) = 0 \end{array}}}
\vnbprop{cc-diff-on-of-eps-bound}{\vnbstmt{\begin{array}{@{}l@{}} \forall s.; \forall g \in (s \to \mathbb{C}), a \in s, l \in \mathbb{C}. \\ \quad \operatorname{is\text{-}open}(\operatorname{nf\text{-}metric\text{-}space}(\operatorname{cc\text{-}normed\text{-}field}), s) \wedge \\ \quad \forall e. \\ \quad \quad \operatorname{pos\text{-}rr}(e) \Rightarrow  \\ \quad \quad \quad \exists d. \\ \quad \quad \quad \quad \operatorname{pos\text{-}rr}(d) \wedge \\ \quad \quad \quad \quad \forall y \in s. \\ \quad \quad \quad \quad \quad (\operatorname{magnitude}((y - a)) < d) \Rightarrow  \\ \quad \quad \quad \quad \quad \quad \operatorname{magnitude}(((g(y) - g(a)) - ((y - a) \cdot l))) \leq (e \cdot \operatorname{magnitude}((y - a))) \Rightarrow \\ \quad \quad \operatorname{is\text{-}diff\text{-}on}(\operatorname{cc\text{-}normed\text{-}field}, s, g, a, l) \end{array}}}
\vnbprop{cc-continuous-at-add-const}{\vnbstmt{\begin{array}{@{}l@{}} \forall u.; \forall f, g \in (u \to \mathbb{C}), k \in \mathbb{C}, y \in u. \\ \quad u \subseteq \mathbb{C} \wedge \\ \quad \forall z \in u.\; (g(z) \simeq (f(z) + k)) \wedge \\ \quad \forall y \in u. \\ \quad \quad \operatorname{is\text{-}continuous\text{-}at}(\operatorname{subspace\text{-}ms}(\operatorname{nf\text{-}metric\text{-}space}(\operatorname{cc\text{-}normed\text{-}field}), u), \\ \quad \quad \quad \operatorname{nf\text{-}metric\text{-}space}(\operatorname{cc\text{-}normed\text{-}field}), f, y) \Rightarrow \\ \quad \quad \operatorname{is\text{-}continuous\text{-}at}(\operatorname{subspace\text{-}ms}(\operatorname{nf\text{-}metric\text{-}space}(\operatorname{cc\text{-}normed\text{-}field}), u), \\ \quad \quad \quad \operatorname{nf\text{-}metric\text{-}space}(\operatorname{cc\text{-}normed\text{-}field}), g, y) \end{array}}}
\vnbprop{cauchy-convex-segments-except}{\vnbstmt{\begin{array}{@{}l@{}} \forall u, f, x.; \forall a, b, c \in u. \\ \quad \operatorname{is\text{-}convex\text{-}cc}(u) \wedge \\ \quad \operatorname{is\text{-}open}(\operatorname{nf\text{-}metric\text{-}space}(\operatorname{cc\text{-}normed\text{-}field}), u) \wedge \\ \quad f \in (u \to \mathbb{C}) \wedge \\ \quad \forall y \in u. \\ \quad \quad \operatorname{is\text{-}continuous\text{-}at}(\operatorname{subspace\text{-}ms}(\operatorname{nf\text{-}metric\text{-}space}(\operatorname{cc\text{-}normed\text{-}field}), u), \\ \quad \quad \quad \operatorname{nf\text{-}metric\text{-}space}(\operatorname{cc\text{-}normed\text{-}field}), f, y) \wedge \\ \quad x \in u \wedge \\ \quad \operatorname{holomorphic\text{-}on}(\operatorname{difference}(u, \operatorname{singleton}(x)), \operatorname{restrict}(f, \operatorname{difference}(u, \operatorname{singleton}(x)))) \Rightarrow \\ \quad \quad \operatorname{seg\text{-}int}(f, a, c) = (\operatorname{seg\text{-}int}(f, a, b) + \operatorname{seg\text{-}int}(f, b, c)) \end{array}}}
\vnbprop{cauchy-convex-primitive-except}{\vnbstmt{\begin{array}{@{}l@{}} \forall u, f, x.; \forall o \in u. \\ \quad \operatorname{is\text{-}convex\text{-}cc}(u) \wedge \\ \quad \operatorname{is\text{-}open}(\operatorname{nf\text{-}metric\text{-}space}(\operatorname{cc\text{-}normed\text{-}field}), u) \wedge \\ \quad f \in (u \to \mathbb{C}) \wedge \\ \quad \forall y \in u. \\ \quad \quad \operatorname{is\text{-}continuous\text{-}at}(\operatorname{subspace\text{-}ms}(\operatorname{nf\text{-}metric\text{-}space}(\operatorname{cc\text{-}normed\text{-}field}), u), \\ \quad \quad \quad \operatorname{nf\text{-}metric\text{-}space}(\operatorname{cc\text{-}normed\text{-}field}), f, y) \wedge \\ \quad x \in u \wedge \\ \quad \operatorname{holomorphic\text{-}on}(\operatorname{difference}(u, \operatorname{singleton}(x)), \operatorname{restrict}(f, \operatorname{difference}(u, \operatorname{singleton}(x)))) \Rightarrow \\ \quad \quad \exists g \in (u \to \mathbb{C}). \\ \quad \quad \quad (\operatorname{holomorphic\text{-}on}(u, g) \wedge \forall z \in u.\; \operatorname{is\text{-}diff\text{-}on}(\operatorname{cc\text{-}normed\text{-}field}, u, g, z, f(z))) \end{array}}}
\vnbprop{cauchy-convex-closed-except}{\vnbstmt{\begin{array}{@{}l@{}} \forall u, f, x, m, c, a, b. \\ \quad \operatorname{is\text{-}convex\text{-}cc}(u) \wedge \\ \quad \operatorname{is\text{-}open}(\operatorname{nf\text{-}metric\text{-}space}(\operatorname{cc\text{-}normed\text{-}field}), u) \wedge \\ \quad f \in (u \to \mathbb{C}) \wedge \\ \quad \forall y \in u. \\ \quad \quad \operatorname{is\text{-}continuous\text{-}at}(\operatorname{subspace\text{-}ms}(\operatorname{nf\text{-}metric\text{-}space}(\operatorname{cc\text{-}normed\text{-}field}), u), \\ \quad \quad \quad \operatorname{nf\text{-}metric\text{-}space}(\operatorname{cc\text{-}normed\text{-}field}), f, y) \wedge \\ \quad x \in u \wedge \\ \quad \operatorname{holomorphic\text{-}on}(\operatorname{difference}(u, \operatorname{singleton}(x)), \operatorname{restrict}(f, \operatorname{difference}(u, \operatorname{singleton}(x)))) \wedge \\ \quad \operatorname{is\text{-}road}(m, c, a, b) \wedge \\ \quad \operatorname{trace}(m, a, b) \subseteq u \wedge \\ \quad m(b) = m(a) \Rightarrow \\ \quad \quad \operatorname{line\text{-}int}(f, m, c, a, b) = 0 \end{array}}}

\vnbbookfile{theorem-library/road-juxtaposition-laws}{theorem-library/road-juxtaposition-laws}
\noindent{\small road-juxtaposition-laws.scm -- THE JUXTAPOSITION OF ROADS (Dieudonne 9.6, the notes' Proposition 3.1), THE TRIANGULAR ROAD, AND THE LENGTH OF A ROAD (the notes' (53)). Batch 37, 2026-09-26.  The vocabulary is structure-library/road-juxtaposition.scm (JUXTA, ROAD-LENGTH, TRI-ROAD, TRI-ROAD-DERIV).}\par
\vnbprop{juxta-unfold}{\vnbstmt{\begin{array}{@{}l@{}} \forall p, q, a, b, c. \\ \quad (\operatorname{juxta}(p, q, a, b, c) \simeq (\lambda t \in \operatorname{ccint}(a, c).\; \left(\text{if }(t \leq b)\text{ then }p(t)\text{ else }q(t)\right))) \end{array}}}
\vnbprop{road-length-unfold}{\vnbstmt{\begin{array}{@{}l@{}} \forall d, a, b. \\ \quad (\operatorname{road\text{-}length}(d, a, b) \simeq \int_{a}^{b} (\lambda t \in \operatorname{ccint}(a, b).\; \operatorname{magnitude}(d(t)))) \end{array}}}
\vnbprop{tri-road-unfold}{\vnbstmt{\begin{array}{@{}l@{}} \forall a, b, c. \\ \quad (\operatorname{tri\text{-}road}(a, b, c) \simeq \operatorname{juxta}(\operatorname{juxta}(\operatorname{seg\text{-}path}(a, b), (\lambda s \in \operatorname{ccint}(1, 2).\; \operatorname{seg\text{-}path}(b, c)((s - 1))), 0, 1, 2), (\lambda s \in \operatorname{ccint}(2, 3).\; \operatorname{seg\text{-}path}(c, a)((s - 2))), 0, 2, 3)) \end{array}}}
\vnbprop{tri-road-deriv-unfold}{\vnbstmt{\begin{array}{@{}l@{}} \forall a, b, c. \\ \quad (\operatorname{tri\text{-}road\text{-}deriv}(a, b, c) \simeq \operatorname{juxta}(\operatorname{juxta}(\operatorname{seg\text{-}deriv}(a, b), (\lambda s \in \operatorname{ccint}(1, 2).\; \operatorname{seg\text{-}deriv}(b, c)((s - 1))), 0, 1, 2), (\lambda s \in \operatorname{ccint}(2, 3).\; \operatorname{seg\text{-}deriv}(c, a)((s - 2))), 0, 2, 3)) \end{array}}}
\vnbprop{juxta-apply-left}{\vnbstmt{\begin{array}{@{}l@{}} \forall p, q.; \forall a, b, c \in \mathbb{R}, t \in \operatorname{ccint}(a, b). \\ \quad (b \leq c) \Rightarrow (\operatorname{juxta}(p, q, a, b, c)(t) \simeq p(t)) \end{array}}}
\vnbprop{juxta-apply-right}{\vnbstmt{\begin{array}{@{}l@{}} \forall p, q.; \forall a, b, c \in \mathbb{R}, t \in \operatorname{ccint}(b, c). \\ \quad ((a \leq b) \wedge (b < t)) \Rightarrow (\operatorname{juxta}(p, q, a, b, c)(t) \simeq q(t)) \end{array}}}
\vnbprop{juxta-apply-right-closed}{\vnbstmt{\begin{array}{@{}l@{}} \forall p, q.; \forall a, b, c \in \mathbb{R}, t \in \operatorname{ccint}(b, c). \\ \quad ((a \leq b) \wedge (p(b) = q(b))) \Rightarrow (\operatorname{juxta}(p, q, a, b, c)(t) \simeq q(t)) \end{array}}}
\vnbprop{juxta-in-fun}{\vnbstmt{\begin{array}{@{}l@{}} \forall p, q.; \forall a, b, c \in \mathbb{R}, d. \\ \quad ((a \leq b) \wedge ((b \leq c) \wedge ((p \in (\operatorname{ccint}(a, b) \to d)) \wedge (q \in (\operatorname{ccint}(b, c) \to d))))) \Rightarrow  \\ \quad \quad \operatorname{juxta}(p, q, a, b, c) \in (\operatorname{ccint}(a, c) \to d) \end{array}}}
\vnbprop{juxta-trace-subset}{\vnbstmt{\begin{array}{@{}l@{}} \forall p, q.; \forall a, b, c \in \mathbb{R}. \\ \quad ((a \leq b) \wedge ((b \leq c) \wedge ((p \in (\operatorname{ccint}(a, b) \to \mathbb{C})) \wedge (q \in (\operatorname{ccint}(b, c) \to \mathbb{C}))))) \Rightarrow  \\ \quad \quad \operatorname{trace}(\operatorname{juxta}(p, q, a, b, c), a, c) \subseteq (\operatorname{trace}(p, a, b) \cup \operatorname{trace}(q, b, c)) \end{array}}}
\vnbprop{juxta-trace}{\vnbstmt{\begin{array}{@{}l@{}} \forall p, q.; \forall a, b, c \in \mathbb{R}. \\ \quad ((a \leq b) \wedge ((b \leq c) \wedge ((p \in (\operatorname{ccint}(a, b) \to \mathbb{C})) \wedge ((q \in (\operatorname{ccint}(b, c) \to \mathbb{C})) \wedge (p(b) = q(b)))))) \Rightarrow  \\ \quad \quad \operatorname{trace}(\operatorname{juxta}(p, q, a, b, c), a, c) = (\operatorname{trace}(p, a, b) \cup \operatorname{trace}(q, b, c)) \end{array}}}
\vnbprop{right-limit-within-local}{\vnbstmt{\begin{array}{@{}l@{}} \forall g, v, f, w, x, l, r. \\ \quad \operatorname{is\text{-}right\text{-}limit\text{-}within}(g, v, x, l) \wedge \\ \quad w \subseteq \mathbb{R} \wedge \\ \quad f \in (w \to \mathbb{R}) \wedge \\ \quad \operatorname{pos\text{-}rr}(r) \wedge \\ \quad \forall t \in w.\; (((x < t) \wedge (t < (x + r))) \Rightarrow (t \in v)) \wedge \\ \quad \forall t \in w.\; (((x < t) \wedge (t < (x + r))) \Rightarrow (f(t) \simeq g(t))) \Rightarrow \\ \quad \quad \operatorname{is\text{-}right\text{-}limit\text{-}within}(f, w, x, l) \end{array}}}
\vnbprop{left-limit-within-local}{\vnbstmt{\begin{array}{@{}l@{}} \forall g, v, f, w, x, l, r. \\ \quad \operatorname{is\text{-}left\text{-}limit\text{-}within}(g, v, x, l) \wedge \\ \quad w \subseteq \mathbb{R} \wedge \\ \quad f \in (w \to \mathbb{R}) \wedge \\ \quad \operatorname{pos\text{-}rr}(r) \wedge \\ \quad \forall t \in w.\; ((((x - r) < t) \wedge (t < x)) \Rightarrow (t \in v)) \wedge \\ \quad \forall t \in w.\; ((((x - r) < t) \wedge (t < x)) \Rightarrow (f(t) \simeq g(t))) \Rightarrow \\ \quad \quad \operatorname{is\text{-}left\text{-}limit\text{-}within}(f, w, x, l) \end{array}}}
\vnbprop{regulated-on-glue}{\vnbstmt{\begin{array}{@{}l@{}} \forall a, c, b \in \mathbb{R}, d, g, f. \\ \quad a < c \wedge \\ \quad c < b \wedge \\ \quad \operatorname{is\text{-}regulated\text{-}on}(d, a, c) \wedge \\ \quad \operatorname{is\text{-}regulated\text{-}on}(g, c, b) \wedge \\ \quad f \in (\operatorname{ccint}(a, b) \to \mathbb{R}) \wedge \\ \quad \forall t \in \operatorname{ccint}(a, c).\; ((t < c) \Rightarrow (f(t) \simeq d(t))) \wedge \\ \quad \forall t \in \operatorname{ccint}(c, b).\; ((c < t) \Rightarrow (f(t) \simeq g(t))) \Rightarrow \\ \quad \quad \operatorname{is\text{-}regulated\text{-}on}(f, a, b) \end{array}}}
\vnbprop{juxta-is-road}{\vnbstmt{\begin{array}{@{}l@{}} \forall m_{1}, d, m_{2}, f, a, b, c. \\ \quad (\operatorname{is\text{-}road}(m_{1}, d, a, b) \wedge (\operatorname{is\text{-}road}(m_{2}, f, b, c) \wedge (\operatorname{m1}(b) = \operatorname{m2}(b)))) \Rightarrow  \\ \quad \quad \operatorname{is\text{-}road}(\operatorname{juxta}(m_{1}, m_{2}, a, b, c), \operatorname{juxta}(d, f, a, b, c), a, c) \end{array}}}
\vnbprop{line-int-interior-congruence}{\vnbstmt{\begin{array}{@{}l@{}} \forall u, f, m, c, a, b.; \forall d, g \in (\operatorname{ccint}(a, b) \to \mathbb{C}). \\ \quad u \subseteq \mathbb{C} \wedge \\ \quad f \in (u \to \mathbb{C}) \wedge \\ \quad \forall y \in u. \\ \quad \quad \operatorname{is\text{-}continuous\text{-}at}(\operatorname{subspace\text{-}ms}(\operatorname{nf\text{-}metric\text{-}space}(\operatorname{cc\text{-}normed\text{-}field}), u), \\ \quad \quad \quad \operatorname{nf\text{-}metric\text{-}space}(\operatorname{cc\text{-}normed\text{-}field}), f, y) \wedge \\ \quad \operatorname{is\text{-}road}(m, c, a, b) \wedge \\ \quad \operatorname{trace}(m, a, b) \subseteq u \wedge \\ \quad \forall t \in \operatorname{ccint}(a, b).\; (d(t) \simeq m(t)) \wedge \\ \quad \forall t \in \operatorname{ooint}(a, b).\; (g(t) \simeq c(t)) \Rightarrow \\ \quad \quad \operatorname{line\text{-}int}(f, d, g, a, b) = \operatorname{line\text{-}int}(f, m, c, a, b) \end{array}}}
\vnbprop{line-int-juxta}{\vnbstmt{\begin{array}{@{}l@{}} \forall u, f, m_{1}, d, m_{2}, g, a, b, c. \\ \quad u \subseteq \mathbb{C} \wedge \\ \quad f \in (u \to \mathbb{C}) \wedge \\ \quad \forall y \in u. \\ \quad \quad \operatorname{is\text{-}continuous\text{-}at}(\operatorname{subspace\text{-}ms}(\operatorname{nf\text{-}metric\text{-}space}(\operatorname{cc\text{-}normed\text{-}field}), u), \\ \quad \quad \quad \operatorname{nf\text{-}metric\text{-}space}(\operatorname{cc\text{-}normed\text{-}field}), f, y) \wedge \\ \quad \operatorname{is\text{-}road}(m_{1}, d, a, b) \wedge \\ \quad \operatorname{is\text{-}road}(m_{2}, g, b, c) \wedge \\ \quad \operatorname{m1}(b) = \operatorname{m2}(b) \wedge \\ \quad \operatorname{trace}(m_{1}, a, b) \subseteq u \wedge \\ \quad \operatorname{trace}(m_{2}, b, c) \subseteq u \Rightarrow \\ \quad \quad \operatorname{line\text{-}int}(f, \operatorname{juxta}(m_{1}, m_{2}, a, b, c), \operatorname{juxta}(d, g, a, b, c), a, c) \\ \quad \quad =\; \operatorname{line\text{-}int}(f, m_{1}, d, a, b) + \operatorname{line\text{-}int}(f, m_{2}, g, b, c) \end{array}}}
\vnbprop{line-int-concat}{\vnbstmt{\begin{array}{@{}l@{}} \forall u, f, m, g, o, h, a, b, c, d. \\ \quad u \subseteq \mathbb{C} \wedge \\ \quad f \in (u \to \mathbb{C}) \wedge \\ \quad \forall y \in u. \\ \quad \quad \operatorname{is\text{-}continuous\text{-}at}(\operatorname{subspace\text{-}ms}(\operatorname{nf\text{-}metric\text{-}space}(\operatorname{cc\text{-}normed\text{-}field}), u), \\ \quad \quad \quad \operatorname{nf\text{-}metric\text{-}space}(\operatorname{cc\text{-}normed\text{-}field}), f, y) \wedge \\ \quad \operatorname{is\text{-}road}(m, g, a, b) \wedge \\ \quad \operatorname{is\text{-}road}(o, h, c, d) \wedge \\ \quad m(b) = o(c) \wedge \\ \quad \operatorname{trace}(m, a, b) \subseteq u \wedge \\ \quad \operatorname{trace}(o, c, d) \subseteq u \Rightarrow \\ \quad \quad \operatorname{line\text{-}int}(f, \\ \quad \quad \quad \operatorname{juxta}(m, (\lambda s \in \operatorname{ccint}(b, (b + (d - c))).\; o((c + (s - b)))), a, b, (b + (d - c))), \\ \quad \quad \quad \operatorname{juxta}(g, (\lambda s \in \operatorname{ccint}(b, (b + (d - c))).\; h((c + (s - b)))), a, b, (b + (d - c))), \\ \quad \quad \quad a, (b + (d - c))) \\ \quad \quad =\; \operatorname{line\text{-}int}(f, m, g, a, b) + \operatorname{line\text{-}int}(f, o, h, c, d) \end{array}}}
\vnbprop{road-length-regulated}{\vnbstmt{\begin{array}{@{}l@{}} \forall m, c, a, b. \\ \quad \operatorname{is\text{-}road}(m, c, a, b) \Rightarrow  \\ \quad \quad \operatorname{is\text{-}regulated\text{-}on}((\lambda t \in \operatorname{ccint}(a, b).\; \operatorname{magnitude}(c(t))), a, b) \end{array}}}
\vnbprop{road-length-primitive}{\vnbstmt{\begin{array}{@{}l@{}} \forall m, c, a, b. \\ \quad \operatorname{is\text{-}road}(m, c, a, b) \Rightarrow  \\ \quad \quad \exists g.\; \operatorname{is\text{-}primitive}(g, (\lambda t \in \operatorname{ccint}(a, b).\; \operatorname{magnitude}(c(t))), a, b) \end{array}}}
\vnbprop{road-length-in-rr}{\vnbstmt{\forall m, c, a, b.\; (\operatorname{is\text{-}road}(m, c, a, b) \Rightarrow (\operatorname{road\text{-}length}(c, a, b) \in \mathbb{R}))}}
\vnbprop{road-length-nonneg}{\vnbstmt{\forall m, c, a, b.\; (\operatorname{is\text{-}road}(m, c, a, b) \Rightarrow (0 \leq \operatorname{road\text{-}length}(c, a, b)))}}
\vnbprop{pw-int-interior-congruence}{\vnbstmt{\begin{array}{@{}l@{}} \forall f, i, \mathit{ppsi}, a, b. \\ \quad \operatorname{is\text{-}primitive}(f, i, a, b) \wedge \\ \quad \mathit{ppsi} \in (\operatorname{ccint}(a, b) \to \mathbb{R}) \wedge \\ \quad \forall y \in \operatorname{ooint}(a, b).\; (i(y) \simeq \operatorname{ppsi}(y)) \Rightarrow \\ \quad \quad \int_{a}^{b} i = \int_{a}^{b} \mathit{ppsi} \end{array}}}
\vnbprop{road-length-juxta}{\vnbstmt{\begin{array}{@{}l@{}} \forall m_{1}, d, m_{2}, f, a, b, c. \\ \quad (\operatorname{is\text{-}road}(m_{1}, d, a, b) \wedge (\operatorname{is\text{-}road}(m_{2}, f, b, c) \wedge (\operatorname{m1}(b) = \operatorname{m2}(b)))) \Rightarrow  \\ \quad \quad \operatorname{road\text{-}length}(\operatorname{juxta}(d, f, a, b, c), a, c) = (\operatorname{road\text{-}length}(d, a, b) + \operatorname{road\text{-}length}(f, b, c)) \end{array}}}
\vnbprop{seg-shift-is-road}{\vnbstmt{\begin{array}{@{}l@{}} \forall p, q \in \mathbb{C}, k, m \in \mathbb{R}. \\ \quad (m = (k + 1)) \Rightarrow  \\ \quad \quad \operatorname{is\text{-}road}((\lambda s \in \operatorname{ccint}(k, m).\; \operatorname{seg\text{-}path}(p, q)((s - k))), \\ \quad \quad \quad (\lambda s \in \operatorname{ccint}(k, m).\; \operatorname{seg\text{-}deriv}(p, q)((s - k))), k, m) \end{array}}}
\vnbprop{seg-shift-trace}{\vnbstmt{\begin{array}{@{}l@{}} \forall p, q \in \mathbb{C}, k, m \in \mathbb{R}. \\ \quad (m = (k + 1)) \Rightarrow  \\ \quad \quad \operatorname{trace}((\lambda s \in \operatorname{ccint}(k, m).\; \operatorname{seg\text{-}path}(p, q)((s - k))), k, m) = \operatorname{trace}(\operatorname{seg\text{-}path}(p, q), 0, 1) \end{array}}}
\vnbprop{road-length-segment}{\vnbstmt{\begin{array}{@{}l@{}} \forall p, q \in \mathbb{C}. \\ \quad \operatorname{road\text{-}length}(\operatorname{seg\text{-}deriv}(p, q), 0, 1) = \operatorname{magnitude}((q - p)) \end{array}}}
\vnbprop{seg-shift-length}{\vnbstmt{\begin{array}{@{}l@{}} \forall p, q \in \mathbb{C}, k, m \in \mathbb{R}. \\ \quad (m = (k + 1)) \Rightarrow  \\ \quad \quad \operatorname{road\text{-}length}((\lambda s \in \operatorname{ccint}(k, m).\; \operatorname{seg\text{-}deriv}(p, q)((s - k))), k, m) = \operatorname{magnitude}((q - p)) \end{array}}}
\vnbprop{seg-shift-line-int}{\vnbstmt{\begin{array}{@{}l@{}} \forall u, f.; \forall p, q \in \mathbb{C}, k, m \in \mathbb{R}. \\ \quad u \subseteq \mathbb{C} \wedge \\ \quad f \in (u \to \mathbb{C}) \wedge \\ \quad \forall y \in u. \\ \quad \quad \operatorname{is\text{-}continuous\text{-}at}(\operatorname{subspace\text{-}ms}(\operatorname{nf\text{-}metric\text{-}space}(\operatorname{cc\text{-}normed\text{-}field}), u), \\ \quad \quad \quad \operatorname{nf\text{-}metric\text{-}space}(\operatorname{cc\text{-}normed\text{-}field}), f, y) \wedge \\ \quad m = (k + 1) \wedge \\ \quad \operatorname{trace}(\operatorname{seg\text{-}path}(p, q), 0, 1) \subseteq u \Rightarrow \\ \quad \quad \operatorname{line\text{-}int}(f, (\lambda s \in \operatorname{ccint}(k, m).\; \operatorname{seg\text{-}path}(p, q)((s - k))), \\ \quad \quad \quad (\lambda s \in \operatorname{ccint}(k, m).\; \operatorname{seg\text{-}deriv}(p, q)((s - k))), k, m) \\ \quad \quad =\; \operatorname{line\text{-}int}(f, \operatorname{seg\text{-}path}(p, q), \operatorname{seg\text{-}deriv}(p, q), 0, 1) \end{array}}}
\vnbprop{tri-road-is-road}{\vnbstmt{\begin{array}{@{}l@{}} \forall a, b, c \in \mathbb{C}. \\ \quad \operatorname{is\text{-}road}(\operatorname{tri\text{-}road}(a, b, c), \operatorname{tri\text{-}road\text{-}deriv}(a, b, c), 0, 3) \end{array}}}
\vnbprop{tri-road-trace}{\vnbstmt{\begin{array}{@{}l@{}} \forall a, b, c \in \mathbb{C}. \\ \quad \operatorname{trace}(\operatorname{tri\text{-}road}(a, b, c), 0, 3) \\ \quad =\; ((\operatorname{trace}(\operatorname{seg\text{-}path}(a, b), 0, 1) \cup \operatorname{trace}(\operatorname{seg\text{-}path}(b, c), 0, 1)) \cup \operatorname{trace}(\operatorname{seg\text{-}path}(c, a), 0, 1)) \end{array}}}
\vnbprop{tri-int-is-line-int}{\vnbstmt{\begin{array}{@{}l@{}} \forall u, f.; \forall a, b, c \in \mathbb{C}. \\ \quad u \subseteq \mathbb{C} \wedge \\ \quad f \in (u \to \mathbb{C}) \wedge \\ \quad \forall y \in u. \\ \quad \quad \operatorname{is\text{-}continuous\text{-}at}(\operatorname{subspace\text{-}ms}(\operatorname{nf\text{-}metric\text{-}space}(\operatorname{cc\text{-}normed\text{-}field}), u), \\ \quad \quad \quad \operatorname{nf\text{-}metric\text{-}space}(\operatorname{cc\text{-}normed\text{-}field}), f, y) \wedge \\ \quad \operatorname{trace}(\operatorname{seg\text{-}path}(a, b), 0, 1) \subseteq u \wedge \\ \quad \operatorname{trace}(\operatorname{seg\text{-}path}(b, c), 0, 1) \subseteq u \wedge \\ \quad \operatorname{trace}(\operatorname{seg\text{-}path}(c, a), 0, 1) \subseteq u \Rightarrow \\ \quad \quad \operatorname{tri\text{-}int}(f, a, b, c) = \operatorname{line\text{-}int}(f, \operatorname{tri\text{-}road}(a, b, c), \operatorname{tri\text{-}road\text{-}deriv}(a, b, c), 0, 3) \end{array}}}
\vnbprop{road-length-tri}{\vnbstmt{\begin{array}{@{}l@{}} \forall a, b, c \in \mathbb{C}. \\ \quad \operatorname{road\text{-}length}(\operatorname{tri\text{-}road\text{-}deriv}(a, b, c), 0, 3) \\ \quad =\; (\operatorname{magnitude}((b - a)) + \operatorname{magnitude}((c - b))) + \operatorname{magnitude}((a - c)) \end{array}}}
\vnbprop{line-int-length-bound}{\vnbstmt{\begin{array}{@{}l@{}} \forall u, f, m, c, a, b.; \forall d \in \mathbb{R}. \\ \quad u \subseteq \mathbb{C} \wedge \\ \quad f \in (u \to \mathbb{C}) \wedge \\ \quad \forall y \in u. \\ \quad \quad \operatorname{is\text{-}continuous\text{-}at}(\operatorname{subspace\text{-}ms}(\operatorname{nf\text{-}metric\text{-}space}(\operatorname{cc\text{-}normed\text{-}field}), u), \\ \quad \quad \quad \operatorname{nf\text{-}metric\text{-}space}(\operatorname{cc\text{-}normed\text{-}field}), f, y) \wedge \\ \quad \operatorname{is\text{-}road}(m, c, a, b) \wedge \\ \quad \operatorname{trace}(m, a, b) \subseteq u \wedge \\ \quad \forall \mathit{bny} \in \operatorname{trace}(m, a, b).\; (\operatorname{magnitude}(f(\mathit{bny})) \leq d) \Rightarrow \\ \quad \quad \operatorname{magnitude}(\operatorname{line\text{-}int}(f, m, c, a, b)) \leq (\operatorname{road\text{-}length}(c, a, b) \cdot d) \end{array}}}
\vnbprop{road-length-affine-reparam}{\vnbstmt{\begin{array}{@{}l@{}} \forall a, b, c, d, l, m \in \mathbb{R}, f, g, k. \\ \quad 0 < l \wedge \\ \quad (m + (l \cdot c)) = a \wedge \\ \quad (m + (l \cdot d)) = b \wedge \\ \quad c < d \wedge \\ \quad \operatorname{is\text{-}road}(f, g, a, b) \wedge \\ \quad k \in (\operatorname{ccint}(c, d) \to \mathbb{C}) \wedge \\ \quad \forall z \in \operatorname{ccint}(c, d).\; (k(z) \simeq (l \cdot g((m + (l \cdot z))))) \Rightarrow \\ \quad \quad \operatorname{road\text{-}length}(k, c, d) = \operatorname{road\text{-}length}(g, a, b) \end{array}}}
\vnbprop{tri-int-is-line-int-conv3}{\vnbstmt{\begin{array}{@{}l@{}} \forall u, f.; \forall a, b, c \in \mathbb{C}. \\ \quad u \subseteq \mathbb{C} \wedge \\ \quad f \in (u \to \mathbb{C}) \wedge \\ \quad \forall y \in u. \\ \quad \quad \operatorname{is\text{-}continuous\text{-}at}(\operatorname{subspace\text{-}ms}(\operatorname{nf\text{-}metric\text{-}space}(\operatorname{cc\text{-}normed\text{-}field}), u), \\ \quad \quad \quad \operatorname{nf\text{-}metric\text{-}space}(\operatorname{cc\text{-}normed\text{-}field}), f, y) \wedge \\ \quad \operatorname{conv3}(a, b, c) \subseteq u \Rightarrow \\ \quad \quad \operatorname{tri\text{-}int}(f, a, b, c) = \operatorname{line\text{-}int}(f, \operatorname{tri\text{-}road}(a, b, c), \operatorname{tri\text{-}road\text{-}deriv}(a, b, c), 0, 3) \end{array}}}

\vnbbookfile{structure-library/cc-elementary}{structure-library/cc-elementary}
\noindent{\small cc-elementary.scm -- THE ELEMENTARY FUNCTIONS exp, cos, sin ON CC, DEFINED BY THEIR POWER SERIES (the notes 2.3.1, 2.3.2). DEFINITIONS ONLY; the laws are PROVEN in theorem-library/cc-exp.scm. Batch 31-A, 2026-09-25.  The user's decision of 2026-09-25: CC-EXP, CC-COS and}\par
\par\medskip\noindent\textbf{Definition} (\texttt{cc-exp}).\;$\exp z :=$
\begin{equation*}
\operatorname{cc\text{-}series\text{-}limit}((\lambda \operatorname{cek\_} \in \mathbb{N}.\; (\operatorname{recip}(\operatorname{factorial}(\operatorname{cek\_})) \cdot \operatorname{power}(z, \operatorname{cek\_}))))
\end{equation*}

\par\medskip\noindent\textbf{Definition} (\texttt{cc-cos}).\;$\cos z :=$
\begin{equation*}
\operatorname{cc\text{-}series\text{-}limit}((\lambda \operatorname{cek\_} \in \mathbb{N}.\; ((\operatorname{power}(-1, \operatorname{cek\_}) \cdot \operatorname{recip}(\operatorname{factorial}((2 \cdot \operatorname{cek\_})))) \cdot \operatorname{power}(z, (2 \cdot \operatorname{cek\_})))))
\end{equation*}

\par\medskip\noindent\textbf{Definition} (\texttt{cc-sin}).\;$\sin z :=$
\begin{equation*}
\operatorname{cc\text{-}series\text{-}limit}((\lambda \operatorname{cek\_} \in \mathbb{N}.\; ((\operatorname{power}(-1, \operatorname{cek\_}) \cdot \operatorname{recip}(\operatorname{factorial}(((2 \cdot \operatorname{cek\_}) + 1)))) \cdot \operatorname{power}(z, ((2 \cdot \operatorname{cek\_}) + 1)))))
\end{equation*}


\vnbbookfile{theorem-library/cc-exp}{theorem-library/cc-exp}
\noindent{\small cc-exp.scm -- THE ELEMENTARY FUNCTIONS exp, cos, sin ON CC (the notes 2.3.1, 2.3.2, everything that does not mention pi). Batch 31-A, 2026-09-25.  The definitions are structure-library/cc-elementary.scm (CC-EXP, CC-COS, CC-SIN by their power series, the user's decision of 2026-09-25).}\par
\vnbprop{cc-exp-unfold}{\vnbstmt{\begin{array}{@{}l@{}} \forall z. \\ \quad (\exp z \simeq \operatorname{cc\text{-}series\text{-}limit}((\lambda k \in \mathbb{N}.\; (\operatorname{recip}(\operatorname{factorial}(k)) \cdot \operatorname{power}(z, k))))) \end{array}}}
\vnbprop{cc-cos-unfold}{\vnbstmt{\begin{array}{@{}l@{}} \forall z. \\ \quad (\cos z \simeq \operatorname{cc\text{-}series\text{-}limit}((\lambda k \in \mathbb{N}.\; ((\operatorname{power}(-1, k) \cdot \operatorname{recip}(\operatorname{factorial}((2 \cdot k)))) \cdot \operatorname{power}(z, (2 \cdot k)))))) \end{array}}}
\vnbprop{cc-sin-unfold}{\vnbstmt{\begin{array}{@{}l@{}} \forall z. \\ \quad (\sin z \simeq \operatorname{cc\text{-}series\text{-}limit}((\lambda k \in \mathbb{N}.\; ((\operatorname{power}(-1, k) \cdot \operatorname{recip}(\operatorname{factorial}(((2 \cdot k) + 1)))) \cdot \operatorname{power}(z, ((2 \cdot k) + 1)))))) \end{array}}}
\vnbprop{cps-radius-pos-inf}{\vnbstmt{\begin{array}{@{}l@{}} \forall f \in (\mathbb{N} \to \mathbb{C}). \\ \quad \forall t \in \mathbb{R}.\; ((0 < t) \Rightarrow \operatorname{series\text{-}converges}((\lambda k \in \mathbb{N}.\; (\operatorname{magnitude}(f(k)) \cdot \operatorname{power}(t, k))))) \Rightarrow  \\ \quad \quad \operatorname{cps\text{-}radius}(f) = \operatorname{pos\text{-}inf} \end{array}}}
\vnbprop{cex-exp-series-real-converges}{\vnbstmt{\begin{array}{@{}l@{}} \forall t \in \mathbb{R}. \\ \quad (0 < t) \Rightarrow \operatorname{series\text{-}converges}((\lambda k \in \mathbb{N}.\; (\operatorname{recip}(\operatorname{factorial}(k)) \cdot \operatorname{power}(t, k)))) \end{array}}}
\vnbprop{cc-exp-radius}{\vnbstmt{\begin{array}{@{}l@{}} \forall f \in (\mathbb{N} \to \mathbb{C}). \\ \quad \forall k \in \mathbb{N}.\; (f(k) = \operatorname{recip}(\operatorname{factorial}(k))) \Rightarrow (\operatorname{cps\text{-}radius}(f) = \operatorname{pos\text{-}inf}) \end{array}}}
\vnbprop{cps-entire-converges}{\vnbstmt{\begin{array}{@{}l@{}} \forall f \in (\mathbb{N} \to \mathbb{C}), z \in \mathbb{C}. \\ \quad (\operatorname{cps\text{-}radius}(f) = \operatorname{pos\text{-}inf}) \Rightarrow  \\ \quad \quad \operatorname{cc\text{-}series\text{-}converges}((\lambda k \in \mathbb{N}.\; (f(k) \cdot \operatorname{power}(z, k)))) \end{array}}}
\vnbprop{cc-diff-on-local}{\vnbstmt{\begin{array}{@{}l@{}} \forall u.; \forall f \in (u \to \mathbb{C}), v, g, a, l. \\ \quad \operatorname{is\text{-}open}(\operatorname{nf\text{-}metric\text{-}space}(\operatorname{cc\text{-}normed\text{-}field}), u) \wedge \\ \quad v \subseteq u \wedge \\ \quad \operatorname{is\text{-}diff\text{-}on}(\operatorname{cc\text{-}normed\text{-}field}, v, g, a, l) \wedge \\ \quad \forall x \in v.\; (f(x) = g(x)) \Rightarrow \\ \quad \quad \operatorname{is\text{-}diff\text{-}on}(\operatorname{cc\text{-}normed\text{-}field}, u, f, a, l) \end{array}}}
\vnbprop{cps-entire-diff}{\vnbstmt{\begin{array}{@{}l@{}} \forall f, a \in (\mathbb{N} \to \mathbb{C}), \mathit{cpda} \in \mathbb{C}. \\ \quad (\forall j \in \mathbb{N}.\; (a(j) = ((j + 1) \cdot f((j + 1)))) \wedge (\operatorname{cps\text{-}radius}(f) = \operatorname{pos\text{-}inf})) \Rightarrow  \\ \quad \quad \operatorname{is\text{-}diff\text{-}on}(\operatorname{cc\text{-}normed\text{-}field}, \mathbb{C}, \\ \quad \quad \quad (\lambda z \in \mathbb{C}.\; \operatorname{cc\text{-}series\text{-}limit}((\lambda k \in \mathbb{N}.\; (f(k) \cdot \operatorname{power}(z, k))))), \\ \quad \quad \quad \mathit{cpda}, \\ \quad \quad \quad \operatorname{cc\text{-}series\text{-}limit}((\lambda k \in \mathbb{N}.\; (a(k) \cdot \operatorname{power}(\mathit{cpda}, k))))) \end{array}}}
\vnbprop{cc-exp-as-power-series}{\vnbstmt{\begin{array}{@{}l@{}} \forall f \in (\mathbb{N} \to \mathbb{C}), z \in \mathbb{C}. \\ \quad \forall k \in \mathbb{N}.\; (f(k) = \operatorname{recip}(\operatorname{factorial}(k))) \Rightarrow  \\ \quad \quad \operatorname{cc\text{-}series\text{-}converges}((\lambda k \in \mathbb{N}.\; (\operatorname{recip}(\operatorname{factorial}(k)) \cdot \operatorname{power}(z, k)))) \wedge \\ \quad \quad \exp z = \operatorname{cc\text{-}series\text{-}limit}((\lambda \mathit{cpsk} \in \mathbb{N}.\; (f(\mathit{cpsk}) \cdot \operatorname{power}(z, \mathit{cpsk})))) \end{array}}}
\vnbprop{cc-exp-converges}{\vnbstmt{\begin{array}{@{}l@{}} \forall z \in \mathbb{C}. \\ \quad (\operatorname{cc\text{-}series\text{-}converges}((\lambda k \in \mathbb{N}.\; (\operatorname{recip}(\operatorname{factorial}(k)) \cdot \operatorname{power}(z, k)))) \wedge (\exp z \in \mathbb{C})) \end{array}}}
\vnbprop{cc-exp-deriv}{\vnbstmt{\begin{array}{@{}l@{}} \forall z \in \mathbb{C}. \\ \quad \operatorname{is\text{-}diff\text{-}on}(\operatorname{cc\text{-}normed\text{-}field}, \mathbb{C}, (\lambda y \in \mathbb{C}.\; \exp y), z, \exp z) \end{array}}}
\vnbprop{cc-exp-holomorphic}{\vnbstmt{\operatorname{holomorphic\text{-}on}(\mathbb{C}, (\lambda y \in \mathbb{C}.\; \exp y))}}
\vnbprop{cc-exp-zero}{\vnbstmt{\exp 0 = 1}}
\vnbprop{cex-sub-deriv}{\vnbstmt{\begin{array}{@{}l@{}} \forall w, z \in \mathbb{C}. \\ \quad \operatorname{is\text{-}diff\text{-}on}(\operatorname{cc\text{-}normed\text{-}field}, \mathbb{C}, (\lambda y \in \mathbb{C}.\; (w - y)), z, -1) \end{array}}}
\vnbprop{cex-shift-deriv-zero}{\vnbstmt{\begin{array}{@{}l@{}} \forall w, z \in \mathbb{C}. \\ \quad \operatorname{is\text{-}diff\text{-}on}(\operatorname{cc\text{-}normed\text{-}field}, \mathbb{C}, (\lambda y \in \mathbb{C}.\; (\exp y \cdot \exp (w - y))), z, 0) \end{array}}}
\vnbprop{cc-exp-shift}{\vnbstmt{\forall w, z \in \mathbb{C}.\; ((\exp z \cdot \exp (w - z)) = \exp w)}}
\vnbprop{cc-exp-add}{\vnbstmt{\forall z, w \in \mathbb{C}.\; (\exp (z + w) = (\exp z \cdot \exp w))}}
\vnbprop{cc-exp-nonzero}{\vnbstmt{\forall z \in \mathbb{C}.\; \neg (\exp z = 0)}}
\vnbprop{cc-exp-maps-to-nonzero}{\vnbstmt{\forall z \in \mathbb{C}.\; (\exp z \in \operatorname{difference}(\mathbb{C}, \operatorname{singleton}(0)))}}
\vnbprop{cc-exp-neg}{\vnbstmt{\forall z \in \mathbb{C}.\; (\exp -z = \operatorname{recip}(\exp z))}}
\vnbprop{cc-exp-continuous}{\vnbstmt{\begin{array}{@{}l@{}} \forall z \in \mathbb{C}. \\ \quad \operatorname{is\text{-}continuous\text{-}at}(\operatorname{subspace\text{-}ms}(\operatorname{nf\text{-}metric\text{-}space}(\operatorname{cc\text{-}normed\text{-}field}), \mathbb{C}), \\ \quad \quad \operatorname{nf\text{-}metric\text{-}space}(\operatorname{cc\text{-}normed\text{-}field}), (\lambda y \in \mathbb{C}.\; \exp y), z) \end{array}}}
\vnbprop{cc-conjugate-sub}{\vnbstmt{\forall a, b \in \mathbb{C}.\; (\operatorname{conjugate}((a - b)) = (\operatorname{conjugate}(a) - \operatorname{conjugate}(b)))}}
\vnbprop{cc-magnitude-conjugate}{\vnbstmt{\forall z \in \mathbb{C}.\; (\operatorname{magnitude}(\operatorname{conjugate}(z)) = \operatorname{magnitude}(z))}}
\vnbprop{cc-conjugate-power}{\vnbstmt{\forall k \in \mathbb{N}, z \in \mathbb{C}.\; (\operatorname{conjugate}(\operatorname{power}(z, k)) = \operatorname{power}(\operatorname{conjugate}(z), k))}}
\vnbprop{cc-series-limit-conjugate}{\vnbstmt{\begin{array}{@{}l@{}} \forall f, g \in (\mathbb{N} \to \mathbb{C}). \\ \quad (\forall j \in \mathbb{N}.\; (g(j) = \operatorname{conjugate}(f(j))) \wedge \operatorname{cc\text{-}series\text{-}converges}(f)) \Rightarrow  \\ \quad \quad \operatorname{cc\text{-}series\text{-}converges}(g) \wedge (\operatorname{cc\text{-}series\text{-}limit}(g) = \operatorname{conjugate}(\operatorname{cc\text{-}series\text{-}limit}(f))) \end{array}}}
\vnbprop{cc-exp-conj}{\vnbstmt{\forall z \in \mathbb{C}.\; (\exp \operatorname{conjugate}(z) = \operatorname{conjugate}(\exp z))}}
\vnbprop{cc-magnitude-sq-conjugate}{\vnbstmt{\forall z \in \mathbb{C}.\; ((\operatorname{magnitude}(z) \cdot \operatorname{magnitude}(z)) = (z \cdot \operatorname{conjugate}(z)))}}
\vnbprop{cc-exp-real-valued}{\vnbstmt{\forall t \in \mathbb{R}.\; (\exp t \in \mathbb{R})}}
\vnbprop{cc-exp-real}{\vnbstmt{\forall t \in \mathbb{R}.\; ((\exp t \in \mathbb{R}) \wedge (0 < \exp t))}}
\vnbprop{cc-exp-magnitude-imaginary}{\vnbstmt{\forall t \in \mathbb{R}.\; (\operatorname{magnitude}(\exp (+i \cdot t)) = 1)}}
\vnbprop{cc-power-add}{\vnbstmt{\forall n, m \in \mathbb{N}, z \in \mathbb{C}.\; (\operatorname{power}(z, (m + n)) = (\operatorname{power}(z, m) \cdot \operatorname{power}(z, n)))}}
\vnbprop{cc-power-mul}{\vnbstmt{\forall k, m \in \mathbb{N}, z \in \mathbb{C}.\; (\operatorname{power}(z, (m \cdot k)) = \operatorname{power}(\operatorname{power}(z, m), k))}}
\vnbprop{cc-power-mul-base}{\vnbstmt{\forall k \in \mathbb{N}, a, b \in \mathbb{C}.\; (\operatorname{power}((a \cdot b), k) = (\operatorname{power}(a, k) \cdot \operatorname{power}(b, k)))}}
\vnbprop{factorial-mono}{\vnbstmt{\forall m, k \in \mathbb{N}.\; (\operatorname{factorial}(k) \leq \operatorname{factorial}((k + m)))}}
\vnbprop{cex-sign-magnitude}{\vnbstmt{\forall k \in \mathbb{N}, a \in \mathbb{R}.\; ((0 \leq a) \Rightarrow (\operatorname{magnitude}((\operatorname{power}(-1, k) \cdot a)) = a))}}
\vnbprop{cc-cos-radius}{\vnbstmt{\begin{array}{@{}l@{}} \forall f \in (\mathbb{N} \to \mathbb{C}). \\ \quad \forall k \in \mathbb{N}.\; (f(k) = (\operatorname{power}(-1, k) \cdot \operatorname{recip}(\operatorname{factorial}((2 \cdot k))))) \Rightarrow  \\ \quad \quad \operatorname{cps\text{-}radius}(f) = \operatorname{pos\text{-}inf} \end{array}}}
\vnbprop{cc-sin-radius}{\vnbstmt{\begin{array}{@{}l@{}} \forall f \in (\mathbb{N} \to \mathbb{C}). \\ \quad \forall k \in \mathbb{N}.\; (f(k) = (\operatorname{power}(-1, k) \cdot \operatorname{recip}(\operatorname{factorial}(((2 \cdot k) + 1))))) \Rightarrow  \\ \quad \quad \operatorname{cps\text{-}radius}(f) = \operatorname{pos\text{-}inf} \end{array}}}
\vnbprop{cc-cos-as-power-series}{\vnbstmt{\begin{array}{@{}l@{}} \forall f \in (\mathbb{N} \to \mathbb{C}), z \in \mathbb{C}. \\ \quad \forall k \in \mathbb{N}.\; (f(k) = (\operatorname{power}(-1, k) \cdot \operatorname{recip}(\operatorname{factorial}((2 \cdot k))))) \Rightarrow  \\ \quad \quad \operatorname{cc\text{-}series\text{-}converges}((\lambda k \in \mathbb{N}.\; ((\operatorname{power}(-1, k) \cdot \operatorname{recip}(\operatorname{factorial}((2 \cdot k)))) \cdot \operatorname{power}(z, (2 \cdot k))))) \wedge \\ \quad \quad \cos z = \operatorname{cc\text{-}series\text{-}limit}((\lambda \mathit{cpsk} \in \mathbb{N}.\; (f(\mathit{cpsk}) \cdot \operatorname{power}(\operatorname{power}(z, 2), \mathit{cpsk})))) \end{array}}}
\vnbprop{cc-sin-as-power-series}{\vnbstmt{\begin{array}{@{}l@{}} \forall f \in (\mathbb{N} \to \mathbb{C}), z \in \mathbb{C}. \\ \quad \forall k \in \mathbb{N}.\; (f(k) = (\operatorname{power}(-1, k) \cdot \operatorname{recip}(\operatorname{factorial}(((2 \cdot k) + 1))))) \Rightarrow  \\ \quad \quad \operatorname{cc\text{-}series\text{-}converges}((\lambda k \in \mathbb{N}.\; ((\operatorname{power}(-1, k) \cdot \operatorname{recip}(\operatorname{factorial}(((2 \cdot k) + 1)))) \cdot \operatorname{power}(z, ((2 \cdot k) + 1))))) \wedge \\ \quad \quad \sin z = (z \cdot \operatorname{cc\text{-}series\text{-}limit}((\lambda \mathit{cpsk} \in \mathbb{N}.\; (f(\mathit{cpsk}) \cdot \operatorname{power}(\operatorname{power}(z, 2), \mathit{cpsk}))))) \end{array}}}
\vnbprop{cc-cos-converges}{\vnbstmt{\begin{array}{@{}l@{}} \forall z \in \mathbb{C}. \\ \quad \operatorname{cc\text{-}series\text{-}converges}((\lambda k \in \mathbb{N}.\; ((\operatorname{power}(-1, k) \cdot \operatorname{recip}(\operatorname{factorial}((2 \cdot k)))) \cdot \operatorname{power}(z, (2 \cdot k))))) \wedge \\ \quad \cos z \in \mathbb{C} \end{array}}}
\vnbprop{cc-sin-converges}{\vnbstmt{\begin{array}{@{}l@{}} \forall z \in \mathbb{C}. \\ \quad \operatorname{cc\text{-}series\text{-}converges}((\lambda k \in \mathbb{N}.\; ((\operatorname{power}(-1, k) \cdot \operatorname{recip}(\operatorname{factorial}(((2 \cdot k) + 1)))) \cdot \operatorname{power}(z, ((2 \cdot k) + 1))))) \wedge \\ \quad \sin z \in \mathbb{C} \end{array}}}
\vnbprop{cc-series-even-odd}{\vnbstmt{\begin{array}{@{}l@{}} \forall f, e, o \in (\mathbb{N} \to \mathbb{C}). \\ \quad \forall j \in \mathbb{N}.\; (e(j) = f((2 \cdot j))) \wedge \\ \quad \forall j \in \mathbb{N}.\; (o(j) = f(((2 \cdot j) + 1))) \wedge \\ \quad \operatorname{cc\text{-}series\text{-}converges}(f) \wedge \\ \quad \operatorname{cc\text{-}series\text{-}converges}(e) \wedge \\ \quad \operatorname{cc\text{-}series\text{-}converges}(o) \Rightarrow \\ \quad \quad \operatorname{cc\text{-}series\text{-}limit}(f) = (\operatorname{cc\text{-}series\text{-}limit}(e) + \operatorname{cc\text{-}series\text{-}limit}(o)) \end{array}}}
\vnbprop{cc-exp-i}{\vnbstmt{\forall z \in \mathbb{C}.\; (\exp (+i \cdot z) = (\cos z + (+i \cdot \sin z)))}}
\vnbprop{cc-cos-even}{\vnbstmt{\forall z \in \mathbb{C}.\; (\cos -z = \cos z)}}
\vnbprop{cc-sin-odd}{\vnbstmt{\forall z \in \mathbb{C}.\; (\sin -z = -\sin z)}}
\vnbprop{cc-sin-exp}{\vnbstmt{\begin{array}{@{}l@{}} \forall z \in \mathbb{C}. \\ \quad \sin z = ((\exp (+i \cdot z) - \exp -(+i \cdot z)) \cdot \operatorname{recip}(+2i)) \end{array}}}
\vnbprop{cc-cos-exp}{\vnbstmt{\forall z \in \mathbb{C}.\; (\cos z = ((\exp (+i \cdot z) + \exp -(+i \cdot z)) \cdot \operatorname{recip}(2)))}}
\vnbprop{cex-lin-deriv}{\vnbstmt{\begin{array}{@{}l@{}} \forall c, z \in \mathbb{C}. \\ \quad \operatorname{is\text{-}diff\text{-}on}(\operatorname{cc\text{-}normed\text{-}field}, \mathbb{C}, (\lambda y \in \mathbb{C}.\; (c \cdot y)), z, c) \end{array}}}
\vnbprop{cc-exp-scaled-deriv}{\vnbstmt{\begin{array}{@{}l@{}} \forall c, z \in \mathbb{C}. \\ \quad \operatorname{is\text{-}diff\text{-}on}(\operatorname{cc\text{-}normed\text{-}field}, \mathbb{C}, (\lambda y \in \mathbb{C}.\; \exp (c \cdot y)), z, (c \cdot \exp (c \cdot z))) \end{array}}}
\vnbprop{cc-diff-on-scale}{\vnbstmt{\begin{array}{@{}l@{}} \forall u, f, a, l.; \forall c \in \mathbb{C}. \\ \quad \operatorname{is\text{-}diff\text{-}on}(\operatorname{cc\text{-}normed\text{-}field}, u, f, a, l) \Rightarrow  \\ \quad \quad \operatorname{is\text{-}diff\text{-}on}(\operatorname{cc\text{-}normed\text{-}field}, u, (\lambda y \in u.\; (c \cdot f(y))), a, (c \cdot l)) \end{array}}}
\vnbprop{cex-exp-comb-deriv}{\vnbstmt{\begin{array}{@{}l@{}} \forall a, b, c, d, z \in \mathbb{C}. \\ \quad \operatorname{is\text{-}diff\text{-}on}(\operatorname{cc\text{-}normed\text{-}field}, \mathbb{C}, \\ \quad \quad (\lambda y \in \mathbb{C}.\; ((a \cdot \exp (c \cdot y)) + (b \cdot \exp (d \cdot y)))), \\ \quad \quad z, ((a \cdot (c \cdot \exp (c \cdot z))) + (b \cdot (d \cdot \exp (d \cdot z))))) \end{array}}}
\vnbprop{cc-cos-deriv}{\vnbstmt{\begin{array}{@{}l@{}} \forall z \in \mathbb{C}. \\ \quad \operatorname{is\text{-}diff\text{-}on}(\operatorname{cc\text{-}normed\text{-}field}, \mathbb{C}, (\lambda y \in \mathbb{C}.\; \cos y), z, -\sin z) \end{array}}}
\vnbprop{cc-sin-deriv}{\vnbstmt{\begin{array}{@{}l@{}} \forall z \in \mathbb{C}. \\ \quad \operatorname{is\text{-}diff\text{-}on}(\operatorname{cc\text{-}normed\text{-}field}, \mathbb{C}, (\lambda y \in \mathbb{C}.\; \sin y), z, \cos z) \end{array}}}
\vnbprop{cc-cos-holomorphic}{\vnbstmt{\operatorname{holomorphic\text{-}on}(\mathbb{C}, (\lambda y \in \mathbb{C}.\; \cos y))}}
\vnbprop{cc-sin-holomorphic}{\vnbstmt{\operatorname{holomorphic\text{-}on}(\mathbb{C}, (\lambda y \in \mathbb{C}.\; \sin y))}}
\vnbprop{cc-cos-add}{\vnbstmt{\begin{array}{@{}l@{}} \forall z, w \in \mathbb{C}. \\ \quad \cos (z + w) = ((\cos z \cdot \cos w) - (\sin z \cdot \sin w)) \end{array}}}
\vnbprop{cc-sin-add}{\vnbstmt{\begin{array}{@{}l@{}} \forall z, w \in \mathbb{C}. \\ \quad \sin (z + w) = ((\cos w \cdot \sin z) + (\cos z \cdot \sin w)) \end{array}}}
\vnbprop{cc-cos-sin-one}{\vnbstmt{\forall z \in \mathbb{C}.\; ((\operatorname{power}(\cos z, 2) + \operatorname{power}(\sin z, 2)) = 1)}}

\vnbbookfile{structure-library/cc-pi}{structure-library/cc-pi}
\noindent{\small cc-pi.scm -- THE CONSTANT pi, DEFINED FROM THE KERNEL OF exp (the notes 2.3.3). DEFINITION ONLY; existence, uniqueness and the characterisation `pi-spec' are PROVEN in theorem-library/cc-exp-kernel.scm. Batch 31-B, 2026-09-25.}\par
\par\medskip\noindent\textbf{Definition} (\texttt{pi}).\;
\begin{equation*}
(\pi \simeq (\iota\, \operatorname{cekp\_}.\;((\operatorname{cekp\_} \in \mathbb{R}) \wedge ((0 < \operatorname{cekp\_}) \wedge \forall \operatorname{cekz\_} \in \mathbb{C}.\; ((\exp \operatorname{cekz\_} = 1) \Leftrightarrow \exists \operatorname{cekk\_} \in \mathbb{Z}.\; (\operatorname{cekz\_} = \cdot(\operatorname{cekk\_}, \cdot(\cdot(2, \operatorname{cekp\_}), +i))))))))
\end{equation*}


\vnbbookfile{structure-library/analytic-log}{structure-library/analytic-log}
\noindent{\small analytic-log.scm -- ANALYTICITY AT A POINT (the notes' Def 2.11) AND THE PRINCIPAL LOGARITHM (the notes 2.3.3.1, after Prop 2.19). DEFINITIONS ONLY; the laws are PROVEN in theorem-library/analytic-log-laws.scm. Batch 38, 2026-09-27 (the user's decision of the same day: IS-ANALYTIC-AT}\par
\par\medskip\noindent\textbf{Definition} (\texttt{is-analytic-at}).\;For all $\operatorname{anlf\_}$, $\operatorname{anlu\_}$ and $\operatorname{anla\_}$, 
\begin{equation*}
\begin{array}{@{}l@{}} \operatorname{anlf\_} \text{ is analytic at } \operatorname{anla\_} \\ \Leftrightarrow\; \operatorname{is\text{-}open}(\operatorname{cc\text{-}ms}, \operatorname{anlu\_}) \wedge \\ \operatorname{anlf\_} \in (\operatorname{anlu\_} \to \mathbb{C}) \wedge \\ \operatorname{anla\_} \in \operatorname{anlu\_} \wedge \\ \exists \operatorname{anlcf\_} \in (\mathbb{N} \to \mathbb{C}), \operatorname{anlr\_} \in \mathbb{R}. \\ \quad 0 < \operatorname{anlr\_} \wedge \\ \quad \operatorname{anlr\_} < \operatorname{cps\text{-}radius}(\operatorname{anlcf\_}) \wedge \\ \quad \operatorname{ball}(\operatorname{cc\text{-}ms}, \operatorname{anla\_}, \operatorname{anlr\_}) \subseteq \operatorname{anlu\_} \wedge \\ \quad \forall \operatorname{anlz\_} \in \operatorname{ball}(\operatorname{cc\text{-}ms}, \operatorname{anla\_}, \operatorname{anlr\_}). \\ \quad \quad \operatorname{anlf\_}(\operatorname{anlz\_}) = \operatorname{cc\text{-}series\text{-}limit}((\lambda \operatorname{anlk\_} \in \mathbb{N}.\; \cdot(\operatorname{anlcf\_}(\operatorname{anlk\_}), \operatorname{power}(-(\operatorname{anlz\_}, \operatorname{anla\_}), \operatorname{anlk\_})))) \end{array}
\end{equation*}

\par\medskip\noindent\textbf{Definition} (\texttt{cc-log}).\;$\log\left(\operatorname{clgw\_}\right) :=$
\begin{equation*}
(\iota\, \operatorname{clgz\_}.\;((\operatorname{clgz\_} \in \mathbb{C}) \wedge ((-\pi < \operatorname{imag\text{-}part}(\operatorname{clgz\_})) \wedge ((\operatorname{imag\text{-}part}(\operatorname{clgz\_}) < \pi) \wedge (\exp \operatorname{clgz\_} = \operatorname{clgw\_})))))
\end{equation*}


\vnbbookfile{theorem-library/cc-exp-kernel}{theorem-library/cc-exp-kernel}
\noindent{\small cc-exp-kernel.scm -- THE KERNEL OF exp AND THE CONSTANT pi (the notes 2.3.3: Lemma 2.18, Prop 2.19, and what they rest on). Batch 31-B, 2026-09-25.  exp, cos, sin are theorem-library/cc-exp.scm's (31-A, definitions structure-library/cc-elementary.scm); PI is defined in}\par
\vnbprop{cek-limit-tail-le}{\vnbstmt{\begin{array}{@{}l@{}} \forall s \in (\mathbb{N} \to \mathbb{C}), v, c \in \mathbb{C}, d \in \mathbb{R}, m \in \mathbb{N}. \\ \quad (\operatorname{converges\text{-}to}(\operatorname{cc\text{-}ms}, s, v) \wedge \forall n \in \mathbb{N}.\; ((m \leq n) \Rightarrow (\operatorname{magnitude}((s(n) - c)) \leq d))) \Rightarrow  \\ \quad \quad \operatorname{magnitude}((v - c)) \leq d \end{array}}}
\vnbprop{cek-half-ratio-partial}{\vnbstmt{\begin{array}{@{}l@{}} \forall n \in \mathbb{N}, g \in (\mathbb{N} \to \mathbb{R}), m \in \mathbb{N}. \\ \quad \forall j \in \mathbb{N}.\; ((m \leq j) \Rightarrow ((2 \cdot g((j + 1))) \leq g(j))) \Rightarrow  \\ \quad \quad ((\operatorname{series\text{-}partial\text{-}sum}(g, (m + n)) - \operatorname{series\text{-}partial\text{-}sum}(g, m)) + (2 \cdot g((m + n)))) \leq (2 \cdot g(m)) \end{array}}}
\vnbprop{cc-series-ratio-tail}{\vnbstmt{\begin{array}{@{}l@{}} \forall f \in (\mathbb{N} \to \mathbb{C}), m \in \mathbb{N}. \\ \quad \operatorname{cc\text{-}series\text{-}converges}(f) \wedge \\ \quad \forall j \in \mathbb{N}.\; ((m \leq j) \Rightarrow ((2 \cdot \operatorname{magnitude}(f((j + 1)))) \leq \operatorname{magnitude}(f(j)))) \Rightarrow \\ \quad \quad \operatorname{magnitude}((\operatorname{cc\text{-}series\text{-}limit}(f) - \operatorname{cc\text{-}series\text{-}partial\text{-}sum}(f, m))) \leq (2 \cdot \operatorname{magnitude}(f(m))) \end{array}}}
\vnbprop{cc-exp-tail}{\vnbstmt{\begin{array}{@{}l@{}} \forall z \in \mathbb{C}, m \in \mathbb{N}. \\ \quad ((2 \cdot \operatorname{magnitude}(z)) \leq (m + 1)) \Rightarrow  \\ \quad \quad \operatorname{magnitude}((\exp z - \operatorname{cc\text{-}series\text{-}partial\text{-}sum}((\lambda k \in \mathbb{N}.\; (\operatorname{recip}(\operatorname{factorial}(k)) \cdot \operatorname{power}(z, k))), m))) \\ \quad \quad \leq\; 2 \cdot (\operatorname{recip}(\operatorname{factorial}(m)) \cdot \operatorname{power}(\operatorname{magnitude}(z), m)) \end{array}}}
\vnbprop{cc-exp-near-one}{\vnbstmt{\begin{array}{@{}l@{}} \forall z \in \mathbb{C}. \\ \quad (\operatorname{magnitude}(z) \leq 1) \Rightarrow (\operatorname{magnitude}(((\exp z - 1) - z)) \leq \operatorname{power}(\operatorname{magnitude}(z), 2)) \end{array}}}
\vnbprop{cc-exp-kernel-discrete}{\vnbstmt{\forall z \in \mathbb{C}.\; (((\exp z = 1) \wedge (\operatorname{magnitude}(z) < 1/2)) \Rightarrow (z = 0))}}
\vnbprop{cek-exp-gt-one-small}{\vnbstmt{\forall t \in \mathbb{R}.\; (((0 < t) \wedge (t \leq 1/2)) \Rightarrow (1 < \exp t))}}
\vnbprop{cek-exp-gt-one-n}{\vnbstmt{\forall n \in \mathbb{N}, t \in \mathbb{R}.\; (((0 < t) \wedge (t \leq (n \cdot 1/2))) \Rightarrow (1 < \exp t))}}
\vnbprop{cc-exp-real-gt-one}{\vnbstmt{\forall t \in \mathbb{R}.\; ((0 < t) \Rightarrow (1 < \exp t))}}
\vnbprop{cc-exp-real-one-iff}{\vnbstmt{\forall t \in \mathbb{R}.\; ((\exp t = 1) \Leftrightarrow (t = 0))}}
\vnbprop{cc-exp-kernel-add}{\vnbstmt{\begin{array}{@{}l@{}} \forall z, w \in \mathbb{C}. \\ \quad ((\exp z = 1) \wedge (\exp w = 1)) \Rightarrow (\exp (z + w) = 1) \end{array}}}
\vnbprop{cc-exp-kernel-neg}{\vnbstmt{\forall z \in \mathbb{C}.\; ((\exp z = 1) \Rightarrow (\exp -z = 1))}}
\vnbprop{cc-exp-kernel-sub}{\vnbstmt{\begin{array}{@{}l@{}} \forall z, w \in \mathbb{C}. \\ \quad ((\exp z = 1) \wedge (\exp w = 1)) \Rightarrow (\exp (z - w) = 1) \end{array}}}
\vnbprop{cc-exp-kernel-imaginary}{\vnbstmt{\forall z \in \mathbb{C}.\; ((\exp z = 1) \Rightarrow (\operatorname{real\text{-}part}(z) = 0))}}
\vnbprop{cc-exp-kernel-on-imaginary-axis}{\vnbstmt{\forall z \in \mathbb{C}.\; ((\exp z = 1) \Rightarrow (z = (+i \cdot \operatorname{imag\text{-}part}(z))))}}
\vnbprop{cc-cos-real-part}{\vnbstmt{\forall t \in \mathbb{R}.\; (\cos t = \operatorname{real\text{-}part}(\exp (+i \cdot t)))}}
\vnbprop{cc-sin-real-part}{\vnbstmt{\forall t \in \mathbb{R}.\; (\sin t = \operatorname{imag\text{-}part}(\exp (+i \cdot t)))}}
\vnbprop{cc-cos-sin-real}{\vnbstmt{\forall t \in \mathbb{R}.\; ((\cos t \in \mathbb{R}) \wedge (\sin t \in \mathbb{R}))}}
\vnbprop{cc-cos-one-bounds}{\vnbstmt{(5/12 \leq \cos 1) \wedge (\cos 1 \leq 7/12)}}
\vnbprop{cc-cos-two-neg}{\vnbstmt{\cos 2 < 0}}
\vnbprop{cc-continuous-restrict-rr}{\vnbstmt{\begin{array}{@{}l@{}} \forall f, g, x. \\ \quad \operatorname{is\text{-}continuous\text{-}at}(\operatorname{subspace\text{-}ms}(\operatorname{nf\text{-}metric\text{-}space}(\operatorname{cc\text{-}normed\text{-}field}), \mathbb{C}), \\ \quad \quad \operatorname{nf\text{-}metric\text{-}space}(\operatorname{cc\text{-}normed\text{-}field}), f, x) \wedge \\ \quad g \in (\mathbb{R} \to \mathbb{R}) \wedge \\ \quad x \in \mathbb{R} \wedge \\ \quad \forall v \in \mathbb{R}.\; (g(v) = f(v)) \Rightarrow \\ \quad \quad \operatorname{is\text{-}continuous\text{-}at}(\operatorname{rr\text{-}ms}, \operatorname{rr\text{-}ms}, g, x) \end{array}}}
\vnbprop{cc-neg-cos-continuous-rr}{\vnbstmt{\begin{array}{@{}l@{}} (\lambda v \in \mathbb{R}.\; (-1 \cdot \cos v)) \in (\mathbb{R} \to \mathbb{R}) \wedge \\ \forall x \in \mathbb{R}. \\ \quad \operatorname{is\text{-}continuous\text{-}at}(\operatorname{rr\text{-}ms}, \operatorname{rr\text{-}ms}, (\lambda v \in \mathbb{R}.\; (-1 \cdot \cos v)), x) \end{array}}}
\vnbprop{cc-cos-real-continuous}{\vnbstmt{\begin{array}{@{}l@{}} (\lambda v \in \mathbb{R}.\; \cos v) \in (\mathbb{R} \to \mathbb{R}) \wedge \\ \forall x \in \mathbb{R}. \\ \quad \operatorname{is\text{-}continuous\text{-}at}(\operatorname{rr\text{-}ms}, \operatorname{rr\text{-}ms}, (\lambda v \in \mathbb{R}.\; \cos v), x) \end{array}}}
\vnbprop{cc-cos-zero}{\vnbstmt{\cos 0 = 1}}
\vnbprop{cc-cos-has-zero}{\vnbstmt{\exists t \in \mathbb{R}.\; ((0 < t) \wedge ((t < 2) \wedge (\cos t = 0)))}}
\vnbprop{cc-exp-imaginary-period-exists}{\vnbstmt{\exists t \in \mathbb{R}.\; ((0 < t) \wedge ((t < 8) \wedge (\exp (+i \cdot t) = 1)))}}
\vnbprop{cc-exp-kernel-nontrivial}{\vnbstmt{\exists z \in \mathbb{C}.\; ((\exp z = 1) \wedge \neg (z = 0))}}
\vnbprop{cek-period-ge-half}{\vnbstmt{\forall t \in \mathbb{R}.\; (((0 < t) \wedge (\exp (+i \cdot t) = 1)) \Rightarrow (1/2 \leq t))}}
\vnbprop{cc-exp-least-period}{\vnbstmt{\begin{array}{@{}l@{}} \exists a \in \mathbb{R}. \\ \quad ((0 < a) \wedge ((\exp (+i \cdot a) = 1) \wedge \forall t \in \mathbb{R}.\; (((0 < t) \wedge (\exp (+i \cdot t) = 1)) \Rightarrow (a \leq t)))) \end{array}}}
\vnbprop{cek-kernel-nn-multiple}{\vnbstmt{\forall n \in \mathbb{N}, w \in \mathbb{C}.\; ((\exp w = 1) \Rightarrow (\exp (n \cdot w) = 1))}}
\vnbprop{cc-exp-kernel-zz-multiple}{\vnbstmt{\forall w \in \mathbb{C}, k \in \mathbb{Z}.\; ((\exp w = 1) \Rightarrow (\exp (k \cdot w) = 1))}}
\vnbprop{cek-period-multiple-nn}{\vnbstmt{\begin{array}{@{}l@{}} \forall n \in \mathbb{N}, a, r \in \mathbb{R}. \\ \quad 0 < a \wedge \\ \quad \exp (+i \cdot a) = 1 \wedge \\ \quad \forall t \in \mathbb{R}.\; (((0 < t) \wedge (\exp (+i \cdot t) = 1)) \Rightarrow (a \leq t)) \wedge \\ \quad 0 \leq r \wedge \\ \quad r < (n \cdot a) \wedge \\ \quad \exp (+i \cdot r) = 1 \Rightarrow \\ \quad \quad \exists k \in \mathbb{N}.\; (r = (k \cdot a)) \end{array}}}
\vnbprop{cek-period-multiple-zz}{\vnbstmt{\begin{array}{@{}l@{}} \forall a, r \in \mathbb{R}. \\ \quad 0 < a \wedge \\ \quad \exp (+i \cdot a) = 1 \wedge \\ \quad \forall t \in \mathbb{R}.\; (((0 < t) \wedge (\exp (+i \cdot t) = 1)) \Rightarrow (a \leq t)) \wedge \\ \quad \exp (+i \cdot r) = 1 \Rightarrow \\ \quad \quad \exists k \in \mathbb{Z}.\; (r = (k \cdot a)) \end{array}}}
\vnbprop{cc-pi-exists}{\vnbstmt{\begin{array}{@{}l@{}} \exists p \in \mathbb{R}. \\ \quad ((0 < p) \wedge \forall z \in \mathbb{C}.\; ((\exp z = 1) \Leftrightarrow \exists k \in \mathbb{Z}.\; (z = (k \cdot ((2 \cdot p) \cdot +i))))) \end{array}}}
\vnbprop{cek-pi-le}{\vnbstmt{\begin{array}{@{}l@{}} \forall p, q. \\ \quad p \in \mathbb{R} \wedge \\ \quad 0 < p \wedge \\ \quad \forall z \in \mathbb{C}.\; ((\exp z = 1) \Leftrightarrow \exists k \in \mathbb{Z}.\; (z = (k \cdot ((2 \cdot p) \cdot +i)))) \wedge \\ \quad q \in \mathbb{R} \wedge \\ \quad 0 < q \wedge \\ \quad \forall z \in \mathbb{C}.\; ((\exp z = 1) \Leftrightarrow \exists k \in \mathbb{Z}.\; (z = (k \cdot ((2 \cdot q) \cdot +i)))) \Rightarrow \\ \quad \quad q \leq p \end{array}}}
\vnbprop{cc-pi-unique}{\vnbstmt{\begin{array}{@{}l@{}} \forall p, q. \\ \quad p \in \mathbb{R} \wedge \\ \quad 0 < p \wedge \\ \quad \forall z \in \mathbb{C}.\; ((\exp z = 1) \Leftrightarrow \exists k \in \mathbb{Z}.\; (z = (k \cdot ((2 \cdot p) \cdot +i)))) \wedge \\ \quad q \in \mathbb{R} \wedge \\ \quad 0 < q \wedge \\ \quad \forall z \in \mathbb{C}.\; ((\exp z = 1) \Leftrightarrow \exists k \in \mathbb{Z}.\; (z = (k \cdot ((2 \cdot q) \cdot +i)))) \Rightarrow \\ \quad \quad p = q \end{array}}}
\vnbprop{pi-spec}{\vnbstmt{\begin{array}{@{}l@{}} \pi \in \mathbb{R} \wedge \\ 0 < \pi \wedge \\ \forall z \in \mathbb{C}.\; ((\exp z = 1) \Leftrightarrow \exists k \in \mathbb{Z}.\; (z = (k \cdot ((2 \cdot \pi) \cdot +i)))) \end{array}}}
\vnbprop{pi-pos}{\vnbstmt{(\pi \in \mathbb{R}) \wedge (0 < \pi)}}
\vnbprop{cc-exp-one-iff}{\vnbstmt{\forall z \in \mathbb{C}.\; ((\exp z = 1) \Leftrightarrow \exists k \in \mathbb{Z}.\; (z = (k \cdot ((2 \cdot \pi) \cdot +i))))}}
\vnbprop{cc-exp-two-pi-i}{\vnbstmt{\exp ((2 \cdot \pi) \cdot +i) = 1}}
\vnbprop{cc-cos-sin-two-pi}{\vnbstmt{(\cos (2 \cdot \pi) = 1) \wedge (\sin (2 \cdot \pi) = 0)}}
\vnbprop{cc-exp-i-pi}{\vnbstmt{\exp (+i \cdot \pi) = -1}}
\vnbprop{cc-cos-sin-pi}{\vnbstmt{(\cos \pi = -1) \wedge (\sin \pi = 0)}}
\vnbprop{cc-cos-sin-half-pi}{\vnbstmt{(\cos (1/2 \cdot \pi) = 0) \wedge (\operatorname{power}(\sin (1/2 \cdot \pi), 2) = 1)}}
\vnbprop{cc-exp-injective-strip}{\vnbstmt{\begin{array}{@{}l@{}} \forall z, w \in \mathbb{C}. \\ \quad -\pi \leq \operatorname{imag\text{-}part}(z) \wedge \\ \quad \operatorname{imag\text{-}part}(z) < \pi \wedge \\ \quad -\pi \leq \operatorname{imag\text{-}part}(w) \wedge \\ \quad \operatorname{imag\text{-}part}(w) < \pi \wedge \\ \quad \exp z = \exp w \Rightarrow \\ \quad \quad z = w \end{array}}}
\vnbprop{cc-exp-kernel-generator}{\vnbstmt{\begin{array}{@{}l@{}} \exists g \in \mathbb{C}. \\ \quad ((0 < \operatorname{imag\text{-}part}(g)) \wedge \forall z \in \mathbb{C}.\; ((\exp z = 1) \Leftrightarrow \exists k \in \mathbb{Z}.\; (z = (k \cdot g)))) \end{array}}}
\vnbprop{cc-exp-kernel-generator-unique}{\vnbstmt{\begin{array}{@{}l@{}} \forall g \in \mathbb{C}. \\ \quad ((0 < \operatorname{imag\text{-}part}(g)) \wedge \forall z \in \mathbb{C}.\; ((\exp z = 1) \Leftrightarrow \exists k \in \mathbb{Z}.\; (z = (k \cdot g)))) \Rightarrow  \\ \quad \quad g = ((2 \cdot \pi) \cdot +i) \end{array}}}
\vnbprop{cc-pi-lt-four}{\vnbstmt{\pi < 4}}
\vnbprop{cc-sin-half-pi}{\vnbstmt{\sin (1/2 \cdot \pi) = 1}}
\vnbprop{cc-cos-two-pi}{\vnbstmt{\cos (2 \cdot \pi) = 1}}
\vnbprop{cc-sin-two-pi}{\vnbstmt{\sin (2 \cdot \pi) = 0}}
\vnbprop{cc-cos-pi}{\vnbstmt{\cos \pi = -1}}
\vnbprop{cc-sin-pi}{\vnbstmt{\sin \pi = 0}}
\vnbprop{cc-cos-half-pi}{\vnbstmt{\cos (1/2 \cdot \pi) = 0}}

\vnbbookfile{theorem-library/cc-exp-surjective}{theorem-library/cc-exp-surjective}
\noindent{\small cc-exp-surjective.scm -- PROPOSITION 2.19 COMPLETED (exp is a bijection of each strip onto its image) AND THE BRIDGE cc-exp|RR = r-exp. Batch 34, 2026-09-26.  exp, cos, sin: theorem-library/cc-exp.scm (31-A); PI and the injectivity half of 2.19: theorem-library/cc-exp-kernel.scm (31-B).}\par
\vnbprop{cc-exp-real-continuous-rr}{\vnbstmt{\begin{array}{@{}l@{}} (\lambda v \in \mathbb{R}.\; \exp v) \in (\mathbb{R} \to \mathbb{R}) \wedge \\ \forall x \in \mathbb{R}. \\ \quad \operatorname{is\text{-}continuous\text{-}at}(\operatorname{rr\text{-}ms}, \operatorname{rr\text{-}ms}, (\lambda v \in \mathbb{R}.\; \exp v), x) \end{array}}}
\vnbprop{cc-diff-on-restrict-rr}{\vnbstmt{\begin{array}{@{}l@{}} \forall f, g, a, l. \\ \quad \operatorname{is\text{-}diff\text{-}on}(\operatorname{cc\text{-}normed\text{-}field}, \mathbb{C}, f, a, l) \wedge \\ \quad g \in (\mathbb{R} \to \mathbb{R}) \wedge \\ \quad a \in \mathbb{R} \wedge \\ \quad l \in \mathbb{R} \wedge \\ \quad \forall v \in \mathbb{R}.\; (g(v) = f(v)) \Rightarrow \\ \quad \quad \operatorname{is\text{-}diff\text{-}at}(g, a, l) \end{array}}}
\vnbprop{cc-exp-real-deriv}{\vnbstmt{\forall t \in \mathbb{R}.\; \operatorname{is\text{-}diff\text{-}at}((\lambda v \in \mathbb{R}.\; \exp v), t, \exp t)}}
\vnbprop{cxs-log-exp-deriv}{\vnbstmt{\forall x \in \mathbb{R}.\; \operatorname{is\text{-}diff\text{-}at}((\lambda h \in \mathbb{R}.\; (\log\left(\exp h\right) - h)), x, 0)}}
\vnbprop{cc-exp-real-is-r-exp}{\vnbstmt{\forall t \in \mathbb{R}.\; (\exp t = \exp\left(t\right))}}
\vnbprop{cc-exp-real-surjective-pos}{\vnbstmt{\forall r \in \mathbb{R}.\; ((0 < r) \Rightarrow \exists s \in \mathbb{R}.\; (\exp s = r))}}
\vnbprop{cc-circle-arg}{\vnbstmt{\begin{array}{@{}l@{}} \forall a, b \in \mathbb{R}. \\ \quad ((\operatorname{power}(a, 2) + \operatorname{power}(b, 2)) = 1) \Rightarrow  \\ \quad \quad \exists h \in \mathbb{R}.\; ((-\pi \leq h) \wedge ((h < \pi) \wedge ((\cos h = a) \wedge (\sin h = b)))) \end{array}}}
\vnbprop{cc-polar-form}{\vnbstmt{\begin{array}{@{}l@{}} \forall w \in \mathbb{C}. \\ \quad \neg (w = 0) \Rightarrow  \\ \quad \quad \exists r, h \in \mathbb{R}. \\ \quad \quad \quad (0 < r) \wedge ((-\pi \leq h) \wedge ((h < \pi) \wedge (w = (r \cdot \exp (+i \cdot h))))) \end{array}}}
\vnbprop{cc-exp-surjective-strip}{\vnbstmt{\begin{array}{@{}l@{}} \forall w \in \mathbb{C}. \\ \quad \neg (w = 0) \Rightarrow  \\ \quad \quad \exists z \in \mathbb{C}.\; ((-\pi \leq \operatorname{imag\text{-}part}(z)) \wedge ((\operatorname{imag\text{-}part}(z) < \pi) \wedge (\exp z = w))) \end{array}}}
\vnbprop{cc-exp-bijection-strip}{\vnbstmt{\begin{array}{@{}l@{}} (\lambda z \in \{q \in \mathbb{C} : ((-\pi \leq \operatorname{imag\text{-}part}(q)) \wedge (\operatorname{imag\text{-}part}(q) < \pi))\}.\; \exp z) \\ \in\; \mathrm{Bij}(\{q \in \mathbb{C} : ((-\pi \leq \operatorname{imag\text{-}part}(q)) \wedge (\operatorname{imag\text{-}part}(q) < \pi))\}, \operatorname{difference}(\mathbb{C}, \operatorname{singleton}(0))) \end{array}}}
\vnbprop{cc-exp-open-strip-avoids-negatives}{\vnbstmt{\begin{array}{@{}l@{}} \forall z \in \mathbb{C}. \\ \quad ((-\pi < \operatorname{imag\text{-}part}(z)) \wedge (\operatorname{imag\text{-}part}(z) < \pi)) \Rightarrow  \\ \quad \quad \neg ((\operatorname{imag\text{-}part}(\exp z) = 0) \wedge (\operatorname{real\text{-}part}(\exp z) \leq 0)) \end{array}}}
\vnbprop{cc-exp-open-strip-onto}{\vnbstmt{\begin{array}{@{}l@{}} \forall w \in \mathbb{C}. \\ \quad \neg ((\operatorname{imag\text{-}part}(w) = 0) \wedge (\operatorname{real\text{-}part}(w) \leq 0)) \Rightarrow  \\ \quad \quad \exists z \in \mathbb{C}.\; ((-\pi < \operatorname{imag\text{-}part}(z)) \wedge ((\operatorname{imag\text{-}part}(z) < \pi) \wedge (\exp z = w))) \end{array}}}
\vnbprop{cc-exp-bijection-open-strip}{\vnbstmt{\begin{array}{@{}l@{}} (\lambda z \in \{q \in \mathbb{C} : ((-\pi < \operatorname{imag\text{-}part}(q)) \wedge (\operatorname{imag\text{-}part}(q) < \pi))\}.\; \exp z) \\ \in\; \mathrm{Bij}(\{q \in \mathbb{C} : ((-\pi < \operatorname{imag\text{-}part}(q)) \wedge (\operatorname{imag\text{-}part}(q) < \pi))\}, \{u \in \mathbb{C} : \neg ((\operatorname{imag\text{-}part}(u) = 0) \wedge (\operatorname{real\text{-}part}(u) \leq 0))\}) \end{array}}}

\vnbbookfile{theorem-library/analytic-log-laws}{theorem-library/analytic-log-laws}
\noindent{\small analytic-log-laws.scm -- ANALYTICITY AT A CENTRE (Def 2.11, Cor 2.14 at a), THE PRINCIPAL LOGARITHM, LEMMA 2.20 AND PROPOSITION 2.21. Batch 38, 2026-09-27.  Definitions: structure-library/analytic-log.scm (IS-ANALYTIC-AT, CC-LOG).  The source: scratchpad/triage/ca.txt 1142-1150}\par
\vnbprop{alg-shift-in-ball}{\vnbstmt{\begin{array}{@{}l@{}} \forall a \in \mathbb{C}, r \in \mathbb{R}, z \in \operatorname{ball}(\operatorname{cc\text{-}ms}, a, r). \\ \quad (z \in \mathbb{C}) \wedge ((z - a) \in \operatorname{ball}(\operatorname{cc\text{-}ms}, 0, r)) \end{array}}}
\vnbprop{alg-shift-diff}{\vnbstmt{\begin{array}{@{}l@{}} \forall a \in \mathbb{C}, r \in \mathbb{R}, w \in \operatorname{ball}(\operatorname{cc\text{-}ms}, a, r). \\ \quad (0 < r) \Rightarrow  \\ \quad \quad \operatorname{is\text{-}diff\text{-}on}(\operatorname{cc\text{-}normed\text{-}field}, \operatorname{ball}(\operatorname{cc\text{-}ms}, a, r), (\lambda x \in \operatorname{ball}(\operatorname{cc\text{-}ms}, a, r).\; (x - a)), w, 1) \end{array}}}
\vnbprop{cps-translate-diff}{\vnbstmt{\begin{array}{@{}l@{}} \forall c, d \in (\mathbb{N} \to \mathbb{C}), a \in \mathbb{C}, r \in \mathbb{R}, w \in \operatorname{ball}(\operatorname{cc\text{-}ms}, a, r). \\ \quad (\forall j \in \mathbb{N}.\; (d(j) = ((j + 1) \cdot c((j + 1)))) \wedge ((0 < r) \wedge (r < \operatorname{cps\text{-}radius}(c)))) \Rightarrow  \\ \quad \quad \operatorname{is\text{-}diff\text{-}on}(\operatorname{cc\text{-}normed\text{-}field}, \operatorname{ball}(\operatorname{cc\text{-}ms}, a, r), \\ \quad \quad \quad (\lambda z \in \operatorname{ball}(\operatorname{cc\text{-}ms}, a, r).\; \operatorname{cc\text{-}series\text{-}limit}((\lambda k \in \mathbb{N}.\; (c(k) \cdot \operatorname{power}((z - a), k))))), \\ \quad \quad \quad w, \\ \quad \quad \quad \operatorname{cc\text{-}series\text{-}limit}((\lambda k \in \mathbb{N}.\; (d(k) \cdot \operatorname{power}((w - a), k))))) \end{array}}}
\vnbprop{cps-translate-holomorphic}{\vnbstmt{\begin{array}{@{}l@{}} \forall c \in (\mathbb{N} \to \mathbb{C}), a \in \mathbb{C}, r \in \mathbb{R}. \\ \quad ((0 < r) \wedge (r < \operatorname{cps\text{-}radius}(c))) \Rightarrow  \\ \quad \quad \operatorname{holomorphic\text{-}on}(\operatorname{ball}(\operatorname{cc\text{-}ms}, a, r), \\ \quad \quad \quad (\lambda z \in \operatorname{ball}(\operatorname{cc\text{-}ms}, a, r).\; \operatorname{cc\text{-}series\text{-}limit}((\lambda k \in \mathbb{N}.\; (c(k) \cdot \operatorname{power}((z - a), k)))))) \end{array}}}
\vnbprop{cps-translate-nth-deriv}{\vnbstmt{\begin{array}{@{}l@{}} \forall c \in (\mathbb{N} \to \mathbb{C}), a \in \mathbb{C}, r \in \mathbb{R}, k \in \mathbb{N}, g \in (\mathbb{N} \to \mathbb{C}). \\ \quad 0 < r \wedge \\ \quad r < \operatorname{cps\text{-}radius}(c) \wedge \\ \quad \forall n \in \mathbb{N}.\; (g(n) = ((\operatorname{factorial}((n + k)) \cdot \operatorname{recip}(\operatorname{factorial}(n))) \cdot c((n + k)))) \Rightarrow \\ \quad \quad (\lambda z \in \operatorname{ball}(\operatorname{cc\text{-}ms}, a, r).\; \operatorname{cc\text{-}series\text{-}limit}((\lambda \mathit{cpsk} \in \mathbb{N}.\; (c(\mathit{cpsk}) \cdot \operatorname{power}((z - a), \mathit{cpsk})))))^{(k)} \\ \quad \quad =\; (\lambda z \in \operatorname{ball}(\operatorname{cc\text{-}ms}, a, r).\; \operatorname{cc\text{-}series\text{-}limit}((\lambda \mathit{cpsk} \in \mathbb{N}.\; (g(\mathit{cpsk}) \cdot \operatorname{power}((z - a), \mathit{cpsk}))))) \end{array}}}
\vnbprop{taylor-coefficients-at}{\vnbstmt{\begin{array}{@{}l@{}} \forall f, u.; \forall c \in (\mathbb{N} \to \mathbb{C}), a \in \mathbb{C}, r \in \mathbb{R}, k \in \mathbb{N}. \\ \quad f \in (u \to \mathbb{C}) \wedge \\ \quad 0 < r \wedge \\ \quad r < \operatorname{cps\text{-}radius}(c) \wedge \\ \quad \operatorname{ball}(\operatorname{cc\text{-}ms}, a, r) \subseteq u \wedge \\ \quad \forall z \in \operatorname{ball}(\operatorname{cc\text{-}ms}, a, r). \\ \quad \quad f(z) = \operatorname{cc\text{-}series\text{-}limit}((\lambda \mathit{cpsk} \in \mathbb{N}.\; (c(\mathit{cpsk}) \cdot \operatorname{power}((z - a), \mathit{cpsk})))) \Rightarrow \\ \quad \quad \operatorname{restrict}(f, \operatorname{ball}(\operatorname{cc\text{-}ms}, a, r))^{(k)}(a) \\ \quad \quad =\; \operatorname{factorial}(k) \cdot c(k) \end{array}}}
\vnbprop{analytic-rep-holomorphic}{\vnbstmt{\begin{array}{@{}l@{}} \forall f, u.; \forall c \in (\mathbb{N} \to \mathbb{C}), a \in \mathbb{C}, r \in \mathbb{R}. \\ \quad f \in (u \to \mathbb{C}) \wedge \\ \quad 0 < r \wedge \\ \quad r < \operatorname{cps\text{-}radius}(c) \wedge \\ \quad \operatorname{ball}(\operatorname{cc\text{-}ms}, a, r) \subseteq u \wedge \\ \quad \forall z \in \operatorname{ball}(\operatorname{cc\text{-}ms}, a, r). \\ \quad \quad f(z) = \operatorname{cc\text{-}series\text{-}limit}((\lambda k \in \mathbb{N}.\; (c(k) \cdot \operatorname{power}((z - a), k)))) \Rightarrow \\ \quad \quad \operatorname{holomorphic\text{-}on}(\operatorname{ball}(\operatorname{cc\text{-}ms}, a, r), \operatorname{restrict}(f, \operatorname{ball}(\operatorname{cc\text{-}ms}, a, r))) \end{array}}}
\vnbprop{analytic-deriv}{\vnbstmt{\begin{array}{@{}l@{}} \forall f, u.; \forall c \in (\mathbb{N} \to \mathbb{C}), a \in \mathbb{C}, r \in \mathbb{R}. \\ \quad f \in (u \to \mathbb{C}) \wedge \\ \quad 0 < r \wedge \\ \quad r < \operatorname{cps\text{-}radius}(c) \wedge \\ \quad \operatorname{ball}(\operatorname{cc\text{-}ms}, a, r) \subseteq u \wedge \\ \quad \forall z \in \operatorname{ball}(\operatorname{cc\text{-}ms}, a, r). \\ \quad \quad f(z) = \operatorname{cc\text{-}series\text{-}limit}((\lambda k \in \mathbb{N}.\; (c(k) \cdot \operatorname{power}((z - a), k)))) \wedge \\ \quad \operatorname{is\text{-}open}(\operatorname{cc\text{-}ms}, u) \Rightarrow \\ \quad \quad \operatorname{is\text{-}diff\text{-}on}(\operatorname{cc\text{-}normed\text{-}field}, u, f, a, c(1)) \end{array}}}
\vnbprop{analytic-at-intro}{\vnbstmt{\begin{array}{@{}l@{}} \forall f, u.; \forall c \in (\mathbb{N} \to \mathbb{C}), a \in \mathbb{C}, r \in \mathbb{R}. \\ \quad f \in (u \to \mathbb{C}) \wedge \\ \quad 0 < r \wedge \\ \quad r < \operatorname{cps\text{-}radius}(c) \wedge \\ \quad \operatorname{ball}(\operatorname{cc\text{-}ms}, a, r) \subseteq u \wedge \\ \quad \forall z \in \operatorname{ball}(\operatorname{cc\text{-}ms}, a, r). \\ \quad \quad f(z) = \operatorname{cc\text{-}series\text{-}limit}((\lambda k \in \mathbb{N}.\; (c(k) \cdot \operatorname{power}((z - a), k)))) \wedge \\ \quad \operatorname{is\text{-}open}(\operatorname{cc\text{-}ms}, u) \Rightarrow \\ \quad \quad f \text{ is analytic at } a \end{array}}}
\vnbprop{analytic-holomorphic-ball}{\vnbstmt{\begin{array}{@{}l@{}} \forall f, u, a. \\ \quad f \text{ is analytic at } a \Rightarrow  \\ \quad \quad \exists r \in \mathbb{R}. \\ \quad \quad \quad 0 < r \wedge \\ \quad \quad \quad \operatorname{ball}(\operatorname{cc\text{-}ms}, a, r) \subseteq u \wedge \\ \quad \quad \quad \operatorname{holomorphic\text{-}on}(\operatorname{ball}(\operatorname{cc\text{-}ms}, a, r), \operatorname{restrict}(f, \operatorname{ball}(\operatorname{cc\text{-}ms}, a, r))) \end{array}}}
\vnbprop{cc-log-unique}{\vnbstmt{\begin{array}{@{}l@{}} \forall a, b \in \mathbb{C}. \\ \quad -\pi < \operatorname{imag\text{-}part}(a) \wedge \\ \quad \operatorname{imag\text{-}part}(a) < \pi \wedge \\ \quad -\pi < \operatorname{imag\text{-}part}(b) \wedge \\ \quad \operatorname{imag\text{-}part}(b) < \pi \wedge \\ \quad \exp a = \exp b \Rightarrow \\ \quad \quad a = b \end{array}}}
\vnbprop{cc-log-spec}{\vnbstmt{\begin{array}{@{}l@{}} \forall w \in \mathbb{C}. \\ \quad \neg ((\operatorname{imag\text{-}part}(w) = 0) \wedge (\operatorname{real\text{-}part}(w) \leq 0)) \Rightarrow  \\ \quad \quad \log\left(w\right) \in \mathbb{C} \wedge \\ \quad \quad -\pi < \operatorname{imag\text{-}part}(\log\left(w\right)) \wedge \\ \quad \quad \operatorname{imag\text{-}part}(\log\left(w\right)) < \pi \wedge \\ \quad \quad \exp \log\left(w\right) = w \end{array}}}
\vnbprop{cc-log-in-slit}{\vnbstmt{\begin{array}{@{}l@{}} \forall w \in \mathbb{C}. \\ \quad \neg ((\operatorname{imag\text{-}part}(w) = 0) \wedge (\operatorname{real\text{-}part}(w) \leq 0)) \Rightarrow  \\ \quad \quad (\log\left(w\right) \in \mathbb{C}) \wedge ((-\pi < \operatorname{imag\text{-}part}(\log\left(w\right))) \wedge (\operatorname{imag\text{-}part}(\log\left(w\right)) < \pi)) \end{array}}}
\vnbprop{cc-exp-log}{\vnbstmt{\begin{array}{@{}l@{}} \forall w \in \mathbb{C}. \\ \quad \neg ((\operatorname{imag\text{-}part}(w) = 0) \wedge (\operatorname{real\text{-}part}(w) \leq 0)) \Rightarrow (\exp \log\left(w\right) = w) \end{array}}}
\vnbprop{cc-log-char}{\vnbstmt{\begin{array}{@{}l@{}} \forall w, z \in \mathbb{C}. \\ \quad ((-\pi < \operatorname{imag\text{-}part}(z)) \wedge ((\operatorname{imag\text{-}part}(z) < \pi) \wedge (\exp z = w))) \Rightarrow (\log\left(w\right) = z) \end{array}}}
\vnbprop{cc-log-exp}{\vnbstmt{\begin{array}{@{}l@{}} \forall z \in \mathbb{C}. \\ \quad ((-\pi < \operatorname{imag\text{-}part}(z)) \wedge (\operatorname{imag\text{-}part}(z) < \pi)) \Rightarrow (\log\left(\exp z\right) = z) \end{array}}}
\vnbprop{cc-log-real}{\vnbstmt{\forall t \in \mathbb{R}.\; ((0 < t) \Rightarrow (\log\left(t\right) = \log\left(t\right)))}}
\vnbprop{cc-log-one}{\vnbstmt{\log\left(1\right) = 0}}
\vnbprop{cc-log-polar}{\vnbstmt{\begin{array}{@{}l@{}} \forall r, h \in \mathbb{R}. \\ \quad ((0 < r) \wedge ((-\pi < h) \wedge (h < \pi))) \Rightarrow  \\ \quad \quad \log\left((r \cdot \exp (+i \cdot h))\right) = (\log\left(r\right) + (+i \cdot h)) \end{array}}}
\vnbprop{cc-log-scale}{\vnbstmt{\begin{array}{@{}l@{}} \forall s \in \mathbb{R}, w \in \mathbb{C}. \\ \quad ((0 < s) \wedge \neg ((\operatorname{imag\text{-}part}(w) = 0) \wedge (\operatorname{real\text{-}part}(w) \leq 0))) \Rightarrow  \\ \quad \quad \log\left((s \cdot w)\right) = (\log\left(s\right) + \log\left(w\right)) \end{array}}}
\vnbprop{cc-geometric-alternating}{\vnbstmt{\begin{array}{@{}l@{}} \forall g \in (\mathbb{N} \to \mathbb{C}), z \in \mathbb{C}. \\ \quad (\forall j \in \mathbb{N}.\; (g(j) = \operatorname{power}(-1, j)) \wedge (\operatorname{magnitude}(z) < 1)) \Rightarrow  \\ \quad \quad \operatorname{cc\text{-}series\text{-}converges}((\lambda k \in \mathbb{N}.\; (g(k) \cdot \operatorname{power}(z, k)))) \wedge \\ \quad \quad \operatorname{cc\text{-}series\text{-}limit}((\lambda k \in \mathbb{N}.\; (g(k) \cdot \operatorname{power}(z, k)))) = \operatorname{recip}((1 + z)) \end{array}}}
\vnbprop{cps-log-radius}{\vnbstmt{\begin{array}{@{}l@{}} \forall c \in (\mathbb{N} \to \mathbb{C}), r \in \mathbb{R}. \\ \quad ((c(0) = 0) \wedge (\forall j \in \mathbb{N}.\; (c((j + 1)) = (\operatorname{power}(-1, j) \cdot \operatorname{recip}((j + 1)))) \wedge ((0 \leq r) \wedge (r < 1)))) \Rightarrow  \\ \quad \quad r < \operatorname{cps\text{-}radius}(c) \end{array}}}
\vnbprop{cps-log-series}{\vnbstmt{\begin{array}{@{}l@{}} \forall c \in (\mathbb{N} \to \mathbb{C}), r \in \mathbb{R}, z \in \operatorname{ball}(\operatorname{cc\text{-}ms}, 0, r). \\ \quad ((c(0) = 0) \wedge (\forall j \in \mathbb{N}.\; (c((j + 1)) = (\operatorname{power}(-1, j) \cdot \operatorname{recip}((j + 1)))) \wedge ((0 < r) \wedge (r < 1)))) \Rightarrow  \\ \quad \quad \operatorname{is\text{-}diff\text{-}on}(\operatorname{cc\text{-}normed\text{-}field}, \operatorname{ball}(\operatorname{cc\text{-}ms}, 0, r), \\ \quad \quad \quad (\lambda \mathit{cpdz} \in \operatorname{ball}(\operatorname{cc\text{-}ms}, 0, r).\; \operatorname{cc\text{-}series\text{-}limit}((\lambda k \in \mathbb{N}.\; (c(k) \cdot \operatorname{power}(\mathit{cpdz}, k))))), \\ \quad \quad \quad z, \operatorname{recip}((1 + z))) \end{array}}}
\vnbprop{alg-zero-deriv-ball}{\vnbstmt{\begin{array}{@{}l@{}} \forall a \in \mathbb{C}, r \in \mathbb{R}, h \in (\operatorname{ball}(\operatorname{cc\text{-}ms}, a, r) \to \mathbb{C}), z \in \operatorname{ball}(\operatorname{cc\text{-}ms}, a, r). \\ \quad 0 < r \wedge \\ \quad \forall \mathit{hbz} \in \operatorname{ball}(\operatorname{cc\text{-}ms}, a, r). \\ \quad \quad \operatorname{is\text{-}diff\text{-}on}(\operatorname{cc\text{-}normed\text{-}field}, \operatorname{ball}(\operatorname{cc\text{-}ms}, a, r), h, \mathit{hbz}, 0) \Rightarrow \\ \quad \quad h(z) = h(a) \end{array}}}
\vnbprop{cps-log-deriv-at-one}{\vnbstmt{\begin{array}{@{}l@{}} \forall c \in (\mathbb{N} \to \mathbb{C}), r \in \mathbb{R}, w \in \operatorname{ball}(\operatorname{cc\text{-}ms}, 1, r). \\ \quad ((c(0) = 0) \wedge (\forall j \in \mathbb{N}.\; (c((j + 1)) = (\operatorname{power}(-1, j) \cdot \operatorname{recip}((j + 1)))) \wedge ((0 < r) \wedge (r < 1)))) \Rightarrow  \\ \quad \quad \operatorname{is\text{-}diff\text{-}on}(\operatorname{cc\text{-}normed\text{-}field}, \operatorname{ball}(\operatorname{cc\text{-}ms}, 1, r), \\ \quad \quad \quad (\lambda z \in \operatorname{ball}(\operatorname{cc\text{-}ms}, 1, r).\; \operatorname{cc\text{-}series\text{-}limit}((\lambda k \in \mathbb{N}.\; (c(k) \cdot \operatorname{power}((z - 1), k))))), \\ \quad \quad \quad w, \operatorname{recip}(w)) \end{array}}}
\vnbprop{log0-exp}{\vnbstmt{\begin{array}{@{}l@{}} \forall c \in (\mathbb{N} \to \mathbb{C}), r \in \mathbb{R}, z \in \operatorname{ball}(\operatorname{cc\text{-}ms}, 1, r). \\ \quad ((c(0) = 0) \wedge (\forall j \in \mathbb{N}.\; (c((j + 1)) = (\operatorname{power}(-1, j) \cdot \operatorname{recip}((j + 1)))) \wedge ((0 < r) \wedge (r < 1)))) \Rightarrow  \\ \quad \quad \exp \operatorname{cc\text{-}series\text{-}limit}((\lambda k \in \mathbb{N}.\; (c(k) \cdot \operatorname{power}((z - 1), k)))) = z \end{array}}}
\vnbprop{log0-bound}{\vnbstmt{\begin{array}{@{}l@{}} \forall c \in (\mathbb{N} \to \mathbb{C}), r \in \mathbb{R}, z \in \operatorname{ball}(\operatorname{cc\text{-}ms}, 1, r). \\ \quad ((c(0) = 0) \wedge (\forall j \in \mathbb{N}.\; (c((j + 1)) = (\operatorname{power}(-1, j) \cdot \operatorname{recip}((j + 1)))) \wedge ((0 < r) \wedge (r < 1)))) \Rightarrow  \\ \quad \quad \operatorname{magnitude}(\operatorname{cc\text{-}series\text{-}limit}((\lambda k \in \mathbb{N}.\; (c(k) \cdot \operatorname{power}((z - 1), k))))) \leq (\operatorname{recip}((1 - \operatorname{magnitude}((z - 1)))) - 1) \end{array}}}
\vnbprop{log0-is-log-on-ball}{\vnbstmt{\begin{array}{@{}l@{}} \forall c \in (\mathbb{N} \to \mathbb{C}), r \in \mathbb{R}, z \in \operatorname{ball}(\operatorname{cc\text{-}ms}, 1, r). \\ \quad c(0) = 0 \wedge \\ \quad \forall j \in \mathbb{N}.\; (c((j + 1)) = (\operatorname{power}(-1, j) \cdot \operatorname{recip}((j + 1)))) \wedge \\ \quad 0 < r \wedge \\ \quad (r \cdot (1 + \pi)) \leq \pi \Rightarrow \\ \quad \quad \log\left(z\right) = \operatorname{cc\text{-}series\text{-}limit}((\lambda k \in \mathbb{N}.\; (c(k) \cdot \operatorname{power}((z - 1), k)))) \end{array}}}
\vnbprop{cc-log-holomorphic-near-one}{\vnbstmt{\begin{array}{@{}l@{}} \forall r \in \mathbb{R}. \\ \quad ((0 < r) \wedge ((r \cdot (1 + \pi)) \leq \pi)) \Rightarrow  \\ \quad \quad \operatorname{holomorphic\text{-}on}(\operatorname{ball}(\operatorname{cc\text{-}ms}, 1, r), (\lambda z \in \operatorname{ball}(\operatorname{cc\text{-}ms}, 1, r).\; \log\left(z\right))) \end{array}}}
\vnbprop{cc-log-one-plus}{\vnbstmt{\begin{array}{@{}l@{}} \forall c \in (\mathbb{N} \to \mathbb{C}), r \in \mathbb{R}, z \in \mathbb{C}. \\ \quad c(0) = 0 \wedge \\ \quad \forall j \in \mathbb{N}.\; (c((j + 1)) = (\operatorname{power}(-1, j) \cdot \operatorname{recip}((j + 1)))) \wedge \\ \quad 0 < r \wedge \\ \quad (r \cdot (1 + \pi)) \leq \pi \wedge \\ \quad \operatorname{magnitude}(z) < r \Rightarrow \\ \quad \quad \log\left((1 + z)\right) = \operatorname{cc\text{-}series\text{-}limit}((\lambda k \in \mathbb{N}.\; (c(k) \cdot \operatorname{power}(z, k)))) \end{array}}}
\vnbprop{holomorphic-compose}{\vnbstmt{\begin{array}{@{}l@{}} \forall u, v, f, g. \\ \quad (\operatorname{holomorphic\text{-}on}(u, f) \wedge (\operatorname{holomorphic\text{-}on}(v, g) \wedge \forall y \in u.\; (f(y) \in v))) \Rightarrow  \\ \quad \quad \operatorname{holomorphic\text{-}on}(u, (\lambda x \in u.\; g(f(x)))) \end{array}}}
\vnbprop{cc-log-holomorphic-scale}{\vnbstmt{\begin{array}{@{}l@{}} \forall r, s \in \mathbb{R}. \\ \quad ((0 < r) \wedge (((r \cdot (1 + \pi)) \leq \pi) \wedge (0 < s))) \Rightarrow  \\ \quad \quad \operatorname{holomorphic\text{-}on}(\operatorname{ball}(\operatorname{cc\text{-}ms}, s, (s \cdot r)), (\lambda z \in \operatorname{ball}(\operatorname{cc\text{-}ms}, s, (s \cdot r)).\; \log\left(z\right))) \end{array}}}
\vnbprop{log0-real}{\vnbstmt{\begin{array}{@{}l@{}} \forall c \in (\mathbb{N} \to \mathbb{C}), r \in \mathbb{R}, z \in \operatorname{ball}(\operatorname{cc\text{-}ms}, 1, r). \\ \quad c(0) = 0 \wedge \\ \quad \forall j \in \mathbb{N}.\; (c((j + 1)) = (\operatorname{power}(-1, j) \cdot \operatorname{recip}((j + 1)))) \wedge \\ \quad 0 < r \wedge \\ \quad r < 1 \wedge \\ \quad z \in \mathbb{R} \Rightarrow \\ \quad \quad \operatorname{cc\text{-}series\text{-}limit}((\lambda k \in \mathbb{N}.\; (c(k) \cdot \operatorname{power}((z - 1), k)))) \in \mathbb{C} \wedge \\ \quad \quad \operatorname{imag\text{-}part}(\operatorname{cc\text{-}series\text{-}limit}((\lambda k \in \mathbb{N}.\; (c(k) \cdot \operatorname{power}((z - 1), k))))) = 0 \end{array}}}
\vnbprop{log0-is-ln}{\vnbstmt{\begin{array}{@{}l@{}} \forall c \in (\mathbb{N} \to \mathbb{C}), t \in \mathbb{R}. \\ \quad ((c(0) = 0) \wedge (\forall j \in \mathbb{N}.\; (c((j + 1)) = (\operatorname{power}(-1, j) \cdot \operatorname{recip}((j + 1)))) \wedge ((0 < t) \wedge (t < 2)))) \Rightarrow  \\ \quad \quad \operatorname{cc\text{-}series\text{-}limit}((\lambda k \in \mathbb{N}.\; (c(k) \cdot \operatorname{power}((t - 1), k)))) = \log\left(t\right) \end{array}}}

\vnbbookfile{theorem-library/cc-log-holomorphic}{theorem-library/cc-log-holomorphic}
\noindent{\small cc-log-holomorphic.scm -- THE SIGN OF cos ON [-pi/2, pi], pi $>$ 3, AND THE HOLOMORPHY OF THE PRINCIPAL LOGARITHM ON THE SLIT PLANE C \textbackslash{} (-inf, 0]; the two loose ends of Lemma 2.20 (the radius is exactly 1) and Remark 2.22 (log(1 - z) = - sum z\textasciicircum{}k/k).}\par
\vnbprop{cc-cos-least-zero}{\vnbstmt{\forall t \in \mathbb{R}.\; (((0 < t) \wedge (t < (1/2 \cdot \pi))) \Rightarrow \neg (\cos t = 0))}}
\vnbprop{cc-cos-pos-right}{\vnbstmt{\forall s \in \mathbb{R}.\; (((0 < s) \wedge (s < (1/2 \cdot \pi))) \Rightarrow (0 < \cos s))}}
\vnbprop{cc-cos-pos-half}{\vnbstmt{\begin{array}{@{}l@{}} \forall t \in \mathbb{R}. \\ \quad ((-(1/2 \cdot \pi) < t) \wedge (t < (1/2 \cdot \pi))) \Rightarrow (0 < \cos t) \end{array}}}
\vnbprop{cc-cos-nonpos-far}{\vnbstmt{\forall t \in \mathbb{R}.\; ((((1/2 \cdot \pi) \leq t) \wedge (t \leq \pi)) \Rightarrow (\cos t \leq 0))}}
\vnbprop{cc-cos-pos-small}{\vnbstmt{\forall t \in \mathbb{R}.\; (((0 < t) \wedge (t < 1)) \Rightarrow (0 < \cos t))}}
\vnbprop{cc-pi-gt-two}{\vnbstmt{2 < \pi}}
\vnbprop{cc-cos-three-halves-pos}{\vnbstmt{0 < \cos 3/2}}
\vnbprop{cc-pi-gt-three}{\vnbstmt{3 < \pi}}
\vnbprop{cc-log-right-half}{\vnbstmt{\begin{array}{@{}l@{}} \forall w \in \mathbb{C}. \\ \quad (0 < \operatorname{real\text{-}part}(w)) \Rightarrow  \\ \quad \quad (-(1/2 \cdot \pi) < \operatorname{imag\text{-}part}(\log\left(w\right))) \wedge (\operatorname{imag\text{-}part}(\log\left(w\right)) < (1/2 \cdot \pi)) \end{array}}}
\vnbprop{cc-log-mul}{\vnbstmt{\begin{array}{@{}l@{}} \forall u, v \in \mathbb{C}. \\ \quad \neg ((\operatorname{imag\text{-}part}(u) = 0) \wedge (\operatorname{real\text{-}part}(u) \leq 0)) \wedge \\ \quad \neg ((\operatorname{imag\text{-}part}(v) = 0) \wedge (\operatorname{real\text{-}part}(v) \leq 0)) \wedge \\ \quad -\pi < (\operatorname{imag\text{-}part}(\log\left(u\right)) + \operatorname{imag\text{-}part}(\log\left(v\right))) \wedge \\ \quad (\operatorname{imag\text{-}part}(\log\left(u\right)) + \operatorname{imag\text{-}part}(\log\left(v\right))) < \pi \Rightarrow \\ \quad \quad \log\left((u \cdot v)\right) = (\log\left(u\right) + \log\left(v\right)) \end{array}}}
\vnbprop{cc-log-rotate}{\vnbstmt{\begin{array}{@{}l@{}} \forall w \in \mathbb{C}. \\ \quad (0 < \operatorname{real\text{-}part}(w)) \Rightarrow (\log\left((+i \cdot w)\right) = (\log\left(w\right) + (+i \cdot (1/2 \cdot \pi)))) \end{array}}}
\vnbprop{cc-log-rotate-neg}{\vnbstmt{\begin{array}{@{}l@{}} \forall w \in \mathbb{C}. \\ \quad (0 < \operatorname{real\text{-}part}(w)) \Rightarrow (\log\left((-+i \cdot w)\right) = (\log\left(w\right) - (+i \cdot (1/2 \cdot \pi)))) \end{array}}}
\vnbprop{cc-log-im-near-one}{\vnbstmt{\begin{array}{@{}l@{}} \forall d, r \in \mathbb{R}, z \in \operatorname{ball}(\operatorname{cc\text{-}ms}, 1, r). \\ \quad ((0 < d) \wedge ((0 < r) \wedge (((r \cdot (1 + \pi)) \leq \pi) \wedge ((r \cdot (1 + d)) \leq d)))) \Rightarrow  \\ \quad \quad \neg ((\operatorname{imag\text{-}part}(z) = 0) \wedge (\operatorname{real\text{-}part}(z) \leq 0)) \wedge \\ \quad \quad -d < \operatorname{imag\text{-}part}(\log\left(z\right)) \wedge \\ \quad \quad \operatorname{imag\text{-}part}(\log\left(z\right)) < d \end{array}}}
\vnbprop{cc-log-holomorphic-at}{\vnbstmt{\begin{array}{@{}l@{}} \forall a \in \mathbb{C}. \\ \quad \neg ((\operatorname{imag\text{-}part}(a) = 0) \wedge (\operatorname{real\text{-}part}(a) \leq 0)) \Rightarrow  \\ \quad \quad \exists p \in \mathbb{R}. \\ \quad \quad \quad 0 < p \wedge \\ \quad \quad \quad \operatorname{ball}(\operatorname{cc\text{-}ms}, a, p) \subseteq \{u \in \mathbb{C} : \neg ((\operatorname{imag\text{-}part}(u) = 0) \wedge (\operatorname{real\text{-}part}(u) \leq 0))\} \wedge \\ \quad \quad \quad \operatorname{holomorphic\text{-}on}(\operatorname{ball}(\operatorname{cc\text{-}ms}, a, p), (\lambda x \in \operatorname{ball}(\operatorname{cc\text{-}ms}, a, p).\; \log\left(x\right))) \end{array}}}
\vnbprop{cc-log-holomorphic}{\vnbstmt{\begin{array}{@{}l@{}} \operatorname{holomorphic\text{-}on}(\{u \in \mathbb{C} : \neg ((\operatorname{imag\text{-}part}(u) = 0) \wedge (\operatorname{real\text{-}part}(u) \leq 0))\}, \\ \quad (\lambda w \in \{u \in \mathbb{C} : \neg ((\operatorname{imag\text{-}part}(u) = 0) \wedge (\operatorname{real\text{-}part}(u) \leq 0))\}.\; \log\left(w\right))) \end{array}}}
\vnbprop{cc-log-deriv}{\vnbstmt{\begin{array}{@{}l@{}} \forall b \in \mathbb{C}. \\ \quad \neg ((\operatorname{imag\text{-}part}(b) = 0) \wedge (\operatorname{real\text{-}part}(b) \leq 0)) \Rightarrow  \\ \quad \quad \operatorname{is\text{-}diff\text{-}on}(\operatorname{cc\text{-}normed\text{-}field}, \\ \quad \quad \quad \{u \in \mathbb{C} : \neg ((\operatorname{imag\text{-}part}(u) = 0) \wedge (\operatorname{real\text{-}part}(u) \leq 0))\}, \\ \quad \quad \quad (\lambda w \in \{u \in \mathbb{C} : \neg ((\operatorname{imag\text{-}part}(u) = 0) \wedge (\operatorname{real\text{-}part}(u) \leq 0))\}.\; \log\left(w\right)), \\ \quad \quad \quad b, \operatorname{recip}(b)) \end{array}}}
\vnbprop{cc-log-one-minus}{\vnbstmt{\begin{array}{@{}l@{}} \forall r \in \mathbb{R}, z \in \mathbb{C}. \\ \quad ((0 < r) \wedge (((r \cdot (1 + \pi)) \leq \pi) \wedge (\operatorname{magnitude}(z) < r))) \Rightarrow  \\ \quad \quad \log\left((1 - z)\right) = -\operatorname{cc\text{-}series\text{-}limit}((\lambda k \in \mathbb{N}.\; ({k}^{\ast\!-1} \cdot \operatorname{power}(z, k)))) \end{array}}}
\vnbprop{cc-series-terms-small}{\vnbstmt{\begin{array}{@{}l@{}} \forall f \in (\mathbb{N} \to \mathbb{C}), e. \\ \quad (\operatorname{cc\text{-}series\text{-}converges}(f) \wedge \operatorname{pos\text{-}rr}(e)) \Rightarrow  \\ \quad \quad \exists n \in \mathbb{N}.\; \forall m \in \mathbb{N}.\; ((n \leq m) \Rightarrow (\operatorname{magnitude}(f(m)) \leq e)) \end{array}}}
\vnbprop{cps-log-radius-le-one}{\vnbstmt{\begin{array}{@{}l@{}} \forall c \in (\mathbb{N} \to \mathbb{C}), r \in \mathbb{R}. \\ \quad ((c(0) = 0) \wedge (\forall j \in \mathbb{N}.\; (c((j + 1)) = (\operatorname{power}(-1, j) \cdot \operatorname{recip}((j + 1)))) \wedge (1 < r))) \Rightarrow  \\ \quad \quad \neg (r < \operatorname{cps\text{-}radius}(c)) \end{array}}}
\vnbprop{cc-sin-pos-half}{\vnbstmt{\forall t \in \mathbb{R}.\; (((0 < t) \wedge (t < \pi)) \Rightarrow (0 < \sin t))}}
\vnbprop{cps-log-radius-one}{\vnbstmt{\begin{array}{@{}l@{}} \forall c \in (\mathbb{N} \to \mathbb{C}). \\ \quad ((c(0) = 0) \wedge \forall j \in \mathbb{N}.\; (c((j + 1)) = (\operatorname{power}(-1, j) \cdot \operatorname{recip}((j + 1))))) \Rightarrow  \\ \quad \quad \operatorname{cps\text{-}radius}(c) = 1 \end{array}}}

\vnbbookfile{structure-library/winding-number}{structure-library/winding-number}
\noindent{\small winding-number.scm -- THE WINDING NUMBER AND THE CIRCLE AS A ROAD (complex-analysis.pdf ch. 3: Definition 3.26, equation (73); Example 3.13). DEFINITIONS ONLY; every law is PROVEN in theorem-library/winding-number-laws.scm. Agent CA-1 of the October roadmap, 2026-10-03.}\par
\par\medskip\noindent\textbf{Definition} (\texttt{winding}).\;$\operatorname{winding}(\operatorname{wnp\_}, \operatorname{wnd\_}, \operatorname{wna\_}, \operatorname{wnb\_}, \operatorname{wnz\_}) :=$
\begin{equation*}
\operatorname{recip}(((2 \cdot \pi) \cdot +i)) \cdot \operatorname{line\text{-}int}((\lambda \operatorname{wnw\_} \in \operatorname{difference}(\mathbb{C}, \operatorname{singleton}(\operatorname{wnz\_})).\; \operatorname{recip}((\operatorname{wnw\_} - \operatorname{wnz\_}))), \operatorname{wnp\_}, \operatorname{wnd\_}, \operatorname{wna\_}, \operatorname{wnb\_})
\end{equation*}

\par\medskip\noindent\textbf{Definition} (\texttt{circle-path}).\;$\operatorname{circle\text{-}path}(\operatorname{wnc\_}, \operatorname{wnr\_}) :=$
\begin{equation*}
(\lambda \operatorname{wnt\_} \in \operatorname{ccint}(0, (2 \cdot \pi)).\; (\operatorname{wnc\_} + (\operatorname{wnr\_} \cdot \exp (+i \cdot \operatorname{wnt\_}))))
\end{equation*}

\par\medskip\noindent\textbf{Definition} (\texttt{circle-deriv}).\;$\operatorname{circle\text{-}deriv}(\operatorname{wnc\_}, \operatorname{wnr\_}) :=$
\begin{equation*}
(\lambda \operatorname{wnt\_} \in \operatorname{ccint}(0, (2 \cdot \pi)).\; (+i \cdot (\operatorname{wnr\_} \cdot \exp (+i \cdot \operatorname{wnt\_}))))
\end{equation*}


\vnbbookfile{theorem-library/winding-number-laws}{theorem-library/winding-number-laws}
\noindent{\small winding-number-laws.scm -- CAUCHY ON A DISC, THE WINDING NUMBER, THE CIRCLE AS A ROAD (complex-analysis.pdf ch. 3: Corollary 3.9 for a disc inside an open set; Definition 3.26 and the disc case of Corollary 3.30; Example 3.13 at the centre). Agent CA-1 of the October roadmap (docs/roadmap-2026-10-03.md, line (a), ...}\par
\vnbprop{cc-ball-is-convex}{\vnbstmt{\forall c \in \mathbb{C}, r \in \mathbb{R}.\; \operatorname{is\text{-}convex\text{-}cc}(\operatorname{ball}(\operatorname{cc\text{-}ms}, c, r))}}
\vnbprop{line-int-domain-congruence}{\vnbstmt{\begin{array}{@{}l@{}} \forall u, f, m, c, a, b, v.; \forall g \in (v \to \mathbb{C}). \\ \quad u \subseteq \mathbb{C} \wedge \\ \quad f \in (u \to \mathbb{C}) \wedge \\ \quad \forall y \in u. \\ \quad \quad \operatorname{is\text{-}continuous\text{-}at}(\operatorname{subspace\text{-}ms}(\operatorname{nf\text{-}metric\text{-}space}(\operatorname{cc\text{-}normed\text{-}field}), u), \\ \quad \quad \quad \operatorname{nf\text{-}metric\text{-}space}(\operatorname{cc\text{-}normed\text{-}field}), f, y) \wedge \\ \quad \operatorname{is\text{-}road}(m, c, a, b) \wedge \\ \quad \operatorname{trace}(m, a, b) \subseteq u \wedge \\ \quad \operatorname{trace}(m, a, b) \subseteq v \wedge \\ \quad \forall x \in \operatorname{trace}(m, a, b).\; (g(x) \simeq f(x)) \Rightarrow \\ \quad \quad \operatorname{line\text{-}int}(g, m, c, a, b) = \operatorname{line\text{-}int}(f, m, c, a, b) \end{array}}}
\vnbprop{holomorphic-on-restrict}{\vnbstmt{\begin{array}{@{}l@{}} \forall u, f, v. \\ \quad (\operatorname{holomorphic\text{-}on}(u, f) \wedge (\operatorname{is\text{-}open}(\operatorname{nf\text{-}metric\text{-}space}(\operatorname{cc\text{-}normed\text{-}field}), v) \wedge (v \subseteq u))) \Rightarrow  \\ \quad \quad \operatorname{holomorphic\text{-}on}(v, \operatorname{restrict}(f, v)) \end{array}}}
\vnbprop{cauchy-convex-sub}{\vnbstmt{\begin{array}{@{}l@{}} \forall u, f, v, m, c, a, b. \\ \quad \operatorname{holomorphic\text{-}on}(u, f) \wedge \\ \quad \operatorname{is\text{-}open}(\operatorname{nf\text{-}metric\text{-}space}(\operatorname{cc\text{-}normed\text{-}field}), v) \wedge \\ \quad \operatorname{is\text{-}convex\text{-}cc}(v) \wedge \\ \quad v \subseteq u \wedge \\ \quad \operatorname{is\text{-}road}(m, c, a, b) \wedge \\ \quad \operatorname{trace}(m, a, b) \subseteq v \wedge \\ \quad m(b) = m(a) \Rightarrow \\ \quad \quad \operatorname{line\text{-}int}(f, m, c, a, b) = 0 \end{array}}}
\vnbprop{cauchy-disc}{\vnbstmt{\begin{array}{@{}l@{}} \forall u, f.; \forall c \in \mathbb{C}, r \in \mathbb{R}, m, d, a, b. \\ \quad \operatorname{holomorphic\text{-}on}(u, f) \wedge \\ \quad 0 < r \wedge \\ \quad \operatorname{ball}(\operatorname{cc\text{-}ms}, c, r) \subseteq u \wedge \\ \quad \operatorname{is\text{-}road}(m, d, a, b) \wedge \\ \quad \operatorname{trace}(m, a, b) \subseteq \operatorname{ball}(\operatorname{cc\text{-}ms}, c, r) \wedge \\ \quad m(b) = m(a) \Rightarrow \\ \quad \quad \operatorname{line\text{-}int}(f, m, d, a, b) = 0 \end{array}}}
\vnbprop{cc-punctured-membership}{\vnbstmt{\begin{array}{@{}l@{}} \forall z \in \mathbb{C}, y. \\ \quad (y \in \operatorname{difference}(\mathbb{C}, \operatorname{singleton}(z))) \Leftrightarrow ((y \in \mathbb{C}) \wedge \neg (y = z)) \end{array}}}
\vnbprop{cc-punctured-open}{\vnbstmt{\forall z \in \mathbb{C}.\; \operatorname{is\text{-}open}(\operatorname{cc\text{-}ms}, \operatorname{difference}(\mathbb{C}, \operatorname{singleton}(z)))}}
\vnbprop{winding-integrand-holomorphic}{\vnbstmt{\begin{array}{@{}l@{}} \forall z \in \mathbb{C}. \\ \quad \operatorname{holomorphic\text{-}on}(\operatorname{difference}(\mathbb{C}, \operatorname{singleton}(z)), \\ \quad \quad (\lambda w \in \operatorname{difference}(\mathbb{C}, \operatorname{singleton}(z)).\; \operatorname{recip}((w - z)))) \end{array}}}
\vnbprop{winding-in-cc}{\vnbstmt{\begin{array}{@{}l@{}} \forall m, c, a, b.; \forall z \in \mathbb{C}. \\ \quad (\operatorname{is\text{-}road}(m, c, a, b) \wedge \neg (z \in \operatorname{trace}(m, a, b))) \Rightarrow  \\ \quad \quad \operatorname{winding}(m, c, a, b, z) \in \mathbb{C} \end{array}}}
\vnbprop{winding-zero-outside-disc}{\vnbstmt{\begin{array}{@{}l@{}} \forall m, d, a, b.; \forall c \in \mathbb{C}, r \in \mathbb{R}, z \in \mathbb{C}. \\ \quad \operatorname{is\text{-}road}(m, d, a, b) \wedge \\ \quad m(b) = m(a) \wedge \\ \quad 0 < r \wedge \\ \quad \operatorname{trace}(m, a, b) \subseteq \operatorname{ball}(\operatorname{cc\text{-}ms}, c, r) \wedge \\ \quad \neg (z \in \operatorname{ball}(\operatorname{cc\text{-}ms}, c, r)) \Rightarrow \\ \quad \quad \operatorname{winding}(m, d, a, b, z) = 0 \end{array}}}
\vnbprop{winding-circle-center}{\vnbstmt{\begin{array}{@{}l@{}} \forall c \in \mathbb{C}, r \in \mathbb{R}. \\ \quad (0 < r) \Rightarrow (\operatorname{winding}(\operatorname{circle\text{-}path}(c, r), \operatorname{circle\text{-}deriv}(c, r), 0, (2 \cdot \pi), c) = 1) \end{array}}}
\vnbprop{road-holomorphic-image}{\vnbstmt{\begin{array}{@{}l@{}} \forall u, f, c, m, d, a, b. \\ \quad \operatorname{holomorphic\text{-}on}(u, f) \wedge \\ \quad c \in (u \to \mathbb{C}) \wedge \\ \quad \forall z \in u.\; \operatorname{is\text{-}diff\text{-}on}(\operatorname{cc\text{-}normed\text{-}field}, u, f, z, c(z)) \wedge \\ \quad \forall y \in u. \\ \quad \quad \operatorname{is\text{-}continuous\text{-}at}(\operatorname{subspace\text{-}ms}(\operatorname{nf\text{-}metric\text{-}space}(\operatorname{cc\text{-}normed\text{-}field}), u), \\ \quad \quad \quad \operatorname{nf\text{-}metric\text{-}space}(\operatorname{cc\text{-}normed\text{-}field}), c, y) \wedge \\ \quad \operatorname{is\text{-}road}(m, d, a, b) \wedge \\ \quad \operatorname{trace}(m, a, b) \subseteq u \Rightarrow \\ \quad \quad \operatorname{is\text{-}road}((\lambda t \in \operatorname{ccint}(a, b).\; f(m(t))), (\lambda t \in \operatorname{ccint}(a, b).\; (c(m(t)) \cdot d(t))), a, b) \end{array}}}
\vnbprop{circle-is-road}{\vnbstmt{\begin{array}{@{}l@{}} \forall c \in \mathbb{C}, r \in \mathbb{R}. \\ \quad \operatorname{is\text{-}road}(\operatorname{circle\text{-}path}(c, r), \operatorname{circle\text{-}deriv}(c, r), 0, (2 \cdot \pi)) \end{array}}}

\vnbbookfile{theorem-library/analytic-from-integral}{theorem-library/analytic-from-integral}
\noindent{\small analytic-from-integral.scm -- THE CAUCHY-TYPE INTEGRAL IS A POWER SERIES (complex-analysis.pdf 3.5: Lemma 3.33 with Remark 3.34, the engine of Theorem 3.35; 3.7: the Cauchy estimate behind Prop 3.47 / Cor 3.48). Agent CA-2 of the October roadmap (docs/roadmap-2026-10-03.md 0(a)), 2026-10-03.}\par
\vnbprop{aif-cont-product}{\vnbstmt{\begin{array}{@{}l@{}} \forall u, f, g, h.; \forall y \in u. \\ \quad u \subseteq \mathbb{C} \wedge \\ \quad f \in (u \to \mathbb{C}) \wedge \\ \quad \forall y \in u. \\ \quad \quad \operatorname{is\text{-}continuous\text{-}at}(\operatorname{subspace\text{-}ms}(\operatorname{nf\text{-}metric\text{-}space}(\operatorname{cc\text{-}normed\text{-}field}), u), \\ \quad \quad \quad \operatorname{nf\text{-}metric\text{-}space}(\operatorname{cc\text{-}normed\text{-}field}), f, y) \wedge \\ \quad g \in (u \to \mathbb{C}) \wedge \\ \quad \forall y \in u. \\ \quad \quad \operatorname{is\text{-}continuous\text{-}at}(\operatorname{subspace\text{-}ms}(\operatorname{nf\text{-}metric\text{-}space}(\operatorname{cc\text{-}normed\text{-}field}), u), \\ \quad \quad \quad \operatorname{nf\text{-}metric\text{-}space}(\operatorname{cc\text{-}normed\text{-}field}), g, y) \wedge \\ \quad h \in (u \to \mathbb{C}) \wedge \\ \quad \forall x \in u.\; (h(x) \simeq (f(x) \cdot g(x))) \Rightarrow \\ \quad \quad \operatorname{is\text{-}continuous\text{-}at}(\operatorname{subspace\text{-}ms}(\operatorname{nf\text{-}metric\text{-}space}(\operatorname{cc\text{-}normed\text{-}field}), u), \\ \quad \quad \quad \operatorname{nf\text{-}metric\text{-}space}(\operatorname{cc\text{-}normed\text{-}field}), h, y) \end{array}}}
\vnbprop{aif-cont-sum}{\vnbstmt{\begin{array}{@{}l@{}} \forall u, f, g, h.; \forall y \in u. \\ \quad u \subseteq \mathbb{C} \wedge \\ \quad f \in (u \to \mathbb{C}) \wedge \\ \quad \forall y \in u. \\ \quad \quad \operatorname{is\text{-}continuous\text{-}at}(\operatorname{subspace\text{-}ms}(\operatorname{nf\text{-}metric\text{-}space}(\operatorname{cc\text{-}normed\text{-}field}), u), \\ \quad \quad \quad \operatorname{nf\text{-}metric\text{-}space}(\operatorname{cc\text{-}normed\text{-}field}), f, y) \wedge \\ \quad g \in (u \to \mathbb{C}) \wedge \\ \quad \forall y \in u. \\ \quad \quad \operatorname{is\text{-}continuous\text{-}at}(\operatorname{subspace\text{-}ms}(\operatorname{nf\text{-}metric\text{-}space}(\operatorname{cc\text{-}normed\text{-}field}), u), \\ \quad \quad \quad \operatorname{nf\text{-}metric\text{-}space}(\operatorname{cc\text{-}normed\text{-}field}), g, y) \wedge \\ \quad h \in (u \to \mathbb{C}) \wedge \\ \quad \forall x \in u.\; (h(x) \simeq (f(x) + g(x))) \Rightarrow \\ \quad \quad \operatorname{is\text{-}continuous\text{-}at}(\operatorname{subspace\text{-}ms}(\operatorname{nf\text{-}metric\text{-}space}(\operatorname{cc\text{-}normed\text{-}field}), u), \\ \quad \quad \quad \operatorname{nf\text{-}metric\text{-}space}(\operatorname{cc\text{-}normed\text{-}field}), h, y) \end{array}}}
\vnbprop{aif-cont-holo-restrict}{\vnbstmt{\begin{array}{@{}l@{}} \forall v, g, u.; \forall h \in (u \to \mathbb{C}), y \in u. \\ \quad (\operatorname{holomorphic\text{-}on}(v, g) \wedge ((u \subseteq v) \wedge \forall x \in u.\; (h(x) \simeq g(x)))) \Rightarrow  \\ \quad \quad \operatorname{is\text{-}continuous\text{-}at}(\operatorname{subspace\text{-}ms}(\operatorname{nf\text{-}metric\text{-}space}(\operatorname{cc\text{-}normed\text{-}field}), u), \\ \quad \quad \quad \operatorname{nf\text{-}metric\text{-}space}(\operatorname{cc\text{-}normed\text{-}field}), h, y) \end{array}}}
\vnbprop{aif-cont-compose}{\vnbstmt{\begin{array}{@{}l@{}} \forall u.; \forall g \in (u \to \mathbb{C}), v, \mathit{aiph}.; \forall h \in (u \to \mathbb{C}), y \in u. \\ \quad u \subseteq \mathbb{C} \wedge \\ \quad \forall y \in u. \\ \quad \quad \operatorname{is\text{-}continuous\text{-}at}(\operatorname{subspace\text{-}ms}(\operatorname{nf\text{-}metric\text{-}space}(\operatorname{cc\text{-}normed\text{-}field}), u), \\ \quad \quad \quad \operatorname{nf\text{-}metric\text{-}space}(\operatorname{cc\text{-}normed\text{-}field}), g, y) \wedge \\ \quad v \subseteq \mathbb{C} \wedge \\ \quad \forall x \in u.\; (g(x) \in v) \wedge \\ \quad \forall y \in v. \\ \quad \quad \operatorname{is\text{-}continuous\text{-}at}(\operatorname{subspace\text{-}ms}(\operatorname{nf\text{-}metric\text{-}space}(\operatorname{cc\text{-}normed\text{-}field}), v), \\ \quad \quad \quad \operatorname{nf\text{-}metric\text{-}space}(\operatorname{cc\text{-}normed\text{-}field}), \mathit{aiph}, y) \wedge \\ \quad \forall x \in u.\; (h(x) \simeq \operatorname{aiph}(g(x))) \Rightarrow \\ \quad \quad \operatorname{is\text{-}continuous\text{-}at}(\operatorname{subspace\text{-}ms}(\operatorname{nf\text{-}metric\text{-}space}(\operatorname{cc\text{-}normed\text{-}field}), u), \\ \quad \quad \quad \operatorname{nf\text{-}metric\text{-}space}(\operatorname{cc\text{-}normed\text{-}field}), h, y) \end{array}}}
\vnbprop{aif-cont-affine}{\vnbstmt{\begin{array}{@{}l@{}} \forall u.; \forall p, q, w \in \mathbb{C}, h \in (u \to \mathbb{C}), y \in u. \\ \quad ((u \subseteq \mathbb{C}) \wedge \forall x \in u.\; (h(x) \simeq (p + (q \cdot (x - w))))) \Rightarrow  \\ \quad \quad \operatorname{is\text{-}continuous\text{-}at}(\operatorname{subspace\text{-}ms}(\operatorname{nf\text{-}metric\text{-}space}(\operatorname{cc\text{-}normed\text{-}field}), u), \\ \quad \quad \quad \operatorname{nf\text{-}metric\text{-}space}(\operatorname{cc\text{-}normed\text{-}field}), h, y) \end{array}}}
\vnbprop{aif-cont-shift}{\vnbstmt{\begin{array}{@{}l@{}} \forall u.; \forall c \in \mathbb{C}, h \in (u \to \mathbb{C}), y \in u. \\ \quad ((u \subseteq \mathbb{C}) \wedge \forall x \in u.\; (h(x) \simeq (x - c))) \Rightarrow  \\ \quad \quad \operatorname{is\text{-}continuous\text{-}at}(\operatorname{subspace\text{-}ms}(\operatorname{nf\text{-}metric\text{-}space}(\operatorname{cc\text{-}normed\text{-}field}), u), \\ \quad \quad \quad \operatorname{nf\text{-}metric\text{-}space}(\operatorname{cc\text{-}normed\text{-}field}), h, y) \end{array}}}
\vnbprop{aif-cont-const}{\vnbstmt{\begin{array}{@{}l@{}} \forall u.; \forall k \in \mathbb{C}, h \in (u \to \mathbb{C}), y \in u. \\ \quad ((u \subseteq \mathbb{C}) \wedge \forall x \in u.\; (h(x) \simeq k)) \Rightarrow  \\ \quad \quad \operatorname{is\text{-}continuous\text{-}at}(\operatorname{subspace\text{-}ms}(\operatorname{nf\text{-}metric\text{-}space}(\operatorname{cc\text{-}normed\text{-}field}), u), \\ \quad \quad \quad \operatorname{nf\text{-}metric\text{-}space}(\operatorname{cc\text{-}normed\text{-}field}), h, y) \end{array}}}
\vnbprop{aif-cont-recip-shift}{\vnbstmt{\begin{array}{@{}l@{}} \forall u.; \forall c \in \mathbb{C}, h \in (u \to \mathbb{C}), y \in u. \\ \quad ((u \subseteq \mathbb{C}) \wedge (\forall x \in u.\; \neg (x = c) \wedge \forall x \in u.\; (h(x) \simeq \operatorname{recip}((x - c))))) \Rightarrow  \\ \quad \quad \operatorname{is\text{-}continuous\text{-}at}(\operatorname{subspace\text{-}ms}(\operatorname{nf\text{-}metric\text{-}space}(\operatorname{cc\text{-}normed\text{-}field}), u), \\ \quad \quad \quad \operatorname{nf\text{-}metric\text{-}space}(\operatorname{cc\text{-}normed\text{-}field}), h, y) \end{array}}}
\vnbprop{aif-cont-power-of}{\vnbstmt{\begin{array}{@{}l@{}} \forall u.; \forall g \in (u \to \mathbb{C}), n \in \mathbb{N}, h \in (u \to \mathbb{C}), y \in u. \\ \quad u \subseteq \mathbb{C} \wedge \\ \quad \forall y \in u. \\ \quad \quad \operatorname{is\text{-}continuous\text{-}at}(\operatorname{subspace\text{-}ms}(\operatorname{nf\text{-}metric\text{-}space}(\operatorname{cc\text{-}normed\text{-}field}), u), \\ \quad \quad \quad \operatorname{nf\text{-}metric\text{-}space}(\operatorname{cc\text{-}normed\text{-}field}), g, y) \wedge \\ \quad \forall x \in u.\; (h(x) \simeq \operatorname{power}(g(x), n)) \Rightarrow \\ \quad \quad \operatorname{is\text{-}continuous\text{-}at}(\operatorname{subspace\text{-}ms}(\operatorname{nf\text{-}metric\text{-}space}(\operatorname{cc\text{-}normed\text{-}field}), u), \\ \quad \quad \quad \operatorname{nf\text{-}metric\text{-}space}(\operatorname{cc\text{-}normed\text{-}field}), h, y) \end{array}}}
\vnbprop{aif-magnitude-recip}{\vnbstmt{\forall d \in \mathbb{C}.\; (\neg (d = 0) \Rightarrow (\operatorname{magnitude}(\operatorname{recip}(d)) = \operatorname{recip}(\operatorname{magnitude}(d))))}}
\vnbprop{aif-geometric-partial}{\vnbstmt{\begin{array}{@{}l@{}} \forall n \in \mathbb{N}, p, v, \mathit{aiqv}, e. \\ \quad ((p \in \mathbb{C}) \wedge ((v \in \mathbb{C}) \wedge ((\mathit{aiqv} \in \mathbb{C}) \wedge ((e \in \mathbb{C}) \wedge ((e \cdot (1 - (\mathit{aiqv} \cdot v))) = v))))) \Rightarrow  \\ \quad \quad \operatorname{cc\text{-}series\text{-}partial\text{-}sum}((\lambda j \in \mathbb{N}.\; ((p \cdot \operatorname{power}(v, (j + 1))) \cdot \operatorname{power}(\mathit{aiqv}, j))), n) \\ \quad \quad =\; ((p \cdot e) - ((p \cdot \operatorname{power}((\mathit{aiqv} \cdot v), n)) \cdot e)) \end{array}}}
\vnbprop{aif-trace-bounded}{\vnbstmt{\begin{array}{@{}l@{}} \forall u.; \forall i \in (u \to \mathbb{C}), m, a, b. \\ \quad u \subseteq \mathbb{C} \wedge \\ \quad \forall y \in u. \\ \quad \quad \operatorname{is\text{-}continuous\text{-}at}(\operatorname{subspace\text{-}ms}(\operatorname{nf\text{-}metric\text{-}space}(\operatorname{cc\text{-}normed\text{-}field}), u), \\ \quad \quad \quad \operatorname{nf\text{-}metric\text{-}space}(\operatorname{cc\text{-}normed\text{-}field}), i, y) \wedge \\ \quad \operatorname{is\text{-}path}(m, a, b) \wedge \\ \quad \operatorname{trace}(m, a, b) \subseteq u \Rightarrow \\ \quad \quad \exists \mathit{aim} \in \mathbb{R}. \\ \quad \quad \quad (0 \leq \mathit{aim}) \wedge \forall \mathit{aiy} \in \operatorname{trace}(m, a, b).\; (\operatorname{magnitude}(i(\mathit{aiy})) \leq \mathit{aim}) \end{array}}}
\vnbprop{aif-scale-small}{\vnbstmt{\begin{array}{@{}l@{}} \forall l \in \mathbb{R}, s. \\ \quad ((0 \leq l) \wedge \operatorname{pos\text{-}rr}(s)) \Rightarrow  \\ \quad \quad \operatorname{pos\text{-}rr}((s \cdot \operatorname{recip}((l + 1)))) \wedge ((l \cdot (s \cdot \operatorname{recip}((l + 1)))) \leq s) \end{array}}}
\vnbprop{aif-cont-scale}{\vnbstmt{\begin{array}{@{}l@{}} \forall u.; \forall k \in \mathbb{C}, g, h \in (u \to \mathbb{C}), y \in u. \\ \quad u \subseteq \mathbb{C} \wedge \\ \quad \forall y \in u. \\ \quad \quad \operatorname{is\text{-}continuous\text{-}at}(\operatorname{subspace\text{-}ms}(\operatorname{nf\text{-}metric\text{-}space}(\operatorname{cc\text{-}normed\text{-}field}), u), \\ \quad \quad \quad \operatorname{nf\text{-}metric\text{-}space}(\operatorname{cc\text{-}normed\text{-}field}), g, y) \wedge \\ \quad \forall x \in u.\; (h(x) \simeq (k \cdot g(x))) \Rightarrow \\ \quad \quad \operatorname{is\text{-}continuous\text{-}at}(\operatorname{subspace\text{-}ms}(\operatorname{nf\text{-}metric\text{-}space}(\operatorname{cc\text{-}normed\text{-}field}), u), \\ \quad \quad \quad \operatorname{nf\text{-}metric\text{-}space}(\operatorname{cc\text{-}normed\text{-}field}), h, y) \end{array}}}
\vnbprop{line-int-of-uniform-limit}{\vnbstmt{\begin{array}{@{}l@{}} \forall u, m, c, a, b.; \forall g \in (u \to \mathbb{C}), d. \\ \quad u \subseteq \mathbb{C} \wedge \\ \quad \operatorname{is\text{-}road}(m, c, a, b) \wedge \\ \quad \operatorname{trace}(m, a, b) \subseteq u \wedge \\ \quad \forall y \in u. \\ \quad \quad \operatorname{is\text{-}continuous\text{-}at}(\operatorname{subspace\text{-}ms}(\operatorname{nf\text{-}metric\text{-}space}(\operatorname{cc\text{-}normed\text{-}field}), u), \\ \quad \quad \quad \operatorname{nf\text{-}metric\text{-}space}(\operatorname{cc\text{-}normed\text{-}field}), g, y) \wedge \\ \quad \forall k \in \mathbb{N}.\; (d(k) \in (u \to \mathbb{C})) \wedge \\ \quad \forall k \in \mathbb{N}, y \in u. \\ \quad \quad \operatorname{is\text{-}continuous\text{-}at}(\operatorname{subspace\text{-}ms}(\operatorname{nf\text{-}metric\text{-}space}(\operatorname{cc\text{-}normed\text{-}field}), u), \\ \quad \quad \quad \operatorname{nf\text{-}metric\text{-}space}(\operatorname{cc\text{-}normed\text{-}field}), d(k), y) \wedge \\ \quad \forall s. \\ \quad \quad \operatorname{pos\text{-}rr}(s) \Rightarrow  \\ \quad \quad \quad \exists p \in \mathbb{N}. \forall k \in \mathbb{N}, \mathit{aiy} \in \operatorname{trace}(m, a, b). \\ \quad \quad \quad \quad (p \leq k) \Rightarrow (\operatorname{magnitude}((d(k)(\mathit{aiy}) - g(\mathit{aiy}))) \leq s) \Rightarrow \\ \quad \quad \operatorname{converges\text{-}to}(\operatorname{cc\text{-}ms}, (\lambda k \in \mathbb{N}.\; \operatorname{line\text{-}int}(d(k), m, c, a, b)), \operatorname{line\text{-}int}(g, m, c, a, b)) \end{array}}}
\vnbprop{aif-geometric-relation}{\vnbstmt{\begin{array}{@{}l@{}} \forall w, c, z \in \mathbb{C}. \\ \quad (\neg (w = c) \wedge \neg (w = z)) \Rightarrow  \\ \quad \quad (\operatorname{recip}((w - z)) \cdot (1 - ((z - c) \cdot \operatorname{recip}((w - c))))) = \operatorname{recip}((w - c)) \end{array}}}
\vnbprop{aif-geometric-remainder}{\vnbstmt{\begin{array}{@{}l@{}} \forall c \in \mathbb{C}, r \in \mathbb{R}, w, z, p \in \mathbb{C}, m \in \mathbb{R}, k \in \mathbb{N}. \\ \quad ((0 < r) \wedge ((r \leq \operatorname{magnitude}((w - c))) \wedge ((\operatorname{magnitude}((z - c)) < r) \wedge (\operatorname{magnitude}(p) \leq m)))) \Rightarrow  \\ \quad \quad \operatorname{magnitude}((\operatorname{cc\text{-}series\text{-}partial\text{-}sum}((\lambda j \in \mathbb{N}.\; ((p \cdot \operatorname{power}(\operatorname{recip}((w - c)), (j + 1))) \cdot \operatorname{power}((z - c), j))), k) - (p \cdot \operatorname{recip}((w - z))))) \\ \quad \quad \in\; \mathbb{R} \wedge \\ \quad \quad \operatorname{magnitude}((\operatorname{cc\text{-}series\text{-}partial\text{-}sum}((\lambda j \in \mathbb{N}.\; ((p \cdot \operatorname{power}(\operatorname{recip}((w - c)), (j + 1))) \cdot \operatorname{power}((z - c), j))), k) - (p \cdot \operatorname{recip}((w - z))))) \\ \quad \quad \leq\; (m \cdot \operatorname{power}((\operatorname{magnitude}((z - c)) \cdot \operatorname{recip}(r)), k)) \cdot \operatorname{recip}((r - \operatorname{magnitude}((z - c)))) \end{array}}}
\vnbprop{geometric-expansion-uniform}{\vnbstmt{\begin{array}{@{}l@{}} \forall u, c, r.; \forall z \in \mathbb{C}, i.; \forall m \in \mathbb{R}, s. \\ \quad u \subseteq \mathbb{C} \wedge \\ \quad c \in \mathbb{C} \wedge \\ \quad r \in \mathbb{R} \wedge \\ \quad 0 < r \wedge \\ \quad \forall x \in u.\; (r \leq \operatorname{magnitude}((x - c))) \wedge \\ \quad \operatorname{magnitude}((z - c)) < r \wedge \\ \quad \forall x \in u.\; (i(x) \in \mathbb{C}) \wedge \\ \quad 0 \leq m \wedge \\ \quad \forall x \in u.\; (\operatorname{magnitude}(i(x)) \leq m) \wedge \\ \quad \operatorname{pos\text{-}rr}(s) \Rightarrow \\ \quad \quad \exists p \in \mathbb{N}. \forall k \in \mathbb{N}, y \in u. \\ \quad \quad \quad (p \leq k) \Rightarrow  \\ \quad \quad \quad \quad \operatorname{magnitude}((\operatorname{cc\text{-}series\text{-}partial\text{-}sum}((\lambda j \in \mathbb{N}.\; ((i(y) \cdot \operatorname{power}(\operatorname{recip}((y - c)), (j + 1))) \cdot \operatorname{power}((z - c), j))), k) - (i(y) \cdot \operatorname{recip}((y - z))))) \\ \quad \quad \quad \quad \leq\; s \end{array}}}
\vnbprop{aif-off-center}{\vnbstmt{\begin{array}{@{}l@{}} \forall c \in \mathbb{C}, r \in \mathbb{R}, w, z \in \mathbb{C}. \\ \quad ((0 < r) \wedge ((r \leq \operatorname{magnitude}((w - c))) \wedge (\operatorname{magnitude}((z - c)) < r))) \Rightarrow  \\ \quad \quad \neg (w = c) \wedge \neg (w = z) \end{array}}}
\vnbprop{aif-partial-integrals}{\vnbstmt{\begin{array}{@{}l@{}} \forall k \in \mathbb{N}, u, c, r, i, m, d, a, b, z. \\ \quad u \subseteq \mathbb{C} \wedge \\ \quad c \in \mathbb{C} \wedge \\ \quad r \in \mathbb{R} \wedge \\ \quad 0 < r \wedge \\ \quad \forall x \in u.\; (r \leq \operatorname{magnitude}((x - c))) \wedge \\ \quad i \in (u \to \mathbb{C}) \wedge \\ \quad \forall y \in u. \\ \quad \quad \operatorname{is\text{-}continuous\text{-}at}(\operatorname{subspace\text{-}ms}(\operatorname{nf\text{-}metric\text{-}space}(\operatorname{cc\text{-}normed\text{-}field}), u), \\ \quad \quad \quad \operatorname{nf\text{-}metric\text{-}space}(\operatorname{cc\text{-}normed\text{-}field}), i, y) \wedge \\ \quad \operatorname{is\text{-}road}(m, d, a, b) \wedge \\ \quad \operatorname{trace}(m, a, b) \subseteq u \wedge \\ \quad z \in \mathbb{C} \wedge \\ \quad \operatorname{magnitude}((z - c)) < r \Rightarrow \\ \quad \quad (\lambda w \in u.\; \\ \quad \quad \quad \operatorname{cc\text{-}series\text{-}partial\text{-}sum}((\lambda j \in \mathbb{N}.\; ((i(w) \cdot \operatorname{power}(\operatorname{recip}((w - c)), (j + 1))) \cdot \operatorname{power}((z - c), j))), k)) \\ \quad \quad \in\; (u \to \mathbb{C}) \wedge \\ \quad \quad \forall y \in u. \\ \quad \quad \quad \operatorname{is\text{-}continuous\text{-}at}(\operatorname{subspace\text{-}ms}(\operatorname{nf\text{-}metric\text{-}space}(\operatorname{cc\text{-}normed\text{-}field}), u), \\ \quad \quad \quad \quad \operatorname{nf\text{-}metric\text{-}space}(\operatorname{cc\text{-}normed\text{-}field}), \\ \quad \quad \quad \quad (\lambda w \in u.\; \\ \quad \quad \quad \quad \quad \operatorname{cc\text{-}series\text{-}partial\text{-}sum}((\lambda j \in \mathbb{N}.\; ((i(w) \cdot \operatorname{power}(\operatorname{recip}((w - c)), (j + 1))) \cdot \operatorname{power}((z - c), j))), k)), \\ \quad \quad \quad \quad y) \wedge \\ \quad \quad \operatorname{line\text{-}int}((\lambda w \in u.\; \\ \quad \quad \quad \quad \operatorname{cc\text{-}series\text{-}partial\text{-}sum}((\lambda j \in \mathbb{N}.\; ((i(w) \cdot \operatorname{power}(\operatorname{recip}((w - c)), (j + 1))) \cdot \operatorname{power}((z - c), j))), k)), \\ \quad \quad \quad m, d, a, b) \\ \quad \quad =\; \operatorname{cc\text{-}series\text{-}partial\text{-}sum}((\lambda \mathit{aitk} \in \mathbb{N}.\; \\ \quad \quad \quad \quad \operatorname{line\text{-}int}((\lambda w \in u.\; (i(w) \cdot \operatorname{power}(\operatorname{recip}((w - c)), (\mathit{aitk} + 1)))), m, d, a, b) \cdot \operatorname{power}((z - c), \mathit{aitk})), \\ \quad \quad \quad k) \end{array}}}
\vnbprop{cauchy-kernel-integral-series}{\vnbstmt{\begin{array}{@{}l@{}} \forall u, c, r, i, m, d, a, b.; \forall z \in \mathbb{C}. \\ \quad u \subseteq \mathbb{C} \wedge \\ \quad c \in \mathbb{C} \wedge \\ \quad r \in \mathbb{R} \wedge \\ \quad 0 < r \wedge \\ \quad \forall x \in u.\; (r \leq \operatorname{magnitude}((x - c))) \wedge \\ \quad i \in (u \to \mathbb{C}) \wedge \\ \quad \forall y \in u. \\ \quad \quad \operatorname{is\text{-}continuous\text{-}at}(\operatorname{subspace\text{-}ms}(\operatorname{nf\text{-}metric\text{-}space}(\operatorname{cc\text{-}normed\text{-}field}), u), \\ \quad \quad \quad \operatorname{nf\text{-}metric\text{-}space}(\operatorname{cc\text{-}normed\text{-}field}), i, y) \wedge \\ \quad \operatorname{is\text{-}road}(m, d, a, b) \wedge \\ \quad \operatorname{trace}(m, a, b) \subseteq u \wedge \\ \quad \operatorname{magnitude}((z - c)) < r \Rightarrow \\ \quad \quad \operatorname{cc\text{-}series\text{-}converges\text{-}to}((\lambda k \in \mathbb{N}.\; \\ \quad \quad \quad \quad \operatorname{line\text{-}int}((\lambda w \in u.\; (i(w) \cdot \operatorname{power}(\operatorname{recip}((w - c)), (k + 1)))), m, d, a, b) \cdot \operatorname{power}((z - c), k)), \\ \quad \quad \quad \operatorname{line\text{-}int}((\lambda w \in u.\; (i(w) \cdot \operatorname{recip}((w - z)))), m, d, a, b)) \end{array}}}
\vnbprop{analytic-from-integral-formula}{\vnbstmt{\begin{array}{@{}l@{}} \forall u, c, r, i, m, d, a, b.; \forall p \in \mathbb{C}, f.; \forall z \in \mathbb{C}. \\ \quad u \subseteq \mathbb{C} \wedge \\ \quad c \in \mathbb{C} \wedge \\ \quad r \in \mathbb{R} \wedge \\ \quad 0 < r \wedge \\ \quad \forall x \in u.\; (r \leq \operatorname{magnitude}((x - c))) \wedge \\ \quad i \in (u \to \mathbb{C}) \wedge \\ \quad \forall y \in u. \\ \quad \quad \operatorname{is\text{-}continuous\text{-}at}(\operatorname{subspace\text{-}ms}(\operatorname{nf\text{-}metric\text{-}space}(\operatorname{cc\text{-}normed\text{-}field}), u), \\ \quad \quad \quad \operatorname{nf\text{-}metric\text{-}space}(\operatorname{cc\text{-}normed\text{-}field}), i, y) \wedge \\ \quad \operatorname{is\text{-}road}(m, d, a, b) \wedge \\ \quad \operatorname{trace}(m, a, b) \subseteq u \wedge \\ \quad \forall z \in \mathbb{C}. \\ \quad \quad (\operatorname{magnitude}((z - c)) < r) \Rightarrow  \\ \quad \quad \quad f(z) = (p \cdot \operatorname{line\text{-}int}((\lambda w \in u.\; (i(w) \cdot \operatorname{recip}((w - z)))), m, d, a, b)) \wedge \\ \quad \operatorname{magnitude}((z - c)) < r \Rightarrow \\ \quad \quad \operatorname{cc\text{-}series\text{-}converges\text{-}to}((\lambda k \in \mathbb{N}.\; \\ \quad \quad \quad \quad (p \cdot \operatorname{line\text{-}int}((\lambda w \in u.\; (i(w) \cdot \operatorname{power}(\operatorname{recip}((w - c)), (k + 1)))), m, d, a, b)) \cdot \operatorname{power}((z - c), k)), \\ \quad \quad \quad f(z)) \end{array}}}
\vnbprop{aif-coefficient-bound}{\vnbstmt{\begin{array}{@{}l@{}} \forall u, c, r, i, m, d, a, b.; \forall \mathit{aim} \in \mathbb{R}, k \in \mathbb{N}. \\ \quad u \subseteq \mathbb{C} \wedge \\ \quad c \in \mathbb{C} \wedge \\ \quad r \in \mathbb{R} \wedge \\ \quad 0 < r \wedge \\ \quad \forall x \in u.\; (r \leq \operatorname{magnitude}((x - c))) \wedge \\ \quad i \in (u \to \mathbb{C}) \wedge \\ \quad \forall y \in u. \\ \quad \quad \operatorname{is\text{-}continuous\text{-}at}(\operatorname{subspace\text{-}ms}(\operatorname{nf\text{-}metric\text{-}space}(\operatorname{cc\text{-}normed\text{-}field}), u), \\ \quad \quad \quad \operatorname{nf\text{-}metric\text{-}space}(\operatorname{cc\text{-}normed\text{-}field}), i, y) \wedge \\ \quad \operatorname{is\text{-}road}(m, d, a, b) \wedge \\ \quad \operatorname{trace}(m, a, b) \subseteq u \wedge \\ \quad \forall \mathit{aiy} \in \operatorname{trace}(m, a, b).\; (\operatorname{magnitude}(i(\mathit{aiy})) \leq \mathit{aim}) \Rightarrow \\ \quad \quad \operatorname{magnitude}(\operatorname{line\text{-}int}((\lambda w \in u.\; (i(w) \cdot \operatorname{power}(\operatorname{recip}((w - c)), (k + 1)))), m, d, a, b)) \\ \quad \quad \leq\; \operatorname{road\text{-}length}(d, a, b) \cdot (\mathit{aim} \cdot \operatorname{power}(\operatorname{recip}(r), (k + 1))) \end{array}}}
\vnbprop{cauchy-estimate-from-formula}{\vnbstmt{\begin{array}{@{}l@{}} \forall u, c, r, i, m, d, a, b.; \forall \mathit{aim} \in \mathbb{R}, p \in \mathbb{C}, k \in \mathbb{N}. \\ \quad u \subseteq \mathbb{C} \wedge \\ \quad c \in \mathbb{C} \wedge \\ \quad r \in \mathbb{R} \wedge \\ \quad 0 < r \wedge \\ \quad \forall x \in u.\; (r \leq \operatorname{magnitude}((x - c))) \wedge \\ \quad i \in (u \to \mathbb{C}) \wedge \\ \quad \forall y \in u. \\ \quad \quad \operatorname{is\text{-}continuous\text{-}at}(\operatorname{subspace\text{-}ms}(\operatorname{nf\text{-}metric\text{-}space}(\operatorname{cc\text{-}normed\text{-}field}), u), \\ \quad \quad \quad \operatorname{nf\text{-}metric\text{-}space}(\operatorname{cc\text{-}normed\text{-}field}), i, y) \wedge \\ \quad \operatorname{is\text{-}road}(m, d, a, b) \wedge \\ \quad \operatorname{trace}(m, a, b) \subseteq u \wedge \\ \quad \forall \mathit{aiy} \in \operatorname{trace}(m, a, b).\; (\operatorname{magnitude}(i(\mathit{aiy})) \leq \mathit{aim}) \Rightarrow \\ \quad \quad \operatorname{magnitude}((p \cdot \operatorname{line\text{-}int}((\lambda w \in u.\; (i(w) \cdot \operatorname{power}(\operatorname{recip}((w - c)), (k + 1)))), m, d, a, b))) \\ \quad \quad \leq\; \operatorname{magnitude}(p) \cdot (\operatorname{road\text{-}length}(d, a, b) \cdot (\mathit{aim} \cdot \operatorname{power}(\operatorname{recip}(r), (k + 1)))) \end{array}}}
\vnbprop{cauchy-kernel-series-absolute}{\vnbstmt{\begin{array}{@{}l@{}} \forall u, c, r, i, m, d, a, b.; \forall z \in \mathbb{C}. \\ \quad u \subseteq \mathbb{C} \wedge \\ \quad c \in \mathbb{C} \wedge \\ \quad r \in \mathbb{R} \wedge \\ \quad 0 < r \wedge \\ \quad \forall x \in u.\; (r \leq \operatorname{magnitude}((x - c))) \wedge \\ \quad i \in (u \to \mathbb{C}) \wedge \\ \quad \forall y \in u. \\ \quad \quad \operatorname{is\text{-}continuous\text{-}at}(\operatorname{subspace\text{-}ms}(\operatorname{nf\text{-}metric\text{-}space}(\operatorname{cc\text{-}normed\text{-}field}), u), \\ \quad \quad \quad \operatorname{nf\text{-}metric\text{-}space}(\operatorname{cc\text{-}normed\text{-}field}), i, y) \wedge \\ \quad \operatorname{is\text{-}road}(m, d, a, b) \wedge \\ \quad \operatorname{trace}(m, a, b) \subseteq u \wedge \\ \quad \operatorname{magnitude}((z - c)) < r \Rightarrow \\ \quad \quad \operatorname{series\text{-}converges}((\lambda k \in \mathbb{N}.\; \\ \quad \quad \quad \quad \operatorname{magnitude}((\operatorname{line\text{-}int}((\lambda w \in u.\; (i(w) \cdot \operatorname{power}(\operatorname{recip}((w - c)), (k + 1)))), m, d, a, b) \cdot \operatorname{power}((z - c), k))))) \end{array}}}
\vnbprop{analytic-at-from-integral-formula}{\vnbstmt{\begin{array}{@{}l@{}} \forall u, c, r, i, m, d, a, b.; \forall p \in \mathbb{C}, \mathit{aiou}.; \forall f \in (\mathit{aiou} \to \mathbb{C}). \\ \quad u \subseteq \mathbb{C} \wedge \\ \quad c \in \mathbb{C} \wedge \\ \quad r \in \mathbb{R} \wedge \\ \quad 0 < r \wedge \\ \quad \forall x \in u.\; (r \leq \operatorname{magnitude}((x - c))) \wedge \\ \quad i \in (u \to \mathbb{C}) \wedge \\ \quad \forall y \in u. \\ \quad \quad \operatorname{is\text{-}continuous\text{-}at}(\operatorname{subspace\text{-}ms}(\operatorname{nf\text{-}metric\text{-}space}(\operatorname{cc\text{-}normed\text{-}field}), u), \\ \quad \quad \quad \operatorname{nf\text{-}metric\text{-}space}(\operatorname{cc\text{-}normed\text{-}field}), i, y) \wedge \\ \quad \operatorname{is\text{-}road}(m, d, a, b) \wedge \\ \quad \operatorname{trace}(m, a, b) \subseteq u \wedge \\ \quad \operatorname{is\text{-}open}(\operatorname{cc\text{-}ms}, \mathit{aiou}) \wedge \\ \quad \operatorname{ball}(\operatorname{cc\text{-}ms}, c, r) \subseteq \mathit{aiou} \wedge \\ \quad \forall z \in \mathbb{C}. \\ \quad \quad (\operatorname{magnitude}((z - c)) < r) \Rightarrow  \\ \quad \quad \quad f(z) = (p \cdot \operatorname{line\text{-}int}((\lambda w \in u.\; (i(w) \cdot \operatorname{recip}((w - z)))), m, d, a, b)) \Rightarrow \\ \quad \quad f \text{ is analytic at } c \end{array}}}

\vnbbookfile{theorem-library/cauchy-integral-formula}{theorem-library/cauchy-integral-formula}
\noindent{\small cauchy-integral-formula.scm -- CAUCHY WITH AN EXCEPTIONAL POINT ON A DISC, THE WINDING NUMBER OF A CIRCLE INSIDE IT, THE CAUCHY INTEGRAL FORMULA ON A DISC, THEOREM 3.35 (HOLOMORPHIC =$>$ ANALYTIC) AND THE HOLOMORPHY OF THE DERIVATIVE. (complex-analysis.pdf ch. 3: Cor 3.9 with one exceptional point, Example 3.13 /}\par
\vnbprop{cif-open-minus-point}{\vnbstmt{\begin{array}{@{}l@{}} \forall v.; \forall x \in \mathbb{C}. \\ \quad \operatorname{is\text{-}open}(\operatorname{nf\text{-}metric\text{-}space}(\operatorname{cc\text{-}normed\text{-}field}), v) \Rightarrow  \\ \quad \quad \operatorname{is\text{-}open}(\operatorname{nf\text{-}metric\text{-}space}(\operatorname{cc\text{-}normed\text{-}field}), \operatorname{difference}(v, \operatorname{singleton}(x))) \end{array}}}
\vnbprop{cauchy-convex-sub-except}{\vnbstmt{\begin{array}{@{}l@{}} \forall u, f, x, v, m, c, a, b. \\ \quad \operatorname{is\text{-}open}(\operatorname{nf\text{-}metric\text{-}space}(\operatorname{cc\text{-}normed\text{-}field}), u) \wedge \\ \quad f \in (u \to \mathbb{C}) \wedge \\ \quad \forall y \in u. \\ \quad \quad \operatorname{is\text{-}continuous\text{-}at}(\operatorname{subspace\text{-}ms}(\operatorname{nf\text{-}metric\text{-}space}(\operatorname{cc\text{-}normed\text{-}field}), u), \\ \quad \quad \quad \operatorname{nf\text{-}metric\text{-}space}(\operatorname{cc\text{-}normed\text{-}field}), f, y) \wedge \\ \quad x \in u \wedge \\ \quad \operatorname{holomorphic\text{-}on}(\operatorname{difference}(u, \operatorname{singleton}(x)), \operatorname{restrict}(f, \operatorname{difference}(u, \operatorname{singleton}(x)))) \wedge \\ \quad \operatorname{is\text{-}open}(\operatorname{nf\text{-}metric\text{-}space}(\operatorname{cc\text{-}normed\text{-}field}), v) \wedge \\ \quad \operatorname{is\text{-}convex\text{-}cc}(v) \wedge \\ \quad v \subseteq u \wedge \\ \quad \operatorname{is\text{-}road}(m, c, a, b) \wedge \\ \quad \operatorname{trace}(m, a, b) \subseteq v \wedge \\ \quad m(b) = m(a) \Rightarrow \\ \quad \quad \operatorname{line\text{-}int}(f, m, c, a, b) = 0 \end{array}}}
\vnbprop{cauchy-disc-except}{\vnbstmt{\begin{array}{@{}l@{}} \forall u, f, x.; \forall c \in \mathbb{C}, r \in \mathbb{R}, m, d, a, b. \\ \quad \operatorname{is\text{-}open}(\operatorname{nf\text{-}metric\text{-}space}(\operatorname{cc\text{-}normed\text{-}field}), u) \wedge \\ \quad f \in (u \to \mathbb{C}) \wedge \\ \quad \forall y \in u. \\ \quad \quad \operatorname{is\text{-}continuous\text{-}at}(\operatorname{subspace\text{-}ms}(\operatorname{nf\text{-}metric\text{-}space}(\operatorname{cc\text{-}normed\text{-}field}), u), \\ \quad \quad \quad \operatorname{nf\text{-}metric\text{-}space}(\operatorname{cc\text{-}normed\text{-}field}), f, y) \wedge \\ \quad x \in u \wedge \\ \quad \operatorname{holomorphic\text{-}on}(\operatorname{difference}(u, \operatorname{singleton}(x)), \operatorname{restrict}(f, \operatorname{difference}(u, \operatorname{singleton}(x)))) \wedge \\ \quad 0 < r \wedge \\ \quad \operatorname{ball}(\operatorname{cc\text{-}ms}, c, r) \subseteq u \wedge \\ \quad \operatorname{is\text{-}road}(m, d, a, b) \wedge \\ \quad \operatorname{trace}(m, a, b) \subseteq \operatorname{ball}(\operatorname{cc\text{-}ms}, c, r) \wedge \\ \quad m(b) = m(a) \Rightarrow \\ \quad \quad \operatorname{line\text{-}int}(f, m, d, a, b) = 0 \end{array}}}
\vnbprop{cif-recip-shift-diff}{\vnbstmt{\begin{array}{@{}l@{}} \forall c \in \mathbb{C}, w \in \operatorname{difference}(\mathbb{C}, \operatorname{singleton}(c)). \\ \quad \operatorname{is\text{-}diff\text{-}on}(\operatorname{cc\text{-}normed\text{-}field}, \operatorname{difference}(\mathbb{C}, \operatorname{singleton}(c)), \\ \quad \quad (\lambda v \in \operatorname{difference}(\mathbb{C}, \operatorname{singleton}(c)).\; \operatorname{recip}((v - c))), \\ \quad \quad w, -(\operatorname{recip}((w - c)) \cdot \operatorname{recip}((w - c)))) \end{array}}}
\vnbprop{cif-recip-power-diff}{\vnbstmt{\begin{array}{@{}l@{}} \forall j \in \mathbb{N}, c \in \mathbb{C}, w \in \operatorname{difference}(\mathbb{C}, \operatorname{singleton}(c)). \\ \quad \operatorname{is\text{-}diff\text{-}on}(\operatorname{cc\text{-}normed\text{-}field}, \operatorname{difference}(\mathbb{C}, \operatorname{singleton}(c)), \\ \quad \quad (\lambda v \in \operatorname{difference}(\mathbb{C}, \operatorname{singleton}(c)).\; \operatorname{power}(\operatorname{recip}((v - c)), (j + 1))), \\ \quad \quad w, (-(j + 1) \cdot \operatorname{power}(\operatorname{recip}((w - c)), ((j + 1) + 1)))) \end{array}}}
\vnbprop{cif-recip-power-closed-zero}{\vnbstmt{\begin{array}{@{}l@{}} \forall j \in \mathbb{N}, c \in \mathbb{C}, m, d, a, b. \\ \quad (\operatorname{is\text{-}road}(m, d, a, b) \wedge ((\operatorname{trace}(m, a, b) \subseteq \operatorname{difference}(\mathbb{C}, \operatorname{singleton}(c))) \wedge (m(b) = m(a)))) \Rightarrow  \\ \quad \quad \operatorname{line\text{-}int}((\lambda v \in \operatorname{difference}(\mathbb{C}, \operatorname{singleton}(c)).\; \operatorname{power}(\operatorname{recip}((v - c)), ((j + 1) + 1))), m, d, a, b) = 0 \end{array}}}
\vnbprop{circle-trace-magnitude}{\vnbstmt{\begin{array}{@{}l@{}} \forall c \in \mathbb{C}, r \in \mathbb{R}, y \in \operatorname{trace}(\operatorname{circle\text{-}path}(c, r), 0, (2 \cdot \pi)). \\ \quad (0 \leq r) \Rightarrow (\operatorname{magnitude}((y - c)) = r) \end{array}}}
\vnbprop{circle-path-closed}{\vnbstmt{\forall c \in \mathbb{C}, r \in \mathbb{R}.\; (\operatorname{circle\text{-}path}(c, r)((2 \cdot \pi)) = \operatorname{circle\text{-}path}(c, r)(0))}}
\vnbprop{winding-circle-inside}{\vnbstmt{\begin{array}{@{}l@{}} \forall c \in \mathbb{C}, r \in \mathbb{R}, z \in \mathbb{C}. \\ \quad ((0 < r) \wedge (\operatorname{magnitude}((z - c)) < r)) \Rightarrow  \\ \quad \quad \operatorname{winding}(\operatorname{circle\text{-}path}(c, r), \operatorname{circle\text{-}deriv}(c, r), 0, (2 \cdot \pi), z) = 1 \end{array}}}
\vnbprop{cauchy-integral-formula-disc}{\vnbstmt{\begin{array}{@{}l@{}} \forall u, f.; \forall c \in \mathbb{C}, s, r \in \mathbb{R}, z \in \mathbb{C}. \\ \quad \operatorname{holomorphic\text{-}on}(u, f) \wedge \\ \quad 0 < s \wedge \\ \quad \operatorname{ball}(\operatorname{cc\text{-}ms}, c, s) \subseteq u \wedge \\ \quad 0 < r \wedge \\ \quad r < s \wedge \\ \quad \operatorname{magnitude}((z - c)) < r \Rightarrow \\ \quad \quad f(z) \\ \quad \quad =\; \operatorname{recip}(((2 \cdot \pi) \cdot +i)) \cdot \operatorname{line\text{-}int}((\lambda w \in \operatorname{trace}(\operatorname{circle\text{-}path}(c, r), 0, (2 \cdot \pi)).\; (f(w) \cdot \operatorname{recip}((w - z)))), \operatorname{circle\text{-}path}(c, r), \operatorname{circle\text{-}deriv}(c, r), 0, (2 \cdot \pi)) \end{array}}}
\vnbprop{cauchy-integral-formula-disc-restrict}{\vnbstmt{\begin{array}{@{}l@{}} \forall u, f.; \forall c \in \mathbb{C}, s, r \in \mathbb{R}, z \in \mathbb{C}. \\ \quad \operatorname{holomorphic\text{-}on}(u, f) \wedge \\ \quad 0 < s \wedge \\ \quad \operatorname{ball}(\operatorname{cc\text{-}ms}, c, s) \subseteq u \wedge \\ \quad 0 < r \wedge \\ \quad r < s \wedge \\ \quad \operatorname{magnitude}((z - c)) < r \Rightarrow \\ \quad \quad f(z) \\ \quad \quad =\; \operatorname{recip}(((2 \cdot \pi) \cdot +i)) \cdot \operatorname{line\text{-}int}((\lambda w \in \operatorname{trace}(\operatorname{circle\text{-}path}(c, r), 0, (2 \cdot \pi)).\; (\operatorname{restrict}(f, \operatorname{trace}(\operatorname{circle\text{-}path}(c, r), 0, (2 \cdot \pi)))(w) \cdot \operatorname{recip}((w - z)))), \operatorname{circle\text{-}path}(c, r), \operatorname{circle\text{-}deriv}(c, r), 0, (2 \cdot \pi)) \end{array}}}
\vnbprop{holomorphic-power-series-disc}{\vnbstmt{\begin{array}{@{}l@{}} \forall u, f.; \forall c \in \mathbb{C}, s, r \in \mathbb{R}, z \in \mathbb{C}. \\ \quad \operatorname{holomorphic\text{-}on}(u, f) \wedge \\ \quad 0 < s \wedge \\ \quad \operatorname{ball}(\operatorname{cc\text{-}ms}, c, s) \subseteq u \wedge \\ \quad 0 < r \wedge \\ \quad r < s \wedge \\ \quad \operatorname{magnitude}((z - c)) < r \Rightarrow \\ \quad \quad \operatorname{cc\text{-}series\text{-}converges\text{-}to}((\lambda k \in \mathbb{N}.\; \\ \quad \quad \quad \quad (\operatorname{recip}(((2 \cdot \pi) \cdot +i)) \cdot \operatorname{line\text{-}int}((\lambda w \in \operatorname{trace}(\operatorname{circle\text{-}path}(c, r), 0, (2 \cdot \pi)).\; (\operatorname{restrict}(f, \operatorname{trace}(\operatorname{circle\text{-}path}(c, r), 0, (2 \cdot \pi)))(w) \cdot \operatorname{power}(\operatorname{recip}((w - c)), (k + 1)))), \operatorname{circle\text{-}path}(c, r), \operatorname{circle\text{-}deriv}(c, r), 0, (2 \cdot \pi))) \cdot \operatorname{power}((z - c), k)), \\ \quad \quad \quad f(z)) \end{array}}}
\vnbprop{holomorphic-analytic-at}{\vnbstmt{\forall u, f.; \forall c \in u.\; (\operatorname{holomorphic\text{-}on}(u, f) \Rightarrow f \text{ is analytic at } c)}}
\vnbprop{holomorphic-deriv-holomorphic}{\vnbstmt{\begin{array}{@{}l@{}} \forall u, f. \\ \quad \operatorname{holomorphic\text{-}on}(u, f) \Rightarrow  \\ \quad \quad \operatorname{holomorphic\text{-}on}(u, (\lambda w \in u.\; \operatorname{deriv\text{-}on}(\operatorname{cc\text{-}normed\text{-}field}, u, f, w))) \end{array}}}

\vnbbookfile{theorem-library/cauchy-estimates}{theorem-library/cauchy-estimates}
\noindent{\small cauchy-estimates.scm -- THE CIRCLE'S LENGTH, CAUCHY'S ESTIMATES, THE DERIVATIVE FORMULA, LIOUVILLE AND THE FUNDAMENTAL THEOREM OF ALGEBRA. (complex-analysis.pdf ch. 3: Prop 3.40 / Cor 3.42 through the coefficients, Prop 3.47 and Cor 3.48 (Cauchy's estimates, Liouville); ch. 1, 1.5 (the fundamental theorem of algebra).)}\par
\vnbprop{cc-magnitude-i}{\vnbstmt{\operatorname{magnitude}(+i) = 1}}
\vnbprop{circle-road-length}{\vnbstmt{\begin{array}{@{}l@{}} \forall c \in \mathbb{C}, r \in \mathbb{R}. \\ \quad (0 \leq r) \Rightarrow (\operatorname{road\text{-}length}(\operatorname{circle\text{-}deriv}(c, r), 0, (2 \cdot \pi)) = ((2 \cdot \pi) \cdot r)) \end{array}}}
\vnbprop{cauchy-kernel-constant-magnitude}{\vnbstmt{\operatorname{magnitude}(\operatorname{recip}(((2 \cdot \pi) \cdot +i))) = \operatorname{recip}((2 \cdot \pi))}}
\vnbprop{cauchy-estimate-circle}{\vnbstmt{\begin{array}{@{}l@{}} \forall u, f.; \forall c \in \mathbb{C}, s, r, d \in \mathbb{R}, k \in \mathbb{N}. \\ \quad \operatorname{holomorphic\text{-}on}(u, f) \wedge \\ \quad 0 < s \wedge \\ \quad \operatorname{ball}(\operatorname{cc\text{-}ms}, c, s) \subseteq u \wedge \\ \quad 0 < r \wedge \\ \quad r < s \wedge \\ \quad \forall y \in \operatorname{trace}(\operatorname{circle\text{-}path}(c, r), 0, (2 \cdot \pi)).\; (\operatorname{magnitude}(f(y)) \leq d) \Rightarrow \\ \quad \quad \operatorname{magnitude}(\operatorname{recip}(((2 \cdot \pi) \cdot +i)) \cdot \operatorname{line\text{-}int}((\lambda w \in \operatorname{trace}(\operatorname{circle\text{-}path}(c, r), 0, (2 \cdot \pi)).\; (\operatorname{restrict}(f, \operatorname{trace}(\operatorname{circle\text{-}path}(c, r), 0, (2 \cdot \pi)))(w) \cdot \operatorname{power}(\operatorname{recip}((w - c)), (k + 1)))), \operatorname{circle\text{-}path}(c, r), \operatorname{circle\text{-}deriv}(c, r), 0, (2 \cdot \pi))) \\ \quad \quad \leq\; d \cdot \operatorname{power}(\operatorname{recip}(r), k) \end{array}}}
\vnbprop{cauchy-estimate-center}{\vnbstmt{\begin{array}{@{}l@{}} \forall u, f.; \forall c \in \mathbb{C}, s, r, d \in \mathbb{R}. \\ \quad \operatorname{holomorphic\text{-}on}(u, f) \wedge \\ \quad 0 < s \wedge \\ \quad \operatorname{ball}(\operatorname{cc\text{-}ms}, c, s) \subseteq u \wedge \\ \quad 0 < r \wedge \\ \quad r < s \wedge \\ \quad \forall y \in \operatorname{trace}(\operatorname{circle\text{-}path}(c, r), 0, (2 \cdot \pi)).\; (\operatorname{magnitude}(f(y)) \leq d) \Rightarrow \\ \quad \quad \operatorname{magnitude}(f(c)) \leq d \end{array}}}
\vnbprop{cauchy-coefficients-radius}{\vnbstmt{\begin{array}{@{}l@{}} \forall u, f.; \forall c \in \mathbb{C}, s, r \in \mathbb{R}, g \in (\mathbb{N} \to \mathbb{C}), o \in \mathbb{R}. \\ \quad \operatorname{holomorphic\text{-}on}(u, f) \wedge \\ \quad 0 < s \wedge \\ \quad \operatorname{ball}(\operatorname{cc\text{-}ms}, c, s) \subseteq u \wedge \\ \quad 0 < r \wedge \\ \quad r < s \wedge \\ \quad \forall k \in \mathbb{N}. \\ \quad \quad g(k) \\ \quad \quad =\; \operatorname{recip}(((2 \cdot \pi) \cdot +i)) \cdot \operatorname{line\text{-}int}((\lambda w \in \operatorname{trace}(\operatorname{circle\text{-}path}(c, r), 0, (2 \cdot \pi)).\; (\operatorname{restrict}(f, \operatorname{trace}(\operatorname{circle\text{-}path}(c, r), 0, (2 \cdot \pi)))(w) \cdot \operatorname{power}(\operatorname{recip}((w - c)), (k + 1)))), \operatorname{circle\text{-}path}(c, r), \operatorname{circle\text{-}deriv}(c, r), 0, (2 \cdot \pi)) \wedge \\ \quad 0 \leq o \wedge \\ \quad o < r \Rightarrow \\ \quad \quad o < \operatorname{cps\text{-}radius}(g) \end{array}}}
\vnbprop{holomorphic-coefficient-series}{\vnbstmt{\begin{array}{@{}l@{}} \forall u, f.; \forall c \in \mathbb{C}, s, r \in \mathbb{R}, g \in (\mathbb{N} \to \mathbb{C}), z \in \mathbb{C}. \\ \quad \operatorname{holomorphic\text{-}on}(u, f) \wedge \\ \quad 0 < s \wedge \\ \quad \operatorname{ball}(\operatorname{cc\text{-}ms}, c, s) \subseteq u \wedge \\ \quad 0 < r \wedge \\ \quad r < s \wedge \\ \quad \forall k \in \mathbb{N}. \\ \quad \quad g(k) \\ \quad \quad =\; \operatorname{recip}(((2 \cdot \pi) \cdot +i)) \cdot \operatorname{line\text{-}int}((\lambda w \in \operatorname{trace}(\operatorname{circle\text{-}path}(c, r), 0, (2 \cdot \pi)).\; (\operatorname{restrict}(f, \operatorname{trace}(\operatorname{circle\text{-}path}(c, r), 0, (2 \cdot \pi)))(w) \cdot \operatorname{power}(\operatorname{recip}((w - c)), (k + 1)))), \operatorname{circle\text{-}path}(c, r), \operatorname{circle\text{-}deriv}(c, r), 0, (2 \cdot \pi)) \wedge \\ \quad \operatorname{magnitude}((z - c)) < r \Rightarrow \\ \quad \quad f(z) = \operatorname{cc\text{-}series\text{-}limit}((\lambda \mathit{cpsk} \in \mathbb{N}.\; (g(\mathit{cpsk}) \cdot \operatorname{power}((z - c), \mathit{cpsk})))) \end{array}}}
\vnbprop{cauchy-derivative-formula}{\vnbstmt{\begin{array}{@{}l@{}} \forall u, f.; \forall c \in \mathbb{C}, s, r, o \in \mathbb{R}, k \in \mathbb{N}. \\ \quad \operatorname{holomorphic\text{-}on}(u, f) \wedge \\ \quad 0 < s \wedge \\ \quad \operatorname{ball}(\operatorname{cc\text{-}ms}, c, s) \subseteq u \wedge \\ \quad 0 < r \wedge \\ \quad r < s \wedge \\ \quad 0 < o \wedge \\ \quad o < r \Rightarrow \\ \quad \quad \operatorname{restrict}(f, \operatorname{ball}(\operatorname{cc\text{-}ms}, c, o))^{(k)}(c) \\ \quad \quad =\; \operatorname{factorial}(k) \cdot (\operatorname{recip}(((2 \cdot \pi) \cdot +i)) \cdot \operatorname{line\text{-}int}((\lambda w \in \operatorname{trace}(\operatorname{circle\text{-}path}(c, r), 0, (2 \cdot \pi)).\; (\operatorname{restrict}(f, \operatorname{trace}(\operatorname{circle\text{-}path}(c, r), 0, (2 \cdot \pi)))(w) \cdot \operatorname{power}(\operatorname{recip}((w - c)), (k + 1)))), \operatorname{circle\text{-}path}(c, r), \operatorname{circle\text{-}deriv}(c, r), 0, (2 \cdot \pi))) \end{array}}}
\vnbprop{cauchy-estimate-derivative}{\vnbstmt{\begin{array}{@{}l@{}} \forall u, f.; \forall c \in \mathbb{C}, s, r, o, d \in \mathbb{R}, k \in \mathbb{N}. \\ \quad \operatorname{holomorphic\text{-}on}(u, f) \wedge \\ \quad 0 < s \wedge \\ \quad \operatorname{ball}(\operatorname{cc\text{-}ms}, c, s) \subseteq u \wedge \\ \quad 0 < r \wedge \\ \quad r < s \wedge \\ \quad 0 < o \wedge \\ \quad o < r \wedge \\ \quad \forall y \in \operatorname{trace}(\operatorname{circle\text{-}path}(c, r), 0, (2 \cdot \pi)).\; (\operatorname{magnitude}(f(y)) \leq d) \Rightarrow \\ \quad \quad \operatorname{magnitude}(\operatorname{restrict}(f, \operatorname{ball}(\operatorname{cc\text{-}ms}, c, o))^{(k)}(c)) \\ \quad \quad \leq\; \operatorname{factorial}(k) \cdot (d \cdot \operatorname{power}(\operatorname{recip}(r), k)) \end{array}}}
\vnbprop{cie-squeeze}{\vnbstmt{\begin{array}{@{}l@{}} \forall x, k, a \in \mathbb{R}. \\ \quad ((0 \leq a) \wedge ((0 \leq k) \wedge \forall q \in \mathbb{R}.\; ((a < q) \Rightarrow (x \leq (k \cdot \operatorname{recip}(q)))))) \Rightarrow  \\ \quad \quad x \leq 0 \end{array}}}
\vnbprop{bounded-entire-deriv-zero}{\vnbstmt{\begin{array}{@{}l@{}} \forall f.; \forall d \in \mathbb{R}, a \in \mathbb{C}. \\ \quad (\operatorname{holomorphic\text{-}on}(\mathbb{C}, f) \wedge \forall y \in \mathbb{C}.\; (\operatorname{magnitude}(f(y)) \leq d)) \Rightarrow  \\ \quad \quad \operatorname{is\text{-}diff\text{-}on}(\operatorname{cc\text{-}normed\text{-}field}, \mathbb{C}, f, a, 0) \end{array}}}
\vnbprop{liouville}{\vnbstmt{\begin{array}{@{}l@{}} \forall f.; \forall d \in \mathbb{R}, z \in \mathbb{C}. \\ \quad (\operatorname{holomorphic\text{-}on}(\mathbb{C}, f) \wedge \forall y \in \mathbb{C}.\; (\operatorname{magnitude}(f(y)) \leq d)) \Rightarrow  \\ \quad \quad f(z) = f(0) \end{array}}}
\vnbprop{polynomial-holomorphic}{\vnbstmt{\begin{array}{@{}l@{}} \forall j \in \mathbb{N}, p \in (\mathbb{N} \to \mathbb{C}). \\ \quad \operatorname{holomorphic\text{-}on}(\mathbb{C}, (\lambda w \in \mathbb{C}.\; \operatorname{cc\text{-}series\text{-}partial\text{-}sum}((\lambda k \in \mathbb{N}.\; (p(k) \cdot \operatorname{power}(w, k))), j))) \end{array}}}
\vnbprop{polynomial-growth-bound}{\vnbstmt{\begin{array}{@{}l@{}} \forall j \in \mathbb{N}, p \in (\mathbb{N} \to \mathbb{C}). \\ \quad \operatorname{cc\text{-}series\text{-}partial\text{-}sum}((\lambda k \in \mathbb{N}.\; \operatorname{magnitude}(p(k))), (j + 1)) \in \mathbb{R} \wedge \\ \quad 0 \leq \operatorname{cc\text{-}series\text{-}partial\text{-}sum}((\lambda k \in \mathbb{N}.\; \operatorname{magnitude}(p(k))), (j + 1)) \wedge \\ \quad \forall w \in \mathbb{C}. \\ \quad \quad (1 \leq \operatorname{magnitude}(w)) \Rightarrow  \\ \quad \quad \quad \operatorname{magnitude}(\operatorname{cc\text{-}series\text{-}partial\text{-}sum}((\lambda k \in \mathbb{N}.\; (p(k) \cdot \operatorname{power}(w, k))), (j + 1))) \\ \quad \quad \quad \leq\; \operatorname{cc\text{-}series\text{-}partial\text{-}sum}((\lambda k \in \mathbb{N}.\; \operatorname{magnitude}(p(k))), (j + 1)) \cdot \operatorname{power}(\operatorname{magnitude}(w), j) \end{array}}}
\vnbprop{fundamental-theorem-of-algebra}{\vnbstmt{\begin{array}{@{}l@{}} \forall p \in (\mathbb{N} \to \mathbb{C}), m \in \mathbb{N}. \\ \quad \neg (p((m + 1)) = 0) \Rightarrow  \\ \quad \quad \exists z \in \mathbb{C}. \\ \quad \quad \quad \operatorname{cc\text{-}series\text{-}partial\text{-}sum}((\lambda k \in \mathbb{N}.\; (p(k) \cdot \operatorname{power}(z, k))), ((m + 1) + 1)) = 0 \end{array}}}

\vnbbookfile{theorem-library/winding-number-integer}{theorem-library/winding-number-integer}
\noindent{\small winding-number-integer.scm -- THE WINDING NUMBER IS AN INTEGER (complex-analysis.pdf ch. 3, Proposition 3.27, first half), with the two pieces of road calculus it rests on. Agent CA-5 of week 3 of the October roadmap, 2026-10-04. Definitions: structure-library/winding-number.scm (WINDING, by the integral (73)).}\par
\vnbprop{road-product-constant}{\vnbstmt{\begin{array}{@{}l@{}} \forall m, c, d, f, a, b. \\ \quad \operatorname{is\text{-}road}(m, c, a, b) \wedge \\ \quad \operatorname{is\text{-}road}(d, f, a, b) \wedge \\ \quad \forall t \in \operatorname{ccint}(a, b).\; (((c(t) \cdot d(t)) + (m(t) \cdot f(t))) = 0) \Rightarrow \\ \quad \quad (m(b) \cdot d(b)) = (m(a) \cdot d(a)) \end{array}}}
\vnbprop{line-int-primitive-road}{\vnbstmt{\begin{array}{@{}l@{}} \forall u, f, m, c, a, b. \\ \quad u \subseteq \mathbb{C} \wedge \\ \quad f \in (u \to \mathbb{C}) \wedge \\ \quad \forall y \in u. \\ \quad \quad \operatorname{is\text{-}continuous\text{-}at}(\operatorname{subspace\text{-}ms}(\operatorname{nf\text{-}metric\text{-}space}(\operatorname{cc\text{-}normed\text{-}field}), u), \\ \quad \quad \quad \operatorname{nf\text{-}metric\text{-}space}(\operatorname{cc\text{-}normed\text{-}field}), f, y) \wedge \\ \quad \operatorname{is\text{-}road}(m, c, a, b) \wedge \\ \quad \operatorname{trace}(m, a, b) \subseteq u \Rightarrow \\ \quad \quad \exists h. \\ \quad \quad \quad \operatorname{is\text{-}road}(h, (\lambda t \in \operatorname{ccint}(a, b).\; (f(m(t)) \cdot c(t))), a, b) \wedge \\ \quad \quad \quad \operatorname{line\text{-}int}(f, m, c, a, b) = (h(b) - h(a)) \end{array}}}
\vnbprop{winding-number-integer}{\vnbstmt{\begin{array}{@{}l@{}} \forall m, c, a, b.; \forall z \in \mathbb{C}. \\ \quad (\operatorname{is\text{-}road}(m, c, a, b) \wedge ((m(b) = m(a)) \wedge \neg (z \in \operatorname{trace}(m, a, b)))) \Rightarrow  \\ \quad \quad \exists k \in \mathbb{Z}.\; (\operatorname{winding}(m, c, a, b, z) = k) \end{array}}}

\vnbbookfile{theorem-library/winding-number-components}{theorem-library/winding-number-components}
\noindent{\small winding-number-components.scm -- THE WINDING NUMBER IS LOCALLY CONSTANT OFF THE TRACE, CONSTANT ON A CONNECTED SET OFF THE TRACE, AND 0 ON THE UNBOUNDED COMPONENT (complex-analysis.pdf ch. 3: Proposition 3.27, second half; Corollary 3.30). Agent CA-5 of week 3 of the October roadmap, 2026-10-04.}\par
\vnbprop{winding-disc-constant}{\vnbstmt{\begin{array}{@{}l@{}} \forall m, d, a, b.; \forall c \in \mathbb{C}, r \in \mathbb{R}, z \in \mathbb{C}. \\ \quad \operatorname{is\text{-}road}(m, d, a, b) \wedge \\ \quad m(b) = m(a) \wedge \\ \quad 0 < r \wedge \\ \quad \forall y \in \operatorname{trace}(m, a, b).\; (r \leq \operatorname{magnitude}((y - c))) \wedge \\ \quad \operatorname{magnitude}((z - c)) < r \Rightarrow \\ \quad \quad \operatorname{winding}(m, d, a, b, z) = \operatorname{winding}(m, d, a, b, c) \end{array}}}
\vnbprop{winding-locally-constant}{\vnbstmt{\begin{array}{@{}l@{}} \forall m, c, a, b.; \forall z \in \mathbb{C}. \\ \quad (\operatorname{is\text{-}road}(m, c, a, b) \wedge ((m(b) = m(a)) \wedge \neg (z \in \operatorname{trace}(m, a, b)))) \Rightarrow  \\ \quad \quad \exists r. \\ \quad \quad \quad \operatorname{pos\text{-}rr}(r) \wedge \\ \quad \quad \quad \forall \mathit{wcz} \in \operatorname{ball}(\operatorname{cc\text{-}ms}, z, r). \\ \quad \quad \quad \quad \neg (\mathit{wcz} \in \operatorname{trace}(m, a, b)) \wedge (\operatorname{winding}(m, c, a, b, \mathit{wcz}) = \operatorname{winding}(m, c, a, b, z)) \end{array}}}
\vnbprop{winding-constant-on-connected}{\vnbstmt{\begin{array}{@{}l@{}} \forall m, c, a, b, \mathit{wca}.; \forall p, q \in \mathit{wca}. \\ \quad \operatorname{is\text{-}road}(m, c, a, b) \wedge \\ \quad m(b) = m(a) \wedge \\ \quad \operatorname{is\text{-}connected}(\operatorname{cc\text{-}ms}, \mathit{wca}) \wedge \\ \quad \forall x \in \mathit{wca}.\; \neg (x \in \operatorname{trace}(m, a, b)) \Rightarrow \\ \quad \quad \operatorname{winding}(m, c, a, b, p) = \operatorname{winding}(m, c, a, b, q) \end{array}}}
\vnbprop{winding-zero-unbounded}{\vnbstmt{\begin{array}{@{}l@{}} \forall m, d, a, b.; \forall c \in \mathbb{C}, r \in \mathbb{R}, \mathit{wca}.; \forall u, p \in \mathit{wca}. \\ \quad \operatorname{is\text{-}road}(m, d, a, b) \wedge \\ \quad m(b) = m(a) \wedge \\ \quad 0 < r \wedge \\ \quad \operatorname{trace}(m, a, b) \subseteq \operatorname{ball}(\operatorname{cc\text{-}ms}, c, r) \wedge \\ \quad \operatorname{is\text{-}connected}(\operatorname{cc\text{-}ms}, \mathit{wca}) \wedge \\ \quad \forall x \in \mathit{wca}.\; \neg (x \in \operatorname{trace}(m, a, b)) \wedge \\ \quad \neg (u \in \operatorname{ball}(\operatorname{cc\text{-}ms}, c, r)) \Rightarrow \\ \quad \quad \operatorname{winding}(m, d, a, b, p) = 0 \end{array}}}
\vnbprop{winding-juxta}{\vnbstmt{\begin{array}{@{}l@{}} \forall m_{1}, d, m_{2}, f, a, b, c.; \forall z \in \mathbb{C}. \\ \quad \operatorname{is\text{-}road}(m_{1}, d, a, b) \wedge \\ \quad \operatorname{is\text{-}road}(m_{2}, f, b, c) \wedge \\ \quad \operatorname{m1}(b) = \operatorname{m2}(b) \wedge \\ \quad \neg (z \in \operatorname{trace}(m_{1}, a, b)) \wedge \\ \quad \neg (z \in \operatorname{trace}(m_{2}, b, c)) \Rightarrow \\ \quad \quad \operatorname{winding}(\operatorname{juxta}(m_{1}, m_{2}, a, b, c), \operatorname{juxta}(d, f, a, b, c), a, c, z) \\ \quad \quad =\; \operatorname{winding}(m_{1}, d, a, b, z) + \operatorname{winding}(m_{2}, f, b, c, z) \end{array}}}
\vnbprop{winding-opposite}{\vnbstmt{\begin{array}{@{}l@{}} \forall m, c, a, b, \mathit{prgb}, \mathit{prdb}.; \forall z \in \mathbb{C}. \\ \quad \operatorname{is\text{-}road}(m, c, a, b) \wedge \\ \quad \mathit{prgb} \in (\operatorname{ccint}(a, b) \to \mathbb{C}) \wedge \\ \quad \mathit{prdb} \in (\operatorname{ccint}(a, b) \to \mathbb{C}) \wedge \\ \quad \forall y \in \operatorname{ccint}(a, b).\; (\operatorname{prgb}(y) \simeq m(((a + b) - y))) \wedge \\ \quad \forall y \in \operatorname{ccint}(a, b).\; (\operatorname{prdb}(y) \simeq (-1 \cdot c(((a + b) - y)))) \wedge \\ \quad \neg (z \in \operatorname{trace}(m, a, b)) \Rightarrow \\ \quad \quad \operatorname{winding}(\mathit{prgb}, \mathit{prdb}, a, b, z) = (-1 \cdot \operatorname{winding}(m, c, a, b, z)) \end{array}}}

\vnbbookfile{theorem-library/cauchy-centre-mean-value}{theorem-library/cauchy-centre-mean-value}
\noindent{\small cauchy-centre-mean-value.scm -- THE INTEGRAL FORMULA AT THE CENTRE, THE MEAN VALUE PROPERTY, THE MEAN VALUE INEQUALITY, THE MAXIMUM MODULUS PRINCIPLE ON A DISC (complex-analysis.pdf 3.7 at z = c; 3.8: equation (92), Definition 3.54, Proposition 3.55 in its disc form).  Agent CA-4 of week 2 of the October roadmap}\par
\vnbprop{cc-continuous-at-local}{\vnbstmt{\begin{array}{@{}l@{}} \forall u.; \forall h \in (u \to \mathbb{C}), v.; \forall w \in v. \\ \quad u \subseteq \mathbb{C} \wedge \\ \quad \operatorname{is\text{-}open}(\operatorname{nf\text{-}metric\text{-}space}(\operatorname{cc\text{-}normed\text{-}field}), v) \wedge \\ \quad v \subseteq u \wedge \\ \quad \operatorname{is\text{-}continuous\text{-}at}(\operatorname{subspace\text{-}ms}(\operatorname{nf\text{-}metric\text{-}space}(\operatorname{cc\text{-}normed\text{-}field}), v), \\ \quad \quad \operatorname{nf\text{-}metric\text{-}space}(\operatorname{cc\text{-}normed\text{-}field}), \operatorname{restrict}(h, v), w) \Rightarrow \\ \quad \quad \operatorname{is\text{-}continuous\text{-}at}(\operatorname{subspace\text{-}ms}(\operatorname{nf\text{-}metric\text{-}space}(\operatorname{cc\text{-}normed\text{-}field}), u), \\ \quad \quad \quad \operatorname{nf\text{-}metric\text{-}space}(\operatorname{cc\text{-}normed\text{-}field}), h, w) \end{array}}}
\vnbprop{holomorphic-difference-quotient}{\vnbstmt{\begin{array}{@{}l@{}} \forall u, f.; \forall c \in u. \\ \quad \operatorname{holomorphic\text{-}on}(u, f) \Rightarrow  \\ \quad \quad \exists q \in (u \to \mathbb{C}). \\ \quad \quad \quad \forall y \in u. \\ \quad \quad \quad \quad \operatorname{is\text{-}continuous\text{-}at}(\operatorname{subspace\text{-}ms}(\operatorname{nf\text{-}metric\text{-}space}(\operatorname{cc\text{-}normed\text{-}field}), u), \\ \quad \quad \quad \quad \quad \operatorname{nf\text{-}metric\text{-}space}(\operatorname{cc\text{-}normed\text{-}field}), q, y) \wedge \\ \quad \quad \quad \operatorname{holomorphic\text{-}on}(\operatorname{difference}(u, \operatorname{singleton}(c)), \operatorname{restrict}(q, \operatorname{difference}(u, \operatorname{singleton}(c)))) \wedge \\ \quad \quad \quad \forall x \in u.\; ((f(x) - f(c)) = (q(x) \cdot (x - c))) \end{array}}}
\vnbprop{circle-trace-in-punctured-ball}{\vnbstmt{\begin{array}{@{}l@{}} \forall c \in \mathbb{C}, r, \mathit{cmvbr} \in \mathbb{R}, y \in \operatorname{trace}(\operatorname{circle\text{-}path}(c, r), 0, (2 \cdot \pi)). \\ \quad ((0 < r) \wedge (r < \mathit{cmvbr})) \Rightarrow ((y \in \operatorname{ball}(\operatorname{cc\text{-}ms}, c, \mathit{cmvbr})) \wedge \neg (y = c)) \end{array}}}
\vnbprop{cauchy-integral-formula-centre}{\vnbstmt{\begin{array}{@{}l@{}} \forall u, f.; \forall c \in \mathbb{C}, r, \mathit{cmvbr} \in \mathbb{R}. \\ \quad (\operatorname{holomorphic\text{-}on}(u, f) \wedge ((0 < r) \wedge ((r < \mathit{cmvbr}) \wedge (\operatorname{ball}(\operatorname{cc\text{-}ms}, c, \mathit{cmvbr}) \subseteq u)))) \Rightarrow  \\ \quad \quad f(c) \\ \quad \quad =\; \operatorname{recip}(((2 \cdot \pi) \cdot +i)) \cdot \operatorname{line\text{-}int}((\lambda w \in \operatorname{difference}(u, \operatorname{singleton}(c)).\; (f(w) \cdot \operatorname{recip}((w - c)))), \operatorname{circle\text{-}path}(c, r), \operatorname{circle\text{-}deriv}(c, r), 0, (2 \cdot \pi)) \end{array}}}
\vnbprop{holomorphic-mean-value}{\vnbstmt{\begin{array}{@{}l@{}} \forall u, f.; \forall c \in \mathbb{C}, r, \mathit{cmvbr} \in \mathbb{R}. \\ \quad (\operatorname{holomorphic\text{-}on}(u, f) \wedge ((0 < r) \wedge ((r < \mathit{cmvbr}) \wedge (\operatorname{ball}(\operatorname{cc\text{-}ms}, c, \mathit{cmvbr}) \subseteq u)))) \Rightarrow  \\ \quad \quad f(c) = (\operatorname{recip}((2 \cdot \pi)) \cdot \int_{0}^{(2 \cdot \pi)} (\lambda t \in \operatorname{ccint}(0, (2 \cdot \pi)).\; f((c + (r \cdot \exp (+i \cdot t)))))) \end{array}}}
\vnbprop{pw-int-nonneg-zero}{\vnbstmt{\begin{array}{@{}l@{}} \forall f, h, a, b.; \forall t \in \operatorname{ccint}(a, b). \\ \quad \operatorname{is\text{-}primitive}(f, h, a, b) \wedge \\ \quad \operatorname{is\text{-}continuous\text{-}on}(h, \operatorname{ccint}(a, b)) \wedge \\ \quad \forall t \in \operatorname{ccint}(a, b).\; (0 \leq h(t)) \wedge \\ \quad \int_{a}^{b} h = 0 \Rightarrow \\ \quad \quad h(t) = 0 \end{array}}}
\vnbprop{holomorphic-circle-max-constant}{\vnbstmt{\begin{array}{@{}l@{}} \forall u, f.; \forall c \in \mathbb{C}, r, \mathit{cmvbr} \in \mathbb{R}, t \in \operatorname{ccint}(0, (2 \cdot \pi)). \\ \quad \operatorname{holomorphic\text{-}on}(u, f) \wedge \\ \quad 0 < r \wedge \\ \quad r < \mathit{cmvbr} \wedge \\ \quad \operatorname{ball}(\operatorname{cc\text{-}ms}, c, \mathit{cmvbr}) \subseteq u \wedge \\ \quad \forall y \in \operatorname{trace}(\operatorname{circle\text{-}path}(c, r), 0, (2 \cdot \pi)). \\ \quad \quad \operatorname{magnitude}(f(y)) \leq \operatorname{magnitude}(f(c)) \Rightarrow \\ \quad \quad f((c + (r \cdot \exp (+i \cdot t)))) = f(c) \end{array}}}
\vnbprop{maximum-modulus-disc}{\vnbstmt{\begin{array}{@{}l@{}} \forall u, f.; \forall c \in \mathbb{C}, r \in \mathbb{R}, z \in \operatorname{ball}(\operatorname{cc\text{-}ms}, c, r). \\ \quad \operatorname{holomorphic\text{-}on}(u, f) \wedge \\ \quad 0 < r \wedge \\ \quad \operatorname{ball}(\operatorname{cc\text{-}ms}, c, r) \subseteq u \wedge \\ \quad \forall z \in \operatorname{ball}(\operatorname{cc\text{-}ms}, c, r).\; (\operatorname{magnitude}(f(z)) \leq \operatorname{magnitude}(f(c))) \Rightarrow \\ \quad \quad f(z) = f(c) \end{array}}}
\vnbprop{holomorphic-mean-value-bound}{\vnbstmt{\begin{array}{@{}l@{}} \forall u, f.; \forall c \in \mathbb{C}, r, \mathit{cmvbr}, m \in \mathbb{R}. \\ \quad \operatorname{holomorphic\text{-}on}(u, f) \wedge \\ \quad 0 < r \wedge \\ \quad r < \mathit{cmvbr} \wedge \\ \quad \operatorname{ball}(\operatorname{cc\text{-}ms}, c, \mathit{cmvbr}) \subseteq u \wedge \\ \quad \forall y \in \operatorname{trace}(\operatorname{circle\text{-}path}(c, r), 0, (2 \cdot \pi)).\; (\operatorname{magnitude}(f(y)) \leq m) \Rightarrow \\ \quad \quad \operatorname{magnitude}(f(c)) \leq m \end{array}}}

\vnbbookfile{theorem-library/morera}{theorem-library/morera}
\noindent{\small morera.scm -- MORERA'S THEOREM AND LOCALLY UNIFORM LIMITS OF HOLOMORPHIC FUNCTIONS (complex-analysis.pdf 3.6: Proposition 3.43, Proposition 3.44), both CONDITIONAL on "the derivative of a holomorphic function is holomorphic" (Corollary 3.37). Agent CA-4 of week 2 of the October roadmap (scratchpad/triage/WEEK2.md), ...}\par
\vnbprop{diff-on-local}{\vnbstmt{\begin{array}{@{}l@{}} \forall u.; \forall f \in (u \to \mathbb{C}), v.; \forall a \in v, l. \\ \quad \operatorname{is\text{-}open}(\operatorname{nf\text{-}metric\text{-}space}(\operatorname{cc\text{-}normed\text{-}field}), u) \wedge \\ \quad \operatorname{is\text{-}open}(\operatorname{nf\text{-}metric\text{-}space}(\operatorname{cc\text{-}normed\text{-}field}), v) \wedge \\ \quad v \subseteq u \wedge \\ \quad \operatorname{is\text{-}diff\text{-}on}(\operatorname{cc\text{-}normed\text{-}field}, v, \operatorname{restrict}(f, v), a, l) \Rightarrow \\ \quad \quad \operatorname{is\text{-}diff\text{-}on}(\operatorname{cc\text{-}normed\text{-}field}, u, f, a, l) \end{array}}}
\vnbprop{cauchy-convex-segments-of-triangles}{\vnbstmt{\begin{array}{@{}l@{}} \forall u, f.; \forall a, b, c \in u. \\ \quad \operatorname{is\text{-}convex\text{-}cc}(u) \wedge \\ \quad \operatorname{is\text{-}open}(\operatorname{nf\text{-}metric\text{-}space}(\operatorname{cc\text{-}normed\text{-}field}), u) \wedge \\ \quad f \in (u \to \mathbb{C}) \wedge \\ \quad \forall y \in u. \\ \quad \quad \operatorname{is\text{-}continuous\text{-}at}(\operatorname{subspace\text{-}ms}(\operatorname{nf\text{-}metric\text{-}space}(\operatorname{cc\text{-}normed\text{-}field}), u), \\ \quad \quad \quad \operatorname{nf\text{-}metric\text{-}space}(\operatorname{cc\text{-}normed\text{-}field}), f, y) \wedge \\ \quad \forall \mathit{mrta}, \mathit{mrtb}, \mathit{mrtc} \in u. \\ \quad \quad (\operatorname{conv3}(\mathit{mrta}, \mathit{mrtb}, \mathit{mrtc}) \subseteq u) \Rightarrow (\operatorname{tri\text{-}int}(f, \mathit{mrta}, \mathit{mrtb}, \mathit{mrtc}) = 0) \Rightarrow \\ \quad \quad \operatorname{seg\text{-}int}(f, a, c) = (\operatorname{seg\text{-}int}(f, a, b) + \operatorname{seg\text{-}int}(f, b, c)) \end{array}}}
\vnbprop{cauchy-convex-primitive-of-triangles}{\vnbstmt{\begin{array}{@{}l@{}} \forall u, f.; \forall o \in u. \\ \quad \operatorname{is\text{-}convex\text{-}cc}(u) \wedge \\ \quad \operatorname{is\text{-}open}(\operatorname{nf\text{-}metric\text{-}space}(\operatorname{cc\text{-}normed\text{-}field}), u) \wedge \\ \quad f \in (u \to \mathbb{C}) \wedge \\ \quad \forall y \in u. \\ \quad \quad \operatorname{is\text{-}continuous\text{-}at}(\operatorname{subspace\text{-}ms}(\operatorname{nf\text{-}metric\text{-}space}(\operatorname{cc\text{-}normed\text{-}field}), u), \\ \quad \quad \quad \operatorname{nf\text{-}metric\text{-}space}(\operatorname{cc\text{-}normed\text{-}field}), f, y) \wedge \\ \quad \forall a, b, c \in u. \\ \quad \quad (\operatorname{conv3}(a, b, c) \subseteq u) \Rightarrow (\operatorname{tri\text{-}int}(f, a, b, c) = 0) \Rightarrow \\ \quad \quad \exists g \in (u \to \mathbb{C}). \\ \quad \quad \quad (\operatorname{holomorphic\text{-}on}(u, g) \wedge \forall z \in u.\; \operatorname{is\text{-}diff\text{-}on}(\operatorname{cc\text{-}normed\text{-}field}, u, g, z, f(z))) \end{array}}}
\vnbprop{morera-conditional}{\vnbstmt{\begin{array}{@{}l@{}} \forall u.; \forall f \in (u \to \mathbb{C}). \\ \quad \forall v, g, d. \\ \quad \quad \operatorname{holomorphic\text{-}on}(v, g) \wedge \\ \quad \quad d \in (v \to \mathbb{C}) \wedge \\ \quad \quad \forall z \in v.\; \operatorname{is\text{-}diff\text{-}on}(\operatorname{cc\text{-}normed\text{-}field}, v, g, z, d(z)) \Rightarrow \\ \quad \quad \quad \operatorname{holomorphic\text{-}on}(v, d) \wedge \\ \quad \operatorname{is\text{-}open}(\operatorname{nf\text{-}metric\text{-}space}(\operatorname{cc\text{-}normed\text{-}field}), u) \wedge \\ \quad \forall y \in u. \\ \quad \quad \operatorname{is\text{-}continuous\text{-}at}(\operatorname{subspace\text{-}ms}(\operatorname{nf\text{-}metric\text{-}space}(\operatorname{cc\text{-}normed\text{-}field}), u), \\ \quad \quad \quad \operatorname{nf\text{-}metric\text{-}space}(\operatorname{cc\text{-}normed\text{-}field}), f, y) \wedge \\ \quad \forall a, b, c \in u. \\ \quad \quad (\operatorname{conv3}(a, b, c) \subseteq u) \Rightarrow (\operatorname{tri\text{-}int}(f, a, b, c) = 0) \Rightarrow \\ \quad \quad \operatorname{holomorphic\text{-}on}(u, f) \end{array}}}
\vnbprop{holomorphic-uniform-limit-conditional}{\vnbstmt{\begin{array}{@{}l@{}} \forall u, m.; \forall f \in (u \to \mathbb{C}). \\ \quad \forall v, g, d. \\ \quad \quad \operatorname{holomorphic\text{-}on}(v, g) \wedge \\ \quad \quad d \in (v \to \mathbb{C}) \wedge \\ \quad \quad \forall z \in v.\; \operatorname{is\text{-}diff\text{-}on}(\operatorname{cc\text{-}normed\text{-}field}, v, g, z, d(z)) \Rightarrow \\ \quad \quad \quad \operatorname{holomorphic\text{-}on}(v, d) \wedge \\ \quad \operatorname{is\text{-}open}(\operatorname{nf\text{-}metric\text{-}space}(\operatorname{cc\text{-}normed\text{-}field}), u) \wedge \\ \quad \forall k \in \mathbb{N}.\; \operatorname{holomorphic\text{-}on}(u, m(k)) \wedge \\ \quad \forall \mathit{mrlz} \in u. \exists r. \\ \quad \quad \operatorname{pos\text{-}rr}(r) \wedge \\ \quad \quad \operatorname{ball}(\operatorname{cc\text{-}ms}, \mathit{mrlz}, r) \subseteq u \wedge \\ \quad \quad \forall s. \\ \quad \quad \quad \operatorname{pos\text{-}rr}(s) \Rightarrow  \\ \quad \quad \quad \quad \exists p \in \mathbb{N}. \forall k \in \mathbb{N}, y \in \operatorname{ball}(\operatorname{cc\text{-}ms}, \mathit{mrlz}, r). \\ \quad \quad \quad \quad \quad (p \leq k) \Rightarrow (\operatorname{magnitude}((m(k)(y) - f(y))) \leq s) \Rightarrow \\ \quad \quad \operatorname{holomorphic\text{-}on}(u, f) \end{array}}}
\vnbprop{deriv-holomorphic-bridge}{\vnbstmt{\begin{array}{@{}l@{}} \forall v, g, d. \\ \quad \forall u, f. \\ \quad \quad \operatorname{holomorphic\text{-}on}(u, f) \Rightarrow  \\ \quad \quad \quad \operatorname{holomorphic\text{-}on}(u, (\lambda w \in u.\; \operatorname{deriv\text{-}on}(\operatorname{cc\text{-}normed\text{-}field}, u, f, w))) \wedge \\ \quad \operatorname{holomorphic\text{-}on}(v, g) \wedge \\ \quad d \in (v \to \mathbb{C}) \wedge \\ \quad \forall z \in v.\; \operatorname{is\text{-}diff\text{-}on}(\operatorname{cc\text{-}normed\text{-}field}, v, g, z, d(z)) \Rightarrow \\ \quad \quad \operatorname{holomorphic\text{-}on}(v, d) \end{array}}}
\vnbprop{morera-of-deriv-holomorphic}{\vnbstmt{\begin{array}{@{}l@{}} \forall u.; \forall f \in (u \to \mathbb{C}). \\ \quad \forall u, f. \\ \quad \quad \operatorname{holomorphic\text{-}on}(u, f) \Rightarrow  \\ \quad \quad \quad \operatorname{holomorphic\text{-}on}(u, (\lambda w \in u.\; \operatorname{deriv\text{-}on}(\operatorname{cc\text{-}normed\text{-}field}, u, f, w))) \wedge \\ \quad \operatorname{is\text{-}open}(\operatorname{nf\text{-}metric\text{-}space}(\operatorname{cc\text{-}normed\text{-}field}), u) \wedge \\ \quad \forall y \in u. \\ \quad \quad \operatorname{is\text{-}continuous\text{-}at}(\operatorname{subspace\text{-}ms}(\operatorname{nf\text{-}metric\text{-}space}(\operatorname{cc\text{-}normed\text{-}field}), u), \\ \quad \quad \quad \operatorname{nf\text{-}metric\text{-}space}(\operatorname{cc\text{-}normed\text{-}field}), f, y) \wedge \\ \quad \forall a, b, c \in u. \\ \quad \quad (\operatorname{conv3}(a, b, c) \subseteq u) \Rightarrow (\operatorname{tri\text{-}int}(f, a, b, c) = 0) \Rightarrow \\ \quad \quad \operatorname{holomorphic\text{-}on}(u, f) \end{array}}}
\vnbprop{holomorphic-uniform-limit-of-deriv-holomorphic}{\vnbstmt{\begin{array}{@{}l@{}} \forall u, m.; \forall f \in (u \to \mathbb{C}). \\ \quad \forall u, f. \\ \quad \quad \operatorname{holomorphic\text{-}on}(u, f) \Rightarrow  \\ \quad \quad \quad \operatorname{holomorphic\text{-}on}(u, (\lambda w \in u.\; \operatorname{deriv\text{-}on}(\operatorname{cc\text{-}normed\text{-}field}, u, f, w))) \wedge \\ \quad \operatorname{is\text{-}open}(\operatorname{nf\text{-}metric\text{-}space}(\operatorname{cc\text{-}normed\text{-}field}), u) \wedge \\ \quad \forall k \in \mathbb{N}.\; \operatorname{holomorphic\text{-}on}(u, m(k)) \wedge \\ \quad \forall z \in u. \exists r. \\ \quad \quad \operatorname{pos\text{-}rr}(r) \wedge \\ \quad \quad \operatorname{ball}(\operatorname{cc\text{-}ms}, z, r) \subseteq u \wedge \\ \quad \quad \forall s. \\ \quad \quad \quad \operatorname{pos\text{-}rr}(s) \Rightarrow  \\ \quad \quad \quad \quad \exists p \in \mathbb{N}. \forall k \in \mathbb{N}, y \in \operatorname{ball}(\operatorname{cc\text{-}ms}, z, r). \\ \quad \quad \quad \quad \quad (p \leq k) \Rightarrow (\operatorname{magnitude}((m(k)(y) - f(y))) \leq s) \Rightarrow \\ \quad \quad \operatorname{holomorphic\text{-}on}(u, f) \end{array}}}
\vnbprop{morera}{\vnbstmt{\begin{array}{@{}l@{}} \forall u.; \forall f \in (u \to \mathbb{C}). \\ \quad \operatorname{is\text{-}open}(\operatorname{nf\text{-}metric\text{-}space}(\operatorname{cc\text{-}normed\text{-}field}), u) \wedge \\ \quad \forall y \in u. \\ \quad \quad \operatorname{is\text{-}continuous\text{-}at}(\operatorname{subspace\text{-}ms}(\operatorname{nf\text{-}metric\text{-}space}(\operatorname{cc\text{-}normed\text{-}field}), u), \\ \quad \quad \quad \operatorname{nf\text{-}metric\text{-}space}(\operatorname{cc\text{-}normed\text{-}field}), f, y) \wedge \\ \quad \forall a, b, c \in u. \\ \quad \quad (\operatorname{conv3}(a, b, c) \subseteq u) \Rightarrow (\operatorname{tri\text{-}int}(f, a, b, c) = 0) \Rightarrow \\ \quad \quad \operatorname{holomorphic\text{-}on}(u, f) \end{array}}}
\vnbprop{holomorphic-uniform-limit}{\vnbstmt{\begin{array}{@{}l@{}} \forall u, m.; \forall f \in (u \to \mathbb{C}). \\ \quad \operatorname{is\text{-}open}(\operatorname{nf\text{-}metric\text{-}space}(\operatorname{cc\text{-}normed\text{-}field}), u) \wedge \\ \quad \forall k \in \mathbb{N}.\; \operatorname{holomorphic\text{-}on}(u, m(k)) \wedge \\ \quad \forall z \in u. \exists r. \\ \quad \quad \operatorname{pos\text{-}rr}(r) \wedge \\ \quad \quad \operatorname{ball}(\operatorname{cc\text{-}ms}, z, r) \subseteq u \wedge \\ \quad \quad \forall s. \\ \quad \quad \quad \operatorname{pos\text{-}rr}(s) \Rightarrow  \\ \quad \quad \quad \quad \exists p \in \mathbb{N}. \forall k \in \mathbb{N}, y \in \operatorname{ball}(\operatorname{cc\text{-}ms}, z, r). \\ \quad \quad \quad \quad \quad (p \leq k) \Rightarrow (\operatorname{magnitude}((m(k)(y) - f(y))) \leq s) \Rightarrow \\ \quad \quad \operatorname{holomorphic\text{-}on}(u, f) \end{array}}}

\vnbbookfile{structure-library/gamma-function}{structure-library/gamma-function}
\noindent{\small gamma-function.scm -- THE IMPROPER INTEGRAL ON (0, +inf) AND THE GAMMA FUNCTION (complex-analysis.pdf ch. 3, section 3.9: equations (96)-(99), Definition 3.58). DEFINITIONS ONLY; every law is PROVEN in theorem-library/gamma-function.scm. Agent CA-6 of week 3 of the October roadmap (scratchpad/triage/WEEK3.md), ...}\par
\par\medskip\noindent\textbf{Definition} (\texttt{improper-int-converges-to}).\;For all $\operatorname{ipf\_}$ and $\operatorname{ipl\_}$, 
\begin{equation*}
\begin{array}{@{}l@{}} \operatorname{improper\text{-}int\text{-}converges\text{-}to}(\operatorname{ipf\_}, \operatorname{ipl\_}) \\ \Leftrightarrow\; \operatorname{ipl\_} \in \mathbb{C} \wedge \\ \forall \operatorname{ipe\_}. \\ \quad \operatorname{pos\text{-}rr}(\operatorname{ipe\_}) \Rightarrow  \\ \quad \quad \exists \operatorname{ipd\_}. \\ \quad \quad \quad \operatorname{pos\text{-}rr}(\operatorname{ipd\_}) \wedge \\ \quad \quad \quad \exists \operatorname{ipm\_} \in \mathbb{R}. \forall \operatorname{ipr\_}, \operatorname{ipq\_} \in \mathbb{R}. \\ \quad \quad \quad \quad (0 < \operatorname{ipr\_}) \Rightarrow  \\ \quad \quad \quad \quad (\operatorname{ipr\_} < \operatorname{ipd\_}) \Rightarrow  \\ \quad \quad \quad \quad (\operatorname{ipm\_} < \operatorname{ipq\_}) \Rightarrow  \\ \quad \quad \quad \quad (\operatorname{ipr\_} < \operatorname{ipq\_}) \Rightarrow (\operatorname{magnitude}(-(\int_{\operatorname{ipr\_}}^{\operatorname{ipq\_}} \operatorname{ipf\_}, \operatorname{ipl\_})) \leq \operatorname{ipe\_}) \end{array}
\end{equation*}

\par\medskip\noindent\textbf{Definition} (\texttt{improper-int}).\;$\int_0^{\infty} \operatorname{ipf\_} :=$ $(\iota\, \operatorname{ipv\_}.\;\operatorname{improper\text{-}int\text{-}converges\text{-}to}(\operatorname{ipf\_}, \operatorname{ipv\_}))$
\par\medskip\noindent\textbf{Definition} (\texttt{gamma-integrand}).\;$\operatorname{gamma\text{-}integrand}(\operatorname{gmz\_}) :=$
\begin{equation*}
\begin{array}{@{}l@{}} (\lambda \operatorname{gmw\_} \in \{\operatorname{gmu\_} \in \mathbb{C} : (0 < \operatorname{real\text{-}part}(\operatorname{gmu\_}))\}.\; \\ \quad \exp ((\operatorname{gmz\_} - 1) \cdot \log\left(\operatorname{gmw\_}\right)) \cdot \exp -\operatorname{gmw\_}) \end{array}
\end{equation*}

\par\medskip\noindent\textbf{Definition} (\texttt{gamma-function}).\;$\Gamma(\operatorname{gmz\_}) :=$ $\int_0^{\infty} \operatorname{gamma\text{-}integrand}(\operatorname{gmz\_})$

\vnbbookfile{theorem-library/gamma-function}{theorem-library/gamma-function}
\noindent{\small gamma-function.scm -- THE IMPROPER INTEGRAL ON (0, +inf), THE GAMMA FUNCTION, ITS CONVERGENCE ON THE RIGHT HALF-PLANE, ITS HOLOMORPHY (PROPOSITION 3.59), GAMMA(1) = 1 AND GAMMA(z + 1) = z GAMMA(z), GAMMA(n + 1) = n!. (complex-analysis.pdf ch. 3, section 3.9: (96)-(101), Definition 3.58, Proposition 3.59.)}\par
\vnbprop{right-half-plane-open}{\vnbstmt{\operatorname{is\text{-}open}(\operatorname{nf\text{-}metric\text{-}space}(\operatorname{cc\text{-}normed\text{-}field}), \{u \in \mathbb{C} : (0 < \operatorname{real\text{-}part}(u))\})}}
\vnbprop{real-interval-is-road}{\vnbstmt{\begin{array}{@{}l@{}} \forall a, b \in \mathbb{R}. \\ \quad (a < b) \Rightarrow  \\ \quad \quad \operatorname{is\text{-}road}((\lambda t \in \operatorname{ccint}(a, b).\; t), (\lambda t \in \operatorname{ccint}(a, b).\; 1), a, b) \end{array}}}
\vnbprop{real-interval-trace}{\vnbstmt{\begin{array}{@{}l@{}} \forall a, b, u. \\ \quad ((a \in \mathbb{R}) \wedge ((b \in \mathbb{R}) \wedge ((a < b) \wedge (\operatorname{ccint}(a, b) \subseteq u)))) \Rightarrow  \\ \quad \quad \operatorname{trace}((\lambda t \in \operatorname{ccint}(a, b).\; t), a, b) \subseteq u \end{array}}}
\vnbprop{cc-int-as-line-int}{\vnbstmt{\begin{array}{@{}l@{}} \forall u, f, a, b. \\ \quad u \subseteq \mathbb{C} \wedge \\ \quad f \in (u \to \mathbb{C}) \wedge \\ \quad \forall y \in u. \\ \quad \quad \operatorname{is\text{-}continuous\text{-}at}(\operatorname{subspace\text{-}ms}(\operatorname{nf\text{-}metric\text{-}space}(\operatorname{cc\text{-}normed\text{-}field}), u), \\ \quad \quad \quad \operatorname{nf\text{-}metric\text{-}space}(\operatorname{cc\text{-}normed\text{-}field}), f, y) \wedge \\ \quad a \in \mathbb{R} \wedge \\ \quad b \in \mathbb{R} \wedge \\ \quad a < b \wedge \\ \quad \operatorname{ccint}(a, b) \subseteq u \Rightarrow \\ \quad \quad \int_{a}^{b} f = \operatorname{line\text{-}int}(f, (\lambda t \in \operatorname{ccint}(a, b).\; t), (\lambda t \in \operatorname{ccint}(a, b).\; 1), a, b) \end{array}}}
\vnbprop{cc-int-in-cc-of-continuous}{\vnbstmt{\begin{array}{@{}l@{}} \forall u, f, a, b. \\ \quad u \subseteq \mathbb{C} \wedge \\ \quad f \in (u \to \mathbb{C}) \wedge \\ \quad \forall y \in u. \\ \quad \quad \operatorname{is\text{-}continuous\text{-}at}(\operatorname{subspace\text{-}ms}(\operatorname{nf\text{-}metric\text{-}space}(\operatorname{cc\text{-}normed\text{-}field}), u), \\ \quad \quad \quad \operatorname{nf\text{-}metric\text{-}space}(\operatorname{cc\text{-}normed\text{-}field}), f, y) \wedge \\ \quad a \in \mathbb{R} \wedge \\ \quad b \in \mathbb{R} \wedge \\ \quad a < b \wedge \\ \quad \operatorname{ccint}(a, b) \subseteq u \Rightarrow \\ \quad \quad \int_{a}^{b} f \in \mathbb{C} \end{array}}}
\vnbprop{cc-int-split}{\vnbstmt{\begin{array}{@{}l@{}} \forall u, f, a, b.; \forall c \in \mathbb{R}. \\ \quad u \subseteq \mathbb{C} \wedge \\ \quad f \in (u \to \mathbb{C}) \wedge \\ \quad \forall y \in u. \\ \quad \quad \operatorname{is\text{-}continuous\text{-}at}(\operatorname{subspace\text{-}ms}(\operatorname{nf\text{-}metric\text{-}space}(\operatorname{cc\text{-}normed\text{-}field}), u), \\ \quad \quad \quad \operatorname{nf\text{-}metric\text{-}space}(\operatorname{cc\text{-}normed\text{-}field}), f, y) \wedge \\ \quad a \in \mathbb{R} \wedge \\ \quad b \in \mathbb{R} \wedge \\ \quad a < b \wedge \\ \quad \operatorname{ccint}(a, b) \subseteq u \wedge \\ \quad a < c \wedge \\ \quad c < b \Rightarrow \\ \quad \quad \int_{a}^{b} f = (\int_{a}^{c} f + \int_{c}^{b} f) \end{array}}}
\vnbprop{cc-int-compare-primitive}{\vnbstmt{\begin{array}{@{}l@{}} \forall u, f, r, c, a, b. \\ \quad u \subseteq \mathbb{C} \wedge \\ \quad f \in (u \to \mathbb{C}) \wedge \\ \quad \forall y \in u. \\ \quad \quad \operatorname{is\text{-}continuous\text{-}at}(\operatorname{subspace\text{-}ms}(\operatorname{nf\text{-}metric\text{-}space}(\operatorname{cc\text{-}normed\text{-}field}), u), \\ \quad \quad \quad \operatorname{nf\text{-}metric\text{-}space}(\operatorname{cc\text{-}normed\text{-}field}), f, y) \wedge \\ \quad a \in \mathbb{R} \wedge \\ \quad b \in \mathbb{R} \wedge \\ \quad a < b \wedge \\ \quad \operatorname{ccint}(a, b) \subseteq u \wedge \\ \quad \operatorname{holomorphic\text{-}on}(u, c) \wedge \\ \quad r \in (u \to \mathbb{C}) \wedge \\ \quad \forall y \in u. \\ \quad \quad \operatorname{is\text{-}continuous\text{-}at}(\operatorname{subspace\text{-}ms}(\operatorname{nf\text{-}metric\text{-}space}(\operatorname{cc\text{-}normed\text{-}field}), u), \\ \quad \quad \quad \operatorname{nf\text{-}metric\text{-}space}(\operatorname{cc\text{-}normed\text{-}field}), r, y) \wedge \\ \quad \forall z \in u.\; \operatorname{is\text{-}diff\text{-}on}(\operatorname{cc\text{-}normed\text{-}field}, u, c, z, r(z)) \wedge \\ \quad \forall \mathit{gmy} \in \operatorname{ccint}(a, b).\; (\operatorname{magnitude}(f(\mathit{gmy})) \leq \operatorname{real\text{-}part}(r(\mathit{gmy}))) \Rightarrow \\ \quad \quad \operatorname{magnitude}(\int_{a}^{b} f) \leq \operatorname{real\text{-}part}((c(b) - c(a))) \end{array}}}
\vnbprop{right-half-plane-in-slit}{\vnbstmt{\begin{array}{@{}l@{}} \forall y \in \{u \in \mathbb{C} : (0 < \operatorname{real\text{-}part}(u))\}. \\ \quad \neg ((\operatorname{imag\text{-}part}(y) = 0) \wedge (\operatorname{real\text{-}part}(y) \leq 0)) \wedge \\ \quad \neg (y = 0) \wedge \\ \quad y \in \{\mathit{cxsu} \in \mathbb{C} : \neg ((\operatorname{imag\text{-}part}(\mathit{cxsu}) = 0) \wedge (\operatorname{real\text{-}part}(\mathit{cxsu}) \leq 0))\} \end{array}}}
\vnbprop{gamma-integrand-holomorphic}{\vnbstmt{\begin{array}{@{}l@{}} \forall z \in \mathbb{C}. \\ \quad \operatorname{holomorphic\text{-}on}(\{u \in \mathbb{C} : (0 < \operatorname{real\text{-}part}(u))\}, \operatorname{gamma\text{-}integrand}(z)) \end{array}}}
\vnbprop{rr-positive-below-two}{\vnbstmt{\begin{array}{@{}l@{}} \forall d, r \in \mathbb{R}. \\ \quad ((0 < d) \wedge (0 < r)) \Rightarrow \exists x \in \mathbb{R}.\; ((0 < x) \wedge ((x < d) \wedge (x < r))) \end{array}}}
\vnbprop{rr-above-two}{\vnbstmt{\forall m, r \in \mathbb{R}.\; \exists x \in \mathbb{R}.\; ((m < x) \wedge (r < x))}}
\vnbprop{improper-int-cauchy}{\vnbstmt{\begin{array}{@{}l@{}} \forall f. \\ \quad \forall r, q \in \mathbb{R}.\; (((0 < r) \wedge (r < q)) \Rightarrow (\int_{r}^{q} f \in \mathbb{C})) \wedge \\ \quad \forall e. \\ \quad \quad \operatorname{pos\text{-}rr}(e) \Rightarrow  \\ \quad \quad \quad \exists d. \\ \quad \quad \quad \quad \operatorname{pos\text{-}rr}(d) \wedge \\ \quad \quad \quad \quad \exists m \in \mathbb{R}. \forall r_{1}, r_{2}, q_{1}, q_{2}. \\ \quad \quad \quad \quad \quad r_{1} \in \mathbb{R} \wedge \\ \quad \quad \quad \quad \quad r_{2} \in \mathbb{R} \wedge \\ \quad \quad \quad \quad \quad q_{1} \in \mathbb{R} \wedge \\ \quad \quad \quad \quad \quad q_{2} \in \mathbb{R} \wedge \\ \quad \quad \quad \quad \quad 0 < r_{1} \wedge \\ \quad \quad \quad \quad \quad r_{1} < d \wedge \\ \quad \quad \quad \quad \quad 0 < r_{2} \wedge \\ \quad \quad \quad \quad \quad r_{2} < d \wedge \\ \quad \quad \quad \quad \quad m < q_{1} \wedge \\ \quad \quad \quad \quad \quad m < q_{2} \wedge \\ \quad \quad \quad \quad \quad r_{1} < q_{1} \wedge \\ \quad \quad \quad \quad \quad r_{2} < q_{2} \Rightarrow \\ \quad \quad \quad \quad \quad \quad \operatorname{magnitude}((\int_{r_{1}}^{q_{1}} f - \int_{r_{2}}^{q_{2}} f)) \leq e \Rightarrow \\ \quad \quad \exists l.\; \operatorname{improper\text{-}int\text{-}converges\text{-}to}(f, l) \end{array}}}
\vnbprop{improper-int-tail-bound}{\vnbstmt{\begin{array}{@{}l@{}} \forall f, l, r, q, t. \\ \quad \forall \mathit{ipr}, \mathit{ipq} \in \mathbb{R}. \\ \quad \quad ((0 < \mathit{ipr}) \wedge (\mathit{ipr} < \mathit{ipq})) \Rightarrow (\int_{\mathit{ipr}}^{\mathit{ipq}} f \in \mathbb{C}) \wedge \\ \quad \operatorname{improper\text{-}int\text{-}converges\text{-}to}(f, l) \wedge \\ \quad r \in \mathbb{R} \wedge \\ \quad q \in \mathbb{R} \wedge \\ \quad t \in \mathbb{R} \wedge \\ \quad 0 < r \wedge \\ \quad r < q \wedge \\ \quad \forall r_{2}, q_{2} \in \mathbb{R}. \\ \quad \quad ((0 < r_{2}) \wedge ((r_{2} < r) \wedge (q < q_{2}))) \Rightarrow  \\ \quad \quad \quad \operatorname{magnitude}((\int_{r_{2}}^{q_{2}} f - \int_{r}^{q} f)) \leq t \Rightarrow \\ \quad \quad \operatorname{magnitude}((l - \int_{r}^{q} f)) \leq t \end{array}}}
\vnbprop{improper-int-unique}{\vnbstmt{\begin{array}{@{}l@{}} \forall f, l_{1}, l_{2}. \\ \quad \forall r, q \in \mathbb{R}.\; (((0 < r) \wedge (r < q)) \Rightarrow (\int_{r}^{q} f \in \mathbb{C})) \wedge \\ \quad \operatorname{improper\text{-}int\text{-}converges\text{-}to}(f, l_{1}) \wedge \\ \quad \operatorname{improper\text{-}int\text{-}converges\text{-}to}(f, l_{2}) \Rightarrow \\ \quad \quad l_{1} = l_{2} \end{array}}}
\vnbprop{improper-int-value}{\vnbstmt{\begin{array}{@{}l@{}} \forall f, l. \\ \quad \forall r, q \in \mathbb{R}.\; (((0 < r) \wedge (r < q)) \Rightarrow (\int_{r}^{q} f \in \mathbb{C})) \wedge \\ \quad \operatorname{improper\text{-}int\text{-}converges\text{-}to}(f, l) \Rightarrow \\ \quad \quad \int_0^{\infty} f = l \end{array}}}
\vnbprop{cc-exp-magnitude}{\vnbstmt{\forall v \in \mathbb{C}.\; (\operatorname{magnitude}(\exp v) = \exp \operatorname{real\text{-}part}(v))}}
\vnbprop{cc-exp-real-mono}{\vnbstmt{\forall u, v \in \mathbb{R}.\; ((u \leq v) \Rightarrow (\exp u \leq \exp v))}}
\vnbprop{cc-exp-log-real}{\vnbstmt{\forall x \in \mathbb{R}.\; ((0 < x) \Rightarrow (\exp \log\left(x\right) = x))}}
\vnbprop{log-cc-exp-real}{\vnbstmt{\forall x \in \mathbb{R}.\; (\log\left(\exp x\right) = x)}}
\vnbprop{gamma-integrand-magnitude}{\vnbstmt{\begin{array}{@{}l@{}} \forall z \in \mathbb{C}, t \in \mathbb{R}. \\ \quad (0 < t) \Rightarrow  \\ \quad \quad \operatorname{magnitude}(\operatorname{gamma\text{-}integrand}(z)(t)) = \exp (((\operatorname{real\text{-}part}(z) - 1) \cdot \log\left(t\right)) - t) \end{array}}}
\vnbprop{log-nonneg-of-ge-one}{\vnbstmt{\forall t \in \mathbb{R}.\; ((1 \leq t) \Rightarrow (0 \leq \log\left(t\right)))}}
\vnbprop{gamma-integrand-bound-small}{\vnbstmt{\begin{array}{@{}l@{}} \forall z, a, t. \\ \quad ((z \in \mathbb{C}) \wedge ((a \in \mathbb{R}) \wedge ((t \in \mathbb{R}) \wedge ((0 < a) \wedge ((a \leq \operatorname{real\text{-}part}(z)) \wedge ((0 < t) \wedge (t \leq 1))))))) \Rightarrow  \\ \quad \quad \operatorname{magnitude}(\operatorname{gamma\text{-}integrand}(z)(t)) \leq \exp ((a - 1) \cdot \log\left(t\right)) \end{array}}}
\vnbprop{gamma-integrand-bound-large}{\vnbstmt{\begin{array}{@{}l@{}} \forall z, b, a, t. \\ \quad z \in \mathbb{C} \wedge \\ \quad b \in \mathbb{R} \wedge \\ \quad a \in \mathbb{R} \wedge \\ \quad t \in \mathbb{R} \wedge \\ \quad \operatorname{real\text{-}part}(z) \leq b \wedge \\ \quad 1 \leq a \wedge \\ \quad (b - 1) \leq a \wedge \\ \quad 1 \leq t \Rightarrow \\ \quad \quad \operatorname{magnitude}(\operatorname{gamma\text{-}integrand}(z)(t)) \leq \exp ((a \cdot (\log\left((2 \cdot a)\right) - 1)) - (1/2 \cdot t)) \end{array}}}
\vnbprop{ccint-in-right-half-plane}{\vnbstmt{\begin{array}{@{}l@{}} \forall a, b \in \mathbb{R}. \\ \quad (0 < a) \Rightarrow (\operatorname{ccint}(a, b) \subseteq \{u \in \mathbb{C} : (0 < \operatorname{real\text{-}part}(u))\}) \end{array}}}
\vnbprop{gamma-integrand-locint}{\vnbstmt{\begin{array}{@{}l@{}} \forall z \in \mathbb{C}, r, q \in \mathbb{R}. \\ \quad ((0 < r) \wedge (r < q)) \Rightarrow (\int_{r}^{q} \operatorname{gamma\text{-}integrand}(z) \in \mathbb{C}) \end{array}}}
\vnbprop{gamma-tail-small}{\vnbstmt{\begin{array}{@{}l@{}} \forall z, a, b, c. \\ \quad z \in \mathbb{C} \wedge \\ \quad a \in \mathbb{R} \wedge \\ \quad b \in \mathbb{R} \wedge \\ \quad c \in \mathbb{R} \wedge \\ \quad 0 < a \wedge \\ \quad a \leq \operatorname{real\text{-}part}(z) \wedge \\ \quad 0 < b \wedge \\ \quad b < c \wedge \\ \quad c \leq 1 \Rightarrow \\ \quad \quad \operatorname{magnitude}(\int_{b}^{c} \operatorname{gamma\text{-}integrand}(z)) \leq (\exp (a \cdot \log\left(c\right)) \cdot \operatorname{recip}(a)) \end{array}}}
\vnbprop{gamma-tail-large}{\vnbstmt{\begin{array}{@{}l@{}} \forall z, b, a, c, d. \\ \quad z \in \mathbb{C} \wedge \\ \quad b \in \mathbb{R} \wedge \\ \quad a \in \mathbb{R} \wedge \\ \quad c \in \mathbb{R} \wedge \\ \quad d \in \mathbb{R} \wedge \\ \quad \operatorname{real\text{-}part}(z) \leq b \wedge \\ \quad 1 \leq a \wedge \\ \quad (b - 1) \leq a \wedge \\ \quad 1 \leq c \wedge \\ \quad c < d \Rightarrow \\ \quad \quad \operatorname{magnitude}(\int_{c}^{d} \operatorname{gamma\text{-}integrand}(z)) \leq (2 \cdot \exp ((a \cdot (\log\left((2 \cdot a)\right) - 1)) - (1/2 \cdot c))) \end{array}}}
\vnbprop{log-monotone}{\vnbstmt{\forall u, v \in \mathbb{R}.\; (((0 < u) \wedge (u \leq v)) \Rightarrow (\log\left(u\right) \leq \log\left(v\right)))}}
\vnbprop{rr-power-small-at-zero}{\vnbstmt{\begin{array}{@{}l@{}} \forall s, e. \\ \quad ((s \in \mathbb{R}) \wedge ((e \in \mathbb{R}) \wedge ((0 < s) \wedge (0 < e)))) \Rightarrow  \\ \quad \quad \exists d \in \mathbb{R}.\; ((0 < d) \wedge ((d \leq 1) \wedge ((\exp (s \cdot \log\left(d\right)) \cdot \operatorname{recip}(s)) \leq e))) \end{array}}}
\vnbprop{rr-power-monotone}{\vnbstmt{\begin{array}{@{}l@{}} \forall s, r, d. \\ \quad ((s \in \mathbb{R}) \wedge ((r \in \mathbb{R}) \wedge ((d \in \mathbb{R}) \wedge ((0 < s) \wedge ((0 < r) \wedge (r \leq d)))))) \Rightarrow  \\ \quad \quad (\exp (s \cdot \log\left(r\right)) \cdot \operatorname{recip}(s)) \leq (\exp (s \cdot \log\left(d\right)) \cdot \operatorname{recip}(s)) \end{array}}}
\vnbprop{exp-tail-small-at-infinity}{\vnbstmt{\begin{array}{@{}l@{}} \forall k, e. \\ \quad ((k \in \mathbb{R}) \wedge ((e \in \mathbb{R}) \wedge (0 < e))) \Rightarrow  \\ \quad \quad \exists m \in \mathbb{R}.\; ((1 \leq m) \wedge ((2 \cdot \exp (k - (1/2 \cdot m))) \leq e)) \end{array}}}
\vnbprop{exp-tail-monotone}{\vnbstmt{\begin{array}{@{}l@{}} \forall k, m, x. \\ \quad ((k \in \mathbb{R}) \wedge ((m \in \mathbb{R}) \wedge ((x \in \mathbb{R}) \wedge (m \leq x)))) \Rightarrow  \\ \quad \quad (2 \cdot \exp (k - (1/2 \cdot x))) \leq (2 \cdot \exp (k - (1/2 \cdot m))) \end{array}}}
\vnbprop{gamma-small-difference}{\vnbstmt{\begin{array}{@{}l@{}} \forall z, a, d, b, c. \\ \quad z \in \mathbb{C} \wedge \\ \quad a \in \mathbb{R} \wedge \\ \quad d \in \mathbb{R} \wedge \\ \quad b \in \mathbb{R} \wedge \\ \quad c \in \mathbb{R} \wedge \\ \quad 0 < a \wedge \\ \quad a \leq \operatorname{real\text{-}part}(z) \wedge \\ \quad 0 < b \wedge \\ \quad b < d \wedge \\ \quad 0 < c \wedge \\ \quad c < d \wedge \\ \quad d \leq 1 \Rightarrow \\ \quad \quad \operatorname{magnitude}((\int_{b}^{1} \operatorname{gamma\text{-}integrand}(z) - \int_{c}^{1} \operatorname{gamma\text{-}integrand}(z))) \\ \quad \quad \leq\; \exp (a \cdot \log\left(d\right)) \cdot \operatorname{recip}(a) \end{array}}}
\vnbprop{gamma-large-difference}{\vnbstmt{\begin{array}{@{}l@{}} \forall z, b, a, m, c, d. \\ \quad z \in \mathbb{C} \wedge \\ \quad b \in \mathbb{R} \wedge \\ \quad a \in \mathbb{R} \wedge \\ \quad m \in \mathbb{R} \wedge \\ \quad c \in \mathbb{R} \wedge \\ \quad d \in \mathbb{R} \wedge \\ \quad \operatorname{real\text{-}part}(z) \leq b \wedge \\ \quad 1 \leq a \wedge \\ \quad (b - 1) \leq a \wedge \\ \quad 1 \leq m \wedge \\ \quad m < c \wedge \\ \quad m < d \Rightarrow \\ \quad \quad \operatorname{magnitude}((\int_{1}^{c} \operatorname{gamma\text{-}integrand}(z) - \int_{1}^{d} \operatorname{gamma\text{-}integrand}(z))) \\ \quad \quad \leq\; 2 \cdot \exp ((a \cdot (\log\left((2 \cdot a)\right) - 1)) - (1/2 \cdot m)) \end{array}}}
\vnbprop{gamma-converges}{\vnbstmt{\begin{array}{@{}l@{}} \forall z \in \mathbb{C}. \\ \quad (0 < \operatorname{real\text{-}part}(z)) \Rightarrow  \\ \quad \quad \operatorname{improper\text{-}int\text{-}converges\text{-}to}(\operatorname{gamma\text{-}integrand}(z), \Gamma(z)) \end{array}}}
\vnbprop{cc-int-linear-combination}{\vnbstmt{\begin{array}{@{}l@{}} \forall u, a, b, p, q.; \forall c \in \mathbb{C}, h \in (u \to \mathbb{C}). \\ \quad u \subseteq \mathbb{C} \wedge \\ \quad a \in (u \to \mathbb{C}) \wedge \\ \quad \forall y \in u. \\ \quad \quad \operatorname{is\text{-}continuous\text{-}at}(\operatorname{subspace\text{-}ms}(\operatorname{nf\text{-}metric\text{-}space}(\operatorname{cc\text{-}normed\text{-}field}), u), \\ \quad \quad \quad \operatorname{nf\text{-}metric\text{-}space}(\operatorname{cc\text{-}normed\text{-}field}), a, y) \wedge \\ \quad p \in \mathbb{R} \wedge \\ \quad q \in \mathbb{R} \wedge \\ \quad p < q \wedge \\ \quad \operatorname{ccint}(p, q) \subseteq u \wedge \\ \quad b \in (u \to \mathbb{C}) \wedge \\ \quad \forall y \in u. \\ \quad \quad \operatorname{is\text{-}continuous\text{-}at}(\operatorname{subspace\text{-}ms}(\operatorname{nf\text{-}metric\text{-}space}(\operatorname{cc\text{-}normed\text{-}field}), u), \\ \quad \quad \quad \operatorname{nf\text{-}metric\text{-}space}(\operatorname{cc\text{-}normed\text{-}field}), b, y) \wedge \\ \quad \forall v \in u.\; (h(v) \simeq (a(v) + (c \cdot b(v)))) \Rightarrow \\ \quad \quad \forall y \in u. \\ \quad \quad \quad \operatorname{is\text{-}continuous\text{-}at}(\operatorname{subspace\text{-}ms}(\operatorname{nf\text{-}metric\text{-}space}(\operatorname{cc\text{-}normed\text{-}field}), u), \\ \quad \quad \quad \quad \operatorname{nf\text{-}metric\text{-}space}(\operatorname{cc\text{-}normed\text{-}field}), h, y) \wedge \\ \quad \quad \int_{p}^{q} h = (\int_{p}^{q} a + (c \cdot \int_{p}^{q} b)) \end{array}}}
\vnbprop{gamma-integrand-expansion}{\vnbstmt{\begin{array}{@{}l@{}} \forall x, a, t. \\ \quad ((x \in \mathbb{C}) \wedge ((a \in \mathbb{C}) \wedge ((t \in \mathbb{R}) \wedge (0 < t)))) \Rightarrow  \\ \quad \quad ((\operatorname{gamma\text{-}integrand}(x)(t) - \operatorname{gamma\text{-}integrand}(a)(t)) - ((x - a) \cdot (\operatorname{gamma\text{-}integrand}(a)(t) \cdot \log\left(t\right)))) \\ \quad \quad =\; \operatorname{gamma\text{-}integrand}(a)(t) \cdot ((\exp ((x - a) \cdot \log\left(t\right)) - 1) - ((x - a) \cdot \log\left(t\right))) \end{array}}}
\vnbprop{log-abs-bound-on-interval}{\vnbstmt{\begin{array}{@{}l@{}} \forall p, q, t. \\ \quad ((p \in \mathbb{R}) \wedge ((q \in \mathbb{R}) \wedge ((t \in \mathbb{R}) \wedge ((0 < p) \wedge ((p \leq t) \wedge (t \leq q)))))) \Rightarrow  \\ \quad \quad \lvert \log\left(t\right) \rvert \leq (\lvert \log\left(p\right) \rvert + \lvert \log\left(q\right) \rvert) \end{array}}}
\vnbprop{gamma-integrand-bound-interval}{\vnbstmt{\begin{array}{@{}l@{}} \forall a, p, q, t. \\ \quad ((a \in \mathbb{C}) \wedge ((p \in \mathbb{R}) \wedge ((q \in \mathbb{R}) \wedge ((t \in \mathbb{R}) \wedge ((0 < p) \wedge ((p \leq t) \wedge (t \leq q))))))) \Rightarrow  \\ \quad \quad \operatorname{magnitude}(\operatorname{gamma\text{-}integrand}(a)(t)) \leq \exp (\lvert (\operatorname{real\text{-}part}(a) - 1) \rvert \cdot (\lvert \log\left(p\right) \rvert + \lvert \log\left(q\right) \rvert)) \end{array}}}
\vnbprop{gamma-integrand-quadratic}{\vnbstmt{\begin{array}{@{}l@{}} \forall a, x, p, q, t. \\ \quad a \in \mathbb{C} \wedge \\ \quad x \in \mathbb{C} \wedge \\ \quad p \in \mathbb{R} \wedge \\ \quad q \in \mathbb{R} \wedge \\ \quad t \in \mathbb{R} \wedge \\ \quad 0 < p \wedge \\ \quad p \leq t \wedge \\ \quad t \leq q \wedge \\ \quad (\operatorname{magnitude}((x - a)) \cdot (\lvert \log\left(p\right) \rvert + \lvert \log\left(q\right) \rvert)) \leq 1 \Rightarrow \\ \quad \quad \operatorname{magnitude}(((\operatorname{gamma\text{-}integrand}(x)(t) - \operatorname{gamma\text{-}integrand}(a)(t)) - ((x - a) \cdot (\operatorname{gamma\text{-}integrand}(a)(t) \cdot \log\left(t\right))))) \\ \quad \quad \leq\; (\exp (\lvert (\operatorname{real\text{-}part}(a) - 1) \rvert \cdot (\lvert \log\left(p\right) \rvert + \lvert \log\left(q\right) \rvert)) \cdot ((\lvert \log\left(p\right) \rvert + \lvert \log\left(q\right) \rvert) \cdot (\lvert \log\left(p\right) \rvert + \lvert \log\left(q\right) \rvert))) \cdot (\operatorname{magnitude}((x - a)) \cdot \operatorname{magnitude}((x - a))) \end{array}}}
\vnbprop{gamma-truncation-quadratic}{\vnbstmt{\begin{array}{@{}l@{}} \forall a, x, p, q. \\ \quad a \in \mathbb{C} \wedge \\ \quad x \in \mathbb{C} \wedge \\ \quad p \in \mathbb{R} \wedge \\ \quad q \in \mathbb{R} \wedge \\ \quad 0 < p \wedge \\ \quad p < q \wedge \\ \quad (\operatorname{magnitude}((x - a)) \cdot (\lvert \log\left(p\right) \rvert + \lvert \log\left(q\right) \rvert)) \leq 1 \Rightarrow \\ \quad \quad \operatorname{magnitude}(((\int_{p}^{q} \operatorname{gamma\text{-}integrand}(x) - \int_{p}^{q} \operatorname{gamma\text{-}integrand}(a)) - ((x - a) \cdot \int_{p}^{q} (\lambda w \in \{u \in \mathbb{C} : (0 < \operatorname{real\text{-}part}(u))\}.\; ((\exp ((a - 1) \cdot \log\left(w\right)) \cdot \exp -w) \cdot \log\left(w\right)))))) \\ \quad \quad \leq\; (q - p) \cdot ((\exp (\lvert (\operatorname{real\text{-}part}(a) - 1) \rvert \cdot (\lvert \log\left(p\right) \rvert + \lvert \log\left(q\right) \rvert)) \cdot ((\lvert \log\left(p\right) \rvert + \lvert \log\left(q\right) \rvert) \cdot (\lvert \log\left(p\right) \rvert + \lvert \log\left(q\right) \rvert))) \cdot (\operatorname{magnitude}((x - a)) \cdot \operatorname{magnitude}((x - a)))) \end{array}}}
\vnbprop{gamma-truncation-holomorphic}{\vnbstmt{\begin{array}{@{}l@{}} \forall p, q. \\ \quad ((p \in \mathbb{R}) \wedge ((q \in \mathbb{R}) \wedge ((0 < p) \wedge (p < q)))) \Rightarrow  \\ \quad \quad \operatorname{holomorphic\text{-}on}(\{u \in \mathbb{C} : (0 < \operatorname{real\text{-}part}(u))\}, \\ \quad \quad \quad (\lambda x \in \{u \in \mathbb{C} : (0 < \operatorname{real\text{-}part}(u))\}.\; \int_{p}^{q} \operatorname{gamma\text{-}integrand}(x))) \end{array}}}
\vnbprop{gamma-truncation-error}{\vnbstmt{\begin{array}{@{}l@{}} \forall z, b, c, a, l, h. \\ \quad z \in \mathbb{C} \wedge \\ \quad b \in \mathbb{R} \wedge \\ \quad c \in \mathbb{R} \wedge \\ \quad a \in \mathbb{R} \wedge \\ \quad l \in \mathbb{R} \wedge \\ \quad h \in \mathbb{R} \wedge \\ \quad 0 < b \wedge \\ \quad b \leq \operatorname{real\text{-}part}(z) \wedge \\ \quad \operatorname{real\text{-}part}(z) \leq c \wedge \\ \quad 1 \leq a \wedge \\ \quad (c - 1) \leq a \wedge \\ \quad 0 < l \wedge \\ \quad l \leq 1 \wedge \\ \quad 1 \leq h \wedge \\ \quad l < h \Rightarrow \\ \quad \quad \operatorname{magnitude}((\Gamma(z) - \int_{l}^{h} \operatorname{gamma\text{-}integrand}(z))) \\ \quad \quad \leq\; (\exp (b \cdot \log\left(l\right)) \cdot \operatorname{recip}(b)) + (2 \cdot \exp ((a \cdot (\log\left((2 \cdot a)\right) - 1)) - (1/2 \cdot h))) \end{array}}}
\vnbprop{gamma-holomorphic}{\vnbstmt{\begin{array}{@{}l@{}} \operatorname{holomorphic\text{-}on}(\{u \in \mathbb{C} : (0 < \operatorname{real\text{-}part}(u))\}, \\ \quad (\lambda x \in \{u \in \mathbb{C} : (0 < \operatorname{real\text{-}part}(u))\}.\; \Gamma(x))) \end{array}}}
\vnbprop{cc-int-of-derivative}{\vnbstmt{\begin{array}{@{}l@{}} \forall u, r, c, a, b. \\ \quad u \subseteq \mathbb{C} \wedge \\ \quad r \in (u \to \mathbb{C}) \wedge \\ \quad \forall y \in u. \\ \quad \quad \operatorname{is\text{-}continuous\text{-}at}(\operatorname{subspace\text{-}ms}(\operatorname{nf\text{-}metric\text{-}space}(\operatorname{cc\text{-}normed\text{-}field}), u), \\ \quad \quad \quad \operatorname{nf\text{-}metric\text{-}space}(\operatorname{cc\text{-}normed\text{-}field}), r, y) \wedge \\ \quad a \in \mathbb{R} \wedge \\ \quad b \in \mathbb{R} \wedge \\ \quad a < b \wedge \\ \quad \operatorname{ccint}(a, b) \subseteq u \wedge \\ \quad \operatorname{holomorphic\text{-}on}(u, c) \wedge \\ \quad \forall z \in u.\; \operatorname{is\text{-}diff\text{-}on}(\operatorname{cc\text{-}normed\text{-}field}, u, c, z, r(z)) \Rightarrow \\ \quad \quad \int_{a}^{b} r = (c(b) - c(a)) \end{array}}}
\vnbprop{gamma-at-one}{\vnbstmt{\Gamma(1) = 1}}
\vnbprop{gamma-product-derivative}{\vnbstmt{\begin{array}{@{}l@{}} \forall z \in \mathbb{C}, \mathit{hbz} \in \{u \in \mathbb{C} : (0 < \operatorname{real\text{-}part}(u))\}. \\ \quad \operatorname{is\text{-}diff\text{-}on}(\operatorname{cc\text{-}normed\text{-}field}, \{u \in \mathbb{C} : (0 < \operatorname{real\text{-}part}(u))\}, \\ \quad \quad (\lambda w \in \{u \in \mathbb{C} : (0 < \operatorname{real\text{-}part}(u))\}.\; -(\exp (z \cdot \log\left(w\right)) \cdot \exp -w)), \\ \quad \quad \mathit{hbz}, \\ \quad \quad (\lambda w \in \{u \in \mathbb{C} : (0 < \operatorname{real\text{-}part}(u))\}.\; (\operatorname{gamma\text{-}integrand}((z + 1))(w) + (-z \cdot \operatorname{gamma\text{-}integrand}(z)(w))))(\mathit{hbz})) \end{array}}}
\vnbprop{rr-power-small-at-zero-plain}{\vnbstmt{\begin{array}{@{}l@{}} \forall s, e. \\ \quad ((s \in \mathbb{R}) \wedge ((e \in \mathbb{R}) \wedge ((0 < s) \wedge (0 < e)))) \Rightarrow  \\ \quad \quad \exists d \in \mathbb{R}. \\ \quad \quad \quad ((0 < d) \wedge ((d \leq 1) \wedge \forall t \in \mathbb{R}.\; (((0 < t) \wedge (t \leq d)) \Rightarrow (\exp (s \cdot \log\left(t\right)) \leq e)))) \end{array}}}
\vnbprop{gamma-recursion}{\vnbstmt{\begin{array}{@{}l@{}} \forall z \in \mathbb{C}. \\ \quad (0 < \operatorname{real\text{-}part}(z)) \Rightarrow (\Gamma((z + 1)) = (z \cdot \Gamma(z))) \end{array}}}
\vnbprop{gamma-factorial}{\vnbstmt{\forall n \in \mathbb{N}.\; (\Gamma((n + 1)) = \operatorname{factorial}(n))}}

\vnbbookfile{structure-library/ca5-function-spaces}{structure-library/ca5-function-spaces}
\noindent{\small ca5-function-spaces.scm -- THE VOCABULARY OF CHAPTER 5 (FUNCTION SPACES) of the user's complex analysis notes (\textasciitilde{}/docs/complex-analysis.pdf, s.5.1), first part. DEFINITIONS ONLY; the laws are PROVEN in theorem-library/ca5-function-spaces.scm. Agent CA-5, 2026-10-14.  Every definition below is a decision for the user,}\par
\par\medskip\noindent\textbf{Definition} (\texttt{cc-exhaust}).\;$\operatorname{cc\text{-}exhaust}(u, n) :=$
\begin{equation*}
\{\operatorname{cxz\_} \in \mathbb{C} : ((\operatorname{magnitude}(\operatorname{cxz\_}) \leq n) \wedge (\operatorname{ball}(\operatorname{cc\text{-}ms}, \operatorname{cxz\_}, \operatorname{recip}((n + 1))) \subseteq u))\}
\end{equation*}
\par\noindent{\small\textit{Comment.} $u$ is an open subset of $\mathbb{C}$ and $n \in \mathbb{N}$; $K_n = \{ z \in \mathbb{C} : |z| \le n,\ B(z, 1/(n+1)) \subseteq u \}$.}\par

\par\medskip\noindent\textbf{Definition} (\texttt{cc-cont-fns}).\;$\operatorname{cc\text{-}cont\text{-}fns}(u) :=$
\begin{equation*}
\{\operatorname{ccf\_} \in (u \to \mathbb{C}) : \forall \operatorname{ccy\_} \in u.\; \operatorname{is\text{-}continuous\text{-}at}(\operatorname{subspace\text{-}ms}(\operatorname{nf\text{-}metric\text{-}space}(\operatorname{cc\text{-}normed\text{-}field}), u), \operatorname{nf\text{-}metric\text{-}space}(\operatorname{cc\text{-}normed\text{-}field}), \operatorname{ccf\_}, \operatorname{ccy\_})\}
\end{equation*}

\par\medskip\noindent\textbf{Definition} (\texttt{cc-hol-fns}).\;$\operatorname{cc\text{-}hol\text{-}fns}(u) :=$ $\{\operatorname{ccf\_} \in (u \to \mathbb{C}) : \operatorname{holomorphic\text{-}on}(u, \operatorname{ccf\_})\}$
\par\medskip\noindent\textbf{Definition} (\texttt{cc-fn-module}).\;$\operatorname{cc\text{-}fn\text{-}module}(u, v) :=$
\begin{equation*}
\langle \operatorname{normed\text{-}field\text{-}as\text{-}commutative\text{-}ring}(\operatorname{cc\text{-}normed\text{-}field}), v, (\lambda \langle \operatorname{cfa\_}, \operatorname{cfb\_} \rangle \in (v \times v).\; (\lambda \operatorname{cfz\_} \in u.\; (\operatorname{cfa\_}(\operatorname{cfz\_}) + \operatorname{cfb\_}(\operatorname{cfz\_})))), (\lambda \operatorname{cfz\_} \in u.\; 0), (\lambda \operatorname{cfa\_} \in v.\; (\lambda \operatorname{cfz\_} \in u.\; -\operatorname{cfa\_}(\operatorname{cfz\_}))), (\lambda \langle \operatorname{cfc\_}, \operatorname{cfa\_} \rangle \in (\mathbb{C} \times v).\; (\lambda \operatorname{cfz\_} \in u.\; (\operatorname{cfc\_} \cdot \operatorname{cfa\_}(\operatorname{cfz\_})))) \rangle
\end{equation*}

\par\medskip\noindent\textbf{Definition} (\texttt{cc-exhaust-seminorms}).\;$\operatorname{cc\text{-}exhaust\text{-}seminorms}(u, v) :=$
\begin{equation*}
(\lambda \operatorname{cfk\_} \in \mathbb{N}.\; (\lambda \operatorname{cfa\_} \in v.\; \operatorname{cc\text{-}sup\text{-}norm}(\operatorname{cfa\_}, \operatorname{cc\text{-}exhaust}(u, \operatorname{cfk\_}))))
\end{equation*}

\par\medskip\noindent\textbf{Definition} (\texttt{cc-hol-bounded}).\;For all $u$ and $a$, 
\begin{equation*}
\begin{array}{@{}l@{}} \operatorname{cc\text{-}hol\text{-}bounded}(u, a) \\ \Leftrightarrow\; a \subseteq \operatorname{cc\text{-}hol\text{-}fns}(u) \wedge \\ \forall \operatorname{chn\_} \in \mathbb{N}. \exists \operatorname{chm\_} \in \mathbb{R}. \forall \operatorname{chf\_} \in a, \operatorname{chz\_} \in \operatorname{cc\text{-}exhaust}(u, \operatorname{chn\_}). \\ \quad \operatorname{magnitude}(\operatorname{chf\_}(\operatorname{chz\_})) \leq \operatorname{chm\_} \end{array}
\end{equation*}


\vnbbookfile{structure-library/ca-ac-vocab}{structure-library/ca-ac-vocab}
\noindent{\small ca-ac-vocab.scm -- UNIFORM CONVERGENCE ON COMPACT SUBSETS of an open set of CC. and ZEROS-ACCUMULATE-AT, an accumulation point of the zeros of a function. DEFINITIONS ONLY; the laws and Proposition 3.44 of complex-analysis.pdf are proven in theorem-library/weierstrass-convergence.scm, Proposition 2.12 and Corollary ...}\par
\par\medskip\noindent\textbf{Definition} (\texttt{zeros-accumulate-at}).\;For all $u$, $f$ and $a$, 
\begin{equation*}
\begin{array}{@{}l@{}} \operatorname{zeros\text{-}accumulate\text{-}at}(u, f, a) \\ \Leftrightarrow\; \forall \operatorname{acr\_}. \\ \quad \operatorname{pos\text{-}rr}(\operatorname{acr\_}) \Rightarrow  \\ \quad \quad \exists \operatorname{acx\_} \in u. \\ \quad \quad \quad (\operatorname{acx\_} \in \operatorname{ball}(\operatorname{cc\text{-}ms}, a, \operatorname{acr\_})) \wedge (\neg (\operatorname{acx\_} = a) \wedge (f(\operatorname{acx\_}) = 0)) \end{array}
\end{equation*}

\par\medskip\noindent\textbf{Definition} (\texttt{unif-on-compacta}).\;For all $u$, $\mathit{fam}$ and $f$, 
\begin{equation*}
\begin{array}{@{}l@{}} \operatorname{unif\text{-}on\text{-}compacta}(u, \mathit{fam}, f) \\ \Leftrightarrow\; \forall \operatorname{wcq\_}. \\ \quad (\operatorname{wcq\_} \subseteq u) \Rightarrow  \\ \quad \operatorname{is\text{-}compact}(\operatorname{subspace\text{-}ms}(\operatorname{cc\text{-}ms}, \operatorname{wcq\_})) \Rightarrow  \\ \quad \quad \forall \operatorname{wce\_}. \\ \quad \quad \quad \operatorname{pos\text{-}rr}(\operatorname{wce\_}) \Rightarrow  \\ \quad \quad \quad \quad \exists \operatorname{wcn\_} \in \mathbb{N}. \forall \operatorname{wck\_} \in \mathbb{N}. \\ \quad \quad \quad \quad \quad (\operatorname{wcn\_} \leq \operatorname{wck\_}) \Rightarrow  \\ \quad \quad \quad \quad \quad \quad \forall \operatorname{wcy\_} \in \operatorname{wcq\_}.\; (\operatorname{magnitude}(-(\operatorname{fam}(\operatorname{wck\_})(\operatorname{wcy\_}), f(\operatorname{wcy\_}))) \leq \operatorname{wce\_}) \end{array}
\end{equation*}


\vnbbookfile{theorem-library/ca5-function-spaces}{theorem-library/ca5-function-spaces}
\noindent{\small ca5-function-spaces.scm -- CHAPTER 5 OF THE COMPLEX ANALYSIS NOTES (FUNCTION SPACES), FIRST PART: the exhaustion of an open subset of CC by compact sets (Lemma 5.1), the sup seminorms on it, and the topology of Proposition 5.2. Agent CA-5, 2026-10-14.  Vocabulary: structure-library/ca5-function-spaces.scm}\par
\vnbprop{cc-exhaust-unfold}{\vnbstmt{\begin{array}{@{}l@{}} \forall u, n. \\ \quad (\operatorname{cc\text{-}exhaust}(u, n) \simeq \{z \in \mathbb{C} : ((\operatorname{magnitude}(z) \leq n) \wedge (\operatorname{ball}(\operatorname{cc\text{-}ms}, z, \operatorname{recip}((n + 1))) \subseteq u))\}) \end{array}}}
\vnbprop{cc-exhaust-membership}{\vnbstmt{\begin{array}{@{}l@{}} \forall u, n, z. \\ \quad (z \in \operatorname{cc\text{-}exhaust}(u, n)) \Leftrightarrow ((z \in \mathbb{C}) \wedge ((\operatorname{magnitude}(z) \leq n) \wedge (\operatorname{ball}(\operatorname{cc\text{-}ms}, z, \operatorname{recip}((n + 1))) \subseteq u))) \end{array}}}
\vnbprop{cc-ms-ball-membership}{\vnbstmt{\begin{array}{@{}l@{}} \forall c \in \mathbb{C}, r \in \mathbb{R}, y. \\ \quad (y \in \operatorname{ball}(\operatorname{cc\text{-}ms}, c, r)) \Leftrightarrow ((y \in \mathbb{C}) \wedge (\operatorname{magnitude}((y - c)) < r)) \end{array}}}
\vnbprop{cc-exhaust-radius}{\vnbstmt{\forall n \in \mathbb{N}.\; ((\operatorname{recip}((n + 1)) \in \mathbb{R}) \wedge (0 < \operatorname{recip}((n + 1))))}}
\vnbprop{cc-exhaust-subset}{\vnbstmt{\forall u.; \forall n \in \mathbb{N}.\; (\operatorname{cc\text{-}exhaust}(u, n) \subseteq u)}}
\vnbprop{cc-magnitude-sub-chain}{\vnbstmt{\begin{array}{@{}l@{}} \forall a, b, c \in \mathbb{C}. \\ \quad \operatorname{magnitude}((a - c)) \leq (\operatorname{magnitude}((a - b)) + \operatorname{magnitude}((b - c))) \end{array}}}
\vnbprop{cc-exhaust-radius-decrease}{\vnbstmt{\forall n \in \mathbb{N}.\; ((\operatorname{recip}(((n + 1) + 1)) < \operatorname{recip}((n + 1))) \wedge (\operatorname{recip}((n + 1)) \leq 1))}}
\vnbprop{cc-ms-ball-centre}{\vnbstmt{\forall c \in \mathbb{C}, r \in \mathbb{R}.\; ((0 < r) \Rightarrow (c \in \operatorname{ball}(\operatorname{cc\text{-}ms}, c, r)))}}
\vnbprop{cc-exhaust-interior}{\vnbstmt{\begin{array}{@{}l@{}} \forall u.; \forall n \in \mathbb{N}. \\ \quad \operatorname{cc\text{-}exhaust}(u, n) \subseteq \operatorname{interior}(\operatorname{cc\text{-}ms}, \operatorname{cc\text{-}exhaust}(u, (n + 1))) \end{array}}}
\vnbprop{cc-ms-ball-mono}{\vnbstmt{\begin{array}{@{}l@{}} \forall c \in \mathbb{C}, r, q \in \mathbb{R}. \\ \quad (r \leq q) \Rightarrow (\operatorname{ball}(\operatorname{cc\text{-}ms}, c, r) \subseteq \operatorname{ball}(\operatorname{cc\text{-}ms}, c, q)) \end{array}}}
\vnbprop{cc-exhaust-mono}{\vnbstmt{\begin{array}{@{}l@{}} \forall u.; \forall a, b \in \mathbb{N}. \\ \quad (a \leq b) \Rightarrow (\operatorname{cc\text{-}exhaust}(u, a) \subseteq \operatorname{cc\text{-}exhaust}(u, b)) \end{array}}}
\vnbprop{cc-exhaust-covers}{\vnbstmt{\begin{array}{@{}l@{}} \forall u.; \forall z \in u. \\ \quad \operatorname{is\text{-}open}(\operatorname{cc\text{-}ms}, u) \Rightarrow \exists n \in \mathbb{N}.\; (z \in \operatorname{cc\text{-}exhaust}(u, n)) \end{array}}}
\vnbprop{not-subset-witness}{\vnbstmt{\forall a, b.\; (\neg (a \subseteq b) \Rightarrow \exists x \in a.\; \neg (x \in b))}}
\vnbprop{cc-exhaust-exhaustive}{\vnbstmt{\begin{array}{@{}l@{}} \forall u, p. \\ \quad (\operatorname{is\text{-}open}(\operatorname{nf\text{-}metric\text{-}space}(\operatorname{cc\text{-}normed\text{-}field}), u) \wedge ((p \subseteq u) \wedge \operatorname{is\text{-}compact}(\operatorname{subspace\text{-}ms}(\operatorname{cc\text{-}ms}, p)))) \Rightarrow  \\ \quad \quad \exists n \in \mathbb{N}.\; (p \subseteq \operatorname{cc\text{-}exhaust}(u, n)) \end{array}}}
\vnbprop{cc-bolzano-weierstrass}{\vnbstmt{\begin{array}{@{}l@{}} \forall h \in (\mathbb{N} \to \mathbb{C}), d \in \mathbb{R}. \\ \quad \forall i \in \mathbb{N}.\; (\operatorname{magnitude}(h(i)) \leq d) \Rightarrow  \\ \quad \quad \exists a. \\ \quad \quad \quad \operatorname{strictly\text{-}mono\text{-}nn}(a) \wedge \exists l \in \mathbb{C}.\; \operatorname{converges\text{-}to}(\operatorname{cc\text{-}ms}, \operatorname{subseq}(h, a), l) \end{array}}}
\vnbprop{cc-exhaust-limit-in}{\vnbstmt{\begin{array}{@{}l@{}} \forall u.; \forall n \in \mathbb{N}, s \in (\mathbb{N} \to \mathbb{C}), l \in \mathbb{C}. \\ \quad (\forall i \in \mathbb{N}.\; (s(i) \in \operatorname{cc\text{-}exhaust}(u, n)) \wedge \operatorname{converges\text{-}to}(\operatorname{cc\text{-}ms}, s, l)) \Rightarrow  \\ \quad \quad l \in \operatorname{cc\text{-}exhaust}(u, n) \end{array}}}
\vnbprop{cc-exhaust-compact}{\vnbstmt{\forall u.; \forall n \in \mathbb{N}.\; \operatorname{is\text{-}compact}(\operatorname{subspace\text{-}ms}(\operatorname{cc\text{-}ms}, \operatorname{cc\text{-}exhaust}(u, n)))}}
\vnbprop{cc-continuous-bounded-on-exhaust}{\vnbstmt{\begin{array}{@{}l@{}} \forall u.; \forall f \in (u \to \mathbb{C}), n \in \mathbb{N}. \\ \quad \operatorname{is\text{-}open}(\operatorname{nf\text{-}metric\text{-}space}(\operatorname{cc\text{-}normed\text{-}field}), u) \wedge \\ \quad \forall y \in u. \\ \quad \quad \operatorname{is\text{-}continuous\text{-}at}(\operatorname{subspace\text{-}ms}(\operatorname{nf\text{-}metric\text{-}space}(\operatorname{cc\text{-}normed\text{-}field}), u), \\ \quad \quad \quad \operatorname{nf\text{-}metric\text{-}space}(\operatorname{cc\text{-}normed\text{-}field}), f, y) \Rightarrow \\ \quad \quad \exists m \in \mathbb{R}. \\ \quad \quad \quad (0 \leq m) \wedge \forall z \in \operatorname{cc\text{-}exhaust}(u, n).\; (\operatorname{magnitude}(f(z)) \leq m) \end{array}}}
\vnbprop{cc-sup-norm-upper}{\vnbstmt{\begin{array}{@{}l@{}} \forall g, s.; \forall x \in s. \\ \quad \forall y \in s.\; (g(y) \in \mathbb{C}) \Rightarrow (\operatorname{magnitude}(g(x)) \leq \operatorname{cc\text{-}sup\text{-}norm}(g, s)) \end{array}}}
\vnbprop{cc-exhaust-seminorm-real}{\vnbstmt{\begin{array}{@{}l@{}} \forall u.; \forall f \in (u \to \mathbb{C}), n \in \mathbb{N}. \\ \quad \operatorname{is\text{-}open}(\operatorname{nf\text{-}metric\text{-}space}(\operatorname{cc\text{-}normed\text{-}field}), u) \wedge \\ \quad \forall y \in u. \\ \quad \quad \operatorname{is\text{-}continuous\text{-}at}(\operatorname{subspace\text{-}ms}(\operatorname{nf\text{-}metric\text{-}space}(\operatorname{cc\text{-}normed\text{-}field}), u), \\ \quad \quad \quad \operatorname{nf\text{-}metric\text{-}space}(\operatorname{cc\text{-}normed\text{-}field}), f, y) \Rightarrow \\ \quad \quad (\operatorname{cc\text{-}sup\text{-}norm}(f, \operatorname{cc\text{-}exhaust}(u, n)) \in \mathbb{R}) \wedge (0 \leq \operatorname{cc\text{-}sup\text{-}norm}(f, \operatorname{cc\text{-}exhaust}(u, n))) \end{array}}}
\vnbprop{cc-exhaust-uniform-iff-compacta}{\vnbstmt{\begin{array}{@{}l@{}} \forall u, m, f. \\ \quad \operatorname{is\text{-}open}(\operatorname{nf\text{-}metric\text{-}space}(\operatorname{cc\text{-}normed\text{-}field}), u) \Rightarrow  \\ \quad \quad \forall n \in \mathbb{N}, s. \\ \quad \quad \quad \operatorname{pos\text{-}rr}(s) \Rightarrow  \\ \quad \quad \quad \quad \exists p \in \mathbb{N}. \forall k \in \mathbb{N}, z \in \operatorname{cc\text{-}exhaust}(u, n). \\ \quad \quad \quad \quad \quad (p \leq k) \Rightarrow (\operatorname{magnitude}((m(k)(z) - f(z))) \leq s) \\ \quad \quad \Leftrightarrow\; \operatorname{unif\text{-}on\text{-}compacta}(u, m, f) \end{array}}}

\vnbbookfile{theorem-library/weierstrass-convergence}{theorem-library/weierstrass-convergence}
\noindent{\small weierstrass-convergence.scm -- UNIFORM CONVERGENCE ON COMPACT SUBSETS AND PROPOSITION 3.44 OF complex-analysis.pdf (with Weierstrass's complement: the derivatives converge too).  Agent CA-AC, 2026-10-14. Vocabulary: structure-library/ca-ac-vocab.scm (UNIF-ON-COMPACTA, new; the user's to confirm).}\par
\vnbprop{cc-closed-ball-mem}{\vnbstmt{\begin{array}{@{}l@{}} \forall c \in \mathbb{C}, r.; \forall z \in \operatorname{closed\text{-}ball}(\operatorname{cc\text{-}ms}, c, r). \\ \quad (z \in \mathbb{C}) \wedge (\operatorname{magnitude}((z - c)) \leq r) \end{array}}}
\vnbprop{cc-closed-ball-mem-intro}{\vnbstmt{\begin{array}{@{}l@{}} \forall c \in \mathbb{C}, r.; \forall z \in \mathbb{C}. \\ \quad (\operatorname{magnitude}((z - c)) \leq r) \Rightarrow (z \in \operatorname{closed\text{-}ball}(\operatorname{cc\text{-}ms}, c, r)) \end{array}}}
\vnbprop{cc-parts-bound-of-dist}{\vnbstmt{\begin{array}{@{}l@{}} \forall c \in \mathbb{C}, r \in \mathbb{R}, z \in \mathbb{C}. \\ \quad (\operatorname{magnitude}((z - c)) \leq r) \Rightarrow  \\ \quad \quad ((\lvert \operatorname{real\text{-}part}(z) \rvert \leq (\lvert \operatorname{real\text{-}part}(c) \rvert + r)) \wedge (\lvert \operatorname{imag\text{-}part}(z) \rvert \leq (\lvert \operatorname{imag\text{-}part}(c) \rvert + r))) \end{array}}}
\vnbprop{rr-subseq-converges-abs}{\vnbstmt{\begin{array}{@{}l@{}} \forall b, g, q, e. \\ \quad (\operatorname{strictly\text{-}mono\text{-}nn}(g) \wedge (\operatorname{converges\text{-}to}(\operatorname{rr\text{-}ms}, \operatorname{subseq}(b, g), q) \wedge \operatorname{pos\text{-}rr}(e))) \Rightarrow  \\ \quad \quad \exists n \in \mathbb{N}.\; \forall j \in \mathbb{N}.\; ((n \leq j) \Rightarrow (\lvert (b(g(j)) - q) \rvert \leq e)) \end{array}}}
\vnbprop{cc-closed-ball-compact}{\vnbstmt{\begin{array}{@{}l@{}} \forall c \in \mathbb{C}, r \in \mathbb{R}. \\ \quad \operatorname{is\text{-}compact}(\operatorname{subspace\text{-}ms}(\operatorname{cc\text{-}ms}, \operatorname{closed\text{-}ball}(\operatorname{cc\text{-}ms}, c, r))) \end{array}}}
\vnbprop{unif-on-compacta-locally}{\vnbstmt{\begin{array}{@{}l@{}} \forall u, m, f.; \forall z \in u. \\ \quad (\operatorname{is\text{-}open}(\operatorname{nf\text{-}metric\text{-}space}(\operatorname{cc\text{-}normed\text{-}field}), u) \wedge \operatorname{unif\text{-}on\text{-}compacta}(u, m, f)) \Rightarrow  \\ \quad \quad \exists r. \\ \quad \quad \quad \operatorname{pos\text{-}rr}(r) \wedge \\ \quad \quad \quad \operatorname{ball}(\operatorname{cc\text{-}ms}, z, r) \subseteq u \wedge \\ \quad \quad \quad \forall s. \\ \quad \quad \quad \quad \operatorname{pos\text{-}rr}(s) \Rightarrow  \\ \quad \quad \quad \quad \quad \exists p \in \mathbb{N}. \forall k \in \mathbb{N}, y \in \operatorname{ball}(\operatorname{cc\text{-}ms}, z, r). \\ \quad \quad \quad \quad \quad \quad (p \leq k) \Rightarrow (\operatorname{magnitude}((m(k)(y) - f(y))) \leq s) \end{array}}}
\vnbprop{holomorphic-limit-on-compacta}{\vnbstmt{\begin{array}{@{}l@{}} \forall u, m.; \forall f \in (u \to \mathbb{C}). \\ \quad \operatorname{is\text{-}open}(\operatorname{nf\text{-}metric\text{-}space}(\operatorname{cc\text{-}normed\text{-}field}), u) \wedge \\ \quad \forall k \in \mathbb{N}.\; \operatorname{holomorphic\text{-}on}(u, m(k)) \wedge \\ \quad \operatorname{unif\text{-}on\text{-}compacta}(u, m, f) \Rightarrow \\ \quad \quad \operatorname{holomorphic\text{-}on}(u, f) \end{array}}}
\vnbprop{compact-cover-by-balls}{\vnbstmt{\begin{array}{@{}l@{}} \forall s.; \forall h \in (s \to \mathbb{R}). \\ \quad ((s \subseteq \mathbb{C}) \wedge (\operatorname{is\text{-}compact}(\operatorname{subspace\text{-}ms}(\operatorname{cc\text{-}ms}, s)) \wedge \forall x \in s.\; (0 < h(x)))) \Rightarrow  \\ \quad \quad \exists p. \\ \quad \quad \quad ((p \subseteq s) \wedge ((|p| \in \mathbb{N}) \wedge \forall y \in s.\; \exists x \in p.\; (y \in \operatorname{ball}(\operatorname{cc\text{-}ms}, x, h(x))))) \end{array}}}
\vnbprop{locally-uniform-on-compacta}{\vnbstmt{\begin{array}{@{}l@{}} \forall u, g, l. \\ \quad u \subseteq \mathbb{C} \wedge \\ \quad \forall z \in u. \exists r. \\ \quad \quad \operatorname{pos\text{-}rr}(r) \wedge \\ \quad \quad \operatorname{ball}(\operatorname{cc\text{-}ms}, z, r) \subseteq u \wedge \\ \quad \quad \forall s. \\ \quad \quad \quad \operatorname{pos\text{-}rr}(s) \Rightarrow  \\ \quad \quad \quad \quad \exists p \in \mathbb{N}. \forall k \in \mathbb{N}, y \in \operatorname{ball}(\operatorname{cc\text{-}ms}, z, r). \\ \quad \quad \quad \quad \quad (p \leq k) \Rightarrow (\operatorname{magnitude}((g(k)(y) - l(y))) \leq s) \Rightarrow \\ \quad \quad \operatorname{unif\text{-}on\text{-}compacta}(u, g, l) \end{array}}}
\vnbprop{holomorphic-limit-derivatives}{\vnbstmt{\begin{array}{@{}l@{}} \forall u, m.; \forall f \in (u \to \mathbb{C}). \\ \quad \operatorname{is\text{-}open}(\operatorname{nf\text{-}metric\text{-}space}(\operatorname{cc\text{-}normed\text{-}field}), u) \wedge \\ \quad \forall k \in \mathbb{N}.\; \operatorname{holomorphic\text{-}on}(u, m(k)) \wedge \\ \quad \operatorname{unif\text{-}on\text{-}compacta}(u, m, f) \Rightarrow \\ \quad \quad \operatorname{unif\text{-}on\text{-}compacta}(u, \\ \quad \quad \quad (\lambda \mathit{wcvk} \in \mathbb{N}.\; (\lambda w \in u.\; \operatorname{deriv\text{-}on}(\operatorname{cc\text{-}normed\text{-}field}, u, m(\mathit{wcvk}), w))), \\ \quad \quad \quad (\lambda w \in u.\; \operatorname{deriv\text{-}on}(\operatorname{cc\text{-}normed\text{-}field}, u, f, w))) \end{array}}}

\vnbbookfile{theorem-library/analytic-continuation}{theorem-library/analytic-continuation}
\noindent{\small analytic-continuation.scm -- PROPOSITION 2.12 (PRINCIPLE OF ANALYTIC CONTINUATION) AND COROLLARY 2.13 (FINITELY MANY ZEROS IN A CLOSED BALL) OF complex-analysis.pdf. Agent CA-AC, 2026-10-14.  Vocabulary: structure-library/ca-ac-vocab.scm (ZEROS-ACCUMULATE-AT, new; the user's to confirm).}\par
\vnbprop{zeros-accumulate-value}{\vnbstmt{\begin{array}{@{}l@{}} \forall v.; \forall h \in (v \to \mathbb{C}), a \in v. \\ \quad \operatorname{is\text{-}open}(\operatorname{nf\text{-}metric\text{-}space}(\operatorname{cc\text{-}normed\text{-}field}), v) \wedge \\ \quad \forall y \in v. \\ \quad \quad \operatorname{is\text{-}continuous\text{-}at}(\operatorname{subspace\text{-}ms}(\operatorname{nf\text{-}metric\text{-}space}(\operatorname{cc\text{-}normed\text{-}field}), v), \\ \quad \quad \quad \operatorname{nf\text{-}metric\text{-}space}(\operatorname{cc\text{-}normed\text{-}field}), h, y) \wedge \\ \quad \operatorname{zeros\text{-}accumulate\text{-}at}(v, h, a) \Rightarrow \\ \quad \quad h(a) = 0 \end{array}}}
\vnbprop{cc-series-of-zeros}{\vnbstmt{\begin{array}{@{}l@{}} \forall t \in (\mathbb{N} \to \mathbb{C}), l. \\ \quad (\forall k \in \mathbb{N}.\; (t(k) = 0) \wedge \operatorname{cc\text{-}series\text{-}converges\text{-}to}(t, l)) \Rightarrow (l = 0) \end{array}}}
\vnbprop{cc-power-nonzero}{\vnbstmt{\forall n \in \mathbb{N}, t \in \mathbb{C}.\; (\neg (t = 0) \Rightarrow \neg (\operatorname{power}(t, n) = 0))}}
\vnbprop{holomorphic-zero-factor}{\vnbstmt{\begin{array}{@{}l@{}} \forall n \in \mathbb{N}, v, h.; \forall a \in v. \\ \quad (\operatorname{holomorphic\text{-}on}(v, h) \wedge \operatorname{zeros\text{-}accumulate\text{-}at}(v, h, a)) \Rightarrow  \\ \quad \quad \exists q. \\ \quad \quad \quad \operatorname{holomorphic\text{-}on}(v, q) \wedge \forall x \in v.\; (h(x) = (\operatorname{power}((x - a), n) \cdot q(x))) \end{array}}}
\vnbprop{holomorphic-zero-coefficients}{\vnbstmt{\begin{array}{@{}l@{}} \forall u, f.; \forall c \in \mathbb{C}, s, r \in \mathbb{R}, k \in \mathbb{N}. \\ \quad \operatorname{holomorphic\text{-}on}(u, f) \wedge \\ \quad 0 < s \wedge \\ \quad \operatorname{ball}(\operatorname{cc\text{-}ms}, c, s) \subseteq u \wedge \\ \quad 0 < r \wedge \\ \quad r < s \wedge \\ \quad \operatorname{zeros\text{-}accumulate\text{-}at}(u, f, c) \Rightarrow \\ \quad \quad \operatorname{recip}(((2 \cdot \pi) \cdot +i)) \cdot \operatorname{line\text{-}int}((\lambda w \in \operatorname{trace}(\operatorname{circle\text{-}path}(c, r), 0, (2 \cdot \pi)).\; (\operatorname{restrict}(f, \operatorname{trace}(\operatorname{circle\text{-}path}(c, r), 0, (2 \cdot \pi)))(w) \cdot \operatorname{power}(\operatorname{recip}((w - c)), (k + 1)))), \operatorname{circle\text{-}path}(c, r), \operatorname{circle\text{-}deriv}(c, r), 0, (2 \cdot \pi)) \\ \quad \quad =\; 0 \end{array}}}
\vnbprop{holomorphic-zeros-local}{\vnbstmt{\begin{array}{@{}l@{}} \forall u, f.; \forall c \in u. \\ \quad (\operatorname{holomorphic\text{-}on}(u, f) \wedge \operatorname{zeros\text{-}accumulate\text{-}at}(u, f, c)) \Rightarrow  \\ \quad \quad \exists r. \\ \quad \quad \quad (\operatorname{pos\text{-}rr}(r) \wedge ((\operatorname{ball}(\operatorname{cc\text{-}ms}, c, r) \subseteq u) \wedge \forall z \in \operatorname{ball}(\operatorname{cc\text{-}ms}, c, r).\; (f(z) = 0))) \end{array}}}
\vnbprop{holomorphic-identity-principle}{\vnbstmt{\begin{array}{@{}l@{}} \forall u, f.; \forall a, z \in u. \\ \quad (\operatorname{is\text{-}connected}(\operatorname{cc\text{-}ms}, u) \wedge (\operatorname{holomorphic\text{-}on}(u, f) \wedge \operatorname{zeros\text{-}accumulate\text{-}at}(u, f, a))) \Rightarrow  \\ \quad \quad f(z) = 0 \end{array}}}
\vnbprop{analytic-on-holomorphic}{\vnbstmt{\begin{array}{@{}l@{}} \forall u.; \forall f \in (u \to \mathbb{C}). \\ \quad (\operatorname{is\text{-}open}(\operatorname{cc\text{-}ms}, u) \wedge \forall b \in u.\; f \text{ is analytic at } b) \Rightarrow  \\ \quad \quad \operatorname{holomorphic\text{-}on}(u, f) \end{array}}}
\vnbprop{analytic-continuation}{\vnbstmt{\begin{array}{@{}l@{}} \forall u.; \forall f \in (u \to \mathbb{C}), z \in u. \\ \quad \operatorname{is\text{-}open}(\operatorname{cc\text{-}ms}, u) \wedge \\ \quad \operatorname{is\text{-}connected}(\operatorname{cc\text{-}ms}, u) \wedge \\ \quad \forall b \in u.\; f \text{ is analytic at } b \wedge \\ \quad \exists a \in u.\; \operatorname{zeros\text{-}accumulate\text{-}at}(u, f, a) \Rightarrow \\ \quad \quad f(z) = 0 \end{array}}}
\vnbprop{analytic-zeros-finite-in-ball}{\vnbstmt{\begin{array}{@{}l@{}} \forall u.; \forall f \in (u \to \mathbb{C}), c \in \mathbb{C}, r \in \mathbb{R}. \\ \quad \operatorname{is\text{-}open}(\operatorname{cc\text{-}ms}, u) \wedge \\ \quad \operatorname{is\text{-}connected}(\operatorname{cc\text{-}ms}, u) \wedge \\ \quad \forall b \in u.\; f \text{ is analytic at } b \wedge \\ \quad \exists e \in u.\; \neg (f(e) = 0) \wedge \\ \quad \operatorname{closed\text{-}ball}(\operatorname{cc\text{-}ms}, c, r) \subseteq u \Rightarrow \\ \quad \quad |\{z \in \operatorname{closed\text{-}ball}(\operatorname{cc\text{-}ms}, c, r) : (f(z) = 0)\}| \in \mathbb{N} \end{array}}}

\vnbbookfile{theorem-library/ca5-frechet}{theorem-library/ca5-frechet}
\noindent{\small ca5-frechet.scm -- CHAPTER 5 OF THE COMPLEX ANALYSIS NOTES, continued: C(U) AND H(U) AS FRECHET SPACES (Proposition 5.2) and H(U) CLOSED IN C(U) (Proposition 5.3). Agent CA-5, 2026-10-14.  Vocabulary: structure-library/ca5-function-spaces.scm (CC-CONT-FNS, CC-HOL-FNS, CC-FN-MODULE, CC-EXHAUST-SEMINORMS); the Frechet ...}\par
\vnbprop{cc-scalar-ring-carr}{\vnbstmt{(\operatorname{carr}(\operatorname{normed\text{-}field\text{-}as\text{-}commutative\text{-}ring}(\operatorname{cc\text{-}normed\text{-}field})) \simeq \mathbb{C})}}
\vnbprop{cc-scalar-ring-add}{\vnbstmt{(\operatorname{add}(\operatorname{normed\text{-}field\text{-}as\text{-}commutative\text{-}ring}(\operatorname{cc\text{-}normed\text{-}field})) \simeq (\lambda \langle \operatorname{x\_}, \operatorname{y\_} \rangle \in (\mathbb{C} \times \mathbb{C}).\; (\operatorname{x\_} + \operatorname{y\_})))}}
\vnbprop{cc-scalar-ring-mul}{\vnbstmt{(\operatorname{mul}(\operatorname{normed\text{-}field\text{-}as\text{-}commutative\text{-}ring}(\operatorname{cc\text{-}normed\text{-}field})) \simeq (\lambda \langle \operatorname{x\_}, \operatorname{y\_} \rangle \in (\mathbb{C} \times \mathbb{C}).\; (\operatorname{x\_} \cdot \operatorname{y\_})))}}
\vnbprop{cc-scalar-ring-neg}{\vnbstmt{(\operatorname{neg}(\operatorname{normed\text{-}field\text{-}as\text{-}commutative\text{-}ring}(\operatorname{cc\text{-}normed\text{-}field})) \simeq (\lambda x \in \mathbb{C}.\; -x))}}
\vnbprop{cc-scalar-ring-zero}{\vnbstmt{(\operatorname{zero}(\operatorname{normed\text{-}field\text{-}as\text{-}commutative\text{-}ring}(\operatorname{cc\text{-}normed\text{-}field})) \simeq 0)}}
\vnbprop{cc-scalar-ring-one}{\vnbstmt{(\operatorname{one}(\operatorname{normed\text{-}field\text{-}as\text{-}commutative\text{-}ring}(\operatorname{cc\text{-}normed\text{-}field})) \simeq 1)}}
\vnbprop{cc-cont-fns-unfold}{\vnbstmt{\begin{array}{@{}l@{}} \forall u. \\ \quad (\operatorname{cc\text{-}cont\text{-}fns}(u) \simeq \{f \in (u \to \mathbb{C}) : \forall y \in u.\; \operatorname{is\text{-}continuous\text{-}at}(\operatorname{subspace\text{-}ms}(\operatorname{nf\text{-}metric\text{-}space}(\operatorname{cc\text{-}normed\text{-}field}), u), \operatorname{nf\text{-}metric\text{-}space}(\operatorname{cc\text{-}normed\text{-}field}), f, y)\}) \end{array}}}
\vnbprop{cc-cont-fns-membership}{\vnbstmt{\begin{array}{@{}l@{}} \forall u, h. \\ \quad h \in \operatorname{cc\text{-}cont\text{-}fns}(u) \\ \quad \Leftrightarrow\; h \in (u \to \mathbb{C}) \wedge \\ \quad \forall y \in u. \\ \quad \quad \operatorname{is\text{-}continuous\text{-}at}(\operatorname{subspace\text{-}ms}(\operatorname{nf\text{-}metric\text{-}space}(\operatorname{cc\text{-}normed\text{-}field}), u), \\ \quad \quad \quad \operatorname{nf\text{-}metric\text{-}space}(\operatorname{cc\text{-}normed\text{-}field}), h, y) \end{array}}}
\vnbprop{cc-fn-module-scal}{\vnbstmt{\begin{array}{@{}l@{}} \forall u, v \in \mathbf{Set}. \\ \quad (\operatorname{scal}(\operatorname{cc\text{-}fn\text{-}module}(u, v)) \simeq \operatorname{normed\text{-}field\text{-}as\text{-}commutative\text{-}ring}(\operatorname{cc\text{-}normed\text{-}field})) \end{array}}}
\vnbprop{cc-fn-module-vec}{\vnbstmt{\forall u, v \in \mathbf{Set}.\; (\operatorname{vec}(\operatorname{cc\text{-}fn\text{-}module}(u, v)) \simeq v)}}
\vnbprop{cc-fn-module-vadd}{\vnbstmt{\begin{array}{@{}l@{}} \forall u, v \in \mathbf{Set}. \\ \quad (\operatorname{vadd}(\operatorname{cc\text{-}fn\text{-}module}(u, v)) \simeq (\lambda \langle \operatorname{cfa\_}, \operatorname{cfb\_} \rangle \in (v \times v).\; (\lambda z \in u.\; (\operatorname{cfa\_}(z) + \operatorname{cfb\_}(z))))) \end{array}}}
\vnbprop{cc-fn-module-vzero}{\vnbstmt{\forall u, v \in \mathbf{Set}.\; (\operatorname{vzero}(\operatorname{cc\text{-}fn\text{-}module}(u, v)) \simeq (\lambda z \in u.\; 0))}}
\vnbprop{cc-fn-module-vneg}{\vnbstmt{\begin{array}{@{}l@{}} \forall u, v \in \mathbf{Set}. \\ \quad (\operatorname{vneg}(\operatorname{cc\text{-}fn\text{-}module}(u, v)) \simeq (\lambda a \in v.\; (\lambda z \in u.\; -a(z)))) \end{array}}}
\vnbprop{cc-fn-module-act}{\vnbstmt{\begin{array}{@{}l@{}} \forall u, v \in \mathbf{Set}. \\ \quad (\operatorname{act}(\operatorname{cc\text{-}fn\text{-}module}(u, v)) \simeq (\lambda \langle \operatorname{cfc\_}, \operatorname{cfa\_} \rangle \in (\mathbb{C} \times v).\; (\lambda z \in u.\; (\operatorname{cfc\_} \cdot \operatorname{cfa\_}(z))))) \end{array}}}
\vnbprop{cc-cont-fns-fun}{\vnbstmt{\forall u.; \forall a \in \operatorname{cc\text{-}cont\text{-}fns}(u).\; ((u \subseteq \mathbb{C}) \Rightarrow (a \in (u \to \mathbb{C})))}}
\vnbprop{cc-cont-fns-add}{\vnbstmt{\begin{array}{@{}l@{}} \forall u.; \forall a, b \in \operatorname{cc\text{-}cont\text{-}fns}(u). \\ \quad (u \subseteq \mathbb{C}) \Rightarrow ((\lambda z \in u.\; (a(z) + b(z))) \in \operatorname{cc\text{-}cont\text{-}fns}(u)) \end{array}}}
\vnbprop{cc-cont-fns-zero}{\vnbstmt{\forall u.\; ((u \subseteq \mathbb{C}) \Rightarrow ((\lambda z \in u.\; 0) \in \operatorname{cc\text{-}cont\text{-}fns}(u)))}}
\vnbprop{cc-cont-fns-scale}{\vnbstmt{\begin{array}{@{}l@{}} \forall u.; \forall c \in \mathbb{C}, a \in \operatorname{cc\text{-}cont\text{-}fns}(u). \\ \quad (u \subseteq \mathbb{C}) \Rightarrow ((\lambda z \in u.\; (c \cdot a(z))) \in \operatorname{cc\text{-}cont\text{-}fns}(u)) \end{array}}}
\vnbprop{cc-cont-fns-neg}{\vnbstmt{\begin{array}{@{}l@{}} \forall u.; \forall a \in \operatorname{cc\text{-}cont\text{-}fns}(u). \\ \quad (u \subseteq \mathbb{C}) \Rightarrow ((\lambda z \in u.\; -a(z)) \in \operatorname{cc\text{-}cont\text{-}fns}(u)) \end{array}}}
\vnbprop{cc-cont-fns-is-set}{\vnbstmt{\forall u.\; ((u \subseteq \mathbb{C}) \Rightarrow (\operatorname{cc\text{-}cont\text{-}fns}(u) \in \mathbf{Set}))}}
\vnbprop{cc-fn-vadd-assoc}{\vnbstmt{\begin{array}{@{}l@{}} \forall u.; \forall a, b, c \in \operatorname{cc\text{-}cont\text{-}fns}(u). \\ \quad (u \subseteq \mathbb{C}) \Rightarrow  \\ \quad \quad (\lambda \langle \operatorname{cfa\_}, \operatorname{cfb\_} \rangle \in (\operatorname{cc\text{-}cont\text{-}fns}(u) \times \operatorname{cc\text{-}cont\text{-}fns}(u)).\; (\lambda z \in u.\; (\operatorname{cfa\_}(z) + \operatorname{cfb\_}(z))))((\lambda \langle \operatorname{cfa\_}, \operatorname{cfb\_} \rangle \in (\operatorname{cc\text{-}cont\text{-}fns}(u) \times \operatorname{cc\text{-}cont\text{-}fns}(u)).\; (\lambda z \in u.\; (\operatorname{cfa\_}(z) + \operatorname{cfb\_}(z))))(a, b), \\ \quad \quad \quad c) \\ \quad \quad =\; (\lambda \langle \operatorname{cfa\_}, \operatorname{cfb\_} \rangle \in (\operatorname{cc\text{-}cont\text{-}fns}(u) \times \operatorname{cc\text{-}cont\text{-}fns}(u)).\; (\lambda z \in u.\; (\operatorname{cfa\_}(z) + \operatorname{cfb\_}(z))))(a, \\ \quad \quad \quad (\lambda \langle \operatorname{cfa\_}, \operatorname{cfb\_} \rangle \in (\operatorname{cc\text{-}cont\text{-}fns}(u) \times \operatorname{cc\text{-}cont\text{-}fns}(u)).\; (\lambda z \in u.\; (\operatorname{cfa\_}(z) + \operatorname{cfb\_}(z))))(b, c)) \end{array}}}
\vnbprop{cc-fn-vadd-comm}{\vnbstmt{\begin{array}{@{}l@{}} \forall u.; \forall a, b \in \operatorname{cc\text{-}cont\text{-}fns}(u). \\ \quad (u \subseteq \mathbb{C}) \Rightarrow  \\ \quad \quad (\lambda \langle \operatorname{cfa\_}, \operatorname{cfb\_} \rangle \in (\operatorname{cc\text{-}cont\text{-}fns}(u) \times \operatorname{cc\text{-}cont\text{-}fns}(u)).\; (\lambda z \in u.\; (\operatorname{cfa\_}(z) + \operatorname{cfb\_}(z))))(a, b) \\ \quad \quad =\; (\lambda \langle \operatorname{cfa\_}, \operatorname{cfb\_} \rangle \in (\operatorname{cc\text{-}cont\text{-}fns}(u) \times \operatorname{cc\text{-}cont\text{-}fns}(u)).\; (\lambda z \in u.\; (\operatorname{cfa\_}(z) + \operatorname{cfb\_}(z))))(b, a) \end{array}}}
\vnbprop{cc-fn-vadd-zero-left}{\vnbstmt{\begin{array}{@{}l@{}} \forall u.; \forall a \in \operatorname{cc\text{-}cont\text{-}fns}(u). \\ \quad (u \subseteq \mathbb{C}) \Rightarrow  \\ \quad \quad (\lambda \langle \operatorname{cfa\_}, \operatorname{cfb\_} \rangle \in (\operatorname{cc\text{-}cont\text{-}fns}(u) \times \operatorname{cc\text{-}cont\text{-}fns}(u)).\; (\lambda z \in u.\; (\operatorname{cfa\_}(z) + \operatorname{cfb\_}(z))))((\lambda z \in u.\; 0), a) \\ \quad \quad =\; a \end{array}}}
\vnbprop{cc-fn-vadd-zero-right}{\vnbstmt{\begin{array}{@{}l@{}} \forall u.; \forall a \in \operatorname{cc\text{-}cont\text{-}fns}(u). \\ \quad (u \subseteq \mathbb{C}) \Rightarrow  \\ \quad \quad (\lambda \langle \operatorname{cfa\_}, \operatorname{cfb\_} \rangle \in (\operatorname{cc\text{-}cont\text{-}fns}(u) \times \operatorname{cc\text{-}cont\text{-}fns}(u)).\; (\lambda z \in u.\; (\operatorname{cfa\_}(z) + \operatorname{cfb\_}(z))))(a, (\lambda z \in u.\; 0)) \\ \quad \quad =\; a \end{array}}}
\vnbprop{cc-fn-vadd-neg-left}{\vnbstmt{\begin{array}{@{}l@{}} \forall u.; \forall a \in \operatorname{cc\text{-}cont\text{-}fns}(u). \\ \quad (u \subseteq \mathbb{C}) \Rightarrow  \\ \quad \quad (\lambda \langle \mathit{cfa}, \operatorname{cfb\_} \rangle \in (\operatorname{cc\text{-}cont\text{-}fns}(u) \times \operatorname{cc\text{-}cont\text{-}fns}(u)).\; (\lambda z \in u.\; (\operatorname{cfa}(z) + \operatorname{cfb\_}(z))))((\lambda \mathit{cfa} \in \operatorname{cc\text{-}cont\text{-}fns}(u).\; (\lambda z \in u.\; -\operatorname{cfa}(z)))(a), \\ \quad \quad \quad a) \\ \quad \quad =\; (\lambda z \in u.\; 0) \end{array}}}
\vnbprop{cc-fn-vadd-neg-right}{\vnbstmt{\begin{array}{@{}l@{}} \forall u.; \forall a \in \operatorname{cc\text{-}cont\text{-}fns}(u). \\ \quad (u \subseteq \mathbb{C}) \Rightarrow  \\ \quad \quad (\lambda \langle \mathit{cfa}, \operatorname{cfb\_} \rangle \in (\operatorname{cc\text{-}cont\text{-}fns}(u) \times \operatorname{cc\text{-}cont\text{-}fns}(u)).\; (\lambda z \in u.\; (\operatorname{cfa}(z) + \operatorname{cfb\_}(z))))(a, \\ \quad \quad \quad (\lambda \mathit{cfa} \in \operatorname{cc\text{-}cont\text{-}fns}(u).\; (\lambda z \in u.\; -\operatorname{cfa}(z)))(a)) \\ \quad \quad =\; (\lambda z \in u.\; 0) \end{array}}}
\vnbprop{cc-fn-act-distrib-vec}{\vnbstmt{\begin{array}{@{}l@{}} \forall u.; \forall a, b \in \operatorname{cc\text{-}cont\text{-}fns}(u), r \in \mathbb{C}. \\ \quad (u \subseteq \mathbb{C}) \Rightarrow  \\ \quad \quad (\lambda \langle \operatorname{cfc\_}, \operatorname{cfa\_} \rangle \in (\mathbb{C} \times \operatorname{cc\text{-}cont\text{-}fns}(u)).\; (\lambda z \in u.\; (\operatorname{cfc\_} \cdot \operatorname{cfa\_}(z))))(r, \\ \quad \quad \quad (\lambda \langle \operatorname{cfa\_}, \operatorname{cfb\_} \rangle \in (\operatorname{cc\text{-}cont\text{-}fns}(u) \times \operatorname{cc\text{-}cont\text{-}fns}(u)).\; (\lambda z \in u.\; (\operatorname{cfa\_}(z) + \operatorname{cfb\_}(z))))(a, b)) \\ \quad \quad =\; (\lambda \langle \operatorname{cfa\_}, \operatorname{cfb\_} \rangle \in (\operatorname{cc\text{-}cont\text{-}fns}(u) \times \operatorname{cc\text{-}cont\text{-}fns}(u)).\; (\lambda z \in u.\; (\operatorname{cfa\_}(z) + \operatorname{cfb\_}(z))))((\lambda \langle \operatorname{cfc\_}, \operatorname{cfa\_} \rangle \in (\mathbb{C} \times \operatorname{cc\text{-}cont\text{-}fns}(u)).\; (\lambda z \in u.\; (\operatorname{cfc\_} \cdot \operatorname{cfa\_}(z))))(r, a), \\ \quad \quad \quad (\lambda \langle \operatorname{cfc\_}, \operatorname{cfa\_} \rangle \in (\mathbb{C} \times \operatorname{cc\text{-}cont\text{-}fns}(u)).\; (\lambda z \in u.\; (\operatorname{cfc\_} \cdot \operatorname{cfa\_}(z))))(r, b)) \end{array}}}
\vnbprop{cc-fn-act-distrib-scalar}{\vnbstmt{\begin{array}{@{}l@{}} \forall u.; \forall a \in \operatorname{cc\text{-}cont\text{-}fns}(u), r, s \in \mathbb{C}. \\ \quad (u \subseteq \mathbb{C}) \Rightarrow  \\ \quad \quad (\lambda \langle \operatorname{cfc\_}, \operatorname{cfa\_} \rangle \in (\mathbb{C} \times \operatorname{cc\text{-}cont\text{-}fns}(u)).\; (\lambda z \in u.\; (\operatorname{cfc\_} \cdot \operatorname{cfa\_}(z))))((\lambda \langle \operatorname{x\_}, \operatorname{y\_} \rangle \in (\mathbb{C} \times \mathbb{C}).\; (\operatorname{x\_} + \operatorname{y\_}))(r, s), \\ \quad \quad \quad a) \\ \quad \quad =\; (\lambda \langle \operatorname{cfa\_}, \operatorname{cfb\_} \rangle \in (\operatorname{cc\text{-}cont\text{-}fns}(u) \times \operatorname{cc\text{-}cont\text{-}fns}(u)).\; (\lambda z \in u.\; (\operatorname{cfa\_}(z) + \operatorname{cfb\_}(z))))((\lambda \langle \operatorname{cfc\_}, \operatorname{cfa\_} \rangle \in (\mathbb{C} \times \operatorname{cc\text{-}cont\text{-}fns}(u)).\; (\lambda z \in u.\; (\operatorname{cfc\_} \cdot \operatorname{cfa\_}(z))))(r, a), \\ \quad \quad \quad (\lambda \langle \operatorname{cfc\_}, \operatorname{cfa\_} \rangle \in (\mathbb{C} \times \operatorname{cc\text{-}cont\text{-}fns}(u)).\; (\lambda z \in u.\; (\operatorname{cfc\_} \cdot \operatorname{cfa\_}(z))))(s, a)) \end{array}}}
\vnbprop{cc-fn-act-mul-compat}{\vnbstmt{\begin{array}{@{}l@{}} \forall u.; \forall a \in \operatorname{cc\text{-}cont\text{-}fns}(u), r, s \in \mathbb{C}. \\ \quad (u \subseteq \mathbb{C}) \Rightarrow  \\ \quad \quad (\lambda \langle \operatorname{cfc\_}, \operatorname{cfa\_} \rangle \in (\mathbb{C} \times \operatorname{cc\text{-}cont\text{-}fns}(u)).\; (\lambda z \in u.\; (\operatorname{cfc\_} \cdot \operatorname{cfa\_}(z))))((\lambda \langle \operatorname{x\_}, \operatorname{y\_} \rangle \in (\mathbb{C} \times \mathbb{C}).\; (\operatorname{x\_} \cdot \operatorname{y\_}))(r, s), \\ \quad \quad \quad a) \\ \quad \quad =\; (\lambda \langle \operatorname{cfc\_}, \operatorname{cfa\_} \rangle \in (\mathbb{C} \times \operatorname{cc\text{-}cont\text{-}fns}(u)).\; (\lambda z \in u.\; (\operatorname{cfc\_} \cdot \operatorname{cfa\_}(z))))(r, \\ \quad \quad \quad (\lambda \langle \operatorname{cfc\_}, \operatorname{cfa\_} \rangle \in (\mathbb{C} \times \operatorname{cc\text{-}cont\text{-}fns}(u)).\; (\lambda z \in u.\; (\operatorname{cfc\_} \cdot \operatorname{cfa\_}(z))))(s, a)) \end{array}}}
\vnbprop{cc-fn-act-unital}{\vnbstmt{\begin{array}{@{}l@{}} \forall u.; \forall a \in \operatorname{cc\text{-}cont\text{-}fns}(u). \\ \quad (u \subseteq \mathbb{C}) \Rightarrow  \\ \quad \quad (\lambda \langle \operatorname{cfc\_}, \operatorname{cfa\_} \rangle \in (\mathbb{C} \times \operatorname{cc\text{-}cont\text{-}fns}(u)).\; (\lambda z \in u.\; (\operatorname{cfc\_} \cdot \operatorname{cfa\_}(z))))(1, a) \\ \quad \quad =\; a \end{array}}}
\vnbprop{cc-fn-module-is-module}{\vnbstmt{\forall u.\; ((u \subseteq \mathbb{C}) \Rightarrow \operatorname{is\text{-}module}(\operatorname{cc\text{-}fn\text{-}module}(u, \operatorname{cc\text{-}cont\text{-}fns}(u))))}}
\vnbprop{cc-exhaust-seminorm-subadditive}{\vnbstmt{\begin{array}{@{}l@{}} \forall u.; \forall a, b \in \operatorname{cc\text{-}cont\text{-}fns}(u), n \in \mathbb{N}. \\ \quad \operatorname{is\text{-}open}(\operatorname{nf\text{-}metric\text{-}space}(\operatorname{cc\text{-}normed\text{-}field}), u) \Rightarrow  \\ \quad \quad \operatorname{cc\text{-}sup\text{-}norm}((\lambda z \in u.\; (a(z) + b(z))), \operatorname{cc\text{-}exhaust}(u, n)) \\ \quad \quad \leq\; \operatorname{cc\text{-}sup\text{-}norm}(a, \operatorname{cc\text{-}exhaust}(u, n)) + \operatorname{cc\text{-}sup\text{-}norm}(b, \operatorname{cc\text{-}exhaust}(u, n)) \end{array}}}
\vnbprop{cc-exhaust-seminorm-homogeneous}{\vnbstmt{\begin{array}{@{}l@{}} \forall u.; \forall c \in \mathbb{C}, a \in \operatorname{cc\text{-}cont\text{-}fns}(u), n \in \mathbb{N}. \\ \quad \operatorname{is\text{-}open}(\operatorname{nf\text{-}metric\text{-}space}(\operatorname{cc\text{-}normed\text{-}field}), u) \Rightarrow  \\ \quad \quad \operatorname{cc\text{-}sup\text{-}norm}((\lambda z \in u.\; (c \cdot a(z))), \operatorname{cc\text{-}exhaust}(u, n)) \\ \quad \quad =\; \operatorname{magnitude}(c) \cdot \operatorname{cc\text{-}sup\text{-}norm}(a, \operatorname{cc\text{-}exhaust}(u, n)) \end{array}}}
\vnbprop{cc-exhaust-seminorms-separate}{\vnbstmt{\begin{array}{@{}l@{}} \forall u.; \forall a \in \operatorname{cc\text{-}cont\text{-}fns}(u). \\ \quad (\operatorname{is\text{-}open}(\operatorname{nf\text{-}metric\text{-}space}(\operatorname{cc\text{-}normed\text{-}field}), u) \wedge \neg (a = (\lambda z \in u.\; 0))) \Rightarrow  \\ \quad \quad \exists k \in \mathbb{N}.\; \neg (\operatorname{cc\text{-}sup\text{-}norm}(a, \operatorname{cc\text{-}exhaust}(u, k)) = 0) \end{array}}}
\vnbprop{cc-exhaust-seminorms-family}{\vnbstmt{\begin{array}{@{}l@{}} \forall u. \\ \quad \operatorname{is\text{-}open}(\operatorname{nf\text{-}metric\text{-}space}(\operatorname{cc\text{-}normed\text{-}field}), u) \Rightarrow  \\ \quad \quad \operatorname{is\text{-}seminorm\text{-}family}(\operatorname{cc\text{-}normed\text{-}field}, \operatorname{cc\text{-}fn\text{-}module}(u, \operatorname{cc\text{-}cont\text{-}fns}(u)), \\ \quad \quad \quad \operatorname{cc\text{-}exhaust\text{-}seminorms}(u, \operatorname{cc\text{-}cont\text{-}fns}(u))) \end{array}}}
\vnbprop{cc-exhaust-seminorms-unfold}{\vnbstmt{\begin{array}{@{}l@{}} \forall u, v. \\ \quad (\operatorname{cc\text{-}exhaust\text{-}seminorms}(u, v) \simeq (\lambda k \in \mathbb{N}.\; (\lambda a \in v.\; \operatorname{cc\text{-}sup\text{-}norm}(a, \operatorname{cc\text{-}exhaust}(u, k))))) \end{array}}}
\vnbprop{cc-cont-fns-diff}{\vnbstmt{\begin{array}{@{}l@{}} \forall u.; \forall a, b \in \operatorname{cc\text{-}cont\text{-}fns}(u). \\ \quad (u \subseteq \mathbb{C}) \Rightarrow ((\lambda z \in u.\; (a(z) - b(z))) \in \operatorname{cc\text{-}cont\text{-}fns}(u)) \end{array}}}
\vnbprop{cc-cont-fns-add-neg}{\vnbstmt{\begin{array}{@{}l@{}} \forall u.; \forall a, b \in \operatorname{cc\text{-}cont\text{-}fns}(u). \\ \quad (u \subseteq \mathbb{C}) \Rightarrow ((\lambda z \in u.\; (a(z) + -b(z))) \in \operatorname{cc\text{-}cont\text{-}fns}(u)) \end{array}}}
\vnbprop{cc-fn-module-seminorm-of-difference}{\vnbstmt{\begin{array}{@{}l@{}} \forall u.; \forall a, b \in \operatorname{cc\text{-}cont\text{-}fns}(u), k \in \mathbb{N}. \\ \quad (u \subseteq \mathbb{C}) \Rightarrow  \\ \quad \quad (\operatorname{cc\text{-}exhaust\text{-}seminorms}(u, \operatorname{cc\text{-}cont\text{-}fns}(u))(k)(\operatorname{vadd}(\operatorname{cc\text{-}fn\text{-}module}(u, \operatorname{cc\text{-}cont\text{-}fns}(u)))(a, \operatorname{vneg}(\operatorname{cc\text{-}fn\text{-}module}(u, \operatorname{cc\text{-}cont\text{-}fns}(u)))(b))) \simeq \operatorname{cc\text{-}sup\text{-}norm}((\lambda z \in u.\; (a(z) - b(z))), \operatorname{cc\text{-}exhaust}(u, k))) \end{array}}}
\vnbprop{cc-fn-module-cauchy-uniform}{\vnbstmt{\begin{array}{@{}l@{}} \forall u.; \forall s \in (\mathbb{N} \to \operatorname{cc\text{-}cont\text{-}fns}(u)), k \in \mathbb{N}, a. \\ \quad \operatorname{is\text{-}open}(\operatorname{nf\text{-}metric\text{-}space}(\operatorname{cc\text{-}normed\text{-}field}), u) \wedge \\ \quad \operatorname{fr\text{-}cauchy}(\operatorname{cc\text{-}fn\text{-}module}(u, \operatorname{cc\text{-}cont\text{-}fns}(u)), \operatorname{cc\text{-}exhaust\text{-}seminorms}(u, \operatorname{cc\text{-}cont\text{-}fns}(u)), s) \wedge \\ \quad \operatorname{pos\text{-}rr}(a) \Rightarrow \\ \quad \quad \exists p \in \mathbb{N}. \forall i, j \in \mathbb{N}, z \in \operatorname{cc\text{-}exhaust}(u, k). \\ \quad \quad \quad ((p \leq i) \wedge (p \leq j)) \Rightarrow (\operatorname{magnitude}((s(i)(z) - s(j)(z))) < a) \end{array}}}
\vnbprop{cc-fn-module-pointwise-limit}{\vnbstmt{\begin{array}{@{}l@{}} \forall u.; \forall s \in (\mathbb{N} \to \operatorname{cc\text{-}cont\text{-}fns}(u)), z \in u. \\ \quad \operatorname{is\text{-}open}(\operatorname{nf\text{-}metric\text{-}space}(\operatorname{cc\text{-}normed\text{-}field}), u) \wedge \\ \quad \forall k \in \mathbb{N}, a. \\ \quad \quad \operatorname{pos\text{-}rr}(a) \Rightarrow  \\ \quad \quad \quad \exists p \in \mathbb{N}. \forall i, j \in \mathbb{N}, z \in \operatorname{cc\text{-}exhaust}(u, k). \\ \quad \quad \quad \quad ((p \leq i) \wedge (p \leq j)) \Rightarrow (\operatorname{magnitude}((s(i)(z) - s(j)(z))) < a) \Rightarrow \\ \quad \quad \exists l \in \mathbb{C}.\; \operatorname{converges\text{-}to}(\operatorname{cc\text{-}ms}, (\lambda \mathit{cfi} \in \mathbb{N}.\; s(\mathit{cfi})(z)), l) \end{array}}}
\vnbprop{cc-fn-module-limit-function}{\vnbstmt{\begin{array}{@{}l@{}} \forall u.; \forall s \in (\mathbb{N} \to \operatorname{cc\text{-}cont\text{-}fns}(u)). \\ \quad \operatorname{is\text{-}open}(\operatorname{nf\text{-}metric\text{-}space}(\operatorname{cc\text{-}normed\text{-}field}), u) \wedge \\ \quad \forall k \in \mathbb{N}, a. \\ \quad \quad \operatorname{pos\text{-}rr}(a) \Rightarrow  \\ \quad \quad \quad \exists p \in \mathbb{N}. \forall i, j \in \mathbb{N}, z \in \operatorname{cc\text{-}exhaust}(u, k). \\ \quad \quad \quad \quad ((p \leq i) \wedge (p \leq j)) \Rightarrow (\operatorname{magnitude}((s(i)(z) - s(j)(z))) < a) \Rightarrow \\ \quad \quad (\lambda \mathit{cfz} \in u.\; \varepsilon(\{l \in \mathbb{C} : \operatorname{converges\text{-}to}(\operatorname{cc\text{-}ms}, (\lambda \mathit{cfi} \in \mathbb{N}.\; s(\mathit{cfi})(\mathit{cfz})), l)\})) \\ \quad \quad \in\; (u \to \mathbb{C}) \wedge \\ \quad \quad \forall z \in u. \\ \quad \quad \quad \operatorname{converges\text{-}to}(\operatorname{cc\text{-}ms}, (\lambda \mathit{cfi} \in \mathbb{N}.\; s(\mathit{cfi})(z)), \\ \quad \quad \quad \quad (\lambda \mathit{cfz} \in u.\; \varepsilon(\{l \in \mathbb{C} : \operatorname{converges\text{-}to}(\operatorname{cc\text{-}ms}, (\lambda \mathit{cfi} \in \mathbb{N}.\; s(\mathit{cfi})(\mathit{cfz})), l)\}))(z)) \end{array}}}
\vnbprop{cc-fn-module-limit-uniform}{\vnbstmt{\begin{array}{@{}l@{}} \forall u.; \forall s \in (\mathbb{N} \to \operatorname{cc\text{-}cont\text{-}fns}(u)), k \in \mathbb{N}, a. \\ \quad \operatorname{is\text{-}open}(\operatorname{nf\text{-}metric\text{-}space}(\operatorname{cc\text{-}normed\text{-}field}), u) \wedge \\ \quad \forall k \in \mathbb{N}, a. \\ \quad \quad \operatorname{pos\text{-}rr}(a) \Rightarrow  \\ \quad \quad \quad \exists p \in \mathbb{N}. \forall i, j \in \mathbb{N}, z \in \operatorname{cc\text{-}exhaust}(u, k). \\ \quad \quad \quad \quad ((p \leq i) \wedge (p \leq j)) \Rightarrow (\operatorname{magnitude}((s(i)(z) - s(j)(z))) < a) \wedge \\ \quad \operatorname{pos\text{-}rr}(a) \Rightarrow \\ \quad \quad \exists p \in \mathbb{N}. \forall i \in \mathbb{N}, z \in \operatorname{cc\text{-}exhaust}(u, k). \\ \quad \quad \quad (p \leq i) \Rightarrow  \\ \quad \quad \quad \quad \operatorname{magnitude}((s(i)(z) - (\lambda \mathit{cfz} \in u.\; \varepsilon(\{l \in \mathbb{C} : \operatorname{converges\text{-}to}(\operatorname{cc\text{-}ms}, (\lambda \mathit{cfi} \in \mathbb{N}.\; s(\mathit{cfi})(\mathit{cfz})), l)\}))(z))) \\ \quad \quad \quad \quad \leq\; a \end{array}}}
\vnbprop{cc-uniform-limit-continuous-at}{\vnbstmt{\begin{array}{@{}l@{}} \forall u.; \forall s \in (\mathbb{N} \to \operatorname{cc\text{-}cont\text{-}fns}(u)), g \in (u \to \mathbb{C}), a \in u, o \in \mathbb{R}. \\ \quad u \subseteq \mathbb{C} \wedge \\ \quad 0 < o \wedge \\ \quad \operatorname{ball}(\operatorname{cc\text{-}ms}, a, o) \subseteq u \wedge \\ \quad \forall b. \\ \quad \quad \operatorname{pos\text{-}rr}(b) \Rightarrow  \\ \quad \quad \quad \exists p \in \mathbb{N}. \forall i \in \mathbb{N}, z \in \operatorname{ball}(\operatorname{cc\text{-}ms}, a, o). \\ \quad \quad \quad \quad (p \leq i) \Rightarrow (\operatorname{magnitude}((s(i)(z) - g(z))) \leq b) \Rightarrow \\ \quad \quad \operatorname{is\text{-}continuous\text{-}at}(\operatorname{subspace\text{-}ms}(\operatorname{nf\text{-}metric\text{-}space}(\operatorname{cc\text{-}normed\text{-}field}), u), \\ \quad \quad \quad \operatorname{nf\text{-}metric\text{-}space}(\operatorname{cc\text{-}normed\text{-}field}), g, a) \end{array}}}
\vnbprop{cc-cont-fns-frechet}{\vnbstmt{\begin{array}{@{}l@{}} \forall u. \\ \quad \operatorname{is\text{-}open}(\operatorname{nf\text{-}metric\text{-}space}(\operatorname{cc\text{-}normed\text{-}field}), u) \Rightarrow  \\ \quad \quad \operatorname{is\text{-}frechet\text{-}structure}(\operatorname{cc\text{-}normed\text{-}field}, \operatorname{cc\text{-}fn\text{-}module}(u, \operatorname{cc\text{-}cont\text{-}fns}(u)), \\ \quad \quad \quad \operatorname{cc\text{-}exhaust\text{-}seminorms}(u, \operatorname{cc\text{-}cont\text{-}fns}(u))) \end{array}}}
\vnbprop{cc-hol-fns-unfold}{\vnbstmt{\forall u.\; (\operatorname{cc\text{-}hol\text{-}fns}(u) \simeq \{f \in (u \to \mathbb{C}) : \operatorname{holomorphic\text{-}on}(u, f)\})}}
\vnbprop{cc-hol-fns-membership}{\vnbstmt{\forall u, h.\; ((h \in \operatorname{cc\text{-}hol\text{-}fns}(u)) \Leftrightarrow ((h \in (u \to \mathbb{C})) \wedge \operatorname{holomorphic\text{-}on}(u, h)))}}
\vnbprop{cc-hol-fns-subset-cont}{\vnbstmt{\forall u.\; (\operatorname{cc\text{-}hol\text{-}fns}(u) \subseteq \operatorname{cc\text{-}cont\text{-}fns}(u))}}
\vnbprop{cc-hol-fns-closed}{\vnbstmt{\begin{array}{@{}l@{}} \forall u.; \forall s \in (\mathbb{N} \to \operatorname{cc\text{-}hol\text{-}fns}(u)), f \in \operatorname{cc\text{-}cont\text{-}fns}(u). \\ \quad \operatorname{is\text{-}open}(\operatorname{nf\text{-}metric\text{-}space}(\operatorname{cc\text{-}normed\text{-}field}), u) \wedge \\ \quad \operatorname{fr\text{-}conv}(\operatorname{cc\text{-}fn\text{-}module}(u, \operatorname{cc\text{-}cont\text{-}fns}(u)), \operatorname{cc\text{-}exhaust\text{-}seminorms}(u, \operatorname{cc\text{-}cont\text{-}fns}(u)), s, f) \Rightarrow \\ \quad \quad f \in \operatorname{cc\text{-}hol\text{-}fns}(u) \end{array}}}

\vnbbookfile{theorem-library/ca5-montel}{theorem-library/ca5-montel}
\noindent{\small ca5-montel.scm -- CHAPTER 5 OF THE COMPLEX ANALYSIS NOTES, section 5.2: BOUNDED SETS OF HOLOMORPHIC FUNCTIONS, Proposition 5.8 (equicontinuity) and MONTEL'S THEOREM (Theorem 5.7 / Corollary 5.19 for sequences).  Agent CA-5, 2026-10-14.}\par
\vnbprop{cc-holomorphic-deriv-estimate}{\vnbstmt{\begin{array}{@{}l@{}} \forall u, f.; \forall y \in \mathbb{C}, q, m \in \mathbb{R}. \\ \quad \operatorname{is\text{-}open}(\operatorname{nf\text{-}metric\text{-}space}(\operatorname{cc\text{-}normed\text{-}field}), u) \wedge \\ \quad \operatorname{holomorphic\text{-}on}(u, f) \wedge \\ \quad 0 < q \wedge \\ \quad \operatorname{ball}(\operatorname{cc\text{-}ms}, y, (2 \cdot q)) \subseteq u \wedge \\ \quad \forall \mathit{cmy} \in \operatorname{ball}(\operatorname{cc\text{-}ms}, y, (2 \cdot q)).\; (\operatorname{magnitude}(f(\mathit{cmy})) \leq m) \Rightarrow \\ \quad \quad \operatorname{magnitude}(\operatorname{deriv\text{-}on}(\operatorname{cc\text{-}normed\text{-}field}, u, f, y)) \leq (m \cdot \operatorname{recip}(q)) \end{array}}}
\vnbprop{cc-holomorphic-lipschitz-at-centre}{\vnbstmt{\begin{array}{@{}l@{}} \forall u, f.; \forall c \in \mathbb{C}, q, m \in \mathbb{R}, z \in \operatorname{ball}(\operatorname{cc\text{-}ms}, c, q). \\ \quad \operatorname{is\text{-}open}(\operatorname{nf\text{-}metric\text{-}space}(\operatorname{cc\text{-}normed\text{-}field}), u) \wedge \\ \quad \operatorname{holomorphic\text{-}on}(u, f) \wedge \\ \quad 0 < q \wedge \\ \quad \operatorname{ball}(\operatorname{cc\text{-}ms}, c, (3 \cdot q)) \subseteq u \wedge \\ \quad \forall y \in \operatorname{ball}(\operatorname{cc\text{-}ms}, c, (3 \cdot q)).\; (\operatorname{magnitude}(f(y)) \leq m) \Rightarrow \\ \quad \quad \operatorname{magnitude}((f(z) - f(c))) \leq ((m \cdot \operatorname{recip}(q)) \cdot \operatorname{magnitude}((z - c))) \end{array}}}
\vnbprop{cc-hol-bounded-equicontinuous}{\vnbstmt{\begin{array}{@{}l@{}} \forall u, a.; \forall b \in u, s. \\ \quad (\operatorname{is\text{-}open}(\operatorname{nf\text{-}metric\text{-}space}(\operatorname{cc\text{-}normed\text{-}field}), u) \wedge (\operatorname{cc\text{-}hol\text{-}bounded}(u, a) \wedge \operatorname{pos\text{-}rr}(s))) \Rightarrow  \\ \quad \quad \exists l. \\ \quad \quad \quad (\operatorname{pos\text{-}rr}(l) \wedge \forall f \in a, z \in u.\; ((\operatorname{magnitude}((z - b)) < l) \Rightarrow (\operatorname{magnitude}((f(z) - f(b))) < s))) \end{array}}}
\vnbprop{cc-re-unif-cauchy-on-compact}{\vnbstmt{\begin{array}{@{}l@{}} \forall a, q, t, s. \\ \quad a \subseteq \mathbb{C} \wedge \\ \quad \operatorname{is\text{-}compact}(\operatorname{subspace\text{-}ms}(\operatorname{cc\text{-}ms}, a)) \wedge \\ \quad \operatorname{is\text{-}dense\text{-}seq}(\operatorname{subspace\text{-}ms}(\operatorname{cc\text{-}ms}, a), q) \wedge \\ \quad \forall k \in \mathbb{N}, z \in a.\; (t(k)(z) \in \mathbb{C}) \wedge \\ \quad \forall x \in a, s. \\ \quad \quad \operatorname{pos\text{-}rr}(s) \Rightarrow  \\ \quad \quad \quad \exists l. \\ \quad \quad \quad \quad \operatorname{pos\text{-}rr}(l) \wedge \\ \quad \quad \quad \quad \forall k \in \mathbb{N}, y \in a. \\ \quad \quad \quad \quad \quad (\operatorname{magnitude}((y - x)) < l) \Rightarrow (\operatorname{magnitude}((t(k)(y) - t(k)(x))) < s) \wedge \\ \quad \forall j \in \mathbb{N}. \exists \ell \in \mathbb{R}. \\ \quad \quad \operatorname{converges\text{-}to}(\operatorname{rr\text{-}ms}, (\lambda \mathit{cmk} \in \mathbb{N}.\; \operatorname{real\text{-}part}(t(\mathit{cmk})(q(j)))), \ell) \wedge \\ \quad \operatorname{pos\text{-}rr}(s) \Rightarrow \\ \quad \quad \exists p \in \mathbb{N}. \forall i, j \in \mathbb{N}, z \in a. \\ \quad \quad \quad ((p \leq i) \wedge (p \leq j)) \Rightarrow (\lvert (\operatorname{real\text{-}part}(t(i)(z)) - \operatorname{real\text{-}part}(t(j)(z))) \rvert < s) \end{array}}}
\vnbprop{cc-im-unif-cauchy-on-compact}{\vnbstmt{\begin{array}{@{}l@{}} \forall a, q, t, s. \\ \quad a \subseteq \mathbb{C} \wedge \\ \quad \operatorname{is\text{-}compact}(\operatorname{subspace\text{-}ms}(\operatorname{cc\text{-}ms}, a)) \wedge \\ \quad \operatorname{is\text{-}dense\text{-}seq}(\operatorname{subspace\text{-}ms}(\operatorname{cc\text{-}ms}, a), q) \wedge \\ \quad \forall k \in \mathbb{N}, z \in a.\; (t(k)(z) \in \mathbb{C}) \wedge \\ \quad \forall x \in a, s. \\ \quad \quad \operatorname{pos\text{-}rr}(s) \Rightarrow  \\ \quad \quad \quad \exists l. \\ \quad \quad \quad \quad \operatorname{pos\text{-}rr}(l) \wedge \\ \quad \quad \quad \quad \forall k \in \mathbb{N}, y \in a. \\ \quad \quad \quad \quad \quad (\operatorname{magnitude}((y - x)) < l) \Rightarrow (\operatorname{magnitude}((t(k)(y) - t(k)(x))) < s) \wedge \\ \quad \forall j \in \mathbb{N}. \exists \ell \in \mathbb{R}. \\ \quad \quad \operatorname{converges\text{-}to}(\operatorname{rr\text{-}ms}, (\lambda \mathit{cmk} \in \mathbb{N}.\; \operatorname{imag\text{-}part}(t(\mathit{cmk})(q(j)))), \ell) \wedge \\ \quad \operatorname{pos\text{-}rr}(s) \Rightarrow \\ \quad \quad \exists p \in \mathbb{N}. \forall i, j \in \mathbb{N}, z \in a. \\ \quad \quad \quad ((p \leq i) \wedge (p \leq j)) \Rightarrow (\lvert (\operatorname{imag\text{-}part}(t(i)(z)) - \operatorname{imag\text{-}part}(t(j)(z))) \rvert < s) \end{array}}}
\vnbprop{cc-unif-cauchy-on-compact}{\vnbstmt{\begin{array}{@{}l@{}} \forall a, q, t, s. \\ \quad a \subseteq \mathbb{C} \wedge \\ \quad \operatorname{is\text{-}compact}(\operatorname{subspace\text{-}ms}(\operatorname{cc\text{-}ms}, a)) \wedge \\ \quad \operatorname{is\text{-}dense\text{-}seq}(\operatorname{subspace\text{-}ms}(\operatorname{cc\text{-}ms}, a), q) \wedge \\ \quad \forall k \in \mathbb{N}, z \in a.\; (t(k)(z) \in \mathbb{C}) \wedge \\ \quad \forall x \in a, s. \\ \quad \quad \operatorname{pos\text{-}rr}(s) \Rightarrow  \\ \quad \quad \quad \exists l. \\ \quad \quad \quad \quad \operatorname{pos\text{-}rr}(l) \wedge \\ \quad \quad \quad \quad \forall k \in \mathbb{N}, y \in a. \\ \quad \quad \quad \quad \quad (\operatorname{magnitude}((y - x)) < l) \Rightarrow (\operatorname{magnitude}((t(k)(y) - t(k)(x))) < s) \wedge \\ \quad \forall j \in \mathbb{N}. \exists \ell \in \mathbb{C}. \\ \quad \quad \operatorname{converges\text{-}to}(\operatorname{cc\text{-}ms}, (\lambda \mathit{cmk} \in \mathbb{N}.\; t(\mathit{cmk})(q(j))), \ell) \wedge \\ \quad \operatorname{pos\text{-}rr}(s) \Rightarrow \\ \quad \quad \exists p \in \mathbb{N}. \forall i, j \in \mathbb{N}, z \in a. \\ \quad \quad \quad ((p \leq i) \wedge (p \leq j)) \Rightarrow (\operatorname{magnitude}((t(i)(z) - t(j)(z))) < s) \end{array}}}
\vnbprop{cc-pointwise-diagonal}{\vnbstmt{\begin{array}{@{}l@{}} \forall u, s, d. \\ \quad u \subseteq \mathbb{C} \wedge \\ \quad \forall k \in \mathbb{N}, z \in u.\; (s(k)(z) \in \mathbb{C}) \wedge \\ \quad \forall z \in u. \exists a \in \mathbb{R}. \forall k \in \mathbb{N}. \\ \quad \quad \operatorname{magnitude}(s(k)(z)) \leq a \wedge \\ \quad \operatorname{is\text{-}dense\text{-}seq}(\operatorname{subspace\text{-}ms}(\operatorname{cc\text{-}ms}, u), d) \Rightarrow \\ \quad \quad \exists i. \\ \quad \quad \quad \operatorname{strictly\text{-}mono\text{-}nn}(i) \wedge \\ \quad \quad \quad \forall j \in \mathbb{N}. \exists \ell \in \mathbb{C}. \\ \quad \quad \quad \quad \operatorname{converges\text{-}to}(\operatorname{cc\text{-}ms}, (\lambda \mathit{cmk} \in \mathbb{N}.\; s(i(\mathit{cmk}))(d(j))), \ell) \end{array}}}
\vnbprop{cc-exhaust-dense-seqs}{\vnbstmt{\begin{array}{@{}l@{}} \forall u.; \forall p \in u. \\ \quad \operatorname{is\text{-}open}(\operatorname{nf\text{-}metric\text{-}space}(\operatorname{cc\text{-}normed\text{-}field}), u) \Rightarrow  \\ \quad \quad \exists a \in \mathbb{N}, q \in (\mathbb{N} \to (\mathbb{N} \to \mathbb{C})). \forall n \in \mathbb{N}. \\ \quad \quad \quad \operatorname{is\text{-}dense\text{-}seq}(\operatorname{subspace\text{-}ms}(\operatorname{cc\text{-}ms}, \operatorname{cc\text{-}exhaust}(u, (n + a))), q(n)) \end{array}}}
\vnbprop{cc-exhaust-dense-in-u}{\vnbstmt{\begin{array}{@{}l@{}} \forall u.; \forall a \in \mathbb{N}, q \in (\mathbb{N} \to (\mathbb{N} \to \mathbb{C})). \\ \quad \operatorname{is\text{-}open}(\operatorname{nf\text{-}metric\text{-}space}(\operatorname{cc\text{-}normed\text{-}field}), u) \wedge \\ \quad \forall n \in \mathbb{N}.\; \operatorname{is\text{-}dense\text{-}seq}(\operatorname{subspace\text{-}ms}(\operatorname{cc\text{-}ms}, \operatorname{cc\text{-}exhaust}(u, (n + a))), q(n)) \Rightarrow \\ \quad \quad \operatorname{is\text{-}dense\text{-}seq}(\operatorname{subspace\text{-}ms}(\operatorname{cc\text{-}ms}, u), (\lambda m \in \mathbb{N}.\; q(\operatorname{nnfst}(m))(\operatorname{nnsnd}(m)))) \end{array}}}
\vnbprop{cc-montel-tail}{\vnbstmt{\begin{array}{@{}l@{}} \forall u.; \forall s \in (\mathbb{N} \to \operatorname{cc\text{-}hol\text{-}fns}(u)). \\ \quad \operatorname{is\text{-}open}(\operatorname{nf\text{-}metric\text{-}space}(\operatorname{cc\text{-}normed\text{-}field}), u) \wedge \\ \quad \forall k \in \mathbb{N}, a. \\ \quad \quad \operatorname{pos\text{-}rr}(a) \Rightarrow  \\ \quad \quad \quad \exists p \in \mathbb{N}. \forall i, j \in \mathbb{N}, z \in \operatorname{cc\text{-}exhaust}(u, k). \\ \quad \quad \quad \quad ((p \leq i) \wedge (p \leq j)) \Rightarrow (\operatorname{magnitude}((s(i)(z) - s(j)(z))) < a) \Rightarrow \\ \quad \quad \exists g \in \operatorname{cc\text{-}hol\text{-}fns}(u).\; \operatorname{unif\text{-}on\text{-}compacta}(u, s, g) \end{array}}}
\vnbprop{cc-montel-subsequence}{\vnbstmt{\begin{array}{@{}l@{}} \forall u.; \forall s \in (\mathbb{N} \to \operatorname{cc\text{-}hol\text{-}fns}(u)). \\ \quad (\operatorname{is\text{-}open}(\operatorname{nf\text{-}metric\text{-}space}(\operatorname{cc\text{-}normed\text{-}field}), u) \wedge \operatorname{cc\text{-}hol\text{-}bounded}(u, \operatorname{image}(s, \mathbb{N}))) \Rightarrow  \\ \quad \quad \exists a. \\ \quad \quad \quad \operatorname{strictly\text{-}mono\text{-}nn}(a) \wedge \\ \quad \quad \quad \forall k \in \mathbb{N}, b. \\ \quad \quad \quad \quad \operatorname{pos\text{-}rr}(b) \Rightarrow  \\ \quad \quad \quad \quad \quad \exists p \in \mathbb{N}. \forall i, j \in \mathbb{N}, z \in \operatorname{cc\text{-}exhaust}(u, k). \\ \quad \quad \quad \quad \quad \quad ((p \leq i) \wedge (p \leq j)) \Rightarrow  \\ \quad \quad \quad \quad \quad \quad \quad \operatorname{magnitude}((\operatorname{subseq}(s, a)(i)(z) - \operatorname{subseq}(s, a)(j)(z))) < b \end{array}}}
\vnbprop{cc-montel}{\vnbstmt{\begin{array}{@{}l@{}} \forall u.; \forall s \in (\mathbb{N} \to \operatorname{cc\text{-}hol\text{-}fns}(u)). \\ \quad (\operatorname{is\text{-}open}(\operatorname{nf\text{-}metric\text{-}space}(\operatorname{cc\text{-}normed\text{-}field}), u) \wedge \operatorname{cc\text{-}hol\text{-}bounded}(u, \operatorname{image}(s, \mathbb{N}))) \Rightarrow  \\ \quad \quad \exists i. \\ \quad \quad \quad (\operatorname{strictly\text{-}mono\text{-}nn}(i) \wedge \exists g \in \operatorname{cc\text{-}hol\text{-}fns}(u).\; \operatorname{unif\text{-}on\text{-}compacta}(u, \operatorname{subseq}(s, i), g)) \end{array}}}
\vnbprop{cc-fr-conv-iff-unif-on-compacta}{\vnbstmt{\begin{array}{@{}l@{}} \forall u.; \forall s \in (\mathbb{N} \to \operatorname{cc\text{-}cont\text{-}fns}(u)), f \in \operatorname{cc\text{-}cont\text{-}fns}(u). \\ \quad \operatorname{is\text{-}open}(\operatorname{nf\text{-}metric\text{-}space}(\operatorname{cc\text{-}normed\text{-}field}), u) \Rightarrow  \\ \quad \quad \operatorname{fr\text{-}conv}(\operatorname{cc\text{-}fn\text{-}module}(u, \operatorname{cc\text{-}cont\text{-}fns}(u)), \operatorname{cc\text{-}exhaust\text{-}seminorms}(u, \operatorname{cc\text{-}cont\text{-}fns}(u)), s, f) \\ \quad \quad \Leftrightarrow\; \operatorname{unif\text{-}on\text{-}compacta}(u, s, f) \end{array}}}
\vnbprop{cc-deriv-continuous-on-hol}{\vnbstmt{\begin{array}{@{}l@{}} \forall u.; \forall s \in (\mathbb{N} \to \operatorname{cc\text{-}hol\text{-}fns}(u)), f \in \operatorname{cc\text{-}hol\text{-}fns}(u). \\ \quad \operatorname{is\text{-}open}(\operatorname{nf\text{-}metric\text{-}space}(\operatorname{cc\text{-}normed\text{-}field}), u) \wedge \\ \quad \operatorname{fr\text{-}conv}(\operatorname{cc\text{-}fn\text{-}module}(u, \operatorname{cc\text{-}cont\text{-}fns}(u)), \operatorname{cc\text{-}exhaust\text{-}seminorms}(u, \operatorname{cc\text{-}cont\text{-}fns}(u)), s, f) \Rightarrow \\ \quad \quad \operatorname{fr\text{-}conv}(\operatorname{cc\text{-}fn\text{-}module}(u, \operatorname{cc\text{-}cont\text{-}fns}(u)), \\ \quad \quad \quad \operatorname{cc\text{-}exhaust\text{-}seminorms}(u, \operatorname{cc\text{-}cont\text{-}fns}(u)), \\ \quad \quad \quad (\lambda k \in \mathbb{N}.\; (\lambda w \in u.\; \operatorname{deriv\text{-}on}(\operatorname{cc\text{-}normed\text{-}field}, u, s(k), w))), \\ \quad \quad \quad (\lambda w \in u.\; \operatorname{deriv\text{-}on}(\operatorname{cc\text{-}normed\text{-}field}, u, f, w))) \end{array}}}

\vnbbookfile{theorem-library/ca5-series}{theorem-library/ca5-series}
\noindent{\small ca5-series.scm -- CHAPTER 5 OF THE COMPLEX ANALYSIS NOTES, section 5.3: DIFFERENTIATION AS A LINEAR OPERATOR, Lemma 5.21, Proposition 5.22 and the linearity half of Proposition 5.20. Agent CA-5, 2026-10-14.}\par
\vnbprop{cc-exhaust-radius-quarter}{\vnbstmt{\forall n \in \mathbb{N}.\; (\operatorname{recip}((((4 \cdot n) + 3) + 1)) = (1/4 \cdot \operatorname{recip}((n + 1))))}}
\vnbprop{cc-exhaust-ball-inside}{\vnbstmt{\begin{array}{@{}l@{}} \forall u.; \forall n \in \mathbb{N}, y \in \operatorname{cc\text{-}exhaust}(u, n). \\ \quad \operatorname{ball}(\operatorname{cc\text{-}ms}, y, (1/2 \cdot \operatorname{recip}((n + 1)))) \subseteq \operatorname{cc\text{-}exhaust}(u, ((4 \cdot n) + 3)) \end{array}}}
\vnbprop{cc-deriv-seminorm-estimate}{\vnbstmt{\begin{array}{@{}l@{}} \forall u, f.; \forall n \in \mathbb{N}, m \in \mathbb{R}, y \in \operatorname{cc\text{-}exhaust}(u, n). \\ \quad \operatorname{is\text{-}open}(\operatorname{nf\text{-}metric\text{-}space}(\operatorname{cc\text{-}normed\text{-}field}), u) \wedge \\ \quad \operatorname{holomorphic\text{-}on}(u, f) \wedge \\ \quad \forall \mathit{cmy} \in \operatorname{cc\text{-}exhaust}(u, ((4 \cdot n) + 3)).\; (\operatorname{magnitude}(f(\mathit{cmy})) \leq m) \Rightarrow \\ \quad \quad \operatorname{magnitude}(\operatorname{deriv\text{-}on}(\operatorname{cc\text{-}normed\text{-}field}, u, f, y)) \leq (m \cdot (4 \cdot (n + 1))) \end{array}}}
\vnbprop{cc-derived-series-normally-convergent}{\vnbstmt{\begin{array}{@{}l@{}} \forall u.; \forall g \in (\mathbb{N} \to \operatorname{cc\text{-}hol\text{-}fns}(u)), n \in \mathbb{N}. \\ \quad \operatorname{is\text{-}open}(\operatorname{nf\text{-}metric\text{-}space}(\operatorname{cc\text{-}normed\text{-}field}), u) \wedge \\ \quad \forall n \in \mathbb{N}.\; \operatorname{is\text{-}normally\text{-}convergent}(g, \operatorname{cc\text{-}exhaust}(u, n)) \Rightarrow \\ \quad \quad \operatorname{is\text{-}normally\text{-}convergent}((\lambda i \in \mathbb{N}.\; (\lambda z \in u.\; \operatorname{deriv\text{-}on}(\operatorname{cc\text{-}normed\text{-}field}, u, g(i), z))), \\ \quad \quad \quad \operatorname{cc\text{-}exhaust}(u, n)) \end{array}}}
\vnbprop{cc-normal-convergence-pointwise}{\vnbstmt{\begin{array}{@{}l@{}} \forall g, s.; \forall z \in s. \\ \quad (\forall i \in \mathbb{N}, y \in s.\; (g(i)(y) \in \mathbb{C}) \wedge \operatorname{is\text{-}normally\text{-}convergent}(g, s)) \Rightarrow  \\ \quad \quad \operatorname{cc\text{-}series\text{-}converges}((\lambda i \in \mathbb{N}.\; g(i)(z))) \end{array}}}
\vnbprop{cc-normal-convergence-uniform}{\vnbstmt{\begin{array}{@{}l@{}} \forall g, s, a. \\ \quad (\forall i \in \mathbb{N}, y \in s.\; (g(i)(y) \in \mathbb{C}) \wedge (\operatorname{is\text{-}normally\text{-}convergent}(g, s) \wedge \operatorname{pos\text{-}rr}(a))) \Rightarrow  \\ \quad \quad \exists p \in \mathbb{N}. \forall n \in \mathbb{N}, z \in s. \\ \quad \quad \quad (p \leq n) \Rightarrow  \\ \quad \quad \quad \quad \operatorname{magnitude}((\operatorname{cc\text{-}series\text{-}limit}((\lambda i \in \mathbb{N}.\; g(i)(z))) - \operatorname{cc\text{-}series\text{-}partial\text{-}sum}((\lambda i \in \mathbb{N}.\; g(i)(z)), n))) \\ \quad \quad \quad \quad \leq\; a \end{array}}}
\vnbprop{cc-partial-sums-diff}{\vnbstmt{\begin{array}{@{}l@{}} \forall n \in \mathbb{N}, u.; \forall g \in (\mathbb{N} \to \operatorname{cc\text{-}hol\text{-}fns}(u)), z \in u. \\ \quad \operatorname{is\text{-}open}(\operatorname{nf\text{-}metric\text{-}space}(\operatorname{cc\text{-}normed\text{-}field}), u) \Rightarrow  \\ \quad \quad \operatorname{is\text{-}diff\text{-}on}(\operatorname{cc\text{-}normed\text{-}field}, u, \\ \quad \quad \quad (\lambda w \in u.\; \operatorname{cc\text{-}series\text{-}partial\text{-}sum}((\lambda i \in \mathbb{N}.\; g(i)(w)), n)), \\ \quad \quad \quad z, \\ \quad \quad \quad \operatorname{cc\text{-}series\text{-}partial\text{-}sum}((\lambda \mathit{cdi} \in \mathbb{N}.\; \operatorname{deriv\text{-}on}(\operatorname{cc\text{-}normed\text{-}field}, u, g(\mathit{cdi}), z)), n)) \end{array}}}
\vnbprop{cc-series-holomorphic}{\vnbstmt{\begin{array}{@{}l@{}} \forall u.; \forall g \in (\mathbb{N} \to \operatorname{cc\text{-}hol\text{-}fns}(u)). \\ \quad \operatorname{is\text{-}open}(\operatorname{nf\text{-}metric\text{-}space}(\operatorname{cc\text{-}normed\text{-}field}), u) \wedge \\ \quad \forall n \in \mathbb{N}.\; \operatorname{is\text{-}normally\text{-}convergent}(g, \operatorname{cc\text{-}exhaust}(u, n)) \Rightarrow \\ \quad \quad \operatorname{holomorphic\text{-}on}(u, (\lambda w \in u.\; \operatorname{cc\text{-}series\text{-}limit}((\lambda i \in \mathbb{N}.\; g(i)(w))))) \wedge \\ \quad \quad \forall z \in u. \\ \quad \quad \quad \operatorname{deriv\text{-}on}(\operatorname{cc\text{-}normed\text{-}field}, u, (\lambda w \in u.\; \operatorname{cc\text{-}series\text{-}limit}((\lambda i \in \mathbb{N}.\; g(i)(w)))), z) \\ \quad \quad \quad =\; \operatorname{cc\text{-}series\text{-}limit}((\lambda \mathit{cdi} \in \mathbb{N}.\; \operatorname{deriv\text{-}on}(\operatorname{cc\text{-}normed\text{-}field}, u, g(\mathit{cdi}), z))) \end{array}}}
\vnbprop{cc-deriv-linear}{\vnbstmt{\begin{array}{@{}l@{}} \forall u, f, g.; \forall a, b \in \mathbb{C}, z \in u. \\ \quad (\operatorname{is\text{-}open}(\operatorname{nf\text{-}metric\text{-}space}(\operatorname{cc\text{-}normed\text{-}field}), u) \wedge (\operatorname{holomorphic\text{-}on}(u, f) \wedge \operatorname{holomorphic\text{-}on}(u, g))) \Rightarrow  \\ \quad \quad \operatorname{deriv\text{-}on}(\operatorname{cc\text{-}normed\text{-}field}, u, (\lambda w \in u.\; ((a \cdot f(w)) + (b \cdot g(w)))), z) \\ \quad \quad =\; (a \cdot \operatorname{deriv\text{-}on}(\operatorname{cc\text{-}normed\text{-}field}, u, f, z)) + (b \cdot \operatorname{deriv\text{-}on}(\operatorname{cc\text{-}normed\text{-}field}, u, g, z)) \end{array}}}

\part{Measure}
\noindent Sigma-algebras, outer measures and the Caratheodory construction; measures and measurable maps; the integral of simple and measurable functions; sums over the extended nonnegative reals.
\par\noindent\textit{Builds on:} Analysis.
\par\noindent{\small 40 files.}

\vnbbookfile{structure-library/sigma-algebra}{structure-library/sigma-algebra}
\noindent{\small sigma-algebra.scm -- ALGEBRAS and SIGMA-ALGEBRAS of sets.  Vocabulary only: the two definitions of Chapter 1 of the measure-theory notes, cross-checked against Rudin.  Nothing is asserted here.}\par
\par\medskip\noindent\textbf{Definition} (\texttt{is-sigma-algebra}).\;For all $\omega$ and $\mathit{ca}$, 
\begin{equation*}
\begin{array}{@{}l@{}} \operatorname{is\text{-}sigma\text{-}algebra}(\omega, \mathit{ca}) \\ \Leftrightarrow\; \omega \in \mathbf{Set} \wedge \\ \mathit{ca} \subseteq \operatorname{power}(\omega) \wedge \\ \omega \in \mathit{ca} \wedge \\ \forall \operatorname{a\_} \in \mathit{ca}.\; (\operatorname{complement\text{-}in}(\omega, \operatorname{a\_}) \in \mathit{ca}) \wedge \\ \forall \operatorname{f\_} \in (\mathbb{N} \to \mathit{ca}).\; (\operatorname{big\text{-}union}(\operatorname{n\_}, \mathbb{N}, \operatorname{f\_}(\operatorname{n\_})) \in \mathit{ca}) \end{array}
\end{equation*}

\par\medskip\noindent\textbf{Definition} (\texttt{is-algebra-of-sets}).\;For all $\omega$ and $\mathit{ca}$, 
\begin{equation*}
\begin{array}{@{}l@{}} \operatorname{is\text{-}algebra\text{-}of\text{-}sets}(\omega, \mathit{ca}) \\ \Leftrightarrow\; \omega \in \mathbf{Set} \wedge \\ \mathit{ca} \subseteq \operatorname{power}(\omega) \wedge \\ \exists \operatorname{a\_}.\; (\operatorname{a\_} \in \mathit{ca}) \wedge \\ \forall \operatorname{a\_} \in \mathit{ca}.\; (\operatorname{complement\text{-}in}(\omega, \operatorname{a\_}) \in \mathit{ca}) \wedge \\ \forall \operatorname{a\_}, \operatorname{b\_} \in \mathit{ca}.\; ((\operatorname{a\_} \cup \operatorname{b\_}) \in \mathit{ca}) \wedge \\ \forall \operatorname{a\_}, \operatorname{b\_} \in \mathit{ca}.\; ((\operatorname{a\_} \cap \operatorname{b\_}) \in \mathit{ca}) \end{array}
\end{equation*}


\vnbbookfile{structure-library/measure}{structure-library/measure}
\noindent{\small measure.scm -- OUTER MEASURES, CARATHEODORY MEASURABILITY, MEASURES. Vocabulary.  The Caratheodory construction's five statements were asserted here with `reference' warrants until 2026-10-07; they are PROVEN, statements unchanged, in theorem-library/caratheodory.scm (MEASURE-1), and this file}\par
\par\medskip\noindent\textbf{Definition} (\texttt{is-measure}).\;For all $\omega$, $\mathit{ca}$ and $\mu$, 
\begin{equation*}
\begin{array}{@{}l@{}} \operatorname{is\text{-}measure}(\omega, \mathit{ca}, \mu) \\ \Leftrightarrow\; \operatorname{is\text{-}sigma\text{-}algebra}(\omega, \mathit{ca}) \wedge \\ \mu \in (\mathit{ca} \to \operatorname{rr\text{-}pos\text{-}star}) \wedge \\ \mu(\emptyset) = 0 \wedge \\ \forall \operatorname{f\_} \in (\mathbb{N} \to \mathit{ca}). \\ \quad \forall \operatorname{i\_}, \operatorname{j\_} \in \mathbb{N}. \\ \quad \quad \neg (\operatorname{i\_} = \operatorname{j\_}) \Rightarrow ((\operatorname{f\_}(\operatorname{i\_}) \cap \operatorname{f\_}(\operatorname{j\_})) = \emptyset) \Rightarrow \\ \quad \quad \mu(\operatorname{big\text{-}union}(\operatorname{n\_}, \mathbb{N}, \operatorname{f\_}(\operatorname{n\_}))) = \operatorname{esum}((\lambda \operatorname{n\_} \in \mathbb{N}.\; \mu(\operatorname{f\_}(\operatorname{n\_})))) \end{array}
\end{equation*}

\par\medskip\noindent\textbf{Definition} (\texttt{is-caratheodory-measurable}).\;For all $\omega$, $\mu$ and $\mathit{es}$, 
\begin{equation*}
\begin{array}{@{}l@{}} \operatorname{is\text{-}caratheodory\text{-}measurable}(\omega, \mu, \mathit{es}) \\ \Leftrightarrow\; \mathit{es} \subseteq \omega \wedge \\ \forall \operatorname{a\_} \in \operatorname{power}(\omega). \\ \quad \mu(\operatorname{a\_}) = \operatorname{eplus}(\mu((\operatorname{a\_} \cap \mathit{es})), \mu(\operatorname{difference}(\operatorname{a\_}, \mathit{es}))) \end{array}
\end{equation*}

\par\medskip\noindent\textbf{Definition} (\texttt{is-outer-measure}).\;For all $\omega$ and $\mu$, 
\begin{equation*}
\begin{array}{@{}l@{}} \operatorname{is\text{-}outer\text{-}measure}(\omega, \mu) \\ \Leftrightarrow\; \omega \in \mathbf{Set} \wedge \\ \mu \in (\operatorname{power}(\omega) \to \operatorname{rr\text{-}pos\text{-}star}) \wedge \\ \mu(\emptyset) = 0 \wedge \\ \forall \operatorname{a\_}, \operatorname{b\_} \in \operatorname{power}(\omega). \\ \quad (\operatorname{a\_} \subseteq \operatorname{b\_}) \Rightarrow (\mu(\operatorname{a\_}) \leq \mu(\operatorname{b\_})) \wedge \\ \forall \operatorname{f\_} \in (\mathbb{N} \to \operatorname{power}(\omega)). \\ \quad \mu(\operatorname{big\text{-}union}(\operatorname{n\_}, \mathbb{N}, \operatorname{f\_}(\operatorname{n\_}))) \leq \operatorname{esum}((\lambda \operatorname{n\_} \in \mathbb{N}.\; \mu(\operatorname{f\_}(\operatorname{n\_})))) \end{array}
\end{equation*}

\par\medskip\noindent\textbf{Definition} (\texttt{sigma-generated}).\;$\operatorname{sigma\text{-}generated}(\omega, \mathit{ce}) :=$
\begin{equation*}
\{\operatorname{a\_} \in \operatorname{power}(\omega) : \forall \operatorname{cb\_}.\; ((\operatorname{is\text{-}sigma\text{-}algebra}(\omega, \operatorname{cb\_}) \wedge (\mathit{ce} \subseteq \operatorname{cb\_})) \Rightarrow (\operatorname{a\_} \in \operatorname{cb\_}))\}
\end{equation*}

\par\medskip\noindent\textbf{Definition} (\texttt{caratheodory-sets}).\;$\operatorname{caratheodory\text{-}sets}(\omega, \mu) :=$
\begin{equation*}
\{\operatorname{es\_} \in \operatorname{power}(\omega) : \operatorname{is\text{-}caratheodory\text{-}measurable}(\omega, \mu, \operatorname{es\_})\}
\end{equation*}

\par\medskip\noindent\textbf{Definition} (\texttt{is-measurable-map}).\;For all $\omega$, $\mathit{ca}$, $\omega_{2}$, $\mathit{cb}$ and $f$, 
\begin{equation*}
\begin{array}{@{}l@{}} \operatorname{is\text{-}measurable\text{-}map}(\omega, \mathit{ca}, \omega_{2}, \mathit{cb}, f) \\ \Leftrightarrow\; \operatorname{is\text{-}sigma\text{-}algebra}(\omega, \mathit{ca}) \wedge \\ \operatorname{is\text{-}sigma\text{-}algebra}(\omega_{2}, \mathit{cb}) \wedge \\ f \in (\omega \to \omega_{2}) \wedge \\ \forall \operatorname{b\_} \in \mathit{cb}.\; (\{\operatorname{x\_} \in \omega : (f(\operatorname{x\_}) \in \operatorname{b\_})\} \in \mathit{ca}) \end{array}
\end{equation*}

\vnbdef{sigma-generated-membership}{\vnbstmt{\begin{array}{@{}l@{}} \forall a, e, \alpha. \\ \quad \alpha \in \operatorname{sigma\text{-}generated}(a, e) \\ \quad \Leftrightarrow\; ((\alpha \in \operatorname{power}(a)) \wedge \forall b.\; ((\operatorname{is\text{-}sigma\text{-}algebra}(a, b) \wedge (e \subseteq b)) \Rightarrow (\alpha \in b))) \end{array}}}
\vnbdef{caratheodory-sets-membership}{\vnbstmt{\begin{array}{@{}l@{}} \forall a, u, s. \\ \quad (s \in \operatorname{caratheodory\text{-}sets}(a, u)) \Leftrightarrow ((s \in \operatorname{power}(a)) \wedge \operatorname{is\text{-}caratheodory\text{-}measurable}(a, u, s)) \end{array}}}

\vnbbookfile{structure-library/measurable-space}{structure-library/measurable-space}
\noindent{\small measurable-space.scm -- MEASURABLE-SPACE and MEASURE-SPACE as declared STRUCTURES.}\par
\vnbstructdef{measurable-space}{A measurable space $s$ is a 2-tuple whose components are listed below.\par
\begin{enumerate}[nosep]
\item $\operatorname{pts}(s) = \operatorname{nth}(1, s)$, a carrier (\texttt{carriers}): $\operatorname{pts}(s) \in \mathbf{Set}$.
\item $\sigma(s) = \operatorname{nth}(2, s)$, a family of subsets (\texttt{family}): $\sigma(s) \in \operatorname{power}(\operatorname{power}(\operatorname{pts}(s)))$.
\end{enumerate}
The laws (the \texttt{property} and \texttt{law} clauses):
\begin{itemize}[nosep]
\item $\operatorname{is\text{-}sigma\text{-}algebra}(\operatorname{pts}(s), \sigma(s))$
\end{itemize}
The characteristic law (the axiom \texttt{is-measurable-space}): $\operatorname{is\text{-}measurable\text{-}space}(s)$ holds if and only if $\operatorname{length}(s) = 2$, every component has the typing listed with it, and every law holds.
}
\vnbstructdef{measure-space}{A measure space $s$ is a 3-tuple whose components are listed below.\par
\begin{enumerate}[nosep]
\item $\operatorname{pts}(s) = \operatorname{nth}(1, s)$, a carrier (\texttt{carriers}): $\operatorname{pts}(s) \in \mathbf{Set}$.
\item $\sigma(s) = \operatorname{nth}(2, s)$, a family of subsets (\texttt{family}): $\sigma(s) \in \operatorname{power}(\operatorname{power}(\operatorname{pts}(s)))$.
\item $\operatorname{meas}(s) = \operatorname{nth}(3, s)$, a function on a family (\texttt{family-fun}): $\operatorname{meas}(s) \in (\sigma(s) \to \operatorname{rr\text{-}pos\text{-}star})$.
\end{enumerate}
The laws (the \texttt{property} and \texttt{law} clauses):
\begin{itemize}[nosep]
\item $\operatorname{is\text{-}measure}(\operatorname{pts}(s), \sigma(s), \operatorname{meas}(s))$
\end{itemize}
The characteristic law (the axiom \texttt{is-measure-space}): $\operatorname{is\text{-}measure\text{-}space}(s)$ holds if and only if $\operatorname{length}(s) = 3$, every component has the typing listed with it, and every law holds.
}
\par\medskip\noindent\textbf{Definition} (\texttt{hom-measurable-space}).\;$\operatorname{hom\text{-}measurable\text{-}space}(a, b) :=$
\begin{equation*}
\{\operatorname{homf\_} \in (\operatorname{pts}(a) \to \operatorname{pts}(b)) : \operatorname{is\text{-}hom\text{-}measurable\text{-}space}(a, b, \operatorname{homf\_})\}
\end{equation*}

\par\medskip\noindent\textbf{Definition} (\texttt{hom-measure-space}).\;$\operatorname{hom\text{-}measure\text{-}space}(a, b) :=$
\begin{equation*}
\{\operatorname{homf\_} \in (\operatorname{pts}(a) \to \operatorname{pts}(b)) : \operatorname{is\text{-}hom\text{-}measure\text{-}space}(a, b, \operatorname{homf\_})\}
\end{equation*}

\vnbdef{measurable-space-class}{\vnbstmt{\forall s.\; ((s \in \operatorname{measurable\text{-}space}) \Leftrightarrow \operatorname{is\text{-}measurable\text{-}space}(s))}}
\vnbdef{measure-space-as-measurable-space-is-measurable-space}{\vnbstmt{\begin{array}{@{}l@{}} \forall r. \\ \quad \operatorname{is\text{-}measure\text{-}space}(r) \Rightarrow  \\ \quad \quad \operatorname{is\text{-}measurable\text{-}space}(\operatorname{measure\text{-}space\text{-}as\text{-}measurable\text{-}space}(r)) \end{array}}}
\vnbdef{is-hom-measurable-space-def}{\vnbstmt{\begin{array}{@{}l@{}} \forall s, t, f. \\ \quad \operatorname{is\text{-}hom\text{-}measurable\text{-}space}(s, t, f) \\ \quad \Leftrightarrow\; \operatorname{is\text{-}measurable\text{-}space}(s) \wedge \\ \quad \operatorname{is\text{-}measurable\text{-}space}(t) \wedge \\ \quad f \in (\operatorname{pts}(s) \to \operatorname{pts}(t)) \wedge \\ \quad \forall u \in \sigma(t).\; (\operatorname{preimage}(s, f, u) \in \sigma(s)) \end{array}}}
\vnbdef{is-hom-measure-space-def}{\vnbstmt{\begin{array}{@{}l@{}} \forall s, t, f. \\ \quad \operatorname{is\text{-}hom\text{-}measure\text{-}space}(s, t, f) \\ \quad \Leftrightarrow\; \operatorname{is\text{-}measure\text{-}space}(s) \wedge \\ \quad \operatorname{is\text{-}measure\text{-}space}(t) \wedge \\ \quad f \in (\operatorname{pts}(s) \to \operatorname{pts}(t)) \wedge \\ \quad \forall u \in \sigma(t).\; (\operatorname{preimage}(s, f, u) \in \sigma(s)) \wedge \\ \quad \forall u \in \sigma(t).\; (\operatorname{meas}(t)(u) = \operatorname{meas}(s)(\operatorname{preimage}(s, f, u))) \end{array}}}
\vnbdef{measure-space-class}{\vnbstmt{\forall s.\; ((s \in \operatorname{measure\text{-}space}) \Leftrightarrow \operatorname{is\text{-}measure\text{-}space}(s))}}

\vnbbookfile{structure-library/integral}{structure-library/integral}
\noindent{\small SINCE 2026-10-04 NOTHING IN THIS FILE IS ASSERTED: the fifteen supports that characterised the integral are THEOREMS -- six in theorem-library/integral-laws.scm (M-2), nine in theorem-library/integral-convergence.scm (M-3) -- and INTEGRAL is DEFINED in structure-library/simple-integral.scm (the supremum of the simple}\par
\par\medskip\noindent\textbf{Definition} (\texttt{is-measurable-fn}).\;For all $\omega$, $\mathit{ca}$ and $f$, 
\begin{equation*}
\begin{array}{@{}l@{}} \operatorname{is\text{-}measurable\text{-}fn}(\omega, \mathit{ca}, f) \\ \Leftrightarrow\; \operatorname{is\text{-}sigma\text{-}algebra}(\omega, \mathit{ca}) \wedge \\ f \in (\omega \to \operatorname{rr\text{-}pos\text{-}star}) \wedge \\ \forall \operatorname{alpha\_} \in \mathbb{R}.\; (\{\operatorname{x\_} \in \omega : (\operatorname{alpha\_} < f(\operatorname{x\_}))\} \in \mathit{ca}) \end{array}
\end{equation*}

\par\medskip\noindent\textbf{Definition} (\texttt{ptwise-eplus}).\;$\operatorname{ptwise\text{-}eplus}(\omega, f, g) :=$ $(\lambda \operatorname{x\_} \in \omega.\; \operatorname{eplus}(f(\operatorname{x\_}), g(\operatorname{x\_})))$
\par\medskip\noindent\textbf{Definition} (\texttt{ptwise-etimes}).\;$\operatorname{ptwise\text{-}etimes}(\omega, f, g) :=$ $(\lambda \operatorname{x\_} \in \omega.\; \operatorname{etimes}(f(\operatorname{x\_}), g(\operatorname{x\_})))$
\par\medskip\noindent\textbf{Definition} (\texttt{ptwise-scale}).\;$\operatorname{ptwise\text{-}scale}(\omega, c, f) :=$ $(\lambda \operatorname{x\_} \in \omega.\; \operatorname{etimes}(c, f(\operatorname{x\_})))$
\par\medskip\noindent\textbf{Definition} (\texttt{ptwise-sup}).\;$\operatorname{ptwise\text{-}sup}(\omega, \mathit{cf}) :=$
\begin{equation*}
(\lambda \operatorname{x\_} \in \omega.\; \operatorname{esup}(\operatorname{image}((\lambda \operatorname{n\_} \in \mathbb{N}.\; \operatorname{cf}(\operatorname{n\_})(\operatorname{x\_})), \mathbb{N})))
\end{equation*}

\par\medskip\noindent\textbf{Definition} (\texttt{ptwise-liminf}).\;$\operatorname{ptwise\text{-}liminf}(\omega, \mathit{cf}) :=$
\begin{equation*}
(\lambda \operatorname{x\_} \in \omega.\; \operatorname{eliminf}((\lambda \operatorname{n\_} \in \mathbb{N}.\; \operatorname{cf}(\operatorname{n\_})(\operatorname{x\_}))))
\end{equation*}

\par\medskip\noindent\textbf{Definition} (\texttt{indicator}).\;$\operatorname{indicator}(\omega, a) :=$ $(\lambda \operatorname{x\_} \in \omega.\; \left(\text{if }(\operatorname{x\_} \in a)\text{ then }1\text{ else }0\right))$
\par\medskip\noindent\textbf{Definition} (\texttt{ptwise-le}).\;For all $\omega$, $f$ and $g$, 
\begin{equation*}
\operatorname{ptwise\text{-}le}(\omega, f, g) \Leftrightarrow \forall \operatorname{x\_} \in \omega.\; (f(\operatorname{x\_}) \leq g(\operatorname{x\_}))
\end{equation*}

\par\medskip\noindent\textbf{Definition} (\texttt{is-simple-fn}).\;For all $\omega$, $\mathit{ca}$ and $f$, 
\begin{equation*}
\begin{array}{@{}l@{}} \operatorname{is\text{-}simple\text{-}fn}(\omega, \mathit{ca}, f) \\ \Leftrightarrow\; \operatorname{is\text{-}measurable\text{-}fn}(\omega, \mathit{ca}, f) \wedge \\ |\operatorname{image}(f, \omega)| \in \mathbb{N} \wedge \\ \forall \operatorname{x\_} \in \omega.\; (f(\operatorname{x\_}) \in \mathbb{R}) \end{array}
\end{equation*}


\vnbbookfile{structure-library/simple-integral}{structure-library/simple-integral}
\noindent{\small simple-integral.scm -- THE INTEGRAL OF A [0,+inf]-VALUED FUNCTION, DEFINED. Definitions only; nothing here is asserted and nothing is proved.}\par
\par\medskip\noindent\textbf{Definition} (\texttt{integral}).\;$\operatorname{integral}(\omega, \mathit{ca}, \mu, f) :=$
\begin{equation*}
\operatorname{esup}(\{\operatorname{iny\_} \in \operatorname{rr\text{-}pos\text{-}star} : \exists \operatorname{ins\_}.\; (\operatorname{is\text{-}simple\text{-}fn}(\omega, \mathit{ca}, \operatorname{ins\_}) \wedge (\operatorname{ptwise\text{-}le}(\omega, \operatorname{ins\_}, f) \wedge (\operatorname{iny\_} = \operatorname{simple\text{-}integral}(\omega, \mu, \operatorname{ins\_}))))\})
\end{equation*}

\par\medskip\noindent\textbf{Definition} (\texttt{simple-integral}).\;$\operatorname{simple\text{-}integral}(\omega, \mu, s) :=$
\begin{equation*}
\begin{array}{@{}l@{}} \operatorname{finsum}(\operatorname{rr\text{-}pos\text{-}star\text{-}add\text{-}monoid}, \\ \quad (\lambda \operatorname{siy\_} \in \operatorname{image}(s, \omega).\; \operatorname{etimes}(\operatorname{siy\_}, \mu(\{\operatorname{six\_} \in \omega : (s(\operatorname{six\_}) = \operatorname{siy\_})\}))), \\ \quad \operatorname{image}(s, \omega)) \end{array}
\end{equation*}


\vnbbookfile{theorem-library/discrete-space}{theorem-library/discrete-space}
\noindent{\small A file shared across books; the entries that belong to this one.}\par
\noindent{\small discrete-space.scm -- THE FIRST WITNESS OF A MEASURABLE SPACE, and the first outright inhabitant of TOP-SPACE.}\par
\vnbprop{discrete-space-is-measurable-space}{\vnbstmt{\forall a \in \mathbf{Set}.\; \operatorname{is\text{-}measurable\text{-}space}(\operatorname{discrete\text{-}space}(a))}}

\vnbbookfile{theorem-library/rake-measure}{theorem-library/rake-measure}
\noindent{\small rake-measure.scm -- rake batch 5, batch W: the light end of measure theory.}\par
\vnbprop{sigma-algebra-ambient-set}{\vnbstmt{\forall a, b.\; (\operatorname{is\text{-}sigma\text{-}algebra}(a, b) \Rightarrow (a \in \mathbf{Set}))}}
\vnbprop{sigma-algebra-whole-in}{\vnbstmt{\forall a, b.\; (\operatorname{is\text{-}sigma\text{-}algebra}(a, b) \Rightarrow (a \in b))}}
\vnbprop{sigma-algebra-complement-closed}{\vnbstmt{\begin{array}{@{}l@{}} \forall a, b, \alpha. \\ \quad (\operatorname{is\text{-}sigma\text{-}algebra}(a, b) \wedge (\alpha \in b)) \Rightarrow (\operatorname{complement\text{-}in}(a, \alpha) \in b) \end{array}}}
\vnbprop{sigma-algebra-union-closed}{\vnbstmt{\begin{array}{@{}l@{}} \forall a, b, f. \\ \quad (\operatorname{is\text{-}sigma\text{-}algebra}(a, b) \wedge (f \in (\mathbb{N} \to b))) \Rightarrow (\operatorname{big\text{-}union}(n, \mathbb{N}, f(n)) \in b) \end{array}}}
\vnbprop{sigma-generated-in-power}{\vnbstmt{\forall a, e.; \forall \alpha \in \operatorname{sigma\text{-}generated}(a, e).\; (\alpha \in \operatorname{power}(a))}}
\vnbprop{sigma-generated-in-every}{\vnbstmt{\begin{array}{@{}l@{}} \forall a, e, \alpha, b. \\ \quad ((\alpha \in \operatorname{sigma\text{-}generated}(a, e)) \wedge (\operatorname{is\text{-}sigma\text{-}algebra}(a, b) \wedge (e \subseteq b))) \Rightarrow  \\ \quad \quad \alpha \in b \end{array}}}
\vnbprop{sigma-generated-intro}{\vnbstmt{\begin{array}{@{}l@{}} \forall a, e.; \forall \alpha \in \operatorname{power}(a). \\ \quad \forall b.\; ((\operatorname{is\text{-}sigma\text{-}algebra}(a, b) \wedge (e \subseteq b)) \Rightarrow (\alpha \in b)) \Rightarrow  \\ \quad \quad \alpha \in \operatorname{sigma\text{-}generated}(a, e) \end{array}}}
\vnbprop{sigma-generated-is-sigma-algebra}{\vnbstmt{\begin{array}{@{}l@{}} \forall a, e. \\ \quad ((a \in \mathbf{Set}) \wedge (e \subseteq \operatorname{power}(a))) \Rightarrow  \\ \quad \quad \operatorname{is\text{-}sigma\text{-}algebra}(a, \operatorname{sigma\text{-}generated}(a, e)) \end{array}}}
\vnbprop{measurable-map-source-algebra}{\vnbstmt{\begin{array}{@{}l@{}} \forall a, b, a_{2}, c, f. \\ \quad \operatorname{is\text{-}measurable\text{-}map}(a, b, a_{2}, c, f) \Rightarrow \operatorname{is\text{-}sigma\text{-}algebra}(a, b) \end{array}}}
\vnbprop{measurable-map-target-algebra}{\vnbstmt{\begin{array}{@{}l@{}} \forall a, b, a_{2}, c, f. \\ \quad \operatorname{is\text{-}measurable\text{-}map}(a, b, a_{2}, c, f) \Rightarrow \operatorname{is\text{-}sigma\text{-}algebra}(a_{2}, c) \end{array}}}
\vnbprop{measurable-map-in-fun}{\vnbstmt{\begin{array}{@{}l@{}} \forall a, b, a_{2}, c, f. \\ \quad \operatorname{is\text{-}measurable\text{-}map}(a, b, a_{2}, c, f) \Rightarrow (f \in (a \to a_{2})) \end{array}}}
\vnbprop{measurable-map-preimage}{\vnbstmt{\begin{array}{@{}l@{}} \forall a, b, a_{2}, c, f, \beta. \\ \quad (\operatorname{is\text{-}measurable\text{-}map}(a, b, a_{2}, c, f) \wedge (\beta \in c)) \Rightarrow  \\ \quad \quad \{x \in a : (f(x) \in \beta)\} \in b \end{array}}}
\vnbprop{measurable-map-intro}{\vnbstmt{\begin{array}{@{}l@{}} \forall a, b, a_{2}, c, f. \\ \quad \operatorname{is\text{-}sigma\text{-}algebra}(a, b) \wedge \\ \quad \operatorname{is\text{-}sigma\text{-}algebra}(a_{2}, c) \wedge \\ \quad f \in (a \to a_{2}) \wedge \\ \quad \forall \beta \in c.\; (\{x \in a : (f(x) \in \beta)\} \in b) \Rightarrow \\ \quad \quad \operatorname{is\text{-}measurable\text{-}map}(a, b, a_{2}, c, f) \end{array}}}
\vnbprop{measurable-map-compose}{\vnbstmt{\begin{array}{@{}l@{}} \forall a, b, a_{2}, c, a_{3}, d, f, g. \\ \quad (\operatorname{is\text{-}measurable\text{-}map}(a, b, a_{2}, c, f) \wedge \operatorname{is\text{-}measurable\text{-}map}(a_{2}, c, a_{3}, d, g)) \Rightarrow  \\ \quad \quad \operatorname{is\text{-}measurable\text{-}map}(a, b, a_{3}, d, \operatorname{compose}(g, f)) \end{array}}}

\vnbbookfile{theorem-library/rake-measure2}{theorem-library/rake-measure2}
\noindent{\small rake-measure2.scm -- rake batch 5b-F: the indicator function, and the four SET-LEVEL bricks sigma-algebra.scm's header lists as "NEXT (not stated)".}\par
\vnbprop{sigma-algebra-empty-in}{\vnbstmt{\forall a, b.\; (\operatorname{is\text{-}sigma\text{-}algebra}(a, b) \Rightarrow (\emptyset \in b))}}
\vnbprop{sigma-algebra-union-closed-2}{\vnbstmt{\begin{array}{@{}l@{}} \forall a, b, \alpha, \beta. \\ \quad (\operatorname{is\text{-}sigma\text{-}algebra}(a, b) \wedge ((\alpha \in b) \wedge (\beta \in b))) \Rightarrow  \\ \quad \quad (\alpha \cup \beta) \in b \end{array}}}
\vnbprop{sigma-algebra-inter-closed-2}{\vnbstmt{\begin{array}{@{}l@{}} \forall a, b, \alpha, \beta. \\ \quad (\operatorname{is\text{-}sigma\text{-}algebra}(a, b) \wedge ((\alpha \in b) \wedge (\beta \in b))) \Rightarrow  \\ \quad \quad (\alpha \cap \beta) \in b \end{array}}}
\vnbprop{sigma-algebra-difference-closed}{\vnbstmt{\begin{array}{@{}l@{}} \forall a, b, \alpha, \beta. \\ \quad (\operatorname{is\text{-}sigma\text{-}algebra}(a, b) \wedge ((\alpha \in b) \wedge (\beta \in b))) \Rightarrow  \\ \quad \quad \operatorname{complement\text{-}in}(\alpha, \beta) \in b \end{array}}}
\vnbprop{measurable-fn-indicator}{\vnbstmt{\begin{array}{@{}l@{}} \forall a, b, \alpha. \\ \quad (\operatorname{is\text{-}sigma\text{-}algebra}(a, b) \wedge (\alpha \in b)) \Rightarrow  \\ \quad \quad \operatorname{is\text{-}measurable\text{-}fn}(a, b, \operatorname{indicator}(a, \alpha)) \end{array}}}

\vnbbookfile{theorem-library/hom-kinds}{theorem-library/hom-kinds}
\noindent{\small A file shared across books; the entries that belong to this one.}\par
\noindent{\small hom-kinds.scm -- the category laws for the arrows the generic driver of hom-laws.scm does not reach: the FAMILY slot kinds, and the RELATION kind (SETOID, since batch 40; two carriers in batch 39).}\par
\vnbprop{hom-measurable-space-member-iff}{\vnbstmt{\begin{array}{@{}l@{}} \forall a, b, f. \\ \quad (f \in \operatorname{hom\text{-}measurable\text{-}space}(a, b)) \Leftrightarrow ((f \in (\operatorname{pts}(a) \to \operatorname{pts}(b))) \wedge \operatorname{is\text{-}hom\text{-}measurable\text{-}space}(a, b, f)) \end{array}}}
\vnbprop{hom-measurable-space-id}{\vnbstmt{\begin{array}{@{}l@{}} \forall a. \\ \quad \operatorname{is\text{-}measurable\text{-}space}(a) \Rightarrow \operatorname{is\text{-}hom\text{-}measurable\text{-}space}(a, a, \operatorname{id\text{-}fun}(\operatorname{pts}(a))) \end{array}}}
\vnbprop{hom-measurable-space-in-set}{\vnbstmt{\begin{array}{@{}l@{}} \forall a, b. \\ \quad (\operatorname{is\text{-}measurable\text{-}space}(a) \wedge \operatorname{is\text{-}measurable\text{-}space}(b)) \Rightarrow  \\ \quad \quad \operatorname{hom\text{-}measurable\text{-}space}(a, b) \in \mathbf{Set} \end{array}}}
\vnbprop{hom-measure-space-member-iff}{\vnbstmt{\begin{array}{@{}l@{}} \forall a, b, f. \\ \quad (f \in \operatorname{hom\text{-}measure\text{-}space}(a, b)) \Leftrightarrow ((f \in (\operatorname{pts}(a) \to \operatorname{pts}(b))) \wedge \operatorname{is\text{-}hom\text{-}measure\text{-}space}(a, b, f)) \end{array}}}
\vnbprop{hom-measure-space-id}{\vnbstmt{\begin{array}{@{}l@{}} \forall a. \\ \quad \operatorname{is\text{-}measure\text{-}space}(a) \Rightarrow \operatorname{is\text{-}hom\text{-}measure\text{-}space}(a, a, \operatorname{id\text{-}fun}(\operatorname{pts}(a))) \end{array}}}
\vnbprop{hom-measure-space-in-set}{\vnbstmt{\begin{array}{@{}l@{}} \forall a, b. \\ \quad (\operatorname{is\text{-}measure\text{-}space}(a) \wedge \operatorname{is\text{-}measure\text{-}space}(b)) \Rightarrow  \\ \quad \quad \operatorname{hom\text{-}measure\text{-}space}(a, b) \in \mathbf{Set} \end{array}}}

\vnbbookfile{theorem-library/hom-functors}{theorem-library/hom-functors}
\noindent{\small A file shared across books; the entries that belong to this one.}\par
\noindent{\small hom-functors.scm -- the Hom functors: Hom(a, -) covariant, Hom(-, b) contravariant, as maps on hom-sets.}\par
\vnbprop{hom-measurable-space-post-type}{\vnbstmt{\begin{array}{@{}l@{}} \forall a, b, c, h. \\ \quad (\operatorname{is\text{-}measurable\text{-}space}(a) \wedge (h \in \operatorname{hom\text{-}measurable\text{-}space}(b, c))) \Rightarrow  \\ \quad \quad (\lambda f \in \operatorname{hom\text{-}measurable\text{-}space}(a, b).\; \operatorname{compose}(h, f)) \\ \quad \quad \in\; (\operatorname{hom\text{-}measurable\text{-}space}(a, b) \to \operatorname{hom\text{-}measurable\text{-}space}(a, c)) \end{array}}}
\vnbprop{hom-measurable-space-pre-type}{\vnbstmt{\begin{array}{@{}l@{}} \forall a, b, c, k. \\ \quad (\operatorname{is\text{-}measurable\text{-}space}(b) \wedge (k \in \operatorname{hom\text{-}measurable\text{-}space}(c, a))) \Rightarrow  \\ \quad \quad (\lambda f \in \operatorname{hom\text{-}measurable\text{-}space}(a, b).\; \operatorname{compose}(f, k)) \\ \quad \quad \in\; (\operatorname{hom\text{-}measurable\text{-}space}(a, b) \to \operatorname{hom\text{-}measurable\text{-}space}(c, b)) \end{array}}}
\vnbprop{hom-measure-space-post-type}{\vnbstmt{\begin{array}{@{}l@{}} \forall a, b, c, h. \\ \quad (\operatorname{is\text{-}measure\text{-}space}(a) \wedge (h \in \operatorname{hom\text{-}measure\text{-}space}(b, c))) \Rightarrow  \\ \quad \quad (\lambda f \in \operatorname{hom\text{-}measure\text{-}space}(a, b).\; \operatorname{compose}(h, f)) \\ \quad \quad \in\; (\operatorname{hom\text{-}measure\text{-}space}(a, b) \to \operatorname{hom\text{-}measure\text{-}space}(a, c)) \end{array}}}
\vnbprop{hom-measure-space-pre-type}{\vnbstmt{\begin{array}{@{}l@{}} \forall a, b, c, k. \\ \quad (\operatorname{is\text{-}measure\text{-}space}(b) \wedge (k \in \operatorname{hom\text{-}measure\text{-}space}(c, a))) \Rightarrow  \\ \quad \quad (\lambda f \in \operatorname{hom\text{-}measure\text{-}space}(a, b).\; \operatorname{compose}(f, k)) \\ \quad \quad \in\; (\operatorname{hom\text{-}measure\text{-}space}(a, b) \to \operatorname{hom\text{-}measure\text{-}space}(c, b)) \end{array}}}

\vnbbookfile{theorem-library/measure-laws}{theorem-library/measure-laws}
\noindent{\small measure-laws.scm -- Chapter 1 of the measure notes (Thayer, Construction of Measures, Sec. 3-4): the elementary laws of a countably additive measure. Agent M-1 of the October roadmap (docs/roadmap-2026-10-03.md, line (b), decision 4).  Helper prefix mlw-.  25 theorems, every one `modulo 0'.}\par
\vnbprop{rps-prefix-zero}{\vnbstmt{\forall f.\; (\operatorname{finsum}(\operatorname{rr\text{-}pos\text{-}star\text{-}add\text{-}monoid}, f, \mathrm{OS}(0)) = 0)}}
\vnbprop{rps-prefix-succ}{\vnbstmt{\begin{array}{@{}l@{}} \forall f \in (\mathbb{N} \to \operatorname{rr\text{-}pos\text{-}star}), n \in \mathbb{N}. \\ \quad \operatorname{finsum}(\operatorname{rr\text{-}pos\text{-}star\text{-}add\text{-}monoid}, f, \mathrm{OS}((n + 1))) \\ \quad =\; \operatorname{eplus}(\operatorname{finsum}(\operatorname{rr\text{-}pos\text{-}star\text{-}add\text{-}monoid}, f, \mathrm{OS}(n)), f(n)) \end{array}}}
\vnbprop{rps-finsum-mono-set}{\vnbstmt{\begin{array}{@{}l@{}} \forall f, s.; \forall t \in \mathbf{Set}. \\ \quad ((|t| \in \mathbb{N}) \wedge ((s \subseteq t) \wedge \forall z \in t.\; (f(z) \in \operatorname{rr\text{-}pos\text{-}star}))) \Rightarrow  \\ \quad \quad \operatorname{finsum}(\operatorname{rr\text{-}pos\text{-}star\text{-}add\text{-}monoid}, f, s) \leq \operatorname{finsum}(\operatorname{rr\text{-}pos\text{-}star\text{-}add\text{-}monoid}, f, t) \end{array}}}
\vnbprop{esum-nn-prefix-le}{\vnbstmt{\begin{array}{@{}l@{}} \forall f \in (\mathbb{N} \to \operatorname{rr\text{-}pos\text{-}star}), n \in \mathbb{N}. \\ \quad \operatorname{finsum}(\operatorname{rr\text{-}pos\text{-}star\text{-}add\text{-}monoid}, f, \mathrm{OS}(n)) \leq \operatorname{esum}(f) \end{array}}}
\vnbprop{esum-nn-le-of-prefixes}{\vnbstmt{\begin{array}{@{}l@{}} \forall f \in (\mathbb{N} \to \operatorname{rr\text{-}pos\text{-}star}), b \in \operatorname{rr\text{-}pos\text{-}star}. \\ \quad \forall n \in \mathbb{N}.\; (\operatorname{finsum}(\operatorname{rr\text{-}pos\text{-}star\text{-}add\text{-}monoid}, f, \mathrm{OS}(n)) \leq b) \Rightarrow  \\ \quad \quad \operatorname{esum}(f) \leq b \end{array}}}
\vnbprop{rps-prefix-type}{\vnbstmt{\begin{array}{@{}l@{}} \forall f \in (\mathbb{N} \to \operatorname{rr\text{-}pos\text{-}star}), n \in \mathbb{N}. \\ \quad \operatorname{finsum}(\operatorname{rr\text{-}pos\text{-}star\text{-}add\text{-}monoid}, f, \mathrm{OS}(n)) \in \operatorname{rr\text{-}pos\text{-}star} \end{array}}}
\vnbprop{rps-prefix-le-termwise}{\vnbstmt{\begin{array}{@{}l@{}} \forall n \in \mathbb{N}, f, g. \\ \quad ((f \in (\mathbb{N} \to \operatorname{rr\text{-}pos\text{-}star})) \wedge ((g \in (\mathbb{N} \to \operatorname{rr\text{-}pos\text{-}star})) \wedge \forall k \in \mathbb{N}.\; (f(k) \leq g(k)))) \Rightarrow  \\ \quad \quad \operatorname{finsum}(\operatorname{rr\text{-}pos\text{-}star\text{-}add\text{-}monoid}, f, \mathrm{OS}(n)) \\ \quad \quad \leq\; \operatorname{finsum}(\operatorname{rr\text{-}pos\text{-}star\text{-}add\text{-}monoid}, g, \mathrm{OS}(n)) \end{array}}}
\vnbprop{esum-nn-le-termwise}{\vnbstmt{\begin{array}{@{}l@{}} \forall f, g. \\ \quad ((f \in (\mathbb{N} \to \operatorname{rr\text{-}pos\text{-}star})) \wedge ((g \in (\mathbb{N} \to \operatorname{rr\text{-}pos\text{-}star})) \wedge \forall k \in \mathbb{N}.\; (f(k) \leq g(k)))) \Rightarrow  \\ \quad \quad \operatorname{esum}(f) \leq \operatorname{esum}(g) \end{array}}}
\vnbprop{esum-nn-two-support}{\vnbstmt{\begin{array}{@{}l@{}} \forall h \in (\mathbb{N} \to \operatorname{rr\text{-}pos\text{-}star}). \\ \quad \forall q \in \mathbb{N}.\; (h(((q + 1) + 1)) = 0) \Rightarrow  \\ \quad \quad \operatorname{esum}(h) = \operatorname{eplus}(h(0), h(1)) \end{array}}}
\vnbprop{measure-finitely-additive}{\vnbstmt{\begin{array}{@{}l@{}} \forall a, b, u, \alpha, \beta. \\ \quad (\operatorname{is\text{-}measure}(a, b, u) \wedge ((\alpha \in b) \wedge ((\beta \in b) \wedge ((\alpha \cap \beta) = \emptyset)))) \Rightarrow  \\ \quad \quad u((\alpha \cup \beta)) = \operatorname{eplus}(u(\alpha), u(\beta)) \end{array}}}
\vnbprop{measure-monotone}{\vnbstmt{\begin{array}{@{}l@{}} \forall a, b, u, \alpha, \beta. \\ \quad (\operatorname{is\text{-}measure}(a, b, u) \wedge ((\alpha \in b) \wedge ((\beta \in b) \wedge (\alpha \subseteq \beta)))) \Rightarrow  \\ \quad \quad u(\alpha) \leq u(\beta) \end{array}}}
\vnbprop{bu-seg-zero}{\vnbstmt{\forall f.\; (\operatorname{big\text{-}union}(j, \mathrm{OS}(0), f(j)) = \emptyset)}}
\vnbprop{bu-seg-succ}{\vnbstmt{\begin{array}{@{}l@{}} \forall k \in \mathbb{N}, c.; \forall f \in (\mathbb{N} \to c). \\ \quad \operatorname{big\text{-}union}(j, \mathrm{OS}((k + 1)), f(j)) = (\operatorname{big\text{-}union}(j, \mathrm{OS}(k), f(j)) \cup f(k)) \end{array}}}
\vnbprop{sigma-algebra-seg-union-in}{\vnbstmt{\begin{array}{@{}l@{}} \forall k \in \mathbb{N}, a, b, f. \\ \quad (\operatorname{is\text{-}sigma\text{-}algebra}(a, b) \wedge (f \in (\mathbb{N} \to b))) \Rightarrow  \\ \quad \quad \operatorname{big\text{-}union}(j, \mathrm{OS}(k), f(j)) \in b \end{array}}}
\vnbprop{disjointify-type}{\vnbstmt{\begin{array}{@{}l@{}} \forall a, b, f. \\ \quad (\operatorname{is\text{-}sigma\text{-}algebra}(a, b) \wedge (f \in (\mathbb{N} \to b))) \Rightarrow  \\ \quad \quad (\lambda k \in \mathbb{N}.\; \operatorname{complement\text{-}in}(f(k), \operatorname{big\text{-}union}(j, \mathrm{OS}(k), f(j)))) \in (\mathbb{N} \to b) \end{array}}}
\vnbprop{disjointify-disjoint}{\vnbstmt{\begin{array}{@{}l@{}} \forall f, c, i, j. \\ \quad ((f \in (\mathbb{N} \to c)) \wedge ((i \in \mathbb{N}) \wedge ((j \in \mathbb{N}) \wedge \neg (i = j)))) \Rightarrow  \\ \quad \quad ((\lambda k \in \mathbb{N}.\; \operatorname{complement\text{-}in}(f(k), \operatorname{big\text{-}union}(\mathit{uj}, \mathrm{OS}(k), f(\mathit{uj}))))(i) \cap (\lambda k \in \mathbb{N}.\; \operatorname{complement\text{-}in}(f(k), \operatorname{big\text{-}union}(\mathit{uj}, \mathrm{OS}(k), f(\mathit{uj}))))(j)) \\ \quad \quad =\; \emptyset \end{array}}}
\vnbprop{disjointify-union}{\vnbstmt{\begin{array}{@{}l@{}} \forall f, c. \\ \quad (f \in (\mathbb{N} \to c)) \Rightarrow  \\ \quad \quad \operatorname{big\text{-}union}(n, \mathbb{N}, (\lambda k \in \mathbb{N}.\; \operatorname{complement\text{-}in}(f(k), \operatorname{big\text{-}union}(j, \mathrm{OS}(k), f(j))))(n)) \\ \quad \quad =\; \operatorname{big\text{-}union}(n, \mathbb{N}, f(n)) \end{array}}}
\vnbprop{measure-countably-subadditive}{\vnbstmt{\begin{array}{@{}l@{}} \forall a, b, u, f. \\ \quad (\operatorname{is\text{-}measure}(a, b, u) \wedge (f \in (\mathbb{N} \to b))) \Rightarrow  \\ \quad \quad u(\operatorname{big\text{-}union}(n, \mathbb{N}, f(n))) \leq \operatorname{esum}((\lambda n \in \mathbb{N}.\; u(f(n)))) \end{array}}}
\vnbprop{disjointify-prefix}{\vnbstmt{\begin{array}{@{}l@{}} \forall n \in \mathbb{N}, a, b, u, f. \\ \quad (\operatorname{is\text{-}measure}(a, b, u) \wedge (f \in (\mathbb{N} \to b))) \Rightarrow  \\ \quad \quad \operatorname{finsum}(\operatorname{rr\text{-}pos\text{-}star\text{-}add\text{-}monoid}, \\ \quad \quad \quad (\lambda c \in \mathbb{N}.\; \\ \quad \quad \quad \quad u((\lambda k \in \mathbb{N}.\; \operatorname{complement\text{-}in}(f(k), \operatorname{big\text{-}union}(j, \mathrm{OS}(k), f(j))))(c))), \\ \quad \quad \quad \mathrm{OS}(n)) \\ \quad \quad =\; u(\operatorname{big\text{-}union}(j, \mathrm{OS}(n), f(j))) \end{array}}}
\vnbprop{nondecreasing-seg-subset}{\vnbstmt{\begin{array}{@{}l@{}} \forall n \in \mathbb{N}, c.; \forall f \in (\mathbb{N} \to c), j \in \mathrm{OS}(n). \\ \quad \forall a \in \mathbb{N}.\; (f(a) \subseteq f((a + 1))) \Rightarrow (f(j) \subseteq f(n)) \end{array}}}
\vnbprop{measure-continuity-from-below}{\vnbstmt{\begin{array}{@{}l@{}} \forall a, b, u, f. \\ \quad (\operatorname{is\text{-}measure}(a, b, u) \wedge ((f \in (\mathbb{N} \to b)) \wedge \forall n \in \mathbb{N}.\; (f(n) \subseteq f((n + 1))))) \Rightarrow  \\ \quad \quad u(\operatorname{big\text{-}union}(n, \mathbb{N}, f(n))) = \operatorname{esup}(\operatorname{image}((\lambda n \in \mathbb{N}.\; u(f(n))), \mathbb{N})) \end{array}}}
\vnbprop{measure-finite-subset}{\vnbstmt{\begin{array}{@{}l@{}} \forall a, b, u, \alpha, \beta. \\ \quad (\operatorname{is\text{-}measure}(a, b, u) \wedge ((\alpha \in b) \wedge ((\beta \in b) \wedge ((\alpha \subseteq \beta) \wedge (u(\beta) \in \mathbb{R}))))) \Rightarrow  \\ \quad \quad u(\alpha) \in \mathbb{R} \end{array}}}
\vnbprop{measure-subtractive}{\vnbstmt{\begin{array}{@{}l@{}} \forall a, b, u, \alpha, \beta. \\ \quad (\operatorname{is\text{-}measure}(a, b, u) \wedge ((\alpha \in b) \wedge ((\beta \in b) \wedge ((\alpha \subseteq \beta) \wedge (u(\beta) \in \mathbb{R}))))) \Rightarrow  \\ \quad \quad u(\operatorname{complement\text{-}in}(\beta, \alpha)) = (u(\beta) - u(\alpha)) \end{array}}}
\vnbprop{measure-continuity-from-above}{\vnbstmt{\begin{array}{@{}l@{}} \forall a, b, u, f. \\ \quad \operatorname{is\text{-}measure}(a, b, u) \wedge \\ \quad f \in (\mathbb{N} \to b) \wedge \\ \quad \forall n \in \mathbb{N}.\; (f((n + 1)) \subseteq f(n)) \wedge \\ \quad u(f(0)) \in \mathbb{R} \Rightarrow \\ \quad \quad u(\operatorname{complement\text{-}in}(a, \operatorname{big\text{-}union}(n, \mathbb{N}, \operatorname{complement\text{-}in}(a, f(n))))) \\ \quad \quad =\; \operatorname{einf}(\operatorname{image}((\lambda n \in \mathbb{N}.\; u(f(n))), \mathbb{N})) \end{array}}}
\vnbprop{measure-pushforward}{\vnbstmt{\begin{array}{@{}l@{}} \forall a, b, u, a_{2}, c, f. \\ \quad (\operatorname{is\text{-}measure}(a, b, u) \wedge \operatorname{is\text{-}measurable\text{-}map}(a, b, a_{2}, c, f)) \Rightarrow  \\ \quad \quad \operatorname{is\text{-}measure}(a_{2}, c, (\lambda \mathit{pb} \in c.\; u(\{x \in a : (f(x) \in \mathit{pb})\}))) \end{array}}}

\vnbbookfile{theorem-library/integral-laws}{theorem-library/integral-laws}
\noindent{\small integral-laws.scm -- THE INTEGRAL OF A [0,+inf]-VALUED FUNCTION BY SIMPLE FUNCTIONS: the laws, proved from the definition, and the bricks they rest on. Agent M-2 of the October roadmap (week 2, scratchpad/triage/WEEK2.md; decision 4 of docs/roadmap-2026-10-03.md).  Helper prefixes ibk- (the}\par
\vnbprop{etimes-closed}{\vnbstmt{\begin{array}{@{}l@{}} \forall x, y. \\ \quad ((x \in \operatorname{rr\text{-}pos\text{-}star}) \wedge (y \in \operatorname{rr\text{-}pos\text{-}star})) \Rightarrow (\operatorname{etimes}(x, y) \in \operatorname{rr\text{-}pos\text{-}star}) \end{array}}}
\vnbprop{etimes-zero-left}{\vnbstmt{\forall x \in \operatorname{rr\text{-}pos\text{-}star}.\; (\operatorname{etimes}(0, x) = 0)}}
\vnbprop{etimes-zero-right}{\vnbstmt{\forall x \in \operatorname{rr\text{-}pos\text{-}star}.\; (\operatorname{etimes}(x, 0) = 0)}}
\vnbprop{etimes-one-left}{\vnbstmt{\forall x \in \operatorname{rr\text{-}pos\text{-}star}.\; (\operatorname{etimes}(1, x) = x)}}
\vnbprop{etimes-comm}{\vnbstmt{\begin{array}{@{}l@{}} \forall x, y. \\ \quad ((x \in \operatorname{rr\text{-}pos\text{-}star}) \wedge (y \in \operatorname{rr\text{-}pos\text{-}star})) \Rightarrow (\operatorname{etimes}(x, y) = \operatorname{etimes}(y, x)) \end{array}}}
\vnbprop{etimes-distrib-left}{\vnbstmt{\begin{array}{@{}l@{}} \forall c \in \mathbb{R}, x, y. \\ \quad ((0 \leq c) \wedge ((x \in \operatorname{rr\text{-}pos\text{-}star}) \wedge (y \in \operatorname{rr\text{-}pos\text{-}star}))) \Rightarrow  \\ \quad \quad \operatorname{etimes}(c, \operatorname{eplus}(x, y)) = \operatorname{eplus}(\operatorname{etimes}(c, x), \operatorname{etimes}(c, y)) \end{array}}}
\vnbprop{etimes-mono-left}{\vnbstmt{\begin{array}{@{}l@{}} \forall c \in \mathbb{R}, x, y. \\ \quad ((0 \leq c) \wedge (((x \in \operatorname{rr\text{-}pos\text{-}star}) \wedge (y \in \operatorname{rr\text{-}pos\text{-}star})) \wedge (x \leq y))) \Rightarrow  \\ \quad \quad \operatorname{etimes}(c, x) \leq \operatorname{etimes}(c, y) \end{array}}}
\vnbprop{etimes-assoc-real}{\vnbstmt{\begin{array}{@{}l@{}} \forall c, d \in \mathbb{R}, x \in \operatorname{rr\text{-}pos\text{-}star}. \\ \quad ((0 \leq c) \wedge (0 \leq d)) \Rightarrow (\operatorname{etimes}(c, \operatorname{etimes}(d, x)) = \operatorname{etimes}((c \cdot d), x)) \end{array}}}
\vnbprop{etimes-distrib-right-real}{\vnbstmt{\begin{array}{@{}l@{}} \forall c, d \in \mathbb{R}, x \in \operatorname{rr\text{-}pos\text{-}star}. \\ \quad ((0 \leq c) \wedge (0 \leq d)) \Rightarrow  \\ \quad \quad \operatorname{etimes}(\operatorname{eplus}(c, d), x) = \operatorname{eplus}(\operatorname{etimes}(c, x), \operatorname{etimes}(d, x)) \end{array}}}
\vnbprop{etimes-mono-real-left}{\vnbstmt{\begin{array}{@{}l@{}} \forall a, b.; \forall m \in \operatorname{rr\text{-}pos\text{-}star}. \\ \quad ((a \in \mathbb{R}) \wedge ((b \in \mathbb{R}) \wedge ((0 \leq a) \wedge (a \leq b)))) \Rightarrow (\operatorname{etimes}(a, m) \leq \operatorname{etimes}(b, m)) \end{array}}}
\vnbprop{finsum-cm-all-iden-ptwise}{\vnbstmt{\begin{array}{@{}l@{}} \forall m.; \forall s \in \mathbf{Set}, f. \\ \quad (\operatorname{is\text{-}comm\text{-}monoid}(m) \wedge ((|s| \in \mathbb{N}) \wedge \forall z \in s.\; (f(z) = \operatorname{iden}(m)))) \Rightarrow  \\ \quad \quad \operatorname{finsum}(m, f, s) = \operatorname{iden}(m) \end{array}}}
\vnbprop{finsum-cm-add-ptwise}{\vnbstmt{\begin{array}{@{}l@{}} \forall m.; \forall s \in \mathbf{Set}, f, g, h. \\ \quad \operatorname{is\text{-}comm\text{-}monoid}(m) \wedge \\ \quad |s| \in \mathbb{N} \wedge \\ \quad \forall z \in s.\; (f(z) \in \operatorname{carr}(m)) \wedge \\ \quad \forall z \in s.\; (g(z) \in \operatorname{carr}(m)) \wedge \\ \quad \forall z \in s.\; (h(z) = \operatorname{opr}(m)(f(z), g(z))) \Rightarrow \\ \quad \quad \operatorname{finsum}(m, h, s) = \operatorname{opr}(m)(\operatorname{finsum}(m, f, s), \operatorname{finsum}(m, g, s)) \end{array}}}
\vnbprop{finsum-cm-remove-ptwise}{\vnbstmt{\begin{array}{@{}l@{}} \forall m.; \forall t \in \mathbf{Set}, k \in t, g. \\ \quad (\operatorname{is\text{-}comm\text{-}monoid}(m) \wedge ((|t| \in \mathbb{N}) \wedge \forall z \in t.\; (g(z) \in \operatorname{carr}(m)))) \Rightarrow  \\ \quad \quad \operatorname{finsum}(m, g, t) = \operatorname{opr}(m)(\operatorname{finsum}(m, g, \operatorname{difference}(t, \{k\})), g(k)) \end{array}}}
\vnbprop{sep-fiber-insert-other}{\vnbstmt{\begin{array}{@{}l@{}} \forall s, x, i, z. \\ \quad ((x \in \mathbf{Set}) \wedge \neg (i(x) = z)) \Rightarrow  \\ \quad \quad \{w \in (s \cup \{x\}) : (i(w) = z)\} = \{w \in s : (i(w) = z)\} \end{array}}}
\vnbprop{sep-fiber-insert-same}{\vnbstmt{\begin{array}{@{}l@{}} \forall s, x, i. \\ \quad ((x \in \mathbf{Set}) \wedge (i(x) \in \mathbf{Set})) \Rightarrow  \\ \quad \quad \{w \in (s \cup \{x\}) : (i(w) = i(x))\} = (\{w \in s : (i(w) = i(x))\} \cup \{x\}) \end{array}}}
\vnbprop{finsum-cm-fiber-ptwise}{\vnbstmt{\begin{array}{@{}l@{}} \forall m.; \forall s, t \in \mathbf{Set}, i, f, q. \\ \quad \operatorname{is\text{-}comm\text{-}monoid}(m) \wedge \\ \quad |s| \in \mathbb{N} \wedge \\ \quad |t| \in \mathbb{N} \wedge \\ \quad \forall z \in s.\; (i(z) \in t) \wedge \\ \quad \forall z \in s.\; (f(z) \in \operatorname{carr}(m)) \wedge \\ \quad \forall z \in t.\; (q(z) \simeq \operatorname{finsum}(m, f, \{w \in s : (i(w) = z)\})) \Rightarrow \\ \quad \quad \operatorname{finsum}(m, f, s) = \operatorname{finsum}(m, q, t) \end{array}}}
\vnbprop{finsum-cm-pad-ptwise}{\vnbstmt{\begin{array}{@{}l@{}} \forall m.; \forall t \in \mathbf{Set}, r, g. \\ \quad \operatorname{is\text{-}comm\text{-}monoid}(m) \wedge \\ \quad |t| \in \mathbb{N} \wedge \\ \quad \forall z \in r.\; (z \in t) \wedge \\ \quad \forall z \in t.\; (g(z) \in \operatorname{carr}(m)) \wedge \\ \quad \forall z \in t.\; (\neg (z \in r) \Rightarrow (g(z) = \operatorname{iden}(m))) \Rightarrow \\ \quad \quad \operatorname{finsum}(m, g, t) = \operatorname{finsum}(m, g, r) \end{array}}}
\vnbprop{rps-finsum-le-ptwise}{\vnbstmt{\begin{array}{@{}l@{}} \forall s \in \mathbf{Set}, f, g. \\ \quad |s| \in \mathbb{N} \wedge \\ \quad \forall z \in s.\; (f(z) \in \operatorname{rr\text{-}pos\text{-}star}) \wedge \\ \quad \forall z \in s.\; (g(z) \in \operatorname{rr\text{-}pos\text{-}star}) \wedge \\ \quad \forall z \in s.\; (f(z) \leq g(z)) \Rightarrow \\ \quad \quad \operatorname{finsum}(\operatorname{rr\text{-}pos\text{-}star\text{-}add\text{-}monoid}, f, s) \leq \operatorname{finsum}(\operatorname{rr\text{-}pos\text{-}star\text{-}add\text{-}monoid}, g, s) \end{array}}}
\vnbprop{rps-finsum-scale-ptwise}{\vnbstmt{\begin{array}{@{}l@{}} \forall c \in \mathbb{R}, s \in \mathbf{Set}, f, h. \\ \quad 0 \leq c \wedge \\ \quad |s| \in \mathbb{N} \wedge \\ \quad \forall z \in s.\; (f(z) \in \operatorname{rr\text{-}pos\text{-}star}) \wedge \\ \quad \forall z \in s.\; (h(z) = \operatorname{etimes}(c, f(z))) \Rightarrow \\ \quad \quad \operatorname{finsum}(\operatorname{rr\text{-}pos\text{-}star\text{-}add\text{-}monoid}, h, s) = \operatorname{etimes}(c, \operatorname{finsum}(\operatorname{rr\text{-}pos\text{-}star\text{-}add\text{-}monoid}, f, s)) \end{array}}}
\vnbprop{sigma-algebra-preimage-finite-in}{\vnbstmt{\begin{array}{@{}l@{}} \forall a, b.; \forall s \in \mathbf{Set}, k. \\ \quad (\operatorname{is\text{-}sigma\text{-}algebra}(a, b) \wedge ((|s| \in \mathbb{N}) \wedge \forall z \in s.\; (\{x \in a : (k(x) = z)\} \in b))) \Rightarrow  \\ \quad \quad \{x \in a : (k(x) \in s)\} \in b \end{array}}}
\vnbprop{measure-preimage-finite-sum}{\vnbstmt{\begin{array}{@{}l@{}} \forall a, b, u.; \forall s \in \mathbf{Set}, k, g. \\ \quad \operatorname{is\text{-}measure}(a, b, u) \wedge \\ \quad |s| \in \mathbb{N} \wedge \\ \quad \forall z \in s.\; (\{x \in a : (k(x) = z)\} \in b) \wedge \\ \quad \forall z \in s.\; (g(z) = u(\{x \in a : (k(x) = z)\})) \Rightarrow \\ \quad \quad u(\{x \in a : (k(x) \in s)\}) = \operatorname{finsum}(\operatorname{rr\text{-}pos\text{-}star\text{-}add\text{-}monoid}, g, s) \end{array}}}
\vnbprop{measurable-fn-ge}{\vnbstmt{\begin{array}{@{}l@{}} \forall b, c, f.; \forall a \in \mathbb{R}. \\ \quad \operatorname{is\text{-}measurable\text{-}fn}(b, c, f) \Rightarrow (\{x \in b : (a \leq f(x))\} \in c) \end{array}}}
\vnbprop{measurable-fn-level-set}{\vnbstmt{\begin{array}{@{}l@{}} \forall b, c, f.; \forall a \in \mathbb{R}. \\ \quad \operatorname{is\text{-}measurable\text{-}fn}(b, c, f) \Rightarrow (\{x \in b : (f(x) = a)\} \in c) \end{array}}}
\vnbprop{simple-integral-unfold}{\vnbstmt{\begin{array}{@{}l@{}} \forall a, u, s. \\ \quad (\operatorname{simple\text{-}integral}(a, u, s) \simeq \operatorname{finsum}(\operatorname{rr\text{-}pos\text{-}star\text{-}add\text{-}monoid}, (\lambda y \in \operatorname{image}(s, a).\; \operatorname{etimes}(y, u(\{x \in a : (s(x) = y)\}))), \operatorname{image}(s, a))) \end{array}}}
\vnbprop{integral-unfold}{\vnbstmt{\begin{array}{@{}l@{}} \forall a, b, u, f. \\ \quad (\operatorname{integral}(a, b, u, f) \simeq \operatorname{esup}(\{y \in \operatorname{rr\text{-}pos\text{-}star} : \exists s.\; (\operatorname{is\text{-}simple\text{-}fn}(a, b, s) \wedge (\operatorname{ptwise\text{-}le}(a, s, f) \wedge (y = \operatorname{simple\text{-}integral}(a, u, s))))\})) \end{array}}}
\vnbprop{integral-in}{\vnbstmt{\begin{array}{@{}l@{}} \forall a, b, u, f. \\ \quad (\operatorname{is\text{-}measure}(a, b, u) \wedge \operatorname{is\text{-}measurable\text{-}fn}(a, b, f)) \Rightarrow  \\ \quad \quad \operatorname{integral}(a, b, u, f) \in \operatorname{rr\text{-}pos\text{-}star} \end{array}}}
\vnbprop{integral-monotone}{\vnbstmt{\begin{array}{@{}l@{}} \forall a, b, u, f, g. \\ \quad \operatorname{is\text{-}measure}(a, b, u) \wedge \\ \quad \operatorname{is\text{-}measurable\text{-}fn}(a, b, f) \wedge \\ \quad \operatorname{is\text{-}measurable\text{-}fn}(a, b, g) \wedge \\ \quad \operatorname{ptwise\text{-}le}(a, f, g) \Rightarrow \\ \quad \quad \operatorname{integral}(a, b, u, f) \leq \operatorname{integral}(a, b, u, g) \end{array}}}
\vnbprop{simple-integral-partition}{\vnbstmt{\begin{array}{@{}l@{}} \forall a, b, u, s.; \forall c \in \mathbf{Set}, k, h, g. \\ \quad \operatorname{is\text{-}measure}(a, b, u) \wedge \\ \quad |c| \in \mathbb{N} \wedge \\ \quad \forall z \in a.\; (k(z) \in c) \wedge \\ \quad \forall z \in c.\; (\{x \in a : (k(x) = z)\} \in b) \wedge \\ \quad \forall z \in c.\; (h(z) \in \mathbb{R}) \wedge \\ \quad \forall z \in c.\; (0 \leq h(z)) \wedge \\ \quad \forall z \in a.\; (s(z) = h(k(z))) \wedge \\ \quad \forall z \in c.\; (g(z) = \operatorname{etimes}(h(z), u(\{x \in a : (k(x) = z)\}))) \Rightarrow \\ \quad \quad \operatorname{simple\text{-}integral}(a, u, s) = \operatorname{finsum}(\operatorname{rr\text{-}pos\text{-}star\text{-}add\text{-}monoid}, g, c) \end{array}}}
\vnbprop{simple-integral-in}{\vnbstmt{\begin{array}{@{}l@{}} \forall a, b, u, s. \\ \quad (\operatorname{is\text{-}measure}(a, b, u) \wedge \operatorname{is\text{-}simple\text{-}fn}(a, b, s)) \Rightarrow  \\ \quad \quad \operatorname{simple\text{-}integral}(a, u, s) \in \operatorname{rr\text{-}pos\text{-}star} \end{array}}}
\vnbprop{simple-integral-mono}{\vnbstmt{\begin{array}{@{}l@{}} \forall a, b, u, s, t. \\ \quad (\operatorname{is\text{-}measure}(a, b, u) \wedge (\operatorname{is\text{-}simple\text{-}fn}(a, b, s) \wedge (\operatorname{is\text{-}simple\text{-}fn}(a, b, t) \wedge \operatorname{ptwise\text{-}le}(a, s, t)))) \Rightarrow  \\ \quad \quad \operatorname{simple\text{-}integral}(a, u, s) \leq \operatorname{simple\text{-}integral}(a, u, t) \end{array}}}
\vnbprop{integral-simple}{\vnbstmt{\begin{array}{@{}l@{}} \forall a, b, u, s. \\ \quad (\operatorname{is\text{-}measure}(a, b, u) \wedge \operatorname{is\text{-}simple\text{-}fn}(a, b, s)) \Rightarrow  \\ \quad \quad \operatorname{integral}(a, b, u, s) = \operatorname{simple\text{-}integral}(a, u, s) \end{array}}}
\vnbprop{integral-sup-of-simple}{\vnbstmt{\begin{array}{@{}l@{}} \forall a, b, u, f. \\ \quad (\operatorname{is\text{-}measure}(a, b, u) \wedge \operatorname{is\text{-}measurable\text{-}fn}(a, b, f)) \Rightarrow  \\ \quad \quad \operatorname{integral}(a, b, u, f) \\ \quad \quad =\; \operatorname{esup}(\{y \in \operatorname{rr\text{-}pos\text{-}star} : \exists s.\; (\operatorname{is\text{-}simple\text{-}fn}(a, b, s) \wedge (\operatorname{ptwise\text{-}le}(a, s, f) \wedge (y = \operatorname{integral}(a, b, u, s))))\}) \end{array}}}
\vnbprop{indicator-value-in}{\vnbstmt{\forall a, \alpha.; \forall x \in a.\; ((x \in \alpha) \Rightarrow (\operatorname{indicator}(a, \alpha)(x) = 1))}}
\vnbprop{indicator-value-out}{\vnbstmt{\forall a, \alpha.; \forall x \in a.\; (\neg (x \in \alpha) \Rightarrow (\operatorname{indicator}(a, \alpha)(x) = 0))}}
\vnbprop{indicator-simple}{\vnbstmt{\begin{array}{@{}l@{}} \forall a, b, \alpha. \\ \quad (\operatorname{is\text{-}sigma\text{-}algebra}(a, b) \wedge (\alpha \in b)) \Rightarrow  \\ \quad \quad \operatorname{is\text{-}simple\text{-}fn}(a, b, \operatorname{indicator}(a, \alpha)) \end{array}}}
\vnbprop{simple-integral-indicator}{\vnbstmt{\begin{array}{@{}l@{}} \forall a, b, u, \alpha. \\ \quad (\operatorname{is\text{-}measure}(a, b, u) \wedge (\alpha \in b)) \Rightarrow  \\ \quad \quad \operatorname{simple\text{-}integral}(a, u, \operatorname{indicator}(a, \alpha)) = u(\alpha) \end{array}}}
\vnbprop{integral-indicator}{\vnbstmt{\begin{array}{@{}l@{}} \forall a, b, u, \alpha. \\ \quad (\operatorname{is\text{-}measure}(a, b, u) \wedge (\alpha \in b)) \Rightarrow  \\ \quad \quad \operatorname{integral}(a, b, u, \operatorname{indicator}(a, \alpha)) = u(\alpha) \end{array}}}
\vnbprop{integral-infinite-on-null}{\vnbstmt{\begin{array}{@{}l@{}} \forall a, b, u, \alpha. \\ \quad (\operatorname{is\text{-}measure}(a, b, u) \wedge ((\alpha \in b) \wedge (u(\alpha) = 0))) \Rightarrow  \\ \quad \quad \operatorname{integral}(a, b, u, \operatorname{ptwise\text{-}scale}(a, \operatorname{pos\text{-}inf}, \operatorname{indicator}(a, \alpha))) = 0 \end{array}}}
\vnbprop{simple-integral-additive}{\vnbstmt{\begin{array}{@{}l@{}} \forall a, b, u, s, t. \\ \quad (\operatorname{is\text{-}measure}(a, b, u) \wedge (\operatorname{is\text{-}simple\text{-}fn}(a, b, s) \wedge \operatorname{is\text{-}simple\text{-}fn}(a, b, t))) \Rightarrow  \\ \quad \quad \operatorname{simple\text{-}integral}(a, u, \operatorname{ptwise\text{-}eplus}(a, s, t)) \\ \quad \quad =\; \operatorname{eplus}(\operatorname{simple\text{-}integral}(a, u, s), \operatorname{simple\text{-}integral}(a, u, t)) \end{array}}}
\vnbprop{simple-integral-homogeneous}{\vnbstmt{\begin{array}{@{}l@{}} \forall a, b, u, s, c. \\ \quad (\operatorname{is\text{-}measure}(a, b, u) \wedge (\operatorname{is\text{-}simple\text{-}fn}(a, b, s) \wedge ((c \in \mathbb{R}) \wedge (0 \leq c)))) \Rightarrow  \\ \quad \quad \operatorname{simple\text{-}integral}(a, u, \operatorname{ptwise\text{-}scale}(a, c, s)) = \operatorname{etimes}(c, \operatorname{simple\text{-}integral}(a, u, s)) \end{array}}}
\vnbprop{measurable-fn-finite-range}{\vnbstmt{\begin{array}{@{}l@{}} \forall a, b, u. \\ \quad \operatorname{is\text{-}sigma\text{-}algebra}(a, b) \wedge \\ \quad u \in (a \to \operatorname{rr\text{-}pos\text{-}star}) \wedge \\ \quad |\operatorname{image}(u, a)| \in \mathbb{N} \wedge \\ \quad \forall z \in \operatorname{image}(u, a).\; (\{x \in a : (u(x) = z)\} \in b) \Rightarrow \\ \quad \quad \operatorname{is\text{-}measurable\text{-}fn}(a, b, u) \end{array}}}
\vnbprop{simple-fn-scale}{\vnbstmt{\begin{array}{@{}l@{}} \forall a, b, s, c. \\ \quad (\operatorname{is\text{-}sigma\text{-}algebra}(a, b) \wedge (\operatorname{is\text{-}simple\text{-}fn}(a, b, s) \wedge ((c \in \mathbb{R}) \wedge (0 \leq c)))) \Rightarrow  \\ \quad \quad \operatorname{is\text{-}simple\text{-}fn}(a, b, \operatorname{ptwise\text{-}scale}(a, c, s)) \end{array}}}
\vnbprop{integral-homogeneous}{\vnbstmt{\begin{array}{@{}l@{}} \forall a, b, u, f, c. \\ \quad (\operatorname{is\text{-}measure}(a, b, u) \wedge (\operatorname{is\text{-}measurable\text{-}fn}(a, b, f) \wedge ((c \in \mathbb{R}) \wedge (0 \leq c)))) \Rightarrow  \\ \quad \quad \operatorname{integral}(a, b, u, \operatorname{ptwise\text{-}scale}(a, c, f)) = \operatorname{etimes}(c, \operatorname{integral}(a, b, u, f)) \end{array}}}

\vnbbookfile{theorem-library/esum-scale}{theorem-library/esum-scale}
\noindent{\small ==================================================================== theorem-library/esum-scale.scm -- MEASURE-5-A, 2026-10-10. Homogeneity of ESUM, and a scaled measure is a measure.}\par
\vnbprop{esum-scale-ptwise}{\vnbstmt{\begin{array}{@{}l@{}} \forall c \in \mathbb{R}, d, g, h. \\ \quad 0 \leq c \wedge \\ \quad g \in (d \to \operatorname{rr\text{-}pos\text{-}star}) \wedge \\ \quad h \in (d \to \operatorname{rr\text{-}pos\text{-}star}) \wedge \\ \quad \forall z \in d.\; (h(z) = \operatorname{etimes}(c, g(z))) \Rightarrow \\ \quad \quad \operatorname{esum}(h) = \operatorname{etimes}(c, \operatorname{esum}(g)) \end{array}}}
\vnbprop{esum-scale}{\vnbstmt{\begin{array}{@{}l@{}} \forall c \in \mathbb{R}, d.; \forall f \in (d \to \operatorname{rr\text{-}pos\text{-}star}). \\ \quad (0 \leq c) \Rightarrow (\operatorname{esum}((\lambda z \in d.\; \operatorname{etimes}(c, f(z)))) = \operatorname{etimes}(c, \operatorname{esum}(f))) \end{array}}}
\vnbprop{measure-scale}{\vnbstmt{\begin{array}{@{}l@{}} \forall a, b, u.; \forall c \in \mathbb{R}. \\ \quad (\operatorname{is\text{-}measure}(a, b, u) \wedge (0 \leq c)) \Rightarrow  \\ \quad \quad \operatorname{is\text{-}measure}(a, b, (\lambda x \in b.\; \operatorname{etimes}(c, u(x)))) \end{array}}}

\vnbbookfile{theorem-library/integral-convergence}{theorem-library/integral-convergence}
\noindent{\small integral-convergence.scm -- MEASURABLE FUNCTIONS CLOSED UNDER +, *, sup and lim inf; APPROXIMATION BY SIMPLE FUNCTIONS; MONOTONE CONVERGENCE; ADDITIVITY OF THE INTEGRAL; FATOU'S LEMMA; DOMINATED CONVERGENCE.  Agent M-3 of the October roadmap (week 3, scratchpad/triage/WEEK3.md).  Helper prefix icv-.  Every theorem ...}\par
\vnbprop{rps-le-total}{\vnbstmt{\forall x, y \in \operatorname{rr\text{-}pos\text{-}star}.\; ((x \leq y) \vee (y \leq x))}}
\vnbprop{rps-not-lt-le}{\vnbstmt{\forall x, y \in \operatorname{rr\text{-}pos\text{-}star}.\; (\neg (x < y) \Rightarrow (y \leq x))}}
\vnbprop{rps-lt-le-trans}{\vnbstmt{\begin{array}{@{}l@{}} \forall x, y, z \in \operatorname{rr\text{-}pos\text{-}star}. \\ \quad ((x < y) \wedge (y \leq z)) \Rightarrow (x < z) \end{array}}}
\vnbprop{esup-lt-witness}{\vnbstmt{\begin{array}{@{}l@{}} \forall s.; \forall a \in \operatorname{rr\text{-}pos\text{-}star}. \\ \quad ((s \subseteq \operatorname{rr\text{-}pos\text{-}star}) \wedge (a < \operatorname{esup}(s))) \Rightarrow \exists y \in s.\; (a < y) \end{array}}}
\vnbprop{esup-gt-of-member}{\vnbstmt{\begin{array}{@{}l@{}} \forall s.; \forall y \in s, a \in \operatorname{rr\text{-}pos\text{-}star}. \\ \quad ((s \subseteq \operatorname{rr\text{-}pos\text{-}star}) \wedge (a < y)) \Rightarrow (a < \operatorname{esup}(s)) \end{array}}}
\vnbprop{measurable-fn-of-nonneg-tails}{\vnbstmt{\begin{array}{@{}l@{}} \forall a, b, f. \\ \quad \operatorname{is\text{-}sigma\text{-}algebra}(a, b) \wedge \\ \quad f \in (a \to \operatorname{rr\text{-}pos\text{-}star}) \wedge \\ \quad \forall l \in \mathbb{R}.\; ((0 \leq l) \Rightarrow (\{x \in a : (l < f(x))\} \in b)) \Rightarrow \\ \quad \quad \operatorname{is\text{-}measurable\text{-}fn}(a, b, f) \end{array}}}
\vnbprop{measurable-fn-sup-of-values}{\vnbstmt{\begin{array}{@{}l@{}} \forall a, b, f, g. \\ \quad \operatorname{is\text{-}sigma\text{-}algebra}(a, b) \wedge \\ \quad f \in (\mathbb{N} \to (a \to \operatorname{rr\text{-}pos\text{-}star})) \wedge \\ \quad \forall n \in \mathbb{N}.\; \operatorname{is\text{-}measurable\text{-}fn}(a, b, f(n)) \wedge \\ \quad g \in (a \to \operatorname{rr\text{-}pos\text{-}star}) \wedge \\ \quad \forall x \in a.\; (g(x) = \operatorname{esup}(\operatorname{image}((\lambda n \in \mathbb{N}.\; f(n)(x)), \mathbb{N}))) \Rightarrow \\ \quad \quad \operatorname{is\text{-}measurable\text{-}fn}(a, b, g) \end{array}}}
\vnbprop{measurable-fn-sup}{\vnbstmt{\begin{array}{@{}l@{}} \forall a, b, f. \\ \quad \operatorname{is\text{-}sigma\text{-}algebra}(a, b) \wedge \\ \quad f \in (\mathbb{N} \to (a \to \operatorname{rr\text{-}pos\text{-}star})) \wedge \\ \quad \forall n \in \mathbb{N}.\; \operatorname{is\text{-}measurable\text{-}fn}(a, b, f(n)) \Rightarrow \\ \quad \quad \operatorname{is\text{-}measurable\text{-}fn}(a, b, \operatorname{ptwise\text{-}sup}(a, f)) \end{array}}}
\vnbprop{measurable-fn-guarded-lt}{\vnbstmt{\begin{array}{@{}l@{}} \forall a, c, f, p, s.; \forall b \in \mathbb{R}. \\ \quad \operatorname{is\text{-}measurable\text{-}fn}(a, c, f) \Rightarrow (\{x \in a : ((p \in s) \wedge (f(x) < b))\} \in c) \end{array}}}
\vnbprop{measurable-fn-lower-bound-set}{\vnbstmt{\begin{array}{@{}l@{}} \forall a, c, f, s, b. \\ \quad \operatorname{is\text{-}sigma\text{-}algebra}(a, c) \wedge \\ \quad f \in (\mathbb{N} \to (a \to \operatorname{rr\text{-}pos\text{-}star})) \wedge \\ \quad \forall n \in \mathbb{N}.\; \operatorname{is\text{-}measurable\text{-}fn}(a, c, f(n)) \wedge \\ \quad s \subseteq \mathbb{N} \wedge \\ \quad b \in \mathbb{R} \Rightarrow \\ \quad \quad \{x \in a : \forall p \in s.\; (b \leq f(p)(x))\} \in c \end{array}}}
\vnbprop{measurable-fn-inf-of-values}{\vnbstmt{\begin{array}{@{}l@{}} \forall a, b, f, s, g. \\ \quad \operatorname{is\text{-}sigma\text{-}algebra}(a, b) \wedge \\ \quad f \in (\mathbb{N} \to (a \to \operatorname{rr\text{-}pos\text{-}star})) \wedge \\ \quad \forall n \in \mathbb{N}.\; \operatorname{is\text{-}measurable\text{-}fn}(a, b, f(n)) \wedge \\ \quad s \subseteq \mathbb{N} \wedge \\ \quad g \in (a \to \operatorname{rr\text{-}pos\text{-}star}) \wedge \\ \quad \forall x \in a.\; (g(x) = \operatorname{einf}(\operatorname{image}((\lambda n \in \mathbb{N}.\; f(n)(x)), s))) \Rightarrow \\ \quad \quad \operatorname{is\text{-}measurable\text{-}fn}(a, b, g) \end{array}}}
\vnbprop{measurable-fn-liminf}{\vnbstmt{\begin{array}{@{}l@{}} \forall a, b, f. \\ \quad \operatorname{is\text{-}sigma\text{-}algebra}(a, b) \wedge \\ \quad f \in (\mathbb{N} \to (a \to \operatorname{rr\text{-}pos\text{-}star})) \wedge \\ \quad \forall n \in \mathbb{N}.\; \operatorname{is\text{-}measurable\text{-}fn}(a, b, f(n)) \Rightarrow \\ \quad \quad \operatorname{is\text{-}measurable\text{-}fn}(a, b, \operatorname{ptwise\text{-}liminf}(a, f)) \end{array}}}
\vnbprop{rr-nn-floor-mul}{\vnbstmt{\begin{array}{@{}l@{}} \forall c \in \mathbb{R}, r. \\ \quad ((0 \leq c) \wedge \operatorname{pos\text{-}rr}(r)) \Rightarrow \exists f \in \mathbb{N}.\; (((f \cdot r) \leq c) \wedge (c < ((f \cdot r) + r))) \end{array}}}
\vnbprop{rr-grid-between}{\vnbstmt{\begin{array}{@{}l@{}} \forall a, b \in \mathbb{R}. \\ \quad ((0 \leq a) \wedge (a < b)) \Rightarrow  \\ \quad \quad \exists n, k \in \mathbb{N}.\; ((a < (k \cdot \operatorname{recip}((n + 1)))) \wedge ((k \cdot \operatorname{recip}((n + 1))) < b)) \end{array}}}
\vnbprop{measurable-fn-grid-row}{\vnbstmt{\begin{array}{@{}l@{}} \forall a, b, f, g, l, r. \\ \quad (\operatorname{is\text{-}measurable\text{-}fn}(a, b, f) \wedge (\operatorname{is\text{-}measurable\text{-}fn}(a, b, g) \wedge ((l \in \mathbb{R}) \wedge (r \in \mathbb{R})))) \Rightarrow  \\ \quad \quad \{x \in a : \exists k \in \mathbb{N}.\; (((k \cdot r) < f(x)) \wedge ((l - (k \cdot r)) < g(x)))\} \in b \end{array}}}
\vnbprop{measurable-fn-eplus-of-values}{\vnbstmt{\begin{array}{@{}l@{}} \forall a, b, f, g, h. \\ \quad \operatorname{is\text{-}measurable\text{-}fn}(a, b, f) \wedge \\ \quad \operatorname{is\text{-}measurable\text{-}fn}(a, b, g) \wedge \\ \quad h \in (a \to \operatorname{rr\text{-}pos\text{-}star}) \wedge \\ \quad \forall x \in a.\; (h(x) = \operatorname{eplus}(f(x), g(x))) \Rightarrow \\ \quad \quad \operatorname{is\text{-}measurable\text{-}fn}(a, b, h) \end{array}}}
\vnbprop{measurable-fn-eplus}{\vnbstmt{\begin{array}{@{}l@{}} \forall a, b, f, g. \\ \quad (\operatorname{is\text{-}measurable\text{-}fn}(a, b, f) \wedge \operatorname{is\text{-}measurable\text{-}fn}(a, b, g)) \Rightarrow  \\ \quad \quad \operatorname{is\text{-}measurable\text{-}fn}(a, b, \operatorname{ptwise\text{-}eplus}(a, f, g)) \end{array}}}
\vnbprop{rr-lt-div-of-lt-mul}{\vnbstmt{\begin{array}{@{}l@{}} \forall l, f, g \in \mathbb{R}. \\ \quad ((0 < g) \wedge (l < (f \cdot g))) \Rightarrow ((l \cdot \operatorname{recip}(g)) < f) \end{array}}}
\vnbprop{rr-recip-swap-lt}{\vnbstmt{\begin{array}{@{}l@{}} \forall l, q, g \in \mathbb{R}. \\ \quad ((0 < g) \wedge ((0 < q) \wedge ((l \cdot \operatorname{recip}(g)) < q))) \Rightarrow ((l \cdot \operatorname{recip}(q)) < g) \end{array}}}
\vnbprop{rr-lt-mul-of-grid}{\vnbstmt{\begin{array}{@{}l@{}} \forall l, q, f, g \in \mathbb{R}. \\ \quad ((0 \leq l) \wedge ((0 < q) \wedge ((q < f) \wedge ((l \cdot \operatorname{recip}(q)) < g)))) \Rightarrow (l < (f \cdot g)) \end{array}}}
\vnbprop{measurable-fn-etimes-cell}{\vnbstmt{\begin{array}{@{}l@{}} \forall a, b, f, g, l, q. \\ \quad (\operatorname{is\text{-}measurable\text{-}fn}(a, b, f) \wedge (\operatorname{is\text{-}measurable\text{-}fn}(a, b, g) \wedge ((l \in \mathbb{R}) \wedge (q \in \mathbb{R})))) \Rightarrow  \\ \quad \quad \{x \in a : ((0 < q) \wedge ((q < f(x)) \wedge ((l \cdot \operatorname{recip}(q)) < g(x))))\} \in b \end{array}}}
\vnbprop{measurable-fn-etimes-row}{\vnbstmt{\begin{array}{@{}l@{}} \forall a, b, f, g, l, r. \\ \quad (\operatorname{is\text{-}measurable\text{-}fn}(a, b, f) \wedge (\operatorname{is\text{-}measurable\text{-}fn}(a, b, g) \wedge ((l \in \mathbb{R}) \wedge (r \in \mathbb{R})))) \Rightarrow  \\ \quad \quad \{x \in a : \exists k \in \mathbb{N}.\; ((0 < (k \cdot r)) \wedge (((k \cdot r) < f(x)) \wedge ((l \cdot \operatorname{recip}((k \cdot r))) < g(x))))\} \in b \end{array}}}
\vnbprop{measurable-fn-etimes-of-values}{\vnbstmt{\begin{array}{@{}l@{}} \forall a, b, f, g, h. \\ \quad \operatorname{is\text{-}measurable\text{-}fn}(a, b, f) \wedge \\ \quad \operatorname{is\text{-}measurable\text{-}fn}(a, b, g) \wedge \\ \quad h \in (a \to \operatorname{rr\text{-}pos\text{-}star}) \wedge \\ \quad \forall x \in a.\; (h(x) = \operatorname{etimes}(f(x), g(x))) \Rightarrow \\ \quad \quad \operatorname{is\text{-}measurable\text{-}fn}(a, b, h) \end{array}}}
\vnbprop{measurable-fn-etimes}{\vnbstmt{\begin{array}{@{}l@{}} \forall a, b, f, g. \\ \quad (\operatorname{is\text{-}measurable\text{-}fn}(a, b, f) \wedge \operatorname{is\text{-}measurable\text{-}fn}(a, b, g)) \Rightarrow  \\ \quad \quad \operatorname{is\text{-}measurable\text{-}fn}(a, b, \operatorname{ptwise\text{-}etimes}(a, f, g)) \end{array}}}
\vnbprop{esup-eplus-le}{\vnbstmt{\begin{array}{@{}l@{}} \forall u, v.; \forall l \in \operatorname{rr\text{-}pos\text{-}star}. \\ \quad u \in (\mathbb{N} \to \operatorname{rr\text{-}pos\text{-}star}) \wedge \\ \quad v \in (\mathbb{N} \to \operatorname{rr\text{-}pos\text{-}star}) \wedge \\ \quad \forall n \in \mathbb{N}.\; (u(n) \leq u((n + 1))) \wedge \\ \quad \forall n \in \mathbb{N}.\; (v(n) \leq v((n + 1))) \wedge \\ \quad \forall n \in \mathbb{N}.\; (\operatorname{eplus}(u(n), v(n)) \leq l) \Rightarrow \\ \quad \quad \operatorname{eplus}(\operatorname{esup}(\operatorname{image}(u, \mathbb{N})), \operatorname{esup}(\operatorname{image}(v, \mathbb{N}))) \leq l \end{array}}}
\vnbprop{rps-finsum-esup-le}{\vnbstmt{\begin{array}{@{}l@{}} \forall s \in \mathbf{Set}, h, a.; \forall l \in \operatorname{rr\text{-}pos\text{-}star}. \\ \quad |s| \in \mathbb{N} \wedge \\ \quad h \in (\mathbb{N} \to (s \to \operatorname{rr\text{-}pos\text{-}star})) \wedge \\ \quad \forall z \in s, n \in \mathbb{N}.\; (h(n)(z) \leq h((n + 1))(z)) \wedge \\ \quad \forall z \in s.\; (a(z) = \operatorname{esup}(\operatorname{image}((\lambda m \in \mathbb{N}.\; h(m)(z)), \mathbb{N}))) \wedge \\ \quad \forall n \in \mathbb{N}.\; (\operatorname{finsum}(\operatorname{rr\text{-}pos\text{-}star\text{-}add\text{-}monoid}, h(n), s) \leq l) \Rightarrow \\ \quad \quad \operatorname{finsum}(\operatorname{rr\text{-}pos\text{-}star\text{-}add\text{-}monoid}, a, s) \leq l \end{array}}}
\vnbprop{simple-fn-of-partition}{\vnbstmt{\begin{array}{@{}l@{}} \forall a, b, s, k, h, t. \\ \quad \operatorname{is\text{-}sigma\text{-}algebra}(a, b) \wedge \\ \quad s \in \mathbf{Set} \wedge \\ \quad |s| \in \mathbb{N} \wedge \\ \quad \forall z \in a.\; (k(z) \in s) \wedge \\ \quad \forall z \in s.\; (\{x \in a : (k(x) = z)\} \in b) \wedge \\ \quad \forall z \in s.\; (h(z) \in \mathbb{R}) \wedge \\ \quad \forall z \in s.\; (0 \leq h(z)) \wedge \\ \quad t \in (a \to \operatorname{rr\text{-}pos\text{-}star}) \wedge \\ \quad \forall z \in a.\; (t(z) = h(k(z))) \Rightarrow \\ \quad \quad \operatorname{is\text{-}simple\text{-}fn}(a, b, t) \end{array}}}
\vnbprop{rr-lt-mul-of-div-lt}{\vnbstmt{\begin{array}{@{}l@{}} \forall l, q, g \in \mathbb{R}. \\ \quad ((0 < g) \wedge ((l \cdot \operatorname{recip}(g)) < q)) \Rightarrow (l < (q \cdot g)) \end{array}}}
\vnbprop{rps-le-of-scaled}{\vnbstmt{\begin{array}{@{}l@{}} \forall u, l \in \operatorname{rr\text{-}pos\text{-}star}. \\ \quad \forall c \in \mathbb{R}.\; (((0 < c) \wedge (c < 1)) \Rightarrow (\operatorname{etimes}(c, u) \leq l)) \Rightarrow (u \leq l) \end{array}}}
\vnbprop{esup-etimes-ge}{\vnbstmt{\begin{array}{@{}l@{}} \forall w \in \mathbb{R}, m \in (\mathbb{N} \to \operatorname{rr\text{-}pos\text{-}star}). \\ \quad (0 \leq w) \Rightarrow  \\ \quad \quad \operatorname{etimes}(w, \operatorname{esup}(\operatorname{image}(m, \mathbb{N}))) \leq \operatorname{esup}(\operatorname{image}((\lambda n \in \mathbb{N}.\; \operatorname{etimes}(w, m(n))), \mathbb{N})) \end{array}}}
\vnbprop{simple-restrict-integral}{\vnbstmt{\begin{array}{@{}l@{}} \forall a, b, u, s, c, t, f. \\ \quad \operatorname{is\text{-}measure}(a, b, u) \wedge \\ \quad \operatorname{is\text{-}simple\text{-}fn}(a, b, s) \wedge \\ \quad c \in \mathbb{R} \wedge \\ \quad 0 \leq c \wedge \\ \quad \forall y \in \operatorname{image}(s, a).\; ((\{\mathit{six} \in a : (s(\mathit{six}) = y)\} \cap t) \in b) \wedge \\ \quad f \in (a \to \operatorname{rr\text{-}pos\text{-}star}) \wedge \\ \quad \forall x \in a.\; (f(x) = \left(\text{if }(x \in t)\text{ then }(c \cdot s(x))\text{ else }0\right)) \Rightarrow \\ \quad \quad \operatorname{is\text{-}simple\text{-}fn}(a, b, f) \wedge \\ \quad \quad \operatorname{simple\text{-}integral}(a, u, f) \\ \quad \quad =\; \operatorname{finsum}(\operatorname{rr\text{-}pos\text{-}star\text{-}add\text{-}monoid}, \\ \quad \quad \quad (\lambda y \in \operatorname{image}(s, a).\; \operatorname{etimes}((c \cdot y), u((\{\mathit{six} \in a : (s(\mathit{six}) = y)\} \cap t)))), \\ \quad \quad \quad \operatorname{image}(s, a)) \end{array}}}
\vnbprop{mct-simple-below}{\vnbstmt{\begin{array}{@{}l@{}} \forall a, b, u, f, s. \\ \quad \operatorname{is\text{-}measure}(a, b, u) \wedge \\ \quad f \in (\mathbb{N} \to (a \to \operatorname{rr\text{-}pos\text{-}star})) \wedge \\ \quad \forall n \in \mathbb{N}.\; \operatorname{is\text{-}measurable\text{-}fn}(a, b, f(n)) \wedge \\ \quad \forall n \in \mathbb{N}.\; \operatorname{ptwise\text{-}le}(a, f(n), f((n + 1))) \wedge \\ \quad \operatorname{is\text{-}simple\text{-}fn}(a, b, s) \wedge \\ \quad \operatorname{ptwise\text{-}le}(a, s, \operatorname{ptwise\text{-}sup}(a, f)) \Rightarrow \\ \quad \quad \operatorname{simple\text{-}integral}(a, u, s) \leq \operatorname{esup}(\operatorname{image}((\lambda n \in \mathbb{N}.\; \operatorname{integral}(a, b, u, f(n))), \mathbb{N})) \end{array}}}
\vnbprop{monotone-convergence}{\vnbstmt{\begin{array}{@{}l@{}} \forall a, b, u, f. \\ \quad \operatorname{is\text{-}measure}(a, b, u) \wedge \\ \quad f \in (\mathbb{N} \to (a \to \operatorname{rr\text{-}pos\text{-}star})) \wedge \\ \quad \forall n \in \mathbb{N}.\; \operatorname{is\text{-}measurable\text{-}fn}(a, b, f(n)) \wedge \\ \quad \forall n \in \mathbb{N}.\; \operatorname{ptwise\text{-}le}(a, f(n), f((n + 1))) \Rightarrow \\ \quad \quad \operatorname{is\text{-}measurable\text{-}fn}(a, b, \operatorname{ptwise\text{-}sup}(a, f)) \wedge \\ \quad \quad \operatorname{integral}(a, b, u, \operatorname{ptwise\text{-}sup}(a, f)) = \operatorname{esup}(\operatorname{image}((\lambda n \in \mathbb{N}.\; \operatorname{integral}(a, b, u, f(n))), \mathbb{N})) \end{array}}}
\vnbprop{fatou}{\vnbstmt{\begin{array}{@{}l@{}} \forall a, b, u, f. \\ \quad \operatorname{is\text{-}measure}(a, b, u) \wedge \\ \quad f \in (\mathbb{N} \to (a \to \operatorname{rr\text{-}pos\text{-}star})) \wedge \\ \quad \forall n \in \mathbb{N}.\; \operatorname{is\text{-}measurable\text{-}fn}(a, b, f(n)) \Rightarrow \\ \quad \quad \operatorname{is\text{-}measurable\text{-}fn}(a, b, \operatorname{ptwise\text{-}liminf}(a, f)) \wedge \\ \quad \quad \operatorname{integral}(a, b, u, \operatorname{ptwise\text{-}liminf}(a, f)) \leq \operatorname{eliminf}((\lambda n \in \mathbb{N}.\; \operatorname{integral}(a, b, u, f(n)))) \end{array}}}
\vnbprop{esup-grid-max}{\vnbstmt{\begin{array}{@{}l@{}} \forall d, y \in \mathbb{R}, m \in \mathbb{N}, a. \\ \quad ((0 < d) \wedge ((a \subseteq \mathbb{N}) \wedge ((m \in a) \wedge ((y < ((m \cdot d) + d)) \wedge \forall k \in a.\; ((k \cdot d) \leq y))))) \Rightarrow  \\ \quad \quad \operatorname{esup}(\operatorname{image}((\lambda \mathit{sk} \in \mathbb{N}.\; (\mathit{sk} \cdot d)), a)) = (m \cdot d) \end{array}}}
\vnbprop{grid-staircase-value}{\vnbstmt{\begin{array}{@{}l@{}} \forall n \in \mathbb{N}, x \in \operatorname{rr\text{-}pos\text{-}star}. \exists m \in \mathbb{N}. \\ \quad m \in \{q \in \mathbb{N} : (((q \cdot \operatorname{power}(1/2, n)) \leq x) \wedge ((q \cdot \operatorname{power}(1/2, n)) \leq n))\} \wedge \\ \quad \operatorname{esup}(\operatorname{image}((\lambda k \in \mathbb{N}.\; (k \cdot \operatorname{power}(1/2, n))), \\ \quad \quad \quad \{q \in \mathbb{N} : (((q \cdot \operatorname{power}(1/2, n)) \leq x) \wedge ((q \cdot \operatorname{power}(1/2, n)) \leq n))\})) \\ \quad =\; m \cdot \operatorname{power}(1/2, n) \wedge \\ \quad \forall z \in \mathbb{R}.\; (((z \leq x) \wedge (z \leq n)) \Rightarrow (z < ((m \cdot \operatorname{power}(1/2, n)) + \operatorname{power}(1/2, n)))) \end{array}}}
\vnbprop{measurable-fn-guarded-ge}{\vnbstmt{\begin{array}{@{}l@{}} \forall c, d, f, a, b, g. \\ \quad (\operatorname{is\text{-}measurable\text{-}fn}(c, d, f) \wedge (b \in \mathbb{R})) \Rightarrow  \\ \quad \quad \{x \in c : (((a < b) \wedge (b \leq g)) \wedge (b \leq f(x)))\} \in d \end{array}}}
\vnbprop{grid-staircase-simple}{\vnbstmt{\begin{array}{@{}l@{}} \forall a, b, f.; \forall n \in \mathbb{N}. \\ \quad \operatorname{is\text{-}measurable\text{-}fn}(a, b, f) \Rightarrow  \\ \quad \quad \operatorname{is\text{-}simple\text{-}fn}(a, b, \\ \quad \quad \quad (\lambda x \in a.\; \\ \quad \quad \quad \quad \operatorname{esup}(\operatorname{image}((\lambda k \in \mathbb{N}.\; (k \cdot \operatorname{power}(1/2, n))), \\ \quad \quad \quad \quad \quad \quad \{q \in \mathbb{N} : (((q \cdot \operatorname{power}(1/2, n)) \leq f(x)) \wedge ((q \cdot \operatorname{power}(1/2, n)) \leq n))\})))) \end{array}}}
\vnbprop{integral-simple-approx}{\vnbstmt{\begin{array}{@{}l@{}} \forall a, b, f. \\ \quad (\operatorname{is\text{-}sigma\text{-}algebra}(a, b) \wedge \operatorname{is\text{-}measurable\text{-}fn}(a, b, f)) \Rightarrow  \\ \quad \quad \exists c \in (\mathbb{N} \to (a \to \operatorname{rr\text{-}pos\text{-}star})). \\ \quad \quad \quad \forall n \in \mathbb{N}.\; \operatorname{is\text{-}simple\text{-}fn}(a, b, c(n)) \wedge \\ \quad \quad \quad \forall n \in \mathbb{N}.\; \operatorname{ptwise\text{-}le}(a, c(n), c((n + 1))) \wedge \\ \quad \quad \quad \forall x \in a.\; (\operatorname{ptwise\text{-}sup}(a, c)(x) = f(x)) \end{array}}}
\vnbprop{esup-eplus-eq}{\vnbstmt{\begin{array}{@{}l@{}} \forall u, v. \\ \quad u \in (\mathbb{N} \to \operatorname{rr\text{-}pos\text{-}star}) \wedge \\ \quad v \in (\mathbb{N} \to \operatorname{rr\text{-}pos\text{-}star}) \wedge \\ \quad \forall n \in \mathbb{N}.\; (u(n) \leq u((n + 1))) \wedge \\ \quad \forall n \in \mathbb{N}.\; (v(n) \leq v((n + 1))) \Rightarrow \\ \quad \quad \operatorname{esup}(\operatorname{image}((\lambda a \in \mathbb{N}.\; \operatorname{eplus}(u(a), v(a))), \mathbb{N})) \\ \quad \quad =\; \operatorname{eplus}(\operatorname{esup}(\operatorname{image}(u, \mathbb{N})), \operatorname{esup}(\operatorname{image}(v, \mathbb{N}))) \end{array}}}
\vnbprop{simple-fn-eplus}{\vnbstmt{\begin{array}{@{}l@{}} \forall a, b, s, f. \\ \quad (\operatorname{is\text{-}simple\text{-}fn}(a, b, s) \wedge \operatorname{is\text{-}simple\text{-}fn}(a, b, f)) \Rightarrow  \\ \quad \quad \operatorname{is\text{-}simple\text{-}fn}(a, b, \operatorname{ptwise\text{-}eplus}(a, s, f)) \end{array}}}
\vnbprop{integral-additive}{\vnbstmt{\begin{array}{@{}l@{}} \forall a, b, u, f, g. \\ \quad (\operatorname{is\text{-}measure}(a, b, u) \wedge (\operatorname{is\text{-}measurable\text{-}fn}(a, b, f) \wedge \operatorname{is\text{-}measurable\text{-}fn}(a, b, g))) \Rightarrow  \\ \quad \quad \operatorname{integral}(a, b, u, \operatorname{ptwise\text{-}eplus}(a, f, g)) = \operatorname{eplus}(\operatorname{integral}(a, b, u, f), \operatorname{integral}(a, b, u, g)) \end{array}}}
\vnbprop{measurable-fn-lt-row}{\vnbstmt{\begin{array}{@{}l@{}} \forall a, b, u, v, r. \\ \quad (\operatorname{is\text{-}measurable\text{-}fn}(a, b, u) \wedge (\operatorname{is\text{-}measurable\text{-}fn}(a, b, v) \wedge (r \in \mathbb{R}))) \Rightarrow  \\ \quad \quad \{x \in a : \exists k \in \mathbb{N}.\; ((v(x) < (k \cdot r)) \wedge ((k \cdot r) < u(x)))\} \in b \end{array}}}
\vnbprop{measurable-fn-lt-set}{\vnbstmt{\begin{array}{@{}l@{}} \forall a, b, u, v. \\ \quad (\operatorname{is\text{-}measurable\text{-}fn}(a, b, u) \wedge \operatorname{is\text{-}measurable\text{-}fn}(a, b, v)) \Rightarrow  \\ \quad \quad \{x \in a : (v(x) < u(x))\} \in b \end{array}}}
\vnbprop{measurable-fn-const}{\vnbstmt{\begin{array}{@{}l@{}} \forall a, b, c. \\ \quad (\operatorname{is\text{-}sigma\text{-}algebra}(a, b) \wedge (c \in \operatorname{rr\text{-}pos\text{-}star})) \Rightarrow  \\ \quad \quad \operatorname{is\text{-}measurable\text{-}fn}(a, b, (\lambda x \in a.\; c)) \end{array}}}
\vnbprop{measurable-fn-finite-part}{\vnbstmt{\forall a, b, u.\; (\operatorname{is\text{-}measurable\text{-}fn}(a, b, u) \Rightarrow (\{p \in a : (u(p) \in \mathbb{R})\} \in b))}}
\vnbprop{integral-finite-null}{\vnbstmt{\begin{array}{@{}l@{}} \forall a, b, c, u. \\ \quad (\operatorname{is\text{-}measure}(a, b, c) \wedge (\operatorname{is\text{-}measurable\text{-}fn}(a, b, u) \wedge (\operatorname{integral}(a, b, c, u) \in \mathbb{R}))) \Rightarrow  \\ \quad \quad c(\operatorname{complement\text{-}in}(a, \{p \in a : (u(p) \in \mathbb{R})\})) = 0 \end{array}}}
\vnbprop{integral-le-off-null}{\vnbstmt{\begin{array}{@{}l@{}} \forall a, b, u, c, d, f. \\ \quad \operatorname{is\text{-}measure}(a, b, u) \wedge \\ \quad c \in b \wedge \\ \quad u(c) = 0 \wedge \\ \quad \operatorname{is\text{-}measurable\text{-}fn}(a, b, d) \wedge \\ \quad \operatorname{is\text{-}measurable\text{-}fn}(a, b, f) \wedge \\ \quad \forall x \in a.\; (\neg (x \in c) \Rightarrow (d(x) \leq f(x))) \Rightarrow \\ \quad \quad \operatorname{integral}(a, b, u, d) \leq \operatorname{integral}(a, b, u, f) \end{array}}}
\vnbprop{difference-value}{\vnbstmt{\begin{array}{@{}l@{}} \forall a, u, v.; \forall x \in a. \\ \quad ((\lambda \mathit{dx} \in a.\; \left(\text{if }(u(\mathit{dx}) \in \mathbb{R})\text{ then }(u(\mathit{dx}) - v(\mathit{dx}))\text{ else }0\right))(x) \simeq \left(\text{if }(u(x) \in \mathbb{R})\text{ then }(u(x) - v(x))\text{ else }0\right)) \end{array}}}
\vnbprop{measurable-fn-sub}{\vnbstmt{\begin{array}{@{}l@{}} \forall a, b, u, v. \\ \quad (\operatorname{is\text{-}measurable\text{-}fn}(a, b, u) \wedge (\operatorname{is\text{-}measurable\text{-}fn}(a, b, v) \wedge \forall x \in a.\; (v(x) \leq u(x)))) \Rightarrow  \\ \quad \quad \operatorname{is\text{-}measurable\text{-}fn}(a, b, (\lambda \mathit{dx} \in a.\; \left(\text{if }(u(\mathit{dx}) \in \mathbb{R})\text{ then }(u(\mathit{dx}) - v(\mathit{dx}))\text{ else }0\right))) \end{array}}}
\vnbprop{integral-sub}{\vnbstmt{\begin{array}{@{}l@{}} \forall a, b, c, u, v. \\ \quad \operatorname{is\text{-}measure}(a, b, c) \wedge \\ \quad \operatorname{is\text{-}measurable\text{-}fn}(a, b, u) \wedge \\ \quad \operatorname{is\text{-}measurable\text{-}fn}(a, b, v) \wedge \\ \quad \forall x \in a.\; (v(x) \leq u(x)) \wedge \\ \quad \operatorname{integral}(a, b, c, u) \in \mathbb{R} \Rightarrow \\ \quad \quad \operatorname{eplus}(\operatorname{integral}(a, b, c, (\lambda \mathit{dx} \in a.\; \left(\text{if }(u(\mathit{dx}) \in \mathbb{R})\text{ then }(u(\mathit{dx}) - v(\mathit{dx}))\text{ else }0\right))), \operatorname{integral}(a, b, c, v)) \\ \quad \quad =\; \operatorname{integral}(a, b, c, u) \end{array}}}
\vnbprop{elimsup-le-bound}{\vnbstmt{\begin{array}{@{}l@{}} \forall a \in (\mathbb{N} \to \operatorname{rr\text{-}pos\text{-}star}), b \in \operatorname{rr\text{-}pos\text{-}star}. \\ \quad \forall n \in \mathbb{N}.\; (a(n) \leq b) \Rightarrow (\operatorname{elimsup}(a) \leq b) \end{array}}}
\vnbprop{eliminf-le-elimsup}{\vnbstmt{\forall a \in (\mathbb{N} \to \operatorname{rr\text{-}pos\text{-}star}).\; (\operatorname{eliminf}(a) \leq \operatorname{elimsup}(a))}}
\vnbprop{einf-reflection-tail}{\vnbstmt{\begin{array}{@{}l@{}} \forall c \in \mathbb{R}, a, b.; \forall m \in \mathbb{N}. \\ \quad a \in (\mathbb{N} \to \operatorname{rr\text{-}pos\text{-}star}) \wedge \\ \quad b \in (\mathbb{N} \to \operatorname{rr\text{-}pos\text{-}star}) \wedge \\ \quad \forall n \in \mathbb{N}.\; (a(n) \leq c) \wedge \\ \quad \forall n \in \mathbb{N}.\; (b(n) = (c - a(n))) \Rightarrow \\ \quad \quad (\operatorname{esup}(\operatorname{etail}(a, m)) \in \mathbb{R}) \wedge (\operatorname{einf}(\operatorname{etail}(b, m)) = (c - \operatorname{esup}(\operatorname{etail}(a, m)))) \end{array}}}
\vnbprop{eliminf-reflection}{\vnbstmt{\begin{array}{@{}l@{}} \forall c \in \mathbb{R}, a, b. \\ \quad a \in (\mathbb{N} \to \operatorname{rr\text{-}pos\text{-}star}) \wedge \\ \quad b \in (\mathbb{N} \to \operatorname{rr\text{-}pos\text{-}star}) \wedge \\ \quad \forall n \in \mathbb{N}.\; (a(n) \leq c) \wedge \\ \quad \forall n \in \mathbb{N}.\; (b(n) = (c - a(n))) \Rightarrow \\ \quad \quad (\operatorname{elimsup}(a) \in \mathbb{R}) \wedge (\operatorname{eliminf}(b) = (c - \operatorname{elimsup}(a))) \end{array}}}
\vnbprop{dominated-convergence}{\vnbstmt{\begin{array}{@{}l@{}} \forall a, b, u, c, f, g. \\ \quad \operatorname{is\text{-}measure}(a, b, u) \wedge \\ \quad c \in (\mathbb{N} \to (a \to \operatorname{rr\text{-}pos\text{-}star})) \wedge \\ \quad \forall n \in \mathbb{N}.\; \operatorname{is\text{-}measurable\text{-}fn}(a, b, c(n)) \wedge \\ \quad \operatorname{is\text{-}measurable\text{-}fn}(a, b, f) \wedge \\ \quad \operatorname{is\text{-}measurable\text{-}fn}(a, b, g) \wedge \\ \quad \forall n \in \mathbb{N}.\; \operatorname{ptwise\text{-}le}(a, c(n), g) \wedge \\ \quad \operatorname{integral}(a, b, u, g) \in \mathbb{R} \wedge \\ \quad \forall x \in a.\; \operatorname{econverges\text{-}to}((\lambda n \in \mathbb{N}.\; c(n)(x)), f(x)) \Rightarrow \\ \quad \quad \operatorname{econverges\text{-}to}((\lambda n \in \mathbb{N}.\; \operatorname{integral}(a, b, u, c(n))), \operatorname{integral}(a, b, u, f)) \end{array}}}

\vnbbookfile{theorem-library/caratheodory}{theorem-library/caratheodory}
\noindent{\small caratheodory.scm -- THE CARATHEODORY CONSTRUCTION, PROVEN. Agent MEASURE-1, 2026-10-07 (scratchpad/triage/MEASURE-1.md).  Library: measure.  Helper prefix cth-.  16 theorems, every one `modulo 0'.}\par
\vnbprop{difference-unfold}{\vnbstmt{\forall u, b.\; (\operatorname{difference}(u, b) \simeq \operatorname{complement\text{-}in}(u, b))}}
\vnbprop{inter-complement-in-shift}{\vnbstmt{\begin{array}{@{}l@{}} \forall a, u, e. \\ \quad (\operatorname{complement\text{-}in}(a, u) \cap e) = (a \cap \operatorname{complement\text{-}in}(e, u)) \end{array}}}
\vnbprop{complement-in-complement-in-union}{\vnbstmt{\begin{array}{@{}l@{}} \forall a, u, e. \\ \quad \operatorname{complement\text{-}in}(\operatorname{complement\text{-}in}(a, u), e) = \operatorname{complement\text{-}in}(a, (u \cup e)) \end{array}}}
\vnbprop{complement-in-empty-set}{\vnbstmt{\forall a.\; (\operatorname{complement\text{-}in}(a, \emptyset) = a)}}
\vnbprop{union-inter-complement-in-split}{\vnbstmt{\forall a, e.\; (a = ((a \cap e) \cup \operatorname{complement\text{-}in}(a, e)))}}
\vnbprop{inter-complement-in-ambient}{\vnbstmt{\begin{array}{@{}l@{}} \forall t, m, e. \\ \quad (t \subseteq m) \Rightarrow ((t \cap \operatorname{complement\text{-}in}(m, e)) = \operatorname{complement\text{-}in}(t, e)) \end{array}}}
\vnbprop{complement-in-complement-in-ambient}{\vnbstmt{\begin{array}{@{}l@{}} \forall t, m, e. \\ \quad (t \subseteq m) \Rightarrow (\operatorname{complement\text{-}in}(t, \operatorname{complement\text{-}in}(m, e)) = (t \cap e)) \end{array}}}
\vnbprop{inter-of-subset}{\vnbstmt{\forall t, m.\; ((t \subseteq m) \Rightarrow ((t \cap m) = t))}}
\vnbprop{complement-in-of-subset-empty}{\vnbstmt{\forall t, m.\; ((t \subseteq m) \Rightarrow (\operatorname{complement\text{-}in}(t, m) = \emptyset))}}
\vnbprop{outer-measure-binary-subadditive}{\vnbstmt{\begin{array}{@{}l@{}} \forall a, u, \alpha, b. \\ \quad (\operatorname{is\text{-}outer\text{-}measure}(a, u) \wedge ((\alpha \in \operatorname{power}(a)) \wedge (b \in \operatorname{power}(a)))) \Rightarrow  \\ \quad \quad u((\alpha \cup b)) \leq \operatorname{eplus}(u(\alpha), u(b)) \end{array}}}
\vnbprop{outer-measure-split-le}{\vnbstmt{\begin{array}{@{}l@{}} \forall a, u, \alpha, s. \\ \quad (\operatorname{is\text{-}outer\text{-}measure}(a, u) \wedge ((\alpha \in \operatorname{power}(a)) \wedge (s \in \operatorname{power}(a)))) \Rightarrow  \\ \quad \quad u(\alpha) \leq \operatorname{eplus}(u((\alpha \cap s)), u(\operatorname{complement\text{-}in}(\alpha, s))) \end{array}}}
\vnbprop{caratheodory-null-measurable}{\vnbstmt{\begin{array}{@{}l@{}} \forall a, u, s. \\ \quad (\operatorname{is\text{-}outer\text{-}measure}(a, u) \wedge ((s \subseteq a) \wedge (u(s) = 0))) \Rightarrow  \\ \quad \quad \operatorname{is\text{-}caratheodory\text{-}measurable}(a, u, s) \end{array}}}
\vnbprop{caratheodory-superadditive-suffices}{\vnbstmt{\begin{array}{@{}l@{}} \forall a, u, s. \\ \quad \operatorname{is\text{-}outer\text{-}measure}(a, u) \wedge \\ \quad s \subseteq a \wedge \\ \quad \forall \alpha \in \operatorname{power}(a). \\ \quad \quad \operatorname{eplus}(u((\alpha \cap s)), u(\operatorname{difference}(\alpha, s))) \leq u(\alpha) \Rightarrow \\ \quad \quad \operatorname{is\text{-}caratheodory\text{-}measurable}(a, u, s) \end{array}}}
\vnbprop{caratheodory-prefix-split}{\vnbstmt{\begin{array}{@{}l@{}} \forall n \in \mathbb{N}, a, u, f, \alpha. \\ \quad (\operatorname{is\text{-}outer\text{-}measure}(a, u) \wedge ((f \in (\mathbb{N} \to \operatorname{caratheodory\text{-}sets}(a, u))) \wedge (\alpha \in \operatorname{power}(a)))) \Rightarrow  \\ \quad \quad u(\alpha) \\ \quad \quad =\; \operatorname{eplus}(\operatorname{finsum}(\operatorname{rr\text{-}pos\text{-}star\text{-}add\text{-}monoid}, \\ \quad \quad \quad \quad (\lambda k \in \mathbb{N}.\; u((\alpha \cap \operatorname{complement\text{-}in}(f(k), \operatorname{big\text{-}union}(j, \mathrm{OS}(k), f(j)))))), \\ \quad \quad \quad \quad \mathrm{OS}(n)), \\ \quad \quad \quad u(\operatorname{complement\text{-}in}(\alpha, \operatorname{big\text{-}union}(j, \mathrm{OS}(n), f(j))))) \end{array}}}
\vnbprop{caratheodory-measurable-sigma-algebra}{\vnbstmt{\begin{array}{@{}l@{}} \forall a, u. \\ \quad \operatorname{is\text{-}outer\text{-}measure}(a, u) \Rightarrow \operatorname{is\text{-}sigma\text{-}algebra}(a, \operatorname{caratheodory\text{-}sets}(a, u)) \end{array}}}
\vnbprop{caratheodory-measure}{\vnbstmt{\begin{array}{@{}l@{}} \forall a, u. \\ \quad \operatorname{is\text{-}outer\text{-}measure}(a, u) \Rightarrow  \\ \quad \quad \operatorname{is\text{-}measure}(a, \operatorname{caratheodory\text{-}sets}(a, u), \operatorname{res}(u, \operatorname{caratheodory\text{-}sets}(a, u))) \end{array}}}

\vnbbookfile{structure-library/outer-measure-construction}{structure-library/outer-measure-construction}
\noindent{\small outer-measure-construction.scm -- METHOD I OUTER MEASURES; null sets, complete and sigma-finite measures.  Vocabulary only: no proof, no support. Agent MEASURE-2, 2026-10-08 (scratchpad/triage/MEASURE-2.md).  The proofs are theorem-library/outer-measure-construction.scm.}\par
\par\medskip\noindent\textbf{Definition} (\texttt{is-premeasure}).\;For all $\omega$, $\mathit{cls}$ and $\mathit{pm}$, 
\begin{equation*}
\begin{array}{@{}l@{}} \operatorname{is\text{-}premeasure}(\omega, \mathit{cls}, \mathit{pm}) \\ \Leftrightarrow\; \omega \in \mathbf{Set} \wedge \\ \mathit{cls} \subseteq \operatorname{power}(\omega) \wedge \\ \emptyset \in \mathit{cls} \wedge \\ \mathit{pm} \in (\mathit{cls} \to \operatorname{rr\text{-}pos\text{-}star}) \wedge \\ \operatorname{pm}(\emptyset) = 0 \end{array}
\end{equation*}

\par\medskip\noindent\textbf{Definition} (\texttt{cover-sums}).\;$\operatorname{cover\text{-}sums}(\omega, \mathit{cls}, \mathit{pm}, \mathit{es}) :=$
\begin{equation*}
\{\operatorname{cvv\_} \in \operatorname{rr\text{-}pos\text{-}star} : \exists \operatorname{cvs\_} \in (\mathbb{N} \to \mathit{cls}).\; ((\mathit{es} \subseteq \operatorname{big\text{-}union}(\operatorname{cvn\_}, \mathbb{N}, \operatorname{cvs\_}(\operatorname{cvn\_}))) \wedge (\operatorname{cvv\_} = \operatorname{esum}((\lambda \operatorname{cvk\_} \in \mathbb{N}.\; \operatorname{pm}(\operatorname{cvs\_}(\operatorname{cvk\_}))))))\}
\end{equation*}

\par\medskip\noindent\textbf{Definition} (\texttt{method-i-outer}).\;$\operatorname{method\text{-}i\text{-}outer}(\omega, \mathit{cls}, \mathit{pm}) :=$
\begin{equation*}
(\lambda \operatorname{cve\_} \in \operatorname{power}(\omega).\; \operatorname{einf}(\operatorname{cover\text{-}sums}(\omega, \mathit{cls}, \mathit{pm}, \operatorname{cve\_})))
\end{equation*}

\par\medskip\noindent\textbf{Definition} (\texttt{is-null-set}).\;For all $\omega$, $\mathit{ca}$, $\mu$ and $\mathit{ns}$, 
\begin{equation*}
\operatorname{is\text{-}null\text{-}set}(\omega, \mathit{ca}, \mu, \mathit{ns}) \Leftrightarrow \exists \operatorname{nla\_} \in \mathit{ca}.\; ((\mathit{ns} \subseteq \operatorname{nla\_}) \wedge (\mu(\operatorname{nla\_}) = 0))
\end{equation*}

\par\medskip\noindent\textbf{Definition} (\texttt{is-complete-measure}).\;For all $\omega$, $\mathit{ca}$ and $\mu$, 
\begin{equation*}
\begin{array}{@{}l@{}} \operatorname{is\text{-}complete\text{-}measure}(\omega, \mathit{ca}, \mu) \\ \Leftrightarrow\; (\operatorname{is\text{-}measure}(\omega, \mathit{ca}, \mu) \wedge \forall \operatorname{nls\_}.\; (\operatorname{is\text{-}null\text{-}set}(\omega, \mathit{ca}, \mu, \operatorname{nls\_}) \Rightarrow (\operatorname{nls\_} \in \mathit{ca}))) \end{array}
\end{equation*}

\par\medskip\noindent\textbf{Definition} (\texttt{is-sigma-finite}).\;For all $\omega$, $\mathit{ca}$ and $\mu$, 
\begin{equation*}
\begin{array}{@{}l@{}} \operatorname{is\text{-}sigma\text{-}finite}(\omega, \mathit{ca}, \mu) \\ \Leftrightarrow\; \exists \operatorname{sff\_} \in (\mathbb{N} \to \mathit{ca}). \\ \quad (\forall \operatorname{sfn\_} \in \mathbb{N}.\; (\mu(\operatorname{sff\_}(\operatorname{sfn\_})) \in \mathbb{R}) \wedge (\operatorname{big\text{-}union}(\operatorname{sfk\_}, \mathbb{N}, \operatorname{sff\_}(\operatorname{sfk\_})) = \omega)) \end{array}
\end{equation*}

\par\medskip\noindent\textbf{Definition} (\texttt{closed-intervals}).\;
\begin{equation*}
(\operatorname{closed\text{-}intervals} \simeq \{\operatorname{oci\_} \in \operatorname{power}(\mathbb{R}) : ((\operatorname{oci\_} = \emptyset) \vee \exists \operatorname{oca\_}, \operatorname{ocb\_} \in \mathbb{R}.\; ((\operatorname{oca\_} \leq \operatorname{ocb\_}) \wedge (\operatorname{oci\_} = \operatorname{ccint}(\operatorname{oca\_}, \operatorname{ocb\_}))))\})
\end{equation*}

\par\medskip\noindent\textbf{Definition} (\texttt{interval-length}).\;$\operatorname{interval\text{-}length}(s) :=$
\begin{equation*}
\operatorname{esup}(\{\operatorname{odv\_} \in \operatorname{rr\text{-}pos\text{-}star} : ((\operatorname{odv\_} = 0) \vee \exists \operatorname{odx\_}, \operatorname{ody\_} \in s.\; ((\operatorname{odx\_} \leq \operatorname{ody\_}) \wedge (\operatorname{odv\_} = (\operatorname{ody\_} - \operatorname{odx\_}))))\})
\end{equation*}


\vnbbookfile{theorem-library/outer-measure-construction}{theorem-library/outer-measure-construction}
\noindent{\small outer-measure-construction.scm -- METHOD I OUTER MEASURES; the outer measure of a measure, its Caratheodory measure, completeness and the completion; Lebesgue outer measure on the line as an instance. Agent MEASURE-2, 2026-10-08 (scratchpad/triage/MEASURE-2.md).  Library:}\par
\vnbprop{cover-sums-unfold}{\vnbstmt{\begin{array}{@{}l@{}} \forall a, s, m, b. \\ \quad (\operatorname{cover\text{-}sums}(a, s, m, b) \simeq \{v \in \operatorname{rr\text{-}pos\text{-}star} : \exists \mathit{cvs} \in (\mathbb{N} \to s).\; ((b \subseteq \operatorname{big\text{-}union}(n, \mathbb{N}, \operatorname{cvs}(n))) \wedge (v = \operatorname{esum}((\lambda k \in \mathbb{N}.\; m(\operatorname{cvs}(k))))))\}) \end{array}}}
\vnbprop{cover-sums-subset}{\vnbstmt{\forall a, s, m, b.\; (\operatorname{cover\text{-}sums}(a, s, m, b) \subseteq \operatorname{rr\text{-}pos\text{-}star})}}
\vnbprop{cover-sums-elim}{\vnbstmt{\begin{array}{@{}l@{}} \forall a, s, m, b.; \forall v \in \operatorname{cover\text{-}sums}(a, s, m, b). \exists \mathit{cvs} \in (\mathbb{N} \to s). \\ \quad (b \subseteq \operatorname{big\text{-}union}(n, \mathbb{N}, \operatorname{cvs}(n))) \wedge (v = \operatorname{esum}((\lambda k \in \mathbb{N}.\; m(\operatorname{cvs}(k))))) \end{array}}}
\vnbprop{cover-sums-intro}{\vnbstmt{\begin{array}{@{}l@{}} \forall a, s, m, b, q. \\ \quad ((m \in (s \to \operatorname{rr\text{-}pos\text{-}star})) \wedge ((q \in (\mathbb{N} \to s)) \wedge (b \subseteq \operatorname{big\text{-}union}(n, \mathbb{N}, q(n))))) \Rightarrow  \\ \quad \quad \operatorname{esum}((\lambda k \in \mathbb{N}.\; m(q(k)))) \in \operatorname{cover\text{-}sums}(a, s, m, b) \end{array}}}
\vnbprop{method-i-outer-value}{\vnbstmt{\begin{array}{@{}l@{}} \forall a, s, m.; \forall b \in \operatorname{power}(a). \\ \quad \operatorname{method\text{-}i\text{-}outer}(a, s, m)(b) = \operatorname{einf}(\operatorname{cover\text{-}sums}(a, s, m, b)) \end{array}}}
\vnbprop{method-i-outer-type}{\vnbstmt{\begin{array}{@{}l@{}} \forall a \in \mathbf{Set}, s, m. \\ \quad \operatorname{method\text{-}i\text{-}outer}(a, s, m) \in (\operatorname{power}(a) \to \operatorname{rr\text{-}pos\text{-}star}) \end{array}}}
\vnbprop{einf-lt-witness}{\vnbstmt{\begin{array}{@{}l@{}} \forall s.; \forall a \in \operatorname{rr\text{-}pos\text{-}star}. \\ \quad ((s \subseteq \operatorname{rr\text{-}pos\text{-}star}) \wedge (\operatorname{einf}(s) < a)) \Rightarrow \exists y \in s.\; (y < a) \end{array}}}
\vnbprop{rps-le-of-forall-eps}{\vnbstmt{\begin{array}{@{}l@{}} \forall x, y. \\ \quad ((x \in \operatorname{rr\text{-}pos\text{-}star}) \wedge ((y \in \mathbb{R}) \wedge ((0 \leq y) \wedge \forall s.\; (\operatorname{pos\text{-}rr}(s) \Rightarrow (x \leq \operatorname{eplus}(y, s)))))) \Rightarrow  \\ \quad \quad x \leq y \end{array}}}
\vnbprop{eplus-swap4}{\vnbstmt{\begin{array}{@{}l@{}} \forall a, b, c, d. \\ \quad ((a \in \operatorname{rr\text{-}pos\text{-}star}) \wedge ((b \in \operatorname{rr\text{-}pos\text{-}star}) \wedge ((c \in \operatorname{rr\text{-}pos\text{-}star}) \wedge (d \in \operatorname{rr\text{-}pos\text{-}star})))) \Rightarrow  \\ \quad \quad \operatorname{eplus}(\operatorname{eplus}(a, b), \operatorname{eplus}(c, d)) = \operatorname{eplus}(\operatorname{eplus}(a, c), \operatorname{eplus}(b, d)) \end{array}}}
\vnbprop{rps-prefix-add}{\vnbstmt{\begin{array}{@{}l@{}} \forall n \in \mathbb{N}, f, g. \\ \quad ((f \in (\mathbb{N} \to \operatorname{rr\text{-}pos\text{-}star})) \wedge (g \in (\mathbb{N} \to \operatorname{rr\text{-}pos\text{-}star}))) \Rightarrow  \\ \quad \quad \operatorname{finsum}(\operatorname{rr\text{-}pos\text{-}star\text{-}add\text{-}monoid}, (\lambda k \in \mathbb{N}.\; \operatorname{eplus}(f(k), g(k))), \mathrm{OS}(n)) \\ \quad \quad =\; \operatorname{eplus}(\operatorname{finsum}(\operatorname{rr\text{-}pos\text{-}star\text{-}add\text{-}monoid}, f, \mathrm{OS}(n)), \\ \quad \quad \quad \operatorname{finsum}(\operatorname{rr\text{-}pos\text{-}star\text{-}add\text{-}monoid}, g, \mathrm{OS}(n))) \end{array}}}
\vnbprop{esum-nn-add-le}{\vnbstmt{\begin{array}{@{}l@{}} \forall f, g. \\ \quad ((f \in (\mathbb{N} \to \operatorname{rr\text{-}pos\text{-}star})) \wedge (g \in (\mathbb{N} \to \operatorname{rr\text{-}pos\text{-}star}))) \Rightarrow  \\ \quad \quad \operatorname{esum}((\lambda k \in \mathbb{N}.\; \operatorname{eplus}(f(k), g(k)))) \leq \operatorname{eplus}(\operatorname{esum}(f), \operatorname{esum}(g)) \end{array}}}
\vnbprop{dyadic-eps-type}{\vnbstmt{\begin{array}{@{}l@{}} \forall s. \\ \quad \operatorname{pos\text{-}rr}(s) \Rightarrow  \\ \quad \quad (\lambda i \in \mathbb{N}.\; (s \cdot \operatorname{recip}(\operatorname{power}(2, (i + 1))))) \in (\mathbb{N} \to \operatorname{rr\text{-}pos\text{-}star}) \wedge \\ \quad \quad \forall k \in \mathbb{N}.\; (0 < (\lambda i \in \mathbb{N}.\; (s \cdot \operatorname{recip}(\operatorname{power}(2, (i + 1)))))(k)) \end{array}}}
\vnbprop{dyadic-eps-prefix}{\vnbstmt{\begin{array}{@{}l@{}} \forall a \in \mathbb{N}, s. \\ \quad \operatorname{pos\text{-}rr}(s) \Rightarrow  \\ \quad \quad \operatorname{finsum}(\operatorname{rr\text{-}pos\text{-}star\text{-}add\text{-}monoid}, (\lambda i \in \mathbb{N}.\; (s \cdot \operatorname{recip}(\operatorname{power}(2, (i + 1))))), \mathrm{OS}(a)) \\ \quad \quad =\; s \cdot \operatorname{series\text{-}partial\text{-}sum}((\lambda n \in \mathbb{N}.\; /(1, \operatorname{power}(2, (n + 1)))), a) \end{array}}}
\vnbprop{dyadic-eps-sum-le}{\vnbstmt{\begin{array}{@{}l@{}} \forall s. \\ \quad \operatorname{pos\text{-}rr}(s) \Rightarrow (\operatorname{esum}((\lambda i \in \mathbb{N}.\; (s \cdot \operatorname{recip}(\operatorname{power}(2, (i + 1)))))) \leq s) \end{array}}}
\vnbprop{trinum-ge}{\vnbstmt{\forall s \in \mathbb{N}.\; (s \leq \operatorname{trinum}(s))}}
\vnbprop{nn-unpair-le}{\vnbstmt{\forall v \in \mathbb{N}.\; ((\operatorname{nnfst}(v) \leq v) \wedge (\operatorname{nnsnd}(v) \leq v))}}
\vnbprop{rps-finsum-fiber-le}{\vnbstmt{\begin{array}{@{}l@{}} \forall s, t, p, f, g. \\ \quad s \in \mathbf{Set} \wedge \\ \quad |s| \in \mathbb{N} \wedge \\ \quad t \in \mathbf{Set} \wedge \\ \quad |t| \in \mathbb{N} \wedge \\ \quad \forall z \in s.\; (p(z) \in t) \wedge \\ \quad \forall z \in s.\; (f(z) \in \operatorname{rr\text{-}pos\text{-}star}) \wedge \\ \quad \forall z \in t.\; (g(z) \in \operatorname{rr\text{-}pos\text{-}star}) \wedge \\ \quad \forall z \in t.\; (\operatorname{finsum}(\operatorname{rr\text{-}pos\text{-}star\text{-}add\text{-}monoid}, f, \{w \in s : (p(w) = z)\}) \leq g(z)) \Rightarrow \\ \quad \quad \operatorname{finsum}(\operatorname{rr\text{-}pos\text{-}star\text{-}add\text{-}monoid}, f, s) \leq \operatorname{finsum}(\operatorname{rr\text{-}pos\text{-}star\text{-}add\text{-}monoid}, g, t) \end{array}}}
\vnbprop{esum-nn-flatten-le}{\vnbstmt{\begin{array}{@{}l@{}} \forall a, l, b. \\ \quad a \in (\mathbb{N} \to (\mathbb{N} \to \operatorname{rr\text{-}pos\text{-}star})) \wedge \\ \quad l \in (\mathbb{N} \to \operatorname{rr\text{-}pos\text{-}star}) \wedge \\ \quad b \in (\mathbb{N} \to \operatorname{rr\text{-}pos\text{-}star}) \wedge \\ \quad \forall m \in \mathbb{N}.\; (l(m) = a(\operatorname{nnfst}(m))(\operatorname{nnsnd}(m))) \wedge \\ \quad \forall i \in \mathbb{N}.\; (b(i) = \operatorname{esum}(a(i))) \Rightarrow \\ \quad \quad \operatorname{esum}(l) \leq \operatorname{esum}(b) \end{array}}}
\vnbprop{esum-nn-term-le}{\vnbstmt{\forall f \in (\mathbb{N} \to \operatorname{rr\text{-}pos\text{-}star}), k \in \mathbb{N}.\; (f(k) \leq \operatorname{esum}(f))}}
\vnbprop{method-i-outer-empty}{\vnbstmt{\begin{array}{@{}l@{}} \forall a, s, m. \\ \quad \operatorname{is\text{-}premeasure}(a, s, m) \Rightarrow (\operatorname{einf}(\operatorname{cover\text{-}sums}(a, s, m, \emptyset)) = 0) \end{array}}}
\vnbprop{method-i-outer-mono}{\vnbstmt{\begin{array}{@{}l@{}} \forall a, s, m, b, t. \\ \quad ((m \in (s \to \operatorname{rr\text{-}pos\text{-}star})) \wedge (b \subseteq t)) \Rightarrow  \\ \quad \quad \operatorname{einf}(\operatorname{cover\text{-}sums}(a, s, m, b)) \leq \operatorname{einf}(\operatorname{cover\text{-}sums}(a, s, m, t)) \end{array}}}
\vnbprop{method-i-outer-subadditive}{\vnbstmt{\begin{array}{@{}l@{}} \forall a, s, m, f. \\ \quad (\operatorname{is\text{-}premeasure}(a, s, m) \wedge (f \in (\mathbb{N} \to \operatorname{power}(a)))) \Rightarrow  \\ \quad \quad \operatorname{einf}(\operatorname{cover\text{-}sums}(a, s, m, \operatorname{big\text{-}union}(n, \mathbb{N}, f(n)))) \\ \quad \quad \leq\; \operatorname{esum}((\lambda k \in \mathbb{N}.\; \operatorname{einf}(\operatorname{cover\text{-}sums}(a, s, m, f(k))))) \end{array}}}
\vnbprop{method-i-outer-is-outer-measure}{\vnbstmt{\begin{array}{@{}l@{}} \forall a, s, m. \\ \quad \operatorname{is\text{-}premeasure}(a, s, m) \Rightarrow \operatorname{is\text{-}outer\text{-}measure}(a, \operatorname{method\text{-}i\text{-}outer}(a, s, m)) \end{array}}}
\vnbprop{method-i-outer-le-single}{\vnbstmt{\begin{array}{@{}l@{}} \forall a, s, m, b, c. \\ \quad (\operatorname{is\text{-}premeasure}(a, s, m) \wedge ((c \in s) \wedge (b \subseteq c))) \Rightarrow  \\ \quad \quad \operatorname{einf}(\operatorname{cover\text{-}sums}(a, s, m, b)) \leq m(c) \end{array}}}
\vnbprop{measure-is-premeasure}{\vnbstmt{\forall a, b, o.\; (\operatorname{is\text{-}measure}(a, b, o) \Rightarrow \operatorname{is\text{-}premeasure}(a, b, o))}}
\vnbprop{cover-sum-dominates}{\vnbstmt{\begin{array}{@{}l@{}} \forall a, b, o, s, y. \\ \quad (\operatorname{is\text{-}measure}(a, b, o) \wedge (y \in \operatorname{cover\text{-}sums}(a, b, o, s))) \Rightarrow  \\ \quad \quad \exists \mathit{oda} \in b.\; ((s \subseteq \mathit{oda}) \wedge (o(\mathit{oda}) \leq y)) \end{array}}}
\vnbprop{method-i-outer-of-measure}{\vnbstmt{\begin{array}{@{}l@{}} \forall a, b, o, s. \\ \quad (\operatorname{is\text{-}measure}(a, b, o) \wedge (s \in \operatorname{power}(a))) \Rightarrow  \\ \quad \quad \operatorname{method\text{-}i\text{-}outer}(a, b, o)(s) \\ \quad \quad =\; \operatorname{einf}(\{v \in \operatorname{rr\text{-}pos\text{-}star} : \exists \mathit{osa} \in b.\; ((s \subseteq \mathit{osa}) \wedge (v = o(\mathit{osa})))\}) \end{array}}}
\vnbprop{method-i-outer-extends}{\vnbstmt{\begin{array}{@{}l@{}} \forall a, b, o, c. \\ \quad (\operatorname{is\text{-}measure}(a, b, o) \wedge (c \in b)) \Rightarrow (\operatorname{method\text{-}i\text{-}outer}(a, b, o)(c) = o(c)) \end{array}}}
\vnbprop{method-i-outer-measurable}{\vnbstmt{\begin{array}{@{}l@{}} \forall a, b, o, c. \\ \quad (\operatorname{is\text{-}measure}(a, b, o) \wedge (c \in b)) \Rightarrow  \\ \quad \quad \operatorname{is\text{-}caratheodory\text{-}measurable}(a, \operatorname{method\text{-}i\text{-}outer}(a, b, o), c) \end{array}}}
\vnbprop{sigma-algebra-countable-inter}{\vnbstmt{\begin{array}{@{}l@{}} \forall a, b, f. \\ \quad (\operatorname{is\text{-}sigma\text{-}algebra}(a, b) \wedge (f \in (\mathbb{N} \to b))) \Rightarrow  \\ \quad \quad \operatorname{complement\text{-}in}(a, \operatorname{big\text{-}union}(n, \mathbb{N}, \operatorname{complement\text{-}in}(a, f(n)))) \in b \end{array}}}
\vnbprop{method-i-outer-attained}{\vnbstmt{\begin{array}{@{}l@{}} \forall a, b, o, s. \\ \quad (\operatorname{is\text{-}measure}(a, b, o) \wedge ((s \in \operatorname{power}(a)) \wedge (\operatorname{method\text{-}i\text{-}outer}(a, b, o)(s) \in \mathbb{R}))) \Rightarrow  \\ \quad \quad \exists z \in b.\; ((s \subseteq z) \wedge (o(z) = \operatorname{method\text{-}i\text{-}outer}(a, b, o)(s))) \end{array}}}
\vnbprop{method-i-outer-null-iff}{\vnbstmt{\begin{array}{@{}l@{}} \forall a, b, o, s. \\ \quad (\operatorname{is\text{-}measure}(a, b, o) \wedge (s \in \operatorname{power}(a))) \Rightarrow  \\ \quad \quad (\operatorname{method\text{-}i\text{-}outer}(a, b, o)(s) = 0) \Leftrightarrow \operatorname{is\text{-}null\text{-}set}(a, b, o, s) \end{array}}}
\vnbprop{method-i-outer-finite-approx}{\vnbstmt{\begin{array}{@{}l@{}} \forall a, b, o, s. \\ \quad \operatorname{is\text{-}measure}(a, b, o) \wedge \\ \quad \operatorname{is\text{-}caratheodory\text{-}measurable}(a, \operatorname{method\text{-}i\text{-}outer}(a, b, o), s) \wedge \\ \quad \operatorname{method\text{-}i\text{-}outer}(a, b, o)(s) \in \mathbb{R} \Rightarrow \\ \quad \quad \exists y \in b.\; ((s \subseteq y) \wedge \operatorname{is\text{-}null\text{-}set}(a, b, o, \operatorname{complement\text{-}in}(y, s))) \end{array}}}
\vnbprop{caratheodory-measure-complete}{\vnbstmt{\begin{array}{@{}l@{}} \forall a, u. \\ \quad \operatorname{is\text{-}outer\text{-}measure}(a, u) \Rightarrow  \\ \quad \quad \operatorname{is\text{-}complete\text{-}measure}(a, \operatorname{caratheodory\text{-}sets}(a, u), \operatorname{res}(u, \operatorname{caratheodory\text{-}sets}(a, u))) \end{array}}}
\vnbprop{method-i-outer-caratheodory-complete}{\vnbstmt{\begin{array}{@{}l@{}} \forall a, b, o. \\ \quad \operatorname{is\text{-}measure}(a, b, o) \Rightarrow  \\ \quad \quad \operatorname{is\text{-}complete\text{-}measure}(a, \operatorname{caratheodory\text{-}sets}(a, \operatorname{method\text{-}i\text{-}outer}(a, b, o)), \\ \quad \quad \quad \operatorname{res}(\operatorname{method\text{-}i\text{-}outer}(a, b, o), \operatorname{caratheodory\text{-}sets}(a, \operatorname{method\text{-}i\text{-}outer}(a, b, o)))) \end{array}}}
\vnbprop{method-i-outer-extends-completion}{\vnbstmt{\begin{array}{@{}l@{}} \forall a, b, o, c, s. \\ \quad (\operatorname{is\text{-}measure}(a, b, o) \wedge ((c \in b) \wedge \operatorname{is\text{-}null\text{-}set}(a, b, o, s))) \Rightarrow  \\ \quad \quad (c \cup s) \in \operatorname{caratheodory\text{-}sets}(a, \operatorname{method\text{-}i\text{-}outer}(a, b, o)) \wedge \\ \quad \quad \operatorname{res}(\operatorname{method\text{-}i\text{-}outer}(a, b, o), \operatorname{caratheodory\text{-}sets}(a, \operatorname{method\text{-}i\text{-}outer}(a, b, o)))((c \cup s)) = o(c) \end{array}}}
\vnbprop{esum-nn-zero}{\vnbstmt{\begin{array}{@{}l@{}} \forall f \in (\mathbb{N} \to \operatorname{rr\text{-}pos\text{-}star}). \\ \quad \forall n \in \mathbb{N}.\; (f(n) = 0) \Rightarrow (\operatorname{esum}(f) = 0) \end{array}}}
\vnbprop{method-i-outer-sigma-finite-approx}{\vnbstmt{\begin{array}{@{}l@{}} \forall a, b, o, s. \\ \quad \operatorname{is\text{-}measure}(a, b, o) \wedge \\ \quad \operatorname{is\text{-}sigma\text{-}finite}(a, b, o) \wedge \\ \quad \operatorname{is\text{-}caratheodory\text{-}measurable}(a, \operatorname{method\text{-}i\text{-}outer}(a, b, o), s) \Rightarrow \\ \quad \quad \exists y \in b.\; ((s \subseteq y) \wedge \operatorname{is\text{-}null\text{-}set}(a, b, o, \operatorname{complement\text{-}in}(y, s))) \end{array}}}
\vnbprop{method-i-outer-sigma-finite-completion}{\vnbstmt{\begin{array}{@{}l@{}} \forall a, b, o, s. \\ \quad \operatorname{is\text{-}measure}(a, b, o) \wedge \\ \quad \operatorname{is\text{-}sigma\text{-}finite}(a, b, o) \wedge \\ \quad s \in \operatorname{caratheodory\text{-}sets}(a, \operatorname{method\text{-}i\text{-}outer}(a, b, o)) \Rightarrow \\ \quad \quad \exists \mathit{ofa} \in b, \mathit{ofb}.\; (\operatorname{is\text{-}null\text{-}set}(a, b, o, \mathit{ofb}) \wedge (s = (\mathit{ofa} \cup \mathit{ofb}))) \end{array}}}
\vnbprop{lebesgue-premeasure}{\vnbstmt{\operatorname{is\text{-}premeasure}(\mathbb{R}, \operatorname{closed\text{-}intervals}, (\lambda s \in \operatorname{closed\text{-}intervals}.\; \operatorname{interval\text{-}length}(s)))}}
\vnbprop{lebesgue-outer-measure}{\vnbstmt{\begin{array}{@{}l@{}} \operatorname{is\text{-}outer\text{-}measure}(\mathbb{R}, \\ \quad \operatorname{method\text{-}i\text{-}outer}(\mathbb{R}, \operatorname{closed\text{-}intervals}, (\lambda s \in \operatorname{closed\text{-}intervals}.\; \operatorname{interval\text{-}length}(s)))) \end{array}}}
\vnbprop{lebesgue-measure-complete}{\vnbstmt{\begin{array}{@{}l@{}} \operatorname{is\text{-}complete\text{-}measure}(\mathbb{R}, \\ \quad \operatorname{caratheodory\text{-}sets}(\mathbb{R}, \\ \quad \quad \operatorname{method\text{-}i\text{-}outer}(\mathbb{R}, \operatorname{closed\text{-}intervals}, (\lambda s \in \operatorname{closed\text{-}intervals}.\; \operatorname{interval\text{-}length}(s)))), \\ \quad \operatorname{res}(\operatorname{method\text{-}i\text{-}outer}(\mathbb{R}, \operatorname{closed\text{-}intervals}, (\lambda s \in \operatorname{closed\text{-}intervals}.\; \operatorname{interval\text{-}length}(s))), \\ \quad \quad \operatorname{caratheodory\text{-}sets}(\mathbb{R}, \\ \quad \quad \quad \operatorname{method\text{-}i\text{-}outer}(\mathbb{R}, \operatorname{closed\text{-}intervals}, (\lambda s \in \operatorname{closed\text{-}intervals}.\; \operatorname{interval\text{-}length}(s)))))) \end{array}}}
\vnbprop{interval-length-ccint}{\vnbstmt{\begin{array}{@{}l@{}} \forall a, b. \\ \quad ((a \in \mathbb{R}) \wedge ((b \in \mathbb{R}) \wedge (a \leq b))) \Rightarrow  \\ \quad \quad (\lambda s \in \operatorname{closed\text{-}intervals}.\; \operatorname{interval\text{-}length}(s))(\operatorname{ccint}(a, b)) = (b - a) \end{array}}}

\vnbbookfile{structure-library/real-integral}{structure-library/real-integral}
\noindent{\small real-integral.scm -- THE INTEGRAL OF A REAL-VALUED FUNCTION: positive and negative parts, real measurability, integrability, the real integral. Definitions only; nothing here is asserted and nothing is proved.  The proofs are theorem-library/real-integral.scm (MEASURE-4-C, 2026-10-10).}\par
\par\medskip\noindent\textbf{Definition} (\texttt{pos-part}).\;$\operatorname{pos\text{-}part}(\omega, f) :=$ $(\lambda \operatorname{ppx\_} \in \omega.\; \left(\text{if }(0 \leq f(\operatorname{ppx\_}))\text{ then }f(\operatorname{ppx\_})\text{ else }0\right))$
\par\medskip\noindent\textbf{Definition} (\texttt{neg-part}).\;$\operatorname{neg\text{-}part}(\omega, f) :=$ $(\lambda \operatorname{npx\_} \in \omega.\; \left(\text{if }(0 \leq f(\operatorname{npx\_}))\text{ then }0\text{ else }-f(\operatorname{npx\_})\right))$
\par\medskip\noindent\textbf{Definition} (\texttt{abs-fn}).\;$\operatorname{abs\text{-}fn}(\omega, f) :=$ $(\lambda \operatorname{afx\_} \in \omega.\; \lvert f(\operatorname{afx\_}) \rvert)$
\par\medskip\noindent\textbf{Definition} (\texttt{is-real-measurable-fn}).\;For all $\omega$, $\mathit{ca}$ and $f$, 
\begin{equation*}
\begin{array}{@{}l@{}} \operatorname{is\text{-}real\text{-}measurable\text{-}fn}(\omega, \mathit{ca}, f) \\ \Leftrightarrow\; f \in (\omega \to \mathbb{R}) \wedge \\ \operatorname{is\text{-}measurable\text{-}fn}(\omega, \mathit{ca}, \operatorname{pos\text{-}part}(\omega, f)) \wedge \\ \operatorname{is\text{-}measurable\text{-}fn}(\omega, \mathit{ca}, \operatorname{neg\text{-}part}(\omega, f)) \end{array}
\end{equation*}

\par\medskip\noindent\textbf{Definition} (\texttt{is-integrable}).\;For all $\omega$, $\mathit{ca}$, $\mu$ and $f$, 
\begin{equation*}
\begin{array}{@{}l@{}} \operatorname{is\text{-}integrable}(\omega, \mathit{ca}, \mu, f) \\ \Leftrightarrow\; \operatorname{is\text{-}real\text{-}measurable\text{-}fn}(\omega, \mathit{ca}, f) \wedge \\ \operatorname{integral}(\omega, \mathit{ca}, \mu, \operatorname{pos\text{-}part}(\omega, f)) \in \mathbb{R} \wedge \\ \operatorname{integral}(\omega, \mathit{ca}, \mu, \operatorname{neg\text{-}part}(\omega, f)) \in \mathbb{R} \end{array}
\end{equation*}

\par\medskip\noindent\textbf{Definition} (\texttt{rintegral}).\;$\operatorname{rintegral}(\omega, \mathit{ca}, \mu, f) :=$
\begin{equation*}
(\operatorname{integral}(\omega, \mathit{ca}, \mu, \operatorname{pos\text{-}part}(\omega, f)) - \operatorname{integral}(\omega, \mathit{ca}, \mu, \operatorname{neg\text{-}part}(\omega, f)))
\end{equation*}


\vnbbookfile{theorem-library/real-integral}{theorem-library/real-integral}
\noindent{\small real-integral.scm -- THE INTEGRAL OF A REAL-VALUED FUNCTION (MEASURE-4-C, 2026-10-10).  Vocabulary: structure-library/real-integral.scm (POS-PART, NEG-PART, ABS-FN, IS-REAL-MEASURABLE-FN, IS-INTEGRABLE, RINTEGRAL).}\par
\vnbprop{pos-part-unfold}{\vnbstmt{\forall a, f.\; (\operatorname{pos\text{-}part}(a, f) \simeq (\lambda x \in a.\; \left(\text{if }(0 \leq f(x))\text{ then }f(x)\text{ else }0\right)))}}
\vnbprop{neg-part-unfold}{\vnbstmt{\forall a, f.\; (\operatorname{neg\text{-}part}(a, f) \simeq (\lambda x \in a.\; \left(\text{if }(0 \leq f(x))\text{ then }0\text{ else }-f(x)\right)))}}
\vnbprop{abs-fn-unfold}{\vnbstmt{\forall a, f.\; (\operatorname{abs\text{-}fn}(a, f) \simeq (\lambda x \in a.\; \lvert f(x) \rvert))}}
\vnbprop{rintegral-unfold}{\vnbstmt{\begin{array}{@{}l@{}} \forall a, b, u, f. \\ \quad (\operatorname{rintegral}(a, b, u, f) \simeq (\operatorname{integral}(a, b, u, \operatorname{pos\text{-}part}(a, f)) - \operatorname{integral}(a, b, u, \operatorname{neg\text{-}part}(a, f)))) \end{array}}}
\vnbprop{pos-part-value}{\vnbstmt{\forall a, f.; \forall x \in a.\; (\operatorname{pos\text{-}part}(a, f)(x) \simeq \left(\text{if }(0 \leq f(x))\text{ then }f(x)\text{ else }0\right))}}
\vnbprop{neg-part-value}{\vnbstmt{\forall a, f.; \forall x \in a.\; (\operatorname{neg\text{-}part}(a, f)(x) \simeq \left(\text{if }(0 \leq f(x))\text{ then }0\text{ else }-f(x)\right))}}
\vnbprop{abs-fn-value}{\vnbstmt{\forall a, f.; \forall x \in a.\; (\operatorname{abs\text{-}fn}(a, f)(x) \simeq \lvert f(x) \rvert)}}
\vnbprop{ptwise-eplus-value-q}{\vnbstmt{\forall a, b, c.; \forall x \in a.\; (\operatorname{ptwise\text{-}eplus}(a, b, c)(x) \simeq \operatorname{eplus}(b(x), c(x)))}}
\vnbprop{ptwise-etimes-value-q}{\vnbstmt{\forall a, b, c.; \forall x \in a.\; (\operatorname{ptwise\text{-}etimes}(a, b, c)(x) \simeq \operatorname{etimes}(b(x), c(x)))}}
\vnbprop{ptwise-scale-value-q}{\vnbstmt{\forall a, k, g.; \forall x \in a.\; (\operatorname{ptwise\text{-}scale}(a, k, g)(x) \simeq \operatorname{etimes}(k, g(x)))}}
\vnbprop{part-values-at}{\vnbstmt{\begin{array}{@{}l@{}} \forall a.; \forall f \in (a \to \mathbb{R}), x \in a. \\ \quad \operatorname{pos\text{-}part}(a, f)(x) \in \mathbb{R} \wedge \\ \quad 0 \leq \operatorname{pos\text{-}part}(a, f)(x) \wedge \\ \quad \operatorname{neg\text{-}part}(a, f)(x) \in \mathbb{R} \wedge \\ \quad 0 \leq \operatorname{neg\text{-}part}(a, f)(x) \wedge \\ \quad f(x) = (\operatorname{pos\text{-}part}(a, f)(x) - \operatorname{neg\text{-}part}(a, f)(x)) \wedge \\ \quad \lvert f(x) \rvert = (\operatorname{pos\text{-}part}(a, f)(x) + \operatorname{neg\text{-}part}(a, f)(x)) \end{array}}}
\vnbprop{pos-part-fun}{\vnbstmt{\forall a.; \forall f \in (a \to \mathbb{R}).\; (\operatorname{pos\text{-}part}(a, f) \in (a \to \operatorname{rr\text{-}pos\text{-}star}))}}
\vnbprop{neg-part-fun}{\vnbstmt{\forall a.; \forall f \in (a \to \mathbb{R}).\; (\operatorname{neg\text{-}part}(a, f) \in (a \to \operatorname{rr\text{-}pos\text{-}star}))}}
\vnbprop{abs-fn-fun}{\vnbstmt{\forall a.; \forall f \in (a \to \mathbb{R}).\; (\operatorname{abs\text{-}fn}(a, f) \in (a \to \operatorname{rr\text{-}pos\text{-}star}))}}
\vnbprop{part-values-nonneg}{\vnbstmt{\begin{array}{@{}l@{}} \forall a.; \forall f \in (a \to \mathbb{R}), x \in a. \\ \quad (0 \leq f(x)) \Rightarrow ((\operatorname{pos\text{-}part}(a, f)(x) = f(x)) \wedge (\operatorname{neg\text{-}part}(a, f)(x) = 0)) \end{array}}}
\vnbprop{part-values-neg}{\vnbstmt{\begin{array}{@{}l@{}} \forall a.; \forall f \in (a \to \mathbb{R}), x \in a. \\ \quad \neg (0 \leq f(x)) \Rightarrow ((\operatorname{pos\text{-}part}(a, f)(x) = 0) \wedge (\operatorname{neg\text{-}part}(a, f)(x) = -f(x))) \end{array}}}
\vnbprop{pos-part-measurable-of-diff}{\vnbstmt{\begin{array}{@{}l@{}} \forall a, b, f, c, d. \\ \quad f \in (a \to \mathbb{R}) \wedge \\ \quad \operatorname{is\text{-}measurable\text{-}fn}(a, b, c) \wedge \\ \quad \operatorname{is\text{-}measurable\text{-}fn}(a, b, d) \wedge \\ \quad \forall x \in a.\; (c(x) \in \mathbb{R}) \wedge \\ \quad \forall x \in a.\; (d(x) \in \mathbb{R}) \wedge \\ \quad \forall x \in a.\; (f(x) = (c(x) - d(x))) \Rightarrow \\ \quad \quad \operatorname{is\text{-}measurable\text{-}fn}(a, b, \operatorname{pos\text{-}part}(a, f)) \end{array}}}
\vnbprop{neg-part-measurable-of-diff}{\vnbstmt{\begin{array}{@{}l@{}} \forall a, b, f, c, d. \\ \quad f \in (a \to \mathbb{R}) \wedge \\ \quad \operatorname{is\text{-}measurable\text{-}fn}(a, b, c) \wedge \\ \quad \operatorname{is\text{-}measurable\text{-}fn}(a, b, d) \wedge \\ \quad \forall x \in a.\; (c(x) \in \mathbb{R}) \wedge \\ \quad \forall x \in a.\; (d(x) \in \mathbb{R}) \wedge \\ \quad \forall x \in a.\; (f(x) = (c(x) - d(x))) \Rightarrow \\ \quad \quad \operatorname{is\text{-}measurable\text{-}fn}(a, b, \operatorname{neg\text{-}part}(a, f)) \end{array}}}
\vnbprop{rintegral-decomposition}{\vnbstmt{\begin{array}{@{}l@{}} \forall a, b, u, f, c, d. \\ \quad \operatorname{is\text{-}measure}(a, b, u) \wedge \\ \quad f \in (a \to \mathbb{R}) \wedge \\ \quad \operatorname{is\text{-}measurable\text{-}fn}(a, b, c) \wedge \\ \quad \operatorname{is\text{-}measurable\text{-}fn}(a, b, d) \wedge \\ \quad \forall x \in a.\; (c(x) \in \mathbb{R}) \wedge \\ \quad \forall x \in a.\; (d(x) \in \mathbb{R}) \wedge \\ \quad \forall x \in a.\; (f(x) = (c(x) - d(x))) \wedge \\ \quad \operatorname{integral}(a, b, u, c) \in \mathbb{R} \wedge \\ \quad \operatorname{integral}(a, b, u, d) \in \mathbb{R} \Rightarrow \\ \quad \quad (\operatorname{is\text{-}integrable}(a, b, u, f) \wedge (\operatorname{rintegral}(a, b, u, f) = (\operatorname{integral}(a, b, u, c) - \operatorname{integral}(a, b, u, d)))) \end{array}}}
\vnbprop{integral-zero-fn}{\vnbstmt{\begin{array}{@{}l@{}} \forall a, b, u. \\ \quad \operatorname{is\text{-}measure}(a, b, u) \Rightarrow (\operatorname{integral}(a, b, u, (\lambda x \in a.\; 0)) = 0) \end{array}}}
\vnbprop{ptwise-scale-measurable}{\vnbstmt{\begin{array}{@{}l@{}} \forall a, b, g, k. \\ \quad (\operatorname{is\text{-}measurable\text{-}fn}(a, b, g) \wedge ((k \in \mathbb{R}) \wedge (0 \leq k))) \Rightarrow  \\ \quad \quad \operatorname{is\text{-}measurable\text{-}fn}(a, b, \operatorname{ptwise\text{-}scale}(a, k, g)) \end{array}}}
\vnbprop{abs-fn-eq-parts}{\vnbstmt{\begin{array}{@{}l@{}} \forall a.; \forall f \in (a \to \mathbb{R}). \\ \quad \operatorname{abs\text{-}fn}(a, f) = \operatorname{ptwise\text{-}eplus}(a, \operatorname{pos\text{-}part}(a, f), \operatorname{neg\text{-}part}(a, f)) \end{array}}}
\vnbprop{integral-abs-fn}{\vnbstmt{\begin{array}{@{}l@{}} \forall a, b, u, f. \\ \quad (\operatorname{is\text{-}measure}(a, b, u) \wedge \operatorname{is\text{-}real\text{-}measurable\text{-}fn}(a, b, f)) \Rightarrow  \\ \quad \quad \operatorname{is\text{-}measurable\text{-}fn}(a, b, \operatorname{abs\text{-}fn}(a, f)) \wedge \\ \quad \quad \operatorname{integral}(a, b, u, \operatorname{abs\text{-}fn}(a, f)) \\ \quad \quad =\; \operatorname{eplus}(\operatorname{integral}(a, b, u, \operatorname{pos\text{-}part}(a, f)), \operatorname{integral}(a, b, u, \operatorname{neg\text{-}part}(a, f))) \end{array}}}
\vnbprop{integrable-iff-abs-finite}{\vnbstmt{\begin{array}{@{}l@{}} \forall a, b, u, f. \\ \quad (\operatorname{is\text{-}measure}(a, b, u) \wedge \operatorname{is\text{-}real\text{-}measurable\text{-}fn}(a, b, f)) \Rightarrow  \\ \quad \quad \operatorname{is\text{-}integrable}(a, b, u, f) \Leftrightarrow (\operatorname{integral}(a, b, u, \operatorname{abs\text{-}fn}(a, f)) \in \mathbb{R}) \end{array}}}
\vnbprop{rintegral-add}{\vnbstmt{\begin{array}{@{}l@{}} \forall a, b, u, f, g. \\ \quad (\operatorname{is\text{-}measure}(a, b, u) \wedge (\operatorname{is\text{-}integrable}(a, b, u, f) \wedge \operatorname{is\text{-}integrable}(a, b, u, g))) \Rightarrow  \\ \quad \quad \operatorname{is\text{-}integrable}(a, b, u, (\lambda x \in a.\; (f(x) + g(x)))) \wedge \\ \quad \quad \operatorname{rintegral}(a, b, u, (\lambda x \in a.\; (f(x) + g(x)))) = (\operatorname{rintegral}(a, b, u, f) + \operatorname{rintegral}(a, b, u, g)) \end{array}}}
\vnbprop{rintegral-scale}{\vnbstmt{\begin{array}{@{}l@{}} \forall a, b, u, f, c. \\ \quad (\operatorname{is\text{-}measure}(a, b, u) \wedge (\operatorname{is\text{-}integrable}(a, b, u, f) \wedge (c \in \mathbb{R}))) \Rightarrow  \\ \quad \quad \operatorname{is\text{-}integrable}(a, b, u, (\lambda x \in a.\; (c \cdot f(x)))) \wedge \\ \quad \quad \operatorname{rintegral}(a, b, u, (\lambda x \in a.\; (c \cdot f(x)))) = (c \cdot \operatorname{rintegral}(a, b, u, f)) \end{array}}}
\vnbprop{rintegral-mono-off-null}{\vnbstmt{\begin{array}{@{}l@{}} \forall a, b, u, f, g, s. \\ \quad \operatorname{is\text{-}measure}(a, b, u) \wedge \\ \quad \operatorname{is\text{-}integrable}(a, b, u, f) \wedge \\ \quad \operatorname{is\text{-}integrable}(a, b, u, g) \wedge \\ \quad \operatorname{is\text{-}null\text{-}set}(a, b, u, s) \wedge \\ \quad \forall x \in a.\; (\neg (x \in s) \Rightarrow (f(x) \leq g(x))) \Rightarrow \\ \quad \quad \operatorname{rintegral}(a, b, u, f) \leq \operatorname{rintegral}(a, b, u, g) \end{array}}}
\vnbprop{rintegral-monotone}{\vnbstmt{\begin{array}{@{}l@{}} \forall a, b, u, f, g. \\ \quad \operatorname{is\text{-}measure}(a, b, u) \wedge \\ \quad \operatorname{is\text{-}integrable}(a, b, u, f) \wedge \\ \quad \operatorname{is\text{-}integrable}(a, b, u, g) \wedge \\ \quad \forall x \in a.\; (f(x) \leq g(x)) \Rightarrow \\ \quad \quad \operatorname{rintegral}(a, b, u, f) \leq \operatorname{rintegral}(a, b, u, g) \end{array}}}
\vnbprop{rintegral-abs-le}{\vnbstmt{\begin{array}{@{}l@{}} \forall a, b, u, f. \\ \quad (\operatorname{is\text{-}measure}(a, b, u) \wedge \operatorname{is\text{-}integrable}(a, b, u, f)) \Rightarrow  \\ \quad \quad \lvert \operatorname{rintegral}(a, b, u, f) \rvert \leq \operatorname{integral}(a, b, u, \operatorname{abs\text{-}fn}(a, f)) \end{array}}}
\vnbprop{rintegral-of-nonneg}{\vnbstmt{\begin{array}{@{}l@{}} \forall a, b, u, f. \\ \quad (\operatorname{is\text{-}measure}(a, b, u) \wedge (\operatorname{is\text{-}integrable}(a, b, u, f) \wedge \forall x \in a.\; (0 \leq f(x)))) \Rightarrow  \\ \quad \quad \operatorname{is\text{-}measurable\text{-}fn}(a, b, f) \wedge (\operatorname{rintegral}(a, b, u, f) = \operatorname{integral}(a, b, u, f)) \end{array}}}
\vnbprop{rintegral-ae-eq}{\vnbstmt{\begin{array}{@{}l@{}} \forall a, b, u, f, g, s. \\ \quad \operatorname{is\text{-}measure}(a, b, u) \wedge \\ \quad \operatorname{is\text{-}integrable}(a, b, u, f) \wedge \\ \quad \operatorname{is\text{-}real\text{-}measurable\text{-}fn}(a, b, g) \wedge \\ \quad \operatorname{is\text{-}null\text{-}set}(a, b, u, s) \wedge \\ \quad \forall x \in a.\; (\neg (x \in s) \Rightarrow (f(x) = g(x))) \Rightarrow \\ \quad \quad \operatorname{is\text{-}integrable}(a, b, u, g) \wedge (\operatorname{rintegral}(a, b, u, g) = \operatorname{rintegral}(a, b, u, f)) \end{array}}}
\vnbprop{rintegral-restrict}{\vnbstmt{\begin{array}{@{}l@{}} \forall a, b, u, f, c. \\ \quad (\operatorname{is\text{-}measure}(a, b, u) \wedge (\operatorname{is\text{-}integrable}(a, b, u, f) \wedge (c \in b))) \Rightarrow  \\ \quad \quad \operatorname{is\text{-}integrable}(a, b, u, (\lambda x \in a.\; (f(x) \cdot \operatorname{indicator}(a, c)(x)))) \wedge \\ \quad \quad \operatorname{rintegral}(a, b, u, (\lambda x \in a.\; (f(x) \cdot \operatorname{indicator}(a, c)(x)))) \\ \quad \quad =\; (\operatorname{integral}(a, b, u, \operatorname{ptwise\text{-}etimes}(a, \operatorname{pos\text{-}part}(a, f), \operatorname{indicator}(a, c))) - \operatorname{integral}(a, b, u, \operatorname{ptwise\text{-}etimes}(a, \operatorname{neg\text{-}part}(a, f), \operatorname{indicator}(a, c)))) \end{array}}}
\vnbprop{rintegral-nonneg-on-sets-null}{\vnbstmt{\begin{array}{@{}l@{}} \forall a, b, u, f. \\ \quad \operatorname{is\text{-}measure}(a, b, u) \wedge \\ \quad \operatorname{is\text{-}integrable}(a, b, u, f) \wedge \\ \quad \forall c \in b. \\ \quad \quad 0 \leq \operatorname{rintegral}(a, b, u, (\lambda x \in a.\; (f(x) \cdot \operatorname{indicator}(a, c)(x)))) \Rightarrow \\ \quad \quad \operatorname{is\text{-}null\text{-}set}(a, b, u, \{\mathit{nx} \in a : (f(\mathit{nx}) < 0)\}) \end{array}}}
\vnbprop{real-measurable-fn-tails}{\vnbstmt{\begin{array}{@{}l@{}} \forall a, b, f.; \forall l \in \mathbb{R}. \\ \quad \operatorname{is\text{-}real\text{-}measurable\text{-}fn}(a, b, f) \Rightarrow (\{x \in a : (l < f(x))\} \in b) \end{array}}}
\vnbprop{pos-part-measurable-of-tails}{\vnbstmt{\begin{array}{@{}l@{}} \forall a, b, f. \\ \quad (\operatorname{is\text{-}sigma\text{-}algebra}(a, b) \wedge ((f \in (a \to \mathbb{R})) \wedge \forall l \in \mathbb{R}.\; (\{x \in a : (l < f(x))\} \in b))) \Rightarrow  \\ \quad \quad \operatorname{is\text{-}measurable\text{-}fn}(a, b, \operatorname{pos\text{-}part}(a, f)) \end{array}}}
\vnbprop{real-measurable-of-tails}{\vnbstmt{\begin{array}{@{}l@{}} \forall a, b, f. \\ \quad (\operatorname{is\text{-}sigma\text{-}algebra}(a, b) \wedge ((f \in (a \to \mathbb{R})) \wedge \forall l \in \mathbb{R}.\; (\{x \in a : (l < f(x))\} \in b))) \Rightarrow  \\ \quad \quad \operatorname{is\text{-}real\text{-}measurable\text{-}fn}(a, b, f) \end{array}}}
\vnbprop{real-measurable-fn-iff-tails}{\vnbstmt{\begin{array}{@{}l@{}} \forall a, b, f. \\ \quad (\operatorname{is\text{-}sigma\text{-}algebra}(a, b) \wedge (f \in (a \to \mathbb{R}))) \Rightarrow  \\ \quad \quad \operatorname{is\text{-}real\text{-}measurable\text{-}fn}(a, b, f) \Leftrightarrow \forall l \in \mathbb{R}.\; (\{x \in a : (l < f(x))\} \in b) \end{array}}}

\vnbbookfile{structure-library/monotone-class}{structure-library/monotone-class}
\noindent{\small monotone-class.scm -- MONOTONE CLASSES.  Vocabulary only: no proof, no support. Agent MEASURE-2, 2026-10-08.  The proofs (Prop. 1.11, Thm. 1.13) are theorem-library/monotone-class.scm.}\par
\par\medskip\noindent\textbf{Definition} (\texttt{monotone-class-generated}).\;$\operatorname{monotone\text{-}class\text{-}generated}(\omega, \mathit{ce}) :=$
\begin{equation*}
\{\operatorname{mca\_} \in \operatorname{power}(\omega) : \forall \operatorname{mcb\_}.\; ((\operatorname{is\text{-}monotone\text{-}class}(\omega, \operatorname{mcb\_}) \wedge (\mathit{ce} \subseteq \operatorname{mcb\_})) \Rightarrow (\operatorname{mca\_} \in \operatorname{mcb\_}))\}
\end{equation*}

\par\medskip\noindent\textbf{Definition} (\texttt{is-monotone-class}).\;For all $\omega$ and $\mathit{mc}$, 
\begin{equation*}
\begin{array}{@{}l@{}} \operatorname{is\text{-}monotone\text{-}class}(\omega, \mathit{mc}) \\ \Leftrightarrow\; \omega \in \mathbf{Set} \wedge \\ \mathit{mc} \subseteq \operatorname{power}(\omega) \wedge \\ \forall \operatorname{mcf\_} \in (\mathbb{N} \to \mathit{mc}). \\ \quad \forall \operatorname{mcn\_} \in \mathbb{N}.\; (\operatorname{mcf\_}(\operatorname{mcn\_}) \subseteq \operatorname{mcf\_}(\operatorname{succ}(\operatorname{mcn\_}))) \Rightarrow  \\ \quad \quad \operatorname{big\text{-}union}(\operatorname{mck\_}, \mathbb{N}, \operatorname{mcf\_}(\operatorname{mck\_})) \in \mathit{mc} \wedge \\ \forall \operatorname{mcg\_} \in (\mathbb{N} \to \mathit{mc}). \\ \quad \forall \operatorname{mcm\_} \in \mathbb{N}.\; (\operatorname{mcg\_}(\operatorname{succ}(\operatorname{mcm\_})) \subseteq \operatorname{mcg\_}(\operatorname{mcm\_})) \Rightarrow  \\ \quad \quad \operatorname{complement\text{-}in}(\omega, \operatorname{big\text{-}union}(\operatorname{mcj\_}, \mathbb{N}, \operatorname{complement\text{-}in}(\omega, \operatorname{mcg\_}(\operatorname{mcj\_})))) \in \mathit{mc} \end{array}
\end{equation*}


\vnbbookfile{theorem-library/monotone-class}{theorem-library/monotone-class}
\noindent{\small monotone-class.scm -- MONOTONE CLASSES: Prop. 1.11 and the monotone class theorem (Thm. 1.13).  Agent MEASURE-2, 2026-10-08 (scratchpad/triage/MEASURE-2.md).  Library: measure.  Helper prefix mcl-. 12 theorems, every one `modulo 0', no oracle.  Vocabulary:}\par
\vnbprop{algebra-complement-closed}{\vnbstmt{\begin{array}{@{}l@{}} \forall a, b, \alpha. \\ \quad (\operatorname{is\text{-}algebra\text{-}of\text{-}sets}(a, b) \wedge (\alpha \in b)) \Rightarrow (\operatorname{complement\text{-}in}(a, \alpha) \in b) \end{array}}}
\vnbprop{algebra-union-closed}{\vnbstmt{\begin{array}{@{}l@{}} \forall a, b, \alpha, \beta. \\ \quad (\operatorname{is\text{-}algebra\text{-}of\text{-}sets}(a, b) \wedge ((\alpha \in b) \wedge (\beta \in b))) \Rightarrow  \\ \quad \quad (\alpha \cup \beta) \in b \end{array}}}
\vnbprop{algebra-inter-closed}{\vnbstmt{\begin{array}{@{}l@{}} \forall a, b, \alpha, \beta. \\ \quad (\operatorname{is\text{-}algebra\text{-}of\text{-}sets}(a, b) \wedge ((\alpha \in b) \wedge (\beta \in b))) \Rightarrow  \\ \quad \quad (\alpha \cap \beta) \in b \end{array}}}
\vnbprop{algebra-whole-empty-in}{\vnbstmt{\forall a, b.\; (\operatorname{is\text{-}algebra\text{-}of\text{-}sets}(a, b) \Rightarrow ((a \in b) \wedge (\emptyset \in b)))}}
\vnbprop{algebra-seg-union-in}{\vnbstmt{\begin{array}{@{}l@{}} \forall k \in \mathbb{N}, a, b, f. \\ \quad (\operatorname{is\text{-}algebra\text{-}of\text{-}sets}(a, b) \wedge (f \in (\mathbb{N} \to b))) \Rightarrow  \\ \quad \quad \operatorname{big\text{-}union}(j, \mathrm{OS}(k), f(j)) \in b \end{array}}}
\vnbprop{algebra-monotone-union-sigma-algebra}{\vnbstmt{\begin{array}{@{}l@{}} \forall a, b. \\ \quad \operatorname{is\text{-}algebra\text{-}of\text{-}sets}(a, b) \wedge \\ \quad \forall f \in (\mathbb{N} \to b). \\ \quad \quad \forall n \in \mathbb{N}.\; (f(n) \subseteq f((n + 1))) \Rightarrow (\operatorname{big\text{-}union}(k, \mathbb{N}, f(k)) \in b) \Rightarrow \\ \quad \quad \operatorname{is\text{-}sigma\text{-}algebra}(a, b) \end{array}}}
\vnbprop{monotone-class-generated-unfold}{\vnbstmt{\begin{array}{@{}l@{}} \forall a, e. \\ \quad (\operatorname{monotone\text{-}class\text{-}generated}(a, e) \simeq \{\mathit{mca} \in \operatorname{power}(a) : \forall b.\; ((\operatorname{is\text{-}monotone\text{-}class}(a, b) \wedge (e \subseteq b)) \Rightarrow (\mathit{mca} \in b))\}) \end{array}}}
\vnbprop{monotone-class-generated-membership}{\vnbstmt{\begin{array}{@{}l@{}} \forall a, e, \alpha. \\ \quad \alpha \in \operatorname{monotone\text{-}class\text{-}generated}(a, e) \\ \quad \Leftrightarrow\; ((\alpha \in \operatorname{power}(a)) \wedge \forall b.\; ((\operatorname{is\text{-}monotone\text{-}class}(a, b) \wedge (e \subseteq b)) \Rightarrow (\alpha \in b))) \end{array}}}
\vnbprop{monotone-class-generated-is-monotone-class}{\vnbstmt{\begin{array}{@{}l@{}} \forall a, e. \\ \quad ((a \in \mathbf{Set}) \wedge (e \subseteq \operatorname{power}(a))) \Rightarrow  \\ \quad \quad \operatorname{is\text{-}monotone\text{-}class}(a, \operatorname{monotone\text{-}class\text{-}generated}(a, e)) \end{array}}}
\vnbprop{monotone-class-union-section}{\vnbstmt{\begin{array}{@{}l@{}} \forall a, m, d. \\ \quad (\operatorname{is\text{-}monotone\text{-}class}(a, m) \wedge (d \in \operatorname{power}(a))) \Rightarrow  \\ \quad \quad \operatorname{is\text{-}monotone\text{-}class}(a, \{e \in \operatorname{power}(a) : ((d \cup e) \in m)\}) \end{array}}}
\vnbprop{monotone-class-complement-section}{\vnbstmt{\begin{array}{@{}l@{}} \forall a, m. \\ \quad \operatorname{is\text{-}monotone\text{-}class}(a, m) \Rightarrow  \\ \quad \quad \operatorname{is\text{-}monotone\text{-}class}(a, \{c \in \operatorname{power}(a) : (\operatorname{complement\text{-}in}(a, c) \in m)\}) \end{array}}}
\vnbprop{monotone-class-theorem}{\vnbstmt{\begin{array}{@{}l@{}} \forall a, c, e. \\ \quad (\operatorname{is\text{-}monotone\text{-}class}(a, c) \wedge (\operatorname{is\text{-}algebra\text{-}of\text{-}sets}(a, e) \wedge (e \subseteq c))) \Rightarrow  \\ \quad \quad \operatorname{sigma\text{-}generated}(a, e) \subseteq c \end{array}}}

\vnbbookfile{theorem-library/lebesgue-interval-length}{theorem-library/lebesgue-interval-length}
\noindent{\small lebesgue-interval-length.scm -- THE LEBESGUE OUTER MEASURE OF AN INTERVAL IS ITS LENGTH.  Agent MEASURE-3-B, 2026-10-09 (scratchpad/triage/MEASURE-3-B.md). Library: measure.  Helper prefix lil-.  No definition, no support, no stamp.}\par
\vnbprop{interval-length-unfold}{\vnbstmt{\begin{array}{@{}l@{}} \forall s. \\ \quad (\operatorname{interval\text{-}length}(s) \simeq \operatorname{esup}(\{v \in \operatorname{rr\text{-}pos\text{-}star} : ((v = 0) \vee \exists x, y \in s.\; ((x \leq y) \wedge (v = (y - x))))\})) \end{array}}}
\vnbprop{interval-length-type}{\vnbstmt{\forall s.\; (\operatorname{interval\text{-}length}(s) \in \operatorname{rr\text{-}pos\text{-}star})}}
\vnbprop{interval-length-mono}{\vnbstmt{\forall a, b.\; ((a \subseteq b) \Rightarrow (\operatorname{interval\text{-}length}(a) \leq \operatorname{interval\text{-}length}(b)))}}
\vnbprop{interval-length-extend}{\vnbstmt{\begin{array}{@{}l@{}} \forall a, b, w, z. \\ \quad b \subseteq \mathbb{R} \wedge \\ \quad a \subseteq b \wedge \\ \quad w \in \mathbb{R} \wedge \\ \quad z \in \mathbb{R} \wedge \\ \quad w \leq z \wedge \\ \quad w \in b \wedge \\ \quad z \in b \wedge \\ \quad \forall x \in a.\; (x \leq w) \Rightarrow \\ \quad \quad \operatorname{eplus}(\operatorname{interval\text{-}length}(a), (z - w)) \leq \operatorname{interval\text{-}length}(b) \end{array}}}
\vnbprop{interval-length-widen}{\vnbstmt{\begin{array}{@{}l@{}} \forall s, r. \\ \quad ((s \subseteq \mathbb{R}) \wedge ((r \in \mathbb{R}) \wedge (0 < r))) \Rightarrow  \\ \quad \quad \operatorname{interval\text{-}length}(\{z \in \mathbb{R} : \exists u \in s.\; (((z - u) < r) \wedge ((u - z) < r))\}) \\ \quad \quad \leq\; \operatorname{eplus}(\operatorname{interval\text{-}length}(s), \operatorname{eplus}(r, r)) \end{array}}}
\vnbprop{rps-prefix-mono-range}{\vnbstmt{\begin{array}{@{}l@{}} \forall f, m, n. \\ \quad ((f \in (\mathbb{N} \to \operatorname{rr\text{-}pos\text{-}star})) \wedge ((m \in \mathbb{N}) \wedge ((n \in \mathbb{N}) \wedge (m \leq n)))) \Rightarrow  \\ \quad \quad \operatorname{finsum}(\operatorname{rr\text{-}pos\text{-}star\text{-}add\text{-}monoid}, f, \mathrm{OS}(m)) \\ \quad \quad \leq\; \operatorname{finsum}(\operatorname{rr\text{-}pos\text{-}star\text{-}add\text{-}monoid}, f, \mathrm{OS}(n)) \end{array}}}
\vnbprop{rps-prefix-bump}{\vnbstmt{\begin{array}{@{}l@{}} \forall n \in \mathbb{N}, k, f, g, c. \\ \quad k \in \mathbb{N} \wedge \\ \quad f \in (\mathbb{N} \to \operatorname{rr\text{-}pos\text{-}star}) \wedge \\ \quad g \in (\mathbb{N} \to \operatorname{rr\text{-}pos\text{-}star}) \wedge \\ \quad c \in \operatorname{rr\text{-}pos\text{-}star} \wedge \\ \quad \forall q \in \mathbb{N}.\; (f(q) \leq g(q)) \wedge \\ \quad \operatorname{eplus}(f(k), c) \leq g(k) \wedge \\ \quad k < n \Rightarrow \\ \quad \quad \operatorname{eplus}(\operatorname{finsum}(\operatorname{rr\text{-}pos\text{-}star\text{-}add\text{-}monoid}, f, \mathrm{OS}(n)), c) \\ \quad \quad \leq\; \operatorname{finsum}(\operatorname{rr\text{-}pos\text{-}star\text{-}add\text{-}monoid}, g, \mathrm{OS}(n)) \end{array}}}
\vnbprop{rr-near-center-half}{\vnbstmt{\begin{array}{@{}l@{}} \forall t, d, r, w. \\ \quad ((t \in \mathbb{R}) \wedge ((d \in \mathbb{R}) \wedge ((r \in \mathbb{R}) \wedge ((w \in \mathbb{R}) \wedge ((0 < d) \wedge (((d + d) = r) \wedge (((t - d) < w) \wedge (w \leq (t + d))))))))) \Rightarrow  \\ \quad \quad ((w - t) < r) \wedge ((t - w) < r) \end{array}}}
\vnbprop{ccint-cover-widened-sum}{\vnbstmt{\begin{array}{@{}l@{}} \forall a, b, s, e. \\ \quad a \in \mathbb{R} \wedge \\ \quad b \in \mathbb{R} \wedge \\ \quad a \leq b \wedge \\ \quad s \in (\mathbb{N} \to \operatorname{power}(\mathbb{R})) \wedge \\ \quad \operatorname{ccint}(a, b) \subseteq \operatorname{big\text{-}union}(n, \mathbb{N}, s(n)) \wedge \\ \quad \operatorname{pos\text{-}rr}(e) \Rightarrow \\ \quad \quad (b - a) \\ \quad \quad \leq\; \operatorname{esum}((\lambda \mathit{lnn} \in \mathbb{N}.\; \\ \quad \quad \quad \quad \operatorname{interval\text{-}length}(\{z \in \mathbb{R} : \exists u \in s(\mathit{lnn}).\; (((z - u) < (e \cdot \operatorname{recip}(\operatorname{power}(2, (\mathit{lnn} + 1))))) \wedge ((u - z) < (e \cdot \operatorname{recip}(\operatorname{power}(2, (\mathit{lnn} + 1))))))\}))) \end{array}}}
\vnbprop{ccint-cover-sum-ge}{\vnbstmt{\begin{array}{@{}l@{}} \forall a, b, s. \\ \quad a \in \mathbb{R} \wedge \\ \quad b \in \mathbb{R} \wedge \\ \quad a \leq b \wedge \\ \quad s \in (\mathbb{N} \to \operatorname{power}(\mathbb{R})) \wedge \\ \quad \operatorname{ccint}(a, b) \subseteq \operatorname{big\text{-}union}(n, \mathbb{N}, s(n)) \Rightarrow \\ \quad \quad (b - a) \leq \operatorname{esum}((\lambda \mathit{lkn} \in \mathbb{N}.\; \operatorname{interval\text{-}length}(s(\mathit{lkn})))) \end{array}}}
\vnbprop{subset-rr-in-power-rr}{\vnbstmt{\forall s.\; ((s \subseteq \mathbb{R}) \Rightarrow (s \in \operatorname{power}(\mathbb{R})))}}
\vnbprop{lebesgue-outer-mono}{\vnbstmt{\begin{array}{@{}l@{}} \forall a, b. \\ \quad ((a \subseteq \mathbb{R}) \wedge ((b \subseteq \mathbb{R}) \wedge (a \subseteq b))) \Rightarrow  \\ \quad \quad \operatorname{method\text{-}i\text{-}outer}(\mathbb{R}, \operatorname{closed\text{-}intervals}, (\lambda s \in \operatorname{closed\text{-}intervals}.\; \operatorname{interval\text{-}length}(s)))(a) \\ \quad \quad \leq\; \operatorname{method\text{-}i\text{-}outer}(\mathbb{R}, \operatorname{closed\text{-}intervals}, (\lambda s \in \operatorname{closed\text{-}intervals}.\; \operatorname{interval\text{-}length}(s)))(b) \end{array}}}
\vnbprop{lebesgue-outer-ccint}{\vnbstmt{\begin{array}{@{}l@{}} \forall a, b. \\ \quad ((a \in \mathbb{R}) \wedge ((b \in \mathbb{R}) \wedge (a \leq b))) \Rightarrow  \\ \quad \quad \operatorname{method\text{-}i\text{-}outer}(\mathbb{R}, \operatorname{closed\text{-}intervals}, (\lambda s \in \operatorname{closed\text{-}intervals}.\; \operatorname{interval\text{-}length}(s)))(\operatorname{ccint}(a, b)) \\ \quad \quad =\; (b - a) \end{array}}}
\vnbprop{lebesgue-outer-point}{\vnbstmt{\begin{array}{@{}l@{}} \forall x \in \mathbb{R}. \\ \quad \operatorname{method\text{-}i\text{-}outer}(\mathbb{R}, \operatorname{closed\text{-}intervals}, (\lambda s \in \operatorname{closed\text{-}intervals}.\; \operatorname{interval\text{-}length}(s)))(\operatorname{singleton}(x)) \\ \quad =\; 0 \end{array}}}
\vnbprop{lebesgue-outer-empty}{\vnbstmt{\begin{array}{@{}l@{}} \operatorname{method\text{-}i\text{-}outer}(\mathbb{R}, \operatorname{closed\text{-}intervals}, (\lambda s \in \operatorname{closed\text{-}intervals}.\; \operatorname{interval\text{-}length}(s)))(\emptyset) \\ =\; 0 \end{array}}}
\vnbprop{lebesgue-outer-between}{\vnbstmt{\begin{array}{@{}l@{}} \forall a, b, s. \\ \quad a \in \mathbb{R} \wedge \\ \quad b \in \mathbb{R} \wedge \\ \quad a \leq b \wedge \\ \quad s \subseteq \operatorname{ccint}(a, b) \wedge \\ \quad \operatorname{ccint}(a, b) \subseteq ((s \cup \operatorname{singleton}(a)) \cup \operatorname{singleton}(b)) \Rightarrow \\ \quad \quad \operatorname{method\text{-}i\text{-}outer}(\mathbb{R}, \operatorname{closed\text{-}intervals}, (\lambda \mathit{ols} \in \operatorname{closed\text{-}intervals}.\; \operatorname{interval\text{-}length}(\mathit{ols})))(s) \\ \quad \quad =\; (b - a) \end{array}}}
\vnbprop{lebesgue-outer-ooint}{\vnbstmt{\begin{array}{@{}l@{}} \forall a, b. \\ \quad ((a \in \mathbb{R}) \wedge ((b \in \mathbb{R}) \wedge (a \leq b))) \Rightarrow  \\ \quad \quad \operatorname{method\text{-}i\text{-}outer}(\mathbb{R}, \operatorname{closed\text{-}intervals}, (\lambda s \in \operatorname{closed\text{-}intervals}.\; \operatorname{interval\text{-}length}(s)))(\operatorname{ooint}(a, b)) \\ \quad \quad =\; (b - a) \end{array}}}
\vnbprop{lebesgue-outer-coint}{\vnbstmt{\begin{array}{@{}l@{}} \forall a, b. \\ \quad ((a \in \mathbb{R}) \wedge ((b \in \mathbb{R}) \wedge (a \leq b))) \Rightarrow  \\ \quad \quad \operatorname{method\text{-}i\text{-}outer}(\mathbb{R}, \operatorname{closed\text{-}intervals}, (\lambda s \in \operatorname{closed\text{-}intervals}.\; \operatorname{interval\text{-}length}(s)))(\operatorname{coint}(a, b)) \\ \quad \quad =\; (b - a) \end{array}}}
\vnbprop{lebesgue-outer-ocint}{\vnbstmt{\begin{array}{@{}l@{}} \forall a, b. \\ \quad ((a \in \mathbb{R}) \wedge ((b \in \mathbb{R}) \wedge (a \leq b))) \Rightarrow  \\ \quad \quad \operatorname{method\text{-}i\text{-}outer}(\mathbb{R}, \operatorname{closed\text{-}intervals}, (\lambda s \in \operatorname{closed\text{-}intervals}.\; \operatorname{interval\text{-}length}(s)))(\operatorname{ocint}(a, b)) \\ \quad \quad =\; (b - a) \end{array}}}
\vnbprop{interval-length-translate-le}{\vnbstmt{\begin{array}{@{}l@{}} \forall s.; \forall c \in \mathbb{R}. \\ \quad \operatorname{interval\text{-}length}(\{z \in \mathbb{R} : ((z - c) \in s)\}) \leq \operatorname{interval\text{-}length}(s) \end{array}}}
\vnbprop{translate-closed-interval}{\vnbstmt{\begin{array}{@{}l@{}} \forall i, c. \\ \quad ((i \in \operatorname{closed\text{-}intervals}) \wedge (c \in \mathbb{R})) \Rightarrow  \\ \quad \quad \{z \in \mathbb{R} : ((z - c) \in i)\} \in \operatorname{closed\text{-}intervals} \end{array}}}
\vnbprop{lebesgue-outer-translate-le}{\vnbstmt{\begin{array}{@{}l@{}} \forall s, c. \\ \quad ((s \subseteq \mathbb{R}) \wedge (c \in \mathbb{R})) \Rightarrow  \\ \quad \quad \operatorname{method\text{-}i\text{-}outer}(\mathbb{R}, \operatorname{closed\text{-}intervals}, (\lambda \mathit{ols} \in \operatorname{closed\text{-}intervals}.\; \operatorname{interval\text{-}length}(\mathit{ols})))(\{z \in \mathbb{R} : ((z - c) \in s)\}) \\ \quad \quad \leq\; \operatorname{method\text{-}i\text{-}outer}(\mathbb{R}, \operatorname{closed\text{-}intervals}, (\lambda \mathit{ols} \in \operatorname{closed\text{-}intervals}.\; \operatorname{interval\text{-}length}(\mathit{ols})))(s) \end{array}}}
\vnbprop{lebesgue-outer-translate}{\vnbstmt{\begin{array}{@{}l@{}} \forall s, c. \\ \quad ((s \subseteq \mathbb{R}) \wedge (c \in \mathbb{R})) \Rightarrow  \\ \quad \quad \operatorname{method\text{-}i\text{-}outer}(\mathbb{R}, \operatorname{closed\text{-}intervals}, (\lambda \mathit{ols} \in \operatorname{closed\text{-}intervals}.\; \operatorname{interval\text{-}length}(\mathit{ols})))(\{z \in \mathbb{R} : ((z - c) \in s)\}) \\ \quad \quad =\; \operatorname{method\text{-}i\text{-}outer}(\mathbb{R}, \operatorname{closed\text{-}intervals}, (\lambda \mathit{ols} \in \operatorname{closed\text{-}intervals}.\; \operatorname{interval\text{-}length}(\mathit{ols})))(s) \end{array}}}

\vnbbookfile{structure-library/metric-outer-measure}{structure-library/metric-outer-measure}
\noindent{\small metric-outer-measure.scm -- the distance between sets, METRIC OUTER MEASURES and the BOREL SETS of a metric space.  Vocabulary only: no proof, no support. Agent MEASURE-3-A, 2026-10-09 (scratchpad/triage/MEASURE-3-A.md).  The proofs are theorem-library/metric-outer-measure.scm.}\par
\par\medskip\noindent\textbf{Definition} (\texttt{pt-set-dist}).\;$\operatorname{pt\text{-}set\text{-}dist}(m, x, \mathit{es}) :=$
\begin{equation*}
\operatorname{einf}(\{\operatorname{mov\_} \in \operatorname{rr\text{-}pos\text{-}star} : \exists \operatorname{moz\_} \in \mathit{es}.\; (\operatorname{mov\_} = \operatorname{dist}(m)(x, \operatorname{moz\_}))\})
\end{equation*}

\par\medskip\noindent\textbf{Definition} (\texttt{set-dist}).\;$\operatorname{set\text{-}dist}(m, \mathit{aa}, \mathit{bb}) :=$
\begin{equation*}
\operatorname{einf}(\{\operatorname{mow\_} \in \operatorname{rr\text{-}pos\text{-}star} : \exists \operatorname{mox\_} \in \mathit{aa}, \operatorname{moy\_} \in \mathit{bb}.\; (\operatorname{mow\_} = \operatorname{dist}(m)(\operatorname{mox\_}, \operatorname{moy\_}))\})
\end{equation*}

\par\medskip\noindent\textbf{Definition} (\texttt{is-metric-outer-measure}).\;For all $m$ and $\mu$, 
\begin{equation*}
\begin{array}{@{}l@{}} \operatorname{is\text{-}metric\text{-}outer\text{-}measure}(m, \mu) \\ \Leftrightarrow\; \operatorname{is\text{-}metric\text{-}space}(m) \wedge \\ \operatorname{is\text{-}outer\text{-}measure}(\operatorname{pts}(m), \mu) \wedge \\ \forall \operatorname{moa\_}, \operatorname{mob\_} \in \operatorname{power}(\operatorname{pts}(m)). \\ \quad \neg (\operatorname{set\text{-}dist}(m, \operatorname{moa\_}, \operatorname{mob\_}) = 0) \Rightarrow  \\ \quad \quad \mu((\operatorname{moa\_} \cup \operatorname{mob\_})) = \operatorname{eplus}(\mu(\operatorname{moa\_}), \mu(\operatorname{mob\_})) \end{array}
\end{equation*}

\par\medskip\noindent\textbf{Definition} (\texttt{borel-sets}).\;$\operatorname{borel\text{-}sets}(m) :=$
\begin{equation*}
\operatorname{sigma\text{-}generated}(\operatorname{pts}(m), \{\operatorname{mou\_} \in \operatorname{power}(\operatorname{pts}(m)) : \operatorname{is\text{-}open}(m, \operatorname{mou\_})\})
\end{equation*}


\vnbbookfile{theorem-library/metric-outer-measure}{theorem-library/metric-outer-measure}
\noindent{\small metric-outer-measure.scm -- METRIC OUTER MEASURES: separated sets, the layers of a closed set, Borel sets are measurable; Lebesgue outer measure on the line is metric. Agent MEASURE-3-A, 2026-10-09 (scratchpad/triage/MEASURE-3-A.md).  Library: measure. Helper prefix mom-.  58 theorems, every one `modulo 0' (oracles ...}\par
\vnbprop{pt-set-dist-unfold}{\vnbstmt{\begin{array}{@{}l@{}} \forall m, x, s. \\ \quad (\operatorname{pt\text{-}set\text{-}dist}(m, x, s) \simeq \operatorname{einf}(\{v \in \operatorname{rr\text{-}pos\text{-}star} : \exists z \in s.\; (v = \operatorname{dist}(m)(x, z))\})) \end{array}}}
\vnbprop{set-dist-unfold}{\vnbstmt{\begin{array}{@{}l@{}} \forall m, a, b. \\ \quad (\operatorname{set\text{-}dist}(m, a, b) \simeq \operatorname{einf}(\{w \in \operatorname{rr\text{-}pos\text{-}star} : \exists x \in a, y \in b.\; (w = \operatorname{dist}(m)(x, y))\})) \end{array}}}
\vnbprop{pt-set-dist-ge-intro}{\vnbstmt{\begin{array}{@{}l@{}} \forall m, x, s.; \forall r \in \operatorname{rr\text{-}pos\text{-}star}. \\ \quad \forall z \in s.\; (r \leq \operatorname{dist}(m)(x, z)) \Rightarrow (r \leq \operatorname{pt\text{-}set\text{-}dist}(m, x, s)) \end{array}}}
\vnbprop{pt-set-dist-ge-elim}{\vnbstmt{\begin{array}{@{}l@{}} \forall m, x, s, r.; \forall z \in s. \\ \quad \operatorname{is\text{-}metric\text{-}space}(m) \wedge \\ \quad x \in \operatorname{pts}(m) \wedge \\ \quad s \subseteq \operatorname{pts}(m) \wedge \\ \quad r \in \operatorname{rr\text{-}pos\text{-}star} \wedge \\ \quad r \leq \operatorname{pt\text{-}set\text{-}dist}(m, x, s) \Rightarrow \\ \quad \quad r \leq \operatorname{dist}(m)(x, z) \end{array}}}
\vnbprop{set-dist-pos-intro}{\vnbstmt{\begin{array}{@{}l@{}} \forall m, a, b, r. \\ \quad (\operatorname{pos\text{-}rr}(r) \wedge \forall x \in a, y \in b.\; (r \leq \operatorname{dist}(m)(x, y))) \Rightarrow  \\ \quad \quad \neg (\operatorname{set\text{-}dist}(m, a, b) = 0) \end{array}}}
\vnbprop{set-dist-pos-elim}{\vnbstmt{\begin{array}{@{}l@{}} \forall m, a, b. \\ \quad (\operatorname{is\text{-}metric\text{-}space}(m) \wedge ((a \subseteq \operatorname{pts}(m)) \wedge ((b \subseteq \operatorname{pts}(m)) \wedge \neg (\operatorname{set\text{-}dist}(m, a, b) = 0)))) \Rightarrow  \\ \quad \quad \exists r.\; (\operatorname{pos\text{-}rr}(r) \wedge \forall x \in a, y \in b.\; (r \leq \operatorname{dist}(m)(x, y))) \end{array}}}
\vnbprop{set-dist-union-pos}{\vnbstmt{\begin{array}{@{}l@{}} \forall m, a, b, s. \\ \quad \operatorname{is\text{-}metric\text{-}space}(m) \wedge \\ \quad a \subseteq \operatorname{pts}(m) \wedge \\ \quad b \subseteq \operatorname{pts}(m) \wedge \\ \quad s \subseteq \operatorname{pts}(m) \wedge \\ \quad \neg (\operatorname{set\text{-}dist}(m, a, s) = 0) \wedge \\ \quad \neg (\operatorname{set\text{-}dist}(m, b, s) = 0) \Rightarrow \\ \quad \quad \neg (\operatorname{set\text{-}dist}(m, (a \cup b), s) = 0) \end{array}}}
\vnbprop{metric-outer-measure-prefix-separated}{\vnbstmt{\begin{array}{@{}l@{}} \forall n \in \mathbb{N}, m, u, s. \\ \quad \operatorname{is\text{-}metric\text{-}outer\text{-}measure}(m, u) \wedge \\ \quad s \in (\mathbb{N} \to \operatorname{power}(\operatorname{pts}(m))) \wedge \\ \quad \forall k, l \in \mathbb{N}.\; (\neg (k = l) \Rightarrow \neg (\operatorname{set\text{-}dist}(m, s(k), s(l)) = 0)) \Rightarrow \\ \quad \quad u(\operatorname{big\text{-}union}(j, \mathrm{OS}(n), s(j))) \\ \quad \quad =\; \operatorname{finsum}(\operatorname{rr\text{-}pos\text{-}star\text{-}add\text{-}monoid}, (\lambda a \in \mathbb{N}.\; u(s(a))), \mathrm{OS}(n)) \wedge \\ \quad \quad \forall \mathit{moj} \in \mathbb{N}. \\ \quad \quad \quad (n \leq \mathit{moj}) \Rightarrow  \\ \quad \quad \quad \quad \neg (\operatorname{set\text{-}dist}(m, \operatorname{big\text{-}union}(j, \mathrm{OS}(n), s(j)), s(\mathit{moj})) = 0) \end{array}}}
\vnbprop{metric-outer-measure-finite-separated}{\vnbstmt{\begin{array}{@{}l@{}} \forall m, u, s.; \forall n \in \mathbb{N}. \\ \quad \operatorname{is\text{-}metric\text{-}outer\text{-}measure}(m, u) \wedge \\ \quad s \in (\mathbb{N} \to \operatorname{power}(\operatorname{pts}(m))) \wedge \\ \quad \forall k, l \in \mathbb{N}.\; (\neg (k = l) \Rightarrow \neg (\operatorname{set\text{-}dist}(m, s(k), s(l)) = 0)) \Rightarrow \\ \quad \quad u(\operatorname{big\text{-}union}(j, \mathrm{OS}(n), s(j))) \\ \quad \quad =\; \operatorname{finsum}(\operatorname{rr\text{-}pos\text{-}star\text{-}add\text{-}monoid}, (\lambda a \in \mathbb{N}.\; u(s(a))), \mathrm{OS}(n)) \end{array}}}
\vnbprop{metric-outer-measure-separated-sum}{\vnbstmt{\begin{array}{@{}l@{}} \forall m, u, s. \\ \quad \operatorname{is\text{-}metric\text{-}outer\text{-}measure}(m, u) \wedge \\ \quad s \in (\mathbb{N} \to \operatorname{power}(\operatorname{pts}(m))) \wedge \\ \quad \forall k, l \in \mathbb{N}.\; (\neg (k = l) \Rightarrow \neg (\operatorname{set\text{-}dist}(m, s(k), s(l)) = 0)) \Rightarrow \\ \quad \quad u(\operatorname{big\text{-}union}(n, \mathbb{N}, s(n))) = \operatorname{esum}((\lambda n \in \mathbb{N}.\; u(s(n)))) \end{array}}}
\vnbprop{rr-mul-succ-distrib}{\vnbstmt{\forall a, c \in \mathbb{R}.\; ((c \cdot (a + 1)) = ((a \cdot c) + c))}}
\vnbprop{nn-recip-succ-gap}{\vnbstmt{\forall j \in \mathbb{N}.\; \operatorname{pos\text{-}rr}((\operatorname{recip}((j + 1)) - \operatorname{recip}(((j + 1) + 1))))}}
\vnbprop{dist-layer-mem-elim}{\vnbstmt{\begin{array}{@{}l@{}} \forall m, l.; \forall n \in \mathbb{N}, x \in \{\mathit{mex} \in \operatorname{pts}(m) : (\operatorname{recip}((n + 1)) \leq \operatorname{pt\text{-}set\text{-}dist}(m, \mathit{mex}, l))\}. \\ \quad (\operatorname{is\text{-}metric\text{-}space}(m) \wedge (l \subseteq \operatorname{pts}(m))) \Rightarrow  \\ \quad \quad (x \in \operatorname{pts}(m)) \wedge \forall z \in l.\; (\operatorname{recip}((n + 1)) \leq \operatorname{dist}(m)(x, z)) \end{array}}}
\vnbprop{dist-layer-mem-intro}{\vnbstmt{\begin{array}{@{}l@{}} \forall m, l.; \forall n \in \mathbb{N}, x \in \operatorname{pts}(m). \\ \quad \forall z \in l.\; (\operatorname{recip}((n + 1)) \leq \operatorname{dist}(m)(x, z)) \Rightarrow  \\ \quad \quad x \in \{\mathit{mex} \in \operatorname{pts}(m) : (\operatorname{recip}((n + 1)) \leq \operatorname{pt\text{-}set\text{-}dist}(m, \mathit{mex}, l))\} \end{array}}}
\vnbprop{dist-layer-mono}{\vnbstmt{\begin{array}{@{}l@{}} \forall m, l, k, n. \\ \quad (\operatorname{is\text{-}metric\text{-}space}(m) \wedge ((l \subseteq \operatorname{pts}(m)) \wedge ((k \in \mathbb{N}) \wedge ((n \in \mathbb{N}) \wedge (k \leq n))))) \Rightarrow  \\ \quad \quad \{x \in \operatorname{pts}(m) : (\operatorname{recip}((k + 1)) \leq \operatorname{pt\text{-}set\text{-}dist}(m, x, l))\} \\ \quad \quad \subseteq\; \{x \in \operatorname{pts}(m) : (\operatorname{recip}((n + 1)) \leq \operatorname{pt\text{-}set\text{-}dist}(m, x, l))\} \end{array}}}
\vnbprop{dist-layer-disjoint}{\vnbstmt{\begin{array}{@{}l@{}} \forall m, l.; \forall n \in \mathbb{N}, x \in \{\mathit{mex} \in \operatorname{pts}(m) : (\operatorname{recip}((n + 1)) \leq \operatorname{pt\text{-}set\text{-}dist}(m, \mathit{mex}, l))\}. \\ \quad (\operatorname{is\text{-}metric\text{-}space}(m) \wedge (l \subseteq \operatorname{pts}(m))) \Rightarrow \neg (x \in l) \end{array}}}
\vnbprop{dist-layer-cover}{\vnbstmt{\begin{array}{@{}l@{}} \forall m, l.; \forall x \in \operatorname{pts}(m). \\ \quad (\operatorname{is\text{-}closed}(m, l) \wedge \neg (x \in l)) \Rightarrow  \\ \quad \quad \exists k \in \mathbb{N}.\; (x \in \{\mathit{mex} \in \operatorname{pts}(m) : (\operatorname{recip}((k + 1)) \leq \operatorname{pt\text{-}set\text{-}dist}(m, \mathit{mex}, l))\}) \end{array}}}
\vnbprop{dist-layer-gap}{\vnbstmt{\begin{array}{@{}l@{}} \forall m, l.; \forall j \in \mathbb{N}, x, y. \\ \quad \operatorname{is\text{-}metric\text{-}space}(m) \wedge \\ \quad l \subseteq \operatorname{pts}(m) \wedge \\ \quad x \in \{\mathit{mex} \in \operatorname{pts}(m) : (\operatorname{recip}((j + 1)) \leq \operatorname{pt\text{-}set\text{-}dist}(m, \mathit{mex}, l))\} \wedge \\ \quad y \in \operatorname{pts}(m) \wedge \\ \quad \neg (y \in \{\mathit{mex} \in \operatorname{pts}(m) : (\operatorname{recip}(((j + 1) + 1)) \leq \operatorname{pt\text{-}set\text{-}dist}(m, \mathit{mex}, l))\}) \Rightarrow \\ \quad \quad (\operatorname{recip}((j + 1)) - \operatorname{recip}(((j + 1) + 1))) \leq \operatorname{dist}(m)(x, y) \end{array}}}
\vnbprop{rps-prefix-mono}{\vnbstmt{\begin{array}{@{}l@{}} \forall k \in \mathbb{N}, f \in (\mathbb{N} \to \operatorname{rr\text{-}pos\text{-}star}), n \in \mathbb{N}. \\ \quad (n \leq k) \Rightarrow  \\ \quad \quad \operatorname{finsum}(\operatorname{rr\text{-}pos\text{-}star\text{-}add\text{-}monoid}, f, \mathrm{OS}(n)) \\ \quad \quad \leq\; \operatorname{finsum}(\operatorname{rr\text{-}pos\text{-}star\text{-}add\text{-}monoid}, f, \mathrm{OS}(k)) \end{array}}}
\vnbprop{rps-prefix-even-odd}{\vnbstmt{\begin{array}{@{}l@{}} \forall n \in \mathbb{N}, f, g, h. \\ \quad f \in (\mathbb{N} \to \operatorname{rr\text{-}pos\text{-}star}) \wedge \\ \quad g \in (\mathbb{N} \to \operatorname{rr\text{-}pos\text{-}star}) \wedge \\ \quad h \in (\mathbb{N} \to \operatorname{rr\text{-}pos\text{-}star}) \wedge \\ \quad \forall q \in \mathbb{N}.\; (g(q) = f((2 \cdot q))) \wedge \\ \quad \forall q \in \mathbb{N}.\; (h(q) = f(((2 \cdot q) + 1))) \Rightarrow \\ \quad \quad \operatorname{finsum}(\operatorname{rr\text{-}pos\text{-}star\text{-}add\text{-}monoid}, f, \mathrm{OS}((2 \cdot n))) \\ \quad \quad =\; \operatorname{eplus}(\operatorname{finsum}(\operatorname{rr\text{-}pos\text{-}star\text{-}add\text{-}monoid}, g, \mathrm{OS}(n)), \\ \quad \quad \quad \operatorname{finsum}(\operatorname{rr\text{-}pos\text{-}star\text{-}add\text{-}monoid}, h, \mathrm{OS}(n))) \end{array}}}
\vnbprop{esum-nn-even-odd-le}{\vnbstmt{\begin{array}{@{}l@{}} \forall f, g, h. \\ \quad f \in (\mathbb{N} \to \operatorname{rr\text{-}pos\text{-}star}) \wedge \\ \quad g \in (\mathbb{N} \to \operatorname{rr\text{-}pos\text{-}star}) \wedge \\ \quad h \in (\mathbb{N} \to \operatorname{rr\text{-}pos\text{-}star}) \wedge \\ \quad \forall q \in \mathbb{N}.\; (g(q) = f((2 \cdot q))) \wedge \\ \quad \forall q \in \mathbb{N}.\; (h(q) = f(((2 \cdot q) + 1))) \Rightarrow \\ \quad \quad \operatorname{esum}(f) \leq \operatorname{eplus}(\operatorname{esum}(g), \operatorname{esum}(h)) \end{array}}}
\vnbprop{rps-prefix-shift-le}{\vnbstmt{\begin{array}{@{}l@{}} \forall k, n \in \mathbb{N}, f, g. \\ \quad ((f \in (\mathbb{N} \to \operatorname{rr\text{-}pos\text{-}star})) \wedge ((g \in (\mathbb{N} \to \operatorname{rr\text{-}pos\text{-}star})) \wedge \forall j \in \mathbb{N}.\; (g(j) \leq f((n + j))))) \Rightarrow  \\ \quad \quad \operatorname{eplus}(\operatorname{finsum}(\operatorname{rr\text{-}pos\text{-}star\text{-}add\text{-}monoid}, f, \mathrm{OS}(n)), \\ \quad \quad \quad \operatorname{finsum}(\operatorname{rr\text{-}pos\text{-}star\text{-}add\text{-}monoid}, g, \mathrm{OS}(k))) \\ \quad \quad \leq\; \operatorname{finsum}(\operatorname{rr\text{-}pos\text{-}star\text{-}add\text{-}monoid}, f, \mathrm{OS}((n + k))) \end{array}}}
\vnbprop{esum-nn-tail-small}{\vnbstmt{\begin{array}{@{}l@{}} \forall f \in (\mathbb{N} \to \operatorname{rr\text{-}pos\text{-}star}), s. \\ \quad ((\operatorname{esum}(f) \in \mathbb{R}) \wedge \operatorname{pos\text{-}rr}(s)) \Rightarrow  \\ \quad \quad \exists n \in \mathbb{N}. \forall g \in (\mathbb{N} \to \operatorname{rr\text{-}pos\text{-}star}). \\ \quad \quad \quad \forall j \in \mathbb{N}.\; (g(j) \leq f((n + j))) \Rightarrow (\operatorname{esum}(g) \leq s) \end{array}}}
\vnbprop{dist-layer-ring-cover}{\vnbstmt{\begin{array}{@{}l@{}} \forall m, l.; \forall n \in \mathbb{N}, x \in \operatorname{pts}(m). \\ \quad \operatorname{is\text{-}closed}(m, l) \wedge \\ \quad \neg (x \in l) \wedge \\ \quad \neg (x \in \{\mathit{mex} \in \operatorname{pts}(m) : (\operatorname{recip}((n + 1)) \leq \operatorname{pt\text{-}set\text{-}dist}(m, \mathit{mex}, l))\}) \Rightarrow \\ \quad \quad \exists j \in \mathbb{N}. \\ \quad \quad \quad x \\ \quad \quad \quad \in\; \operatorname{complement\text{-}in}(\{\mathit{mex} \in \operatorname{pts}(m) : (\operatorname{recip}((((n + j) + 1) + 1)) \leq \operatorname{pt\text{-}set\text{-}dist}(m, \mathit{mex}, l))\}, \\ \quad \quad \quad \quad \{\mathit{mex} \in \operatorname{pts}(m) : (\operatorname{recip}(((n + j) + 1)) \leq \operatorname{pt\text{-}set\text{-}dist}(m, \mathit{mex}, l))\}) \end{array}}}
\vnbprop{dist-layer-set-sep}{\vnbstmt{\begin{array}{@{}l@{}} \forall m, l.; \forall j \in \mathbb{N}, a, b. \\ \quad \operatorname{is\text{-}metric\text{-}space}(m) \wedge \\ \quad l \subseteq \operatorname{pts}(m) \wedge \\ \quad a \subseteq \{x \in \operatorname{pts}(m) : (\operatorname{recip}((j + 1)) \leq \operatorname{pt\text{-}set\text{-}dist}(m, x, l))\} \wedge \\ \quad \forall y \in b. \\ \quad \quad ((y \in \operatorname{pts}(m)) \wedge \neg (y \in \{x \in \operatorname{pts}(m) : (\operatorname{recip}(((j + 1) + 1)) \leq \operatorname{pt\text{-}set\text{-}dist}(m, x, l))\})) \Rightarrow \\ \quad \quad \neg (\operatorname{set\text{-}dist}(m, a, b) = 0) \wedge \neg (\operatorname{set\text{-}dist}(m, b, a) = 0) \end{array}}}
\vnbprop{dist-ring-even-sum-le}{\vnbstmt{\begin{array}{@{}l@{}} \forall m, u, l.; \forall a \in \operatorname{power}(\operatorname{pts}(m)). \\ \quad (\operatorname{is\text{-}metric\text{-}outer\text{-}measure}(m, u) \wedge \operatorname{is\text{-}closed}(m, l)) \Rightarrow  \\ \quad \quad \operatorname{esum}((\lambda q \in \mathbb{N}.\; \\ \quad \quad \quad \quad u((a \cap \operatorname{complement\text{-}in}(\{x \in \operatorname{pts}(m) : (\operatorname{recip}((((2 \cdot q) + 1) + 1)) \leq \operatorname{pt\text{-}set\text{-}dist}(m, x, l))\}, \{x \in \operatorname{pts}(m) : (\operatorname{recip}(((2 \cdot q) + 1)) \leq \operatorname{pt\text{-}set\text{-}dist}(m, x, l))\}))))) \\ \quad \quad \leq\; \operatorname{esup}(\operatorname{image}((\lambda k \in \mathbb{N}.\; u((a \cap \{x \in \operatorname{pts}(m) : (\operatorname{recip}((k + 1)) \leq \operatorname{pt\text{-}set\text{-}dist}(m, x, l))\}))), \\ \quad \quad \quad \quad \mathbb{N})) \end{array}}}
\vnbprop{dist-ring-odd-sum-le}{\vnbstmt{\begin{array}{@{}l@{}} \forall m, u, l.; \forall a \in \operatorname{power}(\operatorname{pts}(m)). \\ \quad (\operatorname{is\text{-}metric\text{-}outer\text{-}measure}(m, u) \wedge \operatorname{is\text{-}closed}(m, l)) \Rightarrow  \\ \quad \quad \operatorname{esum}((\lambda q \in \mathbb{N}.\; \\ \quad \quad \quad \quad u((a \cap \operatorname{complement\text{-}in}(\{x \in \operatorname{pts}(m) : (\operatorname{recip}(((((2 \cdot q) + 1) + 1) + 1)) \leq \operatorname{pt\text{-}set\text{-}dist}(m, x, l))\}, \{x \in \operatorname{pts}(m) : (\operatorname{recip}((((2 \cdot q) + 1) + 1)) \leq \operatorname{pt\text{-}set\text{-}dist}(m, x, l))\}))))) \\ \quad \quad \leq\; \operatorname{esup}(\operatorname{image}((\lambda k \in \mathbb{N}.\; u((a \cap \{x \in \operatorname{pts}(m) : (\operatorname{recip}((k + 1)) \leq \operatorname{pt\text{-}set\text{-}dist}(m, x, l))\}))), \\ \quad \quad \quad \quad \mathbb{N})) \end{array}}}
\vnbprop{metric-outer-layers-limit}{\vnbstmt{\begin{array}{@{}l@{}} \forall m, u, l. \\ \quad (\operatorname{is\text{-}metric\text{-}outer\text{-}measure}(m, u) \wedge \operatorname{is\text{-}closed}(m, l)) \Rightarrow  \\ \quad \quad \forall n \in \mathbb{N}. \\ \quad \quad \quad \{x \in \operatorname{pts}(m) : (\operatorname{recip}((n + 1)) \leq \operatorname{pt\text{-}set\text{-}dist}(m, x, l))\} \\ \quad \quad \quad \subseteq\; \{x \in \operatorname{pts}(m) : (\operatorname{recip}(((n + 1) + 1)) \leq \operatorname{pt\text{-}set\text{-}dist}(m, x, l))\} \wedge \\ \quad \quad \forall a \in \operatorname{power}(\operatorname{pts}(m)). \\ \quad \quad \quad u((a \cap \operatorname{complement\text{-}in}(\operatorname{pts}(m), l))) \\ \quad \quad \quad =\; \operatorname{esup}(\operatorname{image}((\lambda k \in \mathbb{N}.\; u((a \cap \{x \in \operatorname{pts}(m) : (\operatorname{recip}((k + 1)) \leq \operatorname{pt\text{-}set\text{-}dist}(m, x, l))\}))), \\ \quad \quad \quad \quad \quad \mathbb{N})) \end{array}}}
\vnbprop{metric-outer-closed-measurable}{\vnbstmt{\begin{array}{@{}l@{}} \forall m, u, l. \\ \quad (\operatorname{is\text{-}metric\text{-}outer\text{-}measure}(m, u) \wedge \operatorname{is\text{-}closed}(m, l)) \Rightarrow  \\ \quad \quad \operatorname{is\text{-}caratheodory\text{-}measurable}(\operatorname{pts}(m), u, l) \end{array}}}
\vnbprop{metric-outer-open-measurable}{\vnbstmt{\begin{array}{@{}l@{}} \forall m, u, a. \\ \quad (\operatorname{is\text{-}metric\text{-}outer\text{-}measure}(m, u) \wedge \operatorname{is\text{-}open}(m, a)) \Rightarrow  \\ \quad \quad \operatorname{is\text{-}caratheodory\text{-}measurable}(\operatorname{pts}(m), u, a) \end{array}}}
\vnbprop{borel-sets-unfold}{\vnbstmt{\begin{array}{@{}l@{}} \forall m. \\ \quad (\operatorname{borel\text{-}sets}(m) \simeq \operatorname{sigma\text{-}generated}(\operatorname{pts}(m), \{u \in \operatorname{power}(\operatorname{pts}(m)) : \operatorname{is\text{-}open}(m, u)\})) \end{array}}}
\vnbprop{borel-sets-caratheodory-measurable}{\vnbstmt{\begin{array}{@{}l@{}} \forall m, u.; \forall s \in \operatorname{borel\text{-}sets}(m). \\ \quad \operatorname{is\text{-}metric\text{-}outer\text{-}measure}(m, u) \Rightarrow \operatorname{is\text{-}caratheodory\text{-}measurable}(\operatorname{pts}(m), u, s) \end{array}}}
\vnbprop{ccint-inter-left-ray}{\vnbstmt{\begin{array}{@{}l@{}} \forall a, b, c. \\ \quad ((a \in \mathbb{R}) \wedge ((b \in \mathbb{R}) \wedge ((c \in \mathbb{R}) \wedge (c \leq b)))) \Rightarrow  \\ \quad \quad (\operatorname{ccint}(a, b) \cap \{x \in \mathbb{R} : (x \leq c)\}) = \operatorname{ccint}(a, c) \end{array}}}
\vnbprop{ccint-inter-left-ray-all}{\vnbstmt{\begin{array}{@{}l@{}} \forall a, b, c. \\ \quad ((a \in \mathbb{R}) \wedge ((b \in \mathbb{R}) \wedge ((c \in \mathbb{R}) \wedge (b \leq c)))) \Rightarrow  \\ \quad \quad (\operatorname{ccint}(a, b) \cap \{x \in \mathbb{R} : (x \leq c)\}) = \operatorname{ccint}(a, b) \end{array}}}
\vnbprop{ccint-inter-right-ray}{\vnbstmt{\begin{array}{@{}l@{}} \forall a, b, c. \\ \quad ((a \in \mathbb{R}) \wedge ((b \in \mathbb{R}) \wedge ((c \in \mathbb{R}) \wedge (a \leq c)))) \Rightarrow  \\ \quad \quad (\operatorname{ccint}(a, b) \cap \{x \in \mathbb{R} : (c \leq x)\}) = \operatorname{ccint}(c, b) \end{array}}}
\vnbprop{ccint-inter-right-ray-all}{\vnbstmt{\begin{array}{@{}l@{}} \forall a, b, c. \\ \quad ((a \in \mathbb{R}) \wedge ((b \in \mathbb{R}) \wedge ((c \in \mathbb{R}) \wedge (c \leq a)))) \Rightarrow  \\ \quad \quad (\operatorname{ccint}(a, b) \cap \{x \in \mathbb{R} : (c \leq x)\}) = \operatorname{ccint}(a, b) \end{array}}}
\vnbprop{ccint-empty-of-lt}{\vnbstmt{\begin{array}{@{}l@{}} \forall a, b, c. \\ \quad ((a \in \mathbb{R}) \wedge ((b \in \mathbb{R}) \wedge ((c \in \mathbb{R}) \wedge (b < a)))) \Rightarrow (\operatorname{ccint}(a, b) = \emptyset) \end{array}}}
\vnbprop{empty-inter-left-ray}{\vnbstmt{\begin{array}{@{}l@{}} \forall a, b, c. \\ \quad ((a \in \mathbb{R}) \wedge ((b \in \mathbb{R}) \wedge (c \in \mathbb{R}))) \Rightarrow  \\ \quad \quad (\emptyset \cap \{x \in \mathbb{R} : (x \leq c)\}) = \emptyset \end{array}}}
\vnbprop{empty-inter-right-ray}{\vnbstmt{\begin{array}{@{}l@{}} \forall a, b, c. \\ \quad ((a \in \mathbb{R}) \wedge ((b \in \mathbb{R}) \wedge (c \in \mathbb{R}))) \Rightarrow  \\ \quad \quad (\emptyset \cap \{x \in \mathbb{R} : (c \leq x)\}) = \emptyset \end{array}}}
\vnbprop{ccint-in-closed-intervals}{\vnbstmt{\begin{array}{@{}l@{}} \forall a, b. \\ \quad ((a \in \mathbb{R}) \wedge ((b \in \mathbb{R}) \wedge (a \leq b))) \Rightarrow (\operatorname{ccint}(a, b) \in \operatorname{closed\text{-}intervals}) \end{array}}}
\vnbprop{lebesgue-cut}{\vnbstmt{\begin{array}{@{}l@{}} \forall s, c. \\ \quad ((s \in \operatorname{closed\text{-}intervals}) \wedge (c \in \mathbb{R})) \Rightarrow  \\ \quad \quad (s \cap \{x \in \mathbb{R} : (x \leq c)\}) \in \operatorname{closed\text{-}intervals} \wedge \\ \quad \quad (s \cap \{x \in \mathbb{R} : (c \leq x)\}) \in \operatorname{closed\text{-}intervals} \wedge \\ \quad \quad \operatorname{eplus}((\lambda \mathit{ols} \in \operatorname{closed\text{-}intervals}.\; \operatorname{interval\text{-}length}(\mathit{ols}))((s \cap \{x \in \mathbb{R} : (x \leq c)\})), \\ \quad \quad \quad (\lambda \mathit{ols} \in \operatorname{closed\text{-}intervals}.\; \operatorname{interval\text{-}length}(\mathit{ols}))((s \cap \{x \in \mathbb{R} : (c \leq x)\}))) \\ \quad \quad \leq\; (\lambda \mathit{ols} \in \operatorname{closed\text{-}intervals}.\; \operatorname{interval\text{-}length}(\mathit{ols}))(s) \end{array}}}
\vnbprop{esum-nn-superadditive}{\vnbstmt{\begin{array}{@{}l@{}} \forall f, g, h. \\ \quad f \in (\mathbb{N} \to \operatorname{rr\text{-}pos\text{-}star}) \wedge \\ \quad g \in (\mathbb{N} \to \operatorname{rr\text{-}pos\text{-}star}) \wedge \\ \quad h \in (\mathbb{N} \to \operatorname{rr\text{-}pos\text{-}star}) \wedge \\ \quad \forall k \in \mathbb{N}.\; (\operatorname{eplus}(f(k), g(k)) \leq h(k)) \Rightarrow \\ \quad \quad \operatorname{eplus}(\operatorname{esum}(f), \operatorname{esum}(g)) \leq \operatorname{esum}(h) \end{array}}}
\vnbprop{lebesgue-split-le}{\vnbstmt{\begin{array}{@{}l@{}} \forall s, c, u, v. \\ \quad s \in \operatorname{power}(\mathbb{R}) \wedge \\ \quad c \in \mathbb{R} \wedge \\ \quad u \subseteq (s \cap \{x \in \mathbb{R} : (x \leq c)\}) \wedge \\ \quad v \subseteq (s \cap \{x \in \mathbb{R} : (c \leq x)\}) \Rightarrow \\ \quad \quad \operatorname{eplus}(\operatorname{method\text{-}i\text{-}outer}(\mathbb{R}, \operatorname{closed\text{-}intervals}, (\lambda \mathit{ols} \in \operatorname{closed\text{-}intervals}.\; \operatorname{interval\text{-}length}(\mathit{ols})))(u), \\ \quad \quad \quad \operatorname{method\text{-}i\text{-}outer}(\mathbb{R}, \operatorname{closed\text{-}intervals}, (\lambda \mathit{ols} \in \operatorname{closed\text{-}intervals}.\; \operatorname{interval\text{-}length}(\mathit{ols})))(v)) \\ \quad \quad \leq\; \operatorname{method\text{-}i\text{-}outer}(\mathbb{R}, \operatorname{closed\text{-}intervals}, (\lambda \mathit{ols} \in \operatorname{closed\text{-}intervals}.\; \operatorname{interval\text{-}length}(\mathit{ols})))(s) \end{array}}}
\vnbprop{lebesgue-left-ray-measurable}{\vnbstmt{\begin{array}{@{}l@{}} \forall c \in \mathbb{R}. \\ \quad \operatorname{is\text{-}caratheodory\text{-}measurable}(\mathbb{R}, \\ \quad \quad \operatorname{method\text{-}i\text{-}outer}(\mathbb{R}, \operatorname{closed\text{-}intervals}, (\lambda s \in \operatorname{closed\text{-}intervals}.\; \operatorname{interval\text{-}length}(s))), \\ \quad \quad \{x \in \mathbb{R} : (x \leq c)\}) \end{array}}}
\vnbprop{lebesgue-right-ray-measurable}{\vnbstmt{\begin{array}{@{}l@{}} \forall c \in \mathbb{R}. \\ \quad \operatorname{is\text{-}caratheodory\text{-}measurable}(\mathbb{R}, \\ \quad \quad \operatorname{method\text{-}i\text{-}outer}(\mathbb{R}, \operatorname{closed\text{-}intervals}, (\lambda s \in \operatorname{closed\text{-}intervals}.\; \operatorname{interval\text{-}length}(s))), \\ \quad \quad \{x \in \mathbb{R} : (c \leq x)\}) \end{array}}}
\vnbprop{sigma-algebra-binary-union}{\vnbstmt{\begin{array}{@{}l@{}} \forall a, b, \alpha, \beta. \\ \quad (\operatorname{is\text{-}sigma\text{-}algebra}(a, b) \wedge ((\alpha \in b) \wedge (\beta \in b))) \Rightarrow  \\ \quad \quad (\alpha \cup \beta) \in b \end{array}}}
\vnbprop{sigma-algebra-binary-inter}{\vnbstmt{\begin{array}{@{}l@{}} \forall a, b, \alpha, \beta. \\ \quad (\operatorname{is\text{-}sigma\text{-}algebra}(a, b) \wedge ((\alpha \in b) \wedge (\beta \in b))) \Rightarrow  \\ \quad \quad (\alpha \cap \beta) \in b \end{array}}}
\vnbprop{sigma-algebra-sep-meets}{\vnbstmt{\begin{array}{@{}l@{}} \forall a, b, s, c. \\ \quad (\operatorname{is\text{-}sigma\text{-}algebra}(a, b) \wedge (s \in b)) \Rightarrow  \\ \quad \quad \{x \in s : \exists y \in c.\; (y \in s)\} \in b \end{array}}}
\vnbprop{metric-outer-of-separating-sets}{\vnbstmt{\begin{array}{@{}l@{}} \forall m, u. \\ \quad \operatorname{is\text{-}metric\text{-}space}(m) \wedge \\ \quad \operatorname{is\text{-}outer\text{-}measure}(\operatorname{pts}(m), u) \wedge \\ \quad \forall a, b \in \operatorname{power}(\operatorname{pts}(m)). \\ \quad \quad \neg (\operatorname{set\text{-}dist}(m, a, b) = 0) \Rightarrow  \\ \quad \quad \quad \exists s. \\ \quad \quad \quad \quad (\operatorname{is\text{-}caratheodory\text{-}measurable}(\operatorname{pts}(m), u, s) \wedge ((a \subseteq s) \wedge \forall y \in b.\; \neg (y \in s))) \Rightarrow \\ \quad \quad \operatorname{is\text{-}metric\text{-}outer\text{-}measure}(m, u) \end{array}}}
\vnbprop{grid-cell-cover-nonneg}{\vnbstmt{\begin{array}{@{}l@{}} \forall r.; \forall x \in \mathbb{R}. \\ \quad (\operatorname{pos\text{-}rr}(r) \wedge (0 \leq x)) \Rightarrow  \\ \quad \quad \exists n \in \mathbb{N}. \\ \quad \quad \quad x \in \operatorname{complement\text{-}in}(\{\mathit{mhx} \in \mathbb{R} : ((n \cdot r) \leq \mathit{mhx})\}, \{\mathit{mhx} \in \mathbb{R} : (((n \cdot r) + r) \leq \mathit{mhx})\}) \end{array}}}
\vnbprop{grid-cell-cover-nonpos}{\vnbstmt{\begin{array}{@{}l@{}} \forall r.; \forall x \in \mathbb{R}. \\ \quad (\operatorname{pos\text{-}rr}(r) \wedge (x \leq 0)) \Rightarrow  \\ \quad \quad \exists n \in \mathbb{N}. \\ \quad \quad \quad x \in \operatorname{complement\text{-}in}(\{\mathit{mhx} \in \mathbb{R} : (\mathit{mhx} \leq -(n \cdot r))\}, \{\mathit{mhx} \in \mathbb{R} : (\mathit{mhx} \leq -((n \cdot r) + r))\}) \end{array}}}
\vnbprop{grid-cell-close-nonneg}{\vnbstmt{\begin{array}{@{}l@{}} \forall r, n, x, y. \\ \quad \operatorname{pos\text{-}rr}(r) \wedge \\ \quad n \in \mathbb{N} \wedge \\ \quad x \in \operatorname{complement\text{-}in}(\{\mathit{mhx} \in \mathbb{R} : ((n \cdot r) \leq \mathit{mhx})\}, \{\mathit{mhx} \in \mathbb{R} : (((n \cdot r) + r) \leq \mathit{mhx})\}) \wedge \\ \quad y \in \operatorname{complement\text{-}in}(\{\mathit{mhx} \in \mathbb{R} : ((n \cdot r) \leq \mathit{mhx})\}, \{\mathit{mhx} \in \mathbb{R} : (((n \cdot r) + r) \leq \mathit{mhx})\}) \Rightarrow \\ \quad \quad \lvert (x - y) \rvert < r \end{array}}}
\vnbprop{grid-cell-close-nonpos}{\vnbstmt{\begin{array}{@{}l@{}} \forall r, n, x, y. \\ \quad \operatorname{pos\text{-}rr}(r) \wedge \\ \quad n \in \mathbb{N} \wedge \\ \quad x \in \operatorname{complement\text{-}in}(\{\mathit{mhx} \in \mathbb{R} : (\mathit{mhx} \leq -(n \cdot r))\}, \{\mathit{mhx} \in \mathbb{R} : (\mathit{mhx} \leq -((n \cdot r) + r))\}) \wedge \\ \quad y \in \operatorname{complement\text{-}in}(\{\mathit{mhx} \in \mathbb{R} : (\mathit{mhx} \leq -(n \cdot r))\}, \{\mathit{mhx} \in \mathbb{R} : (\mathit{mhx} \leq -((n \cdot r) + r))\}) \Rightarrow \\ \quad \quad \lvert (x - y) \rvert < r \end{array}}}
\vnbprop{lebesgue-metric-outer-measure}{\vnbstmt{\begin{array}{@{}l@{}} \operatorname{is\text{-}metric\text{-}outer\text{-}measure}(\operatorname{rr\text{-}ms}, \\ \quad \operatorname{method\text{-}i\text{-}outer}(\mathbb{R}, \operatorname{closed\text{-}intervals}, (\lambda s \in \operatorname{closed\text{-}intervals}.\; \operatorname{interval\text{-}length}(s)))) \end{array}}}
\vnbprop{lebesgue-borel-measurable}{\vnbstmt{\begin{array}{@{}l@{}} \forall s \in \operatorname{borel\text{-}sets}(\operatorname{rr\text{-}ms}). \\ \quad \operatorname{is\text{-}caratheodory\text{-}measurable}(\mathbb{R}, \\ \quad \quad \operatorname{method\text{-}i\text{-}outer}(\mathbb{R}, \operatorname{closed\text{-}intervals}, (\lambda \mathit{ols} \in \operatorname{closed\text{-}intervals}.\; \operatorname{interval\text{-}length}(\mathit{ols}))), \\ \quad \quad s) \end{array}}}
\vnbprop{lebesgue-open-left-ray-measurable}{\vnbstmt{\begin{array}{@{}l@{}} \forall c \in \mathbb{R}. \\ \quad \{x \in \mathbb{R} : (x < c)\} \\ \quad \in\; \operatorname{caratheodory\text{-}sets}(\mathbb{R}, \\ \quad \quad \operatorname{method\text{-}i\text{-}outer}(\mathbb{R}, \operatorname{closed\text{-}intervals}, (\lambda s \in \operatorname{closed\text{-}intervals}.\; \operatorname{interval\text{-}length}(s)))) \end{array}}}
\vnbprop{lebesgue-open-right-ray-measurable}{\vnbstmt{\begin{array}{@{}l@{}} \forall c \in \mathbb{R}. \\ \quad \{x \in \mathbb{R} : (c < x)\} \\ \quad \in\; \operatorname{caratheodory\text{-}sets}(\mathbb{R}, \\ \quad \quad \operatorname{method\text{-}i\text{-}outer}(\mathbb{R}, \operatorname{closed\text{-}intervals}, (\lambda s \in \operatorname{closed\text{-}intervals}.\; \operatorname{interval\text{-}length}(s)))) \end{array}}}
\vnbprop{lebesgue-intervals-measurable}{\vnbstmt{\begin{array}{@{}l@{}} \forall a, b. \\ \quad ((a \in \mathbb{R}) \wedge (b \in \mathbb{R})) \Rightarrow  \\ \quad \quad \operatorname{is\text{-}caratheodory\text{-}measurable}(\mathbb{R}, \\ \quad \quad \quad \operatorname{method\text{-}i\text{-}outer}(\mathbb{R}, \operatorname{closed\text{-}intervals}, (\lambda s \in \operatorname{closed\text{-}intervals}.\; \operatorname{interval\text{-}length}(s))), \\ \quad \quad \quad \operatorname{ccint}(a, b)) \wedge \\ \quad \quad \operatorname{is\text{-}caratheodory\text{-}measurable}(\mathbb{R}, \\ \quad \quad \quad \operatorname{method\text{-}i\text{-}outer}(\mathbb{R}, \operatorname{closed\text{-}intervals}, (\lambda s \in \operatorname{closed\text{-}intervals}.\; \operatorname{interval\text{-}length}(s))), \\ \quad \quad \quad \operatorname{ooint}(a, b)) \wedge \\ \quad \quad \operatorname{is\text{-}caratheodory\text{-}measurable}(\mathbb{R}, \\ \quad \quad \quad \operatorname{method\text{-}i\text{-}outer}(\mathbb{R}, \operatorname{closed\text{-}intervals}, (\lambda s \in \operatorname{closed\text{-}intervals}.\; \operatorname{interval\text{-}length}(s))), \\ \quad \quad \quad \operatorname{coint}(a, b)) \wedge \\ \quad \quad \operatorname{is\text{-}caratheodory\text{-}measurable}(\mathbb{R}, \\ \quad \quad \quad \operatorname{method\text{-}i\text{-}outer}(\mathbb{R}, \operatorname{closed\text{-}intervals}, (\lambda s \in \operatorname{closed\text{-}intervals}.\; \operatorname{interval\text{-}length}(s))), \\ \quad \quad \quad \operatorname{ocint}(a, b)) \end{array}}}

\vnbbookfile{theorem-library/riesz-markov-cantor}{theorem-library/riesz-markov-cantor}
\noindent{\small theorem-library/riesz-markov-cantor.scm -- RIESZ-MARKOV, STAGE 2 (RM-2, 2026-10-10): a positive linear functional on the continuous functions of the Cantor model K is integration against a measure.  Vocabulary: structure-library/positive-functional.scm (CONT-FNS, IS-POSITIVE-FUNCTIONAL,}\par
\vnbprop{rmc-cont-fns-mem}{\vnbstmt{\begin{array}{@{}l@{}} \forall s, h. \\ \quad (h \in \operatorname{cont\text{-}fns}(s)) \Leftrightarrow ((h \in (\operatorname{pts}(s) \to \mathbb{R})) \wedge \operatorname{is\text{-}continuous}(s, \operatorname{rr\text{-}ms}, h)) \end{array}}}
\vnbprop{rmc-clopen-mem}{\vnbstmt{\begin{array}{@{}l@{}} \forall s, u. \\ \quad (u \in \operatorname{clopen\text{-}sets}(s)) \Leftrightarrow ((u \in \operatorname{power}(\operatorname{pts}(s))) \wedge (\operatorname{is\text{-}open}(s, u) \wedge \operatorname{is\text{-}closed}(s, u))) \end{array}}}
\vnbprop{rmc-clopen-intro}{\vnbstmt{\forall s, u.\; ((\operatorname{is\text{-}open}(s, u) \wedge \operatorname{is\text{-}closed}(s, u)) \Rightarrow (u \in \operatorname{clopen\text{-}sets}(s)))}}
\vnbprop{rmc-clopen-parts}{\vnbstmt{\begin{array}{@{}l@{}} \forall s.; \forall u \in \operatorname{clopen\text{-}sets}(s). \\ \quad \operatorname{is\text{-}open}(s, u) \wedge (\operatorname{is\text{-}closed}(s, u) \wedge (u \subseteq \operatorname{pts}(s))) \end{array}}}
\vnbprop{rmc-cont-fns-parts}{\vnbstmt{\begin{array}{@{}l@{}} \forall s.; \forall h \in \operatorname{cont\text{-}fns}(s). \\ \quad h \in (\operatorname{pts}(s) \to \mathbb{R}) \wedge \\ \quad \operatorname{is\text{-}metric\text{-}space}(s) \wedge \\ \quad \forall a \in \operatorname{pts}(s).\; \operatorname{is\text{-}continuous\text{-}at}(s, \operatorname{rr\text{-}ms}, h, a) \end{array}}}
\vnbprop{rmc-cont-fns-intro}{\vnbstmt{\begin{array}{@{}l@{}} \forall s, h. \\ \quad \operatorname{is\text{-}metric\text{-}space}(s) \wedge \\ \quad h \in (\operatorname{pts}(s) \to \mathbb{R}) \wedge \\ \quad \forall a \in \operatorname{pts}(s).\; \operatorname{is\text{-}continuous\text{-}at}(s, \operatorname{rr\text{-}ms}, h, a) \Rightarrow \\ \quad \quad h \in \operatorname{cont\text{-}fns}(s) \end{array}}}
\vnbprop{rmc-cont-fns-scale}{\vnbstmt{\begin{array}{@{}l@{}} \forall s, c, f. \\ \quad ((c \in \mathbb{R}) \wedge (f \in \operatorname{cont\text{-}fns}(s))) \Rightarrow  \\ \quad \quad (\lambda x \in \operatorname{pts}(s).\; (c \cdot f(x))) \in \operatorname{cont\text{-}fns}(s) \end{array}}}
\vnbprop{rmc-cont-fns-add}{\vnbstmt{\begin{array}{@{}l@{}} \forall s, f, g. \\ \quad ((f \in \operatorname{cont\text{-}fns}(s)) \wedge (g \in \operatorname{cont\text{-}fns}(s))) \Rightarrow  \\ \quad \quad (\lambda x \in \operatorname{pts}(s).\; (f(x) + g(x))) \in \operatorname{cont\text{-}fns}(s) \end{array}}}
\vnbprop{rmc-cont-at-locally-const}{\vnbstmt{\begin{array}{@{}l@{}} \forall s, h, a, w. \\ \quad (\operatorname{is\text{-}open}(s, w) \wedge ((a \in w) \wedge ((h \in (\operatorname{pts}(s) \to \mathbb{R})) \wedge \forall y \in w.\; (h(y) = h(a))))) \Rightarrow  \\ \quad \quad \operatorname{is\text{-}continuous\text{-}at}(s, \operatorname{rr\text{-}ms}, h, a) \end{array}}}
\vnbprop{rmc-indicator-fun}{\vnbstmt{\forall b \in \mathbf{Set}, a.\; (\operatorname{indicator}(b, a) \in (b \to \mathbb{R}))}}
\vnbprop{rmc-indicator-values}{\vnbstmt{\begin{array}{@{}l@{}} \forall b, a.; \forall x \in b. \\ \quad ((\operatorname{indicator}(b, a)(x) \in \mathbb{R}) \wedge ((0 \leq \operatorname{indicator}(b, a)(x)) \wedge (\operatorname{indicator}(b, a)(x) \leq 1))) \end{array}}}
\vnbprop{rmc-indicator-cont}{\vnbstmt{\forall s.; \forall u \in \operatorname{clopen\text{-}sets}(s).\; (\operatorname{indicator}(\operatorname{pts}(s), u) \in \operatorname{cont\text{-}fns}(s))}}
\vnbprop{rmc-cyl-set}{\vnbstmt{\forall s, q, a, b.\; (\operatorname{rm\text{-}cyl}(s, q, a, b) \in \mathbf{Set})}}
\vnbprop{rmc-cyl-clopen}{\vnbstmt{\forall s, q, a.; \forall b \in \mathbb{N}.\; (\operatorname{rm\text{-}cyl}(s, q, a, b) \in \operatorname{clopen\text{-}sets}(\operatorname{rm\text{-}kms}(s, q)))}}
\vnbprop{rmc-cyl-indicator-cont}{\vnbstmt{\begin{array}{@{}l@{}} \forall s, q, a.; \forall b \in \mathbb{N}. \\ \quad \operatorname{indicator}(\operatorname{pts}(\operatorname{rm\text{-}kms}(s, q)), \operatorname{rm\text{-}cyl}(s, q, a, b)) \in \operatorname{cont\text{-}fns}(\operatorname{rm\text{-}kms}(s, q)) \end{array}}}
\vnbprop{rmc-compl-union}{\vnbstmt{\begin{array}{@{}l@{}} \forall a, u, v. \\ \quad \operatorname{complement\text{-}in}(a, (u \cup v)) = (\operatorname{complement\text{-}in}(a, u) \cap \operatorname{complement\text{-}in}(a, v)) \end{array}}}
\vnbprop{rmc-compl-inter}{\vnbstmt{\begin{array}{@{}l@{}} \forall a, u, v. \\ \quad \operatorname{complement\text{-}in}(a, (u \cap v)) = (\operatorname{complement\text{-}in}(a, u) \cup \operatorname{complement\text{-}in}(a, v)) \end{array}}}
\vnbprop{rmc-diff-inter}{\vnbstmt{\begin{array}{@{}l@{}} \forall a, u, v. \\ \quad (u \subseteq a) \Rightarrow (\operatorname{difference}(u, v) = (u \cap \operatorname{complement\text{-}in}(a, v))) \end{array}}}
\vnbprop{rmc-open-union}{\vnbstmt{\begin{array}{@{}l@{}} \forall s, u, v. \\ \quad (\operatorname{is\text{-}open}(s, u) \wedge \operatorname{is\text{-}open}(s, v)) \Rightarrow \operatorname{is\text{-}open}(s, (u \cup v)) \end{array}}}
\vnbprop{rmc-clopen-compl}{\vnbstmt{\forall s.; \forall u \in \operatorname{clopen\text{-}sets}(s).\; (\operatorname{complement\text{-}in}(\operatorname{pts}(s), u) \in \operatorname{clopen\text{-}sets}(s))}}
\vnbprop{rmc-clopen-inter}{\vnbstmt{\begin{array}{@{}l@{}} \forall s, u, v. \\ \quad ((u \in \operatorname{clopen\text{-}sets}(s)) \wedge (v \in \operatorname{clopen\text{-}sets}(s))) \Rightarrow  \\ \quad \quad (u \cap v) \in \operatorname{clopen\text{-}sets}(s) \end{array}}}
\vnbprop{rmc-clopen-union}{\vnbstmt{\begin{array}{@{}l@{}} \forall s, u, v. \\ \quad ((u \in \operatorname{clopen\text{-}sets}(s)) \wedge (v \in \operatorname{clopen\text{-}sets}(s))) \Rightarrow  \\ \quad \quad (u \cup v) \in \operatorname{clopen\text{-}sets}(s) \end{array}}}
\vnbprop{rmc-clopen-diff}{\vnbstmt{\begin{array}{@{}l@{}} \forall s, u, v. \\ \quad ((u \in \operatorname{clopen\text{-}sets}(s)) \wedge (v \in \operatorname{clopen\text{-}sets}(s))) \Rightarrow  \\ \quad \quad \operatorname{difference}(u, v) \in \operatorname{clopen\text{-}sets}(s) \end{array}}}
\vnbprop{rmc-clopen-empty-whole}{\vnbstmt{\begin{array}{@{}l@{}} \forall s. \\ \quad \operatorname{is\text{-}metric\text{-}space}(s) \Rightarrow  \\ \quad \quad (\emptyset \in \operatorname{clopen\text{-}sets}(s)) \wedge (\operatorname{pts}(s) \in \operatorname{clopen\text{-}sets}(s)) \end{array}}}
\vnbprop{rmc-clopen-sub-power}{\vnbstmt{\forall s.\; (\operatorname{is\text{-}metric\text{-}space}(s) \Rightarrow (\operatorname{clopen\text{-}sets}(s) \subseteq \operatorname{power}(\operatorname{pts}(s))))}}
\vnbprop{rmc-clopen-algebra}{\vnbstmt{\begin{array}{@{}l@{}} \forall s. \\ \quad \operatorname{is\text{-}metric\text{-}space}(s) \Rightarrow \operatorname{is\text{-}algebra\text{-}of\text{-}sets}(\operatorname{pts}(s), \operatorname{clopen\text{-}sets}(s)) \end{array}}}
\vnbprop{rmc-pf-add}{\vnbstmt{\begin{array}{@{}l@{}} \forall s, m, f, g. \\ \quad (\operatorname{is\text{-}positive\text{-}functional}(s, m) \wedge ((f \in \operatorname{cont\text{-}fns}(s)) \wedge (g \in \operatorname{cont\text{-}fns}(s)))) \Rightarrow  \\ \quad \quad m((\lambda x \in \operatorname{pts}(s).\; (f(x) + g(x)))) = (m(f) + m(g)) \end{array}}}
\vnbprop{rmc-pf-scale}{\vnbstmt{\begin{array}{@{}l@{}} \forall s, m, c, f. \\ \quad (\operatorname{is\text{-}positive\text{-}functional}(s, m) \wedge ((c \in \mathbb{R}) \wedge (f \in \operatorname{cont\text{-}fns}(s)))) \Rightarrow  \\ \quad \quad m((\lambda x \in \operatorname{pts}(s).\; (c \cdot f(x)))) = (c \cdot m(f)) \end{array}}}
\vnbprop{rmc-pf-mono}{\vnbstmt{\begin{array}{@{}l@{}} \forall s, m, f, g. \\ \quad \operatorname{is\text{-}positive\text{-}functional}(s, m) \wedge \\ \quad f \in \operatorname{cont\text{-}fns}(s) \wedge \\ \quad g \in \operatorname{cont\text{-}fns}(s) \wedge \\ \quad \forall x \in \operatorname{pts}(s).\; (f(x) \leq g(x)) \Rightarrow \\ \quad \quad m(f) \leq m(g) \end{array}}}
\vnbprop{rmc-pf-le-sum}{\vnbstmt{\begin{array}{@{}l@{}} \forall s, m, f, a, b. \\ \quad \operatorname{is\text{-}positive\text{-}functional}(s, m) \wedge \\ \quad f \in \operatorname{cont\text{-}fns}(s) \wedge \\ \quad a \in \operatorname{cont\text{-}fns}(s) \wedge \\ \quad b \in \operatorname{cont\text{-}fns}(s) \wedge \\ \quad \forall x \in \operatorname{pts}(s).\; (f(x) \leq (a(x) + b(x))) \Rightarrow \\ \quad \quad m(f) \leq (m(a) + m(b)) \end{array}}}
\vnbprop{rmc-pf-ge-sum}{\vnbstmt{\begin{array}{@{}l@{}} \forall s, m, f, a, b. \\ \quad \operatorname{is\text{-}positive\text{-}functional}(s, m) \wedge \\ \quad f \in \operatorname{cont\text{-}fns}(s) \wedge \\ \quad a \in \operatorname{cont\text{-}fns}(s) \wedge \\ \quad b \in \operatorname{cont\text{-}fns}(s) \wedge \\ \quad \forall x \in \operatorname{pts}(s).\; ((a(x) + b(x)) \leq f(x)) \Rightarrow \\ \quad \quad (m(a) + m(b)) \leq m(f) \end{array}}}
\vnbprop{rmc-ind-nonneg}{\vnbstmt{\begin{array}{@{}l@{}} \forall s, m, u. \\ \quad (\operatorname{is\text{-}positive\text{-}functional}(s, m) \wedge (u \in \operatorname{clopen\text{-}sets}(s))) \Rightarrow  \\ \quad \quad (m(\operatorname{indicator}(\operatorname{pts}(s), u)) \in \mathbb{R}) \wedge (0 \leq m(\operatorname{indicator}(\operatorname{pts}(s), u))) \end{array}}}
\vnbprop{rmc-ind-split-ge}{\vnbstmt{\begin{array}{@{}l@{}} \forall s, m, c, a. \\ \quad (\operatorname{is\text{-}positive\text{-}functional}(s, m) \wedge ((c \in \operatorname{clopen\text{-}sets}(s)) \wedge (a \in \operatorname{clopen\text{-}sets}(s)))) \Rightarrow  \\ \quad \quad m(\operatorname{indicator}(\operatorname{pts}(s), (c \cap a))) + m(\operatorname{indicator}(\operatorname{pts}(s), \operatorname{difference}(c, a))) \\ \quad \quad \leq\; m(\operatorname{indicator}(\operatorname{pts}(s), c)) \end{array}}}
\vnbprop{rmc-ind-cover-le}{\vnbstmt{\begin{array}{@{}l@{}} \forall s, m, a, c. \\ \quad (\operatorname{is\text{-}positive\text{-}functional}(s, m) \wedge ((a \in \operatorname{clopen\text{-}sets}(s)) \wedge (c \in \operatorname{clopen\text{-}sets}(s)))) \Rightarrow  \\ \quad \quad m(\operatorname{indicator}(\operatorname{pts}(s), a)) \leq (m(\operatorname{indicator}(\operatorname{pts}(s), \operatorname{difference}(a, c))) + m(\operatorname{indicator}(\operatorname{pts}(s), c))) \end{array}}}
\vnbprop{rmc-ind-empty}{\vnbstmt{\begin{array}{@{}l@{}} \forall s, m. \\ \quad (\operatorname{is\text{-}positive\text{-}functional}(s, m) \wedge \operatorname{is\text{-}metric\text{-}space}(s)) \Rightarrow  \\ \quad \quad m(\operatorname{indicator}(\operatorname{pts}(s), \emptyset)) = 0 \end{array}}}
\vnbprop{rmc-ind-additive}{\vnbstmt{\begin{array}{@{}l@{}} \forall s, m, u, v. \\ \quad \operatorname{is\text{-}positive\text{-}functional}(s, m) \wedge \\ \quad u \in \operatorname{clopen\text{-}sets}(s) \wedge \\ \quad v \in \operatorname{clopen\text{-}sets}(s) \wedge \\ \quad (u \cap v) = \emptyset \Rightarrow \\ \quad \quad m(\operatorname{indicator}(\operatorname{pts}(s), (u \cup v))) = (m(\operatorname{indicator}(\operatorname{pts}(s), u)) + m(\operatorname{indicator}(\operatorname{pts}(s), v))) \end{array}}}
\vnbprop{rmc-finite-subcover}{\vnbstmt{\begin{array}{@{}l@{}} \forall m, a, q. \\ \quad \operatorname{is\text{-}compact}(m) \wedge \\ \quad \operatorname{is\text{-}closed}(m, a) \wedge \\ \quad \forall n \in \mathbb{N}.\; \operatorname{is\text{-}open}(m, q(n)) \wedge \\ \quad \forall x \in a.\; \exists n \in \mathbb{N}.\; (x \in q(n)) \Rightarrow \\ \quad \quad \exists f \in \mathbb{N}.\; \forall x \in a.\; \exists n \in \mathbb{N}.\; ((n < f) \wedge (x \in q(n))) \end{array}}}
\vnbprop{rmc-rps-of-nonneg}{\vnbstmt{\forall x \in \mathbb{R}.\; ((0 \leq x) \Rightarrow (x \in \operatorname{rr\text{-}pos\text{-}star}))}}
\vnbprop{rmc-clopen-set}{\vnbstmt{\forall s.\; (\operatorname{is\text{-}metric\text{-}space}(s) \Rightarrow (\operatorname{clopen\text{-}sets}(s) \in \mathbf{Set}))}}
\vnbprop{rmc-mu0-value}{\vnbstmt{\begin{array}{@{}l@{}} \forall s, m, u. \\ \quad (\operatorname{is\text{-}positive\text{-}functional}(s, m) \wedge (u \in \operatorname{clopen\text{-}sets}(s))) \Rightarrow  \\ \quad \quad (\lambda \mathit{rcu} \in \operatorname{clopen\text{-}sets}(s).\; m(\operatorname{indicator}(\operatorname{pts}(s), \mathit{rcu})))(u) = m(\operatorname{indicator}(\operatorname{pts}(s), u)) \end{array}}}
\vnbprop{rmc-fseq-laws}{\vnbstmt{\begin{array}{@{}l@{}} \forall s, m, q. \\ \quad (\operatorname{is\text{-}positive\text{-}functional}(s, m) \wedge (q \in (\mathbb{N} \to \operatorname{clopen\text{-}sets}(s)))) \Rightarrow  \\ \quad \quad (\lambda \mathit{cvk} \in \mathbb{N}.\; (\lambda u \in \operatorname{clopen\text{-}sets}(s).\; m(\operatorname{indicator}(\operatorname{pts}(s), u)))(q(\mathit{cvk}))) \\ \quad \quad \in\; (\mathbb{N} \to \operatorname{rr\text{-}pos\text{-}star}) \wedge \\ \quad \quad \forall k \in \mathbb{N}. \\ \quad \quad \quad (\lambda \mathit{cvk} \in \mathbb{N}.\; (\lambda u \in \operatorname{clopen\text{-}sets}(s).\; m(\operatorname{indicator}(\operatorname{pts}(s), u)))(q(\mathit{cvk})))(k) \\ \quad \quad \quad =\; m(\operatorname{indicator}(\operatorname{pts}(s), q(k))) \end{array}}}
\vnbprop{rmc-prefix-bound}{\vnbstmt{\begin{array}{@{}l@{}} \forall q \in \mathbb{N}, s, m, b.; \forall a \in \operatorname{clopen\text{-}sets}(s). \\ \quad \operatorname{is\text{-}positive\text{-}functional}(s, m) \wedge \\ \quad \operatorname{is\text{-}metric\text{-}space}(s) \wedge \\ \quad b \in (\mathbb{N} \to \operatorname{clopen\text{-}sets}(s)) \wedge \\ \quad \forall x \in a.\; \exists n \in \mathbb{N}.\; ((n < q) \wedge (x \in b(n))) \Rightarrow \\ \quad \quad m(\operatorname{indicator}(\operatorname{pts}(s), a)) \\ \quad \quad \leq\; \operatorname{finsum}(\operatorname{rr\text{-}pos\text{-}star\text{-}add\text{-}monoid}, \\ \quad \quad \quad (\lambda k \in \mathbb{N}.\; (\lambda u \in \operatorname{clopen\text{-}sets}(s).\; m(\operatorname{indicator}(\operatorname{pts}(s), u)))(b(k))), \\ \quad \quad \quad \mathrm{OS}(q)) \end{array}}}
\vnbprop{rmc-countable-subadditive}{\vnbstmt{\begin{array}{@{}l@{}} \forall s, m, q, a. \\ \quad \operatorname{is\text{-}positive\text{-}functional}(s, m) \wedge \\ \quad \operatorname{is\text{-}compact}(s) \wedge \\ \quad q \in (\mathbb{N} \to \operatorname{clopen\text{-}sets}(s)) \wedge \\ \quad a \in \operatorname{clopen\text{-}sets}(s) \wedge \\ \quad a \subseteq \operatorname{big\text{-}union}(n, \mathbb{N}, q(n)) \Rightarrow \\ \quad \quad m(\operatorname{indicator}(\operatorname{pts}(s), a)) \\ \quad \quad \leq\; \operatorname{esum}((\lambda k \in \mathbb{N}.\; (\lambda u \in \operatorname{clopen\text{-}sets}(s).\; m(\operatorname{indicator}(\operatorname{pts}(s), u)))(q(k)))) \end{array}}}
\vnbprop{rmc-premeasure}{\vnbstmt{\begin{array}{@{}l@{}} \forall s, m. \\ \quad (\operatorname{is\text{-}positive\text{-}functional}(s, m) \wedge \operatorname{is\text{-}metric\text{-}space}(s)) \Rightarrow  \\ \quad \quad \operatorname{is\text{-}premeasure}(\operatorname{pts}(s), \operatorname{clopen\text{-}sets}(s), (\lambda u \in \operatorname{clopen\text{-}sets}(s).\; m(\operatorname{indicator}(\operatorname{pts}(s), u)))) \end{array}}}
\vnbprop{rmc-power-subset}{\vnbstmt{\forall a.; \forall x \in \operatorname{power}(a).\; (x \subseteq a)}}
\vnbprop{rmc-outer-on-clopen}{\vnbstmt{\begin{array}{@{}l@{}} \forall s, m, a. \\ \quad (\operatorname{is\text{-}positive\text{-}functional}(s, m) \wedge (\operatorname{is\text{-}compact}(s) \wedge (a \in \operatorname{clopen\text{-}sets}(s)))) \Rightarrow  \\ \quad \quad \operatorname{method\text{-}i\text{-}outer}(\operatorname{pts}(s), \operatorname{clopen\text{-}sets}(s), (\lambda u \in \operatorname{clopen\text{-}sets}(s).\; m(\operatorname{indicator}(\operatorname{pts}(s), u))))(a) \\ \quad \quad =\; m(\operatorname{indicator}(\operatorname{pts}(s), a)) \end{array}}}
\vnbprop{rmc-clopen-measurable}{\vnbstmt{\begin{array}{@{}l@{}} \forall s, m, a. \\ \quad (\operatorname{is\text{-}positive\text{-}functional}(s, m) \wedge (\operatorname{is\text{-}metric\text{-}space}(s) \wedge (a \in \operatorname{clopen\text{-}sets}(s)))) \Rightarrow  \\ \quad \quad \operatorname{is\text{-}caratheodory\text{-}measurable}(\operatorname{pts}(s), \\ \quad \quad \quad \operatorname{method\text{-}i\text{-}outer}(\operatorname{pts}(s), \operatorname{clopen\text{-}sets}(s), (\lambda u \in \operatorname{clopen\text{-}sets}(s).\; m(\operatorname{indicator}(\operatorname{pts}(s), u)))), \\ \quad \quad \quad a) \end{array}}}
\vnbprop{rmc-measure}{\vnbstmt{\begin{array}{@{}l@{}} \forall s, m. \\ \quad (\operatorname{is\text{-}positive\text{-}functional}(s, m) \wedge \operatorname{is\text{-}compact}(s)) \Rightarrow  \\ \quad \quad \operatorname{is\text{-}measure}(\operatorname{pts}(s), \\ \quad \quad \quad \operatorname{caratheodory\text{-}sets}(\operatorname{pts}(s), \\ \quad \quad \quad \quad \operatorname{method\text{-}i\text{-}outer}(\operatorname{pts}(s), \operatorname{clopen\text{-}sets}(s), (\lambda \mathit{rcu} \in \operatorname{clopen\text{-}sets}(s).\; m(\operatorname{indicator}(\operatorname{pts}(s), \mathit{rcu}))))), \\ \quad \quad \quad \operatorname{res}(\operatorname{method\text{-}i\text{-}outer}(\operatorname{pts}(s), \operatorname{clopen\text{-}sets}(s), (\lambda \mathit{rcu} \in \operatorname{clopen\text{-}sets}(s).\; m(\operatorname{indicator}(\operatorname{pts}(s), \mathit{rcu})))), \\ \quad \quad \quad \quad \operatorname{caratheodory\text{-}sets}(\operatorname{pts}(s), \\ \quad \quad \quad \quad \quad \operatorname{method\text{-}i\text{-}outer}(\operatorname{pts}(s), \operatorname{clopen\text{-}sets}(s), (\lambda \mathit{rcu} \in \operatorname{clopen\text{-}sets}(s).\; m(\operatorname{indicator}(\operatorname{pts}(s), \mathit{rcu}))))))) \wedge \\ \quad \quad \operatorname{clopen\text{-}sets}(s) \\ \quad \quad \subseteq\; \operatorname{caratheodory\text{-}sets}(\operatorname{pts}(s), \\ \quad \quad \quad \operatorname{method\text{-}i\text{-}outer}(\operatorname{pts}(s), \operatorname{clopen\text{-}sets}(s), (\lambda \mathit{rcu} \in \operatorname{clopen\text{-}sets}(s).\; m(\operatorname{indicator}(\operatorname{pts}(s), \mathit{rcu}))))) \wedge \\ \quad \quad \forall u \in \operatorname{clopen\text{-}sets}(s). \\ \quad \quad \quad \operatorname{res}(\operatorname{method\text{-}i\text{-}outer}(\operatorname{pts}(s), \operatorname{clopen\text{-}sets}(s), (\lambda \mathit{rcu} \in \operatorname{clopen\text{-}sets}(s).\; m(\operatorname{indicator}(\operatorname{pts}(s), \mathit{rcu})))), \operatorname{caratheodory\text{-}sets}(\operatorname{pts}(s), \operatorname{method\text{-}i\text{-}outer}(\operatorname{pts}(s), \operatorname{clopen\text{-}sets}(s), (\lambda \mathit{rcu} \in \operatorname{clopen\text{-}sets}(s).\; m(\operatorname{indicator}(\operatorname{pts}(s), \mathit{rcu}))))))(u) \\ \quad \quad \quad =\; m(\operatorname{indicator}(\operatorname{pts}(s), u)) \end{array}}}
\vnbprop{rmc-cyl-same}{\vnbstmt{\begin{array}{@{}l@{}} \forall s, q, e, p, b. \\ \quad ((e \in \operatorname{rm\text{-}k}(s, q)) \wedge (p \in \operatorname{rm\text{-}cyl}(s, q, e, b))) \Rightarrow  \\ \quad \quad \operatorname{rm\text{-}cyl}(s, q, p, b) = \operatorname{rm\text{-}cyl}(s, q, e, b) \end{array}}}
\vnbprop{rmc-saturated-open}{\vnbstmt{\begin{array}{@{}l@{}} \forall s, q, w.; \forall b \in \mathbb{N}. \\ \quad ((w \subseteq \operatorname{rm\text{-}k}(s, q)) \wedge \forall y \in w, z \in \operatorname{rm\text{-}cyl}(s, q, y, b).\; (z \in w)) \Rightarrow  \\ \quad \quad \operatorname{is\text{-}open}(\operatorname{rm\text{-}kms}(s, q), w) \end{array}}}
\vnbprop{rmc-saturated-clopen}{\vnbstmt{\begin{array}{@{}l@{}} \forall s, q, w.; \forall b \in \mathbb{N}. \\ \quad ((w \subseteq \operatorname{rm\text{-}k}(s, q)) \wedge \forall y \in w, z \in \operatorname{rm\text{-}cyl}(s, q, y, b).\; (z \in w)) \Rightarrow  \\ \quad \quad w \in \operatorname{clopen\text{-}sets}(\operatorname{rm\text{-}kms}(s, q)) \end{array}}}
\vnbprop{rmc-w-mem}{\vnbstmt{\begin{array}{@{}l@{}} \forall s, q, u, b, x. \\ \quad x \in \{e \in \operatorname{rm\text{-}k}(s, q) : (\operatorname{rm\text{-}cyl}(s, q, e, b) \subseteq u)\} \\ \quad \Leftrightarrow\; (x \in \operatorname{rm\text{-}k}(s, q)) \wedge (\operatorname{rm\text{-}cyl}(s, q, x, b) \subseteq u) \end{array}}}
\vnbprop{rmc-w-clopen}{\vnbstmt{\begin{array}{@{}l@{}} \forall s, q, u.; \forall b \in \mathbb{N}. \\ \quad \{e \in \operatorname{rm\text{-}k}(s, q) : (\operatorname{rm\text{-}cyl}(s, q, e, b) \subseteq u)\} \in \operatorname{clopen\text{-}sets}(\operatorname{rm\text{-}kms}(s, q)) \end{array}}}
\vnbprop{rmc-open-union-w}{\vnbstmt{\begin{array}{@{}l@{}} \forall s, q, u. \\ \quad \operatorname{is\text{-}open}(\operatorname{rm\text{-}kms}(s, q), u) \Rightarrow  \\ \quad \quad u = \operatorname{big\text{-}union}(n, \mathbb{N}, (\lambda k \in \mathbb{N}.\; \{e \in \operatorname{rm\text{-}k}(s, q) : (\operatorname{rm\text{-}cyl}(s, q, e, k) \subseteq u)\})(n)) \end{array}}}
\vnbprop{rmc-k-open-measurable}{\vnbstmt{\begin{array}{@{}l@{}} \forall s, q, m, u. \\ \quad (\operatorname{is\text{-}positive\text{-}functional}(\operatorname{rm\text{-}kms}(s, q), m) \wedge \operatorname{is\text{-}open}(\operatorname{rm\text{-}kms}(s, q), u)) \Rightarrow  \\ \quad \quad u \\ \quad \quad \in\; \operatorname{caratheodory\text{-}sets}(\operatorname{pts}(\operatorname{rm\text{-}kms}(s, q)), \\ \quad \quad \quad \operatorname{method\text{-}i\text{-}outer}(\operatorname{pts}(\operatorname{rm\text{-}kms}(s, q)), \operatorname{clopen\text{-}sets}(\operatorname{rm\text{-}kms}(s, q)), \\ \quad \quad \quad \quad (\lambda \mathit{rcu} \in \operatorname{clopen\text{-}sets}(\operatorname{rm\text{-}kms}(s, q)).\; m(\operatorname{indicator}(\operatorname{pts}(\operatorname{rm\text{-}kms}(s, q)), \mathit{rcu}))))) \end{array}}}
\vnbprop{rmc-k-measure}{\vnbstmt{\begin{array}{@{}l@{}} \forall s, q, m. \\ \quad \operatorname{is\text{-}positive\text{-}functional}(\operatorname{rm\text{-}kms}(s, q), m) \Rightarrow  \\ \quad \quad \operatorname{is\text{-}measure}(\operatorname{pts}(\operatorname{rm\text{-}kms}(s, q)), \\ \quad \quad \quad \operatorname{caratheodory\text{-}sets}(\operatorname{pts}(\operatorname{rm\text{-}kms}(s, q)), \\ \quad \quad \quad \quad \operatorname{method\text{-}i\text{-}outer}(\operatorname{pts}(\operatorname{rm\text{-}kms}(s, q)), \operatorname{clopen\text{-}sets}(\operatorname{rm\text{-}kms}(s, q)), \\ \quad \quad \quad \quad \quad (\lambda \mathit{rcu} \in \operatorname{clopen\text{-}sets}(\operatorname{rm\text{-}kms}(s, q)).\; m(\operatorname{indicator}(\operatorname{pts}(\operatorname{rm\text{-}kms}(s, q)), \mathit{rcu}))))), \\ \quad \quad \quad \operatorname{res}(\operatorname{method\text{-}i\text{-}outer}(\operatorname{pts}(\operatorname{rm\text{-}kms}(s, q)), \operatorname{clopen\text{-}sets}(\operatorname{rm\text{-}kms}(s, q)), \\ \quad \quad \quad \quad \quad (\lambda \mathit{rcu} \in \operatorname{clopen\text{-}sets}(\operatorname{rm\text{-}kms}(s, q)).\; m(\operatorname{indicator}(\operatorname{pts}(\operatorname{rm\text{-}kms}(s, q)), \mathit{rcu})))), \\ \quad \quad \quad \quad \operatorname{caratheodory\text{-}sets}(\operatorname{pts}(\operatorname{rm\text{-}kms}(s, q)), \\ \quad \quad \quad \quad \quad \operatorname{method\text{-}i\text{-}outer}(\operatorname{pts}(\operatorname{rm\text{-}kms}(s, q)), \operatorname{clopen\text{-}sets}(\operatorname{rm\text{-}kms}(s, q)), \\ \quad \quad \quad \quad \quad \quad (\lambda \mathit{rcu} \in \operatorname{clopen\text{-}sets}(\operatorname{rm\text{-}kms}(s, q)).\; m(\operatorname{indicator}(\operatorname{pts}(\operatorname{rm\text{-}kms}(s, q)), \mathit{rcu}))))))) \wedge \\ \quad \quad \operatorname{borel\text{-}sets}(\operatorname{rm\text{-}kms}(s, q)) \\ \quad \quad \subseteq\; \operatorname{caratheodory\text{-}sets}(\operatorname{pts}(\operatorname{rm\text{-}kms}(s, q)), \\ \quad \quad \quad \operatorname{method\text{-}i\text{-}outer}(\operatorname{pts}(\operatorname{rm\text{-}kms}(s, q)), \operatorname{clopen\text{-}sets}(\operatorname{rm\text{-}kms}(s, q)), \\ \quad \quad \quad \quad (\lambda \mathit{rcu} \in \operatorname{clopen\text{-}sets}(\operatorname{rm\text{-}kms}(s, q)).\; m(\operatorname{indicator}(\operatorname{pts}(\operatorname{rm\text{-}kms}(s, q)), \mathit{rcu}))))) \wedge \\ \quad \quad \operatorname{clopen\text{-}sets}(\operatorname{rm\text{-}kms}(s, q)) \\ \quad \quad \subseteq\; \operatorname{caratheodory\text{-}sets}(\operatorname{pts}(\operatorname{rm\text{-}kms}(s, q)), \\ \quad \quad \quad \operatorname{method\text{-}i\text{-}outer}(\operatorname{pts}(\operatorname{rm\text{-}kms}(s, q)), \operatorname{clopen\text{-}sets}(\operatorname{rm\text{-}kms}(s, q)), \\ \quad \quad \quad \quad (\lambda \mathit{rcu} \in \operatorname{clopen\text{-}sets}(\operatorname{rm\text{-}kms}(s, q)).\; m(\operatorname{indicator}(\operatorname{pts}(\operatorname{rm\text{-}kms}(s, q)), \mathit{rcu}))))) \wedge \\ \quad \quad \forall u \in \operatorname{clopen\text{-}sets}(\operatorname{rm\text{-}kms}(s, q)). \\ \quad \quad \quad \operatorname{res}(\operatorname{method\text{-}i\text{-}outer}(\operatorname{pts}(\operatorname{rm\text{-}kms}(s, q)), \operatorname{clopen\text{-}sets}(\operatorname{rm\text{-}kms}(s, q)), (\lambda \mathit{rcu} \in \operatorname{clopen\text{-}sets}(\operatorname{rm\text{-}kms}(s, q)).\; m(\operatorname{indicator}(\operatorname{pts}(\operatorname{rm\text{-}kms}(s, q)), \mathit{rcu})))), \operatorname{caratheodory\text{-}sets}(\operatorname{pts}(\operatorname{rm\text{-}kms}(s, q)), \operatorname{method\text{-}i\text{-}outer}(\operatorname{pts}(\operatorname{rm\text{-}kms}(s, q)), \operatorname{clopen\text{-}sets}(\operatorname{rm\text{-}kms}(s, q)), (\lambda \mathit{rcu} \in \operatorname{clopen\text{-}sets}(\operatorname{rm\text{-}kms}(s, q)).\; m(\operatorname{indicator}(\operatorname{pts}(\operatorname{rm\text{-}kms}(s, q)), \mathit{rcu}))))))(u) \\ \quad \quad \quad =\; m(\operatorname{indicator}(\operatorname{pts}(\operatorname{rm\text{-}kms}(s, q)), u)) \end{array}}}
\vnbprop{rmc-lt-set-open}{\vnbstmt{\begin{array}{@{}l@{}} \forall m, g, c. \\ \quad ((g \in \operatorname{cont\text{-}fns}(m)) \wedge (c \in \mathbb{R})) \Rightarrow \operatorname{is\text{-}open}(m, \{x \in \operatorname{pts}(m) : (g(x) < c)\}) \end{array}}}
\vnbprop{rmc-gt-set-open}{\vnbstmt{\begin{array}{@{}l@{}} \forall m, g, c. \\ \quad ((g \in \operatorname{cont\text{-}fns}(m)) \wedge (c \in \mathbb{R})) \Rightarrow \operatorname{is\text{-}open}(m, \{x \in \operatorname{pts}(m) : (c < g(x))\}) \end{array}}}
\vnbprop{rmc-cont-bounded}{\vnbstmt{\begin{array}{@{}l@{}} \forall m, g. \\ \quad (\operatorname{is\text{-}compact}(m) \wedge (g \in \operatorname{cont\text{-}fns}(m))) \Rightarrow  \\ \quad \quad \exists d \in \mathbb{R}.\; ((0 \leq d) \wedge \forall x \in \operatorname{pts}(m).\; (\lvert g(x) \rvert \leq d)) \end{array}}}
\vnbprop{rmc-cyl-mono}{\vnbstmt{\begin{array}{@{}l@{}} \forall s, q, e, n, b. \\ \quad ((n \in \mathbb{N}) \wedge ((b \in \mathbb{N}) \wedge (n \leq b))) \Rightarrow  \\ \quad \quad \operatorname{rm\text{-}cyl}(s, q, e, b) \subseteq \operatorname{rm\text{-}cyl}(s, q, e, n) \end{array}}}
\vnbprop{rmc-uniform-level}{\vnbstmt{\begin{array}{@{}l@{}} \forall s, q, g, a. \\ \quad ((g \in \operatorname{cont\text{-}fns}(\operatorname{rm\text{-}kms}(s, q))) \wedge \operatorname{pos\text{-}rr}(a)) \Rightarrow  \\ \quad \quad \exists l \in \mathbb{N}. \forall x \in \operatorname{rm\text{-}k}(s, q), y \in \operatorname{rm\text{-}cyl}(s, q, x, l). \\ \quad \quad \quad \lvert (g(y) - g(x)) \rvert \leq a \end{array}}}
\vnbprop{rmc-bounded-integrable}{\vnbstmt{\begin{array}{@{}l@{}} \forall a, b, u, f, d. \\ \quad \operatorname{is\text{-}measure}(a, b, u) \wedge \\ \quad u(a) \in \mathbb{R} \wedge \\ \quad f \in (a \to \mathbb{R}) \wedge \\ \quad \forall l \in \mathbb{R}.\; (\{c \in a : (l < f(c))\} \in b) \wedge \\ \quad d \in \mathbb{R} \wedge \\ \quad 0 \leq d \wedge \\ \quad \forall x \in a.\; (\lvert f(x) \rvert \leq d) \Rightarrow \\ \quad \quad \operatorname{is\text{-}integrable}(a, b, u, f) \end{array}}}
\vnbprop{rmc-good-ind}{\vnbstmt{\begin{array}{@{}l@{}} \forall s, q, m, u. \\ \quad (\operatorname{is\text{-}positive\text{-}functional}(\operatorname{rm\text{-}kms}(s, q), m) \wedge (u \in \operatorname{clopen\text{-}sets}(\operatorname{rm\text{-}kms}(s, q)))) \Rightarrow  \\ \quad \quad \operatorname{indicator}(\operatorname{pts}(\operatorname{rm\text{-}kms}(s, q)), u) \in \operatorname{cont\text{-}fns}(\operatorname{rm\text{-}kms}(s, q)) \wedge \\ \quad \quad \operatorname{is\text{-}integrable}(\operatorname{pts}(\operatorname{rm\text{-}kms}(s, q)), \\ \quad \quad \quad \operatorname{caratheodory\text{-}sets}(\operatorname{pts}(\operatorname{rm\text{-}kms}(s, q)), \\ \quad \quad \quad \quad \operatorname{method\text{-}i\text{-}outer}(\operatorname{pts}(\operatorname{rm\text{-}kms}(s, q)), \operatorname{clopen\text{-}sets}(\operatorname{rm\text{-}kms}(s, q)), \\ \quad \quad \quad \quad \quad (\lambda \mathit{rcu} \in \operatorname{clopen\text{-}sets}(\operatorname{rm\text{-}kms}(s, q)).\; m(\operatorname{indicator}(\operatorname{pts}(\operatorname{rm\text{-}kms}(s, q)), \mathit{rcu}))))), \\ \quad \quad \quad \operatorname{res}(\operatorname{method\text{-}i\text{-}outer}(\operatorname{pts}(\operatorname{rm\text{-}kms}(s, q)), \operatorname{clopen\text{-}sets}(\operatorname{rm\text{-}kms}(s, q)), \\ \quad \quad \quad \quad \quad (\lambda \mathit{rcu} \in \operatorname{clopen\text{-}sets}(\operatorname{rm\text{-}kms}(s, q)).\; m(\operatorname{indicator}(\operatorname{pts}(\operatorname{rm\text{-}kms}(s, q)), \mathit{rcu})))), \\ \quad \quad \quad \quad \operatorname{caratheodory\text{-}sets}(\operatorname{pts}(\operatorname{rm\text{-}kms}(s, q)), \\ \quad \quad \quad \quad \quad \operatorname{method\text{-}i\text{-}outer}(\operatorname{pts}(\operatorname{rm\text{-}kms}(s, q)), \operatorname{clopen\text{-}sets}(\operatorname{rm\text{-}kms}(s, q)), \\ \quad \quad \quad \quad \quad \quad (\lambda \mathit{rcu} \in \operatorname{clopen\text{-}sets}(\operatorname{rm\text{-}kms}(s, q)).\; m(\operatorname{indicator}(\operatorname{pts}(\operatorname{rm\text{-}kms}(s, q)), \mathit{rcu})))))), \\ \quad \quad \quad \operatorname{indicator}(\operatorname{pts}(\operatorname{rm\text{-}kms}(s, q)), u)) \wedge \\ \quad \quad \operatorname{rintegral}(\operatorname{pts}(\operatorname{rm\text{-}kms}(s, q)), \\ \quad \quad \quad \operatorname{caratheodory\text{-}sets}(\operatorname{pts}(\operatorname{rm\text{-}kms}(s, q)), \\ \quad \quad \quad \quad \operatorname{method\text{-}i\text{-}outer}(\operatorname{pts}(\operatorname{rm\text{-}kms}(s, q)), \operatorname{clopen\text{-}sets}(\operatorname{rm\text{-}kms}(s, q)), \\ \quad \quad \quad \quad \quad (\lambda \mathit{rcu} \in \operatorname{clopen\text{-}sets}(\operatorname{rm\text{-}kms}(s, q)).\; m(\operatorname{indicator}(\operatorname{pts}(\operatorname{rm\text{-}kms}(s, q)), \mathit{rcu}))))), \\ \quad \quad \quad \operatorname{res}(\operatorname{method\text{-}i\text{-}outer}(\operatorname{pts}(\operatorname{rm\text{-}kms}(s, q)), \operatorname{clopen\text{-}sets}(\operatorname{rm\text{-}kms}(s, q)), \\ \quad \quad \quad \quad \quad (\lambda \mathit{rcu} \in \operatorname{clopen\text{-}sets}(\operatorname{rm\text{-}kms}(s, q)).\; m(\operatorname{indicator}(\operatorname{pts}(\operatorname{rm\text{-}kms}(s, q)), \mathit{rcu})))), \\ \quad \quad \quad \quad \operatorname{caratheodory\text{-}sets}(\operatorname{pts}(\operatorname{rm\text{-}kms}(s, q)), \\ \quad \quad \quad \quad \quad \operatorname{method\text{-}i\text{-}outer}(\operatorname{pts}(\operatorname{rm\text{-}kms}(s, q)), \operatorname{clopen\text{-}sets}(\operatorname{rm\text{-}kms}(s, q)), \\ \quad \quad \quad \quad \quad \quad (\lambda \mathit{rcu} \in \operatorname{clopen\text{-}sets}(\operatorname{rm\text{-}kms}(s, q)).\; m(\operatorname{indicator}(\operatorname{pts}(\operatorname{rm\text{-}kms}(s, q)), \mathit{rcu})))))), \\ \quad \quad \quad \operatorname{indicator}(\operatorname{pts}(\operatorname{rm\text{-}kms}(s, q)), u)) \\ \quad \quad =\; m(\operatorname{indicator}(\operatorname{pts}(\operatorname{rm\text{-}kms}(s, q)), u)) \end{array}}}
\vnbprop{rmc-good-scale}{\vnbstmt{\begin{array}{@{}l@{}} \forall s, q, m, c, f. \\ \quad \operatorname{is\text{-}positive\text{-}functional}(\operatorname{rm\text{-}kms}(s, q), m) \wedge \\ \quad c \in \mathbb{R} \wedge \\ \quad f \in \operatorname{cont\text{-}fns}(\operatorname{rm\text{-}kms}(s, q)) \wedge \\ \quad \operatorname{is\text{-}integrable}(\operatorname{pts}(\operatorname{rm\text{-}kms}(s, q)), \\ \quad \quad \operatorname{caratheodory\text{-}sets}(\operatorname{pts}(\operatorname{rm\text{-}kms}(s, q)), \\ \quad \quad \quad \operatorname{method\text{-}i\text{-}outer}(\operatorname{pts}(\operatorname{rm\text{-}kms}(s, q)), \operatorname{clopen\text{-}sets}(\operatorname{rm\text{-}kms}(s, q)), \\ \quad \quad \quad \quad (\lambda u \in \operatorname{clopen\text{-}sets}(\operatorname{rm\text{-}kms}(s, q)).\; m(\operatorname{indicator}(\operatorname{pts}(\operatorname{rm\text{-}kms}(s, q)), u))))), \\ \quad \quad \operatorname{res}(\operatorname{method\text{-}i\text{-}outer}(\operatorname{pts}(\operatorname{rm\text{-}kms}(s, q)), \operatorname{clopen\text{-}sets}(\operatorname{rm\text{-}kms}(s, q)), \\ \quad \quad \quad \quad (\lambda u \in \operatorname{clopen\text{-}sets}(\operatorname{rm\text{-}kms}(s, q)).\; m(\operatorname{indicator}(\operatorname{pts}(\operatorname{rm\text{-}kms}(s, q)), u)))), \\ \quad \quad \quad \operatorname{caratheodory\text{-}sets}(\operatorname{pts}(\operatorname{rm\text{-}kms}(s, q)), \\ \quad \quad \quad \quad \operatorname{method\text{-}i\text{-}outer}(\operatorname{pts}(\operatorname{rm\text{-}kms}(s, q)), \operatorname{clopen\text{-}sets}(\operatorname{rm\text{-}kms}(s, q)), \\ \quad \quad \quad \quad \quad (\lambda u \in \operatorname{clopen\text{-}sets}(\operatorname{rm\text{-}kms}(s, q)).\; m(\operatorname{indicator}(\operatorname{pts}(\operatorname{rm\text{-}kms}(s, q)), u)))))), \\ \quad \quad f) \wedge \\ \quad \operatorname{rintegral}(\operatorname{pts}(\operatorname{rm\text{-}kms}(s, q)), \\ \quad \quad \operatorname{caratheodory\text{-}sets}(\operatorname{pts}(\operatorname{rm\text{-}kms}(s, q)), \\ \quad \quad \quad \operatorname{method\text{-}i\text{-}outer}(\operatorname{pts}(\operatorname{rm\text{-}kms}(s, q)), \operatorname{clopen\text{-}sets}(\operatorname{rm\text{-}kms}(s, q)), \\ \quad \quad \quad \quad (\lambda u \in \operatorname{clopen\text{-}sets}(\operatorname{rm\text{-}kms}(s, q)).\; m(\operatorname{indicator}(\operatorname{pts}(\operatorname{rm\text{-}kms}(s, q)), u))))), \\ \quad \quad \operatorname{res}(\operatorname{method\text{-}i\text{-}outer}(\operatorname{pts}(\operatorname{rm\text{-}kms}(s, q)), \operatorname{clopen\text{-}sets}(\operatorname{rm\text{-}kms}(s, q)), \\ \quad \quad \quad \quad (\lambda u \in \operatorname{clopen\text{-}sets}(\operatorname{rm\text{-}kms}(s, q)).\; m(\operatorname{indicator}(\operatorname{pts}(\operatorname{rm\text{-}kms}(s, q)), u)))), \\ \quad \quad \quad \operatorname{caratheodory\text{-}sets}(\operatorname{pts}(\operatorname{rm\text{-}kms}(s, q)), \\ \quad \quad \quad \quad \operatorname{method\text{-}i\text{-}outer}(\operatorname{pts}(\operatorname{rm\text{-}kms}(s, q)), \operatorname{clopen\text{-}sets}(\operatorname{rm\text{-}kms}(s, q)), \\ \quad \quad \quad \quad \quad (\lambda u \in \operatorname{clopen\text{-}sets}(\operatorname{rm\text{-}kms}(s, q)).\; m(\operatorname{indicator}(\operatorname{pts}(\operatorname{rm\text{-}kms}(s, q)), u)))))), \\ \quad \quad f) \\ \quad =\; m(f) \Rightarrow \\ \quad \quad (\lambda x \in \operatorname{pts}(\operatorname{rm\text{-}kms}(s, q)).\; (c \cdot f(x))) \in \operatorname{cont\text{-}fns}(\operatorname{rm\text{-}kms}(s, q)) \wedge \\ \quad \quad \operatorname{is\text{-}integrable}(\operatorname{pts}(\operatorname{rm\text{-}kms}(s, q)), \\ \quad \quad \quad \operatorname{caratheodory\text{-}sets}(\operatorname{pts}(\operatorname{rm\text{-}kms}(s, q)), \\ \quad \quad \quad \quad \operatorname{method\text{-}i\text{-}outer}(\operatorname{pts}(\operatorname{rm\text{-}kms}(s, q)), \operatorname{clopen\text{-}sets}(\operatorname{rm\text{-}kms}(s, q)), \\ \quad \quad \quad \quad \quad (\lambda u \in \operatorname{clopen\text{-}sets}(\operatorname{rm\text{-}kms}(s, q)).\; m(\operatorname{indicator}(\operatorname{pts}(\operatorname{rm\text{-}kms}(s, q)), u))))), \\ \quad \quad \quad \operatorname{res}(\operatorname{method\text{-}i\text{-}outer}(\operatorname{pts}(\operatorname{rm\text{-}kms}(s, q)), \operatorname{clopen\text{-}sets}(\operatorname{rm\text{-}kms}(s, q)), \\ \quad \quad \quad \quad \quad (\lambda u \in \operatorname{clopen\text{-}sets}(\operatorname{rm\text{-}kms}(s, q)).\; m(\operatorname{indicator}(\operatorname{pts}(\operatorname{rm\text{-}kms}(s, q)), u)))), \\ \quad \quad \quad \quad \operatorname{caratheodory\text{-}sets}(\operatorname{pts}(\operatorname{rm\text{-}kms}(s, q)), \\ \quad \quad \quad \quad \quad \operatorname{method\text{-}i\text{-}outer}(\operatorname{pts}(\operatorname{rm\text{-}kms}(s, q)), \operatorname{clopen\text{-}sets}(\operatorname{rm\text{-}kms}(s, q)), \\ \quad \quad \quad \quad \quad \quad (\lambda u \in \operatorname{clopen\text{-}sets}(\operatorname{rm\text{-}kms}(s, q)).\; m(\operatorname{indicator}(\operatorname{pts}(\operatorname{rm\text{-}kms}(s, q)), u)))))), \\ \quad \quad \quad (\lambda x \in \operatorname{pts}(\operatorname{rm\text{-}kms}(s, q)).\; (c \cdot f(x)))) \wedge \\ \quad \quad \operatorname{rintegral}(\operatorname{pts}(\operatorname{rm\text{-}kms}(s, q)), \\ \quad \quad \quad \operatorname{caratheodory\text{-}sets}(\operatorname{pts}(\operatorname{rm\text{-}kms}(s, q)), \\ \quad \quad \quad \quad \operatorname{method\text{-}i\text{-}outer}(\operatorname{pts}(\operatorname{rm\text{-}kms}(s, q)), \operatorname{clopen\text{-}sets}(\operatorname{rm\text{-}kms}(s, q)), \\ \quad \quad \quad \quad \quad (\lambda u \in \operatorname{clopen\text{-}sets}(\operatorname{rm\text{-}kms}(s, q)).\; m(\operatorname{indicator}(\operatorname{pts}(\operatorname{rm\text{-}kms}(s, q)), u))))), \\ \quad \quad \quad \operatorname{res}(\operatorname{method\text{-}i\text{-}outer}(\operatorname{pts}(\operatorname{rm\text{-}kms}(s, q)), \operatorname{clopen\text{-}sets}(\operatorname{rm\text{-}kms}(s, q)), \\ \quad \quad \quad \quad \quad (\lambda u \in \operatorname{clopen\text{-}sets}(\operatorname{rm\text{-}kms}(s, q)).\; m(\operatorname{indicator}(\operatorname{pts}(\operatorname{rm\text{-}kms}(s, q)), u)))), \\ \quad \quad \quad \quad \operatorname{caratheodory\text{-}sets}(\operatorname{pts}(\operatorname{rm\text{-}kms}(s, q)), \\ \quad \quad \quad \quad \quad \operatorname{method\text{-}i\text{-}outer}(\operatorname{pts}(\operatorname{rm\text{-}kms}(s, q)), \operatorname{clopen\text{-}sets}(\operatorname{rm\text{-}kms}(s, q)), \\ \quad \quad \quad \quad \quad \quad (\lambda u \in \operatorname{clopen\text{-}sets}(\operatorname{rm\text{-}kms}(s, q)).\; m(\operatorname{indicator}(\operatorname{pts}(\operatorname{rm\text{-}kms}(s, q)), u)))))), \\ \quad \quad \quad (\lambda x \in \operatorname{pts}(\operatorname{rm\text{-}kms}(s, q)).\; (c \cdot f(x)))) \\ \quad \quad =\; m((\lambda x \in \operatorname{pts}(\operatorname{rm\text{-}kms}(s, q)).\; (c \cdot f(x)))) \end{array}}}
\vnbprop{rmc-good-add}{\vnbstmt{\begin{array}{@{}l@{}} \forall s, q, m, f, g. \\ \quad \operatorname{is\text{-}positive\text{-}functional}(\operatorname{rm\text{-}kms}(s, q), m) \wedge \\ \quad f \in \operatorname{cont\text{-}fns}(\operatorname{rm\text{-}kms}(s, q)) \wedge \\ \quad \operatorname{is\text{-}integrable}(\operatorname{pts}(\operatorname{rm\text{-}kms}(s, q)), \\ \quad \quad \operatorname{caratheodory\text{-}sets}(\operatorname{pts}(\operatorname{rm\text{-}kms}(s, q)), \\ \quad \quad \quad \operatorname{method\text{-}i\text{-}outer}(\operatorname{pts}(\operatorname{rm\text{-}kms}(s, q)), \operatorname{clopen\text{-}sets}(\operatorname{rm\text{-}kms}(s, q)), \\ \quad \quad \quad \quad (\lambda u \in \operatorname{clopen\text{-}sets}(\operatorname{rm\text{-}kms}(s, q)).\; m(\operatorname{indicator}(\operatorname{pts}(\operatorname{rm\text{-}kms}(s, q)), u))))), \\ \quad \quad \operatorname{res}(\operatorname{method\text{-}i\text{-}outer}(\operatorname{pts}(\operatorname{rm\text{-}kms}(s, q)), \operatorname{clopen\text{-}sets}(\operatorname{rm\text{-}kms}(s, q)), \\ \quad \quad \quad \quad (\lambda u \in \operatorname{clopen\text{-}sets}(\operatorname{rm\text{-}kms}(s, q)).\; m(\operatorname{indicator}(\operatorname{pts}(\operatorname{rm\text{-}kms}(s, q)), u)))), \\ \quad \quad \quad \operatorname{caratheodory\text{-}sets}(\operatorname{pts}(\operatorname{rm\text{-}kms}(s, q)), \\ \quad \quad \quad \quad \operatorname{method\text{-}i\text{-}outer}(\operatorname{pts}(\operatorname{rm\text{-}kms}(s, q)), \operatorname{clopen\text{-}sets}(\operatorname{rm\text{-}kms}(s, q)), \\ \quad \quad \quad \quad \quad (\lambda u \in \operatorname{clopen\text{-}sets}(\operatorname{rm\text{-}kms}(s, q)).\; m(\operatorname{indicator}(\operatorname{pts}(\operatorname{rm\text{-}kms}(s, q)), u)))))), \\ \quad \quad f) \wedge \\ \quad \operatorname{rintegral}(\operatorname{pts}(\operatorname{rm\text{-}kms}(s, q)), \\ \quad \quad \operatorname{caratheodory\text{-}sets}(\operatorname{pts}(\operatorname{rm\text{-}kms}(s, q)), \\ \quad \quad \quad \operatorname{method\text{-}i\text{-}outer}(\operatorname{pts}(\operatorname{rm\text{-}kms}(s, q)), \operatorname{clopen\text{-}sets}(\operatorname{rm\text{-}kms}(s, q)), \\ \quad \quad \quad \quad (\lambda u \in \operatorname{clopen\text{-}sets}(\operatorname{rm\text{-}kms}(s, q)).\; m(\operatorname{indicator}(\operatorname{pts}(\operatorname{rm\text{-}kms}(s, q)), u))))), \\ \quad \quad \operatorname{res}(\operatorname{method\text{-}i\text{-}outer}(\operatorname{pts}(\operatorname{rm\text{-}kms}(s, q)), \operatorname{clopen\text{-}sets}(\operatorname{rm\text{-}kms}(s, q)), \\ \quad \quad \quad \quad (\lambda u \in \operatorname{clopen\text{-}sets}(\operatorname{rm\text{-}kms}(s, q)).\; m(\operatorname{indicator}(\operatorname{pts}(\operatorname{rm\text{-}kms}(s, q)), u)))), \\ \quad \quad \quad \operatorname{caratheodory\text{-}sets}(\operatorname{pts}(\operatorname{rm\text{-}kms}(s, q)), \\ \quad \quad \quad \quad \operatorname{method\text{-}i\text{-}outer}(\operatorname{pts}(\operatorname{rm\text{-}kms}(s, q)), \operatorname{clopen\text{-}sets}(\operatorname{rm\text{-}kms}(s, q)), \\ \quad \quad \quad \quad \quad (\lambda u \in \operatorname{clopen\text{-}sets}(\operatorname{rm\text{-}kms}(s, q)).\; m(\operatorname{indicator}(\operatorname{pts}(\operatorname{rm\text{-}kms}(s, q)), u)))))), \\ \quad \quad f) \\ \quad =\; m(f) \wedge \\ \quad g \in \operatorname{cont\text{-}fns}(\operatorname{rm\text{-}kms}(s, q)) \wedge \\ \quad \operatorname{is\text{-}integrable}(\operatorname{pts}(\operatorname{rm\text{-}kms}(s, q)), \\ \quad \quad \operatorname{caratheodory\text{-}sets}(\operatorname{pts}(\operatorname{rm\text{-}kms}(s, q)), \\ \quad \quad \quad \operatorname{method\text{-}i\text{-}outer}(\operatorname{pts}(\operatorname{rm\text{-}kms}(s, q)), \operatorname{clopen\text{-}sets}(\operatorname{rm\text{-}kms}(s, q)), \\ \quad \quad \quad \quad (\lambda u \in \operatorname{clopen\text{-}sets}(\operatorname{rm\text{-}kms}(s, q)).\; m(\operatorname{indicator}(\operatorname{pts}(\operatorname{rm\text{-}kms}(s, q)), u))))), \\ \quad \quad \operatorname{res}(\operatorname{method\text{-}i\text{-}outer}(\operatorname{pts}(\operatorname{rm\text{-}kms}(s, q)), \operatorname{clopen\text{-}sets}(\operatorname{rm\text{-}kms}(s, q)), \\ \quad \quad \quad \quad (\lambda u \in \operatorname{clopen\text{-}sets}(\operatorname{rm\text{-}kms}(s, q)).\; m(\operatorname{indicator}(\operatorname{pts}(\operatorname{rm\text{-}kms}(s, q)), u)))), \\ \quad \quad \quad \operatorname{caratheodory\text{-}sets}(\operatorname{pts}(\operatorname{rm\text{-}kms}(s, q)), \\ \quad \quad \quad \quad \operatorname{method\text{-}i\text{-}outer}(\operatorname{pts}(\operatorname{rm\text{-}kms}(s, q)), \operatorname{clopen\text{-}sets}(\operatorname{rm\text{-}kms}(s, q)), \\ \quad \quad \quad \quad \quad (\lambda u \in \operatorname{clopen\text{-}sets}(\operatorname{rm\text{-}kms}(s, q)).\; m(\operatorname{indicator}(\operatorname{pts}(\operatorname{rm\text{-}kms}(s, q)), u)))))), \\ \quad \quad g) \wedge \\ \quad \operatorname{rintegral}(\operatorname{pts}(\operatorname{rm\text{-}kms}(s, q)), \\ \quad \quad \operatorname{caratheodory\text{-}sets}(\operatorname{pts}(\operatorname{rm\text{-}kms}(s, q)), \\ \quad \quad \quad \operatorname{method\text{-}i\text{-}outer}(\operatorname{pts}(\operatorname{rm\text{-}kms}(s, q)), \operatorname{clopen\text{-}sets}(\operatorname{rm\text{-}kms}(s, q)), \\ \quad \quad \quad \quad (\lambda u \in \operatorname{clopen\text{-}sets}(\operatorname{rm\text{-}kms}(s, q)).\; m(\operatorname{indicator}(\operatorname{pts}(\operatorname{rm\text{-}kms}(s, q)), u))))), \\ \quad \quad \operatorname{res}(\operatorname{method\text{-}i\text{-}outer}(\operatorname{pts}(\operatorname{rm\text{-}kms}(s, q)), \operatorname{clopen\text{-}sets}(\operatorname{rm\text{-}kms}(s, q)), \\ \quad \quad \quad \quad (\lambda u \in \operatorname{clopen\text{-}sets}(\operatorname{rm\text{-}kms}(s, q)).\; m(\operatorname{indicator}(\operatorname{pts}(\operatorname{rm\text{-}kms}(s, q)), u)))), \\ \quad \quad \quad \operatorname{caratheodory\text{-}sets}(\operatorname{pts}(\operatorname{rm\text{-}kms}(s, q)), \\ \quad \quad \quad \quad \operatorname{method\text{-}i\text{-}outer}(\operatorname{pts}(\operatorname{rm\text{-}kms}(s, q)), \operatorname{clopen\text{-}sets}(\operatorname{rm\text{-}kms}(s, q)), \\ \quad \quad \quad \quad \quad (\lambda u \in \operatorname{clopen\text{-}sets}(\operatorname{rm\text{-}kms}(s, q)).\; m(\operatorname{indicator}(\operatorname{pts}(\operatorname{rm\text{-}kms}(s, q)), u)))))), \\ \quad \quad g) \\ \quad =\; m(g) \Rightarrow \\ \quad \quad (\lambda x \in \operatorname{pts}(\operatorname{rm\text{-}kms}(s, q)).\; (f(x) + g(x))) \in \operatorname{cont\text{-}fns}(\operatorname{rm\text{-}kms}(s, q)) \wedge \\ \quad \quad \operatorname{is\text{-}integrable}(\operatorname{pts}(\operatorname{rm\text{-}kms}(s, q)), \\ \quad \quad \quad \operatorname{caratheodory\text{-}sets}(\operatorname{pts}(\operatorname{rm\text{-}kms}(s, q)), \\ \quad \quad \quad \quad \operatorname{method\text{-}i\text{-}outer}(\operatorname{pts}(\operatorname{rm\text{-}kms}(s, q)), \operatorname{clopen\text{-}sets}(\operatorname{rm\text{-}kms}(s, q)), \\ \quad \quad \quad \quad \quad (\lambda u \in \operatorname{clopen\text{-}sets}(\operatorname{rm\text{-}kms}(s, q)).\; m(\operatorname{indicator}(\operatorname{pts}(\operatorname{rm\text{-}kms}(s, q)), u))))), \\ \quad \quad \quad \operatorname{res}(\operatorname{method\text{-}i\text{-}outer}(\operatorname{pts}(\operatorname{rm\text{-}kms}(s, q)), \operatorname{clopen\text{-}sets}(\operatorname{rm\text{-}kms}(s, q)), \\ \quad \quad \quad \quad \quad (\lambda u \in \operatorname{clopen\text{-}sets}(\operatorname{rm\text{-}kms}(s, q)).\; m(\operatorname{indicator}(\operatorname{pts}(\operatorname{rm\text{-}kms}(s, q)), u)))), \\ \quad \quad \quad \quad \operatorname{caratheodory\text{-}sets}(\operatorname{pts}(\operatorname{rm\text{-}kms}(s, q)), \\ \quad \quad \quad \quad \quad \operatorname{method\text{-}i\text{-}outer}(\operatorname{pts}(\operatorname{rm\text{-}kms}(s, q)), \operatorname{clopen\text{-}sets}(\operatorname{rm\text{-}kms}(s, q)), \\ \quad \quad \quad \quad \quad \quad (\lambda u \in \operatorname{clopen\text{-}sets}(\operatorname{rm\text{-}kms}(s, q)).\; m(\operatorname{indicator}(\operatorname{pts}(\operatorname{rm\text{-}kms}(s, q)), u)))))), \\ \quad \quad \quad (\lambda x \in \operatorname{pts}(\operatorname{rm\text{-}kms}(s, q)).\; (f(x) + g(x)))) \wedge \\ \quad \quad \operatorname{rintegral}(\operatorname{pts}(\operatorname{rm\text{-}kms}(s, q)), \\ \quad \quad \quad \operatorname{caratheodory\text{-}sets}(\operatorname{pts}(\operatorname{rm\text{-}kms}(s, q)), \\ \quad \quad \quad \quad \operatorname{method\text{-}i\text{-}outer}(\operatorname{pts}(\operatorname{rm\text{-}kms}(s, q)), \operatorname{clopen\text{-}sets}(\operatorname{rm\text{-}kms}(s, q)), \\ \quad \quad \quad \quad \quad (\lambda u \in \operatorname{clopen\text{-}sets}(\operatorname{rm\text{-}kms}(s, q)).\; m(\operatorname{indicator}(\operatorname{pts}(\operatorname{rm\text{-}kms}(s, q)), u))))), \\ \quad \quad \quad \operatorname{res}(\operatorname{method\text{-}i\text{-}outer}(\operatorname{pts}(\operatorname{rm\text{-}kms}(s, q)), \operatorname{clopen\text{-}sets}(\operatorname{rm\text{-}kms}(s, q)), \\ \quad \quad \quad \quad \quad (\lambda u \in \operatorname{clopen\text{-}sets}(\operatorname{rm\text{-}kms}(s, q)).\; m(\operatorname{indicator}(\operatorname{pts}(\operatorname{rm\text{-}kms}(s, q)), u)))), \\ \quad \quad \quad \quad \operatorname{caratheodory\text{-}sets}(\operatorname{pts}(\operatorname{rm\text{-}kms}(s, q)), \\ \quad \quad \quad \quad \quad \operatorname{method\text{-}i\text{-}outer}(\operatorname{pts}(\operatorname{rm\text{-}kms}(s, q)), \operatorname{clopen\text{-}sets}(\operatorname{rm\text{-}kms}(s, q)), \\ \quad \quad \quad \quad \quad \quad (\lambda u \in \operatorname{clopen\text{-}sets}(\operatorname{rm\text{-}kms}(s, q)).\; m(\operatorname{indicator}(\operatorname{pts}(\operatorname{rm\text{-}kms}(s, q)), u)))))), \\ \quad \quad \quad (\lambda x \in \operatorname{pts}(\operatorname{rm\text{-}kms}(s, q)).\; (f(x) + g(x)))) \\ \quad \quad =\; m((\lambda x \in \operatorname{pts}(\operatorname{rm\text{-}kms}(s, q)).\; (f(x) + g(x)))) \end{array}}}
\vnbprop{rmc-layer}{\vnbstmt{\begin{array}{@{}l@{}} \forall n \in \mathbb{N}, s, q, m, l, c, a, f. \\ \quad \operatorname{is\text{-}positive\text{-}functional}(\operatorname{rm\text{-}kms}(s, q), m) \wedge \\ \quad l \in \mathbb{N} \wedge \\ \quad c \in \mathbb{R} \wedge \\ \quad \operatorname{pos\text{-}rr}(a) \wedge \\ \quad f \in (\operatorname{pts}(\operatorname{rm\text{-}kms}(s, q)) \to \mathbb{R}) \wedge \\ \quad \forall x \in \operatorname{rm\text{-}k}(s, q), y \in \operatorname{rm\text{-}cyl}(s, q, x, l).\; (f(y) = f(x)) \wedge \\ \quad \forall x \in \operatorname{pts}(\operatorname{rm\text{-}kms}(s, q)).\; ((c \leq f(x)) \wedge (f(x) \leq (c + ((n + 1) \cdot a)))) \Rightarrow \\ \quad \quad \exists h. \\ \quad \quad \quad h \in \operatorname{cont\text{-}fns}(\operatorname{rm\text{-}kms}(s, q)) \wedge \\ \quad \quad \quad \operatorname{is\text{-}integrable}(\operatorname{pts}(\operatorname{rm\text{-}kms}(s, q)), \\ \quad \quad \quad \quad \operatorname{caratheodory\text{-}sets}(\operatorname{pts}(\operatorname{rm\text{-}kms}(s, q)), \\ \quad \quad \quad \quad \quad \operatorname{method\text{-}i\text{-}outer}(\operatorname{pts}(\operatorname{rm\text{-}kms}(s, q)), \operatorname{clopen\text{-}sets}(\operatorname{rm\text{-}kms}(s, q)), \\ \quad \quad \quad \quad \quad \quad (\lambda u \in \operatorname{clopen\text{-}sets}(\operatorname{rm\text{-}kms}(s, q)).\; m(\operatorname{indicator}(\operatorname{pts}(\operatorname{rm\text{-}kms}(s, q)), u))))), \\ \quad \quad \quad \quad \operatorname{res}(\operatorname{method\text{-}i\text{-}outer}(\operatorname{pts}(\operatorname{rm\text{-}kms}(s, q)), \operatorname{clopen\text{-}sets}(\operatorname{rm\text{-}kms}(s, q)), \\ \quad \quad \quad \quad \quad \quad (\lambda u \in \operatorname{clopen\text{-}sets}(\operatorname{rm\text{-}kms}(s, q)).\; m(\operatorname{indicator}(\operatorname{pts}(\operatorname{rm\text{-}kms}(s, q)), u)))), \\ \quad \quad \quad \quad \quad \operatorname{caratheodory\text{-}sets}(\operatorname{pts}(\operatorname{rm\text{-}kms}(s, q)), \\ \quad \quad \quad \quad \quad \quad \operatorname{method\text{-}i\text{-}outer}(\operatorname{pts}(\operatorname{rm\text{-}kms}(s, q)), \operatorname{clopen\text{-}sets}(\operatorname{rm\text{-}kms}(s, q)), \\ \quad \quad \quad \quad \quad \quad \quad (\lambda u \in \operatorname{clopen\text{-}sets}(\operatorname{rm\text{-}kms}(s, q)).\; m(\operatorname{indicator}(\operatorname{pts}(\operatorname{rm\text{-}kms}(s, q)), u)))))), \\ \quad \quad \quad \quad h) \wedge \\ \quad \quad \quad \operatorname{rintegral}(\operatorname{pts}(\operatorname{rm\text{-}kms}(s, q)), \\ \quad \quad \quad \quad \operatorname{caratheodory\text{-}sets}(\operatorname{pts}(\operatorname{rm\text{-}kms}(s, q)), \\ \quad \quad \quad \quad \quad \operatorname{method\text{-}i\text{-}outer}(\operatorname{pts}(\operatorname{rm\text{-}kms}(s, q)), \operatorname{clopen\text{-}sets}(\operatorname{rm\text{-}kms}(s, q)), \\ \quad \quad \quad \quad \quad \quad (\lambda u \in \operatorname{clopen\text{-}sets}(\operatorname{rm\text{-}kms}(s, q)).\; m(\operatorname{indicator}(\operatorname{pts}(\operatorname{rm\text{-}kms}(s, q)), u))))), \\ \quad \quad \quad \quad \operatorname{res}(\operatorname{method\text{-}i\text{-}outer}(\operatorname{pts}(\operatorname{rm\text{-}kms}(s, q)), \operatorname{clopen\text{-}sets}(\operatorname{rm\text{-}kms}(s, q)), \\ \quad \quad \quad \quad \quad \quad (\lambda u \in \operatorname{clopen\text{-}sets}(\operatorname{rm\text{-}kms}(s, q)).\; m(\operatorname{indicator}(\operatorname{pts}(\operatorname{rm\text{-}kms}(s, q)), u)))), \\ \quad \quad \quad \quad \quad \operatorname{caratheodory\text{-}sets}(\operatorname{pts}(\operatorname{rm\text{-}kms}(s, q)), \\ \quad \quad \quad \quad \quad \quad \operatorname{method\text{-}i\text{-}outer}(\operatorname{pts}(\operatorname{rm\text{-}kms}(s, q)), \operatorname{clopen\text{-}sets}(\operatorname{rm\text{-}kms}(s, q)), \\ \quad \quad \quad \quad \quad \quad \quad (\lambda u \in \operatorname{clopen\text{-}sets}(\operatorname{rm\text{-}kms}(s, q)).\; m(\operatorname{indicator}(\operatorname{pts}(\operatorname{rm\text{-}kms}(s, q)), u)))))), \\ \quad \quad \quad \quad h) \\ \quad \quad \quad =\; m(h) \wedge \\ \quad \quad \quad \forall x \in \operatorname{pts}(\operatorname{rm\text{-}kms}(s, q)).\; (\lvert (f(x) - h(x)) \rvert \leq a) \end{array}}}
\vnbprop{rmc-approx}{\vnbstmt{\begin{array}{@{}l@{}} \forall s, q, m, g, a. \\ \quad (\operatorname{is\text{-}positive\text{-}functional}(\operatorname{rm\text{-}kms}(s, q), m) \wedge ((g \in \operatorname{cont\text{-}fns}(\operatorname{rm\text{-}kms}(s, q))) \wedge \operatorname{pos\text{-}rr}(a))) \Rightarrow  \\ \quad \quad \exists h. \\ \quad \quad \quad h \in \operatorname{cont\text{-}fns}(\operatorname{rm\text{-}kms}(s, q)) \wedge \\ \quad \quad \quad \operatorname{is\text{-}integrable}(\operatorname{pts}(\operatorname{rm\text{-}kms}(s, q)), \\ \quad \quad \quad \quad \operatorname{caratheodory\text{-}sets}(\operatorname{pts}(\operatorname{rm\text{-}kms}(s, q)), \\ \quad \quad \quad \quad \quad \operatorname{method\text{-}i\text{-}outer}(\operatorname{pts}(\operatorname{rm\text{-}kms}(s, q)), \operatorname{clopen\text{-}sets}(\operatorname{rm\text{-}kms}(s, q)), \\ \quad \quad \quad \quad \quad \quad (\lambda u \in \operatorname{clopen\text{-}sets}(\operatorname{rm\text{-}kms}(s, q)).\; m(\operatorname{indicator}(\operatorname{pts}(\operatorname{rm\text{-}kms}(s, q)), u))))), \\ \quad \quad \quad \quad \operatorname{res}(\operatorname{method\text{-}i\text{-}outer}(\operatorname{pts}(\operatorname{rm\text{-}kms}(s, q)), \operatorname{clopen\text{-}sets}(\operatorname{rm\text{-}kms}(s, q)), \\ \quad \quad \quad \quad \quad \quad (\lambda u \in \operatorname{clopen\text{-}sets}(\operatorname{rm\text{-}kms}(s, q)).\; m(\operatorname{indicator}(\operatorname{pts}(\operatorname{rm\text{-}kms}(s, q)), u)))), \\ \quad \quad \quad \quad \quad \operatorname{caratheodory\text{-}sets}(\operatorname{pts}(\operatorname{rm\text{-}kms}(s, q)), \\ \quad \quad \quad \quad \quad \quad \operatorname{method\text{-}i\text{-}outer}(\operatorname{pts}(\operatorname{rm\text{-}kms}(s, q)), \operatorname{clopen\text{-}sets}(\operatorname{rm\text{-}kms}(s, q)), \\ \quad \quad \quad \quad \quad \quad \quad (\lambda u \in \operatorname{clopen\text{-}sets}(\operatorname{rm\text{-}kms}(s, q)).\; m(\operatorname{indicator}(\operatorname{pts}(\operatorname{rm\text{-}kms}(s, q)), u)))))), \\ \quad \quad \quad \quad h) \wedge \\ \quad \quad \quad \operatorname{rintegral}(\operatorname{pts}(\operatorname{rm\text{-}kms}(s, q)), \\ \quad \quad \quad \quad \operatorname{caratheodory\text{-}sets}(\operatorname{pts}(\operatorname{rm\text{-}kms}(s, q)), \\ \quad \quad \quad \quad \quad \operatorname{method\text{-}i\text{-}outer}(\operatorname{pts}(\operatorname{rm\text{-}kms}(s, q)), \operatorname{clopen\text{-}sets}(\operatorname{rm\text{-}kms}(s, q)), \\ \quad \quad \quad \quad \quad \quad (\lambda u \in \operatorname{clopen\text{-}sets}(\operatorname{rm\text{-}kms}(s, q)).\; m(\operatorname{indicator}(\operatorname{pts}(\operatorname{rm\text{-}kms}(s, q)), u))))), \\ \quad \quad \quad \quad \operatorname{res}(\operatorname{method\text{-}i\text{-}outer}(\operatorname{pts}(\operatorname{rm\text{-}kms}(s, q)), \operatorname{clopen\text{-}sets}(\operatorname{rm\text{-}kms}(s, q)), \\ \quad \quad \quad \quad \quad \quad (\lambda u \in \operatorname{clopen\text{-}sets}(\operatorname{rm\text{-}kms}(s, q)).\; m(\operatorname{indicator}(\operatorname{pts}(\operatorname{rm\text{-}kms}(s, q)), u)))), \\ \quad \quad \quad \quad \quad \operatorname{caratheodory\text{-}sets}(\operatorname{pts}(\operatorname{rm\text{-}kms}(s, q)), \\ \quad \quad \quad \quad \quad \quad \operatorname{method\text{-}i\text{-}outer}(\operatorname{pts}(\operatorname{rm\text{-}kms}(s, q)), \operatorname{clopen\text{-}sets}(\operatorname{rm\text{-}kms}(s, q)), \\ \quad \quad \quad \quad \quad \quad \quad (\lambda u \in \operatorname{clopen\text{-}sets}(\operatorname{rm\text{-}kms}(s, q)).\; m(\operatorname{indicator}(\operatorname{pts}(\operatorname{rm\text{-}kms}(s, q)), u)))))), \\ \quad \quad \quad \quad h) \\ \quad \quad \quad =\; m(h) \wedge \\ \quad \quad \quad \forall x \in \operatorname{pts}(\operatorname{rm\text{-}kms}(s, q)).\; (\lvert (g(x) - h(x)) \rvert \leq (a + a)) \end{array}}}
\vnbprop{rmc-riesz-markov-k}{\vnbstmt{\begin{array}{@{}l@{}} \forall s, q, m, g. \\ \quad (\operatorname{is\text{-}positive\text{-}functional}(\operatorname{rm\text{-}kms}(s, q), m) \wedge (g \in \operatorname{cont\text{-}fns}(\operatorname{rm\text{-}kms}(s, q)))) \Rightarrow  \\ \quad \quad \operatorname{is\text{-}integrable}(\operatorname{pts}(\operatorname{rm\text{-}kms}(s, q)), \\ \quad \quad \quad \operatorname{caratheodory\text{-}sets}(\operatorname{pts}(\operatorname{rm\text{-}kms}(s, q)), \\ \quad \quad \quad \quad \operatorname{method\text{-}i\text{-}outer}(\operatorname{pts}(\operatorname{rm\text{-}kms}(s, q)), \operatorname{clopen\text{-}sets}(\operatorname{rm\text{-}kms}(s, q)), \\ \quad \quad \quad \quad \quad (\lambda u \in \operatorname{clopen\text{-}sets}(\operatorname{rm\text{-}kms}(s, q)).\; m(\operatorname{indicator}(\operatorname{pts}(\operatorname{rm\text{-}kms}(s, q)), u))))), \\ \quad \quad \quad \operatorname{res}(\operatorname{method\text{-}i\text{-}outer}(\operatorname{pts}(\operatorname{rm\text{-}kms}(s, q)), \operatorname{clopen\text{-}sets}(\operatorname{rm\text{-}kms}(s, q)), \\ \quad \quad \quad \quad \quad (\lambda u \in \operatorname{clopen\text{-}sets}(\operatorname{rm\text{-}kms}(s, q)).\; m(\operatorname{indicator}(\operatorname{pts}(\operatorname{rm\text{-}kms}(s, q)), u)))), \\ \quad \quad \quad \quad \operatorname{caratheodory\text{-}sets}(\operatorname{pts}(\operatorname{rm\text{-}kms}(s, q)), \\ \quad \quad \quad \quad \quad \operatorname{method\text{-}i\text{-}outer}(\operatorname{pts}(\operatorname{rm\text{-}kms}(s, q)), \operatorname{clopen\text{-}sets}(\operatorname{rm\text{-}kms}(s, q)), \\ \quad \quad \quad \quad \quad \quad (\lambda u \in \operatorname{clopen\text{-}sets}(\operatorname{rm\text{-}kms}(s, q)).\; m(\operatorname{indicator}(\operatorname{pts}(\operatorname{rm\text{-}kms}(s, q)), u)))))), \\ \quad \quad \quad g) \wedge \\ \quad \quad \operatorname{rintegral}(\operatorname{pts}(\operatorname{rm\text{-}kms}(s, q)), \\ \quad \quad \quad \operatorname{caratheodory\text{-}sets}(\operatorname{pts}(\operatorname{rm\text{-}kms}(s, q)), \\ \quad \quad \quad \quad \operatorname{method\text{-}i\text{-}outer}(\operatorname{pts}(\operatorname{rm\text{-}kms}(s, q)), \operatorname{clopen\text{-}sets}(\operatorname{rm\text{-}kms}(s, q)), \\ \quad \quad \quad \quad \quad (\lambda u \in \operatorname{clopen\text{-}sets}(\operatorname{rm\text{-}kms}(s, q)).\; m(\operatorname{indicator}(\operatorname{pts}(\operatorname{rm\text{-}kms}(s, q)), u))))), \\ \quad \quad \quad \operatorname{res}(\operatorname{method\text{-}i\text{-}outer}(\operatorname{pts}(\operatorname{rm\text{-}kms}(s, q)), \operatorname{clopen\text{-}sets}(\operatorname{rm\text{-}kms}(s, q)), \\ \quad \quad \quad \quad \quad (\lambda u \in \operatorname{clopen\text{-}sets}(\operatorname{rm\text{-}kms}(s, q)).\; m(\operatorname{indicator}(\operatorname{pts}(\operatorname{rm\text{-}kms}(s, q)), u)))), \\ \quad \quad \quad \quad \operatorname{caratheodory\text{-}sets}(\operatorname{pts}(\operatorname{rm\text{-}kms}(s, q)), \\ \quad \quad \quad \quad \quad \operatorname{method\text{-}i\text{-}outer}(\operatorname{pts}(\operatorname{rm\text{-}kms}(s, q)), \operatorname{clopen\text{-}sets}(\operatorname{rm\text{-}kms}(s, q)), \\ \quad \quad \quad \quad \quad \quad (\lambda u \in \operatorname{clopen\text{-}sets}(\operatorname{rm\text{-}kms}(s, q)).\; m(\operatorname{indicator}(\operatorname{pts}(\operatorname{rm\text{-}kms}(s, q)), u)))))), \\ \quad \quad \quad g) \\ \quad \quad =\; m(g) \end{array}}}

\vnbbookfile{theorem-library/positive-extension}{theorem-library/positive-extension}
\noindent{\small theorem-library/positive-extension.scm -- THE POSITIVE EXTENSION (M. Riesz) for C(K), K a compact metric space (RM-3A, 2026-10-10).  Prefix pex-.}\par
\vnbprop{pex-c-set}{\vnbstmt{\forall s.\; (\operatorname{is\text{-}metric\text{-}space}(s) \Rightarrow (\operatorname{cont\text{-}fns}(s) \in \mathbf{Set}))}}
\vnbprop{pex-cf-scal}{\vnbstmt{\begin{array}{@{}l@{}} \forall s. \\ \quad \operatorname{is\text{-}metric\text{-}space}(s) \Rightarrow  \\ \quad \quad (\operatorname{scal}(\langle \operatorname{normed\text{-}field\text{-}as\text{-}commutative\text{-}ring}(\operatorname{rr\text{-}normed\text{-}field}), \operatorname{cont\text{-}fns}(s), (\lambda \langle a, \operatorname{peb\_} \rangle \in (\operatorname{cont\text{-}fns}(s) \times \operatorname{cont\text{-}fns}(s)).\; (\lambda x \in \operatorname{pts}(s).\; (a(x) + \operatorname{peb\_}(x)))), (\lambda x \in \operatorname{pts}(s).\; 0), (\lambda a \in \operatorname{cont\text{-}fns}(s).\; (\lambda x \in \operatorname{pts}(s).\; (-1 \cdot a(x)))), (\lambda \langle \operatorname{pec\_}, a \rangle \in (\mathbb{R} \times \operatorname{cont\text{-}fns}(s)).\; (\lambda x \in \operatorname{pts}(s).\; (\operatorname{pec\_} \cdot a(x)))) \rangle) \simeq \operatorname{normed\text{-}field\text{-}as\text{-}commutative\text{-}ring}(\operatorname{rr\text{-}normed\text{-}field})) \end{array}}}
\vnbprop{pex-cf-vec}{\vnbstmt{\begin{array}{@{}l@{}} \forall s. \\ \quad \operatorname{is\text{-}metric\text{-}space}(s) \Rightarrow  \\ \quad \quad (\operatorname{vec}(\langle \operatorname{normed\text{-}field\text{-}as\text{-}commutative\text{-}ring}(\operatorname{rr\text{-}normed\text{-}field}), \operatorname{cont\text{-}fns}(s), (\lambda \langle a, \operatorname{peb\_} \rangle \in (\operatorname{cont\text{-}fns}(s) \times \operatorname{cont\text{-}fns}(s)).\; (\lambda x \in \operatorname{pts}(s).\; (a(x) + \operatorname{peb\_}(x)))), (\lambda x \in \operatorname{pts}(s).\; 0), (\lambda a \in \operatorname{cont\text{-}fns}(s).\; (\lambda x \in \operatorname{pts}(s).\; (-1 \cdot a(x)))), (\lambda \langle \operatorname{pec\_}, a \rangle \in (\mathbb{R} \times \operatorname{cont\text{-}fns}(s)).\; (\lambda x \in \operatorname{pts}(s).\; (\operatorname{pec\_} \cdot a(x)))) \rangle) \simeq \operatorname{cont\text{-}fns}(s)) \end{array}}}
\vnbprop{pex-cf-vadd}{\vnbstmt{\begin{array}{@{}l@{}} \forall s. \\ \quad \operatorname{is\text{-}metric\text{-}space}(s) \Rightarrow  \\ \quad \quad (\operatorname{vadd}(\langle \operatorname{normed\text{-}field\text{-}as\text{-}commutative\text{-}ring}(\operatorname{rr\text{-}normed\text{-}field}), \operatorname{cont\text{-}fns}(s), (\lambda \langle a, \operatorname{peb\_} \rangle \in (\operatorname{cont\text{-}fns}(s) \times \operatorname{cont\text{-}fns}(s)).\; (\lambda x \in \operatorname{pts}(s).\; (a(x) + \operatorname{peb\_}(x)))), (\lambda x \in \operatorname{pts}(s).\; 0), (\lambda a \in \operatorname{cont\text{-}fns}(s).\; (\lambda x \in \operatorname{pts}(s).\; (-1 \cdot a(x)))), (\lambda \langle \operatorname{pec\_}, a \rangle \in (\mathbb{R} \times \operatorname{cont\text{-}fns}(s)).\; (\lambda x \in \operatorname{pts}(s).\; (\operatorname{pec\_} \cdot a(x)))) \rangle) \simeq (\lambda \langle a, \operatorname{peb\_} \rangle \in (\operatorname{cont\text{-}fns}(s) \times \operatorname{cont\text{-}fns}(s)).\; (\lambda x \in \operatorname{pts}(s).\; (a(x) + \operatorname{peb\_}(x))))) \end{array}}}
\vnbprop{pex-cf-vzero}{\vnbstmt{\begin{array}{@{}l@{}} \forall s. \\ \quad \operatorname{is\text{-}metric\text{-}space}(s) \Rightarrow  \\ \quad \quad (\operatorname{vzero}(\langle \operatorname{normed\text{-}field\text{-}as\text{-}commutative\text{-}ring}(\operatorname{rr\text{-}normed\text{-}field}), \operatorname{cont\text{-}fns}(s), (\lambda \langle a, \operatorname{peb\_} \rangle \in (\operatorname{cont\text{-}fns}(s) \times \operatorname{cont\text{-}fns}(s)).\; (\lambda x \in \operatorname{pts}(s).\; (a(x) + \operatorname{peb\_}(x)))), (\lambda x \in \operatorname{pts}(s).\; 0), (\lambda a \in \operatorname{cont\text{-}fns}(s).\; (\lambda x \in \operatorname{pts}(s).\; (-1 \cdot a(x)))), (\lambda \langle \operatorname{pec\_}, a \rangle \in (\mathbb{R} \times \operatorname{cont\text{-}fns}(s)).\; (\lambda x \in \operatorname{pts}(s).\; (\operatorname{pec\_} \cdot a(x)))) \rangle) \simeq (\lambda x \in \operatorname{pts}(s).\; 0)) \end{array}}}
\vnbprop{pex-cf-vneg}{\vnbstmt{\begin{array}{@{}l@{}} \forall s. \\ \quad \operatorname{is\text{-}metric\text{-}space}(s) \Rightarrow  \\ \quad \quad (\operatorname{vneg}(\langle \operatorname{normed\text{-}field\text{-}as\text{-}commutative\text{-}ring}(\operatorname{rr\text{-}normed\text{-}field}), \operatorname{cont\text{-}fns}(s), (\lambda \langle a, \operatorname{peb\_} \rangle \in (\operatorname{cont\text{-}fns}(s) \times \operatorname{cont\text{-}fns}(s)).\; (\lambda x \in \operatorname{pts}(s).\; (a(x) + \operatorname{peb\_}(x)))), (\lambda x \in \operatorname{pts}(s).\; 0), (\lambda a \in \operatorname{cont\text{-}fns}(s).\; (\lambda x \in \operatorname{pts}(s).\; (-1 \cdot a(x)))), (\lambda \langle \operatorname{pec\_}, a \rangle \in (\mathbb{R} \times \operatorname{cont\text{-}fns}(s)).\; (\lambda x \in \operatorname{pts}(s).\; (\operatorname{pec\_} \cdot a(x)))) \rangle) \simeq (\lambda a \in \operatorname{cont\text{-}fns}(s).\; (\lambda x \in \operatorname{pts}(s).\; (-1 \cdot a(x))))) \end{array}}}
\vnbprop{pex-cf-act}{\vnbstmt{\begin{array}{@{}l@{}} \forall s. \\ \quad \operatorname{is\text{-}metric\text{-}space}(s) \Rightarrow  \\ \quad \quad (\operatorname{act}(\langle \operatorname{normed\text{-}field\text{-}as\text{-}commutative\text{-}ring}(\operatorname{rr\text{-}normed\text{-}field}), \operatorname{cont\text{-}fns}(s), (\lambda \langle a, \operatorname{peb\_} \rangle \in (\operatorname{cont\text{-}fns}(s) \times \operatorname{cont\text{-}fns}(s)).\; (\lambda x \in \operatorname{pts}(s).\; (a(x) + \operatorname{peb\_}(x)))), (\lambda x \in \operatorname{pts}(s).\; 0), (\lambda a \in \operatorname{cont\text{-}fns}(s).\; (\lambda x \in \operatorname{pts}(s).\; (-1 \cdot a(x)))), (\lambda \langle \operatorname{pec\_}, a \rangle \in (\mathbb{R} \times \operatorname{cont\text{-}fns}(s)).\; (\lambda x \in \operatorname{pts}(s).\; (\operatorname{pec\_} \cdot a(x)))) \rangle) \simeq (\lambda \langle \operatorname{pec\_}, a \rangle \in (\mathbb{R} \times \operatorname{cont\text{-}fns}(s)).\; (\lambda x \in \operatorname{pts}(s).\; (\operatorname{pec\_} \cdot a(x))))) \end{array}}}
\vnbprop{pex-zero-in}{\vnbstmt{\forall s.\; (\operatorname{is\text{-}metric\text{-}space}(s) \Rightarrow ((\lambda x \in \operatorname{pts}(s).\; 0) \in \operatorname{cont\text{-}fns}(s)))}}
\vnbprop{pex-cf-module}{\vnbstmt{\begin{array}{@{}l@{}} \forall s. \\ \quad \operatorname{is\text{-}metric\text{-}space}(s) \Rightarrow  \\ \quad \quad \operatorname{is\text{-}module}(\langle \operatorname{normed\text{-}field\text{-}as\text{-}commutative\text{-}ring}(\operatorname{rr\text{-}normed\text{-}field}), \operatorname{cont\text{-}fns}(s), (\lambda \langle a, \operatorname{peb\_} \rangle \in (\operatorname{cont\text{-}fns}(s) \times \operatorname{cont\text{-}fns}(s)).\; (\lambda x \in \operatorname{pts}(s).\; (a(x) + \operatorname{peb\_}(x)))), (\lambda x \in \operatorname{pts}(s).\; 0), (\lambda a \in \operatorname{cont\text{-}fns}(s).\; (\lambda x \in \operatorname{pts}(s).\; (-1 \cdot a(x)))), (\lambda \langle \operatorname{pec\_}, a \rangle \in (\mathbb{R} \times \operatorname{cont\text{-}fns}(s)).\; (\lambda x \in \operatorname{pts}(s).\; (\operatorname{pec\_} \cdot a(x)))) \rangle) \end{array}}}
\vnbprop{pex-vadd-val}{\vnbstmt{\begin{array}{@{}l@{}} \forall s.; \forall f, g \in \operatorname{cont\text{-}fns}(s). \\ \quad \operatorname{is\text{-}metric\text{-}space}(s) \Rightarrow  \\ \quad \quad \operatorname{vadd}(\langle \operatorname{normed\text{-}field\text{-}as\text{-}commutative\text{-}ring}(\operatorname{rr\text{-}normed\text{-}field}), \operatorname{cont\text{-}fns}(s), (\lambda \langle a, \operatorname{peb\_} \rangle \in (\operatorname{cont\text{-}fns}(s) \times \operatorname{cont\text{-}fns}(s)).\; (\lambda x \in \operatorname{pts}(s).\; (a(x) + \operatorname{peb\_}(x)))), (\lambda x \in \operatorname{pts}(s).\; 0), (\lambda a \in \operatorname{cont\text{-}fns}(s).\; (\lambda x \in \operatorname{pts}(s).\; (-1 \cdot a(x)))), (\lambda \langle \operatorname{pec\_}, a \rangle \in (\mathbb{R} \times \operatorname{cont\text{-}fns}(s)).\; (\lambda x \in \operatorname{pts}(s).\; (\operatorname{pec\_} \cdot a(x)))) \rangle)(f, g) \\ \quad \quad =\; (\lambda x \in \operatorname{pts}(s).\; (f(x) + g(x))) \end{array}}}
\vnbprop{pex-act-val}{\vnbstmt{\begin{array}{@{}l@{}} \forall s.; \forall c \in \mathbb{R}, f \in \operatorname{cont\text{-}fns}(s). \\ \quad \operatorname{is\text{-}metric\text{-}space}(s) \Rightarrow  \\ \quad \quad \operatorname{act}(\langle \operatorname{normed\text{-}field\text{-}as\text{-}commutative\text{-}ring}(\operatorname{rr\text{-}normed\text{-}field}), \operatorname{cont\text{-}fns}(s), (\lambda \langle a, \operatorname{peb\_} \rangle \in (\operatorname{cont\text{-}fns}(s) \times \operatorname{cont\text{-}fns}(s)).\; (\lambda x \in \operatorname{pts}(s).\; (a(x) + \operatorname{peb\_}(x)))), (\lambda x \in \operatorname{pts}(s).\; 0), (\lambda a \in \operatorname{cont\text{-}fns}(s).\; (\lambda x \in \operatorname{pts}(s).\; (-1 \cdot a(x)))), (\lambda \langle \operatorname{pec\_}, a \rangle \in (\mathbb{R} \times \operatorname{cont\text{-}fns}(s)).\; (\lambda x \in \operatorname{pts}(s).\; (\operatorname{pec\_} \cdot a(x)))) \rangle)(c, f) \\ \quad \quad =\; (\lambda x \in \operatorname{pts}(s).\; (c \cdot f(x))) \end{array}}}
\vnbprop{pex-cf-pin}{\vnbstmt{\begin{array}{@{}l@{}} \forall s. \\ \quad \operatorname{is\text{-}metric\text{-}space}(s) \Rightarrow  \\ \quad \quad \operatorname{scal}(\langle \operatorname{normed\text{-}field\text{-}as\text{-}commutative\text{-}ring}(\operatorname{rr\text{-}normed\text{-}field}), \operatorname{cont\text{-}fns}(s), (\lambda \langle a, \operatorname{peb\_} \rangle \in (\operatorname{cont\text{-}fns}(s) \times \operatorname{cont\text{-}fns}(s)).\; (\lambda x \in \operatorname{pts}(s).\; (a(x) + \operatorname{peb\_}(x)))), (\lambda x \in \operatorname{pts}(s).\; 0), (\lambda a \in \operatorname{cont\text{-}fns}(s).\; (\lambda x \in \operatorname{pts}(s).\; (-1 \cdot a(x)))), (\lambda \langle \operatorname{pec\_}, a \rangle \in (\mathbb{R} \times \operatorname{cont\text{-}fns}(s)).\; (\lambda x \in \operatorname{pts}(s).\; (\operatorname{pec\_} \cdot a(x)))) \rangle) \\ \quad \quad =\; \operatorname{normed\text{-}field\text{-}as\text{-}commutative\text{-}ring}(\operatorname{rr\text{-}normed\text{-}field}) \end{array}}}
\vnbprop{pex-vs-zero}{\vnbstmt{\begin{array}{@{}l@{}} \forall s, a.; \forall \ell \in (a \to \mathbb{R}). \\ \quad \operatorname{is\text{-}metric\text{-}space}(s) \wedge \\ \quad \operatorname{is\text{-}compact}(s) \wedge \\ \quad a \subseteq \operatorname{cont\text{-}fns}(s) \wedge \\ \quad \forall f, g \in a.\; ((\lambda x \in \operatorname{pts}(s).\; (f(x) + g(x))) \in a) \wedge \\ \quad \forall c \in \mathbb{R}, f \in a.\; ((\lambda x \in \operatorname{pts}(s).\; (c \cdot f(x))) \in a) \wedge \\ \quad (\lambda x \in \operatorname{pts}(s).\; 1) \in a \wedge \\ \quad \forall f, g \in a. \\ \quad \quad \ell((\lambda x \in \operatorname{pts}(s).\; (f(x) + g(x)))) = (\ell(f) + \ell(g)) \wedge \\ \quad \forall c \in \mathbb{R}, f \in a.\; (\ell((\lambda x \in \operatorname{pts}(s).\; (c \cdot f(x)))) = (c \cdot \ell(f))) \wedge \\ \quad \forall f \in a.\; (\forall b \in \operatorname{pts}(s).\; (0 \leq f(b)) \Rightarrow (0 \leq \ell(f))) \Rightarrow \\ \quad \quad ((\lambda x \in \operatorname{pts}(s).\; 0) \in a) \wedge (\ell((\lambda x \in \operatorname{pts}(s).\; 0)) = 0) \end{array}}}
\vnbprop{pex-lm-mono}{\vnbstmt{\begin{array}{@{}l@{}} \forall s, a.; \forall \ell \in (a \to \mathbb{R}), f, g \in a. \\ \quad \operatorname{is\text{-}metric\text{-}space}(s) \wedge \\ \quad \operatorname{is\text{-}compact}(s) \wedge \\ \quad a \subseteq \operatorname{cont\text{-}fns}(s) \wedge \\ \quad \forall f, g \in a.\; ((\lambda x \in \operatorname{pts}(s).\; (f(x) + g(x))) \in a) \wedge \\ \quad \forall c \in \mathbb{R}, f \in a.\; ((\lambda x \in \operatorname{pts}(s).\; (c \cdot f(x))) \in a) \wedge \\ \quad (\lambda x \in \operatorname{pts}(s).\; 1) \in a \wedge \\ \quad \forall f, g \in a. \\ \quad \quad \ell((\lambda x \in \operatorname{pts}(s).\; (f(x) + g(x)))) = (\ell(f) + \ell(g)) \wedge \\ \quad \forall c \in \mathbb{R}, f \in a.\; (\ell((\lambda x \in \operatorname{pts}(s).\; (c \cdot f(x)))) = (c \cdot \ell(f))) \wedge \\ \quad \forall f \in a.\; (\forall b \in \operatorname{pts}(s).\; (0 \leq f(b)) \Rightarrow (0 \leq \ell(f))) \wedge \\ \quad \forall b \in \operatorname{pts}(s).\; (f(b) \leq g(b)) \Rightarrow \\ \quad \quad \ell(f) \leq \ell(g) \end{array}}}
\vnbprop{pex-vs-submodule}{\vnbstmt{\begin{array}{@{}l@{}} \forall s, a.; \forall \ell \in (a \to \mathbb{R}). \\ \quad \operatorname{is\text{-}metric\text{-}space}(s) \wedge \\ \quad \operatorname{is\text{-}compact}(s) \wedge \\ \quad a \subseteq \operatorname{cont\text{-}fns}(s) \wedge \\ \quad \forall f, g \in a.\; ((\lambda x \in \operatorname{pts}(s).\; (f(x) + g(x))) \in a) \wedge \\ \quad \forall c \in \mathbb{R}, f \in a.\; ((\lambda x \in \operatorname{pts}(s).\; (c \cdot f(x))) \in a) \wedge \\ \quad (\lambda x \in \operatorname{pts}(s).\; 1) \in a \wedge \\ \quad \forall f, g \in a. \\ \quad \quad \ell((\lambda x \in \operatorname{pts}(s).\; (f(x) + g(x)))) = (\ell(f) + \ell(g)) \wedge \\ \quad \forall c \in \mathbb{R}, f \in a.\; (\ell((\lambda x \in \operatorname{pts}(s).\; (c \cdot f(x)))) = (c \cdot \ell(f))) \wedge \\ \quad \forall f \in a.\; (\forall b \in \operatorname{pts}(s).\; (0 \leq f(b)) \Rightarrow (0 \leq \ell(f))) \Rightarrow \\ \quad \quad \operatorname{is\text{-}submodule}(\langle \operatorname{normed\text{-}field\text{-}as\text{-}commutative\text{-}ring}(\operatorname{rr\text{-}normed\text{-}field}), \operatorname{cont\text{-}fns}(s), (\lambda \langle \mathit{pea}, \operatorname{peb\_} \rangle \in (\operatorname{cont\text{-}fns}(s) \times \operatorname{cont\text{-}fns}(s)).\; (\lambda x \in \operatorname{pts}(s).\; (\operatorname{pea}(x) + \operatorname{peb\_}(x)))), (\lambda x \in \operatorname{pts}(s).\; 0), (\lambda \mathit{pea} \in \operatorname{cont\text{-}fns}(s).\; (\lambda x \in \operatorname{pts}(s).\; (-1 \cdot \operatorname{pea}(x)))), (\lambda \langle \operatorname{pec\_}, \mathit{pea} \rangle \in (\mathbb{R} \times \operatorname{cont\text{-}fns}(s)).\; (\lambda x \in \operatorname{pts}(s).\; (\operatorname{pec\_} \cdot \operatorname{pea}(x)))) \rangle, \\ \quad \quad \quad a) \end{array}}}
\vnbprop{pex-lm-linear}{\vnbstmt{\begin{array}{@{}l@{}} \forall s, a.; \forall \ell \in (a \to \mathbb{R}). \\ \quad \operatorname{is\text{-}metric\text{-}space}(s) \wedge \\ \quad \operatorname{is\text{-}compact}(s) \wedge \\ \quad a \subseteq \operatorname{cont\text{-}fns}(s) \wedge \\ \quad \forall f, g \in a.\; ((\lambda x \in \operatorname{pts}(s).\; (f(x) + g(x))) \in a) \wedge \\ \quad \forall c \in \mathbb{R}, f \in a.\; ((\lambda x \in \operatorname{pts}(s).\; (c \cdot f(x))) \in a) \wedge \\ \quad (\lambda x \in \operatorname{pts}(s).\; 1) \in a \wedge \\ \quad \forall f, g \in a. \\ \quad \quad \ell((\lambda x \in \operatorname{pts}(s).\; (f(x) + g(x)))) = (\ell(f) + \ell(g)) \wedge \\ \quad \forall c \in \mathbb{R}, f \in a.\; (\ell((\lambda x \in \operatorname{pts}(s).\; (c \cdot f(x)))) = (c \cdot \ell(f))) \wedge \\ \quad \forall f \in a.\; (\forall b \in \operatorname{pts}(s).\; (0 \leq f(b)) \Rightarrow (0 \leq \ell(f))) \Rightarrow \\ \quad \quad \operatorname{is\text{-}linear\text{-}functional\text{-}on}(\langle \operatorname{normed\text{-}field\text{-}as\text{-}commutative\text{-}ring}(\operatorname{rr\text{-}normed\text{-}field}), \operatorname{cont\text{-}fns}(s), (\lambda \langle \mathit{pea}, \operatorname{peb\_} \rangle \in (\operatorname{cont\text{-}fns}(s) \times \operatorname{cont\text{-}fns}(s)).\; (\lambda x \in \operatorname{pts}(s).\; (\operatorname{pea}(x) + \operatorname{peb\_}(x)))), (\lambda x \in \operatorname{pts}(s).\; 0), (\lambda \mathit{pea} \in \operatorname{cont\text{-}fns}(s).\; (\lambda x \in \operatorname{pts}(s).\; (-1 \cdot \operatorname{pea}(x)))), (\lambda \langle \operatorname{pec\_}, \mathit{pea} \rangle \in (\mathbb{R} \times \operatorname{cont\text{-}fns}(s)).\; (\lambda x \in \operatorname{pts}(s).\; (\operatorname{pec\_} \cdot \operatorname{pea}(x)))) \rangle, \\ \quad \quad \quad a, \ell) \end{array}}}
\vnbprop{pex-iif-mem}{\vnbstmt{\begin{array}{@{}l@{}} \forall s, a.; \forall g \in \operatorname{cont\text{-}fns}(s), v. \\ \quad (\operatorname{is\text{-}metric\text{-}space}(s) \wedge (a \subseteq \operatorname{cont\text{-}fns}(s))) \Rightarrow  \\ \quad \quad v \in (\lambda \mathit{peg} \in \operatorname{cont\text{-}fns}(s).\; \{\mathit{pev} \in a : \forall x \in \operatorname{pts}(s).\; (\operatorname{peg}(x) \leq \operatorname{pev}(x))\})(g) \\ \quad \quad \Leftrightarrow\; v \in \{\mathit{pev} \in a : \forall x \in \operatorname{pts}(s).\; (g(x) \leq \operatorname{pev}(x))\} \end{array}}}
\vnbprop{pex-p-exists}{\vnbstmt{\begin{array}{@{}l@{}} \forall s, a.; \forall \ell \in (a \to \mathbb{R}). \\ \quad \operatorname{is\text{-}metric\text{-}space}(s) \wedge \\ \quad \operatorname{is\text{-}compact}(s) \wedge \\ \quad a \subseteq \operatorname{cont\text{-}fns}(s) \wedge \\ \quad \forall f, g \in a.\; ((\lambda \mathit{sx} \in \operatorname{pts}(s).\; (f(\mathit{sx}) + g(\mathit{sx}))) \in a) \wedge \\ \quad \forall c \in \mathbb{R}, f \in a.\; ((\lambda \mathit{sx} \in \operatorname{pts}(s).\; (c \cdot f(\mathit{sx}))) \in a) \wedge \\ \quad (\lambda \mathit{sx} \in \operatorname{pts}(s).\; 1) \in a \wedge \\ \quad \forall f, g \in a. \\ \quad \quad \ell((\lambda \mathit{sx} \in \operatorname{pts}(s).\; (f(\mathit{sx}) + g(\mathit{sx})))) = (\ell(f) + \ell(g)) \wedge \\ \quad \forall c \in \mathbb{R}, f \in a.\; (\ell((\lambda \mathit{sx} \in \operatorname{pts}(s).\; (c \cdot f(\mathit{sx})))) = (c \cdot \ell(f))) \wedge \\ \quad \forall f \in a.\; (\forall b \in \operatorname{pts}(s).\; (0 \leq f(b)) \Rightarrow (0 \leq \ell(f))) \wedge \\ \quad \forall g \in \operatorname{cont\text{-}fns}(s). \exists d \in \mathbb{R}. \\ \quad \quad (0 \leq d) \wedge \forall x \in \operatorname{pts}(s).\; (\lvert g(x) \rvert \leq d) \Rightarrow \\ \quad \quad \exists p \in (\operatorname{cont\text{-}fns}(s) \to \mathbb{R}). \\ \quad \quad \quad \forall g \in \operatorname{cont\text{-}fns}(s), v \in a. \\ \quad \quad \quad \quad \forall \mathit{pex} \in \operatorname{pts}(s).\; (g(\mathit{pex}) \leq v(\mathit{pex})) \Rightarrow (p(g) \leq \ell(v)) \wedge \\ \quad \quad \quad \forall g \in \operatorname{cont\text{-}fns}(s), \beta \in \mathbb{R}. \\ \quad \quad \quad \quad \forall v \in a.\; (\forall \mathit{pex} \in \operatorname{pts}(s).\; (g(\mathit{pex}) \leq v(\mathit{pex})) \Rightarrow (\beta \leq \ell(v))) \Rightarrow  \\ \quad \quad \quad \quad \quad \beta \leq p(g) \end{array}}}
\vnbprop{pex-p-props}{\vnbstmt{\begin{array}{@{}l@{}} \forall s, a.; \forall \ell \in (a \to \mathbb{R}), p \in (\operatorname{cont\text{-}fns}(s) \to \mathbb{R}). \\ \quad \operatorname{is\text{-}metric\text{-}space}(s) \wedge \\ \quad \operatorname{is\text{-}compact}(s) \wedge \\ \quad a \subseteq \operatorname{cont\text{-}fns}(s) \wedge \\ \quad \forall f, g \in a.\; ((\lambda x \in \operatorname{pts}(s).\; (f(x) + g(x))) \in a) \wedge \\ \quad \forall c \in \mathbb{R}, f \in a.\; ((\lambda x \in \operatorname{pts}(s).\; (c \cdot f(x))) \in a) \wedge \\ \quad (\lambda x \in \operatorname{pts}(s).\; 1) \in a \wedge \\ \quad \forall f, g \in a. \\ \quad \quad \ell((\lambda x \in \operatorname{pts}(s).\; (f(x) + g(x)))) = (\ell(f) + \ell(g)) \wedge \\ \quad \forall c \in \mathbb{R}, f \in a.\; (\ell((\lambda x \in \operatorname{pts}(s).\; (c \cdot f(x)))) = (c \cdot \ell(f))) \wedge \\ \quad \forall f \in a.\; (\forall b \in \operatorname{pts}(s).\; (0 \leq f(b)) \Rightarrow (0 \leq \ell(f))) \wedge \\ \quad \forall g \in \operatorname{cont\text{-}fns}(s), v \in a. \\ \quad \quad \forall \mathit{pex} \in \operatorname{pts}(s).\; (g(\mathit{pex}) \leq v(\mathit{pex})) \Rightarrow (p(g) \leq \ell(v)) \wedge \\ \quad \forall g \in \operatorname{cont\text{-}fns}(s), \beta \in \mathbb{R}. \\ \quad \quad \forall v \in a.\; (\forall \mathit{pex} \in \operatorname{pts}(s).\; (g(\mathit{pex}) \leq v(\mathit{pex})) \Rightarrow (\beta \leq \ell(v))) \Rightarrow  \\ \quad \quad \quad \beta \leq p(g) \Rightarrow \\ \quad \quad \operatorname{is\text{-}sublinear}(\langle \operatorname{normed\text{-}field\text{-}as\text{-}commutative\text{-}ring}(\operatorname{rr\text{-}normed\text{-}field}), \operatorname{cont\text{-}fns}(s), (\lambda \langle \mathit{pea}, \operatorname{peb\_} \rangle \in (\operatorname{cont\text{-}fns}(s) \times \operatorname{cont\text{-}fns}(s)).\; (\lambda x \in \operatorname{pts}(s).\; (\operatorname{pea}(x) + \operatorname{peb\_}(x)))), (\lambda x \in \operatorname{pts}(s).\; 0), (\lambda \mathit{pea} \in \operatorname{cont\text{-}fns}(s).\; (\lambda x \in \operatorname{pts}(s).\; (-1 \cdot \operatorname{pea}(x)))), (\lambda \langle \operatorname{pec\_}, \mathit{pea} \rangle \in (\mathbb{R} \times \operatorname{cont\text{-}fns}(s)).\; (\lambda x \in \operatorname{pts}(s).\; (\operatorname{pec\_} \cdot \operatorname{pea}(x)))) \rangle, \\ \quad \quad \quad p) \wedge \\ \quad \quad \forall v \in a.\; (\ell(v) \leq p(v)) \wedge \\ \quad \quad \forall g \in \operatorname{cont\text{-}fns}(s).\; (\forall b \in \operatorname{pts}(s).\; (g(b) \leq 0) \Rightarrow (p(g) \leq 0)) \end{array}}}
\vnbprop{pex-extension-from-bound}{\vnbstmt{\begin{array}{@{}l@{}} \forall s, a.; \forall \ell \in (a \to \mathbb{R}). \\ \quad \operatorname{is\text{-}metric\text{-}space}(s) \wedge \\ \quad \operatorname{is\text{-}compact}(s) \wedge \\ \quad a \subseteq \operatorname{cont\text{-}fns}(s) \wedge \\ \quad \forall f, g \in a.\; ((\lambda \mathit{sx} \in \operatorname{pts}(s).\; (f(\mathit{sx}) + g(\mathit{sx}))) \in a) \wedge \\ \quad \forall c \in \mathbb{R}, f \in a.\; ((\lambda \mathit{sx} \in \operatorname{pts}(s).\; (c \cdot f(\mathit{sx}))) \in a) \wedge \\ \quad (\lambda \mathit{sx} \in \operatorname{pts}(s).\; 1) \in a \wedge \\ \quad \forall f, g \in a. \\ \quad \quad \ell((\lambda \mathit{sx} \in \operatorname{pts}(s).\; (f(\mathit{sx}) + g(\mathit{sx})))) = (\ell(f) + \ell(g)) \wedge \\ \quad \forall c \in \mathbb{R}, f \in a.\; (\ell((\lambda \mathit{sx} \in \operatorname{pts}(s).\; (c \cdot f(\mathit{sx})))) = (c \cdot \ell(f))) \wedge \\ \quad \forall f \in a.\; (\forall b \in \operatorname{pts}(s).\; (0 \leq f(b)) \Rightarrow (0 \leq \ell(f))) \wedge \\ \quad \forall g \in \operatorname{cont\text{-}fns}(s). \exists d \in \mathbb{R}. \\ \quad \quad (0 \leq d) \wedge \forall x \in \operatorname{pts}(s).\; (\lvert g(x) \rvert \leq d) \Rightarrow \\ \quad \quad \exists m.\; (\operatorname{is\text{-}positive\text{-}functional}(s, m) \wedge \operatorname{extends\text{-}on}(a, m, \ell)) \end{array}}}
\vnbprop{pex-positive-extension}{\vnbstmt{\begin{array}{@{}l@{}} \forall s, a.; \forall \ell \in (a \to \mathbb{R}). \\ \quad \operatorname{is\text{-}metric\text{-}space}(s) \wedge \\ \quad \operatorname{is\text{-}compact}(s) \wedge \\ \quad a \subseteq \operatorname{cont\text{-}fns}(s) \wedge \\ \quad \forall f, g \in a.\; ((\lambda x \in \operatorname{pts}(s).\; (f(x) + g(x))) \in a) \wedge \\ \quad \forall c \in \mathbb{R}, f \in a.\; ((\lambda x \in \operatorname{pts}(s).\; (c \cdot f(x))) \in a) \wedge \\ \quad (\lambda x \in \operatorname{pts}(s).\; 1) \in a \wedge \\ \quad \forall f, g \in a. \\ \quad \quad \ell((\lambda x \in \operatorname{pts}(s).\; (f(x) + g(x)))) = (\ell(f) + \ell(g)) \wedge \\ \quad \forall c \in \mathbb{R}, f \in a.\; (\ell((\lambda x \in \operatorname{pts}(s).\; (c \cdot f(x)))) = (c \cdot \ell(f))) \wedge \\ \quad \forall f \in a.\; (\forall b \in \operatorname{pts}(s).\; (0 \leq f(b)) \Rightarrow (0 \leq \ell(f))) \Rightarrow \\ \quad \quad \exists m.\; (\operatorname{is\text{-}positive\text{-}functional}(s, m) \wedge \operatorname{extends\text{-}on}(a, m, \ell)) \end{array}}}

\vnbbookfile{theorem-library/lebesgue-measure-line}{theorem-library/lebesgue-measure-line}
\noindent{\small ==================================================================== theorem-library/lebesgue-measure-line.scm -- MEASURE-3 (the join), 2026-10-09.}\par
\vnbprop{lebesgue-measure-ccint}{\vnbstmt{\begin{array}{@{}l@{}} \forall a, b. \\ \quad ((a \in \mathbb{R}) \wedge ((b \in \mathbb{R}) \wedge (a \leq b))) \Rightarrow  \\ \quad \quad \operatorname{ccint}(a, b) \\ \quad \quad \in\; \operatorname{caratheodory\text{-}sets}(\mathbb{R}, \\ \quad \quad \quad \operatorname{method\text{-}i\text{-}outer}(\mathbb{R}, \operatorname{closed\text{-}intervals}, (\lambda s \in \operatorname{closed\text{-}intervals}.\; \operatorname{interval\text{-}length}(s)))) \wedge \\ \quad \quad \operatorname{res}(\operatorname{method\text{-}i\text{-}outer}(\mathbb{R}, \operatorname{closed\text{-}intervals}, (\lambda s \in \operatorname{closed\text{-}intervals}.\; \operatorname{interval\text{-}length}(s))), \operatorname{caratheodory\text{-}sets}(\mathbb{R}, \operatorname{method\text{-}i\text{-}outer}(\mathbb{R}, \operatorname{closed\text{-}intervals}, (\lambda s \in \operatorname{closed\text{-}intervals}.\; \operatorname{interval\text{-}length}(s)))))(\operatorname{ccint}(a, b)) \\ \quad \quad =\; (b - a) \end{array}}}
\vnbprop{lebesgue-measure-ooint}{\vnbstmt{\begin{array}{@{}l@{}} \forall a, b. \\ \quad ((a \in \mathbb{R}) \wedge ((b \in \mathbb{R}) \wedge (a \leq b))) \Rightarrow  \\ \quad \quad \operatorname{ooint}(a, b) \\ \quad \quad \in\; \operatorname{caratheodory\text{-}sets}(\mathbb{R}, \\ \quad \quad \quad \operatorname{method\text{-}i\text{-}outer}(\mathbb{R}, \operatorname{closed\text{-}intervals}, (\lambda s \in \operatorname{closed\text{-}intervals}.\; \operatorname{interval\text{-}length}(s)))) \wedge \\ \quad \quad \operatorname{res}(\operatorname{method\text{-}i\text{-}outer}(\mathbb{R}, \operatorname{closed\text{-}intervals}, (\lambda s \in \operatorname{closed\text{-}intervals}.\; \operatorname{interval\text{-}length}(s))), \operatorname{caratheodory\text{-}sets}(\mathbb{R}, \operatorname{method\text{-}i\text{-}outer}(\mathbb{R}, \operatorname{closed\text{-}intervals}, (\lambda s \in \operatorname{closed\text{-}intervals}.\; \operatorname{interval\text{-}length}(s)))))(\operatorname{ooint}(a, b)) \\ \quad \quad =\; (b - a) \end{array}}}
\vnbprop{lebesgue-measure-coint}{\vnbstmt{\begin{array}{@{}l@{}} \forall a, b. \\ \quad ((a \in \mathbb{R}) \wedge ((b \in \mathbb{R}) \wedge (a \leq b))) \Rightarrow  \\ \quad \quad \operatorname{coint}(a, b) \\ \quad \quad \in\; \operatorname{caratheodory\text{-}sets}(\mathbb{R}, \\ \quad \quad \quad \operatorname{method\text{-}i\text{-}outer}(\mathbb{R}, \operatorname{closed\text{-}intervals}, (\lambda s \in \operatorname{closed\text{-}intervals}.\; \operatorname{interval\text{-}length}(s)))) \wedge \\ \quad \quad \operatorname{res}(\operatorname{method\text{-}i\text{-}outer}(\mathbb{R}, \operatorname{closed\text{-}intervals}, (\lambda s \in \operatorname{closed\text{-}intervals}.\; \operatorname{interval\text{-}length}(s))), \operatorname{caratheodory\text{-}sets}(\mathbb{R}, \operatorname{method\text{-}i\text{-}outer}(\mathbb{R}, \operatorname{closed\text{-}intervals}, (\lambda s \in \operatorname{closed\text{-}intervals}.\; \operatorname{interval\text{-}length}(s)))))(\operatorname{coint}(a, b)) \\ \quad \quad =\; (b - a) \end{array}}}
\vnbprop{lebesgue-measure-ocint}{\vnbstmt{\begin{array}{@{}l@{}} \forall a, b. \\ \quad ((a \in \mathbb{R}) \wedge ((b \in \mathbb{R}) \wedge (a \leq b))) \Rightarrow  \\ \quad \quad \operatorname{ocint}(a, b) \\ \quad \quad \in\; \operatorname{caratheodory\text{-}sets}(\mathbb{R}, \\ \quad \quad \quad \operatorname{method\text{-}i\text{-}outer}(\mathbb{R}, \operatorname{closed\text{-}intervals}, (\lambda s \in \operatorname{closed\text{-}intervals}.\; \operatorname{interval\text{-}length}(s)))) \wedge \\ \quad \quad \operatorname{res}(\operatorname{method\text{-}i\text{-}outer}(\mathbb{R}, \operatorname{closed\text{-}intervals}, (\lambda s \in \operatorname{closed\text{-}intervals}.\; \operatorname{interval\text{-}length}(s))), \operatorname{caratheodory\text{-}sets}(\mathbb{R}, \operatorname{method\text{-}i\text{-}outer}(\mathbb{R}, \operatorname{closed\text{-}intervals}, (\lambda s \in \operatorname{closed\text{-}intervals}.\; \operatorname{interval\text{-}length}(s)))))(\operatorname{ocint}(a, b)) \\ \quad \quad =\; (b - a) \end{array}}}
\vnbprop{lebesgue-measure-point}{\vnbstmt{\begin{array}{@{}l@{}} \forall x \in \mathbb{R}. \\ \quad \operatorname{singleton}(x) \\ \quad \in\; \operatorname{caratheodory\text{-}sets}(\mathbb{R}, \\ \quad \quad \operatorname{method\text{-}i\text{-}outer}(\mathbb{R}, \operatorname{closed\text{-}intervals}, (\lambda s \in \operatorname{closed\text{-}intervals}.\; \operatorname{interval\text{-}length}(s)))) \wedge \\ \quad \operatorname{res}(\operatorname{method\text{-}i\text{-}outer}(\mathbb{R}, \operatorname{closed\text{-}intervals}, (\lambda s \in \operatorname{closed\text{-}intervals}.\; \operatorname{interval\text{-}length}(s))), \operatorname{caratheodory\text{-}sets}(\mathbb{R}, \operatorname{method\text{-}i\text{-}outer}(\mathbb{R}, \operatorname{closed\text{-}intervals}, (\lambda s \in \operatorname{closed\text{-}intervals}.\; \operatorname{interval\text{-}length}(s)))))(\operatorname{singleton}(x)) \\ \quad =\; 0 \end{array}}}

\vnbbookfile{structure-library/lebesgue-translate}{structure-library/lebesgue-translate}
\noindent{\small lebesgue-translate.scm -- the TRANSLATE of a set of reals.  Vocabulary only: no proof, no support.  Agent MEASURE-4-A, 2026-10-10 (scratchpad/triage/MEASURE-4-A.md). The proofs that use it are theorem-library/lebesgue-measure-unique.scm.}\par
\par\medskip\noindent\textbf{Definition} (\texttt{translate}).\;$(\mathit{es} + c) :=$ $\{\operatorname{ltz\_} \in \mathbb{R} : ((\operatorname{ltz\_} - c) \in \mathit{es})\}$

\vnbbookfile{theorem-library/lebesgue-measure-unique}{theorem-library/lebesgue-measure-unique}
\noindent{\small ==================================================================== theorem-library/lebesgue-measure-unique.scm -- MEASURE-4-A, 2026-10-10. Lebesgue measure on the line: translation invariance and uniqueness.}\par
\vnbprop{translate-unfold}{\vnbstmt{\forall s, c.\; ((s + c) \simeq \{z \in \mathbb{R} : ((z - c) \in s)\})}}
\vnbprop{translate-membership}{\vnbstmt{\forall s, c, z.\; ((z \in (s + c)) \Leftrightarrow ((z \in \mathbb{R}) \wedge ((z - c) \in s)))}}
\vnbprop{translate-subset-rr}{\vnbstmt{\forall s, c.\; ((s + c) \subseteq \mathbb{R})}}
\vnbprop{lebesgue-outer-translate-invariant}{\vnbstmt{\begin{array}{@{}l@{}} \forall s, c. \\ \quad ((s \subseteq \mathbb{R}) \wedge (c \in \mathbb{R})) \Rightarrow  \\ \quad \quad \operatorname{method\text{-}i\text{-}outer}(\mathbb{R}, \operatorname{closed\text{-}intervals}, (\lambda \mathit{ols} \in \operatorname{closed\text{-}intervals}.\; \operatorname{interval\text{-}length}(\mathit{ols})))((s + c)) \\ \quad \quad =\; \operatorname{method\text{-}i\text{-}outer}(\mathbb{R}, \operatorname{closed\text{-}intervals}, (\lambda \mathit{ols} \in \operatorname{closed\text{-}intervals}.\; \operatorname{interval\text{-}length}(\mathit{ols})))(s) \end{array}}}
\vnbprop{translate-intersection}{\vnbstmt{\begin{array}{@{}l@{}} \forall a, b.; \forall c \in \mathbb{R}. \\ \quad ((a \cap b) + c) = ((a + c) \cap (b + c)) \end{array}}}
\vnbprop{translate-difference}{\vnbstmt{\begin{array}{@{}l@{}} \forall a, b.; \forall c \in \mathbb{R}. \\ \quad (\operatorname{difference}(a, b) + c) = \operatorname{difference}((a + c), (b + c)) \end{array}}}
\vnbprop{translate-neg-back}{\vnbstmt{\begin{array}{@{}l@{}} \forall a, c. \\ \quad ((a \subseteq \mathbb{R}) \wedge (c \in \mathbb{R})) \Rightarrow (((a + (0 - c)) + c) = a) \end{array}}}
\vnbprop{lebesgue-measurable-translate}{\vnbstmt{\begin{array}{@{}l@{}} \forall s \in \operatorname{caratheodory\text{-}sets}(\mathbb{R}, \operatorname{method\text{-}i\text{-}outer}(\mathbb{R}, \operatorname{closed\text{-}intervals}, (\lambda \mathit{ols} \in \operatorname{closed\text{-}intervals}.\; \operatorname{interval\text{-}length}(\mathit{ols})))), c \in \mathbb{R}. \\ \quad (s + c) \\ \quad \in\; \operatorname{caratheodory\text{-}sets}(\mathbb{R}, \\ \quad \quad \operatorname{method\text{-}i\text{-}outer}(\mathbb{R}, \operatorname{closed\text{-}intervals}, (\lambda \mathit{ols} \in \operatorname{closed\text{-}intervals}.\; \operatorname{interval\text{-}length}(\mathit{ols})))) \wedge \\ \quad \operatorname{res}(\operatorname{method\text{-}i\text{-}outer}(\mathbb{R}, \operatorname{closed\text{-}intervals}, (\lambda \mathit{ols} \in \operatorname{closed\text{-}intervals}.\; \operatorname{interval\text{-}length}(\mathit{ols}))), \operatorname{caratheodory\text{-}sets}(\mathbb{R}, \operatorname{method\text{-}i\text{-}outer}(\mathbb{R}, \operatorname{closed\text{-}intervals}, (\lambda \mathit{ols} \in \operatorname{closed\text{-}intervals}.\; \operatorname{interval\text{-}length}(\mathit{ols})))))((s + c)) \\ \quad =\; \operatorname{res}(\operatorname{method\text{-}i\text{-}outer}(\mathbb{R}, \operatorname{closed\text{-}intervals}, (\lambda \mathit{ols} \in \operatorname{closed\text{-}intervals}.\; \operatorname{interval\text{-}length}(\mathit{ols}))), \operatorname{caratheodory\text{-}sets}(\mathbb{R}, \operatorname{method\text{-}i\text{-}outer}(\mathbb{R}, \operatorname{closed\text{-}intervals}, (\lambda \mathit{ols} \in \operatorname{closed\text{-}intervals}.\; \operatorname{interval\text{-}length}(\mathit{ols})))))(s) \end{array}}}
\vnbprop{rr-closed-borel}{\vnbstmt{\forall l.\; (\operatorname{is\text{-}closed}(\operatorname{rr\text{-}ms}, l) \Rightarrow (l \in \operatorname{borel\text{-}sets}(\operatorname{rr\text{-}ms})))}}
\vnbprop{ccint-closed-rr-ms}{\vnbstmt{\forall a, b.\; (((a \in \mathbb{R}) \wedge (b \in \mathbb{R})) \Rightarrow \operatorname{is\text{-}closed}(\operatorname{rr\text{-}ms}, \operatorname{ccint}(a, b)))}}
\vnbprop{closed-intervals-borel}{\vnbstmt{\operatorname{closed\text{-}intervals} \subseteq \operatorname{borel\text{-}sets}(\operatorname{rr\text{-}ms})}}
\vnbprop{borel-measure-le-lebesgue-outer}{\vnbstmt{\begin{array}{@{}l@{}} \forall u, s. \\ \quad \operatorname{is\text{-}measure}(\mathbb{R}, \operatorname{borel\text{-}sets}(\operatorname{rr\text{-}ms}), u) \wedge \\ \quad \forall a, b.\; (((a \in \mathbb{R}) \wedge ((b \in \mathbb{R}) \wedge (a \leq b))) \Rightarrow (u(\operatorname{ccint}(a, b)) = (b - a))) \wedge \\ \quad s \in \operatorname{borel\text{-}sets}(\operatorname{rr\text{-}ms}) \Rightarrow \\ \quad \quad u(s) \\ \quad \quad \leq\; \operatorname{method\text{-}i\text{-}outer}(\mathbb{R}, \operatorname{closed\text{-}intervals}, (\lambda \mathit{ols} \in \operatorname{closed\text{-}intervals}.\; \operatorname{interval\text{-}length}(\mathit{ols})))(s) \end{array}}}
\vnbprop{borel-measure-eq-lebesgue-outer-bounded}{\vnbstmt{\begin{array}{@{}l@{}} \forall u, p, q, s. \\ \quad \operatorname{is\text{-}measure}(\mathbb{R}, \operatorname{borel\text{-}sets}(\operatorname{rr\text{-}ms}), u) \wedge \\ \quad \forall a, b.\; (((a \in \mathbb{R}) \wedge ((b \in \mathbb{R}) \wedge (a \leq b))) \Rightarrow (u(\operatorname{ccint}(a, b)) = (b - a))) \wedge \\ \quad p \in \mathbb{R} \wedge \\ \quad q \in \mathbb{R} \wedge \\ \quad p \leq q \wedge \\ \quad s \in \operatorname{borel\text{-}sets}(\operatorname{rr\text{-}ms}) \wedge \\ \quad s \subseteq \operatorname{ccint}(p, q) \Rightarrow \\ \quad \quad u(s) \\ \quad \quad =\; \operatorname{method\text{-}i\text{-}outer}(\mathbb{R}, \operatorname{closed\text{-}intervals}, (\lambda \mathit{ols} \in \operatorname{closed\text{-}intervals}.\; \operatorname{interval\text{-}length}(\mathit{ols})))(s) \end{array}}}
\vnbprop{measure-le-of-ccint-exhaust}{\vnbstmt{\begin{array}{@{}l@{}} \forall a, u, s, d. \\ \quad \operatorname{is\text{-}measure}(\mathbb{R}, a, u) \wedge \\ \quad \forall n \in \mathbb{N}.\; (\operatorname{ccint}((0 - n), n) \in a) \wedge \\ \quad s \in a \wedge \\ \quad d \in \operatorname{rr\text{-}pos\text{-}star} \wedge \\ \quad \forall n \in \mathbb{N}.\; (u((s \cap \operatorname{ccint}((0 - n), n))) \leq d) \Rightarrow \\ \quad \quad u(s) \leq d \end{array}}}
\vnbprop{lebesgue-measure-unique}{\vnbstmt{\begin{array}{@{}l@{}} \forall u, s. \\ \quad \operatorname{is\text{-}measure}(\mathbb{R}, \operatorname{borel\text{-}sets}(\operatorname{rr\text{-}ms}), u) \wedge \\ \quad \forall a, b.\; (((a \in \mathbb{R}) \wedge ((b \in \mathbb{R}) \wedge (a \leq b))) \Rightarrow (u(\operatorname{ccint}(a, b)) = (b - a))) \wedge \\ \quad s \in \operatorname{borel\text{-}sets}(\operatorname{rr\text{-}ms}) \Rightarrow \\ \quad \quad u(s) \\ \quad \quad =\; \operatorname{res}(\operatorname{method\text{-}i\text{-}outer}(\mathbb{R}, \operatorname{closed\text{-}intervals}, (\lambda \mathit{ols} \in \operatorname{closed\text{-}intervals}.\; \operatorname{interval\text{-}length}(\mathit{ols}))), \operatorname{caratheodory\text{-}sets}(\mathbb{R}, \operatorname{method\text{-}i\text{-}outer}(\mathbb{R}, \operatorname{closed\text{-}intervals}, (\lambda \mathit{ols} \in \operatorname{closed\text{-}intervals}.\; \operatorname{interval\text{-}length}(\mathit{ols})))))(s) \end{array}}}

\vnbbookfile{theorem-library/lebesgue-translation-invariant}{theorem-library/lebesgue-translation-invariant}
\noindent{\small ==================================================================== theorem-library/lebesgue-translation-invariant.scm -- MEASURE-5-A, 2026-10-10. A translation-invariant Borel measure on the line with 0 $<$ mu([0, 1]) = k $<$ +inf is k times Lebesgue measure on the Borel sets.}\par
\vnbprop{translate-ccint}{\vnbstmt{\begin{array}{@{}l@{}} \forall a, b, c. \\ \quad ((a \in \mathbb{R}) \wedge ((b \in \mathbb{R}) \wedge (c \in \mathbb{R}))) \Rightarrow  \\ \quad \quad (\operatorname{ccint}(a, b) + c) = \operatorname{ccint}((a + c), (b + c)) \end{array}}}
\vnbprop{invariant-measure-ccint-shift}{\vnbstmt{\begin{array}{@{}l@{}} \forall u, a, b. \\ \quad \operatorname{is\text{-}measure}(\mathbb{R}, \operatorname{borel\text{-}sets}(\operatorname{rr\text{-}ms}), u) \wedge \\ \quad \forall v, c. \\ \quad \quad ((v \in \operatorname{borel\text{-}sets}(\operatorname{rr\text{-}ms})) \wedge ((c \in \mathbb{R}) \wedge ((v + c) \in \operatorname{borel\text{-}sets}(\operatorname{rr\text{-}ms})))) \Rightarrow  \\ \quad \quad \quad u((v + c)) = u(v) \wedge \\ \quad a \in \mathbb{R} \wedge \\ \quad b \in \mathbb{R} \Rightarrow \\ \quad \quad u(\operatorname{ccint}(a, b)) = u(\operatorname{ccint}(0, (b - a))) \end{array}}}
\vnbprop{invariant-measure-points-null}{\vnbstmt{\begin{array}{@{}l@{}} \forall u.; \forall x \in \mathbb{R}. \\ \quad \operatorname{is\text{-}measure}(\mathbb{R}, \operatorname{borel\text{-}sets}(\operatorname{rr\text{-}ms}), u) \wedge \\ \quad \forall v, c. \\ \quad \quad ((v \in \operatorname{borel\text{-}sets}(\operatorname{rr\text{-}ms})) \wedge ((c \in \mathbb{R}) \wedge ((v + c) \in \operatorname{borel\text{-}sets}(\operatorname{rr\text{-}ms})))) \Rightarrow  \\ \quad \quad \quad u((v + c)) = u(v) \wedge \\ \quad u(\operatorname{ccint}(0, 1)) \in \mathbb{R} \Rightarrow \\ \quad \quad u(\operatorname{ccint}(x, x)) = 0 \end{array}}}
\vnbprop{null-points-ccint-split}{\vnbstmt{\begin{array}{@{}l@{}} \forall u, a, c, b. \\ \quad \operatorname{is\text{-}measure}(\mathbb{R}, \operatorname{borel\text{-}sets}(\operatorname{rr\text{-}ms}), u) \wedge \\ \quad \forall x \in \mathbb{R}.\; (u(\operatorname{ccint}(x, x)) = 0) \wedge \\ \quad a \in \mathbb{R} \wedge \\ \quad c \in \mathbb{R} \wedge \\ \quad b \in \mathbb{R} \wedge \\ \quad a \leq c \wedge \\ \quad c \leq b \Rightarrow \\ \quad \quad u(\operatorname{ccint}(a, b)) = \operatorname{eplus}(u(\operatorname{ccint}(a, c)), u(\operatorname{ccint}(c, b))) \end{array}}}
\vnbprop{invariant-measure-ccint-additive}{\vnbstmt{\begin{array}{@{}l@{}} \forall u, s, t. \\ \quad \operatorname{is\text{-}measure}(\mathbb{R}, \operatorname{borel\text{-}sets}(\operatorname{rr\text{-}ms}), u) \wedge \\ \quad \forall v, c. \\ \quad \quad ((v \in \operatorname{borel\text{-}sets}(\operatorname{rr\text{-}ms})) \wedge ((c \in \mathbb{R}) \wedge ((v + c) \in \operatorname{borel\text{-}sets}(\operatorname{rr\text{-}ms})))) \Rightarrow  \\ \quad \quad \quad u((v + c)) = u(v) \wedge \\ \quad u(\operatorname{ccint}(0, 1)) \in \mathbb{R} \wedge \\ \quad s \in \mathbb{R} \wedge \\ \quad t \in \mathbb{R} \wedge \\ \quad 0 \leq s \wedge \\ \quad 0 \leq t \Rightarrow \\ \quad \quad u(\operatorname{ccint}(0, (s + t))) = \operatorname{eplus}(u(\operatorname{ccint}(0, s)), u(\operatorname{ccint}(0, t))) \end{array}}}
\vnbprop{translation-invariant-measure-nn}{\vnbstmt{\begin{array}{@{}l@{}} \forall u.; \forall m \in \mathbb{N}. \\ \quad \operatorname{is\text{-}measure}(\mathbb{R}, \operatorname{borel\text{-}sets}(\operatorname{rr\text{-}ms}), u) \wedge \\ \quad \forall v, c. \\ \quad \quad ((v \in \operatorname{borel\text{-}sets}(\operatorname{rr\text{-}ms})) \wedge ((c \in \mathbb{R}) \wedge ((v + c) \in \operatorname{borel\text{-}sets}(\operatorname{rr\text{-}ms})))) \Rightarrow  \\ \quad \quad \quad u((v + c)) = u(v) \wedge \\ \quad u(\operatorname{ccint}(0, 1)) = 1 \Rightarrow \\ \quad \quad u(\operatorname{ccint}(0, m)) = m \end{array}}}
\vnbprop{translation-invariant-measure-finite}{\vnbstmt{\begin{array}{@{}l@{}} \forall u, t. \\ \quad \operatorname{is\text{-}measure}(\mathbb{R}, \operatorname{borel\text{-}sets}(\operatorname{rr\text{-}ms}), u) \wedge \\ \quad \forall v, c. \\ \quad \quad ((v \in \operatorname{borel\text{-}sets}(\operatorname{rr\text{-}ms})) \wedge ((c \in \mathbb{R}) \wedge ((v + c) \in \operatorname{borel\text{-}sets}(\operatorname{rr\text{-}ms})))) \Rightarrow  \\ \quad \quad \quad u((v + c)) = u(v) \wedge \\ \quad u(\operatorname{ccint}(0, 1)) = 1 \wedge \\ \quad t \in \mathbb{R} \wedge \\ \quad 0 \leq t \Rightarrow \\ \quad \quad u(\operatorname{ccint}(0, t)) \in \mathbb{R} \end{array}}}
\vnbprop{translation-invariant-measure-nn-mul}{\vnbstmt{\begin{array}{@{}l@{}} \forall j \in \mathbb{N}, u, s. \\ \quad \operatorname{is\text{-}measure}(\mathbb{R}, \operatorname{borel\text{-}sets}(\operatorname{rr\text{-}ms}), u) \wedge \\ \quad \forall v, c. \\ \quad \quad ((v \in \operatorname{borel\text{-}sets}(\operatorname{rr\text{-}ms})) \wedge ((c \in \mathbb{R}) \wedge ((v + c) \in \operatorname{borel\text{-}sets}(\operatorname{rr\text{-}ms})))) \Rightarrow  \\ \quad \quad \quad u((v + c)) = u(v) \wedge \\ \quad u(\operatorname{ccint}(0, 1)) = 1 \wedge \\ \quad s \in \mathbb{R} \wedge \\ \quad 0 \leq s \Rightarrow \\ \quad \quad u(\operatorname{ccint}(0, (j \cdot s))) = (j \cdot u(\operatorname{ccint}(0, s))) \end{array}}}
\vnbprop{translation-invariant-measure-ccint-zero}{\vnbstmt{\begin{array}{@{}l@{}} \forall u, t. \\ \quad \operatorname{is\text{-}measure}(\mathbb{R}, \operatorname{borel\text{-}sets}(\operatorname{rr\text{-}ms}), u) \wedge \\ \quad \forall v, c. \\ \quad \quad ((v \in \operatorname{borel\text{-}sets}(\operatorname{rr\text{-}ms})) \wedge ((c \in \mathbb{R}) \wedge ((v + c) \in \operatorname{borel\text{-}sets}(\operatorname{rr\text{-}ms})))) \Rightarrow  \\ \quad \quad \quad u((v + c)) = u(v) \wedge \\ \quad u(\operatorname{ccint}(0, 1)) = 1 \wedge \\ \quad t \in \mathbb{R} \wedge \\ \quad 0 \leq t \Rightarrow \\ \quad \quad u(\operatorname{ccint}(0, t)) = t \end{array}}}
\vnbprop{translation-invariant-measure-ccint}{\vnbstmt{\begin{array}{@{}l@{}} \forall u, a, b. \\ \quad \operatorname{is\text{-}measure}(\mathbb{R}, \operatorname{borel\text{-}sets}(\operatorname{rr\text{-}ms}), u) \wedge \\ \quad \forall v, c. \\ \quad \quad ((v \in \operatorname{borel\text{-}sets}(\operatorname{rr\text{-}ms})) \wedge ((c \in \mathbb{R}) \wedge ((v + c) \in \operatorname{borel\text{-}sets}(\operatorname{rr\text{-}ms})))) \Rightarrow  \\ \quad \quad \quad u((v + c)) = u(v) \wedge \\ \quad u(\operatorname{ccint}(0, 1)) = 1 \wedge \\ \quad a \in \mathbb{R} \wedge \\ \quad b \in \mathbb{R} \wedge \\ \quad a \leq b \Rightarrow \\ \quad \quad u(\operatorname{ccint}(a, b)) = (b - a) \end{array}}}
\vnbprop{translation-invariant-measure-lebesgue}{\vnbstmt{\begin{array}{@{}l@{}} \forall u, s. \\ \quad \operatorname{is\text{-}measure}(\mathbb{R}, \operatorname{borel\text{-}sets}(\operatorname{rr\text{-}ms}), u) \wedge \\ \quad \forall v, c. \\ \quad \quad ((v \in \operatorname{borel\text{-}sets}(\operatorname{rr\text{-}ms})) \wedge ((c \in \mathbb{R}) \wedge ((v + c) \in \operatorname{borel\text{-}sets}(\operatorname{rr\text{-}ms})))) \Rightarrow  \\ \quad \quad \quad u((v + c)) = u(v) \wedge \\ \quad u(\operatorname{ccint}(0, 1)) = 1 \wedge \\ \quad s \in \operatorname{borel\text{-}sets}(\operatorname{rr\text{-}ms}) \Rightarrow \\ \quad \quad u(s) \\ \quad \quad =\; \operatorname{res}(\operatorname{method\text{-}i\text{-}outer}(\mathbb{R}, \operatorname{closed\text{-}intervals}, (\lambda \mathit{ols} \in \operatorname{closed\text{-}intervals}.\; \operatorname{interval\text{-}length}(\mathit{ols}))), \operatorname{caratheodory\text{-}sets}(\mathbb{R}, \operatorname{method\text{-}i\text{-}outer}(\mathbb{R}, \operatorname{closed\text{-}intervals}, (\lambda \mathit{ols} \in \operatorname{closed\text{-}intervals}.\; \operatorname{interval\text{-}length}(\mathit{ols})))))(s) \end{array}}}
\vnbprop{lebesgue-translation-invariant}{\vnbstmt{\begin{array}{@{}l@{}} \forall u, k, s. \\ \quad \operatorname{is\text{-}measure}(\mathbb{R}, \operatorname{borel\text{-}sets}(\operatorname{rr\text{-}ms}), u) \wedge \\ \quad \forall v, c. \\ \quad \quad ((v \in \operatorname{borel\text{-}sets}(\operatorname{rr\text{-}ms})) \wedge ((c \in \mathbb{R}) \wedge ((v + c) \in \operatorname{borel\text{-}sets}(\operatorname{rr\text{-}ms})))) \Rightarrow  \\ \quad \quad \quad u((v + c)) = u(v) \wedge \\ \quad k \in \mathbb{R} \wedge \\ \quad 0 < k \wedge \\ \quad u(\operatorname{ccint}(0, 1)) = k \wedge \\ \quad s \in \operatorname{borel\text{-}sets}(\operatorname{rr\text{-}ms}) \Rightarrow \\ \quad \quad u(s) \\ \quad \quad =\; \operatorname{etimes}(k, \\ \quad \quad \quad \operatorname{res}(\operatorname{method\text{-}i\text{-}outer}(\mathbb{R}, \operatorname{closed\text{-}intervals}, (\lambda \mathit{ols} \in \operatorname{closed\text{-}intervals}.\; \operatorname{interval\text{-}length}(\mathit{ols}))), \operatorname{caratheodory\text{-}sets}(\mathbb{R}, \operatorname{method\text{-}i\text{-}outer}(\mathbb{R}, \operatorname{closed\text{-}intervals}, (\lambda \mathit{ols} \in \operatorname{closed\text{-}intervals}.\; \operatorname{interval\text{-}length}(\mathit{ols})))))(s)) \end{array}}}

\vnbbookfile{structure-library/lp-inequalities}{structure-library/lp-inequalities}
\noindent{\small lp-inequalities.scm -- POWERS IN [0,+inf] AND THE p-TH POWER INTEGRAL I\_p. Definitions only; nothing here is asserted and nothing is proved.  The laws and the inequalities of Hoelder and Minkowski are proved in theorem-library/lp-inequalities.scm (MEASURE-4-B, 2026-10-10).}\par
\par\medskip\noindent\textbf{Definition} (\texttt{epow}).\;${\operatorname{ept\_}}^{\operatorname{epp\_}} :=$ $\left(\text{if }(\operatorname{ept\_} = \operatorname{pos\text{-}inf})\text{ then }\operatorname{pos\text{-}inf}\text{ else }{\operatorname{ept\_}}^{\operatorname{epp\_}}\right)$
\par\medskip\noindent\textbf{Definition} (\texttt{ptwise-epow}).\;$\operatorname{ptwise\text{-}epow}(\operatorname{epo\_}, \operatorname{epf\_}, \operatorname{epp\_}) :=$ $(\lambda \operatorname{epx\_} \in \operatorname{epo\_}.\; {\operatorname{epf\_}(\operatorname{epx\_})}^{\operatorname{epp\_}})$
\par\medskip\noindent\textbf{Definition} (\texttt{ep-norm}).\;$I_{\operatorname{epp\_}}({\operatorname{epf\_}}) :=$
\begin{equation*}
{\operatorname{integral}(\omega, \mathit{ca}, \mu, \operatorname{ptwise\text{-}epow}(\omega, \operatorname{epf\_}, \operatorname{epp\_}))}^{\operatorname{recip}(\operatorname{epp\_})}
\end{equation*}


\vnbbookfile{theorem-library/lp-inequalities}{theorem-library/lp-inequalities}
\noindent{\small theorem-library/lp-inequalities.scm -- POWERS IN [0,+inf], HOELDER AND MINKOWSKI, FUNCTIONS EQUAL ALMOST EVERYWHERE.  MEASURE-4-B, 2026-10-10.}\par
\vnbprop{epow-unfold}{\vnbstmt{\forall t, p.\; ({t}^{p} \simeq \left(\text{if }(t = \operatorname{pos\text{-}inf})\text{ then }\operatorname{pos\text{-}inf}\text{ else }{t}^{p}\right))}}
\vnbprop{ep-norm-unfold}{\vnbstmt{\begin{array}{@{}l@{}} \forall a, b, u, p, f. \\ \quad (I_{p}({f}) \simeq {\operatorname{integral}(a, b, u, \operatorname{ptwise\text{-}epow}(a, f, p))}^{\operatorname{recip}(p)}) \end{array}}}
\vnbprop{rr-lt-pos-inf}{\vnbstmt{\forall x \in \mathbb{R}.\; (x < \operatorname{pos\text{-}inf})}}
\vnbprop{epow-pos-inf}{\vnbstmt{\forall p.\; ({\operatorname{pos\text{-}inf}}^{p} = \operatorname{pos\text{-}inf})}}
\vnbprop{epow-real}{\vnbstmt{\forall t \in \mathbb{R}, p.\; ({t}^{p} \simeq {t}^{p})}}
\vnbprop{epow-zero}{\vnbstmt{\forall p \in \mathbb{R}.\; ((0 < p) \Rightarrow ({0}^{p} = 0))}}
\vnbprop{epow-one-base}{\vnbstmt{\forall p \in \mathbb{R}.\; ({1}^{p} = 1)}}
\vnbprop{epow-pos-real}{\vnbstmt{\forall t, p \in \mathbb{R}.\; ((0 < t) \Rightarrow (({t}^{p} \in \mathbb{R}) \wedge (0 < {t}^{p})))}}
\vnbprop{epow-real-in}{\vnbstmt{\begin{array}{@{}l@{}} \forall t, p \in \mathbb{R}. \\ \quad ((0 \leq t) \wedge (0 < p)) \Rightarrow (({t}^{p} \in \mathbb{R}) \wedge (0 \leq {t}^{p})) \end{array}}}
\vnbprop{epow-in}{\vnbstmt{\forall t \in \operatorname{rr\text{-}pos\text{-}star}, p \in \mathbb{R}.\; ((0 < p) \Rightarrow ({t}^{p} \in \operatorname{rr\text{-}pos\text{-}star}))}}
\vnbprop{epow-strict-mono}{\vnbstmt{\begin{array}{@{}l@{}} \forall s, t \in \operatorname{rr\text{-}pos\text{-}star}, p \in \mathbb{R}. \\ \quad ((0 < p) \wedge (s < t)) \Rightarrow ({s}^{p} < {t}^{p}) \end{array}}}
\vnbprop{epow-mono}{\vnbstmt{\begin{array}{@{}l@{}} \forall s, t \in \operatorname{rr\text{-}pos\text{-}star}, p \in \mathbb{R}. \\ \quad ((0 < p) \wedge (s \leq t)) \Rightarrow ({s}^{p} \leq {t}^{p}) \end{array}}}
\vnbprop{epow-epow}{\vnbstmt{\begin{array}{@{}l@{}} \forall t \in \operatorname{rr\text{-}pos\text{-}star}, a, b \in \mathbb{R}. \\ \quad ((0 < a) \wedge (0 < b)) \Rightarrow ({{t}^{a}}^{b} = {t}^{(a \cdot b)}) \end{array}}}
\vnbprop{epow-one}{\vnbstmt{\forall t \in \operatorname{rr\text{-}pos\text{-}star}.\; ({t}^{1} = t)}}
\vnbprop{epow-inverse}{\vnbstmt{\begin{array}{@{}l@{}} \forall t \in \operatorname{rr\text{-}pos\text{-}star}, p \in \mathbb{R}. \\ \quad (0 < p) \Rightarrow (({{t}^{p}}^{\operatorname{recip}(p)} = t) \wedge ({{t}^{\operatorname{recip}(p)}}^{p} = t)) \end{array}}}
\vnbprop{epow-eq-zero}{\vnbstmt{\forall t \in \operatorname{rr\text{-}pos\text{-}star}, p \in \mathbb{R}.\; (((0 < p) \wedge ({t}^{p} = 0)) \Rightarrow (t = 0))}}
\vnbprop{epow-eq-pos-inf}{\vnbstmt{\begin{array}{@{}l@{}} \forall t \in \operatorname{rr\text{-}pos\text{-}star}, p \in \mathbb{R}. \\ \quad ((0 < p) \wedge ({t}^{p} = \operatorname{pos\text{-}inf})) \Rightarrow (t = \operatorname{pos\text{-}inf}) \end{array}}}
\vnbprop{epow-etimes}{\vnbstmt{\begin{array}{@{}l@{}} \forall c \in \mathbb{R}, u \in \operatorname{rr\text{-}pos\text{-}star}, p \in \mathbb{R}. \\ \quad ((0 < c) \wedge (0 < p)) \Rightarrow ({\operatorname{etimes}(c, u)}^{p} = \operatorname{etimes}({c}^{p}, {u}^{p})) \end{array}}}
\vnbprop{ptwise-epow-value}{\vnbstmt{\begin{array}{@{}l@{}} \forall a.; \forall f \in (a \to \operatorname{rr\text{-}pos\text{-}star}), p \in \mathbb{R}, x \in a. \\ \quad (0 < p) \Rightarrow (\operatorname{ptwise\text{-}epow}(a, f, p)(x) = {f(x)}^{p}) \end{array}}}
\vnbprop{measurable-fn-epow}{\vnbstmt{\begin{array}{@{}l@{}} \forall a, b, f, p. \\ \quad (\operatorname{is\text{-}measurable\text{-}fn}(a, b, f) \wedge ((p \in \mathbb{R}) \wedge (0 < p))) \Rightarrow  \\ \quad \quad \operatorname{is\text{-}measurable\text{-}fn}(a, b, \operatorname{ptwise\text{-}epow}(a, f, p)) \end{array}}}
\vnbprop{ptwise-eplus-value}{\vnbstmt{\begin{array}{@{}l@{}} \forall a, f, g.; \forall x \in a. \\ \quad ((f(x) \in \operatorname{rr\text{-}pos\text{-}star}) \wedge (g(x) \in \operatorname{rr\text{-}pos\text{-}star})) \Rightarrow  \\ \quad \quad \operatorname{ptwise\text{-}eplus}(a, f, g)(x) = \operatorname{eplus}(f(x), g(x)) \end{array}}}
\vnbprop{ptwise-etimes-value}{\vnbstmt{\begin{array}{@{}l@{}} \forall a, f, g.; \forall x \in a. \\ \quad ((f(x) \in \operatorname{rr\text{-}pos\text{-}star}) \wedge (g(x) \in \operatorname{rr\text{-}pos\text{-}star})) \Rightarrow  \\ \quad \quad \operatorname{ptwise\text{-}etimes}(a, f, g)(x) = \operatorname{etimes}(f(x), g(x)) \end{array}}}
\vnbprop{ptwise-scale-value}{\vnbstmt{\begin{array}{@{}l@{}} \forall a, c, f.; \forall x \in a. \\ \quad ((c \in \operatorname{rr\text{-}pos\text{-}star}) \wedge (f(x) \in \operatorname{rr\text{-}pos\text{-}star})) \Rightarrow  \\ \quad \quad \operatorname{ptwise\text{-}scale}(a, c, f)(x) = \operatorname{etimes}(c, f(x)) \end{array}}}
\vnbprop{young-scaled}{\vnbstmt{\begin{array}{@{}l@{}} \forall c, d \in \mathbb{R}, u, v \in \operatorname{rr\text{-}pos\text{-}star}, p, q \in \mathbb{R}. \\ \quad ((0 < c) \wedge ((0 < d) \wedge ((1 < p) \wedge ((1 < q) \wedge ((/(1, p) + /(1, q)) = 1))))) \Rightarrow  \\ \quad \quad \operatorname{etimes}((c \cdot d), \operatorname{etimes}(u, v)) \\ \quad \quad \leq\; \operatorname{eplus}(\operatorname{etimes}((\operatorname{recip}(p) \cdot {c}^{p}), {u}^{p}), \operatorname{etimes}((\operatorname{recip}(q) \cdot {d}^{q}), {v}^{q})) \end{array}}}
\vnbprop{young-extended}{\vnbstmt{\begin{array}{@{}l@{}} \forall a, b \in \operatorname{rr\text{-}pos\text{-}star}, p, q \in \mathbb{R}. \\ \quad ((1 < p) \wedge ((1 < q) \wedge ((/(1, p) + /(1, q)) = 1))) \Rightarrow  \\ \quad \quad \operatorname{etimes}(a, b) \leq \operatorname{eplus}(\operatorname{etimes}(\operatorname{recip}(p), {a}^{p}), \operatorname{etimes}(\operatorname{recip}(q), {b}^{q})) \end{array}}}
\vnbprop{etimes-eq-zero-pos}{\vnbstmt{\forall c \in \mathbb{R}, m \in \operatorname{rr\text{-}pos\text{-}star}.\; (((0 < c) \wedge (\operatorname{etimes}(c, m) = 0)) \Rightarrow (m = 0))}}
\vnbprop{integral-markov}{\vnbstmt{\begin{array}{@{}l@{}} \forall a, b, u, h, c. \\ \quad (\operatorname{is\text{-}measure}(a, b, u) \wedge (\operatorname{is\text{-}measurable\text{-}fn}(a, b, h) \wedge ((c \in \mathbb{R}) \wedge (0 \leq c)))) \Rightarrow  \\ \quad \quad \operatorname{etimes}(c, u(\{x \in a : (c \leq h(x))\})) \leq \operatorname{integral}(a, b, u, h) \end{array}}}
\vnbprop{integral-zero-null}{\vnbstmt{\begin{array}{@{}l@{}} \forall a, b, u, h. \\ \quad (\operatorname{is\text{-}measure}(a, b, u) \wedge (\operatorname{is\text{-}measurable\text{-}fn}(a, b, h) \wedge (\operatorname{integral}(a, b, u, h) = 0))) \Rightarrow  \\ \quad \quad u(\{x \in a : (0 < h(x))\}) = 0 \end{array}}}
\vnbprop{ep-norm-zero-null}{\vnbstmt{\begin{array}{@{}l@{}} \forall a, b, u, f, p. \\ \quad \operatorname{is\text{-}measure}(a, b, u) \wedge \\ \quad \operatorname{is\text{-}measurable\text{-}fn}(a, b, f) \wedge \\ \quad p \in \mathbb{R} \wedge \\ \quad 0 < p \wedge \\ \quad I_{p}({f}) = 0 \Rightarrow \\ \quad \quad (\{x \in a : (0 < f(x))\} \in b) \wedge (u(\{x \in a : (0 < f(x))\}) = 0) \end{array}}}
\vnbprop{integral-ae-eq}{\vnbstmt{\begin{array}{@{}l@{}} \forall a, b, u, c, f, g. \\ \quad \operatorname{is\text{-}measure}(a, b, u) \wedge \\ \quad c \in b \wedge \\ \quad u(c) = 0 \wedge \\ \quad \operatorname{is\text{-}measurable\text{-}fn}(a, b, f) \wedge \\ \quad \operatorname{is\text{-}measurable\text{-}fn}(a, b, g) \wedge \\ \quad \forall x \in a.\; (\neg (x \in c) \Rightarrow (f(x) = g(x))) \Rightarrow \\ \quad \quad \operatorname{integral}(a, b, u, f) = \operatorname{integral}(a, b, u, g) \end{array}}}
\vnbprop{integral-ptwise-eq}{\vnbstmt{\begin{array}{@{}l@{}} \forall a, b, u, f, g. \\ \quad \operatorname{is\text{-}measure}(a, b, u) \wedge \\ \quad \operatorname{is\text{-}measurable\text{-}fn}(a, b, f) \wedge \\ \quad \operatorname{is\text{-}measurable\text{-}fn}(a, b, g) \wedge \\ \quad \forall x \in a.\; (f(x) = g(x)) \Rightarrow \\ \quad \quad \operatorname{integral}(a, b, u, f) = \operatorname{integral}(a, b, u, g) \end{array}}}
\vnbprop{integral-finite-ae-finite}{\vnbstmt{\begin{array}{@{}l@{}} \forall a, b, u, f. \\ \quad (\operatorname{is\text{-}measure}(a, b, u) \wedge (\operatorname{is\text{-}measurable\text{-}fn}(a, b, f) \wedge (\operatorname{integral}(a, b, u, f) < \operatorname{pos\text{-}inf}))) \Rightarrow  \\ \quad \quad (\{x \in a : (f(x) = \operatorname{pos\text{-}inf})\} \in b) \wedge (u(\{x \in a : (f(x) = \operatorname{pos\text{-}inf})\}) = 0) \end{array}}}
\vnbprop{hoelder}{\vnbstmt{\begin{array}{@{}l@{}} \forall a, b, u, f, g.; \forall p, q \in \mathbb{R}. \\ \quad \operatorname{is\text{-}measure}(a, b, u) \wedge \\ \quad \operatorname{is\text{-}measurable\text{-}fn}(a, b, f) \wedge \\ \quad \operatorname{is\text{-}measurable\text{-}fn}(a, b, g) \wedge \\ \quad 1 < p \wedge \\ \quad 1 < q \wedge \\ \quad (/(1, p) + /(1, q)) = 1 \Rightarrow \\ \quad \quad \operatorname{integral}(a, b, u, \operatorname{ptwise\text{-}etimes}(a, f, g)) \leq \operatorname{etimes}(I_{p}({f}), I_{q}({g})) \end{array}}}
\vnbprop{conjugate-exponent}{\vnbstmt{\begin{array}{@{}l@{}} \forall p \in \mathbb{R}. \\ \quad (1 < p) \Rightarrow  \\ \quad \quad (p \cdot \operatorname{recip}((p - 1))) \in \mathbb{R} \wedge \\ \quad \quad 1 < (p \cdot \operatorname{recip}((p - 1))) \wedge \\ \quad \quad (/(1, p) + /(1, (p \cdot \operatorname{recip}((p - 1))))) = 1 \wedge \\ \quad \quad ((p - 1) \cdot (p \cdot \operatorname{recip}((p - 1)))) = p \end{array}}}
\vnbprop{ep-norm-one}{\vnbstmt{\begin{array}{@{}l@{}} \forall a, b, u, h. \\ \quad (\operatorname{is\text{-}measure}(a, b, u) \wedge \operatorname{is\text{-}measurable\text{-}fn}(a, b, h)) \Rightarrow  \\ \quad \quad I_{1}({h}) = \operatorname{integral}(a, b, u, h) \end{array}}}
\vnbprop{epow-split}{\vnbstmt{\begin{array}{@{}l@{}} \forall t \in \operatorname{rr\text{-}pos\text{-}star}, p \in \mathbb{R}. \\ \quad (1 < p) \Rightarrow ({t}^{p} = \operatorname{etimes}({t}^{(p - 1)}, t)) \end{array}}}
\vnbprop{etimes-distrib-left-star}{\vnbstmt{\begin{array}{@{}l@{}} \forall c \in \operatorname{rr\text{-}pos\text{-}star}, x, y. \\ \quad ((x \in \operatorname{rr\text{-}pos\text{-}star}) \wedge (y \in \operatorname{rr\text{-}pos\text{-}star})) \Rightarrow  \\ \quad \quad \operatorname{etimes}(c, \operatorname{eplus}(x, y)) = \operatorname{eplus}(\operatorname{etimes}(c, x), \operatorname{etimes}(c, y)) \end{array}}}
\vnbprop{epow-eplus-le}{\vnbstmt{\begin{array}{@{}l@{}} \forall u, v \in \operatorname{rr\text{-}pos\text{-}star}, p \in \mathbb{R}. \\ \quad (0 < p) \Rightarrow  \\ \quad \quad {\operatorname{eplus}(u, v)}^{p} \leq \operatorname{etimes}({2}^{p}, \operatorname{eplus}({u}^{p}, {v}^{p})) \end{array}}}
\vnbprop{minkowski}{\vnbstmt{\begin{array}{@{}l@{}} \forall a, b, u, f, g.; \forall p \in \mathbb{R}. \\ \quad (\operatorname{is\text{-}measure}(a, b, u) \wedge (\operatorname{is\text{-}measurable\text{-}fn}(a, b, f) \wedge (\operatorname{is\text{-}measurable\text{-}fn}(a, b, g) \wedge (1 \leq p)))) \Rightarrow  \\ \quad \quad I_{p}({\operatorname{ptwise\text{-}eplus}(a, f, g)}) \leq \operatorname{eplus}(I_{p}({f}), I_{p}({g})) \end{array}}}

\vnbbookfile{structure-library/lp-seminorm}{structure-library/lp-seminorm}
\noindent{\small lp-seminorm.scm -- THE L\textasciicircum{}p SEMINORM OF A REAL-VALUED FUNCTION. Definitions only; nothing here is asserted and nothing is proved.  The laws are proved in theorem-library/lp-seminorm.scm (MEASURE-4-C, 2026-10-10).}\par
\par\medskip\noindent\textbf{Definition} (\texttt{lp-norm}).\;$\|{\operatorname{lpf\_}}\|_{\operatorname{lpp\_}} :=$ $I_{\operatorname{lpp\_}}({\operatorname{abs\text{-}fn}(\omega, \operatorname{lpf\_})})$

\vnbbookfile{theorem-library/lp-seminorm}{theorem-library/lp-seminorm}
\noindent{\small lp-seminorm.scm -- THE L\textasciicircum{}p SEMINORM: the seminorm laws (MEASURE-4-C, 2026-10-10).  Vocabulary: structure-library/lp-seminorm.scm (LP-NORM = EP-NORM of ABS-FN).}\par
\vnbprop{lp-norm-unfold}{\vnbstmt{\begin{array}{@{}l@{}} \forall a, b, u, p, f. \\ \quad (\|{f}\|_{p} \simeq I_{p}({\operatorname{abs\text{-}fn}(a, f)})) \end{array}}}
\vnbprop{real-measurable-add}{\vnbstmt{\begin{array}{@{}l@{}} \forall a, b, f, g. \\ \quad (\operatorname{is\text{-}real\text{-}measurable\text{-}fn}(a, b, f) \wedge \operatorname{is\text{-}real\text{-}measurable\text{-}fn}(a, b, g)) \Rightarrow  \\ \quad \quad \operatorname{is\text{-}real\text{-}measurable\text{-}fn}(a, b, (\lambda x \in a.\; (f(x) + g(x)))) \end{array}}}
\vnbprop{real-measurable-sub}{\vnbstmt{\begin{array}{@{}l@{}} \forall a, b, f, g. \\ \quad (\operatorname{is\text{-}real\text{-}measurable\text{-}fn}(a, b, f) \wedge \operatorname{is\text{-}real\text{-}measurable\text{-}fn}(a, b, g)) \Rightarrow  \\ \quad \quad \operatorname{is\text{-}real\text{-}measurable\text{-}fn}(a, b, (\lambda x \in a.\; (f(x) - g(x)))) \end{array}}}
\vnbprop{ep-norm-mono}{\vnbstmt{\begin{array}{@{}l@{}} \forall a, b, u, p, c, d. \\ \quad \operatorname{is\text{-}measure}(a, b, u) \wedge \\ \quad \operatorname{is\text{-}measurable\text{-}fn}(a, b, c) \wedge \\ \quad \operatorname{is\text{-}measurable\text{-}fn}(a, b, d) \wedge \\ \quad \forall x \in a.\; (c(x) \leq d(x)) \wedge \\ \quad p \in \mathbb{R} \wedge \\ \quad 0 < p \Rightarrow \\ \quad \quad I_{p}({c}) \leq I_{p}({d}) \end{array}}}
\vnbprop{ep-norm-scale}{\vnbstmt{\begin{array}{@{}l@{}} \forall a, b, u, p, k, h. \\ \quad (\operatorname{is\text{-}measure}(a, b, u) \wedge (\operatorname{is\text{-}measurable\text{-}fn}(a, b, h) \wedge ((k \in \mathbb{R}) \wedge ((0 \leq k) \wedge ((p \in \mathbb{R}) \wedge (0 < p)))))) \Rightarrow  \\ \quad \quad I_{p}({\operatorname{ptwise\text{-}scale}(a, k, h)}) = \operatorname{etimes}(k, I_{p}({h})) \end{array}}}
\vnbprop{lp-norm-homogeneous}{\vnbstmt{\begin{array}{@{}l@{}} \forall a, b, u, p, f, c. \\ \quad (\operatorname{is\text{-}measure}(a, b, u) \wedge (\operatorname{is\text{-}real\text{-}measurable\text{-}fn}(a, b, f) \wedge ((c \in \mathbb{R}) \wedge ((p \in \mathbb{R}) \wedge (0 < p))))) \Rightarrow  \\ \quad \quad \|{(\lambda x \in a.\; (c \cdot f(x)))}\|_{p} = \operatorname{etimes}(\lvert c \rvert, \|{f}\|_{p}) \end{array}}}
\vnbprop{lp-norm-triangle}{\vnbstmt{\begin{array}{@{}l@{}} \forall a, b, u, p, f, g. \\ \quad \operatorname{is\text{-}measure}(a, b, u) \wedge \\ \quad \operatorname{is\text{-}real\text{-}measurable\text{-}fn}(a, b, f) \wedge \\ \quad \operatorname{is\text{-}real\text{-}measurable\text{-}fn}(a, b, g) \wedge \\ \quad p \in \mathbb{R} \wedge \\ \quad 1 \leq p \Rightarrow \\ \quad \quad \|{(\lambda x \in a.\; (f(x) + g(x)))}\|_{p} \\ \quad \quad \leq\; \operatorname{eplus}(\|{f}\|_{p}, \|{g}\|_{p}) \end{array}}}
\vnbprop{lp-norm-zero-iff}{\vnbstmt{\begin{array}{@{}l@{}} \forall a, b, u, p, f, g. \\ \quad \operatorname{is\text{-}measure}(a, b, u) \wedge \\ \quad \operatorname{is\text{-}real\text{-}measurable\text{-}fn}(a, b, f) \wedge \\ \quad \operatorname{is\text{-}real\text{-}measurable\text{-}fn}(a, b, g) \wedge \\ \quad p \in \mathbb{R} \wedge \\ \quad 0 < p \Rightarrow \\ \quad \quad \|{(\lambda \mathit{sx} \in a.\; (f(\mathit{sx}) - g(\mathit{sx})))}\|_{p} = 0 \\ \quad \quad \Leftrightarrow\; \exists s. \\ \quad \quad \quad \operatorname{is\text{-}null\text{-}set}(a, b, u, s) \wedge \forall x \in a.\; (\neg (x \in s) \Rightarrow (f(x) = g(x))) \end{array}}}

\vnbbookfile{structure-library/product-measure}{structure-library/product-measure}
\noindent{\small product-measure.scm -- PRODUCT SIGMA-ALGEBRAS, SECTIONS, DYNKIN SYSTEMS. Vocabulary only: no proof, no support.  Agent FUBINI-1, 2026-10-10 (scratchpad/triage/FUBINI-1.md).  The proofs are theorem-library/product-measure.scm (prefix pms-) and theorem-library/fubini.scm (prefix fub-).}\par
\par\medskip\noindent\textbf{Definition} (\texttt{rectangles}).\;$\operatorname{rectangles}(\mathit{xs}, \mathit{ca}, \mathit{ys}, \mathit{cb}) :=$
\begin{equation*}
\{\operatorname{rcr\_} \in \operatorname{power}((\mathit{xs} \times \mathit{ys})) : \exists \operatorname{rca\_} \in \mathit{ca}, \operatorname{rcb\_} \in \mathit{cb}.\; (\operatorname{rcr\_} = (\operatorname{rca\_} \times \operatorname{rcb\_}))\}
\end{equation*}

\par\medskip\noindent\textbf{Definition} (\texttt{product-sigma}).\;$\operatorname{product\text{-}sigma}(\mathit{xs}, \mathit{ca}, \mathit{ys}, \mathit{cb}) :=$
\begin{equation*}
\operatorname{sigma\text{-}generated}((\mathit{xs} \times \mathit{ys}), \operatorname{rectangles}(\mathit{xs}, \mathit{ca}, \mathit{ys}, \mathit{cb}))
\end{equation*}

\par\medskip\noindent\textbf{Definition} (\texttt{x-section}).\;$\operatorname{x\text{-}section}(\mathit{ys}, \mathit{cs}, x) :=$ $\{\operatorname{xsy\_} \in \mathit{ys} : (\langle x, \operatorname{xsy\_} \rangle \in \mathit{cs})\}$
\par\medskip\noindent\textbf{Definition} (\texttt{y-section}).\;$\operatorname{y\text{-}section}(\mathit{xs}, \mathit{cs}, y) :=$ $\{\operatorname{ysx\_} \in \mathit{xs} : (\langle \operatorname{ysx\_}, y \rangle \in \mathit{cs})\}$
\par\medskip\noindent\textbf{Definition} (\texttt{fn-x-section}).\;$\operatorname{fn\text{-}x\text{-}section}(\mathit{ys}, f, x) :=$ $(\lambda \operatorname{fsy\_} \in \mathit{ys}.\; f(\langle x, \operatorname{fsy\_} \rangle))$
\par\medskip\noindent\textbf{Definition} (\texttt{fn-y-section}).\;$\operatorname{fn\text{-}y\text{-}section}(\mathit{xs}, f, y) :=$ $(\lambda \operatorname{fsx\_} \in \mathit{xs}.\; f(\langle \operatorname{fsx\_}, y \rangle))$
\par\medskip\noindent\textbf{Definition} (\texttt{is-dynkin-system}).\;For all $\omega$ and $\mathit{dd}$, 
\begin{equation*}
\begin{array}{@{}l@{}} \operatorname{is\text{-}dynkin\text{-}system}(\omega, \mathit{dd}) \\ \Leftrightarrow\; \omega \in \mathbf{Set} \wedge \\ \mathit{dd} \subseteq \operatorname{power}(\omega) \wedge \\ \omega \in \mathit{dd} \wedge \\ \forall \operatorname{dya\_}, \operatorname{dyb\_} \in \mathit{dd}. \\ \quad (\operatorname{dya\_} \subseteq \operatorname{dyb\_}) \Rightarrow (\operatorname{complement\text{-}in}(\operatorname{dyb\_}, \operatorname{dya\_}) \in \mathit{dd}) \wedge \\ \forall \operatorname{dyf\_} \in (\mathbb{N} \to \mathit{dd}). \\ \quad \forall \operatorname{dyn\_} \in \mathbb{N}.\; (\operatorname{dyf\_}(\operatorname{dyn\_}) \subseteq \operatorname{dyf\_}(\operatorname{succ}(\operatorname{dyn\_}))) \Rightarrow  \\ \quad \quad \operatorname{big\text{-}union}(\operatorname{dyk\_}, \mathbb{N}, \operatorname{dyf\_}(\operatorname{dyk\_})) \in \mathit{dd} \end{array}
\end{equation*}

\par\medskip\noindent\textbf{Definition} (\texttt{is-pi-system}).\;For every $\mathit{pp}$, 
\begin{equation*}
\operatorname{is\text{-}pi\text{-}system}(\mathit{pp}) \Leftrightarrow \forall \operatorname{pya\_}, \operatorname{pyb\_} \in \mathit{pp}.\; ((\operatorname{pya\_} \cap \operatorname{pyb\_}) \in \mathit{pp})
\end{equation*}

\par\medskip\noindent\textbf{Definition} (\texttt{dynkin-generated}).\;$\operatorname{dynkin\text{-}generated}(\omega, \mathit{ce}) :=$
\begin{equation*}
\{\operatorname{dga\_} \in \operatorname{power}(\omega) : \forall \operatorname{dgb\_}.\; ((\operatorname{is\text{-}dynkin\text{-}system}(\omega, \operatorname{dgb\_}) \wedge (\mathit{ce} \subseteq \operatorname{dgb\_})) \Rightarrow (\operatorname{dga\_} \in \operatorname{dgb\_}))\}
\end{equation*}

\par\medskip\noindent\textbf{Definition} (\texttt{x-section-measure}).\;$\operatorname{x\text{-}section\text{-}measure}(\mathit{xs}, \mathit{ys}, \nu, \mathit{cs}) :=$ $(\lambda \operatorname{smx\_} \in \mathit{xs}.\; \nu(\operatorname{x\text{-}section}(\mathit{ys}, \mathit{cs}, \operatorname{smx\_})))$
\par\medskip\noindent\textbf{Definition} (\texttt{y-section-measure}).\;$\operatorname{y\text{-}section\text{-}measure}(\mathit{xs}, \mathit{ys}, \mu, \mathit{cs}) :=$ $(\lambda \operatorname{smy\_} \in \mathit{ys}.\; \mu(\operatorname{y\text{-}section}(\mathit{xs}, \mathit{cs}, \operatorname{smy\_})))$
\par\medskip\noindent\textbf{Definition} (\texttt{is-fubini-set}).\;For all $\mathit{xs}$, $\mathit{ca}$, $\mu$, $\mathit{ys}$, $\mathit{cb}$, $\nu$ and $\mathit{cs}$, 
\begin{equation*}
\begin{array}{@{}l@{}} \operatorname{is\text{-}fubini\text{-}set}(\mathit{xs}, \mathit{ca}, \mu, \mathit{ys}, \mathit{cb}, \nu, \mathit{cs}) \\ \Leftrightarrow\; \operatorname{is\text{-}measurable\text{-}fn}(\mathit{xs}, \mathit{ca}, \operatorname{x\text{-}section\text{-}measure}(\mathit{xs}, \mathit{ys}, \nu, \mathit{cs})) \wedge \\ \operatorname{is\text{-}measurable\text{-}fn}(\mathit{ys}, \mathit{cb}, \operatorname{y\text{-}section\text{-}measure}(\mathit{xs}, \mathit{ys}, \mu, \mathit{cs})) \wedge \\ \operatorname{integral}(\mathit{xs}, \mathit{ca}, \mu, \operatorname{x\text{-}section\text{-}measure}(\mathit{xs}, \mathit{ys}, \nu, \mathit{cs})) \\ =\; \operatorname{integral}(\mathit{ys}, \mathit{cb}, \nu, \operatorname{y\text{-}section\text{-}measure}(\mathit{xs}, \mathit{ys}, \mu, \mathit{cs})) \end{array}
\end{equation*}

\par\medskip\noindent\textbf{Definition} (\texttt{x-section-integral}).\;$\operatorname{x\text{-}section\text{-}integral}(\mathit{xs}, \mathit{ys}, \mathit{cb}, \nu, f) :=$
\begin{equation*}
(\lambda \operatorname{six\_} \in \mathit{xs}.\; \operatorname{integral}(\mathit{ys}, \mathit{cb}, \nu, \operatorname{fn\text{-}x\text{-}section}(\mathit{ys}, f, \operatorname{six\_})))
\end{equation*}

\par\medskip\noindent\textbf{Definition} (\texttt{y-section-integral}).\;$\operatorname{y\text{-}section\text{-}integral}(\mathit{xs}, \mathit{ca}, \mu, \mathit{ys}, f) :=$
\begin{equation*}
(\lambda \operatorname{siy\_} \in \mathit{ys}.\; \operatorname{integral}(\mathit{xs}, \mathit{ca}, \mu, \operatorname{fn\text{-}y\text{-}section}(\mathit{xs}, f, \operatorname{siy\_})))
\end{equation*}

\par\medskip\noindent\textbf{Definition} (\texttt{product-measure}).\;$\operatorname{product\text{-}measure}(\mathit{xs}, \mathit{ca}, \mu, \mathit{ys}, \mathit{cb}, \nu) :=$
\begin{equation*}
\begin{array}{@{}l@{}} (\lambda \operatorname{pmc\_} \in \operatorname{product\text{-}sigma}(\mathit{xs}, \mathit{ca}, \mathit{ys}, \mathit{cb}).\; \\ \quad \operatorname{integral}(\mathit{xs}, \mathit{ca}, \mu, \operatorname{x\text{-}section\text{-}measure}(\mathit{xs}, \mathit{ys}, \nu, \operatorname{pmc\_}))) \end{array}
\end{equation*}


\vnbbookfile{theorem-library/product-measure}{theorem-library/product-measure}
\noindent{\small ==================================================================== theorem-library/product-measure.scm -- PRODUCT SIGMA-ALGEBRAS, SECTIONS, DYNKIN'S PI-LAMBDA THEOREM, and FUBINI'S THEOREM for non-negative functions. Agent FUBINI-1, 2026-10-10 (scratchpad/triage/FUBINI-1.md).  Library: measure.}\par
\vnbprop{power-complement-in-member}{\vnbstmt{\forall c, a.; \forall b \in \operatorname{power}(c).\; (\operatorname{complement\text{-}in}(b, a) \in \operatorname{power}(c))}}
\vnbprop{dynkin-generated-unfold}{\vnbstmt{\begin{array}{@{}l@{}} \forall a, e. \\ \quad (\operatorname{dynkin\text{-}generated}(a, e) \simeq \{\mathit{dga} \in \operatorname{power}(a) : \forall b.\; ((\operatorname{is\text{-}dynkin\text{-}system}(a, b) \wedge (e \subseteq b)) \Rightarrow (\mathit{dga} \in b))\}) \end{array}}}
\vnbprop{dynkin-generated-membership}{\vnbstmt{\begin{array}{@{}l@{}} \forall a, e, \alpha. \\ \quad \alpha \in \operatorname{dynkin\text{-}generated}(a, e) \\ \quad \Leftrightarrow\; ((\alpha \in \operatorname{power}(a)) \wedge \forall b.\; ((\operatorname{is\text{-}dynkin\text{-}system}(a, b) \wedge (e \subseteq b)) \Rightarrow (\alpha \in b))) \end{array}}}
\vnbprop{dynkin-generated-is-dynkin}{\vnbstmt{\begin{array}{@{}l@{}} \forall a, e. \\ \quad ((a \in \mathbf{Set}) \wedge (e \subseteq \operatorname{power}(a))) \Rightarrow  \\ \quad \quad \operatorname{is\text{-}dynkin\text{-}system}(a, \operatorname{dynkin\text{-}generated}(a, e)) \end{array}}}
\vnbprop{dynkin-inter-section}{\vnbstmt{\begin{array}{@{}l@{}} \forall a, d, s. \\ \quad (\operatorname{is\text{-}dynkin\text{-}system}(a, d) \wedge (s \in d)) \Rightarrow  \\ \quad \quad \operatorname{is\text{-}dynkin\text{-}system}(a, \{e \in \operatorname{power}(a) : ((s \cap e) \in d)\}) \end{array}}}
\vnbprop{dynkin-pi-is-sigma-algebra}{\vnbstmt{\begin{array}{@{}l@{}} \forall a, d. \\ \quad (\operatorname{is\text{-}dynkin\text{-}system}(a, d) \wedge \operatorname{is\text{-}pi\text{-}system}(d)) \Rightarrow \operatorname{is\text{-}sigma\text{-}algebra}(a, d) \end{array}}}
\vnbprop{dynkin-pi-lambda}{\vnbstmt{\begin{array}{@{}l@{}} \forall a, d, p. \\ \quad (\operatorname{is\text{-}dynkin\text{-}system}(a, d) \wedge (\operatorname{is\text{-}pi\text{-}system}(p) \wedge (p \subseteq d))) \Rightarrow  \\ \quad \quad \operatorname{sigma\text{-}generated}(a, p) \subseteq d \end{array}}}
\vnbprop{cartesian-membership-nth}{\vnbstmt{\begin{array}{@{}l@{}} \forall a, b, w. \\ \quad (w \in (a \times b)) \Leftrightarrow ((w = \langle \operatorname{nth}(1, w), \operatorname{nth}(2, w) \rangle) \wedge ((\operatorname{nth}(1, w) \in a) \wedge (\operatorname{nth}(2, w) \in b))) \end{array}}}
\vnbprop{pair-cartesian-iff}{\vnbstmt{\forall a, b, u, v.\; ((\langle u, v \rangle \in (a \times b)) \Leftrightarrow ((u \in a) \wedge (v \in b)))}}
\vnbprop{x-section-unfold}{\vnbstmt{\forall s, a, x.\; (\operatorname{x\text{-}section}(s, a, x) \simeq \{y \in s : (\langle x, y \rangle \in a)\})}}
\vnbprop{y-section-unfold}{\vnbstmt{\forall s, a, y.\; (\operatorname{y\text{-}section}(s, a, y) \simeq \{x \in s : (\langle x, y \rangle \in a)\})}}
\vnbprop{x-section-membership}{\vnbstmt{\forall s, a, x, z.\; ((z \in \operatorname{x\text{-}section}(s, a, x)) \Leftrightarrow ((z \in s) \wedge (\langle x, z \rangle \in a)))}}
\vnbprop{y-section-membership}{\vnbstmt{\forall s, a, y, z.\; ((z \in \operatorname{y\text{-}section}(s, a, y)) \Leftrightarrow ((z \in s) \wedge (\langle z, y \rangle \in a)))}}
\vnbprop{product-sigma-unfold}{\vnbstmt{\begin{array}{@{}l@{}} \forall s, a, b, c. \\ \quad (\operatorname{product\text{-}sigma}(s, a, b, c) \simeq \operatorname{sigma\text{-}generated}((s \times b), \operatorname{rectangles}(s, a, b, c))) \end{array}}}
\vnbprop{rectangles-unfold}{\vnbstmt{\begin{array}{@{}l@{}} \forall s, a, b, c. \\ \quad (\operatorname{rectangles}(s, a, b, c) \simeq \{r \in \operatorname{power}((s \times b)) : \exists \mathit{rca} \in a, \mathit{rcb} \in c.\; (r = (\mathit{rca} \times \mathit{rcb}))\}) \end{array}}}
\vnbprop{rectangles-membership}{\vnbstmt{\begin{array}{@{}l@{}} \forall s, a, b, c, r. \\ \quad r \in \operatorname{rectangles}(s, a, b, c) \\ \quad \Leftrightarrow\; ((r \in \operatorname{power}((s \times b))) \wedge \exists \mathit{rca} \in a, \mathit{rcb} \in c.\; (r = (\mathit{rca} \times \mathit{rcb}))) \end{array}}}
\vnbprop{cartesian-in-power-product}{\vnbstmt{\begin{array}{@{}l@{}} \forall s, a, b, c. \\ \quad ((b \in \operatorname{power}(s)) \wedge (c \in \operatorname{power}(a))) \Rightarrow  \\ \quad \quad (b \times c) \in \operatorname{power}((s \times a)) \end{array}}}
\vnbprop{cartesian-intersection}{\vnbstmt{\begin{array}{@{}l@{}} \forall a, b, c, d. \\ \quad ((a \times b) \cap (c \times d)) = ((a \cap c) \times (b \cap d)) \end{array}}}
\vnbprop{rectangle-in-rectangles}{\vnbstmt{\begin{array}{@{}l@{}} \forall s, a, b, c, d, f. \\ \quad (\operatorname{is\text{-}sigma\text{-}algebra}(s, a) \wedge (\operatorname{is\text{-}sigma\text{-}algebra}(b, c) \wedge ((d \in a) \wedge (f \in c)))) \Rightarrow  \\ \quad \quad (d \times f) \in \operatorname{rectangles}(s, a, b, c) \end{array}}}
\vnbprop{rectangles-pi-system}{\vnbstmt{\begin{array}{@{}l@{}} \forall s, a, b, c. \\ \quad (\operatorname{is\text{-}sigma\text{-}algebra}(s, a) \wedge \operatorname{is\text{-}sigma\text{-}algebra}(b, c)) \Rightarrow  \\ \quad \quad \operatorname{is\text{-}pi\text{-}system}(\operatorname{rectangles}(s, a, b, c)) \end{array}}}
\vnbprop{product-sigma-is-sigma-algebra}{\vnbstmt{\begin{array}{@{}l@{}} \forall s, a, b, c. \\ \quad (\operatorname{is\text{-}sigma\text{-}algebra}(s, a) \wedge \operatorname{is\text{-}sigma\text{-}algebra}(b, c)) \Rightarrow  \\ \quad \quad \operatorname{is\text{-}sigma\text{-}algebra}((s \times b), \operatorname{product\text{-}sigma}(s, a, b, c)) \end{array}}}
\vnbprop{rectangle-in-product-sigma}{\vnbstmt{\begin{array}{@{}l@{}} \forall s, a, b, c, d, f. \\ \quad (\operatorname{is\text{-}sigma\text{-}algebra}(s, a) \wedge (\operatorname{is\text{-}sigma\text{-}algebra}(b, c) \wedge ((d \in a) \wedge (f \in c)))) \Rightarrow  \\ \quad \quad (d \times f) \in \operatorname{product\text{-}sigma}(s, a, b, c) \end{array}}}
\vnbprop{x-section-measurable}{\vnbstmt{\begin{array}{@{}l@{}} \forall s, a, b, c, d, x. \\ \quad (\operatorname{is\text{-}sigma\text{-}algebra}(s, a) \wedge (\operatorname{is\text{-}sigma\text{-}algebra}(b, c) \wedge ((d \in \operatorname{product\text{-}sigma}(s, a, b, c)) \wedge (x \in s)))) \Rightarrow  \\ \quad \quad \operatorname{x\text{-}section}(b, d, x) \in c \end{array}}}
\vnbprop{y-section-measurable}{\vnbstmt{\begin{array}{@{}l@{}} \forall s, a, b, c, d, y. \\ \quad (\operatorname{is\text{-}sigma\text{-}algebra}(s, a) \wedge (\operatorname{is\text{-}sigma\text{-}algebra}(b, c) \wedge ((d \in \operatorname{product\text{-}sigma}(s, a, b, c)) \wedge (y \in b)))) \Rightarrow  \\ \quad \quad \operatorname{y\text{-}section}(s, d, y) \in a \end{array}}}
\vnbprop{fn-x-section-apply}{\vnbstmt{\forall s, f, x.; \forall z \in s.\; (\operatorname{fn\text{-}x\text{-}section}(s, f, x)(z) \simeq f(\langle x, z \rangle))}}
\vnbprop{fn-y-section-apply}{\vnbstmt{\forall s, f, y.; \forall z \in s.\; (\operatorname{fn\text{-}y\text{-}section}(s, f, y)(z) \simeq f(\langle z, y \rangle))}}
\vnbprop{fn-x-section-measurable}{\vnbstmt{\begin{array}{@{}l@{}} \forall s, a, b, c, f, x. \\ \quad \operatorname{is\text{-}sigma\text{-}algebra}(s, a) \wedge \\ \quad \operatorname{is\text{-}sigma\text{-}algebra}(b, c) \wedge \\ \quad \operatorname{is\text{-}measurable\text{-}fn}((s \times b), \operatorname{product\text{-}sigma}(s, a, b, c), f) \wedge \\ \quad x \in s \Rightarrow \\ \quad \quad \operatorname{is\text{-}measurable\text{-}fn}(b, c, \operatorname{fn\text{-}x\text{-}section}(b, f, x)) \end{array}}}
\vnbprop{fn-y-section-measurable}{\vnbstmt{\begin{array}{@{}l@{}} \forall s, a, b, c, f, y. \\ \quad \operatorname{is\text{-}sigma\text{-}algebra}(s, a) \wedge \\ \quad \operatorname{is\text{-}sigma\text{-}algebra}(b, c) \wedge \\ \quad \operatorname{is\text{-}measurable\text{-}fn}((s \times b), \operatorname{product\text{-}sigma}(s, a, b, c), f) \wedge \\ \quad y \in b \Rightarrow \\ \quad \quad \operatorname{is\text{-}measurable\text{-}fn}(s, a, \operatorname{fn\text{-}y\text{-}section}(s, f, y)) \end{array}}}
\vnbprop{x-section-rectangle-in}{\vnbstmt{\begin{array}{@{}l@{}} \forall s, a, b, x. \\ \quad ((b \subseteq s) \wedge (x \in a)) \Rightarrow (\operatorname{x\text{-}section}(s, (a \times b), x) = b) \end{array}}}
\vnbprop{x-section-rectangle-out}{\vnbstmt{\begin{array}{@{}l@{}} \forall s, a, b, x. \\ \quad \neg (x \in a) \Rightarrow (\operatorname{x\text{-}section}(s, (a \times b), x) = \emptyset) \end{array}}}
\vnbprop{y-section-rectangle-in}{\vnbstmt{\begin{array}{@{}l@{}} \forall s, a, b, y. \\ \quad ((a \subseteq s) \wedge (y \in b)) \Rightarrow (\operatorname{y\text{-}section}(s, (a \times b), y) = a) \end{array}}}
\vnbprop{y-section-rectangle-out}{\vnbstmt{\begin{array}{@{}l@{}} \forall s, a, b, y. \\ \quad \neg (y \in b) \Rightarrow (\operatorname{y\text{-}section}(s, (a \times b), y) = \emptyset) \end{array}}}
\vnbprop{integral-congruence-ptwise}{\vnbstmt{\begin{array}{@{}l@{}} \forall a, b, u, f, g. \\ \quad ((f \in (a \to \operatorname{rr\text{-}pos\text{-}star})) \wedge ((g \in (a \to \operatorname{rr\text{-}pos\text{-}star})) \wedge \forall x \in a.\; (f(x) = g(x)))) \Rightarrow  \\ \quad \quad (\operatorname{integral}(a, b, u, f) \simeq \operatorname{integral}(a, b, u, g)) \end{array}}}
\vnbprop{measurable-fn-congruence-ptwise}{\vnbstmt{\begin{array}{@{}l@{}} \forall a, b, f, g. \\ \quad (\operatorname{is\text{-}measurable\text{-}fn}(a, b, f) \wedge ((g \in (a \to \operatorname{rr\text{-}pos\text{-}star})) \wedge \forall x \in a.\; (f(x) = g(x)))) \Rightarrow  \\ \quad \quad \operatorname{is\text{-}measurable\text{-}fn}(a, b, g) \end{array}}}
\vnbprop{x-section-measure-apply}{\vnbstmt{\begin{array}{@{}l@{}} \forall s, a, u, b.; \forall x \in s. \\ \quad (\operatorname{x\text{-}section\text{-}measure}(s, a, u, b)(x) \simeq u(\operatorname{x\text{-}section}(a, b, x))) \end{array}}}
\vnbprop{y-section-measure-apply}{\vnbstmt{\begin{array}{@{}l@{}} \forall s, a, u, b.; \forall y \in a. \\ \quad (\operatorname{y\text{-}section\text{-}measure}(s, a, u, b)(y) \simeq u(\operatorname{y\text{-}section}(s, b, y))) \end{array}}}
\vnbprop{x-section-measure-type}{\vnbstmt{\begin{array}{@{}l@{}} \forall s, a, u, b, c, d, f. \\ \quad (\operatorname{is\text{-}measure}(s, a, u) \wedge (\operatorname{is\text{-}measure}(b, c, d) \wedge (f \in \operatorname{product\text{-}sigma}(s, a, b, c)))) \Rightarrow  \\ \quad \quad \operatorname{x\text{-}section\text{-}measure}(s, b, d, f) \in (s \to \operatorname{rr\text{-}pos\text{-}star}) \end{array}}}
\vnbprop{y-section-measure-type}{\vnbstmt{\begin{array}{@{}l@{}} \forall s, a, u, b, c, d, f. \\ \quad (\operatorname{is\text{-}measure}(s, a, u) \wedge (\operatorname{is\text{-}measure}(b, c, d) \wedge (f \in \operatorname{product\text{-}sigma}(s, a, b, c)))) \Rightarrow  \\ \quad \quad \operatorname{y\text{-}section\text{-}measure}(s, b, u, f) \in (b \to \operatorname{rr\text{-}pos\text{-}star}) \end{array}}}
\vnbprop{x-section-measure-rectangle}{\vnbstmt{\begin{array}{@{}l@{}} \forall s, a, u, b, c, d, f, g. \\ \quad (\operatorname{is\text{-}measure}(s, a, u) \wedge (\operatorname{is\text{-}measure}(b, c, d) \wedge ((f \in a) \wedge ((g \in c) \wedge (d(g) \in \mathbb{R}))))) \Rightarrow  \\ \quad \quad \operatorname{is\text{-}measurable\text{-}fn}(s, a, \operatorname{x\text{-}section\text{-}measure}(s, b, d, (f \times g))) \wedge \\ \quad \quad \operatorname{integral}(s, a, u, \operatorname{x\text{-}section\text{-}measure}(s, b, d, (f \times g))) = \operatorname{etimes}(d(g), u(f)) \end{array}}}
\vnbprop{y-section-measure-rectangle}{\vnbstmt{\begin{array}{@{}l@{}} \forall s, a, u, b, c, d, f, g. \\ \quad (\operatorname{is\text{-}measure}(s, a, u) \wedge (\operatorname{is\text{-}measure}(b, c, d) \wedge ((f \in a) \wedge ((g \in c) \wedge (u(f) \in \mathbb{R}))))) \Rightarrow  \\ \quad \quad \operatorname{is\text{-}measurable\text{-}fn}(b, c, \operatorname{y\text{-}section\text{-}measure}(s, b, u, (f \times g))) \wedge \\ \quad \quad \operatorname{integral}(b, c, d, \operatorname{y\text{-}section\text{-}measure}(s, b, u, (f \times g))) = \operatorname{etimes}(u(f), d(g)) \end{array}}}
\vnbprop{fubini-set-rectangle}{\vnbstmt{\begin{array}{@{}l@{}} \forall s, a, u, b, c, d, f, g. \\ \quad (\operatorname{is\text{-}measure}(s, a, u) \wedge (\operatorname{is\text{-}measure}(b, c, d) \wedge ((f \in a) \wedge ((g \in c) \wedge ((u(f) \in \mathbb{R}) \wedge (d(g) \in \mathbb{R})))))) \Rightarrow  \\ \quad \quad \operatorname{is\text{-}fubini\text{-}set}(s, a, u, b, c, d, (f \times g)) \end{array}}}
\vnbprop{x-section-mono}{\vnbstmt{\begin{array}{@{}l@{}} \forall s, a, b, x. \\ \quad (a \subseteq b) \Rightarrow (\operatorname{x\text{-}section}(s, a, x) \subseteq \operatorname{x\text{-}section}(s, b, x)) \end{array}}}
\vnbprop{y-section-mono}{\vnbstmt{\begin{array}{@{}l@{}} \forall s, a, b, y. \\ \quad (a \subseteq b) \Rightarrow (\operatorname{y\text{-}section}(s, a, y) \subseteq \operatorname{y\text{-}section}(s, b, y)) \end{array}}}
\vnbprop{x-section-difference}{\vnbstmt{\begin{array}{@{}l@{}} \forall s, a, b, x. \\ \quad \operatorname{x\text{-}section}(s, \operatorname{complement\text{-}in}(b, a), x) = \operatorname{complement\text{-}in}(\operatorname{x\text{-}section}(s, b, x), \operatorname{x\text{-}section}(s, a, x)) \end{array}}}
\vnbprop{y-section-difference}{\vnbstmt{\begin{array}{@{}l@{}} \forall s, a, b, y. \\ \quad \operatorname{y\text{-}section}(s, \operatorname{complement\text{-}in}(b, a), y) = \operatorname{complement\text{-}in}(\operatorname{y\text{-}section}(s, b, y), \operatorname{y\text{-}section}(s, a, y)) \end{array}}}
\vnbprop{x-section-measure-difference}{\vnbstmt{\begin{array}{@{}l@{}} \forall s, a, u, b, c, d, f, g. \\ \quad \operatorname{is\text{-}measure}(s, a, u) \wedge \\ \quad \operatorname{is\text{-}measure}(b, c, d) \wedge \\ \quad f \in \operatorname{product\text{-}sigma}(s, a, b, c) \wedge \\ \quad g \in \operatorname{product\text{-}sigma}(s, a, b, c) \wedge \\ \quad f \subseteq g \wedge \\ \quad \operatorname{is\text{-}measurable\text{-}fn}(s, a, \operatorname{x\text{-}section\text{-}measure}(s, b, d, f)) \wedge \\ \quad \operatorname{is\text{-}measurable\text{-}fn}(s, a, \operatorname{x\text{-}section\text{-}measure}(s, b, d, g)) \wedge \\ \quad \forall z \in s.\; (d(\operatorname{x\text{-}section}(b, g, z)) \in \mathbb{R}) \wedge \\ \quad \operatorname{integral}(s, a, u, \operatorname{x\text{-}section\text{-}measure}(s, b, d, g)) \in \mathbb{R} \Rightarrow \\ \quad \quad \operatorname{is\text{-}measurable\text{-}fn}(s, a, \operatorname{x\text{-}section\text{-}measure}(s, b, d, \operatorname{complement\text{-}in}(g, f))) \wedge \\ \quad \quad \operatorname{integral}(s, a, u, \operatorname{x\text{-}section\text{-}measure}(s, b, d, f)) \in \mathbb{R} \wedge \\ \quad \quad \operatorname{integral}(s, a, u, \operatorname{x\text{-}section\text{-}measure}(s, b, d, \operatorname{complement\text{-}in}(g, f))) \\ \quad \quad =\; (\operatorname{integral}(s, a, u, \operatorname{x\text{-}section\text{-}measure}(s, b, d, g)) - \operatorname{integral}(s, a, u, \operatorname{x\text{-}section\text{-}measure}(s, b, d, f))) \end{array}}}
\vnbprop{y-section-measure-difference}{\vnbstmt{\begin{array}{@{}l@{}} \forall s, a, u, b, c, d, f, g. \\ \quad \operatorname{is\text{-}measure}(s, a, u) \wedge \\ \quad \operatorname{is\text{-}measure}(b, c, d) \wedge \\ \quad f \in \operatorname{product\text{-}sigma}(s, a, b, c) \wedge \\ \quad g \in \operatorname{product\text{-}sigma}(s, a, b, c) \wedge \\ \quad f \subseteq g \wedge \\ \quad \operatorname{is\text{-}measurable\text{-}fn}(b, c, \operatorname{y\text{-}section\text{-}measure}(s, b, u, f)) \wedge \\ \quad \operatorname{is\text{-}measurable\text{-}fn}(b, c, \operatorname{y\text{-}section\text{-}measure}(s, b, u, g)) \wedge \\ \quad \forall z \in b.\; (u(\operatorname{y\text{-}section}(s, g, z)) \in \mathbb{R}) \wedge \\ \quad \operatorname{integral}(b, c, d, \operatorname{y\text{-}section\text{-}measure}(s, b, u, g)) \in \mathbb{R} \Rightarrow \\ \quad \quad \operatorname{is\text{-}measurable\text{-}fn}(b, c, \operatorname{y\text{-}section\text{-}measure}(s, b, u, \operatorname{complement\text{-}in}(g, f))) \wedge \\ \quad \quad \operatorname{integral}(b, c, d, \operatorname{y\text{-}section\text{-}measure}(s, b, u, f)) \in \mathbb{R} \wedge \\ \quad \quad \operatorname{integral}(b, c, d, \operatorname{y\text{-}section\text{-}measure}(s, b, u, \operatorname{complement\text{-}in}(g, f))) \\ \quad \quad =\; (\operatorname{integral}(b, c, d, \operatorname{y\text{-}section\text{-}measure}(s, b, u, g)) - \operatorname{integral}(b, c, d, \operatorname{y\text{-}section\text{-}measure}(s, b, u, f))) \end{array}}}
\vnbprop{fubini-set-difference}{\vnbstmt{\begin{array}{@{}l@{}} \forall s, a, u, b, c, d, f, g. \\ \quad \operatorname{is\text{-}measure}(s, a, u) \wedge \\ \quad \operatorname{is\text{-}measure}(b, c, d) \wedge \\ \quad f \in \operatorname{product\text{-}sigma}(s, a, b, c) \wedge \\ \quad g \in \operatorname{product\text{-}sigma}(s, a, b, c) \wedge \\ \quad f \subseteq g \wedge \\ \quad \operatorname{is\text{-}fubini\text{-}set}(s, a, u, b, c, d, f) \wedge \\ \quad \operatorname{is\text{-}fubini\text{-}set}(s, a, u, b, c, d, g) \wedge \\ \quad \forall z \in s.\; (d(\operatorname{x\text{-}section}(b, g, z)) \in \mathbb{R}) \wedge \\ \quad \forall z \in b.\; (u(\operatorname{y\text{-}section}(s, g, z)) \in \mathbb{R}) \wedge \\ \quad \operatorname{integral}(s, a, u, \operatorname{x\text{-}section\text{-}measure}(s, b, d, g)) \in \mathbb{R} \Rightarrow \\ \quad \quad \operatorname{is\text{-}fubini\text{-}set}(s, a, u, b, c, d, \operatorname{complement\text{-}in}(g, f)) \end{array}}}
\vnbprop{x-section-measure-union}{\vnbstmt{\begin{array}{@{}l@{}} \forall s, a, u, b, c, d, f. \\ \quad \operatorname{is\text{-}measure}(s, a, u) \wedge \\ \quad \operatorname{is\text{-}measure}(b, c, d) \wedge \\ \quad f \in (\mathbb{N} \to \operatorname{product\text{-}sigma}(s, a, b, c)) \wedge \\ \quad \forall n \in \mathbb{N}.\; (f(n) \subseteq f((n + 1))) \wedge \\ \quad \forall n \in \mathbb{N}.\; \operatorname{is\text{-}measurable\text{-}fn}(s, a, \operatorname{x\text{-}section\text{-}measure}(s, b, d, f(n))) \Rightarrow \\ \quad \quad \operatorname{is\text{-}measurable\text{-}fn}(s, a, \operatorname{x\text{-}section\text{-}measure}(s, b, d, \operatorname{big\text{-}union}(g, \mathbb{N}, f(g)))) \wedge \\ \quad \quad \operatorname{integral}(s, a, u, \operatorname{x\text{-}section\text{-}measure}(s, b, d, \operatorname{big\text{-}union}(g, \mathbb{N}, f(g)))) \\ \quad \quad =\; \operatorname{esup}(\operatorname{image}((\lambda n \in \mathbb{N}.\; \operatorname{integral}(s, a, u, \operatorname{x\text{-}section\text{-}measure}(s, b, d, f(n)))), \mathbb{N})) \end{array}}}
\vnbprop{y-section-measure-union}{\vnbstmt{\begin{array}{@{}l@{}} \forall s, a, u, b, c, d, f. \\ \quad \operatorname{is\text{-}measure}(s, a, u) \wedge \\ \quad \operatorname{is\text{-}measure}(b, c, d) \wedge \\ \quad f \in (\mathbb{N} \to \operatorname{product\text{-}sigma}(s, a, b, c)) \wedge \\ \quad \forall n \in \mathbb{N}.\; (f(n) \subseteq f((n + 1))) \wedge \\ \quad \forall n \in \mathbb{N}.\; \operatorname{is\text{-}measurable\text{-}fn}(b, c, \operatorname{y\text{-}section\text{-}measure}(s, b, u, f(n))) \Rightarrow \\ \quad \quad \operatorname{is\text{-}measurable\text{-}fn}(b, c, \operatorname{y\text{-}section\text{-}measure}(s, b, u, \operatorname{big\text{-}union}(g, \mathbb{N}, f(g)))) \wedge \\ \quad \quad \operatorname{integral}(b, c, d, \operatorname{y\text{-}section\text{-}measure}(s, b, u, \operatorname{big\text{-}union}(g, \mathbb{N}, f(g)))) \\ \quad \quad =\; \operatorname{esup}(\operatorname{image}((\lambda n \in \mathbb{N}.\; \operatorname{integral}(b, c, d, \operatorname{y\text{-}section\text{-}measure}(s, b, u, f(n)))), \mathbb{N})) \end{array}}}
\vnbprop{fubini-set-union}{\vnbstmt{\begin{array}{@{}l@{}} \forall s, a, u, b, c, d, f. \\ \quad \operatorname{is\text{-}measure}(s, a, u) \wedge \\ \quad \operatorname{is\text{-}measure}(b, c, d) \wedge \\ \quad f \in (\mathbb{N} \to \operatorname{product\text{-}sigma}(s, a, b, c)) \wedge \\ \quad \forall n \in \mathbb{N}.\; (f(n) \subseteq f((n + 1))) \wedge \\ \quad \forall n \in \mathbb{N}.\; \operatorname{is\text{-}fubini\text{-}set}(s, a, u, b, c, d, f(n)) \Rightarrow \\ \quad \quad \operatorname{is\text{-}fubini\text{-}set}(s, a, u, b, c, d, \operatorname{big\text{-}union}(g, \mathbb{N}, f(g))) \end{array}}}
\vnbprop{x-section-in-cut}{\vnbstmt{\begin{array}{@{}l@{}} \forall s, a, b, c, x. \\ \quad \operatorname{x\text{-}section}(s, (a \cap (b \times c)), x) \subseteq c \end{array}}}
\vnbprop{y-section-in-cut}{\vnbstmt{\begin{array}{@{}l@{}} \forall s, a, b, c, y. \\ \quad \operatorname{y\text{-}section}(s, (a \cap (b \times c)), y) \subseteq b \end{array}}}
\vnbprop{fubini-set-rectangle-cut}{\vnbstmt{\begin{array}{@{}l@{}} \forall s, a, u, b, c, d, f, g, h, k. \\ \quad \operatorname{is\text{-}measure}(s, a, u) \wedge \\ \quad \operatorname{is\text{-}measure}(b, c, d) \wedge \\ \quad f \in a \wedge \\ \quad g \in c \wedge \\ \quad u(f) \in \mathbb{R} \wedge \\ \quad d(g) \in \mathbb{R} \wedge \\ \quad h \in a \wedge \\ \quad k \in c \Rightarrow \\ \quad \quad \operatorname{is\text{-}fubini\text{-}set}(s, a, u, b, c, d, ((h \times k) \cap (f \times g))) \end{array}}}
\vnbprop{fubini-set-cut-bounds}{\vnbstmt{\begin{array}{@{}l@{}} \forall s, a, u, b, c, d, f, g, h. \\ \quad \operatorname{is\text{-}measure}(s, a, u) \wedge \\ \quad \operatorname{is\text{-}measure}(b, c, d) \wedge \\ \quad f \in a \wedge \\ \quad g \in c \wedge \\ \quad u(f) \in \mathbb{R} \wedge \\ \quad d(g) \in \mathbb{R} \wedge \\ \quad h \in \operatorname{product\text{-}sigma}(s, a, b, c) \wedge \\ \quad \operatorname{is\text{-}fubini\text{-}set}(s, a, u, b, c, d, (h \cap (f \times g))) \Rightarrow \\ \quad \quad \forall z \in s.\; (d(\operatorname{x\text{-}section}(b, (h \cap (f \times g)), z)) \in \mathbb{R}) \wedge \\ \quad \quad \forall z \in b.\; (u(\operatorname{y\text{-}section}(s, (h \cap (f \times g)), z)) \in \mathbb{R}) \wedge \\ \quad \quad \operatorname{integral}(s, a, u, \operatorname{x\text{-}section\text{-}measure}(s, b, d, (h \cap (f \times g)))) \in \mathbb{R} \end{array}}}
\vnbprop{fubini-set-in-finite-rectangle}{\vnbstmt{\begin{array}{@{}l@{}} \forall s, a, u, b, c, d, f, g, h. \\ \quad \operatorname{is\text{-}measure}(s, a, u) \wedge \\ \quad \operatorname{is\text{-}measure}(b, c, d) \wedge \\ \quad f \in a \wedge \\ \quad g \in c \wedge \\ \quad u(f) \in \mathbb{R} \wedge \\ \quad d(g) \in \mathbb{R} \wedge \\ \quad h \in \operatorname{product\text{-}sigma}(s, a, b, c) \Rightarrow \\ \quad \quad \operatorname{is\text{-}fubini\text{-}set}(s, a, u, b, c, d, (h \cap (f \times g))) \end{array}}}
\vnbprop{measure-union-finite}{\vnbstmt{\begin{array}{@{}l@{}} \forall c, d, u, a, b. \\ \quad (\operatorname{is\text{-}measure}(c, d, u) \wedge ((a \in d) \wedge ((b \in d) \wedge ((u(a) \in \mathbb{R}) \wedge (u(b) \in \mathbb{R}))))) \Rightarrow  \\ \quad \quad u((a \cup b)) \in \mathbb{R} \end{array}}}
\vnbprop{measure-seg-union-finite}{\vnbstmt{\begin{array}{@{}l@{}} \forall k \in \mathbb{N}, a, b, u, f. \\ \quad (\operatorname{is\text{-}measure}(a, b, u) \wedge ((f \in (\mathbb{N} \to b)) \wedge \forall n \in \mathbb{N}.\; (u(f(n)) \in \mathbb{R}))) \Rightarrow  \\ \quad \quad u(\operatorname{big\text{-}union}(j, \mathrm{OS}(k), f(j))) \in \mathbb{R} \end{array}}}
\vnbprop{sigma-finite-increasing}{\vnbstmt{\begin{array}{@{}l@{}} \forall a, b, u. \\ \quad (\operatorname{is\text{-}measure}(a, b, u) \wedge \operatorname{is\text{-}sigma\text{-}finite}(a, b, u)) \Rightarrow  \\ \quad \quad \exists g \in (\mathbb{N} \to b). \\ \quad \quad \quad \forall n \in \mathbb{N}.\; (u(g(n)) \in \mathbb{R}) \wedge \\ \quad \quad \quad \forall n \in \mathbb{N}.\; (g(n) \subseteq g((n + 1))) \wedge \\ \quad \quad \quad \operatorname{big\text{-}union}(k, \mathbb{N}, g(k)) = a \end{array}}}
\vnbprop{cartesian-mono}{\vnbstmt{\begin{array}{@{}l@{}} \forall a, b, c, d. \\ \quad ((a \subseteq c) \wedge (b \subseteq d)) \Rightarrow ((a \times b) \subseteq (c \times d)) \end{array}}}
\vnbprop{fubini-sets}{\vnbstmt{\begin{array}{@{}l@{}} \forall s, a, u, b, c, d, f. \\ \quad \operatorname{is\text{-}measure}(s, a, u) \wedge \\ \quad \operatorname{is\text{-}measure}(b, c, d) \wedge \\ \quad \operatorname{is\text{-}sigma\text{-}finite}(s, a, u) \wedge \\ \quad \operatorname{is\text{-}sigma\text{-}finite}(b, c, d) \wedge \\ \quad f \in \operatorname{product\text{-}sigma}(s, a, b, c) \Rightarrow \\ \quad \quad \operatorname{is\text{-}fubini\text{-}set}(s, a, u, b, c, d, f) \end{array}}}
\vnbprop{integral-finsum}{\vnbstmt{\begin{array}{@{}l@{}} \forall a, b, u, d, s, m. \\ \quad \operatorname{is\text{-}measure}(a, b, u) \wedge \\ \quad s \in \mathbf{Set} \wedge \\ \quad |s| \in \mathbb{N} \wedge \\ \quad s \subseteq d \wedge \\ \quad \forall v \in s.\; \operatorname{is\text{-}measurable\text{-}fn}(a, b, m(v)) \Rightarrow \\ \quad \quad \operatorname{is\text{-}measurable\text{-}fn}(a, b, \\ \quad \quad \quad (\lambda x \in a.\; \operatorname{finsum}(\operatorname{rr\text{-}pos\text{-}star\text{-}add\text{-}monoid}, (\lambda v \in d.\; m(v)(x)), s))) \wedge \\ \quad \quad \operatorname{integral}(a, b, u, (\lambda x \in a.\; \operatorname{finsum}(\operatorname{rr\text{-}pos\text{-}star\text{-}add\text{-}monoid}, (\lambda v \in d.\; m(v)(x)), s))) \\ \quad \quad =\; \operatorname{finsum}(\operatorname{rr\text{-}pos\text{-}star\text{-}add\text{-}monoid}, (\lambda v \in d.\; \operatorname{integral}(a, b, u, m(v))), s) \end{array}}}
\vnbprop{x-section-simple-integral}{\vnbstmt{\begin{array}{@{}l@{}} \forall a, b, c, d, u, s, x. \\ \quad \operatorname{is\text{-}sigma\text{-}algebra}(a, b) \wedge \\ \quad \operatorname{is\text{-}measure}(c, d, u) \wedge \\ \quad \operatorname{is\text{-}simple\text{-}fn}((a \times c), \operatorname{product\text{-}sigma}(a, b, c, d), s) \wedge \\ \quad x \in a \Rightarrow \\ \quad \quad \operatorname{integral}(c, d, u, \operatorname{fn\text{-}x\text{-}section}(c, s, x)) \\ \quad \quad =\; \operatorname{finsum}(\operatorname{rr\text{-}pos\text{-}star\text{-}add\text{-}monoid}, \\ \quad \quad \quad (\lambda v \in \operatorname{image}(s, (a \times c)).\; \\ \quad \quad \quad \quad \operatorname{etimes}(v, u(\operatorname{x\text{-}section}(c, \{\mathit{fx} \in (a \times c) : (s(\mathit{fx}) = v)\}, x)))), \\ \quad \quad \quad \operatorname{image}(s, (a \times c))) \end{array}}}
\vnbprop{y-section-simple-integral}{\vnbstmt{\begin{array}{@{}l@{}} \forall a, b, u, c, d, s, y. \\ \quad \operatorname{is\text{-}measure}(a, b, u) \wedge \\ \quad \operatorname{is\text{-}sigma\text{-}algebra}(c, d) \wedge \\ \quad \operatorname{is\text{-}simple\text{-}fn}((a \times c), \operatorname{product\text{-}sigma}(a, b, c, d), s) \wedge \\ \quad y \in c \Rightarrow \\ \quad \quad \operatorname{integral}(a, b, u, \operatorname{fn\text{-}y\text{-}section}(a, s, y)) \\ \quad \quad =\; \operatorname{finsum}(\operatorname{rr\text{-}pos\text{-}star\text{-}add\text{-}monoid}, \\ \quad \quad \quad (\lambda v \in \operatorname{image}(s, (a \times c)).\; \\ \quad \quad \quad \quad \operatorname{etimes}(v, u(\operatorname{y\text{-}section}(a, \{x \in (a \times c) : (s(x) = v)\}, y)))), \\ \quad \quad \quad \operatorname{image}(s, (a \times c))) \end{array}}}
\vnbprop{x-section-integral-apply}{\vnbstmt{\begin{array}{@{}l@{}} \forall s, a, b, u, f.; \forall x \in s. \\ \quad (\operatorname{x\text{-}section\text{-}integral}(s, a, b, u, f)(x) \simeq \operatorname{integral}(a, b, u, \operatorname{fn\text{-}x\text{-}section}(a, f, x))) \end{array}}}
\vnbprop{y-section-integral-apply}{\vnbstmt{\begin{array}{@{}l@{}} \forall s, a, u, b, f.; \forall y \in b. \\ \quad (\operatorname{y\text{-}section\text{-}integral}(s, a, u, b, f)(y) \simeq \operatorname{integral}(s, a, u, \operatorname{fn\text{-}y\text{-}section}(s, f, y))) \end{array}}}
\vnbprop{x-section-integral-type}{\vnbstmt{\begin{array}{@{}l@{}} \forall s, a, b, c, u, f. \\ \quad \operatorname{is\text{-}sigma\text{-}algebra}(s, a) \wedge \\ \quad \operatorname{is\text{-}sigma\text{-}algebra}(b, c) \wedge \\ \quad \operatorname{is\text{-}measure}(b, c, u) \wedge \\ \quad \operatorname{is\text{-}measurable\text{-}fn}((s \times b), \operatorname{product\text{-}sigma}(s, a, b, c), f) \Rightarrow \\ \quad \quad \operatorname{x\text{-}section\text{-}integral}(s, b, c, u, f) \in (s \to \operatorname{rr\text{-}pos\text{-}star}) \end{array}}}
\vnbprop{y-section-integral-type}{\vnbstmt{\begin{array}{@{}l@{}} \forall s, a, u, b, c, f. \\ \quad \operatorname{is\text{-}sigma\text{-}algebra}(s, a) \wedge \\ \quad \operatorname{is\text{-}sigma\text{-}algebra}(b, c) \wedge \\ \quad \operatorname{is\text{-}measure}(s, a, u) \wedge \\ \quad \operatorname{is\text{-}measurable\text{-}fn}((s \times b), \operatorname{product\text{-}sigma}(s, a, b, c), f) \Rightarrow \\ \quad \quad \operatorname{y\text{-}section\text{-}integral}(s, a, u, b, f) \in (b \to \operatorname{rr\text{-}pos\text{-}star}) \end{array}}}
\vnbprop{fubini-simple}{\vnbstmt{\begin{array}{@{}l@{}} \forall a, b, u, c, d, f, s. \\ \quad \operatorname{is\text{-}measure}(a, b, u) \wedge \\ \quad \operatorname{is\text{-}measure}(c, d, f) \wedge \\ \quad \operatorname{is\text{-}sigma\text{-}finite}(a, b, u) \wedge \\ \quad \operatorname{is\text{-}sigma\text{-}finite}(c, d, f) \wedge \\ \quad \operatorname{is\text{-}simple\text{-}fn}((a \times c), \operatorname{product\text{-}sigma}(a, b, c, d), s) \Rightarrow \\ \quad \quad \operatorname{is\text{-}measurable\text{-}fn}(a, b, \operatorname{x\text{-}section\text{-}integral}(a, c, d, f, s)) \wedge \\ \quad \quad \operatorname{is\text{-}measurable\text{-}fn}(c, d, \operatorname{y\text{-}section\text{-}integral}(a, b, u, c, s)) \wedge \\ \quad \quad \operatorname{integral}(a, b, u, \operatorname{x\text{-}section\text{-}integral}(a, c, d, f, s)) \\ \quad \quad =\; \operatorname{integral}(c, d, f, \operatorname{y\text{-}section\text{-}integral}(a, b, u, c, s)) \end{array}}}
\vnbprop{ptwise-sup-apply}{\vnbstmt{\begin{array}{@{}l@{}} \forall a, f.; \forall x \in a. \\ \quad (\operatorname{ptwise\text{-}sup}(a, f)(x) \simeq \operatorname{esup}(\operatorname{image}((\lambda n \in \mathbb{N}.\; f(n)(x)), \mathbb{N}))) \end{array}}}
\vnbprop{x-section-integral-limit}{\vnbstmt{\begin{array}{@{}l@{}} \forall s, a, u, b, c, d, f, g. \\ \quad \operatorname{is\text{-}measure}(s, a, u) \wedge \\ \quad \operatorname{is\text{-}measure}(b, c, d) \wedge \\ \quad \operatorname{is\text{-}measurable\text{-}fn}((s \times b), \operatorname{product\text{-}sigma}(s, a, b, c), f) \wedge \\ \quad g \in (\mathbb{N} \to ((s \times b) \to \operatorname{rr\text{-}pos\text{-}star})) \wedge \\ \quad \forall n \in \mathbb{N}. \\ \quad \quad \operatorname{is\text{-}measurable\text{-}fn}((s \times b), \operatorname{product\text{-}sigma}(s, a, b, c), g(n)) \wedge \\ \quad \forall n \in \mathbb{N}.\; \operatorname{ptwise\text{-}le}((s \times b), g(n), g((n + 1))) \wedge \\ \quad \forall w \in (s \times b).\; (\operatorname{ptwise\text{-}sup}((s \times b), g)(w) = f(w)) \wedge \\ \quad \forall n \in \mathbb{N}.\; \operatorname{is\text{-}measurable\text{-}fn}(s, a, \operatorname{x\text{-}section\text{-}integral}(s, b, c, d, g(n))) \Rightarrow \\ \quad \quad \operatorname{is\text{-}measurable\text{-}fn}(s, a, \operatorname{x\text{-}section\text{-}integral}(s, b, c, d, f)) \wedge \\ \quad \quad \operatorname{integral}(s, a, u, \operatorname{x\text{-}section\text{-}integral}(s, b, c, d, f)) \\ \quad \quad =\; \operatorname{esup}(\operatorname{image}((\lambda n \in \mathbb{N}.\; \operatorname{integral}(s, a, u, \operatorname{x\text{-}section\text{-}integral}(s, b, c, d, g(n)))), \mathbb{N})) \end{array}}}
\vnbprop{y-section-integral-limit}{\vnbstmt{\begin{array}{@{}l@{}} \forall s, a, u, b, c, d, f, g. \\ \quad \operatorname{is\text{-}measure}(s, a, u) \wedge \\ \quad \operatorname{is\text{-}measure}(b, c, d) \wedge \\ \quad \operatorname{is\text{-}measurable\text{-}fn}((s \times b), \operatorname{product\text{-}sigma}(s, a, b, c), f) \wedge \\ \quad g \in (\mathbb{N} \to ((s \times b) \to \operatorname{rr\text{-}pos\text{-}star})) \wedge \\ \quad \forall n \in \mathbb{N}. \\ \quad \quad \operatorname{is\text{-}measurable\text{-}fn}((s \times b), \operatorname{product\text{-}sigma}(s, a, b, c), g(n)) \wedge \\ \quad \forall n \in \mathbb{N}.\; \operatorname{ptwise\text{-}le}((s \times b), g(n), g((n + 1))) \wedge \\ \quad \forall w \in (s \times b).\; (\operatorname{ptwise\text{-}sup}((s \times b), g)(w) = f(w)) \wedge \\ \quad \forall n \in \mathbb{N}.\; \operatorname{is\text{-}measurable\text{-}fn}(b, c, \operatorname{y\text{-}section\text{-}integral}(s, a, u, b, g(n))) \Rightarrow \\ \quad \quad \operatorname{is\text{-}measurable\text{-}fn}(b, c, \operatorname{y\text{-}section\text{-}integral}(s, a, u, b, f)) \wedge \\ \quad \quad \operatorname{integral}(b, c, d, \operatorname{y\text{-}section\text{-}integral}(s, a, u, b, f)) \\ \quad \quad =\; \operatorname{esup}(\operatorname{image}((\lambda n \in \mathbb{N}.\; \operatorname{integral}(b, c, d, \operatorname{y\text{-}section\text{-}integral}(s, a, u, b, g(n)))), \mathbb{N})) \end{array}}}
\vnbprop{fubini}{\vnbstmt{\begin{array}{@{}l@{}} \forall s, a, u, b, c, d, f. \\ \quad \operatorname{is\text{-}measure}(s, a, u) \wedge \\ \quad \operatorname{is\text{-}measure}(b, c, d) \wedge \\ \quad \operatorname{is\text{-}sigma\text{-}finite}(s, a, u) \wedge \\ \quad \operatorname{is\text{-}sigma\text{-}finite}(b, c, d) \wedge \\ \quad \operatorname{is\text{-}measurable\text{-}fn}((s \times b), \operatorname{product\text{-}sigma}(s, a, b, c), f) \Rightarrow \\ \quad \quad \operatorname{is\text{-}measurable\text{-}fn}(s, a, \operatorname{x\text{-}section\text{-}integral}(s, b, c, d, f)) \wedge \\ \quad \quad \operatorname{is\text{-}measurable\text{-}fn}(b, c, \operatorname{y\text{-}section\text{-}integral}(s, a, u, b, f)) \wedge \\ \quad \quad \operatorname{integral}(s, a, u, \operatorname{x\text{-}section\text{-}integral}(s, b, c, d, f)) \\ \quad \quad =\; \operatorname{integral}(b, c, d, \operatorname{y\text{-}section\text{-}integral}(s, a, u, b, f)) \end{array}}}
\vnbprop{esum-as-esup-prefix}{\vnbstmt{\begin{array}{@{}l@{}} \forall f \in (\mathbb{N} \to \operatorname{rr\text{-}pos\text{-}star}). \\ \quad \operatorname{esum}(f) = \operatorname{esup}(\operatorname{image}((\lambda n \in \mathbb{N}.\; \operatorname{finsum}(\operatorname{rr\text{-}pos\text{-}star\text{-}add\text{-}monoid}, f, \mathrm{OS}(n))), \mathbb{N})) \end{array}}}
\vnbprop{integral-esum}{\vnbstmt{\begin{array}{@{}l@{}} \forall a, b, u, g, h. \\ \quad \operatorname{is\text{-}measure}(a, b, u) \wedge \\ \quad \forall v \in \mathbb{N}.\; \operatorname{is\text{-}measurable\text{-}fn}(a, b, g(v)) \wedge \\ \quad h \in (a \to \operatorname{rr\text{-}pos\text{-}star}) \wedge \\ \quad \forall x \in a.\; (h(x) = \operatorname{esum}((\lambda v \in \mathbb{N}.\; g(v)(x)))) \Rightarrow \\ \quad \quad \operatorname{is\text{-}measurable\text{-}fn}(a, b, h) \wedge \\ \quad \quad \operatorname{integral}(a, b, u, h) = \operatorname{esum}((\lambda v \in \mathbb{N}.\; \operatorname{integral}(a, b, u, g(v)))) \end{array}}}
\vnbprop{product-measure-apply}{\vnbstmt{\begin{array}{@{}l@{}} \forall s, a, u, b, c, d.; \forall f \in \operatorname{product\text{-}sigma}(s, a, b, c). \\ \quad (\operatorname{product\text{-}measure}(s, a, u, b, c, d)(f) \simeq \operatorname{integral}(s, a, u, \operatorname{x\text{-}section\text{-}measure}(s, b, d, f))) \end{array}}}
\vnbprop{x-section-empty}{\vnbstmt{\forall s, x.\; (\operatorname{x\text{-}section}(s, \emptyset, x) = \emptyset)}}
\vnbprop{product-measure-is-measure}{\vnbstmt{\begin{array}{@{}l@{}} \forall s, a, u, b, c, d. \\ \quad \operatorname{is\text{-}measure}(s, a, u) \wedge \\ \quad \operatorname{is\text{-}measure}(b, c, d) \wedge \\ \quad \operatorname{is\text{-}sigma\text{-}finite}(s, a, u) \wedge \\ \quad \operatorname{is\text{-}sigma\text{-}finite}(b, c, d) \Rightarrow \\ \quad \quad \operatorname{is\text{-}measure}((s \times b), \operatorname{product\text{-}sigma}(s, a, b, c), \operatorname{product\text{-}measure}(s, a, u, b, c, d)) \end{array}}}
\vnbprop{product-measure-symmetric}{\vnbstmt{\begin{array}{@{}l@{}} \forall s, a, u, b, c, d, f. \\ \quad \operatorname{is\text{-}measure}(s, a, u) \wedge \\ \quad \operatorname{is\text{-}measure}(b, c, d) \wedge \\ \quad \operatorname{is\text{-}sigma\text{-}finite}(s, a, u) \wedge \\ \quad \operatorname{is\text{-}sigma\text{-}finite}(b, c, d) \wedge \\ \quad f \in \operatorname{product\text{-}sigma}(s, a, b, c) \Rightarrow \\ \quad \quad \operatorname{product\text{-}measure}(s, a, u, b, c, d)(f) = \operatorname{integral}(b, c, d, \operatorname{y\text{-}section\text{-}measure}(s, b, u, f)) \end{array}}}
\vnbprop{product-measure-rectangle}{\vnbstmt{\begin{array}{@{}l@{}} \forall s, a, u, b, c, d, f, g. \\ \quad (\operatorname{is\text{-}measure}(s, a, u) \wedge (\operatorname{is\text{-}measure}(b, c, d) \wedge ((f \in a) \wedge ((g \in c) \wedge ((u(f) \in \mathbb{R}) \wedge (d(g) \in \mathbb{R})))))) \Rightarrow  \\ \quad \quad \operatorname{product\text{-}measure}(s, a, u, b, c, d)((f \times g)) = \operatorname{etimes}(u(f), d(g)) \end{array}}}

\vnbbookfile{structure-library/push-forward}{structure-library/push-forward}
\noindent{\small structure-library/push-forward.scm -- the push-forward of a measure along a map (RM-3B, 2026-10-10).  DEFINITIONS ONLY; every law is PROVEN in theorem-library/push-forward.scm (prefix pfw-).}\par
\par\medskip\noindent\textbf{Definition} (\texttt{push-forward}).\;$\operatorname{push\text{-}forward}(\omega, \mathit{cb}, \mu, h) :=$ $(\lambda \operatorname{pwb\_} \in \mathit{cb}.\; \mu(\{\operatorname{pwx\_} \in \omega : (h(\operatorname{pwx\_}) \in \operatorname{pwb\_})\}))$

\vnbbookfile{theorem-library/push-forward}{theorem-library/push-forward}
\noindent{\small theorem-library/push-forward.scm -- Borel preimages under a continuous map, the push-forward of a measure, and the change of variables formula (RM-3B, 2026-10-10).  Vocabulary: structure-library/push-forward.scm (PUSH-FORWARD). Source: Thayer, Construction of Measures, Prop. 1.19 and Ch. 2 s. 7,}\par
\vnbprop{pfw-push-forward-unfold}{\vnbstmt{\begin{array}{@{}l@{}} \forall a, b, u, h. \\ \quad (\operatorname{push\text{-}forward}(a, b, u, h) \simeq (\lambda \mathit{pwb} \in b.\; u(\{x \in a : (h(x) \in \mathit{pwb})\}))) \end{array}}}
\vnbprop{pfw-push-forward-value}{\vnbstmt{\begin{array}{@{}l@{}} \forall a, c, u, h.; \forall b \in c. \\ \quad (\operatorname{push\text{-}forward}(a, c, u, h)(b) \simeq u(\{x \in a : (h(x) \in b)\})) \end{array}}}
\vnbprop{pfw-preimage-sigma-algebra}{\vnbstmt{\begin{array}{@{}l@{}} \forall a, b, a_{2}, h. \\ \quad (\operatorname{is\text{-}sigma\text{-}algebra}(a, b) \wedge ((a_{2} \in \mathbf{Set}) \wedge (h \in (a \to a_{2})))) \Rightarrow  \\ \quad \quad \operatorname{is\text{-}sigma\text{-}algebra}(a_{2}, \{e \in \operatorname{power}(a_{2}) : (\{x \in a : (h(x) \in e)\} \in b)\}) \end{array}}}
\vnbprop{pfw-borel-sigma-algebra}{\vnbstmt{\forall s.\; (\operatorname{is\text{-}metric\text{-}space}(s) \Rightarrow \operatorname{is\text{-}sigma\text{-}algebra}(\operatorname{pts}(s), \operatorname{borel\text{-}sets}(s)))}}
\vnbprop{pfw-open-borel}{\vnbstmt{\forall s, u.\; (\operatorname{is\text{-}open}(s, u) \Rightarrow (u \in \operatorname{borel\text{-}sets}(s)))}}
\vnbprop{pfw-borel-preimage}{\vnbstmt{\begin{array}{@{}l@{}} \forall s, t, h, b. \\ \quad (\operatorname{is\text{-}continuous}(s, t, h) \wedge (b \in \operatorname{borel\text{-}sets}(t))) \Rightarrow  \\ \quad \quad \{x \in \operatorname{pts}(s) : (h(x) \in b)\} \in \operatorname{borel\text{-}sets}(s) \end{array}}}
\vnbprop{pfw-continuous-borel-measurable}{\vnbstmt{\begin{array}{@{}l@{}} \forall s, t, h. \\ \quad \operatorname{is\text{-}continuous}(s, t, h) \Rightarrow  \\ \quad \quad \operatorname{is\text{-}measurable\text{-}map}(\operatorname{pts}(s), \operatorname{borel\text{-}sets}(s), \operatorname{pts}(t), \operatorname{borel\text{-}sets}(t), h) \end{array}}}
\vnbprop{pfw-measurable-map-enlarge}{\vnbstmt{\begin{array}{@{}l@{}} \forall a, b, a_{2}, c, h, d. \\ \quad (\operatorname{is\text{-}measurable\text{-}map}(a, b, a_{2}, c, h) \wedge (\operatorname{is\text{-}sigma\text{-}algebra}(a, d) \wedge (b \subseteq d))) \Rightarrow  \\ \quad \quad \operatorname{is\text{-}measurable\text{-}map}(a, d, a_{2}, c, h) \end{array}}}
\vnbprop{pfw-push-forward-measure}{\vnbstmt{\begin{array}{@{}l@{}} \forall a, b, u, a_{2}, c, h. \\ \quad (\operatorname{is\text{-}measure}(a, b, u) \wedge \operatorname{is\text{-}measurable\text{-}map}(a, b, a_{2}, c, h)) \Rightarrow  \\ \quad \quad \operatorname{is\text{-}measure}(a_{2}, c, \operatorname{push\text{-}forward}(a, c, u, h)) \end{array}}}
\vnbprop{pfw-comp-apply}{\vnbstmt{\forall a, f, h.; \forall x \in a.\; ((\lambda z \in a.\; f(h(z)))(x) \simeq f(h(x)))}}
\vnbprop{pfw-comp-fun}{\vnbstmt{\begin{array}{@{}l@{}} \forall a, a_{2}, d, f, h. \\ \quad ((a \in \mathbf{Set}) \wedge ((h \in (a \to a_{2})) \wedge (f \in (a_{2} \to d)))) \Rightarrow  \\ \quad \quad (\lambda z \in a.\; f(h(z))) \in (a \to d) \end{array}}}
\vnbprop{pfw-comp-measurable-fn}{\vnbstmt{\begin{array}{@{}l@{}} \forall a, b, a_{2}, c, h, f. \\ \quad (\operatorname{is\text{-}measurable\text{-}map}(a, b, a_{2}, c, h) \wedge \operatorname{is\text{-}measurable\text{-}fn}(a_{2}, c, f)) \Rightarrow  \\ \quad \quad \operatorname{is\text{-}measurable\text{-}fn}(a, b, (\lambda z \in a.\; f(h(z)))) \end{array}}}
\vnbprop{pfw-simple-change-of-variables}{\vnbstmt{\begin{array}{@{}l@{}} \forall a, b, u, a_{2}, c, h, s. \\ \quad (\operatorname{is\text{-}measure}(a, b, u) \wedge (\operatorname{is\text{-}measurable\text{-}map}(a, b, a_{2}, c, h) \wedge \operatorname{is\text{-}simple\text{-}fn}(a_{2}, c, s))) \Rightarrow  \\ \quad \quad \operatorname{is\text{-}simple\text{-}fn}(a, b, (\lambda z \in a.\; s(h(z)))) \wedge \\ \quad \quad \operatorname{integral}(a_{2}, c, \operatorname{push\text{-}forward}(a, c, u, h), s) = \operatorname{integral}(a, b, u, (\lambda z \in a.\; s(h(z)))) \end{array}}}
\vnbprop{pfw-change-of-variables}{\vnbstmt{\begin{array}{@{}l@{}} \forall a, b, u, a_{2}, c, h, f. \\ \quad (\operatorname{is\text{-}measure}(a, b, u) \wedge (\operatorname{is\text{-}measurable\text{-}map}(a, b, a_{2}, c, h) \wedge \operatorname{is\text{-}measurable\text{-}fn}(a_{2}, c, f))) \Rightarrow  \\ \quad \quad \operatorname{is\text{-}measurable\text{-}fn}(a, b, (\lambda z \in a.\; f(h(z)))) \wedge \\ \quad \quad \operatorname{integral}(a_{2}, c, \operatorname{push\text{-}forward}(a, c, u, h), f) = \operatorname{integral}(a, b, u, (\lambda z \in a.\; f(h(z)))) \end{array}}}
\vnbprop{pfw-pos-part-comp}{\vnbstmt{\begin{array}{@{}l@{}} \forall a, a_{2}.; \forall f \in (a_{2} \to \mathbb{R}), h \in (a \to a_{2}), v \in a. \\ \quad \operatorname{pos\text{-}part}(a, (\lambda z \in a.\; f(h(z))))(v) = (\lambda z \in a.\; \operatorname{pos\text{-}part}(a_{2}, f)(h(z)))(v) \end{array}}}
\vnbprop{pfw-neg-part-comp}{\vnbstmt{\begin{array}{@{}l@{}} \forall a, a_{2}.; \forall f \in (a_{2} \to \mathbb{R}), h \in (a \to a_{2}), v \in a. \\ \quad \operatorname{neg\text{-}part}(a, (\lambda z \in a.\; f(h(z))))(v) = (\lambda z \in a.\; \operatorname{neg\text{-}part}(a_{2}, f)(h(z)))(v) \end{array}}}
\vnbprop{pfw-comp-real-measurable}{\vnbstmt{\begin{array}{@{}l@{}} \forall a, b, a_{2}, c, h, f. \\ \quad (\operatorname{is\text{-}measurable\text{-}map}(a, b, a_{2}, c, h) \wedge \operatorname{is\text{-}real\text{-}measurable\text{-}fn}(a_{2}, c, f)) \Rightarrow  \\ \quad \quad \operatorname{is\text{-}real\text{-}measurable\text{-}fn}(a, b, (\lambda z \in a.\; f(h(z)))) \end{array}}}
\vnbprop{pfw-integrable-iff}{\vnbstmt{\begin{array}{@{}l@{}} \forall a, b, u, a_{2}, c, h, f. \\ \quad (\operatorname{is\text{-}measure}(a, b, u) \wedge (\operatorname{is\text{-}measurable\text{-}map}(a, b, a_{2}, c, h) \wedge \operatorname{is\text{-}real\text{-}measurable\text{-}fn}(a_{2}, c, f))) \Rightarrow  \\ \quad \quad \operatorname{is\text{-}integrable}(a_{2}, c, \operatorname{push\text{-}forward}(a, c, u, h), f) \\ \quad \quad \Leftrightarrow\; \operatorname{is\text{-}integrable}(a, b, u, (\lambda z \in a.\; f(h(z)))) \end{array}}}
\vnbprop{pfw-change-of-variables-real}{\vnbstmt{\begin{array}{@{}l@{}} \forall a, b, u, a_{2}, c, h, f. \\ \quad \operatorname{is\text{-}measure}(a, b, u) \wedge \\ \quad \operatorname{is\text{-}measurable\text{-}map}(a, b, a_{2}, c, h) \wedge \\ \quad \operatorname{is\text{-}real\text{-}measurable\text{-}fn}(a_{2}, c, f) \wedge \\ \quad \operatorname{is\text{-}integrable}(a_{2}, c, \operatorname{push\text{-}forward}(a, c, u, h), f) \Rightarrow \\ \quad \quad \operatorname{is\text{-}integrable}(a, b, u, (\lambda z \in a.\; f(h(z)))) \wedge \\ \quad \quad \operatorname{rintegral}(a_{2}, c, \operatorname{push\text{-}forward}(a, c, u, h), f) = \operatorname{rintegral}(a, b, u, (\lambda z \in a.\; f(h(z)))) \end{array}}}

\vnbbookfile{theorem-library/riesz-markov}{theorem-library/riesz-markov}
\noindent{\small theorem-library/riesz-markov.scm -- RIESZ-MARKOV, STAGE 3B (RM-3B, 2026-10-10): the assembly on X, conditional on the positive extension.  Helper prefix: rmk-. Source: the user's notes-63, form (b).}\par
\vnbprop{rmk-riesz-markov-from-k}{\vnbstmt{\begin{array}{@{}l@{}} \forall s, q, m, x. \\ \quad \operatorname{is\text{-}compact}(s) \wedge \\ \quad \operatorname{is\text{-}dense\text{-}seq}(s, q) \wedge \\ \quad \operatorname{is\text{-}positive\text{-}functional}(\operatorname{rm\text{-}kms}(s, q), m) \wedge \\ \quad \operatorname{is\text{-}positive\text{-}functional}(s, x) \wedge \\ \quad \forall f \in \operatorname{cont\text{-}fns}(s). \\ \quad \quad m((\lambda z \in \operatorname{pts}(\operatorname{rm\text{-}kms}(s, q)).\; f(\operatorname{rm\text{-}pi}(s, q)(z)))) = x(f) \Rightarrow \\ \quad \quad \operatorname{is\text{-}measure}(\operatorname{pts}(s), \operatorname{borel\text{-}sets}(s), \\ \quad \quad \quad \operatorname{push\text{-}forward}(\operatorname{pts}(\operatorname{rm\text{-}kms}(s, q)), \operatorname{borel\text{-}sets}(s), \\ \quad \quad \quad \quad \operatorname{res}(\operatorname{method\text{-}i\text{-}outer}(\operatorname{pts}(\operatorname{rm\text{-}kms}(s, q)), \operatorname{clopen\text{-}sets}(\operatorname{rm\text{-}kms}(s, q)), \\ \quad \quad \quad \quad \quad \quad (\lambda u \in \operatorname{clopen\text{-}sets}(\operatorname{rm\text{-}kms}(s, q)).\; m(\operatorname{indicator}(\operatorname{pts}(\operatorname{rm\text{-}kms}(s, q)), u)))), \\ \quad \quad \quad \quad \quad \operatorname{caratheodory\text{-}sets}(\operatorname{pts}(\operatorname{rm\text{-}kms}(s, q)), \\ \quad \quad \quad \quad \quad \quad \operatorname{method\text{-}i\text{-}outer}(\operatorname{pts}(\operatorname{rm\text{-}kms}(s, q)), \operatorname{clopen\text{-}sets}(\operatorname{rm\text{-}kms}(s, q)), \\ \quad \quad \quad \quad \quad \quad \quad (\lambda u \in \operatorname{clopen\text{-}sets}(\operatorname{rm\text{-}kms}(s, q)).\; m(\operatorname{indicator}(\operatorname{pts}(\operatorname{rm\text{-}kms}(s, q)), u)))))), \\ \quad \quad \quad \quad \operatorname{rm\text{-}pi}(s, q))) \wedge \\ \quad \quad \forall f \in \operatorname{cont\text{-}fns}(s). \\ \quad \quad \quad \operatorname{is\text{-}integrable}(\operatorname{pts}(s), \operatorname{borel\text{-}sets}(s), \\ \quad \quad \quad \quad \operatorname{push\text{-}forward}(\operatorname{pts}(\operatorname{rm\text{-}kms}(s, q)), \operatorname{borel\text{-}sets}(s), \\ \quad \quad \quad \quad \quad \operatorname{res}(\operatorname{method\text{-}i\text{-}outer}(\operatorname{pts}(\operatorname{rm\text{-}kms}(s, q)), \operatorname{clopen\text{-}sets}(\operatorname{rm\text{-}kms}(s, q)), \\ \quad \quad \quad \quad \quad \quad \quad (\lambda u \in \operatorname{clopen\text{-}sets}(\operatorname{rm\text{-}kms}(s, q)).\; m(\operatorname{indicator}(\operatorname{pts}(\operatorname{rm\text{-}kms}(s, q)), u)))), \\ \quad \quad \quad \quad \quad \quad \operatorname{caratheodory\text{-}sets}(\operatorname{pts}(\operatorname{rm\text{-}kms}(s, q)), \\ \quad \quad \quad \quad \quad \quad \quad \operatorname{method\text{-}i\text{-}outer}(\operatorname{pts}(\operatorname{rm\text{-}kms}(s, q)), \operatorname{clopen\text{-}sets}(\operatorname{rm\text{-}kms}(s, q)), \\ \quad \quad \quad \quad \quad \quad \quad \quad (\lambda u \in \operatorname{clopen\text{-}sets}(\operatorname{rm\text{-}kms}(s, q)).\; m(\operatorname{indicator}(\operatorname{pts}(\operatorname{rm\text{-}kms}(s, q)), u)))))), \\ \quad \quad \quad \quad \quad \operatorname{rm\text{-}pi}(s, q)), \\ \quad \quad \quad \quad f) \wedge \\ \quad \quad \quad \operatorname{rintegral}(\operatorname{pts}(s), \operatorname{borel\text{-}sets}(s), \\ \quad \quad \quad \quad \operatorname{push\text{-}forward}(\operatorname{pts}(\operatorname{rm\text{-}kms}(s, q)), \operatorname{borel\text{-}sets}(s), \\ \quad \quad \quad \quad \quad \operatorname{res}(\operatorname{method\text{-}i\text{-}outer}(\operatorname{pts}(\operatorname{rm\text{-}kms}(s, q)), \operatorname{clopen\text{-}sets}(\operatorname{rm\text{-}kms}(s, q)), \\ \quad \quad \quad \quad \quad \quad \quad (\lambda u \in \operatorname{clopen\text{-}sets}(\operatorname{rm\text{-}kms}(s, q)).\; m(\operatorname{indicator}(\operatorname{pts}(\operatorname{rm\text{-}kms}(s, q)), u)))), \\ \quad \quad \quad \quad \quad \quad \operatorname{caratheodory\text{-}sets}(\operatorname{pts}(\operatorname{rm\text{-}kms}(s, q)), \\ \quad \quad \quad \quad \quad \quad \quad \operatorname{method\text{-}i\text{-}outer}(\operatorname{pts}(\operatorname{rm\text{-}kms}(s, q)), \operatorname{clopen\text{-}sets}(\operatorname{rm\text{-}kms}(s, q)), \\ \quad \quad \quad \quad \quad \quad \quad \quad (\lambda u \in \operatorname{clopen\text{-}sets}(\operatorname{rm\text{-}kms}(s, q)).\; m(\operatorname{indicator}(\operatorname{pts}(\operatorname{rm\text{-}kms}(s, q)), u)))))), \\ \quad \quad \quad \quad \quad \operatorname{rm\text{-}pi}(s, q)), \\ \quad \quad \quad \quad f) \\ \quad \quad \quad =\; x(f) \end{array}}}

\vnbbookfile{theorem-library/riesz-markov-compact}{theorem-library/riesz-markov-compact}
\noindent{\small theorem-library/riesz-markov-compact.scm -- RIESZ-MARKOV on a compact metric space (RM-3A follow-up, 2026-10-10).  Prefix rmz-.  Source: the user's notes-63, form (b).}\par
\vnbprop{rmz-comp-cont}{\vnbstmt{\begin{array}{@{}l@{}} \forall s, q.; \forall f \in \operatorname{cont\text{-}fns}(s). \\ \quad (\operatorname{is\text{-}compact}(s) \wedge \operatorname{is\text{-}dense\text{-}seq}(s, q)) \Rightarrow  \\ \quad \quad (\lambda z \in \operatorname{pts}(\operatorname{rm\text{-}kms}(s, q)).\; f(\operatorname{rm\text{-}pi}(s, q)(z))) \in \operatorname{cont\text{-}fns}(\operatorname{rm\text{-}kms}(s, q)) \end{array}}}
\vnbprop{rmz-one-cont}{\vnbstmt{\forall s.\; (\operatorname{is\text{-}metric\text{-}space}(s) \Rightarrow ((\lambda x \in \operatorname{pts}(s).\; 1) \in \operatorname{cont\text{-}fns}(s)))}}
\vnbprop{rmz-lm-val}{\vnbstmt{\begin{array}{@{}l@{}} \forall s, q, x.; \forall f \in \operatorname{cont\text{-}fns}(s). \\ \quad (\operatorname{is\text{-}compact}(s) \wedge (\operatorname{is\text{-}dense\text{-}seq}(s, q) \wedge \operatorname{is\text{-}positive\text{-}functional}(s, x))) \Rightarrow  \\ \quad \quad (\lambda z \in \operatorname{pts}(\operatorname{rm\text{-}kms}(s, q)).\; f(\operatorname{rm\text{-}pi}(s, q)(z))) \\ \quad \quad \in\; \{g \in \operatorname{cont\text{-}fns}(\operatorname{rm\text{-}kms}(s, q)) : \exists \mathit{rzf} \in \operatorname{cont\text{-}fns}(s).\; (g = (\lambda z \in \operatorname{pts}(\operatorname{rm\text{-}kms}(s, q)).\; \operatorname{rzf}(\operatorname{rm\text{-}pi}(s, q)(z))))\} \wedge \\ \quad \quad (\lambda g \in \{g \in \operatorname{cont\text{-}fns}(\operatorname{rm\text{-}kms}(s, q)) : \exists \mathit{rzf} \in \operatorname{cont\text{-}fns}(s).\; (g = (\lambda z \in \operatorname{pts}(\operatorname{rm\text{-}kms}(s, q)).\; \operatorname{rzf}(\operatorname{rm\text{-}pi}(s, q)(z))))\}.\; x((\lambda \mathit{rzx} \in \operatorname{pts}(s).\; g(\operatorname{rm\text{-}code}(s, q, \mathit{rzx})))))((\lambda z \in \operatorname{pts}(\operatorname{rm\text{-}kms}(s, q)).\; f(\operatorname{rm\text{-}pi}(s, q)(z)))) \\ \quad \quad =\; x(f) \end{array}}}
\vnbprop{rmz-riesz-markov-dense}{\vnbstmt{\begin{array}{@{}l@{}} \forall s, q, x. \\ \quad (\operatorname{is\text{-}compact}(s) \wedge (\operatorname{is\text{-}dense\text{-}seq}(s, q) \wedge \operatorname{is\text{-}positive\text{-}functional}(s, x))) \Rightarrow  \\ \quad \quad \exists u. \\ \quad \quad \quad \operatorname{is\text{-}measure}(\operatorname{pts}(s), \operatorname{borel\text{-}sets}(s), u) \wedge \\ \quad \quad \quad \forall f \in \operatorname{cont\text{-}fns}(s). \\ \quad \quad \quad \quad (\operatorname{is\text{-}integrable}(\operatorname{pts}(s), \operatorname{borel\text{-}sets}(s), u, f) \wedge (\operatorname{rintegral}(\operatorname{pts}(s), \operatorname{borel\text{-}sets}(s), u, f) = x(f))) \end{array}}}
\vnbprop{rmz-riesz-markov}{\vnbstmt{\begin{array}{@{}l@{}} \forall s, a. \\ \quad (\operatorname{is\text{-}compact}(s) \wedge (\exists x.\; (x \in \operatorname{pts}(s)) \wedge \operatorname{is\text{-}positive\text{-}functional}(s, a))) \Rightarrow  \\ \quad \quad \exists u. \\ \quad \quad \quad \operatorname{is\text{-}measure}(\operatorname{pts}(s), \operatorname{borel\text{-}sets}(s), u) \wedge \\ \quad \quad \quad \forall f \in \operatorname{cont\text{-}fns}(s). \\ \quad \quad \quad \quad (\operatorname{is\text{-}integrable}(\operatorname{pts}(s), \operatorname{borel\text{-}sets}(s), u, f) \wedge (\operatorname{rintegral}(\operatorname{pts}(s), \operatorname{borel\text{-}sets}(s), u, f) = a(f))) \end{array}}}

\part{Strand Spaces}
\noindent Messages as tokens, lists and encryptions; strands, strand spaces, nodes and the two node relations; bundles and causal precedence; origination and entry points; the adversary and the bound on what it can originate.
\par\noindent\textit{Builds on:} Foundations.
\par\noindent{\small 8 files.}

\vnbbookfile{structure-library/strand-messages}{structure-library/strand-messages}
\noindent{\small strand-messages.scm -- THE MESSAGES OF STRAND SPACES, ON THE USER'S MODEL, WITH ENCRYPTION (\textasciitilde{}/notes/notes-60.txt, item (a): "Let T be a set of tokens.  A message is a nested list of tokens, i.e. a token is a message and any list of messages is a message.  Subterms are defined recursively in the obvious way."  The user, ...}\par
\par\medskip\noindent\textbf{Definition} (\texttt{encr}).\;$\operatorname{encr}(\operatorname{smg\_}, \operatorname{smk\_}) :=$ $\langle 2, \operatorname{smg\_}, \operatorname{smk\_} \rangle$
\par\medskip\noindent\textbf{Definition} (\texttt{lst}).\;$\operatorname{lst}(\operatorname{sml\_}) :=$ $\langle 1, \operatorname{sml\_} \rangle$
\par\medskip\noindent\textbf{Definition} (\texttt{is-token-set}).\;For every $\operatorname{smt\_}$, 
\begin{equation*}
\operatorname{is\text{-}token\text{-}set}(\operatorname{smt\_}) \Leftrightarrow ((\operatorname{smt\_} \in \mathbf{Set}) \wedge \forall \operatorname{smx\_} \in \operatorname{smt\_}.\; \neg (\operatorname{smx\_} \in \operatorname{tuples}(\mathbf{Set})))
\end{equation*}

\par\medskip\noindent\textbf{Definition} (\texttt{msg-set}).\;$\operatorname{msg\text{-}set}(\operatorname{smt\_}, \operatorname{smk\_}) :=$ $\operatorname{big\text{-}union}(\operatorname{smn\_}, \mathbb{N}, \operatorname{msg\text{-}stage}(\operatorname{smt\_}, \operatorname{smk\_}, \operatorname{smn\_}))$
\par\medskip\noindent\textbf{Definition} (\texttt{msg-stage}).\;For all $\operatorname{smt\_}$ and $\operatorname{smk\_}$,  $(\operatorname{msg\text{-}stage}(\operatorname{smt\_}, \operatorname{smk\_}, 0) \simeq \operatorname{smt\_})$\par
Let $\operatorname{smt\_}$ and $\operatorname{smk\_}$ be arbitrary. Let $\operatorname{smn\_} \in \mathbb{N}$. Then
\begin{equation*}
(\operatorname{msg\text{-}stage}(\operatorname{smt\_}, \operatorname{smk\_}, \operatorname{succ}(\operatorname{smn\_})) \simeq (\operatorname{msg\text{-}stage}(\operatorname{smt\_}, \operatorname{smk\_}, \operatorname{smn\_}) \cup ((\{1\} \times \operatorname{tuples}(\operatorname{msg\text{-}stage}(\operatorname{smt\_}, \operatorname{smk\_}, \operatorname{smn\_}))) \cup (\{2\} \times \operatorname{msg\text{-}stage}(\operatorname{smt\_}, \operatorname{smk\_}, \operatorname{smn\_})))))
\end{equation*}

\par\medskip\noindent\textbf{Definition} (\texttt{subterm-rel}).\;$\operatorname{subterm\text{-}rel}(\operatorname{smt\_}, \operatorname{smk\_}) :=$
\begin{equation*}
\{\operatorname{smp\_} \in (\operatorname{msg\text{-}set}(\operatorname{smt\_}, \operatorname{smk\_}) \times \operatorname{msg\text{-}set}(\operatorname{smt\_}, \operatorname{smk\_})) : \forall \operatorname{smq\_} \in \operatorname{power}((\operatorname{msg\text{-}set}(\operatorname{smt\_}, \operatorname{smk\_}) \times \operatorname{msg\text{-}set}(\operatorname{smt\_}, \operatorname{smk\_}))).\; (\operatorname{is\text{-}subterm\text{-}closed}(\operatorname{smt\_}, \operatorname{smk\_}, \operatorname{smq\_}) \Rightarrow (\operatorname{smp\_} \in \operatorname{smq\_}))\}
\end{equation*}

\par\medskip\noindent\textbf{Definition} (\texttt{is-subterm}).\;For all $\operatorname{smt\_}$, $\operatorname{smk\_}$, $\operatorname{sma\_}$ and $\operatorname{smb\_}$, 
\begin{equation*}
\operatorname{is\text{-}subterm}(\operatorname{smt\_}, \operatorname{smk\_}, \operatorname{sma\_}, \operatorname{smb\_}) \Leftrightarrow (\langle \operatorname{sma\_}, \operatorname{smb\_} \rangle \in \operatorname{subterm\text{-}rel}(\operatorname{smt\_}, \operatorname{smk\_}))
\end{equation*}

\par\medskip\noindent\textbf{Definition} (\texttt{msg-rank}).\;$\operatorname{msg\text{-}rank}(\operatorname{smt\_}, \operatorname{smk\_}, \operatorname{smm\_}) :=$
\begin{equation*}
(\iota\, \operatorname{smr\_}.\;((\operatorname{smr\_} \in \mathbb{N}) \wedge ((\operatorname{smm\_} \in \operatorname{msg\text{-}stage}(\operatorname{smt\_}, \operatorname{smk\_}, \operatorname{smr\_})) \wedge \forall \operatorname{smj\_} \in \mathbb{N}.\; ((\operatorname{smm\_} \in \operatorname{msg\text{-}stage}(\operatorname{smt\_}, \operatorname{smk\_}, \operatorname{smj\_})) \Rightarrow (\operatorname{smr\_} \leq \operatorname{smj\_})))))
\end{equation*}

\par\medskip\noindent\textbf{Definition} (\texttt{is-subterm-closed}).\;For all $\operatorname{smt\_}$, $\operatorname{smk\_}$ and $\operatorname{smr\_}$, 
\begin{equation*}
\begin{array}{@{}l@{}} \operatorname{is\text{-}subterm\text{-}closed}(\operatorname{smt\_}, \operatorname{smk\_}, \operatorname{smr\_}) \\ \Leftrightarrow\; \forall \operatorname{sma\_} \in \operatorname{msg\text{-}set}(\operatorname{smt\_}, \operatorname{smk\_}).\; (\langle \operatorname{sma\_}, \operatorname{sma\_} \rangle \in \operatorname{smr\_}) \wedge \\ \forall \operatorname{sma\_}.; \forall \operatorname{sml\_} \in \operatorname{tuples}(\operatorname{msg\text{-}set}(\operatorname{smt\_}, \operatorname{smk\_})), \operatorname{smi\_} \in \mathbb{N}. \\ \quad (1 \leq \operatorname{smi\_}) \Rightarrow  \\ \quad (\operatorname{smi\_} \leq \operatorname{length}(\operatorname{sml\_})) \Rightarrow  \\ \quad (\langle \operatorname{sma\_}, \operatorname{nth}(\operatorname{smi\_}, \operatorname{sml\_}) \rangle \in \operatorname{smr\_}) \Rightarrow (\langle \operatorname{sma\_}, \operatorname{lst}(\operatorname{sml\_}) \rangle \in \operatorname{smr\_}) \wedge \\ \forall \operatorname{sma\_}.; \forall \operatorname{smg\_} \in \operatorname{msg\text{-}set}(\operatorname{smt\_}, \operatorname{smk\_}), \operatorname{smy\_} \in \operatorname{smk\_}. \\ \quad (\langle \operatorname{sma\_}, \operatorname{smg\_} \rangle \in \operatorname{smr\_}) \Rightarrow (\langle \operatorname{sma\_}, \operatorname{encr}(\operatorname{smg\_}, \operatorname{smy\_}) \rangle \in \operatorname{smr\_}) \end{array}
\end{equation*}


\vnbbookfile{theorem-library/strand-messages-laws}{theorem-library/strand-messages-laws}
\noindent{\small strand-messages-laws.scm -- THE MESSAGES OF STRAND SPACES (the user's model, with ENCRYPTION): their laws. (\textasciitilde{}/notes/notes-60.txt item (a), the user's decision of 2026-10-06 16:40Z adding \{m\}\_k; the definitions are structure-library/strand-messages.scm.)}\par
\vnbprop{lst-unfold}{\vnbstmt{\forall l.\; (\operatorname{lst}(l) \simeq \langle 1, l \rangle)}}
\vnbprop{encr-unfold}{\vnbstmt{\forall g, y.\; (\operatorname{encr}(g, y) \simeq \langle 2, g, y \rangle)}}
\vnbprop{encr-ne-list}{\vnbstmt{\forall g, y, l.\; \neg (\operatorname{encr}(g, y) = \operatorname{lst}(l))}}
\vnbprop{lst-ne-encr}{\vnbstmt{\forall l, g, y.\; \neg (\operatorname{lst}(l) = \operatorname{encr}(g, y))}}
\vnbprop{lst-injective}{\vnbstmt{\forall l, a.\; ((\operatorname{lst}(l) = \operatorname{lst}(a)) \Rightarrow (l = a))}}
\vnbprop{encr-injective}{\vnbstmt{\forall g, y, u, z.\; ((\operatorname{encr}(g, y) = \operatorname{encr}(u, z)) \Rightarrow ((g = u) \wedge (y = z)))}}
\vnbprop{msg-stage-token}{\vnbstmt{\forall t, k.; \forall a \in t.\; (a \in \operatorname{msg\text{-}stage}(t, k, 0))}}
\vnbprop{msg-stage-keep}{\vnbstmt{\begin{array}{@{}l@{}} \forall t, k, n, m. \\ \quad ((n \in \mathbb{N}) \wedge (m \in \operatorname{msg\text{-}stage}(t, k, n))) \Rightarrow (m \in \operatorname{msg\text{-}stage}(t, k, (n + 1))) \end{array}}}
\vnbprop{msg-stage-list}{\vnbstmt{\begin{array}{@{}l@{}} \forall t, k, n, l. \\ \quad ((n \in \mathbb{N}) \wedge (l \in \operatorname{tuples}(\operatorname{msg\text{-}stage}(t, k, n)))) \Rightarrow  \\ \quad \quad \operatorname{lst}(l) \in \operatorname{msg\text{-}stage}(t, k, (n + 1)) \end{array}}}
\vnbprop{msg-stage-encr}{\vnbstmt{\begin{array}{@{}l@{}} \forall t, k, n, g, y. \\ \quad ((n \in \mathbb{N}) \wedge ((g \in \operatorname{msg\text{-}stage}(t, k, n)) \wedge (y \in k))) \Rightarrow  \\ \quad \quad \operatorname{encr}(g, y) \in \operatorname{msg\text{-}stage}(t, k, (n + 1)) \end{array}}}
\vnbprop{msg-stage-zero-elim}{\vnbstmt{\forall t, k.; \forall m \in \operatorname{msg\text{-}stage}(t, k, 0).\; (m \in t)}}
\vnbprop{msg-stage-succ-cases}{\vnbstmt{\begin{array}{@{}l@{}} \forall t, k, n, m. \\ \quad ((n \in \mathbb{N}) \wedge (m \in \operatorname{msg\text{-}stage}(t, k, (n + 1)))) \Rightarrow  \\ \quad \quad m \in \operatorname{msg\text{-}stage}(t, k, n) \vee \\ \quad \quad \exists l \in \operatorname{tuples}(\operatorname{msg\text{-}stage}(t, k, n)).\; (m = \operatorname{lst}(l)) \vee \\ \quad \quad \exists g \in \operatorname{msg\text{-}stage}(t, k, n), y \in k.\; (m = \operatorname{encr}(g, y)) \end{array}}}
\vnbprop{msg-stage-add}{\vnbstmt{\begin{array}{@{}l@{}} \forall d \in \mathbb{N}, t, k, j, m. \\ \quad ((j \in \mathbb{N}) \wedge (m \in \operatorname{msg\text{-}stage}(t, k, j))) \Rightarrow (m \in \operatorname{msg\text{-}stage}(t, k, (j + d))) \end{array}}}
\vnbprop{msg-set-unfold}{\vnbstmt{\forall t, k.\; (\operatorname{msg\text{-}set}(t, k) \simeq \operatorname{big\text{-}union}(n, \mathbb{N}, \operatorname{msg\text{-}stage}(t, k, n)))}}
\vnbprop{msg-set-intro}{\vnbstmt{\begin{array}{@{}l@{}} \forall t, k, n, m. \\ \quad ((n \in \mathbb{N}) \wedge (m \in \operatorname{msg\text{-}stage}(t, k, n))) \Rightarrow (m \in \operatorname{msg\text{-}set}(t, k)) \end{array}}}
\vnbprop{msg-set-elim}{\vnbstmt{\forall t, k.; \forall m \in \operatorname{msg\text{-}set}(t, k).\; \exists n \in \mathbb{N}.\; (m \in \operatorname{msg\text{-}stage}(t, k, n))}}
\vnbprop{token-in-msgs}{\vnbstmt{\forall t, k.; \forall a \in t.\; (a \in \operatorname{msg\text{-}set}(t, k))}}
\vnbprop{msg-stage-set}{\vnbstmt{\forall n \in \mathbb{N}, t, k.\; (((t \in \mathbf{Set}) \wedge (k \in \mathbf{Set})) \Rightarrow (\operatorname{msg\text{-}stage}(t, k, n) \in \mathbf{Set}))}}
\vnbprop{msg-set-is-set}{\vnbstmt{\forall t, k.\; (((t \in \mathbf{Set}) \wedge (k \in \mathbf{Set})) \Rightarrow (\operatorname{msg\text{-}set}(t, k) \in \mathbf{Set}))}}
\vnbprop{msg-lists-are-lists}{\vnbstmt{\forall t, k.; \forall l \in \operatorname{tuples}(\operatorname{msg\text{-}set}(t, k)).\; (l \in \operatorname{tuples}(\mathbf{Set}))}}
\vnbprop{list-in-some-stage}{\vnbstmt{\begin{array}{@{}l@{}} \forall t, k.; \forall l \in \operatorname{tuples}(\operatorname{msg\text{-}set}(t, k)). \exists n \in \mathbb{N}. \\ \quad l \in \operatorname{tuples}(\operatorname{msg\text{-}stage}(t, k, n)) \end{array}}}
\vnbprop{list-in-msgs}{\vnbstmt{\forall t, k.; \forall l \in \operatorname{tuples}(\operatorname{msg\text{-}set}(t, k)).\; (\operatorname{lst}(l) \in \operatorname{msg\text{-}set}(t, k))}}
\vnbprop{encr-in-msgs}{\vnbstmt{\begin{array}{@{}l@{}} \forall t, k, g, y. \\ \quad ((g \in \operatorname{msg\text{-}set}(t, k)) \wedge (y \in k)) \Rightarrow (\operatorname{encr}(g, y) \in \operatorname{msg\text{-}set}(t, k)) \end{array}}}
\vnbprop{lst-in-tuples-set}{\vnbstmt{\forall l \in \mathbf{Set}.\; (\operatorname{lst}(l) \in \operatorname{tuples}(\mathbf{Set}))}}
\vnbprop{encr-in-tuples-set}{\vnbstmt{\forall g, y.\; (((g \in \mathbf{Set}) \wedge (y \in \mathbf{Set})) \Rightarrow (\operatorname{encr}(g, y) \in \operatorname{tuples}(\mathbf{Set})))}}
\vnbprop{token-set-is-set}{\vnbstmt{\forall t.\; (\operatorname{is\text{-}token\text{-}set}(t) \Rightarrow (t \in \mathbf{Set}))}}
\vnbprop{token-not-list}{\vnbstmt{\forall t, k.; \forall a \in t.\; (\operatorname{is\text{-}token\text{-}set}(t) \Rightarrow \neg (a \in \operatorname{tuples}(\operatorname{msg\text{-}set}(t, k))))}}
\vnbprop{token-ne-list}{\vnbstmt{\begin{array}{@{}l@{}} \forall t, k, l. \\ \quad (\operatorname{is\text{-}token\text{-}set}(t) \wedge ((k \in t) \wedge (l \in \mathbf{Set}))) \Rightarrow \neg (k = \operatorname{lst}(l)) \end{array}}}
\vnbprop{encr-ne-token}{\vnbstmt{\begin{array}{@{}l@{}} \forall t, k, g, y. \\ \quad (\operatorname{is\text{-}token\text{-}set}(t) \wedge ((k \in t) \wedge ((g \in \mathbf{Set}) \wedge (y \in \mathbf{Set})))) \Rightarrow \neg (k = \operatorname{encr}(g, y)) \end{array}}}
\vnbprop{msg-induction}{\vnbstmt{\begin{array}{@{}l@{}} \forall t, k, c.; \forall m \in \operatorname{msg\text{-}set}(t, k). \\ \quad \forall a \in t.\; (a \in c) \wedge \\ \quad \forall l \in \operatorname{tuples}(\operatorname{msg\text{-}set}(t, k)). \\ \quad \quad \forall i \in \mathbb{N}.\; (((1 \leq i) \wedge (i \leq \operatorname{length}(l))) \Rightarrow (\operatorname{nth}(i, l) \in c)) \Rightarrow  \\ \quad \quad \quad \operatorname{lst}(l) \in c \wedge \\ \quad \forall g, y.\; (((g \in \operatorname{msg\text{-}set}(t, k)) \wedge ((y \in k) \wedge (g \in c))) \Rightarrow (\operatorname{encr}(g, y) \in c)) \Rightarrow \\ \quad \quad m \in c \end{array}}}
\vnbprop{msg-cases}{\vnbstmt{\begin{array}{@{}l@{}} \forall t, k.; \forall m \in \operatorname{msg\text{-}set}(t, k). \\ \quad m \in t \vee \\ \quad \exists l \in \operatorname{tuples}(\operatorname{msg\text{-}set}(t, k)).\; (m = \operatorname{lst}(l)) \vee \\ \quad \exists g \in \operatorname{msg\text{-}set}(t, k), y \in k.\; (m = \operatorname{encr}(g, y)) \end{array}}}
\vnbprop{pair-in-cartesian-parts}{\vnbstmt{\forall x, y, u, v.\; ((\langle u, v \rangle \in (x \times y)) \Rightarrow ((u \in x) \wedge (v \in y)))}}
\vnbprop{subterm-rel-unfold}{\vnbstmt{\begin{array}{@{}l@{}} \forall t, k. \\ \quad (\operatorname{subterm\text{-}rel}(t, k) \simeq \{p \in (\operatorname{msg\text{-}set}(t, k) \times \operatorname{msg\text{-}set}(t, k)) : \forall q \in \operatorname{power}((\operatorname{msg\text{-}set}(t, k) \times \operatorname{msg\text{-}set}(t, k))).\; (\operatorname{is\text{-}subterm\text{-}closed}(t, k, q) \Rightarrow (p \in q))\}) \end{array}}}
\vnbprop{subterm-in-msgs}{\vnbstmt{\begin{array}{@{}l@{}} \forall t, k, a, b. \\ \quad \operatorname{is\text{-}subterm}(t, k, a, b) \Rightarrow ((a \in \operatorname{msg\text{-}set}(t, k)) \wedge (b \in \operatorname{msg\text{-}set}(t, k))) \end{array}}}
\vnbprop{subterm-least}{\vnbstmt{\begin{array}{@{}l@{}} \forall t, k.; \forall r \in \operatorname{power}((\operatorname{msg\text{-}set}(t, k) \times \operatorname{msg\text{-}set}(t, k))), a, b. \\ \quad (\operatorname{is\text{-}subterm\text{-}closed}(t, k, r) \wedge \operatorname{is\text{-}subterm}(t, k, a, b)) \Rightarrow (\langle a, b \rangle \in r) \end{array}}}
\vnbprop{subterm-refl}{\vnbstmt{\forall t, k.; \forall a \in \operatorname{msg\text{-}set}(t, k).\; \operatorname{is\text{-}subterm}(t, k, a, a)}}
\vnbprop{subterm-component}{\vnbstmt{\begin{array}{@{}l@{}} \forall t, k, a, l, i. \\ \quad l \in \operatorname{tuples}(\operatorname{msg\text{-}set}(t, k)) \wedge \\ \quad i \in \mathbb{N} \wedge \\ \quad 1 \leq i \wedge \\ \quad i \leq \operatorname{length}(l) \wedge \\ \quad \operatorname{is\text{-}subterm}(t, k, a, \operatorname{nth}(i, l)) \Rightarrow \\ \quad \quad \operatorname{is\text{-}subterm}(t, k, a, \operatorname{lst}(l)) \end{array}}}
\vnbprop{subterm-encr}{\vnbstmt{\begin{array}{@{}l@{}} \forall t, k, a, g, y. \\ \quad ((g \in \operatorname{msg\text{-}set}(t, k)) \wedge ((y \in k) \wedge \operatorname{is\text{-}subterm}(t, k, a, g))) \Rightarrow  \\ \quad \quad \operatorname{is\text{-}subterm}(t, k, a, \operatorname{encr}(g, y)) \end{array}}}
\vnbprop{subterm-cases}{\vnbstmt{\begin{array}{@{}l@{}} \forall t, k, a, b. \\ \quad ((t \in \mathbf{Set}) \wedge ((k \in \mathbf{Set}) \wedge \operatorname{is\text{-}subterm}(t, k, a, b))) \Rightarrow  \\ \quad \quad a = b \vee \\ \quad \quad \exists l \in \operatorname{tuples}(\operatorname{msg\text{-}set}(t, k)). \\ \quad \quad \quad ((b = \operatorname{lst}(l)) \wedge \exists i \in \mathbb{N}.\; ((1 \leq i) \wedge ((i \leq \operatorname{length}(l)) \wedge \operatorname{is\text{-}subterm}(t, k, a, \operatorname{nth}(i, l))))) \vee \\ \quad \quad \exists g \in \operatorname{msg\text{-}set}(t, k), y \in k.\; ((b = \operatorname{encr}(g, y)) \wedge \operatorname{is\text{-}subterm}(t, k, a, g)) \end{array}}}
\vnbprop{subterm-trans}{\vnbstmt{\begin{array}{@{}l@{}} \forall t, k, a, b, c. \\ \quad ((t \in \mathbf{Set}) \wedge ((k \in \mathbf{Set}) \wedge (\operatorname{is\text{-}subterm}(t, k, a, b) \wedge \operatorname{is\text{-}subterm}(t, k, b, c)))) \Rightarrow  \\ \quad \quad \operatorname{is\text{-}subterm}(t, k, a, c) \end{array}}}
\vnbprop{subterm-of-token}{\vnbstmt{\begin{array}{@{}l@{}} \forall t, k, a.; \forall b \in t. \\ \quad (\operatorname{is\text{-}token\text{-}set}(t) \wedge ((k \in \mathbf{Set}) \wedge \operatorname{is\text{-}subterm}(t, k, a, b))) \Rightarrow (a = b) \end{array}}}
\vnbprop{subterm-of-list}{\vnbstmt{\begin{array}{@{}l@{}} \forall t, k, a.; \forall l \in \operatorname{tuples}(\operatorname{msg\text{-}set}(t, k)). \\ \quad ((t \in \mathbf{Set}) \wedge ((k \in \mathbf{Set}) \wedge \operatorname{is\text{-}subterm}(t, k, a, \operatorname{lst}(l)))) \Rightarrow  \\ \quad \quad ((a = \operatorname{lst}(l)) \vee \exists i \in \mathbb{N}.\; ((1 \leq i) \wedge ((i \leq \operatorname{length}(l)) \wedge \operatorname{is\text{-}subterm}(t, k, a, \operatorname{nth}(i, l))))) \end{array}}}
\vnbprop{subterm-of-encr}{\vnbstmt{\begin{array}{@{}l@{}} \forall t, k, a, g, y. \\ \quad ((t \in \mathbf{Set}) \wedge ((k \in \mathbf{Set}) \wedge ((g \in \operatorname{msg\text{-}set}(t, k)) \wedge ((y \in k) \wedge \operatorname{is\text{-}subterm}(t, k, a, \operatorname{encr}(g, y)))))) \Rightarrow  \\ \quad \quad (a = \operatorname{encr}(g, y)) \vee \operatorname{is\text{-}subterm}(t, k, a, g) \end{array}}}
\vnbprop{token-subterm-of-list}{\vnbstmt{\begin{array}{@{}l@{}} \forall t, k, a, l. \\ \quad \operatorname{is\text{-}token\text{-}set}(t) \wedge \\ \quad k \in \mathbf{Set} \wedge \\ \quad a \in t \wedge \\ \quad l \in \operatorname{tuples}(\operatorname{msg\text{-}set}(t, k)) \wedge \\ \quad \operatorname{is\text{-}subterm}(t, k, a, \operatorname{lst}(l)) \Rightarrow \\ \quad \quad \exists i \in \mathbb{N}.\; ((1 \leq i) \wedge ((i \leq \operatorname{length}(l)) \wedge \operatorname{is\text{-}subterm}(t, k, a, \operatorname{nth}(i, l)))) \end{array}}}
\vnbprop{token-subterm-of-encr}{\vnbstmt{\begin{array}{@{}l@{}} \forall t, k, a, g, y. \\ \quad \operatorname{is\text{-}token\text{-}set}(t) \wedge \\ \quad k \in \mathbf{Set} \wedge \\ \quad a \in t \wedge \\ \quad g \in \operatorname{msg\text{-}set}(t, k) \wedge \\ \quad y \in k \wedge \\ \quad \operatorname{is\text{-}subterm}(t, k, a, \operatorname{encr}(g, y)) \Rightarrow \\ \quad \quad \operatorname{is\text{-}subterm}(t, k, a, g) \end{array}}}
\vnbprop{subterm-encr-other-key}{\vnbstmt{\begin{array}{@{}l@{}} \forall t, k, h, a, y, b. \\ \quad t \in \mathbf{Set} \wedge \\ \quad k \in \mathbf{Set} \wedge \\ \quad h \in \operatorname{msg\text{-}set}(t, k) \wedge \\ \quad a \in \operatorname{msg\text{-}set}(t, k) \wedge \\ \quad y \in k \wedge \\ \quad b \in k \wedge \\ \quad \neg (y = b) \wedge \\ \quad \operatorname{is\text{-}subterm}(t, k, \operatorname{encr}(a, b), \operatorname{encr}(h, y)) \Rightarrow \\ \quad \quad \operatorname{is\text{-}subterm}(t, k, \operatorname{encr}(a, b), h) \end{array}}}
\vnbprop{msg-rank-prop}{\vnbstmt{\begin{array}{@{}l@{}} \forall t, k.; \forall m \in \operatorname{msg\text{-}set}(t, k). \\ \quad \operatorname{msg\text{-}rank}(t, k, m) \in \mathbb{N} \wedge \\ \quad m \in \operatorname{msg\text{-}stage}(t, k, \operatorname{msg\text{-}rank}(t, k, m)) \wedge \\ \quad \forall j \in \mathbb{N}.\; ((m \in \operatorname{msg\text{-}stage}(t, k, j)) \Rightarrow (\operatorname{msg\text{-}rank}(t, k, m) \leq j)) \end{array}}}
\vnbprop{msg-stage-list-inv}{\vnbstmt{\begin{array}{@{}l@{}} \forall n \in \mathbb{N}, t, k.; \forall l \in \mathbf{Set}. \\ \quad (\operatorname{is\text{-}token\text{-}set}(t) \wedge (\operatorname{lst}(l) \in \operatorname{msg\text{-}stage}(t, k, n))) \Rightarrow  \\ \quad \quad \exists j \in \mathbb{N}.\; (((j + 1) \leq n) \wedge (l \in \operatorname{tuples}(\operatorname{msg\text{-}stage}(t, k, j)))) \end{array}}}
\vnbprop{msg-stage-encr-inv}{\vnbstmt{\begin{array}{@{}l@{}} \forall n \in \mathbb{N}, t, k, g, y. \\ \quad (\operatorname{is\text{-}token\text{-}set}(t) \wedge ((g \in \mathbf{Set}) \wedge ((y \in \mathbf{Set}) \wedge (\operatorname{encr}(g, y) \in \operatorname{msg\text{-}stage}(t, k, n))))) \Rightarrow  \\ \quad \quad \exists j \in \mathbb{N}.\; (((j + 1) \leq n) \wedge (g \in \operatorname{msg\text{-}stage}(t, k, j))) \end{array}}}
\vnbprop{msg-rank-token}{\vnbstmt{\forall t, k.; \forall a \in t.\; (\operatorname{msg\text{-}rank}(t, k, a) = 0)}}
\vnbprop{msg-rank-component}{\vnbstmt{\begin{array}{@{}l@{}} \forall t, k, l, i. \\ \quad (\operatorname{is\text{-}token\text{-}set}(t) \wedge ((l \in \operatorname{tuples}(\operatorname{msg\text{-}set}(t, k))) \wedge ((i \in \mathbb{N}) \wedge ((1 \leq i) \wedge (i \leq \operatorname{length}(l)))))) \Rightarrow  \\ \quad \quad (\operatorname{msg\text{-}rank}(t, k, \operatorname{nth}(i, l)) + 1) \leq \operatorname{msg\text{-}rank}(t, k, \operatorname{lst}(l)) \end{array}}}
\vnbprop{msg-rank-encr}{\vnbstmt{\begin{array}{@{}l@{}} \forall t, k, g, y. \\ \quad (\operatorname{is\text{-}token\text{-}set}(t) \wedge ((g \in \operatorname{msg\text{-}set}(t, k)) \wedge (y \in k))) \Rightarrow  \\ \quad \quad (\operatorname{msg\text{-}rank}(t, k, g) + 1) \leq \operatorname{msg\text{-}rank}(t, k, \operatorname{encr}(g, y)) \end{array}}}
\vnbprop{subterm-rank}{\vnbstmt{\begin{array}{@{}l@{}} \forall t, k, a, b. \\ \quad (\operatorname{is\text{-}token\text{-}set}(t) \wedge ((k \in \mathbf{Set}) \wedge \operatorname{is\text{-}subterm}(t, k, a, b))) \Rightarrow  \\ \quad \quad (a = b) \vee ((\operatorname{msg\text{-}rank}(t, k, a) + 1) \leq \operatorname{msg\text{-}rank}(t, k, b)) \end{array}}}
\vnbprop{subterm-antisym}{\vnbstmt{\begin{array}{@{}l@{}} \forall t, k, a, b. \\ \quad (\operatorname{is\text{-}token\text{-}set}(t) \wedge ((k \in \mathbf{Set}) \wedge (\operatorname{is\text{-}subterm}(t, k, a, b) \wedge \operatorname{is\text{-}subterm}(t, k, b, a)))) \Rightarrow  \\ \quad \quad a = b \end{array}}}

\vnbbookfile{structure-library/strand-spaces}{structure-library/strand-spaces}
\noindent{\small strand-spaces.scm -- STRAND SPACES, NODES, BUNDLES ON THE USER'S MODEL (\textasciitilde{}/notes/notes-60.txt items (b)-(g); the user's decisions of 2026-10-06 recorded in scratchpad/triage/WEEK3.md, section S-2; the paper's sections 2.1-2.4 supply the THEOREMS, not the definitions: Thayer, Herzog, Guttman, ...}\par
\vnbstructdef{bundle}{A bundle $s$ is a 3-tuple whose components are listed below.\par
\begin{enumerate}[nosep]
\item $\operatorname{bspace}(s) = \operatorname{nth}(1, s)$, a substructure (\texttt{substructure}): $\operatorname{is\text{-}strand\text{-}space}(\operatorname{bspace}(s))$.
\item $\operatorname{bnodes}(s) = \operatorname{nth}(2, s)$, a carrier (\texttt{carriers}): $\operatorname{bnodes}(s) \in \mathbf{Set}$.
\item $\operatorname{bsend}(s) = \operatorname{nth}(3, s)$, a relation (\texttt{relation}): $\operatorname{bsend}(s) \in \operatorname{power}((\operatorname{bnodes}(s) \times \operatorname{bnodes}(s)))$.
\end{enumerate}
The laws (the \texttt{property} and \texttt{law} clauses):
\begin{itemize}[nosep]
\item $\forall \operatorname{ssx\_} \in \operatorname{bnodes}(s).\; (\operatorname{ssx\_} \in \operatorname{nodes}(\operatorname{bspace}(s)))$
\item $|\operatorname{bnodes}(s)| \in \mathbb{N}$
\item $\forall \operatorname{ssx\_}, \operatorname{ssy\_}.\; ((\langle \operatorname{ssx\_}, \operatorname{ssy\_} \rangle \in \operatorname{bsend}(s)) \Rightarrow \operatorname{sends\text{-}to}(\operatorname{bspace}(s), \operatorname{ssx\_}, \operatorname{ssy\_}))$
\item
\begin{equation*}
\begin{array}{@{}l@{}} \forall \operatorname{ssy\_} \in \operatorname{bnodes}(s). \\ \quad \operatorname{is\text{-}received}(\operatorname{bspace}(s), \operatorname{ssy\_}) \Rightarrow  \\ \quad \quad \exists \operatorname{ssx\_}. \\ \quad \quad \quad ((\langle \operatorname{ssx\_}, \operatorname{ssy\_} \rangle \in \operatorname{bsend}(s)) \wedge \forall \operatorname{ssz\_}.\; ((\langle \operatorname{ssz\_}, \operatorname{ssy\_} \rangle \in \operatorname{bsend}(s)) \Rightarrow (\operatorname{ssz\_} = \operatorname{ssx\_}))) \end{array}
\end{equation*}

\item
\begin{equation*}
\begin{array}{@{}l@{}} \forall \operatorname{ssy\_} \in \operatorname{bnodes}(s), \operatorname{ssx\_}. \\ \quad \operatorname{same\text{-}strand\text{-}before}(\operatorname{bspace}(s), \operatorname{ssx\_}, \operatorname{ssy\_}) \Rightarrow (\operatorname{ssx\_} \in \operatorname{bnodes}(s)) \end{array}
\end{equation*}

\item $\forall \operatorname{ssx\_}.\; \neg (\langle \operatorname{ssx\_}, \operatorname{ssx\_} \rangle \in \operatorname{prec\text{-}strict}(\operatorname{bspace}(s), \operatorname{bnodes}(s), \operatorname{bsend}(s)))$
\end{itemize}
The characteristic law (the axiom \texttt{is-bundle}): $\operatorname{is\text{-}bundle}(s)$ holds if and only if $\operatorname{length}(s) = 3$, every component has the typing listed with it, and every law holds.
}
\vnbstructdef{strand-space}{A strand space $s$ is a 5-tuple whose components are listed below.\par
\begin{enumerate}[nosep]
\item $\operatorname{tokens}(s) = \operatorname{nth}(1, s)$, a carrier (\texttt{carriers}): $\operatorname{tokens}(s) \in \mathbf{Set}$.
\item $\operatorname{ids}(s) = \operatorname{nth}(2, s)$, a carrier (\texttt{carriers}): $\operatorname{ids}(s) \in \mathbf{Set}$.
\item $\operatorname{keys}(s) = \operatorname{nth}(3, s)$, a carrier (\texttt{carriers}): $\operatorname{keys}(s) \in \mathbf{Set}$.
\item $\operatorname{strand\text{-}map}(s) = \operatorname{nth}(4, s)$, an operation (\texttt{op}): $\operatorname{strand\text{-}map}(s) \in (\operatorname{ids}(s) \to \operatorname{strands}(\operatorname{tokens}(s), \operatorname{keys}(s)))$.
\item $\operatorname{kinv}(s) = \operatorname{nth}(5, s)$, an operation (\texttt{op}): $\operatorname{kinv}(s) \in (\operatorname{keys}(s) \to \operatorname{keys}(s))$.
\end{enumerate}
The laws (the \texttt{property} and \texttt{law} clauses):
\begin{itemize}[nosep]
\item $\operatorname{is\text{-}token\text{-}set}(\operatorname{tokens}(s))$
\item $\forall \operatorname{ssx\_} \in \operatorname{keys}(s).\; (\operatorname{ssx\_} \in \operatorname{tokens}(s))$
\item $\forall \operatorname{ssx\_} \in \operatorname{keys}(s).\; (\operatorname{kinv}(s)(\operatorname{kinv}(s)(\operatorname{ssx\_})) = \operatorname{ssx\_})$
\end{itemize}
The characteristic law (the axiom \texttt{is-strand-space}): $\operatorname{is\text{-}strand\text{-}space}(s)$ holds if and only if $\operatorname{length}(s) = 5$, every component has the typing listed with it, and every law holds.
}
\par\medskip\noindent\textbf{Definition} (\texttt{nodes}).\;$\operatorname{nodes}(\operatorname{ssp\_}) :=$
\begin{equation*}
\{\operatorname{ssn\_} \in (\mathbb{N} \times \operatorname{ids}(\operatorname{ssp\_})) : ((1 \leq \operatorname{nth}(1, \operatorname{ssn\_})) \wedge (\operatorname{nth}(1, \operatorname{ssn\_}) \leq \operatorname{length}(\operatorname{strand\text{-}map}(\operatorname{ssp\_})(\operatorname{nth}(2, \operatorname{ssn\_})))))\}
\end{equation*}

\par\medskip\noindent\textbf{Definition} (\texttt{directed}).\;$\operatorname{directed}(\operatorname{sst\_}, \operatorname{ssk\_}) :=$ $(\{\mathit{sent}, \mathit{received}\} \times \operatorname{msg\text{-}set}(\operatorname{sst\_}, \operatorname{ssk\_}))$
\par\medskip\noindent\textbf{Definition} (\texttt{sent}).\; $\mathit{sent} = 0$
\par\medskip\noindent\textbf{Definition} (\texttt{received}).\; $\mathit{received} = 1$
\par\medskip\noindent\textbf{Definition} (\texttt{dir-of}).\;$\operatorname{dir\text{-}of}(\operatorname{ssd\_}) :=$ $\operatorname{nth}(1, \operatorname{ssd\_})$
\par\medskip\noindent\textbf{Definition} (\texttt{content-of}).\;$\operatorname{content\text{-}of}(\operatorname{ssd\_}) :=$ $\operatorname{nth}(2, \operatorname{ssd\_})$
\par\medskip\noindent\textbf{Definition} (\texttt{strands}).\;$\operatorname{strands}(\operatorname{sst\_}, \operatorname{ssk\_}) :=$ $\operatorname{tuples}(\operatorname{directed}(\operatorname{sst\_}, \operatorname{ssk\_}))$
\par\medskip\noindent\textbf{Definition} (\texttt{index-of}).\;$\operatorname{index\text{-}of}(\operatorname{ssn\_}) :=$ $\operatorname{nth}(1, \operatorname{ssn\_})$
\par\medskip\noindent\textbf{Definition} (\texttt{id-of}).\;$\operatorname{id\text{-}of}(\operatorname{ssn\_}) :=$ $\operatorname{nth}(2, \operatorname{ssn\_})$
\par\medskip\noindent\textbf{Definition} (\texttt{node-content}).\;$\operatorname{node\text{-}content}(\operatorname{ssp\_}, \operatorname{ssn\_}) :=$ $\operatorname{nth}(2, \operatorname{nth}(\operatorname{nth}(1, \operatorname{ssn\_}), \operatorname{strand\text{-}map}(\operatorname{ssp\_})(\operatorname{nth}(2, \operatorname{ssn\_}))))$
\par\medskip\noindent\textbf{Definition} (\texttt{node-dir}).\;$\operatorname{node\text{-}dir}(\operatorname{ssp\_}, \operatorname{ssn\_}) :=$ $\operatorname{nth}(1, \operatorname{nth}(\operatorname{nth}(1, \operatorname{ssn\_}), \operatorname{strand\text{-}map}(\operatorname{ssp\_})(\operatorname{nth}(2, \operatorname{ssn\_}))))$
\par\medskip\noindent\textbf{Definition} (\texttt{is-sent}).\;For all $\operatorname{ssp\_}$ and $\operatorname{ssn\_}$,  $\operatorname{is\text{-}sent}(\operatorname{ssp\_}, \operatorname{ssn\_}) \Leftrightarrow (\operatorname{node\text{-}dir}(\operatorname{ssp\_}, \operatorname{ssn\_}) = \mathit{sent})$
\par\medskip\noindent\textbf{Definition} (\texttt{is-received}).\;For all $\operatorname{ssp\_}$ and $\operatorname{ssn\_}$,  $\operatorname{is\text{-}received}(\operatorname{ssp\_}, \operatorname{ssn\_}) \Leftrightarrow (\operatorname{node\text{-}dir}(\operatorname{ssp\_}, \operatorname{ssn\_}) = \mathit{received})$
\par\medskip\noindent\textbf{Definition} (\texttt{same-strand-before}).\;For all $\operatorname{ssp\_}$, $\operatorname{ssn\_}$ and $\operatorname{ssm\_}$, 
\begin{equation*}
\begin{array}{@{}l@{}} \operatorname{same\text{-}strand\text{-}before}(\operatorname{ssp\_}, \operatorname{ssn\_}, \operatorname{ssm\_}) \\ \Leftrightarrow\; \operatorname{ssn\_} \in \operatorname{nodes}(\operatorname{ssp\_}) \wedge \\ \operatorname{ssm\_} \in \operatorname{nodes}(\operatorname{ssp\_}) \wedge \\ \operatorname{nth}(2, \operatorname{ssn\_}) = \operatorname{nth}(2, \operatorname{ssm\_}) \wedge \\ \operatorname{nth}(1, \operatorname{ssn\_}) < \operatorname{nth}(1, \operatorname{ssm\_}) \end{array}
\end{equation*}

\par\medskip\noindent\textbf{Definition} (\texttt{same-strand-upto}).\;For all $\operatorname{ssp\_}$, $\operatorname{ssn\_}$ and $\operatorname{ssm\_}$, 
\begin{equation*}
\begin{array}{@{}l@{}} \operatorname{same\text{-}strand\text{-}upto}(\operatorname{ssp\_}, \operatorname{ssn\_}, \operatorname{ssm\_}) \\ \Leftrightarrow\; \operatorname{ssn\_} \in \operatorname{nodes}(\operatorname{ssp\_}) \wedge \\ \operatorname{ssm\_} \in \operatorname{nodes}(\operatorname{ssp\_}) \wedge \\ \operatorname{nth}(2, \operatorname{ssn\_}) = \operatorname{nth}(2, \operatorname{ssm\_}) \wedge \\ \operatorname{nth}(1, \operatorname{ssn\_}) \leq \operatorname{nth}(1, \operatorname{ssm\_}) \end{array}
\end{equation*}

\par\medskip\noindent\textbf{Definition} (\texttt{sends-to}).\;For all $\operatorname{ssp\_}$, $\operatorname{ssn\_}$ and $\operatorname{ssm\_}$, 
\begin{equation*}
\begin{array}{@{}l@{}} \operatorname{sends\text{-}to}(\operatorname{ssp\_}, \operatorname{ssn\_}, \operatorname{ssm\_}) \\ \Leftrightarrow\; \operatorname{ssn\_} \in \operatorname{nodes}(\operatorname{ssp\_}) \wedge \\ \operatorname{ssm\_} \in \operatorname{nodes}(\operatorname{ssp\_}) \wedge \\ \operatorname{is\text{-}sent}(\operatorname{ssp\_}, \operatorname{ssn\_}) \wedge \\ \operatorname{is\text{-}received}(\operatorname{ssp\_}, \operatorname{ssm\_}) \wedge \\ \operatorname{node\text{-}content}(\operatorname{ssp\_}, \operatorname{ssn\_}) = \operatorname{node\text{-}content}(\operatorname{ssp\_}, \operatorname{ssm\_}) \end{array}
\end{equation*}

\par\medskip\noindent\textbf{Definition} (\texttt{occurs-in-space}).\;For all $\operatorname{ssp\_}$ and $\operatorname{ssu\_}$, 
\begin{equation*}
\begin{array}{@{}l@{}} \operatorname{occurs\text{-}in\text{-}space}(\operatorname{ssp\_}, \operatorname{ssu\_}) \\ \Leftrightarrow\; \exists \operatorname{ssn\_} \in \operatorname{nodes}(\operatorname{ssp\_}). \\ \quad \operatorname{is\text{-}subterm}(\operatorname{tokens}(\operatorname{ssp\_}), \operatorname{keys}(\operatorname{ssp\_}), \operatorname{ssu\_}, \operatorname{node\text{-}content}(\operatorname{ssp\_}, \operatorname{ssn\_})) \end{array}
\end{equation*}

\par\medskip\noindent\textbf{Definition} (\texttt{is-entry-point}).\;For all $\operatorname{ssp\_}$, $\operatorname{ssn\_}$ and $\operatorname{ssi\_}$, 
\begin{equation*}
\begin{array}{@{}l@{}} \operatorname{is\text{-}entry\text{-}point}(\operatorname{ssp\_}, \operatorname{ssn\_}, \operatorname{ssi\_}) \\ \Leftrightarrow\; \operatorname{ssn\_} \in \operatorname{nodes}(\operatorname{ssp\_}) \wedge \\ \operatorname{is\text{-}sent}(\operatorname{ssp\_}, \operatorname{ssn\_}) \wedge \\ \operatorname{node\text{-}content}(\operatorname{ssp\_}, \operatorname{ssn\_}) \in \operatorname{ssi\_} \wedge \\ \forall \operatorname{ssm\_}. \\ \quad \operatorname{same\text{-}strand\text{-}before}(\operatorname{ssp\_}, \operatorname{ssm\_}, \operatorname{ssn\_}) \Rightarrow  \\ \quad \quad \neg (\operatorname{node\text{-}content}(\operatorname{ssp\_}, \operatorname{ssm\_}) \in \operatorname{ssi\_}) \end{array}
\end{equation*}

\par\medskip\noindent\textbf{Definition} (\texttt{originates}).\;For all $\operatorname{ssp\_}$, $\operatorname{ssu\_}$ and $\operatorname{ssn\_}$, 
\begin{equation*}
\begin{array}{@{}l@{}} \operatorname{originates}(\operatorname{ssp\_}, \operatorname{ssu\_}, \operatorname{ssn\_}) \\ \Leftrightarrow\; \operatorname{is\text{-}entry\text{-}point}(\operatorname{ssp\_}, \operatorname{ssn\_}, \\ \quad \{\operatorname{ssw\_} \in \operatorname{msg\text{-}set}(\operatorname{tokens}(\operatorname{ssp\_}), \operatorname{keys}(\operatorname{ssp\_})) : \operatorname{is\text{-}subterm}(\operatorname{tokens}(\operatorname{ssp\_}), \operatorname{keys}(\operatorname{ssp\_}), \operatorname{ssu\_}, \operatorname{ssw\_})\}) \end{array}
\end{equation*}

\par\medskip\noindent\textbf{Definition} (\texttt{uniquely-originates}).\;For all $\operatorname{ssp\_}$ and $\operatorname{ssu\_}$, 
\begin{equation*}
\begin{array}{@{}l@{}} \operatorname{uniquely\text{-}originates}(\operatorname{ssp\_}, \operatorname{ssu\_}) \\ \Leftrightarrow\; \exists \operatorname{ssn\_}. \\ \quad (\operatorname{originates}(\operatorname{ssp\_}, \operatorname{ssu\_}, \operatorname{ssn\_}) \wedge \forall \operatorname{ssm\_}.\; (\operatorname{originates}(\operatorname{ssp\_}, \operatorname{ssu\_}, \operatorname{ssm\_}) \Rightarrow (\operatorname{ssm\_} = \operatorname{ssn\_}))) \end{array}
\end{equation*}

\par\medskip\noindent\textbf{Definition} (\texttt{is-prec-closed}).\;For all $\operatorname{spc\_}$, $\operatorname{snd\_}$, $\operatorname{ssd\_}$ and $\operatorname{srl\_}$, 
\begin{equation*}
\begin{array}{@{}l@{}} \operatorname{is\text{-}prec\text{-}closed}(\operatorname{spc\_}, \operatorname{snd\_}, \operatorname{ssd\_}, \operatorname{srl\_}) \\ \Leftrightarrow\; \forall \operatorname{ssx\_}, \operatorname{ssy\_}.\; ((\langle \operatorname{ssx\_}, \operatorname{ssy\_} \rangle \in \operatorname{ssd\_}) \Rightarrow (\langle \operatorname{ssx\_}, \operatorname{ssy\_} \rangle \in \operatorname{srl\_})) \wedge \\ \forall \operatorname{ssx\_}, \operatorname{ssy\_} \in \operatorname{snd\_}. \\ \quad \operatorname{same\text{-}strand\text{-}before}(\operatorname{spc\_}, \operatorname{ssx\_}, \operatorname{ssy\_}) \Rightarrow (\langle \operatorname{ssx\_}, \operatorname{ssy\_} \rangle \in \operatorname{srl\_}) \wedge \\ \forall \operatorname{ssx\_}, \operatorname{ssy\_}, \operatorname{ssz\_}. \\ \quad (\langle \operatorname{ssx\_}, \operatorname{ssy\_} \rangle \in \operatorname{srl\_}) \Rightarrow ((\langle \operatorname{ssy\_}, \operatorname{ssz\_} \rangle \in \operatorname{srl\_}) \Rightarrow (\langle \operatorname{ssx\_}, \operatorname{ssz\_} \rangle \in \operatorname{srl\_})) \end{array}
\end{equation*}

\par\medskip\noindent\textbf{Definition} (\texttt{prec-strict}).\;$\operatorname{prec\text{-}strict}(\operatorname{spc\_}, \operatorname{snd\_}, \operatorname{ssd\_}) :=$
\begin{equation*}
\{\operatorname{ssp\_} \in (\operatorname{snd\_} \times \operatorname{snd\_}) : \forall \operatorname{ssq\_} \in \operatorname{power}((\operatorname{snd\_} \times \operatorname{snd\_})).\; (\operatorname{is\text{-}prec\text{-}closed}(\operatorname{spc\_}, \operatorname{snd\_}, \operatorname{ssd\_}, \operatorname{ssq\_}) \Rightarrow (\operatorname{ssp\_} \in \operatorname{ssq\_}))\}
\end{equation*}

\par\medskip\noindent\textbf{Definition} (\texttt{precedes-strict}).\;$\operatorname{precedes\text{-}strict}(\operatorname{ssc\_}) :=$ $\operatorname{prec\text{-}strict}(\operatorname{bspace}(\operatorname{ssc\_}), \operatorname{bnodes}(\operatorname{ssc\_}), \operatorname{bsend}(\operatorname{ssc\_}))$
\par\medskip\noindent\textbf{Definition} (\texttt{precedes}).\;$\operatorname{precedes}(\operatorname{ssc\_}) :=$
\begin{equation*}
(\operatorname{precedes\text{-}strict}(\operatorname{ssc\_}) \cup \{\operatorname{ssp\_} \in (\operatorname{bnodes}(\operatorname{ssc\_}) \times \operatorname{bnodes}(\operatorname{ssc\_})) : (\operatorname{nth}(1, \operatorname{ssp\_}) = \operatorname{nth}(2, \operatorname{ssp\_}))\})
\end{equation*}

\par\medskip\noindent\textbf{Definition} (\texttt{hom-bundle}).\;$\operatorname{hom\text{-}bundle}(a, b) :=$
\begin{equation*}
\{\operatorname{homf\_} \in (\operatorname{bnodes}(a) \to \operatorname{bnodes}(b)) : \operatorname{is\text{-}hom\text{-}bundle}(a, b, \operatorname{homf\_})\}
\end{equation*}

\vnbdef{strand-space-class}{\vnbstmt{\forall s.\; ((s \in \operatorname{strand\text{-}space}) \Leftrightarrow \operatorname{is\text{-}strand\text{-}space}(s))}}
\vnbdef{is-hom-strand-space-def}{\vnbstmt{\begin{array}{@{}l@{}} \forall a, b, c, d, f. \\ \quad \operatorname{is\text{-}hom\text{-}strand\text{-}space}(a, b, c, d, f) \\ \quad \Leftrightarrow\; \operatorname{is\text{-}strand\text{-}space}(a) \wedge \\ \quad \operatorname{is\text{-}strand\text{-}space}(b) \wedge \\ \quad c \in (\operatorname{tokens}(a) \to \operatorname{tokens}(b)) \wedge \\ \quad d \in (\operatorname{ids}(a) \to \operatorname{ids}(b)) \wedge \\ \quad f \in (\operatorname{keys}(a) \to \operatorname{keys}(b)) \wedge \\ \quad \forall x_{1} \in \operatorname{ids}(a).\; (\operatorname{strand\text{-}map}(a)(x_{1}) = \operatorname{strand\text{-}map}(b)(d(x_{1}))) \wedge \\ \quad \forall x_{1} \in \operatorname{keys}(a).\; (f(\operatorname{kinv}(a)(x_{1})) = \operatorname{kinv}(b)(f(x_{1}))) \end{array}}}
\vnbdef{is-hom-bundle-def}{\vnbstmt{\begin{array}{@{}l@{}} \forall a, b, f. \\ \quad \operatorname{is\text{-}hom\text{-}bundle}(a, b, f) \\ \quad \Leftrightarrow\; \operatorname{is\text{-}bundle}(a) \wedge \\ \quad \operatorname{is\text{-}bundle}(b) \wedge \\ \quad f \in (\operatorname{bnodes}(a) \to \operatorname{bnodes}(b)) \wedge \\ \quad \operatorname{bspace}(a) = \operatorname{bspace}(b) \wedge \\ \quad \forall x_{1}, x_{2} \in \operatorname{bnodes}(a). \\ \quad \quad (\langle x_{1}, x_{2} \rangle \in \operatorname{bsend}(a)) \Rightarrow (\langle f(x_{1}), f(x_{2}) \rangle \in \operatorname{bsend}(b)) \end{array}}}
\vnbdef{bundle-class}{\vnbstmt{\forall s.\; ((s \in \mathit{bundle}) \Leftrightarrow \operatorname{is\text{-}bundle}(s))}}

\vnbbookfile{theorem-library/strand-space-laws}{theorem-library/strand-space-laws}
\noindent{\small strand-space-laws.scm -- STRAND SPACES, NODES, BUNDLES, CAUSAL PRECEDENCE: the laws, and the paper's Lemmas 2.7, 2.8, 2.9 on the user's model. (\textasciitilde{}/notes/notes-60.txt; the definitions are structure-library/strand-spaces.scm; the paper is Thayer, Herzog, Guttman, \textasciitilde{}/docs/jcs\_strand\_spaces.pdf, sections 2.2-2.4.)}\par
\vnbprop{nodes-unfold}{\vnbstmt{\begin{array}{@{}l@{}} \forall a. \\ \quad (\operatorname{nodes}(a) \simeq \{n \in (\mathbb{N} \times \operatorname{ids}(a)) : ((1 \leq \operatorname{nth}(1, n)) \wedge (\operatorname{nth}(1, n) \leq \operatorname{length}(\operatorname{strand\text{-}map}(a)(\operatorname{nth}(2, n)))))\}) \end{array}}}
\vnbprop{node-content-unfold}{\vnbstmt{\begin{array}{@{}l@{}} \forall a, n. \\ \quad (\operatorname{node\text{-}content}(a, n) \simeq \operatorname{nth}(2, \operatorname{nth}(\operatorname{nth}(1, n), \operatorname{strand\text{-}map}(a)(\operatorname{nth}(2, n))))) \end{array}}}
\vnbprop{node-dir-unfold}{\vnbstmt{\begin{array}{@{}l@{}} \forall a, n. \\ \quad (\operatorname{node\text{-}dir}(a, n) \simeq \operatorname{nth}(1, \operatorname{nth}(\operatorname{nth}(1, n), \operatorname{strand\text{-}map}(a)(\operatorname{nth}(2, n))))) \end{array}}}
\vnbprop{directed-unfold}{\vnbstmt{\forall t, k.\; (\operatorname{directed}(t, k) \simeq (\{\mathit{sent}, \mathit{received}\} \times \operatorname{msg\text{-}set}(t, k)))}}
\vnbprop{strands-unfold}{\vnbstmt{\forall t, k.\; (\operatorname{strands}(t, k) \simeq \operatorname{tuples}(\operatorname{directed}(t, k)))}}
\vnbprop{index-id-unfold}{\vnbstmt{\forall n.\; ((\operatorname{index\text{-}of}(n) \simeq \operatorname{nth}(1, n)) \wedge (\operatorname{id\text{-}of}(n) \simeq \operatorname{nth}(2, n)))}}
\vnbprop{prec-strict-unfold}{\vnbstmt{\begin{array}{@{}l@{}} \forall a, d, b. \\ \quad (\operatorname{prec\text{-}strict}(a, d, b) \simeq \{p \in (d \times d) : \forall q \in \operatorname{power}((d \times d)).\; (\operatorname{is\text{-}prec\text{-}closed}(a, d, b, q) \Rightarrow (p \in q))\}) \end{array}}}
\vnbprop{precedes-strict-unfold}{\vnbstmt{\forall c.\; (\operatorname{precedes\text{-}strict}(c) \simeq \operatorname{prec\text{-}strict}(\operatorname{bspace}(c), \operatorname{bnodes}(c), \operatorname{bsend}(c)))}}
\vnbprop{precedes-unfold}{\vnbstmt{\begin{array}{@{}l@{}} \forall c. \\ \quad (\operatorname{precedes}(c) \simeq (\operatorname{precedes\text{-}strict}(c) \cup \{p \in (\operatorname{bnodes}(c) \times \operatorname{bnodes}(c)) : (\operatorname{nth}(1, p) = \operatorname{nth}(2, p))\})) \end{array}}}
\vnbprop{strand-space-parts}{\vnbstmt{\begin{array}{@{}l@{}} \forall a. \\ \quad \operatorname{is\text{-}strand\text{-}space}(a) \Rightarrow  \\ \quad \quad \operatorname{is\text{-}token\text{-}set}(\operatorname{tokens}(a)) \wedge \\ \quad \quad \operatorname{tokens}(a) \in \mathbf{Set} \wedge \\ \quad \quad \operatorname{ids}(a) \in \mathbf{Set} \wedge \\ \quad \quad \operatorname{keys}(a) \in \mathbf{Set} \wedge \\ \quad \quad \operatorname{strand\text{-}map}(a) \in (\operatorname{ids}(a) \to \operatorname{strands}(\operatorname{tokens}(a), \operatorname{keys}(a))) \end{array}}}
\vnbprop{keys-in-tokens}{\vnbstmt{\forall a.; \forall x \in \operatorname{keys}(a).\; (\operatorname{is\text{-}strand\text{-}space}(a) \Rightarrow (x \in \operatorname{tokens}(a)))}}
\vnbprop{kinv-in-fun}{\vnbstmt{\forall a.\; (\operatorname{is\text{-}strand\text{-}space}(a) \Rightarrow (\operatorname{kinv}(a) \in (\operatorname{keys}(a) \to \operatorname{keys}(a))))}}
\vnbprop{kinv-involution}{\vnbstmt{\begin{array}{@{}l@{}} \forall a.; \forall y \in \operatorname{keys}(a). \\ \quad \operatorname{is\text{-}strand\text{-}space}(a) \Rightarrow (\operatorname{kinv}(a)(\operatorname{kinv}(a)(y)) = y) \end{array}}}
\vnbprop{kinv-injective}{\vnbstmt{\begin{array}{@{}l@{}} \forall a, y, z. \\ \quad (\operatorname{is\text{-}strand\text{-}space}(a) \wedge ((y \in \operatorname{keys}(a)) \wedge ((z \in \operatorname{keys}(a)) \wedge (\operatorname{kinv}(a)(y) = \operatorname{kinv}(a)(z))))) \Rightarrow  \\ \quad \quad y = z \end{array}}}
\vnbprop{strand-at-id}{\vnbstmt{\begin{array}{@{}l@{}} \forall b, a. \\ \quad (\operatorname{is\text{-}strand\text{-}space}(b) \wedge (a \in \operatorname{ids}(b))) \Rightarrow  \\ \quad \quad \operatorname{strand\text{-}map}(b)(a) \in \operatorname{tuples}(\operatorname{directed}(\operatorname{tokens}(b), \operatorname{keys}(b))) \end{array}}}
\vnbprop{node-parts}{\vnbstmt{\begin{array}{@{}l@{}} \forall a.; \forall n \in \operatorname{nodes}(a). \\ \quad \operatorname{nth}(1, n) \in \mathbb{N} \wedge \\ \quad \operatorname{nth}(2, n) \in \operatorname{ids}(a) \wedge \\ \quad 1 \leq \operatorname{nth}(1, n) \wedge \\ \quad \operatorname{nth}(1, n) \leq \operatorname{length}(\operatorname{strand\text{-}map}(a)(\operatorname{nth}(2, n))) \wedge \\ \quad n = \langle \operatorname{nth}(1, n), \operatorname{nth}(2, n) \rangle \end{array}}}
\vnbprop{node-intro}{\vnbstmt{\begin{array}{@{}l@{}} \forall b, k, a. \\ \quad ((k \in \mathbb{N}) \wedge ((a \in \operatorname{ids}(b)) \wedge ((1 \leq k) \wedge (k \leq \operatorname{length}(\operatorname{strand\text{-}map}(b)(a)))))) \Rightarrow  \\ \quad \quad \langle k, a \rangle \in \operatorname{nodes}(b) \end{array}}}
\vnbprop{node-pair-injective}{\vnbstmt{\forall k, a, j, b.\; ((\langle k, a \rangle = \langle j, b \rangle) \Rightarrow ((k = j) \wedge (a = b)))}}
\vnbprop{node-ext}{\vnbstmt{\begin{array}{@{}l@{}} \forall a, x, y. \\ \quad ((x \in \operatorname{nodes}(a)) \wedge ((y \in \operatorname{nodes}(a)) \wedge ((\operatorname{nth}(1, x) = \operatorname{nth}(1, y)) \wedge (\operatorname{nth}(2, x) = \operatorname{nth}(2, y))))) \Rightarrow  \\ \quad \quad x = y \end{array}}}
\vnbprop{directed-parts}{\vnbstmt{\begin{array}{@{}l@{}} \forall t, k.; \forall d \in \operatorname{directed}(t, k). \\ \quad (\operatorname{nth}(1, d) = \mathit{sent}) \vee (\operatorname{nth}(1, d) = \mathit{received}) \wedge \\ \quad \operatorname{nth}(2, d) \in \operatorname{msg\text{-}set}(t, k) \wedge \\ \quad d = \langle \operatorname{nth}(1, d), \operatorname{nth}(2, d) \rangle \end{array}}}
\vnbprop{node-step-in-directed}{\vnbstmt{\begin{array}{@{}l@{}} \forall a, n. \\ \quad (\operatorname{is\text{-}strand\text{-}space}(a) \wedge (n \in \operatorname{nodes}(a))) \Rightarrow  \\ \quad \quad \operatorname{nth}(\operatorname{nth}(1, n), \operatorname{strand\text{-}map}(a)(\operatorname{nth}(2, n))) \in \operatorname{directed}(\operatorname{tokens}(a), \operatorname{keys}(a)) \end{array}}}
\vnbprop{node-content-in-msgs}{\vnbstmt{\begin{array}{@{}l@{}} \forall a, n. \\ \quad (\operatorname{is\text{-}strand\text{-}space}(a) \wedge (n \in \operatorname{nodes}(a))) \Rightarrow  \\ \quad \quad \operatorname{node\text{-}content}(a, n) \in \operatorname{msg\text{-}set}(\operatorname{tokens}(a), \operatorname{keys}(a)) \end{array}}}
\vnbprop{node-dir-cases}{\vnbstmt{\begin{array}{@{}l@{}} \forall a, n. \\ \quad (\operatorname{is\text{-}strand\text{-}space}(a) \wedge (n \in \operatorname{nodes}(a))) \Rightarrow  \\ \quad \quad (\operatorname{node\text{-}dir}(a, n) = \mathit{sent}) \vee (\operatorname{node\text{-}dir}(a, n) = \mathit{received}) \end{array}}}
\vnbprop{sent-ne-received}{\vnbstmt{\neg (\mathit{sent} = \mathit{received})}}
\vnbprop{node-sent-or-received}{\vnbstmt{\begin{array}{@{}l@{}} \forall a, n. \\ \quad (\operatorname{is\text{-}strand\text{-}space}(a) \wedge (n \in \operatorname{nodes}(a))) \Rightarrow (\operatorname{is\text{-}sent}(a, n) \vee \operatorname{is\text{-}received}(a, n)) \end{array}}}
\vnbprop{node-not-sent-and-received}{\vnbstmt{\forall a, n.\; (\operatorname{is\text{-}sent}(a, n) \Rightarrow \neg \operatorname{is\text{-}received}(a, n))}}
\vnbprop{node-content-eq}{\vnbstmt{\begin{array}{@{}l@{}} \forall a, n, k, d. \\ \quad ((\operatorname{nth}(1, n) = k) \wedge (\operatorname{nth}(k, \operatorname{strand\text{-}map}(a)(\operatorname{nth}(2, n))) = d)) \Rightarrow  \\ \quad \quad (\operatorname{node\text{-}content}(a, n) \simeq \operatorname{nth}(2, d)) \wedge (\operatorname{node\text{-}dir}(a, n) \simeq \operatorname{nth}(1, d)) \end{array}}}
\vnbprop{ssb-irreflexive}{\vnbstmt{\forall a, n.\; \neg \operatorname{same\text{-}strand\text{-}before}(a, n, n)}}
\vnbprop{ssb-transitive}{\vnbstmt{\begin{array}{@{}l@{}} \forall a, x, y, z. \\ \quad (\operatorname{same\text{-}strand\text{-}before}(a, x, y) \wedge \operatorname{same\text{-}strand\text{-}before}(a, y, z)) \Rightarrow  \\ \quad \quad \operatorname{same\text{-}strand\text{-}before}(a, x, z) \end{array}}}
\vnbprop{ssu-of-ssb}{\vnbstmt{\forall a, x, y.\; (\operatorname{same\text{-}strand\text{-}before}(a, x, y) \Rightarrow \operatorname{same\text{-}strand\text{-}upto}(a, x, y))}}
\vnbprop{ssu-refl}{\vnbstmt{\forall a.; \forall x \in \operatorname{nodes}(a).\; \operatorname{same\text{-}strand\text{-}upto}(a, x, x)}}
\vnbprop{ssu-cases}{\vnbstmt{\begin{array}{@{}l@{}} \forall a, x, y. \\ \quad \operatorname{same\text{-}strand\text{-}upto}(a, x, y) \Rightarrow (\operatorname{same\text{-}strand\text{-}before}(a, x, y) \vee (x = y)) \end{array}}}
\vnbprop{sends-to-dirs}{\vnbstmt{\begin{array}{@{}l@{}} \forall a, x, y. \\ \quad \operatorname{sends\text{-}to}(a, x, y) \Rightarrow  \\ \quad \quad \operatorname{is\text{-}sent}(a, x) \wedge (\operatorname{is\text{-}received}(a, y) \wedge (\operatorname{node\text{-}content}(a, x) = \operatorname{node\text{-}content}(a, y))) \end{array}}}
\vnbprop{bundle-space}{\vnbstmt{\forall c.\; (\operatorname{is\text{-}bundle}(c) \Rightarrow \operatorname{is\text{-}strand\text{-}space}(\operatorname{bspace}(c)))}}
\vnbprop{bundle-nodes-set}{\vnbstmt{\forall c.\; (\operatorname{is\text{-}bundle}(c) \Rightarrow (\operatorname{bnodes}(c) \in \mathbf{Set}))}}
\vnbprop{bundle-send-typing}{\vnbstmt{\begin{array}{@{}l@{}} \forall c. \\ \quad \operatorname{is\text{-}bundle}(c) \Rightarrow (\operatorname{bsend}(c) \in \operatorname{power}((\operatorname{bnodes}(c) \times \operatorname{bnodes}(c)))) \end{array}}}
\vnbprop{bundle-nodes-in-space}{\vnbstmt{\forall c.; \forall x \in \operatorname{bnodes}(c).\; (\operatorname{is\text{-}bundle}(c) \Rightarrow (x \in \operatorname{nodes}(\operatorname{bspace}(c))))}}
\vnbprop{bundle-finite}{\vnbstmt{\forall c.\; (\operatorname{is\text{-}bundle}(c) \Rightarrow (|\operatorname{bnodes}(c)| \in \mathbb{N}))}}
\vnbprop{bundle-send-sends}{\vnbstmt{\begin{array}{@{}l@{}} \forall c, x, y. \\ \quad (\operatorname{is\text{-}bundle}(c) \wedge (\langle x, y \rangle \in \operatorname{bsend}(c))) \Rightarrow \operatorname{sends\text{-}to}(\operatorname{bspace}(c), x, y) \end{array}}}
\vnbprop{bundle-received-sender}{\vnbstmt{\begin{array}{@{}l@{}} \forall c.; \forall y \in \operatorname{bnodes}(c). \\ \quad (\operatorname{is\text{-}bundle}(c) \wedge \operatorname{is\text{-}received}(\operatorname{bspace}(c), y)) \Rightarrow  \\ \quad \quad \exists x. \\ \quad \quad \quad (\langle x, y \rangle \in \operatorname{bsend}(c)) \wedge \forall z.\; ((\langle z, y \rangle \in \operatorname{bsend}(c)) \Rightarrow (z = x)) \end{array}}}
\vnbprop{bundle-ssb-closed}{\vnbstmt{\begin{array}{@{}l@{}} \forall c.; \forall y \in \operatorname{bnodes}(c), x. \\ \quad (\operatorname{is\text{-}bundle}(c) \wedge \operatorname{same\text{-}strand\text{-}before}(\operatorname{bspace}(c), x, y)) \Rightarrow (x \in \operatorname{bnodes}(c)) \end{array}}}
\vnbprop{bundle-acyclic}{\vnbstmt{\begin{array}{@{}l@{}} \forall c, x. \\ \quad \operatorname{is\text{-}bundle}(c) \Rightarrow \neg (\langle x, x \rangle \in \operatorname{prec\text{-}strict}(\operatorname{bspace}(c), \operatorname{bnodes}(c), \operatorname{bsend}(c))) \end{array}}}
\vnbprop{pair-in-power-box}{\vnbstmt{\begin{array}{@{}l@{}} \forall r, d, x, y. \\ \quad ((r \in \operatorname{power}((d \times d))) \wedge (\langle x, y \rangle \in r)) \Rightarrow ((x \in d) \wedge (y \in d)) \end{array}}}
\vnbprop{bundle-send-nodes}{\vnbstmt{\begin{array}{@{}l@{}} \forall c, x, y. \\ \quad (\operatorname{is\text{-}bundle}(c) \wedge (\langle x, y \rangle \in \operatorname{bsend}(c))) \Rightarrow ((x \in \operatorname{bnodes}(c)) \wedge (y \in \operatorname{bnodes}(c))) \end{array}}}
\vnbprop{ps-in-box}{\vnbstmt{\begin{array}{@{}l@{}} \forall a, d, b, x, y. \\ \quad (\langle x, y \rangle \in \operatorname{prec\text{-}strict}(a, d, b)) \Rightarrow ((x \in d) \wedge (y \in d)) \end{array}}}
\vnbprop{ps-least}{\vnbstmt{\begin{array}{@{}l@{}} \forall a, d, b.; \forall r \in \operatorname{power}((d \times d)), x, y. \\ \quad (\operatorname{is\text{-}prec\text{-}closed}(a, d, b, r) \wedge (\langle x, y \rangle \in \operatorname{prec\text{-}strict}(a, d, b))) \Rightarrow  \\ \quad \quad \langle x, y \rangle \in r \end{array}}}
\vnbprop{ps-send}{\vnbstmt{\begin{array}{@{}l@{}} \forall a, d, b, x, y. \\ \quad ((b \in \operatorname{power}((d \times d))) \wedge (\langle x, y \rangle \in b)) \Rightarrow  \\ \quad \quad \langle x, y \rangle \in \operatorname{prec\text{-}strict}(a, d, b) \end{array}}}
\vnbprop{ps-ssb}{\vnbstmt{\begin{array}{@{}l@{}} \forall a, d, b, x, y. \\ \quad ((x \in d) \wedge ((y \in d) \wedge \operatorname{same\text{-}strand\text{-}before}(a, x, y))) \Rightarrow  \\ \quad \quad \langle x, y \rangle \in \operatorname{prec\text{-}strict}(a, d, b) \end{array}}}
\vnbprop{ps-trans}{\vnbstmt{\begin{array}{@{}l@{}} \forall a, d, b, x, y, z. \\ \quad ((\langle x, y \rangle \in \operatorname{prec\text{-}strict}(a, d, b)) \wedge (\langle y, z \rangle \in \operatorname{prec\text{-}strict}(a, d, b))) \Rightarrow  \\ \quad \quad \langle x, z \rangle \in \operatorname{prec\text{-}strict}(a, d, b) \end{array}}}
\vnbprop{precedes-strict-trans}{\vnbstmt{\begin{array}{@{}l@{}} \forall c, x, y, z. \\ \quad (\operatorname{is\text{-}bundle}(c) \wedge ((\langle x, y \rangle \in \operatorname{precedes\text{-}strict}(c)) \wedge (\langle y, z \rangle \in \operatorname{precedes\text{-}strict}(c)))) \Rightarrow  \\ \quad \quad \langle x, z \rangle \in \operatorname{precedes\text{-}strict}(c) \end{array}}}
\vnbprop{precedes-strict-irreflexive}{\vnbstmt{\forall c, x.\; (\operatorname{is\text{-}bundle}(c) \Rightarrow \neg (\langle x, x \rangle \in \operatorname{precedes\text{-}strict}(c)))}}
\vnbprop{bundle-send-precedes}{\vnbstmt{\begin{array}{@{}l@{}} \forall c, x, y. \\ \quad (\operatorname{is\text{-}bundle}(c) \wedge (\langle x, y \rangle \in \operatorname{bsend}(c))) \Rightarrow (\langle x, y \rangle \in \operatorname{precedes\text{-}strict}(c)) \end{array}}}
\vnbprop{bundle-ssb-precedes}{\vnbstmt{\begin{array}{@{}l@{}} \forall c, x, y. \\ \quad (\operatorname{is\text{-}bundle}(c) \wedge ((y \in \operatorname{bnodes}(c)) \wedge \operatorname{same\text{-}strand\text{-}before}(\operatorname{bspace}(c), x, y))) \Rightarrow  \\ \quad \quad \langle x, y \rangle \in \operatorname{precedes\text{-}strict}(c) \end{array}}}
\vnbprop{precedes-intro-strict}{\vnbstmt{\begin{array}{@{}l@{}} \forall c, x, y. \\ \quad (\operatorname{is\text{-}bundle}(c) \wedge (\langle x, y \rangle \in \operatorname{precedes\text{-}strict}(c))) \Rightarrow (\langle x, y \rangle \in \operatorname{precedes}(c)) \end{array}}}
\vnbprop{precedes-refl}{\vnbstmt{\forall c.; \forall x \in \operatorname{bnodes}(c).\; (\langle x, x \rangle \in \operatorname{precedes}(c))}}
\vnbprop{precedes-elim}{\vnbstmt{\begin{array}{@{}l@{}} \forall c, x, y. \\ \quad (\langle x, y \rangle \in \operatorname{precedes}(c)) \Rightarrow  \\ \quad \quad (\langle x, y \rangle \in \operatorname{precedes\text{-}strict}(c)) \vee ((x = y) \wedge (x \in \operatorname{bnodes}(c))) \end{array}}}
\vnbprop{ps-last-step}{\vnbstmt{\begin{array}{@{}l@{}} \forall a, d, b, x, y. \\ \quad ((d \in \mathbf{Set}) \wedge ((b \in \operatorname{power}((d \times d))) \wedge (\langle x, y \rangle \in \operatorname{prec\text{-}strict}(a, d, b)))) \Rightarrow  \\ \quad \quad \exists z. \\ \quad \quad \quad (\langle z, y \rangle \in b) \vee ((z \in d) \wedge \operatorname{same\text{-}strand\text{-}before}(a, z, y)) \wedge \\ \quad \quad \quad (z = x) \vee (\langle x, z \rangle \in \operatorname{prec\text{-}strict}(a, d, b)) \end{array}}}
\vnbprop{precedes-trans}{\vnbstmt{\begin{array}{@{}l@{}} \forall c, x, y, z. \\ \quad (\operatorname{is\text{-}bundle}(c) \wedge ((\langle x, y \rangle \in \operatorname{precedes}(c)) \wedge (\langle y, z \rangle \in \operatorname{precedes}(c)))) \Rightarrow  \\ \quad \quad \langle x, z \rangle \in \operatorname{precedes}(c) \end{array}}}
\vnbprop{precedes-antisym}{\vnbstmt{\begin{array}{@{}l@{}} \forall c, x, y. \\ \quad (\operatorname{is\text{-}bundle}(c) \wedge ((\langle x, y \rangle \in \operatorname{precedes}(c)) \wedge (\langle y, x \rangle \in \operatorname{precedes}(c)))) \Rightarrow (x = y) \end{array}}}
\vnbprop{precedes-set}{\vnbstmt{\forall c.\; (\operatorname{is\text{-}bundle}(c) \Rightarrow (\operatorname{precedes}(c) \in \mathbf{Set}))}}
\vnbprop{precedes-partial-order}{\vnbstmt{\forall c.\; (\operatorname{is\text{-}bundle}(c) \Rightarrow \operatorname{is\text{-}partial\text{-}order}(\operatorname{bnodes}(c), \operatorname{precedes}(c)))}}
\vnbprop{finite-strict-minimal}{\vnbstmt{\begin{array}{@{}l@{}} \forall r.; \forall s \in \mathbf{Set}. \\ \quad \forall x, y, z.\; (((\langle x, y \rangle \in r) \wedge (\langle y, z \rangle \in r)) \Rightarrow (\langle x, z \rangle \in r)) \wedge \\ \quad \forall x.\; \neg (\langle x, x \rangle \in r) \wedge \\ \quad |s| \in \mathbb{N} \wedge \\ \quad \exists w.\; (w \in s) \Rightarrow \\ \quad \quad \exists m \in s.\; \forall z \in s.\; \neg (\langle z, m \rangle \in r) \end{array}}}
\vnbprop{bundle-minimal}{\vnbstmt{\begin{array}{@{}l@{}} \forall c, s. \\ \quad (\operatorname{is\text{-}bundle}(c) \wedge (\forall z \in s.\; (z \in \operatorname{bnodes}(c)) \wedge \exists w.\; (w \in s))) \Rightarrow  \\ \quad \quad \exists m \in s.\; \forall z \in s.\; \neg (\langle z, m \rangle \in \operatorname{precedes\text{-}strict}(c)) \end{array}}}
\vnbprop{bundle-minimal-precedes}{\vnbstmt{\begin{array}{@{}l@{}} \forall c, s. \\ \quad (\operatorname{is\text{-}bundle}(c) \wedge (\forall z \in s.\; (z \in \operatorname{bnodes}(c)) \wedge \exists w.\; (w \in s))) \Rightarrow  \\ \quad \quad \exists m \in s. \forall z \in s. \\ \quad \quad \quad \neg (z = m) \Rightarrow \neg (\langle z, m \rangle \in \operatorname{precedes}(c)) \end{array}}}
\vnbprop{bundle-minimal-sent}{\vnbstmt{\begin{array}{@{}l@{}} \forall c, s, n. \\ \quad \operatorname{is\text{-}bundle}(c) \wedge \\ \quad \forall z \in s.\; (z \in \operatorname{bnodes}(c)) \wedge \\ \quad \forall x, y. \\ \quad \quad x \in \operatorname{bnodes}(c) \wedge \\ \quad \quad y \in \operatorname{bnodes}(c) \wedge \\ \quad \quad \operatorname{node\text{-}content}(\operatorname{bspace}(c), x) = \operatorname{node\text{-}content}(\operatorname{bspace}(c), y) \wedge \\ \quad \quad y \in s \Rightarrow \\ \quad \quad \quad x \in s \wedge \\ \quad n \in s \wedge \\ \quad \forall z \in s.\; \neg (\langle z, n \rangle \in \operatorname{precedes\text{-}strict}(c)) \Rightarrow \\ \quad \quad \operatorname{is\text{-}sent}(\operatorname{bspace}(c), n) \end{array}}}
\vnbprop{originates-elim}{\vnbstmt{\begin{array}{@{}l@{}} \forall a, u, n. \\ \quad (\operatorname{is\text{-}strand\text{-}space}(a) \wedge \operatorname{originates}(a, u, n)) \Rightarrow  \\ \quad \quad n \in \operatorname{nodes}(a) \wedge \\ \quad \quad \operatorname{is\text{-}sent}(a, n) \wedge \\ \quad \quad \operatorname{is\text{-}subterm}(\operatorname{tokens}(a), \operatorname{keys}(a), u, \operatorname{node\text{-}content}(a, n)) \wedge \\ \quad \quad \forall m. \\ \quad \quad \quad \operatorname{same\text{-}strand\text{-}before}(a, m, n) \Rightarrow  \\ \quad \quad \quad \quad \neg \operatorname{is\text{-}subterm}(\operatorname{tokens}(a), \operatorname{keys}(a), u, \operatorname{node\text{-}content}(a, m)) \end{array}}}
\vnbprop{originates-intro}{\vnbstmt{\begin{array}{@{}l@{}} \forall a, u, n. \\ \quad \operatorname{is\text{-}strand\text{-}space}(a) \wedge \\ \quad n \in \operatorname{nodes}(a) \wedge \\ \quad \operatorname{is\text{-}sent}(a, n) \wedge \\ \quad \operatorname{is\text{-}subterm}(\operatorname{tokens}(a), \operatorname{keys}(a), u, \operatorname{node\text{-}content}(a, n)) \wedge \\ \quad \forall m. \\ \quad \quad \operatorname{same\text{-}strand\text{-}before}(a, m, n) \Rightarrow  \\ \quad \quad \quad \neg \operatorname{is\text{-}subterm}(\operatorname{tokens}(a), \operatorname{keys}(a), u, \operatorname{node\text{-}content}(a, m)) \Rightarrow \\ \quad \quad \operatorname{originates}(a, u, n) \end{array}}}
\vnbprop{originates-sent}{\vnbstmt{\begin{array}{@{}l@{}} \forall a, u, n. \\ \quad (\operatorname{is\text{-}strand\text{-}space}(a) \wedge \operatorname{originates}(a, u, n)) \Rightarrow  \\ \quad \quad \operatorname{is\text{-}sent}(a, n) \wedge \operatorname{is\text{-}subterm}(\operatorname{tokens}(a), \operatorname{keys}(a), u, \operatorname{node\text{-}content}(a, n)) \end{array}}}
\vnbprop{originates-no-earlier}{\vnbstmt{\begin{array}{@{}l@{}} \forall a, u, n, m. \\ \quad (\operatorname{is\text{-}strand\text{-}space}(a) \wedge (\operatorname{originates}(a, u, n) \wedge \operatorname{same\text{-}strand\text{-}before}(a, m, n))) \Rightarrow  \\ \quad \quad \neg \operatorname{is\text{-}subterm}(\operatorname{tokens}(a), \operatorname{keys}(a), u, \operatorname{node\text{-}content}(a, m)) \end{array}}}
\vnbprop{originates-unique}{\vnbstmt{\begin{array}{@{}l@{}} \forall a, u, n, m. \\ \quad (\operatorname{uniquely\text{-}originates}(a, u) \wedge (\operatorname{originates}(a, u, n) \wedge \operatorname{originates}(a, u, m))) \Rightarrow  \\ \quad \quad m = n \end{array}}}
\vnbprop{bundle-minimal-originates}{\vnbstmt{\begin{array}{@{}l@{}} \forall c, u, n. \\ \quad \operatorname{is\text{-}bundle}(c) \wedge \\ \quad n \in \operatorname{bnodes}(c) \wedge \\ \quad \operatorname{is\text{-}subterm}(\operatorname{tokens}(\operatorname{bspace}(c)), \operatorname{keys}(\operatorname{bspace}(c)), u, \operatorname{node\text{-}content}(\operatorname{bspace}(c), n)) \wedge \\ \quad \forall z \in \operatorname{bnodes}(c). \\ \quad \quad \operatorname{is\text{-}subterm}(\operatorname{tokens}(\operatorname{bspace}(c)), \operatorname{keys}(\operatorname{bspace}(c)), u, \operatorname{node\text{-}content}(\operatorname{bspace}(c), z)) \Rightarrow  \\ \quad \quad \quad \neg (\langle z, n \rangle \in \operatorname{precedes\text{-}strict}(c)) \Rightarrow \\ \quad \quad \operatorname{originates}(\operatorname{bspace}(c), u, n) \end{array}}}

\vnbbookfile{structure-library/strand-adversary}{structure-library/strand-adversary}
\noindent{\small strand-adversary.scm -- THE ADVERSARY: its traces, the infiltrated strand space (The paper's Definitions 3.1 and 3.2, Thayer, Herzog, Guttman, \textasciitilde{}/docs/jcs\_strand\_spaces.pdf, section 3.1, on the user's model of notes-60 with ENCRYPTION (the user's decision of 2026-10-06 16:40Z): the binary C and S traces made n-ary, the ...}\par
\vnbstructdef{infiltrated-space}{An infiltrated space $s$ is a 3-tuple whose components are listed below.\par
\begin{enumerate}[nosep]
\item $\operatorname{ispace}(s) = \operatorname{nth}(1, s)$, a substructure (\texttt{substructure}): $\operatorname{is\text{-}strand\text{-}space}(\operatorname{ispace}(s))$.
\item $\operatorname{adv\text{-}ids}(s) = \operatorname{nth}(2, s)$, a carrier (\texttt{carriers}): $\operatorname{adv\text{-}ids}(s) \in \mathbf{Set}$.
\item $\operatorname{adv\text{-}tokens}(s) = \operatorname{nth}(3, s)$, a carrier (\texttt{carriers}): $\operatorname{adv\text{-}tokens}(s) \in \mathbf{Set}$.
\end{enumerate}
The laws (the \texttt{property} and \texttt{law} clauses):
\begin{itemize}[nosep]
\item $\forall \operatorname{sap\_} \in \operatorname{adv\text{-}ids}(s).\; (\operatorname{sap\_} \in \operatorname{ids}(\operatorname{ispace}(s)))$
\item $\forall \operatorname{sat\_} \in \operatorname{adv\text{-}tokens}(s).\; (\operatorname{sat\_} \in \operatorname{tokens}(\operatorname{ispace}(s)))$
\item
\begin{equation*}
\begin{array}{@{}l@{}} \forall \operatorname{sap\_} \in \operatorname{adv\text{-}ids}(s). \\ \quad \operatorname{is\text{-}adversary\text{-}trace}(\operatorname{ispace}(s), \operatorname{adv\text{-}tokens}(s), \operatorname{strand\text{-}map}(\operatorname{ispace}(s))(\operatorname{sap\_})) \end{array}
\end{equation*}

\end{itemize}
The characteristic law (the axiom \texttt{is-infiltrated-space}): $\operatorname{is\text{-}infiltrated\text{-}space}(s)$ holds if and only if $\operatorname{length}(s) = 3$, every component has the typing listed with it, and every law holds.
}
\par\medskip\noindent\textbf{Definition} (\texttt{is-text-trace}).\;For all $\operatorname{sasp\_}$, $\operatorname{sap\_}$ and $\operatorname{sas\_}$, 
\begin{equation*}
\operatorname{is\text{-}text\text{-}trace}(\operatorname{sasp\_}, \operatorname{sap\_}, \operatorname{sas\_}) \Leftrightarrow \exists \operatorname{saa\_} \in \operatorname{sap\_}.\; (\operatorname{sas\_} = \langle \langle \mathit{sent}, \operatorname{saa\_} \rangle \rangle)
\end{equation*}

\par\medskip\noindent\textbf{Definition} (\texttt{is-flush-trace}).\;For all $\operatorname{sasp\_}$ and $\operatorname{sas\_}$, 
\begin{equation*}
\begin{array}{@{}l@{}} \operatorname{is\text{-}flush\text{-}trace}(\operatorname{sasp\_}, \operatorname{sas\_}) \\ \Leftrightarrow\; \exists \operatorname{sag\_} \in \operatorname{msg\text{-}set}(\operatorname{tokens}(\operatorname{sasp\_}), \operatorname{keys}(\operatorname{sasp\_})). \\ \quad \operatorname{sas\_} = \langle \langle \mathit{received}, \operatorname{sag\_} \rangle \rangle \end{array}
\end{equation*}

\par\medskip\noindent\textbf{Definition} (\texttt{is-tee-trace}).\;For all $\operatorname{sasp\_}$ and $\operatorname{sas\_}$, 
\begin{equation*}
\begin{array}{@{}l@{}} \operatorname{is\text{-}tee\text{-}trace}(\operatorname{sasp\_}, \operatorname{sas\_}) \\ \Leftrightarrow\; \exists \operatorname{sag\_} \in \operatorname{msg\text{-}set}(\operatorname{tokens}(\operatorname{sasp\_}), \operatorname{keys}(\operatorname{sasp\_})). \\ \quad \operatorname{sas\_} = \langle \langle \mathit{received}, \operatorname{sag\_} \rangle, \langle \mathit{sent}, \operatorname{sag\_} \rangle, \langle \mathit{sent}, \operatorname{sag\_} \rangle \rangle \end{array}
\end{equation*}

\par\medskip\noindent\textbf{Definition} (\texttt{is-concat-trace}).\;For all $\operatorname{sasp\_}$ and $\operatorname{sas\_}$, 
\begin{equation*}
\begin{array}{@{}l@{}} \operatorname{is\text{-}concat\text{-}trace}(\operatorname{sasp\_}, \operatorname{sas\_}) \\ \Leftrightarrow\; \exists \operatorname{sal\_} \in \operatorname{tuples}(\operatorname{msg\text{-}set}(\operatorname{tokens}(\operatorname{sasp\_}), \operatorname{keys}(\operatorname{sasp\_}))). \\ \quad \operatorname{length}(\operatorname{sas\_}) = \operatorname{succ}(\operatorname{length}(\operatorname{sal\_})) \wedge \\ \quad \forall \operatorname{sai\_} \in \mathbb{N}. \\ \quad \quad (1 \leq \operatorname{sai\_}) \Rightarrow  \\ \quad \quad (\operatorname{sai\_} \leq \operatorname{length}(\operatorname{sal\_})) \Rightarrow (\operatorname{nth}(\operatorname{sai\_}, \operatorname{sas\_}) = \langle \mathit{received}, \operatorname{nth}(\operatorname{sai\_}, \operatorname{sal\_}) \rangle) \wedge \\ \quad \operatorname{nth}(\operatorname{succ}(\operatorname{length}(\operatorname{sal\_})), \operatorname{sas\_}) = \langle \mathit{sent}, \operatorname{lst}(\operatorname{sal\_}) \rangle \end{array}
\end{equation*}

\par\medskip\noindent\textbf{Definition} (\texttt{is-separate-trace}).\;For all $\operatorname{sasp\_}$ and $\operatorname{sas\_}$, 
\begin{equation*}
\begin{array}{@{}l@{}} \operatorname{is\text{-}separate\text{-}trace}(\operatorname{sasp\_}, \operatorname{sas\_}) \\ \Leftrightarrow\; \exists \operatorname{sal\_} \in \operatorname{tuples}(\operatorname{msg\text{-}set}(\operatorname{tokens}(\operatorname{sasp\_}), \operatorname{keys}(\operatorname{sasp\_}))). \\ \quad \operatorname{length}(\operatorname{sas\_}) = \operatorname{succ}(\operatorname{length}(\operatorname{sal\_})) \wedge \\ \quad \operatorname{nth}(1, \operatorname{sas\_}) = \langle \mathit{received}, \operatorname{lst}(\operatorname{sal\_}) \rangle \wedge \\ \quad \forall \operatorname{sai\_} \in \mathbb{N}. \\ \quad \quad (1 \leq \operatorname{sai\_}) \Rightarrow  \\ \quad \quad (\operatorname{sai\_} \leq \operatorname{length}(\operatorname{sal\_})) \Rightarrow (\operatorname{nth}(\operatorname{succ}(\operatorname{sai\_}), \operatorname{sas\_}) = \langle \mathit{sent}, \operatorname{nth}(\operatorname{sai\_}, \operatorname{sal\_}) \rangle) \end{array}
\end{equation*}

\par\medskip\noindent\textbf{Definition} (\texttt{is-encrypt-trace}).\;For all $\operatorname{sasp\_}$ and $\operatorname{sas\_}$, 
\begin{equation*}
\begin{array}{@{}l@{}} \operatorname{is\text{-}encrypt\text{-}trace}(\operatorname{sasp\_}, \operatorname{sas\_}) \\ \Leftrightarrow\; \exists \operatorname{sak\_} \in \operatorname{keys}(\operatorname{sasp\_}), \operatorname{sag\_} \in \operatorname{msg\text{-}set}(\operatorname{tokens}(\operatorname{sasp\_}), \operatorname{keys}(\operatorname{sasp\_})). \\ \quad \operatorname{sas\_} = \langle \langle \mathit{received}, \operatorname{sak\_} \rangle, \langle \mathit{received}, \operatorname{sag\_} \rangle, \langle \mathit{sent}, \operatorname{encr}(\operatorname{sag\_}, \operatorname{sak\_}) \rangle \rangle \end{array}
\end{equation*}

\par\medskip\noindent\textbf{Definition} (\texttt{is-decrypt-trace}).\;For all $\operatorname{sasp\_}$ and $\operatorname{sas\_}$, 
\begin{equation*}
\begin{array}{@{}l@{}} \operatorname{is\text{-}decrypt\text{-}trace}(\operatorname{sasp\_}, \operatorname{sas\_}) \\ \Leftrightarrow\; \exists \operatorname{sak\_} \in \operatorname{keys}(\operatorname{sasp\_}), \operatorname{sag\_} \in \operatorname{msg\text{-}set}(\operatorname{tokens}(\operatorname{sasp\_}), \operatorname{keys}(\operatorname{sasp\_})). \\ \quad \operatorname{sas\_} = \langle \langle \mathit{received}, \operatorname{kinv}(\operatorname{sasp\_})(\operatorname{sak\_}) \rangle, \langle \mathit{received}, \operatorname{encr}(\operatorname{sag\_}, \operatorname{sak\_}) \rangle, \langle \mathit{sent}, \operatorname{sag\_} \rangle \rangle \end{array}
\end{equation*}

\par\medskip\noindent\textbf{Definition} (\texttt{is-adversary-trace}).\;For all $\operatorname{sasp\_}$, $\operatorname{sap\_}$ and $\operatorname{sas\_}$, 
\begin{equation*}
\begin{array}{@{}l@{}} \operatorname{is\text{-}adversary\text{-}trace}(\operatorname{sasp\_}, \operatorname{sap\_}, \operatorname{sas\_}) \\ \Leftrightarrow\; \operatorname{is\text{-}text\text{-}trace}(\operatorname{sasp\_}, \operatorname{sap\_}, \operatorname{sas\_}) \vee \\ \operatorname{is\text{-}flush\text{-}trace}(\operatorname{sasp\_}, \operatorname{sas\_}) \vee \\ \operatorname{is\text{-}tee\text{-}trace}(\operatorname{sasp\_}, \operatorname{sas\_}) \vee \\ \operatorname{is\text{-}concat\text{-}trace}(\operatorname{sasp\_}, \operatorname{sas\_}) \vee \\ \operatorname{is\text{-}separate\text{-}trace}(\operatorname{sasp\_}, \operatorname{sas\_}) \vee \\ \operatorname{is\text{-}encrypt\text{-}trace}(\operatorname{sasp\_}, \operatorname{sas\_}) \vee \\ \operatorname{is\text{-}decrypt\text{-}trace}(\operatorname{sasp\_}, \operatorname{sas\_}) \end{array}
\end{equation*}

\par\medskip\noindent\textbf{Definition} (\texttt{is-adversary-node}).\;For all $\operatorname{sai\_}$ and $\operatorname{san\_}$, 
\begin{equation*}
\operatorname{is\text{-}adversary\text{-}node}(\operatorname{sai\_}, \operatorname{san\_}) \Leftrightarrow (\operatorname{nth}(2, \operatorname{san\_}) \in \operatorname{adv\text{-}ids}(\operatorname{sai\_}))
\end{equation*}

\par\medskip\noindent\textbf{Definition} (\texttt{is-regular-node}).\;For all $\operatorname{sai\_}$ and $\operatorname{san\_}$, 
\begin{equation*}
\operatorname{is\text{-}regular\text{-}node}(\operatorname{sai\_}, \operatorname{san\_}) \Leftrightarrow \neg (\operatorname{nth}(2, \operatorname{san\_}) \in \operatorname{adv\text{-}ids}(\operatorname{sai\_}))
\end{equation*}

\par\medskip\noindent\textbf{Definition} (\texttt{hom-infiltrated-space}).\;$\operatorname{hom\text{-}infiltrated\text{-}space}(a, b) :=$
\begin{equation*}
\{\operatorname{homp\_} \in ((\operatorname{adv\text{-}ids}(a) \to \operatorname{adv\text{-}ids}(b)) \times (\operatorname{adv\text{-}tokens}(a) \to \operatorname{adv\text{-}tokens}(b))) : \operatorname{is\text{-}hom\text{-}infiltrated\text{-}space}(a, b, \operatorname{nth}(1, \operatorname{homp\_}), \operatorname{nth}(2, \operatorname{homp\_}))\}
\end{equation*}

\vnbdef{infiltrated-space-class}{\vnbstmt{\forall s.\; ((s \in \operatorname{infiltrated\text{-}space}) \Leftrightarrow \operatorname{is\text{-}infiltrated\text{-}space}(s))}}
\vnbdef{is-hom-infiltrated-space-def}{\vnbstmt{\begin{array}{@{}l@{}} \forall a, b, c, d. \\ \quad \operatorname{is\text{-}hom\text{-}infiltrated\text{-}space}(a, b, c, d) \\ \quad \Leftrightarrow\; \operatorname{is\text{-}infiltrated\text{-}space}(a) \wedge \\ \quad \operatorname{is\text{-}infiltrated\text{-}space}(b) \wedge \\ \quad c \in (\operatorname{adv\text{-}ids}(a) \to \operatorname{adv\text{-}ids}(b)) \wedge \\ \quad d \in (\operatorname{adv\text{-}tokens}(a) \to \operatorname{adv\text{-}tokens}(b)) \wedge \\ \quad \operatorname{ispace}(a) = \operatorname{ispace}(b) \end{array}}}

\vnbbookfile{theorem-library/strand-adversary-bound}{theorem-library/strand-adversary-bound}
\noindent{\small strand-adversary-bound.scm -- THE ADVERSARY BOUND (the paper's Proposition 3.3 on the user's model with encryption): a token the adversary does not know, and which originates on no regular node of a bundle, is a subterm of the content of NO node of the bundle. (Thayer, Herzog, Guttman, \textasciitilde{}/docs/jcs\_strand\_spaces.pdf, ...}\par
\vnbprop{text-trace-origin}{\vnbstmt{\begin{array}{@{}l@{}} \forall a, p, u, n. \\ \quad \operatorname{is\text{-}strand\text{-}space}(a) \wedge \\ \quad \forall x \in p.\; (x \in \operatorname{tokens}(a)) \wedge \\ \quad n \in \operatorname{nodes}(a) \wedge \\ \quad \operatorname{is\text{-}text\text{-}trace}(a, p, \operatorname{strand\text{-}map}(a)(\operatorname{nth}(2, n))) \wedge \\ \quad \operatorname{originates}(a, u, n) \Rightarrow \\ \quad \quad u \in p \end{array}}}
\vnbprop{flush-trace-no-sent}{\vnbstmt{\begin{array}{@{}l@{}} \forall a, n. \\ \quad (\operatorname{is\text{-}strand\text{-}space}(a) \wedge ((n \in \operatorname{nodes}(a)) \wedge \operatorname{is\text{-}flush\text{-}trace}(a, \operatorname{strand\text{-}map}(a)(\operatorname{nth}(2, n))))) \Rightarrow  \\ \quad \quad \neg \operatorname{is\text{-}sent}(a, n) \end{array}}}
\vnbprop{tee-trace-no-origin}{\vnbstmt{\begin{array}{@{}l@{}} \forall a, u, n. \\ \quad (\operatorname{is\text{-}strand\text{-}space}(a) \wedge ((n \in \operatorname{nodes}(a)) \wedge \operatorname{is\text{-}tee\text{-}trace}(a, \operatorname{strand\text{-}map}(a)(\operatorname{nth}(2, n))))) \Rightarrow  \\ \quad \quad \neg \operatorname{originates}(a, u, n) \end{array}}}
\vnbprop{concat-trace-no-origin}{\vnbstmt{\begin{array}{@{}l@{}} \forall a, u, n. \\ \quad \operatorname{is\text{-}strand\text{-}space}(a) \wedge \\ \quad u \in \operatorname{tokens}(a) \wedge \\ \quad n \in \operatorname{nodes}(a) \wedge \\ \quad \operatorname{is\text{-}concat\text{-}trace}(a, \operatorname{strand\text{-}map}(a)(\operatorname{nth}(2, n))) \Rightarrow \\ \quad \quad \neg \operatorname{originates}(a, u, n) \end{array}}}
\vnbprop{separate-trace-no-origin}{\vnbstmt{\begin{array}{@{}l@{}} \forall a, u, n. \\ \quad (\operatorname{is\text{-}strand\text{-}space}(a) \wedge ((n \in \operatorname{nodes}(a)) \wedge \operatorname{is\text{-}separate\text{-}trace}(a, \operatorname{strand\text{-}map}(a)(\operatorname{nth}(2, n))))) \Rightarrow  \\ \quad \quad \neg \operatorname{originates}(a, u, n) \end{array}}}
\vnbprop{encrypt-trace-no-origin}{\vnbstmt{\begin{array}{@{}l@{}} \forall a, u, n. \\ \quad \operatorname{is\text{-}strand\text{-}space}(a) \wedge \\ \quad u \in \operatorname{tokens}(a) \wedge \\ \quad n \in \operatorname{nodes}(a) \wedge \\ \quad \operatorname{is\text{-}encrypt\text{-}trace}(a, \operatorname{strand\text{-}map}(a)(\operatorname{nth}(2, n))) \Rightarrow \\ \quad \quad \neg \operatorname{originates}(a, u, n) \end{array}}}
\vnbprop{decrypt-trace-no-origin}{\vnbstmt{\begin{array}{@{}l@{}} \forall a, u, n. \\ \quad (\operatorname{is\text{-}strand\text{-}space}(a) \wedge ((n \in \operatorname{nodes}(a)) \wedge \operatorname{is\text{-}decrypt\text{-}trace}(a, \operatorname{strand\text{-}map}(a)(\operatorname{nth}(2, n))))) \Rightarrow  \\ \quad \quad \neg \operatorname{originates}(a, u, n) \end{array}}}
\vnbprop{adversary-node-origin}{\vnbstmt{\begin{array}{@{}l@{}} \forall p, u, n. \\ \quad \operatorname{is\text{-}infiltrated\text{-}space}(p) \wedge \\ \quad u \in \operatorname{tokens}(\operatorname{ispace}(p)) \wedge \\ \quad n \in \operatorname{nodes}(\operatorname{ispace}(p)) \wedge \\ \quad \operatorname{is\text{-}adversary\text{-}node}(p, n) \wedge \\ \quad \operatorname{originates}(\operatorname{ispace}(p), u, n) \Rightarrow \\ \quad \quad u \in \operatorname{adv\text{-}tokens}(p) \end{array}}}
\vnbprop{adversary-bound}{\vnbstmt{\begin{array}{@{}l@{}} \forall p, c, u.; \forall n \in \operatorname{bnodes}(c). \\ \quad \operatorname{is\text{-}infiltrated\text{-}space}(p) \wedge \\ \quad \operatorname{is\text{-}bundle}(c) \wedge \\ \quad \operatorname{bspace}(c) = \operatorname{ispace}(p) \wedge \\ \quad u \in \operatorname{tokens}(\operatorname{ispace}(p)) \wedge \\ \quad \neg (u \in \operatorname{adv\text{-}tokens}(p)) \wedge \\ \quad \forall n \in \operatorname{bnodes}(c). \\ \quad \quad \operatorname{is\text{-}regular\text{-}node}(p, n) \Rightarrow \neg \operatorname{originates}(\operatorname{ispace}(p), u, n) \Rightarrow \\ \quad \quad \neg \operatorname{is\text{-}subterm}(\operatorname{tokens}(\operatorname{ispace}(p)), \operatorname{keys}(\operatorname{ispace}(p)), u, \operatorname{node\text{-}content}(\operatorname{ispace}(p), n)) \end{array}}}
\vnbprop{adversary-bound-adversary-nodes}{\vnbstmt{\begin{array}{@{}l@{}} \forall p, c, u.; \forall n \in \operatorname{bnodes}(c). \\ \quad \operatorname{is\text{-}infiltrated\text{-}space}(p) \wedge \\ \quad \operatorname{is\text{-}bundle}(c) \wedge \\ \quad \operatorname{bspace}(c) = \operatorname{ispace}(p) \wedge \\ \quad u \in \operatorname{tokens}(\operatorname{ispace}(p)) \wedge \\ \quad \neg (u \in \operatorname{adv\text{-}tokens}(p)) \wedge \\ \quad \forall n \in \operatorname{bnodes}(c). \\ \quad \quad \operatorname{is\text{-}regular\text{-}node}(p, n) \Rightarrow \neg \operatorname{originates}(\operatorname{ispace}(p), u, n) \wedge \\ \quad \operatorname{is\text{-}adversary\text{-}node}(p, n) \Rightarrow \\ \quad \quad \neg \operatorname{is\text{-}subterm}(\operatorname{tokens}(\operatorname{ispace}(p)), \operatorname{keys}(\operatorname{ispace}(p)), u, \operatorname{node\text{-}content}(\operatorname{ispace}(p), n)) \end{array}}}
\vnbprop{sent-received-in-pair}{\vnbstmt{(\mathit{sent} \in \{\mathit{sent}, \mathit{received}\}) \wedge (\mathit{received} \in \{\mathit{sent}, \mathit{received}\})}}
\vnbprop{directed-intro}{\vnbstmt{\begin{array}{@{}l@{}} \forall t, k.; \forall g \in \operatorname{msg\text{-}set}(t, k). \\ \quad (\langle \mathit{sent}, g \rangle \in \operatorname{directed}(t, k)) \wedge (\langle \mathit{received}, g \rangle \in \operatorname{directed}(t, k)) \end{array}}}
\vnbprop{separate-trace-exists}{\vnbstmt{\begin{array}{@{}l@{}} \forall a, l. \\ \quad (\operatorname{is\text{-}strand\text{-}space}(a) \wedge (l \in \operatorname{tuples}(\operatorname{msg\text{-}set}(\operatorname{tokens}(a), \operatorname{keys}(a))))) \Rightarrow  \\ \quad \quad \exists s \in \operatorname{strands}(\operatorname{tokens}(a), \operatorname{keys}(a)).\; \operatorname{is\text{-}separate\text{-}trace}(a, s) \end{array}}}
\vnbprop{concat-trace-exists}{\vnbstmt{\begin{array}{@{}l@{}} \forall a, l. \\ \quad (\operatorname{is\text{-}strand\text{-}space}(a) \wedge (l \in \operatorname{tuples}(\operatorname{msg\text{-}set}(\operatorname{tokens}(a), \operatorname{keys}(a))))) \Rightarrow  \\ \quad \quad \exists s \in \operatorname{strands}(\operatorname{tokens}(a), \operatorname{keys}(a)).\; \operatorname{is\text{-}concat\text{-}trace}(a, s) \end{array}}}
\vnbprop{encrypt-trace-strand}{\vnbstmt{\begin{array}{@{}l@{}} \forall a, k, m. \\ \quad (\operatorname{is\text{-}strand\text{-}space}(a) \wedge ((k \in \operatorname{keys}(a)) \wedge (m \in \operatorname{msg\text{-}set}(\operatorname{tokens}(a), \operatorname{keys}(a))))) \Rightarrow  \\ \quad \quad \langle \langle \mathit{received}, k \rangle, \langle \mathit{received}, m \rangle, \langle \mathit{sent}, \operatorname{encr}(m, k) \rangle \rangle \in \operatorname{strands}(\operatorname{tokens}(a), \operatorname{keys}(a)) \wedge \\ \quad \quad \operatorname{is\text{-}encrypt\text{-}trace}(a, \langle \langle \mathit{received}, k \rangle, \langle \mathit{received}, m \rangle, \langle \mathit{sent}, \operatorname{encr}(m, k) \rangle \rangle) \end{array}}}
\vnbprop{decrypt-trace-strand}{\vnbstmt{\begin{array}{@{}l@{}} \forall a, k, m. \\ \quad (\operatorname{is\text{-}strand\text{-}space}(a) \wedge ((k \in \operatorname{keys}(a)) \wedge (m \in \operatorname{msg\text{-}set}(\operatorname{tokens}(a), \operatorname{keys}(a))))) \Rightarrow  \\ \quad \quad \langle \langle \mathit{received}, \operatorname{kinv}(a)(k) \rangle, \langle \mathit{received}, \operatorname{encr}(m, k) \rangle, \langle \mathit{sent}, m \rangle \rangle \in \operatorname{strands}(\operatorname{tokens}(a), \operatorname{keys}(a)) \wedge \\ \quad \quad \operatorname{is\text{-}decrypt\text{-}trace}(a, \langle \langle \mathit{received}, \operatorname{kinv}(a)(k) \rangle, \langle \mathit{received}, \operatorname{encr}(m, k) \rangle, \langle \mathit{sent}, m \rangle \rangle) \end{array}}}

\vnbbookfile{structure-library/nsl}{structure-library/nsl}
\noindent{\small nsl.scm -- THE NEEDHAM-SCHROEDER-LOWE PROTOCOL ON THE USER'S STRAND MODEL: names, public keys, the initiator's and the responder's strands, NSL spaces. (Thayer, Herzog, Guttman, \textasciitilde{}/docs/jcs\_strand\_spaces.pdf, section 5: the paper's Definition 5.1 and the term-algebra specialisation that precedes it; the model is ...}\par
\vnbstructdef{pk-space}{A pk space $s$ is a 3-tuple whose components are listed below.\par
\begin{enumerate}[nosep]
\item $\operatorname{pk\text{-}inf}(s) = \operatorname{nth}(1, s)$, a substructure (\texttt{substructure}): $\operatorname{is\text{-}infiltrated\text{-}space}(\operatorname{pk\text{-}inf}(s))$.
\item $\operatorname{names}(s) = \operatorname{nth}(2, s)$, a carrier (\texttt{carriers}): $\operatorname{names}(s) \in \mathbf{Set}$.
\item $\operatorname{pubkey}(s) = \operatorname{nth}(3, s)$, an operation (\texttt{op}): $\operatorname{pubkey}(s) \in (\operatorname{names}(s) \to \operatorname{keys}(\operatorname{ispace}(\operatorname{pk\text{-}inf}(s))))$.
\end{enumerate}
The laws (the \texttt{property} and \texttt{law} clauses):
\begin{itemize}[nosep]
\item $\forall \operatorname{nlx\_} \in \operatorname{names}(s).\; (\operatorname{nlx\_} \in \operatorname{tokens}(\operatorname{ispace}(\operatorname{pk\text{-}inf}(s))))$
\item $\forall \operatorname{nlx\_} \in \operatorname{names}(s).\; \neg (\operatorname{nlx\_} \in \operatorname{keys}(\operatorname{ispace}(\operatorname{pk\text{-}inf}(s))))$
\item
\begin{equation*}
\begin{array}{@{}l@{}} \forall \operatorname{nlx\_}, \operatorname{nly\_} \in \operatorname{names}(s). \\ \quad (\operatorname{pubkey}(s)(\operatorname{nlx\_}) = \operatorname{pubkey}(s)(\operatorname{nly\_})) \Rightarrow (\operatorname{nlx\_} = \operatorname{nly\_}) \end{array}
\end{equation*}

\end{itemize}
The characteristic law (the axiom \texttt{is-pk-space}): $\operatorname{is\text{-}pk\text{-}space}(s)$ holds if and only if $\operatorname{length}(s) = 3$, every component has the typing listed with it, and every law holds.
}
\par\medskip\noindent\textbf{Definition} (\texttt{is-init-strand}).\;For all $\operatorname{nlp\_}$, $\operatorname{nls\_}$, $\operatorname{nla\_}$, $\operatorname{nlb\_}$, $\operatorname{nlm\_}$ and $\operatorname{nln\_}$, 
\begin{equation*}
\begin{array}{@{}l@{}} \operatorname{is\text{-}init\text{-}strand}(\operatorname{nlp\_}, \operatorname{nls\_}, \operatorname{nla\_}, \operatorname{nlb\_}, \operatorname{nlm\_}, \operatorname{nln\_}) \\ \Leftrightarrow\; \operatorname{nla\_} \in \operatorname{names}(\operatorname{nlp\_}) \wedge \\ \operatorname{nlb\_} \in \operatorname{names}(\operatorname{nlp\_}) \wedge \\ \operatorname{nlm\_} \in \operatorname{tokens}(\operatorname{ispace}(\operatorname{pk\text{-}inf}(\operatorname{nlp\_}))) \wedge \\ \operatorname{nln\_} \in \operatorname{tokens}(\operatorname{ispace}(\operatorname{pk\text{-}inf}(\operatorname{nlp\_}))) \wedge \\ \neg (\operatorname{nlm\_} \in \operatorname{names}(\operatorname{nlp\_})) \wedge \\ \neg (\operatorname{nlm\_} \in \operatorname{keys}(\operatorname{ispace}(\operatorname{pk\text{-}inf}(\operatorname{nlp\_})))) \wedge \\ \neg (\operatorname{nln\_} \in \operatorname{keys}(\operatorname{ispace}(\operatorname{pk\text{-}inf}(\operatorname{nlp\_})))) \wedge \\ \operatorname{nls\_} \\ =\; \langle \langle \mathit{sent}, \operatorname{encr}(\operatorname{lst}(\langle \operatorname{nlm\_}, \operatorname{nla\_} \rangle), \operatorname{pubkey}(\operatorname{nlp\_})(\operatorname{nlb\_})) \rangle, \langle \mathit{received}, \operatorname{encr}(\operatorname{lst}(\langle \operatorname{nlm\_}, \operatorname{nln\_}, \operatorname{nlb\_} \rangle), \operatorname{pubkey}(\operatorname{nlp\_})(\operatorname{nla\_})) \rangle, \langle \mathit{sent}, \operatorname{encr}(\operatorname{nln\_}, \operatorname{pubkey}(\operatorname{nlp\_})(\operatorname{nlb\_})) \rangle \rangle \end{array}
\end{equation*}

\par\medskip\noindent\textbf{Definition} (\texttt{is-resp-strand}).\;For all $\operatorname{nlp\_}$, $\operatorname{nls\_}$, $\operatorname{nla\_}$, $\operatorname{nlb\_}$, $\operatorname{nlm\_}$ and $\operatorname{nln\_}$, 
\begin{equation*}
\begin{array}{@{}l@{}} \operatorname{is\text{-}resp\text{-}strand}(\operatorname{nlp\_}, \operatorname{nls\_}, \operatorname{nla\_}, \operatorname{nlb\_}, \operatorname{nlm\_}, \operatorname{nln\_}) \\ \Leftrightarrow\; \operatorname{nla\_} \in \operatorname{names}(\operatorname{nlp\_}) \wedge \\ \operatorname{nlb\_} \in \operatorname{names}(\operatorname{nlp\_}) \wedge \\ \operatorname{nlm\_} \in \operatorname{tokens}(\operatorname{ispace}(\operatorname{pk\text{-}inf}(\operatorname{nlp\_}))) \wedge \\ \operatorname{nln\_} \in \operatorname{tokens}(\operatorname{ispace}(\operatorname{pk\text{-}inf}(\operatorname{nlp\_}))) \wedge \\ \neg (\operatorname{nln\_} \in \operatorname{names}(\operatorname{nlp\_})) \wedge \\ \neg (\operatorname{nlm\_} \in \operatorname{keys}(\operatorname{ispace}(\operatorname{pk\text{-}inf}(\operatorname{nlp\_})))) \wedge \\ \neg (\operatorname{nln\_} \in \operatorname{keys}(\operatorname{ispace}(\operatorname{pk\text{-}inf}(\operatorname{nlp\_})))) \wedge \\ \operatorname{nls\_} \\ =\; \langle \langle \mathit{received}, \operatorname{encr}(\operatorname{lst}(\langle \operatorname{nlm\_}, \operatorname{nla\_} \rangle), \operatorname{pubkey}(\operatorname{nlp\_})(\operatorname{nlb\_})) \rangle, \langle \mathit{sent}, \operatorname{encr}(\operatorname{lst}(\langle \operatorname{nlm\_}, \operatorname{nln\_}, \operatorname{nlb\_} \rangle), \operatorname{pubkey}(\operatorname{nlp\_})(\operatorname{nla\_})) \rangle, \langle \mathit{received}, \operatorname{encr}(\operatorname{nln\_}, \operatorname{pubkey}(\operatorname{nlp\_})(\operatorname{nlb\_})) \rangle \rangle \end{array}
\end{equation*}

\par\medskip\noindent\textbf{Definition} (\texttt{is-nsl-space}).\;For every $\operatorname{nlp\_}$, 
\begin{equation*}
\begin{array}{@{}l@{}} \operatorname{is\text{-}nsl\text{-}space}(\operatorname{nlp\_}) \\ \Leftrightarrow\; \operatorname{is\text{-}pk\text{-}space}(\operatorname{nlp\_}) \wedge \\ \forall \operatorname{nli\_} \in \operatorname{ids}(\operatorname{ispace}(\operatorname{pk\text{-}inf}(\operatorname{nlp\_}))). \\ \quad \operatorname{nli\_} \in \operatorname{adv\text{-}ids}(\operatorname{pk\text{-}inf}(\operatorname{nlp\_})) \vee \\ \quad \exists \operatorname{nla\_}, \operatorname{nlb\_}, \operatorname{nlm\_}, \operatorname{nln\_}. \\ \quad \quad \operatorname{is\text{-}init\text{-}strand}(\operatorname{nlp\_}, \operatorname{strand\text{-}map}(\operatorname{ispace}(\operatorname{pk\text{-}inf}(\operatorname{nlp\_})))(\operatorname{nli\_}), \operatorname{nla\_}, \operatorname{nlb\_}, \operatorname{nlm\_}, \operatorname{nln\_}) \vee \\ \quad \exists \operatorname{nla\_}, \operatorname{nlb\_}, \operatorname{nlm\_}, \operatorname{nln\_}. \\ \quad \quad \operatorname{is\text{-}resp\text{-}strand}(\operatorname{nlp\_}, \operatorname{strand\text{-}map}(\operatorname{ispace}(\operatorname{pk\text{-}inf}(\operatorname{nlp\_})))(\operatorname{nli\_}), \operatorname{nla\_}, \operatorname{nlb\_}, \operatorname{nlm\_}, \operatorname{nln\_}) \end{array}
\end{equation*}

\par\medskip\noindent\textbf{Definition} (\texttt{hom-pk-space}).\;$\operatorname{hom\text{-}pk\text{-}space}(a, b) :=$
\begin{equation*}
\{\operatorname{homf\_} \in (\operatorname{names}(a) \to \operatorname{names}(b)) : \operatorname{is\text{-}hom\text{-}pk\text{-}space}(a, b, \operatorname{homf\_})\}
\end{equation*}

\vnbdef{is-hom-pk-space-def}{\vnbstmt{\begin{array}{@{}l@{}} \forall a, b, f. \\ \quad \operatorname{is\text{-}hom\text{-}pk\text{-}space}(a, b, f) \\ \quad \Leftrightarrow\; \operatorname{is\text{-}pk\text{-}space}(a) \wedge \\ \quad \operatorname{is\text{-}pk\text{-}space}(b) \wedge \\ \quad f \in (\operatorname{names}(a) \to \operatorname{names}(b)) \wedge \\ \quad \operatorname{pk\text{-}inf}(a) = \operatorname{pk\text{-}inf}(b) \wedge \\ \quad \forall x_{1} \in \operatorname{names}(a).\; (\operatorname{pubkey}(a)(x_{1}) = \operatorname{pubkey}(b)(f(x_{1}))) \end{array}}}
\vnbdef{pk-space-class}{\vnbstmt{\forall s.\; ((s \in \operatorname{pk\text{-}space}) \Leftrightarrow \operatorname{is\text{-}pk\text{-}space}(s))}}

\vnbbookfile{theorem-library/nsl}{theorem-library/nsl}
\noindent{\small nsl.scm -- THE NEEDHAM-SCHROEDER-LOWE PROTOCOL ON THE USER'S STRAND MODEL: the responder's guarantee (agreement, Proposition 5.2) and the secrecy of both nonces, the injectivity of the initiator's strand, the initiator's agreement (Propositions 5.8, 5.10, 5.13, 5.14). (Thayer, Herzog, Guttman, ...}\par
\vnbprop{pk-space-parts}{\vnbstmt{\begin{array}{@{}l@{}} \forall p. \\ \quad \operatorname{is\text{-}pk\text{-}space}(p) \Rightarrow  \\ \quad \quad \operatorname{is\text{-}infiltrated\text{-}space}(\operatorname{pk\text{-}inf}(p)) \wedge \\ \quad \quad \operatorname{is\text{-}strand\text{-}space}(\operatorname{ispace}(\operatorname{pk\text{-}inf}(p))) \wedge \\ \quad \quad \operatorname{names}(p) \in \mathbf{Set} \wedge \\ \quad \quad \operatorname{pubkey}(p) \in (\operatorname{names}(p) \to \operatorname{keys}(\operatorname{ispace}(\operatorname{pk\text{-}inf}(p)))) \wedge \\ \quad \quad \forall x \in \operatorname{names}(p).\; (x \in \operatorname{tokens}(\operatorname{ispace}(\operatorname{pk\text{-}inf}(p)))) \wedge \\ \quad \quad \forall x \in \operatorname{names}(p).\; \neg (x \in \operatorname{keys}(\operatorname{ispace}(\operatorname{pk\text{-}inf}(p)))) \wedge \\ \quad \quad \forall x, y. \\ \quad \quad \quad ((x \in \operatorname{names}(p)) \wedge ((y \in \operatorname{names}(p)) \wedge (\operatorname{pubkey}(p)(x) = \operatorname{pubkey}(p)(y)))) \Rightarrow  \\ \quad \quad \quad \quad x = y \end{array}}}
\vnbprop{three-strand-node}{\vnbstmt{\begin{array}{@{}l@{}} \forall a, n, b, c, d. \\ \quad ((n \in \operatorname{nodes}(a)) \wedge (\operatorname{strand\text{-}map}(a)(\operatorname{nth}(2, n)) = \langle b, c, d \rangle)) \Rightarrow  \\ \quad \quad ((\operatorname{nth}(1, n) = 1) \wedge ((\operatorname{node\text{-}content}(a, n) \simeq \operatorname{nth}(2, b)) \wedge (\operatorname{node\text{-}dir}(a, n) \simeq \operatorname{nth}(1, b)))) \vee \\ \quad \quad ((\operatorname{nth}(1, n) = 2) \wedge ((\operatorname{node\text{-}content}(a, n) \simeq \operatorname{nth}(2, c)) \wedge (\operatorname{node\text{-}dir}(a, n) \simeq \operatorname{nth}(1, c)))) \vee \\ \quad \quad ((\operatorname{nth}(1, n) = 3) \wedge ((\operatorname{node\text{-}content}(a, n) \simeq \operatorname{nth}(2, d)) \wedge (\operatorname{node\text{-}dir}(a, n) \simeq \operatorname{nth}(1, d)))) \end{array}}}
\vnbprop{three-strand-nodes}{\vnbstmt{\begin{array}{@{}l@{}} \forall b, a, c, d, f. \\ \quad ((a \in \operatorname{ids}(b)) \wedge (\operatorname{strand\text{-}map}(b)(a) = \langle c, d, f \rangle)) \Rightarrow  \\ \quad \quad \langle 1, a \rangle \in \operatorname{nodes}(b) \wedge \\ \quad \quad (\operatorname{node\text{-}content}(b, \langle 1, a \rangle) \simeq \operatorname{nth}(2, c)) \wedge \\ \quad \quad (\operatorname{node\text{-}dir}(b, \langle 1, a \rangle) \simeq \operatorname{nth}(1, c)) \wedge \\ \quad \quad \langle 2, a \rangle \in \operatorname{nodes}(b) \wedge \\ \quad \quad (\operatorname{node\text{-}content}(b, \langle 2, a \rangle) \simeq \operatorname{nth}(2, d)) \wedge \\ \quad \quad (\operatorname{node\text{-}dir}(b, \langle 2, a \rangle) \simeq \operatorname{nth}(1, d)) \wedge \\ \quad \quad \langle 3, a \rangle \in \operatorname{nodes}(b) \wedge \\ \quad \quad (\operatorname{node\text{-}content}(b, \langle 3, a \rangle) \simeq \operatorname{nth}(2, f)) \wedge \\ \quad \quad (\operatorname{node\text{-}dir}(b, \langle 3, a \rangle) \simeq \operatorname{nth}(1, f)) \end{array}}}
\vnbprop{init-strand-is-strand}{\vnbstmt{\begin{array}{@{}l@{}} \forall p, s, a, b, c, d. \\ \quad (\operatorname{is\text{-}pk\text{-}space}(p) \wedge \operatorname{is\text{-}init\text{-}strand}(p, s, a, b, c, d)) \Rightarrow  \\ \quad \quad s \in \operatorname{strands}(\operatorname{tokens}(\operatorname{ispace}(\operatorname{pk\text{-}inf}(p))), \operatorname{keys}(\operatorname{ispace}(\operatorname{pk\text{-}inf}(p)))) \end{array}}}
\vnbprop{resp-strand-is-strand}{\vnbstmt{\begin{array}{@{}l@{}} \forall p, s, a, b, c, d. \\ \quad (\operatorname{is\text{-}pk\text{-}space}(p) \wedge \operatorname{is\text{-}resp\text{-}strand}(p, s, a, b, c, d)) \Rightarrow  \\ \quad \quad s \in \operatorname{strands}(\operatorname{tokens}(\operatorname{ispace}(\operatorname{pk\text{-}inf}(p))), \operatorname{keys}(\operatorname{ispace}(\operatorname{pk\text{-}inf}(p)))) \end{array}}}
\vnbprop{init-strand-nodes}{\vnbstmt{\begin{array}{@{}l@{}} \forall p, a, b, c, d, f. \\ \quad \operatorname{is\text{-}pk\text{-}space}(p) \wedge \\ \quad a \in \operatorname{ids}(\operatorname{ispace}(\operatorname{pk\text{-}inf}(p))) \wedge \\ \quad \operatorname{is\text{-}init\text{-}strand}(p, \operatorname{strand\text{-}map}(\operatorname{ispace}(\operatorname{pk\text{-}inf}(p)))(a), b, c, d, f) \Rightarrow \\ \quad \quad \langle 1, a \rangle \in \operatorname{nodes}(\operatorname{ispace}(\operatorname{pk\text{-}inf}(p))) \wedge \\ \quad \quad \operatorname{is\text{-}sent}(\operatorname{ispace}(\operatorname{pk\text{-}inf}(p)), \langle 1, a \rangle) \wedge \\ \quad \quad \operatorname{node\text{-}content}(\operatorname{ispace}(\operatorname{pk\text{-}inf}(p)), \langle 1, a \rangle) = \operatorname{encr}(\operatorname{lst}(\langle d, b \rangle), \operatorname{pubkey}(p)(c)) \wedge \\ \quad \quad \langle 2, a \rangle \in \operatorname{nodes}(\operatorname{ispace}(\operatorname{pk\text{-}inf}(p))) \wedge \\ \quad \quad \operatorname{is\text{-}received}(\operatorname{ispace}(\operatorname{pk\text{-}inf}(p)), \langle 2, a \rangle) \wedge \\ \quad \quad \operatorname{node\text{-}content}(\operatorname{ispace}(\operatorname{pk\text{-}inf}(p)), \langle 2, a \rangle) = \operatorname{encr}(\operatorname{lst}(\langle d, f, c \rangle), \operatorname{pubkey}(p)(b)) \wedge \\ \quad \quad \langle 3, a \rangle \in \operatorname{nodes}(\operatorname{ispace}(\operatorname{pk\text{-}inf}(p))) \wedge \\ \quad \quad \operatorname{is\text{-}sent}(\operatorname{ispace}(\operatorname{pk\text{-}inf}(p)), \langle 3, a \rangle) \wedge \\ \quad \quad \operatorname{node\text{-}content}(\operatorname{ispace}(\operatorname{pk\text{-}inf}(p)), \langle 3, a \rangle) = \operatorname{encr}(f, \operatorname{pubkey}(p)(c)) \end{array}}}
\vnbprop{resp-strand-nodes}{\vnbstmt{\begin{array}{@{}l@{}} \forall p, a, b, c, d, f. \\ \quad \operatorname{is\text{-}pk\text{-}space}(p) \wedge \\ \quad a \in \operatorname{ids}(\operatorname{ispace}(\operatorname{pk\text{-}inf}(p))) \wedge \\ \quad \operatorname{is\text{-}resp\text{-}strand}(p, \operatorname{strand\text{-}map}(\operatorname{ispace}(\operatorname{pk\text{-}inf}(p)))(a), b, c, d, f) \Rightarrow \\ \quad \quad \langle 1, a \rangle \in \operatorname{nodes}(\operatorname{ispace}(\operatorname{pk\text{-}inf}(p))) \wedge \\ \quad \quad \operatorname{is\text{-}received}(\operatorname{ispace}(\operatorname{pk\text{-}inf}(p)), \langle 1, a \rangle) \wedge \\ \quad \quad \operatorname{node\text{-}content}(\operatorname{ispace}(\operatorname{pk\text{-}inf}(p)), \langle 1, a \rangle) = \operatorname{encr}(\operatorname{lst}(\langle d, b \rangle), \operatorname{pubkey}(p)(c)) \wedge \\ \quad \quad \langle 2, a \rangle \in \operatorname{nodes}(\operatorname{ispace}(\operatorname{pk\text{-}inf}(p))) \wedge \\ \quad \quad \operatorname{is\text{-}sent}(\operatorname{ispace}(\operatorname{pk\text{-}inf}(p)), \langle 2, a \rangle) \wedge \\ \quad \quad \operatorname{node\text{-}content}(\operatorname{ispace}(\operatorname{pk\text{-}inf}(p)), \langle 2, a \rangle) = \operatorname{encr}(\operatorname{lst}(\langle d, f, c \rangle), \operatorname{pubkey}(p)(b)) \wedge \\ \quad \quad \langle 3, a \rangle \in \operatorname{nodes}(\operatorname{ispace}(\operatorname{pk\text{-}inf}(p))) \wedge \\ \quad \quad \operatorname{is\text{-}received}(\operatorname{ispace}(\operatorname{pk\text{-}inf}(p)), \langle 3, a \rangle) \wedge \\ \quad \quad \operatorname{node\text{-}content}(\operatorname{ispace}(\operatorname{pk\text{-}inf}(p)), \langle 3, a \rangle) = \operatorname{encr}(f, \operatorname{pubkey}(p)(c)) \end{array}}}
\vnbprop{token-not-encr}{\vnbstmt{\begin{array}{@{}l@{}} \forall t, g, y. \\ \quad (\operatorname{is\text{-}token\text{-}set}(t) \wedge ((g \in \mathbf{Set}) \wedge (y \in \mathbf{Set}))) \Rightarrow \neg (\operatorname{encr}(g, y) \in t) \end{array}}}
\vnbprop{token-not-lst}{\vnbstmt{\forall t, l.\; ((\operatorname{is\text{-}token\text{-}set}(t) \wedge (l \in \mathbf{Set})) \Rightarrow \neg (\operatorname{lst}(l) \in t))}}
\vnbprop{subterm-of-flat-list}{\vnbstmt{\begin{array}{@{}l@{}} \forall t, k, d, w, l. \\ \quad \operatorname{is\text{-}token\text{-}set}(t) \wedge \\ \quad k \in \mathbf{Set} \wedge \\ \quad \forall x \in d.\; (x \in t) \wedge \\ \quad l \in \operatorname{tuples}(d) \wedge \\ \quad \operatorname{is\text{-}subterm}(t, k, w, \operatorname{lst}(l)) \Rightarrow \\ \quad \quad (w = \operatorname{lst}(l)) \vee (w \in d) \end{array}}}
\vnbprop{subterm-of-flat-encr-list}{\vnbstmt{\begin{array}{@{}l@{}} \forall t, k, d, w, l, y. \\ \quad \operatorname{is\text{-}token\text{-}set}(t) \wedge \\ \quad k \in \mathbf{Set} \wedge \\ \quad \forall x \in d.\; (x \in t) \wedge \\ \quad l \in \operatorname{tuples}(d) \wedge \\ \quad y \in k \wedge \\ \quad \operatorname{is\text{-}subterm}(t, k, w, \operatorname{encr}(\operatorname{lst}(l), y)) \Rightarrow \\ \quad \quad (w = \operatorname{encr}(\operatorname{lst}(l), y)) \vee ((w = \operatorname{lst}(l)) \vee (w \in d)) \end{array}}}
\vnbprop{subterm-of-flat-encr-token}{\vnbstmt{\begin{array}{@{}l@{}} \forall t, k, d, w, x, y. \\ \quad \operatorname{is\text{-}token\text{-}set}(t) \wedge \\ \quad k \in \mathbf{Set} \wedge \\ \quad \forall z \in d.\; (z \in t) \wedge \\ \quad x \in d \wedge \\ \quad y \in k \wedge \\ \quad \operatorname{is\text{-}subterm}(t, k, w, \operatorname{encr}(x, y)) \Rightarrow \\ \quad \quad (w = \operatorname{encr}(x, y)) \vee (w \in d) \end{array}}}
\vnbprop{nsl-regular-content}{\vnbstmt{\begin{array}{@{}l@{}} \forall p, n. \\ \quad (\operatorname{is\text{-}nsl\text{-}space}(p) \wedge ((n \in \operatorname{nodes}(\operatorname{ispace}(\operatorname{pk\text{-}inf}(p)))) \wedge \operatorname{is\text{-}regular\text{-}node}(\operatorname{pk\text{-}inf}(p), n))) \Rightarrow  \\ \quad \quad \exists k \in \operatorname{keys}(\operatorname{ispace}(\operatorname{pk\text{-}inf}(p))). \\ \quad \quad \quad \exists x \in \{d \in \operatorname{tokens}(\operatorname{ispace}(\operatorname{pk\text{-}inf}(p))) : \neg (d \in \operatorname{keys}(\operatorname{ispace}(\operatorname{pk\text{-}inf}(p))))\}. \\ \quad \quad \quad \quad \operatorname{node\text{-}content}(\operatorname{ispace}(\operatorname{pk\text{-}inf}(p)), n) = \operatorname{encr}(x, k) \vee \\ \quad \quad \quad \exists l \in \operatorname{tuples}(\{d \in \operatorname{tokens}(\operatorname{ispace}(\operatorname{pk\text{-}inf}(p))) : \neg (d \in \operatorname{keys}(\operatorname{ispace}(\operatorname{pk\text{-}inf}(p))))\}). \\ \quad \quad \quad \quad \operatorname{node\text{-}content}(\operatorname{ispace}(\operatorname{pk\text{-}inf}(p)), n) = \operatorname{encr}(\operatorname{lst}(l), k) \end{array}}}
\vnbprop{nsl-regular-encr-subterm}{\vnbstmt{\begin{array}{@{}l@{}} \forall p, n, g, y. \\ \quad \operatorname{is\text{-}nsl\text{-}space}(p) \wedge \\ \quad n \in \operatorname{nodes}(\operatorname{ispace}(\operatorname{pk\text{-}inf}(p))) \wedge \\ \quad \operatorname{is\text{-}regular\text{-}node}(\operatorname{pk\text{-}inf}(p), n) \wedge \\ \quad g \in \operatorname{msg\text{-}set}(\operatorname{tokens}(\operatorname{ispace}(\operatorname{pk\text{-}inf}(p))), \operatorname{keys}(\operatorname{ispace}(\operatorname{pk\text{-}inf}(p)))) \wedge \\ \quad y \in \operatorname{keys}(\operatorname{ispace}(\operatorname{pk\text{-}inf}(p))) \wedge \\ \quad \operatorname{is\text{-}subterm}(\operatorname{tokens}(\operatorname{ispace}(\operatorname{pk\text{-}inf}(p))), \operatorname{keys}(\operatorname{ispace}(\operatorname{pk\text{-}inf}(p))), \operatorname{encr}(g, y), \\ \quad \quad \operatorname{node\text{-}content}(\operatorname{ispace}(\operatorname{pk\text{-}inf}(p)), n)) \Rightarrow \\ \quad \quad \operatorname{encr}(g, y) = \operatorname{node\text{-}content}(\operatorname{ispace}(\operatorname{pk\text{-}inf}(p)), n) \end{array}}}
\vnbprop{nsl-regular-list-subterm}{\vnbstmt{\begin{array}{@{}l@{}} \forall p, n, m. \\ \quad \operatorname{is\text{-}nsl\text{-}space}(p) \wedge \\ \quad n \in \operatorname{nodes}(\operatorname{ispace}(\operatorname{pk\text{-}inf}(p))) \wedge \\ \quad \operatorname{is\text{-}regular\text{-}node}(\operatorname{pk\text{-}inf}(p), n) \wedge \\ \quad m \in \operatorname{tuples}(\operatorname{msg\text{-}set}(\operatorname{tokens}(\operatorname{ispace}(\operatorname{pk\text{-}inf}(p))), \operatorname{keys}(\operatorname{ispace}(\operatorname{pk\text{-}inf}(p))))) \wedge \\ \quad \operatorname{is\text{-}subterm}(\operatorname{tokens}(\operatorname{ispace}(\operatorname{pk\text{-}inf}(p))), \operatorname{keys}(\operatorname{ispace}(\operatorname{pk\text{-}inf}(p))), \operatorname{lst}(m), \\ \quad \quad \operatorname{node\text{-}content}(\operatorname{ispace}(\operatorname{pk\text{-}inf}(p)), n)) \Rightarrow \\ \quad \quad m \in \operatorname{tuples}(\operatorname{tokens}(\operatorname{ispace}(\operatorname{pk\text{-}inf}(p)))) \end{array}}}
\vnbprop{nsl-regular-no-key}{\vnbstmt{\begin{array}{@{}l@{}} \forall p, n, y. \\ \quad \operatorname{is\text{-}nsl\text{-}space}(p) \wedge \\ \quad n \in \operatorname{nodes}(\operatorname{ispace}(\operatorname{pk\text{-}inf}(p))) \wedge \\ \quad \operatorname{is\text{-}regular\text{-}node}(\operatorname{pk\text{-}inf}(p), n) \wedge \\ \quad y \in \operatorname{keys}(\operatorname{ispace}(\operatorname{pk\text{-}inf}(p))) \Rightarrow \\ \quad \quad \neg \operatorname{is\text{-}subterm}(\operatorname{tokens}(\operatorname{ispace}(\operatorname{pk\text{-}inf}(p))), \operatorname{keys}(\operatorname{ispace}(\operatorname{pk\text{-}inf}(p))), y, \operatorname{node\text{-}content}(\operatorname{ispace}(\operatorname{pk\text{-}inf}(p)), n)) \end{array}}}
\vnbprop{adversary-encr-loss}{\vnbstmt{\begin{array}{@{}l@{}} \forall p, m, x, g, y. \\ \quad \operatorname{is\text{-}infiltrated\text{-}space}(p) \wedge \\ \quad m \in \operatorname{nodes}(\operatorname{ispace}(p)) \wedge \\ \quad \operatorname{is\text{-}adversary\text{-}node}(p, m) \wedge \\ \quad \operatorname{is\text{-}sent}(\operatorname{ispace}(p), m) \wedge \\ \quad \operatorname{same\text{-}strand\text{-}before}(\operatorname{ispace}(p), x, m) \wedge \\ \quad g \in \operatorname{msg\text{-}set}(\operatorname{tokens}(\operatorname{ispace}(p)), \operatorname{keys}(\operatorname{ispace}(p))) \wedge \\ \quad y \in \operatorname{keys}(\operatorname{ispace}(p)) \wedge \\ \quad \operatorname{is\text{-}subterm}(\operatorname{tokens}(\operatorname{ispace}(p)), \operatorname{keys}(\operatorname{ispace}(p)), \operatorname{encr}(g, y), \operatorname{node\text{-}content}(\operatorname{ispace}(p), x)) \wedge \\ \quad \neg \operatorname{is\text{-}subterm}(\operatorname{tokens}(\operatorname{ispace}(p)), \operatorname{keys}(\operatorname{ispace}(p)), \operatorname{encr}(g, y), \operatorname{node\text{-}content}(\operatorname{ispace}(p), m)) \Rightarrow \\ \quad \quad \operatorname{same\text{-}strand\text{-}before}(\operatorname{ispace}(p), \langle 1, \operatorname{nth}(2, m) \rangle, m) \wedge \\ \quad \quad \exists l \in \operatorname{tuples}(\operatorname{msg\text{-}set}(\operatorname{tokens}(\operatorname{ispace}(p)), \operatorname{keys}(\operatorname{ispace}(p)))). \\ \quad \quad \quad \operatorname{node\text{-}content}(\operatorname{ispace}(p), \langle 1, \operatorname{nth}(2, m) \rangle) = \operatorname{lst}(l) \wedge \\ \quad \quad \quad \operatorname{is\text{-}subterm}(\operatorname{tokens}(\operatorname{ispace}(p)), \operatorname{keys}(\operatorname{ispace}(p)), \operatorname{encr}(g, y), \operatorname{lst}(l)) \wedge \\ \quad \quad \quad \exists q \in \mathbb{N}. \\ \quad \quad \quad \quad (1 \leq q) \wedge ((q \leq \operatorname{length}(l)) \wedge (\operatorname{node\text{-}content}(\operatorname{ispace}(p), m) = \operatorname{nth}(q, l))) \vee \\ \quad \quad \operatorname{node\text{-}content}(\operatorname{ispace}(p), \langle 1, \operatorname{nth}(2, m) \rangle) = \operatorname{kinv}(\operatorname{ispace}(p))(y) \end{array}}}
\vnbprop{adversary-sent-source}{\vnbstmt{\begin{array}{@{}l@{}} \forall p, m, w. \\ \quad \operatorname{is\text{-}infiltrated\text{-}space}(p) \wedge \\ \quad m \in \operatorname{nodes}(\operatorname{ispace}(p)) \wedge \\ \quad \operatorname{is\text{-}adversary\text{-}node}(p, m) \wedge \\ \quad \operatorname{is\text{-}sent}(\operatorname{ispace}(p), m) \wedge \\ \quad \operatorname{is\text{-}subterm}(\operatorname{tokens}(\operatorname{ispace}(p)), \operatorname{keys}(\operatorname{ispace}(p)), w, \operatorname{node\text{-}content}(\operatorname{ispace}(p), m)) \Rightarrow \\ \quad \quad \exists x. \\ \quad \quad \quad \operatorname{same\text{-}strand\text{-}before}(\operatorname{ispace}(p), x, m) \wedge \\ \quad \quad \quad \operatorname{is\text{-}subterm}(\operatorname{tokens}(\operatorname{ispace}(p)), \operatorname{keys}(\operatorname{ispace}(p)), w, \operatorname{node\text{-}content}(\operatorname{ispace}(p), x)) \vee \\ \quad \quad w = \operatorname{node\text{-}content}(\operatorname{ispace}(p), m) \wedge \\ \quad \quad w \in \operatorname{adv\text{-}tokens}(p) \vee \\ \quad \quad \exists l \in \operatorname{tuples}(\operatorname{msg\text{-}set}(\operatorname{tokens}(\operatorname{ispace}(p)), \operatorname{keys}(\operatorname{ispace}(p)))). \\ \quad \quad \quad \operatorname{node\text{-}content}(\operatorname{ispace}(p), m) = \operatorname{lst}(l) \wedge \\ \quad \quad \quad \forall i \in \mathbb{N}. \\ \quad \quad \quad \quad ((1 \leq i) \wedge (i \leq \operatorname{length}(l))) \Rightarrow  \\ \quad \quad \quad \quad \quad \exists x. \\ \quad \quad \quad \quad \quad \quad \operatorname{same\text{-}strand\text{-}before}(\operatorname{ispace}(p), x, m) \wedge (\operatorname{node\text{-}content}(\operatorname{ispace}(p), x) = \operatorname{nth}(i, l)) \vee \\ \quad \quad \exists h.; \exists k \in \operatorname{keys}(\operatorname{ispace}(p)). \\ \quad \quad \quad \operatorname{node\text{-}content}(\operatorname{ispace}(p), m) = \operatorname{encr}(h, k) \wedge \\ \quad \quad \quad \exists x. \\ \quad \quad \quad \quad \operatorname{same\text{-}strand\text{-}before}(\operatorname{ispace}(p), x, m) \wedge (\operatorname{node\text{-}content}(\operatorname{ispace}(p), x) = h) \end{array}}}
\vnbprop{bundle-send-content}{\vnbstmt{\begin{array}{@{}l@{}} \forall c, d, a, b. \\ \quad (\operatorname{is\text{-}bundle}(c) \wedge ((\operatorname{bspace}(c) = d) \wedge (\langle a, b \rangle \in \operatorname{bsend}(c)))) \Rightarrow  \\ \quad \quad (a \in \operatorname{bnodes}(c)) \wedge ((b \in \operatorname{bnodes}(c)) \wedge (\operatorname{node\text{-}content}(d, a) = \operatorname{node\text{-}content}(d, b))) \end{array}}}
\vnbprop{bundle-minimal-sent-back}{\vnbstmt{\begin{array}{@{}l@{}} \forall c, s, n. \\ \quad \operatorname{is\text{-}bundle}(c) \wedge \\ \quad \forall z \in s.\; (z \in \operatorname{bnodes}(c)) \wedge \\ \quad \forall x, y.\; (((\langle x, y \rangle \in \operatorname{bsend}(c)) \wedge (y \in s)) \Rightarrow (x \in s)) \wedge \\ \quad n \in s \wedge \\ \quad \forall z \in s.\; \neg (\langle z, n \rangle \in \operatorname{precedes\text{-}strict}(c)) \Rightarrow \\ \quad \quad \operatorname{is\text{-}sent}(\operatorname{bspace}(c), n) \end{array}}}
\vnbprop{bundle-minimal-earlier}{\vnbstmt{\begin{array}{@{}l@{}} \forall c, a, s, u, b, n. \\ \quad \operatorname{is\text{-}bundle}(c) \wedge \\ \quad \operatorname{bspace}(c) = a \wedge \\ \quad \forall z \in s.\; (z \in \operatorname{bnodes}(c)) \wedge \\ \quad \operatorname{uniquely\text{-}originates}(a, u) \wedge \\ \quad \operatorname{originates}(a, u, b) \wedge \\ \quad \neg (b \in s) \wedge \\ \quad n \in s \wedge \\ \quad \operatorname{is\text{-}sent}(a, n) \wedge \\ \quad \operatorname{is\text{-}subterm}(\operatorname{tokens}(a), \operatorname{keys}(a), u, \operatorname{node\text{-}content}(a, n)) \Rightarrow \\ \quad \quad \exists x \in \operatorname{bnodes}(c). \\ \quad \quad \quad \langle x, n \rangle \in \operatorname{precedes\text{-}strict}(c) \wedge \\ \quad \quad \quad \operatorname{same\text{-}strand\text{-}before}(a, x, n) \wedge \\ \quad \quad \quad \operatorname{is\text{-}subterm}(\operatorname{tokens}(a), \operatorname{keys}(a), u, \operatorname{node\text{-}content}(a, x)) \end{array}}}
\vnbprop{nsl-key-secret}{\vnbstmt{\begin{array}{@{}l@{}} \forall p, c, y.; \forall n \in \operatorname{bnodes}(c). \\ \quad \operatorname{is\text{-}nsl\text{-}space}(p) \wedge \\ \quad \operatorname{is\text{-}bundle}(c) \wedge \\ \quad \operatorname{bspace}(c) = \operatorname{ispace}(\operatorname{pk\text{-}inf}(p)) \wedge \\ \quad y \in \operatorname{keys}(\operatorname{ispace}(\operatorname{pk\text{-}inf}(p))) \wedge \\ \quad \neg (y \in \operatorname{adv\text{-}tokens}(\operatorname{pk\text{-}inf}(p))) \Rightarrow \\ \quad \quad \neg \operatorname{is\text{-}subterm}(\operatorname{tokens}(\operatorname{ispace}(\operatorname{pk\text{-}inf}(p))), \operatorname{keys}(\operatorname{ispace}(\operatorname{pk\text{-}inf}(p))), y, \operatorname{node\text{-}content}(\operatorname{ispace}(\operatorname{pk\text{-}inf}(p)), n)) \end{array}}}
\vnbprop{nsl-list-source}{\vnbstmt{\begin{array}{@{}l@{}} \forall p, c, a, l, q. \\ \quad \operatorname{is\text{-}nsl\text{-}space}(p) \wedge \\ \quad \operatorname{is\text{-}bundle}(c) \wedge \\ \quad \operatorname{bspace}(c) = \operatorname{ispace}(\operatorname{pk\text{-}inf}(p)) \wedge \\ \quad a \in \operatorname{bnodes}(c) \wedge \\ \quad l \in \operatorname{tuples}(\operatorname{msg\text{-}set}(\operatorname{tokens}(\operatorname{ispace}(\operatorname{pk\text{-}inf}(p))), \operatorname{keys}(\operatorname{ispace}(\operatorname{pk\text{-}inf}(p))))) \wedge \\ \quad q \in \mathbb{N} \wedge \\ \quad 1 \leq q \wedge \\ \quad q \leq \operatorname{length}(l) \wedge \\ \quad \operatorname{node\text{-}content}(\operatorname{ispace}(\operatorname{pk\text{-}inf}(p)), a) = \operatorname{nth}(q, l) \wedge \\ \quad \neg (l \in \operatorname{tuples}(\operatorname{tokens}(\operatorname{ispace}(\operatorname{pk\text{-}inf}(p))))) \wedge \\ \quad \exists x \in \operatorname{bnodes}(c). \\ \quad \quad \langle x, a \rangle \in \operatorname{precedes\text{-}strict}(c) \wedge \\ \quad \quad \operatorname{is\text{-}subterm}(\operatorname{tokens}(\operatorname{ispace}(\operatorname{pk\text{-}inf}(p))), \operatorname{keys}(\operatorname{ispace}(\operatorname{pk\text{-}inf}(p))), \operatorname{lst}(l), \\ \quad \quad \quad \operatorname{node\text{-}content}(\operatorname{ispace}(\operatorname{pk\text{-}inf}(p)), x)) \Rightarrow \\ \quad \quad \exists z \in \operatorname{bnodes}(c). \\ \quad \quad \quad \langle z, a \rangle \in \operatorname{precedes\text{-}strict}(c) \wedge \\ \quad \quad \quad \operatorname{node\text{-}content}(\operatorname{ispace}(\operatorname{pk\text{-}inf}(p)), z) = \operatorname{node\text{-}content}(\operatorname{ispace}(\operatorname{pk\text{-}inf}(p)), a) \end{array}}}
\vnbprop{nsl-adversary-loss}{\vnbstmt{\begin{array}{@{}l@{}} \forall p, c, m, x, g, y. \\ \quad \operatorname{is\text{-}nsl\text{-}space}(p) \wedge \\ \quad \operatorname{is\text{-}bundle}(c) \wedge \\ \quad \operatorname{bspace}(c) = \operatorname{ispace}(\operatorname{pk\text{-}inf}(p)) \wedge \\ \quad m \in \operatorname{bnodes}(c) \wedge \\ \quad \operatorname{is\text{-}adversary\text{-}node}(\operatorname{pk\text{-}inf}(p), m) \wedge \\ \quad \operatorname{is\text{-}sent}(\operatorname{ispace}(\operatorname{pk\text{-}inf}(p)), m) \wedge \\ \quad \operatorname{same\text{-}strand\text{-}before}(\operatorname{ispace}(\operatorname{pk\text{-}inf}(p)), x, m) \wedge \\ \quad g \in \operatorname{msg\text{-}set}(\operatorname{tokens}(\operatorname{ispace}(\operatorname{pk\text{-}inf}(p))), \operatorname{keys}(\operatorname{ispace}(\operatorname{pk\text{-}inf}(p)))) \wedge \\ \quad y \in \operatorname{keys}(\operatorname{ispace}(\operatorname{pk\text{-}inf}(p))) \wedge \\ \quad \neg (\operatorname{kinv}(\operatorname{ispace}(\operatorname{pk\text{-}inf}(p)))(y) \in \operatorname{adv\text{-}tokens}(\operatorname{pk\text{-}inf}(p))) \wedge \\ \quad \operatorname{is\text{-}subterm}(\operatorname{tokens}(\operatorname{ispace}(\operatorname{pk\text{-}inf}(p))), \operatorname{keys}(\operatorname{ispace}(\operatorname{pk\text{-}inf}(p))), \operatorname{encr}(g, y), \\ \quad \quad \operatorname{node\text{-}content}(\operatorname{ispace}(\operatorname{pk\text{-}inf}(p)), x)) \wedge \\ \quad \neg \operatorname{is\text{-}subterm}(\operatorname{tokens}(\operatorname{ispace}(\operatorname{pk\text{-}inf}(p))), \operatorname{keys}(\operatorname{ispace}(\operatorname{pk\text{-}inf}(p))), \operatorname{encr}(g, y), \operatorname{node\text{-}content}(\operatorname{ispace}(\operatorname{pk\text{-}inf}(p)), m)) \Rightarrow \\ \quad \quad \exists z \in \operatorname{bnodes}(c). \\ \quad \quad \quad \langle z, m \rangle \in \operatorname{precedes\text{-}strict}(c) \wedge \\ \quad \quad \quad \operatorname{node\text{-}content}(\operatorname{ispace}(\operatorname{pk\text{-}inf}(p)), z) = \operatorname{node\text{-}content}(\operatorname{ispace}(\operatorname{pk\text{-}inf}(p)), m) \end{array}}}
\vnbprop{nsl-resp-nonce-originates}{\vnbstmt{\begin{array}{@{}l@{}} \forall p, s, a, b, c, d. \\ \quad \operatorname{is\text{-}pk\text{-}space}(p) \wedge \\ \quad s \in \operatorname{ids}(\operatorname{ispace}(\operatorname{pk\text{-}inf}(p))) \wedge \\ \quad \operatorname{is\text{-}resp\text{-}strand}(p, \operatorname{strand\text{-}map}(\operatorname{ispace}(\operatorname{pk\text{-}inf}(p)))(s), a, b, c, d) \wedge \\ \quad \neg (c = d) \Rightarrow \\ \quad \quad \operatorname{originates}(\operatorname{ispace}(\operatorname{pk\text{-}inf}(p)), d, \langle 2, s \rangle) \end{array}}}
\vnbprop{nsl-init-nonce-originates}{\vnbstmt{\begin{array}{@{}l@{}} \forall p, t, a, b, c, d. \\ \quad \operatorname{is\text{-}pk\text{-}space}(p) \wedge \\ \quad t \in \operatorname{ids}(\operatorname{ispace}(\operatorname{pk\text{-}inf}(p))) \wedge \\ \quad \operatorname{is\text{-}init\text{-}strand}(p, \operatorname{strand\text{-}map}(\operatorname{ispace}(\operatorname{pk\text{-}inf}(p)))(t), a, b, c, d) \Rightarrow \\ \quad \quad \operatorname{originates}(\operatorname{ispace}(\operatorname{pk\text{-}inf}(p)), c, \langle 1, t \rangle) \end{array}}}
\vnbprop{nsl-init-unique}{\vnbstmt{\begin{array}{@{}l@{}} \forall p, a, b, c, d, f, g, a_{2}, b_{2}, h. \\ \quad \operatorname{is\text{-}pk\text{-}space}(p) \wedge \\ \quad \operatorname{uniquely\text{-}originates}(\operatorname{ispace}(\operatorname{pk\text{-}inf}(p)), f) \wedge \\ \quad a \in \operatorname{ids}(\operatorname{ispace}(\operatorname{pk\text{-}inf}(p))) \wedge \\ \quad b \in \operatorname{ids}(\operatorname{ispace}(\operatorname{pk\text{-}inf}(p))) \wedge \\ \quad \operatorname{is\text{-}init\text{-}strand}(p, \operatorname{strand\text{-}map}(\operatorname{ispace}(\operatorname{pk\text{-}inf}(p)))(a), c, d, f, g) \wedge \\ \quad \operatorname{is\text{-}init\text{-}strand}(p, \operatorname{strand\text{-}map}(\operatorname{ispace}(\operatorname{pk\text{-}inf}(p)))(b), a_{2}, b_{2}, f, h) \Rightarrow \\ \quad \quad a = b \end{array}}}
\vnbprop{nsl-lemma-5-4}{\vnbstmt{\begin{array}{@{}l@{}} \forall p, c, s, a, b, d, f. \\ \quad \operatorname{is\text{-}nsl\text{-}space}(p) \wedge \\ \quad \operatorname{is\text{-}bundle}(c) \wedge \\ \quad \operatorname{bspace}(c) = \operatorname{ispace}(\operatorname{pk\text{-}inf}(p)) \wedge \\ \quad s \in \operatorname{ids}(\operatorname{ispace}(\operatorname{pk\text{-}inf}(p))) \wedge \\ \quad \operatorname{is\text{-}resp\text{-}strand}(p, \operatorname{strand\text{-}map}(\operatorname{ispace}(\operatorname{pk\text{-}inf}(p)))(s), a, b, d, f) \wedge \\ \quad \langle 3, s \rangle \in \operatorname{bnodes}(c) \wedge \\ \quad \neg (\operatorname{kinv}(\operatorname{ispace}(\operatorname{pk\text{-}inf}(p)))(\operatorname{pubkey}(p)(a)) \in \operatorname{adv\text{-}tokens}(\operatorname{pk\text{-}inf}(p))) \wedge \\ \quad \neg (d = f) \wedge \\ \quad \operatorname{uniquely\text{-}originates}(\operatorname{ispace}(\operatorname{pk\text{-}inf}(p)), f) \Rightarrow \\ \quad \quad \exists g \in \operatorname{bnodes}(c). \\ \quad \quad \quad \operatorname{is\text{-}subterm}(\operatorname{tokens}(\operatorname{ispace}(\operatorname{pk\text{-}inf}(p))), \operatorname{keys}(\operatorname{ispace}(\operatorname{pk\text{-}inf}(p))), f, \\ \quad \quad \quad \quad \operatorname{node\text{-}content}(\operatorname{ispace}(\operatorname{pk\text{-}inf}(p)), g)) \wedge \\ \quad \quad \quad \neg \operatorname{is\text{-}subterm}(\operatorname{tokens}(\operatorname{ispace}(\operatorname{pk\text{-}inf}(p))), \operatorname{keys}(\operatorname{ispace}(\operatorname{pk\text{-}inf}(p))), \operatorname{encr}(\operatorname{lst}(\langle d, f, b \rangle), \operatorname{pubkey}(p)(a)), \operatorname{node\text{-}content}(\operatorname{ispace}(\operatorname{pk\text{-}inf}(p)), g)) \wedge \\ \quad \quad \quad \forall z \in \operatorname{bnodes}(c). \\ \quad \quad \quad \quad \operatorname{is\text{-}subterm}(\operatorname{tokens}(\operatorname{ispace}(\operatorname{pk\text{-}inf}(p))), \operatorname{keys}(\operatorname{ispace}(\operatorname{pk\text{-}inf}(p))), f, \\ \quad \quad \quad \quad \quad \operatorname{node\text{-}content}(\operatorname{ispace}(\operatorname{pk\text{-}inf}(p)), z)) \wedge \\ \quad \quad \quad \quad \neg \operatorname{is\text{-}subterm}(\operatorname{tokens}(\operatorname{ispace}(\operatorname{pk\text{-}inf}(p))), \operatorname{keys}(\operatorname{ispace}(\operatorname{pk\text{-}inf}(p))), \operatorname{encr}(\operatorname{lst}(\langle d, f, b \rangle), \operatorname{pubkey}(p)(a)), \operatorname{node\text{-}content}(\operatorname{ispace}(\operatorname{pk\text{-}inf}(p)), z)) \Rightarrow \\ \quad \quad \quad \quad \quad \neg (\langle z, g \rangle \in \operatorname{precedes\text{-}strict}(c)) \wedge \\ \quad \quad \quad \operatorname{is\text{-}regular\text{-}node}(\operatorname{pk\text{-}inf}(p), g) \wedge \\ \quad \quad \quad \operatorname{is\text{-}sent}(\operatorname{ispace}(\operatorname{pk\text{-}inf}(p)), g) \end{array}}}
\vnbprop{nsl-lemma-5-6}{\vnbstmt{\begin{array}{@{}l@{}} \forall p, c, s, a, b, d, f.; \forall g \in \operatorname{bnodes}(c). \\ \quad \operatorname{is\text{-}nsl\text{-}space}(p) \wedge \\ \quad \operatorname{is\text{-}bundle}(c) \wedge \\ \quad \operatorname{bspace}(c) = \operatorname{ispace}(\operatorname{pk\text{-}inf}(p)) \wedge \\ \quad s \in \operatorname{ids}(\operatorname{ispace}(\operatorname{pk\text{-}inf}(p))) \wedge \\ \quad \operatorname{is\text{-}resp\text{-}strand}(p, \operatorname{strand\text{-}map}(\operatorname{ispace}(\operatorname{pk\text{-}inf}(p)))(s), a, b, d, f) \wedge \\ \quad \langle 3, s \rangle \in \operatorname{bnodes}(c) \wedge \\ \quad \neg (\operatorname{kinv}(\operatorname{ispace}(\operatorname{pk\text{-}inf}(p)))(\operatorname{pubkey}(p)(a)) \in \operatorname{adv\text{-}tokens}(\operatorname{pk\text{-}inf}(p))) \wedge \\ \quad \neg (d = f) \wedge \\ \quad \operatorname{uniquely\text{-}originates}(\operatorname{ispace}(\operatorname{pk\text{-}inf}(p)), f) \wedge \\ \quad \operatorname{is\text{-}subterm}(\operatorname{tokens}(\operatorname{ispace}(\operatorname{pk\text{-}inf}(p))), \operatorname{keys}(\operatorname{ispace}(\operatorname{pk\text{-}inf}(p))), f, \\ \quad \quad \operatorname{node\text{-}content}(\operatorname{ispace}(\operatorname{pk\text{-}inf}(p)), g)) \wedge \\ \quad \neg \operatorname{is\text{-}subterm}(\operatorname{tokens}(\operatorname{ispace}(\operatorname{pk\text{-}inf}(p))), \operatorname{keys}(\operatorname{ispace}(\operatorname{pk\text{-}inf}(p))), \operatorname{encr}(\operatorname{lst}(\langle d, f, b \rangle), \operatorname{pubkey}(p)(a)), \operatorname{node\text{-}content}(\operatorname{ispace}(\operatorname{pk\text{-}inf}(p)), g)) \wedge \\ \quad \forall z \in \operatorname{bnodes}(c). \\ \quad \quad \operatorname{is\text{-}subterm}(\operatorname{tokens}(\operatorname{ispace}(\operatorname{pk\text{-}inf}(p))), \operatorname{keys}(\operatorname{ispace}(\operatorname{pk\text{-}inf}(p))), f, \\ \quad \quad \quad \operatorname{node\text{-}content}(\operatorname{ispace}(\operatorname{pk\text{-}inf}(p)), z)) \wedge \\ \quad \quad \neg \operatorname{is\text{-}subterm}(\operatorname{tokens}(\operatorname{ispace}(\operatorname{pk\text{-}inf}(p))), \operatorname{keys}(\operatorname{ispace}(\operatorname{pk\text{-}inf}(p))), \operatorname{encr}(\operatorname{lst}(\langle d, f, b \rangle), \operatorname{pubkey}(p)(a)), \operatorname{node\text{-}content}(\operatorname{ispace}(\operatorname{pk\text{-}inf}(p)), z)) \Rightarrow \\ \quad \quad \quad \neg (\langle z, g \rangle \in \operatorname{precedes\text{-}strict}(c)) \wedge \\ \quad \operatorname{is\text{-}regular\text{-}node}(\operatorname{pk\text{-}inf}(p), g) \wedge \\ \quad \operatorname{is\text{-}sent}(\operatorname{ispace}(\operatorname{pk\text{-}inf}(p)), g) \Rightarrow \\ \quad \quad \exists h. \\ \quad \quad \quad \operatorname{same\text{-}strand\text{-}before}(\operatorname{ispace}(\operatorname{pk\text{-}inf}(p)), h, g) \wedge \\ \quad \quad \quad \operatorname{node\text{-}content}(\operatorname{ispace}(\operatorname{pk\text{-}inf}(p)), h) = \operatorname{encr}(\operatorname{lst}(\langle d, f, b \rangle), \operatorname{pubkey}(p)(a)) \end{array}}}
\vnbprop{nsl-lemma-5-7}{\vnbstmt{\begin{array}{@{}l@{}} \forall p, a, b, c, d, f, g. \\ \quad \operatorname{is\text{-}nsl\text{-}space}(p) \wedge \\ \quad c \in \operatorname{names}(p) \wedge \\ \quad b \in \operatorname{nodes}(\operatorname{ispace}(\operatorname{pk\text{-}inf}(p))) \wedge \\ \quad \operatorname{is\text{-}regular\text{-}node}(\operatorname{pk\text{-}inf}(p), b) \wedge \\ \quad \operatorname{is\text{-}sent}(\operatorname{ispace}(\operatorname{pk\text{-}inf}(p)), b) \wedge \\ \quad \operatorname{same\text{-}strand\text{-}before}(\operatorname{ispace}(\operatorname{pk\text{-}inf}(p)), a, b) \wedge \\ \quad \operatorname{node\text{-}content}(\operatorname{ispace}(\operatorname{pk\text{-}inf}(p)), a) = \operatorname{encr}(\operatorname{lst}(\langle f, g, d \rangle), \operatorname{pubkey}(p)(c)) \Rightarrow \\ \quad \quad (\operatorname{is\text{-}init\text{-}strand}(p, \operatorname{strand\text{-}map}(\operatorname{ispace}(\operatorname{pk\text{-}inf}(p)))(\operatorname{nth}(2, b)), c, d, f, g) \wedge (\operatorname{nth}(1, b) = 3)) \end{array}}}
\vnbprop{nsl-responder-agreement}{\vnbstmt{\begin{array}{@{}l@{}} \forall p, c, s, a, b, d, f. \\ \quad \operatorname{is\text{-}nsl\text{-}space}(p) \wedge \\ \quad \operatorname{is\text{-}bundle}(c) \wedge \\ \quad \operatorname{bspace}(c) = \operatorname{ispace}(\operatorname{pk\text{-}inf}(p)) \wedge \\ \quad s \in \operatorname{ids}(\operatorname{ispace}(\operatorname{pk\text{-}inf}(p))) \wedge \\ \quad \operatorname{is\text{-}resp\text{-}strand}(p, \operatorname{strand\text{-}map}(\operatorname{ispace}(\operatorname{pk\text{-}inf}(p)))(s), a, b, d, f) \wedge \\ \quad \langle 3, s \rangle \in \operatorname{bnodes}(c) \wedge \\ \quad \neg (\operatorname{kinv}(\operatorname{ispace}(\operatorname{pk\text{-}inf}(p)))(\operatorname{pubkey}(p)(a)) \in \operatorname{adv\text{-}tokens}(\operatorname{pk\text{-}inf}(p))) \wedge \\ \quad \neg (d = f) \wedge \\ \quad \operatorname{uniquely\text{-}originates}(\operatorname{ispace}(\operatorname{pk\text{-}inf}(p)), f) \Rightarrow \\ \quad \quad \exists t \in \operatorname{ids}(\operatorname{ispace}(\operatorname{pk\text{-}inf}(p))). \\ \quad \quad \quad (\operatorname{is\text{-}init\text{-}strand}(p, \operatorname{strand\text{-}map}(\operatorname{ispace}(\operatorname{pk\text{-}inf}(p)))(t), a, b, d, f) \wedge (\langle 3, t \rangle \in \operatorname{bnodes}(c))) \end{array}}}
\vnbprop{nsl-lemma-5-12}{\vnbstmt{\begin{array}{@{}l@{}} \forall p, c, s, a, b, d, f.; \forall g \in \operatorname{bnodes}(c). \\ \quad \operatorname{is\text{-}nsl\text{-}space}(p) \wedge \\ \quad \operatorname{is\text{-}bundle}(c) \wedge \\ \quad \operatorname{bspace}(c) = \operatorname{ispace}(\operatorname{pk\text{-}inf}(p)) \wedge \\ \quad s \in \operatorname{ids}(\operatorname{ispace}(\operatorname{pk\text{-}inf}(p))) \wedge \\ \quad \operatorname{is\text{-}resp\text{-}strand}(p, \operatorname{strand\text{-}map}(\operatorname{ispace}(\operatorname{pk\text{-}inf}(p)))(s), a, b, d, f) \wedge \\ \quad \neg (\operatorname{kinv}(\operatorname{ispace}(\operatorname{pk\text{-}inf}(p)))(\operatorname{pubkey}(p)(a)) \in \operatorname{adv\text{-}tokens}(\operatorname{pk\text{-}inf}(p))) \wedge \\ \quad \neg (\operatorname{kinv}(\operatorname{ispace}(\operatorname{pk\text{-}inf}(p)))(\operatorname{pubkey}(p)(b)) \in \operatorname{adv\text{-}tokens}(\operatorname{pk\text{-}inf}(p))) \wedge \\ \quad \neg (d = f) \wedge \\ \quad \operatorname{uniquely\text{-}originates}(\operatorname{ispace}(\operatorname{pk\text{-}inf}(p)), f) \wedge \\ \quad \operatorname{is\text{-}subterm}(\operatorname{tokens}(\operatorname{ispace}(\operatorname{pk\text{-}inf}(p))), \operatorname{keys}(\operatorname{ispace}(\operatorname{pk\text{-}inf}(p))), f, \\ \quad \quad \operatorname{node\text{-}content}(\operatorname{ispace}(\operatorname{pk\text{-}inf}(p)), g)) \wedge \\ \quad \neg \operatorname{is\text{-}subterm}(\operatorname{tokens}(\operatorname{ispace}(\operatorname{pk\text{-}inf}(p))), \operatorname{keys}(\operatorname{ispace}(\operatorname{pk\text{-}inf}(p))), \operatorname{encr}(\operatorname{lst}(\langle d, f, b \rangle), \operatorname{pubkey}(p)(a)), \operatorname{node\text{-}content}(\operatorname{ispace}(\operatorname{pk\text{-}inf}(p)), g)) \wedge \\ \quad \neg \operatorname{is\text{-}subterm}(\operatorname{tokens}(\operatorname{ispace}(\operatorname{pk\text{-}inf}(p))), \operatorname{keys}(\operatorname{ispace}(\operatorname{pk\text{-}inf}(p))), \operatorname{encr}(f, \operatorname{pubkey}(p)(b)), \operatorname{node\text{-}content}(\operatorname{ispace}(\operatorname{pk\text{-}inf}(p)), g)) \wedge \\ \quad \forall z \in \operatorname{bnodes}(c). \\ \quad \quad \operatorname{is\text{-}subterm}(\operatorname{tokens}(\operatorname{ispace}(\operatorname{pk\text{-}inf}(p))), \operatorname{keys}(\operatorname{ispace}(\operatorname{pk\text{-}inf}(p))), f, \\ \quad \quad \quad \operatorname{node\text{-}content}(\operatorname{ispace}(\operatorname{pk\text{-}inf}(p)), z)) \wedge \\ \quad \quad \neg \operatorname{is\text{-}subterm}(\operatorname{tokens}(\operatorname{ispace}(\operatorname{pk\text{-}inf}(p))), \operatorname{keys}(\operatorname{ispace}(\operatorname{pk\text{-}inf}(p))), \operatorname{encr}(\operatorname{lst}(\langle d, f, b \rangle), \operatorname{pubkey}(p)(a)), \operatorname{node\text{-}content}(\operatorname{ispace}(\operatorname{pk\text{-}inf}(p)), z)) \wedge \\ \quad \quad \neg \operatorname{is\text{-}subterm}(\operatorname{tokens}(\operatorname{ispace}(\operatorname{pk\text{-}inf}(p))), \operatorname{keys}(\operatorname{ispace}(\operatorname{pk\text{-}inf}(p))), \operatorname{encr}(f, \operatorname{pubkey}(p)(b)), \operatorname{node\text{-}content}(\operatorname{ispace}(\operatorname{pk\text{-}inf}(p)), z)) \Rightarrow \\ \quad \quad \quad \neg (\langle z, g \rangle \in \operatorname{precedes\text{-}strict}(c)) \Rightarrow \\ \quad \quad \neg \operatorname{is\text{-}adversary\text{-}node}(\operatorname{pk\text{-}inf}(p), g) \end{array}}}
\vnbprop{nsl-lemma-5-11}{\vnbstmt{\begin{array}{@{}l@{}} \forall p, c, s, a, b, d, f.; \forall g \in \operatorname{bnodes}(c). \\ \quad \operatorname{is\text{-}nsl\text{-}space}(p) \wedge \\ \quad \operatorname{is\text{-}bundle}(c) \wedge \\ \quad \operatorname{bspace}(c) = \operatorname{ispace}(\operatorname{pk\text{-}inf}(p)) \wedge \\ \quad s \in \operatorname{ids}(\operatorname{ispace}(\operatorname{pk\text{-}inf}(p))) \wedge \\ \quad \operatorname{is\text{-}resp\text{-}strand}(p, \operatorname{strand\text{-}map}(\operatorname{ispace}(\operatorname{pk\text{-}inf}(p)))(s), a, b, d, f) \wedge \\ \quad \neg (\operatorname{kinv}(\operatorname{ispace}(\operatorname{pk\text{-}inf}(p)))(\operatorname{pubkey}(p)(a)) \in \operatorname{adv\text{-}tokens}(\operatorname{pk\text{-}inf}(p))) \wedge \\ \quad \neg (\operatorname{kinv}(\operatorname{ispace}(\operatorname{pk\text{-}inf}(p)))(\operatorname{pubkey}(p)(b)) \in \operatorname{adv\text{-}tokens}(\operatorname{pk\text{-}inf}(p))) \wedge \\ \quad \neg (d = f) \wedge \\ \quad \operatorname{uniquely\text{-}originates}(\operatorname{ispace}(\operatorname{pk\text{-}inf}(p)), f) \wedge \\ \quad \operatorname{is\text{-}subterm}(\operatorname{tokens}(\operatorname{ispace}(\operatorname{pk\text{-}inf}(p))), \operatorname{keys}(\operatorname{ispace}(\operatorname{pk\text{-}inf}(p))), f, \\ \quad \quad \operatorname{node\text{-}content}(\operatorname{ispace}(\operatorname{pk\text{-}inf}(p)), g)) \wedge \\ \quad \neg \operatorname{is\text{-}subterm}(\operatorname{tokens}(\operatorname{ispace}(\operatorname{pk\text{-}inf}(p))), \operatorname{keys}(\operatorname{ispace}(\operatorname{pk\text{-}inf}(p))), \operatorname{encr}(\operatorname{lst}(\langle d, f, b \rangle), \operatorname{pubkey}(p)(a)), \operatorname{node\text{-}content}(\operatorname{ispace}(\operatorname{pk\text{-}inf}(p)), g)) \wedge \\ \quad \neg \operatorname{is\text{-}subterm}(\operatorname{tokens}(\operatorname{ispace}(\operatorname{pk\text{-}inf}(p))), \operatorname{keys}(\operatorname{ispace}(\operatorname{pk\text{-}inf}(p))), \operatorname{encr}(f, \operatorname{pubkey}(p)(b)), \operatorname{node\text{-}content}(\operatorname{ispace}(\operatorname{pk\text{-}inf}(p)), g)) \wedge \\ \quad \forall z \in \operatorname{bnodes}(c). \\ \quad \quad \operatorname{is\text{-}subterm}(\operatorname{tokens}(\operatorname{ispace}(\operatorname{pk\text{-}inf}(p))), \operatorname{keys}(\operatorname{ispace}(\operatorname{pk\text{-}inf}(p))), f, \\ \quad \quad \quad \operatorname{node\text{-}content}(\operatorname{ispace}(\operatorname{pk\text{-}inf}(p)), z)) \wedge \\ \quad \quad \neg \operatorname{is\text{-}subterm}(\operatorname{tokens}(\operatorname{ispace}(\operatorname{pk\text{-}inf}(p))), \operatorname{keys}(\operatorname{ispace}(\operatorname{pk\text{-}inf}(p))), \operatorname{encr}(\operatorname{lst}(\langle d, f, b \rangle), \operatorname{pubkey}(p)(a)), \operatorname{node\text{-}content}(\operatorname{ispace}(\operatorname{pk\text{-}inf}(p)), z)) \wedge \\ \quad \quad \neg \operatorname{is\text{-}subterm}(\operatorname{tokens}(\operatorname{ispace}(\operatorname{pk\text{-}inf}(p))), \operatorname{keys}(\operatorname{ispace}(\operatorname{pk\text{-}inf}(p))), \operatorname{encr}(f, \operatorname{pubkey}(p)(b)), \operatorname{node\text{-}content}(\operatorname{ispace}(\operatorname{pk\text{-}inf}(p)), z)) \Rightarrow \\ \quad \quad \quad \neg (\langle z, g \rangle \in \operatorname{precedes\text{-}strict}(c)) \Rightarrow \\ \quad \quad \neg \operatorname{is\text{-}regular\text{-}node}(\operatorname{pk\text{-}inf}(p), g) \end{array}}}
\vnbprop{nsl-responder-secrecy}{\vnbstmt{\begin{array}{@{}l@{}} \forall p, c, s, a, b, d, f.; \forall m \in \operatorname{bnodes}(c). \\ \quad \operatorname{is\text{-}nsl\text{-}space}(p) \wedge \\ \quad \operatorname{is\text{-}bundle}(c) \wedge \\ \quad \operatorname{bspace}(c) = \operatorname{ispace}(\operatorname{pk\text{-}inf}(p)) \wedge \\ \quad s \in \operatorname{ids}(\operatorname{ispace}(\operatorname{pk\text{-}inf}(p))) \wedge \\ \quad \operatorname{is\text{-}resp\text{-}strand}(p, \operatorname{strand\text{-}map}(\operatorname{ispace}(\operatorname{pk\text{-}inf}(p)))(s), a, b, d, f) \wedge \\ \quad \neg (\operatorname{kinv}(\operatorname{ispace}(\operatorname{pk\text{-}inf}(p)))(\operatorname{pubkey}(p)(a)) \in \operatorname{adv\text{-}tokens}(\operatorname{pk\text{-}inf}(p))) \wedge \\ \quad \neg (\operatorname{kinv}(\operatorname{ispace}(\operatorname{pk\text{-}inf}(p)))(\operatorname{pubkey}(p)(b)) \in \operatorname{adv\text{-}tokens}(\operatorname{pk\text{-}inf}(p))) \wedge \\ \quad \neg (d = f) \wedge \\ \quad \operatorname{uniquely\text{-}originates}(\operatorname{ispace}(\operatorname{pk\text{-}inf}(p)), f) \wedge \\ \quad \operatorname{is\text{-}subterm}(\operatorname{tokens}(\operatorname{ispace}(\operatorname{pk\text{-}inf}(p))), \operatorname{keys}(\operatorname{ispace}(\operatorname{pk\text{-}inf}(p))), f, \\ \quad \quad \operatorname{node\text{-}content}(\operatorname{ispace}(\operatorname{pk\text{-}inf}(p)), m)) \Rightarrow \\ \quad \quad \operatorname{is\text{-}subterm}(\operatorname{tokens}(\operatorname{ispace}(\operatorname{pk\text{-}inf}(p))), \operatorname{keys}(\operatorname{ispace}(\operatorname{pk\text{-}inf}(p))), \\ \quad \quad \quad \operatorname{encr}(\operatorname{lst}(\langle d, f, b \rangle), \operatorname{pubkey}(p)(a)), \\ \quad \quad \quad \operatorname{node\text{-}content}(\operatorname{ispace}(\operatorname{pk\text{-}inf}(p)), m)) \vee \\ \quad \quad \operatorname{is\text{-}subterm}(\operatorname{tokens}(\operatorname{ispace}(\operatorname{pk\text{-}inf}(p))), \operatorname{keys}(\operatorname{ispace}(\operatorname{pk\text{-}inf}(p))), \\ \quad \quad \quad \operatorname{encr}(f, \operatorname{pubkey}(p)(b)), \operatorname{node\text{-}content}(\operatorname{ispace}(\operatorname{pk\text{-}inf}(p)), m)) \end{array}}}
\vnbprop{nsl-responder-nonce-not-sent}{\vnbstmt{\begin{array}{@{}l@{}} \forall p, c, s, a, b, d, f.; \forall m \in \operatorname{bnodes}(c). \\ \quad \operatorname{is\text{-}nsl\text{-}space}(p) \wedge \\ \quad \operatorname{is\text{-}bundle}(c) \wedge \\ \quad \operatorname{bspace}(c) = \operatorname{ispace}(\operatorname{pk\text{-}inf}(p)) \wedge \\ \quad s \in \operatorname{ids}(\operatorname{ispace}(\operatorname{pk\text{-}inf}(p))) \wedge \\ \quad \operatorname{is\text{-}resp\text{-}strand}(p, \operatorname{strand\text{-}map}(\operatorname{ispace}(\operatorname{pk\text{-}inf}(p)))(s), a, b, d, f) \wedge \\ \quad \neg (\operatorname{kinv}(\operatorname{ispace}(\operatorname{pk\text{-}inf}(p)))(\operatorname{pubkey}(p)(a)) \in \operatorname{adv\text{-}tokens}(\operatorname{pk\text{-}inf}(p))) \wedge \\ \quad \neg (\operatorname{kinv}(\operatorname{ispace}(\operatorname{pk\text{-}inf}(p)))(\operatorname{pubkey}(p)(b)) \in \operatorname{adv\text{-}tokens}(\operatorname{pk\text{-}inf}(p))) \wedge \\ \quad \neg (d = f) \wedge \\ \quad \operatorname{uniquely\text{-}originates}(\operatorname{ispace}(\operatorname{pk\text{-}inf}(p)), f) \Rightarrow \\ \quad \quad \neg (\operatorname{node\text{-}content}(\operatorname{ispace}(\operatorname{pk\text{-}inf}(p)), m) = f) \end{array}}}
\vnbprop{nsl-initiator-secrecy-adversary}{\vnbstmt{\begin{array}{@{}l@{}} \forall p, c, s, a, b, d, f.; \forall g \in \operatorname{bnodes}(c). \\ \quad \operatorname{is\text{-}nsl\text{-}space}(p) \wedge \\ \quad \operatorname{is\text{-}bundle}(c) \wedge \\ \quad \operatorname{bspace}(c) = \operatorname{ispace}(\operatorname{pk\text{-}inf}(p)) \wedge \\ \quad s \in \operatorname{ids}(\operatorname{ispace}(\operatorname{pk\text{-}inf}(p))) \wedge \\ \quad \operatorname{is\text{-}init\text{-}strand}(p, \operatorname{strand\text{-}map}(\operatorname{ispace}(\operatorname{pk\text{-}inf}(p)))(s), a, b, d, f) \wedge \\ \quad \langle 3, s \rangle \in \operatorname{bnodes}(c) \wedge \\ \quad \neg (\operatorname{kinv}(\operatorname{ispace}(\operatorname{pk\text{-}inf}(p)))(\operatorname{pubkey}(p)(a)) \in \operatorname{adv\text{-}tokens}(\operatorname{pk\text{-}inf}(p))) \wedge \\ \quad \neg (\operatorname{kinv}(\operatorname{ispace}(\operatorname{pk\text{-}inf}(p)))(\operatorname{pubkey}(p)(b)) \in \operatorname{adv\text{-}tokens}(\operatorname{pk\text{-}inf}(p))) \wedge \\ \quad \neg (d = f) \wedge \\ \quad \operatorname{uniquely\text{-}originates}(\operatorname{ispace}(\operatorname{pk\text{-}inf}(p)), d) \wedge \\ \quad \operatorname{is\text{-}subterm}(\operatorname{tokens}(\operatorname{ispace}(\operatorname{pk\text{-}inf}(p))), \operatorname{keys}(\operatorname{ispace}(\operatorname{pk\text{-}inf}(p))), d, \\ \quad \quad \operatorname{node\text{-}content}(\operatorname{ispace}(\operatorname{pk\text{-}inf}(p)), g)) \wedge \\ \quad \neg \operatorname{is\text{-}subterm}(\operatorname{tokens}(\operatorname{ispace}(\operatorname{pk\text{-}inf}(p))), \operatorname{keys}(\operatorname{ispace}(\operatorname{pk\text{-}inf}(p))), \operatorname{encr}(\operatorname{lst}(\langle d, a \rangle), \operatorname{pubkey}(p)(b)), \operatorname{node\text{-}content}(\operatorname{ispace}(\operatorname{pk\text{-}inf}(p)), g)) \wedge \\ \quad \forall w \in \operatorname{tokens}(\operatorname{ispace}(\operatorname{pk\text{-}inf}(p))). \\ \quad \quad \neg \operatorname{is\text{-}subterm}(\operatorname{tokens}(\operatorname{ispace}(\operatorname{pk\text{-}inf}(p))), \operatorname{keys}(\operatorname{ispace}(\operatorname{pk\text{-}inf}(p))), \operatorname{encr}(\operatorname{lst}(\langle d, w, b \rangle), \operatorname{pubkey}(p)(a)), \operatorname{node\text{-}content}(\operatorname{ispace}(\operatorname{pk\text{-}inf}(p)), g)) \wedge \\ \quad \forall z \in \operatorname{bnodes}(c). \\ \quad \quad \operatorname{is\text{-}subterm}(\operatorname{tokens}(\operatorname{ispace}(\operatorname{pk\text{-}inf}(p))), \operatorname{keys}(\operatorname{ispace}(\operatorname{pk\text{-}inf}(p))), d, \\ \quad \quad \quad \operatorname{node\text{-}content}(\operatorname{ispace}(\operatorname{pk\text{-}inf}(p)), z)) \wedge \\ \quad \quad \neg \operatorname{is\text{-}subterm}(\operatorname{tokens}(\operatorname{ispace}(\operatorname{pk\text{-}inf}(p))), \operatorname{keys}(\operatorname{ispace}(\operatorname{pk\text{-}inf}(p))), \operatorname{encr}(\operatorname{lst}(\langle d, a \rangle), \operatorname{pubkey}(p)(b)), \operatorname{node\text{-}content}(\operatorname{ispace}(\operatorname{pk\text{-}inf}(p)), z)) \wedge \\ \quad \quad \forall w \in \operatorname{tokens}(\operatorname{ispace}(\operatorname{pk\text{-}inf}(p))). \\ \quad \quad \quad \neg \operatorname{is\text{-}subterm}(\operatorname{tokens}(\operatorname{ispace}(\operatorname{pk\text{-}inf}(p))), \operatorname{keys}(\operatorname{ispace}(\operatorname{pk\text{-}inf}(p))), \operatorname{encr}(\operatorname{lst}(\langle d, w, b \rangle), \operatorname{pubkey}(p)(a)), \operatorname{node\text{-}content}(\operatorname{ispace}(\operatorname{pk\text{-}inf}(p)), z)) \Rightarrow \\ \quad \quad \quad \neg (\langle z, g \rangle \in \operatorname{precedes\text{-}strict}(c)) \Rightarrow \\ \quad \quad \neg \operatorname{is\text{-}adversary\text{-}node}(\operatorname{pk\text{-}inf}(p), g) \end{array}}}
\vnbprop{nsl-initiator-secrecy-regular}{\vnbstmt{\begin{array}{@{}l@{}} \forall p, c, s, a, b, d, f.; \forall g \in \operatorname{bnodes}(c). \\ \quad \operatorname{is\text{-}nsl\text{-}space}(p) \wedge \\ \quad \operatorname{is\text{-}bundle}(c) \wedge \\ \quad \operatorname{bspace}(c) = \operatorname{ispace}(\operatorname{pk\text{-}inf}(p)) \wedge \\ \quad s \in \operatorname{ids}(\operatorname{ispace}(\operatorname{pk\text{-}inf}(p))) \wedge \\ \quad \operatorname{is\text{-}init\text{-}strand}(p, \operatorname{strand\text{-}map}(\operatorname{ispace}(\operatorname{pk\text{-}inf}(p)))(s), a, b, d, f) \wedge \\ \quad \langle 3, s \rangle \in \operatorname{bnodes}(c) \wedge \\ \quad \neg (\operatorname{kinv}(\operatorname{ispace}(\operatorname{pk\text{-}inf}(p)))(\operatorname{pubkey}(p)(a)) \in \operatorname{adv\text{-}tokens}(\operatorname{pk\text{-}inf}(p))) \wedge \\ \quad \neg (\operatorname{kinv}(\operatorname{ispace}(\operatorname{pk\text{-}inf}(p)))(\operatorname{pubkey}(p)(b)) \in \operatorname{adv\text{-}tokens}(\operatorname{pk\text{-}inf}(p))) \wedge \\ \quad \neg (d = f) \wedge \\ \quad \operatorname{uniquely\text{-}originates}(\operatorname{ispace}(\operatorname{pk\text{-}inf}(p)), d) \wedge \\ \quad \operatorname{is\text{-}subterm}(\operatorname{tokens}(\operatorname{ispace}(\operatorname{pk\text{-}inf}(p))), \operatorname{keys}(\operatorname{ispace}(\operatorname{pk\text{-}inf}(p))), d, \\ \quad \quad \operatorname{node\text{-}content}(\operatorname{ispace}(\operatorname{pk\text{-}inf}(p)), g)) \wedge \\ \quad \neg \operatorname{is\text{-}subterm}(\operatorname{tokens}(\operatorname{ispace}(\operatorname{pk\text{-}inf}(p))), \operatorname{keys}(\operatorname{ispace}(\operatorname{pk\text{-}inf}(p))), \operatorname{encr}(\operatorname{lst}(\langle d, a \rangle), \operatorname{pubkey}(p)(b)), \operatorname{node\text{-}content}(\operatorname{ispace}(\operatorname{pk\text{-}inf}(p)), g)) \wedge \\ \quad \forall w \in \operatorname{tokens}(\operatorname{ispace}(\operatorname{pk\text{-}inf}(p))). \\ \quad \quad \neg \operatorname{is\text{-}subterm}(\operatorname{tokens}(\operatorname{ispace}(\operatorname{pk\text{-}inf}(p))), \operatorname{keys}(\operatorname{ispace}(\operatorname{pk\text{-}inf}(p))), \operatorname{encr}(\operatorname{lst}(\langle d, w, b \rangle), \operatorname{pubkey}(p)(a)), \operatorname{node\text{-}content}(\operatorname{ispace}(\operatorname{pk\text{-}inf}(p)), g)) \wedge \\ \quad \forall z \in \operatorname{bnodes}(c). \\ \quad \quad \operatorname{is\text{-}subterm}(\operatorname{tokens}(\operatorname{ispace}(\operatorname{pk\text{-}inf}(p))), \operatorname{keys}(\operatorname{ispace}(\operatorname{pk\text{-}inf}(p))), d, \\ \quad \quad \quad \operatorname{node\text{-}content}(\operatorname{ispace}(\operatorname{pk\text{-}inf}(p)), z)) \wedge \\ \quad \quad \neg \operatorname{is\text{-}subterm}(\operatorname{tokens}(\operatorname{ispace}(\operatorname{pk\text{-}inf}(p))), \operatorname{keys}(\operatorname{ispace}(\operatorname{pk\text{-}inf}(p))), \operatorname{encr}(\operatorname{lst}(\langle d, a \rangle), \operatorname{pubkey}(p)(b)), \operatorname{node\text{-}content}(\operatorname{ispace}(\operatorname{pk\text{-}inf}(p)), z)) \wedge \\ \quad \quad \forall w \in \operatorname{tokens}(\operatorname{ispace}(\operatorname{pk\text{-}inf}(p))). \\ \quad \quad \quad \neg \operatorname{is\text{-}subterm}(\operatorname{tokens}(\operatorname{ispace}(\operatorname{pk\text{-}inf}(p))), \operatorname{keys}(\operatorname{ispace}(\operatorname{pk\text{-}inf}(p))), \operatorname{encr}(\operatorname{lst}(\langle d, w, b \rangle), \operatorname{pubkey}(p)(a)), \operatorname{node\text{-}content}(\operatorname{ispace}(\operatorname{pk\text{-}inf}(p)), z)) \Rightarrow \\ \quad \quad \quad \neg (\langle z, g \rangle \in \operatorname{precedes\text{-}strict}(c)) \Rightarrow \\ \quad \quad \neg \operatorname{is\text{-}regular\text{-}node}(\operatorname{pk\text{-}inf}(p), g) \end{array}}}
\vnbprop{nsl-initiator-secrecy}{\vnbstmt{\begin{array}{@{}l@{}} \forall p, c, s, a, b, d, f.; \forall m \in \operatorname{bnodes}(c). \\ \quad \operatorname{is\text{-}nsl\text{-}space}(p) \wedge \\ \quad \operatorname{is\text{-}bundle}(c) \wedge \\ \quad \operatorname{bspace}(c) = \operatorname{ispace}(\operatorname{pk\text{-}inf}(p)) \wedge \\ \quad s \in \operatorname{ids}(\operatorname{ispace}(\operatorname{pk\text{-}inf}(p))) \wedge \\ \quad \operatorname{is\text{-}init\text{-}strand}(p, \operatorname{strand\text{-}map}(\operatorname{ispace}(\operatorname{pk\text{-}inf}(p)))(s), a, b, d, f) \wedge \\ \quad \langle 3, s \rangle \in \operatorname{bnodes}(c) \wedge \\ \quad \neg (\operatorname{kinv}(\operatorname{ispace}(\operatorname{pk\text{-}inf}(p)))(\operatorname{pubkey}(p)(a)) \in \operatorname{adv\text{-}tokens}(\operatorname{pk\text{-}inf}(p))) \wedge \\ \quad \neg (\operatorname{kinv}(\operatorname{ispace}(\operatorname{pk\text{-}inf}(p)))(\operatorname{pubkey}(p)(b)) \in \operatorname{adv\text{-}tokens}(\operatorname{pk\text{-}inf}(p))) \wedge \\ \quad \neg (d = f) \wedge \\ \quad \operatorname{uniquely\text{-}originates}(\operatorname{ispace}(\operatorname{pk\text{-}inf}(p)), d) \wedge \\ \quad \operatorname{is\text{-}subterm}(\operatorname{tokens}(\operatorname{ispace}(\operatorname{pk\text{-}inf}(p))), \operatorname{keys}(\operatorname{ispace}(\operatorname{pk\text{-}inf}(p))), d, \\ \quad \quad \operatorname{node\text{-}content}(\operatorname{ispace}(\operatorname{pk\text{-}inf}(p)), m)) \Rightarrow \\ \quad \quad \operatorname{is\text{-}subterm}(\operatorname{tokens}(\operatorname{ispace}(\operatorname{pk\text{-}inf}(p))), \operatorname{keys}(\operatorname{ispace}(\operatorname{pk\text{-}inf}(p))), \\ \quad \quad \quad \operatorname{encr}(\operatorname{lst}(\langle d, a \rangle), \operatorname{pubkey}(p)(b)), \\ \quad \quad \quad \operatorname{node\text{-}content}(\operatorname{ispace}(\operatorname{pk\text{-}inf}(p)), m)) \vee \\ \quad \quad \exists w \in \operatorname{tokens}(\operatorname{ispace}(\operatorname{pk\text{-}inf}(p))). \\ \quad \quad \quad \operatorname{is\text{-}subterm}(\operatorname{tokens}(\operatorname{ispace}(\operatorname{pk\text{-}inf}(p))), \operatorname{keys}(\operatorname{ispace}(\operatorname{pk\text{-}inf}(p))), \\ \quad \quad \quad \quad \operatorname{encr}(\operatorname{lst}(\langle d, w, b \rangle), \operatorname{pubkey}(p)(a)), \\ \quad \quad \quad \quad \operatorname{node\text{-}content}(\operatorname{ispace}(\operatorname{pk\text{-}inf}(p)), m)) \end{array}}}
\vnbprop{nsl-initiator-nonce-not-sent}{\vnbstmt{\begin{array}{@{}l@{}} \forall p, c, s, a, b, d, f.; \forall m \in \operatorname{bnodes}(c). \\ \quad \operatorname{is\text{-}nsl\text{-}space}(p) \wedge \\ \quad \operatorname{is\text{-}bundle}(c) \wedge \\ \quad \operatorname{bspace}(c) = \operatorname{ispace}(\operatorname{pk\text{-}inf}(p)) \wedge \\ \quad s \in \operatorname{ids}(\operatorname{ispace}(\operatorname{pk\text{-}inf}(p))) \wedge \\ \quad \operatorname{is\text{-}init\text{-}strand}(p, \operatorname{strand\text{-}map}(\operatorname{ispace}(\operatorname{pk\text{-}inf}(p)))(s), a, b, d, f) \wedge \\ \quad \langle 3, s \rangle \in \operatorname{bnodes}(c) \wedge \\ \quad \neg (\operatorname{kinv}(\operatorname{ispace}(\operatorname{pk\text{-}inf}(p)))(\operatorname{pubkey}(p)(a)) \in \operatorname{adv\text{-}tokens}(\operatorname{pk\text{-}inf}(p))) \wedge \\ \quad \neg (\operatorname{kinv}(\operatorname{ispace}(\operatorname{pk\text{-}inf}(p)))(\operatorname{pubkey}(p)(b)) \in \operatorname{adv\text{-}tokens}(\operatorname{pk\text{-}inf}(p))) \wedge \\ \quad \neg (d = f) \wedge \\ \quad \operatorname{uniquely\text{-}originates}(\operatorname{ispace}(\operatorname{pk\text{-}inf}(p)), d) \Rightarrow \\ \quad \quad \neg (\operatorname{node\text{-}content}(\operatorname{ispace}(\operatorname{pk\text{-}inf}(p)), m) = d) \end{array}}}
\vnbprop{nsl-initiator-agreement}{\vnbstmt{\begin{array}{@{}l@{}} \forall p, c, s, a, b, d, f. \\ \quad \operatorname{is\text{-}nsl\text{-}space}(p) \wedge \\ \quad \operatorname{is\text{-}bundle}(c) \wedge \\ \quad \operatorname{bspace}(c) = \operatorname{ispace}(\operatorname{pk\text{-}inf}(p)) \wedge \\ \quad s \in \operatorname{ids}(\operatorname{ispace}(\operatorname{pk\text{-}inf}(p))) \wedge \\ \quad \operatorname{is\text{-}init\text{-}strand}(p, \operatorname{strand\text{-}map}(\operatorname{ispace}(\operatorname{pk\text{-}inf}(p)))(s), a, b, d, f) \wedge \\ \quad \langle 3, s \rangle \in \operatorname{bnodes}(c) \wedge \\ \quad \neg (\operatorname{kinv}(\operatorname{ispace}(\operatorname{pk\text{-}inf}(p)))(\operatorname{pubkey}(p)(a)) \in \operatorname{adv\text{-}tokens}(\operatorname{pk\text{-}inf}(p))) \wedge \\ \quad \neg (\operatorname{kinv}(\operatorname{ispace}(\operatorname{pk\text{-}inf}(p)))(\operatorname{pubkey}(p)(b)) \in \operatorname{adv\text{-}tokens}(\operatorname{pk\text{-}inf}(p))) \wedge \\ \quad \neg (d = f) \wedge \\ \quad \operatorname{uniquely\text{-}originates}(\operatorname{ispace}(\operatorname{pk\text{-}inf}(p)), d) \Rightarrow \\ \quad \quad \exists t \in \operatorname{ids}(\operatorname{ispace}(\operatorname{pk\text{-}inf}(p))). \\ \quad \quad \quad (\operatorname{is\text{-}resp\text{-}strand}(p, \operatorname{strand\text{-}map}(\operatorname{ispace}(\operatorname{pk\text{-}inf}(p)))(t), a, b, d, f) \wedge (\langle 2, t \rangle \in \operatorname{bnodes}(c))) \end{array}}}

\part{The Kernel Encoded}
\noindent Expressions of the VNB language encoded as sets (syntax, free variables, substitution, rank); the core calculus as an inductively defined set DERIVABLE(T) of sequents; contextual rewriting and the rewriting rules as admissible rules of the core; the induction rules and the definedness clauses of the kernel against the core (K2, docs/meta-k2-admissibility-2026-10-08.md).
\par\noindent\textit{Builds on:} Foundations.
\par\noindent{\small 10 files.}

\vnbbookfile{structure-library/meta-syntax}{structure-library/meta-syntax}
\noindent{\small meta-syntax.scm -- THE SYNTAX OF THE KERNEL AS A VNB OBJECT: locally nameless expressions, their stages, and the four structural operations FV, SUBST, OPEN, LC. (docs/meta-k2-admissibility-2026-10-08.md, sections 5 and 6, phase 1: the cost probe.) DEFINITIONS ONLY; the laws are in theorem-library/meta-syntax-laws.scm.}\par
\par\medskip\noindent\textbf{Definition} (\texttt{efvar}).\;$\operatorname{efvar}(\operatorname{mxn\_}) :=$ $\langle 0, \operatorname{mxn\_} \rangle$
\par\medskip\noindent\textbf{Definition} (\texttt{ebvar}).\;$\operatorname{ebvar}(\operatorname{mxk\_}) :=$ $\langle 1, \operatorname{mxk\_} \rangle$
\par\medskip\noindent\textbf{Definition} (\texttt{econst}).\;$\operatorname{econst}(\operatorname{mxc\_}) :=$ $\langle 2, \operatorname{mxc\_} \rangle$
\par\medskip\noindent\textbf{Definition} (\texttt{eapp}).\;$\operatorname{eapp}(\operatorname{mxh\_}, \operatorname{mxl\_}) :=$ $\langle 3, \operatorname{mxh\_}, \operatorname{mxl\_} \rangle$
\par\medskip\noindent\textbf{Definition} (\texttt{ebind}).\;$\operatorname{ebind}(\operatorname{mxh\_}, \operatorname{mxn\_}, \operatorname{mxa\_}, \operatorname{mxb\_}) :=$ $\langle 4, \operatorname{mxh\_}, \operatorname{mxn\_}, \operatorname{mxa\_}, \operatorname{mxb\_} \rangle$
\par\medskip\noindent\textbf{Definition} (\texttt{expr-fv}).\;$\operatorname{expr\text{-}fv}(\operatorname{mxe\_}) :=$ $\operatorname{fv\text{-}stage}(\operatorname{expr\text{-}rank}(\operatorname{mxe\_}))(\operatorname{mxe\_})$
\par\medskip\noindent\textbf{Definition} (\texttt{expr-subst}).\;$\operatorname{expr\text{-}subst}(\operatorname{mxe\_}, \operatorname{mxm\_}, \operatorname{mxu\_}) :=$ $\operatorname{subst\text{-}stage}(\operatorname{mxm\_}, \operatorname{mxu\_}, \operatorname{expr\text{-}rank}(\operatorname{mxe\_}))(\operatorname{mxe\_})$
\par\medskip\noindent\textbf{Definition} (\texttt{expr-open}).\;$\operatorname{expr\text{-}open}(\operatorname{mxe\_}, \operatorname{mxk\_}, \operatorname{mxu\_}) :=$ $\operatorname{open\text{-}stage}(\operatorname{mxu\_}, \operatorname{expr\text{-}rank}(\operatorname{mxe\_}))(\operatorname{mxe\_}, \operatorname{mxk\_})$
\par\medskip\noindent\textbf{Definition} (\texttt{expr-lc}).\;For every $\operatorname{mxe\_}$,  $\operatorname{expr\text{-}lc}(\operatorname{mxe\_}) \Leftrightarrow (0 \in \operatorname{expr\text{-}lcd}(\operatorname{mxe\_}))$
\par\medskip\noindent\textbf{Definition} (\texttt{expr-set}).\; $(\operatorname{expr\text{-}set} \simeq \operatorname{big\text{-}union}(\operatorname{mxn\_}, \mathbb{N}, \operatorname{expr\text{-}stage}(\operatorname{mxn\_})))$
\par\medskip\noindent\textbf{Definition} (\texttt{expr-stage}).\;
\begin{equation*}
(\operatorname{expr\text{-}stage}(0) \simeq ((\{0\} \times \mathbb{N}) \cup ((\{1\} \times \mathbb{N}) \cup (\{2\} \times \mathbb{N}))))
\end{equation*}
\par
Suppose $\operatorname{mxn\_} \in \mathbb{N}$. Then
\begin{equation*}
(\operatorname{expr\text{-}stage}(\operatorname{succ}(\operatorname{mxn\_})) \simeq (\operatorname{expr\text{-}stage}(\operatorname{mxn\_}) \cup ((\{3\} \times \mathbb{N}) \cup (\{4\} \times \mathbb{N}))))
\end{equation*}

\par\medskip\noindent\textbf{Definition} (\texttt{expr-rank}).\;$\operatorname{expr\text{-}rank}(\operatorname{mxe\_}) :=$
\begin{equation*}
(\iota\, \operatorname{mxr\_}.\;((\operatorname{mxr\_} \in \mathbb{N}) \wedge ((\operatorname{mxe\_} \in \operatorname{expr\text{-}stage}(\operatorname{mxr\_})) \wedge \forall \operatorname{mxj\_} \in \mathbb{N}.\; ((\operatorname{mxe\_} \in \operatorname{expr\text{-}stage}(\operatorname{mxj\_})) \Rightarrow (\operatorname{mxr\_} \leq \operatorname{mxj\_})))))
\end{equation*}

\par\medskip\noindent\textbf{Definition} (\texttt{mapl}).\;$\operatorname{mapl}(\operatorname{mxf\_}, \operatorname{mxl\_}) :=$
\begin{equation*}
(\iota\, \operatorname{mxr\_}.\;((\operatorname{mxr\_} \in \operatorname{tuples}(\mathbf{Set})) \wedge ((\operatorname{length}(\operatorname{mxr\_}) = \operatorname{length}(\operatorname{mxl\_})) \wedge \forall \operatorname{mxi\_} \in \mathbb{N}.\; ((1 \leq \operatorname{mxi\_}) \Rightarrow ((\operatorname{mxi\_} \leq \operatorname{length}(\operatorname{mxl\_})) \Rightarrow (\operatorname{nth}(\operatorname{mxi\_}, \operatorname{mxr\_}) = \operatorname{mxf\_}(\operatorname{nth}(\operatorname{mxi\_}, \operatorname{mxl\_}))))))))
\end{equation*}

\par\medskip\noindent\textbf{Definition} (\texttt{maplk}).\;$\operatorname{maplk}(\operatorname{mxf\_}, \operatorname{mxl\_}, \operatorname{mxk\_}) :=$
\begin{equation*}
(\iota\, \operatorname{mxr\_}.\;((\operatorname{mxr\_} \in \operatorname{tuples}(\mathbf{Set})) \wedge ((\operatorname{length}(\operatorname{mxr\_}) = \operatorname{length}(\operatorname{mxl\_})) \wedge \forall \operatorname{mxi\_} \in \mathbb{N}.\; ((1 \leq \operatorname{mxi\_}) \Rightarrow ((\operatorname{mxi\_} \leq \operatorname{length}(\operatorname{mxl\_})) \Rightarrow (\operatorname{nth}(\operatorname{mxi\_}, \operatorname{mxr\_}) = \operatorname{mxf\_}(\operatorname{nth}(\operatorname{mxi\_}, \operatorname{mxl\_}), \operatorname{mxk\_})))))))
\end{equation*}

\par\medskip\noindent\textbf{Definition} (\texttt{fv-stage}).\;
\begin{equation*}
(\operatorname{fv\text{-}stage}(0) \simeq (\lambda \operatorname{mxe\_} \in \operatorname{expr\text{-}stage}(0).\; \left(\text{if }(\operatorname{nth}(1, \operatorname{mxe\_}) = 0)\text{ then }\{\operatorname{nth}(2, \operatorname{mxe\_})\}\text{ else }\emptyset\right)))
\end{equation*}
\par
Suppose $\operatorname{mxn\_} \in \mathbb{N}$. Then
\begin{equation*}
(\operatorname{fv\text{-}stage}(\operatorname{succ}(\operatorname{mxn\_})) \simeq (\lambda \operatorname{mxe\_} \in \operatorname{expr\text{-}stage}(\operatorname{succ}(\operatorname{mxn\_})).\; \left(\text{if }(\operatorname{mxe\_} \in \operatorname{expr\text{-}stage}(\operatorname{mxn\_}))\text{ then }\operatorname{fv\text{-}stage}(\operatorname{mxn\_})(\operatorname{mxe\_})\text{ else }\left(\text{if }(\operatorname{nth}(1, \operatorname{mxe\_}) = 3)\text{ then }\operatorname{big\text{-}union}(\operatorname{mxw\_}, \operatorname{make\text{-}set}(\operatorname{nth}(3, \operatorname{mxe\_})), \operatorname{fv\text{-}stage}(\operatorname{mxn\_})(\operatorname{mxw\_}))\text{ else }(\operatorname{big\text{-}union}(\operatorname{mxw\_}, \operatorname{make\text{-}set}(\operatorname{nth}(4, \operatorname{mxe\_})), \operatorname{fv\text{-}stage}(\operatorname{mxn\_})(\operatorname{mxw\_})) \cup \operatorname{fv\text{-}stage}(\operatorname{mxn\_})(\operatorname{nth}(5, \operatorname{mxe\_})))\right)\right)))
\end{equation*}

\par\medskip\noindent\textbf{Definition} (\texttt{subst-stage}).\;For all $\operatorname{mxm\_}$ and $\operatorname{mxu\_}$, 
\begin{equation*}
(\operatorname{subst\text{-}stage}(\operatorname{mxm\_}, \operatorname{mxu\_}, 0) \simeq (\lambda \operatorname{mxe\_} \in \operatorname{expr\text{-}stage}(0).\; \left(\text{if }(\operatorname{mxe\_} = \langle 0, \operatorname{mxm\_} \rangle)\text{ then }\operatorname{mxu\_}\text{ else }\operatorname{mxe\_}\right)))
\end{equation*}
\par
Let $\operatorname{mxm\_}$ and $\operatorname{mxu\_}$ be arbitrary. Let $\operatorname{mxn\_} \in \mathbb{N}$. Then
\begin{equation*}
(\operatorname{subst\text{-}stage}(\operatorname{mxm\_}, \operatorname{mxu\_}, \operatorname{succ}(\operatorname{mxn\_})) \simeq (\lambda \operatorname{mxe\_} \in \operatorname{expr\text{-}stage}(\operatorname{succ}(\operatorname{mxn\_})).\; \left(\text{if }(\operatorname{mxe\_} \in \operatorname{expr\text{-}stage}(\operatorname{mxn\_}))\text{ then }\operatorname{subst\text{-}stage}(\operatorname{mxm\_}, \operatorname{mxu\_}, \operatorname{mxn\_})(\operatorname{mxe\_})\text{ else }\left(\text{if }(\operatorname{nth}(1, \operatorname{mxe\_}) = 3)\text{ then }\langle 3, \operatorname{nth}(2, \operatorname{mxe\_}), \operatorname{mapl}(\operatorname{subst\text{-}stage}(\operatorname{mxm\_}, \operatorname{mxu\_}, \operatorname{mxn\_}), \operatorname{nth}(3, \operatorname{mxe\_})) \rangle\text{ else }\langle 4, \operatorname{nth}(2, \operatorname{mxe\_}), \operatorname{nth}(3, \operatorname{mxe\_}), \operatorname{mapl}(\operatorname{subst\text{-}stage}(\operatorname{mxm\_}, \operatorname{mxu\_}, \operatorname{mxn\_}), \operatorname{nth}(4, \operatorname{mxe\_})), \operatorname{subst\text{-}stage}(\operatorname{mxm\_}, \operatorname{mxu\_}, \operatorname{mxn\_})(\operatorname{nth}(5, \operatorname{mxe\_})) \rangle\right)\right)))
\end{equation*}

\par\medskip\noindent\textbf{Definition} (\texttt{open-stage}).\;For every $\operatorname{mxu\_}$, 
\begin{equation*}
(\operatorname{open\text{-}stage}(\operatorname{mxu\_}, 0) \simeq (\lambda \langle \operatorname{mxe\_}, \operatorname{mxk\_} \rangle \in (\operatorname{expr\text{-}stage}(0) \times \mathbb{N}).\; \left(\text{if }(\operatorname{mxe\_} = \langle 1, \operatorname{mxk\_} \rangle)\text{ then }\operatorname{mxu\_}\text{ else }\operatorname{mxe\_}\right)))
\end{equation*}
\par
Let $\operatorname{mxu\_}$ be arbitrary. Let $\operatorname{mxn\_} \in \mathbb{N}$. Then
\begin{equation*}
(\operatorname{open\text{-}stage}(\operatorname{mxu\_}, \operatorname{succ}(\operatorname{mxn\_})) \simeq (\lambda \langle \operatorname{mxe\_}, \operatorname{mxk\_} \rangle \in (\operatorname{expr\text{-}stage}(\operatorname{succ}(\operatorname{mxn\_})) \times \mathbb{N}).\; \left(\text{if }(\operatorname{mxe\_} \in \operatorname{expr\text{-}stage}(\operatorname{mxn\_}))\text{ then }\operatorname{open\text{-}stage}(\operatorname{mxu\_}, \operatorname{mxn\_})(\operatorname{mxe\_}, \operatorname{mxk\_})\text{ else }\left(\text{if }(\operatorname{nth}(1, \operatorname{mxe\_}) = 3)\text{ then }\langle 3, \operatorname{nth}(2, \operatorname{mxe\_}), \operatorname{maplk}(\operatorname{open\text{-}stage}(\operatorname{mxu\_}, \operatorname{mxn\_}), \operatorname{nth}(3, \operatorname{mxe\_}), \operatorname{mxk\_}) \rangle\text{ else }\langle 4, \operatorname{nth}(2, \operatorname{mxe\_}), \operatorname{nth}(3, \operatorname{mxe\_}), \operatorname{maplk}(\operatorname{open\text{-}stage}(\operatorname{mxu\_}, \operatorname{mxn\_}), \operatorname{nth}(4, \operatorname{mxe\_}), \operatorname{mxk\_}), \operatorname{open\text{-}stage}(\operatorname{mxu\_}, \operatorname{mxn\_})(\operatorname{nth}(5, \operatorname{mxe\_}), +(\operatorname{mxk\_}, \operatorname{nth}(3, \operatorname{mxe\_}))) \rangle\right)\right)))
\end{equation*}

\par\medskip\noindent\textbf{Definition} (\texttt{expr-lcd}).\;$\operatorname{expr\text{-}lcd}(\operatorname{mxe\_}) :=$ $\operatorname{lcd\text{-}stage}(\operatorname{expr\text{-}rank}(\operatorname{mxe\_}))(\operatorname{mxe\_})$
\par\medskip\noindent\textbf{Definition} (\texttt{lcd-stage}).\;
\begin{equation*}
(\operatorname{lcd\text{-}stage}(0) \simeq (\lambda \operatorname{mxe\_} \in \operatorname{expr\text{-}stage}(0).\; \left(\text{if }(\operatorname{nth}(1, \operatorname{mxe\_}) = 1)\text{ then }\{\operatorname{mxd\_} \in \mathbb{N} : (\operatorname{nth}(2, \operatorname{mxe\_}) < \operatorname{mxd\_})\}\text{ else }\mathbb{N}\right)))
\end{equation*}
\par
Suppose $\operatorname{mxn\_} \in \mathbb{N}$. Then
\begin{equation*}
(\operatorname{lcd\text{-}stage}(\operatorname{succ}(\operatorname{mxn\_})) \simeq (\lambda \operatorname{mxe\_} \in \operatorname{expr\text{-}stage}(\operatorname{succ}(\operatorname{mxn\_})).\; \left(\text{if }(\operatorname{mxe\_} \in \operatorname{expr\text{-}stage}(\operatorname{mxn\_}))\text{ then }\operatorname{lcd\text{-}stage}(\operatorname{mxn\_})(\operatorname{mxe\_})\text{ else }\left(\text{if }(\operatorname{nth}(1, \operatorname{mxe\_}) = 3)\text{ then }\{\operatorname{mxd\_} \in \mathbb{N} : \forall \operatorname{mxw\_} \in \operatorname{make\text{-}set}(\operatorname{nth}(3, \operatorname{mxe\_})).\; (\operatorname{mxd\_} \in \operatorname{lcd\text{-}stage}(\operatorname{mxn\_})(\operatorname{mxw\_}))\}\text{ else }\{\operatorname{mxd\_} \in \mathbb{N} : (\forall \operatorname{mxw\_} \in \operatorname{make\text{-}set}(\operatorname{nth}(4, \operatorname{mxe\_})).\; (\operatorname{mxd\_} \in \operatorname{lcd\text{-}stage}(\operatorname{mxn\_})(\operatorname{mxw\_})) \wedge (+(\operatorname{mxd\_}, \operatorname{nth}(3, \operatorname{mxe\_})) \in \operatorname{lcd\text{-}stage}(\operatorname{mxn\_})(\operatorname{nth}(5, \operatorname{mxe\_}))))\}\right)\right)))
\end{equation*}

\par\medskip\noindent\textbf{Definition} (\texttt{subst-fun}).\;$\operatorname{subst\text{-}fun}(\operatorname{mxm\_}, \operatorname{mxu\_}) :=$ $(\lambda \operatorname{mxw\_} \in \operatorname{expr\text{-}set}.\; \operatorname{expr\text{-}subst}(\operatorname{mxw\_}, \operatorname{mxm\_}, \operatorname{mxu\_}))$
\par\medskip\noindent\textbf{Definition} (\texttt{open-fun}).\;$\operatorname{open\text{-}fun}(\operatorname{mxu\_}) :=$
\begin{equation*}
(\lambda \langle \operatorname{mxw\_}, \operatorname{mxj\_} \rangle \in (\operatorname{expr\text{-}set} \times \mathbb{N}).\; \operatorname{expr\text{-}open}(\operatorname{mxw\_}, \operatorname{mxj\_}, \operatorname{mxu\_}))
\end{equation*}


\vnbbookfile{theorem-library/meta-syntax-laws}{theorem-library/meta-syntax-laws}
\noindent{\small meta-syntax-laws.scm -- THE SYNTAX OF THE KERNEL AS A VNB OBJECT: the laws of the locally nameless expressions (structure-library/meta-syntax.scm). (docs/meta-k2-admissibility-2026-10-08.md, section 6, phase 1: the cost probe.) Agent META-1, 2026-10-08.  Template: theorem-library/strand-messages-laws.scm, whose local}\par
\vnbprop{efvar-unfold}{\vnbstmt{\forall n.\; (\operatorname{efvar}(n) \simeq \langle 0, n \rangle)}}
\vnbprop{ebvar-unfold}{\vnbstmt{\forall k.\; (\operatorname{ebvar}(k) \simeq \langle 1, k \rangle)}}
\vnbprop{econst-unfold}{\vnbstmt{\forall c.\; (\operatorname{econst}(c) \simeq \langle 2, c \rangle)}}
\vnbprop{eapp-unfold}{\vnbstmt{\forall h, l.\; (\operatorname{eapp}(h, l) \simeq \langle 3, h, l \rangle)}}
\vnbprop{ebind-unfold}{\vnbstmt{\forall h, c, a, b.\; (\operatorname{ebind}(h, c, a, b) \simeq \langle 4, h, c, a, b \rangle)}}
\vnbprop{eapp-injective}{\vnbstmt{\forall h, l, a, b.\; ((\operatorname{eapp}(h, l) = \operatorname{eapp}(a, b)) \Rightarrow ((h = a) \wedge (l = b)))}}
\vnbprop{ebind-injective}{\vnbstmt{\begin{array}{@{}l@{}} \forall h, c, a, b, d, f, g, k. \\ \quad (\operatorname{ebind}(h, c, a, b) = \operatorname{ebind}(d, f, g, k)) \Rightarrow  \\ \quad \quad (h = d) \wedge ((c = f) \wedge ((a = g) \wedge (b = k))) \end{array}}}
\vnbprop{expr-stage-fvar}{\vnbstmt{\forall n \in \mathbb{N}.\; (\operatorname{efvar}(n) \in \operatorname{expr\text{-}stage}(0))}}
\vnbprop{expr-stage-bvar}{\vnbstmt{\forall k \in \mathbb{N}.\; (\operatorname{ebvar}(k) \in \operatorname{expr\text{-}stage}(0))}}
\vnbprop{expr-stage-const}{\vnbstmt{\forall c \in \mathbb{N}.\; (\operatorname{econst}(c) \in \operatorname{expr\text{-}stage}(0))}}
\vnbprop{expr-stage-keep}{\vnbstmt{\forall n, m.\; (((n \in \mathbb{N}) \wedge (m \in \operatorname{expr\text{-}stage}(n))) \Rightarrow (m \in \operatorname{expr\text{-}stage}((n + 1))))}}
\vnbprop{expr-stage-app}{\vnbstmt{\begin{array}{@{}l@{}} \forall n, h, l. \\ \quad ((n \in \mathbb{N}) \wedge ((h \in \mathbb{N}) \wedge (l \in \operatorname{tuples}(\operatorname{expr\text{-}stage}(n))))) \Rightarrow  \\ \quad \quad \operatorname{eapp}(h, l) \in \operatorname{expr\text{-}stage}((n + 1)) \end{array}}}
\vnbprop{expr-stage-bind}{\vnbstmt{\begin{array}{@{}l@{}} \forall n, h, c, a, b. \\ \quad ((n \in \mathbb{N}) \wedge ((h \in \mathbb{N}) \wedge ((c \in \mathbb{N}) \wedge ((a \in \operatorname{tuples}(\operatorname{expr\text{-}stage}(n))) \wedge (b \in \operatorname{expr\text{-}stage}(n)))))) \Rightarrow  \\ \quad \quad \operatorname{ebind}(h, c, a, b) \in \operatorname{expr\text{-}stage}((n + 1)) \end{array}}}
\vnbprop{expr-stage-zero-cases}{\vnbstmt{\begin{array}{@{}l@{}} \forall m \in \operatorname{expr\text{-}stage}(0). \\ \quad \exists c \in \mathbb{N}.\; (m = \operatorname{efvar}(c)) \vee \\ \quad \exists c \in \mathbb{N}.\; (m = \operatorname{ebvar}(c)) \vee \\ \quad \exists c \in \mathbb{N}.\; (m = \operatorname{econst}(c)) \end{array}}}
\vnbprop{expr-stage-succ-cases}{\vnbstmt{\begin{array}{@{}l@{}} \forall n, m. \\ \quad ((n \in \mathbb{N}) \wedge (m \in \operatorname{expr\text{-}stage}((n + 1)))) \Rightarrow  \\ \quad \quad m \in \operatorname{expr\text{-}stage}(n) \vee \\ \quad \quad \exists h \in \mathbb{N}, l \in \operatorname{tuples}(\operatorname{expr\text{-}stage}(n)).\; (m = \operatorname{eapp}(h, l)) \vee \\ \quad \quad \exists h, c \in \mathbb{N}, a \in \operatorname{tuples}(\operatorname{expr\text{-}stage}(n)), b \in \operatorname{expr\text{-}stage}(n). \\ \quad \quad \quad m = \operatorname{ebind}(h, c, a, b) \end{array}}}
\vnbprop{expr-stage-add}{\vnbstmt{\begin{array}{@{}l@{}} \forall d \in \mathbb{N}, j, m. \\ \quad ((j \in \mathbb{N}) \wedge (m \in \operatorname{expr\text{-}stage}(j))) \Rightarrow (m \in \operatorname{expr\text{-}stage}((j + d))) \end{array}}}
\vnbprop{expr-set-intro}{\vnbstmt{\forall n, m.\; (((n \in \mathbb{N}) \wedge (m \in \operatorname{expr\text{-}stage}(n))) \Rightarrow (m \in \operatorname{expr\text{-}set}))}}
\vnbprop{expr-set-elim}{\vnbstmt{\forall m \in \operatorname{expr\text{-}set}.\; \exists n \in \mathbb{N}.\; (m \in \operatorname{expr\text{-}stage}(n))}}
\vnbprop{cartesian5-set}{\vnbstmt{\begin{array}{@{}l@{}} \forall a_{1}, a_{2}, a_{3}, a_{4}, a_{5}. \\ \quad ((a_{1} \in \mathbf{Set}) \wedge ((a_{2} \in \mathbf{Set}) \wedge ((a_{3} \in \mathbf{Set}) \wedge ((a_{4} \in \mathbf{Set}) \wedge (a_{5} \in \mathbf{Set}))))) \Rightarrow  \\ \quad \quad (a_{1} \times a_{2}) \in \mathbf{Set} \end{array}}}
\vnbprop{expr-stage-set}{\vnbstmt{\forall n \in \mathbb{N}.\; (\operatorname{expr\text{-}stage}(n) \in \mathbf{Set})}}
\vnbprop{expr-set-is-set}{\vnbstmt{\operatorname{expr\text{-}set} \in \mathbf{Set}}}
\vnbprop{expr-list-in-some-stage}{\vnbstmt{\forall l \in \operatorname{tuples}(\operatorname{expr\text{-}set}).\; \exists n \in \mathbb{N}.\; (l \in \operatorname{tuples}(\operatorname{expr\text{-}stage}(n)))}}
\vnbprop{efvar-in-expr}{\vnbstmt{\forall n \in \mathbb{N}.\; (\operatorname{efvar}(n) \in \operatorname{expr\text{-}set})}}
\vnbprop{ebvar-in-expr}{\vnbstmt{\forall k \in \mathbb{N}.\; (\operatorname{ebvar}(k) \in \operatorname{expr\text{-}set})}}
\vnbprop{econst-in-expr}{\vnbstmt{\forall c \in \mathbb{N}.\; (\operatorname{econst}(c) \in \operatorname{expr\text{-}set})}}
\vnbprop{eapp-in-expr}{\vnbstmt{\forall h, l.\; (((h \in \mathbb{N}) \wedge (l \in \operatorname{tuples}(\operatorname{expr\text{-}set}))) \Rightarrow (\operatorname{eapp}(h, l) \in \operatorname{expr\text{-}set}))}}
\vnbprop{ebind-in-expr}{\vnbstmt{\begin{array}{@{}l@{}} \forall h, c, a, b. \\ \quad ((h \in \mathbb{N}) \wedge ((c \in \mathbb{N}) \wedge ((a \in \operatorname{tuples}(\operatorname{expr\text{-}set})) \wedge (b \in \operatorname{expr\text{-}set})))) \Rightarrow  \\ \quad \quad \operatorname{ebind}(h, c, a, b) \in \operatorname{expr\text{-}set} \end{array}}}
\vnbprop{expr-induction}{\vnbstmt{\begin{array}{@{}l@{}} \forall s.; \forall e \in \operatorname{expr\text{-}set}. \\ \quad \forall n \in \mathbb{N}.\; (\operatorname{efvar}(n) \in s) \wedge \\ \quad \forall k \in \mathbb{N}.\; (\operatorname{ebvar}(k) \in s) \wedge \\ \quad \forall c \in \mathbb{N}.\; (\operatorname{econst}(c) \in s) \wedge \\ \quad \forall h, l. \\ \quad \quad h \in \mathbb{N} \wedge \\ \quad \quad l \in \operatorname{tuples}(\operatorname{expr\text{-}set}) \wedge \\ \quad \quad \forall i \in \mathbb{N}.\; (((1 \leq i) \wedge (i \leq \operatorname{length}(l))) \Rightarrow (\operatorname{nth}(i, l) \in s)) \Rightarrow \\ \quad \quad \quad \operatorname{eapp}(h, l) \in s \wedge \\ \quad \forall h, b, a, \mathit{mxjb}. \\ \quad \quad h \in \mathbb{N} \wedge \\ \quad \quad b \in \mathbb{N} \wedge \\ \quad \quad a \in \operatorname{tuples}(\operatorname{expr\text{-}set}) \wedge \\ \quad \quad \mathit{mxjb} \in \operatorname{expr\text{-}set} \wedge \\ \quad \quad \forall i \in \mathbb{N}.\; (((1 \leq i) \wedge (i \leq \operatorname{length}(a))) \Rightarrow (\operatorname{nth}(i, a) \in s)) \wedge \\ \quad \quad \mathit{mxjb} \in s \Rightarrow \\ \quad \quad \quad \operatorname{ebind}(h, b, a, \mathit{mxjb}) \in s \Rightarrow \\ \quad \quad e \in s \end{array}}}
\vnbprop{expr-rank-prop}{\vnbstmt{\begin{array}{@{}l@{}} \forall e \in \operatorname{expr\text{-}set}. \\ \quad \operatorname{expr\text{-}rank}(e) \in \mathbb{N} \wedge \\ \quad e \in \operatorname{expr\text{-}stage}(\operatorname{expr\text{-}rank}(e)) \wedge \\ \quad \forall j \in \mathbb{N}.\; ((e \in \operatorname{expr\text{-}stage}(j)) \Rightarrow (\operatorname{expr\text{-}rank}(e) \leq j)) \end{array}}}
\vnbprop{expr-rank-fvar}{\vnbstmt{\forall n \in \mathbb{N}.\; (\operatorname{expr\text{-}rank}(\operatorname{efvar}(n)) = 0)}}
\vnbprop{expr-rank-bvar}{\vnbstmt{\forall n \in \mathbb{N}.\; (\operatorname{expr\text{-}rank}(\operatorname{ebvar}(n)) = 0)}}
\vnbprop{expr-rank-const}{\vnbstmt{\forall n \in \mathbb{N}.\; (\operatorname{expr\text{-}rank}(\operatorname{econst}(n)) = 0)}}
\vnbprop{expr-rank-app}{\vnbstmt{\begin{array}{@{}l@{}} \forall h, l. \\ \quad ((h \in \mathbb{N}) \wedge (l \in \operatorname{tuples}(\operatorname{expr\text{-}set}))) \Rightarrow  \\ \quad \quad \exists k \in \mathbb{N}. \\ \quad \quad \quad ((\operatorname{expr\text{-}rank}(\operatorname{eapp}(h, l)) = (k + 1)) \wedge ((l \in \operatorname{tuples}(\operatorname{expr\text{-}stage}(k))) \wedge \neg (\operatorname{eapp}(h, l) \in \operatorname{expr\text{-}stage}(k)))) \end{array}}}
\vnbprop{expr-rank-bind}{\vnbstmt{\begin{array}{@{}l@{}} \forall h, c, a, b. \\ \quad ((h \in \mathbb{N}) \wedge ((c \in \mathbb{N}) \wedge ((a \in \operatorname{tuples}(\operatorname{expr\text{-}set})) \wedge (b \in \operatorname{expr\text{-}set})))) \Rightarrow  \\ \quad \quad \exists k \in \mathbb{N}. \\ \quad \quad \quad \operatorname{expr\text{-}rank}(\operatorname{ebind}(h, c, a, b)) = (k + 1) \wedge \\ \quad \quad \quad a \in \operatorname{tuples}(\operatorname{expr\text{-}stage}(k)) \wedge \\ \quad \quad \quad b \in \operatorname{expr\text{-}stage}(k) \wedge \\ \quad \quad \quad \neg (\operatorname{ebind}(h, c, a, b) \in \operatorname{expr\text{-}stage}(k)) \end{array}}}
\vnbprop{makeset-member-in}{\vnbstmt{\forall s, l, w.\; (((l \in \operatorname{tuples}(s)) \wedge (w \in \operatorname{make\text{-}set}(l))) \Rightarrow (w \in s))}}
\vnbprop{fv-stage-step}{\vnbstmt{\begin{array}{@{}l@{}} \forall k, e. \\ \quad ((k \in \mathbb{N}) \wedge (e \in \operatorname{expr\text{-}stage}(k))) \Rightarrow  \\ \quad \quad (\operatorname{fv\text{-}stage}((k + 1))(e) \simeq \operatorname{fv\text{-}stage}(k)(e)) \end{array}}}
\vnbprop{fv-stage-coh}{\vnbstmt{\begin{array}{@{}l@{}} \forall d \in \mathbb{N}, j, e. \\ \quad ((j \in \mathbb{N}) \wedge (e \in \operatorname{expr\text{-}stage}(j))) \Rightarrow (\operatorname{fv\text{-}stage}((j + d))(e) \simeq \operatorname{fv\text{-}stage}(j)(e)) \end{array}}}
\vnbprop{expr-fv-at}{\vnbstmt{\begin{array}{@{}l@{}} \forall j, e. \\ \quad ((j \in \mathbb{N}) \wedge (e \in \operatorname{expr\text{-}stage}(j))) \Rightarrow (\operatorname{expr\text{-}fv}(e) \simeq \operatorname{fv\text{-}stage}(j)(e)) \end{array}}}
\vnbprop{fv-efvar}{\vnbstmt{\forall n \in \mathbb{N}.\; (\operatorname{expr\text{-}fv}(\operatorname{efvar}(n)) \simeq \{n\})}}
\vnbprop{fv-ebvar}{\vnbstmt{\forall k \in \mathbb{N}.\; (\operatorname{expr\text{-}fv}(\operatorname{ebvar}(k)) \simeq \emptyset)}}
\vnbprop{fv-econst}{\vnbstmt{\forall c \in \mathbb{N}.\; (\operatorname{expr\text{-}fv}(\operatorname{econst}(c)) \simeq \emptyset)}}
\vnbprop{fv-eapp}{\vnbstmt{\begin{array}{@{}l@{}} \forall h, l. \\ \quad ((h \in \mathbb{N}) \wedge (l \in \operatorname{tuples}(\operatorname{expr\text{-}set}))) \Rightarrow  \\ \quad \quad (\operatorname{expr\text{-}fv}(\operatorname{eapp}(h, l)) \simeq \operatorname{big\text{-}union}(w, \operatorname{make\text{-}set}(l), \operatorname{expr\text{-}fv}(w))) \end{array}}}
\vnbprop{fv-ebind}{\vnbstmt{\begin{array}{@{}l@{}} \forall h, c, a, b. \\ \quad ((h \in \mathbb{N}) \wedge ((c \in \mathbb{N}) \wedge ((a \in \operatorname{tuples}(\operatorname{expr\text{-}set})) \wedge (b \in \operatorname{expr\text{-}set})))) \Rightarrow  \\ \quad \quad (\operatorname{expr\text{-}fv}(\operatorname{ebind}(h, c, a, b)) \simeq (\operatorname{big\text{-}union}(w, \operatorname{make\text{-}set}(a), \operatorname{expr\text{-}fv}(w)) \cup \operatorname{expr\text{-}fv}(b))) \end{array}}}
\vnbprop{bu-finite-nn}{\vnbstmt{\begin{array}{@{}l@{}} \forall s \in \mathbf{Set}, m. \\ \quad ((|s| \in \mathbb{N}) \wedge \forall v_{5} \in s.\; ((m(v_{5}) \in \mathbf{Set}) \wedge ((m(v_{5}) \subseteq \mathbb{N}) \wedge (|m(v_{5})| \in \mathbb{N})))) \Rightarrow  \\ \quad \quad \operatorname{big\text{-}union}(w, s, m(w)) \in \mathbf{Set} \wedge \\ \quad \quad \operatorname{big\text{-}union}(w, s, m(w)) \subseteq \mathbb{N} \wedge \\ \quad \quad |\operatorname{big\text{-}union}(w, s, m(w))| \in \mathbb{N} \end{array}}}
\vnbprop{fv-finite}{\vnbstmt{\begin{array}{@{}l@{}} \forall e \in \operatorname{expr\text{-}set}. \\ \quad (\operatorname{expr\text{-}fv}(e) \in \mathbf{Set}) \wedge ((\operatorname{expr\text{-}fv}(e) \subseteq \mathbb{N}) \wedge (|\operatorname{expr\text{-}fv}(e)| \in \mathbb{N})) \end{array}}}
\vnbprop{mapl-prop}{\vnbstmt{\begin{array}{@{}l@{}} \forall s, g, l. \\ \quad ((\operatorname{length}(l) \in \mathbb{N}) \wedge \forall i \in \mathbb{N}.\; (((1 \leq i) \wedge (i \leq \operatorname{length}(l))) \Rightarrow (g(\operatorname{nth}(i, l)) \in s))) \Rightarrow  \\ \quad \quad \operatorname{mapl}(g, l) \in \operatorname{tuples}(s) \wedge \\ \quad \quad \operatorname{length}(\operatorname{mapl}(g, l)) = \operatorname{length}(l) \wedge \\ \quad \quad \forall i \in \mathbb{N}. \\ \quad \quad \quad ((1 \leq i) \wedge (i \leq \operatorname{length}(l))) \Rightarrow (\operatorname{nth}(i, \operatorname{mapl}(g, l)) = g(\operatorname{nth}(i, l))) \end{array}}}
\vnbprop{maplk-prop}{\vnbstmt{\begin{array}{@{}l@{}} \forall s, g, l, x. \\ \quad ((\operatorname{length}(l) \in \mathbb{N}) \wedge \forall i \in \mathbb{N}.\; (((1 \leq i) \wedge (i \leq \operatorname{length}(l))) \Rightarrow (g(\operatorname{nth}(i, l), x) \in s))) \Rightarrow  \\ \quad \quad \operatorname{maplk}(g, l, x) \in \operatorname{tuples}(s) \wedge \\ \quad \quad \operatorname{length}(\operatorname{maplk}(g, l, x)) = \operatorname{length}(l) \wedge \\ \quad \quad \forall i \in \mathbb{N}. \\ \quad \quad \quad ((1 \leq i) \wedge (i \leq \operatorname{length}(l))) \Rightarrow (\operatorname{nth}(i, \operatorname{maplk}(g, l, x)) = g(\operatorname{nth}(i, l), x)) \end{array}}}
\vnbprop{subst-stage-fun}{\vnbstmt{\begin{array}{@{}l@{}} \forall n \in \mathbb{N}, m.; \forall u \in \operatorname{expr\text{-}set}. \\ \quad \operatorname{subst\text{-}stage}(m, u, n) \in (\operatorname{expr\text{-}stage}(n) \to \operatorname{expr\text{-}set}) \end{array}}}
\vnbprop{subst-stage-step}{\vnbstmt{\begin{array}{@{}l@{}} \forall m, u, k, e. \\ \quad ((k \in \mathbb{N}) \wedge (e \in \operatorname{expr\text{-}stage}(k))) \Rightarrow  \\ \quad \quad (\operatorname{subst\text{-}stage}(m, u, (k + 1))(e) \simeq \operatorname{subst\text{-}stage}(m, u, k)(e)) \end{array}}}
\vnbprop{subst-stage-coh}{\vnbstmt{\begin{array}{@{}l@{}} \forall d \in \mathbb{N}, m, u, j, e. \\ \quad ((j \in \mathbb{N}) \wedge (e \in \operatorname{expr\text{-}stage}(j))) \Rightarrow  \\ \quad \quad (\operatorname{subst\text{-}stage}(m, u, (j + d))(e) \simeq \operatorname{subst\text{-}stage}(m, u, j)(e)) \end{array}}}
\vnbprop{expr-subst-at}{\vnbstmt{\begin{array}{@{}l@{}} \forall m, u, j, e. \\ \quad ((j \in \mathbb{N}) \wedge (e \in \operatorname{expr\text{-}stage}(j))) \Rightarrow  \\ \quad \quad (\operatorname{expr\text{-}subst}(e, m, u) \simeq \operatorname{subst\text{-}stage}(m, u, j)(e)) \end{array}}}
\vnbprop{subst-in-expr}{\vnbstmt{\forall m.; \forall u, e \in \operatorname{expr\text{-}set}.\; (\operatorname{expr\text{-}subst}(e, m, u) \in \operatorname{expr\text{-}set})}}
\vnbprop{subst-efvar-hit}{\vnbstmt{\forall m \in \mathbb{N}, u \in \operatorname{expr\text{-}set}.\; (\operatorname{expr\text{-}subst}(\operatorname{efvar}(m), m, u) \simeq u)}}
\vnbprop{subst-efvar-miss}{\vnbstmt{\begin{array}{@{}l@{}} \forall m, u, n. \\ \quad ((u \in \operatorname{expr\text{-}set}) \wedge ((n \in \mathbb{N}) \wedge \neg (n = m))) \Rightarrow  \\ \quad \quad (\operatorname{expr\text{-}subst}(\operatorname{efvar}(n), m, u) \simeq \operatorname{efvar}(n)) \end{array}}}
\vnbprop{subst-ebvar}{\vnbstmt{\begin{array}{@{}l@{}} \forall m, u, k. \\ \quad ((u \in \operatorname{expr\text{-}set}) \wedge (k \in \mathbb{N})) \Rightarrow (\operatorname{expr\text{-}subst}(\operatorname{ebvar}(k), m, u) \simeq \operatorname{ebvar}(k)) \end{array}}}
\vnbprop{subst-econst}{\vnbstmt{\begin{array}{@{}l@{}} \forall m, u, c. \\ \quad ((u \in \operatorname{expr\text{-}set}) \wedge (c \in \mathbb{N})) \Rightarrow (\operatorname{expr\text{-}subst}(\operatorname{econst}(c), m, u) \simeq \operatorname{econst}(c)) \end{array}}}
\vnbprop{subst-eapp}{\vnbstmt{\begin{array}{@{}l@{}} \forall m, u, h, l. \\ \quad ((u \in \operatorname{expr\text{-}set}) \wedge ((h \in \mathbb{N}) \wedge (l \in \operatorname{tuples}(\operatorname{expr\text{-}set})))) \Rightarrow  \\ \quad \quad (\operatorname{expr\text{-}subst}(\operatorname{eapp}(h, l), m, u) \simeq \operatorname{eapp}(h, \operatorname{mapl}(\operatorname{subst\text{-}fun}(m, u), l))) \end{array}}}
\vnbprop{subst-ebind}{\vnbstmt{\begin{array}{@{}l@{}} \forall m, u, h, c, a, b. \\ \quad ((u \in \operatorname{expr\text{-}set}) \wedge ((h \in \mathbb{N}) \wedge ((c \in \mathbb{N}) \wedge ((a \in \operatorname{tuples}(\operatorname{expr\text{-}set})) \wedge (b \in \operatorname{expr\text{-}set}))))) \Rightarrow  \\ \quad \quad (\operatorname{expr\text{-}subst}(\operatorname{ebind}(h, c, a, b), m, u) \simeq \operatorname{ebind}(h, c, \operatorname{mapl}(\operatorname{subst\text{-}fun}(m, u), a), \operatorname{expr\text{-}subst}(b, m, u))) \end{array}}}
\vnbprop{subst-fresh}{\vnbstmt{\begin{array}{@{}l@{}} \forall m.; \forall u, e \in \operatorname{expr\text{-}set}. \\ \quad \neg (m \in \operatorname{expr\text{-}fv}(e)) \Rightarrow (\operatorname{expr\text{-}subst}(e, m, u) = e) \end{array}}}
\vnbprop{subst-fun-unfold}{\vnbstmt{\forall m, u.\; (\operatorname{subst\text{-}fun}(m, u) \simeq (\lambda w \in \operatorname{expr\text{-}set}.\; \operatorname{expr\text{-}subst}(w, m, u)))}}
\vnbprop{fv-subst}{\vnbstmt{\begin{array}{@{}l@{}} \forall m.; \forall u, e \in \operatorname{expr\text{-}set}, x \in \operatorname{expr\text{-}fv}(\operatorname{expr\text{-}subst}(e, m, u)). \\ \quad ((x \in \operatorname{expr\text{-}fv}(e)) \wedge \neg (x = m)) \vee (x \in \operatorname{expr\text{-}fv}(u)) \end{array}}}
\vnbprop{open-stage-fun}{\vnbstmt{\begin{array}{@{}l@{}} \forall n \in \mathbb{N}, u \in \operatorname{expr\text{-}set}. \\ \quad \operatorname{open\text{-}stage}(u, n) \in ((\operatorname{expr\text{-}stage}(n) \times \mathbb{N}) \to \operatorname{expr\text{-}set}) \end{array}}}
\vnbprop{open-stage-step}{\vnbstmt{\begin{array}{@{}l@{}} \forall u, k, e, x. \\ \quad ((k \in \mathbb{N}) \wedge ((e \in \operatorname{expr\text{-}stage}(k)) \wedge (x \in \mathbb{N}))) \Rightarrow  \\ \quad \quad (\operatorname{open\text{-}stage}(u, (k + 1))(e, x) \simeq \operatorname{open\text{-}stage}(u, k)(e, x)) \end{array}}}
\vnbprop{open-stage-coh}{\vnbstmt{\begin{array}{@{}l@{}} \forall d \in \mathbb{N}, u, j, e, x. \\ \quad ((j \in \mathbb{N}) \wedge ((e \in \operatorname{expr\text{-}stage}(j)) \wedge (x \in \mathbb{N}))) \Rightarrow  \\ \quad \quad (\operatorname{open\text{-}stage}(u, (j + d))(e, x) \simeq \operatorname{open\text{-}stage}(u, j)(e, x)) \end{array}}}
\vnbprop{expr-open-at}{\vnbstmt{\begin{array}{@{}l@{}} \forall u, j, e, x. \\ \quad ((j \in \mathbb{N}) \wedge ((e \in \operatorname{expr\text{-}stage}(j)) \wedge (x \in \mathbb{N}))) \Rightarrow  \\ \quad \quad (\operatorname{expr\text{-}open}(e, x, u) \simeq \operatorname{open\text{-}stage}(u, j)(e, x)) \end{array}}}
\vnbprop{open-in-expr}{\vnbstmt{\forall u \in \operatorname{expr\text{-}set}, x \in \mathbb{N}, e \in \operatorname{expr\text{-}set}.\; (\operatorname{expr\text{-}open}(e, x, u) \in \operatorname{expr\text{-}set})}}
\vnbprop{open-ebvar-hit}{\vnbstmt{\forall u, x.\; (((u \in \operatorname{expr\text{-}set}) \wedge (x \in \mathbb{N})) \Rightarrow (\operatorname{expr\text{-}open}(\operatorname{ebvar}(x), x, u) \simeq u))}}
\vnbprop{open-ebvar-miss}{\vnbstmt{\begin{array}{@{}l@{}} \forall u, x, j. \\ \quad ((u \in \operatorname{expr\text{-}set}) \wedge ((x \in \mathbb{N}) \wedge ((j \in \mathbb{N}) \wedge \neg (j = x)))) \Rightarrow  \\ \quad \quad (\operatorname{expr\text{-}open}(\operatorname{ebvar}(j), x, u) \simeq \operatorname{ebvar}(j)) \end{array}}}
\vnbprop{open-efvar}{\vnbstmt{\begin{array}{@{}l@{}} \forall u, x, n. \\ \quad ((u \in \operatorname{expr\text{-}set}) \wedge ((x \in \mathbb{N}) \wedge (n \in \mathbb{N}))) \Rightarrow  \\ \quad \quad (\operatorname{expr\text{-}open}(\operatorname{efvar}(n), x, u) \simeq \operatorname{efvar}(n)) \end{array}}}
\vnbprop{open-econst}{\vnbstmt{\begin{array}{@{}l@{}} \forall u, x, c. \\ \quad ((u \in \operatorname{expr\text{-}set}) \wedge ((x \in \mathbb{N}) \wedge (c \in \mathbb{N}))) \Rightarrow  \\ \quad \quad (\operatorname{expr\text{-}open}(\operatorname{econst}(c), x, u) \simeq \operatorname{econst}(c)) \end{array}}}
\vnbprop{open-eapp}{\vnbstmt{\begin{array}{@{}l@{}} \forall u, x, h, l. \\ \quad ((u \in \operatorname{expr\text{-}set}) \wedge ((x \in \mathbb{N}) \wedge ((h \in \mathbb{N}) \wedge (l \in \operatorname{tuples}(\operatorname{expr\text{-}set}))))) \Rightarrow  \\ \quad \quad (\operatorname{expr\text{-}open}(\operatorname{eapp}(h, l), x, u) \simeq \operatorname{eapp}(h, \operatorname{maplk}(\operatorname{open\text{-}fun}(u), l, x))) \end{array}}}
\vnbprop{open-ebind}{\vnbstmt{\begin{array}{@{}l@{}} \forall u, x, h, c, a, b. \\ \quad ((u \in \operatorname{expr\text{-}set}) \wedge ((x \in \mathbb{N}) \wedge ((h \in \mathbb{N}) \wedge ((c \in \mathbb{N}) \wedge ((a \in \operatorname{tuples}(\operatorname{expr\text{-}set})) \wedge (b \in \operatorname{expr\text{-}set})))))) \Rightarrow  \\ \quad \quad (\operatorname{expr\text{-}open}(\operatorname{ebind}(h, c, a, b), x, u) \simeq \operatorname{ebind}(h, c, \operatorname{maplk}(\operatorname{open\text{-}fun}(u), a, x), \operatorname{expr\text{-}open}(b, (x + c), u))) \end{array}}}
\vnbprop{lcd-stage-step}{\vnbstmt{\begin{array}{@{}l@{}} \forall k, e. \\ \quad ((k \in \mathbb{N}) \wedge (e \in \operatorname{expr\text{-}stage}(k))) \Rightarrow  \\ \quad \quad (\operatorname{lcd\text{-}stage}((k + 1))(e) \simeq \operatorname{lcd\text{-}stage}(k)(e)) \end{array}}}
\vnbprop{lcd-stage-coh}{\vnbstmt{\begin{array}{@{}l@{}} \forall d \in \mathbb{N}, j, e. \\ \quad ((j \in \mathbb{N}) \wedge (e \in \operatorname{expr\text{-}stage}(j))) \Rightarrow  \\ \quad \quad (\operatorname{lcd\text{-}stage}((j + d))(e) \simeq \operatorname{lcd\text{-}stage}(j)(e)) \end{array}}}
\vnbprop{expr-lcd-at}{\vnbstmt{\begin{array}{@{}l@{}} \forall j, e. \\ \quad ((j \in \mathbb{N}) \wedge (e \in \operatorname{expr\text{-}stage}(j))) \Rightarrow (\operatorname{expr\text{-}lcd}(e) \simeq \operatorname{lcd\text{-}stage}(j)(e)) \end{array}}}
\vnbprop{lcd-efvar}{\vnbstmt{\forall n \in \mathbb{N}.\; (\operatorname{expr\text{-}lcd}(\operatorname{efvar}(n)) \simeq \mathbb{N})}}
\vnbprop{lcd-econst}{\vnbstmt{\forall c \in \mathbb{N}.\; (\operatorname{expr\text{-}lcd}(\operatorname{econst}(c)) \simeq \mathbb{N})}}
\vnbprop{lcd-ebvar}{\vnbstmt{\forall k \in \mathbb{N}.\; (\operatorname{expr\text{-}lcd}(\operatorname{ebvar}(k)) \simeq \{d \in \mathbb{N} : (k < d)\})}}
\vnbprop{lcd-eapp}{\vnbstmt{\begin{array}{@{}l@{}} \forall h, l. \\ \quad ((h \in \mathbb{N}) \wedge (l \in \operatorname{tuples}(\operatorname{expr\text{-}set}))) \Rightarrow  \\ \quad \quad (\operatorname{expr\text{-}lcd}(\operatorname{eapp}(h, l)) \simeq \{d \in \mathbb{N} : \forall w \in \operatorname{make\text{-}set}(l).\; (d \in \operatorname{expr\text{-}lcd}(w))\}) \end{array}}}
\vnbprop{lcd-ebind}{\vnbstmt{\begin{array}{@{}l@{}} \forall h, c, a, b. \\ \quad ((h \in \mathbb{N}) \wedge ((c \in \mathbb{N}) \wedge ((a \in \operatorname{tuples}(\operatorname{expr\text{-}set})) \wedge (b \in \operatorname{expr\text{-}set})))) \Rightarrow  \\ \quad \quad (\operatorname{expr\text{-}lcd}(\operatorname{ebind}(h, c, a, b)) \simeq \{d \in \mathbb{N} : (\forall w \in \operatorname{make\text{-}set}(a).\; (d \in \operatorname{expr\text{-}lcd}(w)) \wedge ((d + c) \in \operatorname{expr\text{-}lcd}(b)))\}) \end{array}}}
\vnbprop{lcd-open}{\vnbstmt{\begin{array}{@{}l@{}} \forall u, e \in \operatorname{expr\text{-}set}, x, d. \\ \quad ((x \in \mathbb{N}) \wedge ((d \in \mathbb{N}) \wedge ((d \leq x) \wedge (d \in \operatorname{expr\text{-}lcd}(e))))) \Rightarrow  \\ \quad \quad \operatorname{expr\text{-}open}(e, x, u) = e \end{array}}}
\vnbprop{lc-open}{\vnbstmt{\begin{array}{@{}l@{}} \forall u, e \in \operatorname{expr\text{-}set}, x \in \mathbb{N}. \\ \quad \operatorname{expr\text{-}lc}(e) \Rightarrow (\operatorname{expr\text{-}open}(e, x, u) = e) \end{array}}}
\vnbprop{subst-open-commute}{\vnbstmt{\begin{array}{@{}l@{}} \forall n, u, v.; \forall e \in \operatorname{expr\text{-}set}, x \in \mathbb{N}. \\ \quad ((u \in \operatorname{expr\text{-}set}) \wedge ((v \in \operatorname{expr\text{-}set}) \wedge \operatorname{expr\text{-}lc}(v))) \Rightarrow  \\ \quad \quad \operatorname{expr\text{-}subst}(\operatorname{expr\text{-}open}(e, x, u), n, v) = \operatorname{expr\text{-}open}(\operatorname{expr\text{-}subst}(e, n, v), x, \operatorname{expr\text{-}subst}(u, n, v)) \end{array}}}

\vnbbookfile{structure-library/meta-core}{structure-library/meta-core}
\noindent{\small meta-core.scm -- THE CORE OF K2: well-formedness, sequents, the core rule relations, DERIVABLE, and the eight kernel relations transcribed (three in phase 2, five by META-KIT). (docs/meta-k2-admissibility-2026-10-08.md, sections 3, 4 and 5, phase 2.) DEFINITIONS ONLY; the laws are in theorem-library/meta-core-laws.scm.}\par
\par\medskip\noindent\textbf{Definition} (\texttt{derivable}).\;$\operatorname{derivable}(\operatorname{mct\_}) :=$
\begin{equation*}
\{\operatorname{mcq\_} \in \operatorname{seq\text{-}set} : \forall \operatorname{mcs\_} \in \operatorname{power}(\operatorname{seq\text{-}set}).\; (\operatorname{core\text{-}closed}(\operatorname{mct\_}, \operatorname{mcs\_}) \Rightarrow (\operatorname{mcq\_} \in \operatorname{mcs\_}))\}
\end{equation*}

\par\medskip\noindent\textbf{Definition} (\texttt{enc-num}).\;$\operatorname{enc\text{-}num}(\operatorname{mck\_}) :=$ $\operatorname{econst}((2 \cdot \operatorname{mck\_}))$
\par\medskip\noindent\textbf{Definition} (\texttt{enc-setc}).\; $(\operatorname{enc\text{-}setc} \simeq \operatorname{econst}(1))$
\par\medskip\noindent\textbf{Definition} (\texttt{pred-code}).\;For every $\operatorname{mcpk\_}$, 
\begin{equation*}
\operatorname{pred\text{-}code}(\operatorname{mcpk\_}) \Leftrightarrow ((\operatorname{mcpk\_} = 7) \vee ((\operatorname{mcpk\_} = 9) \vee ((\operatorname{mcpk\_} = 11) \vee (\operatorname{mcpk\_} = 13))))
\end{equation*}

\par\medskip\noindent\textbf{Definition} (\texttt{hd-res}).\;$\operatorname{hd\text{-}res}(\operatorname{mch\_}) :=$ $\left(\text{if }(\operatorname{mch\_} \leq 10)\text{ then }2\text{ else }1\right)$
\par\medskip\noindent\textbf{Definition} (\texttt{app-ar-ok}).\;For all $\operatorname{mch\_}$ and $\operatorname{mcn\_}$, 
\begin{equation*}
\begin{array}{@{}l@{}} \operatorname{app\text{-}ar\text{-}ok}(\operatorname{mch\_}, \operatorname{mcn\_}) \\ \Leftrightarrow\; ((10 \leq \operatorname{mch\_}) \vee (((\operatorname{mch\_} \leq 1) \wedge (\operatorname{mcn\_} = 0)) \vee (((\operatorname{mch\_} = 2) \wedge (\operatorname{mcn\_} = 1)) \vee ((3 \leq \operatorname{mch\_}) \wedge (\operatorname{mcn\_} = 2))))) \end{array}
\end{equation*}

\par\medskip\noindent\textbf{Definition} (\texttt{hd-argk}).\;$\operatorname{hd\text{-}argk}(\operatorname{mch\_}) :=$ $\left(\text{if }(\operatorname{mch\_} \leq 6)\text{ then }2\text{ else }1\right)$
\par\medskip\noindent\textbf{Definition} (\texttt{hd-bres}).\;$\operatorname{hd\text{-}bres}(\operatorname{mch\_}) :=$ $\left(\text{if }(\operatorname{mch\_} \leq 21)\text{ then }2\text{ else }1\right)$
\par\medskip\noindent\textbf{Definition} (\texttt{bind-sig}).\;For all $\operatorname{mch\_}$, $\operatorname{mcnb\_}$ and $\operatorname{mcn\_}$, 
\begin{equation*}
\begin{array}{@{}l@{}} \operatorname{bind\text{-}sig}(\operatorname{mch\_}, \operatorname{mcnb\_}, \operatorname{mcn\_}) \\ \Leftrightarrow\; (20 \leq \operatorname{mch\_}) \wedge ((\operatorname{mch\_} \leq 21) \wedge ((\operatorname{mcnb\_} = 1) \wedge (\operatorname{mcn\_} = 0))) \vee \\ ((\operatorname{mch\_} = 22) \vee ((\operatorname{mch\_} = 24) \vee (\operatorname{mch\_} = 26))) \wedge ((\operatorname{mcnb\_} = 1) \wedge (\operatorname{mcn\_} = 1)) \vee \\ ((\operatorname{mch\_} = 23) \vee (\operatorname{mch\_} = 25)) \wedge ((\operatorname{mcnb\_} = 1) \wedge (\operatorname{mcn\_} = 0)) \vee \\ (\operatorname{mch\_} = 27) \wedge ((\operatorname{mcnb\_} = 0) \wedge (\operatorname{mcn\_} = 2)) \end{array}
\end{equation*}

\par\medskip\noindent\textbf{Definition} (\texttt{hd-body}).\;$\operatorname{hd\text{-}body}(\operatorname{mch\_}) :=$ $\left(\text{if }((\operatorname{mch\_} = 24) \vee (\operatorname{mch\_} = 26))\text{ then }1\text{ else }2\right)$
\par\medskip\noindent\textbf{Definition} (\texttt{wf-term}).\;For every $\operatorname{mce\_}$,  $\operatorname{wf\text{-}term}(\operatorname{mce\_}) \Leftrightarrow (1 \in \operatorname{expr\text{-}wks}(\operatorname{mce\_}))$
\par\medskip\noindent\textbf{Definition} (\texttt{wf-fml}).\;For every $\operatorname{mce\_}$,  $\operatorname{wf\text{-}fml}(\operatorname{mce\_}) \Leftrightarrow (2 \in \operatorname{expr\text{-}wks}(\operatorname{mce\_}))$
\par\medskip\noindent\textbf{Definition} (\texttt{clf}).\; $(\mathit{clf} \simeq \{\operatorname{mce\_} \in \operatorname{expr\text{-}set} : (\operatorname{wf\text{-}fml}(\operatorname{mce\_}) \wedge \operatorname{expr\text{-}lc}(\operatorname{mce\_}))\})$
\par\medskip\noindent\textbf{Definition} (\texttt{clt}).\; $(\mathit{clt} \simeq \{\operatorname{mce\_} \in \operatorname{expr\text{-}set} : (\operatorname{wf\text{-}term}(\operatorname{mce\_}) \wedge \operatorname{expr\text{-}lc}(\operatorname{mce\_}))\})$
\par\medskip\noindent\textbf{Definition} (\texttt{ctx-set}).\; $(\operatorname{ctx\text{-}set} \simeq \{\operatorname{mcg\_} \in \operatorname{power}(\mathit{clf}) : (|\operatorname{mcg\_}| \in \mathbb{N})\})$
\par\medskip\noindent\textbf{Definition} (\texttt{seq-set}).\; $(\operatorname{seq\text{-}set} \simeq (\operatorname{ctx\text{-}set} \times \mathit{clf}))$
\par\medskip\noindent\textbf{Definition} (\texttt{ctx-add}).\;$\operatorname{ctx\text{-}add}(\operatorname{mcg\_}, \operatorname{mca\_}) :=$ $(\operatorname{mcg\_} \cup \{\operatorname{mca\_}\})$
\par\medskip\noindent\textbf{Definition} (\texttt{ctx-fv}).\;$\operatorname{ctx\text{-}fv}(\operatorname{mcg\_}) :=$ $\operatorname{big\text{-}union}(\operatorname{mcw\_}, \operatorname{mcg\_}, \operatorname{expr\text{-}fv}(\operatorname{mcw\_}))$
\par\medskip\noindent\textbf{Definition} (\texttt{core-closed}).\;For all $\operatorname{mct\_}$ and $\operatorname{mcs\_}$, 
\begin{equation*}
\begin{array}{@{}l@{}} \operatorname{core\text{-}closed}(\operatorname{mct\_}, \operatorname{mcs\_}) \\ \Leftrightarrow\; \forall \operatorname{mchs\_}, \operatorname{mcc\_}. \\ \quad \operatorname{core\text{-}rule}(\operatorname{mct\_}, \operatorname{mchs\_}, \operatorname{mcc\_}) \Rightarrow  \\ \quad \operatorname{hyps\text{-}in}(\operatorname{mchs\_}, \operatorname{mcs\_}) \Rightarrow ((\operatorname{mcc\_} \in \operatorname{seq\text{-}set}) \Rightarrow (\operatorname{mcc\_} \in \operatorname{mcs\_})) \end{array}
\end{equation*}

\par\medskip\noindent\textbf{Definition} (\texttt{core-rule}).\;For all $\operatorname{mct\_}$, $\operatorname{mchs\_}$ and $\operatorname{mcc\_}$, 
\begin{equation*}
\begin{array}{@{}l@{}} \operatorname{core\text{-}rule}(\operatorname{mct\_}, \operatorname{mchs\_}, \operatorname{mcc\_}) \\ \Leftrightarrow\; \operatorname{cr\text{-}and\text{-}intro}(\operatorname{mct\_}, \operatorname{mchs\_}, \operatorname{mcc\_}) \vee \\ \operatorname{cr\text{-}and\text{-}elim}(\operatorname{mct\_}, \operatorname{mchs\_}, \operatorname{mcc\_}) \vee \\ \operatorname{cr\text{-}or\text{-}intro\text{-}left}(\operatorname{mct\_}, \operatorname{mchs\_}, \operatorname{mcc\_}) \vee \\ \operatorname{cr\text{-}or\text{-}intro\text{-}right}(\operatorname{mct\_}, \operatorname{mchs\_}, \operatorname{mcc\_}) \vee \\ \operatorname{cr\text{-}or\text{-}elim}(\operatorname{mct\_}, \operatorname{mchs\_}, \operatorname{mcc\_}) \vee \\ \operatorname{cr\text{-}implies\text{-}intro}(\operatorname{mct\_}, \operatorname{mchs\_}, \operatorname{mcc\_}) \vee \\ \operatorname{cr\text{-}backchain}(\operatorname{mct\_}, \operatorname{mchs\_}, \operatorname{mcc\_}) \vee \\ \operatorname{cr\text{-}not\text{-}intro}(\operatorname{mct\_}, \operatorname{mchs\_}, \operatorname{mcc\_}) \vee \\ \operatorname{cr\text{-}not\text{-}elim}(\operatorname{mct\_}, \operatorname{mchs\_}, \operatorname{mcc\_}) \vee \\ \operatorname{cr\text{-}iff\text{-}intro}(\operatorname{mct\_}, \operatorname{mchs\_}, \operatorname{mcc\_}) \vee \\ \operatorname{cr\text{-}iff\text{-}elim}(\operatorname{mct\_}, \operatorname{mchs\_}, \operatorname{mcc\_}) \vee \\ \operatorname{cr\text{-}forall\text{-}intro}(\operatorname{mct\_}, \operatorname{mchs\_}, \operatorname{mcc\_}) \vee \\ \operatorname{cr\text{-}forall\text{-}elim}(\operatorname{mct\_}, \operatorname{mchs\_}, \operatorname{mcc\_}) \vee \\ \operatorname{cr\text{-}forsome\text{-}intro}(\operatorname{mct\_}, \operatorname{mchs\_}, \operatorname{mcc\_}) \vee \\ \operatorname{cr\text{-}forsome\text{-}elim}(\operatorname{mct\_}, \operatorname{mchs\_}, \operatorname{mcc\_}) \vee \\ \operatorname{cr\text{-}assumption}(\operatorname{mct\_}, \operatorname{mchs\_}, \operatorname{mcc\_}) \vee \\ \operatorname{cr\text{-}theorem\text{-}assumption}(\operatorname{mct\_}, \operatorname{mchs\_}, \operatorname{mcc\_}) \vee \\ \operatorname{cr\text{-}cut}(\operatorname{mct\_}, \operatorname{mchs\_}, \operatorname{mcc\_}) \vee \\ \operatorname{cr\text{-}weakening}(\operatorname{mct\_}, \operatorname{mchs\_}, \operatorname{mcc\_}) \vee \\ \operatorname{cr\text{-}contradiction}(\operatorname{mct\_}, \operatorname{mchs\_}, \operatorname{mcc\_}) \vee \\ \operatorname{cr\text{-}truth\text{-}intro}(\operatorname{mct\_}, \operatorname{mchs\_}, \operatorname{mcc\_}) \vee \\ \operatorname{cr\text{-}eq\text{-}subst}(\operatorname{mct\_}, \operatorname{mchs\_}, \operatorname{mcc\_}) \vee \\ \operatorname{cr\text{-}eq\text{-}subst\text{-}sym}(\operatorname{mct\_}, \operatorname{mchs\_}, \operatorname{mcc\_}) \vee \\ \operatorname{cr\text{-}reflexivity}(\operatorname{mct\_}, \operatorname{mchs\_}, \operatorname{mcc\_}) \vee \\ \operatorname{cr\text{-}quasi\text{-}reflexivity}(\operatorname{mct\_}, \operatorname{mchs\_}, \operatorname{mcc\_}) \vee \\ \operatorname{cr\text{-}union\text{-}mem}(\operatorname{mct\_}, \operatorname{mchs\_}, \operatorname{mcc\_}) \vee \\ \operatorname{cr\text{-}intersection\text{-}mem}(\operatorname{mct\_}, \operatorname{mchs\_}, \operatorname{mcc\_}) \vee \\ \operatorname{cr\text{-}cartesian\text{-}mem}(\operatorname{mct\_}, \operatorname{mchs\_}, \operatorname{mcc\_}) \vee \\ \operatorname{cr\text{-}tuples\text{-}mem}(\operatorname{mct\_}, \operatorname{mchs\_}, \operatorname{mcc\_}) \vee \\ \operatorname{cr\text{-}nth\text{-}eq}(\operatorname{mct\_}, \operatorname{mchs\_}, \operatorname{mcc\_}) \vee \\ \operatorname{cr\text{-}length\text{-}eq}(\operatorname{mct\_}, \operatorname{mchs\_}, \operatorname{mcc\_}) \vee \\ \operatorname{cr\text{-}tuple\text{-}eq}(\operatorname{mct\_}, \operatorname{mchs\_}, \operatorname{mcc\_}) \vee \\ \operatorname{cr\text{-}if\text{-}true}(\operatorname{mct\_}, \operatorname{mchs\_}, \operatorname{mcc\_}) \vee \\ \operatorname{cr\text{-}if\text{-}false}(\operatorname{mct\_}, \operatorname{mchs\_}, \operatorname{mcc\_}) \vee \\ \operatorname{cr\text{-}sep\text{-}sethood}(\operatorname{mct\_}, \operatorname{mchs\_}, \operatorname{mcc\_}) \vee \\ \operatorname{cr\text{-}sep\text{-}mem}(\operatorname{mct\_}, \operatorname{mchs\_}, \operatorname{mcc\_}) \vee \\ \operatorname{cr\text{-}comp\text{-}mem}(\operatorname{mct\_}, \operatorname{mchs\_}, \operatorname{mcc\_}) \vee \\ \operatorname{cr\text{-}big\text{-}union\text{-}sethood}(\operatorname{mct\_}, \operatorname{mchs\_}, \operatorname{mcc\_}) \vee \\ \operatorname{cr\text{-}big\text{-}union\text{-}mem}(\operatorname{mct\_}, \operatorname{mchs\_}, \operatorname{mcc\_}) \vee \\ \operatorname{cr\text{-}iota\text{-}def}(\operatorname{mct\_}, \operatorname{mchs\_}, \operatorname{mcc\_}) \vee \\ \operatorname{cr\text{-}iota\text{-}in\text{-}elim}(\operatorname{mct\_}, \operatorname{mchs\_}, \operatorname{mcc\_}) \vee \\ \operatorname{cr\text{-}lambda\text{-}type}(\operatorname{mct\_}, \operatorname{mchs\_}, \operatorname{mcc\_}) \vee \\ \operatorname{cr\text{-}lambda\text{-}beta\text{-}eq}(\operatorname{mct\_}, \operatorname{mchs\_}, \operatorname{mcc\_}) \vee \\ \operatorname{cr\text{-}def\text{-}strict\text{-}in\text{-}l}(\operatorname{mct\_}, \operatorname{mchs\_}, \operatorname{mcc\_}) \vee \\ \operatorname{cr\text{-}def\text{-}strict\text{-}in\text{-}r}(\operatorname{mct\_}, \operatorname{mchs\_}, \operatorname{mcc\_}) \vee \\ \operatorname{cr\text{-}def\text{-}total}(\operatorname{mct\_}, \operatorname{mchs\_}, \operatorname{mcc\_}) \vee \\ \operatorname{cr\text{-}def\text{-}binder}(\operatorname{mct\_}, \operatorname{mchs\_}, \operatorname{mcc\_}) \vee \\ \operatorname{cr\text{-}def\text{-}comp}(\operatorname{mct\_}, \operatorname{mchs\_}, \operatorname{mcc\_}) \vee \\ \operatorname{cr\text{-}def\text{-}list}(\operatorname{mct\_}, \operatorname{mchs\_}, \operatorname{mcc\_}) \vee \\ \operatorname{cr\text{-}def\text{-}strict\text{-}arg}(\operatorname{mct\_}, \operatorname{mchs\_}, \operatorname{mcc\_}) \vee \\ \operatorname{cr\text{-}def\text{-}strict\text{-}dom}(\operatorname{mct\_}, \operatorname{mchs\_}, \operatorname{mcc\_}) \vee \\ \operatorname{cr\text{-}def\text{-}strict\text{-}pred}(\operatorname{mct\_}, \operatorname{mchs\_}, \operatorname{mcc\_}) \end{array}
\end{equation*}

\par\medskip\noindent\textbf{Definition} (\texttt{wks-stage}).\; $(\operatorname{wks\text{-}stage}(0) \simeq (\lambda \operatorname{mce\_} \in \operatorname{expr\text{-}stage}(0).\; \{1\}))$\par
Suppose $\operatorname{mcn\_} \in \mathbb{N}$. Then
\begin{equation*}
(\operatorname{wks\text{-}stage}(\operatorname{succ}(\operatorname{mcn\_})) \simeq (\lambda \operatorname{mce\_} \in \operatorname{expr\text{-}stage}(\operatorname{succ}(\operatorname{mcn\_})).\; \left(\text{if }(\operatorname{mce\_} \in \operatorname{expr\text{-}stage}(\operatorname{mcn\_}))\text{ then }\operatorname{wks\text{-}stage}(\operatorname{mcn\_})(\operatorname{mce\_})\text{ else }\left(\text{if }(\operatorname{nth}(1, \operatorname{mce\_}) = 3)\text{ then }\{\operatorname{mckk\_} \in \mathbb{N} : ((\operatorname{mckk\_} = \operatorname{hd\text{-}res}(\operatorname{nth}(2, \operatorname{mce\_}))) \wedge (\operatorname{app\text{-}ar\text{-}ok}(\operatorname{nth}(2, \operatorname{mce\_}), \operatorname{length}(\operatorname{nth}(3, \operatorname{mce\_}))) \wedge \forall \operatorname{mcw\_} \in \operatorname{make\text{-}set}(\operatorname{nth}(3, \operatorname{mce\_})).\; (\operatorname{hd\text{-}argk}(\operatorname{nth}(2, \operatorname{mce\_})) \in \operatorname{wks\text{-}stage}(\operatorname{mcn\_})(\operatorname{mcw\_}))))\}\text{ else }\{\operatorname{mckk\_} \in \mathbb{N} : ((\operatorname{mckk\_} = \operatorname{hd\text{-}bres}(\operatorname{nth}(2, \operatorname{mce\_}))) \wedge (\operatorname{bind\text{-}sig}(\operatorname{nth}(2, \operatorname{mce\_}), \operatorname{nth}(3, \operatorname{mce\_}), \operatorname{length}(\operatorname{nth}(4, \operatorname{mce\_}))) \wedge (\forall \operatorname{mcw\_} \in \operatorname{make\text{-}set}(\operatorname{nth}(4, \operatorname{mce\_})).\; (1 \in \operatorname{wks\text{-}stage}(\operatorname{mcn\_})(\operatorname{mcw\_})) \wedge (\operatorname{hd\text{-}body}(\operatorname{nth}(2, \operatorname{mce\_})) \in \operatorname{wks\text{-}stage}(\operatorname{mcn\_})(\operatorname{nth}(5, \operatorname{mce\_}))))))\}\right)\right)))
\end{equation*}

\par\medskip\noindent\textbf{Definition} (\texttt{expr-wks}).\;$\operatorname{expr\text{-}wks}(\operatorname{mce\_}) :=$ $\operatorname{wks\text{-}stage}(\operatorname{expr\text{-}rank}(\operatorname{mce\_}))(\operatorname{mce\_})$
\par\medskip\noindent\textbf{Definition} (\texttt{enc-truth}).\; $(\operatorname{enc\text{-}truth} \simeq \operatorname{eapp}(0, \langle  \rangle))$
\par\medskip\noindent\textbf{Definition} (\texttt{enc-falsity}).\; $(\operatorname{enc\text{-}falsity} \simeq \operatorname{eapp}(1, \langle  \rangle))$
\par\medskip\noindent\textbf{Definition} (\texttt{enc-not}).\;$\operatorname{enc\text{-}not}(\operatorname{mcp\_}) :=$ $\operatorname{eapp}(2, \langle \operatorname{mcp\_} \rangle)$
\par\medskip\noindent\textbf{Definition} (\texttt{enc-and}).\;$\operatorname{enc\text{-}and}(\operatorname{mcp\_}, \operatorname{mcq\_}) :=$ $\operatorname{eapp}(3, \langle \operatorname{mcp\_}, \operatorname{mcq\_} \rangle)$
\par\medskip\noindent\textbf{Definition} (\texttt{enc-or}).\;$\operatorname{enc\text{-}or}(\operatorname{mcp\_}, \operatorname{mcq\_}) :=$ $\operatorname{eapp}(4, \langle \operatorname{mcp\_}, \operatorname{mcq\_} \rangle)$
\par\medskip\noindent\textbf{Definition} (\texttt{enc-imp}).\;$\operatorname{enc\text{-}imp}(\operatorname{mcp\_}, \operatorname{mcq\_}) :=$ $\operatorname{eapp}(5, \langle \operatorname{mcp\_}, \operatorname{mcq\_} \rangle)$
\par\medskip\noindent\textbf{Definition} (\texttt{enc-iff}).\;$\operatorname{enc\text{-}iff}(\operatorname{mcp\_}, \operatorname{mcq\_}) :=$ $\operatorname{eapp}(6, \langle \operatorname{mcp\_}, \operatorname{mcq\_} \rangle)$
\par\medskip\noindent\textbf{Definition} (\texttt{enc-in}).\;$\operatorname{enc\text{-}in}(\operatorname{mcp\_}, \operatorname{mcq\_}) :=$ $\operatorname{eapp}(7, \langle \operatorname{mcp\_}, \operatorname{mcq\_} \rangle)$
\par\medskip\noindent\textbf{Definition} (\texttt{enc-eq}).\;$\operatorname{enc\text{-}eq}(\operatorname{mcp\_}, \operatorname{mcq\_}) :=$ $\operatorname{eapp}(8, \langle \operatorname{mcp\_}, \operatorname{mcq\_} \rangle)$
\par\medskip\noindent\textbf{Definition} (\texttt{enc-qeq}).\;$\operatorname{enc\text{-}qeq}(\operatorname{mcp\_}, \operatorname{mcq\_}) :=$ $\operatorname{eapp}(9, \langle \operatorname{mcp\_}, \operatorname{mcq\_} \rangle)$
\par\medskip\noindent\textbf{Definition} (\texttt{enc-union}).\;$\operatorname{enc\text{-}union}(\operatorname{mcl\_}) :=$ $\operatorname{eapp}(11, \operatorname{mcl\_})$
\par\medskip\noindent\textbf{Definition} (\texttt{enc-inter}).\;$\operatorname{enc\text{-}inter}(\operatorname{mcl\_}) :=$ $\operatorname{eapp}(12, \operatorname{mcl\_})$
\par\medskip\noindent\textbf{Definition} (\texttt{enc-cart}).\;$\operatorname{enc\text{-}cart}(\operatorname{mcl\_}) :=$ $\operatorname{eapp}(13, \operatorname{mcl\_})$
\par\medskip\noindent\textbf{Definition} (\texttt{enc-tuples}).\;$\operatorname{enc\text{-}tuples}(\operatorname{mcp\_}) :=$ $\operatorname{eapp}(14, \langle \operatorname{mcp\_} \rangle)$
\par\medskip\noindent\textbf{Definition} (\texttt{enc-nth}).\;$\operatorname{enc\text{-}nth}(\operatorname{mcp\_}, \operatorname{mcq\_}) :=$ $\operatorname{eapp}(15, \langle \operatorname{mcp\_}, \operatorname{mcq\_} \rangle)$
\par\medskip\noindent\textbf{Definition} (\texttt{enc-length}).\;$\operatorname{enc\text{-}length}(\operatorname{mcp\_}) :=$ $\operatorname{eapp}(16, \langle \operatorname{mcp\_} \rangle)$
\par\medskip\noindent\textbf{Definition} (\texttt{enc-list}).\;$\operatorname{enc\text{-}list}(\operatorname{mcl\_}) :=$ $\operatorname{eapp}(17, \operatorname{mcl\_})$
\par\medskip\noindent\textbf{Definition} (\texttt{enc-apply}).\;$\operatorname{enc\text{-}apply}(\operatorname{mcp\_}, \operatorname{mcq\_}) :=$ $\operatorname{eapp}(18, \langle \operatorname{mcp\_}, \operatorname{mcq\_} \rangle)$
\par\medskip\noindent\textbf{Definition} (\texttt{enc-fun}).\;$\operatorname{enc\text{-}fun}(\operatorname{mcp\_}, \operatorname{mcq\_}) :=$ $\operatorname{eapp}(19, \langle \operatorname{mcp\_}, \operatorname{mcq\_} \rangle)$
\par\medskip\noindent\textbf{Definition} (\texttt{enc-all}).\;$\operatorname{enc\text{-}all}(\operatorname{mcb\_}) :=$ $\operatorname{ebind}(20, 1, \langle  \rangle, \operatorname{mcb\_})$
\par\medskip\noindent\textbf{Definition} (\texttt{enc-ex}).\;$\operatorname{enc\text{-}ex}(\operatorname{mcb\_}) :=$ $\operatorname{ebind}(21, 1, \langle  \rangle, \operatorname{mcb\_})$
\par\medskip\noindent\textbf{Definition} (\texttt{enc-sep}).\;$\operatorname{enc\text{-}sep}(\operatorname{mca\_}, \operatorname{mcb\_}) :=$ $\operatorname{ebind}(22, 1, \langle \operatorname{mca\_} \rangle, \operatorname{mcb\_})$
\par\medskip\noindent\textbf{Definition} (\texttt{enc-comp}).\;$\operatorname{enc\text{-}comp}(\operatorname{mcb\_}) :=$ $\operatorname{ebind}(23, 1, \langle  \rangle, \operatorname{mcb\_})$
\par\medskip\noindent\textbf{Definition} (\texttt{enc-bu}).\;$\operatorname{enc\text{-}bu}(\operatorname{mca\_}, \operatorname{mcb\_}) :=$ $\operatorname{ebind}(24, 1, \langle \operatorname{mca\_} \rangle, \operatorname{mcb\_})$
\par\medskip\noindent\textbf{Definition} (\texttt{enc-iota}).\;$\operatorname{enc\text{-}iota}(\operatorname{mcb\_}) :=$ $\operatorname{ebind}(25, 1, \langle  \rangle, \operatorname{mcb\_})$
\par\medskip\noindent\textbf{Definition} (\texttt{enc-lam}).\;$\operatorname{enc\text{-}lam}(\operatorname{mca\_}, \operatorname{mcb\_}) :=$ $\operatorname{ebind}(26, 1, \langle \operatorname{mca\_} \rangle, \operatorname{mcb\_})$
\par\medskip\noindent\textbf{Definition} (\texttt{enc-if}).\;$\operatorname{enc\text{-}if}(\operatorname{mcp\_}, \operatorname{mca\_}, \operatorname{mcb\_}) :=$ $\operatorname{ebind}(27, 0, \langle \operatorname{mca\_}, \operatorname{mcb\_} \rangle, \operatorname{mcp\_})$
\par\medskip\noindent\textbf{Definition} (\texttt{enc-orl}).\;For all $\operatorname{mcx\_}$ and $\operatorname{mcl\_}$,  $(\operatorname{enc\text{-}orl}(\operatorname{mcx\_}, \operatorname{mcl\_}, 0) \simeq \operatorname{enc\text{-}falsity})$\par
Let $\operatorname{mcx\_}$ and $\operatorname{mcl\_}$ be arbitrary. Let $\operatorname{mcn\_} \in \mathbb{N}$. Then
\begin{equation*}
(\operatorname{enc\text{-}orl}(\operatorname{mcx\_}, \operatorname{mcl\_}, \operatorname{succ}(\operatorname{mcn\_})) \simeq \operatorname{enc\text{-}or}(\operatorname{enc\text{-}orl}(\operatorname{mcx\_}, \operatorname{mcl\_}, \operatorname{mcn\_}), \operatorname{enc\text{-}in}(\operatorname{mcx\_}, \operatorname{nth}(\operatorname{succ}(\operatorname{mcn\_}), \operatorname{mcl\_}))))
\end{equation*}

\par\medskip\noindent\textbf{Definition} (\texttt{enc-andl}).\;For all $\operatorname{mcx\_}$ and $\operatorname{mcl\_}$,  $(\operatorname{enc\text{-}andl}(\operatorname{mcx\_}, \operatorname{mcl\_}, 0) \simeq \operatorname{enc\text{-}truth})$\par
Let $\operatorname{mcx\_}$ and $\operatorname{mcl\_}$ be arbitrary. Let $\operatorname{mcn\_} \in \mathbb{N}$. Then
\begin{equation*}
(\operatorname{enc\text{-}andl}(\operatorname{mcx\_}, \operatorname{mcl\_}, \operatorname{succ}(\operatorname{mcn\_})) \simeq \operatorname{enc\text{-}and}(\operatorname{enc\text{-}andl}(\operatorname{mcx\_}, \operatorname{mcl\_}, \operatorname{mcn\_}), \operatorname{enc\text{-}in}(\operatorname{mcx\_}, \operatorname{nth}(\operatorname{succ}(\operatorname{mcn\_}), \operatorname{mcl\_}))))
\end{equation*}

\par\medskip\noindent\textbf{Definition} (\texttt{enc-ande}).\;For all $\operatorname{mcl\_}$ and $\operatorname{mca\_}$,  $(\operatorname{enc\text{-}ande}(\operatorname{mcl\_}, \operatorname{mca\_}, 0) \simeq \operatorname{enc\text{-}truth})$\par
Let $\operatorname{mcl\_}$ and $\operatorname{mca\_}$ be arbitrary. Let $\operatorname{mcn\_} \in \mathbb{N}$. Then
\begin{equation*}
(\operatorname{enc\text{-}ande}(\operatorname{mcl\_}, \operatorname{mca\_}, \operatorname{succ}(\operatorname{mcn\_})) \simeq \operatorname{enc\text{-}and}(\operatorname{enc\text{-}ande}(\operatorname{mcl\_}, \operatorname{mca\_}, \operatorname{mcn\_}), \operatorname{enc\text{-}in}(\operatorname{nth}(\operatorname{succ}(\operatorname{mcn\_}), \operatorname{mcl\_}), \operatorname{mca\_})))
\end{equation*}

\par\medskip\noindent\textbf{Definition} (\texttt{enc-cartl}).\;For all $\operatorname{mcx\_}$ and $\operatorname{mcl\_}$,  $(\operatorname{enc\text{-}cartl}(\operatorname{mcx\_}, \operatorname{mcl\_}, 0) \simeq \operatorname{enc\text{-}truth})$\par
Let $\operatorname{mcx\_}$ and $\operatorname{mcl\_}$ be arbitrary. Let $\operatorname{mcn\_} \in \mathbb{N}$. Then
\begin{equation*}
(\operatorname{enc\text{-}cartl}(\operatorname{mcx\_}, \operatorname{mcl\_}, \operatorname{succ}(\operatorname{mcn\_})) \simeq \operatorname{enc\text{-}and}(\operatorname{enc\text{-}cartl}(\operatorname{mcx\_}, \operatorname{mcl\_}, \operatorname{mcn\_}), \operatorname{enc\text{-}in}(\operatorname{enc\text{-}nth}(\operatorname{enc\text{-}num}(\operatorname{succ}(\operatorname{mcn\_})), \operatorname{mcx\_}), \operatorname{nth}(\operatorname{succ}(\operatorname{mcn\_}), \operatorname{mcl\_}))))
\end{equation*}

\par\medskip\noindent\textbf{Definition} (\texttt{enc-eql}).\;For all $\operatorname{mcl\_}$ and $\operatorname{mca\_}$,  $(\operatorname{enc\text{-}eql}(\operatorname{mcl\_}, \operatorname{mca\_}, 0) \simeq \operatorname{enc\text{-}truth})$\par
Let $\operatorname{mcl\_}$ and $\operatorname{mca\_}$ be arbitrary. Let $\operatorname{mcn\_} \in \mathbb{N}$. Then
\begin{equation*}
(\operatorname{enc\text{-}eql}(\operatorname{mcl\_}, \operatorname{mca\_}, \operatorname{succ}(\operatorname{mcn\_})) \simeq \operatorname{enc\text{-}and}(\operatorname{enc\text{-}eql}(\operatorname{mcl\_}, \operatorname{mca\_}, \operatorname{mcn\_}), \operatorname{enc\text{-}eq}(\operatorname{nth}(\operatorname{succ}(\operatorname{mcn\_}), \operatorname{mcl\_}), \operatorname{nth}(\operatorname{succ}(\operatorname{mcn\_}), \operatorname{mca\_}))))
\end{equation*}

\par\medskip\noindent\textbf{Definition} (\texttt{cr-and-intro}).\;For all $\operatorname{mct\_}$, $\operatorname{mchs\_}$ and $\operatorname{mcc\_}$, 
\begin{equation*}
\begin{array}{@{}l@{}} \operatorname{cr\text{-}and\text{-}intro}(\operatorname{mct\_}, \operatorname{mchs\_}, \operatorname{mcc\_}) \\ \Leftrightarrow\; \exists \operatorname{mcg\_}, \operatorname{mca\_}, \operatorname{mcb\_}. \\ \quad (\operatorname{mcc\_} = \langle \operatorname{mcg\_}, \operatorname{enc\text{-}and}(\operatorname{mca\_}, \operatorname{mcb\_}) \rangle) \wedge (\operatorname{mchs\_} = \langle \langle \operatorname{mcg\_}, \operatorname{mca\_} \rangle, \langle \operatorname{mcg\_}, \operatorname{mcb\_} \rangle \rangle) \end{array}
\end{equation*}

\par\medskip\noindent\textbf{Definition} (\texttt{cr-and-elim}).\;For all $\operatorname{mct\_}$, $\operatorname{mchs\_}$ and $\operatorname{mcc\_}$, 
\begin{equation*}
\begin{array}{@{}l@{}} \operatorname{cr\text{-}and\text{-}elim}(\operatorname{mct\_}, \operatorname{mchs\_}, \operatorname{mcc\_}) \\ \Leftrightarrow\; \exists \operatorname{mcg\_}, \operatorname{mcf\_}, \operatorname{mca\_}, \operatorname{mcb\_}. \\ \quad \operatorname{mcc\_} = \langle \operatorname{mcg\_}, \operatorname{mcf\_} \rangle \wedge \\ \quad \operatorname{enc\text{-}and}(\operatorname{mca\_}, \operatorname{mcb\_}) \in \operatorname{mcg\_} \wedge \\ \quad \operatorname{mchs\_} = \langle \langle \operatorname{ctx\text{-}add}(\operatorname{ctx\text{-}add}(\operatorname{mcg\_}, \operatorname{mca\_}), \operatorname{mcb\_}), \operatorname{mcf\_} \rangle \rangle \end{array}
\end{equation*}

\par\medskip\noindent\textbf{Definition} (\texttt{cr-or-intro-left}).\;For all $\operatorname{mct\_}$, $\operatorname{mchs\_}$ and $\operatorname{mcc\_}$, 
\begin{equation*}
\begin{array}{@{}l@{}} \operatorname{cr\text{-}or\text{-}intro\text{-}left}(\operatorname{mct\_}, \operatorname{mchs\_}, \operatorname{mcc\_}) \\ \Leftrightarrow\; \exists \operatorname{mcg\_}, \operatorname{mca\_}, \operatorname{mcb\_}. \\ \quad (\operatorname{mcc\_} = \langle \operatorname{mcg\_}, \operatorname{enc\text{-}or}(\operatorname{mca\_}, \operatorname{mcb\_}) \rangle) \wedge (\operatorname{mchs\_} = \langle \langle \operatorname{mcg\_}, \operatorname{mca\_} \rangle \rangle) \end{array}
\end{equation*}

\par\medskip\noindent\textbf{Definition} (\texttt{cr-or-intro-right}).\;For all $\operatorname{mct\_}$, $\operatorname{mchs\_}$ and $\operatorname{mcc\_}$, 
\begin{equation*}
\begin{array}{@{}l@{}} \operatorname{cr\text{-}or\text{-}intro\text{-}right}(\operatorname{mct\_}, \operatorname{mchs\_}, \operatorname{mcc\_}) \\ \Leftrightarrow\; \exists \operatorname{mcg\_}, \operatorname{mca\_}, \operatorname{mcb\_}. \\ \quad (\operatorname{mcc\_} = \langle \operatorname{mcg\_}, \operatorname{enc\text{-}or}(\operatorname{mca\_}, \operatorname{mcb\_}) \rangle) \wedge (\operatorname{mchs\_} = \langle \langle \operatorname{mcg\_}, \operatorname{mcb\_} \rangle \rangle) \end{array}
\end{equation*}

\par\medskip\noindent\textbf{Definition} (\texttt{cr-or-elim}).\;For all $\operatorname{mct\_}$, $\operatorname{mchs\_}$ and $\operatorname{mcc\_}$, 
\begin{equation*}
\begin{array}{@{}l@{}} \operatorname{cr\text{-}or\text{-}elim}(\operatorname{mct\_}, \operatorname{mchs\_}, \operatorname{mcc\_}) \\ \Leftrightarrow\; \exists \operatorname{mcg\_}, \operatorname{mcf\_}, \operatorname{mca\_}, \operatorname{mcb\_}. \\ \quad \operatorname{mcc\_} = \langle \operatorname{mcg\_}, \operatorname{mcf\_} \rangle \wedge \\ \quad \operatorname{enc\text{-}or}(\operatorname{mca\_}, \operatorname{mcb\_}) \in \operatorname{mcg\_} \wedge \\ \quad \operatorname{mchs\_} = \langle \langle \operatorname{ctx\text{-}add}(\operatorname{mcg\_}, \operatorname{mca\_}), \operatorname{mcf\_} \rangle, \langle \operatorname{ctx\text{-}add}(\operatorname{mcg\_}, \operatorname{mcb\_}), \operatorname{mcf\_} \rangle \rangle \end{array}
\end{equation*}

\par\medskip\noindent\textbf{Definition} (\texttt{cr-implies-intro}).\;For all $\operatorname{mct\_}$, $\operatorname{mchs\_}$ and $\operatorname{mcc\_}$, 
\begin{equation*}
\begin{array}{@{}l@{}} \operatorname{cr\text{-}implies\text{-}intro}(\operatorname{mct\_}, \operatorname{mchs\_}, \operatorname{mcc\_}) \\ \Leftrightarrow\; \exists \operatorname{mcg\_}, \operatorname{mca\_}, \operatorname{mcb\_}. \\ \quad (\operatorname{mcc\_} = \langle \operatorname{mcg\_}, \operatorname{enc\text{-}imp}(\operatorname{mca\_}, \operatorname{mcb\_}) \rangle) \wedge (\operatorname{mchs\_} = \langle \langle \operatorname{ctx\text{-}add}(\operatorname{mcg\_}, \operatorname{mca\_}), \operatorname{mcb\_} \rangle \rangle) \end{array}
\end{equation*}

\par\medskip\noindent\textbf{Definition} (\texttt{cr-backchain}).\;For all $\operatorname{mct\_}$, $\operatorname{mchs\_}$ and $\operatorname{mcc\_}$, 
\begin{equation*}
\begin{array}{@{}l@{}} \operatorname{cr\text{-}backchain}(\operatorname{mct\_}, \operatorname{mchs\_}, \operatorname{mcc\_}) \\ \Leftrightarrow\; \exists \operatorname{mcg\_}, \operatorname{mcf\_}, \operatorname{mca\_}. \\ \quad (\operatorname{mcc\_} = \langle \operatorname{mcg\_}, \operatorname{mcf\_} \rangle) \wedge ((\operatorname{enc\text{-}imp}(\operatorname{mca\_}, \operatorname{mcf\_}) \in \operatorname{mcg\_}) \wedge (\operatorname{mchs\_} = \langle \langle \operatorname{mcg\_}, \operatorname{mca\_} \rangle \rangle)) \end{array}
\end{equation*}

\par\medskip\noindent\textbf{Definition} (\texttt{cr-not-intro}).\;For all $\operatorname{mct\_}$, $\operatorname{mchs\_}$ and $\operatorname{mcc\_}$, 
\begin{equation*}
\begin{array}{@{}l@{}} \operatorname{cr\text{-}not\text{-}intro}(\operatorname{mct\_}, \operatorname{mchs\_}, \operatorname{mcc\_}) \\ \Leftrightarrow\; \exists \operatorname{mcg\_}, \operatorname{mca\_}. \\ \quad (\operatorname{mcc\_} = \langle \operatorname{mcg\_}, \operatorname{enc\text{-}not}(\operatorname{mca\_}) \rangle) \wedge (\operatorname{mchs\_} = \langle \langle \operatorname{ctx\text{-}add}(\operatorname{mcg\_}, \operatorname{mca\_}), \operatorname{enc\text{-}falsity} \rangle \rangle) \end{array}
\end{equation*}

\par\medskip\noindent\textbf{Definition} (\texttt{cr-not-elim}).\;For all $\operatorname{mct\_}$, $\operatorname{mchs\_}$ and $\operatorname{mcc\_}$, 
\begin{equation*}
\begin{array}{@{}l@{}} \operatorname{cr\text{-}not\text{-}elim}(\operatorname{mct\_}, \operatorname{mchs\_}, \operatorname{mcc\_}) \\ \Leftrightarrow\; \exists \operatorname{mcg\_}, \operatorname{mcf\_}, \operatorname{mca\_}. \\ \quad (\operatorname{mcc\_} = \langle \operatorname{mcg\_}, \operatorname{mcf\_} \rangle) \wedge ((\operatorname{enc\text{-}not}(\operatorname{mca\_}) \in \operatorname{mcg\_}) \wedge ((\operatorname{mca\_} \in \operatorname{mcg\_}) \wedge (\operatorname{mchs\_} = \langle  \rangle))) \end{array}
\end{equation*}

\par\medskip\noindent\textbf{Definition} (\texttt{cr-iff-intro}).\;For all $\operatorname{mct\_}$, $\operatorname{mchs\_}$ and $\operatorname{mcc\_}$, 
\begin{equation*}
\begin{array}{@{}l@{}} \operatorname{cr\text{-}iff\text{-}intro}(\operatorname{mct\_}, \operatorname{mchs\_}, \operatorname{mcc\_}) \\ \Leftrightarrow\; \exists \operatorname{mcg\_}, \operatorname{mca\_}, \operatorname{mcb\_}. \\ \quad \operatorname{mcc\_} = \langle \operatorname{mcg\_}, \operatorname{enc\text{-}iff}(\operatorname{mca\_}, \operatorname{mcb\_}) \rangle \wedge \\ \quad \operatorname{mchs\_} = \langle \langle \operatorname{ctx\text{-}add}(\operatorname{mcg\_}, \operatorname{mca\_}), \operatorname{mcb\_} \rangle, \langle \operatorname{ctx\text{-}add}(\operatorname{mcg\_}, \operatorname{mcb\_}), \operatorname{mca\_} \rangle \rangle \end{array}
\end{equation*}

\par\medskip\noindent\textbf{Definition} (\texttt{cr-iff-elim}).\;For all $\operatorname{mct\_}$, $\operatorname{mchs\_}$ and $\operatorname{mcc\_}$, 
\begin{equation*}
\begin{array}{@{}l@{}} \operatorname{cr\text{-}iff\text{-}elim}(\operatorname{mct\_}, \operatorname{mchs\_}, \operatorname{mcc\_}) \\ \Leftrightarrow\; \exists \operatorname{mcg\_}, \operatorname{mcf\_}, \operatorname{mca\_}, \operatorname{mcb\_}. \\ \quad \operatorname{mcc\_} = \langle \operatorname{mcg\_}, \operatorname{mcf\_} \rangle \wedge \\ \quad \operatorname{enc\text{-}iff}(\operatorname{mca\_}, \operatorname{mcb\_}) \in \operatorname{mcg\_} \wedge \\ \quad \operatorname{mchs\_} = \langle \langle \operatorname{ctx\text{-}add}(\operatorname{ctx\text{-}add}(\operatorname{mcg\_}, \operatorname{enc\text{-}imp}(\operatorname{mca\_}, \operatorname{mcb\_})), \operatorname{enc\text{-}imp}(\operatorname{mcb\_}, \operatorname{mca\_})), \operatorname{mcf\_} \rangle \rangle \end{array}
\end{equation*}

\par\medskip\noindent\textbf{Definition} (\texttt{cr-forall-intro}).\;For all $\operatorname{mct\_}$, $\operatorname{mchs\_}$ and $\operatorname{mcc\_}$, 
\begin{equation*}
\begin{array}{@{}l@{}} \operatorname{cr\text{-}forall\text{-}intro}(\operatorname{mct\_}, \operatorname{mchs\_}, \operatorname{mcc\_}) \\ \Leftrightarrow\; \exists \operatorname{mcg\_}, \operatorname{mcb\_}, \operatorname{mcy\_}. \\ \quad \operatorname{mcc\_} = \langle \operatorname{mcg\_}, \operatorname{enc\text{-}all}(\operatorname{mcb\_}) \rangle \wedge \\ \quad \operatorname{mcy\_} \in \mathbb{N} \wedge \\ \quad \neg (\operatorname{mcy\_} \in (\operatorname{ctx\text{-}fv}(\operatorname{mcg\_}) \cup \operatorname{expr\text{-}fv}(\operatorname{mcb\_}))) \wedge \\ \quad \operatorname{mchs\_} = \langle \langle \operatorname{mcg\_}, \operatorname{expr\text{-}open}(\operatorname{mcb\_}, 0, \operatorname{efvar}(\operatorname{mcy\_})) \rangle \rangle \end{array}
\end{equation*}

\par\medskip\noindent\textbf{Definition} (\texttt{cr-forall-elim}).\;For all $\operatorname{mct\_}$, $\operatorname{mchs\_}$ and $\operatorname{mcc\_}$, 
\begin{equation*}
\begin{array}{@{}l@{}} \operatorname{cr\text{-}forall\text{-}elim}(\operatorname{mct\_}, \operatorname{mchs\_}, \operatorname{mcc\_}) \\ \Leftrightarrow\; \exists \operatorname{mcg\_}, \operatorname{mcf\_}, \operatorname{mcb\_}, \operatorname{mctm\_}. \\ \quad \operatorname{mcc\_} = \langle \operatorname{mcg\_}, \operatorname{mcf\_} \rangle \wedge \\ \quad \operatorname{enc\text{-}all}(\operatorname{mcb\_}) \in \operatorname{mcg\_} \wedge \\ \quad \operatorname{mctm\_} \in \mathit{clt} \wedge \\ \quad \operatorname{mchs\_} = \langle \langle \operatorname{ctx\text{-}add}(\operatorname{mcg\_}, \operatorname{expr\text{-}open}(\operatorname{mcb\_}, 0, \operatorname{mctm\_})), \operatorname{mcf\_} \rangle, \langle \operatorname{mcg\_}, \operatorname{enc\text{-}eq}(\operatorname{mctm\_}, \operatorname{mctm\_}) \rangle \rangle \end{array}
\end{equation*}

\par\medskip\noindent\textbf{Definition} (\texttt{cr-forsome-intro}).\;For all $\operatorname{mct\_}$, $\operatorname{mchs\_}$ and $\operatorname{mcc\_}$, 
\begin{equation*}
\begin{array}{@{}l@{}} \operatorname{cr\text{-}forsome\text{-}intro}(\operatorname{mct\_}, \operatorname{mchs\_}, \operatorname{mcc\_}) \\ \Leftrightarrow\; \exists \operatorname{mcg\_}, \operatorname{mcb\_}, \operatorname{mctm\_}. \\ \quad \operatorname{mcc\_} = \langle \operatorname{mcg\_}, \operatorname{enc\text{-}ex}(\operatorname{mcb\_}) \rangle \wedge \\ \quad \operatorname{mctm\_} \in \mathit{clt} \wedge \\ \quad \operatorname{mchs\_} = \langle \langle \operatorname{mcg\_}, \operatorname{expr\text{-}open}(\operatorname{mcb\_}, 0, \operatorname{mctm\_}) \rangle, \langle \operatorname{mcg\_}, \operatorname{enc\text{-}eq}(\operatorname{mctm\_}, \operatorname{mctm\_}) \rangle \rangle \end{array}
\end{equation*}

\par\medskip\noindent\textbf{Definition} (\texttt{cr-forsome-elim}).\;For all $\operatorname{mct\_}$, $\operatorname{mchs\_}$ and $\operatorname{mcc\_}$, 
\begin{equation*}
\begin{array}{@{}l@{}} \operatorname{cr\text{-}forsome\text{-}elim}(\operatorname{mct\_}, \operatorname{mchs\_}, \operatorname{mcc\_}) \\ \Leftrightarrow\; \exists \operatorname{mcg\_}, \operatorname{mcf\_}, \operatorname{mcb\_}, \operatorname{mcy\_}. \\ \quad \operatorname{mcc\_} = \langle \operatorname{mcg\_}, \operatorname{mcf\_} \rangle \wedge \\ \quad \operatorname{enc\text{-}ex}(\operatorname{mcb\_}) \in \operatorname{mcg\_} \wedge \\ \quad \operatorname{mcy\_} \in \mathbb{N} \wedge \\ \quad \neg (\operatorname{mcy\_} \in (\operatorname{ctx\text{-}fv}(\operatorname{mcg\_}) \cup \operatorname{expr\text{-}fv}(\operatorname{mcf\_}))) \wedge \\ \quad \operatorname{mchs\_} = \langle \langle \operatorname{ctx\text{-}add}(\operatorname{mcg\_}, \operatorname{expr\text{-}open}(\operatorname{mcb\_}, 0, \operatorname{efvar}(\operatorname{mcy\_}))), \operatorname{mcf\_} \rangle \rangle \end{array}
\end{equation*}

\par\medskip\noindent\textbf{Definition} (\texttt{cr-assumption}).\;For all $\operatorname{mct\_}$, $\operatorname{mchs\_}$ and $\operatorname{mcc\_}$, 
\begin{equation*}
\begin{array}{@{}l@{}} \operatorname{cr\text{-}assumption}(\operatorname{mct\_}, \operatorname{mchs\_}, \operatorname{mcc\_}) \\ \Leftrightarrow\; \exists \operatorname{mcg\_}, \operatorname{mca\_}.\; ((\operatorname{mcc\_} = \langle \operatorname{mcg\_}, \operatorname{mca\_} \rangle) \wedge ((\operatorname{mca\_} \in \operatorname{mcg\_}) \wedge (\operatorname{mchs\_} = \langle  \rangle))) \end{array}
\end{equation*}

\par\medskip\noindent\textbf{Definition} (\texttt{cr-theorem-assumption}).\;For all $\operatorname{mct\_}$, $\operatorname{mchs\_}$ and $\operatorname{mcc\_}$, 
\begin{equation*}
\begin{array}{@{}l@{}} \operatorname{cr\text{-}theorem\text{-}assumption}(\operatorname{mct\_}, \operatorname{mchs\_}, \operatorname{mcc\_}) \\ \Leftrightarrow\; \exists \operatorname{mcg\_}, \operatorname{mca\_}.\; ((\operatorname{mcc\_} = \langle \operatorname{mcg\_}, \operatorname{mca\_} \rangle) \wedge ((\operatorname{mca\_} \in \operatorname{mct\_}) \wedge (\operatorname{mchs\_} = \langle  \rangle))) \end{array}
\end{equation*}

\par\medskip\noindent\textbf{Definition} (\texttt{cr-cut}).\;For all $\operatorname{mct\_}$, $\operatorname{mchs\_}$ and $\operatorname{mcc\_}$, 
\begin{equation*}
\begin{array}{@{}l@{}} \operatorname{cr\text{-}cut}(\operatorname{mct\_}, \operatorname{mchs\_}, \operatorname{mcc\_}) \\ \Leftrightarrow\; \exists \operatorname{mcg\_}, \operatorname{mcf\_}, \operatorname{mca\_}. \\ \quad (\operatorname{mcc\_} = \langle \operatorname{mcg\_}, \operatorname{mcf\_} \rangle) \wedge (\operatorname{mchs\_} = \langle \langle \operatorname{mcg\_}, \operatorname{mca\_} \rangle, \langle \operatorname{ctx\text{-}add}(\operatorname{mcg\_}, \operatorname{mca\_}), \operatorname{mcf\_} \rangle \rangle) \end{array}
\end{equation*}

\par\medskip\noindent\textbf{Definition} (\texttt{cr-weakening}).\;For all $\operatorname{mct\_}$, $\operatorname{mchs\_}$ and $\operatorname{mcc\_}$, 
\begin{equation*}
\begin{array}{@{}l@{}} \operatorname{cr\text{-}weakening}(\operatorname{mct\_}, \operatorname{mchs\_}, \operatorname{mcc\_}) \\ \Leftrightarrow\; \exists \operatorname{mcg\_}, \operatorname{mcf\_}, \operatorname{mcd\_}. \\ \quad (\operatorname{mcc\_} = \langle \operatorname{mcg\_}, \operatorname{mcf\_} \rangle) \wedge ((\operatorname{mcd\_} \subseteq \operatorname{mcg\_}) \wedge (\operatorname{mchs\_} = \langle \langle \operatorname{mcd\_}, \operatorname{mcf\_} \rangle \rangle)) \end{array}
\end{equation*}

\par\medskip\noindent\textbf{Definition} (\texttt{cr-contradiction}).\;For all $\operatorname{mct\_}$, $\operatorname{mchs\_}$ and $\operatorname{mcc\_}$, 
\begin{equation*}
\begin{array}{@{}l@{}} \operatorname{cr\text{-}contradiction}(\operatorname{mct\_}, \operatorname{mchs\_}, \operatorname{mcc\_}) \\ \Leftrightarrow\; \exists \operatorname{mcg\_}, \operatorname{mcf\_}. \\ \quad (\operatorname{mcc\_} = \langle \operatorname{mcg\_}, \operatorname{mcf\_} \rangle) \wedge (\operatorname{mchs\_} = \langle \langle \operatorname{ctx\text{-}add}(\operatorname{mcg\_}, \operatorname{enc\text{-}not}(\operatorname{mcf\_})), \operatorname{enc\text{-}falsity} \rangle \rangle) \end{array}
\end{equation*}

\par\medskip\noindent\textbf{Definition} (\texttt{cr-truth-intro}).\;For all $\operatorname{mct\_}$, $\operatorname{mchs\_}$ and $\operatorname{mcc\_}$, 
\begin{equation*}
\operatorname{cr\text{-}truth\text{-}intro}(\operatorname{mct\_}, \operatorname{mchs\_}, \operatorname{mcc\_}) \Leftrightarrow \exists \operatorname{mcg\_}.\; ((\operatorname{mcc\_} = \langle \operatorname{mcg\_}, \operatorname{enc\text{-}truth} \rangle) \wedge (\operatorname{mchs\_} = \langle  \rangle))
\end{equation*}

\par\medskip\noindent\textbf{Definition} (\texttt{cr-eq-subst}).\;For all $\operatorname{mct\_}$, $\operatorname{mchs\_}$ and $\operatorname{mcc\_}$, 
\begin{equation*}
\begin{array}{@{}l@{}} \operatorname{cr\text{-}eq\text{-}subst}(\operatorname{mct\_}, \operatorname{mchs\_}, \operatorname{mcc\_}) \\ \Leftrightarrow\; \exists \operatorname{mcg\_}, \operatorname{mcp\_}, \operatorname{mcz\_}, \operatorname{mcs\_}, \operatorname{mctm\_}. \\ \quad \operatorname{mcc\_} = \langle \operatorname{mcg\_}, \operatorname{expr\text{-}subst}(\operatorname{mcp\_}, \operatorname{mcz\_}, \operatorname{mcs\_}) \rangle \wedge \\ \quad \operatorname{mcz\_} \in \mathbb{N} \wedge \\ \quad \operatorname{mcp\_} \in \operatorname{expr\text{-}set} \wedge \\ \quad (\operatorname{enc\text{-}eq}(\operatorname{mcs\_}, \operatorname{mctm\_}) \in \operatorname{mcg\_}) \vee (\operatorname{enc\text{-}qeq}(\operatorname{mcs\_}, \operatorname{mctm\_}) \in \operatorname{mcg\_}) \wedge \\ \quad \operatorname{mchs\_} = \langle \langle \operatorname{mcg\_}, \operatorname{expr\text{-}subst}(\operatorname{mcp\_}, \operatorname{mcz\_}, \operatorname{mctm\_}) \rangle \rangle \end{array}
\end{equation*}

\par\medskip\noindent\textbf{Definition} (\texttt{cr-eq-subst-sym}).\;For all $\operatorname{mct\_}$, $\operatorname{mchs\_}$ and $\operatorname{mcc\_}$, 
\begin{equation*}
\begin{array}{@{}l@{}} \operatorname{cr\text{-}eq\text{-}subst\text{-}sym}(\operatorname{mct\_}, \operatorname{mchs\_}, \operatorname{mcc\_}) \\ \Leftrightarrow\; \exists \operatorname{mcg\_}, \operatorname{mcp\_}, \operatorname{mcz\_}, \operatorname{mcs\_}, \operatorname{mctm\_}. \\ \quad \operatorname{mcc\_} = \langle \operatorname{mcg\_}, \operatorname{expr\text{-}subst}(\operatorname{mcp\_}, \operatorname{mcz\_}, \operatorname{mcs\_}) \rangle \wedge \\ \quad \operatorname{mcz\_} \in \mathbb{N} \wedge \\ \quad \operatorname{mcp\_} \in \operatorname{expr\text{-}set} \wedge \\ \quad (\operatorname{enc\text{-}eq}(\operatorname{mctm\_}, \operatorname{mcs\_}) \in \operatorname{mcg\_}) \vee (\operatorname{enc\text{-}qeq}(\operatorname{mctm\_}, \operatorname{mcs\_}) \in \operatorname{mcg\_}) \wedge \\ \quad \operatorname{mchs\_} = \langle \langle \operatorname{mcg\_}, \operatorname{expr\text{-}subst}(\operatorname{mcp\_}, \operatorname{mcz\_}, \operatorname{mctm\_}) \rangle \rangle \end{array}
\end{equation*}

\par\medskip\noindent\textbf{Definition} (\texttt{cr-reflexivity}).\;For all $\operatorname{mct\_}$, $\operatorname{mchs\_}$ and $\operatorname{mcc\_}$, 
\begin{equation*}
\begin{array}{@{}l@{}} \operatorname{cr\text{-}reflexivity}(\operatorname{mct\_}, \operatorname{mchs\_}, \operatorname{mcc\_}) \\ \Leftrightarrow\; \exists \operatorname{mcg\_}, \operatorname{mctm\_}, \operatorname{mcn\_}. \\ \quad \operatorname{mcc\_} = \langle \operatorname{mcg\_}, \operatorname{enc\text{-}eq}(\operatorname{mctm\_}, \operatorname{mctm\_}) \rangle \wedge \\ \quad \operatorname{mcn\_} \in \mathbb{N} \wedge \\ \quad (\operatorname{mctm\_} = \operatorname{efvar}(\operatorname{mcn\_})) \vee (\operatorname{mctm\_} = \operatorname{econst}(\operatorname{mcn\_})) \wedge \\ \quad \operatorname{mchs\_} = \langle  \rangle \end{array}
\end{equation*}

\par\medskip\noindent\textbf{Definition} (\texttt{cr-quasi-reflexivity}).\;For all $\operatorname{mct\_}$, $\operatorname{mchs\_}$ and $\operatorname{mcc\_}$, 
\begin{equation*}
\begin{array}{@{}l@{}} \operatorname{cr\text{-}quasi\text{-}reflexivity}(\operatorname{mct\_}, \operatorname{mchs\_}, \operatorname{mcc\_}) \\ \Leftrightarrow\; \exists \operatorname{mcg\_}, \operatorname{mctm\_}.\; ((\operatorname{mcc\_} = \langle \operatorname{mcg\_}, \operatorname{enc\text{-}qeq}(\operatorname{mctm\_}, \operatorname{mctm\_}) \rangle) \wedge (\operatorname{mchs\_} = \langle  \rangle)) \end{array}
\end{equation*}

\par\medskip\noindent\textbf{Definition} (\texttt{cr-union-mem}).\;For all $\operatorname{mct\_}$, $\operatorname{mchs\_}$ and $\operatorname{mcc\_}$, 
\begin{equation*}
\begin{array}{@{}l@{}} \operatorname{cr\text{-}union\text{-}mem}(\operatorname{mct\_}, \operatorname{mchs\_}, \operatorname{mcc\_}) \\ \Leftrightarrow\; \exists \operatorname{mcg\_}, \operatorname{mcx\_}, \operatorname{mcl\_}. \\ \quad \operatorname{mcc\_} = \langle \operatorname{mcg\_}, \operatorname{enc\text{-}iff}(\operatorname{enc\text{-}in}(\operatorname{mcx\_}, \operatorname{enc\text{-}union}(\operatorname{mcl\_})), \operatorname{enc\text{-}orl}(\operatorname{mcx\_}, \operatorname{mcl\_}, \operatorname{length}(\operatorname{mcl\_}))) \rangle \wedge \\ \quad \operatorname{mchs\_} = \langle  \rangle \end{array}
\end{equation*}

\par\medskip\noindent\textbf{Definition} (\texttt{cr-intersection-mem}).\;For all $\operatorname{mct\_}$, $\operatorname{mchs\_}$ and $\operatorname{mcc\_}$, 
\begin{equation*}
\begin{array}{@{}l@{}} \operatorname{cr\text{-}intersection\text{-}mem}(\operatorname{mct\_}, \operatorname{mchs\_}, \operatorname{mcc\_}) \\ \Leftrightarrow\; \exists \operatorname{mcg\_}, \operatorname{mcx\_}, \operatorname{mcl\_}. \\ \quad \operatorname{mcc\_} = \langle \operatorname{mcg\_}, \operatorname{enc\text{-}iff}(\operatorname{enc\text{-}in}(\operatorname{mcx\_}, \operatorname{enc\text{-}inter}(\operatorname{mcl\_})), \operatorname{enc\text{-}andl}(\operatorname{mcx\_}, \operatorname{mcl\_}, \operatorname{length}(\operatorname{mcl\_}))) \rangle \wedge \\ \quad 1 \leq \operatorname{length}(\operatorname{mcl\_}) \wedge \\ \quad \operatorname{mchs\_} = \langle  \rangle \end{array}
\end{equation*}

\par\medskip\noindent\textbf{Definition} (\texttt{cr-cartesian-mem}).\;For all $\operatorname{mct\_}$, $\operatorname{mchs\_}$ and $\operatorname{mcc\_}$, 
\begin{equation*}
\begin{array}{@{}l@{}} \operatorname{cr\text{-}cartesian\text{-}mem}(\operatorname{mct\_}, \operatorname{mchs\_}, \operatorname{mcc\_}) \\ \Leftrightarrow\; \exists \operatorname{mcg\_}, \operatorname{mcx\_}, \operatorname{mcl\_}. \\ \quad \operatorname{mcc\_} \\ \quad =\; \langle \operatorname{mcg\_}, \operatorname{enc\text{-}iff}(\operatorname{enc\text{-}in}(\operatorname{mcx\_}, \operatorname{enc\text{-}cart}(\operatorname{mcl\_})), \operatorname{enc\text{-}and}(\operatorname{enc\text{-}in}(\operatorname{mcx\_}, \operatorname{enc\text{-}tuples}(\operatorname{enc\text{-}setc})), \operatorname{enc\text{-}and}(\operatorname{enc\text{-}eq}(\operatorname{enc\text{-}length}(\operatorname{mcx\_}), \operatorname{enc\text{-}num}(\operatorname{length}(\operatorname{mcl\_}))), \operatorname{enc\text{-}cartl}(\operatorname{mcx\_}, \operatorname{mcl\_}, \operatorname{length}(\operatorname{mcl\_}))))) \rangle \wedge \\ \quad 1 \leq \operatorname{length}(\operatorname{mcl\_}) \wedge \\ \quad \operatorname{mchs\_} = \langle  \rangle \end{array}
\end{equation*}

\par\medskip\noindent\textbf{Definition} (\texttt{cr-tuples-mem}).\;For all $\operatorname{mct\_}$, $\operatorname{mchs\_}$ and $\operatorname{mcc\_}$, 
\begin{equation*}
\begin{array}{@{}l@{}} \operatorname{cr\text{-}tuples\text{-}mem}(\operatorname{mct\_}, \operatorname{mchs\_}, \operatorname{mcc\_}) \\ \Leftrightarrow\; \exists \operatorname{mcg\_}, \operatorname{mcl\_}, \operatorname{mca\_}. \\ \quad \operatorname{mcc\_} = \langle \operatorname{mcg\_}, \operatorname{enc\text{-}iff}(\operatorname{enc\text{-}in}(\operatorname{enc\text{-}list}(\operatorname{mcl\_}), \operatorname{enc\text{-}tuples}(\operatorname{mca\_})), \operatorname{enc\text{-}ande}(\operatorname{mcl\_}, \operatorname{mca\_}, \operatorname{length}(\operatorname{mcl\_}))) \rangle \wedge \\ \quad \operatorname{mchs\_} = \langle  \rangle \end{array}
\end{equation*}

\par\medskip\noindent\textbf{Definition} (\texttt{cr-nth-eq}).\;For all $\operatorname{mct\_}$, $\operatorname{mchs\_}$ and $\operatorname{mcc\_}$, 
\begin{equation*}
\begin{array}{@{}l@{}} \operatorname{cr\text{-}nth\text{-}eq}(\operatorname{mct\_}, \operatorname{mchs\_}, \operatorname{mcc\_}) \\ \Leftrightarrow\; \exists \operatorname{mcg\_}, \operatorname{mcl\_}, \operatorname{mck\_}. \\ \quad \operatorname{mcc\_} = \langle \operatorname{mcg\_}, \operatorname{enc\text{-}qeq}(\operatorname{enc\text{-}nth}(\operatorname{enc\text{-}num}(\operatorname{mck\_}), \operatorname{enc\text{-}list}(\operatorname{mcl\_})), \operatorname{nth}(\operatorname{mck\_}, \operatorname{mcl\_})) \rangle \wedge \\ \quad \operatorname{mck\_} \in \mathbb{N} \wedge \\ \quad 1 \leq \operatorname{mck\_} \wedge \\ \quad \operatorname{mck\_} \leq \operatorname{length}(\operatorname{mcl\_}) \wedge \\ \quad \operatorname{mchs\_} = \langle \langle \operatorname{mcg\_}, \operatorname{enc\text{-}eq}(\operatorname{enc\text{-}list}(\operatorname{mcl\_}), \operatorname{enc\text{-}list}(\operatorname{mcl\_})) \rangle \rangle \end{array}
\end{equation*}

\par\medskip\noindent\textbf{Definition} (\texttt{cr-length-eq}).\;For all $\operatorname{mct\_}$, $\operatorname{mchs\_}$ and $\operatorname{mcc\_}$, 
\begin{equation*}
\begin{array}{@{}l@{}} \operatorname{cr\text{-}length\text{-}eq}(\operatorname{mct\_}, \operatorname{mchs\_}, \operatorname{mcc\_}) \\ \Leftrightarrow\; \exists \operatorname{mcg\_}, \operatorname{mcl\_}. \\ \quad \operatorname{mcc\_} = \langle \operatorname{mcg\_}, \operatorname{enc\text{-}qeq}(\operatorname{enc\text{-}length}(\operatorname{enc\text{-}list}(\operatorname{mcl\_})), \operatorname{enc\text{-}num}(\operatorname{length}(\operatorname{mcl\_}))) \rangle \wedge \\ \quad \operatorname{mchs\_} = \langle \langle \operatorname{mcg\_}, \operatorname{enc\text{-}eq}(\operatorname{enc\text{-}list}(\operatorname{mcl\_}), \operatorname{enc\text{-}list}(\operatorname{mcl\_})) \rangle \rangle \end{array}
\end{equation*}

\par\medskip\noindent\textbf{Definition} (\texttt{cr-tuple-eq}).\;For all $\operatorname{mct\_}$, $\operatorname{mchs\_}$ and $\operatorname{mcc\_}$, 
\begin{equation*}
\begin{array}{@{}l@{}} \operatorname{cr\text{-}tuple\text{-}eq}(\operatorname{mct\_}, \operatorname{mchs\_}, \operatorname{mcc\_}) \\ \Leftrightarrow\; \exists \operatorname{mcg\_}, \operatorname{mcl\_}, \operatorname{mca\_}. \\ \quad \operatorname{mcc\_} = \langle \operatorname{mcg\_}, \operatorname{enc\text{-}iff}(\operatorname{enc\text{-}eq}(\operatorname{enc\text{-}list}(\operatorname{mcl\_}), \operatorname{enc\text{-}list}(\operatorname{mca\_})), \operatorname{enc\text{-}eql}(\operatorname{mcl\_}, \operatorname{mca\_}, \operatorname{length}(\operatorname{mcl\_}))) \rangle \wedge \\ \quad \operatorname{length}(\operatorname{mcl\_}) = \operatorname{length}(\operatorname{mca\_}) \wedge \\ \quad \operatorname{mchs\_} \\ \quad =\; \langle \langle \operatorname{mcg\_}, \operatorname{enc\text{-}eq}(\operatorname{enc\text{-}list}(\operatorname{mcl\_}), \operatorname{enc\text{-}list}(\operatorname{mcl\_})) \rangle, \langle \operatorname{mcg\_}, \operatorname{enc\text{-}eq}(\operatorname{enc\text{-}list}(\operatorname{mca\_}), \operatorname{enc\text{-}list}(\operatorname{mca\_})) \rangle \rangle \end{array}
\end{equation*}

\par\medskip\noindent\textbf{Definition} (\texttt{cr-if-true}).\;For all $\operatorname{mct\_}$, $\operatorname{mchs\_}$ and $\operatorname{mcc\_}$, 
\begin{equation*}
\begin{array}{@{}l@{}} \operatorname{cr\text{-}if\text{-}true}(\operatorname{mct\_}, \operatorname{mchs\_}, \operatorname{mcc\_}) \\ \Leftrightarrow\; \exists \operatorname{mcg\_}, \operatorname{mcf\_}, \operatorname{mcp\_}, \operatorname{mca\_}, \operatorname{mcb\_}. \\ \quad \operatorname{mcc\_} = \langle \operatorname{mcg\_}, \operatorname{mcf\_} \rangle \wedge \\ \quad \operatorname{mchs\_} = \langle \langle \operatorname{mcg\_}, \operatorname{mcp\_} \rangle, \langle \operatorname{ctx\text{-}add}(\operatorname{mcg\_}, \operatorname{enc\text{-}qeq}(\operatorname{enc\text{-}if}(\operatorname{mcp\_}, \operatorname{mca\_}, \operatorname{mcb\_}), \operatorname{mca\_})), \operatorname{mcf\_} \rangle \rangle \end{array}
\end{equation*}

\par\medskip\noindent\textbf{Definition} (\texttt{cr-if-false}).\;For all $\operatorname{mct\_}$, $\operatorname{mchs\_}$ and $\operatorname{mcc\_}$, 
\begin{equation*}
\begin{array}{@{}l@{}} \operatorname{cr\text{-}if\text{-}false}(\operatorname{mct\_}, \operatorname{mchs\_}, \operatorname{mcc\_}) \\ \Leftrightarrow\; \exists \operatorname{mcg\_}, \operatorname{mcf\_}, \operatorname{mcp\_}, \operatorname{mca\_}, \operatorname{mcb\_}. \\ \quad \operatorname{mcc\_} = \langle \operatorname{mcg\_}, \operatorname{mcf\_} \rangle \wedge \\ \quad \operatorname{mchs\_} = \langle \langle \operatorname{mcg\_}, \operatorname{enc\text{-}not}(\operatorname{mcp\_}) \rangle, \langle \operatorname{ctx\text{-}add}(\operatorname{mcg\_}, \operatorname{enc\text{-}qeq}(\operatorname{enc\text{-}if}(\operatorname{mcp\_}, \operatorname{mca\_}, \operatorname{mcb\_}), \operatorname{mcb\_})), \operatorname{mcf\_} \rangle \rangle \end{array}
\end{equation*}

\par\medskip\noindent\textbf{Definition} (\texttt{cr-sep-sethood}).\;For all $\operatorname{mct\_}$, $\operatorname{mchs\_}$ and $\operatorname{mcc\_}$, 
\begin{equation*}
\begin{array}{@{}l@{}} \operatorname{cr\text{-}sep\text{-}sethood}(\operatorname{mct\_}, \operatorname{mchs\_}, \operatorname{mcc\_}) \\ \Leftrightarrow\; \exists \operatorname{mcg\_}, \operatorname{mca\_}, \operatorname{mcp\_}. \\ \quad ((\operatorname{mcc\_} = \langle \operatorname{mcg\_}, \operatorname{enc\text{-}in}(\operatorname{enc\text{-}sep}(\operatorname{mca\_}, \operatorname{mcp\_}), \operatorname{enc\text{-}setc}) \rangle) \wedge (\operatorname{mchs\_} = \langle \langle \operatorname{mcg\_}, \operatorname{enc\text{-}in}(\operatorname{mca\_}, \operatorname{enc\text{-}setc}) \rangle \rangle)) \end{array}
\end{equation*}

\par\medskip\noindent\textbf{Definition} (\texttt{cr-sep-mem}).\;For all $\operatorname{mct\_}$, $\operatorname{mchs\_}$ and $\operatorname{mcc\_}$, 
\begin{equation*}
\begin{array}{@{}l@{}} \operatorname{cr\text{-}sep\text{-}mem}(\operatorname{mct\_}, \operatorname{mchs\_}, \operatorname{mcc\_}) \\ \Leftrightarrow\; \exists \operatorname{mcg\_}, \operatorname{mca\_}, \operatorname{mcp\_}, \operatorname{mctm\_}. \\ \quad \operatorname{mcc\_} \\ \quad =\; \langle \operatorname{mcg\_}, \operatorname{enc\text{-}iff}(\operatorname{enc\text{-}in}(\operatorname{mctm\_}, \operatorname{enc\text{-}sep}(\operatorname{mca\_}, \operatorname{mcp\_})), \operatorname{enc\text{-}and}(\operatorname{enc\text{-}in}(\operatorname{mctm\_}, \operatorname{mca\_}), \operatorname{expr\text{-}open}(\operatorname{mcp\_}, 0, \operatorname{mctm\_}))) \rangle \wedge \\ \quad \operatorname{mchs\_} = \langle  \rangle \end{array}
\end{equation*}

\par\medskip\noindent\textbf{Definition} (\texttt{cr-comp-mem}).\;For all $\operatorname{mct\_}$, $\operatorname{mchs\_}$ and $\operatorname{mcc\_}$, 
\begin{equation*}
\begin{array}{@{}l@{}} \operatorname{cr\text{-}comp\text{-}mem}(\operatorname{mct\_}, \operatorname{mchs\_}, \operatorname{mcc\_}) \\ \Leftrightarrow\; \exists \operatorname{mcg\_}, \operatorname{mcp\_}, \operatorname{mctm\_}. \\ \quad \operatorname{mcc\_} \\ \quad =\; \langle \operatorname{mcg\_}, \operatorname{enc\text{-}iff}(\operatorname{enc\text{-}in}(\operatorname{mctm\_}, \operatorname{enc\text{-}comp}(\operatorname{mcp\_})), \operatorname{enc\text{-}and}(\operatorname{enc\text{-}in}(\operatorname{mctm\_}, \operatorname{enc\text{-}setc}), \operatorname{expr\text{-}open}(\operatorname{mcp\_}, 0, \operatorname{mctm\_}))) \rangle \wedge \\ \quad \operatorname{mchs\_} = \langle  \rangle \end{array}
\end{equation*}

\par\medskip\noindent\textbf{Definition} (\texttt{cr-big-union-sethood}).\;For all $\operatorname{mct\_}$, $\operatorname{mchs\_}$ and $\operatorname{mcc\_}$, 
\begin{equation*}
\begin{array}{@{}l@{}} \operatorname{cr\text{-}big\text{-}union\text{-}sethood}(\operatorname{mct\_}, \operatorname{mchs\_}, \operatorname{mcc\_}) \\ \Leftrightarrow\; \exists \operatorname{mcg\_}, \operatorname{mca\_}, \operatorname{mcb\_}. \\ \quad \operatorname{mcc\_} = \langle \operatorname{mcg\_}, \operatorname{enc\text{-}in}(\operatorname{enc\text{-}bu}(\operatorname{mca\_}, \operatorname{mcb\_}), \operatorname{enc\text{-}setc}) \rangle \wedge \\ \quad \operatorname{mchs\_} \\ \quad =\; \langle \langle \operatorname{mcg\_}, \operatorname{enc\text{-}in}(\operatorname{mca\_}, \operatorname{enc\text{-}setc}) \rangle, \langle \operatorname{mcg\_}, \operatorname{enc\text{-}all}(\operatorname{enc\text{-}imp}(\operatorname{enc\text{-}in}(\operatorname{ebvar}(0), \operatorname{mca\_}), \operatorname{enc\text{-}in}(\operatorname{mcb\_}, \operatorname{enc\text{-}setc}))) \rangle \rangle \end{array}
\end{equation*}

\par\medskip\noindent\textbf{Definition} (\texttt{cr-big-union-mem}).\;For all $\operatorname{mct\_}$, $\operatorname{mchs\_}$ and $\operatorname{mcc\_}$, 
\begin{equation*}
\begin{array}{@{}l@{}} \operatorname{cr\text{-}big\text{-}union\text{-}mem}(\operatorname{mct\_}, \operatorname{mchs\_}, \operatorname{mcc\_}) \\ \Leftrightarrow\; \exists \operatorname{mcg\_}, \operatorname{mca\_}, \operatorname{mcb\_}, \operatorname{mcx\_}. \\ \quad \operatorname{mcc\_} \\ \quad =\; \langle \operatorname{mcg\_}, \operatorname{enc\text{-}iff}(\operatorname{enc\text{-}in}(\operatorname{mcx\_}, \operatorname{enc\text{-}bu}(\operatorname{mca\_}, \operatorname{mcb\_})), \operatorname{enc\text{-}ex}(\operatorname{enc\text{-}and}(\operatorname{enc\text{-}in}(\operatorname{ebvar}(0), \operatorname{mca\_}), \operatorname{enc\text{-}in}(\operatorname{mcx\_}, \operatorname{mcb\_})))) \rangle \wedge \\ \quad \operatorname{mchs\_} = \langle  \rangle \end{array}
\end{equation*}

\par\medskip\noindent\textbf{Definition} (\texttt{cr-iota-def}).\;For all $\operatorname{mct\_}$, $\operatorname{mchs\_}$ and $\operatorname{mcc\_}$, 
\begin{equation*}
\begin{array}{@{}l@{}} \operatorname{cr\text{-}iota\text{-}def}(\operatorname{mct\_}, \operatorname{mchs\_}, \operatorname{mcc\_}) \\ \Leftrightarrow\; \exists \operatorname{mcg\_}, \operatorname{mcf\_}, \operatorname{mcp\_}. \\ \quad \operatorname{mcc\_} = \langle \operatorname{mcg\_}, \operatorname{mcf\_} \rangle \wedge \\ \quad \operatorname{mchs\_} \\ \quad =\; \langle \langle \operatorname{mcg\_}, \operatorname{enc\text{-}ex}(\operatorname{enc\text{-}and}(\operatorname{mcp\_}, \operatorname{enc\text{-}all}(\operatorname{enc\text{-}imp}(\operatorname{mcp\_}, \operatorname{enc\text{-}eq}(\operatorname{ebvar}(1), \operatorname{ebvar}(0)))))) \rangle, \langle \operatorname{ctx\text{-}add}(\operatorname{mcg\_}, \operatorname{expr\text{-}open}(\operatorname{mcp\_}, 0, \operatorname{enc\text{-}iota}(\operatorname{mcp\_}))), \operatorname{mcf\_} \rangle \rangle \end{array}
\end{equation*}

\par\medskip\noindent\textbf{Definition} (\texttt{cr-iota-in-elim}).\;For all $\operatorname{mct\_}$, $\operatorname{mchs\_}$ and $\operatorname{mcc\_}$, 
\begin{equation*}
\begin{array}{@{}l@{}} \operatorname{cr\text{-}iota\text{-}in\text{-}elim}(\operatorname{mct\_}, \operatorname{mchs\_}, \operatorname{mcc\_}) \\ \Leftrightarrow\; \exists \operatorname{mcg\_}, \operatorname{mcf\_}, \operatorname{mcp\_}. \\ \quad \operatorname{mcc\_} = \langle \operatorname{mcg\_}, \operatorname{mcf\_} \rangle \wedge \\ \quad \operatorname{mchs\_} \\ \quad =\; \langle \langle \operatorname{mcg\_}, \operatorname{enc\text{-}eq}(\operatorname{enc\text{-}iota}(\operatorname{mcp\_}), \operatorname{enc\text{-}iota}(\operatorname{mcp\_})) \rangle, \langle \operatorname{ctx\text{-}add}(\operatorname{mcg\_}, \operatorname{expr\text{-}open}(\operatorname{mcp\_}, 0, \operatorname{enc\text{-}iota}(\operatorname{mcp\_}))), \operatorname{mcf\_} \rangle \rangle \end{array}
\end{equation*}

\par\medskip\noindent\textbf{Definition} (\texttt{cr-lambda-type}).\;For all $\operatorname{mct\_}$, $\operatorname{mchs\_}$ and $\operatorname{mcc\_}$, 
\begin{equation*}
\begin{array}{@{}l@{}} \operatorname{cr\text{-}lambda\text{-}type}(\operatorname{mct\_}, \operatorname{mchs\_}, \operatorname{mcc\_}) \\ \Leftrightarrow\; \exists \operatorname{mcg\_}, \operatorname{mca\_}, \operatorname{mcb\_}, \operatorname{mcd\_}. \\ \quad \operatorname{mcc\_} = \langle \operatorname{mcg\_}, \operatorname{enc\text{-}in}(\operatorname{enc\text{-}lam}(\operatorname{mca\_}, \operatorname{mcb\_}), \operatorname{enc\text{-}fun}(\operatorname{mca\_}, \operatorname{mcd\_})) \rangle \wedge \\ \quad \operatorname{mchs\_} \\ \quad =\; \langle \langle \operatorname{mcg\_}, \operatorname{enc\text{-}all}(\operatorname{enc\text{-}imp}(\operatorname{enc\text{-}in}(\operatorname{ebvar}(0), \operatorname{mca\_}), \operatorname{enc\text{-}in}(\operatorname{mcb\_}, \operatorname{mcd\_}))) \rangle, \langle \operatorname{mcg\_}, \operatorname{enc\text{-}in}(\operatorname{mca\_}, \operatorname{enc\text{-}setc}) \rangle \rangle \end{array}
\end{equation*}

\par\medskip\noindent\textbf{Definition} (\texttt{cr-lambda-beta-eq}).\;For all $\operatorname{mct\_}$, $\operatorname{mchs\_}$ and $\operatorname{mcc\_}$, 
\begin{equation*}
\begin{array}{@{}l@{}} \operatorname{cr\text{-}lambda\text{-}beta\text{-}eq}(\operatorname{mct\_}, \operatorname{mchs\_}, \operatorname{mcc\_}) \\ \Leftrightarrow\; \exists \operatorname{mcg\_}, \operatorname{mca\_}, \operatorname{mcb\_}, \operatorname{mctm\_}. \\ \quad \operatorname{mcc\_} = \langle \operatorname{mcg\_}, \operatorname{enc\text{-}qeq}(\operatorname{enc\text{-}apply}(\operatorname{enc\text{-}lam}(\operatorname{mca\_}, \operatorname{mcb\_}), \operatorname{mctm\_}), \operatorname{expr\text{-}open}(\operatorname{mcb\_}, 0, \operatorname{mctm\_})) \rangle \wedge \\ \quad \operatorname{mchs\_} = \langle \langle \operatorname{mcg\_}, \operatorname{enc\text{-}in}(\operatorname{mctm\_}, \operatorname{mca\_}) \rangle \rangle \end{array}
\end{equation*}

\par\medskip\noindent\textbf{Definition} (\texttt{strict-head}).\;For every $\operatorname{mchd\_}$, 
\begin{equation*}
\begin{array}{@{}l@{}} \operatorname{strict\text{-}head}(\operatorname{mchd\_}) \\ \Leftrightarrow\; (((7 \leq \operatorname{mchd\_}) \wedge (\operatorname{mchd\_} \leq 8)) \vee (((11 \leq \operatorname{mchd\_}) \wedge (\operatorname{mchd\_} \leq 19)) \vee ((30 \leq \operatorname{mchd\_}) \wedge (\operatorname{mchd\_} \leq 31)))) \end{array}
\end{equation*}

\par\medskip\noindent\textbf{Definition} (\texttt{cr-def-total}).\;For all $\operatorname{mct\_}$, $\operatorname{mchs\_}$ and $\operatorname{mcc\_}$, 
\begin{equation*}
\begin{array}{@{}l@{}} \operatorname{cr\text{-}def\text{-}total}(\operatorname{mct\_}, \operatorname{mchs\_}, \operatorname{mcc\_}) \\ \Leftrightarrow\; \exists \operatorname{mcg\_}, \operatorname{mchd\_}, \operatorname{mcl\_}. \\ \quad \operatorname{mcc\_} = \langle \operatorname{mcg\_}, \operatorname{enc\text{-}eq}(\operatorname{eapp}(\operatorname{mchd\_}, \operatorname{mcl\_}), \operatorname{eapp}(\operatorname{mchd\_}, \operatorname{mcl\_})) \rangle \wedge \\ \quad \operatorname{total\text{-}head}(\operatorname{mchd\_}) \wedge \\ \quad \operatorname{length}(\operatorname{mchs\_}) = \operatorname{length}(\operatorname{mcl\_}) \wedge \\ \quad \forall \operatorname{mci\_} \in \mathbb{N}. \\ \quad \quad (1 \leq \operatorname{mci\_}) \Rightarrow  \\ \quad \quad (\operatorname{mci\_} \leq \operatorname{length}(\operatorname{mcl\_})) \Rightarrow  \\ \quad \quad \quad \operatorname{nth}(\operatorname{mci\_}, \operatorname{mchs\_}) = \langle \operatorname{mcg\_}, \operatorname{enc\text{-}eq}(\operatorname{nth}(\operatorname{mci\_}, \operatorname{mcl\_}), \operatorname{nth}(\operatorname{mci\_}, \operatorname{mcl\_})) \rangle \end{array}
\end{equation*}

\par\medskip\noindent\textbf{Definition} (\texttt{total-head}).\;For every $\operatorname{mchd\_}$, 
\begin{equation*}
\operatorname{total\text{-}head}(\operatorname{mchd\_}) \Leftrightarrow (((11 \leq \operatorname{mchd\_}) \wedge (\operatorname{mchd\_} \leq 14)) \vee ((\operatorname{mchd\_} = 19) \vee ((30 \leq \operatorname{mchd\_}) \wedge (\operatorname{mchd\_} \leq 31))))
\end{equation*}

\par\medskip\noindent\textbf{Definition} (\texttt{cr-def-strict-in-l}).\;For all $\operatorname{mct\_}$, $\operatorname{mchs\_}$ and $\operatorname{mcc\_}$, 
\begin{equation*}
\begin{array}{@{}l@{}} \operatorname{cr\text{-}def\text{-}strict\text{-}in\text{-}l}(\operatorname{mct\_}, \operatorname{mchs\_}, \operatorname{mcc\_}) \\ \Leftrightarrow\; \exists \operatorname{mcg\_}, \operatorname{mcw\_}, \operatorname{mca\_}. \\ \quad (\operatorname{mcc\_} = \langle \operatorname{mcg\_}, \operatorname{enc\text{-}eq}(\operatorname{mcw\_}, \operatorname{mcw\_}) \rangle) \wedge (\operatorname{mchs\_} = \langle \langle \operatorname{mcg\_}, \operatorname{enc\text{-}in}(\operatorname{mcw\_}, \operatorname{mca\_}) \rangle \rangle) \end{array}
\end{equation*}

\par\medskip\noindent\textbf{Definition} (\texttt{cr-def-strict-in-r}).\;For all $\operatorname{mct\_}$, $\operatorname{mchs\_}$ and $\operatorname{mcc\_}$, 
\begin{equation*}
\begin{array}{@{}l@{}} \operatorname{cr\text{-}def\text{-}strict\text{-}in\text{-}r}(\operatorname{mct\_}, \operatorname{mchs\_}, \operatorname{mcc\_}) \\ \Leftrightarrow\; \exists \operatorname{mcg\_}, \operatorname{mcw\_}, \operatorname{mca\_}. \\ \quad (\operatorname{mcc\_} = \langle \operatorname{mcg\_}, \operatorname{enc\text{-}eq}(\operatorname{mca\_}, \operatorname{mca\_}) \rangle) \wedge (\operatorname{mchs\_} = \langle \langle \operatorname{mcg\_}, \operatorname{enc\text{-}in}(\operatorname{mcw\_}, \operatorname{mca\_}) \rangle \rangle) \end{array}
\end{equation*}

\par\medskip\noindent\textbf{Definition} (\texttt{cr-def-binder}).\;For all $\operatorname{mct\_}$, $\operatorname{mchs\_}$ and $\operatorname{mcc\_}$, 
\begin{equation*}
\begin{array}{@{}l@{}} \operatorname{cr\text{-}def\text{-}binder}(\operatorname{mct\_}, \operatorname{mchs\_}, \operatorname{mcc\_}) \\ \Leftrightarrow\; \exists \operatorname{mcg\_}, \operatorname{mchd\_}, \operatorname{mca\_}, \operatorname{mcb\_}. \\ \quad \operatorname{mcc\_} = \langle \operatorname{mcg\_}, \operatorname{enc\text{-}eq}(\operatorname{ebind}(\operatorname{mchd\_}, 1, \langle \operatorname{mca\_} \rangle, \operatorname{mcb\_}), \operatorname{ebind}(\operatorname{mchd\_}, 1, \langle \operatorname{mca\_} \rangle, \operatorname{mcb\_})) \rangle \wedge \\ \quad (\operatorname{mchd\_} = 22) \vee ((\operatorname{mchd\_} = 24) \vee (\operatorname{mchd\_} = 26)) \wedge \\ \quad \operatorname{mchs\_} = \langle \langle \operatorname{mcg\_}, \operatorname{enc\text{-}eq}(\operatorname{mca\_}, \operatorname{mca\_}) \rangle \rangle \end{array}
\end{equation*}

\par\medskip\noindent\textbf{Definition} (\texttt{cr-def-comp}).\;For all $\operatorname{mct\_}$, $\operatorname{mchs\_}$ and $\operatorname{mcc\_}$, 
\begin{equation*}
\begin{array}{@{}l@{}} \operatorname{cr\text{-}def\text{-}comp}(\operatorname{mct\_}, \operatorname{mchs\_}, \operatorname{mcc\_}) \\ \Leftrightarrow\; \exists \operatorname{mcg\_}, \operatorname{mcp\_}. \\ \quad (\operatorname{mcc\_} = \langle \operatorname{mcg\_}, \operatorname{enc\text{-}eq}(\operatorname{enc\text{-}comp}(\operatorname{mcp\_}), \operatorname{enc\text{-}comp}(\operatorname{mcp\_})) \rangle) \wedge (\operatorname{mchs\_} = \langle  \rangle) \end{array}
\end{equation*}

\par\medskip\noindent\textbf{Definition} (\texttt{cr-def-list}).\;For all $\operatorname{mct\_}$, $\operatorname{mchs\_}$ and $\operatorname{mcc\_}$, 
\begin{equation*}
\begin{array}{@{}l@{}} \operatorname{cr\text{-}def\text{-}list}(\operatorname{mct\_}, \operatorname{mchs\_}, \operatorname{mcc\_}) \\ \Leftrightarrow\; \exists \operatorname{mcg\_}, \operatorname{mcl\_}. \\ \quad \operatorname{mcc\_} = \langle \operatorname{mcg\_}, \operatorname{enc\text{-}eq}(\operatorname{enc\text{-}list}(\operatorname{mcl\_}), \operatorname{enc\text{-}list}(\operatorname{mcl\_})) \rangle \wedge \\ \quad \operatorname{length}(\operatorname{mchs\_}) = \operatorname{length}(\operatorname{mcl\_}) \wedge \\ \quad \forall \operatorname{mci\_} \in \mathbb{N}. \\ \quad \quad (1 \leq \operatorname{mci\_}) \Rightarrow  \\ \quad \quad (\operatorname{mci\_} \leq \operatorname{length}(\operatorname{mcl\_})) \Rightarrow  \\ \quad \quad \quad \operatorname{nth}(\operatorname{mci\_}, \operatorname{mchs\_}) = \langle \operatorname{mcg\_}, \operatorname{enc\text{-}in}(\operatorname{nth}(\operatorname{mci\_}, \operatorname{mcl\_}), \operatorname{enc\text{-}setc}) \rangle \end{array}
\end{equation*}

\par\medskip\noindent\textbf{Definition} (\texttt{cr-def-strict-arg}).\;For all $\operatorname{mct\_}$, $\operatorname{mchs\_}$ and $\operatorname{mcc\_}$, 
\begin{equation*}
\begin{array}{@{}l@{}} \operatorname{cr\text{-}def\text{-}strict\text{-}arg}(\operatorname{mct\_}, \operatorname{mchs\_}, \operatorname{mcc\_}) \\ \Leftrightarrow\; \exists \operatorname{mcg\_}, \operatorname{mchd\_}, \operatorname{mcl\_}, \operatorname{mci\_}. \\ \quad \operatorname{mcc\_} = \langle \operatorname{mcg\_}, \operatorname{enc\text{-}eq}(\operatorname{nth}(\operatorname{mci\_}, \operatorname{mcl\_}), \operatorname{nth}(\operatorname{mci\_}, \operatorname{mcl\_})) \rangle \wedge \\ \quad \operatorname{strict\text{-}head}(\operatorname{mchd\_}) \wedge \\ \quad \operatorname{mci\_} \in \mathbb{N} \wedge \\ \quad 1 \leq \operatorname{mci\_} \wedge \\ \quad \operatorname{mci\_} \leq \operatorname{length}(\operatorname{mcl\_}) \wedge \\ \quad \operatorname{mchs\_} = \langle \langle \operatorname{mcg\_}, \operatorname{enc\text{-}eq}(\operatorname{eapp}(\operatorname{mchd\_}, \operatorname{mcl\_}), \operatorname{eapp}(\operatorname{mchd\_}, \operatorname{mcl\_})) \rangle \rangle \end{array}
\end{equation*}

\par\medskip\noindent\textbf{Definition} (\texttt{cr-def-strict-dom}).\;For all $\operatorname{mct\_}$, $\operatorname{mchs\_}$ and $\operatorname{mcc\_}$, 
\begin{equation*}
\begin{array}{@{}l@{}} \operatorname{cr\text{-}def\text{-}strict\text{-}dom}(\operatorname{mct\_}, \operatorname{mchs\_}, \operatorname{mcc\_}) \\ \Leftrightarrow\; \exists \operatorname{mcg\_}, \operatorname{mchd\_}, \operatorname{mca\_}, \operatorname{mcb\_}. \\ \quad \operatorname{mcc\_} = \langle \operatorname{mcg\_}, \operatorname{enc\text{-}eq}(\operatorname{mca\_}, \operatorname{mca\_}) \rangle \wedge \\ \quad (\operatorname{mchd\_} = 22) \vee ((\operatorname{mchd\_} = 24) \vee (\operatorname{mchd\_} = 26)) \wedge \\ \quad \operatorname{mchs\_} = \langle \langle \operatorname{mcg\_}, \operatorname{enc\text{-}eq}(\operatorname{ebind}(\operatorname{mchd\_}, 1, \langle \operatorname{mca\_} \rangle, \operatorname{mcb\_}), \operatorname{ebind}(\operatorname{mchd\_}, 1, \langle \operatorname{mca\_} \rangle, \operatorname{mcb\_})) \rangle \rangle \end{array}
\end{equation*}

\par\medskip\noindent\textbf{Definition} (\texttt{cr-def-strict-pred}).\;For all $\operatorname{mct\_}$, $\operatorname{mchs\_}$ and $\operatorname{mcc\_}$, 
\begin{equation*}
\begin{array}{@{}l@{}} \operatorname{cr\text{-}def\text{-}strict\text{-}pred}(\operatorname{mct\_}, \operatorname{mchs\_}, \operatorname{mcc\_}) \\ \Leftrightarrow\; \exists \operatorname{mcg\_}, \operatorname{mcl\_}, \operatorname{mck\_}, \operatorname{mci\_}. \\ \quad \operatorname{mcc\_} = \langle \operatorname{mcg\_}, \operatorname{enc\text{-}eq}(\operatorname{nth}(\operatorname{mci\_}, \operatorname{mcl\_}), \operatorname{nth}(\operatorname{mci\_}, \operatorname{mcl\_})) \rangle \wedge \\ \quad \operatorname{nth}(1, \operatorname{mcl\_}) = \operatorname{econst}(\operatorname{mck\_}) \wedge \\ \quad \operatorname{pred\text{-}code}(\operatorname{mck\_}) \wedge \\ \quad \operatorname{mci\_} \in \mathbb{N} \wedge \\ \quad 2 \leq \operatorname{mci\_} \wedge \\ \quad \operatorname{mci\_} \leq \operatorname{length}(\operatorname{mcl\_}) \wedge \\ \quad \operatorname{mchs\_} = \langle \langle \operatorname{mcg\_}, \operatorname{eapp}(10, \operatorname{mcl\_}) \rangle \rangle \end{array}
\end{equation*}

\par\medskip\noindent\textbf{Definition} (\texttt{hyps-in}).\;For all $\operatorname{mchs\_}$ and $\operatorname{mcs\_}$, 
\begin{equation*}
\begin{array}{@{}l@{}} \operatorname{hyps\text{-}in}(\operatorname{mchs\_}, \operatorname{mcs\_}) \\ \Leftrightarrow\; \forall \operatorname{mci\_} \in \mathbb{N}. \\ \quad (1 \leq \operatorname{mci\_}) \Rightarrow ((\operatorname{mci\_} \leq \operatorname{length}(\operatorname{mchs\_})) \Rightarrow (\operatorname{nth}(\operatorname{mci\_}, \operatorname{mchs\_}) \in \operatorname{mcs\_})) \end{array}
\end{equation*}

\par\medskip\noindent\textbf{Definition} (\texttt{rel-detach}).\;For all $\operatorname{mct\_}$, $\operatorname{mchs\_}$ and $\operatorname{mcc\_}$, 
\begin{equation*}
\begin{array}{@{}l@{}} \operatorname{rel\text{-}detach}(\operatorname{mct\_}, \operatorname{mchs\_}, \operatorname{mcc\_}) \\ \Leftrightarrow\; \operatorname{mcc\_} \in \operatorname{seq\text{-}set} \wedge \\ \exists \operatorname{mcg\_}, \operatorname{mcf\_}, \operatorname{mcp\_}, \operatorname{mcq\_}. \\ \quad \operatorname{mcc\_} = \langle \operatorname{mcg\_}, \operatorname{mcf\_} \rangle \wedge \\ \quad \operatorname{enc\text{-}imp}(\operatorname{mcp\_}, \operatorname{mcq\_}) \in \operatorname{mcg\_} \wedge \\ \quad \operatorname{mcp\_} \in \operatorname{mcg\_} \wedge \\ \quad \operatorname{mchs\_} = \langle \langle \operatorname{ctx\text{-}add}(\operatorname{mcg\_}, \operatorname{mcq\_}), \operatorname{mcf\_} \rangle \rangle \end{array}
\end{equation*}

\par\medskip\noindent\textbf{Definition} (\texttt{rel-contraposition}).\;For all $\operatorname{mct\_}$, $\operatorname{mchs\_}$ and $\operatorname{mcc\_}$, 
\begin{equation*}
\begin{array}{@{}l@{}} \operatorname{rel\text{-}contraposition}(\operatorname{mct\_}, \operatorname{mchs\_}, \operatorname{mcc\_}) \\ \Leftrightarrow\; \operatorname{mcc\_} \in \operatorname{seq\text{-}set} \wedge \\ \exists \operatorname{mcg\_}, \operatorname{mcp\_}, \operatorname{mcq\_}. \\ \quad ((\operatorname{mcc\_} = \langle \operatorname{mcg\_}, \operatorname{enc\text{-}not}(\operatorname{mcp\_}) \rangle) \wedge ((\operatorname{enc\text{-}imp}(\operatorname{mcp\_}, \operatorname{mcq\_}) \in \operatorname{mcg\_}) \wedge (\operatorname{mchs\_} = \langle \langle \operatorname{mcg\_}, \operatorname{enc\text{-}not}(\operatorname{mcq\_}) \rangle \rangle))) \end{array}
\end{equation*}

\par\medskip\noindent\textbf{Definition} (\texttt{rel-union-intro}).\;For all $\operatorname{mct\_}$, $\operatorname{mchs\_}$ and $\operatorname{mcc\_}$, 
\begin{equation*}
\begin{array}{@{}l@{}} \operatorname{rel\text{-}union\text{-}intro}(\operatorname{mct\_}, \operatorname{mchs\_}, \operatorname{mcc\_}) \\ \Leftrightarrow\; \operatorname{mcc\_} \in \operatorname{seq\text{-}set} \wedge \\ \exists \operatorname{mcg\_}, \operatorname{mcx\_}, \operatorname{mcl\_}, \operatorname{mck\_}. \\ \quad \operatorname{mcc\_} = \langle \operatorname{mcg\_}, \operatorname{enc\text{-}in}(\operatorname{mcx\_}, \operatorname{enc\text{-}union}(\operatorname{mcl\_})) \rangle \wedge \\ \quad \operatorname{mck\_} \in \mathbb{N} \wedge \\ \quad 1 \leq \operatorname{mck\_} \wedge \\ \quad \operatorname{mck\_} \leq \operatorname{length}(\operatorname{mcl\_}) \wedge \\ \quad \operatorname{mchs\_} = \langle \langle \operatorname{mcg\_}, \operatorname{enc\text{-}in}(\operatorname{mcx\_}, \operatorname{nth}(\operatorname{mck\_}, \operatorname{mcl\_})) \rangle \rangle \end{array}
\end{equation*}

\par\medskip\noindent\textbf{Definition} (\texttt{rel-union-elim}).\;For all $\operatorname{mct\_}$, $\operatorname{mchs\_}$ and $\operatorname{mcc\_}$, 
\begin{equation*}
\begin{array}{@{}l@{}} \operatorname{rel\text{-}union\text{-}elim}(\operatorname{mct\_}, \operatorname{mchs\_}, \operatorname{mcc\_}) \\ \Leftrightarrow\; \operatorname{mcc\_} \in \operatorname{seq\text{-}set} \wedge \\ \exists \operatorname{mcg\_}, \operatorname{mcf\_}, \operatorname{mcx\_}, \operatorname{mcl\_}. \\ \quad \operatorname{mcc\_} = \langle \operatorname{mcg\_}, \operatorname{mcf\_} \rangle \wedge \\ \quad \operatorname{enc\text{-}in}(\operatorname{mcx\_}, \operatorname{enc\text{-}union}(\operatorname{mcl\_})) \in \operatorname{mcg\_} \wedge \\ \quad \operatorname{length}(\operatorname{mchs\_}) = \operatorname{length}(\operatorname{mcl\_}) \wedge \\ \quad \forall \operatorname{mci\_} \in \mathbb{N}. \\ \quad \quad (1 \leq \operatorname{mci\_}) \Rightarrow  \\ \quad \quad (\operatorname{mci\_} \leq \operatorname{length}(\operatorname{mcl\_})) \Rightarrow  \\ \quad \quad \quad \operatorname{nth}(\operatorname{mci\_}, \operatorname{mchs\_}) \\ \quad \quad \quad =\; \langle \operatorname{ctx\text{-}add}(\operatorname{difference}(\operatorname{mcg\_}, \{\operatorname{enc\text{-}in}(\operatorname{mcx\_}, \operatorname{enc\text{-}union}(\operatorname{mcl\_}))\}), \operatorname{enc\text{-}in}(\operatorname{mcx\_}, \operatorname{nth}(\operatorname{mci\_}, \operatorname{mcl\_}))), \operatorname{mcf\_} \rangle \end{array}
\end{equation*}

\par\medskip\noindent\textbf{Definition} (\texttt{rel-intersection-intro}).\;For all $\operatorname{mct\_}$, $\operatorname{mchs\_}$ and $\operatorname{mcc\_}$, 
\begin{equation*}
\begin{array}{@{}l@{}} \operatorname{rel\text{-}intersection\text{-}intro}(\operatorname{mct\_}, \operatorname{mchs\_}, \operatorname{mcc\_}) \\ \Leftrightarrow\; \operatorname{mcc\_} \in \operatorname{seq\text{-}set} \wedge \\ \exists \operatorname{mcg\_}, \operatorname{mcx\_}, \operatorname{mcl\_}. \\ \quad \operatorname{mcc\_} = \langle \operatorname{mcg\_}, \operatorname{enc\text{-}in}(\operatorname{mcx\_}, \operatorname{enc\text{-}inter}(\operatorname{mcl\_})) \rangle \wedge \\ \quad 1 \leq \operatorname{length}(\operatorname{mcl\_}) \wedge \\ \quad \operatorname{length}(\operatorname{mchs\_}) = \operatorname{length}(\operatorname{mcl\_}) \wedge \\ \quad \forall \operatorname{mci\_} \in \mathbb{N}. \\ \quad \quad (1 \leq \operatorname{mci\_}) \Rightarrow  \\ \quad \quad (\operatorname{mci\_} \leq \operatorname{length}(\operatorname{mcl\_})) \Rightarrow  \\ \quad \quad \quad \operatorname{nth}(\operatorname{mci\_}, \operatorname{mchs\_}) = \langle \operatorname{mcg\_}, \operatorname{enc\text{-}in}(\operatorname{mcx\_}, \operatorname{nth}(\operatorname{mci\_}, \operatorname{mcl\_})) \rangle \end{array}
\end{equation*}

\par\medskip\noindent\textbf{Definition} (\texttt{rel-intersection-elim}).\;For all $\operatorname{mct\_}$, $\operatorname{mchs\_}$ and $\operatorname{mcc\_}$, 
\begin{equation*}
\begin{array}{@{}l@{}} \operatorname{rel\text{-}intersection\text{-}elim}(\operatorname{mct\_}, \operatorname{mchs\_}, \operatorname{mcc\_}) \\ \Leftrightarrow\; \operatorname{mcc\_} \in \operatorname{seq\text{-}set} \wedge \\ \exists \operatorname{mcg\_}, \operatorname{mcf\_}, \operatorname{mcx\_}, \operatorname{mcl\_}, \operatorname{mck\_}. \\ \quad \operatorname{mcc\_} = \langle \operatorname{mcg\_}, \operatorname{mcf\_} \rangle \wedge \\ \quad \operatorname{enc\text{-}in}(\operatorname{mcx\_}, \operatorname{enc\text{-}inter}(\operatorname{mcl\_})) \in \operatorname{mcg\_} \wedge \\ \quad \operatorname{mck\_} \in \mathbb{N} \wedge \\ \quad 1 \leq \operatorname{mck\_} \wedge \\ \quad \operatorname{mck\_} \leq \operatorname{length}(\operatorname{mcl\_}) \wedge \\ \quad \operatorname{mchs\_} = \langle \langle \operatorname{ctx\text{-}add}(\operatorname{mcg\_}, \operatorname{enc\text{-}in}(\operatorname{mcx\_}, \operatorname{nth}(\operatorname{mck\_}, \operatorname{mcl\_}))), \operatorname{mcf\_} \rangle \rangle \end{array}
\end{equation*}

\par\medskip\noindent\textbf{Definition} (\texttt{rel-tuples-intro}).\;For all $\operatorname{mct\_}$, $\operatorname{mchs\_}$ and $\operatorname{mcc\_}$, 
\begin{equation*}
\begin{array}{@{}l@{}} \operatorname{rel\text{-}tuples\text{-}intro}(\operatorname{mct\_}, \operatorname{mchs\_}, \operatorname{mcc\_}) \\ \Leftrightarrow\; \operatorname{mcc\_} \in \operatorname{seq\text{-}set} \wedge \\ \exists \operatorname{mcg\_}, \operatorname{mcl\_}, \operatorname{mca\_}. \\ \quad \operatorname{mcc\_} = \langle \operatorname{mcg\_}, \operatorname{enc\text{-}in}(\operatorname{enc\text{-}list}(\operatorname{mcl\_}), \operatorname{enc\text{-}tuples}(\operatorname{mca\_})) \rangle \wedge \\ \quad \operatorname{length}(\operatorname{mchs\_}) = \operatorname{length}(\operatorname{mcl\_}) \wedge \\ \quad \forall \operatorname{mci\_} \in \mathbb{N}. \\ \quad \quad (1 \leq \operatorname{mci\_}) \Rightarrow  \\ \quad \quad (\operatorname{mci\_} \leq \operatorname{length}(\operatorname{mcl\_})) \Rightarrow  \\ \quad \quad \quad \operatorname{nth}(\operatorname{mci\_}, \operatorname{mchs\_}) = \langle \operatorname{mcg\_}, \operatorname{enc\text{-}in}(\operatorname{nth}(\operatorname{mci\_}, \operatorname{mcl\_}), \operatorname{mca\_}) \rangle \end{array}
\end{equation*}

\par\medskip\noindent\textbf{Definition} (\texttt{rel-tuples-elim}).\;For all $\operatorname{mct\_}$, $\operatorname{mchs\_}$ and $\operatorname{mcc\_}$, 
\begin{equation*}
\begin{array}{@{}l@{}} \operatorname{rel\text{-}tuples\text{-}elim}(\operatorname{mct\_}, \operatorname{mchs\_}, \operatorname{mcc\_}) \\ \Leftrightarrow\; \operatorname{mcc\_} \in \operatorname{seq\text{-}set} \wedge \\ \exists \operatorname{mcg\_}, \operatorname{mcf\_}, \operatorname{mcl\_}, \operatorname{mca\_}, \operatorname{mck\_}. \\ \quad \operatorname{mcc\_} = \langle \operatorname{mcg\_}, \operatorname{mcf\_} \rangle \wedge \\ \quad \operatorname{enc\text{-}in}(\operatorname{enc\text{-}list}(\operatorname{mcl\_}), \operatorname{enc\text{-}tuples}(\operatorname{mca\_})) \in \operatorname{mcg\_} \wedge \\ \quad \operatorname{mck\_} \in \mathbb{N} \wedge \\ \quad 1 \leq \operatorname{mck\_} \wedge \\ \quad \operatorname{mck\_} \leq \operatorname{length}(\operatorname{mcl\_}) \wedge \\ \quad \operatorname{mchs\_} = \langle \langle \operatorname{ctx\text{-}add}(\operatorname{mcg\_}, \operatorname{enc\text{-}in}(\operatorname{nth}(\operatorname{mck\_}, \operatorname{mcl\_}), \operatorname{mca\_})), \operatorname{mcf\_} \rangle \rangle \end{array}
\end{equation*}


\vnbbookfile{theorem-library/meta-core-laws}{theorem-library/meta-core-laws}
\noindent{\small meta-core-laws.scm -- THE CORE OF K2: the laws of well-formedness, sequents, DERIVABLE, and the admissibility of eight kernel rules (structure-library/meta-core.scm): detach, contraposition, union-intro (section H, phase 2; their proofs rewritten as dk-derive! trees by META-KIT), and union-elim, intersection-intro, ...}\par
\vnbprop{enc-not-unfold}{\vnbstmt{\forall p.\; (\operatorname{enc\text{-}not}(p) \simeq \operatorname{eapp}(2, \langle p \rangle))}}
\vnbprop{enc-and-unfold}{\vnbstmt{\forall p, q.\; (\operatorname{enc\text{-}and}(p, q) \simeq \operatorname{eapp}(3, \langle p, q \rangle))}}
\vnbprop{enc-or-unfold}{\vnbstmt{\forall p, q.\; (\operatorname{enc\text{-}or}(p, q) \simeq \operatorname{eapp}(4, \langle p, q \rangle))}}
\vnbprop{enc-imp-unfold}{\vnbstmt{\forall p, q.\; (\operatorname{enc\text{-}imp}(p, q) \simeq \operatorname{eapp}(5, \langle p, q \rangle))}}
\vnbprop{enc-iff-unfold}{\vnbstmt{\forall p, q.\; (\operatorname{enc\text{-}iff}(p, q) \simeq \operatorname{eapp}(6, \langle p, q \rangle))}}
\vnbprop{enc-in-unfold}{\vnbstmt{\forall p, q.\; (\operatorname{enc\text{-}in}(p, q) \simeq \operatorname{eapp}(7, \langle p, q \rangle))}}
\vnbprop{enc-eq-unfold}{\vnbstmt{\forall p, q.\; (\operatorname{enc\text{-}eq}(p, q) \simeq \operatorname{eapp}(8, \langle p, q \rangle))}}
\vnbprop{enc-qeq-unfold}{\vnbstmt{\forall p, q.\; (\operatorname{enc\text{-}qeq}(p, q) \simeq \operatorname{eapp}(9, \langle p, q \rangle))}}
\vnbprop{enc-union-unfold}{\vnbstmt{\forall l.\; (\operatorname{enc\text{-}union}(l) \simeq \operatorname{eapp}(11, l))}}
\vnbprop{ctx-add-unfold}{\vnbstmt{\forall g, a.\; (\operatorname{ctx\text{-}add}(g, a) \simeq (g \cup \{a\}))}}
\vnbprop{ctx-fv-unfold}{\vnbstmt{\forall g.\; (\operatorname{ctx\text{-}fv}(g) \simeq \operatorname{big\text{-}union}(w, g, \operatorname{expr\text{-}fv}(w)))}}
\vnbprop{derivable-unfold}{\vnbstmt{\begin{array}{@{}l@{}} \forall t. \\ \quad (\operatorname{derivable}(t) \simeq \{q \in \operatorname{seq\text{-}set} : \forall s \in \operatorname{power}(\operatorname{seq\text{-}set}).\; (\operatorname{core\text{-}closed}(t, s) \Rightarrow (q \in s))\}) \end{array}}}
\vnbprop{wks-stage-step}{\vnbstmt{\begin{array}{@{}l@{}} \forall k, e. \\ \quad ((k \in \mathbb{N}) \wedge (e \in \operatorname{expr\text{-}stage}(k))) \Rightarrow  \\ \quad \quad (\operatorname{wks\text{-}stage}((k + 1))(e) \simeq \operatorname{wks\text{-}stage}(k)(e)) \end{array}}}
\vnbprop{wks-stage-coh}{\vnbstmt{\begin{array}{@{}l@{}} \forall d \in \mathbb{N}, j, e. \\ \quad ((j \in \mathbb{N}) \wedge (e \in \operatorname{expr\text{-}stage}(j))) \Rightarrow  \\ \quad \quad (\operatorname{wks\text{-}stage}((j + d))(e) \simeq \operatorname{wks\text{-}stage}(j)(e)) \end{array}}}
\vnbprop{expr-wks-at}{\vnbstmt{\begin{array}{@{}l@{}} \forall j, e. \\ \quad ((j \in \mathbb{N}) \wedge (e \in \operatorname{expr\text{-}stage}(j))) \Rightarrow (\operatorname{expr\text{-}wks}(e) \simeq \operatorname{wks\text{-}stage}(j)(e)) \end{array}}}
\vnbprop{wks-efvar}{\vnbstmt{\forall n \in \mathbb{N}.\; (\operatorname{expr\text{-}wks}(\operatorname{efvar}(n)) \simeq \{1\})}}
\vnbprop{wks-ebvar}{\vnbstmt{\forall n \in \mathbb{N}.\; (\operatorname{expr\text{-}wks}(\operatorname{ebvar}(n)) \simeq \{1\})}}
\vnbprop{wks-econst}{\vnbstmt{\forall n \in \mathbb{N}.\; (\operatorname{expr\text{-}wks}(\operatorname{econst}(n)) \simeq \{1\})}}
\vnbprop{wks-eapp}{\vnbstmt{\begin{array}{@{}l@{}} \forall h, l. \\ \quad ((h \in \mathbb{N}) \wedge (l \in \operatorname{tuples}(\operatorname{expr\text{-}set}))) \Rightarrow  \\ \quad \quad (\operatorname{expr\text{-}wks}(\operatorname{eapp}(h, l)) \simeq \{k \in \mathbb{N} : ((k = \operatorname{hd\text{-}res}(h)) \wedge (\operatorname{app\text{-}ar\text{-}ok}(h, \operatorname{length}(l)) \wedge \forall w \in \operatorname{make\text{-}set}(l).\; (\operatorname{hd\text{-}argk}(h) \in \operatorname{expr\text{-}wks}(w))))\}) \end{array}}}
\vnbprop{wks-ebind}{\vnbstmt{\begin{array}{@{}l@{}} \forall h, c, a, b. \\ \quad ((h \in \mathbb{N}) \wedge ((c \in \mathbb{N}) \wedge ((a \in \operatorname{tuples}(\operatorname{expr\text{-}set})) \wedge (b \in \operatorname{expr\text{-}set})))) \Rightarrow  \\ \quad \quad (\operatorname{expr\text{-}wks}(\operatorname{ebind}(h, c, a, b)) \simeq \{k \in \mathbb{N} : ((k = \operatorname{hd\text{-}bres}(h)) \wedge (\operatorname{bind\text{-}sig}(h, c, \operatorname{length}(a)) \wedge (\forall w \in \operatorname{make\text{-}set}(a).\; (1 \in \operatorname{expr\text{-}wks}(w)) \wedge (\operatorname{hd\text{-}body}(h) \in \operatorname{expr\text{-}wks}(b)))))\}) \end{array}}}
\vnbprop{eapp-stage-inv}{\vnbstmt{\begin{array}{@{}l@{}} \forall n \in \mathbb{N}, h, l. \\ \quad (\operatorname{eapp}(h, l) \in \operatorname{expr\text{-}stage}(n)) \Rightarrow ((h \in \mathbb{N}) \wedge (l \in \operatorname{tuples}(\operatorname{expr\text{-}set}))) \end{array}}}
\vnbprop{eapp-expr-inv}{\vnbstmt{\forall h, l.\; ((\operatorname{eapp}(h, l) \in \operatorname{expr\text{-}set}) \Rightarrow ((h \in \mathbb{N}) \wedge (l \in \operatorname{tuples}(\operatorname{expr\text{-}set}))))}}
\vnbprop{ms2-intro}{\vnbstmt{\forall a, b.\; (((a \in \mathbf{Set}) \wedge (b \in \mathbf{Set})) \Rightarrow ((a \in \operatorname{make\text{-}set}(\langle a, b \rangle)) \wedge (b \in \operatorname{make\text{-}set}(\langle a, b \rangle))))}}
\vnbprop{ms1-intro}{\vnbstmt{\forall a \in \mathbf{Set}.\; (a \in \operatorname{make\text{-}set}(\langle a \rangle))}}
\vnbprop{ms2-cases}{\vnbstmt{\forall a, b.; \forall w \in \operatorname{make\text{-}set}(\langle a, b \rangle).\; ((w = a) \vee (w = b))}}
\vnbprop{ms1-cases}{\vnbstmt{\forall a.; \forall w \in \operatorname{make\text{-}set}(\langle a \rangle).\; (w = a)}}
\vnbprop{clf-iff}{\vnbstmt{\forall e.\; ((e \in \mathit{clf}) \Leftrightarrow ((e \in \operatorname{expr\text{-}set}) \wedge (\operatorname{wf\text{-}fml}(e) \wedge \operatorname{expr\text{-}lc}(e))))}}
\vnbprop{clt-iff}{\vnbstmt{\forall e.\; ((e \in \mathit{clt}) \Leftrightarrow ((e \in \operatorname{expr\text{-}set}) \wedge (\operatorname{wf\text{-}term}(e) \wedge \operatorname{expr\text{-}lc}(e))))}}
\vnbprop{wks-eapp-iff}{\vnbstmt{\begin{array}{@{}l@{}} \forall h, l, k. \\ \quad ((h \in \mathbb{N}) \wedge ((l \in \operatorname{tuples}(\operatorname{expr\text{-}set})) \wedge (k \in \mathbb{N}))) \Rightarrow  \\ \quad \quad k \in \operatorname{expr\text{-}wks}(\operatorname{eapp}(h, l)) \\ \quad \quad \Leftrightarrow\; ((k = \operatorname{hd\text{-}res}(h)) \wedge (\operatorname{app\text{-}ar\text{-}ok}(h, \operatorname{length}(l)) \wedge \forall w \in \operatorname{make\text{-}set}(l).\; (\operatorname{hd\text{-}argk}(h) \in \operatorname{expr\text{-}wks}(w)))) \end{array}}}
\vnbprop{lc-eapp-iff}{\vnbstmt{\begin{array}{@{}l@{}} \forall h, l. \\ \quad ((h \in \mathbb{N}) \wedge (l \in \operatorname{tuples}(\operatorname{expr\text{-}set}))) \Rightarrow  \\ \quad \quad \operatorname{expr\text{-}lc}(\operatorname{eapp}(h, l)) \Leftrightarrow \forall w \in \operatorname{make\text{-}set}(l).\; \operatorname{expr\text{-}lc}(w) \end{array}}}
\vnbprop{clf-conn2}{\vnbstmt{\begin{array}{@{}l@{}} \forall h, a, b. \\ \quad ((h \in \mathbb{N}) \wedge ((3 \leq h) \wedge ((h \leq 6) \wedge ((a \in \operatorname{expr\text{-}set}) \wedge (b \in \operatorname{expr\text{-}set}))))) \Rightarrow  \\ \quad \quad (\operatorname{eapp}(h, \langle a, b \rangle) \in \mathit{clf}) \Leftrightarrow ((a \in \mathit{clf}) \wedge (b \in \mathit{clf})) \end{array}}}
\vnbprop{clf-rel2}{\vnbstmt{\begin{array}{@{}l@{}} \forall h, a, b. \\ \quad ((h \in \mathbb{N}) \wedge ((7 \leq h) \wedge ((h \leq 9) \wedge ((a \in \operatorname{expr\text{-}set}) \wedge (b \in \operatorname{expr\text{-}set}))))) \Rightarrow  \\ \quad \quad (\operatorname{eapp}(h, \langle a, b \rangle) \in \mathit{clf}) \Leftrightarrow ((a \in \mathit{clt}) \wedge (b \in \mathit{clt})) \end{array}}}
\vnbprop{clf-not1}{\vnbstmt{\forall a \in \operatorname{expr\text{-}set}.\; ((\operatorname{eapp}(2, \langle a \rangle) \in \mathit{clf}) \Leftrightarrow (a \in \mathit{clf}))}}
\vnbprop{clt-union}{\vnbstmt{\begin{array}{@{}l@{}} \forall l \in \operatorname{tuples}(\operatorname{expr\text{-}set}). \\ \quad (\operatorname{eapp}(11, l) \in \mathit{clt}) \Leftrightarrow \forall u \in \operatorname{make\text{-}set}(l).\; (u \in \mathit{clt}) \end{array}}}
\vnbprop{clf-falsity0}{\vnbstmt{\operatorname{eapp}(1, \langle  \rangle) \in \mathit{clf}}}
\vnbprop{clf-truth0}{\vnbstmt{\operatorname{eapp}(0, \langle  \rangle) \in \mathit{clf}}}
\vnbprop{clf-imp-elim}{\vnbstmt{\forall p, q.\; ((\operatorname{enc\text{-}imp}(p, q) \in \mathit{clf}) \Rightarrow ((p \in \mathit{clf}) \wedge (q \in \mathit{clf})))}}
\vnbprop{clf-imp-intro}{\vnbstmt{\forall p, q.\; (((p \in \mathit{clf}) \wedge (q \in \mathit{clf})) \Rightarrow (\operatorname{enc\text{-}imp}(p, q) \in \mathit{clf}))}}
\vnbprop{clf-or-intro}{\vnbstmt{\forall p, q.\; (((p \in \mathit{clf}) \wedge (q \in \mathit{clf})) \Rightarrow (\operatorname{enc\text{-}or}(p, q) \in \mathit{clf}))}}
\vnbprop{clf-iff-intro}{\vnbstmt{\forall p, q.\; (((p \in \mathit{clf}) \wedge (q \in \mathit{clf})) \Rightarrow (\operatorname{enc\text{-}iff}(p, q) \in \mathit{clf}))}}
\vnbprop{clf-not-intro}{\vnbstmt{\forall p \in \mathit{clf}.\; (\operatorname{enc\text{-}not}(p) \in \mathit{clf})}}
\vnbprop{clf-in-elim}{\vnbstmt{\forall x, y.\; ((\operatorname{enc\text{-}in}(x, y) \in \mathit{clf}) \Rightarrow ((x \in \mathit{clt}) \wedge (y \in \mathit{clt})))}}
\vnbprop{clf-in-intro}{\vnbstmt{\forall x, y.\; (((x \in \mathit{clt}) \wedge (y \in \mathit{clt})) \Rightarrow (\operatorname{enc\text{-}in}(x, y) \in \mathit{clf}))}}
\vnbprop{clt-union-elim}{\vnbstmt{\begin{array}{@{}l@{}} \forall l. \\ \quad (\operatorname{enc\text{-}union}(l) \in \mathit{clt}) \Rightarrow  \\ \quad \quad (l \in \operatorname{tuples}(\operatorname{expr\text{-}set})) \wedge \forall u \in \operatorname{make\text{-}set}(l).\; (u \in \mathit{clt}) \end{array}}}
\vnbprop{clf-falsity}{\vnbstmt{\operatorname{enc\text{-}falsity} \in \mathit{clf}}}
\vnbprop{clf-and-intro}{\vnbstmt{\forall p, q.\; (((p \in \mathit{clf}) \wedge (q \in \mathit{clf})) \Rightarrow (\operatorname{enc\text{-}and}(p, q) \in \mathit{clf}))}}
\vnbprop{clf-and-elim}{\vnbstmt{\forall p, q.\; ((\operatorname{enc\text{-}and}(p, q) \in \mathit{clf}) \Rightarrow ((p \in \mathit{clf}) \wedge (q \in \mathit{clf})))}}
\vnbprop{clf-or-elim}{\vnbstmt{\forall p, q.\; ((\operatorname{enc\text{-}or}(p, q) \in \mathit{clf}) \Rightarrow ((p \in \mathit{clf}) \wedge (q \in \mathit{clf})))}}
\vnbprop{clf-iff-elim}{\vnbstmt{\forall p, q.\; ((\operatorname{enc\text{-}iff}(p, q) \in \mathit{clf}) \Rightarrow ((p \in \mathit{clf}) \wedge (q \in \mathit{clf})))}}
\vnbprop{clf-not-elim}{\vnbstmt{\forall p.\; ((\operatorname{enc\text{-}not}(p) \in \mathit{clf}) \Rightarrow (p \in \mathit{clf}))}}
\vnbprop{clf-truth}{\vnbstmt{\operatorname{enc\text{-}truth} \in \mathit{clf}}}
\vnbprop{seq-iff}{\vnbstmt{\forall g, f.\; ((\langle g, f \rangle \in \operatorname{seq\text{-}set}) \Leftrightarrow ((g \in \operatorname{ctx\text{-}set}) \wedge (f \in \mathit{clf})))}}
\vnbprop{ctx-iff}{\vnbstmt{\forall g.\; ((g \in \operatorname{ctx\text{-}set}) \Leftrightarrow ((g \in \operatorname{power}(\mathit{clf})) \wedge (|g| \in \mathbb{N})))}}
\vnbprop{ctx-mem}{\vnbstmt{\forall g, a.\; (((g \in \operatorname{ctx\text{-}set}) \wedge (a \in g)) \Rightarrow (a \in \mathit{clf}))}}
\vnbprop{ctx-add-has}{\vnbstmt{\forall g.; \forall a \in \mathbf{Set}.\; (a \in \operatorname{ctx\text{-}add}(g, a))}}
\vnbprop{ctx-add-keeps}{\vnbstmt{\forall g, a, b.\; (((a \in \mathbf{Set}) \wedge (b \in g)) \Rightarrow (b \in \operatorname{ctx\text{-}add}(g, a)))}}
\vnbprop{ctx-add-ctx}{\vnbstmt{\forall g, a.\; (((g \in \operatorname{ctx\text{-}set}) \wedge (a \in \mathit{clf})) \Rightarrow (\operatorname{ctx\text{-}add}(g, a) \in \operatorname{ctx\text{-}set}))}}
\vnbprop{ctx-add-sub}{\vnbstmt{\forall g.; \forall a \in \mathbf{Set}.\; (g \subseteq \operatorname{ctx\text{-}add}(g, a))}}
\vnbprop{ctx-sub-cons}{\vnbstmt{\forall d, g.; \forall a \in \mathbf{Set}.\; (((d \subseteq g) \wedge (a \in g)) \Rightarrow (\operatorname{ctx\text{-}add}(d, a) \subseteq g))}}
\vnbprop{ctx-sub-snoc}{\vnbstmt{\forall d, g.; \forall b \in \mathbf{Set}.\; ((d \subseteq g) \Rightarrow (d \subseteq \operatorname{ctx\text{-}add}(g, b)))}}
\vnbprop{der-in-seq}{\vnbstmt{\forall t.; \forall s \in \operatorname{derivable}(t).\; (s \in \operatorname{seq\text{-}set})}}
\vnbprop{derivable-induction}{\vnbstmt{\begin{array}{@{}l@{}} \forall t, \mathit{mczs}, s. \\ \quad ((\mathit{mczs} \in \operatorname{power}(\operatorname{seq\text{-}set})) \wedge (\operatorname{core\text{-}closed}(t, \mathit{mczs}) \wedge (s \in \operatorname{derivable}(t)))) \Rightarrow  \\ \quad \quad s \in \mathit{mczs} \end{array}}}
\vnbprop{core-closed-apply}{\vnbstmt{\begin{array}{@{}l@{}} \forall t, s, a, c. \\ \quad (\operatorname{core\text{-}closed}(t, s) \wedge (\operatorname{core\text{-}rule}(t, a, c) \wedge (\operatorname{hyps\text{-}in}(a, s) \wedge (c \in \operatorname{seq\text{-}set})))) \Rightarrow  \\ \quad \quad c \in s \end{array}}}
\vnbprop{derivable-closed}{\vnbstmt{\begin{array}{@{}l@{}} \forall t, s, c. \\ \quad (\operatorname{core\text{-}rule}(t, s, c) \wedge (\operatorname{hyps\text{-}in}(s, \operatorname{derivable}(t)) \wedge (c \in \operatorname{seq\text{-}set}))) \Rightarrow  \\ \quad \quad c \in \operatorname{derivable}(t) \end{array}}}
\vnbprop{hyps-in-0}{\vnbstmt{\forall s.\; \operatorname{hyps\text{-}in}(\langle  \rangle, s)}}
\vnbprop{hyps-in-1}{\vnbstmt{\forall a, s.\; ((a \in s) \Rightarrow \operatorname{hyps\text{-}in}(\langle a \rangle, s))}}
\vnbprop{hyps-in-2}{\vnbstmt{\forall a, b, s.\; (((a \in s) \wedge (b \in s)) \Rightarrow \operatorname{hyps\text{-}in}(\langle a, b \rangle, s))}}
\vnbprop{hyps-in-1-elim}{\vnbstmt{\forall a, s.\; (((\langle a \rangle = \langle a \rangle) \wedge \operatorname{hyps\text{-}in}(\langle a \rangle, s)) \Rightarrow (a \in s))}}
\vnbprop{derivable-and-intro}{\vnbstmt{\begin{array}{@{}l@{}} \forall t, s, c. \\ \quad (\operatorname{cr\text{-}and\text{-}intro}(t, s, c) \wedge (\operatorname{hyps\text{-}in}(s, \operatorname{derivable}(t)) \wedge (c \in \operatorname{seq\text{-}set}))) \Rightarrow  \\ \quad \quad c \in \operatorname{derivable}(t) \end{array}}}
\vnbprop{derivable-and-elim}{\vnbstmt{\begin{array}{@{}l@{}} \forall t, s, c. \\ \quad (\operatorname{cr\text{-}and\text{-}elim}(t, s, c) \wedge (\operatorname{hyps\text{-}in}(s, \operatorname{derivable}(t)) \wedge (c \in \operatorname{seq\text{-}set}))) \Rightarrow  \\ \quad \quad c \in \operatorname{derivable}(t) \end{array}}}
\vnbprop{derivable-or-intro-left}{\vnbstmt{\begin{array}{@{}l@{}} \forall t, s, c. \\ \quad (\operatorname{cr\text{-}or\text{-}intro\text{-}left}(t, s, c) \wedge (\operatorname{hyps\text{-}in}(s, \operatorname{derivable}(t)) \wedge (c \in \operatorname{seq\text{-}set}))) \Rightarrow  \\ \quad \quad c \in \operatorname{derivable}(t) \end{array}}}
\vnbprop{derivable-or-intro-right}{\vnbstmt{\begin{array}{@{}l@{}} \forall t, s, c. \\ \quad (\operatorname{cr\text{-}or\text{-}intro\text{-}right}(t, s, c) \wedge (\operatorname{hyps\text{-}in}(s, \operatorname{derivable}(t)) \wedge (c \in \operatorname{seq\text{-}set}))) \Rightarrow  \\ \quad \quad c \in \operatorname{derivable}(t) \end{array}}}
\vnbprop{derivable-or-elim}{\vnbstmt{\begin{array}{@{}l@{}} \forall t, s, c. \\ \quad (\operatorname{cr\text{-}or\text{-}elim}(t, s, c) \wedge (\operatorname{hyps\text{-}in}(s, \operatorname{derivable}(t)) \wedge (c \in \operatorname{seq\text{-}set}))) \Rightarrow  \\ \quad \quad c \in \operatorname{derivable}(t) \end{array}}}
\vnbprop{derivable-implies-intro}{\vnbstmt{\begin{array}{@{}l@{}} \forall t, s, c. \\ \quad (\operatorname{cr\text{-}implies\text{-}intro}(t, s, c) \wedge (\operatorname{hyps\text{-}in}(s, \operatorname{derivable}(t)) \wedge (c \in \operatorname{seq\text{-}set}))) \Rightarrow  \\ \quad \quad c \in \operatorname{derivable}(t) \end{array}}}
\vnbprop{derivable-backchain}{\vnbstmt{\begin{array}{@{}l@{}} \forall t, s, c. \\ \quad (\operatorname{cr\text{-}backchain}(t, s, c) \wedge (\operatorname{hyps\text{-}in}(s, \operatorname{derivable}(t)) \wedge (c \in \operatorname{seq\text{-}set}))) \Rightarrow  \\ \quad \quad c \in \operatorname{derivable}(t) \end{array}}}
\vnbprop{derivable-not-intro}{\vnbstmt{\begin{array}{@{}l@{}} \forall t, s, c. \\ \quad (\operatorname{cr\text{-}not\text{-}intro}(t, s, c) \wedge (\operatorname{hyps\text{-}in}(s, \operatorname{derivable}(t)) \wedge (c \in \operatorname{seq\text{-}set}))) \Rightarrow  \\ \quad \quad c \in \operatorname{derivable}(t) \end{array}}}
\vnbprop{derivable-not-elim}{\vnbstmt{\begin{array}{@{}l@{}} \forall t, s, c. \\ \quad (\operatorname{cr\text{-}not\text{-}elim}(t, s, c) \wedge (\operatorname{hyps\text{-}in}(s, \operatorname{derivable}(t)) \wedge (c \in \operatorname{seq\text{-}set}))) \Rightarrow  \\ \quad \quad c \in \operatorname{derivable}(t) \end{array}}}
\vnbprop{derivable-iff-intro}{\vnbstmt{\begin{array}{@{}l@{}} \forall t, s, c. \\ \quad (\operatorname{cr\text{-}iff\text{-}intro}(t, s, c) \wedge (\operatorname{hyps\text{-}in}(s, \operatorname{derivable}(t)) \wedge (c \in \operatorname{seq\text{-}set}))) \Rightarrow  \\ \quad \quad c \in \operatorname{derivable}(t) \end{array}}}
\vnbprop{derivable-iff-elim}{\vnbstmt{\begin{array}{@{}l@{}} \forall t, s, c. \\ \quad (\operatorname{cr\text{-}iff\text{-}elim}(t, s, c) \wedge (\operatorname{hyps\text{-}in}(s, \operatorname{derivable}(t)) \wedge (c \in \operatorname{seq\text{-}set}))) \Rightarrow  \\ \quad \quad c \in \operatorname{derivable}(t) \end{array}}}
\vnbprop{derivable-forall-intro}{\vnbstmt{\begin{array}{@{}l@{}} \forall t, s, c. \\ \quad (\operatorname{cr\text{-}forall\text{-}intro}(t, s, c) \wedge (\operatorname{hyps\text{-}in}(s, \operatorname{derivable}(t)) \wedge (c \in \operatorname{seq\text{-}set}))) \Rightarrow  \\ \quad \quad c \in \operatorname{derivable}(t) \end{array}}}
\vnbprop{derivable-forall-elim}{\vnbstmt{\begin{array}{@{}l@{}} \forall t, s, c. \\ \quad (\operatorname{cr\text{-}forall\text{-}elim}(t, s, c) \wedge (\operatorname{hyps\text{-}in}(s, \operatorname{derivable}(t)) \wedge (c \in \operatorname{seq\text{-}set}))) \Rightarrow  \\ \quad \quad c \in \operatorname{derivable}(t) \end{array}}}
\vnbprop{derivable-forsome-intro}{\vnbstmt{\begin{array}{@{}l@{}} \forall t, s, c. \\ \quad (\operatorname{cr\text{-}forsome\text{-}intro}(t, s, c) \wedge (\operatorname{hyps\text{-}in}(s, \operatorname{derivable}(t)) \wedge (c \in \operatorname{seq\text{-}set}))) \Rightarrow  \\ \quad \quad c \in \operatorname{derivable}(t) \end{array}}}
\vnbprop{derivable-forsome-elim}{\vnbstmt{\begin{array}{@{}l@{}} \forall t, s, c. \\ \quad (\operatorname{cr\text{-}forsome\text{-}elim}(t, s, c) \wedge (\operatorname{hyps\text{-}in}(s, \operatorname{derivable}(t)) \wedge (c \in \operatorname{seq\text{-}set}))) \Rightarrow  \\ \quad \quad c \in \operatorname{derivable}(t) \end{array}}}
\vnbprop{derivable-assumption}{\vnbstmt{\begin{array}{@{}l@{}} \forall t, s, c. \\ \quad (\operatorname{cr\text{-}assumption}(t, s, c) \wedge (\operatorname{hyps\text{-}in}(s, \operatorname{derivable}(t)) \wedge (c \in \operatorname{seq\text{-}set}))) \Rightarrow  \\ \quad \quad c \in \operatorname{derivable}(t) \end{array}}}
\vnbprop{derivable-theorem-assumption}{\vnbstmt{\begin{array}{@{}l@{}} \forall t, s, c. \\ \quad (\operatorname{cr\text{-}theorem\text{-}assumption}(t, s, c) \wedge (\operatorname{hyps\text{-}in}(s, \operatorname{derivable}(t)) \wedge (c \in \operatorname{seq\text{-}set}))) \Rightarrow  \\ \quad \quad c \in \operatorname{derivable}(t) \end{array}}}
\vnbprop{derivable-cut}{\vnbstmt{\begin{array}{@{}l@{}} \forall t, s, c. \\ \quad (\operatorname{cr\text{-}cut}(t, s, c) \wedge (\operatorname{hyps\text{-}in}(s, \operatorname{derivable}(t)) \wedge (c \in \operatorname{seq\text{-}set}))) \Rightarrow  \\ \quad \quad c \in \operatorname{derivable}(t) \end{array}}}
\vnbprop{derivable-weakening}{\vnbstmt{\begin{array}{@{}l@{}} \forall t, s, c. \\ \quad (\operatorname{cr\text{-}weakening}(t, s, c) \wedge (\operatorname{hyps\text{-}in}(s, \operatorname{derivable}(t)) \wedge (c \in \operatorname{seq\text{-}set}))) \Rightarrow  \\ \quad \quad c \in \operatorname{derivable}(t) \end{array}}}
\vnbprop{derivable-contradiction}{\vnbstmt{\begin{array}{@{}l@{}} \forall t, s, c. \\ \quad (\operatorname{cr\text{-}contradiction}(t, s, c) \wedge (\operatorname{hyps\text{-}in}(s, \operatorname{derivable}(t)) \wedge (c \in \operatorname{seq\text{-}set}))) \Rightarrow  \\ \quad \quad c \in \operatorname{derivable}(t) \end{array}}}
\vnbprop{derivable-truth-intro}{\vnbstmt{\begin{array}{@{}l@{}} \forall t, s, c. \\ \quad (\operatorname{cr\text{-}truth\text{-}intro}(t, s, c) \wedge (\operatorname{hyps\text{-}in}(s, \operatorname{derivable}(t)) \wedge (c \in \operatorname{seq\text{-}set}))) \Rightarrow  \\ \quad \quad c \in \operatorname{derivable}(t) \end{array}}}
\vnbprop{derivable-eq-subst}{\vnbstmt{\begin{array}{@{}l@{}} \forall t, s, c. \\ \quad (\operatorname{cr\text{-}eq\text{-}subst}(t, s, c) \wedge (\operatorname{hyps\text{-}in}(s, \operatorname{derivable}(t)) \wedge (c \in \operatorname{seq\text{-}set}))) \Rightarrow  \\ \quad \quad c \in \operatorname{derivable}(t) \end{array}}}
\vnbprop{derivable-reflexivity}{\vnbstmt{\begin{array}{@{}l@{}} \forall t, s, c. \\ \quad (\operatorname{cr\text{-}reflexivity}(t, s, c) \wedge (\operatorname{hyps\text{-}in}(s, \operatorname{derivable}(t)) \wedge (c \in \operatorname{seq\text{-}set}))) \Rightarrow  \\ \quad \quad c \in \operatorname{derivable}(t) \end{array}}}
\vnbprop{derivable-quasi-reflexivity}{\vnbstmt{\begin{array}{@{}l@{}} \forall t, s, c. \\ \quad (\operatorname{cr\text{-}quasi\text{-}reflexivity}(t, s, c) \wedge (\operatorname{hyps\text{-}in}(s, \operatorname{derivable}(t)) \wedge (c \in \operatorname{seq\text{-}set}))) \Rightarrow  \\ \quad \quad c \in \operatorname{derivable}(t) \end{array}}}
\vnbprop{derivable-union-mem}{\vnbstmt{\begin{array}{@{}l@{}} \forall t, s, c. \\ \quad (\operatorname{cr\text{-}union\text{-}mem}(t, s, c) \wedge (\operatorname{hyps\text{-}in}(s, \operatorname{derivable}(t)) \wedge (c \in \operatorname{seq\text{-}set}))) \Rightarrow  \\ \quad \quad c \in \operatorname{derivable}(t) \end{array}}}
\vnbprop{derivable-intersection-mem}{\vnbstmt{\begin{array}{@{}l@{}} \forall t, s, c. \\ \quad (\operatorname{cr\text{-}intersection\text{-}mem}(t, s, c) \wedge (\operatorname{hyps\text{-}in}(s, \operatorname{derivable}(t)) \wedge (c \in \operatorname{seq\text{-}set}))) \Rightarrow  \\ \quad \quad c \in \operatorname{derivable}(t) \end{array}}}
\vnbprop{derivable-cartesian-mem}{\vnbstmt{\begin{array}{@{}l@{}} \forall t, s, c. \\ \quad (\operatorname{cr\text{-}cartesian\text{-}mem}(t, s, c) \wedge (\operatorname{hyps\text{-}in}(s, \operatorname{derivable}(t)) \wedge (c \in \operatorname{seq\text{-}set}))) \Rightarrow  \\ \quad \quad c \in \operatorname{derivable}(t) \end{array}}}
\vnbprop{derivable-tuples-mem}{\vnbstmt{\begin{array}{@{}l@{}} \forall t, s, c. \\ \quad (\operatorname{cr\text{-}tuples\text{-}mem}(t, s, c) \wedge (\operatorname{hyps\text{-}in}(s, \operatorname{derivable}(t)) \wedge (c \in \operatorname{seq\text{-}set}))) \Rightarrow  \\ \quad \quad c \in \operatorname{derivable}(t) \end{array}}}
\vnbprop{derivable-nth-eq}{\vnbstmt{\begin{array}{@{}l@{}} \forall t, s, c. \\ \quad (\operatorname{cr\text{-}nth\text{-}eq}(t, s, c) \wedge (\operatorname{hyps\text{-}in}(s, \operatorname{derivable}(t)) \wedge (c \in \operatorname{seq\text{-}set}))) \Rightarrow  \\ \quad \quad c \in \operatorname{derivable}(t) \end{array}}}
\vnbprop{derivable-length-eq}{\vnbstmt{\begin{array}{@{}l@{}} \forall t, s, c. \\ \quad (\operatorname{cr\text{-}length\text{-}eq}(t, s, c) \wedge (\operatorname{hyps\text{-}in}(s, \operatorname{derivable}(t)) \wedge (c \in \operatorname{seq\text{-}set}))) \Rightarrow  \\ \quad \quad c \in \operatorname{derivable}(t) \end{array}}}
\vnbprop{derivable-tuple-eq}{\vnbstmt{\begin{array}{@{}l@{}} \forall t, s, c. \\ \quad (\operatorname{cr\text{-}tuple\text{-}eq}(t, s, c) \wedge (\operatorname{hyps\text{-}in}(s, \operatorname{derivable}(t)) \wedge (c \in \operatorname{seq\text{-}set}))) \Rightarrow  \\ \quad \quad c \in \operatorname{derivable}(t) \end{array}}}
\vnbprop{derivable-if-true}{\vnbstmt{\begin{array}{@{}l@{}} \forall t, s, c. \\ \quad (\operatorname{cr\text{-}if\text{-}true}(t, s, c) \wedge (\operatorname{hyps\text{-}in}(s, \operatorname{derivable}(t)) \wedge (c \in \operatorname{seq\text{-}set}))) \Rightarrow  \\ \quad \quad c \in \operatorname{derivable}(t) \end{array}}}
\vnbprop{derivable-if-false}{\vnbstmt{\begin{array}{@{}l@{}} \forall t, s, c. \\ \quad (\operatorname{cr\text{-}if\text{-}false}(t, s, c) \wedge (\operatorname{hyps\text{-}in}(s, \operatorname{derivable}(t)) \wedge (c \in \operatorname{seq\text{-}set}))) \Rightarrow  \\ \quad \quad c \in \operatorname{derivable}(t) \end{array}}}
\vnbprop{derivable-sep-sethood}{\vnbstmt{\begin{array}{@{}l@{}} \forall t, s, c. \\ \quad (\operatorname{cr\text{-}sep\text{-}sethood}(t, s, c) \wedge (\operatorname{hyps\text{-}in}(s, \operatorname{derivable}(t)) \wedge (c \in \operatorname{seq\text{-}set}))) \Rightarrow  \\ \quad \quad c \in \operatorname{derivable}(t) \end{array}}}
\vnbprop{derivable-sep-mem}{\vnbstmt{\begin{array}{@{}l@{}} \forall t, s, c. \\ \quad (\operatorname{cr\text{-}sep\text{-}mem}(t, s, c) \wedge (\operatorname{hyps\text{-}in}(s, \operatorname{derivable}(t)) \wedge (c \in \operatorname{seq\text{-}set}))) \Rightarrow  \\ \quad \quad c \in \operatorname{derivable}(t) \end{array}}}
\vnbprop{derivable-comp-mem}{\vnbstmt{\begin{array}{@{}l@{}} \forall t, s, c. \\ \quad (\operatorname{cr\text{-}comp\text{-}mem}(t, s, c) \wedge (\operatorname{hyps\text{-}in}(s, \operatorname{derivable}(t)) \wedge (c \in \operatorname{seq\text{-}set}))) \Rightarrow  \\ \quad \quad c \in \operatorname{derivable}(t) \end{array}}}
\vnbprop{derivable-big-union-sethood}{\vnbstmt{\begin{array}{@{}l@{}} \forall t, s, c. \\ \quad (\operatorname{cr\text{-}big\text{-}union\text{-}sethood}(t, s, c) \wedge (\operatorname{hyps\text{-}in}(s, \operatorname{derivable}(t)) \wedge (c \in \operatorname{seq\text{-}set}))) \Rightarrow  \\ \quad \quad c \in \operatorname{derivable}(t) \end{array}}}
\vnbprop{derivable-big-union-mem}{\vnbstmt{\begin{array}{@{}l@{}} \forall t, s, c. \\ \quad (\operatorname{cr\text{-}big\text{-}union\text{-}mem}(t, s, c) \wedge (\operatorname{hyps\text{-}in}(s, \operatorname{derivable}(t)) \wedge (c \in \operatorname{seq\text{-}set}))) \Rightarrow  \\ \quad \quad c \in \operatorname{derivable}(t) \end{array}}}
\vnbprop{derivable-iota-def}{\vnbstmt{\begin{array}{@{}l@{}} \forall t, s, c. \\ \quad (\operatorname{cr\text{-}iota\text{-}def}(t, s, c) \wedge (\operatorname{hyps\text{-}in}(s, \operatorname{derivable}(t)) \wedge (c \in \operatorname{seq\text{-}set}))) \Rightarrow  \\ \quad \quad c \in \operatorname{derivable}(t) \end{array}}}
\vnbprop{derivable-iota-in-elim}{\vnbstmt{\begin{array}{@{}l@{}} \forall t, s, c. \\ \quad (\operatorname{cr\text{-}iota\text{-}in\text{-}elim}(t, s, c) \wedge (\operatorname{hyps\text{-}in}(s, \operatorname{derivable}(t)) \wedge (c \in \operatorname{seq\text{-}set}))) \Rightarrow  \\ \quad \quad c \in \operatorname{derivable}(t) \end{array}}}
\vnbprop{derivable-lambda-type}{\vnbstmt{\begin{array}{@{}l@{}} \forall t, s, c. \\ \quad (\operatorname{cr\text{-}lambda\text{-}type}(t, s, c) \wedge (\operatorname{hyps\text{-}in}(s, \operatorname{derivable}(t)) \wedge (c \in \operatorname{seq\text{-}set}))) \Rightarrow  \\ \quad \quad c \in \operatorname{derivable}(t) \end{array}}}
\vnbprop{derivable-lambda-beta-eq}{\vnbstmt{\begin{array}{@{}l@{}} \forall t, s, c. \\ \quad (\operatorname{cr\text{-}lambda\text{-}beta\text{-}eq}(t, s, c) \wedge (\operatorname{hyps\text{-}in}(s, \operatorname{derivable}(t)) \wedge (c \in \operatorname{seq\text{-}set}))) \Rightarrow  \\ \quad \quad c \in \operatorname{derivable}(t) \end{array}}}
\vnbprop{derivable-eq-subst-sym}{\vnbstmt{\begin{array}{@{}l@{}} \forall t, s, c. \\ \quad (\operatorname{cr\text{-}eq\text{-}subst\text{-}sym}(t, s, c) \wedge (\operatorname{hyps\text{-}in}(s, \operatorname{derivable}(t)) \wedge (c \in \operatorname{seq\text{-}set}))) \Rightarrow  \\ \quad \quad c \in \operatorname{derivable}(t) \end{array}}}
\vnbprop{derivable-def-strict-in-l}{\vnbstmt{\begin{array}{@{}l@{}} \forall t, s, c. \\ \quad (\operatorname{cr\text{-}def\text{-}strict\text{-}in\text{-}l}(t, s, c) \wedge (\operatorname{hyps\text{-}in}(s, \operatorname{derivable}(t)) \wedge (c \in \operatorname{seq\text{-}set}))) \Rightarrow  \\ \quad \quad c \in \operatorname{derivable}(t) \end{array}}}
\vnbprop{derivable-def-strict-in-r}{\vnbstmt{\begin{array}{@{}l@{}} \forall t, s, c. \\ \quad (\operatorname{cr\text{-}def\text{-}strict\text{-}in\text{-}r}(t, s, c) \wedge (\operatorname{hyps\text{-}in}(s, \operatorname{derivable}(t)) \wedge (c \in \operatorname{seq\text{-}set}))) \Rightarrow  \\ \quad \quad c \in \operatorname{derivable}(t) \end{array}}}
\vnbprop{derivable-def-total}{\vnbstmt{\begin{array}{@{}l@{}} \forall t, s, c. \\ \quad (\operatorname{cr\text{-}def\text{-}total}(t, s, c) \wedge (\operatorname{hyps\text{-}in}(s, \operatorname{derivable}(t)) \wedge (c \in \operatorname{seq\text{-}set}))) \Rightarrow  \\ \quad \quad c \in \operatorname{derivable}(t) \end{array}}}
\vnbprop{derivable-def-binder}{\vnbstmt{\begin{array}{@{}l@{}} \forall t, s, c. \\ \quad (\operatorname{cr\text{-}def\text{-}binder}(t, s, c) \wedge (\operatorname{hyps\text{-}in}(s, \operatorname{derivable}(t)) \wedge (c \in \operatorname{seq\text{-}set}))) \Rightarrow  \\ \quad \quad c \in \operatorname{derivable}(t) \end{array}}}
\vnbprop{derivable-def-comp}{\vnbstmt{\begin{array}{@{}l@{}} \forall t, s, c. \\ \quad (\operatorname{cr\text{-}def\text{-}comp}(t, s, c) \wedge (\operatorname{hyps\text{-}in}(s, \operatorname{derivable}(t)) \wedge (c \in \operatorname{seq\text{-}set}))) \Rightarrow  \\ \quad \quad c \in \operatorname{derivable}(t) \end{array}}}
\vnbprop{derivable-def-list}{\vnbstmt{\begin{array}{@{}l@{}} \forall t, s, c. \\ \quad (\operatorname{cr\text{-}def\text{-}list}(t, s, c) \wedge (\operatorname{hyps\text{-}in}(s, \operatorname{derivable}(t)) \wedge (c \in \operatorname{seq\text{-}set}))) \Rightarrow  \\ \quad \quad c \in \operatorname{derivable}(t) \end{array}}}
\vnbprop{derivable-def-strict-arg}{\vnbstmt{\begin{array}{@{}l@{}} \forall t, s, c. \\ \quad (\operatorname{cr\text{-}def\text{-}strict\text{-}arg}(t, s, c) \wedge (\operatorname{hyps\text{-}in}(s, \operatorname{derivable}(t)) \wedge (c \in \operatorname{seq\text{-}set}))) \Rightarrow  \\ \quad \quad c \in \operatorname{derivable}(t) \end{array}}}
\vnbprop{derivable-def-strict-dom}{\vnbstmt{\begin{array}{@{}l@{}} \forall t, s, c. \\ \quad (\operatorname{cr\text{-}def\text{-}strict\text{-}dom}(t, s, c) \wedge (\operatorname{hyps\text{-}in}(s, \operatorname{derivable}(t)) \wedge (c \in \operatorname{seq\text{-}set}))) \Rightarrow  \\ \quad \quad c \in \operatorname{derivable}(t) \end{array}}}
\vnbprop{derivable-def-strict-pred}{\vnbstmt{\begin{array}{@{}l@{}} \forall t, s, c. \\ \quad (\operatorname{cr\text{-}def\text{-}strict\text{-}pred}(t, s, c) \wedge (\operatorname{hyps\text{-}in}(s, \operatorname{derivable}(t)) \wedge (c \in \operatorname{seq\text{-}set}))) \Rightarrow  \\ \quad \quad c \in \operatorname{derivable}(t) \end{array}}}
\vnbprop{derivable-assumption-c}{\vnbstmt{\forall t, g, a.\; (((\langle g, a \rangle \in \operatorname{seq\text{-}set}) \wedge (a \in g)) \Rightarrow (\langle g, a \rangle \in \operatorname{derivable}(t)))}}
\vnbprop{derivable-backchain-c}{\vnbstmt{\begin{array}{@{}l@{}} \forall t, g, f, a. \\ \quad ((\langle g, f \rangle \in \operatorname{seq\text{-}set}) \wedge ((\operatorname{enc\text{-}imp}(a, f) \in g) \wedge (\langle g, a \rangle \in \operatorname{derivable}(t)))) \Rightarrow  \\ \quad \quad \langle g, f \rangle \in \operatorname{derivable}(t) \end{array}}}
\vnbprop{derivable-cut-c}{\vnbstmt{\begin{array}{@{}l@{}} \forall t, g, f, a. \\ \quad ((\langle g, f \rangle \in \operatorname{seq\text{-}set}) \wedge ((\langle g, a \rangle \in \operatorname{derivable}(t)) \wedge (\langle \operatorname{ctx\text{-}add}(g, a), f \rangle \in \operatorname{derivable}(t)))) \Rightarrow  \\ \quad \quad \langle g, f \rangle \in \operatorname{derivable}(t) \end{array}}}
\vnbprop{derivable-weakening-c}{\vnbstmt{\begin{array}{@{}l@{}} \forall t, g, f, d. \\ \quad ((\langle g, f \rangle \in \operatorname{seq\text{-}set}) \wedge ((d \subseteq g) \wedge (\langle d, f \rangle \in \operatorname{derivable}(t)))) \Rightarrow  \\ \quad \quad \langle g, f \rangle \in \operatorname{derivable}(t) \end{array}}}
\vnbprop{derivable-not-intro-c}{\vnbstmt{\begin{array}{@{}l@{}} \forall t, g, a. \\ \quad ((\langle g, \operatorname{enc\text{-}not}(a) \rangle \in \operatorname{seq\text{-}set}) \wedge (\langle \operatorname{ctx\text{-}add}(g, a), \operatorname{enc\text{-}falsity} \rangle \in \operatorname{derivable}(t))) \Rightarrow  \\ \quad \quad \langle g, \operatorname{enc\text{-}not}(a) \rangle \in \operatorname{derivable}(t) \end{array}}}
\vnbprop{derivable-not-elim-c}{\vnbstmt{\begin{array}{@{}l@{}} \forall t, g, f, a. \\ \quad ((\langle g, f \rangle \in \operatorname{seq\text{-}set}) \wedge ((\operatorname{enc\text{-}not}(a) \in g) \wedge (a \in g))) \Rightarrow (\langle g, f \rangle \in \operatorname{derivable}(t)) \end{array}}}
\vnbprop{derivable-or-intro-left-c}{\vnbstmt{\begin{array}{@{}l@{}} \forall t, g, a, b. \\ \quad ((\langle g, \operatorname{enc\text{-}or}(a, b) \rangle \in \operatorname{seq\text{-}set}) \wedge (\langle g, a \rangle \in \operatorname{derivable}(t))) \Rightarrow  \\ \quad \quad \langle g, \operatorname{enc\text{-}or}(a, b) \rangle \in \operatorname{derivable}(t) \end{array}}}
\vnbprop{derivable-or-intro-right-c}{\vnbstmt{\begin{array}{@{}l@{}} \forall t, g, a, b. \\ \quad ((\langle g, \operatorname{enc\text{-}or}(a, b) \rangle \in \operatorname{seq\text{-}set}) \wedge (\langle g, b \rangle \in \operatorname{derivable}(t))) \Rightarrow  \\ \quad \quad \langle g, \operatorname{enc\text{-}or}(a, b) \rangle \in \operatorname{derivable}(t) \end{array}}}
\vnbprop{derivable-iff-elim-c}{\vnbstmt{\begin{array}{@{}l@{}} \forall t, g, f, a, b. \\ \quad \langle g, f \rangle \in \operatorname{seq\text{-}set} \wedge \\ \quad \operatorname{enc\text{-}iff}(a, b) \in g \wedge \\ \quad \langle \operatorname{ctx\text{-}add}(\operatorname{ctx\text{-}add}(g, \operatorname{enc\text{-}imp}(a, b)), \operatorname{enc\text{-}imp}(b, a)), f \rangle \in \operatorname{derivable}(t) \Rightarrow \\ \quad \quad \langle g, f \rangle \in \operatorname{derivable}(t) \end{array}}}
\vnbprop{derivable-union-mem-c}{\vnbstmt{\begin{array}{@{}l@{}} \forall t, g, x, l. \\ \quad (\langle g, \operatorname{enc\text{-}iff}(\operatorname{enc\text{-}in}(x, \operatorname{enc\text{-}union}(l)), \operatorname{enc\text{-}orl}(x, l, \operatorname{length}(l))) \rangle \in \operatorname{seq\text{-}set}) \Rightarrow  \\ \quad \quad \langle g, \operatorname{enc\text{-}iff}(\operatorname{enc\text{-}in}(x, \operatorname{enc\text{-}union}(l)), \operatorname{enc\text{-}orl}(x, l, \operatorname{length}(l))) \rangle \in \operatorname{derivable}(t) \end{array}}}
\vnbprop{adm-detach}{\vnbstmt{\begin{array}{@{}l@{}} \forall t, s, c. \\ \quad (\operatorname{rel\text{-}detach}(t, s, c) \wedge \operatorname{hyps\text{-}in}(s, \operatorname{derivable}(t))) \Rightarrow (c \in \operatorname{derivable}(t)) \end{array}}}
\vnbprop{adm-contraposition}{\vnbstmt{\begin{array}{@{}l@{}} \forall t, s, c. \\ \quad (\operatorname{rel\text{-}contraposition}(t, s, c) \wedge \operatorname{hyps\text{-}in}(s, \operatorname{derivable}(t))) \Rightarrow  \\ \quad \quad c \in \operatorname{derivable}(t) \end{array}}}
\vnbprop{orl-clf}{\vnbstmt{\begin{array}{@{}l@{}} \forall n \in \mathbb{N}, x, l. \\ \quad ((x \in \mathit{clt}) \wedge ((l \in \operatorname{tuples}(\operatorname{expr\text{-}set})) \wedge (\forall u \in \operatorname{make\text{-}set}(l).\; (u \in \mathit{clt}) \wedge (n \leq \operatorname{length}(l))))) \Rightarrow  \\ \quad \quad \operatorname{enc\text{-}orl}(x, l, n) \in \mathit{clf} \end{array}}}
\vnbprop{orl-der}{\vnbstmt{\begin{array}{@{}l@{}} \forall n \in \mathbb{N}, t, g, x, l, k. \\ \quad \langle g, \operatorname{enc\text{-}in}(x, \operatorname{nth}(k, l)) \rangle \in \operatorname{derivable}(t) \wedge \\ \quad g \in \operatorname{ctx\text{-}set} \wedge \\ \quad x \in \mathit{clt} \wedge \\ \quad l \in \operatorname{tuples}(\operatorname{expr\text{-}set}) \wedge \\ \quad \forall u \in \operatorname{make\text{-}set}(l).\; (u \in \mathit{clt}) \wedge \\ \quad k \in \mathbb{N} \wedge \\ \quad 1 \leq k \wedge \\ \quad k \leq n \wedge \\ \quad n \leq \operatorname{length}(l) \Rightarrow \\ \quad \quad \langle g, \operatorname{enc\text{-}orl}(x, l, n) \rangle \in \operatorname{derivable}(t) \end{array}}}
\vnbprop{adm-union-intro}{\vnbstmt{\begin{array}{@{}l@{}} \forall t, s, c. \\ \quad (\operatorname{rel\text{-}union\text{-}intro}(t, s, c) \wedge \operatorname{hyps\text{-}in}(s, \operatorname{derivable}(t))) \Rightarrow (c \in \operatorname{derivable}(t)) \end{array}}}
\vnbprop{clf-is-set}{\vnbstmt{\mathit{clf} \in \mathbf{Set}}}
\vnbprop{seq-set-is-set}{\vnbstmt{\operatorname{seq\text{-}set} \in \mathbf{Set}}}
\vnbprop{ctx-fv-finite}{\vnbstmt{\begin{array}{@{}l@{}} \forall g \in \operatorname{ctx\text{-}set}. \\ \quad (\operatorname{ctx\text{-}fv}(g) \in \mathbf{Set}) \wedge ((\operatorname{ctx\text{-}fv}(g) \subseteq \mathbb{N}) \wedge (|\operatorname{ctx\text{-}fv}(g)| \in \mathbb{N})) \end{array}}}
\vnbprop{fresh-name-exists}{\vnbstmt{\begin{array}{@{}l@{}} \forall g, f. \\ \quad ((g \in \operatorname{ctx\text{-}set}) \wedge (f \in \operatorname{expr\text{-}set})) \Rightarrow  \\ \quad \quad \exists n \in \mathbb{N}.\; \neg (n \in (\operatorname{ctx\text{-}fv}(g) \cup \operatorname{expr\text{-}fv}(f))) \end{array}}}
\vnbprop{fv-app2}{\vnbstmt{\begin{array}{@{}l@{}} \forall h, a, b. \\ \quad ((h \in \mathbb{N}) \wedge ((a \in \operatorname{expr\text{-}set}) \wedge (b \in \operatorname{expr\text{-}set}))) \Rightarrow  \\ \quad \quad \operatorname{expr\text{-}fv}(\operatorname{eapp}(h, \langle a, b \rangle)) = (\operatorname{expr\text{-}fv}(a) \cup \operatorname{expr\text{-}fv}(b)) \end{array}}}
\vnbprop{fv-app1}{\vnbstmt{\begin{array}{@{}l@{}} \forall h, a. \\ \quad ((h \in \mathbb{N}) \wedge (a \in \operatorname{expr\text{-}set})) \Rightarrow (\operatorname{expr\text{-}fv}(\operatorname{eapp}(h, \langle a \rangle)) = \operatorname{expr\text{-}fv}(a)) \end{array}}}
\vnbprop{fv-app0}{\vnbstmt{\forall h \in \mathbb{N}.\; (\operatorname{expr\text{-}fv}(\operatorname{eapp}(h, \langle  \rangle)) = \emptyset)}}
\vnbprop{fv-bind0}{\vnbstmt{\begin{array}{@{}l@{}} \forall h, a, b. \\ \quad ((h \in \mathbb{N}) \wedge ((a \in \mathbb{N}) \wedge (b \in \operatorname{expr\text{-}set}))) \Rightarrow  \\ \quad \quad \operatorname{expr\text{-}fv}(\operatorname{ebind}(h, a, \langle  \rangle, b)) = \operatorname{expr\text{-}fv}(b) \end{array}}}
\vnbprop{wks-eapp-transfer}{\vnbstmt{\begin{array}{@{}l@{}} \forall h, l, a. \\ \quad h \in \mathbb{N} \wedge \\ \quad l \in \operatorname{tuples}(\operatorname{expr\text{-}set}) \wedge \\ \quad a \in \operatorname{tuples}(\operatorname{expr\text{-}set}) \wedge \\ \quad \operatorname{length}(a) = \operatorname{length}(l) \wedge \\ \quad \forall i \in \mathbb{N}. \\ \quad \quad ((1 \leq i) \wedge (i \leq \operatorname{length}(l))) \Rightarrow (\operatorname{expr\text{-}wks}(\operatorname{nth}(i, a)) = \operatorname{expr\text{-}wks}(\operatorname{nth}(i, l))) \Rightarrow \\ \quad \quad \operatorname{expr\text{-}wks}(\operatorname{eapp}(h, a)) = \operatorname{expr\text{-}wks}(\operatorname{eapp}(h, l)) \end{array}}}
\vnbprop{wks-ebind-transfer}{\vnbstmt{\begin{array}{@{}l@{}} \forall h, c, a, d, b, f. \\ \quad h \in \mathbb{N} \wedge \\ \quad c \in \mathbb{N} \wedge \\ \quad a \in \operatorname{tuples}(\operatorname{expr\text{-}set}) \wedge \\ \quad d \in \operatorname{tuples}(\operatorname{expr\text{-}set}) \wedge \\ \quad b \in \operatorname{expr\text{-}set} \wedge \\ \quad f \in \operatorname{expr\text{-}set} \wedge \\ \quad \operatorname{length}(d) = \operatorname{length}(a) \wedge \\ \quad \forall i \in \mathbb{N}. \\ \quad \quad ((1 \leq i) \wedge (i \leq \operatorname{length}(a))) \Rightarrow (\operatorname{expr\text{-}wks}(\operatorname{nth}(i, d)) = \operatorname{expr\text{-}wks}(\operatorname{nth}(i, a))) \wedge \\ \quad \operatorname{expr\text{-}wks}(f) = \operatorname{expr\text{-}wks}(b) \Rightarrow \\ \quad \quad \operatorname{expr\text{-}wks}(\operatorname{ebind}(h, c, d, f)) = \operatorname{expr\text{-}wks}(\operatorname{ebind}(h, c, a, b)) \end{array}}}
\vnbprop{wks-subst}{\vnbstmt{\begin{array}{@{}l@{}} \forall m.; \forall u, e \in \operatorname{expr\text{-}set}. \\ \quad (\operatorname{expr\text{-}wks}(u) = \{1\}) \Rightarrow (\operatorname{expr\text{-}wks}(\operatorname{expr\text{-}subst}(e, m, u)) = \operatorname{expr\text{-}wks}(e)) \end{array}}}
\vnbprop{wks-open}{\vnbstmt{\begin{array}{@{}l@{}} \forall u, e \in \operatorname{expr\text{-}set}, x \in \mathbb{N}. \\ \quad (\operatorname{expr\text{-}wks}(u) = \{1\}) \Rightarrow (\operatorname{expr\text{-}wks}(\operatorname{expr\text{-}open}(e, x, u)) = \operatorname{expr\text{-}wks}(e)) \end{array}}}
\vnbprop{enc-inter-unfold}{\vnbstmt{\forall l.\; (\operatorname{enc\text{-}inter}(l) \simeq \operatorname{eapp}(12, l))}}
\vnbprop{enc-tuples-unfold}{\vnbstmt{\forall p.\; (\operatorname{enc\text{-}tuples}(p) \simeq \operatorname{eapp}(14, \langle p \rangle))}}
\vnbprop{enc-list-unfold}{\vnbstmt{\forall l.\; (\operatorname{enc\text{-}list}(l) \simeq \operatorname{eapp}(17, l))}}
\vnbprop{clt-inter}{\vnbstmt{\begin{array}{@{}l@{}} \forall l \in \operatorname{tuples}(\operatorname{expr\text{-}set}). \\ \quad (\operatorname{eapp}(12, l) \in \mathit{clt}) \Leftrightarrow \forall u \in \operatorname{make\text{-}set}(l).\; (u \in \mathit{clt}) \end{array}}}
\vnbprop{clt-list}{\vnbstmt{\begin{array}{@{}l@{}} \forall l \in \operatorname{tuples}(\operatorname{expr\text{-}set}). \\ \quad (\operatorname{eapp}(17, l) \in \mathit{clt}) \Leftrightarrow \forall u \in \operatorname{make\text{-}set}(l).\; (u \in \mathit{clt}) \end{array}}}
\vnbprop{clt-tuples1}{\vnbstmt{\forall a \in \operatorname{expr\text{-}set}.\; ((\operatorname{eapp}(14, \langle a \rangle) \in \mathit{clt}) \Leftrightarrow (a \in \mathit{clt}))}}
\vnbprop{clt-inter-elim}{\vnbstmt{\begin{array}{@{}l@{}} \forall l. \\ \quad (\operatorname{enc\text{-}inter}(l) \in \mathit{clt}) \Rightarrow  \\ \quad \quad (l \in \operatorname{tuples}(\operatorname{expr\text{-}set})) \wedge \forall u \in \operatorname{make\text{-}set}(l).\; (u \in \mathit{clt}) \end{array}}}
\vnbprop{clt-list-elim}{\vnbstmt{\begin{array}{@{}l@{}} \forall l. \\ \quad (\operatorname{enc\text{-}list}(l) \in \mathit{clt}) \Rightarrow  \\ \quad \quad (l \in \operatorname{tuples}(\operatorname{expr\text{-}set})) \wedge \forall u \in \operatorname{make\text{-}set}(l).\; (u \in \mathit{clt}) \end{array}}}
\vnbprop{clt-tuples-elim}{\vnbstmt{\forall p.\; ((\operatorname{enc\text{-}tuples}(p) \in \mathit{clt}) \Rightarrow (p \in \mathit{clt}))}}
\vnbprop{ctx-diff-sub}{\vnbstmt{\forall g, b.\; (\operatorname{difference}(g, b) \subseteq g)}}
\vnbprop{enc-andl-clf}{\vnbstmt{\begin{array}{@{}l@{}} \forall n \in \mathbb{N}, x, l. \\ \quad ((x \in \mathit{clt}) \wedge ((l \in \operatorname{tuples}(\operatorname{expr\text{-}set})) \wedge (\forall u \in \operatorname{make\text{-}set}(l).\; (u \in \mathit{clt}) \wedge (n \leq \operatorname{length}(l))))) \Rightarrow  \\ \quad \quad \operatorname{enc\text{-}andl}(x, l, n) \in \mathit{clf} \end{array}}}
\vnbprop{enc-ande-clf}{\vnbstmt{\begin{array}{@{}l@{}} \forall n \in \mathbb{N}, l.; \forall a \in \mathit{clt}. \\ \quad ((l \in \operatorname{tuples}(\operatorname{expr\text{-}set})) \wedge (\forall u \in \operatorname{make\text{-}set}(l).\; (u \in \mathit{clt}) \wedge (n \leq \operatorname{length}(l)))) \Rightarrow  \\ \quad \quad \operatorname{enc\text{-}ande}(l, a, n) \in \mathit{clf} \end{array}}}
\vnbprop{enc-andl-der}{\vnbstmt{\begin{array}{@{}l@{}} \forall n \in \mathbb{N}, t, g, x, l. \\ \quad g \in \operatorname{ctx\text{-}set} \wedge \\ \quad x \in \mathit{clt} \wedge \\ \quad l \in \operatorname{tuples}(\operatorname{expr\text{-}set}) \wedge \\ \quad \forall u \in \operatorname{make\text{-}set}(l).\; (u \in \mathit{clt}) \wedge \\ \quad \forall k \in \mathbb{N}. \\ \quad \quad ((1 \leq k) \wedge (k \leq n)) \Rightarrow (\langle g, \operatorname{enc\text{-}in}(x, \operatorname{nth}(k, l)) \rangle \in \operatorname{derivable}(t)) \wedge \\ \quad n \leq \operatorname{length}(l) \Rightarrow \\ \quad \quad \langle g, \operatorname{enc\text{-}andl}(x, l, n) \rangle \in \operatorname{derivable}(t) \end{array}}}
\vnbprop{enc-ande-der}{\vnbstmt{\begin{array}{@{}l@{}} \forall n \in \mathbb{N}, t, g, l, a. \\ \quad g \in \operatorname{ctx\text{-}set} \wedge \\ \quad a \in \mathit{clt} \wedge \\ \quad l \in \operatorname{tuples}(\operatorname{expr\text{-}set}) \wedge \\ \quad \forall u \in \operatorname{make\text{-}set}(l).\; (u \in \mathit{clt}) \wedge \\ \quad \forall k \in \mathbb{N}. \\ \quad \quad ((1 \leq k) \wedge (k \leq n)) \Rightarrow (\langle g, \operatorname{enc\text{-}in}(\operatorname{nth}(k, l), a) \rangle \in \operatorname{derivable}(t)) \wedge \\ \quad n \leq \operatorname{length}(l) \Rightarrow \\ \quad \quad \langle g, \operatorname{enc\text{-}ande}(l, a, n) \rangle \in \operatorname{derivable}(t) \end{array}}}
\vnbprop{enc-andl-elim-der}{\vnbstmt{\begin{array}{@{}l@{}} \forall n \in \mathbb{N}, t, g, x, l, k. \\ \quad \langle g, \operatorname{enc\text{-}andl}(x, l, n) \rangle \in \operatorname{derivable}(t) \wedge \\ \quad g \in \operatorname{ctx\text{-}set} \wedge \\ \quad x \in \mathit{clt} \wedge \\ \quad l \in \operatorname{tuples}(\operatorname{expr\text{-}set}) \wedge \\ \quad \forall u \in \operatorname{make\text{-}set}(l).\; (u \in \mathit{clt}) \wedge \\ \quad k \in \mathbb{N} \wedge \\ \quad 1 \leq k \wedge \\ \quad k \leq n \wedge \\ \quad n \leq \operatorname{length}(l) \Rightarrow \\ \quad \quad \langle g, \operatorname{enc\text{-}in}(x, \operatorname{nth}(k, l)) \rangle \in \operatorname{derivable}(t) \end{array}}}
\vnbprop{enc-ande-elim-der}{\vnbstmt{\begin{array}{@{}l@{}} \forall n \in \mathbb{N}, t, g, l, a, k. \\ \quad \langle g, \operatorname{enc\text{-}ande}(l, a, n) \rangle \in \operatorname{derivable}(t) \wedge \\ \quad g \in \operatorname{ctx\text{-}set} \wedge \\ \quad a \in \mathit{clt} \wedge \\ \quad l \in \operatorname{tuples}(\operatorname{expr\text{-}set}) \wedge \\ \quad \forall u \in \operatorname{make\text{-}set}(l).\; (u \in \mathit{clt}) \wedge \\ \quad k \in \mathbb{N} \wedge \\ \quad 1 \leq k \wedge \\ \quad k \leq n \wedge \\ \quad n \leq \operatorname{length}(l) \Rightarrow \\ \quad \quad \langle g, \operatorname{enc\text{-}in}(\operatorname{nth}(k, l), a) \rangle \in \operatorname{derivable}(t) \end{array}}}
\vnbprop{orl-elim-der}{\vnbstmt{\begin{array}{@{}l@{}} \forall n \in \mathbb{N}, t, g, f, x, l. \\ \quad g \in \operatorname{ctx\text{-}set} \wedge \\ \quad f \in \mathit{clf} \wedge \\ \quad x \in \mathit{clt} \wedge \\ \quad l \in \operatorname{tuples}(\operatorname{expr\text{-}set}) \wedge \\ \quad \forall u \in \operatorname{make\text{-}set}(l).\; (u \in \mathit{clt}) \wedge \\ \quad \forall k \in \mathbb{N}. \\ \quad \quad ((1 \leq k) \wedge (k \leq n)) \Rightarrow (\langle \operatorname{ctx\text{-}add}(g, \operatorname{enc\text{-}in}(x, \operatorname{nth}(k, l))), f \rangle \in \operatorname{derivable}(t)) \wedge \\ \quad n \leq \operatorname{length}(l) \Rightarrow \\ \quad \quad \langle \operatorname{ctx\text{-}add}(g, \operatorname{enc\text{-}orl}(x, l, n)), f \rangle \in \operatorname{derivable}(t) \end{array}}}
\vnbprop{adm-union-elim}{\vnbstmt{\begin{array}{@{}l@{}} \forall t, s, c. \\ \quad (\operatorname{rel\text{-}union\text{-}elim}(t, s, c) \wedge \operatorname{hyps\text{-}in}(s, \operatorname{derivable}(t))) \Rightarrow (c \in \operatorname{derivable}(t)) \end{array}}}
\vnbprop{adm-intersection-intro}{\vnbstmt{\begin{array}{@{}l@{}} \forall t, s, c. \\ \quad (\operatorname{rel\text{-}intersection\text{-}intro}(t, s, c) \wedge \operatorname{hyps\text{-}in}(s, \operatorname{derivable}(t))) \Rightarrow  \\ \quad \quad c \in \operatorname{derivable}(t) \end{array}}}
\vnbprop{adm-intersection-elim}{\vnbstmt{\begin{array}{@{}l@{}} \forall t, s, c. \\ \quad (\operatorname{rel\text{-}intersection\text{-}elim}(t, s, c) \wedge \operatorname{hyps\text{-}in}(s, \operatorname{derivable}(t))) \Rightarrow  \\ \quad \quad c \in \operatorname{derivable}(t) \end{array}}}
\vnbprop{adm-tuples-intro}{\vnbstmt{\begin{array}{@{}l@{}} \forall t, s, c. \\ \quad (\operatorname{rel\text{-}tuples\text{-}intro}(t, s, c) \wedge \operatorname{hyps\text{-}in}(s, \operatorname{derivable}(t))) \Rightarrow (c \in \operatorname{derivable}(t)) \end{array}}}
\vnbprop{adm-tuples-elim}{\vnbstmt{\begin{array}{@{}l@{}} \forall t, s, c. \\ \quad (\operatorname{rel\text{-}tuples\text{-}elim}(t, s, c) \wedge \operatorname{hyps\text{-}in}(s, \operatorname{derivable}(t))) \Rightarrow (c \in \operatorname{derivable}(t)) \end{array}}}

\vnbbookfile{structure-library/meta-cr}{structure-library/meta-cr}
\noindent{\small meta-cr.scm -- CONTEXTUAL REWRITING (CR) over the encoded core: the one-step rewrite relation with its local contexts, the equivalence it is to be proved inside, and the kernel relations of the seven rewriting rules, transcribed. (docs/meta-k2-admissibility-2026-10-08.md, sections 4.5 and 6, phase 3.)}\par
\par\medskip\noindent\textbf{Definition} (\texttt{rewr}).\;$\operatorname{rewr}(\operatorname{mrt\_}, \operatorname{mrk\_}) :=$
\begin{equation*}
\{\operatorname{mrx\_} \in \operatorname{rw\text{-}univ} : \forall \operatorname{mrs\_} \in \operatorname{power}(\operatorname{rw\text{-}univ}).\; (\operatorname{rw\text{-}closed}(\operatorname{mrt\_}, \operatorname{mrk\_}, \operatorname{mrs\_}) \Rightarrow (\operatorname{mrx\_} \in \operatorname{mrs\_}))\}
\end{equation*}

\par\medskip\noindent\textbf{Definition} (\texttt{eqv}).\;$\operatorname{eqv}(\operatorname{mrt\_}) :=$ $\{\operatorname{mrx\_} \in \operatorname{rw\text{-}univ} : \operatorname{eqv\text{-}ok}(\operatorname{mrt\_}, \operatorname{mrx\_})\}$
\par\medskip\noindent\textbf{Definition} (\texttt{rw-univ}).\; $(\operatorname{rw\text{-}univ} \simeq (\mathbb{N} \times \operatorname{ctx\text{-}set}))$
\par\medskip\noindent\textbf{Definition} (\texttt{eqv-ok}).\;For all $\operatorname{mrt\_}$ and $\operatorname{mrx\_}$, 
\begin{equation*}
\begin{array}{@{}l@{}} \operatorname{eqv\text{-}ok}(\operatorname{mrt\_}, \operatorname{mrx\_}) \\ \Leftrightarrow\; \exists \operatorname{mrg\_}, \operatorname{mre\_}, \operatorname{mrf\_}. \\ \quad \operatorname{mrx\_} = \langle 2, \operatorname{mrg\_}, \operatorname{mre\_}, \operatorname{mrf\_} \rangle \wedge \\ \quad \operatorname{mre\_} \in \mathit{clf} \wedge \\ \quad \operatorname{mrf\_} \in \mathit{clf} \wedge \\ \quad \langle \operatorname{mrg\_}, \operatorname{enc\text{-}iff}(\operatorname{mre\_}, \operatorname{mrf\_}) \rangle \in \operatorname{derivable}(\operatorname{mrt\_}) \vee \\ \quad \operatorname{mrx\_} = \langle 1, \operatorname{mrg\_}, \operatorname{mre\_}, \operatorname{mrf\_} \rangle \wedge \\ \quad \operatorname{mre\_} \in \mathit{clt} \wedge \\ \quad \operatorname{mrf\_} \in \mathit{clt} \wedge \\ \quad \langle \operatorname{mrg\_}, \operatorname{enc\text{-}qeq}(\operatorname{mre\_}, \operatorname{mrf\_}) \rangle \in \operatorname{derivable}(\operatorname{mrt\_}) \end{array}
\end{equation*}

\par\medskip\noindent\textbf{Definition} (\texttt{rw-rule}).\;For all $\operatorname{mrt\_}$, $\operatorname{mrk\_}$, $\operatorname{mrs\_}$ and $\operatorname{mrx\_}$, 
\begin{equation*}
\begin{array}{@{}l@{}} \operatorname{rw\text{-}rule}(\operatorname{mrt\_}, \operatorname{mrk\_}, \operatorname{mrs\_}, \operatorname{mrx\_}) \\ \Leftrightarrow\; (1 \leq \operatorname{mrk\_}) \wedge \operatorname{rw\text{-}top}(\operatorname{mrt\_}, \operatorname{mrs\_}, \operatorname{mrx\_}) \vee \\ (1 \leq \operatorname{mrk\_}) \wedge \operatorname{rw\text{-}not}(\operatorname{mrt\_}, \operatorname{mrs\_}, \operatorname{mrx\_}) \vee \\ (1 \leq \operatorname{mrk\_}) \wedge \operatorname{rw\text{-}and\text{-}l}(\operatorname{mrt\_}, \operatorname{mrs\_}, \operatorname{mrx\_}) \vee \\ (1 \leq \operatorname{mrk\_}) \wedge \operatorname{rw\text{-}and\text{-}r}(\operatorname{mrt\_}, \operatorname{mrs\_}, \operatorname{mrx\_}) \vee \\ (1 \leq \operatorname{mrk\_}) \wedge \operatorname{rw\text{-}or\text{-}l}(\operatorname{mrt\_}, \operatorname{mrs\_}, \operatorname{mrx\_}) \vee \\ (1 \leq \operatorname{mrk\_}) \wedge \operatorname{rw\text{-}or\text{-}r}(\operatorname{mrt\_}, \operatorname{mrs\_}, \operatorname{mrx\_}) \vee \\ (1 \leq \operatorname{mrk\_}) \wedge \operatorname{rw\text{-}imp\text{-}l}(\operatorname{mrt\_}, \operatorname{mrs\_}, \operatorname{mrx\_}) \vee \\ (1 \leq \operatorname{mrk\_}) \wedge \operatorname{rw\text{-}imp\text{-}r}(\operatorname{mrt\_}, \operatorname{mrs\_}, \operatorname{mrx\_}) \vee \\ (1 \leq \operatorname{mrk\_}) \wedge \operatorname{rw\text{-}iff\text{-}l}(\operatorname{mrt\_}, \operatorname{mrs\_}, \operatorname{mrx\_}) \vee \\ (1 \leq \operatorname{mrk\_}) \wedge \operatorname{rw\text{-}iff\text{-}r}(\operatorname{mrt\_}, \operatorname{mrs\_}, \operatorname{mrx\_}) \vee \\ (1 \leq \operatorname{mrk\_}) \wedge \operatorname{rw\text{-}rel\text{-}l}(\operatorname{mrt\_}, \operatorname{mrs\_}, \operatorname{mrx\_}) \vee \\ (1 \leq \operatorname{mrk\_}) \wedge \operatorname{rw\text{-}rel\text{-}r}(\operatorname{mrt\_}, \operatorname{mrs\_}, \operatorname{mrx\_}) \vee \\ (2 \leq \operatorname{mrk\_}) \wedge \operatorname{rw\text{-}all}(\operatorname{mrt\_}, \operatorname{mrs\_}, \operatorname{mrx\_}) \vee \\ (2 \leq \operatorname{mrk\_}) \wedge \operatorname{rw\text{-}ex}(\operatorname{mrt\_}, \operatorname{mrs\_}, \operatorname{mrx\_}) \vee \\ (4 \leq \operatorname{mrk\_}) \wedge \operatorname{rw\text{-}pred}(\operatorname{mrt\_}, \operatorname{mrs\_}, \operatorname{mrx\_}) \vee \\ (4 \leq \operatorname{mrk\_}) \wedge \operatorname{rw\text{-}tapp}(\operatorname{mrt\_}, \operatorname{mrs\_}, \operatorname{mrx\_}) \vee \\ (5 \leq \operatorname{mrk\_}) \wedge \operatorname{rw\text{-}bind\text{-}dom}(\operatorname{mrt\_}, \operatorname{mrs\_}, \operatorname{mrx\_}) \vee \\ (6 \leq \operatorname{mrk\_}) \wedge \operatorname{rw\text{-}sep\text{-}body}(\operatorname{mrt\_}, \operatorname{mrs\_}, \operatorname{mrx\_}) \vee \\ (6 \leq \operatorname{mrk\_}) \wedge \operatorname{rw\text{-}tbind\text{-}body}(\operatorname{mrt\_}, \operatorname{mrs\_}, \operatorname{mrx\_}) \vee \\ (7 \leq \operatorname{mrk\_}) \wedge \operatorname{rw\text{-}desc\text{-}body}(\operatorname{mrt\_}, \operatorname{mrs\_}, \operatorname{mrx\_}) \vee \\ (3 \leq \operatorname{mrk\_}) \wedge \operatorname{rw\text{-}if\text{-}cond}(\operatorname{mrt\_}, \operatorname{mrs\_}, \operatorname{mrx\_}) \vee \\ (3 \leq \operatorname{mrk\_}) \wedge \operatorname{rw\text{-}if\text{-}then}(\operatorname{mrt\_}, \operatorname{mrs\_}, \operatorname{mrx\_}) \vee \\ (3 \leq \operatorname{mrk\_}) \wedge \operatorname{rw\text{-}if\text{-}else}(\operatorname{mrt\_}, \operatorname{mrs\_}, \operatorname{mrx\_}) \end{array}
\end{equation*}

\par\medskip\noindent\textbf{Definition} (\texttt{rw-closed}).\;For all $\operatorname{mrt\_}$, $\operatorname{mrk\_}$ and $\operatorname{mrs\_}$, 
\begin{equation*}
\begin{array}{@{}l@{}} \operatorname{rw\text{-}closed}(\operatorname{mrt\_}, \operatorname{mrk\_}, \operatorname{mrs\_}) \\ \Leftrightarrow\; \forall \operatorname{mrx\_}. \\ \quad \operatorname{rw\text{-}rule}(\operatorname{mrt\_}, \operatorname{mrk\_}, \operatorname{mrs\_}, \operatorname{mrx\_}) \Rightarrow ((\operatorname{mrx\_} \in \operatorname{rw\text{-}univ}) \Rightarrow (\operatorname{mrx\_} \in \operatorname{mrs\_})) \end{array}
\end{equation*}

\par\medskip\noindent\textbf{Definition} (\texttt{rw-top}).\;For all $\operatorname{mrt\_}$, $\operatorname{mrs\_}$ and $\operatorname{mrx\_}$,  $\operatorname{rw\text{-}top}(\operatorname{mrt\_}, \operatorname{mrs\_}, \operatorname{mrx\_}) \Leftrightarrow (\operatorname{mrx\_} \in \operatorname{mrt\_})$
\par\medskip\noindent\textbf{Definition} (\texttt{rw-not}).\;For all $\operatorname{mrt\_}$, $\operatorname{mrs\_}$ and $\operatorname{mrx\_}$, 
\begin{equation*}
\begin{array}{@{}l@{}} \operatorname{rw\text{-}not}(\operatorname{mrt\_}, \operatorname{mrs\_}, \operatorname{mrx\_}) \\ \Leftrightarrow\; \exists \operatorname{mrg\_}, \operatorname{mra\_}, \operatorname{mrb\_}. \\ \quad (\operatorname{mrx\_} = \langle 2, \operatorname{mrg\_}, \operatorname{enc\text{-}not}(\operatorname{mra\_}), \operatorname{enc\text{-}not}(\operatorname{mrb\_}) \rangle) \wedge (\langle 2, \operatorname{mrg\_}, \operatorname{mra\_}, \operatorname{mrb\_} \rangle \in \operatorname{mrs\_}) \end{array}
\end{equation*}

\par\medskip\noindent\textbf{Definition} (\texttt{rw-and-l}).\;For all $\operatorname{mrt\_}$, $\operatorname{mrs\_}$ and $\operatorname{mrx\_}$, 
\begin{equation*}
\begin{array}{@{}l@{}} \operatorname{rw\text{-}and\text{-}l}(\operatorname{mrt\_}, \operatorname{mrs\_}, \operatorname{mrx\_}) \\ \Leftrightarrow\; \exists \operatorname{mrg\_}, \operatorname{mra\_}, \operatorname{mrb\_}, \operatorname{mrc\_}. \\ \quad \operatorname{mrx\_} = \langle 2, \operatorname{mrg\_}, \operatorname{enc\text{-}and}(\operatorname{mra\_}, \operatorname{mrc\_}), \operatorname{enc\text{-}and}(\operatorname{mrb\_}, \operatorname{mrc\_}) \rangle \wedge \\ \quad \operatorname{mrc\_} \in \mathit{clf} \wedge \\ \quad \langle 2, \operatorname{ctx\text{-}add}(\operatorname{mrg\_}, \operatorname{mrc\_}), \operatorname{mra\_}, \operatorname{mrb\_} \rangle \in \operatorname{mrs\_} \end{array}
\end{equation*}

\par\medskip\noindent\textbf{Definition} (\texttt{rw-and-r}).\;For all $\operatorname{mrt\_}$, $\operatorname{mrs\_}$ and $\operatorname{mrx\_}$, 
\begin{equation*}
\begin{array}{@{}l@{}} \operatorname{rw\text{-}and\text{-}r}(\operatorname{mrt\_}, \operatorname{mrs\_}, \operatorname{mrx\_}) \\ \Leftrightarrow\; \exists \operatorname{mrg\_}, \operatorname{mra\_}, \operatorname{mrb\_}, \operatorname{mrc\_}. \\ \quad \operatorname{mrx\_} = \langle 2, \operatorname{mrg\_}, \operatorname{enc\text{-}and}(\operatorname{mrc\_}, \operatorname{mra\_}), \operatorname{enc\text{-}and}(\operatorname{mrc\_}, \operatorname{mrb\_}) \rangle \wedge \\ \quad \operatorname{mrc\_} \in \mathit{clf} \wedge \\ \quad \langle 2, \operatorname{ctx\text{-}add}(\operatorname{mrg\_}, \operatorname{mrc\_}), \operatorname{mra\_}, \operatorname{mrb\_} \rangle \in \operatorname{mrs\_} \end{array}
\end{equation*}

\par\medskip\noindent\textbf{Definition} (\texttt{rw-or-l}).\;For all $\operatorname{mrt\_}$, $\operatorname{mrs\_}$ and $\operatorname{mrx\_}$, 
\begin{equation*}
\begin{array}{@{}l@{}} \operatorname{rw\text{-}or\text{-}l}(\operatorname{mrt\_}, \operatorname{mrs\_}, \operatorname{mrx\_}) \\ \Leftrightarrow\; \exists \operatorname{mrg\_}, \operatorname{mra\_}, \operatorname{mrb\_}, \operatorname{mrc\_}. \\ \quad \operatorname{mrx\_} = \langle 2, \operatorname{mrg\_}, \operatorname{enc\text{-}or}(\operatorname{mra\_}, \operatorname{mrc\_}), \operatorname{enc\text{-}or}(\operatorname{mrb\_}, \operatorname{mrc\_}) \rangle \wedge \\ \quad \operatorname{mrc\_} \in \mathit{clf} \wedge \\ \quad \langle 2, \operatorname{ctx\text{-}add}(\operatorname{mrg\_}, \operatorname{enc\text{-}not}(\operatorname{mrc\_})), \operatorname{mra\_}, \operatorname{mrb\_} \rangle \in \operatorname{mrs\_} \end{array}
\end{equation*}

\par\medskip\noindent\textbf{Definition} (\texttt{rw-or-r}).\;For all $\operatorname{mrt\_}$, $\operatorname{mrs\_}$ and $\operatorname{mrx\_}$, 
\begin{equation*}
\begin{array}{@{}l@{}} \operatorname{rw\text{-}or\text{-}r}(\operatorname{mrt\_}, \operatorname{mrs\_}, \operatorname{mrx\_}) \\ \Leftrightarrow\; \exists \operatorname{mrg\_}, \operatorname{mra\_}, \operatorname{mrb\_}, \operatorname{mrc\_}. \\ \quad \operatorname{mrx\_} = \langle 2, \operatorname{mrg\_}, \operatorname{enc\text{-}or}(\operatorname{mrc\_}, \operatorname{mra\_}), \operatorname{enc\text{-}or}(\operatorname{mrc\_}, \operatorname{mrb\_}) \rangle \wedge \\ \quad \operatorname{mrc\_} \in \mathit{clf} \wedge \\ \quad \langle 2, \operatorname{ctx\text{-}add}(\operatorname{mrg\_}, \operatorname{enc\text{-}not}(\operatorname{mrc\_})), \operatorname{mra\_}, \operatorname{mrb\_} \rangle \in \operatorname{mrs\_} \end{array}
\end{equation*}

\par\medskip\noindent\textbf{Definition} (\texttt{rw-imp-l}).\;For all $\operatorname{mrt\_}$, $\operatorname{mrs\_}$ and $\operatorname{mrx\_}$, 
\begin{equation*}
\begin{array}{@{}l@{}} \operatorname{rw\text{-}imp\text{-}l}(\operatorname{mrt\_}, \operatorname{mrs\_}, \operatorname{mrx\_}) \\ \Leftrightarrow\; \exists \operatorname{mrg\_}, \operatorname{mra\_}, \operatorname{mrb\_}, \operatorname{mrc\_}. \\ \quad \operatorname{mrx\_} = \langle 2, \operatorname{mrg\_}, \operatorname{enc\text{-}imp}(\operatorname{mra\_}, \operatorname{mrc\_}), \operatorname{enc\text{-}imp}(\operatorname{mrb\_}, \operatorname{mrc\_}) \rangle \wedge \\ \quad \operatorname{mrc\_} \in \mathit{clf} \wedge \\ \quad \langle 2, \operatorname{mrg\_}, \operatorname{mra\_}, \operatorname{mrb\_} \rangle \in \operatorname{mrs\_} \end{array}
\end{equation*}

\par\medskip\noindent\textbf{Definition} (\texttt{rw-imp-r}).\;For all $\operatorname{mrt\_}$, $\operatorname{mrs\_}$ and $\operatorname{mrx\_}$, 
\begin{equation*}
\begin{array}{@{}l@{}} \operatorname{rw\text{-}imp\text{-}r}(\operatorname{mrt\_}, \operatorname{mrs\_}, \operatorname{mrx\_}) \\ \Leftrightarrow\; \exists \operatorname{mrg\_}, \operatorname{mra\_}, \operatorname{mrb\_}, \operatorname{mrc\_}. \\ \quad \operatorname{mrx\_} = \langle 2, \operatorname{mrg\_}, \operatorname{enc\text{-}imp}(\operatorname{mrc\_}, \operatorname{mra\_}), \operatorname{enc\text{-}imp}(\operatorname{mrc\_}, \operatorname{mrb\_}) \rangle \wedge \\ \quad \operatorname{mrc\_} \in \mathit{clf} \wedge \\ \quad \langle 2, \operatorname{ctx\text{-}add}(\operatorname{mrg\_}, \operatorname{mrc\_}), \operatorname{mra\_}, \operatorname{mrb\_} \rangle \in \operatorname{mrs\_} \end{array}
\end{equation*}

\par\medskip\noindent\textbf{Definition} (\texttt{rw-iff-l}).\;For all $\operatorname{mrt\_}$, $\operatorname{mrs\_}$ and $\operatorname{mrx\_}$, 
\begin{equation*}
\begin{array}{@{}l@{}} \operatorname{rw\text{-}iff\text{-}l}(\operatorname{mrt\_}, \operatorname{mrs\_}, \operatorname{mrx\_}) \\ \Leftrightarrow\; \exists \operatorname{mrg\_}, \operatorname{mra\_}, \operatorname{mrb\_}, \operatorname{mrc\_}. \\ \quad \operatorname{mrx\_} = \langle 2, \operatorname{mrg\_}, \operatorname{enc\text{-}iff}(\operatorname{mra\_}, \operatorname{mrc\_}), \operatorname{enc\text{-}iff}(\operatorname{mrb\_}, \operatorname{mrc\_}) \rangle \wedge \\ \quad \operatorname{mrc\_} \in \mathit{clf} \wedge \\ \quad \langle 2, \operatorname{mrg\_}, \operatorname{mra\_}, \operatorname{mrb\_} \rangle \in \operatorname{mrs\_} \end{array}
\end{equation*}

\par\medskip\noindent\textbf{Definition} (\texttt{rw-iff-r}).\;For all $\operatorname{mrt\_}$, $\operatorname{mrs\_}$ and $\operatorname{mrx\_}$, 
\begin{equation*}
\begin{array}{@{}l@{}} \operatorname{rw\text{-}iff\text{-}r}(\operatorname{mrt\_}, \operatorname{mrs\_}, \operatorname{mrx\_}) \\ \Leftrightarrow\; \exists \operatorname{mrg\_}, \operatorname{mra\_}, \operatorname{mrb\_}, \operatorname{mrc\_}. \\ \quad \operatorname{mrx\_} = \langle 2, \operatorname{mrg\_}, \operatorname{enc\text{-}iff}(\operatorname{mrc\_}, \operatorname{mra\_}), \operatorname{enc\text{-}iff}(\operatorname{mrc\_}, \operatorname{mrb\_}) \rangle \wedge \\ \quad \operatorname{mrc\_} \in \mathit{clf} \wedge \\ \quad \langle 2, \operatorname{mrg\_}, \operatorname{mra\_}, \operatorname{mrb\_} \rangle \in \operatorname{mrs\_} \end{array}
\end{equation*}

\par\medskip\noindent\textbf{Definition} (\texttt{rw-rel-l}).\;For all $\operatorname{mrt\_}$, $\operatorname{mrs\_}$ and $\operatorname{mrx\_}$, 
\begin{equation*}
\begin{array}{@{}l@{}} \operatorname{rw\text{-}rel\text{-}l}(\operatorname{mrt\_}, \operatorname{mrs\_}, \operatorname{mrx\_}) \\ \Leftrightarrow\; \exists \operatorname{mrg\_}, \operatorname{mrh\_}, \operatorname{mra\_}, \operatorname{mrb\_}, \operatorname{mrc\_}. \\ \quad \operatorname{mrx\_} = \langle 2, \operatorname{mrg\_}, \operatorname{eapp}(\operatorname{mrh\_}, \langle \operatorname{mra\_}, \operatorname{mrc\_} \rangle), \operatorname{eapp}(\operatorname{mrh\_}, \langle \operatorname{mrb\_}, \operatorname{mrc\_} \rangle) \rangle \wedge \\ \quad \operatorname{mrh\_} \in \mathbb{N} \wedge \\ \quad 7 \leq \operatorname{mrh\_} \wedge \\ \quad \operatorname{mrh\_} \leq 9 \wedge \\ \quad \operatorname{mrc\_} \in \mathit{clt} \wedge \\ \quad \langle 1, \operatorname{mrg\_}, \operatorname{mra\_}, \operatorname{mrb\_} \rangle \in \operatorname{mrs\_} \end{array}
\end{equation*}

\par\medskip\noindent\textbf{Definition} (\texttt{rw-rel-r}).\;For all $\operatorname{mrt\_}$, $\operatorname{mrs\_}$ and $\operatorname{mrx\_}$, 
\begin{equation*}
\begin{array}{@{}l@{}} \operatorname{rw\text{-}rel\text{-}r}(\operatorname{mrt\_}, \operatorname{mrs\_}, \operatorname{mrx\_}) \\ \Leftrightarrow\; \exists \operatorname{mrg\_}, \operatorname{mrh\_}, \operatorname{mra\_}, \operatorname{mrb\_}, \operatorname{mrc\_}. \\ \quad \operatorname{mrx\_} = \langle 2, \operatorname{mrg\_}, \operatorname{eapp}(\operatorname{mrh\_}, \langle \operatorname{mrc\_}, \operatorname{mra\_} \rangle), \operatorname{eapp}(\operatorname{mrh\_}, \langle \operatorname{mrc\_}, \operatorname{mrb\_} \rangle) \rangle \wedge \\ \quad \operatorname{mrh\_} \in \mathbb{N} \wedge \\ \quad 7 \leq \operatorname{mrh\_} \wedge \\ \quad \operatorname{mrh\_} \leq 9 \wedge \\ \quad \operatorname{mrc\_} \in \mathit{clt} \wedge \\ \quad \langle 1, \operatorname{mrg\_}, \operatorname{mra\_}, \operatorname{mrb\_} \rangle \in \operatorname{mrs\_} \end{array}
\end{equation*}

\par\medskip\noindent\textbf{Definition} (\texttt{rw-all}).\;For all $\operatorname{mrt\_}$, $\operatorname{mrs\_}$ and $\operatorname{mrx\_}$, 
\begin{equation*}
\begin{array}{@{}l@{}} \operatorname{rw\text{-}all}(\operatorname{mrt\_}, \operatorname{mrs\_}, \operatorname{mrx\_}) \\ \Leftrightarrow\; \exists \operatorname{mrg\_}, \operatorname{mra\_}, \operatorname{mrb\_}. \\ \quad \operatorname{mrx\_} = \langle 2, \operatorname{mrg\_}, \operatorname{enc\text{-}all}(\operatorname{mra\_}), \operatorname{enc\text{-}all}(\operatorname{mrb\_}) \rangle \wedge \\ \quad \operatorname{enc\text{-}all}(\operatorname{mra\_}) \in \mathit{clf} \wedge \\ \quad \operatorname{enc\text{-}all}(\operatorname{mrb\_}) \in \mathit{clf} \wedge \\ \quad \exists \operatorname{mrfs\_} \in \mathbf{Set}. \\ \quad \quad \operatorname{mrfs\_} \subseteq \mathbb{N} \wedge \\ \quad \quad |\operatorname{mrfs\_}| \in \mathbb{N} \wedge \\ \quad \quad \forall \operatorname{mry\_} \in \mathbb{N}. \\ \quad \quad \quad \neg (\operatorname{mry\_} \in \operatorname{mrfs\_}) \Rightarrow  \\ \quad \quad \quad \quad \langle 2, \operatorname{mrg\_}, \operatorname{expr\text{-}open}(\operatorname{mra\_}, 0, \operatorname{efvar}(\operatorname{mry\_})), \operatorname{expr\text{-}open}(\operatorname{mrb\_}, 0, \operatorname{efvar}(\operatorname{mry\_})) \rangle \in \operatorname{mrs\_} \end{array}
\end{equation*}

\par\medskip\noindent\textbf{Definition} (\texttt{rw-ex}).\;For all $\operatorname{mrt\_}$, $\operatorname{mrs\_}$ and $\operatorname{mrx\_}$, 
\begin{equation*}
\begin{array}{@{}l@{}} \operatorname{rw\text{-}ex}(\operatorname{mrt\_}, \operatorname{mrs\_}, \operatorname{mrx\_}) \\ \Leftrightarrow\; \exists \operatorname{mrg\_}, \operatorname{mra\_}, \operatorname{mrb\_}. \\ \quad \operatorname{mrx\_} = \langle 2, \operatorname{mrg\_}, \operatorname{enc\text{-}ex}(\operatorname{mra\_}), \operatorname{enc\text{-}ex}(\operatorname{mrb\_}) \rangle \wedge \\ \quad \operatorname{enc\text{-}ex}(\operatorname{mra\_}) \in \mathit{clf} \wedge \\ \quad \operatorname{enc\text{-}ex}(\operatorname{mrb\_}) \in \mathit{clf} \wedge \\ \quad \exists \operatorname{mrfs\_} \in \mathbf{Set}. \\ \quad \quad \operatorname{mrfs\_} \subseteq \mathbb{N} \wedge \\ \quad \quad |\operatorname{mrfs\_}| \in \mathbb{N} \wedge \\ \quad \quad \forall \operatorname{mry\_} \in \mathbb{N}. \\ \quad \quad \quad \neg (\operatorname{mry\_} \in \operatorname{mrfs\_}) \Rightarrow  \\ \quad \quad \quad \quad \langle 2, \operatorname{mrg\_}, \operatorname{expr\text{-}open}(\operatorname{mra\_}, 0, \operatorname{efvar}(\operatorname{mry\_})), \operatorname{expr\text{-}open}(\operatorname{mrb\_}, 0, \operatorname{efvar}(\operatorname{mry\_})) \rangle \in \operatorname{mrs\_} \end{array}
\end{equation*}

\par\medskip\noindent\textbf{Definition} (\texttt{rw-pred}).\;For all $\operatorname{mrt\_}$, $\operatorname{mrs\_}$ and $\operatorname{mrx\_}$, 
\begin{equation*}
\begin{array}{@{}l@{}} \operatorname{rw\text{-}pred}(\operatorname{mrt\_}, \operatorname{mrs\_}, \operatorname{mrx\_}) \\ \Leftrightarrow\; \exists \operatorname{mrg\_}, \operatorname{mrh\_}, \operatorname{mra\_}, \operatorname{mrb\_}, \operatorname{mri\_}. \\ \quad \operatorname{mrx\_} = \langle 2, \operatorname{mrg\_}, \operatorname{eapp}(\operatorname{mrh\_}, \operatorname{mra\_}), \operatorname{eapp}(\operatorname{mrh\_}, \operatorname{mrb\_}) \rangle \wedge \\ \quad \operatorname{mrh\_} \in \mathbb{N} \wedge \\ \quad \operatorname{mrh\_} = 10 \wedge \\ \quad \operatorname{eapp}(\operatorname{mrh\_}, \operatorname{mra\_}) \in \mathit{clf} \wedge \\ \quad \operatorname{eapp}(\operatorname{mrh\_}, \operatorname{mrb\_}) \in \mathit{clf} \wedge \\ \quad \operatorname{mra\_} \in \operatorname{tuples}(\operatorname{expr\text{-}set}) \wedge \\ \quad \operatorname{mrb\_} \in \operatorname{tuples}(\operatorname{expr\text{-}set}) \wedge \\ \quad \operatorname{length}(\operatorname{mrb\_}) = \operatorname{length}(\operatorname{mra\_}) \wedge \\ \quad \operatorname{mri\_} \in \mathbb{N} \wedge \\ \quad 1 \leq \operatorname{mri\_} \wedge \\ \quad \operatorname{mri\_} \leq \operatorname{length}(\operatorname{mra\_}) \wedge \\ \quad \forall \operatorname{mrj\_} \in \mathbb{N}. \\ \quad \quad (1 \leq \operatorname{mrj\_}) \Rightarrow  \\ \quad \quad (\operatorname{mrj\_} \leq \operatorname{length}(\operatorname{mra\_})) \Rightarrow  \\ \quad \quad \neg (\operatorname{mrj\_} = \operatorname{mri\_}) \Rightarrow (\operatorname{nth}(\operatorname{mrj\_}, \operatorname{mrb\_}) = \operatorname{nth}(\operatorname{mrj\_}, \operatorname{mra\_})) \wedge \\ \quad \langle 1, \operatorname{mrg\_}, \operatorname{nth}(\operatorname{mri\_}, \operatorname{mra\_}), \operatorname{nth}(\operatorname{mri\_}, \operatorname{mrb\_}) \rangle \in \operatorname{mrs\_} \end{array}
\end{equation*}

\par\medskip\noindent\textbf{Definition} (\texttt{rw-tapp}).\;For all $\operatorname{mrt\_}$, $\operatorname{mrs\_}$ and $\operatorname{mrx\_}$, 
\begin{equation*}
\begin{array}{@{}l@{}} \operatorname{rw\text{-}tapp}(\operatorname{mrt\_}, \operatorname{mrs\_}, \operatorname{mrx\_}) \\ \Leftrightarrow\; \exists \operatorname{mrg\_}, \operatorname{mrh\_}, \operatorname{mra\_}, \operatorname{mrb\_}, \operatorname{mri\_}. \\ \quad \operatorname{mrx\_} = \langle 1, \operatorname{mrg\_}, \operatorname{eapp}(\operatorname{mrh\_}, \operatorname{mra\_}), \operatorname{eapp}(\operatorname{mrh\_}, \operatorname{mrb\_}) \rangle \wedge \\ \quad \operatorname{mrh\_} \in \mathbb{N} \wedge \\ \quad 11 \leq \operatorname{mrh\_} \wedge \\ \quad \operatorname{eapp}(\operatorname{mrh\_}, \operatorname{mra\_}) \in \mathit{clt} \wedge \\ \quad \operatorname{eapp}(\operatorname{mrh\_}, \operatorname{mrb\_}) \in \mathit{clt} \wedge \\ \quad \operatorname{mra\_} \in \operatorname{tuples}(\operatorname{expr\text{-}set}) \wedge \\ \quad \operatorname{mrb\_} \in \operatorname{tuples}(\operatorname{expr\text{-}set}) \wedge \\ \quad \operatorname{length}(\operatorname{mrb\_}) = \operatorname{length}(\operatorname{mra\_}) \wedge \\ \quad \operatorname{mri\_} \in \mathbb{N} \wedge \\ \quad 1 \leq \operatorname{mri\_} \wedge \\ \quad \operatorname{mri\_} \leq \operatorname{length}(\operatorname{mra\_}) \wedge \\ \quad \forall \operatorname{mrj\_} \in \mathbb{N}. \\ \quad \quad (1 \leq \operatorname{mrj\_}) \Rightarrow  \\ \quad \quad (\operatorname{mrj\_} \leq \operatorname{length}(\operatorname{mra\_})) \Rightarrow  \\ \quad \quad \neg (\operatorname{mrj\_} = \operatorname{mri\_}) \Rightarrow (\operatorname{nth}(\operatorname{mrj\_}, \operatorname{mrb\_}) = \operatorname{nth}(\operatorname{mrj\_}, \operatorname{mra\_})) \wedge \\ \quad \langle 1, \operatorname{mrg\_}, \operatorname{nth}(\operatorname{mri\_}, \operatorname{mra\_}), \operatorname{nth}(\operatorname{mri\_}, \operatorname{mrb\_}) \rangle \in \operatorname{mrs\_} \end{array}
\end{equation*}

\par\medskip\noindent\textbf{Definition} (\texttt{rw-bind-dom}).\;For all $\operatorname{mrt\_}$, $\operatorname{mrs\_}$ and $\operatorname{mrx\_}$, 
\begin{equation*}
\begin{array}{@{}l@{}} \operatorname{rw\text{-}bind\text{-}dom}(\operatorname{mrt\_}, \operatorname{mrs\_}, \operatorname{mrx\_}) \\ \Leftrightarrow\; \exists \operatorname{mrg\_}, \operatorname{mrh\_}, \operatorname{mra\_}, \operatorname{mrb\_}, \operatorname{mrc\_}. \\ \quad \operatorname{mrx\_} = \langle 1, \operatorname{mrg\_}, \operatorname{ebind}(\operatorname{mrh\_}, 1, \langle \operatorname{mra\_} \rangle, \operatorname{mrc\_}), \operatorname{ebind}(\operatorname{mrh\_}, 1, \langle \operatorname{mrb\_} \rangle, \operatorname{mrc\_}) \rangle \wedge \\ \quad (\operatorname{mrh\_} = 22) \vee ((\operatorname{mrh\_} = 24) \vee (\operatorname{mrh\_} = 26)) \wedge \\ \quad \operatorname{ebind}(\operatorname{mrh\_}, 1, \langle \operatorname{mra\_} \rangle, \operatorname{mrc\_}) \in \mathit{clt} \wedge \\ \quad \operatorname{ebind}(\operatorname{mrh\_}, 1, \langle \operatorname{mrb\_} \rangle, \operatorname{mrc\_}) \in \mathit{clt} \wedge \\ \quad \langle 1, \operatorname{mrg\_}, \operatorname{mra\_}, \operatorname{mrb\_} \rangle \in \operatorname{mrs\_} \end{array}
\end{equation*}

\par\medskip\noindent\textbf{Definition} (\texttt{rw-sep-body}).\;For all $\operatorname{mrt\_}$, $\operatorname{mrs\_}$ and $\operatorname{mrx\_}$, 
\begin{equation*}
\begin{array}{@{}l@{}} \operatorname{rw\text{-}sep\text{-}body}(\operatorname{mrt\_}, \operatorname{mrs\_}, \operatorname{mrx\_}) \\ \Leftrightarrow\; \exists \operatorname{mrg\_}, \operatorname{mra\_}, \operatorname{mrb\_}, \operatorname{mrc\_}. \\ \quad \operatorname{mrx\_} = \langle 1, \operatorname{mrg\_}, \operatorname{enc\text{-}sep}(\operatorname{mrc\_}, \operatorname{mra\_}), \operatorname{enc\text{-}sep}(\operatorname{mrc\_}, \operatorname{mrb\_}) \rangle \wedge \\ \quad \operatorname{enc\text{-}sep}(\operatorname{mrc\_}, \operatorname{mra\_}) \in \mathit{clt} \wedge \\ \quad \operatorname{enc\text{-}sep}(\operatorname{mrc\_}, \operatorname{mrb\_}) \in \mathit{clt} \wedge \\ \quad \exists \operatorname{mrfs\_} \in \mathbf{Set}. \\ \quad \quad \operatorname{mrfs\_} \subseteq \mathbb{N} \wedge \\ \quad \quad |\operatorname{mrfs\_}| \in \mathbb{N} \wedge \\ \quad \quad \forall \operatorname{mry\_} \in \mathbb{N}. \\ \quad \quad \quad \neg (\operatorname{mry\_} \in \operatorname{mrfs\_}) \Rightarrow  \\ \quad \quad \quad \quad \langle 2, \operatorname{mrg\_}, \operatorname{expr\text{-}open}(\operatorname{mra\_}, 0, \operatorname{efvar}(\operatorname{mry\_})), \operatorname{expr\text{-}open}(\operatorname{mrb\_}, 0, \operatorname{efvar}(\operatorname{mry\_})) \rangle \in \operatorname{mrs\_} \end{array}
\end{equation*}

\par\medskip\noindent\textbf{Definition} (\texttt{rw-tbind-body}).\;For all $\operatorname{mrt\_}$, $\operatorname{mrs\_}$ and $\operatorname{mrx\_}$, 
\begin{equation*}
\begin{array}{@{}l@{}} \operatorname{rw\text{-}tbind\text{-}body}(\operatorname{mrt\_}, \operatorname{mrs\_}, \operatorname{mrx\_}) \\ \Leftrightarrow\; \exists \operatorname{mrg\_}, \operatorname{mrh\_}, \operatorname{mra\_}, \operatorname{mrb\_}, \operatorname{mrc\_}. \\ \quad \operatorname{mrx\_} = \langle 1, \operatorname{mrg\_}, \operatorname{ebind}(\operatorname{mrh\_}, 1, \langle \operatorname{mrc\_} \rangle, \operatorname{mra\_}), \operatorname{ebind}(\operatorname{mrh\_}, 1, \langle \operatorname{mrc\_} \rangle, \operatorname{mrb\_}) \rangle \wedge \\ \quad (\operatorname{mrh\_} = 24) \vee (\operatorname{mrh\_} = 26) \wedge \\ \quad \operatorname{ebind}(\operatorname{mrh\_}, 1, \langle \operatorname{mrc\_} \rangle, \operatorname{mra\_}) \in \mathit{clt} \wedge \\ \quad \operatorname{ebind}(\operatorname{mrh\_}, 1, \langle \operatorname{mrc\_} \rangle, \operatorname{mrb\_}) \in \mathit{clt} \wedge \\ \quad \exists \operatorname{mrfs\_} \in \mathbf{Set}. \\ \quad \quad \operatorname{mrfs\_} \subseteq \mathbb{N} \wedge \\ \quad \quad |\operatorname{mrfs\_}| \in \mathbb{N} \wedge \\ \quad \quad \forall \operatorname{mry\_} \in \mathbb{N}. \\ \quad \quad \quad \neg (\operatorname{mry\_} \in \operatorname{mrfs\_}) \Rightarrow  \\ \quad \quad \quad \quad \langle 1, \operatorname{mrg\_}, \operatorname{expr\text{-}open}(\operatorname{mra\_}, 0, \operatorname{efvar}(\operatorname{mry\_})), \operatorname{expr\text{-}open}(\operatorname{mrb\_}, 0, \operatorname{efvar}(\operatorname{mry\_})) \rangle \in \operatorname{mrs\_} \end{array}
\end{equation*}

\par\medskip\noindent\textbf{Definition} (\texttt{rw-desc-body}).\;For all $\operatorname{mrt\_}$, $\operatorname{mrs\_}$ and $\operatorname{mrx\_}$, 
\begin{equation*}
\begin{array}{@{}l@{}} \operatorname{rw\text{-}desc\text{-}body}(\operatorname{mrt\_}, \operatorname{mrs\_}, \operatorname{mrx\_}) \\ \Leftrightarrow\; \exists \operatorname{mrg\_}, \operatorname{mrh\_}, \operatorname{mra\_}, \operatorname{mrb\_}. \\ \quad \operatorname{mrx\_} = \langle 1, \operatorname{mrg\_}, \operatorname{ebind}(\operatorname{mrh\_}, 1, \langle  \rangle, \operatorname{mra\_}), \operatorname{ebind}(\operatorname{mrh\_}, 1, \langle  \rangle, \operatorname{mrb\_}) \rangle \wedge \\ \quad (\operatorname{mrh\_} = 23) \vee (\operatorname{mrh\_} = 25) \wedge \\ \quad \operatorname{ebind}(\operatorname{mrh\_}, 1, \langle  \rangle, \operatorname{mra\_}) \in \mathit{clt} \wedge \\ \quad \operatorname{ebind}(\operatorname{mrh\_}, 1, \langle  \rangle, \operatorname{mrb\_}) \in \mathit{clt} \wedge \\ \quad \exists \operatorname{mrfs\_} \in \mathbf{Set}. \\ \quad \quad \operatorname{mrfs\_} \subseteq \mathbb{N} \wedge \\ \quad \quad |\operatorname{mrfs\_}| \in \mathbb{N} \wedge \\ \quad \quad \forall \operatorname{mry\_} \in \mathbb{N}. \\ \quad \quad \quad \neg (\operatorname{mry\_} \in \operatorname{mrfs\_}) \Rightarrow  \\ \quad \quad \quad \quad \langle 2, \operatorname{mrg\_}, \operatorname{expr\text{-}open}(\operatorname{mra\_}, 0, \operatorname{efvar}(\operatorname{mry\_})), \operatorname{expr\text{-}open}(\operatorname{mrb\_}, 0, \operatorname{efvar}(\operatorname{mry\_})) \rangle \in \operatorname{mrs\_} \end{array}
\end{equation*}

\par\medskip\noindent\textbf{Definition} (\texttt{rw-if-cond}).\;For all $\operatorname{mrt\_}$, $\operatorname{mrs\_}$ and $\operatorname{mrx\_}$, 
\begin{equation*}
\begin{array}{@{}l@{}} \operatorname{rw\text{-}if\text{-}cond}(\operatorname{mrt\_}, \operatorname{mrs\_}, \operatorname{mrx\_}) \\ \Leftrightarrow\; \exists \operatorname{mrg\_}, \operatorname{mra\_}, \operatorname{mrb\_}, \operatorname{mrc\_}, \operatorname{mrd\_}. \\ \quad \operatorname{mrx\_} = \langle 1, \operatorname{mrg\_}, \operatorname{enc\text{-}if}(\operatorname{mra\_}, \operatorname{mrc\_}, \operatorname{mrd\_}), \operatorname{enc\text{-}if}(\operatorname{mrb\_}, \operatorname{mrc\_}, \operatorname{mrd\_}) \rangle \wedge \\ \quad \operatorname{enc\text{-}if}(\operatorname{mra\_}, \operatorname{mrc\_}, \operatorname{mrd\_}) \in \mathit{clt} \wedge \\ \quad \operatorname{enc\text{-}if}(\operatorname{mrb\_}, \operatorname{mrc\_}, \operatorname{mrd\_}) \in \mathit{clt} \wedge \\ \quad \operatorname{mrc\_} \in \mathit{clt} \wedge \\ \quad \operatorname{mrd\_} \in \mathit{clt} \wedge \\ \quad \langle 2, \operatorname{mrg\_}, \operatorname{mra\_}, \operatorname{mrb\_} \rangle \in \operatorname{mrs\_} \end{array}
\end{equation*}

\par\medskip\noindent\textbf{Definition} (\texttt{rw-if-then}).\;For all $\operatorname{mrt\_}$, $\operatorname{mrs\_}$ and $\operatorname{mrx\_}$, 
\begin{equation*}
\begin{array}{@{}l@{}} \operatorname{rw\text{-}if\text{-}then}(\operatorname{mrt\_}, \operatorname{mrs\_}, \operatorname{mrx\_}) \\ \Leftrightarrow\; \exists \operatorname{mrg\_}, \operatorname{mra\_}, \operatorname{mrb\_}, \operatorname{mrc\_}, \operatorname{mrd\_}. \\ \quad \operatorname{mrx\_} = \langle 1, \operatorname{mrg\_}, \operatorname{enc\text{-}if}(\operatorname{mrc\_}, \operatorname{mra\_}, \operatorname{mrd\_}), \operatorname{enc\text{-}if}(\operatorname{mrc\_}, \operatorname{mrb\_}, \operatorname{mrd\_}) \rangle \wedge \\ \quad \operatorname{enc\text{-}if}(\operatorname{mrc\_}, \operatorname{mra\_}, \operatorname{mrd\_}) \in \mathit{clt} \wedge \\ \quad \operatorname{enc\text{-}if}(\operatorname{mrc\_}, \operatorname{mrb\_}, \operatorname{mrd\_}) \in \mathit{clt} \wedge \\ \quad \operatorname{mrc\_} \in \mathit{clf} \wedge \\ \quad \operatorname{mrd\_} \in \mathit{clt} \wedge \\ \quad \langle 1, \operatorname{mrg\_}, \operatorname{mra\_}, \operatorname{mrb\_} \rangle \in \operatorname{mrs\_} \end{array}
\end{equation*}

\par\medskip\noindent\textbf{Definition} (\texttt{rw-if-else}).\;For all $\operatorname{mrt\_}$, $\operatorname{mrs\_}$ and $\operatorname{mrx\_}$, 
\begin{equation*}
\begin{array}{@{}l@{}} \operatorname{rw\text{-}if\text{-}else}(\operatorname{mrt\_}, \operatorname{mrs\_}, \operatorname{mrx\_}) \\ \Leftrightarrow\; \exists \operatorname{mrg\_}, \operatorname{mra\_}, \operatorname{mrb\_}, \operatorname{mrc\_}, \operatorname{mrd\_}. \\ \quad \operatorname{mrx\_} = \langle 1, \operatorname{mrg\_}, \operatorname{enc\text{-}if}(\operatorname{mrc\_}, \operatorname{mrd\_}, \operatorname{mra\_}), \operatorname{enc\text{-}if}(\operatorname{mrc\_}, \operatorname{mrd\_}, \operatorname{mrb\_}) \rangle \wedge \\ \quad \operatorname{enc\text{-}if}(\operatorname{mrc\_}, \operatorname{mrd\_}, \operatorname{mra\_}) \in \mathit{clt} \wedge \\ \quad \operatorname{enc\text{-}if}(\operatorname{mrc\_}, \operatorname{mrd\_}, \operatorname{mrb\_}) \in \mathit{clt} \wedge \\ \quad \operatorname{mrc\_} \in \mathit{clf} \wedge \\ \quad \operatorname{mrd\_} \in \mathit{clt} \wedge \\ \quad \langle 1, \operatorname{mrg\_}, \operatorname{mra\_}, \operatorname{mrb\_} \rangle \in \operatorname{mrs\_} \end{array}
\end{equation*}

\par\medskip\noindent\textbf{Definition} (\texttt{rel-rewrite-1}).\;For all $\operatorname{mrt\_}$, $\operatorname{mrhs\_}$ and $\operatorname{mrc\_}$, 
\begin{equation*}
\begin{array}{@{}l@{}} \operatorname{rel\text{-}rewrite\text{-}1}(\operatorname{mrt\_}, \operatorname{mrhs\_}, \operatorname{mrc\_}) \\ \Leftrightarrow\; \operatorname{mrc\_} \in \operatorname{seq\text{-}set} \wedge \\ \exists \operatorname{mrg\_}, \operatorname{mra\_}, \operatorname{mrb\_}. \\ \quad ((\operatorname{mrc\_} = \langle \operatorname{mrg\_}, \operatorname{mra\_} \rangle) \wedge ((\operatorname{mrhs\_} = \langle \langle \operatorname{mrg\_}, \operatorname{mrb\_} \rangle \rangle) \wedge (\langle 2, \operatorname{mrg\_}, \operatorname{mra\_}, \operatorname{mrb\_} \rangle \in \operatorname{rewr}(\operatorname{mrt\_}, 5)))) \end{array}
\end{equation*}

\par\medskip\noindent\textbf{Definition} (\texttt{rel-rewrite-hyp-1}).\;For all $\operatorname{mrt\_}$, $\operatorname{mrhs\_}$ and $\operatorname{mrc\_}$, 
\begin{equation*}
\begin{array}{@{}l@{}} \operatorname{rel\text{-}rewrite\text{-}hyp\text{-}1}(\operatorname{mrt\_}, \operatorname{mrhs\_}, \operatorname{mrc\_}) \\ \Leftrightarrow\; \operatorname{mrc\_} \in \operatorname{seq\text{-}set} \wedge \\ \exists \operatorname{mrg\_}, \operatorname{mrf\_}, \operatorname{mra\_}, \operatorname{mrb\_}. \\ \quad \operatorname{mrc\_} = \langle \operatorname{mrg\_}, \operatorname{mrf\_} \rangle \wedge \\ \quad \operatorname{mra\_} \in \operatorname{mrg\_} \wedge \\ \quad \operatorname{mrhs\_} = \langle \langle \operatorname{ctx\text{-}add}(\operatorname{difference}(\operatorname{mrg\_}, \{\operatorname{mra\_}\}), \operatorname{mrb\_}), \operatorname{mrf\_} \rangle \rangle \wedge \\ \quad \langle 2, \operatorname{mrg\_}, \operatorname{mra\_}, \operatorname{mrb\_} \rangle \in \operatorname{rewr}(\operatorname{mrt\_}, 5) \end{array}
\end{equation*}

\par\medskip\noindent\textbf{Definition} (\texttt{j-nth}).\;$\operatorname{j\text{-}nth}(\operatorname{mrt\_}) :=$
\begin{equation*}
\{\operatorname{mrx\_} \in \operatorname{rw\text{-}univ} : \exists \operatorname{mrg\_}, \operatorname{mrl\_}, \operatorname{mrk\_}.\; ((\operatorname{mrx\_} = \langle 1, \operatorname{mrg\_}, \operatorname{enc\text{-}nth}(\operatorname{enc\text{-}num}(\operatorname{mrk\_}), \operatorname{enc\text{-}list}(\operatorname{mrl\_})), \operatorname{nth}(\operatorname{mrk\_}, \operatorname{mrl\_}) \rangle) \wedge ((\operatorname{mrk\_} \in \mathbb{N}) \wedge ((1 \leq \operatorname{mrk\_}) \wedge ((\operatorname{mrk\_} \leq \operatorname{length}(\operatorname{mrl\_})) \wedge ((\operatorname{enc\text{-}nth}(\operatorname{enc\text{-}num}(\operatorname{mrk\_}), \operatorname{enc\text{-}list}(\operatorname{mrl\_})) \in \mathit{clt}) \wedge ((\operatorname{nth}(\operatorname{mrk\_}, \operatorname{mrl\_}) \in \mathit{clt}) \wedge (\langle \operatorname{mrg\_}, \operatorname{enc\text{-}eq}(\operatorname{enc\text{-}list}(\operatorname{mrl\_}), \operatorname{enc\text{-}list}(\operatorname{mrl\_})) \rangle \in \operatorname{derivable}(\operatorname{mrt\_}))))))))\}
\end{equation*}

\par\medskip\noindent\textbf{Definition} (\texttt{j-len}).\;$\operatorname{j\text{-}len}(\operatorname{mrt\_}) :=$
\begin{equation*}
\{\operatorname{mrx\_} \in \operatorname{rw\text{-}univ} : \exists \operatorname{mrg\_}, \operatorname{mrl\_}.\; ((\operatorname{mrx\_} = \langle 1, \operatorname{mrg\_}, \operatorname{enc\text{-}length}(\operatorname{enc\text{-}list}(\operatorname{mrl\_})), \operatorname{enc\text{-}num}(\operatorname{length}(\operatorname{mrl\_})) \rangle) \wedge ((\operatorname{enc\text{-}length}(\operatorname{enc\text{-}list}(\operatorname{mrl\_})) \in \mathit{clt}) \wedge ((\operatorname{enc\text{-}num}(\operatorname{length}(\operatorname{mrl\_})) \in \mathit{clt}) \wedge (\langle \operatorname{mrg\_}, \operatorname{enc\text{-}eq}(\operatorname{enc\text{-}list}(\operatorname{mrl\_}), \operatorname{enc\text{-}list}(\operatorname{mrl\_})) \rangle \in \operatorname{derivable}(\operatorname{mrt\_})))))\}
\end{equation*}

\par\medskip\noindent\textbf{Definition} (\texttt{j-lb}).\;$\operatorname{j\text{-}lb}(\operatorname{mrt\_}) :=$
\begin{equation*}
\{\operatorname{mrx\_} \in \operatorname{rw\text{-}univ} : \exists \operatorname{mrg\_}, \operatorname{mra\_}, \operatorname{mrb\_}, \operatorname{mre\_}.\; ((\operatorname{mrx\_} = \langle 1, \operatorname{mrg\_}, \operatorname{enc\text{-}apply}(\operatorname{enc\text{-}lam}(\operatorname{mra\_}, \operatorname{mrb\_}), \operatorname{mre\_}), \operatorname{expr\text{-}open}(\operatorname{mrb\_}, 0, \operatorname{mre\_}) \rangle) \wedge ((\operatorname{enc\text{-}apply}(\operatorname{enc\text{-}lam}(\operatorname{mra\_}, \operatorname{mrb\_}), \operatorname{mre\_}) \in \mathit{clt}) \wedge ((\operatorname{expr\text{-}open}(\operatorname{mrb\_}, 0, \operatorname{mre\_}) \in \mathit{clt}) \wedge (\langle \operatorname{mrg\_}, \operatorname{enc\text{-}in}(\operatorname{mre\_}, \operatorname{mra\_}) \rangle \in \operatorname{derivable}(\operatorname{mrt\_})))))\}
\end{equation*}

\par\medskip\noindent\textbf{Definition} (\texttt{j-mac}).\;$\operatorname{j\text{-}mac}(\operatorname{mrt\_}) :=$
\begin{equation*}
\{\operatorname{mrx\_} \in \operatorname{rw\text{-}univ} : \exists \operatorname{mrg\_}, \operatorname{mrc\_}, \operatorname{mra\_}, \operatorname{mrb\_}.; \exists \operatorname{mre\_} \in \mathit{clt}.\; ((\langle \operatorname{mrg\_}, \operatorname{expr\text{-}open}(\operatorname{mrc\_}, 0, \operatorname{mre\_}) \rangle \in \operatorname{derivable}(\operatorname{mrt\_})) \wedge ((\langle \operatorname{mrg\_}, \operatorname{enc\text{-}eq}(\operatorname{mre\_}, \operatorname{mre\_}) \rangle \in \operatorname{derivable}(\operatorname{mrt\_})) \wedge (((\operatorname{mrx\_} = \langle 1, \operatorname{mrg\_}, \operatorname{expr\text{-}open}(\operatorname{mra\_}, 0, \operatorname{mre\_}), \operatorname{expr\text{-}open}(\operatorname{mrb\_}, 0, \operatorname{mre\_}) \rangle) \wedge ((\operatorname{enc\text{-}all}(\operatorname{enc\text{-}imp}(\operatorname{mrc\_}, \operatorname{enc\text{-}qeq}(\operatorname{mra\_}, \operatorname{mrb\_}))) \in \operatorname{mrt\_}) \wedge ((\operatorname{enc\text{-}all}(\operatorname{enc\text{-}imp}(\operatorname{mrc\_}, \operatorname{enc\text{-}qeq}(\operatorname{mra\_}, \operatorname{mrb\_}))) \in \mathit{clf}) \wedge ((\operatorname{expr\text{-}open}(\operatorname{mra\_}, 0, \operatorname{mre\_}) \in \mathit{clt}) \wedge (\operatorname{expr\text{-}open}(\operatorname{mrb\_}, 0, \operatorname{mre\_}) \in \mathit{clt}))))) \vee ((\operatorname{mrx\_} = \langle 2, \operatorname{mrg\_}, \operatorname{expr\text{-}open}(\operatorname{mra\_}, 0, \operatorname{mre\_}), \operatorname{expr\text{-}open}(\operatorname{mrb\_}, 0, \operatorname{mre\_}) \rangle) \wedge ((\operatorname{enc\text{-}all}(\operatorname{enc\text{-}imp}(\operatorname{mrc\_}, \operatorname{enc\text{-}iff}(\operatorname{mra\_}, \operatorname{mrb\_}))) \in \operatorname{mrt\_}) \wedge ((\operatorname{enc\text{-}all}(\operatorname{enc\text{-}imp}(\operatorname{mrc\_}, \operatorname{enc\text{-}iff}(\operatorname{mra\_}, \operatorname{mrb\_}))) \in \mathit{clf}) \wedge ((\operatorname{expr\text{-}open}(\operatorname{mra\_}, 0, \operatorname{mre\_}) \in \mathit{clf}) \wedge (\operatorname{expr\text{-}open}(\operatorname{mrb\_}, 0, \operatorname{mre\_}) \in \mathit{clf}))))))))\}
\end{equation*}

\par\medskip\noindent\textbf{Definition} (\texttt{rel-macete-1}).\;For all $\operatorname{mrt\_}$, $\operatorname{mrhs\_}$ and $\operatorname{mrc\_}$, 
\begin{equation*}
\operatorname{rel\text{-}macete\text{-}1}(\operatorname{mrt\_}, \operatorname{mrhs\_}, \operatorname{mrc\_}) \Leftrightarrow \operatorname{rel\text{-}rewrite\text{-}1}(\operatorname{j\text{-}mac}(\operatorname{mrt\_}), \operatorname{mrhs\_}, \operatorname{mrc\_})
\end{equation*}

\par\medskip\noindent\textbf{Definition} (\texttt{rel-nth-reduce-1}).\;For all $\operatorname{mrt\_}$, $\operatorname{mrhs\_}$ and $\operatorname{mrc\_}$, 
\begin{equation*}
\operatorname{rel\text{-}nth\text{-}reduce\text{-}1}(\operatorname{mrt\_}, \operatorname{mrhs\_}, \operatorname{mrc\_}) \Leftrightarrow \operatorname{rel\text{-}rewrite\text{-}1}(\operatorname{j\text{-}nth}(\operatorname{mrt\_}), \operatorname{mrhs\_}, \operatorname{mrc\_})
\end{equation*}

\par\medskip\noindent\textbf{Definition} (\texttt{rel-length-reduce-1}).\;For all $\operatorname{mrt\_}$, $\operatorname{mrhs\_}$ and $\operatorname{mrc\_}$, 
\begin{equation*}
\operatorname{rel\text{-}length\text{-}reduce\text{-}1}(\operatorname{mrt\_}, \operatorname{mrhs\_}, \operatorname{mrc\_}) \Leftrightarrow \operatorname{rel\text{-}rewrite\text{-}1}(\operatorname{j\text{-}len}(\operatorname{mrt\_}), \operatorname{mrhs\_}, \operatorname{mrc\_})
\end{equation*}

\par\medskip\noindent\textbf{Definition} (\texttt{rel-lambda-beta-1}).\;For all $\operatorname{mrt\_}$, $\operatorname{mrhs\_}$ and $\operatorname{mrc\_}$, 
\begin{equation*}
\operatorname{rel\text{-}lambda\text{-}beta\text{-}1}(\operatorname{mrt\_}, \operatorname{mrhs\_}, \operatorname{mrc\_}) \Leftrightarrow \operatorname{rel\text{-}rewrite\text{-}1}(\operatorname{j\text{-}lb}(\operatorname{mrt\_}), \operatorname{mrhs\_}, \operatorname{mrc\_})
\end{equation*}

\par\medskip\noindent\textbf{Definition} (\texttt{rel-lambda-beta-hyp-1}).\;For all $\operatorname{mrt\_}$, $\operatorname{mrhs\_}$ and $\operatorname{mrc\_}$, 
\begin{equation*}
\operatorname{rel\text{-}lambda\text{-}beta\text{-}hyp\text{-}1}(\operatorname{mrt\_}, \operatorname{mrhs\_}, \operatorname{mrc\_}) \Leftrightarrow \operatorname{rel\text{-}rewrite\text{-}hyp\text{-}1}(\operatorname{j\text{-}lb}(\operatorname{mrt\_}), \operatorname{mrhs\_}, \operatorname{mrc\_})
\end{equation*}

\par\medskip\noindent\textbf{Definition} (\texttt{rel-macete-hyp-1}).\;For all $\operatorname{mrt\_}$, $\operatorname{mrhs\_}$ and $\operatorname{mrc\_}$, 
\begin{equation*}
\operatorname{rel\text{-}macete\text{-}hyp\text{-}1}(\operatorname{mrt\_}, \operatorname{mrhs\_}, \operatorname{mrc\_}) \Leftrightarrow \operatorname{rel\text{-}rewrite\text{-}hyp\text{-}1}(\operatorname{j\text{-}mac}(\operatorname{mrt\_}), \operatorname{mrhs\_}, \operatorname{mrc\_})
\end{equation*}


\vnbbookfile{theorem-library/meta-cr-laws}{theorem-library/meta-cr-laws}
\noindent{\small meta-cr-laws.scm -- CONTEXTUAL REWRITING (CR) as an admissibility theorem over the encoded core, and the admissibility of the rewriting rules at one position (structure-library/meta-cr.scm). (docs/meta-k2-admissibility-2026-10-08.md, section 6, phase 3.) Agent META-3, 2026-10-09.  Builds on ...}\par
\vnbprop{cartesian4-set}{\vnbstmt{\begin{array}{@{}l@{}} \forall a_{1}, a_{2}, a_{3}, a_{4}. \\ \quad ((a_{1} \in \mathbf{Set}) \wedge ((a_{2} \in \mathbf{Set}) \wedge ((a_{3} \in \mathbf{Set}) \wedge (a_{4} \in \mathbf{Set})))) \Rightarrow  \\ \quad \quad (a_{1} \times a_{2}) \in \mathbf{Set} \end{array}}}
\vnbprop{rw-univ-is-set}{\vnbstmt{\operatorname{rw\text{-}univ} \in \mathbf{Set}}}
\vnbprop{eqv-unfold}{\vnbstmt{\forall t.\; (\operatorname{eqv}(t) \simeq \{x \in \operatorname{rw\text{-}univ} : \operatorname{eqv\text{-}ok}(t, x)\})}}
\vnbprop{rewr-unfold}{\vnbstmt{\begin{array}{@{}l@{}} \forall t, k. \\ \quad (\operatorname{rewr}(t, k) \simeq \{x \in \operatorname{rw\text{-}univ} : \forall s \in \operatorname{power}(\operatorname{rw\text{-}univ}).\; (\operatorname{rw\text{-}closed}(t, k, s) \Rightarrow (x \in s))\}) \end{array}}}
\vnbprop{eqv-in-power}{\vnbstmt{\forall t.\; (\operatorname{eqv}(t) \in \operatorname{power}(\operatorname{rw\text{-}univ}))}}
\vnbprop{rewr-least}{\vnbstmt{\begin{array}{@{}l@{}} \forall t, k.; \forall s \in \operatorname{power}(\operatorname{rw\text{-}univ}), x \in \operatorname{rewr}(t, k). \\ \quad \operatorname{rw\text{-}closed}(t, k, s) \Rightarrow (x \in s) \end{array}}}
\vnbprop{eqv-f-open}{\vnbstmt{\begin{array}{@{}l@{}} \forall t, g, a, b. \\ \quad (\langle 2, g, a, b \rangle \in \operatorname{eqv}(t)) \Rightarrow  \\ \quad \quad (a \in \mathit{clf}) \wedge ((b \in \mathit{clf}) \wedge (\langle g, \operatorname{enc\text{-}iff}(a, b) \rangle \in \operatorname{derivable}(t))) \end{array}}}
\vnbprop{eqv-t-open}{\vnbstmt{\begin{array}{@{}l@{}} \forall t, g, a, b. \\ \quad (\langle 1, g, a, b \rangle \in \operatorname{eqv}(t)) \Rightarrow  \\ \quad \quad (a \in \mathit{clt}) \wedge ((b \in \mathit{clt}) \wedge (\langle g, \operatorname{enc\text{-}qeq}(a, b) \rangle \in \operatorname{derivable}(t))) \end{array}}}
\vnbprop{eqv-f-intro}{\vnbstmt{\begin{array}{@{}l@{}} \forall t, g, a, b. \\ \quad ((a \in \mathit{clf}) \wedge ((b \in \mathit{clf}) \wedge (\langle g, \operatorname{enc\text{-}iff}(a, b) \rangle \in \operatorname{derivable}(t)))) \Rightarrow  \\ \quad \quad \langle 2, g, a, b \rangle \in \operatorname{eqv}(t) \end{array}}}
\vnbprop{eqv-t-intro}{\vnbstmt{\begin{array}{@{}l@{}} \forall t, g, a, b. \\ \quad ((a \in \mathit{clt}) \wedge ((b \in \mathit{clt}) \wedge (\langle g, \operatorname{enc\text{-}qeq}(a, b) \rangle \in \operatorname{derivable}(t)))) \Rightarrow  \\ \quad \quad \langle 1, g, a, b \rangle \in \operatorname{eqv}(t) \end{array}}}
\vnbprop{rw-univ-ctx}{\vnbstmt{\forall k, g, a, b.\; ((\langle k, g, a, b \rangle \in \operatorname{rw\text{-}univ}) \Rightarrow (g \in \operatorname{ctx\text{-}set}))}}
\vnbprop{der-iff-r}{\vnbstmt{\begin{array}{@{}l@{}} \forall t, g, a, b. \\ \quad ((\langle g, \operatorname{enc\text{-}iff}(a, b) \rangle \in \operatorname{derivable}(t)) \wedge (\langle g, a \rangle \in \operatorname{derivable}(t))) \Rightarrow  \\ \quad \quad \langle g, b \rangle \in \operatorname{derivable}(t) \end{array}}}
\vnbprop{der-iff-l}{\vnbstmt{\begin{array}{@{}l@{}} \forall t, g, a, b. \\ \quad ((\langle g, \operatorname{enc\text{-}iff}(a, b) \rangle \in \operatorname{derivable}(t)) \wedge (\langle g, b \rangle \in \operatorname{derivable}(t))) \Rightarrow  \\ \quad \quad \langle g, a \rangle \in \operatorname{derivable}(t) \end{array}}}
\vnbprop{eqv-top}{\vnbstmt{\begin{array}{@{}l@{}} \forall j, t, x. \\ \quad ((j \subseteq \operatorname{eqv}(t)) \wedge (\operatorname{rw\text{-}top}(j, \operatorname{eqv}(t), x) \wedge (x \in \operatorname{rw\text{-}univ}))) \Rightarrow (x \in \operatorname{eqv}(t)) \end{array}}}
\vnbprop{eqv-not}{\vnbstmt{\forall j, t, x.\; ((\operatorname{rw\text{-}not}(j, \operatorname{eqv}(t), x) \wedge (x \in \operatorname{rw\text{-}univ})) \Rightarrow (x \in \operatorname{eqv}(t)))}}
\vnbprop{eqv-and-l}{\vnbstmt{\forall j, t, x.\; ((\operatorname{rw\text{-}and\text{-}l}(j, \operatorname{eqv}(t), x) \wedge (x \in \operatorname{rw\text{-}univ})) \Rightarrow (x \in \operatorname{eqv}(t)))}}
\vnbprop{eqv-and-r}{\vnbstmt{\forall j, t, x.\; ((\operatorname{rw\text{-}and\text{-}r}(j, \operatorname{eqv}(t), x) \wedge (x \in \operatorname{rw\text{-}univ})) \Rightarrow (x \in \operatorname{eqv}(t)))}}
\vnbprop{eqv-or-l}{\vnbstmt{\forall j, t, x.\; ((\operatorname{rw\text{-}or\text{-}l}(j, \operatorname{eqv}(t), x) \wedge (x \in \operatorname{rw\text{-}univ})) \Rightarrow (x \in \operatorname{eqv}(t)))}}
\vnbprop{eqv-or-r}{\vnbstmt{\forall j, t, x.\; ((\operatorname{rw\text{-}or\text{-}r}(j, \operatorname{eqv}(t), x) \wedge (x \in \operatorname{rw\text{-}univ})) \Rightarrow (x \in \operatorname{eqv}(t)))}}
\vnbprop{eqv-imp-l}{\vnbstmt{\forall j, t, x.\; ((\operatorname{rw\text{-}imp\text{-}l}(j, \operatorname{eqv}(t), x) \wedge (x \in \operatorname{rw\text{-}univ})) \Rightarrow (x \in \operatorname{eqv}(t)))}}
\vnbprop{eqv-imp-r}{\vnbstmt{\forall j, t, x.\; ((\operatorname{rw\text{-}imp\text{-}r}(j, \operatorname{eqv}(t), x) \wedge (x \in \operatorname{rw\text{-}univ})) \Rightarrow (x \in \operatorname{eqv}(t)))}}
\vnbprop{eqv-iff-l}{\vnbstmt{\forall j, t, x.\; ((\operatorname{rw\text{-}iff\text{-}l}(j, \operatorname{eqv}(t), x) \wedge (x \in \operatorname{rw\text{-}univ})) \Rightarrow (x \in \operatorname{eqv}(t)))}}
\vnbprop{eqv-iff-r}{\vnbstmt{\forall j, t, x.\; ((\operatorname{rw\text{-}iff\text{-}r}(j, \operatorname{eqv}(t), x) \wedge (x \in \operatorname{rw\text{-}univ})) \Rightarrow (x \in \operatorname{eqv}(t)))}}
\vnbprop{subst-app2}{\vnbstmt{\begin{array}{@{}l@{}} \forall m, u, h, a, b. \\ \quad ((u \in \operatorname{expr\text{-}set}) \wedge ((h \in \mathbb{N}) \wedge ((a \in \operatorname{expr\text{-}set}) \wedge (b \in \operatorname{expr\text{-}set})))) \Rightarrow  \\ \quad \quad \operatorname{expr\text{-}subst}(\operatorname{eapp}(h, \langle a, b \rangle), m, u) = \operatorname{eapp}(h, \langle \operatorname{expr\text{-}subst}(a, m, u), \operatorname{expr\text{-}subst}(b, m, u) \rangle) \end{array}}}
\vnbprop{subst-hole-l}{\vnbstmt{\begin{array}{@{}l@{}} \forall z, s, c, h. \\ \quad ((z \in \mathbb{N}) \wedge ((s \in \operatorname{expr\text{-}set}) \wedge ((c \in \operatorname{expr\text{-}set}) \wedge ((h \in \mathbb{N}) \wedge \neg (z \in \operatorname{expr\text{-}fv}(c)))))) \Rightarrow  \\ \quad \quad \operatorname{expr\text{-}subst}(\operatorname{eapp}(h, \langle \operatorname{efvar}(z), c \rangle), z, s) = \operatorname{eapp}(h, \langle s, c \rangle) \end{array}}}
\vnbprop{subst-hole-r}{\vnbstmt{\begin{array}{@{}l@{}} \forall z, s, c, h. \\ \quad ((z \in \mathbb{N}) \wedge ((s \in \operatorname{expr\text{-}set}) \wedge ((c \in \operatorname{expr\text{-}set}) \wedge ((h \in \mathbb{N}) \wedge \neg (z \in \operatorname{expr\text{-}fv}(c)))))) \Rightarrow  \\ \quad \quad \operatorname{expr\text{-}subst}(\operatorname{eapp}(h, \langle c, \operatorname{efvar}(z) \rangle), z, s) = \operatorname{eapp}(h, \langle c, s \rangle) \end{array}}}
\vnbprop{der-atom-l}{\vnbstmt{\begin{array}{@{}l@{}} \forall a, g, h, s, t, c. \\ \quad h \in \mathbb{N} \wedge \\ \quad s \in \mathit{clt} \wedge \\ \quad t \in \mathit{clt} \wedge \\ \quad c \in \mathit{clt} \wedge \\ \quad \operatorname{eapp}(h, \langle s, c \rangle) \in \mathit{clf} \wedge \\ \quad \langle g, \operatorname{enc\text{-}qeq}(s, t) \rangle \in \operatorname{derivable}(a) \wedge \\ \quad \langle g, \operatorname{eapp}(h, \langle t, c \rangle) \rangle \in \operatorname{derivable}(a) \Rightarrow \\ \quad \quad \langle g, \operatorname{eapp}(h, \langle s, c \rangle) \rangle \in \operatorname{derivable}(a) \end{array}}}
\vnbprop{der-atom-r}{\vnbstmt{\begin{array}{@{}l@{}} \forall a, g, h, s, t, c. \\ \quad h \in \mathbb{N} \wedge \\ \quad s \in \mathit{clt} \wedge \\ \quad t \in \mathit{clt} \wedge \\ \quad c \in \mathit{clt} \wedge \\ \quad \operatorname{eapp}(h, \langle c, s \rangle) \in \mathit{clf} \wedge \\ \quad \langle g, \operatorname{enc\text{-}qeq}(s, t) \rangle \in \operatorname{derivable}(a) \wedge \\ \quad \langle g, \operatorname{eapp}(h, \langle c, t \rangle) \rangle \in \operatorname{derivable}(a) \Rightarrow \\ \quad \quad \langle g, \operatorname{eapp}(h, \langle c, s \rangle) \rangle \in \operatorname{derivable}(a) \end{array}}}
\vnbprop{der-qeq-sym}{\vnbstmt{\begin{array}{@{}l@{}} \forall a, g, s, t. \\ \quad ((s \in \mathit{clt}) \wedge ((t \in \mathit{clt}) \wedge (\langle g, \operatorname{enc\text{-}qeq}(s, t) \rangle \in \operatorname{derivable}(a)))) \Rightarrow  \\ \quad \quad \langle g, \operatorname{enc\text{-}qeq}(t, s) \rangle \in \operatorname{derivable}(a) \end{array}}}
\vnbprop{eqv-rel-l}{\vnbstmt{\forall j, t, x.\; ((\operatorname{rw\text{-}rel\text{-}l}(j, \operatorname{eqv}(t), x) \wedge (x \in \operatorname{rw\text{-}univ})) \Rightarrow (x \in \operatorname{eqv}(t)))}}
\vnbprop{eqv-rel-r}{\vnbstmt{\forall j, t, x.\; ((\operatorname{rw\text{-}rel\text{-}r}(j, \operatorname{eqv}(t), x) \wedge (x \in \operatorname{rw\text{-}univ})) \Rightarrow (x \in \operatorname{eqv}(t)))}}
\vnbprop{enc-all-unfold}{\vnbstmt{\forall p.\; (\operatorname{enc\text{-}all}(p) \simeq \operatorname{ebind}(20, 1, \langle  \rangle, p))}}
\vnbprop{enc-ex-unfold}{\vnbstmt{\forall p.\; (\operatorname{enc\text{-}ex}(p) \simeq \operatorname{ebind}(21, 1, \langle  \rangle, p))}}
\vnbprop{clt-efvar}{\vnbstmt{\forall n \in \mathbb{N}.\; (\operatorname{efvar}(n) \in \mathit{clt})}}
\vnbprop{clf-eq-intro}{\vnbstmt{\forall x, y.\; (((x \in \mathit{clt}) \wedge (y \in \mathit{clt})) \Rightarrow (\operatorname{enc\text{-}eq}(x, y) \in \mathit{clf}))}}
\vnbprop{ebind-stage-inv}{\vnbstmt{\begin{array}{@{}l@{}} \forall n \in \mathbb{N}, h, c, a, b. \\ \quad (\operatorname{ebind}(h, c, a, b) \in \operatorname{expr\text{-}stage}(n)) \Rightarrow (b \in \operatorname{expr\text{-}set}) \end{array}}}
\vnbprop{ebind-body-expr}{\vnbstmt{\forall h, c, a, b.\; ((\operatorname{ebind}(h, c, a, b) \in \operatorname{expr\text{-}set}) \Rightarrow (b \in \operatorname{expr\text{-}set}))}}
\vnbprop{mr-fresh}{\vnbstmt{\begin{array}{@{}l@{}} \forall s, g, f. \\ \quad ((s \in \mathbf{Set}) \wedge ((s \subseteq \mathbb{N}) \wedge ((|s| \in \mathbb{N}) \wedge ((g \in \operatorname{ctx\text{-}set}) \wedge (f \in \operatorname{expr\text{-}set}))))) \Rightarrow  \\ \quad \quad \exists n \in \mathbb{N}.\; (\neg (n \in s) \wedge \neg (n \in (\operatorname{ctx\text{-}fv}(g) \cup \operatorname{expr\text{-}fv}(f)))) \end{array}}}
\vnbprop{eqv-all}{\vnbstmt{\forall j, t, x.\; ((\operatorname{rw\text{-}all}(j, \operatorname{eqv}(t), x) \wedge (x \in \operatorname{rw\text{-}univ})) \Rightarrow (x \in \operatorname{eqv}(t)))}}
\vnbprop{eqv-ex}{\vnbstmt{\forall j, t, x.\; ((\operatorname{rw\text{-}ex}(j, \operatorname{eqv}(t), x) \wedge (x \in \operatorname{rw\text{-}univ})) \Rightarrow (x \in \operatorname{eqv}(t)))}}
\vnbprop{clf-qeq-intro}{\vnbstmt{\forall x, y.\; (((x \in \mathit{clt}) \wedge (y \in \mathit{clt})) \Rightarrow (\operatorname{enc\text{-}qeq}(x, y) \in \mathit{clf}))}}
\vnbprop{der-qeq-join}{\vnbstmt{\begin{array}{@{}l@{}} \forall t, g, x, y, c. \\ \quad x \in \mathit{clt} \wedge \\ \quad y \in \mathit{clt} \wedge \\ \quad c \in \mathit{clt} \wedge \\ \quad \langle g, \operatorname{enc\text{-}qeq}(x, c) \rangle \in \operatorname{derivable}(t) \wedge \\ \quad \langle g, \operatorname{enc\text{-}qeq}(y, c) \rangle \in \operatorname{derivable}(t) \Rightarrow \\ \quad \quad \langle g, \operatorname{enc\text{-}qeq}(x, y) \rangle \in \operatorname{derivable}(t) \end{array}}}
\vnbprop{eqv-if-cond}{\vnbstmt{\forall j, t, x.\; ((\operatorname{rw\text{-}if\text{-}cond}(j, \operatorname{eqv}(t), x) \wedge (x \in \operatorname{rw\text{-}univ})) \Rightarrow (x \in \operatorname{eqv}(t)))}}
\vnbprop{eqv-if-then}{\vnbstmt{\forall j, t, x.\; ((\operatorname{rw\text{-}if\text{-}then}(j, \operatorname{eqv}(t), x) \wedge (x \in \operatorname{rw\text{-}univ})) \Rightarrow (x \in \operatorname{eqv}(t)))}}
\vnbprop{eqv-if-else}{\vnbstmt{\forall j, t, x.\; ((\operatorname{rw\text{-}if\text{-}else}(j, \operatorname{eqv}(t), x) \wedge (x \in \operatorname{rw\text{-}univ})) \Rightarrow (x \in \operatorname{eqv}(t)))}}
\vnbprop{cr-closed-1}{\vnbstmt{\forall t, j.\; ((j \subseteq \operatorname{eqv}(t)) \Rightarrow \operatorname{rw\text{-}closed}(j, 1, \operatorname{eqv}(t)))}}
\vnbprop{cr-closed-2}{\vnbstmt{\forall t, j.\; ((j \subseteq \operatorname{eqv}(t)) \Rightarrow \operatorname{rw\text{-}closed}(j, 2, \operatorname{eqv}(t)))}}
\vnbprop{cr-closed-3}{\vnbstmt{\forall t, j.\; ((j \subseteq \operatorname{eqv}(t)) \Rightarrow \operatorname{rw\text{-}closed}(j, 3, \operatorname{eqv}(t)))}}
\vnbprop{cr-1}{\vnbstmt{\forall t, j, x.\; (((j \subseteq \operatorname{eqv}(t)) \wedge (x \in \operatorname{rewr}(j, 1))) \Rightarrow (x \in \operatorname{eqv}(t)))}}
\vnbprop{cr-2}{\vnbstmt{\forall t, j, x.\; (((j \subseteq \operatorname{eqv}(t)) \wedge (x \in \operatorname{rewr}(j, 2))) \Rightarrow (x \in \operatorname{eqv}(t)))}}
\vnbprop{cr-3}{\vnbstmt{\forall t, j, x.\; (((j \subseteq \operatorname{eqv}(t)) \wedge (x \in \operatorname{rewr}(j, 3))) \Rightarrow (x \in \operatorname{eqv}(t)))}}
\vnbprop{cr-goal-1}{\vnbstmt{\begin{array}{@{}l@{}} \forall t, j, g, o, o_{2}. \\ \quad ((j \subseteq \operatorname{eqv}(t)) \wedge ((\langle 2, g, o, o_{2} \rangle \in \operatorname{rewr}(j, 1)) \wedge (\langle g, o_{2} \rangle \in \operatorname{derivable}(t)))) \Rightarrow  \\ \quad \quad \langle g, o \rangle \in \operatorname{derivable}(t) \end{array}}}
\vnbprop{cr-goal-2}{\vnbstmt{\begin{array}{@{}l@{}} \forall t, j, g, o, o_{2}. \\ \quad ((j \subseteq \operatorname{eqv}(t)) \wedge ((\langle 2, g, o, o_{2} \rangle \in \operatorname{rewr}(j, 2)) \wedge (\langle g, o_{2} \rangle \in \operatorname{derivable}(t)))) \Rightarrow  \\ \quad \quad \langle g, o \rangle \in \operatorname{derivable}(t) \end{array}}}
\vnbprop{cr-goal-3}{\vnbstmt{\begin{array}{@{}l@{}} \forall t, j, g, o, o_{2}. \\ \quad ((j \subseteq \operatorname{eqv}(t)) \wedge ((\langle 2, g, o, o_{2} \rangle \in \operatorname{rewr}(j, 3)) \wedge (\langle g, o_{2} \rangle \in \operatorname{derivable}(t)))) \Rightarrow  \\ \quad \quad \langle g, o \rangle \in \operatorname{derivable}(t) \end{array}}}
\vnbprop{fv-entry}{\vnbstmt{\begin{array}{@{}l@{}} \forall h, l, k, z. \\ \quad h \in \mathbb{N} \wedge \\ \quad l \in \operatorname{tuples}(\operatorname{expr\text{-}set}) \wedge \\ \quad k \in \mathbb{N} \wedge \\ \quad 1 \leq k \wedge \\ \quad k \leq \operatorname{length}(l) \wedge \\ \quad z \in \operatorname{expr\text{-}fv}(\operatorname{nth}(k, l)) \Rightarrow \\ \quad \quad z \in \operatorname{expr\text{-}fv}(\operatorname{eapp}(h, l)) \end{array}}}
\vnbprop{hole-list}{\vnbstmt{\begin{array}{@{}l@{}} \forall z, l, i. \\ \quad ((z \in \mathbb{N}) \wedge ((l \in \operatorname{tuples}(\operatorname{expr\text{-}set})) \wedge ((i \in \mathbb{N}) \wedge ((1 \leq i) \wedge (i \leq \operatorname{length}(l)))))) \Rightarrow  \\ \quad \quad \exists \mathit{mrlz} \in \operatorname{tuples}(\operatorname{expr\text{-}set}). \\ \quad \quad \quad \operatorname{length}(\mathit{mrlz}) = \operatorname{length}(l) \wedge \\ \quad \quad \quad \operatorname{nth}(i, \mathit{mrlz}) = \operatorname{efvar}(z) \wedge \\ \quad \quad \quad \forall j \in \mathbb{N}. \\ \quad \quad \quad \quad ((1 \leq j) \wedge ((j \leq \operatorname{length}(l)) \wedge \neg (j = i))) \Rightarrow (\operatorname{nth}(j, \mathit{mrlz}) = \operatorname{nth}(j, l)) \end{array}}}
\vnbprop{subst-hole-list}{\vnbstmt{\begin{array}{@{}l@{}} \forall z, h, l, a, b, i. \\ \quad z \in \mathbb{N} \wedge \\ \quad h \in \mathbb{N} \wedge \\ \quad l \in \operatorname{tuples}(\operatorname{expr\text{-}set}) \wedge \\ \quad a \in \operatorname{tuples}(\operatorname{expr\text{-}set}) \wedge \\ \quad b \in \operatorname{tuples}(\operatorname{expr\text{-}set}) \wedge \\ \quad \operatorname{length}(a) = \operatorname{length}(l) \wedge \\ \quad \operatorname{length}(b) = \operatorname{length}(l) \wedge \\ \quad i \in \mathbb{N} \wedge \\ \quad 1 \leq i \wedge \\ \quad i \leq \operatorname{length}(l) \wedge \\ \quad \operatorname{nth}(i, a) = \operatorname{efvar}(z) \wedge \\ \quad \forall j \in \mathbb{N}. \\ \quad \quad ((1 \leq j) \wedge ((j \leq \operatorname{length}(l)) \wedge \neg (j = i))) \Rightarrow (\operatorname{nth}(j, a) = \operatorname{nth}(j, l)) \wedge \\ \quad \forall j \in \mathbb{N}. \\ \quad \quad ((1 \leq j) \wedge ((j \leq \operatorname{length}(l)) \wedge \neg (j = i))) \Rightarrow (\operatorname{nth}(j, b) = \operatorname{nth}(j, l)) \wedge \\ \quad \neg (z \in \operatorname{expr\text{-}fv}(\operatorname{eapp}(h, l))) \Rightarrow \\ \quad \quad \operatorname{expr\text{-}subst}(\operatorname{eapp}(h, a), z, \operatorname{nth}(i, b)) = \operatorname{eapp}(h, b) \end{array}}}
\vnbprop{agree-sym}{\vnbstmt{\begin{array}{@{}l@{}} \forall l, a, i.; \forall j \in \mathbb{N}. \\ \quad \operatorname{length}(a) = \operatorname{length}(l) \wedge \\ \quad \forall j \in \mathbb{N}. \\ \quad \quad ((1 \leq j) \wedge ((j \leq \operatorname{length}(l)) \wedge \neg (j = i))) \Rightarrow (\operatorname{nth}(j, a) = \operatorname{nth}(j, l)) \wedge \\ \quad 1 \leq j \wedge \\ \quad j \leq \operatorname{length}(a) \wedge \\ \quad \neg (j = i) \Rightarrow \\ \quad \quad \operatorname{nth}(j, l) = \operatorname{nth}(j, a) \end{array}}}
\vnbprop{der-list-hole}{\vnbstmt{\begin{array}{@{}l@{}} \forall t, g, h, l, a, i. \\ \quad h \in \mathbb{N} \wedge \\ \quad l \in \operatorname{tuples}(\operatorname{expr\text{-}set}) \wedge \\ \quad a \in \operatorname{tuples}(\operatorname{expr\text{-}set}) \wedge \\ \quad \operatorname{length}(a) = \operatorname{length}(l) \wedge \\ \quad i \in \mathbb{N} \wedge \\ \quad 1 \leq i \wedge \\ \quad i \leq \operatorname{length}(l) \wedge \\ \quad \forall j \in \mathbb{N}. \\ \quad \quad ((1 \leq j) \wedge ((j \leq \operatorname{length}(l)) \wedge \neg (j = i))) \Rightarrow (\operatorname{nth}(j, a) = \operatorname{nth}(j, l)) \wedge \\ \quad \operatorname{eapp}(h, l) \in \mathit{clf} \wedge \\ \quad \operatorname{eapp}(h, a) \in \mathit{clf} \wedge \\ \quad \langle g, \operatorname{enc\text{-}qeq}(\operatorname{nth}(i, a), \operatorname{nth}(i, l)) \rangle \in \operatorname{derivable}(t) \wedge \\ \quad \langle g, \operatorname{eapp}(h, l) \rangle \in \operatorname{derivable}(t) \Rightarrow \\ \quad \quad \langle g, \operatorname{eapp}(h, a) \rangle \in \operatorname{derivable}(t) \end{array}}}
\vnbprop{der-tapp}{\vnbstmt{\begin{array}{@{}l@{}} \forall t, g, h, l, a, i. \\ \quad h \in \mathbb{N} \wedge \\ \quad l \in \operatorname{tuples}(\operatorname{expr\text{-}set}) \wedge \\ \quad a \in \operatorname{tuples}(\operatorname{expr\text{-}set}) \wedge \\ \quad \operatorname{length}(a) = \operatorname{length}(l) \wedge \\ \quad i \in \mathbb{N} \wedge \\ \quad 1 \leq i \wedge \\ \quad i \leq \operatorname{length}(l) \wedge \\ \quad \forall j \in \mathbb{N}. \\ \quad \quad ((1 \leq j) \wedge ((j \leq \operatorname{length}(l)) \wedge \neg (j = i))) \Rightarrow (\operatorname{nth}(j, a) = \operatorname{nth}(j, l)) \wedge \\ \quad \operatorname{eapp}(h, l) \in \mathit{clt} \wedge \\ \quad \operatorname{eapp}(h, a) \in \mathit{clt} \wedge \\ \quad \langle g, \operatorname{enc\text{-}qeq}(\operatorname{nth}(i, l), \operatorname{nth}(i, a)) \rangle \in \operatorname{derivable}(t) \Rightarrow \\ \quad \quad \langle g, \operatorname{enc\text{-}qeq}(\operatorname{eapp}(h, l), \operatorname{eapp}(h, a)) \rangle \in \operatorname{derivable}(t) \end{array}}}
\vnbprop{eqv-tapp}{\vnbstmt{\forall j, t, x.\; ((\operatorname{rw\text{-}tapp}(j, \operatorname{eqv}(t), x) \wedge (x \in \operatorname{rw\text{-}univ})) \Rightarrow (x \in \operatorname{eqv}(t)))}}
\vnbprop{eqv-pred}{\vnbstmt{\forall j, t, x.\; ((\operatorname{rw\text{-}pred}(j, \operatorname{eqv}(t), x) \wedge (x \in \operatorname{rw\text{-}univ})) \Rightarrow (x \in \operatorname{eqv}(t)))}}
\vnbprop{cr-closed-4}{\vnbstmt{\forall t, j.\; ((j \subseteq \operatorname{eqv}(t)) \Rightarrow \operatorname{rw\text{-}closed}(j, 4, \operatorname{eqv}(t)))}}
\vnbprop{cr-4}{\vnbstmt{\forall t, j, x.\; (((j \subseteq \operatorname{eqv}(t)) \wedge (x \in \operatorname{rewr}(j, 4))) \Rightarrow (x \in \operatorname{eqv}(t)))}}
\vnbprop{cr-goal-4}{\vnbstmt{\begin{array}{@{}l@{}} \forall t, j, g, o, o_{2}. \\ \quad ((j \subseteq \operatorname{eqv}(t)) \wedge ((\langle 2, g, o, o_{2} \rangle \in \operatorname{rewr}(j, 4)) \wedge (\langle g, o_{2} \rangle \in \operatorname{derivable}(t)))) \Rightarrow  \\ \quad \quad \langle g, o \rangle \in \operatorname{derivable}(t) \end{array}}}
\vnbprop{subst-bind1}{\vnbstmt{\begin{array}{@{}l@{}} \forall m, u, h, c, a, b. \\ \quad ((u \in \operatorname{expr\text{-}set}) \wedge ((h \in \mathbb{N}) \wedge ((c \in \mathbb{N}) \wedge ((a \in \operatorname{expr\text{-}set}) \wedge (b \in \operatorname{expr\text{-}set}))))) \Rightarrow  \\ \quad \quad \operatorname{expr\text{-}subst}(\operatorname{ebind}(h, c, \langle a \rangle, b), m, u) = \operatorname{ebind}(h, c, \langle \operatorname{expr\text{-}subst}(a, m, u) \rangle, \operatorname{expr\text{-}subst}(b, m, u)) \end{array}}}
\vnbprop{fv-bind-body}{\vnbstmt{\begin{array}{@{}l@{}} \forall h, c, a, b, z. \\ \quad ((h \in \mathbb{N}) \wedge ((c \in \mathbb{N}) \wedge ((a \in \operatorname{tuples}(\operatorname{expr\text{-}set})) \wedge ((b \in \operatorname{expr\text{-}set}) \wedge (z \in \operatorname{expr\text{-}fv}(b)))))) \Rightarrow  \\ \quad \quad z \in \operatorname{expr\text{-}fv}(\operatorname{ebind}(h, c, a, b)) \end{array}}}
\vnbprop{der-bind-dom}{\vnbstmt{\begin{array}{@{}l@{}} \forall t, g, h, a, b, c. \\ \quad h \in \mathbb{N} \wedge \\ \quad a \in \mathit{clt} \wedge \\ \quad b \in \mathit{clt} \wedge \\ \quad \operatorname{ebind}(h, 1, \langle a \rangle, c) \in \mathit{clt} \wedge \\ \quad \operatorname{ebind}(h, 1, \langle b \rangle, c) \in \mathit{clt} \wedge \\ \quad \langle g, \operatorname{enc\text{-}qeq}(a, b) \rangle \in \operatorname{derivable}(t) \Rightarrow \\ \quad \quad \langle g, \operatorname{enc\text{-}qeq}(\operatorname{ebind}(h, 1, \langle a \rangle, c), \operatorname{ebind}(h, 1, \langle b \rangle, c)) \rangle \in \operatorname{derivable}(t) \end{array}}}
\vnbprop{eqv-bind-dom}{\vnbstmt{\begin{array}{@{}l@{}} \forall j, t, x. \\ \quad (\operatorname{rw\text{-}bind\text{-}dom}(j, \operatorname{eqv}(t), x) \wedge (x \in \operatorname{rw\text{-}univ})) \Rightarrow (x \in \operatorname{eqv}(t)) \end{array}}}
\vnbprop{cr-closed-5}{\vnbstmt{\forall t, j.\; ((j \subseteq \operatorname{eqv}(t)) \Rightarrow \operatorname{rw\text{-}closed}(j, 5, \operatorname{eqv}(t)))}}
\vnbprop{cr-5}{\vnbstmt{\forall t, j, x.\; (((j \subseteq \operatorname{eqv}(t)) \wedge (x \in \operatorname{rewr}(j, 5))) \Rightarrow (x \in \operatorname{eqv}(t)))}}
\vnbprop{cr-goal-5}{\vnbstmt{\begin{array}{@{}l@{}} \forall t, j, g, o, o_{2}. \\ \quad ((j \subseteq \operatorname{eqv}(t)) \wedge ((\langle 2, g, o, o_{2} \rangle \in \operatorname{rewr}(j, 5)) \wedge (\langle g, o_{2} \rangle \in \operatorname{derivable}(t)))) \Rightarrow  \\ \quad \quad \langle g, o \rangle \in \operatorname{derivable}(t) \end{array}}}
\vnbprop{adm-rewrite-1}{\vnbstmt{\begin{array}{@{}l@{}} \forall t, j, s, c. \\ \quad ((j \subseteq \operatorname{eqv}(t)) \wedge (\operatorname{rel\text{-}rewrite\text{-}1}(j, s, c) \wedge \operatorname{hyps\text{-}in}(s, \operatorname{derivable}(t)))) \Rightarrow  \\ \quad \quad c \in \operatorname{derivable}(t) \end{array}}}
\vnbprop{adm-rewrite-hyp-1}{\vnbstmt{\begin{array}{@{}l@{}} \forall t, j, s, c. \\ \quad ((j \subseteq \operatorname{eqv}(t)) \wedge (\operatorname{rel\text{-}rewrite\text{-}hyp\text{-}1}(j, s, c) \wedge \operatorname{hyps\text{-}in}(s, \operatorname{derivable}(t)))) \Rightarrow  \\ \quad \quad c \in \operatorname{derivable}(t) \end{array}}}
\vnbprop{j-nth-sub}{\vnbstmt{\forall t.\; (\operatorname{j\text{-}nth}(t) \subseteq \operatorname{eqv}(t))}}
\vnbprop{j-len-sub}{\vnbstmt{\forall t.\; (\operatorname{j\text{-}len}(t) \subseteq \operatorname{eqv}(t))}}
\vnbprop{j-lb-sub}{\vnbstmt{\forall t.\; (\operatorname{j\text{-}lb}(t) \subseteq \operatorname{eqv}(t))}}
\vnbprop{open-fun-unfold}{\vnbstmt{\begin{array}{@{}l@{}} \forall u. \\ \quad (\operatorname{open\text{-}fun}(u) \simeq (\lambda \langle \operatorname{mxw\_}, \operatorname{mxj\_} \rangle \in (\operatorname{expr\text{-}set} \times \mathbb{N}).\; \operatorname{expr\text{-}open}(\operatorname{mxw\_}, \operatorname{mxj\_}, u))) \end{array}}}
\vnbprop{open-app2}{\vnbstmt{\begin{array}{@{}l@{}} \forall u, k, h, a, b. \\ \quad ((u \in \operatorname{expr\text{-}set}) \wedge ((k \in \mathbb{N}) \wedge ((h \in \mathbb{N}) \wedge ((a \in \operatorname{expr\text{-}set}) \wedge (b \in \operatorname{expr\text{-}set}))))) \Rightarrow  \\ \quad \quad \operatorname{expr\text{-}open}(\operatorname{eapp}(h, \langle a, b \rangle), k, u) = \operatorname{eapp}(h, \langle \operatorname{expr\text{-}open}(a, k, u), \operatorname{expr\text{-}open}(b, k, u) \rangle) \end{array}}}
\vnbprop{j-mac-sub}{\vnbstmt{\forall t.\; (\operatorname{j\text{-}mac}(t) \subseteq \operatorname{eqv}(t))}}
\vnbprop{adm-nth-reduce-1}{\vnbstmt{\begin{array}{@{}l@{}} \forall t, s, c. \\ \quad (\operatorname{rel\text{-}nth\text{-}reduce\text{-}1}(t, s, c) \wedge \operatorname{hyps\text{-}in}(s, \operatorname{derivable}(t))) \Rightarrow (c \in \operatorname{derivable}(t)) \end{array}}}
\vnbprop{adm-length-reduce-1}{\vnbstmt{\begin{array}{@{}l@{}} \forall t, s, c. \\ \quad (\operatorname{rel\text{-}length\text{-}reduce\text{-}1}(t, s, c) \wedge \operatorname{hyps\text{-}in}(s, \operatorname{derivable}(t))) \Rightarrow  \\ \quad \quad c \in \operatorname{derivable}(t) \end{array}}}
\vnbprop{adm-lambda-beta-1}{\vnbstmt{\begin{array}{@{}l@{}} \forall t, s, c. \\ \quad (\operatorname{rel\text{-}lambda\text{-}beta\text{-}1}(t, s, c) \wedge \operatorname{hyps\text{-}in}(s, \operatorname{derivable}(t))) \Rightarrow  \\ \quad \quad c \in \operatorname{derivable}(t) \end{array}}}
\vnbprop{adm-lambda-beta-hyp-1}{\vnbstmt{\begin{array}{@{}l@{}} \forall t, s, c. \\ \quad (\operatorname{rel\text{-}lambda\text{-}beta\text{-}hyp\text{-}1}(t, s, c) \wedge \operatorname{hyps\text{-}in}(s, \operatorname{derivable}(t))) \Rightarrow  \\ \quad \quad c \in \operatorname{derivable}(t) \end{array}}}
\vnbprop{adm-macete-1}{\vnbstmt{\begin{array}{@{}l@{}} \forall t, s, c. \\ \quad (\operatorname{rel\text{-}macete\text{-}1}(t, s, c) \wedge \operatorname{hyps\text{-}in}(s, \operatorname{derivable}(t))) \Rightarrow (c \in \operatorname{derivable}(t)) \end{array}}}
\vnbprop{adm-macete-hyp-1}{\vnbstmt{\begin{array}{@{}l@{}} \forall t, s, c. \\ \quad (\operatorname{rel\text{-}macete\text{-}hyp\text{-}1}(t, s, c) \wedge \operatorname{hyps\text{-}in}(s, \operatorname{derivable}(t))) \Rightarrow (c \in \operatorname{derivable}(t)) \end{array}}}

\vnbbookfile{structure-library/meta-ind}{structure-library/meta-ind}
\noindent{\small meta-ind.scm -- K2 PHASE 4: the kernel relations of the induction rules, the set-formation rules and the structural rules, TRANSCRIBED from their checkers (rule-checkers-logic.scm, rule-checkers-schema.scm, rule-checkers-rewrite.scm of tree 95bf8fe, 2026-10-12), and the encoded catalog axioms the induction rules are ...}\par
\par\medskip\noindent\textbf{Definition} (\texttt{rel-sep-sethood}).\;For all $\operatorname{mit\_}$, $\operatorname{mihs\_}$ and $\operatorname{mic\_}$, 
\begin{equation*}
\begin{array}{@{}l@{}} \operatorname{rel\text{-}sep\text{-}sethood}(\operatorname{mit\_}, \operatorname{mihs\_}, \operatorname{mic\_}) \\ \Leftrightarrow\; \operatorname{mic\_} \in \operatorname{seq\text{-}set} \wedge \\ \exists \operatorname{mig\_}, \operatorname{mia\_}, \operatorname{mip\_}. \\ \quad ((\operatorname{mic\_} = \langle \operatorname{mig\_}, \operatorname{enc\text{-}in}(\operatorname{enc\text{-}sep}(\operatorname{mia\_}, \operatorname{mip\_}), \operatorname{enc\text{-}setc}) \rangle) \wedge (\operatorname{mihs\_} = \langle \langle \operatorname{mig\_}, \operatorname{enc\text{-}in}(\operatorname{mia\_}, \operatorname{enc\text{-}setc}) \rangle \rangle)) \end{array}
\end{equation*}

\par\medskip\noindent\textbf{Definition} (\texttt{rel-sep-mem-intro}).\;For all $\operatorname{mit\_}$, $\operatorname{mihs\_}$ and $\operatorname{mic\_}$, 
\begin{equation*}
\begin{array}{@{}l@{}} \operatorname{rel\text{-}sep\text{-}mem\text{-}intro}(\operatorname{mit\_}, \operatorname{mihs\_}, \operatorname{mic\_}) \\ \Leftrightarrow\; \operatorname{mic\_} \in \operatorname{seq\text{-}set} \wedge \\ \exists \operatorname{mig\_}, \operatorname{mia\_}, \operatorname{mip\_}, \operatorname{miy\_}. \\ \quad \operatorname{mic\_} = \langle \operatorname{mig\_}, \operatorname{enc\text{-}in}(\operatorname{miy\_}, \operatorname{enc\text{-}sep}(\operatorname{mia\_}, \operatorname{mip\_})) \rangle \wedge \\ \quad \operatorname{mihs\_} = \langle \langle \operatorname{mig\_}, \operatorname{enc\text{-}in}(\operatorname{miy\_}, \operatorname{mia\_}) \rangle, \langle \operatorname{mig\_}, \operatorname{expr\text{-}open}(\operatorname{mip\_}, 0, \operatorname{miy\_}) \rangle \rangle \end{array}
\end{equation*}

\par\medskip\noindent\textbf{Definition} (\texttt{rel-sep-mem-elim}).\;For all $\operatorname{mit\_}$, $\operatorname{mihs\_}$ and $\operatorname{mic\_}$, 
\begin{equation*}
\begin{array}{@{}l@{}} \operatorname{rel\text{-}sep\text{-}mem\text{-}elim}(\operatorname{mit\_}, \operatorname{mihs\_}, \operatorname{mic\_}) \\ \Leftrightarrow\; \operatorname{mic\_} \in \operatorname{seq\text{-}set} \wedge \\ \exists \operatorname{mig\_}, \operatorname{mif\_}, \operatorname{mia\_}, \operatorname{mip\_}, \operatorname{miy\_}. \\ \quad \operatorname{mic\_} = \langle \operatorname{mig\_}, \operatorname{mif\_} \rangle \wedge \\ \quad \operatorname{enc\text{-}in}(\operatorname{miy\_}, \operatorname{enc\text{-}sep}(\operatorname{mia\_}, \operatorname{mip\_})) \in \operatorname{mig\_} \wedge \\ \quad \operatorname{mihs\_} \\ \quad =\; \langle \langle \operatorname{ctx\text{-}add}(\operatorname{ctx\text{-}add}(\operatorname{difference}(\operatorname{mig\_}, \{\operatorname{enc\text{-}in}(\operatorname{miy\_}, \operatorname{enc\text{-}sep}(\operatorname{mia\_}, \operatorname{mip\_}))\}), \operatorname{enc\text{-}in}(\operatorname{miy\_}, \operatorname{mia\_})), \operatorname{expr\text{-}open}(\operatorname{mip\_}, 0, \operatorname{miy\_})), \operatorname{mif\_} \rangle \rangle \end{array}
\end{equation*}

\par\medskip\noindent\textbf{Definition} (\texttt{rel-comp-mem-intro}).\;For all $\operatorname{mit\_}$, $\operatorname{mihs\_}$ and $\operatorname{mic\_}$, 
\begin{equation*}
\begin{array}{@{}l@{}} \operatorname{rel\text{-}comp\text{-}mem\text{-}intro}(\operatorname{mit\_}, \operatorname{mihs\_}, \operatorname{mic\_}) \\ \Leftrightarrow\; \operatorname{mic\_} \in \operatorname{seq\text{-}set} \wedge \\ \exists \operatorname{mig\_}, \operatorname{mip\_}, \operatorname{miy\_}. \\ \quad \operatorname{mic\_} = \langle \operatorname{mig\_}, \operatorname{enc\text{-}in}(\operatorname{miy\_}, \operatorname{enc\text{-}comp}(\operatorname{mip\_})) \rangle \wedge \\ \quad \operatorname{mihs\_} = \langle \langle \operatorname{mig\_}, \operatorname{enc\text{-}in}(\operatorname{miy\_}, \operatorname{enc\text{-}setc}) \rangle, \langle \operatorname{mig\_}, \operatorname{expr\text{-}open}(\operatorname{mip\_}, 0, \operatorname{miy\_}) \rangle \rangle \end{array}
\end{equation*}

\par\medskip\noindent\textbf{Definition} (\texttt{rel-comp-mem-elim}).\;For all $\operatorname{mit\_}$, $\operatorname{mihs\_}$ and $\operatorname{mic\_}$, 
\begin{equation*}
\begin{array}{@{}l@{}} \operatorname{rel\text{-}comp\text{-}mem\text{-}elim}(\operatorname{mit\_}, \operatorname{mihs\_}, \operatorname{mic\_}) \\ \Leftrightarrow\; \operatorname{mic\_} \in \operatorname{seq\text{-}set} \wedge \\ \exists \operatorname{mig\_}, \operatorname{mif\_}, \operatorname{mip\_}, \operatorname{miy\_}. \\ \quad \operatorname{mic\_} = \langle \operatorname{mig\_}, \operatorname{mif\_} \rangle \wedge \\ \quad \operatorname{enc\text{-}in}(\operatorname{miy\_}, \operatorname{enc\text{-}comp}(\operatorname{mip\_})) \in \operatorname{mig\_} \wedge \\ \quad \operatorname{mihs\_} \\ \quad =\; \langle \langle \operatorname{ctx\text{-}add}(\operatorname{ctx\text{-}add}(\operatorname{difference}(\operatorname{mig\_}, \{\operatorname{enc\text{-}in}(\operatorname{miy\_}, \operatorname{enc\text{-}comp}(\operatorname{mip\_}))\}), \operatorname{enc\text{-}in}(\operatorname{miy\_}, \operatorname{enc\text{-}setc})), \operatorname{expr\text{-}open}(\operatorname{mip\_}, 0, \operatorname{miy\_})), \operatorname{mif\_} \rangle \rangle \end{array}
\end{equation*}

\par\medskip\noindent\textbf{Definition} (\texttt{rel-big-union-sethood}).\;For all $\operatorname{mit\_}$, $\operatorname{mihs\_}$ and $\operatorname{mic\_}$, 
\begin{equation*}
\begin{array}{@{}l@{}} \operatorname{rel\text{-}big\text{-}union\text{-}sethood}(\operatorname{mit\_}, \operatorname{mihs\_}, \operatorname{mic\_}) \\ \Leftrightarrow\; \operatorname{mic\_} \in \operatorname{seq\text{-}set} \wedge \\ \exists \operatorname{mig\_}, \operatorname{mia\_}, \operatorname{mib\_}. \\ \quad \operatorname{mic\_} = \langle \operatorname{mig\_}, \operatorname{enc\text{-}in}(\operatorname{enc\text{-}bu}(\operatorname{mia\_}, \operatorname{mib\_}), \operatorname{enc\text{-}setc}) \rangle \wedge \\ \quad \operatorname{mihs\_} \\ \quad =\; \langle \langle \operatorname{mig\_}, \operatorname{enc\text{-}in}(\operatorname{mia\_}, \operatorname{enc\text{-}setc}) \rangle, \langle \operatorname{mig\_}, \operatorname{enc\text{-}all}(\operatorname{enc\text{-}imp}(\operatorname{enc\text{-}in}(\operatorname{ebvar}(0), \operatorname{mia\_}), \operatorname{enc\text{-}in}(\operatorname{mib\_}, \operatorname{enc\text{-}setc}))) \rangle \rangle \end{array}
\end{equation*}

\par\medskip\noindent\textbf{Definition} (\texttt{rel-big-union-mem-intro}).\;For all $\operatorname{mit\_}$, $\operatorname{mihs\_}$ and $\operatorname{mic\_}$, 
\begin{equation*}
\begin{array}{@{}l@{}} \operatorname{rel\text{-}big\text{-}union\text{-}mem\text{-}intro}(\operatorname{mit\_}, \operatorname{mihs\_}, \operatorname{mic\_}) \\ \Leftrightarrow\; \operatorname{mic\_} \in \operatorname{seq\text{-}set} \wedge \\ \exists \operatorname{mig\_}, \operatorname{mia\_}, \operatorname{mib\_}, \operatorname{mix\_}, \operatorname{miw\_}. \\ \quad \operatorname{mic\_} = \langle \operatorname{mig\_}, \operatorname{enc\text{-}in}(\operatorname{mix\_}, \operatorname{enc\text{-}bu}(\operatorname{mia\_}, \operatorname{mib\_})) \rangle \wedge \\ \quad \operatorname{mihs\_} = \langle \langle \operatorname{mig\_}, \operatorname{enc\text{-}in}(\operatorname{miw\_}, \operatorname{mia\_}) \rangle, \langle \operatorname{mig\_}, \operatorname{enc\text{-}in}(\operatorname{mix\_}, \operatorname{expr\text{-}open}(\operatorname{mib\_}, 0, \operatorname{miw\_})) \rangle \rangle \end{array}
\end{equation*}

\par\medskip\noindent\textbf{Definition} (\texttt{rel-big-union-mem-elim}).\;For all $\operatorname{mit\_}$, $\operatorname{mihs\_}$ and $\operatorname{mic\_}$, 
\begin{equation*}
\begin{array}{@{}l@{}} \operatorname{rel\text{-}big\text{-}union\text{-}mem\text{-}elim}(\operatorname{mit\_}, \operatorname{mihs\_}, \operatorname{mic\_}) \\ \Leftrightarrow\; \operatorname{mic\_} \in \operatorname{seq\text{-}set} \wedge \\ \exists \operatorname{mig\_}, \operatorname{mif\_}, \operatorname{mia\_}, \operatorname{mib\_}, \operatorname{mix\_}, \operatorname{mie\_}. \\ \quad \operatorname{mic\_} = \langle \operatorname{mig\_}, \operatorname{mif\_} \rangle \wedge \\ \quad \operatorname{enc\text{-}in}(\operatorname{mix\_}, \operatorname{enc\text{-}bu}(\operatorname{mia\_}, \operatorname{mib\_})) \in \operatorname{mig\_} \wedge \\ \quad \operatorname{mie\_} \in \mathbb{N} \wedge \\ \quad \neg (\operatorname{mie\_} \in (\operatorname{ctx\text{-}fv}(\operatorname{mig\_}) \cup \operatorname{expr\text{-}fv}(\operatorname{mif\_}))) \wedge \\ \quad \operatorname{mihs\_} \\ \quad =\; \langle \langle \operatorname{ctx\text{-}add}(\operatorname{ctx\text{-}add}(\operatorname{difference}(\operatorname{mig\_}, \{\operatorname{enc\text{-}in}(\operatorname{mix\_}, \operatorname{enc\text{-}bu}(\operatorname{mia\_}, \operatorname{mib\_}))\}), \operatorname{enc\text{-}in}(\operatorname{efvar}(\operatorname{mie\_}), \operatorname{mia\_})), \operatorname{enc\text{-}in}(\operatorname{mix\_}, \operatorname{expr\text{-}open}(\operatorname{mib\_}, 0, \operatorname{efvar}(\operatorname{mie\_})))), \operatorname{mif\_} \rangle \rangle \end{array}
\end{equation*}

\par\medskip\noindent\textbf{Definition} (\texttt{enc-nnc}).\; $(\operatorname{enc\text{-}nnc} \simeq \operatorname{econst}(3))$
\par\medskip\noindent\textbf{Definition} (\texttt{enc-ordc}).\; $(\operatorname{enc\text{-}ordc} \simeq \operatorname{econst}(5))$
\par\medskip\noindent\textbf{Definition} (\texttt{enc-succ}).\;$\operatorname{enc\text{-}succ}(\operatorname{mip\_}) :=$ $\operatorname{eapp}(30, \langle \operatorname{mip\_} \rangle)$
\par\medskip\noindent\textbf{Definition} (\texttt{enc-succo}).\;$\operatorname{enc\text{-}succo}(\operatorname{mip\_}) :=$ $\operatorname{eapp}(31, \langle \operatorname{mip\_} \rangle)$
\par\medskip\noindent\textbf{Definition} (\texttt{enc-ordlt}).\;$\operatorname{enc\text{-}ordlt}(\operatorname{mip\_}, \operatorname{miq\_}) :=$ $\operatorname{eapp}(10, \langle \operatorname{econst}(7), \operatorname{mip\_}, \operatorname{miq\_} \rangle)$
\par\medskip\noindent\textbf{Definition} (\texttt{enc-limit}).\;$\operatorname{enc\text{-}limit}(\operatorname{mip\_}) :=$ $\operatorname{eapp}(10, \langle \operatorname{econst}(9), \operatorname{mip\_} \rangle)$
\par\medskip\noindent\textbf{Definition} (\texttt{ax-nn-ind}).\;
\begin{equation*}
(\operatorname{ax\text{-}nn\text{-}ind} \simeq \operatorname{enc\text{-}all}(\operatorname{enc\text{-}imp}(\operatorname{enc\text{-}and}(\operatorname{enc\text{-}in}(\operatorname{enc\text{-}num}(0), \operatorname{ebvar}(0)), \operatorname{enc\text{-}all}(\operatorname{enc\text{-}imp}(\operatorname{enc\text{-}and}(\operatorname{enc\text{-}in}(\operatorname{ebvar}(0), \operatorname{enc\text{-}nnc}), \operatorname{enc\text{-}in}(\operatorname{ebvar}(0), \operatorname{ebvar}(1))), \operatorname{enc\text{-}in}(\operatorname{enc\text{-}succ}(\operatorname{ebvar}(0)), \operatorname{ebvar}(1))))), \operatorname{enc\text{-}all}(\operatorname{enc\text{-}imp}(\operatorname{enc\text{-}in}(\operatorname{ebvar}(0), \operatorname{enc\text{-}nnc}), \operatorname{enc\text{-}in}(\operatorname{ebvar}(0), \operatorname{ebvar}(1)))))))
\end{equation*}

\par\medskip\noindent\textbf{Definition} (\texttt{ax-tfi}).\;
\begin{equation*}
(\operatorname{ax\text{-}tfi} \simeq \operatorname{enc\text{-}all}(\operatorname{enc\text{-}imp}(\operatorname{enc\text{-}all}(\operatorname{enc\text{-}imp}(\operatorname{enc\text{-}and}(\operatorname{enc\text{-}in}(\operatorname{ebvar}(0), \operatorname{enc\text{-}ordc}), \operatorname{enc\text{-}all}(\operatorname{enc\text{-}imp}(\operatorname{enc\text{-}ordlt}(\operatorname{ebvar}(0), \operatorname{ebvar}(1)), \operatorname{enc\text{-}in}(\operatorname{ebvar}(0), \operatorname{ebvar}(2))))), \operatorname{enc\text{-}in}(\operatorname{ebvar}(0), \operatorname{ebvar}(1)))), \operatorname{enc\text{-}all}(\operatorname{enc\text{-}imp}(\operatorname{enc\text{-}in}(\operatorname{ebvar}(0), \operatorname{enc\text{-}ordc}), \operatorname{enc\text{-}in}(\operatorname{ebvar}(0), \operatorname{ebvar}(1)))))))
\end{equation*}

\par\medskip\noindent\textbf{Definition} (\texttt{ax-mem-set}).\;
\begin{equation*}
(\operatorname{ax\text{-}mem\text{-}set} \simeq \operatorname{enc\text{-}all}(\operatorname{enc\text{-}all}(\operatorname{enc\text{-}imp}(\operatorname{enc\text{-}in}(\operatorname{ebvar}(1), \operatorname{ebvar}(0)), \operatorname{enc\text{-}in}(\operatorname{ebvar}(1), \operatorname{enc\text{-}setc})))))
\end{equation*}

\par\medskip\noindent\textbf{Definition} (\texttt{rel-nn-induction}).\;For all $\operatorname{mit\_}$, $\operatorname{mihs\_}$ and $\operatorname{mic\_}$, 
\begin{equation*}
\begin{array}{@{}l@{}} \operatorname{rel\text{-}nn\text{-}induction}(\operatorname{mit\_}, \operatorname{mihs\_}, \operatorname{mic\_}) \\ \Leftrightarrow\; \operatorname{mic\_} \in \operatorname{seq\text{-}set} \wedge \\ \exists \operatorname{mig\_}, \operatorname{mip\_}, \operatorname{mis\_}, \operatorname{mifs\_}. \\ \quad \operatorname{mic\_} = \langle \operatorname{mig\_}, \operatorname{enc\text{-}all}(\operatorname{enc\text{-}imp}(\operatorname{enc\text{-}in}(\operatorname{ebvar}(0), \operatorname{enc\text{-}nnc}), \operatorname{mip\_})) \rangle \wedge \\ \quad \operatorname{mihs\_} = \langle \langle \operatorname{mig\_}, \operatorname{expr\text{-}open}(\operatorname{mip\_}, 0, \operatorname{enc\text{-}num}(0)) \rangle, \langle \operatorname{mig\_}, \operatorname{enc\text{-}all}(\operatorname{mis\_}) \rangle \rangle \wedge \\ \quad \operatorname{mifs\_} \in \mathbf{Set} \wedge \\ \quad \operatorname{mifs\_} \subseteq \mathbb{N} \wedge \\ \quad |\operatorname{mifs\_}| \in \mathbb{N} \wedge \\ \quad \forall \operatorname{miy\_} \in \mathbb{N}. \\ \quad \quad \neg (\operatorname{miy\_} \in \operatorname{mifs\_}) \Rightarrow  \\ \quad \quad \quad \operatorname{expr\text{-}open}(\operatorname{mis\_}, 0, \operatorname{efvar}(\operatorname{miy\_})) \\ \quad \quad \quad =\; \operatorname{enc\text{-}imp}(\operatorname{enc\text{-}in}(\operatorname{efvar}(\operatorname{miy\_}), \operatorname{enc\text{-}nnc}), \\ \quad \quad \quad \quad \operatorname{enc\text{-}imp}(\operatorname{expr\text{-}open}(\operatorname{mip\_}, 0, \operatorname{efvar}(\operatorname{miy\_})), \operatorname{expr\text{-}open}(\operatorname{mip\_}, 0, \operatorname{enc\text{-}succ}(\operatorname{efvar}(\operatorname{miy\_}))))) \end{array}
\end{equation*}

\par\medskip\noindent\textbf{Definition} (\texttt{rel-transfinite-induction}).\;For all $\operatorname{mit\_}$, $\operatorname{mihs\_}$ and $\operatorname{mic\_}$, 
\begin{equation*}
\begin{array}{@{}l@{}} \operatorname{rel\text{-}transfinite\text{-}induction}(\operatorname{mit\_}, \operatorname{mihs\_}, \operatorname{mic\_}) \\ \Leftrightarrow\; \operatorname{mic\_} \in \operatorname{seq\text{-}set} \wedge \\ \exists \operatorname{mig\_}, \operatorname{mip\_}. \\ \quad \operatorname{mic\_} = \langle \operatorname{mig\_}, \operatorname{enc\text{-}all}(\operatorname{enc\text{-}imp}(\operatorname{enc\text{-}in}(\operatorname{ebvar}(0), \operatorname{enc\text{-}ordc}), \operatorname{mip\_})) \rangle \wedge \\ \quad \operatorname{mihs\_} \\ \quad =\; \langle \langle \operatorname{mig\_}, \operatorname{enc\text{-}all}(\operatorname{enc\text{-}imp}(\operatorname{enc\text{-}and}(\operatorname{enc\text{-}in}(\operatorname{ebvar}(0), \operatorname{enc\text{-}ordc}), \operatorname{enc\text{-}all}(\operatorname{enc\text{-}imp}(\operatorname{enc\text{-}ordlt}(\operatorname{ebvar}(0), \operatorname{ebvar}(1)), \operatorname{mip\_}))), \operatorname{mip\_})) \rangle \rangle \end{array}
\end{equation*}

\par\medskip\noindent\textbf{Definition} (\texttt{rel-tfi-3cases}).\;For all $\operatorname{mit\_}$, $\operatorname{mihs\_}$ and $\operatorname{mic\_}$, 
\begin{equation*}
\begin{array}{@{}l@{}} \operatorname{rel\text{-}tfi\text{-}3cases}(\operatorname{mit\_}, \operatorname{mihs\_}, \operatorname{mic\_}) \\ \Leftrightarrow\; \operatorname{mic\_} \in \operatorname{seq\text{-}set} \wedge \\ \exists \operatorname{mig\_}, \operatorname{mip\_}, \operatorname{mis\_}, \operatorname{mifs\_}. \\ \quad \operatorname{mic\_} = \langle \operatorname{mig\_}, \operatorname{enc\text{-}all}(\operatorname{enc\text{-}imp}(\operatorname{enc\text{-}in}(\operatorname{ebvar}(0), \operatorname{enc\text{-}ordc}), \operatorname{mip\_})) \rangle \wedge \\ \quad \operatorname{mihs\_} \\ \quad =\; \langle \langle \operatorname{mig\_}, \operatorname{expr\text{-}open}(\operatorname{mip\_}, 0, \operatorname{enc\text{-}num}(0)) \rangle, \langle \operatorname{mig\_}, \operatorname{enc\text{-}all}(\operatorname{mis\_}) \rangle, \langle \operatorname{mig\_}, \operatorname{enc\text{-}all}(\operatorname{enc\text{-}imp}(\operatorname{enc\text{-}and}(\operatorname{enc\text{-}limit}(\operatorname{ebvar}(0)), \operatorname{enc\text{-}all}(\operatorname{enc\text{-}imp}(\operatorname{enc\text{-}ordlt}(\operatorname{ebvar}(0), \operatorname{ebvar}(1)), \operatorname{mip\_}))), \operatorname{mip\_})) \rangle \rangle \wedge \\ \quad \operatorname{mifs\_} \in \mathbf{Set} \wedge \\ \quad \operatorname{mifs\_} \subseteq \mathbb{N} \wedge \\ \quad |\operatorname{mifs\_}| \in \mathbb{N} \wedge \\ \quad \forall \operatorname{miy\_} \in \mathbb{N}. \\ \quad \quad \neg (\operatorname{miy\_} \in \operatorname{mifs\_}) \Rightarrow  \\ \quad \quad \quad \operatorname{expr\text{-}open}(\operatorname{mis\_}, 0, \operatorname{efvar}(\operatorname{miy\_})) \\ \quad \quad \quad =\; \operatorname{enc\text{-}imp}(\operatorname{enc\text{-}and}(\operatorname{enc\text{-}in}(\operatorname{efvar}(\operatorname{miy\_}), \operatorname{enc\text{-}ordc}), \operatorname{expr\text{-}open}(\operatorname{mip\_}, 0, \operatorname{efvar}(\operatorname{miy\_}))), \\ \quad \quad \quad \quad \operatorname{expr\text{-}open}(\operatorname{mip\_}, 0, \operatorname{enc\text{-}succo}(\operatorname{efvar}(\operatorname{miy\_})))) \end{array}
\end{equation*}

\par\medskip\noindent\textbf{Definition} (\texttt{rel-cartesian-intro}).\;For all $\operatorname{mit\_}$, $\operatorname{mihs\_}$ and $\operatorname{mic\_}$, 
\begin{equation*}
\begin{array}{@{}l@{}} \operatorname{rel\text{-}cartesian\text{-}intro}(\operatorname{mit\_}, \operatorname{mihs\_}, \operatorname{mic\_}) \\ \Leftrightarrow\; \operatorname{mic\_} \in \operatorname{seq\text{-}set} \wedge \\ \exists \operatorname{mig\_}, \operatorname{mie\_}, \operatorname{mia\_}. \\ \quad \operatorname{mic\_} = \langle \operatorname{mig\_}, \operatorname{enc\text{-}in}(\operatorname{enc\text{-}list}(\operatorname{mie\_}), \operatorname{enc\text{-}cart}(\operatorname{mia\_})) \rangle \wedge \\ \quad \operatorname{length}(\operatorname{mie\_}) = \operatorname{length}(\operatorname{mia\_}) \wedge \\ \quad 1 \leq \operatorname{length}(\operatorname{mia\_}) \wedge \\ \quad \operatorname{length}(\operatorname{mihs\_}) = \operatorname{length}(\operatorname{mie\_}) \wedge \\ \quad \forall \operatorname{mii\_} \in \mathbb{N}. \\ \quad \quad (1 \leq \operatorname{mii\_}) \Rightarrow  \\ \quad \quad (\operatorname{mii\_} \leq \operatorname{length}(\operatorname{mie\_})) \Rightarrow  \\ \quad \quad \quad \operatorname{nth}(\operatorname{mii\_}, \operatorname{mihs\_}) = \langle \operatorname{mig\_}, \operatorname{enc\text{-}in}(\operatorname{nth}(\operatorname{mii\_}, \operatorname{mie\_}), \operatorname{nth}(\operatorname{mii\_}, \operatorname{mia\_})) \rangle \end{array}
\end{equation*}

\par\medskip\noindent\textbf{Definition} (\texttt{rel-cartesian-elim}).\;For all $\operatorname{mit\_}$, $\operatorname{mihs\_}$ and $\operatorname{mic\_}$, 
\begin{equation*}
\begin{array}{@{}l@{}} \operatorname{rel\text{-}cartesian\text{-}elim}(\operatorname{mit\_}, \operatorname{mihs\_}, \operatorname{mic\_}) \\ \Leftrightarrow\; \operatorname{mic\_} \in \operatorname{seq\text{-}set} \wedge \\ \exists \operatorname{mig\_}, \operatorname{mif\_}, \operatorname{mie\_}, \operatorname{mia\_}, \operatorname{mik\_}. \\ \quad \operatorname{mic\_} = \langle \operatorname{mig\_}, \operatorname{mif\_} \rangle \wedge \\ \quad \operatorname{enc\text{-}in}(\operatorname{enc\text{-}list}(\operatorname{mie\_}), \operatorname{enc\text{-}cart}(\operatorname{mia\_})) \in \operatorname{mig\_} \wedge \\ \quad \operatorname{length}(\operatorname{mie\_}) = \operatorname{length}(\operatorname{mia\_}) \wedge \\ \quad \operatorname{mik\_} \in \mathbb{N} \wedge \\ \quad 1 \leq \operatorname{mik\_} \wedge \\ \quad \operatorname{mik\_} \leq \operatorname{length}(\operatorname{mie\_}) \wedge \\ \quad \operatorname{mihs\_} = \langle \langle \operatorname{ctx\text{-}add}(\operatorname{mig\_}, \operatorname{enc\text{-}in}(\operatorname{nth}(\operatorname{mik\_}, \operatorname{mie\_}), \operatorname{nth}(\operatorname{mik\_}, \operatorname{mia\_}))), \operatorname{mif\_} \rangle \rangle \end{array}
\end{equation*}

\par\medskip\noindent\textbf{Definition} (\texttt{idx-list}).\;$\operatorname{idx\text{-}list}(\operatorname{min\_}) :=$
\begin{equation*}
(\iota\, \operatorname{mir\_}.\;((\operatorname{mir\_} \in \operatorname{tuples}(\mathbf{Set})) \wedge ((\operatorname{length}(\operatorname{mir\_}) = \operatorname{min\_}) \wedge \forall \operatorname{mii\_} \in \mathbb{N}.\; ((1 \leq \operatorname{mii\_}) \Rightarrow ((\operatorname{mii\_} \leq \operatorname{min\_}) \Rightarrow (\operatorname{nth}(\operatorname{mii\_}, \operatorname{mir\_}) = \operatorname{ebvar}((\operatorname{min\_} - \operatorname{mii\_}))))))))
\end{equation*}

\par\medskip\noindent\textbf{Definition} (\texttt{enc-cch}).\;For all $\operatorname{mix\_}$, $\operatorname{mia\_}$ and $\operatorname{min\_}$, 
\begin{equation*}
(\operatorname{enc\text{-}cch}(\operatorname{mix\_}, \operatorname{mia\_}, \operatorname{min\_}, 0) \simeq \operatorname{enc\text{-}eq}(\operatorname{mix\_}, \operatorname{enc\text{-}list}(\operatorname{idx\text{-}list}(\operatorname{min\_}))))
\end{equation*}
\par
Let $\operatorname{mix\_}$, $\operatorname{mia\_}$ and $\operatorname{min\_}$ be arbitrary. Let $\operatorname{mik\_} \in \mathbb{N}$. Then
\begin{equation*}
(\operatorname{enc\text{-}cch}(\operatorname{mix\_}, \operatorname{mia\_}, \operatorname{min\_}, \operatorname{succ}(\operatorname{mik\_})) \simeq \operatorname{enc\text{-}ex}(\operatorname{enc\text{-}and}(\operatorname{enc\text{-}in}(\operatorname{ebvar}(0), \operatorname{nth}(-(\operatorname{min\_}, \operatorname{mik\_}), \operatorname{mia\_})), \operatorname{enc\text{-}cch}(\operatorname{mix\_}, \operatorname{mia\_}, \operatorname{min\_}, \operatorname{mik\_}))))
\end{equation*}

\par\medskip\noindent\textbf{Definition} (\texttt{j-cart}).\;$\operatorname{j\text{-}cart}(\operatorname{mit\_}) :=$
\begin{equation*}
\{\operatorname{mix\_} \in \operatorname{rw\text{-}univ} : \exists \operatorname{mig\_}, \operatorname{miz\_}, \operatorname{mia\_}.\; ((\operatorname{mix\_} = \langle 2, \operatorname{mig\_}, \operatorname{enc\text{-}in}(\operatorname{miz\_}, \operatorname{enc\text{-}cart}(\operatorname{mia\_})), \operatorname{enc\text{-}cch}(\operatorname{miz\_}, \operatorname{mia\_}, \operatorname{length}(\operatorname{mia\_}), \operatorname{length}(\operatorname{mia\_})) \rangle) \wedge ((\operatorname{mig\_} \in \operatorname{ctx\text{-}set}) \wedge (1 \leq \operatorname{length}(\operatorname{mia\_}))))\}
\end{equation*}

\par\medskip\noindent\textbf{Definition} (\texttt{rel-cartesian-decompose-1}).\;For all $\operatorname{mit\_}$, $\operatorname{mihs\_}$ and $\operatorname{mic\_}$, 
\begin{equation*}
\operatorname{rel\text{-}cartesian\text{-}decompose\text{-}1}(\operatorname{mit\_}, \operatorname{mihs\_}, \operatorname{mic\_}) \Leftrightarrow \operatorname{rel\text{-}rewrite\text{-}1}(\operatorname{j\text{-}cart}(\operatorname{mit\_}), \operatorname{mihs\_}, \operatorname{mic\_})
\end{equation*}

\par\medskip\noindent\textbf{Definition} (\texttt{enc-tcj}).\;For all $\operatorname{mia\_}$, $\operatorname{mib\_}$ and $\operatorname{min\_}$,  $(\operatorname{enc\text{-}tcj}(\operatorname{mia\_}, \operatorname{mib\_}, \operatorname{min\_}, 0) \simeq \operatorname{enc\text{-}truth})$\par
Let $\operatorname{mia\_}$, $\operatorname{mib\_}$ and $\operatorname{min\_}$ be arbitrary. Let $\operatorname{mik\_} \in \mathbb{N}$. Then
\begin{equation*}
(\operatorname{enc\text{-}tcj}(\operatorname{mia\_}, \operatorname{mib\_}, \operatorname{min\_}, \operatorname{succ}(\operatorname{mik\_})) \simeq \left(\text{if }(\operatorname{mik\_} = 0)\text{ then }\operatorname{enc\text{-}eq}(\operatorname{nth}(\operatorname{min\_}, \operatorname{mia\_}), \operatorname{nth}(\operatorname{min\_}, \operatorname{mib\_}))\text{ else }\operatorname{enc\text{-}and}(\operatorname{enc\text{-}eq}(\operatorname{nth}(-(\operatorname{min\_}, \operatorname{mik\_}), \operatorname{mia\_}), \operatorname{nth}(-(\operatorname{min\_}, \operatorname{mik\_}), \operatorname{mib\_})), \operatorname{enc\text{-}tcj}(\operatorname{mia\_}, \operatorname{mib\_}, \operatorname{min\_}, \operatorname{mik\_}))\right))
\end{equation*}

\par\medskip\noindent\textbf{Definition} (\texttt{j-teq}).\;$\operatorname{j\text{-}teq}(\operatorname{mit\_}) :=$
\begin{equation*}
\{\operatorname{mix\_} \in \operatorname{rw\text{-}univ} : \exists \operatorname{mig\_}, \operatorname{mia\_}, \operatorname{mib\_}.\; ((\operatorname{mix\_} = \langle 2, \operatorname{mig\_}, \operatorname{enc\text{-}eq}(\operatorname{enc\text{-}list}(\operatorname{mia\_}), \operatorname{enc\text{-}list}(\operatorname{mib\_})), \operatorname{enc\text{-}tcj}(\operatorname{mia\_}, \operatorname{mib\_}, \operatorname{length}(\operatorname{mia\_}), \operatorname{length}(\operatorname{mia\_})) \rangle) \wedge ((\operatorname{mig\_} \in \operatorname{ctx\text{-}set}) \wedge (\operatorname{length}(\operatorname{mia\_}) = \operatorname{length}(\operatorname{mib\_}))))\}
\end{equation*}

\par\medskip\noindent\textbf{Definition} (\texttt{rel-tuple-eq-decompose-1}).\;For all $\operatorname{mit\_}$, $\operatorname{mihs\_}$ and $\operatorname{mic\_}$, 
\begin{equation*}
\operatorname{rel\text{-}tuple\text{-}eq\text{-}decompose\text{-}1}(\operatorname{mit\_}, \operatorname{mihs\_}, \operatorname{mic\_}) \Leftrightarrow \operatorname{rel\text{-}rewrite\text{-}1}(\operatorname{j\text{-}teq}(\operatorname{mit\_}), \operatorname{mihs\_}, \operatorname{mic\_})
\end{equation*}

\par\medskip\noindent\textbf{Definition} (\texttt{def-comp-closed}).\;For every $\operatorname{mit\_}$, 
\begin{equation*}
\begin{array}{@{}l@{}} \operatorname{def\text{-}comp\text{-}closed}(\operatorname{mit\_}) \\ \Leftrightarrow\; \forall \operatorname{mig\_} \in \operatorname{ctx\text{-}set}, \operatorname{miq\_}. \\ \quad (\operatorname{enc\text{-}comp}(\operatorname{miq\_}) \in \mathit{clt}) \Rightarrow  \\ \quad \quad \langle \operatorname{mig\_}, \operatorname{enc\text{-}eq}(\operatorname{enc\text{-}comp}(\operatorname{miq\_}), \operatorname{enc\text{-}comp}(\operatorname{miq\_})) \rangle \in \operatorname{derivable}(\operatorname{mit\_}) \end{array}
\end{equation*}

\par\medskip\noindent\textbf{Definition} (\texttt{def-in-closed}).\;For every $\operatorname{mit\_}$, 
\begin{equation*}
\begin{array}{@{}l@{}} \operatorname{def\text{-}in\text{-}closed}(\operatorname{mit\_}) \\ \Leftrightarrow\; \forall \operatorname{mig\_}, \operatorname{miw\_}, \operatorname{mia\_}. \\ \quad (\langle \operatorname{mig\_}, \operatorname{enc\text{-}in}(\operatorname{miw\_}, \operatorname{mia\_}) \rangle \in \operatorname{derivable}(\operatorname{mit\_})) \Rightarrow  \\ \quad \quad \langle \operatorname{mig\_}, \operatorname{enc\text{-}eq}(\operatorname{miw\_}, \operatorname{miw\_}) \rangle \in \operatorname{derivable}(\operatorname{mit\_}) \end{array}
\end{equation*}

\par\medskip\noindent\textbf{Definition} (\texttt{ax-zero-nn}).\; $(\operatorname{ax\text{-}zero\text{-}nn} \simeq \operatorname{enc\text{-}in}(\operatorname{enc\text{-}num}(0), \operatorname{enc\text{-}nnc}))$
\par\medskip\noindent\textbf{Definition} (\texttt{ax-succ-nn}).\;
\begin{equation*}
(\operatorname{ax\text{-}succ\text{-}nn} \simeq \operatorname{enc\text{-}all}(\operatorname{enc\text{-}imp}(\operatorname{enc\text{-}in}(\operatorname{ebvar}(0), \operatorname{enc\text{-}nnc}), \operatorname{enc\text{-}in}(\operatorname{enc\text{-}succ}(\operatorname{ebvar}(0)), \operatorname{enc\text{-}nnc}))))
\end{equation*}

\par\medskip\noindent\textbf{Definition} (\texttt{ax-limit-iff}).\;
\begin{equation*}
(\operatorname{ax\text{-}limit\text{-}iff} \simeq \operatorname{enc\text{-}all}(\operatorname{enc\text{-}iff}(\operatorname{enc\text{-}limit}(\operatorname{ebvar}(0)), \operatorname{enc\text{-}and}(\operatorname{enc\text{-}in}(\operatorname{ebvar}(0), \operatorname{enc\text{-}ordc}), \operatorname{enc\text{-}and}(\operatorname{enc\text{-}not}(\operatorname{enc\text{-}eq}(\operatorname{ebvar}(0), \operatorname{enc\text{-}num}(0))), \operatorname{enc\text{-}not}(\operatorname{enc\text{-}ex}(\operatorname{enc\text{-}and}(\operatorname{enc\text{-}in}(\operatorname{ebvar}(0), \operatorname{enc\text{-}ordc}), \operatorname{enc\text{-}eq}(\operatorname{ebvar}(1), \operatorname{enc\text{-}succo}(\operatorname{ebvar}(0)))))))))))
\end{equation*}

\par\medskip\noindent\textbf{Definition} (\texttt{ax-ord-succ-above}).\;
\begin{equation*}
(\operatorname{ax\text{-}ord\text{-}succ\text{-}above} \simeq \operatorname{enc\text{-}all}(\operatorname{enc\text{-}imp}(\operatorname{enc\text{-}in}(\operatorname{ebvar}(0), \operatorname{enc\text{-}ordc}), \operatorname{enc\text{-}ordlt}(\operatorname{ebvar}(0), \operatorname{enc\text{-}succo}(\operatorname{ebvar}(0))))))
\end{equation*}


\vnbbookfile{structure-library/meta-def}{structure-library/meta-def}
\noindent{\small meta-def.scm -- K2 PHASE 5: the clauses of the kernel's definedness certificate pi--defined? (primitive-inferences.scm 1164-1235 and its helpers, tree f7e2246, 2026-10-12), each TRANSCRIBED as a rule DEF-k over the encoding: hypotheses the context facts the clause reads, conclusion [g, t = t]. ...}\par
\par\medskip\noindent\textbf{Definition} (\texttt{def-strict-in-l}).\;For all $\operatorname{mdt\_}$, $\operatorname{mdhs\_}$ and $\operatorname{mdc\_}$, 
\begin{equation*}
\begin{array}{@{}l@{}} \operatorname{def\text{-}strict\text{-}in\text{-}l}(\operatorname{mdt\_}, \operatorname{mdhs\_}, \operatorname{mdc\_}) \\ \Leftrightarrow\; \operatorname{mdc\_} \in \operatorname{seq\text{-}set} \wedge \\ \exists \operatorname{mdg\_}, \operatorname{mdw\_}, \operatorname{mda\_}. \\ \quad (\operatorname{mdc\_} = \langle \operatorname{mdg\_}, \operatorname{enc\text{-}eq}(\operatorname{mdw\_}, \operatorname{mdw\_}) \rangle) \wedge (\operatorname{mdhs\_} = \langle \langle \operatorname{mdg\_}, \operatorname{enc\text{-}in}(\operatorname{mdw\_}, \operatorname{mda\_}) \rangle \rangle) \end{array}
\end{equation*}

\par\medskip\noindent\textbf{Definition} (\texttt{def-strict-eq-l}).\;For all $\operatorname{mdt\_}$, $\operatorname{mdhs\_}$ and $\operatorname{mdc\_}$, 
\begin{equation*}
\begin{array}{@{}l@{}} \operatorname{def\text{-}strict\text{-}eq\text{-}l}(\operatorname{mdt\_}, \operatorname{mdhs\_}, \operatorname{mdc\_}) \\ \Leftrightarrow\; \operatorname{mdc\_} \in \operatorname{seq\text{-}set} \wedge \\ \exists \operatorname{mdg\_}, \operatorname{mds\_}, \operatorname{mdu\_}. \\ \quad (\operatorname{mdc\_} = \langle \operatorname{mdg\_}, \operatorname{enc\text{-}eq}(\operatorname{mds\_}, \operatorname{mds\_}) \rangle) \wedge (\operatorname{mdhs\_} = \langle \langle \operatorname{mdg\_}, \operatorname{enc\text{-}eq}(\operatorname{mds\_}, \operatorname{mdu\_}) \rangle \rangle) \end{array}
\end{equation*}

\par\medskip\noindent\textbf{Definition} (\texttt{def-strict-eq-r}).\;For all $\operatorname{mdt\_}$, $\operatorname{mdhs\_}$ and $\operatorname{mdc\_}$, 
\begin{equation*}
\begin{array}{@{}l@{}} \operatorname{def\text{-}strict\text{-}eq\text{-}r}(\operatorname{mdt\_}, \operatorname{mdhs\_}, \operatorname{mdc\_}) \\ \Leftrightarrow\; \operatorname{mdc\_} \in \operatorname{seq\text{-}set} \wedge \\ \exists \operatorname{mdg\_}, \operatorname{mds\_}, \operatorname{mdu\_}. \\ \quad (\operatorname{mdc\_} = \langle \operatorname{mdg\_}, \operatorname{enc\text{-}eq}(\operatorname{mdu\_}, \operatorname{mdu\_}) \rangle) \wedge (\operatorname{mdhs\_} = \langle \langle \operatorname{mdg\_}, \operatorname{enc\text{-}eq}(\operatorname{mds\_}, \operatorname{mdu\_}) \rangle \rangle) \end{array}
\end{equation*}

\par\medskip\noindent\textbf{Definition} (\texttt{def-strict-arg}).\;For all $\operatorname{mdt\_}$, $\operatorname{mdhs\_}$ and $\operatorname{mdc\_}$, 
\begin{equation*}
\begin{array}{@{}l@{}} \operatorname{def\text{-}strict\text{-}arg}(\operatorname{mdt\_}, \operatorname{mdhs\_}, \operatorname{mdc\_}) \\ \Leftrightarrow\; \operatorname{mdc\_} \in \operatorname{seq\text{-}set} \wedge \\ \exists \operatorname{mdg\_}, \operatorname{mdh\_}, \operatorname{mdl\_}, \operatorname{mdi\_}. \\ \quad \operatorname{mdc\_} = \langle \operatorname{mdg\_}, \operatorname{enc\text{-}eq}(\operatorname{nth}(\operatorname{mdi\_}, \operatorname{mdl\_}), \operatorname{nth}(\operatorname{mdi\_}, \operatorname{mdl\_})) \rangle \wedge \\ \quad \operatorname{strict\text{-}head}(\operatorname{mdh\_}) \wedge \\ \quad \operatorname{mdi\_} \in \mathbb{N} \wedge \\ \quad 1 \leq \operatorname{mdi\_} \wedge \\ \quad \operatorname{mdi\_} \leq \operatorname{length}(\operatorname{mdl\_}) \wedge \\ \quad \operatorname{mdhs\_} = \langle \langle \operatorname{mdg\_}, \operatorname{enc\text{-}eq}(\operatorname{eapp}(\operatorname{mdh\_}, \operatorname{mdl\_}), \operatorname{eapp}(\operatorname{mdh\_}, \operatorname{mdl\_})) \rangle \rangle \end{array}
\end{equation*}

\par\medskip\noindent\textbf{Definition} (\texttt{def-strict-hyp}).\;For all $\operatorname{mdt\_}$, $\operatorname{mdhs\_}$ and $\operatorname{mdc\_}$, 
\begin{equation*}
\begin{array}{@{}l@{}} \operatorname{def\text{-}strict\text{-}hyp}(\operatorname{mdt\_}, \operatorname{mdhs\_}, \operatorname{mdc\_}) \\ \Leftrightarrow\; \operatorname{mdc\_} \in \operatorname{seq\text{-}set} \wedge \\ \exists \operatorname{mdg\_}, \operatorname{mdf\_}, \operatorname{md7h2\_}, \operatorname{mdw\_}. \\ \quad \operatorname{mdc\_} = \langle \operatorname{mdg\_}, \operatorname{enc\text{-}eq}(\operatorname{mdw\_}, \operatorname{mdw\_}) \rangle \wedge \\ \quad \operatorname{mdhs\_} = \langle \langle \operatorname{mdg\_}, \operatorname{mdf\_} \rangle \rangle \wedge \\ \quad \operatorname{md7h2\_} \in \operatorname{true\text{-}reach}(\operatorname{mdf\_}) \wedge \\ \quad \operatorname{strict\text{-}atom}(\operatorname{md7h2\_}, \operatorname{mdw\_}) \end{array}
\end{equation*}

\par\medskip\noindent\textbf{Definition} (\texttt{def-strict-pred}).\;For all $\operatorname{mdt\_}$, $\operatorname{mdhs\_}$ and $\operatorname{mdc\_}$, 
\begin{equation*}
\begin{array}{@{}l@{}} \operatorname{def\text{-}strict\text{-}pred}(\operatorname{mdt\_}, \operatorname{mdhs\_}, \operatorname{mdc\_}) \\ \Leftrightarrow\; \operatorname{mdc\_} \in \operatorname{seq\text{-}set} \wedge \\ \exists \operatorname{mdg\_}, \operatorname{mdl\_}, \operatorname{mdk\_}, \operatorname{mdi\_}. \\ \quad \operatorname{mdc\_} = \langle \operatorname{mdg\_}, \operatorname{enc\text{-}eq}(\operatorname{nth}(\operatorname{mdi\_}, \operatorname{mdl\_}), \operatorname{nth}(\operatorname{mdi\_}, \operatorname{mdl\_})) \rangle \wedge \\ \quad \operatorname{nth}(1, \operatorname{mdl\_}) = \operatorname{econst}(\operatorname{mdk\_}) \wedge \\ \quad \operatorname{pred\text{-}code}(\operatorname{mdk\_}) \wedge \\ \quad \operatorname{mdi\_} \in \mathbb{N} \wedge \\ \quad 2 \leq \operatorname{mdi\_} \wedge \\ \quad \operatorname{mdi\_} \leq \operatorname{length}(\operatorname{mdl\_}) \wedge \\ \quad \operatorname{mdhs\_} = \langle \langle \operatorname{mdg\_}, \operatorname{eapp}(10, \operatorname{mdl\_}) \rangle \rangle \end{array}
\end{equation*}

\par\medskip\noindent\textbf{Definition} (\texttt{def-struct-hyp}).\;For all $\operatorname{mdt\_}$, $\operatorname{mdhs\_}$ and $\operatorname{mdc\_}$, 
\begin{equation*}
\begin{array}{@{}l@{}} \operatorname{def\text{-}struct\text{-}hyp}(\operatorname{mdt\_}, \operatorname{mdhs\_}, \operatorname{mdc\_}) \\ \Leftrightarrow\; \operatorname{mdc\_} \in \operatorname{seq\text{-}set} \wedge \\ \exists \operatorname{mdg\_}, \operatorname{mdk\_}, \operatorname{mds\_}. \\ \quad \operatorname{mdc\_} = \langle \operatorname{mdg\_}, \operatorname{enc\text{-}eq}(\operatorname{mds\_}, \operatorname{mds\_}) \rangle \wedge \\ \quad \operatorname{pred\text{-}code}(\operatorname{mdk\_}) \wedge \\ \quad \operatorname{mdhs\_} = \langle \langle \operatorname{mdg\_}, \operatorname{enc\text{-}pred1}(\operatorname{econst}(\operatorname{mdk\_}), \operatorname{mds\_}) \rangle \rangle \end{array}
\end{equation*}

\par\medskip\noindent\textbf{Definition} (\texttt{def-list}).\;For all $\operatorname{mdt\_}$, $\operatorname{mdhs\_}$ and $\operatorname{mdc\_}$, 
\begin{equation*}
\begin{array}{@{}l@{}} \operatorname{def\text{-}list}(\operatorname{mdt\_}, \operatorname{mdhs\_}, \operatorname{mdc\_}) \\ \Leftrightarrow\; \operatorname{mdc\_} \in \operatorname{seq\text{-}set} \wedge \\ \exists \operatorname{mdg\_}, \operatorname{mdl\_}. \\ \quad \operatorname{mdc\_} = \langle \operatorname{mdg\_}, \operatorname{enc\text{-}eq}(\operatorname{enc\text{-}list}(\operatorname{mdl\_}), \operatorname{enc\text{-}list}(\operatorname{mdl\_})) \rangle \wedge \\ \quad \operatorname{length}(\operatorname{mdhs\_}) = \operatorname{length}(\operatorname{mdl\_}) \wedge \\ \quad \forall \operatorname{mdi\_} \in \mathbb{N}. \\ \quad \quad (1 \leq \operatorname{mdi\_}) \Rightarrow  \\ \quad \quad (\operatorname{mdi\_} \leq \operatorname{length}(\operatorname{mdl\_})) \Rightarrow  \\ \quad \quad \quad \operatorname{nth}(\operatorname{mdi\_}, \operatorname{mdhs\_}) = \langle \operatorname{mdg\_}, \operatorname{enc\text{-}in}(\operatorname{nth}(\operatorname{mdi\_}, \operatorname{mdl\_}), \operatorname{enc\text{-}setc}) \rangle \end{array}
\end{equation*}

\par\medskip\noindent\textbf{Definition} (\texttt{def-total}).\;For all $\operatorname{mdt\_}$, $\operatorname{mdhs\_}$ and $\operatorname{mdc\_}$, 
\begin{equation*}
\begin{array}{@{}l@{}} \operatorname{def\text{-}total}(\operatorname{mdt\_}, \operatorname{mdhs\_}, \operatorname{mdc\_}) \\ \Leftrightarrow\; \operatorname{mdc\_} \in \operatorname{seq\text{-}set} \wedge \\ \exists \operatorname{mdg\_}, \operatorname{mdh\_}, \operatorname{mdl\_}. \\ \quad \operatorname{mdc\_} = \langle \operatorname{mdg\_}, \operatorname{enc\text{-}eq}(\operatorname{eapp}(\operatorname{mdh\_}, \operatorname{mdl\_}), \operatorname{eapp}(\operatorname{mdh\_}, \operatorname{mdl\_})) \rangle \wedge \\ \quad \operatorname{total\text{-}head}(\operatorname{mdh\_}) \wedge \\ \quad \operatorname{length}(\operatorname{mdhs\_}) = \operatorname{length}(\operatorname{mdl\_}) \wedge \\ \quad \forall \operatorname{mdi\_} \in \mathbb{N}. \\ \quad \quad (1 \leq \operatorname{mdi\_}) \Rightarrow  \\ \quad \quad (\operatorname{mdi\_} \leq \operatorname{length}(\operatorname{mdl\_})) \Rightarrow  \\ \quad \quad \quad \operatorname{nth}(\operatorname{mdi\_}, \operatorname{mdhs\_}) = \langle \operatorname{mdg\_}, \operatorname{enc\text{-}eq}(\operatorname{nth}(\operatorname{mdi\_}, \operatorname{mdl\_}), \operatorname{nth}(\operatorname{mdi\_}, \operatorname{mdl\_})) \rangle \end{array}
\end{equation*}

\par\medskip\noindent\textbf{Definition} (\texttt{def-binder}).\;For all $\operatorname{mdt\_}$, $\operatorname{mdhs\_}$ and $\operatorname{mdc\_}$, 
\begin{equation*}
\begin{array}{@{}l@{}} \operatorname{def\text{-}binder}(\operatorname{mdt\_}, \operatorname{mdhs\_}, \operatorname{mdc\_}) \\ \Leftrightarrow\; \operatorname{mdc\_} \in \operatorname{seq\text{-}set} \wedge \\ \exists \operatorname{mdg\_}, \operatorname{mdh\_}, \operatorname{mda\_}, \operatorname{mdb\_}. \\ \quad \operatorname{mdc\_} = \langle \operatorname{mdg\_}, \operatorname{enc\text{-}eq}(\operatorname{ebind}(\operatorname{mdh\_}, 1, \langle \operatorname{mda\_} \rangle, \operatorname{mdb\_}), \operatorname{ebind}(\operatorname{mdh\_}, 1, \langle \operatorname{mda\_} \rangle, \operatorname{mdb\_})) \rangle \wedge \\ \quad (\operatorname{mdh\_} = 22) \vee ((\operatorname{mdh\_} = 24) \vee (\operatorname{mdh\_} = 26)) \wedge \\ \quad \operatorname{mdhs\_} = \langle \langle \operatorname{mdg\_}, \operatorname{enc\text{-}eq}(\operatorname{mda\_}, \operatorname{mda\_}) \rangle \rangle \end{array}
\end{equation*}

\par\medskip\noindent\textbf{Definition} (\texttt{def-comp}).\;For all $\operatorname{mdt\_}$, $\operatorname{mdhs\_}$ and $\operatorname{mdc\_}$, 
\begin{equation*}
\begin{array}{@{}l@{}} \operatorname{def\text{-}comp}(\operatorname{mdt\_}, \operatorname{mdhs\_}, \operatorname{mdc\_}) \\ \Leftrightarrow\; \operatorname{mdc\_} \in \operatorname{seq\text{-}set} \wedge \\ \exists \operatorname{mdg\_}, \operatorname{mdp\_}. \\ \quad (\operatorname{mdc\_} = \langle \operatorname{mdg\_}, \operatorname{enc\text{-}eq}(\operatorname{enc\text{-}comp}(\operatorname{mdp\_}), \operatorname{enc\text{-}comp}(\operatorname{mdp\_})) \rangle) \wedge (\operatorname{mdhs\_} = \langle  \rangle) \end{array}
\end{equation*}

\par\medskip\noindent\textbf{Definition} (\texttt{def-nth-lit}).\;For all $\operatorname{mdt\_}$, $\operatorname{mdhs\_}$ and $\operatorname{mdc\_}$, 
\begin{equation*}
\begin{array}{@{}l@{}} \operatorname{def\text{-}nth\text{-}lit}(\operatorname{mdt\_}, \operatorname{mdhs\_}, \operatorname{mdc\_}) \\ \Leftrightarrow\; \operatorname{mdc\_} \in \operatorname{seq\text{-}set} \wedge \\ \exists \operatorname{mdg\_}, \operatorname{mdl\_}, \operatorname{mdk\_}. \\ \quad \operatorname{mdc\_} \\ \quad =\; \langle \operatorname{mdg\_}, \operatorname{enc\text{-}eq}(\operatorname{enc\text{-}nth}(\operatorname{enc\text{-}num}(\operatorname{mdk\_}), \operatorname{enc\text{-}list}(\operatorname{mdl\_})), \operatorname{enc\text{-}nth}(\operatorname{enc\text{-}num}(\operatorname{mdk\_}), \operatorname{enc\text{-}list}(\operatorname{mdl\_}))) \rangle \wedge \\ \quad \operatorname{mdk\_} \in \mathbb{N} \wedge \\ \quad 1 \leq \operatorname{mdk\_} \wedge \\ \quad \operatorname{mdk\_} \leq \operatorname{length}(\operatorname{mdl\_}) \wedge \\ \quad \operatorname{mdhs\_} = \langle \langle \operatorname{mdg\_}, \operatorname{enc\text{-}eq}(\operatorname{nth}(\operatorname{mdk\_}, \operatorname{mdl\_}), \operatorname{nth}(\operatorname{mdk\_}, \operatorname{mdl\_})) \rangle \rangle \end{array}
\end{equation*}

\par\medskip\noindent\textbf{Definition} (\texttt{def-nth-list}).\;For all $\operatorname{mdt\_}$, $\operatorname{mdhs\_}$ and $\operatorname{mdc\_}$, 
\begin{equation*}
\begin{array}{@{}l@{}} \operatorname{def\text{-}nth\text{-}list}(\operatorname{mdt\_}, \operatorname{mdhs\_}, \operatorname{mdc\_}) \\ \Leftrightarrow\; \operatorname{mdc\_} \in \operatorname{seq\text{-}set} \wedge \\ \exists \operatorname{mdg\_}, \operatorname{mdl\_}, \operatorname{mdk\_}. \\ \quad \operatorname{mdc\_} \\ \quad =\; \langle \operatorname{mdg\_}, \operatorname{enc\text{-}eq}(\operatorname{enc\text{-}nth}(\operatorname{enc\text{-}num}(\operatorname{mdk\_}), \operatorname{enc\text{-}list}(\operatorname{mdl\_})), \operatorname{enc\text{-}nth}(\operatorname{enc\text{-}num}(\operatorname{mdk\_}), \operatorname{enc\text{-}list}(\operatorname{mdl\_}))) \rangle \wedge \\ \quad \operatorname{mdk\_} \in \mathbb{N} \wedge \\ \quad 1 \leq \operatorname{mdk\_} \wedge \\ \quad \operatorname{mdk\_} \leq \operatorname{length}(\operatorname{mdl\_}) \wedge \\ \quad \operatorname{length}(\operatorname{mdhs\_}) = \operatorname{length}(\operatorname{mdl\_}) \wedge \\ \quad \forall \operatorname{mdi\_} \in \mathbb{N}. \\ \quad \quad (1 \leq \operatorname{mdi\_}) \Rightarrow  \\ \quad \quad (\operatorname{mdi\_} \leq \operatorname{length}(\operatorname{mdl\_})) \Rightarrow  \\ \quad \quad \quad \operatorname{nth}(\operatorname{mdi\_}, \operatorname{mdhs\_}) = \langle \operatorname{mdg\_}, \operatorname{enc\text{-}in}(\operatorname{nth}(\operatorname{mdi\_}, \operatorname{mdl\_}), \operatorname{enc\text{-}setc}) \rangle \end{array}
\end{equation*}

\par\medskip\noindent\textbf{Definition} (\texttt{def-app}).\;For all $\operatorname{mdt\_}$, $\operatorname{mdhs\_}$ and $\operatorname{mdc\_}$, 
\begin{equation*}
\begin{array}{@{}l@{}} \operatorname{def\text{-}app}(\operatorname{mdt\_}, \operatorname{mdhs\_}, \operatorname{mdc\_}) \\ \Leftrightarrow\; \operatorname{mdc\_} \in \operatorname{seq\text{-}set} \wedge \\ \exists \operatorname{mdg\_}, \operatorname{mdf\_}, \operatorname{mdd\_}, \operatorname{mda\_}. \\ \quad \operatorname{mdc\_} = \langle \operatorname{mdg\_}, \operatorname{enc\text{-}eq}(\operatorname{enc\text{-}apply}(\operatorname{mdf\_}, \operatorname{mda\_}), \operatorname{enc\text{-}apply}(\operatorname{mdf\_}, \operatorname{mda\_})) \rangle \wedge \\ \quad \operatorname{mdhs\_} = \langle \langle \operatorname{mdg\_}, \operatorname{enc\text{-}in}(\operatorname{mdf\_}, \operatorname{enc\text{-}fun1}(\operatorname{mdd\_})) \rangle, \langle \operatorname{mdg\_}, \operatorname{enc\text{-}in}(\operatorname{mda\_}, \operatorname{mdd\_}) \rangle \rangle \end{array}
\end{equation*}

\par\medskip\noindent\textbf{Definition} (\texttt{def-pred-unfold}).\;For all $\operatorname{mdt\_}$, $\operatorname{mdhs\_}$ and $\operatorname{mdc\_}$, 
\begin{equation*}
\begin{array}{@{}l@{}} \operatorname{def\text{-}pred\text{-}unfold}(\operatorname{mdt\_}, \operatorname{mdhs\_}, \operatorname{mdc\_}) \\ \Leftrightarrow\; \operatorname{mdc\_} \in \operatorname{seq\text{-}set} \wedge \\ \exists \operatorname{mdg\_}, \operatorname{mdk\_}, \operatorname{mds\_}, \operatorname{mdb\_}. \\ \quad \operatorname{mdc\_} = \langle \operatorname{mdg\_}, \operatorname{expr\text{-}open}(\operatorname{mdb\_}, 0, \operatorname{mds\_}) \rangle \wedge \\ \quad \operatorname{pred\text{-}code}(\operatorname{mdk\_}) \wedge \\ \quad \operatorname{mdhs\_} \\ \quad =\; \langle \langle \operatorname{mdg\_}, \operatorname{enc\text{-}pred1}(\operatorname{econst}(\operatorname{mdk\_}), \operatorname{mds\_}) \rangle, \langle \operatorname{mdg\_}, \operatorname{enc\text{-}all}(\operatorname{enc\text{-}iff}(\operatorname{enc\text{-}pred1}(\operatorname{econst}(\operatorname{mdk\_}), \operatorname{ebvar}(0)), \operatorname{mdb\_})) \rangle \rangle \end{array}
\end{equation*}

\par\medskip\noindent\textbf{Definition} (\texttt{def-pred-hyp}).\;For all $\operatorname{mdt\_}$, $\operatorname{mdhs\_}$ and $\operatorname{mdc\_}$, 
\begin{equation*}
\begin{array}{@{}l@{}} \operatorname{def\text{-}pred\text{-}hyp}(\operatorname{mdt\_}, \operatorname{mdhs\_}, \operatorname{mdc\_}) \\ \Leftrightarrow\; \operatorname{mdc\_} \in \operatorname{seq\text{-}set} \wedge \\ \exists \operatorname{mdg\_}, \operatorname{mdk\_}, \operatorname{mds\_}, \operatorname{mdb\_}, \operatorname{md10f\_}, \operatorname{mdw\_}. \\ \quad \operatorname{mdc\_} = \langle \operatorname{mdg\_}, \operatorname{enc\text{-}eq}(\operatorname{mdw\_}, \operatorname{mdw\_}) \rangle \wedge \\ \quad \operatorname{pred\text{-}code}(\operatorname{mdk\_}) \wedge \\ \quad \operatorname{mdhs\_} \\ \quad =\; \langle \langle \operatorname{mdg\_}, \operatorname{enc\text{-}pred1}(\operatorname{econst}(\operatorname{mdk\_}), \operatorname{mds\_}) \rangle, \langle \operatorname{mdg\_}, \operatorname{enc\text{-}all}(\operatorname{enc\text{-}iff}(\operatorname{enc\text{-}pred1}(\operatorname{econst}(\operatorname{mdk\_}), \operatorname{ebvar}(0)), \operatorname{mdb\_})) \rangle \rangle \wedge \\ \quad \operatorname{md10f\_} \in \operatorname{true\text{-}reach}(\operatorname{expr\text{-}open}(\operatorname{mdb\_}, 0, \operatorname{mds\_})) \wedge \\ \quad \operatorname{strict\text{-}atom}(\operatorname{md10f\_}, \operatorname{mdw\_}) \end{array}
\end{equation*}

\par\medskip\noindent\textbf{Definition} (\texttt{def-strict-in-r}).\;For all $\operatorname{mdt\_}$, $\operatorname{mdhs\_}$ and $\operatorname{mdc\_}$, 
\begin{equation*}
\begin{array}{@{}l@{}} \operatorname{def\text{-}strict\text{-}in\text{-}r}(\operatorname{mdt\_}, \operatorname{mdhs\_}, \operatorname{mdc\_}) \\ \Leftrightarrow\; \operatorname{mdc\_} \in \operatorname{seq\text{-}set} \wedge \\ \exists \operatorname{mdg\_}, \operatorname{mdw\_}, \operatorname{mda\_}. \\ \quad (\operatorname{mdc\_} = \langle \operatorname{mdg\_}, \operatorname{enc\text{-}eq}(\operatorname{mda\_}, \operatorname{mda\_}) \rangle) \wedge (\operatorname{mdhs\_} = \langle \langle \operatorname{mdg\_}, \operatorname{enc\text{-}in}(\operatorname{mdw\_}, \operatorname{mda\_}) \rangle \rangle) \end{array}
\end{equation*}

\par\medskip\noindent\textbf{Definition} (\texttt{def-strict-dom}).\;For all $\operatorname{mdt\_}$, $\operatorname{mdhs\_}$ and $\operatorname{mdc\_}$, 
\begin{equation*}
\begin{array}{@{}l@{}} \operatorname{def\text{-}strict\text{-}dom}(\operatorname{mdt\_}, \operatorname{mdhs\_}, \operatorname{mdc\_}) \\ \Leftrightarrow\; \operatorname{mdc\_} \in \operatorname{seq\text{-}set} \wedge \\ \exists \operatorname{mdg\_}, \operatorname{mdh\_}, \operatorname{mda\_}, \operatorname{mdb\_}. \\ \quad \operatorname{mdc\_} = \langle \operatorname{mdg\_}, \operatorname{enc\text{-}eq}(\operatorname{mda\_}, \operatorname{mda\_}) \rangle \wedge \\ \quad (\operatorname{mdh\_} = 22) \vee ((\operatorname{mdh\_} = 24) \vee (\operatorname{mdh\_} = 26)) \wedge \\ \quad \operatorname{mdhs\_} = \langle \langle \operatorname{mdg\_}, \operatorname{enc\text{-}eq}(\operatorname{ebind}(\operatorname{mdh\_}, 1, \langle \operatorname{mda\_} \rangle, \operatorname{mdb\_}), \operatorname{ebind}(\operatorname{mdh\_}, 1, \langle \operatorname{mda\_} \rangle, \operatorname{mdb\_})) \rangle \rangle \end{array}
\end{equation*}

\par\medskip\noindent\textbf{Definition} (\texttt{strict-atom}).\;For all $\operatorname{md7af\_}$ and $\operatorname{md7at\_}$, 
\begin{equation*}
\begin{array}{@{}l@{}} \operatorname{strict\text{-}atom}(\operatorname{md7af\_}, \operatorname{md7at\_}) \\ \Leftrightarrow\; \exists \operatorname{md7ax\_}, \operatorname{md7ay\_}. \\ \quad (\operatorname{md7af\_} = \operatorname{enc\text{-}in}(\operatorname{md7ax\_}, \operatorname{md7ay\_})) \vee (\operatorname{md7af\_} = \operatorname{enc\text{-}eq}(\operatorname{md7ax\_}, \operatorname{md7ay\_})) \wedge \\ \quad (\operatorname{md7at\_} \in \operatorname{strict\text{-}reach}(\operatorname{md7ax\_})) \vee (\operatorname{md7at\_} \in \operatorname{strict\text{-}reach}(\operatorname{md7ay\_})) \vee \\ \exists \operatorname{md7pl\_}, \operatorname{md7pk\_}, \operatorname{md7pi\_}. \\ \quad \operatorname{md7af\_} = \operatorname{eapp}(10, \operatorname{md7pl\_}) \wedge \\ \quad \operatorname{nth}(1, \operatorname{md7pl\_}) = \operatorname{econst}(\operatorname{md7pk\_}) \wedge \\ \quad \operatorname{pred\text{-}code}(\operatorname{md7pk\_}) \wedge \\ \quad \operatorname{md7pi\_} \in \mathbb{N} \wedge \\ \quad 2 \leq \operatorname{md7pi\_} \wedge \\ \quad \operatorname{md7pi\_} \leq \operatorname{length}(\operatorname{md7pl\_}) \wedge \\ \quad \operatorname{md7at\_} \in \operatorname{strict\text{-}reach}(\operatorname{nth}(\operatorname{md7pi\_}, \operatorname{md7pl\_})) \end{array}
\end{equation*}

\par\medskip\noindent\textbf{Definition} (\texttt{enc-le}).\;$\operatorname{enc\text{-}le}(\operatorname{md10x\_}, \operatorname{md10y\_}) :=$ $\operatorname{eapp}(10, \langle \operatorname{econst}(11), \operatorname{md10x\_}, \operatorname{md10y\_} \rangle)$
\par\medskip\noindent\textbf{Definition} (\texttt{enc-lt}).\;$\operatorname{enc\text{-}lt}(\operatorname{md10x\_}, \operatorname{md10y\_}) :=$ $\operatorname{eapp}(10, \langle \operatorname{econst}(13), \operatorname{md10x\_}, \operatorname{md10y\_} \rangle)$
\par\medskip\noindent\textbf{Definition} (\texttt{true-reach}).\;$\operatorname{true\text{-}reach}(\operatorname{md7j\_}) :=$
\begin{equation*}
\{\operatorname{md7zz\_} \in \operatorname{expr\text{-}set} : \forall \operatorname{md7ww\_} \in \operatorname{power}(\operatorname{expr\text{-}set}).\; ((\operatorname{md7j\_} \in \operatorname{md7ww\_}) \Rightarrow (\operatorname{true\text{-}closed}(\operatorname{md7ww\_}) \Rightarrow (\operatorname{md7zz\_} \in \operatorname{md7ww\_})))\}
\end{equation*}

\par\medskip\noindent\textbf{Definition} (\texttt{strict-step}).\;For all $\operatorname{md7s\_}$ and $\operatorname{md7t\_}$, 
\begin{equation*}
\begin{array}{@{}l@{}} \operatorname{strict\text{-}step}(\operatorname{md7s\_}, \operatorname{md7t\_}) \\ \Leftrightarrow\; \exists \operatorname{md7h\_}, \operatorname{md7l\_}, \operatorname{md7i\_}. \\ \quad \operatorname{md7s\_} = \operatorname{eapp}(\operatorname{md7h\_}, \operatorname{md7l\_}) \wedge \\ \quad \operatorname{strict\text{-}head}(\operatorname{md7h\_}) \wedge \\ \quad \operatorname{md7i\_} \in \mathbb{N} \wedge \\ \quad 1 \leq \operatorname{md7i\_} \wedge \\ \quad \operatorname{md7i\_} \leq \operatorname{length}(\operatorname{md7l\_}) \wedge \\ \quad \operatorname{md7t\_} = \operatorname{nth}(\operatorname{md7i\_}, \operatorname{md7l\_}) \vee \\ \exists \operatorname{md7k\_}, \operatorname{md7a\_}, \operatorname{md7b\_}. \\ \quad (\operatorname{md7k\_} = 22) \vee ((\operatorname{md7k\_} = 24) \vee (\operatorname{md7k\_} = 26)) \wedge \\ \quad \operatorname{md7s\_} = \operatorname{ebind}(\operatorname{md7k\_}, 1, \langle \operatorname{md7a\_} \rangle, \operatorname{md7b\_}) \wedge \\ \quad \operatorname{md7t\_} = \operatorname{md7a\_} \end{array}
\end{equation*}

\par\medskip\noindent\textbf{Definition} (\texttt{strict-reach}).\;$\operatorname{strict\text{-}reach}(\operatorname{md7e\_}) :=$
\begin{equation*}
\{\operatorname{md7z\_} \in \operatorname{expr\text{-}set} : \forall \operatorname{md7w\_} \in \operatorname{power}(\operatorname{expr\text{-}set}).\; ((\operatorname{md7e\_} \in \operatorname{md7w\_}) \Rightarrow (\operatorname{strict\text{-}closed}(\operatorname{md7w\_}) \Rightarrow (\operatorname{md7z\_} \in \operatorname{md7w\_})))\}
\end{equation*}

\par\medskip\noindent\textbf{Definition} (\texttt{enc-fun1}).\;$\operatorname{enc\text{-}fun1}(\operatorname{mdp\_}) :=$ $\operatorname{eapp}(19, \langle \operatorname{mdp\_} \rangle)$
\par\medskip\noindent\textbf{Definition} (\texttt{adm-def-strict-in-l}).\;For every $\operatorname{mdt\_}$, 
\begin{equation*}
\begin{array}{@{}l@{}} \operatorname{adm\text{-}def\text{-}strict\text{-}in\text{-}l}(\operatorname{mdt\_}) \\ \Leftrightarrow\; \forall \operatorname{mdhs\_}, \operatorname{mdc\_}. \\ \quad \operatorname{def\text{-}strict\text{-}in\text{-}l}(\operatorname{mdt\_}, \operatorname{mdhs\_}, \operatorname{mdc\_}) \Rightarrow  \\ \quad \operatorname{hyps\text{-}in}(\operatorname{mdhs\_}, \operatorname{derivable}(\operatorname{mdt\_})) \Rightarrow (\operatorname{mdc\_} \in \operatorname{derivable}(\operatorname{mdt\_})) \end{array}
\end{equation*}

\par\medskip\noindent\textbf{Definition} (\texttt{adm-def-comp}).\;For every $\operatorname{mdt\_}$, 
\begin{equation*}
\begin{array}{@{}l@{}} \operatorname{adm\text{-}def\text{-}comp}(\operatorname{mdt\_}) \\ \Leftrightarrow\; \forall \operatorname{mdhs\_}, \operatorname{mdc\_}. \\ \quad \operatorname{def\text{-}comp}(\operatorname{mdt\_}, \operatorname{mdhs\_}, \operatorname{mdc\_}) \Rightarrow  \\ \quad \operatorname{hyps\text{-}in}(\operatorname{mdhs\_}, \operatorname{derivable}(\operatorname{mdt\_})) \Rightarrow (\operatorname{mdc\_} \in \operatorname{derivable}(\operatorname{mdt\_})) \end{array}
\end{equation*}

\par\medskip\noindent\textbf{Definition} (\texttt{adm-def-strict-in-r}).\;For every $\operatorname{mdt\_}$, 
\begin{equation*}
\begin{array}{@{}l@{}} \operatorname{adm\text{-}def\text{-}strict\text{-}in\text{-}r}(\operatorname{mdt\_}) \\ \Leftrightarrow\; \forall \operatorname{mdhs\_}, \operatorname{mdc\_}. \\ \quad \operatorname{def\text{-}strict\text{-}in\text{-}r}(\operatorname{mdt\_}, \operatorname{mdhs\_}, \operatorname{mdc\_}) \Rightarrow  \\ \quad \operatorname{hyps\text{-}in}(\operatorname{mdhs\_}, \operatorname{derivable}(\operatorname{mdt\_})) \Rightarrow (\operatorname{mdc\_} \in \operatorname{derivable}(\operatorname{mdt\_})) \end{array}
\end{equation*}

\par\medskip\noindent\textbf{Definition} (\texttt{adm-def-strict-arg}).\;For every $\operatorname{mdt\_}$, 
\begin{equation*}
\begin{array}{@{}l@{}} \operatorname{adm\text{-}def\text{-}strict\text{-}arg}(\operatorname{mdt\_}) \\ \Leftrightarrow\; \forall \operatorname{mdhs\_}, \operatorname{mdc\_}. \\ \quad \operatorname{def\text{-}strict\text{-}arg}(\operatorname{mdt\_}, \operatorname{mdhs\_}, \operatorname{mdc\_}) \Rightarrow  \\ \quad \operatorname{hyps\text{-}in}(\operatorname{mdhs\_}, \operatorname{derivable}(\operatorname{mdt\_})) \Rightarrow (\operatorname{mdc\_} \in \operatorname{derivable}(\operatorname{mdt\_})) \end{array}
\end{equation*}

\par\medskip\noindent\textbf{Definition} (\texttt{adm-def-strict-dom}).\;For every $\operatorname{mdt\_}$, 
\begin{equation*}
\begin{array}{@{}l@{}} \operatorname{adm\text{-}def\text{-}strict\text{-}dom}(\operatorname{mdt\_}) \\ \Leftrightarrow\; \forall \operatorname{mdhs\_}, \operatorname{mdc\_}. \\ \quad \operatorname{def\text{-}strict\text{-}dom}(\operatorname{mdt\_}, \operatorname{mdhs\_}, \operatorname{mdc\_}) \Rightarrow  \\ \quad \operatorname{hyps\text{-}in}(\operatorname{mdhs\_}, \operatorname{derivable}(\operatorname{mdt\_})) \Rightarrow (\operatorname{mdc\_} \in \operatorname{derivable}(\operatorname{mdt\_})) \end{array}
\end{equation*}

\par\medskip\noindent\textbf{Definition} (\texttt{ax-fun-app}).\;
\begin{equation*}
(\operatorname{ax\text{-}fun\text{-}app} \simeq \operatorname{enc\text{-}all}(\operatorname{enc\text{-}all}(\operatorname{enc\text{-}all}(\operatorname{enc\text{-}imp}(\operatorname{enc\text{-}in}(\operatorname{ebvar}(1), \operatorname{enc\text{-}fun1}(\operatorname{ebvar}(2))), \operatorname{enc\text{-}iff}(\operatorname{enc\text{-}eq}(\operatorname{enc\text{-}apply}(\operatorname{ebvar}(1), \operatorname{ebvar}(0)), \operatorname{enc\text{-}apply}(\operatorname{ebvar}(1), \operatorname{ebvar}(0))), \operatorname{enc\text{-}in}(\operatorname{ebvar}(0), \operatorname{ebvar}(2))))))))
\end{equation*}

\par\medskip\noindent\textbf{Definition} (\texttt{strict-closed}).\;For every $\operatorname{md7x\_}$, 
\begin{equation*}
\begin{array}{@{}l@{}} \operatorname{strict\text{-}closed}(\operatorname{md7x\_}) \\ \Leftrightarrow\; \forall \operatorname{md7p\_} \in \operatorname{md7x\_}, \operatorname{md7q\_}. \\ \quad \operatorname{strict\text{-}step}(\operatorname{md7p\_}, \operatorname{md7q\_}) \Rightarrow (\operatorname{md7q\_} \in \operatorname{md7x\_}) \end{array}
\end{equation*}

\par\medskip\noindent\textbf{Definition} (\texttt{true-step}).\;For all $\operatorname{md7f\_}$ and $\operatorname{md7g\_}$, 
\begin{equation*}
\begin{array}{@{}l@{}} \operatorname{true\text{-}step}(\operatorname{md7f\_}, \operatorname{md7g\_}) \\ \Leftrightarrow\; \exists \operatorname{md7c\_}, \operatorname{md7d\_}. \\ \quad (\operatorname{md7f\_} = \operatorname{enc\text{-}and}(\operatorname{md7c\_}, \operatorname{md7d\_})) \wedge ((\operatorname{md7g\_} = \operatorname{md7c\_}) \vee (\operatorname{md7g\_} = \operatorname{md7d\_})) \vee \\ \exists \operatorname{md7y\_}, \operatorname{md7o\_}, \operatorname{md7r\_}. \\ \quad ((\operatorname{md7f\_} = \operatorname{enc\text{-}in}(\operatorname{md7y\_}, \operatorname{enc\text{-}sep}(\operatorname{md7o\_}, \operatorname{md7r\_}))) \wedge (\operatorname{md7g\_} = \operatorname{expr\text{-}open}(\operatorname{md7r\_}, 0, \operatorname{md7y\_}))) \end{array}
\end{equation*}

\par\medskip\noindent\textbf{Definition} (\texttt{true-closed}).\;For every $\operatorname{md7v\_}$, 
\begin{equation*}
\begin{array}{@{}l@{}} \operatorname{true\text{-}closed}(\operatorname{md7v\_}) \\ \Leftrightarrow\; \forall \operatorname{md7m\_} \in \operatorname{md7v\_}, \operatorname{md7n\_}. \\ \quad \operatorname{true\text{-}step}(\operatorname{md7m\_}, \operatorname{md7n\_}) \Rightarrow (\operatorname{md7n\_} \in \operatorname{md7v\_}) \end{array}
\end{equation*}

\par\medskip\noindent\textbf{Definition} (\texttt{adm-def-strict-hyp}).\;For every $\operatorname{mdt\_}$, 
\begin{equation*}
\begin{array}{@{}l@{}} \operatorname{adm\text{-}def\text{-}strict\text{-}hyp}(\operatorname{mdt\_}) \\ \Leftrightarrow\; \forall \operatorname{mdhs\_}, \operatorname{mdc\_}. \\ \quad \operatorname{def\text{-}strict\text{-}hyp}(\operatorname{mdt\_}, \operatorname{mdhs\_}, \operatorname{mdc\_}) \Rightarrow  \\ \quad \operatorname{hyps\text{-}in}(\operatorname{mdhs\_}, \operatorname{derivable}(\operatorname{mdt\_})) \Rightarrow (\operatorname{mdc\_} \in \operatorname{derivable}(\operatorname{mdt\_})) \end{array}
\end{equation*}

\par\medskip\noindent\textbf{Definition} (\texttt{adm-def-nth-list}).\;For every $\operatorname{mdt\_}$, 
\begin{equation*}
\begin{array}{@{}l@{}} \operatorname{adm\text{-}def\text{-}nth\text{-}list}(\operatorname{mdt\_}) \\ \Leftrightarrow\; \forall \operatorname{mdhs\_}, \operatorname{mdc\_}. \\ \quad \operatorname{def\text{-}nth\text{-}list}(\operatorname{mdt\_}, \operatorname{mdhs\_}, \operatorname{mdc\_}) \Rightarrow  \\ \quad \operatorname{hyps\text{-}in}(\operatorname{mdhs\_}, \operatorname{derivable}(\operatorname{mdt\_})) \Rightarrow (\operatorname{mdc\_} \in \operatorname{derivable}(\operatorname{mdt\_})) \end{array}
\end{equation*}

\par\medskip\noindent\textbf{Definition} (\texttt{adm-def-strict-pred}).\;For every $\operatorname{mdt\_}$, 
\begin{equation*}
\begin{array}{@{}l@{}} \operatorname{adm\text{-}def\text{-}strict\text{-}pred}(\operatorname{mdt\_}) \\ \Leftrightarrow\; \forall \operatorname{mdhs\_}, \operatorname{mdc\_}. \\ \quad \operatorname{def\text{-}strict\text{-}pred}(\operatorname{mdt\_}, \operatorname{mdhs\_}, \operatorname{mdc\_}) \Rightarrow  \\ \quad \operatorname{hyps\text{-}in}(\operatorname{mdhs\_}, \operatorname{derivable}(\operatorname{mdt\_})) \Rightarrow (\operatorname{mdc\_} \in \operatorname{derivable}(\operatorname{mdt\_})) \end{array}
\end{equation*}

\par\medskip\noindent\textbf{Definition} (\texttt{enc-pred1}).\;$\operatorname{enc\text{-}pred1}(\operatorname{md10n\_}, \operatorname{md10p\_}) :=$ $\operatorname{eapp}(10, \langle \operatorname{md10n\_}, \operatorname{md10p\_} \rangle)$
\par\medskip\noindent\textbf{Definition} (\texttt{adm-def-struct-hyp}).\;For every $\operatorname{mdt\_}$, 
\begin{equation*}
\begin{array}{@{}l@{}} \operatorname{adm\text{-}def\text{-}struct\text{-}hyp}(\operatorname{mdt\_}) \\ \Leftrightarrow\; \forall \operatorname{mdhs\_}, \operatorname{mdc\_}. \\ \quad \operatorname{def\text{-}struct\text{-}hyp}(\operatorname{mdt\_}, \operatorname{mdhs\_}, \operatorname{mdc\_}) \Rightarrow  \\ \quad \operatorname{hyps\text{-}in}(\operatorname{mdhs\_}, \operatorname{derivable}(\operatorname{mdt\_})) \Rightarrow (\operatorname{mdc\_} \in \operatorname{derivable}(\operatorname{mdt\_})) \end{array}
\end{equation*}

\par\medskip\noindent\textbf{Definition} (\texttt{adm-def-pred-unfold}).\;For every $\operatorname{mdt\_}$, 
\begin{equation*}
\begin{array}{@{}l@{}} \operatorname{adm\text{-}def\text{-}pred\text{-}unfold}(\operatorname{mdt\_}) \\ \Leftrightarrow\; \forall \operatorname{mdhs\_}, \operatorname{mdc\_}. \\ \quad \operatorname{def\text{-}pred\text{-}unfold}(\operatorname{mdt\_}, \operatorname{mdhs\_}, \operatorname{mdc\_}) \Rightarrow  \\ \quad \operatorname{hyps\text{-}in}(\operatorname{mdhs\_}, \operatorname{derivable}(\operatorname{mdt\_})) \Rightarrow (\operatorname{mdc\_} \in \operatorname{derivable}(\operatorname{mdt\_})) \end{array}
\end{equation*}

\par\medskip\noindent\textbf{Definition} (\texttt{adm-def-pred-hyp}).\;For every $\operatorname{mdt\_}$, 
\begin{equation*}
\begin{array}{@{}l@{}} \operatorname{adm\text{-}def\text{-}pred\text{-}hyp}(\operatorname{mdt\_}) \\ \Leftrightarrow\; \forall \operatorname{mdhs\_}, \operatorname{mdc\_}. \\ \quad \operatorname{def\text{-}pred\text{-}hyp}(\operatorname{mdt\_}, \operatorname{mdhs\_}, \operatorname{mdc\_}) \Rightarrow  \\ \quad \operatorname{hyps\text{-}in}(\operatorname{mdhs\_}, \operatorname{derivable}(\operatorname{mdt\_})) \Rightarrow (\operatorname{mdc\_} \in \operatorname{derivable}(\operatorname{mdt\_})) \end{array}
\end{equation*}


\vnbbookfile{theorem-library/meta-ind-laws}{theorem-library/meta-ind-laws}
\noindent{\small meta-ind-laws.scm -- K2 PHASE 4: the admissibility of the kernel's set-formation rules and of the three induction rules over the core, the typing machinery an instance at a compound term needs, and (section J) the one phase-5 derivation that needs the OPEN calculator (docs/meta-k2-admissibility-2026-10-08.md, section ...}\par
\vnbprop{hyps-in-2-elim}{\vnbstmt{\begin{array}{@{}l@{}} \forall a, b, s. \\ \quad ((\langle a, b \rangle = \langle a, b \rangle) \wedge \operatorname{hyps\text{-}in}(\langle a, b \rangle, s)) \Rightarrow ((a \in s) \wedge (b \in s)) \end{array}}}
\vnbprop{adm-sep-sethood}{\vnbstmt{\begin{array}{@{}l@{}} \forall t, s, c. \\ \quad (\operatorname{rel\text{-}sep\text{-}sethood}(t, s, c) \wedge \operatorname{hyps\text{-}in}(s, \operatorname{derivable}(t))) \Rightarrow (c \in \operatorname{derivable}(t)) \end{array}}}
\vnbprop{adm-big-union-sethood}{\vnbstmt{\begin{array}{@{}l@{}} \forall t, s, c. \\ \quad (\operatorname{rel\text{-}big\text{-}union\text{-}sethood}(t, s, c) \wedge \operatorname{hyps\text{-}in}(s, \operatorname{derivable}(t))) \Rightarrow  \\ \quad \quad c \in \operatorname{derivable}(t) \end{array}}}
\vnbprop{adm-sep-mem-intro}{\vnbstmt{\begin{array}{@{}l@{}} \forall t, s, c. \\ \quad (\operatorname{rel\text{-}sep\text{-}mem\text{-}intro}(t, s, c) \wedge \operatorname{hyps\text{-}in}(s, \operatorname{derivable}(t))) \Rightarrow  \\ \quad \quad c \in \operatorname{derivable}(t) \end{array}}}
\vnbprop{adm-comp-mem-intro}{\vnbstmt{\begin{array}{@{}l@{}} \forall t, s, c. \\ \quad (\operatorname{rel\text{-}comp\text{-}mem\text{-}intro}(t, s, c) \wedge \operatorname{hyps\text{-}in}(s, \operatorname{derivable}(t))) \Rightarrow  \\ \quad \quad c \in \operatorname{derivable}(t) \end{array}}}
\vnbprop{enc-sep-unfold}{\vnbstmt{\forall a, p.\; (\operatorname{enc\text{-}sep}(a, p) \simeq \operatorname{ebind}(22, 1, \langle a \rangle, p))}}
\vnbprop{enc-comp-unfold}{\vnbstmt{\forall p.\; (\operatorname{enc\text{-}comp}(p) \simeq \operatorname{ebind}(23, 1, \langle  \rangle, p))}}
\vnbprop{enc-bu-unfold}{\vnbstmt{\forall a, p.\; (\operatorname{enc\text{-}bu}(a, p) \simeq \operatorname{ebind}(24, 1, \langle a \rangle, p))}}
\vnbprop{adm-sep-mem-elim}{\vnbstmt{\begin{array}{@{}l@{}} \forall t, s, c. \\ \quad (\operatorname{rel\text{-}sep\text{-}mem\text{-}elim}(t, s, c) \wedge \operatorname{hyps\text{-}in}(s, \operatorname{derivable}(t))) \Rightarrow (c \in \operatorname{derivable}(t)) \end{array}}}
\vnbprop{adm-comp-mem-elim}{\vnbstmt{\begin{array}{@{}l@{}} \forall t, s, c. \\ \quad (\operatorname{rel\text{-}comp\text{-}mem\text{-}elim}(t, s, c) \wedge \operatorname{hyps\text{-}in}(s, \operatorname{derivable}(t))) \Rightarrow  \\ \quad \quad c \in \operatorname{derivable}(t) \end{array}}}
\vnbprop{wks-single}{\vnbstmt{\forall e \in \operatorname{expr\text{-}set}.\; ((1 \in \operatorname{expr\text{-}wks}(e)) \Rightarrow (\operatorname{expr\text{-}wks}(e) = \{1\}))}}
\vnbprop{lcd-mono}{\vnbstmt{\begin{array}{@{}l@{}} \forall e \in \operatorname{expr\text{-}set}, d, a. \\ \quad ((d \in \mathbb{N}) \wedge ((a \in \mathbb{N}) \wedge ((d \leq a) \wedge (d \in \operatorname{expr\text{-}lcd}(e))))) \Rightarrow (a \in \operatorname{expr\text{-}lcd}(e)) \end{array}}}
\vnbprop{lcd-open-closes}{\vnbstmt{\begin{array}{@{}l@{}} \forall u, e \in \operatorname{expr\text{-}set}, x \in \mathbb{N}. \\ \quad (((x + 1) \in \operatorname{expr\text{-}lcd}(e)) \wedge (x \in \operatorname{expr\text{-}lcd}(u))) \Rightarrow  \\ \quad \quad x \in \operatorname{expr\text{-}lcd}(\operatorname{expr\text{-}open}(e, x, u)) \end{array}}}
\vnbprop{enc-succ-unfold}{\vnbstmt{\forall p.\; (\operatorname{enc\text{-}succ}(p) \simeq \operatorname{eapp}(30, \langle p \rangle))}}
\vnbprop{enc-succo-unfold}{\vnbstmt{\forall p.\; (\operatorname{enc\text{-}succo}(p) \simeq \operatorname{eapp}(31, \langle p \rangle))}}
\vnbprop{enc-ordlt-unfold}{\vnbstmt{\forall p, q.\; (\operatorname{enc\text{-}ordlt}(p, q) \simeq \operatorname{eapp}(10, \langle \operatorname{econst}(7), p, q \rangle))}}
\vnbprop{enc-limit-unfold}{\vnbstmt{\forall p.\; (\operatorname{enc\text{-}limit}(p) \simeq \operatorname{eapp}(10, \langle \operatorname{econst}(9), p \rangle))}}
\vnbprop{enc-num-unfold}{\vnbstmt{\forall k \in \mathbb{N}.\; (\operatorname{enc\text{-}num}(k) = \operatorname{econst}((2 \cdot k)))}}
\vnbprop{ebind-args-stage}{\vnbstmt{\begin{array}{@{}l@{}} \forall n \in \mathbb{N}, h, c, a, b. \\ \quad (\operatorname{ebind}(h, c, a, b) \in \operatorname{expr\text{-}stage}(n)) \Rightarrow (a \in \operatorname{tuples}(\operatorname{expr\text{-}set})) \end{array}}}
\vnbprop{ebind-args-expr}{\vnbstmt{\begin{array}{@{}l@{}} \forall h, c, a, b. \\ \quad (\operatorname{ebind}(h, c, a, b) \in \operatorname{expr\text{-}set}) \Rightarrow (a \in \operatorname{tuples}(\operatorname{expr\text{-}set})) \end{array}}}
\vnbprop{comp-of-body-clt}{\vnbstmt{\forall q, p.\; ((\operatorname{enc\text{-}all}(\operatorname{enc\text{-}imp}(q, p)) \in \mathit{clf}) \Rightarrow (\operatorname{enc\text{-}comp}(p) \in \mathit{clt}))}}
\vnbprop{open-app1}{\vnbstmt{\begin{array}{@{}l@{}} \forall u, k, h, a_{1}. \\ \quad ((u \in \operatorname{expr\text{-}set}) \wedge ((k \in \mathbb{N}) \wedge ((h \in \mathbb{N}) \wedge (a_{1} \in \operatorname{expr\text{-}set})))) \Rightarrow  \\ \quad \quad \operatorname{expr\text{-}open}(\operatorname{eapp}(h, \langle a_{1} \rangle), k, u) = \operatorname{eapp}(h, \langle \operatorname{expr\text{-}open}(a_{1}, k, u) \rangle) \end{array}}}
\vnbprop{open-app3}{\vnbstmt{\begin{array}{@{}l@{}} \forall u, k, h, a_{1}, a_{2}, a_{3}. \\ \quad ((u \in \operatorname{expr\text{-}set}) \wedge ((k \in \mathbb{N}) \wedge ((h \in \mathbb{N}) \wedge ((a_{1} \in \operatorname{expr\text{-}set}) \wedge ((a_{2} \in \operatorname{expr\text{-}set}) \wedge (a_{3} \in \operatorname{expr\text{-}set})))))) \Rightarrow  \\ \quad \quad \operatorname{expr\text{-}open}(\operatorname{eapp}(h, \langle a_{1}, a_{2}, a_{3} \rangle), k, u) \\ \quad \quad =\; \operatorname{eapp}(h, \langle \operatorname{expr\text{-}open}(a_{1}, k, u), \operatorname{expr\text{-}open}(a_{2}, k, u), \operatorname{expr\text{-}open}(a_{3}, k, u) \rangle) \end{array}}}
\vnbprop{maplk-nil}{\vnbstmt{\forall u, k.\; (((u \in \operatorname{expr\text{-}set}) \wedge (k \in \mathbb{N})) \Rightarrow (\operatorname{maplk}(\operatorname{open\text{-}fun}(u), \langle  \rangle, k) = \langle  \rangle))}}
\vnbprop{der-comp-out}{\vnbstmt{\begin{array}{@{}l@{}} \forall a, h, q, t. \\ \quad (\langle h, \operatorname{enc\text{-}in}(t, \operatorname{enc\text{-}comp}(q)) \rangle \in \operatorname{derivable}(a)) \Rightarrow  \\ \quad \quad \langle h, \operatorname{expr\text{-}open}(q, 0, t) \rangle \in \operatorname{derivable}(a) \end{array}}}
\vnbprop{der-comp-in}{\vnbstmt{\begin{array}{@{}l@{}} \forall a, h, q, t. \\ \quad \operatorname{enc\text{-}comp}(q) \in \mathit{clt} \wedge \\ \quad \langle h, \operatorname{enc\text{-}in}(t, \operatorname{enc\text{-}setc}) \rangle \in \operatorname{derivable}(a) \wedge \\ \quad \langle h, \operatorname{expr\text{-}open}(q, 0, t) \rangle \in \operatorname{derivable}(a) \Rightarrow \\ \quad \quad \langle h, \operatorname{enc\text{-}in}(t, \operatorname{enc\text{-}comp}(q)) \rangle \in \operatorname{derivable}(a) \end{array}}}
\vnbprop{der-name-set}{\vnbstmt{\begin{array}{@{}l@{}} \forall t, h, n, m. \\ \quad ((\operatorname{ax\text{-}mem\text{-}set} \in t) \wedge ((n \in \mathbb{N}) \wedge ((m \in \mathbb{N}) \wedge (\langle h, \operatorname{enc\text{-}in}(\operatorname{efvar}(n), \operatorname{econst}(m)) \rangle \in \operatorname{derivable}(t))))) \Rightarrow  \\ \quad \quad \langle h, \operatorname{enc\text{-}in}(\operatorname{efvar}(n), \operatorname{enc\text{-}setc}) \rangle \in \operatorname{derivable}(t) \end{array}}}
\vnbprop{adm-big-union-mem-intro}{\vnbstmt{\begin{array}{@{}l@{}} \forall t, s, c. \\ \quad (\operatorname{def\text{-}in\text{-}closed}(t) \wedge (\operatorname{rel\text{-}big\text{-}union\text{-}mem\text{-}intro}(t, s, c) \wedge \operatorname{hyps\text{-}in}(s, \operatorname{derivable}(t)))) \Rightarrow  \\ \quad \quad c \in \operatorname{derivable}(t) \end{array}}}
\vnbprop{clf-open-inst}{\vnbstmt{\forall b, t.\; (((\operatorname{enc\text{-}all}(b) \in \mathit{clf}) \wedge (t \in \mathit{clt})) \Rightarrow (\operatorname{expr\text{-}open}(b, 0, t) \in \mathit{clf}))}}
\vnbprop{clf-open-inst-ex}{\vnbstmt{\forall b, t.\; (((\operatorname{enc\text{-}ex}(b) \in \mathit{clf}) \wedge (t \in \mathit{clt})) \Rightarrow (\operatorname{expr\text{-}open}(b, 0, t) \in \mathit{clf}))}}
\vnbprop{adm-transfinite-induction}{\vnbstmt{\begin{array}{@{}l@{}} \forall t, s, c. \\ \quad t \subseteq \mathit{clf} \wedge \\ \quad \operatorname{ax\text{-}tfi} \in t \wedge \\ \quad \operatorname{ax\text{-}mem\text{-}set} \in t \wedge \\ \quad \operatorname{def\text{-}comp\text{-}closed}(t) \wedge \\ \quad \operatorname{rel\text{-}transfinite\text{-}induction}(t, s, c) \wedge \\ \quad \operatorname{hyps\text{-}in}(s, \operatorname{derivable}(t)) \Rightarrow \\ \quad \quad c \in \operatorname{derivable}(t) \end{array}}}
\vnbprop{der-mem-set}{\vnbstmt{\begin{array}{@{}l@{}} \forall t, h, x, m. \\ \quad t \subseteq \mathit{clf} \wedge \\ \quad \operatorname{ax\text{-}mem\text{-}set} \in t \wedge \\ \quad m \in \mathbb{N} \wedge \\ \quad \langle h, \operatorname{enc\text{-}in}(x, \operatorname{econst}(m)) \rangle \in \operatorname{derivable}(t) \wedge \\ \quad \langle h, \operatorname{enc\text{-}eq}(x, x) \rangle \in \operatorname{derivable}(t) \Rightarrow \\ \quad \quad \langle h, \operatorname{enc\text{-}in}(x, \operatorname{enc\text{-}setc}) \rangle \in \operatorname{derivable}(t) \end{array}}}
\vnbprop{adm-nn-induction}{\vnbstmt{\begin{array}{@{}l@{}} \forall t, s, c. \\ \quad t \subseteq \mathit{clf} \wedge \\ \quad \operatorname{ax\text{-}nn\text{-}ind} \in t \wedge \\ \quad \operatorname{ax\text{-}mem\text{-}set} \in t \wedge \\ \quad \operatorname{ax\text{-}zero\text{-}nn} \in t \wedge \\ \quad \operatorname{ax\text{-}succ\text{-}nn} \in t \wedge \\ \quad \operatorname{def\text{-}comp\text{-}closed}(t) \wedge \\ \quad \operatorname{def\text{-}in\text{-}closed}(t) \wedge \\ \quad \operatorname{rel\text{-}nn\text{-}induction}(t, s, c) \wedge \\ \quad \operatorname{hyps\text{-}in}(s, \operatorname{derivable}(t)) \Rightarrow \\ \quad \quad c \in \operatorname{derivable}(t) \end{array}}}
\vnbprop{der-em}{\vnbstmt{\begin{array}{@{}l@{}} \forall t, h, x. \\ \quad ((h \in \operatorname{ctx\text{-}set}) \wedge (x \in \mathit{clf})) \Rightarrow (\langle h, \operatorname{enc\text{-}or}(x, \operatorname{enc\text{-}not}(x)) \rangle \in \operatorname{derivable}(t)) \end{array}}}
\vnbprop{subst-open-name}{\vnbstmt{\begin{array}{@{}l@{}} \forall y, e, v. \\ \quad ((y \in \mathbb{N}) \wedge ((e \in \operatorname{expr\text{-}set}) \wedge ((v \in \operatorname{expr\text{-}set}) \wedge (\operatorname{expr\text{-}lc}(v) \wedge \neg (y \in \operatorname{expr\text{-}fv}(e)))))) \Rightarrow  \\ \quad \quad \operatorname{expr\text{-}subst}(\operatorname{expr\text{-}open}(e, 0, \operatorname{efvar}(y)), y, v) = \operatorname{expr\text{-}open}(e, 0, v) \end{array}}}
\vnbprop{hyps-in-3-elim}{\vnbstmt{\begin{array}{@{}l@{}} \forall a, b, c, s. \\ \quad ((\langle a, b, c \rangle = \langle a, b, c \rangle) \wedge \operatorname{hyps\text{-}in}(\langle a, b, c \rangle, s)) \Rightarrow  \\ \quad \quad (a \in s) \wedge ((b \in s) \wedge (c \in s)) \end{array}}}
\vnbprop{adm-tfi-3cases}{\vnbstmt{\begin{array}{@{}l@{}} \forall t, s, c. \\ \quad t \subseteq \mathit{clf} \wedge \\ \quad \operatorname{ax\text{-}tfi} \in t \wedge \\ \quad \operatorname{ax\text{-}mem\text{-}set} \in t \wedge \\ \quad \operatorname{ax\text{-}limit\text{-}iff} \in t \wedge \\ \quad \operatorname{ax\text{-}ord\text{-}succ\text{-}above} \in t \wedge \\ \quad \operatorname{def\text{-}comp\text{-}closed}(t) \wedge \\ \quad \operatorname{rel\text{-}tfi\text{-}3cases}(t, s, c) \wedge \\ \quad \operatorname{hyps\text{-}in}(s, \operatorname{derivable}(t)) \Rightarrow \\ \quad \quad c \in \operatorname{derivable}(t) \end{array}}}
\vnbprop{enc-apply-unfold}{\vnbstmt{\forall p, q.\; (\operatorname{enc\text{-}apply}(p, q) \simeq \operatorname{eapp}(18, \langle p, q \rangle))}}
\vnbprop{enc-fun1-unfold}{\vnbstmt{\forall p.\; (\operatorname{enc\text{-}fun1}(p) \simeq \operatorname{eapp}(19, \langle p \rangle))}}
\vnbprop{adm-def-app}{\vnbstmt{\begin{array}{@{}l@{}} \forall t, s, c. \\ \quad t \subseteq \mathit{clf} \wedge \\ \quad \operatorname{ax\text{-}fun\text{-}app} \in t \wedge \\ \quad \operatorname{adm\text{-}def\text{-}strict\text{-}in\text{-}l}(t) \wedge \\ \quad \operatorname{adm\text{-}def\text{-}strict\text{-}in\text{-}r}(t) \wedge \\ \quad \operatorname{def\text{-}app}(t, s, c) \wedge \\ \quad \operatorname{hyps\text{-}in}(s, \operatorname{derivable}(t)) \Rightarrow \\ \quad \quad c \in \operatorname{derivable}(t) \end{array}}}
\vnbprop{enc-pred1-unfold}{\vnbstmt{\forall p, q.\; (\operatorname{enc\text{-}pred1}(p, q) \simeq \operatorname{eapp}(10, \langle p, q \rangle))}}
\vnbprop{enc-le-unfold}{\vnbstmt{\forall p, q.\; (\operatorname{enc\text{-}le}(p, q) \simeq \operatorname{eapp}(10, \langle \operatorname{econst}(11), p, q \rangle))}}
\vnbprop{enc-lt-unfold}{\vnbstmt{\forall p, q.\; (\operatorname{enc\text{-}lt}(p, q) \simeq \operatorname{eapp}(10, \langle \operatorname{econst}(13), p, q \rangle))}}
\vnbprop{pred-code-nn}{\vnbstmt{\forall k.\; (\operatorname{pred\text{-}code}(k) \Rightarrow (k \in \mathbb{N}))}}
\vnbprop{clf-pred-arg}{\vnbstmt{\begin{array}{@{}l@{}} \forall l, i. \\ \quad ((\operatorname{eapp}(10, l) \in \mathit{clf}) \wedge ((i \in \mathbb{N}) \wedge ((1 \leq i) \wedge (i \leq \operatorname{length}(l))))) \Rightarrow  \\ \quad \quad \operatorname{nth}(i, l) \in \mathit{clt} \end{array}}}
\vnbprop{adm-pred1-arg}{\vnbstmt{\begin{array}{@{}l@{}} \forall t, g, k, s. \\ \quad (\operatorname{pred\text{-}code}(k) \wedge (\langle g, \operatorname{enc\text{-}pred1}(\operatorname{econst}(k), s) \rangle \in \operatorname{derivable}(t))) \Rightarrow  \\ \quad \quad \langle g, \operatorname{enc\text{-}eq}(s, s) \rangle \in \operatorname{derivable}(t) \end{array}}}
\vnbprop{adm-pred1-unfold}{\vnbstmt{\begin{array}{@{}l@{}} \forall t, g, k, s, b. \\ \quad \operatorname{pred\text{-}code}(k) \wedge \\ \quad \langle g, \operatorname{enc\text{-}pred1}(\operatorname{econst}(k), s) \rangle \in \operatorname{derivable}(t) \wedge \\ \quad \langle g, \operatorname{enc\text{-}all}(\operatorname{enc\text{-}iff}(\operatorname{enc\text{-}pred1}(\operatorname{econst}(k), \operatorname{ebvar}(0)), b)) \rangle \in \operatorname{derivable}(t) \Rightarrow \\ \quad \quad \langle g, \operatorname{expr\text{-}open}(b, 0, s) \rangle \in \operatorname{derivable}(t) \end{array}}}
\vnbprop{def-pred-unfold-derived}{\vnbstmt{\begin{array}{@{}l@{}} \forall t, s, c. \\ \quad (\operatorname{def\text{-}pred\text{-}unfold}(t, s, c) \wedge \operatorname{hyps\text{-}in}(s, \operatorname{derivable}(t))) \Rightarrow (c \in \operatorname{derivable}(t)) \end{array}}}

\vnbbookfile{theorem-library/meta-def-laws}{theorem-library/meta-def-laws}
\noindent{\small meta-def-laws.scm -- K2 PHASE 5: the definedness rules of structure-library/meta-def.scm against the core (docs/meta-k2-admissibility-2026-10-08.md, section 4.7; report scratchpad/triage/META-5-REPORT.md). Agent META-5, 2026-10-12.  Every theorem here is (qed ...); helper procedures are local (prefix md-). (1) the two ...}\par
\vnbprop{def-in-closed-from-strict}{\vnbstmt{\forall t.\; (\operatorname{adm\text{-}def\text{-}strict\text{-}in\text{-}l}(t) \Rightarrow \operatorname{def\text{-}in\text{-}closed}(t))}}
\vnbprop{def-comp-closed-from-schema}{\vnbstmt{\forall t.\; (\operatorname{adm\text{-}def\text{-}comp}(t) \Rightarrow \operatorname{def\text{-}comp\text{-}closed}(t))}}
\vnbprop{clf-eq-elim}{\vnbstmt{\forall x, y.\; ((\operatorname{enc\text{-}eq}(x, y) \in \mathit{clf}) \Rightarrow ((x \in \mathit{clt}) \wedge (y \in \mathit{clt})))}}
\vnbprop{adm-def-strict-eq-l}{\vnbstmt{\begin{array}{@{}l@{}} \forall t, s, c. \\ \quad (\operatorname{def\text{-}strict\text{-}eq\text{-}l}(t, s, c) \wedge \operatorname{hyps\text{-}in}(s, \operatorname{derivable}(t))) \Rightarrow (c \in \operatorname{derivable}(t)) \end{array}}}
\vnbprop{adm-def-strict-eq-r}{\vnbstmt{\begin{array}{@{}l@{}} \forall t, s, c. \\ \quad (\operatorname{def\text{-}strict\text{-}eq\text{-}r}(t, s, c) \wedge \operatorname{hyps\text{-}in}(s, \operatorname{derivable}(t))) \Rightarrow (c \in \operatorname{derivable}(t)) \end{array}}}
\vnbprop{def-strict-in-l-admitted}{\vnbstmt{\begin{array}{@{}l@{}} \forall t, s, c. \\ \quad (\operatorname{def\text{-}strict\text{-}in\text{-}l}(t, s, c) \wedge \operatorname{hyps\text{-}in}(s, \operatorname{derivable}(t))) \Rightarrow (c \in \operatorname{derivable}(t)) \end{array}}}
\vnbprop{def-strict-in-r-admitted}{\vnbstmt{\begin{array}{@{}l@{}} \forall t, s, c. \\ \quad (\operatorname{def\text{-}strict\text{-}in\text{-}r}(t, s, c) \wedge \operatorname{hyps\text{-}in}(s, \operatorname{derivable}(t))) \Rightarrow (c \in \operatorname{derivable}(t)) \end{array}}}
\vnbprop{def-total-admitted}{\vnbstmt{\begin{array}{@{}l@{}} \forall t, s, c. \\ \quad (\operatorname{def\text{-}total}(t, s, c) \wedge \operatorname{hyps\text{-}in}(s, \operatorname{derivable}(t))) \Rightarrow (c \in \operatorname{derivable}(t)) \end{array}}}
\vnbprop{def-binder-admitted}{\vnbstmt{\begin{array}{@{}l@{}} \forall t, s, c. \\ \quad (\operatorname{def\text{-}binder}(t, s, c) \wedge \operatorname{hyps\text{-}in}(s, \operatorname{derivable}(t))) \Rightarrow (c \in \operatorname{derivable}(t)) \end{array}}}
\vnbprop{def-comp-admitted}{\vnbstmt{\begin{array}{@{}l@{}} \forall t, s, c. \\ \quad (\operatorname{def\text{-}comp}(t, s, c) \wedge \operatorname{hyps\text{-}in}(s, \operatorname{derivable}(t))) \Rightarrow (c \in \operatorname{derivable}(t)) \end{array}}}
\vnbprop{def-list-admitted}{\vnbstmt{\begin{array}{@{}l@{}} \forall t, s, c. \\ \quad (\operatorname{def\text{-}list}(t, s, c) \wedge \operatorname{hyps\text{-}in}(s, \operatorname{derivable}(t))) \Rightarrow (c \in \operatorname{derivable}(t)) \end{array}}}
\vnbprop{adm-def-strict-in-l-holds}{\vnbstmt{\forall t.\; \operatorname{adm\text{-}def\text{-}strict\text{-}in\text{-}l}(t)}}
\vnbprop{adm-def-strict-in-r-holds}{\vnbstmt{\forall t.\; \operatorname{adm\text{-}def\text{-}strict\text{-}in\text{-}r}(t)}}
\vnbprop{adm-def-comp-holds}{\vnbstmt{\forall t.\; \operatorname{adm\text{-}def\text{-}comp}(t)}}
\vnbprop{def-in-closed-holds}{\vnbstmt{\forall t.\; \operatorname{def\text{-}in\text{-}closed}(t)}}
\vnbprop{def-comp-closed-holds}{\vnbstmt{\forall t.\; \operatorname{def\text{-}comp\text{-}closed}(t)}}
\vnbprop{adm-big-union-mem-intro-core}{\vnbstmt{\begin{array}{@{}l@{}} \forall t, s, c. \\ \quad (\operatorname{rel\text{-}big\text{-}union\text{-}mem\text{-}intro}(t, s, c) \wedge \operatorname{hyps\text{-}in}(s, \operatorname{derivable}(t))) \Rightarrow  \\ \quad \quad c \in \operatorname{derivable}(t) \end{array}}}
\vnbprop{adm-transfinite-induction-core}{\vnbstmt{\begin{array}{@{}l@{}} \forall t, s, c. \\ \quad t \subseteq \mathit{clf} \wedge \\ \quad \operatorname{ax\text{-}tfi} \in t \wedge \\ \quad \operatorname{ax\text{-}mem\text{-}set} \in t \wedge \\ \quad \operatorname{rel\text{-}transfinite\text{-}induction}(t, s, c) \wedge \\ \quad \operatorname{hyps\text{-}in}(s, \operatorname{derivable}(t)) \Rightarrow \\ \quad \quad c \in \operatorname{derivable}(t) \end{array}}}
\vnbprop{adm-nn-induction-core}{\vnbstmt{\begin{array}{@{}l@{}} \forall t, s, c. \\ \quad t \subseteq \mathit{clf} \wedge \\ \quad \operatorname{ax\text{-}nn\text{-}ind} \in t \wedge \\ \quad \operatorname{ax\text{-}mem\text{-}set} \in t \wedge \\ \quad \operatorname{ax\text{-}zero\text{-}nn} \in t \wedge \\ \quad \operatorname{ax\text{-}succ\text{-}nn} \in t \wedge \\ \quad \operatorname{rel\text{-}nn\text{-}induction}(t, s, c) \wedge \\ \quad \operatorname{hyps\text{-}in}(s, \operatorname{derivable}(t)) \Rightarrow \\ \quad \quad c \in \operatorname{derivable}(t) \end{array}}}
\vnbprop{adm-tfi-3cases-core}{\vnbstmt{\begin{array}{@{}l@{}} \forall t, s, c. \\ \quad t \subseteq \mathit{clf} \wedge \\ \quad \operatorname{ax\text{-}tfi} \in t \wedge \\ \quad \operatorname{ax\text{-}mem\text{-}set} \in t \wedge \\ \quad \operatorname{ax\text{-}limit\text{-}iff} \in t \wedge \\ \quad \operatorname{ax\text{-}ord\text{-}succ\text{-}above} \in t \wedge \\ \quad \operatorname{rel\text{-}tfi\text{-}3cases}(t, s, c) \wedge \\ \quad \operatorname{hyps\text{-}in}(s, \operatorname{derivable}(t)) \Rightarrow \\ \quad \quad c \in \operatorname{derivable}(t) \end{array}}}
\vnbprop{def-strict-arg-admitted}{\vnbstmt{\begin{array}{@{}l@{}} \forall t, s, c. \\ \quad (\operatorname{def\text{-}strict\text{-}arg}(t, s, c) \wedge \operatorname{hyps\text{-}in}(s, \operatorname{derivable}(t))) \Rightarrow (c \in \operatorname{derivable}(t)) \end{array}}}
\vnbprop{def-strict-dom-admitted}{\vnbstmt{\begin{array}{@{}l@{}} \forall t, s, c. \\ \quad (\operatorname{def\text{-}strict\text{-}dom}(t, s, c) \wedge \operatorname{hyps\text{-}in}(s, \operatorname{derivable}(t))) \Rightarrow (c \in \operatorname{derivable}(t)) \end{array}}}
\vnbprop{adm-def-strict-arg-holds}{\vnbstmt{\forall t.\; \operatorname{adm\text{-}def\text{-}strict\text{-}arg}(t)}}
\vnbprop{adm-def-strict-dom-holds}{\vnbstmt{\forall t.\; \operatorname{adm\text{-}def\text{-}strict\text{-}dom}(t)}}
\vnbprop{clt-eapp-arg}{\vnbstmt{\begin{array}{@{}l@{}} \forall h, l, i. \\ \quad ((\operatorname{eapp}(h, l) \in \mathit{clt}) \wedge ((i \in \mathbb{N}) \wedge ((1 \leq i) \wedge (i \leq \operatorname{length}(l))))) \Rightarrow  \\ \quad \quad \operatorname{nth}(i, l) \in \mathit{clt} \end{array}}}
\vnbprop{clt-ebind-dom}{\vnbstmt{\forall h, a, b.\; (((h \in \mathbb{N}) \wedge (\operatorname{ebind}(h, 1, \langle a \rangle, b) \in \mathit{clt})) \Rightarrow (a \in \mathit{clt}))}}
\vnbprop{clf-sep-open}{\vnbstmt{\begin{array}{@{}l@{}} \forall y, d, p. \\ \quad (\operatorname{enc\text{-}in}(y, \operatorname{enc\text{-}sep}(d, p)) \in \mathit{clf}) \Rightarrow (\operatorname{expr\text{-}open}(p, 0, y) \in \mathit{clf}) \end{array}}}
\vnbprop{strict-reach-unfold}{\vnbstmt{\begin{array}{@{}l@{}} \forall e. \\ \quad (\operatorname{strict\text{-}reach}(e) \simeq \{z \in \operatorname{expr\text{-}set} : \forall w \in \operatorname{power}(\operatorname{expr\text{-}set}).\; (((e \in w) \wedge \operatorname{strict\text{-}closed}(w)) \Rightarrow (z \in w))\}) \end{array}}}
\vnbprop{true-reach-unfold}{\vnbstmt{\begin{array}{@{}l@{}} \forall j. \\ \quad (\operatorname{true\text{-}reach}(j) \simeq \{z \in \operatorname{expr\text{-}set} : \forall w \in \operatorname{power}(\operatorname{expr\text{-}set}).\; (((j \in w) \wedge \operatorname{true\text{-}closed}(w)) \Rightarrow (z \in w))\}) \end{array}}}
\vnbprop{strict-reach-least}{\vnbstmt{\begin{array}{@{}l@{}} \forall e.; \forall s \in \operatorname{power}(\operatorname{expr\text{-}set}), x \in \operatorname{strict\text{-}reach}(e). \\ \quad ((e \in s) \wedge \operatorname{strict\text{-}closed}(s)) \Rightarrow (x \in s) \end{array}}}
\vnbprop{true-reach-least}{\vnbstmt{\begin{array}{@{}l@{}} \forall e.; \forall s \in \operatorname{power}(\operatorname{expr\text{-}set}), x \in \operatorname{true\text{-}reach}(e). \\ \quad ((e \in s) \wedge \operatorname{true\text{-}closed}(s)) \Rightarrow (x \in s) \end{array}}}
\vnbprop{adm-strict-walk}{\vnbstmt{\begin{array}{@{}l@{}} \forall t, g, s, x. \\ \quad ((\langle g, \operatorname{enc\text{-}eq}(s, s) \rangle \in \operatorname{derivable}(t)) \wedge (x \in \operatorname{strict\text{-}reach}(s))) \Rightarrow  \\ \quad \quad \langle g, \operatorname{enc\text{-}eq}(x, x) \rangle \in \operatorname{derivable}(t) \end{array}}}
\vnbprop{adm-true-walk}{\vnbstmt{\begin{array}{@{}l@{}} \forall t, g, f, x. \\ \quad ((\langle g, f \rangle \in \operatorname{derivable}(t)) \wedge (x \in \operatorname{true\text{-}reach}(f))) \Rightarrow (\langle g, x \rangle \in \operatorname{derivable}(t)) \end{array}}}
\vnbprop{adm-strict-atom}{\vnbstmt{\begin{array}{@{}l@{}} \forall t, g, f, w. \\ \quad ((\langle g, f \rangle \in \operatorname{derivable}(t)) \wedge \operatorname{strict\text{-}atom}(f, w)) \Rightarrow  \\ \quad \quad \langle g, \operatorname{enc\text{-}eq}(w, w) \rangle \in \operatorname{derivable}(t) \end{array}}}
\vnbprop{def-strict-hyp-derived}{\vnbstmt{\begin{array}{@{}l@{}} \forall t, s, c. \\ \quad (\operatorname{def\text{-}strict\text{-}hyp}(t, s, c) \wedge \operatorname{hyps\text{-}in}(s, \operatorname{derivable}(t))) \Rightarrow (c \in \operatorname{derivable}(t)) \end{array}}}
\vnbprop{adm-def-strict-hyp-holds}{\vnbstmt{\forall t.\; \operatorname{adm\text{-}def\text{-}strict\text{-}hyp}(t)}}
\vnbprop{enc-nth-unfold}{\vnbstmt{\forall p_{7}, q_{7}.\; (\operatorname{enc\text{-}nth}(p_{7}, q_{7}) \simeq \operatorname{eapp}(15, \langle p_{7}, q_{7} \rangle))}}
\vnbprop{def-nth-list-derived}{\vnbstmt{\begin{array}{@{}l@{}} \forall t, s, c. \\ \quad (\operatorname{def\text{-}nth\text{-}list}(t, s, c) \wedge \operatorname{hyps\text{-}in}(s, \operatorname{derivable}(t))) \Rightarrow (c \in \operatorname{derivable}(t)) \end{array}}}
\vnbprop{adm-def-nth-list-holds}{\vnbstmt{\forall t.\; \operatorname{adm\text{-}def\text{-}nth\text{-}list}(t)}}
\vnbprop{def-strict-pred-admitted}{\vnbstmt{\begin{array}{@{}l@{}} \forall t, s, c. \\ \quad (\operatorname{def\text{-}strict\text{-}pred}(t, s, c) \wedge \operatorname{hyps\text{-}in}(s, \operatorname{derivable}(t))) \Rightarrow (c \in \operatorname{derivable}(t)) \end{array}}}
\vnbprop{adm-def-strict-pred-holds}{\vnbstmt{\forall t.\; \operatorname{adm\text{-}def\text{-}strict\text{-}pred}(t)}}
\vnbprop{def-strict-pred-refuses-qeq}{\vnbstmt{\begin{array}{@{}l@{}} \forall t, g, a, b, c. \\ \quad ((g \in \operatorname{ctx\text{-}set}) \wedge ((a \in \mathit{clt}) \wedge (b \in \mathit{clt}))) \Rightarrow  \\ \quad \quad \neg \operatorname{cr\text{-}def\text{-}strict\text{-}pred}(t, \langle \langle g, \operatorname{enc\text{-}qeq}(a, b) \rangle \rangle, c) \end{array}}}
\vnbprop{def-strict-pred-refuses-not}{\vnbstmt{\begin{array}{@{}l@{}} \forall t, g, a, c. \\ \quad ((g \in \operatorname{ctx\text{-}set}) \wedge (a \in \mathit{clt})) \Rightarrow  \\ \quad \quad \neg \operatorname{cr\text{-}def\text{-}strict\text{-}pred}(t, \langle \langle g, \operatorname{enc\text{-}not}(\operatorname{enc\text{-}eq}(a, a)) \rangle \rangle, c) \end{array}}}
\vnbprop{strict-atom-pred2}{\vnbstmt{\begin{array}{@{}l@{}} \forall k, x, y, w. \\ \quad (\operatorname{pred\text{-}code}(k) \wedge ((x \in \operatorname{expr\text{-}set}) \wedge ((y \in \operatorname{expr\text{-}set}) \wedge ((w \in \operatorname{strict\text{-}reach}(x)) \vee (w \in \operatorname{strict\text{-}reach}(y)))))) \Rightarrow  \\ \quad \quad \operatorname{strict\text{-}atom}(\operatorname{eapp}(10, \langle \operatorname{econst}(k), x, y \rangle), w) \end{array}}}
\vnbprop{strict-atom-le}{\vnbstmt{\begin{array}{@{}l@{}} \forall x, y, w. \\ \quad ((x \in \operatorname{expr\text{-}set}) \wedge ((y \in \operatorname{expr\text{-}set}) \wedge ((w \in \operatorname{strict\text{-}reach}(x)) \vee (w \in \operatorname{strict\text{-}reach}(y))))) \Rightarrow  \\ \quad \quad \operatorname{strict\text{-}atom}(\operatorname{enc\text{-}le}(x, y), w) \end{array}}}
\vnbprop{strict-atom-lt}{\vnbstmt{\begin{array}{@{}l@{}} \forall x, y, w. \\ \quad ((x \in \operatorname{expr\text{-}set}) \wedge ((y \in \operatorname{expr\text{-}set}) \wedge ((w \in \operatorname{strict\text{-}reach}(x)) \vee (w \in \operatorname{strict\text{-}reach}(y))))) \Rightarrow  \\ \quad \quad \operatorname{strict\text{-}atom}(\operatorname{enc\text{-}lt}(x, y), w) \end{array}}}
\vnbprop{def-struct-hyp-derived}{\vnbstmt{\begin{array}{@{}l@{}} \forall t, s, c. \\ \quad (\operatorname{def\text{-}struct\text{-}hyp}(t, s, c) \wedge \operatorname{hyps\text{-}in}(s, \operatorname{derivable}(t))) \Rightarrow (c \in \operatorname{derivable}(t)) \end{array}}}
\vnbprop{adm-def-struct-hyp-holds}{\vnbstmt{\forall t.\; \operatorname{adm\text{-}def\text{-}struct\text{-}hyp}(t)}}
\vnbprop{adm-def-pred-unfold-holds}{\vnbstmt{\forall t.\; \operatorname{adm\text{-}def\text{-}pred\text{-}unfold}(t)}}
\vnbprop{def-pred-hyp-derived}{\vnbstmt{\begin{array}{@{}l@{}} \forall t, s, c. \\ \quad (\operatorname{def\text{-}pred\text{-}hyp}(t, s, c) \wedge \operatorname{hyps\text{-}in}(s, \operatorname{derivable}(t))) \Rightarrow (c \in \operatorname{derivable}(t)) \end{array}}}
\vnbprop{adm-def-pred-hyp-holds}{\vnbstmt{\forall t.\; \operatorname{adm\text{-}def\text{-}pred\text{-}hyp}(t)}}



\end{document}
